\IfFileExists{stacks-project-book.cls}{%
\documentclass{stacks-project-book}
}{%
\documentclass{amsbook}
}

\usepackage{verbatim}
\newenvironment{reference}{\comment}{\endcomment}
\newenvironment{slogan}{\comment}{\endcomment}
\newenvironment{history}{\comment}{\endcomment}

\usepackage[all]{xy}

\xyoption{2cell}
\UseAllTwocells

\usepackage{multicol}

\usepackage{lmodern}
\usepackage[T1]{fontenc}
\usepackage[french]{babel}
\AtBeginDocument{\shorthandoff{?}}


\usepackage{hyperref}


\theoremstyle{plain}
\newtheorem{theorem}[subsection]{Théorème}
\newtheorem{proposition}[subsection]{Proposition}
\newtheorem{lemma}[subsection]{Lemme}

\theoremstyle{definition}
\newtheorem{definition}[subsection]{Définition}
\newtheorem{example}[subsection]{Exemple}
\newtheorem{exercise}[subsection]{Exercice}
\newtheorem{situation}[subsection]{Situation}

\theoremstyle{remark}
\newtheorem{remark}[subsection]{Remarque}
\newtheorem{remarks}[subsection]{Remarques}

\renewcommand{\proofname}{Démonstration}

\numberwithin{equation}{subsection}

\def\lim{\mathop{\mathrm{lim}}\nolimits}
\def\colim{\mathop{\mathrm{colim}}\nolimits}
\def\Spec{\mathop{\mathrm{Spec}}}
\def\Hom{\mathop{\mathrm{Hom}}\nolimits}
\def\Ext{\mathop{\mathrm{Ext}}\nolimits}
\def\SheafHom{\mathop{\mathcal{H}\!\mathit{om}}\nolimits}
\def\SheafExt{\mathop{\mathcal{E}\!\mathit{xt}}\nolimits}
\def\Sch{\mathit{Sch}}
\def\Mor{\mathop{\mathrm{Mor}}\nolimits}
\def\Ob{\mathop{\mathrm{Ob}}\nolimits}
\def\Sh{\mathop{\mathit{Sh}}\nolimits}
\def\NL{\mathop{N\!L}\nolimits}
\def\CH{\mathop{\mathrm{CH}}\nolimits}
\def\proetale{{pro\text{-}\acute{e}tale}}
\def\etale{{\acute{e}tale}}
\def\QCoh{\mathit{QCoh}}
\def\Ker{\mathop{\mathrm{Ker}}}
\def\Im{\mathop{\mathrm{Im}}}
\def\Coker{\mathop{\mathrm{Coker}}}
\def\Coim{\mathop{\mathrm{Coim}}}

\DeclareMathSymbol{\boxtimes}{\mathbin}{AMSa}{"02}

\def\QCohstack{\mathcal{QC}\!\mathit{oh}}
\def\Cohstack{\mathcal{C}\!\mathit{oh}}
\def\Spacesstack{\mathcal{S}\!\mathit{paces}}
\def\Quotfunctor{\mathrm{Quot}}
\def\Hilbfunctor{\mathrm{Hilb}}
\def\Curvesstack{\mathcal{C}\!\mathit{urves}}
\def\Polarizedstack{\mathcal{P}\!\mathit{olarized}}
\def\Complexesstack{\mathcal{C}\!\mathit{omplexes}}
\def\Pic{\mathop{\mathrm{Pic}}\nolimits}
\def\Picardstack{\mathcal{P}\!\mathit{ic}}
\def\Picardfunctor{\mathrm{Pic}}
\def\Deformationcategory{\mathcal{D}\!\mathit{ef}}
\begin{document}
\title{Le projet Stacks --- édition française, partie 3}
\maketitle
\tableofcontents

% Bon, commençons ici.
%

\chapter{Lissage des morphismes d'anneaux}



\phantomsection
\label{smoothing-section-phantom}


\section{Introduction}
\label{smoothing-section-introduction}

\noindent
Le résultat principal de ce chapitre est le suivant :
$$
\fbox{Un morphisme d'anneaux régulier entre anneaux noethériens est une
colimite filtrante de morphismes lisses.}
$$
Ce théorème est dû à Popescu, voir \cite{popescu-letter}.
Richard Swan a donné une présentation accessible de la démonstration de Popescu,
voir \cite{swan} ; il s'est appuyé sur des notes d'Andr\'e et sur un article d'Ogoma, voir
\cite{Ogoma}.

\medskip\noindent
Notre présentation suit celle de Swan, mais nous démontrons d'abord un résultat
intermédiaire qui permet de nous placer dans une situation un peu plus simple.
En voici un aperçu. Nous résolvons d'abord le ``problème de relèvement'' suivant :
une déformation infinitésimale plate d'une colimite filtrante d'algèbres lisses
est une colimite filtrante d'algèbres lisses. Ce résultat dit essentiellement
qu'il suffit de démontrer le théorème principal pour les morphismes entre
anneaux noethériens réduits. Nous démontrons ensuite deux lemmes très ingénieux,
appelés le ``lemme de relèvement'' et le ``lemme de désingularisation''.
Nous montrons que, combinés, ces lemmes ramènent le théorème principal à prouver
qu'une algèbre noethérienne $\Lambda$, géométriquement régulière sur un corps $k$,
est une colimite filtrante de $k$-algèbres lisses. Nous examinons ensuite les
arguments locaux nécessaires à la démonstration de Popescu-Ogoma-Swan-Andr\'e.
Enfin, nous donnons la démonstration dans les trois dernières sections.

\medskip\noindent
Nous terminons cette introduction par quelques indications bibliographiques.
Soit $A$ un anneau local noethérien hensélien.
Nous disons que $A$ possède la {\it propriété d'approximation} si, quels que soient
$f_1, \ldots, f_m \in A[x_1, \ldots, x_n]$, le système d'équations
$f_1 = 0, \ldots, f_m = 0$ admet une solution dans le complété de $A$ si et
seulement s'il admet une solution dans $A$. Cette définition est due à Artin.
Artin a d'abord démontré la propriété d'approximation pour les systèmes
d'équations analytiques, voir \cite{Artin-Analytic-Approximation}.
Dans \cite{Artin-Algebraic-Approximation}, Artin a démontré la propriété
d'approximation pour les anneaux locaux essentiellement de type fini sur un
anneau de valuation discrète excellent.
Artin a conjecturé (page 26 de \cite{Artin-Algebraic-Approximation}) que tout
anneau local hensélien excellent devait posséder la propriété d'approximation.

\medskip\noindent
À un certain moment, on a conjecturé que tout morphisme d'anneaux régulier entre
anneaux noethériens est une colimite filtrante d'algèbres lisses (voir par exemple
\cite{Raynaud-Rennes}, \cite{popescu-global}, \cite{Artin-power-series},
\cite{Artin-Denef}). Nous ne savons pas à qui cette conjecture\footnote{La
question/conjecture formulée dans \cite{Artin-power-series},
\cite{Artin-Denef} et \cite{popescu-global} est plus forte ; il a été démontré
dans \cite{Cipu} qu'elle est équivalente à la version originale.}
est due. Le lien avec la propriété d'approximation est le suivant : si
$A \to A^\wedge$ est une colimite d'algèbres lisses, avec $A$ comme ci-dessus,
alors la propriété d'approximation est satisfaite (insérer ici une référence
ultérieure). En outre, le théorème principal s'applique au morphisme
$A \to A^\wedge$ si $A$ est un anneau local excellent, car l'une des conditions
définissant un anneau local excellent est que ses fibres formelles soient
géométriquement régulières. Notons que les anneaux locaux excellents ont été
définis par Grothendieck et que leur définition a été publiée en 1965.

\medskip\noindent
Dans \cite{Artin-power-series}, il a été démontré que
$R \to R^\wedge$ est une colimite filtrante d'algèbres lisses pour tout anneau
local $R$ essentiellement de type fini sur un corps.
Dans \cite{Rotthaus-Artin}, il a été démontré que $R \to R^\wedge$ est une
colimite filtrante d'algèbres lisses pour tout anneau local $R$ essentiellement
de type fini sur un anneau de valuation discrète excellent.
Enfin, le théorème principal a été démontré dans
\cite{popescu-GND}, \cite{popescu-GNDA}, \cite{popescu-letter}, 
\cite{Ogoma} et \cite{swan}, comme indiqué ci-dessus.

\medskip\noindent
Réciproquement, à l'aide de certains des résultats précédents, il a été démontré
dans \cite{Rotthaus-excellent} que tout anneau local noethérien possédant la
propriété d'approximation est excellent.

\medskip\noindent
L'article \cite{Spivakovsky} propose une autre approche du théorème principal,
mais il semble difficile à lire (par exemple, \cite[lemme 5.2]{Spivakovsky}
semble être une reformulation incorrecte de \cite[lemme 3]{Elkik}). Il existe
également un exposé Bourbaki sur ce sujet, voir \cite{Teissier}.












\section{Idéaux singuliers}
\label{smoothing-section-singular-ideal}

\noindent
Soit $R \to A$ un morphisme d'anneaux. L'idéal singulier de $A$ sur $R$ est
l'idéal radical de $A$ qui définit le lieu singulier du morphisme
$\Spec(A) \to \Spec(R)$. En voici une définition précise.

\begin{definition}
\label{smoothing-definition-singular-ideal}
Soit $R \to A$ un morphisme d'anneaux. L'{\it idéal singulier de $A$ sur $R$},
noté $H_{A/R}$, est l'unique idéal radical $H_{A/R} \subset A$ tel que
$$
V(H_{A/R}) = \{\mathfrak q \in \Spec(A) \mid R \to A
\text{ non lisse en }\mathfrak q\}
$$
\end{definition}

\noindent
Cette définition a un sens car l'ensemble des idéaux premiers où $R \to A$
est lisse est ouvert, voir
Algèbre, Définition \ref{algebra-definition-smooth-at-prime}.
Afin de trouver un ensemble explicite de générateurs de l'idéal singulier,
nous démontrons d'abord le lemme suivant.

\begin{lemma}
\label{smoothing-lemma-find-strictly-standard}
Soit $R$ un anneau. Soit $A = R[x_1, \ldots, x_n]/(f_1, \ldots, f_m)$.
Soit $\mathfrak q \subset A$ un idéal premier. Supposons que $R \to A$ soit
lisse en l'idéal premier $\mathfrak q$. Il existe alors un $a \in A$, $a \not \in \mathfrak q$,
un entier $c$, $0 \leq c \leq \min(n, m)$, et des sous-ensembles
$U \subset \{1, \ldots, n\}$, $V \subset \{1, \ldots, m\}$
de cardinal $c$ tels que
$$
a = a' \det(\partial f_j/\partial x_i)_{j \in V, i \in U}
$$
pour un certain $a' \in A$, et
$$
a f_\ell \in (f_j, j \in V) + (f_1, \ldots, f_m)^2
$$
pour tout $\ell \in \{1, \ldots, m\}$.
\end{lemma}

\begin{proof}
Posons $I = (f_1, \ldots, f_m)$, de sorte que le complexe cotangent naïf de
$A$ sur $R$ soit homotopiquement équivalent à
$I/I^2 \to \bigoplus A\text{d}x_i$, voir
Algèbre, Lemme \ref{algebra-lemma-NL-homotopy}.
Nous utiliserons le fait que la formation du complexe cotangent naïf commute à
la localisation, voir Algèbre, section \ref{algebra-section-netherlander},
et en particulier Algèbre, Lemme \ref{algebra-lemma-localize-NL}.
D'après Algèbre, Définitions \ref{algebra-definition-smooth} et
\ref{algebra-definition-smooth-at-prime}
nous voyons que $(I/I^2)_a \to \bigoplus A_a\text{d}x_i$ est une injection
scindée pour un certain $a \in A$, $a \not \in \mathfrak q$.
Après avoir renuméroté $x_1, \ldots, x_n$ et $f_1, \ldots, f_m$, nous pouvons
supposer que $f_1, \ldots, f_c$ forment une base de l'espace vectoriel
$I/I^2 \otimes_A \kappa(\mathfrak q)$ et que les images de
$\text{d}x_{c + 1}, \ldots, \text{d}x_n$ forment une base de
$\Omega_{A/R} \otimes_A \kappa(\mathfrak q)$. Ainsi, après avoir remplacé $a$
par $aa'$ pour un certain $a' \in A$, $a' \not \in \mathfrak q$, nous pouvons
supposer que $f_1, \ldots, f_c$ forment une base de $(I/I^2)_a$ et que les images
de $\text{d}x_{c + 1}, \ldots, \text{d}x_n$ forment une base de
$(\Omega_{A/R})_a$. Dans cette situation, $a^N$ satisfait les conditions du
lemme pour un entier $N$ suffisamment grand (avec $U = V = \{1, \ldots, c\}$).
\end{proof}

\noindent
Nous utiliserons la notion d'élément {\it strictement standard} de $A$ sur $R$.
Notre notion est légèrement plus faible que celle de l'article de Swan
\cite{swan}. Nous définissons aussi un élément {\it standard élémentaire}
comme étant du type obtenu dans le lemme précédent. Nous comparons les
différents types d'éléments dans le lemme \ref{smoothing-lemma-compare-standard}.

\begin{definition}
\label{smoothing-definition-strictly-standard}
Soit $R \to A$ un morphisme d'anneaux de présentation finie.
Nous disons qu'un élément $a \in A$ est {\it standard élémentaire dans $A$ sur $R$}
s'il existe une présentation
$A = R[x_1, \ldots, x_n]/(f_1, \ldots, f_m)$
et un entier $0 \leq c \leq \min(n, m)$ tels que
\begin{equation}
\label{smoothing-equation-elementary-standard-one}
a = a' \det(\partial f_j/\partial x_i)_{i, j = 1, \ldots, c}
\end{equation}
pour un certain $a' \in A$, et
\begin{equation}
\label{smoothing-equation-elementary-standard-two}
a f_{c + j} \in (f_1, \ldots, f_c) + (f_1, \ldots, f_m)^2
\end{equation}
pour $j = 1, \ldots, m - c$. Nous disons que $a \in A$ est
{\it strictement standard dans $A$ sur $R$} s'il existe une présentation
$A = R[x_1, \ldots, x_n]/(f_1, \ldots, f_m)$
et un entier $0 \leq c \leq \min(n, m)$ tels que
\begin{equation}
\label{smoothing-equation-strictly-standard-one}
a = \sum\nolimits_{I \subset \{1, \ldots, n\},\ |I| = c}
a_I \det(\partial f_j/\partial x_i)_{j = 1, \ldots, c,\ i \in I}
\end{equation}
pour certains $a_I \in A$, et
\begin{equation}
\label{smoothing-equation-strictly-standard-two}
a f_{c + j} \in (f_1, \ldots, f_c) + (f_1, \ldots, f_m)^2
\end{equation}
pour $j = 1, \ldots, m - c$.
\end{definition}

\noindent
Le lemme suivant est utile pour dégager des conséquences de
(\ref{smoothing-equation-strictly-standard-one}).

\begin{lemma}
\label{smoothing-lemma-parse-equation-strictly-standard-one}
Soit $R$ un anneau. Soit $A = R[x_1, \ldots, x_n]/(f_1, \ldots, f_m)$
et posons $I = (f_1, \ldots, f_m)$. Soit $a \in A$. Alors
(\ref{smoothing-equation-strictly-standard-one}) implique qu'il existe un morphisme
$A$-linéaire
$\psi : \bigoplus\nolimits_{i = 1, \ldots, n} A \text{d}x_i \to A^{\oplus c}$
tel que la composée
$$
A^{\oplus c} \xrightarrow{(f_1, \ldots, f_c)}
I/I^2 \xrightarrow{f \mapsto \text{d}f}
\bigoplus\nolimits_{i = 1, \ldots, n} A \text{d}x_i
\xrightarrow{\psi}
A^{\oplus c}
$$
soit la multiplication par $a$. Réciproquement, si un tel $\psi$ existe,
alors $a^c$ satisfait (\ref{smoothing-equation-strictly-standard-one}).
\end{lemma}

\begin{proof}
Il s'agit d'un cas particulier de
Algèbre, Lemme \ref{algebra-lemma-matrix-left-inverse}.
\end{proof}

\begin{lemma}[Elkik]
\label{smoothing-lemma-elkik}
Soit $R \to A$ un morphisme d'anneaux de présentation finie.
L'idéal singulier $H_{A/R}$ est le radical de l'idéal engendré par les éléments
strictement standard de $A$ sur $R$, et aussi le radical de l'idéal engendré
par les éléments standard élémentaires de $A$ sur $R$.
\end{lemma}

\begin{proof}
Supposons que $a$ soit strictement standard dans $A$ sur $R$. Nous affirmons
que $A_a$ est lisse sur $R$, ce qui prouve que $a \in H_{A/R}$. En effet,
soient $A = R[x_1, \ldots, x_n]/(f_1, \ldots, f_m)$, $c$ et $a' \in A$
comme dans la définition \ref{smoothing-definition-strictly-standard}.
Posons $I = (f_1, \ldots, f_m)$, de sorte que le complexe cotangent naïf de
$A$ sur $R$ soit donné par $I/I^2 \to \bigoplus A\text{d}x_i$.
L'hypothèse (\ref{smoothing-equation-strictly-standard-two}) implique que $(I/I^2)_a$
est engendré par les classes de $f_1, \ldots, f_c$.
L'hypothèse (\ref{smoothing-equation-strictly-standard-one}) implique que la différentielle
$(I/I^2)_a \to \bigoplus A_a\text{d}x_i$ admet un inverse à gauche, voir
le lemme \ref{smoothing-lemma-parse-equation-strictly-standard-one}.
Ainsi, $R \to A_a$ est lisse par définition et d'après
Algèbre, Lemme \ref{algebra-lemma-localize-NL}.

\medskip\noindent
Soient $H_e, H_s \subset A$ les radicaux des idéaux engendrés par les éléments
standard élémentaires, resp.\ strictement standard, de $A$ sur $R$.
Par définition et d'après ce que nous venons de démontrer, nous avons
$H_e \subset H_s \subset H_{A/R}$. L'inclusion $H_{A/R} \subset H_e$ résulte
du lemme \ref{smoothing-lemma-find-strictly-standard}.
\end{proof}

\begin{example}
\label{smoothing-example-not-quasi-compact}
L'ensemble des points où un morphisme d'anneaux de présentation finie est lisse
n'est pas nécessairement un ouvert quasi-compact. Par exemple, soient
$R = k[x, y_1, y_2, y_3, \ldots]/(xy_i)$ et $A = R/(x)$.
Alors le lieu de lissité de $R \to A$ est $\bigcup D(y_i)$, qui n'est pas
quasi-compact.
\end{example}

\begin{lemma}
\label{smoothing-lemma-strictly-standard-base-change}
Soit $R \to A$ un morphisme d'anneaux de présentation finie.
Soit $R \to R'$ un morphisme d'anneaux. Si $a \in A$ est standard élémentaire,
resp.\ strictement standard, dans $A$ sur $R$, alors $a \otimes 1$ est standard
élémentaire, resp.\ strictement standard, dans $A \otimes_R R'$ sur $R'$.
\end{lemma}

\begin{proof}
Si $A = R[x_1, \ldots, x_n]/(f_1, \ldots, f_m)$ est une présentation de $A$
sur $R$, alors
$A \otimes_R R' = R'[x_1, \ldots, x_n]/(f'_1, \ldots, f'_m)$
est une présentation de $A \otimes_R R'$ sur $R'$. Ici, $f'_j$ est l'image de
$f_j$ dans $R'[x_1, \ldots, x_n]$. Le résultat découle donc des définitions.
\end{proof}

\begin{lemma}
\label{smoothing-lemma-final-solve}
Soient $R \to A \to \Lambda$ des morphismes d'anneaux, avec $A$ de présentation
finie sur $R$. Supposons que $H_{A/R} \Lambda = \Lambda$. Il existe alors une
factorisation $A \to B \to \Lambda$ telle que $B$ soit lisse sur $R$.
\end{lemma}

\begin{proof}
Choisissons $f_1, \ldots, f_r \in H_{A/R}$ et
$\lambda_1, \ldots, \lambda_r \in \Lambda$ tels que
$\sum f_i\lambda_i = 1$ dans $\Lambda$. Posons
$B = A[x_1, \ldots, x_r]/(f_1x_1 + \ldots + f_rx_r - 1)$
et définissons $B \to \Lambda$ en envoyant $x_i$ sur $\lambda_i$.
Pour vérifier que $B$ est lisse sur $R$, utilisons le fait que $A_{f_i}$ est
lisse sur $R$ par définition de $H_{A/R}$ et que $B_{f_i}$ est lisse sur
$A_{f_i}$. Détails omis.
\end{proof}





\section{Présentations d'algèbres}
\label{smoothing-section-presentations}

\noindent
Certains des résultats de cette section sont dus à Elkik. Notons que l'algèbre
$C$ du lemme suivant est une algèbre symétrique sur $A$. De plus, si $R$ est
noethérien, alors $C$ est de présentation finie sur $R$.

\begin{lemma}
\label{smoothing-lemma-improve-presentation}
Soit $R$ un anneau et soit $A$ une $R$-algèbre de présentation finie.
Il existe un morphisme de $R$-algèbres de type fini $A \to C$ qui admet une
rétraction et possède les deux propriétés suivantes :
\begin{enumerate}
\item pour tout $a \in A$ tel que $R \to A_a$ soit une intersection complète
locale (Compléments d'algèbre, Définition
\ref{more-algebra-definition-local-complete-intersection}),
l'anneau $C_a$ est lisse sur $A_a$ et admet une présentation
$C_a = R[y_1, \ldots, y_m]/J$ telle que $J/J^2$ soit libre sur $C_a$, et
\item pour tout $a \in A$ tel que $A_a$ soit lisse sur $R$, le module
$\Omega_{C_a/R}$ est libre sur $C_a$.
\end{enumerate}
\end{lemma}

\begin{proof}
Choisissons une présentation $A = R[x_1, \ldots, x_n]/I$ et écrivons
$I = (f_1, \ldots, f_m)$. Définissons le $A$-module $K$ par la suite exacte courte
$$
0 \to K \to A^{\oplus m} \to I/I^2 \to 0
$$
où le $j$-ième vecteur de base $e_j$ du terme central est envoyé sur la classe de
$f_j$ dans le terme de droite. Posons
$$
C = \text{Sym}^*_A(I/I^2).
$$
La rétraction est simplement la projection sur la composante de degré $0$ de
$C$. Nous avons une surjection
$R[x_1, \ldots, x_n, y_1, \ldots, y_m] \to C$ qui envoie $y_j$ sur la classe de
$f_j$ dans $I/I^2$. Le noyau $J$ de ce morphisme est engendré par les éléments
$f_1, \ldots, f_m$ et par les éléments $\sum h_j y_j$, où
$h_j \in R[x_1, \ldots, x_n]$, tels que $\sum h_j e_j$ définisse un élément de
$K$. D'après
Algèbre, Lemme \ref{algebra-lemma-exact-sequence-NL}
appliqué à $R \to A \to C$ et aux présentations ci-dessus, ainsi que
Compléments d'algèbre, Lemme
\ref{more-algebra-lemma-cotangent-complex-symmetric-algebra}
il existe une suite exacte
\begin{equation}
\label{smoothing-equation-sequence}
I/I^2 \otimes_A C \to J/J^2 \to K \otimes_A C \to 0
\end{equation}
de $C$-modules. Soit $h \in R[x_1, \ldots, x_n]$ un élément d'image $a \in A$.
Nous utiliserons, pour les anneaux localisés, les présentations
$$
A_a = R[x_0, x_1, \ldots, x_n]/I'
\quad\text{et}\quad
C_a = R[x_0, x_1, \ldots, x_n, y_1, \ldots, y_m]/J'
$$
où $I' = (hx_0 - 1, I)$ et $J' = (hx_0 - 1, J)$. Ainsi,
$I'/(I')^2 = A_a \oplus (I/I^2)_a$ comme $A_a$-modules et
$J'/(J')^2 = C_a \oplus (J/J^2)_a$ comme $C_a$-modules.
Nous obtenons donc
\begin{equation}
\label{smoothing-equation-sequence-localized}
C_a \oplus I/I^2 \otimes_A C_a \to
C_a \oplus (J/J^2)_a \to
K \otimes_A C_a \to 0
\end{equation}
qui est la suite fournie par
Algèbre, lemme \ref{algebra-lemma-exact-sequence-NL},
pour $R \to A_a \to C_a$ et les présentations ci-dessus.

\medskip\noindent
Supposons maintenant que $a \in A$ soit tel que $A_a$ soit une intersection
complète locale sur $R$. Alors $(I/I^2)_a$ est projectif de type fini sur $A_a$,
voir
Compléments d'algèbre, Lemme
\ref{more-algebra-lemma-quasi-regular-ideal-finite-projective}.
Nous voyons donc que $K_a \oplus (I/I^2)_a \cong A_a^{\oplus m}$ est libre.
En particulier, $K_a$ est lui aussi projectif de type fini.
D'après Compléments d'algèbre, Lemme \ref{more-algebra-lemma-transitive-lci-at-end},
la suite (\ref{smoothing-equation-sequence-localized}) est exacte à gauche. Ainsi,
$$
J'/(J')^2 \cong
C_a \oplus I/I^2 \otimes_A C_a \oplus K \otimes_A C_a \cong
C_a^{\oplus m + 1}
$$
Cela démontre (1). Supposons enfin que, de plus, $A_a$ soit lisse sur $R$.
La même présentation montre alors que le module des différentielles de Kähler
$\Omega_{C_a/R}$ est le conoyau du morphisme
$$
J'/(J')^2 \longrightarrow
\bigoplus\nolimits_i C_a\text{d}x_i \oplus \bigoplus\nolimits_j C_a\text{d}y_j
$$
Dans la décomposition ci-dessus, le facteur direct $C_a$ de $J'/(J')^2$
correspond à $hx_0 - 1$ et s'envoie donc isomorphiquement sur le facteur direct
$C_a\text{d}x_0$. Le facteur direct $I/I^2 \otimes_A C_a$ de $J'/(J')^2$
s'injecte dans $\bigoplus_{i = 1, \ldots, n} C_a\text{d}x_i$, avec pour quotient
$\Omega_{A_a/R} \otimes_{A_a} C_a$. Le facteur direct $K \otimes_A C_a$
s'injecte dans $\bigoplus_{j \geq 1} C_a\text{d}y_j$, avec un quotient isomorphe
à $I/I^2 \otimes_A C_a$. Ainsi, le conoyau du dernier morphisme affiché est le
module
$I/I^2 \otimes_A C_a \oplus \Omega_{A_a/R} \otimes_{A_a} C_a$.
Puisque $(I/I^2)_a \oplus \Omega_{A_a/R}$ est libre (d'après la définition des
morphismes d'anneaux lisses), nous voyons que (2) est satisfaite.
\end{proof}

\noindent
La proposition suivante a été démontrée par Elkik dans \cite{Elkik} pour les
morphismes d'anneaux lisses sur des couples henséliens. Pour les morphismes
d'anneaux lisses, on la trouve dans \cite{Arabia}, où il est également démontré
que les morphismes d'anneaux entre algèbres lisses peuvent être relevés.

\begin{proposition}
\label{smoothing-proposition-lift-smooth}
\begin{slogan}
Les algèbres lisses et syntomiques se relèvent à travers les surjections
\end{slogan}
Soit $R \to R_0$ un morphisme d'anneaux surjectif de noyau $I$.
\begin{enumerate}
\item Si $R_0 \to A_0$ est un morphisme d'anneaux syntomique, alors il existe
un morphisme d'anneaux syntomique $R \to A$ tel que $A/IA \cong A_0$.
\item Si $R_0 \to A_0$ est un morphisme d'anneaux lisse, alors il existe un
morphisme d'anneaux lisse $R \to A$ tel que $A/IA \cong A_0$.
\end{enumerate}
\end{proposition}

\begin{proof}
Supposons que $R_0 \to A_0$ soit syntomique, donc en particulier une
intersection complète locale (Compléments d'algèbre, Lemme
\ref{more-algebra-lemma-syntomic-lci}).
Choisissons une présentation $A_0 = R_0[x_1, \ldots, x_n]/J_0$. Posons
$C_0 = \text{Sym}^*_{A_0}(J_0/J_0^2)$. Notons que $J_0/J_0^2$ est un
$A_0$-module projectif de type fini (Algèbre, Lemme
\ref{algebra-lemma-syntomic-presentation-ideal-mod-squares}).
D'après le lemme \ref{smoothing-lemma-improve-presentation}, le morphisme d'anneaux
$A_0 \to C_0$ est lisse et nous pouvons trouver une présentation
$C_0 = R_0[y_1, \ldots, y_m]/K_0$ dans laquelle $K_0/K_0^2$ est libre sur
$C_0$. D'après Algèbre, Lemme \ref{algebra-lemma-huber}, nous pouvons supposer
$C_0 = R_0[y_1, \ldots, y_m]/(\overline{f}_1, \ldots, \overline{f}_c)$
où les images de $\overline{f}_1, \ldots, \overline{f}_c$ forment une base de
$K_0/K_0^2$ sur $C_0$. Choisissons des relèvements
$f_1, \ldots, f_c \in R[y_1, \ldots, y_m]$ de
$\overline{f}_1, \ldots, \overline{f}_c$ et posons
$$
C = R[y_1, \ldots, y_m]/(f_1, \ldots, f_c)
$$
Par construction, $C_0 = C/IC$. D'après Algèbre, Lemme
\ref{algebra-lemma-localize-relative-complete-intersection},
quitte à remplacer $C$ par $C_g$, nous pouvons supposer que $C$ est une
intersection complète globale relative sur $R$.
Nous en concluons qu'il existe un $A_0$-module projectif de type fini $P_0$ tel
que $C_0 = \text{Sym}^*_{A_0}(P_0)$ soit isomorphe à $C/IC$ pour une certaine
$R$-algèbre syntomique $C$.

\medskip\noindent
Choisissons un entier $n$ et une décomposition en somme directe
$A_0^{\oplus n} = P_0 \oplus Q_0$.
D'après Compléments d'algèbre, Lemme \ref{more-algebra-lemma-lift-projective-module},
nous pouvons trouver un morphisme d'anneaux \'etale $C \to C'$ qui induit un
isomorphisme $C/IC \to C'/IC'$ et un $C'$-module projectif de type fini $Q$ tel
que $Q/IQ$ soit isomorphe à $Q_0 \otimes_{A_0} C/IC$.
Alors $D = \text{Sym}_{C'}^*(Q)$ est une $C'$-algèbre lisse (voir
Compléments d'algèbre, Lemme \ref{more-algebra-lemma-symmetric-algebra-smooth}).
Considérons le diagramme
$$
\xymatrix{
R \ar[d] \ar[rr] & &
C \ar[r] \ar[d] &
C' \ar[r] \ar[d] &
D \ar[d] \\
R/I \ar[r] &
A_0 \ar[r] &
C/IC \ar[r]^{\cong} &
C'/IC' \ar[r] &
D/ID
}
$$
Observons que notre choix de $Q$ donne
\begin{align*}
D/ID & =
\text{Sym}_{C/IC}^*(Q_0 \otimes_{A_0} C/IC) \\
& =
\text{Sym}_{A_0}^*(Q_0) \otimes_{A_0} C/IC \\
& =
\text{Sym}_{A_0}^*(Q_0) \otimes_{A_0}
\text{Sym}_{A_0}^*(P_0) \\
& =
\text{Sym}_{A_0}^*(Q_0 \oplus P_0) \\
& =
\text{Sym}_{A_0}^*(A_0^{\oplus n}) \\
& =
A_0[x_1, \ldots, x_n]
\end{align*}
Choisissons $f_1, \ldots, f_n \in D$ dont les images sont
$x_1, \ldots, x_n$ dans $D/ID = A_0[x_1, \ldots, x_n]$. Posons
$A = D/(f_1, \ldots, f_n)$. Notons que $A_0 = A/IA$. Nous affirmons que
$R \to A$ est syntomique dans un voisinage de $V(IA)$. Si cette affirmation est
vraie, nous pouvons trouver un $f \in A$ d'image $1 \in A_0$ tel que $A_f$ soit
syntomique sur $R$, ce qui achève la démonstration de (1).

\medskip\noindent
Démonstration de l'affirmation. Observons que $R \to D$ est syntomique en tant
que composée du morphisme d'anneaux syntomique $R \to C$, du morphisme
d'anneaux \'etale $C \to C'$ et du morphisme d'anneaux lisse $C' \to D$
(Algèbre, Lemmes
\ref{algebra-lemma-composition-syntomic} et
\ref{algebra-lemma-smooth-syntomic}).
La question est locale sur $\Spec(D)$ ; nous pouvons donc supposer que $D$ est
une intersection complète globale relative
(Algèbre, Lemme \ref{algebra-lemma-syntomic}).
Écrivons $D = R[y_1, \ldots, y_m]/(g_1, \ldots, g_s)$.
Soient $f'_1, \ldots, f'_n \in R[y_1, \ldots, y_m]$ des relèvements de
$f_1, \ldots, f_n$. Nous pouvons alors appliquer
Algèbre, Lemme \ref{algebra-lemma-localize-relative-complete-intersection}
pour obtenir l'affirmation.

\medskip\noindent
Démonstration de (2). Comme tout morphisme d'anneaux lisse est syntomique, nous
pouvons trouver un morphisme d'anneaux syntomique $R \to A$ tel que
$A_0 = A/IA$. Par hypothèse, les fibres de $R \to A$ sont lisses au-dessus des
idéaux premiers de $V(I)$ ; ainsi, $R \to A$ est lisse dans un voisinage ouvert
de $V(IA)$
(Algèbre, Lemme \ref{algebra-lemma-flat-fibre-smooth}).
Nous pouvons donc remplacer $A$ par une localisation pour obtenir le résultat
voulu.
\end{proof}

\noindent
Nous savons que tout morphisme d'anneaux syntomique $R \to A$ est localement
une intersection complète globale relative, voir
Algèbre, Lemme \ref{algebra-lemma-syntomic}.
Le lemme suivant affirme qu'un fibré vectoriel sur $\Spec(A)$ est une
intersection complète globale relative.

\begin{lemma}
\label{smoothing-lemma-syntomic-complete-intersection}
Soit $R \to A$ un morphisme d'anneaux syntomique. Il existe alors un morphisme
lisse de $R$-algèbres $A \to C$ admettant une rétraction tel que $C$ soit une
intersection complète globale relative sur $R$, c'est-à-dire
$$
C \cong R[x_1, \ldots, x_n]/(f_1, \ldots, f_c)
$$
plate sur $R$ et dont toutes les fibres sont de dimension $n - c$.
\end{lemma}

\begin{proof}
Appliquons le lemme \ref{smoothing-lemma-improve-presentation} pour obtenir $A \to C$.
D'après Algèbre, Lemme \ref{algebra-lemma-huber}, nous pouvons écrire
$C = R[x_1, \ldots, x_n]/(f_1, \ldots, f_c)$ de sorte que les images des $f_i$
forment une base de $J/J^2$.
Le morphisme d'anneaux $R \to C$ est syntomique (donc plat), car il est la
composée d'un morphisme d'anneaux syntomique et d'un morphisme d'anneaux lisse.
La dimension des fibres est $n - c$ d'après
Algèbre, Lemme \ref{algebra-lemma-lci}
(les fibres sont des intersections complètes locales, le lemme s'applique donc).
\end{proof}

\begin{lemma}
\label{smoothing-lemma-smooth-standard-smooth}
Soit $R \to A$ un morphisme d'anneaux lisse. Il existe alors un morphisme lisse
de $R$-algèbres $A \to B$ admettant une rétraction tel que $B$ soit standard
lisse sur $R$, c'est-à-dire
$$
B \cong R[x_1, \ldots, x_n]/(f_1, \ldots, f_c)
$$
et $\det(\partial f_j/\partial x_i)_{i, j = 1, \ldots, c}$ est inversible dans
$B$.
\end{lemma}

\begin{proof}
Appliquons le lemme \ref{smoothing-lemma-syntomic-complete-intersection} pour obtenir un
morphisme lisse de $R$-algèbres $A \to C$ admettant une rétraction, tel que
$C = R[x_1, \ldots, x_n]/(f_1, \ldots, f_c)$ soit une intersection complète
globale relative sur $R$. Comme $C$ est lisse sur $R$, nous avons une suite
exacte courte
$$
0 \to
\bigoplus\nolimits_{j = 1, \ldots, c} C f_j \to
\bigoplus\nolimits_{i = 1, \ldots, n} C\text{d}x_i \to
\Omega_{C/R} \to 0
$$
Puisque $\Omega_{C/R}$ est un $C$-module projectif, cette suite est scindée.
Choisissons un inverse à gauche $t$ du premier morphisme. Écrivons
$t(\text{d}x_i) = \sum c_{ij} f_j$
de sorte que $\sum_i \frac{\partial f_j}{\partial x_i} c_{i\ell} = \delta_{j\ell}$
(symbole de Kronecker). Soit
$$
B' = C[y_1, \ldots, y_c] =
R[x_1, \ldots, x_n, y_1, \ldots, y_c]/(f_1, \ldots, f_c)
$$
Le morphisme de $R$-algèbres $C \to B'$ admet la rétraction qui envoie $y_j$
sur zéro. Nous affirmons que le morphisme
$$
R[z_1, \ldots, z_n] \longrightarrow B',\quad
z_i \longmapsto x_i - \sum\nolimits_j c_{ij} y_j
$$
est \'etale en tout point de l'image de $\Spec(C) \to \Spec(B')$.
Dans $\Omega_{B'/R[z_1, \ldots, z_n]}$, nous avons
$$
0 =
\text{d}f_j - \sum\nolimits_i \frac{\partial f_j}{\partial x_i} \text{d}z_i
\equiv
\sum\nolimits_{i, \ell}
\frac{\partial f_j}{\partial x_i} c_{i\ell} \text{d}y_\ell
\equiv
\text{d}y_j \bmod (y_1, \ldots, y_c)\Omega_{B'/R[z_1, \ldots, z_n]}
$$
Puisque $0 = \text{d}z_i = \text{d}x_i$ modulo
$\sum B'\text{d}y_j + (y_1, \ldots, y_c)\Omega_{B'/R[z_1, \ldots, z_n]}$,
nous en concluons que
$$
\Omega_{B'/R[z_1, \ldots, z_n]}/
(y_1, \ldots, y_c)\Omega_{B'/R[z_1, \ldots, z_n]} = 0.
$$
Comme $\Omega_{B'/R[z_1, \ldots, z_n]}$ est un $B'$-module de type fini, le
lemme de Nakayama montre qu'il existe un $g \in 1 + (y_1, \ldots, y_c)$ tel que
$(\Omega_{B'/R[z_1, \ldots, z_n]})_g = 0$. Cela démontre que
$R[z_1, \ldots, z_n] \to B'_g$ est non ramifié, voir
Algèbre, Définition \ref{algebra-definition-unramified}.
Pour tout morphisme d'anneaux $R \to k$, où $k$ est un corps, nous obtenons un
morphisme d'anneaux non ramifié
$k[z_1, \ldots, z_n] \to (B'_g) \otimes_R k$ entre $k$-algèbres lisses de
dimension $n$. Il en résulte que
$k[z_1, \ldots, z_n] \to (B'_g) \otimes_R k$ est plat d'après
Algèbre, Lemmes \ref{algebra-lemma-CM-over-regular-flat} et
\ref{algebra-lemma-characterize-smooth-kbar}. Par le crit\`ere
de platitude par fibre
(Algèbre, Lemme \ref{algebra-lemma-criterion-flatness-fibre}),
nous concluons que $R[z_1, \ldots, z_n] \to B'_g$ est plat.
Enfin, Algèbre, Lemme \ref{algebra-lemma-characterize-etale} implique que
$R[z_1, \ldots, z_n] \to B'_g$ est \'etale.
Posons $B = B'_g$. Notons que $C \to B$ est lisse et admet une rétraction ; il
en est donc de même de $A \to B$. De plus,
$R[z_1, \ldots, z_n] \to B$ est \'etale.
D'après Algèbre, Lemme \ref{algebra-lemma-etale-standard-smooth}, nous pouvons
écrire
$$
B = R[z_1, \ldots, z_n, w_1, \ldots, w_m]/(g_1, \ldots, g_m)
$$
avec $\det(\partial g_j/\partial w_i)$ inversible dans $B$.
Cela démontre le lemme.
\end{proof}

\begin{lemma}
\label{smoothing-lemma-colimit-standard-smooth}
Soit $R \to \Lambda$ un morphisme d'anneaux. Si $\Lambda$ est une colimite
filtrante de $R$-algèbres lisses, alors $\Lambda$ est une colimite filtrante de
$R$-algèbres standard lisses.
\end{lemma}

\begin{proof}
Soit $A \to \Lambda$ un morphisme de $R$-algèbres, avec $A$ de présentation
finie sur $R$. D'après
Algèbre, Lemme \ref{algebra-lemma-when-colimit}
nous devons factoriser ce morphisme par une algèbre standard lisse, et nous
savons que nous pouvons le factoriser sous la forme $A \to B \to \Lambda$, où
$B$ est lisse sur $R$. Choisissons un morphisme de $R$-algèbres $B \to C$
admettant une rétraction $C \to B$, tel que $C$ soit standard lisse sur $R$, voir
Lemme \ref{smoothing-lemma-smooth-standard-smooth}.
La factorisation cherchée est alors $A \to B \to C \to B \to \Lambda$.
\end{proof}

\begin{lemma}
\label{smoothing-lemma-standard-smooth-include-generators}
Soit $R \to A$ un morphisme d'anneaux standard lisse.
Soit $E \subset A$ un sous-ensemble fini de cardinal $|E| = n$.
Il existe alors une présentation
$A = R[x_1, \ldots, x_{n + m}]/(f_1, \ldots, f_c)$ avec $c \geq n$, dans
laquelle $\det(\partial f_j/\partial x_i)_{i, j = 1, \ldots, c}$ est inversible
dans $A$ et telle que $E$ soit l'ensemble des classes, dans ce quotient, de
$x_1, \ldots, x_n$.
\end{lemma}

\begin{proof}
Choisissons une présentation $A = R[y_1, \ldots, y_m]/(g_1, \ldots, g_d)$
telle que l'image de $\det(\partial g_j/\partial y_i)_{i, j = 1, \ldots, d}$
soit inversible dans $A$. Choisissons une énumération
$E = \{a_1, \ldots, a_n\}$ et des éléments $h_i \in R[y_1, \ldots, y_m]$ dont
l'image dans $A$ est $a_i$. Considérons la présentation
$$
A = R[x_1, \ldots, x_n, y_1, \ldots, y_m]/
(x_1 - h_1, \ldots, x_n - h_n, g_1, \ldots, g_d)
$$
et posons $c = n + d$.
\end{proof}

\begin{lemma}
\label{smoothing-lemma-compare-standard}
Soit $R \to A$ un morphisme d'anneaux de présentation finie.
Soit $a \in A$. Considérons les conditions suivantes sur $a$ :
\begin{enumerate}
\item $A_a$ est lisse sur $R$,
\item $A_a$ est lisse sur $R$ et $\Omega_{A_a/R}$ est stablement libre,
\item $A_a$ est lisse sur $R$ et $\Omega_{A_a/R}$ est libre,
\item $A_a$ est standard lisse sur $R$,
\item $a$ est strictement standard dans $A$ sur $R$,
\item $a$ est standard élémentaire dans $A$ sur $R$.
\end{enumerate}
Alors :
\begin{enumerate}
\item[(a)] (4) $\Rightarrow$ (3) $\Rightarrow$ (2) $\Rightarrow$ (1),
\item[(b)] (6) $\Rightarrow$ (5),
\item[(c)] (6) $\Rightarrow$ (4),
\item[(d)] (5) $\Rightarrow$ (2),
\item[(e)] (2) $\Rightarrow$ les éléments $a^e$, $e \geq e_0$, sont
strictement standard dans $A$ sur $R$,
\item[(f)] (4) $\Rightarrow$ les éléments $a^e$, $e \geq e_0$, sont standard
élémentaires dans $A$ sur $R$.
\end{enumerate}
\end{lemma}

\begin{proof}
D'après les définitions et
Algèbre, Lemme \ref{algebra-lemma-standard-smooth},
la partie (a) est claire. La partie (b) résulte immédiatement de
la définition \ref{smoothing-definition-strictly-standard}.

\medskip\noindent
Démonstration de (c). Choisissons une présentation
$A = R[x_1, \ldots, x_n]/(f_1, \ldots, f_m)$ telle que
(\ref{smoothing-equation-elementary-standard-one}) et
(\ref{smoothing-equation-elementary-standard-two}) soient satisfaites.
Choisissons $h \in R[x_1, \ldots, x_n]$ d'image $a$. Alors
$$
A_a = R[x_0, x_1, \ldots, x_n]/(x_0h - 1, f_1, \ldots, f_m).
$$
Écrivons $J = (x_0h - 1, f_1, \ldots, f_m)$.
D'après (\ref{smoothing-equation-elementary-standard-two}), nous voyons que le
$A_a$-module $J/J^2$ est engendré par $x_0h - 1, f_1, \ldots, f_c$ sur
$A_a$. Ainsi, comme dans la démonstration de
Algèbre, Lemme \ref{algebra-lemma-huber}, nous pouvons choisir un $g \in 1 + J$
tel que
$$
A_a = R[x_0, \ldots, x_n, x_{n + 1}]/
(x_0h - 1, f_1, \ldots, f_c, gx_{n + 1} - 1).
$$
À ce stade, (\ref{smoothing-equation-elementary-standard-one}) implique que $R \to A_a$
est standard lisse (utiliser les coordonnées
$x_0, x_1, \ldots, x_c, x_{n + 1}$ pour calculer les dérivées).

\medskip\noindent
Démonstration de (d). Choisissons une présentation
$A = R[x_1, \ldots, x_n]/(f_1, \ldots, f_m)$ telle que
(\ref{smoothing-equation-strictly-standard-one}) et
(\ref{smoothing-equation-strictly-standard-two}) soient satisfaites.
Écrivons $I = (f_1, \ldots, f_m)$.
Nous savons déjà que $A_a$ est lisse sur $R$, voir
Lemme \ref{smoothing-lemma-elkik}.
D'après le lemme \ref{smoothing-lemma-parse-equation-strictly-standard-one}, nous voyons que
$(I/I^2)_a$ est libre de base $f_1, \ldots, f_c$ et s'envoie isomorphiquement
sur un facteur direct de $\bigoplus A_a \text{d}x_i$. Puisque
$\Omega_{A_a/R} = (\Omega_{A/R})_a$
est le conoyau du morphisme
$(I/I^2)_a \to \bigoplus A_a \text{d}x_i$,
nous concluons qu'il est stablement libre.

\medskip\noindent
Démonstration de (e). Choisissons une présentation
$A = R[x_1, \ldots, x_n]/I$ avec $I$ de type fini.
Par hypothèse, nous avons une suite exacte courte
$$
0 \to (I/I^2)_a \to \bigoplus\nolimits_{i = 1, \ldots, n} A_a\text{d}x_i \to
\Omega_{A_a/R} \to 0
$$
qui est exacte scindée. Nous voyons donc que
$(I/I^2)_a \oplus \Omega_{A_a/R}$ est un $A_a$-module libre.
Puisque $\Omega_{A_a/R}$ est stablement libre, nous voyons que $(I/I^2)_a$
l'est également. Ainsi, en remplaçant la présentation choisie ci-dessus par
$A = R[x_1, \ldots, x_n, x_{n + 1}, \ldots, x_{n + r}]/J$, avec
$J = (I, x_{n + 1}, \ldots, x_{n + r})$ pour un certain $r$, nous obtenons que
$(J/J^2)_a$ est libre de type fini. Choisissons
$f_1, \ldots, f_c \in J$ dont les images forment une base de $(J/J^2)_a$.
Complétons cette famille en une famille de générateurs
$f_1, \ldots, f_m \in J$. Considérons la présentation
$A = R[x_1, \ldots, x_{n + r}]/(f_1, \ldots, f_m)$.
Alors, par construction, (\ref{smoothing-equation-strictly-standard-two}) est satisfaite
pour $a^e$ dès que $e$ est suffisamment grand. De plus, puisque
$(J/J^2)_a \to \bigoplus\nolimits_{i = 1, \ldots, n + r} A_a\text{d}x_i$
est une injection scindée, nous pouvons trouver un inverse à gauche
$A_a$-linéaire. En exprimant cet inverse à gauche dans la base
$f_1, \ldots, f_c$ et en chassant les dénominateurs, nous trouvons un morphisme
linéaire
$\psi_0 : A^{\oplus n + r} \to A^{\oplus c}$ tel que
$$
A^{\oplus c} \xrightarrow{(f_1, \ldots, f_c)}
J/J^2 \xrightarrow{f \mapsto \text{d}f}
\bigoplus\nolimits_{i = 1, \ldots, n + r} A \text{d}x_i
\xrightarrow{\psi_0}
A^{\oplus c}
$$
soit la multiplication par $a^{e_0}$ pour un certain $e_0 \geq 1$. D'après
Lemme \ref{smoothing-lemma-parse-equation-strictly-standard-one},
nous voyons que (\ref{smoothing-equation-strictly-standard-one}) est satisfaite pour
$a^{ce_0}$, et donc pour $a^e$ pour tout $e$ tel que $e \geq ce_0$.

\medskip\noindent
Démonstration de (f). Choisissons une présentation
$A_a = R[x_1, \ldots, x_n]/(f_1, \ldots, f_c)$ telle que
$\det(\partial f_j/\partial x_i)_{i, j = 1, \ldots, c}$
soit inversible dans $A_a$. Nous pouvons supposer que, pour un certain $m < n$,
les classes des éléments $x_1, \ldots, x_m$ correspondent aux $a_i/1$, où
$a_1, \ldots, a_m \in A$ sont des générateurs de $A$ sur $R$, voir
Lemme \ref{smoothing-lemma-standard-smooth-include-generators}.
Après avoir remplacé $x_i$ par $a^Nx_i$ pour $m < i \leq n$, nous pouvons
supposer que la classe de $x_i$ est $a_i/1 \in A_a$ pour un certain $a_i \in A$.
Considérons le morphisme d'anneaux
$$
\Psi : R[x_1, \ldots, x_n] \longrightarrow A,\quad
x_i \longmapsto a_i.
$$
Il s'agit d'un morphisme d'anneaux surjectif. En remplaçant $f_j$ par $a^Nf_j$,
nous pouvons supposer que $f_j \in R[x_1, \ldots, x_n]$ et que
$\Psi(f_j) = 0$ (car, après tout, $f_j(a_1/1, \ldots, a_n/1) = 0$ dans $A_a$).
Soit $J = \Ker(\Psi)$. Alors $A = R[x_1, \ldots, x_n]/J$ est une présentation,
et $f_1, \ldots, f_c \in J$ sont des éléments tels que $(J/J^2)_a$ soit libre
de base $f_1, \ldots, f_c$ et que
$\det(\partial f_j/\partial x_i)_{i, j = 1, \ldots, c}$ s'envoie sur un élément
inversible de $A_a$. Il en résulte que
(\ref{smoothing-equation-elementary-standard-one}) et
(\ref{smoothing-equation-elementary-standard-two})
sont satisfaites pour $a^e$ et tout $e$ suffisamment grand, comme souhaité.
\end{proof}
\section{Intermède : désingularisation de Néron}
\label{smoothing-section-neron}

\noindent
Nous interrompons l'étude du cas général du théorème de Popescu pour
examiner un cas plus facile mais déjà très intéressant, à savoir celui où
$R \to \Lambda$ est un homomorphisme d'anneaux de valuation discrète.
Ce cas est étudié dans
\cite[section 4]{Artin-Algebraic-Approximation}.

\begin{situation}
\label{smoothing-situation-neron}
Ici, $R \subset \Lambda$ est une extension d'anneaux de valuation discrète
d'indice de ramification $1$ (Compléments d'algèbre, définition
\ref{more-algebra-definition-extension-discrete-valuation-rings}).
Nous nous donnons une factorisation
$$
R \to A \xrightarrow{\varphi} \Lambda
$$
où $R \to A$ est plat et de type fini. Posons $\mathfrak q = \Ker(\varphi)$
et $\mathfrak p = \varphi^{-1}(\mathfrak m_\Lambda)$.
\end{situation}

\noindent
Dans la situation \ref{smoothing-situation-neron}, soit $\pi \in R$ une uniformisante.
Rappelons que la platitude de $A$ sur $R$ signifie que $\pi$
est un non-diviseur de zéro dans $A$
(Compléments d'algèbre, lemme
\ref{more-algebra-lemma-valuation-ring-torsion-free-flat}).
Notre hypothèse sur $R \subset \Lambda$ montre que l'image de $\pi$
est une uniformisante de $\Lambda$. Comme $\pi \in \mathfrak p$, nous pouvons considérer
l'algèbre d'éclatement affine de Néron (voir
Algèbre, section \ref{algebra-section-blow-up})
$$
\varphi' :
A' = A[\textstyle{\frac{\mathfrak p}{\pi}}]
\longrightarrow
\Lambda
$$
munie du morphisme induit vers $\Lambda$ qui envoie
$a/\pi^n$, $a \in \mathfrak p^n$, sur $\pi^{-n}\varphi(a)$ dans $\Lambda$.
Nous noterons $\mathfrak q' \subset \mathfrak p' \subset A'$
les idéaux premiers correspondants de $A'$. Remarquons que la classe
d'isomorphisme de $A'$ ne dépend pas du choix de l'uniformisante.
En répétant cette construction, nous obtenons une suite
$$
A \to A' \to A'' \to \ldots \to \Lambda
$$

\begin{lemma}
\label{smoothing-lemma-neron-functorial}
Dans la situation \ref{smoothing-situation-neron}, l'éclatement de Néron est fonctoriel
au sens suivant :
\begin{enumerate}
\item si $a \in A$, $a \not \in \mathfrak p$, alors l'éclatement de Néron
de $A_a$ est $A'_a$, et
\item si $B \to A$ est une surjection de $R$-algèbres plates de type fini
de noyau $I$, alors $A'$ est le quotient de $B'/IB'$ par sa
torsion annulée par une puissance de $\pi$.
\end{enumerate}
\end{lemma}

\begin{proof}
Les assertions (1) et (2) sont des cas particuliers de
Algèbre, lemme \ref{algebra-lemma-blowup-base-change}.
En effet, dès que nous avons $A_1 \to A_2 \to \Lambda$ avec
$\mathfrak p_1 A_2 = \mathfrak p_2$, l'algèbre $A_2'$ est
le quotient de $A_1' \otimes_{A_1} A_2$ par sa torsion annulée par une puissance de $\pi$.
\end{proof}

\begin{lemma}
\label{smoothing-lemma-neron-blowup-smooth}
Dans la situation \ref{smoothing-situation-neron}, supposons que $R \to A$ soit lisse
en $\mathfrak p$ et que $R/\pi R \subset \Lambda/\pi \Lambda$
soit une extension séparable de corps. Alors $R \to A'$ est lisse en
$\mathfrak p'$ et il existe une suite exacte courte
$$
0 \to
\Omega_{A/R} \otimes_A A'_{\mathfrak p'} \to
\Omega_{A'/R, \mathfrak p'} \to
(A'/\pi A')_{\mathfrak p'}^{\oplus c} \to 0
$$
où $c = \dim((A/\pi A)_\mathfrak p)$.
\end{lemma}

\begin{proof}
D'après le lemme \ref{smoothing-lemma-neron-functorial}, nous pouvons remplacer $A$ par
un localisé en un élément qui n'appartient pas à $\mathfrak p$ ; nous le ferons
désormais sans autre mention. Posons $\kappa = R/\pi R$. Comme la lissité est
stable par changement de base
(Algèbre, lemme \ref{algebra-lemma-base-change-smooth}),
nous voyons que $A/\pi A$ est lisse sur $\kappa$ en $\mathfrak p$.
Par conséquent, $(A/\pi A)_\mathfrak p$ est un anneau local régulier
(Algèbre, lemme \ref{algebra-lemma-characterize-smooth-over-field}).
Choisissons $g_1, \ldots, g_c \in \mathfrak p$ dont les images forment
un système régulier de paramètres de $(A/\pi A)_\mathfrak p$.
Nous voyons alors que $\mathfrak p = (\pi, g_1, \ldots, g_c)$,
quitte à remplacer $A$ par un localisé.
Remarquons que $\pi, g_1, \ldots, g_c$ est une suite régulière
dans $A_\mathfrak p$ (d'abord $\pi$ est un non-diviseur de zéro, puis on applique
Algèbre, lemme \ref{algebra-lemma-regular-ring-CM}
au reste de la suite).
Après avoir remplacé $A$ par un localisé, nous pouvons supposer que
$\pi, g_1, \ldots, g_c$ est une suite régulière dans $A$
(Algèbre, lemme \ref{algebra-lemma-regular-sequence-in-neighbourhood}).
Il s'ensuit que
$$
A' = A[y_1, \ldots, y_c]/(\pi y_1 - g_1, \ldots, \pi y_c - g_c) =
A[y_1, \ldots, y_c]/I
$$
d'après Compléments d'algèbre, lemme \ref{more-algebra-lemma-blowup-regular-sequence}.
Dans la suite, nous utiliserons sans autre mention la définition de la lissité
au moyen du complexe cotangent naïf (Algèbre, définition \ref{algebra-definition-smooth})
et le critère d'Algèbre, lemme \ref{algebra-lemma-smooth-at-point}.
La suite exacte d'Algèbre, lemme \ref{algebra-lemma-exact-sequence-NL},
pour $R \to A[y_1, \ldots, y_c] \to A'$, s'écrit
$$
0 \to H_1(\NL_{A'/R}) \to I/I^2 \to
\Omega_{A/R} \otimes_A A' \oplus
\bigoplus\nolimits_{i = 1, \ldots, c} A' \text{d}y_i \to
\Omega_{A'/R} \to 0
$$
où la classe de $\pi y_i - g_i$ dans $I/I^2$ est envoyée
sur $- \text{d}g_i + \pi \text{d}y_i$ dans le terme suivant.
Ici, nous avons utilisé Algèbre, lemme \ref{algebra-lemma-NL-surjection},
pour calculer $\NL_{A'/A[y_1, \ldots, y_c]}$, ainsi que le fait que
$R \to A[y_1, \ldots, y_c]$ est lisse, de sorte que
$H_1(\NL_{A[y_1, \ldots, y_c]/R}) = 0$ et
$\Omega_{A[y_1, \ldots, y_c]/R}$ est un
$A[y_1, \ldots, y_c]$-module projectif fini (donc, a fortiori, plat),
qui est en fait la somme directe de
$\Omega_{A/R} \otimes_A A[y_1, \ldots, y_c]$ et d'un module libre
de base les $\text{d}y_i$. Pour achever la démonstration, il suffit de
montrer que $\text{d}g_1, \ldots, \text{d}g_c$ font partie d'une base du
module libre fini $\Omega_{A/R, \mathfrak p}$.
En effet, cela montrera que $(I/I^2)_{\mathfrak p'}$ est libre sur les
$\pi y_i - g_i$, que le localisé en $\mathfrak p'$ du morphisme central
de la suite est injectif, donc que $H_1(\NL_{A'/R})_{\mathfrak p'} = 0$, et que
le conoyau $\Omega_{A'/R, \mathfrak p'}$ est libre fini.
Pour cela, il suffit de montrer que les images des $\text{d}g_i$ sont
linéairement indépendantes sur $\kappa(\mathfrak p)$ dans
$\Omega_{A/R, \mathfrak p}/\pi = \Omega_{(A/\pi A)/\kappa, \mathfrak p}$
(l'égalité résulte d'Algèbre, lemme \ref{algebra-lemma-differentials-base-change}).
Comme $\kappa \subset \kappa(\mathfrak p) \subset \Lambda/\pi \Lambda$,
nous voyons que $\kappa(\mathfrak p)$ est séparable sur $\kappa$
(Algèbre, définition \ref{algebra-definition-separable-field-extension}).
L'indépendance linéaire recherchée résulte alors d'Algèbre,
lemme \ref{algebra-lemma-computation-differential}.
\end{proof}

\begin{lemma}
\label{smoothing-lemma-neron-when-smooth}
Dans la situation \ref{smoothing-situation-neron}, supposons que $R \to A$ soit lisse
en $\mathfrak q$ et que nous ayons une surjection de $R$-algèbres
$B \to A$ de noyau $I$. Supposons $R \to B$ lisse en
$\mathfrak p_B = (B \to A)^{-1}\mathfrak p$. Si le conoyau de
$$
I/I^2 \otimes_A \Lambda \to \Omega_{B/R} \otimes_B \Lambda
$$
est un $\Lambda$-module libre, alors $R \to A$ est lisse en $\mathfrak p$.
\end{lemma}

\begin{proof}
Le conoyau du morphisme $I/I^2 \to \Omega_{B/R} \otimes_B A$
est $\Omega_{A/R}$, voir Algèbre, lemme \ref{algebra-lemma-differential-seq}.
Soit $d = \dim_\mathfrak q(A/R)$ la dimension relative de $R \to A$
en $\mathfrak q$, c'est-à-dire la dimension de $\Spec(A[1/\pi])$ en $\mathfrak q$.
Voir Algèbre, définition \ref{algebra-definition-relative-dimension}.
Alors $\Omega_{A/R, \mathfrak q}$ est libre sur $A_\mathfrak q$ et de rang $d$
(Algèbre, lemme \ref{algebra-lemma-characterize-smooth-over-field}).
Ainsi, si l'hypothèse du lemme est satisfaite,
$\Omega_{A/R} \otimes_A \Lambda$ est libre de rang $d$.
Il s'ensuit que $\Omega_{A/R} \otimes_A \kappa(\mathfrak p)$
est de dimension $d$ (car c'est le cas après tensorisation par $\Lambda/\pi \Lambda$).
Comme $R \to A$ est plat et que $\mathfrak p$
est une spécialisation de $\mathfrak q$, nous voyons que
$\dim_\mathfrak p(A/R) \geq d$ d'après Algèbre, lemme
\ref{algebra-lemma-dimension-fibres-bounded-open-upstairs}.
Il en résulte que $R \to A$ est lisse en $\mathfrak p$ d'après
Algèbre, lemmes \ref{algebra-lemma-flat-fibre-smooth} et
\ref{algebra-lemma-characterize-smooth-over-field}.
\end{proof}

\begin{lemma}
\label{smoothing-lemma-neron-desingularization}
Dans la situation \ref{smoothing-situation-neron},
supposons que $R \to A$ soit lisse en $\mathfrak q$
et que $R/\pi R \subset \Lambda/\pi \Lambda$ soit une extension
séparable de corps. Après un nombre fini d'éclatements affines de Néron,
l'algèbre $A$ devient alors lisse sur $R$ en $\mathfrak p$.
\end{lemma}

\begin{proof}
Choisissons une $R$-algèbre $B$ et une surjection $B \to A$. Posons
$\mathfrak p_B = (B \to A)^{-1}(\mathfrak p)$ et notons $r$
la dimension relative de $R \to B$ en $\mathfrak p_B$. Nous choisissons $B$
de sorte que $R \to B$ soit lisse en $\mathfrak p_B$.
Par exemple, nous pouvons prendre pour $B$ une algèbre de polynômes en $r$
variables sur $R$. Considérons le complexe
$$
I/I^2 \otimes_A \Lambda \longrightarrow \Omega_{B/R} \otimes_B \Lambda
$$
du lemme \ref{smoothing-lemma-neron-when-smooth}. La structure des modules finis sur
$\Lambda$ (Compléments d'algèbre, lemme \ref{more-algebra-lemma-modules-PID})
montre que le conoyau est de la forme
$$
\Lambda^{\oplus d} \oplus
\bigoplus\nolimits_{i = 1, \ldots, n} \Lambda/\pi^{e_i} \Lambda
$$
pour certains $d \geq 0$, $n \geq 0$ et $e_i \geq 1$. Remarquons que $d$
est la dimension relative de $A/R$ en $\mathfrak q$
(Algèbre, lemme \ref{algebra-lemma-characterize-smooth-over-field}).
Si le défaut $e = \sum_{i = 1, \ldots, n} e_i$ est nul, la conclusion résulte du
lemme \ref{smoothing-lemma-neron-when-smooth}.

\medskip\noindent
Examinons maintenant ce qui se passe lorsque nous effectuons l'éclatement de Néron.
Rappelons que $A'$ est le quotient de $B'/IB'$ par sa torsion annulée par
une puissance de $\pi$ (lemme \ref{smoothing-lemma-neron-functorial}) et que $R \to B'$ est lisse en
$\mathfrak p_{B'}$ (lemme \ref{smoothing-lemma-neron-blowup-smooth}).
Après éclatement, nous nous retrouvons donc exactement dans la même situation. Représentons-la par
$$
\xymatrix{
0 \ar[r] & I' \ar[r] & B' \ar[r] & A' \ar[r] & 0 \\
0 \ar[r] & I \ar[u] \ar[r] & B \ar[u] \ar[r] & A \ar[r] \ar[u] & 0
}
$$
Comme $I \subset \mathfrak p_B$, nous voyons que $I \to I'$
se factorise par $\pi I'$. En considérant le morphisme induit de
complexes, nous obtenons
$$
\xymatrix{
I'/(I')^2 \otimes_{A'} \Lambda \ar[r] &
\Omega_{B'/R} \otimes_{B'} \Lambda \ar@{=}[r] & M' \\
I/I^2 \otimes_A \Lambda \ar[r] \ar[u] &
\Omega_{B/R} \otimes_B \Lambda \ar[u] \ar@{=}[r] & M
}
$$
Alors $M \subset M'$ sont des $\Lambda$-modules libres finis, et le quotient
$M'/M$ est annulé par $\pi$, voir le lemme \ref{smoothing-lemma-neron-blowup-smooth}.
Soient $N \subset M$ et $N' \subset M'$ les images des morphismes horizontaux,
et posons $Q = M/N$ et $Q' = M'/N'$.
Nous obtenons un diagramme commutatif
$$
\xymatrix{
0 \ar[r] &
N' \ar[r] &
M' \ar[r] &
Q' \ar[r] &
0 \\
0 \ar[r] &
N \ar[r] \ar[u] &
M \ar[r] \ar[u] &
Q \ar[r] \ar[u] &
0
}
$$
Alors $N \subset N'$ sont des $\Lambda$-modules libres de rang $r - d$.
Comme $I$ s'envoie dans $\pi I'$, nous voyons que $N \subset \pi N'$.

\medskip\noindent
Soit $K = \Lambda_\pi$ le corps des fractions de $\Lambda$.
Nous avons un diagramme commutatif
$$
\xymatrix{
0 \ar[r] &
N' \ar[r] &
N'_K \cap M' \ar[r] &
Q'_{tor} \ar[r] &
0 \\
0 \ar[r] &
N \ar[r] \ar[u] &
N_K \cap M \ar[r] \ar[u] &
Q_{tor} \ar[r] \ar[u] &
0
}
$$
dont les lignes sont des suites exactes courtes. Cela montre que la variation du défaut
est donnée par
$$
e - e' =
\text{longueur}(Q_{tor}) - \text{longueur}(Q'_{tor})
=
\text{longueur}(N'/N) - \text{longueur}(N'_K \cap M' / N_K \cap M)
$$
Comme $M'/M$ est annulé par $\pi$, il en va de même de $N'_K \cap M' / N_K \cap M$,
et sa longueur est au plus $\dim_K(N_K)$.
Comme $N \subset \pi N'$, nous obtenons $\text{longueur}(N'/N) \ge \dim_K(N_K)$,
avec égalité si et seulement si $N = \pi N'$.

\medskip\noindent
Pour achever la démonstration, nous devons montrer que $N$ est strictement contenu
dans $\pi N'$ lorsque $A$ n'est pas lisse en $\mathfrak p$ ;
c'est le calcul clé de l'argument de Néron.
Pour cela, considérons la suite exacte
$$
I/I^2 \otimes_B \kappa(\mathfrak p_B)
\to \Omega_{B/R} \otimes_B \kappa(\mathfrak p_B)
\to \Omega_{A/R} \otimes_A \kappa(\mathfrak p) \to 0
$$
(qui résulte d'Algèbre, lemme \ref{algebra-lemma-differential-seq}).
Comme $R \to A$ n'est pas lisse en $\mathfrak p$, nous voyons que la dimension $s$ de
$\Omega_{A/R} \otimes_A \kappa(\mathfrak p)$
est strictement supérieure à $d$. D'autre part,
la première flèche se factorise par le morphisme injectif
$$
\mathfrak p_B B_{\mathfrak p_B}/\mathfrak p_B^2 B_{\mathfrak p_B}
\to \Omega_{B/R} \otimes_B \kappa(\mathfrak p_B)
$$
d'Algèbre, lemme \ref{algebra-lemma-computation-differential} ;
remarquons que $\kappa(\mathfrak p)$ est séparable sur $\kappa$
d'après notre hypothèse sur $R/\pi R \subset \Lambda/\pi \Lambda$.
Nous en concluons qu'il existe des générateurs
$g_1, \ldots, g_t \in I$ tels que $g_j \in \mathfrak p^2$
pour $j > r - s$. Les images des $g_j$ dans $A'$ appartiennent alors à $\pi^2 I'$
pour $j > r - s$. Comme $r - s < r - d$,
au moins un des générateurs minimaux
de $N$ devient divisible par $\pi^2$ dans $N'$.
Ainsi, $e$ diminue d'au moins $1$, ce qui conclut la démonstration.
\end{proof}

\noindent
Si $R \to \Lambda$ est une extension d'anneaux de valuation discrète,
alors $R \to \Lambda$ est régulier si et seulement si
(a) l'indice de ramification vaut $1$,
(b) l'extension des corps de fractions est séparable, et
(c) $R/\mathfrak m_R \subset \Lambda/\mathfrak m_\Lambda$
est séparable. Le résultat suivant est donc un cas particulier
de la désingularisation générale de Néron donnée dans le
théorème \ref{smoothing-theorem-popescu}.

\begin{lemma}
\label{smoothing-lemma-neron-colimit}
\begin{slogan}
Les extensions non ramifiées d'anneaux de valuation discrète sont ind-lisses : désingularisation de Néron
\end{slogan}
Soit $R \subset \Lambda$ une extension d'anneaux de valuation
discrète d'indice de ramification $1$ qui induit une extension séparable
des corps résiduels et des corps de fractions.
Alors $\Lambda$ est une colimite filtrante de $R$-algèbres lisses.
\end{lemma}

\begin{proof}
D'après Algèbre, lemme \ref{algebra-lemma-when-colimit}, il suffit de montrer
que tout $R \to A \to \Lambda$ comme dans la situation \ref{smoothing-situation-neron}
se factorise sous la forme $A \to B \to \Lambda$, où $B$ est une
$R$-algèbre lisse. Après avoir remplacé $A$ par son image dans $\Lambda$,
nous pouvons supposer que $A$ est un domaine dont le corps des fractions $K$
est un sous-corps du corps des fractions de $\Lambda$.
En particulier, $K$ est séparable sur le corps des fractions de $R$
d'après nos hypothèses. Alors $R \to A$ est lisse en $\mathfrak q = (0)$ d'après
Algèbre, lemme \ref{algebra-lemma-smooth-at-generic-point}.
Après un nombre fini d'éclatements de Néron, nous pouvons supposer que $R \to A$
est lisse en $\mathfrak p$, voir le lemme \ref{smoothing-lemma-neron-desingularization}.
Enfin, après avoir remplacé $A$ par son localisé
en un élément $a \in A$, $a \not \in \mathfrak p$, elle devient
lisse sur $R$, ce qui démontre le lemme.
\end{proof}












\section{Le problème de relèvement}
\label{smoothing-section-lifting}

\noindent
Le but de cette section est de démontrer (proposition \ref{smoothing-proposition-lift})
que la classe des algèbres qui sont des colimites filtrantes d'algèbres lisses
est stable par déformations infinitésimales plates. La démonstration est élémentaire
et n'utilise que les résultats sur les présentations d'algèbres lisses de la
section \ref{smoothing-section-presentations}.

\begin{lemma}
\label{smoothing-lemma-lift-once}
Soit $R \to \Lambda$ un homomorphisme d'anneaux. Soit $I \subset R$ un idéal.
Supposons que
\begin{enumerate}
\item $I^2 = 0$, et
\item $\Lambda/I\Lambda$ soit une colimite filtrante de $R/I$-algèbres lisses.
\end{enumerate}
Soit $\varphi : A \to \Lambda$ un morphisme de $R$-algèbres, où $A$ est de
présentation finie sur $R$. Il existe alors une factorisation
$$
A \to B/J \to \Lambda
$$
où $B$ est une $R$-algèbre lisse et $J \subset IB$ un idéal de type fini.
\end{lemma}

\begin{proof}
Choisissons une factorisation
$$
A/IA \to \bar B \to \Lambda/I\Lambda
$$
où $\bar B$ est standard lisse sur $R/I$ ; cela est possible d'après
l'hypothèse et le lemme \ref{smoothing-lemma-colimit-standard-smooth}. Écrivons
$$
\bar B = A/IA[t_1, \ldots, t_r]/(\bar g_1, \ldots, \bar g_s)
$$
et supposons que $\bar B \to \Lambda/I\Lambda$ envoie $t_i$ sur la classe
de $\lambda_i$ modulo $I\Lambda$. Choisissons
$g_1, \ldots, g_s \in A[t_1, \ldots, t_r]$ relevant
$\bar g_1, \ldots, \bar g_s$. Écrivons
$\varphi(g_i)(\lambda_1, \ldots, \lambda_r) =
\sum \epsilon_{ij} \mu_{ij}$
pour certains $\epsilon_{ij} \in I$ et $\mu_{ij} \in \Lambda$. Définissons
$$
A' = A[t_1, \ldots, t_r, \delta_{i, j}]/
(g_i - \sum \epsilon_{ij} \delta_{ij})
$$
et considérons le morphisme
$$
A' \longrightarrow \Lambda,\quad
a \longmapsto \varphi(a),\quad
t_i \longmapsto \lambda_i,\quad
\delta_{ij} \longmapsto \mu_{ij}
$$
Nous avons
$$
A'/IA' = A/IA[t_1, \ldots, t_r]/(\bar g_1, \ldots, \bar g_s)[\delta_{ij}]
\cong \bar B[\delta_{ij}]
$$
C'est une algèbre standard lisse sur $R/I$, puisque $\bar B$ est standard
lisse. Choisissons une présentation
$A'/IA' = R/I[x_1, \ldots, x_n]/(\bar f_1, \ldots, \bar f_c)$ telle que
$\det(\partial \bar f_j/\partial x_i)_{i, j = 1, \ldots, c}$ soit inversible dans
$A'/IA'$. Choisissons des relèvements $f_1, \ldots, f_c \in R[x_1, \ldots, x_n]$ de
$\bar f_1, \ldots, \bar f_c$. Alors
$$
B = R[x_1, \ldots, x_n, x_{n + 1}]/
(f_1, \ldots, f_c,
x_{n + 1}\det(\partial f_j/\partial x_i)_{i, j = 1, \ldots, c} - 1)
$$
est lisse sur $R$. Comme les morphismes d'anneaux lisses sont formellement lisses
(Algèbre, proposition \ref{algebra-proposition-smooth-formally-smooth}),
il existe un morphisme de $R$-algèbres $B \to A'$ qui est un isomorphisme
modulo $I$. Le morphisme $B \to A'$ est alors surjectif par le lemme de Nakayama
(Algèbre, lemme \ref{algebra-lemma-NAK}).
Ainsi $A' = B/J$, où $J \subset IB$ est de type fini (voir
Algèbre, lemme \ref{algebra-lemma-finite-presentation-independent}).
\end{proof}

\begin{lemma}
\label{smoothing-lemma-lift-twice}
Soit $R \to \Lambda$ un homomorphisme d'anneaux. Soit $I \subset R$ un idéal.
Supposons que
\begin{enumerate}
\item $I^2 = 0$,
\item $\Lambda/I\Lambda$ soit une colimite filtrante de $R/I$-algèbres lisses, et
\item $R \to \Lambda$ soit plat.
\end{enumerate}
Soit $\varphi : B \to \Lambda$ un morphisme de $R$-algèbres, où $B$
est lisse sur $R$. Soit $J \subset IB$ un idéal de type fini
tel que $\varphi(J) = 0$.
Il existe alors des morphismes de $R$-algèbres
$$
B \xrightarrow{\alpha} B' \xrightarrow{\beta} \Lambda
$$
tels que $B'$ soit lisse sur $R$, que $\alpha(J) = 0$ et
que $\beta \circ \alpha = \varphi$.
\end{lemma}

\begin{proof}
Si nous savons démontrer le lemme dans le cas $J = (h)$, nous pouvons le
démontrer par récurrence sur le nombre de générateurs de $J$. Supposons en effet
que $J$ puisse être engendré par $n$ éléments $h_1, \ldots, h_n$ et que le
lemme soit vrai dans tous les cas où $J$ est engendré par $n - 1$ éléments.
Nous appliquons d'abord le cas $n = 1$ pour construire $B \to B' \to \Lambda$
tel que le premier morphisme annule $h_n$. Soit ensuite $J'$ l'idéal de $B'$
engendré par les images de $h_1, \ldots, h_{n - 1}$ ; nous appliquons le cas
$n - 1$ pour construire $B' \to B'' \to \Lambda$.
Il est immédiat de vérifier que $B \to B'' \to \Lambda$ convient.

\medskip\noindent
Supposons $J = (h)$ et écrivons $h = \sum \epsilon_i b_i$
pour certains $\epsilon_i \in I$ et $b_i \in B$. Remarquons que
$0 = \varphi(h) = \sum \epsilon_i \varphi(b_i)$.
Comme $\Lambda$ est plat sur $R$, le critère équationnel de platitude
(Algèbre, lemme \ref{algebra-lemma-flat-eq}) implique qu'il existe
$\lambda_j \in \Lambda$, $j = 1, \ldots, m$, et $a_{ij} \in R$ tels que
$\varphi(b_i) = \sum_j a_{ij} \lambda_j$ et $\sum_i \epsilon_i a_{ij} = 0$.
Posons
$$
C = B[x_1, \ldots, x_m]/(b_i - \sum a_{ij} x_j)
$$
où $C \to \Lambda$ est défini par $\varphi$ et $x_j \mapsto \lambda_j$.
Choisissons une factorisation
$$
C \to B'/J' \to \Lambda
$$
comme dans le lemme \ref{smoothing-lemma-lift-once}. Comme $B$ est lisse sur $R$, nous pouvons
relever le morphisme $B \to C \to B'/J'$ en un morphisme $\alpha : B \to B'$. Alors
$\varphi = \beta \circ \alpha$. Pour achever la démonstration, vérifions que
$\alpha(h) = 0$. En effet, le fait que $\alpha$ relève
$B \to C \to B'/J'$ implique que
$$
\alpha(b_i) = \sum a_{ij} \xi_j + \theta_i
$$
pour certains $\xi_j \in B'$ et $\theta_i \in J' \subset IB'$.
Par conséquent,
$$
\alpha(h) = \alpha(\sum \epsilon_i b_i) =
\sum \epsilon_i a_{ij} \xi_j + \sum \epsilon_i \theta_i = 0
$$
en vertu des relations ci-dessus et du fait que $I^2 = 0$.
\end{proof}

\begin{proposition}
\label{smoothing-proposition-lift}
\begin{slogan}
L'ind-lissité d'une algèbre est stable par déformations infinitésimales
\end{slogan}
Soit $R \to \Lambda$ un homomorphisme d'anneaux. Soit $I \subset R$ un idéal.
Supposons que
\begin{enumerate}
\item $I$ soit nilpotent,
\item $\Lambda/I\Lambda$ soit une colimite filtrante de $R/I$-algèbres lisses, et
\item $R \to \Lambda$ soit plat.
\end{enumerate}
Alors $\Lambda$ est une colimite filtrante de $R$-algèbres lisses.
\end{proposition}

\begin{proof}
Comme $I^n = 0$ pour un certain $n$, une récurrence sur $n$ montre qu'il
suffit de considérer le cas où $I^2 = 0$. Soit
$\varphi : A \to \Lambda$ un morphisme de $R$-algèbres, où $A$ est de
présentation finie sur $R$. Nous devons trouver une factorisation $A \to B \to \Lambda$
avec $B$ lisse sur $R$, voir Algèbre, lemme \ref{algebra-lemma-when-colimit}.
D'après le lemme \ref{smoothing-lemma-lift-once}, nous pouvons supposer que
$A = B/J$, où $B$ est lisse sur $R$ et $J \subset IB$
est un idéal de type fini. D'après le
lemme \ref{smoothing-lemma-lift-twice},
il existe un diagramme commutatif
$$
\xymatrix{
B \ar[rr]_\alpha \ar[rd]_\varphi & & B' \ar[ld]^\beta \\
& \Lambda
}
$$
de $R$-algèbres, où $B'$ est lisse sur $R$ et $\alpha(J) = 0$.
Ainsi, $\alpha$ se factorise sous la forme $B \to A \to B'$, ce qui achève la démonstration.
\end{proof}






\section{Le lemme de relèvement}
\label{smoothing-section-lifting-lemma}

\noindent
Voici un lemme diaboliquement ingénieux.

\begin{lemma}
\label{smoothing-lemma-lifting}
Soit $R$ un anneau noethérien et soit $\Lambda$ une $R$-algèbre.
Soit $\pi \in R$ et supposons que $\text{Ann}_R(\pi) = \text{Ann}_R(\pi^2)$ et
$\text{Ann}_\Lambda(\pi) = \text{Ann}_\Lambda(\pi^2)$.
Supposons donnés des morphismes de $R$-algèbres
$R/\pi^2R \to \bar C \to \Lambda/\pi^2\Lambda$
avec $\bar C$ de présentation finie.
Il existe alors un homomorphisme de $R$-algèbres
$D \to \Lambda$ et un diagramme commutatif
$$
\xymatrix{
R/\pi^2R \ar[r] \ar[d] &
\bar C \ar[r] \ar[d] &
\Lambda/\pi^2\Lambda \ar[d] \\
R/\pi R \ar[r] &
D/\pi D \ar[r] &
\Lambda/\pi \Lambda
}
$$
ayant les propriétés suivantes :
\begin{enumerate}
\item[(a)] $D$ est de présentation finie,
\item[(b)] $R \to D$ est lisse en tout idéal premier $\mathfrak q$ tel que
$\pi \not \in \mathfrak q$,
\item[(c)] $R \to D$ est lisse en tout idéal premier $\mathfrak q$ tel que
$\pi \in \mathfrak q$ et situé au-dessus d'un idéal premier de $\bar C$ où
$R/\pi^2 R \to \bar C$ est lisse, et
\item[(d)] $\bar C/\pi \bar C \to D/\pi D$ est lisse en tout idéal premier
situé au-dessus d'un idéal premier de $\bar C$ où $R/\pi^2R \to \bar C$ est lisse.
\end{enumerate}
\end{lemma}

\begin{proof}
Choisissons une présentation
$$
\bar C = R[x_1, \ldots, x_n]/(f_1, \ldots, f_m)
$$
Notons également $I = (f_1, \ldots, f_m)$ et $\bar I$ l'image de
$I$ dans $R/\pi^2R[x_1, \ldots, x_n]$. Puisque $R$ est noethérien,
$\bar C$ l'est aussi. Le lieu lisse de $R/\pi^2 R \to \bar C$
est donc quasi-compact, voir
Topologie, lemme \ref{topology-lemma-Noetherian}.
En appliquant le
lemme \ref{smoothing-lemma-find-strictly-standard},
nous pouvons choisir une famille finie d'éléments
$a_1, \ldots, a_r \in R[x_1, \ldots, x_n]$ telle que
\begin{enumerate}
\item la réunion des sous-espaces ouverts
$\Spec(\bar C_{a_k}) \subset \Spec(\bar C)$
recouvre le lieu lisse de $R/\pi^2 R \to \bar C$, et
\item pour chaque $k = 1, \ldots, r$, il existe une partie finie
$E_k \subset \{1, \ldots, m\}$ telle que
$(\bar I/\bar I^2)_{a_k}$ soit librement engendré par les classes des
$f_j$, $j \in E_k$.
\end{enumerate}
Posons $I_k = (f_j, j \in E_k) \subset I$ et notons $\bar I_k$
l'image de $I_k$ dans $R/\pi^2R[x_1, \ldots, x_n]$.
D'après (2) et le lemme de Nakayama, $(\bar I/\bar I_k)_{a_k}$
est annulé par $1 + b'_k$ pour un certain $b'_k \in \bar I_{a_k}$.
Supposons que $b'_k$ soit l'image de $b_k/(a_k)^N$ pour un certain $b_k \in I$
et un certain entier $N$. En remplaçant $a_k$ par $(a_k)^N + b_k$, nous obtenons
\begin{enumerate}
\item[(3)] $(\bar I_k)_{a_k} = (\bar I)_{a_k}$.
\end{enumerate}
Ainsi, quitte à remplacer $a_k$ par une puissance suffisamment grande, nous pouvons écrire
\begin{enumerate}
\item[(4)]
$a_k f_\ell = \sum\nolimits_{j \in E_k} h_{k, \ell}^jf_j + \pi^2 g_{k, \ell}$
\end{enumerate}
pour tout $\ell \in \{1, \ldots, m\}$ et certains
$h_{k, \ell}^j, g_{k, \ell} \in R[x_1, \ldots, x_n]$.
Si $\ell \in E_k$, nous choisissons $h_{k, \ell}^j = a_k\delta_{\ell, j}$
(symbole de Kronecker) et $g_{k, \ell} = 0$. Posons
$$
D = R[x_1, \ldots, x_n, z_1, \ldots, z_m]/
(f_j - \pi z_j, p_{k, \ell}).
$$
Ici, $j \in \{1, \ldots, m\}$, $k \in \{1, \ldots, r\}$,
$\ell \in \{1, \ldots, m\}$, et
$$
p_{k, \ell} = a_k z_\ell - \sum\nolimits_{j \in E_k} h_{k, \ell}^j z_j
- \pi g_{k, \ell}.
$$
Remarquons que, pour $\ell \in E_k$, nos choix ci-dessus donnent $p_{k, \ell} = 0$.

\medskip\noindent
Le morphisme $R \to D$ est le morphisme structural.
Supposons que $\bar C \to \Lambda/\pi^2\Lambda$ envoie $x_i$
sur la classe de $\lambda_i$ modulo $\pi^2$. Pour un élément
$f \in R[x_1, \ldots, x_n]$, nous notons $f(\lambda) \in \Lambda$
le résultat de la substitution de $\lambda_i$ à $x_i$. Nous savons alors que
$f_j(\lambda) = \pi^2 \mu_j$ pour un certain $\mu_j \in \Lambda$.
Définissons $D \to \Lambda$ par les règles $x_i \mapsto \lambda_i$ et
$z_j \mapsto \pi\mu_j$. Ce morphisme est bien défini car
\begin{align*}
p_{k, \ell} & \mapsto
a_k(\lambda) \pi \mu_\ell -
\sum\nolimits_{j \in E_k} h_{k, \ell}^j(\lambda) \pi \mu_j
- \pi g_{k, \ell}(\lambda) \\
& =
\pi\left(a_k(\lambda) \mu_\ell -
\sum\nolimits_{j \in E_k} h_{k, \ell}^j(\lambda) \mu_j
- g_{k, \ell}(\lambda)\right)
\end{align*}
En substituant $x_i = \lambda_i$ dans (4), nous voyons que l'expression
entre parenthèses est annulée par $\pi^2$, donc par $\pi$, puisque nous avons supposé
$\text{Ann}_\Lambda(\pi) = \text{Ann}_\Lambda(\pi^2)$.
Le morphisme $\bar C \to D/\pi D$ est déterminé par $x_i \mapsto x_i$
(il est manifestement bien défini). Il reste donc à démontrer (b), (c) et (d).

\medskip\noindent
À l'aide de (4), nous obtenons l'égalité clé suivante :
\begin{align*}
\pi p_{k, \ell} & =
\pi a_k z_\ell - \sum\nolimits_{j \in E_k} \pi h_{k, \ell}^jz_j
- \pi^2 g_{k, \ell} \\
& =
- a_k (f_\ell - \pi z_\ell) + a_k f_\ell +
\sum\nolimits_{j \in E_k} h_{k, \ell}^j (f_j - \pi z_j) -
\sum\nolimits_{j \in E_k} h_{k, \ell}^j f_j - \pi^2 g_{k, \ell} \\
& =
-a_k(f_\ell - \pi z_\ell) +
\sum\nolimits_{j \in E_k} h_{k, \ell}^j(f_j - \pi z_j)
\end{align*}
Le résultat est un élément de l'idéal engendré par les $f_j - \pi z_j$.
En particulier, nous voyons que $D[1/\pi]$ est isomorphe à
$R[1/\pi][x_1, \ldots, x_n, z_1, \ldots, z_m]/(f_j - \pi z_j)$
qui est isomorphe à $R[1/\pi][x_1, \ldots, x_n]$, donc lisse
sur $R$. Cela démontre (b).

\medskip\noindent
Pour $k \in \{1, \ldots, r\}$ fixé, considérons l'anneau
$$
D_k = R[x_1, \ldots, x_n, z_1, \ldots, z_m]/
(f_j - \pi z_j, j \in E_k, p_{k, \ell})
$$
Le nombre d'équations est $m = |E_k| + (m - |E_k|)$, puisque $p_{k, \ell}$
est nul si $\ell \in E_k$. Remarquons aussi que
\begin{align*}
(D_k/\pi D_k)_{a_k}
& =
R/\pi R[x_1, \ldots, x_n, 1/a_k, z_1, \ldots, z_m]/
(f_j, j \in E_k, p_{k, \ell}) \\
& =
(\bar C/\pi \bar C)_{a_k}[z_1, \ldots, z_m]/
(a_kz_\ell - \sum\nolimits_{j \in E_k} h_{k, \ell}^j z_j) \\
& \cong
(\bar C/\pi \bar C)_{a_k}[z_j, j \in E_k]
\end{align*}
En particulier, $(D_k/\pi D_k)_{a_k}$ est lisse sur $(\bar C/\pi \bar C)_{a_k}$.
Par notre choix de $a_k$, l'anneau $(\bar C/\pi \bar C)_{a_k}$ est lisse
sur $R/\pi R$ de dimension relative $n - |E_k|$, voir (2). Par conséquent,
pour un idéal premier $\mathfrak q_k \subset D_k$ contenant $\pi$ et situé au-dessus de
$\Spec(\bar C_{a_k})$, l'anneau de la fibre de $R \to D_k$
est lisse en $\mathfrak q_k$ de dimension $n$. Ainsi, $R \to D_k$ est syntomique
en $\mathfrak q_k$ d'après le décompte ci-dessus du nombre d'équations, voir
Algèbre, lemme \ref{algebra-lemma-localize-relative-complete-intersection}.
Il s'ensuit que $R \to D_k$ est lisse en $\mathfrak q_k$, voir
Algèbre, lemme \ref{algebra-lemma-flat-fibre-smooth}.

\medskip\noindent
Pour achever la démonstration, soit $\mathfrak q \subset D$ un idéal premier
contenant $\pi$ et situé au-dessus d'un idéal premier où $R/\pi^2 R \to \bar C$
est lisse. D'après (1), $a_k \not \in \mathfrak q$ pour un certain $k$.
Nous allons montrer que la surjection $D_k \to D$ induit
un isomorphisme sur les anneaux locaux en $\mathfrak q$. Nous savons, par le
paragraphe précédent, que les morphismes d'anneaux
$\bar C/\pi \bar C \to D_k/\pi D_k$ et $R \to D_k$ sont lisses en
l'idéal premier correspondant $\mathfrak q_k$ ; cela démontrera (c) et (d),
et achèvera donc la démonstration.

\medskip\noindent
Tout d'abord, pour tout $\ell$, l'équation démontrée ci-dessus
$\pi p_{k, \ell} = -a_k(f_\ell - \pi z_\ell) +
\sum_{j \in E_k} h_{k, \ell}^j (f_j - \pi z_j)$ montre que
$f_\ell - \pi z_\ell$ a une image nulle dans $(D_k)_{a_k}$, et donc dans
$(D_k)_{\mathfrak q_k}$.
Les relations (4) impliquent que $a_k f_\ell =
\sum_{j \in E_k} h_{k, \ell}^j f_j$ dans $I/I^2$.
Comme $(\bar I_k/\bar I_k^2)_{a_k}$ est libre sur les $f_j$, $j \in E_k$,
nous voyons que
$$
a_{k'} h_{k, \ell}^j -
\sum\nolimits_{j' \in E_{k'}} h_{k', \ell}^{j'} h_{k, j'}^j
$$
est nul dans $\bar C_{a_k}$ pour tous $k, k', \ell$ et $j \in E_k$.
Il existe donc un entier $N$ suffisamment grand tel que
$$
a_k^N\left(
a_{k'} h_{k, \ell}^j -
\sum\nolimits_{j' \in E_{k'}} h_{k', \ell}^{j'} h_{k, j'}^j
\right)
$$
appartienne à $I_k + \pi^2R[x_1, \ldots, x_n]$. En calculant modulo $\pi$, nous obtenons
\begin{align*}
&
a_kp_{k', \ell} - a_{k'}p_{k, \ell} + \sum h_{k', \ell}^{j'} p_{k, j'}
\\
&
=
- a_k \sum h_{k', \ell}^{j'} z_{j'}
+ a_{k'} \sum h_{k, \ell}^j z_j
+ \sum h_{k', \ell}^{j'} a_k z_{j'}
- \sum \sum h_{k', \ell}^{j'} h_{k, j'}^j z_j \\
&
=
\sum \left(
a_{k'} h_{k, \ell}^j
- \sum h_{k', \ell}^{j'} h_{k, j'}^j
\right) z_j
\end{align*}
avec la convention de sommation d'Einstein. En combinant ce calcul avec ce qui précède,
nous voyons que $a_k^{N + 1} p_{k', \ell}$ appartient à l'idéal engendré
par $I_k$ et $\pi$ dans $R[x_1, \ldots, x_n, z_1, \ldots, z_m]$.
Ainsi, $p_{k', \ell}$ s'envoie dans $\pi (D_k)_{a_k}$. D'autre part, l'équation
$$
\pi p_{k', \ell} =
-a_{k'} (f_\ell - \pi z_\ell) +
\sum\nolimits_{j' \in E_{k'}} h_{k', \ell}^{j'}(f_{j'} - \pi z_{j'})
$$
montre que $\pi p_{k', \ell}$ est nul dans $(D_k)_{a_k}$.
Comme nous avons supposé que $\text{Ann}_R(\pi) = \text{Ann}_R(\pi^2)$
et que $(D_k)_{\mathfrak q_k}$ est lisse, donc plat, sur $R$,
nous voyons que
$\text{Ann}_{(D_k)_{\mathfrak q_k}}(\pi) =
\text{Ann}_{(D_k)_{\mathfrak q_k}}(\pi^2)$.
Nous en concluons que $p_{k', \ell}$ s'annule également, et donc que
$D_{\mathfrak q} = (D_k)_{\mathfrak q_k}$, ce qui conclut la démonstration.
\end{proof}
\section{Le lemme de désingularisation}
\label{smoothing-section-desingularization-lemma}

\noindent
Voici un autre lemme d'une ingéniosité diabolique.

\begin{lemma}
\label{smoothing-lemma-desingularize}
Soit $R$ un anneau noethérien.
Soit $\Lambda$ une $R$-algèbre. Soit $\pi \in R$ et
supposons que $\text{Ann}_\Lambda(\pi) = \text{Ann}_\Lambda(\pi^2)$. Soit
$A \to \Lambda$ un morphisme de $R$-algèbres, avec $A$ de
présentation finie. Supposons que
\begin{enumerate}
\item l'image de $\pi$ est strictement standard dans $A$ sur $R$, et
\item il existe une section $\rho : A/\pi^4 A \to R/\pi^4 R$
compatible avec le morphisme vers $\Lambda/\pi^4 \Lambda$.
\end{enumerate}
Alors il existe des morphismes de $R$-algèbres $A \to B \to \Lambda$, avec
$B$ de présentation finie, tels que $\mathfrak a B \subset H_{B/R}$, où
$\mathfrak a = \text{Ann}_R(\text{Ann}_R(\pi^2)/\text{Ann}_R(\pi))$.
\end{lemma}

\begin{proof}
Choisissons une présentation
$$
A = R[x_1, \ldots, x_n]/(f_1, \ldots, f_m)
$$
et un entier $0 \leq c \leq \min(n, m)$ tels que
(\ref{smoothing-equation-strictly-standard-one}) soit satisfaite pour $\pi$ et que
\begin{equation}
\label{smoothing-equation-star}
\pi f_{c + j} \in (f_1, \ldots, f_c) + (f_1, \ldots, f_m)^2
\end{equation}
pour $j = 1, \ldots, m - c$. Disons que $\rho$ envoie $x_i$ sur la classe de
$r_i \in R$. Nous pouvons alors remplacer $x_i$ par $x_i - r_i$. Ainsi, nous pouvons
supposer que $\rho(x_i) = 0$ dans $R/\pi^4 R$. Il s'ensuit que
$f_j(0) \in \pi^4R$ et que $A \to \Lambda$ envoie $x_i$
sur $\pi^4\lambda_i$ pour un certain $\lambda_i \in \Lambda$. Écrivons
$$
f_j = f_j(0) + \sum\nolimits_{i = 1, \ldots, n} r_{ji} x_i + \text{termes d'ordre supérieur}
$$
Il en résulte que le terme constant de $\partial f_j/\partial x_i$ est
$r_{ji}$. Appliquons $\rho$ à (\ref{smoothing-equation-strictly-standard-one})
pour $\pi$ ; nous obtenons
$$
\pi = \sum\nolimits_{I \subset \{1, \ldots, n\},\ |I| = c}
r_I \det(r_{ji})_{j = 1, \ldots, c,\ i \in I} \bmod \pi^4R
$$
pour certains $r_I \in R$. Nous avons donc
$$
u\pi = \sum\nolimits_{I \subset \{1, \ldots, n\},\ |I| = c}
r_I \det(r_{ji})_{j = 1, \ldots, c,\ i \in I}
$$
pour un certain $u \in 1 + \pi^3R$. D'après
Algèbre, lemme \ref{algebra-lemma-matrix-left-inverse},
il existe alors une matrice $n \times c$ $(s_{ik})$ telle que
$$
u\pi \delta_{jk} = \sum\nolimits_{i = 1, \ldots, n} r_{ji}s_{ik}\quad
\text{pour tous } j, k = 1, \ldots, c
$$
(symbole de Kronecker). Introduisons des variables auxiliaires
$v_1, \ldots, v_c, w_1, \ldots, w_n$ et posons
$$
h_i = x_i - \pi^2 \sum\nolimits_{j = 1, \ldots c} s_{ij} v_j - \pi^3 w_i
$$
Dans la suite, nous utiliserons que
$$
R[x_1, \ldots, x_n, v_1, \ldots, v_c, w_1, \ldots, w_n]/
(h_1, \ldots, h_n) = R[v_1, \ldots, v_c, w_1, \ldots, w_n]
$$
sans autre mention. Dans
$R[x_1, \ldots, x_n, v_1, \ldots, v_c, w_1, \ldots, w_n]/
(h_1, \ldots, h_n)$, nous avons
\begin{align*}
f_j & = f_j(x_1 - h_1, \ldots, x_n - h_n) \\
& =
\pi^2 \sum\nolimits_{k = 1}^c
\left(\sum\nolimits_{i = 1}^n r_{ji} s_{ik}\right) v_k
+
\pi^3 \sum\nolimits_{i = 1}^n r_{ji}w_i \bmod \pi^4 \\
& =
\pi^3 v_j + \pi^3 \sum\nolimits_{i = 1}^n r_{ji}w_i \bmod \pi^4
\end{align*}
pour $1 \leq j \leq c$. Nous pouvons donc choisir des éléments
$g_j \in R[v_1, \ldots, v_c, w_1, \ldots, w_n]$
tels que $g_j = v_j + \sum r_{ji}w_i \bmod \pi$
et que $f_j = \pi^3 g_j$ dans la $R$-algèbre :
$$
R[x_1, \ldots, x_n, v_1, \ldots, v_c, w_1, \ldots, w_n]/
(h_1, \ldots, h_n)
$$
Posons
$$
B = R[x_1, \ldots, x_n, v_1, \ldots, v_c, w_1, \ldots, w_n]/
(f_1, \ldots, f_m, h_1, \ldots, h_n, g_1, \ldots, g_c).
$$
Le morphisme $A \to B$ est évident. Définissons $B \to \Lambda$ en envoyant
$x_i \to \pi^4\lambda_i$, $v_i \mapsto 0$, et $w_i \mapsto \pi \lambda_i$.
Il est alors clair que les éléments $f_j$ et $h_i$ ont pour image zéro
dans $\Lambda$. De plus, $g_i$ a manifestement pour image un élément
$t$ de $\pi\Lambda$ tel que $\pi^3t = 0$ (car $f_i = \pi^3 g_i$ modulo
l'idéal engendré par les $h$). Notre hypothèse
$\text{Ann}_\Lambda(\pi) = \text{Ann}_\Lambda(\pi^2)$ entraîne donc que $t = 0$.
Il reste ainsi à démontrer l'assertion concernant la lissité.

\medskip\noindent
Remarquons que $B_\pi \cong A_\pi[v_1, \ldots, v_c]$, car les équations
$g_i = 0$ résultent de $f_i = 0$. Ainsi, $B_\pi$ est lisse sur $R$
puisque $A_\pi$ est lisse sur $R$ par l'hypothèse selon laquelle $\pi$ est
strictement standard dans $A$ sur $R$ ; voir le
Lemme \ref{smoothing-lemma-elkik}.

\medskip\noindent
Posons $B' = R[v_1, \ldots, v_c, w_1, \ldots, w_n]/(g_1, \ldots, g_c)$.
Comme $g_i = v_i + \sum r_{ji}w_i \bmod \pi$, nous voyons que
$B'/\pi B' = R/\pi R[w_1, \ldots, w_n]$. Ainsi,
$R \to B'$ est lisse de dimension relative $n$ en tout
point de $V(\pi)$ d'après
Algèbre, lemmes
\ref{algebra-lemma-localize-relative-complete-intersection} et
\ref{algebra-lemma-flat-fibre-smooth}
(le premier lemme montre qu'il est syntomique en ces idéaux premiers, donc en particulier
plat, après quoi le second lemme montre qu'il est lisse).

\medskip\noindent
Soit $\mathfrak q \subset B$ un idéal premier tel que $\pi \in \mathfrak q$ et
tel que, pour un certain $r \in \mathfrak a$, $r \not \in \mathfrak q$.
Notons $\mathfrak q' = B' \cap \mathfrak q$.
Nous affirmons que la surjection $B' \to B$ induit un isomorphisme d'anneaux
locaux $(B')_{\mathfrak q'} \to B_\mathfrak q$. Cela
achèvera la démonstration du lemme. Remarquons que $B_\mathfrak q$ est le
quotient de $(B')_{\mathfrak q'}$ par l'idéal engendré par
$f_{c + j}$, $j = 1, \ldots, m - c$. Nous observons deux faits :
premièrement, l'image de $f_{c + j}$ dans $(B')_{\mathfrak q'}$ est
divisible par $\pi^2$ ;
deuxièmement, l'image de $\pi f_{c + j}$ dans $(B')_{\mathfrak q'}$
peut s'écrire $\sum b_{j_1 j_2} f_{c + j_1}f_{c + j_2}$ d'après
(\ref{smoothing-equation-star}). Nous voyons donc que l'image de chaque $\pi f_{c + j}$
appartient à l'idéal engendré par les éléments $\pi^2 f_{c + j'}$.
Ainsi, $\pi f_{c + j} = 0$ dans $(B')_{\mathfrak q'}$, puisque celui-ci est un
anneau local noethérien ; voir
Algèbre, lemme \ref{algebra-lemma-intersect-powers-ideal-module-zero}.
Comme $R \to (B')_{\mathfrak q'}$ est plat, nous avons
$$
\left(\text{Ann}_R(\pi^2)/\text{Ann}_R(\pi)\right)
\otimes_R (B')_{\mathfrak q'}
=
\text{Ann}_{(B')_{\mathfrak q'}}(\pi^2)/\text{Ann}_{(B')_{\mathfrak q'}}(\pi)
$$
Puisque $r \in \mathfrak a$ est inversible dans
$(B')_{\mathfrak q'}$, nous voyons que ce module est nul.
Il s'ensuit que l'image de $f_{c + j}$ est nulle dans
$(B')_{\mathfrak q'}$, comme souhaité.
\end{proof}

\begin{lemma}
\label{smoothing-lemma-desingularize-strictly-standard}
Soit $R$ un anneau noethérien. Soit $\Lambda$ une $R$-algèbre.
Soit $\pi \in R$ et supposons que $\text{Ann}_R(\pi) = \text{Ann}_R(\pi^2)$ et
$\text{Ann}_\Lambda(\pi) = \text{Ann}_\Lambda(\pi^2)$.
Soient $A \to \Lambda$ et $D \to \Lambda$ des morphismes de $R$-algèbres, avec
$A$ et $D$ de présentation finie. Supposons que
\begin{enumerate}
\item $\pi$ est strictement standard dans $A$ sur $R$, et
\item il existe un morphisme de $R$-algèbres $A/\pi^4 A \to D/\pi^4 D$ compatible
avec les morphismes vers $\Lambda/\pi^4 \Lambda$.
\end{enumerate}
Alors il existe un morphisme de $R$-algèbres $B \to \Lambda$, avec $B$ de
présentation finie, ainsi que des morphismes de $R$-algèbres $A \to B$ et $D \to B$
compatibles avec les morphismes vers $\Lambda$, tels que $H_{D/R}B \subset H_{B/D}$
et $H_{D/R}B \subset H_{B/R}$.
\end{lemma}

\begin{proof}
Appliquons le Lemme \ref{smoothing-lemma-desingularize} à
$$
D \longrightarrow A \otimes_R D \longrightarrow \Lambda
$$
et à l'image de $\pi$ dans $D$. D'après le
Lemme \ref{smoothing-lemma-strictly-standard-base-change},
$\pi$ est strictement standard dans $A \otimes_R D$ sur $D$.
Comme section $\rho : (A \otimes_R D)/\pi^4 (A \otimes_R D) \to D/\pi^4 D$,
prenons le morphisme induit par celui de (2). Ainsi, le
Lemme \ref{smoothing-lemma-desingularize} s'applique et fournit une factorisation
$A \otimes_R D \to B \to \Lambda$, avec $B$ de présentation finie,
et $\mathfrak a B \subset H_{B/D}$, où
$$
\mathfrak a = \text{Ann}_D(\text{Ann}_D(\pi^2)/\text{Ann}_D(\pi)).
$$
Pour tout idéal premier $\mathfrak q$ de $D$ tel que $D_\mathfrak q$ soit plat sur $R$,
nous avons
$\text{Ann}_{D_\mathfrak q}(\pi^2)/\text{Ann}_{D_\mathfrak q}(\pi) = 0$
car la formation des annulateurs d'éléments commute au changement de base plat et
nous avons supposé $\text{Ann}_R(\pi) = \text{Ann}_R(\pi^2)$. Comme $D$ est
noethérien, $\text{Ann}_D(\pi^2)/\text{Ann}_D(\pi)$ est un
$D$-module de type fini ; la formation de son annulateur commute donc à la localisation.
Nous voyons ainsi que $\mathfrak a \not \subset \mathfrak q$. Par conséquent,
$D \to B$ est lisse en tout idéal premier de $B$ situé au-dessus de $\mathfrak q$.
Puisque, en tout idéal premier de $D$ où $R \to D$ est lisse,
$D_\mathfrak q$ est plat sur $R$, nous concluons que $H_{D/R}B \subset H_{B/D}$.
La dernière inclusion $H_{D/R}B \subset H_{B/R}$ résulte du fait que les composées
de morphismes d'anneaux lisses sont lisses
(Algèbre, lemme \ref{algebra-lemma-compose-smooth}).
\end{proof}

\begin{lemma}
\label{smoothing-lemma-desingularize-lifting-apply}
Soit $R$ un anneau noethérien. Soit $\Lambda$ une $R$-algèbre.
Soit $\pi \in R$ et supposons que $\text{Ann}_R(\pi) = \text{Ann}_R(\pi^2)$ et
$\text{Ann}_\Lambda(\pi) = \text{Ann}_\Lambda(\pi^2)$.
Soit $A \to \Lambda$ un morphisme de $R$-algèbres, avec
$A$ de présentation finie, et supposons que $\pi$ est strictement standard
dans $A$ sur $R$. Soit
$$
A/\pi^8A \to \bar C \to \Lambda/\pi^8\Lambda
$$
une factorisation, avec $\bar C$ de présentation finie.
Alors il existe une factorisation $A \to B \to \Lambda$, avec $B$ de
présentation finie, telle que $R_\pi \to B_\pi$ soit lisse et que
$$
H_{\bar C/(R/\pi^8 R)} \cdot \Lambda/\pi^8\Lambda
\subset
\sqrt{H_{B/R} \Lambda} \bmod \pi^8\Lambda.
$$
\end{lemma}

\begin{proof}
Appliquons le Lemme \ref{smoothing-lemma-lifting} pour obtenir $R \to D \to \Lambda$
avec une factorisation
$\bar C/\pi^4\bar C \to D/\pi^4 D \to \Lambda/\pi^4\Lambda$
telle que $R \to D$ soit lisse en tout idéal premier ne contenant pas $\pi$
et en tout idéal premier situé au-dessus d'un idéal premier de $\bar C/\pi^4\bar C$
où $R/\pi^8 R \to \bar C$ est lisse.
D'après le Lemme \ref{smoothing-lemma-desingularize-strictly-standard},
il existe une $R$-algèbre $B$ de présentation finie et des
factorisations $A \to B \to \Lambda$ et $D \to B \to \Lambda$
telles que $H_{D/R}B \subset H_{B/R}$. Nous omettons de vérifier qu'il
s'agit bien d'une solution au problème posé par le lemme.
\end{proof}












\section{Préliminaire : réduction à un corps de base}
\label{smoothing-section-reduction}

\noindent
Dans cette section, nous appliquons les lemmes des sections précédentes
pour montrer qu'il suffit de démontrer le résultat principal lorsque l'anneau
de base est un corps ; voir le Lemme \ref{smoothing-lemma-reduce-to-field}.

\begin{situation}
\label{smoothing-situation-global}
Ici, $R \to \Lambda$ est un morphisme d'anneaux régulier entre anneaux noethériens.
\end{situation}

\noindent
Soit $R \to \Lambda$ comme dans la Situation \ref{smoothing-situation-global}.
Nous disons que {\it PT est vraie pour $R \to \Lambda$} si $\Lambda$ est une
colimite filtrante de $R$-algèbres lisses.

\begin{lemma}
\label{smoothing-lemma-product}
Soient $R_i \to \Lambda_i$, $i = 1, 2$, comme dans la Situation \ref{smoothing-situation-global}.
Si PT est vraie pour $R_i \to \Lambda_i$, $i = 1, 2$, alors PT est vraie pour
$R_1 \times R_2 \to \Lambda_1 \times \Lambda_2$.
\end{lemma}

\begin{proof}
Omis. Indication : un produit de colimites filtrantes est une colimite filtrante.
\end{proof}

\begin{lemma}
\label{smoothing-lemma-delocalize-base}
Soient $R \to A \to \Lambda$ des morphismes d'anneaux, avec $A$ de présentation finie
sur $R$. Soit $S \subset R$ une partie
multiplicative. Soit $S^{-1}A \to B' \to S^{-1}\Lambda$ une factorisation telle que
$B'$ soit lisse sur $S^{-1}R$. Alors il existe une factorisation
$A \to B \to \Lambda$ telle que l'image d'un certain $s \in S$ soit un
élément standard élémentaire (Définition \ref{smoothing-definition-strictly-standard})
dans $B$ sur $R$.
\end{lemma}

\begin{proof}
Appliquons d'abord le Lemme \ref{smoothing-lemma-smooth-standard-smooth} à $S^{-1}R \to B'$.
Nous pouvons donc supposer que $B'$ est standard lisse sur $S^{-1}R$.
Écrivons $A = R[x_1, \ldots, x_n]/(g_1, \ldots, g_t)$ et supposons que
$x_i \mapsto \lambda_i$ dans $\Lambda$. Nous pouvons écrire
$B' = S^{-1}R[x_1, \ldots, x_{n + m}]/(f_1, \ldots, f_c)$
pour un certain $c \geq n$, où
$\det(\partial f_j/\partial x_i)_{i, j = 1, \ldots, c}$
est inversible dans $B'$ et tel que $A \to B'$ soit donné par $x_i \mapsto x_i$ ;
voir le Lemme \ref{smoothing-lemma-standard-smooth-include-generators}.
Après avoir multiplié $x_i$, $i > n$, par un élément de $S$ et modifié en conséquence
les équations $f_j$, nous pouvons supposer que $B' \to S^{-1}\Lambda$ envoie
$x_i$ sur $\lambda_i/1$ pour un certain $\lambda_i \in \Lambda$, lorsque $i > n$.
Choisissons une relation
$$
1 =
a_0 \det(\partial f_j/\partial x_i)_{i, j = 1, \ldots, c}
+
\sum\nolimits_{j = 1, \ldots, c} a_jf_j
$$
pour certains $a_j \in S^{-1}R[x_1, \ldots, x_{n + m}]$. Puisque tout élément de $S$
est inversible dans $B'$, nous pouvons, en chassant les dénominateurs, supposer que
$f_j, a_j \in R[x_1, \ldots, x_{n + m}]$ et que
$$
s_0 = a_0 \det(\partial f_j/\partial x_i)_{i, j = 1, \ldots, c}
+
\sum\nolimits_{j = 1, \ldots, c} a_jf_j
$$
pour un certain $s_0 \in S$. Puisque $g_j$ a pour image zéro dans
$S^{-1}R[x_1, \ldots, x_{n + m}]/(f_1, \ldots, f_c)$
il existe des éléments $s_j \in S$ tels que $s_j g_j = 0$ dans
$R[x_1, \ldots, x_{n + m}]/(f_1, \ldots, f_c)$.
Puisque $f_j$ a pour image zéro dans $S^{-1}\Lambda$, il existe $s'_j \in S$
tel que $s'_j f_j(\lambda_1, \ldots, \lambda_{n + m}) = 0$ dans
$\Lambda$. Considérons l'anneau
$$
B = R[x_1, \ldots, x_{n + m}]/
(s'_1f_1, \ldots, s'_cf_c, g_1, \ldots, g_t)
$$
et la factorisation $A \to B \to \Lambda$, où $B \to \Lambda$ est donné par
$x_i \mapsto \lambda_i$. Nous affirmons que $s = s_0s_1 \ldots s_ts'_1 \ldots s'_c$
est standard élémentaire dans $B$ sur $R$, ce qui achève la démonstration.
En effet, $s_j g_j \in (f_1, \ldots, f_c)$, et donc
$sg_j \in (s'_1f_1, \ldots, s'_cf_c)$. Enfin, nous avons
$$
a_0\det(\partial s'_jf_j/\partial x_i)_{i, j = 1, \ldots, c}
+
\sum\nolimits_{j = 1, \ldots, c}
(s'_1 \ldots \hat{s'_j} \ldots s'_c) a_j s'_jf_j
=
s_0s'_1\ldots s'_c
$$
qui divise $s$, comme souhaité.
\end{proof}

\begin{lemma}
\label{smoothing-lemma-reduce-to-field}
\begin{slogan}
Démontrer l'approximation de Popescu se ramène aux algèbres sur un corps
\end{slogan}
Si PT est vraie pour toute Situation \ref{smoothing-situation-global} où $R$
est un corps, alors PT est vraie en général.
\end{lemma}

\begin{proof}
Supposons que PT soit vraie pour toute Situation \ref{smoothing-situation-global} où $R$ est un corps.
Soit $R \to \Lambda$ une Situation \ref{smoothing-situation-global} arbitraire.
Remarquons que $R/I \to \Lambda/I\Lambda$ est encore un morphisme d'anneaux régulier
entre anneaux noethériens ; voir
Compléments d'algèbre, lemme \ref{more-algebra-lemma-regular-base-change}.
Considérons l'ensemble d'idéaux
$$
\mathcal{I} = \{I \subset R \mid R/I \to \Lambda/I\Lambda
\text{ ne vérifie pas PT}\}
$$
Nous devons montrer que $\mathcal{I}$ est vide. Si cet ensemble est non vide,
il contient un élément maximal puisque $R$ est noethérien.
En remplaçant $R$ par $R/I$ et $\Lambda$ par $\Lambda/I$, nous obtenons une
situation où PT est vraie pour $R/I \to \Lambda/I\Lambda$ pour tout
idéal non nul de $R$. En particulier, en appliquant la
Proposition \ref{smoothing-proposition-lift},
nous voyons que $R$ est un anneau réduit.

\medskip\noindent
Soit $A \to \Lambda$ un morphisme de $R$-algèbres, avec $A$ de
présentation finie. Nous devons trouver une factorisation $A \to B \to \Lambda$
où $B$ est lisse sur $R$ ; voir Algèbre, lemme \ref{algebra-lemma-when-colimit}.

\medskip\noindent
Soit $S \subset R$ l'ensemble des non-diviseurs de zéro et
considérons l'anneau total des fractions $Q = S^{-1}R$ de $R$. Nous savons que
$Q = K_1 \times \ldots \times K_n$ est un produit de corps ; voir
Algèbre, lemmes \ref{algebra-lemma-total-ring-fractions-no-embedded-points} et
\ref{algebra-lemma-Noetherian-irreducible-components}.
D'après le Lemme \ref{smoothing-lemma-product} et notre hypothèse,
PT est vraie pour le morphisme d'anneaux $S^{-1}R \to S^{-1}\Lambda$.
Il existe donc une factorisation $S^{-1}A \to B' \to S^{-1}\Lambda$
où $B'$ est lisse sur $S^{-1}R$.

\medskip\noindent
Appliquons le Lemme \ref{smoothing-lemma-delocalize-base} ;
nous obtenons une factorisation $A \to B \to \Lambda$ telle qu'un
certain $\pi \in S$ soit standard élémentaire dans $B$ sur $R$.
Après avoir remplacé $A$ par $B$, nous pouvons supposer que $\pi$ est
standard élémentaire, donc strictement standard dans $A$. Nous savons que
$R/\pi^8R \to \Lambda/\pi^8\Lambda$ vérifie PT.
Il existe donc une factorisation
$R/\pi^8 R \to A/\pi^8A \to \bar C \to \Lambda/\pi^8\Lambda$
où $R/\pi^8 R \to \bar C$ est lisse. D'après le
Lemme \ref{smoothing-lemma-lifting},
il existe un morphisme de $R$-algèbres $D \to \Lambda$, avec $D$ lisse sur $R$,
et une factorisation
$R/\pi^4 R \to A/\pi^4A \to D/\pi^4D \to \Lambda/\pi^4\Lambda$.
D'après le Lemme \ref{smoothing-lemma-desingularize-strictly-standard},
il existe $A \to B \to \Lambda$ avec $B$ lisse sur $R$,
ce qui achève la démonstration.
\end{proof}










\section{Arguments locaux}
\label{smoothing-section-local}


\begin{situation}
\label{smoothing-situation-local}
On se donne un anneau noethérien $R$, un morphisme de $R$-algèbres $A \to \Lambda$
et un idéal premier $\mathfrak q \subset \Lambda$. Nous supposons que $A$ est de
présentation finie sur $R$. Dans cette situation, nous notons
$\mathfrak h_A = \sqrt{H_{A/R} \Lambda}$.
\end{situation}

\noindent
Soit $R \to A \to \Lambda \supset \mathfrak q$ comme dans la
Situation \ref{smoothing-situation-local}. Nous disons que
{\it $R \to A \to \Lambda \supset \mathfrak q$
peut être résolu} s'il existe une factorisation $A \to B \to \Lambda$
avec $B$ de présentation finie et
$\mathfrak h_A \subset \mathfrak h_B \not \subset \mathfrak q$.
Dans ce cas, nous appelons la factorisation $A \to B \to \Lambda$
une {\it résolution de $R \to A \to \Lambda \supset \mathfrak q$}.

\begin{lemma}
\label{smoothing-lemma-lift-solution}
Soit $R \to A \to \Lambda \supset \mathfrak q$ comme dans la
Situation \ref{smoothing-situation-local}. Soit $r \geq 1$, et supposons que
$\pi_1, \ldots, \pi_r \in R$ ont leurs images dans $\mathfrak q$. Supposons que
\begin{enumerate}
\item pour $i = 1, \ldots, r$, nous avons
$$
\text{Ann}_{R/(\pi_1^8, \ldots, \pi_{i - 1}^8)R}(\pi_i)
=
\text{Ann}_{R/(\pi_1^8, \ldots, \pi_{i - 1}^8)R}(\pi_i^2)
$$
et
$$
\text{Ann}_{\Lambda/(\pi_1^8, \ldots, \pi_{i - 1}^8)\Lambda}(\pi_i)
=
\text{Ann}_{\Lambda/(\pi_1^8, \ldots, \pi_{i - 1}^8)\Lambda}(\pi_i^2)
$$
\item pour $i = 1, \ldots, r$, l'élément $\pi_i$ a pour image un élément
strictement standard dans $A$ sur $R$.
\end{enumerate}
Alors, si
$$
R/(\pi_1^8, \ldots, \pi_r^8)R \to A/(\pi_1^8, \ldots, \pi_r^8)A
\to \Lambda/(\pi_1^8, \ldots, \pi_r^8)\Lambda \supset
\mathfrak q/(\pi_1^8, \ldots, \pi_r^8)\Lambda
$$
peut être résolu, il en va de même de $R \to A \to \Lambda \supset \mathfrak q$.
\end{lemma}

\begin{proof}
Nous allons le démontrer par récurrence sur $r$.

\medskip\noindent
Le cas $r = 1$. L'hypothèse affirme ici qu'il existe une
factorisation $A/\pi_1^8 \to \bar C \to \Lambda/\pi_1^8$
qui résout la situation modulo $\pi_1^8$. Les conditions (1) et
(2) sont les hypothèses nécessaires pour appliquer le
Lemme \ref{smoothing-lemma-desingularize-lifting-apply}.
Nous pouvons donc ``relever'' la résolution $\bar C$
en une résolution de $R \to A \to \Lambda \supset \mathfrak q$.

\medskip\noindent
Le cas $r > 1$. Dans ce cas, appliquons l'hypothèse de récurrence pour $r - 1$
à la situation
$R/\pi_1^8 \to A/\pi_1^8 \to \Lambda/\pi_1^8
\supset \mathfrak q/\pi_1^8\Lambda$.
Remarquons que la propriété (2) est préservée par le
Lemme \ref{smoothing-lemma-strictly-standard-base-change}.
\end{proof}

\begin{lemma}
\label{smoothing-lemma-delocalize-weak}
\begin{reference}
\cite[Lemma 12.2]{swan} ou
\cite[Lemma 2]{popescu-GND}
\end{reference}
Soit $R \to A \to \Lambda \supset \mathfrak q$ comme dans la
Situation \ref{smoothing-situation-local}. Soit $\mathfrak p = R \cap \mathfrak q$.
Supposons que $\mathfrak q$ soit minimal au-dessus de $\mathfrak h_A$ et que
$R_\mathfrak p \to A_\mathfrak p \to \Lambda_\mathfrak q
\supset \mathfrak q\Lambda_\mathfrak q$ puisse être résolu.
Alors il existe une factorisation $A \to C \to \Lambda$, avec $C$ de
présentation finie, telle que $H_{C/R} \Lambda \not \subset \mathfrak q$.
\end{lemma}

\begin{proof}
Soit $A_\mathfrak p \to C \to \Lambda_\mathfrak q$ une résolution de
$R_\mathfrak p \to A_\mathfrak p \to \Lambda_\mathfrak q
\supset \mathfrak q\Lambda_\mathfrak q$. Comme $\mathfrak q$ est, par hypothèse,
minimal au-dessus de $\mathfrak h_A$, cela
signifie que $H_{C/R_\mathfrak p} \Lambda_\mathfrak q = \Lambda_\mathfrak q$.
D'après le Lemme \ref{smoothing-lemma-final-solve},
nous pouvons supposer que $C$ est lisse sur $R_\mathfrak p$.
D'après le Lemme \ref{smoothing-lemma-smooth-standard-smooth}, nous pouvons supposer que
$C$ est standard lisse sur $R_\mathfrak p$.
Écrivons $A = R[x_1, \ldots, x_n]/(g_1, \ldots, g_t)$ et supposons que
$A \to \Lambda$ est donné par $x_i \mapsto \lambda_i$.
Écrivons $C = R_\mathfrak p[x_1, \ldots, x_{n + m}]/(f_1, \ldots, f_c)$
pour un certain $c \geq n$, de sorte que $A \to C$ envoie $x_i$ sur $x_i$ et que
$\det(\partial f_j/\partial x_i)_{i, j = 1, \ldots, c}$
soit inversible dans $C$ ; voir le
Lemme \ref{smoothing-lemma-standard-smooth-include-generators}.
Après avoir chassé les dénominateurs, nous pouvons supposer que
$f_1, \ldots, f_c$ sont des éléments de $R[x_1, \ldots, x_{n + m}]$.
Bien entendu,
$\det(\partial f_j/\partial x_i)_{i, j = 1, \ldots, c}$
n'est pas inversible dans $R[x_1, \ldots, x_{n + m}]/(f_1, \ldots, f_c)$,
mais il le devient après inversion d'un certain élément $s_0 \in R$,
$s_0 \not \in \mathfrak p$.
Comme $g_j$ a pour image zéro par $R[x_1, \ldots, x_n] \to A \to C$,
il existe $s_j \in R$, $s_j \not \in \mathfrak p$, tel que
$s_j g_j$ soit nul dans $R[x_1, \ldots, x_{n + m}]/(f_1, \ldots, f_c)$.
Écrivons $f_j = F_j(x_1, \ldots, x_{n + m}, 1)$
pour un polynôme
$F_j \in R[x_1, \ldots, x_n, X_{n + 1}, \ldots, X_{n + m + 1}]$
homogène en $X_{n + 1}, \ldots, X_{n + m + 1}$.
Choisissons $\lambda_{n + i} \in \Lambda$, $i = 1, \ldots, m + 1$, avec
$\lambda_{n + m + 1} \not \in \mathfrak q$, de sorte que $x_{n + i}$ ait pour image
$\lambda_{n + i}/\lambda_{n + m + 1}$ dans $\Lambda_\mathfrak q$.
Alors
\begin{align*}
F_j(\lambda_1, \ldots, \lambda_{n + m + 1})
& =
(\lambda_{n + m + 1})^{\deg(F_j)} F_j(\lambda_1, \ldots, \lambda_n,
\frac{\lambda_{n + 1}}{\lambda_{n + m + 1}}, \ldots,
\frac{\lambda_{n + m}}{\lambda_{n + m + 1}}, 1) \\
& =
(\lambda_{n + m + 1})^{\deg(F_j)} f_j(\lambda_1, \ldots, \lambda_n,
\frac{\lambda_{n + 1}}{\lambda_{n + m + 1}}, \ldots,
\frac{\lambda_{n + m}}{\lambda_{n + m + 1}}) \\
& = 0
\end{align*}
dans $\Lambda_\mathfrak q$. Il existe donc
$\lambda_0 \in \Lambda$, $\lambda_0 \not \in \mathfrak q$, tel que
$\lambda_0 F_j(\lambda_1, \ldots, \lambda_{n + m + 1}) = 0$
dans $\Lambda$. Posons maintenant $B$ égal à
$$
R[x_0, \ldots, x_{n + m + 1}]/
(g_1, \ldots, g_t, x_0F_1(x_1, \ldots, x_{n + m + 1}), \ldots,
x_0F_c(x_1, \ldots, x_{n + m + 1}))
$$
que nous envoyons dans $\Lambda$ en envoyant $x_i$ sur $\lambda_i$.
Soit $b$ l'image de $x_0 x_{n + m + 1} s_0 s_1 \ldots s_t$ dans $B$.
Alors $B_b$ est isomorphe à
$$
R_{s_0s_1 \ldots s_t}[x_0, x_1, \ldots, x_{n + m + 1}, 1/x_0x_{n + m + 1}]/
(f_1, \ldots, f_c)
$$
qui est lisse sur $R$ par construction.
Puisque l'image de $b$ n'appartient pas à $\mathfrak q$, le résultat s'ensuit.
\end{proof}

\begin{lemma}
\label{smoothing-lemma-delocalize-height-zero}
Soit $R \to A \to \Lambda \supset \mathfrak q$ comme dans la
Situation \ref{smoothing-situation-local}. Soit $\mathfrak p = R \cap \mathfrak q$.
Supposons que
\begin{enumerate}
\item $\mathfrak q$ est minimal au-dessus de $\mathfrak h_A$,
\item $R_\mathfrak p \to A_\mathfrak p \to \Lambda_\mathfrak q
\supset \mathfrak q\Lambda_\mathfrak q$ peut être résolu, et
\item $\dim(\Lambda_\mathfrak q) = 0$.
\end{enumerate}
Alors $R \to A \to \Lambda \supset \mathfrak q$ peut être résolu.
\end{lemma}

\begin{proof}
Par (3), l'anneau $\Lambda_\mathfrak q$ est local artinien ; ainsi,
$\mathfrak q\Lambda_\mathfrak q$ est nilpotent. Par conséquent,
$(\mathfrak h_A)^N \Lambda_\mathfrak q = 0$ pour un certain $N > 0$.
Il existe donc $\lambda \in \Lambda$, $\lambda \not \in \mathfrak q$,
tel que $\lambda (\mathfrak h_A)^N = 0$ dans $\Lambda$.
Écrivons $H_{A/R} = (a_1, \ldots, a_r)$, de sorte que $\lambda a_i^N = 0$
dans $\Lambda$. D'après le Lemme \ref{smoothing-lemma-delocalize-weak}, il existe une factorisation
$A \to C \to \Lambda$, avec $C$ de présentation finie, telle que
$\mathfrak h_C \not \subset \mathfrak q$.
Écrivons $C = A[x_1, \ldots, x_n]/(f_1, \ldots, f_m)$.
Posons
$$
B = A[x_1, \ldots, x_n, y_1, \ldots, y_r, z, t_{ij}]/
(f_j - \sum y_i t_{ij}, zy_i)
$$
où les $t_{ij}$ forment un ensemble de $rm$ variables.
Remarquons qu'il existe un morphisme $B \to C[y_i, z]/(y_iz)$ obtenu en posant les $t_{ij}$
égaux à zéro. Le morphisme $B \to \Lambda$ est la composée
$B \to C[y_i, z]/(y_iz) \to \Lambda$, où $C[y_i, z]/(y_iz) \to \Lambda$
est le morphisme donné $C \to \Lambda$, envoie $z$ sur $\lambda$ et envoie
$y_i$ sur l'image de $a_i^N$ dans $\Lambda$.

\medskip\noindent
Nous affirmons que $B$ est une solution pour $R \to A \to \Lambda \supset \mathfrak q$.
Remarquons d'abord que $B_z$ est isomorphe à $C[z, z^{-1}, t_{ij}]$.
Choisissons $c \in H_{C/R}$ dont l'image dans $\Lambda$ n'appartient pas à
$\mathfrak q$. Alors $B_{zc}$ est lisse sur $R$. D'autre part,
$B_{y_\ell} \cong A[x_i, y_i, y_\ell^{-1}, t_{ij}, i \not = \ell]$
qui est lisse sur $A$. Nous voyons ainsi que $zc$ et les $a_\ell y_\ell$
(les composées de morphismes lisses sont lisses) sont tous
des éléments de $H_{B/R}$. Cela démontre le lemme.
\end{proof}




\section{Corps résiduels séparables}
\label{smoothing-section-separable}

\noindent
Dans cette section, nous expliquons comment résoudre un problème local dans le cas
d'une extension séparable de corps résiduels.

\begin{lemma}[Ogoma]
\label{smoothing-lemma-ogoma}
Soit $A$ un anneau noethérien et soit $M$ un $A$-module de type fini.
Soit $S \subset A$ une partie multiplicative. Si $\pi \in A$ et
$\Ker(\pi : S^{-1}M \to S^{-1}M) =
\Ker(\pi^2 : S^{-1}M \to S^{-1}M)$
alors il existe $s \in S$ tel que, pour tout $n > 0$, nous ayons
$\Ker(s^n\pi : M \to M) = \Ker((s^n\pi)^2 : M \to M)$.
\end{lemma}

\begin{proof}
Posons $K = \Ker(\pi : M \to M)$ et
$K' = \{m \in M \mid \pi^2 m = 0\text{ dans }S^{-1}M\}$ et
$Q = K'/K$. Remarquons que $S^{-1}Q = 0$ par hypothèse. Comme $A$
est noethérien, $Q$ est un $A$-module de type fini.
Il existe donc $s \in S$ tel que $s$ annule $Q$.
Alors $s$ convient.
\end{proof}

\begin{lemma}
\label{smoothing-lemma-find-sequence}
Soit $\Lambda$ un anneau noethérien. Soit $I \subset \Lambda$ un idéal.
Soit $I \subset \mathfrak q$, où ce dernier est un idéal premier. Soient $n, e$ des entiers positifs.
Supposons que $\mathfrak q^n\Lambda_\mathfrak q \subset I\Lambda_\mathfrak q$
et que $\Lambda_\mathfrak q$ soit un anneau local régulier de dimension $d$.
Alors il existe des éléments $\pi_1, \ldots, \pi_d \in \Lambda$ tels que
\begin{enumerate}
\item $(\pi_1, \ldots, \pi_d)\Lambda_\mathfrak q =
\mathfrak q\Lambda_\mathfrak q$,
\item $\pi_1^n, \ldots, \pi_d^n \in I$, et
\item pour $i = 1, \ldots, d$, nous avons
$$
\text{Ann}_{\Lambda/(\pi_1^e, \ldots, \pi_{i - 1}^e)\Lambda}(\pi_i) =
\text{Ann}_{\Lambda/(\pi_1^e, \ldots, \pi_{i - 1}^e)\Lambda}(\pi_i^2).
$$
\end{enumerate}
\end{lemma}

\begin{proof}
Posons $S = \Lambda \setminus \mathfrak q$, de sorte que
$\Lambda_\mathfrak q = S^{-1}\Lambda$.
Choisissons d'abord $\pi_1, \ldots, \pi_d$ satisfaisant (1), ce qui est possible
puisque $\Lambda_\mathfrak q$ est régulier. Par hypothèse,
$\pi_i^n \in I\Lambda_\mathfrak q$. Il existe donc
$s_1, \ldots, s_d \in S$ tels que $s_i\pi_i^n \in I$.
En remplaçant $\pi_i$ par $s_i\pi_i$, nous obtenons (2).
Remarquons que (1) et (2) sont préservées par toute multiplication ultérieure par des éléments de $S$.
Supposons que (3) soit satisfaite pour $i = 1, \ldots, t$, pour un certain
$t \in \{0, \ldots, d\}$. Remarquons que
$\pi_1, \ldots, \pi_d$ est une suite régulière dans $S^{-1}\Lambda$ ; voir
Algèbre, lemme \ref{algebra-lemma-regular-ring-CM}.
En particulier, $\pi_1^e, \ldots, \pi_t^e, \pi_{t + 1}$ est une
suite régulière dans $S^{-1}\Lambda = \Lambda_\mathfrak q$ d'après
Algèbre, lemme \ref{algebra-lemma-regular-sequence-powers}.
Nous obtenons donc
$$
\text{Ann}_{S^{-1}\Lambda/(\pi_1^e, \ldots, \pi_{i - 1}^e)}(\pi_i) =
\text{Ann}_{S^{-1}\Lambda/(\pi_1^e, \ldots, \pi_{i - 1}^e)}(\pi_i^2).
$$
Nous obtenons ainsi (3) pour $i = t + 1$ après avoir remplacé $\pi_{t + 1}$ par $s\pi_{t + 1}$
pour un certain $s \in S$, d'après le Lemme \ref{smoothing-lemma-ogoma}. Par récurrence sur $t$, ceci
produit une suite satisfaisant (1), (2) et (3).
\end{proof}

\begin{lemma}
\label{smoothing-lemma-resolve-special}
Soit $k \to A \to \Lambda \supset \mathfrak q$ comme dans la
Situation \ref{smoothing-situation-local}, où
\begin{enumerate}
\item $k$ est un corps,
\item $\Lambda$ est noethérien,
\item $\mathfrak q$ est minimal au-dessus de $\mathfrak h_A$,
\item $\Lambda_\mathfrak q$ est un anneau local régulier, et
\item l'extension de corps $\kappa(\mathfrak q)/k$ est séparable.
\end{enumerate}
Alors $k \to A \to \Lambda \supset \mathfrak q$ peut être résolu.
\end{lemma}

\begin{proof}
Posons $d = \dim \Lambda_\mathfrak q$. Posons $R = k[x_1, \ldots, x_d]$.
Choisissons $n > 0$ tel que
$\mathfrak q^n\Lambda_\mathfrak q \subset \mathfrak h_A\Lambda_\mathfrak q$
ce qui est possible puisque $\mathfrak q$ est minimal au-dessus de $\mathfrak h_A$.
Choisissons des générateurs $a_1, \ldots, a_r$ de $H_{A/k}$. Posons
$$
B = A[x_1, \ldots, x_d, z_{ij}]/(x_i^n - \sum z_{ij}a_j)
$$
Chaque $B_{a_j}$ est lisse sur $R$ : c'est une algèbre de polynômes
sur $A_{a_j}[x_1, \ldots, x_d]$, et $A_{a_j}$ est lisse sur $k$.
Par conséquent, $B_{x_i}$ est lisse sur $R$. Soit $B \to C$ le morphisme de $R$-algèbres
construit dans le Lemme \ref{smoothing-lemma-improve-presentation},
qui est muni d'une rétraction de $R$-algèbres $C \to B$. En particulier,
nous obtenons un morphisme $C \to \Lambda$ s'insérant dans le diagramme ci-dessous.
Par construction, $C_{x_i}$ est une $R$-algèbre lisse dont
$\Omega_{C_{x_i}/R}$ est libre. Il existe donc $c > 0$
tel que $x_i^c$ soit strictement standard dans $C/R$ ; voir le
Lemme \ref{smoothing-lemma-compare-standard}.
Choisissons maintenant $\pi_1, \ldots, \pi_d \in \Lambda$ comme dans le
Lemme \ref{smoothing-lemma-find-sequence},
où $n = n$, $e = 8c$, $\mathfrak q = \mathfrak q$ et $I = \mathfrak h_A$.
Écrivons $\pi_i^n = \sum \lambda_{ij} a_j$ pour certains $\lambda_{ij} \in \Lambda$.
Il existe un morphisme $B \to \Lambda$ donné par $x_i \mapsto \pi_i$
et $z_{ij} \mapsto \lambda_{ij}$. Posons $R = k[x_1, \ldots, x_d]$.
Considérons le diagramme
$$
\xymatrix{
R \ar[r] & B \ar[rd] \\
k \ar[u] \ar[r] & A \ar[u] \ar[r] & \Lambda
}
$$
Appliquons maintenant le
Lemme \ref{smoothing-lemma-lift-solution}
à $R \to C \to \Lambda \supset \mathfrak q$
et à la suite d'éléments $x_1^c, \ldots, x_d^c$ de $R$.
L'hypothèse (2) est claire. L'hypothèse (1) est satisfaite pour $R$
par inspection et pour $\Lambda$ grâce à notre choix de
$\pi_1, \ldots, \pi_d$. (Remarquons que si
$\text{Ann}_\Lambda(\pi) = \text{Ann}_\Lambda(\pi^2)$, alors nous avons
$\text{Ann}_\Lambda(\pi) = \text{Ann}_\Lambda(\pi^c)$ pour tout $c > 0$.)
Il suffit donc de résoudre
$$
R/(x_1^e, \ldots, x_d^e) \to C/(x_1^e, \ldots, x_d^e) \to
\Lambda/(\pi_1^e, \ldots, \pi_d^e) \supset
\mathfrak q/(\pi_1^e, \ldots, \pi_d^e)
$$
pour $e = 8c$. D'après le
Lemme \ref{smoothing-lemma-delocalize-height-zero},
il suffit de résoudre cette situation après localisation en $\mathfrak q$.
Mais puisque $x_1, \ldots, x_d$ ont pour images une suite régulière
dans $\Lambda_\mathfrak q$, nous voyons que $R_\mathfrak p \to \Lambda_\mathfrak q$
est plat ; voir Algèbre, lemme \ref{algebra-lemma-flat-over-regular}. Par conséquent,
$$
R_\mathfrak p/(x_1^e, \ldots, x_d^e) \to
\Lambda_\mathfrak q/(\pi_1^e, \ldots, \pi_d^e)
$$
est un morphisme d'anneaux plat entre anneaux locaux artiniens.
De plus, ce morphisme induit par hypothèse une extension séparable
de corps résiduels. Ce morphisme est donc une colimite filtrante
d'algèbres lisses d'après
Algèbre, lemme \ref{algebra-lemma-colimit-syntomic},
et la Proposition \ref{smoothing-proposition-lift}.
L'existence de la solution souhaitée résulte de
Algèbre, lemme \ref{algebra-lemma-when-colimit}.
\end{proof}
\section{Extensions inséparables de corps résiduels}
\label{smoothing-section-inseparable}

\noindent
Dans cette section, nous expliquons comment résoudre un problème local dans le cas
d'une extension inséparable de corps résiduels.

\begin{lemma}
\label{smoothing-lemma-helper}
Soit $k$ un corps de caractéristique $p > 0$.
Soit $(\Lambda, \mathfrak m, K)$ une $k$-algèbre locale artinienne.
Supposons que $\dim H_1(L_{K/k}) < \infty$.
Alors $\Lambda$ est une colimite filtrante de $k$-algèbres
locales artiniennes $A$ telles que chaque morphisme $A \to \Lambda$ soit plat, que
$\mathfrak m_A \Lambda = \mathfrak m$ et que
$A$ soit essentiellement de type fini sur $k$.
\end{lemma}

\begin{proof}
Remarquons que la platitude de $A \to \Lambda$ implique que $A \to \Lambda$
est injectif ; le lemme affirme donc en réalité que $\Lambda$ est une
réunion filtrante de sous-anneaux $A \subset \Lambda$ de ce type.
Soit $n$ le plus petit entier tel que $\mathfrak m^n = 0$.
Nous démontrons ce lemme par récurrence sur $n$. Le cas $n = 1$ est clair,
car une extension de corps est la réunion de ses sous-extensions de type fini.

\medskip\noindent
Choisissons $\lambda_1, \ldots, \lambda_d \in \mathfrak m$ qui engendrent
$\mathfrak m$. Comme $K$ est formellement lisse sur $\mathbf{F}_p$ (voir
Algèbre, lemme \ref{algebra-lemma-formally-smooth-extensions-easy}), nous pouvons
trouver un morphisme d'anneaux $\sigma : K \to \Lambda$ qui est une section du
morphisme quotient $\Lambda \to K$. En général, $\sigma$ n'est {\bf pas}
un morphisme de $k$-algèbres. Étant donné $\sigma$, définissons
$$
\Psi_\sigma : K[x_1, \ldots, x_d] \longrightarrow \Lambda
$$
en utilisant $\sigma$ sur les éléments de $K$ et en envoyant $x_i$ sur $\lambda_i$.
Assertion : il existe un $\sigma : K \to \Lambda$
et un sous-corps $k \subset F \subset K$ de type fini sur $k$
tels que l'image de $k$ dans $\Lambda$ soit contenue dans
$\Psi_\sigma(F[x_1, \ldots, x_d])$.

\medskip\noindent
Nous démontrons l'assertion par récurrence sur le plus petit entier $n$ tel que
$\mathfrak m^n = 0$. Elle est claire pour $n = 1$. Si $n > 1$, posons
$I = \mathfrak m^{n - 1}$ et $\Lambda' = \Lambda/I$.
Par récurrence, nous pouvons supposer donnés
$\sigma' : K \to \Lambda'$ et $k \subset F' \subset K$, le corps intermédiaire
étant de type fini sur le corps de base, tels que l'image de $k \to \Lambda \to \Lambda'$
soit contenue dans $A' = \Psi_{\sigma'}(F'[x_1, \ldots, x_d])$.
Notons $\tau' : k \to A'$ le morphisme induit.
Choisissons un relèvement $\sigma : K \to \Lambda$ de $\sigma'$ (cela est possible
par la lissité formelle de $K/\mathbf{F}_p$ mentionnée ci-dessus).
Pour usage ultérieur, remarquons que nous pouvons remplacer $\sigma$ par
$\sigma + D$ pour une certaine dérivation $D : K \to I$.
Posons $A = F'[x_1, \ldots, x_d]/(x_1, \ldots, x_d)^n$.
Alors $\Psi_\sigma$ induit un morphisme d'anneaux
$\Psi_\sigma : A \to \Lambda$. La composition avec le
morphisme quotient $\Lambda \to \Lambda'$ induit un morphisme surjectif
$A \to A'$ de noyau nilpotent.
Choisissons un relèvement $\tau : k \to A$ de $\tau'$ (possible puisque $k/\mathbf{F}_p$
est formellement lisse). Nous obtenons ainsi deux morphismes $k \to \Lambda$, à savoir
$\Psi_\sigma \circ \tau : k \to \Lambda$ et le morphisme donné $i : k \to \Lambda$.
Ces morphismes coïncident modulo $I$ ; leur différence est donc une
dérivation $\theta = i - \Psi_\sigma \circ \tau : k \to I$.
Remarquons que si nous remplaçons $\sigma$ par $\sigma + D$, alors nous remplaçons
$\theta$ par $\theta - D|_k$.

\medskip\noindent
Choisissons une famille d'éléments $\{y_j\}_{j \in J}$ de $k$ dont les différentielles
$\text{d}y_j$ forment une base de $\Omega_{k/\mathbf{F}_p}$. La suite de Jacobi-Zariski
pour $\mathbf{F}_p \subset k \subset K$ est
$$
0 \to H_1(L_{K/k}) \to \Omega_{k/\mathbf{F}_p} \otimes K \to
\Omega_{K/\mathbf{F}_p} \to \Omega_{K/k} \to 0
$$
Comme $\dim H_1(L_{K/k}) < \infty$, nous pouvons trouver un sous-ensemble fini $J_0 \subset J$
tel que l'image du premier morphisme soit contenue dans
$\bigoplus_{j \in J_0} K\text{d}y_j$. Par conséquent, les images des éléments
$\text{d}y_j$, $j \in J \setminus J_0$, sont $K$-linéairement indépendantes
dans $\Omega_{K/\mathbf{F}_p}$. Nous pouvons donc choisir
$D : K \to I$ tel que $\theta - D|_k = \xi \circ \text{d}$,
où $\xi$ est la composition
$$
\Omega_{k/\mathbf{F}_p} = \bigoplus\nolimits_{j \in J} k \text{d}y_j
\longrightarrow \bigoplus\nolimits_{j \in J_0} k \text{d}y_j
\longrightarrow I
$$
Posons $f_j = \xi(\text{d}y_j) \in I$ pour $j \in J_0$.
Remplaçons $\sigma$ par $\sigma + D$ comme ci-dessus. Nous voyons alors que
$\theta(a) = \sum_{j \in J_0} a_j f_j$ pour $a \in k$, où
$\text{d}a = \sum a_j \text{d}y_j$ dans $\Omega_{k/\mathbf{F}_p}$.
Remarquons que $I$ est engendré par les monômes
$\lambda^E = \lambda_1^{e_1} \ldots \lambda_d^{e_d}$ de
degré total $|E| = \sum e_i = n - 1$ en $\lambda_1, \ldots, \lambda_d$.
Écrivons $f_j = \sum_E c_{j, E} \lambda^E$ avec $c_{j, E} \in K$.
Remplaçons $F'$ par $F = F'(c_{j, E})$. L'assertion est alors démontrée.

\medskip\noindent
Choisissons $\sigma$ et $F$ comme dans l'assertion. Le noyau de $\Psi_\sigma$ est
engendré par un nombre fini de polynômes
$g_1, \ldots, g_t \in K[x_1, \ldots, x_d]$ et, après avoir agrandi $F$ en lui adjoignant
un nombre fini d'éléments, nous pouvons supposer que leurs coefficients appartiennent à $F$.
Dans ce cas, il est clair que le morphisme
$A = F[x_1, \ldots, x_d]/(g_1, \ldots, g_t) \to
K[x_1, \ldots, x_d]/(g_1, \ldots, g_t) = \Lambda$ est plat.
D'après l'assertion, $A$ est une sous-$k$-algèbre de $\Lambda$.
Il est clair que $\Lambda$ est la colimite filtrante de ces
algèbres, puisque $K$ est la réunion filtrante des sous-corps $F$.
Enfin, ces algèbres sont essentiellement de type fini sur $k$ d'après
Algèbre, lemme
\ref{algebra-lemma-essentially-of-finite-type-into-artinian-local}.
\end{proof}

\begin{lemma}
\label{smoothing-lemma-solution-modulo}
Soit $k$ un corps de caractéristique $p > 0$.
Soit $\Lambda$ une $k$-algèbre noethérienne géométriquement régulière.
Soit $\mathfrak q \subset \Lambda$ un idéal premier.
Soit $n \geq 1$ un entier et soit
$E \subset \Lambda_\mathfrak q/\mathfrak q^n\Lambda_\mathfrak q$
un sous-ensemble fini.
Alors nous pouvons trouver $m \geq 0$ et
$\varphi : k[y_1, \ldots, y_m] \to \Lambda$ ayant les propriétés suivantes :
\begin{enumerate}
\item en posant $\mathfrak p = \varphi^{-1}(\mathfrak q)$, nous avons
$\mathfrak q\Lambda_\mathfrak q = \mathfrak p \Lambda_\mathfrak q$
et $k[y_1, \ldots, y_m]_\mathfrak p \to \Lambda_\mathfrak q$ est plat,
\item il existe une factorisation par des morphismes d'anneaux locaux artiniens
$$
k[y_1, \ldots, y_m]_\mathfrak p/\mathfrak p^n k[y_1, \ldots, y_m]_\mathfrak p
\to D \to
\Lambda_\mathfrak q/\mathfrak q^n\Lambda_\mathfrak q
$$
où la première flèche est essentiellement lisse et la seconde est plate,
\item $E$ est contenu dans $D$ modulo $\mathfrak q^n\Lambda_\mathfrak q$.
\end{enumerate}
\end{lemma}

\begin{proof}
Posons $\bar \Lambda = \Lambda_\mathfrak q/\mathfrak q^n\Lambda_\mathfrak q$.
Remarquons que $\dim H_1(L_{\kappa(\mathfrak q)/k}) < \infty$ d'après
Compléments d'algèbre, proposition
\ref{more-algebra-proposition-characterization-geometrically-regular}.
Choisissons $A \subset \bar \Lambda$ contenant $E$, telle que $A$ soit locale
artinienne, essentiellement de type fini sur $k$, que le morphisme
$A \to \bar \Lambda$ soit plat et que $\mathfrak m_A$ engendre l'idéal maximal
de $\bar \Lambda$ ; voir le Lemme \ref{smoothing-lemma-helper}.
Notons $F = A/\mathfrak m_A$ le corps résiduel, de sorte que $k \subset F \subset K$.
Choisissons $\lambda_1, \ldots, \lambda_t \in \Lambda$ dont les images
appartiennent à $A$ dans $\bar \Lambda$ et tels que, de plus, les images
de $\text{d}\lambda_1, \ldots, \text{d}\lambda_t$ forment une base
de $\Omega_{F/k}$. Considérons le morphisme
$\varphi' : k[y_1, \ldots, y_t] \to \Lambda$ qui envoie $y_j$ sur $\lambda_j$.
Posons $\mathfrak p' = (\varphi')^{-1}(\mathfrak q)$. D'après
Compléments d'algèbre, lemme
\ref{more-algebra-lemma-geometrically-regular-over-field}
le morphisme d'anneaux $k[y_1, \ldots, y_t]_{\mathfrak p'} \to \Lambda_\mathfrak q$
est plat et $\Lambda_\mathfrak q/\mathfrak p' \Lambda_\mathfrak q$ est
régulier. Nous pouvons donc choisir d'autres éléments
$\lambda_{t + 1}, \ldots, \lambda_m \in \Lambda$
dont les images appartiennent à $A \subset \bar \Lambda$ et dont les
images forment un système régulier de paramètres de
$\Lambda_\mathfrak q/\mathfrak p' \Lambda_\mathfrak q$.
Nous obtenons $\varphi : k[y_1, \ldots, y_m] \to \Lambda$ ayant
la propriété (1), tel que
$k[y_1, \ldots, y_m]_\mathfrak p/\mathfrak p^n k[y_1, \ldots, y_m]_\mathfrak p
\to \bar\Lambda$
se factorise par $A$. Ainsi,
$k[y_1, \ldots, y_m]_\mathfrak p/\mathfrak p^n k[y_1, \ldots, y_m]_\mathfrak p
\to A$ est plat d'après
Algèbre, lemme \ref{algebra-lemma-flatness-descends-more-general}.
Par construction, l'extension de corps résiduels $F/\kappa(\mathfrak p)$
est de type fini et $\Omega_{F/\kappa(\mathfrak p)} = 0$. Cette extension est
donc finie séparable d'après
Compléments d'algèbre, lemme \ref{more-algebra-lemma-cartier-equality}.
Ainsi,
$k[y_1, \ldots, y_m]_\mathfrak p/\mathfrak p^n k[y_1, \ldots, y_m]_\mathfrak p
\to A$
est fini d'après Algèbre, lemme
\ref{algebra-lemma-essentially-of-finite-type-into-artinian-local}.
Enfin, nous concluons qu'il est étale d'après
Algèbre, lemme \ref{algebra-lemma-characterize-etale}.
Comme un morphisme d'anneaux étale est certainement essentiellement lisse, la démonstration est achevée.
\end{proof}

\begin{lemma}
\label{smoothing-lemma-enlarge-solution-modulo}
Soient $\varphi : k[y_1, \ldots, y_m] \to \Lambda$, $n$, $\mathfrak q$,
$\mathfrak p$ et
$$
k[y_1, \ldots, y_m]_\mathfrak p/\mathfrak p^n \to
D \to \Lambda_\mathfrak q/\mathfrak q^n \Lambda_\mathfrak q
$$
comme dans le Lemme \ref{smoothing-lemma-solution-modulo}. Alors, pour tout
$\lambda \in \Lambda \setminus \mathfrak q$
il existe un entier $q > 0$ et une factorisation
$$
k[y_1, \ldots, y_m]_\mathfrak p/\mathfrak p^n \to
D \to D' \to \Lambda_\mathfrak q/\mathfrak q^n \Lambda_\mathfrak q
$$
tels que $D \to D'$ soit un morphisme essentiellement lisse d'anneaux locaux artiniens,
que la dernière flèche soit plate et que $\lambda^q$ appartienne à $D'$.
\end{lemma}

\begin{proof}
Posons $\bar \Lambda = \Lambda_\mathfrak q/\mathfrak q^n\Lambda_\mathfrak q$.
Soit $\bar \lambda$ l'image de $\lambda$ dans $\bar \Lambda$.
Soit $\alpha \in \kappa(\mathfrak q)$ l'image de $\lambda$ dans le
corps résiduel.
Soit $k \subset F \subset \kappa(\mathfrak q)$ le corps résiduel de $D$.
Si $\alpha$ appartient à $F$, nous pouvons trouver
$x \in D$ tel que $x \bar\lambda = 1 \bmod \mathfrak q$. Par conséquent,
$(x \bar \lambda)^q = 1 \bmod (\mathfrak q)^q$ si $q$ est divisible par $p$.
Ainsi, $\bar\lambda^q$ appartient à $D$. Si $\alpha$ est
transcendant sur $F$, nous pouvons prendre $D' = (D[\bar \lambda])_\mathfrak m$
égal au sous-anneau engendré par $D$ et $\bar \lambda$, localisé
en $\mathfrak m = D[\bar \lambda] \cap \mathfrak q \bar \Lambda$.
Cela convient car, dans ce cas, $D[\bar \lambda]$ est en fait une algèbre de polynômes
sur $D$. Enfin, si $\lambda \bmod \mathfrak q$ est
algébrique sur $F$, nous pouvons trouver une puissance de $p$, notée $q$, telle que
$\alpha^q$ soit séparable algébrique sur $F$ ; voir
Corps, section \ref{fields-section-algebraic}.
Remarquons que $D$ et $\bar\Lambda$ sont des anneaux locaux henséliens ; voir
Algèbre, lemme \ref{algebra-lemma-local-dimension-zero-henselian}.
Soit $D \to D'$ une extension étale finie
dont l'extension de corps résiduels est $F(\alpha^q)/F$ ; voir
Algèbre, lemme \ref{algebra-lemma-henselian-cat-finite-etale}.
Comme $\bar\Lambda$ est hensélien et que $F(\alpha^q)$ est contenu
dans son corps résiduel, nous pouvons trouver une factorisation
$D' \to \bar \Lambda$. La première partie de l'argument montre que
$\bar\lambda^{qq'} \in D'$ pour un certain $q' > 0$.
\end{proof}

\begin{lemma}
\label{smoothing-lemma-resolve-general}
Soit $k \to A \to \Lambda \supset \mathfrak q$ comme dans la
Situation \ref{smoothing-situation-local}, où
\begin{enumerate}
\item $k$ est un corps de caractéristique $p > 0$,
\item $\Lambda$ est noethérien et géométriquement régulier sur $k$,
\item $\mathfrak q$ est minimal au-dessus de $\mathfrak h_A$.
\end{enumerate}
Alors $k \to A \to \Lambda \supset \mathfrak q$ peut être résolu.
\end{lemma}

\begin{proof}
Le lemme se démontre par les étapes suivantes, dans l'ordre indiqué.
Nous justifierons chacune de ces étapes ci-dessous.
\begin{enumerate}
\item
\label{smoothing-item-power}
Choisissons un entier $N > 0$ tel que
$\mathfrak q^N\Lambda_\mathfrak q \subset H_{A/k}\Lambda_\mathfrak q$.
\item
\label{smoothing-item-generators}
Choisissons des générateurs $a_1, \ldots, a_t \in A$ de l'idéal $H_{A/k}$.
\item
\label{smoothing-item-dimension}
Posons $d = \dim(\Lambda_\mathfrak q)$.
\item
\label{smoothing-item-standardizer}
Posons $B = A[x_1, \ldots, x_d, z_{ij}]/(x_i^{2N} - \sum z_{ij}a_j)$.
\item
\label{smoothing-item-elkik}
Considérons $B$ comme une $k[x_1, \ldots, x_d]$-algèbre et soit
$B \to C$ comme dans le
Lemme \ref{smoothing-lemma-improve-presentation}.
Nous obtenons aussi une section $C \to B$.
\item
\label{smoothing-item-strictly-standard}
Choisissons $c > 0$ tel que chaque $x_i^c$
soit strictement standard dans $C$ sur $k[x_1, \ldots, x_d]$.
\item
\label{smoothing-item-set-n}
Posons $n = N + dc$ et $e = 8c$.
\item
\label{smoothing-item-set-E}
Soit $E \subset \Lambda_\mathfrak q/\mathfrak q^n\Lambda_\mathfrak q$
l'ensemble des images de générateurs de $A$ comme $k$-algèbre.
\item
\label{smoothing-item-NP}
Choisissons un entier $m$, un morphisme de $k$-algèbres
$\varphi : k[y_1, \ldots, y_m] \to \Lambda$
et une factorisation par des anneaux locaux artiniens
$$
k[y_1, \ldots, y_m]_\mathfrak p/\mathfrak p^n k[y_1, \ldots, y_m]_\mathfrak p
\to D \to
\Lambda_\mathfrak q/\mathfrak q^n\Lambda_\mathfrak q
$$
telle que la première flèche soit essentiellement lisse, la seconde plate,
que $E$ soit contenu dans $D$, et, avec $\mathfrak p = \varphi^{-1}(\mathfrak q)$,
que le morphisme $k[y_1, \ldots, y_m]_\mathfrak p \to \Lambda_\mathfrak q$ soit
plat et que $\mathfrak p \Lambda_\mathfrak q = \mathfrak q \Lambda_\mathfrak q$.
\item
\label{smoothing-item-choose-pii}
Choisissons $\pi_1, \ldots, \pi_d \in \mathfrak p$ dont les images forment un
système régulier de paramètres de $k[y_1, \ldots, y_m]_\mathfrak p$.
\item
\label{smoothing-item-set-R}
Soit $R = k[y_1, \ldots, y_m, t_1, \ldots, t_d]$ et $\gamma_i = \pi_i t_i$.
\item
\label{smoothing-item-modify-pii}
Si nécessaire, modifions le choix des $\pi_i$ de sorte que,
pour $i = 1, \ldots, d$, nous ayons
$$
\text{Ann}_{R/(\gamma_1^e, \ldots, \gamma_{i - 1}^e)R}(\gamma_i)
=
\text{Ann}_{R/(\gamma_1^e, \ldots, \gamma_{i - 1}^e)R}(\gamma_i^2)
$$
\item
\label{smoothing-item-choose-deltai}
Il existe $\delta_1, \ldots, \delta_d \in \Lambda$,
$\delta_i \not \in \mathfrak q$, et une factorisation
$D \to D' \to \Lambda_\mathfrak q/\mathfrak q^n\Lambda_\mathfrak q$
avec $D'$ local artinien, $D \to D'$ essentiellement lisse, la dernière flèche
$D' \to \Lambda_\mathfrak q/\mathfrak q^n\Lambda_\mathfrak q$
étant plate, et tels que, si l'on pose $\pi_i' = \delta_i \pi_i$, on ait pour
$i = 1, \ldots, d$ :
\begin{enumerate}
\item $(\pi_i')^{2N} = \sum a_j\lambda_{ij}$ dans $\Lambda$, où
$\lambda_{ij} \bmod \mathfrak q^n\Lambda_\mathfrak q$ est un élément de $D'$,
\item $\text{Ann}_{\Lambda/({\pi'}_1^e, \ldots, {\pi'}_{i - 1}^e)}({\pi'}_i) =
\text{Ann}_{\Lambda/({\pi'}_1^e, \ldots, {\pi'}_{i - 1}^e)}({\pi'}_i^2)$,
\item $\delta_i \bmod \mathfrak q^n\Lambda_\mathfrak q$ est un élément de $D'$.
\end{enumerate}
\item
\label{smoothing-item-map-B-Lambda}
Définissons $B \to \Lambda$ en envoyant $x_i$ sur $\pi'_i$ et
$z_{ij}$ sur le $\lambda_{ij}$ trouvé ci-dessus. Définissons $C \to \Lambda$
en composant le morphisme $B \to \Lambda$ avec la rétraction $C \to B$.
\item
\label{smoothing-item-set-map-R}
Définissons $R \to \Lambda$ par $\varphi$ sur $k[y_1, \ldots, y_m]$
et en envoyant $t_i$ sur $\delta_i$. Introduisons en outre un morphisme
$$
k[x_1, \ldots, x_d]
\longrightarrow
R = k[y_1, \ldots, y_m, t_1, \ldots, t_d]
$$
en envoyant $x_i$ sur $\gamma_i = \pi_i t_i$.
\item
\label{smoothing-item-first-solve}
Il suffit de résoudre
$$
R
\to
C \otimes_{k[x_1, \ldots, x_d]} R
\to
\Lambda \supset \mathfrak q
$$
\item
\label{smoothing-item-set-I}
Posons $I = (\gamma_1^e, \ldots, \gamma_d^e) \subset R$.
\item
\label{smoothing-item-second-resolve}
Il suffit de résoudre
$$
R/I
\to
C \otimes_{k[x_1, \ldots, x_d]} R/I
\to
\Lambda/I\Lambda \supset \mathfrak q/I\Lambda
$$
\item
\label{smoothing-item-set-r}
Notons $\mathfrak r \subset R = k[y_1, \ldots, y_m, t_1, \ldots, t_d]$
l'image inverse de $\mathfrak q$.
\item
\label{smoothing-item-third-resolve}
Il suffit de résoudre
$$
(R/I)_\mathfrak r \to
C \otimes_{k[x_1, \ldots, x_d]} (R/I)_\mathfrak r \to
\Lambda_\mathfrak q/I\Lambda_\mathfrak q
\supset
\mathfrak q\Lambda_\mathfrak q/I\Lambda_\mathfrak q
$$
\item
\label{smoothing-item-fourth-resolve}
Posons $J = (\pi_1^e, \ldots, \pi_d^e)$ dans $k[y_1, \ldots, y_m]$.
\item
\label{smoothing-item-fifth-resolve}
Il suffit de résoudre
$$
(R/JR)_\mathfrak p \to
C \otimes_{k[x_1, \ldots, x_d]} (R/JR)_\mathfrak p \to
\Lambda_\mathfrak q/J\Lambda_\mathfrak q
\supset
\mathfrak q\Lambda_\mathfrak q/J\Lambda_\mathfrak q
$$
\item
\label{smoothing-item-sixth-resolve}
Il suffit de résoudre
$$
(R/\mathfrak p^nR)_\mathfrak p \to
C \otimes_{k[x_1, \ldots, x_d]} (R/\mathfrak p^nR)_\mathfrak p \to
\Lambda_\mathfrak q/\mathfrak q^n\Lambda_\mathfrak q
\supset
\mathfrak q\Lambda_\mathfrak q/\mathfrak q^n\Lambda_\mathfrak q
$$
\item
\label{smoothing-item-seventh-resolve}
Il suffit de résoudre
$$
(R/\mathfrak p^nR)_\mathfrak p \to
B \otimes_{k[x_1, \ldots, x_d]} (R/\mathfrak p^nR)_\mathfrak p \to
\Lambda_\mathfrak q/\mathfrak q^n\Lambda_\mathfrak q
\supset
\mathfrak q\Lambda_\mathfrak q/\mathfrak q^n\Lambda_\mathfrak q
$$
\item
\label{smoothing-item-eighth-resolve}
On munit l'anneau $D'[t_1, \ldots, t_d]$ d'une structure de
$R_\mathfrak p/\mathfrak p^nR_\mathfrak p$-algèbre au moyen du morphisme donné
$k[y_1, \ldots, y_m]_\mathfrak p/\mathfrak p^n k[y_1, \ldots, y_m]_\mathfrak p
\to D'$ et en envoyant $t_i$ sur $t_i$. Il suffit de trouver une factorisation
$$
B \otimes_{k[x_1, \ldots, x_d]} (R/\mathfrak p^nR)_\mathfrak p
\to D'[t_1, \ldots, t_d] \to
\Lambda_\mathfrak q/\mathfrak q^n\Lambda_\mathfrak q
$$
où la seconde flèche envoie $t_i$ sur $\delta_i$ et induit le
morphisme donné $D' \to \Lambda_\mathfrak q/\mathfrak q^n\Lambda_\mathfrak q$.
\item
\label{smoothing-item-done}
Une telle factorisation existe grâce au choix de $D'$ fait ci-dessus.
\end{enumerate}
Nous donnons maintenant la justification de chacune des étapes, en omettant
celles qui ne font qu'introduire des notations.

\medskip\noindent
Pour (\ref{smoothing-item-power}). Cela est possible puisque $\mathfrak q$ est minimal
au-dessus de $\mathfrak h_A = \sqrt{H_{A/k}\Lambda}$.

\medskip\noindent
Pour (\ref{smoothing-item-strictly-standard}). Remarquons que $A_{a_j}$ est lisse
sur $k$. Ainsi $B_{a_j}$, qui est isomorphe à une algèbre de polynômes
sur $A_{a_j}[x_1, \ldots, x_d]$, est lisse sur
$k[x_1, \ldots, x_d]$. Par conséquent, $B_{x_i}$ est lisse sur $k[x_1, \ldots, x_d]$.
D'après le Lemme \ref{smoothing-lemma-improve-presentation},
nous voyons que $C_{x_i}$ est lisse sur $k[x_1, \ldots, x_d]$
et que son module des différentielles de Kähler est libre de type fini. Une certaine puissance de
$x_i$ est donc strictement standard dans $C$ sur $k[x_1, \ldots, x_d]$
d'après le Lemme \ref{smoothing-lemma-compare-standard}.

\medskip\noindent
Pour (\ref{smoothing-item-NP}). Cela résulte de l'application du Lemme \ref{smoothing-lemma-solution-modulo}.

\medskip\noindent
Pour (\ref{smoothing-item-choose-pii}). Puisque
$k[y_1, \ldots, y_m]_\mathfrak p \to \Lambda_\mathfrak q$
est plat et que $\mathfrak p \Lambda_\mathfrak q = \mathfrak q \Lambda_\mathfrak q$
par construction,
nous avons $\dim(k[y_1, \ldots, y_m]_\mathfrak p) = d$ d'après
Algèbre, lemme \ref{algebra-lemma-dimension-base-fibre-equals-total}.
Nous pouvons donc trouver $\pi_1, \ldots, \pi_d \in \mathfrak p$ formant
un système régulier de paramètres de $k[y_1, \ldots, y_m]_\mathfrak p$.

\medskip\noindent
Pour (\ref{smoothing-item-modify-pii}). D'après
Algèbre, lemme \ref{algebra-lemma-regular-ring-CM},
toute permutation de la suite $\pi_1, \ldots, \pi_d$ est une
suite régulière dans $k[y_1, \ldots, y_m]_\mathfrak p$. Par conséquent,
$\gamma_1 = \pi_1 t_1, \ldots, \gamma_d = \pi_d t_d$ est une suite
régulière dans
$R_\mathfrak p = k[y_1, \ldots, y_m]_\mathfrak p[t_1, \ldots, t_d]$ ; voir
Algèbre, lemme \ref{algebra-lemma-regular-sequence-in-polynomial-ring}.
Soit $S = k[y_1, \ldots, y_m] \setminus \mathfrak p$, de sorte que
$R_\mathfrak p = S^{-1}R$. Remarquons que $\pi_1, \ldots, \pi_d$
et $\gamma_1, \ldots, \gamma_d$
restent des suites régulières si nous multiplions nos $\pi_i$ par des éléments de $S$.
Supposons que
$$
\text{Ann}_{R/(\gamma_1^e, \ldots, \gamma_{i - 1}^e)R}(\gamma_i)
=
\text{Ann}_{R/(\gamma_1^e, \ldots, \gamma_{i - 1}^e)R}(\gamma_i^2)
$$
soit satisfaite pour $i = 1, \ldots, t$, pour un certain $t \in \{0, \ldots, d\}$. Remarquons que
$\gamma_1^e, \ldots, \gamma_t^e, \gamma_{t + 1}$ est une
suite régulière dans $S^{-1}R$ d'après
Algèbre, lemme \ref{algebra-lemma-regular-sequence-powers}.
Nous voyons donc que
$$
\text{Ann}_{S^{-1}R/(\gamma_1^e, \ldots, \gamma_{i - 1}^e)}(\gamma_i) =
\text{Ann}_{S^{-1}R/(\gamma_1^e, \ldots, \gamma_{i - 1}^e)}(\gamma_i^2).
$$
Nous obtenons ainsi
$$
\text{Ann}_{R/(\gamma_1^e, \ldots, \gamma_t^e)R}(\gamma_{t + 1})
=
\text{Ann}_{R/(\gamma_1^e, \ldots, \gamma_t^e)R}(\gamma_{t + 1}^2)
$$
après avoir remplacé $\pi_{t + 1}$ par $s\pi_{t + 1}$ pour un certain $s \in S$, d'après le
Lemme \ref{smoothing-lemma-ogoma}. Par récurrence sur $t$, ceci produit la
suite souhaitée.

\medskip\noindent
Pour (\ref{smoothing-item-choose-deltai}). Posons $S = \Lambda \setminus \mathfrak q$,
de sorte que $\Lambda_\mathfrak q = S^{-1}\Lambda$. Posons
$\bar \Lambda = \Lambda_\mathfrak q/\mathfrak q^n \Lambda_\mathfrak q$.
Supposons que nous ayons un $t \in \{0, \ldots, d\}$, des éléments
$\delta_1, \ldots, \delta_t \in S$ et une factorisation
$D \to D' \to \bar \Lambda$ comme dans (\ref{smoothing-item-choose-deltai}),
tels que (a), (b), (c) soient satisfaites pour $i = 1, \ldots, t$. Nous avons
$\pi_{t + 1}^N \in H_{A/k}\Lambda_\mathfrak q$
car $\mathfrak q^N \Lambda_\mathfrak q \subset H_{A/k}\Lambda_\mathfrak q$
d'après (\ref{smoothing-item-power}). Par conséquent,
$\pi_{t + 1}^N \in H_{A/k} \bar\Lambda$, puis
$\pi_{t + 1}^N \in H_{A/k}D'$, puisque $D' \to \bar \Lambda$
est fidèlement plat ; voir
Algèbre, lemme \ref{algebra-lemma-faithfully-flat-universally-injective}.
Rappelons que $H_{A/k} = (a_1, \ldots, a_t)$.
Écrivons $\pi_{t + 1}^N = \sum a_j d_j$ dans $D'$ et choisissons
$c_j \in \Lambda_\mathfrak q$ relevant $d_j \in D'$. Alors
$\pi_{t + 1}^N = \sum c_j a_j + \epsilon$, avec
$\epsilon \in \mathfrak q^n\Lambda_\mathfrak q \subset
\mathfrak q^{n - N}H_{A/k}\Lambda_\mathfrak q$.
Écrivons $\epsilon = \sum a_j c'_j$ pour certains
$c'_j \in \mathfrak q^{n - N}\Lambda_\mathfrak q$.
Ainsi, $\pi_{t + 1}^{2N} = \sum (\pi_{t + 1}^N c_j + \pi_{t + 1}^N c'_j) a_j$.
Remarquons que $\pi_{t + 1}^Nc'_j$ a pour image zéro dans $\bar \Lambda$ ; cette observation
élémentaire mais essentielle garantira plus loin que (a) est satisfaite.
Choisissons maintenant $s \in S$ tel qu'il existe des
$\mu_{t + 1j} \in \Lambda$ vérifiant, d'une part,
$\pi_{t + 1}^N c_j + \pi_{t + 1}^N c'_j = \mu_{t + 1j}/s^{2N}$
dans $S^{-1}\Lambda$ et, d'autre part,
$(s \pi_{t + 1})^{2N} = \sum \mu_{t + 1j}a_j$
dans $\Lambda$ (détail mineur omis). Nous pouvons en outre remplacer $s$ par
une puissance et agrandir $D'$ de sorte que $s$ ait pour image un élément de $D'$.
Avec ces choix, $\mu_{t + 1j}$ a pour image $s^{2N}d_j$, qui est
un élément de $D'$. Remarquons que $\pi_1, \ldots, \pi_d$ forment un système
régulier de paramètres dans $S^{-1}\Lambda$ par notre
choix de $\varphi$. Ainsi, $\pi_1, \ldots, \pi_d$ forment une suite régulière
dans $\Lambda_\mathfrak q$ d'après
Algèbre, lemme \ref{algebra-lemma-regular-ring-CM}.
Il s'ensuit que ${\pi'}_1^e, \ldots, {\pi'}_t^e, s\pi_{t + 1}$ est une
suite régulière dans $S^{-1}\Lambda$ d'après
Algèbre, lemme \ref{algebra-lemma-regular-sequence-powers}.
Nous obtenons donc
$$
\text{Ann}_{S^{-1}\Lambda/({\pi'}_1^e, \ldots, {\pi'}_t^e)}(s\pi_{t + 1}) =
\text{Ann}_{S^{-1}\Lambda/({\pi'}_1^e, \ldots, {\pi'}_t^e)}((s\pi_{t + 1})^2).
$$
Nous pouvons donc appliquer le Lemme \ref{smoothing-lemma-ogoma} afin de trouver $s' \in S$
tel que
$$
\text{Ann}_{\Lambda/({\pi'}_1^e, \ldots, {\pi'}_t^e)}((s')^qs\pi_{t + 1})
=
\text{Ann}_{\Lambda/({\pi'}_1^e, \ldots, {\pi'}_t^e)}(((s')^qs\pi_{t + 1})^2).
$$
pour tout $q > 0$. D'après le Lemme \ref{smoothing-lemma-enlarge-solution-modulo},
nous pouvons choisir $q$ et agrandir $D'$ de sorte que $(s')^q$ ait pour image un élément
de $D'$. En posant $\delta_{t + 1} = (s')^qs$, nous concluons que
(a), (b), (c) sont satisfaites pour $i = 1, \ldots, t + 1$. Pour (a), remarquons que
$\lambda_{t + 1j} = (s')^{2Nq}\mu_{t + 1j}$ convient.
Une récurrence sur $t$ achève la démonstration.

\medskip\noindent
Pour (\ref{smoothing-item-first-solve}). Par construction, le radical de
$H_{(C \otimes_{k[x_1, \ldots, x_d]} R)/R} \Lambda$ contient
$\mathfrak h_A$. En effet, les éléments $a_j \in H_{A/k}$
ont pour images des éléments de $H_{B/k[x_1, \ldots, x_d]}$, puis des éléments
de $H_{C/k[x_1, \ldots, x_d]}$ ; par conséquent, les $a_j \otimes 1$ ont pour images des éléments de
$H_{C \otimes_{k[x_1, \ldots, x_d]} R/R}$. De plus, si nous avons une solution
$C \otimes_{k[x_1, \ldots, x_d]} R \to T \to \Lambda$ de
$$
R
\to
C \otimes_{k[x_1, \ldots, x_d]} R
\to
\Lambda \supset \mathfrak q
$$
alors $H_{T/R} \subset H_{T/k}$, puisque $R$ est lisse sur $k$.
Ainsi, $T$ sera aussi une solution de
la situation initiale $k \to A \to \Lambda \supset \mathfrak q$.

\medskip\noindent
Pour (\ref{smoothing-item-second-resolve}). Cela résulte de l'application du
Lemme \ref{smoothing-lemma-lift-solution} à
$R \to C \otimes_{k[x_1, \ldots, x_d]} R
\to \Lambda \supset \mathfrak q$ et à la suite d'éléments
$\gamma_1^c, \ldots, \gamma_d^c$. Remarquons que, puisque les $x_i^c$
sont strictement standard dans $C$ sur $k[x_1, \ldots, x_d]$, les éléments
$\gamma_i^c$ sont strictement standard dans $C \otimes_{k[x_1, \ldots, x_d]} R$
sur $R$ d'après le Lemme \ref{smoothing-lemma-strictly-standard-base-change}.
L'autre hypothèse du Lemme \ref{smoothing-lemma-lift-solution} est satisfaite grâce aux étapes
(\ref{smoothing-item-modify-pii}) et (\ref{smoothing-item-choose-deltai}).

\medskip\noindent
Pour (\ref{smoothing-item-third-resolve}). Appliquons le Lemme \ref{smoothing-lemma-delocalize-height-zero}
à la situation de (\ref{smoothing-item-second-resolve}). Dans la suite des
arguments, l'anneau d'arrivée est local artinien ; nous cherchons donc
une factorisation par une algèbre $T$ lisse sur l'anneau de départ.

\medskip\noindent
Pour (\ref{smoothing-item-fifth-resolve}).
Supposons que $C \otimes_{k[x_1, \ldots, x_d]} (R/JR)_\mathfrak p \to
T \to \Lambda_\mathfrak q/J\Lambda_\mathfrak q$ soit une solution de
$$
(R/JR)_\mathfrak p \to
C \otimes_{k[x_1, \ldots, x_d]} (R/JR)_\mathfrak p \to
\Lambda_\mathfrak q/J\Lambda_\mathfrak q
\supset
\mathfrak q\Lambda_\mathfrak q/J\Lambda_\mathfrak q
$$
Alors $C \otimes_{k[x_1, \ldots, x_d]} (R/I)_\mathfrak r \to T_\mathfrak r \to
\Lambda_\mathfrak q/I\Lambda_\mathfrak q$
est une solution de la situation de (\ref{smoothing-item-third-resolve}).

\medskip\noindent
Pour (\ref{smoothing-item-sixth-resolve}). Notre entier $n = N + dc$ est assez grand pour que
$\mathfrak p^nk[y_1, \ldots, y_m]_\mathfrak p \subset J_\mathfrak p$
et $\mathfrak q^n \Lambda_\mathfrak q \subset J\Lambda_\mathfrak q$.
Par conséquent, si nous avons une solution
$C \otimes_{k[x_1, \ldots, x_d]} (R/\mathfrak p^nR)_\mathfrak p \to
T \to \Lambda_\mathfrak q/\mathfrak q^n\Lambda_\mathfrak q$
de (\ref{smoothing-item-fifth-resolve},
nous pouvons prendre $T/JT$ comme solution de
(\ref{smoothing-item-sixth-resolve}).

\medskip\noindent
Pour (\ref{smoothing-item-seventh-resolve}). Cela est vrai parce que nous avons une
section $C \to B$ dans la catégorie des $R$-algèbres.

\medskip\noindent
Pour (\ref{smoothing-item-eighth-resolve}). Cela est vrai parce que $D'$ est
essentiellement lisse sur l'anneau local artinien
$k[y_1, \ldots, y_m]_\mathfrak p/\mathfrak p^n k[y_1, \ldots, y_m]_\mathfrak p$
et que
$$
R_\mathfrak p/\mathfrak p^nR_\mathfrak p =
k[y_1, \ldots, y_m]_\mathfrak p/
\mathfrak p^n k[y_1, \ldots, y_m]_\mathfrak p[t_1, \ldots, t_d].
$$
Ainsi, $D'[t_1, \ldots, t_d]$ est une colimite filtrante de
$R_\mathfrak p/\mathfrak p^nR_\mathfrak p$-algèbres lisses et
$B \otimes_{k[x_1, \ldots, x_d]} (R_\mathfrak p/\mathfrak p^nR_\mathfrak p)$
se factorise par l'une d'elles.

\medskip\noindent
Pour (\ref{smoothing-item-done}). La dernière subtilité de la démonstration est que nous ne pouvons
pas simplement utiliser le morphisme $B \to D'$ qui envoie $x_i$ sur l'image de $\pi_i'$
dans $D'$ et $z_{ij}$ sur l'image de $\lambda_{ij}$ dans $D'$, 
car nous avons besoin que le diagramme
$$
\xymatrix{
B \ar[r] & D'[t_1, \ldots, t_d] \\
k[x_1, \ldots, x_d] \ar[r] \ar[u] &
R_\mathfrak p/\mathfrak p^nR_\mathfrak p \ar[u]
}
$$
soit commutatif et que la composée
$B \to D'[t_1, \ldots, t_d] \to
\Lambda_\mathfrak q/\mathfrak q^n\Lambda_\mathfrak q$
soit le morphisme de (\ref{smoothing-item-map-B-Lambda}).
Cela nous impose d'envoyer $x_i$ sur l'image de
$\pi_i t_i$ dans $D'[t_1, \ldots, t_d]$.
Nous envoyons donc $z_{ij}$ sur l'image de
$\lambda_{ij} t_i^{2N} / \delta_i^{2N}$ dans $D'[t_1, \ldots, t_d]$
et tout est clair.
\end{proof}







\section{Le théorème principal}
\label{smoothing-section-main}

\noindent
Dans cette section, nous concluons la discussion.

\begin{theorem}[Popescu]
\label{smoothing-theorem-popescu}
Tout morphisme régulier entre anneaux noethériens est une colimite filtrante
de morphismes d'anneaux lisses.
\end{theorem}

\begin{proof}
D'après le lemme \ref{smoothing-lemma-reduce-to-field},
il suffit de démontrer l'énoncé pour $k \to \Lambda$,
où $\Lambda$ est noethérien et géométriquement régulier sur $k$.
Soit $k \to A \to \Lambda$ une factorisation, avec $A$ une
$k$-algèbre de type fini. Il suffit de construire une factorisation
$A \to B \to \Lambda$, avec $B$ de type fini, telle que
$\mathfrak h_B = \Lambda$ ; voir le lemme \ref{smoothing-lemma-final-solve}.
Nous pouvons donc procéder par récurrence noethérienne sur l'idéal
$\mathfrak h_A$. Choisissons un idéal premier
$\mathfrak q \supset \mathfrak h_A$ tel que
$\mathfrak q$ soit minimal au-dessus de $\mathfrak h_A$.
Il suffit alors de résoudre $k \to A \to \Lambda \supset \mathfrak q$
(au sens défini dans le texte qui suit la situation \ref{smoothing-situation-local}).
Si la caractéristique de $k$ est nulle, cela résulte du
lemme \ref{smoothing-lemma-resolve-special}.
Si la caractéristique de $k$ est $p > 0$, cela résulte du
lemme \ref{smoothing-lemma-resolve-general}.
\end{proof}




\section{La propriété d'approximation pour les G-anneaux}
\label{smoothing-section-approximation-G-rings}

\noindent
Soit $R$ un anneau local noethérien. Alors $R$ est un G-anneau si et
seulement si le morphisme d'anneaux $R \to R^\wedge$ est régulier ; voir
Compléments d'algèbre, lemme
\ref{more-algebra-lemma-check-G-ring-maximal-ideals}.
Dans ce cas, l'hensélisation $R^h$ et l'hensélisation stricte
$R^{sh}$ de $R$ sont des G-anneaux ; voir
Compléments d'algèbre, lemme
\ref{more-algebra-lemma-henselization-G-ring}.
En outre, toute algèbre essentiellement de type fini sur un corps, sur un
anneau local complet, sur $\mathbf{Z}$ ou sur un anneau de Dedekind de
caractéristique nulle est un G-anneau ; voir
Compléments d'algèbre, proposition
\ref{more-algebra-proposition-ubiquity-G-ring}.
Nous disposons ainsi d'un vaste choix d'anneaux auxquels s'applique le
résultat ci-dessous.

\medskip\noindent
Soit $R$ un anneau. Soient $f_1, \ldots, f_m \in R[x_1, \ldots, x_n]$.
Soit $S$ une $R$-algèbre. Dans cette situation, nous disons qu'un vecteur
$(a_1, \ldots, a_n) \in S^n$ est une {\it solution dans $S$}
si et seulement si
$$
f_j(a_1, \ldots, a_n) = 0 \text{ dans } S, \text{ pour }
j = 1, \ldots, m
$$
L'une des questions importantes de la géométrie algébrique consiste
naturellement à déterminer quand des systèmes d'équations polynomiales
ont des solutions. Le théorème suivant affirme que l'existence de solutions
dans le complété d'un anneau local noethérien suffit souvent à montrer
qu'il en existe dans l'hensélisation de cet anneau.

\begin{theorem}
\label{smoothing-theorem-approximation-property}
Soit $R$ un anneau local noethérien. Soient
$f_1, \ldots, f_m \in R[x_1, \ldots, x_n]$.
Supposons que $(a_1, \ldots, a_n) \in (R^\wedge)^n$ soit une solution
dans $R^\wedge$. Si $R$ est un G-anneau hensélien, alors, pour tout entier
$N$, il existe une solution $(b_1, \ldots, b_n) \in R^n$ dans $R$ telle que
$a_i - b_i \in \mathfrak m^NR^\wedge$.
\end{theorem}

\begin{proof}
Choisissons $c_i \in R$ tel que $a_i - c_i \in \mathfrak m^N$.
Choisissons des générateurs $\mathfrak m^N = (d_1, \ldots, d_M)$.
Écrivons $a_i = c_i + \sum a_{i, l} d_l$.
Considérons l'anneau de polynômes $R[x_{i, l}]$ et les éléments
$$
g_j = f_j(c_1 + \sum x_{1, l} d_l , \ldots, c_n + \sum x_{n, l} d_l)
\in R[x_{i, l}]
$$
Le système d'équations $g_j = 0$ admet la solution $(a_{i, l})$.
Supposons que nous puissions montrer que $g_j$ admet une solution
$(b_{i, l})$ dans $R$.
Il s'ensuit que $b_i = c_i + \sum b_{i, l}d_l$ est une solution
de $f_j = 0$ congrue à $a_i$ modulo $\mathfrak m^N$.
Il suffit donc de montrer que l'existence d'une solution sur $R^\wedge$
entraîne l'existence d'une solution sur $R$.

\medskip\noindent
Soit $A \subset R^\wedge$ la sous-$R$-algèbre engendrée par
$a_1, \ldots, a_n$. Puisque nous avons supposé que $R$ est un G-anneau,
c'est-à-dire que $R \to R^\wedge$ est régulier, il existe une factorisation
$$
A \to B \to R^\wedge
$$
où $B$ est lisse sur $R$ ; voir le théorème \ref{smoothing-theorem-popescu}.
Notons $\kappa = R/\mathfrak m$ le corps résiduel. C'est aussi
le corps résiduel de $R^\wedge$, si bien que nous obtenons un diagramme
commutatif
$$
\xymatrix{
B \ar[rd] \ar@{..>}[r] & R' \ar@{..>}[d] \\
R \ar[r] \ar[u] & \kappa
}
$$
Comme la flèche verticale est lisse,
Compléments d'algèbre, lemme
\ref{more-algebra-lemma-lift-section-smooth-morphism},
il existe un morphisme d'anneaux \'etale $R \to R'$ qui induit un isomorphisme
$R/\mathfrak m \to R'/\mathfrak mR'$ et un morphisme de $R$-algèbres
$B \to R'$ rendant commutatif le diagramme ci-dessus.
Puisque $R$ est hensélien, le morphisme $R \to R'$ possède une section ; voir
Algèbre, lemme \ref{algebra-lemma-characterize-henselian}.
Soit $b_i \in R$ l'image de $a_i$ par les morphismes d'anneaux
$A \to B \to R' \to R$. Puisque tous ces morphismes sont des morphismes de
$R$-algèbres, $(b_1, \ldots, b_n)$ est une solution dans $R$.
\end{proof}

\noindent
Étant donnés un anneau local noethérien $(R, \mathfrak m)$, un morphisme
d'anneaux \'etale $R \to R'$ et un idéal maximal
$\mathfrak m' \subset R'$ au-dessus de $\mathfrak m$, tels que
$\kappa(\mathfrak m) = \kappa(\mathfrak m')$. Nous avons alors les inclusions
$$
R \subset R_{\mathfrak m'} \subset R^h \subset R^\wedge,
$$
d'après Algèbre, lemme \ref{algebra-lemma-henselian-functorial-prepare}, et
Compléments d'algèbre, lemme
\ref{more-algebra-lemma-henselization-noetherian}.

\begin{theorem}
\label{smoothing-theorem-approximation-property-variant}
Soit $R$ un anneau local noethérien. Soient
$f_1, \ldots, f_m \in R[x_1, \ldots, x_n]$.
Supposons que $(a_1, \ldots, a_n) \in (R^\wedge)^n$ soit une solution.
Si $R$ est un G-anneau, alors, pour tout entier $N$, il existe
\begin{enumerate}
\item un morphisme d'anneaux \'etale $R \to R'$,
\item un idéal maximal $\mathfrak m' \subset R'$ au-dessus de $\mathfrak m$,
\item une solution $(b_1, \ldots, b_n) \in (R')^n$ dans $R'$,
\end{enumerate}
tels que $\kappa(\mathfrak m) = \kappa(\mathfrak m')$ et
$a_i - b_i \in (\mathfrak m')^NR^\wedge$.
\end{theorem}

\begin{proof}
Nous pourrions déduire ce théorème du théorème
\ref{smoothing-theorem-approximation-property}, en utilisant le fait que
l'hensélisation $R^h$ est un G-anneau d'après
Compléments d'algèbre, lemme
\ref{more-algebra-lemma-henselization-G-ring}, et en écrivant $R^h$ comme une
colimite filtrante d'extensions \'etales $R'$. Nous préférons en donner une
démonstration en reprenant, dans ce cas, celle du théorème précédent.

\medskip\noindent
Choisissons $c_i \in R$ tel que $a_i - c_i \in \mathfrak m^N$.
Choisissons des générateurs $\mathfrak m^N = (d_1, \ldots, d_M)$.
Écrivons $a_i = c_i + \sum a_{i, l} d_l$.
Considérons l'anneau de polynômes $R[x_{i, l}]$ et les éléments
$$
g_j = f_j(c_1 + \sum x_{1, l} d_l , \ldots, c_n + \sum x_{n, l} d_l)
\in R[x_{i, l}]
$$
Le système d'équations $g_j = 0$ admet la solution $(a_{i, l})$.
Supposons que nous puissions montrer que $g_j$ admet une solution
$(b_{i, l})$ dans $R'$, pour un morphisme d'anneaux \'etale $R \to R'$ muni
d'un idéal maximal $\mathfrak m'$ tel que
$\kappa(\mathfrak m) = \kappa(\mathfrak m')$.
Il s'ensuit que $b_i = c_i + \sum b_{i, l}d_l$ est une solution
de $f_j = 0$ congrue à $a_i$ modulo $(\mathfrak m')^N$.
Il suffit donc de montrer que l'existence d'une solution sur $R^\wedge$
entraîne l'existence d'une solution sur une extension d'anneaux \'etale
qui induit une extension triviale des corps résiduels en un idéal premier
au-dessus de $\mathfrak m$.

\medskip\noindent
Soit $A \subset R^\wedge$ la sous-$R$-algèbre engendrée par
$a_1, \ldots, a_n$. Puisque nous avons supposé que $R$ est un G-anneau,
c'est-à-dire que $R \to R^\wedge$ est régulier, il existe une factorisation
$$
A \to B \to R^\wedge
$$
où $B$ est lisse sur $R$ ; voir le théorème \ref{smoothing-theorem-popescu}.
Notons $\kappa = R/\mathfrak m$ le corps résiduel. C'est aussi
le corps résiduel de $R^\wedge$, si bien que nous obtenons un diagramme
commutatif
$$
\xymatrix{
B \ar[rd] \ar@{..>}[r] & R' \ar@{..>}[d] \\
R \ar[r] \ar[u] & \kappa
}
$$
Comme la flèche verticale est lisse,
Compléments d'algèbre, lemme
\ref{more-algebra-lemma-lift-section-smooth-morphism},
il existe un morphisme d'anneaux \'etale $R \to R'$ qui induit un isomorphisme
$R/\mathfrak m \to R'/\mathfrak mR'$ et un morphisme de $R$-algèbres
$B \to R'$ rendant commutatif le diagramme ci-dessus.
Soit $b_i \in R'$ l'image de $a_i$ par les morphismes d'anneaux
$A \to B \to R'$. Puisque tous ces morphismes sont des morphismes de
$R$-algèbres, $(b_1, \ldots, b_n)$ est une solution dans $R'$.
\end{proof}

\begin{example}
\label{smoothing-example-describe-henselian}
Soit $(R, \mathfrak m)$ un anneau local noethérien d'hensélisation $R^h$.
Le morphisme entre les complétés $R^\wedge \to (R^h)^\wedge$ est un
isomorphisme ; voir Compléments d'algèbre, lemme
\ref{more-algebra-lemma-henselization-noetherian}.
Comme $R^h$ est lui aussi noethérien (ibid.), nous pouvons considérer $R^h$
comme un sous-anneau de son complété (car la complétion est fidèlement
plate). Nous pouvons ainsi identifier $R^h$ à un sous-anneau de $R^\wedge$.

\medskip\noindent
Cherchons à comprendre quels sont les éléments de $R^\wedge$ qui
appartiennent à $R^h$. Pour simplifier, supposons que $R$ soit un domaine
ayant $K$ pour corps des fractions. Il est clair que tout élément $f$ de $R^h$
est algébrique sur $R$, au sens où il existe une équation de la forme
$a_n f^n + \ldots + a_1 f + a_0 = 0$, pour certains $a_i \in R$ tels que
$n > 0$ et $a_n \not = 0$.

\medskip\noindent
Réciproquement, soient $f \in R^\wedge$, $n \in \mathbf{N}$ et
$a_0, \ldots, a_n \in R$, avec $a_n \not = 0$, tels que
$a_n f^n + \ldots + a_1 f + a_0 = 0$. Si $R$ est un G-anneau, alors, pour
tout $N > 0$, il existe un élément $g \in R^h$ tel que
$a_n g^n + \ldots + a_1 g + a_0 = 0$ et $f - g \in \mathfrak m^N R^\wedge$,
voir le théorème \ref{smoothing-theorem-approximation-property-variant}.
Nous voudrions en déduire que $f = g$ lorsque $N \gg 0$.
Si ce n'est pas le cas, nous obtenons une infinité de racines $g$ de $P(T)$
dans $R^h$. C'est impossible, car (1) $R^h \subset R^h \otimes_R K$
et (2) $R^h \otimes_R K$ est un produit fini d'extensions de corps de $K$.
En effet, $R \to K$ est injectif et $R \to R^h$ est plat ; par conséquent,
$R^h \to R^h \otimes_R K$ est injectif, et (2) résulte de
Compléments d'algèbre, lemme
\ref{more-algebra-lemma-fibres-henselization}.

\medskip\noindent
Conclusion : si $R$ est un domaine local noethérien ayant $K$ pour corps des
fractions et un G-anneau, alors $R^h \subset R^\wedge$ est l'ensemble des
éléments algébriques sur $K$.
\end{example}

\noindent
Voici une autre variante du théorème principal de cette section.

\begin{lemma}
\label{smoothing-lemma-approximation-property-variant}
Soit $R$ un anneau noethérien. Soit $\mathfrak p \subset R$ un idéal
premier. Soient
$f_1, \ldots, f_m \in R[x_1, \ldots, x_n]$.
Supposons que $(a_1, \ldots, a_n) \in ((R_\mathfrak p)^\wedge)^n$ soit une
solution. Si $R_\mathfrak p$ est un G-anneau, alors, pour tout entier $N$,
il existe
\begin{enumerate}
\item un morphisme d'anneaux \'etale $R \to R'$,
\item un idéal premier $\mathfrak p' \subset R'$ au-dessus de $\mathfrak p$,
\item une solution $(b_1, \ldots, b_n) \in (R')^n$ dans $R'$,
\end{enumerate}
tels que $\kappa(\mathfrak p) = \kappa(\mathfrak p')$ et
$a_i - b_i \in (\mathfrak p')^N(R'_{\mathfrak p'})^\wedge$.
\end{lemma}

\begin{proof}
D'après le théorème \ref{smoothing-theorem-approximation-property-variant},
nous pouvons trouver une solution $(b'_1, \ldots, b'_n)$ dans un anneau
$R''$ \'etale sur $R_\mathfrak p$, muni d'un idéal premier
$\mathfrak p''$ au-dessus de $\mathfrak p$, tel que
$\kappa(\mathfrak p) = \kappa(\mathfrak p'')$ et
$a_i - b'_i \in (\mathfrak p'')^N(R''_{\mathfrak p''})^\wedge$.
Nous pouvons écrire
$R'' = R' \otimes_R R_\mathfrak p$
pour une $R$-algèbre \'etale $R'$ (voir Algèbre, lemme
\ref{algebra-lemma-etale}). Après avoir remplacé $R'$ par une localisation
principale si nécessaire, nous pouvons supposer que
$(b'_1, \ldots, b'_n)$ provient d'une solution $(b_1, \ldots, b_n)$
dans $R'$. Posons $\mathfrak p' = R' \cap \mathfrak p''$ ; alors
$R''_{\mathfrak p''} = R'_{\mathfrak p'}$, ce qui achève la démonstration.
\end{proof}



\section{Approximation pour les couples henséliens}
\label{smoothing-section-approximation-pairs}

\noindent
Nous pouvons généraliser la discussion de la
section \ref{smoothing-section-approximation-G-rings} au cas des couples henséliens.
Les couples henséliens ont été définis dans
Compléments d'algèbre, section
\ref{more-algebra-section-henselian-pairs}.

\begin{lemma}
\label{smoothing-lemma-henselian-pair}
\begin{slogan}
Approximation pour les couples henséliens.
\end{slogan}
Soit $(A, I)$ un couple hensélien, avec $A$ noethérien.
Soit $A^\wedge$ le complété $I$-adique de $A$. Supposons que l'une au
moins des conditions suivantes soit satisfaite :
\begin{enumerate}
\item $A \to A^\wedge$ est un morphisme d'anneaux régulier ;
\item $A$ est un G-anneau noethérien ;
\item $(A, I)$ est l'hensélisation
(Compléments d'algèbre, lemme \ref{more-algebra-lemma-henselization})
d'un couple $(B, J)$, où $B$ est un G-anneau noethérien.
\end{enumerate}
Soient $f_1, \ldots, f_m \in A[x_1, \ldots, x_n]$
et $\hat{a}_1, \ldots, \hat{a}_n \in A^\wedge$ tels que
$f_j(\hat{a}_1, \ldots, \hat{a}_n) = 0$
pour $j = 1, \ldots, m$. Alors, pour tout $N \geq 1$, il existe
$a_1, \ldots, a_n \in A$ tels que
$\hat{a}_i - a_i \in I^N$ et $f_j(a_1, \ldots, a_n) = 0$
pour $j = 1, \ldots, m$.
\end{lemma}

\begin{proof}
D'après Compléments d'algèbre, lemme
\ref{more-algebra-lemma-henselization-pair-G-ring}
nous voyons que (3) implique (2). D'après Compléments d'algèbre, lemme
\ref{more-algebra-lemma-map-G-ring-to-completion-regular}
nous voyons que (2) implique (1). Il suffit donc de démontrer le lemme
dans le cas où $A \to A^\wedge$ est un morphisme d'anneaux régulier.

\medskip\noindent
Soient $\hat{a}_1, \ldots, \hat{a}_n$ comme dans l'énoncé du lemme.
D'après le théorème \ref{smoothing-theorem-popescu}, nous pouvons trouver une
factorisation $A \to B \to A^\wedge$ telle que $A \to B$ soit lisse, ainsi que
des éléments $b_1, \ldots, b_n \in B$ tels que
$f_j(b_1, \ldots, b_n) = 0$ dans $B$.
Notons $\sigma : B \to A^\wedge \to A/I^N$ la composée.
D'après Compléments d'algèbre, lemme
\ref{more-algebra-lemma-lift-section-smooth-morphism}, il existe un morphisme
d'anneaux \'etale $A \to A'$ qui induit un isomorphisme
$A/I^N \to A'/I^NA'$ et un morphisme de $A$-algèbres
$\tilde \sigma : B \to A'$ relevant $\sigma$.
Comme $(A, I)$ est hensélien, il existe un morphisme de $A$-algèbres
$\chi : A' \to A$ ; voir Compléments d'algèbre, lemme
\ref{more-algebra-lemma-characterize-henselian-pair}.
En posant alors $a_i = \chi(\tilde \sigma(b_i))$, nous obtenons une solution.
\end{proof}







% OK, start here.
%

\chapter{Faisceaux de modules}



\phantomsection
\label{modules-section-phantom}


\section{Introduction}
\label{modules-section-introduction}

\noindent
Dans ce chapitre, nous développons les notions de base relatives aux faisceaux
de modules. Cela comprend en particulier le cas des faisceaux abéliens, puisque
ceux-ci peuvent être considérés comme des faisceaux de
$\underline{\mathbf{Z}}$-modules. Les références de base sont
\cite{FAC}, \cite{EGA} et \cite{SGA4}.

\medskip\noindent
Nous étudions dans un autre chapitre le cas des faisceaux de modules sur les
topos annelés (voir Modules sur les sites, section
\ref{sites-modules-section-introduction}), bien que nous y reprenions
essentiellement les considérations du présent chapitre.





\section{Pathologie}
\label{modules-section-pathology}

\noindent
Un espace annelé est un couple formé d'un espace topologique $X$ et d'un
faisceau d'anneaux $\mathcal{O}$. Nous autorisons $\mathcal{O} = 0$ dans
la définition. Dans ce cas, la catégorie des modules possède un seul objet
(à savoir $0$). Elle demeure une catégorie abélienne, etc., mais elle est
quelque peu dégénérée. De même, le faisceau $\mathcal{O}$ peut être nul
sur certains ouverts de $X$, etc.

\medskip\noindent
Ce phénomène ne se produit pas pour les espaces localement annelés (que nous
considérerons plus loin).








\section{La catégorie abélienne des faisceaux de modules}
\label{modules-section-kernels}

\noindent
Soit $(X, \mathcal{O}_X)$ un espace annelé, voir Faisceaux, Définition
\ref{sheaves-definition-ringed-space}. Soient $\mathcal{F}$ et $\mathcal{G}$
des faisceaux de $\mathcal{O}_X$-modules, voir Faisceaux, Définition
\ref{sheaves-definition-sheaf-modules}. Soient
$\varphi, \psi : \mathcal{F} \to \mathcal{G}$ des morphismes de faisceaux
de $\mathcal{O}_X$-modules. Nous définissons
$\varphi + \psi : \mathcal{F} \to \mathcal{G}$ comme l'application qui,
sur chaque ouvert $U \subset X$, est la somme des applications induites par
$\varphi$ et $\psi$. Il s'agit manifestement encore d'une application de
faisceaux de $\mathcal{O}_X$-modules. Il est également clair que la composition
des applications de $\mathcal{O}_X$-modules est bilinéaire pour cette addition.
Ainsi, $\textit{Mod}(\mathcal{O}_X)$ est une catégorie préadditive, voir
Homologie, Définition \ref{homology-definition-preadditive}.

\medskip\noindent
Nous noterons $0$ le faisceau de $\mathcal{O}_X$-modules qui prend la valeur
constante $\{0\}$ sur tout ouvert $U \subset X$. Il s'agit manifestement à la
fois d'un objet final et d'un objet initial de
$\textit{Mod}(\mathcal{O}_X)$. Étant donné un morphisme de
$\mathcal{O}_X$-modules $\varphi : \mathcal{F} \to \mathcal{G}$, les
assertions suivantes sont équivalentes :
(a) $\varphi$ est nul, (b) $\varphi$ se factorise par $0$,
(c) $\varphi$ est nul sur les sections au-dessus de tout ouvert $U$, et
(d) $\varphi_x = 0$ pour tout $x \in X$. Voir Faisceaux, Lemme
\ref{sheaves-lemma-points-exactness}.

\medskip\noindent
En outre, étant donnés deux faisceaux $\mathcal{F}$, $\mathcal{G}$ de
$\mathcal{O}_X$-modules, nous pouvons définir leur somme directe par
$$
\mathcal{F} \oplus \mathcal{G} = \mathcal{F} \times \mathcal{G}
$$
munie des applications évidentes $(i, j, p, q)$ comme dans Homologie,
Définition \ref{homology-definition-direct-sum}. Ainsi
$\textit{Mod}(\mathcal{O}_X)$ est une catégorie additive, voir Homologie,
Définition \ref{homology-definition-additive-category}.

\medskip\noindent
Soit $\varphi : \mathcal{F} \to \mathcal{G}$ un morphisme de
$\mathcal{O}_X$-modules. Nous pouvons définir $\Ker(\varphi)$ comme le
sous-faisceau de $\mathcal{F}$ dont les sections sont
$$
\Ker(\varphi)(U) =
\{ s \in \mathcal{F}(U) \mid \varphi(s) = 0 \text{ dans } \mathcal{G}(U)\}
$$
pour tout ouvert $U \subset X$. On vérifie aisément qu'il s'agit bien d'un
noyau dans la catégorie des $\mathcal{O}_X$-modules. Autrement dit, un
morphisme $\alpha : \mathcal{H} \to \mathcal{F}$ se factorise par
$\Ker(\varphi)$ si et seulement si $\varphi \circ \alpha = 0$. De plus, au
niveau des germes, nous avons $\Ker(\varphi)_x = \Ker(\varphi_x)$.

\medskip\noindent
D'autre part, nous définissons $\Coker(\varphi)$ comme le faisceau de
$\mathcal{O}_X$-modules associé au préfaisceau de
$\mathcal{O}_X$-modules défini par la règle
$$
U
\longmapsto
\Coker(\mathcal{F}(U)\to \mathcal{G}(U)) =
\mathcal{G}(U)/\varphi(\mathcal{F}(U)).
$$
Comme la prise des germes commute au passage au faisceau associé, voir
Faisceaux, Lemme \ref{sheaves-lemma-stalk-sheafification}, nous obtenons
$\Coker(\varphi)_x = \Coker(\varphi_x)$. L'application
$\mathcal{G} \to \Coker(\varphi)$ est donc surjective (comme application de
faisceaux d'ensembles), voir Faisceaux, section
\ref{sheaves-section-exactness-points}. Pour montrer qu'il s'agit d'un
conoyau, remarquons que si
$\beta : \mathcal{G} \to \mathcal{H}$ est un morphisme de
$\mathcal{O}_X$-modules tel que $\beta \circ \varphi$ soit nul, alors, pour
tout ouvert $U \subset X$, $\beta$ induit une application de
$\mathcal{G}(U)/\varphi(\mathcal{F}(U))$ dans $\mathcal{H}(U)$. Par la
propriété universelle du faisceau associé (voir Faisceaux, Lemme
\ref{sheaves-lemma-sheafification-presheaf-modules}), nous obtenons une
application canonique $\Coker(\varphi) \to \mathcal{H}$ telle que
l'application $\beta$ initiale soit égale à la composée
$\mathcal{G} \to \Coker(\varphi) \to \mathcal{H}$. Le morphisme
$\Coker(\varphi) \to \mathcal{H}$ est unique en vertu de la surjectivité
mentionnée ci-dessus.

\begin{lemma}
\label{modules-lemma-abelian}
Soit $(X, \mathcal{O}_X)$ un espace annelé. La catégorie
$\textit{Mod}(\mathcal{O}_X)$ est une catégorie abélienne. De plus, un complexe
$$
\mathcal{F} \to \mathcal{G} \to \mathcal{H}
$$
est exact en $\mathcal{G}$ si et seulement si, pour tout $x \in X$, le complexe
$$
\mathcal{F}_x \to \mathcal{G}_x \to \mathcal{H}_x
$$
est exact en $\mathcal{G}_x$.
\end{lemma}

\begin{proof}
D'après Homologie, Définition \ref{homology-definition-abelian-category},
il faut montrer que l'image et la coimage coïncident. D'après Faisceaux,
Lemme \ref{sheaves-lemma-points-exactness}, il suffit de montrer qu'elles ont
le même germe en tout $x \in X$. Par les constructions des noyaux et conoyaux
ci-dessus, ces germes sont respectivement la coimage et l'image dans la
catégorie des $\mathcal{O}_{X, x}$-modules. Le résultat découle donc du fait
que la catégorie des modules sur un anneau est abélienne.
\end{proof}

\noindent
En fait, la catégorie $\textit{Mod}(\mathcal{O}_X)$ possède bien d'autres
propriétés. Voici deux constructions que nous pouvons effectuer.
\begin{enumerate}
\item Étant donnés un ensemble $I$ et, pour chaque $i \in I$, un
$\mathcal{O}_X$-module, nous pouvons former le produit
$$
\prod\nolimits_{i \in I} \mathcal{F}_i
$$
qui est le faisceau associant à chaque ouvert $U$ le produit des modules
$\mathcal{F}_i(U)$. C'est aussi le produit catégorique, comme dans Catégories,
Définition \ref{categories-definition-product}.
\item Étant donnés un ensemble $I$ et, pour chaque $i \in I$, un
$\mathcal{O}_X$-module, nous pouvons former la somme directe
$$
\bigoplus\nolimits_{i \in I} \mathcal{F}_i
$$
qui est le {\it faisceau associé} au préfaisceau associant à chaque ouvert
$U$ la somme directe des modules $\mathcal{F}_i(U)$. C'est aussi le coproduit
catégorique, comme dans Catégories, Définition
\ref{categories-definition-coproduct}. Cela résulte de la propriété
universelle du faisceau associé.
\end{enumerate}
Nous en déduisons que toutes les limites projectives et inductives existent
dans $\textit{Mod}(\mathcal{O}_X)$.

\begin{lemma}
\label{modules-lemma-limits-colimits}
Soit $(X, \mathcal{O}_X)$ un espace annelé.
\begin{enumerate}
\item Toutes les limites projectives existent dans
$\textit{Mod}(\mathcal{O}_X)$. Elles coïncident avec les limites projectives
correspondantes de préfaisceaux de $\mathcal{O}_X$-modules (autrement dit,
elles commutent à la prise des sections sur les ouverts).
\item Toutes les limites inductives existent dans
$\textit{Mod}(\mathcal{O}_X)$. Elles s'obtiennent en prenant le faisceau
associé à la limite inductive correspondante dans la catégorie des
préfaisceaux. La prise des limites inductives commute à la prise des germes.
\item Les limites inductives filtrantes sont exactes.
\item Les sommes directes finies coïncident avec les sommes directes finies
correspondantes de préfaisceaux de $\mathcal{O}_X$-modules.
\end{enumerate}
\end{lemma}

\begin{proof}
Comme $\textit{Mod}(\mathcal{O}_X)$ est abélienne (Lemme
\ref{modules-lemma-abelian}), elle admet toutes les limites projectives et inductives
finies (Homologie, Lemme \ref{homology-lemma-colimit-abelian-category}).
L'existence et la description des limites projectives et inductives découlent
donc de celles des produits et coproduits (voir la discussion ci-dessus) et
des lemmes de Catégories \ref{categories-lemma-limits-products-equalizers} et
\ref{categories-lemma-colimits-coproducts-coequalizers}. Puisque le passage au
faisceau associé commute à la prise des germes, les limites inductives
commutent à la prise des germes. L'assertion (3) signifie que, pour un système
$0 \to \mathcal{F}_i \to \mathcal{G}_i \to \mathcal{H}_i \to 0$ de suites
exactes de $\mathcal{O}_X$-modules indexé par un ensemble filtrant $I$, la
suite $0 \to \colim \mathcal{F}_i \to \colim \mathcal{G}_i \to
\colim \mathcal{H}_i \to 0$ est encore exacte. Comme l'exactitude se vérifie
sur les germes (Lemme \ref{modules-lemma-abelian}), cela découle du cas des modules,
qui est Algèbre, Lemme \ref{algebra-lemma-directed-colimit-exact}. Nous
omettons la démonstration de (4).
\end{proof}

\noindent
L'existence des limites projectives et inductives permet d'exprimer les
propriétés d'exactitude des foncteurs définis sur la catégorie des
$\mathcal{O}$-modules en termes de ces limites, comme dans Catégories,
section \ref{categories-section-exact-functor}. Voir Homologie, Lemme
\ref{homology-lemma-exact-functor}, pour une description des propriétés
d'exactitude en termes de suites exactes courtes.

\begin{lemma}
\label{modules-lemma-exactness-pushforward-pullback}
Soit $f : (X, \mathcal{O}_X) \to (Y, \mathcal{O}_Y)$ un morphisme
d'espaces annelés.
\begin{enumerate}
\item Le foncteur
$f_* : \textit{Mod}(\mathcal{O}_X) \to \textit{Mod}(\mathcal{O}_Y)$
est exact à gauche. En fait, il commute à toutes les limites projectives.
\item Le foncteur
$f^* : \textit{Mod}(\mathcal{O}_Y) \to \textit{Mod}(\mathcal{O}_X)$
est exact à droite. En fait, il commute à toutes les limites inductives.
\item Le foncteur image inverse
$f^{-1} : \textit{Ab}(Y) \to \textit{Ab}(X)$ sur les faisceaux abéliens est
exact.
\end{enumerate}
\end{lemma}

\begin{proof}
Les assertions (1) et (2) résultent de ce que $(f^*, f_*)$ est un couple de
foncteurs adjoints, voir Faisceaux, Lemme
\ref{sheaves-lemma-adjoint-pullback-pushforward-modules}, et Catégories,
section \ref{categories-section-adjoint}. L'assertion (3) résulte de ce que
l'exactitude se vérifie sur les germes (Lemme \ref{modules-lemma-abelian}) et de la
description des germes de l'image inverse, voir Faisceaux, Lemme
\ref{sheaves-lemma-pullback-abelian-stalk}.
\end{proof}

\begin{lemma}
\label{modules-lemma-j-shriek-exact}
Soit $j : U \to X$ une immersion ouverte d'espaces topologiques. Le foncteur
$j_! : \textit{Ab}(U) \to \textit{Ab}(X)$ est exact.
\end{lemma}

\begin{proof}
Cela découle de la description des germes donnée dans Faisceaux, Lemme
\ref{sheaves-lemma-j-shriek-abelian}.
\end{proof}

\begin{lemma}
\label{modules-lemma-section-direct-sum-quasi-compact}
Soit $(X, \mathcal{O}_X)$ un espace annelé. Soit $I$ un ensemble. Pour
$i \in I$, soit $\mathcal{F}_i$ un faisceau de $\mathcal{O}_X$-modules.
Pour tout ouvert quasi-compact $U \subset X$, l'application
$$
\bigoplus\nolimits_{i \in I} \mathcal{F}_i(U)
\longrightarrow
\left(\bigoplus\nolimits_{i \in I} \mathcal{F}_i\right)(U)
$$
est bijective.
\end{lemma}

\begin{proof}
Si $s$ est un élément du membre de droite, il existe un recouvrement ouvert
$U = \bigcup_{j \in J} U_j$ tel que $s|_{U_j}$ soit une somme finie
$\sum_{i \in I_j} s_{ji}$, où $s_{ji} \in \mathcal{F}_i(U_j)$. Comme $U$
est quasi-compact, nous pouvons supposer le recouvrement fini, c'est-à-dire
$J$ fini. Alors $I' = \bigcup_{j \in J} I_j$ est une partie finie de $I$.
Manifestement, $s$ est une section du sous-faisceau
$\bigoplus_{i \in I'} \mathcal{F}_i$. Le résultat découle du fait que, pour
une somme directe finie, le passage au faisceau associé n'est pas nécessaire ;
voir le Lemme \ref{modules-lemma-limits-colimits} ci-dessus.
\end{proof}



\section{sections des faisceaux de modules}
\label{modules-section-sections}

\noindent
Soit $(X, \mathcal{O}_X)$ un espace annelé. Soit $\mathcal{F}$ un faisceau de
$\mathcal{O}_X$-modules. Soit
$s \in \Gamma(X, \mathcal{F}) = \mathcal{F}(X)$ une section globale. Il existe
une unique {\it application de $\mathcal{O}_X$-modules}
$$
\mathcal{O}_X \longrightarrow \mathcal{F}, \ f \longmapsto fs
$$
{\it associée à $s$}. La notation ci-dessus signifie qu'une section locale
$f$ de $\mathcal{O}_X$, c'est-à-dire une section $f$ sur un ouvert $U$, est
envoyée sur le produit de $f$ par la restriction de $s$ à $U$.
Réciproquement, toute application
$\varphi : \mathcal{O}_X \to \mathcal{F}$ donne une section
$s = \varphi(1)$ telle que $\varphi$ soit le morphisme associé à $s$.

\begin{definition}
\label{modules-definition-globally-generated}
Soit $(X, \mathcal{O}_X)$ un espace annelé. Soit $\mathcal{F}$ un faisceau de
$\mathcal{O}_X$-modules. Nous disons que $\mathcal{F}$ est {\it engendré par
ses sections globales} s'il existe un ensemble $I$ et des sections globales
$s_i \in \Gamma(X, \mathcal{F})$, $i \in I$, tels que l'application
$$
\bigoplus\nolimits_{i \in I}
\mathcal{O}_X \longrightarrow \mathcal{F}
$$
qui est l'application associée à $s_i$ sur le facteur correspondant à $i$,
soit surjective. Dans ce cas, nous disons que les sections $s_i$
{\it engendrent} $\mathcal{F}$.
\end{definition}

\noindent
Nous emploierons souvent l'abus de notation introduit dans Faisceaux,
section \ref{sheaves-section-stalks} : si $s$ est une section locale de
$\mathcal{F}$ définie dans un voisinage ouvert d'un point $x \in X$, nous
notons $s_x$, voire $s$, l'image de $s$ dans le germe $\mathcal{F}_x$.

\begin{lemma}
\label{modules-lemma-globally-generated}
Soit $(X, \mathcal{O}_X)$ un espace annelé. Soit $\mathcal{F}$ un faisceau de
$\mathcal{O}_X$-modules. Soit $I$ un ensemble. Soient
$s_i \in \Gamma(X, \mathcal{F})$, $i \in I$, des sections globales. Les
sections $s_i$ engendrent $\mathcal{F}$ si et seulement si, pour tout
$x\in X$, les éléments $s_{i, x} \in \mathcal{F}_x$ engendrent le
$\mathcal{O}_{X, x}$-module $\mathcal{F}_x$.
\end{lemma}

\begin{proof}
Omis.
\end{proof}

\begin{lemma}
\label{modules-lemma-tensor-product-globally-generated}
\begin{slogan}
Le produit tensoriel de faisceaux de modules engendrés par leurs sections
globales est engendré par ses sections globales.
\end{slogan}
Soit $(X, \mathcal{O}_X)$ un espace annelé. Soient $\mathcal{F}$ et
$\mathcal{G}$ des faisceaux de $\mathcal{O}_X$-modules. Si $\mathcal{F}$ et
$\mathcal{G}$ sont engendrés par leurs sections globales, il en est de même de
$\mathcal{F} \otimes_{\mathcal{O}_X} \mathcal{G}$.
\end{lemma}

\begin{proof}
Omis.
\end{proof}

\begin{lemma}
\label{modules-lemma-generated-by-local-sections}
Soit $(X, \mathcal{O}_X)$ un espace annelé. Soit $\mathcal{F}$ un faisceau de
$\mathcal{O}_X$-modules. Soit $I$ un ensemble. Soient $s_i$, $i \in I$, une
famille de sections locales de $\mathcal{F}$, c'est-à-dire
$s_i \in \mathcal{F}(U_i)$ pour certains ouverts $U_i \subset X$. Il existe
un unique plus petit sous-faisceau de $\mathcal{O}_X$-modules $\mathcal{G}$
tel que chaque $s_i$ corresponde à une section locale de $\mathcal{G}$.
\end{lemma}

\begin{proof}
Considérons le sous-préfaisceau de $\mathcal{O}_X$-modules défini par la règle
$$
U
\longmapsto
\{
\text{sommes } \sum\nolimits_{i \in J} f_i (s_i|_U)
\text{ où } J \text{ est fini, }
U \subset U_i \text{ pour } i\in J, \text{ et }
f_i \in \mathcal{O}_X(U)
\}
$$
Soit $\mathcal{G}$ le faisceau associé à ce sous-préfaisceau. C'est un
sous-faisceau de $\mathcal{F}$ d'après Faisceaux, Lemme
\ref{sheaves-lemma-sheafify-epi-mono}. Puisque toutes les sommes finies
doivent manifestement appartenir à $\mathcal{G}$, c'est bien le plus petit
sous-faisceau voulu.
\end{proof}

\begin{definition}
\label{modules-definition-generated-by-local-sections}
Soit $(X, \mathcal{O}_X)$ un espace annelé. Soit $\mathcal{F}$ un faisceau de
$\mathcal{O}_X$-modules. Étant donnés un ensemble $I$ et des sections locales
$s_i$, $i \in I$, de $\mathcal{F}$, nous disons que le sous-faisceau
$\mathcal{G}$ du Lemme \ref{modules-lemma-generated-by-local-sections} ci-dessus est
le {\it sous-faisceau engendré par les $s_i$}.
\end{definition}

\begin{lemma}
\label{modules-lemma-generated-by-local-sections-stalk}
Soit $(X, \mathcal{O}_X)$ un espace annelé. Soit $\mathcal{F}$ un faisceau de
$\mathcal{O}_X$-modules. Soient un ensemble $I$ et des sections locales
$s_i$, $i \in I$, de $\mathcal{F}$. Soit $\mathcal{G}$ le sous-faisceau
engendré par les $s_i$, et soit $x\in X$. Alors $\mathcal{G}_x$ est le
sous-$\mathcal{O}_{X, x}$-module de $\mathcal{F}_x$ engendré par les éléments
$s_{i, x}$ pour les indices $i$ tels que $s_i$ soit définie en $x$.
\end{lemma}

\begin{proof}
Cela résulte immédiatement de la construction de $\mathcal{G}$ dans la
démonstration du Lemme \ref{modules-lemma-generated-by-local-sections}.
\end{proof}










\section{Supports des modules et des sections}
\label{modules-section-support}

\begin{definition}
\label{modules-definition-support}
Soit $(X, \mathcal{O}_X)$ un espace annelé. Soit $\mathcal{F}$ un faisceau de
$\mathcal{O}_X$-modules.
\begin{enumerate}
\item Le {\it support de $\mathcal{F}$} est l'ensemble des points
$x \in X$ tels que $\mathcal{F}_x \not = 0$.
\item Nous notons $\text{Supp}(\mathcal{F})$ le support de $\mathcal{F}$.
\item Soit $s \in \Gamma(X, \mathcal{F})$ une section globale. Le
{\it support de $s$} est l'ensemble des points $x \in X$ tels que l'image
$s_x \in \mathcal{F}_x$ de $s$ soit non nulle.
\end{enumerate}
\end{definition}

\noindent
Le support d'une section locale se trouve bien entendu ainsi défini, puisqu'une
section locale est une section globale de la restriction de $\mathcal{F}$.

\begin{lemma}
\label{modules-lemma-support-section-closed}
Soit $(X, \mathcal{O}_X)$ un espace annelé. Soit $\mathcal{F}$ un faisceau de
$\mathcal{O}_X$-modules. Soit $U \subset X$ un ouvert.
\begin{enumerate}
\item Le support de $s \in \mathcal{F}(U)$ est fermé dans $U$.
\item Le support de $fs$ est contenu dans l'intersection des supports de
$f \in \mathcal{O}_X(U)$ et de $s \in \mathcal{F}(U)$.
\item Le support de $s + s'$ est contenu dans la réunion des supports de
$s, s' \in \mathcal{F}(U)$.
\item Le support de $\mathcal{F}$ est la réunion des supports de toutes les
sections locales de $\mathcal{F}$.
\item Si $\varphi : \mathcal{F} \to \mathcal{G}$ est un morphisme de
$\mathcal{O}_X$-modules, alors le support de $\varphi(s)$ est contenu dans
celui de $s \in \mathcal{F}(U)$.
\end{enumerate}
\end{lemma}

\begin{proof}
En effet, si $s_x = 0$, alors $s$ est nul dans un voisinage ouvert de $x$ par
définition des germes. Il en va de même pour $f$. Nous omettons les détails.
\end{proof}

\noindent
En général, le support d'un faisceau de modules n'est pas fermé. En effet, on
peut prendre sur $\mathbf{R}$ (muni de sa topologie archimédienne usuelle) le
faisceau abélien somme directe d'une infinité de faisceaux gratte-ciel non nuls,
chacun supporté en un seul point $p_i$ de $\mathbf{R}$. Son support est alors
l'ensemble des points $p_i$, qui peut ne pas être fermé.

\medskip\noindent
Un autre exemple s'obtient en considérant l'immersion ouverte
$j : U = (0 , \infty) \to \mathbf{R} = X$ et le faisceau abélien
$j_!\underline{\mathbf{Z}}_U$. D'après Faisceaux, section
\ref{sheaves-section-open-immersions}, le support de ce faisceau est exactement
$U$.

\begin{lemma}
\label{modules-lemma-support-sheaf-rings-closed}
Soit $X$ un espace topologique. Le support d'un faisceau d'anneaux est fermé.
\end{lemma}

\begin{proof}
Cela résulte de ce que, selon nos conventions, un anneau est $0$ si et seulement
si $1 = 0$ ; le support d'un faisceau d'anneaux est donc le support de la
section unité.
\end{proof}






\section{Immersions fermées et faisceaux abéliens}
\label{modules-section-closed-immersions}

\noindent
Rappelons que nous considérons un faisceau abélien sur un espace topologique
$X$ comme un faisceau de $\underline{\mathbf{Z}}_X$-modules. Nous pouvons donc
appliquer aux faisceaux abéliens tous les résultats et définitions concernant
les faisceaux de modules.

\begin{lemma}
\label{modules-lemma-i-star-exact}
Soit $X$ un espace topologique et soit $Z \subset X$ une partie fermée.
Notons $i : Z \to X$ l'inclusion. Le foncteur
$$
i_* : \textit{Ab}(Z) \longrightarrow \textit{Ab}(X)
$$
est exact et pleinement fidèle ; son image essentielle est exactement
constituée des faisceaux abéliens dont le support est contenu dans $Z$. Le
foncteur $i^{-1}$ est un inverse à gauche de $i_*$.
\end{lemma}

\begin{proof}
L'exactitude résulte de la description des germes dans Faisceaux, Lemme
\ref{sheaves-lemma-stalks-closed-pushforward}, et du Lemme
\ref{modules-lemma-abelian}. Le reste a été démontré dans Faisceaux, Lemme
\ref{sheaves-lemma-equivalence-categories-closed-abelian}.
\end{proof}

\noindent
Soit $\mathcal{F}$ un faisceau abélien sur l'espace topologique $X$. Étant
donnée une partie fermée $Z$, il existe un sous-faisceau abélien canonique de
$\mathcal{F}$ formé exactement des sections dont le support est contenu dans
$Z$. Voici l'énoncé précis.

\begin{remark}
\label{modules-remark-sections-support-in-closed}
Soit $X$ un espace topologique, soit $Z \subset X$ une partie fermée et soit
$\mathcal{F}$ un faisceau abélien sur $X$. Pour tout ouvert $U \subset X$,
posons
$$
\mathcal{H}_Z(\mathcal{F})(U) =
\{s \in \mathcal{F}(U) \mid
\text{ le support de }s\text{ est contenu dans }Z \cap U\}.
$$
Alors $\mathcal{H}_Z(\mathcal{F})$ est un sous-faisceau abélien de
$\mathcal{F}$. C'est le plus grand sous-faisceau abélien de $\mathcal{F}$ dont
le support est contenu dans $Z$. D'après le Lemme \ref{modules-lemma-i-star-exact},
nous pouvons (et allons) considérer $\mathcal{H}_Z(\mathcal{F})$ comme un
faisceau abélien sur $Z$. Nous obtenons ainsi un foncteur exact à gauche
$$
\textit{Ab}(X) \longrightarrow \textit{Ab}(Z),\quad
\mathcal{F} \longmapsto \mathcal{H}_Z(\mathcal{F})
\text{ considéré comme faisceau abélien sur }Z
$$
Toutes les assertions précédentes découlent directement du Lemme
\ref{modules-lemma-support-section-closed}.
\end{remark}

\noindent
C'est l'occasion de montrer que le foncteur $i_*$ possède un adjoint à droite
sur les faisceaux abéliens.

\begin{lemma}
\label{modules-lemma-i-star-right-adjoint}
Soit $i : Z \to X$ l'inclusion d'une partie fermée dans l'espace topologique
$X$. Le foncteur $\textit{Ab}(X) \to \textit{Ab}(Z)$,
$\mathcal{F} \mapsto \mathcal{H}_Z(\mathcal{F})$ de la Remarque
\ref{modules-remark-sections-support-in-closed} est adjoint à droite de
$i_* : \textit{Ab}(Z) \to \textit{Ab}(X)$. En particulier, $i_*$ commute aux
limites inductives arbitraires.
\end{lemma}

\begin{proof}
Il faut montrer que, pour tout faisceau abélien $\mathcal{F}$ sur $X$ et tout
faisceau abélien $\mathcal{G}$ sur $Z$, nous avons
$$
\Hom_{\textit{Ab}(X)}(i_*\mathcal{G}, \mathcal{F}) =
\Hom_{\textit{Ab}(Z)}(\mathcal{G}, \mathcal{H}_Z(\mathcal{F}))
$$
C'est clair puisque toute section de $i_*\mathcal{G}$ a son support dans $Z$.
Nous omettons les détails.
\end{proof}

\begin{remark}
\label{modules-remark-i-star-right-adjoint}
Dans Faisceaux, Remarque \ref{sheaves-remark-i-star-not-exact}, nous avons
montré que $i_*$, comme foncteur entre catégories de faisceaux d'ensembles,
n'admet pas d'adjoint à droite, simplement parce qu'il n'est pas exact.
Toutefois, cette assertion est presque vraie : le foncteur $i_*$ est en effet
exact sur les faisceaux d'ensembles pointés, on peut définir pour ceux-ci les
sections à support dans $Z$, et $\mathcal{H}_Z$ a un sens et est adjoint à
droite de $i_*$.
\end{remark}









\section{Une suite exacte canonique}
\label{modules-section-canonical-exact-sequence}

\noindent
Nous consacrons une section à cette suite exacte.

\begin{lemma}
\label{modules-lemma-canonical-exact-sequence}
Soit $X$ un espace topologique. Soit $U \subset X$ une partie ouverte de
complémentaire $Z \subset X$. Notons $j : U \to X$ l'immersion ouverte et
$i : Z \to X$ l'immersion fermée. Pour tout faisceau de groupes abéliens
$\mathcal{F}$ sur $X$, les morphismes d'adjonction
$j_{!}j^{-1}\mathcal{F} \to \mathcal{F}$ et
$\mathcal{F} \to i_*i^{-1}\mathcal{F}$ donnent une suite exacte courte
$$
0 \to j_{!}j^{-1}\mathcal{F} \to \mathcal{F} \to i_*i^{-1}\mathcal{F} \to 0
$$
de faisceaux de groupes abéliens. Pour tout morphisme
$\varphi : \mathcal{F} \to \mathcal{G}$ de faisceaux abéliens sur $X$, nous
obtenons un morphisme de suites exactes courtes
$$
\xymatrix{
0 \ar[r] &
j_{!}j^{-1}\mathcal{F} \ar[r] \ar[d] &
\mathcal{F} \ar[r] \ar[d] &
i_*i^{-1}\mathcal{F} \ar[r] \ar[d] &
0 \\
0 \ar[r] &
j_{!}j^{-1}\mathcal{G} \ar[r] &
\mathcal{G} \ar[r] &
i_*i^{-1}\mathcal{G} \ar[r] &
0
}
$$
\end{lemma}

\begin{proof}
La fonctorialité de la suite exacte courte résulte immédiatement de la
naturalité des morphismes d'adjonction. Nous pouvons vérifier l'exactitude sur
les germes (Lemme \ref{modules-lemma-abelian}). Pour une description des germes en
question, voir Faisceaux, Lemmes \ref{sheaves-lemma-j-shriek-abelian} et
\ref{sheaves-lemma-stalks-closed-pushforward}.
\end{proof}








\section{Modules localement engendrés par des sections}
\label{modules-section-locally-generated}

\noindent
Soit $(X, \mathcal{O}_X)$ un espace annelé. Dans cette section et la suivante,
nous restreindrons souvent les faisceaux aux sous-espaces ouverts
$U \subset X$, voir Faisceaux, section
\ref{sheaves-section-open-immersions}. En particulier, nous noterons souvent
le sous-espace ouvert $(U, \mathcal{O}_U)$ plutôt que d'employer la notation
plus correcte $(U, \mathcal{O}_X|_U)$, voir Faisceaux, Définition
\ref{sheaves-definition-restriction}.

\medskip\noindent
Considérons l'immersion ouverte
$j : U = (0 , \infty) \to \mathbf{R} = X$ et le faisceau abélien
$j_!\underline{\mathbf{Z}}_U$. D'après Faisceaux, section
\ref{sheaves-section-open-immersions}, le germe de
$j_!\underline{\mathbf{Z}}_U$ en $x = 0$ est $0$. En fait, les sections de ce
faisceau sur tout intervalle ouvert contenant $0$ sont $0$. Il n'existe donc
aucun voisinage ouvert du point $0$ sur lequel le faisceau puisse être engendré
par des sections.

\begin{definition}
\label{modules-definition-locally-generated}
Soit $(X, \mathcal{O}_X)$ un espace annelé. Soit $\mathcal{F}$ un faisceau de
$\mathcal{O}_X$-modules. Nous disons que $\mathcal{F}$ est
{\it localement engendré par des sections} si, pour tout $x \in X$, il existe
un voisinage ouvert $U$ de $x$ tel que $\mathcal{F}|_U$ soit engendré par ses
sections globales comme faisceau de $\mathcal{O}_U$-modules.
\end{definition}

\noindent
Autrement dit, il existe un ensemble $I$ et, pour chaque $i$, une section
$s_i \in \mathcal{F}(U)$ telle que l'application associée
$$
\bigoplus\nolimits_{i \in I} \mathcal{O}_U
\longrightarrow
\mathcal{F}|_U
$$
soit surjective.

\begin{lemma}
\label{modules-lemma-pullback-locally-generated}
Soit $f : (X, \mathcal{O}_X) \to (Y, \mathcal{O}_Y)$ un morphisme d'espaces
annelés. L'image inverse $f^*\mathcal{G}$ est localement engendrée par des
sections si $\mathcal{G}$ l'est.
\end{lemma}

\begin{proof}
Étant donné un sous-espace ouvert $V$ de $Y$, nous pouvons considérer le
diagramme commutatif d'espaces annelés
$$
\xymatrix{
(f^{-1}V, \mathcal{O}_{f^{-1}V}) \ar[r]_{j'} \ar[d]_{f'} &
(X, \mathcal{O}_X) \ar[d]^f \\
(V, \mathcal{O}_V) \ar[r]^j &
(Y, \mathcal{O}_Y)
}
$$
Nous savons que
$f^*\mathcal{G}|_{f^{-1}V} \cong (f')^*(\mathcal{G}|_V)$, voir Faisceaux,
Lemme \ref{sheaves-lemma-push-pull-composition-modules}. Nous pouvons donc
supposer que $\mathcal{G}$ est engendré par ses sections globales.

\medskip\noindent
Nous avons vu que $f^*$ commute à toutes les limites inductives et est exact à
droite, voir le Lemme \ref{modules-lemma-exactness-pushforward-pullback}. Ainsi, si
nous avons une surjection
$$
\bigoplus\nolimits_{i \in I}
\mathcal{O}_Y
\to
\mathcal{G}
\to
0
$$
alors, en appliquant $f^*$, nous obtenons la surjection
$$
\bigoplus\nolimits_{i \in I}
\mathcal{O}_X
\to
f^*\mathcal{G}
\to
0.
$$
Cela démontre le lemme.
\end{proof}












\section{Modules de type fini}
\label{modules-section-finite-type}

\begin{definition}
\label{modules-definition-finite-type}
Soit $(X, \mathcal{O}_X)$ un espace annelé. Soit $\mathcal{F}$ un faisceau de
$\mathcal{O}_X$-modules. Nous disons que $\mathcal{F}$ est de
{\it type fini} si, pour tout $x \in X$, il existe un voisinage ouvert $U$ tel
que $\mathcal{F}|_U$ soit engendré par un nombre fini de sections.
\end{definition}

\begin{lemma}
\label{modules-lemma-pullback-finite-type}
Soit $f : (X, \mathcal{O}_X) \to (Y, \mathcal{O}_Y)$ un morphisme d'espaces
annelés. L'image inverse $f^*\mathcal{G}$ d'un $\mathcal{O}_Y$-module de type
fini est un $\mathcal{O}_X$-module de type fini.
\end{lemma}

\begin{proof}
En raisonnant comme dans la démonstration du Lemme
\ref{modules-lemma-pullback-locally-generated}, nous pouvons supposer que
$\mathcal{G}$ est engendré par un nombre fini de sections globales. Nous avons
vu que $f^*$ commute à toutes les limites inductives et est exact à droite,
voir le Lemme \ref{modules-lemma-exactness-pushforward-pullback}. Ainsi, si nous avons
une surjection
$$
\bigoplus\nolimits_{i = 1, \ldots, n}
\mathcal{O}_Y
\to
\mathcal{G}
\to
0
$$
alors, en appliquant $f^*$, nous obtenons la surjection
$$
\bigoplus\nolimits_{i = 1, \ldots, n}
\mathcal{O}_X
\to
f^*\mathcal{G}
\to
0.
$$
Cela démontre le lemme.
\end{proof}

\begin{lemma}
\label{modules-lemma-extension-finite-type}
Soit $X$ un espace annelé. L'image d'un morphisme de
$\mathcal{O}_X$-modules de type fini est de type fini. Soit
$0 \to \mathcal{F}_1 \to \mathcal{F}_2 \to \mathcal{F}_3 \to 0$
une suite exacte courte de $\mathcal{O}_X$-modules. Si $\mathcal{F}_1$ et
$\mathcal{F}_3$ sont de type fini, alors $\mathcal{F}_2$ l'est aussi.
\end{lemma}

\begin{proof}
L'assertion sur les images est triviale. Celle sur les suites exactes courtes
résulte de ce que les sections de $\mathcal{F}_3$ se relèvent localement en
des sections de $\mathcal{F}_2$, ainsi que du résultat correspondant dans la
catégorie des modules sur un anneau (appliqué, par exemple, aux germes).
\end{proof}

\begin{lemma}
\label{modules-lemma-finite-type-surjective-on-stalk}
Soit $X$ un espace annelé. Soit
$\varphi : \mathcal{G} \to \mathcal{F}$ un homomorphisme de
$\mathcal{O}_X$-modules. Soit $x \in X$. Supposons $\mathcal{F}$ de type fini
et l'application sur les germes
$\varphi_x : \mathcal{G}_x \to \mathcal{F}_x$ surjective. Il existe alors un
voisinage ouvert $x \in U \subset X$ tel que $\varphi|_U$ soit surjective.
\end{lemma}

\begin{proof}
Choisissons un voisinage ouvert $U \subset X$ de $x$ tel que $\mathcal{F}$
soit engendré par $s_1, \ldots, s_n \in \mathcal{F}(U)$ sur $U$. Par
l'hypothèse de surjectivité de $\varphi_x$, après avoir rétréci $U$, nous
pouvons supposer que $s_i = \varphi(t_i)$ pour certains
$t_i \in \mathcal{G}(U)$. Alors $U$ convient.
\end{proof}

\begin{lemma}
\label{modules-lemma-finite-type-stalk-zero}
Soit $X$ un espace annelé. Soit $\mathcal{F}$ un
$\mathcal{O}_X$-module et soit $x \in X$. Supposons $\mathcal{F}$ de type
fini et $\mathcal{F}_x = 0$. Il existe alors un voisinage ouvert
$x \in U \subset X$ tel que $\mathcal{F}|_U$ soit nul.
\end{lemma}

\begin{proof}
C'est un cas particulier du Lemme
\ref{modules-lemma-finite-type-surjective-on-stalk}, appliqué au morphisme
$0 \to \mathcal{F}$.
\end{proof}

\begin{lemma}
\label{modules-lemma-support-finite-type-closed}
\begin{slogan}
Sur tout espace annelé, les faisceaux de modules de type fini ont un support
fermé.
\end{slogan}
Soit $(X, \mathcal{O}_X)$ un espace annelé. Soit $\mathcal{F}$ un faisceau de
$\mathcal{O}_X$-modules. Si $\mathcal{F}$ est de type fini, le support de
$\mathcal{F}$ est fermé.
\end{lemma}

\begin{proof}
C'est une reformulation du Lemme \ref{modules-lemma-finite-type-stalk-zero}.
\end{proof}

\begin{lemma}
\label{modules-lemma-finite-type-quasi-compact-colimit}
Soit $X$ un espace annelé. Soit $I$ un ensemble préordonné et soit
$(\mathcal{F}_i, f_{ii'})$ un système indexé par $I$ et formé de faisceaux de
$\mathcal{O}_X$-modules (voir Catégories, section
\ref{categories-section-posets-limits}). Soit
$\mathcal{F} = \colim \mathcal{F}_i$ sa limite inductive. Supposons que
(a) $I$ est filtrant, (b) $\mathcal{F}$ est un $\mathcal{O}_X$-module de type
fini, et (c) $X$ est quasi-compact. Il existe alors un $i$ tel que
$\mathcal{F}_i \to \mathcal{F}$ soit surjective. Si les applications de
transition $f_{ii'}$ sont injectives, nous en déduisons que
$\mathcal{F} = \mathcal{F}_i$ pour un certain $i \in I$.
\end{lemma}

\begin{proof}
Soit $x \in X$. Il existe un voisinage ouvert $U \subset X$ de $x$ et un
nombre fini de sections $s_j \in \mathcal{F}(U)$, $j = 1, \ldots, m$, tels
que $s_1, \ldots, s_m$ engendrent $\mathcal{F}$ comme
$\mathcal{O}_U$-module. Après avoir éventuellement rétréci $U$ en un voisinage
ouvert plus petit de $x$, nous pouvons supposer que chaque $s_j$ provient
d'une section de $\mathcal{F}_i$ pour un certain $i \in I$. Puisque $X$ est
quasi-compact, nous pouvons donc trouver un recouvrement ouvert fini
$X = \bigcup_{j = 1, \ldots, m} U_j$ et, pour chaque $j$, un indice $i_j$
ainsi qu'un nombre fini de sections
$s_{jl} \in \mathcal{F}_{i_j}(U_j)$ dont les images engendrent la restriction
de $\mathcal{F}$ à $U_j$. Le lemme vaut manifestement pour tout indice
$i \in I$ qui est $\geq$ à tous les $i_j$.
\end{proof}

\begin{lemma}
\label{modules-lemma-set-isomorphism-classes-finite-type-modules}
Soit $X$ un espace annelé. Il existe un ensemble de
$\mathcal{O}_X$-modules de type fini $\{\mathcal{F}_i\}_{i \in I}$ tel que
chaque $\mathcal{O}_X$-module de type fini sur $X$ soit isomorphe à exactement
un des $\mathcal{F}_i$.
\end{lemma}

\begin{proof}
Pour chaque recouvrement ouvert $\mathcal{U} : X = \bigcup U_j$, considérons
les faisceaux de $\mathcal{O}_X$-modules $\mathcal{F}$ tels que chaque
restriction $\mathcal{F}|_{U_j}$ soit un quotient de
$\mathcal{O}_{U_j}^{\oplus r_j}$ pour un certain $r_j \geq 0$. Ils sont
paramétrés par des sous-faisceaux
$\mathcal{K}_j \subset \mathcal{O}_{U_j}^{\oplus r_j}$ et des données de
recollement
$$
\varphi_{jj'} :
\mathcal{O}_{U_j \cap U_{j'}}^{\oplus r_j}/
(\mathcal{K}_j|_{U_j \cap U_{j'}})
\longrightarrow
\mathcal{O}_{U_j \cap U_{j'}}^{\oplus r_{j'}}/
(\mathcal{K}_{j'}|_{U_j \cap U_{j'}})
$$
voir Faisceaux, section \ref{sheaves-section-glueing-sheaves}. Remarquons que
la collection de toutes les données de recollement forme un ensemble. La
collection de tous les recouvrements
$\mathcal{U} : X = \bigcup_{j \in J} U_i$ pour lesquels
$J \to \mathcal{P}(X)$, $j \mapsto U_j$, est injective forme également un
ensemble. Par conséquent, la collection de tous les faisceaux de
$\mathcal{O}_X$-modules obtenus par recollement de quotients comme ci-dessus
forme un ensemble $\mathcal{I}$. Par définition, tout
$\mathcal{O}_X$-module de type fini est isomorphe à un élément de
$\mathcal{I}$. En choisissant un élément dans chaque classe d'isomorphisme de
$\mathcal{I}$, on obtient l'ensemble de faisceaux voulu (ce qui utilise
l'axiome du choix).
\end{proof}













\section{Modules quasi-cohérents}
\label{modules-section-quasi-coherent}

\noindent
Dans cette section, nous introduisons une notion abstraite de
$\mathcal{O}_X$-module quasi-cohérent. Cette notion est très utile en géométrie
algébrique, puisque les modules quasi-cohérents sur un schéma admettent une
bonne description sur tout ouvert affine. Nous avertissons toutefois le
lecteur que, dans le cadre général des espaces (localement) annelés, cette
notion se comporte très mal. La catégorie des faisceaux quasi-cohérents n'est
pas abélienne en général, les sommes directes infinies de faisceaux
quasi-cohérents ne sont pas quasi-cohérentes, etc.

\begin{definition}
\label{modules-definition-quasi-coherent}
Soit $(X, \mathcal{O}_X)$ un espace annelé. Soit $\mathcal{F}$ un faisceau de
$\mathcal{O}_X$-modules. Nous disons que $\mathcal{F}$ est un
{\it faisceau quasi-cohérent de $\mathcal{O}_X$-modules} si, pour tout point
$x \in X$, il existe un voisinage ouvert $x\in U \subset X$ tel que
$\mathcal{F}|_U$ soit isomorphe au conoyau d'une application
$$
\bigoplus\nolimits_{j \in J}
\mathcal{O}_U
\longrightarrow
\bigoplus\nolimits_{i \in I}
\mathcal{O}_U
$$
La catégorie des $\mathcal{O}_X$-modules quasi-cohérents est notée
$\QCoh(\mathcal{O}_X)$.
\end{definition}

\noindent
La définition signifie que $X$ est recouvert par des ouverts $U$ tels que
$\mathcal{F}|_U$ admette une {\it présentation} de la forme
$$
\bigoplus\nolimits_{j \in J}
\mathcal{O}_U
\longrightarrow
\bigoplus\nolimits_{i \in I}
\mathcal{O}_U
\longrightarrow
\mathcal{F}|_U
\longrightarrow
0.
$$
Ici, « présentation » signifie que la suite affichée est exacte. Autrement dit,
\begin{enumerate}
\item pour tout point $x$ de $X$, il existe un voisinage ouvert tel que
$\mathcal{F}|_U$ soit engendré par ses sections globales, et
\item pour un choix convenable de ces sections, le noyau de la surjection
associée est lui aussi engendré par ses sections globales.
\end{enumerate}

\begin{lemma}
\label{modules-lemma-direct-sum-quasi-coherent}
Soit $(X, \mathcal{O}_X)$ un espace annelé. La somme directe de deux
$\mathcal{O}_X$-modules quasi-cohérents est un $\mathcal{O}_X$-module
quasi-cohérent.
\end{lemma}

\begin{proof}
Omis.
\end{proof}

\begin{remark}
\label{modules-remark-infinite-direct-sum-quasi-coherent-not}
Attention : il n'est pas vrai, en général, qu'une somme directe infinie de
$\mathcal{O}_X$-modules quasi-cohérents soit quasi-cohérente. Pour des
comportements plus exotiques des modules quasi-cohérents, voir l'Exemple
\ref{modules-example-quasi-coherent}.
\end{remark}

\begin{lemma}
\label{modules-lemma-pullback-quasi-coherent}
Soit $f : (X, \mathcal{O}_X) \to (Y, \mathcal{O}_Y)$ un morphisme d'espaces
annelés. L'image inverse $f^*\mathcal{G}$ d'un $\mathcal{O}_Y$-module
quasi-cohérent est quasi-cohérente.
\end{lemma}

\begin{proof}
En raisonnant comme dans la démonstration du Lemme
\ref{modules-lemma-pullback-locally-generated}, nous pouvons supposer que
$\mathcal{G}$ admet une présentation globale par des sommes directes de copies
de $\mathcal{O}_Y$. Nous avons vu que $f^*$ commute à toutes les limites
inductives et est exact à droite, voir le Lemme
\ref{modules-lemma-exactness-pushforward-pullback}. Ainsi, si nous avons une suite
exacte
$$
\bigoplus\nolimits_{j \in J}
\mathcal{O}_Y
\longrightarrow
\bigoplus\nolimits_{i \in I}
\mathcal{O}_Y
\longrightarrow
\mathcal{G}
\longrightarrow
0
$$
alors, en appliquant $f^*$, nous obtenons la suite exacte
$$
\bigoplus\nolimits_{j \in J}
\mathcal{O}_X
\longrightarrow
\bigoplus\nolimits_{i \in I}
\mathcal{O}_X
\longrightarrow
f^*\mathcal{G}
\longrightarrow
0.
$$
Cela démontre le lemme.
\end{proof}

\noindent
Cela fournit de nombreux exemples de faisceaux quasi-cohérents.

\begin{lemma}
\label{modules-lemma-construct-quasi-coherent-sheaves}
Soit $(X, \mathcal{O}_X)$ un espace annelé.
Soit $\alpha : R \to \Gamma(X, \mathcal{O}_X)$ un homomorphisme d'un anneau
$R$ vers l'anneau des sections globales sur $X$.
Soit $M$ un $R$-module.
Les trois constructions suivantes donnent des faisceaux de
$\mathcal{O}_X$-modules canoniquement isomorphes :
\begin{enumerate}
\item Soit $\pi : (X, \mathcal{O}_X) \longrightarrow (\{*\}, R)$
le morphisme d'espaces annelés dont $\pi : X \to \{*\}$ est l'unique
application et dont l'application associée à $\pi$ est $\pi^\sharp$,
l'application donnée
$\alpha : R \to \Gamma(X, \mathcal{O}_X)$. Posons $\mathcal{F}_1 = \pi^*M$.
\item Choisissons une présentation
$\bigoplus_{j \in J} R \to \bigoplus_{i \in I} R \to M \to 0$.
Posons
$$
\mathcal{F}_2 = \Coker\left(
\bigoplus\nolimits_{j \in J} \mathcal{O}_X
\to
\bigoplus\nolimits_{i \in I} \mathcal{O}_X
\right).
$$
Ici, l'application sur le facteur $\mathcal{O}_X$ correspondant à $j \in J$
est donnée par la section $\sum_i \alpha(r_{ij})$, où les $r_{ij}$
sont les coefficients matriciels de l'application dans la présentation de $M$.
\item Posons $\mathcal{F}_3$ égal au faisceau associé au préfaisceau
$U \mapsto \mathcal{O}_X(U) \otimes_R M$, où l'application
$R \to \mathcal{O}_X(U)$ est la composée de $\alpha$ et de
l'application de restriction $\mathcal{O}_X(X) \to \mathcal{O}_X(U)$.
\end{enumerate}
Cette construction possède les propriétés suivantes :
\begin{enumerate}
\item Le faisceau de $\mathcal{O}_X$-modules ainsi obtenu
$\mathcal{F}_M = \mathcal{F}_1 = \mathcal{F}_2 = \mathcal{F}_3$
est quasi-cohérent.
\item La construction définit un foncteur de la catégorie des $R$-modules
vers la catégorie des faisceaux quasi-cohérents sur $X$, qui commute aux
limites inductives arbitraires.
\item Pour tout $x \in X$, nous avons
$\mathcal{F}_{M, x} = \mathcal{O}_{X, x} \otimes_R M$
fonctoriellement en $M$.
\item Pour tout $\mathcal{O}_X$-module
$\mathcal{G}$, nous avons
$$
\Mor_{\mathcal{O}_X}(\mathcal{F}_M, \mathcal{G})
=
\Hom_R(M, \Gamma(X, \mathcal{G}))
$$
où la structure de $R$-module sur $\Gamma(X, \mathcal{G})$
provient de la structure de $\Gamma(X, \mathcal{O}_X)$-module via
$\alpha$.
\end{enumerate}
\end{lemma}

\begin{proof}
L'isomorphisme entre $\mathcal{F}_1$ et $\mathcal{F}_3$
provient du fait que $\pi^*$ est défini comme la faisceautisation
du préfaisceau de (3), voir Faisceaux, section
\ref{sheaves-section-ringed-spaces-functoriality-modules}.
L'isomorphisme entre les constructions (2) et (1) provient
du fait que le foncteur $\pi^*$ est exact à droite, de sorte que
$\pi^*(\bigoplus_{j \in J} R) \to \pi^*(\bigoplus_{i \in I} R) \to
\pi^*M \to 0$ est exacte, que $\pi^*$ commute aux sommes directes
arbitraires, voir le Lemme \ref{modules-lemma-exactness-pushforward-pullback},
et enfin du fait que $\pi^*(R) = \mathcal{O}_X$.

\medskip\noindent
L'assertion (1) est claire d'après la construction (2).
L'assertion (2) est claire puisque $\pi^*$ possède ces propriétés.
L'assertion (3) résulte de la description des germes des faisceaux images
inverses, voir Faisceaux, Lemme
\ref{sheaves-lemma-stalk-pullback-modules}.
L'assertion (4) résulte de l'adjonction entre $\pi_*$ et
$\pi^*$.
\end{proof}

\begin{definition}
\label{modules-definition-sheaf-associated}
Dans la situation du Lemme \ref{modules-lemma-construct-quasi-coherent-sheaves},
nous disons que $\mathcal{F}_M$ est le {\it faisceau associé au module $M$
et à l'homomorphisme d'anneaux $\alpha$}. Si $R = \Gamma(X, \mathcal{O}_X)$
et $\alpha = \text{id}_R$, nous disons simplement que $\mathcal{F}_M$ est le
{\it faisceau associé au module $M$}.
\end{definition}


\begin{lemma}
\label{modules-lemma-restrict-quasi-coherent}
Soit $(X, \mathcal{O}_X)$ un espace annelé.
Posons $R = \Gamma(X, \mathcal{O}_X)$.
Soit $M$ un $R$-module.
Soit $\mathcal{F}_M$ le faisceau quasi-cohérent de
$\mathcal{O}_X$-modules associé à $M$.
Si $g : (Y, \mathcal{O}_Y) \to (X, \mathcal{O}_X)$
est un morphisme d'espaces annelés, alors
$g^*\mathcal{F}_M$ est le faisceau associé
au $\Gamma(Y, \mathcal{O}_Y)$-module
$\Gamma(Y, \mathcal{O}_Y) \otimes_R M$.
\end{lemma}

\begin{proof}
L'assertion résulte de la première description de
$\mathcal{F}_M$ dans le Lemme \ref{modules-lemma-construct-quasi-coherent-sheaves}
comme $\pi^*M$, ainsi que du diagramme commutatif suivant
d'espaces annelés
$$
\xymatrix{
(Y, \mathcal{O}_Y) \ar[r]_-\pi \ar[d]_g &
(\{*\}, \Gamma(Y, \mathcal{O}_Y)) \ar[d]^{\text{induite par }g^\sharp} \\
(X, \mathcal{O}_X) \ar[r]^-\pi &
(\{*\}, \Gamma(X, \mathcal{O}_X))
}
$$
(On utilise aussi Faisceaux, Lemme
\ref{sheaves-lemma-push-pull-composition-modules}.)
\end{proof}

\begin{lemma}
\label{modules-lemma-quasi-coherent-module}
Soit $(X, \mathcal{O}_X)$ un espace annelé.
Soit $x \in X$ un point.
Supposons que $x$ possède un système fondamental de voisinages quasi-compacts.
Considérons un $\mathcal{O}_X$-module quasi-cohérent quelconque $\mathcal{F}$.
Il existe alors un voisinage ouvert $U$ de $x$
tel que $\mathcal{F}|_U$ soit isomorphe au
faisceau de modules $\mathcal{F}_M$ sur $(U, \mathcal{O}_U)$
associé à un certain $\Gamma(U, \mathcal{O}_U)$-module $M$.
\end{lemma}

\begin{proof}
Nous pouvons d'abord remplacer $X$ par un voisinage ouvert de $x$
et supposer que $\mathcal{F}$ est isomorphe au
conoyau d'une application
$$
\Psi :
\bigoplus\nolimits_{j \in J}
\mathcal{O}_X
\longrightarrow
\bigoplus\nolimits_{i \in I}
\mathcal{O}_X.
$$
Le problème est que cette application n'est pas nécessairement donnée par
une « matrice », car le module des sections globales d'une somme directe
diffère en général de la somme directe des modules
de sections globales.

\medskip\noindent
Soit $x \in E \subset X$ un voisinage quasi-compact
de $x$ (remarquons que $E$ n'est pas nécessairement ouvert).
Soit $x \in U \subset E$ un voisinage ouvert de
$x$ contenu dans $E$.
Nous procédons ensuite comme dans la démonstration du
Lemme \ref{modules-lemma-section-direct-sum-quasi-compact}.
Pour chaque $j \in J$, notons
$s_j \in \Gamma(X, \bigoplus\nolimits_{i \in I} \mathcal{O}_X)$
l'image de la section $1$ dans le facteur $\mathcal{O}_X$
correspondant à $j$. Il existe une famille finie d'ouverts
$U_{jk}$, $k \in K_j$, telle que $E \subset \bigcup_{k \in K_j} U_{jk}$
et telle que chaque restriction $s_j|_{U_{jk}}$
soit une somme finie $\sum_{i \in I_{jk}} f_{jki}$
avec $I_{jk} \subset I$, et $f_{jki}$ dans le facteur
$\mathcal{O}_X$ correspondant à $i \in I$. Posons
$I_j = \bigcup_{k \in K_j} I_{jk}$. C'est un ensemble fini.
Puisque $U \subset E \subset \bigcup_{k \in K_j} U_{jk}$,
la section $s_j|_U$ est une section de la somme directe finie
$\bigoplus_{i \in I_j} \mathcal{O}_X$.
D'après le Lemme \ref{modules-lemma-limits-colimits},
nous voyons qu'en fait $s_j|_U$ est une somme
$\sum_{i \in I_j} f_{ij}$ et que
$f_{ij} \in \mathcal{O}_X(U) = \Gamma(U, \mathcal{O}_U)$.

\medskip\noindent
Nous pouvons alors définir un module $M$ comme le conoyau de l'application
$$
\bigoplus\nolimits_{j \in J}
\Gamma(U, \mathcal{O}_U)
\longrightarrow
\bigoplus\nolimits_{i \in I}
\Gamma(U, \mathcal{O}_U)
$$
dont la matrice est $(f_{ij})$. D'après la construction (2) du
Lemme \ref{modules-lemma-construct-quasi-coherent-sheaves}, nous voyons que
$\mathcal{F}_M$ possède la même présentation que $\mathcal{F}|_U$
et donc que $\mathcal{F}_M \cong \mathcal{F}|_U$.
\end{proof}

\begin{example}
\label{modules-example-quasi-coherent}
Soit $X$ l'espace formé d'une infinité dénombrable de copies
$L_1, L_2, L_3, \ldots$ de la droite réelle, toutes recollées en $0$ ;
un système fondamental de voisinages de $0$ est la famille
$\{U_n\}_{n \in \mathbf{N}}$, où $U_n \cap L_i = (-1/n, 1/n)$.
Soit $\mathcal{O}_X$ le faisceau des fonctions continues à valeurs réelles.
Soit $f : \mathbf{R} \to \mathbf{R}$ une fonction continue
identiquement nulle sur $(-1, 1)$ et identiquement égale à $1$
sur $(-\infty, -2) \cup (2, \infty)$. Notons $f_n$ la fonction continue
sur $X$ qui vaut $x \mapsto f(nx)$ sur chaque
$L_j = \mathbf{R}$. Soit $1_{L_j}$ la fonction caractéristique
de $L_j$. Considérons l'application
$$
\bigoplus\nolimits_{j \in \mathbf{N}}
\mathcal{O}_X
\longrightarrow
\bigoplus\nolimits_{j, i \in \mathbf{N}}
\mathcal{O}_X, \quad
e_j \longmapsto \sum\nolimits_{i \in \mathbf{N}} f_j 1_{L_i} e_{ij}
$$
avec les notations évidentes. Cela a un sens, car cette somme est localement
finie puisque $f_j$ est nulle dans un voisinage de $0$. Sur $U_n$,
l'image de $e_j$, pour $j > 2n$, n'est pas une combinaison linéaire finie
$\sum g_{ij} e_{ij}$ avec les $g_{ij}$ continues. Il n'existe donc
aucun voisinage de $0 \in X$ sur lequel l'application affichée soit donnée par
une « matrice » comme dans la démonstration du Lemme
\ref{modules-lemma-quasi-coherent-module} ci-dessus.

\medskip\noindent
Remarquons que $\bigoplus\nolimits_{j \in \mathbf{N}} \mathcal{O}_X$
est le faisceau associé au module libre de base $e_j$,
et de même pour l'autre somme directe. Nous voyons donc qu'en général,
même localement sur $X$, un morphisme entre faisceaux associés à des modules
ne provient pas d'un morphisme de modules.
Il devrait de même exister un exemple d'espace annelé $X$
et de $\mathcal{O}_X$-module quasi-cohérent $\mathcal{F}$
tel que $\mathcal{F}$ ne soit pas localement de la forme $\mathcal{F}_M$.
(Merci de nous écrire si vous en trouvez un.)
En outre, il devrait exister des exemples d'espaces localement compacts
$X$ et d'applications $\mathcal{F}_M \to \mathcal{F}_N$ qui, elles non plus,
ne proviennent pas localement d'applications de modules (la démonstration du
Lemme \ref{modules-lemma-quasi-coherent-module} montre que cela ne peut se produire
si $N$ est libre).
\end{example}












\section{Modules de présentation finie}
\label{modules-section-finite-presentation}

\noindent
Voici la définition.

\begin{definition}
\label{modules-definition-finite-presentation}
Soit $(X, \mathcal{O}_X)$ un espace annelé.
Soit $\mathcal{F}$ un faisceau de $\mathcal{O}_X$-modules.
Nous disons que $\mathcal{F}$ est de {\it présentation finie}
si, pour tout point $x \in X$, il existe un voisinage ouvert
$x\in U \subset X$ et des entiers $n, m \in \mathbf{N}$ tels que
$\mathcal{F}|_U$ soit isomorphe au conoyau d'une application
$$
\bigoplus\nolimits_{j = 1, \ldots, m}
\mathcal{O}_U
\longrightarrow
\bigoplus\nolimits_{i = 1, \ldots, n}
\mathcal{O}_U
$$
\end{definition}

\noindent
Cela signifie que $X$ est recouvert par des ouverts $U$
tels que $\mathcal{F}|_U$ possède une {\it présentation}
de la forme
$$
\bigoplus\nolimits_{j = 1, \ldots, m}
\mathcal{O}_U
\longrightarrow
\bigoplus\nolimits_{i = 1, \ldots, n}
\mathcal{O}_U
\to
\mathcal{F}|_U
\to
0.
$$
Ici, « présentation » signifie que la suite affichée
est exacte. Autrement dit,
\begin{enumerate}
\item pour tout point $x$ de $X$, il existe
un voisinage ouvert tel que $\mathcal{F}|_U$
soit engendré par un nombre fini de sections globales, et
\item pour un choix convenable de ces sections,
le noyau de la surjection associée est lui aussi
engendré par un nombre fini de sections globales.
\end{enumerate}

\begin{lemma}
\label{modules-lemma-finite-presentation-quasi-coherent}
Soit $(X, \mathcal{O}_X)$ un espace annelé.
Tout $\mathcal{O}_X$-module de présentation finie
est quasi-cohérent.
\end{lemma}

\begin{proof}
Cela découle immédiatement des définitions.
\end{proof}

\begin{lemma}
\label{modules-lemma-cokernel-finite-finite-presentation}
Soit $(X,\mathcal{O}_X)$ un espace annelé. Soit $\mathcal{F}$ un
$\mathcal{O}_X$-module de présentation finie. Soit
$\varphi : \mathcal{G} \to \mathcal{F}$ un morphisme de
$\mathcal{O}_X$-modules. Si $\mathcal{G}$ est de type fini, alors
$\Coker(\varphi)$ est de présentation finie.
\end{lemma}

\begin{proof}
Localement sur $X$, nous pouvons écrire
$\mathcal{F} = \mathcal{O}_X^{\oplus n} / \mathcal{M}$, où
$\mathcal{M} \subset \mathcal{O}_X^{\oplus n}$ est un sous-module de
$\mathcal{O}_X$ de type fini. Par conséquent,
$\Im(\varphi) = \mathcal{N} / \mathcal{M}$, où
$\mathcal{N} \subset \mathcal{O}_X^{\oplus n}$ est un sous-module de
$\mathcal{O}_X$ contenant $\mathcal{M}$.
Le $\mathcal{O}_X$-module $\Im(\varphi)$ est de type fini
parce que $\mathcal{G} \to \Im(\varphi)$ est surjective et que
$\mathcal{G}$ est de type fini. D'après le Lemme
\ref{modules-lemma-extension-finite-type}, $\mathcal{N}$ est de type fini. Ainsi,
$\Coker(\varphi) = \mathcal{O}_X^{\oplus n} / \mathcal{N}$
est de présentation finie.
\end{proof}

\begin{lemma}
\label{modules-lemma-kernel-surjection-finite-free-onto-finite-presentation}
Soit $(X, \mathcal{O}_X)$ un espace annelé.
Soit $\mathcal{F}$ un $\mathcal{O}_X$-module de présentation finie.
\begin{enumerate}
\item Si $\psi : \mathcal{O}_X^{\oplus r} \to \mathcal{F}$ est surjective,
alors $\Ker(\psi)$ est de type fini.
\item Si $\theta : \mathcal{G} \to \mathcal{F}$ est surjective et si
$\mathcal{G}$ est de type fini, alors $\Ker(\theta)$ est de type fini.
\end{enumerate}
\end{lemma}

\begin{proof}
Démonstration de (1).
Soit $x \in X$. Choisissons un voisinage ouvert $U \subset X$ de $x$
sur lequel il existe une présentation
$$
\mathcal{O}_U^{\oplus m}
\xrightarrow{\chi}
\mathcal{O}_U^{\oplus n}
\xrightarrow{\varphi}
\mathcal{F}|_U
\to
0.
$$
Soit $e_k$ la section qui engendre le $k$-ième facteur de
$\mathcal{O}_X^{\oplus r}$. Pour tout $k = 1, \ldots, r$,
nous pouvons, après avoir rétréci $U$ en un petit voisinage de $x$,
relever $\psi(e_k)$ en une section $\tilde e_k$ de
$\mathcal{O}_U^{\oplus n}$ sur $U$. Cela définit un morphisme
de faisceaux $\alpha : \mathcal{O}_U^{\oplus r} \to \mathcal{O}_U^{\oplus n}$
tel que $\varphi \circ \alpha = \psi$.
De même, après avoir rétréci $U$, nous pouvons trouver un morphisme
$\beta : \mathcal{O}_U^{\oplus n} \to \mathcal{O}_U^{\oplus r}$
tel que $\psi \circ \beta = \varphi$. Alors l'application
$$
\mathcal{O}_U^{\oplus m} \oplus
\mathcal{O}_U^{\oplus r}
\xrightarrow{\beta \circ \chi, 1 - \beta \circ \alpha}
\mathcal{O}_U^{\oplus r}
$$
est une surjection sur le noyau de $\psi$.

\medskip\noindent
Pour démontrer (2), nous pouvons choisir localement une surjection
$\eta : \mathcal{O}_X^{\oplus r} \to \mathcal{G}$.
D'après (1), $\Ker(\theta \circ \eta)$ est de type fini.
Comme $\Ker(\theta) = \eta(\Ker(\theta \circ \eta))$, le résultat s'ensuit.
\end{proof}

\begin{lemma}
\label{modules-lemma-pullback-finite-presentation}
Soit $f : (X, \mathcal{O}_X) \to (Y, \mathcal{O}_Y)$
un morphisme d'espaces annelés.
L'image inverse $f^*\mathcal{G}$ d'un module de présentation finie
est de présentation finie.
\end{lemma}

\begin{proof}
La démonstration est exactement la même que celle du Lemme
\ref{modules-lemma-pullback-quasi-coherent}, mais avec des ensembles d'indices finis.
\end{proof}

\begin{lemma}
\label{modules-lemma-quasi-coherent-limit-finite-presentation}
Soit $(X, \mathcal{O}_X)$ un espace annelé.
Posons $R = \Gamma(X, \mathcal{O}_X)$.
Soit $M$ un $R$-module.
Le $\mathcal{O}_X$-module $\mathcal{F}_M$ associé à $M$
est une limite inductive filtrante de $\mathcal{O}_X$-modules de présentation
finie.
\end{lemma}

\begin{proof}
Cela résulte immédiatement du
Lemme \ref{modules-lemma-construct-quasi-coherent-sheaves}
et du fait que tout module est une limite inductive filtrante
de modules de présentation finie, voir
Algèbre, Lemme \ref{algebra-lemma-module-colimit-fp}.
\end{proof}

\begin{lemma}
\label{modules-lemma-finite-presentation-stalk-free}
Soit $(X, \mathcal{O}_X)$ un espace annelé. Soit $\mathcal{F}$ un
$\mathcal{O}_X$-module de présentation finie. Soit $x \in X$ tel que
$\mathcal{F}_x \cong \mathcal{O}_{X, x}^{\oplus r}$. Il existe alors
un voisinage ouvert $U$ de $x$ tel que
$\mathcal{F}|_U \cong \mathcal{O}_U^{\oplus r}$.
\end{lemma}

\begin{proof}
Choisissons $s_1, \ldots, s_r \in \mathcal{F}_x$ qui correspondent, par
l'isomorphisme, à une base de $\mathcal{O}_{X, x}^{\oplus r}$. Choisissons un
voisinage ouvert $U$ de $x$ tel que $s_i$ se relève en
$s_i \in \mathcal{F}(U)$. Après avoir rétréci $U$, nous voyons que
l'application induite $\psi : \mathcal{O}_U^{\oplus r} \to \mathcal{F}|_U$
est surjective (Lemme \ref{modules-lemma-finite-type-surjective-on-stalk}).
D'après le Lemme
\ref{modules-lemma-kernel-surjection-finite-free-onto-finite-presentation},
$\Ker(\psi)$ est de type fini.
Alors $\Ker(\psi)_x = 0$ implique que $\Ker(\psi)$ devient nul
après avoir encore rétréci $U$ (Lemme \ref{modules-lemma-finite-type-stalk-zero}).
\end{proof}
\section{Modules cohérents}
\label{modules-section-coherent}

\noindent
Une référence pour cette section est \cite{FAC}.

\medskip\noindent
La catégorie des faisceaux cohérents sur un espace annelé $X$
est un objet plus raisonnable
que la catégorie des faisceaux quasi-cohérents, en ce sens
qu'elle est au moins une sous-catégorie abélienne de $\textit{Mod}(\mathcal{O}_X)$
quel que soit $X$. En revanche, l'image réciproque d'un
module cohérent n'est « presque jamais » cohérente dans le cadre général
des espaces annelés.

\begin{definition}
\label{modules-definition-coherent}
Soit $(X, \mathcal{O}_X)$ un espace annelé.
Soit $\mathcal{F}$ un faisceau de $\mathcal{O}_X$-modules.
On dit que $\mathcal{F}$ est un {\it $\mathcal{O}_X$-module cohérent}
si les deux conditions suivantes sont satisfaites :
\begin{enumerate}
\item $\mathcal{F}$ est de type fini, et
\item pour tout ouvert $U \subset X$ et toute famille finie
$s_i \in \mathcal{F}(U)$, $i = 1, \ldots, n$,
le noyau de l'application associée
$\bigoplus_{i = 1, \ldots, n} \mathcal{O}_U \to \mathcal{F}|_U$
est de type fini.
\end{enumerate}
La catégorie des $\mathcal{O}_X$-modules cohérents est notée
$\textit{Coh}(\mathcal{O}_X)$.
\end{definition}

\begin{lemma}
\label{modules-lemma-coherent-finite-presentation}
Soit $(X, \mathcal{O}_X)$ un espace annelé.
Tout $\mathcal{O}_X$-module cohérent est de présentation finie
et, par conséquent, quasi-cohérent.
\end{lemma}

\begin{proof}
Soit $\mathcal{F}$ un faisceau cohérent sur $X$.
Choisissons un point $x \in X$.
D'après (1) de la définition de la cohérence, il existe un voisinage ouvert $U$
et des sections $s_i$, $i = 1, \ldots, n$, de $\mathcal{F}$ sur $U$
tels que $\Psi : \bigoplus_{i = 1, \ldots, n} \mathcal{O}_U \to \mathcal{F}$
soit surjective. D'après (2) de la définition de la cohérence, il existe
un voisinage ouvert $V$, $x \in V \subset U$, et des sections
$t_1, \ldots, t_m$ de $\bigoplus_{i = 1, \ldots, n} \mathcal{O}_V$
qui engendrent le noyau de $\Psi|_V$. On obtient alors sur $V$ la
présentation
$$
\bigoplus\nolimits_{j = 1, \ldots, m}
\mathcal{O}_V
\longrightarrow
\bigoplus\nolimits_{i = 1, \ldots, n}
\mathcal{O}_V
\to
\mathcal{F}|_V
\to
0
$$
voulue.
\end{proof}

\begin{example}
\label{modules-example-coherent-not-Noetherian}
Supposons que $X$ soit un point. Dans ce cas, la définition
ci-dessus donne une notion pour les modules sur les anneaux.
Que signifie la définition de la cohérence ?
Elle est étroitement liée à la notion de noethérianité,
mais ne lui est pas identique : en effet, l'anneau
$R = \mathbf{C}[x_1, x_2, x_3, \ldots]$ est cohérent
comme module sur lui-même, mais n'est pas noethérien comme module
sur lui-même. Voir Algèbre, section \ref{algebra-section-coherent},
pour une discussion plus approfondie.
\end{example}

\begin{lemma}
\label{modules-lemma-coherent-abelian}
Soit $(X, \mathcal{O}_X)$ un espace annelé.
\begin{enumerate}
\item Tout sous-faisceau de type fini d'un faisceau cohérent est cohérent.
\item Soit $\varphi : \mathcal{F} \to \mathcal{G}$
un morphisme d'un faisceau de type fini $\mathcal{F}$
vers un faisceau cohérent $\mathcal{G}$. Alors $\Ker(\varphi)$ est de type fini.
\item Soit $\varphi : \mathcal{F} \to \mathcal{G}$ un morphisme
de $\mathcal{O}_X$-modules cohérents. Alors
$\Ker(\varphi)$ et
$\Coker(\varphi)$ sont cohérents.
\item Étant donnée une suite exacte courte de $\mathcal{O}_X$-modules
$0 \to \mathcal{F}_1 \to \mathcal{F}_2 \to \mathcal{F}_3 \to 0$
si deux des trois termes sont cohérents, le troisième l'est aussi.
\item La catégorie $\textit{Coh}(\mathcal{O}_X)$ est une sous-catégorie de Serre faible
de $\textit{Mod}(\mathcal{O}_X)$. En particulier, la catégorie des
modules cohérents est abélienne et le foncteur d'inclusion
$\textit{Coh}(\mathcal{O}_X) \to \textit{Mod}(\mathcal{O}_X)$
est exact.
\end{enumerate}
\end{lemma}

\begin{proof}
La condition (2) de la définition \ref{modules-definition-coherent}
est satisfaite par tout sous-faisceau d'un faisceau cohérent. On obtient donc (1).

\medskip\noindent
Supposons vérifiées les hypothèses de (2).
Montrons que $\Ker(\varphi)$ est de type fini. Choisissons $x \in X$.
Choisissons un voisinage ouvert $U$ de $x$ dans $X$ tel
que $\mathcal{F}|_U$ soit engendré par $s_1, \ldots, s_n$.
D'après la définition \ref{modules-definition-coherent}, le noyau $\mathcal{K}$
de l'application induite
$\bigoplus_{i = 1}^n \mathcal{O}_U \to \mathcal{G}$,
$e_i \mapsto \varphi(s_i)$ est de type fini.
Par conséquent, $\Ker(\varphi)$, qui est l'image de la
composée
$\mathcal{K} \to \bigoplus_{i = 1}^n \mathcal{O}_U \to \mathcal{F}$
est de type fini.

\medskip\noindent
Supposons vérifiées les hypothèses de (3).
D'après (2), le noyau de $\varphi$ est de type fini et
il est donc cohérent d'après (1).

\medskip\noindent
Sous les mêmes hypothèses, montrons que $\Coker(\varphi)$ est cohérent.
Puisque $\mathcal{G}$ est de type fini, $\Coker(\varphi)$ l'est aussi.
Soit $U \subset X$ un ouvert et soient $\overline{s}_i \in \Coker(\varphi)(U)$,
$i = 1, \ldots, n$, des sections. Nous devons montrer que le noyau du
morphisme associé
$\overline{\Psi} : \bigoplus_{i = 1}^n \mathcal{O}_U \to \Coker(\varphi)$
est de type fini. Il existe un recouvrement ouvert
de $U$ tel que, sur chaque ouvert, toutes les sections $\overline{s}_i$
se relèvent en des sections $s_i$ de $\mathcal{G}$. Nous pouvons donc supposer
que tel est le cas sur $U$. Nous pouvons en outre supposer qu'il existe des
sections $t_j$, $j = 1, \ldots, m$, de $\Im(\varphi)$ sur $U$
qui engendrent $\Im(\varphi)$ sur $U$.
Soit $\Phi : \bigoplus_{j = 1}^m \mathcal{O}_U \to \Im(\varphi)$
définie à l'aide des $t_j$, et
$\Psi :
\bigoplus_{j = 1}^m \mathcal{O}_U \oplus
\bigoplus_{i = 1}^n \mathcal{O}_U \to \mathcal{G}$
définie à l'aide des $t_j$ et des $s_i$.
Considérons le diagramme commutatif suivant
$$
\xymatrix{
0 \ar[r] &
\bigoplus_{j = 1}^m \mathcal{O}_U \ar[d]_\Phi \ar[r] &
\bigoplus_{j = 1}^m \mathcal{O}_U \oplus
\bigoplus_{i = 1}^n \mathcal{O}_U \ar[d]_\Psi \ar[r] &
\bigoplus_{i = 1}^n \mathcal{O}_U \ar[d]_{\overline{\Psi}} \ar[r] &
0 \\
0 \ar[r] &
\Im(\varphi) \ar[r] &
\mathcal{G} \ar[r] &
\Coker(\varphi) \ar[r] &
0
}
$$
Le lemme du serpent donne une suite exacte
$\Ker(\Psi) \to \Ker(\overline{\Psi}) \to 0$. Puisque $\Ker(\Psi)$
est un module de type fini, on voit que $\Ker(\overline{\Psi})$ est de type fini.

\medskip\noindent
Démonstration de la partie (4).
Soit $0 \to \mathcal{F}_1 \to \mathcal{F}_2 \to \mathcal{F}_3 \to 0$
une suite exacte courte de $\mathcal{O}_X$-modules. D'après la partie
(3), il suffit
de prouver que, si $\mathcal{F}_1$ et $\mathcal{F}_3$ sont cohérents,
$\mathcal{F}_2$ l'est aussi. Le lemme \ref{modules-lemma-extension-finite-type}
montre que $\mathcal{F}_2$ est de type fini. Soient
$s_1, \ldots, s_n$ un nombre fini de sections locales
de $\mathcal{F}_2$ définies sur un même ouvert $U$ de $X$.
Nous devons montrer que le module des relations $\mathcal{K}$
entre celles-ci est de type fini.
Considérons le diagramme commutatif suivant
$$
\xymatrix{
0 \ar[r] &
0 \ar[r] \ar[d] &
\bigoplus_{i = 1}^{n} \mathcal{O}_U \ar[r] \ar[d] &
\bigoplus_{i = 1}^{n} \mathcal{O}_U \ar[r] \ar[d] &
0 \\
0 \ar[r] &
\mathcal{F}_1 \ar[r] &
\mathcal{F}_2 \ar[r] &
\mathcal{F}_3 \ar[r] &
0
}
$$
avec les notations évidentes. Le lemme du serpent
donne une suite exacte courte
$0 \to \mathcal{K} \to \mathcal{K}_3 \to \mathcal{F}_1$
où $\mathcal{K}_3$ est le module des relations entre
les images des sections $s_i$ dans $\mathcal{F}_3$.
Puisque $\mathcal{F}_1$ est cohérent, on voit que
$\mathcal{K}$ est le noyau d'une application d'un module de type fini
vers un module cohérent, et est donc de type fini d'après (2).

\medskip\noindent
Démonstration de (5). Cela résulte de (3) et (4), qui montrent que
le lemme \ref{homology-lemma-characterize-weak-serre-subcategory} d'Homologie
s'applique.
\end{proof}

\begin{lemma}
\label{modules-lemma-coherent-structure-sheaf}
Soit $(X, \mathcal{O}_X)$ un espace annelé.
Soit $\mathcal{F}$ un $\mathcal{O}_X$-module.
Supposons que $\mathcal{O}_X$ soit un $\mathcal{O}_X$-module cohérent.
Alors $\mathcal{F}$ est cohérent si et seulement s'il est
de présentation finie.
\end{lemma}

\begin{proof}
Omis.
\end{proof}

\begin{lemma}
\label{modules-lemma-finite-type-to-coherent-injective-on-stalk}
Soit $X$ un espace annelé.
Soit $\varphi : \mathcal{G} \to \mathcal{F}$ un homomorphisme
de $\mathcal{O}_X$-modules.
Soit $x \in X$. Supposons $\mathcal{G}$ de type fini,
$\mathcal{F}$ cohérent et l'application sur les germes
$\varphi_x : \mathcal{G}_x \to \mathcal{F}_x$ injective.
Alors il existe un voisinage ouvert
$x \in U \subset X$ tel que $\varphi|_U$ soit injective.
\end{lemma}

\begin{proof}
Notons $\mathcal{K} \subset \mathcal{G}$ le noyau de $\varphi$.
Le lemme \ref{modules-lemma-coherent-abelian} montre que $\mathcal{K}$ est
un $\mathcal{O}_X$-module de type fini. Notre hypothèse est que
$\mathcal{K}_x = 0$. D'après le lemme \ref{modules-lemma-finite-type-stalk-zero},
il existe un voisinage ouvert $U$ de $x$ tel que $\mathcal{K}|_U = 0$.
Cet ouvert $U$ convient.
\end{proof}













\section{Immersions fermées d'espaces annelés}
\label{modules-section-closed-immersion}

\noindent
Quand déclarons-nous qu'un morphisme d'espaces annelés
$i : (Z, \mathcal{O}_Z) \to (X, \mathcal{O}_X)$
est une immersion fermée ?

\medskip\noindent
Motivés par l'exemple d'une immersion fermée d'espaces topologiques normaux
(annelés par le faisceau des fonctions continues), ou de variétés différentielles
(annelées par le faisceau des fonctions différentiables), il semble naturel de
supposer au moins que :
\begin{enumerate}
\item L'application $i$ est une immersion fermée d'espaces topologiques.
\item L'application associée $\mathcal{O}_X \to i_*\mathcal{O}_Z$
est surjective. Notons son noyau $\mathcal{I}$.
\end{enumerate}
Ces conditions impliquent déjà plusieurs résultats agréables : par exemple,
nous démontrons que la catégorie des $\mathcal{O}_Z$-modules est équivalente à
la catégorie des $\mathcal{O}_X$-modules annulés par $\mathcal{I}$,
ce qui généralise le résultat sur les faisceaux abéliens de la
section \ref{modules-section-closed-immersions}.

\medskip\noindent
Cependant, dans le projet Champs nous choisissons la
définition qui garantit que, si $i$ est une immersion fermée
et $(X, \mathcal{O}_X)$ un schéma, alors $(Z, \mathcal{O}_Z)$
est également un schéma. De plus, dans cette situation, nous voulons que $i_*$ et $i^*$
induisent une équivalence entre la catégorie des
$\mathcal{O}_Z$-modules quasi-cohérents et la catégorie des
$\mathcal{O}_X$-modules quasi-cohérents annulés par $\mathcal{I}$.
Une condition minimale est que $i_*\mathcal{O}_Z$ soit un
faisceau quasi-cohérent de $\mathcal{O}_X$-modules.
Une bonne manière de garantir que $i_*\mathcal{O}_Z$ est un
$\mathcal{O}_X$-module quasi-cohérent consiste à supposer que
$\mathcal{I}$ est localement engendré par des sections.
On peut interpréter cette condition en disant que « $(Z, \mathcal{O}_Z)$ est
localement défini sur $(X, \mathcal{O}_X)$ en annulant certaines fonctions régulières
$f_i$, c'est-à-dire des sections locales de $\mathcal{O}_X$ ».
Ceci conduit à la définition suivante.

\begin{definition}
\label{modules-definition-closed-immersion}
Une {\it immersion fermée d'espaces annelés}\footnote{Cette terminologie n'est
pas standard ; voir la discussion ci-dessus.} est un morphisme
$i : (Z, \mathcal{O}_Z) \to (X, \mathcal{O}_X)$
ayant les propriétés suivantes :
\begin{enumerate}
\item L'application $i$ est une immersion fermée d'espaces topologiques.
\item L'application associée $\mathcal{O}_X \to i_*\mathcal{O}_Z$
est surjective. Notons son noyau $\mathcal{I}$.
\item Le $\mathcal{O}_X$-module $\mathcal{I}$ est localement
engendré par des sections.
\end{enumerate}
\end{definition}

\noindent
En fait, cette définition ne garantit toujours pas que
l'image directe par $i_*$ d'un $\mathcal{O}_Z$-module quasi-cohérent soit un
$\mathcal{O}_X$-module quasi-cohérent. Le problème est que
l'on ne voit pas clairement comment transformer une présentation locale d'un
$\mathcal{O}_Z$-module quasi-cohérent en une présentation locale
de son image directe. Toutefois, l'énoncé suivant
est immédiat.

\begin{lemma}
\label{modules-lemma-i-star-quasi-coherent}
Soit $i : (Z, \mathcal{O}_Z) \to (X, \mathcal{O}_X)$
une immersion fermée d'espaces annelés.
Soit $\mathcal{F}$ un $\mathcal{O}_Z$-module quasi-cohérent.
Alors $i_*\mathcal{F}$ est localement sur $X$ le conoyau d'une
application de $\mathcal{O}_X$-modules quasi-cohérents.
\end{lemma}

\begin{proof}
Ceci est vrai parce que $i_*\mathcal{O}_Z$ est quasi-cohérent
par définition. Et, localement sur $Z$, le faisceau $\mathcal{F}$
est le conoyau d'une application entre sommes directes de copies
de $\mathcal{O}_Z$. De plus, toute somme directe de copies d'un
{\it même} faisceau quasi-cohérent est quasi-cohérente.
Enfin, $i_*$ commute aux colimites arbitraires,
voir le lemme \ref{modules-lemma-i-star-right-adjoint}. Certains détails sont omis.
\end{proof}

\begin{lemma}
\label{modules-lemma-i-star-reflects-finite-type}
Soit $i : (Z, \mathcal{O}_Z) \to (X, \mathcal{O}_X)$ un morphisme d'espaces
annelés. Supposons que $i$ soit un homéomorphisme sur un fermé de $X$ et
que $\mathcal{O}_X \to i_*\mathcal{O}_Z$ soit surjective.
Soit $\mathcal{F}$ un $\mathcal{O}_Z$-module.
Alors $i_*\mathcal{F}$ est de type fini si et seulement si
$\mathcal{F}$ est de type fini.
\end{lemma}

\begin{proof}
Supposons que $\mathcal{F}$ soit de type fini.
Choisissons $x \in X$. Si $x \not \in Z$, alors $i_*\mathcal{F}$
est nul dans un voisinage de $x$, et donc engendré par un nombre fini d'éléments dans
un voisinage de $x$. Si $x = i(z)$, choisissons un voisinage ouvert
$z \in V \subset Z$ et des sections $s_1, \ldots, s_n \in \mathcal{F}(V)$
qui engendrent $\mathcal{F}$ sur $V$. Écrivons $V = Z \cap U$ pour un ouvert
$U \subset X$. Remarquons que $U$ est un voisinage de $x$. Les
sections $s_i$ donnent clairement des sections $s_i$ de $i_*\mathcal{F}$ sur $U$.
L'application qui en résulte
$$
\bigoplus\nolimits_{i = 1, \ldots, n} \mathcal{O}_U
\longrightarrow
i_*\mathcal{F}|_U
$$
est surjective, comme on le voit sur les germes
(on utilise ici la surjectivité de $\mathcal{O}_X \to i_*\mathcal{O}_Z$).
Par conséquent, $i_*\mathcal{F}$
est de type fini.

\medskip\noindent
Réciproquement, supposons que $i_*\mathcal{F}$ soit de type fini.
Choisissons $z \in Z$. Posons $x = i(z)$. Par hypothèse, il existe
un voisinage ouvert $U \subset X$ de $x$, et des sections
$s_1, \ldots, s_n \in (i_*\mathcal{F})(U)$ qui engendrent $i_*\mathcal{F}$
sur $U$. Posons $V = Z \cap U$. Par définition de $i_*$, les sections
$s_i$ correspondent à des sections $s_i$ de $\mathcal{F}$ sur $V$.
L'application qui en résulte
$$
\bigoplus\nolimits_{i = 1, \ldots, n} \mathcal{O}_V
\longrightarrow
\mathcal{F}|_V
$$
est surjective, comme on le voit sur les germes. Par conséquent,
$\mathcal{F}$ est de type fini.
\end{proof}

\begin{lemma}
\label{modules-lemma-i-star-equivalence}
Soit $i : (Z, \mathcal{O}_Z) \to (X, \mathcal{O}_X)$ un morphisme
d'espaces annelés. Supposons que $i$ soit un homéomorphisme sur un fermé de $X$
et que $i^\sharp : \mathcal{O}_X \to i_*\mathcal{O}_Z$ soit surjective.
Notons $\mathcal{I} \subset \mathcal{O}_X$ le noyau de $i^\sharp$.
Le foncteur
$$
i_* :
\textit{Mod}(\mathcal{O}_Z)
\longrightarrow
\textit{Mod}(\mathcal{O}_X)
$$
est exact, pleinement fidèle, et a pour image essentielle les
$\mathcal{O}_X$-modules $\mathcal{G}$ tels que $\mathcal{I}\mathcal{G} = 0$.
\end{lemma}

\begin{proof}
Nous affirmons que, pour un $\mathcal{O}_Z$-module $\mathcal{F}$, l'application canonique
$$
i^*i_*\mathcal{F} \longrightarrow \mathcal{F}
$$
est un isomorphisme. Vérifions-le sur les germes. Soient $z \in Z$ et $x = i(z)$.
Nous avons
$$
(i^*i_*\mathcal{F})_z =
(i_*\mathcal{F})_x \otimes_{\mathcal{O}_{X, x}} \mathcal{O}_{Z, z} =
\mathcal{F}_z \otimes_{\mathcal{O}_{X, x}} \mathcal{O}_{Z, z} =
\mathcal{F}_z
$$
d'après le lemme \ref{sheaves-lemma-stalk-pullback-modules} de Faisceaux,
le fait que $\mathcal{O}_{Z, z}$ est un quotient de $\mathcal{O}_{X, x}$, et
le lemme \ref{sheaves-lemma-stalks-closed-pushforward} de Faisceaux.
Il s'ensuit que $i_*$ est pleinement fidèle.

\medskip\noindent
Soit $\mathcal{G}$ un $\mathcal{O}_X$-module tel que
$\mathcal{I}\mathcal{G} = 0$. Nous allons démontrer que l'application canonique
$$
\mathcal{G} \longrightarrow i_*i^*\mathcal{G}
$$
est un isomorphisme. Ceci prouve que $\mathcal{G} = i_*\mathcal{F}$
avec $\mathcal{F} = i^*\mathcal{G}$, ce qui achève la démonstration.
Vérifions que l'application affichée induit un isomorphisme sur les germes.
Si $x \in X$, $x \not \in i(Z)$, alors $\mathcal{G}_x = 0$
car $\mathcal{I}_x = \mathcal{O}_{X, x}$ dans ce
cas. Comme ci-dessus, $(i_*i^*\mathcal{G})_x = 0$ d'après le
lemme \ref{sheaves-lemma-stalks-closed-pushforward} de Faisceaux.
En revanche, si $x \in Z$, nous obtenons l'application
$$
\mathcal{G}_x
\longrightarrow
\mathcal{G}_x \otimes_{\mathcal{O}_{X, x}} \mathcal{O}_{Z, x}
$$
d'après les lemmes \ref{sheaves-lemma-stalk-pullback-modules} et
\ref{sheaves-lemma-stalks-closed-pushforward} de Faisceaux. Cette application est un isomorphisme
car $\mathcal{O}_{Z, x} = \mathcal{O}_{X, x}/\mathcal{I}_x$
et parce que $\mathcal{G}_x$ est annulé par $\mathcal{I}_x$ par hypothèse.
\end{proof}

\begin{remark}
\label{modules-remark-sections-support-in-closed-modules}
Soit $(X, \mathcal{O}_X)$ un espace annelé. Soit $Z \subset X$
un fermé. Pour un $\mathcal{O}_X$-module $\mathcal{F}$, on peut
considérer le {\it sous-module des sections à support dans $Z$}, noté
$\mathcal{H}_Z(\mathcal{F})$, défini par la règle
$$
\mathcal{H}_Z(\mathcal{F})(U) =
\{s \in \mathcal{F}(U) \mid \text{Supp}(s) \subset U \cap Z\}
$$
Observons que $\mathcal{H}_Z(\mathcal{F})(U)$ est un module sur
$\mathcal{O}_X(U)$, c'est-à-dire que $\mathcal{H}_Z(\mathcal{F})$ est un
$\mathcal{O}_X$-module. Par construction, $\mathcal{H}_Z(\mathcal{F})$
est le plus grand $\mathcal{O}_X$-sous-module de $\mathcal{F}$ dont le support est
contenu dans $Z$.
En appliquant le lemme \ref{modules-lemma-i-star-equivalence} au
morphisme d'espaces annelés $(Z, \mathcal{O}_X|_Z) \to (X, \mathcal{O}_X)$, nous
pouvons (et nous allons) considérer
$\mathcal{H}_Z(\mathcal{F})$ comme un $\mathcal{O}_X|_Z$-module sur $Z$.
Nous obtenons ainsi un foncteur
$$
\textit{Mod}(\mathcal{O}_X) \longrightarrow \textit{Mod}(\mathcal{O}_X|_Z),
\quad
\mathcal{F} \longmapsto \mathcal{H}_Z(\mathcal{F})
\text{ considéré comme }\mathcal{O}_X|_Z\text{-module sur }Z
$$
Ce foncteur est exact à gauche, mais n'est pas exact en général.
Tous les énoncés ci-dessus résultent directement du
lemme \ref{modules-lemma-support-section-closed}.
Cette construction est manifestement compatible avec celle de la
remarque \ref{modules-remark-sections-support-in-closed}.
\end{remark}

\begin{lemma}
\label{modules-lemma-adjoint-section-with-support}
Soit $(X, \mathcal{O}_X)$ un espace annelé. Soit $i : Z \to X$
l'inclusion d'un fermé. Le foncteur
$\mathcal{H}_Z : \textit{Mod}(\mathcal{O}_X) \to
\textit{Mod}(\mathcal{O}_X|_Z)$ de la
remarque \ref{modules-remark-sections-support-in-closed-modules}
est adjoint à droite de
$i_* : \textit{Mod}(\mathcal{O}_X|_Z) \to \textit{Mod}(\mathcal{O}_X)$.
\end{lemma}

\begin{proof}
Nous devons montrer que, pour tout $\mathcal{O}_X$-module $\mathcal{F}$
et tout $\mathcal{O}_X|_Z$-module $\mathcal{G}$, nous avons
$$
\Hom_{\mathcal{O}_X|_Z}(\mathcal{G}, \mathcal{H}_Z(\mathcal{F})) =
\Hom_{\mathcal{O}_X}(i_*\mathcal{G}, \mathcal{F})
$$
C'est clair,
car toute section de $i_*\mathcal{G}$ est après tout à support dans $Z$.
Les détails sont omis.
\end{proof}










\section{Faisceaux localement libres}
\label{modules-section-locally-free}

\noindent
Soit $(X, \mathcal{O}_X)$ un espace annelé.
Nos conventions permettent que (certains) germes $\mathcal{O}_{X, x}$
soient l'anneau nul. Cela signifie qu'il faut prendre quelques précautions
lors de la définition du rang d'un faisceau localement libre.

\begin{definition}
\label{modules-definition-locally-free}
Soit $(X, \mathcal{O}_X)$ un espace annelé.
Soit $\mathcal{F}$ un faisceau de $\mathcal{O}_X$-modules.
\begin{enumerate}
\item On dit que $\mathcal{F}$ est {\it localement libre} si, pour tout
point $x \in X$, il existe un ensemble $I$ et un
voisinage ouvert $x \in U \subset X$
tels que $\mathcal{F}|_U$ soit isomorphe à
$\bigoplus_{i \in I} \mathcal{O}_X|_U$ comme $\mathcal{O}_X|_U$-module.
\item On dit que $\mathcal{F}$ est {\it localement libre de type fini} si l'on peut
choisir les ensembles d'indices $I$ finis.
\item On dit que $\mathcal{F}$ est {\it localement libre de type fini et de rang $r$}
si l'on peut choisir les ensembles d'indices $I$ de cardinal $r$.
\end{enumerate}
\end{definition}

\noindent
Une somme directe finie de faisceaux localement libres (de type fini) est
localement libre (de type fini). Toutefois, une somme directe infinie
de faisceaux localement libres peut ne pas être localement libre.

\begin{lemma}
\label{modules-lemma-locally-free-quasi-coherent}
Soit $(X, \mathcal{O}_X)$ un espace annelé.
Soit $\mathcal{F}$ un faisceau de $\mathcal{O}_X$-modules.
Si $\mathcal{F}$ est localement libre, alors il est quasi-cohérent.
\end{lemma}

\begin{proof}
Omis.
\end{proof}

\begin{lemma}
\label{modules-lemma-pullback-locally-free}
Soit $f : (X, \mathcal{O}_X) \to (Y, \mathcal{O}_Y)$
un morphisme d'espaces annelés. Si $\mathcal{G}$ est
un $\mathcal{O}_Y$-module localement libre, alors
$f^*\mathcal{G}$ est un $\mathcal{O}_X$-module localement libre.
\end{lemma}

\begin{proof}
Omis.
\end{proof}

\begin{lemma}
\label{modules-lemma-rank}
Soit $(X, \mathcal{O}_X)$ un espace annelé.
Supposons que le support de $\mathcal{O}_X$ soit $X$,
c'est-à-dire que tous les germes de $\mathcal{O}_X$ soient des anneaux non nuls.
Soit $\mathcal{F}$ un faisceau localement libre de $\mathcal{O}_X$-modules.
Il existe une fonction localement constante
$$
\text{rang}_\mathcal{F} :
X \longrightarrow \{0, 1, 2, \ldots\}\cup\{\infty\}
$$
telle que, pour tout point $x \in X$, le cardinal de tout
ensemble $I$ tel que $\mathcal{F}$ soit isomorphe à
$\bigoplus_{i\in I} \mathcal{O}_X$ dans un voisinage
de $x$ soit $\text{rang}_\mathcal{F}(x)$.
\end{lemma}

\begin{proof}
Sous l'hypothèse du lemme, le cardinal de $I$ se lit
sur le rang du module libre $\mathcal{F}_x$ sur l'anneau non nul
$\mathcal{O}_{X, x}$, et il est constant dans un voisinage de $x$.
\end{proof}

\begin{lemma}
\label{modules-lemma-map-finite-locally-free}
Soit $(X, \mathcal{O}_X)$ un espace annelé. Soit $r \geq 0$.
Soit $\varphi : \mathcal{F} \to \mathcal{G}$ une application de
$\mathcal{O}_X$-modules localement libres de type fini et de rang $r$.
Alors $\varphi$ est un isomorphisme si et seulement si $\varphi$
est surjective.
\end{lemma}

\begin{proof}
Supposons $\varphi$ surjective. Choisissons $x \in X$.
Il existe un voisinage ouvert $U$ de $x$
tel que $\mathcal{F}|_U$ et $\mathcal{G}|_U$ soient tous deux
isomorphes à $\mathcal{O}_U^{\oplus r}$.
Choisissons des relèvements des générateurs libres de $\mathcal{G}|_U$
pour obtenir une application $\psi : \mathcal{G}|_U \to \mathcal{F}|_U$
telle que $\varphi|_U \circ \psi = \text{id}$.
Nous en déduisons que l'application
$\Gamma(U, \mathcal{F}) \to \Gamma(U, \mathcal{G})$
induite par $\varphi$ est surjective. Puisque
$\Gamma(U, \mathcal{F})$ et $\Gamma(U, \mathcal{G})$
sont isomorphes à $\Gamma(U, \mathcal{O}_U)^{\oplus r}$
comme $\Gamma(U, \mathcal{O}_U)$-modules, nous pouvons
appliquer le lemme \ref{algebra-lemma-fun} d'Algèbre pour voir que
$\Gamma(U, \mathcal{F}) \to \Gamma(U, \mathcal{G})$
est injective. Ceci achève la démonstration.
\end{proof}

\begin{lemma}
\label{modules-lemma-direct-summand-of-locally-free-is-locally-free}
Soit $(X, \mathcal{O}_X)$ un espace annelé. Si tous les germes $\mathcal{O}_{X, x}$
sont des anneaux locaux, alors tout facteur direct d'un
$\mathcal{O}_X$-module localement libre de type fini est localement libre de type fini.
\end{lemma}

\begin{proof}
Supposons que $\mathcal{F}$ soit un facteur direct du
$\mathcal{O}_X$-module localement libre de type fini $\mathcal{H}$.
Soit $x \in X$ un point. Alors $\mathcal{H}_x$ est un
$\mathcal{O}_{X, x}$-module libre de type fini.
Comme $\mathcal{O}_{X, x}$ est local, on voit que
$\mathcal{F}_x \cong \mathcal{O}_{X, x}^{\oplus r}$ pour un certain $r$, voir le
lemme \ref{algebra-lemma-finite-projective} d'Algèbre.
D'après le lemme \ref{modules-lemma-finite-presentation-stalk-free},
$\mathcal{F}$ est libre de rang $r$ dans un voisinage ouvert de $x$.
(Remarquons que $\mathcal{F}$ est de présentation finie comme facteur direct de
$\mathcal{H}$.)
\end{proof}





\section{Applications bilinéaires}
\label{modules-section-bilinear}


\noindent
Soit $(X, \mathcal{O}_X)$ un espace annelé.
Soient $\mathcal{F}$, $\mathcal{G}$ et $\mathcal{H}$ des $\mathcal{O}_X$-modules.
Une {\it application bilinéaire} $f : \mathcal{F} \times \mathcal{G} \to \mathcal{H}$
de faisceaux de $\mathcal{O}_X$-modules est une application de faisceaux d'ensembles comme indiqué
telle que, pour tout ouvert $U \subset X$, l'application induite
$$
\mathcal{F}(U) \times \mathcal{G}(U) \to \mathcal{H}(U)
$$
soit une application bilinéaire de $\mathcal{O}_X(U)$-modules. De manière équivalente, on peut demander
que certains diagrammes d'applications de faisceaux d'ensembles commutent, en imitant les
axiomes usuels des applications bilinéaires de modules. Par exemple, l'axiome
$f(x + y, z) = f(x, z) + f(y, z)$ est représenté par la commutativité
du diagramme
$$
\xymatrix{
\mathcal{F} \times \mathcal{F} \times \mathcal{G}
\ar[rrr]_{(f \circ \text{pr}_{13}, f \circ \text{pr}_{23})}
\ar[d]_{(+ \circ \text{pr}_{12}, \text{pr}_3)} & & &
\mathcal{H} \times \mathcal{H}
\ar[d]^{+} \\
\mathcal{F} \times \mathcal{G} \ar[rrr]^f & & &
\mathcal{H}
}
$$
Une autre caractérisation est la suivante : si
$f : \mathcal{F} \times \mathcal{G} \to \mathcal{H}$
est une application de faisceaux d'ensembles et
induit une application bilinéaire de modules sur les germes en tout point de $X$, alors
$f$ est une application bilinéaire de faisceaux de modules. En effet, on peut vérifier
l'égalité de sections locales sur les germes.

\medskip\noindent
Notons $\text{Mor}( - , - )$ les morphismes dans la catégorie des faisceaux
d'ensembles sur $X$. Voici une autre caractérisation d'une application bilinéaire :
une application de faisceaux d'ensembles
$f : \mathcal{F} \times \mathcal{G} \to \mathcal{H}$
est bilinéaire si, pour tout faisceau d'ensembles $\mathcal{S}$, la règle
$$
\text{Mor}(\mathcal{S}, \mathcal{F}) \times
\text{Mor}(\mathcal{S}, \mathcal{G}) \to
\text{Mor}(\mathcal{S}, \mathcal{H}),\quad
(a, b) \mapsto f \circ (a \times b)
$$
est une application bilinéaire de modules sur l'anneau
$\text{Mor}(\mathcal{S}, \mathcal{O}_X)$.
Nous n'adoptons généralement pas ce point de vue, car il est plus facile de raisonner sur les
ensembles de sections locales, et les deux points de vue sont manifestement équivalents.

\medskip\noindent
Enfin, voici encore une autre formulation de la définition :
$\mathcal{O}_X$ est un objet en anneaux dans la catégorie des faisceaux d'ensembles,
et $\mathcal{F}$, $\mathcal{G}$, $\mathcal{H}$ sont des objets en modules
sur cet anneau. On peut alors définir une application bilinéaire entre objets en modules
sur un objet en anneaux dans toute catégorie.
Pour formuler ce qu'est un objet en anneaux, ce qu'est un objet en modules
sur un objet en anneaux, et ce qu'est une application bilinéaire entre de tels objets dans une
catégorie, il est commode (mais pas strictement nécessaire)
de supposer que la catégorie possède des produits finis ;
c'est le cas de la catégorie des faisceaux d'ensembles.







\section{Produit tensoriel}
\label{modules-section-tensor-product}

\noindent
Nous avons déjà brièvement abordé le produit tensoriel dans le cadre
du changement d'anneaux dans Faisceaux, sections
\ref{sheaves-section-presheaves-modules} et
\ref{sheaves-section-sheafification-presheaves-modules}.
Généralisons cette discussion aux produits tensoriels de modules.

\medskip\noindent
Soit $(X, \mathcal{O}_X)$ un espace annelé et soient
$\mathcal{F}$ et $\mathcal{G}$ des $\mathcal{O}_X$-modules.
Nous définissons d'abord le {\it préfaisceau produit tensoriel}
$$
\mathcal{F} \otimes_{p, \mathcal{O}_X} \mathcal{G}
$$
comme la règle qui associe à tout ouvert $U \subset X$ le
$\mathcal{O}_X(U)$-module
$\mathcal{F}(U) \otimes_{\mathcal{O}_X(U)} \mathcal{G}(U)$.
Cela étant défini, nous appelons {\it faisceau produit tensoriel}
le faisceau associé au préfaisceau précédent :
$$
\mathcal{F} \otimes_{\mathcal{O}_X} \mathcal{G}
=
(\mathcal{F} \otimes_{p, \mathcal{O}_X} \mathcal{G})^\#
$$
On peut le caractériser comme le faisceau de
$\mathcal{O}_X$-modules tel que, pour tout troisième
faisceau de $\mathcal{O}_X$-modules $\mathcal{H}$,
on ait
$$
\Hom_{\mathcal{O}_X}
(\mathcal{F} \otimes_{\mathcal{O}_X} \mathcal{G}, \mathcal{H})
=
\text{Bilin}_{\mathcal{O}_X}(\mathcal{F} \times \mathcal{G}, \mathcal{H}).
$$
Ici, le membre de droite désigne l'ensemble des applications bilinéaires de faisceaux
de $\mathcal{O}_X$-modules définies dans la section \ref{modules-section-bilinear}.

\medskip\noindent
Le produit tensoriel de modules $M, N$ sur un anneau $R$ est symétrique,
à savoir $M \otimes_R N = N \otimes_R M$ ; il en va donc de même pour
les produits tensoriels de faisceaux de modules, c'est-à-dire que l'on a
$$
\mathcal{F} \otimes_{\mathcal{O}_X} \mathcal{G}
=
\mathcal{G} \otimes_{\mathcal{O}_X} \mathcal{F}
$$
fonctoriellement en $\mathcal{F}$, $\mathcal{G}$.
Puisque le produit tensoriel de modules est associatif, nous
obtenons aussi des isomorphismes canoniques fonctoriels
$$
(\mathcal{F} \otimes_{\mathcal{O}_X} \mathcal{G})
\otimes_{\mathcal{O}_X} \mathcal{H}
=
\mathcal{F} \otimes_{\mathcal{O}_X}
(\mathcal{G} \otimes_{\mathcal{O}_X} \mathcal{H})
$$
en $\mathcal{F}$, $\mathcal{G}$ et $\mathcal{H}$.

\begin{lemma}
\label{modules-lemma-stalk-tensor-product}
Soit $(X, \mathcal{O}_X)$ un espace annelé.
Soient $\mathcal{F}$, $\mathcal{G}$ des $\mathcal{O}_X$-modules.
Soit $x \in X$. Il existe un isomorphisme canonique
de $\mathcal{O}_{X, x}$-modules
$$
(\mathcal{F} \otimes_{\mathcal{O}_X} \mathcal{G})_x
=
\mathcal{F}_x \otimes_{\mathcal{O}_{X, x}} \mathcal{G}_x
$$
fonctoriel en $\mathcal{F}$ et $\mathcal{G}$.
\end{lemma}

\begin{proof}
La démonstration est omise.
\end{proof}

\begin{lemma}
\label{modules-lemma-tensor-product-sheafification}
Soit $(X, \mathcal{O}_X)$ un espace annelé.
Soient $\mathcal{F}'$, $\mathcal{G}'$ des préfaisceaux de $\mathcal{O}_X$-modules
dont les faisceaux associés sont $\mathcal{F}$, $\mathcal{G}$. Alors
$\mathcal{F} \otimes_{\mathcal{O}_X} \mathcal{G} =
(\mathcal{F}' \otimes_{p, \mathcal{O}_X} \mathcal{G}')^\#$.
\end{lemma}

\begin{proof}
La démonstration est omise.
\end{proof}


\begin{lemma}
\label{modules-lemma-tensor-product-exact}
Soit $(X, \mathcal{O}_X)$ un espace annelé.
Soit $\mathcal{G}$ un $\mathcal{O}_X$-module.
Si
$\mathcal{F}_1
\to \mathcal{F}_2
\to \mathcal{F}_3
\to 0$
est une suite exacte de $\mathcal{O}_X$-modules, alors
la suite induite
$$
\mathcal{F}_1 \otimes_{\mathcal{O}_X} \mathcal{G} \to
\mathcal{F}_2 \otimes_{\mathcal{O}_X} \mathcal{G} \to
\mathcal{F}_3 \otimes_{\mathcal{O}_X} \mathcal{G} \to
0
$$
est exacte.
\end{lemma}

\begin{proof}
Cela résulte du fait que l'exactitude peut se vérifier sur les germes
(lemme \ref{modules-lemma-abelian}), de la description des germes
(lemme \ref{modules-lemma-stalk-tensor-product}) et du résultat correspondant
pour les produits tensoriels de modules
(Algèbre, lemme \ref{algebra-lemma-tensor-product-exact}).
\end{proof}

\begin{lemma}
\label{modules-lemma-tensor-product-pullback}
Soit $f : (X, \mathcal{O}_X) \to (Y, \mathcal{O}_Y)$
un morphisme d'espaces annelés. Soient $\mathcal{F}$, $\mathcal{G}$
des $\mathcal{O}_Y$-modules. Alors
$f^*(\mathcal{F} \otimes_{\mathcal{O}_Y} \mathcal{G})
= f^*\mathcal{F} \otimes_{\mathcal{O}_X} f^*\mathcal{G}$
fonctoriellement en $\mathcal{F}$, $\mathcal{G}$.
\end{lemma}

\begin{proof}
La démonstration est omise.
\end{proof}

\begin{lemma}
\label{modules-lemma-tensor-commute-colimits}
Soit $(X, \mathcal{O}_X)$ un espace annelé.
Pour tout $\mathcal{O}_X$-module $\mathcal{F}$, le foncteur
$$
\textit{Mod}(\mathcal{O}_X) \longrightarrow \textit{Mod}(\mathcal{O}_X)
, \quad
\mathcal{G} \longmapsto \mathcal{F} \otimes_{\mathcal{O}_X} \mathcal{G}
$$
commute aux colimites quelconques.
\end{lemma}

\begin{proof}
Soit $I$ un ensemble préordonné et soit $\{\mathcal{G}_i\}$
un système indexé par $I$. Posons $\mathcal{G} = \colim_i \mathcal{G}_i$.
Rappelons que $\mathcal{G}$ est le faisceau associé au préfaisceau
$\mathcal{G}' : U \mapsto \colim_i \mathcal{G}_i(U)$ ; voir
Faisceaux, section \ref{sheaves-section-limits-sheaves}.
D'après le
lemme \ref{modules-lemma-tensor-product-sheafification},
le produit tensoriel $\mathcal{F} \otimes_{\mathcal{O}_X} \mathcal{G}$
est le faisceau associé au préfaisceau
$$
U \longmapsto
\mathcal{F}(U) \otimes_{\mathcal{O}_X(U)} \colim_i \mathcal{G}_i(U) =
\colim_i \mathcal{F}(U) \otimes_{\mathcal{O}_X(U)} \mathcal{G}_i(U)
$$
où l'égalité est donnée par
Algèbre, lemme \ref{algebra-lemma-tensor-products-commute-with-limits}.
Le lemme résulte donc de la description des colimites dans
$\textit{Mod}(\mathcal{O}_X)$ ; voir le lemme \ref{modules-lemma-limits-colimits}.
\end{proof}

\begin{lemma}
\label{modules-lemma-tensor-product-permanence}
Soit $(X, \mathcal{O}_X)$ un espace annelé.
Soient $\mathcal{F}$, $\mathcal{G}$ des $\mathcal{O}_X$-modules.
\begin{enumerate}
\item Si $\mathcal{F}$, $\mathcal{G}$ sont localement engendrés
par leurs sections, il en va de même pour $\mathcal{F} \otimes_{\mathcal{O}_X} \mathcal{G}$.
\item Si $\mathcal{F}$, $\mathcal{G}$ sont de type fini,
il en va de même pour $\mathcal{F} \otimes_{\mathcal{O}_X} \mathcal{G}$.
\item Si $\mathcal{F}$, $\mathcal{G}$ sont quasi-cohérents,
il en va de même pour $\mathcal{F} \otimes_{\mathcal{O}_X} \mathcal{G}$.
\item Si $\mathcal{F}$, $\mathcal{G}$ sont de présentation finie,
il en va de même pour $\mathcal{F} \otimes_{\mathcal{O}_X} \mathcal{G}$.
\item Si $\mathcal{F}$ est de présentation finie et $\mathcal{G}$ est cohérent,
alors $\mathcal{F} \otimes_{\mathcal{O}_X} \mathcal{G}$ est cohérent.
\item Si $\mathcal{F}$, $\mathcal{G}$ sont cohérents,
il en va de même pour $\mathcal{F} \otimes_{\mathcal{O}_X} \mathcal{G}$.
\item Si $\mathcal{F}$, $\mathcal{G}$ sont localement libres,
il en va de même pour $\mathcal{F} \otimes_{\mathcal{O}_X} \mathcal{G}$.
\end{enumerate}
\end{lemma}

\begin{proof}
Montrons d'abord que le produit tensoriel de
$\mathcal{O}_X$-modules localement libres est localement libre. Il suffit de montrer
que
$(\bigoplus_{i \in I} \mathcal{O}_X) \otimes_{\mathcal{O}_X}
(\bigoplus_{j \in J} \mathcal{O}_X) \cong
\bigoplus_{(i, j) \in I \times J} \mathcal{O}_X$.
Le faisceau $\bigoplus_{i \in I} \mathcal{O}_X$ est le faisceau associé
au préfaisceau $U \mapsto \bigoplus_{i \in I} \mathcal{O}_X(U)$.
Par conséquent, le produit tensoriel est le faisceau associé
au préfaisceau
$$
U \longmapsto
(\bigoplus\nolimits_{i \in I} \mathcal{O}_X(U))
\otimes_{\mathcal{O}_X(U)}
(\bigoplus\nolimits_{j \in J} \mathcal{O}_X(U)).
$$
On en déduit le résultat voulu, puisque pour tout anneau $R$ on a
$(\bigoplus_{i \in I} R) \otimes_R (\bigoplus_{j \in J} R) =
\bigoplus_{(i, j) \in I \times J} R$.

\medskip\noindent
Si $\mathcal{F}_2 \to \mathcal{F}_1 \to \mathcal{F} \to 0$
est exacte, alors, d'après le lemme \ref{modules-lemma-tensor-product-exact},
le complexe
$\mathcal{F}_2 \otimes_{\mathcal{O}_X} \mathcal{G} \to
\mathcal{F}_1 \otimes_{\mathcal{O}_X} \mathcal{G} \to
\mathcal{F} \otimes_{\mathcal{O}_X} \mathcal{G} \to 0$
est exact. Cela permet de démontrer (5). En effet, dans ce cas, il existe
localement une telle suite exacte où les $\mathcal{F}_i$, $i = 1, 2$, sont
libres de type fini. Les deux termes
$\mathcal{F}_2 \otimes_{\mathcal{O}_X} \mathcal{G}$ et
$\mathcal{F}_1 \otimes_{\mathcal{O}_X} \mathcal{G}$
sont donc isomorphes à des sommes directes finies de $\mathcal{G}$ (par exemple
d'après le lemme \ref{modules-lemma-tensor-commute-colimits}).
Comme les sommes directes finies sont des faisceaux cohérents, ces termes sont cohérents,
ainsi que le conoyau de l'application ; voir le lemme \ref{modules-lemma-coherent-abelian}.

\medskip\noindent
Et si, de plus,
$\mathcal{G}_2 \to \mathcal{G}_1 \to \mathcal{G} \to 0$
est exacte, alors on voit que
$$
\mathcal{F}_2 \otimes_{\mathcal{O}_X} \mathcal{G}_1
\oplus
\mathcal{F}_1 \otimes_{\mathcal{O}_X} \mathcal{G}_2
\to
\mathcal{F}_1 \otimes_{\mathcal{O}_X} \mathcal{G}_1
\to
\mathcal{F} \otimes_{\mathcal{O}_X} \mathcal{G}
\to 0
$$
est exacte. Cela permet, par exemple, de démontrer (3).
En effet, l'hypothèse signifie que l'on peut trouver localement des présentations
comme ci-dessus où les $\mathcal{F}_i$ et les $\mathcal{G}_i$ sont des
$\mathcal{O}_X$-modules libres. La présentation affichée est donc aussi
une présentation du produit tensoriel par des faisceaux libres.

\medskip\noindent
La démonstration des autres assertions est omise.
\end{proof}





\section{Modules plats}
\label{modules-section-flat}

\noindent
On peut définir les modules plats exactement comme dans le cas des modules sur un anneau.

\begin{definition}
\label{modules-definition-flat}
Soit $(X, \mathcal{O}_X)$ un espace annelé.
Un $\mathcal{O}_X$-module $\mathcal{F}$ est {\it plat} si le foncteur
$$
\textit{Mod}(\mathcal{O}_X)
\longrightarrow
\textit{Mod}(\mathcal{O}_X), \quad
\mathcal{G} \mapsto \mathcal{G} \otimes_{\mathcal{O}_X} \mathcal{F}
$$
est exact.
\end{definition}

\noindent
On peut caractériser la platitude au moyen des germes.

\begin{lemma}
\label{modules-lemma-flat-stalks-flat}
Soit $(X, \mathcal{O}_X)$ un espace annelé.
Un $\mathcal{O}_X$-module $\mathcal{F}$ est plat
si et seulement si le germe $\mathcal{F}_x$ est un
$\mathcal{O}_{X, x}$-module plat pour tout $x \in X$.
\end{lemma}

\begin{proof}
Supposons que $\mathcal{F}_x$ soit un $\mathcal{O}_{X, x}$-module plat pour tout
$x \in X$. Dans ce cas, si $\mathcal{G} \to \mathcal{H} \to \mathcal{K}$
est exacte, alors
$\mathcal{G} \otimes_{\mathcal{O}_X} \mathcal{F} \to
\mathcal{H} \otimes_{\mathcal{O}_X} \mathcal{F} \to
\mathcal{K} \otimes_{\mathcal{O}_X} \mathcal{F}$
est également exacte, car on peut vérifier l'exactitude sur les germes et
le produit tensoriel commute à la prise des germes ; voir le
lemme \ref{modules-lemma-stalk-tensor-product}.
Réciproquement, supposons que $\mathcal{F}$ soit plat, et soit $x \in X$.
Considérons les faisceaux gratte-ciel $i_{x, *} M$, où $M$ est un
$\mathcal{O}_{X, x}$-module. Remarquons que
$$
M \otimes_{\mathcal{O}_{X, x}} \mathcal{F}_x =
\left(i_{x, *} M \otimes_{\mathcal{O}_X} \mathcal{F}\right)_x
$$
de nouveau d'après le
lemme \ref{modules-lemma-stalk-tensor-product}.
Puisque $i_{x, *}$ est exact, la platitude de $\mathcal{F}$ implique que
$M \mapsto M \otimes_{\mathcal{O}_{X, x}} \mathcal{F}_x$
est exact. Ainsi $\mathcal{F}_x$ est un $\mathcal{O}_{X, x}$-module plat.
\end{proof}

\noindent
La définition suivante a donc un sens.

\begin{definition}
\label{modules-definition-flat-at-point}
Soit $(X, \mathcal{O}_X)$ un espace annelé. Soit $x \in X$.
Un $\mathcal{O}_X$-module $\mathcal{F}$ est
{\it plat en $x$} si $\mathcal{F}_x$ est un
$\mathcal{O}_{X, x}$-module plat.
\end{definition}

\noindent
Ainsi, $\mathcal{F}$ est un $\mathcal{O}_X$-module plat
si et seulement s'il est plat en tout point.

\begin{lemma}
\label{modules-lemma-base-change-flat}
Soit $f : (X, \mathcal{O}_X) \to (Y, \mathcal{O}_Y)$
un morphisme d'espaces annelés. Si $\mathcal{G}$ est un
$\mathcal{O}_Y$-module plat, alors $f^*\mathcal{G}$ est un
$\mathcal{O}_X$-module plat.
\end{lemma}

\begin{proof}
Combiner le lemme \ref{modules-lemma-flat-stalks-flat} avec
Faisceaux, lemme \ref{sheaves-lemma-stalk-pullback-modules}, et
Algèbre, lemme \ref{algebra-lemma-flat-base-change}.
\end{proof}

\begin{lemma}
\label{modules-lemma-colimits-flat}
Soit $(X, \mathcal{O}_X)$ un espace annelé.
Une colimite filtrante de $\mathcal{O}_X$-modules plats est plate.
Une somme directe de $\mathcal{O}_X$-modules plats est plate.
\end{lemma}

\begin{proof}
Cela résulte du
lemme \ref{modules-lemma-tensor-commute-colimits}, du
lemme \ref{modules-lemma-stalk-tensor-product}, d'
Algèbre, lemme \ref{algebra-lemma-directed-colimit-exact},
et du fait que l'on peut vérifier l'exactitude sur les germes.
\end{proof}

\begin{lemma}
\label{modules-lemma-j-shriek-flat}
Soit $(X, \mathcal{O}_X)$ un espace annelé.
Soit $U \subset X$ un ouvert. Le faisceau $j_{U!}\mathcal{O}_U$
est un faisceau plat de $\mathcal{O}_X$-modules.
\end{lemma}

\begin{proof}
Les germes de $j_{U!}\mathcal{O}_U$ sont soit nuls, soit égaux
à $\mathcal{O}_{X, x}$. Appliquer le
lemme \ref{modules-lemma-flat-stalks-flat}.
\end{proof}

\begin{lemma}
\label{modules-lemma-module-quotient-flat}
Soit $(X, \mathcal{O}_X)$ un espace annelé.
\begin{enumerate}
\item Tout faisceau de $\mathcal{O}_X$-modules est quotient d'une
somme directe $\bigoplus j_{U_i!}\mathcal{O}_{U_i}$.
\item Tout $\mathcal{O}_X$-module est quotient d'un
$\mathcal{O}_X$-module plat.
\end{enumerate}
\end{lemma}

\begin{proof}
Soit $\mathcal{F}$ un $\mathcal{O}_X$-module.
Pour tout ouvert $U \subset X$ et tout
$s \in \mathcal{F}(U)$, on obtient un morphisme
$j_{U!}\mathcal{O}_U \to \mathcal{F}$, à savoir l'adjoint du
morphisme $\mathcal{O}_U \to \mathcal{F}|_U$, $1 \mapsto s$.
L'application
$$
\bigoplus\nolimits_{(U, s)} j_{U!}\mathcal{O}_U
\longrightarrow
\mathcal{F}
$$
est manifestement surjective, et sa source est plate d'après les lemmes
\ref{modules-lemma-colimits-flat} et \ref{modules-lemma-j-shriek-flat}.
\end{proof}

\begin{lemma}
\label{modules-lemma-flat-tor-zero}
Soit $(X, \mathcal{O}_X)$ un espace annelé.
Soit
$$
0 \to \mathcal{F}'' \to \mathcal{F}' \to \mathcal{F} \to 0
$$
une suite exacte courte de $\mathcal{O}_X$-modules.
Supposons $\mathcal{F}$ plat. Alors, pour tout $\mathcal{O}_X$-module
$\mathcal{G}$, la suite
$$
0 \to
\mathcal{F}'' \otimes_\mathcal{O} \mathcal{G} \to
\mathcal{F}' \otimes_\mathcal{O} \mathcal{G} \to
\mathcal{F} \otimes_\mathcal{O} \mathcal{G} \to 0
$$
est exacte.
\end{lemma}

\begin{proof}
Comme $\mathcal{F}_x$ est un $\mathcal{O}_{X, x}$-module plat
pour tout $x \in X$ et que l'exactitude se vérifie sur les germes, le résultat
découle d'
Algèbre, lemme \ref{algebra-lemma-flat-tor-zero}.
\end{proof}

\begin{lemma}
\label{modules-lemma-flat-ses}
\begin{slogan}
Les noyaux d'épimorphismes et les extensions de faisceaux plats de modules sur
un espace annelé sont encore plats.
\end{slogan}
Soit $(X, \mathcal{O}_X)$ un espace annelé.
Soit
$$
0 \to
\mathcal{F}_2 \to
\mathcal{F}_1 \to
\mathcal{F}_0 \to 0
$$
une suite exacte courte de $\mathcal{O}_X$-modules.
\begin{enumerate}
\item Si $\mathcal{F}_2$ et $\mathcal{F}_0$ sont plats, il en va de même pour
$\mathcal{F}_1$.
\item Si $\mathcal{F}_1$ et $\mathcal{F}_0$ sont plats, il en va de même pour
$\mathcal{F}_2$.
\end{enumerate}
\end{lemma}

\begin{proof}
Comme l'exactitude et la platitude se vérifient sur les germes,
cela résulte d'
Algèbre, lemme \ref{algebra-lemma-flat-ses}.
\end{proof}

\begin{lemma}
\label{modules-lemma-flat-resolution-of-flat}
Soit $(X, \mathcal{O}_X)$ un espace annelé.
Soit
$$
\ldots \to
\mathcal{F}_2 \to
\mathcal{F}_1 \to
\mathcal{F}_0 \to
\mathcal{Q} \to 0
$$
un complexe exact de $\mathcal{O}_X$-modules.
Si $\mathcal{Q}$ et tous les $\mathcal{F}_i$ sont des $\mathcal{O}_X$-modules plats,
alors, pour tout $\mathcal{O}_X$-module $\mathcal{G}$, le complexe
$$
\ldots \to
\mathcal{F}_2 \otimes_{\mathcal{O}_X} \mathcal{G} \to
\mathcal{F}_1 \otimes_{\mathcal{O}_X} \mathcal{G} \to
\mathcal{F}_0 \otimes_{\mathcal{O}_X} \mathcal{G} \to
\mathcal{Q} \otimes_{\mathcal{O}_X} \mathcal{G} \to 0
$$
est également exact.
\end{lemma}

\begin{proof}
Cela résulte du lemme \ref{modules-lemma-flat-tor-zero}, en décomposant le complexe
en suites exactes courtes et en utilisant le lemme \ref{modules-lemma-flat-ses} pour
montrer par récurrence que $\Im(\mathcal{F}_{i + 1} \to \mathcal{F}_i)$
est plat.
\end{proof}

\noindent
Le lemme suivant donne une implication du critère équationnel de
platitude (Algèbre, lemme \ref{algebra-lemma-flat-eq}).

\begin{lemma}
\label{modules-lemma-flat-eq}
Soit $(X, \mathcal{O}_X)$ un espace annelé. Soit $\mathcal{F}$ un
$\mathcal{O}_X$-module plat. Soit $U \subset X$ un ouvert et soit
$$
\mathcal{O}_U \xrightarrow{(f_1, \ldots, f_n)}
\mathcal{O}_U^{\oplus n} \xrightarrow{(s_1, \ldots, s_n)}
\mathcal{F}|_U
$$
un complexe de $\mathcal{O}_U$-modules. Pour tout $x \in U$, il
existe un voisinage ouvert $V \subset U$ de $x$ et une factorisation
$$
\mathcal{O}_V^{\oplus n}
\xrightarrow{A}
\mathcal{O}_V^{\oplus m} \xrightarrow{(t_1, \ldots, t_m)}
\mathcal{F}|_V
$$
de $(s_1, \ldots, s_n)|_V$ telle que $A \circ (f_1, \ldots, f_n)|_V = 0$.
\end{lemma}

\begin{proof}
Soit $\mathcal{I} \subset \mathcal{O}_U$ le faisceau d'idéaux
engendré par $f_1, \ldots, f_n$. Alors $\sum f_i \otimes s_i$ est
une section de $\mathcal{I} \otimes_{\mathcal{O}_U} \mathcal{F}|_U$
dont l'image dans $\mathcal{F}|_U$ est nulle. Comme $\mathcal{F}|_U$ est plat,
l'application
$\mathcal{I} \otimes_{\mathcal{O}_U} \mathcal{F}|_U \to \mathcal{F}|_U$
est injective. Puisque $\mathcal{I} \otimes_{\mathcal{O}_U} \mathcal{F}|_U$
est le faisceau associé au préfaisceau produit tensoriel, il existe
un voisinage ouvert $V \subset U$ de $x$ tel
que $\sum f_i|_V \otimes s_i|_V$ soit nul dans
$\mathcal{I}(V) \otimes_{\mathcal{O}(V)} \mathcal{F}(V)$.
En explicitant les définitions à l'aide d'Algèbre, lemme \ref{algebra-lemma-relations},
on trouve $t_1, \ldots, t_m \in \mathcal{F}(V)$ et $a_{ij} \in \mathcal{O}(V)$
tels que $\sum a_{ij}f_i|_V = 0$ et $s_i|_V = \sum a_{ij}t_j$.
\end{proof}














\section{Duaux}
\label{modules-section-duals}

\noindent
Soit $(X, \mathcal{O}_X)$ un espace annelé. La catégorie des
$\mathcal{O}_X$-modules, munie du produit tensoriel
construit dans la section \ref{modules-section-tensor-product},
est une catégorie monoïdale symétrique. Pour un $\mathcal{O}_X$-module
$\mathcal{F}$, les conditions suivantes sont équivalentes :
\begin{enumerate}
\item $\mathcal{F}$ admet un dual à gauche dans la catégorie monoïdale
des $\mathcal{O}_X$-modules ;
\item $\mathcal{F}$ est localement facteur direct d'un
$\mathcal{O}_X$-module libre de type fini ;
\item $\mathcal{F}$ est de présentation finie et plat comme
$\mathcal{O}_X$-module.
\end{enumerate}
Cela est démontré dans l'exemple \ref{modules-example-dual} et les
lemmes \ref{modules-lemma-left-dual-module} et
\ref{modules-lemma-flat-locally-finite-presentation} de cette section.

\begin{example}
\label{modules-example-dual}
Soit $(X, \mathcal{O}_X)$ un espace annelé. Soit $\mathcal{F}$
un $\mathcal{O}_X$-module qui est localement facteur direct d'un
$\mathcal{O}_X$-module libre de type fini. Alors l'application
$$
\mathcal{F} \otimes_{\mathcal{O}_X}
\SheafHom_{\mathcal{O}_X}(\mathcal{F}, \mathcal{O}_X)
\longrightarrow
\SheafHom_{\mathcal{O}_X}(\mathcal{F}, \mathcal{F})
$$
est un isomorphisme. En effet, la question est locale, l'assertion est vraie
si $\mathcal{F}$ est libre de type fini, et elle vaut pour tout facteur direct
d'un module pour lequel elle est vraie. Notons
$$
\eta :
\mathcal{O}_X
\longrightarrow
\mathcal{F} \otimes_{\mathcal{O}_X}
\SheafHom_{\mathcal{O}_X}(\mathcal{F}, \mathcal{O}_X)
$$
l'application qui envoie $1$ sur la section correspondant à
$\text{id}_\mathcal{F}$ par l'isomorphisme ci-dessus. Notons
$$
\epsilon : 
\SheafHom_{\mathcal{O}_X}(\mathcal{F}, \mathcal{O}_X)
\otimes_{\mathcal{O}_X} \mathcal{F}
\longrightarrow
\mathcal{O}_X
$$
l'application d'évaluation. Alors
$\SheafHom_{\mathcal{O}_X}(\mathcal{F}, \mathcal{O}_X), \eta, \epsilon$
est un dual à gauche de $\mathcal{F}$ au sens de
Catégories, définition \ref{categories-definition-dual}.
Nous omettons la vérification des égalités
$(1 \otimes \epsilon) \circ (\eta \otimes 1) = \text{id}_\mathcal{F}$
et
$(\epsilon \otimes 1) \circ (1 \otimes \eta) =
\text{id}_{\SheafHom_{\mathcal{O}_X}(\mathcal{F}, \mathcal{O}_X)}$.
\end{example}

\begin{lemma}
\label{modules-lemma-left-dual-module}
Soit $(X, \mathcal{O}_X)$ un espace annelé. Soit $\mathcal{F}$ un
$\mathcal{O}_X$-module. Soit $\mathcal{G}, \eta, \epsilon$
un dual à gauche de $\mathcal{F}$ dans la catégorie monoïdale des
$\mathcal{O}_X$-modules ; voir
Catégories, définition \ref{categories-definition-dual}. Alors :
\begin{enumerate}
\item $\mathcal{F}$ est localement facteur direct d'un
$\mathcal{O}_X$-module libre de type fini ;
\item l'application
$e : \SheafHom_{\mathcal{O}_X}(\mathcal{F}, \mathcal{O}_X) \to \mathcal{G}$
qui envoie une section locale $\lambda$ sur $(\lambda \otimes 1)(\eta)$
est un isomorphisme ;
\item on a $\epsilon(f, g) = e^{-1}(g)(f)$ pour des sections locales
$f$ et $g$ de $\mathcal{F}$ et $\mathcal{G}$.
\end{enumerate}
\end{lemma}

\begin{proof}
Les hypothèses signifient que
$$
\mathcal{F} \xrightarrow{\eta \otimes 1}
\mathcal{F} \otimes_{\mathcal{O}_X} \mathcal{G}
\otimes_{\mathcal{O}_X} \mathcal{F}
\xrightarrow{1 \otimes \epsilon} \mathcal{F}
\quad\text{et}\quad
\mathcal{G} \xrightarrow{1 \otimes \eta}
\mathcal{G} \otimes_{\mathcal{O}_X} \mathcal{F}
\otimes_{\mathcal{O}_X} \mathcal{G}
\xrightarrow{\epsilon \otimes 1} \mathcal{G}
$$
sont les applications identiques. Soit $x \in X$. On peut trouver un voisinage ouvert
$U$ de $x$, un nombre fini de sections $f_1, \ldots, f_n$
et $g_1, \ldots, g_n$ de
$\mathcal{F}$ et $\mathcal{G}$ sur $U$ telles que
$\eta(1) = \sum f_i g_i$. Notons
$$
\mathcal{O}_U^{\oplus n} \to \mathcal{F}|_U
$$
l'application qui envoie le $i$-ième vecteur de base sur $f_i$. On peut alors
factoriser l'application $\eta|_U$ par une application
$\tilde \eta : \mathcal{O}_U \to
\mathcal{O}_U^{\oplus n} \otimes_{\mathcal{O}_U} \mathcal{G}|_U$.
On obtient un diagramme commutatif
$$
\xymatrix{
\mathcal{F}|_U
\ar[rr]_-{\eta \otimes 1} \ar[rrd]_-{\tilde \eta \otimes 1} & &
\mathcal{F}|_U \otimes \mathcal{G}|_U \otimes \mathcal{F}|_U
\ar[r]_-{1 \otimes \epsilon} &
\mathcal{F}|_U \\
& &
\mathcal{O}_U^{\oplus n} \otimes \mathcal{G}|_U \otimes \mathcal{F}|_U
\ar[u] \ar[r]^-{1 \otimes \epsilon} &
\mathcal{O}_U^{\oplus n} \ar[u]
}
$$
Cela montre que l'identité de $\mathcal{F}$, localement sur $X$,
se factorise par un module libre de type fini. Ceci démontre (1). L'assertion (2) résulte de
Catégories, lemme \ref{categories-lemma-left-dual}, et de sa démonstration.
L'assertion (3) résulte de la première égalité de la démonstration.
On peut aussi déduire (2) et (3) de l'unicité des duaux à gauche
(Catégories, remarque \ref{categories-remark-left-dual-adjoint})
et de la construction du dual à gauche dans
l'exemple \ref{modules-example-dual}.
\end{proof}

\begin{lemma}
\label{modules-lemma-flat-locally-finite-presentation}
Soit $(X, \mathcal{O}_X)$ un espace annelé. Soit $\mathcal{F}$ un
$\mathcal{O}_X$-module plat de présentation finie. Alors $\mathcal{F}$ est
localement facteur direct d'un $\mathcal{O}_X$-module libre de type fini.
\end{lemma}

\begin{proof}
Après avoir remplacé $X$ par les membres d'un recouvrement ouvert, on peut
supposer qu'il existe une présentation
$$
\mathcal{O}_X^{\oplus r} \to
\mathcal{O}_X^{\oplus n} \to \mathcal{F} \to 0
$$
Soit $x \in X$. D'après le lemme \ref{modules-lemma-flat-eq},
on peut, après avoir restreint $X$ à un voisinage
ouvert de $x$, supposer qu'il existe une factorisation
$$
\mathcal{O}_X^{\oplus n} \to
\mathcal{O}_X^{\oplus n_1} \to \mathcal{F}
$$
telle que la composée
$\mathcal{O}_X^{\oplus r} \to \mathcal{O}_X^{\oplus n} \to
\mathcal{O}_X^{\oplus n_1}$
annule le premier facteur de $\mathcal{O}_X^{\oplus r}$.
En répétant cet argument encore $r - 1$ fois, on obtient une factorisation
$$
\mathcal{O}_X^{\oplus n} \to
\mathcal{O}_X^{\oplus n_r} \to \mathcal{F}
$$
telle que la composée
$\mathcal{O}_X^{\oplus r} \to \mathcal{O}_X^{\oplus n}
\to \mathcal{O}_X^{\oplus n_r}$ est nulle.
Cela signifie que la surjection $\mathcal{O}_X^{\oplus n_r} \to \mathcal{F}$
admet une section, ce qui conclut.
\end{proof}







\section{Faisceaux constructibles d'ensembles}
\label{modules-section-constructible}

\noindent
Soit $X$ un espace topologique. Étant donné un ensemble $S$, rappelons que $\underline{S}$
ou $\underline{S}_X$ désigne le faisceau constant de valeur $S$ ; voir
Faisceaux, définition \ref{sheaves-definition-constant-sheaf}.
Soit $U \subset X$ un ouvert d'un espace topologique $X$.
Nous noterons $j_U$ le morphisme d'inclusion et
$j_{U!} : \Sh(U) \to \Sh(X)$ le prolongement par l'ensemble vide
décrit dans Faisceaux, section \ref{sheaves-section-open-immersions}.

\begin{lemma}
\label{modules-lemma-surjection}
Soit $X$ un espace topologique. Soit $\mathcal{B}$ une base de la
topologie de $X$. Soit $\mathcal{F}$ un faisceau d'ensembles sur $X$.
Il existe un ensemble $I$ et, pour tout $i \in I$, un élément
$U_i \in \mathcal{B}$ et un ensemble fini $S_i$ tels qu'il existe
une surjection $\coprod_{i \in I} j_{U_i!}\underline{S_i} \to \mathcal{F}$.
\end{lemma}

\begin{proof}
Soit $S$ un singleton. Nous démontrerons le résultat avec $S_i = S$.
Pour tout $x \in X$ et tout élément $s \in \mathcal{F}_x$, on peut choisir
$U(x, s) \in \mathcal{B}$ et $s(x, s) \in \mathcal{F}(U(x, s))$
dont l'image est $s$ dans $\mathcal{F}_x$. D'après
Faisceaux, lemme \ref{sheaves-lemma-j-shriek},
la section $s(x, s)$
correspond à une application de faisceaux $j_{U(x, s)!}\underline{S} \to \mathcal{F}$.
Alors
$$
\coprod\nolimits_{(x, s)} j_{U(x, s)!}\underline{S} \to \mathcal{F}
$$
est surjective sur les germes, donc surjective.
\end{proof}

\begin{lemma}
\label{modules-lemma-filtered-colimit-constructibles}
Soit $X$ un espace topologique. Soit $\mathcal{B}$ une base de la
topologie de $X$, et supposons que tout $U \in \mathcal{B}$ soit quasi-compact.
Alors tout faisceau d'ensembles sur $X$ est une colimite filtrante de faisceaux de la forme
\begin{equation}
\label{modules-equation-towards-constructible-sets}
\text{Coégalisateur}\left(
\xymatrix{
\coprod\nolimits_{b = 1, \ldots, m} j_{V_b!}\underline{S_b}
\ar@<1ex>[r] \ar@<-1ex>[r] &
\coprod\nolimits_{a = 1, \ldots, n} j_{U_a!}\underline{S_a}
}
\right)
\end{equation}
où $U_a$ et $V_b$ appartiennent à $\mathcal{B}$, et où $S_a$ et $S_b$ sont des ensembles finis.
\end{lemma}

\begin{proof}
Par le lemme \ref{modules-lemma-surjection}, tout faisceau d'ensembles $\mathcal{F}$
est le but d'une surjection dont la source $\mathcal{F}_0$ est un coproduit
de faisceaux de la forme $j_{U!}\underline{S}$, avec $U \in \mathcal{B}$
et $S$ fini. En appliquant ceci à
$\mathcal{F}_0 \times_\mathcal{F} \mathcal{F}_0$
on trouve que $\mathcal{F}$ est le coégalisateur d'une paire d'applications
$$
\xymatrix{
\coprod\nolimits_{b \in B} j_{V_b!}\underline{S_b}
\ar@<1ex>[r] \ar@<-1ex>[r] &
\coprod\nolimits_{a \in A} j_{U_a!}\underline{S_a}
}
$$
pour certains ensembles d'indices $A$, $B$, avec $V_b$ et $U_a$ dans $\mathcal{B}$ et
$S_a$ et $S_b$ finis. Pour toute partie finie $B' \subset B$,
il existe une partie finie $A' \subset A$ telle que, par les deux applications, le coproduit
indexé par $b \in B'$ s'envoie dans le coproduit indexé par $a \in A'$.
En effet, on peut voir le membre de droite comme une colimite filtrante dont les
applications de transition sont injectives. Par conséquent, la prise des sections sur les ouverts
quasi-compacts $V_b$, $b \in B'$, commute à ce coproduit ;
voir Faisceaux, lemme \ref{sheaves-lemma-directed-colimits-sections}.
Ainsi notre faisceau est la colimite des conoyaux de ces applications
entre coproduits finis.
\end{proof}

\begin{lemma}
\label{modules-lemma-constructible-comes-from-finite}
Soit $X$ un espace topologique spectral. Soit $\mathcal{B}$
l'ensemble des ouverts quasi-compacts de $X$.
Soit $\mathcal{F}$ un faisceau d'ensembles comme dans
l'équation (\ref{modules-equation-towards-constructible-sets}).
Alors il existe une application spectrale continue $f : X \to Y$
vers un espace topologique sobre fini $Y$ et un faisceau
d'ensembles $\mathcal{G}$ sur $Y$, à germes finis,
tels que $f^{-1}\mathcal{G} \cong \mathcal{F}$.
\end{lemma}

\begin{proof}
On peut écrire $X = \lim X_i$ comme une limite projective
d'espaces sobres finis ; voir Topologie, lemme
\ref{topology-lemma-spectral-inverse-limit-finite-sober-spaces}.
Bien sûr, les applications de transition $X_{i'} \to X_i$ sont spectrales ; par conséquent,
d'après Topologie, lemme \ref{topology-lemma-directed-inverse-limit-spectral-spaces},
les applications $p_i : X \to X_i$ sont spectrales.
Pour un certain $i$, on peut trouver des ouverts $U_{a, i}$ et $V_{b, i}$
de $X_i$ dont les images inverses sont $U_a$ et $V_b$ ; voir
Topologie, lemme \ref{topology-lemma-descend-opens}.
Les deux applications
$$
\beta, \gamma :
\coprod\nolimits_{b \in B} j_{V_b!}\underline{S_b}
\longrightarrow
\coprod\nolimits_{a \in A} j_{U_a!}\underline{S_a}
$$
dont le coégalisateur est $\mathcal{F}$ correspondent par adjonction à deux familles
$$
\beta_b, \gamma_b :
S_b
\longrightarrow
\Gamma(V_b, \coprod\nolimits_{a \in A} j_{U_a!}\underline{S_a}), \quad
b \in B
$$
d'applications d'ensembles. Observons que
$p_i^{-1}(j_{U_{a, i}!}\underline{S_a}) = j_{U_a!}\underline{S_a}$
et $(X_{i'} \to X_i)^{-1}(j_{U_{a, i}!}\underline{S_a}) =
j_{U_{a, i'}!}\underline{S_a}$. Il résulte de
Faisceaux, lemme \ref{sheaves-lemma-compute-pullback-to-limit}
(en utilisant aussi que $S_b$ et $B$ sont des ensembles finis) que,
après avoir augmenté $i$, on trouve des applications
$$
\beta_{b, i}, \gamma_{b, i} :
S_b
\longrightarrow
\Gamma(V_{b, i}, \coprod\nolimits_{a \in A} j_{U_{a, i}!}\underline{S_a})
, \quad b \in B
$$
qui donnent les applications $\beta_b$ et $\gamma_b$ après image inverse
par $p_i$. Ces applications correspondent à leur tour à des applications de faisceaux
$$
\beta_i, \gamma_i :
\coprod\nolimits_{b \in B} j_{V_{b, i}!}\underline{S_b}
\longrightarrow
\coprod\nolimits_{a \in A} j_{U_{a, i}!}\underline{S_a}
$$
sur $X_i$. On peut alors prendre $Y = X_i$ et
$$
\mathcal{G} =
\text{Coégalisateur}\left(
\xymatrix{
\coprod\nolimits_{b = 1, \ldots, m} j_{V_{b, i}!}\underline{S_b}
\ar@<1ex>[r] \ar@<-1ex>[r] &
\coprod\nolimits_{a = 1, \ldots, n} j_{U_{a, i}!}\underline{S_a}
}
\right)
$$
Nous omettons quelques détails.
\end{proof}

\begin{lemma}
\label{modules-lemma-constructible-in-constant}
Soit $X$ un espace topologique spectral. Soit $\mathcal{B}$
l'ensemble des ouverts quasi-compacts de $X$.
Soit $\mathcal{F}$ un faisceau d'ensembles comme dans
l'équation (\ref{modules-equation-towards-constructible-sets}).
Alors il existe un nombre fini de fermés constructibles
$Z_1, \ldots, Z_n \subset X$ et d'ensembles finis $S_i$
tels que $\mathcal{F}$ soit isomorphe à un sous-faisceau de
$\prod (Z_i \to X)_*\underline{S_i}$.
\end{lemma}

\begin{proof}
Par le lemme \ref{modules-lemma-constructible-comes-from-finite},
on se ramène au cas d'un espace topologique sobre fini
et d'un faisceau à germes finis.
Dans ce cas, $\mathcal{F} \subset \prod_{x \in X} i_{x, *}\mathcal{F}_x$,
où $i_x : \{x\} \to X$ est l'inclusion. Nous omettons la démonstration
du fait que $i_{x, *}\mathcal{F}_x$ est un faisceau constant sur $\overline{\{x\}}$.
\end{proof}
\section{Morphismes plats d'espaces annelés}
\label{modules-section-flat-morphisms}

\noindent
La définition ponctuelle est motivée par le
Lemme \ref{modules-lemma-flat-stalks-flat}
et la
Définition \ref{modules-definition-flat-at-point}
ci-dessus.

\begin{definition}
\label{modules-definition-flat-morphism}
Soit $f : X \to Y$ un morphisme d'espaces annelés.
Soit $x \in X$. Nous disons que $f$ est {\it plat en $x$}
si le morphisme d'anneaux $\mathcal{O}_{Y, f(x)} \to \mathcal{O}_{X, x}$ est plat.
Nous disons que $f$ est {\it plat} si $f$ est plat en tout $x \in X$.
\end{definition}

\noindent
Considérons le morphisme de faisceaux d'anneaux
$f^\sharp : f^{-1}\mathcal{O}_Y \to \mathcal{O}_X$.
Nous voyons que son germe en $x$ est le morphisme d'anneaux
$f^\sharp_x : \mathcal{O}_{Y, f(x)} \to \mathcal{O}_{X, x}$.
Ainsi, $f$ est plat en $x$ si et seulement si $\mathcal{O}_X$ est plat en $x$
comme $f^{-1}\mathcal{O}_Y$-module. De plus, $f$ est plat si et seulement si
$\mathcal{O}_X$ est plat comme $f^{-1}\mathcal{O}_Y$-module.
Les immersions ouvertes fournissent un cas particulier de morphismes plats.

\begin{lemma}
\label{modules-lemma-pullback-flat}
Soit $f : X \to Y$ un morphisme plat d'espaces annelés.
Alors le foncteur image inverse
$f^* : \textit{Mod}(\mathcal{O}_Y) \to \textit{Mod}(\mathcal{O}_X)$
est exact.
\end{lemma}

\begin{proof}
Le foncteur $f^*$ est la composée du foncteur exact
$f^{-1} : \textit{Mod}(\mathcal{O}_Y) \to \textit{Mod}(f^{-1}\mathcal{O}_Y)$
et du foncteur de changement d'anneaux
$$
\textit{Mod}(f^{-1}\mathcal{O}_Y) \to \textit{Mod}(\mathcal{O}_X), \quad
\mathcal{F} \longmapsto
\mathcal{F} \otimes_{f^{-1}\mathcal{O}_Y} \mathcal{O}_X.
$$
Le résultat découle donc de la discussion qui suit la
Définition \ref{modules-definition-flat-morphism}.
\end{proof}

\begin{definition}
\label{modules-definition-flat-module}
Soit $f : (X, \mathcal{O}_X) \to (Y, \mathcal{O}_Y)$ un morphisme
d'espaces annelés. Soit $\mathcal{F}$ un faisceau de $\mathcal{O}_X$-modules.
\begin{enumerate}
\item Nous disons que $\mathcal{F}$ est {\it plat sur $Y$ en un point $x \in X$}
si le germe $\mathcal{F}_x$ est un $\mathcal{O}_{Y, f(x)}$-module plat.
\item Nous disons que $\mathcal{F}$ est {\it plat sur $Y$} si
$\mathcal{F}$ est plat sur $Y$ en tout point $x$ de $X$.
\end{enumerate}
\end{definition}

\noindent
Avec cette définition, nous voyons que $\mathcal{F}$ est plat sur $Y$ en $x$
si et seulement si $\mathcal{F}$ est plat en $x$ comme
$f^{-1}\mathcal{O}_Y$-module, car
$(f^{-1}\mathcal{O}_Y)_x = \mathcal{O}_{Y, f(x)}$ d'après
Faisceaux, lemme \ref{sheaves-lemma-stalk-pullback}.

\begin{lemma}
\label{modules-lemma-pullback-tensor-flat-module}
Soit $f : X \to Y$ un morphisme d'espaces annelés. Soit $\mathcal{F}$
un $\mathcal{O}_X$-module plat sur $Y$. Alors le foncteur
$$
\textit{Mod}(\mathcal{O}_Y) \to \textit{Mod}(\mathcal{O}_X),\quad
\mathcal{G} \longmapsto f^*\mathcal{G} \otimes_{\mathcal{O}_X} \mathcal{F}
$$
est exact.
\end{lemma}

\begin{proof}
Cela est vrai parce que
$f^*\mathcal{G} \otimes_{\mathcal{O}_X} \mathcal{F} =
f^{-1}\mathcal{G} \otimes_{f^{-1}\mathcal{O}_Y} \mathcal{F}$,
le foncteur $f^{-1}$ est exact et $\mathcal{F}$ est un
$f^{-1}\mathcal{O}_Y$-module plat.
\end{proof}














\section{Puissances symétriques et extérieures}
\label{modules-section-symmetric-exterior}

\noindent
Soit $(X, \mathcal{O}_X)$ un espace annelé.
Soit $\mathcal{F}$ un $\mathcal{O}_X$-module.
Nous définissons l'{\it algèbre tensorielle} de $\mathcal{F}$ comme
le faisceau de $\mathcal{O}_X$-algèbres non commutatives
$$
\text{T}(\mathcal{F})
=
\text{T}_{\mathcal{O}_X}(\mathcal{F})
= \bigoplus\nolimits_{n \geq 0} \text{T}^n(\mathcal{F}).
$$
Ici, $\text{T}^0(\mathcal{F}) = \mathcal{O}_X$,
$\text{T}^1(\mathcal{F}) = \mathcal{F}$
et, pour $n \geq 2$, nous avons
$$
\text{T}^n(\mathcal{F}) =
\mathcal{F} \otimes_{\mathcal{O}_X} \ldots \otimes_{\mathcal{O}_X} \mathcal{F}
\ \ (n\text{ facteurs})
$$
Nous définissons $\wedge(\mathcal{F})$ comme le quotient de
$\text{T}(\mathcal{F})$ par l'idéal bilatère engendré par les
sections locales $s \otimes s$ de $\text{T}^2(\mathcal{F})$, où
$s$ est une section locale de $\mathcal{F}$. On l'appelle
l'{\it algèbre extérieure} de $\mathcal{F}$.
De même, nous définissons $\text{Sym}(\mathcal{F})$ comme
le quotient de $\text{T}(\mathcal{F})$ par l'idéal
bilatère engendré par les sections locales de la forme
$s \otimes t - t \otimes s$ de $\text{T}^2(\mathcal{F})$.

\medskip\noindent
Les algèbres $\text{T}(\mathcal{F})$, $\wedge(\mathcal{F})$ et
$\text{Sym}(\mathcal{F})$ sont des $\mathcal{O}_X$-algèbres graduées
au sens de Faisceaux différentiels gradués, définition
\ref{sdga-definition-ga}.
De plus, $\text{Sym}(\mathcal{F})$
est commutative et $\wedge(\mathcal{F})$ est commutative au sens gradué.

\begin{lemma}
\label{modules-lemma-local-tensor-algebra}
Dans la situation décrite ci-dessus,
le faisceau $\wedge^n\mathcal{F}$ est le faisceau associé au
préfaisceau
$$
U \longmapsto \wedge^n_{\mathcal{O}_X(U)}(\mathcal{F}(U)).
$$
Voir Algèbre, section \ref{algebra-section-tensor-algebra}.
De même, le faisceau $\text{Sym}^n\mathcal{F}$ est le faisceau associé
au préfaisceau
$$
U \longmapsto \text{Sym}^n_{\mathcal{O}_X(U)}(\mathcal{F}(U)).
$$
\end{lemma}

\begin{proof}
Omis. Il peut être plus efficace de définir $\text{Sym}(\mathcal{F})$
et $\wedge(\mathcal{F})$ de cette manière plutôt que par la méthode
donnée ci-dessus.
\end{proof}

\begin{lemma}
\label{modules-lemma-stalk-tensor-algebra}
Dans la situation décrite ci-dessus, soit $x \in X$.
Il existe des isomorphismes canoniques de $\mathcal{O}_{X, x}$-modules
$\text{T}(\mathcal{F})_x = \text{T}(\mathcal{F}_x)$,
$\text{Sym}(\mathcal{F})_x = \text{Sym}(\mathcal{F}_x)$ et
$\wedge(\mathcal{F})_x = \wedge(\mathcal{F}_x)$.
\end{lemma}

\begin{proof}
Cela résulte immédiatement du Lemme \ref{modules-lemma-local-tensor-algebra} ci-dessus et de
Algèbre, lemme \ref{algebra-lemma-colimit-tensor-algebra}.
\end{proof}

\begin{lemma}
\label{modules-lemma-pullback-tensor-algebra}
Soit $f : (X, \mathcal{O}_X) \to (Y, \mathcal{O}_Y)$ un morphisme
d'espaces annelés. Soit $\mathcal{F}$ un faisceau de $\mathcal{O}_Y$-modules.
Alors $f^*\text{T}(\mathcal{F}) = \text{T}(f^*\mathcal{F})$,
et il en va de même des algèbres extérieure et symétrique associées
à $\mathcal{F}$.
\end{lemma}

\begin{proof}
Omis.
\end{proof}

\begin{lemma}
\label{modules-lemma-presentation-sym-exterior}
Soit $(X, \mathcal{O}_X)$ un espace annelé.
Soit $\mathcal{F}_2 \to \mathcal{F}_1 \to \mathcal{F} \to 0$
une suite exacte de faisceaux de $\mathcal{O}_X$-modules.
Pour tout $n \geq 1$, il existe une suite exacte
$$
\mathcal{F}_2 \otimes_{\mathcal{O}_X} \text{Sym}^{n - 1}(\mathcal{F}_1)
\to
\text{Sym}^n(\mathcal{F}_1)
\to
\text{Sym}^n(\mathcal{F})
\to
0
$$
et, de même, une suite exacte
$$
\mathcal{F}_2 \otimes_{\mathcal{O}_X} \wedge^{n - 1}(\mathcal{F}_1)
\to
\wedge^n(\mathcal{F}_1)
\to
\wedge^n(\mathcal{F})
\to
0
$$
\end{lemma}

\begin{proof}
Voir Algèbre, lemme \ref{algebra-lemma-presentation-sym-exterior}.
\end{proof}

\begin{lemma}
\label{modules-lemma-tensor-algebra-permanence}
Soit $(X, \mathcal{O}_X)$ un espace annelé.
Soit $\mathcal{F}$ un faisceau de $\mathcal{O}_X$-modules.
\begin{enumerate}
\item Si $\mathcal{F}$ est localement engendré par des sections,
il en va de même de chaque $\text{T}^n(\mathcal{F})$,
$\wedge^n(\mathcal{F})$ et $\text{Sym}^n(\mathcal{F})$.
\item Si $\mathcal{F}$ est de type fini,
il en va de même de chaque $\text{T}^n(\mathcal{F})$,
$\wedge^n(\mathcal{F})$ et $\text{Sym}^n(\mathcal{F})$.
\item Si $\mathcal{F}$ est de présentation finie,
il en va de même de chaque $\text{T}^n(\mathcal{F})$,
$\wedge^n(\mathcal{F})$ et $\text{Sym}^n(\mathcal{F})$.
\item Si $\mathcal{F}$ est cohérent,
alors, pour $n > 0$, chacun de $\text{T}^n(\mathcal{F})$,
$\wedge^n(\mathcal{F})$ et $\text{Sym}^n(\mathcal{F})$
est cohérent.
\item Si $\mathcal{F}$ est quasi-cohérent,
il en va de même de chaque $\text{T}^n(\mathcal{F})$,
$\wedge^n(\mathcal{F})$ et $\text{Sym}^n(\mathcal{F})$.
\item Si $\mathcal{F}$ est localement libre,
il en va de même de chaque $\text{T}^n(\mathcal{F})$,
$\wedge^n(\mathcal{F})$ et $\text{Sym}^n(\mathcal{F})$.
\end{enumerate}
\end{lemma}

\begin{proof}
Ces assertions pour $\text{T}^n(\mathcal{F})$ résultent du
Lemme \ref{modules-lemma-tensor-product-permanence}.

\medskip\noindent
Les assertions (1) et (2) résultent du fait que
$\wedge^n(\mathcal{F})$ et $\text{Sym}^n(\mathcal{F})$
sont des quotients de $\text{T}^n(\mathcal{F})$.

\medskip\noindent
L'assertion (6) résulte de
Algèbre, lemme \ref{algebra-lemma-free-tensor-algebra}.

\medskip\noindent
Pour (3) et (5), nous utiliserons le
Lemme \ref{modules-lemma-presentation-sym-exterior} ci-dessus.
En choisissant localement une présentation
$\mathcal{F}_2 \to \mathcal{F}_1 \to \mathcal{F} \to 0$
où les $\mathcal{F}_i$ sont libres, ou libres de type fini, puis en appliquant le
lemme, nous voyons que $\text{Sym}^n(\mathcal{F})$ et $\wedge^n(\mathcal{F})$
admettent une présentation analogue ; nous utilisons ici (6) et le
Lemme \ref{modules-lemma-tensor-product-permanence}.

\medskip\noindent
Pour démontrer (4), nous utiliserons
Algèbre, lemme \ref{algebra-lemma-present-sym-wedge}.
Nous pouvons localiser sur $X$ et supposer que
$\mathcal{F}$ est engendré par une famille finie
$(s_i)_{i \in I}$ de sections globales.
Le lemme mentionné ci-dessus,
combiné au Lemme \ref{modules-lemma-local-tensor-algebra} ci-dessus,
implique que, pour $n \geq 2$,
il existe une suite exacte
$$
\bigoplus\nolimits_{j \in J}
\text{T}^{n - 2}(\mathcal{F})
\to
\text{T}^n(\mathcal{F})
\to
\text{Sym}^n(\mathcal{F})
\to
0
$$
où l'ensemble d'indices $J$ est fini. Or nous savons que
$\text{T}^{n - 2}(\mathcal{F})$ est de type fini ;
l'image de la première flèche est donc un sous-faisceau cohérent
de $\text{T}^n(\mathcal{F})$ ; voir le Lemme \ref{modules-lemma-coherent-abelian}.
Ce même lemme permet de conclure que $\text{Sym}^n(\mathcal{F})$ est
cohérent.
\end{proof}

\begin{lemma}
\label{modules-lemma-whole-tensor-algebra-permanence}
Soit $(X, \mathcal{O}_X)$ un espace annelé.
Soit $\mathcal{F}$ un faisceau de $\mathcal{O}_X$-modules.
\begin{enumerate}
\item Si $\mathcal{F}$ est quasi-cohérent,
il en va de même de chacun de $\text{T}(\mathcal{F})$,
$\wedge(\mathcal{F})$ et $\text{Sym}(\mathcal{F})$.
\item Si $\mathcal{F}$ est localement libre,
il en va de même de chacun de $\text{T}(\mathcal{F})$,
$\wedge(\mathcal{F})$ et $\text{Sym}(\mathcal{F})$.
\end{enumerate}
\end{lemma}

\begin{proof}
Il n'est pas vrai qu'une somme directe infinie $\bigoplus \mathcal{G}_i$ de
modules localement libres soit localement libre, ni qu'une
somme directe infinie de modules quasi-cohérents
soit quasi-cohérente. Le problème est que, pour un
point $x \in X$, les voisinages ouverts $U_i$ de $x$ sur lesquels $\mathcal{G}_i$
devient libre (resp.\ admet une présentation convenable) peuvent avoir une intersection
qui n'est pas un voisinage ouvert de $x$. Cependant, dans la
démonstration du Lemme \ref{modules-lemma-tensor-algebra-permanence}, nous avons vu que,
dès qu'un voisinage ouvert convenable pour $\mathcal{F}$ est choisi,
ce voisinage ouvert convient à chacun des faisceaux
$\text{T}^n(\mathcal{F})$, $\wedge^n(\mathcal{F})$ et
$\text{Sym}^n(\mathcal{F})$.
Le lemme en résulte.
\end{proof}







\section{Hom interne}
\label{modules-section-internal-hom}

\noindent
Soit $(X, \mathcal{O}_X)$ un espace annelé.
Soient $\mathcal{F}$ et $\mathcal{G}$ des $\mathcal{O}_X$-modules.
Considérons la règle
$$
U \longmapsto \Hom_{\mathcal{O}_X|_U}(\mathcal{F}|_U, \mathcal{G}|_U).
$$
Il résulte de la discussion de Faisceaux, section
\ref{sheaves-section-glueing-sheaves}, qu'il s'agit d'un faisceau de
groupes abéliens. De plus, étant donnés un élément
$\varphi \in \Hom_{\mathcal{O}_X|_U}(\mathcal{F}|_U, \mathcal{G}|_U)$
et une section $f \in \mathcal{O}_X(U)$, nous pouvons définir
$f\varphi \in \Hom_{\mathcal{O}_X|_U}(\mathcal{F}|_U, \mathcal{G}|_U)$
soit en précomposant avec la multiplication par $f$ sur $\mathcal{F}|_U$,
soit en postcomposant avec la multiplication par $f$ sur $\mathcal{G}|_U$ (le résultat
est le même). Nous obtenons donc en fait un faisceau de $\mathcal{O}_X$-modules.
Nous noterons ce faisceau
$\SheafHom_{\mathcal{O}_X}(\mathcal{F}, \mathcal{G})$.
Il existe un morphisme canonique d'``évaluation''
$$
\mathcal{F}
\otimes_{\mathcal{O}_X}
\SheafHom_{\mathcal{O}_X}(\mathcal{F}, \mathcal{G})
\longrightarrow
\mathcal{G}.
$$
Pour tout $x \in X$, il existe aussi un morphisme canonique
$$
\SheafHom_{\mathcal{O}_X}(\mathcal{F}, \mathcal{G})_x
\to
\Hom_{\mathcal{O}_{X, x}}(\mathcal{F}_x, \mathcal{G}_x)
$$
qui est rarement un isomorphisme.

\begin{lemma}
\label{modules-lemma-internal-hom}
Soit $(X, \mathcal{O}_X)$ un espace annelé.
Soient $\mathcal{F}$, $\mathcal{G}$ et $\mathcal{H}$ des $\mathcal{O}_X$-modules.
Il existe un isomorphisme canonique
$$
\SheafHom_{\mathcal{O}_X}
(\mathcal{F} \otimes_{\mathcal{O}_X} \mathcal{G}, \mathcal{H})
\longrightarrow
\SheafHom_{\mathcal{O}_X}
(\mathcal{F}, \SheafHom_{\mathcal{O}_X}(\mathcal{G}, \mathcal{H}))
$$
qui est fonctoriel en ses trois arguments (il s'agit du Hom de faisceaux
aux trois emplacements). En particulier, se donner un
morphisme $\mathcal{F} \otimes_{\mathcal{O}_X} \mathcal{G} \to \mathcal{H}$
équivaut à se donner un morphisme
$\mathcal{F} \to \SheafHom_{\mathcal{O}_X}(\mathcal{G}, \mathcal{H})$.
\end{lemma}

\begin{proof}
C'est l'analogue de
Algèbre, lemme \ref{algebra-lemma-hom-from-tensor-product}.
La démonstration, identique, est omise.
\end{proof}

\begin{lemma}
\label{modules-lemma-internal-hom-exact}
Soit $(X, \mathcal{O}_X)$ un espace annelé.
Soient $\mathcal{F}$ et $\mathcal{G}$ des $\mathcal{O}_X$-modules.
\begin{enumerate}
\item Si $\mathcal{F}_2 \to \mathcal{F}_1 \to \mathcal{F} \to 0$
est une suite exacte de $\mathcal{O}_X$-modules, alors
$$
0 \to
\SheafHom_{\mathcal{O}_X}(\mathcal{F}, \mathcal{G}) \to
\SheafHom_{\mathcal{O}_X}(\mathcal{F}_1, \mathcal{G}) \to
\SheafHom_{\mathcal{O}_X}(\mathcal{F}_2, \mathcal{G})
$$
est exacte.
\item Si $0 \to \mathcal{G} \to \mathcal{G}_1 \to \mathcal{G}_2$
est une suite exacte de $\mathcal{O}_X$-modules, alors
$$
0 \to
\SheafHom_{\mathcal{O}_X}(\mathcal{F}, \mathcal{G}) \to
\SheafHom_{\mathcal{O}_X}(\mathcal{F}, \mathcal{G}_1) \to
\SheafHom_{\mathcal{O}_X}(\mathcal{F}, \mathcal{G}_2)
$$
est exacte.
\end{enumerate}
\end{lemma}

\begin{proof}
Soit $\mathcal{F}_2 \to \mathcal{F}_1 \to \mathcal{F} \to 0$
comme en (1). Pour tout ouvert $U \subset X$, la suite
$$
0 \to \Hom_{\mathcal{O}_U}(\mathcal{F}|_U, \mathcal{G}|_U) \to
\Hom_{\mathcal{O}_U}(\mathcal{F}_1|_U, \mathcal{G}|_U) \to
\Hom_{\mathcal{O}_U}(\mathcal{F}_2|_U, \mathcal{G}|_U)
$$
est exacte d'après Homologie, lemme \ref{homology-lemma-check-exactness}.
Cela signifie que
prendre les sections sur $U$ de la suite de faisceaux de (1) donne
une suite exacte de groupes abéliens. La suite de (1) est donc exacte
par définition. La démonstration de (2) est exactement la même.
\end{proof}

\begin{lemma}
\label{modules-lemma-adjoint-tensor-restrict}
Soit $X$ un espace topologique. Soit $\mathcal{O}_1 \to \mathcal{O}_2$
un morphisme de faisceaux d'anneaux. Alors nous avons
$$
\Hom_{\mathcal{O}_1}(\mathcal{F}_{\mathcal{O}_1}, \mathcal{G}) =
\Hom_{\mathcal{O}_2}(\mathcal{F},
\SheafHom_{\mathcal{O}_1}(\mathcal{O}_2, \mathcal{G}))
$$
bifonctoriellement en $\mathcal{F} \in \textit{Mod}(\mathcal{O}_2)$
et $\mathcal{G} \in \textit{Mod}(\mathcal{O}_1)$.
\end{lemma}

\begin{proof}
Omis. C'est l'analogue de
Algèbre, lemme \ref{algebra-lemma-adjoint-hom-restrict},
et la démonstration est exactement la même.
\end{proof}

\begin{lemma}
\label{modules-lemma-stalk-internal-hom}
Soit $(X, \mathcal{O}_X)$ un espace annelé.
Soient $\mathcal{F}$ et $\mathcal{G}$ des $\mathcal{O}_X$-modules.
Si $\mathcal{F}$ est de type fini, alors le morphisme canonique
$$
\SheafHom_{\mathcal{O}_X}(\mathcal{F}, \mathcal{G})_x
\to
\Hom_{\mathcal{O}_{X, x}}(\mathcal{F}_x, \mathcal{G}_x)
$$
est injectif.
Si $\mathcal{F}$ est de présentation finie,
ce morphisme canonique est un isomorphisme.
\end{lemma}

\begin{proof}
Le morphisme envoie la classe d'équivalence de
$(U, \varphi)$ dans $\SheafHom_{\mathcal{O}_X}(\mathcal{F}, \mathcal{G})_x$,
où $x \in U \subset X$, l'ensemble intermédiaire étant ouvert, et
$\varphi \in \Hom_{\mathcal{O}_U}(\mathcal{F}|_U, \mathcal{G}|_U)$,
sur le morphisme induit sur les germes en $x$, à savoir
$\varphi_x : \mathcal{F}_x \to \mathcal{G}_x$.

\medskip\noindent
Supposons $\mathcal{F}$ de type fini.
Choisissons un représentant $(U, \varphi)$ d'un élément $\sigma$
du noyau du morphisme, c'est-à-dire tel que $\varphi_x = 0$.
En rétrécissant $U$ si nécessaire, choisissons des sections
$s^1, \ldots, s^n \in \mathcal{F}(U)$ qui engendrent $\mathcal{F}|_U$.
Comme $\varphi_x(s^i_x) = 0$
et que nous n'avons qu'un nombre fini de sections,
nous pouvons trouver un voisinage ouvert
$V \subset U$ de $x$ tel que $\varphi_V(s^i|_V)=0$
pour tout $i = 1, \ldots, n$. Puisque les $s^i|_V$, $i = 1, \ldots, n$, engendrent
$\mathcal{F}|_V$, cela signifie que $\varphi|_V = 0$.
Comme $(U, \varphi)$ est équivalent à $(V, \varphi|_V)$,
nous concluons que $\sigma = 0$, d'où l'injectivité du morphisme.

\medskip\noindent
Supposons maintenant $\mathcal{F}$ de présentation finie.
En localisant sur $X$, nous pouvons supposer que $\mathcal{F}$ admet une présentation
$$
\bigoplus\nolimits_{j = 1, \ldots, m}
\mathcal{O}_X
\longrightarrow
\bigoplus\nolimits_{i = 1, \ldots, n}
\mathcal{O}_X
\to
\mathcal{F}
\to
0.
$$
D'après le Lemme \ref{modules-lemma-internal-hom-exact}, cela donne une suite exacte
$
0 \to
\SheafHom_{\mathcal{O}_X}(\mathcal{F}, \mathcal{G}) \to
\bigoplus\nolimits_{i = 1, \ldots, n} \mathcal{G}
\longrightarrow
\bigoplus\nolimits_{j = 1, \ldots, m} \mathcal{G}.
$
En prenant les germes, nous obtenons une suite exacte
$
0 \to
\SheafHom_{\mathcal{O}_X}(\mathcal{F}, \mathcal{G})_x \to
\bigoplus\nolimits_{i = 1, \ldots, n} \mathcal{G}_x
\longrightarrow
\bigoplus\nolimits_{j = 1, \ldots, m} \mathcal{G}_x
$
et le résultat en découle puisque $\mathcal{F}_x$ s'insère dans
une suite exacte
$
\bigoplus\nolimits_{j = 1, \ldots, m}
\mathcal{O}_{X, x}
\longrightarrow
\bigoplus\nolimits_{i = 1, \ldots, n}
\mathcal{O}_{X, x}
\to
\mathcal{F}_x
\to
0
$
qui induit la suite exacte
$
0 \to
\Hom_{\mathcal{O}_{X, x}}(\mathcal{F}_x, \mathcal{G}_x) \to
\bigoplus\nolimits_{i = 1, \ldots, n} \mathcal{G}_x
\longrightarrow
\bigoplus\nolimits_{j = 1, \ldots, m} \mathcal{G}_x
$
qui est la même que celle ci-dessus.
\end{proof}

\begin{lemma}
\label{modules-lemma-pullback-internal-hom}
Soit $f : (X, \mathcal{O}_X) \to (Y, \mathcal{O}_Y)$ un morphisme
d'espaces annelés. Soient $\mathcal{F}$ et $\mathcal{G}$ des $\mathcal{O}_Y$-modules.
Si $\mathcal{F}$ est de présentation finie et si $f$ est plat,
alors le morphisme canonique
$$
f^*\SheafHom_{\mathcal{O}_Y}(\mathcal{F}, \mathcal{G})
\longrightarrow
\SheafHom_{\mathcal{O}_X}(f^*\mathcal{F}, f^*\mathcal{G})
$$
est un isomorphisme.
\end{lemma}

\begin{proof}
Remarquons que $f^*\mathcal{F}$ est aussi de présentation finie
(Lemme \ref{modules-lemma-pullback-finite-presentation}).
Soit $x \in X$ d'image $y \in Y$. En considérant les germes
en $x$, nous obtenons un isomorphisme d'après le
Lemme \ref{modules-lemma-stalk-internal-hom} et
Compléments d'algèbre, lemme
\ref{more-algebra-lemma-pseudo-coherence-and-base-change-ext}
qui montre que, dans ce cas, $\Hom$ commute au changement de base par
$\mathcal{O}_{Y, y} \to \mathcal{O}_{X, x}$.
Seconde démonstration : utiliser exactement le même argument que dans
la démonstration du Lemme \ref{modules-lemma-stalk-internal-hom}.
\end{proof}

\begin{lemma}
\label{modules-lemma-internal-hom-locally-kernel-direct-sum}
Soit $(X, \mathcal{O}_X)$ un espace annelé.
Soient $\mathcal{F}$ et $\mathcal{G}$ des $\mathcal{O}_X$-modules.
Si $\mathcal{F}$ est de présentation finie, alors le faisceau
$\SheafHom_{\mathcal{O}_X}(\mathcal{F}, \mathcal{G})$
est localement le noyau d'un morphisme entre des sommes directes finies
de copies de $\mathcal{G}$.
En particulier, si $\mathcal{G}$ est cohérent, alors
$\SheafHom_{\mathcal{O}_X}(\mathcal{F}, \mathcal{G})$
est lui aussi cohérent.
\end{lemma}

\begin{proof}
Nous avons vu la première assertion
dans la démonstration du Lemme \ref{modules-lemma-stalk-internal-hom}.
Le résultat pour les faisceaux cohérents découle alors du
Lemme \ref{modules-lemma-coherent-abelian}.
\end{proof}

\begin{lemma}
\label{modules-lemma-hom-finite-presentation-colimit}
Soit $X$ un espace annelé. Soit $\mathcal{F}$ un $\mathcal{O}_X$-module
de présentation finie. Soit
$\mathcal{G} = \colim_{\lambda \in \Lambda} \mathcal{G}_\lambda$
une colimite filtrante de $\mathcal{O}_X$-modules. Alors le morphisme canonique
$$
\colim_\lambda \SheafHom_{\mathcal{O}_X}(\mathcal{F}, \mathcal{G}_\lambda)
\longrightarrow
\SheafHom_{\mathcal{O}_X}(\mathcal{F}, \mathcal{G})
$$
est un isomorphisme.
\end{lemma}

\begin{proof}
La formation des colimites de faisceaux de modules commute à la restriction aux ouverts ; voir
Faisceaux, section \ref{sheaves-section-limits-sheaves}. Nous pouvons donc supposer que
$\mathcal{F}$ admet une présentation globale
$$
\bigoplus\nolimits_{j = 1, \ldots, m}
\mathcal{O}_X
\longrightarrow
\bigoplus\nolimits_{i = 1, \ldots, n}
\mathcal{O}_X
\to
\mathcal{F}
\to
0
$$
Le foncteur $\SheafHom_{\mathcal{O}_X}(-, -)$ commute aux sommes directes finies
en chacune des variables et $\SheafHom_{\mathcal{O}_X}(\mathcal{O}_X, -)$
est le foncteur identité. D'après cela et le Lemme \ref{modules-lemma-internal-hom-exact},
nous obtenons une suite exacte
$$
0 \to
\SheafHom_{\mathcal{O}_X}(\mathcal{F}, \mathcal{G}) \to
\bigoplus\nolimits_{i = 1, \ldots, n} \mathcal{G} \to
\bigoplus\nolimits_{j = 1, \ldots, m} \mathcal{G}
$$
Puisque les colimites filtrantes sont exactes dans $\textit{Mod}(\mathcal{O}_X)$,
la ligne supérieure du diagramme commutatif suivant est elle aussi exacte
$$
\xymatrix{
0 \ar[r] &
\colim_\lambda \SheafHom_{\mathcal{O}_X}(\mathcal{F}, \mathcal{G}_\lambda)
\ar[r] \ar[d] &
\colim_\lambda \bigoplus\nolimits_{i = 1, \ldots, n} \mathcal{G}_\lambda
\ar[r] \ar[d] &
\colim_\lambda \bigoplus\nolimits_{j = 1, \ldots, m} \mathcal{G}_\lambda
\ar[d] \\
0 \ar[r] &
\SheafHom_{\mathcal{O}_X}(\mathcal{F}, \mathcal{G}) \ar[r] &
\bigoplus\nolimits_{i = 1, \ldots, n} \mathcal{G} \ar[r] &
\bigoplus\nolimits_{j = 1, \ldots, m} \mathcal{G}
}
$$
Comme les deux flèches verticales de droite sont des isomorphismes, nous concluons.
\end{proof}

\begin{lemma}
\label{modules-lemma-finite-presentation-quasi-compact-colimit}
Soit $X$ un espace annelé.
Soit $I$ un ensemble préordonné et
soit $(\mathcal{F}_i, \varphi_{ii'})$ un système sur $I$
formé de faisceaux de $\mathcal{O}_X$-modules
(voir Catégories, section \ref{categories-section-posets-limits}).
Supposons que
\begin{enumerate}
\item $I$ est filtrant,
\item $\mathcal{G}$ est un $\mathcal{O}_X$-module de présentation finie, et
\item $X$ admet un système cofinal de recouvrements ouverts
$\mathcal{U} : X = \bigcup_{j\in J} U_j$, avec
$J$ fini et $U_j \cap U_{j'}$ quasi-compact
pour tous $j, j' \in J$.
\end{enumerate}
Alors nous avons
$$
\colim_i \Hom_X(\mathcal{G}, \mathcal{F}_i)
=
\Hom_X(\mathcal{G}, \colim_i \mathcal{F}_i).
$$
\end{lemma}

\begin{proof}
Posons
$\mathcal{H} = \SheafHom_{\mathcal{O}_X}(\mathcal{G}, \colim \mathcal{F}_i)$
et $\mathcal{H}_i = \SheafHom_{\mathcal{O}_X}(\mathcal{G}, \mathcal{F}_i)$.
Rappelons que
$$
\Hom_X(\mathcal{G}, \mathcal{F}) = \Gamma(X, \mathcal{H})
\quad\text{et}\quad
\Hom_X(\mathcal{G}, \mathcal{F}_i) = \Gamma(X, \mathcal{H}_i)
$$
par construction. D'après le Lemme \ref{modules-lemma-hom-finite-presentation-colimit}, nous avons
$\mathcal{H} = \colim \mathcal{H}_i$. Le lemme résulte donc de
Faisceaux, lemme \ref{sheaves-lemma-directed-colimits-sections}.
\end{proof}

\begin{remark}
\label{modules-remark-condition-necessary}
Dans le lemme ci-dessus, il faut une condition supplémentaire à la
quasi-compacité de $X$. Voir
Faisceaux, exemple \ref{sheaves-example-conditions-needed-colimit}.
\end{remark}




\section{L'annulateur d'un faisceau de modules}
\label{modules-section-annihilator}

\noindent
Soit $(X, \mathcal{O}_X)$ un espace annelé.
Soit $\mathcal{F}$ un $\mathcal{O}_X$-module.
Il existe un morphisme canonique de faisceaux de $\mathcal{O}_X$-modules
$$
\mathcal{O}_X
\longrightarrow
\SheafHom_{\mathcal{O}_X}(\mathcal{F}, \mathcal{F})
$$
qui envoie une section locale $f \in \mathcal{O}_X(U)$ sur le morphisme
$f : \mathcal{F}|_U \to \mathcal{F}|_U$ donné par la multiplication par $f$.

\begin{definition}
\label{modules-definition-annihilator-sheaf}
Soit $(X, \mathcal{O}_X)$ un espace annelé et soit $\mathcal{F}$
un $\mathcal{O}_X$-module. L'{\it annulateur} de $\mathcal{F}$,
noté $\text{Ann}_{\mathcal{O}_X}(\mathcal{F})$,
est le noyau du morphisme
$\mathcal{O}_X \to \SheafHom_{\mathcal{O}_X}(\mathcal{F}, \mathcal{F})$
considéré ci-dessus.
\end{definition}

\noindent
Pour tout $x\in X$, on a une inclusion d'idéaux de $\mathcal{O}_{X, x}$ :
\begin{equation}
\label{modules-equation-inclusion-of-annihilator-ideals}
(\text{Ann}_{\mathcal{O}_X}(\mathcal{F}))_x
\subset
\text{Ann}_{\mathcal{O}_{X, x}}(\mathcal{F}_x)
\end{equation}
car toute section de $\text{Ann}_{\mathcal{O}_X}(\mathcal{F})$
annule les germes de $\mathcal{F}$ en tout point où elle est
définie. Voici une situation simple dans laquelle
(\ref{modules-equation-inclusion-of-annihilator-ideals}) est une égalité.

\begin{lemma}
\label{modules-lemma-finite-type-annihilator-stalk}
Soit $(X, \mathcal{O}_X)$ un espace annelé et soit
$\mathcal{F}$ un faisceau de $\mathcal{O}_X$-modules.
Si $\mathcal{F}$ est de type fini, alors
$(\text{Ann}_{\mathcal{O}_X}(\mathcal{F}))_x =
\text{Ann}_{\mathcal{O}_{X, x}}(\mathcal{F}_x)$.
\end{lemma}

\begin{proof}
D'après le lemme \ref{modules-lemma-stalk-internal-hom}, le morphisme
$$
\SheafHom_{\mathcal{O}_X}(\mathcal{F}, \mathcal{F})_x \longrightarrow
\Hom_{\mathcal{O}_{X,x}}(\mathcal{F}_x, \mathcal{F}_x)
$$
est injectif. Ainsi, toute section $f$ de $\mathcal{O}_X$ sur un voisinage
ouvert $U$ de $x$ qui agit par zéro sur $\mathcal{F}_x$
agit par zéro sur $\mathcal{F}|_V$ pour un ouvert $U \supset V \ni x$ convenable.
L'inclusion (\ref{modules-equation-inclusion-of-annihilator-ideals})
est donc une égalité.
\end{proof}

\begin{lemma}
\label{modules-lemma-induced-quotient-ring-sheaf-module-structure}
Soit $(X, \mathcal{O}_X)$ un espace annelé, soit $\mathcal{F}$ un
$\mathcal{O}_X$-module et soit $\mathcal{I} \subset \mathcal{O}_X$
un faisceau d'idéaux. Si
$\mathcal{I} \subset \text{Ann}_{\mathcal{O}_X}(\mathcal{F})$,
alors $\mathcal{F}$ possède une structure naturelle de
$\mathcal{O}_X/\mathcal{I}$-module qui, sur les germes, coïncide avec
la construction usuelle d'algèbre commutative.
\end{lemma}

\begin{proof}
En appliquant la propriété universelle du conoyau de l'inclusion
$\mathcal{I} \to \mathcal{O}_X$, nous obtenons un diagramme commutatif
$$
\xymatrix{
\mathcal{O}_X \ar[r] \ar[d] &
\SheafHom_{\mathcal{O}_X}(\mathcal{F}, \mathcal{F}) \\
\mathcal{O}_X/\mathcal{I} \ar@{-->}[ur]
}
$$
de $\mathcal{O}_X$-modules. D'après le lemme \ref{modules-lemma-internal-hom},
le morphisme obtenu $\mathcal{O}_X/\mathcal{I} \to
\SheafHom_{\mathcal{O}_X}(\mathcal{F}, \mathcal{F})$
correspond à un morphisme de $\mathcal{O}_X$-modules
$$
\mathcal{O}_X/\mathcal{I} \otimes_{\mathcal{O}_X} \mathcal{F}
\longrightarrow
\mathcal{F}
$$
ce qui nous donne une structure de $\mathcal{O}_X/\mathcal{I}$-module
sur $\mathcal{F}$, compatible avec sa structure donnée de $\mathcal{O}_X$-module.
Nous omettons la vérification de l'assertion sur les germes.
\end{proof}

\begin{lemma}
\label{modules-lemma-coherent-annihilator}
Soit $(X,\mathcal{O}_X)$ un espace annelé.
Si $\mathcal{O}_X$ et $\mathcal{F}$ sont cohérents,
alors $\text{Ann}_{\mathcal{O}_X}(\mathcal{F})$ l'est aussi.
\end{lemma}

\begin{proof}
Puisque $\text{Ann}_{\mathcal{O}_X}(\mathcal{F})$ est le noyau de
$\mathcal{O}_X \to \SheafHom_{\mathcal{O}_X}(\mathcal{F}, \mathcal{F})$
et d'après le lemme \ref{modules-lemma-coherent-abelian}, il suffit
de montrer que $\SheafHom_{\mathcal{O}_X}(\mathcal{F}, \mathcal{F})$ est
cohérent. Cela résulte du
lemme \ref{modules-lemma-internal-hom-locally-kernel-direct-sum}
et du fait que $\mathcal{F}$ est cohérent, donc a fortiori de présentation finie
(lemme \ref{modules-lemma-coherent-finite-presentation}).
\end{proof}





\section{Complexes de Koszul}
\label{modules-section-koszul-complex}

\noindent
Nous suggérons de lire d'abord la section consacrée aux complexes de Koszul dans
Compléments d'algèbre, section \ref{more-algebra-section-koszul}.
Nous définissons comme suit le complexe de Koszul dans la catégorie des
$\mathcal{O}_X$-modules.

\begin{definition}
\label{modules-definition-koszul}
Soit $X$ un espace annelé. Soit $\varphi : \mathcal{E} \to \mathcal{O}_X$
un morphisme de $\mathcal{O}_X$-modules. Le
{\it complexe de Koszul} $K_\bullet(\varphi)$ associé à $\varphi$
est le faisceau d'algèbres différentielles graduées commutatives défini comme suit :
\begin{enumerate}
\item l'algèbre graduée sous-jacente est l'algèbre extérieure
$K_\bullet(\varphi) = \wedge(\mathcal{E})$,
\item la différentielle $d : K_\bullet(\varphi) \to K_\bullet(\varphi)$
est l'unique dérivation telle que $d(e) = \varphi(e)$ pour toute
section locale $e$ de $\mathcal{E} = K_1(\varphi)$.
\end{enumerate}
\end{definition}

\noindent
Explicitement, si $e_1 \wedge \ldots \wedge e_n$ est un produit extérieur de
sections locales de $\mathcal{E}$, alors
$$
d(e_1 \wedge \ldots \wedge e_n) =
\sum\nolimits_{i = 1, \ldots, n} (-1)^{i + 1}
\varphi(e_i)e_1 \wedge \ldots \wedge \widehat{e_i} \wedge \ldots \wedge e_n.
$$
On vérifie immédiatement que cette formule définit une dérivation bien définie
sur l'algèbre tensorielle, qui annule $e \wedge e$ et se factorise donc
par l'algèbre extérieure.

\begin{definition}
\label{modules-definition-koszul-complex}
Soit $X$ un espace annelé et soient
$f_1, \ldots, f_n \in \Gamma(X, \mathcal{O}_X)$. Le
{\it complexe de Koszul associé à $f_1, \ldots, f_n$} est le complexe de Koszul
associé au morphisme
$(f_1, \ldots, f_n) : \mathcal{O}_X^{\oplus n} \to \mathcal{O}_X$.
On le note $K_\bullet(\mathcal{O}_X, f_1, \ldots, f_n)$
ou $K_\bullet(\mathcal{O}_X, f_\bullet)$.
\end{definition}

\noindent
Bien entendu, étant donné un morphisme de $\mathcal{O}_X$-modules
$\varphi : \mathcal{E} \to \mathcal{O}_X$,
si $\mathcal{E}$ est localement libre de type fini, alors
$K_\bullet(\varphi)$ est localement sur $X$ isomorphe à un complexe de Koszul
$K_\bullet(\mathcal{O}_X, f_1, \ldots, f_n)$.



\section{Modules inversibles}
\label{modules-section-invertible}

\noindent
Comme dans le cas des modules sur un anneau
(Compléments d'algèbre, section \ref{more-algebra-section-picard}),
nous avons la définition suivante.

\begin{definition}
\label{modules-definition-invertible}
Soit $(X, \mathcal{O}_X)$ un espace annelé. Un
{\it $\mathcal{O}_X$-module inversible} est un faisceau
de $\mathcal{O}_X$-modules $\mathcal{L}$ tel que
le foncteur
$$
\textit{Mod}(\mathcal{O}_X) \longrightarrow \textit{Mod}(\mathcal{O}_X),\quad
\mathcal{F} \longmapsto \mathcal{L} \otimes_{\mathcal{O}_X} \mathcal{F}
$$
est une équivalence de catégories. Nous disons que $\mathcal{L}$ est
{\it trivial} s'il est isomorphe à $\mathcal{O}_X$ en tant que
$\mathcal{O}_X$-module.
\end{definition}

\noindent
Le lemme \ref{modules-lemma-invertible-is-locally-free-rank-1}
ci-dessous explique le lien avec les modules localement libres
de rang $1$.

\begin{lemma}
\label{modules-lemma-invertible}
Soit $(X, \mathcal{O}_X)$ un espace annelé. Soit $\mathcal{L}$
un $\mathcal{O}_X$-module. Les conditions suivantes sont équivalentes :
\begin{enumerate}
\item $\mathcal{L}$ est inversible ;
\item il existe un $\mathcal{O}_X$-module $\mathcal{N}$
tel que
$\mathcal{L} \otimes_{\mathcal{O}_X} \mathcal{N} \cong \mathcal{O}_X$.
\end{enumerate}
Dans ce cas, $\mathcal{L}$ est localement facteur direct d'un
$\mathcal{O}_X$-module libre de type fini, et le module $\mathcal{N}$ de (2) est isomorphe à
$\SheafHom_{\mathcal{O}_X}(\mathcal{L}, \mathcal{O}_X)$.
\end{lemma}

\begin{proof}
Supposons (1). Le foncteur $- \otimes_{\mathcal{O}_X} \mathcal{L}$
est alors essentiellement surjectif ; il existe donc un $\mathcal{O}_X$-module
$\mathcal{N}$ comme en (2). Si (2) est vérifiée, le foncteur
$- \otimes_{\mathcal{O}_X} \mathcal{N}$ est un quasi-inverse
du foncteur $- \otimes_{\mathcal{O}_X} \mathcal{L}$ ;
on en déduit (1).

\medskip\noindent
Supposons (1) et (2) vérifiées. Notons
$\psi : \mathcal{L} \otimes_{\mathcal{O}_X} \mathcal{N} \to \mathcal{O}_X$
l'isomorphisme donné. Soit $x \in X$. Choisissons un voisinage ouvert
$U$ de $x$, un entier $n \geq 1$ et des sections $s_i \in \mathcal{L}(U)$,
$t_i \in \mathcal{N}(U)$ telles que $\psi(\sum s_i \otimes t_i) = 1$.
Considérons les isomorphismes
$$
\mathcal{L}|_U \to
\mathcal{L}|_U \otimes_{\mathcal{O}_U}
\mathcal{L}|_U \otimes_{\mathcal{O}_U} \mathcal{N}|_U \to \mathcal{L}|_U
$$
où la première flèche envoie $s$ sur $\sum s_i \otimes s \otimes t_i$
et la seconde envoie $s \otimes s' \otimes t$ sur $\psi(s' \otimes t)s$.
Nous en concluons que $s \mapsto \sum \psi(s \otimes t_i)s_i$ est
un automorphisme de $\mathcal{L}|_U$. Cet automorphisme se factorise sous la forme
$$
\mathcal{L}|_U \to \mathcal{O}_U^{\oplus n} \to \mathcal{L}|_U
$$
où la première flèche est donnée par
$s \mapsto (\psi(s \otimes t_1), \ldots, \psi(s \otimes t_n))$
et la seconde par $(a_1, \ldots, a_n) \mapsto \sum a_i s_i$.
Nous en concluons que $\mathcal{L}|_U$ est facteur direct
d'un $\mathcal{O}_U$-module libre de type fini.

\medskip\noindent
Supposons (1) et (2) vérifiées. Considérons le morphisme d'évaluation
$$
\mathcal{L} \otimes_{\mathcal{O}_X}
\SheafHom_{\mathcal{O}_X}(\mathcal{L}, \mathcal{O}_X)
\longrightarrow \mathcal{O}_X
$$
Pour achever la démonstration du lemme,
montrons que c'est un isomorphisme en vérifiant qu'il induit
des isomorphismes sur les germes. Soit $x \in X$.
Comme nous savons, d'après le paragraphe précédent,
que $\mathcal{L}$ est un $\mathcal{O}_X$-module
de présentation finie,
le lemme \ref{modules-lemma-stalk-internal-hom}
montre qu'il suffit d'établir que
$$
\mathcal{L}_x \otimes_{\mathcal{O}_{X, x}}
\Hom_{\mathcal{O}_{X, x}}(\mathcal{L}_x, \mathcal{O}_{X, x})
\longrightarrow \mathcal{O}_{X, x}
$$
est un isomorphisme. Puisque
$\mathcal{L}_x \otimes_{\mathcal{O}_{X, x}} \mathcal{N}_x =
(\mathcal{L} \otimes_{\mathcal{O}_X} \mathcal{N})_x =
\mathcal{O}_{X, x}$ (lemme \ref{modules-lemma-stalk-tensor-product}),
le résultat voulu découle de
Compléments d'algèbre, lemme \ref{more-algebra-lemma-invertible}.
\end{proof}

\begin{lemma}
\label{modules-lemma-pullback-invertible}
Soit $f : (X, \mathcal{O}_X) \to (Y, \mathcal{O}_Y)$ un
morphisme d'espaces annelés. L'image inverse $f^*\mathcal{L}$ d'un
$\mathcal{O}_Y$-module inversible est inversible.
\end{lemma}

\begin{proof}
D'après le lemme \ref{modules-lemma-invertible},
il existe un $\mathcal{O}_Y$-module $\mathcal{N}$ tel que
$\mathcal{L} \otimes_{\mathcal{O}_Y} \mathcal{N} \cong \mathcal{O}_Y$.
En prenant l'image inverse, nous obtenons
$f^*\mathcal{L} \otimes_{\mathcal{O}_X} f^*\mathcal{N} \cong \mathcal{O}_X$
d'après le lemme \ref{modules-lemma-tensor-product-pullback}.
Ainsi $f^*\mathcal{L}$ est inversible d'après le lemme \ref{modules-lemma-invertible}.
\end{proof}

\begin{lemma}
\label{modules-lemma-invertible-is-locally-free-rank-1}
Soit $(X, \mathcal{O}_X)$ un espace annelé. Tout
$\mathcal{O}_X$-module localement libre de rang $1$ est inversible.
Si tous les anneaux $\mathcal{O}_{X, x}$ sont locaux, alors
la réciproque est également vraie (ce qui n'est pas le cas en général).
\end{lemma}

\begin{proof}
L'assertion entre parenthèses s'obtient en considérant un espace réduit à un point
$X$, muni du faisceau d'anneaux $\mathcal{O}_X$ défini par un anneau $R$.
Les $\mathcal{O}_X$-modules inversibles correspondent alors aux
$R$-modules inversibles ; dès que $\Pic(R)$ n'est pas le groupe trivial,
nous obtenons donc un exemple.

\medskip\noindent
Supposons $\mathcal{L}$ localement libre de rang $1$ et considérons le
morphisme d'évaluation
$$
\mathcal{L} \otimes_{\mathcal{O}_X}
\SheafHom_{\mathcal{O}_X}(\mathcal{L}, \mathcal{O}_X)
\longrightarrow \mathcal{O}_X
$$
Sur un recouvrement ouvert qui trivialise $\mathcal{L}$, nous voyons
que ce morphisme est un isomorphisme. Ainsi $\mathcal{L}$ est inversible
d'après le lemme \ref{modules-lemma-invertible}.

\medskip\noindent
Supposons que tous les anneaux $\mathcal{O}_{X, x}$ soient locaux et que $\mathcal{L}$
soit inversible. Dans la démonstration du lemme \ref{modules-lemma-invertible},
nous avons vu que $\mathcal{L}_x$ est un
$\mathcal{O}_{X, x}$-module inversible pour tout $x \in X$. Comme
$\mathcal{O}_{X, x}$ est local, nous voyons que
$\mathcal{L}_x \cong \mathcal{O}_{X, x}$
(Compléments d'algèbre, section \ref{more-algebra-section-picard}).
Puisque $\mathcal{L}$ est de présentation finie d'après le
lemme \ref{modules-lemma-invertible}, nous concluons que $\mathcal{L}$
est localement libre de rang $1$ d'après le
lemme \ref{modules-lemma-finite-presentation-stalk-free}.
\end{proof}

\begin{lemma}
\label{modules-lemma-constructions-invertible}
Soit $(X, \mathcal{O}_X)$ un espace annelé.
\begin{enumerate}
\item Si $\mathcal{L}$ et $\mathcal{N}$ sont des
$\mathcal{O}_X$-modules inversibles, alors il en va de même de
$\mathcal{L} \otimes_{\mathcal{O}_X} \mathcal{N}$.
\item Si $\mathcal{L}$ est un $\mathcal{O}_X$-module inversible, alors
$\SheafHom_{\mathcal{O}_X}(\mathcal{L}, \mathcal{O}_X)$ l'est aussi, et le morphisme d'évaluation
$\mathcal{L} \otimes_{\mathcal{O}_X}
\SheafHom_{\mathcal{O}_X}(\mathcal{L}, \mathcal{O}_X) \to \mathcal{O}_X$
est un isomorphisme.
\end{enumerate}
\end{lemma}

\begin{proof}
L'assertion (1) est immédiate par définition, et (2) résulte du
lemme \ref{modules-lemma-invertible} et de sa démonstration.
\end{proof}

\begin{definition}
\label{modules-definition-powers}
Soit $(X, \mathcal{O}_X)$ un espace annelé. Étant donné un faisceau inversible
$\mathcal{L}$ sur $X$ et $n \in \mathbf{Z}$, nous définissons la
$n$-ième {\it puissance tensorielle} $\mathcal{L}^{\otimes n}$ de $\mathcal{L}$
comme l'image de $\mathcal{O}_X$ par l'application de l'équivalence
$\mathcal{F} \mapsto\mathcal{F} \otimes_{\mathcal{O}_X} \mathcal{L}$
exactement $n$ fois.
\end{definition}

\noindent
Cette définition a également un sens pour les entiers négatifs $n$, puisque nous avons
défini un $\mathcal{O}_X$-module inversible comme un module pour lequel le produit tensoriel
est une équivalence. Plus explicitement, nous avons
$$
\mathcal{L}^{\otimes n} =
\left\{
\begin{matrix}
\mathcal{O}_X & \text{si} & n = 0 \\
\SheafHom_{\mathcal{O}_X}(\mathcal{L}, \mathcal{O}_X) & \text{si} & n = -1\\
\mathcal{L} \otimes_{\mathcal{O}_X} \ldots \otimes_{\mathcal{O}_X} \mathcal{L}
& \text{si} & n > 0 \\
\mathcal{L}^{\otimes -1} \otimes_{\mathcal{O}_X} \ldots
\otimes_{\mathcal{O}_X} \mathcal{L}^{\otimes -1}
& \text{si} & n < -1
\end{matrix}
\right.
$$
voir le lemme \ref{modules-lemma-constructions-invertible}.
Avec cette définition, nous avons des isomorphismes canoniques
$\mathcal{L}^{\otimes n} \otimes_{\mathcal{O}_X}
\mathcal{L}^{\otimes m} \to
\mathcal{L}^{\otimes n + m}$, et ces isomorphismes
satisfont à des contraintes de commutativité et d'associativité
(formulation omise).

\medskip\noindent
Soit $(X, \mathcal{O}_X)$ un espace annelé. Nous pouvons définir une structure
d'anneau $\mathbf{Z}$-gradué sur $\bigoplus \Gamma(X, \mathcal{L}^{\otimes n})$ en envoyant
$s \in \Gamma(X, \mathcal{L}^{\otimes n})$ et
$t \in \Gamma(X, \mathcal{L}^{\otimes m})$ sur la
section correspondant à $s \otimes t$ dans
$\Gamma(X, \mathcal{L}^{\otimes n + m})$.
Nous omettons de vérifier que cela définit un anneau commutatif
et associatif, d'unité $1$. Toutefois, selon nos conventions dans
Algèbre, section \ref{algebra-section-graded},
un anneau gradué ne possède aucun élément non nul de degré négatif.
Cela conduit à la définition suivante.

\begin{definition}
\label{modules-definition-gamma-star}
Soit $(X, \mathcal{O}_X)$ un espace annelé.
Étant donné un faisceau inversible $\mathcal{L}$ sur $X$, nous définissons
l'{\it anneau gradué associé} par
$$
\Gamma_*(X, \mathcal{L})
=
\bigoplus\nolimits_{n \geq 0} \Gamma(X, \mathcal{L}^{\otimes n})
$$
Étant donné un faisceau de $\mathcal{O}_X$-modules $\mathcal{F}$, nous posons
$$
\Gamma_*(X, \mathcal{L}, \mathcal{F})
=
\bigoplus\nolimits_{n \in \mathbf{Z}} \Gamma(X,
\mathcal{F} \otimes_{\mathcal{O}_X} \mathcal{L}^{\otimes n})
$$
que nous considérons comme un $\Gamma_*(X, \mathcal{L})$-module gradué.
\end{definition}

\noindent
Nous écrivons souvent simplement $\Gamma_*(\mathcal{L})$ et $\Gamma_*(\mathcal{F})$
(bien que cette notation soit ambiguë si $\mathcal{F}$ est inversible).
La multiplication de $\Gamma_*(\mathcal{L})$ sur
$\Gamma_*(\mathcal{F})$ est définie au moyen des isomorphismes
ci-dessus. Si $\gamma : \mathcal{F} \to \mathcal{G}$ est un morphisme de
$\mathcal{O}_X$-modules, nous obtenons un morphisme de $\Gamma_*(\mathcal{L})$-modules
$\gamma : \Gamma_*(\mathcal{F}) \to \Gamma_*(\mathcal{G})$.
Si $\alpha : \mathcal{L} \to \mathcal{N}$ est un morphisme de
$\mathcal{O}_X$-modules entre des $\mathcal{O}_X$-modules inversibles, nous obtenons
un morphisme d'anneaux gradués $\Gamma_*(\mathcal{L}) \to \Gamma_*(\mathcal{N})$.
Si $f : (Y, \mathcal{O}_Y) \to (X, \mathcal{O}_X)$
est un morphisme d'espaces annelés et si $\mathcal{L}$ est inversible
sur $X$, nous obtenons un faisceau inversible $f^*\mathcal{L}$ sur $Y$
(lemme \ref{modules-lemma-pullback-invertible})
et un morphisme induit d'anneaux gradués
$$
f^* :
\Gamma_*(X, \mathcal{L})
\longrightarrow
\Gamma_*(Y, f^*\mathcal{L})
$$
Il existe en outre certaines compatibilités entre les constructions
ci-dessus, dont nous omettons les énoncés.

\begin{lemma}
\label{modules-lemma-pic-set}
Soit $(X, \mathcal{O}_X)$ un espace annelé.
Il existe un ensemble de modules inversibles $\{\mathcal{L}_i\}_{i \in I}$
tel que tout module inversible sur $X$ soit isomorphe à exactement
l'un des $\mathcal{L}_i$.
\end{lemma}

\begin{proof}
Rappelons que tout $\mathcal{O}_X$-module inversible est localement
facteur direct d'un $\mathcal{O}_X$-module libre de type fini ; voir le
lemme \ref{modules-lemma-invertible}.
Pour tout recouvrement ouvert $\mathcal{U} : X = \bigcup_{j \in J} U_j$
et toute application $r : J \to \mathbf{N}$, considérons les faisceaux de
$\mathcal{O}_X$-modules $\mathcal{F}$ tels que
$\mathcal{F}_j = \mathcal{F}|_{U_j}$ soit facteur direct de
$\mathcal{O}_{U_j}^{\oplus r(j)}$.
La collection des classes d'isomorphisme de $\mathcal{F}_j$ est un ensemble, car
$\Hom_{\mathcal{O}_U}(\mathcal{O}_U^{\oplus r}, \mathcal{O}_U^{\oplus r})$
est un ensemble. Le faisceau $\mathcal{F}$ s'obtient par recollement des $\mathcal{F}_j$ ;
voir Faisceaux, section
\ref{sheaves-section-glueing-sheaves}. Remarquons que la collection de
toutes les données de recollement forme un ensemble. La collection de tous les recouvrements
$\mathcal{U} : X = \bigcup_{j \in J} U_j$ pour lesquels $J \to \mathcal{P}(X)$,
$j \mapsto U_j$ est injective forme elle aussi un ensemble. Pour chaque recouvrement,
il existe un ensemble d'applications $r : J \to \mathbf{N}$. La collection
de tous les $\mathcal{F}$ forme donc un ensemble.
\end{proof}

\noindent
Ce lemme affirme, en substance, que la collection des
classes d'isomorphisme de faisceaux inversibles forme un ensemble.
Le lemme \ref{modules-lemma-constructions-invertible} affirme que
le produit tensoriel définit sur cet ensemble une structure de groupe
abélien.

\begin{definition}
\label{modules-definition-pic}
Soit $(X, \mathcal{O}_X)$ un espace annelé.
Le {\it groupe de Picard} $\Pic(X)$ de $X$ est le
groupe abélien dont les éléments sont les classes d'isomorphisme de
$\mathcal{O}_X$-modules inversibles, l'addition
correspondant au produit tensoriel.
\end{definition}

\begin{lemma}
\label{modules-lemma-s-open}
\begin{slogan}
Une trivialisation (locale) d'un fibré en droites
équivaut à une section (locale) ne s'annulant pas.
\end{slogan}
Soit $X$ un espace annelé. Supposons que chaque anneau local $\mathcal{O}_{X, x}$
soit un anneau local d'idéal maximal $\mathfrak m_x$.
Soit $\mathcal{L}$ un $\mathcal{O}_X$-module inversible.
Pour toute section $s \in \Gamma(X, \mathcal{L})$, l'ensemble
$$
X_s = \{x \in X \mid \text{image de }s \not\in \mathfrak m_x\mathcal{L}_x\}
$$
est ouvert dans $X$. Le morphisme $s : \mathcal{O}_{X_s} \to \mathcal{L}|_{X_s}$
est un isomorphisme, et il existe une section $s'$
de $\mathcal{L}^{\otimes -1}$ sur $X_s$ telle que $s' (s|_{X_s}) = 1$.
\end{lemma}

\begin{proof}
Supposons $x \in X_s$.
Nous avons un isomorphisme
$$
\mathcal{L}_x \otimes_{\mathcal{O}_{X, x}} (\mathcal{L}^{\otimes -1})_x
\longrightarrow
\mathcal{O}_{X, x}
$$
d'après le lemme \ref{modules-lemma-constructions-invertible}.
Les modules $\mathcal{L}_x$ et $(\mathcal{L}^{\otimes -1})_x$
sont tous deux des $\mathcal{O}_{X, x}$-modules libres de rang $1$. Le
lemme de Nakayama (Algèbre, lemme \ref{algebra-lemma-NAK}) montre que
$s_x$ est une base de $\mathcal{L}_x$. Il existe donc
un élément de base $t_x \in (\mathcal{L}^{\otimes -1})_x$
tel que $s_x \otimes t_x$ s'envoie sur $1$.
Choisissons un voisinage ouvert $U$ de
$x$ tel que $t_x$ provienne d'une section $t$
de $\mathcal{L}^{\otimes -1}$ sur $U$ et que
$s \otimes t$ s'envoie sur $1 \in \mathcal{O}_X(U)$.
Il est clair que, pour tout $x' \in U$, la section $s$ engendre
le module $\mathcal{L}_{x'}$. Ainsi $U \subset X_s$.
Ceci prouve que $X_s$ est ouvert. De plus, la section
$t$ construite ci-dessus sur $U$ est unique ; ces sections
se recollent donc en la section $s'$ du lemme.
\end{proof}

\begin{remark}
\label{modules-remark-pullback-s-open}
Étant donné un morphisme d'espaces localement annelés $f : Y \to X$
(voir Schémas, définition \ref{schemes-definition-locally-ringed-space}),
l'image inverse $f^{-1}(X_s)$ est égale à $Y_{f^*s}$, où
$f^*s \in \Gamma(Y, f^*\mathcal{L})$ est l'image inverse de $s$.
\end{remark}


\section{Rang et déterminant}
\label{modules-section-rank-and-det}

\noindent
Soit $(X, \mathcal{O}_X)$ un espace annelé. Considérons la catégorie
$\textit{Vect}(X)$ des $\mathcal{O}_X$-modules localement libres de type fini.
C'est une catégorie exacte
(voir Injectifs, remarque \ref{injectives-remark-embed-exact-category})
dont les épimorphismes admissibles sont les
surjections et dont les monomorphismes admissibles sont les noyaux de
surjections. De plus, les classes d'isomorphisme des objets de
$\textit{Vect}(X)$ forment un ensemble (démonstration omise). Nous pouvons donc former
le $K$-groupe de Grothendieck d'indice zéro $K_0(\textit{Vect}(X))$.
Explicitement, dans ce cas, $K_0(\textit{Vect}(X))$
est le groupe abélien engendré par $[\mathcal{E}]$, où $\mathcal{E}$
est un $\mathcal{O}_X$-module localement libre de type fini, soumis aux relations
$$
[\mathcal{E}] = [\mathcal{E}'] + [\mathcal{E}'']
$$
pour toute suite exacte courte
$0 \to \mathcal{E}' \to \mathcal{E} \to \mathcal{E}'' \to 0$
de $\mathcal{O}_X$-modules localement libres de type fini.

\medskip\noindent
{\bf Rangs.} Supposons que tous les anneaux locaux $\mathcal{O}_{X, x}$ soient non nuls.
Étant donné un $\mathcal{O}_X$-module localement libre de type fini $\mathcal{E}$,
le {\it rang} est la fonction localement constante
$$
\text{rang}_\mathcal{E} : X \longrightarrow \mathbf{Z}_{\geq 0},\quad
x \longmapsto \text{rang}_{\mathcal{O}_{X, x}} \mathcal{E}_x
$$
Voir le lemme \ref{modules-lemma-rank}. Par définition des modules localement
libres, la fonction $\text{rang}_\mathcal{E}$ est localement constante. Si
$0 \to \mathcal{E}' \to \mathcal{E} \to \mathcal{E}'' \to 0$
est une suite exacte courte de $\mathcal{O}_X$-modules localement libres de type fini,
alors $\text{rang}_\mathcal{E} = \text{rang}_{\mathcal{E}'} +
\text{rang}_{\mathcal{E}''}$.
Le rang définit donc un morphisme
$$
K_0(\textit{Vect}(X)) \longrightarrow \text{Map}_{cont}(X, \mathbf{Z}),\quad
[\mathcal{E}] \longmapsto \text{rang}_\mathcal{E}
$$

\medskip\noindent
{\bf Déterminants.} Étant donné un
$\mathcal{O}_X$-module localement libre de type fini $\mathcal{E}$, nous obtenons une
décomposition en union disjointe
$$
X = X_0 \amalg X_1 \amalg X_2 \amalg \ldots
$$
où les $X_i$ sont ouverts et fermés, de telle sorte que $\mathcal{E}$ soit localement
libre de type fini et de rang $i$ sur $X_i$ (cela revient exactement à dire que
$\text{rang}_\mathcal{E}$ est localement constante). Dans ce cas, nous définissons
$\det(\mathcal{E})$ comme le faisceau inversible sur $X$ qui est égal à
$\wedge^i(\mathcal{E}|_{X_i})$ sur $X_i$ pour tout $i \geq 0$.
La décomposition ci-dessus étant disjointe, il n'y a aucune condition de recollement
à vérifier. D'après le lemme \ref{modules-lemma-det-ses} ci-dessous,
ceci définit un morphisme
$$
\det : K_0(\textit{Vect}(X)) \longrightarrow \Pic(X),\quad
[\mathcal{E}] \longmapsto \det(\mathcal{E})
$$
de groupes abéliens. Les éléments de $\Pic(X)$ ainsi obtenus
sont localement libres de rang $1$ (voir, après le lemme, une généralisation).

\begin{lemma}
\label{modules-lemma-det-ses}
Soit $X$ un espace annelé. Soit
$0 \to \mathcal{E}' \to \mathcal{E} \to \mathcal{E}'' \to 0$
une suite exacte courte de $\mathcal{O}_X$-modules localement libres de type fini.
Il existe alors un isomorphisme canonique
$$
\det(\mathcal{E}') \otimes_{\mathcal{O}_X}\det(\mathcal{E}'')
\longrightarrow
\det(\mathcal{E})
$$
de $\mathcal{O}_X$-modules.
\end{lemma}

\begin{proof}
Nous pouvons décomposer $X$ en sous-ensembles ouverts et fermés disjoints sur lesquels
$\mathcal{E}'$ et $\mathcal{E}''$ sont tous deux de rang constant.
Nous nous ramenons donc au cas où $\mathcal{E}'$ et $\mathcal{E}''$
sont de rang constant, disons $r'$ et $r''$. Dans cette situation, nous
définissons
$$
\wedge^{r'}(\mathcal{E}') \otimes_{\mathcal{O}_X} \wedge^{r''}(\mathcal{E}'')
\longrightarrow
\wedge^{r' + r''}(\mathcal{E})
$$
comme suit. Étant données des sections locales $s'_1, \ldots, s'_{r'}$ de $\mathcal{E}'$
et des sections locales $s''_1, \ldots, s''_{r''}$ de $\mathcal{E}''$,
nous envoyons
$$
s'_1 \wedge \ldots \wedge s'_{r'} \otimes
s''_1 \wedge \ldots \wedge s''_{r''}
\quad\text{sur}\quad
s'_1 \wedge \ldots \wedge s'_{r'} \wedge
\tilde s''_1 \wedge \ldots \wedge \tilde s''_{r''}
$$
où $\tilde s''_i$ est un relèvement local de la section
$s''_i$ en une section de $\mathcal{E}$. Nous omettons les détails.
\end{proof}

\noindent
Soit $(X, \mathcal{O}_X)$ un espace annelé. Au lieu de considérer les
$\mathcal{O}_X$-modules localement libres de type fini, nous pouvons considérer les
$\mathcal{O}_X$-modules $\mathcal{F}$ qui sont localement sur $X$
facteurs directs d'un $\mathcal{O}_X$-module libre de type fini. Cela revient à
demander que $\mathcal{F}$ soit un $\mathcal{O}_X$-module plat et de
présentation finie ; voir le lemme \ref{modules-lemma-flat-locally-finite-presentation}.
Si tous les anneaux locaux $\mathcal{O}_{X, x}$ sont locaux, alors un tel module
$\mathcal{F}$ est localement libre de type fini ; voir le
lemme \ref{modules-lemma-direct-summand-of-locally-free-is-locally-free}.
Ce n'est toutefois pas le cas en général ; par
exemple, $X$ pourrait être un point, $\Gamma(X, \mathcal{O}_X)$ pourrait
être le produit $A \times B$ de deux anneaux non nuls et $\mathcal{F}$
pourrait correspondre à $A \times 0$.
Pour un tel module, la fonction rang n'est donc pas définie.
Néanmoins, il est encore possible de définir $\det(\mathcal{F})$,
qui sera un $\mathcal{O}_X$-module inversible au sens de la
définition \ref{modules-definition-invertible} (sans être nécessairement localement
libre de rang $1$). Notre construction coïncidera avec celle donnée ci-dessus
lorsque $\mathcal{F}$ est localement libre de type fini.
Nous recommandons au lecteur de sauter le reste de cette section.

\begin{lemma}
\label{modules-lemma-determinant-as-socle}
Soit $(X, \mathcal{O}_X)$ un espace annelé. Soit $\mathcal{F}$
un $\mathcal{O}_X$-module plat et de présentation finie.
Notons
$$
\det(\mathcal{F}) \subset
\wedge^*_{\mathcal{O}_X}(\mathcal{F})
$$
l'annulateur de $\mathcal{F} \subset \wedge^*_{\mathcal{O}_X}(\mathcal{F})$.
Alors $\det(\mathcal{F})$ est un $\mathcal{O}_X$-module inversible.
\end{lemma}

\begin{proof}
Pour le démontrer, nous pouvons travailler localement sur $X$. Nous pouvons donc supposer que $\mathcal{F}$
est facteur direct d'un module libre de type fini ; voir le
lemme \ref{modules-lemma-flat-locally-finite-presentation}. Écrivons
$\mathcal{F} \oplus \mathcal{G} = \mathcal{O}_X^{\oplus n}$.
Posons $R = \mathcal{O}_X(X)$. Nous avons alors
$\mathcal{F}(X) \oplus \mathcal{G}(X) = R^{\oplus n}$
et, de même,
$\mathcal{F}(U) \oplus \mathcal{G}(U) = \mathcal{O}_X(U)^{\oplus n}$
pour tout ouvert $U \subset X$.
Nous en concluons que $\mathcal{F} = \mathcal{F}_M$ comme dans le
lemme \ref{modules-lemma-construct-quasi-coherent-sheaves},
où $M = \mathcal{F}(X)$ est un $R$-module projectif de type fini.
Autrement dit, nous avons $\mathcal{F}(U) = M \otimes_R \mathcal{O}_X(U)$.
Ceci implique que
$\det(M) \otimes_R \mathcal{O}_X(U) = \det(\mathcal{F}(U))$
pour tout ouvert $U \subset X$, où $\det$ est celui de
Compléments d'algèbre, section \ref{more-algebra-section-determinants}. D'après
Compléments d'algèbre, remarque \ref{more-algebra-remark-determinant-as-socle},
nous voyons que
$$
\det(M) \otimes_R \mathcal{O}_X(U) =
\det(\mathcal{F}(U)) \subset
\wedge^*_{\mathcal{O}_X(U)}(\mathcal{F}(U))
$$
est l'annulateur de $\mathcal{F}(U)$. Nous en concluons que
$\det(\mathcal{F})$, tel qu'il est défini dans l'énoncé du lemme,
est égal à $\mathcal{F}_{\det(M)}$. Nous omettons certains détails : il faut
prendre garde au fait que les annulateurs ne peuvent pas être définis en faisceautisant
la prise des annulateurs sur les sections au-dessus des ouverts. Ainsi $\det(\mathcal{F})$
est l'image inverse d'un module inversible, ce qui conclut la démonstration.
\end{proof}
\section{Localisation des faisceaux d'anneaux}
\label{modules-section-localizing-sheaves-rings}

\noindent
Soit $X$ un espace topologique et soit $\mathcal{O}_X$ un
préfaisceau d'anneaux. Soit $\mathcal{S} \subset \mathcal{O}_X$
un préfaisceau d'ensembles contenu dans $\mathcal{O}_X$.
Supposons que, pour tout ouvert $U \subset X$, l'ensemble
$\mathcal{S}(U) \subset \mathcal{O}_X(U)$ soit une partie multiplicative;
voir Algèbre, définition \ref{algebra-definition-multiplicative-subset}.
On peut alors considérer le préfaisceau d'anneaux
$$
\mathcal{S}^{-1}\mathcal{O}_X :
U \longmapsto \mathcal{S}(U)^{-1}\mathcal{O}_X(U).
$$
L'application de restriction envoie la section $f/s$, où
$f \in \mathcal{O}_X(U)$ et $s \in \mathcal{S}(U)$, sur
$(f|_V)/(s|_V)$ lorsque $V \subset U$ sont des ouverts de $X$.

\begin{lemma}
\label{modules-lemma-simple-invert}
Soit $X$ un espace topologique et soit $\mathcal{O}_X$ un
préfaisceau d'anneaux. Soit $\mathcal{S} \subset \mathcal{O}_X$
un préfaisceau d'ensembles contenu dans $\mathcal{O}_X$.
Supposons que, pour tout ouvert $U \subset X$, l'ensemble
$\mathcal{S}(U) \subset \mathcal{O}_X(U)$ soit une partie multiplicative.
\begin{enumerate}
\item Il existe un morphisme de préfaisceaux d'anneaux
$\mathcal{O}_X \to \mathcal{S}^{-1}\mathcal{O}_X$
tel que toute section locale de $\mathcal{S}$ ait pour image une section
inversible de $\mathcal{S}^{-1}\mathcal{O}_X$.
\item Pour tout homomorphisme de préfaisceaux d'anneaux
$\mathcal{O}_X \to \mathcal{A}$ tel que toute section locale
de $\mathcal{S}$ ait pour image une section inversible de $\mathcal{A}$,
il existe une unique factorisation
$\mathcal{S}^{-1}\mathcal{O}_X \to \mathcal{A}$.
\item Pour tout $x \in X$, on a
$$
(\mathcal{S}^{-1}\mathcal{O}_X)_x = \mathcal{S}_x^{-1} \mathcal{O}_{X, x}.
$$
\item Le faisceau associé $(\mathcal{S}^{-1}\mathcal{O}_X)^\#$ est un
faisceau d'anneaux muni d'un morphisme de faisceaux d'anneaux
$(\mathcal{O}_X)^\# \to (\mathcal{S}^{-1}\mathcal{O}_X)^\#$
qui est universel parmi les morphismes de $(\mathcal{O}_X)^\#$ vers des
faisceaux d'anneaux envoyant toute section locale de $\mathcal{S}$ sur une
section inversible.
\item Pour tout $x \in X$, on a
$$
(\mathcal{S}^{-1}\mathcal{O}_X)^\#_x = \mathcal{S}_x^{-1} \mathcal{O}_{X, x}.
$$
\end{enumerate}
\end{lemma}

\begin{proof}
Démonstration omise.
\end{proof}

\noindent
Soit $X$ un espace topologique et soit $\mathcal{O}_X$ un
préfaisceau d'anneaux. Soit $\mathcal{S} \subset \mathcal{O}_X$
un préfaisceau d'ensembles contenu dans $\mathcal{O}_X$.
Supposons que, pour tout ouvert $U \subset X$, l'ensemble
$\mathcal{S}(U) \subset \mathcal{O}_X(U)$ soit une partie multiplicative.
Soit $\mathcal{F}$ un préfaisceau de $\mathcal{O}_X$-modules.
On peut alors considérer le préfaisceau de
$\mathcal{S}^{-1}\mathcal{O}_X$-modules
$$
\mathcal{S}^{-1}\mathcal{F} :
U \longmapsto \mathcal{S}(U)^{-1}\mathcal{F}(U).
$$
L'application de restriction envoie la section $t/s$, où
$t \in \mathcal{F}(U)$ et $s \in \mathcal{S}(U)$, sur
$(t|_V)/(s|_V)$ lorsque $V \subset U$ sont des ouverts de $X$.

\begin{lemma}
\label{modules-lemma-simple-invert-module}
Soit $X$ un espace topologique.
Soit $\mathcal{O}_X$ un préfaisceau d'anneaux.
Soit $\mathcal{S} \subset \mathcal{O}_X$ un préfaisceau d'ensembles contenu
dans $\mathcal{O}_X$. Supposons que, pour tout ouvert $U \subset X$,
l'ensemble $\mathcal{S}(U) \subset \mathcal{O}_X(U)$ soit une partie
multiplicative. Pour tout préfaisceau de $\mathcal{O}_X$-modules
$\mathcal{F}$, on a
$$
\mathcal{S}^{-1}\mathcal{F}
=
\mathcal{S}^{-1}\mathcal{O}_X \otimes_{p, \mathcal{O}_X} \mathcal{F}
$$
(voir Faisceaux, section \ref{sheaves-section-presheaves-modules} pour la
notation) et, si $\mathcal{F}$ et $\mathcal{O}_X$ sont des faisceaux, alors
$$
(\mathcal{S}^{-1}\mathcal{F})^\#
=
(\mathcal{S}^{-1}\mathcal{O}_X)^\# \otimes_{\mathcal{O}_X} \mathcal{F}
$$
(voir Faisceaux, section
\ref{sheaves-section-sheafification-presheaves-modules} pour la notation).
\end{lemma}

\begin{proof}
Démonstration omise.
\end{proof}








\section{Modules de différentielles}
\label{modules-section-differentials}

\noindent
Dans cette section, nous expliquons brièvement comment définir le module des
différentielles relatives d'un morphisme d'espaces annelés. Nous conseillons
au lecteur de consulter la section correspondante du chapitre d'algèbre
commutative (Algèbre, section \ref{algebra-section-differentials}).

\begin{definition}
\label{modules-definition-derivation}
Soit $X$ un espace topologique. Soit
$\varphi : \mathcal{O}_1 \to \mathcal{O}_2$ un homomorphisme de faisceaux
d'anneaux. Soit $\mathcal{F}$ un $\mathcal{O}_2$-module. Une
{\it $\mathcal{O}_1$-dérivation}, ou plus précisément une
{\it $\varphi$-dérivation}, à valeurs dans $\mathcal{F}$ est une application
$D : \mathcal{O}_2 \to \mathcal{F}$ additive, nulle sur l'image de
$\mathcal{O}_1 \to \mathcal{O}_2$, et satisfaisant la
{\it règle de Leibniz}
$$
D(ab) = aD(b) + D(a)b
$$
pour toutes sections locales $a, b$ de $\mathcal{O}_2$ (sur tout ouvert où
elles sont toutes deux définies). On note
$\text{Der}_{\mathcal{O}_1}(\mathcal{O}_2, \mathcal{F})$ l'ensemble des
$\varphi$-dérivations à valeurs dans $\mathcal{F}$.
\end{definition}

\noindent
C'est l'analogue faisceautique d'Algèbre, définition
\ref{algebra-definition-derivation}. Étant donnée une dérivation
$D : \mathcal{O}_2 \to \mathcal{F}$ comme dans la définition,
l'application induite sur les sections globales
$$
D : \Gamma(X, \mathcal{O}_2) \longrightarrow \Gamma(X, \mathcal{F})
$$
est une $\Gamma(X, \mathcal{O}_1)$-dérivation au sens de la définition
algébrique. Remarquons que, si
$\alpha : \mathcal{F} \to \mathcal{G}$ est un morphisme de
$\mathcal{O}_2$-modules, il existe une application induite
$$
\text{Der}_{\mathcal{O}_1}(\mathcal{O}_2, \mathcal{F})
\longrightarrow
\text{Der}_{\mathcal{O}_1}(\mathcal{O}_2, \mathcal{G})
$$
donnée par $D \mapsto \alpha \circ D$. Autrement dit, on obtient un
foncteur.

\begin{lemma}
\label{modules-lemma-universal-module}
Soit $X$ un espace topologique. Soit
$\varphi : \mathcal{O}_1 \to \mathcal{O}_2$ un homomorphisme de faisceaux
d'anneaux. Le foncteur
$$
\textit{Mod}(\mathcal{O}_2) \longrightarrow \textit{Ab}, \quad
\mathcal{F} \longmapsto \text{Der}_{\mathcal{O}_1}(\mathcal{O}_2, \mathcal{F})
$$
est représentable.
\end{lemma}

\begin{proof}
La démonstration est exactement la même que celle de l'énoncé analogue en
algèbre. Au cours de cette démonstration, pour tout faisceau d'ensembles
$\mathcal{F}$ sur $X$, notons $\mathcal{O}_2[\mathcal{F}]$ le faisceau
associé au préfaisceau $U \mapsto \mathcal{O}_2(U)[\mathcal{F}(U)]$, où
cette dernière expression désigne le $\mathcal{O}_2(U)$-module libre de base
l'ensemble $\mathcal{F}(U)$. Pour $s \in \mathcal{F}(U)$, on note $[s]$ la
section correspondante de $\mathcal{O}_2[\mathcal{F}]$ sur $U$. Si
$\mathcal{F}$ est un faisceau de $\mathcal{O}_2$-modules, il existe un
morphisme canonique
$$
c : \mathcal{O}_2[\mathcal{F}] \longrightarrow \mathcal{F}
$$
qui, au niveau des préfaisceaux, est donné par la règle
$\sum f_s[s] \mapsto \sum f_s s$. Nous emploierons l'abréviation
$[s] \mapsto s$ pour décrire ce morphisme, et des abréviations analogues
pour les morphismes ci-dessous. Considérons le morphisme de
$\mathcal{O}_2$-modules
\begin{equation}
\label{modules-equation-define-module-differentials}
\begin{matrix}
\mathcal{O}_2[\mathcal{O}_2 \times \mathcal{O}_2] \oplus
\mathcal{O}_2[\mathcal{O}_2 \times \mathcal{O}_2] \oplus
\mathcal{O}_2[\mathcal{O}_1] &
\longrightarrow &
\mathcal{O}_2[\mathcal{O}_2] \\
[(a, b)] \oplus [(f, g)] \oplus [h] & \longmapsto & [a + b] - [a] - [b] + \\
& & [fg] - g[f] - f[g] + \\
& & [\varphi(h)]
\end{matrix}
\end{equation}
avec la notation abrégée ci-dessus. Posons
$\Omega_{\mathcal{O}_2/\mathcal{O}_1}$ égal au conoyau de ce morphisme.
Il est alors clair qu'il existe un morphisme de faisceaux d'ensembles
$$
\text{d} : \mathcal{O}_2 \longrightarrow \Omega_{\mathcal{O}_2/\mathcal{O}_1}
$$
qui envoie une section locale $f$ sur l'image de $[f]$ dans
$\Omega_{\mathcal{O}_2/\mathcal{O}_1}$. Par construction, $\text{d}$ est une
$\mathcal{O}_1$-dérivation. Soit ensuite $\mathcal{F}$ un faisceau de
$\mathcal{O}_2$-modules et soit
$D : \mathcal{O}_2 \to \mathcal{F}$ une $\mathcal{O}_1$-dérivation.
On peut considérer le morphisme $\mathcal{O}_2$-linéaire
$\mathcal{O}_2[\mathcal{O}_2] \to \mathcal{F}$ qui envoie $[g]$ sur $D(g)$.
La définition d'une dérivation montre que ce morphisme annule les sections
appartenant à l'image du morphisme
(\ref{modules-equation-define-module-differentials}); il définit donc un morphisme
$$
\alpha_D : \Omega_{\mathcal{O}_2/\mathcal{O}_1} \longrightarrow \mathcal{F}
$$
Comme $D = \alpha_D \circ \text{d}$, le lemme est démontré.
\end{proof}

\begin{definition}
\label{modules-definition-module-differentials}
Soit $X$ un espace topologique. Soit
$\varphi : \mathcal{O}_1 \to \mathcal{O}_2$ un homomorphisme de faisceaux
d'anneaux sur $X$. Le {\it module des différentielles} de $\varphi$ est
l'objet représentant le foncteur
$\mathcal{F} \mapsto \text{Der}_{\mathcal{O}_1}(\mathcal{O}_2, \mathcal{F})$
dont l'existence résulte du lemme \ref{modules-lemma-universal-module}.
On le note $\Omega_{\mathcal{O}_2/\mathcal{O}_1}$, et la
{\it $\varphi$-dérivation universelle} se note
$\text{d} : \mathcal{O}_2 \to \Omega_{\mathcal{O}_2/\mathcal{O}_1}$.
\end{definition}

\noindent
Remarquons que $\Omega_{\mathcal{O}_2/\mathcal{O}_1}$ est le conoyau du
morphisme de $\mathcal{O}_2$-modules
(\ref{modules-equation-define-module-differentials}). De plus, le morphisme
$\text{d}$ est caractérisé par le fait que $\text{d}f$ est l'image de la
section locale $[f]$.

\begin{lemma}
\label{modules-lemma-differentials-sheafify}
Soit $X$ un espace topologique. Soit
$\varphi : \mathcal{O}_1 \to \mathcal{O}_2$ un homomorphisme de faisceaux
d'anneaux sur $X$. Alors $\Omega_{\mathcal{O}_2/\mathcal{O}_1}$ est le
faisceau associé au préfaisceau
$U \mapsto \Omega_{\mathcal{O}_2(U)/\mathcal{O}_1(U)}$.
\end{lemma}

\begin{proof}
Considérons le morphisme (\ref{modules-equation-define-module-differentials}). Il
existe un morphisme analogue de préfaisceaux dont la valeur sur l'ouvert
$U$ est
$$
\mathcal{O}_2(U)[\mathcal{O}_2(U) \times \mathcal{O}_2(U)] \oplus
\mathcal{O}_2(U)[\mathcal{O}_2(U) \times \mathcal{O}_2(U)] \oplus
\mathcal{O}_2(U)[\mathcal{O}_1(U)]
\longrightarrow
\mathcal{O}_2(U)[\mathcal{O}_2(U)]
$$
Par la construction du module des différentielles dans Algèbre, définition
\ref{algebra-definition-differentials}, le conoyau de ce morphisme a pour
valeur $\Omega_{\mathcal{O}_2(U)/\mathcal{O}_1(U)}$ sur $U$. D'autre part,
les faisceaux figurant dans (\ref{modules-equation-define-module-differentials}) sont
les faisceaux associés aux préfaisceaux ci-dessus. Le résultat découle donc
de l'exactitude de la faisceautisation.
\end{proof}

\begin{lemma}
\label{modules-lemma-localize-differentials}
Soit $X$ un espace topologique. Soit
$\varphi : \mathcal{O}_1 \to \mathcal{O}_2$ un homomorphisme de faisceaux
d'anneaux. Pour tout ouvert $U \subset X$, il existe un isomorphisme canonique
$$
\Omega_{\mathcal{O}_2/\mathcal{O}_1}|_U =
\Omega_{(\mathcal{O}_2|_U)/(\mathcal{O}_1|_U)}
$$
compatible aux dérivations universelles.
\end{lemma}

\begin{proof}
Cela résulte du fait que $\Omega_{\mathcal{O}_2/\mathcal{O}_1}$ est le
conoyau du morphisme (\ref{modules-equation-define-module-differentials}).
\end{proof}

\begin{lemma}
\label{modules-lemma-pullback-differentials}
Soit $f : Y \to X$ une application continue d'espaces topologiques.
Soit $\varphi : \mathcal{O}_1 \to \mathcal{O}_2$ un homomorphisme de
faisceaux d'anneaux sur $X$. Il existe alors une identification canonique
$f^{-1}\Omega_{\mathcal{O}_2/\mathcal{O}_1} =
\Omega_{f^{-1}\mathcal{O}_2/f^{-1}\mathcal{O}_1}$
compatible aux dérivations universelles.
\end{lemma}

\begin{proof}
Cela résulte du fait que le faisceau
$\Omega_{\mathcal{O}_2/\mathcal{O}_1}$ est le conoyau du morphisme
(\ref{modules-equation-define-module-differentials}) et qu'un énoncé analogue vaut
pour $\Omega_{f^{-1}\mathcal{O}_2/f^{-1}\mathcal{O}_1}$, puisque le foncteur
$f^{-1}$ est exact et que
$f^{-1}(\mathcal{O}_2[\mathcal{O}_2]) =
f^{-1}\mathcal{O}_2[f^{-1}\mathcal{O}_2]$,
$f^{-1}(\mathcal{O}_2[\mathcal{O}_2 \times \mathcal{O}_2]) =
f^{-1}\mathcal{O}_2[f^{-1}\mathcal{O}_2 \times f^{-1}\mathcal{O}_2]$, et
$f^{-1}(\mathcal{O}_2[\mathcal{O}_1]) =
f^{-1}\mathcal{O}_2[f^{-1}\mathcal{O}_1]$.
\end{proof}

\begin{lemma}
\label{modules-lemma-stalk-module-differentials}
Soit $X$ un espace topologique. Soit
$\mathcal{O}_1 \to \mathcal{O}_2$ un homomorphisme de faisceaux d'anneaux
sur $X$. Pour $x \in X$, on a
$\Omega_{\mathcal{O}_2/\mathcal{O}_1, x} =
\Omega_{\mathcal{O}_{2, x}/\mathcal{O}_{1, x}}$.
\end{lemma}

\begin{proof}
C'est le cas particulier du lemme \ref{modules-lemma-pullback-differentials}
correspondant à l'inclusion $\{x\} \to X$. On peut aussi utiliser le lemme
\ref{modules-lemma-differentials-sheafify}, Faisceaux, lemme
\ref{sheaves-lemma-stalk-sheafification}, et Algèbre, lemme
\ref{algebra-lemma-colimit-differentials}.
\end{proof}

\begin{lemma}
\label{modules-lemma-functoriality-differentials}
Soit $X$ un espace topologique. Soit
$$
\xymatrix{
\mathcal{O}_2 \ar[r]_\varphi & \mathcal{O}_2' \\
\mathcal{O}_1 \ar[r] \ar[u] & \mathcal{O}'_1 \ar[u]
}
$$
un diagramme commutatif de faisceaux d'anneaux sur $X$. La composée de
$\mathcal{O}_2 \to \mathcal{O}'_2$ avec le morphisme
$\text{d} : \mathcal{O}'_2 \to \Omega_{\mathcal{O}'_2/\mathcal{O}'_1}$
est une $\mathcal{O}_1$-dérivation. On obtient donc un morphisme canonique
de $\mathcal{O}_2$-modules
$\Omega_{\mathcal{O}_2/\mathcal{O}_1} \to
\Omega_{\mathcal{O}'_2/\mathcal{O}'_1}$.
Il est uniquement caractérisé par la propriété
$\text{d}(f) \mapsto \text{d}(\varphi(f))$
pour toute section locale $f$ de $\mathcal{O}_2$. Ainsi,
$\Omega_{-/-}$ devient un foncteur sur la catégorie des flèches de faisceaux
d'anneaux.
\end{lemma}

\begin{proof}
L'énoncé découle immédiatement des définitions.
\end{proof}

\begin{lemma}
\label{modules-lemma-differential-seq}
Dans le lemme \ref{modules-lemma-functoriality-differentials}, supposons que
$\mathcal{O}_2 \to \mathcal{O}'_2$ soit surjectif, de noyau
$\mathcal{I} \subset \mathcal{O}_2$, et que
$\mathcal{O}_1 = \mathcal{O}'_1$. Il existe alors une suite exacte canonique
de $\mathcal{O}'_2$-modules
$$
\mathcal{I}/\mathcal{I}^2
\longrightarrow
\Omega_{\mathcal{O}_2/\mathcal{O}_1} \otimes_{\mathcal{O}_2} \mathcal{O}'_2
\longrightarrow
\Omega_{\mathcal{O}'_2/\mathcal{O}_1}
\longrightarrow
0
$$
Le morphisme de gauche est caractérisé par le fait qu'une section locale
$f$ de $\mathcal{I}$ est envoyée sur $\text{d}f \otimes 1$.
\end{lemma}

\begin{proof}
Pour une section locale $f$ de $\mathcal{I}$, notons $\overline{f}$ l'image
de $f$ dans $\mathcal{I}/\mathcal{I}^2$. Pour montrer que l'application
$\overline{f} \mapsto \text{d}f \otimes 1$ est bien définie, il suffit de
vérifier que $\text{d} f_1f_2 \otimes 1 = 0$ lorsque $f_1, f_2$ sont des
sections locales de $\mathcal{I}$. Cela découle de la règle de Leibniz :
$\text{d} f_1f_2 \otimes 1 =
(f_1 \text{d}f_2 + f_2 \text{d} f_1 )\otimes 1 =
\text{d}f_2 \otimes f_1 + \text{d}f_1 \otimes f_2 = 0$.
Un calcul analogue montre que cette application est
$\mathcal{O}'_2 = \mathcal{O}_2/\mathcal{I}$-linéaire. Le morphisme de droite
est celui du lemme \ref{modules-lemma-functoriality-differentials}. Pour vérifier
l'exactitude de la suite, on peut raisonner sur les fibres (lemme
\ref{modules-lemma-abelian}). Le lemme \ref{modules-lemma-stalk-module-differentials} ramène
alors l'assertion à Algèbre, lemme \ref{algebra-lemma-differential-seq}.
\end{proof}

\begin{definition}
\label{modules-definition-differentials}
Soit $(f, f^\sharp) : (X, \mathcal{O}_X) \to (S, \mathcal{O}_S)$
un morphisme d'espaces annelés.
\begin{enumerate}
\item Soit $\mathcal{F}$ un $\mathcal{O}_X$-module. Une
{\it $S$-dérivation} à valeurs dans $\mathcal{F}$ est une
$f^{-1}\mathcal{O}_S$-dérivation, ou plus précisément une
$f^\sharp$-dérivation au sens de la définition
\ref{modules-definition-derivation}. On note
$\text{Der}_S(\mathcal{O}_X, \mathcal{F})$ l'ensemble des $S$-dérivations
à valeurs dans $\mathcal{F}$.
\item Le {\it faisceau des différentielles $\Omega_{X/S}$ de $X$ sur $S$}
est le module des différentielles
$\Omega_{\mathcal{O}_X/f^{-1}\mathcal{O}_S}$ muni de sa $S$-dérivation
universelle $\text{d}_{X/S} : \mathcal{O}_X \to \Omega_{X/S}$.
\end{enumerate}
\end{definition}

\noindent
Voici une situation particulière dans laquelle les dérivations apparaissent
naturellement.

\begin{lemma}
\label{modules-lemma-double-structure-gives-derivation}
Soit $(f, f^\sharp) : (X, \mathcal{O}_X) \to (S, \mathcal{O}_S)$
un morphisme d'espaces annelés. Considérons une suite exacte courte
$$
0 \to \mathcal{I} \to \mathcal{A} \to \mathcal{O}_X \to 0
$$
Ici, $\mathcal{A}$ est un faisceau de $f^{-1}\mathcal{O}_S$-algèbres,
$\pi : \mathcal{A} \to \mathcal{O}_X$ est un morphisme surjectif de
faisceaux de $f^{-1}\mathcal{O}_S$-algèbres, et
$\mathcal{I} = \Ker(\pi)$ est son noyau. Supposons que $\mathcal{I}$ soit
un faisceau d'idéaux de carré nul dans $\mathcal{A}$. Alors $\mathcal{I}$
possède une structure naturelle de $\mathcal{O}_X$-module.
Une section $s : \mathcal{O}_X \to \mathcal{A}$ de $\pi$ est un morphisme
de $f^{-1}\mathcal{O}_S$-algèbres tel que $\pi \circ s = \text{id}$.
Étant données une section $s : \mathcal{O}_X \to \mathcal{A}$ de $\pi$
et une $S$-dérivation $D : \mathcal{O}_X \to \mathcal{I}$, l'application
$$
s + D : \mathcal{O}_X \to \mathcal{A}
$$
est une section de $\pi$, et toute section $s'$ est de la forme $s + D$ pour
une unique $S$-dérivation $D$.
\end{lemma}

\begin{proof}
Rappelons que la structure de $\mathcal{O}_X$-module sur $\mathcal{I}$ est
donnée par $h \tau = \tilde h \tau$ (produit dans $\mathcal{A}$), où $h$
est une section locale de $\mathcal{O}_X$, où $\tilde h$ est un relevé
local de $h$ à $\mathcal{A}$, et où $\tau$ est une section locale de
$\mathcal{I}$. En particulier, une fois $s$ fixée, on peut prendre
$\tilde h = s(h)$. Pour vérifier que $s + D$ est un homomorphisme de
faisceaux d'anneaux, calculons
\begin{eqnarray*}
(s + D)(ab) & = & s(ab) + D(ab) \\
& = & s(a)s(b) + aD(b) + D(a)b \\
& = & s(a) s(b) + s(a)D(b) + D(a)s(b) \\
& = & (s(a) + D(a))(s(b) + D(b))
\end{eqnarray*}
par la règle de Leibniz. De même, on montre que $s + D$ est un morphisme de
$f^{-1}\mathcal{O}_S$-algèbres, puisque $D$ est une $S$-dérivation.
Réciproquement, étant donnée $s'$, on pose $D = s' - s$. Les détails sont omis.
\end{proof}

\begin{lemma}
\label{modules-lemma-functoriality-differentials-ringed-spaces}
Soit
$$
\xymatrix{
X' \ar[d]_{h'} \ar[r]_f & X \ar[d]^h \\
S' \ar[r]^g & S
}
$$
un diagramme commutatif d'espaces annelés.
\begin{enumerate}
\item La composée du morphisme canonique
$\mathcal{O}_X \to f_*\mathcal{O}_{X'}$ avec
$f_*\text{d}_{X'/S'} : f_*\mathcal{O}_{X'} \to f_*\Omega_{X'/S'}$ est une
$S$-dérivation; on obtient donc un morphisme canonique de
$\mathcal{O}_X$-modules
$\Omega_{X/S} \to f_*\Omega_{X'/S'}$.
\item Le diagramme commutatif
$$
\xymatrix{
f^{-1}\mathcal{O}_X \ar[r] & \mathcal{O}_{X'} \\
f^{-1}h^{-1}\mathcal{O}_S \ar[u] \ar[r] & (h')^{-1}\mathcal{O}_{S'} \ar[u]
}
$$
induit, par les lemmes \ref{modules-lemma-pullback-differentials} et
\ref{modules-lemma-functoriality-differentials}, un morphisme canonique
$f^{-1}\Omega_{X/S} \to \Omega_{X'/S'}$.
\end{enumerate}
Ces deux morphismes correspondent (par adjonction entre $f_*$ et $f^*$, et
en utilisant $f^*\Omega_{X/S} =
f^{-1}\Omega_{X/S} \otimes_{f^{-1}\mathcal{O}_X} \mathcal{O}_{X'}$ et
Faisceaux, lemme \ref{sheaves-lemma-adjointness-tensor-restrict})
au même homomorphisme de $\mathcal{O}_{X'}$-modules
$$
c_f : f^*\Omega_{X/S} \longrightarrow \Omega_{X'/S'}
$$
uniquement caractérisé par la propriété que
$f^*\text{d}_{X/S}(a)$ soit envoyé sur $\text{d}_{X'/S'}(f^*a)$
pour toute section locale $a$ de $\mathcal{O}_X$.
\end{lemma}

\begin{proof}
Démonstration omise.
\end{proof}

\begin{lemma}
\label{modules-lemma-check-functoriality-differentials}
Soit
$$
\xymatrix{
X'' \ar[d] \ar[r]_g & X' \ar[d] \ar[r]_f & X \ar[d] \\
S'' \ar[r] & S' \ar[r] & S
}
$$
un diagramme commutatif d'espaces annelés. Avec les notations du lemme
\ref{modules-lemma-functoriality-differentials-ringed-spaces}, on a
$$
c_{f \circ g} = c_g \circ g^* c_f
$$
comme morphismes $(f \circ g)^*\Omega_{X/S} \to \Omega_{X''/S''}$.
\end{lemma}

\begin{proof}
Démonstration omise.
\end{proof}

\begin{lemma}
\label{modules-lemma-triangle-differentials}
Soient $f : X \to Y$ et $g : Y \to S$ des morphismes d'espaces annelés.
Il existe alors une suite exacte canonique
$$
f^*\Omega_{Y/S} \to \Omega_{X/S} \to \Omega_{X/Y} \to 0
$$
dont les flèches proviennent d'applications du lemme
\ref{modules-lemma-functoriality-differentials-ringed-spaces}.
\end{lemma}

\begin{proof}
En prenant les morphismes induits sur les fibres en $x \in X$ et en utilisant
le lemme \ref{modules-lemma-stalk-module-differentials}, on obtient la suite
$$
\mathcal{O}_{X, x} \otimes_{\mathcal{O}_{Y, f(x)}}
\Omega_{\mathcal{O}_{Y, f(x)}/\mathcal{O}_{S, g(f(x))}} \to
\Omega_{\mathcal{O}_{X, x}/\mathcal{O}_{S, g(f(x))}} \to
\Omega_{\mathcal{O}_{X, x}/\mathcal{O}_{Y, f(x)}} \to 0
$$
Il suffit de vérifier que les flèches de cette suite sont celles d'Algèbre,
lemme \ref{algebra-lemma-exact-sequence-differentials}. Cela résulte de la
caractérisation des flèches dans le lemme
\ref{modules-lemma-functoriality-differentials-ringed-spaces} et du fait que,
via l'isomorphisme de Faisceaux, lemme
\ref{sheaves-lemma-stalk-pullback-modules}, on a
$(f^*s)_x = s_x \otimes 1$ pour toute section locale $s$ d'un faisceau de
$\mathcal{O}_Y$-modules.
\end{proof}


















\section{Opérateurs différentiels d'ordre fini}
\label{modules-section-differential-operators}

\noindent
Dans cette section, nous introduisons les opérateurs différentiels d'ordre
fini. Nous conseillons au lecteur de consulter la section correspondante du
chapitre d'algèbre commutative (Algèbre, section
\ref{algebra-section-differential-operators}).

\begin{definition}
\label{modules-definition-differential-operators}
Soit $X$ un espace topologique. Soit
$\varphi : \mathcal{O}_1 \to \mathcal{O}_2$ un homomorphisme de faisceaux
d'anneaux sur $X$. Soit $k \geq 0$ un entier. Soient $\mathcal{F}$ et
$\mathcal{G}$ des faisceaux de $\mathcal{O}_2$-modules. Un
{\it opérateur différentiel $D : \mathcal{F} \to \mathcal{G}$ d'ordre $k$}
est une application $\mathcal{O}_1$-linéaire telle que, pour toute section
locale $g$ de $\mathcal{O}_2$, l'application
$s \mapsto D(gs) - gD(s)$ soit un opérateur différentiel d'ordre $k - 1$.
Dans le cas initial $k = 0$, on définit un opérateur différentiel d'ordre $0$
comme une application $\mathcal{O}_2$-linéaire.
\end{definition}

\noindent
Si $D : \mathcal{F} \to \mathcal{G}$ est un opérateur différentiel d'ordre
$k$, alors, pour toute section locale $g$ de $\mathcal{O}_2$, l'application
$gD$ est encore un opérateur différentiel d'ordre $k$. La somme de deux
opérateurs différentiels d'ordre $k$ en est encore un. Par conséquent,
l'ensemble de tous ces opérateurs
$$
\text{Diff}^k(\mathcal{F}, \mathcal{G}) =
\text{Diff}^k_{\mathcal{O}_2/\mathcal{O}_1}(\mathcal{F}, \mathcal{G})
$$
est un $\Gamma(X, \mathcal{O}_2)$-module. On a
$$
\text{Diff}^0(\mathcal{F}, \mathcal{G}) \subset
\text{Diff}^1(\mathcal{F}, \mathcal{G}) \subset
\text{Diff}^2(\mathcal{F}, \mathcal{G}) \subset \ldots
$$
La règle qui, à tout ouvert $U \subset X$, associe le module des opérateurs
différentiels $D : \mathcal{F}|_U \to \mathcal{G}|_U$ d'ordre $k$ définit
un faisceau de $\mathcal{O}_2$-modules sur $X$. On obtient ainsi un faisceau
d'opérateurs différentiels (si nous en avons besoin, nous ajouterons ici une
définition).

\begin{lemma}
\label{modules-lemma-composition-differential-operators}
Soit $X$ un espace topologique. Soit
$\mathcal{O}_1 \to \mathcal{O}_2$ un morphisme de faisceaux d'anneaux sur
$X$. Soient $\mathcal{E}, \mathcal{F}, \mathcal{G}$ des faisceaux de
$\mathcal{O}_2$-modules. Si $D : \mathcal{E} \to \mathcal{F}$ et
$D' : \mathcal{F} \to \mathcal{G}$ sont des opérateurs différentiels
d'ordres respectifs $k$ et $k'$, alors $D' \circ D$ est un opérateur
différentiel d'ordre $k + k'$.
\end{lemma}

\begin{proof}
Soit $g$ une section locale de $\mathcal{O}_2$. L'application qui envoie une
section locale $x$ de $\mathcal{E}$ sur
$$
D'(D(gx)) - gD'(D(x)) = D'(D(gx)) - D'(gD(x)) + D'(gD(x)) - gD'(D(x))
$$
est la somme de deux composées d'opérateurs différentiels d'ordre inférieur.
Le lemme en résulte par récurrence sur $k + k'$.
\end{proof}

\begin{lemma}
\label{modules-lemma-module-principal-parts}
Soit $X$ un espace topologique. Soit
$\mathcal{O}_1 \to \mathcal{O}_2$ un morphisme de faisceaux d'anneaux sur
$X$. Soit $\mathcal{F}$ un faisceau de $\mathcal{O}_2$-modules et soit
$k \geq 0$. Il existe un faisceau de $\mathcal{O}_2$-modules
$\mathcal{P}^k_{\mathcal{O}_2/\mathcal{O}_1}(\mathcal{F})$
et un isomorphisme canonique
$$
\text{Diff}^k_{\mathcal{O}_2/\mathcal{O}_1}(\mathcal{F}, \mathcal{G}) =
\Hom_{\mathcal{O}_2}(
\mathcal{P}^k_{\mathcal{O}_2/\mathcal{O}_1}(\mathcal{F}), \mathcal{G})
$$
fonctoriel par rapport au $\mathcal{O}_2$-module $\mathcal{G}$.
\end{lemma}

\begin{proof}
L'existence résulte d'arguments généraux de théorie des catégories
(insérer ici une référence ultérieure), mais nous donnerons aussi une
construction directe, qui sera utile dans les démonstrations à venir. Nous
emploierons librement les notations introduites dans la démonstration du
lemme \ref{modules-lemma-universal-module}. À tout opérateur différentiel
$D : \mathcal{F} \to \mathcal{G}$ est associé un morphisme
$\mathcal{O}_2$-linéaire
$L_D : \mathcal{O}_2[\mathcal{F}] \to \mathcal{G}$
envoyant $[m]$ sur $D(m)$. Si $D$ est d'ordre $0$, alors $L_D$ annule les
sections locales
$$
[m + m'] - [m] - [m'],\quad
g_0[m] - [g_0m]
$$
où $g_0$ est une section locale de $\mathcal{O}_2$ et $m, m'$ sont des
sections locales de $\mathcal{F}$. Si $D$ est d'ordre $1$, alors $L_D$
annule les sections locales
$$
[m + m'] - [m] - [m'],\quad
f[m] - [fm], \quad
g_0g_1[m] - g_0[g_1m] - g_1[g_0m] + [g_1g_0m]
$$
où $f$ est une section locale de $\mathcal{O}_1$, où $g_0, g_1$ sont des
sections locales de $\mathcal{O}_2$, et où $m, m'$ sont des sections locales
de $\mathcal{F}$. Si $D$ est d'ordre $k$, alors $L_D$ annule les sections
locales $[m + m'] - [m] - [m']$, $f[m] - [fm]$, ainsi que les sections locales
$$
g_0g_1\ldots g_k[m] - \sum g_0 \ldots \hat g_i \ldots g_k[g_im] + \ldots
+(-1)^{k + 1}[g_0\ldots g_km]
$$
Réciproquement, si $L : \mathcal{O}_2[\mathcal{F}] \to \mathcal{G}$ est un
morphisme $\mathcal{O}_2$-linéaire annulant toutes les sections locales
énumérées dans la phrase précédente, alors $m \mapsto L([m])$ est un
opérateur différentiel d'ordre $k$. Ainsi,
$\mathcal{P}^k_{\mathcal{O}_2/\mathcal{O}_1}(\mathcal{F})$
est le quotient de $\mathcal{O}_2[\mathcal{F}]$ par le
$\mathcal{O}_2$-sous-module engendré par ces sections locales.
\end{proof}

\begin{definition}
\label{modules-definition-module-principal-parts}
Soit $X$ un espace topologique. Soit
$\mathcal{O}_1 \to \mathcal{O}_2$ un morphisme de faisceaux d'anneaux sur
$X$. Soit $\mathcal{F}$ un faisceau de $\mathcal{O}_2$-modules. Le module
$\mathcal{P}^k_{\mathcal{O}_2/\mathcal{O}_1}(\mathcal{F})$ construit dans le
lemme \ref{modules-lemma-module-principal-parts} est appelé le
{\it module des parties principales d'ordre $k$} de $\mathcal{F}$.
\end{definition}

\noindent
Remarquons que les inclusions
$$
\text{Diff}^0(\mathcal{F}, \mathcal{G}) \subset
\text{Diff}^1(\mathcal{F}, \mathcal{G}) \subset
\text{Diff}^2(\mathcal{F}, \mathcal{G}) \subset \ldots
$$
correspondent, par le lemme de Yoneda (Catégories, lemme
\ref{categories-lemma-yoneda}), aux surjections
$$
\ldots \to \mathcal{P}^2_{\mathcal{O}_2/\mathcal{O}_1}(\mathcal{F})
\to \mathcal{P}^1_{\mathcal{O}_2/\mathcal{O}_1}(\mathcal{F})
\to \mathcal{P}^0_{\mathcal{O}_2/\mathcal{O}_1}(\mathcal{F}) = \mathcal{F}
$$

\begin{lemma}
\label{modules-lemma-differential-operators-sheafify}
Soit $X$ un espace topologique. Soit
$\mathcal{O}_1 \to \mathcal{O}_2$ un homomorphisme de préfaisceaux d'anneaux
sur $X$. Soit $\mathcal{F}$ un préfaisceau de $\mathcal{O}_2$-modules. Alors
$\mathcal{P}^k_{\mathcal{O}_2^\#/\mathcal{O}_1^\#}(\mathcal{F}^\#)$
est le faisceau associé au préfaisceau
$U \mapsto P^k_{\mathcal{O}_2(U)/\mathcal{O}_1(U)}(\mathcal{F}(U))$.
\end{lemma}

\begin{proof}
On peut le démontrer exactement comme pour le faisceau des différentielles
dans le lemme \ref{modules-lemma-differentials-sheafify}. Une méthode peut-être plus
satisfaisante consiste à utiliser directement la propriété universelle du
lemme \ref{modules-lemma-module-principal-parts}. Nous omettons les détails.
\end{proof}

\begin{lemma}
\label{modules-lemma-sequence-of-principal-parts}
Soit $X$ un espace topologique. Soit
$\mathcal{O}_1 \to \mathcal{O}_2$ un homomorphisme de faisceaux d'anneaux
sur $X$. Soit $\mathcal{F}$ un faisceau de $\mathcal{O}_2$-modules. Il
existe une suite exacte courte canonique
$$
0 \to
\Omega_{\mathcal{O}_2/\mathcal{O}_1} \otimes_{\mathcal{O}_2} \mathcal{F} \to
\mathcal{P}^1_{\mathcal{O}_2/\mathcal{O}_1}(\mathcal{F}) \to
\mathcal{F} \to 0
$$
fonctorielle en $\mathcal{F}$, appelée la
{\it suite des parties principales}.
\end{lemma}

\begin{proof}
Cela résulte de la version en algèbre commutative (Algèbre, lemme
\ref{algebra-lemma-sequence-of-principal-parts}) et des lemmes
\ref{modules-lemma-differentials-sheafify} et
\ref{modules-lemma-differential-operators-sheafify}.
\end{proof}

\begin{remark}
\label{modules-remark-functoriality-principal-parts}
Soit $X$ un espace topologique. Supposons donné un diagramme commutatif de
faisceaux d'anneaux
$$
\xymatrix{
\mathcal{B} \ar[r] & \mathcal{B}' \\
\mathcal{A} \ar[u] \ar[r] & \mathcal{A}' \ar[u]
}
$$
sur $X$, ainsi qu'un $\mathcal{B}$-module $\mathcal{F}$, un
$\mathcal{B}'$-module $\mathcal{F}'$ et un morphisme $\mathcal{B}$-linéaire
$\mathcal{F} \to \mathcal{F}'$. On obtient alors un système compatible de
morphismes de modules
$$
\xymatrix{
\ldots \ar[r] &
\mathcal{P}^2_{\mathcal{B}'/\mathcal{A}'}(\mathcal{F}') \ar[r] &
\mathcal{P}^1_{\mathcal{B}'/\mathcal{A}'}(\mathcal{F}') \ar[r] &
\mathcal{P}^0_{\mathcal{B}'/\mathcal{A}'}(\mathcal{F}') \\
\ldots \ar[r] &
\mathcal{P}^2_{\mathcal{B}/\mathcal{A}}(\mathcal{F}) \ar[r] \ar[u] &
\mathcal{P}^1_{\mathcal{B}/\mathcal{A}}(\mathcal{F}) \ar[r] \ar[u] &
\mathcal{P}^0_{\mathcal{B}/\mathcal{A}}(\mathcal{F}) \ar[u]
}
$$
Ces morphismes sont compatibles à toute composition ultérieure de
morphismes de ce type. Le moyen le plus simple de le voir consiste à utiliser
la description des modules
$\mathcal{P}^k_{\mathcal{B}/\mathcal{A}}(\mathcal{M})$ par générateurs
(locaux) et relations donnée dans la démonstration du lemme
\ref{modules-lemma-module-principal-parts}; on peut aussi le déduire directement de
leur propriété universelle. De plus, ces morphismes sont compatibles aux
suites exactes courtes du lemme \ref{modules-lemma-sequence-of-principal-parts}.
\end{remark}

\noindent
Étendons maintenant notre définition aux morphismes d'espaces annelés.

\begin{definition}
\label{modules-definition-relative-differential-operators}
Soit $(f, f^\sharp) : (X, \mathcal{O}_X) \to (S, \mathcal{O}_S)$
un morphisme d'espaces annelés. Soient $\mathcal{F}$ et $\mathcal{G}$ des
$\mathcal{O}_X$-modules et soit $k \geq 0$ un entier. Un
{\it opérateur différentiel d'ordre $k$ sur $X/S$} est un opérateur
différentiel $D : \mathcal{F} \to \mathcal{G}$ relativement à
$f^\sharp : f^{-1}\mathcal{O}_S \to \mathcal{O}_X$. On note
$\text{Diff}^k_{X/S}(\mathcal{F}, \mathcal{G})$ l'ensemble de ces opérateurs
différentiels.
\end{definition}












\section{Le complexe de de Rham}
\label{modules-section-de-rham-complex}

\noindent
Cette section est l'analogue d'Algèbre, section
\ref{algebra-section-de-rham-complex}, pour les morphismes d'espaces annelés.
Nous recommandons vivement au lecteur de commencer par cette section.

\medskip\noindent
Soit $X$ un espace topologique. Soit $\mathcal{A} \to \mathcal{B}$ un
homomorphisme de faisceaux d'anneaux. Notons
$\text{d} : \mathcal{B} \to \Omega_{\mathcal{B}/\mathcal{A}}$
le module des différentielles muni de sa $\mathcal{A}$-dérivation universelle,
construit dans la section \ref{modules-section-differentials}. Posons
$$
\Omega_{\mathcal{B}/\mathcal{A}}^i =
\wedge^i_\mathcal{B}(\Omega_{\mathcal{B}/\mathcal{A}})
$$
pour $i \geq 0$; c'est la $i$-ième puissance extérieure définie dans la
section \ref{modules-section-symmetric-exterior}.

\begin{definition}
\label{modules-definition-de-rham-complex}
Dans la situation ci-dessus, le
{\it complexe de de Rham de $\mathcal{B}$ sur $\mathcal{A}$} est l'unique
complexe
$$
\Omega_{\mathcal{B}/\mathcal{A}}^0 \to
\Omega_{\mathcal{B}/\mathcal{A}}^1 \to
\Omega_{\mathcal{B}/\mathcal{A}}^2 \to \ldots
$$
de faisceaux de $\mathcal{A}$-modules dont la différentielle en degré $0$
est $\text{d} : \mathcal{B} \to \Omega_{\mathcal{B}/\mathcal{A}}$ et dont
les différentielles en degrés supérieurs satisfont la propriété suivante :
\begin{equation}
\label{modules-equation-rule}
\text{d}\left(b_0\text{d}b_1 \wedge \ldots \wedge \text{d}b_p\right) =
\text{d}b_0 \wedge \text{d}b_1 \wedge \ldots \wedge \text{d}b_p
\end{equation}
où $b_0, \ldots, b_p \in \mathcal{B}(U)$ sont des sections définies sur un même
ouvert $U \subset X$.
\end{definition}

\noindent
On pourrait construire ce complexe en reprenant les arguments assez lourds
d'Algèbre, section \ref{algebra-section-de-rham-complex}. Rappelons plutôt
que $\Omega_{\mathcal{B}/\mathcal{A}}$ est le faisceau associé au préfaisceau
$U \mapsto \Omega_{\mathcal{B}(U)/\mathcal{A}(U)}$; voir le lemme
\ref{modules-lemma-differentials-sheafify}. Ainsi,
$\Omega_{\mathcal{B}/\mathcal{A}}^i$ est le faisceau associé au préfaisceau
$U \mapsto \Omega^i_{\mathcal{B}(U)/\mathcal{A}(U)}$; voir le lemme
\ref{modules-lemma-local-tensor-algebra}. On peut donc définir le complexe de de Rham
en faisceautisant la règle
$$
U \longmapsto
\Omega^\bullet_{\mathcal{B}(U)/\mathcal{A}(U)}
$$

\begin{lemma}
\label{modules-lemma-pullback-de-rham-complex}
Soit $f : Y \to X$ une application continue d'espaces topologiques. Soit
$\mathcal{A} \to \mathcal{B}$ un homomorphisme de faisceaux d'anneaux sur
$X$. Il existe alors une identification canonique
$f^{-1}\Omega^\bullet_{\mathcal{B}/\mathcal{A}} =
\Omega^\bullet_{f^{-1}\mathcal{B}/f^{-1}\mathcal{A}}$
de complexes de de Rham.
\end{lemma}

\begin{proof}
Démonstration omise. Indication : comparer avec le lemme
\ref{modules-lemma-pullback-differentials}.
\end{proof}

\begin{lemma}
\label{modules-lemma-differentials-de-rham-complex-order-1}
Soit $X$ un espace topologique. Soit $\mathcal{A} \to \mathcal{B}$ un
homomorphisme de faisceaux d'anneaux sur $X$. Les différentielles
$\text{d} : \Omega^i_{\mathcal{B}/\mathcal{A}}  \to
\Omega^{i + 1}_{\mathcal{B}/\mathcal{A}}$
sont des opérateurs différentiels d'ordre $1$.
\end{lemma}

\begin{proof}
Compte tenu de la construction ci-dessus du complexe de de Rham comme
faisceau associé à la règle
$U \mapsto \Omega^\bullet_{\mathcal{B}(U)/\mathcal{A}(U)}$, cela résulte
d'Algèbre, lemme
\ref{algebra-lemma-differentials-de-rham-complex-order-1}.
\end{proof}

\noindent
Soit $X$ un espace topologique. Soit
$$
\xymatrix{
\mathcal{B} \ar[r] & \mathcal{B}' \\
\mathcal{A} \ar[r] \ar[u] & \mathcal{A}' \ar[u]
}
$$
un diagramme commutatif de faisceaux d'anneaux sur $X$. Il existe un
morphisme naturel de complexes de de Rham
$$
\Omega^\bullet_{\mathcal{B}/\mathcal{A}} \longrightarrow
\Omega^\bullet_{\mathcal{B}'/\mathcal{A}'}
$$
En effet, en degré $0$, c'est le morphisme
$\mathcal{B} \to \mathcal{B}'$; en degré $1$, c'est le morphisme
$\Omega_{\mathcal{B}/\mathcal{A}} \to \Omega_{\mathcal{B}'/\mathcal{A}'}$
construit dans la section \ref{modules-section-differentials}; pour $p \geq 2$, c'est
le morphisme induit
$\Omega^p_{\mathcal{B}/\mathcal{A}} =
\wedge^p_\mathcal{B}(\Omega_{\mathcal{B}/\mathcal{A}}) \to
\wedge^p_{\mathcal{B}'}(\Omega_{\mathcal{B}'/\mathcal{A}'}) =
\Omega^p_{\mathcal{B}'/\mathcal{A}'}$.
La compatibilité aux différentielles résulte de leur caractérisation par la
formule (\ref{modules-equation-rule}).

\begin{definition}
\label{modules-definition-de-rham-complex-morphism-ringed-spaces}
Soit $f : (X, \mathcal{O}_X) \to (Y, \mathcal{O}_Y)$ un morphisme
d'espaces annelés. Le {\it complexe de de Rham} de $f$, ou de $X$ sur $Y$,
est le complexe
$$
\Omega^\bullet_{X/Y} = \Omega^\bullet_{\mathcal{O}_X/f^{-1}\mathcal{O}_Y}
$$
\end{definition}

\noindent
Considérons un diagramme commutatif d'espaces annelés
$$
\xymatrix{
X' \ar[d]_{h'} \ar[r]_f & X \ar[d]^h \\
S' \ar[r]^g & S
}
$$
On obtient alors un morphisme canonique
$$
\Omega^\bullet_{X/S} \to f_*\Omega^\bullet_{X'/S'}
$$
de complexes de de Rham. En effet, le diagramme commutatif de faisceaux
d'anneaux
$$
\xymatrix{
f^{-1}\mathcal{O}_X \ar[r] & \mathcal{O}_{X'} \\
f^{-1}h^{-1}\mathcal{O}_S \ar[u] \ar[r] & (h')^{-1}\mathcal{O}_{S'} \ar[u]
}
$$
sur $X'$ fournit un morphisme de complexes (voir ci-dessus)
$$
f^{-1}\Omega^\bullet_{X/S} =
\Omega^\bullet_{f^{-1}\mathcal{O}_X/f^{-1}h^{-1}\mathcal{O}_S}
\longrightarrow
\Omega^\bullet_{\mathcal{O}_{X'}/(h')^{-1}\mathcal{O}_{S'}} =
\Omega^\bullet_{X'/S'}
$$
(en utilisant le lemme \ref{modules-lemma-pullback-de-rham-complex} pour la première
égalité), puis on applique l'adjonction.

\begin{lemma}
\label{modules-lemma-differentials-relative-de-rham-complex-order-1}
Soit $f : X \to Y$ un morphisme d'espaces annelés. Les différentielles
$\text{d} : \Omega^i_{X/Y}  \to \Omega^{i + 1}_{X/Y}$
sont des opérateurs différentiels d'ordre $1$ sur $X/Y$.
\end{lemma}

\begin{proof}
Cela résulte immédiatement du lemme
\ref{modules-lemma-differentials-de-rham-complex-order-1} et de la définition.
\end{proof}





















\section{Le complexe cotangent naïf}
\label{modules-section-netherlander}

\noindent
Cette section est l'analogue d'Algèbre, section
\ref{algebra-section-netherlander}, pour les morphismes d'espaces annelés.
Nous recommandons vivement au lecteur de commencer par cette section.

\medskip\noindent
Soit $X$ un espace topologique. Soit $\mathcal{A} \to \mathcal{B}$ un
homomorphisme de faisceaux d'anneaux. Dans cette section, pour tout faisceau
d'ensembles $\mathcal{E}$ sur $X$, on note $\mathcal{A}[\mathcal{E}]$ le
faisceau associé au préfaisceau
$U \mapsto \mathcal{A}(U)[\mathcal{E}(U)]$. Ici,
$\mathcal{A}(U)[\mathcal{E}(U)]$
désigne l'algèbre de polynômes sur $\mathcal{A}(U)$ dont les variables
correspondent aux éléments de $\mathcal{E}(U)$. On note
$[e] \in \mathcal{A}(U)[\mathcal{E}(U)]$ la variable correspondant à
$e \in \mathcal{E}(U)$. Il existe une surjection canonique de
$\mathcal{A}$-algèbres
\begin{equation}
\label{modules-equation-canonical-presentation}
\mathcal{A}[\mathcal{B}] \longrightarrow \mathcal{B},\quad [b] \longmapsto b
\end{equation}
dont on note le noyau $\mathcal{I} \subset \mathcal{A}[\mathcal{B}]$.
On observe immédiatement que $\mathcal{I}$ est engendré par les sections
locales $[b][b'] - [bb']$ et $[a] - a$. D'après le lemme
\ref{modules-lemma-differential-seq}, il existe un morphisme canonique
\begin{equation}
\label{modules-equation-naive-cotangent-complex}
\mathcal{I}/\mathcal{I}^2
\longrightarrow
\Omega_{\mathcal{A}[\mathcal{B}]/\mathcal{A}}
\otimes_{\mathcal{A}[\mathcal{B}]} \mathcal{B}
\end{equation}
dont le conoyau est canoniquement isomorphe à
$\Omega_{\mathcal{B}/\mathcal{A}}$.

\begin{definition}
\label{modules-definition-naive-cotangent-complex}
Soit $X$ un espace topologique. Soit $\mathcal{A} \to \mathcal{B}$ un
homomorphisme de faisceaux d'anneaux. Le
{\it complexe cotangent naïf} $\NL_{\mathcal{B}/\mathcal{A}}$ est le
complexe de chaînes
(\ref{modules-equation-naive-cotangent-complex})
$$
\NL_{\mathcal{B}/\mathcal{A}} =
\left(\mathcal{I}/\mathcal{I}^2
\longrightarrow
\Omega_{\mathcal{A}[\mathcal{B}]/\mathcal{A}}
\otimes_{\mathcal{A}[\mathcal{B}]} \mathcal{B}\right)
$$
où $\mathcal{I}/\mathcal{I}^2$ est placé en degré $-1$ et où
$\Omega_{\mathcal{A}[\mathcal{B}]/\mathcal{A}}
\otimes_{\mathcal{A}[\mathcal{B}]} \mathcal{B}$
est placé en degré $0$.
\end{definition}

\noindent
Cette construction possède une fonctorialité analogue à celle étudiée dans
le lemme \ref{modules-lemma-functoriality-differentials} pour les modules de
différentielles. Plus précisément, étant donné un diagramme commutatif
\begin{equation}
\label{modules-equation-commutative-square-sheaves}
\vcenter{
\xymatrix{
\mathcal{B} \ar[r] & \mathcal{B}' \\
\mathcal{A} \ar[u] \ar[r] & \mathcal{A}' \ar[u]
}
}
\end{equation}
de faisceaux d'anneaux sur $X$, il existe un morphisme canonique
$\mathcal{B}$-linéaire de complexes
$$
\NL_{\mathcal{B}/\mathcal{A}} \longrightarrow \NL_{\mathcal{B}'/\mathcal{A}'}
$$
En effet, les morphismes du diagramme commutatif donnent un morphisme
canonique $\mathcal{A}[\mathcal{B}] \to \mathcal{A}'[\mathcal{B}']$ qui
envoie $\mathcal{I}$ dans
$\mathcal{I}' = \Ker(\mathcal{A}'[\mathcal{B}'] \to \mathcal{B}')$.
On obtient donc un morphisme
$\mathcal{I}/\mathcal{I}^2 \to \mathcal{I}'/(\mathcal{I}')^2$ et un
morphisme entre les modules de différentielles; ensemble, ils donnent le
morphisme voulu entre les complexes cotangents naïfs. Ce morphisme est
compatible aux compositions au sens suivant : étant donné un diagramme
commutatif
$$
\xymatrix{
\mathcal{B} \ar[r] &
\mathcal{B}' \ar[r] &
\mathcal{B}'' \\
\mathcal{A} \ar[u] \ar[r] &
\mathcal{A}' \ar[u] \ar[r] &
\mathcal{A}'' \ar[u]
}
$$
de faisceaux d'anneaux, la composée
$$
\NL_{\mathcal{B}/\mathcal{A}} \longrightarrow
\NL_{\mathcal{B}'/\mathcal{A}'} \longrightarrow
\NL_{\mathcal{B}''/\mathcal{A}''}
$$
est le morphisme associé au rectangle extérieur.

\medskip\noindent
On peut choisir une autre présentation de $\mathcal{B}$ comme quotient d'une
algèbre de polynômes sur $\mathcal{A}$ et obtenir néanmoins le même objet de
$D(\mathcal{B})$. Pour l'expliquer, soit $\mathcal{E}$ un faisceau
d'ensembles sur $X$ et soit $\alpha : \mathcal{E} \to \mathcal{B}$ un
morphisme de faisceaux d'ensembles. On obtient alors un homomorphisme de
$\mathcal{A}$-algèbres $\mathcal{A}[\mathcal{E}] \to \mathcal{B}$. Si ce
morphisme est surjectif, c'est-à-dire si $\alpha(\mathcal{E})$ engendre
$\mathcal{B}$ comme $\mathcal{A}$-algèbre, on pose
$$
\NL(\alpha) = \left(
\mathcal{J}/\mathcal{J}^2
\longrightarrow
\Omega_{\mathcal{A}[\mathcal{E}]/\mathcal{A}}
\otimes_{\mathcal{A}[\mathcal{E}]} \mathcal{B}\right)
$$
où $\mathcal{J} \subset \mathcal{A}[\mathcal{E}]$ est le noyau de la
surjection $\mathcal{A}[\mathcal{E}] \to \mathcal{B}$. Voici le résultat.

\begin{lemma}
\label{modules-lemma-NL-up-to-qis}
Dans la situation ci-dessus, il existe un isomorphisme canonique
$\NL(\alpha) = \NL_{\mathcal{B}/\mathcal{A}}$ in $D(\mathcal{B})$.
\end{lemma}

\begin{proof}
Remarquons que
$\NL_{\mathcal{B}/\mathcal{A}} = \NL(\text{id}_\mathcal{B})$. Il suffit
donc de montrer que, pour deux morphismes
$\alpha_i : \mathcal{E}_i \to \mathcal{B}$ comme ci-dessus, il existe un
quasi-isomorphisme canonique
$\NL(\alpha_1) = \NL(\alpha_2)$ dans $D(\mathcal{B})$. Pour cela, posons
$\mathcal{E} = \mathcal{E}_1 \amalg \mathcal{E}_2$ et
$\alpha = \alpha_1 \amalg \alpha_2 : \mathcal{E} \to \mathcal{B}$.
Posons
$\mathcal{J}_i = \Ker(\mathcal{A}[\mathcal{E}_i] \to \mathcal{B})$
et
$\mathcal{J} = \Ker(\mathcal{A}[\mathcal{E}] \to \mathcal{B})$.
On obtient des morphismes
$\mathcal{A}[\mathcal{E}_i] \to \mathcal{A}[\mathcal{E}]$ qui envoient
$\mathcal{J}_i$ dans $\mathcal{J}$. On obtient donc des morphismes
canoniques de complexes
$$
\NL(\alpha_i) \longrightarrow \NL(\alpha)
$$
et il suffit de montrer que ce sont des quasi-isomorphismes. Il suffit pour
cela de le vérifier sur les fibres (lemme \ref{modules-lemma-abelian}). Si $x \in X$,
la fibre de $\NL(\alpha)$ est le complexe $\NL(\alpha_x)$ d'Algèbre,
section \ref{algebra-section-netherlander}, associé à la présentation
$\mathcal{A}_x[\mathcal{E}_x] \to \mathcal{B}_x$ provenant du morphisme
$\alpha_x : \mathcal{E}_x \to \mathcal{B}_x$. (Certains détails sont omis;
utiliser le lemme \ref{modules-lemma-stalk-module-differentials} pour vérifier la
compatibilité entre la formation des différentielles et le passage aux
fibres.) Le résultat découle d'Algèbre, lemme
\ref{algebra-lemma-NL-homotopy}.
\end{proof}

\begin{lemma}
\label{modules-lemma-pullback-NL}
Soit $f : X \to Y$ une application continue d'espaces topologiques. Soit
$\mathcal{A} \to \mathcal{B}$ un homomorphisme de faisceaux d'anneaux sur
$Y$. Alors $f^{-1}\NL_{\mathcal{B}/\mathcal{A}} =
\NL_{f^{-1}\mathcal{B}/f^{-1}\mathcal{A}}$.
\end{lemma}

\begin{proof}
Démonstration omise. Indication : utiliser le lemme
\ref{modules-lemma-pullback-differentials}.
\end{proof}

\begin{lemma}
\label{modules-lemma-stalk-NL}
Soit $X$ un espace topologique. Soit $\mathcal{A} \to \mathcal{B}$ un
homomorphisme de faisceaux d'anneaux sur $X$. Pour $x \in X$, on a
$\NL_{\mathcal{B}/\mathcal{A}, x} =
\NL_{\mathcal{B}_x/\mathcal{A}_x}$.
\end{lemma}

\begin{proof}
C'est le cas particulier du lemme \ref{modules-lemma-pullback-NL} correspondant à
l'inclusion $\{x\} \to X$.
\end{proof}

\begin{lemma}
\label{modules-lemma-exact-sequence-NL}
Soit $X$ un espace topologique. Soient
$\mathcal{A} \to \mathcal{B} \to \mathcal{C}$ des morphismes de faisceaux
d'anneaux. Soit $C$ le cône (Catégories dérivées, définition
\ref{derived-definition-cone}) du morphisme de complexes
$\NL_{\mathcal{C}/\mathcal{A}} \to \NL_{\mathcal{C}/\mathcal{B}}$.
Il existe un morphisme canonique
$$
c :
\NL_{\mathcal{B}/\mathcal{A}} \otimes_\mathcal{B} \mathcal{C}
\longrightarrow
C[-1]
$$
de complexes de $\mathcal{C}$-modules qui fournit une suite exacte canonique
à six termes
$$
\xymatrix{
H^0(\NL_{\mathcal{B}/\mathcal{A}} \otimes_\mathcal{B} \mathcal{C}) \ar[r] &
H^0(\NL_{\mathcal{C}/\mathcal{A}}) \ar[r] &
H^0(\NL_{\mathcal{C}/\mathcal{B}}) \ar[r] &
0 \\
H^{-1}(\NL_{\mathcal{B}/\mathcal{A}} \otimes_\mathcal{B} \mathcal{C}) \ar[r] &
H^{-1}(\NL_{\mathcal{C}/\mathcal{A}}) \ar[r] &
H^{-1}(\NL_{\mathcal{C}/\mathcal{B}}) \ar[llu]
}
$$
de faisceaux de cohomologie.
\end{lemma}

\begin{proof}
Pour définir le morphisme $c$, il faut définir un morphisme
$c_1 : \NL_{\mathcal{B}/\mathcal{A}} \otimes_\mathcal{B} \mathcal{C}
\to \NL_{\mathcal{C}/\mathcal{A}}$ ainsi qu'une homotopie explicite entre
la composée
$$
\NL_{\mathcal{B}/\mathcal{A}} \otimes_\mathcal{B} \mathcal{C} \to
\NL_{\mathcal{C}/\mathcal{A}} \to \NL_{\mathcal{C}/\mathcal{B}}
$$
et le morphisme nul; voir Catégories dérivées, lemme
\ref{derived-lemma-map-from-cone}. Pour $c_1$, on utilise la fonctorialité
décrite ci-dessus appliquée au diagramme évident. Pour l'homotopie, on
utilise le morphisme
$$
\NL_{\mathcal{B}/\mathcal{A}}^0 \otimes_\mathcal{B} \mathcal{C}
\longrightarrow
\NL_{\mathcal{C}/\mathcal{B}}^{-1},\quad
\text{d}[b] \otimes 1 \longmapsto [\varphi(b)] - b[1]
$$
où $\varphi : \mathcal{B} \to \mathcal{C}$ est le morphisme donné. Comparer
avec Algèbre, remarque
\ref{algebra-remark-composition-homotopy-equivalent-to-zero}. Pour en déduire
l'assertion sur les faisceaux de cohomologie, il suffit de montrer que
$H^0(c)$ est un isomorphisme et que $H^{-1}(c)$ est surjectif. On peut le
vérifier sur les fibres, en utilisant le lemme \ref{modules-lemma-stalk-NL}, puis
appliquer le résultat correspondant d'algèbre commutative; voir Algèbre,
lemme \ref{algebra-lemma-exact-sequence-NL}. Certains détails sont omis.
\end{proof}

\noindent
Le complexe cotangent naïf d'un morphisme d'espaces annelés se définit au
moyen du complexe construit ci-dessus.

\begin{definition}
\label{modules-definition-cotangent-complex-morphism-ringed-topoi}
Le {\it complexe cotangent naïf} $\NL_f = \NL_{X/Y}$ d'un morphisme d'espaces
annelés $f : (X, \mathcal{O}_X) \to (Y, \mathcal{O}_Y)$ est
$\NL_{\mathcal{O}_X/f^{-1}\mathcal{O}_Y}$.
\end{definition}

\noindent
Étant donné un diagramme commutatif
$$
\xymatrix{
X' \ar[r]_g \ar[d]_{f'} & X \ar[d]^f \\
Y' \ar[r]^h & Y
}
$$
d'espaces annelés, il existe un morphisme canonique
$c : g^*\NL_{X/Y} \to \NL_{X'/Y'}$. C'est le morphisme
$$
g^*\NL_{X/Y} =
\mathcal{O}_{X'} \otimes_{g^{-1}\mathcal{O}_X}
\NL_{g^{-1}\mathcal{O}_X/g^{-1}f^{-1}\mathcal{O}_Y}
\longrightarrow
\NL_{\mathcal{O}_{X'}/(f')^{-1}\mathcal{O}_{Y'}} = \NL_{X'/Y'}
$$
où la flèche provient du diagramme commutatif de faisceaux d'anneaux
$$
\xymatrix{
g^{-1}\mathcal{O}_X \ar[r]_{g^\sharp} &
\mathcal{O}_{X'} \\
g^{-1}f^{-1}\mathcal{O}_Y
\ar[r]^{g^{-1}h^\sharp}
\ar[u]^{g^{-1}f^\sharp} &
(f')^{-1}\mathcal{O}_{Y'} \ar[u]_{(f')^\sharp}
}
$$
comme dans (\ref{modules-equation-commutative-square-sheaves}) ci-dessus. Étant
donné un second diagramme de ce type
$$
\xymatrix{
X'' \ar[r]_{g'} \ar[d] & X' \ar[d] \\
Y'' \ar[r] & Y'
}
$$
la composée de $(g')^*c$ et du morphisme
$c' : (g')^*\NL_{X'/Y'} \to \NL_{X''/Y''}$ est le morphisme
$(g \circ g')^*\NL_{X/Y} \to \NL_{X''/Y''}$.

\begin{lemma}
\label{modules-lemma-exact-sequence-NL-ringed-topoi}
Soient $f : X \to Y$ et $g : Y \to Z$ des morphismes d'espaces annelés.
Soit $C$ le cône du morphisme $\NL_{X/Z} \to \NL_{X/Y}$ de complexes de
$\mathcal{O}_X$-modules. Il existe un morphisme canonique
$$
f^*\NL_{Y/Z} \to C[-1]
$$
qui fournit une suite exacte canonique à six termes
$$
\xymatrix{
H^0(f^*\NL_{Y/Z}) \ar[r] &
H^0(\NL_{X/Z}) \ar[r] &
H^0(\NL_{X/Y}) \ar[r] &
0 \\
H^{-1}(f^*\NL_{Y/Z}) \ar[r] &
H^{-1}(\NL_{X/Z}) \ar[r] &
H^{-1}(\NL_{X/Y}) \ar[llu]
}
$$
de faisceaux de cohomologie.
\end{lemma}

\begin{proof}
Considérons les morphismes de faisceaux d'anneaux
$$
(g \circ f)^{-1}\mathcal{O}_Z \to f^{-1}\mathcal{O}_Y \to \mathcal{O}_X
$$
et appliquons le lemme \ref{modules-lemma-exact-sequence-NL}.
\end{proof}











% OK, start here.
%

\chapter{Modules sur les sites}



\phantomsection
\label{sites-modules-section-phantom}


\section{Introduction}
\label{sites-modules-section-introduction}

\noindent
Dans ce chapitre, nous développons les notions de base relatives aux faisceaux
de modules sur les topos annelés ou les sites annelés. Nous établissons d'abord
quelques faits fondamentaux sur les faisceaux abéliens, puis nous introduisons
les sites annelés et les topos annelés. Nous étudions ensuite certaines notions
très élémentaires concernant les (pré)faisceaux de $\mathcal{O}$-modules,
parallèlement au traitement des (pré)faisceaux de $\mathcal{O}$-modules dans le
chapitre consacré aux faisceaux sur les espaces. Cela fait, nous reprenons une
grande partie de la discussion du chapitre sur les faisceaux de modules (voir
Faisceaux de modules, section \ref{modules-section-introduction}). Les
références de base sont \cite{FAC}, \cite{EGA} et \cite{SGA4}.






\section{Préfaisceaux abéliens}
\label{sites-modules-section-abelian-pre-sheaves}

\noindent
Soit $\mathcal{C}$ une catégorie. Les préfaisceaux abéliens ont été introduits
dans Sites, sections \ref{sites-section-presheaves} et
\ref{sites-section-sheaves}, puis étudiés plus avant dans Sites, section
\ref{sites-section-sheaves-algebraic-structures}. Nous suivrons la convention
de cette dernière référence : un préfaisceau abélien est un préfaisceau
d'ensembles muni, sur chacun de ses ensembles de sections, d'une loi d'addition
compatible avec les applications de restriction. Rappelons que la catégorie
des préfaisceaux abéliens sur $\mathcal{C}$ est notée
$\textit{PAb}(\mathcal{C})$.

\medskip\noindent
La catégorie $\textit{PAb}(\mathcal{C})$ est abélienne au sens d'Homologie,
Définition \ref{homology-definition-abelian-category}. Étant donnée une
application de préfaisceaux $\varphi : \mathcal{G}_1 \to \mathcal{G}_2$, le
noyau de $\varphi$ est le préfaisceau abélien
$U \mapsto \Ker(\mathcal{G}_1(U) \to \mathcal{G}_2(U))$, et le conoyau de
$\varphi$ est
le préfaisceau $U \mapsto
\Coker(\mathcal{G}_1(U) \to \mathcal{G}_2(U))$. Puisque la catégorie des
groupes abéliens est abélienne, il s'ensuit que $\Coim = \Im$, cette égalité
étant vraie au-dessus de chaque $U$. Une suite de préfaisceaux abéliens
$$
\mathcal{G}_1 \longrightarrow
\mathcal{G}_2 \longrightarrow
\mathcal{G}_3
$$
est exacte si et seulement si
$\mathcal{G}_1(U) \to \mathcal{G}_2(U) \to \mathcal{G}_3(U)$
est une suite exacte de groupes abéliens pour tout $U \in \Ob(\mathcal{C})$.
Nous laissons les vérifications au lecteur.

\begin{lemma}
\label{sites-modules-lemma-limits-colimits-abelian-presheaves}
Soit $\mathcal{C}$ une catégorie.
\begin{enumerate}
\item Toutes les limites et toutes les colimites existent dans
$\textit{PAb}(\mathcal{C})$.
\item Toutes les limites et toutes les colimites commutent à la prise des
sections au-dessus des objets de $\mathcal{C}$.
\end{enumerate}
\end{lemma}

\begin{proof}
Soit $\mathcal{I} \to \textit{PAb}(\mathcal{C})$, $i \mapsto \mathcal{F}_i$,
un diagramme. Nous pouvons simplement définir des préfaisceaux abéliens $L$ et
$C$ par les règles
$$
L : U \longmapsto \lim_i \mathcal{F}_i(U)
$$
et
$$
C : U \longmapsto \colim_i \mathcal{F}_i(U).
$$
Il existe manifestement des applications de préfaisceaux abéliens
$L \to \mathcal{F}_i$ et $\mathcal{F}_i \to C$, obtenues à partir des
applications correspondantes sur les groupes de sections au-dessus de chaque
$U$. On vérifie immédiatement que $L$ et $C$, munis de ces applications, sont
respectivement la limite et la colimite du diagramme dans
$\textit{PAb}(\mathcal{C})$. Cela démontre (1) et (2). Nous omettons les
détails.
\end{proof}


\section{Faisceaux abéliens}
\label{sites-modules-section-abelian-sheaves}

\noindent
Soit $\mathcal{C}$ un site. La catégorie des faisceaux abéliens sur
$\mathcal{C}$ est notée $\textit{Ab}(\mathcal{C})$. C'est la sous-catégorie
pleine de $\textit{PAb}(\mathcal{C})$ formée des préfaisceaux abéliens dont le
préfaisceau d'ensembles sous-jacent est un faisceau. Les propriétés
($\alpha$) -- ($\zeta$) de Sites, section
\ref{sites-section-sheaves-algebraic-structures}, sont satisfaites ; voir
Sites, Proposition
\ref{sites-proposition-functoriality-algebraic-structures-topoi}. En
particulier, le foncteur d'inclusion
$\textit{Ab}(\mathcal{C}) \to \textit{PAb}(\mathcal{C})$ admet un adjoint à
gauche, à savoir le foncteur de faisceautisation
$\mathcal{G} \mapsto \mathcal{G}^\#$.

\medskip\noindent
Nous suggérons au lecteur de démontrer le lemme sur une feuille de brouillon
plutôt que d'en lire la démonstration.

\begin{lemma}
\label{sites-modules-lemma-abelian-abelian}
Soit $\mathcal{C}$ un site. Soit $\varphi : \mathcal{F} \to \mathcal{G}$
un morphisme de faisceaux abéliens sur $\mathcal{C}$.
\begin{enumerate}
\item La catégorie $\textit{Ab}(\mathcal{C})$ est une catégorie abélienne.
\item Le noyau $\Ker(\varphi)$ de $\varphi$ coïncide avec le noyau de
$\varphi$ comme morphisme de préfaisceaux.
\item Le morphisme $\varphi$ est injectif (Homologie, Définition
\ref{homology-definition-injective-surjective}) si et seulement si $\varphi$ est
injectif comme application de préfaisceaux (Sites, Définition
\ref{sites-definition-presheaves-injective-surjective}), si et seulement si
$\varphi$ est injectif comme application de faisceaux (Sites, Définition
\ref{sites-definition-sheaves-injective-surjective}).
\item Le conoyau $\Coker(\varphi)$ de $\varphi$ est le faisceau associé au
conoyau de $\varphi$ comme morphisme de préfaisceaux.
\item Le morphisme $\varphi$ est surjectif (Homologie, Définition
\ref{homology-definition-injective-surjective}) si et seulement si $\varphi$ est
surjectif comme application de faisceaux (Sites, Définition
\ref{sites-definition-sheaves-injective-surjective}).
\item Un complexe de faisceaux abéliens
$$
\mathcal{F} \to \mathcal{G} \to \mathcal{H}
$$
est exact en $\mathcal{G}$ si et seulement si, pour tout
$U \in \Ob(\mathcal{C})$ et tout $s \in \mathcal{G}(U)$ dont l'image dans
$\mathcal{H}(U)$ est nulle, il existe un recouvrement
$\{U_i \to U\}_{i \in I}$ dans $\mathcal{C}$ tel que chaque $s|_{U_i}$
appartienne à l'image de $\mathcal{F}(U_i) \to \mathcal{G}(U_i)$.
\end{enumerate}
\end{lemma}

\begin{proof}
Nous affirmons qu'Homologie, Lemme
\ref{homology-lemma-adjoint-get-abelian}, s'applique aux catégories
$\mathcal{A} = \textit{Ab}(\mathcal{C})$ et
$\mathcal{B} = \textit{PAb}(\mathcal{C})$, ainsi qu'aux foncteurs
$a : \mathcal{A} \to \mathcal{B}$ (l'inclusion) et
$b : \mathcal{B} \to \mathcal{A}$ (la faisceautisation). Vérifions les
hypothèses d'Homologie, Lemme \ref{homology-lemma-adjoint-get-abelian}.
L'hypothèse (1) dit que $\mathcal{A}$ et $\mathcal{B}$ sont des
catégories additives, que $a$ et $b$ sont des foncteurs additifs, et que $a$
est adjoint à droite de $b$. Les deux premières assertions sont claires, et
l'adjonction est Sites, section
\ref{sites-section-sheaves-algebraic-structures} ($\epsilon$). L'hypothèse
(2) dit que $\textit{PAb}(\mathcal{C})$ est abélienne, ce que nous avons vu à
la section \ref{sites-modules-section-abelian-pre-sheaves}, et que la faisceautisation est
exacte à gauche, ce qui est Sites, section
\ref{sites-section-sheaves-algebraic-structures} ($\zeta$). Enfin, la dernière
hypothèse est $ba \cong \text{id}_\mathcal{A}$, ce qui est Sites, section
\ref{sites-section-sheaves-algebraic-structures} ($\delta$). Ainsi Homologie,
Lemme \ref{homology-lemma-adjoint-get-abelian}, s'applique et nous concluons
que $\textit{Ab}(\mathcal{C})$ est abélienne.

\medskip\noindent
La démonstration d'Homologie, Lemme
\ref{homology-lemma-adjoint-get-abelian}, montre que $\Ker(\varphi)$ et
$\Coker(\varphi)$ sont les faisceaux associés respectivement au noyau et au
conoyau de $\varphi$ comme morphisme de préfaisceaux abéliens. Cela démontre
(4). Puisque le noyau est un égalisateur (donc une limite) et que la
faisceautisation commute aux limites finies, nous en déduisons (2).

\medskip\noindent
L'assertion (2) implique (3). L'assertion (4) implique (5) d'après notre
description de la faisceautisation. La caractérisation de l'exactitude en (6)
résulte de (2), de (5), et du fait que la suite est exacte si et seulement si
$\Im(\mathcal{F} \to \mathcal{G}) =
\Ker(\mathcal{G} \to \mathcal{H})$.
\end{proof}

\noindent
On peut encore formuler la partie (6) du lemme en disant qu'une suite de
faisceaux abéliens
$$
\mathcal{F}_1 \longrightarrow
\mathcal{F}_2 \longrightarrow
\mathcal{F}_3
$$
est exacte si et seulement si le faisceau associé à
$U \mapsto \Im(\mathcal{F}_1(U) \to \mathcal{F}_2(U))$
est égal au noyau de $\mathcal{F}_2 \to \mathcal{F}_3$.

\begin{lemma}
\label{sites-modules-lemma-limits-colimits-abelian-sheaves}
Soit $\mathcal{C}$ un site.
\begin{enumerate}
\item Toutes les limites et toutes les colimites existent dans
$\textit{Ab}(\mathcal{C})$.
\item Les limites coïncident avec les limites correspondantes de préfaisceaux
abéliens sur $\mathcal{C}$ (autrement dit, elles commutent à la prise des
sections au-dessus des objets de $\mathcal{C}$).
\item Les sommes directes finies coïncident avec les sommes directes finies
correspondantes dans la catégorie des préfaisceaux abéliens sur $\mathcal{C}$.
\item Une colimite est le faisceau associé à la colimite correspondante dans la
catégorie des préfaisceaux abéliens.
\item Les colimites filtrantes sont exactes.
\end{enumerate}
\end{lemma}

\begin{proof}
D'après le Lemme \ref{sites-modules-lemma-limits-colimits-abelian-presheaves}, les limites et
les colimites de préfaisceaux abéliens existent et s'obtiennent en prenant les
limites et les colimites au niveau des sections au-dessus des objets.

\medskip\noindent
Soit $\mathcal{I} \to \textit{Ab}(\mathcal{C})$,
$i \mapsto \mathcal{F}_i$, un diagramme. Soit $\lim_i \mathcal{F}_i$ sa
limite comme préfaisceau abélien. D'après Sites, Lemme
\ref{sites-lemma-limit-sheaf}, c'est un faisceau abélien. On voit alors sans
difficulté que $\lim_i \mathcal{F}_i$ est la limite du diagramme dans
$\textit{Ab}(\mathcal{C})$. Cela démontre l'existence des limites et (2).

\medskip\noindent
D'après Catégories, Lemme \ref{categories-lemma-adjoint-exact}, et puisque la
faisceautisation est adjointe à gauche au foncteur d'inclusion, nous voyons que
$\colim_i \mathcal{F}$ existe et est le faisceau associé à la colimite dans
$\textit{PAb}(\mathcal{C})$. Cela démontre l'existence des colimites et (4).

\medskip\noindent
Dans toute catégorie abélienne, les sommes directes finies coïncident avec les
produits finis. Ainsi (3) résulte de (2).

\medskip\noindent
Preuve de (5). L'assertion signifie qu'étant donné un système
$0 \to \mathcal{F}_i \to \mathcal{G}_i \to \mathcal{H}_i \to 0$
de suites exactes de faisceaux abéliens indexé par un ensemble filtrant $I$, la
suite
$0 \to \colim \mathcal{F}_i \to \colim \mathcal{G}_i \to
\colim \mathcal{H}_i \to 0$ est également exacte. Un argument formel utilisant
Homologie, Lemme \ref{homology-lemma-check-exactness}, et la définition des
colimites montre que la suite
$\colim \mathcal{F}_i \to \colim \mathcal{G}_i \to \colim \mathcal{H}_i \to 0$
est exacte. Remarquons que
$\colim \mathcal{F}_i \to \colim \mathcal{G}_i$ est le faisceau associé à
l'application entre les colimites de préfaisceaux ; celle-ci est injective
puisque chacune des applications
$\mathcal{F}_i \to \mathcal{G}_i$ l'est. La faisceautisation étant exacte, on
conclut.
\end{proof}







\section{Préfaisceaux abéliens libres}
\label{sites-modules-section-free-abelian-presheaf}

\noindent
Afin de préparer la notation de la définition suivante, convenons de noter le
groupe abélien libre sur un ensemble $S$ par\footnote{Dans d'autres chapitres,
la notation $\mathbf{Z}[S]$ désigne parfois l'anneau de polynômes sur
$\mathbf{Z}$ en les éléments de $S$.}
$\mathbf{Z}[S] = \bigoplus_{s \in S} \mathbf{Z}$. Il est caractérisé par la
propriété
$$
\Mor_{\textit{Ab}}(\mathbf{Z}[S], A)
=
\Mor_{\textit{Ensembles}}(S, A)
$$
Autrement dit, la construction $S \mapsto \mathbf{Z}[S]$ est adjointe à
gauche au foncteur d'oubli $\textit{Ab} \to \textit{Ensembles}$.

\begin{definition}
\label{sites-modules-definition-free-abelian-presheaf-on}
Soit $\mathcal{C}$ une catégorie. Soit $\mathcal{G}$ un préfaisceau
d'ensembles. Le {\it préfaisceau abélien libre}
$\mathbf{Z}_\mathcal{G}$ sur $\mathcal{G}$ est le préfaisceau abélien défini
par la règle
$$
U \longmapsto \mathbf{Z}[\mathcal{G}(U)].
$$
Dans le cas particulier $\mathcal{G} = h_X$ d'un préfaisceau représentable
associé à un objet $X$ de $\mathcal{C}$, nous utilisons la notation
$\mathbf{Z}_X = \mathbf{Z}_{h_X}$. Autrement dit,
$$
\mathbf{Z}_X(U) = \mathbf{Z}[\Mor_\mathcal{C}(U, X)].
$$
\end{definition}

\noindent
Cette construction est manifestement fonctorielle en le préfaisceau
$\mathcal{G}$. Elle est en fait adjointe au foncteur d'oubli
$\textit{PAb}(\mathcal{C}) \to \textit{PSh}(\mathcal{C})$. Voici l'énoncé
précis.

\begin{lemma}
\label{sites-modules-lemma-obvious-adjointness}
Soit $\mathcal{C}$ une catégorie. Soient $\mathcal{G}$ et $\mathcal{F}$ des
préfaisceaux d'ensembles. Soit $\mathcal{A}$ un préfaisceau abélien. Soit $U$
un objet de $\mathcal{C}$. Alors nous avons
\begin{align*}
\Mor_{\textit{PSh}(\mathcal{C})}(h_U, \mathcal{F})
& =
\mathcal{F}(U), \\
\Mor_{\textit{PAb}(\mathcal{C})}(\mathbf{Z}_\mathcal{G}, \mathcal{A})
& =
\Mor_{\textit{PSh}(\mathcal{C})}(\mathcal{G}, \mathcal{A}), \\
\Mor_{\textit{PAb}(\mathcal{C})}(\mathbf{Z}_U, \mathcal{A})
& =
\mathcal{A}(U).
\end{align*}
Toutes ces égalités sont fonctorielles.
\end{lemma}

\begin{proof}
Omis.
\end{proof}

\begin{lemma}
\label{sites-modules-lemma-coproduct-sum-free-abelian-presheaf}
Soit $\mathcal{C}$ une catégorie. Soit $I$ un ensemble. Pour chaque
$i \in I$, soit $\mathcal{G}_i$ un préfaisceau d'ensembles. Alors
$$
\mathbf{Z}_{\coprod_i \mathcal{G}_i}
=
\bigoplus\nolimits_{i \in I} \mathbf{Z}_{\mathcal{G}_i}
$$
dans $\textit{PAb}(\mathcal{C})$.
\end{lemma}

\begin{proof}
Omis.
\end{proof}



\section{Faisceaux abéliens libres}
\label{sites-modules-section-free-abelian-sheaf}

\noindent
Voici la notion de faisceau abélien libre sur un faisceau d'ensembles.

\begin{definition}
\label{sites-modules-definition-free-abelian-sheaf-on}
Soit $\mathcal{C}$ un site. Soit $\mathcal{G}$ un préfaisceau d'ensembles. Le
{\it faisceau abélien libre} $\mathbf{Z}_\mathcal{G}^\#$ sur $\mathcal{G}$
est le faisceau abélien $\mathbf{Z}_\mathcal{G}^\#$ associé au préfaisceau
abélien libre sur $\mathcal{G}$. Dans le cas particulier
$\mathcal{G} = h_X$ d'un préfaisceau représentable associé à un objet $X$ de
$\mathcal{C}$, nous utilisons la notation $\mathbf{Z}_X^\#$.
\end{definition}

\noindent
Cette construction est manifestement fonctorielle en le préfaisceau
$\mathcal{G}$. Elle fournit en fait un adjoint au foncteur d'oubli
$\textit{Ab}(\mathcal{C}) \to \Sh(\mathcal{C})$. Voici l'énoncé précis.

\begin{lemma}
\label{sites-modules-lemma-obvious-adjointness-sheaves}
Soit $\mathcal{C}$ un site. Soient $\mathcal{G}$ et $\mathcal{F}$ des
faisceaux d'ensembles. Soit $\mathcal{A}$ un faisceau abélien. Soit $U$ un
objet de $\mathcal{C}$. Alors nous avons
\begin{align*}
\Mor_{\Sh(\mathcal{C})}(h_U^\#, \mathcal{F})
& =
\mathcal{F}(U), \\
\Mor_{\textit{Ab}(\mathcal{C})}(\mathbf{Z}_\mathcal{G}^\#,
\mathcal{A})
& =
\Mor_{\Sh(\mathcal{C})}(\mathcal{G}, \mathcal{A}), \\
\Mor_{\textit{Ab}(\mathcal{C})}(\mathbf{Z}_U^\#, \mathcal{A})
& =
\mathcal{A}(U).
\end{align*}
Toutes ces égalités sont fonctorielles.
\end{lemma}

\begin{proof}
Omis.
\end{proof}

\begin{lemma}
\label{sites-modules-lemma-may-sheafify-before-abelianize}
Soit $\mathcal{C}$ un site. Soit $\mathcal{G}$ un préfaisceau d'ensembles.
Alors $\mathbf{Z}_\mathcal{G}^\# = (\mathbf{Z}_{\mathcal{G}^\#})^\#$.
\end{lemma}

\begin{proof}
Omis.
\end{proof}








\section{Sites annelés}
\label{sites-modules-section-ringed-sites}

\noindent
Dans ce chapitre, nous travaillons principalement avec des faisceaux de modules
sur un site annelé. Il nous faut donc définir cette notion.

\begin{definition}
\label{sites-modules-definition-ringed-site}
Sites annelés.
\begin{enumerate}
\item Un {\it site annelé} est un couple $(\mathcal{C}, \mathcal{O})$, où
$\mathcal{C}$ est un site et $\mathcal{O}$ un faisceau d'anneaux sur
$\mathcal{C}$. Le faisceau $\mathcal{O}$ est appelé le {\it faisceau
structural} du site annelé.
\item Soient $(\mathcal{C}, \mathcal{O})$ et
$(\mathcal{C}', \mathcal{O}')$ des sites annelés. Un {\it morphisme de sites
annelés}
$$
(f, f^\sharp) :
(\mathcal{C}, \mathcal{O})
\longrightarrow
(\mathcal{C}', \mathcal{O}')
$$
est la donnée d'un morphisme de sites $f : \mathcal{C} \to \mathcal{C}'$ (voir
Sites, Définition \ref{sites-definition-morphism-sites}) et d'une application
de faisceaux d'anneaux $f^\sharp : f^{-1}\mathcal{O}' \to \mathcal{O}$, ce qui,
par adjonction, revient à donner une application de faisceaux d'anneaux
$f^\sharp : \mathcal{O}' \to f_*\mathcal{O}$.
\item Soient
$(f, f^\sharp) :
(\mathcal{C}_1, \mathcal{O}_1) \to (\mathcal{C}_2, \mathcal{O}_2)$ et
$(g, g^\sharp) :
(\mathcal{C}_2, \mathcal{O}_2) \to (\mathcal{C}_3, \mathcal{O}_3)$
des morphismes de sites annelés. Nous définissons alors la {\it composition des
morphismes de sites annelés} par la règle
$$
(g, g^\sharp) \circ (f, f^\sharp) = (g \circ f, f^\sharp \circ g^\sharp).
$$
Nous utilisons ici la composition des morphismes de sites définie dans Sites,
Définition \ref{sites-definition-composition-morphisms-sites}, et
$f^\sharp \circ g^\sharp$ désigne le morphisme de faisceaux d'anneaux
$$
\mathcal{O}_3 \xrightarrow{g^\sharp} g_*\mathcal{O}_2
\xrightarrow{g_*f^\sharp} g_*f_*\mathcal{O}_1 = (g \circ f)_*\mathcal{O}_1
$$
\end{enumerate}
\end{definition}






\section{Topos annelés}
\label{sites-modules-section-ringed-topoi}

\noindent
Un topos annelé est simplement un site annelé, à ceci près que la notion de
morphisme de topos annelés diffère de celle de morphisme de sites annelés.

\begin{definition}
\label{sites-modules-definition-ringed-topos}
Topos annelés.
\begin{enumerate}
\item Un {\it topos annelé} est un couple
$(\Sh(\mathcal{C}), \mathcal{O})$
où $\mathcal{C}$ est un site et $\mathcal{O}$ un faisceau d'anneaux sur
$\mathcal{C}$. Le faisceau $\mathcal{O}$ est appelé le {\it faisceau
structural} du topos annelé.
\item Soient $(\Sh(\mathcal{C}), \mathcal{O})$ et
$(\Sh(\mathcal{C}'), \mathcal{O}')$ des topos annelés. Un {\it morphisme de
topos annelés}
$$
(f, f^\sharp) :
(\Sh(\mathcal{C}), \mathcal{O})
\longrightarrow
(\Sh(\mathcal{C}'), \mathcal{O}')
$$
est la donnée d'un morphisme de topos
$f : \Sh(\mathcal{C}) \to \Sh(\mathcal{C}')$ (voir Sites, Définition
\ref{sites-definition-topos}) et d'une application de faisceaux d'anneaux
$f^\sharp : f^{-1}\mathcal{O}' \to \mathcal{O}$, ce qui, par adjonction,
revient à donner une application de faisceaux d'anneaux
$f^\sharp : \mathcal{O}' \to f_*\mathcal{O}$.
\item Soient
$(f, f^\sharp) :
(\Sh(\mathcal{C}_1), \mathcal{O}_1)
\to (\Sh(\mathcal{C}_2), \mathcal{O}_2)$ et
$(g, g^\sharp) :
(\Sh(\mathcal{C}_2), \mathcal{O}_2) \to
(\Sh(\mathcal{C}_3), \mathcal{O}_3)$
des morphismes de topos annelés. Nous définissons alors la {\it composition des
morphismes de topos annelés} par la règle
$$
(g, g^\sharp) \circ (f, f^\sharp) = (g \circ f, f^\sharp \circ g^\sharp).
$$
Nous utilisons ici la composition des morphismes de topos définie dans Sites,
Définition \ref{sites-definition-topos}, et $f^\sharp \circ g^\sharp$ désigne
le morphisme de faisceaux d'anneaux
$$
\mathcal{O}_3 \xrightarrow{g^\sharp} g_*\mathcal{O}_2
\xrightarrow{g_*f^\sharp} g_*f_*\mathcal{O}_1 = (g \circ f)_*\mathcal{O}_1
$$
\end{enumerate}
\end{definition}

\noindent
Tout morphisme de topos annelés est la composée d'une équivalence de topos
annelés et d'un morphisme de topos annelés associé à un morphisme de sites
annelés. Voici l'énoncé précis.

\begin{lemma}
\label{sites-modules-lemma-morphism-ringed-topoi-comes-from-morphism-ringed-sites}
Soit $(f, f^\sharp) :
(\Sh(\mathcal{C}), \mathcal{O}_\mathcal{C})
\to (\Sh(\mathcal{D}), \mathcal{O}_\mathcal{D})$
un morphisme de topos annelés. Il existe une factorisation
$$
\xymatrix{
(\Sh(\mathcal{C}), \mathcal{O}_\mathcal{C})
\ar[rr]_{(f, f^\sharp)}
\ar[d]_{(g, g^\sharp)}
& &
(\Sh(\mathcal{D}), \mathcal{O}_\mathcal{D}) \ar[d]^{(e, e^\sharp)}
\\
(\Sh(\mathcal{C}'), \mathcal{O}_{\mathcal{C}'})
\ar[rr]^{(h, h^\sharp)} & &
(\Sh(\mathcal{D}'), \mathcal{O}_{\mathcal{D}'})
}
$$
où
\begin{enumerate}
\item $g : \Sh(\mathcal{C}) \to \Sh(\mathcal{C}')$ est une équivalence de
topos induite par un foncteur cocontinu spécial
$\mathcal{C} \to \mathcal{C}'$ (voir Sites, Définition
\ref{sites-definition-special-cocontinuous-functor}),
\item $e : \Sh(\mathcal{D}) \to \Sh(\mathcal{D}')$ est une équivalence de
topos induite par un foncteur cocontinu spécial
$\mathcal{D} \to \mathcal{D}'$ (voir Sites, Définition
\ref{sites-definition-special-cocontinuous-functor}),
\item $\mathcal{O}_{\mathcal{C}'} = g_*\mathcal{O}_\mathcal{C}$
et $g^\sharp$ est l'application évidente,
\item $\mathcal{O}_{\mathcal{D}'} = e_*\mathcal{O}_\mathcal{D}$
et $e^\sharp$ est l'application évidente,
\item les sites $\mathcal{C}'$ et $\mathcal{D}'$ possèdent des objets finaux
et des produits fibrés (autrement dit, toutes les limites finies),
\item $h$ est un morphisme de sites induit par un foncteur continu
$u : \mathcal{D}' \to \mathcal{C}'$ qui commute à toutes les limites finies
(autrement dit, qui satisfait les hypothèses de Sites, Proposition
\ref{sites-proposition-get-morphism}), et
\item étant données des familles quelconques de faisceaux $\mathcal{F}_i$
(resp.\ $\mathcal{G}_j$) sur $\mathcal{C}$ (resp.\ $\mathcal{D}$), nous pouvons
supposer que chacun est un faisceau représentable sur $\mathcal{C}'$
(resp.\ $\mathcal{D}'$).
\end{enumerate}
De plus, si $(f, f^\sharp)$ est une équivalence de topos annelés, nous pouvons
choisir le diagramme de sorte que
$\mathcal{C}' = \mathcal{D}'$,
$\mathcal{O}_{\mathcal{C}'} = \mathcal{O}_{\mathcal{D}'}$
et que $(h, h^\sharp)$ soit l'identité.
\end{lemma}

\begin{proof}
Cela résulte de Sites, Lemme
\ref{sites-lemma-morphism-topoi-comes-from-morphism-sites}, et de Sites,
Remarques \ref{sites-remark-morphism-topoi-comes-from-morphism-sites} et
\ref{sites-remark-equivalence-topoi-comes-from-morphism-sites}. Il suffit de
transporter simultanément les faisceaux d'anneaux. Nous omettons certains
détails.
\end{proof}







\section{2-morphismes de topos annelés}
\label{sites-modules-section-2-category}

\noindent
Cette brève section concerne la notion de $2$-morphisme de topos annelés.

\begin{definition}
\label{sites-modules-definition-2-morphism-ringed-topoi}
Soient
$f, g :
(\Sh(\mathcal{C}), \mathcal{O}_\mathcal{C})
\to
(\Sh(\mathcal{D}), \mathcal{O}_\mathcal{D})$
deux morphismes de topos annelés. Un {\it 2-morphisme de $f$ vers $g$} est
la donnée d'une transformation de foncteurs $t : f_* \to g_*$ telle que
$$
\xymatrix{
& \mathcal{O}_\mathcal{D}
\ar[ld]_{f^\sharp}
\ar[rd]^{g^\sharp} \\
f_*\mathcal{O}_\mathcal{C} \ar[rr]^t & &
g_*\mathcal{O}_\mathcal{C}
}
$$
soit commutatif.
\end{definition}

\noindent
Nous représentons parfois $t$ graphiquement de la manière suivante :
$$
\xymatrix{
(\Sh(\mathcal{C}), \mathcal{O}_\mathcal{C})
\rrtwocell^f_g{t}
&
&
(\Sh(\mathcal{D}), \mathcal{O}_\mathcal{D})
}
$$
Comme dans
Sites, section \ref{sites-section-2-category}
donner un 2-morphisme $t : f_* \to g_*$ équivaut à donner
$t : g^{-1} \to f^{-1}$ (habituellement noté par le même symbole), de sorte
que le diagramme
$$
\xymatrix{
f^{-1}\mathcal{O}_\mathcal{D}
\ar[rd]_{f^\sharp}  & &
g^{-1}\mathcal{O}_\mathcal{D} \ar[ll]^t \ar[ld]^{g^\sharp} \\
& \mathcal{O}_\mathcal{C}
}
$$
soit commutatif. Comme dans
Sites, section \ref{sites-section-2-category}
les axiomes d'une 2-catégorie stricte sont satisfaits pour les compositions
horizontale et verticale définies comme dans la référence citée.










\section{Préfaisceaux de modules}
\label{sites-modules-section-presheaves-modules}

\noindent
Soit $\mathcal{C}$ une catégorie. Soit $\mathcal{O}$ un préfaisceau d'anneaux
sur $\mathcal{C}$. Nous n'avons pas encore défini les préfaisceaux de
$\mathcal{O}$-modules ; faisons-le maintenant.

\begin{definition}
\label{sites-modules-definition-presheaf-modules}
Soit $\mathcal{C}$ une catégorie et soit $\mathcal{O}$ un préfaisceau
d'anneaux sur $\mathcal{C}$.
\begin{enumerate}
\item Un {\it préfaisceau de $\mathcal{O}$-modules} est la donnée d'un
préfaisceau abélien $\mathcal{F}$ et d'une application de préfaisceaux
d'ensembles
$$
\mathcal{O} \times \mathcal{F} \longrightarrow \mathcal{F}
$$
tel que, pour tout objet $U$ de $\mathcal{C}$, l'application
$\mathcal{O}(U) \times \mathcal{F}(U) \to \mathcal{F}(U)$
définisse une structure de $\mathcal{O}(U)$-module sur le groupe abélien
$\mathcal{F}(U)$.
\item Un {\it morphisme $\varphi : \mathcal{F} \to \mathcal{G}$ de
préfaisceaux de $\mathcal{O}$-modules} est un morphisme de préfaisceaux
abéliens $\varphi : \mathcal{F} \to \mathcal{G}$ tel que le diagramme
$$
\xymatrix{
\mathcal{O} \times \mathcal{F} \ar[r] \ar[d]_{\text{id} \times \varphi} &
\mathcal{F} \ar[d]^{\varphi} \\
\mathcal{O} \times \mathcal{G} \ar[r] &
\mathcal{G}
}
$$
soit commutatif.
\item L'ensemble des morphismes de $\mathcal{O}$-modules ci-dessus est noté
$\Hom_\mathcal{O}(\mathcal{F}, \mathcal{G})$.
\item La catégorie des préfaisceaux de $\mathcal{O}$-modules est notée
$\textit{PMod}(\mathcal{O})$.
\end{enumerate}
\end{definition}

\noindent
Supposons que $\mathcal{O}_1 \to \mathcal{O}_2$ soit un morphisme de
préfaisceaux d'anneaux sur la catégorie $\mathcal{C}$. Si $\mathcal{F}$ est un
préfaisceau de $\mathcal{O}_2$-modules, nous pouvons alors considérer
$\mathcal{F}$ comme un préfaisceau de $\mathcal{O}_1$-modules au moyen de la composée
$$
\mathcal{O}_1 \times \mathcal{F}
\to
\mathcal{O}_2 \times \mathcal{F}
\to
\mathcal{F}.
$$
Nous notons parfois ce préfaisceau
$\mathcal{F}_{\mathcal{O}_1}$ pour signaler la restriction des scalaires. Nous
appelons ce préfaisceau la {\it restriction de $\mathcal{F}$}. Nous obtenons le
foncteur de restriction des scalaires
$$
\textit{PMod}(\mathcal{O}_2)
\longrightarrow
\textit{PMod}(\mathcal{O}_1)
$$

\medskip\noindent
Réciproquement, étant donné un préfaisceau de $\mathcal{O}_1$-modules
$\mathcal{G}$, nous pouvons construire un préfaisceau de
$\mathcal{O}_2$-modules
$\mathcal{O}_2 \otimes_{p, \mathcal{O}_1} \mathcal{G}$
par la règle
$$
U \longmapsto
\left(\mathcal{O}_2 \otimes_{p, \mathcal{O}_1} \mathcal{G}\right)(U)
=
\mathcal{O}_2(U) \otimes_{\mathcal{O}_1(U)} \mathcal{G}(U)
$$
pour $U \in \Ob(\mathcal{C})$, avec les applications de restriction évidentes.
L'indice $p$ signifie « préfaisceau » et non « point ». Ce préfaisceau est
appelé le préfaisceau produit tensoriel. Nous obtenons le foncteur
d'{\it extension des scalaires}
$$
\textit{PMod}(\mathcal{O}_1)
\longrightarrow
\textit{PMod}(\mathcal{O}_2)
$$

\begin{lemma}
\label{sites-modules-lemma-adjointness-tensor-restrict-presheaves}
Avec $\mathcal{C}$, $\mathcal{O}_1 \to \mathcal{O}_2$, $\mathcal{F}$ et
$\mathcal{G}$ comme ci-dessus, il existe une bijection canonique
$$
\Hom_{\mathcal{O}_1}(\mathcal{G}, \mathcal{F}_{\mathcal{O}_1})
=
\Hom_{\mathcal{O}_2}(
\mathcal{O}_2 \otimes_{p, \mathcal{O}_1} \mathcal{G},
\mathcal{F}
)
$$
Autrement dit, les foncteurs de restriction et d'extension des scalaires
définis ci-dessus sont adjoints l'un de l'autre.
\end{lemma}

\begin{proof}
Cela résulte du fait que, pour un homomorphisme d'anneaux $A \to B$, les
foncteurs de restriction et d'extension des scalaires sont adjoints l'un de
l'autre.
\end{proof}


\section{Faisceaux de modules}
\label{sites-modules-section-sheaves-modules}

\begin{definition}
\label{sites-modules-definition-sheaf-modules}
Soit $\mathcal{C}$ un site. Soit $\mathcal{O}$ un faisceau d'anneaux sur
$\mathcal{C}$.
\begin{enumerate}
\item Un {\it faisceau de $\mathcal{O}$-modules} est un préfaisceau de
$\mathcal{O}$-modules $\mathcal{F}$, voir Définition
\ref{sites-modules-definition-presheaf-modules}, tel que le préfaisceau de groupes abéliens
sous-jacent à $\mathcal{F}$ soit un faisceau.
\item Un {\it morphisme de faisceaux de $\mathcal{O}$-modules} est un
morphisme de préfaisceaux de $\mathcal{O}$-modules.
\item Étant donnés des faisceaux de $\mathcal{O}$-modules
$\mathcal{F}$ et $\mathcal{G}$, nous notons
$\Hom_\mathcal{O}(\mathcal{F}, \mathcal{G})$ l'ensemble de leurs morphismes de
faisceaux de $\mathcal{O}$-modules.
\item La catégorie des faisceaux de $\mathcal{O}$-modules est notée
$\textit{Mod}(\mathcal{O})$.
\end{enumerate}
\end{definition}

\noindent
Cette définition garde un certain sens lorsque $\mathcal{O}$ n'est qu'un
préfaisceau d'anneaux, bien que nous ne connaissions aucun exemple où cela soit
utile. Nous éviterons donc l'expression « faisceaux de $\mathcal{O}$-modules »
lorsque $\mathcal{O}$ n'est pas un faisceau d'anneaux.



\section{Faisceautisation des préfaisceaux de modules}
\label{sites-modules-section-sheafification-presheaves-modules}

\begin{lemma}
\label{sites-modules-lemma-sheafification-presheaf-modules}
Soit $\mathcal{C}$ un site. Soit $\mathcal{O}$ un préfaisceau d'anneaux sur
$\mathcal{C}$ et soit $\mathcal{F}$ un préfaisceau de $\mathcal{O}$-modules.
Soit $\mathcal{O}^\#$ le faisceau associé à $\mathcal{O}$ comme préfaisceau
d'anneaux, voir Sites, section
\ref{sites-section-sheaves-algebraic-structures}. Soit $\mathcal{F}^\#$ le
faisceau associé à $\mathcal{F}$ comme préfaisceau de groupes abéliens. Il
existe une unique application de faisceaux d'ensembles
$$
\mathcal{O}^\# \times \mathcal{F}^\#
\longrightarrow
\mathcal{F}^\#
$$
qui rende commutatif le diagramme
$$
\xymatrix{
\mathcal{O} \times \mathcal{F} \ar[r] \ar[d] &
\mathcal{F} \ar[d] \\
\mathcal{O}^\# \times \mathcal{F}^\# \ar[r] &
\mathcal{F}^\#
}
$$
et qui munisse $\mathcal{F}^\#$ d'une structure de faisceau de
$\mathcal{O}^\#$-modules. De plus, si $\mathcal{G}$ est un faisceau de
$\mathcal{O}^\#$-modules, tout morphisme de préfaisceaux de
$\mathcal{O}$-modules $\mathcal{F} \to \mathcal{G}$ (où $\mathcal{G}$ est
considéré par restriction comme un $\mathcal{O}$-module) se factorise de façon
unique sous la forme $\mathcal{F} \to \mathcal{F}^\# \to \mathcal{G}$, où
$\mathcal{F}^\# \to \mathcal{G}$ est un morphisme de
$\mathcal{O}^\#$-modules.
\end{lemma}

\begin{proof}
Omis.
\end{proof}

\noindent
Cela signifie précisément que le foncteur
$i : \textit{Mod}(\mathcal{O}^\#) \to \textit{PMod}(\mathcal{O})$
(qui combine la restriction des scalaires et l'inclusion des faisceaux dans
les préfaisceaux) et le foncteur de faisceautisation du lemme
${}^\# : \textit{PMod}(\mathcal{O}) \to \textit{Mod}(\mathcal{O}^\#)$
sont adjoints. En formule,
$$
\Mor_{\textit{PMod}(\mathcal{O})}(\mathcal{F}, i\mathcal{G})
=
\Mor_{\textit{Mod}(\mathcal{O}^\#)}(\mathcal{F}^\#, \mathcal{G})
$$
Un cas important se présente lorsque $\mathcal{O}$ est déjà un faisceau
d'anneaux. La formule devient alors
$$
\Mor_{\textit{PMod}(\mathcal{O})}(\mathcal{F}, i\mathcal{G})
=
\Mor_{\textit{Mod}(\mathcal{O})}(\mathcal{F}^\#, \mathcal{G})
$$
puisque $\mathcal{O} = \mathcal{O}^\#$ dans ce cas.

\begin{lemma}
\label{sites-modules-lemma-sheafification-exact}
Soit $\mathcal{C}$ un site. Soit $\mathcal{O}$ un préfaisceau d'anneaux sur
$\mathcal{C}$. Le foncteur de faisceautisation
$$
\textit{PMod}(\mathcal{O}) \longrightarrow \textit{Mod}(\mathcal{O}^\#), \quad
\mathcal{F} \longmapsto \mathcal{F}^\#
$$
est exact.
\end{lemma}

\begin{proof}
Cela est vrai parce que la faisceautisation
$\textit{PAb}(\mathcal{C}) \to \textit{Ab}(\mathcal{C})$ est exacte. Voir la
discussion de la section \ref{sites-modules-section-abelian-sheaves}.
\end{proof}

\noindent
Soit $\mathcal{C}$ un site. Soit
$\mathcal{O}_1 \to \mathcal{O}_2$ un morphisme de faisceaux d'anneaux sur
$\mathcal{C}$. À la section \ref{sites-modules-section-presheaves-modules}, nous avons défini
les foncteurs de restriction et d'extension des scalaires correspondants sur
les préfaisceaux de modules.

\medskip\noindent
Si $\mathcal{F}$ est un faisceau de $\mathcal{O}_2$-modules, la restriction
$\mathcal{F}_{\mathcal{O}_1}$ de $\mathcal{F}$ est manifestement un faisceau de
$\mathcal{O}_1$-modules. Nous obtenons le foncteur de restriction des scalaires
$$
\textit{Mod}(\mathcal{O}_2)
\longrightarrow
\textit{Mod}(\mathcal{O}_1)
$$

\medskip\noindent
Réciproquement, étant donné un faisceau de $\mathcal{O}_1$-modules
$\mathcal{G}$, le préfaisceau de $\mathcal{O}_2$-modules
$\mathcal{O}_2 \otimes_{p, \mathcal{O}_1} \mathcal{G}$
n'est en général pas un faisceau. Nous définissons donc le {\it faisceau produit
tensoriel}
$\mathcal{O}_2 \otimes_{\mathcal{O}_1} \mathcal{G}$
par la formule
$$
\mathcal{O}_2 \otimes_{\mathcal{O}_1} \mathcal{G}
=
(\mathcal{O}_2 \otimes_{p, \mathcal{O}_1} \mathcal{G})^\#
$$
comme le faisceau associé à la construction précédente sur les préfaisceaux.
Nous obtenons le foncteur d'{\it extension des scalaires}
$$
\textit{Mod}(\mathcal{O}_1)
\longrightarrow
\textit{Mod}(\mathcal{O}_2)
$$

\begin{lemma}
\label{sites-modules-lemma-adjointness-tensor-restrict}
Avec $X$, $\mathcal{O}_1$, $\mathcal{O}_2$, $\mathcal{F}$ et $\mathcal{G}$
comme ci-dessus, il existe une bijection canonique
$$
\Hom_{\mathcal{O}_1}(\mathcal{G}, \mathcal{F}_{\mathcal{O}_1})
=
\Hom_{\mathcal{O}_2}(
\mathcal{O}_2 \otimes_{\mathcal{O}_1} \mathcal{G},
\mathcal{F}
)
$$
Autrement dit, les foncteurs de restriction et d'extension des scalaires sont
adjoints l'un de l'autre.
\end{lemma}

\begin{proof}
Cela résulte du Lemme \ref{sites-modules-lemma-adjointness-tensor-restrict-presheaves} et du
fait que
$\Hom_{\mathcal{O}_2}(
\mathcal{O}_2 \otimes_{\mathcal{O}_1} \mathcal{G},
\mathcal{F}
)
=
\Hom_{\mathcal{O}_2}(
\mathcal{O}_2 \otimes_{p, \mathcal{O}_1} \mathcal{G},
\mathcal{F}
)$
puisque $\mathcal{F}$ est un faisceau.
\end{proof}

\begin{lemma}
\label{sites-modules-lemma-epimorphism-modules}
Soit $\mathcal{C}$ un site. Soit $\mathcal{O} \to \mathcal{O}'$ un
épimorphisme de faisceaux d'anneaux. Soient
$\mathcal{G}_1, \mathcal{G}_2$ des $\mathcal{O}'$-modules. Alors
$$
\Hom_{\mathcal{O}'}(\mathcal{G}_1, \mathcal{G}_2) =
\Hom_\mathcal{O}(\mathcal{G}_1, \mathcal{G}_2).
$$
Autrement dit, le foncteur de restriction
$\textit{Mod}(\mathcal{O}') \to \textit{Mod}(\mathcal{O})$ est pleinement
fidèle.
\end{lemma}

\begin{proof}
C'est la version faisceautique d'Algèbre, Lemme
\ref{algebra-lemma-epimorphism-modules}, et elle se démontre exactement de la
même manière.
\end{proof}




\section{Morphismes de topos et faisceaux de modules}
\label{sites-modules-section-sheaves-modules-functorial}

\noindent
Tout ce qui suit est immédiat. Nous formulons les résultats pour les morphismes
de topos, mais ils valent bien entendu aussi pour les morphismes de sites.

\begin{lemma}
\label{sites-modules-lemma-pushforward-module}
Soient $\mathcal{C}$ et $\mathcal{D}$ des sites. Soit
$f : \Sh(\mathcal{C}) \to \Sh(\mathcal{D})$ un morphisme de topos. Soit
$\mathcal{O}$ un faisceau d'anneaux sur $\mathcal{C}$ et soit $\mathcal{F}$ un
faisceau de $\mathcal{O}$-modules. Il existe une application naturelle de
faisceaux d'ensembles
$$
f_*\mathcal{O} \times f_*\mathcal{F}
\longrightarrow
f_*\mathcal{F}
$$
qui munit $f_*\mathcal{F}$ d'une structure de faisceau de
$f_*\mathcal{O}$-modules. Cette construction est fonctorielle en $\mathcal{F}$.
\end{lemma}

\begin{proof}
Notons $\mu : \mathcal{O} \times \mathcal{F} \to \mathcal{F}$ l'application
de multiplication. Rappelons que $f_*$ (sur les faisceaux d'ensembles) est
exact à gauche et commute donc aux produits. Ainsi $f_*\mu$ est une application
du type indiqué, ce qui démontre le lemme.
\end{proof}

\begin{lemma}
\label{sites-modules-lemma-pullback-module}
Soient $\mathcal{C}$ et $\mathcal{D}$ des sites. Soit
$f : \Sh(\mathcal{C}) \to \Sh(\mathcal{D})$ un morphisme de topos. Soit
$\mathcal{O}$ un faisceau d'anneaux sur $\mathcal{D}$ et soit $\mathcal{G}$ un
faisceau de $\mathcal{O}$-modules. Il existe une application naturelle de
faisceaux d'ensembles
$$
f^{-1}\mathcal{O} \times f^{-1}\mathcal{G}
\longrightarrow
f^{-1}\mathcal{G}
$$
qui munit $f^{-1}\mathcal{G}$ d'une structure de faisceau de
$f^{-1}\mathcal{O}$-modules. Cette construction est fonctorielle en
$\mathcal{G}$.
\end{lemma}

\begin{proof}
Notons $\mu : \mathcal{O} \times \mathcal{G} \to \mathcal{G}$ l'application
de multiplication. Rappelons que $f^{-1}$ (sur les faisceaux d'ensembles) est
exact et commute donc aux produits. Ainsi $f^{-1}\mu$ est une application du
type indiqué, ce qui démontre le lemme.
\end{proof}

\begin{lemma}
\label{sites-modules-lemma-adjoint-push-pull-modules}
Soient $\mathcal{C}$ et $\mathcal{D}$ des sites. Soit
$f : \Sh(\mathcal{C}) \to \Sh(\mathcal{D})$ un morphisme de topos. Soit
$\mathcal{O}$ un faisceau d'anneaux sur $\mathcal{D}$. Soit $\mathcal{G}$ un
faisceau de $\mathcal{O}$-modules et soit $\mathcal{F}$ un faisceau de
$f^{-1}\mathcal{O}$-modules. Alors
$$
\Mor_{\textit{Mod}(f^{-1}\mathcal{O})}(f^{-1}\mathcal{G}, \mathcal{F})
=
\Mor_{\textit{Mod}(\mathcal{O})}(\mathcal{G}, f_*\mathcal{F}).
$$
Nous utilisons ici les Lemmes \ref{sites-modules-lemma-pullback-module} et
\ref{sites-modules-lemma-pushforward-module}, et nous considérons $f_*\mathcal{F}$ comme un
$\mathcal{O}$-module par restriction des scalaires suivant
$\mathcal{O} \to f_*f^{-1}\mathcal{O}$.
\end{lemma}

\begin{proof}
Remarquons d'abord que
$$
\Mor_{\textit{Ab}(\mathcal{C})}(f^{-1}\mathcal{G}, \mathcal{F})
=
\Mor_{\textit{Ab}(\mathcal{D})}(\mathcal{G}, f_*\mathcal{F}).
$$
d'après Sites, Proposition
\ref{sites-proposition-functoriality-algebraic-structures-topoi}. Supposons que
$\alpha : f^{-1}\mathcal{G} \to \mathcal{F}$ et
$\beta : \mathcal{G} \to f_*\mathcal{F}$ soient des morphismes de faisceaux
abéliens qui se correspondent par la formule ci-dessus. Il faut montrer que
$\alpha$ est $f^{-1}\mathcal{O}$-linéaire si et seulement si $\beta$ est
$\mathcal{O}$-linéaire. Supposons par exemple que $\alpha$ soit
$f^{-1}\mathcal{O}$-linéaire. Alors $f_*\alpha$ est manifestement
$f_*f^{-1}\mathcal{O}$-linéaire, donc $\mathcal{O}$-linéaire puisque la
restriction des scalaires est fonctorielle. Il suffit donc de montrer que
l'application d'adjonction $\mathcal{G} \to f_*f^{-1}\mathcal{G}$ est
$\mathcal{O}$-linéaire. Comme $f_*$ et $f^{-1}$ commutent tous deux aux
produits (sur les faisceaux d'ensembles), cela revient à montrer que
$$
\xymatrix{
\mathcal{O} \times \mathcal{G} \ar[r] \ar[d] &
f_*f^{-1}(\mathcal{O} \times \mathcal{G}) \ar[d] \\
\mathcal{G} \ar[r] & f_*f^{-1}\mathcal{G}
}
$$
est commutatif. Cela résulte de ce que l'application d'adjonction
$\text{id}_{\Sh(\mathcal{D})} \to f_*f^{-1}$ est une transformation de
foncteurs. Nous omettons la démonstration de l'implication « $\beta$ linéaire
$\Rightarrow$ $\alpha$ linéaire ».
\end{proof}

\begin{lemma}
\label{sites-modules-lemma-adjoint-pull-push-modules}
Soient $\mathcal{C}$ et $\mathcal{D}$ des sites. Soit
$f : \Sh(\mathcal{C}) \to \Sh(\mathcal{D})$ un morphisme de topos. Soit
$\mathcal{O}$ un faisceau d'anneaux sur $\mathcal{C}$. Soit $\mathcal{F}$ un
faisceau de $\mathcal{O}$-modules et soit $\mathcal{G}$ un faisceau de
$f_*\mathcal{O}$-modules. Alors
$$
\Mor_{\textit{Mod}(\mathcal{O})}(
\mathcal{O} \otimes_{f^{-1}f_*\mathcal{O}} f^{-1}\mathcal{G}, \mathcal{F})
=
\Mor_{\textit{Mod}(f_*\mathcal{O})}(\mathcal{G}, f_*\mathcal{F}).
$$
Nous utilisons ici les Lemmes \ref{sites-modules-lemma-pullback-module} et
\ref{sites-modules-lemma-pushforward-module}, ainsi que l'application canonique
$f^{-1}f_*\mathcal{O} \to \mathcal{O}$ dans la définition du produit tensoriel.
\end{lemma}

\begin{proof}
Remarquons que
$$
\Mor_{\textit{Mod}(\mathcal{O})}(
\mathcal{O} \otimes_{f^{-1}f_*\mathcal{O}} f^{-1}\mathcal{G}, \mathcal{F})
=
\Mor_{\textit{Mod}(f^{-1}f_*\mathcal{O})}(
f^{-1}\mathcal{G}, \mathcal{F}_{f^{-1}f_*\mathcal{O}})
$$
d'après le Lemme \ref{sites-modules-lemma-adjointness-tensor-restrict}. Le résultat découle
donc du Lemme \ref{sites-modules-lemma-adjoint-push-pull-modules}.
\end{proof}






\section{Morphismes de topos annelés et modules}
\label{sites-modules-section-functoriality-modules}

\noindent
Nous avons maintenant introduit assez de notations pour définir l'image inverse
et l'image directe des modules le long d'un morphisme de topos annelés.

\begin{definition}
\label{sites-modules-definition-pushforward}
Soit
$(f, f^\sharp) :
(\Sh(\mathcal{C}), \mathcal{O}_\mathcal{C})
\to
(\Sh(\mathcal{D}), \mathcal{O}_\mathcal{D})$
un morphisme de topos annelés ou de sites annelés.
\begin{enumerate}
\item Soit $\mathcal{F}$ un faisceau de $\mathcal{O}_\mathcal{C}$-modules. Nous
définissons l'{\it image directe} de $\mathcal{F}$ comme le faisceau de
$\mathcal{O}_\mathcal{D}$-modules dont le faisceau de groupes abéliens
sous-jacent est $f_*\mathcal{F}$ et dont la structure de module est obtenue par
restriction des scalaires suivant
$f^\sharp : \mathcal{O}_\mathcal{D} \to f_*\mathcal{O}_\mathcal{C}$ à partir
de la structure de module
$$
f_*\mathcal{O}_\mathcal{C} \times f_*\mathcal{F}
\longrightarrow
f_*\mathcal{F}
$$
du Lemme \ref{sites-modules-lemma-pushforward-module}.
\item Soit $\mathcal{G}$ un faisceau de $\mathcal{O}_\mathcal{D}$-modules. Nous
définissons l'{\it image inverse} $f^*\mathcal{G}$ comme le faisceau de
$\mathcal{O}_\mathcal{C}$-modules donné par la formule
$$
f^*\mathcal{G}
=
\mathcal{O}_\mathcal{C} \otimes_{f^{-1}\mathcal{O}_\mathcal{D}}
f^{-1}\mathcal{G}
$$
où l'homomorphisme d'anneaux
$f^{-1}\mathcal{O}_\mathcal{D} \to \mathcal{O}_\mathcal{C}$
est $f^\sharp$, et où la structure de module est celle du Lemme
\ref{sites-modules-lemma-pullback-module}.
\end{enumerate}
\end{definition}

\noindent
Nous avons ainsi défini des foncteurs
\begin{eqnarray*}
f_* : \textit{Mod}(\mathcal{O}_\mathcal{C})
& \longrightarrow &
\textit{Mod}(\mathcal{O}_\mathcal{D}) \\
f^* : \textit{Mod}(\mathcal{O}_\mathcal{D})
& \longrightarrow &
\textit{Mod}(\mathcal{O}_\mathcal{C})
\end{eqnarray*}
Le dernier résultat sur ces foncteurs affirme qu'ils sont bien adjoints, comme
prévu.

\begin{lemma}
\label{sites-modules-lemma-adjoint-pullback-pushforward-modules}
Soit
$(f, f^\sharp) :
(\Sh(\mathcal{C}), \mathcal{O}_\mathcal{C})
\to
(\Sh(\mathcal{D}), \mathcal{O}_\mathcal{D})$
un morphisme de topos annelés ou de sites annelés. Soit $\mathcal{F}$ un
faisceau de $\mathcal{O}_\mathcal{C}$-modules et soit $\mathcal{G}$ un faisceau
de $\mathcal{O}_\mathcal{D}$-modules. Il existe une bijection canonique
$$
\Hom_{\mathcal{O}_\mathcal{C}}(f^*\mathcal{G}, \mathcal{F})
=
\Hom_{\mathcal{O}_\mathcal{D}}(\mathcal{G}, f_*\mathcal{F}).
$$
Autrement dit, le foncteur $f^*$ est adjoint à gauche de $f_*$.
\end{lemma}

\begin{proof}
Cela résulte des calculs précédents :
\begin{eqnarray*}
\Hom_{\mathcal{O}_\mathcal{C}}(f^*\mathcal{G}, \mathcal{F})
& = &
\Mor_{\textit{Mod}(\mathcal{O}_\mathcal{C})}(
\mathcal{O}_\mathcal{C}
\otimes_{f^{-1}\mathcal{O}_\mathcal{D}} f^{-1}\mathcal{G},
\mathcal{F}) \\
& = &
\Mor_{\textit{Mod}(f^{-1}\mathcal{O}_\mathcal{D})}(
f^{-1}\mathcal{G}, \mathcal{F}_{f^{-1}\mathcal{O}_\mathcal{D}}) \\
& = &
\Hom_{\mathcal{O}_\mathcal{D}}(\mathcal{G}, f_*\mathcal{F}).
\end{eqnarray*}
Nous utilisons ici les Lemmes \ref{sites-modules-lemma-adjointness-tensor-restrict} et
\ref{sites-modules-lemma-adjoint-push-pull-modules}.
\end{proof}

\begin{lemma}
\label{sites-modules-lemma-push-pull-composition-modules}
$(f, f^\sharp) :
(\Sh(\mathcal{C}_1), \mathcal{O}_1)
\to (\Sh(\mathcal{C}_2), \mathcal{O}_2)$ et
$(g, g^\sharp) :
(\Sh(\mathcal{C}_2), \mathcal{O}_2) \to
(\Sh(\mathcal{C}_3), \mathcal{O}_3)$
soient des morphismes de topos annelés. Il existe des isomorphismes canoniques
de foncteurs
$(g \circ f)_* \cong g_* \circ f_*$ et
$(g \circ f)^* \cong f^* \circ g^*$.
\end{lemma}

\begin{proof}
Cela résulte immédiatement des définitions.
\end{proof}





\section{La catégorie abélienne des faisceaux de modules}
\label{sites-modules-section-kernels}

\noindent
Soit $(\Sh(\mathcal{C}), \mathcal{O})$ un topos annelé. Soient
$\mathcal{F}$ et $\mathcal{G}$ des faisceaux de $\mathcal{O}$-modules, voir
Faisceaux, Définition \ref{sheaves-definition-sheaf-modules}. Soient
$\varphi, \psi : \mathcal{F} \to \mathcal{G}$ des morphismes de faisceaux de
$\mathcal{O}$-modules. Nous définissons
$\varphi + \psi : \mathcal{F} \to \mathcal{G}$ comme la somme de $\varphi$ et
$\psi$ en tant que morphismes de faisceaux abéliens. C'est manifestement encore
une application de $\mathcal{O}$-modules. La composition des applications de
$\mathcal{O}$-modules est en outre bilinéaire pour cette addition. Ainsi,
$\textit{Mod}(\mathcal{O})$ est une catégorie préadditive, voir Homologie,
Définition \ref{homology-definition-preadditive}.

\medskip\noindent
Nous noterons $0$ le faisceau de $\mathcal{O}$-modules qui prend la valeur
constante $\{0\}$ sur tout objet $U$ de $\mathcal{C}$. C'est manifestement à
la fois un objet final et un objet initial de $\textit{Mod}(\mathcal{O})$.
Étant donné un morphisme de $\mathcal{O}$-modules
$\varphi : \mathcal{F} \to \mathcal{G}$, les assertions suivantes sont
équivalentes : (a) $\varphi$ est nul, (b) $\varphi$ se factorise par $0$, et
(c) $\varphi$ est nul sur les sections au-dessus de tout objet $U$.

\medskip\noindent
En outre, étant donnés deux faisceaux $\mathcal{F}$ et $\mathcal{G}$ de
$\mathcal{O}$-modules, nous pouvons définir leur somme directe par
$$
\mathcal{F} \oplus \mathcal{G} = \mathcal{F} \times \mathcal{G}
$$
munie des applications évidentes $(i, j, p, q)$ comme dans Homologie,
Définition \ref{homology-definition-direct-sum}. Ainsi
$\textit{Mod}(\mathcal{O})$ est une catégorie additive, voir Homologie,
Définition \ref{homology-definition-additive-category}.

\medskip\noindent
Soit $\varphi : \mathcal{F} \to \mathcal{G}$ un morphisme de
$\mathcal{O}$-modules. Nous pouvons définir $\Ker(\varphi)$ comme le noyau de
$\varphi$ en tant qu'application de faisceaux abéliens. D'après la section
\ref{sites-modules-section-abelian-sheaves}, c'est le sous-faisceau de $\mathcal{F}$ dont les
sections sont
$$
\Ker(\varphi)(U) =
\{ s \in \mathcal{F}(U) \mid \varphi(s) = 0 \text{ dans } \mathcal{G}(U)\}
$$
pour tout objet $U$ de $\mathcal{C}$. On vérifie aisément qu'il s'agit bien
d'un noyau dans la catégorie des $\mathcal{O}$-modules. Autrement dit, un
morphisme $\alpha : \mathcal{H} \to \mathcal{F}$ se factorise par
$\Ker(\varphi)$ si et seulement si $\varphi \circ \alpha = 0$.

\medskip\noindent
De même, nous définissons $\Coker(\varphi)$ comme le conoyau de $\varphi$ en
tant qu'application de faisceaux abéliens. Il existe une unique application de
multiplication
$$
\mathcal{O} \times \Coker(\varphi) \longrightarrow \Coker(\varphi)
$$
qui fait de $\mathcal{G} \to \Coker(\varphi)$ un morphisme de
$\mathcal{O}$-modules (vérification omise). L'application
$\mathcal{G} \to \Coker(\varphi)$ est surjective (comme application de
faisceaux d'ensembles, voir section \ref{sites-modules-section-abelian-sheaves}). Pour montrer
que $\Coker(\varphi)$ est un conoyau dans $\textit{Mod}(\mathcal{O})$,
remarquons que si $\beta : \mathcal{G} \to \mathcal{H}$ est un morphisme de
$\mathcal{O}$-modules tel que $\beta \circ \varphi$ soit nul, alors, pour tout
objet $U$ de $\mathcal{C}$, $\beta$ induit une application de
$\mathcal{G}(U)/\varphi(\mathcal{F}(U))$ vers $\mathcal{H}(U)$. Par la
propriété universelle de la faisceautisation (voir Faisceaux, Lemme
\ref{sheaves-lemma-sheafification-presheaf-modules}), nous obtenons une
application canonique $\Coker(\varphi) \to \mathcal{H}$ telle que $\beta$ soit
la composée $\mathcal{G} \to \Coker(\varphi) \to \mathcal{H}$. Le morphisme
$\Coker(\varphi) \to \mathcal{H}$ est unique en raison de la surjectivité
mentionnée ci-dessus.

\begin{lemma}
\label{sites-modules-lemma-abelian}
Soit $(\Sh(\mathcal{C}), \mathcal{O})$ un topos annelé. La catégorie
$\textit{Mod}(\mathcal{O})$ est abélienne. Le foncteur d'oubli
$\textit{Mod}(\mathcal{O}) \to \textit{Ab}(\mathcal{C})$
est exact ; les noyaux, les conoyaux et l'exactitude des
$\mathcal{O}$-modules correspondent donc aux notions analogues pour les
faisceaux abéliens.
\end{lemma}

\begin{proof}
Nous avons vu ci-dessus que $\textit{Mod}(\mathcal{O})$ est une catégorie
additive, qu'elle possède noyaux et conoyaux, et que
$\textit{Mod}(\mathcal{O}) \to \textit{Ab}(\mathcal{C})$ les préserve. D'après
Homologie, Définition \ref{homology-definition-abelian-category}, il reste à
montrer que l'image et la coimage coïncident. Cela est clair puisque c'est vrai
dans $\textit{Ab}(\mathcal{C})$. Le lemme en résulte.
\end{proof}

\begin{lemma}
\label{sites-modules-lemma-limits-colimits}
Soit $(\Sh(\mathcal{C}), \mathcal{O})$ un topos annelé. Toutes les limites et
toutes les colimites existent dans $\textit{Mod}(\mathcal{O})$, et le foncteur
d'oubli
$\textit{Mod}(\mathcal{O}) \to \textit{Ab}(\mathcal{C})$
commute avec elles. De plus, les colimites filtrantes sont exactes.
\end{lemma}

\begin{proof}
La dernière assertion résulte de la première, car les colimites filtrantes sont
exactes dans $\textit{Ab}(\mathcal{C})$ d'après le Lemme
\ref{sites-modules-lemma-limits-colimits-abelian-sheaves}. Soit
$\mathcal{I} \to \textit{Mod}(\mathcal{C})$, $i \mapsto \mathcal{F}_i$, un
diagramme. Soit $\lim_i \mathcal{F}_i$ sa limite dans
$\textit{Ab}(\mathcal{C})$. La description de cette limite au Lemme
\ref{sites-modules-lemma-limits-colimits-abelian-sheaves} montre immédiatement qu'il existe
une multiplication
$$
\mathcal{O} \times \lim_i \mathcal{F}_i
\longrightarrow
\lim_i \mathcal{F}_i
$$
qui munit $\lim_i \mathcal{F}_i$ d'une structure de faisceau de
$\mathcal{O}$-modules. On voit aisément qu'il s'agit de la limite du diagramme
dans $\textit{Mod}(\mathcal{C})$. Soit $\colim_i \mathcal{F}_i$ la colimite du
diagramme dans $\textit{PAb}(\mathcal{C})$. La description de cette colimite
dans la démonstration du Lemme
\ref{sites-modules-lemma-limits-colimits-abelian-presheaves} montre immédiatement qu'il
existe une multiplication
$$
\mathcal{O} \times \colim_i \mathcal{F}_i
\longrightarrow
\colim_i \mathcal{F}_i
$$
qui munit $\colim_i \mathcal{F}_i$ d'une structure de préfaisceau de
$\mathcal{O}$-modules. En faisceautisant, nous obtenons le faisceau de
$\mathcal{O}$-modules $(\colim_i \mathcal{F}_i)^\#$, voir Lemme
\ref{sites-modules-lemma-sheafification-presheaf-modules}. On voit aisément que
$(\colim_i \mathcal{F}_i)^\#$ est la colimite du diagramme dans
$\textit{Mod}(\mathcal{O})$ et que, d'après le Lemme
\ref{sites-modules-lemma-limits-colimits-abelian-sheaves}, l'oubli de sa structure de
$\mathcal{O}$-module donne la colimite dans $\textit{Ab}(\mathcal{C})$.
\end{proof}

\noindent
L'existence des limites et des colimites permet d'exprimer par celles-ci les
propriétés d'exactitude des foncteurs définis sur la catégorie des
$\mathcal{O}$-modules, comme dans Catégories, section
\ref{categories-section-exact-functor}. Voir Homologie, Lemme
\ref{homology-lemma-exact-functor}, pour une description de ces propriétés en
termes de suites exactes courtes.

\begin{lemma}
\label{sites-modules-lemma-exactness-pushforward-pullback}
Soit $f : (\Sh(\mathcal{C}), \mathcal{O}_\mathcal{C})
\to (\Sh(\mathcal{D}), \mathcal{O}_\mathcal{D})$
un morphisme de topos annelés.
\begin{enumerate}
\item Le foncteur $f_*$ est exact à gauche. En fait, il commute à toutes les
limites.
\item Le foncteur $f^*$ est exact à droite. En fait, il commute à toutes les
colimites.
\end{enumerate}
\end{lemma}

\begin{proof}
Cela est vrai parce que $(f^*, f_*)$ est un couple de foncteurs adjoints, voir
Lemme \ref{sites-modules-lemma-adjoint-pullback-pushforward-modules}. Voir Catégories,
section \ref{categories-section-adjoint}.
\end{proof}

\begin{lemma}
\label{sites-modules-lemma-check-exactness-stalks}
Soit $\mathcal{C}$ un site. Si $\{p_i\}_{i \in I}$ est une famille conservative
de points, l'exactitude d'une suite de faisceaux abéliens peut être vérifiée sur
les germes aux points $p_i$, $i \in I$. Si $\mathcal{C}$ possède suffisamment
de points, l'exactitude d'une suite de faisceaux abéliens peut être vérifiée sur
les germes.
\end{lemma}

\begin{proof}
Cela résulte immédiatement de Sites, Lemme
\ref{sites-lemma-exactness-stalks}.
\end{proof}





\section{Exactitude de l'image directe}
\label{sites-modules-section-pushforward}

\noindent
Quelques lemmes techniques sur les propriétés d'exactitude de l'image directe.

\begin{lemma}
\label{sites-modules-lemma-reflect-surjections}
Soit $f : \Sh(\mathcal{C}) \to \Sh(\mathcal{D})$ un morphisme de topos. Les
assertions suivantes sont équivalentes :
\begin{enumerate}
\item $f^{-1}f_*\mathcal{F} \to \mathcal{F}$ est surjectif pour tout
$\mathcal{F}$ dans $\textit{Ab}(\mathcal{C})$, et
\item $f_* : \textit{Ab}(\mathcal{C}) \to \textit{Ab}(\mathcal{D})$
reflète les surjections.
\end{enumerate}
Dans ce cas, le foncteur
$f_* : \textit{Ab}(\mathcal{C}) \to \textit{Ab}(\mathcal{D})$
est fidèle.
\end{lemma}

\begin{proof}
Supposons (1). Soit $a : \mathcal{F} \to \mathcal{F}'$ une application de
faisceaux abéliens sur $\mathcal{C}$ telle que $f_*a$ soit surjective. Puisque
$f^{-1}$ est exact, il s'ensuit que
$f^{-1}f_*a : f^{-1}f_*\mathcal{F} \to f^{-1}f_*\mathcal{F}'$
est surjectif. Avec (1), cela implique que $a$ est surjectif. Ainsi (2) est
satisfaite.

\medskip\noindent
Supposons (2). Soit $\mathcal{F}$ un faisceau abélien sur $\mathcal{C}$. Nous
devons montrer que $f^{-1}f_*\mathcal{F} \to \mathcal{F}$ est surjectif.
D'après (2), il suffit de montrer que
$f_*f^{-1}f_*\mathcal{F} \to f_*\mathcal{F}$ est surjectif. Or cela est vrai,
car il existe une application canonique
$f_*\mathcal{F} \to f_*f^{-1}f_*\mathcal{F}$ qui en est un inverse d'un côté.

\medskip\noindent
Nous omettons la démonstration de la dernière assertion.
\end{proof}

\begin{lemma}
\label{sites-modules-lemma-exactness}
Soit $f : \Sh(\mathcal{C}) \to \Sh(\mathcal{D})$ un morphisme de topos.
Supposons qu'au moins l'une des propriétés suivantes soit satisfaite :
\begin{enumerate}
\item $f_*$ transforme les surjections de faisceaux d'ensembles en
surjections,
\item $f_*$ transforme les surjections de faisceaux abéliens en surjections,
\item $f_*$ commute aux coégalisateurs de faisceaux d'ensembles,
\item $f_*$ commute aux sommes amalgamées de faisceaux d'ensembles,
\end{enumerate}
Alors $f_* : \textit{Ab}(\mathcal{C}) \to \textit{Ab}(\mathcal{D})$ est
exact.
\end{lemma}

\begin{proof}
Puisque $f_* : \textit{Ab}(\mathcal{C}) \to \textit{Ab}(\mathcal{D})$ est un
adjoint à droite, nous savons déjà qu'il transforme une suite exacte courte
$0 \to \mathcal{F}_1 \to \mathcal{F}_2 \to \mathcal{F}_3 \to 0$
de faisceaux abéliens sur $\mathcal{C}$ en une suite exacte
$$
0 \to f_*\mathcal{F}_1 \to f_*\mathcal{F}_2 \to f_*\mathcal{F}_3
$$
voir Catégories, sections \ref{categories-section-exact-functor} et
\ref{categories-section-adjoint}, ainsi qu'Homologie, section
\ref{homology-section-functors}. Il suffit donc de démontrer que l'application
$f_*\mathcal{F}_2 \to f_*\mathcal{F}_3$ est surjective. Si (1) ou (2) est
satisfaite, cela découle des définitions. D'après Sites, Lemme
\ref{sites-lemma-exactness-properties}, chacune des propriétés (3) et (4)
implique formellement (1), ce qui conclut également dans ces cas.
\end{proof}

\begin{lemma}
\label{sites-modules-lemma-morphism-ringed-sites-almost-cocontinuous}
Soit $f : \mathcal{D} \to \mathcal{C}$ un morphisme de sites associé au
foncteur continu $u : \mathcal{C} \to \mathcal{D}$. Supposons que $u$ soit
presque cocontinu. Alors
\begin{enumerate}
\item $f_* : \textit{Ab}(\mathcal{D}) \to \textit{Ab}(\mathcal{C})$ est
exact.
\item si l'on se donne
$f^\sharp : f^{-1}\mathcal{O}_\mathcal{C} \to \mathcal{O}_\mathcal{D}$ de
sorte que $f$ devienne un morphisme de sites annelés, alors
$f_* : \textit{Mod}(\mathcal{O}_\mathcal{D}) \to
\textit{Mod}(\mathcal{O}_\mathcal{C})$ est exact.
\end{enumerate}
\end{lemma}

\begin{proof}
La partie (2) résulte de la partie (1) par le Lemme
\ref{sites-modules-lemma-limits-colimits}. La partie (1) résulte de Sites, Lemmes
\ref{sites-lemma-morphism-of-sites-almost-cocontinuous} et
\ref{sites-lemma-exactness-properties}.
\end{proof}





\section{Exactitude de l'adjoint à gauche}
\label{sites-modules-section-exactness-lower-shriek}

\noindent
Soit $u : \mathcal{C} \to \mathcal{D}$ un foncteur entre sites. Supposons que
\begin{enumerate}
\item[(a)] $u$ soit cocontinu, et
\item[(b)] $u$ soit continu.
\end{enumerate}
Soit $g : \Sh(\mathcal{C}) \to \Sh(\mathcal{D})$ le morphisme de topos associé
à $u$, voir Sites, Lemme \ref{sites-lemma-cocontinuous-morphism-topoi}.
Rappelons que $g^{-1} = u^p$, c'est-à-dire que $g^{-1}$ est donné par la
formule simple
$(g^{-1}\mathcal{G})(U) = \mathcal{G}(u(U))$, voir Sites, Lemme
\ref{sites-lemma-when-shriek}.
Nous voulons montrer que
$g^{-1} : \textit{Ab}(\mathcal{D}) \to \textit{Ab}(\mathcal{C})$
admet un adjoint à gauche $g_!$. D'après Sites, Lemme
\ref{sites-lemma-when-shriek}, le foncteur
$g^{Sh}_! = (u_p\ )^\#$ est un adjoint à gauche sur les faisceaux d'ensembles.
De plus, nous savons que $g^{Sh}_!\mathcal{F}$ est le faisceau associé au
préfaisceau
$$
V \longmapsto \colim_{V \to u(U)} \mathcal{F}(U)
$$
où la colimite est indexée par $(\mathcal{I}_V^u)^{opp}$ et prise dans la
catégorie des ensembles. La définition suivante est donc naturelle.

\begin{definition}
\label{sites-modules-definition-g-shriek}
Soit $u : \mathcal{C} \to \mathcal{D}$ satisfaisant (a) et (b) ci-dessus. Pour
$\mathcal{F} \in \textit{PAb}(\mathcal{C})$, nous définissons
{\it $g_{p!}\mathcal{F}$} comme le préfaisceau
$$
V \longmapsto \colim_{V \to u(U)} \mathcal{F}(U)
$$
où les colimites sur $(\mathcal{I}_V^u)^{opp}$ sont prises dans
$\textit{Ab}$. Pour $\mathcal{F} \in \textit{PAb}(\mathcal{C})$, nous posons
{\it $g_!\mathcal{F} = (g_{p!}\mathcal{F})^\#$}.
\end{definition}

\noindent
Nous sommes si explicites parce que les foncteurs $g^{Sh}_!$ et $g_!$ sont
différents. Lorsque nous les utilisons tous deux, il faut soigneusement les
distinguer.

\begin{lemma}
\label{sites-modules-lemma-g-shriek-adjoint}
Le foncteur $g_{p!}$ est adjoint à gauche du foncteur $u^p$. Le foncteur $g_!$
est adjoint à gauche du foncteur $g^{-1}$. Autrement dit, les formules
\begin{align*}
\Mor_{\textit{PAb}(\mathcal{C})}(\mathcal{F}, u^p\mathcal{G})
& =
\Mor_{\textit{PAb}(\mathcal{D})}(g_{p!}\mathcal{F}, \mathcal{G}), \\
\Mor_{\textit{Ab}(\mathcal{C})}(\mathcal{F}, g^{-1}\mathcal{G})
& =
\Mor_{\textit{Ab}(\mathcal{D})}(g_!\mathcal{F}, \mathcal{G})
\end{align*}
sont bifonctorielles en $\mathcal{F}$ et $\mathcal{G}$.
\end{lemma}

\begin{proof}
La seconde formule résulte formellement de la première, car si $\mathcal{F}$ et
$\mathcal{G}$ sont des faisceaux abéliens, alors
\begin{align*}
\Mor_{\textit{Ab}(\mathcal{C})}(\mathcal{F}, g^{-1}\mathcal{G})
& =
\Mor_{\textit{PAb}(\mathcal{D})}(g_{p!}\mathcal{F}, \mathcal{G}) \\
& =
\Mor_{\textit{Ab}(\mathcal{D})}(g_!\mathcal{F}, \mathcal{G})
\end{align*}
par la propriété universelle de la faisceautisation.

\medskip\noindent
Pour démontrer la première formule, soient $\mathcal{F}$ et $\mathcal{G}$ des
préfaisceaux abéliens. Nous construirons des applications du groupe de gauche
vers celui de droite et omettrons de vérifier qu'elles sont inverses l'une de
l'autre.

\medskip\noindent
Il existe une application canonique de préfaisceaux abéliens
$\mathcal{F} \to u^pg_{p!}\mathcal{F}$ qui, sur les sections au-dessus de $U$,
est l'application naturelle
$\mathcal{F}(U) \to \colim_{u(U) \to u(U')} \mathcal{F}(U')$, voir Sites,
Lemme \ref{sites-lemma-recover}.
Étant donnée une application $\alpha : g_{p!}\mathcal{F} \to \mathcal{G}$,
nous obtenons $u^p\alpha : u^pg_{p!}\mathcal{F} \to u^p\mathcal{G}$, que nous
pouvons précomposer par l'application
$\mathcal{F} \to u^pg_{p!}\mathcal{F}$.

\medskip\noindent
Il existe une application canonique de préfaisceaux abéliens
$g_{p!}u^p\mathcal{G} \to \mathcal{G}$ qui, sur les sections au-dessus de $V$,
est l'application naturelle
$\colim_{V \to u(U)} \mathcal{G}(u(U)) \to \mathcal{G}(V)$.
Elle envoie une section $s \in \mathcal{G}(u(U))$ dans le facteur correspondant
à $t : V \to u(U)$ sur $t^*s \in \mathcal{G}(V)$. Ainsi, étant donnée une
application $\beta : \mathcal{F} \to u^p\mathcal{G}$, nous obtenons une
application $g_{p!}\beta : g_{p!}\mathcal{F} \to g_{p!}u^p\mathcal{G}$, que
nous pouvons postcomposer par l'application
$g_{p!}u^p\mathcal{G} \to \mathcal{G}$ ci-dessus.
\end{proof}

\begin{lemma}
\label{sites-modules-lemma-exactness-lower-shriek}
Soient $\mathcal{C}$ et $\mathcal{D}$ des sites. Soit
$u : \mathcal{C} \to \mathcal{D}$ un foncteur. Supposons que
\begin{enumerate}
\item[(a)] $u$ soit cocontinu,
\item[(b)] $u$ soit continu, et
\item[(c)] les produits fibrés et les égalisateurs existent dans
$\mathcal{C}$ et $u$ commute à ces constructions.
\end{enumerate}
Dans ce cas, le foncteur
$g_! : \textit{Ab}(\mathcal{C}) \to \textit{Ab}(\mathcal{D})$
est exact.
\end{lemma}

\begin{proof}
Comparer avec Sites, Lemme \ref{sites-lemma-preserve-equalizers}. Supposons
(a), (b) et (c). Nous savons déjà que $g_!$ est exact à droite puisqu'il est un
adjoint à gauche, voir Catégories, Lemme
\ref{categories-lemma-exact-adjoint}, et Homologie, section
\ref{homology-section-functors}. Nous avons $g_! = (g_{p!}\ )^\#$. Il faut
montrer que $g_!$ transforme les applications injectives de faisceaux abéliens
en applications injectives de préfaisceaux abéliens. Rappelons que la
faisceautisation des préfaisceaux abéliens est exacte, voir Lemme
\ref{sites-modules-lemma-limits-colimits-abelian-sheaves}. Il suffit donc de montrer que
$g_{p!}$ transforme les applications injectives de préfaisceaux abéliens en
applications injectives de préfaisceaux abéliens. Pour cela, il suffit que les
colimites sur les catégories
$(\mathcal{I}_V^u)^{opp}$ de Sites, section
\ref{sites-section-functoriality-PSh}
transforment les applications injectives entre diagrammes en injections. Cela
résulte de Sites, Lemme \ref{sites-lemma-almost-directed}, et d'Algèbre, Lemme
\ref{algebra-lemma-almost-directed-colimit-exact}.
\end{proof}

\begin{lemma}
\label{sites-modules-lemma-back-and-forth}
Soient $\mathcal{C}$ et $\mathcal{D}$ des sites. Soit
$u : \mathcal{C} \to \mathcal{D}$ un foncteur. Supposons que
\begin{enumerate}
\item[(a)] $u$ soit cocontinu,
\item[(b)] $u$ soit continu, et
\item[(c)] $u$ soit pleinement fidèle.
\end{enumerate}
Pour $g_!, g^{-1}, g_*$ comme ci-dessus, les applications canoniques
$\mathcal{F} \to g^{-1}g_!\mathcal{F}$ et
$g^{-1}g_*\mathcal{F} \to \mathcal{F}$
sont des isomorphismes pour tout faisceau abélien $\mathcal{F}$ sur
$\mathcal{C}$.
\end{lemma}

\begin{proof}
L'application $g^{-1}g_*\mathcal{F} \to \mathcal{F}$ est un isomorphisme
d'après Sites, Lemme \ref{sites-lemma-back-and-forth}, et puisque les images
inverse et directe des faisceaux abéliens coïncident avec celles de leurs
faisceaux d'ensembles sous-jacents.

\medskip\noindent
Choisissons $U \in \Ob(\mathcal{C})$. Nous allons montrer que
$g^{-1}g_!\mathcal{F}(U) = \mathcal{F}(U)$. Remarquons d'abord que
$g^{-1}g_!\mathcal{F}(U) = g_!\mathcal{F}(u(U))$. Il suffit donc de montrer que
$g_!\mathcal{F}(u(U)) = \mathcal{F}(U)$. Nous savons que $g_!\mathcal{F}$ est
le faisceau abélien associé au préfaisceau $g_{p!}\mathcal{F}$ défini par la
règle
$$
V \longmapsto \colim_{V \to u(U')} \mathcal{F}(U')
$$
où la colimite est prise dans $\textit{Ab}$. Si $V = u(U)$, alors, puisque $u$
est pleinement fidèle, cette colimite est indexée par les $U \to U'$. Nous en
concluons que $g_{p!}\mathcal{F}(u(U)) = \mathcal{F}(U)$. Comme $u$ est
cocontinu et continu, tout recouvrement de $u(U)$ dans $\mathcal{D}$ peut être
raffiné par un recouvrement (!) $\{u(U_i) \to u(U)\}$ de $\mathcal{D}$, où
$\{U_i \to U\}$ est un recouvrement dans $\mathcal{C}$. Il s'ensuit également
que $(g_{p!}\mathcal{F})^+(u(U)) = \mathcal{F}(U)$, car, dans la colimite qui
définit la valeur de $(g_{p!}\mathcal{F})^+$ en $u(U)$, nous pouvons nous
restreindre au système cofinal des recouvrements
$\{u(U_i) \to u(U)\}$ ci-dessus. Ainsi
$(g_{p!}\mathcal{F})^+(u(U)) = \mathcal{F}(U)$ pour tout objet $U$ de
$\mathcal{C}$. En répétant encore une fois cet argument, nous obtenons
$(g_{p!}\mathcal{F})^\#(u(U)) = \mathcal{F}(U)$ pour tout objet $U$ de
$\mathcal{C}$, d'où l'égalité cherchée $g^{-1}g_!\mathcal{F} = \mathcal{F}$.
\end{proof}

\begin{remark}
\label{sites-modules-remark-no-extension}
En général, le foncteur $g_!$ ne s'étend pas aux catégories de modules lorsque
$g$ est (une partie d')un morphisme de topos annelés. En effet, pour tout
homomorphisme d'anneaux $A \to B$, le foncteur
$M \mapsto B \otimes_A M$ admet un adjoint à droite (la restriction des
scalaires), mais n'admet pas en général d'adjoint à gauche, car l'existence de
celui-ci impliquerait que $A \to B$ est plat. Nous verrons à la section
\ref{sites-modules-section-localize} qu'il est possible de définir $j_!$ sur les faisceaux de
modules dans le cas d'une localisation de sites. Nous en discuterons avec une
plus grande généralité à la section \ref{sites-modules-section-lower-shriek-modules}.
\end{remark}

\begin{lemma}
\label{sites-modules-lemma-have-left-adjoint}
Soient $\mathcal{C}$ et $\mathcal{D}$ des sites. Soit
$g : \Sh(\mathcal{C}) \to \Sh(\mathcal{D})$ le morphisme de topos associé à un
foncteur continu et cocontinu $u : \mathcal{C} \to \mathcal{D}$.
\begin{enumerate}
\item Si $u$ admet un adjoint à gauche $w$, alors $g_!$ et $g_!^{\Sh}$
coïncident sur les faisceaux d'ensembles sous-jacents et $g_!$ est exact.
\item Si, de plus, $w$ est cocontinu, alors $g_! = h^{-1}$ et
$g^{-1} = h_*$, où $h : \Sh(\mathcal{D}) \to \Sh(\mathcal{C})$ est le
morphisme de topos associé à $w$.
\end{enumerate}
\end{lemma}

\begin{proof}
Ce lemme est l'analogue de Sites, Lemme
\ref{sites-lemma-have-left-adjoint}. D'après Sites, Lemme
\ref{sites-lemma-adjoint-functors}, les catégories $\mathcal{I}_V^u$ possèdent
un objet initial. L'ensemble sous-jacent à la colimite d'un système de groupes
abéliens sur $(\mathcal{I}_V^u)^{opp}$ est donc la colimite des ensembles
sous-jacents. Il s'ensuit que $g_!^{\Sh}$ et $g_!$ coïncident, compte tenu de
notre construction de $g_!$ dans la Définition \ref{sites-modules-definition-g-shriek}.
L'exactitude et (2) résultent immédiatement des assertions correspondantes de
Sites, Lemme \ref{sites-lemma-have-left-adjoint}.
\end{proof}





\section{Types globaux de modules}
\label{sites-modules-section-global}

\begin{definition}
\label{sites-modules-definition-global}
Soit $(\Sh(\mathcal{C}), \mathcal{O})$ un topos annelé. Soit $\mathcal{F}$ un
faisceau de $\mathcal{O}$-modules.
\begin{enumerate}
\item Nous disons que $\mathcal{F}$ est un {\it $\mathcal{O}$-module libre}
si $\mathcal{F}$ est isomorphe, comme $\mathcal{O}$-module, à un faisceau de la
forme
$\bigoplus_{i \in I} \mathcal{O}$.
\item Nous disons que $\mathcal{F}$ est {\it libre de rang fini} si
$\mathcal{F}$ est isomorphe, comme $\mathcal{O}$-module, à un faisceau de la forme
$\bigoplus_{i \in I} \mathcal{O}$ avec $I$ fini.
\item Nous disons que $\mathcal{F}$ est {\it engendré par ses sections
globales} s'il existe une surjection
$$
\bigoplus\nolimits_{i \in I} \mathcal{O} \longrightarrow \mathcal{F}
$$
depuis un $\mathcal{O}$-module libre vers $\mathcal{F}$.
\item Pour $r \geq 0$, nous disons que $\mathcal{F}$ est {\it engendré par
$r$ sections globales} s'il existe une surjection
$\mathcal{O}^{\oplus r} \to \mathcal{F}$.
\item Nous disons que $\mathcal{F}$ est {\it engendré par un nombre fini de
sections globales} s'il est engendré par $r$ sections globales pour un certain
$r \geq 0$.
\item Nous disons que $\mathcal{F}$ admet une {\it présentation globale} s'il
existe une suite exacte
$$
\bigoplus\nolimits_{j \in J} \mathcal{O} \longrightarrow
\bigoplus\nolimits_{i \in I} \mathcal{O} \longrightarrow
\mathcal{F} \longrightarrow 0
$$
de $\mathcal{O}$-modules.
\item Nous disons que $\mathcal{F}$ admet une {\it présentation globale finie}
s'il existe une suite exacte
$$
\bigoplus\nolimits_{j \in J} \mathcal{O} \longrightarrow
\bigoplus\nolimits_{i \in I} \mathcal{O} \longrightarrow
\mathcal{F} \longrightarrow 0
$$
de $\mathcal{O}$-modules avec $I$ et $J$ finis.
\end{enumerate}
\end{definition}

\noindent
Remarquons que, pour tout ensemble $I$, la somme directe
$\bigoplus_{i \in I} \mathcal{O}$ existe (Lemme
\ref{sites-modules-lemma-limits-colimits}) et est le faisceau associé au préfaisceau
$U \mapsto \bigoplus_{i \in I} \mathcal{O}(U)$. Ce module est appelé le
{\it $\mathcal{O}$-module libre sur l'ensemble $I$}.

\begin{lemma}
\label{sites-modules-lemma-global-pullback}
Soit
$(f, f^\sharp) :
(\Sh(\mathcal{C}), \mathcal{O}_\mathcal{C})
\to
(\Sh(\mathcal{D}), \mathcal{O}_\mathcal{D})$
un morphisme de topos annelés. Soit $\mathcal{F}$ un
$\mathcal{O}_\mathcal{D}$-module.
\begin{enumerate}
\item Si $\mathcal{F}$ est libre, alors $f^*\mathcal{F}$ est libre.
\item Si $\mathcal{F}$ est libre de rang fini, alors $f^*\mathcal{F}$ est
libre de rang fini.
\item Si $\mathcal{F}$ est engendré par ses sections globales, alors
$f^*\mathcal{F}$ l'est aussi.
\item Si $r \geq 0$ et si $\mathcal{F}$ est engendré par $r$ sections
globales, alors $f^*\mathcal{F}$ est engendré par $r$ sections globales.
\item Si $\mathcal{F}$ est engendré par un nombre fini de sections globales,
alors $f^*\mathcal{F}$ l'est aussi.
\item Si $\mathcal{F}$ admet une présentation globale, alors
$f^*\mathcal{F}$ en admet une.
\item Si $\mathcal{F}$ admet une présentation globale finie, alors
$f^*\mathcal{F}$ en admet une.
\end{enumerate}
\end{lemma}

\begin{proof}
Cela est vrai parce que $f^*$ commute aux colimites quelconques (Lemme
\ref{sites-modules-lemma-exactness-pushforward-pullback}) et que
$f^*\mathcal{O}_\mathcal{D} = \mathcal{O}_\mathcal{C}$.
\end{proof}






\section{Propriétés intrinsèques des modules}
\label{sites-modules-section-intrinsic}

\noindent
Soit $\mathcal{P}$ une propriété des faisceaux de modules sur les topos
annelés. Nous disons que $\mathcal{P}$ est une {\it propriété intrinsèque} si
$\mathcal{P}(\mathcal{F}) \Leftrightarrow \mathcal{P}(f^*\mathcal{F})$
chaque fois que $(f, f^\sharp) :
(\Sh(\mathcal{C}'), \mathcal{O}')
\to
(\Sh(\mathcal{C}), \mathcal{O})$
est une équivalence de topos annelés. Par exemple, la propriété d'être libre est
intrinsèque. En effet, le $\mathcal{O}$-module libre sur l'ensemble $I$ est
caractérisé par la propriété
$$
\Mor_{\textit{Mod}(\mathcal{O})}(
\bigoplus\nolimits_{i \in I} \mathcal{O},
\mathcal{F})
=
\prod\nolimits_{i \in I} \Mor_{\Sh(\mathcal{C})}(\{*\},
\mathcal{F})
$$
pour $\mathcal{F}$ variable dans $\textit{Mod}(\mathcal{O})$. On peut aussi
utiliser le Lemme \ref{sites-modules-lemma-global-pullback} pour voir que la liberté est
intrinsèque. En fait, chacune des propriétés définies dans la Définition
\ref{sites-modules-definition-global} est intrinsèque pour la même raison. Comment
définirons-nous d'autres propriétés intrinsèques des $\mathcal{O}$-modules ?

\medskip\noindent
La conséquence du Lemme
\ref{sites-modules-lemma-morphism-ringed-topoi-comes-from-morphism-ringed-sites} est la
suivante. Supposons que l'on veuille définir une propriété intrinsèque
$\mathcal{P}$ d'un $\mathcal{O}$-module sur un topos. On peut procéder comme
suit :
\begin{enumerate}
\item Pour tout site $\mathcal{C}$, tout faisceau d'anneaux $\mathcal{O}$ sur
$\mathcal{C}$ et tout $\mathcal{O}$-module $\mathcal{F}$, définir la propriété
correspondante $\mathcal{P}(\mathcal{C}, \mathcal{O}, \mathcal{F})$.
\item Pour toute paire de sites $\mathcal{C}$, $\mathcal{C}'$, tout foncteur
cocontinu spécial $u : \mathcal{C} \to \mathcal{C}'$, tout faisceau d'anneaux
$\mathcal{O}$ sur $\mathcal{C}$ et tout $\mathcal{O}$-module $\mathcal{F}$,
montrer que
$$
\mathcal{P}(\mathcal{C}, \mathcal{O}, \mathcal{F})
\Leftrightarrow
\mathcal{P}(\mathcal{C}', g_*\mathcal{O}, g_*\mathcal{F})
$$
où $g : \Sh(\mathcal{C}) \to \Sh(\mathcal{C}')$ est l'équivalence de topos
associée à $u$.
\end{enumerate}
Dans ce cas, étant donnés un topos annelé
$(\Sh(\mathcal{C}), \mathcal{O})$ et un faisceau de $\mathcal{O}$-modules
$\mathcal{F}$, nous disons simplement que $\mathcal{F}$ possède la
propriété $\mathcal{P}$ si
$\mathcal{P}(\mathcal{C}, \mathcal{O}, \mathcal{F})$ est vraie. Le Lemme
\ref{sites-modules-lemma-morphism-ringed-topoi-comes-from-morphism-ringed-sites}, combiné à
(2), garantit que cette définition est indépendante des choix.

\medskip\noindent
De plus, le même Lemme
\ref{sites-modules-lemma-morphism-ringed-topoi-comes-from-morphism-ringed-sites} garantit que
si, en outre,
\begin{enumerate}
\item[(3)] pour tout morphisme de sites annelés
$(f, f^\sharp) :
(\mathcal{C}, \mathcal{O}_\mathcal{C})
\to
(\mathcal{D}, \mathcal{O}_\mathcal{D})$
tel que $f$ soit donné par un foncteur
$u : \mathcal{D} \to \mathcal{C}$ satisfaisant les hypothèses de Sites,
Proposition \ref{sites-proposition-get-morphism}, et pour tout
$\mathcal{O}_\mathcal{D}$-module $\mathcal{G}$, on ait
$$
\mathcal{P}(\mathcal{D}, \mathcal{O}_\mathcal{D}, \mathcal{G})
\Rightarrow
\mathcal{P}(\mathcal{C}, \mathcal{O}_\mathcal{C}, f^*\mathcal{G})
$$
\end{enumerate}
alors $\mathcal{P}$ est préservée par image inverse des modules suivant les
morphismes arbitraires de topos annelés.

\medskip\noindent
Dans les sections suivantes, nous utiliserons cette méthode pour montrer que les
propriétés suivantes sont intrinsèques : être localement libre, être localement
engendré par des sections, être localement engendré par $r$ sections, être de
type fini, être de présentation finie, être quasi-cohérent, et être cohérent.

\medskip\noindent
Une méthode peut-être plus satisfaisante consisterait à donner une définition
intrinsèque de ces notions, plutôt qu'à suivre la procédure laborieuse esquissée
ici. Cependant, dans de nombreuses situations géométriques où nous voulons
appliquer ces définitions, un site annelé précis et un faisceau de modules précis
nous sont donnés ; il est alors utile de disposer d'une définition déjà adaptée
à ce langage.
\section{Localisation de sites annelés}
\label{sites-modules-section-localize}

\noindent
Soit $(\mathcal{C}, \mathcal{O})$ un site annelé.
Soit $U \in \Ob(\mathcal{C})$.
Nous exposons dans ce cadre les analogues des résultats de
Sites, section \ref{sites-section-localize}.

\medskip\noindent
Notons
$\mathcal{O}_U = j_U^{-1}\mathcal{O}$ la restriction de $\mathcal{O}$
au site $\mathcal{C}/U$. Elle est décrite par la règle simple
$\mathcal{O}_U(V/U) = \mathcal{O}(V)$. Avec cette notation,
le morphisme de localisation $j_U$ devient un morphisme de topos annelés
$$
(j_U, j_U^\sharp) :
(\Sh(\mathcal{C}/U), \mathcal{O}_U)
\longrightarrow
(\Sh(\mathcal{C}), \mathcal{O})
$$
à savoir que l'on prend pour
$j_U^\sharp : j_U^{-1}\mathcal{O} \to \mathcal{O}_U$
l'application identique.
On obtient en outre les descriptions suivantes de l'image directe
et de l'image inverse des modules.

\begin{definition}
\label{sites-modules-definition-localize-ringed-site}
Soit $(\mathcal{C}, \mathcal{O})$ un site annelé.
Soit $U \in \Ob(\mathcal{C})$.
\begin{enumerate}
\item Le site annelé $(\mathcal{C}/U, \mathcal{O}_U)$ est appelé la
{\it localisation du site annelé $(\mathcal{C}, \mathcal{O})$
en l'objet $U$}.
\item Le morphisme de topos annelés
$(j_U, j_U^\sharp) :
(\Sh(\mathcal{C}/U), \mathcal{O}_U)
\to
(\Sh(\mathcal{C}), \mathcal{O})$
est appelé le {\it morphisme de localisation}.
\item Le foncteur
$j_{U*} : \textit{Mod}(\mathcal{O}_U) \to \textit{Mod}(\mathcal{O})$
est appelé le {\it foncteur image directe}.
\item Pour un faisceau de $\mathcal{O}$-modules $\mathcal{F}$ sur $\mathcal{C}$,
le faisceau $j_U^*\mathcal{F}$ est appelé la
{\it restriction de $\mathcal{F}$ à $\mathcal{C}/U$}.
Nous le noterons parfois
$\mathcal{F}|_{\mathcal{C}/U}$, voire $\mathcal{F}|_U$.
Il est décrit par la règle simple
$j_U^*(\mathcal{F})(X/U) = \mathcal{F}(X)$.
\item L'adjoint à gauche
$j_{U!} : \textit{Mod}(\mathcal{O}_U) \to \textit{Mod}(\mathcal{O})$
de la restriction est appelé {\it prolongement par zéro}. Il existe et est
exact d'après les
Lemmes \ref{sites-modules-lemma-extension-by-zero} et
\ref{sites-modules-lemma-extension-by-zero-exact}.
\end{enumerate}
\end{definition}

\noindent
Comme dans le cas topologique, voir
Faisceaux, section \ref{sheaves-section-open-immersions},
le foncteur de prolongement par zéro $j_{U!}$ diffère du
prolongement par l'ensemble vide $j_{U!}$ défini sur les faisceaux
d'ensembles. Voici le lemme qui définit le prolongement par zéro.

\begin{lemma}
\label{sites-modules-lemma-extension-by-zero}
Soit $(\mathcal{C}, \mathcal{O})$ un site annelé.
Soit $U \in \Ob(\mathcal{C})$.
Le foncteur de restriction
$j_U^* : \textit{Mod}(\mathcal{O}) \to \textit{Mod}(\mathcal{O}_U)$
a un adjoint à gauche
$j_{U!} : \textit{Mod}(\mathcal{O}_U) \to \textit{Mod}(\mathcal{O})$.
Ainsi,
$$
\Mor_{\textit{Mod}(\mathcal{O}_U)}(\mathcal{G}, j_U^*\mathcal{F})
=
\Mor_{\textit{Mod}(\mathcal{O})}(j_{U!}\mathcal{G}, \mathcal{F})
$$
pour $\mathcal{F} \in \Ob(\textit{Mod}(\mathcal{O}))$
et $\mathcal{G} \in \Ob(\textit{Mod}(\mathcal{O}_U))$.
De plus, le prolongement par zéro $j_{U!}\mathcal{G}$ de $\mathcal{G}$
est le faisceau associé au préfaisceau
$$
V
\longmapsto
\bigoplus\nolimits_{\varphi \in \Mor_\mathcal{C}(V, U)}
\mathcal{G}(V \xrightarrow{\varphi} U)
$$
muni des applications de restriction évidentes et de la structure
évidente de $\mathcal{O}$-module.
\end{lemma}

\begin{proof}
La structure de $\mathcal{O}$-module sur le préfaisceau est définie
comme suit. Si $f \in \mathcal{O}(V)$ et
$s \in \mathcal{G}(V \xrightarrow{\varphi} U)$, on pose
$f \cdot s = fs$, où
$f \in \mathcal{O}_U(\varphi : V \to U) = \mathcal{O}(V)$
(car $\mathcal{O}_U$ est la restriction de $\mathcal{O}$ à
$\mathcal{C}/U$).

\medskip\noindent
De même, soit $\alpha : \mathcal{G} \to \mathcal{F}|_U$ un
morphisme de $\mathcal{O}_U$-modules. On peut alors définir une
application du préfaisceau du lemme dans $\mathcal{F}$ en envoyant
$$
\bigoplus\nolimits_{\varphi \in \Mor_\mathcal{C}(V, U)}
\mathcal{G}(V \xrightarrow{\varphi} U)
\longrightarrow
\mathcal{F}(V)
$$
selon la règle qui à $s \in \mathcal{G}(V \xrightarrow{\varphi} U)$
associe $\alpha(s) \in \mathcal{F}(V)$. Cette application est clairement
$\mathcal{O}$-linéaire et induit donc un morphisme de
$\mathcal{O}$-modules $\alpha' : j_{U!}\mathcal{G} \to \mathcal{F}$
par les propriétés de la faisceautisation des modules
(Lemme \ref{sites-modules-lemma-sheafification-presheaf-modules}).

\medskip\noindent
Réciproquement, soit $\beta : j_{U!}\mathcal{G} \to \mathcal{F}$
un morphisme de $\mathcal{O}$-modules.
Rappelons d'après Sites, section \ref{sites-section-localize},
qu'il existe un prolongement par l'ensemble vide
$j^{Sh}_{U!} : \Sh(\mathcal{C}/U) \to \Sh(\mathcal{C})$
sur les faisceaux d'ensembles, adjoint à gauche de $j_U^{-1}$.
De plus, $j^{Sh}_{U!}\mathcal{G}$ est le faisceau associé au préfaisceau
$$
V
\longmapsto
\coprod\nolimits_{\varphi \in \Mor_\mathcal{C}(V, U)}
\mathcal{G}(V \xrightarrow{\varphi} U)
$$
Il existe donc une application naturelle
$j^{Sh}_{U!}\mathcal{G} \to j_{U!}\mathcal{G}$ de faisceaux d'ensembles.
En précomposant $\beta$ avec cette application, on obtient une
application de faisceaux d'ensembles
$j^{Sh}_{U!}\mathcal{G} \to \mathcal{F}$ qui correspond par adjonction
à une application de faisceaux d'ensembles
$\beta' : \mathcal{G} \to \mathcal{F}|_U$.
Nous affirmons que $\beta'$ est $\mathcal{O}_U$-linéaire. Supposons en effet
que $\varphi : V \to U$ soit un objet de $\mathcal{C}/U$ et que
$s, s' \in \mathcal{G}(\varphi : V \to U)$, et
$f \in \mathcal{O}(V) = \mathcal{O}_U(\varphi : V \to U)$.
La discussion précédente montre alors que
$\beta'(s + s')$, resp.\ $\beta'(fs)$ dans
$\mathcal{F}|_U(\varphi : V \to U)$, correspond à
$\beta(s + s')$, resp.\ $\beta(fs)$ dans
$\mathcal{F}(V)$. On conclut puisque $\beta$ est un homomorphisme.

\medskip\noindent
Pour achever la preuve du lemme, il reste à montrer que les constructions
$\alpha \mapsto \alpha'$ et $\beta \mapsto \beta'$ sont inverses l'une de
l'autre. Nous omettons les vérifications.
\end{proof}

\noindent
Notons que, dans la situation de la
Définition \ref{sites-modules-definition-localize-ringed-site}, on a
\begin{equation}
\label{sites-modules-equation-map-lower-shriek-OU-into-module}
\Hom_\mathcal{O}(j_{U!}\mathcal{O}_U, \mathcal{F}) =
\Hom_{\mathcal{O}_U}(\mathcal{O}_U, j_U^*\mathcal{F}) =
\mathcal{F}(U)
\end{equation}
pour tout $\mathcal{O}$-module $\mathcal{F}$. En effet, la première égalité
résulte de l'adjonction entre $j_{U!}$ et $j_U^*$, et la seconde du fait que
$\Hom_{\mathcal{O}_U}(\mathcal{O}_U, j_U^*\mathcal{F}) =
j_U^*\mathcal{F}(U/U) = \mathcal{F}|_U(U/U) = \mathcal{F}(U)$.

\begin{lemma}
\label{sites-modules-lemma-extension-by-zero-exact}
Soit $(\mathcal{C}, \mathcal{O})$ un site annelé.
Soit $U \in \Ob(\mathcal{C})$.
Le foncteur
$j_{U!} : \textit{Mod}(\mathcal{O}_U) \to \textit{Mod}(\mathcal{O})$
est exact.
\end{lemma}

\begin{proof}
Puisque $j_{U!}$ est adjoint à gauche de $j_U^*$, il est exact à droite
(voir
Catégories, Lemme \ref{categories-lemma-exact-adjoint}
et
Homology, section \ref{homology-section-functors}).
Il suffit donc de montrer que, si $\mathcal{G}_1 \to \mathcal{G}_2$
est un morphisme injectif de $\mathcal{O}_U$-modules, alors
$j_{U!}\mathcal{G}_1 \to j_{U!}\mathcal{G}_2$ est injectif.
Le morphisme sur les sections des préfaisceaux au-dessus d'un objet $V$
(comme dans le Lemme \ref{sites-modules-lemma-extension-by-zero}) est
$$
\bigoplus\nolimits_{\varphi \in \Mor_\mathcal{C}(V, U)}
\mathcal{G}_1(V \xrightarrow{\varphi} U)
\longrightarrow
\bigoplus\nolimits_{\varphi \in \Mor_\mathcal{C}(V, U)}
\mathcal{G}_2(V \xrightarrow{\varphi} U)
$$
qui est injectif par hypothèse. La faisceautisation étant exacte d'après le
Lemme \ref{sites-modules-lemma-sheafification-exact},
on en conclut que $j_{U!}\mathcal{G}_1 \to j_{U!}\mathcal{G}_2$ est injectif.
\end{proof}

\begin{lemma}
\label{sites-modules-lemma-j-shriek-reflects-exactness}
Soit $(\mathcal{C}, \mathcal{O})$ un site annelé.
Soit $U \in \Ob(\mathcal{C})$. Un complexe de $\mathcal{O}_U$-modules
$\mathcal{G}_1 \to \mathcal{G}_2 \to \mathcal{G}_3$ est exact
si et seulement si
$j_{U!}\mathcal{G}_1 \to j_{U!}\mathcal{G}_2 \to j_{U!}\mathcal{G}_3$
est exact comme suite de $\mathcal{O}$-modules.
\end{lemma}

\begin{proof}
Nous savons déjà que $j_{U!}$ est exact, voir le
Lemme \ref{sites-modules-lemma-extension-by-zero-exact}.
Il suffit donc de montrer que
$j_{U!} :  \textit{Mod}(\mathcal{O}_U) \to \textit{Mod}(\mathcal{O})$
reflète les injections et les surjections.

\medskip\noindent
Pour tout $\mathcal{G}$ dans $\textit{Mod}(\mathcal{O}_U)$,
on dispose de l'unité $\mathcal{G} \to j_U^*j_{U!}\mathcal{G}$
de l'adjonction. Nous affirmons que cette application est une injection
de faisceaux. En examinant la construction du
Lemme \ref{sites-modules-lemma-extension-by-zero}, on voit en effet qu'elle est la
faisceautisation de la règle qui associe à l'objet $V/U$ de
$\mathcal{C}/U$ l'application injective
$$
\mathcal{G}(V/U) \longrightarrow
\bigoplus\nolimits_{\varphi \in \Mor_\mathcal{C}(V, U)}
\mathcal{G}(V \xrightarrow{\varphi} U)
$$
donnée par l'inclusion du facteur correspondant au morphisme structural
$V \to U$. L'affirmation résulte alors de l'exactitude de la faisceautisation.
Nous omettons quelques détails.

\medskip\noindent
Si $\mathcal{G} \to \mathcal{G}'$ est un morphisme de
$\mathcal{O}_U$-modules tel que
$j_{U!}\mathcal{G} \to j_{U!}\mathcal{G}'$ soit injectif, alors
$j_U^*j_{U!}\mathcal{G} \to j_U^*j_{U!}\mathcal{G}'$ est injectif
(la restriction est exacte), donc
$\mathcal{G} \to j_U^*j_{U!}\mathcal{G}'$ est injectif, et par suite
$\mathcal{G} \to \mathcal{G}'$ est injectif.
Ainsi $j_{U!}$ reflète les injections.

\medskip\noindent
Soit $a : \mathcal{G} \to \mathcal{G}'$ un morphisme de
$\mathcal{O}_U$-modules tel que
$j_{U!}\mathcal{G} \to j_{U!}\mathcal{G}'$ soit surjectif.
Soit $\mathcal{H}$ le conoyau de $a$.
Alors $j_{U!}\mathcal{H} = 0$, puisque $j_{U!}$ est exact.
D'après ce qui précède, l'application
$\mathcal{H} \to j^*_U j_{U!}\mathcal{H}$ est injective.
Ainsi $\mathcal{H} = 0$, comme voulu.
\end{proof}

\begin{lemma}
\label{sites-modules-lemma-relocalize}
Soit $(\mathcal{C}, \mathcal{O})$ un site annelé.
Soit $f : V \to U$ un morphisme de $\mathcal{C}$.
Il existe alors un diagramme commutatif
$$
\xymatrix{
(\Sh(\mathcal{C}/V), \mathcal{O}_V)
\ar[rd]_{(j_V, j_V^\sharp)} \ar[rr]_{(j, j^\sharp)} & &
(\Sh(\mathcal{C}/U), \mathcal{O}_U)
\ar[ld]^{(j_U, j_U^\sharp)} \\
& (\Sh(\mathcal{C}), \mathcal{O}) &
}
$$
dans la catégorie des topos annelés. Ici $(j, j^\sharp)$ est le morphisme
de localisation associé à l'objet $V/U$ du site annelé
$(\mathcal{C}/V, \mathcal{O}_V)$.
\end{lemma}

\begin{proof}
Le seul point à vérifier est que
$j_V^\sharp = j^\sharp \circ j^{-1}(j_U^\sharp)$,
car tout le reste résulte directement de
Sites, Lemme \ref{sites-lemma-relocalize}, et
Sites, Équation (\ref{sites-equation-relocalize}).
Nous omettons la vérification de cette égalité.
\end{proof}

\begin{remark}
\label{sites-modules-remark-localize-shriek-equal}
Dans la situation du Lemme \ref{sites-modules-lemma-extension-by-zero},
le diagramme
$$
\xymatrix{
\textit{Mod}(\mathcal{O}_U) \ar[r]_{j_{U!}} \ar[d]_{forget} &
\textit{Mod}(\mathcal{O}_\mathcal{C}) \ar[d]^{forget} \\
\textit{Ab}(\mathcal{C}/U) \ar[r]^{j^{Ab}_{U!}} &
\textit{Ab}(\mathcal{C})
}
$$
est commutatif. Cela résulte clairement de la description explicite du
foncteur $j_{U!}$ donnée dans le lemme.
\end{remark}

\begin{remark}
\label{sites-modules-remark-localize-presheaves}
Localisation et préfaisceaux de modules ; voir
Sites, Remarque \ref{sites-remark-localize-presheaves}.
Soit $\mathcal{C}$ une catégorie.
Soit $\mathcal{O}$ un préfaisceau d'anneaux.
Soit $U$ un objet de $\mathcal{C}$.
À proprement parler, les foncteurs $j_U^*$, $j_{U*}$ et $j_{U!}$
n'ont pas été définis pour les préfaisceaux de $\mathcal{O}$-modules.
On peut cependant considérer un préfaisceau comme un faisceau pour la
topologie chaotique sur $\mathcal{C}$ (voir
Sites, Exemples \ref{sites-example-indiscrete}).
On obtient donc aussi un foncteur
$$
j_U^* :
\textit{PMod}(\mathcal{O})
\longrightarrow
\textit{PMod}(\mathcal{O}_U)
$$
et des foncteurs
$$
j_{U*}, j_{U!} :
\textit{PMod}(\mathcal{O}_U)
\longrightarrow
\textit{PMod}(\mathcal{O})
$$
qui sont respectivement adjoint à droite et adjoint à gauche de $j_U^*$.
En examinant la preuve du Lemme \ref{sites-modules-lemma-extension-by-zero}, on voit que
$j_{U!}\mathcal{G}$ est le préfaisceau
$$
V \longmapsto
\bigoplus\nolimits_{\varphi \in \Mor_\mathcal{C}(V, U)}
\mathcal{G}(V \xrightarrow{\varphi} U)
$$
En outre, le foncteur $j_{U!}$ est exact (d'après le
Lemme \ref{sites-modules-lemma-extension-by-zero-exact}, appliqué aux topologies
discrètes). Si de plus $\mathcal{C}$ est un site et $\mathcal{O}$ un faisceau
d'anneaux, alors le diagramme
$$
\xymatrix{
\textit{Mod}(\mathcal{O}_U) \ar[r]_{j_{U!}} \ar[d]_{forget} &
\textit{Mod}(\mathcal{O}) \\
\textit{PMod}(\mathcal{O}_U) \ar[r]^{j_{U!}} &
\textit{PMod}(\mathcal{O}) \ar[u]_{(\ )^\#}
}
$$
est commutatif.
\end{remark}

\begin{lemma}
\label{sites-modules-lemma-restrict-back}
Soit $\mathcal{C}$ un site et $U \in \Ob(\mathcal{C})$.
Supposons que tout $X$ de $\mathcal{C}$ admette au plus un morphisme vers $U$.
Soit $\mathcal{F}$ un faisceau abélien sur $\mathcal{C}/U$.
Les applications canoniques $\mathcal{F} \to j_U^{-1}j_{U!}\mathcal{F}$
et $j_U^{-1}j_{U*}\mathcal{F} \to \mathcal{F}$ sont des isomorphismes.
\end{lemma}

\begin{proof}
C'est un cas particulier du Lemme \ref{sites-modules-lemma-back-and-forth}, car l'hypothèse
sur $U$ équivaut à la pleine fidélité du foncteur de localisation
$\mathcal{C}/U \to \mathcal{C}$.
\end{proof}









\section{Localisation des morphismes de sites annelés}
\label{sites-modules-section-localize-morphisms}

\noindent
Cette section est l'analogue de
Sites, section \ref{sites-section-localize-morphisms}.

\begin{lemma}
\label{sites-modules-lemma-localize-morphism-ringed-sites}
Soit
$(f, f^\sharp) :
(\mathcal{C}, \mathcal{O})
\longrightarrow
(\mathcal{D}, \mathcal{O}')$
un morphisme de sites annelés, où $f$ est donné par le foncteur continu
$u : \mathcal{D} \to \mathcal{C}$.
Soit $V$ un objet de $\mathcal{D}$ et posons $U = u(V)$.
Il existe alors une application canonique de faisceaux d'anneaux
$(f')^\sharp$ telle que le diagramme du
Lemme \ref{sites-lemma-localize-morphism} de Sites
devienne un diagramme commutatif de topos annelés
$$
\xymatrix{
(\Sh(\mathcal{C}/U), \mathcal{O}_U)
\ar[rr]_{(j_U, j_U^\sharp)} \ar[d]_{(f', (f')^\sharp)} & &
(\Sh(\mathcal{C}), \mathcal{O})
\ar[d]^{(f, f^\sharp)} \\
(\Sh(\mathcal{D}/V), \mathcal{O}'_V)
\ar[rr]^{(j_V, j_V^\sharp)} & &
(\Sh(\mathcal{D}), \mathcal{O}').
}
$$
De plus, dans cette situation, on a $f'_*j_U^{-1} = j_V^{-1}f_*$
et $f'_*j_U^* = j_V^*f_*$.
\end{lemma}

\begin{proof}
Il suffit de prendre pour $(f')^\sharp$ l'application
$$
(f')^{-1}\mathcal{O}'_V =
(f')^{-1}j_V^{-1}\mathcal{O}' =
j_U^{-1}f^{-1}\mathcal{O}' \xrightarrow{j_U^{-1}f^\sharp}
j_U^{-1}\mathcal{O} = \mathcal{O}_U
$$
et tout le reste résulte de
Sites, Lemme \ref{sites-lemma-localize-morphism}.
(Notons que $j^{-1} = j^*$ sur les faisceaux de modules lorsque $j$ est un
morphisme de localisation ; la première égalité de foncteurs implique donc
la seconde.)
\end{proof}

\begin{lemma}
\label{sites-modules-lemma-relocalize-morphism-ringed-sites}
Soit
$(f, f^\sharp) :
(\mathcal{C}, \mathcal{O})
\longrightarrow
(\mathcal{D}, \mathcal{O}')$
un morphisme de sites annelés, où $f$ est donné par le foncteur continu
$u : \mathcal{D} \to \mathcal{C}$.
Soient $V \in \Ob(\mathcal{D})$, $U \in \Ob(\mathcal{C})$ et
$c : U \to u(V)$ un morphisme de $\mathcal{C}$.
Il existe un diagramme commutatif de topos annelés
$$
\xymatrix{
(\Sh(\mathcal{C}/U), \mathcal{O}_U)
\ar[rr]_{(j_U, j_U^\sharp)} \ar[d]_{(f_c, f_c^\sharp)} & &
(\Sh(\mathcal{C}), \mathcal{O}) \ar[d]^{(f, f^\sharp)} \\
(\Sh(\mathcal{D}/V), \mathcal{O}'_V)
\ar[rr]^{(j_V, j_V^\sharp)} & &
(\Sh(\mathcal{D}), \mathcal{O}').
}
$$
Le morphisme $(f_c, f_c^\sharp)$ est égal au composé du morphisme
$$
(f', (f')^\sharp) :
(\Sh(\mathcal{C}/u(V)), \mathcal{O}_{u(V)})
\longrightarrow
(\Sh(\mathcal{D}/V), \mathcal{O}'_V)
$$
du Lemme \ref{sites-modules-lemma-localize-morphism-ringed-sites} et du morphisme
$$
(j, j^\sharp) :
(\Sh(\mathcal{C}/U), \mathcal{O}_U)
\to
(\Sh(\mathcal{C}/u(V)), \mathcal{O}_{u(V)})
$$
du Lemme \ref{sites-modules-lemma-relocalize}.
Étant donnés des morphismes $b : V' \to V$, $a : U' \to U$ et
$c' : U' \to u(V')$ tels que
$$
\xymatrix{
U' \ar[r]_-{c'} \ar[d]_a & u(V') \ar[d]^{u(b)} \\
U \ar[r]^-c & u(V)
}
$$
soit commutatif, le diagramme suivant de topos annelés
$$
\xymatrix{
(\Sh(\mathcal{C}/U'), \mathcal{O}_{U'})
\ar[rr]_{(j_{U'/U}, j_{U'/U}^\sharp)} \ar[d]_{(f_{c'}, f_{c'}^\sharp)} & &
(\Sh(\mathcal{C}/U), \mathcal{O}_U)
\ar[d]^{(f_c, f_c^\sharp)} \\
(\Sh(\mathcal{D}/V'), \mathcal{O}'_{V'})
\ar[rr]^{(j_{V'/V}, j_{V'/V}^\sharp)} & &
(\Sh(\mathcal{D}/V), \mathcal{O}'_{V'})
}
$$
est commutatif.
\end{lemma}

\begin{proof}
Au niveau des morphismes de topos, c'est le
Lemme \ref{sites-lemma-relocalize-morphism} de Sites.
Pour vérifier que les diagrammes commutent comme diagrammes de morphismes de
topos annelés, on utilise les Lemmes \ref{sites-modules-lemma-relocalize} et
\ref{sites-modules-lemma-localize-morphism-ringed-sites}, exactement comme dans la preuve du
Lemme \ref{sites-lemma-relocalize-morphism} de Sites.
\end{proof}
















\section{Localisation de topos annelés}
\label{sites-modules-section-localize-ringed-topoi}

\noindent
Cette section est l'analogue de
Sites, section \ref{sites-section-localize-topoi}
dans le cadre des topos annelés.

\begin{lemma}
\label{sites-modules-lemma-localize-ringed-topos}
Soit $(\Sh(\mathcal{C}), \mathcal{O})$ un topos annelé.
Soit $\mathcal{F} \in \Sh(\mathcal{C})$ un faisceau.
Pour tout faisceau $\mathcal{H}$ sur $\mathcal{C}$, notons
$\mathcal{H}_\mathcal{F}$ le faisceau $\mathcal{H} \times \mathcal{F}$,
considéré comme objet de la catégorie
$\Sh(\mathcal{C})/\mathcal{F}$. Le couple
$(\Sh(\mathcal{C})/\mathcal{F}, \mathcal{O}_\mathcal{F})$
est un topos annelé et il existe un morphisme canonique de topos annelés
$$
(j_\mathcal{F}, j_\mathcal{F}^\sharp) :
(\Sh(\mathcal{C})/\mathcal{F}, \mathcal{O}_\mathcal{F})
\longrightarrow
(\Sh(\mathcal{C}), \mathcal{O})
$$
qui est une localisation au sens de la
section \ref{sites-modules-section-localize} et tel que
\begin{enumerate}
\item le foncteur $j_\mathcal{F}^{-1}$ est le foncteur
$\mathcal{H} \mapsto \mathcal{H}_\mathcal{F}$,
\item le foncteur $j_\mathcal{F}^*$ est le foncteur
$\mathcal{H} \mapsto \mathcal{H}_\mathcal{F}$,
\item le foncteur $j_{\mathcal{F}!}$ sur les faisceaux d'ensembles est le
foncteur d'oubli $\mathcal{G}/\mathcal{F} \mapsto \mathcal{G}$,
\item le foncteur $j_{\mathcal{F}!}$ sur les faisceaux de modules associe au
$\mathcal{O}_\mathcal{F}$-module
$\varphi : \mathcal{G} \to \mathcal{F}$ le $\mathcal{O}$-module qui est la
faisceautisation du préfaisceau
$$
V \longmapsto
\bigoplus\nolimits_{s \in \mathcal{F}(V)}
\{\sigma \in \mathcal{G}(V) \mid \varphi(\sigma) = s \}
$$
pour $V \in \Ob(\mathcal{C})$.
\end{enumerate}
\end{lemma}

\begin{proof}
D'après le Lemme \ref{sites-lemma-localize-topos} de Sites,
$\Sh(\mathcal{C})/\mathcal{F}$ est un topos et (1) et (3) sont vrais.
En particulier, cela montre que
$j_\mathcal{F}^{-1}\mathcal{O} = \mathcal{O}_\mathcal{F}$ et que
$\mathcal{O}_\mathcal{F}$ est un faisceau d'anneaux.
On peut donc choisir pour $j_\mathcal{F}^\sharp$ l'application identique ;
en particulier, (2) est vrai. De plus, la preuve du
Lemme \ref{sites-lemma-localize-topos} de Sites montre que l'on peut supposer
que $\mathcal{C}$ est un site admettant toutes les limites finies et muni
d'une topologie sous-canonique, et que $\mathcal{F} = h_U$ pour un objet
$U$ de $\mathcal{C}$. Alors (4) résulte de la description de
$j_{\mathcal{F}!}$ au Lemme \ref{sites-modules-lemma-extension-by-zero}.
On pourrait aussi montrer directement que le foncteur décrit en (4) est
adjoint à gauche de $j_\mathcal{F}^*$.
\end{proof}

\begin{definition}
\label{sites-modules-definition-localize-ringed-topos}
Soit $(\Sh(\mathcal{C}), \mathcal{O})$ un topos annelé.
Soit $\mathcal{F} \in \Sh(\mathcal{C})$.
\begin{enumerate}
\item Le topos annelé
$(\Sh(\mathcal{C})/\mathcal{F}, \mathcal{O}_\mathcal{F})$
est appelé la
{\it localisation du topos annelé
$(\Sh(\mathcal{C}), \mathcal{O})$ en $\mathcal{F}$}.
\item Le morphisme de topos annelés
$(j_\mathcal{F}, j_\mathcal{F}^\sharp) :
(\Sh(\mathcal{C})/\mathcal{F}, \mathcal{O}_\mathcal{F})
\to
(\Sh(\mathcal{C}), \mathcal{O})$ of
Lemme \ref{sites-modules-lemma-localize-ringed-topos}
est appelé le {\it morphisme de localisation}.
\end{enumerate}
\end{definition}

\noindent
Conformément à l'usage établi dans le chapitre sur les sites, nous
vérifions, chaque fois que cela a un sens, que les constructions de
localisation des topos sont compatibles avec celles des sites.

\begin{lemma}
\label{sites-modules-lemma-localize-compare}
Soient
$(\Sh(\mathcal{C}), \mathcal{O})$ et
$\mathcal{F} \in \Sh(\mathcal{C})$ comme dans le
Lemme \ref{sites-modules-lemma-localize-ringed-topos}.
Si $\mathcal{F} = h_U^\#$ pour un objet $U$ de $\mathcal{C}$, alors, via
l'identification
$\Sh(\mathcal{C}/U) = \Sh(\mathcal{C})/h_U^\#$ of
Sites, Lemme \ref{sites-lemma-essential-image-j-shriek}
on a
\begin{enumerate}
\item canoniquement $\mathcal{O}_U = \mathcal{O}_\mathcal{F}$, et
\item sous ces identifications,
$(j_\mathcal{F}, j_\mathcal{F}^\sharp) = (j_U, j_U^\sharp)$.
\end{enumerate}
\end{lemma}

\begin{proof}
L'assertion concernant les topos sous-jacents est le
Lemme \ref{sites-lemma-localize-compare} de Sites.
Notons que $\mathcal{O}_U$ est la restriction de $\mathcal{O}$ qui, d'après
Sites, Lemme \ref{sites-lemma-compute-j-shriek-restrict}
correspond à $\mathcal{O} \times h_U^\#$ sous l'équivalence fournie par
Sites, Lemme \ref{sites-lemma-essential-image-j-shriek}.
La définition de $\mathcal{O}_\mathcal{F}$ donne (1).
Il reste à prouver que $j_\mathcal{F}^\sharp = j_U^\sharp$ sous cette
identification. Nous omettons la vérification.
\end{proof}

\noindent
La localisation est fonctorielle de deux manières : on peut
« relocaliser » une localisation (voir le
Lemme \ref{sites-modules-lemma-relocalize-ringed-topos}) ou, étant donné un morphisme de
topos annelés, localiser en haut en l'image inverse d'un faisceau d'en bas et
obtenir un diagramme commutatif de topos annelés (voir le
Lemme \ref{sites-modules-lemma-localize-morphism-ringed-topoi}).

\begin{lemma}
\label{sites-modules-lemma-relocalize-ringed-topos}
Soit $(\Sh(\mathcal{C}), \mathcal{O})$ un topos annelé.
Si $s : \mathcal{G} \to \mathcal{F}$ est un morphisme de faisceaux sur
$\mathcal{C}$, il existe un diagramme commutatif naturel de morphismes de
topos annelés
$$
\xymatrix{
(\Sh(\mathcal{C})/\mathcal{G}, \mathcal{O}_\mathcal{G})
\ar[rd]_{(j_\mathcal{G}, j_\mathcal{G}^\sharp)} \ar[rr]_{(j, j^\sharp)} & &
(\Sh(\mathcal{C})/\mathcal{F}, \mathcal{O}_\mathcal{F})
\ar[ld]^{(j_\mathcal{F}, j_\mathcal{F}^\sharp)} \\
& (\Sh(\mathcal{C}), \mathcal{O}) &
}
$$
où $(j, j^\sharp)$ est le morphisme de localisation du topos annelé
$(\Sh(\mathcal{C})/\mathcal{F}, \mathcal{O}_\mathcal{F})$
en l'objet $\mathcal{G}/\mathcal{F}$.
\end{lemma}

\begin{proof}
Toutes les assertions résultent de
Sites, Lemme \ref{sites-lemma-relocalize-topos}
excepté l'assertion
$j_\mathcal{G}^\sharp = j^\sharp \circ j^{-1}(j_\mathcal{F}^\sharp)$.
Nous omettons la vérification.
\end{proof}

\begin{lemma}
\label{sites-modules-lemma-relocalize-compare}
Soient $(\Sh(\mathcal{C}), \mathcal{O})$ et
$s : \mathcal{G} \to \mathcal{F}$ comme dans le
Lemme \ref{sites-modules-lemma-relocalize-ringed-topos}.
S'il existe un morphisme $f : V \to U$ de $\mathcal{C}$ tel que
$\mathcal{G} = h_V^\#$, $\mathcal{F} = h_U^\#$ et que $s$ soit induit par
$f$, alors les diagrammes du
Lemme \ref{sites-modules-lemma-relocalize}
et du
Lemme \ref{sites-modules-lemma-relocalize-ringed-topos}
coïncident via les identifications
$(j_\mathcal{F}, j_\mathcal{F}^\sharp) = (j_U, j_U^\sharp)$
et
$(j_\mathcal{G}, j_\mathcal{G}^\sharp) = (j_V, j_V^\sharp)$
du Lemme \ref{sites-modules-lemma-localize-compare}.
\end{lemma}

\begin{proof}
Toutes les assertions résultent de
Sites, Lemme \ref{sites-lemma-relocalize-compare}
excepté l'assertion selon laquelle les deux applications $j^\sharp$
coïncident. Elle est vraie puisque, dans les deux cas, $j^\sharp$ est
simplement l'identité. Nous omettons quelques détails.
\end{proof}










\section{Localisation des morphismes de topos annelés}
\label{sites-modules-section-localize-morphisms-ringed-topoi}

\noindent
Cette section est l'analogue de
Sites, section \ref{sites-section-localize-morphisms-topoi}.

\begin{lemma}
\label{sites-modules-lemma-localize-morphism-ringed-topoi}
Soit
$$
f :
(\Sh(\mathcal{C}), \mathcal{O})
\longrightarrow
(\Sh(\mathcal{D}), \mathcal{O}')
$$
soit un morphisme de topos annelés. Soit $\mathcal{G}$ un faisceau sur
$\mathcal{D}$. Posons $\mathcal{F} = f^{-1}\mathcal{G}$.
Il existe alors un diagramme commutatif de topos annelés
$$
\xymatrix{
(\Sh(\mathcal{C})/\mathcal{F}, \mathcal{O}_\mathcal{F})
\ar[rr]_{(j_\mathcal{F}, j_\mathcal{F}^\sharp)}
\ar[d]_{(f', (f')^\sharp)} & &
(\Sh(\mathcal{C}), \mathcal{O}) \ar[d]^{(f, f^\sharp)} \\
(\Sh(\mathcal{D})/\mathcal{G}, \mathcal{O}'_\mathcal{G})
\ar[rr]^{(j_\mathcal{G}, j_\mathcal{G}^\sharp)} & &
(\Sh(\mathcal{D}), \mathcal{O}')
}
$$
On a $f'_*j_\mathcal{F}^{-1} = j_\mathcal{G}^{-1}f_*$ et
$f'_*j_\mathcal{F}^* = j_\mathcal{G}^*f_*$. De plus, le morphisme $f'$ est
caractérisé par la règle
$$
(f')^{-1}(\mathcal{H} \xrightarrow{\varphi} \mathcal{G})
=
(f^{-1}\mathcal{H} \xrightarrow{f^{-1}\varphi} \mathcal{F}).
$$
\end{lemma}

\begin{proof}
Le Lemme \ref{sites-lemma-localize-morphism-topoi} de Sites donne le
diagramme des topos sous-jacents, l'égalité
$f'_*j_\mathcal{F}^{-1} = j_\mathcal{G}^{-1}f_*$ et la description de
$(f')^{-1}$. Pour définir $(f')^\sharp$, on utilise l'application
$$
(f')^\sharp :
\mathcal{O}'_\mathcal{G} =
j_\mathcal{G}^{-1} \mathcal{O}'
\xrightarrow{j_\mathcal{G}^{-1}f^\sharp}
j_\mathcal{G}^{-1} f_*\mathcal{O} =
f'_* j_\mathcal{F}^{-1}\mathcal{O} =
f'_* \mathcal{O}_\mathcal{F}
$$
ou, de manière équivalente, l'application
$$
(f')^\sharp :
(f')^{-1}\mathcal{O}'_\mathcal{G} =
(f')^{-1}j_\mathcal{G}^{-1} \mathcal{O}' =
j_\mathcal{F}^{-1}f^{-1}\mathcal{O}'
\xrightarrow{j_\mathcal{F}^{-1}f^\sharp}
j_\mathcal{F}^{-1} \mathcal{O} =
\mathcal{O}_\mathcal{F}.
$$
Nous omettons la vérification que ces deux applications sont bien adjointes.
La seconde construction de $(f')^\sharp$ montre que le diagramme commute
dans la $2$-catégorie des topos annelés (car les applications
$j_\mathcal{F}^\sharp$ et $j_\mathcal{G}^\sharp$ sont des identités).
Enfin, l'égalité $f'_*j_\mathcal{F}^* = j_\mathcal{G}^*f_*$ résulte de
$f'_*j_\mathcal{F}^{-1} = j_\mathcal{G}^{-1}f_*$
et du fait que les images inverses de faisceaux de modules et de faisceaux
d'ensembles coïncident, voir le Lemme \ref{sites-modules-lemma-localize-ringed-topos}.
\end{proof}

\begin{lemma}
\label{sites-modules-lemma-localize-morphism-compare}
Soit
$$
f :
(\Sh(\mathcal{C}), \mathcal{O})
\longrightarrow
(\Sh(\mathcal{D}), \mathcal{O}')
$$
soit un morphisme de topos annelés.
Soit $\mathcal{G}$ un faisceau sur $\mathcal{D}$.
Posons $\mathcal{F} = f^{-1}\mathcal{G}$.
Si $f$ est donné par un foncteur continu $u : \mathcal{D} \to \mathcal{C}$
et $\mathcal{G} = h_V^\#$, alors les diagrammes commutatifs du
Lemme \ref{sites-modules-lemma-localize-morphism-ringed-sites}
et du
Lemme \ref{sites-modules-lemma-localize-morphism-ringed-topoi}
coïncident via les identifications du
Lemme \ref{sites-modules-lemma-localize-compare}.
\end{lemma}

\begin{proof}
Au niveau des morphismes de topos, c'est le
Lemme \ref{sites-lemma-localize-morphism-compare} de Sites.
Cela vaut aussi au niveau des morphismes de topos annelés puisque les
formules qui définissent $(f')^\sharp$ dans les preuves du
Lemme \ref{sites-modules-lemma-localize-morphism-ringed-sites}
et du
Lemme \ref{sites-modules-lemma-localize-morphism-ringed-topoi}
coïncident.
\end{proof}

\begin{lemma}
\label{sites-modules-lemma-relocalize-morphism-ringed-topoi}
Soit
$(f, f^\sharp) :
(\Sh(\mathcal{C}), \mathcal{O})
\to
(\Sh(\mathcal{D}), \mathcal{O}')$
soit un morphisme de topos annelés.
Soient $\mathcal{G}$ un faisceau sur $\mathcal{D}$,
$\mathcal{F}$ un faisceau sur $\mathcal{C}$ et
$s : \mathcal{F} \to f^{-1}\mathcal{G}$ un morphisme de faisceaux.
Il existe un diagramme commutatif de topos annelés
$$
\xymatrix{
(\Sh(\mathcal{C})/\mathcal{F}, \mathcal{O}_\mathcal{F})
\ar[rr]_{(j_\mathcal{F}, j_\mathcal{F}^\sharp)}
\ar[d]_{(f_c, f_c^\sharp)} & &
(\Sh(\mathcal{C}), \mathcal{O})
\ar[d]^{(f, f^\sharp)} \\
(\Sh(\mathcal{D})/\mathcal{G}, \mathcal{O}'_\mathcal{G})
\ar[rr]^{(j_\mathcal{G}, j_\mathcal{G}^\sharp)} & &
(\Sh(\mathcal{D}), \mathcal{O}').
}
$$
Le morphisme $(f_s, f_s^\sharp)$ est égal au composé du morphisme
$$
(f', (f')^\sharp) :
(\Sh(\mathcal{C})/f^{-1}\mathcal{G}, \mathcal{O}_{f^{-1}\mathcal{G}})
\longrightarrow
(\Sh(\mathcal{D})/{\mathcal{G}}, \mathcal{O}'_\mathcal{G})
$$
du Lemme \ref{sites-modules-lemma-localize-morphism-ringed-topoi} et du morphisme
$$
(j, j^\sharp) :
(\Sh(\mathcal{C})/\mathcal{F}, \mathcal{O}_\mathcal{F})
\to
(\Sh(\mathcal{C})/f^{-1}\mathcal{G}, \mathcal{O}_{f^{-1}\mathcal{G}})
$$
du Lemme \ref{sites-modules-lemma-relocalize-ringed-topos}.
Étant donnés des morphismes $b : \mathcal{G}' \to \mathcal{G}$,
$a : \mathcal{F}' \to \mathcal{F}$ et
$s' : \mathcal{F}' \to f^{-1}\mathcal{G}'$ tels que
$$
\xymatrix{
\mathcal{F}' \ar[r]_-{s'} \ar[d]_a &
f^{-1}\mathcal{G}' \ar[d]^{f^{-1}b} \\
\mathcal{F} \ar[r]^-s &
f^{-1}\mathcal{G}
}
$$
soit commutatif, le diagramme suivant de topos annelés
$$
\xymatrix{
(\Sh(\mathcal{C})/\mathcal{F}', \mathcal{O}_{\mathcal{F}'})
\ar[rr]_{(j_{\mathcal{F}'/\mathcal{F}}, j_{\mathcal{F}'/\mathcal{F}}^\sharp)}
\ar[d]_{(f_{s'}, f_{s'}^\sharp)} & &
(\Sh(\mathcal{C})/\mathcal{F}, \mathcal{O}_\mathcal{F})
\ar[d]^{(f_s, f_s^\sharp)} \\
(\Sh(\mathcal{D})/\mathcal{G}', \mathcal{O}'_{\mathcal{G}'})
\ar[rr]^{(j_{\mathcal{G}'/\mathcal{G}}, j_{\mathcal{G}'/\mathcal{G}}^\sharp)}
& &
(\Sh(\mathcal{D})/\mathcal{G}, \mathcal{O}'_{\mathcal{G}'})
}
$$
est commutatif.
\end{lemma}

\begin{proof}
Au niveau des morphismes de topos, c'est le
Lemme \ref{sites-lemma-relocalize-morphism-topoi} de Sites.
Pour vérifier que les diagrammes commutent comme diagrammes de morphismes de
topos annelés, on utilise les diagrammes commutatifs des
Lemmes \ref{sites-modules-lemma-relocalize-ringed-topos} et
\ref{sites-modules-lemma-localize-morphism-ringed-topoi}.
\end{proof}

\begin{lemma}
\label{sites-modules-lemma-relocalize-morphism-compare}
Soient
$(f, f^\sharp) :
(\Sh(\mathcal{C}), \mathcal{O})
\to
(\Sh(\mathcal{D}), \mathcal{O}')$,
$s : \mathcal{F} \to f^{-1}\mathcal{G}$ étant comme dans le
Lemme \ref{sites-modules-lemma-relocalize-morphism-ringed-topoi}.
Si $f$ est donné par un foncteur continu
$u : \mathcal{D} \to \mathcal{C}$
et si $\mathcal{G} = h_V^\#$, $\mathcal{F} = h_U^\#$ et $s$ provient d'un
morphisme $c : U \to u(V)$, alors les diagrammes commutatifs du
Lemme \ref{sites-modules-lemma-relocalize-morphism-ringed-sites}
et du
Lemme \ref{sites-modules-lemma-relocalize-morphism-ringed-topoi}
coïncident via les identifications du
Lemme \ref{sites-modules-lemma-localize-compare}.
\end{lemma}

\begin{proof}
Cela résulte formellement des
Lemmes \ref{sites-modules-lemma-relocalize-compare} et
\ref{sites-modules-lemma-localize-morphism-compare}.
\end{proof}

















\section{Propriétés locales des modules}
\label{sites-modules-section-local}

\noindent
Conformément à la stratégie générale expliquée à la
section \ref{sites-modules-section-intrinsic}, nous définissons d'abord les propriétés
locales des faisceaux de modules sur un site annelé, puis montrons aussitôt
qu'elles sont intrinsèques et ont donc un sens pour les faisceaux de modules
sur les topos annelés.

\begin{definition}
\label{sites-modules-definition-site-local}
Soit $(\mathcal{C}, \mathcal{O})$ un site annelé.
Soit $\mathcal{F}$ un faisceau de $\mathcal{O}$-modules.
Nous utiliserons librement les notions définies dans la
Définition \ref{sites-modules-definition-global}.
\begin{enumerate}
\item On dit que $\mathcal{F}$ est {\it localement libre} si, pour tout objet
$U$ de $\mathcal{C}$, il existe un recouvrement
$\{U_i \to U\}_{i \in I}$ de $\mathcal{C}$ tel que chaque restriction
$\mathcal{F}|_{\mathcal{C}/U_i}$ soit un
$\mathcal{O}_{U_i}$-module libre.
\item On dit que $\mathcal{F}$ est {\it localement libre de rang fini} si,
pour tout objet $U$ de $\mathcal{C}$, il existe un recouvrement
$\{U_i \to U\}_{i \in I}$ de $\mathcal{C}$ tel que chaque restriction
$\mathcal{F}|_{\mathcal{C}/U_i}$ soit un
$\mathcal{O}_{U_i}$-module libre de rang fini.
\item On dit que $\mathcal{F}$ est {\it localement engendré par des sections}
si, pour tout objet $U$ de $\mathcal{C}$, il existe un recouvrement
$\{U_i \to U\}_{i \in I}$ de $\mathcal{C}$ tel que chaque restriction
$\mathcal{F}|_{\mathcal{C}/U_i}$ soit un $\mathcal{O}_{U_i}$-module engendré
par ses sections globales.
\item Étant donné $r \geq 0$, on dit que $\mathcal{F}$ est
{\it localement engendré par $r$ sections} si, pour tout objet $U$ de
$\mathcal{C}$, il existe un recouvrement
$\{U_i \to U\}_{i \in I}$ de $\mathcal{C}$ tel que chaque restriction
$\mathcal{F}|_{\mathcal{C}/U_i}$ soit un $\mathcal{O}_{U_i}$-module engendré
par $r$ sections globales.
\item On dit que $\mathcal{F}$ est {\it de type fini} si, pour tout objet $U$
de $\mathcal{C}$, il existe un recouvrement
$\{U_i \to U\}_{i \in I}$ de $\mathcal{C}$ tel que chaque restriction
$\mathcal{F}|_{\mathcal{C}/U_i}$ soit un $\mathcal{O}_{U_i}$-module engendré
par un nombre fini de sections globales.
\item On dit que $\mathcal{F}$ est {\it quasi-cohérent} si, pour tout objet
$U$ de $\mathcal{C}$, il existe un recouvrement
$\{U_i \to U\}_{i \in I}$ de $\mathcal{C}$ tel que chaque restriction
$\mathcal{F}|_{\mathcal{C}/U_i}$ soit un $\mathcal{O}_{U_i}$-module admettant
une présentation globale.
\item On dit que $\mathcal{F}$ est {\it de présentation finie} si, pour tout
objet $U$ de $\mathcal{C}$, il existe un recouvrement
$\{U_i \to U\}_{i \in I}$ de $\mathcal{C}$ tel que chaque restriction
$\mathcal{F}|_{\mathcal{C}/U_i}$ soit un $\mathcal{O}_{U_i}$-module admettant
une présentation globale finie.
\item On dit que $\mathcal{F}$ est {\it cohérent} si et seulement si
$\mathcal{F}$ est de type fini et si, pour tout objet $U$ de $\mathcal{C}$
et tous $s_1, \ldots, s_n \in \mathcal{F}(U)$, le noyau de l'application
$\bigoplus_{i = 1, \ldots, n} \mathcal{O}_U \to \mathcal{F}|_U$
est de type fini sur $(\mathcal{C}/U, \mathcal{O}_U)$.
\end{enumerate}
\end{definition}

\begin{lemma}
\label{sites-modules-lemma-special-locally-free}
Chacune des propriétés (1) -- (8) de la
Définition \ref{sites-modules-definition-site-local} est intrinsèque (voir la discussion
de la section \ref{sites-modules-section-intrinsic}).
\end{lemma}

\begin{proof}
Soient $\mathcal{C}$ et $\mathcal{D}$ des sites.
Soit $u : \mathcal{C} \to \mathcal{D}$ un foncteur cocontinu spécial.
Soit $\mathcal{O}$ un faisceau d'anneaux sur $\mathcal{C}$.
Soit $\mathcal{F}$ un faisceau de $\mathcal{O}$-modules sur $\mathcal{C}$.
Soit $g : \Sh(\mathcal{C}) \to \Sh(\mathcal{D})$ l'équivalence de topos
associée à $u$.
Posons $\mathcal{O}' = g_*\mathcal{O}$ et prenons pour
$g^\sharp : \mathcal{O}' \to g_*\mathcal{O}$ l'identité.
Enfin, posons $\mathcal{F}' = g_*\mathcal{F}$.
Soit $\mathcal{P}_l$ l'une des propriétés (1) -- (7) de la
Définition \ref{sites-modules-definition-site-local}.
(Nous traiterons le cas cohérent à la fin de la preuve.)
Notons $\mathcal{P}_g$ la propriété correspondante de la
Définition \ref{sites-modules-definition-global}. Nous avons déjà vu que
$\mathcal{P}_g$ est intrinsèque. Il faut montrer que
$\mathcal{P}_l(\mathcal{C}, \mathcal{O}, \mathcal{F})$
est vraie si et seulement si
$\mathcal{P}_l(\mathcal{D}, \mathcal{O}', \mathcal{F}')$
est vraie.

\medskip\noindent
Supposons que $\mathcal{F}$ possède la propriété $\mathcal{P}_l$.
Soit $V$ un objet de $\mathcal{D}$.
L'une des propriétés d'un foncteur cocontinu spécial assure l'existence
d'un recouvrement $\{u(U_i) \to V\}_{i \in I}$ dans le site $\mathcal{D}$.
Par hypothèse, pour chaque $i$, il existe un recouvrement
$\{U_{ij} \to U_i\}_{j \in J_i}$ dans $\mathcal{C}$ tel que chaque restriction
$\mathcal{F}|_{U_{ij}}$ possède la propriété $\mathcal{P}_g$. D'après
Sites, Lemme \ref{sites-lemma-localize-special-cocontinuous}
on a des diagrammes commutatifs de topos annelés
$$
\xymatrix{
(\Sh(\mathcal{C}/U_{ij}), \mathcal{O}_{U_{ij}}) \ar[r] \ar[d] &
(\Sh(\mathcal{C}), \mathcal{O}) \ar[d] \\
(\Sh(\mathcal{D}/u(U_{ij})), \mathcal{O}'_{u(U_{ij})}) \ar[r] &
(\Sh(\mathcal{D}), \mathcal{O}')
}
$$
où les flèches verticales sont des équivalences. Ainsi
$\mathcal{F}'|_{u(U_{ij})}$ possède aussi la propriété $\mathcal{P}_g$.
De plus, $\{u(U_{ij}) \to V\}_{i \in I, j \in J_i}$ est un recouvrement du
site $\mathcal{D}$. Par conséquent, $\mathcal{F}'$ possède la propriété
$\mathcal{P}_l$.

\medskip\noindent
Supposons que $\mathcal{F}'$ possède la propriété $\mathcal{P}_l$.
Soit $U$ un objet de $\mathcal{C}$.
Par hypothèse, il existe un recouvrement
$\{V_i \to u(U)\}_{i \in I}$ tel que $\mathcal{F}'|_{V_i}$ possède la
propriété $\mathcal{P}_g$. Puisque $u$ est cocontinu, on peut raffiner ce
recouvrement par une famille $\{u(U_j) \to u(U)\}_{j \in J}$, où
$\{U_j \to U\}_{j \in J}$ est un recouvrement dans $\mathcal{C}$.
Notons que le raffinement est donné par $\alpha : J \to I$ et par des
morphismes $u(U_j) \to V_{\alpha(j)}$.
La restriction est transitive, c'est-à-dire
$(\mathcal{F}'|_{V_{\alpha(j)}})|_{u(U_j)} = \mathcal{F}'|_{u(U_j)}$.
Le Lemme \ref{sites-modules-lemma-global-pullback} montre donc que
$\mathcal{F}'|_{u(U_j)}$ possède la propriété $\mathcal{P}_g$.
Ainsi, le diagramme
$$
\xymatrix{
(\Sh(\mathcal{C}/U_j), \mathcal{O}_{U_j}) \ar[r] \ar[d] &
(\Sh(\mathcal{C}), \mathcal{O}) \ar[d] \\
(\Sh(\mathcal{D}/u(U_j)), \mathcal{O}'_{u(U_j)})
\ar[r] &
(\Sh(\mathcal{D}), \mathcal{O}')
}
$$
où les flèches verticales sont des équivalences montre que
$\mathcal{F}|_{U_j}$ possède aussi la propriété $\mathcal{P}_g$.
Ainsi $\mathcal{F}$ possède la propriété $\mathcal{P}_l$, comme voulu.

\medskip\noindent
Enfin, prouvons le lemme dans le cas
$\mathcal{P}_l = coherent$\footnote{Les détails techniques sont quelque peu
malcommodes ; nous conseillons au lecteur de sauter cette partie de la preuve.}.
Supposons $\mathcal{F}$ cohérent. Alors $\mathcal{F}$ est de type fini, et la
première partie de la preuve montre que $\mathcal{F}'$ l'est aussi.
Soit $V$ un objet de $\mathcal{D}$ et soient
$s_1, \ldots, s_n \in \mathcal{F}'(V)$. Il faut montrer que le noyau
$\mathcal{K}'$ of
$\bigoplus_{j = 1, \ldots, n} \mathcal{O}_V \to \mathcal{F}'|_V$
soit de type fini sur $\mathcal{D}/V$. Il faut donc montrer que, pour tout
$V'/V$, il existe un recouvrement $\{V'_i \to V'\}$ tel que
$\mathcal{F}'|_{V'_i}$ soit engendré par un nombre fini de sections.
En remplaçant $V$ par $V'$ (et en restreignant les sections $s_j$ à $V'$),
on se ramène au cas $V' = V$. Puisque $u$ est un foncteur cocontinu spécial,
il existe un recouvrement $\{u(U_i) \to V\}_{i \in I}$ dans le site
$\mathcal{D}$. En utilisant l'isomorphisme de topos
$\Sh(\mathcal{C}/U_i) = \Sh(\mathcal{D}/u(U_i))$
on voit que $\mathcal{K}'|_{u(U_i)}$ correspond au noyau
$\mathcal{K}_i$ d'un morphisme
$\bigoplus_{j = 1, \ldots, n} \mathcal{O}_{U_i} \to \mathcal{F}|_{U_i}$.
La cohérence de $\mathcal{F}$ montre que $\mathcal{K}_i$ est de type fini.
On en conclut (en appliquant à nouveau la première partie de la preuve) que
$\mathcal{K}|_{u(U_i)}$ est de type fini. Il existe donc des recouvrements
$\{V_{il} \to u(U_i)\}$ tels que $\mathcal{K}|_{V_{il}}$ soit engendré par un
nombre fini de sections globales. Comme $\{V_{il} \to V\}$ est un
recouvrement de $\mathcal{D}$, on conclut que $\mathcal{K}$ est de type fini,
comme voulu.

\medskip\noindent
Supposons $\mathcal{F}'$ cohérent. Alors $\mathcal{F}'$ est de type fini et,
d'après la première partie de la preuve, $\mathcal{F}$ l'est aussi.
Soit $U$ un objet de $\mathcal{C}$ et soient
$s_1, \ldots, s_n \in \mathcal{F}(U)$. Il faut montrer que le noyau
$\mathcal{K}$ of
$\bigoplus_{j = 1, \ldots, n} \mathcal{O}_U \to \mathcal{F}|_U$
est de type fini sur $\mathcal{C}/U$. En utilisant l'isomorphisme de topos
$\Sh(\mathcal{C}/U) = \Sh(\mathcal{D}/u(U))$
on voit que $\mathcal{K}|_U$ correspond au noyau
$\mathcal{K}'$ d'un morphisme
$\bigoplus_{j = 1, \ldots, n} \mathcal{O}_{u(U)} \to \mathcal{F}'|_{u(U)}$.
Comme $\mathcal{F}'$ est cohérent, $\mathcal{K}'$ est de type fini.
La première partie de la preuve montre alors que $\mathcal{K}$ est de type
fini.
\end{proof}

\noindent
Désormais, nous pouvons donc parler des propriétés des
$\mathcal{O}$-modules définies dans la Définition \ref{sites-modules-definition-site-local}
sans préciser de site.

\begin{lemma}
\label{sites-modules-lemma-local-final-object}
Soit $(\Sh(\mathcal{C}), \mathcal{O})$ un topos annelé.
Soit $\mathcal{F}$ un $\mathcal{O}$-module.
Supposons que le site $\mathcal{C}$ admette un objet final $X$.
Alors
\begin{enumerate}
\item Les conditions suivantes sont équivalentes :
\begin{enumerate}
\item $\mathcal{F}$ est localement libre ;
\item il existe un recouvrement $\{X_i \to X\}$ dans $\mathcal{C}$ tel que
chaque restriction $\mathcal{F}|_{\mathcal{C}/X_i}$ soit un
$\mathcal{O}_{X_i}$-module localement libre ;
\item il existe un recouvrement $\{X_i \to X\}$ dans $\mathcal{C}$ tel que
chaque restriction $\mathcal{F}|_{\mathcal{C}/X_i}$ soit un
$\mathcal{O}_{X_i}$-module libre.
\end{enumerate}
\item Les conditions suivantes sont équivalentes :
\begin{enumerate}
\item $\mathcal{F}$ est localement libre de rang fini ;
\item il existe un recouvrement $\{X_i \to X\}$ dans $\mathcal{C}$ tel que
chaque restriction $\mathcal{F}|_{\mathcal{C}/X_i}$ soit un
$\mathcal{O}_{X_i}$-module localement libre de rang fini ;
\item il existe un recouvrement $\{X_i \to X\}$ dans $\mathcal{C}$ tel que
chaque restriction $\mathcal{F}|_{\mathcal{C}/X_i}$ soit un
$\mathcal{O}_{X_i}$-module libre de rang fini.
\end{enumerate}
\item Les conditions suivantes sont équivalentes :
\begin{enumerate}
\item $\mathcal{F}$ est localement engendré par des sections ;
\item il existe un recouvrement $\{X_i \to X\}$ dans $\mathcal{C}$ tel que
chaque restriction $\mathcal{F}|_{\mathcal{C}/X_i}$ soit un
$\mathcal{O}_{X_i}$-module localement engendré par des sections ;
\item il existe un recouvrement $\{X_i \to X\}$ dans $\mathcal{C}$ tel que
chaque restriction $\mathcal{F}|_{\mathcal{C}/X_i}$ soit un
$\mathcal{O}_{X_i}$-module engendré par ses sections globales.
\end{enumerate}
\item Étant donné $r \geq 0$, les conditions suivantes sont équivalentes :
\begin{enumerate}
\item $\mathcal{F}$ est localement engendré par $r$ sections ;
\item il existe un recouvrement $\{X_i \to X\}$ dans $\mathcal{C}$ tel que
chaque restriction $\mathcal{F}|_{\mathcal{C}/X_i}$ soit un
$\mathcal{O}_{X_i}$-module localement engendré par $r$ sections ;
\item il existe un recouvrement $\{X_i \to X\}$ dans $\mathcal{C}$ tel que
chaque restriction $\mathcal{F}|_{\mathcal{C}/X_i}$ soit un
$\mathcal{O}_{X_i}$-module engendré par $r$ sections globales.
\end{enumerate}
\item Les conditions suivantes sont équivalentes :
\begin{enumerate}
\item $\mathcal{F}$ est de type fini ;
\item il existe un recouvrement $\{X_i \to X\}$ dans $\mathcal{C}$ tel que
chaque restriction $\mathcal{F}|_{\mathcal{C}/X_i}$ soit un
$\mathcal{O}_{X_i}$-module de type fini ;
\item il existe un recouvrement $\{X_i \to X\}$ dans $\mathcal{C}$ tel que
chaque restriction $\mathcal{F}|_{\mathcal{C}/X_i}$ soit un
$\mathcal{O}_{X_i}$-module engendré par un nombre fini de sections globales.
\end{enumerate}
\item Les conditions suivantes sont équivalentes :
\begin{enumerate}
\item $\mathcal{F}$ est quasi-cohérent ;
\item il existe un recouvrement $\{X_i \to X\}$ dans $\mathcal{C}$ tel que
chaque restriction $\mathcal{F}|_{\mathcal{C}/X_i}$ soit un
$\mathcal{O}_{X_i}$-module quasi-cohérent ;
\item il existe un recouvrement $\{X_i \to X\}$ dans $\mathcal{C}$ tel que
chaque restriction $\mathcal{F}|_{\mathcal{C}/X_i}$ soit un
$\mathcal{O}_{X_i}$-module admettant une présentation globale.
\end{enumerate}
\item Les conditions suivantes sont équivalentes :
\begin{enumerate}
\item $\mathcal{F}$ est de présentation finie ;
\item il existe un recouvrement $\{X_i \to X\}$ dans $\mathcal{C}$ tel que
chaque restriction $\mathcal{F}|_{\mathcal{C}/X_i}$ soit un
$\mathcal{O}_{X_i}$-module de présentation finie ;
\item il existe un recouvrement $\{X_i \to X\}$ dans $\mathcal{C}$ tel que
chaque restriction $\mathcal{F}|_{\mathcal{C}/X_i}$ soit un
$\mathcal{O}_{X_i}$-module admettant une présentation globale finie.
\end{enumerate}
\item Les conditions suivantes sont équivalentes :
\begin{enumerate}
\item $\mathcal{F}$ est cohérent ;
\item il existe un recouvrement $\{X_i \to X\}$ dans $\mathcal{C}$ tel que
chaque restriction $\mathcal{F}|_{\mathcal{C}/X_i}$ soit un
$\mathcal{O}_{X_i}$-module cohérent.
\end{enumerate}
\end{enumerate}
\end{lemma}

\begin{proof}
Dans chaque cas, (a) $\Rightarrow (b)$. Dans chacun des cas (1) -- (6),
la condition (b) implique la condition (c) par l'axiome (2) d'un site
(voir Sites, Définition \ref{sites-definition-site}) et par la définition
des propriétés locales des modules.
Supposons que $\{X_i \to X\}$ soit un recouvrement.
Pour tout objet $U$ de $\mathcal{C}$, on obtient alors le recouvrement induit
$\{X_i \times_X U \to U\}$. De plus, la propriété globale de
$\mathcal{F}|_{\mathcal{C}/X_i}$ dans (c) implique, d'après le
Lemme \ref{sites-modules-lemma-global-pullback}, la propriété globale correspondante de
$\mathcal{F}|_{\mathcal{C}/X_i \times_X U}$ ; le faisceau possède donc par
définition la propriété (a). Nous omettons la preuve de
(b) $\Rightarrow$ (a) dans le cas (7).
\end{proof}

\begin{lemma}
\label{sites-modules-lemma-local-pullback}
Soit
$(f, f^\sharp) :
(\Sh(\mathcal{C}), \mathcal{O}_\mathcal{C})
\to
(\Sh(\mathcal{D}), \mathcal{O}_\mathcal{D})$
soit un morphisme de topos annelés.
Soit $\mathcal{F}$ un $\mathcal{O}_\mathcal{D}$-module.
\begin{enumerate}
\item Si $\mathcal{F}$ est localement libre, alors $f^*\mathcal{F}$ est
localement libre.
\item Si $\mathcal{F}$ est localement libre de rang fini, alors
$f^*\mathcal{F}$ est localement libre de rang fini.
\item Si $\mathcal{F}$ est localement engendré par des sections, alors
$f^*\mathcal{F}$ est localement engendré par des sections.
\item Si $\mathcal{F}$ est localement engendré par $r$ sections, alors
$f^*\mathcal{F}$ est localement engendré par $r$ sections.
\item Si $\mathcal{F}$ est de type fini, alors $f^*\mathcal{F}$ est de type
fini.
\item Si $\mathcal{F}$ est quasi-cohérent, alors $f^*\mathcal{F}$ est
quasi-cohérent.
\item Si $\mathcal{F}$ est de présentation finie, alors $f^*\mathcal{F}$ est
de présentation finie.
\end{enumerate}
\end{lemma}

\begin{proof}
Selon la discussion de la section \ref{sites-modules-section-intrinsic}, il suffit de
vérifier la conservation par image inverse pour un morphisme de sites annelés
$(f, f^\sharp) :
(\mathcal{C}, \mathcal{O}_\mathcal{C})
\to
(\mathcal{D}, \mathcal{O}_\mathcal{D})$
tel que $f$ soit donné par un foncteur exact à gauche et continu
$u : \mathcal{D} \to \mathcal{C}$ entre des sites admettant toutes les
limites finies.
Soit $\mathcal{G}$ un faisceau de $\mathcal{O}_\mathcal{D}$-modules possédant
l'une des propriétés (1) -- (6) de la
Définition \ref{sites-modules-definition-site-local}.
Nous savons que $\mathcal{D}$ admet un objet final $Y$ et que $X = u(Y)$ est
un objet final de $\mathcal{C}$. Par hypothèse, il existe un recouvrement
$\{Y_i \to Y\}$ tel que $\mathcal{G}|_{\mathcal{D}/Y_i}$ possède la
propriété globale correspondante. Posons $X_i = u(Y_i)$ ; alors
$\{X_i \to X\}$ est un recouvrement dans $\mathcal{C}$.
On obtient un diagramme commutatif de morphismes de sites annelés
$$
\xymatrix{
(\mathcal{C}/X_i, \mathcal{O}_\mathcal{C}|_{X_i}) \ar[r] \ar[d] &
(\mathcal{C}, \mathcal{O}_\mathcal{C}) \ar[d] \\
(\mathcal{D}/Y_i, \mathcal{O}_\mathcal{D}|_{Y_i}) \ar[r] &
(\mathcal{D}, \mathcal{O}_\mathcal{D})
}
$$
d'après Sites, Lemme \ref{sites-lemma-localize-morphism-strong}.
Le Lemme \ref{sites-modules-lemma-global-pullback} montre donc que
$f^*\mathcal{G}|_{X_i}$ possède la propriété globale correspondante.
Le Lemme \ref{sites-modules-lemma-local-final-object} permet alors de conclure que
$\mathcal{G}$ possède la propriété locale considérée.
\end{proof}







\section{Résultats élémentaires sur les propriétés locales des modules}
\label{sites-modules-section-basics}

\noindent
Voici quelques lemmes élémentaires relatifs aux définitions précédentes.

\begin{lemma}
\label{sites-modules-lemma-cokernel-finite-finite-presentation}
Soit $(\mathcal{C},\mathcal{O})$ un site annelé. Soit $\mathcal{F}$ un
$\mathcal{O}$-module de présentation finie. Soit
$\varphi : \mathcal{G} \to \mathcal{F}$ un morphisme de
$\mathcal{O}$-modules. Si $\mathcal{G}$ est de type fini, alors
$\Coker(\varphi)$ est de présentation finie.
\end{lemma}

\begin{proof}
La preuve est omise. Indication : voir Modules, Lemme
\ref{modules-lemma-cokernel-finite-finite-presentation}.
\end{proof}

\begin{lemma}
\label{sites-modules-lemma-kernel-surjection-finite-onto-finite-presentation}
Soit $(\mathcal{C}, \mathcal{O})$ un site annelé.
Soit $\theta : \mathcal{G} \to \mathcal{F}$ un morphisme surjectif de
$\mathcal{O}$-modules, où $\mathcal{F}$ est de présentation finie et
$\mathcal{G}$ de type fini. Alors $\Ker(\theta)$ est de type fini.
\end{lemma}

\begin{proof}
La preuve est omise. Indication : voir Modules, Lemme
\ref{modules-lemma-kernel-surjection-finite-free-onto-finite-presentation}.
\end{proof}

\begin{lemma}
\label{sites-modules-lemma-chaotic-quasi-coherent}
Soit $\mathcal{C}$ une catégorie considérée comme un site muni de la
topologie chaotique, voir Sites, Exemple \ref{sites-example-indiscrete}.
Soient $\mathcal{O}$ un faisceau d'anneaux sur $\mathcal{C}$ et
$\mathcal{F}$ un faisceau de $\mathcal{O}$-modules. Alors $\mathcal{F}$ est
quasi-cohérent si et seulement si, pour tout $U \to V$ dans $\mathcal{C}$,
l'application canonique
$$
\mathcal{F}(V) \otimes_{\mathcal{O}(V)} \mathcal{O}(U)
\longrightarrow
\mathcal{F}(U)
$$
est un isomorphisme.
\end{lemma}

\begin{proof}
Supposons $\mathcal{F}$ quasi-cohérent et soit $U \to V$ un morphisme de
$\mathcal{C}$. Comme tout recouvrement de $V$ est donné par un isomorphisme,
la Définition \ref{sites-modules-definition-site-local} donne une présentation
$$
\bigoplus\nolimits_{j \in J} \mathcal{O}_V \longrightarrow
\bigoplus\nolimits_{i \in I} \mathcal{O}_V \longrightarrow
\mathcal{F}|_{\mathcal{C}/V} \longrightarrow 0
$$
La topologie sur $\mathcal{C}$ étant chaotique, la prise des sections sur
tout objet de $\mathcal{C}$ est exacte. On obtient donc une présentation
$$
\bigoplus\nolimits_{j \in J} \mathcal{O}(V) \longrightarrow
\bigoplus\nolimits_{i \in I} \mathcal{O}(V) \longrightarrow
\mathcal{F}(V) \longrightarrow 0
$$
de $\mathcal{F}(V)$ comme $\mathcal{O}(V)$-module, et de même pour
$\mathcal{F}(U)$. Cela montre aisément que l'application affichée dans
l'énoncé du lemme est un isomorphisme.

\medskip\noindent
Supposons que l'application affichée dans l'énoncé du lemme soit un
isomorphisme pour tout morphisme $U \to V$ de $\mathcal{C}$. Fixons $V$ et
choisissons une présentation
$$
\bigoplus\nolimits_{j \in J} \mathcal{O}(V) \longrightarrow
\bigoplus\nolimits_{i \in I} \mathcal{O}(V) \longrightarrow
\mathcal{F}(V) \longrightarrow 0
$$
de $\mathcal{F}(V)$ comme $\mathcal{O}(V)$-module. L'hypothèse sur
$\mathcal{F}$ signifie exactement que la suite correspondante
$$
\bigoplus\nolimits_{j \in J} \mathcal{O}_V \longrightarrow
\bigoplus\nolimits_{i \in I} \mathcal{O}_V \longrightarrow
\mathcal{F}|_{\mathcal{C}/V} \longrightarrow 0
$$
est exacte, d'où l'on conclut que $\mathcal{F}$ est quasi-cohérent.
\end{proof}

\begin{lemma}
\label{sites-modules-lemma-chaotic-flat-quasi-coherent}
Soit $\mathcal{C}$ une catégorie considérée comme un site muni de la
topologie chaotique, voir Sites, Exemple \ref{sites-example-indiscrete}.
Soit $\mathcal{O}$ un faisceau d'anneaux sur $\mathcal{C}$.
Supposons que, pour tout $U \to V$ dans $\mathcal{C}$, l'application de
restriction $\mathcal{O}(V) \to \mathcal{O}(U)$ soit un homomorphisme plat
d'anneaux. Alors la catégorie des $\mathcal{O}$-modules quasi-cohérents est
une sous-catégorie de Serre faible de $\textit{Mod}(\mathcal{O})$.
\end{lemma}

\begin{proof}
Nous allons vérifier la définition d'une sous-catégorie de Serre faible,
voir Homology, Définition \ref{homology-definition-serre-subcategory}.
Pour cela, nous utiliserons la caractérisation des modules quasi-cohérents
donnée au Lemme \ref{sites-modules-lemma-chaotic-quasi-coherent}.
Considérons une suite exacte
$$
\mathcal{F}_0 \to \mathcal{F}_1 \to
\mathcal{F}_2 \to \mathcal{F}_3 \to
\mathcal{F}_4
$$
dans $\textit{Mod}(\mathcal{O})$, où $\mathcal{F}_0$, $\mathcal{F}_1$,
$\mathcal{F}_3$ et $\mathcal{F}_4$ sont quasi-cohérents. Soit $U \to V$ un
morphisme de $\mathcal{C}$ et considérons le diagramme commutatif
$$
\xymatrix@C=1em{
{\scriptstyle \mathcal{F}_0(V) \otimes_{\mathcal{O}(V)} \mathcal{O}(U)} \ar[r] \ar[d] &
{\scriptstyle \mathcal{F}_1(V) \otimes_{\mathcal{O}(V)} \mathcal{O}(U)} \ar[r] \ar[d] &
{\scriptstyle \mathcal{F}_2(V) \otimes_{\mathcal{O}(V)} \mathcal{O}(U)} \ar[r] \ar[d] &
{\scriptstyle \mathcal{F}_3(V) \otimes_{\mathcal{O}(V)} \mathcal{O}(U)} \ar[r] \ar[d] &
{\scriptstyle \mathcal{F}_4(V) \otimes_{\mathcal{O}(V)} \mathcal{O}(U)} \ar[d] \\
{\scriptstyle \mathcal{F}_0(U)} \ar[r] &
{\scriptstyle \mathcal{F}_1(U)} \ar[r] &
{\scriptstyle \mathcal{F}_2(U)} \ar[r] &
{\scriptstyle \mathcal{F}_3(U)} \ar[r] &
{\scriptstyle \mathcal{F}_4(U)}
}
$$
Par hypothèse, les flèches verticales d'indices $0$, $1$, $3$, $4$ sont des
isomorphismes. Comme la topologie sur $\mathcal{C}$ est chaotique, la prise
des sections sur un objet de $\mathcal{C}$ est exacte ; la ligne inférieure
est donc exacte. Comme $\mathcal{O}(V) \to \mathcal{O}(U)$ est plat, la ligne
supérieure est elle aussi exacte. Le lemme des $5$ permet alors de conclure
que la flèche du milieu est un isomorphisme
(Homology, Lemme \ref{homology-lemma-five-lemma}).
\end{proof}






\section{Immersions fermées de topos annelés}
\label{sites-modules-section-closed-immersion}

\noindent
Quand doit-on déclarer qu'un morphisme de topos annelés
$i : (\Sh(\mathcal{C}), \mathcal{O}) \to (\Sh(\mathcal{D}), \mathcal{O}')$
est une immersion fermée ? Par analogie avec la discussion de
Modules, section \ref{modules-section-closed-immersion}, il semble naturel de
supposer au moins que :
\begin{enumerate}
\item le foncteur $i$ est une immersion fermée de topos
(Sites, Définition \ref{sites-definition-immersion-topoi}).
\item l'application associée $\mathcal{O}' \to i_*\mathcal{O}$ est surjective.
\end{enumerate}
Ces conditions impliquent déjà plusieurs résultats satisfaisants, que nous
examinons dans cette section. Il paraît toutefois prudent de ne pas définir
formellement la notion d'immersion fermée de topos annelés, car plusieurs
définitions différentes seraient possibles.

\begin{lemma}
\label{sites-modules-lemma-i-star-equivalence}
Soit $i : (\Sh(\mathcal{C}), \mathcal{O}) \to (\Sh(\mathcal{D}), \mathcal{O}')$
un morphisme de topos annelés. Supposons que $i$ soit une immersion fermée
de topos et que $i^\sharp : \mathcal{O}' \to i_*\mathcal{O}$ soit surjectif.
Notons $\mathcal{I} \subset \mathcal{O}'$ le noyau de $i^\sharp$.
Le foncteur
$$
i_* :
\textit{Mod}(\mathcal{O})
\longrightarrow
\textit{Mod}(\mathcal{O}')
$$
est exact et pleinement fidèle ; son image essentielle est formée des
$\mathcal{O}'$-modules $\mathcal{G}$ tels que $\mathcal{I}\mathcal{G} = 0$.
\end{lemma}

\begin{proof}
Le Lemme \ref{sites-modules-lemma-exactness} et Sites, Lemme
\ref{sites-lemma-closed-immersion} montrent que $i_*$ est exact.
Puisque $i_*$ est pleinement fidèle sur les faisceaux d'ensembles et que
$i^\sharp$ est surjectif, $i_*$ est pleinement fidèle comme foncteur
$\textit{Mod}(\mathcal{O}) \to \textit{Mod}(\mathcal{O}')$.
En effet, supposons que
$\alpha : i_*\mathcal{F}_1 \to i_*\mathcal{F}_2$ soit un morphisme de
$\mathcal{O}'$-modules. La pleine fidélité de $i_*$ fournit une application
$\beta : \mathcal{F}_1 \to \mathcal{F}_2$ de faisceaux d'ensembles.
Pour montrer que $\beta$ est un morphisme de modules, il faut vérifier que
$$
\xymatrix{
\mathcal{O} \times \mathcal{F}_1 \ar[r] \ar[d] &
\mathcal{F}_1 \ar[d] \\
\mathcal{O} \times \mathcal{F}_2 \ar[r] &
\mathcal{F}_2
}
$$
est commutatif. Il suffit de le vérifier après application de $i_*$.
Considérons
$$
\xymatrix{
\mathcal{O}' \times i_*\mathcal{F}_1 \ar[r] \ar[d] &
i_*\mathcal{O} \times i_*\mathcal{F}_1 \ar[r] \ar[d] &
i_*\mathcal{F}_1 \ar[d] \\
\mathcal{O}' \times i_*\mathcal{F}_2 \ar[r] &
i_*\mathcal{O} \times i_*\mathcal{F}_2 \ar[r] &
i_*\mathcal{F}_2
}
$$
Le rectangle extérieur est commutatif. La conclusion résulte de la
surjectivité de $i^\sharp$.

\medskip\noindent
Pour achever la preuve, il reste à établir l'assertion sur l'image
essentielle de $i_*$. Il est clair que $i_*\mathcal{F}$ est annulé par
$\mathcal{I}$ pour tout $\mathcal{O}$-module $\mathcal{F}$. Réciproquement,
soit $\mathcal{G}$ un $\mathcal{O}'$-module tel que
$\mathcal{I}\mathcal{G} = 0$. Par définition d'un sous-topos fermé, il existe
un sous-faisceau $\mathcal{U}$ de l'objet final de $\mathcal{D}$ tel que
l'image essentielle de $i_*$ sur les faisceaux d'ensembles soit la classe des
faisceaux d'ensembles $\mathcal{H}$ pour lesquels
$\mathcal{H} \times \mathcal{U} \to \mathcal{U}$ est un isomorphisme.
En particulier, $i_*\mathcal{O} \times \mathcal{U} = \mathcal{U}$, ce qui
implique $\mathcal{I} \times \mathcal{U} = \mathcal{O} \times \mathcal{U}$.
Notre module $\mathcal{G}$ satisfait donc
$\mathcal{G} \times \mathcal{U} = \{0\} \times \mathcal{U} = \mathcal{U}$
(car le module nul est isomorphe à l'objet final de la catégorie des
faisceaux d'ensembles). Il existe donc un faisceau d'ensembles $\mathcal{F}$
sur $\mathcal{C}$ tel que $i_*\mathcal{F} = \mathcal{G}$. Comme $i_*$ est
pleinement fidèle sur les faisceaux d'ensembles, pour définir l'addition
$\mathcal{F} \times \mathcal{F} \to \mathcal{F}$ et la multiplication
$\mathcal{O} \times \mathcal{F} \to \mathcal{F}$, il suffit d'utiliser
l'addition
$$
\mathcal{G} \times \mathcal{G} \longrightarrow \mathcal{G}
$$
(dont on dispose puisque $\mathcal{G}$ est un $\mathcal{O}'$-module) et la
multiplication
$$
i_*\mathcal{O} \times \mathcal{G} \to \mathcal{G}
$$
dont on dispose car la multiplication par $\mathcal{O}'$ annule
$\mathcal{I}$ par hypothèse et
$i_*\mathcal{O} = \mathcal{O}'/\mathcal{I}$. Par construction,
$\mathcal{G}$ est isomorphe à l'image directe du $\mathcal{O}$-module
$\mathcal{F}$ ainsi construit.
\end{proof}












\section{Produit tensoriel}
\label{sites-modules-section-tensor-product}

\noindent
Dans les sections \ref{sites-modules-section-presheaves-modules} et
\ref{sites-modules-section-sheafification-presheaves-modules}, nous avons défini le
foncteur de changement d'anneaux au moyen d'un produit tensoriel.
Cette construction permet aussi de définir le produit tensoriel de
préfaisceaux de $\mathcal{O}$-modules. Plus précisément, supposons que
$\mathcal{C}$ soit une catégorie, $\mathcal{O}$ un préfaisceau d'anneaux et
$\mathcal{F}$, $\mathcal{G}$ des préfaisceaux de $\mathcal{O}$-modules.
On définit alors $\mathcal{F} \otimes_{p, \mathcal{O}} \mathcal{G}$ comme le
préfaisceau
$$
U
\longmapsto
(\mathcal{F} \otimes_{p, \mathcal{O}} \mathcal{G})(U)
=
\mathcal{F}(U) \otimes_{\mathcal{O}(U)} \mathcal{G}(U)
$$
Si $\mathcal{C}$ est un site, $\mathcal{O}$ un faisceau d'anneaux et
$\mathcal{F}$, $\mathcal{G}$ des faisceaux de $\mathcal{O}$-modules, on pose
$$
\mathcal{F} \otimes_\mathcal{O} \mathcal{G}
=
(\mathcal{F} \otimes_{p, \mathcal{O}} \mathcal{G})^\#
$$
qui est le faisceau de $\mathcal{O}$-modules associé au préfaisceau
$\mathcal{F} \otimes_{p, \mathcal{O}} \mathcal{G}$.

\medskip\noindent
Voici quelques formules que nous utiliserons plus loin sans autre mention :
$$
(\mathcal{F}
\otimes_{p, \mathcal{O}} \mathcal{G})
\otimes_{p, \mathcal{O}} \mathcal{H}
=
\mathcal{F}
\otimes_{p, \mathcal{O}} (\mathcal{G}
\otimes_{p, \mathcal{O}} \mathcal{H}),
$$
et de même pour les faisceaux.
Si $\mathcal{O}_1 \to \mathcal{O}_2$ est un morphisme de préfaisceaux
d'anneaux, alors
$$
(\mathcal{F} \otimes_{p, \mathcal{O}_1} \mathcal{G})
\otimes_{p, \mathcal{O}_1} \mathcal{O}_2 =
(\mathcal{F} \otimes_{p, \mathcal{O}_1} \mathcal{O}_2)
\otimes_{p, \mathcal{O}_2}
(\mathcal{G} \otimes_{p, \mathcal{O}_1} \mathcal{O}_2),
$$
et de même pour les faisceaux.
Ces formules résultent de leurs analogues algébriques et de la
faisceautisation.

\begin{lemma}
\label{sites-modules-lemma-sheafification-tensor}
Soit $\mathcal{C}$ un site. Soit $\mathcal{O}$ un préfaisceau d'anneaux.
Soient $\mathcal{F}$ et $\mathcal{G}$ des préfaisceaux de
$\mathcal{O}$-modules. Alors
$\mathcal{F}^\# \otimes_{\mathcal{O}^\#} \mathcal{G}^\#$
est égal à
$(\mathcal{F} \otimes_{p, \mathcal{O}} \mathcal{G})^\#$.
\end{lemma}

\begin{proof}
La preuve est omise. Indication : utiliser la caractérisation du produit
tensoriel par les applications bilinéaires donnée ci-dessous, ainsi que la
propriété universelle de la faisceautisation.
\end{proof}

\noindent
Soit $\mathcal{C}$ un site, soit $\mathcal{O}$ un faisceau d'anneaux et soient
$\mathcal{F}$, $\mathcal{G}$, $\mathcal{H}$ des faisceaux de
$\mathcal{O}$-modules. On pose
$$
\text{Bilin}_\mathcal{O}(\mathcal{F} \times \mathcal{G}, \mathcal{H})
=
\{\varphi \in
\Mor_{\Sh(\mathcal{C})}(
\mathcal{F} \times \mathcal{G}, \mathcal{H}) \mid
\varphi \text{ est }\mathcal{O}\text{-bilinéaire}\}.
$$
Avec cette définition, on a
$$
\Hom_\mathcal{O}
(\mathcal{F} \otimes_\mathcal{O} \mathcal{G}, \mathcal{H})
=
\text{Bilin}_\mathcal{O}(\mathcal{F} \times \mathcal{G}, \mathcal{H}).
$$
Autrement dit, $\mathcal{F} \otimes_\mathcal{O} \mathcal{G}$ représente le
foncteur qui associe à $\mathcal{H}$ l'ensemble des applications bilinéaires
$\mathcal{F} \times \mathcal{G} \to \mathcal{H}$.
En particulier, puisque la notion d'application bilinéaire a un sens pour
une paire de modules sur un topos annelé, le produit tensoriel de faisceaux de
modules est intrinsèque au topos (comparer la discussion de la
section \ref{sites-modules-section-intrinsic}). On a en fait le résultat suivant.

\begin{lemma}
\label{sites-modules-lemma-tensor-product-pullback}
Soit $f : (\Sh(\mathcal{C}), \mathcal{O}_\mathcal{C})
\to (\Sh(\mathcal{D}), \mathcal{O}_\mathcal{D})$ un morphisme de topos
annelés. Soient $\mathcal{F}$ et $\mathcal{G}$ des
$\mathcal{O}_\mathcal{D}$-modules. Alors
$f^*(\mathcal{F} \otimes_{\mathcal{O}_\mathcal{D}} \mathcal{G})
= f^*\mathcal{F} \otimes_{\mathcal{O}_\mathcal{C}} f^*\mathcal{G}$
fonctoriellement en $\mathcal{F}$ et $\mathcal{G}$.
\end{lemma}

\begin{proof}
Pour tout faisceau $\mathcal{H}$ de $\mathcal{O}_\mathcal{C}$-modules, on a
\begin{align*}
\Hom_{\mathcal{O}_\mathcal{C}}(
f^*(\mathcal{F} \otimes_\mathcal{O} \mathcal{G}), \mathcal{H})
& =
\Hom_{\mathcal{O}_\mathcal{D}}(
\mathcal{F} \otimes_\mathcal{O} \mathcal{G}, f_*\mathcal{H}) \\
& =
\text{Bilin}_{\mathcal{O}_\mathcal{D}}(
\mathcal{F} \times \mathcal{G}, f_*\mathcal{H}) \\
& =
\text{Bilin}_{f^{-1}\mathcal{O}_\mathcal{D}}(
f^{-1}\mathcal{F} \times f^{-1}\mathcal{G}, \mathcal{H}) \\
& =
\Hom_{f^{-1}\mathcal{O}_\mathcal{D}}(
f^{-1}\mathcal{F} \otimes_{f^{-1}\mathcal{O}_\mathcal{D}} f^{-1}\mathcal{G},
\mathcal{H}) \\
& =
\Hom_{\mathcal{O}_\mathcal{C}}(
f^*\mathcal{F} \otimes_{f^*\mathcal{O}_\mathcal{D}} f^*\mathcal{G},
\mathcal{H})
\end{align*}
Le symbole « $=$ » intéressant de cette chaîne correspond à la
troisième égalité. Elle résulte
directement de la définition et de l'adjonction entre $f_*$ et $f^{-1}$
(discutée dans les sections précédentes).
\end{proof}

\begin{lemma}
\label{sites-modules-lemma-tensor-product-permanence}
Soit $(\mathcal{C}, \mathcal{O})$ un site annelé.
Soient $\mathcal{F}$ et $\mathcal{G}$ des faisceaux de $\mathcal{O}$-modules.
\begin{enumerate}
\item Si $\mathcal{F}$ et $\mathcal{G}$ sont localement libres, alors
$\mathcal{F} \otimes_\mathcal{O} \mathcal{G}$ l'est aussi.
\item Si $\mathcal{F}$ et $\mathcal{G}$ sont localement libres de rang fini,
alors $\mathcal{F} \otimes_\mathcal{O} \mathcal{G}$ l'est aussi.
\item Si $\mathcal{F}$ et $\mathcal{G}$ sont localement engendrés par des
sections, alors $\mathcal{F} \otimes_\mathcal{O} \mathcal{G}$ l'est aussi.
\item Si $\mathcal{F}$ et $\mathcal{G}$ sont de type fini, alors
$\mathcal{F} \otimes_\mathcal{O} \mathcal{G}$ l'est aussi.
\item Si $\mathcal{F}$ et $\mathcal{G}$ sont quasi-cohérents, alors
$\mathcal{F} \otimes_\mathcal{O} \mathcal{G}$ l'est aussi.
\item Si $\mathcal{F}$ et $\mathcal{G}$ sont de présentation finie, alors
$\mathcal{F} \otimes_\mathcal{O} \mathcal{G}$ l'est aussi.
\item Si $\mathcal{F}$ est de présentation finie et $\mathcal{G}$ cohérent,
alors $\mathcal{F} \otimes_\mathcal{O} \mathcal{G}$ est cohérent.
\item Si $\mathcal{F}$ et $\mathcal{G}$ sont cohérents, alors
$\mathcal{F} \otimes_\mathcal{O} \mathcal{G}$ l'est aussi.
\end{enumerate}
\end{lemma}

\begin{proof}
La preuve est omise. Indication : comparer avec
Faisceaux de modules, Lemme \ref{modules-lemma-tensor-product-permanence}.
\end{proof}



\section{Hom interne}
\label{sites-modules-section-internal-hom}

\noindent
Soit $\mathcal{C}$ une catégorie et soit $\mathcal{O}$ un préfaisceau
d'anneaux. Soient $\mathcal{F}$ et $\mathcal{G}$ des préfaisceaux de
$\mathcal{O}$-modules. Considérons la règle
$$
U \longmapsto \Hom_{\mathcal{O}_U}(\mathcal{F}|_U, \mathcal{G}|_U).
$$
Pour $\varphi : V \to U$ dans $\mathcal{C}$, on définit une application de
restriction
$$
\Hom_{\mathcal{O}_U}(\mathcal{F}|_U, \mathcal{G}|_U)
\longrightarrow
\Hom_{\mathcal{O}_V}(\mathcal{F}|_V, \mathcal{G}|_V)
$$
en restreignant par le morphisme de relocalisation
$j : \mathcal{C}/V \to \mathcal{C}/U$, voir Sites, Lemme
\ref{sites-lemma-relocalize}. Cela définit donc un préfaisceau
$\SheafHom_\mathcal{O}(\mathcal{F}, \mathcal{G})$.
De plus, étant donnés
$\varphi \in \Hom_{\mathcal{O}|_U}(\mathcal{F}|_U, \mathcal{G}|_U)$ et une
section $f \in \mathcal{O}(U)$, on peut définir
$f\varphi \in \Hom_{\mathcal{O}|_U}(\mathcal{F}|_U, \mathcal{G}|_U)$ soit
en précomposant par la multiplication par $f$ sur $\mathcal{F}|_U$, soit en
postcomposant par la multiplication par $f$ sur $\mathcal{G}|_U$ (le résultat
est le même). On obtient donc en fait un préfaisceau de
$\mathcal{O}$-modules. Il existe un morphisme canonique d'évaluation
$$
\mathcal{F}
\otimes_{p, \mathcal{O}}
\SheafHom_\mathcal{O}(\mathcal{F}, \mathcal{G})
\longrightarrow
\mathcal{G}.
$$

\begin{lemma}
\label{sites-modules-lemma-internal-hom}
Si $\mathcal{C}$ est un site, $\mathcal{O}$ un faisceau d'anneaux,
$\mathcal{F}$ un préfaisceau de $\mathcal{O}$-modules et
$\mathcal{G}$ un faisceau de $\mathcal{O}$-modules, alors
$\SheafHom_\mathcal{O}(\mathcal{F}, \mathcal{G})$
est un faisceau de $\mathcal{O}$-modules.
\end{lemma}

\begin{proof}
La preuve est omise. Indications : noter d'abord que
$\SheafHom_\mathcal{O}(\mathcal{F}, \mathcal{G})
= \SheafHom_\mathcal{O}(\mathcal{F}^\#, \mathcal{G})$, ce qui ramène
au cas où $\mathcal{F}$ et $\mathcal{G}$ sont tous deux des
faisceaux. Le résultat pour les faisceaux d'ensembles est Sites, Lemme
\ref{sites-lemma-glue-maps}.
\end{proof}

\begin{lemma}
\label{sites-modules-lemma-internal-hom-restriction}
Soit $(\mathcal{C}, \mathcal{O})$ un site annelé.
Soient $\mathcal{F}, \mathcal{G}$ des faisceaux de $\mathcal{O}$-modules.
Alors la formation de $\SheafHom_\mathcal{O}(\mathcal{F}, \mathcal{G})$
commute à la restriction à $U$, pour $U \in \Ob(\mathcal{C})$.
\end{lemma}

\begin{proof}
Cela résulte immédiatement de la définition.
\end{proof}

\begin{remark}
\label{sites-modules-remark-pullback-internal-hom}
Soit $f : (\mathcal{C}, \mathcal{O}_\mathcal{C}) \to
(\mathcal{D}, \mathcal{O}_\mathcal{D})$ un morphisme de sites annelés.
Soient $\mathcal{F}, \mathcal{G}$ des faisceaux de
$\mathcal{O}_\mathcal{D}$-modules. Il existe une application canonique
$$
f^*\SheafHom_{\mathcal{O}_\mathcal{D}}(\mathcal{F}, \mathcal{G})
\longrightarrow
\SheafHom_{\mathcal{O}_\mathcal{C}}(f^*\mathcal{F}, f^*\mathcal{G})
$$
Cette application est en effet adjointe à l'application
$$
\SheafHom_{\mathcal{O}_\mathcal{D}}(\mathcal{F}, \mathcal{G})
\longrightarrow
f_*\SheafHom_{\mathcal{O}_\mathcal{C}}(f^*\mathcal{F}, f^*\mathcal{G})
$$
définie comme suit. Supposons que $f$ soit donné par le foncteur continu
$u : \mathcal{D} \to \mathcal{C}$. Sur les sections au-dessus de
$V \in \Ob(\mathcal{D})$, on utilise l'application
\begin{align*}
\Gamma(V, \SheafHom_{\mathcal{O}_\mathcal{D}}(\mathcal{F}, \mathcal{G}))
& =
\Hom_{\mathcal{O}_V}(\mathcal{F}|_V, \mathcal{G}|_V) \\
& \longrightarrow
\Hom_{\mathcal{O}_{u(V)}}(f^*\mathcal{F}|_{u(V)}, f^*\mathcal{G}|_{u(V)}) \\
& =
\Gamma(u(V), \SheafHom_{\mathcal{O}_\mathcal{C}}(f^*\mathcal{F},
f^*\mathcal{G})) \\
& =
\Gamma(V, f_*\SheafHom_{\mathcal{O}_\mathcal{C}}(f^*\mathcal{F},
f^*\mathcal{G}))
\end{align*}
où, pour la flèche, on utilise l'image inverse par le morphisme
$(\mathcal{C}/u(V), \mathcal{O}_{u(V)}) \to
(\mathcal{D}/V, \mathcal{O}_V)$ induit par $f$.
\end{remark}

\noindent
Dans la situation du Lemme \ref{sites-modules-lemma-internal-hom}, le morphisme
d'évaluation se factorise par le produit tensoriel de faisceaux de
modules
$$
\mathcal{F}
\otimes_\mathcal{O}
\SheafHom_\mathcal{O}(\mathcal{F}, \mathcal{G})
\longrightarrow
\mathcal{G}.
$$

\begin{lemma}
\label{sites-modules-lemma-internal-hom-commute-limits}
Hom interne et (co)limites.
Soit $\mathcal{C}$ une catégorie et soit $\mathcal{O}$ un préfaisceau
d'anneaux.
\begin{enumerate}
\item Pour tout préfaisceau de $\mathcal{O}$-modules $\mathcal{F}$, le foncteur
$$
\textit{PMod}(\mathcal{O}) \longrightarrow \textit{PMod}(\mathcal{O})
, \quad
\mathcal{G} \longmapsto \SheafHom_\mathcal{O}(\mathcal{F}, \mathcal{G})
$$
commute aux limites arbitraires.
\item Pour tout préfaisceau de $\mathcal{O}$-modules $\mathcal{G}$, le foncteur
$$
\textit{PMod}(\mathcal{O}) \longrightarrow \textit{PMod}(\mathcal{O})^{opp}
, \quad
\mathcal{F} \longmapsto \SheafHom_\mathcal{O}(\mathcal{F}, \mathcal{G})
$$
commute aux colimites arbitraires ; en formule,
$$
\SheafHom_\mathcal{O}(\colim_i \mathcal{F}_i, \mathcal{G})
=
\lim_i \SheafHom_\mathcal{O}(\mathcal{F}_i, \mathcal{G}).
$$
\end{enumerate}
Supposons que $\mathcal{C}$ soit un site et $\mathcal{O}$ un faisceau d'anneaux.
\begin{enumerate}
\item[(3)] Pour tout faisceau de $\mathcal{O}$-modules $\mathcal{F}$, le foncteur
$$
\textit{Mod}(\mathcal{O}) \longrightarrow \textit{Mod}(\mathcal{O})
, \quad
\mathcal{G} \longmapsto \SheafHom_\mathcal{O}(\mathcal{F}, \mathcal{G})
$$
commute aux limites arbitraires.
\item[(4)] Pour tout faisceau de $\mathcal{O}$-modules $\mathcal{G}$, le foncteur
$$
\textit{Mod}(\mathcal{O}) \longrightarrow \textit{Mod}(\mathcal{O})^{opp}
, \quad
\mathcal{F} \longmapsto \SheafHom_\mathcal{O}(\mathcal{F}, \mathcal{G})
$$
commute aux colimites arbitraires ; en formule,
$$
\SheafHom_\mathcal{O}(\colim_i \mathcal{F}_i, \mathcal{G})
=
\lim_i \SheafHom_\mathcal{O}(\mathcal{F}_i, \mathcal{G}).
$$
\end{enumerate}
\end{lemma}

\begin{proof}
Soit $\mathcal{I} \to \textit{PMod}(\mathcal{O})$,
$i \mapsto \mathcal{G}_i$, un diagramme. Soit $U$ un objet de la catégorie
$\mathcal{C}$. Comme $j_U^*$ est à la fois adjoint à gauche et à droite,
on a
$\lim_i j_U^*\mathcal{G}_i = j_U^* \lim_i \mathcal{G}_i$.
Par conséquent,
\begin{align*}
\SheafHom_\mathcal{O}(\mathcal{F}, \lim_i \mathcal{G}_i)(U)
& =
\Hom_{\mathcal{O}_U}(\mathcal{F}|_U, \lim_i \mathcal{G}_i|_U) \\
& =
\lim_i \Hom_{\mathcal{O}_U}(\mathcal{F}|_U, \mathcal{G}_i|_U) \\
& = \lim_i \SheafHom_\mathcal{O}(\mathcal{F}, \mathcal{G}_i)(U)
\end{align*}
par définition d'une limite. Cela prouve (1). La preuve de (2) est exactement
la même. L'assertion (3) résulte de (1), car la limite d'un diagramme de
faisceaux est la même que sa limite dans la catégorie des préfaisceaux.
Enfin, (4) résulte du fait que, dans la formule, on a
$$
\Mor_{\textit{Mod}(\mathcal{O})}(
\colim_i \mathcal{F}_i, \mathcal{G})
=
\Mor_{\textit{PMod}(\mathcal{O})}(
\colim^{PSh}_i \mathcal{F}_i, \mathcal{G})
$$
car la colimite $\colim_i \mathcal{F}_i$ est la faisceautisation de la
colimite $\colim^{PSh}_i \mathcal{F}_i$ dans
$\textit{PMod}(\mathcal{O})$. Ainsi (4) résulte de (2) (en utilisant à
nouveau la remarque précédente sur les limites).
\end{proof}

\begin{lemma}
\label{sites-modules-lemma-internal-hom-exact}
Soit $(\mathcal{C}, \mathcal{O})$ un site annelé.
Soient $\mathcal{F}$ et $\mathcal{G}$ des $\mathcal{O}$-modules.
\begin{enumerate}
\item Si $\mathcal{F}_2 \to \mathcal{F}_1 \to \mathcal{F} \to 0$ est une
suite exacte de $\mathcal{O}$-modules, alors
$$
0 \to
\SheafHom_\mathcal{O}(\mathcal{F}, \mathcal{G}) \to
\SheafHom_\mathcal{O}(\mathcal{F}_1, \mathcal{G}) \to
\SheafHom_\mathcal{O}(\mathcal{F}_2, \mathcal{G})
$$
est exacte.
\item Si $0 \to \mathcal{G} \to \mathcal{G}_1 \to \mathcal{G}_2$ est une
suite exacte de $\mathcal{O}$-modules, alors
$$
0 \to
\SheafHom_\mathcal{O}(\mathcal{F}, \mathcal{G}) \to
\SheafHom_\mathcal{O}(\mathcal{F}, \mathcal{G}_1) \to
\SheafHom_\mathcal{O}(\mathcal{F}, \mathcal{G}_2)
$$
est exacte.
\end{enumerate}
\end{lemma}

\begin{proof}
Cela résulte du Lemme \ref{sites-modules-lemma-internal-hom-commute-limits} et de
Homology, Lemme \ref{homology-lemma-exact-functor}.
\end{proof}

\begin{lemma}
\label{sites-modules-lemma-internal-hom-adjoint-tensor}
Soit $\mathcal{C}$ une catégorie. Soit $\mathcal{O}$ un préfaisceau
d'anneaux.
\begin{enumerate}
\item Soient $\mathcal{F}$, $\mathcal{G}$ et $\mathcal{H}$ des préfaisceaux
de $\mathcal{O}$-modules. Il existe un isomorphisme canonique
$$
\SheafHom_\mathcal{O}
(\mathcal{F} \otimes_{p, \mathcal{O}} \mathcal{G}, \mathcal{H})
\longrightarrow
\SheafHom_\mathcal{O}
(\mathcal{F}, \SheafHom_\mathcal{O}(\mathcal{G}, \mathcal{H}))
$$
fonctoriel en ses trois arguments (le Hom est celui des préfaisceaux dans les
trois termes). En particulier,
$$
\Mor_{\textit{PMod}(\mathcal{O})}(
\mathcal{F} \otimes_{p, \mathcal{O}} \mathcal{G}, \mathcal{H})
=
\Mor_{\textit{PMod}(\mathcal{O})}(
\mathcal{F}, \SheafHom_\mathcal{O}(\mathcal{G}, \mathcal{H}))
$$
\item
Supposons que $\mathcal{C}$ soit un site, que $\mathcal{O}$ soit un faisceau
d'anneaux et que $\mathcal{F}$, $\mathcal{G}$, $\mathcal{H}$ soient des
faisceaux de $\mathcal{O}$-modules. Il existe un isomorphisme canonique
$$
\SheafHom_\mathcal{O}
(\mathcal{F} \otimes_\mathcal{O} \mathcal{G}, \mathcal{H})
\longrightarrow
\SheafHom_\mathcal{O}
(\mathcal{F}, \SheafHom_\mathcal{O}(\mathcal{G}, \mathcal{H}))
$$
fonctoriel en ses trois arguments (le Hom est celui des faisceaux dans les
trois termes). En particulier,
$$
\Mor_{\textit{Mod}(\mathcal{O})}(
\mathcal{F} \otimes_\mathcal{O} \mathcal{G}, \mathcal{H})
=
\Mor_{\textit{Mod}(\mathcal{O})}(
\mathcal{F}, \SheafHom_\mathcal{O}(\mathcal{G}, \mathcal{H}))
$$
\end{enumerate}
\end{lemma}

\begin{proof}
Il s'agit de l'analogue de
Algèbre, Lemme \ref{algebra-lemma-hom-from-tensor-product}.
La preuve est la même et est omise.
\end{proof}

\begin{lemma}
\label{sites-modules-lemma-tensor-commute-colimits}
Produit tensoriel et colimites.
Soit $\mathcal{C}$ une catégorie et soit $\mathcal{O}$ un préfaisceau
d'anneaux.
\begin{enumerate}
\item Pour tout préfaisceau de $\mathcal{O}$-modules $\mathcal{F}$, le foncteur
$$
\textit{PMod}(\mathcal{O}) \longrightarrow \textit{PMod}(\mathcal{O})
, \quad
\mathcal{G} \longmapsto \mathcal{F} \otimes_{p, \mathcal{O}} \mathcal{G}
$$
commute aux colimites arbitraires.
\item
Supposons que $\mathcal{C}$ soit un site et $\mathcal{O}$ un faisceau
d'anneaux. Pour tout faisceau de $\mathcal{O}$-modules $\mathcal{F}$, le
foncteur
$$
\textit{Mod}(\mathcal{O}) \longrightarrow \textit{Mod}(\mathcal{O})
, \quad
\mathcal{G} \longmapsto \mathcal{F} \otimes_\mathcal{O} \mathcal{G}
$$
commute aux colimites arbitraires.
\end{enumerate}
\end{lemma}

\begin{proof}
Cela tient à ce que le produit tensoriel est adjoint au Hom interne d'après
le Lemme \ref{sites-modules-lemma-internal-hom-adjoint-tensor}.
Voir Catégories, Lemme \ref{categories-lemma-adjoint-exact}.
\end{proof}

\begin{lemma}
\label{sites-modules-lemma-adjoint-hom-restrict}
Soit $\mathcal{C}$ une catégorie, resp.\ un site.
Soit $\mathcal{O} \to \mathcal{O}'$ un morphisme de préfaisceaux, resp.\ de
faisceaux d'anneaux. Alors
$$
\Hom_\mathcal{O}(\mathcal{G}, \mathcal{F}) =
\Hom_{\mathcal{O}'}(\mathcal{G},
\SheafHom_\mathcal{O}(\mathcal{O}', \mathcal{F}))
$$
pour tout $\mathcal{O}'$-module $\mathcal{G}$ et tout $\mathcal{O}$-module
$\mathcal{F}$.
\end{lemma}

\begin{proof}
Il s'agit de l'analogue de
Algèbre, Lemme \ref{algebra-lemma-adjoint-hom-restrict}.
La preuve est la même et est omise.
\end{proof}

\begin{lemma}
\label{sites-modules-lemma-j-shriek-and-tensor}
Soit $(\mathcal{C}, \mathcal{O})$ un site annelé.
Soit $U \in \Ob(\mathcal{C})$.
Pour $\mathcal{G}$ dans $\textit{Mod}(\mathcal{O}_U)$ et $\mathcal{F}$ dans
$\textit{Mod}(\mathcal{O})$, on a
$j_{U!}\mathcal{G} \otimes_\mathcal{O} \mathcal{F} =
j_{U!}(\mathcal{G} \otimes_{\mathcal{O}_U} \mathcal{F}|_U)$.
\end{lemma}

\begin{proof}
Soit $\mathcal{H}$ un objet de $\textit{Mod}(\mathcal{O})$. Alors
\begin{align*}
\Hom_\mathcal{O}(j_{U!}(\mathcal{G} \otimes_{\mathcal{O}_U} \mathcal{F}|_U),
\mathcal{H})
& =
\Hom_{\mathcal{O}_U}(\mathcal{G} \otimes_{\mathcal{O}_U} \mathcal{F}|_U,
\mathcal{H}|_U) \\
& =
\Hom_{\mathcal{O}_U}(\mathcal{G},
\SheafHom_{\mathcal{O}_U}(\mathcal{F}|_U, \mathcal{H}|_U)) \\
& =
\Hom_{\mathcal{O}_U}(\mathcal{G},
\SheafHom_\mathcal{O}(\mathcal{F}, \mathcal{H})|_U) \\
& =
\Hom_\mathcal{O}(j_{U!}\mathcal{G},
\SheafHom_\mathcal{O}(\mathcal{F}, \mathcal{H})) \\
& =
\Hom_\mathcal{O}(j_{U!}\mathcal{G} \otimes_\mathcal{O} \mathcal{F},
\mathcal{H})
\end{align*}
La première égalité vient de ce que $j_{U!}$ est adjoint à gauche de la
restriction des modules. La deuxième résulte du
Lemme \ref{sites-modules-lemma-internal-hom-adjoint-tensor}.
La troisième résulte du Lemme \ref{sites-modules-lemma-internal-hom-restriction}.
La quatrième vient à nouveau de ce que $j_{U!}$ est adjoint à gauche de la
restriction des modules. La cinquième résulte du
Lemme \ref{sites-modules-lemma-internal-hom-adjoint-tensor}.
Le lemme découle de ces égalités et du lemme de Yoneda.
\end{proof}

\begin{remark}
\label{sites-modules-remark-j-shriek-tensor}
Soit $\mathcal{C}$ un site. Soit $\mathcal{F}$ un faisceau d'ensembles sur
$\mathcal{C}$ et considérons le morphisme de localisation
$j : \Sh(\mathcal{C})/\mathcal{F} \to \Sh(\mathcal{C})$.
Voir Sites, Définition \ref{sites-definition-localize-topos}.
Nous affirmons que (a) $j_!\mathbf{Z} = \mathbf{Z}_\mathcal{F}^\#$ et que
(b) $j_!(j^{-1}\mathcal{H}) = j_!\mathbf{Z} \otimes_\mathbf{Z} \mathcal{H}$
pour tout faisceau abélien $\mathcal{H}$ sur $\mathcal{C}$.
Soit $\mathcal{G}$ un faisceau abélien sur $\mathcal{C}$.
L'assertion (a) résulte du lemme de Yoneda, car
$$
\Hom(j_!\mathbf{Z}, \mathcal{G}) =
\Hom(\mathbf{Z}, j^{-1}\mathcal{G}) =
\Hom(\mathbf{Z}_\mathcal{F}^\#, \mathcal{G})
$$
où la seconde égalité vaut parce que ses deux membres s'identifient à
l'ensemble des applications $\mathcal{F} \to \mathcal{G}$, considéré comme
groupe abélien. Pour (b), on utilise le lemme de Yoneda et
\begin{align*}
\Hom(j_!(j^{-1}\mathcal{H}), \mathcal{G})
& =
\Hom(j^{-1}\mathcal{H}, j^{-1}\mathcal{G}) \\
& =
\Hom(\mathbf{Z}, \SheafHom(j^{-1}\mathcal{H}, j^{-1}\mathcal{G})) \\
& =
\Hom(\mathbf{Z}, j^{-1}\SheafHom(\mathcal{H}, \mathcal{G})) \\
& =
\Hom(j_!\mathbf{Z}, \SheafHom(\mathcal{H}, \mathcal{G})) \\
& =
\Hom(j_!\mathbf{Z} \otimes_\mathbf{Z} \mathcal{H}, \mathcal{G})
\end{align*}
On utilise ici l'adjonction, le fait que la formation de $\SheafHom$ commute à
la localisation, et le Lemme \ref{sites-modules-lemma-internal-hom-adjoint-tensor}.
\end{remark}

\begin{lemma}
\label{sites-modules-lemma-hom-finite-presentation-colimit}
Soit $(\mathcal{C}, \mathcal{O})$ un site annelé.
Soit $\mathcal{F}$ un $\mathcal{O}$-module de présentation finie. Soit
$\mathcal{G} = \colim_{\lambda \in \Lambda} \mathcal{G}_\lambda$ une
colimite filtrante de $\mathcal{O}$-modules. Alors l'application canonique
$$
\colim_\lambda \SheafHom_\mathcal{O}(\mathcal{F}, \mathcal{G}_\lambda)
\longrightarrow
\SheafHom_\mathcal{O}(\mathcal{F}, \mathcal{G})
$$
est un isomorphisme.
\end{lemma}

\begin{proof}
Il suffit de montrer que la flèche devient un isomorphisme après restriction
à $U$, pour tout $U$ de $\mathcal{C}$. La formation des colimites de faisceaux de
modules et celle du Hom interne commutent toutes deux à la restriction à
$U$. Voir par exemple les Lemmes \ref{sites-modules-lemma-exactness-pushforward-pullback}
et \ref{sites-modules-lemma-internal-hom-restriction}. Fixons $U$. Étant donné un
recouvrement $\{U_i \to U\}_{i \in I}$, il suffit de prouver que la
restriction à chaque $U_i$ est un isomorphisme. On peut donc supposer que
$\mathcal{F}$ admet une présentation globale
$$
\bigoplus\nolimits_{j = 1, \ldots, m}
\mathcal{O}
\longrightarrow
\bigoplus\nolimits_{i = 1, \ldots, n}
\mathcal{O}
\to
\mathcal{F}
\to
0
$$
Le foncteur $\SheafHom_\mathcal{O}(-, -)$ commute aux sommes directes finies
en chacune des variables, et $\SheafHom_\mathcal{O}(\mathcal{O}, -)$ est le
foncteur identité. Avec le Lemme \ref{sites-modules-lemma-internal-hom-exact}, cela donne
une suite exacte
$$
0 \to
\SheafHom_\mathcal{O}(\mathcal{F}, \mathcal{G}) \to
\bigoplus\nolimits_{i = 1, \ldots, n} \mathcal{G} \to
\bigoplus\nolimits_{j = 1, \ldots, m} \mathcal{G}
$$
Comme les colimites filtrantes sont exactes dans
$\textit{Mod}(\mathcal{O})$ d'après le Lemme \ref{sites-modules-lemma-limits-colimits},
la ligne supérieure du diagramme commutatif suivant est aussi exacte
$$
\xymatrix{
0 \ar[r] &
\colim_\lambda \SheafHom_\mathcal{O}(\mathcal{F}, \mathcal{G}_\lambda)
\ar[r] \ar[d] &
\colim_\lambda \bigoplus\nolimits_{i = 1, \ldots, n} \mathcal{G}_\lambda
\ar[r] \ar[d] &
\colim_\lambda \bigoplus\nolimits_{j = 1, \ldots, m} \mathcal{G}_\lambda
\ar[d] \\
0 \ar[r] &
\SheafHom_\mathcal{O}(\mathcal{F}, \mathcal{G}) \ar[r] &
\bigoplus\nolimits_{i = 1, \ldots, n} \mathcal{G} \ar[r] &
\bigoplus\nolimits_{j = 1, \ldots, m} \mathcal{G}
}
$$
On conclut puisque les deux flèches verticales de droite sont des
isomorphismes.
\end{proof}

\begin{lemma}
\label{sites-modules-lemma-finite-presentation-quasi-compact-colimit}
Soit $(\mathcal{C}, \mathcal{O})$ un site annelé.
Soit $\mathcal{G} = \colim_{\lambda \in \Lambda} \mathcal{G}_\lambda$ une
colimite filtrante de $\mathcal{O}$-modules.
Soit $\mathcal{F}$ un $\mathcal{O}$-module de présentation finie.
Alors on a
$$
\colim_\lambda \Hom_\mathcal{O}(\mathcal{F}, \mathcal{G}_\lambda)
=
\Hom_\mathcal{O}(\mathcal{F}, \mathcal{G}).
$$
si les hypothèses de Sites, Lemme
\ref{sites-lemma-directed-colimits-global-sections}, partie (4), sont
satisfaites pour le site $\mathcal{C}$ ; voir Sites, Remarque
\ref{sites-remark-stronger-conditions}.
\end{lemma}

\begin{proof}
Posons $\mathcal{H} =
\SheafHom_\mathcal{O}(\mathcal{F}, \colim \mathcal{G}_\lambda)$
et $\mathcal{H}_\lambda =
\SheafHom_\mathcal{O}(\mathcal{F}, \mathcal{G}_\lambda)$.
Rappelons que
$$
\Hom_\mathcal{O}(\mathcal{F}, \mathcal{G}) = \Gamma(\mathcal{C}, \mathcal{H})
\quad\text{et}\quad
\Hom_\mathcal{O}(\mathcal{F}, \mathcal{G}_\lambda) =
\Gamma(\mathcal{C}, \mathcal{H}_\lambda)
$$
par construction. Le Lemme \ref{sites-modules-lemma-hom-finite-presentation-colimit} donne
$\mathcal{H} = \colim \mathcal{H}_\lambda$. Le lemme résulte donc de Sites,
Lemme \ref{sites-lemma-directed-colimits-global-sections}.
\end{proof}
\section{Modules plats}

\label{sites-modules-section-flat}

\noindent
On peut définir les modules plats exactement comme dans le cas des modules sur un anneau.

\begin{definition}

\label{sites-modules-definition-flat}

Soit $\mathcal{C}$ une catégorie et soit $\mathcal{O}$ un préfaisceau d'anneaux.

\begin{enumerate}

\item Un préfaisceau $\mathcal{F}$ de $\mathcal{O}$-modules est dit
{\it plat} si le foncteur
$$
\textit{PMod}(\mathcal{O})
\longrightarrow
\textit{PMod}(\mathcal{O}), \quad
\mathcal{G} \mapsto \mathcal{G} \otimes_{p, \mathcal{O}} \mathcal{F}
$$
est exact.
\item Un morphisme $\mathcal{O} \to \mathcal{O}'$ de préfaisceaux d'anneaux
est dit {\it plat} si $\mathcal{O}'$ est plat comme préfaisceau de
$\mathcal{O}$-modules.
\item Si $\mathcal{C}$ est un site, si $\mathcal{O}$ est un faisceau
d'anneaux et si $\mathcal{F}$ est un faisceau de $\mathcal{O}$-modules,
on dit que $\mathcal{F}$ est {\it plat} si le foncteur
$$
\textit{Mod}(\mathcal{O})
\longrightarrow
\textit{Mod}(\mathcal{O}), \quad
\mathcal{G} \mapsto \mathcal{G} \otimes_\mathcal{O} \mathcal{F}
$$
est exact.
\item Un morphisme $\mathcal{O} \to \mathcal{O}'$ de faisceaux d'anneaux
sur un site est dit {\it plat} si $\mathcal{O}'$ est plat comme faisceau de
$\mathcal{O}$-modules.

\end{enumerate}

\end{definition}

\noindent
La notion de module plat ou de morphisme plat d'anneaux est intrinsèque
(section \ref{sites-modules-section-intrinsic}).

\begin{lemma}

\label{sites-modules-lemma-flatness-presheaves}
Soit $\mathcal{C}$ une catégorie, soit $\mathcal{O}$ un préfaisceau
d'anneaux et soit $\mathcal{F}$ un préfaisceau de $\mathcal{O}$-modules.
Si chaque $\mathcal{F}(U)$ est un $\mathcal{O}(U)$-module plat,
alors $\mathcal{F}$ est plat.
\end{lemma}

\begin{proof}
Cela résulte immédiatement des définitions.
\end{proof}

\begin{lemma}

\label{sites-modules-lemma-flatness-sheafification}
Soit $\mathcal{C}$ un site, soit $\mathcal{O}$ un préfaisceau d'anneaux
et soit $\mathcal{F}$ un préfaisceau de $\mathcal{O}$-modules.
Si $\mathcal{F}$ est plat sur $\mathcal{O}$, alors $\mathcal{F}^\#$
est plat sur $\mathcal{O}^\#$.
\end{lemma}

\begin{proof}

Démonstration omise. (Indication : la faisceautisation est exacte.)

\end{proof}

\begin{lemma}

\label{sites-modules-lemma-flatness-sheafification-refined}
Soit $\mathcal{C}$ un site, soit $\mathcal{O}$ un préfaisceau d'anneaux
et soit $\mathcal{F}$ un préfaisceau de $\mathcal{O}$-modules. Supposons
que tout objet $U$ de $\mathcal{C}$ admette un recouvrement
$\{U_i \to U\}_{i \in I}$ tel que $\mathcal{F}(U_i)$ soit un
$\mathcal{O}(U_i)$-module plat. Alors $\mathcal{F}^\#$ est un
$\mathcal{O}^\#$-module plat.
\end{lemma}

\begin{proof}

Soit $\mathcal{G} \subset \mathcal{G}'$ une inclusion de
$\mathcal{O}^\#$-modules. Il faut montrer que
$$
\mathcal{G} \otimes_{\mathcal{O}^\#} \mathcal{F}^\#
\longrightarrow
\mathcal{G}' \otimes_{\mathcal{O}^\#} \mathcal{F}^\#
$$
est injective. D'après le lemme \ref{sites-modules-lemma-sheafification-tensor},
la source de cette flèche est le faisceau associé au préfaisceau
$\mathcal{G} \otimes_{p, \mathcal{O}} \mathcal{F}$,
et il en va de même pour la cible. Si $U$ est un objet de $\mathcal{C}$ tel que
$\mathcal{F}(U)$ soit un $\mathcal{O}(U)$-module plat, alors
$$
(\mathcal{G} \otimes_{p, \mathcal{O}} \mathcal{F})(U) =
\mathcal{G}(U) \otimes_{\mathcal{O}(U)} \mathcal{F}(U)
\longrightarrow
\mathcal{G}'(U) \otimes_{\mathcal{O}(U)} \mathcal{F}(U) =
(\mathcal{G}' \otimes_{p, \mathcal{O}} \mathcal{F})(U)
$$
est injective. Il suffit donc d'établir l'assertion suivante : si un morphisme
de préfaisceaux $f : \mathcal{H} \to \mathcal{H}'$ sur $\mathcal{C}$ est tel
que tout objet $U$ de $\mathcal{C}$ admette un recouvrement
$\{U_i \to U\}_{i \in I}$ pour lequel
$\mathcal{H}(U_i) \to \mathcal{H}'(U_i)$ est injectif, alors $f^\#$ est
injectif. Nous laissons cette vérification au lecteur.

\end{proof}

\begin{lemma}

\label{sites-modules-lemma-colimits-flat}
Colimites et produit tensoriel.
\begin{enumerate}
\item Une colimite filtrante de préfaisceaux de modules plats est plate.
Une somme directe de préfaisceaux de modules plats est plate.
\item Une colimite filtrante de faisceaux de modules plats est plate.
Une somme directe de faisceaux de modules plats est plate.
\end{enumerate}

\end{lemma}

\begin{proof}

La partie (1) découle du lemme \ref{sites-modules-lemma-tensor-commute-colimits} et
d'Algèbre, lemme \ref{algebra-lemma-directed-colimit-exact}, en évaluant
sur les objets. Pour la partie (2), on utilise le lemme
\ref{sites-modules-lemma-tensor-commute-colimits} et le fait qu'une colimite filtrante de
complexes exacts est exacte ; cela repose sur l'exactitude de la
faisceautisation et sur sa commutation aux colimites. Certains détails sont
omis.

\end{proof}

\begin{lemma}

\label{sites-modules-lemma-restriction-flat}
Soit $(\mathcal{C}, \mathcal{O})$ un site annelé et soit $U$ un objet de
$\mathcal{C}$. Si $\mathcal{F}$ est un $\mathcal{O}$-module plat, alors
$\mathcal{F}|_U$ est un $\mathcal{O}_U$-module plat.
\end{lemma}

\begin{proof}
Soit $\mathcal{G}_1 \to \mathcal{G}_2 \to \mathcal{G}_3$ un complexe exact
de $\mathcal{O}_U$-modules. Comme $j_{U!}$ est exact (lemme
\ref{sites-modules-lemma-extension-by-zero-exact}) et que $\mathcal{F}$ est plat comme
$\mathcal{O}$-module, le complexe formé des modules
$$
j_{U!}(\mathcal{G}_i \otimes_{\mathcal{O}_U} \mathcal{F}|_U) =
j_{U!}\mathcal{G}_i \otimes_\mathcal{O} \mathcal{F}
$$
est exact d'après le lemme \ref{sites-modules-lemma-j-shriek-and-tensor}. On en déduit que
$\mathcal{G}_1 \otimes_{\mathcal{O}_U} \mathcal{F}|_U \to
\mathcal{G}_2 \otimes_{\mathcal{O}_U} \mathcal{F}|_U \to
\mathcal{G}_3 \otimes_{\mathcal{O}_U} \mathcal{F}|_U$ est exact par le lemme \ref{sites-modules-lemma-j-shriek-reflects-exactness}.
\end{proof}

\begin{lemma}

\label{sites-modules-lemma-j-shriek-flat}

Soient $\mathcal{C}$ une catégorie, $\mathcal{O}$ un préfaisceau d'anneaux
et $U$ un objet de $\mathcal{C}$. Considérons le foncteur
$j_U : \mathcal{C}/U \to \mathcal{C}$.

\begin{enumerate}

\item Le préfaisceau de $\mathcal{O}$-modules
$j_{U!}\mathcal{O}_U$ (voir la remarque
\ref{sites-modules-remark-localize-presheaves}) est plat.
\item Si $\mathcal{C}$ est un site, $\mathcal{O}$ est un faisceau d'anneaux,
$j_{U!}\mathcal{O}_U$ est un faisceau plat de $\mathcal{O}$-modules.

\end{enumerate}

\end{lemma}

\begin{proof}
Preuve de (1). La discussion de la remarque
\ref{sites-modules-remark-localize-presheaves} donne
$$
j_{U!}\mathcal{O}_U(V)
=
\bigoplus\nolimits_{\varphi \in \Mor_\mathcal{C}(V, U)}
\mathcal{O}(V)
$$
qui est un $\mathcal{O}(V)$-module plat. La partie (1) découle donc du lemme
\ref{sites-modules-lemma-flatness-presheaves}. La partie (2) résulte ensuite de l'égalité
$j_{U!}\mathcal{O}_U = (j_{U!}\mathcal{O}_U)^\#$ (le premier $j_{U!}$
étant pris sur les faisceaux et le second sur les préfaisceaux) et du lemme
\ref{sites-modules-lemma-flatness-sheafification}.
\end{proof}

\begin{lemma}

\label{sites-modules-lemma-module-quotient-flat}

Soient $\mathcal{C}$ une catégorie et $\mathcal{O}$ un préfaisceau d'anneaux.

\begin{enumerate}

\item Tout préfaisceau de $\mathcal{O}$-modules est quotient d'une somme
directe $\bigoplus j_{U_i!}\mathcal{O}_{U_i}$.
\item Tout préfaisceau de $\mathcal{O}$-modules est quotient d'un préfaisceau
plat de $\mathcal{O}$-modules.
\item Si $\mathcal{C}$ est un site et $\mathcal{O}$ un faisceau d'anneaux,
tout faisceau de $\mathcal{O}$-modules est quotient d'une somme directe
$\bigoplus j_{U_i!}\mathcal{O}_{U_i}$.
\item Si $\mathcal{C}$ est un site et $\mathcal{O}$ un faisceau d'anneaux,
tout faisceau de $\mathcal{O}$-modules est quotient d'un faisceau plat de
$\mathcal{O}$-modules.

\end{enumerate}

\end{lemma}

\begin{proof}
Preuve de (1). Pour tout objet $U$ de $\mathcal{C}$ et toute section
$s \in \mathcal{F}(U)$, on obtient un morphisme
$j_{U!}\mathcal{O}_U \to \mathcal{F}$, adjoint du morphisme
$\mathcal{O}_U \to \mathcal{F}|_U$, $1 \mapsto s$. L'application
$$
\bigoplus\nolimits_{(U, s)} j_{U!}\mathcal{O}_U
\longrightarrow
\mathcal{F}
$$
est surjective. Les lemmes \ref{sites-modules-lemma-colimits-flat} et
\ref{sites-modules-lemma-j-shriek-flat} montrent que sa source est plate, ce qui prouve
(2). Le cas des faisceaux s'en déduit soit par faisceautisation, soit en
répétant le même argument.
\end{proof}

\begin{lemma}

\label{sites-modules-lemma-flat-tor-zero}

Soient $\mathcal{C}$ une catégorie et $\mathcal{O}$ un préfaisceau d'anneaux.
$$
0 \to \mathcal{F}'' \to \mathcal{F}' \to \mathcal{F} \to 0
$$
une suite exacte courte de préfaisceaux de $\mathcal{O}$-modules, et soit
$\mathcal{G}$ un préfaisceau de $\mathcal{O}$-modules.

\begin{enumerate}

\item Si $\mathcal{F}$ est un préfaisceau plat de modules, alors la suite
$$
0 \to
\mathcal{F}'' \otimes_{p, \mathcal{O}} \mathcal{G} \to
\mathcal{F}' \otimes_{p, \mathcal{O}} \mathcal{G} \to
\mathcal{F} \otimes_{p, \mathcal{O}} \mathcal{G} \to 0
$$
est exacte.
\item Si $\mathcal{C}$ est un site, $\mathcal{O}$,
$\mathcal{F}$, $\mathcal{F}'$, $\mathcal{F}''$,
$\mathcal{G}$ sont des faisceaux et si $\mathcal{F}$ est plat comme faisceau
de modules, alors la suite
$$
0 \to
\mathcal{F}'' \otimes_\mathcal{O} \mathcal{G} \to
\mathcal{F}' \otimes_\mathcal{O} \mathcal{G} \to
\mathcal{F} \otimes_\mathcal{O} \mathcal{G} \to 0
$$
est exacte.

\end{enumerate}

\end{lemma}

\begin{proof}
Choisissons un préfaisceau plat de $\mathcal{O}$-modules $\mathcal{G}'$
surjectant sur $\mathcal{G}$ ; c'est possible d'après le lemme
\ref{sites-modules-lemma-module-quotient-flat}. Posons
$\mathcal{G}'' = \Ker(\mathcal{G}' \to \mathcal{G})$. Le résultat s'obtient
en appliquant le lemme du serpent au diagramme suivant :
$$
\begin{matrix}
 & & 0 & & 0 & & 0 & & \\
 & & \uparrow & & \uparrow & & \uparrow & & \\
 & & \mathcal{F}'' \otimes_{p, \mathcal{O}} \mathcal{G} & \to &
     \mathcal{F}' \otimes_{p, \mathcal{O}} \mathcal{G} & \to &
     \mathcal{F} \otimes_{p, \mathcal{O}} \mathcal{G} & \to & 0 \\
 & & \uparrow & & \uparrow & & \uparrow & & \\
0 & \to & \mathcal{F}'' \otimes_{p, \mathcal{O}} \mathcal{G}' & \to &
          \mathcal{F}' \otimes_{p, \mathcal{O}} \mathcal{G}' & \to &
          \mathcal{F} \otimes_{p, \mathcal{O}} \mathcal{G}' & \to & 0 \\
 & & \uparrow & & \uparrow & & \uparrow & & \\
 & & \mathcal{F}'' \otimes_{p, \mathcal{O}} \mathcal{G}'' & \to &
     \mathcal{F}' \otimes_{p, \mathcal{O}} \mathcal{G}'' & \to &
     \mathcal{F} \otimes_{p, \mathcal{O}} \mathcal{G}'' & \to & 0 \\
 & & & & & & \uparrow & & \\
 & & & & & & 0 & &
\end{matrix}
$$
dont les lignes et les colonnes sont exactes. La ligne médiane est exacte
parce que le produit tensoriel avec le module plat $\mathcal{G}'$ est exact.
La démonstration dans le cas des faisceaux est identique.
\end{proof}

\begin{lemma}

\label{sites-modules-lemma-flat-ses}

Soient $\mathcal{C}$ une catégorie et $\mathcal{O}$ un préfaisceau d'anneaux.
$$
0 \to
\mathcal{F}_2 \to
\mathcal{F}_1 \to
\mathcal{F}_0 \to 0
$$
une suite exacte courte de préfaisceaux de $\mathcal{O}$-modules.

\begin{enumerate}

\item Si $\mathcal{F}_2$ et $\mathcal{F}_0$ sont plats, alors
$\mathcal{F}_1$ l'est aussi.
\item Si $\mathcal{F}_1$ et $\mathcal{F}_0$ sont plats, alors
$\mathcal{F}_2$ l'est aussi.

\end{enumerate}

Si $\mathcal{C}$ est un site et $\mathcal{O}$ un faisceau d'anneaux, le même
résultat vaut dans $\textit{Mod}(\mathcal{O})$.

\end{lemma}

\begin{proof}
Soit $\mathcal{G}^\bullet$ un complexe exact quelconque de préfaisceaux de
$\mathcal{O}$-modules. Supposons $\mathcal{F}_0$ plat. Le lemme
\ref{sites-modules-lemma-flat-tor-zero} montre que
$$
0 \to
\mathcal{G}^\bullet \otimes_{p, \mathcal{O}} \mathcal{F}_2 \to
\mathcal{G}^\bullet \otimes_{p, \mathcal{O}} \mathcal{F}_1 \to
\mathcal{G}^\bullet \otimes_{p, \mathcal{O}} \mathcal{F}_0 \to 0
$$
est une suite exacte courte de complexes de préfaisceaux de
$\mathcal{O}$-modules. Les assertions (1) et (2) découlent donc du lemme du
serpent. Le cas des faisceaux de modules se démontre de la même manière.
\end{proof}

\begin{lemma}

\label{sites-modules-lemma-flat-resolution-of-flat}
Soient $\mathcal{C}$ une catégorie et $\mathcal{O}$ un préfaisceau d'anneaux.
Soit
$$
\ldots \to
\mathcal{F}_2 \to
\mathcal{F}_1 \to
\mathcal{F}_0 \to
\mathcal{Q} \to 0
$$
un complexe exact de préfaisceaux de $\mathcal{O}$-modules. Si
$\mathcal{Q}$ et tous les $\mathcal{F}_i$ sont des $\mathcal{O}$-modules
plats, alors, pour tout préfaisceau $\mathcal{G}$ de $\mathcal{O}$-modules,
le complexe
$$
\ldots \to
\mathcal{F}_2 \otimes_{p, \mathcal{O}} \mathcal{G} \to
\mathcal{F}_1 \otimes_{p, \mathcal{O}} \mathcal{G} \to
\mathcal{F}_0 \otimes_{p, \mathcal{O}} \mathcal{G} \to
\mathcal{Q} \otimes_{p, \mathcal{O}} \mathcal{G} \to 0
$$
est également exact. Si $\mathcal{C}$ est un site et $\mathcal{O}$ un
faisceau d'anneaux, le même résultat vaut dans
$\textit{Mod}(\mathcal{O})$.
\end{lemma}

\begin{proof}

Cela découle du lemme \ref{sites-modules-lemma-flat-tor-zero}, en décomposant le complexe
en suites exactes courtes et en appliquant le lemme \ref{sites-modules-lemma-flat-ses}
pour démontrer par récurrence que
$\Im(\mathcal{F}_{i + 1} \to \mathcal{F}_i)$ est plat.

\end{proof}

\begin{lemma}

\label{sites-modules-lemma-tensor-flats}
Soit $(\mathcal{C}, \mathcal{O})$ un site annelé. Si $\mathcal{G}$ et
$\mathcal{F}$ sont des $\mathcal{O}$-modules plats, alors
$\mathcal{G} \otimes_\mathcal{O} \mathcal{F}$ est un
$\mathcal{O}$-module plat.
\end{lemma}

\begin{proof}

C'est vrai parce que
$$
(\mathcal{G} \otimes_\mathcal{O} \mathcal{F}) \otimes_\mathcal{O} \mathcal{H}
=
\mathcal{G} \otimes_\mathcal{O} (\mathcal{F} \otimes_\mathcal{O} \mathcal{H})
$$
et une composition de foncteurs exacts est exacte.

\end{proof}

\begin{lemma}

\label{sites-modules-lemma-flat-change-of-rings}
Soit $\mathcal{O}_1 \to \mathcal{O}_2$ un morphisme de faisceaux d'anneaux
sur un site $\mathcal{C}$. Si $\mathcal{G}$ est un
$\mathcal{O}_1$-module plat, alors
$\mathcal{G} \otimes_{\mathcal{O}_1} \mathcal{O}_2$ est un
$\mathcal{O}_2$-module plat.
\end{lemma}

\begin{proof}

C'est vrai parce que
$$
(\mathcal{G} \otimes_{\mathcal{O}_1} \mathcal{O}_2)
\otimes_{\mathcal{O}_2} \mathcal{H}
=
\mathcal{G} \otimes_{\mathcal{O}_1} \mathcal{H}
$$
(par exemple comme faisceaux de groupes abéliens).

\end{proof}

\noindent
Le lemme suivant est l'analogue du critère d'équation de la planéité (Algèbre, lemme \ref{algebra-lemma-flat-eq}).

\begin{lemma}

\label{sites-modules-lemma-flat-eq}

Soit $(\mathcal{C}, \mathcal{O})$ un site annelé et soit $\mathcal{F}$ un
$\mathcal{O}$-module. Les conditions suivantes sont équivalentes :

\begin{enumerate}

\item $\mathcal{F}$ est un $\mathcal{O}$-module plat.
\item Soit $U$ un objet de $\mathcal{C}$ et soit
$$
\mathcal{O}_U \xrightarrow{(f_1, \ldots, f_n)}
\mathcal{O}_U^{\oplus n} \xrightarrow{(s_1, \ldots, s_n)}
\mathcal{F}|_U
$$
un complexe de $\mathcal{O}_U$-modules. Il existe alors un recouvrement
$\{U_i \to U\}$ et, pour chaque $i$, une factorisation
$$
\mathcal{O}_{U_i}^{\oplus n}
\xrightarrow{B_i}
\mathcal{O}_{U_i}^{\oplus l_i} \xrightarrow{(t_{i1}, \ldots, t_{il_i})}
\mathcal{F}|_{U_i}
$$
de $(s_1, \ldots, s_n)|_{U_i}$ telle que
$B_i \circ (f_1, \ldots, f_n)|_{U_i} = 0$.
\item Soit $U$ un objet de $\mathcal{C}$ et soit
$$
\mathcal{O}_U^{\oplus m} \xrightarrow{A}
\mathcal{O}_U^{\oplus n} \xrightarrow{(s_1, \ldots, s_n)}
\mathcal{F}|_U
$$
un complexe de $\mathcal{O}_U$-modules. Il existe alors un recouvrement
$\{U_i \to U\}$ et, pour chaque $i$, une factorisation
$$
\mathcal{O}_{U_i}^{\oplus n}
\xrightarrow{B_i}
\mathcal{O}_{U_i}^{\oplus l_i} \xrightarrow{(t_{i1}, \ldots, t_{il_i})}
\mathcal{F}|_{U_i}
$$
de $(s_1, \ldots, s_n)|_{U_i}$ telle que
$B_i \circ A|_{U_i} = 0$.

\end{enumerate}

\end{lemma}

\begin{proof}

Supposons (1). Soit $\mathcal{I} \subset \mathcal{O}_U$ le faisceau
d'idéaux engendré par $f_1, \ldots, f_n$. Alors
$\sum f_j \otimes s_j$ est une section de
$\mathcal{I} \otimes_{\mathcal{O}_U} \mathcal{F}|_U$
dont l'image dans $\mathcal{F}|_U$ est nulle. Comme $\mathcal{F}|_U$ est
plat (lemme \ref{sites-modules-lemma-restriction-flat}), l'application
$\mathcal{I} \otimes_{\mathcal{O}_U} \mathcal{F}|_U \to \mathcal{F}|_U$
est injective. $\mathcal{I} \otimes_{\mathcal{O}_U} \mathcal{F}|_U$
est le faisceau associé au préfaisceau produit tensoriel, il existe un
recouvrement $\{U_i \to U\}$ tel que
$\sum f_j|_{U_i} \otimes s_j|_{U_i}$ soit nul dans
$\mathcal{I}(U_i) \otimes_{\mathcal{O}(U_i)} \mathcal{F}(U_i)$. En
explicitant les définitions à l'aide d'Algèbre, lemme
\ref{algebra-lemma-relations}, on trouve
$t_{i1}, \ldots, t_{i l_i} \in \mathcal{F}(U_i)$ et
$a_{ijk} \in \mathcal{O}(U_i)$
de sorte que $\sum_j a_{ijk}f_j|_{U_i} = 0$ et
$s_j|_{U_i} = \sum_k a_{ijk}t_{ik}$. Ainsi (2) tient.

\medskip\noindent

Supposons (2), et donnons-nous $U$, $n$, $m$, $A$ et
$s_1, \ldots, s_n$ comme en (3). La matrice $A$ possède $m$ colonnes.
Nous allons démontrer l'assertion (3) par récurrence sur $m$. Pour
$m = 0$, on peut prendre le recouvrement $\{U \to U\}$ et la factorisation
$\mathcal{O}_U^n \to \mathcal{O}_U^n \to \mathcal{F}|_U$, la première
flèche étant l'identité et la seconde $(s_1, \ldots, s_n)$. Supposons
$m > 0$. Soit $(f_1, \ldots, f_n)$ la dernière colonne de $A$ et appliquons
(2). On obtient les nouveaux diagrammes
$$
\mathcal{O}_{U_i}^{\oplus m} \xrightarrow{B_i \circ A|_{U_i}}
\mathcal{O}_{U_i}^{\oplus l_i} \xrightarrow{(t_{i1}, \ldots, t_{il_i})}
\mathcal{F}|_{U_i}
$$
dont la dernière colonne de $A_i = B_i \circ A|_{U_i}$ est nulle.
On peut donc appliquer l'hypothèse de récurrence à
$U_i$, $l_i$, $m - 1$, la matrice composée de la première $m - 1$ colonnes $A_i$, $t_{i1}, \ldots, t_{il_i}$
pour obtenir des recouvrements $\{U_{ij} \to U_i\}$ et des factorisations
$$
\mathcal{O}_{U_{ij}}^{\oplus l_i}
\xrightarrow{C_{ij}}
\mathcal{O}_{U_{ij}}^{\oplus k_{ij}}
\xrightarrow{(v_{ij1}, \ldots, v_{ij k_{ij}})}
\mathcal{F}|_{U_{ij}}
$$
de $(t_{i1}, \ldots, t_{il_i})|_{U_{ij}}$ telles que
$C_{ij} \circ B_i|_{U_{ij}} \circ A|_{U_{ij}} = 0$. Alors
$\{U_{ij} \to U\}$ est un recouvrement, et les factorisations cherchées
s'obtiennent en prenant $B_{ij} = C_{ij} \circ B_i|_{U_{ij}}$ et
$v_{ija}$. Cela montre que (2) implique (3).

\medskip\noindent

Supposons (3). Soit $\mathcal{G} \to \mathcal{H}$ un morphisme injectif
de $\mathcal{O}$-modules. Il faut montrer que
$\mathcal{G} \otimes_\mathcal{O} \mathcal{F} \to
\mathcal{H} \otimes_\mathcal{O} \mathcal{F}$
est injective. Soit $U$ un objet de $\mathcal{C}$ et soit
$s \in (\mathcal{G} \otimes_\mathcal{O} \mathcal{F})(U)$ une section dont
l'image dans $\mathcal{H} \otimes_\mathcal{O} \mathcal{F}$ est nulle.
Il faut montrer que $s$ est nulle. Comme
$\mathcal{G} \otimes_\mathcal{O} \mathcal{F}$
est un faisceau, il suffit de trouver un recouvrement
$\{U_i \to U\}_{i \in I}$ de $\mathcal{C}$ tel que $s|_{U_i}$ soit nulle
pour tout $i \in I$. On peut donc remplacer $U$ par les membres d'un
recouvrement. En particulier, puisque
$\mathcal{G} \otimes_\mathcal{O} \mathcal{F}$
est la faisceautisation de $\mathcal{G} \otimes_{p, \mathcal{O}} \mathcal{F}$
on peut supposer que $s$ est l'image d'un élément $s' \in 
\mathcal{G}(U) \otimes_{\mathcal{O}(U)} \mathcal{F}(U)$. En raisonnant de
même pour $\mathcal{H} \otimes_\mathcal{O} \mathcal{F}$, on peut supposer
que l'image de $s'$ dans
$\mathcal{H}(U) \otimes_{\mathcal{O}(U)} \mathcal{F}(U)$ est nulle.
Écrivons $\mathcal{F}(U) = \colim M_\alpha$ comme colimite filtrante de
$\mathcal{O}(U)$-modules de présentation finie $M_\alpha$ (Algèbre, lemme
\ref{algebra-lemma-module-colimit-fp}). Puisque le produit tensoriel commute
aux colimites filtrantes (Algèbre, lemme
\ref{algebra-lemma-tensor-products-commute-with-limits}), on peut choisir un
$\alpha$ tel que $s'$ provienne d'un élément
$s'' \in \mathcal{G}(U) \otimes_{\mathcal{O}(U)} M_\alpha$ et que l'image
de $s''$ dans
$\mathcal{H}(U) \otimes_{\mathcal{O}(U)} M_\alpha$ soit nulle. Fixons
$\alpha$ et $s''$, puis choisissons une présentation
$$
\mathcal{O}(U)^{\oplus m} \xrightarrow{A} \mathcal{O}(U)^{\oplus n}
\to M_\alpha \to 0
$$
Appliquons (3) au complexe correspondant de $\mathcal{O}_U$-modules
$$
\mathcal{O}_U^{\oplus m} \xrightarrow{A}
\mathcal{O}_U^{\oplus n} \xrightarrow{(s_1, \ldots, s_n)}
\mathcal{F}|_U
$$
Après avoir remplacé $U$ par les membres du recouvrement $U_i$, on obtient
que l'application
$$
M_\alpha \to \mathcal{F}(U)
$$
se factorise par un module libre $\mathcal{O}(U)^{\oplus l}$ pour un certain
$l$. Comme $\mathcal{G}(U) \to \mathcal{H}(U)$ est injective, on en déduit que
$$
\mathcal{G}(U) \otimes_{\mathcal{O}(U)} \mathcal{O}(U)^{\oplus l}
\to
\mathcal{H}(U) \otimes_{\mathcal{O}(U)} \mathcal{O}(U)^{\oplus l}
$$
est également injective. Puisque l'image de $s''$ dans le module de droite
est nulle, son image dans celui de gauche l'est aussi ; autrement dit,
$s$ est nulle, comme voulu.

\end{proof}

\begin{lemma}

\label{sites-modules-lemma-flat-over-thickening}

Soit $\mathcal{C}$ un site. Soit $\mathcal{O}' \to \mathcal{O}$ une
surjection de faisceaux d'anneaux dont le noyau $\mathcal{I}$ est un idéal
de carré nul. Soit $\mathcal{F}'$ un $\mathcal{O}'$-module et posons
$\mathcal{F} = \mathcal{F}'/\mathcal{I}\mathcal{F}'$. Les conditions
suivantes sont équivalentes :

\begin{enumerate}

\item $\mathcal{F}'$ est un $\mathcal{O}'$-module plat ;
\item $\mathcal{F}$ est un $\mathcal{O}$-module plat et
$\mathcal{I} \otimes_\mathcal{O} \mathcal{F} \to \mathcal{F}'$
est injective.

\end{enumerate}

\end{lemma}

\begin{proof}

Si (1) est satisfaite, alors
$\mathcal{F} = \mathcal{F}' \otimes_{\mathcal{O}'} \mathcal{O}$ est plat
sur $\mathcal{O}$ d'après le lemme \ref{sites-modules-lemma-flat-change-of-rings}, et
l'application
$\mathcal{I} \otimes_\mathcal{O} \mathcal{F} \to \mathcal{F}'$
est injective lorsque l'on applique $- \otimes_{\mathcal{O}'} \mathcal{F}'$
à la suite exacte
$0 \to \mathcal{I} \to \mathcal{O}' \to \mathcal{O} \to 0$ ; voir le lemme
\ref{sites-modules-lemma-flat-tor-zero}. Supposons maintenant (2). Dans la suite, nous
utiliserons sans le rappeler l'égalité
$\mathcal{K} \otimes_{\mathcal{O}'} \mathcal{F}' =
\mathcal{K} \otimes_\mathcal{O} \mathcal{F}$
pour tout $\mathcal{O}'$-module $\mathcal{K}$ annulé par $\mathcal{I}$.
Soit $\alpha : \mathcal{G}' \to \mathcal{H}'$ un morphisme injectif de
$\mathcal{O}'$-modules. Soient $\mathcal{G} \subset \mathcal{G}'$,
respectivement $\mathcal{H} \subset \mathcal{H}'$, les sous-faisceaux des
sections annulées par $\mathcal{I}$. Considérons le diagramme
$$
\xymatrix{
\mathcal{G} \otimes_{\mathcal{O}'} \mathcal{F}' \ar[r] \ar[d] &
\mathcal{G}' \otimes_{\mathcal{O}'} \mathcal{F}' \ar[r] \ar[d] &
\mathcal{G}'/\mathcal{G} \otimes_{\mathcal{O}'} \mathcal{F}' \ar[r] \ar[d] &
0 \\
\mathcal{H} \otimes_{\mathcal{O}'} \mathcal{F}' \ar[r] &
\mathcal{H}' \otimes_{\mathcal{O}'} \mathcal{F}' \ar[r] &
\mathcal{H}'/\mathcal{H} \otimes_{\mathcal{O}'} \mathcal{F}' \ar[r] & 0
}
$$
Notez que $\mathcal{G}'/\mathcal{G}$ et $\mathcal{H}'/\mathcal{H}$
sont annulés par $\mathcal{I}$ et que
$\mathcal{G}'/\mathcal{G} \to \mathcal{H}'/\mathcal{H}$ est injectif.
La flèche verticale de droite est donc injective puisque $\mathcal{F}$ est
plat sur $\mathcal{O}$. Il en va de même de celle de gauche. La flèche
verticale médiane est alors injective, et $\mathcal{F}'$ est plat.

\end{proof}

\begin{lemma}

\label{sites-modules-lemma-neighbourhood-isomorphism}
Soit $\mathcal{C}$ un site. Soit $\mathcal{O} \to \mathcal{O}'$ un
morphisme plat de faisceaux d'anneaux. Soit
$\mathcal{I} \subset \mathcal{O}$ un faisceau d'idéaux tel que le morphisme
induit $\mathcal{O}/\mathcal{I} \to
\mathcal{O}'/\mathcal{I}\mathcal{O}'$ soit un isomorphisme. Pour tout
$\mathcal{O}$-module $\mathcal{F}$ annulé par $\mathcal{I}^n$ pour un
$n \geq 0$, l'application $\text{id} \otimes 1 :
\mathcal{F} \to \mathcal{F} \otimes_\mathcal{O} \mathcal{O}'$ est un isomorphisme.
\end{lemma}

\begin{proof}
Démonstration omise. Indication : voir Compléments sur l'algèbre, lemme
\ref{more-algebra-lemma-neighbourhood-isomorphism}.
\end{proof}





\section{Duaux}

\label{sites-modules-section-duals}

\noindent

Soit $(\mathcal{C}, \mathcal{O})$ un site annelé. La catégorie des
$\mathcal{O}$-modules, munie du produit tensoriel construit à la section
\ref{sites-modules-section-tensor-product}, est une catégorie monoïdale symétrique. Pour
un $\mathcal{O}$-module $\mathcal{F}$, les conditions suivantes sont
équivalentes :

\begin{enumerate}

\item $\mathcal{F}$ admet un dual à gauche dans la catégorie monoïdale des
$\mathcal{O}$-modules ;
\item pour tout objet $U$ de $\mathcal{C}$, il existe un recouvrement
$\{U_i \to U\}$ tel que $\mathcal{F}|_{U_i}$ soit facteur direct d'un
$\mathcal{O}|_{U_i}$-module libre de rang fini ;
\item $\mathcal{F}$ est de présentation finie et plate comme
$\mathcal{O}$-module.

\end{enumerate}

Ces équivalences sont établies par l'exemple \ref{sites-modules-example-dual} et les
lemmes \ref{sites-modules-lemma-left-dual-module} et
\ref{sites-modules-lemma-flat-locally-finite-presentation} de la présente section.

\begin{example}

\label{sites-modules-example-dual}

Soit $(\mathcal{C}, \mathcal{O})$ un site annelé. Soit $\mathcal{F}$ un
$\mathcal{O}$-module tel que, pour tout objet $U$ de $\mathcal{C}$, il
existe un recouvrement $\{U_i \to U\}$ pour lequel
$\mathcal{F}|_{U_i}$ est facteur direct d'un
$\mathcal{O}|_{U_i}$-module libre de rang fini. Alors l'application
$$
\mathcal{F} \otimes_\mathcal{O}
\SheafHom_\mathcal{O}(\mathcal{F}, \mathcal{O})
\longrightarrow
\SheafHom_\mathcal{O}(\mathcal{F}, \mathcal{F})
$$
est un isomorphisme. En effet, cette assertion est locale ; elle est vraie
lorsque $\mathcal{F}$ est libre de rang fini et reste vraie pour tout
facteur direct d'un module pour lequel elle est satisfaite. Nous omettons
les détails. Notons
$$
\eta :
\mathcal{O}
\longrightarrow
\mathcal{F} \otimes_\mathcal{O}
\SheafHom_\mathcal{O}(\mathcal{F}, \mathcal{O})
$$
le morphisme qui envoie $1$ sur la section correspondant à
$\text{id}_\mathcal{F}$ sous l'isomorphisme ci-dessus.
Notons également
$$
\epsilon : 
\SheafHom_\mathcal{O}(\mathcal{F}, \mathcal{O})
\otimes_\mathcal{O} \mathcal{F}
\longrightarrow
\mathcal{O}
$$
le morphisme d'évaluation. On vérifie alors que
$\SheafHom_\mathcal{O}(\mathcal{F}, \mathcal{O}), \eta, \epsilon$
est un dual à gauche de $\mathcal{F}$ au sens de Catégories, définition
\ref{categories-definition-dual}. Nous omettons la vérification des
identités
$(1 \otimes \epsilon) \circ (\eta \otimes 1) = \text{id}_\mathcal{F}$
et
$(\epsilon \otimes 1) \circ (1 \otimes \eta) =
\text{id}_{\SheafHom_\mathcal{O}(\mathcal{F}, \mathcal{O})}$.

\end{example}

\begin{lemma}

\label{sites-modules-lemma-left-dual-module}

Soit $(\mathcal{C}, \mathcal{O})$ un site annelé et soit $\mathcal{F}$ un
$\mathcal{O}$-module. Supposons que $\mathcal{G}, \eta, \epsilon$
soit un dual à gauche de $\mathcal{F}$ dans la catégorie monoïdale des
$\mathcal{O}$-modules ; voir Catégories, définition
\ref{categories-definition-dual}. Alors :

\begin{enumerate}

\item pour tout objet $U$ de $\mathcal{C}$, il existe un recouvrement
$\{U_i \to U\}$ tel que $\mathcal{F}|_{U_i}$ soit facteur direct d'un
$\mathcal{O}|_{U_i}$-module libre de rang fini ;
\item l'application
$e : \SheafHom_\mathcal{O}(\mathcal{F}, \mathcal{O}) \to \mathcal{G}$
qui envoie une section locale $\lambda$ sur $(\lambda \otimes 1)(\eta)$
est un isomorphisme ;
\item on a $\epsilon(f, g) = e^{-1}(g)(f)$ pour les sections locales
$f$ et $g$ de $\mathcal{F}$ et $\mathcal{G}$.

\end{enumerate}

\end{lemma}

\begin{proof}

Les hypothèses signifient que
$$
\mathcal{F} \xrightarrow{\eta \otimes 1}
\mathcal{F} \otimes_\mathcal{O} \mathcal{G}
\otimes_\mathcal{O} \mathcal{F}
\xrightarrow{1 \otimes \epsilon} \mathcal{F}
\quad\text{et}\quad
\mathcal{G} \xrightarrow{1 \otimes \eta}
\mathcal{G} \otimes_\mathcal{O} \mathcal{F}
\otimes_\mathcal{O} \mathcal{G}
\xrightarrow{\epsilon \otimes 1} \mathcal{G}
$$
sont les applications identiques. Soit $U$ un objet de $\mathcal{C}$.
Après avoir remplacé $U$ par les membres d'un recouvrement de $U$, on peut trouver
un nombre fini de sections $f_1, \ldots, f_n$ et
$g_1, \ldots, g_n$ de $\mathcal{F}$ et $\mathcal{G}$ sur $U$ telles que
$\eta(1) = \sum f_i g_i$. Considérons
$$
\mathcal{O}_U^{\oplus n} \to \mathcal{F}|_U
$$
l'application qui envoie le $i$-ième vecteur de base sur $f_i$. On peut
alors factoriser $\eta|_U$ par une application
$\tilde \eta : \mathcal{O}_U \to
\mathcal{O}_U^{\oplus n} \otimes_{\mathcal{O}_U} \mathcal{G}|_U$. Nous obtenons un diagramme commutatif
$$
\xymatrix{
\mathcal{F}|_U
\ar[rr]_-{\eta \otimes 1} \ar[rrd]_-{\tilde \eta \otimes 1} & &
\mathcal{F}|_U \otimes \mathcal{G}|_U \otimes \mathcal{F}|_U
\ar[r]_-{1 \otimes \epsilon} &
\mathcal{F}|_U \\
& &
\mathcal{O}_U^{\oplus n} \otimes \mathcal{G}|_U \otimes \mathcal{F}|_U
\ar[u] \ar[r]^-{1 \otimes \epsilon} &
\mathcal{O}_U^{\oplus n} \ar[u]
}
$$
Ce diagramme montre que l'identité de $\mathcal{F}|_U$ se factorise par un
$\mathcal{O}_U$-module libre de rang fini, ce qui prouve (1). La partie (2)
découle de Catégories, lemme \ref{categories-lemma-left-dual}, et de sa
démonstration. La partie (3) découle de la première identité de la preuve.
On peut aussi déduire (2) et (3) de l'unicité des duaux à gauche
(Catégories, remarque \ref{categories-remark-left-dual-adjoint}) et de la
construction du dual à gauche dans l'exemple \ref{sites-modules-example-dual}.

\end{proof}

\begin{lemma}

\label{sites-modules-lemma-flat-locally-finite-presentation}
Soit $(\mathcal{C}, \mathcal{O})$ un site annelé et soit $\mathcal{F}$
localement de présentation finie et plat. Pour tout
objet $U$ de $\mathcal{C}$, il existe alors un recouvrement
$\{U_i \to U\}$ tel que $\mathcal{F}|_{U_i}$ soit facteur direct d'un
$\mathcal{O}_{U_i}$-module libre de rang fini.
\end{lemma}

\begin{proof}

Choisissons un objet $U$ de $\mathcal{C}$. Après avoir remplacé $U$ par les
membres d'un recouvrement, on peut supposer qu'il existe une présentation
$$
\mathcal{O}_U^{\oplus r} \to
\mathcal{O}_U^{\oplus n} \to \mathcal{F}|_U \to 0
$$
D'après le lemme \ref{sites-modules-lemma-flat-eq}, après avoir de nouveau remplacé $U$
par les membres d'un recouvrement, on peut supposer qu'il existe une
factorisation
$$
\mathcal{O}_U^{\oplus n} \to
\mathcal{O}_U^{\oplus n_1} \to \mathcal{F}|_U
$$
tel que le composé
$\mathcal{O}_U^{\oplus r} \to \mathcal{O}_U^{\oplus n}
\to \mathcal{O}_U^{\oplus n_r}$ est zéro. Cela signifie que la surjection $\mathcal{O}_U^{\oplus n_r} \to \mathcal{F}|_U$
admet une section, ce qui conclut la preuve.

\end{proof}











\section{Vers les modules constructibles}

\label{sites-modules-section-constructible}

\noindent
Rappelons qu'un objet quasi-compact d'un site est, grossièrement, un objet
dont tout recouvrement admet un raffinement fini ; la définition précise est
un peu plus élaborée, voir Sites, section
\ref{sites-section-quasi-compact}. Si tout objet d'un site admet un
recouvrement par des objets quasi-compacts, les modules
$j_!\mathcal{O}_U$, avec $U$ quasi-compact, forment un ensemble de
générateurs particulièrement commode pour la catégorie de tous les modules.

\begin{lemma}

\label{sites-modules-lemma-covering-gives-surjection}
Soit $(\mathcal{C}, \mathcal{O})$ un site annelé et soit
$\{U_i \to U\}$ un recouvrement de $\mathcal{C}$. Alors la suite
$$
\bigoplus j_{U_i \times_U U_j!}\mathcal{O}_{U_i \times_U U_j} \to
\bigoplus j_{U_i!}\mathcal{O}_{U_i} \to j_!\mathcal{O}_U \to 0
$$ est exacte.
\end{lemma}

\begin{proof}
Pour tout $\mathcal{O}$-module $\mathcal{F}$, le foncteur
$\Hom_\mathcal{O}(-, \mathcal{F})$ transforme la suite de l'énoncé en la
suite exacte $0 \to \mathcal{F}(U) \to \prod \mathcal{F}(U_i) \to
\prod \mathcal{F}(U_i \times_U U_j)$ ; voir
(\ref{sites-modules-equation-map-lower-shriek-OU-into-module}). Le résultat découle alors
d'Homologie, lemme \ref{homology-lemma-check-exactness}.
\end{proof}

\begin{lemma}

\label{sites-modules-lemma-silly-quasi-compact}
Soit $(\mathcal{C}, \mathcal{O})$ un site annelé. Soit
$\mathcal{U} = \{U_i \to U\}_{i \in I}$ un recouvrement dans
$\mathcal{C}$. Si $U$ est quasi-compact, il existe un sous-ensemble fini
$I' \subset I$ tel que la suite
$$
\bigoplus\nolimits_{i, i' \in I'}
j_{U_i \times_U U_{i'}!}\mathcal{O}_{U_i \times_U U_{i'}} \to
\bigoplus\nolimits_{i \in I'}
j_{U_i!}\mathcal{O}_{U_i} \to
j_!\mathcal{O}_U \to 0
$$ est exacte.
\end{lemma}

\begin{proof}

Le lemme découle immédiatement du lemme
\ref{sites-modules-lemma-covering-gives-surjection} lorsque $U$ satisfait la condition
(3) de Sites, lemme \ref{sites-lemma-conclude-quasi-compact}. Le lecteur
peut ignorer la démonstration du cas général. Par définition, il existe un
recouvrement
$\mathcal{V} = \{V_j \to U\}_{j \in J}$ et un morphisme
$\mathcal{V} \to \mathcal{U}$ de familles de morphismes de but fixé, donné
par $\text{id} : U \to U$, $\alpha : J \to I$ et
$f_j : V_j \to U_{\alpha(j)}$
(voir Sites, définition \ref{sites-definition-morphism-coverings}), de sorte
que l'image $I' \subset I$ de $\alpha$ soit finie. D'après Homologie, lemme
\ref{homology-lemma-check-exactness}, il suffit de montrer que, pour tout
faisceau de $\mathcal{O}$-modules $\mathcal{F}$, le foncteur
$\Hom_\mathcal{O}(-, \mathcal{F})$ transforme la suite de l'énoncé en une
suite exacte. Par (\ref{sites-modules-equation-map-lower-shriek-OU-into-module}), on
obtient la suite usuelle
$$
0 \to
\mathcal{F}(U) \to
\prod\nolimits_{i \in I'} \mathcal{F}(U_i) \to
\prod\nolimits_{i, i' \in I'} \mathcal{F}(U_i \times_U U_{i'})
$$
Cette suite est exacte d'après Sites, lemme
\ref{sites-lemma-compare-sheaf-condition},
appliqué à la famille de cartes $\{U_i \to U\}_{i \in I'}$
qui est raffinée par le recouvrement $\mathcal{V}$.

\end{proof}

\begin{lemma}

\label{sites-modules-lemma-sections-over-quasi-compact}
Soit $\mathcal{C}$ un site et soit $W$ un objet quasi-compact de
$\mathcal{C}$.
\begin{enumerate}
\item Le foncteur $\Sh(\mathcal{C}) \to \textit{Ensembles}$,
$\mathcal{F} \mapsto \mathcal{F}(W)$ commute aux coproduits.
\item Soit $\mathcal{O}$ un faisceau d'anneaux sur $\mathcal{C}$. Le
foncteur $\textit{Mod}(\mathcal{O}) \to \textit{Ab}$,
$\mathcal{F} \mapsto \mathcal{F}(W)$ commute aux sommes directes.
\end{enumerate}

\end{lemma}

\begin{proof}
Preuve de (1). D'après Sites, lemme
\ref{sites-lemma-directed-colimits-sections}, la prise des sections sur $W$
commute aux colimites filtrantes dont les morphismes de transition sont
injectifs. Soit $\mathcal{F}_i$ une famille de faisceaux d'ensembles
indexée par un ensemble $I$. Le coproduit $\coprod \mathcal{F}_i$ est la
colimite filtrante, sur l'ensemble ordonné des sous-ensembles finis
$E \subset I$, des coproduits
$\mathcal{F}_E = \coprod_{i \in E} \mathcal{F}_i$. Les morphismes de
transition étant injectifs, on conclut.

\medskip\noindent

Preuve de (2). Soit $\mathcal{F}_i$ une famille de faisceaux de
$\mathcal{O}$-modules indexée par un ensemble $I$. Alors
$\bigoplus \mathcal{F}_i$ est la colimite filtrante, sur l'ensemble ordonné
des sous-ensembles finis $E \subset I$, des sommes directes
$\mathcal{F}_E = \bigoplus_{i \in E} \mathcal{F}_i$. Une colimite filtrante
de faisceaux abéliens se calcule dans la catégorie des faisceaux d'ensembles.
Si $E \subset E'$, le morphisme de transition
$\mathcal{F}_E \to \mathcal{F}_{E'}$ est injectif, car la faisceautisation
est exacte et l'injectivité est claire sur les préfaisceaux sous-jacents.
Sites, lemme \ref{sites-lemma-directed-colimits-sections}, ramène donc au
cas d'un ensemble fini d'indices. Ce cas est traité par le lemme
\ref{sites-modules-lemma-limits-colimits-abelian-sheaves}, qui vaut sur tout objet de
$\mathcal{C}$.

\end{proof}

\begin{lemma}

\label{sites-modules-lemma-quasi-compact-hom-from}
Soit $(\mathcal{C}, \mathcal{O})$ un site annelé et soit $U$ un objet
quasi-compact de $\mathcal{C}$. Alors le foncteur
$\Hom_\mathcal{O}(j_!\mathcal{O}_U, -)$ commute aux sommes directes.
\end{lemma}

\begin{proof}
En effet, $\Hom_\mathcal{O}(j_!\mathcal{O}_U, \mathcal{F}) =
\mathcal{F}(U)$ d'après
(\ref{sites-modules-equation-map-lower-shriek-OU-into-module}), et le foncteur
$\mathcal{F} \mapsto \mathcal{F}(U)$ commute aux sommes directes par le
lemme \ref{sites-modules-lemma-sections-over-quasi-compact}.
\end{proof}

\noindent
Afin d'indiquer les résultats les plus nets possibles dans ce qui suit, nous introduisons une notation.

\begin{situation}

\label{sites-modules-situation-quasi-compact-objects}
Soit $\mathcal{C}$ un site et soit
$\mathcal{B} \subset \text{Ob}(\mathcal{C})$ un ensemble d'objets.
Considérons les conditions suivantes :
\begin{enumerate}
\item
\label{sites-modules-item-enough}
Tout objet de $\mathcal{C}$ admet un recouvrement par des éléments de
$\mathcal{B}$. \item
\label{sites-modules-item-enough-qc}
Tout $U \in \mathcal{B}$ est quasi-compact (Sites, section
\ref{sites-section-quasi-compact}). \item
\label{sites-modules-item-enough-qc-qs}
Pour tout recouvrement $\{U_i \to U\}$ avec
$U_i, U \in \mathcal{B}$, les produits fibrés
$U_i \times_U U_j$ sont quasi-compacts.
\end{enumerate}

\end{situation}

\begin{lemma}

\label{sites-modules-lemma-module-quotient-direct-sum}
Dans la situation \ref{sites-modules-situation-quasi-compact-objects}, supposons
(\ref{sites-modules-item-enough}) satisfaite.
\begin{enumerate}
\item Tout faisceau d'ensembles est le but d'un morphisme surjectif dont la
source est un coproduit $\coprod h_{U_i}^\#$, avec $U_i$ dans
$\mathcal{B}$.
\item Si $\mathcal{O}$ est un faisceau d'anneaux, tout
$\mathcal{O}$-module est quotient d'une somme directe
$\bigoplus\nolimits j_{U_i!}\mathcal{O}_{U_i}$, avec $U_i$ dans
$\mathcal{B}$.
\end{enumerate}

\end{lemma}

\begin{proof}

La partie (1) découle de Sites, lemmes
\ref{sites-lemma-sheaf-coequalizer-representable} et
\ref{sites-lemma-covering-surjective-after-sheafification}. La partie (2)
découle des lemmes \ref{sites-modules-lemma-module-quotient-flat} et
\ref{sites-modules-lemma-covering-gives-surjection}.

\end{proof}

\begin{lemma}

\label{sites-modules-lemma-module-filtered-colimit-constructibles}

Dans la situation \ref{sites-modules-situation-quasi-compact-objects}, supposons
(\ref{sites-modules-item-enough}) et (\ref{sites-modules-item-enough-qc}) satisfaites.

\begin{enumerate}

\item Tout faisceau d'ensembles est une colimite filtrante de faisceaux de
la forme

\begin{equation}
\label{sites-modules-equation-towards-constructible-sets}
\text{Coégalisateur}\left(
\xymatrix{
\coprod\nolimits_{j = 1, \ldots, m} h_{V_j}^\#
\ar@<1ex>[r] \ar@<-1ex>[r] &
\coprod\nolimits_{i = 1, \ldots, n} h_{U_i}^\#
}
\right)
\end{equation}

avec $U_i$ et $V_j$ dans $\mathcal{B}$.
\item Si $\mathcal{O}$ est un faisceau d'anneaux, tout
$\mathcal{O}$-module est une colimite filtrante de faisceaux de la forme

\begin{equation}
\label{sites-modules-equation-towards-constructible}
\Coker\left(
\bigoplus\nolimits_{j = 1, \ldots, m} j_{V_j!}\mathcal{O}_{V_j}
\longrightarrow
\bigoplus\nolimits_{i = 1, \ldots, n} j_{U_i!}\mathcal{O}_{U_i}
\right)
\end{equation}

avec $U_i$ et $V_j$ dans $\mathcal{B}$.

\end{enumerate}

\end{lemma}

\begin{proof}

Preuve de (1). D'après le lemme
\ref{sites-modules-lemma-module-quotient-direct-sum}, tout faisceau d'ensembles
$\mathcal{F}$ est le but d'une surjection dont la source est un coproduit
$\mathcal{F}_0$ de faisceaux de la forme $h_{U}^\#$, avec
$U \in \mathcal{B}$. En appliquant ce résultat à
$\mathcal{F}_0 \times_\mathcal{F} \mathcal{F}_0$, on obtient que
$\mathcal{F}$ est le coégalisateur d'une paire de morphismes
$$
\xymatrix{
\coprod\nolimits_{j \in J} h_{V_j}^\#
\ar@<1ex>[r] \ar@<-1ex>[r] &
\coprod\nolimits_{i \in I} h_{U_i}^\#
}
$$
pour certains ensembles d'indices $I$, $J$, et certains $V_j$, $U_i$ dans
$\mathcal{B}$. Pour tout sous-ensemble fini $J' \subset J$, il existe un
sous-ensemble fini $I' \subset I$ tel que les deux morphismes envoient le
coproduit indexé par $j \in J'$ dans celui indexé par $i \in I'$ ; voir
Sites, lemme \ref{sites-lemma-directed-colimits-sections}. Nous omettons les
détails ; on utilise qu'un coproduit infini est la colimite filtrante de ses
sous-coproduits finis. Notre faisceau est donc la colimite des conoyaux de
ces morphismes entre coproduits finis.

\medskip\noindent

Preuve de (2). Par le lemme \ref{sites-modules-lemma-module-quotient-direct-sum}, tout
module est quotient d'une somme directe de modules de la forme
$j_{U!}\mathcal{O}_U$, avec $U \in \mathcal{B}$. Tout module est donc un
conoyau
$$
\Coker\left(
\bigoplus\nolimits_{j \in J} j_{V_j!}\mathcal{O}_{V_j}
\longrightarrow
\bigoplus\nolimits_{i \in I} j_{U_i!}\mathcal{O}_{U_i}
\right)
$$
pour certains ensembles d'indices $I$, $J$, et certains $V_j$, $U_i$ dans
$\mathcal{B}$. Pour tout sous-ensemble fini $J' \subset J$, il existe un
sous-ensemble fini $I' \subset I$ tel que la somme directe indexée par
$j \in J'$ s'envoie dans celle indexée par $i \in I'$ ; voir le lemme
\ref{sites-modules-lemma-quasi-compact-hom-from}. Notre module est donc la colimite des
conoyaux de ces morphismes entre sommes directes finies.

\end{proof}

\begin{lemma}

\label{sites-modules-lemma-cokernel-map-towards-constructibles}

Dans la situation \ref{sites-modules-situation-quasi-compact-objects}, supposons
(\ref{sites-modules-item-enough}) et (\ref{sites-modules-item-enough-qc}) satisfaites. Soit
$\mathcal{O}$ un faisceau d'anneaux. Le conoyau d'un morphisme entre deux
modules de la forme (\ref{sites-modules-equation-towards-constructible}) est encore un
module de la forme (\ref{sites-modules-equation-towards-constructible}).

\end{lemma}

\begin{proof}

Soit $\mathcal{F} = \Coker(\bigoplus j_{V_j!}\mathcal{O}_{V_j} \to
\bigoplus j_{U_i!}\mathcal{O}_{U_i})$
comme dans (\ref{sites-modules-equation-towards-constructible}). Il suffit de montrer que
le conoyau d'un morphisme
$\varphi : j_{W!}\mathcal{O}_W \to \mathcal{F}$, avec
$W \in \mathcal{B}$, est encore un module du même type. Le morphisme
$\varphi$ correspond à un élément $s \in \mathcal{F}(W)$. Comme
$\bigoplus j_{U_i!}\mathcal{O}_{U_i} \to \mathcal{F}$ est surjectif, on
peut, par (\ref{sites-modules-item-enough}), choisir un recouvrement
$\{W_k \to W\}_{k \in K}$ avec $W_k \in \mathcal{B}$
tel que $s|_{W_k}$ soit l'image d'une section
$s_k$ de $\bigoplus j_{U_i!}\mathcal{O}_{U_i})$. Par
(\ref{sites-modules-item-enough-qc}), l'objet $W$ est quasi-compact. D'après le lemme
\ref{sites-modules-lemma-silly-quasi-compact}, il existe un sous-ensemble fini
$K' \subset K$ tel que
$\bigoplus_{k \in K'} j_{W_k!}\mathcal{O}_{W_k} \to j_{W!}\mathcal{O}_W$
soit surjectif. On en déduit que $\Coker(\varphi)$ est égal à
$$
\Coker\left(
\bigoplus\nolimits_{k \in K'} j_{W_k!}\mathcal{O}_{W_k} \oplus
\bigoplus j_{V_j!}\mathcal{O}_{V_j}
\longrightarrow
\bigoplus j_{U_i!}\mathcal{O}_{U_i}
\right)
$$
où le morphisme $\bigoplus_{k \in K'} j_{W_k!}\mathcal{O}_{W_k}
\to \bigoplus j_{U_i!}\mathcal{O}_{U_i}$ correspond à
$\sum_{k \in K'} s_k$. Cela termine la preuve.

\end{proof}

\begin{lemma}

\label{sites-modules-lemma-change-presentation-towards-constructibles}

Dans la situation \ref{sites-modules-situation-quasi-compact-objects}, supposons
(\ref{sites-modules-item-enough}), (\ref{sites-modules-item-enough-qc}) et
(\ref{sites-modules-item-enough-qc-qs}) satisfaites. Soit $\mathcal{O}$ un faisceau
d'anneaux. Donnons-nous un morphisme
$$
\bigoplus\nolimits_{j = 1, \ldots, m} j_{V_j!}\mathcal{O}_{V_j}
\longrightarrow
\bigoplus\nolimits_{i = 1, \ldots, n} j_{U_i!}\mathcal{O}_{U_i}
$$
avec $U_i$ et $V_j$ dans $\mathcal{B}$, ainsi que des recouvrements
$\{U_{ik} \to U_i\}_{k \in K_i}$ avec $U_{ik} \in \mathcal{B}$. Il existe
alors des sous-ensembles finis $K'_i \subset K_i$, un ensemble fini $L$
d'objets $W_l \in \mathcal{B}$ et un diagramme commutatif
$$
\xymatrix{
\bigoplus_{l \in L} j_{W_l!}\mathcal{O}_{W_l} \ar[d] \ar[r] &
\bigoplus_{i = 1, \ldots, n} \bigoplus_{k \in K'_i}
j_{U_{ik}!}\mathcal{O}_{U_{ik}} \ar[d] \\
\bigoplus_{j = 1, \ldots, m} j_{V_j!}\mathcal{O}_{V_j} \ar[r] &
\bigoplus_{i = 1, \ldots, n} j_{U_i!}\mathcal{O}_{U_i}
}
$$
induisant un isomorphisme sur les conoyaux des morphismes horizontaux.

\end{lemma}

\begin{proof}

Puisque $U_i$ est quasi-compact, on peut choisir des sous-ensembles finis
$K'_i \subset K_i$ comme dans le lemme
\ref{sites-modules-lemma-silly-quasi-compact}. Puisque
$\bigoplus_{i = 1, \ldots, n}
\bigoplus_{k \in K'_i} j_{U_{ik}!}\mathcal{O}_{U_{ik}} \to
\bigoplus_{i = 1, \ldots, n} j_{U_i!}\mathcal{O}_{U_i}$ est surjectif, on
peut trouver des recouvrements
$\{V_{jm} \to V_j\}_{m \in M_j}$ avec $V_{jm} \in \mathcal{B}$ pour
lesquels il existe un diagramme commutatif
$$
\xymatrix{
\bigoplus_{j = 1, \ldots, m} \bigoplus_{m \in M_j}
j_{V_{jm}!}\mathcal{O}_{V_{jm}} \ar[d] \ar[r] &
\bigoplus_{i = 1, \ldots n} \bigoplus_{k \in K'_i}
j_{U_{ik}!}\mathcal{O}_{U_{ik}} \ar[d] \\
\bigoplus_{j = 1, \ldots, m} j_{V_j!}\mathcal{O}_{V_j} \ar[r] &
\bigoplus_{i = 1, \ldots, n} j_{U_i!}\mathcal{O}_{U_i}
}
$$
Puisque $V_j$ est quasi-compact, on peut choisir des sous-ensembles finis
$M'_j \subset M_j$ comme dans le lemme
\ref{sites-modules-lemma-silly-quasi-compact}. Posons
$$
L = \left(\coprod\nolimits_{i = 1, \ldots, n} K'_i \times K'_i \right)
\coprod
\left(\coprod\nolimits_{j = 1, \ldots, m} M'_j\right)
$$
et, pour $l = (k, k') \in K'_i \times K'_i \subset L$, posons
$W_l = U_{ik} \times_{U_i} U_{ik'}$ ; pour
$l = m \in M'_j \subset L$, posons $W_l = V_{jm}$. Les suites exactes du
lemme \ref{sites-modules-lemma-silly-quasi-compact}
pour les familles $\{U_{ik} \to U_i\}_{k \in K'_i}$
montrent que l'on obtient un diagramme comme dans l'énoncé du lemme ; nous
omettons les détails. Il reste seulement à assurer que
$W_l \in \mathcal{B}$. Or $W_l$ est quasi-compact pour tout $l \in L$.
Une nouvelle application du lemme \ref{sites-modules-lemma-silly-quasi-compact} fournit
des familles finies de morphismes
$\{W_{lt} \to W_l\}_{t \in T_l}$ avec $W_{lt} \in \mathcal{B}$ telles que
$\bigoplus j_{W_{lt}!}\mathcal{O}_{W_{lt}} \to j_{W_l!}\mathcal{O}_{W_l}$
soit surjectif. On remplace alors $L$ par $\coprod_{l \in L} T_l$, ce qui
achève la démonstration.

\end{proof}

\begin{lemma}

\label{sites-modules-lemma-extension-towards-constructibles}
Dans la situation \ref{sites-modules-situation-quasi-compact-objects}, supposons
(\ref{sites-modules-item-enough}), (\ref{sites-modules-item-enough-qc}) et
(\ref{sites-modules-item-enough-qc-qs}) satisfaites. Soit $\mathcal{O}$ un faisceau
d'anneaux. Toute extension de modules de la forme
(\ref{sites-modules-equation-towards-constructible}) est encore un module de la forme
(\ref{sites-modules-equation-towards-constructible}).
\end{lemma}

\begin{proof}

Soit $0 \to \mathcal{F}_1 \to \mathcal{F}_2 \to \mathcal{F}_3 \to 0$
une suite exacte courte de $\mathcal{O}$-modules dans laquelle
$\mathcal{F}_1$ et $\mathcal{F}_3$ sont de la forme
(\ref{sites-modules-equation-towards-constructible}). Choisissons des présentations
$$
\bigoplus A_{V_j} \to \bigoplus A_{U_i} \to \mathcal{F}_1 \to 0
\quad\text{et}\quad
\bigoplus A_{T_j} \to \bigoplus A_{W_i} \to \mathcal{F}_3 \to 0
$$
Dans cette preuve, toutes les sommes directes sont finies, et nous écrivons
$A_U = j_{U!}\mathcal{O}_U$ pour $U \in \mathcal{B}$. Puisque
$\mathcal{F}_2 \to \mathcal{F}_3$ est surjectif, on peut choisir des
recouvrements $\{W_{ik} \to W_i\}$ avec $W_{ik} \in \mathcal{B}$ tels que
$A_{W_{ik}} \to \mathcal{F}_3$ se relève en un morphisme
$A_{W_{ik}} \to \mathcal{F}_2$. D'après le lemme
\ref{sites-modules-lemma-change-presentation-towards-constructibles}, on peut remplacer
la famille $\{W_i\}$ par une sous-famille finie de la famille
$\{W_{ik}\}$ et supposer que le morphisme
$\bigoplus A_{W_i} \to \mathcal{F}_3$ se relève dans $\mathcal{F}_2$.
Considérons le noyau
$$
\mathcal{K}_2 = \Ker(\bigoplus A_{U_i} \oplus \bigoplus A_{W_i}
\longrightarrow \mathcal{F}_2)
$$
Le lemme du serpent montre que ce noyau surjecte sur
$\mathcal{K}_3 = \Ker(\bigoplus A_{W_i} \to \mathcal{F}_3)$. En raisonnant
comme ci-dessus, et après avoir remplacé chaque $T_j$ par une famille finie
d'éléments de $\mathcal{B}$, ce que permet le lemme
\ref{sites-modules-lemma-silly-quasi-compact}, on peut supposer qu'il existe un morphisme
$\bigoplus A_{T_j} \to \mathcal{K}_2$ relevant le morphisme donné
$\bigoplus A_{T_j} \to \mathcal{K}_3$. Alors $\bigoplus A_{V_j} \oplus \bigoplus A_{T_j} \to \mathcal{K}_2$
est surjectif, ce qui achève la démonstration.

\end{proof}

\begin{lemma}

\label{sites-modules-lemma-towards-constructible-when-serre-subcategory}

Dans la situation \ref{sites-modules-situation-quasi-compact-objects}, supposons
(\ref{sites-modules-item-enough}), (\ref{sites-modules-item-enough-qc}) et
(\ref{sites-modules-item-enough-qc-qs}) satisfaites. Soit $\mathcal{O}$ un faisceau
d'anneaux. Soit $\mathcal{A} \subset \textit{Mod}(\mathcal{O})$ la
sous-catégorie pleine des modules isomorphes à un conoyau de la forme
(\ref{sites-modules-equation-towards-constructible}). Si le noyau de tout morphisme de
$\mathcal{O}$-modules de la forme
$$
\bigoplus\nolimits_{j = 1, \ldots, m} j_{V_j!}\mathcal{O}_{V_j}
\longrightarrow
\bigoplus\nolimits_{i = 1, \ldots, n} j_{U_i!}\mathcal{O}_{U_i}
$$
avec $U_i$ et $V_j$ dans $\mathcal{B}$ appartient à $\mathcal{A}$, alors
$\mathcal{A}$ est une sous-catégorie de Serre faible de
$\textit{Mod}(\mathcal{O})$.

\end{lemma}

\begin{proof}

Nous utiliserons le critère d'Homologie, lemme
\ref{homology-lemma-characterize-weak-serre-subcategory}. D'après les
lemmes \ref{sites-modules-lemma-cokernel-map-towards-constructibles} et
\ref{sites-modules-lemma-extension-towards-constructibles}
il suffit de montrer que le noyau d'un morphisme
$\mathcal{F} \to \mathcal{G}$ entre objets de $\mathcal{A}$ appartient à
$\mathcal{A}$. Pour cela, choisissons des présentations
$$
\bigoplus A_{V_j} \to \bigoplus A_{U_i} \to \mathcal{F} \to 0
\quad\text{et}\quad
\bigoplus A_{T_j} \to \bigoplus A_{W_i} \to \mathcal{G} \to 0
$$
Dans cette preuve, toutes les sommes directes sont finies, et nous écrivons
$A_U = j_{U!}\mathcal{O}_U$ pour $U \in \mathcal{B}$. En appliquant les
lemmes \ref{sites-modules-lemma-covering-gives-surjection} et
\ref{sites-modules-lemma-change-presentation-towards-constructibles}
et en raisonnant comme dans la preuve du lemme
\ref{sites-modules-lemma-extension-towards-constructibles}, on peut supposer que le
morphisme $\mathcal{F} \to \mathcal{G}$ se relève en un morphisme de
présentations
$$
\xymatrix{
\bigoplus A_{V_j} \ar[r] \ar[d] &
\bigoplus A_{U_i} \ar[r] \ar[d] &
\mathcal{F} \ar[r] \ar[d] & 0 \\
\bigoplus A_{T_j} \ar[r] &
\bigoplus A_{W_i} \ar[r] &
\mathcal{G} \ar[r] & 0
}
$$
On obtient alors
$$
\Ker(\mathcal{F} \to \mathcal{G}) =
\Coker\left(\bigoplus A_{V_j} \to
\Ker\left(
\bigoplus A_{T_j} \oplus \bigoplus A_{U_i} \to \bigoplus A_{W_i}\right)\right)
$$
et le résultat découle de l'hypothèse et du lemme
\ref{sites-modules-lemma-cokernel-map-towards-constructibles}.

\end{proof}







\section{Morphismes plats}

\label{sites-modules-section-flat-morphisms}

\begin{definition}

\label{sites-modules-definition-flat-morphism}

Soit
$(f, f^\sharp) :
(\Sh(\mathcal{C}), \mathcal{O})
\longrightarrow
(\Sh(\mathcal{C}'), \mathcal{O}')$
un morphisme de topos annelés. On dit que $(f, f^\sharp)$ est
{\it plat} si le morphisme d'anneaux $f^\sharp : f^{-1}\mathcal{O}' \to \mathcal{O}$
est plat. Nous disons qu'un morphisme de sites annelés est {\it plat}
si le morphisme de topos annelés associé est plat.

\end{definition}

\begin{lemma}

\label{sites-modules-lemma-flat-pullback-exact}
Soit $f : \Sh(\mathcal{C}) \to \Sh(\mathcal{C}')$ un morphisme de topos
annelés. Alors
$$
f^{-1} : \textit{Ab}(\mathcal{C}') \longrightarrow \textit{Ab}(\mathcal{C}),
\quad
\mathcal{F} \longmapsto f^{-1}\mathcal{F}
$$ est exact. Si $(f, f^\sharp) :
(\Sh(\mathcal{C}), \mathcal{O})
\to
(\Sh(\mathcal{C}'), \mathcal{O}')$ est un morphisme de topos annelés plat, alors $$
f^* : \textit{Mod}(\mathcal{O}') \longrightarrow \textit{Mod}(\mathcal{O}),
\quad
\mathcal{F} \longmapsto f^*\mathcal{F}
$$ est exact.
\end{lemma}

\begin{proof}
Pour un faisceau abélien $\mathcal{G}$ sur $\mathcal{C}'$, le faisceau
d'ensembles sous-jacent à $f^{-1}\mathcal{G}$ est l'image inverse par
$f^{-1}$ du faisceau d'ensembles sous-jacent à $\mathcal{G}$ ; voir Sites,
section \ref{sites-section-sheaves-algebraic-structures}. L'exactitude de
$f^{-1}$ sur les faisceaux d'ensembles, qui fait partie de la définition
d'un morphisme de topos (voir Sites, définition
\ref{sites-definition-topos}), implique donc l'exactitude de $f^{-1}$ sur les
faisceaux abéliens.

\medskip\noindent

Pour voir l'énoncé sur les modules, rappelons que $f^*\mathcal{F}$ est défini comme le produit tensoriel
$f^{-1}\mathcal{F} \otimes_{f^{-1}\mathcal{O}', f^\sharp} \mathcal{O}$.
Ainsi, $f^*$ est le composé de deux foncteurs exacts.

\end{proof}

\begin{definition}

\label{sites-modules-definition-flat-module}
Soit $f : (\Sh(\mathcal{C}), \mathcal{O}) \to
(\Sh(\mathcal{D}), \mathcal{O}')$ un morphisme de topos annelés et soit
$\mathcal{F}$ un faisceau de $\mathcal{O}$-modules. On dit que
$\mathcal{F}$ est {\it plat sur
$(\Sh(\mathcal{D}), \mathcal{O}')$} si $\mathcal{F}$ est plat comme
$f^{-1}\mathcal{O}'$-module.
\end{definition}

\noindent
Ceci est compatible avec la notion telle que définie pour les morphismes des espaces annelés, voir Modules, définition \ref{modules-definition-flat-module} et la discussion suivante.

\begin{lemma}

\label{sites-modules-lemma-pullback-internal-hom}
Soit $f : (\mathcal{C}, \mathcal{O}_\mathcal{C}) \to
(\mathcal{D}, \mathcal{O}_\mathcal{D})$ un morphisme de sites annelés.
Soient $\mathcal{F}$ et $\mathcal{G}$ des
$\mathcal{O}_\mathcal{D}$-modules. Si $\mathcal{F}$ est de présentation
finie et si $f$ est plat, alors le morphisme canonique
$$
f^*\SheafHom_{\mathcal{O}_\mathcal{D}}(\mathcal{F}, \mathcal{G})
\longrightarrow
\SheafHom_{\mathcal{O}_\mathcal{C}}(f^*\mathcal{F}, f^*\mathcal{G})
$$
de la remarque \ref{sites-modules-remark-pullback-internal-hom} est un isomorphisme.
\end{lemma}

\begin{proof}

Supposons que $f$ soit donné par le foncteur continu
$u : \mathcal{D} \to \mathcal{C}$. Il faut montrer que la restriction du
morphisme à $\mathcal{C}/U$ est un isomorphisme pour tout
$U \in \Ob(\mathcal{C})$. On peut remplacer $U$ par les membres
d'un recouvrement de $U$. D'après Sites, lemme
\ref{sites-lemma-morphism-of-sites-covering}, on peut donc supposer qu'il
existe un morphisme $U \to u(V)$ pour un certain
$V \in \Ob(\mathcal{D})$. On peut alors remplacer $U$ par $u(V)$.
Puisque $u$ est continu, on peut remplacer $V$ par les membres d'un
recouvrement et supposer qu'il existe une présentation
$\mathcal{O}_V^{\oplus m} \to \mathcal{O}_V^{\oplus n} \to
\mathcal{F}|_V \to 0$
sur $\mathcal{D}/V$. Comme la formation de $\SheafHom$ commute à la
localisation (lemme \ref{sites-modules-lemma-internal-hom-restriction}), on peut remplacer
$f$ par le morphisme
$(\mathcal{C}/u(V), \mathcal{O}_{u(V)}) \to
(\mathcal{D}/V, \mathcal{O}_V)$ induit par $f$. On se ramène ainsi au cas
où $\mathcal{F}$ admet une présentation globale
$\mathcal{O}_\mathcal{D}^{\oplus m} \to \mathcal{O}_\mathcal{D}^{\oplus n} \to
\mathcal{F} \to 0$. Puisque $f$ est plat et que
$f^*\mathcal{O}_\mathcal{D} = \mathcal{O}_\mathcal{C}$, on obtient la
présentation correspondante
$\mathcal{O}_\mathcal{C}^{\oplus m} \to \mathcal{O}_\mathcal{C}^{\oplus n} \to
f^*\mathcal{F} \to 0$ ; voir le lemme
\ref{sites-modules-lemma-flat-pullback-exact}. En utilisant la commutation de
$\SheafHom$ aux sommes directes finies dans la première variable, le fait
que les deux foncteurs
$\SheafHom_{\mathcal{O}_\mathcal{C}}(\mathcal{O}_\mathcal{C}, -)$ et
$\SheafHom_{\mathcal{O}_\mathcal{D}}(\mathcal{O}_\mathcal{D}, -)$
sont le foncteur identité, et la fonctorialité de la construction de la
remarque \ref{sites-modules-remark-pullback-internal-hom}, on obtient un diagramme
commutatif
$$
\xymatrix{
0 \ar[r] &
f^*\SheafHom_{\mathcal{O}_\mathcal{D}}(\mathcal{F}, \mathcal{G})
\ar[d] \ar[r] &
f^*\mathcal{G}^{\oplus n} \ar[d] \ar[r] &
f^*\mathcal{G}^{\oplus m} \ar[d] \\
0 \ar[r] &
\SheafHom_{\mathcal{O}_\mathcal{C}}(f^*\mathcal{F}, f^*\mathcal{G})
\ar[r] &
f^*\mathcal{G}^{\oplus n} \ar[r] &
f^*\mathcal{G}^{\oplus m}
}
$$
où les deux flèches verticales de droite sont des isomorphismes. Les lignes
sont exactes d'après le lemme \ref{sites-modules-lemma-internal-hom-exact}. Le lemme des
cinq permet de conclure.

\end{proof}







\section{Modules inversibles}

\label{sites-modules-section-invertible}

\noindent
Voici la définition.

\begin{definition}

\label{sites-modules-definition-invertible-sheaf}

Soit $(\mathcal{C}, \mathcal{O})$ un site annelé.

\begin{enumerate}

\item Un $\mathcal{O}$-module $\mathcal{F}$ localement libre de rang fini
est dit de {\it rang $r$} si, pour tout objet $U$ de $\mathcal{C}$, il
existe un recouvrement $\{U_i \to U\}$ de $U$ tel que
$\mathcal{F}|_{U_i}$ soit isomorphe à
$\mathcal{O}_{U_i}^{\oplus r}$ comme $\mathcal{O}_{U_i}$-module.
\item Un $\mathcal{O}$-module $\mathcal{L}$ est dit {\it inversible}
si le foncteur
$$
\textit{Mod}(\mathcal{O}) \longrightarrow \textit{Mod}(\mathcal{O}),\quad
\mathcal{F}  \longmapsto \mathcal{F} \otimes_\mathcal{O} \mathcal{L}
$$
est une équivalence.
\item Le faisceau {\it $\mathcal{O}^*$} est le sous-faisceau de
$\mathcal{O}$ défini par
$$
U \longmapsto \mathcal{O}^*(U) = \{f \in \mathcal{O}(U) \mid
\exists g \in \mathcal{O}(U)\text{ tel que }fg = 1\}
$$
C'est un faisceau de groupes abéliens pour la multiplication.

\end{enumerate}

\end{definition}

\noindent
Le lemme \ref{sites-modules-lemma-invertible-is-locally-free-rank-1} ci-dessous précise
le lien avec les modules localement libres de rang $1$.

\begin{lemma}

\label{sites-modules-lemma-invertible}

Soit $(\mathcal{C}, \mathcal{O})$ un site annelé et soit $\mathcal{L}$ un
$\mathcal{O}$-module. Les conditions suivantes sont équivalentes :

\begin{enumerate}

\item $\mathcal{L}$ est inversible ;
\item il existe un $\mathcal{O}$-module $\mathcal{N}$
tel que
$\mathcal{L} \otimes_\mathcal{O} \mathcal{N} \cong \mathcal{O}$.

\end{enumerate}

Dans ce cas, nous avons

\begin{enumerate}

\item[(a)] $\mathcal{L}$ est un $\mathcal{O}$-module plat de présentation
finie ;
\item[b)] pour tout objet $U$ de $\mathcal{C}$, il existe un recouvrement
$\{U_i \to U\}$ tel que $\mathcal{L}|_{U_i}$ soit facteur direct d'un
module libre de type fini ;
\item[c)] le module $\mathcal{N}$ en (2) est isomorphe à
$\SheafHom_\mathcal{O}(\mathcal{L}, \mathcal{O})$.

\end{enumerate}

\end{lemma}

\begin{proof}
Supposons (1). Le foncteur $- \otimes_\mathcal{O} \mathcal{L}$ est
essentiellement surjectif ; il existe donc un $\mathcal{O}$-module
$\mathcal{N}$ comme en (2). Réciproquement, si (2) est satisfaite, le
foncteur $- \otimes_\mathcal{O} \mathcal{N}$ est un quasi-inverse du
foncteur $- \otimes_\mathcal{O} \mathcal{L}$, ce qui donne (1).

\medskip\noindent

Supposons (1) et (2). Puisque $- \otimes_\mathcal{O} \mathcal{L}$ est une
équivalence, ce foncteur est exact et $\mathcal{L}$ est donc plat. Soit
$\psi : \mathcal{L} \otimes_\mathcal{O} \mathcal{N} \to \mathcal{O}$
l'isomorphisme donné, et soit $U$ un objet de $\mathcal{C}$. Nous allons
montrer que la restriction de $\mathcal{L}$ aux membres d'un recouvrement de
$U$ est facteur direct d'un module libre ; cela impliquera en
particulier que $\mathcal{L}$ est de présentation finie. Par construction
du produit $\otimes$, après avoir
remplacé $U$ par les membres d'un recouvrement, on peut supposer qu'il
existe un entier $n \geq 1$, des sections $x_i \in \mathcal{L}(U)$ et
$y_i \in \mathcal{N}(U)$ telles que
$\psi(\sum x_i \otimes y_i) = 1$. Considérons les isomorphismes
$$
\mathcal{L}|_U \to
\mathcal{L}|_U \otimes_{\mathcal{O}_U}
\mathcal{L}|_U \otimes_{\mathcal{O}_U} \mathcal{N}|_U \to \mathcal{L}|_U
$$
où la première flèche envoie $x$ sur
$\sum x_i \otimes x \otimes y_i$ et la seconde envoie
$x \otimes x' \otimes y$ sur $\psi(x' \otimes y)x$. On en déduit que
$x \mapsto \sum \psi(x \otimes y_i)x_i$ est un automorphisme de
$\mathcal{L}|_U$. Cet automorphisme se factorise comme
$$
\mathcal{L}|_U \to \mathcal{O}_U^{\oplus n} \to \mathcal{L}|_U
$$
où la première flèche est donnée
$x \mapsto (\psi(x \otimes y_1), \ldots, \psi(x \otimes y_n))$
et la seconde par $(a_1, \ldots, a_n) \mapsto \sum a_i x_i$. Ainsi,
$\mathcal{L}|_U$ est facteur direct d'un $\mathcal{O}_U$-module libre de
rang fini.

\medskip\noindent

Toujours sous les hypothèses (1) et (2), considérons le morphisme
d'évaluation
$$
\mathcal{L} \otimes_\mathcal{O}
\SheafHom_\mathcal{O}(\mathcal{L}, \mathcal{O}_X)
\longrightarrow \mathcal{O}_X
$$
Pour achever la démonstration, montrons qu'il s'agit d'un isomorphisme. Le
lemme \ref{sites-modules-lemma-internal-hom-adjoint-tensor} donne
$$
\Hom_\mathcal{O}(\mathcal{O}, \mathcal{O}) =
\Hom_\mathcal{O}
(\mathcal{N} \otimes_\mathcal{O} \mathcal{L}, \mathcal{O})
\longrightarrow
\Hom_\mathcal{O}
(\mathcal{N}, \SheafHom_\mathcal{O}(\mathcal{L}, \mathcal{O}))
$$
L'image de $1$ définit un morphisme
$\mathcal{N} \to \SheafHom_\mathcal{O}(\mathcal{L}, \mathcal{O})$. En
prenant le produit tensoriel avec $\mathcal{L}$, on obtient
$$
\mathcal{O} = \mathcal{L} \otimes_\mathcal{O} \mathcal{N}
\longrightarrow
\mathcal{L} \otimes_\mathcal{O} \SheafHom_\mathcal{O}(\mathcal{L}, \mathcal{O})
$$
Ce morphisme est l'inverse du morphisme d'évaluation ; nous omettons le
calcul.

\end{proof}

\begin{lemma}

\label{sites-modules-lemma-pullback-invertible}

Soit $f : (\Sh(\mathcal{C}), \mathcal{O}_\mathcal{C}) \to
(\Sh(\mathcal{D}), \mathcal{O}_\mathcal{D})$ un morphisme de topos
annelés. L'image inverse $f^*\mathcal{L}$ d'un
$\mathcal{O}_\mathcal{D}$-module inversible est inversible.

\end{lemma}

\begin{proof}
D'après le lemme \ref{sites-modules-lemma-invertible}, il existe un
$\mathcal{O}_\mathcal{D}$-module $\mathcal{N}$ tel que
$\mathcal{L} \otimes_{\mathcal{O}_\mathcal{D}} \mathcal{N} \cong
\mathcal{O}_\mathcal{D}$. Par image inverse, on obtient
$f^*\mathcal{L} \otimes_{\mathcal{O}_\mathcal{C}} f^*\mathcal{N} \cong
\mathcal{O}_\mathcal{C}$ d'après le lemme
\ref{sites-modules-lemma-tensor-product-pullback}. Le lemme \ref{sites-modules-lemma-invertible} montre
donc que $f^*\mathcal{L}$ est inversible.
\end{proof}

\begin{lemma}

\label{sites-modules-lemma-constructions-invertible}
Soit $(\mathcal{C}, \mathcal{O})$ un site annelé.
\begin{enumerate}
\item Si $\mathcal{L}$ et $\mathcal{N}$ sont des $\mathcal{O}$-modules
inversibles, alors $\mathcal{L} \otimes_\mathcal{O} \mathcal{N}$ est
inversible.
\item Si $\mathcal{L}$ est un $\mathcal{O}$-module inversible, alors
$\SheafHom_\mathcal{O}(\mathcal{L}, \mathcal{O})$ est inversible et le
morphisme d'évaluation $\mathcal{L} \otimes_\mathcal{O}
\SheafHom_\mathcal{O}(\mathcal{L}, \mathcal{O}) \to \mathcal{O}$ est un isomorphisme.
\end{enumerate}

\end{lemma}

\begin{proof}

La partie (1) est immédiate d'après la définition, et la partie (2) découle
du lemme \ref{sites-modules-lemma-invertible} et de sa démonstration.

\end{proof}

\begin{lemma}

\label{sites-modules-lemma-pic-set}

Soit $(\mathcal{C}, \mathcal{O})$ un site annelé. Il existe un ensemble de
modules inversibles $\{\mathcal{L}_i\}_{i \in I}$ tel que tout module
inversible sur $(\mathcal{C}, \mathcal{O})$
est isomorphe à exactement l'un des $\mathcal{L}_i$.

\end{lemma}

\begin{proof}
Démonstration omise ; voir Faisceaux de modules, lemme
\ref{modules-lemma-pic-set}.
\end{proof}

\noindent
Le lemme \ref{sites-modules-lemma-pic-set} affirme que les classes d'isomorphisme des
faisceaux inversibles forment un ensemble. Le lemme
\ref{sites-modules-lemma-constructions-invertible} montre que le produit tensoriel munit
cet ensemble d'une structure de groupe abélien ; l'inverse de
$\mathcal{L}$ y est représenté par
$\SheafHom_\mathcal{O}(\mathcal{L}, \mathcal{O})$.

\medskip\noindent
Plus précisément, pour un $\mathcal{O}$-module inversible $\mathcal{L}$ et
$n \in \mathbf{Z}$, on définit la {\it $n$-ième puissance tensorielle}
$\mathcal{L}^{\otimes n}$ de $\mathcal{L}$ comme l'image de $\mathcal{O}$
par l'équivalence $\mathcal{F} \mapsto
\mathcal{F} \otimes_\mathcal{O} \mathcal{L}$ itérée $n$ fois. Cette
définition a également un sens pour $n$ négatif, puisqu'un
$\mathcal{O}$-module est inversible précisément lorsque le produit tensoriel
définit une équivalence.
Plus explicitement :
$$
\mathcal{L}^{\otimes n} =
\left\{
\begin{matrix}
\mathcal{O} & \text{si} & n = 0 \\
\SheafHom_\mathcal{O}(\mathcal{L}, \mathcal{O}) & \text{si} & n = -1\\
\mathcal{L} \otimes_\mathcal{O} \ldots \otimes_\mathcal{O} \mathcal{L}
& \text{si} & n > 0 \\
\mathcal{L}^{\otimes -1} \otimes_\mathcal{O} \ldots
\otimes_\mathcal{O} \mathcal{L}^{\otimes -1}
& \text{si} & n < -1
\end{matrix}
\right.
$$
voir le lemme \ref{sites-modules-lemma-constructions-invertible}. Avec cette définition,
on dispose d'isomorphismes canoniques
$\mathcal{L}^{\otimes n} \otimes_\mathcal{O}
\mathcal{L}^{\otimes m} \to
\mathcal{L}^{\otimes n + m}$ satisfaisant aux contraintes de commutativité
et d'associativité ; nous en omettons la formulation.

\begin{definition}

\label{sites-modules-definition-pic}
Soit $(\mathcal{C}, \mathcal{O})$ un site annelé. Le {\it groupe de Picard}
$\Pic(\mathcal{O})$ de ce site annelé est le groupe abélien des classes
d'isomorphisme de $\mathcal{O}$-modules inversibles, dont l'addition
correspond au produit tensoriel.
\end{definition}









\section{Modules de différentielles}

\label{sites-modules-section-differentials}

\noindent
Dans cette section, nous expliquons brièvement comment définir le module des
différentielles relatives d'un morphisme de topos annelés. Nous invitons le
lecteur à consulter la section correspondante du chapitre d'algèbre
commutative (Algèbre, section \ref{algebra-section-differentials}).

\begin{definition}

\label{sites-modules-definition-derivation}

Soit $\mathcal{C}$ un site. Soit
$\varphi : \mathcal{O}_1 \to \mathcal{O}_2$ un morphisme de faisceaux
d'anneaux, et soit $\mathcal{F}$ un $\mathcal{O}_2$-module. Une
{\it $\mathcal{O}_1$-dérivation}, ou plus précisément une
{\it $\varphi$-dérivation} à valeurs dans $\mathcal{F}$, est une application
$D : \mathcal{O}_2 \to \mathcal{F}$ additive, qui annule l'image de
$\mathcal{O}_1 \to \mathcal{O}_2$ et satisfait la {\it règle de Leibniz}

$$
D(ab) = aD(b) + D(a)b
$$
pour toutes sections locales $a, b$ de $\mathcal{O}_2$, là où elles sont
toutes deux définies. On note
$\text{Der}_{\mathcal{O}_1}(\mathcal{O}_2, \mathcal{F})$
l'ensemble des $\varphi$-dérivations à valeurs dans $\mathcal{F}$.

\end{definition}

\noindent
C'est l'analogue, pour les faisceaux, d'Algèbre, définition
\ref{sites-modules-definition-derivation}. Pour une dérivation
$D : \mathcal{O}_2 \to \mathcal{F}$ comme dans la définition, l'application
sur les sections globales
$$
D : \Gamma(\mathcal{O}_2) \longrightarrow \Gamma(\mathcal{F})
$$
est clairement une $\Gamma(\mathcal{O}_1)$-dérivation au sens algébrique.
Si $\alpha : \mathcal{F} \to \mathcal{G}$ est un morphisme de
$\mathcal{O}_2$-modules, on obtient un morphisme induit
$$
\text{Der}_{\mathcal{O}_1}(\mathcal{O}_2, \mathcal{F})
\longrightarrow
\text{Der}_{\mathcal{O}_1}(\mathcal{O}_2, \mathcal{G})
$$
donné par $D \mapsto \alpha \circ D$. On obtient ainsi un foncteur.

\begin{lemma}

\label{sites-modules-lemma-universal-module}

Soit $\mathcal{C}$ un site et soit
$\varphi : \mathcal{O}_1 \to \mathcal{O}_2$ un morphisme de faisceaux
d'anneaux. Le foncteur
$$
\textit{Mod}(\mathcal{O}_2) \longrightarrow \textit{Ab}, \quad
\mathcal{F} \longmapsto \text{Der}_{\mathcal{O}_1}(\mathcal{O}_2, \mathcal{F})
$$
est représentable.

\end{lemma}

\begin{proof}

La démonstration est identique à celle de l'énoncé analogue en algèbre.
Au cours de celle-ci, pour tout faisceau d'ensembles $\mathcal{F}$ sur
$\mathcal{C}$, notons $\mathcal{O}_2[\mathcal{F}]$ la faisceautisation du
préfaisceau $U \mapsto \mathcal{O}_2(U)[\mathcal{F}(U)]$, où le membre de
droite désigne le $\mathcal{O}_2(U)$-module libre sur l'ensemble
$\mathcal{F}(U)$. Pour $s \in \mathcal{F}(U)$, notons $[s]$ la section
correspondante de $\mathcal{O}_2[\mathcal{F}]$ au-dessus de $U$. Si
$\mathcal{F}$ est un faisceau de $\mathcal{O}_2$-modules, il existe un
morphisme canonique
$$
c : \mathcal{O}_2[\mathcal{F}] \longrightarrow \mathcal{F}
$$
qui, au niveau des préfaisceaux, est donné par la règle
$\sum f_s[s] \mapsto \sum f_s s$. Nous emploierons l'abréviation
$[s] \mapsto s$ pour décrire ce morphisme, et procéderons de même pour les
autres morphismes ci-dessous. Considérons le morphisme de
$\mathcal{O}_2$-modules

\begin{equation}
\label{sites-modules-equation-define-module-differentials}
\begin{matrix}
\mathcal{O}_2[\mathcal{O}_2 \times  \mathcal{O}_2] \oplus
\mathcal{O}_2[\mathcal{O}_2 \times  \mathcal{O}_2] \oplus
\mathcal{O}_2[\mathcal{O}_1] &
\longrightarrow &
\mathcal{O}_2[\mathcal{O}_2] \\
[(a, b)] \oplus [(f, g)] \oplus [h] & \longmapsto & [a + b] - [a] - [b] + \\
& & [fg] - g[f] - f[g] + \\
& & [\varphi(h)]
\end{matrix}
\end{equation}

où nous utilisons la notation abrégée précédente. Posons
$\Omega_{\mathcal{O}_2/\mathcal{O}_1}$ égal au conoyau de ce morphisme. Il
existe alors manifestement un morphisme de faisceaux d'ensembles
$$
\text{d} : \mathcal{O}_2 \longrightarrow \Omega_{\mathcal{O}_2/\mathcal{O}_1}
$$
qui envoie une section locale $f$ sur l'image de $[f]$ dans
$\Omega_{\mathcal{O}_2/\mathcal{O}_1}$. Par construction, $\text{d}$ est une
$\mathcal{O}_1$-dérivation. Soit maintenant $\mathcal{F}$ un faisceau de
$\mathcal{O}_2$-modules et soit $D : \mathcal{O}_2 \to \mathcal{F}$ une
$\mathcal{O}_1$-dérivation. On peut considérer le morphisme
$\mathcal{O}_2$-linéaire $\mathcal{O}_2[\mathcal{O}_2] \to \mathcal{F}$ qui
envoie $[g]$ sur $D(g)$. Par définition d'une dérivation, ce morphisme annule
les sections appartenant à l'image du morphisme
(\ref{sites-modules-equation-define-module-differentials}) et induit donc un morphisme
$$
\alpha_D : \Omega_{\mathcal{O}_2/\mathcal{O}_1} \longrightarrow \mathcal{F}
$$
Comme $D = \alpha_D \circ \text{d}$, le lemme est démontré.

\end{proof}

\begin{definition}

\label{sites-modules-definition-module-differentials}
Soit $\mathcal{C}$ un site. Soit
$\varphi : \mathcal{O}_1 \to \mathcal{O}_2$ un morphisme de faisceaux
d'anneaux. Le {\it module des différentielles} du morphisme d'anneaux
$\varphi$ est l'objet représentant le foncteur
$\mathcal{F} \mapsto \text{Der}_{\mathcal{O}_1}(\mathcal{O}_2, \mathcal{F})$,
dont l'existence résulte du lemme \ref{sites-modules-lemma-universal-module}. On le note
$\Omega_{\mathcal{O}_2/\mathcal{O}_1}$, et la {\it $\varphi$-dérivation
universelle} est notée
$\text{d} : \mathcal{O}_2 \to \Omega_{\mathcal{O}_2/\mathcal{O}_1}$.
\end{definition}

\noindent
Comme ce module et cette dérivation forment l'objet universel représentant un
foncteur, la notion est intrinsèque : elle ne dépend pas du choix du site
sous-jacent au topos annelé ; voir la section \ref{sites-modules-section-intrinsic}. Le
module $\Omega_{\mathcal{O}_2/\mathcal{O}_1}$ est le conoyau du morphisme
(\ref{sites-modules-equation-define-module-differentials}) de $\mathcal{O}_2$-modules. De plus,
$\text{d}$ est caractérisée par le fait que $\text{d}f$ est l'image de la
section locale $[f]$.

\begin{lemma}

\label{sites-modules-lemma-differentials-sheafify}
Soit $\mathcal{C}$ un site. Soit
$\varphi : \mathcal{O}_1 \to \mathcal{O}_2$ un morphisme de préfaisceaux
d'anneaux. Alors $\Omega_{\mathcal{O}_2^\#/\mathcal{O}_1^\#}$ est le
faisceau associé au préfaisceau
$U \mapsto \Omega_{\mathcal{O}_2(U)/\mathcal{O}_1(U)}$.
\end{lemma}

\begin{proof}
Considérons le morphisme (\ref{sites-modules-equation-define-module-differentials}). Il
existe un morphisme analogue de préfaisceaux dont la valeur en
$U \in \Ob(\mathcal{C})$ est
$$
\mathcal{O}_2(U)[\mathcal{O}_2(U) \times \mathcal{O}_2(U)] \oplus
\mathcal{O}_2(U)[\mathcal{O}_2(U) \times \mathcal{O}_2(U)] \oplus
\mathcal{O}_2(U)[\mathcal{O}_1(U)]
\longrightarrow
\mathcal{O}_2(U)[\mathcal{O}_2(U)]
$$
Le conoyau de ce morphisme vaut
$\Omega_{\mathcal{O}_2(U)/\mathcal{O}_1(U)}$ en $U$, par la construction du
module des différentielles dans Algèbre, définition
\ref{algebra-definition-differentials}. D'autre part, les faisceaux de
(\ref{sites-modules-equation-define-module-differentials}) sont les faisceautisations des
préfaisceaux ci-dessus. Le résultat découle donc de l'exactitude de la
faisceautisation.
\end{proof}

\begin{lemma}

\label{sites-modules-lemma-pullback-differentials}

Soit $f : \Sh(\mathcal{D}) \to \Sh(\mathcal{C})$ un morphisme de topos. Soit
$\varphi : \mathcal{O}_1 \to \mathcal{O}_2$ un morphisme de faisceaux
d'anneaux sur $\mathcal{C}$. Il existe alors une identification canonique
$f^{-1}\Omega_{\mathcal{O}_2/\mathcal{O}_1} =
\Omega_{f^{-1}\mathcal{O}_2/f^{-1}\mathcal{O}_1}$
compatible avec les dérivations universelles.

\end{lemma}

\begin{proof}
Cela résulte du fait que le faisceau
$\Omega_{\mathcal{O}_2/\mathcal{O}_1}$ est le conoyau du morphisme
(\ref{sites-modules-equation-define-module-differentials}), qu'il en va de même pour
$\Omega_{f^{-1}\mathcal{O}_2/f^{-1}\mathcal{O}_1}$, que le foncteur
$f^{-1}$ est exact et que $f^{-1}(\mathcal{O}_2[\mathcal{O}_2]) =
f^{-1}\mathcal{O}_2[f^{-1}\mathcal{O}_2]$, $f^{-1}(\mathcal{O}_2[\mathcal{O}_2 \times \mathcal{O}_2]) =
f^{-1}\mathcal{O}_2[f^{-1}\mathcal{O}_2 \times f^{-1}\mathcal{O}_2]$, et $f^{-1}(\mathcal{O}_2[\mathcal{O}_1]) =
f^{-1}\mathcal{O}_2[f^{-1}\mathcal{O}_1]$.
\end{proof}

\begin{lemma}

\label{sites-modules-lemma-localize-differentials}

Soit $\mathcal{C}$ un site. Soit
$\varphi : \mathcal{O}_1 \to \mathcal{O}_2$ un morphisme de faisceaux
d'anneaux. Pour tout objet $U$ de $\mathcal{C}$, il existe un isomorphisme
canonique
$$
\Omega_{\mathcal{O}_2/\mathcal{O}_1}|_U =
\Omega_{(\mathcal{O}_2|_U)/(\mathcal{O}_1|_U)}
$$
compatible avec les dérivations universelles.

\end{lemma}

\begin{proof}
C'est un cas particulier du lemme \ref{sites-modules-lemma-pullback-differentials}.
\end{proof}

\begin{lemma}

\label{sites-modules-lemma-functoriality-differentials}

Soit $\mathcal{C}$ un site. Soit
$$
\xymatrix{
\mathcal{O}_2 \ar[r]_\varphi & \mathcal{O}_2' \\
\mathcal{O}_1 \ar[r] \ar[u] & \mathcal{O}'_1 \ar[u]
}
$$
un diagramme commutatif de faisceaux d'anneaux sur $\mathcal{C}$. Le morphisme
$\mathcal{O}_2 \to \mathcal{O}'_2$, composé avec le morphisme
$\text{d} : \mathcal{O}'_2 \to \Omega_{\mathcal{O}'_2/\mathcal{O}'_1}$
est une $\mathcal{O}_1$-dérivation. On obtient donc un morphisme canonique de
$\mathcal{O}_2$-modules
$\Omega_{\mathcal{O}_2/\mathcal{O}_1} \to
\Omega_{\mathcal{O}'_2/\mathcal{O}'_1}$. Il est uniquement caractérisé par
la propriété que $\text{d}(f)$ s'envoie sur $\text{d}(\varphi(f))$ pour toute
section locale $f$ de $\mathcal{O}_2$. De cette manière, $\Omega_{-/-}$
devient un foncteur sur la catégorie des flèches de faisceaux d'anneaux.

\end{lemma}

\begin{proof}
Le lemme découle directement des définitions.
\end{proof}

\begin{lemma}

\label{sites-modules-lemma-differential-seq}
Dans le lemme \ref{sites-modules-lemma-functoriality-differentials}, supposons que
$\mathcal{O}_2 \to \mathcal{O}'_2$ soit surjectif, de noyau
$\mathcal{I} \subset \mathcal{O}_2$, et que
$\mathcal{O}_1 = \mathcal{O}'_1$. Il existe alors une suite exacte canonique
de $\mathcal{O}'_2$-modules
$$
\mathcal{I}/\mathcal{I}^2
\longrightarrow
\Omega_{\mathcal{O}_2/\mathcal{O}_1} \otimes_{\mathcal{O}_2} \mathcal{O}'_2
\longrightarrow
\Omega_{\mathcal{O}'_2/\mathcal{O}_1}
\longrightarrow
0
$$
Le morphisme de gauche est caractérisé par le fait qu'une section locale $f$
de $\mathcal{I}$ s'envoie sur $\text{d}f \otimes 1$.
\end{lemma}

\begin{proof}

Pour une section locale $f$ de $\mathcal{I}$, notons $\overline{f}$ l'image de
$f$ dans $\mathcal{I}/\mathcal{I}^2$. Pour montrer que le morphisme
$\overline{f} \mapsto \text{d}f \otimes 1$ est bien défini, il suffit de
vérifier que $\text{d} f_1f_2 \otimes 1 = 0$ lorsque $f_1, f_2$ sont des
sections locales de $\mathcal{I}$. Cela résulte immédiatement de la règle de
Leibniz :
$\text{d} f_1f_2 \otimes 1 =
(f_1 \text{d}f_2 + f_2 \text{d} f_1 )\otimes 1 =
\text{d}f_2 \otimes f_1 + \text{d}f_1 \otimes f_2 = 0$. Un calcul analogue
montre que ce morphisme est $\mathcal{O}'_2 = \mathcal{O}_2/\mathcal{I}$-linéaire.
Le morphisme de droite est celui du lemme
\ref{sites-modules-lemma-functoriality-differentials}.

\medskip\noindent
Pour établir l'exactitude, raisonnons comme suit. Soit
$\mathcal{O}''_2 \subset \mathcal{O}'_2$ le préfaisceau de
$\mathcal{O}_1$-algèbres dont la valeur en $U$ est l'image de
$\mathcal{O}_2(U) \to \mathcal{O}'_2(U)$. D'après Algèbre, lemme
\ref{algebra-lemma-differential-seq}, les suites
$$
\mathcal{I}(U)/\mathcal{I}(U)^2
\longrightarrow
\Omega_{\mathcal{O}_2(U)/\mathcal{O}_1(U)}
\otimes_{\mathcal{O}_2(U)} \mathcal{O}''_2(U)
\longrightarrow
\Omega_{\mathcal{O}''_2(U)/\mathcal{O}_1(U)}
\longrightarrow
0
$$
sont exactes pour tout objet $U$ de $\mathcal{C}$. L'exactitude de la
faisceautisation donne donc une suite exacte de faisceaux de
$(\mathcal{O}'_2)^\#$-modules. On conclut par le lemme
\ref{sites-modules-lemma-differentials-sheafify} et l'égalité
$(\mathcal{O}''_2)^\# = \mathcal{O}'_2$.
\end{proof}

\noindent
Voici une situation particulière où les dérivations apparaissent naturellement.

\begin{lemma}

\label{sites-modules-lemma-double-structure-gives-derivation}

Soit $\mathcal{C}$ un site. Soit
$\varphi : \mathcal{O}_1 \to \mathcal{O}_2$ un morphisme de faisceaux
d'anneaux. Considérons une suite exacte courte
$$
0 \to \mathcal{F} \to \mathcal{A} \to \mathcal{O}_2 \to 0
$$
Ici, $\mathcal{A}$ est un faisceau de $\mathcal{O}_1$-algèbres,
$\pi : \mathcal{A} \to \mathcal{O}_2$ est une surjection de faisceaux de
$\mathcal{O}_1$-algèbres et $\mathcal{F} = \Ker(\pi)$ est son noyau.
Supposons que $\mathcal{F}$ soit un faisceau d'idéaux de carré nul dans
$\mathcal{A}$. Le faisceau $\mathcal{F}$ possède alors une structure naturelle de
$\mathcal{O}_2$-module. Une section $s : \mathcal{O}_2 \to \mathcal{A}$ de
$\pi$ est un morphisme de $\mathcal{O}_1$-algèbres tel que
$\pi \circ s = \text{id}$. Étant données une section quelconque
$s : \mathcal{O}_2 \to \mathcal{A}$ de $\pi$ et une $\varphi$-dérivation
$D : \mathcal{O}_2 \to \mathcal{F}$, le morphisme
$$
s + D : \mathcal{O}_2 \to \mathcal{A}
$$
est une section de $\pi$, et toute section $s'$ s'écrit de manière unique
$s + D$ pour une $\varphi$-dérivation $D$.

\end{lemma}

\begin{proof}

Rappelons que la structure de $\mathcal{O}_2$-module sur $\mathcal{F}$ est
donnée par $h \tau = \tilde h \tau$ (le produit étant pris dans
$\mathcal{A}$), où $h$ est une section locale de $\mathcal{O}_2$,
$\tilde h$ un relèvement local de $h$ en une section locale de $\mathcal{A}$,
et $\tau$ une section locale de $\mathcal{F}$. En présence de $s$, on peut
prendre $\tilde h = s(h)$. Pour vérifier que $s + D$ est un morphisme de
faisceaux d'anneaux, calculons

\begin{eqnarray*}
(s + D)(ab) & = & s(ab) + D(ab) \\
& = & s(a)s(b) + aD(b) + D(a)b \\
& = & s(a) s(b) + s(a)D(b) + D(a)s(b) \\
& = & (s(a) + D(a))(s(b) + D(b))
\end{eqnarray*}

par la règle de Leibniz. De même, $s + D$ est un morphisme de
$\mathcal{O}_1$-algèbres parce que $D$ est une $\mathcal{O}_1$-dérivation.
Réciproquement, pour une section $s'$, posons $D = s' - s$. Les détails sont
laissés au lecteur.

\end{proof}

\begin{definition}

\label{sites-modules-definition-sheaf-differentials}

Soient $X = (\Sh(\mathcal{C}), \mathcal{O})$ et
$Y = (\Sh(\mathcal{C}'), \mathcal{O}')$ des topos annelés. Soit
$(f, f^\sharp) : X \to Y$ un morphisme de topos annelés. Dans cette situation,

\begin{enumerate}

\item pour un faisceau $\mathcal{F}$ de $\mathcal{O}$-modules, une
{\it $Y$-dérivation} $D : \mathcal{O} \to \mathcal{F}$ est simplement une
$f^\sharp$-dérivation, et
\item le {\it faisceau des différentielles $\Omega_{X/Y}$ de $X$ sur $Y$}
est le module des différentielles du morphisme
$f^\sharp : f^{-1}\mathcal{O}' \to \mathcal{O}$, voir définition \ref{sites-modules-definition-module-differentials}.

\end{enumerate}

Ainsi, $\Omega_{X/Y}$ est muni d'une {\it dérivation $Y$-universelle}

$\text{d}_{X/Y} : \mathcal{O} \longrightarrow \Omega_{X/Y}$. Nous écrivons parfois $\Omega_{X/Y} = \Omega_f$.

\end{definition}

\noindent
Rappelons que $f^\sharp : f^{-1}\mathcal{O}' \to \mathcal{O}$, ce qui donne
un sens à cette définition.

\begin{lemma}

\label{sites-modules-lemma-functoriality-differentials-ringed-topoi}

Soient
$X = (\Sh(\mathcal{C}_X), \mathcal{O}_X)$,
$Y = (\Sh(\mathcal{C}_Y), \mathcal{O}_Y)$,
$X' = (\Sh(\mathcal{C}_{X'}), \mathcal{O}_{X'})$,
$Y' = (\Sh(\mathcal{C}_{Y'}), \mathcal{O}_{Y'})$ des topos annelés. Soit
$$
\xymatrix{
X' \ar[d] \ar[r]_f & X \ar[d] \\
Y' \ar[r] & Y
}
$$
un diagramme commutatif de morphismes de topos annelés. Le morphisme
$f^\sharp : \mathcal{O}_X \to f_*\mathcal{O}_{X'}$, composé avec le morphisme
$f_*\text{d}_{X'/Y'} : f_*\mathcal{O}_{X'} \to f_*\Omega_{X'/Y'}$ est un
$Y$-dérivation. On obtient donc un morphisme canonique de
$\mathcal{O}_X$-modules $\Omega_{X/Y} \to f_*\Omega_{X'/Y'}$ et, par
adjonction entre $f_*$ et $f^*$, un homomorphisme canonique de
$\mathcal{O}_{X'}$-modules
$$
c_f : f^*\Omega_{X/Y} \longrightarrow \Omega_{X'/Y'}.
$$
Il est uniquement caractérisé par la propriété que
$f^*\text{d}_{X/Y}(t)$ s'envoie sur $\text{d}_{X'/Y'}(f^* t)$ pour toute
section locale $t$ de $\mathcal{O}_X$.

\end{lemma}

\begin{proof}

Tout est clair, sauf la dernière assertion. Précisons son sens. Soient
$U \in \Ob(\mathcal{C}_X)$ et $t \in \mathcal{O}_X(U)$ ; c'est en ce sens que
$t$ est une section locale de $\mathcal{O}_X$. On peut voir $t$ comme un
morphisme de faisceaux d'ensembles $t : h_U^\# \to \mathcal{O}_X$. On a alors
$f^{-1}t : f^{-1}h_U^\# \to f^{-1}\mathcal{O}_X$. Par $f^*t$, on entend le
composé
$$
\xymatrix{
f^{-1}h_U^\# \ar[rr]^{f^{-1}t} \ar@/^4ex/[rrrr]^{f^*t} & &
f^{-1}\mathcal{O}_X \ar[rr]^{f^\sharp} & &
\mathcal{O}_{X'}
}
$$
Notons que $\text{d}_{X/Y}(t) \in \Omega_{X/Y}(U)$. On peut donc voir
$\text{d}_{X/Y}(t)$ comme un morphisme
$\text{d}_{X/Y}(t) : h_U^\# \to \Omega_{X/Y}$. On a alors
$f^{-1}\text{d}_{X/Y}(t) : f^{-1}h_U^\# \to f^{-1}\Omega_{X/Y}$. Par
$f^*\text{d}_{X/Y}(t)$, on entend le composé
$$
\xymatrix{
f^{-1}h_U^\#
\ar[rr]^{f^{-1}\text{d}_{X/Y}(t)}
\ar@/^4ex/[rrrr]^{f^*\text{d}_{X/Y}(t)} & &
f^{-1}\Omega_{X/Y} \ar[rr]^{1 \otimes \text{id}} & &
f^*\Omega_{X/Y}
}
$$
L'énoncé du lemme signifie alors que l'on a
$$
c_f \circ f^*\text{d}_{X/Y}(t) = \text{d}_{X'/Y'}(f^*t)
$$
comme morphismes de $f^{-1}h_U^\#$ vers $\Omega_{X'/Y'}$. Nous omettons la
vérification que le morphisme $c_f$ défini dans le lemme satisfait cette
propriété. Indication : le premier morphisme
$c'_f : \Omega_{X/Y} \to f_*\Omega_{X'/Y'}$ satisfait
$c'_f(\text{d}_{X/Y}(t)) = f_*\text{d}_{X'/Y'}(f^\sharp(t))$ en tant que
sections de $f_*\Omega_{X'/Y'}$ au-dessus de $U$ ; il suffit de traduire cette
égalité au moyen de l'adjonction. Cette propriété caractérise uniquement
$c_f$, car les images de $f^*\text{d}_{X/Y}(t)$ engendrent le
$\mathcal{O}_{X'}$-module $f^*\Omega_{X/Y}$, puisque les sections locales
$\text{d}_{X/Y}(t)$ engendrent le $\mathcal{O}_X$-module $\Omega_{X/Y}$.

\end{proof}
\section{Opérateurs différentiels d'ordre fini}
\label{sites-modules-section-differential-operators}

\noindent
Dans cette section, nous introduisons les opérateurs différentiels d'ordre fini.
Nous conseillons au lecteur de consulter la section correspondante
du chapitre consacré à l'algèbre commutative
(Algèbre, section \ref{algebra-section-differential-operators}).

\begin{definition}
\label{sites-modules-definition-differential-operators}
Soit $\mathcal{C}$ un site. Soit $\varphi : \mathcal{O}_1 \to \mathcal{O}_2$
un homomorphisme de faisceaux d'anneaux. Soit $k \geq 0$ un entier.
Soient $\mathcal{F}$ et $\mathcal{G}$ des faisceaux de $\mathcal{O}_2$-modules.
Un {\it opérateur différentiel $D : \mathcal{F} \to \mathcal{G}$ d'ordre $k$}
est une application $\mathcal{O}_1$-linéaire telle que, pour toute section locale
$g$ de $\mathcal{O}_2$, l'application $s \mapsto D(gs) - gD(s)$ soit un
opérateur différentiel d'ordre $k - 1$. Pour le cas initial $k = 0$,
nous définissons un opérateur différentiel d'ordre $0$ comme une application
$\mathcal{O}_2$-linéaire.
\end{definition}

\noindent
Si $D : \mathcal{F} \to \mathcal{G}$ est un opérateur différentiel d'ordre $k$,
alors, pour toute section locale $g$ de $\mathcal{O}_2$, l'application $gD$ est
un opérateur différentiel d'ordre $k$. La somme de deux opérateurs différentiels
d'ordre $k$ en est encore un. Par conséquent, l'ensemble de tous ces opérateurs
$$
\text{Diff}^k(\mathcal{F}, \mathcal{G}) =
\text{Diff}^k_{\mathcal{O}_2/\mathcal{O}_1}(\mathcal{F}, \mathcal{G})
$$
est un $\Gamma(\mathcal{C}, \mathcal{O}_2)$-module. On a
$$
\text{Diff}^0(\mathcal{F}, \mathcal{G}) \subset
\text{Diff}^1(\mathcal{F}, \mathcal{G}) \subset
\text{Diff}^2(\mathcal{F}, \mathcal{G}) \subset \ldots
$$
La règle qui associe à $U \in \Ob(\mathcal{C})$ le module des
opérateurs différentiels $D : \mathcal{F}|_U \to \mathcal{G}|_U$
d'ordre $k$ est un faisceau de $\mathcal{O}_2$-modules sur le site $\mathcal{C}$.
On obtient ainsi un faisceau d'opérateurs différentiels (si nous en avons un
jour besoin, nous ajouterons ici une définition).

\begin{lemma}
\label{sites-modules-lemma-composition-differential-operators}
Soit $\mathcal{C}$ un site.
Soit $\mathcal{O}_1 \to \mathcal{O}_2$ un morphisme de faisceaux d'anneaux.
Soient $\mathcal{E}, \mathcal{F}, \mathcal{G}$ des faisceaux de
$\mathcal{O}_2$-modules.
Si $D : \mathcal{E} \to \mathcal{F}$ et $D' : \mathcal{F} \to \mathcal{G}$
sont des opérateurs différentiels d'ordres $k$ et $k'$, alors $D' \circ D$ est
un opérateur différentiel d'ordre $k + k'$.
\end{lemma}

\begin{proof}
Soit $g$ une section locale de $\mathcal{O}_2$.
L'application qui envoie une section locale $x$ de $\mathcal{E}$ sur
$$
D'(D(gx)) - gD'(D(x)) = D'(D(gx)) - D'(gD(x)) + D'(gD(x)) - gD'(D(x))
$$
est une somme de deux composées d'opérateurs différentiels d'ordre inférieur.
Le lemme en résulte par récurrence sur $k + k'$.
\end{proof}

\begin{lemma}
\label{sites-modules-lemma-module-principal-parts}
Soit $\mathcal{C}$ un site.
Soit $\mathcal{O}_1 \to \mathcal{O}_2$ un morphisme de faisceaux d'anneaux.
Soit $\mathcal{F}$ un faisceau de $\mathcal{O}_2$-modules.
Soit $k \geq 0$. Il existe un faisceau de $\mathcal{O}_2$-modules
$\mathcal{P}^k_{\mathcal{O}_2/\mathcal{O}_1}(\mathcal{F})$
et un isomorphisme canonique
$$
\text{Diff}^k_{\mathcal{O}_2/\mathcal{O}_1}(\mathcal{F}, \mathcal{G}) =
\Hom_{\mathcal{O}_2}(
\mathcal{P}^k_{\mathcal{O}_2/\mathcal{O}_1}(\mathcal{F}), \mathcal{G})
$$
fonctoriel par rapport au $\mathcal{O}_2$-module $\mathcal{G}$.
\end{lemma}

\begin{proof}
L'existence résulte d'arguments généraux de théorie des catégories
(insérer ici une référence ultérieure), mais nous donnerons également une
construction directe, car elle sera utile dans les démonstrations à venir.
Nous emploierons librement les notations introduites dans la démonstration du
lemme \ref{sites-modules-lemma-universal-module}.
Étant donné un opérateur différentiel $D : \mathcal{F} \to \mathcal{G}$,
nous obtenons une application $\mathcal{O}_2$-linéaire
$L_D : \mathcal{O}_2[\mathcal{F}] \to \mathcal{G}$
qui envoie $[m]$ sur $D(m)$. Si $D$ est d'ordre $0$,
alors $L_D$ annule les sections locales
$$
[m + m'] - [m] - [m'],\quad
g_0[m] - [g_0m]
$$
où $g_0$ est une section locale de $\mathcal{O}_2$ et $m, m'$
sont des sections locales de $\mathcal{F}$. Si $D$ est d'ordre $1$, alors $L_D$
annule les sections locales
$$
[m + m'] - [m] - [m'],\quad
f[m] - [fm], \quad
g_0g_1[m] - g_0[g_1m] - g_1[g_0m] + [g_1g_0m]
$$
où $f$ est une section locale de $\mathcal{O}_1$,
$g_0, g_1$ sont des sections locales de $\mathcal{O}_2$, et
$m, m'$ sont des sections locales de $\mathcal{F}$.
Si $D$ est d'ordre $k$, alors $L_D$ annule les sections locales
$[m + m'] - [m] - [m']$, $f[m] - [fm]$, ainsi que les sections locales
$$
g_0g_1\ldots g_k[m] - \sum g_0 \ldots \hat g_i \ldots g_k[g_im] + \ldots
+(-1)^{k + 1}[g_0\ldots g_km]
$$
Réciproquement, si $L : \mathcal{O}_2[\mathcal{F}] \to \mathcal{G}$ est une
application $\mathcal{O}_2$-linéaire qui annule toutes les sections locales
énumérées dans la phrase précédente, alors $m \mapsto L([m])$ est un
opérateur différentiel d'ordre $k$. On voit donc que
$\mathcal{P}^k_{\mathcal{O}_2/\mathcal{O}_1}(\mathcal{F})$
est le quotient de $\mathcal{O}_2[\mathcal{F}]$ par le
sous-$\mathcal{O}_2$-module engendré par ces sections locales.
\end{proof}

\begin{definition}
\label{sites-modules-definition-module-principal-parts}
Soit $\mathcal{C}$ un site.
Soit $\mathcal{O}_1 \to \mathcal{O}_2$ un morphisme de faisceaux d'anneaux.
Soit $\mathcal{F}$ un faisceau de $\mathcal{O}_2$-modules.
Le module $\mathcal{P}^k_{\mathcal{O}_2/\mathcal{O}_1}(\mathcal{F})$
construit dans le lemme \ref{sites-modules-lemma-module-principal-parts}
est appelé le {\it module des parties principales d'ordre $k$} de $\mathcal{F}$.
\end{definition}

\noindent
Remarquons que les inclusions
$$
\text{Diff}^0(\mathcal{F}, \mathcal{G}) \subset
\text{Diff}^1(\mathcal{F}, \mathcal{G}) \subset
\text{Diff}^2(\mathcal{F}, \mathcal{G}) \subset \ldots
$$
correspondent, par le lemme de Yoneda (Catégories, lemme
\ref{categories-lemma-yoneda}), aux surjections
$$
\ldots \to \mathcal{P}^2_{\mathcal{O}_2/\mathcal{O}_1}(\mathcal{F})
\to \mathcal{P}^1_{\mathcal{O}_2/\mathcal{O}_1}(\mathcal{F})
\to \mathcal{P}^0_{\mathcal{O}_2/\mathcal{O}_1}(\mathcal{F}) = \mathcal{F}
$$

\begin{lemma}
\label{sites-modules-lemma-differential-operators-sheafify}
Soit $\mathcal{C}$ un site. Soit $\mathcal{O}_1 \to \mathcal{O}_2$
un homomorphisme de préfaisceaux d'anneaux. Soit $\mathcal{F}$ un préfaisceau
de $\mathcal{O}_2$-modules. Alors
$\mathcal{P}^k_{\mathcal{O}_2^\#/\mathcal{O}_1^\#}(\mathcal{F}^\#)$
est le faisceau associé au préfaisceau
$U \mapsto P^k_{\mathcal{O}_2(U)/\mathcal{O}_1(U)}(\mathcal{F}(U))$.
\end{lemma}

\begin{proof}
On peut le démontrer exactement comme pour le faisceau des différentielles
dans le lemme \ref{sites-modules-lemma-differentials-sheafify}.
Une approche peut-être plus élégante consiste à appliquer directement la
propriété universelle du lemme \ref{sites-modules-lemma-module-principal-parts} pour obtenir
l'égalité. Nous omettons les détails.
\end{proof}

\begin{lemma}
\label{sites-modules-lemma-sequence-of-principal-parts}
Soit $\mathcal{C}$ un site. Soit $\mathcal{O}_1 \to \mathcal{O}_2$
un homomorphisme de faisceaux d'anneaux. Soit $\mathcal{F}$ un faisceau
de $\mathcal{O}_2$-modules. Il existe une suite exacte courte canonique
$$
0 \to
\Omega_{\mathcal{O}_2/\mathcal{O}_1} \otimes_{\mathcal{O}_2} \mathcal{F} \to
\mathcal{P}^1_{\mathcal{O}_2/\mathcal{O}_1}(\mathcal{F}) \to
\mathcal{F} \to 0
$$
fonctorielle en $\mathcal{F}$, appelée la {\it suite des parties principales}.
\end{lemma}

\begin{proof}
Cela résulte de la version d'algèbre commutative
(Algèbre, lemme \ref{algebra-lemma-sequence-of-principal-parts})
et des lemmes \ref{sites-modules-lemma-differentials-sheafify} et
\ref{sites-modules-lemma-differential-operators-sheafify}.
\end{proof}

\begin{remark}
\label{sites-modules-remark-functoriality-principal-parts}
Soit $\mathcal{C}$ un site. Supposons donné un diagramme commutatif de
faisceaux d'anneaux
$$
\xymatrix{
\mathcal{B} \ar[r] & \mathcal{B}' \\
\mathcal{A} \ar[u] \ar[r] & \mathcal{A}' \ar[u]
}
$$
un $\mathcal{B}$-module $\mathcal{F}$, un $\mathcal{B}'$-module $\mathcal{F}'$,
et une application $\mathcal{B}$-linéaire $\mathcal{F} \to \mathcal{F}'$.
On obtient alors un système compatible d'applications de modules
$$
\xymatrix{
\ldots \ar[r] &
\mathcal{P}^2_{\mathcal{B}'/\mathcal{A}'}(\mathcal{F}') \ar[r] &
\mathcal{P}^1_{\mathcal{B}'/\mathcal{A}'}(\mathcal{F}') \ar[r] &
\mathcal{P}^0_{\mathcal{B}'/\mathcal{A}'}(\mathcal{F}') \\
\ldots \ar[r] &
\mathcal{P}^2_{\mathcal{B}/\mathcal{A}}(\mathcal{F}) \ar[r] \ar[u] &
\mathcal{P}^1_{\mathcal{B}/\mathcal{A}}(\mathcal{F}) \ar[r] \ar[u] &
\mathcal{P}^0_{\mathcal{B}/\mathcal{A}}(\mathcal{F}) \ar[u]
}
$$
Ces applications sont compatibles avec toute composition ultérieure
d'applications de ce type. Le moyen le plus simple de le voir consiste à
utiliser la description des modules
$\mathcal{P}^k_{\mathcal{B}/\mathcal{A}}(\mathcal{M})$ en termes de
générateurs et relations (locaux) donnée dans la démonstration du
lemme \ref{sites-modules-lemma-module-principal-parts}; cela se voit aussi directement à
partir de la propriété universelle de ces modules. De plus, ces applications
sont compatibles avec les suites exactes courtes du
lemme \ref{sites-modules-lemma-sequence-of-principal-parts}.
\end{remark}








\section{Le complexe cotangent naïf}
\label{sites-modules-section-netherlander}

\noindent
Cette section est l'analogue de
Algèbre, section \ref{algebra-section-netherlander},
et de
Modules, section \ref{modules-section-netherlander}.
Nous conseillons au lecteur de lire d'abord ces sections.

\medskip\noindent
Soit $\mathcal{C}$ un site. Soit $\mathcal{A} \to \mathcal{B}$ un
homomorphisme de faisceaux d'anneaux sur $\mathcal{C}$. Dans cette section,
pour tout faisceau d'ensembles $\mathcal{E}$ sur $\mathcal{C}$, nous notons
$\mathcal{A}[\mathcal{E}]$ le faisceautisé du préfaisceau
$U \mapsto \mathcal{A}(U)[\mathcal{E}(U)]$. Ici,
$\mathcal{A}(U)[\mathcal{E}(U)]$
désigne l'algèbre de polynômes sur $\mathcal{A}(U)$ dont les variables
correspondent aux éléments de $\mathcal{E}(U)$.
Nous notons $[e] \in \mathcal{A}(U)[\mathcal{E}(U)]$ la variable
correspondant à $e \in \mathcal{E}(U)$.
Il existe une surjection canonique de $\mathcal{A}$-algèbres
\begin{equation}
\label{sites-modules-equation-canonical-presentation}
\mathcal{A}[\mathcal{B}] \longrightarrow \mathcal{B},\quad [b] \longmapsto b
\end{equation}
dont nous notons le noyau $\mathcal{I} \subset \mathcal{A}[\mathcal{B}]$.
On observe immédiatement que $\mathcal{I}$ est engendré par les
sections locales $[b][b'] - [bb']$ et $[a] - a$. D'après le
lemme \ref{sites-modules-lemma-differential-seq}, il existe une application canonique
\begin{equation}
\label{sites-modules-equation-naive-cotangent-complex}
\mathcal{I}/\mathcal{I}^2
\longrightarrow
\Omega_{\mathcal{A}[\mathcal{B}]/\mathcal{A}}
\otimes_{\mathcal{A}[\mathcal{B}]} \mathcal{B}
\end{equation}
dont le conoyau est canoniquement isomorphe à $\Omega_{\mathcal{B}/\mathcal{A}}$.

\begin{definition}
\label{sites-modules-definition-naive-cotangent-complex}
Soit $\mathcal{C}$ un site. Soit $\mathcal{A} \to \mathcal{B}$ un
homomorphisme de faisceaux d'anneaux sur $\mathcal{C}$.
Le {\it complexe cotangent naïf} $\NL_{\mathcal{B}/\mathcal{A}}$
est le complexe de chaînes (\ref{sites-modules-equation-naive-cotangent-complex})
$$
\NL_{\mathcal{B}/\mathcal{A}} =
\left(\mathcal{I}/\mathcal{I}^2
\longrightarrow
\Omega_{\mathcal{A}[\mathcal{B}]/\mathcal{A}}
\otimes_{\mathcal{A}[\mathcal{B}]} \mathcal{B}\right)
$$
où $\mathcal{I}/\mathcal{I}^2$ est placé en degré $-1$ et
$\Omega_{\mathcal{A}[\mathcal{B}]/\mathcal{A}}
\otimes_{\mathcal{A}[\mathcal{B}]} \mathcal{B}$
est placé en degré $0$.
\end{definition}

\noindent
Cette construction possède une fonctorialité analogue à celle examinée dans le
lemme \ref{sites-modules-lemma-functoriality-differentials} pour les modules de
différentielles. Plus précisément, étant donné un diagramme commutatif
\begin{equation}
\label{sites-modules-equation-commutative-square-sheaves}
\vcenter{
\xymatrix{
\mathcal{B} \ar[r] & \mathcal{B}' \\
\mathcal{A} \ar[u] \ar[r] & \mathcal{A}' \ar[u]
}
}
\end{equation}
de faisceaux d'anneaux sur $\mathcal{C}$, il existe une application canonique
$\mathcal{B}$-linéaire de complexes
$$
\NL_{\mathcal{B}/\mathcal{A}} \longrightarrow \NL_{\mathcal{B}'/\mathcal{A}'}
$$
En effet, les applications du diagramme commutatif donnent une application canonique
$\mathcal{A}[\mathcal{B}] \to \mathcal{A}'[\mathcal{B}']$
qui envoie $\mathcal{I}$ dans
$\mathcal{I}' = \Ker(\mathcal{A}'[\mathcal{B}'] \to \mathcal{B}')$.
On obtient ainsi une application
$\mathcal{I}/\mathcal{I}^2 \to \mathcal{I}'/(\mathcal{I}')^2$ et une
application entre les modules de différentielles; ensemble, elles donnent
l'application cherchée entre les complexes cotangents naïfs.

\medskip\noindent
On peut choisir une autre présentation de $\mathcal{B}$ comme quotient d'une
algèbre de polynômes sur $\mathcal{A}$ et obtenir néanmoins le même objet de
$D(\mathcal{B})$. Pour l'expliquer, supposons que $\mathcal{E}$ soit un
faisceau d'ensembles sur $\mathcal{C}$ et que
$\alpha : \mathcal{E} \to \mathcal{B}$ soit une application de faisceaux
d'ensembles. On obtient alors un homomorphisme de $\mathcal{A}$-algèbres
$\mathcal{A}[\mathcal{E}] \to \mathcal{B}$. Supposons cette application
surjective et notons $\mathcal{J} \subset \mathcal{A}[\mathcal{E}]$ son noyau.
Posons
$$
\NL(\alpha) = \left(
\mathcal{J}/\mathcal{J}^2
\longrightarrow
\Omega_{\mathcal{A}[\mathcal{E}]/\mathcal{A}}
\otimes_{\mathcal{A}[\mathcal{E}]} \mathcal{B}\right)
$$
Voici le résultat.

\begin{lemma}
\label{sites-modules-lemma-NL-up-to-qis}
Dans la situation ci-dessus, il existe un isomorphisme canonique
$\NL(\alpha) = \NL_{\mathcal{B}/\mathcal{A}}$ dans $D(\mathcal{B})$.
\end{lemma}

\begin{proof}
Remarquons que
$\NL_{\mathcal{B}/\mathcal{A}} = \NL(\text{id}_\mathcal{B})$.
Il suffit donc de montrer que, pour deux applications
$\alpha_i : \mathcal{E}_i \to \mathcal{B}$ comme ci-dessus, il existe un
quasi-isomorphisme canonique $\NL(\alpha_1) = \NL(\alpha_2)$ dans
$D(\mathcal{B})$. Pour cela, posons
$\mathcal{E} = \mathcal{E}_1 \amalg \mathcal{E}_2$ et
$\alpha = \alpha_1 \amalg \alpha_2 : \mathcal{E} \to \mathcal{B}$.
Posons
$\mathcal{J}_i = \Ker(\mathcal{A}[\mathcal{E}_i] \to \mathcal{B})$
et
$\mathcal{J} = \Ker(\mathcal{A}[\mathcal{E}] \to \mathcal{B})$.
On obtient des applications
$\mathcal{A}[\mathcal{E}_i] \to \mathcal{A}[\mathcal{E}]$ qui envoient
$\mathcal{J}_i$ dans $\mathcal{J}$, donc des applications canoniques de complexes
$$
\NL(\alpha_i) \longrightarrow \NL(\alpha)
$$
et il suffit de montrer que ces applications sont des quasi-isomorphismes.
Pour cela, raisonnons comme suit. Tout d'abord, remarquons que
$H^0(\NL(\alpha_i)) = \Omega_{\mathcal{B}/\mathcal{A}}$ et
$H^0(\NL(\alpha)) = \Omega_{\mathcal{B}/\mathcal{A}}$ d'après
le lemme \ref{sites-modules-lemma-differential-seq}; l'application est donc un isomorphisme
sur les faisceaux de cohomologie en degré $0$.
De même, nous affirmons que $H^{-1}(\NL(\alpha_i))$ et
$H^{-1}(\NL(\alpha))$ sont les faisceaux associés au préfaisceau
$U \mapsto H_1(L_{\mathcal{B}(U)/\mathcal{A}(U)})$, où
où $H_1(L_{-/-})$ est défini comme dans
Algèbre, définition \ref{algebra-definition-naive-cotangent-complex}.
Cette assertion achève la démonstration.

\medskip\noindent
Démonstration de l'assertion. Soit
$\alpha : \mathcal{E} \to \mathcal{B}$ comme ci-dessus. Soit
$\mathcal{B}' \subset \mathcal{B}$ le sous-préfaisceau de
$\mathcal{A}$-algèbres dont la valeur sur $U$ est l'image de
$\mathcal{A}(U)[\mathcal{E}(U)] \to \mathcal{B}(U)$. Soit $\mathcal{I}'$
le préfaisceau dont la valeur sur $U$ est le noyau de
$\mathcal{A}(U)[\mathcal{E}(U)] \to \mathcal{B}(U)$. Alors $\mathcal{I}$
est le faisceautisé de $\mathcal{I}'$, et $\mathcal{B}$ celui de
$\mathcal{B}'$. De même, $H^{-1}(\NL(\alpha))$ est le faisceautisé du préfaisceau
$$
U \longmapsto
\Ker(\mathcal{I}'(U)/\mathcal{I}'(U)^2 \to
\Omega_{\mathcal{A}(U)[\mathcal{E}(U)]/\mathcal{A}(U)}
\otimes_{\mathcal{A}(U)[\mathcal{E}(U)]} \mathcal{B}'(U))
$$
d'après le lemme \ref{sites-modules-lemma-differentials-sheafify}.
Par Algèbre, lemme \ref{algebra-lemma-NL-homotopy}, on en déduit que
$H^{-1}(\NL(\alpha))$ est le faisceau associé au préfaisceau
$U \mapsto H_1(L_{\mathcal{B}'(U)/\mathcal{A}(U)})$. Il reste donc à montrer
que les applications
$H_1(L_{\mathcal{B}'(U)/\mathcal{A}(U)}) \to
H_1(L_{\mathcal{B}(U)/\mathcal{A}(U)})$ induisent un isomorphisme
$\mathcal{H}'_1 \to \mathcal{H}_1$ entre faisceautisés.

\medskip\noindent
Injectivité de $\mathcal{H}'_1 \to \mathcal{H}_1$. Soit
$f \in H_1(L_{\mathcal{B}'(U)/\mathcal{A}(U)})$ d'image nulle dans
$\mathcal{H}_1(U)$. Il faut montrer que $f$ a une image nulle dans
$\mathcal{H}'_1(U)$ est nulle. L'hypothèse signifie qu'il existe un recouvrement
$\{U_i \to U\}$ tel que $f$ s'annule dans
$H_1(L_{\mathcal{B}(U_i)/\mathcal{A}(U_i)})$ pour tout $i$.
En remplaçant $U$ par $U_i$, on se ramène au cas où $f$ s'annule dans
$H_1(L_{\mathcal{B}(U)/\mathcal{A}(U)})$.
Par Algèbre, lemme \ref{algebra-lemma-colimits-NL}, on peut trouver une
sous-algèbre de type fini $\mathcal{B}'(U) \subset B \subset \mathcal{B}(U)$
telle que $f$ s'annule dans $H_1(L_{B/\mathcal{A}(U)})$.
Comme $\mathcal{B} = (\mathcal{B}')^\#$, il existe un recouvrement
$\{U_i \to U\}$ tel que $B \to \mathcal{B}(U_i)$ se factorise par
$\mathcal{B}'(U_i)$. Ainsi $f$ s'annule dans
$H_1(L_{\mathcal{B}'(U_i)/\mathcal{A}(U_i)})$, comme voulu.

\medskip\noindent
La surjectivité de $\mathcal{H}'_1 \to \mathcal{H}_1$ se démontre
exactement de la même manière.
\end{proof}

\begin{lemma}
\label{sites-modules-lemma-pullback-NL}
Soit $f : \Sh(\mathcal{C}) \to \Sh(\mathcal{D})$ un morphisme de topos.
Soit $\mathcal{A} \to \mathcal{B}$ un homomorphisme de faisceaux d'anneaux
sur $\mathcal{D}$. Alors $f^{-1}\NL_{\mathcal{B}/\mathcal{A}} =
\NL_{f^{-1}\mathcal{B}/f^{-1}\mathcal{A}}$.
\end{lemma}

\begin{proof}
La démonstration est omise. Indication : appliquer le
lemme \ref{sites-modules-lemma-pullback-differentials}.
\end{proof}

\noindent
Le complexe cotangent d'un morphisme de topos annelés se définit à l'aide du
complexe cotangent introduit ci-dessus.

\begin{definition}
\label{sites-modules-definition-cotangent-complex-morphism-ringed-topoi}
Soient $X = (\Sh(\mathcal{C}), \mathcal{O})$ et
$Y = (\Sh(\mathcal{C}'), \mathcal{O}')$ des topos annelés.
Soit $(f, f^\sharp) : X \to Y$ un morphisme de topos annelés.
Le {\it complexe cotangent naïf} $\NL_f = \NL_{X/Y}$ de ce morphisme de
topos annelés est $\NL_{\mathcal{O}/f^{-1}\mathcal{O}'}$.
Nous écrivons parfois $\NL_{X/Y} = \NL_{\mathcal{O}/\mathcal{O}'}$.
\end{definition}









\section{Germes des modules}
\label{sites-modules-section-stalks}

\noindent
Il faut être quelque peu prudent lorsqu'on prend les germes aux points, car la
colimite qui définit un germe (voir Sites, équation
\ref{sites-equation-stalk}) peut ne pas être filtrante\footnote{Bien entendu,
dans presque tous les cas naturels, la colimite est filtrante et une partie de
la discussion de cette section peut être simplifiée.}. En revanche, par
définition d'un point d'un site, le foncteur des germes est exact et commute aux
colimites arbitraires. Autrement dit, il se comporte exactement comme si la
colimite était filtrante.

\begin{lemma}
\label{sites-modules-lemma-stalk-exact}
Soit $\mathcal{C}$ un site.
Soit $p$ un point de $\mathcal{C}$.
\begin{enumerate}
\item On a $(\mathcal{F}^\#)_p = \mathcal{F}_p$
pour tout préfaisceau d'ensembles sur $\mathcal{C}$.
\item Le foncteur des germes
$\Sh(\mathcal{C}) \to \textit{Ensembles}$,
$\mathcal{F} \mapsto \mathcal{F}_p$ est exact (voir
Catégories, définition \ref{categories-definition-exact})
et commute aux colimites arbitraires.
\item Le foncteur des germes
$\textit{PSh}(\mathcal{C}) \to \textit{Ensembles}$,
$\mathcal{F} \mapsto \mathcal{F}_p$ est exact (voir
Catégories, définition \ref{categories-definition-exact})
et commute aux colimites arbitraires.
\end{enumerate}
\end{lemma}

\begin{proof}
Par Sites, lemme \ref{sites-lemma-point-pushforward-sheaf}, on a (1).
Par Sites, lemme \ref{sites-lemma-adjoint-point-push-stalk}, on voit que
$\textit{PSh}(\mathcal{C}) \to \textit{Ensembles}$,
$\mathcal{F} \mapsto \mathcal{F}_p$ est un adjoint à gauche, et, par
Sites, lemme \ref{sites-lemma-point-pushforward-sheaf}, il en va de même pour
$\Sh(\mathcal{C}) \to \textit{Ensembles}$,
$\mathcal{F} \mapsto \mathcal{F}_p$.
Le foncteur des germes commute donc aux colimites arbitraires (voir
Catégories, lemme \ref{categories-lemma-adjoint-exact}).
Il résulte de la définition d'un point d'un site, voir
Sites, définition \ref{sites-definition-point}, que
$\Sh(\mathcal{C}) \to \textit{Ensembles}$,
$\mathcal{F} \mapsto \mathcal{F}_p$
est exact. Comme la faisceautisation est exacte
(Sites, lemme \ref{sites-lemma-sheafification-exact}), il s'ensuit que
$\textit{PSh}(\mathcal{C}) \to \textit{Ensembles}$,
$\mathcal{F} \mapsto \mathcal{F}_p$
est exact.
\end{proof}

\noindent
En particulier, puisque le foncteur des germes
$\mathcal{F} \mapsto \mathcal{F}_p$ sur les préfaisceaux commute à toutes les
limites et colimites finies, on peut appliquer le raisonnement de la
démonstration de Sites, proposition
\ref{sites-proposition-functoriality-algebraic-structures-topoi}.
On obtient ainsi que, si $\mathcal{F}$ est un (pré)faisceau de structures
algébriques figurant dans Sites, proposition
\ref{sites-proposition-functoriality-algebraic-structures-topoi}, alors le germe
$\mathcal{F}_p$ porte naturellement une structure algébrique du même type.
Détaillons le cas où $\mathcal{F}$ est un préfaisceau abélien. L'addition
$+ : \mathcal{F} \times \mathcal{F} \to \mathcal{F}$ induit une application
$$
+ :
\mathcal{F}_p \times \mathcal{F}_p
=
(\mathcal{F} \times \mathcal{F})_p
\longrightarrow
\mathcal{F}_p
$$
où l'égalité utilise le fait que le foncteur des germes sur les préfaisceaux
d'ensembles commute aux limites finies. Ceci définit une structure de groupe
sur le germe $\mathcal{F}_p$. On obtient ainsi le foncteur des germes
$$
\textit{PAb}(\mathcal{C}) \longrightarrow \textit{Ab}, \quad
\mathcal{F} \longmapsto \mathcal{F}_p
$$
Par construction, l'ensemble sous-jacent à $\mathcal{F}_p$ est le germe du
préfaisceau d'ensembles sous-jacent. Cela définit aussi le foncteur des germes pour
les faisceaux de groupes abéliens, par précomposition avec l'inclusion
$\textit{Ab}(\mathcal{C}) \subset \textit{PAb}(\mathcal{C})$.

\begin{lemma}
\label{sites-modules-lemma-stalk-exact-abelian}
Soit $\mathcal{C}$ un site.
Soit $p$ un point de $\mathcal{C}$.
\begin{enumerate}
\item Le foncteur
$\textit{Ab}(\mathcal{C}) \to \textit{Ab}$,
$\mathcal{F} \mapsto \mathcal{F}_p$ est exact.
\item Le foncteur des germes
$\textit{PAb}(\mathcal{C}) \to \textit{Ab}$,
$\mathcal{F}  \mapsto  \mathcal{F}_p$
est exact.
\item Pour $\mathcal{F} \in \Ob(\textit{PAb}(\mathcal{C}))$, on a
$\mathcal{F}_p = \mathcal{F}^\#_p$.
\end{enumerate}
\end{lemma}

\begin{proof}
Cela résulte formellement du lemme \ref{sites-modules-lemma-stalk-exact} et de la
construction précédente du foncteur des germes.
\end{proof}

\noindent
Passons maintenant au cas des faisceaux de modules.
Soit $(\mathcal{C}, \mathcal{O})$ un site annelé.
(Pour cette discussion, il suffit que $\mathcal{O}$ soit un préfaisceau
d'anneaux.) Soit $\mathcal{F}$ un préfaisceau de $\mathcal{O}$-modules.
Soit $p$ un point de $\mathcal{C}$. On obtient alors une application
$$
\cdot :
\mathcal{O}_p \times \mathcal{O}_p
=
(\mathcal{O} \times \mathcal{O})_p
\longrightarrow
\mathcal{O}_p
$$
qui est le germe de l'application de multiplication, et
$$
\cdot :
\mathcal{O}_p \times \mathcal{F}_p
=
(\mathcal{O} \times \mathcal{F})_p
\longrightarrow
\mathcal{F}_p
$$
qui est le germe de l'application de multiplication. Nous omettons la
vérification que ceci définit une structure d'anneau sur $\mathcal{O}_p$ et une
structure de $\mathcal{O}_p$-module sur $\mathcal{F}_p$. On obtient ainsi un foncteur
$$
\textit{PMod}(\mathcal{O}) \longrightarrow \textit{Mod}(\mathcal{O}_p), \quad
\mathcal{F} \longmapsto \mathcal{F}_p
$$
Par construction, l'ensemble sous-jacent à $\mathcal{F}_p$ est le germe du
préfaisceau d'ensembles sous-jacent. Cela définit aussi le foncteur des germes pour
les faisceaux de $\mathcal{O}$-modules, par précomposition avec l'inclusion
$\textit{Mod}(\mathcal{O}) \subset \textit{PMod}(\mathcal{O})$.

\begin{lemma}
\label{sites-modules-lemma-stalk-exact-modules}
Soit $(\mathcal{C}, \mathcal{O})$ un site annelé.
Soit $p$ un point de $\mathcal{C}$.
\begin{enumerate}
\item Le foncteur
$\textit{Mod}(\mathcal{O}) \to \textit{Mod}(\mathcal{O}_p)$,
$\mathcal{F} \mapsto \mathcal{F}_p$ est exact.
\item Le foncteur des germes
$\textit{PMod}(\mathcal{O}) \to \textit{Mod}(\mathcal{O}_p)$,
$\mathcal{F} \mapsto \mathcal{F}_p$
est exact.
\item Pour $\mathcal{F} \in \Ob(\textit{PMod}(\mathcal{O}))$, on a
$\mathcal{F}_p = \mathcal{F}^\#_p$.
\end{enumerate}
\end{lemma}

\begin{proof}
Cela résulte formellement du lemme \ref{sites-modules-lemma-stalk-exact-abelian}, de la
construction précédente du foncteur des germes et du lemme \ref{sites-modules-lemma-abelian}.
\end{proof}

\begin{lemma}
\label{sites-modules-lemma-pullback-stalk}
Soit
$(f, f^\sharp) :
(\Sh(\mathcal{C}), \mathcal{O}_\mathcal{C})
\to
(\Sh(\mathcal{D}), \mathcal{O}_\mathcal{D})$
soit un morphisme de topos annelés ou de sites annelés.
Soit $p$ un point de $\mathcal{C}$ ou de $\Sh(\mathcal{C})$, et posons
$q = f \circ p$. Alors
$$
(f^*\mathcal{F})_p =
\mathcal{F}_q \otimes_{\mathcal{O}_{\mathcal{D}, q}}
\mathcal{O}_{\mathcal{C}, p}
$$
pour tout $\mathcal{O}_\mathcal{D}$-module $\mathcal{F}$.
\end{lemma}

\begin{proof}
On a
$$
f^*\mathcal{F} =
f^{-1}\mathcal{F} \otimes_{f^{-1}\mathcal{O}_\mathcal{D}}
\mathcal{O}_\mathcal{C}
$$
par définition. Comme la prise du germe en $p$ (c'est-à-dire l'application de
$p^{-1}$) commute à $\otimes$ par le
lemme \ref{sites-modules-lemma-tensor-product-pullback}, le résultat découle de la relation
entre le germe en $p$ d'une image inverse et les germes en $q$, exposée dans
Sites, lemme \ref{sites-lemma-point-morphism-sites}, ou
Sites, lemme \ref{sites-lemma-point-morphism-topoi}.
\end{proof}






\section{Faisceaux gratte-ciel}
\label{sites-modules-section-skyscraper}

\noindent
Soit $p$ un point d'un site $\mathcal{C}$ ou d'un topos
$\Sh(\mathcal{C})$. Dans cette section, nous étudions les propriétés
d'exactitude du foncteur qui associe à un groupe abélien $A$ le faisceau
gratte-ciel $p_*A$. Rappelons d'abord que
$p_* : \textit{Ensembles} \to \Sh(\mathcal{C})$ possède de nombreuses propriétés
d'exactitude; voir Sites, lemmes \ref{sites-lemma-stalk-skyscraper} et
\ref{sites-lemma-skyscraper-functor-exact}.

\begin{lemma}
\label{sites-modules-lemma-skyscraper-exact}
Soit $\mathcal{C}$ un site. Soit $p$ un point de $\mathcal{C}$ ou de son
topos associé.
\begin{enumerate}
\item Le foncteur $p_* : \textit{Ab} \to \textit{Ab}(\mathcal{C})$,
$A \mapsto p_*A$, est exact.
\item Il existe une décomposition fonctorielle en somme directe
$$
p^{-1}p_*A = A \oplus I(A)
$$
pour $A \in \Ob(\textit{Ab})$.
\end{enumerate}
\end{lemma}

\begin{proof}
Par Sites, lemme \ref{sites-lemma-stalk-skyscraper}, il existe des
applications fonctorielles $A \to p^{-1}p_*A \to A$ dont la composée est
$\text{id}_A$. On obtient donc une décomposition fonctorielle en somme directe
comme en (2), où $I(A)$ est le noyau de l'application d'adjonction
$p^{-1}p_*A \to A$. Le foncteur $p_*$ est exact à gauche par le
lemme \ref{sites-modules-lemma-exactness-pushforward-pullback}. Le foncteur $p_*$ transforme
les surjections en surjections par Sites,
lemme \ref{sites-lemma-skyscraper-functor-exact}. D'où (1).
\end{proof}

\noindent
Pour procéder de même avec les faisceaux de modules, soit $p$ un point d'un
topos annelé $(\Sh(\mathcal{C}), \mathcal{O})$. Rappelons que $p^{-1}$ est
précisément le foncteur des germes. On peut donc considérer $p$ comme un morphisme
de topos annelés
$$
(p, \text{id}_{\mathcal{O}_p}) :
(\Sh(pt), \mathcal{O}_p)
\longrightarrow
(\Sh(\mathcal{C}), \mathcal{O}).
$$
On obtient ainsi un foncteur image inverse
$p^* : \textit{Mod}(\mathcal{O}) \to \textit{Mod}(\mathcal{O}_p)$
qui coïncide avec le foncteur des germes étudié dans le
lemme \ref{sites-modules-lemma-stalk-exact-modules}. Dans cette section, nous considérons le foncteur
$p_* : \textit{Mod}(\mathcal{O}_p) \to \textit{Mod}(\mathcal{O})$.

\begin{lemma}
\label{sites-modules-lemma-skyscraper-modules-exact}
Soit $(\Sh(\mathcal{C}), \mathcal{O})$ un topos annelé.
Soit $p$ un point du topos $\Sh(\mathcal{C})$.
\begin{enumerate}
\item Le foncteur
$p_* : \textit{Mod}(\mathcal{O}_p) \to \textit{Mod}(\mathcal{O})$,
$M \mapsto p_*M$ est exact.
\item La surjection canonique $p^{-1}p_*M \to M$ est $\mathcal{O}_p$-linéaire.
\item La décomposition fonctorielle en somme directe
$p^{-1}p_*M = M \oplus I(M)$ du lemme \ref{sites-modules-lemma-skyscraper-exact}
n'est {\bf pas} $\mathcal{O}_p$-linéaire en général.
\end{enumerate}
\end{lemma}

\begin{proof}
La partie (1) et la surjectivité dans (2) résultent immédiatement du résultat
correspondant pour les faisceaux abéliens dans le
lemme \ref{sites-modules-lemma-skyscraper-exact}. Comme
$p^{-1}\mathcal{O} = \mathcal{O}_p$, on a $p^{-1} = p^*$; par conséquent,
$p^{-1}p_*M \to M$ est la co-unité $p^*p_*M \to M$ de l'adjonction pour les
modules, et est donc linéaire.

\medskip\noindent
Démonstration de (3). Soit $G$ un groupe. Considérons le topos
$G\textit{-Ensembles} = \Sh(\mathcal{T}_G)$
et le point $p : \textit{Ensembles} \to G\textit{-Ensembles}$. Voir Sites,
section \ref{sites-section-example-sheaf-G-sets}, et
exemple \ref{sites-example-point-G-sets}. Ici, $p^{-1}$ est le foncteur
d'oubli de l'action de $G$, et $p_*$ est son adjoint à droite, qui envoie
$M$ sur $\text{Map}(G, M)$. Les applications de la décomposition en somme
directe sont
$$
M \to \text{Map}(G, M) \to M
$$
où la première envoie $m \in M$ sur l'application constante de valeur $m$, et
la seconde est l'évaluation en l'élément neutre $1$ de $G$.
Supposons ensuite que $R$ soit un anneau muni d'une action de $G$. Celle-ci
détermine un faisceau d'anneaux $\mathcal{O}$ sur $\mathcal{T}_G$. La catégorie
des $\mathcal{O}$-modules est celle des $R$-modules $M$ munis d'une action de
$G$ compatible à celle sur $R$. La structure de $R$-module sur
$\text{Map}(G, M)$ est donnée par
$$
( r f ) (\sigma) = \sigma(r) f(\sigma)
$$
pour $r \in R$ et $f \in \text{Map}(G, M)$. C'est en effet l'unique structure
$G$-invariante de $R$-module compatible avec l'évaluation en $1$. Le lecteur
observera qu'en général l'image de $M \to \text{Map}(G, M)$ n'est pas un
sous-$R$-module (prendre par exemple $M = R$ et supposer l'action de $G$ non
triviale), ce qui achève la démonstration.
\end{proof}

\begin{example}
\label{sites-modules-example-ring-with-group-action}
Soit $G$ un groupe. Considérons le site $\mathcal{T}_G$ et son point $p$;
voir Sites, exemple \ref{sites-example-point-G-sets}. Soit $R$ un anneau muni
d'une action de $G$, correspondant à un faisceau d'anneaux $\mathcal{O}$ sur
$\mathcal{T}_G$. Alors $\mathcal{O}_p = R$, où l'on oublie l'action de $G$.
Dans ce cas, $p^{-1}p_*M = \text{Map}(G, M)$,
$I(M) = \{f : G \to M \mid f(1_G) = 0\}$, et
$M \to \text{Map}(G, M)$ associe à $m \in M$ la fonction constante de valeur $m$.
\end{example}




\section{Localisation et points}
\label{sites-modules-section-localize-points}

\begin{lemma}
\label{sites-modules-lemma-stalk-j-shriek}
Soit $(\mathcal{C}, \mathcal{O})$ un site annelé.
Soit $p$ un point de $\mathcal{C}$. Soit $U$ un objet de $\mathcal{C}$.
Pour $\mathcal{G}$ dans $\textit{Mod}(\mathcal{O}_U)$, on a
$$
(j_{U!}\mathcal{G})_p =
\bigoplus\nolimits_q \mathcal{G}_q
$$
où la somme directe porte sur les points $q$ de $\mathcal{C}/U$ situés
au-dessus de $p$; voir Sites, lemme \ref{sites-lemma-points-above-point}.
\end{lemma}

\begin{proof}
Nous utilisons la description de $j_{U!}\mathcal{G}$ comme faisceau associé au préfaisceau
$V \mapsto
\bigoplus\nolimits_{\varphi \in \Mor_\mathcal{C}(V, U)}
\mathcal{G}(V/_\varphi U)$
du lemme \ref{sites-modules-lemma-extension-by-zero}. Le germe de $j_{U!}\mathcal{G}$ en
$p$ est égale à celle de ce préfaisceau; voir le
lemme \ref{sites-modules-lemma-stalk-exact-modules}. Soit
$u : \mathcal{C} \to \textit{Ensembles}$ le foncteur correspondant à $p$
(voir Sites, section \ref{sites-section-points}). On voit alors que
$$
(j_{U!}\mathcal{G})_p = \colim_{(V, y)}
\bigoplus\nolimits_{\varphi : V \to U} \mathcal{G}(V/_\varphi U)
$$
où la colimite est prise dans la catégorie des groupes abéliens.
À un quadruplet $(V, y, \varphi, s)$ intervenant dans cette colimite, on peut
associer $x = u(\varphi)(y) \in u(U)$. On obtient ainsi
$$
(j_{U!}\mathcal{G})_p =
\bigoplus\nolimits_{x \in u(U)}
\colim_{(\varphi : V \to U, y), \ u(\varphi)(y) = x} \mathcal{G}(V/_\varphi U).
$$
C'est l'expression du lemme, d'après la description des points $q$ situés
au-dessus de $x$ donnée dans Sites,
lemme \ref{sites-lemma-points-above-point}.
\end{proof}

\begin{remark}
\label{sites-modules-remark-not-pushforward}
Attention : le résultat du lemme \ref{sites-modules-lemma-stalk-j-shriek} n'a pas d'analogue
pour $j_{U, *}$.
\end{remark}












\section{Images inverses de modules plats}
\label{sites-modules-section-pullback-flat}

\noindent
L'image inverse d'un module plat par un morphisme de topos annelés est plate.
La démonstration demande quelque soin.

\begin{lemma}
\label{sites-modules-lemma-pullback-flat}
\begin{reference}
\cite[Expos\'e V, Corollary 1.7.1]{SGA4}
\end{reference}
Soit
$(f, f^\sharp) :
(\Sh(\mathcal{C}), \mathcal{O}_\mathcal{C})
\to
(\Sh(\mathcal{D}), \mathcal{O}_\mathcal{D})$
soit un morphisme de topos annelés ou de sites annelés.
Alors $f^*\mathcal{F}$ est un $\mathcal{O}_\mathcal{C}$-module plat lorsque
$\mathcal{F}$ est un $\mathcal{O}_\mathcal{D}$-module plat.
\end{lemma}

\begin{proof}
Choisissons un diagramme comme dans le
lemme \ref{sites-modules-lemma-morphism-ringed-topoi-comes-from-morphism-ringed-sites}.
Rappelons que la propriété d'être un module plat est intrinsèque
(voir la section \ref{sites-modules-section-intrinsic} et la
définition \ref{sites-modules-definition-flat}). Il suffit donc de démontrer le lemme pour
le morphisme $(h, h^\sharp) :
(\Sh(\mathcal{C}'), \mathcal{O}_{\mathcal{C}'})
\to
(\Sh(\mathcal{D}'), \mathcal{O}_{\mathcal{D}'})$.
Autrement dit, on peut supposer que nos sites $\mathcal{C}$ et $\mathcal{D}$
admettent toutes les limites finies et que $f$ est un morphisme de sites induit
par un foncteur continu $u : \mathcal{D} \to \mathcal{C}$ qui commute aux
limites finies.

\medskip\noindent
Rappelons que $f^*\mathcal{F} =
\mathcal{O}_\mathcal{C} \otimes_{f^{-1}\mathcal{O}_\mathcal{D}}
f^{-1}\mathcal{F}$ (définition \ref{sites-modules-definition-pushforward}).
Par le lemme \ref{sites-modules-lemma-flat-change-of-rings}, il suffit de montrer que
$f^{-1}\mathcal{F}$ est un $f^{-1}\mathcal{O}_\mathcal{D}$-module plat.
Avec le paragraphe précédent, cela nous ramène à la situation du paragraphe suivant.

\medskip\noindent
Supposons que $\mathcal{C}$ et $\mathcal{D}$ soient des sites admettant toutes
les limites finies, et que $u : \mathcal{D} \to \mathcal{C}$ soit un foncteur
continu commutant aux limites finies. Soit $\mathcal{O}$ un faisceau d'anneaux
sur $\mathcal{D}$ et soit $\mathcal{F}$ un $\mathcal{O}$-module plat. Alors $u$
définit un morphisme de sites $f : \mathcal{C} \to \mathcal{D}$
(Sites, proposition \ref{sites-proposition-get-morphism}). Il faut montrer que
$f^{-1}\mathcal{F}$ est un $f^{-1}\mathcal{O}$-module plat. Soit $U$ un objet
de $\mathcal{C}$ et soit
$$
f^{-1}\mathcal{O}|_U \xrightarrow{(f_1, \ldots, f_n)}
f^{-1}\mathcal{O}|_U^{\oplus n} \xrightarrow{(s_1, \ldots, s_n)}
f^{-1}\mathcal{F}|_U
$$
un complexe de $f^{-1}\mathcal{O}|_U$-modules. Notre but est de construire une
factorisation de $(s_1, \ldots, s_n)$ sur les membres d'un recouvrement de $U$,
comme dans le lemme \ref{sites-modules-lemma-flat-eq}, partie (2).
Considérons les éléments $s_a \in f^{-1}\mathcal{F}(U)$ et
$f_a \in f^{-1}\mathcal{O}(U)$. Comme $f^{-1}\mathcal{F}$ et, de même,
$f^{-1}\mathcal{O}$ s'obtiennent par faisceautisation de $u_p\mathcal{F}$, on
peut, après avoir remplacé $U$ par les membres d'un recouvrement, supposer que
$s_a$ est l'image d'un élément
$s'_a \in u_p\mathcal{F}(U)$ et que $f_a$ est l'image d'un élément
$f'_a \in u_p\mathcal{O}(U)$. Après un nouveau remplacement de $U$ par les
membres d'un recouvrement, on peut supposer que $\sum f'_as'_a$ est nul dans
$u_p\mathcal{F}(U)$. Rappelons que la catégorie
$(\mathcal{I}_U^u)^{opp}$ est filtrante
(Sites, lemme \ref{sites-lemma-directed}) et que
$u_p\mathcal{F}(U) = \colim_{(\mathcal{I}_U^u)^{opp}} \mathcal{F}(V)$ et
$u_p\mathcal{O}(U) = \colim_{(\mathcal{I}_U^u)^{opp}} \mathcal{O}(V)$.
On peut donc supposer qu'il existe une paire
$(V, \phi) \in \Ob(\mathcal{I}_U^u)$, où $V$ est un objet de $\mathcal{D}$ et
$\phi$ est un morphisme $\phi : U \to u(V)$ de $\mathcal{D}$, ainsi que des éléments
$s''_a \in \mathcal{F}(V)$ et $f''_a \in \mathcal{O}(V)$ dont les images dans
$u_p\mathcal{F}(U)$ et $u_p\mathcal{O}(U)$ sont $s'_a$ et $f'_a$, et tels que
$\sum f''_a s''_a = 0$ dans $\mathcal{F}(V)$. On obtient alors un complexe
$$
\mathcal{O}|_V \xrightarrow{(f''_1, \ldots, f''_n)}
\mathcal{O}|_V^{\oplus n} \xrightarrow{(s''_1, \ldots, s''_n)}
\mathcal{F}|_V
$$
et l'autre implication du lemme \ref{sites-modules-lemma-flat-eq} montre qu'il existe un
recouvrement $\{V_i \to V\}$ dans $\mathcal{D}$ et, pour chaque $i$, une factorisation
$$
\mathcal{O}|_{V_i}^{\oplus n}
\xrightarrow{B''_i}
\mathcal{O}|_{V_i}^{\oplus l_i} \xrightarrow{(t''_{i1}, \ldots, t''_{il_i})}
\mathcal{F}|_{V_i}
$$
de $(s''_1, \ldots, s''_n)|_{V_i}$ telle que
$B''_i \circ (f''_1, \ldots, f''_n)|_{V_i} = 0$.
Posons $U_i = U \times_{\phi, u(V)} u(V_i)$, notons
$B_i \in \text{Mat}(l_i \times n, f^{-1}\mathcal{O}(U_i))$ l'image de $B''_i$
et $t_{ij} \in f^{-1}\mathcal{F}(U_i)$ l'image de $t''_{ij}$. On obtient une factorisation
$$
f^{-1}\mathcal{O}|_{U_i}^{\oplus n}
\xrightarrow{B_i}
f^{-1}\mathcal{O}|_{U_i}^{\oplus l_i}
\xrightarrow{(t_{i1}, \ldots, t_{il_i})}
\mathcal{F}|_{U_i}
$$
de $(s_1, \ldots, s_n)|_{U_i}$ telle que
$B_i \circ (f_1, \ldots, f_n)|_{U_i} = 0$.
Cela achève la démonstration.
\end{proof}

\begin{lemma}
\label{sites-modules-lemma-stalk-flat}
Soit $(\mathcal{C}, \mathcal{O})$ un site annelé et soit $p$ un point de
$\mathcal{C}$. Si $\mathcal{F}$ est un $\mathcal{O}$-module plat, alors
$\mathcal{F}_p$ est un $\mathcal{O}_p$-module plat.
\end{lemma}

\begin{proof}
Dans la section \ref{sites-modules-section-skyscraper}, nous avons vu que l'on peut
considérer $p$ comme un morphisme de topos annelés
$$
(p, \text{id}_{\mathcal{O}_p}) :
(\Sh(pt), \mathcal{O}_p)
\longrightarrow
(\Sh(\mathcal{C}), \mathcal{O}).
$$
tel que le foncteur image inverse
$p^* : \textit{Mod}(\mathcal{O}) \to \textit{Mod}(\mathcal{O}_p)$
coïncide avec le foncteur des germes. Le lemme résulte donc du
lemme \ref{sites-modules-lemma-pullback-flat}.
\end{proof}

\begin{lemma}
\label{sites-modules-lemma-check-flat-stalks}
Soit $(\mathcal{C}, \mathcal{O})$ un site annelé. Soit $\mathcal{F}$ un
faisceau de $\mathcal{O}$-modules. Soit $\{p_i\}_{i \in I}$ une famille
conservative de points de $\mathcal{C}$. Alors $\mathcal{F}$ est plat si et
seulement si $\mathcal{F}_{p_i}$ est un $\mathcal{O}_{p_i}$-module plat pour
tout $i \in I$.
\end{lemma}

\begin{proof}
Le lemme \ref{sites-modules-lemma-stalk-flat} donne une implication. Pour la réciproque,
utilisons l'égalité
$(\mathcal{F} \otimes_\mathcal{O} \mathcal{G})_p =
\mathcal{F}_p \otimes_{\mathcal{O}_p} \mathcal{G}_p$
fournie par le lemme \ref{sites-modules-lemma-tensor-product-pullback} (puisque le germe en
$p$ est donnée par $p^{-1}$), ainsi que le
lemme \ref{sites-modules-lemma-check-exactness-stalks}.
\end{proof}

\begin{lemma}
\label{sites-modules-lemma-pullback-ses}
Soit
$f : (\Sh(\mathcal{C}'), \mathcal{O}') \to (\Sh(\mathcal{C}), \mathcal{O})$
soit un morphisme de topos annelés. Soit
$0 \to \mathcal{F} \to \mathcal{G} \to \mathcal{H} \to 0$
une suite exacte courte de $\mathcal{O}$-modules, où $\mathcal{H}$ est un
$\mathcal{O}$-module plat. Alors la suite
$0 \to f^*\mathcal{F} \to f^*\mathcal{G} \to f^*\mathcal{H} \to 0$
est elle aussi exacte.
\end{lemma}

\begin{proof}
Comme $f^{-1}$ est exact, on a la suite exacte courte
$0 \to f^{-1}\mathcal{F} \to f^{-1}\mathcal{G} \to f^{-1}\mathcal{H} \to 0$
de $f^{-1}\mathcal{O}$-modules. Par le lemme \ref{sites-modules-lemma-pullback-flat}, le
$f^{-1}\mathcal{O}$-module $f^{-1}\mathcal{H}$ est plat. Le
lemme \ref{sites-modules-lemma-flat-tor-zero} montre donc que la suite reste exacte après
tensorisation sur $f^{-1}\mathcal{O}$ par $\mathcal{O}'$. Comme
$f^*\mathcal{F} = f^{-1}\mathcal{F} \otimes_{f^{-1}\mathcal{O}} \mathcal{O}'$
et de même pour $\mathcal{G}$ et $\mathcal{H}$, on conclut.
\end{proof}





\section{Topos localement annelés}
\label{sites-modules-section-locally-ringed}


\noindent
Une référence pour cette section est
\cite[Expos\'e IV, Exercice 13.9]{SGA4}.

\begin{lemma}
\label{sites-modules-lemma-locally-ringed}
Soit $(\mathcal{C}, \mathcal{O})$ un site annelé. Les conditions suivantes
sont équivalentes :
\begin{enumerate}
\item Pour tout objet $U$ de $\mathcal{C}$ et tout $f \in \mathcal{O}(U)$,
il existe un recouvrement $\{U_j \to U\}$ tel que, pour chaque $j$, soit
$f|_{U_j}$, soit $(1 - f)|_{U_j}$ soit inversible.
\item Pour $U \in \Ob(\mathcal{C})$, $n \geq 1$, et
$f_1, \ldots, f_n \in \mathcal{O}(U)$ engendrant l'idéal unité de
$\mathcal{O}(U)$, il existe un recouvrement $\{U_j \to U\}$ tel que, pour
chaque $j$, il existe $i$ pour lequel $f_i|_{U_j}$ est inversible.
\item L'application de faisceaux d'ensembles
$$
(\mathcal{O} \times \mathcal{O})
\amalg
(\mathcal{O} \times \mathcal{O})
\longrightarrow
\mathcal{O} \times \mathcal{O}
$$
qui envoie $(f, a)$ dans la première composante sur $(f, af)$ et $(f, b)$ dans
la seconde sur $(f, b(1 - f))$ est surjective.
\end{enumerate}
\end{lemma}

\begin{proof}
Il est clair que (2) implique (1). Pour montrer que (1) implique (2),
raisonnons par récurrence sur $n$. Le premier cas est $n = 2$ (le cas $n = 1$
étant trivial). On a alors $a_1f_1 + a_2f_2 = 1$ pour certains
$a_1, a_2 \in \mathcal{O}(U)$. Par hypothèse, il existe un recouvrement
$\{U_j \to U\}$ tel que, pour chaque $j$, soit $a_1f_1|_{U_j}$, soit
$a_2f_2|_{U_j}$ soit inversible. Ainsi $f_1|_{U_j}$ ou $f_2|_{U_j}$ est
inversible, comme voulu. Pour $n > 2$, on a
$a_1f_1 + \ldots + a_nf_n = 1$ pour certains
$a_1, \ldots, a_n \in \mathcal{O}(U)$.
Le cas $n = 2$ donne un recouvrement $\{U_j \to U\}_{j \in J}$ tel que, pour
chaque $j$, soit $f_n|_{U_j}$, soit
$a_1f_1 + \ldots + a_{n - 1}f_{n - 1}|_{U_j}$ soit inversible.
Disons que le premier cas se produit pour $j \in J_n$ et posons
$J' = J \setminus J_n$. Par l'hypothèse de récurrence, pour chaque
$j \in J'$, il existe un recouvrement
$\{U_{jk} \to U_j\}_{k \in K_j}$ tel que, pour chaque $k \in K_j$, il existe
un $i \in \{1, \ldots, n - 1\}$ tel que $f_i|_{U_{jk}}$ soit inversible. Par les
axiomes d'un site, la famille de morphismes
$\{U_j \to U\}_{j \in J_n} \cup \{U_{jk} \to U\}_{j \in J', k \in K_j}$
est un recouvrement possédant la propriété voulue.

\medskip\noindent
Supposons (1). Pour voir que l'application de (3) est surjective, soit
$(f, c)$ une section de $\mathcal{O} \times \mathcal{O}$ sur $U$. Par
hypothèse, il existe un recouvrement $\{U_j \to U\}$ tel que, pour chaque $j$,
$f$ ou $1 - f$ se restreigne en une section inversible. Dans le premier cas, on
prend $a = c|_{U_j} (f|_{U_j})^{-1}$; dans le second,
$b = c|_{U_j} (1 - f|_{U_j})^{-1}$. Ainsi $(f, c)$ appartient à l'image sur
chaque membre du recouvrement. Réciproquement, supposons (3). Pour tout $U$ et
$f \in \mathcal{O}(U)$, il existe un recouvrement $\{U_j \to U\}$ de $U$ tel que la
section $(f, 1)|_{U_j}$ appartienne à l'image de l'application de (3) sur
$U_j$. Cela signifie précisément que $f$ ou $1 - f$ se restreint en une section
inversible sur $U_j$, d'où (1).
\end{proof}

\begin{lemma}
\label{sites-modules-lemma-locally-ringed-stalk}
Soit $(\mathcal{C}, \mathcal{O})$ un site annelé. Considérons les conditions suivantes :
\begin{enumerate}
\item Pour tout objet $U$ de $\mathcal{C}$ et tout $f \in \mathcal{O}(U)$,
il existe un recouvrement $\{U_j \to U\}$ tel que, pour chaque $j$, soit
$f|_{U_j}$, soit $(1 - f)|_{U_j}$ soit inversible.
\item Pour tout point $p$ de $\mathcal{C}$, le germe $\mathcal{O}_p$ est
soit l'anneau nul, soit un anneau local.
\end{enumerate}
On a toujours (1) $\Rightarrow$ (2). Si $\mathcal{C}$ a assez de points,
alors (1) et (2) sont équivalentes.
\end{lemma}

\begin{proof}
Supposons (1). Soit $p$ un point de $\mathcal{C}$ donné par un foncteur
$u : \mathcal{C} \to \textit{Ensembles}$, et soit $f_p \in \mathcal{O}_p$.
Puisque $\mathcal{O}_p$ est calculé par Sites, équation
(\ref{sites-equation-stalk}), on peut représenter $f_p$ par un triplet
$(U, x, f)$, où $x \in U(U)$ et $f \in \mathcal{O}(U)$. Par hypothèse, il existe
un recouvrement $\{U_i \to U\}$ tel que, pour chaque $i$, $f$ ou $1 - f$ soit
inversible sur $U_i$. Comme $u$ définit un point du site, il existe, pour un
certain $i$, un $x_i \in u(U_i)$ qui s'envoie sur $x \in u(U)$. La discussion
autour de Sites, équation (\ref{sites-equation-stalk}), montre que $(U, x, f)$ et
$(U_i, x_i, f|_{U_i})$ définissent le même élément de $\mathcal{O}_p$. Ainsi
$f_p$ ou $1 - f_p$ est inversible. Par conséquent, dans l'anneau
$\mathcal{O}_p$, pour tout élément $a$, $a$ ou $1 - a$ est inversible. Cela
signifie que $\mathcal{O}_p$ est nul ou local; voir Algèbre,
lemme \ref{algebra-lemma-characterize-local-ring}.

\medskip\noindent
Supposons (2) et que $\mathcal{C}$ ait assez de points. Considérons
l'application de faisceaux d'ensembles
$$
\mathcal{O} \times \mathcal{O} \amalg \mathcal{O} \times \mathcal{O}
\longrightarrow
\mathcal{O} \times \mathcal{O}
$$
du lemme \ref{sites-modules-lemma-locally-ringed}, partie (3). Pour tout anneau local $R$,
l'application correspondante
$(R \times R) \amalg (R \times R) \to R \times R$
est surjective; voir par exemple Algèbre,
lemme \ref{algebra-lemma-characterize-local-ring}. Comme chaque
$\mathcal{O}_p$ est un anneau local ou nul, l'application est surjective sur
les germes. Puisque $\mathcal{C}$ a assez de points, elle est donc surjective,
ce qui conclut.
\end{proof}

\noindent
Dans Modules, section \ref{modules-section-pathology}, nous avons signalé
qu'un espace annelé $(X, \mathcal{O}_X)$ peut avoir un sous-espace ouvert sur
lequel le faisceau structural est nul. Pour l'exclure, on peut imposer que les
sections $1$ et $0$ aient des valeurs distinctes dans tout germe de $X$.
Pour les topos et sites annelés, la condition est que
\begin{equation}
\label{sites-modules-equation-one-is-never-zero}
\emptyset^\# \longrightarrow
\text{Equalizer}(0, 1 : * \longrightarrow \mathcal{O})
\end{equation}
soit un isomorphisme de faisceaux. Ici, $*$ est le faisceau singleton et
$\emptyset^\#$ le « faisceau vide », c'est-à-dire respectivement l'objet final
et l'objet initial de la catégorie des faisceaux; voir Sites,
exemple \ref{sites-example-singleton-sheaf}, respectivement
section \ref{sites-section-almost-cocontinuous}. Autrement dit, si
$U \in \Ob(\mathcal{C})$ n'est pas vide au sens faisceautique, alors
$1, 0 \in \mathcal{O}(U)$ ne sont pas égaux. Énonçons le lemme obligé.

\begin{lemma}
\label{sites-modules-lemma-ringed-stalk-not-zero}
Soit $(\mathcal{C}, \mathcal{O})$ un site annelé. Considérons les assertions :
\begin{enumerate}
\item (\ref{sites-modules-equation-one-is-never-zero}) est un isomorphisme, et
\item pour tout point $p$ de $\mathcal{C}$, le germe $\mathcal{O}_p$ n'est
pas l'anneau nul.
\end{enumerate}
On a toujours (1) $\Rightarrow$ (2), et, si $\mathcal{C}$ a assez de points,
alors (1) $\Leftrightarrow$ (2).
\end{lemma}

\begin{proof}
La démonstration est omise.
\end{proof}

\noindent
Les lemmes \ref{sites-modules-lemma-locally-ringed},
\ref{sites-modules-lemma-locally-ringed-stalk} et
\ref{sites-modules-lemma-ringed-stalk-not-zero} motivent la définition suivante.

\begin{definition}
\label{sites-modules-definition-locally-ringed}
Un site annelé $(\mathcal{C}, \mathcal{O})$ est dit
{\it localement annelé} si (\ref{sites-modules-equation-one-is-never-zero}) est un
isomorphisme et si les propriétés équivalentes du
lemme \ref{sites-modules-lemma-locally-ringed} sont satisfaites.
\end{definition}

\noindent
Dans \cite[Expos\'e IV, Exercice 13.9]{SGA4}, la condition que
(\ref{sites-modules-equation-one-is-never-zero}) soit un isomorphisme manque, ce qui conduit
à une notion légèrement différente de site et de topos localement annelés.
Comme nous sommes guidés par la notion d'espace localement annelé, nous avons
choisi d'ajouter cette condition (voir l'explication ci-dessus).

\begin{lemma}
\label{sites-modules-lemma-locally-ringed-intrinsic}
Être un site localement annelé est une propriété intrinsèque. Plus précisément :
\begin{enumerate}
\item si $f : \Sh(\mathcal{C}') \to \Sh(\mathcal{C})$ est un morphisme de
topos et si $(\mathcal{C}, \mathcal{O})$ est un site localement annelé, alors
$(\mathcal{C}', f^{-1}\mathcal{O})$ est un site localement annelé; et
\item si
$(f, f^\sharp) : (\Sh(\mathcal{C}'), \mathcal{O}')
\to (\Sh(\mathcal{C}), \mathcal{O})$
est une équivalence de topos annelés, alors $(\mathcal{C}, \mathcal{O})$ est
localement annelé si et seulement si $(\mathcal{C}', \mathcal{O}')$ l'est.
\end{enumerate}
\end{lemma}

\begin{proof}
Il est clair que (2) résulte de (1). Pour démontrer (1), remarquons que, comme
$f^{-1}$ est exact, on a $f^{-1}* = *$ et
$f^{-1}\emptyset^\# = \emptyset^\#$; de plus, $f^{-1}$ commute aux produits
et aux égalisateurs, et transforme isomorphismes et surjections en
isomorphismes et surjections. Ainsi, $f^{-1}$ transforme l'isomorphisme
(\ref{sites-modules-equation-one-is-never-zero}) en son analogue pour
$f^{-1}\mathcal{O}$, et la surjection du
lemme \ref{sites-modules-lemma-locally-ringed}, partie (3), en la surjection correspondante
pour $f^{-1}\mathcal{O}$.
\end{proof}

\noindent
En fait, le lemme \ref{sites-modules-lemma-locally-ringed-intrinsic}, partie (2), est
l'analogue de Schémas, lemme
\ref{schemes-lemma-isomorphism-locally-ringed}. Il garantit que la définition
suivante a un sens.

\begin{definition}
\label{sites-modules-definition-locally-ringed-topos}
Un topos annelé $(\Sh(\mathcal{C}), \mathcal{O})$ est dit
{\it localement annelé} si le site annelé sous-jacent
$(\mathcal{C}, \mathcal{O})$ est localement annelé.
\end{definition}

\noindent
Voici un exemple de conséquence de la propriété d'être localement annelé.

\begin{lemma}
\label{sites-modules-lemma-invertible-is-locally-free-rank-1}
Soit $(\Sh(\mathcal{C}), \mathcal{O})$ un topos annelé. Tout
$\mathcal{O}$-module localement libre de rang $1$ est inversible.
Si $(\mathcal{C}, \mathcal{O})$ est localement annelé, la réciproque est
également vraie (ce qui n'est pas le cas en général).
\end{lemma}

\begin{proof}
Supposons $\mathcal{L}$ localement libre de rang $1$ et considérons
l'application d'évaluation
$$
\mathcal{L} \otimes_\mathcal{O}
\SheafHom_\mathcal{O}(\mathcal{L}, \mathcal{O})
\longrightarrow \mathcal{O}
$$
Pour tout objet $U$ de $\mathcal{C}$, en se restreignant aux membres d'un
recouvrement qui trivialise $\mathcal{L}$, on voit que cette application est
un isomorphisme (détails omis). Le lemme \ref{sites-modules-lemma-invertible} montre donc que
$\mathcal{L}$ est inversible.

\medskip\noindent
Supposons $(\Sh(\mathcal{C}), \mathcal{O})$ localement annelé et soit $U$ un
objet de $\mathcal{C}$. Dans la démonstration du lemme \ref{sites-modules-lemma-invertible},
nous avons vu qu'il existe un recouvrement $\{U_i \to U\}$ tel que
$\mathcal{L}|_{\mathcal{C}/U_i}$ soit facteur direct d'un
$\mathcal{O}_{U_i}$-module libre fini. Après avoir remplacé $U$ par $U_i$,
soit
$p : \mathcal{O}_U^{\oplus r} \to \mathcal{O}_U^{\oplus r}$
un projecteur dont l'image est isomorphe à $\mathcal{L}|_{\mathcal{C}/U}$.
Alors $p$ correspond à une matrice
$$
P = (p_{ij}) \in \text{Mat}(r \times r, \mathcal{O}(U))
$$
qui est un projecteur : $P^2 = P$. Posons $A = \mathcal{O}(U)$, de sorte que
$P \in \text{Mat}(r \times r, A)$. Par Algèbre,
lemme \ref{algebra-lemma-finite-projective}, l'image de $P$ est un module
$M$ fini localement libre sur $A$. Il existe donc
$f_1, \ldots, f_t \in A$ engendrant l'idéal unité tels que $M_{f_i}$ soit libre
fini. Par le lemme \ref{sites-modules-lemma-locally-ringed}, après avoir remplacé $U$ par les
membres d'un recouvrement ouvert, on peut supposer $M$ libre. Cela signifie
que $\mathcal{L}|_U$ est libre (détails omis). Puisque $\mathcal{L}$ est
inversible, cela n'est possible que si le rang de $\mathcal{L}|_U$ vaut $1$,
ce qui achève la démonstration.
\end{proof}

\noindent
Nous voulons maintenant préciser ce que signifie être un morphisme d'espaces
localement annelés. Nous utiliserons le lemme suivant.

\begin{lemma}
\label{sites-modules-lemma-locally-ringed-morphism}
Soit
$(f, f^\sharp) : (\Sh(\mathcal{C}), \mathcal{O}_\mathcal{C})
\to (\Sh(\mathcal{D}), \mathcal{O}_\mathcal{D})$
soit un morphisme de topos annelés. Considérons les conditions suivantes :
\begin{enumerate}
\item Le diagramme de faisceaux
$$
\xymatrix{
f^{-1}(\mathcal{O}^*_\mathcal{D}) \ar[r]_-{f^\sharp} \ar[d] &
\mathcal{O}^*_\mathcal{C} \ar[d] \\
f^{-1}(\mathcal{O}_\mathcal{D}) \ar[r]^-{f^\sharp} &
\mathcal{O}_\mathcal{C}
}
$$
est cartésien.
\item Pour tout point $p$ de $\mathcal{C}$, en posant $q = f \circ p$, le diagramme
$$
\xymatrix{
\mathcal{O}^*_{\mathcal{D}, q} \ar[r] \ar[d] &
\mathcal{O}^*_{\mathcal{C}, p} \ar[d] \\
\mathcal{O}_{\mathcal{D}, q} \ar[r] &
\mathcal{O}_{\mathcal{C}, p}
}
$$
d'ensembles est cartésien.
\end{enumerate}
On a toujours (1) $\Rightarrow$ (2). Si $\mathcal{C}$ a assez de points,
alors (1) et (2) sont équivalentes. Si
$(\Sh(\mathcal{C}), \mathcal{O}_\mathcal{C})$
et
$(\Sh(\mathcal{D}), \mathcal{O}_\mathcal{D})$
sont des topos localement annelés, alors (2) équivaut à :
\begin{enumerate}
\item[(3)] Pour tout point $p$ de $\mathcal{C}$, en posant $q = f \circ p$,
l'homomorphisme d'anneaux
$\mathcal{O}_{\mathcal{D}, q} \to \mathcal{O}_{\mathcal{C}, p}$ est local.
\end{enumerate}
En fait, la propriété (2), ou la propriété (3) pour une famille conservative
de points, implique (1).
\end{lemma}

\begin{proof}
Le lemme se démontre de lui-même : il suffit de développer les définitions.
\end{proof}

\begin{definition}
\label{sites-modules-definition-morphism-locally-ringed-topoi}
Soit $(f, f^\sharp) : (\Sh(\mathcal{C}), \mathcal{O}_\mathcal{C})
\to (\Sh(\mathcal{D}), \mathcal{O}_\mathcal{D})$
soit un morphisme de topos annelés. Supposons que
$(\Sh(\mathcal{C}), \mathcal{O}_\mathcal{C})$
et
$(\Sh(\mathcal{D}), \mathcal{O}_\mathcal{D})$
soient des topos localement annelés. On dit que $(f, f^\sharp)$ est un
{\it morphisme de topos localement annelés} si et seulement si le diagramme de faisceaux
$$
\xymatrix{
f^{-1}(\mathcal{O}^*_\mathcal{D}) \ar[r]_-{f^\sharp} \ar[d] &
\mathcal{O}^*_\mathcal{C} \ar[d] \\
f^{-1}(\mathcal{O}_\mathcal{D}) \ar[r]^-{f^\sharp} &
\mathcal{O}_\mathcal{C}
}
$$
(voir le lemme \ref{sites-modules-lemma-locally-ringed-morphism}) est cartésien. Si
$(f, f^\sharp)$ est un morphisme de sites annelés, on dit que c'est un
{\it morphisme de sites localement annelés} si le morphisme de topos annelés
associé est un morphisme de topos localement annelés.
\end{definition}

\noindent
Il est clair qu'un isomorphisme de topos annelés entre topos localement
annelés est automatiquement un isomorphisme de topos localement annelés.

\begin{lemma}
\label{sites-modules-lemma-composition-morphisms-locally-ringed-topoi}
Soient
$(f, f^\sharp) :
(\Sh(\mathcal{C}_1), \mathcal{O}_1)
\to (\Sh(\mathcal{C}_2), \mathcal{O}_2)$ et
$(g, g^\sharp) :
(\Sh(\mathcal{C}_2), \mathcal{O}_2) \to
(\Sh(\mathcal{C}_3), \mathcal{O}_3)$
soient des morphismes de topos localement annelés. Alors leur composée
$(g, g^\sharp) \circ (f, f^\sharp)$ (voir la
définition \ref{sites-modules-definition-ringed-topos}) est encore un morphisme de topos
localement annelés.
\end{lemma}

\begin{proof}
La démonstration est omise.
\end{proof}

\begin{lemma}
\label{sites-modules-lemma-locally-ringed-intrinsic-morphism}
Soit $f : \Sh(\mathcal{C}') \to \Sh(\mathcal{C})$ un morphisme de topos.
Si $\mathcal{O}$ est un faisceau d'anneaux sur $\mathcal{C}$, alors
$$
f^{-1}(\mathcal{O}^*) = (f^{-1}\mathcal{O})^*.
$$
En particulier, si $\mathcal{O}$ fait de $\mathcal{C}$ un site localement
annelé, alors, en posant $f^\sharp = \text{id}$, le morphisme de topos annelés
$$
(f, f^\sharp) :
(\Sh(\mathcal{C}'), f^{-1}\mathcal{O})
\to
(\Sh(\mathcal{C}), \mathcal{O})
$$
est un morphisme de topos localement annelés.
\end{lemma}

\begin{proof}
Remarquons que le diagramme
$$
\xymatrix{
\mathcal{O}^* \ar[rr] \ar[d]_{u \mapsto (u, u^{-1})} & &
{*} \ar[d]^{1} \\
\mathcal{O} \times \mathcal{O} \ar[rr]^-{(a, b) \mapsto ab} & &
\mathcal{O}
}
$$
est cartésien. Comme $f^{-1}$ est exact, on en déduit que
$$
\xymatrix{
f^{-1}(\mathcal{O}^*)
\ar[d]_{u \mapsto (u, u^{-1})} \ar[rr] & &
{*} \ar[d]^{1} \\
f^{-1}\mathcal{O} \times f^{-1}\mathcal{O} \ar[rr]^-{(a, b) \mapsto ab} & &
f^{-1}\mathcal{O}
}
$$
est cartésien, ce qui prouve la première assertion. Pour la seconde,
remarquons que $(\mathcal{C}', f^{-1}\mathcal{O})$ est un site localement annelé
par le lemme \ref{sites-modules-lemma-locally-ringed-intrinsic}; l'assertion a donc un sens.
La première partie montre alors que le morphisme est un morphisme de topos
localement annelés.
\end{proof}

\begin{lemma}
\label{sites-modules-lemma-localize-morphism-locally-ringed-topoi}
Localisation des sites et topos localement annelés.
\begin{enumerate}
\item Soit $(\mathcal{C}, \mathcal{O})$ un site localement annelé et soit $U$
un objet de $\mathcal{C}$. Alors la localisation
$(\mathcal{C}/U, \mathcal{O}_U)$ est un site localement annelé, et le morphisme
de localisation
$$
(j_U, j_U^\sharp) :
(\Sh(\mathcal{C}/U), \mathcal{O}_U)
\to
(\Sh(\mathcal{C}), \mathcal{O})
$$
est un morphisme de topos localement annelés.
\item Soit $(\mathcal{C}, \mathcal{O})$ un site localement annelé et soit
$f : V \to U$ un morphisme de $\mathcal{C}$. Alors le morphisme
$$
(j, j^\sharp) :
(\Sh(\mathcal{C}/V), \mathcal{O}_V)
\to
(\Sh(\mathcal{C}/U), \mathcal{O}_U)
$$
du lemme \ref{sites-modules-lemma-relocalize} est un morphisme de topos localement annelés.
\item Soit
$(f, f^\sharp) :
(\mathcal{C}, \mathcal{O})
\longrightarrow
(\mathcal{D}, \mathcal{O}')$
un morphisme de sites localement annelés, où $f$ est donné par le foncteur
continu $u : \mathcal{D} \to \mathcal{C}$. Soit $V$ un objet de $\mathcal{D}$
et posons $U = u(V)$. Alors le morphisme
$$
(f', (f')^\sharp) :
(\Sh(\mathcal{C}/U), \mathcal{O}_U)
\to
(\Sh(\mathcal{D}/V), \mathcal{O}'_V)
$$
du lemme \ref{sites-modules-lemma-localize-morphism-ringed-sites} est un morphisme de sites
localement annelés.
\item Soit
$(f, f^\sharp) :
(\mathcal{C}, \mathcal{O})
\longrightarrow
(\mathcal{D}, \mathcal{O}')$
un morphisme de sites localement annelés, où $f$ est donné par le foncteur
continu $u : \mathcal{D} \to \mathcal{C}$. Soient
$V \in \Ob(\mathcal{D})$, $U \in \Ob(\mathcal{C})$ et $c : U \to u(V)$.
Alors le morphisme
$$
(f_c, (f_c)^\sharp) :
(\Sh(\mathcal{C}/U), \mathcal{O}_U)
\to
(\Sh(\mathcal{D}/V), \mathcal{O}'_V)
$$
du lemme \ref{sites-modules-lemma-relocalize-morphism-ringed-sites} est un morphisme de topos
localement annelés.
\item Soit $(\Sh(\mathcal{C}), \mathcal{O})$ un topos localement annelé et
soit $\mathcal{F}$ un faisceau sur $\mathcal{C}$. Alors la localisation
$(\Sh(\mathcal{C})/\mathcal{F}, \mathcal{O}_\mathcal{F})$
est un topos localement annelé, et le morphisme de localisation
$$
(j_\mathcal{F}, j_\mathcal{F}^\sharp) :
(\Sh(\mathcal{C})/\mathcal{F}, \mathcal{O}_\mathcal{F})
\to
(\Sh(\mathcal{C}), \mathcal{O})
$$
est un morphisme de topos localement annelés.
\item Soit $(\Sh(\mathcal{C}), \mathcal{O})$ un topos localement annelé.
Soit $s : \mathcal{G} \to \mathcal{F}$ une application de faisceaux sur
$\mathcal{C}$. Alors le morphisme
$$
(j, j^\sharp) :
(\Sh(\mathcal{C})/\mathcal{G}, \mathcal{O}_\mathcal{G})
\longrightarrow
(\Sh(\mathcal{C})/\mathcal{F}, \mathcal{O}_\mathcal{F})
$$
du lemme \ref{sites-modules-lemma-relocalize-ringed-topos} est un morphisme de topos
localement annelés.
\item Soit
$f :
(\Sh(\mathcal{C}), \mathcal{O})
\longrightarrow
(\Sh(\mathcal{D}), \mathcal{O}')$
un morphisme de topos localement annelés. Soit $\mathcal{G}$ un faisceau sur
$\mathcal{D}$ et posons $\mathcal{F} = f^{-1}\mathcal{G}$. Alors le morphisme
$$
(f', (f')^\sharp) :
(\Sh(\mathcal{C})/\mathcal{F}, \mathcal{O}_\mathcal{F})
\longrightarrow
(\Sh(\mathcal{D})/\mathcal{G}, \mathcal{O}'_\mathcal{G})
$$
du lemme \ref{sites-modules-lemma-localize-morphism-ringed-topoi} est un morphisme de topos
localement annelés.
\item Soit
$f :
(\Sh(\mathcal{C}), \mathcal{O})
\longrightarrow
(\Sh(\mathcal{D}), \mathcal{O}')$
un morphisme de topos localement annelés. Soient $\mathcal{G}$ un faisceau sur
$\mathcal{D}$ et $\mathcal{F}$ un faisceau sur $\mathcal{C}$, et soit
$s : \mathcal{F} \to f^{-1}\mathcal{G}$ un morphisme de faisceaux. Alors le morphisme
$$
(f_s, (f_s)^\sharp) :
(\Sh(\mathcal{C})/\mathcal{F}, \mathcal{O}_\mathcal{F})
\longrightarrow
(\Sh(\mathcal{D})/\mathcal{G}, \mathcal{O}'_\mathcal{G})
$$
du lemme \ref{sites-modules-lemma-relocalize-morphism-ringed-topoi} est un morphisme de topos
localement annelés.
\end{enumerate}
\end{lemma}

\begin{proof}
La partie (1) est claire, puisque $\mathcal{O}_U$ n'est autre que la
restriction de $\mathcal{O}$; les lemmes
\ref{sites-modules-lemma-locally-ringed-intrinsic} et
\ref{sites-modules-lemma-locally-ringed-intrinsic-morphism} s'appliquent. La partie (2) est
claire, car $(j, j^\sharp)$ est une localisation d'un site localement annelé,
donc (1) s'applique. La partie (3) est également claire, puisque
$(f')^\sharp$ est la restriction de $f^\sharp$ au topos
$\Sh(\mathcal{C}/U)$; voir la démonstration du
lemme \ref{sites-modules-lemma-localize-morphism-ringed-topoi}. Le diagramme de la
définition \ref{sites-modules-definition-morphism-locally-ringed-topoi} pour $f'$ est donc
la restriction du diagramme correspondant pour $f$, et la restriction est un
foncteur exact. La partie (4) résulte formellement de (2) et (3). Les parties
(5), (6), (7) et (8) résultent respectivement de (1), (2), (3) et (4), en
agrandissant les sites comme dans le
lemme \ref{sites-modules-lemma-morphism-ringed-topoi-comes-from-morphism-ringed-sites} et en
traduisant tout en termes de sites et de morphismes de sites au moyen des comparaisons des lemmes
\ref{sites-modules-lemma-localize-compare},
\ref{sites-modules-lemma-relocalize-compare},
\ref{sites-modules-lemma-localize-morphism-compare} et
\ref{sites-modules-lemma-relocalize-morphism-compare}.
(On pourrait aussi appliquer directement aux parties (5), (6), (7) et (8) les
mêmes arguments que dans les démonstrations de (1), (2), (3) et (4).)
\end{proof}






\section{Image directe exceptionnelle pour les modules}
\label{sites-modules-section-lower-shriek-modules}

\noindent
Dans cette section, nous étendons au cas des faisceaux de modules la
construction de $g_!$ étudiée dans la section
\ref{sites-modules-section-exactness-lower-shriek}.

\begin{lemma}
\label{sites-modules-lemma-lower-shriek-modules}
Soit $u : \mathcal{C} \to \mathcal{D}$ un foncteur continu et cocontinu entre
sites. Notons $g : \Sh(\mathcal{C}) \to \Sh(\mathcal{D})$ le morphisme de
topos associé. Soit $\mathcal{O}_\mathcal{D}$ un faisceau d'anneaux sur
$\mathcal{D}$ et posons
$\mathcal{O}_\mathcal{C} = g^{-1}\mathcal{O}_\mathcal{D}$. Ainsi $g$ devient
un morphisme de topos annelés tel que $g^* = g^{-1}$. Il existe alors un foncteur
$$
g_! :
\textit{Mod}(\mathcal{O}_\mathcal{C})
\longrightarrow
\textit{Mod}(\mathcal{O}_\mathcal{D})
$$
adjoint à gauche de $g^*$.
\end{lemma}

\begin{proof}
Soit $U$ un objet de $\mathcal{C}$. Pour tout
$\mathcal{O}_\mathcal{D}$-module $\mathcal{G}$, on a
\begin{align*}
\Hom_{\mathcal{O}_\mathcal{C}}(j_{U!}\mathcal{O}_U, g^{-1}\mathcal{G})
& =
g^{-1}\mathcal{G}(U) \\
& =
\mathcal{G}(u(U)) \\
& =
\Hom_{\mathcal{O}_\mathcal{D}}(j_{u(U)!}\mathcal{O}_{u(U)}, \mathcal{G})
\end{align*}
car $g^{-1}$ se décrit par restriction; voir Sites,
lemme \ref{sites-lemma-when-shriek}. Une formule analogue vaut évidemment
pour une somme directe de modules de la forme $j_{U!}\mathcal{O}_U$. Par
Homologie, lemme \ref{homology-lemma-partially-defined-adjoint}, et par le
lemme \ref{sites-modules-lemma-module-quotient-flat}, on voit que $g_!$ existe.
\end{proof}

\begin{remark}
\label{sites-modules-remark-when-shriek-equal}
Attention ! Soient $u : \mathcal{C} \to \mathcal{D}$, $g$,
$\mathcal{O}_\mathcal{D}$ et $\mathcal{O}_\mathcal{C}$ comme dans le
lemme \ref{sites-modules-lemma-lower-shriek-modules}. En général, le diagramme
$$
\xymatrix{
\textit{Mod}(\mathcal{O}_\mathcal{C}) \ar[r]_{g_!} \ar[d]_{forget} &
\textit{Mod}(\mathcal{O}_\mathcal{D}) \ar[d]^{forget} \\
\textit{Ab}(\mathcal{C}) \ar[r]^{g^{Ab}_!} &
\textit{Ab}(\mathcal{D})
}
$$
ne commute {\bf pas} (ici, $g^{Ab}_!$ est celui du
lemme \ref{sites-modules-lemma-g-shriek-adjoint}). Il existe une transformation de foncteurs
$$
g_!^{Ab} \circ forget \longrightarrow forget \circ g_!
$$
La démonstration du lemme \ref{sites-modules-lemma-lower-shriek-modules} montre que c'est un
isomorphisme si et seulement si
$g^{Ab}_!j_{U!}\mathcal{O}_U \to g_!j_{U!}\mathcal{O}_U$
est un isomorphisme pour tout objet $U$ de $\mathcal{C}$. Comme
$g_!j_{U!}\mathcal{O}_U = j_{u(U)!}\mathcal{O}_{u(U)}$, cette condition équivaut à ce que
$$
g^{Ab}_!j_{U!}\mathcal{O}_U \longrightarrow j_{u(U)!}\mathcal{O}_{u(U)}
$$
soit un isomorphisme pour tout objet $U$ de $\mathcal{C}$. Pour un tel $U$,
on obtient un diagramme commutatif
$$
\xymatrix{
\mathcal{C}/U \ar[r]_-{u'} \ar[d]_{j_U} & \mathcal{D}/u(U) \ar[d]^{j_{u(U)}} \\
\mathcal{C} \ar[r]^u & \mathcal{D}
}
$$
de foncteurs cocontinus de sites; voir Sites,
lemme \ref{sites-lemma-localize-cocontinuous}. Par conséquent,
$g^{Ab}_!j_{U!} = j_{u(U)!}(g')^{Ab}_!$, où
$g' : \Sh(\mathcal{C}/U) \to \Sh(\mathcal{D}/u(U))$ est le morphisme de
topos induit par le foncteur cocontinu $u'$. On voit donc que
$g_! = g^{Ab}_!$ si l'application canonique
\begin{equation}
\label{sites-modules-equation-compare-on-localizations}
(g')^{Ab}_!\mathcal{O}_U \longrightarrow \mathcal{O}_{u(U)}
\end{equation}
est un isomorphisme pour tout objet $U$ de $\mathcal{C}$.
\end{remark}

\noindent
Les deux résultats suivants sont de nature légèrement différente.

\begin{lemma}
\label{sites-modules-lemma-special-square-cocontinuous}
Supposons donné un diagramme commutatif
$$
\xymatrix{
(\Sh(\mathcal{C}'), \mathcal{O}_{\mathcal{C}'})
\ar[r]_{(g', (g')^\sharp)} \ar[d]_{(f', (f')^\sharp)} &
(\Sh(\mathcal{C}), \mathcal{O}_\mathcal{C}) \ar[d]^{(f, f^\sharp)} \\
(\Sh(\mathcal{D}'), \mathcal{O}_{\mathcal{D}'}) \ar[r]^{(g, g^\sharp)} &
(\Sh(\mathcal{D}), \mathcal{O}_\mathcal{D})
}
$$
de topos annelés. Supposons que :
\begin{enumerate}
\item $f$, $f'$, $g$ et $g'$ correspondent aux foncteurs cocontinus
$u$, $u'$, $v$ et $v'$ comme dans Sites,
lemme \ref{sites-lemma-cocontinuous-morphism-topoi};
\item $v \circ u' = u \circ v'$,
\item $v$ et $v'$ sont à la fois continus et cocontinus;
\item pour tout objet $V'$ de $\mathcal{D}'$, le foncteur
${}^{u'}_{V'}\mathcal{I} \to {}^{\ \ \ u}_{v(V')}\mathcal{I}$
induit par $v$ est cofinal; et
\item $g^{-1}\mathcal{O}_{\mathcal{D}} = \mathcal{O}_{\mathcal{D}'}$
et $(g')^{-1}\mathcal{O}_{\mathcal{C}} = \mathcal{O}_{\mathcal{C}'}$.
\end{enumerate}
Alors, sur les modules, on a $f'_* \circ (g')^* = g^* \circ f_*$ et
$g'_! \circ (f')^{-1} = f^{-1} \circ g_!$.
\end{lemma}

\begin{proof}
On a $(g')^*\mathcal{F} = (g')^{-1}\mathcal{F}$ et
$g^*\mathcal{G} = g^{-1}\mathcal{G}$ en vertu de (5). La première égalité
résulte donc immédiatement de l'égalité correspondante dans Sites,
lemme \ref{sites-lemma-special-square-cocontinuous}. Comme les adjoints à
gauche $g_!$ et $g'_!$ de $g^*$ et $(g')^*$ existent par le
lemme \ref{sites-modules-lemma-lower-shriek-modules}, la seconde égalité résulte de l'unicité
des foncteurs adjoints.
\end{proof}

\begin{lemma}
\label{sites-modules-lemma-special-square-continuous}
Considérons un diagramme commutatif
$$
\xymatrix{
(\Sh(\mathcal{C}'), \mathcal{O}_{\mathcal{C}'})
\ar[r]_{(g', (g')^\sharp)} \ar[d]_{(f', (f')^\sharp)} &
(\Sh(\mathcal{C}), \mathcal{O}_\mathcal{C}) \ar[d]^{(f, f^\sharp)} \\
(\Sh(\mathcal{D}'), \mathcal{O}_{\mathcal{D}'}) \ar[r]^{(g, g^\sharp)} &
(\Sh(\mathcal{D}), \mathcal{O}_\mathcal{D})
}
$$
de topos annelés et supposons donnés des foncteurs
$$
\xymatrix{
\mathcal{C}' \ar[r]_{v'} &
\mathcal{C} \\
\mathcal{D}' \ar[r]^v \ar[u]^{u'} &
\mathcal{D} \ar[u]_u
}
$$
tels que (avec les notations de Sites, sections
\ref{sites-section-morphism-sites} et
\ref{sites-section-cocontinuous-morphism-topoi}) :
\begin{enumerate}
\item $u$ et $u'$ sont continus et donnent les morphismes $f$ et $f'$;
\item $v$ et $v'$ sont cocontinus et donnent les morphismes $g$ et $g'$;
\item $u \circ v = v' \circ u'$,
\item $v$ et $v'$ sont à la fois continus et cocontinus; et
\item $g^{-1}\mathcal{O}_{\mathcal{D}} = \mathcal{O}_{\mathcal{D}'}$
et $(g')^{-1}\mathcal{O}_{\mathcal{C}} = \mathcal{O}_{\mathcal{C}'}$.
\end{enumerate}
Alors, sur les modules, $f'_* \circ (g')^* = g^* \circ f_*$ et
$g'_! \circ (f')^{-1} = f^{-1} \circ g_!$.
\end{lemma}

\begin{proof}
On a $(g')^*\mathcal{F} = (g')^{-1}\mathcal{F}$ et
$g^*\mathcal{G} = g^{-1}\mathcal{G}$ en vertu de (5). La première égalité
résulte immédiatement de l'égalité correspondante dans Sites,
lemme \ref{sites-lemma-special-square-continuous}. Comme les adjoints à gauche
$g_!$ et $g'_!$ de $g^*$ et $(g')^*$ existent par le
lemme \ref{sites-modules-lemma-lower-shriek-modules}, la seconde égalité résulte de l'unicité
des foncteurs adjoints.
\end{proof}







\section{Faisceaux constants}
\label{sites-modules-section-constant}

\noindent
Soit $E$ un ensemble et soit $\mathcal{C}$ un site. Le faisceau constant de
valeur $E$ est noté $\underline{E}$; voir Sites,
exemple \ref{sites-example-constant-sheaf}. Si $E$ est un groupe abélien, un
anneau, un module, etc., alors $\underline{E}$ est un faisceau de groupes
abéliens, d'anneaux, de modules, etc.

\begin{lemma}
\label{sites-modules-lemma-constant-exact}
Soit $\mathcal{C}$ un site. Si $0 \to A \to B \to C \to 0$ est une suite
exacte courte de groupes abéliens, alors
$0 \to \underline{A} \to \underline{B} \to \underline{C} \to 0$
est une suite exacte de faisceaux abéliens; elle est même exacte comme suite
de préfaisceaux abéliens.
\end{lemma}

\begin{proof}
Comme la faisceautisation est exacte, il est clair que
$0 \to \underline{A} \to \underline{B} \to \underline{C} \to 0$
est une suite exacte de faisceaux abéliens. Ainsi
$0 \to \underline{A} \to \underline{B} \to \underline{C}$
est une suite exacte de préfaisceaux abéliens. Pour voir que
$\underline{B} \to \underline{C}$ est surjective, choisissons une section
ensembliste $s : C \to B$. Elle induit une section
$\underline{s} : \underline{C} \to \underline{B}$ de faisceaux d'ensembles,
inverse à droite de la surjection $\underline{B} \to \underline{C}$.
\end{proof}

\begin{lemma}
\label{sites-modules-lemma-tensor-with-finitely-presented}
Soit $\mathcal{C}$ un site. Soit $\Lambda$ un anneau et soient $M$ et $Q$ des
$\Lambda$-modules. Si $Q$ est un $\Lambda$-module de présentation finie, alors
$\underline{M \otimes_\Lambda Q}(U) = \underline{M}(U) \otimes_\Lambda Q$
pour tout $U \in \Ob(\mathcal{C})$.
\end{lemma}

\begin{proof}
Choisissons une présentation
$\Lambda^{\oplus m} \to \Lambda^{\oplus n} \to Q \to 0$. Elle donne une
suite exacte $M^{\oplus m} \to M^{\oplus n} \to M \otimes Q \to 0$.
Par le lemme \ref{sites-modules-lemma-constant-exact}, on obtient la suite exacte
$$
\underline{M}(U)^{\oplus m} \to
\underline{M}(U)^{\oplus n} \to
\underline{M \otimes Q}(U) \to 0
$$
ce qui démontre le lemme. (Remarquons que la prise de sections sur $U$ commute
toujours aux sommes directes finies, mais non aux sommes directes arbitraires.)
\end{proof}

\begin{lemma}
\label{sites-modules-lemma-flat-sections}
Soit $\mathcal{C}$ un site. Soit $\Lambda$ un anneau cohérent et soit $M$ un
$\Lambda$-module plat. Pour $U \in \Ob(\mathcal{C})$, le module
$\underline{M}(U)$ est un $\Lambda$-module plat.
\end{lemma}

\begin{proof}
Soit $I \subset \Lambda$ un idéal de type fini. Par Algèbre,
lemme \ref{algebra-lemma-flat}, il suffit de montrer que
$\underline{M}(U) \otimes_\Lambda I \to \underline{M}(U)$
est injective. Comme $\Lambda$ est cohérent, $I$ est un $\Lambda$-module de
présentation finie. Par le lemme \ref{sites-modules-lemma-tensor-with-finitely-presented},
$\underline{M}(U) \otimes I = \underline{M \otimes I}$. Puisque $M$ est plat,
l'application $M \otimes I \to M$ est injective; il en va donc de même de
$\underline{M \otimes I} \to \underline{M}$.
\end{proof}

\begin{lemma}
\label{sites-modules-lemma-completion-flat}
Soit $\mathcal{C}$ un site. Soit $\Lambda$ un anneau noethérien et soit
$I \subset \Lambda$ un idéal. Le faisceau
$\underline{\Lambda}^\wedge = \lim \underline{\Lambda/I^n}$
est une $\underline{\Lambda}$-algèbre plate. De plus, on a les
identifications canoniques
$$
\underline{\Lambda}/I\underline{\Lambda} =
\underline{\Lambda}/\underline{I} =
\underline{\Lambda}^\wedge/I\underline{\Lambda}^\wedge =
\underline{\Lambda}^\wedge/\underline{I} \cdot \underline{\Lambda}^\wedge =
\underline{\Lambda}^\wedge/\underline{I}^\wedge =
\underline{\Lambda/I}
$$
où $\underline{I}^\wedge = \lim \underline{I/I^n}$.
\end{lemma}

\begin{proof}
Pour montrer que $\underline{\Lambda}^\wedge$ est plat, il suffit de vérifier
que $\underline{\Lambda}^\wedge(U)$ est un $\Lambda$-module plat pour tout
$U \in \Ob(\mathcal{C})$; voir les lemmes
\ref{sites-modules-lemma-flatness-presheaves} et \ref{sites-modules-lemma-flatness-sheafification}.
Par le lemme \ref{sites-modules-lemma-flat-sections},
$$
\underline{\Lambda}^\wedge(U) = \lim \underline{\Lambda/I^n}(U)
$$
est la limite d'un système de $\Lambda/I^n$-modules plats. Le
lemme \ref{sites-modules-lemma-constant-exact} montre que les applications de transition
sont surjectives. On conclut par Compléments d'algèbre,
lemme \ref{more-algebra-lemma-limit-flat}.

\medskip\noindent
Pour établir les égalités, remarquons que
$\underline{\Lambda}(U)/I\underline{\Lambda}(U) = \underline{\Lambda/I}(U)$
par le lemme \ref{sites-modules-lemma-tensor-with-finitely-presented}. Il s'ensuit que
$\underline{\Lambda}/I\underline{\Lambda} =
\underline{\Lambda}/\underline{I} = \underline{\Lambda/I}$. Le système de
suites exactes courtes
$$
0 \to \underline{I/I^n}(U) \to \underline{\Lambda/I^n}(U) \to
\underline{\Lambda/I}(U) \to 0
$$
a des applications de transition surjectives et donne donc une suite exacte courte
$$
0 \to \lim \underline{I/I^n}(U) \to \lim \underline{\Lambda/I^n}(U) \to
\lim \underline{\Lambda/I}(U) \to 0
$$
voir Homologie, lemme \ref{homology-lemma-Mittag-Leffler}. Ainsi,
$\underline{\Lambda}^\wedge/\underline{I}^\wedge =
\underline{\Lambda/I}$. Comme
$$
I \underline{\Lambda}^\wedge \subset
\underline{I} \cdot \underline{\Lambda}^\wedge \subset
\underline{I}^\wedge
$$
il suffit de montrer que
$I \underline{\Lambda}^\wedge(U) = \underline{I}^\wedge(U)$
pour tout $U$. Choisissons des générateurs $I = (f_1, \ldots, f_r)$.
Pour tout $n$, on obtient une suite exacte courte
$$
0 \to K_n/(I^n)^{\oplus r} \to (\Lambda/I^n)^{\oplus r}
\xrightarrow{(f_1, \ldots, f_r)}
I/I^{n + 1} \to 0
$$
où $K_n =
\{(x_1, \ldots, x_r) \in \Lambda^{\oplus r} \mid \sum x_i f_i \in I^{n + 1}\}$.
On obtient des suites exactes courtes
$$
0 \to \underline{K_n/(I^n)^{\oplus r}}(U) \to
\underline{(\Lambda/I^n)^{\oplus r}}(U) \to
\underline{I/I^{n + 1}}(U) \to 0
$$
Un calcul donne $K_n = K + (I^n)^{\oplus r}$; les applications de transition
$K_{n + 1}/(I^{n + 1})^{\oplus r} \to K_n/(I^n)^{\oplus r}$ sont donc
surjectives. Le système de modules de gauche a ainsi des applications de
transition surjectives et vérifie a fortiori la condition de Mittag--Leffler.
On voit donc que
$(f_1, \ldots, f_r) :
(\underline{\Lambda}^\wedge)^{\oplus r}(U) \to \underline{I}^\wedge(U)$
est surjective par Homologie, lemme \ref{homology-lemma-Mittag-Leffler}, ce
qu'il fallait démontrer.
\end{proof}

\begin{lemma}
\label{sites-modules-lemma-locally-constant-finite-type}
Soit $\mathcal{C}$ un site. Soit $\Lambda$ un anneau et soit $M$ un
$\Lambda$-module. Supposons que $\Sh(\mathcal{C})$ ne soit pas le topos vide. Alors
\begin{enumerate}
\item $\underline{M}$ est un faisceau de $\underline{\Lambda}$-modules de type
fini si et seulement si $M$ est un $\Lambda$-module fini; et
\item $\underline{M}$ est un faisceau de $\underline{\Lambda}$-modules de
présentation finie si et seulement si $M$ est un $\Lambda$-module de
présentation finie.
\end{enumerate}
\end{lemma}

\begin{proof}
Démonstration de (1). Si $M$ est engendré par $x_1, \ldots, x_r$, alors
$x_1, \ldots, x_r$ définissent des sections globales de $\underline{M}$ qui
l'engendrent; ainsi
$\underline{M}$ est de type fini. Réciproquement, supposons
$\underline{M}$ de type fini. Soit $U \in \mathcal{C}$ un objet qui n'est pas
vide au sens faisceautique (Sites, définition \ref{sites-definition-empty}).
Un tel objet existe puisque $\Sh(\mathcal{C})$ n'est pas le topos vide. Il
existe alors un recouvrement $\{U_i \to U\}$ et un nombre fini de sections
$s_{ij} \in \underline{M}(U_i)$ engendrant $\underline{M}|_{U_i}$. Après
raffinement du recouvrement, on peut supposer que les $s_{ij}$ proviennent
d'éléments $x_{ij}$ de $M$. Les $x_{ij}$ définissent des sections globales de
$\underline{M}$ dont les restrictions à $U$ engendrent $\underline{M}$.

\medskip\noindent
Supposons qu'il existe des éléments $x_1, \ldots, x_r$ de $M$ définissant des
sections globales de $\underline{M}$ qui engendrent $\underline{M}$ comme
faisceau de $\underline{\Lambda}$-modules. Montrons que $x_1, \ldots, x_r$
engendrent $M$ comme $\Lambda$-module. Soit $x \in M$. Il existe un recouvrement
$\{U_i \to U\}_{i \in I}$ et des $f_{i, j} \in \underline{\Lambda}(U_i)$ tels
que $x|_{U_i} = \sum f_{i, j} x_j|_{U_i}$. Après raffinement, on peut supposer
$f_{i, j} \in \Lambda$. Puisque $U$ n'est pas vide au sens faisceautique, il existe
au moins un $i \in I$ tel que $U_i$ ne soit pas non plus vide au sens
faisceautique. L'application
$M \to \underline{M}(U_i)$ est alors injective (détails omis). On conclut que
$x = \sum f_{i, j}x_j$ dans $M$, comme voulu.

\medskip\noindent
Démonstration de (2). Supposons $\underline{M}$ de présentation finie comme
$\underline{\Lambda}$-module. Par (1), $M$ est de type fini. Choisissons des
générateurs $x_1, \ldots, x_r$ de $M$ comme $\Lambda$-module. On obtient une
suite exacte courte $0 \to K \to \Lambda^{\oplus r} \to M \to 0$, qui devient
$$
0 \to \underline{K} \to \underline{\Lambda}^{\oplus r} \to \underline{M} \to 0
$$
par le lemme \ref{sites-modules-lemma-constant-exact}. Le
lemme \ref{sites-modules-lemma-kernel-surjection-finite-onto-finite-presentation} montre que
$\underline{K}$ est de type fini. Par (1), $K$ est donc un
$\Lambda$-module fini, et $M$ est donc un $\Lambda$-module de présentation finie.
\end{proof}




\section{Faisceaux localement constants}
\label{sites-modules-section-locally-constant}

\noindent
Voici la définition générale.

\begin{definition}
\label{sites-modules-definition-locally-constant}
Soit $\mathcal{C}$ un site. Soit $\mathcal{F}$ un faisceau d'ensembles, de
groupes, de groupes abéliens, d'anneaux, de modules sur un anneau fixé
$\Lambda$, etc.
\begin{enumerate}
\item On dit que $\mathcal{F}$ est un {\it faisceau constant} d'ensembles, de
groupes, de groupes abéliens, d'anneaux, de modules sur $\Lambda$, etc., s'il
est isomorphe, comme faisceau d'ensembles, de groupes, de groupes abéliens,
d'anneaux ou de modules sur $\Lambda$, à un faisceau
constant $\underline{E}$ comme dans la section \ref{sites-modules-section-constant}.
\item On dit que $\mathcal{F}$ est {\it localement constant} si, pour tout
objet $U$ de $\mathcal{C}$, il existe un recouvrement $\{U_i \to U\}$ tel que
$\mathcal{F}|_{U_i}$ soit un faisceau constant.
\item Si $\mathcal{F}$ est un faisceau d'ensembles ou de groupes, on dit que $\mathcal{F}$ est
{\it fini localement constant} si les valeurs constantes sont des ensembles
finis ou des groupes finis.
\end{enumerate}
\end{definition}

\begin{lemma}
\label{sites-modules-lemma-pullback-locally-constant}
Soit $f : \Sh(\mathcal{C}) \to \Sh(\mathcal{D})$ un morphisme de topos.
Si $\mathcal{G}$ est un faisceau localement constant d'ensembles, de groupes,
de groupes abéliens, d'anneaux, de modules sur un anneau fixé $\Lambda$, etc.,
sur $\mathcal{D}$, alors $f^{-1}\mathcal{G}$ possède la même propriété sur
$\mathcal{C}$.
\end{lemma}

\begin{proof}
La démonstration est omise.
\end{proof}

\begin{lemma}
\label{sites-modules-lemma-morphism-locally-constant}
Soit $\mathcal{C}$ un site possédant un objet final $X$.
\begin{enumerate}
\item Soit $\varphi : \mathcal{F} \to \mathcal{G}$ une application de
faisceaux d'ensembles localement constants sur $\mathcal{C}$. Si
$\mathcal{F}$ est fini localement constant, il existe un recouvrement
$\{U_i \to X\}$ tel que $\varphi|_{U_i}$ soit l'application de faisceaux
constants associée à une application d'ensembles.
\item Soit $\varphi : \mathcal{F} \to \mathcal{G}$ une application de
faisceaux de groupes abéliens localement constants sur $\mathcal{C}$. Si
$\mathcal{F}$ est fini localement constant, il existe un recouvrement
$\{U_i \to X\}$ tel que $\varphi|_{U_i}$ soit l'application de faisceaux
abéliens constants associée à un homomorphisme de groupes abéliens.
\item Soit $\Lambda$ un anneau et soit
$\varphi : \mathcal{F} \to \mathcal{G}$ une application de faisceaux de
$\Lambda$-modules localement constants sur $\mathcal{C}$. Si $\mathcal{F}$
est de type fini, il existe un recouvrement $\{U_i \to X\}$ tel que
$\varphi|_{U_i}$ soit l'application de faisceaux constants de
$\Lambda$-modules associée à une application de $\Lambda$-modules.
\end{enumerate}
\end{lemma}

\begin{proof}
La démonstration est omise.
\end{proof}

\begin{lemma}
\label{sites-modules-lemma-locally-constant}
Soit $\mathcal{C}$ un site. Soit $\Lambda$ un anneau. Soient $M$, $N$ des
$\Lambda$-modules et $\mathcal{F}, \mathcal{G}$ des faisceaux localement
constants de $\Lambda$-modules.
\begin{enumerate}
\item Si $M$ est de présentation finie, alors
$$
\underline{\Hom_\Lambda(M, N)} =
\SheafHom_{\underline{\Lambda}}(\underline{M}, \underline{N})
$$
\item Si $M$ et $N$ sont tous deux de présentation finie, alors
$$
\underline{\text{Isom}_\Lambda(M, N)} =
\mathit{Isom}_{\underline{\Lambda}}(\underline{M}, \underline{N})
$$
\item Si $\mathcal{F}$ est de présentation finie, alors
$\SheafHom_{\underline{\Lambda}}(\mathcal{F}, \mathcal{G})$
est un faisceau localement constant de $\Lambda$-modules.
\item Si $\mathcal{F}$ et $\mathcal{G}$ sont tous deux de présentation finie, alors
$\mathit{Isom}_{\underline{\Lambda}}(\mathcal{F}, \mathcal{G})$
est un faisceau d'ensembles localement constant.
\end{enumerate}
\end{lemma}

\begin{proof}
Démonstration de (1). Posons $E = \Hom_\Lambda(M, N)$. Nous voulons montrer que
l'application canonique
$$
\underline{E}
\longrightarrow
\SheafHom_{\underline{\Lambda}}(\underline{M}, \underline{N})
$$
est un isomorphisme. Le module $M$ admet une présentation
$\Lambda^{\oplus s} \to \Lambda^{\oplus t} \to M \to 0$.
Alors $E$ s'insère dans une suite exacte
$$
0 \to E \to \Hom_\Lambda(\Lambda^{\oplus t}, N) \to
\Hom_\Lambda(\Lambda^{\oplus s}, N)
$$
et l'on a de même
$$
0 \to
\SheafHom_{\underline{\Lambda}}(\underline{M}, \underline{N}) \to
\SheafHom_{\underline{\Lambda}}(\underline{\Lambda^{\oplus t}}, \underline{N})
\to
\SheafHom_{\underline{\Lambda}}(\underline{\Lambda^{\oplus s}}, \underline{N})
$$
Cela ramène la question au cas où $M$ est libre fini, dans lequel le résultat
est clair.

\medskip\noindent
Démonstration de (3). La question est locale sur $\mathcal{C}$; on peut donc
supposer $\mathcal{F} = \underline{M}$ et $\mathcal{G} = \underline{N}$ pour des
$\Lambda$-modules $M$ et $N$. Par le
lemme \ref{sites-modules-lemma-locally-constant-finite-type}, $M$ est de présentation finie.
Le résultat découle alors de (1).

\medskip\noindent
Les parties (2) et (4) résultent de (1) et (3), ainsi que du fait que
$\mathit{Isom}$ peut être vu comme le sous-faisceau des sections de
$\SheafHom_{\underline{\Lambda}}(\mathcal{F}, \mathcal{G})$
qui admettent un inverse dans
$\SheafHom_{\underline{\Lambda}}(\mathcal{G}, \mathcal{F})$.
\end{proof}

\begin{lemma}
\label{sites-modules-lemma-kernel-finite-locally-constant}
Soit $\mathcal{C}$ un site.
\begin{enumerate}
\item La catégorie des faisceaux d'ensembles finis localement constants est
stable par limites et colimites finies dans $\Sh(\mathcal{C})$.
\item La catégorie des faisceaux abéliens finis localement constants est une
sous-catégorie de Serre faible de $\textit{Ab}(\mathcal{C})$.
\item Soit $\Lambda$ un anneau noethérien. La catégorie des faisceaux de
$\Lambda$-modules localement constants et de type fini sur $\mathcal{C}$ est
une sous-catégorie de Serre faible de $\textit{Mod}(\mathcal{C}, \Lambda)$.
\end{enumerate}
\end{lemma}

\begin{proof}
Démonstration de (1). On peut travailler localement sur $\mathcal{C}$. Par le
lemme \ref{sites-modules-lemma-morphism-locally-constant}, on peut donc supposer donné un
diagramme fini d'ensembles finis tel que notre diagramme de faisceaux soit le
diagramme associé de faisceaux constants. Il suffit alors de prendre la limite
ou la colimite dans la catégorie des ensembles, puis le faisceau constant
associé. Certains détails sont omis.

\medskip\noindent
Pour démontrer (2) et (3), appliquons le critère d'Homologie,
lemme \ref{homology-lemma-characterize-weak-serre-subcategory}. L'existence
des noyaux et conoyaux s'établit comme ci-dessus. L'hypothèse noethérienne dans
(3) garantit que le noyau d'une application de $\Lambda$-modules finis est un
$\Lambda$-module fini. Pour voir que la catégorie est stable par extensions
(dans le cas des faisceaux de $\Lambda$-modules), considérons une extension de
faisceaux de $\Lambda$-modules
$$
0 \to \mathcal{F} \to \mathcal{E} \to \mathcal{G} \to 0
$$
sur $\mathcal{C}$, où $\mathcal{F}$ et $\mathcal{G}$ sont de type fini et
localement constants. Après localisation sur $\mathcal{C}$, on peut supposer
$\mathcal{F}$ et $\mathcal{G}$ constants, c'est-à-dire que l'on obtient
$$
0 \to \underline{M} \to \mathcal{E} \to \underline{N} \to 0
$$
pour certains $\Lambda$-modules $M, N$. Choisissons des générateurs
$y_1, \ldots, y_m$ de $N$, d'où une suite exacte courte
$0 \to K \to \Lambda^{\oplus m} \to N \to 0$ de $\Lambda$-modules. Après une nouvelle
localisation, on peut supposer que $y_j$ se relève en une section $s_j$ de
$\mathcal{E}$. Ainsi $\mathcal{E}$ est un poussé comme dans le diagramme suivant
$$
\xymatrix{
0 \ar[r] &
\underline{K} \ar[d] \ar[r] &
\underline{\Lambda^{\oplus m}} \ar[d] \ar[r] &
\underline{N} \ar[d] \ar[r] & 0 \\
0 \ar[r] &
\underline{M} \ar[r] &
\mathcal{E} \ar[r] &
\underline{N} \ar[r] & 0
}
$$
Par une nouvelle application du lemme \ref{sites-modules-lemma-morphism-locally-constant}
(et parce que $K$ est un $\Lambda$-module fini puisque $\Lambda$ est
noethérien), on voit que l'application
$\underline{K} \to \underline{M}$ est localement constante, ce qui permet de conclure.
\end{proof}

\begin{lemma}
\label{sites-modules-lemma-tensor-product-locally-constant}
Soit $\mathcal{C}$ un site et soit $\Lambda$ un anneau. Le produit tensoriel de
deux faisceaux localement constants de $\Lambda$-modules sur $\mathcal{C}$ est
un faisceau localement constant de $\Lambda$-modules.
\end{lemma}

\begin{proof}
La démonstration est omise.
\end{proof}











\section{Localisation des faisceaux d'anneaux}
\label{sites-modules-section-localizing-sheaves-rings}

\noindent
Soit $(\mathcal{C}, \mathcal{O})$ un site annelé. Soit
$\mathcal{S} \subset \mathcal{O}$ un sous-préfaisceau d'ensembles tel que,
pour tout $U \in \Ob(\mathcal{C})$, l'ensemble
$\mathcal{S}(U) \subset \mathcal{O}(U)$ soit une partie multiplicative; voir
Algèbre, définition \ref{algebra-definition-multiplicative-subset}. On peut
alors considérer le préfaisceau d'anneaux
$$
\mathcal{S}^{-1}\mathcal{O} :
U \longmapsto \mathcal{S}(U)^{-1}\mathcal{O}(U).
$$
L'application de restriction envoie la section $f/s$, avec
$f \in \mathcal{O}(U)$ et $s \in \mathcal{S}(U)$, sur
$(f|_V)/(s|_V)$ pour $V \to U$ dans $\mathcal{C}$.

\begin{lemma}
\label{sites-modules-lemma-simple-invert}
Dans la situation ci-dessus, l'application vers le faisceautisé
$$
\mathcal{O} \longrightarrow (\mathcal{S}^{-1}\mathcal{O})^\#
$$
est un homomorphisme de faisceaux d'anneaux ayant la propriété universelle
suivante : pour tout homomorphisme de faisceaux d'anneaux
$\mathcal{O} \to \mathcal{A}$ qui envoie chaque section locale de
$\mathcal{S}$ sur une section inversible de $\mathcal{A}$, il existe une
unique factorisation $(\mathcal{S}^{-1}\mathcal{O})^\# \to \mathcal{A}$.
\end{lemma}

\begin{proof}
La démonstration est omise.
\end{proof}

\noindent
Soit $(\mathcal{C}, \mathcal{O})$ un site annelé. Soit
$\mathcal{S} \subset \mathcal{O}$ un sous-préfaisceau d'ensembles tel que,
pour tout $U \in \mathcal{C}$, l'ensemble
$\mathcal{S}(U) \subset \mathcal{O}(U)$ soit une partie multiplicative. Soit
$\mathcal{F}$ un faisceau de $\mathcal{O}$-modules. On peut alors considérer
le préfaisceau de $\mathcal{S}^{-1}\mathcal{O}$-modules
$$
\mathcal{S}^{-1}\mathcal{F} :
U \longmapsto \mathcal{S}(U)^{-1}\mathcal{F}(U).
$$
L'application de restriction envoie la section $t/s$, avec
$t \in \mathcal{F}(U)$ et $s \in \mathcal{S}(U)$, sur
$(t|_V)/(s|_V)$ si $V \to U$ est un morphisme de $\mathcal{C}$. Ainsi,
$\mathcal{S}^{-1}\mathcal{F}$ est un préfaisceau de
$\mathcal{S}^{-1}\mathcal{O}$-modules.

\begin{lemma}
\label{sites-modules-lemma-simple-invert-module}
Dans la situation ci-dessus, l'application vers le faisceautisé
$$
\mathcal{F} \longrightarrow (\mathcal{S}^{-1}\mathcal{F})^\#
$$
possède la propriété universelle suivante : pour tout homomorphisme de
$\mathcal{O}$-modules $\mathcal{F} \to \mathcal{G}$ tel que chaque section
locale de $\mathcal{S}$ agisse de façon inversible sur $\mathcal{G}$, il
existe une unique factorisation
$(\mathcal{S}^{-1}\mathcal{F})^\# \to \mathcal{G}$. De plus, on a
$$
(\mathcal{S}^{-1}\mathcal{F})^\#
=
(\mathcal{S}^{-1}\mathcal{O})^\# \otimes_\mathcal{O} \mathcal{F}
$$
comme faisceaux de $(\mathcal{S}^{-1}\mathcal{O})^\#$-modules.
\end{lemma}

\begin{proof}
La démonstration est omise.
\end{proof}









\section{Faisceaux d'ensembles pointés}
\label{sites-modules-section-pointed}

\noindent
Dans cette section, nous rassemblons quelques faits sur les faisceaux
d'ensembles pointés, que nous n'avions jusqu'ici mentionnés que pour les
faisceaux abéliens.

\medskip\noindent
Un ensemble pointé est une paire $(S, 0)$, où $S$ est un ensemble et
$0 \in S$ est un élément de $S$. Un morphisme $(S, 0) \to (S', 0')$ d'ensembles pointés est simplement
une application $S \to S'$ envoyant $0$ sur $0'$. Par abus de notation, nous
dirons « soit $S$ un ensemble pointé » pour signifier que $S$ est muni d'un
élément distingué $0 \in S$. Un faisceau d'ensembles pointés est la même chose
qu'un faisceau d'ensembles $\mathcal{F}$ muni d'un « pointage »
$0 : * \to \mathcal{F}$, où $*$ est le faisceau final
(Sites, exemple \ref{sites-example-singleton-sheaf}).

\medskip\noindent
Pour un morphisme de sites ou de topos, il existe des foncteurs image directe
et image inverse sur les catégories de faisceaux d'ensembles pointés; voir
Sites, section \ref{sites-section-sheaves-algebraic-structures}. Ils se
construisent en prenant l'image directe, respectivement l'image inverse, du
faisceau d'ensembles sous-jacent et en le pointant convenablement (en utilisant
que l'image inverse du faisceau final est le faisceau final).

\medskip\noindent
Soit $u : \mathcal{C} \to \mathcal{D}$ un foncteur continu et cocontinu entre
sites. Soit $g : \Sh(\mathcal{C}) \to \Sh(\mathcal{D})$ le morphisme de topos
associé à $u$; voir Sites,
lemme \ref{sites-lemma-cocontinuous-morphism-topoi}. Sur les faisceaux
d'ensembles pointés, $g^{-1}$ admet également un adjoint à gauche $g_!$. Sa
construction est entièrement analogue à celle de $g_!$ sur les faisceaux
abéliens dans la section \ref{sites-modules-section-exactness-lower-shriek}.

\medskip\noindent
De même, si $j : \mathcal{C}/U \to \mathcal{C}$ est comme dans la
section \ref{sites-modules-section-localize}, alors il existe un adjoint à gauche $j_!$ du
foncteur $j^{-1}$ sur les faisceaux d'ensembles pointés.

\medskip\noindent
Si nous avons un jour besoin de ces faits et constructions, nous énoncerons et
démontrerons ici précisément les lemmes correspondants.










% OK, start here.
%

\chapter{Injectifs}



\phantomsection
\label{injectives-section-phantom}


\section{Introduction}
\label{injectives-section-introduction}

\noindent
Dans des chapitres ultérieurs, nous utiliserons l'existence d'injectifs et
de complexes K-injectifs pour étudier la cohomologie des faisceaux de modules
sur les sites annelés. Dans ce chapitre, nous expliquons comment construire
des injectifs et des complexes K-injectifs, d'abord pour les modules sur les
sites, puis, plus généralement, pour les catégories abéliennes de Grothendieck.

\medskip\noindent
Observons que nous savons déjà que la catégorie des groupes abéliens et la
catégorie des modules sur un anneau possèdent assez d'injectifs ; voir
Compléments d'algèbre, sections
\ref{more-algebra-section-abelian-groups} et
\ref{more-algebra-section-injectives-modules}





\section{L'argument de Baer pour les modules}
\label{injectives-section-baer}

\noindent
Il existe une autre manière, plus ensembliste, de montrer que tout $R$-module
$M$ se plonge dans un module injectif. Cette méthode construit le module
injectif comme une colimite transfinie de sommes amalgamées. Bien qu'elle soit
assez abstraite et plus compliquée que celle de
Compléments d'algèbre, section \ref{more-algebra-section-injectives-modules},
elle est aussi plus générale. Cette méthode semble remonter à Baer ; elle a
été reprise par Cartan et Eilenberg dans \cite{Cartan-Eilenberg}, puis par
Grothendieck dans \cite{Tohoku}. Grothendieck y montre que bien d'autres
catégories abéliennes possèdent assez d'injectifs. Nous reviendrons plus loin
au cas général (section \ref{injectives-section-grothendieck-categories}).

\medskip\noindent
Commençons par quelques remarques ensemblistes.
Soit $\{B_{\beta}\}_{\beta \in \alpha}$ un système inductif d'objets d'une
catégorie $\mathcal{C}$, indexé par un ordinal $\alpha$. Supposons que
$\colim_{\beta \in \alpha} B_\beta$ existe dans $\mathcal{C}$. Si $A$ est un
objet de $\mathcal{C}$, il existe alors une application naturelle
\begin{equation}
\label{injectives-equation-compare}
\colim_{\beta \in \alpha} \Mor_\mathcal{C}(A, B_\beta)
\longrightarrow
\Mor_\mathcal{C}(A, \colim_{\beta \in \alpha} B_\beta).
\end{equation}
En effet, un morphisme $A \to B_\beta$ donné pour un certain $\beta$ fournit
naturellement un morphisme de $A$ vers la colimite, par composition avec
$B_\beta \to \colim_{\beta \in \alpha} B_\alpha$.
Notons que la colimite de gauche est une colimite d'ensembles ! En général,
(\ref{injectives-equation-compare}) n'est ni injective ni surjective.

\begin{example}
\label{injectives-example-not-surjective}
Considérons la catégorie des ensembles. Posons $A = \mathbf{N}$ et prenons
pour $B_n = \{1, \ldots, n\}$ le système inductif indexé par les entiers
naturels, où $B_n \to B_m$, pour $n \leq m$, est l'application évidente.
Alors $\colim B_n = \mathbf{N}$ ; il existe donc une application
$A \to \colim B_n$ qui ne se factorise par aucune application $A \to B_m$,
pour tout $m$.
Par conséquent,
$\colim \Mor(A, B_n) \to \Mor(A, \colim B_n)$
n'est pas surjective.
\end{example}

\begin{example}
\label{injectives-example-not-injective}
Donnons ensuite un exemple où l'application n'est pas injective. Posons
$B_n = \mathbf{N}/\{1, 2, \ldots, n\}$, c'est-à-dire l'ensemble quotient de
$\mathbf{N}$ dans lequel les $n$ premiers éléments sont identifiés à un seul.
Il existe des applications naturelles $B_n \to B_m$ pour $n \leq m$ ; les
$\{B_n\}$ forment donc un système d'ensembles indexé par $\mathbf{N}$. On voit
aisément que $\colim B_n = \{*\}$ : c'est l'ensemble à un élément. Il s'ensuit
que $\Mor(A, \colim B_n)$ est un ensemble à un élément pour tout ensemble $A$.
Cependant, $\colim \Mor(A , B_n)$ n'est {\bf pas} un ensemble à un élément.
Considérons la famille d'applications $A \to B_n$ formée des projections
naturelles $\mathbf{N} \to \mathbf{N}/\{1, 2, \ldots, n\}$, et la famille
d'applications $A \to B_n$ qui envoient tout $A$ sur la classe de $1$.
Ces deux familles d'applications sont distinctes à chaque étape, donc distinctes
dans $\colim \Mor(A, B_n)$, mais elles induisent la même application
$A \to \colim B_n$.
\end{example}

\noindent
Néanmoins, si l'ensemble source est fini, (\ref{injectives-equation-compare}) est toujours
un isomorphisme.

\begin{lemma}
\label{injectives-lemma-out-of-finite}
Supposons que, dans (\ref{injectives-equation-compare}), $\mathcal{C}$ soit la catégorie
des ensembles et que $A$ soit un {\it ensemble fini}. Alors l'application est
une bijection.
\end{lemma}

\begin{proof}
Soit $f : A \to \colim B_\beta$. L'image de $f$ est finie ; notons ses éléments
$c_1, \ldots, c_r \in \colim B_\beta$. Par définition de la colimite, ils
proviennent tous d'éléments d'un même $B_\beta$, pour $\beta \in \alpha$
assez grand. On peut donc définir $\widetilde{f} : A \to B_\beta$ qui relève
$f$ à une étape finie. Cela prouve que (\ref{injectives-equation-compare}) est surjective.
Supposons maintenant que deux applications
$f : A \to B_\gamma, f' : A \to B_{\gamma'}$ définissent la même application
$A \to \colim B_\beta$. Chacun des éléments, en nombre fini, de $A$ a alors
la même image dans la colimite. Par définition de la colimite d'ensembles, il
existe $\beta \geq \gamma, \gamma'$ tel que les éléments de $A$ aient les mêmes images
dans $B_\beta$ par $f$ et $f'$. Cela prouve que (\ref{injectives-equation-compare}) est
injective.
\end{proof}

\noindent
Le cas le plus intéressant du lemme est celui où $\alpha = \omega$, autrement
dit où le système $\{B_\beta\}$ est un système
$\{B_n\}_{n \in \mathbf{N}}$ indexé par les entiers naturels, comme dans les
Exemples \ref{injectives-example-not-surjective} et \ref{injectives-example-not-injective}.
L'idée essentielle est que $A$ est « petit » relativement à la longue chaîne
de composées $B_1 \to B_2 \to \ldots$, de sorte que son morphisme doit se
factoriser par une étape finie. On trouvera une version plus générale de ce
lemme dans Ensembles, lemme \ref{sets-lemma-map-from-set-lifts}. Généralisons
maintenant ce résultat à la catégorie des modules.

\begin{definition}
\label{injectives-definition-small}
Soient $\mathcal{C}$ une catégorie, $I \subset \text{Arrows}(\mathcal{C})$
et $\alpha$ un ordinal. Un objet $A$ de $\mathcal{C}$ est dit
{\it $\alpha$-petit relativement à $I$} si, pour tout système
$\{B_\beta\}$ indexé par $\alpha$ dont les morphismes de transition
appartiennent à $I$, l'application (\ref{injectives-equation-compare}) est un
isomorphisme.
\end{definition}

\noindent
Dans la suite de cette section, nous nous restreindrons à la catégorie des
$R$-modules, où $R$ est un anneau commutatif fixé. Nous prendrons aussi pour
$I$ la famille des applications injectives, c'est-à-dire des
{\it monomorphismes} dans la catégorie des modules sur $R$. Dans ce cas, pour
tout système $\{B_\beta\}$ comme dans la définition, chacune des applications
$$
B_\beta \to \colim_{\beta \in \alpha} B_\beta
$$
est injective. Il s'ensuit que l'application (\ref{injectives-equation-compare}) est une
{\it injection}. On peut en fait considérer les $B_\beta$ comme des
sous-modules du module $B = \colim_{\beta \in \alpha} B_\beta$ ; on a alors
$B = \bigcup_{\beta \in \alpha} B_\beta$. Il n'y a pas abus de notation si
l'on identifie $B_\alpha$ à son image dans la colimite. Nous voulons maintenant
montrer que les modules sont toujours petits pour des ordinaux $\alpha$
« grands ».

\begin{proposition}
\label{injectives-proposition-modules-are-small}
Soient $R$ un anneau et $M$ un $R$-module. Soit $\kappa$ le cardinal de
l'ensemble des sous-modules de $M$. Si $\alpha$ est un ordinal dont la
cofinalité est strictement supérieure à $\kappa$, alors $M$ est
$\alpha$-petit relativement aux injections.
\end{proposition}

\begin{proof}
La démonstration est immédiate, mais considérons d'abord un cas particulier.
Si $M$ est fini, l'assertion est que, pour tout système inductif
$\{B_\beta\}$ à morphismes de transition injectifs et indexé par un ordinal
limite, toute application $M \to \colim B_\beta$ se factorise par l'un des
$B_\beta$. C'est précisément ce que nous avons démontré dans le
Lemme \ref{injectives-lemma-out-of-finite}.

\medskip\noindent
Passons au cas général. Il suffit de montrer que l'application
(\ref{injectives-equation-compare}) est surjective. Soit
$f : M \to \colim B_\beta$ une application. Considérons les sous-objets
$\{f^{-1}(B_\beta)\}$ de $M$, où $B_\beta$ est vu comme un sous-objet de la
colimite $B = \bigcup_\beta B_\beta$. Si l'un d'eux, disons
$f^{-1}(B_\beta)$, est égal à $M$, alors l'application se factorise par
$B_\beta$.

\medskip\noindent
Supposons donc, par l'absurde, que tous les $f^{-1}(B_\beta)$ soient des
sous-objets propres de $M$. Or nous savons que
$$
\bigcup f^{-1}(B_\beta) = f^{-1}\left(\bigcup B_\beta\right) = M.
$$
Par hypothèse, il y a au plus $\kappa$ sous-objets distincts de $M$ parmi les
$f^{-1}(B_\alpha)$. On peut donc trouver un sous-ensemble $S \subset \alpha$
de cardinal au plus $\kappa$ tel que, lorsque $\beta'$ parcourt $S$, les
$f^{-1}(B_{\beta'})$ parcourent \emph{tous} les $f^{-1}(B_\alpha)$.

\medskip\noindent
Cependant, $S$ possède un majorant $\widetilde{\alpha} < \alpha$, puisque
$\alpha$ est de cofinalité strictement supérieure à $\kappa$. En particulier,
tous les $f^{-1}(B_{\beta'})$, pour $\beta' \in S$, sont contenus dans
$f^{-1}(B_{\widetilde{\alpha}})$. Il en résulte que
$f^{-1}(B_{\widetilde{\alpha}}) = M$. Ainsi, l'application $f$ se factorise
par $B_{\widetilde{\alpha}}$.
\end{proof}

\noindent
Ce lemme nous permettra de déduire l'existence de nombreux injectifs.
Rappelons le critère de Baer.

\begin{lemma}[Critère de Baer]
\label{injectives-lemma-criterion-baer}
\begin{reference}
\cite[Theorem 1]{Baer}
\end{reference}
Soit $R$ un anneau. Un $R$-module $Q$ est injectif si et seulement si, dans
tout diagramme commutatif
$$
\xymatrix{
\mathfrak{a} \ar[d] \ar[r] &  Q \\
R \ar@{-->}[ru]
}
$$
où $\mathfrak{a} \subset R$ est un idéal, la flèche pointillée existe.
\end{lemma}

\begin{proof}
Il s'agit de l'équivalence entre (1) et (3) dans
Compléments d'algèbre, lemme
\ref{more-algebra-lemma-characterize-injective-bis};
notons que la démonstration qui y est donnée est élémentaire (elle n'utilise ni
les groupes $\text{Ext}$, ni l'existence d'injectifs ou de projectifs dans la
catégorie des $R$-modules).
\end{proof}

\noindent
Si $M$ est un $R$-module, nous pouvons en général disposer d'un diagramme
partiellement complété comme dans le Lemme \ref{injectives-lemma-criterion-baer}. Nous
pouvons y former la \emph{somme amalgamée}
$$
\xymatrix{
\mathfrak{a} \ar[d]  \ar[r] &  M \ar[d] \\
R \ar[r] &  R \oplus_{\mathfrak{a}} M.
}
$$
Ici, le morphisme vertical est injectif et le diagramme est commutatif. Le point
essentiel est que nous pouvons prolonger $\mathfrak{a} \to M$ à $R$ \emph{si}
nous remplaçons $M$ par le module plus grand
$R \oplus_{\mathfrak{a}} M$.

\medskip\noindent
L'idée fondamentale de l'argument de Baer consiste à répéter cette procédure
un nombre transfini de fois. Pour cela, étant donné un $R$-module $M$, nous
commençons par définir l'immense somme amalgamée suivante
\begin{equation}
\label{injectives-equation-huge-diagram}
\vcenter{
\xymatrix{
\bigoplus_{\mathfrak a}
\bigoplus_{\varphi \in \Hom_R(\mathfrak a, M)}
\mathfrak{a} \ar[r] \ar[d] & M \ar[d] \\
\bigoplus_{\mathfrak a}
\bigoplus_{\varphi \in \Hom_R(\mathfrak a, M)}
R \ar[r] &  \mathbf{M}(M).
}
}
\end{equation}
La flèche horizontale supérieure envoie l'élément $a \in \mathfrak a$ du
facteur correspondant à $\varphi$ sur l'élément $\varphi(a) \in M$. La flèche
verticale de gauche envoie simplement $a \in \mathfrak a$, dans le facteur
correspondant à $\varphi$, sur l'élément $a \in R$ du facteur correspondant à
$\varphi$. Les propriétés fondamentales de cette construction sont énoncées
dans le lemme suivant.

\begin{lemma}
\label{injectives-lemma-construction}
Soit $R$ un anneau.
\begin{enumerate}
\item La construction $M \mapsto (M \to \mathbf{M}(M))$ est fonctorielle
en $M$.
\item L'application $M \to \mathbf{M}(M)$ est injective.
\item Pour tout idéal $\mathfrak{a}$ et tout morphisme de $R$-modules
$\varphi : \mathfrak a \to M$, il existe un morphisme de $R$-modules
$\varphi' : R \to \mathbf{M}(M)$ tel que le diagramme
$$
\xymatrix{
\mathfrak{a} \ar[d] \ar[r]_\varphi &  M \ar[d] \\
R \ar[r]^{\varphi'} & \mathbf{M}(M)
}
$$
soit commutatif.
\end{enumerate}
\end{lemma}

\begin{proof}
Les assertions (2) et (3) résultent immédiatement de la construction. Pour
voir (1), soit $\chi : M \to N$ un morphisme de $R$-modules. Nous affirmons
qu'il existe un diagramme commutatif canonique
$$
\xymatrix{
\bigoplus_{\mathfrak a}
\bigoplus_{\varphi \in \Hom_R(\mathfrak a, M)}
\mathfrak{a} \ar[r] \ar[d] \ar[rrd] & M \ar[rrd]^\chi \\
\bigoplus_{\mathfrak a}
\bigoplus_{\varphi \in \Hom_R(\mathfrak a, M)}
R \ar[rrd] & &
\bigoplus_{\mathfrak a}
\bigoplus_{\psi \in \Hom_R(\mathfrak a, N)}
\mathfrak{a} \ar[r] \ar[d] & N \\
& & \bigoplus_{\mathfrak a}
\bigoplus_{\psi \in \Hom_R(\mathfrak a, N)}
R
}
$$
qui induit l'application voulue $\mathbf{M}(M) \to \mathbf{M}(N)$. La flèche
centrale dirigée vers l'est-sud-est envoie le facteur $\mathfrak a$
correspondant à $\varphi$, par $\text{id}_{\mathfrak a}$, sur le facteur
$\mathfrak a$ correspondant à
$\psi = \chi \circ \varphi$. Il en va de même de la flèche inférieure dirigée
vers l'est-sud-est. Nous omettons les détails.
\end{proof}

\noindent
Nous allons maintenant appliquer le foncteur $\mathbf{M}$ un nombre transfini
de fois. Pour chaque ordinal $\alpha$, nous définissons un foncteur
$\mathbf{M}_\alpha$ sur la catégorie des $R$-modules, ainsi qu'une injection
naturelle $N \to \mathbf{M}_\alpha(N)$, par récurrence transfinie. D'abord,
$\mathbf{M}_1 = \mathbf{M}$ est le foncteur défini ci-dessus. Soit ensuite un
ordinal $\alpha$, et supposons $\mathbf{M}_{\alpha'}$ défini pour
$\alpha' < \alpha$. Si $\alpha$ admet un prédécesseur immédiat
$\widetilde{\alpha}$, posons
$$
\mathbf{M}_\alpha = \mathbf{M} \circ \mathbf{M}_{\widetilde{\alpha}}.
$$
Sinon, c'est-à-dire si $\alpha$ est un ordinal limite, posons
$$
\mathbf{M}_{\alpha}(N) =
\colim_{\alpha' < \alpha} \mathbf{M}_{\alpha'}(N).
$$
Il est clair, par exemple par récurrence, que les
$\mathbf{M}_{\alpha}(N)$ forment un système inductif indexé par les ordinaux ;
cette définition a donc bien un sens.

\begin{theorem}
\label{injectives-theorem-baer-grothendieck}
Soit $\kappa$ le cardinal de l'ensemble des idéaux de $R$, et soit $\alpha$
un ordinal de cofinalité strictement supérieure à $\kappa$. Alors
$\mathbf{M}_\alpha(N)$ est un $R$-module injectif, et
$N \to \mathbf{M}_\alpha(N)$ est un plongement fonctoriel dans un injectif.
\end{theorem}

\begin{proof}
D'après le critère de Baer, Lemme \ref{injectives-lemma-criterion-baer}, il suffit de
montrer que, pour tout idéal $\mathfrak{a} \subset R$, toute application
$f : \mathfrak{a} \to \mathbf{M}_\alpha(N)$ se prolonge en une application
$R \to \mathbf{M}_\alpha(N)$. Or, puisque $\alpha$ est un ordinal limite, nous
savons que
$$
\mathbf{M}_{\alpha}(N) =
\colim_{\beta < \alpha} \mathbf{M}_{\beta}(N),
$$
donc, par la Proposition \ref{injectives-proposition-modules-are-small}, nous obtenons
$$
\Hom_R(\mathfrak{a}, \mathbf{M}_{\alpha}(N)) =
\colim_{\beta < \alpha} \Hom_R(\mathfrak a, \mathbf{M}_{\beta}(N)).
$$
Cela signifie en particulier qu'il existe $\beta' < \alpha$ tel que $f$ se
factorise par le sous-module $\mathbf{M}_{\beta'}(N)$, sous la forme
$$
f : \mathfrak{a} \to \mathbf{M}_{\beta'}(N) \to
\mathbf{M}_{\alpha}(N).
$$
Cependant, d'après la propriété fondamentale du foncteur $\mathbf{M}$, voir le
Lemme \ref{injectives-lemma-construction}, assertion (3), l'application
$\mathfrak{a} \to \mathbf{M}_{\beta'}(N)$ se prolonge en
$$
R \to \mathbf{M}( \mathbf{M}_{\beta'}(N)) =
\mathbf{M}_{\beta' + 1}(N),
$$
et ce dernier objet se plonge dans $\mathbf{M}_{\alpha}(N)$, car
$\beta' + 1 < \alpha$ puisque $\alpha$ est un ordinal limite. Par conséquent,
$f$ se prolonge à $\mathbf{M}_{\alpha}(N)$.
\end{proof}




\section{G-modules}
\label{injectives-section-G-modules}

\noindent
Nous verrons plus loin (Algèbre différentielle graduée, section
\ref{dga-section-modules-noncommutative}) que la catégorie des modules sur une
algèbre admet des plongements fonctoriels dans des injectifs. La construction est
exactement la même que celle de
Compléments d'algèbre, section
\ref{more-algebra-section-injectives-modules}.

\begin{lemma}
\label{injectives-lemma-G-modules}
Soient $G$ un groupe topologique et $R$ un anneau. La catégorie
$\text{Mod}_{R, G}$ des $R\text{-}G$-modules, voir Cohomologie étale,
définition \ref{etale-cohomology-definition-G-module-continuous}, admet des
plongements fonctoriels dans des injectifs. Il en est ainsi, en particulier, de la
catégorie des $G$-modules discrets.
\end{lemma}

\begin{proof}
La remarque qui précède le lemme montre que la catégorie
$\text{Mod}_{R[G]}$ admet des plongements fonctoriels dans des injectifs. Considérons le
foncteur d'oubli
$v : \text{Mod}_{R, G} \to \text{Mod}_{R[G]}$.
Ce foncteur est pleinement fidèle, transforme les applications injectives en
applications injectives et possède un adjoint à droite, à savoir
$$
u : M \mapsto u(M) = \{x \in M \mid \text{le stabilisateur de }x\text{ est ouvert}\}
$$
Comme $v(M) = 0 \Rightarrow M = 0$, on conclut par
Homologie, lemme \ref{homology-lemma-adjoint-functorial-injectives}.
\end{proof}



\section{Faisceaux abéliens sur un espace}
\label{injectives-section-abelian-sheaves-space}


\begin{lemma}
\label{injectives-lemma-abelian-sheaves-space}
Soit $X$ un espace topologique. La catégorie des faisceaux abéliens sur $X$
possède assez d'injectifs. Elle admet même des plongements fonctoriels dans des
injectifs.
\end{lemma}

\begin{proof}
Pour tout groupe abélien $A$, notons $j : A \to J(A)$ le plongement fonctoriel
dans un injectif construit dans
Compléments d'algèbre, section
\ref{more-algebra-section-injectives-modules}.
Soit $\mathcal{F}$ un faisceau abélien sur $X$. D'après Faisceaux, exemple
\ref{sheaves-example-sheaf-product-pointwise}, la donnée
$$
\mathcal{I} : U \mapsto
\mathcal{I}(U) = \prod\nolimits_{x\in U} J(\mathcal{F}_x)
$$
est un faisceau abélien. Il existe une application canonique
$\mathcal{F} \to \mathcal{I}$ qui envoie $s \in \mathcal{F}(U)$ sur
$\prod_{x \in U} j(s_x)$, où $s_x \in \mathcal{F}_x$ désigne le germe de $s$
en $x$. Cette application est injective ; voir, par exemple, Faisceaux, lemme
\ref{sheaves-lemma-sheaf-subset-stalks}.

\medskip\noindent
Il reste à démontrer l'assertion suivante : si une donnée $x \mapsto I_x$
associe à chaque point $x \in X$ un groupe abélien injectif, alors le faisceau
$\mathcal{I} : U \mapsto \prod_{x \in U} I_x$ est injectif. Notons que
$$
\mathcal{I} = \prod\nolimits_{x \in X} i_{x, *}I_x
$$
est le produit des faisceaux gratte-ciel $i_{x, *}I_x$ (pour la notation, voir
Faisceaux, section \ref{sheaves-section-skyscraper-sheaves}). On a
$$
\Mor_{\textit{Ab}}(\mathcal{F}_x, I_x)
=
\Mor_{\textit{Ab}(X)}(\mathcal{F}, i_{x, *}I_x).
$$
voir Faisceaux, lemme \ref{sheaves-lemma-stalk-skyscraper-adjoint}. Il est donc
clair que chaque $i_{x, *}I_x$ est injectif. L'injectivité de $\mathcal{I}$
résulte alors de
Homologie, lemme \ref{homology-lemma-product-injectives}.
\end{proof}









\section{Faisceaux de modules sur un espace annelé}
\label{injectives-section-sheaves-modules-space}


\begin{lemma}
\label{injectives-lemma-sheaves-modules-space}
Soit $(X, \mathcal{O}_X)$ un espace annelé ; voir Faisceaux, section
\ref{sheaves-section-ringed-spaces}. La catégorie des faisceaux de
$\mathcal{O}_X$-modules sur $X$ possède assez d'injectifs. Elle admet même des
plongements fonctoriels dans des injectifs.
\end{lemma}

\begin{proof}
Pour tout anneau $R$ et tout $R$-module $M$, notons
$j : M \to J_R(M)$ le plongement fonctoriel dans un injectif construit dans
Compléments d'algèbre, section
\ref{more-algebra-section-injectives-modules}.
Soit $\mathcal{F}$ un faisceau de $\mathcal{O}_X$-modules sur $X$. D'après
Faisceaux, exemples \ref{sheaves-example-sheaf-product-pointwise} et
\ref{sheaves-example-sheaf-product-pointwise-algebraic-structure}, la donnée
$$
\mathcal{I} : U \mapsto
\mathcal{I}(U) = \prod\nolimits_{x\in U} J_{\mathcal{O}_{X, x}}(\mathcal{F}_x)
$$
est un faisceau abélien. Il existe une application canonique
$\mathcal{F} \to \mathcal{I}$ qui envoie $s \in \mathcal{F}(U)$ sur
$\prod_{x \in U} j(s_x)$, où $s_x \in \mathcal{F}_x$ désigne le germe de $s$
en $x$. Cette application est injective ; voir, par exemple, Faisceaux, lemme
\ref{sheaves-lemma-sheaf-subset-stalks}.

\medskip\noindent
Il reste à démontrer l'assertion suivante : si une donnée $x \mapsto I_x$
associe à chaque point $x \in X$ un $\mathcal{O}_{X, x}$-module injectif, alors
le faisceau $\mathcal{I} : U \mapsto \prod_{x \in U} I_x$ est injectif.
Notons que
$$
\mathcal{I} = \prod\nolimits_{x \in X} i_{x, *}I_x
$$
est le produit des faisceaux gratte-ciel $i_{x, *}I_x$ (pour la notation, voir
Faisceaux, section \ref{sheaves-section-skyscraper-sheaves}). On a
$$
\Hom_{\mathcal{O}_{X, x}}(\mathcal{F}_x, I_x)
=
\Hom_{\mathcal{O}_X}(\mathcal{F}, i_{x, *}I_x).
$$
voir Faisceaux, lemme \ref{sheaves-lemma-stalk-skyscraper-adjoint}. Il est donc
clair que chaque $i_{x, *}I_x$ est un $\mathcal{O}_X$-module injectif (voir
Homologie, lemme \ref{homology-lemma-adjoint-preserve-injectives}, ou raisonner
directement). L'injectivité de $\mathcal{I}$ résulte alors de
Homologie, lemme \ref{homology-lemma-product-injectives}.
\end{proof}













\section{Préfaisceaux abéliens sur une catégorie}
\label{injectives-section-injectives-presheaves}

\noindent
Soit $\mathcal{C}$ une catégorie. Rappelons que cela signifie que
$\Ob(\mathcal{C})$ est un ensemble. Considérons, d'une part, les préfaisceaux
abéliens sur $\mathcal{C}$, voir Sites, section
\ref{sites-section-presheaves}, et, d'autre part, les familles de groupes
abéliens indexées par les éléments de $\Ob(\mathcal{C})$ ; autrement dit, les
préfaisceaux sur la catégorie discrète dont l'ensemble sous-jacent d'objets est
$\Ob(\mathcal{C})$. Notons simplement $\Ob(\mathcal{C})$ cette catégorie
discrète. Il existe un foncteur naturel
$$
i : \Ob(\mathcal{C}) \longrightarrow \mathcal{C}
$$
et, par conséquent, un foncteur naturel de restriction, ou foncteur d'oubli,
$$
v = i^p :
\textit{PAb}(\mathcal{C})
\longrightarrow
\textit{PAb}(\Ob(\mathcal{C}))
$$
à comparer avec Sites, section \ref{sites-section-functoriality-PSh}. Nous
noterons les préfaisceaux sur $\mathcal{C}$ par $B$ et les préfaisceaux sur
$\Ob(\mathcal{C})$ par $A$.

\medskip\noindent
Il existe aussi deux foncteurs, à savoir $i_p$ et ${}_pi$, qui associent un
préfaisceau abélien sur $\mathcal{C}$ à un préfaisceau abélien sur
$\Ob(\mathcal{C})$ ; voir Sites, sections
\ref{sites-section-functoriality-PSh} et
\ref{sites-section-more-functoriality-PSh}. Nous utiliserons ici
$u = {}_pi$, qui, dans le cas présent, est défini par
$$
uA(U) = \prod\nolimits_{U' \to U} A(U').
$$
Ainsi, un élément est une famille $(a_\phi)_\phi$, où $\phi$ parcourt tous les
morphismes de $\mathcal{C}$ de but $U$. L'application de restriction de $uA$
correspondant à $g : V \to U$ envoie l'élément $(a_\phi)_\phi$ sur
$(a_{g \circ \psi})_\psi$.

\medskip\noindent
Il existe une application surjective canonique $vuA \to A$ et une application
injective canonique $B \to uvB$. Nous laissons au lecteur le soin de montrer
que
$$
\Mor_{\textit{PAb}(\mathcal{C})}(B, uA)
=
\Mor_{\textit{PAb}(\Ob(\mathcal{C}))}(vB, A).
$$
dans ce cas simple ; le cas général est traité dans Sites, section
\ref{sites-section-functoriality-PSh}. Ainsi, le couple $(u, v)$ est un exemple
de couple de foncteurs adjoints ; voir Catégories, section
\ref{categories-section-adjoint}.

\medskip\noindent
Nous pouvons maintenant énumérer les faits suivants concernant la situation
ci-dessus.
\begin{enumerate}
\item Les foncteurs $u$ et $v$ sont exacts. Cela résulte de leur description
explicite donnée ci-dessus.
\item En particulier, le foncteur $v$ transforme les applications injectives
en applications injectives.
\item La catégorie $\textit{PAb}(\Ob(\mathcal{C}))$ possède assez d'injectifs.
\item Il existe même un plongement fonctoriel dans un injectif
$A \mapsto \big(A \to J(A)\big)$ comme dans
Homologie, définition \ref{homology-definition-functorial-injective-embedding}.
On peut en effet prendre pour $J(A)$ le préfaisceau
$U\mapsto J(A(U))$, où $J(-)$ est le foncteur construit dans
Compléments d'algèbre, section
\ref{more-algebra-section-injectives-modules}
pour l'anneau $\mathbf{Z}$.
\end{enumerate}
En réunissant ces faits, on obtient la procédure suivante pour plonger les
objets $B$ de $\textit{PAb}(\mathcal{C})$ dans un objet injectif :
$B \to uJ(vB)$. Voir
Homologie, lemme \ref{homology-lemma-adjoint-functorial-injectives}.

\begin{proposition}
\label{injectives-proposition-presheaves-injectives}
Les préfaisceaux abéliens sur une catégorie admettent un plongement fonctoriel
dans un injectif.
\end{proposition}

\begin{proof}
Voir la discussion ci-dessus.
\end{proof}

\section{Faisceaux abéliens sur un site}
\label{injectives-section-injectives-sheaves}

\noindent
Soit $\mathcal{C}$ un site. Dans cette section, nous démontrons que la
catégorie des faisceaux abéliens sur $\mathcal{C}$ a assez d'injectifs.

\medskip\noindent
Notons
$i : \textit{Ab}(\mathcal{C}) \longrightarrow \textit{PAb}(\mathcal{C})$
le foncteur d'oubli des faisceaux abéliens vers les préfaisceaux abéliens.
Notons
${}^\# : \textit{PAb}(\mathcal{C}) \longrightarrow \textit{Ab}(\mathcal{C})$
le foncteur de faisceautisation. Rappelons que ${}^\#$ est adjoint à gauche
de $i$, que ${}^\#$ est exact et que $i\mathcal{F}^\# = \mathcal{F}$
pour tout faisceau abélien $\mathcal{F}$. Enfin, notons
$\mathcal{G} \to J(\mathcal{G})$ le plongement canonique dans un
préfaisceau injectif construit à la
section \ref{injectives-section-injectives-presheaves}.

\medskip\noindent
Pour tout faisceau $\mathcal{F}$ de $\textit{Ab}(\mathcal{C})$ et tout
ordinal $\beta$, nous définissons un faisceau $J_\beta(\mathcal{F})$ par
récurrence transfinie. Posons $J_0(\mathcal{F}) = \mathcal{F}$ et
$J_1(\mathcal{F}) = J(i\mathcal{F})^\#$. La faisceautisation de
l'application canonique $i\mathcal{F} \to J(i\mathcal{F})$ donne une
application fonctorielle
$$
\mathcal{F} \longrightarrow J_1(\mathcal{F})
$$
qui est injective puisque $\#$ est exact. Posons
$J_{\alpha + 1}(\mathcal{F}) = J_1(J_\alpha(\mathcal{F}))$. On obtient ainsi
des applications injectives canoniques
$J_\alpha(\mathcal{F}) \to J_{\alpha + 1}(\mathcal{F})$. Pour un ordinal
limite $\beta$, nous posons
$$
J_\beta(\mathcal{F}) = \colim_{\alpha < \beta} J_\alpha(\mathcal{F}).
$$
Remarquons qu'il s'agit d'une colimite filtrante. Ainsi, pour tous ordinaux
$\alpha < \beta$, on dispose d'une application injective
$J_\alpha(\mathcal{F}) \to J_\beta(\mathcal{F})$.

\begin{lemma}
\label{injectives-lemma-map-into-next-one}
Avec les notations ci-dessus, supposons que
$\mathcal{G}_1 \to \mathcal{G}_2$ soit une application injective de faisceaux
abéliens sur $\mathcal{C}$. Soit $\alpha$ un ordinal et soit
$\mathcal{G}_1 \to J_\alpha(\mathcal{F})$ un morphisme de faisceaux. Il existe
un morphisme $\mathcal{G}_2 \to J_{\alpha + 1}(\mathcal{F})$ tel que le
diagramme suivant soit commutatif
$$
\xymatrix{
\mathcal{G}_1 \ar[d] \ar[r] & \mathcal{G}_2 \ar[d] \\
J_{\alpha}(\mathcal{F}) \ar[r] & J_{\alpha + 1}(\mathcal{F}) }
$$
\end{lemma}

\begin{proof}
En effet, l'application $i\mathcal{G}_1 \to i\mathcal{G}_2$ est injective ;
par conséquent, $i\mathcal{G}_1 \to iJ_\alpha(\mathcal{F})$ se prolonge en
$i\mathcal{G}_2 \to J(iJ_\alpha(\mathcal{F}))$, ce qui donne l'application
voulue après application du foncteur de faisceautisation.
\end{proof}

\noindent
Ce lemme affirme qu'en un certain sens le système
$\{J_{\alpha}(\mathcal{F})\}$ est un plongement injectif de $\mathcal{F}$.
Bien entendu, on ne peut pas prendre la limite sur tous les $\alpha$, car
ceux-ci forment une classe et non un ensemble. L'idée est cependant qu'il
n'est pas nécessaire de vérifier l'injectivité sur toutes les injections
$\mathcal{G}_1 \to \mathcal{G}_2$, compte tenu du lemme suivant.

\begin{lemma}
\label{injectives-lemma-map-into-smaller}
Supposons que les $\mathcal{G}_i$, $i\in I$, forment un ensemble de faisceaux
abéliens sur $\mathcal{C}$. Il existe un ordinal $\beta$ tel que, pour tout
faisceau $\mathcal{F}$, tout $i\in I$ et toute application
$\varphi : \mathcal{G}_i \to J_\beta(\mathcal{F})$, il existe un
$\alpha < \beta$ tel que $ \varphi $ se factorise par
$J_\alpha(\mathcal{F})$.
\end{lemma}

\begin{proof}
On se ramène au cas d'un seul faisceau $\mathcal{G}$ en prenant la somme
directe de tous les $\mathcal{G}_i$.

\medskip\noindent
Considérons les ensembles
$$
S = \coprod\nolimits_{U \in \Ob(\mathcal{C})} \mathcal{G}(U).
$$
et
$$
T_\beta
=
\coprod\nolimits_{U \in \Ob(\mathcal{C})} J_\beta(\mathcal{F})(U)
$$
Les applications de transition entre les ensembles $T_\beta$ sont injectives.
Si la cofinalité de $\beta$ est assez grande, alors
$T_\beta = \colim_{\alpha < \beta} T_\alpha$ ; voir
Sites, lemme \ref{sites-lemma-colimit-over-ordinal-sections}. Un morphisme
$\mathcal{G} \to J_\beta(\mathcal{F})$ se factorise par
$J_\alpha(\mathcal{F})$ si et seulement si l'application associée
$S \to T_\beta$ se factorise par $T_\alpha$. D'après
Ensembles, lemme \ref{sets-lemma-map-from-set-lifts}, si la cofinalité de $\beta$
est strictement supérieure au cardinal de $S$, la conclusion du lemme est
vraie. Le lemme résulte donc du fait qu'il existe des ordinaux de cofinalité
arbitrairement grande ; voir
Ensembles, proposition \ref{sets-proposition-exist-ordinals-large-cofinality}.
\end{proof}

\noindent
Rappelons que, pour un objet $X$ de $\mathcal{C}$, on note $\mathbf{Z}_X$ le
préfaisceau de groupes abéliens $\Gamma(U, \mathbf{Z}_X) =
\oplus_{U \to X} \mathbf{Z}$ ; voir
Modules sur les sites, section \ref{sites-modules-section-free-abelian-presheaf}.
Le faisceau associé à ce préfaisceau est noté $\mathbf{Z}_X^\#$ ; voir
Modules sur les sites, section \ref{sites-modules-section-free-abelian-sheaf}.
Il se caractérise par la propriété
\begin{equation}
\label{injectives-equation-free-sheaf-on}
\Mor_{\textit{Ab}(\mathcal{C})}(\mathbf{Z}_X^\#, \mathcal{G})
=
\mathcal{G}(X)
\end{equation}
où l'élément $\varphi$ du membre de gauche est envoyé sur
$\varphi(1 \cdot \text{id}_X)$ dans celui de droite. Ces faisceaux permettent
de caractériser les faisceaux abéliens injectifs.

\begin{lemma}
\label{injectives-lemma-characterize-injectives}
Soit $\mathcal{J}$ un faisceau de groupes abéliens ayant la propriété
suivante : pour tout $X\in \Ob(\mathcal{C})$, tout sous-faisceau abélien
$\mathcal{S} \subset \mathbf{Z}_X^\#$ et tout morphisme
$\varphi : \mathcal{S} \to \mathcal{J}$, il existe un morphisme
$\mathbf{Z}_X^\# \to \mathcal{J}$ prolongeant $\varphi$. Alors $\mathcal{J}$
est un faisceau injectif de groupes abéliens.
\end{lemma}

\begin{proof}
Soit $\mathcal{F} \to \mathcal{G}$ une application injective de faisceaux
abéliens et soit $\varphi : \mathcal{F} \to \mathcal{J}$ un morphisme.
En raisonnant comme dans la preuve de
Compléments d'algèbre, lemme \ref{more-algebra-lemma-injective-abelian}, on voit qu'il
suffit de démontrer que, si $\mathcal{F} \not = \mathcal{G}$, on peut trouver
un faisceau abélien $\mathcal{F}'$ tel que
$\mathcal{F} \subset \mathcal{F}' \subset \mathcal{G}$, l'inclusion
$\mathcal{F} \subset \mathcal{F}'$ étant stricte, et tel que $\varphi$ se
prolonge à $\mathcal{F}'$. Pour trouver $\mathcal{F}'$, choisissons un objet
$X$ de $\mathcal{C}$ pour lequel l'inclusion
$\mathcal{F}(X) \subset \mathcal{G}(X)$ est stricte, puis
$s \in \mathcal{G}(X)$, $s \not \in \mathcal{F}(X)$. Soit
$\psi : \mathbf{Z}_X^\# \to \mathcal{G}$ le morphisme correspondant à la
section $s$ par (\ref{injectives-equation-free-sheaf-on}). Posons
$\mathcal{S} = \psi^{-1}(\mathcal{F})$. Par hypothèse, le morphisme
$$
\mathcal{S} \xrightarrow{\psi} \mathcal{F} \xrightarrow{\varphi} \mathcal{J}
$$
se prolonge en un morphisme $\varphi' : \mathbf{Z}_X^\# \to \mathcal{J}$.
Remarquons que $\varphi'$ annule le noyau de $\psi$ (puisque c'est le cas de
$\varphi$). Ainsi, $\varphi'$ induit un morphisme
$\varphi'' : \Im(\psi) \to \mathcal{J}$ qui coïncide par construction avec
$\varphi$ sur l'intersection $\mathcal{F} \cap \Im(\psi)$. Les morphismes
$\varphi$ et $\varphi''$ se recollent donc en un prolongement de $\varphi$ au
sous-faisceau strictement plus grand
$\mathcal{F}' = \mathcal{F} + \Im(\psi)$.
\end{proof}

\begin{theorem}
\label{injectives-theorem-sheaves-injectives}
La catégorie des faisceaux de groupes abéliens sur un site a assez
d'injectifs. Plus précisément, il existe un plongement injectif fonctoriel ;
voir Homology, définition
\ref{homology-definition-functorial-injective-embedding}.
\end{theorem}

\begin{proof}
Soit $\mathcal{G}_i$, $i \in I$, un ensemble de faisceaux abéliens tel que
tout sous-faisceau de tout $\mathbf{Z}_X^\#$ soit l'un des $\mathcal{G}_i$.
Appliquons le lemme \ref{injectives-lemma-map-into-smaller} à cette famille et obtenons
ainsi un ordinal $\beta$. Nous affirmons que, pour tout faisceau de groupes
abéliens $\mathcal{F}$, l'application
$\mathcal{F} \to J_\beta(\mathcal{F})$ plonge $\mathcal{F}$ dans un injectif.
Remarquons que, par construction, l'association
$\mathcal{F} \mapsto \big(\mathcal{F} \to J_\beta(\mathcal{F})\big)$ est bien
fonctorielle.

\medskip\noindent
D'après le lemme \ref{injectives-lemma-characterize-injectives}, il suffit, pour prouver
cette affirmation, de prolonger tout morphisme
$\gamma : \mathcal{G} \to J_\beta(\mathcal{F})$ défini sur un sous-faisceau
$\mathcal{G}$ d'un $\mathbf{Z}_X^\#$ à tout $\mathbf{Z}_X^\#$. Or, d'après le
lemme \ref{injectives-lemma-map-into-smaller}, l'application $\gamma$ se relève dans
$J_\alpha(\mathcal{F})$ pour un certain $\alpha < \beta$. Enfin, le
lemme \ref{injectives-lemma-map-into-next-one} fournit le prolongement voulu de $\gamma$
en un morphisme à valeurs dans
$J_{\alpha + 1}(\mathcal{F}) \to J_\beta(\mathcal{F})$.
\end{proof}






\section{Modules sur un site annelé}
\label{injectives-section-sheaves-modules}

\noindent
Soit $\mathcal{C}$ un site et soit $\mathcal{O}$ un faisceau d'anneaux sur
$\mathcal{C}$. Par analogie avec
Compléments d'algèbre, section \ref{more-algebra-section-injectives-modules},
cherchons à démontrer qu'il existe assez de $\mathcal{O}$-modules injectifs.
Choisissons tout d'abord un plongement injectif
$$
\bigoplus\nolimits_{U, \mathcal{I}}
j_{U!}\mathcal{O}_U/\mathcal{I}
\longrightarrow
\mathcal{J}
$$
où $\mathcal{J}$ est un faisceau abélien injectif, dont l'existence résulte de
la section précédente. La somme directe porte ici sur tous les objets $U$ de
$\mathcal{C}$ et tous les sous-$\mathcal{O}$-modules
$\mathcal{I} \subset j_{U!}\mathcal{O}_U$. Pour les foncteurs de restriction
et de prolongement par $0$ associés au foncteur de localisation
$j_U : \mathcal{C}/U \to \mathcal{C}$, voir
Modules sur les sites, section \ref{sites-modules-section-localize}.

\medskip\noindent
Pour tout faisceau de $\mathcal{O}$-modules $\mathcal{F}$, notons
$$
\mathcal{F}^\vee
=
\SheafHom(\mathcal{F}, \mathcal{J})
$$
muni de sa structure naturelle de $\mathcal{O}$-module.
Insérer ici une future référence au Hom interne.
Nous aurons également besoin d'une résolution plate canonique d'un faisceau de
$\mathcal{O}$-modules. Nous la construisons comme suit : pour tout
$\mathcal{O}$-module $\mathcal{F}$, notons
$$
F(\mathcal{F})
=
\bigoplus\nolimits_{U \in \Ob(\mathcal{C}), s \in \mathcal{F}(U)}
j_{U!}\mathcal{O}_U.
$$
C'est un faisceau plat de $\mathcal{O}$-modules, muni d'une surjection canonique
$F(\mathcal{F}) \to \mathcal{F}$ ; voir
Modules sur les sites, lemme \ref{sites-modules-lemma-module-quotient-flat}.
En outre, la construction $\mathcal{F} \mapsto F(\mathcal{F})$ est
fonctorielle en $\mathcal{F}$.

\begin{lemma}
\label{injectives-lemma-vee-exact-sheaves}
Le foncteur $\mathcal{F} \mapsto \mathcal{F}^\vee$ est exact.
\end{lemma}

\begin{proof}
Cela résulte de ce que $\mathcal{J}$ est un faisceau abélien injectif.
\end{proof}

\noindent
Il existe une application canonique d'évaluation
$ev : \mathcal{F} \to (\mathcal{F}^\vee)^\vee$ : étant donné
$x \in \mathcal{F}(U)$, on définit
$ev(x) \in (\mathcal{F}^\vee)^\vee =
\SheafHom(\mathcal{F}^\vee, \mathcal{J})$ comme l'application
$\varphi \mapsto \varphi(x)$.

\begin{lemma}
\label{injectives-lemma-ev-injective-sheaves}
Pour tout $\mathcal{O}$-module $\mathcal{F}$, l'application d'évaluation
$ev : \mathcal{F} \to (\mathcal{F}^\vee)^\vee$ est injective.
\end{lemma}

\begin{proof}
On le vérifie à l'aide de la définition de $\mathcal{J}$. En effet, si
$s \in \mathcal{F}(U)$ est non nul, soit
$j_{U!}\mathcal{O}_U \to \mathcal{F}$ l'application de
$\mathcal{O}$-modules qui lui correspond par adjonction, et soit
$\mathcal{I}$ son noyau. Il existe une application non nulle
$\mathcal{F} \supset j_{U!}\mathcal{O}_U/\mathcal{I} \to \mathcal{J}$ qui
n'annule pas $s$. Comme $\mathcal{J}$ est un faisceau abélien injectif, cette
application sous-jacente de $\mathcal{O}$-modules se prolonge en une application
$\varphi : \mathcal{F} \to \mathcal{J}$.
Alors $ev(s)(\varphi) = \varphi(s) \not = 0$, ce qu'il fallait démontrer.
\end{proof}

\noindent
La surjection canonique $F(\mathcal{F}) \to \mathcal{F}$ de
$\mathcal{O}$-modules donne, d'après ce qui précède, une injection canonique de
$\mathcal{O}$-modules
$$
(\mathcal{F}^\vee)^\vee \longrightarrow (F(\mathcal{F}^\vee))^\vee.
$$
Posons $J(\mathcal{F}) = (F(\mathcal{F}^\vee))^\vee$. La composée de $ev$
avec l'application affichée donne
$\mathcal{F} \to J(\mathcal{F})$, fonctoriellement en $\mathcal{F}$.

\begin{lemma}
\label{injectives-lemma-JM-injective-sheaves}
Soit $\mathcal{O}$ un faisceau d'anneaux. Pour tout $\mathcal{O}$-module
$\mathcal{F}$, le $\mathcal{O}$-module $J(\mathcal{F})$ est injectif.
\end{lemma}

\begin{proof}
Il faut montrer que le foncteur
$\Hom_\mathcal{O}(\mathcal{G}, J(\mathcal{F}))$ est exact. Remarquons que
\begin{eqnarray*}
\Hom_\mathcal{O}(\mathcal{G}, J(\mathcal{F}))
& = &
\Hom_\mathcal{O}(\mathcal{G}, (F(\mathcal{F}^\vee))^\vee) \\
& = &
\Hom_\mathcal{O}
(\mathcal{G}, \SheafHom(F(\mathcal{F}^\vee), \mathcal{J})) \\
& = &
\Hom(\mathcal{G} \otimes_\mathcal{O} F(\mathcal{F}^\vee), \mathcal{J})
\end{eqnarray*}
La conclusion résulte donc de ce que $F(\mathcal{F}^\vee)$ est plat et
$\mathcal{J}$ injectif.
\end{proof}

\begin{theorem}
\label{injectives-theorem-sheaves-modules-injectives}
Soit $\mathcal{C}$ un site et soit $\mathcal{O}$ un faisceau d'anneaux sur
$\mathcal{C}$. La catégorie des faisceaux de $\mathcal{O}$-modules sur ce site
a assez d'injectifs. Plus précisément, il existe un plongement injectif
fonctoriel ; voir Homology, définition
\ref{homology-definition-functorial-injective-embedding}.
\end{theorem}

\begin{proof}
Cela résulte de la discussion de cette section.
\end{proof}

\begin{proposition}
\label{injectives-proposition-presheaves-modules}
Soit $\mathcal{C}$ une catégorie et soit $\mathcal{O}$ un préfaisceau
d'anneaux sur $\mathcal{C}$. La catégorie
$\textit{PMod}(\mathcal{O})$ des préfaisceaux de $\mathcal{O}$-modules possède
des plongements injectifs fonctoriels.
\end{proposition}

\begin{proof}
On pourrait procéder comme dans la discussion de la
section \ref{injectives-section-injectives-presheaves}. Nous utiliserons plutôt le
théorème précédent. Munissons $\mathcal{C}$ de la structure de site pour
laquelle les recouvrements d'un objet $U$ sont tous les singletons
$\{f : V \to U\}$ où $f$ est un isomorphisme. Nous omettons la vérification
qu'il s'agit bien d'un site. Pour cette topologie, un faisceau est la même chose
qu'un préfaisceau (preuve omise). Le théorème s'applique donc.
\end{proof}








\section{Plongement des catégories abéliennes}
\label{injectives-section-embedding}

\noindent
Dans cette section, nous montrons qu'une catégorie abélienne se plonge dans la
catégorie des faisceaux abéliens sur un site ayant assez de points. Le site
sera celui décrit dans le lemme suivant.

\begin{lemma}
\label{injectives-lemma-site-abelian-category}
Soit $\mathcal{A}$ une catégorie abélienne. Posons
$$
\text{Cov} = \{\{f : V \to U\} \mid f\text{ est surjectif}\}.
$$
Alors $(\mathcal{A}, \text{Cov})$ est un site ; voir
Sites, définition \ref{sites-definition-site}.
\end{lemma}

\begin{proof}
Selon nos conventions sur les catégories, $\Ob(\mathcal{A})$ est un ensemble.
Un isomorphisme est un morphisme surjectif, la composée de morphismes
surjectifs est surjective, et le changement de base d'un morphisme surjectif
dans $\mathcal{A}$ est surjectif ; voir
Homology, lemme \ref{homology-lemma-epimorphism-universal-abelian-category}.
\end{proof}

\noindent
Soit $\mathcal{A}$ une catégorie préadditive. Dans ce cas, le plongement de
Yoneda $\mathcal{A} \to \textit{PSh}(\mathcal{A})$, $X \mapsto h_X$, se
factorise par un foncteur $\mathcal{A} \to \textit{PAb}(\mathcal{A})$.

\begin{lemma}
\label{injectives-lemma-embedding}
Soit $\mathcal{A}$ une catégorie abélienne et soit
$\mathcal{C} = (\mathcal{A}, \text{Cov})$ le site défini au
lemme \ref{injectives-lemma-site-abelian-category}. Alors $X \mapsto h_X$ définit un
foncteur exact et pleinement fidèle
$$
\mathcal{A} \longrightarrow \textit{Ab}(\mathcal{C}).
$$
De plus, le site $\mathcal{C}$ a assez de points.
\end{lemma}

\begin{proof}
Supposons que $f : V \to U$ soit un morphisme surjectif de $\mathcal{A}$ et
posons $K = \Ker(f)$. Rappelons que
$V \times_U V = \Ker((f, -f) : V \oplus V \to U)$ ; voir
Homology, exemple \ref{homology-example-fibre-product-pushouts}. Il existe en
particulier une injection $K \oplus K \to V \times_U V$. Soient
$p, q : V \times_U V \to V$ les deux projections. Le morphisme
$p - q : V \times_U V \to V$ vérifie $f \circ (p  - q) = 0$ ; ainsi,
$p - q$ se factorise par $K \to V$. Notons
$c : V \times_U V \to K$ le morphisme ainsi
obtenu. Comme la composée $K \oplus K \to V \times_U V \to K$ est
surjective, $c$ est surjectif. Par conséquent,
$$
V \times_U V \xrightarrow{p - q} V \to U \to 0
$$
est une suite exacte de $\mathcal{A}$. Ainsi, pour un objet $X$ de
$\mathcal{A}$, la suite
$$
0 \to
\Hom_\mathcal{A}(U, X) \to
\Hom_\mathcal{A}(V, X) \to
\Hom_\mathcal{A}(V \times_U V, X)
$$
est une suite exacte de groupes abéliens ; voir
Homology, lemme \ref{homology-lemma-check-exactness}. Cela signifie que $h_X$
satisfait à la condition de faisceau sur $\mathcal{C}$.

\medskip\noindent
Le foncteur est pleinement fidèle d'après
Catégories, lemme \ref{categories-lemma-yoneda}. Il est exact à gauche entre
catégories abéliennes d'après
Homology, lemme \ref{homology-lemma-check-exactness}. Pour montrer qu'il est
exact à droite, soit $X \to Y$ un morphisme surjectif de $\mathcal{A}$. Soit
$U$ un objet de $\mathcal{A}$ et soit
$s \in h_Y(U) = \Mor_\mathcal{A}(U, Y)$ une section de $h_Y$ au-dessus de
$U$. D'après
Homology, lemme \ref{homology-lemma-epimorphism-universal-abelian-category},
la projection $U \times_Y X \to U$ est surjective. Ainsi,
$\{V = U \times_Y X \to U\}$ est un recouvrement de $U$ tel que $s|_V$ se
relève en une section de $h_X$. Ceci montre que $h_X \to h_Y$ est une
surjection de faisceaux abéliens ; voir
Sites, lemme \ref{sites-lemma-mono-epi-sheaves}.

\medskip\noindent
Le site $\mathcal{C}$ a assez de points d'après
Sites, proposition \ref{sites-proposition-criterion-points}.
\end{proof}

\begin{remark}
\label{injectives-remark-embedding}
Le théorème de plongement de Freyd--Mitchell affirme qu'il existe, de toute
catégorie abélienne $\mathcal{A}$ vers la catégorie des modules sur un anneau,
un foncteur exact et pleinement fidèle. Le lemme \ref{injectives-lemma-embedding} n'est
pas tout à fait aussi fort. Son résultat convient toutefois au Champs Project,
puisque nous devons de toute façon étudier en détail les faisceaux de groupes
abéliens sur des sites. En outre, la « chasse au diagramme » fonctionne dans la
catégorie des faisceaux abéliens sur $\mathcal{C}$, par exemple en travaillant
avec des sections au-dessus des objets, ou au niveau des germes en utilisant le
fait que $\mathcal{C}$ a assez de points. Pour déduire le théorème de plongement
de Freyd--Mitchell du lemme \ref{injectives-lemma-embedding}, voir la
remarque \ref{injectives-remark-embedding-freyd}.
\end{remark}

\begin{remark}
\label{injectives-remark-embedding-big}
Si $\mathcal{A}$ est une « grande » catégorie abélienne, c'est-à-dire si
$\mathcal{A}$ possède une classe d'objets, le lemme \ref{injectives-lemma-embedding} ne
s'applique pas. Dans ce cas, pour tout ensemble d'objets
$E \subset \Ob(\mathcal{A})$, il existe une sous-catégorie abélienne pleine
$\mathcal{A}' \subset \mathcal{A}$ telle que $\Ob(\mathcal{A}')$ soit un
ensemble et $E \subset \Ob(\mathcal{A}')$. On peut alors appliquer le
lemme \ref{injectives-lemma-embedding} à $\mathcal{A}'$. Cela permet de prouver que les
résultats fondés sur une chasse au diagramme restent valables dans
$\mathcal{A}$.
\end{remark}

\begin{remark}
\label{injectives-remark-embedding-freyd}
Soit $\mathcal{C}$ un site. La catégorie $\textit{Ab}(\mathcal{C})$ a assez
d'injectifs ; voir le théorème \ref{injectives-theorem-sheaves-injectives}. Lorsque
$\mathcal{C}$ a assez de points, cela est immédiat, car $p_*I$ est un faisceau
injectif si $I$ est un $\mathbf{Z}$-module injectif et $p$ un point. De plus,
$\textit{Ab}(\mathcal{C})$ possède un cogénérateur (détails omis). Le
lemme \ref{injectives-lemma-embedding} donne donc un plongement exact et pleinement fidèle
$\mathcal{A} \to \mathcal{B}$, où $\mathcal{B}$ possède un cogénérateur et
assez d'injectifs. En l'appliquant à $\mathcal{A}^{opp}$, on obtient un foncteur
exact et pleinement fidèle
$i : \mathcal{A} \to \mathcal{D} = \mathcal{B}^{opp}$, où $\mathcal{D}$ a
assez de projectifs et possède un générateur. Ainsi, $\mathcal{D}$ possède un
générateur projectif $P$. Posons $R = \Mor_\mathcal{D}(P, P)$. Alors
$$
\mathcal{A} \longrightarrow \text{Mod}_R, \quad
X \longmapsto \Hom_\mathcal{D}(P, X).
$$
On vérifie que ce foncteur est exact et pleinement fidèle. Autrement dit, on
retrouve le théorème de Freyd--Mitchell mentionné dans la
remarque \ref{injectives-remark-embedding}.
\end{remark}

\begin{remark}
\label{injectives-remark-embed-exact-category}
Les arguments qui démontrent les lemmes
\ref{injectives-lemma-site-abelian-category} et \ref{injectives-lemma-embedding} valent aussi pour
les {\it catégories exactes} ; voir \cite[Appendix A]{Buhler} et
\cite[1.1.4]{BBD}. Nous les rappelons brièvement ici et ajouterons des détails
si le Champs Project en a besoin.

\medskip\noindent
Soit $\mathcal{A}$ une catégorie additive. Une paire {\it noyau-conoyau} est
une paire $(i, p)$ de morphismes de $\mathcal{A}$,
$i : A \to B$, $p : B \to C$, telle que $i$ soit le noyau de $p$ et $p$ le
conoyau de $i$. Étant donné un ensemble $\mathcal{E}$ de paires
noyau-conoyau, on dit que $i : A \to B$ est un
{\it monomorphisme admissible} si $(i, p) \in \mathcal{E}$ pour un certain
morphisme $p$. De même, on dit qu'un morphisme $p : B \to C$ est un
{\it épimorphisme admissible} si $(i, p) \in \mathcal{E}$ pour un certain
morphisme $i$. La paire $(\mathcal{A}, \mathcal{E})$ est appelée
{\it catégorie exacte} si les axiomes suivants sont satisfaits
\begin{enumerate}
\item $\mathcal{E}$ est stable par isomorphismes de paires noyau-conoyau ;
\item pour tout objet $A$, le morphisme $1_A$ est à la fois un épimorphisme
admissible et un monomorphisme admissible ;
\item les monomorphismes admissibles sont stables par composition ;
\item les épimorphismes admissibles sont stables par composition ;
\item la somme amalgamée d'un monomorphisme admissible $i : A \to B$ par tout
morphisme $A \to A'$ existe et le morphisme induit $i' : A' \to B'$ est un
monomorphisme admissible ;
\item le changement de base d'un épimorphisme admissible $p : B \to C$ par
tout morphisme $C' \to C$ existe et le morphisme induit $p' : B' \to C'$ est
un épimorphisme admissible.
\end{enumerate}
Étant donnée une telle structure, posons
$\mathcal{C} = (\mathcal{A}, \text{Cov})$, où les recouvrements, c'est-à-dire
les éléments de $\text{Cov}$, sont les épimorphismes admissibles. Les axiomes
ci-dessus entraînent immédiatement qu'il s'agit d'un site. Considérons le
foncteur
$$
F : \mathcal{A} \longrightarrow \textit{Ab}(\mathcal{C}), \quad
X \longmapsto h_X
$$
exactement comme au lemme \ref{injectives-lemma-embedding}. Ce foncteur est pleinement
fidèle, exact, et reflète l'exactitude. En outre, toute extension d'objets de
l'image essentielle de $F$ appartient à l'image essentielle de $F$.
\end{remark}






\section{Conditions AB de Grothendieck}
\label{injectives-section-grothendieck-conditions}

\noindent
Cette section et les quelques suivantes concernent surtout les « grandes »
catégories abéliennes, c'est-à-dire les catégories énumérées dans
Catégories, remarque \ref{categories-remark-big-categories}. L'exemple à
garder à l'esprit est la catégorie des faisceaux de modules sur un site annelé.

\medskip\noindent
Grothendieck a démontré l'existence d'injectifs dans une grande généralité dans
l'article \cite{Tohoku}. Il a utilisé les conditions suivantes pour distinguer
les catégories abéliennes dotées de propriétés particulières.

\begin{definition}
\label{injectives-definition-grothendieck-conditions}
Soit $\mathcal{A}$ une catégorie abélienne. Nous nommons quelques conditions :
\begin{enumerate}
\item[AB3] $\mathcal{A}$ possède des sommes directes ;
\item[AB4] $\mathcal{A}$ satisfait AB3 et les sommes directes sont exactes ;
\item[AB5] $\mathcal{A}$ satisfait AB3 et les colimites filtrantes sont
exactes.
\end{enumerate}
Voici les notions duales :
\begin{enumerate}
\item[AB3*] $\mathcal{A}$ possède des produits ;
\item[AB4*] $\mathcal{A}$ satisfait AB3* et les produits sont exacts ;
\item[AB5*] $\mathcal{A}$ satisfait AB3* et les limites cofiltrantes sont
exactes.
\end{enumerate}
On dit qu'un objet $U$ de $\mathcal{A}$ est un {\it générateur} si, pour tous
$N \subset M$, $N \not = M$, dans $\mathcal{A}$, il existe un morphisme
$U \to M$ qui ne se factorise pas par $N$. On dit que $\mathcal{A}$ est une
{\it catégorie abélienne de Grothendieck} si elle satisfait AB5 et possède un
générateur.
\end{definition}

\noindent
Discussion : une somme directe dans une catégorie abélienne est un coproduit.
Si une catégorie abélienne possède des sommes directes, c'est-à-dire satisfait
AB3, alors elle possède des colimites ; voir
Catégories, lemme \ref{categories-lemma-colimits-coproducts-coequalizers}. 
De même, si $\mathcal{A}$ satisfait AB3*, alors elle possède des limites ; voir
Catégories, lemme \ref{categories-lemma-limits-products-equalizers}. 
L'exactitude des sommes directes signifie ceci : étant donnés un ensemble
d'indices $I$ et des suites exactes courtes
$$
0 \to A_i \to B_i \to C_i \to 0,\quad i \in I
$$
dans $\mathcal{A}$, la suite
$$
0 \to
\bigoplus\nolimits_{i \in I} A_i \to
\bigoplus\nolimits_{i \in I} B_i \to
\bigoplus\nolimits_{i \in I} C_i \to 0
$$
est elle aussi exacte. Sans supposer AB4, on peut seulement affirmer en général
que cette suite est exacte à droite, c'est-à-dire que la prise des sommes
directes est un foncteur exact à droite lorsqu'elles existent. De même,
l'exactitude des colimites filtrantes signifie ceci : étant donnés un ensemble
filtrant $I$ et un système de suites exactes courtes
$$
0 \to A_i \to B_i \to C_i \to 0
$$
indexé par $I$ dans $\mathcal{A}$, la suite
$$
0 \to
\colim_{i \in I} A_i \to
\colim_{i \in I} B_i \to
\colim_{i \in I} C_i \to 0
$$
est elle aussi exacte. Sans supposer AB5, on peut seulement affirmer en général
que cette suite est exacte à droite, c'est-à-dire que la prise des colimites
est un foncteur exact à droite lorsqu'elles existent. Une explication analogue
vaut pour AB4* et AB5*.
\section{Objets injectifs dans les catégories de Grothendieck}
\label{injectives-section-grothendieck-categories}

\noindent
L'existence d'un générateur implique que, pour tout objet $M$ d'une
catégorie abélienne de Grothendieck $\mathcal{A}$, ses sous-objets
forment un ensemble. (Cela peut être faux pour une catégorie abélienne
``grande'' quelconque.)

\begin{lemma}
\label{injectives-lemma-set-of-subobjects}
Soient $\mathcal{A}$ une catégorie abélienne munie d'un générateur $U$ et
$X$ un objet de $\mathcal{A}$. Si $\kappa$ est le cardinal de
$\Mor(U, X)$, alors
\begin{enumerate}
\item Il n'existe aucune chaîne strictement croissante
(ni strictement décroissante) de sous-objets de $X$ indexée par
un cardinal strictement supérieur à $\kappa$.
\item Si $\alpha$ est un ordinal de cofinalité $> \kappa$,
toute suite croissante (ou décroissante) de sous-objets de $X$
indexée par $\alpha$ est stationnaire à partir d'un certain rang.
\item Le cardinal de l'ensemble des sous-objets de $X$
est $\leq 2^\kappa$.
\end{enumerate}
\end{lemma}

\begin{proof}
Pour (1), soit $\kappa' > \kappa$ un cardinal et supposons que
la famille $X_i$, $i \in \kappa'$, soit strictement croissante. Pour tout
$i$, choisissons alors $\phi_i \in \Mor(U, X)$ tel que $\phi_i$ se factorise
par $X_{i + 1}$, mais non par $X_i$. Les morphismes $\phi_i$ sont donc
deux à deux distincts, ce qui contredit la définition de $\kappa$.

\medskip\noindent
L'assertion (2) résulte de la définition de la cofinalité et de (1).

\medskip\noindent
Démonstration de (3). Pour tout sous-objet $Y \subset X$, définissons
$S_Y \in \mathcal{P}(\Mor(U, X))$ (l'ensemble des parties) par
$S_Y = \{\phi \in \Mor(U,X) : \phi\text{ se factorise par }Y\}$.
Alors $Y = Y'$ si et seulement si $S_Y = S_{Y'}$. Le cardinal de
l'ensemble des sous-objets est donc au plus celui de cet ensemble des parties.
\end{proof}

\noindent
D'après le lemme \ref{injectives-lemma-set-of-subobjects}, la définition suivante
a un sens.

\begin{definition}
\label{injectives-definition-size}
Soit $\mathcal{A}$ une catégorie abélienne de Grothendieck.
Soit $M$ un objet de $\mathcal{A}$.
La {\it taille} $|M|$ de $M$ est le cardinal de l'ensemble des sous-objets
de $M$.
\end{definition}

\begin{lemma}
\label{injectives-lemma-size-goes-down}
Soit $\mathcal{A}$ une catégorie abélienne de Grothendieck.
Si $0 \to M' \to M \to M'' \to 0$ est une suite exacte courte dans
$\mathcal{A}$, alors $|M'|, |M''| \leq |M|$.
\end{lemma}

\begin{proof}
Cela résulte immédiatement des définitions.
\end{proof}

\begin{lemma}
\label{injectives-lemma-set-iso-classes-bounded-size}
Soit $\mathcal{A}$ une catégorie abélienne de Grothendieck munie d'un
générateur $U$.
\begin{enumerate}
\item Si $|M| \leq \kappa$, alors $M$ est un quotient d'une somme directe
d'au plus $\kappa$ copies de $U$.
\item Pour tout cardinal $\kappa$, les classes d'isomorphisme des objets
$M$ tels que $|M| \leq \kappa$ forment un ensemble.
\end{enumerate}
\end{lemma}

\begin{proof}
Pour (1), choisissons, pour tout sous-objet propre $M' \subset M$, un morphisme
$\varphi_{M'} : U \to M$ dont l'image n'est pas contenue dans $M'$. Alors
$\bigoplus_{M' \subset M} \varphi_{M'} : \bigoplus_{M' \subset M} U \to M$
est surjectif. Il est clair que (1) implique (2).
\end{proof}

\begin{proposition}
\label{injectives-proposition-objects-are-small}
Soit $\mathcal{A}$ une catégorie abélienne de Grothendieck. Soit $M$ un
objet de $\mathcal{A}$. Posons $\kappa = |M|$.
Si $\alpha$ est un ordinal dont la cofinalité est strictement supérieure
à $\kappa$, alors $M$ est $\alpha$-petit relativement aux monomorphismes.
\end{proposition}

\begin{proof}
Comparons avec la proposition \ref{injectives-proposition-modules-are-small}.
Il suffit de montrer que l'application (\ref{injectives-equation-compare}) est surjective.
Soit $f : M \to \colim B_\beta$ une application.
Considérons les sous-objets $\{f^{-1}(B_\beta)\}$ de $M$, où $B_\beta$
est vu comme un sous-objet de la colimite $B = \bigcup_\beta B_\beta$.
Si l'un d'eux, disons $f^{-1}(B_\beta)$, est égal à $M$,
alors l'application se factorise par $B_\beta$.

\medskip\noindent
Puisque $\mathcal{A}$ vérifie AB5, on a
$$
\colim f^{-1}(B_\beta) = f^{-1}\left(\colim B_\beta\right) = M.
$$
Par hypothèse, il y a au plus $\kappa$ sous-objets distincts de $M$ parmi
les $f^{-1}(B_\alpha)$.
On peut donc trouver une partie $S \subset \alpha$ de cardinal au plus
$\kappa$ telle que, lorsque $\beta'$ parcourt $S$, les
$f^{-1}(B_{\beta'})$ parcourent \emph{tous} les $f^{-1}(B_\alpha)$.

\medskip\noindent
Toutefois, $S$ admet un majorant $\widetilde{\alpha} < \alpha$, puisque
$\alpha$ est de cofinalité strictement supérieure à $\kappa$. En particulier,
tous les $f^{-1}(B_{\beta'})$, $\beta' \in S$, sont contenus dans
$f^{-1}(B_{\widetilde{\alpha}})$.
Il s'ensuit que $f^{-1}(B_{\widetilde{\alpha}}) = M$.
L'application $f$ se factorise donc par $B_{\widetilde{\alpha}}$.
\end{proof}

\begin{lemma}
\label{injectives-lemma-characterize-injective}
\begin{slogan}
Pour vérifier qu'un objet est injectif, il suffit de vérifier la propriété
de relèvement pour les sous-objets d'un générateur.
\end{slogan}
Soit $\mathcal{A}$ une catégorie abélienne de Grothendieck munie d'un
générateur $U$.
Un objet $I$ de $\mathcal{A}$ est injectif si et seulement si, dans tout
diagramme commutatif
$$
\xymatrix{
M \ar[d] \ar[r] &  I \\
U \ar@{-->}[ru]
}
$$
où $M \subset U$ est un sous-objet, la flèche pointillée existe.
\end{lemma}

\begin{proof}
Voir le lemme \ref{injectives-lemma-criterion-baer} pour le cas des modules.
Soient un monomorphisme $A \subset B$ et un morphisme $\varphi : A \to I$.
Considérons l'ensemble $S$ des couples $(A', \varphi')$ formés de
sous-objets $A \subset A' \subset B$ et d'un morphisme $\varphi' : A' \to I$
prolongeant $\varphi$. Munissons cet ensemble de l'ordre partiel évident.
Soit $T \subset S$ une partie totalement ordonnée. Alors
$$
A' = \colim_{t \in T} A_t \xrightarrow{\colim_{t \in T} \varphi_t} I
$$
est un majorant. Le lemme de Zorn assure donc que $S$ possède un élément
maximal $(A', \varphi')$. Nous affirmons que $A' = B$. Sinon, choisissons
un morphisme $\psi : U \to B$ qui ne se factorise pas par $A'$. Posons
$N = A' \cap \psi(U)$ et $M = \psi^{-1}(N)$. L'application
$$
M \to N \to A' \xrightarrow{\varphi'} I
$$
se prolonge en un morphisme $\chi : U \to I$. Comme
$\chi|_{\Ker(\psi)} = 0$, le morphisme $\chi$ se factorise sous la forme
$$
U \to \Im(\psi) \xrightarrow{\varphi''} I
$$
Puisque $\varphi'$ et $\varphi''$ coïncident sur
$N = A' \cap \Im(\psi)$, ils définissent ensemble un morphisme
$A' + \Im(\psi) \to I$, ce qui contredit la maximalité supposée de $A'$.
\end{proof}

\begin{theorem}
\label{injectives-theorem-injective-embedding-grothendieck}
Soit $\mathcal{A}$ une catégorie abélienne de Grothendieck.
Alors $\mathcal{A}$ admet des plongements injectifs fonctoriels.
\end{theorem}

\begin{proof}
Comparons avec la démonstration du
théorème \ref{injectives-theorem-baer-grothendieck}.
Choisissons un générateur $U$ de $\mathcal{A}$. Pour un objet $M$, on définit
$\mathbf{M}(M)$ par le diagramme cocartésien suivant :
$$
\xymatrix{
\bigoplus_{N \subset U}
\bigoplus_{\varphi \in \Hom(N, M)}
N \ar[r] \ar[d] & M \ar[d] \\
\bigoplus_{N \subset U}
\bigoplus_{\varphi \in \Hom(N, M)}
U \ar[r] &  \mathbf{M}(M).
}
$$
Notons que $M \mapsto \mathbf{M}(M)$ est un foncteur et qu'il existe des
applications injectives fonctorielles $M \to \mathbf{M}(M)$. On définit par
récurrence transfinie des foncteurs $\mathbf{M}_\alpha(M)$ pour tout ordinal
$\alpha$. On pose $\mathbf{M}_0(M) = M$. Une fois $\mathbf{M}_\alpha(M)$
défini, on pose
$\mathbf{M}_{\alpha + 1}(M) = \mathbf{M}(\mathbf{M}_\alpha(M))$.
Pour un ordinal limite $\beta$, on pose
$$
\mathbf{M}_\beta(M) = \colim_{\alpha < \beta} \mathbf{M}_\alpha(M).
$$
Enfin, choisissons un ordinal $\alpha$ dont la cofinalité est strictement
supérieure à $|U|$. Un tel ordinal existe d'après
Ensembles, proposition \ref{sets-proposition-exist-ordinals-large-cofinality}.
Nous affirmons que $M \to \mathbf{M}_\alpha(M)$ est le plongement injectif
fonctoriel recherché. En effet, si $N \subset U$ est un sous-objet et si
$\varphi : N \to \mathbf{M}_\alpha(M)$ est un morphisme, alors
$\varphi$ se factorise par $\mathbf{M}_{\alpha'}(M)$ pour un certain
$\alpha' < \alpha$, d'après la
proposition \ref{injectives-proposition-objects-are-small}.
Par construction de $\mathbf{M}(-)$, le morphisme $\varphi$ se prolonge en
un morphisme de $U$ dans $\mathbf{M}_{\alpha' + 1}(M)$, donc dans
$\mathbf{M}_\alpha(M)$. Le
lemme \ref{injectives-lemma-characterize-injective}
montre alors que $\mathbf{M}_\alpha(M)$ est injectif.
\end{proof}













\section{Complexes K-injectifs dans les catégories de Grothendieck}
\label{injectives-section-K-injective}

\noindent
Le contenu de cette section provient de l'article \cite{serpe} de Serp\'e.
Cet article généralise certains résultats de Spaltenstein
\cite{Spaltenstein} à des catégories abéliennes de Grothendieck quelconques.
Notre lemme \ref{injectives-lemma-characterize-K-injective} n'apparaît
qu'implicitement dans l'article de Serp\'e. Notre démarche consiste à imiter
la démonstration de Grothendieck du
théorème \ref{injectives-theorem-injective-embedding-grothendieck}.

\begin{lemma}
\label{injectives-lemma-surjection-bounded-size}
Soit $\mathcal{A}$ une catégorie abélienne de Grothendieck munie d'un
générateur $U$. Soit $c$ la fonction sur les cardinaux définie par
$c(\kappa) = |\bigoplus_{\alpha \in \kappa} U|$. Si $\pi : M \to N$ est
surjectif, il existe un sous-objet $M' \subset M$ qui se surjecte sur $N$
et tel que $|M'| \leq c(|N|)$.
\end{lemma}

\begin{proof}
Pour tout sous-objet propre $N' \subset N$, choisissons un morphisme
$\varphi_{N'} : U \to M$ tel que $U \to M \to N$ ne se factorise pas
par $N'$. Posons
$$
M' = \Im\left(
\bigoplus\nolimits_{N' \subset N} \varphi_{N'} :
\bigoplus\nolimits_{N' \subset N} U \longrightarrow M\right)
$$
Alors $M'$ convient.
\end{proof}

\begin{lemma}
\label{injectives-lemma-acyclic-quotient-complexes-bounded-size}
Soit $\mathcal{A}$ une catégorie abélienne de Grothendieck. Il existe un
cardinal $\kappa$ tel que, pour tout complexe acyclique $M^\bullet$, on ait :
\begin{enumerate}
\item si $M^\bullet$ est non nul, il existe un sous-complexe non nul
$N^\bullet$, borné supérieurement et acyclique, tel que
$|N^n| \leq \kappa$ ;
\item il existe une surjection de complexes
$$
\bigoplus\nolimits_{i \in I} M_i^\bullet \longrightarrow M^\bullet
$$
où $M_i^\bullet$ est borné supérieurement, acyclique, et
$|M_i^n| \leq \kappa$.
\end{enumerate}
\end{lemma}

\begin{proof}
Choisissons un générateur $U$ de $\mathcal{A}$. Notons $c$ la fonction du
lemme \ref{injectives-lemma-surjection-bounded-size}.
Posons $\kappa = \sup \{c^n(|U|), n = 1, 2, 3, \ldots\}$.
Soient $n \in \mathbf{Z}$ et un morphisme $\psi : U \to M^n$.
Pour démontrer (1) et (2), il suffit d'établir qu'il existe un sous-complexe
$N^\bullet \subset M^\bullet$, borné supérieurement et acyclique, tel que
$|N^m| \leq \kappa$ et que $\psi$ se factorise par $N^n$.
À cette fin, posons $N^n = \Im(\psi)$,
$N^{n + 1} = \Im(U \to M^n \to M^{n + 1})$, et $N^m = 0$ pour
$m \geq n + 2$.
Supposons que, pour un certain $k \leq n$, on ait construit
$N^m \subset M^m$ pour tout $n > m \geq k$, de telle sorte que
\begin{enumerate}
\item $\text{d}(N^m) \subset N^{m + 1}$ pour tout $m \geq k$ ;
\item $\Im(N^{m - 1} \to N^m) = \Ker(N^m \to N^{m + 1})$ pour
tout $m \geq k + 1$ ;
\item $|N^m| \leq c^{\max\{n - m, 0\}}(|U|)$ pour tout $m \geq k$.
\end{enumerate}
Comme $M^\bullet$ est acyclique, le sous-objet
$\text{d}^{-1}(\Ker(N^k \to N^{k + 1})) \subset M^{k - 1}$ se surjecte sur
$\Ker(N^k \to N^{k + 1})$. Le
lemme \ref{injectives-lemma-surjection-bounded-size} permet donc de choisir
$N^{k - 1} \subset M^{k - 1}$ se surjectant sur
$\Ker(N^k \to N^{k + 1})$ et vérifiant
$|N^{k - 1}| \leq c^{n - k + 1}(|U|)$.
On conclut par récurrence sur $k$.
\end{proof}

\begin{lemma}
\label{injectives-lemma-characterize-K-injective}
Soit $\mathcal{A}$ une catégorie abélienne de Grothendieck.
Soit $\kappa$ un cardinal comme dans le
lemme \ref{injectives-lemma-acyclic-quotient-complexes-bounded-size}.
Supposons que $I^\bullet$ soit un complexe tel que
\begin{enumerate}
\item chaque $I^j$ soit injectif ;
\item pour tout complexe acyclique $M^\bullet$, borné supérieurement,
tel que $|M^n| \leq \kappa$, on ait
$\Hom_{K(\mathcal{A})}(M^\bullet, I^\bullet) = 0$.
\end{enumerate}
Alors $I^\bullet$ est un complexe $K$-injectif.
\end{lemma}

\begin{proof}
Soit $M^\bullet$ un complexe acyclique. Nous allons construire, par
récurrence sur l'ordinal $\alpha$, un sous-complexe acyclique
$K_\alpha^\bullet \subset M^\bullet$ de la manière suivante.
Pour $\alpha = 0$, posons $K_0^\bullet = 0$. Pour $\alpha > 0$,
procédons comme suit :
\begin{enumerate}
\item Si $\alpha = \beta + 1$ et $K_\beta^\bullet = M^\bullet$,
on choisit $K_\alpha^\bullet = K_\beta^\bullet$.
\item Si $\alpha = \beta + 1$ et $K_\beta^\bullet \not = M^\bullet$,
alors $M^\bullet/K_\beta^\bullet$ est un complexe acyclique non nul.
On choisit un sous-complexe
$N_\alpha^\bullet \subset M^\bullet/K_\beta^\bullet$ comme dans le
lemme \ref{injectives-lemma-acyclic-quotient-complexes-bounded-size}.
Enfin, on prend pour $K_\alpha^\bullet \subset M^\bullet$
l'image inverse de $N_\alpha^\bullet$.
\item Si $\alpha$ est un ordinal limite, on pose
$K_\alpha^\bullet = \colim_{\beta < \alpha} K_\beta^\bullet$.
\end{enumerate}
Il est clair que $M^\bullet = K_\alpha^\bullet$ pour un ordinal $\alpha$
assez grand. Nous allons démontrer que
$$
\Hom_{K(\mathcal{A})}(K_\alpha^\bullet, I^\bullet)
$$
est nul par récurrence transfinie sur $\alpha$. C'est vrai pour $\alpha = 0$,
car $K_0^\bullet$ est nul. Supposons le résultat vrai pour $\beta$ et
$\alpha = \beta + 1$. Dans le cas (1) de la liste ci-dessus, le résultat est
clair. Dans le cas (2), on dispose d'une suite exacte courte de complexes
$$
0 \to K_\beta^\bullet \to K_\alpha^\bullet \to N_\alpha^\bullet \to 0
$$
Puisque chaque composante de $I^\bullet$ est injective, on obtient une suite
exacte
$$
\Hom_{K(\mathcal{A})}(K_\beta^\bullet, I^\bullet) \leftarrow
\Hom_{K(\mathcal{A})}(K_\alpha^\bullet, I^\bullet) \leftarrow
\Hom_{K(\mathcal{A})}(N_\alpha^\bullet, I^\bullet)
$$
Par récurrence, le terme de gauche est nul et, par hypothèse sur
$I^\bullet$, celui de droite est nul. Le groupe du milieu est donc lui aussi
nul. Supposons enfin que $\alpha$ soit un ordinal limite. Alors
$$
\Hom^\bullet(K_\alpha^\bullet, I^\bullet) =
\lim_{\beta < \alpha} \Hom^\bullet(K_\beta^\bullet, I^\bullet)
$$
avec les notations de
Compléments d'algèbre, section \ref{more-algebra-section-hom-complexes}.
Ces complexes calculent les morphismes dans $K(\mathcal{A})$ d'après
Compléments d'algèbre, équation
(\ref{more-algebra-equation-cohomology-hom-complex}).
Les applications de transition du système sont surjectives, puisque $I^j$
est injectif pour tout $j$. En outre, pour un ordinal limite $\alpha$, la
limite est égale à la valeur (voir la formule affichée ci-dessus). On peut donc
conclure à l'aide de
Homologie, lemme \ref{homology-lemma-ML-over-ordinals}.
\end{proof}

\begin{lemma}
\label{injectives-lemma-functorial-homotopies}
Soit $\mathcal{A}$ une catégorie abélienne de Grothendieck.
Soit $(K_i^\bullet)_{i \in I}$ un ensemble de complexes acycliques.
Il existe un foncteur $M^\bullet \mapsto \mathbf{M}^\bullet(M^\bullet)$
et une transformation naturelle
$j_{M^\bullet} : M^\bullet \to \mathbf{M}^\bullet(M^\bullet)$
tels que
\begin{enumerate}
\item $j_{M^\bullet}$ soit un quasi-isomorphisme injectif (terme à terme) ;
\item pour tout $i \in I$ et tout $w : K_i^\bullet \to M^\bullet$,
le morphisme $j_{M^\bullet} \circ w$ soit homotope à zéro.
\end{enumerate}
\end{lemma}

\begin{proof}
Pour tout $i \in I$, choisissons une application de complexes injective
terme à terme $K_i^\bullet \to L_i^\bullet$, homotope à zéro, où
$L_i^\bullet$ est quasi-isomorphe à zéro. On peut par exemple prendre pour
$L_i^\bullet$ le cône de l'identité de $K_i^\bullet$.
On définit $\mathbf{M}^\bullet(M^\bullet)$ par le diagramme cocartésien suivant :
$$
\xymatrix{
\bigoplus_{i \in I}
\bigoplus_{w : K_i^\bullet \to M^\bullet}
K_i^\bullet \ar[r] \ar[d] & M^\bullet \ar[d] \\
\bigoplus_{i \in I}
\bigoplus_{w : K_i^\bullet \to M^\bullet}
L_i^\bullet \ar[r] &  \mathbf{M}^\bullet(M^\bullet).
}
$$
Alors $M^\bullet \mapsto \mathbf{M}^\bullet(M^\bullet)$ est un foncteur. La flèche
verticale de droite définit l'application injective fonctorielle
$j_{M^\bullet}$. Le conoyau de $j_{M^\bullet}$ est isomorphe à la somme directe
des conoyaux des
applications $K_i^\bullet \to L_i^\bullet$, et est donc acyclique. Ainsi,
$j_{M^\bullet}$ est un quasi-isomorphisme. L'assertion (2) est vraie par
construction.
\end{proof}

\begin{lemma}
\label{injectives-lemma-functorial-injective}
Soit $\mathcal{A}$ une catégorie abélienne de Grothendieck.
Il existe un foncteur $M^\bullet \mapsto \mathbf{N}^\bullet(M^\bullet)$
et une transformation naturelle
$j_{M^\bullet} : M^\bullet \to \mathbf{N}^\bullet(M^\bullet)$
tels que
\begin{enumerate}
\item $j_{M^\bullet}$ soit un quasi-isomorphisme injectif (terme à terme) ;
\item pour tout $n \in \mathbf{Z}$, l'application
$M^n \to \mathbf{N}^n(M^\bullet)$ se factorise par un sous-objet
$I^n \subset \mathbf{N}^n(M^\bullet)$, où $I^n$ est un objet injectif de
$\mathcal{A}$.
\end{enumerate}
\end{lemma}

\begin{proof}
Choisissons des plongements injectifs fonctoriels $i_M : M \to I(M)$, voir le
théorème \ref{injectives-theorem-injective-embedding-grothendieck}.
Pour tout complexe $M^\bullet$, notons $J^\bullet(M^\bullet)$ le complexe
dont les termes sont $J^n(M^\bullet) = I(M^n) \oplus I(M^{n + 1})$ et dont
la différentielle est
$$
d_{J^\bullet(M^\bullet)} =
\left(
\begin{matrix}
0 & 1 \\
0 & 0
\end{matrix}
\right)
$$
Il existe une application injective canonique de complexes
$u_{M^\bullet} : M^\bullet \to J^\bullet(M^\bullet)$, obtenue en envoyant
$M^n$ dans $I(M^n) \oplus I(M^{n + 1})$ au moyen des applications
$i_{M^n} : M^n \to I(M^n)$ et
$i_{M^{n + 1}} \circ d : M^n \to M^{n + 1} \to I(M^{n + 1})$. On dispose donc
d'une suite exacte courte de complexes
$$
0 \to M^\bullet \xrightarrow{u_{M^\bullet}}
J^\bullet(M^\bullet) \xrightarrow{v_{M^\bullet}}
Q^\bullet(M^\bullet) \to 0
$$
fonctorielle en $M^\bullet$. Posons
$$
\mathbf{N}^\bullet(M^\bullet) = C(v_{M^\bullet})^\bullet[-1].
$$
Remarquons que
$$
\mathbf{N}^n(M^\bullet) = Q^{n - 1}(M^\bullet) \oplus J^n(M^\bullet)
$$
dont la différentielle est
$$
\left(
\begin{matrix}
- d^{n - 1}_{Q^\bullet(M^\bullet)} & - v^n_{M^\bullet} \\
0 & d^n_{J^\bullet(M^\bullet)}
\end{matrix}
\right)
$$
Il existe donc une application de complexes
$j_{M^\bullet} : M^\bullet \to \mathbf{N}^\bullet(M^\bullet)$ induite par $u$.
Elle est injective et se factorise par un sous-objet injectif par construction.
L'application $j_{M^\bullet}$ est un quasi-isomorphisme, comme le montre la
suite exacte longue de cohomologie associée aux suites exactes courtes de
complexes ci-dessus.
\end{proof}

\begin{theorem}
\label{injectives-theorem-K-injective-embedding-grothendieck}
\begin{slogan}
Existence de complexes K-injectifs dans les catégories abéliennes de
Grothendieck.
\end{slogan}
Soit $\mathcal{A}$ une catégorie abélienne de Grothendieck.
Pour tout complexe $M^\bullet$, il existe un quasi-isomorphisme
$M^\bullet \to I^\bullet$ tel que $M^n \to I^n$ soit injectif, que $I^n$
soit un objet injectif de $\mathcal{A}$ pour tout $n$, et que $I^\bullet$
soit un complexe K-injectif. En outre, cette construction est fonctorielle
en $M^\bullet$.
\end{theorem}

\begin{proof}
Comparons avec les démonstrations des
théorèmes \ref{injectives-theorem-baer-grothendieck} et
\ref{injectives-theorem-injective-embedding-grothendieck}.
Choisissons un cardinal $\kappa$ comme dans les
lemmes \ref{injectives-lemma-acyclic-quotient-complexes-bounded-size} et
\ref{injectives-lemma-characterize-K-injective}.
Choisissons un ensemble $(K_i^\bullet)_{i \in I}$ de complexes acycliques
bornés supérieurement tel que tout complexe acyclique $K^\bullet$, borné
supérieurement et vérifiant $|K^n| \leq \kappa$, soit isomorphe à
$K_i^\bullet$ pour un certain $i \in I$. Cela est possible d'après le
lemme \ref{injectives-lemma-set-iso-classes-bounded-size}.
Notons $\mathbf{M}^\bullet(-)$ le foncteur construit dans le
lemme \ref{injectives-lemma-functorial-homotopies}, et
$\mathbf{N}^\bullet(-)$ celui construit dans le
lemme \ref{injectives-lemma-functorial-injective}.
Ces deux foncteurs sont munis de transformations injectives
$\text{id} \to \mathbf{M}$ et $\text{id} \to \mathbf{N}$.

\medskip\noindent
Par récurrence transfinie, on définit une suite de foncteurs
$\mathbf{T}_\alpha(-)$ et les transformations correspondantes
$\text{id} \to \mathbf{T}_\alpha$. On pose
$\mathbf{T}_0(M^\bullet) = M^\bullet$. Si $\mathbf{T}_\alpha$ est défini,
on pose
$$
\mathbf{T}_{\alpha + 1}(M^\bullet) =
\mathbf{N}^\bullet(\mathbf{M}^\bullet(\mathbf{T}_\alpha(M^\bullet)))
$$
Si $\beta$ est un ordinal limite, on pose
$$
\mathbf{T}_\beta(M^\bullet) =
\colim_{\alpha < \beta} \mathbf{T}_\alpha(M^\bullet)
$$
Les applications de transition du système sont des quasi-isomorphismes
injectifs. D'après AB5, la colimite est encore quasi-isomorphe à $M^\bullet$.
Nous affirmons que $M^\bullet \to \mathbf{T}_\alpha(M^\bullet)$ convient si
la cofinalité de $\alpha$ est strictement supérieure à
$\max(\kappa, |U|)$, où $U$ est un générateur de $\mathcal{A}$.
Il suffit en effet de vérifier les conditions (1) et (2) du
lemme \ref{injectives-lemma-characterize-K-injective}.

\medskip\noindent
Pour (1), utilisons le critère du
lemme \ref{injectives-lemma-characterize-injective}.
Supposons que $M \subset U$ et que
$\varphi : M \to \mathbf{T}^n_\alpha(M^\bullet)$ soit un morphisme pour un
certain $n \in \mathbf{Z}$. La
proposition \ref{injectives-proposition-objects-are-small} montre que $\varphi$ se
factorise par $\mathbf{T}^n_{\alpha'}(M^\bullet)$ pour un certain
$\alpha' < \alpha$. En particulier, par construction du foncteur
$\mathbf{N}^\bullet(-)$, le morphisme $\varphi$ se factorise par un objet
injectif de $\mathcal{A}$ ; cela montre que $\varphi$ se relève en un morphisme
défini sur $U$.

\medskip\noindent
Pour (2), soit $w : K^\bullet \to \mathbf{T}_\alpha(M^\bullet)$ un morphisme
de complexes, où $K^\bullet$ est un complexe acyclique borné supérieurement
tel que $|K^n| \leq \kappa$. Alors $K^\bullet \cong K_i^\bullet$ pour un
certain $i \in I$. De nouveau, la
proposition \ref{injectives-proposition-objects-are-small} montre que $w$ se factorise
par $\mathbf{T}_{\alpha'}(M^\bullet)$ pour un certain $\alpha' < \alpha$.
Par construction du foncteur $\mathbf{M}^\bullet(-)$, il s'ensuit que $w$
est homotope à zéro. Cela achève la démonstration.
\end{proof}







\section{Remarques complémentaires sur les catégories abéliennes de
Grothendieck}
\label{injectives-section-additional-Grothendieck}

\noindent
Nous rassemblons dans cette section quelques résultats classiques sur les
catégories abéliennes de Grothendieck.

\begin{lemma}
\label{injectives-lemma-grothendieck-brown}
Soit $\mathcal{A}$ une catégorie abélienne de Grothendieck.
Soit $F : \mathcal{A}^{opp} \to \textit{Ensembles}$ un foncteur.
Alors $F$ est représentable si et seulement si $F$ transforme les colimites
en limites, c'est-à-dire si
$$
F(\colim_i N_i) = \lim F(N_i)
$$
pour tout diagramme $\mathcal{I} \to \mathcal{A}$, $i \in \mathcal{I}$.
\end{lemma}

\begin{proof}
Si $F$ est représentable, il transforme les colimites en limites par
définition des colimites.

\medskip\noindent
Supposons que $F$ transforme les colimites en limites. Alors
$F(M \oplus N) = F(M) \times F(N)$, ce qui permet de munir $F(M)$ d'une
structure de groupe. On obtient ainsi un foncteur
$F : \mathcal{A}^{opp} \to \textit{Ab}$ additif et exact à droite au sens
contravariant, c'est-à-dire qu'il transforme une suite exacte courte
$0 \to K \to L \to M \to 0$ en une suite exacte
$F(K) \leftarrow F(L) \leftarrow F(M) \leftarrow 0$ (comparer avec
Homologie, section \ref{homology-section-functors}).

\medskip\noindent
Soit $U$ un générateur de $\mathcal{A}$. Posons
$A = \bigoplus_{s \in F(U)} U$ et
$s_{univ} = (s)_{s \in F(U)} \in F(A) = \prod_{s \in F(U)} F(U)$.
Soit $A' \subset A$ le plus grand sous-objet tel que la restriction de
$s_{univ}$ à $A'$ soit nulle. Il existe parce que $\mathcal{A}$ est une
catégorie de Grothendieck et que $F$ transforme les colimites en limites.
Puisque $F$ transforme les colimites en limites, il existe aussi un unique
élément $\overline{s}_{univ} \in F(A/A')$ qui s'envoie sur $s_{univ}$ dans
$F(A)$.
Nous affirmons que $A/A'$ représente $F$ ; autrement dit, l'application de
Yoneda
$$
\overline{s}_{univ} : h_{A/A'} \longrightarrow F
$$
est un isomorphisme. Soient $M \in \Ob(\mathcal{A})$ et $s \in F(M)$.
Considérons la surjection
$$
c_M :
A_M = \bigoplus\nolimits_{\varphi \in \Hom_\mathcal{A}(U, M)} U
\longrightarrow
M.
$$
Elle donne $F(c_M)(s) = (s_\varphi) \in \prod_\varphi F(U)$.
Considérons l'application
$$
\psi :
A_M = \bigoplus\nolimits_{\varphi \in \Hom_\mathcal{A}(U, M)} U
\longrightarrow
\bigoplus\nolimits_{s \in F(U)} U = A
$$
qui envoie le facteur correspondant à $\varphi$ sur celui correspondant à
$s_\varphi$ par l'application identique de $U$. Par construction, $s_{univ}$
a pour image $(s_\varphi)_\varphi$. Autrement dit, le carré de droite du
diagramme
$$
\xymatrix{
A' \ar[r] &
A \ar@{..>}[r]_{s_{univ}} & F \\
K \ar[r] \ar[u]^{?} & A_M \ar[u]^\psi \ar[r] &
M \ar@{..>}[u]_s
}
$$
est commutatif. Posons $K = \Ker(A_M \to M)$. Comme la restriction de $s$
à $K$ est nulle, on a $\psi(K) \subset A'$ par définition de $A'$. Il existe
donc un morphisme induit $M \to A/A'$. Cette construction fournit un inverse
de l'application $h_{A/A'}(M) \to F(M)$ (nous omettons les détails).
\end{proof}

\begin{lemma}
\label{injectives-lemma-grothendieck-products}
Toute catégorie abélienne de Grothendieck vérifie Ab3*.
\end{lemma}

\begin{proof}
Soit $M_i$, $i \in I$, une famille d'objets de $\mathcal{A}$ indexée par un
ensemble $I$. Le foncteur $F = \prod_{i \in I} h_{M_i}$ transforme les
colimites en limites. Le
lemme \ref{injectives-lemma-grothendieck-brown} s'applique donc.
\end{proof}

\begin{remark}
\label{injectives-remark-existence-D}
Dans le chapitre consacré aux catégories dérivées, nous travaillons
systématiquement avec des catégories abéliennes ``petites'' (conformément à
la convention du Champs project). Pour une catégorie abélienne ``grande''
$\mathcal{A}$, l'existence de la catégorie dérivée $D(\mathcal{A})$ n'est pas
évidente, car il n'est pas clair que les morphismes de la catégorie dérivée
forment des ensembles. En général, c'est faux ; voir
Exemples, lemme \ref{examples-lemma-big-abelian-category}.
En revanche, si $\mathcal{A}$ est une catégorie abélienne de Grothendieck et si
$K^\bullet, L^\bullet$ sont dans $K(\mathcal{A})$, alors le
théorème \ref{injectives-theorem-K-injective-embedding-grothendieck} donne un
quasi-isomorphisme $L^\bullet \to I^\bullet$ vers un complexe K-injectif
$I^\bullet$, et Catégories dérivées,
lemme \ref{derived-lemma-K-injective}, montre que
$$
\Hom_{D(\mathcal{A})}(K^\bullet, L^\bullet) =
\Hom_{K(\mathcal{A})}(K^\bullet, I^\bullet)
$$
qui est un ensemble. La catégorie des modules sur un anneau et, plus
généralement, celle des faisceaux de modules sur un site annelé sont des
exemples de catégories abéliennes de Grothendieck.
\end{remark}

\begin{lemma}
\label{injectives-lemma-derived-products}
Soit $\mathcal{A}$ une catégorie abélienne de Grothendieck. Alors :
\begin{enumerate}
\item $D(\mathcal{A})$ admet des sommes directes et des produits ;
\item les sommes directes s'obtiennent en prenant terme à terme les sommes
directes de complexes quelconques ;
\item les produits s'obtiennent en prenant terme à terme les produits de
complexes K-injectifs.
\end{enumerate}
\end{lemma}

\begin{proof}
Soit $K^\bullet_i$, $i \in I$, une famille d'objets de $D(\mathcal{A})$
indexée par un ensemble $I$. Nous affirmons que la somme directe terme à
terme $\bigoplus_{i \in I} K^\bullet_i$ est une somme directe dans
$D(\mathcal{A})$. Soit en effet $I^\bullet$ un complexe K-injectif. On a
\begin{align*}
\Hom_{D(\mathcal{A})}(\bigoplus\nolimits_{i \in I} K^\bullet_i, I^\bullet)
& =
\Hom_{K(\mathcal{A})}(\bigoplus\nolimits_{i \in I} K^\bullet_i, I^\bullet) \\
& =
\prod\nolimits_{i \in I} \Hom_{K(\mathcal{A})}(K^\bullet_i, I^\bullet) \\
& =
\prod\nolimits_{i \in I} \Hom_{D(\mathcal{A})}(K^\bullet_i, I^\bullet)
\end{align*}
comme souhaité. Cela suffit, puisque tout complexe peut être représenté par
un complexe K-injectif d'après le
théorème \ref{injectives-theorem-K-injective-embedding-grothendieck}.
Pour construire le produit, choisissons une résolution K-injective
$K_i^\bullet \to I_i^\bullet$ pour chaque $i$. Nous affirmons alors que
$\prod_{i \in I} I_i^\bullet$ est un produit dans $D(\mathcal{A})$.
Cela résulte de Catégories dérivées,
lemme \ref{derived-lemma-product-K-injective}.
\end{proof}

\begin{remark}
\label{injectives-remark-direct-sum-product-derived}
Soit $R$ un anneau. Supposons que $M_n$, $n \in \mathbf{Z}$, soient des
$R$-modules. Notons $E_n = M_n[-n] \in D(R)$. Nous affirmons que
$E = \bigoplus M_n[-n]$ est {\it à la fois} la somme directe et le produit des
objets $E_n$ dans $D(R)$. Pour la somme directe, il suffit de consulter la
démonstration du lemme \ref{injectives-lemma-derived-products}.
Pour le produit direct, choisissons des résolutions injectives
$M_n \to I_n^\bullet$. La démonstration du
lemme \ref{injectives-lemma-derived-products} donne
$$
\prod E_n = \prod I_n^\bullet[-n]
$$
dans $D(R)$. Comme les produits sont exacts dans $\text{Mod}_R$,
$\prod I_n^\bullet[-n]$ est quasi-isomorphe à $E$. Le même raisonnement vaut
plus généralement dans $D(\mathcal{A})$, lorsque $\mathcal{A}$ est une
catégorie abélienne de Grothendieck vérifiant Ab4*.
\end{remark}

\begin{lemma}
\label{injectives-lemma-RF-commutes-with-Rlim}
Soit $F : \mathcal{A} \to \mathcal{B}$ un foncteur additif entre catégories
abéliennes. Supposons que
\begin{enumerate}
\item $\mathcal{A}$ soit une catégorie abélienne de Grothendieck ;
\item les produits dénombrables soient exacts dans $\mathcal{B}$ ;
\item $F$ commute aux produits dénombrables.
\end{enumerate}
Alors $RF : D(\mathcal{A}) \to D(\mathcal{B})$ commute aux limites dérivées.
\end{lemma}

\begin{proof}
Remarquons que $RF$ existe, car $\mathcal{A}$ possède suffisamment de
complexes K-injectifs (théorème
\ref{injectives-theorem-K-injective-embedding-grothendieck} et
Catégories dérivées, lemme \ref{derived-lemma-K-injective-defined}).
L'énoncé signifie que, si $K = R\lim K_n$, alors
$RF(K) = R\lim RF(K_n)$. Pour les notations, voir
Catégories dérivées, définition \ref{derived-definition-derived-limit}.
Comme $RF$ est un foncteur exact entre catégories triangulées, il suffit de
vérifier que $RF$ commute aux produits dénombrables d'objets de
$D(\mathcal{A})$. La démonstration du
lemme \ref{injectives-lemma-derived-products} montre que les produits dans
$D(\mathcal{A})$ se calculent en prenant les produits de complexes
K-injectifs, et qu'un produit de complexes K-injectifs est K-injectif.
En outre, Catégories dérivées,
lemme \ref{derived-lemma-products}, montre que les produits dans
$D(\mathcal{B})$ se calculent terme à terme.
Or $RF$ se calcule en appliquant $F$ à un représentant K-injectif, et nous
avons supposé que $F$ commute aux produits dénombrables ; le résultat en
découle.
\end{proof}
\noindent
Le lemme suivant constitue une sorte de généralisation de
l'existence des résolutions de Cartan--Eilenberg
(Catégories dérivées, section \ref{derived-section-cartan-eilenberg}).

\begin{lemma}
\label{injectives-lemma-K-injective-embedding-filtration}
Soit $\mathcal{A}$ une catégorie abélienne de Grothendieck.
Soit $K^\bullet$ un complexe filtré de $\mathcal{A}$, voir
Homologie, définition \ref{homology-definition-filtered-complex}.
Il existe alors un morphisme $j : K^\bullet \to J^\bullet$
de complexes filtrés de $\mathcal{A}$ tel que
\begin{enumerate}
\item $J^n$, $F^pJ^n$, $J^n/F^pJ^n$ et $F^pJ^n/F^{p'}J^n$ sont des
objets injectifs de $\mathcal{A}$,
\item $J^\bullet$, $F^pJ^\bullet$, $J^\bullet/F^pJ^\bullet$ et
$F^pJ^\bullet/F^{p'}J^\bullet$ sont des complexes K-injectifs,
\item $j$ induit les quasi-isomorphismes
$K^\bullet \to J^\bullet$,
$F^pK^\bullet \to F^pJ^\bullet$,
$K^\bullet/F^pK^\bullet \to J^\bullet/F^pJ^\bullet$, et
$F^pK^\bullet/F^{p'}K^\bullet \to F^pJ^\bullet/F^{p'}J^\bullet$.
\end{enumerate}
\end{lemma}

\begin{proof}
D'après le théorème \ref{injectives-theorem-K-injective-embedding-grothendieck},
on obtient des quasi-isomorphismes $i : K^\bullet \to I^\bullet$ et
$i^p : F^pK^\bullet \to I^{p, \bullet}$ ainsi que les diagrammes commutatifs
$$
\vcenter{
\xymatrix{
K^\bullet \ar[d]_i & F^pK^\bullet \ar[l] \ar[d]_{i^p} \\
I^\bullet & I^{p, \bullet} \ar[l]_{\alpha^p}
}
}
\quad\text{et}\quad
\vcenter{
\xymatrix{
F^{p'}K^\bullet \ar[d]_{i^{p'}} &
F^pK^\bullet \ar[l] \ar[d]_{i^p} \\
I^{p', \bullet} &
I^{p, \bullet} \ar[l]_{\alpha^{p p'}}
}
}
\quad\text{pour }p' \leq p
$$
de sorte que $\alpha^{p'} \circ \alpha^{p p'} = \alpha^p$
et $\alpha^{p'p''} \circ \alpha^{pp'} = \alpha^{pp''}$.
Le problème est que les applications $\alpha^p : I^{p, \bullet} \to I^\bullet$
ne sont pas nécessairement injectives. Pour chaque $p$, choisissons une
injection $t^p : I^{p, \bullet} \to J^{p, \bullet}$ dans un complexe
K-injectif acyclique $J^{p, \bullet}$ dont les termes sont des objets
injectifs de $\mathcal{A}$ (on l'envoie d'abord dans le cône de l'identité,
puis on applique le théorème).
Choisissons un morphisme de complexes $s^p : I^\bullet \to J^{p, \bullet}$
tel que le diagramme suivant soit commutatif
$$
\xymatrix{
K^\bullet \ar[d]_i & F^pK^\bullet \ar[l] \ar[d]_{i^p} \\
I^\bullet \ar[rd]_{s^p} & I^{p, \bullet} \ar[d]^{t^p} \\
& J^{p, \bullet}
}
$$
C'est possible : la composée $F^pK^\bullet \to J^{p, \bullet}$
est homotope à zéro, car $J^{p, \bullet}$ est acyclique et K-injectif
(Catégories dérivées, lemme \ref{derived-lemma-K-injective}).
Puisque les objets $J^{p, n - 1}$ sont injectifs et que
$F^pK^n \to K^n \to I^n$ sont des morphismes injectifs, on peut
prolonger les applications $F^pK^n \to J^{p, n - 1}$ qui définissent
l'homotopie en une application $h^n : I^n \to J^{p, n - 1}$.
On pose alors $s^p$ égal à $h \circ \text{d} + \text{d} \circ h$.
(Attention : on n'aura pas $t^p = s^p \circ \alpha^p$ ; il faudra donc
veiller à ne pas utiliser cette égalité ci-dessous.)

\medskip\noindent
Considérons
$$
J^\bullet = I^\bullet \times \prod\nolimits_p J^{p, \bullet}
$$
Comme les produits dans $D(\mathcal{A})$ s'obtiennent en prenant
les produits de complexes K-injectifs
(lemme \ref{injectives-lemma-derived-products}) et que $J^{p, \bullet}$
est isomorphe à $0$ dans $D(\mathcal{A})$, le morphisme
$J^\bullet \to I^\bullet$ est un isomorphisme dans $D(\mathcal{A})$.
Considérons l'application
$$
j = i \times (s^p \circ i)_{p \in \mathbf{Z}} :
K^\bullet
\longrightarrow
I^\bullet \times \prod\nolimits_p J^{p, \bullet} = J^\bullet
$$
D'après les remarques précédentes, c'est un quasi-isomorphisme. Il est aussi
injectif. Pour $p \in \mathbf{Z}$, on définit $F^pJ^\bullet \subset J^\bullet$ par
$$
\Im\left(
\alpha^p \times (t^{p'} \circ \alpha^{pp'})_{p' \leq p} :
I^{p, \bullet}
\to
I^\bullet \times \prod\nolimits_{p' \leq p} J^{p', \bullet}
\right)
\times \prod\nolimits_{p' > p} J^{p', \bullet}
$$
Ce complexe est isomorphe au complexe
$I^{p, \bullet} \times \prod_{p' > p} J^{p', \bullet}$
car $\alpha^{pp} = \text{id}$ et $t^p$ est injective.
Par conséquent, $F^pJ^\bullet$ est quasi-isomorphe à $I^{p, \bullet}$
(on raisonne comme ci-dessus). La commutativité du diagramme précédent donne
$j(F^pK^\bullet) \subset F^pJ^\bullet$. Le morphisme de complexes correspondant
$F^pK^\bullet \to F^pJ^\bullet$ est donc un quasi-isomorphisme.
Enfin, pour voir que
$F^{p + 1}J^\bullet \subset F^pJ^\bullet$
on utilise l'égalité $\alpha^{pp'} \circ \alpha^{p + 1p} = \alpha^{p + 1p'}$
et la commutativité du premier diagramme affiché au premier paragraphe
de la démonstration.

\medskip\noindent
Affirmons que $j : K^\bullet \to J^\bullet$ résout le problème posé par le
lemme. En effet, $F^pJ^n$ est un objet injectif de $\mathcal{A}$ puisqu'il est
isomorphe à $I^{p, n} \times \prod_{p' > p} J^{p', n}$ et que les produits
d'injectifs sont injectifs. L'application injective $F^pJ^n \to J^n$ est alors
scindée ; le quotient $J^n/F^pJ^n$ est donc lui aussi injectif, comme facteur
direct de l'objet injectif $J^n$. Il en va de même de
$F^pJ^n/F^{p'}J^n$. En particulier,
$0 \to F^pJ^\bullet \to J^\bullet \to J^\bullet/F^pJ^\bullet \to 0$
est une suite exacte courte de complexes scindée terme à terme ; elle définit
donc, par convention, un triangle distingué dans $K(\mathcal{A})$.
Puisque $J^\bullet$ et $F^pJ^\bullet$ sont des complexes K-injectifs, il en va
de même de $J^\bullet/F^pJ^\bullet$ d'après Catégories dérivées, lemme
\ref{derived-lemma-triangle-K-injective}. Un raisonnement semblable montre que
$F^pJ^\bullet/F^{p'}J^\bullet$ est K-injectif. Par construction,
$j : K^\bullet \to J^\bullet$ et les applications induites
$F^pK^\bullet \to F^pJ^\bullet$ sont des quasi-isomorphismes. Les longues
suites exactes de cohomologie des complexes considérés montrent qu'il en va
de même pour
$K^\bullet/F^pK^\bullet \to J^\bullet/F^pJ^\bullet$ et
$F^pK^\bullet/F^{p'}K^\bullet \to F^pJ^\bullet/F^{p'}J^\bullet$.
\end{proof}

\begin{remark}
\label{injectives-remark-ext-into-filtered-complex}
Soit $\mathcal{A}$ une catégorie abélienne de Grothendieck.
Soit $K^\bullet$ un complexe filtré de $\mathcal{A}$, voir
Homologie, définition \ref{homology-definition-filtered-complex}.
Pour alléger les notations, désignons par $K$, $F^pK$, $\text{gr}^pK$
les objets de $D(\mathcal{A})$ représentés respectivement par $K^\bullet$,
$F^pK^\bullet$, $\text{gr}^pK^\bullet$. Soit $M \in D(\mathcal{A})$.
Le lemme \ref{injectives-lemma-K-injective-embedding-filtration} permet de construire
une suite spectrale $(E_r, d_r)_{r \geq 1}$ d'objets bigradués de
$\mathcal{A}$, où $d_r$ est de bidegré $(r, -r + 1)$ et
$$
E_1^{p, q} = \Ext^{p + q}(M, \text{gr}^pK)
$$
Si, pour tout $n$, on a
$$
\Ext^n(M, F^pK) = 0 \text{ pour } p \gg 0
\quad\text{et}\quad
\Ext^n(M, F^pK) = \Ext^n(M, K) \text{ pour } p \ll 0
$$
alors la suite spectrale est bornée et converge vers $\Ext^{p + q}(M, K)$.
En effet, choisissons un complexe quelconque $M^\bullet$ représentant $M$,
un morphisme $j : K^\bullet \to J^\bullet$ comme dans le lemme, et considérons
le complexe
$$
\Hom^\bullet(M^\bullet, J^\bullet)
$$
défini exactement comme dans Compléments d'algèbre, section
\ref{more-algebra-section-hom-complexes}. En posant
$F^p\Hom^\bullet(M^\bullet, J^\bullet) =
\Hom^\bullet(M^\bullet, F^pJ^\bullet)$, on obtient un complexe filtré.
La suite spectrale d'Homologie, section
\ref{homology-section-filtered-complex}, possède les différentielles et les
termes décrits ci-dessus ; nous omettons les détails. Le caractère borné et la
convergence résultent d'Homologie, lemme
\ref{homology-lemma-ss-converges-trivial}.
\end{remark}

\begin{remark}
\label{injectives-remark-spectral-sequences-ext}
Soit $\mathcal{A}$ une catégorie abélienne de Grothendieck.
Soient $M, K$ des objets de $D(\mathcal{A})$.
Pour tout choix d'un complexe $K^\bullet$ représentant $K$, on peut utiliser
la filtration $F^pK^\bullet = \tau_{\leq -p}K^\bullet$ et la discussion de la
remarque \ref{injectives-remark-ext-into-filtered-complex} afin d'obtenir une suite
spectrale telle que
$$
E_1^{p, q} = \Ext^{2p + q}(M, H^{-p}(K))
$$
Cette suite spectrale ne dépend pas du choix du complexe $K^\bullet$
représentant $K$. Après la renumérotation $p = -j$ et $q = i + 2j$, on
obtient une suite spectrale $(E'_r, d'_r)_{r \geq 2}$ où $d'_r$ est de
bidegré $(r, -r + 1)$ et
$$
(E'_2)^{i, j} = \Ext^i(M, H^j(K))
$$
Si $M \in D^-(\mathcal{A})$ et $K \in D^+(\mathcal{A})$, alors $E_r$ et
$E'_r$ sont toutes deux bornées et convergent vers $\Ext^{p + q}(M, K)$.
Si l'on utilise la filtration $F^pK^\bullet = \sigma_{\geq p}K^\bullet$,
on obtient
$$
E_1^{p, q} = \Ext^q(M, K^p)
$$
Si $M \in D^-(\mathcal{A})$ et si $K^\bullet$ est borné inférieurement,
cette suite spectrale est bornée et converge vers $\Ext^{p + q}(M, K)$.
\end{remark}

\begin{remark}
\label{injectives-remark-ext-from-filtered-complex}
Soit $\mathcal{A}$ une catégorie abélienne de Grothendieck. Soit
$K \in D(\mathcal{A})$. Soit $M^\bullet$ un complexe filtré de
$\mathcal{A}$, voir Homologie, définition
\ref{homology-definition-filtered-complex}. Pour alléger les notations,
désignons par $M$, $M/F^pM$, $\text{gr}^pM$ les objets de
$D(\mathcal{A})$ représentés respectivement par $M^\bullet$,
$M^\bullet/F^pM^\bullet$, $\text{gr}^pM^\bullet$. Par dualité avec la
remarque \ref{injectives-remark-ext-into-filtered-complex}, on peut construire une suite
spectrale $(E_r, d_r)_{r \geq 1}$ d'objets bigradués de $\mathcal{A}$, où
$d_r$ est de bidegré $(r, -r + 1)$ et
$$
E_1^{p, q} = \Ext^{p + q}(\text{gr}^{-p}M, K)
$$
Si, pour tout $n$, on a
$$
\Ext^n(M/F^pM, K) = 0 \text{ pour } p \ll 0
\quad\text{et}\quad
\Ext^n(M/F^pM, K) = \Ext^n(M, K) \text{ pour } p \gg 0
$$
alors la suite spectrale est bornée et converge vers $\Ext^{p + q}(M, K)$.
En effet, choisissons un complexe K-injectif $I^\bullet$ à termes injectifs
représentant $K$, voir le théorème
\ref{injectives-theorem-K-injective-embedding-grothendieck}. Considérons le complexe
$$
\Hom^\bullet(M^\bullet, I^\bullet)
$$
défini exactement comme dans Compléments d'algèbre, section
\ref{more-algebra-section-hom-complexes}. En posant
$$
F^p\Hom^\bullet(M^\bullet, I^\bullet) =
\Hom^\bullet(M^\bullet/F^{-p + 1}M^\bullet, I^\bullet)
$$
on obtient un complexe filtré (remarquer le signe et le décalage de la
filtration). La suite spectrale d'Homologie, section
\ref{homology-section-filtered-complex}, possède les différentielles et les
termes décrits ci-dessus ; nous omettons les détails. Le caractère borné et la
convergence résultent d'Homologie, lemme
\ref{homology-lemma-ss-converges-trivial}.
\end{remark}

\begin{remark}
\label{injectives-remark-spectral-sequences-ext-variant}
Soit $\mathcal{A}$ une catégorie abélienne de Grothendieck.
Soient $M, K$ des objets de $D(\mathcal{A})$.
Pour tout choix d'un complexe $M^\bullet$ représentant $M$, on peut utiliser
la filtration $F^pM^\bullet = \tau_{\leq -p}M^\bullet$ et la discussion de la
remarque \ref{injectives-remark-ext-into-filtered-complex} afin d'obtenir une suite
spectrale telle que
$$
E_1^{p, q} = \Ext^{2p + q}(H^p(M), K)
$$
Cette suite spectrale ne dépend pas du choix du complexe $M^\bullet$
représentant $M$. Après la renumérotation $p = -j$ et $q = i + 2j$, on
obtient une suite spectrale $(E'_r, d'_r)_{r \geq 2}$ où $d'_r$ est de
bidegré $(r, -r + 1)$ et
$$
(E'_2)^{i, j} = \Ext^i(H^{-j}(M), K)
$$
Si $M \in D^-(\mathcal{A})$ et $K \in D^+(\mathcal{A})$, alors $E_r$ et
$E'_r$ sont bornées et convergent vers $\Ext^{p + q}(M, K)$.
Si l'on utilise la filtration $F^pM^\bullet = \sigma_{\geq p}M^\bullet$,
on obtient
$$
E_1^{p, q} = \Ext^q(M^{-p}, K)
$$
Si $K \in D^+(\mathcal{A})$ et si $M^\bullet$ est borné supérieurement,
cette suite spectrale est bornée et converge vers $\Ext^{p + q}(M, K)$.
\end{remark}

\begin{lemma}
\label{injectives-lemma-represent-by-filtered-complex}
Soit $\mathcal{A}$ une catégorie abélienne de Grothendieck. Supposons donnés
un objet $E \in D(\mathcal{A})$, un système projectif
$\{E^i\}_{i \in \mathbf{Z}}$ d'objets de $D(\mathcal{A})$ indexé par
$\mathbf{Z}$, ainsi qu'un système compatible de morphismes $E^i \to E$.
On a le diagramme suivant :
$$
\ldots \to E^{i + 1} \to E^i \to E^{i - 1} \to \ldots \to E
$$
Il existe alors un complexe filtré $K^\bullet$ de $\mathcal{A}$
(Homologie, définition \ref{homology-definition-filtered-complex}) tel que
$K^\bullet$ représente $E$ et que $F^iK^\bullet$ représente $E^i$, de façon
compatible avec les morphismes donnés.
\end{lemma}

\begin{proof}
D'après le théorème \ref{injectives-theorem-K-injective-embedding-grothendieck}, on peut
choisir un complexe K-injectif $I^\bullet$ représentant $E$, dont tous les
termes $I^n$ sont des objets injectifs de $\mathcal{A}$.
Choisissons un complexe $G^{0, \bullet}$ représentant $E^0$, puis un
morphisme de complexes $\varphi^0 : G^{0, \bullet} \to I^\bullet$
représentant $E^0 \to E$.
Pour $i > 0$, représentons par récurrence $E^i \to E^{i - 1}$ par un
morphisme de complexes
$\delta : G^{i, \bullet} \to G^{i - 1, \bullet}$
et posons $\varphi^i = \varphi^{i - 1} \circ \delta$.
Pour $i < 0$, représentons par récurrence $E^{i + 1} \to E^i$ par un
morphisme de complexes injectif terme à terme
$\delta : G^{i + 1, \bullet} \to G^{i, \bullet}$
(on peut, par exemple, utiliser Catégories dérivées, lemme
\ref{derived-lemma-make-injective}).
Affirmons qu'il existe un morphisme de complexes
$\varphi^i : G^{i, \bullet} \to I^\bullet$ représentant le morphisme
$E^i \to E$ et s'insérant dans le diagramme commutatif
$$
\xymatrix{
G^{i + 1, \bullet} \ar[r]_\delta \ar[d]_{\varphi^{i + 1}} &
G^{i, \bullet} \ar[ld]^{\varphi^i} \\
I^\bullet
}
$$
En effet, choisissons d'abord un morphisme de complexes quelconque
$\varphi : G^{i, \bullet} \to I^\bullet$ représentant le morphisme
$E^i \to E$. Les morphismes $\varphi \circ \delta$ et
$\varphi^{i + 1}$ sont alors homotopes par une homotopie
$h^p : G^{i + 1, p} \to I^{p - 1}$. Comme les termes de $I^\bullet$ sont
injectifs et que $\delta$ est injectif terme à terme, on peut prolonger
$h^p$ en $(h')^p : G^{i, p} \to I^{p - 1}$. En posant
$\varphi^i = \varphi + h' \circ d + d \circ h'$, on obtient ce qui était
annoncé.

\medskip\noindent
Choisissons ensuite, pour tout $i$, un morphisme de complexes injectif terme
à terme $a^i : G^{i, \bullet} \to J^{i, \bullet}$, où $J^{i, \bullet}$ est
acyclique et K-injectif et où les $J^{i, p}$ sont des objets injectifs de
$\mathcal{A}$. Pour cela, on envoie d'abord $G^{i, \bullet}$ dans le cône de
l'identité, puis on applique le théorème cité ci-dessus.
En raisonnant comme précédemment, on trouve des morphismes de complexes
$\delta' : J^{i, \bullet} \to J^{i - 1, \bullet}$ tels que les diagrammes
$$
\xymatrix{
G^{i, \bullet} \ar[r]_\delta \ar[d]_{a^i} &
G^{i - 1, \bullet} \ar[d]^{a^{i - 1}} \\
J^{i, \bullet} \ar[r]^{\delta'} &
J^{i - 1, \bullet}
}
$$
soient commutatifs. (On pourrait aussi utiliser la fonctorialité des cônes et
celle du théorème.) Considérons alors les morphismes
$$
\xymatrix{
G^{i + 1, \bullet} \times \prod\nolimits_{p > i + 1} J^{p, \bullet}
\ar[r] \ar[rd] &
G^{i, \bullet} \times \prod\nolimits_{p > i} J^{p, \bullet}
\ar[r] \ar[d] &
G^{i - 1, \bullet} \times \prod\nolimits_{p > i - 1} J^{p, \bullet}
\ar[ld] \\
& I^\bullet \times \prod\nolimits_p J^{p, \bullet}
}
$$
Ici, les flèches portant sur $J^{p, \bullet}$ sont les flèches évidentes
(l'identité ou le morphisme nul). Sur le facteur $G^{i, \bullet}$, on utilise
$\delta : G^{i, \bullet} \to G^{i - 1, \bullet}$, le morphisme
$\varphi^i : G^{i, \bullet} \to I^\bullet$, le morphisme nul
$0 : G^{i, \bullet} \to J^{p, \bullet}$ pour $p > i$, le morphisme
$a^i : G^{i, \bullet} \to J^{p, \bullet}$ pour $p = i$, et
$(\delta')^{i - p} \circ a^i = a^p \circ \delta^{i - p} :
G^{i, \bullet} \to J^{p, \bullet}$ pour $p < i$.
Nous omettons de vérifier que toutes les flèches du diagramme sont injectives
terme à terme. On obtient ainsi un complexe filtré. Comme les produits dans
$D(\mathcal{A})$ s'obtiennent en prenant les produits de complexes K-injectifs
(lemme \ref{injectives-lemma-derived-products}) et que $J^{p, \bullet}$ est nul dans
$D(\mathcal{A})$, ce diagramme représente le diagramme donné dans la catégorie
dérivée. Cela achève la démonstration.
\end{proof}

\begin{lemma}
\label{injectives-lemma-represent-by-filtered-complex-bis}
Dans la situation du lemme \ref{injectives-lemma-represent-by-filtered-complex},
supposons donné un second système projectif
$\{(E')^i\}_{i \in \mathbf{Z}}$, ainsi qu'un système compatible de morphismes
$(E')^i \to E$. Il existe alors un complexe bifiltré $K^\bullet$ de
$\mathcal{A}$ tel que $K^\bullet$ représente $E$, que $F^iK^\bullet$
représente $E^i$ et que $(F')^iK^\bullet$ représente $(E')^i$, de façon
compatible avec les morphismes donnés.
\end{lemma}

\begin{proof}
Le lemme permet d'abord de choisir $K^\bullet$ et $F$. On choisit ensuite
$(K')^\bullet$ et $F'$ convenant à $\{(E')^i\}_{i \in \mathbf{Z}}$ et aux
morphismes $(E')^i \to E$. D'après le lemme
\ref{injectives-lemma-K-injective-embedding-filtration}, on peut supposer que
$K^\bullet$ est un complexe K-injectif. On choisit alors un morphisme de
complexes $(K')^\bullet \to K^\bullet$ correspondant aux identifications
données $(K')^\bullet \cong E \cong K^\bullet$. On peut en outre choisir un
morphisme $(K')^\bullet \to J^\bullet$ injectif terme à terme, où
$J^\bullet$ est acyclique et K-injectif. (Pour cela, on envoie d'abord
$(K')^\bullet$ dans le cône de l'identité, puis on applique le théorème
\ref{injectives-theorem-K-injective-embedding-grothendieck}.) Les morphismes
$(K')^\bullet \to K^\bullet \times J^\bullet$ et
$K^\bullet \to K^\bullet \times J^\bullet$ sont alors tous deux injectifs
terme à terme et sont des quasi-isomorphismes (car le produit représente $E$
d'après le lemme \ref{injectives-lemma-derived-products}). Il suffit enfin de prendre les
images des filtrations de $K^\bullet$ et $(K')^\bullet$ par ces morphismes.
\end{proof}





\section{Le théorème de Gabriel--Popescu}
\label{injectives-section-gabriel-popescu}

\noindent
Dans cette section, nous étudions le théorème principal de \cite{GP}. La
méthode de démonstration suit des notes de Jacob Lurie et d'Akhil Mathew, qui
suivent à leur tour la présentation de Kuhn dans \cite{Kuhn}. Voir aussi
\cite{Takeuchi}.

\medskip\noindent
Soit $\mathcal{A}$ une catégorie abélienne de Grothendieck et soit $U$ un
générateur de $\mathcal{A}$, voir la définition
\ref{injectives-definition-grothendieck-conditions}. Posons
$R = \Hom_\mathcal{A}(U, U)$. Considérons le foncteur
$G : \mathcal{A} \to \text{Mod}_R$ défini par
$$
G(A) = \Hom_\mathcal{A}(U, A)
$$
muni de sa structure canonique de $R$-module à droite.

\begin{lemma}
\label{injectives-lemma-gabriel-popescu-left-adjoint}
Le foncteur $G$ ci-dessus possède un adjoint à gauche
$F : \text{Mod}_R \to \mathcal{A}$.
\end{lemma}

\begin{proof}
Nous donnons deux démonstrations de ce lemme.

\medskip\noindent
La première utilise le théorème du foncteur adjoint, voir Catégories,
théorème \ref{categories-theorem-adjoint-functor}. Remarquons que
$G : \mathcal{A} \to \text{Mod}_R$ est exact à gauche et envoie les produits
sur les produits. Le foncteur $G$ commute donc aux limites. Pour vérifier la condition
ensembliste du théorème, soit $M$ un objet de $\text{Mod}_R$. Choisissons un
cardinal $\kappa$ suffisamment grand et désignons par $E$ un ensemble d'objets
de $\mathcal{A}$ tel que tout objet $A$ vérifiant $|A| \leq \kappa$ soit
isomorphe à un élément de $E$. C'est possible d'après le lemme
\ref{injectives-lemma-set-iso-classes-bounded-size}. Posons
$I = \coprod_{A \in E} \Hom_R(M, G(A))$. Un élément $i \in I$ sera considéré
comme un couple $(A_i, f_i)$. Enfin, soient $A$ un objet quelconque de
$\mathcal{A}$ et $f : M \to G(A)$ un morphisme quelconque. Interprétons les
éléments de $\Im(f) \subset G(A) = \Hom_\mathcal{A}(U, A)$ comme des
morphismes $u : U \to A$. Posons
$$
A' = \Im(\bigoplus\nolimits_{u \in \Im(f)} U \xrightarrow{u} A)
$$
Puisque $G$ est exact à gauche, $G(A') \subset G(A)$ contient $\Im(f)$ ; on
obtient donc $f' : M \to G(A')$ par lequel $f$ se factorise. D'autre part,
l'objet $A'$ est quotient d'une somme directe d'au plus $|M|$ copies de $U$.
Ainsi, si $\kappa = |\bigoplus_{|M|} U|$, le couple $(A', f')$ est isomorphe
à un élément $(A_i, f_i)$ de $I$, et $f$ se factorise, comme souhaité, sous
la forme $M \xrightarrow{f_i} G(A_i) \to G(A)$.

\medskip\noindent
La seconde démonstration construit $F$ et montre, en un certain sens, que
``$F(M) = M \otimes_R U$''. En effet, pour tout $R$-module $M$, on peut
choisir une présentation
$$
\bigoplus\nolimits_{j \in J} R \to
\bigoplus\nolimits_{i \in I} R \to
M \to 0
$$
On définit alors $F(M)$ par la suite exacte correspondante
$$
\bigoplus\nolimits_{j \in J} U \to
\bigoplus\nolimits_{i \in I} U \to
F(M) \to 0
$$
Cette construction ne dépend pas du choix de la présentation et elle est
fonctorielle ; nous omettons les détails. Pour tout $A$ dans $\mathcal{A}$,
on obtient une suite exacte
$$
0 \to \Hom_\mathcal{A}(F(M), A) \to
\prod\nolimits_{i \in I} G(A) \to
\prod\nolimits_{j \in J} G(A)
$$
qui est isomorphe à la suite
$$
0 \to \Hom_R(M, G(A)) \to
\Hom_R(\bigoplus\nolimits_{i \in I} R, G(A)) \to
\Hom_R(\bigoplus\nolimits_{j \in J} R, G(A))
$$
ce qui montre que $F$ est l'adjoint à gauche de $G$.
\end{proof}

\begin{lemma}
\label{injectives-lemma-F-G-monos}
Soit $f : M \to G(A)$ un morphisme injectif dans $\text{Mod}_R$.
Alors le morphisme adjoint $f' : F(M) \to A$ est lui aussi injectif.
\end{lemma}

\begin{proof}
Choisissons un morphisme $R^{\oplus n} \to M$ et considérons le morphisme
correspondant $U^{\oplus n} \to F(M)$. Soit
$v : U \to U^{\oplus n}$ un morphisme tel que la composée
$U \to U^{\oplus n} \to F(M) \to A$ soit $0$. La flèche
$v : U \to U^{\oplus n}$ est alors un élément $v$ de $R^{\oplus n}$ dont
l'image dans $G(A)$ est nulle. Puisque $f$ est injectif, l'image de $v$ dans
$M$ est nulle ; cela signifie que $U \to U^{\oplus n} \to F(M)$ est nul,
par la construction de $F(M)$ donnée dans la démonstration du
lemme \ref{injectives-lemma-gabriel-popescu-left-adjoint}. Comme $U$ est un générateur, on en
déduit que
$$
\Ker(U^{\oplus n} \to F(M) \to A) = \Ker(U^{\oplus n} \to F(M))
$$
Pour achever la démonstration, choisissons une surjection
$\bigoplus_{i \in I} R \to M$ et considérons la surjection correspondante
$$
\pi : \bigoplus\nolimits_{i \in I} U \longrightarrow F(M)
$$
Pour montrer que $f'$ est injectif, il suffit d'établir l'égalité
$\Ker(\pi) = \Ker(f' \circ \pi)$ entre sous-objets de
$\bigoplus_{i \in I} U$. Or on peut écrire $\bigoplus_{i \in I} U$ comme la
colimite filtrante de ses sous-objets $\bigoplus_{i \in I'} U$, où
$I' \subset I$ parcourt les sous-ensembles finis. Comme les colimites
filtrantes sont exactes dans $\mathcal{A}$ par AB5, on a
$$
\Ker(\pi) =
\colim_{I' \subset I\text{ fini}}
\left(\bigoplus\nolimits_{i \in I'} U\right)
\bigcap \Ker(\pi)
$$
et
$$
\Ker(f' \circ \pi) =
\colim_{I' \subset I\text{ fini}}
\left(\bigoplus\nolimits_{i \in I'} U\right)
\bigcap \Ker(f' \circ \pi)
$$
et l'on obtient l'égalité, car la première égalité affichée ci-dessus
l'établit pour chaque $I'$.
\end{proof}

\begin{theorem}
\label{injectives-theorem-gabriel-popescu}
Soit $\mathcal{A}$ une catégorie abélienne de Grothendieck. Il existe alors
un anneau (non commutatif) $R$ et des foncteurs
$G : \mathcal{A} \to \text{Mod}_R$ et
$F : \text{Mod}_R \to \mathcal{A}$ tels que
\begin{enumerate}
\item $F$ est l'adjoint à gauche de $G$,
\item $G$ est pleinement fidèle, et
\item $F$ est exact.
\end{enumerate}
Ces foncteurs sont en outre ceux construits ci-dessus.
\end{theorem}

\begin{proof}
Montrons d'abord que $G$ est pleinement fidèle, ou, ce qui revient au même,
que $F \circ G \to \text{id}$ est un isomorphisme, voir Catégories, lemme
\ref{categories-lemma-adjoint-fully-faithful}. Pour tout objet $A$, le
morphisme $F(G(A)) \to A$ est surjectif, car, par construction, tout morphisme
$U \to A$ se factorise par $F(G(A))$. D'autre part, le morphisme
$F(G(A)) \to A$ est l'adjoint du morphisme
$\text{id} : G(A) \to G(A)$ ; il est donc injectif d'après le lemme
\ref{injectives-lemma-F-G-monos}.

\medskip\noindent
Le foncteur $F$ est exact à droite puisqu'il est adjoint à gauche.
Comme $\text{Mod}_R$ possède assez de projectifs, pour montrer que $F$ est
exact, il suffit de montrer que son premier foncteur dérivé gauche $L_1F$ est
nul. Pour établir $L_1F(M) = 0$ pour un $R$-module $M$, choisissons une suite
exacte $0 \to K \to P \to M \to 0$ de $R$-modules où $P$ est libre. Il suffit
de montrer que $F(K) \to F(P)$ est injectif. Écrivons cette suite comme
colimite filtrante de suites $0 \to K_i \to P_i \to M_i \to 0$ où $P_i$ est
un $R$-module libre de type fini : il suffit d'écrire ainsi $P$, puis de poser
$K_i = K \cap P_i$ et $M_i = \Im(P_i \to M)$. Puisque $F$ est adjoint à
gauche, il commute aux colimites ; et puisque $\mathcal{A}$ est une catégorie
abélienne de Grothendieck, $F(K) \to F(P)$ est injectif dès que chaque
$F(K_i) \to F(P_i)$ l'est. Il suffit donc de vérifier que
$F(K) \to F(P)$ est injectif lorsque $K \subset P = R^{\oplus n}$.
Le lemme \ref{injectives-lemma-F-G-monos} montre précisément que
$F(K) \to U^{\oplus n}$ est injectif.
\end{proof}

\begin{lemma}
\label{injectives-lemma-gabriel-popescu}
\begin{reference}
\cite[Corollary 4.1]{serpe}
\end{reference}
Soit $\mathcal{A}$ une catégorie abélienne de Grothendieck. Soient
$R$, $F$, $G$ comme dans le théorème de Gabriel--Popescu
(théorème \ref{injectives-theorem-gabriel-popescu}). On obtient alors les foncteurs
dérivés
$$
RG : D(\mathcal{A}) \to D(\text{Mod}_R)
\quad\text{et}\quad
F : D(\text{Mod}_R) \to D(\mathcal{A})
$$
tels que $F$ est adjoint à gauche de $RG$, que $RG$ est pleinement fidèle et
que $F \circ RG = \text{id}$.
\end{lemma}

\begin{proof}
L'existence et l'adjonction de ces foncteurs résultent des théorèmes
\ref{injectives-theorem-gabriel-popescu} et
\ref{injectives-theorem-K-injective-embedding-grothendieck}
et
des lemmes de Catégories dérivées
\ref{derived-lemma-K-injective-defined},
\ref{derived-lemma-right-derived-exact-functor}, et
\ref{derived-lemma-derived-adjoint-functors}.
L'égalité $F \circ RG = \text{id}$ vient du fait qu'on peut calculer $RG$ sur
un objet de $D(\mathcal{A})$ en appliquant $G$ à un complexe représentant
convenable $I^\bullet$ (par exemple un complexe K-injectif) ; on a alors
$F(G(I^\bullet)) = I^\bullet$ puisque $F \circ G = \text{id}$. La pleine
fidélité de $RG$ en résulte d'après Catégories, lemme
\ref{categories-lemma-adjoint-fully-faithful}.
\end{proof}





\section{Représentabilité de Brown et catégories abéliennes de Grothendieck}
\label{injectives-section-brown}

\noindent
Dans cette section, nous démontrons rapidement un théorème de représentabilité
pour les catégories dérivées de catégories abéliennes de Grothendieck. Le
lecteur devrait d'abord étudier le cas des catégories triangulées compactement
engendrées dans Catégories dérivées, section \ref{derived-section-brown}.
Il est ensuite tout indiqué de consulter, plutôt que cette section, la
littérature pour des résultats plus généraux de ce type ; voir par exemple
\cite{Franke}, \cite{Neeman}, \cite{Krause}, ou Catégories dérivées, section
\ref{derived-section-brown-bis}.

\begin{lemma}
\label{injectives-lemma-brown}
Soit $\mathcal{A}$ une catégorie abélienne de Grothendieck.
Soit $H : D(\mathcal{A}) \to \textit{Ab}$ un foncteur cohomologique
contravariant qui transforme les sommes directes en produits.
Alors $H$ est représentable.
\end{lemma}

\begin{proof}
Soient $R, F, G, RG$ comme dans le lemme \ref{injectives-lemma-gabriel-popescu}, et
considérons le foncteur $H \circ F : D(\text{Mod}_R) \to \textit{Ab}$.
Comme $F$ est adjoint à gauche, il envoie les sommes directes sur les sommes
directes ; $H \circ F$ transforme donc les sommes directes en produits.
D'autre part, la catégorie dérivée $D(\text{Mod}_R)$ est engendrée par un
unique objet compact, à savoir $R$. D'après Catégories dérivées, lemme
\ref{derived-lemma-brown}, le foncteur $H \circ F$ est représentable, disons
par $L \in D(\text{Mod}_R)$. Choisissons un triangle distingué
$$
M \to L \to RG(F(L)) \to M[1]
$$
dans $D(\text{Mod}_R)$. Alors $F(M) = 0$ puisque
$F \circ RG = \text{id}$. Par conséquent $H(F(M)) = 0$, puis
$\Hom(M, L) = 0$. Il s'ensuit que $L \to RG(F(L))$ est l'inclusion d'un
facteur direct, voir Catégories dérivées, lemme \ref{derived-lemma-split}.
Pour $A$ dans $D(\mathcal{A})$, on obtient
\begin{align*}
H(A)
& =
H(F(RG(A))) \\
& =
\Hom(RG(A), L) \\
& \to
\Hom(RG(A), RG(F(L))) \\
& =
\Hom(F(RG(A)), F(L)) \\
& =
\Hom(A, F(L))
\end{align*}
où la flèche possède un inverse à gauche fonctoriel en $A$. Autrement dit,
$H$ est facteur direct d'un foncteur représentable. Comme $D(\mathcal{A})$
est karoubienne (Catégories dérivées, lemme
\ref{derived-lemma-projectors-have-images-triangulated}), on conclut.
\end{proof}

\begin{proposition}
\label{injectives-proposition-brown}
Soit $\mathcal{A}$ une catégorie abélienne de Grothendieck. Soit $\mathcal{D}$
une catégorie triangulée. Soit $F : D(\mathcal{A}) \to \mathcal{D}$ un
foncteur exact de catégories triangulées qui transforme les sommes directes en
sommes directes. Alors $F$ possède un adjoint à droite exact.
\end{proposition}

\begin{proof}
Pour un objet $Y$ de $\mathcal{D}$, considérons le foncteur contravariant
$$
D(\mathcal{A}) \to \textit{Ab},\quad
W \mapsto \Hom_\mathcal{D}(F(W), Y)
$$
C'est un foncteur cohomologique puisque $F$ est exact, et il transforme les
sommes directes en produits puisque $F$ transforme les sommes directes en
sommes directes. Le lemme \ref{injectives-lemma-brown} fournit donc un objet $X$ de
$D(\mathcal{A})$ tel que
$\Hom_{D(\mathcal{A})}(W, X) = \Hom_\mathcal{D}(F(W), Y)$.
L'existence de l'adjoint résulte de Catégories, lemme
\ref{categories-lemma-adjoint-exists}. Son exactitude résulte de Catégories
dérivées, lemme \ref{derived-lemma-adjoint-is-exact}.
\end{proof}
















% OK, start here.
%

\chapter{Cohomologie des faisceaux}



\phantomsection
\label{cohomology-section-phantom}


\section{Introduction}

\label{cohomology-section-introduction}

\noindent
Dans ce document, nous étudions quelques thèmes de cohomologie des faisceaux sur les espaces topologiques. Nous travaillons principalement dans le cadre général des modules sur un faisceau d'anneaux et des morphismes d'espaces annelés. Pour le cas des faisceaux sur les sites, voir le chapitre Cohomologie sur les sites, section \ref{sites-cohomology-section-introduction}. Les références de base sont \cite{Godement} et \cite{Iversen}.





\section{Cohomologie des faisceaux}

\label{cohomology-section-cohomology-sheaves}

\noindent

Soit $X$ un espace topologique. Soit $\mathcal{F}$ un faisceau abélien. La catégorie des faisceaux abéliens sur $X$ possède suffisamment d'objets injectifs, voir Injectifs, lemme \ref{injectives-lemma-abelian-sheaves-space}. Nous pouvons donc choisir une résolution injective
$\mathcal{F}[0] \to \mathcal{I}^\bullet$. Comme il est d'usage, nous définissons

\begin{equation}
\label{cohomology-equation-cohomology}
H^i(X, \mathcal{F}) = H^i(\Gamma(X, \mathcal{I}^\bullet))
\end{equation}

le membre de gauche comme le {\it $i$-ième groupe de cohomologie du faisceau abélien $\mathcal{F}$}. La famille des foncteurs $H^i(X, -)$ forme un $\delta$-foncteur universel de $\textit{Ab}(X) \to \textit{Ab}$.

\medskip\noindent

Soit $f : X \to Y$ une application continue d'espaces topologiques. Avec
$\mathcal{F}[0] \to \mathcal{I}^\bullet$ comme ci-dessus, nous définissons

\begin{equation}
\label{cohomology-equation-higher-direct-image}
R^if_*\mathcal{F} = H^i(f_*\mathcal{I}^\bullet)
\end{equation}

le membre de gauche comme la {\it $i$-ième image directe supérieure de $\mathcal{F}$}. La famille des foncteurs $R^if_*$ forme un $\delta$-foncteur universel de $\textit{Ab}(X) \to \textit{Ab}(Y)$.

\medskip\noindent

Soit $(X, \mathcal{O}_X)$ un espace annelé. Soit $\mathcal{F}$ un
$\mathcal{O}_X$-module. La catégorie des $\mathcal{O}_X$-modules sur $X$ possède suffisamment d'objets injectifs, voir Injectifs, lemme \ref{injectives-lemma-sheaves-modules-space}. Nous pouvons donc choisir une résolution injective
$\mathcal{F}[0] \to \mathcal{I}^\bullet$. Comme il est d'usage, nous définissons

\begin{equation}
\label{cohomology-equation-cohomology-modules}
H^i(X, \mathcal{F}) = H^i(\Gamma(X, \mathcal{I}^\bullet))
\end{equation}

le membre de gauche comme le {\it $i$-ième groupe de cohomologie de $\mathcal{F}$}. La famille des foncteurs $H^i(X, -)$ forme un $\delta$-foncteur universel de $\textit{Mod}(\mathcal{O}_X) \to \text{Mod}_{\mathcal{O}_X(X)}$.

\medskip\noindent

Soit $f : (X, \mathcal{O}_X) \to (Y, \mathcal{O}_Y)$ un morphisme d'espaces annelés. Avec $\mathcal{F}[0] \to \mathcal{I}^\bullet$ comme ci-dessus, nous définissons

\begin{equation}
\label{cohomology-equation-higher-direct-image-modules}
R^if_*\mathcal{F} = H^i(f_*\mathcal{I}^\bullet)
\end{equation}

le membre de gauche comme la {\it $i$-ième image directe supérieure de $\mathcal{F}$}. La famille des foncteurs $R^if_*$ forme un $\delta$-foncteur universel de $\textit{Mod}(\mathcal{O}_X) \to \textit{Mod}(\mathcal{O}_Y)$.





\section{Foncteurs dérivés}

\label{cohomology-section-derived-functors}

\noindent
Nous expliquons brièvement comment obtenir les foncteurs dérivés à droite à l'aide de foncteurs de résolution. Pour les foncteurs dérivés non bornés, voir la section \ref{cohomology-section-unbounded}.

\medskip\noindent

Soit $(X, \mathcal{O}_X)$ un espace annelé. La catégorie
$\textit{Mod}(\mathcal{O}_X)$ est abélienne, voir Faisceaux de modules, lemme \ref{modules-lemma-abelian}. Dans ce chapitre, nous écrirons
$$
K(\mathcal{O}_X) = K(\textit{Mod}(\mathcal{O}_X))
\quad
\text{et}
\quad
D(\mathcal{O}_X) = D(\textit{Mod}(\mathcal{O}_X)).
$$
et de même pour les versions bornées des catégories triangulées introduites dans Catégories dérivées, définition \ref{derived-definition-complexes-notation} et définition \ref{derived-definition-unbounded-derived-category}. D'après Catégories dérivées, remarque \ref{derived-remark-big-abelian-category},
il existe un foncteur de résolution
$$
j = j_X :
K^{+}(\textit{Mod}(\mathcal{O}_X))
\longrightarrow
K^{+}(\mathcal{I})
$$
où $\mathcal{I}$ est la sous-catégorie additive strictement pleine de
$\textit{Mod}(\mathcal{O}_X)$ formée des faisceaux injectifs. Pour tout foncteur exact à gauche
$F : \textit{Mod}(\mathcal{O}_X) \to \mathcal{B}$
à valeurs dans une catégorie abélienne $\mathcal{B}$, nous noterons $RF$ le foncteur dérivé à droite décrit dans Catégories dérivées, section \ref{derived-section-right-derived-functor},
et construit à l'aide du foncteur de résolution $j_X$ que nous venons de décrire :

\begin{equation}
\label{cohomology-equation-RF}
RF = F \circ j_X' : D^{+}(X) \longrightarrow D^{+}(\mathcal{B})
\end{equation}

Voir Catégories dérivées, lemme \ref{derived-lemma-right-derived-functor},
pour cette notation. Remarquons que l'on peut considérer $RF$ comme défini sur
$\textit{Mod}(\mathcal{O}_X)$,
$\text{Comp}^{+}(\textit{Mod}(\mathcal{O}_X))$,
$K^{+}(X)$ ou $D^{+}(X)$,
selon la situation. D'après Catégories dérivées, définition \ref{derived-definition-higher-derived-functors},
nous obtenons le $i$-ième foncteur dérivé à droite

\begin{equation}
\label{cohomology-equation-RFi}
R^iF = H^i \circ RF : \textit{Mod}(\mathcal{O}_X) \longrightarrow \mathcal{B}
\end{equation}

de sorte que $R^0F = F$ et que $\{R^iF, \delta\}_{i \geq 0}$ est un
$\delta$-foncteur universel, voir Catégories dérivées, lemme \ref{derived-lemma-higher-derived-functors}.

\medskip\noindent

Voici deux cas particuliers de cette construction. Pour un anneau $R$, nous écrivons $K(R) = K(\text{Mod}_R)$ et
$D(R) = D(\text{Mod}_R)$, et de même pour les versions bornées. Pour tout ouvert $U \subset X$, nous avons un foncteur exact à gauche
$
\Gamma(U, -) :
\textit{Mod}(\mathcal{O}_X)
\longrightarrow
\text{Mod}_{\mathcal{O}_X(U)}
$
qui donne lieu à

\begin{equation}
\label{cohomology-equation-total-derived-cohomology}
R\Gamma(U, -) :
D^{+}(X)
\longrightarrow
D^{+}(\mathcal{O}_X(U))
\end{equation}

d'après la discussion précédente. Nous posons $H^i(U, -) = R^i\Gamma(U, -)$. Si $U = X$, nous retrouvons (\ref{cohomology-equation-cohomology-modules}). Si $f : X \to Y$ est un morphisme d'espaces annelés, nous avons le foncteur exact à gauche
$
f_* :
\textit{Mod}(\mathcal{O}_X)
\longrightarrow
\textit{Mod}(\mathcal{O}_Y)
$
qui donne lieu à l'{\it image directe dérivée}

\begin{equation}
\label{cohomology-equation-total-derived-direct-image}
Rf_* :
D^{+}(X)
\longrightarrow
D^{+}(Y)
\end{equation}

Le $i$-ième faisceau de cohomologie de $Rf_*\mathcal{F}^\bullet$ est noté
$R^if_*\mathcal{F}^\bullet$ et appelé la $i$-ième {\it image directe supérieure},
conformément à (\ref{cohomology-equation-higher-direct-image-modules}). Les deux foncteurs affichés ci-dessus sont des foncteurs exacts de catégories dérivées.

\medskip\noindent
{\bf Abus de notation:} Lorsque le foncteur $Rf_*$, ou tout autre foncteur dérivé, est appliqué à un faisceau $\mathcal{F}$ sur $X$ ou à un complexe de faisceaux, il est entendu que $\mathcal{F}$ a été remplacé par une résolution appropriée de $\mathcal{F}$. Afin de faciliter ce type d'opération, nous dirons, étant donné un objet $\mathcal{F}^\bullet \in D(\mathcal{O}_X)$, qu'un complexe borné inférieurement $\mathcal{I}^\bullet$ d'objets injectifs de $\textit{Mod}(\mathcal{O}_X)$ {\it représente $\mathcal{F}^\bullet$ dans la catégorie dérivée} s'il existe un quasi-isomorphisme $\mathcal{F}^\bullet \to \mathcal{I}^\bullet$. Dans le même esprit, la phrase « soit $\alpha : \mathcal{F}^\bullet \to \mathcal{G}^\bullet$ un morphisme de $D(\mathcal{O}_X)$ » ne signifie pas que $\alpha$ est représenté par un morphisme de complexes. S'il s'agit effectivement d'un morphisme de complexes, nous le préciserons.









\section{Première cohomologie et torseurs}

\label{cohomology-section-h1-torsors}

\begin{definition}

\label{cohomology-definition-torsor}

Soit $X$ un espace topologique. Soit $\mathcal{G}$ un faisceau de groupes (éventuellement non commutatifs) sur $X$. Un {\it pseudo-torseur}, ou plus précisément un {\it pseudo-$\mathcal{G}$-torseur}, est un faisceau d'ensembles $\mathcal{F}$ sur $X$ muni d'une action
$\mathcal{G} \times \mathcal{F} \to \mathcal{F}$ telle que

\begin{enumerate}

\item dès que $\mathcal{F}(U)$ est non vide, l'action
$\mathcal{G}(U) \times \mathcal{F}(U) \to \mathcal{F}(U)$
est simplement transitive

\end{enumerate}

Un {\it morphisme de pseudo-$\mathcal{G}$-torseurs}  $\mathcal{F} \to \mathcal{F}'$
est un morphisme de faisceaux d'ensembles compatible avec les
actions de $\mathcal{G}$. Un {\it torseur}, ou plus précisément un
{\it $\mathcal{G}$-torseur}, est un pseudo-$\mathcal{G}$-torseur tel que, de plus,

\begin{enumerate}

\item[(2)] pour tout $x \in X$, la fibre $\mathcal{F}_x$ est non vide.

\end{enumerate}

Un {\it morphisme de $\mathcal{G}$-torseurs} est un morphisme de pseudo-$\mathcal{G}$-torseurs. Le {\it $\mathcal{G}$-torseur trivial}
est le faisceau $\mathcal{G}$ muni de l'action à gauche évidente de
$\mathcal{G}$.

\end{definition}

\noindent
Il est clair qu'un morphisme de torseurs est automatiquement un isomorphisme.

\begin{lemma}

\label{cohomology-lemma-trivial-torsor}
Soit $X$ un espace topologique. Soit $\mathcal{G}$ un faisceau de groupes (éventuellement non commutatifs) sur $X$. Un $\mathcal{G}$-torseur $\mathcal{F}$ est trivial si et seulement si $\mathcal{F}(X) \not = \emptyset$.
\end{lemma}

\begin{proof}
Démonstration omise.
\end{proof}

\begin{lemma}

\label{cohomology-lemma-torsors-h1}
Soit $X$ un espace topologique. Soit $\mathcal{H}$ un faisceau abélien sur $X$. Il existe une bijection canonique entre l'ensemble des classes d'isomorphisme de $\mathcal{H}$-torseurs et $H^1(X, \mathcal{H})$.
\end{lemma}

\begin{proof}

Soit $\mathcal{F}$ un $\mathcal{H}$-torseur. Considérons le faisceau abélien libre $\mathbf{Z}[\mathcal{F}]$
sur $\mathcal{F}$. C'est la faisceautisation du préfaisceau qui associe à tout ouvert $U \subset X$ l'ensemble des sommes formelles finies $\sum n_i[s_i]$ avec $n_i \in \mathbf{Z}$
et $s_i \in \mathcal{F}(U)$. Il existe un morphisme naturel
$$
\sigma : \mathbf{Z}[\mathcal{F}] \longrightarrow \underline{\mathbf{Z}}
$$
qui, à une section locale $\sum n_i[s_i]$, associe $\sum n_i$. Le noyau de $\sigma$ est engendré par les sections locales de la forme
$[s] - [s']$. Il y a une application canonique
$a : \Ker(\sigma) \to \mathcal{H}$
qui envoie $[s] - [s'] \mapsto h$, où $h$ est la section locale de
$\mathcal{H}$ telle que $h \cdot s' = s$. Considérons le diagramme cocartésien
$$
\xymatrix{
0 \ar[r] &
\Ker(\sigma) \ar[r] \ar[d]^a &
\mathbf{Z}[\mathcal{F}] \ar[r] \ar[d] &
\underline{\mathbf{Z}} \ar[r] \ar[d] &
0 \\
0 \ar[r] &
\mathcal{H} \ar[r] &
\mathcal{E} \ar[r] &
\underline{\mathbf{Z}} \ar[r] &
0
}
$$
Ici, $\mathcal{E}$ est l'extension obtenue par somme amalgamée. La longue suite exacte de cohomologie associée à la suite exacte courte inférieure fournit un élément
$\xi = \xi_\mathcal{F} \in H^1(X, \mathcal{H})$
en appliquant le morphisme de bord à $1 \in H^0(X, \underline{\mathbf{Z}})$.

\medskip\noindent

Réciproquement, étant donné $\xi \in H^1(X, \mathcal{H})$, nous pouvons associer à $\xi$
un torseur comme suit. Choisissons un plongement $\mathcal{H} \to \mathcal{I}$
de $\mathcal{H}$ dans un faisceau abélien injectif $\mathcal{I}$. Posons
$\mathcal{Q} = \mathcal{I}/\mathcal{H}$, de sorte que nous avons une suite exacte courte
$$
\xymatrix{
0 \ar[r] &
\mathcal{H} \ar[r] &
\mathcal{I} \ar[r] &
\mathcal{Q} \ar[r] &
0
}
$$
L'élément $\xi$ est l'image d'une section globale $q \in H^0(X, \mathcal{Q})$
car $H^1(X, \mathcal{I}) = 0$ (voir Catégories dérivées, lemme \ref{derived-lemma-higher-derived-functors}). Soit $\mathcal{F} \subset \mathcal{I}$ le sous-faisceau (d'ensembles) des sections qui s'envoient sur $q$ dans le faisceau $\mathcal{Q}$. Il est facile de vérifier que
$\mathcal{F}$ est un torseur.

\medskip\noindent

Nous omettons de vérifier que les deux constructions données ci-dessus sont inverses l'une de l'autre.

\end{proof}







\section{Première cohomologie et extensions}

\label{cohomology-section-h1-extensions}

\begin{lemma}

\label{cohomology-lemma-h1-extensions}
Soit $(X, \mathcal{O}_X)$ un espace annelé. Soit $\mathcal{F}$ un faisceau de $\mathcal{O}_X$-modules. Il existe une bijection canonique $$
\Ext^1_{\textit{Mod}(\mathcal{O}_X)}(\mathcal{O}_X, \mathcal{F})
\longrightarrow
H^1(X, \mathcal{F})
$$ qui associe à l'extension $$
0 \to \mathcal{F} \to \mathcal{E} \to \mathcal{O}_X \to 0
$$ l'image de $1 \in \Gamma(X, \mathcal{O}_X)$ dans $H^1(X, \mathcal{F})$.
\end{lemma}

\begin{proof}
Construisons l'inverse de l'application donnée dans le lemme. Soit $\xi \in H^1(X, \mathcal{F})$. Choisissons une injection $\mathcal{F} \subset \mathcal{I}$, où $\mathcal{I}$ est injectif dans $\textit{Mod}(\mathcal{O}_X)$. Posons $\mathcal{Q} = \mathcal{I}/\mathcal{F}$. Par la longue suite exacte de cohomologie, $\xi$ est l'image d'une section $\tilde \xi \in \Gamma(X, \mathcal{Q}) =
\Hom_{\mathcal{O}_X}(\mathcal{O}_X, \mathcal{Q})$. Formons alors le produit fibré $$
\xymatrix{
0 \ar[r] &
\mathcal{F} \ar[r] \ar@{=}[d] &
\mathcal{E} \ar[r] \ar[d] &
\mathcal{O}_X \ar[r] \ar[d]^{\tilde \xi} &
0 \\
0 \ar[r] &
\mathcal{F} \ar[r] &
\mathcal{I} \ar[r] &
\mathcal{Q} \ar[r] &
0
}
$$ voir Homologie, section \ref{homology-section-extensions}.
\end{proof}







\section{Première cohomologie et faisceaux inversibles}

\label{cohomology-section-invertible-sheaves}

\noindent
Le groupe de Picard d'un espace annelé est défini dans Faisceaux de modules, section \ref{modules-section-invertible}.

\begin{lemma}

\label{cohomology-lemma-h1-invertible}

Soit $(X, \mathcal{O}_X)$ un espace annelé. Si tous les anneaux locaux
$\mathcal{O}_{X, x}$ sont locaux, alors il existe un isomorphisme canonique
$$
H^1(X, \mathcal{O}_X^*) = \Pic(X).
$$
de groupes abéliens.

\end{lemma}

\begin{proof}

Soit $\mathcal{L}$ un $\mathcal{O}_X$-module inversible. Considérons le préfaisceau $\mathcal{L}^*$ défini par la règle
$$
U \longmapsto \{s \in \mathcal{L}(U)
\text{ tel que } \mathcal{O}_U \xrightarrow{s \cdot -} \mathcal{L}_U
\text{ est un isomorphisme}\}
$$
Ce préfaisceau satisfait à la condition de faisceau. De plus, si
$f \in \mathcal{O}_X^*(U)$ et $s \in \mathcal{L}^*(U)$, alors il est clair que
$fs \in \mathcal{L}^*(U)$. De la même façon, si $s, s' \in \mathcal{L}^*(U)$
alors il existe un unique $f \in \mathcal{O}_X^*(U)$ tel que
$fs = s'$. De plus, le faisceau $\mathcal{L}^*$ possède localement des sections d'après Faisceaux de modules, lemme \ref{modules-lemma-invertible-is-locally-free-rank-1}. Ainsi, $\mathcal{L}^*$ est un $\mathcal{O}_X^*$-torseur. Nous obtenons donc une application
$$
\begin{matrix}
\text{faisceaux inversibles sur }(X, \mathcal{O}_X) \\
\text{ à isomorphisme près}
\end{matrix}
\longrightarrow
\begin{matrix}
\mathcal{O}_X^*\text{-torsors} \\
\text{ à isomorphisme près}
\end{matrix}
$$
Nous omettons de vérifier qu'il s'agit d'un homomorphisme de groupes abéliens. D'après le lemme \ref{cohomology-lemma-torsors-h1},
le membre de droite est canoniquement en bijection avec $H^1(X, \mathcal{O}_X^*)$. Il reste donc à montrer que cette application est injective et surjective.

\medskip\noindent
Injectivité. Si le torseur $\mathcal{L}^*$ est trivial, le lemme \ref{cohomology-lemma-trivial-torsor} montre que $\mathcal{L}^*$ possède une section globale. Cela signifie précisément que $\mathcal{L} \cong \mathcal{O}_X$, qui est l'élément neutre de $\Pic(X)$.

\medskip\noindent

Surjectivité. Soit $\mathcal{F}$ un $\mathcal{O}_X^*$-torseur. Considérons le préfaisceau d'ensembles
$$
\mathcal{L}_1 : U \longmapsto
(\mathcal{F}(U) \times \mathcal{O}_X(U))/\mathcal{O}_X^*(U)
$$
où l'action de $f \in \mathcal{O}_X^*(U)$ sur
$(s, g)$ est donnée par $(fs, f^{-1}g)$. Alors $\mathcal{L}_1$ est un préfaisceau de $\mathcal{O}_X$-modules si l'on pose
$(s, g) + (s', g') = (s, g + (s'/s)g')$, où $s'/s$ est la section locale $f$ de $\mathcal{O}_X^*$ telle que $fs = s'$,
$h(s, g) = (s, hg)$ pour toute section locale $h$ de $\mathcal{O}_X$. Nous omettons de vérifier que la faisceautisation
$\mathcal{L} = \mathcal{L}_1^\#$ est un $\mathcal{O}_X$-module inversible dont le $\mathcal{O}_X^*$-torseur associé $\mathcal{L}^*$ est isomorphe à $\mathcal{F}$.

\end{proof}













\section{Caractère local de la cohomologie}

\label{cohomology-section-locality}

\noindent
Le lemme suivant montre qu'il n'y a aucune ambiguïté dans la définition de la cohomologie d'un faisceau $\mathcal{F}$ sur un ouvert.

\begin{lemma}

\label{cohomology-lemma-cohomology-of-open}
Soit $X$ un espace annelé. Soit $U \subset X$ un sous-espace ouvert.
\begin{enumerate}
\item Si $\mathcal{I}$ est un $\mathcal{O}_X$-module injectif, alors $\mathcal{I}|_U$ est un $\mathcal{O}_U$-module injectif.
\item Pour tout faisceau de $\mathcal{O}_X$-modules $\mathcal{F}$, nous avons $H^p(U, \mathcal{F}) = H^p(U, \mathcal{F}|_U)$.
\end{enumerate}

\end{lemma}

\begin{proof}
Notons $j : U \to X$ l'immersion ouverte. Rappelons que le foncteur de restriction $j^{-1}$ à $U$ est adjoint à droite du foncteur $j_!$ d'extension par $0$, voir Faisceaux, lemme \ref{sheaves-lemma-j-shriek-modules}. De plus, $j_!$ est exact. Par conséquent, (1) découle d'Homologie, lemme \ref{homology-lemma-adjoint-preserve-injectives}.

\medskip\noindent
Par définition, $H^p(U, \mathcal{F}) = H^p(\Gamma(U, \mathcal{I}^\bullet))$, où $\mathcal{F} \to \mathcal{I}^\bullet$ est une résolution injective dans $\textit{Mod}(\mathcal{O}_X)$. D'après ce qui précède, $\mathcal{F}|_U \to \mathcal{I}^\bullet|_U$ est une résolution injective dans $\textit{Mod}(\mathcal{O}_U)$. Par conséquent, $H^p(U, \mathcal{F}|_U)$ est égal à $H^p(\Gamma(U, \mathcal{I}^\bullet|_U))$. Bien entendu, $\Gamma(U, \mathcal{F}) = \Gamma(U, \mathcal{F}|_U)$ pour tout faisceau $\mathcal{F}$ sur $X$. On obtient ainsi l'égalité de (2).
\end{proof}

\noindent

Soit $X$ un espace annelé. Soit $\mathcal{F}$ un faisceau de $\mathcal{O}_X$-modules. Soient $U \subset V \subset X$ des ouverts. Il existe alors une {\it application de restriction} canonique

\begin{equation}
\label{cohomology-equation-restriction-mapping}
H^n(V, \mathcal{F})
\longrightarrow
H^n(U, \mathcal{F}), \quad
\xi \longmapsto \xi|_U
\end{equation}

fonctorielle en $\mathcal{F}$. En effet, choisissons une résolution injective $\mathcal{F} \to \mathcal{I}^\bullet$. Les applications de restriction des faisceaux $\mathcal{I}^p$ donnent un morphisme de complexes
$$
\Gamma(V, \mathcal{I}^\bullet)
\longrightarrow
\Gamma(U, \mathcal{I}^\bullet)
$$
Le membre de gauche est un complexe représentant $R\Gamma(V, \mathcal{F})$,
et celui de droite un complexe représentant $R\Gamma(U, \mathcal{F})$. En appliquant le foncteur $H^n$, nous obtenons l'application sur les groupes de cohomologie. Comme indiqué, nous la noterons $\xi \mapsto \xi|_U$. Ainsi, la règle $U \mapsto H^n(U, \mathcal{F})$ définit un préfaisceau de
$\mathcal{O}_X$-modules. Ce préfaisceau est habituellement noté
$\underline{H}^n(\mathcal{F})$. Nous en donnerons une autre interprétation dans le lemme \ref{cohomology-lemma-include}.

\begin{lemma}

\label{cohomology-lemma-kill-cohomology-class-on-covering}
Soit $X$ un espace annelé. Soit $\mathcal{F}$ un faisceau de $\mathcal{O}_X$-modules. Soit $U \subset X$ un sous-espace ouvert. Soient $n > 0$ et $\xi \in H^n(U, \mathcal{F})$. Il existe alors un recouvrement ouvert $U = \bigcup_{i\in I} U_i$ tel que $\xi|_{U_i} = 0$ pour tout $i \in I$.
\end{lemma}

\begin{proof}

Soit $\mathcal{F} \to \mathcal{I}^\bullet$ une résolution injective. Alors
$$
H^n(U, \mathcal{F}) =
\frac{\Ker(\mathcal{I}^n(U) \to \mathcal{I}^{n + 1}(U))}
{\Im(\mathcal{I}^{n - 1}(U) \to \mathcal{I}^n(U))}.
$$
Choisissons un élément $\tilde \xi \in \mathcal{I}^n(U)$ représentant la classe de cohomologie dans la présentation ci-dessus. Puisque $\mathcal{I}^\bullet$
est une résolution injective de $\mathcal{F}$ et que $n > 0$, le complexe $\mathcal{I}^\bullet$ est exact en degré $n$. Par conséquent,
$\Im(\mathcal{I}^{n - 1} \to \mathcal{I}^n) =
\Ker(\mathcal{I}^n \to \mathcal{I}^{n + 1})$ comme faisceaux. Puisque $\tilde \xi$ est une section du faisceau noyau sur $U$,
il existe un recouvrement ouvert $U = \bigcup_{i \in I} U_i$
tel que $\tilde \xi|_{U_i}$ soit l'image par $d$ d'une section
$\xi_i \in \mathcal{I}^{n - 1}(U_i)$. Par définition de la restriction, $\xi|_{U_i}$ correspond à la classe de
$\tilde \xi|_{U_i}$, ce qui conclut la preuve.

\end{proof}

\begin{lemma}

\label{cohomology-lemma-describe-higher-direct-images}

Soit $f : X \to Y$ un morphisme d'espaces annelés. Soit $\mathcal{F}$ un $\mathcal{O}_X$-module. Les faisceaux $R^if_*\mathcal{F}$ sont les faisceaux associés aux préfaisceaux
$$
V \longmapsto H^i(f^{-1}(V), \mathcal{F})
$$
munis des applications de restriction de l'équation (\ref{cohomology-equation-restriction-mapping}). Il existe un énoncé analogue pour $R^if_*$ appliqué à un complexe borné inférieurement $\mathcal{F}^\bullet$.

\end{lemma}

\begin{proof}

Soit $\mathcal{F} \to \mathcal{I}^\bullet$ une résolution injective. Par définition, $R^if_*\mathcal{F}$ est le $i$-ième faisceau de cohomologie du complexe
$$
f_*\mathcal{I}^0 \to f_*\mathcal{I}^1 \to f_*\mathcal{I}^2 \to \ldots
$$
Par définition de la structure de catégorie abélienne sur les $\mathcal{O}_Y$-modules, ce faisceau de cohomologie est le faisceau associé au préfaisceau
$$
V
\longmapsto
\frac{\Ker(f_*\mathcal{I}^i(V) \to f_*\mathcal{I}^{i + 1}(V))}
{\Im(f_*\mathcal{I}^{i - 1}(V) \to f_*\mathcal{I}^i(V))}
$$
et cela est évidemment égal à
$$
\frac{\Ker(\mathcal{I}^i(f^{-1}(V)) \to \mathcal{I}^{i + 1}(f^{-1}(V)))}
{\Im(\mathcal{I}^{i - 1}(f^{-1}(V)) \to \mathcal{I}^i(f^{-1}(V)))}
$$
qui est égal à $H^i(f^{-1}(V), \mathcal{F})$,
ce qui démontre le résultat.

\end{proof}

\begin{lemma}

\label{cohomology-lemma-localize-higher-direct-images}
Soit $f : X \to Y$ un morphisme d'espaces annelés. Soit $\mathcal{F}$ un $\mathcal{O}_X$-module. Soit $V \subset Y$ un sous-espace ouvert. Notons $g : f^{-1}(V) \to V$ la restriction de $f$. Alors $$
R^pg_*(\mathcal{F}|_{f^{-1}(V)}) = (R^pf_*\mathcal{F})|_V
$$ Il existe un énoncé analogue pour l'image dérivée $Rf_*\mathcal{F}^\bullet$, où $\mathcal{F}^\bullet$ est un complexe borné inférieurement de $\mathcal{O}_X$-modules.
\end{lemma}

\begin{proof}
Première démonstration. Appliquer les lemmes \ref{cohomology-lemma-describe-higher-direct-images} et \ref{cohomology-lemma-cohomology-of-open} donne l'égalité affichée. Deuxième démonstration. Choisissons une résolution injective $\mathcal{F} \to \mathcal{I}^\bullet$ et utilisons le fait que $\mathcal{F}|_{f^{-1}(V)} \to \mathcal{I}^\bullet|_{f^{-1}(V)}$ est également une résolution injective.
\end{proof}

\begin{remark}

\label{cohomology-remark-daniel}

Voici une autre approche des démonstrations des lemmes \ref{cohomology-lemma-kill-cohomology-class-on-covering} et
\ref{cohomology-lemma-describe-higher-direct-images} ci-dessus. Soit $(X, \mathcal{O}_X)$ un espace annelé. Soit $i_X : \textit{Mod}(\mathcal{O}_X) \to \textit{PMod}(\mathcal{O}_X)$
le foncteur d'inclusion, et soit $\#$ le foncteur de faisceautisation. Rappelons que $i_X$ est exact à gauche et que $\#$ est exact.

\begin{enumerate}

\item Démontrons d'abord le lemme \ref{cohomology-lemma-include} ci-dessous, qui affirme que les foncteurs dérivés à droite de $i_X$ sont donnés par
$R^pi_X\mathcal{F} = \underline{H}^p(\mathcal{F})$. En voici une autre démonstration : l'égalité est claire pour $p = 0$. Les deux familles $(R^pi_X)_{p \geq 0}$ et $(\underline{H}^p)_{p \geq 0}$
sont des foncteurs delta qui s'annulent sur les objets injectifs ; ils sont donc universels, et par conséquent isomorphes. Voir Homologie, section \ref{homology-section-cohomological-delta-functor}.
\item Une reformulation du lemme \ref{cohomology-lemma-kill-cohomology-class-on-covering}
est que $(\underline{H}^p(\mathcal{F}))^\# = 0$, $p > 0$, pour tout faisceau
de $\mathcal{O}_X$-modules $\mathcal{F}$. Pour le vérifier, utilisons l'exactitude de ${}^\#$ :
$$
(\underline{H}^p(\mathcal{F}))^\# =
(R^pi_X\mathcal{F})^\# =
R^p(\# \circ i_X)(\mathcal{F}) = 0
$$
car $\# \circ i_X$ est le foncteur identité.
\item Soit $f : X \to Y$ un morphisme d'espaces annelés. Soit $\mathcal{F}$ un $\mathcal{O}_X$-module. Le préfaisceau
$V \mapsto H^p(f^{-1}V, \mathcal{F})$ est égal à
$R^p (i_Y \circ f_*)\mathcal{F}$. On peut le démontrer en remarquant que les deux constructions donnent des foncteurs delta universels, comme dans l'argument de (1). Le lemme \ref{cohomology-lemma-describe-higher-direct-images}
affirme que $R^p f_* \mathcal{F}= (R^p (i_Y \circ f_*)\mathcal{F})^\#$. En utilisant à nouveau l'exactitude de $\#$ et le fait que $\# \circ i_Y$ est le foncteur identité, on obtient
$$
R^p f_* \mathcal{F} =
R^p(\# \circ i_Y \circ f_*)\mathcal{F} =
(R^p (i_Y \circ f_*)\mathcal{F})^\#
$$
comme souhaité.

\end{enumerate}

\end{remark}





\section{Mayer--Vietoris}

\label{cohomology-section-mayer-vietoris}

\noindent
Nous construirons plus loin la suite spectrale reliant la cohomologie de {\v C}ech à la cohomologie des faisceaux, voir le lemme \ref{cohomology-lemma-cech-spectral-sequence}. La longue suite exacte de Mayer--Vietoris en est un cas particulier. Comme il s'agit d'une variante fondamentale, utile et facile à comprendre de cette suite spectrale, nous la traitons ici séparément.

\begin{lemma}

\label{cohomology-lemma-injective-restriction-surjective}

\begin{slogan}
Les injectifs sont flasques.
\end{slogan}
Soit $X$ un espace annelé. Soient $U' \subset U \subset X$ des sous-espaces ouverts. Pour tout $\mathcal{O}_X$-module injectif $\mathcal{I}$, l'application de restriction $\mathcal{I}(U) \to \mathcal{I}(U')$ est surjective.
\end{lemma}

\begin{proof}

Soient $j : U \to X$ et $j' : U' \to X$ les immersions ouvertes. Rappelons que $j_!\mathcal{O}_U$ est l'extension par zéro de
$\mathcal{O}_U = \mathcal{O}_X|_U$, voir Faisceaux, section \ref{sheaves-section-open-immersions}. Puisque $j_!$ est adjoint à gauche de la restriction, pour tout faisceau $\mathcal{F}$ de $\mathcal{O}_X$-modules nous avons
$$
\Hom_{\mathcal{O}_X}(j_!\mathcal{O}_U, \mathcal{F})
=
\Hom_{\mathcal{O}_U}(\mathcal{O}_U, \mathcal{F}|_U)
=
\mathcal{F}(U)
$$
voir Faisceaux, lemme \ref{sheaves-lemma-j-shriek-modules}. De même, le faisceau $j'_!\mathcal{O}_{U'}$ représente le foncteur $\mathcal{F} \mapsto \mathcal{F}(U')$. En outre, il existe une application canonique évidente de $\mathcal{O}_X$-modules
$$
j'_!\mathcal{O}_{U'} \longrightarrow j_!\mathcal{O}_U
$$
qui correspond à l'application de restriction
$\mathcal{F}(U) \to \mathcal{F}(U')$ par le lemme de Yoneda (Catégories, lemme \ref{categories-lemma-yoneda}). D'après la description des fibres de
$j'_!\mathcal{O}_{U'}$, $j_!\mathcal{O}_U$
nous voyons que l'application affichée ci-dessus est injective (voir le lemme cité ci-dessus). Si $\mathcal{I}$ est un $\mathcal{O}_X$-module injectif, alors l'application
$$
\Hom_{\mathcal{O}_X}(j_!\mathcal{O}_U, \mathcal{I})
\longrightarrow
\Hom_{\mathcal{O}_X}(j'_!\mathcal{O}_{U'}, \mathcal{I})
$$
est surjective, voir Homologie, lemme \ref{homology-lemma-characterize-injectives}. En réunissant ces observations, nous obtenons le lemme.

\end{proof}

\begin{lemma}
[Mayer--Vietoris]
\label{cohomology-lemma-mayer-vietoris}
Soit $X$ un espace annelé. Supposons que $X = U \cup V$ soit la réunion de deux ouverts. Pour tout $\mathcal{O}_X$-module $\mathcal{F}$, il existe une longue suite exacte de cohomologie $$
0 \to
H^0(X, \mathcal{F}) \to
H^0(U, \mathcal{F}) \oplus H^0(V, \mathcal{F}) \to
H^0(U \cap V, \mathcal{F}) \to
H^1(X, \mathcal{F}) \to \ldots
$$ Cette longue suite exacte est fonctorielle en $\mathcal{F}$.
\end{lemma}

\begin{proof}

La condition de faisceau affirme que le noyau de
$(1, -1) : \mathcal{F}(U) \oplus \mathcal{F}(V) \to \mathcal{F}(U \cap V)$
est égal à l'image de $\mathcal{F}(X)$ par la première application, pour tout faisceau abélien $\mathcal{F}$. Le lemme \ref{cohomology-lemma-injective-restriction-surjective} ci-dessus implique que l'application
$(1, -1) : \mathcal{I}(U) \oplus \mathcal{I}(V) \to \mathcal{I}(U \cap V)$
est surjective dès que $\mathcal{I}$ est un $\mathcal{O}_X$-module injectif. Ainsi, si $\mathcal{F} \to \mathcal{I}^\bullet$ est une résolution injective de $\mathcal{F}$, nous obtenons une suite exacte courte de complexes
$$
0 \to
\mathcal{I}^\bullet(X) \to
\mathcal{I}^\bullet(U) \oplus \mathcal{I}^\bullet(V) \to
\mathcal{I}^\bullet(U \cap V) \to
0.
$$
Le passage à la cohomologie donne le résultat (utiliser Homologie, lemme \ref{homology-lemma-long-exact-sequence-cochain}). Nous omettons la vérification de la fonctorialité de la suite.

\end{proof}

\begin{lemma}
[Relative Mayer--Vietoris]
\label{cohomology-lemma-relative-mayer-vietoris}
Soit $f : X \to Y$ un morphisme d'espaces annelés. Supposons que $X = U \cup V$ soit la réunion de deux ouverts. Notons $a = f|_U : U \to Y$, $b = f|_V : V \to Y$ et $c = f|_{U \cap V} : U \cap V \to Y$. Pour tout $\mathcal{O}_X$-module $\mathcal{F}$, il existe une longue suite exacte $$
0 \to
f_*\mathcal{F} \to
a_*(\mathcal{F}|_U) \oplus b_*(\mathcal{F}|_V) \to
c_*(\mathcal{F}|_{U \cap V}) \to
R^1f_*\mathcal{F} \to \ldots
$$ Cette longue suite exacte est fonctorielle en $\mathcal{F}$.
\end{lemma}

\begin{proof}

Soit $\mathcal{F} \to \mathcal{I}^\bullet$ une résolution injective de $\mathcal{F}$. Nous affirmons que nous obtenons une suite exacte courte de complexes
$$
0 \to
f_*\mathcal{I}^\bullet \to
a_*\mathcal{I}^\bullet|_U \oplus b_*\mathcal{I}^\bullet|_V \to
c_*\mathcal{I}^\bullet|_{U \cap V} \to
0.
$$
En effet, pour tout ouvert $W \subset Y$ et tout $n \geq 0$, la suite correspondante des groupes de sections sur $W$

$$
0 \to
\mathcal{I}^n(f^{-1}(W)) \to
\mathcal{I}^n(U \cap f^{-1}(W))
\oplus \mathcal{I}^n(V \cap f^{-1}(W)) \to
\mathcal{I}^n(U \cap V \cap f^{-1}(W)) \to
0
$$
est exacte courte d'après la démonstration du lemme \ref{cohomology-lemma-mayer-vietoris}. Le lemme s'obtient en prenant les faisceaux de cohomologie et en utilisant le fait que
$\mathcal{I}^\bullet|_U$ est une résolution injective de $\mathcal{F}|_U$
et de même pour $\mathcal{I}^\bullet|_V$ et $\mathcal{I}^\bullet|_{U \cap V}$,
voir le lemme \ref{cohomology-lemma-cohomology-of-open}.

\end{proof}



















\section{Le complexe de {\v C}
ech et la cohomologie de {\v C}ech}
\label{cohomology-section-cech}

\noindent

Soit $X$ un espace topologique. Soit $\mathcal{U} : U = \bigcup_{i \in I} U_i$ un recouvrement ouvert, voir Topologie, notion de base (\ref{topology-item-covering}). Comme il est d'usage, nous notons
$U_{i_0\ldots i_p} = U_{i_0} \cap \ldots \cap U_{i_p}$ pour
l'intersection de $(p + 1)$ membres de $\mathcal{U}$. Soit $\mathcal{F}$ un préfaisceau abélien sur $X$. Posons
$$
\check{\mathcal{C}}^p(\mathcal{U}, \mathcal{F})
=
\prod\nolimits_{(i_0, \ldots, i_p) \in I^{p + 1}}
\mathcal{F}(U_{i_0\ldots i_p}).
$$
C'est un groupe abélien. Pour
$s \in \check{\mathcal{C}}^p(\mathcal{U}, \mathcal{F})$, nous notons
$s_{i_0\ldots i_p}$ sa valeur dans $\mathcal{F}(U_{i_0\ldots i_p})$. Remarquons que si $s \in \check{\mathcal{C}}^1(\mathcal{U}, \mathcal{F})$
et $i, j \in I$, alors $s_{ij}$ et $s_{ji}$ appartiennent tous deux à $\mathcal{F}(U_i \cap U_j)$, mais aucune relation n'est imposée entre $s_{ij}$ et $s_{ji}$. Autrement dit, nous ne travaillons {\it pas}
avec des cochaînes alternées (celles-ci seront définies dans la section \ref{cohomology-section-alternating-cech}). Nous définissons
$$
d : \check{\mathcal{C}}^p(\mathcal{U}, \mathcal{F})
\longrightarrow
\check{\mathcal{C}}^{p + 1}(\mathcal{U}, \mathcal{F})
$$
par la formule

\begin{equation}
\label{cohomology-equation-d-cech}
d(s)_{i_0\ldots i_{p + 1}}
=
\sum\nolimits_{j = 0}^{p + 1}
(-1)^j
s_{i_0\ldots \hat i_j \ldots i_{p + 1}}|_{U_{i_0\ldots i_{p + 1}}}
\end{equation}

On vérifie immédiatement que $d \circ d = 0$. Autrement dit,
$\check{\mathcal{C}}^\bullet(\mathcal{U}, \mathcal{F})$ est un complexe.

\begin{definition}

\label{cohomology-definition-cech-complex}
Soit $X$ un espace topologique. Soit $\mathcal{U} : U = \bigcup_{i \in I} U_i$ un recouvrement ouvert. Soit $\mathcal{F}$ un préfaisceau abélien sur $X$. Le complexe $\check{\mathcal{C}}^\bullet(\mathcal{U}, \mathcal{F})$ est le {\it complexe de {\v C}ech} associé à $\mathcal{F}$ et au recouvrement ouvert $\mathcal{U}$. Ses groupes de cohomologie $H^i(\check{\mathcal{C}}^\bullet(\mathcal{U}, \mathcal{F}))$ sont appelés les {\it groupes de cohomologie de {\v C}ech} associés à $\mathcal{F}$ et au recouvrement $\mathcal{U}$. Ils sont notés $\check H^i(\mathcal{U}, \mathcal{F})$.
\end{definition}

\begin{lemma}

\label{cohomology-lemma-cech-h0}

Soit $X$ un espace topologique. Soit $\mathcal{F}$ un préfaisceau abélien sur $X$. Les assertions suivantes sont équivalentes :

\begin{enumerate}

\item  $\mathcal{F}$ est un faisceau abélien, et
\item pour chaque recouvrement ouvert $\mathcal{U} : U = \bigcup_{i \in I} U_i$
l'application naturelle
$$
\mathcal{F}(U) \to \check{H}^0(\mathcal{U}, \mathcal{F})
$$
est bijective.

\end{enumerate}

\end{lemma}

\begin{proof}

Cela résulte du fait que la condition de faisceau affirme précisément que
$\mathcal{F}(U) \to \check{H}^0(\mathcal{U}, \mathcal{F})$
est bijectif pour chaque recouvrement ouvert.

\end{proof}

\begin{lemma}

\label{cohomology-lemma-cech-trivial}
Soit $X$ un espace topologique. Soit $\mathcal{F}$ un préfaisceau abélien sur $X$. Soit $\mathcal{U} : U = \bigcup_{i \in I} U_i$ un recouvrement ouvert. Si $U_i = U$ pour un certain $i \in I$, alors le complexe de {\v C}ech augmenté $$
\mathcal{F}(U) \to \check{\mathcal{C}}^\bullet(\mathcal{U}, \mathcal{F})
$$ obtenu en plaçant $\mathcal{F}(U)$ en degré $-1$, avec pour différentielle l'application canonique de $\mathcal{F}(U)$ dans $\check{\mathcal{C}}^0(\mathcal{U}, \mathcal{F})$, est homotopiquement équivalent à $0$.
\end{lemma}

\begin{proof}

Fixons un élément $i \in I$ tel que $U = U_i$. Remarquons que
$U_{i_0 \ldots i_p} = U_{i_0 \ldots \hat i_j \ldots i_p}$ si $i_j = i$. Définissons une homotopie
$$
h :
\prod\nolimits_{i_0 \ldots i_{p + 1}} \mathcal{F}(U_{i_0 \ldots i_{p + 1}})
\longrightarrow
\prod\nolimits_{i_0 \ldots i_p} \mathcal{F}(U_{i_0 \ldots i_p})
$$
par la règle
$$
h(s)_{i_0 \ldots i_p} = s_{i i_0 \ldots i_p}
$$
En d'autres termes, $h : \prod_{i_0} \mathcal{F}(U_{i_0}) \to \mathcal{F}(U)$
est la projection sur le facteur $\mathcal{F}(U_i) = \mathcal{F}(U)$
et en général l'application $h$ égale la projection sur les facteurs
$\mathcal{F}(U_{i i_1 \ldots i_{p + 1}}) =
\mathcal{F}(U_{i_1 \ldots i_{p + 1}})$. Nous calculons

\begin{align*}
(dh + hd)(s)_{i_0 \ldots i_p}
& =
\sum\nolimits_{j = 0}^p
(-1)^j
h(s)_{i_0 \ldots \hat i_j \ldots i_p}
+
d(s)_{i i_0 \ldots i_p}\\
& =
\sum\nolimits_{j = 0}^p
(-1)^j
s_{i i_0 \ldots \hat i_j \ldots i_p}
+
s_{i_0 \ldots i_p}
+
\sum\nolimits_{j = 0}^p
(-1)^{j + 1}
s_{i i_0 \ldots \hat i_j \ldots i_p} \\
& =
s_{i_0 \ldots i_p}
\end{align*}

Ceci prouve que l'application identité est homotope à zéro, comme souhaité.

\end{proof}





\section{La cohomologie de {\v C}
ech comme foncteur sur les préfaisceaux}
\sectionmark{Cohomologie de {\v C}ech des préfaisceaux}
\label{cohomology-section-cech-functor}

\noindent
Attention : dans cette section, nous travaillons presque exclusivement avec des préfaisceaux et des catégories de préfaisceaux ; les résultats sont entièrement faux dans le cadre des faisceaux et des catégories de faisceaux !

\medskip\noindent

Soit $X$ un espace annelé. Soit $\mathcal{U} : U = \bigcup_{i \in I} U_i$ un recouvrement ouvert. Soit $\mathcal{F}$ un préfaisceau de $\mathcal{O}_X$-modules. Nous disposons du complexe de {\v C}ech
$\check{\mathcal{C}}^\bullet(\mathcal{U}, \mathcal{F})$
de $\mathcal{F}$ en considérant simplement $\mathcal{F}$
comme un préfaisceau de groupes abéliens. Cependant, chaque terme
$\check{\mathcal{C}}^p(\mathcal{U}, \mathcal{F})$ est naturellement muni d'une structure de $\mathcal{O}_X(U)$-module, et la différentielle est donnée par des
applications $\mathcal{O}_X(U)$-linéaires. De plus, il est clair que la construction
$$
\mathcal{F} \longmapsto \check{\mathcal{C}}^\bullet(\mathcal{U}, \mathcal{F})
$$
est fonctorielle en $\mathcal{F}$. Il s'agit en fait d'un foncteur

\begin{equation}
\label{cohomology-equation-cech-functor}
\check{\mathcal{C}}^\bullet(\mathcal{U}, -) :
\textit{PMod}(\mathcal{O}_X)
\longrightarrow
\text{Comp}^{+}(\text{Mod}_{\mathcal{O}_X(U)})
\end{equation}

Voir Catégories dérivées, définition \ref{derived-definition-complexes-notation},
pour la notation. Rappelons que la catégorie des complexes bornés inférieurement d'une catégorie abélienne est elle-même abélienne, voir Homologie, lemme \ref{homology-lemma-cat-cochain-abelian}.

\begin{lemma}

\label{cohomology-lemma-cech-exact-presheaves}

Le foncteur donné par l'équation (\ref{cohomology-equation-cech-functor}) est exact (voir Homologie, \ref{homology-lemma-exact-functor}).

\end{lemma}

\begin{proof}

Pour tout ouvert $W \subset U$, le foncteur
$\mathcal{F} \mapsto \mathcal{F}(W)$ est un foncteur additif exact de $\textit{PMod}(\mathcal{O}_X)$ vers $\text{Mod}_{\mathcal{O}_X(U)}$. Les termes
$\check{\mathcal{C}}^p(\mathcal{U}, \mathcal{F})$
du complexe sont des produits de ces foncteurs exacts, donc sont exacts. En outre, une suite de complexes est exacte si et seulement si la suite des termes de chaque degré est exacte.

\end{proof}

\begin{lemma}

\label{cohomology-lemma-cech-cohomology-delta-functor-presheaves}

Soit $X$ un espace annelé. Soit $\mathcal{U} : U = \bigcup_{i \in I} U_i$ un recouvrement ouvert. Les foncteurs $\mathcal{F} \mapsto \check{H}^n(\mathcal{U}, \mathcal{F})$
forment un $\delta$-foncteur de la catégorie abélienne des préfaisceaux de $\mathcal{O}_X$-modules vers la catégorie des $\mathcal{O}_X(U)$-modules (voir Homologie, définition \ref{homology-definition-cohomological-delta-functor}).

\end{lemma}

\begin{proof}
Par le lemme \ref{cohomology-lemma-cech-exact-presheaves}, une suite exacte courte de préfaisceaux de $\mathcal{O}_X$-modules $0 \to \mathcal{F}_1 \to \mathcal{F}_2 \to \mathcal{F}_3 \to 0$ donne une suite exacte courte de complexes de $\mathcal{O}_X(U)$-modules. Nous pouvons donc appliquer Homologie, lemme \ref{homology-lemma-long-exact-sequence-cochain}, pour obtenir les morphismes de bord $\delta_{\mathcal{F}_1 \to \mathcal{F}_2 \to \mathcal{F}_3} :
\check{H}^n(\mathcal{U}, \mathcal{F}_3) \to
\check{H}^{n + 1}(\mathcal{U}, \mathcal{F}_1)$ et la longue suite exacte correspondante. Nous omettons de vérifier que ces applications sont compatibles avec les morphismes entre suites exactes courtes de préfaisceaux.
\end{proof}

\noindent

Dans la formulation du lemme suivant, nous utilisons le foncteur $j_{p!}$ d'extension par $0$ pour les préfaisceaux de modules relativement à une immersion ouverte $j : U \to X$. Voir Faisceaux, section \ref{sheaves-section-open-immersions}. Pour tout ouvert
$W \subset X$ et tout préfaisceau $\mathcal{G}$ de $\mathcal{O}_X|_U$-modules, nous avons
$$
(j_{p!}\mathcal{G})(W) =
\left\{
\begin{matrix}
\mathcal{G}(W) & \text{si } W \subset U \\
0 & \text{sinon.}
\end{matrix}
\right.
$$
De plus, le foncteur $j_{p!}$ est adjoint à gauche du foncteur de restriction, voir Faisceaux, lemme \ref{sheaves-lemma-j-shriek-modules}. En particulier, nous avons la formule
$$
\Hom_{\mathcal{O}_X}(j_{p!}\mathcal{O}_U, \mathcal{F})
=
\Hom_{\mathcal{O}_U}(\mathcal{O}_U, \mathcal{F}|_U)
=
\mathcal{F}(U).
$$
Puisque le foncteur $\mathcal{F} \mapsto \mathcal{F}(U)$ est un foncteur exact sur la catégorie des préfaisceaux, nous concluons que le préfaisceau
$j_{p!}\mathcal{O}_U$ est un objet projectif dans la catégorie
$\textit{PMod}(\mathcal{O}_X)$, voir Homologie, lemme \ref{homology-lemma-characterize-projectives}.

\medskip\noindent
Notons que, si l'on se donne des ouverts $U \subset V \subset X$ et les immersions ouvertes associées $j_U, j_V$, alors on dispose d'une application canonique $(j_U)_{p!}\mathcal{O}_U \to (j_V)_{p!}\mathcal{O}_V$. Sur les sections au-dessus de tout ouvert $W \subset U$, c'est l'identité, et c'est $0$ sinon. Sous l'identification $\Hom_{\mathcal{O}_X}((j_U)_{p!}\mathcal{O}_U, (j_V)_{p!}\mathcal{O}_V) =
(j_V)_{p!}\mathcal{O}_V(U) = \mathcal{O}_V(U)$, elle correspond à l'élément $1 \in \mathcal{O}_V(U)$.

\begin{lemma}

\label{cohomology-lemma-cech-map-into}
Soit $X$ un espace annelé. Soit $\mathcal{U} : U = \bigcup_{i \in I} U_i$ un recouvrement ouvert. Notons $j_{i_0\ldots i_p} : U_{i_0 \ldots i_p} \to X$ l'immersion ouverte. Considérons le complexe de chaînes $K(\mathcal{U})_\bullet$ de préfaisceaux de $\mathcal{O}_X$-modules $$
\ldots
\to
\bigoplus_{i_0i_1i_2} (j_{i_0i_1i_2})_{p!}\mathcal{O}_{U_{i_0i_1i_2}}
\to
\bigoplus_{i_0i_1} (j_{i_0i_1})_{p!}\mathcal{O}_{U_{i_0i_1}}
\to
\bigoplus_{i_0} (j_{i_0})_{p!}\mathcal{O}_{U_{i_0}}
\to 0 \to \ldots
$$ où le dernier terme non nul est placé en degré $0$ et où l'application $$
(j_{i_0\ldots i_{p + 1}})_{p!}\mathcal{O}_{U_{i_0\ldots i_{p + 1}}}
\longrightarrow
(j_{i_0\ldots \hat i_j \ldots i_{p + 1}})_{p!}
\mathcal{O}_{U_{i_0\ldots \hat i_j \ldots i_{p + 1}}}
$$ est donnée par $(-1)^j$ fois l'application canonique. Alors il existe un isomorphisme $$
\Hom_{\mathcal{O}_X}(K(\mathcal{U})_\bullet, \mathcal{F})
=
\check{\mathcal{C}}^\bullet(\mathcal{U}, \mathcal{F})
$$ fonctoriel en $\mathcal{F} \in \Ob(\textit{PMod}(\mathcal{O}_X))$.
\end{lemma}

\begin{proof}

Nous avons vu dans la discussion qui précède le lemme que
$$
\Hom_{\mathcal{O}_X}(
(j_{i_0\ldots i_p})_{p!}\mathcal{O}_{U_{i_0\ldots i_p}},
\mathcal{F})
=
\mathcal{F}(U_{i_0\ldots i_p}).
$$
Par conséquent, on voit bien que la somme directe
$$
\bigoplus\nolimits_{i_0 \ldots i_p}
(j_{i_0 \ldots i_p})_{p!}\mathcal{O}_{U_{i_0 \ldots i_p}}
$$
représente le foncteur
$$
\mathcal{F}
\longmapsto
\prod\nolimits_{i_0\ldots i_p} \mathcal{F}(U_{i_0\ldots i_p}).
$$
Par conséquent, d'après Catégories, lemme de Yoneda \ref{categories-lemma-yoneda},
nous voyons qu'il existe un complexe $K(\mathcal{U})_\bullet$ dont les termes sont ceux indiqués. Il est immédiat que les applications sont celles données dans le lemme.

\end{proof}

\begin{lemma}

\label{cohomology-lemma-homology-complex}
Soit $X$ un espace annelé. Soit $\mathcal{U} : U = \bigcup_{i \in I} U_i$ un recouvrement ouvert. Soit $\mathcal{O}_\mathcal{U} \subset \mathcal{O}_X$ le préfaisceau image de l'application $\bigoplus j_{p!}\mathcal{O}_{U_i} \to \mathcal{O}_X$. Le complexe de chaînes $K(\mathcal{U})_\bullet$ de préfaisceaux du lemme \ref{cohomology-lemma-cech-map-into} ci-dessus a pour préfaisceaux d'homologie $$
H_i(K(\mathcal{U})_\bullet) =
\left\{
\begin{matrix}
0 & \text{si} & i \not = 0 \\
\mathcal{O}_\mathcal{U} & \text{si} & i = 0
\end{matrix}
\right.
$$
\end{lemma}

\begin{proof}
Considérons le complexe étendu $K^{ext}_\bullet$ obtenu en plaçant $\mathcal{O}_\mathcal{U}$ en degré $-1$, muni de l'application évidente $K(\mathcal{U})_0 =
\bigoplus_{i_0} (j_{i_0})_{p!}\mathcal{O}_{U_{i_0}} \to
\mathcal{O}_\mathcal{U}$. Il suffit de montrer que le complexe obtenu en prenant les sections de ce complexe étendu sur tout ouvert $W \subset X$ est acyclique. En fait, nous affirmons que, pour chaque $W \subset X$, le complexe $K^{ext}_\bullet(W)$ est homotopiquement équivalent au complexe nul. Écrivons $I = I_1 \amalg I_2$, où $W \subset U_i$ si et seulement si $i \in I_1$.

\medskip\noindent
Si $I_1 = \emptyset$, alors le complexe $K^{ext}_\bullet(W) = 0$, donc il n'y a rien à prouver.

\medskip\noindent

Si $I_1 \not = \emptyset$, alors
$\mathcal{O}_\mathcal{U}(W) = \mathcal{O}_X(W)$
et
$$
K^{ext}_p(W) =
\bigoplus\nolimits_{i_0 \ldots i_p \in I_1} \mathcal{O}_X(W).
$$
Cela résulte de la description simple des préfaisceaux
$(j_{i_0 \ldots i_p})_{p!}\mathcal{O}_{U_{i_0 \ldots i_p}}$. En outre, le différentiel du complexe $K^{ext}_\bullet(W)$
est donné par
$$
d(s)_{i_0 \ldots i_p} =
\sum\nolimits_{j = 0, \ldots, p + 1} \sum\nolimits_{i \in I_1}
(-1)^j s_{i_0 \ldots i_{j - 1} i i_j \ldots i_p}.
$$
La somme est finie, car l'élément $s$ est à support fini. Fixons un élément $i_{\text{fix}} \in I_1$. Définissons une application
$$
h : K^{ext}_p(W) \longrightarrow K^{ext}_{p + 1}(W)
$$
par la règle
$$
h(s)_{i_0 \ldots i_{p + 1}} =
\left\{
\begin{matrix}
0 & \text{si} & i_0 \not = i_{\text{fix}} \\
s_{i_1 \ldots i_{p + 1}} & \text{si} & i_0 = i_{\text{fix}}
\end{matrix}
\right.
$$
Nous utiliserons la notation abrégée
$h(s)_{i_0 \ldots i_{p + 1}} = (i_0 = i_{\text{fix}}) s_{i_1 \ldots i_p}$
pour cela. Alors nous calculons

\begin{eqnarray*}
& & (dh + hd)(s)_{i_0 \ldots i_p} \\
& = &
\sum_j \sum_{i \in I_1} (-1)^j h(s)_{i_0 \ldots i_{j - 1} i i_j \ldots i_p}
+
(i = i_0) d(s)_{i_1 \ldots i_p} \\
& = &
s_{i_0 \ldots i_p} +
\sum_{j \geq 1}\sum_{i \in I_1}
(-1)^j (i_0 = i_{\text{fix}}) s_{i_1 \ldots i_{j - 1} i i_j \ldots i_p}
+
(i_0 = i_{\text{fix}}) d(s)_{i_1 \ldots i_p}
\end{eqnarray*}

qui est égal à $s_{i_0 \ldots i_p}$, comme voulu.

\end{proof}

\begin{lemma}

\label{cohomology-lemma-cech-cohomology-derived-presheaves}
Soit $X$ un espace annelé. Soit $\mathcal{U} : U = \bigcup_{i \in I} U_i$ un recouvrement ouvert de $U \subset X$. Les foncteurs de cohomologie de {\v C}ech $\check{H}^p(\mathcal{U}, -)$ sont canoniquement isomorphes, en tant que $\delta$-foncteur, aux foncteurs dérivés à droite du foncteur $$
\check{H}^0(\mathcal{U}, -) :
\textit{PMod}(\mathcal{O}_X)
\longrightarrow
\text{Mod}_{\mathcal{O}_X(U)}.
$$ En outre, il existe un quasi-isomorphisme fonctoriel $$
\check{\mathcal{C}}^\bullet(\mathcal{U}, \mathcal{F})
\longrightarrow
R\check{H}^0(\mathcal{U}, \mathcal{F})
$$ où le membre de droite désigne le foncteur dérivé à droite $$
R\check{H}^0(\mathcal{U}, -) :
D^{+}(\textit{PMod}(\mathcal{O}_X))
\longrightarrow
D^{+}(\mathcal{O}_X(U))
$$ du foncteur exact à gauche $\check{H}^0(\mathcal{U}, -)$.
\end{lemma}


\begin{proof}
Notons que la catégorie des préfaisceaux de $\mathcal{O}_X$-modules a suffisamment d'injectifs, voir Injectifs, proposition \ref{injectives-proposition-presheaves-modules}. Notons aussi que $\check{H}^0(\mathcal{U}, -)$ est un foncteur exact à gauche de la catégorie des préfaisceaux de $\mathcal{O}_X$-modules vers celle des $\mathcal{O}_X(U)$-modules. Le foncteur dérivé et le foncteur dérivé à droite existent donc, voir Catégories dérivées, section \ref{derived-section-right-derived-functor}.

\medskip\noindent

Soit $\mathcal{I}$ un préfaisceau injectif de $\mathcal{O}_X$-modules. Dans ce cas, le foncteur $\Hom_{\mathcal{O}_X}(-, \mathcal{I})$
est exact sur $\textit{PMod}(\mathcal{O}_X)$. D'après le lemme \ref{cohomology-lemma-cech-map-into}, nous avons
$$
\Hom_{\mathcal{O}_X}(K(\mathcal{U})_\bullet, \mathcal{I})
=
\check{\mathcal{C}}^\bullet(\mathcal{U}, \mathcal{I}).
$$
D'après le lemme \ref{cohomology-lemma-homology-complex}, nous savons que $K(\mathcal{U})_\bullet$ est quasi-isomorphe à $\mathcal{O}_\mathcal{U}[0]$. Par l'exactitude du foncteur Hom à valeurs dans $\mathcal{I}$ mentionnée ci-dessus, nous obtenons donc $\check{H}^i(\mathcal{U}, \mathcal{I}) = 0$ pour tous
$i > 0$. Ainsi, le $\delta$-foncteur $(\check{H}^n, \delta)$
(voir lemme \ref{cohomology-lemma-cech-cohomology-delta-functor-presheaves}) satisfait aux hypothèses de Homologie, lemme \ref{homology-lemma-efface-implies-universal}, et est donc un $\delta$-foncteur universel.

\medskip\noindent
D'après Catégories dérivées, lemme \ref{derived-lemma-higher-derived-functors}, la suite $R^i\check{H}^0(\mathcal{U}, -)$ forme également un $\delta$-foncteur universel. Par l'unicité des $\delta$-foncteurs universels, voir Homologie, lemme \ref{homology-lemma-uniqueness-universal-delta-functor}, nous concluons que $R^i\check{H}^0(\mathcal{U}, -) = \check{H}^i(\mathcal{U}, -)$. Cela suffit pour la plupart des applications, et l'on suggère au lecteur de sauter le reste de la démonstration.

\medskip\noindent

Soit $\mathcal{F}$ un préfaisceau quelconque de $\mathcal{O}_X$-modules. Choisissons une résolution injective $\mathcal{F} \to \mathcal{I}^\bullet$
dans la catégorie $\textit{PMod}(\mathcal{O}_X)$. Considérons le double complexe
$\check{\mathcal{C}}^\bullet(\mathcal{U}, \mathcal{I}^\bullet)$ avec des termes
$\check{\mathcal{C}}^p(\mathcal{U}, \mathcal{I}^q)$. Considérons le complexe total associé
$\text{Tot}(\check{\mathcal{C}}^\bullet(\mathcal{U}, \mathcal{I}^\bullet))$, voir Homologie, définition \ref{homology-definition-associated-simple-complex}. Il existe un morphisme de complexes
$$
\check{\mathcal{C}}^\bullet(\mathcal{U}, \mathcal{F})
\longrightarrow
\text{Tot}(\check{\mathcal{C}}^\bullet(\mathcal{U}, \mathcal{I}^\bullet))
$$
provenant des applications
$\check{\mathcal{C}}^p(\mathcal{U}, \mathcal{F})
\to \check{\mathcal{C}}^p(\mathcal{U}, \mathcal{I}^0)$
et il existe un morphisme de complexes
$$
\check{H}^0(\mathcal{U}, \mathcal{I}^\bullet)
\longrightarrow
\text{Tot}(\check{\mathcal{C}}^\bullet(\mathcal{U}, \mathcal{I}^\bullet))
$$
provenant des applications
$\check{H}^0(\mathcal{U}, \mathcal{I}^q) \to
\check{\mathcal{C}}^0(\mathcal{U}, \mathcal{I}^q)$. Ces deux applications sont des quasi-isomorphismes d'après Homologie, lemme \ref{homology-lemma-double-complex-gives-resolution}. En effet, les colonnes du double complexe sont exactes en degrés positifs parce que le complexe de {\v C}ech est exact en tant que foncteur (lemme \ref{cohomology-lemma-cech-exact-presheaves}), et les lignes du double complexe sont exactes en degrés positifs puisque, comme nous venons de le voir, les groupes de cohomologie supérieure de {\v C}ech des préfaisceaux injectifs $\mathcal{I}^q$ sont nuls. Comme les quasi-isomorphismes deviennent inversibles dans $D^{+}(\mathcal{O}_X(U))$, on obtient le dernier morphisme affiché du lemme. Nous omettons de vérifier que ce morphisme est fonctoriel.

\end{proof}





\section{Cohomologie de {\v C}
ech et cohomologie}
\label{cohomology-section-cech-cohomology-cohomology}

\begin{lemma}

\label{cohomology-lemma-injective-trivial-cech}
Soit $X$ un espace annelé. Soit $\mathcal{U} : U = \bigcup_{i \in I} U_i$ un recouvrement ouvert. Soit $\mathcal{I}$ un $\mathcal{O}_X$-module injectif. Alors $$
\check{H}^p(\mathcal{U}, \mathcal{I}) =
\left\{
\begin{matrix}
\mathcal{I}(U) & \text{si} & p = 0 \\
0 & \text{si} & p > 0
\end{matrix}
\right.
$$
\end{lemma}

\begin{proof}

Un $\mathcal{O}_X$-module injectif est également injectif comme objet de la catégorie $\textit{PMod}(\mathcal{O}_X)$ (par exemple, puisque la faisceautisation est un adjoint à gauche exact du foncteur d'inclusion, en utilisant Homologie, lemme \ref{homology-lemma-adjoint-preserve-injectives}). Nous pouvons donc appliquer le lemme \ref{cohomology-lemma-cech-cohomology-derived-presheaves}
(ou sa preuve) pour voir le résultat.

\end{proof}

\begin{lemma}

\label{cohomology-lemma-cech-cohomology}
Soit $X$ un espace annelé. Soit $\mathcal{U} : U = \bigcup_{i \in I} U_i$ un recouvrement ouvert. Il existe une transformation $$
\check{\mathcal{C}}^\bullet(\mathcal{U}, -)
\longrightarrow
R\Gamma(U, -)
$$ des foncteurs $\textit{Mod}(\mathcal{O}_X) \to D^{+}(\mathcal{O}_X(U))$. En particulier, cela fournit des applications canoniques $\check{H}^p(\mathcal{U}, \mathcal{F}) \to H^p(U, \mathcal{F})$ lorsque $\mathcal{F}$ parcourt $\textit{Mod}(\mathcal{O}_X)$.
\end{lemma}

\begin{proof}

Soit $\mathcal{F}$ un $\mathcal{O}_X$-module. Choisissons une résolution injective
$\mathcal{F} \to \mathcal{I}^\bullet$. Considérons le double complexe
$\check{\mathcal{C}}^\bullet(\mathcal{U}, \mathcal{I}^\bullet)$ avec des termes
$\check{\mathcal{C}}^p(\mathcal{U}, \mathcal{I}^q)$. Il existe un morphisme de complexes
$$
\alpha :
\Gamma(U, \mathcal{I}^\bullet)
\longrightarrow
\text{Tot}(\check{\mathcal{C}}^\bullet(\mathcal{U}, \mathcal{I}^\bullet))
$$
provenant des applications
$\mathcal{I}^q(U) \to \check{H}^0(\mathcal{U}, \mathcal{I}^q)$
et un morphisme de complexes
$$
\beta :
\check{\mathcal{C}}^\bullet(\mathcal{U}, \mathcal{F})
\longrightarrow
\text{Tot}(\check{\mathcal{C}}^\bullet(\mathcal{U}, \mathcal{I}^\bullet))
$$
provenant de l'application $\mathcal{F} \to \mathcal{I}^0$. Le lemme \ref{homology-lemma-double-complex-gives-resolution} d'Homologie montre que
$\alpha$ est un quasi-isomorphisme. En effet, le lemme \ref{cohomology-lemma-injective-trivial-cech} implique que la $q$-ième ligne du double complexe
$\check{\mathcal{C}}^\bullet(\mathcal{U}, \mathcal{I}^\bullet)$ est une résolution de $\Gamma(U, \mathcal{I}^q)$. Ainsi, $\alpha$ devient inversible dans $D^{+}(\mathcal{O}_X(U))$ et la transformation du lemme est la composition de $\beta$
suivie de l'inverse de $\alpha$. Nous omettons de vérifier que cette construction est fonctorielle.

\end{proof}

\begin{lemma}

\label{cohomology-lemma-cech-h1}
Soit $X$ un espace topologique. Soit $\mathcal{H}$ un faisceau abélien sur $X$. Soit $\mathcal{U} : X = \bigcup_{i \in I} U_i$ un recouvrement ouvert. L'application $$
\check{H}^1(\mathcal{U}, \mathcal{H}) \longrightarrow H^1(X, \mathcal{H})
$$ est injective et identifie $\check{H}^1(\mathcal{U}, \mathcal{H})$, au moyen de la bijection du lemme \ref{cohomology-lemma-torsors-h1}, à l'ensemble des classes d'isomorphisme de $\mathcal{H}$-torseurs dont la restriction à chaque $U_i$ est triviale.
\end{lemma}

\begin{proof}

Pour le voir, construisons une application inverse. Soit $\mathcal{F}$ un
$\mathcal{H}$-torseur dont la restriction à $U_i$ est triviale. D'après le lemme \ref{cohomology-lemma-trivial-torsor}, cela signifie qu'il existe une section $s_i \in \mathcal{F}(U_i)$. Sur $U_{i_0} \cap U_{i_1}$
il existe une section unique $s_{i_0i_1}$ de $\mathcal{H}$ telle que
$s_{i_0i_1} \cdot s_{i_0}|_{U_{i_0} \cap U_{i_1}} =
s_{i_1}|_{U_{i_0} \cap U_{i_1}}$. Un calcul montre que $s_{i_0i_1}$ est un cocycle de {\v C}ech et que sa classe est bien définie (c.-à-d. ne dépend pas du choix des sections $s_i$). L'application inverse envoie la classe d'isomorphisme de $\mathcal{F}$ sur la classe de cohomologie du cocycle $(s_{i_0i_1})$. Nous omettons de vérifier que cette application est bien un inverse.

\end{proof}

\begin{lemma}

\label{cohomology-lemma-include}
Soit $X$ un espace annelé. Considérons le foncteur $i : \textit{Mod}(\mathcal{O}_X) \to \textit{PMod}(\mathcal{O}_X)$. C'est un foncteur exact à gauche dont les foncteurs dérivés à droite sont donnés par $$
R^pi(\mathcal{F}) = \underline{H}^p(\mathcal{F}) :
U \longmapsto H^p(U, \mathcal{F})
$$ voir la discussion dans la section \ref{cohomology-section-locality}.
\end{lemma}

\begin{proof}
Il est clair que $i$ est exact à gauche. Choisissons une résolution injective $\mathcal{F} \to \mathcal{I}^\bullet$. Par définition, $R^pi$ est le $p$-ième {\it préfaisceau} de cohomologie du complexe $\mathcal{I}^\bullet$. En d'autres termes, les sections de $R^pi(\mathcal{F})$ sur un ouvert $U$ sont données par $$
\frac{\Ker(\mathcal{I}^p(U) \to \mathcal{I}^{p + 1}(U))}
{\Im(\mathcal{I}^{p - 1}(U) \to \mathcal{I}^p(U))}.
$$ ce qui est la définition de $H^p(U, \mathcal{F})$.
\end{proof}

\begin{lemma}

\label{cohomology-lemma-cech-spectral-sequence}
Soit $X$ un espace annelé. Soit $\mathcal{U} : U = \bigcup_{i \in I} U_i$ un recouvrement ouvert. Pour tout faisceau de $\mathcal{O}_X$-modules $\mathcal{F}$, il existe une suite spectrale $(E_r, d_r)_{r \geq 0}$ avec $$
E_2^{p, q} = \check{H}^p(\mathcal{U}, \underline{H}^q(\mathcal{F}))
$$ qui converge vers $H^{p + q}(U, \mathcal{F})$. Cette suite spectrale est fonctorielle en $\mathcal{F}$.
\end{lemma}

\begin{proof}
C'est une suite spectrale de Grothendieck (voir Catégories dérivées, lemme \ref{derived-lemma-grothendieck-spectral-sequence}) pour les foncteurs $$
i :  \textit{Mod}(\mathcal{O}_X) \to \textit{PMod}(\mathcal{O}_X)
\quad\text{et}\quad
\check{H}^0(\mathcal{U}, - ) : \textit{PMod}(\mathcal{O}_X)
\to \text{Mod}_{\mathcal{O}_X(U)}.
$$ Nous avons $\check{H}^0(\mathcal{U}, i(\mathcal{F})) = \mathcal{F}(U)$ d'après le lemme \ref{cohomology-lemma-cech-h0}. Le préfaisceau $i(\mathcal{I})$ est acyclique pour la cohomologie de {\v C}ech d'après le lemme \ref{cohomology-lemma-injective-trivial-cech}. Enfin, $\check{H}^p(\mathcal{U}, -) = R^p\check{H}^0(\mathcal{U}, -)$ comme foncteurs sur $\textit{PMod}(\mathcal{O}_X)$ d'après le lemme \ref{cohomology-lemma-cech-cohomology-derived-presheaves}. L'assemblage de ces faits donne le lemme.
\end{proof}

\begin{lemma}

\label{cohomology-lemma-cech-spectral-sequence-application}
Soit $X$ un espace annelé. Soit $\mathcal{U} : U = \bigcup_{i \in I} U_i$ un recouvrement ouvert. Soit $\mathcal{F}$ un $\mathcal{O}_X$-module. Supposons que $H^i(U_{i_0 \ldots i_p}, \mathcal{F}) = 0$ pour tout $i > 0$, tout $p \geq 0$ et tous $i_0, \ldots, i_p \in I$. Alors $\check{H}^p(\mathcal{U}, \mathcal{F}) = H^p(U, \mathcal{F})$ comme $\mathcal{O}_X(U)$-modules.
\end{lemma}

\begin{proof}

Nous utiliserons la suite spectrale du lemme \ref{cohomology-lemma-cech-spectral-sequence}. Les hypothèses signifient que $E_2^{p, q} = 0$ pour tous $(p, q)$ avec
$q \not = 0$. La suite spectrale dégénère donc en $E_2$,
et le résultat s'ensuit.

\end{proof}

\begin{lemma}

\label{cohomology-lemma-ses-cech-h1}
Soit $X$ un espace annelé. Soit $$
0 \to \mathcal{F} \to \mathcal{G} \to \mathcal{H} \to 0
$$ une suite exacte courte de $\mathcal{O}_X$-modules. Soit $U \subset X$ un ouvert. S'il existe un système cofinal de recouvrements ouverts $\mathcal{U}$ de $U$ tel que $\check{H}^1(\mathcal{U}, \mathcal{F}) = 0$, alors l'application $\mathcal{G}(U) \to \mathcal{H}(U)$ est surjective.
\end{lemma}

\begin{proof}

Prenons un élément $s \in \mathcal{H}(U)$. Choisissons un recouvrement ouvert
$\mathcal{U} : U = \bigcup_{i \in I} U_i$ tel que a) $\check{H}^1(\mathcal{U}, \mathcal{F}) = 0$ et b)
$s|_{U_i}$ soit l'image d'une section $s_i \in \mathcal{G}(U_i)$. Puisque nous pouvons certainement trouver $\mathcal{U}$ satisfaisant à (b), les hypothèses du lemme montrent que nous pouvons trouver
$\mathcal{U}$ satisfaisant à la fois à a) et à b). Considérons les sections
$$
s_{i_0i_1} = s_{i_1}|_{U_{i_0i_1}} - s_{i_0}|_{U_{i_0i_1}}.
$$
Puisque $s_i$ relève $s$, nous voyons que $s_{i_0i_1} \in \mathcal{F}(U_{i_0i_1})$. Par l'annulation de $\check{H}^1(\mathcal{U}, \mathcal{F})$, nous pouvons trouver des sections $t_i \in \mathcal{F}(U_i)$ telles que
$$
s_{i_0i_1} = t_{i_1}|_{U_{i_0i_1}} - t_{i_0}|_{U_{i_0i_1}}.
$$
Alors les sections $s_i - t_i$ satisfont clairement à la condition de recollement et se recollent en une section de $\mathcal{G}$ sur $U$ dont l'image est $s$. D'où le résultat.

\end{proof}

\begin{lemma}

\label{cohomology-lemma-cech-vanish}

\begin{slogan}
Si la cohomologie supérieure de {\v C}ech d'un faisceau abélien s'annule pour tous les recouvrements ouverts, alors la cohomologie supérieure s'annule.
\end{slogan}
Soit $X$ un espace annelé. Soit $\mathcal{F}$ un $\mathcal{O}_X$-module tel que $$
\check{H}^p(\mathcal{U}, \mathcal{F}) = 0
$$ pour tout $p > 0$ et tout recouvrement ouvert $\mathcal{U} : U = \bigcup_{i \in I} U_i$ d'un ouvert de $X$. Alors $H^p(U, \mathcal{F}) = 0$ pour tout $p > 0$ et tout ouvert $U \subset X$.
\end{lemma}

\begin{proof}

Soit $\mathcal{F}$ un faisceau satisfaisant à l'hypothèse du lemme. Nous exprimerons cette propriété en disant que « la cohomologie supérieure de {\v C}ech de $\mathcal{F}$ s'annule pour tout recouvrement ouvert ». Choisissons un plongement $\mathcal{F} \to \mathcal{I}$ dans un $\mathcal{O}_X$-module injectif. D'après le lemme \ref{cohomology-lemma-injective-trivial-cech}, la cohomologie supérieure de {\v C}ech de $\mathcal{I}$ s'annule pour tout recouvrement ouvert. Posons $\mathcal{Q} = \mathcal{I}/\mathcal{F}$
de sorte que nous avons une suite exacte courte
$$
0 \to \mathcal{F} \to \mathcal{I} \to \mathcal{Q} \to 0.
$$
D'après le lemme \ref{cohomology-lemma-ses-cech-h1} et nos hypothèses, cette suite est même exacte comme suite de préfaisceaux ! En particulier, nous disposons d'une suite exacte longue de groupes de cohomologie de {\v C}ech pour tout recouvrement ouvert $\mathcal{U}$ ; voir, par exemple, le lemme \ref{cohomology-lemma-cech-cohomology-delta-functor-presheaves}.
Il en résulte que $\mathcal{Q}$ est lui aussi un $\mathcal{O}_X$-module dont la cohomologie supérieure de {\v C}ech s'annule pour tout recouvrement ouvert.

\medskip\noindent

Ensuite, nous examinons la suite exacte longue de cohomologie
$$
\xymatrix{
0 \ar[r] &
H^0(U, \mathcal{F}) \ar[r] &
H^0(U, \mathcal{I}) \ar[r] &
H^0(U, \mathcal{Q}) \ar[lld] \\
&
H^1(U, \mathcal{F}) \ar[r] &
H^1(U, \mathcal{I}) \ar[r] &
H^1(U, \mathcal{Q}) \ar[lld] \\
&
\ldots & \ldots & \ldots \\
}
$$
pour tout ouvert $U \subset X$. Puisque $\mathcal{I}$ est injectif, nous avons $H^n(U, \mathcal{I}) = 0$ pour $n > 0$ (voir Catégories dérivées, lemme \ref{derived-lemma-higher-derived-functors}). D'après ce qui précède, l'application $H^0(U, \mathcal{I}) \to H^0(U, \mathcal{Q})$
est surjective, donc $H^1(U, \mathcal{F}) = 0$. Comme $\mathcal{F}$ était un $\mathcal{O}_X$-module arbitraire dont la cohomologie supérieure de {\v C}ech s'annule, nous en déduisons également que
$H^1(U, \mathcal{Q}) = 0$, puisque $\mathcal{Q}$ est un autre faisceau de ce type (voir ci-dessus). La suite exacte longue entraîne alors $H^2(U, \mathcal{F}) = 0$, et l'on poursuit de même.

\end{proof}

\begin{lemma}

\label{cohomology-lemma-cech-vanish-basis}

(Variante du lemme \ref{cohomology-lemma-cech-vanish}.) Soit $X$ un espace annelé. Soit $\mathcal{B}$ une base de la topologie de $X$. Soit $\mathcal{F}$ un $\mathcal{O}_X$-module. Supposons qu'il existe un ensemble de recouvrements ouverts $\text{Cov}$
ayant les propriétés suivantes :

\begin{enumerate}

\item Pour chaque $\mathcal{U} \in \text{Cov}$
avec $\mathcal{U} : U = \bigcup_{i \in I} U_i$, nous avons
$U, U_i \in \mathcal{B}$ et chaque $U_{i_0 \ldots i_p} \in \mathcal{B}$.
\item Pour chaque $U \in \mathcal{B}$, les recouvrements ouverts de $U$
qui appartiennent à $\text{Cov}$ forment un système cofinal de recouvrements ouverts de $U$.
\item Pour chaque $\mathcal{U} \in \text{Cov}$, nous avons
$\check{H}^p(\mathcal{U}, \mathcal{F}) = 0$ pour tout $p > 0$.

\end{enumerate}

Alors $H^p(U, \mathcal{F}) = 0$ pour tout $p > 0$ et tout $U \in \mathcal{B}$.

\end{lemma}

\begin{proof}

Soient $\mathcal{F}$ et $\text{Cov}$ comme dans le lemme. Nous exprimerons cette propriété en disant que « la cohomologie supérieure de {\v C}ech de $\mathcal{F}$ s'annule pour tout $\mathcal{U} \in \text{Cov}$ ». Choisissons un plongement $\mathcal{F} \to \mathcal{I}$ dans un $\mathcal{O}_X$-module injectif. D'après le lemme \ref{cohomology-lemma-injective-trivial-cech}, la cohomologie supérieure de {\v C}ech de $\mathcal{I}$
s'annule pour tout $\mathcal{U} \in \text{Cov}$. Posons $\mathcal{Q} = \mathcal{I}/\mathcal{F}$
de sorte que nous avons une suite exacte courte
$$
0 \to \mathcal{F} \to \mathcal{I} \to \mathcal{Q} \to 0.
$$
D'après le lemme \ref{cohomology-lemma-ses-cech-h1} et notre hypothèse (2), cette suite donne lieu à une suite exacte
$$
0 \to \mathcal{F}(U) \to \mathcal{I}(U) \to \mathcal{Q}(U) \to 0.
$$
pour chaque $U \in \mathcal{B}$. Par conséquent, pour tout $\mathcal{U} \in \text{Cov}$,
nous obtenons une suite exacte courte de complexes de {\v C}ech
$$
0 \to
\check{\mathcal{C}}^\bullet(\mathcal{U}, \mathcal{F}) \to
\check{\mathcal{C}}^\bullet(\mathcal{U}, \mathcal{I}) \to
\check{\mathcal{C}}^\bullet(\mathcal{U}, \mathcal{Q}) \to 0
$$
car chaque terme du complexe de {\v C}ech est constitué d'un produit de valeurs sur des éléments de $\mathcal{B}$, d'après l'hypothèse (1). En particulier, nous avons une suite exacte longue de groupes de cohomologie de {\v C}ech pour tout recouvrement ouvert $\mathcal{U} \in \text{Cov}$. Cela implique que $\mathcal{Q}$ est également un $\mathcal{O}_X$-module dont la cohomologie supérieure de {\v C}ech s'annule pour tous
$\mathcal{U} \in \text{Cov}$.

\medskip\noindent

Ensuite, nous examinons la suite exacte longue de cohomologie
$$
\xymatrix{
0 \ar[r] &
H^0(U, \mathcal{F}) \ar[r] &
H^0(U, \mathcal{I}) \ar[r] &
H^0(U, \mathcal{Q}) \ar[lld] \\
&
H^1(U, \mathcal{F}) \ar[r] &
H^1(U, \mathcal{I}) \ar[r] &
H^1(U, \mathcal{Q}) \ar[lld] \\
&
\ldots & \ldots & \ldots \\
}
$$
pour tout $U \in \mathcal{B}$. Puisque $\mathcal{I}$ est injectif, nous avons $H^n(U, \mathcal{I}) = 0$ pour $n > 0$ (voir Catégories dérivées, lemme \ref{derived-lemma-higher-derived-functors}). D'après ce qui précède, l'application $H^0(U, \mathcal{I}) \to H^0(U, \mathcal{Q})$
est surjective, donc $H^1(U, \mathcal{F}) = 0$. Comme $\mathcal{F}$ était un $\mathcal{O}_X$-module arbitraire dont la cohomologie supérieure de {\v C}ech s'annule pour tout $\mathcal{U} \in \text{Cov}$,
nous en déduisons que $H^1(U, \mathcal{Q}) = 0$, puisque $\mathcal{Q}$ est un autre faisceau de ce type (voir ci-dessus). La suite exacte longue entraîne alors $H^2(U, \mathcal{F}) = 0$, et l'on poursuit de même.

\end{proof}

\begin{lemma}

\label{cohomology-lemma-pushforward-injective}

Soit $f : X \to Y$ un morphisme d'espaces annelés. Soit $\mathcal{I}$ un $\mathcal{O}_X$-module injectif. Alors :

\begin{enumerate}

\item  $\check{H}^p(\mathcal{V}, f_*\mathcal{I}) = 0$
pour tout $p > 0$ et tout recouvrement ouvert
$\mathcal{V} : V = \bigcup_{j \in J} V_j$ de $Y$.
\item  $H^p(V, f_*\mathcal{I}) = 0$ pour tout $p > 0$ et tout ouvert $V \subset Y$.

\end{enumerate}

En d'autres termes, $f_*\mathcal{I}$ est acyclique à droite pour $\Gamma(V, -)$
(voir Catégories dérivées, définition \ref{derived-definition-derived-functor}) pour tout ouvert $V \subset Y$.

\end{lemma}

\begin{proof}
Posons $\mathcal{U} : f^{-1}(V) = \bigcup_{j \in J} f^{-1}(V_j)$. Il s'agit d'un recouvrement par des ouverts de $X$ et $$
\check{\mathcal{C}}^\bullet(\mathcal{V}, f_*\mathcal{I}) =
\check{\mathcal{C}}^\bullet(\mathcal{U}, \mathcal{I}).
$$ Ceci est vrai parce que $$
f_*\mathcal{I}(V_{j_0 \ldots j_p})
= \mathcal{I}(f^{-1}(V_{j_0 \ldots j_p})) =
\mathcal{I}(f^{-1}(V_{j_0}) \cap \ldots \cap f^{-1}(V_{j_p}))
= \mathcal{I}(U_{j_0 \ldots j_p}).
$$ Ainsi, la première affirmation du lemme résulte du lemme \ref{cohomology-lemma-injective-trivial-cech}. La seconde résulte de la première et du lemme \ref{cohomology-lemma-cech-vanish}.
\end{proof}

\noindent
Le lemme suivant implique en particulier que $f_* : \textit{Ab}(X) \to \textit{Ab}(Y)$ transforme les faisceaux abéliens injectifs en faisceaux abéliens injectifs.

\begin{lemma}

\label{cohomology-lemma-pushforward-injective-flat}
Soit $f : X \to Y$ un morphisme d'espaces annelés. Supposons que $f$ soit plat. Alors $f_*\mathcal{I}$ est un $\mathcal{O}_Y$-module injectif pour tout $\mathcal{O}_X$-module injectif $\mathcal{I}$.
\end{lemma}

\begin{proof}

Dans ce cas, le foncteur $f^*$ transforme les monomorphismes en monomorphismes (Faisceaux de modules, lemme \ref{modules-lemma-pullback-flat}). Le résultat découle alors d'Homologie, lemme \ref{homology-lemma-adjoint-preserve-injectives}.

\end{proof}

\begin{lemma}

\label{cohomology-lemma-cohomology-products}
Soit $(X, \mathcal{O}_X)$ un espace annelé et soit $I$ un ensemble. Pour $i \in I$, soit $\mathcal{F}_i$ un $\mathcal{O}_X$-module. Soit $U \subset X$ un ouvert. L'application canonique $$
H^p(U, \prod\nolimits_{i \in I} \mathcal{F}_i)
\longrightarrow
\prod\nolimits_{i \in I} H^p(U, \mathcal{F}_i)
$$ est un isomorphisme pour $p = 0$ et injective pour $p = 1$.
\end{lemma}

\begin{proof}

L'assertion pour $p = 0$ est vraie parce que le produit des faisceaux est égal au produit des préfaisceaux sous-jacents ; voir Faisceaux, section \ref{sheaves-section-limits-sheaves}. Prouvons l'assertion pour $p = 1$. Posons $\mathcal{F} = \prod \mathcal{F}_i$. Soit $\xi \in H^1(U, \mathcal{F})$ dont l'image est nulle dans
$\prod H^1(U, \mathcal{F}_i)$. Par le caractère local de la cohomologie, voir lemme \ref{cohomology-lemma-kill-cohomology-class-on-covering}, il existe un recouvrement ouvert $\mathcal{U} : U = \bigcup U_j$ tel que
$\xi|_{U_j} = 0$ pour tout $j$. D'après le lemme \ref{cohomology-lemma-cech-h1}, cela signifie que
$\xi$ provient d'un élément
$\check \xi \in \check H^1(\mathcal{U}, \mathcal{F})$. Puisque les applications
$\check H^1(\mathcal{U}, \mathcal{F}_i) \to H^1(U, \mathcal{F}_i)$
sont injectives pour tout $i$ (par le lemme \ref{cohomology-lemma-cech-h1}) et, puisque l'image de $\xi$ est nulle dans $\prod H^1(U, \mathcal{F}_i)$, nous voyons que l'image
$\check \xi_i = 0$ dans $\check H^1(\mathcal{U}, \mathcal{F}_i)$. Cependant, puisque $\mathcal{F} = \prod \mathcal{F}_i$, nous voyons que
$\check{\mathcal{C}}^\bullet(\mathcal{U}, \mathcal{F})$ est le produit des complexes
$\check{\mathcal{C}}^\bullet(\mathcal{U}, \mathcal{F}_i)$, donc par Homologie, lemme \ref{homology-lemma-product-abelian-groups-exact}
nous concluons que $\check \xi = 0$, comme voulu.

\end{proof}







\section{Faisceaux flasques}

\label{cohomology-section-flasque}

\noindent
Voici la définition.

\begin{definition}

\label{cohomology-definition-flasque}
Soit $X$ un espace topologique. On dit qu'un préfaisceau d'ensembles $\mathcal{F}$ est {\it flasque}, ou {\it mou}, si, pour tous ouverts $U \subset V$ de $X$, l'application de restriction $\mathcal{F}(V) \to \mathcal{F}(U)$ est surjective.
\end{definition}

\noindent
Nous utiliserons également cette terminologie pour les faisceaux abéliens et les faisceaux de modules lorsque $X$ est un espace annelé. Il suffit clairement de supposer que les applications de restriction $\mathcal{F}(X) \to \mathcal{F}(U)$ sont surjectives pour chaque ouvert $U \subset X$.

\begin{lemma}

\label{cohomology-lemma-injective-flasque}
Soit $(X, \mathcal{O}_X)$ un espace annelé. Alors tout $\mathcal{O}_X$-module injectif est flasque.
\end{lemma}

\begin{proof}
C'est une reformulation du lemme \ref{cohomology-lemma-injective-restriction-surjective}.
\end{proof}

\begin{lemma}

\label{cohomology-lemma-flasque-acyclic}
Soit $(X, \mathcal{O}_X)$ un espace annelé. Tout $\mathcal{O}_X$-module flasque est acyclique pour $R\Gamma(X, -)$ ainsi que pour $R\Gamma(U, -)$, pour tout ouvert $U$ de $X$.
\end{lemma}

\begin{proof}

Nous utiliserons Catégories dérivées, lemme \ref{derived-lemma-subcategory-right-acyclics}. Comme tout module injectif est flasque, chaque $\mathcal{O}_X$-module se plonge dans un module flasque ; voir Injectifs, lemme \ref{injectives-lemma-abelian-sheaves-space}. Il suffit donc de montrer que, pour toute suite exacte courte
$$
0 \to \mathcal{F} \to \mathcal{G} \to \mathcal{H} \to 0
$$
avec $\mathcal{F}$ et $\mathcal{G}$ flasques, $\mathcal{H}$
est flasque et la suite reste exacte après prise de sections sur tout ouvert de $X$. En fait, la seconde assertion implique la première. Soit $U \subset X$ un sous-espace ouvert. Soit $s \in \mathcal{H}(U)$. Nous montrerons que nous pouvons relever $s$ en une section de $\mathcal{G}$
sur $U$. Pour ce faire, considérons l'ensemble $T$ des couples $(V, t)$
où $V \subset U$ est ouvert et $t \in \mathcal{G}(V)$ est une section envoyée sur $s|_V$ dans $\mathcal{H}$. Munissons $T$ de l'ordre partiel défini par
$(V, t) \leq (V', t')$ si et seulement si $V \subset V'$ et $t'|_V = t$. Si $(V_\alpha, t_\alpha)$, $\alpha \in A$,
est un sous-ensemble totalement ordonné de $T$, alors $V = \bigcup V_\alpha$
est ouvert et il existe une section unique $t \in \mathcal{G}(V)$
se restreignant à $t_\alpha$ sur $V_\alpha$, grâce à la condition de recollement pour
$\mathcal{G}$. Ainsi, le lemme de Zorn fournit un élément maximal
$(V, t)$ de $T$. Nous allons montrer que $V = U$, ce qui achèvera la démonstration. En effet, choisissons $x \in U$. Nous pouvons trouver un petit voisinage ouvert
$W \subset U$ de $x$ et $t' \in \mathcal{G}(W)$ dont l'image est $s|_W$
dans $\mathcal{H}$. Alors $t'|_{W \cap V} - t|_{W \cap V}$ a pour image zéro dans $\mathcal{H}$ et provient donc d'une section
$r' \in \mathcal{F}(W \cap V)$. Comme $\mathcal{F}$ est flasque, nous trouvons une section $r \in \mathcal{F}(W)$ dont la restriction est $r'$
sur $W \cap V$. En modifiant $t'$ par l'image de $r$, nous pouvons supposer que $t$ et $t'$ se restreignent à la même section sur
$W \cap V$. Grâce à la condition de recollement pour $\mathcal{G}$,
nous pouvons trouver une section $\tilde t$ de $\mathcal{G}$ sur
$W \cup V$ qui se restreint à $t$ et à $t'$. Par maximalité de $(V, t)$, nous voyons que $V \cup W = V$. Ainsi, $x \in V$ et la démonstration est achevée.

\end{proof}

\noindent
Le lemme suivant ne vaut pas pour les préfaisceaux flasques.

\begin{lemma}

\label{cohomology-lemma-flasque-acyclic-cech}
Soit $(X, \mathcal{O}_X)$ un espace annelé. Soit $\mathcal{F}$ un faisceau de $\mathcal{O}_X$-modules. Soit $\mathcal{U} : U = \bigcup U_i$ un recouvrement ouvert. Si $\mathcal{F}$ est flasque, alors $\check{H}^p(\mathcal{U}, \mathcal{F}) = 0$ pour $p > 0$.
\end{lemma}

\begin{proof}

Les préfaisceaux $\underline{H}^q(\mathcal{F})$ utilisés dans l'énoncé du lemme \ref{cohomology-lemma-cech-spectral-sequence} sont nuls d'après le lemme \ref{cohomology-lemma-flasque-acyclic}. Par conséquent, $\check{H}^p(U, \mathcal{F}) = H^p(U, \mathcal{F}) = 0$
d'après le lemme \ref{cohomology-lemma-flasque-acyclic}, encore une fois.

\end{proof}

\begin{lemma}

\label{cohomology-lemma-flasque-acyclic-pushforward}
Soit $f : (X, \mathcal{O}_X) \to (Y, \mathcal{O}_Y)$ un morphisme d'espaces annelés. Soit $\mathcal{F}$ un faisceau de $\mathcal{O}_X$-modules. Si $\mathcal{F}$ est flasque, alors $R^pf_*\mathcal{F} = 0$ pour $p > 0$.
\end{lemma}

\begin{proof}
Cela résulte immédiatement des lemmes \ref{cohomology-lemma-describe-higher-direct-images} et \ref{cohomology-lemma-flasque-acyclic}.
\end{proof}

\noindent
Le lemme suivant peut être démontré par un argument élémentaire d'induction pour les recouvrements finis ; comparer avec la discussion de la cohomologie de {\v C}ech dans \cite{FOAG}.

\begin{lemma}

\label{cohomology-lemma-vanishing-ravi}
Soit $X$ un espace topologique. Soit $\mathcal{F}$ un faisceau abélien sur $X$. Soit $\mathcal{U} : U = \bigcup_{i \in I} U_i$ un recouvrement ouvert. Supposons que les applications de restriction $\mathcal{F}(U) \to \mathcal{F}(U')$ soient surjectives pour toute union arbitraire $U'$ d'ouverts de la forme $U_{i_0 \ldots i_p}$. Alors $\check{H}^p(\mathcal{U}, \mathcal{F})$ s'annule pour $p > 0$.
\end{lemma}

\begin{proof}

Soit $Y$ l'ensemble de sous-ensembles non vides de $I$. Nous utiliserons les lettres
$A, B, C, \ldots$ pour désigner les éléments de $Y$, c'est-à-dire les sous-ensembles non vides de $I$. Pour un sous-ensemble fini non vide $J \subset I$, posons
$$
V_J = \{A \in Y \mid J \subset A\}
$$
Cela signifie que $V_{\{i\}} = \{A \in Y \mid i \in A\}$ et
$V_J = \bigcap_{j \in J} V_{\{j\}}$. Alors $V_J \subset V_K$ si et seulement si $J \supset K$. Il existe une topologie unique sur $Y$ telle que la collection des sous-ensembles $V_J$ forme une base de la topologie de $Y$. Tout ouvert est de la forme
$$
V = \bigcup\nolimits_{t \in T} V_{J_t}
$$
pour une famille de sous-ensembles finis $J_t$. Si $J_t \subset J_{t'}$,
alors nous pouvons retirer $J_{t'}$ de la famille sans modifier $V$. Ainsi, nous pouvons supposer qu'il n'y a pas d'inclusions parmi les $J_t$. Dans ce cas, les éléments minimaux de $V$ sont les ensembles $A = J_t$. Par conséquent, nous pouvons retrouver la famille $(J_t)_{t \in T}$ à partir de l'ouvert $V$.

\medskip\noindent

Nous pouvons décrire complètement les recouvrements ouverts de $Y$. Premièrement, puisque les éléments $A \in Y$ sont des sous-ensembles non vides de $I$, nous avons
$$
Y = \bigcup\nolimits_{i \in I} V_{\{i\}}
$$
Pour comprendre les autres recouvrements, soit $V$ comme ci-dessus et soit $V_s \subset Y$
un ouvert correspondant à la famille $(J_{s, t})_{t \in T_s}$. Alors
$$
V = \bigcup\nolimits_{s \in S} V_s
$$
si et seulement si, pour chaque $t \in T$, il existe un $s \in S$ et
$t_s \in T_s$ tel que $J_t = J_{s, t_s}$. Comme la famille
$(J_t)_{t \in T}$ est minimale, l'élément minimal $A = J_t$
doit appartenir à $V_s$ pour un certain $s$, donc $A \in V_{J_{t_s}}$ pour un certain
$t_s \in T_s$. Mais, comme $A$ est également minimal dans $V_s$, nous concluons que $J_{t_s} = J_t$.

\medskip\noindent

Ensuite, nous associons à chaque ouvert de $Y$ un ouvert de $X$. Plus précisément, nous envoyons
$Y$ sur $U$ et adoptons la règle
$$
V_J \mapsto U_J = \bigcap\nolimits_{i \in J} U_i
$$
sur les ouverts $V_J$, et nous l'étendons à des ouverts arbitraires $V$ par la règle
$$
V = \bigcup\nolimits_{t \in T} V_{J_t}
\mapsto
\bigcup\nolimits_{t \in T} U_{J_t}
$$
La classification des recouvrements ouverts de $Y$ ci-dessus montre que cette règle transforme les recouvrements ouverts en recouvrements ouverts. Nous obtenons ainsi un faisceau abélien $\mathcal{G}$ sur $Y$ en posant
$\mathcal{G}(Y) = \mathcal{F}(U)$ et pour
$V = \bigcup\nolimits_{t \in T} V_{J_t}$, en posant
$$
\mathcal{G}(V) = \mathcal{F}\left(\bigcup\nolimits_{t \in T} U_{J_t}\right)
$$
et en utilisant les applications de restriction de $\mathcal{F}$.

\medskip\noindent
Ces préliminaires étant établis, nous pouvons démontrer le lemme comme suit. Nous avons un recouvrement ouvert $\mathcal{V} : Y = \bigcup_{i \in I} V_{\{i\}}$ de $Y$. Par construction, nous avons une égalité $$
\check{C}^\bullet(\mathcal{V}, \mathcal{G}) =
\check{C}^\bullet(\mathcal{U}, \mathcal{F})
$$ de complexes de {\v C}ech. Puisque le faisceau $\mathcal{G}$ est flasque sur $Y$ (d'après notre hypothèse sur $\mathcal{F}$ dans l'énoncé du lemme), l'annulation résulte du lemme \ref{cohomology-lemma-flasque-acyclic-cech}.
\end{proof}



\section{La suite spectrale de Leray}

\label{cohomology-section-Leray}

\begin{lemma}

\label{cohomology-lemma-before-Leray}

Soit $f : X \to Y$ un morphisme d'espaces annelés. Il existe un diagramme commutatif
$$
\xymatrix{
D^{+}(X) \ar[rr]_-{R\Gamma(X, -)} \ar[d]_{Rf_*} & &
D^{+}(\mathcal{O}_X(X)) \ar[d]^{\text{restriction}} \\
D^{+}(Y) \ar[rr]^-{R\Gamma(Y, -)} & &
D^{+}(\mathcal{O}_Y(Y))
}
$$
Plus généralement, pour tout ouvert $V \subset Y$ et $U = f^{-1}(V)$, il existe un diagramme commutatif
$$
\xymatrix{
D^{+}(X) \ar[rr]_-{R\Gamma(U, -)} \ar[d]_{Rf_*} & &
D^{+}(\mathcal{O}_X(U)) \ar[d]^{\text{restriction}} \\
D^{+}(Y) \ar[rr]^-{R\Gamma(V, -)} & &
D^{+}(\mathcal{O}_Y(V))
}
$$
Voir aussi la remarque \ref{cohomology-remark-elucidate-lemma} pour plus d'explications.

\end{lemma}

\begin{proof}

Soit $\Gamma_{res} : \textit{Mod}(\mathcal{O}_X) \to \text{Mod}_{\mathcal{O}_Y(Y)}$
le foncteur qui associe à un $\mathcal{O}_X$-module $\mathcal{F}$
les sections globales de $\mathcal{F}$, considérées comme un $\mathcal{O}_Y(Y)$-module via l'application $f^\sharp : \mathcal{O}_Y(Y) \to \mathcal{O}_X(X)$. Soit $restriction : \text{Mod}_{\mathcal{O}_X(X)} \to \text{Mod}_{\mathcal{O}_Y(Y)}$
le foncteur de restriction des scalaires induit par
$f^\sharp : \mathcal{O}_Y(Y) \to \mathcal{O}_X(X)$. Notons que $restriction$
est exact, si bien que son foncteur dérivé à droite se calcule en appliquant simplement le foncteur de restriction, voir Catégories dérivées, lemme \ref{derived-lemma-right-derived-exact-functor}. Il est clair que
$$
\Gamma_{res}
=
restriction \circ \Gamma(X, -)
=
\Gamma(Y, -) \circ f_*
$$
Nous affirmons que le lemme \ref{derived-lemma-compose-derived-functors} de Catégories dérivées s'applique aux deux compositions.
Pour la première, cela résulte clairement des remarques ci-dessus. Pour la seconde, cela découle du lemme \ref{cohomology-lemma-pushforward-injective}, qui implique que les $\mathcal{O}_X$-modules injectifs sont envoyés sur des faisceaux acycliques pour $\Gamma(Y, -)$ dans la catégorie des faisceaux sur $Y$.

\end{proof}

\begin{remark}

\label{cohomology-remark-elucidate-lemma}
Voici une explication concrète du lemme \ref{cohomology-lemma-before-Leray}. Étant donnés $f : X \to Y$, $\mathcal{F} \in \textit{Mod}(\mathcal{O}_X)$ et une résolution injective $\mathcal{F} \to \mathcal{I}^\bullet$, on a $$
\begin{matrix}
R\Gamma(X, \mathcal{F}) & \text{est représenté par} &
\Gamma(X, \mathcal{I}^\bullet) \\
Rf_*\mathcal{F} & \text{est représenté par} & f_*\mathcal{I}^\bullet \\
R\Gamma(Y, Rf_*\mathcal{F}) & \text{est représenté par} &
\Gamma(Y, f_*\mathcal{I}^\bullet)
\end{matrix}
$$ où la dernière assertion résulte du lemme d'acyclicité de Leray (Catégories dérivées, lemme \ref{derived-lemma-leray-acyclicity}) et du lemme \ref{cohomology-lemma-pushforward-injective}. Enfin, combinons cela avec l'observation immédiate que $$
\Gamma(X, \mathcal{I}^\bullet)
=
\Gamma(Y, f_*\mathcal{I}^\bullet).
$$ ce qui donne la commutativité du diagramme du lemme.
\end{remark}

\begin{lemma}

\label{cohomology-lemma-modules-abelian}

Soit $X$ un espace annelé. Soit $\mathcal{F}$ un $\mathcal{O}_X$-module.

\begin{enumerate}

\item Les groupes de cohomologie $H^i(U, \mathcal{F})$, pour tout ouvert $U \subset X$, calculés en considérant $\mathcal{F}$ comme un $\mathcal{O}_X$-module ou comme un faisceau abélien sont identiques.
\item Soit $f : X \to Y$ un morphisme d'espaces annelés. Les images directes supérieures $R^if_*\mathcal{F}$ de $\mathcal{F}$
calculées en considérant celui-ci comme un $\mathcal{O}_X$-module ou comme un faisceau abélien sont identiques.

\end{enumerate}

Des énoncés analogues valent pour les complexes bornés inférieurement de $\mathcal{O}_X$-modules.

\end{lemma}

\begin{proof}
Considérons le morphisme d'espaces annelés $(X, \mathcal{O}_X) \to (X, \underline{\mathbf{Z}}_X)$ donné par l'identité sur l'espace topologique sous-jacent et par l'unique morphisme de faisceaux d'anneaux $\underline{\mathbf{Z}}_X \to \mathcal{O}_X$. Soit $\mathcal{F}$ un $\mathcal{O}_X$-module. Notons $\mathcal{F}_{ab}$ le même faisceau, vu comme un $\underline{\mathbf{Z}}_X$-module, c'est-à-dire comme un faisceau de groupes abéliens. Soit $\mathcal{F} \to \mathcal{I}^\bullet$ une résolution injective. D'après la remarque \ref{cohomology-remark-elucidate-lemma}, $\Gamma(X, \mathcal{I}^\bullet)$ calcule à la fois $R\Gamma(X, \mathcal{F})$ et $R\Gamma(X, \mathcal{F}_{ab})$. Ceci prouve (1).

\medskip\noindent

Pour prouver (2), nous utilisons (1) et le lemme \ref{cohomology-lemma-describe-higher-direct-images}. Le résultat suit immédiatement.

\end{proof}

\begin{lemma}
[suite spectrale de Leray]
\label{cohomology-lemma-Leray}
Soit $f : X \to Y$ un morphisme d'espaces annelés. Soit $\mathcal{F}^\bullet$ un complexe borné inférieurement de $\mathcal{O}_X$-modules. Il existe une suite spectrale $$
E_2^{p, q} = H^p(Y, R^qf_*(\mathcal{F}^\bullet))
$$ convergeant vers $H^{p + q}(X, \mathcal{F}^\bullet)$.
\end{lemma}

\begin{proof}
Il s'agit simplement de la suite spectrale de Grothendieck des Catégories dérivées, lemme \ref{derived-lemma-grothendieck-spectral-sequence}, issue de la composition des foncteurs $\Gamma_{res} = \Gamma(Y, -) \circ f_*$, où $\Gamma_{res}$ est celui de la démonstration du lemme \ref{cohomology-lemma-before-Leray}. Pour vérifier les hypothèses des Catégories dérivées, lemme \ref{derived-lemma-grothendieck-spectral-sequence}, voir la démonstration du lemme \ref{cohomology-lemma-before-Leray} ou la remarque \ref{cohomology-remark-elucidate-lemma}.
\end{proof}

\begin{remark}

\label{cohomology-remark-Leray-ss-more-structure}

La suite spectrale de Leray, telle que nous l'avons construite dans le lemme \ref{cohomology-lemma-Leray},
est une suite spectrale de $\Gamma(Y, \mathcal{O}_Y)$-modules. Il est toutefois assez facile de voir qu'il s'agit en fait d'une suite spectrale de
$\Gamma(X, \mathcal{O}_X)$-modules. Par exemple $f$ donne lieu à un morphisme d'espaces annelés
$f' :  (X, \mathcal{O}_X) \to (Y, f_*\mathcal{O}_X)$. D'après le lemme \ref{cohomology-lemma-modules-abelian}, les termes $E_r^{p, q}$ de la suite spectrale de Leray associés au $\mathcal{O}_X$-module $\mathcal{F}$
et $f$ sont identiques à ceux associés à $\mathcal{F}$ et $f'$
au moins pour $r \geq 2$. En effet, ils coïncident tous deux avec les termes de la suite spectrale de Leray de $\mathcal{F}$ considéré comme faisceau abélien. Puisque $(f_*\mathcal{O}_X)(Y) = \mathcal{O}_X(X)$, le résultat en découle. La suite spectrale de Leray possède souvent une structure supplémentaire.

\end{remark}

\begin{lemma}

\label{cohomology-lemma-apply-Leray}
Soit $f : X \to Y$ un morphisme d'espaces annelés. Soit $\mathcal{F}$ un $\mathcal{O}_X$-module.
\begin{enumerate}
\item Si $R^qf_*\mathcal{F} = 0$ pour $q > 0$, alors $H^p(X, \mathcal{F}) = H^p(Y, f_*\mathcal{F})$ pour tous $p$.
\item Si $H^p(Y, R^qf_*\mathcal{F}) = 0$ pour tous $q$ et $p > 0$, alors $H^q(X, \mathcal{F}) = H^0(Y, R^qf_*\mathcal{F})$ pour tous $q$.
\end{enumerate}

\end{lemma}

\begin{proof}

Ce sont deux conditions simples qui forcent la suite spectrale de Leray à dégénérer en $E_2$. On peut aussi démontrer ces faits directement, sans utiliser la suite spectrale ; c'est un bon exercice de cohomologie des faisceaux.

\end{proof}

\begin{lemma}

\label{cohomology-lemma-higher-direct-images-compose}

\begin{slogan}
Le foncteur dérivé total d'une composition est la composition des foncteurs dérivés totaux.
\end{slogan}
Soient $f : X \to Y$ et $g : Y \to Z$ des morphismes d'espaces annelés. Alors $Rg_* \circ Rf_* = R(g \circ f)_*$ comme foncteurs de $D^{+}(X) \to D^{+}(Z)$.
\end{lemma}

\begin{proof}
Nous allons appliquer Catégories dérivées, lemme \ref{derived-lemma-compose-derived-functors}. Il est clair que $g_* \circ f_* = (g \circ f)_*$ ; voir Faisceaux, lemme \ref{sheaves-lemma-pushforward-composition}. Il reste à montrer que $f_*\mathcal{I}$ est $g_*$-acyclique. Ceci résulte du lemme \ref{cohomology-lemma-pushforward-injective} et de la description des images directes supérieures $R^ig_*$ dans le lemme \ref{cohomology-lemma-describe-higher-direct-images}.
\end{proof}

\begin{lemma}
[Suite spectrale de Leray relative]
\label{cohomology-lemma-relative-Leray}
Soient $f : X \to Y$ et $g : Y \to Z$ des morphismes d'espaces annelés. Soit $\mathcal{F}$ un $\mathcal{O}_X$-module. Il existe une suite spectrale telle que $$
E_2^{p, q} = R^pg_*(R^qf_*\mathcal{F})
$$ convergeant vers $R^{p + q}(g \circ f)_*\mathcal{F}$. Cette suite spectrale est fonctorielle en $\mathcal{F}$, et il existe une version pour les complexes de $\mathcal{O}_X$-modules.
\end{lemma}

\begin{proof}
Il s'agit de la suite spectrale de Grothendieck associée à la composition des foncteurs ; elle résulte du lemme \ref{cohomology-lemma-higher-direct-images-compose} et des Catégories dérivées, lemme \ref{derived-lemma-grothendieck-spectral-sequence}.
\end{proof}













\section{Fonctorialité de la cohomologie}

\label{cohomology-section-functoriality}

\begin{lemma}

\label{cohomology-lemma-functoriality}

Soit $f : X \to Y$ un morphisme d'espaces annelés. Soient $\mathcal{G}^\bullet$, respectivement $\mathcal{F}^\bullet$, des complexes bornés inférieurement de $\mathcal{O}_Y$-modules, respectivement de $\mathcal{O}_X$-modules.
Soit $\varphi : \mathcal{G}^\bullet \to f_*\mathcal{F}^\bullet$ un
morphisme de complexes. Il existe un morphisme canonique
$$
\mathcal{G}^\bullet
\longrightarrow
Rf_*(\mathcal{F}^\bullet)
$$
dans $D^{+}(Y)$. De plus, cette construction est fonctorielle par rapport au triple
$(\mathcal{G}^\bullet, \mathcal{F}^\bullet, \varphi)$.

\end{lemma}

\begin{proof}
Choisissons une résolution injective $\mathcal{F}^\bullet \to \mathcal{I}^\bullet$. Par définition, $Rf_*(\mathcal{F}^\bullet)$ est représenté par $f_*\mathcal{I}^\bullet$ dans $K^{+}(\mathcal{O}_Y)$. La composée $$
\mathcal{G}^\bullet \to f_*\mathcal{F}^\bullet \to f_*\mathcal{I}^\bullet
$$ est un morphisme dans $K^{+}(Y)$ qui se transforme en morphisme du lemme en appliquant le foncteur de localisation $j_Y : K^{+}(Y) \to D^{+}(Y)$.
\end{proof}

\noindent

Soit $f : X \to Y$ un morphisme d'espaces annelés. Soit $\mathcal{G}$ un $\mathcal{O}_Y$-module et soit
$\mathcal{F}$ un $\mathcal{O}_X$-module. Rappelons qu'un
$f$-morphisme $\varphi$ de $\mathcal{G}$ vers $\mathcal{F}$ est un morphisme
$\varphi : \mathcal{G} \to f_*\mathcal{F}$, ou, ce qui revient au même, un morphisme $\varphi : f^*\mathcal{G} \to \mathcal{F}$. Voir Faisceaux, définition \ref{sheaves-definition-f-map}. Un tel $f$-morphisme donne lieu à un morphisme de complexes

\begin{equation}
\label{cohomology-equation-functorial-derived}
\varphi :
R\Gamma(Y, \mathcal{G})
\longrightarrow
R\Gamma(X, \mathcal{F})
\end{equation}

dans $D^{+}(\mathcal{O}_Y(Y))$. Plus précisément, nous utilisons le morphisme
$\mathcal{G} \to Rf_*\mathcal{F}$ dans $D^{+}(Y)$ du lemme \ref{cohomology-lemma-functoriality}, puis nous appliquons $R\Gamma(Y, -)$. D'après le lemme \ref{cohomology-lemma-before-Leray}, nous avons
$R\Gamma(X, \mathcal{F}) = R\Gamma(Y, Rf_*\mathcal{F})$
et nous obtenons la flèche affichée. Nous la décrivons complètement dans la remarque \ref{cohomology-remark-explain-arrow} ci-dessous. Elle induit en particulier des morphismes en cohomologie

\begin{equation}
\label{cohomology-equation-functorial}
\varphi : H^i(Y, \mathcal{G}) \longrightarrow H^i(X, \mathcal{F}).
\end{equation}

\begin{remark}

\label{cohomology-remark-explain-arrow}

Soit $f : X \to Y$ un morphisme d'espaces annelés. Soit $\mathcal{G}$ un $\mathcal{O}_Y$-module. Soit $\mathcal{F}$ un $\mathcal{O}_X$-module. Soit $\varphi$ un $f$-morphisme de $\mathcal{G}$ vers $\mathcal{F}$. Choisissons une résolution $\mathcal{F} \to \mathcal{I}^\bullet$
par un complexe de $\mathcal{O}_X$-modules injectifs. Choisissons des résolutions $\mathcal{G} \to \mathcal{J}^\bullet$ et
$f_*\mathcal{I}^\bullet \to (\mathcal{J}')^\bullet$ par des complexes de $\mathcal{O}_Y$-modules injectifs. D'après Catégories dérivées, lemme \ref{derived-lemma-morphisms-lift},
il existe un morphisme de complexes
$\beta$ tel que le diagramme

\begin{equation}
\label{cohomology-equation-choice}
\xymatrix{
\mathcal{G} \ar[d] \ar[r] &
f_*\mathcal{F} \ar[r] &
f_*\mathcal{I}^\bullet \ar[d] \\
\mathcal{J}^\bullet \ar[rr]^\beta & &
(\mathcal{J}')^\bullet
}
\end{equation}

En appliquant les foncteurs de sections globales, nous obtenons un diagramme
$$
\xymatrix{
 & & \Gamma(Y, f_*\mathcal{I}^\bullet) \ar[d]_{qis} \ar@{=}[r] &
\Gamma(X, \mathcal{I}^\bullet) \\
\Gamma(Y, \mathcal{J}^\bullet) \ar[rr]^\beta & &
\Gamma(Y, (\mathcal{J}')^\bullet) &
}
$$
Le complexe en bas à gauche représente $R\Gamma(Y, \mathcal{G})$
et le complexe en haut à droite représente $R\Gamma(X, \mathcal{F})$. La flèche verticale est un quasi-isomorphisme d'après le lemme \ref{cohomology-lemma-before-Leray} ; elle devient inversible après application du foncteur de localisation
$K^{+}(\mathcal{O}_Y(Y)) \to D^{+}(\mathcal{O}_Y(Y))$. La flèche (\ref{cohomology-equation-functorial-derived}) est la composée de la flèche horizontale avec l'inverse de la flèche verticale.

\end{remark}





\section{Raffinements et cohomologie de {\v C}ech}

\label{cohomology-section-refinements-cech}

\noindent

Soit $(X, \mathcal{O}_X)$ un espace annelé.
$\mathcal{U} : X = \bigcup_{i \in I} U_i$ et
$\mathcal{V} : X = \bigcup_{j \in J} V_j$ soient des recouvrements ouverts. Supposons que $\mathcal{U}$ soit un raffinement de $\mathcal{V}$. Choisissons une application $c : I \to J$ telle que $U_i \subset V_{c(i)}$
pour tous $i \in I$. Cela induit un morphisme de complexes de {\v C}ech
$$
\gamma :
\check{\mathcal{C}}^\bullet(\mathcal{V}, \mathcal{F})
\longrightarrow
\check{\mathcal{C}}^\bullet(\mathcal{U}, \mathcal{F}),
\quad
(\xi_{j_0 \ldots j_p})
\longmapsto
(\xi_{c(i_0) \ldots c(i_p)}|_{U_{i_0 \ldots i_p}})
$$
fonctoriel par rapport au faisceau de $\mathcal{O}_X$-modules $\mathcal{F}$. Supposons que $c' : I \to J$ soit une seconde application telle que
$U_i \subset V_{c'(i)}$ pour tous $i \in I$. Alors les morphismes correspondants
$\gamma$ et $\gamma'$ sont homotopes.
$\gamma - \gamma' = \text{d} \circ h + h \circ \text{d}$
avec
$h : \check{\mathcal{C}}^{p + 1}(\mathcal{V}, \mathcal{F}) \to
\check{\mathcal{C}}^p(\mathcal{U}, \mathcal{F})$
donné par la formule
$$
h(\alpha)_{i_0 \ldots i_p} =
\sum\nolimits_{a = 0}^{p}
(-1)^a
\alpha_{c(i_0)\ldots c(i_a) c'(i_a) \ldots c'(i_p)}
$$
Nous omettons le calcul qui le vérifie ; voir la discussion qui suit (\ref{cohomology-equation-transformation}) pour une démonstration dans un cadre plus général. En particulier, le morphisme induit sur les groupes de cohomologie de {\v C}ech est indépendant du choix de $c$. En outre, si
$\mathcal{W} : X = \bigcup_{k \in K} W_k$ est un troisième recouvrement ouvert et si $\mathcal{V}$ est un raffinement de $\mathcal{W}$, alors la composée des morphismes
$$
\check{\mathcal{C}}^\bullet(\mathcal{W}, \mathcal{F})
\longrightarrow
\check{\mathcal{C}}^\bullet(\mathcal{V}, \mathcal{F})
\longrightarrow
\check{\mathcal{C}}^\bullet(\mathcal{U}, \mathcal{F})
$$
associés aux applications $I \to J$ et $J \to K$ est le morphisme associé à la composée $I \to K$. Nous pouvons donc définir les groupes de cohomologie de {\v C}ech
$$
\check{H}^p(X, \mathcal{F}) =
\colim_\mathcal{U} \check{H}^p(\mathcal{U}, \mathcal{F})
$$
où la colimite porte sur tous les recouvrements ouverts de $X$, ordonnés par raffinement.

\medskip\noindent

Les morphismes $\gamma$ définis ci-dessus sont compatibles avec le morphisme vers la cohomologie ; autrement dit, la composée
$$
\check{H}^p(\mathcal{V}, \mathcal{F}) \to
\check{H}^p(\mathcal{U}, \mathcal{F})
\xrightarrow{\text{Lemme \ref{cohomology-lemma-cech-cohomology}}}
H^p(X, \mathcal{F})
$$
est le morphisme canonique du premier groupe vers la cohomologie du lemme \ref{cohomology-lemma-cech-cohomology}. Nous le démontrerons ci-dessous dans un cadre légèrement plus général. Il en résulte un morphisme bien défini

\begin{equation}
\label{cohomology-equation-cech-to-cohomology}
\check{H}^p(X, \mathcal{F}) =
\colim_\mathcal{U} \check{H}^p(\mathcal{U}, \mathcal{F})
\longrightarrow
H^p(X, \mathcal{F})
\end{equation}

de la cohomologie de {\v C}ech vers la cohomologie.

\begin{lemma}

\label{cohomology-lemma-functoriality-cech}

Soit $f : X \to Y$ un morphisme d'espaces annelés. Soit $\varphi : f^*\mathcal{G} \to \mathcal{F}$ un $f$-morphisme d'un $\mathcal{O}_Y$-module $\mathcal{G}$ vers un
$\mathcal{O}_X$-module $\mathcal{F}$. Soient $\mathcal{U} : X = \bigcup_{i \in I} U_i$ et
$\mathcal{V} : Y = \bigcup_{j \in J} V_j$ des recouvrements ouverts. Supposons que $\mathcal{U}$ soit un raffinement de
$f^{-1}\mathcal{V} : X = \bigcup_{j \in J} f^{-1}(V_j)$. Dans ce cas, il existe un diagramme commutatif
$$
\xymatrix{
\check{\mathcal{C}}^\bullet(\mathcal{U}, \mathcal{F}) \ar[r] &
R\Gamma(X, \mathcal{F}) \\
\check{\mathcal{C}}^\bullet(\mathcal{V}, \mathcal{G}) \ar[r]
\ar[u]^\gamma &
R\Gamma(Y, \mathcal{G}) \ar[u]
}
$$
dans $D^{+}(\mathcal{O}_X(X))$, où les flèches horizontales sont celles du lemme \ref{cohomology-lemma-cech-cohomology} et la flèche verticale de droite est donnée par (\ref{cohomology-equation-functorial-derived}). En particulier, nous obtenons des diagrammes commutatifs de groupes de cohomologie
$$
\xymatrix{
\check{H}^p(\mathcal{U}, \mathcal{F}) \ar[r] &
H^p(X, \mathcal{F}) \\
\check{H}^p(\mathcal{V}, \mathcal{G}) \ar[r]
\ar[u]^\gamma &
H^p(Y, \mathcal{G}) \ar[u]
}
$$
où la flèche verticale de droite est celle de (\ref{cohomology-equation-functorial}).

\end{lemma}

\begin{proof}

Définissons d'abord la flèche verticale de gauche. Choisissons une application
$c : I \to J$ telle que $U_i \subset f^{-1}(V_{c(i)})$ pour tous
$i \in I$. En degré $p$ nous définissons l'application par la règle
$$
\gamma(s)_{i_0 \ldots i_p} = \varphi(s)_{c(i_0) \ldots c(i_p)}
$$
Ceci est bien défini, car $\varphi$ induit des morphismes
$\mathcal{G}(V_{c(i_0) \ldots c(i_p)}) \to \mathcal{F}(U_{i_0 \ldots i_p})$
par hypothèse. Cela définit clairement un morphisme de complexes. Choisissons des résolutions injectives
$\mathcal{F} \to \mathcal{I}^\bullet$ sur $X$ et
$\mathcal{G} \to J^\bullet$ sur $Y$. Comme dans la démonstration du lemme \ref{cohomology-lemma-cech-cohomology}, introduisons les bicomplexes $A^{\bullet, \bullet}$ et $B^{\bullet, \bullet}$
avec des termes
$$
B^{p, q} = \check{\mathcal{C}}^p(\mathcal{V}, \mathcal{J}^q)
\quad
\text{et}
\quad
A^{p, q} = \check{\mathcal{C}}^p(\mathcal{U}, \mathcal{I}^q).
$$
Comme dans la remarque \ref{cohomology-remark-explain-arrow} ci-dessus, choisissons également une résolution injective
$f_*\mathcal{I} \to (\mathcal{J}')^\bullet$ sur $Y$ et un morphisme de complexes $\beta : \mathcal{J} \to (\mathcal{J}')^\bullet$
rendant (\ref{cohomology-equation-choice}) commutatif. Introduisons deux autres bicomplexes, $(B')^{\bullet, \bullet}$ et
$(B'')^{\bullet, \bullet}$ avec
$$
(B')^{p, q} = \check{\mathcal{C}}^p(\mathcal{V}, (\mathcal{J}')^q)
\quad
\text{et}
\quad
(B'')^{p, q} = \check{\mathcal{C}}^p(\mathcal{V}, f_*\mathcal{I}^q).
$$
Il existe un $f$-morphisme de complexes de
$f_*\mathcal{I}^\bullet$ vers $\mathcal{I}^\bullet$. La même formule que ci-dessus définit donc un morphisme de bicomplexes
$$
\gamma : (B'')^{\bullet, \bullet} \longrightarrow A^{\bullet, \bullet}.
$$
Considérons le diagramme de complexes
$$
\xymatrix{
\check{\mathcal{C}}^\bullet(\mathcal{U}, \mathcal{F})
\ar[r] &
\text{Tot}(A^{\bullet, \bullet}) & & &
\Gamma(X, \mathcal{I}^\bullet) \ar[lll]^{qis}
\ar@{=}[ddl]\\
\check{\mathcal{C}}^\bullet(\mathcal{V}, \mathcal{G})
\ar[r] \ar[u]^\gamma &
\text{Tot}(B^{\bullet, \bullet}) \ar[r]^\beta &
\text{Tot}((B')^{\bullet, \bullet}) &
\text{Tot}((B'')^{\bullet, \bullet}) \ar[l] \ar[llu]_{s\gamma} \\
& \Gamma(Y, \mathcal{J}^\bullet) \ar[u]^{qis} \ar[r]^\beta &
\Gamma(Y, (\mathcal{J}')^\bullet) \ar[u] &
\Gamma(Y, f_*\mathcal{I}^\bullet) \ar[u] \ar[l]_{qis}
}
$$
Les deux flèches horizontales de cibles $\text{Tot}(A^{\bullet, \bullet})$ et
$\text{Tot}(B^{\bullet, \bullet})$
sont celles expliquées dans le lemme \ref{cohomology-lemma-cech-cohomology}. La partie supérieure gauche, un pentagone, est commutative parce que (\ref{cohomology-equation-choice}) l'est. Les deux carrés inférieurs sont manifestement commutatifs. Les définitions montrent aussi immédiatement que la partie supérieure droite, un carré, est commutative. Le lemme résulte alors des définitions et du fait que les deux chemins extérieurs allant de $\check{\mathcal{C}}^\bullet(\mathcal{V}, \mathcal{G})$
à $\Gamma(X, \mathcal{I}^\bullet)$, par le haut ou par le bas, coïncident après inversion des quasi-isomorphismes rencontrés.

\end{proof}





\section{Cohomologie sur les espaces séparés quasi compacts}

\label{cohomology-section-cohomology-LC}

\noindent
Sur un tel espace, la cohomologie de {\v C}ech coïncide avec la cohomologie usuelle.

\begin{lemma}

\label{cohomology-lemma-cech-always}

Soit $X$ un espace topologique et soit $\mathcal{F}$ un faisceau abélien sur cet espace. Alors l'application $\check{H}^1(X, \mathcal{F}) \to H^1(X, \mathcal{F})$ définie dans (\ref{cohomology-equation-cech-to-cohomology}) est un isomorphisme.

\end{lemma}

\begin{proof}
Soit $\mathcal{U}$ un recouvrement ouvert de $X$. D'après le lemme \ref{cohomology-lemma-cech-spectral-sequence}, il existe une suite exacte $$
0 \to \check{H}^1(\mathcal{U}, \mathcal{F}) \to H^1(X, \mathcal{F})
\to \check{H}^0(\mathcal{U}, \underline{H}^1(\mathcal{F}))
$$ L'application est donc injective. Pour établir la surjectivité, il suffit de montrer que tout élément de $\check{H}^0(\mathcal{U}, \underline{H}^1(\mathcal{F}))$ s'envoie sur zéro après remplacement de $\mathcal{U}$ par un raffinement. Cela résulte immédiatement des définitions et du fait que $\underline{H}^1(\mathcal{F})$ est un préfaisceau de groupes abéliens dont la faisceautisation est nulle par le caractère local de la cohomologie ; voir le lemme \ref{cohomology-lemma-kill-cohomology-class-on-covering}.
\end{proof}

\begin{lemma}

\label{cohomology-lemma-cech-Hausdorff-quasi-compact}
Soit $X$ un espace topologique séparé et quasi-compact. Soit $\mathcal{F}$ un faisceau abélien sur $X$. Alors l'application $\check{H}^n(X, \mathcal{F}) \to H^n(X, \mathcal{F})$ définie dans (\ref{cohomology-equation-cech-to-cohomology}) est un isomorphisme pour tout $n$.
\end{lemma}

\begin{proof}
Nous savons déjà que $\check{H}^n(X, -) \to H^n(X, -)$ est un isomorphisme de foncteurs pour $n = 0, 1$ ; voir le lemme \ref{cohomology-lemma-cech-always}. Les foncteurs $H^n(X, -)$ forment un $\delta$-foncteur universel ; voir Catégories dérivées, lemme \ref{derived-lemma-higher-derived-functors}. Si nous montrons que $\check{H}^n(X, -)$ forme un $\delta$-foncteur universel et que $\check{H}^n(X, -) \to H^n(X, -)$ est compatible avec les morphismes de bord, alors l'application sera automatiquement un isomorphisme par unicité des $\delta$-foncteurs universels ; voir Homologie, lemme \ref{homology-lemma-uniqueness-universal-delta-functor}.

\medskip\noindent

Soit $0 \to \mathcal{F} \to \mathcal{G} \to \mathcal{H} \to 0$
une suite exacte courte de faisceaux abéliens sur $X$. Soit $\mathcal{U} : X = \bigcup_{i \in I} U_i$ un recouvrement ouvert. On obtient alors un complexe de complexes
$$
0 \to \check{\mathcal{C}}^\bullet(\mathcal{U}, \mathcal{F}) \to
\check{\mathcal{C}}^\bullet(\mathcal{U}, \mathcal{G}) \to
\check{\mathcal{C}}^\bullet(\mathcal{U}, \mathcal{H}) \to 0
$$
qui n'est en général pas exact à droite. Cette suite définit les morphismes
$$
\check{H}^n(\mathcal{U}, \mathcal{F}) \to
\check{H}^n(\mathcal{U}, \mathcal{G}) \to
\check{H}^n(\mathcal{U}, \mathcal{H})
$$
mais ne suffit pas à définir un morphisme de bord
$\delta : \check{H}^n(\mathcal{U}, \mathcal{H}) \to
\check{H}^{n + 1}(\mathcal{U}, \mathcal{F})$. En effet une telle chose n'existera pas en général. Cependant, compte tenu d'un élément $\overline{h} \in \check{H}^n(\mathcal{U}, \mathcal{H})$
représenté par le cocycle
$h = (h_{i_0 \ldots i_n})$
nous pouvons choisir des recouvrements ouverts
$$
U_{i_0 \ldots i_n} = \bigcup W_{i_0 \ldots i_n, k}
$$
de sorte que $h_{i_0 \ldots i_n}|_{W_{i_0 \ldots i_n, k}}$
se relève en une section de $\mathcal{G}$ sur $W_{i_0 \ldots i_n, k}$. D'après Topologie, lemme \ref{topology-lemma-refine-covering},
(c'est ici que nous utilisons l'hypothèse que $X$ est séparé et quasi-compact), nous pouvons choisir un recouvrement ouvert $\mathcal{V} : X = \bigcup_{j \in J} V_j$
et $\alpha : J \to I$ tels que $V_j \subset U_{\alpha(j)}$
(c'est un raffinement) et tel que pour tous $j_0, \ldots, j_n \in J$
il existe un $k$ tel que
$V_{j_0 \ldots j_n} \subset W_{\alpha(j_0) \ldots \alpha(j_n), k}$. Nous obtenons des morphismes de complexes
$$
\xymatrix{
0 \ar[r] &
\check{\mathcal{C}}^\bullet(\mathcal{U}, \mathcal{F}) \ar[d] \ar[r] &
\check{\mathcal{C}}^\bullet(\mathcal{U}, \mathcal{G}) \ar[d] \ar[r] &
\check{\mathcal{C}}^\bullet(\mathcal{U}, \mathcal{H}) \ar[d] \ar[r] &
0 \\
0 \ar[r] &
\check{\mathcal{C}}^\bullet(\mathcal{V}, \mathcal{F}) \ar[r] &
\check{\mathcal{C}}^\bullet(\mathcal{V}, \mathcal{G}) \ar[r] &
\check{\mathcal{C}}^\bullet(\mathcal{V}, \mathcal{H}) \ar[r] &
0
}
$$
En fait, les flèches verticales sont les morphismes de complexes utilisés pour définir les morphismes de transition entre les groupes de cohomologie de {\v C}ech. Notre choix de raffinement montre que nous pouvons choisir
$$
g_{j_0 \ldots j_n} \in
\mathcal{G}(V_{j_0 \ldots j_n}),\quad
g_{j_0 \ldots j_n} \longmapsto
h_{\alpha(j_0) \ldots \alpha(j_n)}|_{V_{j_0 \ldots j_n}}
$$
La cochaîne $g = (g_{j_0 \ldots j_n})$ n'est pas nécessairement un cocycle, mais son cobord de {\v C}ech $\text{d}(g)$
s'envoie sur zéro dans $\check{\mathcal{C}}^{n + 1}(\mathcal{V}, \mathcal{H})$
(par le diagramme commutatif ci-dessus et parce que $h$ est un cocycle). Ainsi, $\text{d}(g)$ est un cocycle dans
$\check{\mathcal{C}}^\bullet(\mathcal{V}, \mathcal{F})$. Cela nous permet de définir
$$
\delta(\overline{h}) = \text{classe de }\text{d}(g)\text{ dans }
\check{H}^{n + 1}(\mathcal{V}, \mathcal{F})
$$
Étant donné un élément $\xi \in \check{H}^n(X, \mathcal{G})$,
nous choisissons un recouvrement ouvert $\mathcal{U}$ et un élément
$\overline{h} \in \check{H}^n(\mathcal{U}, \mathcal{G})$
qui s'envoie sur $\xi$ dans la colimite définissant la cohomologie de {\v C}ech. Choisissons ensuite $\mathcal{V}$ et $g$ comme ci-dessus, et posons
$\delta(\xi)$ égal à l'image de $\delta(\overline{h})$
dans $\check{H}^n(X, \mathcal{F})$. Il reste à vérifier plusieurs propriétés, toutes immédiates. Il faut notamment montrer que la construction est indépendante du choix de
$\mathcal{U}, \overline{h}, \mathcal{V}, \alpha : J \to I, g$. La classe de $\text{d}(g)$ est indépendante du choix des relèvements
$g_{i_0 \ldots i_n}$, car la différence est un cobord. L'application ne dépend pas de $\alpha$\footnote{Cette vérification est importante : la non-unicité de $\alpha$ est le seul obstacle à prendre la colimite des complexes de {\v C}ech sur tous les recouvrements de $X$ afin d'obtenir une suite exacte courte de complexes calculant la cohomologie de {\v C}ech.}, car deux choix de $\alpha$ donnent des morphismes verticaux homotopes entre les complexes du diagramme ci-dessus ; voir la section \ref{cohomology-section-refinements-cech}. Pour les autres choix intervenant dans la construction, toute famille finie de recouvrements ouverts de $X$ admet un raffinement ouvert commun. L'additivité se vérifie de la même manière. Enfin, il faut vérifier que les applications
$\check{H}^n(X, -) \to H^n(X, -)$ sont compatibles avec les morphismes de bord. Pour ce faire, nous choisissons des résolutions injectives
$$
\xymatrix{
0 \ar[r] &
\mathcal{F} \ar[r] \ar[d] &
\mathcal{G} \ar[r] \ar[d] &
\mathcal{H} \ar[r] \ar[d] &
0 \\
0 \ar[r] &
\mathcal{I}_1^\bullet \ar[r] &
\mathcal{I}_2^\bullet \ar[r] &
\mathcal{I}_3^\bullet \ar[r] &
0
}
$$
comme dans Catégories dérivées, lemme \ref{derived-lemma-injective-resolution-ses}. Ceci donne un diagramme commutatif
$$
\xymatrix{
0 \ar[r] &
\check{\mathcal{C}}^\bullet(\mathcal{U}, \mathcal{F}) \ar[r] \ar[d] &
\check{\mathcal{C}}^\bullet(\mathcal{U}, \mathcal{F}) \ar[r] \ar[d] &
\check{\mathcal{C}}^\bullet(\mathcal{U}, \mathcal{F}) \ar[r] \ar[d] &
0 \\
0 \ar[r] &
\text{Tot}(\check{\mathcal{C}}^\bullet(\mathcal{U}, \mathcal{I}_1^\bullet))
\ar[r] &
\text{Tot}(\check{\mathcal{C}}^\bullet(\mathcal{U}, \mathcal{I}_2^\bullet))
\ar[r] &
\text{Tot}(\check{\mathcal{C}}^\bullet(\mathcal{U}, \mathcal{I}_3^\bullet))
\ar[r] &
0
}
$$
Ici $\mathcal{U}$ est un recouvrement ouvert comme ci-dessus et les flèches verticales sont celles utilisées pour définir les morphismes
$\check{H}^n(\mathcal{U}, -) \to H^n(X, -)$ ; voir le lemme \ref{cohomology-lemma-cech-cohomology}. Le complexe inférieur est exact, car la suite de complexes injectifs est scindée exacte terme à terme. Le morphisme de bord en cohomologie est donc calculé par la procédure usuelle appliquée à cette suite exacte inférieure ; voir Homologie, lemme \ref{homology-lemma-long-exact-sequence-cochain}. Il en va de même après passage au raffinement
$\mathcal{V}$, où le morphisme de bord en cohomologie de {\v C}ech a été défini. Les morphismes de bord coïncident donc, puisqu'ils reposent sur la même construction (lorsque le premier est défini sur un élément de la cohomologie de {\v C}ech associé à un recouvrement ouvert donné). Ceci achève la construction de la structure de $\delta$-foncteur sur la cohomologie de {\v C}ech et établit sa compatibilité avec la
structure de $\delta$-foncteur sur la cohomologie usuelle.

\medskip\noindent
Enfin, le lemme \ref{cohomology-lemma-injective-trivial-cech} montre que la cohomologie supérieure de {\v C}ech est nulle sur les faisceaux injectifs. La cohomologie de {\v C}ech forme donc un $\delta$-foncteur universel, d'après Homologie, lemme \ref{homology-lemma-efface-implies-universal}.
\end{proof}

\begin{lemma}

\label{cohomology-lemma-cohomology-of-closed}

\begin{reference}
\cite[Expose V bis, 4.1.3]{SGA4}
\end{reference}
Soit $X$ un espace topologique. Soit $Z \subset X$ un sous-ensemble quasi-compact tel que deux points quelconques de $Z$ possèdent des voisinages ouverts disjoints dans $X$. Pour tout faisceau abélien $\mathcal{F}$ sur $X$, le morphisme canonique $$
\colim H^p(U, \mathcal{F})
\longrightarrow
H^p(Z, \mathcal{F}|_Z)
$$ où la colimite porte sur les voisinages ouverts $U$ de $Z$ dans $X$, est un isomorphisme.
\end{lemma}

\begin{proof}

Commençons par le cas $p = 0$. L'injectivité découle de la définition de $\mathcal{F}|_Z$ et vaut en général (pour tout sous-ensemble de tout espace topologique $X$). Supposons ensuite que
$s \in H^0(Z, \mathcal{F}|_Z)$. Nous pouvons trouver des ouverts $U_i \subset X$
tels que $Z \subset \bigcup U_i$ et tels que $s|_{Z \cap U_i}$
vient de $s_i \in \mathcal{F}(U_i)$. Il s'ensuit qu'il existe des ouvertures $W_{ij} \subset U_i \cap U_j$ avec
$W_{ij} \cap Z = U_i \cap U_j \cap Z$ tels que
$s_i|_{W_{ij}} = s_j|_{W_{ij}}$. En appliquant Topologie, lemme
\ref{topology-lemma-lift-covering-of-quasi-compact-hausdorff-subset}
nous trouvons des ouverts $V_i$ de $X$ tels que $V_i \subset U_i$ et $V_i \cap V_j \subset W_{ij}$. Nous voyons donc que
$s_i|_{V_i}$ colle sur une section de $\mathcal{F}$ sur le voisinage ouvert $\bigcup V_i$ de $Z$.

\medskip\noindent
Pour achever la démonstration, il suffit de montrer que, si $\mathcal{I}$ est un faisceau abélien injectif sur $X$, alors $H^p(Z, \mathcal{I}|_Z) = 0$ pour $p > 0$. Cela résulte des suites exactes courtes et du décalage des degrés ; nous omettons les détails. Supposons donc que $\overline{\xi}$ soit un élément de $H^p(Z, \mathcal{I}|_Z)$ pour un certain $p > 0$. D'après le lemme \ref{cohomology-lemma-cech-Hausdorff-quasi-compact}, l'élément $\overline{\xi}$ provient de $\check{H}^p(\mathcal{V}, \mathcal{I}|_Z)$ pour un recouvrement ouvert $\mathcal{V} : Z = \bigcup V_i$ de $Z$. Écrivons $\overline{\xi}$ comme l'image de la classe d'un cocycle $\xi = (\xi_{i_0 \ldots i_p})$ de $\check{\mathcal{C}}^p(\mathcal{V}, \mathcal{I}|_Z)$.

\medskip\noindent
Soit $\mathcal{I}' \subset \mathcal{I}|_Z$ le sous-préfaisceau défini par la formule $$
\mathcal{I}'(V) =
\{s \in \mathcal{I}|_Z(V) \mid
\exists (U, t),\ U \subset X\text{ ouvert},
\ t \in \mathcal{I}(U),\ V = Z \cap U,\ s = t|_{Z \cap U} \}
$$ Alors $\mathcal{I}|_Z$ est la faisceautisation de $\mathcal{I}'$. Ainsi, pour chaque $(p + 1)$-uplet $i_0 \ldots i_p$, nous pouvons trouver un recouvrement ouvert $V_{i_0 \ldots i_p} = \bigcup W_{i_0 \ldots i_p, k}$ tel que $\xi_{i_0 \ldots i_p}|_{W_{i_0 \ldots i_p, k}}$ soit une section de $\mathcal{I}'$. En appliquant Topologie, lemme \ref{topology-lemma-refine-covering}, nous pouvons, après avoir raffiné $\mathcal{V}$, supposer que chaque $\xi_{i_0 \ldots i_p}$ est une section du préfaisceau $\mathcal{I}'$.

\medskip\noindent

Écrivons $V_i = Z \cap U_i$ pour certains ouverts $U_i \subset X$. Puisque $\mathcal{I}$ est flasque (lemme \ref{cohomology-lemma-injective-flasque}) et que $\xi_{i_0 \ldots i_p}$ est une section de $\mathcal{I}'$,
pour chaque $(p + 1)$-uplet $i_0 \ldots i_p$, nous pouvons choisir une section $s_{i_0 \ldots i_p} \in \mathcal{I}(U_{i_0 \ldots i_p})$
dont la restriction est $\xi_{i_0 \ldots i_p}$ sur
$V_{i_0 \ldots i_p} = Z \cap U_{i_0 \ldots i_p}$. (On peut éviter ici d'utiliser le fait que les injectifs sont flasques en appliquant une fois de plus Topologie, lemme
\ref{topology-lemma-lift-covering-of-quasi-compact-hausdorff-subset}.) Soit $s = (s_{i_0 \ldots i_p})$ la cochaîne correspondante pour le recouvrement ouvert $U = \bigcup U_i$. Puisque $\text{d}(\xi) = 0$, les sections
$\text{d}(s)_{i_0 \ldots i_{p + 1}}$ se restreignent à zéro sur $Z \cap U_{i_0 \ldots i_{p + 1}}$. Par conséquent, d'après les remarques initiales de la démonstration, il existe des ouverts
$W_{i_0 \ldots i_{p + 1}} \subset U_{i_0 \ldots i_{p + 1}}$
avec $Z \cap W_{i_0 \ldots i_{p + 1}} = Z \cap U_{i_0 \ldots i_{p + 1}}$
tels que $\text{d}(s)_{i_0 \ldots i_{p + 1}}|_{W_{i_0 \ldots i_{p + 1}}} = 0$. D'après Topologie, lemme
\ref{topology-lemma-lift-covering-of-quasi-compact-hausdorff-subset}
nous pouvons trouver $U'_i \subset U_i$ tels que $Z \subset \bigcup U'_i$
et $U'_{i_0 \ldots i_{p + 1}} \subset W_{i_0 \ldots i_{p + 1}}$. Alors $s' = (s'_{i_0 \ldots i_p})$ avec
$s'_{i_0 \ldots i_p} = s_{i_0 \ldots i_p}|_{U'_{i_0 \ldots i_p}}$
est un cocycle à valeurs dans $\mathcal{I}$ pour le recouvrement ouvert
$U' = \bigcup U'_i$ d'un voisinage ouvert de $Z$. Puisque les groupes de cohomologie supérieure de {\v C}ech de $\mathcal{I}$ sont nuls (lemme \ref{cohomology-lemma-injective-trivial-cech}), nous concluons que $s'$ est un cobord. Il s'ensuit que l'image de
$\xi$ dans le complexe de {\v C}ech associé au recouvrement ouvert
$Z = \bigcup Z \cap U'_i$ est un cobord, ce qui achève la démonstration.

\end{proof}





\section{L'application de changement de base}

\label{cohomology-section-base-change-map}

\noindent

Nous devons savoir construire le morphisme de changement de base dans certains cas. Comme nous n'avons pas encore étudié l'image inverse dérivée, nous ne considérons ici que le changement de base par un morphisme plat d'espaces annelés. Avant d'énoncer le résultat, examinons l'image inverse plate sur la catégorie dérivée. Soit $g : X \to Y$ un morphisme plat d'espaces annelés. D'après Faisceaux de modules, lemme \ref{modules-lemma-pullback-flat},
le foncteur $g^* : \textit{Mod}(\mathcal{O}_Y) \to
\textit{Mod}(\mathcal{O}_X)$ est exact. Il admet donc un foncteur dérivé
$$
g^* : D^{+}(Y) \to D^{+}(X)
$$
qui se calcule simplement en appliquant l'image inverse à un représentant de l'objet de $D^{+}(Y)$ considéré ; voir Catégories dérivées, lemme \ref{derived-lemma-right-derived-exact-functor}. Ainsi, comme annoncé, nous notons ce foncteur $g^*$ plutôt que
$Lg^*$.

\begin{lemma}

\label{cohomology-lemma-base-change-map-flat-case}
Considérons le diagramme commutatif suivant d'espaces annelés : $$
\xymatrix{
X' \ar[r]_{g'} \ar[d]_{f'} &
X \ar[d]^f \\
S' \ar[r]^g &
S
}
$$ Soit $\mathcal{F}^\bullet$ un complexe borné inférieurement de $\mathcal{O}_X$-modules. Supposons $g$ et $g'$ plats. Il existe alors un morphisme canonique de changement de base $$
g^*Rf_*\mathcal{F}^\bullet
\longrightarrow
R(f')_*(g')^*\mathcal{F}^\bullet
$$ dans $D^{+}(S')$.
\end{lemma}

\begin{proof}

Choisissons des résolutions injectives $\mathcal{F}^\bullet \to \mathcal{I}^\bullet$
et $(g')^*\mathcal{F}^\bullet \to \mathcal{J}^\bullet$. D'après le lemme \ref{cohomology-lemma-pushforward-injective-flat},
$(g')_*\mathcal{J}^\bullet$ est un complexe de modules injectifs représentant
$R(g')_*(g')^*\mathcal{F}^\bullet$. D'où par catégories dérivées, lemmes \ref{derived-lemma-morphisms-lift}
et \ref{derived-lemma-morphisms-equal-up-to-homotopy}
le morphisme $\beta$ du diagramme
$$
\xymatrix{
(g')_*(g')^*\mathcal{F}^\bullet \ar[r] &
(g')_*\mathcal{J}^\bullet \\
\mathcal{F}^\bullet \ar[u]^{adjunction} \ar[r] &
\mathcal{I}^\bullet \ar[u]_\beta
}
$$
existe et est unique à homotopie près. En poussant vers $S$, nous obtenons
$$
f_*\beta :
f_*\mathcal{I}^\bullet
\longrightarrow
f_*(g')_*\mathcal{J}^\bullet
=
g_*(f')_*\mathcal{J}^\bullet
$$
Par adjonction entre $g^*$ et $g_*$, nous obtenons un morphisme de complexes
$g^*f_*\mathcal{I}^\bullet \to (f')_*\mathcal{J}^\bullet$. Ce morphisme est unique à homotopie près, puisque le seul choix effectué au cours de la construction est celui du morphisme $\beta$
et tout a été fait au niveau des complexes.

\end{proof}

\begin{remark}

\label{cohomology-remark-correct-version-base-change-map}
La version « correcte » du morphisme de changement de base est le morphisme $$
Lg^* Rf_* \mathcal{F}^\bullet
\longrightarrow
R(f')_* L(g')^*\mathcal{F}^\bullet.
$$ Sa construction fait intervenir des complexes non bornés ; voir la remarque \ref{cohomology-remark-base-change}.
\end{remark}





\section{Changement de base propre en topologie}

\label{cohomology-section-proper-base-change}

\noindent
Dans cette section, nous démontrons une version très générale du théorème de changement de base propre en topologie. Elle affirme que les germes des images directes supérieures $R^pf_*$ peuvent se calculer sur la fibre.

\begin{lemma}

\label{cohomology-lemma-proper-base-change}
Soit $f : (X, \mathcal{O}_X) \to (Y, \mathcal{O}_Y)$ un morphisme d'espaces annelés. Soit $y \in Y$. Supposons que
\begin{enumerate}
\item $f$ est fermé,
\item $f$ est séparé, et
\item $f^{-1}(y)$ est quasi-compact.
\end{enumerate}
Alors, pour $E$ dans $D^+(\mathcal{O}_X)$, nous avons $(Rf_*E)_y = R\Gamma(f^{-1}(y), E|_{f^{-1}(y)})$ dans $D^+(\mathcal{O}_{Y, y})$.
\end{lemma}

\begin{proof}

Le morphisme de changement de base du lemme \ref{cohomology-lemma-base-change-map-flat-case}
donne un morphisme canonique $(Rf_*E)_y \to R\Gamma(f^{-1}(y), E|_{f^{-1}(y)})$. Pour montrer que ce morphisme est un isomorphisme, représentons $E$ par un complexe borné inférieurement de modules injectifs $\mathcal{I}^\bullet$. Posons $Z = f^{-1}(\{y\})$. Les hypothèses du lemme \ref{cohomology-lemma-cohomology-of-closed}
sont satisfaites ; voir Topologie, lemme \ref{topology-lemma-separated}. Les restrictions
$\mathcal{I}^n|_Z$ sont acycliques pour $\Gamma(Z, -)$. Ainsi $R\Gamma(Z, E|_Z)$ est représenté par le complexe $\Gamma(Z, \mathcal{I}^\bullet|_Z)$ ; voir Catégories dérivées, lemme \ref{derived-lemma-leray-acyclicity}. En d'autres termes, nous devons montrer que le morphisme
$$
\colim_V \mathcal{I}^\bullet(f^{-1}(V))
\longrightarrow
\Gamma(Z, \mathcal{I}^\bullet|_Z)
$$
est un isomorphisme. D'après le lemme \ref{cohomology-lemma-cohomology-of-closed},
nous voyons qu'il suffit de montrer que la collection de voisinages ouverts
$f^{-1}(V)$ de $Z = f^{-1}(\{y\})$
est cofinal dans le système de tous les voisinages ouverts. Si $f^{-1}(\{y\}) \subset U$ est un voisinage ouvert, alors, puisque $f$ est fermé, l'ensemble $V = Y \setminus f(X \setminus U)$ est un voisinage ouvert de $y$ tel que $f^{-1}(V) \subset U$. Cela démontre le lemme.

\end{proof}

\begin{theorem}
[changement de base propre]
\label{cohomology-theorem-proper-base-change}

\begin{reference}
\cite[Expose V bis, 4.1.1]{SGA4}
\end{reference}
Considérez un carré cartésien d'espaces topologiques $$
\xymatrix{
X' = Y' \times_Y X \ar[d]_{f'} \ar[r]_-{g'} & X \ar[d]^f \\
Y' \ar[r]^g & Y
}
$$ Supposons que $f$ soit propre. Soit $E$ un objet de $D^+(X)$. Alors le morphisme de changement de base $$
g^{-1}Rf_*E \longrightarrow Rf'_*(g')^{-1}E
$$ du lemme \ref{cohomology-lemma-base-change-map-flat-case} est un isomorphisme dans $D^+(Y')$.
\end{theorem}

\begin{proof}

Soit $y' \in Y'$ un point d'image $y \in Y$. Il suffit de montrer que le morphisme de changement de base induit un isomorphisme sur les germes en $y'$. Puisque $f$ est propre, il en va de même de $f'$ ; les fibres de $f$ et de $f'$ sont quasi-compactes, et $f$ et $f'$ sont fermés. Voir Topologie, théorème \ref{topology-theorem-characterize-proper} et lemme \ref{topology-lemma-base-change-separated}. Nous pouvons donc appliquer deux fois le lemme \ref{cohomology-lemma-proper-base-change} pour obtenir
$$
(Rf'_*(g')^{-1}E)_{y'} = R\Gamma((f')^{-1}(y'), (g')^{-1}E|_{(f')^{-1}(y')})
$$
et
$$
(Rf_*E)_y = R\Gamma(f^{-1}(y), E|_{f^{-1}(y)})
$$
Le morphisme induit sur les fibres $(f')^{-1}(y') \to f^{-1}(y)$ est un homéomorphisme d'espaces topologiques, et l'image inverse de
$E|_{f^{-1}(y)}$ est $(g')^{-1}E|_{(f')^{-1}(y')}$. Le résultat en découle.

\end{proof}

\begin{lemma}
[changement de base propre pour faisceaux d'ensembles]
\label{cohomology-lemma-proper-base-change-sheaves-of-sets}
Considérez un carré cartésien d'espaces topologiques $$
\xymatrix{
X' \ar[d]_{f'} \ar[r]_-{g'} & X \ar[d]^f \\
Y' \ar[r]^g & Y
}
$$ Supposons que $f$ soit propre. Alors $g^{-1}f_*\mathcal{F} = f'_*(g')^{-1}\mathcal{F}$ pour tout faisceau d'ensembles $\mathcal{F}$ sur $X$.
\end{lemma}

\begin{proof}
On raisonne exactement comme dans la démonstration du théorème \ref{cohomology-theorem-proper-base-change} : il suffit d'établir $(f_*\mathcal{F})_y = \Gamma(X_y, \mathcal{F}|_{X_y})$. Comme dans le lemme \ref{cohomology-lemma-proper-base-change}, on se ramène ensuite au cas $p = 0$ du lemme \ref{cohomology-lemma-cohomology-of-closed} pour les faisceaux d'ensembles. La première partie de la démonstration du lemme \ref{cohomology-lemma-cohomology-of-closed} s'applique aux faisceaux d'ensembles et achève la preuve. Nous omettons quelques détails.
\end{proof}



\section{Cohomologie et colimites}

\label{cohomology-section-limits}

\noindent

Soit $X$ un espace annelé. Soit $(\mathcal{F}_i, \varphi_{ii'})$ un système de faisceaux de $\mathcal{O}_X$-modules indexé par l'ensemble dirigé $I$ ; voir Catégories, section \ref{categories-section-posets-limits}. Pour chaque $i$, le morphisme canonique
$\mathcal{F}_i \to \colim_i \mathcal{F}_i$ induit un morphisme canonique
$$
\colim_i H^p(X, \mathcal{F}_i)
\longrightarrow
H^p(X, \colim_i \mathcal{F}_i)
$$
pour chaque $p \geq 0$. Il existe bien sûr un morphisme analogue pour chaque $U \subset X$. En général, ces morphismes ne sont pas des isomorphismes, même pour $p = 0$. Dans cette section, nous généralisons le résultat de Faisceaux, \ref{sheaves-lemma-directed-colimits-sections}. Voir aussi Faisceaux de modules, lemme \ref{modules-lemma-finite-presentation-quasi-compact-colimit}
(dans le cas spécial $\mathcal{G} = \mathcal{O}_X$).

\begin{lemma}

\label{cohomology-lemma-quasi-separated-cohomology-colimit}
Soit $X$ un espace annelé. Supposons que l'espace topologique sous-jacent à $X$ possède les propriétés suivantes :
\begin{enumerate}
\item il existe une base de sous-ensembles ouverts quasi compacts, et
\item l'intersection de deux ouverts quasi compacts est quasi-compacte.
\end{enumerate}
Alors, pour tout système dirigé $(\mathcal{F}_i, \varphi_{ii'})$ de faisceaux de $\mathcal{O}_X$-modules et pour tout ouvert quasi compact $U \subset X$, le morphisme canonique $$
\colim_i H^q(U, \mathcal{F}_i)
\longrightarrow
H^q(U, \colim_i \mathcal{F}_i)
$$ est un isomorphisme pour chaque $q \geq 0$.
\end{lemma}

\begin{proof}
Il importe de raisonner simultanément pour tous les ouverts quasi compacts $U \subset X$. Pour $q = 0$ et tout ouvert quasi compact $U \subset X$, le résultat découle de Faisceaux, lemme \ref{sheaves-lemma-directed-colimits-sections}, combiné avec Topologie, lemme \ref{topology-lemma-topology-quasi-separated-scheme}. Supposons le résultat établi pour tout $q \leq q_0$ et démontrons-le pour $q = q_0 + 1$.

\medskip\noindent

Par convention, l'ensemble d'indices $I$ est dirigé, et tout système de $\mathcal{O}_X$-modules $(\mathcal{F}_i, \varphi_{ii'})$
indexé par $I$ est un système dirigé. D'après Injectifs, lemme \ref{injectives-lemma-sheaves-modules-space}, la catégorie des $\mathcal{O}_X$-modules possède des plongements injectifs fonctoriels. Ainsi, pour tout système $(\mathcal{F}_i, \varphi_{ii'})$, il existe un système $(\mathcal{I}_i, \varphi_{ii'})$ tel que chaque $\mathcal{I}_i$ soit un $\mathcal{O}_X$-module injectif et un morphisme de systèmes formé de morphismes injectifs de $\mathcal{O}_X$-modules
$\mathcal{F}_i \to \mathcal{I}_i$. Notons $\mathcal{Q}_i$ le conoyau ; nous obtenons ainsi des suites exactes courtes
$$
0 \to
\mathcal{F}_i \to
\mathcal{I}_i \to
\mathcal{Q}_i \to 0.
$$
Nous affirmons que la séquence
$$
0 \to
\colim_i \mathcal{F}_i \to
\colim_i \mathcal{I}_i \to
\colim_i \mathcal{Q}_i \to 0.
$$
est également une suite exacte courte de $\mathcal{O}_X$-modules ; cela se vérifie sur les germes. Voir \ref{sheaves-section-limits-presheaves}
et \ref{sheaves-section-limits-sheaves}
puisque la colimite dirigée d'une suite exacte courte de groupes abéliens est exacte (voir Algèbre, lemme \ref{algebra-lemma-directed-colimit-exact}), le résultat en découle. Nous affirmons que
$H^q(U, \colim_i \mathcal{I}_i) = 0$ pour tout ouvert quasi compact $U \subset X$ et tout $q \geq 1$. Admettons provisoirement cette assertion et considérons le diagramme
$$
\xymatrix{
\colim_i H^{q_0}(U, \mathcal{I}_i) \ar[d] \ar[r] &
\colim_i H^{q_0}(U, \mathcal{Q}_i) \ar[d] \ar[r] &
\colim_i H^{q_0 + 1}(U, \mathcal{F}_i) \ar[d] \ar[r] &
0 \ar[d] \\
H^{q_0}(U, \colim_i \mathcal{I}_i) \ar[r] &
H^{q_0}(U, \colim_i \mathcal{Q}_i) \ar[r] &
H^{q_0 + 1}(U, \colim_i \mathcal{F}_i) \ar[r] &
0
}
$$
Le zéro dans le coin inférieur droit provient de l'assertion, et celui du coin supérieur droit du fait que les faisceaux
$\mathcal{I}_i$ sont injectifs. La ligne supérieure est exacte par Algèbre, lemme \ref{algebra-lemma-directed-colimit-exact}. Le lemme du serpent donne alors le résultat pour $q = q_0 + 1$.

\medskip\noindent

Il reste à démontrer l'assertion. Nous appliquons le lemme \ref{cohomology-lemma-cech-vanish-basis}. Soit $\mathcal{B}$ l'ensemble de tous les ouverts quasi compacts de $X$ ; par hypothèse, c'est une base de la topologie de $X$. Soit $\text{Cov}$ l'ensemble des recouvrements ouverts finis
$\mathcal{U} : U = \bigcup_{j = 1, \ldots, m} U_j$ tels que $U$ et chaque $U_j$ soient des ouverts quasi compacts de $X$. D'après le cas $q = 0$,
pour $\mathcal{U} \in \text{Cov}$, nous avons
$$
\check{\mathcal{C}}^\bullet(\mathcal{U}, \colim_i \mathcal{I}_i)
=
\colim_i \check{\mathcal{C}}^\bullet(\mathcal{U}, \mathcal{I}_i)
$$
parce que toutes les intersections multiples $U_{j_0 \ldots j_p}$
sont quasi-compactes. D'après le lemme \ref{cohomology-lemma-injective-trivial-cech},
chacun des complexes de la colimite de complexes de {\v C}ech est acyclique en degré $\geq 1$. Par Algèbre, lemme \ref{algebra-lemma-directed-colimit-exact},
le complexe de {\v C}ech
$\check{\mathcal{C}}^\bullet(\mathcal{U}, \colim_i \mathcal{I}_i)$
est acyclique en degrés $\geq 1$. En d'autres termes, nous voyons que
$\check{H}^p(\mathcal{U},  \colim_i \mathcal{I}_i) = 0$
pour tout $p \geq 1$. Les hypothèses du lemme \ref{cohomology-lemma-cech-vanish-basis} sont donc satisfaites, ce qui démontre l'assertion.

\end{proof}

\begin{lemma}

\label{cohomology-lemma-higher-direct-image-colimit}
Soit $f : X \to Y$ une application continue d'espaces topologiques. Soit $(\mathcal{F}_i, \varphi_{ii'})$ un système de faisceaux abéliens sur $X$. Posons $\mathcal{F} = \colim \mathcal{F}_i$. Soit $p \geq 0$ un entier. Supposons que les ouverts $V \subset Y$ tels que $H^p(f^{-1}(V), \mathcal{F}) = \colim H^p(f^{-1}(V), \mathcal{F}_i)$ forment une base de la topologie de $Y$. Alors $R^pf_*\mathcal{F} = \colim R^pf_*\mathcal{F}_i$.
\end{lemma}

\begin{proof}

Rappelons que $R^pf_*\mathcal{F}$ est la faisceautisation du préfaisceau
$\mathcal{G}$ qui associe à $V$ le groupe $H^p(f^{-1}(V), \mathcal{F})$ ; voir le lemme \ref{cohomology-lemma-describe-higher-direct-images}. De même, 
$R^pf_*\mathcal{F}_i$ est la faisceautisation du préfaisceau
$\mathcal{G}_i$ qui associe à $V$ le groupe $H^p(f^{-1}(V), \mathcal{F}_i)$. Rappelons que la faisceautisation est adjointe à gauche au foncteur d'inclusion des faisceaux dans les préfaisceaux ; voir Faisceaux, section
\ref{sheaves-section-sheafification}. La faisceautisation commute donc aux colimites ; voir Catégories, lemme \ref{categories-lemma-adjoint-exact}. Il suffit ainsi de montrer que le morphisme de préfaisceaux, la colimite étant prise dans la catégorie des préfaisceaux,
$$
\colim \mathcal{G}_i \longrightarrow \mathcal{G}
$$
induit un isomorphisme sur les faisceautisés. Il suffit pour cela de montrer que les préfaisceaux $\mathcal{G}$ et $\colim \mathcal{G}_i$
coïncident sur une base de la topologie de $Y$. En effet, les germes de leurs faisceautisés, calculables directement à partir des valeurs des préfaisceaux sur les éléments de la base, coïncident alors. L'égalité requise est précisément l'hypothèse du lemme.

\end{proof}

\noindent

Formulons maintenant l'analogue, pour la cohomologie, de Faisceaux, lemme \ref{sheaves-lemma-descend-opens}.
Soit $X$ un espace spectral, écrit comme limite cofiltrée d'espaces spectraux $X_i$ pour un diagramme dont les morphismes de transition sont spectraux, comme dans Topologie, lemme \ref{topology-lemma-directed-inverse-limit-spectral-spaces}. Supposons donnés

\begin{enumerate}

\item un faisceau abélien $\mathcal{F}_i$ sur $X_i$ pour tout
$i \in \Ob(\mathcal{I})$,
\item pour $a : j \to i$, un $f_a$-morphisme
$\varphi_a : \mathcal{F}_i \to \mathcal{F}_j$ de faisceaux abéliens (voir Faisceaux, définition \ref{sheaves-definition-f-map})

\end{enumerate}

tels que $\varphi_c = \varphi_b \circ \varphi_a$
dès que $c = a \circ b$. Posons $\mathcal{F} = \colim p_i^{-1}\mathcal{F}_i$
sur $X$.

\begin{lemma}

\label{cohomology-lemma-colimit}
Dans la situation ci-dessus, soit $i \in \Ob(\mathcal{I})$ et soit $U_i \subset X_i$ un ouvert quasi compact. Alors $$
\colim_{a : j \to i} H^p(f_a^{-1}(U_i), \mathcal{F}_j) =
H^p(p_i^{-1}(U_i), \mathcal{F})
$$ pour tous $p \geq 0$. En particulier, nous avons $H^p(X, \mathcal{F}) = \colim H^p(X_i, \mathcal{F}_i)$.
\end{lemma}

\begin{proof}
Le cas $p = 0$ est Faisceaux, lemme \ref{sheaves-lemma-descend-opens}.

\medskip\noindent

Dans ce paragraphe, montrons que l'on peut trouver un morphisme de systèmes
$(\gamma_i) : (\mathcal{F}_i, \varphi_a) \to (\mathcal{G}_i, \psi_a)$
où $\mathcal{G}_i$ est un faisceau abélien injectif et $\gamma_i$ est injectif. Pour chaque $i$, choisissons une injection $\mathcal{F}_i \to \mathcal{I}_i$,
où $\mathcal{I}_i$ est un faisceau abélien injectif sur $X_i$. Nous pouvons alors considérer la famille de morphismes
$$
\gamma_i :
\mathcal{F}_i
\longrightarrow
\prod\nolimits_{b : k \to i} f_{b, *}\mathcal{I}_k = \mathcal{G}_i
$$
où les morphismes composantes sont adjoints aux morphismes
$f_b^{-1}\mathcal{F}_i \to \mathcal{F}_k \to \mathcal{I}_k$. Pour $a : j \to i$ dans $\mathcal{I}$, il existe un morphisme canonique
$$
\psi_a : f_a^{-1}\mathcal{G}_i \to \mathcal{G}_j
$$
dont les composantes sont les morphismes canoniques
$f_b^{-1}f_{a \circ b, *}\mathcal{I}_k \to f_{b, *}\mathcal{I}_k$
pour $b : k \to j$. Nous obtenons ainsi une injection
$\{\gamma_i\} : \{\mathcal{F}_i, \varphi_a) \to (\mathcal{G}_i, \psi_a)$
de systèmes de faisceaux abéliens. Le faisceau de groupes abéliens $\mathcal{G}_i$ est injectif sur $X_i$ ; voir le lemme \ref{cohomology-lemma-pushforward-injective-flat} et Homologie, lemme \ref{homology-lemma-product-injectives}. Ceci termine la construction.

\medskip\noindent
En raisonnant exactement comme dans la démonstration du lemme \ref{cohomology-lemma-quasi-separated-cohomology-colimit}, il suffit de montrer que $H^p(X, \colim f_i^{-1}\mathcal{G}_i) = 0$ pour $p > 0$.

\medskip\noindent

Posons $\mathcal{G} = \colim f_i^{-1}\mathcal{G}_i$. Pour établir l'annulation de la cohomologie de $\mathcal{G}$ sur tout ouvert quasi compact de $X$, il suffit de montrer que la cohomologie de {\v C}ech de
$\mathcal{G}$, pour tout recouvrement ouvert $\mathcal{U}$ d'un ouvert quasi compact de
$X$ par un nombre fini d'ouverts quasi compacts, est nulle ; voir le lemme \ref{cohomology-lemma-cech-vanish-basis}. Un tel recouvrement est l'image inverse par $p_i$ d'un recouvrement $\mathcal{U}_i$
de l'espace $X_i$ pour un certain $i$, d'après Topologie, lemme \ref{topology-lemma-descend-opens}. Nous avons
$$
\check{\mathcal{C}}^\bullet(\mathcal{U}, \mathcal{G}) =
\colim_{a : j \to i}
\check{\mathcal{C}}^\bullet(f_a^{-1}(\mathcal{U}_i), \mathcal{G}_j)
$$
d'après le cas $p = 0$. Le membre de droite est une colimite filtrante de complexes, chacun acyclique en degrés positifs d'après le lemme \ref{cohomology-lemma-injective-trivial-cech}. On conclut alors par Algèbre, lemme \ref{algebra-lemma-directed-colimit-exact}.

\end{proof}
\section{Annulation sur les espaces topologiques noethériens}

\label{cohomology-section-vanishing-Noetherian}

\noindent
L'objectif est de prouver un théorème de Grothendieck, à savoir la proposition \ref{cohomology-proposition-vanishing-Noetherian}. Voir \cite{Tohoku}.

\begin{lemma}

\label{cohomology-lemma-cohomology-and-closed-immersions}
Soit $i : Z \to X$ une immersion fermée d'espaces topologiques. Pour tout faisceau abélien $\mathcal{F}$ sur $Z$, on a $H^p(Z, \mathcal{F}) = H^p(X, i_*\mathcal{F})$.
\end{lemma}

\begin{proof}
Cela résulte de l'exactitude de $i_*$ (voir Modules, lemme \ref{modules-lemma-i-star-exact}); par conséquent $R^pi_* = 0$ en tant que foncteur (Catégories dérivées, lemme \ref{derived-lemma-right-derived-exact-functor}). Nous pouvons donc appliquer le lemme \ref{cohomology-lemma-apply-Leray}.
\end{proof}

\begin{lemma}

\label{cohomology-lemma-irreducible-constant-cohomology-zero}
Soit $X$ un espace topologique irréductible. Alors $H^p(X, \underline{A}) = 0$ pour tout $p > 0$ et tout groupe abélien $A$.
\end{lemma}

\begin{proof}

Rappelons que $\underline{A}$ est le faisceau constant défini dans Faisceaux, définition \ref{sheaves-definition-constant-sheaf}. Comme $X$ est irréductible, tout ouvert non vide $U$ est irréductible et, a fortiori, connexe. Si $U \subset X$
est un ouvert non vide, alors $\underline{A}(U) = A$. De plus, $\underline{A}(\emptyset) = 0$. Ainsi $\underline{A}$
est un faisceau abélien flasque sur $X$. L'annulation résulte du lemme \ref{cohomology-lemma-flasque-acyclic}.

\end{proof}

\begin{lemma}

\label{cohomology-lemma-subsheaf-of-constant-sheaf}

\begin{reference}

\cite[Page 168]{Tohoku}.

\end{reference}

Soit $X$ un espace topologique tel que l'intersection de deux ouverts quasi-compacts est quasi-compact.
$\mathcal{F} \subset \underline{\mathbf{Z}}$
un sous-faisceau engendré par un nombre fini de sections sur des ouverts quasi-compacts.
$$
(0) = \mathcal{F}_0 \subset \mathcal{F}_1 \subset \ldots \subset
\mathcal{F}_n = \mathcal{F}
$$
par des sous-faisceaux abéliens tels que, pour chaque $0 < i \leq n$,
il existe une suite exacte courte
$$
0 \to j'_!\underline{\mathbf{Z}}_V \to j_!\underline{\mathbf{Z}}_U \to
\mathcal{F}_i/\mathcal{F}_{i - 1} \to 0
$$
avec $j : U \to X$ et $j' : V \to X$ l'inclusion des ouverts quasi-compacts dans $X$.

\end{lemma}

\begin{proof}

Supposons que $\mathcal{F}$ soit engendré par les sections $s_1, \ldots, s_t$ sur les ouverts quasi-compacts $U_1, \ldots, U_t$. Comme $U_i$ est quasi-compact et que
$s_i$ est une fonction localement constante à valeurs dans $\mathbf{Z}$, nous pouvons, quitte à remplacer $U_i$ par les morceaux d'une décomposition finie en parties ouvertes et fermées, supposer que $s_i$ est constante, disons $s_i = n_i$ avec $n_i \in \mathbf{Z}$. Nous pouvons supprimer
$(U_i, n_i)$ de la liste lorsque $n_i = 0$. En changeant les signes si nécessaire, nous pouvons aussi supposer $n_i > 0$. Ensuite, pour tout sous-ensemble $I \subset \{1, \ldots, t\}$,
nous pouvons ajouter à la liste $\bigcap_{i \in I} U_i$ et $\gcd(n_i, i \in I)$. Après cette opération, la liste $(U_1, n_1), \ldots, (U_t, n_t)$
possède la propriété suivante : pour $x \in X$, posons $I_x = \{i \in \{1, \ldots, t\} \mid x \in U_i\}$. Alors $\gcd(n_i, i \in I_x)$ est égal à l'un des $n_i$, pour un certain $i \in I_x$.

\medskip\noindent
Pour filtration, prenons $\mathcal{F}_0 = (0)$ et $\mathcal{F}_n$ le sous-faisceau engendré par les sections $n_i$ sur $U_i$ pour les $i$ tels que $n_i \leq n$. Il est clair que $\mathcal{F}_n = \mathcal{F}$ pour $n \gg 0$. Le quotient $\mathcal{F}_n/\mathcal{F}_{n - 1}$ est engendré par la section $n$ sur $U = \bigcup_{n_i \leq n} U_i$, et le noyau de $j_!\underline{\mathbf{Z}}_U \to \mathcal{F}_n/\mathcal{F}_{n - 1}$ est engendré par la section $n$ sur $V = \bigcup_{n_i \leq n - 1} U_i$. On obtient ainsi une suite exacte courte comme dans l'énoncé du lemme.
\end{proof}

\begin{lemma}

\label{cohomology-lemma-vanishing-generated-one-section}

\begin{reference}

Il s'agit d'un cas particulier: \cite[proposition 3.6.1]{Tohoku}.

\end{reference}

Soit $X$ un espace topologique et soit $d \geq 0$ un entier. Supposons que

\begin{enumerate}

\item  $X$ est quasi-compact,
\item les ouverts quasi-compacts constituent une base pour $X$,
\item l'intersection de deux ouverts quasi-compacts est quasi-compact.
\item  $H^p(X, j_!\underline{\mathbf{Z}}_U) = 0$ pour tous $p > d$
et tout ouvert quasi-compact $j : U \to X$.

\end{enumerate}

Alors $H^p(X, \mathcal{F}) = 0$ pour tout $p > d$
et tout faisceau abélien $\mathcal{F}$ sur $X$.

\end{lemma}

\begin{proof}

Posons $S = \coprod_{U \subset X} \mathcal{F}(U)$, où $U$ parcourt les ouverts quasi-compacts de $X$. Pour tout sous-ensemble fini $A = \{s_1, \ldots, s_n\} \subset S$, soit $\mathcal{F}_A$ le sous-faisceau de $\mathcal{F}$ engendré par les $s_i$ (voir Modules, définition \ref{modules-definition-generated-by-local-sections}). Si $A \subset A'$, alors $\mathcal{F}_A \subset \mathcal{F}_{A'}$. Ainsi, les $\{\mathcal{F}_A\}$ forment un système indexé par l'ensemble ordonné filtrant des sous-ensembles finis de $S$. D'après Modules, lemme \ref{modules-lemma-generated-by-local-sections-stalk},
il est clair que
$$
\colim_A \mathcal{F}_A = \mathcal{F}
$$
comme on le voit sur les germes. D'après le lemme \ref{cohomology-lemma-quasi-separated-cohomology-colimit}, on a
$$
H^p(X, \mathcal{F}) =
\colim_A H^p(X, \mathcal{F}_A)
$$
Il suffit donc de prouver l'annulation pour les faisceaux abéliens
$\mathcal{F}_A$. Autrement dit, il suffit de traiter le cas où $\mathcal{F}$ est engendré par un nombre fini de sections locales sur des ouverts quasi-compacts de $X$.

\medskip\noindent
Supposons que $\mathcal{F}$ soit engendré par les sections locales $s_1, \ldots, s_n$. Soit $\mathcal{F}' \subset \mathcal{F}$ le sous-faisceau engendré par $s_1, \ldots, s_{n - 1}$. Nous avons alors une suite exacte courte $$
0 \to \mathcal{F}' \to \mathcal{F} \to \mathcal{F}/\mathcal{F}' \to 0
$$ La suite exacte longue de cohomologie montre qu'il suffit d'établir l'annulation pour $\mathcal{F}'$ et $\mathcal{F}/\mathcal{F}'$, qui sont engendrés par moins de $n$ sections locales. Il suffit donc de traiter les faisceaux engendrés par une seule section locale au plus. Ce sont précisément les quotients des faisceaux $j_!\underline{\mathbf{Z}}_U$, où $U$ est un ouvert quasi-compact de $X$.

\medskip\noindent

Supposons maintenant donnée une suite exacte courte
$$
0 \to \mathcal{K} \to j_!\underline{\mathbf{Z}}_U \to \mathcal{F} \to 0
$$
où $U$ est un ouvert quasi-compact de $X$. Il suffit de montrer que $H^q(X, \mathcal{K})$ est nul pour $q \geq d + 1$. Comme ci-dessus, écrivons $\mathcal{K}$ comme colimite filtrante de sous-faisceaux $\mathcal{K}'$ engendrés par un nombre fini de sections sur des ouverts quasi-compacts. Alors $\mathcal{F}$ est la colimite filtrante des faisceaux $j_!\underline{\mathbf{Z}}_U/\mathcal{K}'$. Nous nous ramenons ainsi au cas où $\mathcal{K}$ est engendré par un nombre fini de sections sur des ouverts quasi-compacts. Remarquons que $\mathcal{K}$
est un sous-faisceau de $\underline{\mathbf{Z}}_X$. Le lemme \ref{cohomology-lemma-subsheaf-of-constant-sheaf} fournit donc une filtration finie de $\mathcal{K}$ dont chaque quotient successif $\mathcal{Q}$ s'insère dans une suite exacte courte
$$
0 \to j''_!\underline{\mathbf{Z}}_W \to
j'_!\underline{\mathbf{Z}}_V \to \mathcal{Q} \to 0
$$
avec $j'' : W \to X$ et $j' : V \to X$ les inclusions des ouverts quasi-compacts. L'annulation de $H^p(X, \mathcal{Q})$ pour $p > d$ découle alors de l'hypothèse du lemme appliquée aux groupes de cohomologie de $j''_!\underline{\mathbf{Z}}_W$ et $j'_!\underline{\mathbf{Z}}_V$. En revenant à $\mathcal{K}$, un raisonnement par récurrence utilisant la suite exacte longue de cohomologie donne l'annulation cherchée.

\end{proof}

\begin{example}

\label{cohomology-example-datta}

Soit $X = \mathbf{N}$ doté de la topologie dont les ouverts sont
$\emptyset$, $X$, $U_n = \{i \mid i \leq n\}$ pour $n \geq 1$. Un faisceau abélien $\mathcal{F}$ sur $X$ équivaut à la donnée d'un système projectif de groupes abéliens $A_n = \mathcal{F}(U_n)$ et
$\Gamma(X, \mathcal{F}) = \lim A_n$. Comme le foncteur de limite projective n'est pas exact sur la catégorie des systèmes projectifs, il existe un faisceau abélien dont $H^1$ est non nul. Enfin, le lecteur peut vérifier que $H^p(X, j_!\mathbf{Z}_U) = 0$,
$p \geq 1$ si $j : U = U_n \to X$ est l'inclusion. Ainsi, nous voyons que $X$ est un exemple d'espace satisfaisant aux conditions (2), (3) et (4) du lemme \ref{cohomology-lemma-vanishing-generated-one-section} pour $d = 0$
mais non à sa conclusion.

\end{example}

\begin{lemma}

\label{cohomology-lemma-subsheaf-irreducible}
Soit $X$ un espace topologique irréductible. Soit $\mathcal{H} \subset \underline{\mathbf{Z}}$ un sous-faisceau abélien du faisceau constant. Il existe alors un ouvert non vide $U \subset X$ tel que $\mathcal{H}|_U = \underline{d\mathbf{Z}}_U$ pour un certain $d \in \mathbf{Z}$.
\end{lemma}

\begin{proof}
Rappelons que $\underline{\mathbf{Z}}(V) = \mathbf{Z}$ pour tout ouvert non vide $V$ de $X$ (voir la preuve du lemme \ref{cohomology-lemma-irreducible-constant-cohomology-zero}). Si $\mathcal{H} = 0$, le lemme vaut avec $d = 0$. Si $\mathcal{H} \not = 0$, il existe un ouvert non vide $U \subset X$ tel que $\mathcal{H}(U) \not = 0$. Écrivons $\mathcal{H}(U) = n\mathbf{Z}$ pour un certain $n \geq 1$. On a donc $\underline{n\mathbf{Z}}_U
\subset \mathcal{H}|_U \subset
\underline{\mathbf{Z}}_U$. Si la première inclusion est stricte, nous pouvons trouver un $U' \subset U$ non vide et un entier $1 \leq n' < n$ tel que $\underline{n'\mathbf{Z}}_{U'}
\subset \mathcal{H}|_{U'} \subset
\underline{\mathbf{Z}}_{U'}$. Ce processus s'arrête après un nombre fini d'étapes, ce qui prouve le lemme.
\end{proof}

\begin{proposition}
[Grothendieck]
\label{cohomology-proposition-vanishing-Noetherian}

\begin{reference}
\cite[théorème 3.6.5]{Tohoku}.
\end{reference}
Soit $X$ un espace topologique noethérien. Si $\dim(X) \leq d$, alors $H^p(X, \mathcal{F}) = 0$ pour tous les $p > d$ et tout faisceau abélien $\mathcal{F}$ sur $X$.
\end{proposition}

\begin{proof}
Nous prouvons ce lemme par induction sur $d$. Fixez donc $d$ et supposez que le lemme tient pour tous les espaces topologiques noethériens de dimension $< d$.

\medskip\noindent
Soit $\mathcal{F}$ un faisceau abélien sur $X$. Supposons que $U \subset X$ soit un ouvert et notons $Z \subset X$ son complémentaire fermé. Notons $j : U \to X$ et $i : Z \to X$ les applications d'inclusion. On a alors une suite exacte courte $$
0 \to j_{!}j^*\mathcal{F} \to \mathcal{F} \to i_*i^*\mathcal{F} \to 0
$$ voir Modules, lemme \ref{modules-lemma-canonical-exact-sequence}. Notez que $j_!j^*\mathcal{F}$ est pris en charge sur la fermeture topologique $Z'$ de $U$, c'est-à-dire qu'il est de la forme $i'_*\mathcal{F}'$ pour certains faisceau abélien $\mathcal{F}'$ sur $Z'$, où $i' : Z' \to X$ est l'inclusion.

\medskip\noindent

Nous pouvons utiliser cela pour nous ramener au cas où $X$ est irréductible. En effet, d'après Topologie, lemme \ref{topology-lemma-Noetherian},

$X$ possède un nombre fini de composantes irréductibles. Si $X$ en possède plus d'une, soit $Z \subset X$ une composante irréductible de $X$
et posons $U = X \setminus Z$. D'après ce qui précède et la suite exacte longue de cohomologie, il suffit de prouver l'annulation de
$H^p(X, i_*i^*\mathcal{F})$ et $H^p(X, i'_*\mathcal{F}')$ pour $p > d$. D'après le lemme \ref{cohomology-lemma-cohomology-and-closed-immersions} il suffit de prouver
$H^p(Z, i^*\mathcal{F})$ et de $H^p(Z', \mathcal{F}')$ pour $p > d$. Comme $Z'$ et $Z$ ont moins de composantes irréductibles, nous nous ramenons bien au cas où $X$ est irréductible.

\medskip\noindent
Si $d = 0$ et si $X$ est irréductible, alors $X$ est le seul ouvert non vide de $X$. Par conséquent, tout faisceau est constant et les groupes de cohomologie supérieurs s'annulent (par exemple d'après le lemme \ref{cohomology-lemma-irreducible-constant-cohomology-zero}).

\medskip\noindent
Supposons que $X$ soit irréductible de dimension $d > 0$. D'après le lemme \ref{cohomology-lemma-vanishing-generated-one-section}, nous nous ramenons au cas où $\mathcal{F} = j_!\underline{\mathbf{Z}}_U$ pour un ouvert $U \subset X$. Dans ce cas, considérons la suite exacte courte $$
0 \to j_!(\underline{\mathbf{Z}}_U) \to
\underline{\mathbf{Z}}_X \to i_*\underline{\mathbf{Z}}_Z \to 0
$$ où $Z = X \setminus U$. D'après le lemme \ref{cohomology-lemma-irreducible-constant-cohomology-zero}, on a l'annulation de $H^p(X, \underline{\mathbf{Z}}_X)$ pour tout $p \geq 1$. Par récurrence, $H^p(X, i_*\underline{\mathbf{Z}}_Z) = H^p(Z, \underline{\mathbf{Z}}_Z) = 0$ pour $p \geq d$. Le résultat découle donc de la suite exacte longue de cohomologie.
\end{proof}





\section{Cohomologie à l'appui dans une partie fermée}

\label{cohomology-section-cohomology-support}

\noindent
Cette section traite juste du minimum -- la discussion se poursuivra dans la section \ref{cohomology-section-cohomology-support-bis}.

\medskip\noindent
Soit $X$ un espace topologique, soit $Z \subset X$ une partie fermée et soit $\mathcal{F}$ un faisceau abélien sur $X$. Posons $$
\Gamma_Z(X, \mathcal{F}) =
\{s \in \mathcal{F}(X) \mid \text{Supp}(s) \subset Z\}
$$ le sous-ensemble des sections dont le support est contenu dans $Z$. Le support d'une section est défini dans Modules, définition \ref{modules-definition-support}. Modules, lemme \ref{modules-lemma-support-section-closed} implique que $\Gamma_Z(X, \mathcal{F})$ est un sous-groupe de $\Gamma(X, \mathcal{F})$. Le même lemme garantit que l'association $\mathcal{F} \mapsto \Gamma_Z(X, \mathcal{F})$ est fonctorielle en $\mathcal{F}$. Ce foncteur est exact à gauche, mais n'est pas exact en général.

\medskip\noindent

Comme la catégorie des faisceaux abéliens possède assez d'injectifs (Injectifs, lemme \ref{injectives-lemma-abelian-sheaves-space}), on obtient un foncteur dérivé à droite
$$
R\Gamma_Z(X, -) : D^+(X) \longrightarrow D^+(\textit{Ab})
$$
d'après Catégories dérivées, lemme \ref{derived-lemma-enough-injectives-right-derived}. La valeur de $R\Gamma_Z(X, -)$ sur un objet $K$ se calcule en représentant
$K$ par un complexe borné inférieurement $\mathcal{I}^\bullet$ de faisceaux abéliens injectifs, puis en prenant $\Gamma_Z(X, \mathcal{I}^\bullet)$; voir Catégories dérivées, lemme \ref{derived-lemma-injective-acyclic}. Les groupes de cohomologie d'un faisceau abélien $\mathcal{F}$
à support dans $Z$ sont définis par
$H^q_Z(X, \mathcal{F}) = R^q\Gamma_Z(X, \mathcal{F})$.

\medskip\noindent
Soit $\mathcal{I}$ un faisceau abélien injectif sur $X$, et posons $U = X \setminus Z$. L'application de restriction $\mathcal{I}(X) \to \mathcal{I}(U)$ est surjective (lemme \ref{cohomology-lemma-injective-restriction-surjective}) et a pour noyau $\Gamma_Z(X, \mathcal{I})$. Il s'ensuit immédiatement que, pour $K \in D^+(X)$, on a un triangle distingué $$
R\Gamma_Z(X, K) \to R\Gamma(X, K) \to R\Gamma(U, K) \to R\Gamma_Z(X, K)[1]
$$ dans $D^+(\textit{Ab})$. Par conséquent, nous obtenons une suite exacte longue de cohomologie $$
\ldots \to H^i_Z(X, K) \to H^i(X, K) \to H^i(U, K) \to
H^{i + 1}_Z(X, K) \to \ldots
$$ pour tout $K$ dans $D^+(X)$.

\medskip\noindent

Pour un faisceau abélien $\mathcal{F}$ sur $X$, nous pouvons considérer le
{\it sous-faisceau des sections à support dans $Z$}, dénoté
$\mathcal{H}_Z(\mathcal{F})$, défini par la règle
$$
\mathcal{H}_Z(\mathcal{F})(U) =
\{s \in \mathcal{F}(U) \mid \text{Supp}(s) \subset U \cap Z\} =
\Gamma_{Z \cap U}(U, \mathcal{F}|_U)
$$
En utilisant l'équivalence de Modules, lemme \ref{modules-lemma-i-star-exact},
nous pouvons considérer $\mathcal{H}_Z(\mathcal{F})$ comme un faisceau abélien sur $Z$; voir Modules, remarque \ref{modules-remark-sections-support-in-closed}. Ainsi, nous obtenons un foncteur
$$
\textit{Ab}(X) \longrightarrow \textit{Ab}(Z),\quad
\mathcal{F} \longmapsto
\mathcal{H}_Z(\mathcal{F})\text{ considéré comme faisceau sur }Z
$$
Ce foncteur est exact à gauche, mais n'est pas exact en général. Comme ci-dessus, nous obtenons un foncteur dérivé à droite
$$
R\mathcal{H}_Z : D^+(X) \longrightarrow D^+(Z)
$$
et nous posons
$\mathcal{H}^q_Z(\mathcal{F}) = R^q\mathcal{H}_Z(\mathcal{F})$, de sorte que
$\mathcal{H}^0_Z(\mathcal{F}) = \mathcal{H}_Z(\mathcal{F})$.

\medskip\noindent
On a $\Gamma_Z(X, \mathcal{F}) = \Gamma(Z, \mathcal{H}_Z(\mathcal{F}))$ pour tout faisceau abélien $\mathcal{F}$. D'après le lemme \ref{cohomology-lemma-sections-with-support-acyclic}, le foncteur $\mathcal{H}_Z$ transforme les faisceaux abéliens injectifs en faisceaux acycliques pour $\Gamma(Z, -)$. Le lemme de Catégories dérivées \ref{derived-lemma-grothendieck-spectral-sequence} fournit donc une suite spectrale de Grothendieck convergente $$
E_2^{p, q} = H^p(Z, \mathcal{H}^q_Z(K)) \Rightarrow H^{p + q}_Z(X, K)
$$ fonctorielle en $K$ dans $D^+(X)$.

\begin{lemma}

\label{cohomology-lemma-sections-with-support-acyclic}
Soit $i : Z \to X$ l'inclusion d'une partie fermée, et soit $\mathcal{I}$ un faisceau abélien injectif sur $X$. Alors $\mathcal{H}_Z(\mathcal{I})$ est un faisceau abélien injectif sur $Z$.
\end{lemma}

\begin{proof}
Cela résulte du lemme d'Homologie \ref{homology-lemma-adjoint-preserve-injectives}, car $\mathcal{H}_Z(-)$ est adjoint à droite du foncteur exact $i_*$. Voir Modules, lemmes \ref{modules-lemma-i-star-exact} et \ref{modules-lemma-i-star-right-adjoint}.
\end{proof}



\section{Cohomologie sur les espaces spectraux}

\label{cohomology-section-spectral}

\noindent
Un résultat clé sur la cohomologie des espaces spectraux est le lemme \ref{cohomology-lemma-colimit}, qui affirme, en termes informels, que la cohomologie commute aux limites cofiltrantes dans la catégorie des espaces spectraux définie dans Topologie, définition \ref{topology-definition-spectral-space}. On peut l'appliquer pour obtenir les analogues suivants des lemmes \ref{cohomology-lemma-cohomology-of-closed} et \ref{cohomology-lemma-proper-base-change}.

\begin{lemma}

\label{cohomology-lemma-cohomology-of-neighbourhoods-of-closed}
Soit $X$ un espace spectral, soit $\mathcal{F}$ un faisceau abélien sur $X$, soit $E \subset X$ un sous-ensemble quasi-compact et soit $W \subset X$ l'ensemble des points de $X$ qui se spécialisent en un point de $E$.
\begin{enumerate}
\item $H^p(W, \mathcal{F}|_W) = \colim H^p(U, \mathcal{F})$ où la colimite est sur les voisinages quasi-compacts ouverts de $E$,
\item $H^p(W \setminus E, \mathcal{F}|_{W \setminus E}) =
\colim H^p(U \setminus E, \mathcal{F}|_{U \setminus E})$ si $E$ est un sous-ensemble constructible.
\end{enumerate}

\end{lemma}

\begin{proof}
D'après Topologie, lemme \ref{topology-lemma-make-spectral-space}, on a $W = \lim U$, où la limite porte sur les ouverts quasi-compacts contenant $E$. Chaque $U$ est un espace spectral d'après Topologie, lemme \ref{topology-lemma-spectral-sub}. Le lemme \ref{cohomology-lemma-colimit} donne donc (1). La même démonstration vaut pour (2), en utilisant Topologie, lemme \ref{topology-lemma-make-spectral-space-minus}.
\end{proof}

\begin{lemma}

\label{cohomology-lemma-proper-base-change-spectral}
Soit $f : X \to Y$ une application spectrale entre espaces spectraux. Soit $y \in Y$, et soit $E \subset Y$ l'ensemble des points qui se spécialisent en $y$. Pour tout faisceau abélien $\mathcal{F}$ sur $X$, on a $(R^pf_*\mathcal{F})_y = H^p(f^{-1}(E), \mathcal{F}|_{f^{-1}(E)})$.
\end{lemma}

\begin{proof}

On a $E = \bigcap V$, où $V$ parcourt les voisinages ouverts quasi-compacts de $y$ dans $Y$. Ainsi $f^{-1}(E) = \bigcap f^{-1}(V)$, donc $f^{-1}(E) = \lim f^{-1}(V)$ comme espace topologique. Puisque $f$ est spectrale, chaque $f^{-1}(V)$ est encore un espace spectral (Topologie, lemme \ref{topology-lemma-spectral-sub}). On en déduit que $f^{-1}(E)$ est un espace spectral et que
$$
H^p(f^{-1}(E), \mathcal{F}|_{f^{-1}(E)}) =
\colim H^p(f^{-1}(V), \mathcal{F})
$$
par lemme \ref{cohomology-lemma-colimit}. D'autre part, le germe de
$R^pf_*\mathcal{F}$ en $y$ est donné par la colimite du membre de droite.

\end{proof}

\begin{lemma}

\label{cohomology-lemma-vanishing-for-profinite}
Soit $X$ un espace topologique profini. Alors $H^q(X, \mathcal{F}) = 0$ pour tout $q > 0$ et tout faisceau abélien $\mathcal{F}$.
\end{lemma}

\begin{proof}
Tout recouvrement ouvert de $X$ admet un raffinement qui est une décomposition finie en parties ouvertes disjointes; voir Topologie, lemme \ref{topology-lemma-profinite-refine-open-covering}. Ainsi, si $\mathcal{F} \to \mathcal{G}$ est une surjection de faisceaux abéliens sur $X$, l'application $\mathcal{F}(X) \to \mathcal{G}(X)$ est surjective. Autrement dit, le foncteur des sections globales est exact. Ses foncteurs dérivés supérieurs sont donc nuls; voir Catégories dérivées, lemme \ref{derived-lemma-right-derived-exact-functor}.
\end{proof}

\noindent
Le résultat suivant sur l'annulation cohomologique améliore le résultat de Grothendieck (proposition \ref{cohomology-proposition-vanishing-Noetherian}) et se trouve dans \cite{Scheiderer}.

\begin{proposition}

\label{cohomology-proposition-cohomological-dimension-spectral}

\begin{reference}

La partie 1 est le théorème principal de \cite{Scheiderer}.

\end{reference}

Soit $X$ un espace spectral de dimension de Krull $d$. Soit $\mathcal{F}$ un faisceau abélien sur $X$.

\begin{enumerate}

\item  $H^q(X, \mathcal{F}) = 0$ pour $q > d$,
\item  $H^d(X, \mathcal{F}) \to H^d(U, \mathcal{F})$ est surjectif pour chaque ouvert quasi-compact $U \subset X$,
\item  $H^q_Z(X, \mathcal{F}) = 0$ pour $q > d$ et toute partie fermée constructible $Z \subset X$.

\end{enumerate}

\end{proposition}

\begin{proof}
Nous prouvons ce résultat par induction sur $d$.

\medskip\noindent

Si $d = 0$, alors $X$ est un espace profini; voir Topologie, lemme \ref{topology-lemma-characterize-profinite-spectral}. Le point (1) résulte donc du lemme \ref{cohomology-lemma-vanishing-for-profinite}. Si $U \subset X$ est un ouvert quasi-compact, alors $U$ est aussi fermé, puisqu'il est quasi-compact dans un espace séparé. Ainsi $X = U \amalg (X \setminus U)$, ce qui prouve (2). Pour $Z$ comme dans (3), considérons la suite exacte longue
$$
H^{q - 1}(X, \mathcal{F}) \to
H^{q - 1}(X \setminus Z, \mathcal{F}) \to
H^q_Z(X, \mathcal{F}) \to H^q(X, \mathcal{F})
$$
Comme $X$ et $U = X \setminus Z$ sont profinis (en effet, $U$ est quasi-compact puisque $Z$ est constructible), les points (2) et (1) donnent l'annulation cherchée des groupes de cohomologie à support dans $Z$.

\medskip\noindent

Étape d'induction. Supposons $d \geq 1$ et la proposition vraie pour tout espace spectral de dimension $< d$. Prouvons d'abord (2) pour $X$. Soit $U$ un ouvert quasi-compact et soit $\xi \in H^d(U, \mathcal{F})$. Posons $Z = X \setminus U$, et soit $W \subset X$ l'ensemble des points qui se spécialisent dans $Z$. D'après le lemme \ref{cohomology-lemma-cohomology-of-neighbourhoods-of-closed}, on a
$$
H^d(W \setminus Z, \mathcal{F}|_{W \setminus Z}) =
\colim_{Z \subset V} H^d(V \setminus Z, \mathcal{F})
$$
où la colimite porte sur les voisinages ouverts quasi-compacts $V$ de $Z$ dans $X$.
D'après Topologie, lemme \ref{topology-lemma-make-spectral-space},
$W \setminus Z$ est un espace spectral. Comme tout point de $W$ se spécialise en un point de $Z$,
$W \setminus Z$ est de dimension de Krull $< d$. Par l'hypothèse d'induction, l'image de $\xi$ dans
$H^d(W \setminus Z, \mathcal{F}|_{W \setminus Z})$ est nulle. Par la formule affichée, il existe un ouvert quasi-compact $Z \subset V \subset X$
tel que $\xi|_{V \setminus Z} = 0$. Comme $V \setminus Z = V \cap U$, la suite de Mayer–Vietoris (lemme \ref{cohomology-lemma-mayer-vietoris}) pour le recouvrement $X = U \cup V$
fournit un élément $\tilde \xi \in H^d(X, \mathcal{F})$ qui se restreint à $\xi$ sur $U$ et à zéro sur $V$. Le point (2) est donc établi.

\medskip\noindent

Preuve de (1), en supposant (2). Choisissons une résolution injective
$\mathcal{F} \to \mathcal{I}^\bullet$. Posons
$$
\mathcal{G} = \Im(\mathcal{I}^{d - 1} \to \mathcal{I}^d) =
\Ker(\mathcal{I}^d \to \mathcal{I}^{d + 1})
$$
Pour tout ouvert quasi-compact $U \subset X$, nous avons le morphisme suivant de suites exactes :
$$
\xymatrix{
\mathcal{I}^{d - 1}(X) \ar[r] \ar[d] &
\mathcal{G}(X) \ar[r] \ar[d] &
H^d(X, \mathcal{F}) \ar[d] \ar[r] & 0 \\
\mathcal{I}^{d - 1}(U) \ar[r] &
\mathcal{G}(U) \ar[r] &
H^d(U, \mathcal{F}) \ar[r] & 0
}
$$
Le faisceau $\mathcal{I}^{d - 1}$ est flasque, par le lemme \ref{cohomology-lemma-injective-flasque} et puisque $d \geq 1$. Le point (2) montre que la flèche verticale de droite est surjective. Une chasse au diagramme montre alors que l'application
$\mathcal{G}(X) \to \mathcal{G}(U)$ est surjective. D'après le lemme \ref{cohomology-lemma-vanishing-ravi}, on en déduit que
$\check{H}^q(\mathcal{U}, \mathcal{G}) = 0$ pour $q > 0$ et tout recouvrement fini $\mathcal{U} : U = U_1 \cup \ldots \cup U_n$
d'un ouvert quasi-compact par des ouverts quasi-compacts. D'après le lemme \ref{cohomology-lemma-cech-vanish-basis}, on a $H^q(U, \mathcal{G}) = 0$
pour tout $q > 0$ et tout ouvert quasi-compact $U$ de $X$. Le lemme d'acyclicité de Leray (Catégories dérivées, lemme \ref{derived-lemma-leray-acyclicity}) donne alors
$$
H^q(X, \mathcal{F}) =
H^q\left(
\Gamma(X, \mathcal{I}^0) \to \ldots \to
\Gamma(X, \mathcal{I}^{d - 1}) \to \Gamma(X, \mathcal{G})
\right)
$$
En particulier, ce groupe de cohomologie est nul si $q > d$.

\medskip\noindent

Démonstration de (3). Pour $Z$ comme dans (3), considérons la suite exacte longue
$$
H^{q - 1}(X, \mathcal{F}) \to
H^{q - 1}(X \setminus Z, \mathcal{F}) \to
H^q_Z(X, \mathcal{F}) \to H^q(X, \mathcal{F})
$$
Comme $X$ et $U = X \setminus Z$ sont des espaces spectraux (Topologie, lemme \ref{topology-lemma-spectral-sub}) de dimension $\leq d$,
les points (2) et (1) donnent l'annulation cherchée.

\end{proof}













\section{Le complexe de {\v C}
ech modifié}
\label{cohomology-section-alternating-cech}

\noindent
Cette section compare le complexe {\v C}ech avec le complexe {\v C}ech alternatif et certains complexes connexes.

\medskip\noindent
Soit $X$ un espace topologique et soit $\mathcal{U} : U = \bigcup_{i \in I} U_i$ un recouvrement ouvert. Pour $p \geq 0$, posons $$
\check{\mathcal{C}}_{alt}^p(\mathcal{U}, \mathcal{F})
=
\left\{
\begin{matrix}
s \in  \check{\mathcal{C}}^p(\mathcal{U}, \mathcal{F})
\text{ tel que }
s_{i_0 \ldots i_p} = 0 \text{ si } i_n = i_m \text{ pour certains } n \not = m\\
\text{ et }
s_{i_0\ldots i_n \ldots i_m \ldots i_p}
=
-s_{i_0\ldots i_m \ldots i_n \ldots i_p}
\text{ dans tous les cas.}
\end{matrix}
\right\}
$$ Nous omettons de vérifier que le différentiel $d$ de l'équation (\ref{cohomology-equation-d-cech}) envoie $\check{\mathcal{C}}^p_{alt}(\mathcal{U}, \mathcal{F})$ dans $\check{\mathcal{C}}^{p + 1}_{alt}(\mathcal{U}, \mathcal{F})$.

\begin{definition}

\label{cohomology-definition-alternating-cech-complex}
Soit $X$ un espace topologique, soit $\mathcal{U} : U = \bigcup_{i \in I} U_i$ un recouvrement ouvert, et soit $\mathcal{F}$ un préfaisceau abélien sur $X$. Le complexe $\check{\mathcal{C}}_{alt}^\bullet(\mathcal{U}, \mathcal{F})$ est le {\it complexe de {\v C}ech alterné} associé à $\mathcal{F}$ et au recouvrement $\mathcal{U}$.
\end{definition}

\noindent
Il y a donc un morphisme canonique des complexes $$
\check{\mathcal{C}}_{alt}^\bullet(\mathcal{U}, \mathcal{F})
\longrightarrow
\check{\mathcal{C}}^\bullet(\mathcal{U}, \mathcal{F})
$$ à savoir l'inclusion du complexe {\v C}ech alternatif dans le complexe {\v C}ech habituel.

\medskip\noindent
Supposons le recouvrement $\mathcal{U} : U = \bigcup_{i \in I} U_i$ muni d'un ordre total $<$ sur $I$. Posons alors $$
\check{\mathcal{C}}_{ord}^p(\mathcal{U}, \mathcal{F})
=
\prod\nolimits_{(i_0, \ldots, i_p) \in I^{p + 1}, i_0 < \ldots < i_p}
\mathcal{F}(U_{i_0\ldots i_p}).
$$ C'est un groupe abélien. Pour $s \in \check{\mathcal{C}}_{ord}^p(\mathcal{U}, \mathcal{F})$, nous notons $s_{i_0\ldots i_p}$ sa valeur dans $\mathcal{F}(U_{i_0\ldots i_p})$. Nous définissons $$
d : \check{\mathcal{C}}_{ord}^p(\mathcal{U}, \mathcal{F})
\longrightarrow
\check{\mathcal{C}}_{ord}^{p + 1}(\mathcal{U}, \mathcal{F})
$$ par la formule $$
d(s)_{i_0\ldots i_{p + 1}}
=
\sum\nolimits_{j = 0}^{p + 1}
(-1)^j
s_{i_0\ldots \hat i_j \ldots i_{p + 1}}|_{U_{i_0\ldots i_{p + 1}}}
$$ pour n'importe quelle $i_0 < \ldots < i_{p + 1}$. Notez que cette formule est identique à l'équation (\ref{cohomology-equation-d-cech}). Il est simple de voir que $d \circ d = 0$. En d'autres termes, $\check{\mathcal{C}}_{ord}^\bullet(\mathcal{U}, \mathcal{F})$ est un complexe.

\begin{definition}

\label{cohomology-definition-ordered-cech-complex}
Soit $X$ un espace topologique, soit $\mathcal{U} : U = \bigcup_{i \in I} U_i$ un recouvrement ouvert, et supposons $I$ muni d'un ordre total. Soit $\mathcal{F}$ un préfaisceau abélien sur $X$. Le complexe $\check{\mathcal{C}}_{ord}^\bullet(\mathcal{U}, \mathcal{F})$ est le {\it complexe de {\v C}ech ordonné} associé à $\mathcal{F}$, au recouvrement $\mathcal{U}$ et à l'ordre total donné sur $I$.
\end{definition}

\noindent
Ce complexe est parfois appelé le complexe alterné de {\v C}ech. La raison est qu'il existe une application de comparaison évidente entre le complexe ordonné de {\v C}ech et le complexe alterné de {\v C}ech. Plus précisément, considérons l'application $$
c :
\check{\mathcal{C}}_{ord}^\bullet(\mathcal{U}, \mathcal{F})
\longrightarrow
\check{\mathcal{C}}^\bullet(\mathcal{U}, \mathcal{F})
$$ donnée par la règle $$
c(s)_{i_0\ldots i_p} =
\left\{
\begin{matrix}
0 &
\text{si} &
i_n = i_m \text{ pour certains } n \not = m\\
\text{sgn}(\sigma) s_{i_{\sigma(0)}\ldots i_{\sigma(p)}} &
\text{si} &
i_{\sigma(0)} < i_{\sigma(1)} < \ldots < i_{\sigma(p)}
\end{matrix}
\right.
$$ Ici, $\sigma$ désigne une permutation de $\{0, \ldots, p\}$ et $\text{sgn}(\sigma)$ désigne son signe. Les complexes alterné et ordonné de {\v C}ech sont souvent identifiés dans la littérature au moyen de l'application $c$. Nous avons en effet le lemme élémentaire suivant.

\begin{lemma}

\label{cohomology-lemma-ordered-alternating}
Soit $X$ un espace topologique, soit $\mathcal{U} : U = \bigcup_{i \in I} U_i$ un recouvrement ouvert, et supposons $I$ muni d'un ordre total. L'application $c$ est un morphisme de complexes. Elle induit même un isomorphisme $$
c : \check{\mathcal{C}}_{ord}^\bullet(\mathcal{U}, \mathcal{F})
\to \check{\mathcal{C}}_{alt}^\bullet(\mathcal{U}, \mathcal{F})
$$ de complexes.
\end{lemma}

\begin{proof}
Démonstration omise.
\end{proof}

\noindent
Il y a aussi une application $$
\pi :
\check{\mathcal{C}}^\bullet(\mathcal{U}, \mathcal{F})
\longrightarrow
\check{\mathcal{C}}_{ord}^\bullet(\mathcal{U}, \mathcal{F})
$$ qui est décrite par la règle $$
\pi(s)_{i_0\ldots i_p} = s_{i_0\ldots i_p}
$$ chaque fois que $i_0 < i_1 < \ldots < i_p$.

\begin{lemma}

\label{cohomology-lemma-project-to-ordered}
Soit $X$ un espace topologique, soit $\mathcal{U} : U = \bigcup_{i \in I} U_i$ un recouvrement ouvert, et supposons $I$ muni d'un ordre total. La projection $\pi : \check{\mathcal{C}}^\bullet(\mathcal{U}, \mathcal{F})
\to \check{\mathcal{C}}_{ord}^\bullet(\mathcal{U}, \mathcal{F})$ est un morphisme de complexes. Il induit un isomorphisme $$
\pi : \check{\mathcal{C}}_{alt}^\bullet(\mathcal{U}, \mathcal{F})
\to \check{\mathcal{C}}_{ord}^\bullet(\mathcal{U}, \mathcal{F})
$$ de complexes qui est un inverse gauche au morphisme $c$.
\end{lemma}

\begin{proof}
Démonstration omise.
\end{proof}

\begin{remark}

\label{cohomology-remark-compared-ordered-complexes}
Cela signifie que si nous avons deux commandes totales $<_1$ et $<_2$ sur le jeu d'index $I$, alors nous obtenons un isomorphisme des complexes $\tau = \pi_2 \circ c_1 :
\check{\mathcal{C}}_{ord\text{-}1}(\mathcal{U}, \mathcal{F}) \to
\check{\mathcal{C}}_{ord\text{-}2}(\mathcal{U}, \mathcal{F})$. Il est clair que $$
\tau(s)_{i_0 \ldots i_p} =
\text{sign}(\sigma) s_{i_{\sigma(0)} \ldots i_{\sigma(p)}}
$$ où $i_0 <_1 i_1 <_1 \ldots <_1 i_p$ et $i_{\sigma(0)} <_2 i_{\sigma(1)} <_2 \ldots <_2 i_{\sigma(p)}$. C'est en ce sens que le complexe de {\v C}ech ordonné ne dépend pas de l'ordre total choisi.
\end{remark}

\begin{lemma}

\label{cohomology-lemma-alternating-usual}
Soit $X$ un espace topologique, soit $\mathcal{U} : U = \bigcup_{i \in I} U_i$ un recouvrement ouvert, et supposons $I$ muni d'un ordre total. L'application $c \circ \pi$ est homotope à l'identité de $\check{\mathcal{C}}^\bullet(\mathcal{U}, \mathcal{F})$. En particulier, l'inclusion $\check{\mathcal{C}}_{alt}^\bullet(\mathcal{U}, \mathcal{F}) \to
\check{\mathcal{C}}^\bullet(\mathcal{U}, \mathcal{F})$ est une équivalence d'homotopie.
\end{lemma}

\begin{proof}
Pour tout multi-index $(i_0, \ldots, i_p) \in I^{p + 1}$ il existe une permutation unique $\sigma : \{0, \ldots, p\} \to \{0, \ldots, p\}$ telle que $$
i_{\sigma(0)} \leq i_{\sigma(1)} \leq \ldots \leq i_{\sigma(p)}
\quad
\text{et}
\quad
\sigma(j) < \sigma(j + 1)
\quad
\text{si}
\quad
i_{\sigma(j)} = i_{\sigma(j + 1)}.
$$ Nous dénotons cette permutation $\sigma = \sigma^{i_0 \ldots i_p}$.

\medskip\noindent
Pour toute permutation $\sigma : \{0, \ldots, p\} \to \{0, \ldots, p\}$ et toute $a$, $0 \leq a \leq p$ nous dénotons $\sigma_a$ la permutation unique de $\{0, \ldots, p\}$ telle que $\sigma_a(j) = \sigma(j)$ pour $0 \leq j < a$ et telle que $\sigma_a(a) < \sigma_a(a + 1) < \ldots < \sigma_a(p)$. Donc si $p = 3$ et $\sigma$, $\tau$ sont donnés par $$
\begin{matrix}
\text{id} & 0 & 1 & 2 & 3 \\
\sigma & 3 & 2 & 1 & 0
\end{matrix}
\quad \text{et} \quad
\begin{matrix}
\text{id} & 0 & 1 & 2 & 3 \\
\tau & 3 & 0 & 2 & 1
\end{matrix}
$$ alors nous avons $$
\begin{matrix}
\text{id} & 0 & 1 & 2 & 3 \\
\sigma_0 & 0 & 1 & 2 & 3 \\
\sigma_1 & 3 & 0 & 1 & 2 \\
\sigma_2 & 3 & 2 & 0 & 1 \\
\sigma_3 & 3 & 2 & 1 & 0 \\
\end{matrix}
\quad \text{et} \quad
\begin{matrix}
\text{id} & 0 & 1 & 2 & 3 \\
\tau_0 & 0 & 1 & 2 & 3 \\
\tau_1 & 3 & 0 & 1 & 2 \\
\tau_2 & 3 & 0 & 1 & 2 \\
\tau_3 & 3 & 0 & 2 & 1 \\
\end{matrix}
$$ Il est clair que toujours $\sigma_0 = \text{id}$ et $\sigma_p = \sigma$.

\medskip\noindent

Après avoir introduit cette notation, nous définissons pour
$s \in \check{\mathcal{C}}^{p + 1}(\mathcal{U}, \mathcal{F})$
l'élément $h(s) \in \check{\mathcal{C}}^p(\mathcal{U}, \mathcal{F})$
comme l'élément dont les composantes sont

\begin{equation}
\label{cohomology-equation-first-homotopy}
h(s)_{i_0\ldots i_p} =
\sum\nolimits_{0 \leq a \leq p}
(-1)^a \text{sign}(\sigma_a)
s_{i_{\sigma(0)} \ldots i_{\sigma(a)} i_{\sigma_a(a)} \ldots i_{\sigma_a(p)}}
\end{equation}

où $\sigma = \sigma^{i_0 \ldots i_p}$. L'indice
$i_{\sigma(a)}$ survient deux fois dans
$i_{\sigma(0)} \ldots i_{\sigma(a)} i_{\sigma_a(a)} \ldots i_{\sigma_a(p)}$
une fois dans le premier groupe de $a + 1$ indices et une fois dans le second groupe de $p - a + 1$ indices, puisque $\sigma_a(j) = \sigma(a)$ pour un certain
$j \geq a$ par définition de $\sigma_a$. La somme a donc un sens, car chacun des éléments
$s_{i_{\sigma(0)} \ldots i_{\sigma(a)} i_{\sigma_a(a)} \ldots i_{\sigma_a(p)}}$
est défini sur l'ouvert $U_{i_0 \ldots i_p}$. Notons également que, pour $a = 0$, nous obtenons $s_{i_0 \ldots i_p}$ et que, pour $a = p$, nous obtenons
$(-1)^p \text{sign}(\sigma) s_{i_{\sigma(0)} \ldots i_{\sigma(p)}}$.

\medskip\noindent

Nous affirmons que :
$$
(dh + hd)(s)_{i_0 \ldots i_p} =
s_{i_0 \ldots i_p} -
\text{sign}(\sigma) s_{i_{\sigma(0)} \ldots i_{\sigma(p)}}
$$
où $\sigma = \sigma^{i_0 \ldots i_p}$. Nous omettons la vérification de cette revendication. (Il y a un script PARI/gp appelé first-homotopy.gp dans les scripts de sous-répertoire stacks-project qui peuvent être utilisés pour vérifier de nombreuses instances finies de cette revendication. Nous avons écrit ce script pour nous assurer que les signes sont corrects.)
$$
\kappa :
\check{\mathcal{C}}^\bullet(\mathcal{U}, \mathcal{F})
\longrightarrow
\check{\mathcal{C}}^\bullet(\mathcal{U}, \mathcal{F})
$$
pour l'opérateur donné par la règle
$$
\kappa(s)_{i_0 \ldots i_p} =
\text{sign}(\sigma^{i_0 \ldots i_p}) s_{i_{\sigma(0)} \ldots i_{\sigma(p)}}.
$$
L'allégation ci-dessus implique que: $\kappa$ est un morphisme de complexes et que
$\kappa$ est homotope à l'identité du complexe de {\v C}ech. Cela n'implique pas immédiatement le lemme, car l'image de l'opérateur $\kappa$ n'est pas le sous-complexe alterné. Plus précisément, l'image de $\kappa$ est le complexe « semi-alterné »
$\check{\mathcal{C}}_{semi\text{-}alt}^p(\mathcal{U}, \mathcal{F})$
où $s$ est une $p$-cochaîne de ce complexe si et seulement si
$$
s_{i_0 \ldots i_p} = \text{sign}(\sigma) s_{i_{\sigma(0)} \ldots i_{\sigma(p)}}
$$
pour tout $(i_0, \ldots, i_p) \in I^{p + 1}$ avec
$\sigma = \sigma^{i_0 \ldots i_p}$. Nous introduisons une autre variante du complexe de {\v C}ech, à savoir le complexe de {\v C}ech semi-ordonné défini par
$$
\check{\mathcal{C}}_{semi\text{-}ord}^p(\mathcal{U}, \mathcal{F})
=
\prod\nolimits_{i_0 \leq i_1 \leq \ldots \leq i_p}
\mathcal{F}(U_{i_0 \ldots i_p})
$$
Il est facile de voir que l'équation (\ref{cohomology-equation-d-cech}) définit encore une différentielle, donc un complexe. Il est également clair, comme dans le lemme \ref{cohomology-lemma-project-to-ordered}, que l'application de projection
$$
\check{\mathcal{C}}_{semi\text{-}alt}^\bullet(\mathcal{U}, \mathcal{F})
\longrightarrow
\check{\mathcal{C}}_{semi\text{-}ord}^\bullet(\mathcal{U}, \mathcal{F})
$$
est un isomorphisme des complexes.

\medskip\noindent

Par conséquent, le lemme suit si nous pouvons montrer que l'application d'inclusion évidente
$$
\check{\mathcal{C}}_{ord}^p(\mathcal{U}, \mathcal{F})
\longrightarrow
\check{\mathcal{C}}_{semi\text{-}ord}^p(\mathcal{U}, \mathcal{F})
$$
est une équivalence d'homotopie. Pour cela nous utilisons l'homotopie

\begin{equation}
\label{cohomology-equation-second-homotopy}
h(s)_{i_0 \ldots i_p} =
\left\{
\begin{matrix}
0 & \text{si} & i_0 < i_1 < \ldots < i_p \\
(-1)^a s_{i_0 \ldots i_{a - 1} i_a i_a i_{a + 1} \ldots i_p}
& \text{si} & i_0 < i_1 < \ldots < i_{a - 1} < i_a = i_{a + 1}
\end{matrix}
\right.
\end{equation}

Nous affirmons que :
$$
(dh + hd)(s)_{i_0 \ldots i_p} =
\left\{
\begin{matrix}
0 & \text{si} & i_0 < i_1 < \ldots < i_p \\
s_{i_0 \ldots i_p}
& \text{sinon} &
\end{matrix}
\right.
$$
Nous omettons la vérification. (Il y a un script PARI/gp appelé second-homotopy.gp dans les scripts de sous-répertoire stacks-project qui peuvent être utilisés pour vérifier de nombreuses instances de cette revendication. Nous avons écrit ce script pour nous assurer que les signes sont corrects.) La revendication montre clairement que la composition
$$
\check{\mathcal{C}}_{semi\text{-}ord}^\bullet(\mathcal{U}, \mathcal{F})
\longrightarrow
\check{\mathcal{C}}_{ord}^\bullet(\mathcal{U}, \mathcal{F})
\longrightarrow
\check{\mathcal{C}}_{semi\text{-}ord}^\bullet(\mathcal{U}, \mathcal{F})
$$
de la projection avec l'inclusion naturelle est homotopique à l'application d'identité comme souhaité.

\end{proof}

\begin{lemma}

\label{cohomology-lemma-alternating-cech-trivial}

Soit $X$ un espace topologique, soit $\mathcal{F}$ un préfaisceau abélien sur $X$ et soit $\mathcal{U} : U = \bigcup_{i \in I} U_i$ un recouvrement ouvert. Si
$U_i = U$ pour un certain $i \in I$, alors le complexe de {\v C}ech alterné augmenté
$$
\mathcal{F}(U) \to \check{\mathcal{C}}_{alt}^\bullet(\mathcal{U}, \mathcal{F})
$$
obtenu en plaçant $\mathcal{F}(U)$ en degré $-1$, avec pour différentielle l'application canonique de $\mathcal{F}(U)$ dans
$\check{\mathcal{C}}^0(\mathcal{U}, \mathcal{F})$
est homotopiquement équivalent à $0$. De même, pour tout ordre total sur $I$,
le complexe de {\v C}ech ordonné augmenté
$$
\mathcal{F}(U) \to
\check{\mathcal{C}}_{ord}^\bullet(\mathcal{U}, \mathcal{F})
$$
est homotopiquement équivalent à $0$.

\end{lemma}

\begin{proof}
[Première preuve] Combiner les lemmes \ref{cohomology-lemma-cech-trivial} et \ref{cohomology-lemma-alternating-usual}.
\end{proof}

\begin{proof}
[Deuxième preuve] Puisque les complexes de {\v C}ech alterné et ordonné sont isomorphes, il suffit de traiter le complexe ordonné. Nous emploierons la notation usuelle : une cochaîne $s$ de degré $p$
du complexe de {\v C}ech ordonné augmenté est de la forme
$s = (s_{i_0 \ldots i_p})$ où $s_{i_0 \ldots i_p}$ est dans
$\mathcal{F}(U_{i_0 \ldots i_p})$ et $i_0 < \ldots < i_p$. Avec cette notation nous avons
$$
d(x)_{i_0 \ldots i_{p + 1}} =
\sum\nolimits_j (-1)^j x_{i_0 \ldots \hat i_j \ldots i_p}
$$
Fixons un indice $i \in I$ tel que $U = U_i$. Nous utilisons comme homotopie les applications
$$
h : \text{cochaînes de degré }p + 1 \to \text{cochaînes de degré }p
$$
données par la règle
$$
h(s)_{i_0 \ldots i_p} = 0 \text{ si } i \in \{i_0, \ldots, i_p\}
\text{ et }
h(s)_{i_0 \ldots i_p} = 
(-1)^j s_{i_0 \ldots i_j i i_{j + 1} \ldots i_p} \text{ sinon}
$$
Ici $j$ est l'unique indice tel que $i_j < i < i_{j + 1}$ dans le second cas; puisque $U = U_i$, nous avons aussi l'égalité
$$
\mathcal{F}(U_{i_0 \ldots i_p}) =
\mathcal{F}(U_{i_0 \ldots i_j i i_{j + 1} \ldots i_p})
$$
qui permet de considérer
$(-1)^j s_{i_0 \ldots i_j i i_{j + 1} \ldots i_p}$
comme un élément de $\mathcal{F}(U_{i_0 \ldots i_p})$. Nous allons montrer par un calcul que $d h + h d$ est l'opposé de l'identité, ce qui achèvera la preuve. Fixons une cochaîne $s$ de degré $p$ et des éléments
$i_0 < \ldots < i_p$ de $I$.

\medskip\noindent
Cas I : $i \in \{i_0, \ldots, i_p\}$. Écrivons $i = i_t$. Alors $h(d(s))_{i_0 \ldots i_p} = 0$. D'autre part, on a $$
d(h(s))_{i_0 \ldots i_p} =
\sum (-1)^j h(s)_{i_0 \ldots \hat i_j \ldots i_p} =
(-1)^t h(s)_{i_0 \ldots \hat i \ldots i_p} =
(-1)^t (-1)^{t - 1} s_{i_0 \ldots i_p}
$$ Ainsi $(dh + hd)(s)_{i_0 \ldots i_p} = -s_{i_0 \ldots i_p}$ comme désiré.

\medskip\noindent

Cas II : $i \not \in \{i_0, \ldots, i_p\}$. Soit $j$ tel que
$i_j < i < i_{j + 1}$. Alors nous voyons que

\begin{align*}
h(d(s))_{i_0 \ldots i_p}
& =
(-1)^j d(s)_{i_0 \ldots i_j i i_{j + 1} \ldots i_p} \\
& =
\sum\nolimits_{j' \leq j} (-1)^{j + j'}
s_{i_0 \ldots \hat i_{j'} \ldots i_j i i_{j + 1} \ldots i_p} -
s_{i_0 \ldots i_p} \\
&
+ \sum\nolimits_{j' > j} (-1)^{j + j' + 1}
s_{i_0 \ldots i_j i i_{j + 1} \ldots \hat i_{j'} \ldots i_p}
\end{align*}

D'autre part, nous avons

\begin{align*}
d(h(s))_{i_0 \ldots i_p}
& =
\sum\nolimits_{j'} (-1)^{j'} h(s)_{i_0 \ldots \hat i_{j'} \ldots i_p} \\
& =
\sum\nolimits_{j' \leq j} (-1)^{j' + j - 1}
s_{i_0 \ldots \hat i_{j'} \ldots i_j i i_{j + 1} \ldots i_p} \\
& +
\sum\nolimits_{j' > j} (-1)^{j' + j}
s_{i_0 \ldots i_j i i_{j + 1} \ldots \hat i_{j'} \ldots i_p}
\end{align*}

En les ajoutant, nous obtenons
$(dh + hd)(s)_{i_0 \ldots i_p} = - s_{i_0 \ldots i_p}$
comme désiré.

\end{proof}





\section{Un autre point de vue sur le complexe de {\v C}
ech}
\label{cohomology-section-locally-finite-cech}

\noindent
Dans cette section, nous discutons d'une autre façon d'établir la relation entre le complexe {\v C}ech et la cohomologie.

\begin{lemma}

\label{cohomology-lemma-covering-resolution}

Soit $X$ un espace annelé. Soit $\mathcal{U} : X = \bigcup_{i \in I} U_i$
un recouvrement ouvert de $X$. Soit $\mathcal{F}$ un $\mathcal{O}_X$-module. Notons $\mathcal{F}_{i_0 \ldots i_p}$ la restriction de
$\mathcal{F}$ à $U_{i_0 \ldots i_p}$. Il existe un complexe
${\mathfrak C}^\bullet(\mathcal{U}, \mathcal{F})$
de $\mathcal{O}_X$-modules avec
$$
{\mathfrak C}^p(\mathcal{U}, \mathcal{F}) =
\prod\nolimits_{i_0 \ldots i_p}
(j_{i_0 \ldots i_p})_* \mathcal{F}_{i_0 \ldots i_p}
$$
et différentiel
$d : {\mathfrak C}^p(\mathcal{U}, \mathcal{F})
\to {\mathfrak C}^{p + 1}(\mathcal{U}, \mathcal{F})$
comme dans l'équation (\ref{cohomology-equation-d-cech}). En outre, il existe une application canonique
$$
\mathcal{F} \to {\mathfrak C}^\bullet(\mathcal{U}, \mathcal{F})
$$
qui est un quasi-isomorphisme, c'est-à-dire,
${\mathfrak C}^\bullet(\mathcal{U}, \mathcal{F})$
est une résolution de $\mathcal{F}$.

\end{lemma}

\begin{proof}

Nous vérifions
$$
0 \to \mathcal{F} \to \mathfrak{C}^0(\mathcal{U}, \mathcal{F}) \to
\mathfrak{C}^1(\mathcal{U}, \mathcal{F}) \to  \ldots
$$
est exact sur les germes. $x \in X$ et choisir $i_{\text{fix}} \in I$
de telle manière que $x \in U_{i_{\text{fix}}}$. Ensuite, définir 
$$
h : \mathfrak{C}^p(\mathcal{U}, \mathcal{F})_x
\to \mathfrak{C}^{p - 1}(\mathcal{U}, \mathcal{F})_x
$$
comme suit: $s \in \mathfrak{C}^p(\mathcal{U}, \mathcal{F})_x$, prendre un représentant
$$
\widetilde{s} \in
\mathfrak{C}^p(\mathcal{U}, \mathcal{F})(V) =
\prod\nolimits_{i_0 \ldots i_p}
\mathcal{F}(V \cap U_{i_0} \cap \ldots \cap U_{i_p})
$$
défini sur un voisinage $V$ de $x$, et posons
$$
h(s)_{i_0 \ldots i_{p - 1}} =
\widetilde{s}_{i_{\text{fix}} i_0 \ldots i_{p - 1}, x}.
$$
Par la même formule (pour $p = 0$) nous obtenons une application
$\mathfrak{C}^{0}(\mathcal{U},\mathcal{F})_x \to \mathcal{F}_x$. Nous calculons formellement comme suit :

\begin{align*}
(dh + hd)(s)_{i_0 \ldots i_p}
& =
\sum\nolimits_{j = 0}^p
(-1)^j
h(s)_{i_0 \ldots \hat i_j \ldots i_p}
+
d(s)_{i_{\text{fix}} i_0 \ldots i_p}\\
& =
\sum\nolimits_{j = 0}^p
(-1)^j
s_{i_{\text{fix}} i_0 \ldots \hat i_j \ldots i_p}
+
s_{i_0 \ldots i_p}
+
\sum\nolimits_{j = 0}^p
(-1)^{j + 1}
s_{i_{\text{fix}} i_0 \ldots \hat i_j \ldots i_p} \\
& =
s_{i_0 \ldots i_p}
\end{align*}

Ceci montre $h$ est une homotopie de l'application d'identité du complexe étendu
$$
0 \to \mathcal{F}_x \to \mathfrak{C}^0(\mathcal{U}, \mathcal{F})_x
\to \mathfrak{C}^1(\mathcal{U}, \mathcal{F})_x \to \ldots
$$
à zéro et nous concluons.

\end{proof}

\noindent

Ce lemme permet de redémontrer facilement la suite spectrale reliant la cohomologie de {\v C}ech à la cohomologie du lemme \ref{cohomology-lemma-cech-spectral-sequence}. Soient $X$, $\mathcal{U}$ et $\mathcal{F}$ comme dans le lemme \ref{cohomology-lemma-covering-resolution},
et soit $\mathcal{F} \to \mathcal{I}^\bullet$ une résolution injective. Considérons alors le bicomplexe
$$
A^{\bullet, \bullet} =
\Gamma(X, {\mathfrak C}^\bullet(\mathcal{U}, \mathcal{I}^\bullet)).
$$
Par construction nous avons
$$
A^{p, q} = \prod\nolimits_{i_0 \ldots i_p} \mathcal{I}^q(U_{i_0 \ldots i_p})
$$
Considérons les deux suites spectrales associées à ce bicomplexe dans Homologie, section \ref{homology-section-double-complex}; voir notamment Homologie, lemme \ref{homology-lemma-ss-double-complex}. Pour la suite spectrale $({}'E_r, {}'d_r)_{r \geq 0}$, nous obtenons
${}'E_2^{p, q} = \check{H}^p(\mathcal{U}, \underline{H}^q(\mathcal{F}))$
car la formation des produits est exacte (Homologie, lemme \ref{homology-lemma-product-abelian-groups-exact}). Pour la suite spectrale $({}''E_r, {}''d_r)_{r \geq 0}$, nous obtenons
${}''E_2^{p, q} = 0$ si $p > 0$ et ${}''E_2^{0, q} = H^q(X, \mathcal{F})$. En effet, pour $q$ fixé, le complexe de faisceaux
${\mathfrak C}^\bullet(\mathcal{U}, \mathcal{I}^q)$
est une résolution (lemme \ref{cohomology-lemma-covering-resolution}) du faisceau injectif $\mathcal{I}^q$
par des faisceaux injectifs (d'après les lemmes \ref{cohomology-lemma-cohomology-of-open} et
\ref{cohomology-lemma-pushforward-injective-flat}
et Homologie, lemme \ref{homology-lemma-product-injectives}). Par conséquent, la cohomologie de
$\Gamma(X, {\mathfrak C}^\bullet(\mathcal{U}, \mathcal{I}^q))$
est nulle en degrés positifs et égale à $\Gamma(X, \mathcal{I}^q)$
en degré $0$. En prenant la cohomologie pour la différentielle suivante, on obtient l'assertion sur la suite spectrale $({}''E_r, {}''d_r)_{r \geq 0}$. Le résultat s'ensuit, puisque les deux suites spectrales convergent vers la cohomologie du complexe total associé à $A^{\bullet, \bullet}$.

\begin{definition}

\label{cohomology-definition-covering-locally-finite}

Soit $X$ un espace topologique. Un recouvrement ouvert $X = \bigcup_{i \in I} U_i$ est dit
{\it localement fini} si, pour tout $x \in X$, il existe un voisinage ouvert
$W$ de $x$ tel que $\{i \in I \mid W \cap U_i \not = \emptyset\}$ soit fini.

\end{definition}

\begin{remark}

\label{cohomology-remark-locally-finite-sections}

Soit $X = \bigcup_{i \in I} U_i$ un recouvrement ouvert localement fini. Notons $j_i : U_i \to X$ l'application d'inclusion. Supposons que, pour chaque $i$,
on se donne un faisceau abélien $\mathcal{F}_i$ sur $U_i$. Considérons le faisceau abélien $\mathcal{G} = \bigoplus_{i \in I} (j_i)_*\mathcal{F}_i$. Alors, pour tout ouvert $V \subset X$, on a
$$
\Gamma(V, \mathcal{G}) = \prod\nolimits_{i \in I} \mathcal{F}_i(V \cap U_i).
$$
En d'autres termes, nous avons
$$
\bigoplus\nolimits_{i \in I} (j_i)_*\mathcal{F}_i =
\prod\nolimits_{i \in I} (j_i)_*\mathcal{F}_i
$$
Cela peut sembler étrange jusqu'à ce que l'on se rappelle que la somme directe d'une famille de faisceaux est le faisceau associé au préfaisceau somme directe. Voir la discussion dans Modules, section \ref{modules-section-kernels}. Nous concluons donc que, dans ce cas, les termes du complexe du lemme \ref{cohomology-lemma-covering-resolution} sont
$$
{\mathfrak C}^p(\mathcal{U}, \mathcal{F}) =
\bigoplus\nolimits_{i_0 \ldots i_p}
(j_{i_0 \ldots i_p})_* \mathcal{F}_{i_0 \ldots i_p}
$$
Ce qui est parfois utile.

\end{remark}











\section{Cohomologie de {\v C}
ech des complexes}
\label{cohomology-section-cech-cohomology-of-complexes}

\noindent
En général, pour des faisceaux de groupes abéliens ${\mathcal F}$ et ${\mathcal G}$ sur $X$, il existe une application produit cup $$
H^i(X, {\mathcal F}) \times H^j(X, {\mathcal G})
\longrightarrow
H^{i + j}(X, {\mathcal F} \otimes_{\mathbf Z} {\mathcal G}).
$$ Dans cette section, nous la définissons au moyen de cocycles de {\v C}ech, par une formule explicite pour le produit cup. Si l'on craint que la cohomologie ne coïncide pas avec la cohomologie de {\v C}ech, on peut utiliser des hyperrecouvrements tout en conservant la notation des cocycles. Cette méthode a en outre l'avantage de définir le produit cup en hypercohomologie sur tout topos (insérer ici une référence ultérieure).

\medskip\noindent

Soit ${\mathcal F}^\bullet$ un complexe borné inférieurement de préfaisceaux de groupes abéliens sur $X$. On peut souvent calculer $H^n(X, {\mathcal F}^\bullet)$
au moyen de cocycles de {\v C}ech. Soit donc
${\mathcal U} : X = \bigcup_{i \in I} U_i$
un recouvrement ouvert de $X$. Comme le complexe de {\v C}ech
$\check{\mathcal{C}}^\bullet(\mathcal{U}, \mathcal{F})$
(définition \ref{cohomology-definition-cech-complex}) est fonctoriel en le préfaisceau $\mathcal{F}$, on obtient un bicomplexe
$\check{\mathcal{C}}^\bullet(\mathcal{U}, \mathcal{F}^\bullet)$. Le complexe associé total à
$\check{\mathcal{C}}^\bullet({\mathcal U}, {\mathcal F}^\bullet)$
est le complexe dont le terme de degré $n$ est
$$
\text{Tot}^n(\check{\mathcal{C}}^\bullet({\mathcal U}, {\mathcal F}^\bullet))
=
\bigoplus\nolimits_{p + q = n}
\prod\nolimits_{i_0\ldots i_p} {\mathcal F}^q(U_{i_0\ldots i_p})
$$
Voir Homologie, définition \ref{homology-definition-associated-simple-complex}. Un élément typique de $\text{Tot}^n$ sera noté
$\alpha = \{\alpha_{i_0\ldots i_p}\}$ où
$\alpha_{i_0 \ldots i_p} \in \mathcal{F}^q(U_{i_0\ldots i_p})$. Autrement dit, le degré en $\mathcal{F}$ de $\alpha_{i_0\ldots i_p}$ est
$q = n - p$. Cette notation suppose que l'on garde à l'esprit le degré de $\alpha$;
nous indiquons cette situation par la formule
$\deg_{\mathcal F}(\alpha_{i_0\ldots i_p}) = q$. Selon nos conventions dans Homologie, définition \ref{homology-definition-associated-simple-complex}
la différentielle d'un élément $\alpha$ de degré $n$ est donnée par
$$
d(\alpha)_{i_0\ldots i_{p + 1}}
=
\sum\nolimits_{j = 0}^{p + 1}
(-1)^j \alpha_{i_0 \ldots \hat i_j \ldots i_{p + 1}} + 
(-1)^{p + 1}d_{{\mathcal F}}(\alpha_{i_0 \ldots i_{p + 1}})
$$
où $d_\mathcal{F}$ désigne le différentiel sur le complexe
$\mathcal{F}^\bullet$. L'expression $\alpha_{i_0 \ldots \hat i_j \ldots i_{p + 1}}$ désigne la restriction de $\alpha_{i_0 \ldots \hat i_j \ldots i_{p + 1}}
\in {\mathcal F}(U_{i_0\ldots\hat i_j\ldots i_{p + 1}})$ à
$U_{i_0 \ldots i_{p + 1}}$.

\medskip\noindent

La construction de
$\text{Tot}(\check{\mathcal{C}}^\bullet({\mathcal U}, {\mathcal F}^\bullet))$
est fonctorielle en ${\mathcal F}^\bullet$. Il existe également une transformation fonctorielle

\begin{equation}
\label{cohomology-equation-global-sections-to-cech}
\Gamma(X, {\mathcal F}^\bullet)
\longrightarrow
\text{Tot}(\check{\mathcal{C}}^\bullet({\mathcal U}, {\mathcal F}^\bullet))
\end{equation}

de complexes définie par la règle suivante :
$s\in \Gamma(X, {\mathcal F}^n)$
est envoyé sur l'élément $\alpha = \{\alpha_{i_0\ldots i_p}\}$
avec $\alpha_{i_0} = s|_{U_{i_0}}$ et $\alpha_{i_0\ldots i_p} = 0$
pour $p > 0$.

\medskip\noindent

Raffinements. Soit ${\mathcal V} = \{ V_j \}_{j\in J}$ un raffinement de ${\mathcal U}$. Cela signifie qu'il existe une application $t: J \to I$
telle que $V_j \subset U_{t(j)}$ pour tout $j\in J$. Elle induit une transformation fonctorielle

\begin{equation}
\label{cohomology-equation-transformation}
T_t :
\text{Tot}(\check{\mathcal{C}}^\bullet({\mathcal U}, {\mathcal F}^\bullet))
\longrightarrow
\text{Tot}(\check{\mathcal{C}}^\bullet({\mathcal V}, {\mathcal F}^\bullet)).
\end{equation}

défini par la règle
$$
T_t(\alpha)_{j_0\ldots j_p}
=
\alpha_{t(j_0)\ldots t(j_p)}|_{V_{j_0\ldots j_p}}.
$$
Pour deux applications $t, t' : J \to I$ comme ci-dessus, les applications
$T_t$ et $T_{t'}$ ainsi construites sont homotopes. Une homotopie est donnée par
$$
h(\alpha)_{j_0\ldots j_p}
=
\sum\nolimits_{a = 0}^{p}
(-1)^a
\alpha_{t(j_0)\ldots t(j_a) t'(j_a) \ldots t'(j_p)}
$$
pour un élément $\alpha$ de degré $n$. Cela résulte du calcul suivant, où
$\alpha$ est encore un élément de degré $n$ (ainsi $d(\alpha)$
est de degré $n + 1$ et $h(\alpha)$ de degré $n - 1$) :

\begin{align*}
(
d(h(\alpha)) + h(d(\alpha))
)_{j_0\ldots j_p}
= &
\sum\nolimits_{k = 0}^p
(-1)^k
h(\alpha)_{j_0 \ldots \hat j_k \ldots j_p}
+ \\
&
(-1)^p
d_{\mathcal F}(h(\alpha)_{j_0 \ldots j_p})
+ \\
&
\sum\nolimits_{a = 0}^p
(-1)^a
d(\alpha)_{t(j_0) \ldots t(j_a) t'(j_a) \ldots t'(j_p)}
\\
= &
\sum\nolimits_{k = 0}^p
\sum\nolimits_{a = 0}^{k - 1}
(-1)^{k + a}
\alpha_{t(j_0)\ldots t(j_a)t'(j_a)\ldots \hat{t'(j_k)}\ldots t'(j_p)}
+ \\
&
\sum\nolimits_{k = 0}^p
\sum\nolimits_{a = k + 1}^p
(-1)^{k + a - 1}
\alpha_{t(j_0)\ldots \hat{t(j_k)}\ldots t(j_a)t'(j_a)\ldots t'(j_p)}
+ \\
&
\sum\nolimits_{a = 0}^p
(-1)^{p + a}
d_{\mathcal F}(\alpha_{t(j_0)\ldots t(j_a) t'(j_a) \ldots t'(j_p)})
+ \\
&
\sum\nolimits_{a = 0}^p
\sum\nolimits_{k = 0}^a
(-1)^{a + k}
\alpha_{t(j_0)\ldots\hat{t(j_k)}\ldots t(j_a)t'(j_a)\ldots t'(j_p)}
+ \\
&
\sum\nolimits_{a = 0}^p
\sum\nolimits_{k = a}^p
(-1)^{a + k + 1}
\alpha_{t(j_0) \ldots t(j_a) t'(j_a) \ldots \hat{t'(j_k)} \ldots t'(j_p)}
+ \\
&
\sum\nolimits_{a = 0}^p
(-1)^{a + p + 1}
d_{\mathcal F}(\alpha_{t(j_0)\ldots t(j_a) t'(j_a) \ldots t'(j_p)})
\\
= &
\alpha_{t'(j_0)\ldots t'(j_p)} +
(-1)^{2p + 1}\alpha_{t(j_0)\ldots t(j_p)}
\\
= &
T_{t'}(\alpha)_{j_0\ldots j_p} - T_t(\alpha)_{j_0\ldots j_p}
\end{align*}

Nous laissons le lecteur vérifier les annulations. (Notez que les termes ayant les deux $k$ et $a$ dans les 1er, 2e et 4e, 5e summands annuler, sauf ceux où $a = k$ qui ne se produisent que dans les 4ème et 5ème et qui s'annulent l'un contre l'autre, sauf pour les deux termes souhaités.) Il s'ensuit que l'application induite
$$
H^n(T_t) :
H^n(
\text{Tot}(\check{\mathcal{C}}^\bullet({\mathcal U}, {\mathcal F}^\bullet))
)
\to
H^n(
\text{Tot}(\check{\mathcal{C}}^\bullet({\mathcal V}, {\mathcal F}^\bullet))
)
$$
est indépendant du choix de $t$. Nous définissons
{\it hypercohomologie de {\v C}ech} comme la colimite des groupes de cohomologie de {\v C}ech sur tous les raffinements, relativement aux applications $H^\bullet(T_t)$.

\medskip\noindent
En passant à la colimite sur tous les recouvrements ouverts de $X$, le lemme suivant fournit une application de l'hypercohomologie de {\v C}ech vers la cohomologie. Celle-ci est souvent un isomorphisme et l'est toujours si l'on utilise des hyperrecouvrements.

\begin{lemma}

\label{cohomology-lemma-cech-complex-complex}
Soit $(X, \mathcal{O}_X)$ un espace annelé, soit $\mathcal{U} : X = \bigcup_{i \in I} U_i$ un recouvrement ouvert. Pour tout complexe borné inférieurement $\mathcal{F}^\bullet$ de $\mathcal{O}_X$-modules, il existe une application canonique $$
\text{Tot}(\check{\mathcal{C}}^\bullet(\mathcal{U}, \mathcal{F}^\bullet))
\longrightarrow
R\Gamma(X, \mathcal{F}^\bullet)
$$ fonctorielle en $\mathcal{F}^\bullet$ et compatible avec (\ref{cohomology-equation-global-sections-to-cech}) et (\ref{cohomology-equation-transformation}). Il existe une suite spectrale $(E_r, d_r)_{r \geq 0}$ telle que $$
E_2^{p, q} =
H^p(\text{Tot}(\check{\mathcal{C}}^\bullet(\mathcal{U},
\underline{H}^q(\mathcal{F}^\bullet)))
$$ convergeant vers $H^{p + q}(X, \mathcal{F}^\bullet)$.
\end{lemma}

\begin{proof}

Soit ${\mathcal I}^\bullet$ un complexe borné inférieurement de faisceaux injectifs. L'application (\ref{cohomology-equation-global-sections-to-cech}) pour
$\mathcal{I}^\bullet$ est une application
$\Gamma(X, {\mathcal I}^\bullet) \to
\text{Tot}(\check{\mathcal{C}}^\bullet({\mathcal U}, {\mathcal I}^\bullet))$. C'est un quasi-isomorphisme de complexes de groupes abéliens, comme il résulte d'Homologie, lemme \ref{homology-lemma-double-complex-gives-resolution},
appliqué au complexe double
$\check{\mathcal{C}}^\bullet({\mathcal U}, {\mathcal I}^\bullet)$, en utilisant le lemme \ref{cohomology-lemma-injective-trivial-cech}. Choisissons un quasi-isomorphisme ${\mathcal F}^\bullet \to {\mathcal I}^\bullet$ de ${\mathcal F}^\bullet$ vers un complexe borné inférieurement de faisceaux injectifs. Le complexe représentant $R\Gamma(X, {\mathcal F}^\bullet)$ est
$\Gamma(X, {\mathcal I}^\bullet)$; on obtient l'application du lemme au moyen de
$$
\text{Tot}(\check{\mathcal{C}}^\bullet({\mathcal U}, {\mathcal F}^\bullet))
\longrightarrow
\text{Tot}(\check{\mathcal{C}}^\bullet({\mathcal U}, {\mathcal I}^\bullet)).
$$
Nous omettons la vérification de la fonctorialité et des compatibilités. Pour construire la suite spectrale du lemme, choisissons une résolution de Cartan–Eilenberg $\mathcal{F}^\bullet \to \mathcal{I}^{\bullet, \bullet}$; voir Catégories dérivées, lemme \ref{derived-lemma-cartan-eilenberg}. Alors $\mathcal{F}^\bullet \to \text{Tot}(\mathcal{I}^{\bullet, \bullet})$
est une résolution injective, et donc
$$
\text{Tot}(\check{\mathcal{C}}^\bullet({\mathcal U},
\text{Tot}({\mathcal I}^{\bullet, \bullet})))
$$
calcule $R\Gamma(X, \mathcal{F}^\bullet)$, comme nous l'avons vu plus haut. D'après Homologie, remarque \ref{homology-remark-triple-complex},
nous pouvons le considérer comme le complexe total associé au triple complexe
$\check{\mathcal{C}}^\bullet({\mathcal U}, {\mathcal I}^{\bullet, \bullet})$
et, en utilisant la même remarque, comme le complexe total associé au bicomplexe $A^{\bullet, \bullet}$ dont les termes sont
$$
A^{n, m} =
\bigoplus\nolimits_{p + q = n}
\check{\mathcal{C}}^p({\mathcal U}, \mathcal{I}^{q, m})
$$
Comme $\mathcal{I}^{q, \bullet}$ est une résolution injective de
$\mathcal{F}^q$, nous pouvons appliquer la première suite spectrale associée à
$A^{\bullet, \bullet}$
(Homologie, lemme \ref{homology-lemma-ss-double-complex}) pour obtenir une suite spectrale
$$
E_1^{n, m} =
\bigoplus\nolimits_{p + q = n}
\check{\mathcal{C}}^p(\mathcal{U}, \underline{H}^m(\mathcal{F}^q))
$$
qui est le terme de degré $n$ du complexe
$\text{Tot}(\check{\mathcal{C}}^\bullet(\mathcal{U},
\underline{H}^m(\mathcal{F}^\bullet))$. On obtient donc
les termes $E_2$ décrits dans le lemme. La convergence résulte d'Homologie, lemme \ref{homology-lemma-first-quadrant-ss}.

\end{proof}

\begin{lemma}

\label{cohomology-lemma-cech-complex-complex-computes}
Soit $(X, \mathcal{O}_X)$ un espace annelé. Soit $\mathcal{U} : X = \bigcup_{i \in I} U_i$ un recouvrement ouvert. Soit $\mathcal{F}^\bullet$ un complexe borné inférieurement de $\mathcal{O}_X$-modules. Si $H^i(U_{i_0 \ldots i_p}, \mathcal{F}^q) = 0$ pour tout $i > 0$ et tous $p, i_0, \ldots, i_p, q$, alors l'application $
\text{Tot}(\check{\mathcal{C}}^\bullet(\mathcal{U}, \mathcal{F}^\bullet))
\to
R\Gamma(X, \mathcal{F}^\bullet)
$ du lemme \ref{cohomology-lemma-cech-complex-complex} est un isomorphisme.
\end{lemma}

\begin{proof}
Cela résulte immédiatement de la suite spectrale du lemme \ref{cohomology-lemma-cech-complex-complex}.
\end{proof}

\begin{remark}

\label{cohomology-remark-shift-complex-cech-complex}

Soit $(X, \mathcal{O}_X)$ un espace annelé.
Soit $\mathcal{U} : X = \bigcup_{i \in I} U_i$ un recouvrement ouvert, soit $\mathcal{F}^\bullet$ un complexe borné inférieurement de $\mathcal{O}_X$-modules et soit $b$ un entier. Nous affirmons qu'il existe un diagramme commutatif
$$
\xymatrix{
\text{Tot}(\check{\mathcal{C}}^\bullet(\mathcal{U}, \mathcal{F}^\bullet))[b]
\ar[r] \ar[d]_\gamma &
R\Gamma(X, \mathcal{F}^\bullet)[b] \ar[d] \\
\text{Tot}(\check{\mathcal{C}}^\bullet(\mathcal{U}, \mathcal{F}^\bullet[b]))
\ar[r] &
R\Gamma(X, \mathcal{F}^\bullet[b])
}
$$
dans la catégorie dérivée, où $\gamma$ est l'application de complexes construite dans Homologie, remarque \ref{homology-remark-shift-double-complex}. Cela a un sens puisque le bicomplexe
$\check{\mathcal{C}}^\bullet(\mathcal{U}, \mathcal{F}^\bullet[b])$
est manifestement le même que le bicomplexe
$\check{\mathcal{C}}^\bullet(\mathcal{U}, \mathcal{F}^\bullet)[0, b]$
introduit dans Homologie, remarque \ref{homology-remark-shift-double-complex}. Pour vérifier que le diagramme commute, choisissons une résolution injective
$\mathcal{F}^\bullet \to \mathcal{I}^\bullet$ comme dans la preuve du lemme \ref{cohomology-lemma-cech-complex-complex}. Une chasse au diagramme montre qu'il suffit de vérifier la commutativité après avoir remplacé $\mathcal{F}^\bullet$
par $\mathcal{I}^\bullet$. Ensuite, nous considérons le diagramme étendu
$$
\xymatrix{
\Gamma(X, \mathcal{I}^\bullet)[b] \ar[r] \ar[d] &
\text{Tot}(\check{\mathcal{C}}^\bullet(\mathcal{U}, \mathcal{I}^\bullet))[b]
\ar[r] \ar[d]_\gamma &
R\Gamma(X, \mathcal{I}^\bullet)[b] \ar[d] \\
\Gamma(X, \mathcal{I}^\bullet[b]) \ar[r] &
\text{Tot}(\check{\mathcal{C}}^\bullet(\mathcal{U}, \mathcal{I}^\bullet[b]))
\ar[r] &
R\Gamma(X, \mathcal{I}^\bullet[b])
}
$$
où les flèches horizontales de gauche sont celles de (\ref{cohomology-equation-global-sections-to-cech}). Puisque les flèches horizontales sont ici des isomorphismes dans la catégorie dérivée (voir la preuve du lemme \ref{cohomology-lemma-cech-complex-complex}), il suffit de montrer que le carré de gauche commute. C'est le cas parce que l'application $\gamma$ utilise le signe $1$ sur les facteurs directs
$\check{\mathcal{C}}^0(\mathcal{U}, \mathcal{I}^{q + b})$, voir la formule dans Homologie, remarque \ref{homology-remark-shift-double-complex}.

\end{remark}

\noindent

Soit $X$ un espace topologique, soit $\mathcal{U} : X = \bigcup_{i \in I} U_i$
un recouvrement ouvert et soit $\mathcal{F}^\bullet$ un complexe borné inférieurement de préfaisceaux de groupes abéliens.
$\tau :
\text{Tot}(\check{\mathcal{C}}^\bullet({\mathcal U}, {\mathcal F}^\bullet))
\to
\text{Tot}(\check{\mathcal{C}}^\bullet({\mathcal U}, {\mathcal F}^\bullet))$
défini par
$$
\tau(\alpha)_{i_0 \ldots i_p} = (-1)^{p(p + 1)/2} \alpha_{i_p \ldots i_0}.
$$
Alors, pour tout élément $\alpha$ de degré $n$, on a

\begin{align*}
& d(\tau(\alpha))_{i_0 \ldots i_{p + 1}} \\
& =
\sum\nolimits_{j = 0}^{p + 1}
(-1)^j
\tau(\alpha)_{i_0 \ldots \hat i_j \ldots i_{p + 1}}
+
(-1)^{p + 1}
d_{\mathcal F}(\tau(\alpha)_{i_0 \ldots i_{p + 1}})
\\
& =
\sum\nolimits_{j = 0}^{p + 1}
(-1)^{j + \frac{p(p + 1)}{2}}
\alpha_{i_{p + 1} \ldots \hat i_j \ldots i_0}
+
(-1)^{p + 1 + \frac{(p + 1)(p + 2)}{2}}
d_{\mathcal F}(\alpha_{i_{p + 1} \ldots i_0})
\end{align*}

D'autre part, nous avons

\begin{align*}
& \tau(d(\alpha))_{i_0\ldots i_{p + 1}} \\
& =
(-1)^{\frac{(p + 1)(p + 2)}{2}} d(\alpha)_{i_{p + 1} \ldots i_0}
\\
& =
(-1)^{\frac{(p + 1)(p + 2)}{2}}
\left(
\sum\nolimits_{j = 0}^{p + 1}
(-1)^j
\alpha_{i_{p + 1}\ldots \hat i_{p + 1 - j} \ldots i_0}
+
(-1)^{p + 1}
d_{\mathcal F}(\alpha_{i_{p + 1}\ldots i_0})
\right)
\end{align*}

Ainsi, nous concluons que $d(\tau(\alpha)) = \tau(d(\alpha))$
parce que $p(p + 1)/2 \equiv (p + 1)(p + 2)/2 + p + 1 \bmod 2$. En d'autres termes
$\tau$ est un endomorphisme du complexe
$\text{Tot}(\check{\mathcal{C}}^\bullet({\mathcal U}, {\mathcal F}^\bullet))$. Notez que le schéma
$$
\begin{matrix}
\Gamma(X, {\mathcal F}^\bullet) &
\longrightarrow &
\text{Tot}(\check{\mathcal{C}}^\bullet({\mathcal U}, {\mathcal F}^\bullet)) \\
\downarrow \text{id} & & \downarrow \tau \\
\Gamma(X, {\mathcal F}^\bullet) &
\longrightarrow &
\text{Tot}(\check{\mathcal{C}}^\bullet({\mathcal U}, {\mathcal F}^\bullet))
\end{matrix}
$$
commute. En outre, $\tau$ est manifestement compatible avec les raffinements. Cela suggère que $\tau$ agit comme l'identité sur la cohomologie de {\v C}ech, c'est-à-dire dans la colimite, dès que l'hypercohomologie de {\v C}ech coïncide avec l'hypercohomologie; c'est toujours le cas si l'on utilise des hyperrecouvrements. En fait, $\tau$ est homotope à l'identité sur le complexe total de {\v C}ech
$\text{Tot}(\check{\mathcal{C}}^\bullet({\mathcal U}, {\mathcal F}^\bullet))$. Pour prouver cela, nous utilisons comme homotopie
$$
h(\alpha)_{i_0 \ldots i_p}
=
\sum\nolimits_{a = 0}^p
\epsilon_p(a)
\alpha_{i_0 \ldots i_a i_p \ldots i_a}
\quad\text{avec}\quad
\epsilon_p(a) = (-1)^{\frac{(p - a)(p - a - 1)}{2} + p}
$$
pour $\alpha$ du degré $n$. Comme d'habitude nous omettons d'écrire
$|_{U_{i_0 \ldots i_p}}$. Cela fonctionne à cause du calcul suivant, encore avec
$\alpha$ un élément de degré $n$:

\begin{align*}
(d(h(\alpha)) + h(d(\alpha)))_{i_0 \ldots i_p}
= &
\sum\nolimits_{k = 0}^p
(-1)^k
h(\alpha)_{i_0 \ldots \hat i_k \ldots i_p}
+ \\
&
(-1)^p
d_{\mathcal F}(h(\alpha)_{i_0 \ldots i_p})
+ \\
&
\sum\nolimits_{a = 0}^p
\epsilon_p(a)
d(\alpha)_{i_0 \ldots i_a i_p \ldots i_a}
\\
= &
\sum\nolimits_{k = 0}^p
\sum\nolimits_{a = 0}^{k - 1}
(-1)^k \epsilon_{p - 1}(a)
\alpha_{i_0 \ldots i_a i_p \ldots \hat{i_k} \ldots i_a}
+ \\
&
\sum\nolimits_{k = 0}^p
\sum\nolimits_{a = k + 1}^p
(-1)^k \epsilon_{p - 1}(a - 1)
\alpha_{i_0 \ldots \hat{i_k} \ldots i_a i_p \ldots i_a}
+ \\
&
\sum\nolimits_{a = 0}^p
(-1)^p \epsilon_p(a)
d_{\mathcal F}(\alpha_{i_0 \ldots i_a i_p \ldots i_a})
+ \\
&
\sum\nolimits_{a = 0}^p
\sum\nolimits_{k = 0}^a
\epsilon_p(a) (-1)^k
\alpha_{i_0 \ldots \hat{i_k} \ldots i_a i_p \ldots i_a}
+ \\
&
\sum\nolimits_{a = 0}^p
\sum\nolimits_{k = a}^p
\epsilon_p(a) (-1)^{p + a + 1 - k}
\alpha_{i_0 \ldots i_a i_p \ldots \hat{i_k} \ldots i_a}
+ \\
&
\sum\nolimits_{a = 0}^p
\epsilon_p(a) (-1)^{p + 1}
d_{\mathcal F}(\alpha_{i_0 \ldots i_a i_p \ldots i_a})
\\
= &
\epsilon_p(0) \alpha_{i_p \ldots i_0} +
\epsilon_p(p) (-1)^{p + 1} \alpha_{i_0 \ldots i_p} \\
= &
(-1)^{\frac{p(p + 1)}{2}}\alpha_{i_p \ldots i_0}
- \alpha_{i_0 \ldots i_p}
\end{align*}

Les annulations suivent parce que
$$
(-1)^k \epsilon_{p - 1}(a) + \epsilon_p(a)(-1)^{p + a + 1 - k} = 0
\quad\text{et}\quad
(-1)^k\epsilon_{p - 1}(a - 1) + \epsilon_p(a) (-1)^k = 0
$$
Nous laissons le lecteur vérifier les annulations.

\medskip\noindent

Supposons que nous ayons deux complexes bornés inférieurement de faisceaux abéliens
${\mathcal F}^\bullet$ et ${\mathcal G}^\bullet$. Nous définissons le complexe
$\text{Tot}({\mathcal F}^\bullet\otimes_{\mathbf Z} {\mathcal G}^\bullet)$
être complexe avec les termes
$\bigoplus_{p + q = n} {\mathcal F}^p \otimes {\mathcal G}^q$
et différentiel selon la règle

\begin{equation}
\label{cohomology-equation-differential-tensor-product-complexes}
d(\alpha \otimes \beta) =
d(\alpha)\otimes \beta + (-1)^{\deg(\alpha)} \alpha \otimes d(\beta)
\end{equation}

lorsque $\alpha$ et $\beta$ homogène, voir Homologie, définition \ref{homology-definition-associated-simple-complex}.

\medskip\noindent
Supposons que $M^\bullet$ et $N^\bullet$ soient deux complexes bornés inférieurement de groupes abéliens. Si $m$, respectivement $n$, est un cocycle de $M^\bullet$, respectivement de $N^\bullet$, il est immédiat que $m \otimes n$ est un cocycle de $\text{Tot}(M^\bullet\otimes N^\bullet)$. On obtient ainsi un produit cup $$
H^i(M^\bullet) \times H^j(N^\bullet)
\longrightarrow
H^{i + j}(Tot(M^\bullet\otimes N^\bullet)).
$$ Ceci est discuté également dans Compléments sur l'algèbre, section \ref{more-algebra-section-products-tor}.

\medskip\noindent

Ainsi, la construction du produit de la coupe dans l'hypercohomologie des complexes repose sur une construction d'une application des complexes

\begin{equation}
\label{cohomology-equation-needs-signs}
\text{Tot}\left(
\text{Tot}(\check{\mathcal{C}}^\bullet({\mathcal U}, {\mathcal F}^\bullet))
\otimes_{\mathbf Z}
\text{Tot}(\check{\mathcal{C}}^\bullet({\mathcal U}, {\mathcal G}^\bullet))
\right)
\longrightarrow
\text{Tot}(
\check{\mathcal{C}}^\bullet({\mathcal U},
\text{Tot}({\mathcal F}^\bullet\otimes {\mathcal G}^\bullet)
))
\end{equation}

Cette application est dénotée $\cup$ et est donné par la règle
$$
(\alpha \cup \beta)_{i_0 \ldots i_p}
=
\sum\nolimits_{r = 0}^p
\epsilon(n, m, p, r)
\alpha_{i_0 \ldots i_r} \otimes \beta_{i_r \ldots i_p}.
$$
où $\alpha$ est de degré $n$ et $\beta$ de degré $m$
et avec
$$
\epsilon(n, m, p, r) = (-1)^{(p + r)n + rp + r}.
$$
Notez que $\epsilon(n, m, p, n) = 1$. Par conséquent, si
$\mathcal{F}^\bullet = \mathcal{F}[0]$ est le complexe constitué d'un seul faisceau abélien $\mathcal{F}$ placé dans le degré $0$, puis il n'y a pas de signes dans la formule pour $\cup$ (comme dans ce cas
$\alpha_{i_0 \ldots i_r} = 0$ sauf $r = n$). Pour une explication de pourquoi il doit y avoir un signe et comment le calculer voir \cite[Exposé XVII]{SGA4} par Deligne. Pour vérifier (\ref{cohomology-equation-needs-signs}) est une application des complexes que nous devons montrer que
$$
d(\alpha \cup \beta) =
d(\alpha) \cup \beta +
(-1)^{\deg(\alpha)} \alpha \cup d(\beta)
$$
par la définition du différentiel
$\text{Tot}(
\text{Tot}(\check{\mathcal{C}}^\bullet({\mathcal U}, {\mathcal F}^\bullet))
\otimes_{\mathbf Z}
\text{Tot}(\check{\mathcal{C}}^\bullet({\mathcal U}, {\mathcal G}^\bullet))
)$
tel qu'indiqué dans Homologie, définition \ref{homology-definition-associated-simple-complex}. Nous calculons d'abord

\begin{align*}
d(\alpha \cup \beta)_{i_0 \ldots i_{p + 1}}
= &
\sum\nolimits_{j = 0}^{p + 1}
(-1)^j
(\alpha \cup \beta)_{i_0 \ldots \hat i_j \ldots i_{p + 1}}
+
(-1)^{p + 1}
d_{{\mathcal F} \otimes {\mathcal G}}
((\alpha \cup \beta)_{i_0 \ldots i_{p + 1}})
\\
= &
\sum\nolimits_{j = 0}^{p + 1}
\sum\nolimits_{r = 0}^{j - 1}
(-1)^j \epsilon(n, m, p, r)
\alpha_{i_0 \ldots i_r} \otimes \beta_{i_r \ldots \hat i_j \ldots i_{p + 1}}
+ \\
&
\sum\nolimits_{j = 0}^{p + 1}
\sum\nolimits_{r = j + 1}^{p + 1}
(-1)^j \epsilon(n, m, p, r - 1)
\alpha_{i_0 \ldots \hat i_j \ldots i_r} \otimes \beta_{i_r \ldots i_{p + 1}}
+ \\
&
\sum\nolimits_{r = 0}^{p + 1}
(-1)^{p + 1} \epsilon(n, m, p + 1, r)
d_{{\mathcal F} \otimes {\mathcal G}}
(\alpha_{i_0 \ldots i_r} \otimes \beta_{i_r \ldots i_{p + 1}})
\end{align*}

et noter que les sommes dans le dernier terme égale
$$
(-1)^{p + 1} \epsilon(n, m, p + 1, r)
\left(
d_{\mathcal F}(\alpha_{i_0 \ldots i_r}) \otimes 
\beta_{i_r \ldots i_{p + 1}} +
(-1)^{n - r}
\alpha_{i_0 \ldots i_r} \otimes d_{\mathcal G}(\beta_{i_r \ldots i_{p + 1}})
\right).
$$
parce que $\deg_\mathcal{F}(\alpha_{i_0 \ldots i_r}) = n - r$. D'un autre côté

\begin{align*}
(d(\alpha) \cup \beta)_{i_0\ldots i_{p + 1}}
= &
\sum\nolimits_{r = 0}^{p + 1}
\epsilon(n + 1, m, p + 1, r)
d(\alpha)_{i_0\ldots i_r} \otimes \beta_{i_r\ldots i_{p + 1}}
\\
= &
\sum\nolimits_{r = 0}^{p + 1}
\sum\nolimits_{j = 0}^{r}
\epsilon(n + 1, m, p + 1, r) (-1)^j
\alpha_{i_0\ldots\hat{i_j}\ldots i_r} \otimes \beta_{i_r\ldots i_{p + 1}}
+ \\
&
\sum\nolimits_{r = 0}^{p + 1}
\epsilon(n + 1, m, p + 1, r) (-1)^r
d_{\mathcal F}(\alpha_{i_0 \ldots i_r}) \otimes \beta_{i_r\ldots i_{p + 1}}
\end{align*}

et

\begin{align*}
(\alpha \cup d(\beta))_{i_0\ldots i_{p + 1}}
= &
\sum\nolimits_{r = 0}^{p + 1}
\epsilon(n, m + 1, p + 1, r)
\alpha_{i_0 \ldots i_r} \otimes d(\beta)_{i_r \ldots i_{p + 1}}
\\
= &
\sum\nolimits_{r = 0}^{p + 1}
\sum\nolimits_{j = r}^{p + 1}
\epsilon(n, m + 1, p + 1, r) (-1)^{j - r}
\alpha_{i_0 \ldots i_r} \otimes \beta_{i_r \ldots \hat{i_j}\ldots i_{p + 1}}
+ \\
&
\sum\nolimits_{r = 0}^{p + 1}
\epsilon(n, m + 1, p + 1, r) (-1)^{p + 1 - r}
\alpha_{i_0 \ldots i_r} \otimes d_{\mathcal G}(\beta_{i_r \ldots i_{p + 1}})
\end{align*}

L'égalité souhaitée tient si nous avons

\begin{align*}
(-1)^{p + 1} \epsilon(n, m, p + 1, r)
& =
\epsilon(n + 1, m, p + 1, r) (-1)^r \\
(-1)^{p + 1} \epsilon(n, m, p + 1, r) (-1)^{n - r}
& =
(-1)^n \epsilon(n, m + 1, p + 1, r) (-1)^{p + 1 - r} \\
\epsilon(n + 1, m, p + 1, r) (-1)^r
& =
(-1)^{1 + n} \epsilon(n, m + 1, p + 1, r - 1) \\
(-1)^j \epsilon(n, m, p, r)
& =
(-1)^n \epsilon(n, m + 1, p + 1, r) (-1)^{j - r} \\
(-1)^j \epsilon(n, m, p, r - 1)
& =
\epsilon(n + 1, m, p + 1, r) (-1)^j
\end{align*}

(La troisième égalité est nécessaire pour obtenir les $r = j$
à partir $d(\alpha) \cup \beta$ et $(-1)^n \alpha \cup d(\beta)$
Nous laissons les vérifications au lecteur. (Alternativement, vérifiez les script signs.gp dans le sous-répertoire scripts du projet Champs.)

\medskip\noindent

Associativité du produit de la tasse. Supposons que ${\mathcal F}^\bullet$,
${\mathcal G}^\bullet$ et ${\mathcal H}^\bullet$ sont nés complexes d'intérieur de groupes abéliens sur $X$. La carte évidente (sans intervention des signes) est un isomorphisme des complexes
$$
\text{Tot}(
\text{Tot}({\mathcal F}^\bullet \otimes_{\mathbf Z} {\mathcal G}^\bullet)
\otimes_{\mathbf Z} {\mathcal H}^\bullet
)
\longrightarrow
\text{Tot}(
{\mathcal F}^\bullet \otimes_{\mathbf Z}
\text{Tot}({\mathcal G}^\bullet \otimes_{\mathbf Z} {\mathcal H}^\bullet)
).
$$
Une autre façon de dire ceci est que le triple complexe
${\mathcal F}^\bullet \otimes_{\mathbf Z} {\mathcal G}^\bullet
\otimes_{\mathbf Z} {\mathcal H}^\bullet$ donne lieu à un complexe bien défini total avec différentiel satisfaisant
$$
d(\alpha \otimes \beta \otimes \gamma) =
d(\alpha) \otimes \beta \otimes \gamma +
(-1)^{\deg(\alpha)} \alpha \otimes d(\beta) \otimes \gamma +
(-1)^{\deg(\alpha) + \deg(\beta)} \alpha \otimes \beta \otimes d(\gamma)
$$
pour les éléments homogènes. En utilisant cette application, il est facile de vérifier que
$$
(\alpha \cup \beta) \cup \gamma = \alpha \cup ( \beta \cup \gamma)
$$
à savoir, si $\alpha$ est de degré $a$, $\beta$ de degré $b$ et
$\gamma$ de degré $c$, alors

\begin{align*}
((\alpha \cup \beta) \cup \gamma)_{i_0 \ldots i_p}
= &
\sum\nolimits_{r = 0}^p
\epsilon(a + b, c, p, r)
(\alpha \cup \beta)_{i_0 \ldots i_r} \otimes \gamma_{i_r \ldots i_p}
\\
= &
\sum\nolimits_{r = 0}^p
\sum\nolimits_{s = 0}^r
\epsilon(a + b, c, p, r) \epsilon(a, b, r, s)
\alpha_{i_0 \ldots i_s} \otimes
\beta_{i_s \ldots i_r} \otimes
\gamma_{i_r \ldots i_p}
\end{align*}

et

\begin{align*}
(\alpha \cup (\beta \cup \gamma)_{i_0\ldots i_p}
= &
\sum\nolimits_{s = 0}^p
\epsilon(a, b + c, p, s)
\alpha_{i_0 \ldots i_s} \otimes (\beta \cup \gamma)_{i_s \ldots i_p}
\\
= &
\sum\nolimits_{s = 0}^p
\sum\nolimits_{r = s}^p
\epsilon(a, b + c, p, s) \epsilon(b, c, p - s, r - s)
\alpha_{i_0 \ldots i_s} \otimes \beta_{i_s \ldots i_r} \otimes
\gamma_{i_r \ldots i_p}
\end{align*}

et un mod trivial $2$ (Alternativement, vérifiez les script signs.gp dans le sous-répertoire scripts du projet Champs.)

\medskip\noindent

Enfin, nous indiquons pourquoi le produit de la coupe conserve une structure commutative gradée, au moins au niveau cohomologique.
$\tau$ présenté ci-dessus. Soit ${\mathcal F}^\bullet$ un complexe borné inférieurement de groupes abéliens, muni d'une multiplication commutative graduée
$$
\wedge^\bullet :
\text{Tot}({\mathcal F}^\bullet\otimes {\mathcal F}^\bullet)
\longrightarrow
{\mathcal F}^\bullet.
$$
Cela signifie ce qui suit : $s$ une section locale de
${\mathcal F}^a$, $t$ une section locale de ${\mathcal F}^b$
Nous avons $s \wedge t$ une section locale de ${\mathcal F}^{a + b}$. Commutatif graduel signifie que nous avons
$s \wedge t = (-1)^{ab} t \wedge s$Depuis... $\wedge$ est une application des complexes que nous avons
$d(s\wedge t) = d(s) \wedge t + (-1)^a s \wedge d(t)$. La composition
$$
\text{Tot}(
\text{Tot}(\check{\mathcal{C}}^\bullet({\mathcal U}, {\mathcal F}^\bullet))
\otimes
\text{Tot}(\check{\mathcal{C}}^\bullet({\mathcal U}, {\mathcal F}^\bullet))
)
\to
\text{Tot}(
\check{\mathcal{C}}^\bullet({\mathcal U},
\text{Tot}({\mathcal F}^\bullet\otimes_{\mathbf Z}{\mathcal F}^\bullet))
)
\to
\text{Tot}(\check{\mathcal{C}}^\bullet({\mathcal U}, {\mathcal F}^\bullet))
$$
induit une tasse produit sur la cohomologie
$$
H^n(
\text{Tot}(\check{\mathcal{C}}^\bullet({\mathcal U}, {\mathcal F}^\bullet))
)
\times
H^m(
\text{Tot}(\check{\mathcal{C}}^\bullet({\mathcal U}, {\mathcal F}^\bullet))
)
\longrightarrow
H^{n + m}(
\text{Tot}(\check{\mathcal{C}}^\bullet({\mathcal U}, {\mathcal F}^\bullet))
)
$$
et donc, après passage à la limite inductive, également un produit sur la cohomologie de {\v C}ech et, par conséquent (en utilisant des hyperrecouvrements si nécessaire), un produit sur la cohomologie de ${\mathcal F}^\bullet$. Nous affirmons que ce produit (sur la cohomologie) est lui aussi commutatif gradué. Pour le prouver, considérons d'abord un élément $\alpha$ de degré $n$ dans
$\text{Tot}(\check{\mathcal{C}}^\bullet({\mathcal U}, {\mathcal F}^\bullet))$
et un élément $\beta$ de degré $m$ dans
$\text{Tot}(\check{\mathcal{C}}^\bullet({\mathcal U}, {\mathcal F}^\bullet))$
et calculons

\begin{align*}
\wedge^\bullet(\alpha \cup \beta)_{i_0 \ldots i_p}
= &
\sum\nolimits_{r = 0}^p
\epsilon(n, m, p, r)
\alpha_{i_0 \ldots i_r} \wedge \beta_{i_r \ldots i_p} \\
= &
\sum\nolimits_{r = 0}^p
\epsilon(n, m, p, r)
(-1)^{\deg(\alpha_{i_0 \ldots i_r})\deg(\beta_{i_r \ldots i_p})}
\beta_{i_r \ldots i_p} \wedge \alpha_{i_0 \ldots i_r}
\end{align*}

parce que $\wedge$ est classé comme commutatif. D'autre part, nous avons

\begin{align*}
\tau(\wedge^\bullet(\tau(\beta) \cup \tau(\alpha)))_{i_0 \ldots i_p}
= &
\chi(p)
\sum\nolimits_{r = 0}^p
\epsilon(m, n, p, r)
\tau(\beta)_{i_p \ldots i_{p - r}} \wedge \tau(\alpha)_{i_{p - r} \ldots i_0}
\\
= &
\chi(p)
\sum\nolimits_{r = 0}^p
\epsilon(m, n, p, r) \chi(r) \chi(p - r)
\beta_{i_{p - r} \ldots i_p} \wedge \alpha_{i_0 \ldots i_{p - r}}
\\
= &
\chi(p)
\sum\nolimits_{r = 0}^p
\epsilon(m, n, p, p - r) \chi(r) \chi(p - r)
\beta_{i_r \ldots i_p} \wedge \alpha_{i_0 \ldots i_r}
\end{align*}

où $\chi(t) = (-1)^{\frac{t(t + 1)}{2}}$. Puisque nous avons prouvé plus tôt que
$\tau$ agit comme l'identité sur la cohomologie, nous devons vérifier que
$$
\epsilon(n, m, p, r)
(-1)^{(n - r)(m - (p - r))}
=
(-1)^{nm} \chi(p)\epsilon(m, n, p, p - r) \chi(r) \chi(p - r)
$$
Un calcul immédiat modulo $2$ le vérifie. On peut aussi consulter le script signs.gp dans le sous-répertoire scripts du projet Champs.

\medskip\noindent
Enfin, étudions la compatibilité du produit cup avec les applications de bord. Supposons que $$
0
\to
{\mathcal F}_1^\bullet
\to
{\mathcal F}_2^\bullet
\to
{\mathcal F}_3^\bullet
\to
0
\quad\text{et}\quad
0
\leftarrow
{\mathcal G}_1^\bullet
\leftarrow
{\mathcal G}_2^\bullet
\leftarrow
{\mathcal G}_3^\bullet
\leftarrow
0
$$ soient des suites exactes courtes de complexes bornés inférieurement de faisceaux abéliens sur $X$. Soit ${\mathcal H}^\bullet$ un autre complexe borné inférieurement de faisceaux abéliens, et supposons que nous ayons des morphismes de complexes $$
\gamma_i :
\text{Tot}({\mathcal F}_i^\bullet \otimes_{\mathbf Z} {\mathcal G}_i^\bullet)
\longrightarrow
{\mathcal H}^\bullet
$$ compatibles avec les morphismes entre ces complexes, c'est-à-dire tels que les diagrammes $$
\xymatrix{
\text{Tot}({\mathcal F}_1^\bullet \otimes_{\mathbf Z} {\mathcal G}_1^\bullet)
\ar[d]_{\gamma_1}
&
\text{Tot}({\mathcal F}_1^\bullet \otimes_{\mathbf Z} {\mathcal G}_2^\bullet)
\ar[l] \ar[d]
\\
\mathcal{H}^\bullet &
\text{Tot}({\mathcal F}_2^\bullet \otimes_{\mathbf Z} {\mathcal G}_2^\bullet)
\ar[l]_-{\gamma_2}
}
$$ et $$
\xymatrix{
\text{Tot}({\mathcal F}_2^\bullet \otimes_{\mathbf Z} {\mathcal G}_2^\bullet)
\ar[d]_{\gamma_2}
&
\text{Tot}({\mathcal F}_2^\bullet \otimes_{\mathbf Z} {\mathcal G}_3^\bullet)
\ar[l] \ar[d]
\\
\mathcal{H}^\bullet &
\text{Tot}({\mathcal F}_3^\bullet \otimes_{\mathbf Z} {\mathcal G}_3^\bullet)
\ar[l]_-{\gamma_3}
}
$$ sont commutatifs.

\begin{lemma}

\label{cohomology-lemma-compute-sign-cup-product-boundaries}
Dans la situation ci-dessus, supposons que la cohomologie de {\v C}ech coïncide avec la cohomologie pour les faisceaux $\mathcal{F}_i^p$ et $\mathcal{G}_j^q$. Soient $a_3 \in H^n(X, \mathcal{F}_3^\bullet)$ et $b_1 \in H^m(X, \mathcal{G}_1^\bullet)$. Alors $$
\gamma_1( \partial a_3 \cup b_1) =
(-1)^{n + 1} \gamma_3( a_3 \cup \partial b_1)
$$ dans $H^{n + m + 1}(X, \mathcal{H}^\bullet)$, où $\partial$ désigne l'application de bord en cohomologie associée aux suites exactes courtes de complexes ci-dessus.
\end{lemma}

\begin{proof}

Nous utiliserons les conventions et notations suivantes. Nous considérons
${\mathcal F}_1^p$ comme un sous-faisceau de ${\mathcal F}_2^p$ et
${\mathcal G}_3^q$ comme un sous-faisceau de ${\mathcal G}_2^q$. Ainsi, si $s$ est une section locale de ${\mathcal F}_1^p$, nous notons encore $s$ la section correspondante de ${\mathcal F}_2^p$. De même pour les sections locales de ${\mathcal G}_3^q$. En outre, si $s$ est une section locale de ${\mathcal F}_2^p$, nous notons
$\bar s$ son image dans ${\mathcal F}_3^p$. On adopte la même convention pour l'application ${\mathcal G}_2^q \to {\mathcal G}^q_1$. En particulier, si
$s$ est une section locale de ${\mathcal F}_2^p$ et $\bar s = 0$
alors $s$ est une section locale de ${\mathcal F}_1^p$. La commutativité des diagrammes ci-dessus implique, pour les sections locales $s$ de
${\mathcal F}_2^p$ et $t$ de ${\mathcal G}_3^q$ que
$\gamma_2(s \otimes t) = \gamma_3(\bar s \otimes t)$ en tant que sections
${\mathcal H}^{p + q}$.

\medskip\noindent

Soit ${\mathcal U} : X =  \bigcup_{i \in I} U_i$ un recouvrement ouvert de $X$. Supposons que $\alpha_3$, respectivement $\beta_1$, soit un cocycle de degré $n$, respectivement $m$, de
$\text{Tot}(
\check{\mathcal{C}}^\bullet({\mathcal U}, {\mathcal F}_3^\bullet))$, respectivement de $\text{Tot}(
\check{\mathcal{C}}^\bullet({\mathcal U}, {\mathcal G}_1^\bullet))$
représentant $a_3$, respectivement $b_1$. Après avoir raffiné $\mathcal{U}$ si nécessaire, nous pouvons trouver des cochaînes $\alpha_2$, respectivement $\beta_2$, de degré $n$, respectivement $m$, dans
$\text{Tot}(
\check{\mathcal{C}}^\bullet({\mathcal U}, {\mathcal F}_2^\bullet))$, respectivement dans $\text{Tot}(
\check{\mathcal{C}}^\bullet({\mathcal U}, {\mathcal G}_2^\bullet))$
qui s'envoient sur $\alpha_3$, respectivement $\beta_1$. On a alors
$$
\overline{d(\alpha_2)} = d(\bar \alpha_2) = 0
\quad\text{et}\quad
\overline{d(\beta_2)} = d(\bar \beta_2) = 0.
$$
Cela signifie que $\alpha_1 = d(\alpha_2)$ est un cocycle de degré $n + 1$ dans
$\text{Tot}(\check{\mathcal{C}}^\bullet({\mathcal U}, {\mathcal F}_1^\bullet))$
représentant $\partial a_3$. De même, $\beta_3 = d(\beta_2)$ est un cocycle de degré $m + 1$ dans
$\text{Tot}(\check{\mathcal{C}}^\bullet({\mathcal U}, {\mathcal G}_3^\bullet))$
représentant $\partial b_1$. Ainsi nous pouvons calculer

\begin{align*}
d(\gamma_2(\alpha_2 \cup \beta_2))
& =
\gamma_2(d(\alpha_2 \cup \beta_2))
\\
& =
\gamma_2(d(\alpha_2) \cup \beta_2 + (-1)^n \alpha_2 \cup d(\beta_2) )
\\
& =
\gamma_2( \alpha_1 \cup \beta_2)  + (-1)^n \gamma_2( \alpha_2 \cup \beta_3)
\\
& =
\gamma_1(\alpha_1 \cup \beta_1) + (-1)^n \gamma_3(\alpha_3 \cup \beta_3)
\end{align*}

Donc cela nous dit même que le signe est $(-1)^{n + 1}$ comme indiqué dans le lemme\footnote{Le signe dépend de la convention pour les signes dans la suite exacte longue en cohomologie associée à un triangle dans $D(X)$. Les conventions dans le projet Champs sont : a) les triangles distincts correspondent à des séquences exactes par terme et b) les cartes limites dans la suite exacte longue sont données par les cartes dans la lemme du serpent sans intervention de signes. Voir Catégories dérivées, section \ref{derived-section-homotopy-triangulated}.}.

\end{proof}

\begin{lemma}

\label{cohomology-lemma-boundary-derivation-over-cup-product}

Soit $X$ un espace topologique et soit $\mathcal{O}' \to \mathcal{O}$ une surjection de faisceaux d'anneaux dont le noyau $\mathcal{I} \subset \mathcal{O}'$
est de carré nul. Alors $M = H^1(X, \mathcal{I})$ est un
$R = H^0(X, \mathcal{O})$-module et l'application de bord
$\partial : R \to M$ associée à la suite exacte courte
$$
0 \to \mathcal{I} \to \mathcal{O}' \to \mathcal{O} \to 0
$$
est une dérivation (Algèbre, définition \ref{algebra-definition-derivation}).

\end{lemma}

\begin{proof}

L'application $\mathcal{O}' \to \SheafHom(\mathcal{I}, \mathcal{I})$
se factorise par $\mathcal{O}$ puisque $\mathcal{I} \cdot \mathcal{I} = 0$
par hypothèse. Ainsi, $\mathcal{I}$ est un faisceau de $\mathcal{O}$-modules, ce qui définit la structure de $R$-module sur $M$. L'application de bord est additive; il suffit donc de prouver la règle de Leibniz. Soit $f \in R$. Choisissons un recouvrement ouvert
$\mathcal{U} : X = \bigcup U_i$ tel qu'il existe des éléments
$f_i \in \mathcal{O}'(U_i)$ relevant $f|_{U_i} \in \mathcal{O}(U_i)$. Remarquons que $f_i - f_j$ appartient à $\mathcal{I}(U_i \cap U_j)$. Alors $\partial(f)$ correspond à la classe de cohomologie de {\v C}ech du $1$-cocycle $\alpha$ défini par $\alpha_{i_0i_1} = f_{i_0} - f_{i_1}$. (D'après le lemme \ref{cohomology-lemma-cech-h1}, le premier groupe de cohomologie de {\v C}ech relatif à $\mathcal{U}$ est un sous-module de $M$.) Soit ensuite $g \in R$ un second élément et supposons, après avoir éventuellement raffiné le recouvrement ouvert, que $g_i \in \mathcal{O}'(U_i)$ relève
$g|_{U_i} \in \mathcal{O}(U_i)$. Alors nous voyons que
$\partial(g)$ est donné par le cocycle $\beta$ avec
$\beta_{i_0i_1} = g_{i_0} - g_{i_1}$. Comme $f_ig_i \in \mathcal{O}'(U_i)$
relève $fg|_{U_i}$, on voit que
$\partial(fg)$ est donné par le cocycle $\gamma$ avec
$$
\gamma_{i_0i_1} = f_{i_0}g_{i_0} - f_{i_1}g_{i_1} =
(f_{i_0} - f_{i_1})g_{i_0} + f_{i_1}(g_{i_0} - g_{i_1}) =
\alpha_{i_0i_1}g + f\beta_{i_0i_1}
$$
par notre définition du $\mathcal{O}$-module structure sur $\mathcal{I}$. Cela prouve la règle de Leibniz et la preuve est complète.

\end{proof}









\section{Résolutions plates}

\label{cohomology-section-flat}

\noindent
Une référence pour le contenu de cette section est \cite{Spaltenstein}. Soit $(X, \mathcal{O}_X)$ un espace annelé. D'après Modules, lemme \ref{modules-lemma-module-quotient-flat}, tout $\mathcal{O}_X$-module est quotient d'un $\mathcal{O}_X$-module plat. D'après Catégories dérivées, lemme \ref{derived-lemma-subcategory-left-resolution}, tout complexe borné supérieurement de $\mathcal{O}_X$-modules admet une résolution à gauche par un complexe borné supérieurement de $\mathcal{O}_X$-modules plats. Pour les complexes non bornés, toutefois, les résolutions plates ne suffisent pas.

\begin{lemma}

\label{cohomology-lemma-derived-tor-exact}
Soit $(X, \mathcal{O}_X)$ un espace annelé. Soit $\mathcal{G}^\bullet$ un complexe de $\mathcal{O}_X$-modules. Les foncteurs $$
K(\textit{Mod}(\mathcal{O}_X))
\longrightarrow
K(\textit{Mod}(\mathcal{O}_X)),
\quad
\mathcal{F}^\bullet \longmapsto
\text{Tot}(\mathcal{G}^\bullet \otimes_{\mathcal{O}_X} \mathcal{F}^\bullet)
$$ et $$
K(\textit{Mod}(\mathcal{O}_X))
\longrightarrow
K(\textit{Mod}(\mathcal{O}_X)),
\quad
\mathcal{F}^\bullet \longmapsto
\text{Tot}(\mathcal{F}^\bullet \otimes_{\mathcal{O}_X} \mathcal{G}^\bullet)
$$ sont des foncteurs exacts de catégories triangulées.
\end{lemma}

\begin{proof}
Cela résulte de Catégories dérivées, remarque \ref{derived-remark-double-complex-as-tensor-product-of}.
\end{proof}

\begin{definition}

\label{cohomology-definition-K-flat}

Soit $(X, \mathcal{O}_X)$ un espace annelé. Un complexe $\mathcal{K}^\bullet$ de $\mathcal{O}_X$-modules est dit {\it K-plat} si, pour tout complexe acyclique $\mathcal{F}^\bullet$
de $\mathcal{O}_X$-modules, le complexe
$$
\text{Tot}(\mathcal{F}^\bullet \otimes_{\mathcal{O}_X} \mathcal{K}^\bullet)
$$
est acyclique.

\end{definition}

\begin{lemma}

\label{cohomology-lemma-K-flat-quasi-isomorphism}

Soit $(X, \mathcal{O}_X)$ un espace annelé. Soit $\mathcal{K}^\bullet$ un complexe K-plat. Alors le foncteur
$$
K(\textit{Mod}(\mathcal{O}_X))
\longrightarrow
K(\textit{Mod}(\mathcal{O}_X)), \quad
\mathcal{F}^\bullet
\longmapsto
\text{Tot}(\mathcal{F}^\bullet \otimes_{\mathcal{O}_X} \mathcal{K}^\bullet)
$$
transforme les quasi-isomorphismes en quasi-isomorphismes.

\end{lemma}

\begin{proof}

Cela résulte du lemme \ref{cohomology-lemma-derived-tor-exact}
et le fait que les quasi-isomorphismes sont caractérisés par des cônes acycliques.

\end{proof}

\begin{lemma}

\label{cohomology-lemma-check-K-flat-stalks}

Soit $(X, \mathcal{O}_X)$ un espace annelé. Soit $\mathcal{K}^\bullet$
un complexe de $\mathcal{O}_X$-modules. Le complexe $\mathcal{K}^\bullet$
est K-plat si et seulement si, pour tout $x \in X$, le complexe
$\mathcal{K}_x^\bullet$ de $\mathcal{O}_{X, x}$-modules est K-plat (Compléments sur l'algèbre, définition \ref{more-algebra-definition-K-flat}).

\end{lemma}

\begin{proof}

Si $\mathcal{K}_x^\bullet$ est K-plat pour tout $x \in X$, alors $\mathcal{K}^\bullet$ est K-plat, car $\otimes$ et les sommes directes commutent à la prise des germes et l'exactitude se vérifie sur les germes; voir Modules, lemme \ref{modules-lemma-abelian}. Réciproquement, supposons $\mathcal{K}^\bullet$ K-plat. Fixons $x \in X$
et soit
$M^\bullet$ un complexe acyclique de $\mathcal{O}_{X, x}$-modules. Alors $i_{x, *}M^\bullet$ est un complexe acyclique de $\mathcal{O}_X$-modules, donc $\text{Tot}(i_{x, *}M^\bullet \otimes_{\mathcal{O}_X} \mathcal{K}^\bullet)$
est acyclique. En prenant le germe en $x$, on voit que
$\text{Tot}(M^\bullet \otimes_{\mathcal{O}_{X, x}} \mathcal{K}_x^\bullet)$
est acyclique.

\end{proof}

\begin{lemma}

\label{cohomology-lemma-tensor-product-K-flat}

Soit $(X, \mathcal{O}_X)$ un espace annelé. Si $\mathcal{K}^\bullet$ et $\mathcal{L}^\bullet$ sont des complexes K-plats de $\mathcal{O}_X$-modules, alors
$\text{Tot}(\mathcal{K}^\bullet \otimes_{\mathcal{O}_X} \mathcal{L}^\bullet)$
est un complexe K-plat de $\mathcal{O}_X$-modules.

\end{lemma}

\begin{proof}

Cela résulte de l'isomorphisme
$$
\text{Tot}(\mathcal{M}^\bullet \otimes_{\mathcal{O}_X}
\text{Tot}(\mathcal{K}^\bullet \otimes_{\mathcal{O}_X} \mathcal{L}^\bullet))
=
\text{Tot}(\text{Tot}(\mathcal{M}^\bullet \otimes_{\mathcal{O}_X}
\mathcal{K}^\bullet) \otimes_{\mathcal{O}_X} \mathcal{L}^\bullet)
$$
et la définition.

\end{proof}

\begin{lemma}

\label{cohomology-lemma-K-flat-two-out-of-three}
Soit $(X, \mathcal{O}_X)$ un espace annelé. Soit $(\mathcal{K}_1^\bullet, \mathcal{K}_2^\bullet, \mathcal{K}_3^\bullet)$ un triangle distingué dans $K(\textit{Mod}(\mathcal{O}_X))$. Si deux des trois complexes $\mathcal{K}_i^\bullet$ sont K-plats, alors le troisième l'est aussi.
\end{lemma}

\begin{proof}
Cela résulte du lemme \ref{cohomology-lemma-derived-tor-exact} et du fait que, dans un triangle distingué de $K(\textit{Mod}(\mathcal{O}_X))$, si deux des trois objets sont acycliques, alors le troisième l'est aussi.
\end{proof}

\begin{lemma}

\label{cohomology-lemma-K-flat-two-out-of-three-ses}
Soit $(X, \mathcal{O}_X)$ un espace annelé. Soit $0 \to \mathcal{K}_1^\bullet \to \mathcal{K}_2^\bullet \to
\mathcal{K}_3^\bullet \to 0$ une suite exacte courte de complexes, dont les termes de $\mathcal{K}_3^\bullet$ sont des $\mathcal{O}_X$-modules plats. Si deux des trois complexes $\mathcal{K}_i^\bullet$ sont K-plats, alors le troisième l'est aussi.
\end{lemma}

\begin{proof}
D'après Modules, lemme \ref{modules-lemma-flat-tor-zero}, pour tout complexe $\mathcal{L}^\bullet$, nous obtenons une suite exacte courte $$
0 \to
\text{Tot}(\mathcal{L}^\bullet \otimes_{\mathcal{O}_X} \mathcal{K}_1^\bullet)
\to
\text{Tot}(\mathcal{L}^\bullet \otimes_{\mathcal{O}_X} \mathcal{K}_1^\bullet)
\to
\text{Tot}(\mathcal{L}^\bullet \otimes_{\mathcal{O}_X} \mathcal{K}_1^\bullet)
\to 0
$$ de complexes. Le lemme résulte donc de la suite exacte longue des faisceaux de cohomologie et de la définition des complexes K-plats.
\end{proof}

\begin{lemma}

\label{cohomology-lemma-pullback-K-flat}

Soit $f : (X, \mathcal{O}_X) \to (Y, \mathcal{O}_Y)$ un morphisme d'espaces annelés. L'image inverse d'un complexe K-plat de $\mathcal{O}_Y$-modules est un complexe K-plat de $\mathcal{O}_X$-modules.

\end{lemma}

\begin{proof}
Nous pouvons vérifier ceci sur les germes, voir lemme \ref{cohomology-lemma-check-K-flat-stalks}. C'est ainsi que ceci résulte de Faisceaux, lemme \ref{sheaves-lemma-stalk-pullback-modules} et Compléments sur l'algèbre, lemme \ref{more-algebra-lemma-base-change-K-flat}.
\end{proof}

\begin{lemma}

\label{cohomology-lemma-bounded-flat-K-flat}
Soit $(X, \mathcal{O}_X)$ un espace annelé. Tout complexe borné supérieurement de $\mathcal{O}_X$-modules plats est K-plat.
\end{lemma}

\begin{proof}
Nous pouvons vérifier ceci sur les germes, voir lemme \ref{cohomology-lemma-check-K-flat-stalks}. Ainsi ce lemme suit des Modules, lemme \ref{modules-lemma-flat-stalks-flat} et Compléments sur l'algèbre, lemme \ref{more-algebra-lemma-derived-tor-quasi-isomorphism}.
\end{proof}

\noindent
Dans le lemme suivant, la colimite d'un système de complexes désigne la colimite terme à terme.

\begin{lemma}

\label{cohomology-lemma-colimit-K-flat}

Soit $(X, \mathcal{O}_X)$ un espace annelé. Soit $\mathcal{K}_1^\bullet \to \mathcal{K}_2^\bullet \to \ldots$
un système de complexes K-plats. Alors $\colim_i \mathcal{K}_i^\bullet$ est K-plat.

\end{lemma}

\begin{proof}

Puisque nous prenons les colimites terme à terme, il est clair que
$$
\colim_i \text{Tot}(
\mathcal{F}^\bullet \otimes_{\mathcal{O}_X} \mathcal{K}_i^\bullet)
=
\text{Tot}(\mathcal{F}^\bullet \otimes_{\mathcal{O}_X}
\colim_i \mathcal{K}_i^\bullet)
$$
Le lemme découle donc de l'exactitude des colimites filtrantes.

\end{proof}

\begin{lemma}

\label{cohomology-lemma-resolution-by-direct-sums-extensions-by-zero}

Soit $(X, \mathcal{O}_X)$ un espace annelé. Pour tout complexe $\mathcal{G}^\bullet$ de $\mathcal{O}_X$-modules il existe un diagramme commutatif des complexes de $\mathcal{O}_X$-modules
$$
\xymatrix{
\mathcal{K}_1^\bullet \ar[d] \ar[r] &
\mathcal{K}_2^\bullet \ar[d] \ar[r] & \ldots \\
\tau_{\leq 1}\mathcal{G}^\bullet \ar[r] &
\tau_{\leq 2}\mathcal{G}^\bullet \ar[r] & \ldots
}
$$
avec les propriétés suivantes : (1) les flèches verticales sont des quasi-isomorphismes et sont surjectives terme à terme, (2) chaque $\mathcal{K}_n^\bullet$ est un complexe borné supérieurement dont les termes sont des sommes directes de $\mathcal{O}_X$-modules de la forme
$j_{U!}\mathcal{O}_U$, et (3) les applications $\mathcal{K}_n^\bullet \to \mathcal{K}_{n + 1}^\bullet$ sont des injections scindées terme à terme dont les conoyaux sont des sommes directes de
$\mathcal{O}_X$-modules de la forme $j_{U!}\mathcal{O}_U$. En outre, l'application
$\colim \mathcal{K}_n^\bullet \to \mathcal{G}^\bullet$ est un quasi-isomorphisme.

\end{lemma}

\begin{proof}
L'existence du diagramme et les propriétés (1), (2), (3) résultent immédiatement de Modules, lemme \ref{modules-lemma-module-quotient-flat}, et de Catégories dérivées, lemme \ref{derived-lemma-special-direct-system}. L'application induite $\colim \mathcal{K}_n^\bullet \to \mathcal{G}^\bullet$ est un quasi-isomorphisme parce que les colimites filtrantes sont exactes.
\end{proof}

\begin{lemma}

\label{cohomology-lemma-K-flat-resolution}

Soit $(X, \mathcal{O}_X)$ un espace annelé. Pour tout complexe $\mathcal{G}^\bullet$, il existe un complexe $K$-plat
$\mathcal{K}^\bullet$ dont les termes sont des $\mathcal{O}_X$-modules plats, ainsi qu'un quasi-isomorphisme $\mathcal{K}^\bullet \to \mathcal{G}^\bullet$
surjectif terme à terme.

\end{lemma}

\begin{proof}

Choisissez un diagramme comme dans lemme \ref{cohomology-lemma-resolution-by-direct-sums-extensions-by-zero}. Chaque complexe $\mathcal{K}_n^\bullet$ est un complexe borné supérieurement de modules plats, voir Modules, lemme \ref{modules-lemma-j-shriek-flat}. D'où $\mathcal{K}_n^\bullet$ est K-plat par lemme \ref{cohomology-lemma-bounded-flat-K-flat}. Ainsi $\colim \mathcal{K}_n^\bullet$ est K-plat par lemme \ref{cohomology-lemma-colimit-K-flat}. La carte induite
$\colim \mathcal{K}_n^\bullet \to \mathcal{G}^\bullet$
est un quasi-isomorphisme et surjectif par construction. Propriété (3) de lemme \ref{cohomology-lemma-resolution-by-direct-sums-extensions-by-zero}
montre que $\colim \mathcal{K}_n^m$ est une somme directe de modules plats et donc plats qui prouve l'affirmation finale.

\end{proof}

\begin{lemma}

\label{cohomology-lemma-derived-tor-quasi-isomorphism-other-side}

Soit $(X, \mathcal{O}_X)$ un espace annelé.
Soit $\alpha : \mathcal{P}^\bullet \to \mathcal{Q}^\bullet$ un quasi-isomorphisme de complexes K-plats de $\mathcal{O}_X$-modules. Pour tout complexe $\mathcal{F}^\bullet$ de $\mathcal{O}_X$-modules, l'application induite
$$
\text{Tot}(\text{id}_{\mathcal{F}^\bullet} \otimes \alpha) :
\text{Tot}(\mathcal{F}^\bullet \otimes_{\mathcal{O}_X} \mathcal{P}^\bullet)
\longrightarrow
\text{Tot}(\mathcal{F}^\bullet \otimes_{\mathcal{O}_X} \mathcal{Q}^\bullet)
$$
est un quasi-isomorphisme.

\end{lemma}

\begin{proof}

Choisissez un quasi-isomorphisme $\mathcal{K}^\bullet \to \mathcal{F}^\bullet$
où $\mathcal{K}^\bullet$ est un complexe K-plat; voir le lemme \ref{cohomology-lemma-K-flat-resolution}. Considérons le diagramme commutatif
$$
\xymatrix{
\text{Tot}(\mathcal{K}^\bullet
\otimes_{\mathcal{O}_X} \mathcal{P}^\bullet) \ar[r] \ar[d] &
\text{Tot}(\mathcal{K}^\bullet
\otimes_{\mathcal{O}_X} \mathcal{Q}^\bullet) \ar[d] \\
\text{Tot}(\mathcal{F}^\bullet
\otimes_{\mathcal{O}_X} \mathcal{P}^\bullet) \ar[r] &
\text{Tot}(\mathcal{F}^\bullet
\otimes_{\mathcal{O}_X} \mathcal{Q}^\bullet)
}
$$
Le résultat découle du lemme \ref{cohomology-lemma-K-flat-quasi-isomorphism} :
les flèches verticales et la flèche horizontale supérieure sont des quasi-isomorphismes.

\end{proof}

\noindent
Soit $(X, \mathcal{O}_X)$ un espace annelé et soit $\mathcal{F}^\bullet$ un objet de $D(\mathcal{O}_X)$. Choisissons une résolution K-plate $\mathcal{K}^\bullet \to \mathcal{F}^\bullet$; voir le lemme \ref{cohomology-lemma-K-flat-resolution}. Le lemme \ref{cohomology-lemma-derived-tor-exact} fournit un foncteur exact entre catégories triangulées $$
K(\mathcal{O}_X)
\longrightarrow
K(\mathcal{O}_X),
\quad
\mathcal{G}^\bullet
\longmapsto
\text{Tot}(\mathcal{G}^\bullet \otimes_{\mathcal{O}_X} \mathcal{K}^\bullet)
$$ D'après le lemme \ref{cohomology-lemma-K-flat-quasi-isomorphism}, ce foncteur induit un foncteur $D(\mathcal{O}_X) \to D(\mathcal{O}_X)$, puisque $D(\mathcal{O}_X)$ est la localisation de $K(\mathcal{O}_X)$ relativement aux quasi-isomorphismes. Comme la catégorie des résolutions $K$-plates de $\mathcal{F}^\bullet$ est cofiltrante et compte tenu du lemme \ref{cohomology-lemma-derived-tor-quasi-isomorphism-other-side}, le foncteur obtenu, à isomorphisme près, ne dépend pas du choix de $\mathcal{K}^\bullet$.

\begin{definition}

\label{cohomology-definition-derived-tor}
Soit $(X, \mathcal{O}_X)$ un espace annelé. Soit $\mathcal{F}^\bullet$ un objet de $D(\mathcal{O}_X)$. Le {\it produit tensoriel dérivé} $$
- \otimes_{\mathcal{O}_X}^{\mathbf{L}} \mathcal{F}^\bullet :
D(\mathcal{O}_X)
\longrightarrow
D(\mathcal{O}_X)
$$ est le foncteur exact des catégories triangulées décrites ci-dessus.
\end{definition}

\noindent
Il ressort clairement de nos constructions explicites qu'il y a un isomorphisme canonique $$
\mathcal{F}^\bullet \otimes_{\mathcal{O}_X}^{\mathbf{L}} \mathcal{G}^\bullet
\cong
\mathcal{G}^\bullet \otimes_{\mathcal{O}_X}^{\mathbf{L}} \mathcal{F}^\bullet
$$ pour $\mathcal{G}^\bullet$ et $\mathcal{F}^\bullet$ dans $D(\mathcal{O}_X)$. Ici nous utilisons les règles des signes comme indiqué dans Compléments sur l'algèbre, section \ref{more-algebra-section-sign-rules}. Par conséquent, lorsque nous écrivons $\mathcal{F}^\bullet \otimes_{\mathcal{O}_X}^{\mathbf{L}} \mathcal{G}^\bullet$ nous serons généralement agnostiques sur quelle variable nous utilisons pour définir le produit tensoriel dérivé avec.

\begin{definition}

\label{cohomology-definition-tor}
Soit $(X, \mathcal{O}_X)$ un espace annelé. Soit $\mathcal{F}$, $\mathcal{G}$ $\mathcal{O}_X$-modules. Les {\it Tor} de $\mathcal{F}$ et $\mathcal{G}$ sont définis par la formule $$
\text{Tor}_p^{\mathcal{O}_X}(\mathcal{F}, \mathcal{G}) =
H^{-p}(\mathcal{F} \otimes_{\mathcal{O}_X}^\mathbf{L} \mathcal{G})
$$ avec produit tensoriel dérivé comme défini ci-dessus.
\end{definition}

\noindent
Cette définition implique que, pour toute suite exacte courte de $\mathcal{O}_X$-modules $0 \to \mathcal{F}_1 \to \mathcal{F}_2 \to \mathcal{F}_3 \to 0$, nous avons une suite exacte longue de cohomologie $$
\xymatrix{
\mathcal{F}_1 \otimes_{\mathcal{O}_X} \mathcal{G} \ar[r] &
\mathcal{F}_2 \otimes_{\mathcal{O}_X} \mathcal{G} \ar[r] &
\mathcal{F}_3 \otimes_{\mathcal{O}_X} \mathcal{G} \ar[r] & 0 \\
\text{Tor}_1^{\mathcal{O}_X}(\mathcal{F}_1, \mathcal{G}) \ar[r] &
\text{Tor}_1^{\mathcal{O}_X}(\mathcal{F}_2, \mathcal{G}) \ar[r] &
\text{Tor}_1^{\mathcal{O}_X}(\mathcal{F}_3, \mathcal{G}) \ar[ull]
}
$$ pour chaque $\mathcal{O}_X$-module $\mathcal{G}$. Ceci sera appelé la suite exacte longue de $\text{Tor}$ associée à la situation.

\begin{lemma}

\label{cohomology-lemma-flat-tor-zero}

\begin{slogan}

Tor mesure l'écart de planéité.

\end{slogan}

Soit $(X, \mathcal{O}_X)$ un espace annelé. Soit $\mathcal{F}$ un $\mathcal{O}_X$-module. Les conditions suivantes sont équivalentes :

\begin{enumerate}

\item  $\mathcal{F}$ est un $\mathcal{O}_X$-module plat, et
\item  $\text{Tor}_1^{\mathcal{O}_X}(\mathcal{F}, \mathcal{G}) = 0$
pour tout $\mathcal{O}_X$-module $\mathcal{G}$.

\end{enumerate}

\end{lemma}

\begin{proof}
Si $\mathcal{F}$ est plat, alors $\mathcal{F} \otimes_{\mathcal{O}_X} -$ est un foncteur exact et ses foncteurs dérivés s'annulent. Réciproquement, supposons (2). Si $\mathcal{G} \to \mathcal{H}$ est injectif, de conoyau $\mathcal{Q}$, la suite exacte longue de $\text{Tor}$ montre que le noyau de $\mathcal{F} \otimes_{\mathcal{O}_X} \mathcal{G} \to
\mathcal{F} \otimes_{\mathcal{O}_X} \mathcal{H}$ est un quotient de $\text{Tor}_1^{\mathcal{O}_X}(\mathcal{F}, \mathcal{Q})$, qui est nul par hypothèse. Par conséquent, $\mathcal{F}$ est plat.
\end{proof}

\begin{lemma}

\label{cohomology-lemma-factor-through-K-flat}

Soit $(X, \mathcal{O}_X)$ un espace annelé. Soit $a : \mathcal{K}^\bullet \to \mathcal{L}^\bullet$ un morphisme de complexes de $\mathcal{O}_X$-modules. Si $\mathcal{K}^\bullet$ est K-plat, alors il existe un complexe $\mathcal{N}^\bullet$ et des morphismes de complexes
$b : \mathcal{K}^\bullet \to \mathcal{N}^\bullet$
et $c : \mathcal{N}^\bullet \to \mathcal{L}^\bullet$ de telle manière que

\begin{enumerate}

\item  $\mathcal{N}^\bullet$ est K-plat,
\item  $c$ est un quasi-isomorphisme,
\item  $a$ est homotopique à $c \circ b$.

\end{enumerate}

Si les termes de $\mathcal{K}^\bullet$ sont plats, alors nous pouvons choisir
$\mathcal{N}^\bullet$, $b$, $c$
tel que la même chose est vraie pour $\mathcal{N}^\bullet$.

\end{lemma}

\begin{proof}

Nous utiliserons que la catégorie homotopie $K(\textit{Mod}(\mathcal{O}_X))$
est une catégorie triangulée, voir Catégories dérivées, proposition
\ref{derived-proposition-homotopy-category-triangulated}. Choisissez un triangle distingué
$\mathcal{K}^\bullet \to \mathcal{L}^\bullet \to
\mathcal{C}^\bullet \to \mathcal{K}^\bullet[1]$. Choisissez un quasi-isomorphisme $\mathcal{M}^\bullet \to \mathcal{C}^\bullet$ avec
$\mathcal{M}^\bullet$ K-plat avec des termes plats, voir lemme \ref{cohomology-lemma-K-flat-resolution}. Par les axiomes des catégories triangulées, nous pouvons correspondre à la composition
$\mathcal{M}^\bullet \to \mathcal{C}^\bullet \to \mathcal{K}^\bullet[1]$
dans un triangle distingué
$\mathcal{K}^\bullet \to \mathcal{N}^\bullet \to
\mathcal{M}^\bullet \to \mathcal{K}^\bullet[1]$. D'après le lemme \ref{cohomology-lemma-K-flat-two-out-of-three} Nous voyons que
$\mathcal{N}^\bullet$ est K-plat. En utilisant de nouveau les axiomes des catégories triangulées, nous pouvons choisir une application $\mathcal{N}^\bullet \to \mathcal{L}^\bullet$ qui définit un morphisme entre les triangles distingués suivants
$$
\xymatrix{
\mathcal{K}^\bullet \ar[r] \ar[d] &
\mathcal{N}^\bullet \ar[r] \ar[d] &
\mathcal{M}^\bullet \ar[r] \ar[d] &
\mathcal{K}^\bullet[1] \ar[d] \\
\mathcal{K}^\bullet \ar[r] &
\mathcal{L}^\bullet \ar[r] &
\mathcal{C}^\bullet \ar[r] &
\mathcal{K}^\bullet[1]
}
$$
Puisque deux des trois flèches sont quasi-isomorphismes, ainsi est la troisième flèche $\mathcal{N}^\bullet \to \mathcal{L}^\bullet$
par les suites exécute des longues de cohomologie associées à ces triangles distingués (ou vous pouvez regarder l'image de ce diagramme en $D(\mathcal{O}_X)$ et utiliser Derived Catégories, lemme \ref{derived-lemma-third-isomorphism-triangle}
si vous le souhaitez). Ceci termine la preuve de (1), (2), et (3). Pour prouver l'affirmation finale, nous pouvons choisir $\mathcal{N}^\bullet$
de telle manière que $\mathcal{N}^n \cong \mathcal{M}^n \oplus \mathcal{K}^n$, voir Catégories dérivées, lemme
\ref{derived-lemma-improve-distinguished-triangle-homotopy}. Nous obtenons donc la planéité désirée si les termes de $\mathcal{K}^\bullet$ Ils sont plats.

\end{proof}











\section{Image inverse dérivée}

\label{cohomology-section-derived-pullback}

\noindent
Soit $f : (X, \mathcal{O}_X) \to (Y, \mathcal{O}_Y)$ un morphisme d'espaces annelés. Les résolutions K-plates permettent de définir le foncteur image inverse dérivée $$
Lf^* : D(\mathcal{O}_Y) \to D(\mathcal{O}_X)
$$ Plus précisément, pour tout complexe de $\mathcal{O}_Y$-modules $\mathcal{G}^\bullet$, choisissons une résolution K-plate $\mathcal{K}^\bullet \to \mathcal{G}^\bullet$ et posons $Lf^*\mathcal{G}^\bullet = f^*\mathcal{K}^\bullet$. Les lemmes \ref{cohomology-lemma-pullback-K-flat}, \ref{cohomology-lemma-K-flat-resolution} et \ref{cohomology-lemma-derived-tor-quasi-isomorphism-other-side} montrent que cette définition est indépendante des choix. Pour régler tous les détails, il est toutefois plus commode d'employer une théorie générale.

\begin{lemma}

\label{cohomology-lemma-derived-base-change}
La construction ci-dessus est indépendante des choix et définit un foncteur exact des catégories triangulées $Lf^* : D(\mathcal{O}_Y) \to D(\mathcal{O}_X)$.
\end{lemma}

\begin{proof}

Pour cela, nous utilisons la théorie générale développée dans Catégories dérivées, section \ref{derived-section-derived-functors}. Posons $\mathcal{D} = K(\mathcal{O}_Y)$ et $\mathcal{D}' = D(\mathcal{O}_X)$. Notons $F : \mathcal{D} \to \mathcal{D}'$ le foncteur exact entre catégories triangulées défini par la règle
$F(\mathcal{G}^\bullet) = f^*\mathcal{G}^\bullet$. Soit $S$ l'ensemble des quasi-isomorphismes de
$\mathcal{D} = K(\mathcal{O}_Y)$. On obtient une situation comme dans Catégories dérivées, situation \ref{derived-situation-derived-functor},
de sorte que Catégories dérivées, définition
\ref{derived-definition-right-derived-functor-defined}
est applicable. Nous affirmons que $LF$ est défini partout. Cela résulte de Catégories dérivées, lemme \ref{derived-lemma-find-existence-computes},
en prenant pour $\mathcal{P} \subset \Ob(\mathcal{D})$ la collection des complexes $K$-plats : la condition (1) résulte du lemme \ref{cohomology-lemma-K-flat-resolution};
pour vérifier (2), il faut montrer que, pour tout quasi-isomorphisme
$\mathcal{K}_1^\bullet  \to \mathcal{K}_2^\bullet$ entre complexes K-plats de $\mathcal{O}_Y$-modules, l'application
$f^*\mathcal{K}_1^\bullet  \to f^*\mathcal{K}_2^\bullet$ est un quasi-isomorphisme. Écrivons cette application sous la forme
$$
f^{-1}\mathcal{K}_1^\bullet \otimes_{f^{-1}\mathcal{O}_Y} \mathcal{O}_X
\longrightarrow
f^{-1}\mathcal{K}_2^\bullet \otimes_{f^{-1}\mathcal{O}_Y} \mathcal{O}_X
$$
Le foncteur $f^{-1}$ est exact; l'application
$f^{-1}\mathcal{K}_1^\bullet  \to f^{-1}\mathcal{K}_2^\bullet$ est donc un quasi-isomorphisme. Le lemme \ref{cohomology-lemma-pullback-K-flat},
appliqué au morphisme $(X, f^{-1}\mathcal{O}_Y) \to (Y, \mathcal{O}_Y)$
montre que les complexes $f^{-1}\mathcal{K}_1^\bullet$ et $f^{-1}\mathcal{K}_2^\bullet$
sont des complexes K-plats de $f^{-1}\mathcal{O}_Y$-modules. Le lemme \ref{cohomology-lemma-derived-tor-quasi-isomorphism-other-side}
garantit que l'application affichée est un quasi-isomorphisme.
$$
LF :
D(\mathcal{O}_Y) = S^{-1}\mathcal{D}
\longrightarrow
\mathcal{D}' = D(\mathcal{O}_X)
$$
Voir Catégories dérivées, Équation (\ref{derived-equation-everywhere}). Enfin, les catégories dérivées, lemme \ref{derived-lemma-find-existence-computes}
garantit également que
$LF(\mathcal{K}^\bullet) = F(\mathcal{K}^\bullet) = f^*\mathcal{K}^\bullet$
lorsque $\mathcal{K}^\bullet$ est K-plat, c'est-à-dire, $Lf^* = LF$ est effectivement calculé de la manière décrite ci-dessus.

\end{proof}

\begin{lemma}

\label{cohomology-lemma-derived-pullback-composition}
Soient $f : X \to Y$ et $g : Y \to Z$ des morphismes d'espaces annelés. Alors $Lf^* \circ Lg^* = L(g \circ f)^*$ comme foncteurs $D(\mathcal{O}_Z) \to D(\mathcal{O}_X)$.
\end{lemma}

\begin{proof}
Soit $E$ un objet de $D(\mathcal{O}_Z)$. Par construction, $Lg^*E$ se calcule en choisissant un complexe K-plat $\mathcal{K}^\bullet$ représentant $E$ sur $Z$ et en posant $Lg^*E = g^*\mathcal{K}^\bullet$. D'après le lemme \ref{cohomology-lemma-pullback-K-flat}, $g^*\mathcal{K}^\bullet$ est K-plat sur $Y$. Ainsi, $Lf^*Lg^*E$ est représenté par $f^*g^*\mathcal{K}^\bullet = (g \circ f)^*\mathcal{K}^\bullet$, qui représente également $L(g \circ f)^*E$.
\end{proof}

\begin{lemma}

\label{cohomology-lemma-pullback-tensor-product}
Soit $f : (X, \mathcal{O}_X) \to (Y, \mathcal{O}_Y)$ un morphisme d'espaces annelés. Il existe un isomorphisme canonique bifonctoriel $$
Lf^*(
\mathcal{F}^\bullet \otimes_{\mathcal{O}_Y}^{\mathbf{L}} \mathcal{G}^\bullet
) =
Lf^*\mathcal{F}^\bullet 
\otimes_{\mathcal{O}_X}^{\mathbf{L}}
Lf^*\mathcal{G}^\bullet 
$$ pour $\mathcal{F}^\bullet, \mathcal{G}^\bullet \in \Ob(D(\mathcal{O}_Y))$.
\end{lemma}

\begin{proof}
On peut supposer que $\mathcal{F}^\bullet$ et $\mathcal{G}^\bullet$ sont des complexes K-plats. Dans ce cas, $\mathcal{F}^\bullet \otimes_{\mathcal{O}_Y}^{\mathbf{L}} \mathcal{G}^\bullet$ est le complexe total associé au bicomplexe $\mathcal{F}^\bullet \otimes_{\mathcal{O}_Y} \mathcal{G}^\bullet$. D'après le lemme \ref{cohomology-lemma-tensor-product-K-flat}, $\text{Tot}(\mathcal{F}^\bullet \otimes_{\mathcal{O}_Y} \mathcal{G}^\bullet)$ est lui aussi K-plat. L'isomorphisme du lemme provient donc de l'isomorphisme $$
\text{Tot}(f^*\mathcal{F}^\bullet \otimes_{\mathcal{O}_X}
f^*\mathcal{G}^\bullet)
\longrightarrow
f^*\text{Tot}(\mathcal{F}^\bullet \otimes_{\mathcal{O}_Y} \mathcal{G}^\bullet)
$$ dont les constituants sont les isomorphismes $f^*\mathcal{F}^p \otimes_{\mathcal{O}_X} f^*\mathcal{G}^q \to
f^*(\mathcal{F}^p \otimes_{\mathcal{O}_Y} \mathcal{G}^q)$ des modules, lemme \ref{modules-lemma-tensor-product-pullback}.
\end{proof}

\begin{lemma}

\label{cohomology-lemma-variant-derived-pullback}
Soit $f : (X, \mathcal{O}_X) \to (Y, \mathcal{O}_Y)$ un morphisme d'espaces annelés. Il y a un isomorphisme bifonctionnel canonique $$
\mathcal{F}^\bullet
\otimes_{\mathcal{O}_X}^{\mathbf{L}}
Lf^*\mathcal{G}^\bullet
=
\mathcal{F}^\bullet 
\otimes_{f^{-1}\mathcal{O}_Y}^{\mathbf{L}}
f^{-1}\mathcal{G}^\bullet 
$$ pour $\mathcal{F}^\bullet$ dans $D(\mathcal{O}_X)$ et $\mathcal{G}^\bullet$ dans $D(\mathcal{O}_Y)$.
\end{lemma}

\begin{proof}
Soit $\mathcal{F}$ un $\mathcal{O}_X$-module et soit $\mathcal{G}$ un $\mathcal{O}_Y$-module. Alors $\mathcal{F} \otimes_{\mathcal{O}_X} f^*\mathcal{G} =
\mathcal{F} \otimes_{f^{-1}\mathcal{O}_Y} f^{-1}\mathcal{G}$ parce que $f^*\mathcal{G} =
\mathcal{O}_X \otimes_{f^{-1}\mathcal{O}_Y} f^{-1}\mathcal{G}$. Le lemme suit de ceci et les définitions.
\end{proof}

\begin{lemma}

\label{cohomology-lemma-tensor-pull-compatibility}
Soit $f : (X, \mathcal{O}_X) \to (Y, \mathcal{O}_Y)$ un morphisme des espaces annelés. Soit $\mathcal{K}^\bullet$ et $\mathcal{M}^\bullet$ des complexes des modules $\mathcal{O}_Y$. Le schéma $$
\xymatrix{
Lf^*(\mathcal{K}^\bullet
\otimes_{\mathcal{O}_Y}^\mathbf{L}
\mathcal{M}^\bullet) \ar[r] \ar[d] &
Lf^*\text{Tot}(\mathcal{K}^\bullet
\otimes_{\mathcal{O}_Y}
\mathcal{M}^\bullet) \ar[d] \\
Lf^*\mathcal{K}^\bullet \otimes_{\mathcal{O}_X}^\mathbf{L}
Lf^*\mathcal{M}^\bullet \ar[d] &
f^*\text{Tot}(\mathcal{K}^\bullet
\otimes_{\mathcal{O}_Y}
\mathcal{M}^\bullet) \ar[d] \\
f^*\mathcal{K}^\bullet \otimes_{\mathcal{O}_X}^\mathbf{L}
f^*\mathcal{M}^\bullet \ar[r] &
\text{Tot}(f^*\mathcal{K}^\bullet \otimes_{\mathcal{O}_X}
f^*\mathcal{M}^\bullet)
}
$$ se déplace.
\end{lemma}

\begin{proof}
Nous utiliserons l'existence de résolutions K-plates comme dans le lemme \ref{cohomology-lemma-pullback-K-flat}. Si nous choisissons ces résolutions $\mathcal{P}^\bullet \to \mathcal{K}^\bullet$ et $\mathcal{Q}^\bullet \to \mathcal{M}^\bullet$, alors nous voyons que $$
\xymatrix{
Lf^*\text{Tot}(\mathcal{P}^\bullet
\otimes_{\mathcal{O}_Y}
\mathcal{Q}^\bullet) \ar[r] \ar[d] &
Lf^*\text{Tot}(\mathcal{K}^\bullet
\otimes_{\mathcal{O}_Y}
\mathcal{M}^\bullet) \ar[d] \\
f^*\text{Tot}(\mathcal{P}^\bullet
\otimes_{\mathcal{O}_Y}
\mathcal{Q}^\bullet) \ar[d] \ar[r] &
f^*\text{Tot}(\mathcal{K}^\bullet
\otimes_{\mathcal{O}_Y}
\mathcal{M}^\bullet) \ar[d] \\
\text{Tot}(f^*\mathcal{P}^\bullet \otimes_{\mathcal{O}_X}
f^*\mathcal{Q}^\bullet) \ar[r] &
\text{Tot}(f^*\mathcal{K}^\bullet \otimes_{\mathcal{O}_X}
f^*\mathcal{M}^\bullet)
}
$$ commute. Or le côté gauche de ce diagramme coïncide avec celui du diagramme de l'énoncé, par notre choix de $\mathcal{P}^\bullet$ et $\mathcal{Q}^\bullet$ et par le lemme \ref{cohomology-lemma-tensor-product-K-flat}.
\end{proof}









\section{Cohomologie des complexes non bornés}

\label{cohomology-section-unbounded}

\noindent

Soit $(X, \mathcal{O}_X)$ un espace annelé. La catégorie $\textit{Mod}(\mathcal{O}_X)$ est une catégorie abélienne de Grothendieck : elle admet toutes les colimites, les colimites filtrantes y sont exactes et elle possède un générateur, à savoir
$$
\bigoplus\nolimits_{U \subset X\text{ open}} j_{U!}\mathcal{O}_U,
$$
voir Modules, section \ref{modules-section-kernels} et les lemmes \ref{modules-lemma-j-shriek-flat} et
\ref{modules-lemma-module-quotient-flat}. Par Injectifs, théorème
\ref{injectives-theorem-K-injective-embedding-grothendieck}
pour tout complexe $\mathcal{F}^\bullet$ de $\mathcal{O}_X$-modules, il existe un quasi-isomorphisme
$\mathcal{F}^\bullet \to \mathcal{I}^\bullet$ vers un complexe K-injectif de $\mathcal{O}_X$-modules dont tous les termes sont injectifs;
$\mathcal{O}_X$-modules; de plus, cette inclusion peut être choisie fonctoriellement en le complexe $\mathcal{F}^\bullet$. Il résulte alors de Catégories dérivées, lemme \ref{derived-lemma-enough-K-injectives-implies},
que

\begin{enumerate}

\item tout foncteur exact $F : K(\textit{Mod}(\mathcal{O}_X)) \to \mathcal{D}$
à valeurs dans une catégorie triangulée $\mathcal{D}$ admet un foncteur dérivé à droite
$RF : D(\mathcal{O}_X) \to \mathcal{D}$,
\item pour tout foncteur additif
$F : \textit{Mod}(\mathcal{O}_X) \to \mathcal{A}$
à valeurs dans une catégorie abélienne $\mathcal{A}$, on considère le foncteur exact
$F : K(\textit{Mod}(\mathcal{O}_X)) \to K(\mathcal{A})$ induit par $F$
et nous obtenons un foncteur dérivé à droite
$RF : D(\mathcal{O}_X) \to D(\mathcal{A})$.

\end{enumerate}

Par construction, on a $RF(\mathcal{F}^\bullet) = F(\mathcal{I}^\bullet)$,
où $\mathcal{F}^\bullet \to \mathcal{I}^\bullet$ est comme ci-dessus.

\medskip\noindent
Voici quelques exemples de ce qui précède :
\begin{enumerate}
\item Le foncteur $\Gamma(X, -) : \textit{Mod}(\mathcal{O}_X) \to
\text{Mod}_{\Gamma(X, \mathcal{O}_X)}$ donne lieu à $$
R\Gamma(X, -) : D(\mathcal{O}_X) \to D(\Gamma(X, \mathcal{O}_X))
$$ Nous utiliserons la notation $H^i(X, K) = H^i(R\Gamma(X, K))$ pour la cohomologie.
\item Pour un $U \subset X$ ouvert nous considérons le foncteur $\Gamma(U, -) : \textit{Mod}(\mathcal{O}_X) \to
\text{Mod}_{\Gamma(U, \mathcal{O}_X)}$. Ceci donne lieu à $$
R\Gamma(U, -) : D(\mathcal{O}_X) \to D(\Gamma(U, \mathcal{O}_X))
$$ Nous utiliserons la notation $H^i(U, K) = H^i(R\Gamma(U, K))$ pour la cohomologie.
\item Pour un morphisme d'espaces annelés $f : (X, \mathcal{O}_X) \to (Y, \mathcal{O}_Y)$ nous considérons le foncteur $f_* : \textit{Mod}(\mathcal{O}_X) \to \textit{Mod}(\mathcal{O}_Y)$ qui donne lieu à l'image directe totale $$
Rf_* : D(\mathcal{O}_X) \longrightarrow D(\mathcal{O}_Y)
$$ sur des catégories dérivées non délimitées.
\end{enumerate}

\begin{lemma}

\label{cohomology-lemma-adjoint}
Soit $f : (X, \mathcal{O}_X) \to (Y, \mathcal{O}_Y)$ un morphisme des espaces annelés. Le foncteur $Rf_*$ défini ci-dessus et le foncteur $Lf^*$ défini dans le lemme \ref{cohomology-lemma-derived-base-change} sont adjoints: $$
\Hom_{D(\mathcal{O}_X)}(Lf^*\mathcal{G}^\bullet, \mathcal{F}^\bullet)
=
\Hom_{D(\mathcal{O}_Y)}(\mathcal{G}^\bullet, Rf_*\mathcal{F}^\bullet)
$$ bifonctoriellement en $\mathcal{F}^\bullet \in \Ob(D(\mathcal{O}_X))$ et $\mathcal{G}^\bullet \in \Ob(D(\mathcal{O}_Y))$.
\end{lemma}

\begin{proof}
Ceci résulte formellement du fait que $Rf_*$ et $Lf^*$ existent, voir Catégories dérivées, lemme \ref{derived-lemma-derived-adjoint-functors}.
\end{proof}

\begin{lemma}

\label{cohomology-lemma-derived-pushforward-composition}
Soient $f : X \to Y$ et $g : Y \to Z$ des morphismes d'espaces annelés. Alors $Rg_* \circ Rf_* = R(g \circ f)_*$ comme foncteurs $D(\mathcal{O}_X) \to D(\mathcal{O}_Z)$.
\end{lemma}

\begin{proof}
Par lemme \ref{cohomology-lemma-adjoint} nous voyons que $Rg_* \circ Rf_*$ est adjoint à $Lf^* \circ Lg^*$. Nous avons $Lf^* \circ Lg^* = L(g \circ f)^*$ par lemme \ref{cohomology-lemma-derived-pullback-composition} et donc par unicité des foncteurs adjoints nous avons $Rg_* \circ Rf_* = R(g \circ f)_*$.
\end{proof}

\begin{remark}

\label{cohomology-remark-base-change}

La construction des foncteurs dérivés $Lf^*$ et $Rf_*$
permet de construire l'application de changement de base en pleine généralité.
$$
\xymatrix{
X' \ar[r]_{g'} \ar[d]_{f'} &
X \ar[d]^f \\
S' \ar[r]^g &
S
}
$$
soit un diagramme commutatif d'espaces annelés. Soit $K$ un objet de
$D(\mathcal{O}_X)$. Il existe alors une application canonique de changement de base
$$
Lg^*Rf_*K \longrightarrow R(f')_*L(g')^*K
$$
dans $D(\mathcal{O}_{S'})$. Cette application est, par adjonction, associée à une application
$L(f')^*Lg^*Rf_*K \to L(g')^*K$
Puisque $L(f')^*Lg^* = L(g')^*Lf^*$, il s'agit de la même chose que de donner une application
$L(g')^*Lf^*Rf_*K \to L(g')^*K$
que nous pouvons prendre pour être $L(g')^*$ de l'application d'adjonction
$Lf^*Rf_*K \to K$.

\end{remark}

\begin{remark}

\label{cohomology-remark-compose-base-change}

Considérez un diagramme commutatif
$$
\xymatrix{
X' \ar[r]_k \ar[d]_{f'} & X \ar[d]^f \\
Y' \ar[r]^l \ar[d]_{g'} & Y \ar[d]^g \\
Z' \ar[r]^m & Z
}
$$
d'espaces annelés. Les applications de changement de base de la remarque \ref{cohomology-remark-base-change}
associées aux deux carrés se composent en l'application de changement de base du rectangle extérieur. Plus précisément, la composée

\begin{align*}
Lm^* \circ R(g \circ f)_*
& =
Lm^* \circ Rg_* \circ Rf_* \\
& \to Rg'_* \circ Ll^* \circ Rf_* \\
& \to Rg'_* \circ Rf'_* \circ Lk^* \\
& = R(g' \circ f')_* \circ Lk^*
\end{align*}

est l'application de changement de base pour le rectangle. Nous omettons la vérification.

\end{remark}

\begin{remark}

\label{cohomology-remark-compose-base-change-horizontal}

Considérez un diagramme commutatif
$$
\xymatrix{
X'' \ar[r]_{g'} \ar[d]_{f''} & X' \ar[r]_g \ar[d]_{f'} & X \ar[d]^f \\
Y'' \ar[r]^{h'} & Y' \ar[r]^h & Y
}
$$
d'espaces annelés. Les applications de changement de base de la remarque \ref{cohomology-remark-base-change}
associées aux deux carrés se composent en l'application de changement de base du rectangle extérieur. Plus précisément, la composée

\begin{align*}
L(h \circ h')^* \circ Rf_*
& =
L(h')^* \circ Lh_* \circ Rf_* \\
& \to L(h')^* \circ Rf'_* \circ Lg^* \\
& \to Rf''_* \circ L(g')^* \circ Lg^* \\
& = Rf''_* \circ L(g \circ g')^*
\end{align*}

est l'application de changement de base pour le rectangle. Nous omettons la vérification.

\end{remark}

\begin{lemma}

\label{cohomology-lemma-adjoints-push-pull-compatibility}
Soit $f : (X, \mathcal{O}_X) \to (Y, \mathcal{O}_Y)$ un morphisme d'espaces annelés. Soit $\mathcal{K}^\bullet$ un complexe de $\mathcal{O}_X$-modules. Le diagramme $$
\xymatrix{
Lf^*f_*\mathcal{K}^\bullet \ar[r] \ar[d] &
f^*f_*\mathcal{K}^\bullet \ar[d] \\
Lf^*Rf_*\mathcal{K}^\bullet \ar[r] &
\mathcal{K}^\bullet
}
$$ provenant de $Lf^* \to f^*$ sur les complexes, $f_* \to Rf_*$ sur les complexes, et l'adjonction $Lf^* \circ Rf_* \to \text{id}$ se déplace dans $D(\mathcal{O}_X)$.
\end{lemma}

\begin{proof}

Nous utiliserons l'existence de résolutions K-plates et de résolutions K-injectives; voir le lemme \ref{cohomology-lemma-pullback-K-flat}
et la discussion ci-dessus. Choisissez un quasi-isomorphisme
$\mathcal{K}^\bullet \to \mathcal{I}^\bullet$ où $\mathcal{I}^\bullet$
est K-injective comme un complexe de $\mathcal{O}_X$-modules. Choisissez un quasi-isomorphisme $\mathcal{Q}^\bullet \to f_*\mathcal{I}^\bullet$
où $\mathcal{Q}^\bullet$ est K-plat comme un complexe de
$\mathcal{O}_Y$-modules. Nous pouvons choisir un complexe K-plat de
$\mathcal{O}_Y$-modules $\mathcal{P}^\bullet$
et un diagramme de morphismes de complexes
$$
\xymatrix{
\mathcal{P}^\bullet \ar[r] \ar[d] &
f_*\mathcal{K}^\bullet \ar[d] \\
\mathcal{Q}^\bullet \ar[r] & f_*\mathcal{I}^\bullet
}
$$
commutatif à homotopie près, dont la flèche horizontale supérieure est un quasi-isomorphisme. Un tel $\mathcal{P}^\bullet$ existe parce que les quasi-isomorphismes forment un système multiplicatif dans la catégorie homotopique des complexes; on peut ensuite remplacer
$\mathcal{P}^\bullet$ par un complexe K-plat. En appliquant l'image inverse, on obtient un diagramme de morphismes de complexes
$$
\xymatrix{
f^*\mathcal{P}^\bullet \ar[r] \ar[d] &
f^*f_*\mathcal{K}^\bullet \ar[d] \ar[r] &
\mathcal{K}^\bullet \ar[d] \\
f^*\mathcal{Q}^\bullet \ar[r] &
f^*f_*\mathcal{I}^\bullet \ar[r] &
\mathcal{I}^\bullet
}
$$
Le rectangle extérieur prouve l'assertion du lemme.

\end{proof}

\begin{remark}

\label{cohomology-remark-cup-product}
Soit $f : (X, \mathcal{O}_X) \to (Y, \mathcal{O}_Y)$ un morphisme d'espaces annelés. L'adjonction entre $Lf^*$ et $Rf_*$ permet de construire un produit cup relatif $$
Rf_*K \otimes_{\mathcal{O}_Y}^\mathbf{L} Rf_*L
\longrightarrow
Rf_*(K \otimes_{\mathcal{O}_X}^\mathbf{L} L)
$$ dans $D(\mathcal{O}_Y)$ pour tous $K, L$ dans $D(\mathcal{O}_X)$. Autrement dit, cette application est adjointe à une application $Lf^*(Rf_*K \otimes_{\mathcal{O}_Y}^\mathbf{L} Rf_*L) \to
K \otimes_{\mathcal{O}_X}^\mathbf{L} L$ pour laquelle nous pouvons prendre la composition de l'isomorphisme $Lf^*(Rf_*K \otimes_{\mathcal{O}_Y}^\mathbf{L} Rf_*L) =
Lf^*Rf_*K \otimes_{\mathcal{O}_X}^\mathbf{L} Lf^*Rf_*L$ (lemme \ref{cohomology-lemma-pullback-tensor-product}) avec l'application $Lf^*Rf_*K \otimes_{\mathcal{O}_X}^\mathbf{L} Lf^*Rf_*L
\to K \otimes_{\mathcal{O}_X}^\mathbf{L} L$ provenant de la co-unité $Lf^* \circ Rf_* \to \text{id}$.
\end{remark}





\section{Cohomologie des complexes filtrés}

\label{cohomology-section-cohomology-filtered-object}

\noindent
Les complexes filtrés de faisceaux apparaissent souvent de manière naturelle dans l'étude de la cohomologie des variétés algébriques; par exemple, le complexe de de Rham est muni de sa filtration de Hodge. Dans cette section, nous utilisons le très général lemme \ref{injectives-lemma-K-injective-embedding-filtration} d'Injectifs pour construire des suites spectrales en cohomologie et les relier à celles construites précédemment.

\begin{lemma}

\label{cohomology-lemma-spectral-sequence-filtered-object}

Soit $(X, \mathcal{O}_X)$ un espace annelé. Soit $\mathcal{F}^\bullet$ un complexe filtré de $\mathcal{O}_X$-modules. Il existe une suite spectrale canonique $(E_r, \text{d}_r)_{r \geq 1}$ de
$\Gamma(X, \mathcal{O}_X)$-modules bigradués, où $d_r$ est de bidegré $(r, -r + 1)$, et
$$
E_1^{p, q} = H^{p + q}(X, \text{gr}^p\mathcal{F}^\bullet)
$$
Si, pour tout $n$, on a
$$
H^n(X, F^p\mathcal{F}^\bullet) = 0\text{ pour }p \gg 0
\quad\text{et}\quad
H^n(X, F^p\mathcal{F}^\bullet) = H^n(X, \mathcal{F}^\bullet)\text{ pour }p \ll 0
$$
alors la suite spectrale est bornée et converge vers
$H^*(X, \mathcal{F}^\bullet)$.

\end{lemma}

\begin{proof}
(Pour une preuve dans le cas d'un complexe borné inférieurement de modules munis de filtrations finies, voir la remarque ci-dessous.) Choisissons un morphisme de complexes filtrés $j : \mathcal{F}^\bullet \to \mathcal{J}^\bullet$ comme dans Injectifs, lemme \ref{injectives-lemma-K-injective-embedding-filtration}. La suite spectrale est celle d'Homologie, section \ref{homology-section-filtered-complex}, associée au complexe filtré $$
\Gamma(X, \mathcal{J}^\bullet)
\quad\text{avec}\quad
F^p\Gamma(X, \mathcal{J}^\bullet) = \Gamma(X, F^p\mathcal{J}^\bullet)
$$ Comme la cohomologie se calcule sur des représentants K-injectifs, la page $E_1$ est celle indiquée dans le lemme. Sous les hypothèses énoncées, la convergence et le caractère borné résultent d'Homologie, lemme \ref{homology-lemma-ss-converges-trivial}.
\end{proof}

\begin{remark}

\label{cohomology-remark-spectral-sequence-filtered-object}

Soit $(X, \mathcal{O}_X)$ un espace annelé. Soit $\mathcal{F}^\bullet$ un complexe filtré de $\mathcal{O}_X$-modules. Si $\mathcal{F}^\bullet$
est borné inférieurement et si, pour tout $n$, la filtration de $\mathcal{F}^n$
est finie, on peut construire la suite spectrale du lemme \ref{cohomology-lemma-spectral-sequence-filtered-object}
sans recourir à Injectifs, lemme
\ref{injectives-lemma-K-injective-embedding-filtration}. En effet, d'après Catégories dérivées, lemme
\ref{derived-lemma-right-resolution-by-filtered-injectives}
il y a un quasi-isomorphisme filtré
$i : \mathcal{F}^\bullet \to \mathcal{I}^\bullet$
de complexes filtrés, où $\mathcal{I}^\bullet$ est borné inférieurement, où la filtration de $\mathcal{I}^n$ est finie pour tout $n$ et où chaque $\text{gr}^p\mathcal{I}^n$ est un $\mathcal{O}_X$-module injectif. Nous prenons alors la suite spectrale associée à
$$
\Gamma(X, \mathcal{I}^\bullet)
\quad\text{avec}\quad
F^p\Gamma(X, \mathcal{I}^\bullet) = \Gamma(X, F^p\mathcal{I}^\bullet)
$$
Comme la cohomologie se calcule sur des complexes bornés inférieurement d'injectifs, la page $E_1$ est celle indiquée dans le lemme. Sous les hypothèses énoncées, la convergence et le caractère borné résultent d'Homologie, lemme \ref{homology-lemma-biregular-ss-converges}. Il s'agit en fait d'un cas particulier de la suite spectrale de Catégories dérivées, lemme \ref{derived-lemma-ss-filtered-derived}.

\end{remark}

\begin{example}

\label{cohomology-example-spectral-sequence}

Soit $(X, \mathcal{O}_X)$ un espace annelé. Soit $\mathcal{F}^\bullet$ un complexe de $\mathcal{O}_X$-modules. Nous pouvons appliquer le lemme \ref{cohomology-lemma-spectral-sequence-filtered-object}
avec $F^p\mathcal{F}^\bullet = \tau_{\leq -p}\mathcal{F}^\bullet$. (Si $\mathcal{F}^\bullet$ est borné inférieurement, nous pouvons utiliser la remarque \ref{cohomology-remark-spectral-sequence-filtered-object}.) Nous obtenons alors une suite spectrale
$$
E_1^{p, q} = H^{p + q}(X, H^{-p}(\mathcal{F}^\bullet)[p]) =
H^{2p + q}(X, H^{-p}(\mathcal{F}^\bullet))
$$
Après renumérotation $p = -j$ et $q = i + 2j$, on voit que, pour tout
$K \in D(\mathcal{O}_X)$, il existe une suite spectrale
$(E'_r, d'_r)_{r \geq 2}$ de modules bigradués, où
$d'_r$ est de bidegré $(r, -r + 1)$, et telle que
$$
(E'_2)^{i, j} = H^i(X, H^j(K))
$$
Si $K$ est borné inférieurement, par exemple, cette suite spectrale est bornée et converge vers $H^{i + j}(X, K)$. Dans le cas borné inférieurement, elle est un exemple de la seconde suite spectrale de Catégories dérivées, lemme \ref{derived-lemma-two-ss-complex-functor}
(construit en utilisant les résolutions de Cartan-Eilenberg).

\end{example}

\begin{example}

\label{cohomology-example-spectral-sequence-bis}

Soit $(X, \mathcal{O}_X)$ un espace annelé. Soit $\mathcal{F}^\bullet$ un complexe de $\mathcal{O}_X$-modules. Nous pouvons appliquer le lemme \ref{cohomology-lemma-spectral-sequence-filtered-object}
avec $F^p\mathcal{F}^\bullet = \sigma_{\geq p}\mathcal{F}^\bullet$. Nous obtenons alors une suite spectrale
$$
E_1^{p, q} = H^{p + q}(X, \mathcal{F}^p[-p]) = H^q(X, \mathcal{F}^p)
$$
Si $\mathcal{F}^\bullet$ est borné inférieurement, alors

\begin{enumerate}

\item Nous pouvons utiliser la remarque \ref{cohomology-remark-spectral-sequence-filtered-object}
pour construire cette suite spectrale,
\item la suite spectrale est bornée et converge vers
$H^{i + j}(X, \mathcal{F}^\bullet)$,
\item la suite spectrale est égale à la première suite spectrale de Catégories dérivées, lemme \ref{derived-lemma-two-ss-complex-functor}
(construit en utilisant les résolutions de Cartan-Eilenberg).

\end{enumerate}

\end{example}

\begin{lemma}

\label{cohomology-lemma-relative-spectral-sequence-filtered-object}
Soit $f : (X, \mathcal{O}_X) \to (Y, \mathcal{O}_Y)$ un morphisme d'espaces annelés, et soit $\mathcal{F}^\bullet$ un complexe filtré de $\mathcal{O}_X$-modules. Il existe une suite spectrale canonique $(E_r, \text{d}_r)_{r \geq 1}$ de $\mathcal{O}_Y$-modules bigradués, où $d_r$ est de bidegré $(r, -r + 1)$, et $$
E_1^{p, q} = R^{p + q}f_*\text{gr}^p\mathcal{F}^\bullet
$$ Si pour chaque $n$ nous avons $$
R^nf_*F^p\mathcal{F}^\bullet = 0 \text{ pour }p \gg 0
\quad\text{et}\quad
R^nf_*F^p\mathcal{F}^\bullet = R^nf_*\mathcal{F}^\bullet \text{ pour }p \ll 0
$$, alors la suite spectrale est bornée et converge vers $Rf_*\mathcal{F}^\bullet$.
\end{lemma}

\begin{proof}

La preuve est exactement la même que celle du lemme \ref{cohomology-lemma-spectral-sequence-filtered-object}.

\end{proof}





\section{Résolution de Godement}

\label{cohomology-section-godement}

\noindent
Une référence est \cite{Godement}.

\medskip\noindent

Soit $(X, \mathcal{O}_X)$ un espace annelé. Notons $X_{disc}$ l'espace topologique discret ayant les mêmes points que $X$, et $f : X_{disc} \to X$
l'application continue évidente. Posons $\mathcal{O}_{X_{disc}} = f^{-1}\mathcal{O}_X$. Alors $f : (X_{disc}, \mathcal{O}_{X_{disc}}) \to (X, \mathcal{O}_X)$
est un morphisme plat d'espaces annelés. Appliquons la construction
{\it duale} de Simplicial, section \ref{simplicial-section-standard}, à la paire de foncteurs adjoints $f^*, f_*$ sur les faisceaux de modules. On obtient ainsi un objet cosimplicial augmenté
$$
\xymatrix{
\text{id} \ar[r] &
f_*f^* \ar@<1ex>[r] \ar@<-1ex>[r] &
f_*f^*f_*f^* \ar@<0ex>[l] \ar@<-2ex>[r] \ar@<0ex>[r] \ar@<2ex>[r] &
f_*f^*f_*f^*f_*f^* \ar@<1ex>[l] \ar@<-1ex>[l]
}
$$
dans la catégorie des endofoncteurs de $\textit{Mod}(\mathcal{O}_X)$;
voir Simplicial, lemme \ref{simplicial-lemma-standard-simplicial}. De plus, l'augmentation
$$
\xymatrix{
f^* \ar[r] &
f^*f_*f^* \ar@<1ex>[r] \ar@<-1ex>[r] &
f^*f_*f^*f_*f^* \ar@<0ex>[l] \ar@<-2ex>[r] \ar@<0ex>[r] \ar@<2ex>[r] &
f^*f_*f^*f_*f^*f_*f^* \ar@<1ex>[l] \ar@<-1ex>[l]
}
$$
est une équivalence d'homotopie, voir Simplicial, lemme \ref{simplicial-lemma-standard-simplicial-homotopy}.

\begin{lemma}

\label{cohomology-lemma-godement-resolution}
Soit $(X, \mathcal{O}_X)$ un espace annelé. Pour tout faisceau de $\mathcal{O}_X$-modules $\mathcal{F}$, il existe une résolution $$
0 \to
\mathcal{F} \to
f_*f^*\mathcal{F} \to
f_*f^*f_*f^*\mathcal{F} \to
f_*f^*f_*f^*f_*f^*\mathcal{F} \to \ldots
$$ fonctorielle en $\mathcal{F}$, telle que chaque terme $f_*f^* \ldots f_*f^*\mathcal{F}$ est un $\mathcal{O}_X$-module flasque et que, pour tout $x \in X$, l'application $$
\mathcal{F}_x[0] \to \Big(
(f_*f^*\mathcal{F})_x \to
(f_*f^*f_*f^*\mathcal{F})_x \to
(f_*f^*f_*f^*f_*f^*\mathcal{F})_x \to \ldots
\Big)
$$ est une équivalence d'homotopie dans la catégorie des complexes de $\mathcal{O}_{X, x}$-modules.
\end{lemma}

\begin{proof}

Le complexe $f_*f^*\mathcal{F} \to  f_*f^*f_*f^*\mathcal{F} \to
f_*f^*f_*f^*f_*f^*\mathcal{F} \to \ldots$ est le complexe associé à l'objet cosimplicial avec les termes
$f_*f^*\mathcal{F}, f_*f^*f_*f^*\mathcal{F},
f_*f^*f_*f^*f_*f^*\mathcal{F}, \ldots$ décrit ci-dessus; voir Simplicial, section \ref{simplicial-section-dold-kan-cosimplicial}. L'augmentation fournit l'application $\mathcal{F} \to f_*f^*\mathcal{F}$
indiquée. Pour tout faisceau abélien $\mathcal{H}$ sur $X_{disc}$, l'image directe $f_*\mathcal{H}$ est flasque, car $X_{disc}$
est discret et l'image directe d'un faisceau flasque est flasque. La même observation vaut pour les $\mathcal{O}_X$-modules.

\medskip\noindent

Si $x \in X_{disc} = X$ est un point, alors $(f^*\mathcal{G})_x = \mathcal{G}_x$
pour tout $\mathcal{O}_X$-module $\mathcal{G}$. Ainsi, $f^*$ est un foncteur exact, et un complexe de $\mathcal{O}_X$-modules
$\mathcal{G}_1 \to \mathcal{G}_2 \to \mathcal{G}_3$
est exact si et seulement si
$f^*\mathcal{G}_1 \to f^*\mathcal{G}_2 \to f^*\mathcal{G}_3$
est exact (voir Modules, lemme \ref{modules-lemma-abelian}). Le résultat mentionné au début de cette section montre que l'image inverse par $f^*$ fournit une équivalence d'homotopie de l'objet cosimplicial constant $f^*\mathcal{F}$ vers l'objet cosimplicial dont les termes sont
$f_*f^*\mathcal{F}, f_*f^*f_*f^*\mathcal{F},
f_*f^*f_*f^*f_*f^*\mathcal{F}, \ldots$. Par Simplicial, lemme \ref{simplicial-lemma-homotopy-equivalence-s-Q}
nous obtenons que
$$
f^*\mathcal{F}[0] \to \Big(
f^*f_*f^*\mathcal{F} \to
f^*f_*f^*f_*f^*\mathcal{F} \to
f^*f_*f^*f_*f^*f_*f^*\mathcal{F} \to \ldots
\Big)
$$
est une équivalence d'homotopie. Les deux dernières assertions du lemme en résultent immédiatement.

\end{proof}

\begin{lemma}

\label{cohomology-lemma-godement-resolution-bounded-below}
Soit $(X, \mathcal{O}_X)$ un espace annelé. Soit $\mathcal{F}^\bullet$ un complexe borné inférieurement de $\mathcal{O}_X$-modules. Il existe un quasi-isomorphisme $\mathcal{F}^\bullet \to \mathcal{G}^\bullet$, où $\mathcal{G}^\bullet$ est un complexe borné inférieurement de $\mathcal{O}_X$-modules flasques et où, pour tout $x \in X$, l'application $\mathcal{F}^\bullet_x \to \mathcal{G}^\bullet_x$ est une équivalence d'homotopie dans la catégorie des complexes de $\mathcal{O}_{X, x}$-modules.
\end{lemma}

\begin{proof}

Soit $\mathcal{A}$ la catégorie de complexes de $\mathcal{O}_X$-modules et soit $\mathcal{B}$ la catégorie de complexes de $\mathcal{O}_X$-modules. Nous pouvons alors appliquer la discussion ci-dessus aux foncteurs
$f^*$ et $f_*$ entre $\mathcal{A}$ et $\mathcal{B}$. En raisonnant exactement comme dans la preuve du lemme \ref{cohomology-lemma-godement-resolution},
nous obtenons une résolution
$$
0 \to
\mathcal{F}^\bullet \to
f_*f^*\mathcal{F}^\bullet \to
f_*f^*f_*f^*\mathcal{F}^\bullet \to
f_*f^*f_*f^*f_*f^*\mathcal{F}^\bullet \to \ldots
$$
dans la catégorie abélienne $\mathcal{A}$, de telle sorte que chaque terme de chaque
$f_*f^*\ldots f_*f^*\mathcal{F}^\bullet$ soit un
$\mathcal{O}_X$-module flasque et que, pour tout $x \in X$, l'application
$$
\mathcal{F}^\bullet_x[0] \to \Big(
(f_*f^*\mathcal{F}^\bullet)_x \to
(f_*f^*f_*f^*\mathcal{F}^\bullet)_x \to
(f_*f^*f_*f^*f_*f^*\mathcal{F}^\bullet)_x \to \ldots
\Big)
$$
est une équivalence d'homotopie dans la catégorie des complexes de complexes de $\mathcal{O}_{X, x}$-modules. Comme un complexe de complexes est un bicomplexe, nous pouvons considérer l'application induite
$$
\mathcal{F}^\bullet \to
\mathcal{G}^\bullet =
\text{Tot}(
f_*f^*\mathcal{F}^\bullet \to
f_*f^*f_*f^*\mathcal{F}^\bullet \to
f_*f^*f_*f^*f_*f^*\mathcal{F}^\bullet \to \ldots
)
$$
Comme le complexe $\mathcal{F}^\bullet$ est borné inférieurement, il en est de même de $\mathcal{G}^\bullet$; en fait, chaque terme de $\mathcal{G}^\bullet$ est une somme directe finie de termes des complexes $f_*f^*\ldots f_*f^*\mathcal{F}^\bullet$
et est donc flasque. La dernière assertion du lemme résulte alors d'Homologie, lemme \ref{homology-lemma-homotopy-complex-complexes}. Comme cela montre en particulier que
$\mathcal{F}^\bullet \to \mathcal{G}^\bullet$
est un quasi-isomorphisme, la preuve est complète.

\end{proof}
\section{Cup-produit}

\label{cohomology-section-cup-product}

\noindent

Soit $(X, \mathcal{O}_X)$ un espace annelé. Soient $K, M$ des objets de $D(\mathcal{O}_X)$. Posons $A = \Gamma(X, \mathcal{O}_X)$. Le cup-produit (global) dans ce cadre est un morphisme
$$
\mu :
R\Gamma(X, K) \otimes_A^\mathbf{L} R\Gamma(X, M)
\longrightarrow
R\Gamma(X, K \otimes_{\mathcal{O}_X}^\mathbf{L} M)
$$
dans $D(A)$. Nous le définissons comme le cup-produit relatif associé au morphisme d'espaces annelés $f : (X, \mathcal{O}_X) \to (pt, A)$
de la remarque \ref{cohomology-remark-cup-product}, via $D(pt, A) = D(A)$. Ce morphisme définit en particulier les accouplements
$$
\cup :
H^i(X, K) \times H^j(X, M)
\longrightarrow
H^{i + j}(X, K \otimes_{\mathcal{O}_X}^\mathbf{L} M)
$$
Plus précisément, étant donnés $\xi \in H^i(X, K) = H^i(R\Gamma(X, K))$ et
$\eta \in H^j(X, M) = H^j(R\Gamma(X, M))$, nous pouvons d'abord les ``tensoriser'' pour obtenir un élément $\xi \otimes \eta$ de
$H^{i + j}(R\Gamma(X, K) \otimes_A^\mathbf{L} R\Gamma(X, M))$; voir Compléments sur l'algèbre, section \ref{more-algebra-section-products-tor}. Nous pouvons ensuite appliquer $\mu$ pour obtenir l'élément souhaité
$\xi \cup \eta = \mu(\xi \otimes \eta)$
de $H^{i + j}(X, K \otimes_{\mathcal{O}_X}^\mathbf{L} M)$.

\medskip\noindent
Voici une autre manière d'interpréter le cup-produit de $\xi$ et $\eta$. En effet, nous pouvons écrire $$
H^i(R\Gamma(X, K)) = \Hom_{D(\mathcal{O}_X)}(\mathcal{O}_X[-i], K)
\quad\text{et}\quad
H^j(R\Gamma(X, M)) = \Hom_{D(\mathcal{O}_X)}(\mathcal{O}_X[-j], M)
$$ puisque $\Hom(\mathcal{O}_X, -) = \Gamma(X, -)$. Ainsi, $\xi$ et $\eta$ sont la ``même'' chose que les morphismes $$
\tilde \xi : \mathcal{O}_X[-i] \to K
\quad\text{et}\quad
\tilde \eta : \mathcal{O}_X[-j] \to M
$$ En combinant ceci avec la fonctorialité du produit tensoriel dérivé, nous obtenons $$
\mathcal{O}_X[-i - j] =
\mathcal{O}_X[-i] \otimes_{\mathcal{O}_X}^\mathbf{L} \mathcal{O}_X[-j]
\xrightarrow{\tilde \xi \otimes \tilde \eta}
K \otimes_{\mathcal{O}_X}^\mathbf{L} M
$$ qui, par le même raisonnement que ci-dessus, est un élément de $H^{i + j}(X, K \otimes_{\mathcal{O}_X}^\mathbf{L} M)$.

\begin{lemma}

\label{cohomology-lemma-second-cup-equals-first}

Cette construction donne le cup-produit.

\end{lemma}

\begin{proof}

Avec $f : (X, \mathcal{O}_X) \to (pt, A)$ comme ci-dessus nous avons
$Rf_*(-) = R\Gamma(X, -)$ et notre morphisme $\mu$ est adjoint du morphisme
$$
Lf^*(Rf_*K \otimes_A^\mathbf{L} Rf_*M) =
Lf^*Rf_*K \otimes_{\mathcal{O}_X}^\mathbf{L} Lf^*Rf_*M
\xrightarrow{\epsilon_K \otimes \epsilon_M}
K \otimes_{\mathcal{O}_X}^\mathbf{L} M
$$
où $\epsilon$ est la counité de l'adjonction entre
$Lf^*$ et $Rf_*$. Si nous considérons $\xi$ et $\eta$ comme des morphismes $\xi : A[-i] \to R\Gamma(X, K)$
et $\eta : A[-j] \to R\Gamma(X, M)$, alors le tenseur $\xi \otimes \eta$ correspond au morphisme\footnote{Un signe est caché ici : l'égalité est définie par la composition
$$
A[-i - j] \to (A \otimes_A^\mathbf{L} A)[-i - j] \to
A[-i] \otimes_A^\mathbf{L} A[-j]
$$
où, dans la deuxième étape, nous utilisons l'identification de Compléments sur l'algèbre, point (\ref{more-algebra-item-shift-tensor}), qui fait en principe intervenir un signe. Dans ce cas, ce signe vaut toutefois $+1$ par convention; même s'il ne valait pas
$+1$, cela n'aurait aucune importance puisque nous avons utilisé le même signe dans l'identification
$\mathcal{O}_X[-i - j] =
\mathcal{O}_X[-i] \otimes_{\mathcal{O}_X}^\mathbf{L} \mathcal{O}_X[-j]$.}
$$
A[-i - j] = A[-i] \otimes_A^\mathbf{L} A[-j]
\xrightarrow{\xi \otimes \eta}
R\Gamma(X, K) \otimes_A^\mathbf{L} R\Gamma(X, M)
$$
Par définition, le cup-produit $\xi \cup \eta$ est le morphisme
$A[-i - j] \to R\Gamma(X, K \otimes_{\mathcal{O}_X}^\mathbf{L} M)$
dont l'adjoint est
$$
(\epsilon_K \otimes \epsilon_M) \circ Lf^*(\xi \otimes \eta) =
(\epsilon_K \circ Lf^*\xi) \otimes (\epsilon_M \circ Lf^*\eta)
$$
Cependant, il est facile de voir que
$\epsilon_K \circ Lf^*\xi = \tilde \xi$ et
$\epsilon_M \circ Lf^*\eta = \tilde \eta$. Nous concluons que $\widetilde{\xi \cup \eta} = \tilde \xi \otimes \tilde \eta$,
ce qui donne l'identification souhaitée.

\end{proof}

\begin{remark}

\label{cohomology-remark-cup-with-element-map-total-cohomology}
Soit $(X, \mathcal{O}_X)$ un espace annelé. Soient $K, M$ des objets de $D(\mathcal{O}_X)$. Posons $A = \Gamma(X, \mathcal{O}_X)$. Pour $\xi \in H^i(X, K)$, nous obtenons le morphisme associé $$
\xi = ``\xi \cup -'' :
R\Gamma(X, M)[-i]
\to
R\Gamma(X, K \otimes_{\mathcal{O}_X}^\mathbf{L} M)
$$ en représentant $\xi$ par un morphisme $\xi : A[-i] \to R\Gamma(X, K)$ comme dans la preuve du lemme \ref{cohomology-lemma-second-cup-equals-first}, puis en utilisant la composition $$
R\Gamma(X, M)[-i] = A[-i] \otimes_A^\mathbf{L} R\Gamma(X, M)
\xrightarrow{\xi \otimes 1}
R\Gamma(X, K) \otimes_A^\mathbf{L} R\Gamma(X, M)
\to
R\Gamma(X, K \otimes_{\mathcal{O}_X}^\mathbf{L} M)
$$ où la deuxième flèche est le cup-produit global $\mu$ ci-dessus. En cohomologie, on retrouve ainsi le cup-produit par $\xi$, comme il ressort du lemme \ref{cohomology-lemma-second-cup-equals-first} et de sa preuve.
\end{remark}

\noindent
Formulons et démontrons une compatibilité naturelle du cup-produit relatif. Supposons donné un morphisme $f : (X, \mathcal{O}_X) \to (Y, \mathcal{O}_Y)$ d'espaces annelés. Soient $\mathcal{K}^\bullet$ et $\mathcal{M}^\bullet$ des complexes de $\mathcal{O}_X$-modules. Il existe un cup-produit naïf $$
\text{Tot}(
f_*\mathcal{K}^\bullet
\otimes_{\mathcal{O}_Y}
f_*\mathcal{M}^\bullet)
\longrightarrow
f_*\text{Tot}(\mathcal{K}^\bullet
\otimes_{\mathcal{O}_X}
\mathcal{M}^\bullet)
$$ Nous affirmons que celui-ci est relié au cup-produit relatif.

\begin{lemma}

\label{cohomology-lemma-cup-compatible-with-naive}
Dans la situation ci-dessus, le diagramme ci-dessous commute $$
\xymatrix{
f_*\mathcal{K}^\bullet
\otimes_{\mathcal{O}_Y}^\mathbf{L}
f_*\mathcal{M}^\bullet \ar[r] \ar[d]
&
Rf_*\mathcal{K}^\bullet
\otimes_{\mathcal{O}_Y}^\mathbf{L}
Rf_*\mathcal{M}^\bullet \ar[d]^{\text{Remarque \ref{cohomology-remark-cup-product}}} \\
\text{Tot}(
f_*\mathcal{K}^\bullet
\otimes_{\mathcal{O}_Y}
f_*\mathcal{M}^\bullet) \ar[d]_{\text{cup-produit naïf}} &
Rf_*(\mathcal{K}^\bullet
\otimes_{\mathcal{O}_X}^\mathbf{L}
\mathcal{M}^\bullet) \ar[d] \\
f_*\text{Tot}(\mathcal{K}^\bullet
\otimes_{\mathcal{O}_X}
\mathcal{M}^\bullet) \ar[r] &
Rf_*\text{Tot}(\mathcal{K}^\bullet
\otimes_{\mathcal{O}_X}
\mathcal{M}^\bullet)
}
$$
\end{lemma}

\begin{proof}

D'après la construction de la remarque \ref{cohomology-remark-cup-product}, le parcours du diagramme dans le sens horaire donne un morphisme
$$
f_*\mathcal{K}^\bullet
\otimes_{\mathcal{O}_Y}^\mathbf{L}
f_*\mathcal{M}^\bullet 
\longrightarrow
Rf_*\text{Tot}(\mathcal{K}^\bullet
\otimes_{\mathcal{O}_X}
\mathcal{M}^\bullet)
$$
dont l'adjoint est le morphisme

\begin{align*}
Lf^*(f_*\mathcal{K}^\bullet
\otimes_{\mathcal{O}_Y}^\mathbf{L}
f_*\mathcal{M}^\bullet)
& =
Lf^*f_*\mathcal{K}^\bullet
\otimes_{\mathcal{O}_Y}^\mathbf{L}
Lf^*f_*\mathcal{M}^\bullet \\
& \to
Lf^*Rf_*\mathcal{K}^\bullet
\otimes_{\mathcal{O}_Y}^\mathbf{L}
Lf^*Rf_*\mathcal{M}^\bullet \\
& \to
\mathcal{K}^\bullet
\otimes_{\mathcal{O}_Y}^\mathbf{L}
\mathcal{M}^\bullet \\
& \to
\text{Tot}(\mathcal{K}^\bullet
\otimes_{\mathcal{O}_X}
\mathcal{M}^\bullet)
\end{align*}

D'après le lemme \ref{cohomology-lemma-adjoints-push-pull-compatibility} c'est aussi égal à

\begin{align*}
Lf^*(f_*\mathcal{K}^\bullet
\otimes_{\mathcal{O}_Y}^\mathbf{L}
f_*\mathcal{M}^\bullet)
& =
Lf^*f_*\mathcal{K}^\bullet
\otimes_{\mathcal{O}_Y}^\mathbf{L}
Lf^*f_*\mathcal{M}^\bullet \\
& \to
f^*f_*\mathcal{K}^\bullet
\otimes_{\mathcal{O}_Y}^\mathbf{L}
f^*f_*\mathcal{M}^\bullet \\
& \to
\mathcal{K}^\bullet
\otimes_{\mathcal{O}_Y}^\mathbf{L}
\mathcal{M}^\bullet \\
& \to
\text{Tot}(\mathcal{K}^\bullet
\otimes_{\mathcal{O}_X}
\mathcal{M}^\bullet)
\end{align*}

En parcourant le diagramme dans le sens antihoraire, nous obtenons le morphisme adjoint du morphisme

\begin{align*}
Lf^*(f_*\mathcal{K}^\bullet
\otimes_{\mathcal{O}_Y}^\mathbf{L}
f_*\mathcal{M}^\bullet)
& \to
Lf^*\text{Tot}(
f_*\mathcal{K}^\bullet
\otimes_{\mathcal{O}_Y}
f_*\mathcal{M}^\bullet) \\
& \to
Lf^*f_*\text{Tot}(\mathcal{K}^\bullet
\otimes_{\mathcal{O}_X}
\mathcal{M}^\bullet) \\
& \to
Lf^*Rf_*\text{Tot}(\mathcal{K}^\bullet
\otimes_{\mathcal{O}_X}
\mathcal{M}^\bullet) \\
& \to
\text{Tot}(\mathcal{K}^\bullet
\otimes_{\mathcal{O}_X}
\mathcal{M}^\bullet)
\end{align*}

D'après le lemme \ref{cohomology-lemma-adjoints-push-pull-compatibility} c'est aussi égal à

\begin{align*}
Lf^*(f_*\mathcal{K}^\bullet
\otimes_{\mathcal{O}_Y}^\mathbf{L}
f_*\mathcal{M}^\bullet)
& \to
Lf^*\text{Tot}(
f_*\mathcal{K}^\bullet
\otimes_{\mathcal{O}_Y}
f_*\mathcal{M}^\bullet) \\
& \to
Lf^*f_*\text{Tot}(\mathcal{K}^\bullet
\otimes_{\mathcal{O}_X}
\mathcal{M}^\bullet) \\
& \to
f^*f_*\text{Tot}(\mathcal{K}^\bullet
\otimes_{\mathcal{O}_X}
\mathcal{M}^\bullet) \\
& \to
\text{Tot}(\mathcal{K}^\bullet
\otimes_{\mathcal{O}_X}
\mathcal{M}^\bullet)
\end{align*}

La démonstration s'achève alors par l'examen du diagramme
$$
\xymatrix{
Lf^*(f_*\mathcal{K}^\bullet
\otimes_{\mathcal{O}_Y}^\mathbf{L}
f_*\mathcal{M}^\bullet) \ar[d] \ar[rr] & &
Lf^*f_*\mathcal{K}^\bullet \otimes_{\mathcal{O}_X}^\mathbf{L}
Lf^*f_*\mathcal{M}^\bullet \ar[d] \\
Lf^*\text{Tot}(
f_*\mathcal{K}^\bullet
\otimes_{\mathcal{O}_Y}
f_*\mathcal{M}^\bullet) \ar[d]_{naive} \ar[r] &
f^*\text{Tot}(
f_*\mathcal{K}^\bullet
\otimes_{\mathcal{O}_Y}
f_*\mathcal{M}^\bullet) \ar[ldd]^{naive} \ar[dd] &
f^*f_*\mathcal{K}^\bullet \otimes_{\mathcal{O}_X}^\mathbf{L}
f^*f_*\mathcal{M}^\bullet \ar[dd] \ar[ldd] \\
Lf^*f_*\text{Tot}(\mathcal{K}^\bullet
\otimes_{\mathcal{O}_X}
\mathcal{M}^\bullet) \ar[d] \\
f^*f_*\text{Tot}(\mathcal{K}^\bullet \otimes_{\mathcal{O}_X}
\mathcal{M}^\bullet) \ar[rd] &
\text{Tot}(f^*f_*\mathcal{K}^\bullet \otimes_{\mathcal{O}_X}
f^*f_*\mathcal{M}^\bullet) \ar[d] &
\mathcal{K}^\bullet \otimes_{\mathcal{O}_X}^\mathbf{L}
\mathcal{M}^\bullet \ar[ld] \\
& \text{Tot}(\mathcal{K}^\bullet
\otimes_{\mathcal{O}_X}
\mathcal{M}^\bullet)
}
$$
Tous les polygones de ce diagramme commutent. Cela résulte du lemme \ref{cohomology-lemma-tensor-pull-compatibility}. Le carré comportant les deux cup-produits naïfs commute parce que
$Lf^* \to f^*$ est fonctoriel dans le complexe de modules. De même avec le carré impliquant les deux cartes
$\mathcal{A}^\bullet \otimes^\mathbf{L} \mathcal{B}^\bullet \to
\text{Tot}(\mathcal{A}^\bullet \otimes \mathcal{B}^\bullet)$. Enfin, la commutativité du reste du carré vaut au niveau des complexes et peut être considérée comme la définition du cup-produit naïf (par l'adjonction de $f^*$ et $f_*$). La preuve est achevée, car les deux morphismes décrits ci-dessus font le tour du bord extérieur du diagramme.

\end{proof}

\noindent
Soit $(X, \mathcal{O}_X)$ un espace annelé. Soient $\mathcal{K}^\bullet$ et $\mathcal{M}^\bullet$ des complexes de $\mathcal{O}_X$-modules. Nous avons alors un cup-produit naïf $$
\mu' :
\text{Tot}(
\Gamma(X, \mathcal{K}^\bullet) \otimes_A \Gamma(X, \mathcal{M}^\bullet))
\longrightarrow
\Gamma(X, \text{Tot}(
\mathcal{K}^\bullet \otimes_{\mathcal{O}_X} \mathcal{M}^\bullet))
$$ D'après le lemme \ref{cohomology-lemma-cup-compatible-with-naive}, appliqué au morphisme $(X, \mathcal{O}_X) \to (pt, A)$, ce cup-produit naïf est relié au cup-produit $\mu$ défini au premier paragraphe de cette section par le diagramme commutatif suivant $$
\xymatrix{
\Gamma(X, \mathcal{K}^\bullet)
\otimes_A^\mathbf{L}
\Gamma(X, \mathcal{M}^\bullet) \ar[d] \ar[r] &
R\Gamma(X, \mathcal{K}^\bullet)
\otimes_A^\mathbf{L}
R\Gamma(X, \mathcal{M}^\bullet) \ar[d]^-\mu \\
\text{Tot}(\Gamma(X, \mathcal{K}^\bullet)
\otimes_A
\Gamma(X, \mathcal{M}^\bullet)) \ar[d]_-{\mu'} &
R\Gamma(X, \mathcal{K}^\bullet
\otimes_{\mathcal{O}_X}^\mathbf{L}
\mathcal{M}^\bullet) \ar[d] \\
\Gamma(X, \text{Tot}(\mathcal{K}^\bullet
\otimes_{\mathcal{O}_X}
\mathcal{M}^\bullet)) \ar[r] &
R\Gamma(X, \text{Tot}(\mathcal{K}^\bullet
\otimes_{\mathcal{O}_X}
\mathcal{M}^\bullet))
}
$$ dans $D(A)$. Sur la cohomologie, nous obtenons le diagramme commutatif $$
\xymatrix{
H^i(\Gamma(X, \mathcal{K}^\bullet)) \times
H^j(\Gamma(X, \mathcal{M}^\bullet)) \ar[d] \ar[r] &
H^{i + j}(X,
\text{Tot}(\mathcal{K}^\bullet \otimes_{\mathcal{O}_X} \mathcal{M}^\bullet)) \\
H^i(X, \mathcal{K}^\bullet) \times
H^j(X, \mathcal{M}^\bullet) \ar[r]^\cup &
H^{i + j}(X, \mathcal{K}^\bullet \otimes_{\mathcal{O}_X}^\mathbf{L}
\mathcal{M}^\bullet) \ar[u]
}
$$ qui relie le cup-produit naïf au véritable cup-produit.

\begin{lemma}

\label{cohomology-lemma-diagrams-commute}
Soit $(X, \mathcal{O}_X)$ un espace annelé. Soient $\mathcal{K}^\bullet$ et $\mathcal{M}^\bullet$ des complexes bornés inférieurement de $\mathcal{O}_X$-modules. Soit $\mathcal{U} : X = \bigcup_{i \in I} U_i$ un recouvrement ouvert. Alors $$
\xymatrix{
\text{Tot}(\check{\mathcal{C}}^\bullet(\mathcal{U}, \mathcal{K}^\bullet))
\otimes_A^\mathbf{L}
\text{Tot}(\check{\mathcal{C}}^\bullet(\mathcal{U}, \mathcal{M}^\bullet))
\ar[d] \ar[r] &
R\Gamma(X, \mathcal{K}^\bullet)
\otimes_A^\mathbf{L}
R\Gamma(X, \mathcal{M}^\bullet) \ar[d]^\mu \\
\text{Tot}(
\text{Tot}(\check{\mathcal{C}}^\bullet(\mathcal{U}, \mathcal{K}^\bullet))
\otimes_A
\text{Tot}(\check{\mathcal{C}}^\bullet(\mathcal{U}, \mathcal{M}^\bullet)))
\ar[d]^{(\ref{cohomology-equation-needs-signs})} &
R\Gamma(X,
\mathcal{K}^\bullet \otimes_{\mathcal{O}_X}^\mathbf{L} \mathcal{M}^\bullet)
\ar[d] \\
\text{Tot}(
\check{\mathcal{C}}^\bullet({\mathcal U},
\text{Tot}(\mathcal{K}^\bullet \otimes_{\mathcal{O}_X} \mathcal{M}^\bullet)
)) \ar[r] &
R\Gamma(X,
\text{Tot}(\mathcal{K}^\bullet \otimes_{\mathcal{O}_X} \mathcal{M}^\bullet))
}
$$ où les flèches horizontales sont celles du lemme \ref{cohomology-lemma-cech-complex-complex}, commute dans $D(A)$.
\end{lemma}

\begin{proof}

Choisissez des quasi-isomorphismes de complexes
$a : \mathcal{K}^\bullet \to \mathcal{K}_1^\bullet$ et
$b : \mathcal{M}^\bullet \to \mathcal{M}_1^\bullet$
comme dans le lemme \ref{cohomology-lemma-godement-resolution-bounded-below}. Puisque les morphismes $a$ et $b$ induisent des équivalences d'homotopie sur les germes, nous voyons que le morphisme induit
$$
\text{Tot}(\mathcal{K}^\bullet \otimes_{\mathcal{O}_X} \mathcal{M}^\bullet)
\to
\text{Tot}(\mathcal{K}_1^\bullet \otimes_{\mathcal{O}_X} \mathcal{M}_1^\bullet)
$$
est lui aussi une équivalence d'homotopie sur les germes (Compléments sur l'algèbre, lemme
\ref{more-algebra-lemma-derived-tor-homotopy}) et donc un quasi-isomorphisme. Ainsi les cibles
$$
R\Gamma(X,
\text{Tot}(\mathcal{K}^\bullet
\otimes_{\mathcal{O}_X} \mathcal{M}^\bullet)) =
R\Gamma(X,
\text{Tot}(\mathcal{K}_1^\bullet
\otimes_{\mathcal{O}_X} \mathcal{M}_1^\bullet))
$$
des deux diagrammes coïncident dans $D(A)$. Il suffit donc de démontrer l'assertion après avoir remplacé $\mathcal{K}$ et $\mathcal{M}$
par $\mathcal{K}_1$ et $\mathcal{M}_1$. Nous sommes ainsi ramenés au cas traité dans le paragraphe suivant.

\medskip\noindent

Supposons que $\mathcal{K}^\bullet$ et $\mathcal{M}^\bullet$ soient des complexes bornés inférieurement de $\mathcal{O}_X$-modules flasques et considérons le diagramme qui relie le cup-produit au cup-produit (\ref{cohomology-equation-needs-signs}) sur les complexes de {\v C}ech. Nous pouvons alors considérer le diagramme commutatif
$$
\xymatrix{
\Gamma(X, \mathcal{K}^\bullet)
\otimes_A^\mathbf{L}
\Gamma(X, \mathcal{M}^\bullet) \ar[d] \ar[r] &
\text{Tot}(\check{\mathcal{C}}^\bullet(\mathcal{U}, \mathcal{K}^\bullet))
\otimes_A^\mathbf{L}
\text{Tot}(\check{\mathcal{C}}^\bullet(\mathcal{U}, \mathcal{M}^\bullet))
\ar[d] \\
\text{Tot}(\Gamma(X, \mathcal{K}^\bullet)
\otimes_A
\Gamma(X, \mathcal{M}^\bullet)) \ar[d] \ar[r] &
\text{Tot}(
\text{Tot}(\check{\mathcal{C}}^\bullet(\mathcal{U}, \mathcal{K}^\bullet))
\otimes_A
\text{Tot}(\check{\mathcal{C}}^\bullet(\mathcal{U}, \mathcal{M}^\bullet)))
\ar[d]^{(\ref{cohomology-equation-needs-signs})} \\
\Gamma(X, \text{Tot}(\mathcal{K}^\bullet
\otimes_{\mathcal{O}_X}
\mathcal{M}^\bullet)) \ar[r] &
\text{Tot}(
\check{\mathcal{C}}^\bullet({\mathcal U},
\text{Tot}(\mathcal{K}^\bullet \otimes_{\mathcal{O}_X} \mathcal{M}^\bullet)
))
}
$$
Dans ce diagramme, les flèches horizontales sont des isomorphismes dans $D(A)$, car, pour un complexe borné inférieurement de modules flasques tel que $\mathcal{K}^\bullet$,
nous avons
$$
\Gamma(X, \mathcal{K}^\bullet) =
\text{Tot}(\check{\mathcal{C}}^\bullet(\mathcal{U}, \mathcal{K}^\bullet)) =
R\Gamma(X, \mathcal{K}^\bullet)
$$
dans $D(A)$. Cela résulte du lemme \ref{cohomology-lemma-flasque-acyclic}, de Catégories dérivées, lemme \ref{derived-lemma-leray-acyclicity}, et du lemme \ref{cohomology-lemma-cech-complex-complex-computes}. La commutativité du diagramme du lemme faisant intervenir (\ref{cohomology-equation-needs-signs}) résulte donc de celle, déjà démontrée, du lemme \ref{cohomology-lemma-cup-compatible-with-naive},
où $f$ est le morphisme vers un point (voir la discussion qui suit \ref{cohomology-lemma-cup-compatible-with-naive}).

\end{proof}

\begin{lemma}

\label{cohomology-lemma-cup-product-associative}
Soit $f : (X, \mathcal{O}_X) \to (Y, \mathcal{O}_Y)$ un morphisme d'espaces annelés. Le cup-produit relatif de la remarque \ref{cohomology-remark-cup-product} est associatif au sens où le diagramme $$
\xymatrix{
Rf_*K \otimes_{\mathcal{O}_Y}^\mathbf{L}
Rf_*L \otimes_{\mathcal{O}_Y}^\mathbf{L}
Rf_*M \ar[r] \ar[d] &
Rf_*(K \otimes_{\mathcal{O}_X}^\mathbf{L} L)
\otimes_{\mathcal{O}_Y}^\mathbf{L} Rf_*M \ar[d] \\
Rf_*K \otimes_{\mathcal{O}_Y}^\mathbf{L}
Rf_*(L \otimes_{\mathcal{O}_X}^\mathbf{L} M) \ar[r] &
Rf_*(K \otimes_{\mathcal{O}_X}^\mathbf{L} 
L \otimes_{\mathcal{O}_X}^\mathbf{L} M)
}
$$ est commutatif dans $D(\mathcal{O}_Y)$ pour tous $K, L, M$ dans $D(\mathcal{O}_X)$.
\end{lemma}

\begin{proof}

En suivant chacun des deux chemins, nous obtenons le morphisme adjoint du morphisme évident

\begin{align*}
Lf^*(Rf_*K \otimes_{\mathcal{O}_Y}^\mathbf{L}
Rf_*L \otimes_{\mathcal{O}_Y}^\mathbf{L}
Rf_*M) & =
Lf^*(Rf_*K) \otimes_{\mathcal{O}_X}^\mathbf{L}
Lf^*(Rf_*L) \otimes_{\mathcal{O}_X}^\mathbf{L}
Lf^*(Rf_*M) \\
& \to
K \otimes_{\mathcal{O}_X}^\mathbf{L} 
L \otimes_{\mathcal{O}_X}^\mathbf{L} M
\end{align*}

en $D(\mathcal{O}_X)$.

\end{proof}

\begin{lemma}

\label{cohomology-lemma-cup-product-commutative}

Soit $f : (X, \mathcal{O}_X) \to (Y, \mathcal{O}_Y)$
un morphisme d'espaces annelés. Le cup-produit relatif de la remarque \ref{cohomology-remark-cup-product} est commutatif au sens où le diagramme
$$
\xymatrix{
Rf_*K \otimes_{\mathcal{O}_Y}^\mathbf{L} Rf_*L \ar[r] \ar[d]_\psi &
Rf_*(K \otimes_{\mathcal{O}_X}^\mathbf{L} L) \ar[d]^{Rf_*\psi} \\
Rf_*L \otimes_{\mathcal{O}_Y}^\mathbf{L} Rf_*K \ar[r] &
Rf_*(L \otimes_{\mathcal{O}_X}^\mathbf{L} K)
}
$$
est commutatif dans $D(\mathcal{O}_Y)$ pour tous $K, L$ dans $D(\mathcal{O}_X)$. Ici, $\psi$ est la contrainte de commutativité de la catégorie dérivée (lemme \ref{cohomology-lemma-symmetric-monoidal-derived}).

\end{lemma}

\begin{proof}
Démonstration omise.
\end{proof}

\begin{lemma}

\label{cohomology-lemma-compose-cup-product}
Soient $f : (X, \mathcal{O}_X) \to (Y, \mathcal{O}_Y)$ et $g : (Y, \mathcal{O}_Y) \to (Z, \mathcal{O}_Z)$ des morphismes d'espaces annelés. Le cup-produit relatif de la remarque \ref{cohomology-remark-cup-product} est compatible aux compositions au sens où le diagramme $$
\xymatrix{
R(g \circ f)_*K \otimes_{\mathcal{O}_Z}^\mathbf{L} R(g \circ f)_*L
\ar@{=}[rr] \ar[d] & &
Rg_*Rf_*K \otimes_{\mathcal{O}_Z}^\mathbf{L} Rg_*Rf_*L \ar[d] \\
R(g \circ f)_*(K \otimes_{\mathcal{O}_X}^\mathbf{L} L) \ar@{=}[r] &
Rg_*Rf_*(K \otimes_{\mathcal{O}_X}^\mathbf{L} L) &
Rg_*(Rf_*K \otimes_{\mathcal{O}_Y}^\mathbf{L}  Rf_*L) \ar[l]
}
$$ est commutatif en $D(\mathcal{O}_Z)$ pour tous $K, L$ en $D(\mathcal{O}_X)$.
\end{lemma}

\begin{proof}

En effet, chacun des deux parcours du diagramme donne le morphisme adjoint du morphisme

\begin{align*}
& L(g \circ f)^*\left(R(g \circ f)_*K
\otimes_{\mathcal{O}_Z}^\mathbf{L}
R(g \circ f)_*L\right) \\
& =
L(g \circ f)^*R(g \circ f)_*K
\otimes_{\mathcal{O}_X}^\mathbf{L}
L(g \circ f)^*R(g \circ f)_*L) \\
& \to
K \otimes_{\mathcal{O}_X}^\mathbf{L} L
\end{align*}

dans $D(\mathcal{O}_X)$. Pour le voir, on utilise que la composée des counités
$$
L(g \circ f)^*R(g \circ f)_* =
Lf^* Lg^* Rg_* Rf_*  \to
Lf^* Rf_* \to \text{id}
$$
est la counité pour $L(g \circ f)^*$ et $R(g \circ f)_*$. Voir Catégories, lemme \ref{categories-lemma-compose-counits}.

\end{proof}

\begin{lemma}

\label{cohomology-lemma-base-change-cup-product}
Considérez un carré commutatif $$
\xymatrix{
(X', \mathcal{O}_{X'}) \ar[r]_{g'} \ar[d]_{f'} &
(X, \mathcal{O}_X) \ar[d]^f \\
(Y', \mathcal{O}_{Y'}) \ar[r]^g & (Y, \mathcal{O}_Y) 
}
$$ d'espaces annelés. Soient $K, L$ dans $D(\mathcal{O}_X)$. Le cup-produit relatif est compatible avec ce carré au sens où le diagramme $$
\xymatrix{
Lg^*(Rf_*K \otimes_{\mathcal{O}_Y}^\mathbf{L} Rf_*L)
\ar[r] \ar@{=}[d] &
Lg^*(Rf_*(K \otimes_{\mathcal{O}_X}^\mathbf{L} L)) \ar[d] \\
Lg^*Rf_*K \otimes_{\mathcal{O}_{Y'}}^\mathbf{L} Lg^*Rf_*L \ar[d] &
R(f')_*L(g')^*(K \otimes_{\mathcal{O}_X}^\mathbf{L} L) \ar@{=}[d] \\
R(f')_*(L(g')^*K \otimes_{\mathcal{O}_{Y'}} R(f')_*(L(g')^*L \ar[r] &
R(f')_*(L(g')^*K \otimes_{\mathcal{O}_{X'}}^\mathbf{L} L(g')^*L)
}
$$ est commutatif dans $D(\mathcal{O}_{Y'})$. Les flèches horizontales sont données par le cup-produit relatif (remarque \ref{cohomology-remark-cup-product}) et les flèches verticales par le morphisme de changement de base (remarque \ref{cohomology-remark-base-change}) et le lemme \ref{cohomology-lemma-pullback-tensor-product}.
\end{lemma}

\begin{proof}
Démonstration omise.
\end{proof}













\section{Quelques propriétés des complexes K-injectifs}

\label{cohomology-section-properties-K-injective}

\noindent

Soit $(X, \mathcal{O}_X)$ un espace annelé. Soit $U \subset X$ un sous-ensemble ouvert. Notons $j : (U, \mathcal{O}_U) \to (X, \mathcal{O}_X)$
l'immersion ouverte correspondante. Le foncteur d'image inverse $j^*$ est exact, puisqu'il n'est autre que le foncteur de restriction. Ainsi, l'image inverse dérivée $Lj^*$ se calcule sur tout complexe en restreignant simplement ce complexe. Nous notons souvent simplement le foncteur correspondant
$$
D(\mathcal{O}_X) \to D(\mathcal{O}_U),
\quad
E \mapsto j^*E = E|_U
$$
De même, l'extension de zéro
$j_! : \textit{Mod}(\mathcal{O}_U) \to \textit{Mod}(\mathcal{O}_X)$
(voir Faisceaux, section \ref{sheaves-section-open-immersions}) est un foncteur exact (Modules, lemme \ref{modules-lemma-j-shriek-exact}). Il induit donc un foncteur
$$
j_! : D(\mathcal{O}_U) \to D(\mathcal{O}_X),\quad
F \mapsto j_!F
$$
en appliquant simplement $j_!$ à tout complexe représentant l'objet $F$.

\begin{lemma}

\label{cohomology-lemma-restrict-K-injective-to-open}

Soit $X$ un espace annelé. Soit $U \subset X$ un sous-espace ouvert. La restriction d'un complexe K-injectif de $\mathcal{O}_X$-modules à $U$ est un complexe K-injectif de $\mathcal{O}_U$-modules.

\end{lemma}

\begin{proof}
Cela résulte de Catégories dérivées, lemme \ref{derived-lemma-adjoint-preserve-K-injectives}, et du fait que le foncteur de restriction admet pour adjoint à gauche le foncteur exact $j_!$. Pour la construction de $j_!$, voir Faisceaux, section \ref{sheaves-section-open-immersions}, et, pour son exactitude, voir Modules, lemme \ref{modules-lemma-j-shriek-exact}.
\end{proof}

\begin{lemma}

\label{cohomology-lemma-unbounded-cohomology-of-open}
Soit $X$ un espace annelé. Soit $U \subset X$ un sous-espace ouvert. Pour $K$ dans $D(\mathcal{O}_X)$, nous avons $H^p(U, K) = H^p(U, K|_U)$.
\end{lemma}

\begin{proof}

Soit $\mathcal{I}^\bullet$ un complexe K-injectif de $\mathcal{O}_X$-modules représentant $K$. Alors
$$
H^q(U, K) = H^q(\Gamma(U, \mathcal{I}^\bullet)) =
H^q(\Gamma(U, \mathcal{I}^\bullet|_U))
$$
par construction de la cohomologie. D'après le lemme \ref{cohomology-lemma-restrict-K-injective-to-open},
le complexe $\mathcal{I}^\bullet|_U$ est un complexe K-injectif représentant $K|_U$ et le lemme suit.

\end{proof}

\begin{lemma}

\label{cohomology-lemma-sheafification-cohomology}
Soit $(X, \mathcal{O}_X)$ un espace annelé. Soit $K$ un objet de $D(\mathcal{O}_X)$. La faisceautisation de $$
U \mapsto H^q(U, K) = H^q(U, K|_U)
$$ est le $q$-ième faisceau de cohomologie $H^q(K)$ de $K$.
\end{lemma}

\begin{proof}

L'égalité $H^q(U, K) = H^q(U, K|_U)$ résulte du lemme \ref{cohomology-lemma-unbounded-cohomology-of-open}. Choisissons un complexe K-injectif $\mathcal{I}^\bullet$ représentant $K$. Alors
$$
H^q(U, K) =
\frac{\Ker(\mathcal{I}^q(U) \to \mathcal{I}^{q + 1}(U))}
{\Im(\mathcal{I}^{q - 1}(U) \to \mathcal{I}^q(U))}.
$$
par notre construction de la cohomologie. Comme
$H^q(K) = \Ker(\mathcal{I}^q \to \mathcal{I}^{q + 1})/
\Im(\mathcal{I}^{q - 1} \to \mathcal{I}^q)$, le résultat est clair.

\end{proof}

\begin{lemma}

\label{cohomology-lemma-restrict-direct-image-open}
Soit $f : (X, \mathcal{O}_X) \to (Y, \mathcal{O}_Y)$ un morphisme d'espaces annelés. Étant donné un sous-espace ouvert $V \subset Y$, posons $U = f^{-1}(V)$ et notons $g : U \to V$ le morphisme induit. Alors $(Rf_*E)|_V = Rg_*(E|_U)$ pour $E$ dans $D(\mathcal{O}_X)$.
\end{lemma}

\begin{proof}
Représentons $E$ par un complexe K-injectif $\mathcal{I}^\bullet$ de $\mathcal{O}_X$-modules. Alors $Rf_*(E) = f_*\mathcal{I}^\bullet$ et $Rg_*(E|_U) = g_*(\mathcal{I}^\bullet|_U)$ d'après le lemme \ref{cohomology-lemma-restrict-K-injective-to-open}. Comme il est clair que $(f_*\mathcal{F})|_V = g_*(\mathcal{F}|_U)$ pour tout faisceau $\mathcal{F}$ sur $X$, le résultat suit.
\end{proof}

\begin{lemma}

\label{cohomology-lemma-Leray-unbounded}
Soit $f : X \to Y$ un morphisme d'espaces annelés. Alors $R\Gamma(Y, -) \circ Rf_* = R\Gamma(X, -)$ comme foncteurs $D(\mathcal{O}_X) \to D(\Gamma(Y, \mathcal{O}_Y))$. Plus généralement, pour $V \subset Y$ ouvert et $U = f^{-1}(V)$, nous avons $R\Gamma(U, -) = R\Gamma(V, -) \circ Rf_*$.
\end{lemma}

\begin{proof}
Soit $Z$ l'espace annelé dont l'espace sous-jacent est réduit à un point et tel que $\Gamma(Z, \mathcal{O}_Z) = \Gamma(Y, \mathcal{O}_Y)$. Il existe un morphisme canonique $Y \to Z$ d'espaces annelés qui induit l'identification sur les sections globales des faisceaux structuraux. Alors $D(\mathcal{O}_Z) = D(\Gamma(Y, \mathcal{O}_Y))$. L'égalité $R\Gamma(Y, -) \circ Rf_* = R\Gamma(X, -)$ découle donc du lemme \ref{cohomology-lemma-derived-pushforward-composition} appliqué à $X \to Y \to Z$.

\medskip\noindent

Le second énoncé, plus général, découle du premier après application du lemme \ref{cohomology-lemma-restrict-direct-image-open}.

\end{proof}

\begin{lemma}

\label{cohomology-lemma-unbounded-describe-higher-direct-images}
Soit $f : (X, \mathcal{O}_X) \to (Y, \mathcal{O}_Y)$ un morphisme d'espaces annelés. Soit $K$ dans $D(\mathcal{O}_X)$. Alors $H^i(Rf_*K)$ est le faisceau associé au préfaisceau $$
V \mapsto H^i(f^{-1}(V), K) = H^i(V, Rf_*K)
$$
\end{lemma}

\begin{proof}
L'égalité $H^i(f^{-1}(V), K) = H^i(V, Rf_*K)$ s'obtient en prenant la cohomologie du second énoncé du lemme \ref{cohomology-lemma-Leray-unbounded}. L'énoncé sur la faisceautisation découle alors du lemme \ref{cohomology-lemma-sheafification-cohomology}.
\end{proof}

\begin{lemma}

\label{cohomology-lemma-modules-abelian-unbounded}
Soit $X$ un espace annelé. Soit $K$ un objet de $D(\mathcal{O}_X)$ et notons $K_{ab}$ son image dans $D(\underline{\mathbf{Z}}_X)$.
\begin{enumerate}
\item Pour tout $U \subset X$ ouvert, il y a une application canonique $R\Gamma(U, K) \to R\Gamma(U, K_{ab})$ qui est un isomorphisme dans $D(\textit{Ab})$.
\item Soit $f : X \to Y$ un morphisme d'espaces annelés. Il y a une application canonique $Rf_*K \to Rf_*(K_{ab})$ qui est un isomorphisme dans $D(\underline{\mathbf{Z}}_Y)$.
\end{enumerate}

\end{lemma}

\begin{proof}

Le morphisme est construit comme suit. Choisissons un complexe K-injectif
$\mathcal{I}^\bullet$ représentant $K$. Choisissez un quasi-isomorphisme
$\mathcal{I}^\bullet \to \mathcal{J}^\bullet$ où $\mathcal{J}^\bullet$
est un complexe K-injectif de groupes abéliens. Alors le morphisme de (1) est donné par $\Gamma(U, \mathcal{I}^\bullet) \to \Gamma(U, \mathcal{J}^\bullet)$
et celui de (2) par
$f_*\mathcal{I}^\bullet \to f_*\mathcal{J}^\bullet$. Pour montrer que ces morphismes sont des isomorphismes, il suffit de prouver qu'ils induisent des isomorphismes sur les groupes et les faisceaux de cohomologie. D'après les lemmes \ref{cohomology-lemma-unbounded-cohomology-of-open} et
\ref{cohomology-lemma-unbounded-describe-higher-direct-images}
il suffit de montrer que l'application
$$
H^0(X, K) \longrightarrow H^0(X, K_{ab})
$$
est un isomorphisme. Observez que
$$
H^0(X, K) = \Hom_{D(\mathcal{O}_X)}(\mathcal{O}_X, K)
$$
et de même pour l'autre groupe. Choisissons un complexe quelconque $\mathcal{K}^\bullet$
de $\mathcal{O}_X$-modules représentant $K$. Par la construction de la catégorie dérivée comme une localisation nous avons
$$
\Hom_{D(\mathcal{O}_X)}(\mathcal{O}_X, K) =
\colim_{s : \mathcal{F}^\bullet \to \mathcal{O}_X}
\Hom_{K(\mathcal{O}_X)}(\mathcal{F}^\bullet, \mathcal{K}^\bullet)
$$
où la colimite est prise sur les quasi-isomorphismes $s$ de complexes de
$\mathcal{O}_X$-modules. De même, nous avons
$$
\Hom_{D(\underline{\mathbf{Z}}_X)}(\underline{\mathbf{Z}}_X, K) =
\colim_{s : \mathcal{G}^\bullet \to \underline{\mathbf{Z}}_X}
\Hom_{K(\underline{\mathbf{Z}}_X)}(\mathcal{G}^\bullet, \mathcal{K}^\bullet)
$$
Ensuite, nous observons que les quasi-isomorphismes
$s : \mathcal{G}^\bullet \to \underline{\mathbf{Z}}_X$
avec $\mathcal{G}^\bullet$ complexe borné supérieurement de
$\underline{\mathbf{Z}}_X$-modules plats forment une famille cofinale dans le système. (Ceci découle de Modules, lemme \ref{modules-lemma-module-quotient-flat}, et de Catégories dérivées, lemme \ref{derived-lemma-subcategory-left-resolution}; voir la discussion de la section \ref{cohomology-section-flat}.) Nous pouvons donc construire un inverse au morphisme
$H^0(X, K) \longrightarrow H^0(X, K_{ab})$
en représentant un élément $\xi \in H^0(X, K_{ab})$ par une paire
$$
(s : \mathcal{G}^\bullet \to \underline{\mathbf{Z}}_X,
a : \mathcal{G}^\bullet \to \mathcal{K}^\bullet)
$$
où $\mathcal{G}^\bullet$ est un complexe borné supérieurement de
$\underline{\mathbf{Z}}_X$-modules plats, et envoyer cet élément sur
$$
(\mathcal{G}^\bullet \otimes_{\underline{\mathbf{Z}}_X} \mathcal{O}_X
\to \mathcal{O}_X,
\mathcal{G}^\bullet  \otimes_{\underline{\mathbf{Z}}_X} \mathcal{O}_X
\to \mathcal{K}^\bullet)
$$
Le seul point à noter est que la première flèche est un quasi-isomorphisme d'après les lemmes \ref{cohomology-lemma-derived-tor-quasi-isomorphism-other-side} et
\ref{cohomology-lemma-bounded-flat-K-flat}. Nous omettons la vérification détaillée que cette construction est effectivement un inverse.

\end{proof}

\begin{lemma}

\label{cohomology-lemma-adjoint-lower-shriek-restrict}
Soit $(X, \mathcal{O}_X)$ un espace annelé. Soit $U \subset X$ un sous-ensemble ouvert. Notons $j : (U, \mathcal{O}_U) \to (X, \mathcal{O}_X)$ l'immersion ouverte correspondante. Le foncteur de restriction $D(\mathcal{O}_X) \to D(\mathcal{O}_U)$ est adjoint à droite du foncteur d'extension par zéro $j_! : D(\mathcal{O}_U) \to D(\mathcal{O}_X)$.
\end{lemma}

\begin{proof}
Ceci résulte formellement du fait que $j_!$ et $j^*$ sont adjoints et exacts (et donc $Lj_! = j_!$ et $Rj^* = j^*$ existent), voir Catégories dérivées, lemme \ref{derived-lemma-derived-adjoint-functors}.
\end{proof}

\begin{lemma}

\label{cohomology-lemma-K-injective-flat}

Soit $f : X \to Y$ un morphisme plat d'espaces annelés. Si $\mathcal{I}^\bullet$ est un complexe K-injectif de $\mathcal{O}_X$-modules, alors $f_*\mathcal{I}^\bullet$ est K-injectif comme complexe de
$\mathcal{O}_Y$-modules.

\end{lemma}

\begin{proof}
Ceci est vrai parce que $$
\Hom_{K(\mathcal{O}_Y)}(\mathcal{F}^\bullet, f_*\mathcal{I}^\bullet)
=
\Hom_{K(\mathcal{O}_X)}(f^*\mathcal{F}^\bullet, \mathcal{I}^\bullet)
$$ par Faisceaux, lemme \ref{sheaves-lemma-adjoint-pullback-pushforward-modules} et le fait que $f^*$ est exact comme $f$ est supposé être plat.
\end{proof}





\section{Mayer--Vietoris non borné}

\label{cohomology-section-unbounded-mayer-vietoris}

\noindent
Il existe également une suite de Mayer--Vietoris pour la cohomologie non bornée.

\begin{lemma}

\label{cohomology-lemma-exact-sequence-lower-shriek}
Soit $(X, \mathcal{O}_X)$ un espace annelé. Soit $X = U \cup V$ l'union de deux sous-espaces ouverts. Pour tout objet $E$ de $D(\mathcal{O}_X)$ nous avons un triangle distingué $$
j_{U \cap V!}E|_{U \cap V} \to
j_{U!}E|_U \oplus j_{V!}E|_V \to E \to 
j_{U \cap V!}E|_{U \cap V}[1]
$$ dans $D(\mathcal{O}_X)$.
\end{lemma}

\begin{proof}

Nous avons vu dans la section \ref{cohomology-section-properties-K-injective}
que les foncteurs de restriction et d'extension par zéro se calculent en appliquant simplement les foncteurs correspondants à chaque complexe. Soit $\mathcal{E}^\bullet$ un complexe de $\mathcal{O}_X$-modules représentant $E$. Le triangle distingué du lemme est le triangle distingué associé (par Catégories dérivées, section
\ref{derived-section-canonical-delta-functor} et en particulier \ref{derived-lemma-derived-canonical-delta-functor}) à la suite exacte courte de complexes de $\mathcal{O}_X$-modules
$$
0 \to j_{U \cap V!}\mathcal{E}^\bullet|_{U \cap V} \to
j_{U!}\mathcal{E}^\bullet|_U \oplus j_{V!}\mathcal{E}^\bullet|_V
\to \mathcal{E}^\bullet \to 0
$$
L'exactitude de cette suite se vérifie sur les germes à l'aide de Faisceaux, lemme \ref{sheaves-lemma-j-shriek-modules}
(calcul omis).

\end{proof}

\begin{lemma}

\label{cohomology-lemma-exact-sequence-j-star}
Soit $(X, \mathcal{O}_X)$ un espace annelé. Soit $X = U \cup V$ l'union de deux sous-espaces ouverts. Pour tout objet $E$ de $D(\mathcal{O}_X)$ nous avons un triangle distingué $$
E \to 
Rj_{U, *}E|_U \oplus Rj_{V, *}E|_V \to
Rj_{U \cap V, *}E|_{U \cap V} \to
E[1]
$$ dans $D(\mathcal{O}_X)$.
\end{lemma}

\begin{proof}

Choisissons un complexe K-injectif $\mathcal{I}^\bullet$ représentant $E$
dont les termes $\mathcal{I}^n$ sont des objets injectifs de
$\textit{Mod}(\mathcal{O}_X)$; voir Injectifs, théorème
\ref{injectives-theorem-K-injective-embedding-grothendieck}. Nous avons vu que $\mathcal{I}^\bullet|U$ est également un complexe K-injectif (lemme \ref{cohomology-lemma-restrict-K-injective-to-open}). D'où :
$Rj_{U, *}E|_U$ est représenté par $j_{U, *}\mathcal{I}^\bullet|_U$. De même pour $V$ et $U \cap V$. Par conséquent, le triangle distingué du lemme est le triangle distingué associé (par Catégories dérivées, section
\ref{derived-section-canonical-delta-functor} et en particulier \ref{derived-lemma-derived-canonical-delta-functor}) à la suite exacte courte des complexes
$$
0 \to
\mathcal{I}^\bullet \to
j_{U, *}\mathcal{I}^\bullet|_U \oplus j_{V, *}\mathcal{I}^\bullet|_V \to
j_{U \cap V, *}\mathcal{I}^\bullet|_{U \cap V} \to
0.
$$
Cette suite est exacte parce que, pour tout ouvert $W \subset X$ et tout $n$, la suite
$$
0 \to
\mathcal{I}^n(W) \to
\mathcal{I}^n(W \cap U) \oplus \mathcal{I}^n(W \cap V) \to
\mathcal{I}^n(W \cap U \cap V) \to
0
$$
est exacte (voir la preuve du lemme \ref{cohomology-lemma-mayer-vietoris}).

\end{proof}

\begin{lemma}

\label{cohomology-lemma-mayer-vietoris-hom}
Soit $(X, \mathcal{O}_X)$ un espace annelé. Soit $X = U \cup V$ l'union de deux sous-espaces ouverts de $X$. Pour les objets $E$, $F$ de $D(\mathcal{O}_X)$ nous avons une suite de Mayer--Vietoris $$
\xymatrix{
& \ldots \ar[r] & \Ext^{-1}(E_{U \cap V}, F_{U \cap V}) \ar[lld] \\
\Hom(E, F) \ar[r] &
\Hom(E_U, F_U) \oplus
\Hom(E_V, F_V) \ar[r] &
\Hom(E_{U \cap V}, F_{U \cap V})
}
$$ où les indices désignent les restrictions aux ouverts concernés, et où les $\Hom$ et $\Ext$ sont pris dans les catégories dérivées correspondantes.
\end{lemma}

\begin{proof}
Appliquons le triangle distingué du lemme \ref{cohomology-lemma-exact-sequence-lower-shriek} pour obtenir une suite exacte longue de $\Hom$ (Catégories dérivées, lemme \ref{derived-lemma-representable-homological}) et utilisons l'égalité $$
\Hom_{D(\mathcal{O}_X)}(j_{U!}E|_U, F) =
\Hom_{D(\mathcal{O}_U)}(E|_U, F|_U)
$$ d'après le lemme \ref{cohomology-lemma-adjoint-lower-shriek-restrict}.
\end{proof}

\begin{lemma}

\label{cohomology-lemma-unbounded-mayer-vietoris}
Soit $(X, \mathcal{O}_X)$ un espace annelé. Supposons que $X = U \cup V$ soit la réunion de deux sous-ensembles ouverts. Pour un objet $E$ de $D(\mathcal{O}_X)$, nous avons un triangle distingué $$
R\Gamma(X, E) \to R\Gamma(U, E) \oplus R\Gamma(V, E) \to
R\Gamma(U \cap V, E) \to R\Gamma(X, E)[1]
$$ et, en particulier, une suite exacte longue de cohomologie $$
\ldots \to
H^n(X, E) \to
H^n(U, E) \oplus H^0(V, E) \to
H^n(U \cap V, E) \to
H^{n + 1}(X, E) \to \ldots
$$ La construction du triangle distingué et de la suite exacte longue est fonctorielle en $E$.
\end{lemma}

\begin{proof}

Choisissez un complexe K-injectif $\mathcal{I}^\bullet$
représentant $E$. Nous pouvons supposer que $\mathcal{I}^n$ est un objet injectif de $\textit{Mod}(\mathcal{O}_X)$ pour tout $n$; voir Injectifs, théorème
\ref{injectives-theorem-K-injective-embedding-grothendieck}. Alors $R\Gamma(X, E)$ est calculé par $\Gamma(X, \mathcal{I}^\bullet)$. De même pour $U$, $V$, $U \cap V$ d'après le lemme \ref{cohomology-lemma-restrict-K-injective-to-open}. Par conséquent, le triangle distingué du lemme est le triangle distingué associé (par Catégories dérivées, section
\ref{derived-section-canonical-delta-functor} et en particulier \ref{derived-lemma-derived-canonical-delta-functor}) à la suite exacte courte des complexes
$$
0 \to
\mathcal{I}^\bullet(X) \to
\mathcal{I}^\bullet(U) \oplus \mathcal{I}^\bullet(V) \to
\mathcal{I}^\bullet(U \cap V) \to
0.
$$
Nous avons vu qu'il s'agit d'une suite exacte courte dans la preuve du lemme \ref{cohomology-lemma-mayer-vietoris}. L'assertion finale découle de la fonctorielité de la construction dans Injectifs, théorème
\ref{injectives-theorem-K-injective-embedding-grothendieck}.

\end{proof}

\begin{lemma}

\label{cohomology-lemma-unbounded-relative-mayer-vietoris}
Soit $f : X \to Y$ un morphisme d'espaces annelés. Supposons que $X = U \cup V$ soit la réunion de deux sous-ensembles ouverts. Notons $a = f|_U : U \to Y$, $b = f|_V : V \to Y$ et $c = f|_{U \cap V} : U \cap V \to Y$. Pour chaque objet $E$ de $D(\mathcal{O}_X)$, il existe un triangle distingué $$
Rf_*E \to
Ra_*(E|_U) \oplus Rb_*(E|_V) \to
Rc_*(E|_{U \cap V}) \to
Rf_*E[1]
$$ Ce triangle est fonctoriel dans $E$.
\end{lemma}

\begin{proof}

Choisissez un complexe K-injectif $\mathcal{I}^\bullet$
représentant $E$. Nous pouvons supposer que $\mathcal{I}^n$ est un objet injectif de $\textit{Mod}(\mathcal{O}_X)$ pour tout $n$; voir Injectifs, théorème
\ref{injectives-theorem-K-injective-embedding-grothendieck}. Alors $Rf_*E$ est calculé par $f_*\mathcal{I}^\bullet$. De même pour $U$, $V$, $U \cap V$ d'après le lemme \ref{cohomology-lemma-restrict-K-injective-to-open}. Par conséquent, le triangle distingué du lemme est le triangle distingué associé (par Catégories dérivées, section
\ref{derived-section-canonical-delta-functor} et en particulier \ref{derived-lemma-derived-canonical-delta-functor}) à la suite exacte courte des complexes
$$
0 \to
f_*\mathcal{I}^\bullet \to
a_*\mathcal{I}^\bullet|_U \oplus b_*\mathcal{I}^\bullet|_V \to
c_*\mathcal{I}^\bullet|_{U \cap V} \to
0.
$$
C'est une suite exacte courte de complexes d'après le lemme \ref{cohomology-lemma-relative-mayer-vietoris}
et le fait que $R^1f_*\mathcal{I} = 0$
pour un objet injectif $\mathcal{I}$ de $\textit{Mod}(\mathcal{O}_X)$. L'assertion finale découle de la fonctorielité de la construction dans Injectifs, théorème
\ref{injectives-theorem-K-injective-embedding-grothendieck}.

\end{proof}

\begin{lemma}

\label{cohomology-lemma-pushforward-restriction}
Soit $(X, \mathcal{O}_X)$ un espace annelé. Soit $j : U \to X$ l'inclusion d'un sous-espace ouvert. Soit $T \subset X$ une partie fermée contenue dans $U$.
\begin{enumerate}
\item Si $E$ est un objet de $D(\mathcal{O}_X)$ dont les faisceaux de cohomologie sont à support dans $T$, alors $E \to Rj_*(E|_U)$ est un isomorphisme.
\item Si $F$ est un objet de $D(\mathcal{O}_U)$ dont les faisceaux de cohomologie sont à support dans $T$, alors $j_!F \to Rj_*F$ est un isomorphisme.
\end{enumerate}

\end{lemma}

\begin{proof}

Posons $V = X \setminus T$ et $W = U \cap V$. Remarquons que $X = U \cup V$ est un recouvrement ouvert de $X$. Notons $j_W : W \to V$ l'immersion ouverte. Soit $E$ un objet de $D(\mathcal{O}_X)$ dont les faisceaux de cohomologie sont à support dans $T$. D'après le lemme \ref{cohomology-lemma-restrict-direct-image-open}, nous avons
$(Rj_*E|_U)|_V = Rj_{W, *}(E|_W) = 0$, car $E|_W = 0$ par hypothèse. D'autre part, $Rj_*(E|_U)|_U = E|_U$. Ainsi, (1) est clair. Soit $F$ un objet de $D(\mathcal{O}_U)$ dont les faisceaux de cohomologie sont à support dans $T$. D'après le lemme \ref{cohomology-lemma-restrict-direct-image-open}, nous avons
$(Rj_*F)|_V = Rj_{W, *}(F|_W) = 0$, car $F|_W = 0$ par hypothèse. Nous avons aussi $(j_!F)|_V = j_{W!}(F|_W) = 0$ (la première égalité résulte immédiatement de la définition de l'extension par zéro).
Comme $(Rj_*F)|_U = F$ et $(j_!F)|_U = F$, nous voyons que (2) est vérifié.

\end{proof}

\begin{lemma}

\label{cohomology-lemma-mayer-vietoris-cup}

Soit $(X, \mathcal{O}_X)$ un espace annelé. Posons $A = \Gamma(X, \mathcal{O}_X)$. Supposons que $X = U \cup V$ soit la réunion de deux sous-ensembles ouverts.
Pour des objets $K$ et $M$ de $D(\mathcal{O}_X)$, nous avons un morphisme de triangles distingués
$$
\xymatrix{
R\Gamma(X, K) \otimes_A^\mathbf{L} R\Gamma(X, M) \ar[r] \ar[d] &
R\Gamma(X, K \otimes_{\mathcal{O}_X}^\mathbf{L} M) \ar[d] \\
R\Gamma(X, K) \otimes_A^\mathbf{L}
(R\Gamma(U, M) \oplus R\Gamma(V, M)) \ar[r] \ar[d] &
R\Gamma(U, K \otimes_{\mathcal{O}_X}^\mathbf{L} M)
\oplus R\Gamma(V, K \otimes_{\mathcal{O}_X}^\mathbf{L} M)) \ar[d] \\
R\Gamma(X, K) \otimes_A^\mathbf{L} R\Gamma(U \cap V, M) \ar[r] \ar[d] &
R\Gamma(U \cap V, K \otimes_{\mathcal{O}_X}^\mathbf{L} M) \ar[d] \\
R\Gamma(X, K) \otimes_A^\mathbf{L} R\Gamma(X, M)[1] \ar[r] &
R\Gamma(X, K \otimes_{\mathcal{O}_X}^\mathbf{L} M)[1]
}
$$
où

\begin{enumerate}

\item les flèches horizontales sont données par le cup-produit,
\item sur le côté droit nous avons le triangle distingué du lemme \ref{cohomology-lemma-unbounded-mayer-vietoris} pour
$K \otimes_{\mathcal{O}_X}^\mathbf{L} M$,
\item sur le côté gauche nous avons le foncteur exact
$R\Gamma(X, K) \otimes_A^\mathbf{L} - $ appliqué au triangle distingué du lemme \ref{cohomology-lemma-unbounded-mayer-vietoris} pour $M$.

\end{enumerate}

\end{lemma}

\begin{proof}

Choisissons un complexe K-plat $T^\bullet$ de $A$-modules plats représentant
$R\Gamma(X, K)$, voir Compléments sur l'algèbre, lemme \ref{more-algebra-lemma-K-flat-resolution}. Notons $T^\bullet \otimes_A \mathcal{O}_X$ l'image inverse de $T^\bullet$
par le morphisme d'espaces annelés $(X, \mathcal{O}_X) \to (pt, A)$. Il existe un morphisme d'adjonction naturel
$\epsilon : T^\bullet \otimes_A \mathcal{O}_X \to K$ dans $D(\mathcal{O}_X)$. Remarquons que $T^\bullet \otimes_A \mathcal{O}_X$ est un complexe K-plat de $\mathcal{O}_X$-modules à termes plats; voir le lemme \ref{cohomology-lemma-pullback-K-flat} et Modules, lemme \ref{modules-lemma-pullback-flat}. D'après le lemme \ref{cohomology-lemma-factor-through-K-flat}, nous pouvons trouver un morphisme de complexes
$$
T^\bullet \otimes_A \mathcal{O}_X \longrightarrow \mathcal{K}^\bullet
$$
de $\mathcal{O}_X$-modules représentant $\epsilon$
tel que $\mathcal{K}^\bullet$ soit un complexe K-plat à termes plats. En effet, par la construction de
$D(\mathcal{O}_X)$, nous pouvons d'abord représenter $\epsilon$ par un morphisme de complexes
$e : T^\bullet \otimes_A \mathcal{O}_X \to \mathcal{L}^\bullet$
de $\mathcal{O}_X$-modules représentant $\epsilon$
puis appliquer le lemme à $e$. Choisissons un complexe K-injectif $\mathcal{I}^\bullet$ dont les termes sont des $\mathcal{O}_X$-modules injectifs et qui représente $M$. Enfin, choisissons un quasi-isomorphisme
$$
\text{Tot}(\mathcal{K}^\bullet \otimes_\mathcal{O} \mathcal{I}^\bullet)
\longrightarrow
\mathcal{J}^\bullet
$$
vers un complexe K-injectif dont les termes sont des $\mathcal{O}_X$-modules injectifs. Remarquons que la source et la cible de cette flèche représentent
$K \otimes_{\mathcal{O}_X}^\mathbf{L} M$ dans $D(\mathcal{O}_X)$. À ce stade, pour tout ouvert $W \subset X$, nous obtenons un morphisme de complexes
$$
\text{Tot}(T^\bullet \otimes_A \mathcal{I}^\bullet(W))
\to
\text{Tot}(\mathcal{K}^\bullet(W) \otimes_A \mathcal{I}^\bullet(W))
\to
\mathcal{J}^\bullet(W)
$$
de $A$-modules dont la composition représente l'application
$$
R\Gamma(X, K) \otimes_A^\mathbf{L} R\Gamma(W, M)
\longrightarrow
R\Gamma(W, K \otimes_{\mathcal{O}_X}^\mathbf{L} M)
$$
dans $D(A)$. Ces morphismes sont manifestement compatibles aux applications de restriction. Nous pouvons donc considérer le diagramme commutatif suivant de complexes de $A$-modules
$$
\xymatrix{
0 \ar[d] & 0 \ar[d] \\
\text{Tot}(T^\bullet \otimes_A \mathcal{I}^\bullet(X)) \ar[d] \ar[r] &
\mathcal{J}^\bullet(X) \ar[d] \\
\text{Tot}(T^\bullet \otimes_A
(\mathcal{I}^\bullet(U) \oplus \mathcal{I}^\bullet(V)) \ar[d] \ar[r] &
\mathcal{J}^\bullet(U) \oplus \mathcal{J}^\bullet(V) \ar[d] \\
\text{Tot}(T^\bullet \otimes_A \mathcal{I}^\bullet(U \cap V)) \ar[r] \ar[d] &
\mathcal{J}^\bullet(U \cap V) \ar[d] \\
0 & 0
}
$$
D'après la preuve du lemme \ref{cohomology-lemma-mayer-vietoris}, les colonnes sont des suites exactes de complexes de $A$-modules (ce qui utilise aussi le fait que
$\text{Tot}(T^\bullet \otimes_A -)$ transforme les suites exactes courtes de complexes de $A$-modules en suites exactes courtes, puisque les termes de $T^\bullet$ sont des $A$-modules plats). Comme les triangles distingués du lemme \ref{cohomology-lemma-unbounded-mayer-vietoris}
sont les triangles distingués associés à ces suites exactes courtes de complexes, le résultat souhaité découle de la fonctorialité de la construction du triangle distingué associé, étudiée dans Catégories dérivées, section \ref{derived-section-canonical-delta-functor}.

\end{proof}











\section{Cohomologie à support dans une partie fermée, II}

\label{cohomology-section-cohomology-support-bis}

\noindent
Nous poursuivons la discussion commencée dans la section \ref{cohomology-section-cohomology-support}.

\medskip\noindent
Soit $(X, \mathcal{O}_X)$ un espace annelé. Soit $Z \subset X$ une partie fermée. Dans cette situation, nous pouvons considérer le foncteur $\textit{Mod}(\mathcal{O}_X) \to \textit{Mod}(\mathcal{O}_X(X))$ donné par $\mathcal{F} \mapsto \Gamma_Z(X, \mathcal{F})$. Voir Modules, définition \ref{modules-definition-support} et Modules, lemme \ref{modules-lemma-support-section-closed}. En utilisant des résolutions K-injectives (voir la section \ref{cohomology-section-unbounded}), nous obtenons le foncteur dérivé à droite $$
R\Gamma_Z(X, - ) : D(\mathcal{O}_X) \to D(\mathcal{O}_X(X))
$$ Pour un objet $K$ de $D(\mathcal{O}_X)$, nous notons $H^q_Z(X, K) = H^q(R\Gamma_Z(X, K))$ le module de cohomologie à support dans $Z$. Nous verrons plus loin (lemme \ref{cohomology-lemma-sections-support-abelian-unbounded}) que cette notation coïncide avec la construction de la section \ref{cohomology-section-cohomology-support}.

\medskip\noindent
Pour un $\mathcal{O}_X$-module $\mathcal{F}$, nous pouvons considérer le {\it sous-faisceau des sections à support dans $Z$}, noté $\mathcal{H}_Z(\mathcal{F})$, défini par la formule $$
\mathcal{H}_Z(\mathcal{F})(U) =
\{s \in \mathcal{F}(U) \mid \text{Supp}(s) \subset U \cap Z\} =
\Gamma_{Z \cap U}(U, \mathcal{F}|_U)
$$ Comme expliqué dans Modules, remarque \ref{modules-remark-sections-support-in-closed-modules}, nous pouvons voir $\mathcal{H}_Z(\mathcal{F})$ comme un $\mathcal{O}_X|_Z$-module sur $Z$, et nous obtenons un foncteur $$
\textit{Mod}(\mathcal{O}_X) \longrightarrow \textit{Mod}(\mathcal{O}_X|_Z),
\quad
\mathcal{F} \longmapsto \mathcal{H}_Z(\mathcal{F})
\text{ considéré comme }\mathcal{O}_X|_Z\text{-module sur }Z
$$ Ce foncteur est exact à gauche, mais n'est pas exact en général. Exactement comme ci-dessus, nous obtenons un foncteur dérivé à droite $$
R\mathcal{H}_Z : D(\mathcal{O}_X) \longrightarrow D(\mathcal{O}_X|_Z)
$$ Nous posons $\mathcal{H}^q_Z(K) = H^q(R\mathcal{H}_Z(K))$, de sorte que $\mathcal{H}^0_Z(\mathcal{F}) = \mathcal{H}_Z(\mathcal{F})$ pour tout faisceau de $\mathcal{O}_X$-modules $\mathcal{F}$.

\begin{lemma}

\label{cohomology-lemma-cohomology-with-support-sheaf-on-support}
Soit $(X, \mathcal{O}_X)$ un espace annelé. Soit $i : Z \to X$ l'inclusion d'une partie fermée.
\begin{enumerate}
\item $R\mathcal{H}_Z : D(\mathcal{O}_X) \to D(\mathcal{O}_X|_Z)$ est adjoint à droite de $i_* : D(\mathcal{O}_X|_Z) \to D(\mathcal{O}_X)$.
\item Pour $K$ dans $D(\mathcal{O}_X|_Z)$ nous avons $R\mathcal{H}_Z(i_*K) = K$.
\item Soit $\mathcal{G}$ un faisceau de $\mathcal{O}_X|_Z$-modules sur $Z$. Alors $\mathcal{H}^p_Z(i_*\mathcal{G}) = 0$ pour $p > 0$.
\end{enumerate}

\end{lemma}

\begin{proof}

Le foncteur $i_*$ est exact, donc $i_* = Ri_* = Li_*$. Par conséquent, la partie (1) du lemme découle de Modules, lemme \ref{modules-lemma-adjoint-section-with-support},
et de Catégories dérivées, lemme \ref{derived-lemma-derived-adjoint-functors}. Soit $K$ comme dans (2). Nous pouvons représenter $K$ par un complexe K-injectif
$\mathcal{I}^\bullet$ de $\mathcal{O}_X|_Z$-modules. D'après le lemme \ref{cohomology-lemma-K-injective-flat},
le complexe $i_*\mathcal{I}^\bullet$, qui représente $i_*K$, est un complexe K-injectif de $\mathcal{O}_X$-modules. Ainsi,
$R\mathcal{H}_Z(i_*K)$ est calculé par
$\mathcal{H}_Z(i_*\mathcal{I}^\bullet) = \mathcal{I}^\bullet$
ce qui prouve (2). La partie (3) est un cas particulier de (2).

\end{proof}

\noindent
Soit $(X, \mathcal{O}_X)$ un espace annelé et soit $Z \subset X$ une partie fermée. La catégorie des $\mathcal{O}_X$-modules dont le support est contenu dans $Z$ est une sous-catégorie de Serre de la catégorie de tous les $\mathcal{O}_X$-modules, voir Homologie, définition \ref{homology-definition-serre-subcategory} et Modules, lemme \ref{modules-lemma-support-section-closed}. Nous notons $D_Z(\mathcal{O}_X)$ la sous-catégorie triangulée strictement saturée de $D(\mathcal{O}_X)$ formée des complexes dont les faisceaux de cohomologie sont à support dans $Z$, voir Catégories dérivées, section \ref{derived-section-triangulated-sub}.

\begin{lemma}

\label{cohomology-lemma-complexes-with-support-on-closed}
Soit $(X, \mathcal{O}_X)$ un espace annelé. Soit $i : Z \to X$ l'inclusion d'une partie fermée.
\begin{enumerate}
\item Pour $K$ dans $D(\mathcal{O}_X|_Z)$ nous avons $i_*K$ dans $D_Z(\mathcal{O}_X)$.
\item Le foncteur $i_* : D(\mathcal{O}_X|_Z) \to D_Z(\mathcal{O}_X)$ est une équivalence, de quasi-inverse $i^{-1}|_{D_Z(\mathcal{O}_X)} = R\mathcal{H}_Z|_{D_Z(\mathcal{O}_X)}$.
\item Le foncteur $i_* \circ R\mathcal{H}_Z : D(\mathcal{O}_X) \to D_Z(\mathcal{O}_X)$ est adjoint à droite du foncteur d'inclusion $D_Z(\mathcal{O}_X) \to D(\mathcal{O}_X)$.
\end{enumerate}

\end{lemma}

\begin{proof}
La partie (1) résulte immédiatement des définitions. La partie (3) est une conséquence formelle de la partie (2) et du lemme \ref{cohomology-lemma-cohomology-with-support-sheaf-on-support}. Dans la suite de la preuve, nous démontrons la partie (2).

\medskip\noindent

Considérons $i$ comme le morphisme d'espaces annelés
$i : (Z, \mathcal{O}_X|_Z) \to (X, \mathcal{O}_X)$. Rappelons que $i^*$ et $i_*$ forment une paire de foncteurs adjoints. Puisque $i$ est une immersion fermée, $i_*$ est exact. Comme $i^{-1}\mathcal{O}_X = \mathcal{O}_X|_Z$ est le faisceau structural de $(Z, \mathcal{O}_X|_Z)$, nous voyons que $i^* = i^{-1}$
est exact et que $i^*i_* = i^{-1}i_*$
est isomorphe au foncteur identité. Voir Modules, lemmes \ref{modules-lemma-exactness-pushforward-pullback} et
\ref{modules-lemma-i-star-exact}. Ainsi
$i_* : D(\mathcal{O}_X|_Z) \to D_Z(\mathcal{O}_X)$
est pleinement fidèle et $i^{-1}$ fournit un inverse à gauche. D'autre part, supposons que $K$ soit un objet de
$D_Z(\mathcal{O}_X)$ et considérons le morphisme d'adjonction
$K \to i_*i^{-1}K$. L'exactitude de $i_*$ et de $i^{-1}$ montre que celui-ci induit les morphismes d'adjonction
$H^n(K) \to i_*i^{-1}H^n(K)$ sur les faisceaux de cohomologie. Comme ces faisceaux sont à support dans $Z$, ces morphismes d'adjonction sont des isomorphismes; nous concluons donc que $i_* : D(\mathcal{O}_X|_Z) \to D_Z(\mathcal{O}_X)$
est une équivalence.

\medskip\noindent

Pour terminer la preuve, il suffit de montrer que $R\mathcal{H}_Z(K) = i^{-1}K$ lorsque
$K$ est un objet de $D_Z(\mathcal{O}_X)$. Pour cela, nous pouvons utiliser l'égalité
$K = i_*i^{-1}K$, que nous venons de démontrer. Le lemme \ref{cohomology-lemma-cohomology-with-support-sheaf-on-support}
donne alors le résultat souhaité.

\end{proof}

\begin{lemma}

\label{cohomology-lemma-sections-with-support-K-injective}
Soit $(X, \mathcal{O}_X)$ un espace annelé. Soit $i : Z \to X$ l'inclusion d'une partie fermée. Si $\mathcal{I}^\bullet$ est un complexe K-injectif de $\mathcal{O}_X$-modules, alors $\mathcal{H}_Z(\mathcal{I}^\bullet)$ est un complexe K-injectif de $\mathcal{O}_X|_Z$-modules.
\end{lemma}

\begin{proof}
Puisque $i_* : \textit{Mod}(\mathcal{O}_X|_Z) \to \textit{Mod}(\mathcal{O}_X)$ est exact et adjoint à gauche de $\mathcal{H}_Z$ (Modules, lemme \ref{modules-lemma-adjoint-section-with-support}), ceci découle de Catégories dérivées, lemme \ref{derived-lemma-adjoint-preserve-K-injectives}.
\end{proof}

\begin{lemma}

\label{cohomology-lemma-local-to-global-sections-with-support}
Soit $(X, \mathcal{O}_X)$ un espace annelé. Soit $i : Z \to X$ l'inclusion d'une partie fermée. Alors $R\Gamma(Z, - ) \circ R\mathcal{H}_Z = R\Gamma_Z(X, - )$ comme foncteurs $D(\mathcal{O}_X) \to D(\mathcal{O}_X(X))$.
\end{lemma}

\begin{proof}
Cela résulte de la construction des foncteurs dérivés à droite au moyen de résolutions K-injectives, du lemme \ref{cohomology-lemma-sections-with-support-K-injective} et de l'égalité $\Gamma_Z(X, -) = \Gamma(Z, -) \circ \mathcal{H}_Z$.
\end{proof}

\begin{lemma}

\label{cohomology-lemma-triangle-sections-with-support}
Soit $(X, \mathcal{O}_X)$ un espace annelé. Soit $i : Z \to X$ l'inclusion d'une partie fermée. Posons $U = X \setminus Z$. Il existe un triangle distingué $$
R\Gamma_Z(X, K) \to R\Gamma(X, K) \to R\Gamma(U, K) \to
R\Gamma_Z(X, K)[1]
$$ dans $D(\mathcal{O}_X(X))$, fonctoriellement en $K$ dans $D(\mathcal{O}_X)$.
\end{lemma}

\begin{proof}

Choisissons un complexe K-injectif $\mathcal{I}^\bullet$ dont tous les termes sont des $\mathcal{O}_X$-modules injectifs et qui représente $K$. Voir la section \ref{cohomology-section-unbounded}. Rappelons que $\mathcal{I}^\bullet|_U$
est un complexe K-injectif de $\mathcal{O}_U$-modules; voir le lemme \ref{cohomology-lemma-restrict-K-injective-to-open}. Par conséquent, chacun des foncteurs dérivés du triangle distingué s'obtient en appliquant le foncteur sous-jacent à $\mathcal{I}^\bullet$. Il suffit donc de prouver que, pour tout $\mathcal{O}_X$-module injectif $\mathcal{I}$, nous avons une suite exacte courte
$$
0 \to \Gamma_Z(X, \mathcal{I}) \to \Gamma(X, \mathcal{I})
\to \Gamma(U, \mathcal{I}) \to 0
$$
Cela découle du lemme \ref{cohomology-lemma-injective-restriction-surjective}
et des définitions.

\end{proof}

\begin{lemma}

\label{cohomology-lemma-triangle-sections-with-support-sheaves}
Soit $(X, \mathcal{O}_X)$ un espace annelé. Soit $i : Z \to X$ l'inclusion d'une partie fermée. Notons $j : U = X \setminus Z \to X$ l'inclusion du complémentaire. Il existe un triangle distingué $$
i_*R\mathcal{H}_Z(K) \to K \to Rj_*(K|_U) \to
i_*R\mathcal{H}_Z(K)[1]
$$ dans $D(\mathcal{O}_X)$, fonctoriellement en $K$ dans $D(\mathcal{O}_X)$.
\end{lemma}

\begin{proof}

Choisissons un complexe K-injectif $\mathcal{I}^\bullet$ dont tous les termes sont des $\mathcal{O}_X$-modules injectifs et qui représente $K$. Voir la section \ref{cohomology-section-unbounded}. Rappelons que $\mathcal{I}^\bullet|_U$
est un complexe K-injectif de $\mathcal{O}_U$-modules; voir le lemme \ref{cohomology-lemma-restrict-K-injective-to-open}. Par conséquent, chacun des foncteurs dérivés du triangle distingué s'obtient en appliquant le foncteur sous-jacent à $\mathcal{I}^\bullet$. Il suffit donc de prouver que, pour tout $\mathcal{O}_X$-module injectif $\mathcal{I}$, nous avons une suite exacte courte
$$
0 \to i_*\mathcal{H}_Z(\mathcal{I}) \to \mathcal{I}
\to j_*(\mathcal{I}|_U) \to 0
$$
Cela découle du lemme \ref{cohomology-lemma-injective-restriction-surjective}
et des définitions.

\end{proof}

\begin{lemma}

\label{cohomology-lemma-sections-support-in-closed-disjoint-open}
Soit $(X, \mathcal{O}_X)$ un espace annelé. Soit $Z \subset X$ une partie fermée. Soit $j : U \to X$ l'inclusion d'un sous-ensemble ouvert tel que $U \cap Z = \emptyset$. Alors $R\mathcal{H}_Z(Rj_*K) = 0$ pour tout $K$ dans $D(\mathcal{O}_U)$.
\end{lemma}

\begin{proof}

Choisissons un complexe K-injectif $\mathcal{I}^\bullet$ de $\mathcal{O}_U$-modules représentant $K$. Alors $j_*\mathcal{I}^\bullet$ représente $Rj_*K$. D'après le lemme \ref{cohomology-lemma-K-injective-flat}, le complexe $j_*\mathcal{I}^\bullet$ est un complexe K-injectif de $\mathcal{O}_X$-modules. Ainsi,
$\mathcal{H}_Z(j_*\mathcal{I}^\bullet)$ représente $R\mathcal{H}_Z(Rj_*K)$. Il suffit donc de montrer que $\mathcal{H}_Z(j_*\mathcal{G}) = 0$
pour tout faisceau abélien $\mathcal{G}$ sur $U$. Il faut donc montrer qu'une section $s$ de $j_*\mathcal{G}$ sur un ouvert $W$, à support dans $W \cap Z$, est nulle. La condition sur le support signifie que
$s|_{W \setminus W \cap Z} = 0$. Comme $j_*\mathcal{G}(W) =
\mathcal{G}(U \cap W) = j_*\mathcal{G}(W \setminus W \cap Z)$
cela implique que $s$ est nulle, comme souhaité.

\end{proof}

\begin{lemma}

\label{cohomology-lemma-sections-support-abelian-unbounded}
Soit $(X, \mathcal{O}_X)$ un espace annelé. Soit $Z \subset X$ une partie fermée. Soit $K$ un objet de $D(\mathcal{O}_X)$ et notons $K_{ab}$ son image dans $D(\underline{\mathbf{Z}}_X)$.
\begin{enumerate}
\item Il y a une application canonique $R\Gamma_Z(X, K) \to R\Gamma_Z(X, K_{ab})$ qui est un isomorphisme dans $D(\textit{Ab})$.
\item Il y a une application canonique $R\mathcal{H}_Z(K) \to R\mathcal{H}_Z(K_{ab})$ qui est un isomorphisme dans $D(\underline{\mathbf{Z}}_Z)$.
\end{enumerate}

\end{lemma}

\begin{proof}

Démonstration de (1). Le morphisme est construit comme suit. Choisissons un complexe K-injectif de $\mathcal{O}_X$-modules $\mathcal{I}^\bullet$ représentant $K$. Choisissons un quasi-isomorphisme
$\mathcal{I}^\bullet \to \mathcal{J}^\bullet$ où $\mathcal{J}^\bullet$
est un complexe K-injectif de groupes abéliens. Alors le morphisme de (1) est donné par
$$
\Gamma_Z(X, \mathcal{I}^\bullet) \to \Gamma_Z(X, \mathcal{J}^\bullet)
$$
déterminé par le fait que $\Gamma_Z$ est un foncteur sur les faisceaux abéliens. Une vérification immédiate montre que le morphisme obtenu et les morphismes canoniques du lemme \ref{cohomology-lemma-modules-abelian-unbounded}
s'insèrent dans un morphisme de triangles distingués
$$
\xymatrix{
R\Gamma_Z(X, K) \ar[r] \ar[d] &
R\Gamma(X, K) \ar[r] \ar[d] &
R\Gamma(U, K) \ar[d] \\
R\Gamma_Z(X, K_{ab}) \ar[r] &
R\Gamma(X, K_{ab}) \ar[r] &
R\Gamma(U, K_{ab})
}
$$
du lemme \ref{cohomology-lemma-triangle-sections-with-support}. Comme deux des trois flèches sont des isomorphismes d'après le lemme cité, nous concluons par Catégories dérivées, lemme
\ref{derived-lemma-third-isomorphism-triangle}.

\medskip\noindent

La preuve de (2) est omise. Indication : utiliser le même argument avec le lemme \ref{cohomology-lemma-triangle-sections-with-support-sheaves}
pour le triangle distingué.

\end{proof}

\begin{remark}

\label{cohomology-remark-support-cup-product}

Soit $(X, \mathcal{O}_X)$ un espace annelé. Soit $i : Z \to X$
l'inclusion d'une partie fermée. Étant donnés $K$ et $M$ dans
$D(\mathcal{O}_X)$, il existe un morphisme canonique
$$
K|_Z \otimes_{\mathcal{O}_X|_Z}^\mathbf{L} R\mathcal{H}_Z(M)
\longrightarrow
R\mathcal{H}_Z(K \otimes_{\mathcal{O}_X}^\mathbf{L} M)
$$
dans $D(\mathcal{O}_X|_Z)$. Ici, $K|_Z = i^{-1}K$ est la restriction de
$K$ à $Z$, considérée comme un objet de $D(\mathcal{O}_X|_Z)$. Par l'adjonction entre $i_*$ et $R\mathcal{H}_Z$ du lemme \ref{cohomology-lemma-cohomology-with-support-sheaf-on-support}, il suffit, pour construire ce morphisme, de produire un morphisme canonique
$$
i_*\left(K|_Z \otimes_{\mathcal{O}_X|_Z}^\mathbf{L} R\mathcal{H}_Z(M)\right)
\longrightarrow
K \otimes_{\mathcal{O}_X}^\mathbf{L} M
$$
Pour construire ce morphisme, choisissons un complexe K-injectif $\mathcal{I}^\bullet$
de $\mathcal{O}_X$-modules représentant $M$ et un complexe K-plat
$\mathcal{K}^\bullet$ de $\mathcal{O}_X$-modules représentant $K$. Remarquons que $\mathcal{K}^\bullet|_Z$ est un complexe K-plat de
$\mathcal{O}_X|_Z$-modules représentant $K|_Z$; voir le lemme \ref{cohomology-lemma-pullback-K-flat}. Nous devons donc produire un morphisme de complexes
$$
i_*\text{Tot}\left(
\mathcal{K}^\bullet|_Z \otimes_{\mathcal{O}_X|_Z}
\mathcal{H}_Z(\mathcal{I}^\bullet)\right)
\longrightarrow
\text{Tot}(\mathcal{K}^\bullet \otimes_{\mathcal{O}_X} \mathcal{I}^\bullet)
$$
de $\mathcal{O}_X$-modules. Pour cela, il suffit de produire des morphismes
$$
i_*(\mathcal{K}^a|_Z \otimes_{\mathcal{O}_X|_Z}
\mathcal{H}_Z(\mathcal{I}^b))
\longrightarrow
\mathcal{K}^a \otimes_{\mathcal{O}_X} \mathcal{I}^b
$$
En regardant les germes (par exemple), nous voyons que le côté gauche de cette formule est égal à
$\mathcal{K}^a \otimes_{\mathcal{O}_X} i_*\mathcal{H}_Z(\mathcal{I}^b)$
et nous pouvons utiliser l'inclusion
$\mathcal{H}_Z(\mathcal{I}^b) \to \mathcal{I}^b$ pour obtenir notre morphisme.

\end{remark}

\begin{remark}

\label{cohomology-remark-support-cup-product-global}

Avec les notations de la remarque \ref{cohomology-remark-support-cup-product},
nous obtenons un cup-produit canonique

\begin{align*}
H^a(X, K) \times H^b_Z(X, M)
& =
H^a(X, K) \times H^b(Z, R\mathcal{H}_Z(M)) \\
& \to
H^a(Z, K|_Z) \times H^b(Z, R\mathcal{H}_Z(M)) \\
& \to
H^{a + b}(Z, K|_Z \otimes_{\mathcal{O}_X|_Z}^\mathbf{L} R\mathcal{H}_Z(M)) \\
& \to
H^{a + b}(Z, R\mathcal{H}_Z(K \otimes_{\mathcal{O}_X}^\mathbf{L} M)) \\
& =
H^{a + b}_Z(X, K \otimes_{\mathcal{O}_X}^\mathbf{L} M)
\end{align*}

Ici, les signes d'égalité sont donnés par le lemme \ref{cohomology-lemma-local-to-global-sections-with-support}; la première flèche est la restriction à $Z$, la deuxième est le cup-produit (section \ref{cohomology-section-cup-product}) et la troisième est le morphisme de la remarque \ref{cohomology-remark-support-cup-product}.

\end{remark}

\begin{lemma}

\label{cohomology-lemma-support-cup-product}
Avec les notations de la remarque \ref{cohomology-remark-support-cup-product}, le diagramme $$
\xymatrix{
H^i(X, K) \times H^j_Z(X, M) \ar[r] \ar[d] &
H^{i + j}_Z(X, K \otimes_{\mathcal{O}_X}^\mathbf{L} M) \ar[d] \\
H^i(X, K) \times H^j(X, M) \ar[r] &
H^{i + j}(X, K \otimes_{\mathcal{O}_X}^\mathbf{L} M)
}
$$ commute, la flèche horizontale supérieure étant le cup-produit de la remarque \ref{cohomology-remark-support-cup-product-global}.
\end{lemma}

\begin{proof}
Démonstration omise.
\end{proof}

\begin{remark}

\label{cohomology-remark-support-functorial}

Soit $f : (X', \mathcal{O}_{X'}) \to (X, \mathcal{O}_X)$ un morphisme d'espaces annelés. Soit $Z \subset X$ une partie fermée et posons $Z' = f^{-1}(Z)$. Notons $f|_{Z'} : (Z', \mathcal{O}_{X'}|_{Z'}) \to (Z, \mathcal{O}_X|Z)$ le morphisme induit d'espaces annelés. Pour tout $K$ dans $D(\mathcal{O}_X)$, il existe un morphisme canonique
$$
L(f|_{Z'})^*R\mathcal{H}_Z(K) \longrightarrow R\mathcal{H}_{Z'}(Lf^*K)
$$
dans $D(\mathcal{O}_{X'}|_{Z'})$. Notons $i : Z \to X$ et $i' : Z' \to X'$
les morphismes d'inclusion. D'après le lemme \ref{cohomology-lemma-complexes-with-support-on-closed}, partie (2), appliqué à $i'$, cela revient à donner un morphisme
$$
i'_* L(f|_{Z'})^* R\mathcal{H}_Z(K)
\longrightarrow
i'_*R\mathcal{H}_{Z'}(Lf^*K)
$$
dans $D_{Z'}(\mathcal{O}_{X'})$. Le morphisme de foncteurs
$Lf^* \circ i_* \to i'_* \circ L(f|_{Z'})^*$ de la remarque \ref{cohomology-remark-base-change} est un isomorphisme dans ce cas (cela se vérifie sur les germes en utilisant que $i_*$ et $i'_*$
sont exactes et $\mathcal{O}_{Z, z} = \mathcal{O}_{X, z}$
et de même pour $Z'$). Il suffit donc de construire la flèche horizontale supérieure dans le diagramme suivant
$$
\xymatrix{
Lf^* i_* R\mathcal{H}_Z(K) \ar[rr] \ar[rd] & &
i'_* R\mathcal{H}_{Z'}(Lf^*K) \ar[ld] \\
& Lf^*K
}
$$
Le complexe $Lf^* i_* R\mathcal{H}_Z(K)$ est à support dans $Z'$. La flèche sud-est provient du morphisme d'adjonction $i_*R\mathcal{H}_Z(K) \to K$
(lemme \ref{cohomology-lemma-cohomology-with-support-sheaf-on-support}). Puisque le morphisme d'adjonction $i'_* R\mathcal{H}_{Z'}(Lf^*K) \to Lf^*K$ est universel d'après le lemme \ref{cohomology-lemma-complexes-with-support-on-closed}, partie (3), la flèche sud-est se factorise de manière unique par la flèche sud-ouest, ce qui fournit la flèche souhaitée.

\end{remark}

\begin{lemma}

\label{cohomology-lemma-support-functorial}
Avec les notations et les hypothèses de la remarque \ref{cohomology-remark-support-functorial}, le diagramme $$
\xymatrix{
H^p_Z(X, K) \ar[r] \ar[d] & H^p_{Z'}(X', Lf^*K) \ar[d] \\
H^p(X, K) \ar[r] & H^p(X', Lf^*K)
}
$$ commute. Ici, la flèche horizontale supérieure provient des identifications $H^p_Z(X, K) = H^p(Z, R\mathcal{H}_Z(K))$ et $H^p_{Z'}(X', Lf^*K) = H^p(Z', R\mathcal{H}_{Z'}(K'))$, du morphisme d'image inverse $H^p(Z, R\mathcal{H}_Z(K)) \to H^p(Z', L(f|_{Z'})^*R\mathcal{H}_Z(K))$ et du morphisme construit dans la remarque \ref{cohomology-remark-support-functorial}.
\end{lemma}

\begin{proof}
Démonstration omise. Indication : en utilisant $H^p(Z, R\mathcal{H}_Z(K)) = H^p(X, i_*R\mathcal{H}_Z(K))$, et de même pour $R\mathcal{H}_{Z'}(Lf^*K)$, le résultat découle de la fonctorialité des morphismes d'image inverse et du diagramme commutatif utilisé pour définir le morphisme de la remarque \ref{cohomology-remark-support-functorial}.
\end{proof}





\section{Systèmes projectifs et cohomologie, I}

\label{cohomology-section-inverse-systems}

\noindent
Soit $A$ un anneau et soit $I \subset A$ un idéal. Nous démontrons quelques résultats sur des systèmes projectifs de faisceaux de $A/I^n$-modules.

\begin{lemma}

\label{cohomology-lemma-ML-general}

Soit $I$ un idéal d'un anneau $A$. Soit $X$ un espace topologique. Soit
$$
\ldots \to \mathcal{F}_3 \to \mathcal{F}_2 \to \mathcal{F}_1
$$
soit un système projectif de faisceaux de $A$-modules sur $X$
tel que $\mathcal{F}_n = \mathcal{F}_{n + 1}/I^n\mathcal{F}_{n + 1}$. Soit $p \geq 0$. Supposons que
$$
\bigoplus\nolimits_{n \geq 0} H^{p + 1}(X, I^n\mathcal{F}_{n + 1})
$$
satisfait à la condition de chaîne ascendante comme module gradué sur
$\bigoplus_{n \geq 0} I^n/I^{n + 1}$. Alors le système projectif $M_n = H^p(X, \mathcal{F}_n)$ satisfait à la condition de Mittag--Leffler\footnote{En fait, il existe un $c \geq 0$ tel que $\Im(M_n \to M_{n - c})$ soit l'image stable pour tout $n \geq c$.}.

\end{lemma}

\begin{proof}

Posons $N_n = H^{p + 1}(X, I^n\mathcal{F}_{n + 1})$ et soit
$\delta_n : M_n \to N_n$ l'application de bord en cohomologie provenant de la suite exacte courte
$0 \to I^n\mathcal{F}_{n + 1} \to \mathcal{F}_{n + 1} \to \mathcal{F}_n \to 0$. Alors $\bigoplus \Im(\delta_n) \subset \bigoplus N_n$ est un sous-module gradué. Soient $s \in M_n$ et $f \in I^m$; nous avons alors un diagramme commutatif
$$
\xymatrix{
0 \ar[r] &
I^n\mathcal{F}_{n + 1} \ar[d]_f \ar[r] &
\mathcal{F}_{n + 1} \ar[d]_f \ar[r] &
\mathcal{F}_n \ar[d]_f \ar[r] & 0 \\
0 \ar[r] &
I^{n + m}\mathcal{F}_{n + m + 1} \ar[r] &
\mathcal{F}_{n + m + 1} \ar[r] &
\mathcal{F}_{n + m} \ar[r] & 0
}
$$
Le morphisme vertical du milieu s'obtient en relevant une section locale de
$\mathcal{F}_{n + 1}$ en une section de $\mathcal{F}_{n + m + 1}$,
puis en la multipliant par $f$; il en va de même pour les autres flèches verticales. Nous en déduisons que $\delta_{n + m}(fs) = f \delta_n(s)$. Par hypothèse, nous pouvons trouver $s_j \in M_{n_j}$, $j = 1, \ldots, N$,
tels que les $\delta_{n_j}(s_j)$
engendrent $\bigoplus \Im(\delta_n)$ comme module gradué. Soit $n > c = \max(n_j)$ et soit $s \in M_n$. Nous pouvons alors trouver $f_j \in I^{n - n_j}$ tels que
$\delta_n(s) = \sum f_j \delta_{n_j}(s_j)$. Nous concluons que
$\delta(s - \sum f_j s_j) = 0$, c'est-à-dire, nous pouvons trouver $s' \in M_{n + 1}$
s'envoyant sur $s - \sum f_js_j$ dans $M_n$. Il s'ensuit que
$$
\Im(M_{n + 1} \to M_{n - c}) = \Im(M_n \to M_{n - c})
$$
En effet, les éléments $f_js_j$ s'envoient sur zéro dans $M_{n - c}$. Cela prouve le lemme.

\end{proof}

\begin{lemma}

\label{cohomology-lemma-ML-general-better}

Soit $I$ un idéal d'un anneau $A$. Soit $X$ un espace topologique. Soit
$$
\ldots \to \mathcal{F}_3 \to \mathcal{F}_2 \to \mathcal{F}_1
$$
un système projectif de faisceaux de $A$-modules sur $X$
tel que $\mathcal{F}_n = \mathcal{F}_{n + 1}/I^n\mathcal{F}_{n + 1}$. Soit $p \geq 0$. Pour $n$, posons
$$
N_n =
\bigcap\nolimits_{m \geq n}
\Im\left(
H^{p + 1}(X, I^n\mathcal{F}_{m + 1}) \to H^{p + 1}(X, I^n\mathcal{F}_{n + 1})
\right)
$$
Si $\bigoplus N_n$ satisfait à la condition de chaîne ascendante comme module gradué sur
$\bigoplus_{n \geq 0} I^n/I^{n + 1}$, alors le système projectif
$M_n = H^p(X, \mathcal{F}_n)$ satisfait à la condition de Mittag--Leffler\footnote{En fait, il existe un $c \geq 0$ tel que $\Im(M_n \to M_{n - c})$ soit l'image stable pour tout $n \geq c$.}.

\end{lemma}

\begin{proof}
La preuve est exactement la même que celle du lemme \ref{cohomology-lemma-ML-general}. En effet, les arguments donnés là prouvent le résultat dès que l'on montre que $\bigoplus N_n$ est un sous-module gradué sur $\bigoplus_{n \geq 0} I^n/I^{n + 1}$ de $\bigoplus H^{p + 1}(X, I^n\mathcal{F}_{n + 1})$ et que les applications de bord $\delta_n : M_n \to H^{p + 1}(X, I^n\mathcal{F}_{n + 1})$ ont leur image contenue dans $N_n$.

\medskip\noindent
Supposons que $\xi \in N_n$ et $f \in I^k$. Choisissons $m \gg n + k$, puis $\xi' \in H^{p + 1}(X, I^n\mathcal{F}_{m + 1})$ relevant $\xi$. Considérons le diagramme $$
\xymatrix{
0 \ar[r] &
I^n\mathcal{F}_{m + 1} \ar[d]_f \ar[r] &
\mathcal{F}_{m + 1} \ar[d]_f \ar[r] &
\mathcal{F}_n \ar[d]_f \ar[r] & 0 \\
0 \ar[r] &
I^{n + k}\mathcal{F}_{m + 1} \ar[r] &
\mathcal{F}_{m + 1} \ar[r] &
\mathcal{F}_{n + k} \ar[r] & 0
}
$$ construit comme dans la preuve du lemme \ref{cohomology-lemma-ML-general}. Nous obtenons un morphisme induit en cohomologie et nous voyons que $f \xi' \in H^{p + 1}(X, I^{n + k}\mathcal{F}_{m + 1})$ s'envoie sur $f \xi$. Puisque ceci vaut pour tout $m \gg n + k$, nous voyons que $f\xi$ appartient à $N_{n + k}$, comme souhaité.

\medskip\noindent

Pour voir que les applications de bord $\delta_n$ ont leur image contenue dans $N_n$,
nous considérons les diagrammes
$$
\xymatrix{
0 \ar[r] &
I^n\mathcal{F}_{m + 1} \ar[d] \ar[r] &
\mathcal{F}_{m + 1} \ar[d] \ar[r] &
\mathcal{F}_n \ar[d] \ar[r] & 0 \\
0 \ar[r] &
I^n\mathcal{F}_{n + 1} \ar[r] &
\mathcal{F}_{n + 1} \ar[r] &
\mathcal{F}_n \ar[r] & 0
}
$$
pour $m \geq n$. Le résultat s'obtient en examinant les morphismes induits en cohomologie.

\end{proof}

\begin{lemma}

\label{cohomology-lemma-topology-I-adic-general}
Soit $I$ un idéal d'un anneau $A$. Soit $X$ un espace topologique. Soit $$
\ldots \to \mathcal{F}_3 \to \mathcal{F}_2 \to \mathcal{F}_1
$$ un système projectif de faisceaux de $A$-modules sur $X$ tel que $\mathcal{F}_n = \mathcal{F}_{n + 1}/I^n\mathcal{F}_{n + 1}$. Soit $p \geq 0$. Supposons que $$
\bigoplus\nolimits_{n \geq 0} H^p(X, I^n\mathcal{F}_{n + 1})
$$ satisfasse à la condition de chaîne ascendante comme module gradué sur $\bigoplus_{n \geq 0} I^n/I^{n + 1}$. Alors la topologie limite sur $M = \lim H^p(X, \mathcal{F}_n)$ est la topologie $I$-adique.
\end{lemma}

\begin{proof}

Posons $F^n = \Ker(M \to H^p(X, \mathcal{F}_n))$ pour $n \geq 1$ et $F^0 = M$. Observons que $I F^n \subset F^{n + 1}$. En particulier, $I^n M \subset F^n$. Ainsi, la topologie $I$-adique est plus fine que la topologie limite. Réciproquement, nous montrerons que, pour tout $n$,
il existe un $m \geq n$ tel que $F^m \subset I^nM$.\footnote{En fait, il existe un $c \geq 0$ tel que $F^{n + c} \subset I^nM$ pour tout $n$.} Nous avons des applications injectives
$$
F^n/F^{n + 1} \longrightarrow H^p(X, \mathcal{F}_{n + 1})
$$
dont l'image est contenue dans
$H^p(X, I^n\mathcal{F}_{n + 1}) \to H^p(X, \mathcal{F}_{n + 1})$. Notons
$$
E_n \subset H^p(X, I^n\mathcal{F}_{n + 1})
$$
l'image inverse de $F^n/F^{n + 1}$. Alors $\bigoplus E_n$ est un sous-module gradué sur $\bigoplus I^n/I^{n + 1}$ de
$\bigoplus H^p(X, I^n\mathcal{F}_{n + 1})$ et
$\bigoplus E_n \to \bigoplus F^n/F^{n + 1}$ est un homomorphisme de modules gradués; nous omettons les détails. Par hypothèse, $\bigoplus E_n$ est engendré par un nombre fini d'éléments homogènes sur $\bigoplus I^n/I^{n + 1}$. Comme $E_n \to F^n/F^{n + 1}$ est surjectif, il en va de même pour $\bigoplus F^n/F^{n + 1}$. Nous pouvons donc trouver $r$, des entiers $c_1, \ldots, c_r \geq 0$ et
des éléments $a_i \in F^{c_i}$ dont les images engendrent $\bigoplus F^n/F^{n + 1}$. Posons $c = \max(c_i)$.

\medskip\noindent

Pour $n \geq c$, nous affirmons que $I F^n = F^{n + 1}$. Cette assertion montre que
$F^{n + c} = I^nF^c \subset I^nM$, comme souhaité. Pour la démontrer, supposons que $a \in F^{n + 1}$. L'image de
$a$ en $F^{n + 1}/F^{n + 2}$ est une combinaison linéaire de notre $a_i$. Par conséquent $a  - \sum f_i a_i \in F^{n + 2}$
pour certains $f_i \in I^{n + 1 - c_i}$. Comme
$I^{n + 1 - c_i} = I \cdot I^{n - c_i}$ et $n \geq c_i$, nous pouvons écrire
$f_i = \sum g_{i, j} h_{i, j}$ avec $g_{i, j} \in I$
et $h_{i, j}a_i \in F^n$. Ainsi, nous voyons que
$F^{n + 1} = F^{n + 2} + IF^n$. Un simple argument d'induction donne $F^{n + 1} = F^{n + e} + IF^n$
pour tout $e > 0$. Il s'ensuit que $IF^n$ est dense dans $F^{n + 1}$. Choisissons des générateurs $k_1, \ldots, k_r$ de $I$ et considérons l'application continue
$$
u : (F^n)^{\oplus r} \longrightarrow F^{n + 1},\quad
(x_1, \ldots, x_r) \mapsto \sum k_i x_i
$$
(pour la topologie limite). D'après ce qui précède, l'image de $(F^m)^{\oplus r}$ par $u$ est dense dans
$F^{m + 1}$ pour tout $m \geq n$. Le lemme de l'application ouverte (Compléments sur l'algèbre, lemme \ref{more-algebra-lemma-open-mapping}) montre que $u$ est ouverte. L'application $u$ est donc surjective. Par conséquent, $IF^n = F^{n + 1}$
pour $n \geq c$. Ceci conclut la preuve.

\end{proof}

\begin{lemma}

\label{cohomology-lemma-topology-I-adic-general-better}
Soit $I$ un idéal d'un anneau $A$. Soit $X$ un espace topologique. Soit $$
\ldots \to \mathcal{F}_3 \to \mathcal{F}_2 \to \mathcal{F}_1
$$ un système projectif de faisceaux de $A$-modules sur $X$ tel que $\mathcal{F}_n = \mathcal{F}_{n + 1}/I^n\mathcal{F}_{n + 1}$. Soit $p \geq 0$. Pour $n$, posons $$
N_n =
\bigcap\nolimits_{m \geq n}
\Im\left(
H^p(X, I^n\mathcal{F}_{m + 1}) \to H^p(X, I^n\mathcal{F}_{n + 1})
\right)
$$ Si $\bigoplus N_n$ satisfait à la condition de chaîne ascendante comme module gradué sur $\bigoplus_{n \geq 0} I^n/I^{n + 1}$, alors la topologie limite sur $M = \lim H^p(X, \mathcal{F}_n)$ est la topologie $I$-adique.
\end{lemma}

\begin{proof}
La preuve est exactement la même que celle du lemme \ref{cohomology-lemma-topology-I-adic-general}. En effet, les arguments donnés là prouvent le résultat dès que l'on montre que $\bigoplus N_n$ est un sous-module gradué sur $\bigoplus_{n \geq 0} I^n/I^{n + 1}$ de $\bigoplus H^{p + 1}(X, I^n\mathcal{F}_{n + 1})$ et que $F^n/F^{n + 1} \subset H^p(X, \mathcal{F}_{n + 1})$ est contenu dans l'image de $N_n \to H^p(X, \mathcal{F}_{n + 1})$. Dans la preuve du lemme \ref{cohomology-lemma-ML-general-better}, nous avons établi l'énoncé sur la structure de module.

\medskip\noindent

Soit $t \in F^n$. Choisissons un élément $s \in H^p(X, I^n\mathcal{F}_{n + 1})$
qui s'envoie sur l'image de $t$ dans $H^p(X, \mathcal{F}_{n + 1})$. Nous devons montrer que $s$ appartient à $N_n$. Or $F^n$ est le noyau de l'application
$M \to H^p(X, \mathcal{F}_n)$; par conséquent, pour tout $m \geq n$, nous pouvons envoyer $t$
vers un élément $t_m \in H^p(X, \mathcal{F}_{m + 1})$ qui s'envoie sur zéro dans $H^p(X, \mathcal{F}_n)$. Considérez la séquence de cohomologie
$$
H^{p - 1}(X, \mathcal{F}_n) \to
H^p(X, I^n\mathcal{F}_{m + 1}) \to
H^p(X, \mathcal{F}_{m + 1}) \to
H^p(X, \mathcal{F}_n)
$$
provenant de la suite exacte courte
$0 \to I^n\mathcal{F}_{m + 1} \to \mathcal{F}_{m + 1} \to \mathcal{F}_n \to 0$. Nous pouvons choisir $s_m \in H^p(X, I^n\mathcal{F}_{m + 1})$ s'envoyant sur $t_m$. En comparant la suite ci-dessus à celle correspondant à $m = n$, nous voyons que
$s_m$ s'envoie sur $s$, à un élément de l'image de
$H^{p - 1}(X, \mathcal{F}_n) \to H^p(X, I^n\mathcal{F}_{n + 1})$. Cependant, cette application se factorise par l'application
$H^p(X, I^n\mathcal{F}_{m + 1}) \to H^p(X, I^n\mathcal{F}_{n + 1})$
et nous voyons que $s$ appartient à l'image, comme souhaité.

\end{proof}





\section{Systèmes projectifs et cohomologie, II}

\label{cohomology-section-inverse-systems-bis}

\noindent
Cette section poursuit la discussion dans la section \ref{cohomology-section-inverse-systems} dans le cadre où l'idéal est principal.

\begin{lemma}

\label{cohomology-lemma-equivalent-f-good}

Soit $(X, \mathcal{O}_X)$ un espace annelé et soit $f \in \Gamma(X, \mathcal{O}_X)$. Soit
$$
\ldots \to \mathcal{F}_3 \to \mathcal{F}_2 \to \mathcal{F}_1
$$
un système projectif de $\mathcal{O}_X$-modules. Considérons les conditions

\begin{enumerate}

\item pour tout $n \geq 1$, le morphisme
$f : \mathcal{F}_{n + 1} \to \mathcal{F}_{n + 1}$ se factorise par $\mathcal{F}_{n + 1} \to \mathcal{F}_n$ de manière à donner une suite exacte courte
$0 \to \mathcal{F}_n \to \mathcal{F}_{n + 1} \to \mathcal{F}_1 \to 0$,
\item pour tout $n \geq 1$, le morphisme
$f^n : \mathcal{F}_{n + 1} \to \mathcal{F}_{n + 1}$
se factorise par $\mathcal{F}_{n + 1} \to \mathcal{F}_1$
de manière à donner une suite exacte courte
$0 \to \mathcal{F}_1 \to \mathcal{F}_{n + 1} \to \mathcal{F}_n \to 0$

\item il existe un $\mathcal{O}_X$-module $\mathcal{G}$
qui est $f$-divisible et tel que $\mathcal{F}_n = \mathcal{G}[f^n]$,
\item il existe un $\mathcal{O}_X$-module $\mathcal{F}$
sans $f$-torsion et tel que
$\mathcal{F}_n = \mathcal{F}/f^n\mathcal{F}$.

\end{enumerate}

Alors (4) $\Rightarrow$ (3) $\Leftrightarrow$ (2) $\Leftrightarrow$ (1).

\end{lemma}

\begin{proof}
Nous omettons la preuve de l'équivalence de (1) et (2), ainsi que celle du fait que (3) implique (1). Étant donné $\mathcal{F}_n$ comme dans (1), posons, pour démontrer (3),
$\mathcal{G} = \colim \mathcal{F}_n$, où les morphismes
$\mathcal{F}_1 \to \mathcal{F}_2 \to \mathcal{F}_3 \to \ldots$
sont ceux de (1). Le morphisme $f : \mathcal{G} \to \mathcal{G}$
est surjectif, car l'image de $\mathcal{F}_{n + 1} \subset \mathcal{G}$
est $\mathcal{F}_n \subset \mathcal{G}$ d'après la suite exacte courte de (1).
Ainsi, $\mathcal{G}$ est $f$-divisible comme $\mathcal{O}_X$-module et
$\mathcal{F}_n = \mathcal{G}[f^n]$.

\medskip\noindent
Supposons donné $\mathcal{F}$ comme dans (4). Le morphisme
$\mathcal{F}/f^{n + 1}\mathcal{F} \to \mathcal{F}/f^n\mathcal{F}$
est toujours surjectif, et son noyau est l'image du morphisme
$\mathcal{F}/f\mathcal{F} \to \mathcal{F}/f^{n + 1}\mathcal{F}$
induit par la multiplication par $f^n$. Pour vérifier (2), il suffit de voir que le noyau de
$f^n : \mathcal{F} \to \mathcal{F}/f^{n + 1}\mathcal{F}$
est $f\mathcal{F}$. Il suffit donc de montrer que, pour des sections
$s, t$ de $\mathcal{F}$ sur un ouvert $U \subset X$
telles que $f^ns = f^{n + 1}t$, nous avons $s = ft$. Cela est clair, car
$f : \mathcal{F} \to \mathcal{F}$ est injectif puisque $\mathcal{F}$
est sans $f$-torsion.
\end{proof}

\begin{lemma}

\label{cohomology-lemma-topology-I-adic-f}
Supposons que $X$, $f$ et $(\mathcal{F}_n)$ satisfassent à la condition (1) du lemme \ref{cohomology-lemma-equivalent-f-good}. Soit $p \geq 0$ et posons
$H^p = \lim H^p(X, \mathcal{F}_n)$. Alors $f^cH^p$ est le noyau de $H^p \to H^p(X, \mathcal{F}_c)$ pour tout
$c \geq 1$. Ainsi, la topologie limite sur $H^p$ est la topologie $f$-adique.
\end{lemma}

\begin{proof}

Soit $c \geq 1$. Il est clair que $f^c H^p$ s'envoie sur zéro dans
$H^p(X, \mathcal{F}_c)$. Si $\xi = (\xi_n) \in H^p$ est suffisamment petit pour la topologie limite, alors $\xi_c = 0$, et donc $\xi_n$
s'envoie sur zéro dans $H^p(X, \mathcal{F}_c)$ pour $n \geq c$. Considérons le système projectif de suites exactes courtes
$$
0 \to \mathcal{F}_{n - c} \xrightarrow{f^c} \mathcal{F}_n \to
\mathcal{F}_c \to 0
$$
et le système projectif correspondant de suites exactes longues de cohomologie
$$
H^{p - 1}(X, \mathcal{F}_c) \to
H^p(X, \mathcal{F}_{n - c}) \to
H^p(X, \mathcal{F}_n) \to
H^p(X, \mathcal{F}_c)
$$
Puisque le terme $H^{p - 1}(X, \mathcal{F}_c)$ est indépendant de
$n$, nous pouvons choisir une suite compatible d'éléments
$\xi'_n \in H^p(X, \mathcal{F}_{n - c})$
relevant $\xi_n$. En posant $\xi' = (\xi'_n)$, nous obtenons $\xi = f^c \xi'$,
comme souhaité.

\end{proof}

\begin{lemma}

\label{cohomology-lemma-limit-finite}

Soit $A$ un anneau noethérien complet pour un idéal principal $(f)$. Soit $X$ un espace topologique. Soit
$$
\ldots \to \mathcal{F}_3 \to \mathcal{F}_2 \to \mathcal{F}_1
$$
un système projectif de faisceaux de $A$-modules. Supposons que

\begin{enumerate}

\item $\Gamma(X, \mathcal{F}_1)$ est un $A$-module de type fini,
\item $X$, $f$ et $(\mathcal{F}_n)$ satisfont à la condition (1) du lemme \ref{cohomology-lemma-equivalent-f-good}.

\end{enumerate}

Alors
$$
M = \lim \Gamma(X, \mathcal{F}_n)
$$
est un $A$-module de type fini, $f$ est un non-diviseur de zéro sur $M$, et
$M/fM$ est l'image de $M$ dans $\Gamma(X, \mathcal{F}_1)$.

\end{lemma}

\begin{proof}

D'après le lemme \ref{cohomology-lemma-topology-I-adic-f}, nous avons
$M/fM \subset H^0(X, \mathcal{F}_1)$. D'après (1) et la noethérianité de $A$, nous obtenons que $M/fM$ est un $A$-module de type fini. De plus, $\bigcap f^nM = 0$, car $f^nM$ s'envoie sur zéro dans
$H^0(X, \mathcal{F}_n)$. Par Algèbre, lemme \ref{algebra-lemma-finite-over-complete-ring},
nous concluons que $M$ est de type fini sur $A$. Enfin, supposons que $s = (s_n) \in M = \lim H^0(X, \mathcal{F}_n)$
vérifie $fs = 0$. Alors $s_{n + 1}$ appartient au noyau de
$\mathcal{F}_{n + 1} \to \mathcal{F}_n$ d'après la condition (1) du lemme \ref{cohomology-lemma-equivalent-f-good}. Ainsi, $s_n = 0$. Comme $n$ est arbitraire, nous obtenons $s = 0$. Par conséquent, $f$ est un non-diviseur de zéro sur $M$.

\end{proof}

\begin{lemma}

\label{cohomology-lemma-ML}

Soit $A$ un anneau. Soit $f \in A$. Soit $X$ un espace topologique. Soit
$$
\ldots \to \mathcal{F}_3 \to \mathcal{F}_2 \to \mathcal{F}_1
$$
un système projectif de faisceaux de $A$-modules. Soit $p \geq 0$. Supposons que

\begin{enumerate}

\item soit $H^{p + 1}(X, \mathcal{F}_1)$ est un $A$-module de longueur finie, soit $A$ est noethérien et $H^{p + 1}(X, \mathcal{F}_1)$ est un $A$-module de type fini,
\item $X$, $f$ et $(\mathcal{F}_n)$ satisfont à la condition (1) du lemme \ref{cohomology-lemma-equivalent-f-good}.

\end{enumerate}

Alors le système projectif $M_n = H^p(X, \mathcal{F}_n)$ satisfait à la condition de Mittag--Leffler.

\end{lemma}

\begin{proof}

Posons $I = (f)$. Nous utiliserons le critère du lemme \ref{cohomology-lemma-ML-general}. Remarquons que
$f^n : \mathcal{F}_1 \to I^n\mathcal{F}_{n + 1}$
est un isomorphisme pour tout $n \geq 0$. Il suffit donc de montrer que
$$
\bigoplus\nolimits_{n \geq 1} H^{p + 1}(X, \mathcal{F}_1) \cdot f^{n + 1}
$$
est un module gradué sur $S = \bigoplus_{n \geq 0} A/(f) \cdot f^n$ satisfaisant à la condition de chaîne ascendante. Si $A$ n'est pas noethérien, alors
$H^{p + 1}(X, \mathcal{F}_1)$ est de longueur finie et le résultat est acquis. Si $A$ est noethérien, alors $S$ est un anneau noethérien et le résultat découle de ce que le module est de type fini sur $S$, par l'hypothèse de finitude sur $H^{p + 1}(X, \mathcal{F}_1)$. Nous omettons certains détails.

\end{proof}

\begin{lemma}

\label{cohomology-lemma-ML-better}

Soit $A$ un anneau. Soit $f \in A$. Soit $X$ un espace topologique. Soit
$$
\ldots \to \mathcal{F}_3 \to \mathcal{F}_2 \to \mathcal{F}_1
$$
un système projectif de faisceaux de $A$-modules. Soit $p \geq 0$. Supposons que

\begin{enumerate}

\item soit il existe $m \geq 1$ tel que l'image de
$H^{p + 1}(X, \mathcal{F}_m) \to H^{p + 1}(X, \mathcal{F}_1)$
soit un $A$-module de longueur finie, soit $A$ est noethérien et l'intersection des images de
$H^{p + 1}(X, \mathcal{F}_m) \to H^{p + 1}(X, \mathcal{F}_1)$
est un $A$-module de type fini,
\item $X$, $f$ et $(\mathcal{F}_n)$ satisfont à la condition (1) du lemme \ref{cohomology-lemma-equivalent-f-good}.

\end{enumerate}

Alors le système projectif $M_n = H^p(X, \mathcal{F}_n)$ satisfait à la condition de Mittag--Leffler.

\end{lemma}

\begin{proof}
Posons $I = (f)$. Nous appliquerons le critère du lemme \ref{cohomology-lemma-ML-general-better} aux modules $N_n$. Pour $m \geq n$, nous avons
$I^n\mathcal{F}_{m + 1} = \mathcal{F}_{m + 1 - n}$. Ainsi, nous voyons que
$$
N_n = \bigcap\nolimits_{m \geq 1} \Im\left(
H^{p + 1}(X, \mathcal{F}_m) \to H^{p + 1}(X, \mathcal{F}_1)
\right)
$$
est indépendant de $n$ et que $\bigoplus N_n = \bigoplus N_1 \cdot f^{n + 1}$. Nous concluons donc exactement comme dans la preuve du lemme \ref{cohomology-lemma-ML}.
\end{proof}

\begin{remark}

\label{cohomology-remark-compare-ML}

Soit $(X, \mathcal{O}_X)$ un espace annelé. Soit $f \in \Gamma(X, \mathcal{O}_X)$. Soit $\mathcal{F}$ un $\mathcal{O}_X$-module. Si $\mathcal{F}$ est sans $f$-torsion, alors, pour tout $p \geq 0$, nous avons une suite exacte courte de systèmes projectifs
$$
0 \to \{H^p(X, \mathcal{F})/f^nH^p(X, \mathcal{F})\}
\to \{H^p(X, \mathcal{F}/f^n\mathcal{F})\}
\to \{H^{p + 1}(X, \mathcal{F})[f^n]\} \to 0
$$
Comme le premier système projectif satisfait à la condition de Mittag--Leffler (ML), nous en déduisons trois faits :

\begin{enumerate}

\item Il existe une suite exacte courte
$$
0 \to \widehat{H^p(X, \mathcal{F})} \to
\lim H^p(X, \mathcal{F}/f^n\mathcal{F}) \to
T_f(H^{p + 1}(X, \mathcal{F})) \to 0
$$
où $\widehat{\ }$ désigne la complétion $f$-adique usuelle et
$T_f( - )$ le module de Tate $f$-adique de Compléments sur l'algèbre, exemple
\ref{more-algebra-example-spectral-sequence-principal}.
\item Nous avons
$R^1\lim H^p(X, \mathcal{F}/f^n\mathcal{F}) =
R^1\lim H^{p + 1}(X, \mathcal{F})[f^n]$.
\item Le système $\{H^{p + 1}(X, \mathcal{F})[f^n]\}$ est ML si et seulement si
$\{H^p(X, \mathcal{F}/f^n\mathcal{F})\}$ est ML.

\end{enumerate}

Voir Homologie, lemme \ref{homology-lemma-Mittag-Leffler}, et Compléments sur l'algèbre, lemmes \ref{more-algebra-lemma-six-term-Rlim} et
\ref{more-algebra-lemma-Mittag-Leffler}.

\end{remark}









\section{Limites dérivées}

\label{cohomology-section-derived-limits}

\noindent
Soit $(X, \mathcal{O}_X)$ un espace annelé. Puisque la catégorie triangulée $D(\mathcal{O}_X)$ admet des produits (Injectifs, lemme \ref{injectives-lemma-derived-products}), il s'ensuit que $D(\mathcal{O}_X)$ admet des limites dérivées; voir Catégories dérivées, définition \ref{derived-definition-derived-limit}. Si $(K_n)$ est un système projectif dans $D(\mathcal{O}_X)$, nous notons $R\lim K_n$ sa limite dérivée.

\begin{lemma}

\label{cohomology-lemma-RGamma-commutes-with-Rlim}
Soit $(X, \mathcal{O}_X)$ un espace annelé. Pour tout ouvert $U \subset X$, le foncteur $R\Gamma(U, -)$ commute avec $R\lim$. De plus, il existe des suites exactes courtes $$
0 \to
R^1\lim H^{m - 1}(U, K_n) \to H^m(U, R\lim K_n) \to
\lim H^m(U, K_n) \to 0
$$ pour tout système projectif $(K_n)$ dans $D(\mathcal{O}_X)$ et tout $m \in \mathbf{Z}$.
\end{lemma}

\begin{proof}

Le premier énoncé découle d'Injectifs, lemme \ref{injectives-lemma-RF-commutes-with-Rlim}. Nous pouvons ensuite appliquer Compléments sur l'algèbre, remarque \ref{more-algebra-remark-compare-derived-limit},
à $R\lim R\Gamma(U, K_n) = R\Gamma(U, R\lim K_n)$ pour obtenir les suites exactes courtes.

\end{proof}

\begin{lemma}

\label{cohomology-lemma-Rf-commutes-with-Rlim}

Soit $f : (X, \mathcal{O}_X) \to (Y, \mathcal{O}_Y)$ un morphisme d'espaces annelés. Alors $Rf_*$ commute avec $R\lim$, c'est-à-dire que $Rf_*$ commute avec les limites dérivées.

\end{lemma}

\begin{proof}

Soit $(K_n)$ un système projectif dans $D(\mathcal{O}_X)$. Considérons le triangle distingué qui définit la limite dérivée
$$
R\lim K_n \to \prod K_n \to \prod K_n
$$
dans $D(\mathcal{O}_X)$. En appliquant le foncteur exact $Rf_*$, nous obtenons le triangle distingué
$$
Rf_*(R\lim K_n) \to Rf_*\left(\prod K_n\right) \to Rf_*\left(\prod K_n\right)
$$
dans $D(\mathcal{O}_Y)$. Il suffit donc de montrer que
$Rf_*$ commute aux produits de la catégorie dérivée (lesquels ne sont pas simplement donnés par les produits de complexes; voir Injectifs, lemme \ref{injectives-lemma-derived-products}). Or, puisque $Rf_*$ est adjoint à droite d'après le lemme \ref{cohomology-lemma-adjoint},
cela résulte formellement de Catégories, lemme \ref{categories-lemma-adjoint-exact}. Attention : on ne peut pas appliquer directement Catégories, lemme \ref{categories-lemma-adjoint-exact},
car $R\lim K_n$ n'est pas une limite dans $D(\mathcal{O}_X)$.

\end{proof}

\begin{remark}

\label{cohomology-remark-discuss-derived-limit}

Soit $(X, \mathcal{O}_X)$ un espace annelé. Soit $(K_n)$ un système projectif dans $D(\mathcal{O}_X)$. Posons $K = R\lim K_n$. Pour tout $n$ et tout $m$,
soit $\mathcal{H}^m_n = H^m(K_n)$ le $m$-ième faisceau de cohomologie de
$K_n$, et posons de même $\mathcal{H}^m = H^m(K)$. Notons
$\underline{\mathcal{H}}^m_n$ le préfaisceau
$$
U \longmapsto \underline{\mathcal{H}}^m_n(U) = H^m(U, K_n)
$$
De même, posons $\underline{\mathcal{H}}^m(U) = H^m(U, K)$. D'après le lemme \ref{cohomology-lemma-sheafification-cohomology}, nous voyons que
$\mathcal{H}^m_n$ est la faisceautisation de
$\underline{\mathcal{H}}^m_n$ et que $\mathcal{H}^m$ est la faisceautisation de $\underline{\mathcal{H}}^m$. Voici un diagramme
$$
\xymatrix{
K \ar@{=}[d] &
\underline{\mathcal{H}}^m \ar[d] \ar[r] & 
\mathcal{H}^m \ar[d] \\
R\lim K_n &
\lim \underline{\mathcal{H}}^m_n \ar[r] & 
\lim \mathcal{H}^m_n
}
$$
En général, il n'est pas vrai que
$\lim \mathcal{H}^m_n$ soit le faisceau associé à
$\lim \underline{\mathcal{H}}^m_n$. Si $U \subset X$ est un ouvert, nous avons alors des suites exactes courtes

\begin{equation}
\label{cohomology-equation-ses-Rlim-over-U}
0 \to
R^1\lim \underline{\mathcal{H}}^{m - 1}_n(U) \to
\underline{\mathcal{H}}^m(U) \to
\lim \underline{\mathcal{H}}^m_n(U) \to 0
\end{equation}

d'après le lemme \ref{cohomology-lemma-RGamma-commutes-with-Rlim}.

\end{remark}

\noindent
Le lemme suivant s'applique à un système projectif de modules quasi-cohérents dont les morphismes de transition sont surjectifs sur un schéma.

\begin{lemma}

\label{cohomology-lemma-inverse-limit-is-derived-limit}

Soit $(X, \mathcal{O}_X)$ un espace annelé. Soit $(\mathcal{F}_n)$ un système projectif de $\mathcal{O}_X$-modules. Soit $\mathcal{B}$ un ensemble d'ouverts de $X$. Supposons que

\begin{enumerate}

\item tout ouvert de $X$ admet un recouvrement dont les membres appartiennent à
$\mathcal{B}$,
\item  $H^p(U, \mathcal{F}_n) = 0$ pour $p > 0$ et $U \in \mathcal{B}$,
\item le système projectif $\mathcal{F}_n(U)$ a un $R^1\lim$ nul
pour $U \in \mathcal{B}$.

\end{enumerate}

Alors $R\lim \mathcal{F}_n = \lim \mathcal{F}_n$ et nous avons
$H^p(U, \lim \mathcal{F}_n) = 0$ pour $p > 0$ et $U \in \mathcal{B}$.

\end{lemma}

\begin{proof}

Posons $K_n = \mathcal{F}_n$ et $K = R\lim \mathcal{F}_n$. Avec les notations de la remarque \ref{cohomology-remark-discuss-derived-limit} et l'hypothèse (2), nous voyons que, pour
$U \in \mathcal{B}$, nous avons $\underline{\mathcal{H}}_n^m(U) = 0$
lorsque $m \not = 0$ et $\underline{\mathcal{H}}_n^0(U) = \mathcal{F}_n(U)$. De l'équation (\ref{cohomology-equation-ses-Rlim-over-U}) et l'hypothèse (3) nous voyons que $\underline{\mathcal{H}}^m(U) = 0$
lorsque $m \not = 0$ et égal à $\lim \mathcal{F}_n(U)$
lorsque $m = 0$. En faisceautisant et en utilisant (1), nous obtenons
$\mathcal{H}^m = 0$ lorsque $m \not = 0$ et égal à
$\lim \mathcal{F}_n$ lorsque $m = 0$. Ainsi, $K = \lim \mathcal{F}_n$. Comme $H^m(U, K) = \underline{\mathcal{H}}^m(U) = 0$ pour $m > 0$
(voir ci-dessus), la seconde assertion est vérifiée.

\end{proof}

\begin{lemma}

\label{cohomology-lemma-cohomology-derived-limit-injective}

Soit $(X, \mathcal{O}_X)$ un espace annelé. Soit $(K_n)$ un système projectif dans $D(\mathcal{O}_X)$. Soient $x \in X$ et $m \in \mathbf{Z}$. Supposons qu'il existe un entier $n(x)$ et un système fondamental $\mathfrak{U}_x$
de voisinages ouverts de $x$ tels que, pour $U \in \mathfrak{U}_x$,

\begin{enumerate}

\item  $R^1\lim H^{m - 1}(U, K_n) = 0$,
\item  $H^m(U, K_n) \to H^m(U, K_{n(x)})$ est injectif pour $n \geq n(x)$.

\end{enumerate}

Alors l'application sur les germes $H^m(R\lim K_n)_x \to H^m(K_{n(x)})_x$ est injective.

\end{lemma}

\begin{proof}

Soit $\gamma$ un élément de $H^m(R\lim K_n)_x$ qui s'envoie sur zéro dans $H^m(K_{n(x)})_x$. Comme $H^m(R\lim K_n)$ est la faisceautisation de $U \mapsto H^m(U, R\lim K_n)$,
(d'après le lemme \ref{cohomology-lemma-sheafification-cohomology}) nous pouvons choisir $U \in \mathfrak{U}_x$
et un élément $\tilde \gamma \in H^m(U, R\lim K_n)$ s'envoyant sur $\gamma$. Alors $\tilde\gamma$ s'envoie sur $\tilde\gamma_{n(x)} \in H^m(U, K_{n(x)})$. Puisque $H^m(K_{n(x)})$ est la faisceautisation de
$U \mapsto H^m(U, K_{n(x)})$
(à nouveau d'après le lemme \ref{cohomology-lemma-sheafification-cohomology}), nous voyons qu'après avoir rétréci $U$, nous pouvons supposer que $\tilde\gamma_{n(x)} = 0$. Pour cet ouvert $U$, considérons la suite exacte courte
$$
0 \to
R^1\lim H^{m - 1}(U, K_n) \to H^m(U, R\lim K_n) \to
\lim H^m(U, K_n) \to 0
$$
du lemme \ref{cohomology-lemma-RGamma-commutes-with-Rlim}. Par l'hypothèse (1), le groupe de gauche est nul et, par l'hypothèse (2), celui de droite s'injecte dans $H^m(U, K_{n(x)})$. Nous concluons que $\tilde\gamma = 0$
et donc que $\gamma = 0$, comme souhaité.

\end{proof}

\begin{lemma}

\label{cohomology-lemma-is-limit-per-point}

Soit $(X, \mathcal{O}_X)$ un espace annelé et soit $E \in D(\mathcal{O}_X)$. Supposons que, pour tout $x \in X$, il existe une fonction $p(x, -) : \mathbf{Z} \to \mathbf{Z}$ et un système fondamental $\mathfrak{U}_x$ de voisinages ouverts de $x$
tels que
$$
H^p(U, H^{m - p}(E)) = 0 \text{ pour }
U \in \mathfrak{U}_x \text{ et } p > p(x, m)
$$
Alors l'application $E \to R\lim \tau_{\geq -n} E$ de Catégories dérivées, remarque
\ref{derived-remark-map-into-derived-limit-truncations}
est un isomorphisme dans $D(\mathcal{O}_X)$.

\end{lemma}

\begin{proof}

Posons $K_n = \tau_{\geq -n}E$ et $K = R\lim K_n$. Le morphisme canonique $E \to K$
provient des morphismes canoniques $E \to K_n = \tau_{\geq -n}E$. Nous devons montrer que $E \to K$ induit un isomorphisme
$H^m(E) \to H^m(K)$ de faisceaux de cohomologie. Dans la suite de la preuve, fixons $m$. Si $n \geq -m$, alors le morphisme $E \to \tau_{\geq -n}E = K_n$ induit un isomorphisme
$H^m(E) \to H^m(K_n)$. Pour finir la preuve il suffit de montrer que pour chaque $x \in X$
il existe un entier $n(x) \geq -m$ de sorte que l'application
$H^m(K)_x \to H^m(K_{n(x)})_x$ est injective. En effet, la composée
$$
H^m(E)_x \to H^m(K)_x \to H^m(K_{n(x)})_x
$$
est une bijection et la deuxième flèche est injective; la première est donc bijective. Posons
$$
n(x) = 1 + \max\{-m, p(x, m - 1) - m, -1 + p(x, m) - m, -2 + p(x, m + 1) - m\}.
$$
de sorte que, dans tous les cas, $n(x) \geq -m$. Assertion : les morphismes
$$
H^{m - 1}(U, K_{n + 1}) \to H^{m - 1}(U, K_n)
\quad\text{et}\quad
H^m(U, K_{n + 1}) \to H^m(U, K_n)
$$
sont des isomorphismes pour $n \geq n(x)$ et $U \in \mathfrak{U}_x$. L'assertion implique que les conditions (1) et (2) du lemme \ref{cohomology-lemma-cohomology-derived-limit-injective}
sont satisfaites et fournit donc l'injectivité souhaitée. Rappelons (Catégories dérivées, remarque
\ref{derived-remark-truncation-distinguished-triangle}) que nous avons des triangles distingués
$$
H^{-n - 1}(E)[n + 1] \to
K_{n + 1} \to K_n \to H^{-n - 1}(E)[n + 2]
$$
En examinant la suite exacte longue de cohomologie associée, l'assertion découle de la nullité des groupes
$$
H^{m + n}(U, H^{-n - 1}(E)),\quad
H^{m + n + 1}(U, H^{-n - 1}(E)),\quad
H^{m + n + 2}(U, H^{-n - 1}(E))
$$
pour $n \geq n(x)$ et $U \in \mathfrak{U}_x$. Cette nullité découle de notre choix de $n(x)$
et de l'hypothèse du lemme.

\end{proof}

\begin{lemma}

\label{cohomology-lemma-is-limit-spaltenstein}

\begin{reference}
\cite[proposition 3.13]{Spaltenstein}
\end{reference}
Soit $(X, \mathcal{O}_X)$ un espace annelé. Soit $E \in D(\mathcal{O}_X)$. Supposons que, pour tout $x \in X$, il existe un entier $d_x \geq 0$ et un système fondamental $\mathfrak{U}_x$ de voisinages ouverts de $x$ tels que $$
H^p(U, H^q(E)) = 0 \text{ pour }
U \in \mathfrak{U}_x,\ p > d_x, \text{ et }q < 0
$$ Alors l'application $E \to R\lim \tau_{\geq -n} E$ de Catégories dérivées, remarque \ref{derived-remark-map-into-derived-limit-truncations}, est un isomorphisme dans $D(\mathcal{O}_X)$.
\end{lemma}

\begin{proof}
Ceci résulte du lemme \ref{cohomology-lemma-is-limit-per-point} avec $p(x, m) = d_x + \max(0, m)$.
\end{proof}

\begin{lemma}

\label{cohomology-lemma-is-limit}

Soit $(X, \mathcal{O}_X)$ un espace annelé. Soit $E \in D(\mathcal{O}_X)$. Supposons qu'il existe une fonction $p(-) : \mathbf{Z} \to \mathbf{Z}$
et un ensemble $\mathcal{B}$ d'ouverts de $X$ tels que

\begin{enumerate}

\item tout ouvert de $X$ admet un recouvrement dont les membres appartiennent à $\mathcal{B}$,
\item  $H^p(U, H^{m - p}(E)) = 0$ pour $p > p(m)$ et $U \in \mathcal{B}$.

\end{enumerate}

Alors l'application $E \to R\lim \tau_{\geq -n} E$ de Catégories dérivées, remarque
\ref{derived-remark-map-into-derived-limit-truncations}
est un isomorphisme dans $D(\mathcal{O}_X)$.

\end{lemma}

\begin{proof}
Appliquer le lemme \ref{cohomology-lemma-is-limit-per-point} avec $p(x, m) = p(m)$ et $\mathfrak{U}_x = \{U \in \mathcal{B} \mid x \in U\}$.
\end{proof}

\begin{lemma}

\label{cohomology-lemma-is-limit-dimension}
Soit $(X, \mathcal{O}_X)$ un espace annelé. Soit $E \in D(\mathcal{O}_X)$. Supposons qu'il existe un entier $d \geq 0$ et une base $\mathcal{B}$ de la topologie de $X$ telle que $$
H^p(U, H^q(E)) = 0 \text{ pour }
U \in \mathcal{B},\ p > d, \text{ et }q < 0
$$ Alors l'application $E \to R\lim \tau_{\geq -n} E$ de Catégories dérivées, remarque \ref{derived-remark-map-into-derived-limit-truncations}, est un isomorphisme dans $D(\mathcal{O}_X)$.
\end{lemma}

\begin{proof}
Appliquer le lemme \ref{cohomology-lemma-is-limit-spaltenstein} avec $d_x = d$ et $\mathfrak{U}_x = \{U \in \mathcal{B} \mid x \in U\}$.
\end{proof}

\noindent
Les lemmes ci-dessus peuvent être utilisés pour calculer la cohomologie dans certaines situations.

\begin{lemma}

\label{cohomology-lemma-cohomology-over-U-trivial}

Soit $(X, \mathcal{O}_X)$ un espace annelé. Soit $K$
un objet de $D(\mathcal{O}_X)$. Soit $\mathcal{B}$ un ensemble d'ouverts de $X$. Supposons que

\begin{enumerate}

\item tout ouvert de $X$ admet un recouvrement dont les membres appartiennent à $\mathcal{B}$,
\item  $H^p(U, H^q(K)) = 0$ pour tous $p > 0$, $q \in \mathbf{Z}$,
$U \in \mathcal{B}$.

\end{enumerate}

Alors $H^q(U, K) = H^0(U, H^q(K))$ pour $q \in \mathbf{Z}$
et $U \in \mathcal{B}$.

\end{lemma}

\begin{proof}

Remarquons que $K = R\lim \tau_{\geq -n} K$ d'après le lemme \ref{cohomology-lemma-is-limit-dimension} avec $d = 0$. Soit $U \in \mathcal{B}$. Par l'équation (\ref{cohomology-equation-ses-Rlim-over-U}), nous obtenons une suite exacte courte
$$
0 \to R^1\lim H^{q - 1}(U, \tau_{\geq -n}K) \to
H^q(U, K) \to \lim H^q(U, \tau_{\geq -n}K) \to 0
$$
La condition (2) implique
$H^q(U, \tau_{\geq -n} K) = H^0(U, H^q(\tau_{\geq -n} K))$
pour tout $q$, à l'aide de la suite spectrale de l'exemple \ref{cohomology-example-spectral-sequence}. Cette suite spectrale converge parce que $\tau_{\geq -n}K$ est borné inférieurement. Si $n > -q$, alors $H^q(\tau_{\geq -n}K) = H^q(K)$. Ainsi, les systèmes de gauche et de droite de la suite exacte courte affichée sont stationnaires, de valeurs respectives
$H^0(U, H^{q - 1}(K))$ et $H^0(U, H^q(K))$. Le lemme suit.

\end{proof}

\noindent
Voici un autre cas où nous pouvons décrire la limite dérivée.

\begin{lemma}

\label{cohomology-lemma-derived-limit-suitable-system}

Soit $(X, \mathcal{O}_X)$ un espace annelé. Soit $(K_n)$
un système projectif d'objets de $D(\mathcal{O}_X)$. Soit $\mathcal{B}$ un ensemble d'ouverts de $X$. Supposons que

\begin{enumerate}

\item tout ouvert de $X$ admet un recouvrement dont les membres appartiennent à $\mathcal{B}$,
\item pour tout $U \in \mathcal{B}$ et tout $q \in \mathbf{Z}$, nous avons

\begin{enumerate}

\item  $H^p(U, H^q(K_n)) = 0$ pour $p > 0$,
\item le système projectif $H^0(U, H^q(K_n))$ a un $R^1\lim$ nul.

\end{enumerate}

\end{enumerate}

Alors $H^q(R\lim K_n) = \lim H^q(K_n)$ pour $q \in \mathbf{Z}$.

\end{lemma}

\begin{proof}

Posons $K = R\lim K_n$. Soit $U \in \mathcal{B}$. D'après le lemme \ref{cohomology-lemma-cohomology-over-U-trivial} et (2) a), nous avons $H^q(U, K_n) = H^0(U, H^q(K_n))$. D'après le lemme \ref{cohomology-lemma-RGamma-commutes-with-Rlim} et (2) b), nous avons 
$H^q(U, K) = \lim H^0(U, H^q(K_n))$. Ainsi, $H^q(U, K)$ est la limite projective des sections des faisceaux $H^q(K_n)$ sur $U$. Comme
$\lim H^q(K_n)$ est un faisceau, l'hypothèse (1) montre que $H^q(K)$, qui est la faisceautisation du préfaisceau $U \mapsto H^q(U, K)$, est égal à $\lim H^q(K_n)$. Cela prouve le lemme.

\end{proof}







\section{Construction de résolutions K-injectives}

\label{cohomology-section-K-injective}

\noindent

Soit $(X, \mathcal{O}_X)$ un espace annelé. Soit $\mathcal{F}^\bullet$ un complexe de $\mathcal{O}_X$-modules. La catégorie $\textit{Mod}(\mathcal{O}_X)$ a assez d'injectifs; nous pouvons donc utiliser Catégories dérivées, lemme \ref{derived-lemma-special-inverse-system},
pour produire un diagramme
$$
\xymatrix{
\ldots \ar[r] &
\tau_{\geq -2}\mathcal{F}^\bullet \ar[r] \ar[d] &
\tau_{\geq -1}\mathcal{F}^\bullet \ar[d] \\
\ldots \ar[r] & \mathcal{I}_2^\bullet \ar[r] & \mathcal{I}_1^\bullet
}
$$
dans la catégorie des complexes de $\mathcal{O}_X$-modules, tel que

\begin{enumerate}

\item les flèches verticales sont des quasi-isomorphismes,
\item $\mathcal{I}_n^\bullet$ est un complexe borné inférieurement d'objets injectifs,
\item les flèches $\mathcal{I}_{n + 1}^\bullet \to \mathcal{I}_n^\bullet$
sont des surjections scindées terme à terme.

\end{enumerate}

La catégorie des $\mathcal{O}_X$-modules admet des limites (elles se calculent au niveau des préfaisceaux); nous pouvons donc former la limite
$\mathcal{I}^\bullet = \lim_n \mathcal{I}_n^\bullet$. Par Catégories dérivées, lemmes
\ref{derived-lemma-bounded-below-injectives-K-injective} et
\ref{derived-lemma-limit-K-injectives}
il s'agit d'un complexe K-injectif. En général, l'application canonique

\begin{equation}
\label{cohomology-equation-into-candidate-K-injective}
\mathcal{F}^\bullet \to \mathcal{I}^\bullet
\end{equation}

peut ne pas être un quasi-isomorphisme. Dans le lemme suivant, nous donnons des conditions sous lesquelles il en est un.

\begin{lemma}

\label{cohomology-lemma-K-injective}

Dans la situation décrite ci-dessus, soit $\mathcal{H}^m = H^m(\mathcal{F}^\bullet)$ le faisceau de cohomologie de degré $m$. Soit $\mathcal{B}$ un ensemble de sous-ensembles ouverts de $X$. Soit $d \in \mathbf{N}$. Supposons que

\begin{enumerate}

\item tout ouvert de $X$ admet un recouvrement dont les membres appartiennent à $\mathcal{B}$,
\item pour tout $U \in \mathcal{B}$, nous avons $H^p(U, \mathcal{H}^q) = 0$
pour $p > d$ et $q < 0$\footnote{Il suffit que
$\forall m$, $\exists p(m)$, $H^p(U. \mathcal{H}^{m - p}) = 0$ pour
$p > p(m)$, voir lemme \ref{cohomology-lemma-is-limit}.}.

\end{enumerate}

Alors (\ref{cohomology-equation-into-candidate-K-injective}) est un quasi-isomorphisme.

\end{lemma}

\begin{proof}
D'après Catégories dérivées, lemme \ref{derived-lemma-difficulty-K-injectives}, il suffit de montrer que l'application $\mathcal{F}^\bullet \to R\lim \tau_{\geq -n} \mathcal{F}^\bullet$ est un isomorphisme. C'est le lemme \ref{cohomology-lemma-is-limit-dimension}.
\end{proof}

\noindent
Voici un lemme technique sur les faisceaux de cohomologie de la limite projective d'un système de complexes de faisceaux. En un certain sens, ce n'est pas le bon énoncé à chercher à démontrer, car il faudrait prendre des limites dérivées plutôt que des limites projectives ordinaires.

\begin{lemma}

\label{cohomology-lemma-inverse-limit-complexes}

Soit $(X, \mathcal{O}_X)$ un espace annelé. Soit $(\mathcal{F}_n^\bullet)$
un système projectif de complexes de $\mathcal{O}_X$-modules. Soit $m \in \mathbf{Z}$. Supposons qu'il existe un ensemble $\mathcal{B}$ de sous-ensembles ouverts de $X$ et un entier $n_0$ tels que

\begin{enumerate}

\item tout ouvert de $X$ admet un recouvrement dont les membres appartiennent à $\mathcal{B}$,
\item pour chaque $U \in \mathcal{B}$

\begin{enumerate}

\item les systèmes des groupes abéliens
$\mathcal{F}_n^{m - 2}(U)$ et $\mathcal{F}_n^{m - 1}(U)$
ont un $R^1\lim$ nul (par exemple, ils satisfont à la condition de Mittag--Leffler),
\item le système des groupes abéliens $H^{m - 1}(\mathcal{F}_n^\bullet(U))$
a un $R^1\lim$ nul (par exemple, il satisfait à la condition de Mittag--Leffler), et
\item Nous avons
$H^m(\mathcal{F}_n^\bullet(U)) = H^m(\mathcal{F}_{n_0}^\bullet(U))$
pour tous $n \geq n_0$.

\end{enumerate}

\end{enumerate}

Alors les morphismes
$H^m(\mathcal{F}^\bullet) \to \lim H^m(\mathcal{F}_n^\bullet) \to
H^m(\mathcal{F}_{n_0}^\bullet)$
sont des isomorphismes de faisceaux, où
$\mathcal{F}^\bullet = \lim \mathcal{F}_n^\bullet$ est la limite projective par terme.

\end{lemma}

\begin{proof}

Soit $U \in \mathcal{B}$. Notons que $H^m(\mathcal{F}^\bullet(U))$ est la cohomologie de
$$
\lim_n \mathcal{F}_n^{m - 2}(U) \to
\lim_n \mathcal{F}_n^{m - 1}(U) \to
\lim_n \mathcal{F}_n^m(U) \to
\lim_n \mathcal{F}_n^{m + 1}(U)
$$
Sous les hypothèses (2)a) et (2)b), nous pouvons appliquer Compléments sur l'algèbre, lemme \ref{more-algebra-lemma-apply-Mittag-Leffler-again},
pour conclure que
$$
H^m(\mathcal{F}^\bullet(U)) = \lim H^m(\mathcal{F}_n^\bullet(U))
$$
Par hypothèse (2) c), nous concluons
$$
H^m(\mathcal{F}^\bullet(U)) = H^m(\mathcal{F}_n^\bullet(U))
$$
pour tout $n \geq n_0$. Par l'hypothèse (1), nous concluons que la faisceautisation de
$U \mapsto H^m(\mathcal{F}^\bullet(U))$ est égale à celle de $U \mapsto H^m(\mathcal{F}_n^\bullet(U))$ pour tout $n \geq n_0$. Ainsi, le système projectif des faisceaux $H^m(\mathcal{F}_n^\bullet)$ est constant pour $n \geq n_0$, de valeur $H^m(\mathcal{F}^\bullet)$. Cela prouve le lemme.

\end{proof}











\section{Systèmes projectifs et cohomologie, III}

\label{cohomology-section-inverse-systems-tri}

\noindent
Cette section poursuit la discussion dans la section \ref{cohomology-section-inverse-systems-bis} en utilisant des limites dérivées.

\begin{lemma}

\label{cohomology-lemma-formal-functions-principal}

Soit $(X, \mathcal{O}_X)$ un espace annelé. Soit $A \to \Gamma(X, \mathcal{O}_X)$
un homomorphisme d'anneaux et soit $f \in A$. Soit $E$ un objet de $D(\mathcal{O}_X)$. Posons
$$
E_n = E \otimes_{\mathcal{O}_X} (\mathcal{O}_X \xrightarrow{f^n} \mathcal{O}_X)
$$
et définissons $E^\wedge = R\lim E_n$. Pour $p \in \mathbf{Z}$,
il existe un diagramme commutatif canonique
$$
\xymatrix{
& 0 & 0 \\
0 \ar[r] &
\widehat{H^p(X, E)} \ar[r] \ar[u] &
\lim H^p(X, E_n) \ar[r] \ar[u] &
T_f(H^{p + 1}(X, E)) \ar[r] &
0 \\
0 \ar[r] &
H^0(H^p(X, E)^\wedge) \ar[r] \ar[u] &
H^p(X, E^\wedge) \ar[r] \ar[u] &
T_f(H^{p + 1}(X, E)) \ar[r] \ar@{=}[u] &
0 \\
&
R^1\lim H^p(X, E)[f^n] \ar[u] \ar[r]^\cong &
R^1\lim H^{p - 1}(X, E_n) \ar[u] \\
& 0 \ar[u] & 0 \ar[u]
}
$$
dont les lignes et les colonnes sont exactes, où $\widehat{H^p(X, E)} = \lim H^p(X, E)/f^n H^p(X, E)$ est la complétion
$f$-adique, $H^p(X, E)^\wedge$ est la complétion $f$-adique dérivée, et $T_f(H^{p + 1}(X, E))$ est le module de Tate $f$-adique; voir Compléments sur l'algèbre, exemple
\ref{more-algebra-example-spectral-sequence-principal}. Enfin, nous avons $H^p(X, E^\wedge) = H^p(R\Gamma(X, E)^\wedge)$.

\end{lemma}

\begin{proof}
Remarquons que $R\Gamma(X, E^\wedge) = R\lim R\Gamma(X, E_n)$ d'après le lemme \ref{cohomology-lemma-Rf-commutes-with-Rlim}. D'autre part, nous avons $$
R\Gamma(X, E_n) =
R\Gamma(X, E) \otimes_A^\mathbf{L} (A \xrightarrow{f^n} A)
$$ (détails omis). Nous en déduisons que $R\Gamma(X, E^\wedge)$ est la complétion $f$-adique dérivée $R\Gamma(X, E)^\wedge$. Le diagramme découle alors de Compléments sur l'algèbre, lemme \ref{more-algebra-lemma-compare-completion-derived-completion}.
\end{proof}

\begin{lemma}

\label{cohomology-lemma-stabilizes}
Soit $\mathcal{A}$ une catégorie abélienne. Soit $f : M \to M$ un morphisme de $\mathcal{A}$. Si $M[f^n] = \Ker(f^n : M \to M)$ se stabilise, alors les systèmes projectifs $$
(M \xrightarrow{f^n} M)
\quad\text{et}\quad
\Coker(f^n : M \to M)
$$ sont pro-isomorphes dans $D(\mathcal{A})$.
\end{lemma}

\begin{proof}
Il existe manifestement un morphisme du premier système projectif vers le second. Supposons que $M[f^c] = M[f^{c + 1}] = M[f^{c + 2}] = \ldots$. Nous pouvons alors définir un morphisme de systèmes projectifs dans $D(\mathcal{A})$ en sens inverse au moyen des diagrammes $$
\xymatrix{
M/M[f^c] \ar[r]_-{f^{n + c}} \ar[d]_{f^c} &
M \ar[d]^1 \\
M \ar[r]^{f^n} & M
}
$$ Puisque la flèche horizontale supérieure est injective, le complexe de la ligne supérieure est quasi-isomorphe à $\Coker(f^{n + c} : M \to M)$. Nous omettons quelques détails.
\end{proof}

\begin{example}

\label{cohomology-example-formal-functions-principal}

Soit $(X, \mathcal{O}_X)$ un espace annelé. Soit $A \to \Gamma(X, \mathcal{O}_X)$
un homomorphisme d'anneaux et soit $f \in A$. Soit $\mathcal{F}$ un $\mathcal{O}_X$-module. Supposons qu'il existe un entier $c$ tel que $\mathcal{F}[f^c] = \mathcal{F}[f^n]$
pour tout $n \geq c$. Appliquons le lemme \ref{cohomology-lemma-formal-functions-principal} avec
$E = \mathcal{F}$. D'après le lemme \ref{cohomology-lemma-stabilizes}, le système projectif $(E_n)$ est pro-isomorphe au système projectif
$(\mathcal{F}/f^n\mathcal{F})$. Nous en concluons que, pour $p \in \mathbf{Z}$,
nous obtenons le diagramme commutatif
$$
\xymatrix{
& 0 & 0 \\
0 \ar[r] &
\widehat{H^p(X, \mathcal{F})} \ar[r] \ar[u] &
\lim H^p(X, \mathcal{F}/f^n\mathcal{F}) \ar[r] \ar[u] &
T_f(H^{p + 1}(X, \mathcal{F})) \ar[r] &
0 \\
0 \ar[r] &
H^0(H^p(X, \mathcal{F})^\wedge) \ar[r] \ar[u] &
H^p(R\Gamma(X, \mathcal{F})^\wedge) \ar[r] \ar[u] &
T_f(H^{p + 1}(X, \mathcal{F})) \ar[r] \ar@{=}[u] &
0 \\
&
R^1\lim H^p(X, \mathcal{F})[f^n] \ar[u] \ar[r]^\cong &
R^1\lim H^{p - 1}(X, \mathcal{F}/f^n\mathcal{F}) \ar[u] \\
& 0 \ar[u] & 0 \ar[u]
}
$$
dont les lignes et les colonnes sont exactes, où
$\widehat{H^p(X, \mathcal{F})} =
\lim H^p(X, \mathcal{F})/f^n H^p(X, \mathcal{F})$
est la complétion $f$-adique usuelle et où $M^\wedge$ désigne la complétion $f$-adique dérivée de $M$ dans $D(A)$.

\end{example}















\section{Cohomologie de {\v C}ech des complexes non bornés}

\label{cohomology-section-cech-cohomology-of-unbounded-complexes}

\noindent
La construction de la section \ref{cohomology-section-cech-cohomology-of-complexes} n'est pas la bonne pour les complexes non bornés. Le problème est que, dans le projet Champs, nous utilisons des sommes directes dans la totalisation d'un complexe double et que nous devrions les remplacer par un produit. Au lieu d'effectuer ce remplacement dans cette section, nous supposons que le recouvrement est fini et nous utilisons le complexe de {\v C}ech alterné.

\medskip\noindent

Soit $(X, \mathcal{O}_X)$ un espace annelé. Soit ${\mathcal F}^\bullet$ un complexe de préfaisceaux de
$\mathcal{O}_X$-modules. Soit ${\mathcal U} : X = \bigcup_{i \in I} U_i$
un {\bf recouvrement ouvert fini} de $X$. Comme le complexe de {\v C}ech alterné
$\check{\mathcal{C}}_{alt}^\bullet(\mathcal{U}, \mathcal{F})$
(section \ref{cohomology-section-alternating-cech}) est fonctoriel en $\mathcal{F}$, nous obtenons un complexe double
$\check{\mathcal{C}}^\bullet_{alt}(\mathcal{U}, \mathcal{F}^\bullet)$. Dans cette section, nous travaillons avec le complexe total associé.
La construction de $\text{Tot}(\check{\mathcal{C}}^\bullet_{alt}({\mathcal U},
{\mathcal F}^\bullet))$
est fonctorielle en ${\mathcal F}^\bullet$. Il existe également une transformation fonctorielle

\begin{equation}
\label{cohomology-equation-global-sections-to-alternating-cech}
\Gamma(X, {\mathcal F}^\bullet)
\longrightarrow
\text{Tot}(\check{\mathcal{C}}^\bullet_{alt}({\mathcal U},
{\mathcal F}^\bullet))
\end{equation}

de complexes définie par la règle suivante : une section
$s\in \Gamma(X, {\mathcal F}^n)$
est envoyée vers l'élément $\alpha = \{\alpha_{i_0\ldots i_p}\}$
avec $\alpha_{i_0} = s|_{U_{i_0}}$ et $\alpha_{i_0\ldots i_p} = 0$
pour $p > 0$.

\begin{lemma}

\label{cohomology-lemma-alternating-cech-complex-complex}

Soit $(X, \mathcal{O}_X)$ un espace annelé. Soit $\mathcal{U} : X = \bigcup_{i \in I} U_i$ un recouvrement fini ouvert. Pour un complexe $\mathcal{F}^\bullet$
de $\mathcal{O}_X$-modules, il existe un morphisme canonique
$$
\text{Tot}(\check{\mathcal{C}}^\bullet_{alt}(\mathcal{U}, \mathcal{F}^\bullet))
\longrightarrow
R\Gamma(X, \mathcal{F}^\bullet)
$$
fonctoriel en $\mathcal{F}^\bullet$ et compatible avec (\ref{cohomology-equation-global-sections-to-alternating-cech}).

\end{lemma}

\begin{proof}

Soit ${\mathcal I}^\bullet$ un complexe K-injectif dont les termes sont des $\mathcal{O}_X$-modules injectifs. Le morphisme (\ref{cohomology-equation-global-sections-to-alternating-cech}) pour
$\mathcal{I}^\bullet$ est un morphisme
$\Gamma(X, {\mathcal I}^\bullet) \to
\text{Tot}(\check{\mathcal{C}}^\bullet_{alt}({\mathcal U},
{\mathcal I}^\bullet))$. Il s'agit d'un quasi-isomorphisme de complexes de groupes abéliens comme il résulte de Homologie, lemme \ref{homology-lemma-double-complex-gives-resolution}
appliqué au complexe double
$\check{\mathcal{C}}^\bullet_{alt}({\mathcal U},
{\mathcal I}^\bullet)$, au moyen des lemmes \ref{cohomology-lemma-injective-trivial-cech} et \ref{cohomology-lemma-alternating-usual}. Supposons que ${\mathcal F}^\bullet \to {\mathcal I}^\bullet$ soit un quasi-isomorphisme de ${\mathcal F}^\bullet$ dans un complexe K-injectif dont les termes sont injectifs (Injectifs, théorème
\ref{injectives-theorem-K-injective-embedding-grothendieck}). Puisque $R\Gamma(X, {\mathcal F}^\bullet)$ est représenté par le complexe
$\Gamma(X, {\mathcal I}^\bullet)$, nous obtenons le morphisme du lemme au moyen de
$$
\text{Tot}(\check{\mathcal{C}}^\bullet_{alt}({\mathcal U},
{\mathcal F}^\bullet))
\longrightarrow
\text{Tot}(\check{\mathcal{C}}^\bullet_{alt}({\mathcal U},
{\mathcal I}^\bullet)).
$$
Nous omettons la vérification de la fonctorialité et des compatibilités.

\end{proof}

\begin{lemma}

\label{cohomology-lemma-alternating-cech-complex-complex-ss}

Soit $(X, \mathcal{O}_X)$ un espace annelé.
Soit $\mathcal{U} : X = \bigcup_{i \in I} U_i$ un recouvrement fini ouvert.
Soit $\mathcal{F}^\bullet$ un complexe de $\mathcal{O}_X$-modules. Soit $\mathcal{B}$ un ensemble d'ouverts de $X$. Supposons que

\begin{enumerate}

\item tout ouvert de $X$ possède un recouvrement dont les membres appartiennent à $\mathcal{B}$,
\item on a $U_{i_0\ldots i_p} \in \mathcal{B}$ pour tous
$i_0, \ldots, i_p \in I$,
\item pour tout $U \in \mathcal{B}$ et tout $p > 0$, on a

\begin{enumerate}

\item  $H^p(U, \mathcal{F}^q) = 0$,
\item  $H^p(U, \Coker(\mathcal{F}^{q - 1} \to \mathcal{F}^q)) = 0$,
\item  $H^p(U, H^q(\mathcal{F})) = 0$.

\end{enumerate}

\end{enumerate}

Alors le morphisme
$$
\text{Tot}(\check{\mathcal{C}}^\bullet_{alt}(\mathcal{U}, \mathcal{F}^\bullet))
\longrightarrow
R\Gamma(X, \mathcal{F}^\bullet)
$$
du lemme \ref{cohomology-lemma-alternating-cech-complex-complex}
est un isomorphisme dans $D(\textit{Ab})$.

\end{lemma}

\begin{proof}
Supposons d'abord que $\mathcal{F}^\bullet$ soit borné inférieurement. Dans ce cas, le morphisme $$
\text{Tot}(\check{\mathcal{C}}^\bullet_{alt}(\mathcal{U}, \mathcal{F}^\bullet))
\longrightarrow
\text{Tot}(\check{\mathcal{C}}^\bullet(\mathcal{U}, \mathcal{F}^\bullet))
$$ est un quasi-isomorphisme d'après le lemme \ref{cohomology-lemma-alternating-usual}. En effet, le morphisme de complexes doubles $\check{\mathcal{C}}^\bullet_{alt}(\mathcal{U}, \mathcal{F}^\bullet) \to
\check{\mathcal{C}}^\bullet(\mathcal{U}, \mathcal{F}^\bullet)$ induit un isomorphisme entre les premières pages des secondes suites spectrales associées à ces complexes (par Homologie, lemme \ref{homology-lemma-ss-double-complex}), et ces suites spectrales convergent (Homologie, lemme \ref{homology-lemma-first-quadrant-ss}). La conclusion découle donc, dans ce cas, du lemme \ref{cohomology-lemma-cech-complex-complex-computes} et de l'hypothèse (3)(a).

\medskip\noindent

En général, par l'hypothèse (3)c), nous pouvons choisir une résolution
$\mathcal{F}^\bullet \to \mathcal{I}^\bullet = \lim \mathcal{I}_n^\bullet$
comme dans le lemme \ref{cohomology-lemma-K-injective}. Le morphisme du lemme devient alors
$$
\lim_n 
\text{Tot}(\check{\mathcal{C}}^\bullet_{alt}(\mathcal{U},
\tau_{\geq -n}\mathcal{F}^\bullet))
\longrightarrow
\Gamma(X, \mathcal{I}^\bullet) =
\lim_n \Gamma(X, \mathcal{I}_n^\bullet)
$$
Ici, la flèche appartient à la catégorie dérivée, mais l'égalité de droite vaut au niveau des complexes. Remarquons que l'hypothèse (3)(b) montre que $\tau_{\geq -n}\mathcal{F}^\bullet$
est un complexe borné inférieurement satisfaisant aux hypothèses du lemme. Ainsi, le cas des complexes bornés inférieurement montre que chacun des morphismes
$$
\text{Tot}(\check{\mathcal{C}}^\bullet_{alt}(\mathcal{U},
\tau_{\geq -n}\mathcal{F}^\bullet))
\longrightarrow
\Gamma(X, \mathcal{I}_n^\bullet)
$$
est un quasi-isomorphisme. En chaque degré donné, les cohomologies des complexes du membre de gauche sont stationnaires (car le complexe de {\v C}ech alterné est fini). Il en va donc de même pour le membre de droite. La cohomologie de la limite du membre de droite est par conséquent égale à cette valeur stationnaire, d'après Homologie, lemme \ref{homology-lemma-apply-Mittag-Leffler-again}
(ou, plus précisément, Compléments sur l'algèbre, lemme
\ref{more-algebra-lemma-apply-Mittag-Leffler-again}), ce qui démontre le résultat.

\end{proof}





\section{Complexes de Hom}

\label{cohomology-section-hom-complexes}

\noindent

Soit $(X, \mathcal{O}_X)$ un espace annelé. Soient
$\mathcal{L}^\bullet$ et $\mathcal{M}^\bullet$ deux complexes de $\mathcal{O}_X$-modules. Nous construisons un complexe de $\mathcal{O}_X$-modules
$\SheafHom^\bullet(\mathcal{L}^\bullet, \mathcal{M}^\bullet)$. Plus précisément, pour tout $n$, posons
$$
\SheafHom^n(\mathcal{L}^\bullet, \mathcal{M}^\bullet) =
\prod\nolimits_{n = p + q}
\SheafHom_{\mathcal{O}_X}(\mathcal{L}^{-q}, \mathcal{M}^p)
$$
On peut voir $\SheafHom^n$ comme le faisceau de $\mathcal{O}_X$-modules formé des applications $\mathcal{O}_X$-linéaires de $\mathcal{L}^\bullet$ dans $\mathcal{M}^\bullet$
(considérés comme des $\mathcal{O}_X$-modules gradués) qui sont homogènes de degré $n$. Dans cette terminologie, nous définissons le différentiel par la formule
$$
\text{d}(f) =
\text{d}_\mathcal{M} \circ f - (-1)^n f \circ \text{d}_\mathcal{L}
$$
pour
$f \in \SheafHom^n_{\mathcal{O}_X}(\mathcal{L}^\bullet, \mathcal{M}^\bullet)$. Nous omettons la vérification que $\text{d}^2 = 0$. Cette construction est un cas particulier d'Algèbre différentielle graduée, exemple \ref{dga-example-category-complexes}. Il résulte immédiatement de la construction que

\begin{equation}
\label{cohomology-equation-cohomology-hom-complex}
H^n(\Gamma(U, \SheafHom^\bullet(\mathcal{L}^\bullet, \mathcal{M}^\bullet))) =
\Hom_{K(\mathcal{O}_U)}(\mathcal{L}^\bullet, \mathcal{M}^\bullet[n])
\end{equation}

pour tous $n \in \mathbf{Z}$ et tous les ouverts $U \subset X$.

\begin{lemma}

\label{cohomology-lemma-compose}
Soit $(X, \mathcal{O}_X)$ un espace annelé. Étant donnés des complexes $\mathcal{K}^\bullet, \mathcal{L}^\bullet, \mathcal{M}^\bullet$ de $\mathcal{O}_X$-modules, il existe un isomorphisme $$
\SheafHom^\bullet(\mathcal{K}^\bullet,
\SheafHom^\bullet(\mathcal{L}^\bullet, \mathcal{M}^\bullet))
=
\SheafHom^\bullet(\text{Tot}(\mathcal{K}^\bullet \otimes_{\mathcal{O}_X}
\mathcal{L}^\bullet), \mathcal{M}^\bullet)
$$ de complexes de $\mathcal{O}_X$-modules, fonctoriel en $\mathcal{K}^\bullet, \mathcal{L}^\bullet, \mathcal{M}^\bullet$.
\end{lemma}

\begin{proof}
Démonstration omise. Indication : cela se démontre exactement de la même manière que Compléments sur l'algèbre, lemme \ref{more-algebra-lemma-compose}.
\end{proof}

\begin{lemma}

\label{cohomology-lemma-composition}
Soit $(X, \mathcal{O}_X)$ un espace annelé. Étant donnés des complexes $\mathcal{K}^\bullet, \mathcal{L}^\bullet, \mathcal{M}^\bullet$ de $\mathcal{O}_X$-modules, il existe un morphisme canonique $$
\text{Tot}\left(
\SheafHom^\bullet(\mathcal{L}^\bullet, \mathcal{M}^\bullet)
\otimes_{\mathcal{O}_X}
\SheafHom^\bullet(\mathcal{K}^\bullet, \mathcal{L}^\bullet)
\right)
\longrightarrow
\SheafHom^\bullet(\mathcal{K}^\bullet, \mathcal{M}^\bullet)
$$ de complexes de $\mathcal{O}_X$-modules.
\end{lemma}

\begin{proof}
Démonstration omise. Indication : cela se démontre exactement de la même manière que Compléments sur l'algèbre, lemme \ref{more-algebra-lemma-composition}.
\end{proof}

\begin{lemma}

\label{cohomology-lemma-diagonal-better}

Soit $(X, \mathcal{O}_X)$ un espace annelé. Étant donnés des complexes
$\mathcal{K}^\bullet, \mathcal{L}^\bullet, \mathcal{M}^\bullet$
de $\mathcal{O}_X$-modules, il existe un morphisme canonique
$$
\text{Tot}\left(
\mathcal{K}^\bullet \otimes_{\mathcal{O}_X}
\SheafHom^\bullet(\mathcal{M}^\bullet, \mathcal{L}^\bullet)
\right)
\longrightarrow
\SheafHom^\bullet(\mathcal{M}^\bullet,
\text{Tot}(\mathcal{K}^\bullet \otimes_{\mathcal{O}_X} \mathcal{L}^\bullet))
$$
de complexes de $\mathcal{O}_X$-modules, fonctoriel en les trois complexes.

\end{lemma}

\begin{proof}
Démonstration omise. Indication : cela se démontre exactement de la même manière que Compléments sur l'algèbre, lemme \ref{more-algebra-lemma-diagonal-better}.
\end{proof}

\begin{lemma}

\label{cohomology-lemma-diagonal}

Soit $(X, \mathcal{O}_X)$ un espace annelé. Étant donnés des complexes
$\mathcal{K}^\bullet, \mathcal{L}^\bullet$
de $\mathcal{O}_X$-modules, il existe un morphisme canonique
$$
\mathcal{K}^\bullet
\longrightarrow
\SheafHom^\bullet(\mathcal{L}^\bullet,
\text{Tot}(\mathcal{K}^\bullet \otimes_{\mathcal{O}_X} \mathcal{L}^\bullet))
$$
de complexes de $\mathcal{O}_X$-modules, fonctoriel en les deux complexes.

\end{lemma}

\begin{proof}
Démonstration omise. Indication : cela se démontre exactement de la même manière que Compléments sur l'algèbre, lemme \ref{more-algebra-lemma-diagonal}.
\end{proof}

\begin{lemma}

\label{cohomology-lemma-evaluate-and-more}

Soit $(X, \mathcal{O}_X)$ un espace annelé. Étant donnés des complexes
$\mathcal{K}^\bullet, \mathcal{L}^\bullet, \mathcal{M}^\bullet$
de $\mathcal{O}_X$-modules, il existe un morphisme canonique
$$
\text{Tot}(\SheafHom^\bullet(\mathcal{L}^\bullet,
\mathcal{M}^\bullet) \otimes_{\mathcal{O}_X} \mathcal{K}^\bullet)
\longrightarrow
\SheafHom^\bullet(\SheafHom^\bullet(\mathcal{K}^\bullet,
\mathcal{L}^\bullet), \mathcal{M}^\bullet)
$$
de complexes de $\mathcal{O}_X$-modules, fonctoriel en les trois complexes.

\end{lemma}

\begin{proof}
Démonstration omise. Indication : cela se démontre exactement de la même manière que Compléments sur l'algèbre, lemme \ref{more-algebra-lemma-evaluate-and-more}.
\end{proof}

\begin{lemma}

\label{cohomology-lemma-RHom-into-K-injective}
Soit $(X, \mathcal{O}_X)$ un espace annelé. Soient $L$ et $M$ des objets de $D(\mathcal{O}_X)$. Soit $\mathcal{I}^\bullet$ un complexe K-injectif de $\mathcal{O}_X$-modules représentant $M$. Soit $\mathcal{L}^\bullet$ un complexe de $\mathcal{O}_X$-modules représentant $L$. Alors $$
H^0(\Gamma(U, \SheafHom^\bullet(\mathcal{L}^\bullet, \mathcal{I}^\bullet))) =
\Hom_{D(\mathcal{O}_U)}(L|_U, M|_U)
$$ pour tout ouvert $U \subset X$.
\end{lemma}

\begin{proof}

Nous avons

\begin{align*}
H^0(\Gamma(U, \SheafHom^\bullet(\mathcal{L}^\bullet, \mathcal{I}^\bullet)))
& =
\Hom_{K(\mathcal{O}_U)}(\mathcal{L}^\bullet|_U, \mathcal{I}^\bullet|_U) \\
& =
\Hom_{D(\mathcal{O}_U)}(L|_U, M|_U)
\end{align*}

La première égalité est (\ref{cohomology-equation-cohomology-hom-complex}). La seconde est vraie parce que $\mathcal{I}^\bullet|_U$
est K-injectif d'après le lemme \ref{cohomology-lemma-restrict-K-injective-to-open}.

\end{proof}

\begin{lemma}

\label{cohomology-lemma-RHom-well-defined}

Soit $(X, \mathcal{O}_X)$ un espace annelé. Soit
$(\mathcal{I}')^\bullet \to \mathcal{I}^\bullet$
un quasi-isomorphisme de complexes K-injectifs de $\mathcal{O}_X$-modules. Soit $(\mathcal{L}')^\bullet \to \mathcal{L}^\bullet$
un quasi-isomorphisme de complexes de $\mathcal{O}_X$-modules. Alors
$$
\SheafHom^\bullet(\mathcal{L}^\bullet, (\mathcal{I}')^\bullet)
\longrightarrow
\SheafHom^\bullet((\mathcal{L}')^\bullet, \mathcal{I}^\bullet)
$$
est un quasi-isomorphisme.

\end{lemma}

\begin{proof}
Soit $M$ l'objet de $D(\mathcal{O}_X)$ représenté par $\mathcal{I}^\bullet$ et $(\mathcal{I}')^\bullet$. Soit $L$ l'objet de $D(\mathcal{O}_X)$ représenté par $\mathcal{L}^\bullet$ et $(\mathcal{L}')^\bullet$. D'après le lemme \ref{cohomology-lemma-RHom-into-K-injective}, les faisceaux $$
H^0(\SheafHom^\bullet(\mathcal{L}^\bullet, (\mathcal{I}')^\bullet))
\quad\text{et}\quad
H^0(\SheafHom^\bullet((\mathcal{L}')^\bullet, \mathcal{I}^\bullet))
$$ sont tous deux égaux au faisceau associé au préfaisceau $$
U \longmapsto \Hom_{D(\mathcal{O}_U)}(L|_U, M|_U)
$$ Le morphisme est donc un quasi-isomorphisme.
\end{proof}

\begin{lemma}

\label{cohomology-lemma-RHom-from-K-flat-into-K-injective}
Soit $(X, \mathcal{O}_X)$ un espace annelé. Soit $\mathcal{I}^\bullet$ un complexe K-injectif de $\mathcal{O}_X$-modules. Soit $\mathcal{L}^\bullet$ un complexe K-plat de $\mathcal{O}_X$-modules. Alors $\SheafHom^\bullet(\mathcal{L}^\bullet, \mathcal{I}^\bullet)$ est un complexe K-injectif de $\mathcal{O}_X$-modules.
\end{lemma}

\begin{proof}

En effet, si $\mathcal{K}^\bullet$ est un complexe acyclique de
$\mathcal{O}_X$-modules, alors

\begin{align*}
\Hom_{K(\mathcal{O}_X)}(\mathcal{K}^\bullet,
\SheafHom^\bullet(\mathcal{L}^\bullet, \mathcal{I}^\bullet))
& =
H^0(\Gamma(X,
\SheafHom^\bullet(\mathcal{K}^\bullet,
\SheafHom^\bullet(\mathcal{L}^\bullet, \mathcal{I}^\bullet)))) \\
& =
H^0(\Gamma(X, \SheafHom^\bullet(\text{Tot}(
\mathcal{K}^\bullet \otimes_{\mathcal{O}_X} \mathcal{L}^\bullet),
\mathcal{I}^\bullet))) \\
& =
\Hom_{K(\mathcal{O}_X)}(
\text{Tot}(\mathcal{K}^\bullet \otimes_{\mathcal{O}_X} \mathcal{L}^\bullet),
\mathcal{I}^\bullet) \\
& =
0
\end{align*}

La première égalité découle de (\ref{cohomology-equation-cohomology-hom-complex}), la deuxième du lemme \ref{cohomology-lemma-compose} et la troisième de (\ref{cohomology-equation-cohomology-hom-complex}). La dernière égalité vaut parce que
$\text{Tot}(\mathcal{K}^\bullet \otimes_{\mathcal{O}_X} \mathcal{L}^\bullet)$
est acyclique puisque $\mathcal{L}^\bullet$ est K-plat (définition \ref{cohomology-definition-K-flat}) et puisque $\mathcal{I}^\bullet$
est K-injectif.

\end{proof}







\section{Hom interne dans la catégorie dérivée}

\label{cohomology-section-internal-hom}

\noindent

Soit $(X, \mathcal{O}_X)$ un espace annelé. Soient $L, M$ des objets de $D(\mathcal{O}_X)$. Nous voulons construire un objet
$R\SheafHom(L, M)$ de $D(\mathcal{O}_X)$ tel que, pour tout autre objet $K$ de $D(\mathcal{O}_X)$, il existe une bijection canonique

\begin{equation}
\label{cohomology-equation-internal-hom}
\Hom_{D(\mathcal{O}_X)}(K, R\SheafHom(L, M))
=
\Hom_{D(\mathcal{O}_X)}(K \otimes_{\mathcal{O}_X}^\mathbf{L} L, M)
\end{equation}

Remarquons que cette formule définit $R\SheafHom(L, M)$ à isomorphisme unique près, par le lemme de Yoneda (Catégories, lemme \ref{categories-lemma-yoneda}).

\medskip\noindent
Pour construire un tel objet, choisissons un complexe K-injectif $\mathcal{I}^\bullet$ représentant $M$ et un complexe quelconque de $\mathcal{O}_X$-modules $\mathcal{L}^\bullet$ représentant $L$. Posons alors $$
R\SheafHom(L, M) = \SheafHom^\bullet(\mathcal{L}^\bullet, \mathcal{I}^\bullet)
$$ où le membre de droite est le complexe de $\mathcal{O}_X$-modules construit dans la section \ref{cohomology-section-hom-complexes}. Cette définition est bien posée d'après le lemme \ref{cohomology-lemma-RHom-well-defined}. Nous obtenons un foncteur $$
D(\mathcal{O}_X)^{opp} \times D(\mathcal{O}_X) \longrightarrow D(\mathcal{O}_X),
\quad
(K, L) \longmapsto R\SheafHom(K, L)
$$ Avant de démontrer (\ref{cohomology-equation-internal-hom}), calculons les groupes de cohomologie de $R\SheafHom(K, L)$.

\begin{lemma}

\label{cohomology-lemma-section-RHom-over-U}
Soit $(X, \mathcal{O}_X)$ un espace annelé. Soient $L, M$ des objets de $D(\mathcal{O}_X)$. Pour tout ouvert $U$, on a $$
H^0(U, R\SheafHom(L, M)) =
\Hom_{D(\mathcal{O}_U)}(L|_U, M|_U)
$$ et en particulier $H^0(X, R\SheafHom(L, M)) = \Hom_{D(\mathcal{O}_X)}(L, M)$.
\end{lemma}

\begin{proof}
Choisissons un complexe K-injectif $\mathcal{I}^\bullet$ de $\mathcal{O}_X$-modules représentant $M$ et un complexe K-plat $\mathcal{L}^\bullet$ représentant $L$. Alors $\SheafHom^\bullet(\mathcal{L}^\bullet, \mathcal{I}^\bullet)$ est K-injectif d'après le lemme \ref{cohomology-lemma-RHom-from-K-flat-into-K-injective}. Nous pouvons donc calculer la cohomologie sur $U$ en prenant simplement les sections sur $U$, et le résultat découle du lemme \ref{cohomology-lemma-RHom-into-K-injective}.
\end{proof}

\begin{lemma}

\label{cohomology-lemma-internal-hom}
Soit $(X, \mathcal{O}_X)$ un espace annelé. Soient $K, L, M$ des objets de $D(\mathcal{O}_X)$. Avec la construction décrite ci-dessus, il existe un isomorphisme canonique $$
R\SheafHom(K, R\SheafHom(L, M)) =
R\SheafHom(K \otimes_{\mathcal{O}_X}^\mathbf{L} L, M)
$$ dans $D(\mathcal{O}_X)$, fonctoriel en $K, L, M$, qui redonne (\ref{cohomology-equation-internal-hom}) après application de $H^0(X, -)$.
\end{lemma}

\begin{proof}

Choisissons un complexe K-injectif $\mathcal{I}^\bullet$ représentant
$M$ et un complexe K-plat de $\mathcal{O}_X$-modules $\mathcal{L}^\bullet$
représentant $L$. Soit $\mathcal{K}^\bullet$ un complexe quelconque de
$\mathcal{O}_X$-modules représentant $K$. Alors
$$
\SheafHom^\bullet(\mathcal{K}^\bullet,
\SheafHom^\bullet(\mathcal{L}^\bullet, \mathcal{I}^\bullet))
=
\SheafHom^\bullet(
\text{Tot}(\mathcal{K}^\bullet \otimes_{\mathcal{O}_X} \mathcal{L}^\bullet),
\mathcal{I}^\bullet)
$$
d'après le lemme \ref{cohomology-lemma-compose}. Remarquons que le membre de gauche représente
$R\SheafHom(K, R\SheafHom(L, M))$ (utiliser le lemme \ref{cohomology-lemma-RHom-from-K-flat-into-K-injective}) et que le membre de droite représente
$R\SheafHom(K \otimes_{\mathcal{O}_X}^\mathbf{L} L, M)$. Cela démontre la formule du lemme. En prenant les sections globales et en utilisant le lemme \ref{cohomology-lemma-section-RHom-over-U},
nous obtenons (\ref{cohomology-equation-internal-hom}).

\end{proof}

\begin{lemma}

\label{cohomology-lemma-restriction-RHom-to-U}
Soit $(X, \mathcal{O}_X)$ un espace annelé. Soient $K, L$ des objets de $D(\mathcal{O}_X)$. La construction de $R\SheafHom(K, L)$ commute à la restriction aux ouverts, c'est-à-dire que, pour tout ouvert $U$, on a $R\SheafHom(K|_U, L|_U) = R\SheafHom(K, L)|_U$.
\end{lemma}

\begin{proof}
Cela résulte immédiatement de la construction et du lemme \ref{cohomology-lemma-restrict-K-injective-to-open}.
\end{proof}

\begin{lemma}

\label{cohomology-lemma-RHom-triangulated}

Soit $(X, \mathcal{O}_X)$ un espace annelé. Le bifoncteur $R\SheafHom(- , -)$
transforme les triangles distingués en triangles distingués dans les deux variables.

\end{lemma}

\begin{proof}

Cela résulte du fait que l'application
$$
(\mathcal{L}^\bullet, \mathcal{M}^\bullet) \longmapsto
\SheafHom^\bullet(\mathcal{L}^\bullet, \mathcal{M}^\bullet)
$$
transforme une suite exacte courte de complexes scindée terme à terme, dans l'une ou l'autre variable, en une suite exacte courte scindée terme à terme. Nous omettons les détails.

\end{proof}

\begin{lemma}

\label{cohomology-lemma-internal-hom-composition}
Soit $(X, \mathcal{O}_X)$ un espace annelé. Étant donnés $K, L, M$ dans $D(\mathcal{O}_X)$, il existe un morphisme canonique $$
R\SheafHom(L, M) \otimes_{\mathcal{O}_X}^\mathbf{L} R\SheafHom(K, L)
\longrightarrow R\SheafHom(K, M)
$$ dans $D(\mathcal{O}_X)$, fonctoriel en $K, L, M$.
\end{lemma}

\begin{proof}

Choisissons un complexe K-injectif $\mathcal{I}^\bullet$ représentant $M$, un complexe K-injectif $\mathcal{J}^\bullet$ représentant $L$, et un complexe quelconque de $\mathcal{O}_X$-modules $\mathcal{K}^\bullet$ représentant $K$. D'après le lemme \ref{cohomology-lemma-composition}, il existe un morphisme de complexes
$$
\text{Tot}\left(
\SheafHom^\bullet(\mathcal{J}^\bullet, \mathcal{I}^\bullet)
\otimes_{\mathcal{O}_X}
\SheafHom^\bullet(\mathcal{K}^\bullet, \mathcal{J}^\bullet)
\right)
\longrightarrow
\SheafHom^\bullet(\mathcal{K}^\bullet, \mathcal{I}^\bullet)
$$
Les complexes de $\mathcal{O}_X$-modules
$\SheafHom^\bullet(\mathcal{J}^\bullet, \mathcal{I}^\bullet)$,
$\SheafHom^\bullet(\mathcal{K}^\bullet, \mathcal{J}^\bullet)$,
$\SheafHom^\bullet(\mathcal{K}^\bullet, \mathcal{I}^\bullet)$
représentent respectivement $R\SheafHom(L, M)$, $R\SheafHom(K, L)$ et $R\SheafHom(K, M)$. Si nous choisissons un complexe K-plat $\mathcal{H}^\bullet$ et un quasi-isomorphisme
$\mathcal{H}^\bullet \to
\SheafHom^\bullet(\mathcal{K}^\bullet, \mathcal{J}^\bullet)$, alors il existe un morphisme
$$
\text{Tot}\left(
\SheafHom^\bullet(\mathcal{J}^\bullet, \mathcal{I}^\bullet)
\otimes_{\mathcal{O}_X} \mathcal{H}^\bullet
\right)
\longrightarrow
\text{Tot}\left(
\SheafHom^\bullet(\mathcal{J}^\bullet, \mathcal{I}^\bullet)
\otimes_{\mathcal{O}_X}
\SheafHom^\bullet(\mathcal{K}^\bullet, \mathcal{J}^\bullet)
\right)
$$
dont la source représente
$R\SheafHom(L, M) \otimes_{\mathcal{O}_X}^\mathbf{L} R\SheafHom(K, L)$. En composant les deux flèches affichées, on obtient le morphisme souhaité. Nous omettons la preuve de la fonctorialité de la construction.

\end{proof}

\begin{lemma}

\label{cohomology-lemma-internal-hom-diagonal-better}
Soit $(X, \mathcal{O}_X)$ un espace annelé. Étant donnés $K, L, M$ dans $D(\mathcal{O}_X)$, il existe un morphisme canonique $$
K \otimes_{\mathcal{O}_X}^\mathbf{L} R\SheafHom(M, L)
\longrightarrow
R\SheafHom(M, K \otimes_{\mathcal{O}_X}^\mathbf{L} L)
$$ dans $D(\mathcal{O}_X)$, fonctoriel en $K, L, M$.
\end{lemma}

\begin{proof}
Choisissons un complexe K-plat $\mathcal{K}^\bullet$ représentant $K$, un complexe K-injectif $\mathcal{I}^\bullet$ représentant $L$, et un complexe quelconque de $\mathcal{O}_X$-modules $\mathcal{M}^\bullet$ représentant $M$. Choisissons un quasi-isomorphisme $\text{Tot}(\mathcal{K}^\bullet \otimes_{\mathcal{O}_X} \mathcal{I}^\bullet)
\to \mathcal{J}^\bullet$, où $\mathcal{J}^\bullet$ est K-injectif. Nous utilisons alors le morphisme $$
\text{Tot}\left(
\mathcal{K}^\bullet \otimes_{\mathcal{O}_X}
\SheafHom^\bullet(\mathcal{M}^\bullet, \mathcal{I}^\bullet)
\right)
\to
\SheafHom^\bullet(\mathcal{M}^\bullet,
\text{Tot}(\mathcal{K}^\bullet \otimes_{\mathcal{O}_X} \mathcal{I}^\bullet))
\to
\SheafHom^\bullet(\mathcal{M}^\bullet, \mathcal{J}^\bullet)
$$ où la première flèche est le morphisme du lemme \ref{cohomology-lemma-diagonal-better}.
\end{proof}

\begin{lemma}

\label{cohomology-lemma-internal-hom-diagonal}
Soit $(X, \mathcal{O}_X)$ un espace annelé. Étant donnés $K, L$ dans $D(\mathcal{O}_X)$, il existe un morphisme canonique $$
K \longrightarrow R\SheafHom(L, K \otimes_{\mathcal{O}_X}^\mathbf{L} L)
$$ dans $D(\mathcal{O}_X)$, fonctoriel en $K$ et $L$.
\end{lemma}

\begin{proof}
Choisissons un complexe K-plat $\mathcal{K}^\bullet$ représentant $K$ et un complexe quelconque de $\mathcal{O}_X$-modules $\mathcal{L}^\bullet$ représentant $L$. Choisissons un complexe K-injectif $\mathcal{J}^\bullet$ et un quasi-isomorphisme $\text{Tot}(\mathcal{K}^\bullet \otimes_{\mathcal{O}_X} \mathcal{L}^\bullet)
\to \mathcal{J}^\bullet$. Nous utilisons alors $$
\mathcal{K}^\bullet \to
\SheafHom^\bullet(\mathcal{L}^\bullet,
\text{Tot}(\mathcal{K}^\bullet \otimes_{\mathcal{O}_X} \mathcal{L}^\bullet))
\to
\SheafHom^\bullet(\mathcal{L}^\bullet, \mathcal{J}^\bullet)
$$ où la première flèche provient du lemme \ref{cohomology-lemma-diagonal}.
\end{proof}

\begin{lemma}

\label{cohomology-lemma-dual}

Soit $(X, \mathcal{O}_X)$ un espace annelé. Soit $L$ un objet de $D(\mathcal{O}_X)$. Posons $L^\vee = R\SheafHom(L, \mathcal{O}_X)$. Pour $M$ dans $D(\mathcal{O}_X)$, il existe un morphisme canonique

\begin{equation}
\label{cohomology-equation-eval}
M \otimes^\mathbf{L}_{\mathcal{O}_X} L^\vee
\longrightarrow
R\SheafHom(L, M)
\end{equation}

qui induit un morphisme canonique
$$
H^0(X, M \otimes^\mathbf{L}_{\mathcal{O}_X} L^\vee)
\longrightarrow
\Hom_{D(\mathcal{O}_X)}(L, M)
$$
fonctoriel en $M$ dans $D(\mathcal{O}_X)$.

\end{lemma}

\begin{proof}
Le morphisme (\ref{cohomology-equation-eval}) est un cas particulier du lemme \ref{cohomology-lemma-internal-hom-composition}, en utilisant l'identification $M = R\SheafHom(\mathcal{O}_X, M)$.
\end{proof}

\begin{lemma}

\label{cohomology-lemma-internal-hom-evaluate}
Soit $(X, \mathcal{O}_X)$ un espace annelé. Soient $K, L, M$ des objets de $D(\mathcal{O}_X)$. Il existe un morphisme canonique $$
R\SheafHom(L, M) \otimes_{\mathcal{O}_X}^\mathbf{L} K
\longrightarrow
R\SheafHom(R\SheafHom(K, L), M)
$$ dans $D(\mathcal{O}_X)$, fonctoriel en $K, L, M$.
\end{lemma}

\begin{proof}
Choisissons un complexe K-injectif $\mathcal{I}^\bullet$ représentant $M$, un complexe K-injectif $\mathcal{J}^\bullet$ représentant $L$, et un complexe K-plat $\mathcal{K}^\bullet$ représentant $K$. Le morphisme est défini au moyen du morphisme $$
\text{Tot}(\SheafHom^\bullet(\mathcal{J}^\bullet,
\mathcal{I}^\bullet) \otimes_{\mathcal{O}_X} \mathcal{K}^\bullet)
\longrightarrow
\SheafHom^\bullet(\SheafHom^\bullet(\mathcal{K}^\bullet,
\mathcal{J}^\bullet), \mathcal{I}^\bullet)
$$ du lemme \ref{cohomology-lemma-evaluate-and-more}. Par notre choix particulier des complexes, le membre de gauche représente $R\SheafHom(L, M) \otimes_{\mathcal{O}_X}^\mathbf{L} K$ et le membre de droite représente $R\SheafHom(R\SheafHom(K, L), M)$. Nous omettons la preuve de la fonctorialité en les trois objets de $D(\mathcal{O}_X)$.
\end{proof}

\begin{remark}

\label{cohomology-remark-tensor-internal-hom}

Soit $(X, \mathcal{O}_X)$ un espace annelé. Pour $K, K', M, M'$ dans
$D(\mathcal{O}_X)$, il existe un morphisme canonique
$$
R\SheafHom(K, K') \otimes_{\mathcal{O}_X}^\mathbf{L}
R\SheafHom(M, M')
\longrightarrow
R\SheafHom(K \otimes_{\mathcal{O}_X}^\mathbf{L} M,
K' \otimes_{\mathcal{O}_X}^\mathbf{L} M')
$$
En effet, d'après (\ref{cohomology-equation-internal-hom}), cela revient à construire un morphisme
$$
R\SheafHom(K, K') \otimes_{\mathcal{O}_X}^\mathbf{L}
R\SheafHom(M, M') \otimes_{\mathcal{O}_X}^\mathbf{L}
K \otimes_{\mathcal{O}_X}^\mathbf{L} M
\longrightarrow
K' \otimes_{\mathcal{O}_X}^\mathbf{L} M'
$$
Pour cela, nous pouvons d'abord permuter les deux facteurs du milieu (avec les règles de signes de Compléments sur l'algèbre, section
\ref{more-algebra-section-sign-rules}), puis utiliser les morphismes
$$
R\SheafHom(K, K') \otimes_{\mathcal{O}_X}^\mathbf{L} K \to K'
\quad\text{et}\quad
R\SheafHom(M, M') \otimes_{\mathcal{O}_X}^\mathbf{L} M \to M'
$$
du lemme \ref{cohomology-lemma-internal-hom-composition}, en considérant $K = R\SheafHom(\mathcal{O}_X, K)$ et de même pour
$K'$, $M$ et $M'$.

\end{remark}

\begin{remark}

\label{cohomology-remark-projection-formula-for-internal-hom}

Soit $f : X \to Y$ un morphisme d'espaces annelés. Soient $K, L$ des objets de $D(\mathcal{O}_X)$. Nous affirmons qu'il existe un morphisme canonique
$$
Rf_*R\SheafHom(L, K) \longrightarrow R\SheafHom(Rf_*L, Rf_*K)
$$
En effet, d'après (\ref{cohomology-equation-internal-hom}), cela revient à construire un morphisme
$Rf_*R\SheafHom(L, K) \otimes_{\mathcal{O}_Y}^\mathbf{L} Rf_*L \to Rf_*K$. Pour cela, nous pouvons utiliser la composition
$$
Rf_*R\SheafHom(L, K) \otimes_{\mathcal{O}_Y}^\mathbf{L} Rf_*L \to
Rf_*(R\SheafHom(L, K) \otimes_{\mathcal{O}_X}^\mathbf{L} L) \to
Rf_*K
$$
où la première flèche est le cup-produit relatif (remarque \ref{cohomology-remark-cup-product}) et la seconde est $Rf_*$ appliqué au morphisme canonique
$R\SheafHom(L, K) \otimes_{\mathcal{O}_X}^\mathbf{L} L \to K$
provenant du lemme \ref{cohomology-lemma-internal-hom-composition}
(avec $\mathcal{O}_X$ dans l'un des arguments).

\end{remark}

\begin{remark}

\label{cohomology-remark-relative-cup-and-composition}

Soit $h : X \to Y$ un morphisme d'espaces annelés. Soient $K, M$ des objets de $D(\mathcal{O}_Y)$. Le schéma
$$
\xymatrix{
Rf_*R\SheafHom_{\mathcal{O}_X}(K, M)
\otimes_{\mathcal{O}_Y}^\mathbf{L} Rf_*K
\ar[r] \ar[d] &
Rf_*\left(R\SheafHom_{\mathcal{O}_X}(K, M)
\otimes_{\mathcal{O}_X}^\mathbf{L} K\right)
\ar[d] \\
R\SheafHom_{\mathcal{O}_Y}(Rf_*K, Rf_*M) \otimes_{\mathcal{O}_Y}^\mathbf{L}
Rf_*K \ar[r] &
Rf_*M
}
$$
est commutatif. Ici, la flèche verticale de gauche provient de la remarque \ref{cohomology-remark-projection-formula-for-internal-hom}. La flèche horizontale supérieure est celle de la remarque \ref{cohomology-remark-cup-product}. Les deux autres flèches sont des cas particuliers du morphisme du lemme \ref{cohomology-lemma-internal-hom-composition} (où l'une des entrées est remplacée par $\mathcal{O}_X$ ou $\mathcal{O}_Y$).

\end{remark}

\begin{remark}

\label{cohomology-remark-prepare-fancy-base-change}

Soit $h : X \to Y$ un morphisme d'espaces annelés. Soient $K, L$ des objets de $D(\mathcal{O}_Y)$. Nous affirmons qu'il existe un morphisme canonique
$$
Lh^*R\SheafHom(K, L) \longrightarrow R\SheafHom(Lh^*K, Lh^*L)
$$
dans $D(\mathcal{O}_X)$. En effet, d'après (\ref{cohomology-equation-internal-hom}), démontré dans le lemme \ref{cohomology-lemma-internal-hom},
un tel morphisme revient à construire un morphisme
$$
Lh^*R\SheafHom(K, L) \otimes^\mathbf{L} Lh^*K \longrightarrow Lh^*L
$$
La source de cette flèche est $Lh^*(\SheafHom(K, L) \otimes^\mathbf{L} K)$
d'après le lemme \ref{cohomology-lemma-pullback-tensor-product}
donc il suffit de construire un morphisme canonique
$$
R\SheafHom(K, L) \otimes^\mathbf{L} K \longrightarrow L.
$$
Pour cela, nous prenons la flèche correspondant à
$$
\text{id} :
R\SheafHom(K, L)
\longrightarrow
R\SheafHom(K, L)
$$
via (\ref{cohomology-equation-internal-hom}).

\end{remark}

\begin{remark}

\label{cohomology-remark-fancy-base-change}

Supposons que
$$
\xymatrix{
X' \ar[r]_h \ar[d]_{f'} &
X \ar[d]^f \\
S' \ar[r]^g &
S
}
$$
est un diagramme commutatif d'espaces annelés. Soient $K, L$ des objets de $D(\mathcal{O}_X)$. Nous affirmons qu'il existe un morphisme canonique de changement de base
$$
Lg^*Rf_*R\SheafHom(K, L)
\longrightarrow
R(f')_*R\SheafHom(Lh^*K, Lh^*L)
$$
dans $D(\mathcal{O}_{S'})$. Plus précisément, nous prenons le morphisme adjoint de la composition

\begin{align*}
L(f')^*Lg^*Rf_*R\SheafHom(K, L)
& =
Lh^*Lf^*Rf_*R\SheafHom(K, L) \\
& \to
Lh^*R\SheafHom(K, L) \\
& \to
R\SheafHom(Lh^*K, Lh^*L)
\end{align*}

où la première flèche utilise le morphisme d'adjonction
$Lf^*Rf_* \to \text{id}$ et la seconde est le morphisme canonique construit dans la remarque \ref{cohomology-remark-prepare-fancy-base-change}.

\end{remark}
\section{Faisceaux d'Ext}

\label{cohomology-section-ext}

\noindent

Soit $(X, \mathcal{O}_X)$ un espace annelé et soient $K, L \in D(\mathcal{O}_X)$. La construction du Hom interne dans la catégorie dérivée fournit un faisceau bien défini de $\mathcal{O}_X$-modules
$$
\SheafExt^n(K, L) = H^n(R\SheafHom(K, L))
$$
obtenu en prenant le $n$-ième faisceau de cohomologie de l'objet $R\SheafHom(K, L)$
de $D(\mathcal{O}_X)$. Nous noterons parfois
$\SheafExt^n_{\mathcal{O}_X}(K, L)$ cet objet. Par le lemme \ref{cohomology-lemma-section-RHom-over-U},
ce faisceau $\SheafExt^n$ est le faisceau associé au préfaisceau
$$
U \longmapsto \Ext^n_{D(\mathcal{O}_U)}(K|_U, L|_U)
$$
Par l'exemple \ref{cohomology-example-spectral-sequence}, il existe toujours une suite spectrale
$$
E_2^{p, q} = H^p(X, \SheafExt^q(K, L))
$$
convergeant vers $\Ext^{p + q}_{D(\mathcal{O}_X)}(K, L)$
dans des situations favorables (par exemple si $L$ est borné inférieurement et
$K$ borné supérieurement).





\section{Hom dérivé global}

\label{cohomology-section-global-RHom}

\noindent

Soit $(X, \mathcal{O}_X)$ un espace annelé et soient $K, L \in D(\mathcal{O}_X)$. La construction du Hom interne dans la catégorie dérivée fournit un objet bien défini
$$
R\Hom_X(K, L) = R\Gamma(X, R\SheafHom(K, L))
$$
dans $D(\Gamma(X, \mathcal{O}_X))$. Nous noterons parfois
$R\Hom_{\mathcal{O}_X}(K, L)$ cet objet. Par le lemme \ref{cohomology-lemma-section-RHom-over-U},
nous avons
$$
H^0(R\Hom_X(K, L)) = \Hom_{D(\mathcal{O}_X)}(K, L),
\quad
H^p(R\Hom_X(K, L)) = \Ext_{D(\mathcal{O}_X)}^p(K, L)
$$
Si $f : Y \to X$ est un morphisme d'espaces annelés, il existe une application canonique
$$
R\Hom_X(K, L) \longrightarrow R\Hom_Y(Lf^*K, Lf^*L)
$$
dans $D(\Gamma(X, \mathcal{O}_X))$, obtenue en prenant les sections globales de l'application définie dans la remarque \ref{cohomology-remark-prepare-fancy-base-change}.







\section{Collage de complexes}

\label{cohomology-section-glueing-complexes}

\noindent
Nous pouvons recoller des complexes ! Plus précisément, dans certaines circonstances, des objets de la catégorie dérivée donnés localement se recollent en un objet global. Nous démontrons d'abord quelques cas simples, puis le très général \cite[théorème 3.2.4]{BBD} dans le cadre des espaces topologiques et des recouvrements ouverts.

\begin{lemma}

\label{cohomology-lemma-glue}

Soit $(X, \mathcal{O}_X)$ un espace annelé. Supposons que $X = U \cup V$ soit la réunion de deux sous-espaces ouverts de $X$ et que soient donnés

\begin{enumerate}

\item un objet $A$ de $D(\mathcal{O}_U)$,
\item un objet $B$ de $D(\mathcal{O}_V)$, et
\item un isomorphisme $c : A|_{U \cap V} \to B|_{U \cap V}$.

\end{enumerate}

Alors il existe un objet $F$ de $D(\mathcal{O}_X)$
et des isomorphismes $f : F|_U \to A$, $g : F|_V \to B$ tels que $c = g|_{U \cap V} \circ f^{-1}|_{U \cap V}$. De plus, étant donnés

\begin{enumerate}

\item un objet $E$ de $D(\mathcal{O}_X)$,
\item un morphisme $a : A \to E|_U$ de $D(\mathcal{O}_U)$,
\item un morphisme $b : B \to E|_V$ de $D(\mathcal{O}_V)$, 

\end{enumerate}

tels que
$$
a|_{U \cap V}  = b|_{U \cap V} \circ c.
$$
il existe un morphisme $F \to E$ dans $D(\mathcal{O}_X)$
dont la restriction à $U$ est $a \circ f$
et dont la restriction à $V$ est $b \circ g$.

\end{lemma}

\begin{proof}

Notons $j_U$, $j_V$, $j_{U \cap V}$ les immersions ouvertes correspondantes. Choisissons un triangle distingué
$$
F \to Rj_{U, *}A \oplus Rj_{V, *}B \to Rj_{U \cap V, *}(B|_{U \cap V})
\to F[1]
$$
où l'application $Rj_{V, *}B \to Rj_{U \cap V, *}(B|_{U \cap V})$ est l'application évidente et où
$Rj_{U, *}A \to Rj_{U \cap V, *}(B|_{U \cap V})$
est la composition de
$Rj_{U, *}A \to Rj_{U \cap V, *}(A|_{U \cap V})$
avec $Rj_{U \cap V, *}c$. En restreignant à $U$, nous obtenons
$$
F|_U \to A \oplus (Rj_{V, *}B)|_U \to (Rj_{U \cap V, *}(B|_{U \cap V}))|_U
\to F|_U[1]
$$
Notons $j : U \cap V \to U$. La compatibilité de la restriction aux ouverts avec la cohomologie montre que
$(Rj_{V, *}B)|_U$ et $(Rj_{U \cap V, *}(B|_{U \cap V}))|_U$
sont tous deux canoniquement isomorphes à $Rj_*(B|_{U \cap V})$. La deuxième flèche du dernier diagramme admet donc une section, et le morphisme $F|_U \to A$ est un isomorphisme. De même, le morphisme $F|_V \to B$ est un isomorphisme. L'existence du morphisme $F \to E$ résulte de la suite de Mayer--Vietoris pour $\Hom$ ; voir le lemme \ref{cohomology-lemma-mayer-vietoris-hom}.

\end{proof}

\begin{lemma}

\label{cohomology-lemma-vanishing-and-glueing}

Soit $f : (X, \mathcal{O}_X) \to (Y, \mathcal{O}_Y)$ un morphisme d'espaces annelés. Soit $\mathcal{B}$ une base de la topologie de $Y$.

\begin{enumerate}

\item Supposons que $K$ soit un objet de $D(\mathcal{O}_X)$ tel que, pour $V \in \mathcal{B}$, nous ayons $H^i(f^{-1}(V), K) = 0$ pour $i < 0$. Alors les faisceaux de cohomologie de $Rf_*K$ sont nuls en degrés négatifs,
$H^i(f^{-1}(V), K) = 0$ pour $i < 0$ pour tout ouvert $V \subset Y$, et la règle $V \mapsto H^0(f^{-1}V, K)$ définit un faisceau sur $Y$.
\item Supposons que $K, L$ soient des objets de $D(\mathcal{O}_X)$ tels que, pour $V \in \mathcal{B}$, nous ayons
$\Ext^i(K|_{f^{-1}V}, L|_{f^{-1}V}) = 0$ pour $i < 0$. Alors $\Ext^i(K|_{f^{-1}V}, L|_{f^{-1}V}) = 0$ pour $i < 0$
pour tout ouvert $V \subset Y$, et la règle $V \mapsto \Hom(K|_{f^{-1}V}, L|_{f^{-1}V})$ définit un faisceau sur $Y$.

\end{enumerate}

\end{lemma}

\begin{proof}
Le lemme \ref{cohomology-lemma-unbounded-describe-higher-direct-images} affirme que $H^i(Rf_*K)$ est le faisceau associé au préfaisceau $V \mapsto H^i(f^{-1}(V), K) = H^i(V, Rf_*K)$. Les hypothèses de (1) impliquent que les faisceaux de cohomologie de $Rf_*K$ sont nuls en degrés $< 0$. Pour tout ouvert $V \subset Y$, le groupe de cohomologie $H^i(V, Rf_*K)$ est donc nul pour $i < 0$ et égal à $H^0(V, H^0(Rf_*K))$ pour $i = 0$. Cela démontre (1).

\medskip\noindent

Pour démontrer (2), appliquons (1) au complexe $R\SheafHom(K, L)$ et utilisons le lemme \ref{cohomology-lemma-section-RHom-over-U} pour effectuer la traduction.

\end{proof}

\begin{situation}

\label{cohomology-situation-locally-given}

Soit $(X, \mathcal{O}_X)$ un espace annelé. On se donne

\begin{enumerate}

\item une collection d'ouverts $\mathcal{B}$ de $X$,
\item pour $U \in \mathcal{B}$ un objet $K_U$ en $D(\mathcal{O}_U)$,
\item pour $V \subset U$ avec $V, U \in \mathcal{B}$ un isomorphisme
$\rho^U_V : K_U|_V \to K_V$ en $D(\mathcal{O}_V)$,

\end{enumerate}

de sorte que, pour $W \subset V \subset U$ avec $U, V, W$ dans
$\mathcal{B}$, on ait $\rho^U_W = \rho^V_W \circ \rho ^U_V|_W$.

\end{situation}

\noindent
Nous ne pourrons rien démontrer à ce sujet sans hypothèses supplémentaires. Un cas intéressant est celui où $\mathcal{B}$ est une base de la topologie de $X$. Un autre est celui où l'on dispose d'un morphisme $f : X \to Y$ d'espaces topologiques et où les éléments de $\mathcal{B}$ sont les images réciproques des éléments d'une base de la topologie de $Y$.

\medskip\noindent
Dans la situation \ref{cohomology-situation-locally-given}, une {\it solution} est un couple $(K, \rho_U)$, où $K$ est un objet de $D(\mathcal{O}_X)$ et où les $\rho_U : K|_U \to K_U$, pour $U \in \mathcal{B}$, sont des isomorphismes tels que $\rho^U_V \circ \rho_U|_V = \rho_V$ pour tous $V \subset U$, $U, V \in \mathcal{B}$. Dans certains cas, les solutions sont uniques.

\begin{lemma}

\label{cohomology-lemma-uniqueness}
Dans la situation \ref{cohomology-situation-locally-given}, supposons que
\begin{enumerate}
\item $X = \bigcup_{U \in \mathcal{B}} U$ et pour $U, V \in \mathcal{B}$ nous avons $U \cap V = \bigcup_{W \in \mathcal{B}, W \subset U \cap V} W$,
\item pour tout $U \in \mathcal{B}$, nous ayons $\Ext^i(K_U, K_U) = 0$ pour $i < 0$.
\end{enumerate}
Si une solution $(K, \rho_U)$ existe, elle est unique à un unique isomorphisme près ; de plus, $\Ext^i(K, K) = 0$ pour $i < 0$.
\end{lemma}

\begin{proof}

Soient $(K, \rho_U)$ et $(K', \rho'_U)$ deux solutions. Soit $f : X \to Y$ l'application continue construite dans Topologie, lemme \ref{topology-lemma-create-map-from-subcollection}. Posons $\mathcal{O}_Y = f_*\mathcal{O}_X$. Alors $K, K'$ et $\mathcal{B}$ satisfont les hypothèses de la partie (2) du lemme \ref{cohomology-lemma-vanishing-and-glueing}. Nous obtenons donc l'annulation des Ext négatifs de $K$, et la règle
$$
V \longmapsto \Hom(K|_{f^{-1}V}, K'|_{f^{-1}V})
$$
définit un faisceau sur $Y$. Comme $(K, \rho_U)$ et $(K', \rho'_U)$ sont deux solutions, les applications
$$
(\rho'_U)^{-1} \circ \rho_U : K|_U \longrightarrow K'|_U
$$
sur $U = f^{-1}(f(U))$ coïncident sur les intersections. Nous obtenons donc une unique section globale du faisceau ci-dessus ; elle définit l'isomorphisme recherché
$K \to K'$, compatible avec toutes les données.

\end{proof}

\begin{remark}

\label{cohomology-remark-uniqueness}

Conservons les notations et les hypothèses du lemme \ref{cohomology-lemma-uniqueness}. Soient $U, V \in \mathcal{B}$. Soit $\mathcal{B}'$ l'ensemble des éléments de $\mathcal{B}$ contenus dans $U \cap V$. Alors
$$
(\{K_{U'}\}_{U' \in \mathcal{B}'},
\{\rho_{V'}^{U'}\}_{V' \subset U'\text{ avec }U', V' \in \mathcal{B}'})
$$
est un système sur l'espace annelé $U \cap V$
satisfaisant aux hypothèses du lemme \ref{cohomology-lemma-uniqueness}. De plus, $(K_U|_{U \cap V}, \rho^U_{U'})$ et
$(K_V|_{U \cap V}, \rho^V_{U'})$ sont deux solutions de ce système. Le lemme fournit donc un unique isomorphisme
$$
\rho_{U, V} : K_U|_{U \cap V} \longrightarrow K_V|_{U \cap V}
$$
tel que, pour tout $U' \subset U \cap V$, $U' \in \mathcal{B}$, le diagramme
$$
\xymatrix{
K_U|_{U'} \ar[rr]_{\rho_{U, V}|_{U'}} \ar[rd]_{\rho^U_{U'}} & &
K_V|_{U'} \ar[ld]^{\rho^V_{U'}} \\
& K_{U'}
}
$$
soit commutatif. Choisissons un troisième élément $W \in \mathcal{B}$. Nous obtenons des isomorphismes
$\rho_{U, W} : K_U|_{U \cap W} \to K_W|_{U \cap W}$ et
$\rho_{V, W} : K_U|_{V \cap W} \to K_W|_{V \cap W}$ satisfaisant des propriétés analogues à celles de $\rho_{U, V}$. Enfin, nous avons
$$
\rho_{U, W}|_{U \cap V \cap W} =
\rho_{V, W}|_{U \cap V \cap W} \circ \rho_{U, V}|_{U \cap V \cap W}
$$
Cela résulte de l'unicité dans le lemme, car les deux membres de l'égalité sont l'unique isomorphisme compatible avec les applications $\rho^U_{U''}$ et $\rho^W_{U''}$
pour $U'' \subset U \cap V \cap W$, $U'' \in \mathcal{B}$. Nous omettons quelques détails mineurs. La famille $(K_U, \rho_{U, V})$ est une donnée de descente dans la catégorie dérivée pour le recouvrement ouvert
$\mathcal{U} : X = \bigcup_{U \in \mathcal{B}} U$ de $X$. Dans ce langage, chercher une solution revient à chercher l'« effectivité de la donnée de descente ».

\end{remark}

\begin{lemma}

\label{cohomology-lemma-solution-in-finite-case}

Dans la situation \ref{cohomology-situation-locally-given}, supposons que

\begin{enumerate}

\item  $X = U_1 \cup \ldots \cup U_n$ avec $U_i \in \mathcal{B}$,
\item pour $U, V \in \mathcal{B}$, nous ayons
$U \cap V = \bigcup_{W \in \mathcal{B}, W \subset U \cap V} W$,
\item pour tout $U \in \mathcal{B}$, nous ayons $\Ext^i(K_U, K_U) = 0$
pour $i < 0$.

\end{enumerate}

Alors une solution existe et elle est unique à un unique isomorphisme près.

\end{lemma}

\begin{proof}
L'unicité a été établie dans le lemme \ref{cohomology-lemma-uniqueness}. Démontrons le lemme par récurrence sur $n$. Le cas $n = 1$ est immédiat.

\medskip\noindent

Le cas $n = 2$. Considérons l'isomorphisme
$\rho_{U_1, U_2} : K_{U_1}|_{U_1 \cap U_2} \to K_{U_2}|_{U_1 \cap U_2}$
construit dans la remarque \ref{cohomology-remark-uniqueness}. Par le lemme \ref{cohomology-lemma-glue}, nous obtenons un objet $K$ de $D(\mathcal{O}_X)$
et des isomorphismes $\rho_{U_1} : K|_{U_1} \to K_{U_1}$ et
$\rho_{U_2} : K|_{U_2} \to K_{U_2}$ compatibles avec $\rho_{U_1, U_2}$. Soit $U \in \mathcal{B}$. Nous allons construire un isomorphisme
$\rho_U : K|_U \to K_U$ et laisserons au lecteur le soin de vérifier que $(K, \rho_U)$ est une solution. Considérons l'ensemble $\mathcal{B}'$
des éléments de $\mathcal{B}$ contenus dans $U \cap U_1$ ou dans
$U \cap U_2$. Alors $(K_U, \rho^U_{U'})$ est une solution du système
$(\{K_{U'}\}_{U' \in \mathcal{B}'},
\{\rho_{V'}^{U'}\}_{V' \subset U'\text{ avec }U', V' \in \mathcal{B}'})$
sur l'espace annelé $U$. Nous affirmons que $(K|_U, \tau_{U'})$ est une autre solution, où
$\tau_{U'}$, pour $U' \in \mathcal{B}'$, est défini comme suit : si $U' \subset U_1$, nous prenons la composée
$$
K|_{U'} \xrightarrow{\rho_{U_1}|_{U'}}
K_{U_1}|_{U'} \xrightarrow{\rho^{U_1}_{U'}}
K_{U'}
$$
et, si $U' \subset U_2$, la composée
$$
K|_{U'} \xrightarrow{\rho_{U_2}|_{U'}}
K_{U_2}|_{U'} \xrightarrow{\rho^{U_2}_{U'}}
K_{U'}.
$$
Pour vérifier qu'il s'agit d'une solution, utilisons la propriété de l'application $\rho_{U_1, U_2}$
décrite dans la remarque \ref{cohomology-remark-uniqueness} et la compatibilité de
$\rho_{U_1}$ et $\rho_{U_2}$ avec $\rho_{U_1, U_2}$. Le lemme \ref{cohomology-lemma-uniqueness} fournit alors un unique isomorphisme $K|_{U'} \to K_{U'}$ compatible avec les applications $\tau_{U'}$ et
$\rho^U_{U'}$ pour $U' \in \mathcal{B}'$.

\medskip\noindent

Le cas $n > 2$. Considérons le sous-espace ouvert
$X' = U_1 \cup \ldots \cup U_{n - 1}$ et soit $\mathcal{B}'$ l'ensemble des éléments de $\mathcal{B}$ contenus dans $X'$. Nous obtenons un système
$(\{K_U\}_{U \in \mathcal{B}'}, \{\rho_V^U\}_{U, V \in \mathcal{B}'})$
sur l'espace annelé $X'$, auquel nous pouvons appliquer l'hypothèse de récurrence. Nous obtenons une solution $(K_{X'}, \rho^{X'}_U)$. Considérons alors la famille
$\mathcal{B}^* = \mathcal{B} \cup \{X'\}$ des ouvertures de $X$ et nous voyons que nous obtenons un système
$(\{K_U\}_{U \in \mathcal{B}^*},
\{\rho_V^U\}_{V \subset U\text{ avec }U, V \in \mathcal{B}^*})$. Ce nouveau système satisfait encore la condition (3), par le lemme \ref{cohomology-lemma-uniqueness} appliqué à la solution $K_{X'}$. Pour ce système, $X = X' \cup U_n$, ce qui nous ramène au cas $n = 2$ traité ci-dessus.

\end{proof}

\begin{lemma}

\label{cohomology-lemma-glueing-increasing-union}

Soit $X$ un espace annelé. Soit $E$ un ensemble bien ordonné et soit
$$
X = \bigcup\nolimits_{\alpha \in E} W_\alpha
$$
être un recouvrement ouvert avec $W_\alpha \subset W_{\alpha + 1}$
et $W_\alpha = \bigcup_{\beta < \alpha} W_\beta$ si $\alpha$ n'est pas un successeur. Soit $K_\alpha$ un objet de $D(\mathcal{O}_{W_\alpha})$
tel que $\Ext^i(K_\alpha, K_\alpha) = 0$ pour $i < 0$. Supposons donnés des isomorphismes
$\rho_\beta^\alpha :  K_\alpha|_{W_\beta} \to K_\beta$ en
$D(\mathcal{O}_{W_\beta})$ pour tous $\beta < \alpha$ avec
$\rho_\gamma^\alpha = \rho_\gamma^\beta \circ \rho^\alpha_\beta|_{W_\gamma}$
pour $\gamma < \beta < \alpha$. Alors il existe un objet
$K$ en $D(\mathcal{O}_X)$ et les isomorphismes
$K|_{W_\alpha} \to K_\alpha$ pour $\alpha \in E$
compatible avec les isomorphismes $\rho_\beta^\alpha$.

\end{lemma}

\begin{proof}

Dans cette preuve, $\alpha, \beta, \gamma, \ldots$ désignent des éléments de $E$. Choisissons un complexe K-injectif
$I_\alpha^\bullet$ sur $W_\alpha$ représentant $K_\alpha$. Pour $\beta < \alpha$, notons $j_{\beta, \alpha} : W_\beta \to W_\alpha$
le morphisme d'inclusion. Par récurrence transfinie, nous allons construire, pour tout
$\beta < \alpha$, un morphisme de complexes
$$
\tau_{\beta, \alpha} :
(j_{\beta, \alpha})_!I_\beta^\bullet
\longrightarrow
I_\alpha^\bullet
$$
représentant l'adjoint de l'inverse de l'isomorphisme
$\rho^\alpha_\beta : K_\alpha|_{W_\beta} \to K_\beta$. Nous effectuerons de plus cette construction de sorte que, pour
$\gamma < \beta < \alpha$, nous ayons
$$
\tau_{\gamma, \alpha} = \tau_{\beta, \alpha} \circ
(j_{\beta, \alpha})_!\tau_{\gamma, \beta}
$$
comme morphismes de complexes. Supposons en effet les $\tau_{\gamma, \beta}$ déjà construits
et se composant correctement pour tous $\gamma < \beta < \alpha$. Si $\alpha = \alpha' + 1$ est un successeur, choisissons un morphisme de complexes
$$
(j_{\alpha', \alpha})_!I_{\alpha'}^\bullet \to I_\alpha^\bullet
$$
adjoint de l'inverse de l'isomorphisme
$\rho^\alpha_{\alpha'} : K_\alpha|_{W_{\alpha'}} \to K_{\alpha'}$
(ce qui est possible puisque $I_\alpha^\bullet$ est K-injectif) et, pour tout $\beta < \alpha'$, posons
$$
\tau_{\beta, \alpha} = \tau_{\alpha', \alpha} \circ
(j_{\alpha', \alpha})_!\tau_{\beta, \alpha'}
$$
Si $\alpha$ n'est pas un successeur, considérons sur $W_\alpha$ le complexe
$$
C^\bullet = \colim_{\beta < \alpha} (j_{\beta, \alpha})_!I_\beta^\bullet
$$
(colimite terme à terme), où les morphismes de transition sont les $\tau_{\beta', \beta}$ pour
$\beta' < \beta < \alpha$. Nous affirmons que $C^\bullet$
représente $K_\alpha$. En effet, pour $\beta < \alpha$, la restriction de la coprojection $(j_{\beta, \alpha})_!I_\beta^\bullet \to C^\bullet$
donne un morphisme
$$
\sigma_\beta : I_\beta^\bullet \longrightarrow C^\bullet|_{W_\beta}
$$
qui est un quasi-isomorphisme : si $x \in W_\beta$, le passage aux germes donne
$$
(C^\bullet)_x =
\colim_{\beta' < \alpha}
\left((j_{\beta', \alpha})_!I_{\beta'}^\bullet\right)_x =
\colim_{\beta \leq \beta' < \alpha} (I_{\beta'}^\bullet)_x
\longleftarrow
(I_\beta^\bullet)_x
$$
qui est un quasi-isomorphisme. Nous avons utilisé ici le fait que le passage aux germes commute aux colimites, que les colimites filtrantes sont exactes et que les morphismes $(I_\beta^\bullet)_x \to (I_{\beta'}^\bullet)_x$
sont des quasi-isomorphismes pour $\beta \leq \beta' < \alpha$. Ainsi, $(C^\bullet, \sigma_\beta^{-1})$ est une solution du système $(\{K_\beta\}_{\beta < \alpha},
\{\rho^\beta_{\beta'}\}_{\beta' < \beta < \alpha})$. Comme $(K_\alpha, \rho^\alpha_\beta)$ est une autre solution, nous obtenons un unique isomorphisme $\sigma : K_\alpha \to C^\bullet$
dans $D(\mathcal{O}_{W_\alpha})$, compatible avec tous nos morphismes ; voir le lemme \ref{cohomology-lemma-solution-in-finite-case}
(c'est ici que nous utilisons l'annulation des groupes d'Ext négatifs). Choisissons un morphisme de complexes $\tau : C^\bullet \to I_\alpha^\bullet$
représentant $\sigma$. Posons alors
$$
\tau_{\beta, \alpha} = \tau|_{W_\beta} \circ \sigma_\beta
$$
pour obtenir les morphismes recherchés. Enfin, prenons pour $K$ l'objet de la catégorie dérivée représenté par le complexe
$$
K^\bullet = \colim_{\alpha \in E} (W_\alpha \to X)_!I_\alpha^\bullet
$$
où les morphismes de transition sont les morphismes $\tau_{\beta, \alpha}$ soigneusement construits pour $\beta < \alpha$. Le même raisonnement que ci-dessus montre que, pour tout $\alpha$,
la restriction de la coprojection détermine un isomorphisme
$$
K|_{W_\alpha} \longrightarrow K_\alpha
$$
compatible avec les morphismes donnés $\rho^\alpha_\beta$.

\end{proof}

\noindent
Une récurrence transfinie permet de démontrer le résultat dans le cas général.

\begin{theorem}
[Lemme de recollement de BBD]

\label{cohomology-theorem-glueing-bbd-general}

\begin{reference}

Cas particulier de \cite[théorème 3.2.4]{BBD}, sans hypothèse de bornitude.

\end{reference}

Dans la situation \ref{cohomology-situation-locally-given}, supposons que

\begin{enumerate}

\item  $X = \bigcup_{U \in \mathcal{B}} U$,
\item pour $U, V \in \mathcal{B}$, nous ayons
$U \cap V = \bigcup_{W \in \mathcal{B}, W \subset U \cap V} W$,
\item pour tout $U \in \mathcal{B}$, nous ayons $\Ext^i(K_U, K_U) = 0$
pour $i < 0$.

\end{enumerate}

Alors il existe un objet $K$ de $D(\mathcal{O}_X)$
et des isomorphismes $\rho_U : K|_U \to K_U$ dans $D(\mathcal{O}_U)$ pour
$U \in \mathcal{B}$, tels que $\rho^U_V \circ \rho_U|_V = \rho_V$
pour tous $V \subset U$ avec $U, V \in \mathcal{B}$. Le couple $(K, \rho_U)$ est unique à un unique isomorphisme près.

\end{theorem}

\begin{proof}
Un couple $(K, \rho_U)$ est appelé une solution dans le texte ci-dessus. L'unicité résulte du lemme \ref{cohomology-lemma-uniqueness}. Si $X$ admet un recouvrement fini par des éléments de $\mathcal{B}$ (par exemple si $X$ est quasi-compact), le théorème résulte du lemme \ref{cohomology-lemma-solution-in-finite-case}. Dans le cas général, nous raisonnons de la même manière, en utilisant une récurrence transfinie et le lemme \ref{cohomology-lemma-glueing-increasing-union}.

\medskip\noindent

Utilisons d'abord la récurrence transfinie pour choisir un ouvert $W_\alpha \subset X$
pour tout ordinal $\alpha$. Posons $W_0 = \emptyset$. Si $\alpha = \beta + 1$ est un successeur, ou bien $W_\beta = X$
et nous posons $W_\alpha = X$, ou bien $W_\beta \not = X$ et nous posons
$W_\alpha = W_\beta \cup U_\alpha$ où
$U_\alpha \in \mathcal{B}$ n'est pas contenu dans $W_\beta$. Si $\alpha$ est un ordinal limite, posons
$W_\alpha = \bigcup_{\beta < \alpha} W_\beta$. Alors, pour $\alpha$ assez grand, on a $W_\alpha = X$. Pour tout $\alpha$, l'ouvert $W_\alpha$ est une réunion d'éléments de $\mathcal{B}$. Par conséquent, si
$\mathcal{B}_\alpha = \{U \in \mathcal{B}, U \subset W_\alpha\}$, alors
$$
S_\alpha = (\{K_U\}_{U \in \mathcal{B}_\alpha},
\{\rho_V^U\}_{V \subset U\text{ avec }U, V \in \mathcal{B}_\alpha})
$$
est un système comme dans le lemme \ref{cohomology-lemma-uniqueness} sur l'espace annelé $W_\alpha$.

\medskip\noindent
Nous allons montrer par récurrence transfinie que, pour tout $\alpha$, le système $S_\alpha$ admet une solution. Cela démontrera le théorème, car pour $\alpha$ assez grand ce système est celui de l'énoncé.

\medskip\noindent
Supposons que $\alpha = \beta + 1$ soit un ordinal successeur. (Ce cas a déjà été traité dans la preuve du lemme précédent, mais nous répétons l'argument par souci de clarté.) Rappelons qu'alors $W_\alpha = W_\beta \cup U_\alpha$ pour un certain $U_\alpha \in \mathcal{B}$. Par hypothèse de récurrence, nous disposons d'une solution $(K_{W_\beta}, \{\rho^{W_\beta}_U\}_{U \in \mathcal{B}_\beta})$ du système $S_\beta$. Considérons la famille $\mathcal{B}_\alpha^* = \mathcal{B}_\alpha \cup \{W_\beta\}$ d'ouverts de $W_\alpha$ ; elle fournit un système $(\{K_U\}_{U \in \mathcal{B}_\alpha^*},
\{\rho_V^U\}_{V \subset U\text{ avec }U, V \in \mathcal{B}_\alpha^*})$. Ce nouveau système satisfait encore la condition (3), par le lemme \ref{cohomology-lemma-uniqueness} appliqué à la solution $K_{W_\beta}$. Comme $W_\alpha = W_\beta \cup U_\alpha$, nous sommes ramenés au cas traité dans le lemme \ref{cohomology-lemma-solution-in-finite-case}.

\medskip\noindent

Supposons que $\alpha$ soit un ordinal limite. Rappelons qu'alors
$W_\alpha = \bigcup_{\beta < \alpha} W_\beta$. Pour $\beta < \alpha$, soit
$(K_{W_\beta}, \{\rho^{W_\beta}_U\}_{U \in \mathcal{B}_\beta})$
la solution de $S_\beta$. Pour $\gamma < \beta < \alpha$, la restriction
$K_{W_\beta}|_{W_\gamma}$, munie des morphismes
$\rho^{W_\beta}_U$, $U \in \mathcal{B}_\gamma$
est une solution de $S_\gamma$. Par unicité, nous obtenons des isomorphismes uniques
$\rho_{W_\gamma}^{W_\beta} : K_{W_\beta}|_{W_\gamma} \to K_{W_\gamma}$
compatibles avec les morphismes $\rho^{W_\beta}_U$ et $\rho^{W_\gamma}_U$
pour $U \in \mathcal{B}_\gamma$. Ces morphismes se composent correctement, c'est-à-dire que $\rho_{W_\delta}^{W_\gamma} \circ \rho_{W_\gamma}^{W_\beta}|_{W_\delta}
= \rho^{W_\delta}_{W_\beta}$ pour $\delta < \gamma < \beta < \alpha$. Nous pouvons donc appliquer le lemme \ref{cohomology-lemma-glueing-increasing-union}
(notons que les Ext négatifs de
$K_{W_\beta}$ sont nuls par le lemme \ref{cohomology-lemma-uniqueness} appliqué à la solution $K_{W_\beta}$) afin d'obtenir $K_{W_\alpha}$ et des isomorphismes
$$
\rho_{W_\beta}^{W_\alpha} :
K_{W_\alpha}|_{W_\beta}
\longrightarrow
K_{W_\beta}
$$
compatibles avec les morphismes $\rho_{W_\gamma}^{W_\beta}$ pour
$\gamma < \beta < \alpha$.

\medskip\noindent

Pour montrer que $K_{W_\alpha}$ est une solution, il reste à construire des isomorphismes $\rho_U^{W_\alpha} : K_{W_\alpha}|_U \to K_U$ pour
$U \in \mathcal{B}_\alpha$, satisfaisant certaines compatibilités. Choisissons pour $\rho_U^{W_\alpha}$ l'unique morphisme tel que, pour $\beta < \alpha$ et tout $V \in \mathcal{B}_\beta$
avec $V \subset U$, le diagramme
$$
\xymatrix{
K_{W_\alpha}|_V \ar[r]_{\rho_U^{W_\alpha}|_V}
\ar[d]_{\rho_{W_\beta}^{W_\alpha}|_V}
& K_U|_V \ar[d]^{\rho_U^V} \\
K_{W_\beta} \ar[r]^{\rho_V^{W_\beta}}
& K_V
}
$$
soit commutatif. Ce choix est possible parce que
$$
(\{K_V\}_{V \subset U, V \in \mathcal{B}_\beta\text{ pour un certain }\beta < \alpha},
\{\rho_V^{V'}\}_{V \subset V'\text{ avec }V, V' \subset U
\text{ et }V, V' \in \mathcal{B}_\beta\text{ pour un certain }\beta < \alpha})
$$
est un système comme dans le lemme \ref{cohomology-lemma-uniqueness} sur l'espace annelé $U$
et que $(K_U, \rho^U_V)$ et
$(K_{W_\alpha}|_U,  \rho_V^{W_\beta}\circ \rho_{W_\beta}^{W_\alpha}|_V)$
en sont deux solutions. Cela donne l'existence et l'unicité. Nous omettons de vérifier que ces morphismes satisfont les compatibilités voulues (il ne s'agit que d'une vérification formelle).

\end{proof}



\section{Complexes strictement parfaits}

\label{cohomology-section-strictly-perfect}

\noindent
Les complexes strictement parfaits de modules serviront plus loin à définir les notions de complexe pseudo-cohérent et de complexe parfait. Les voici.

\begin{definition}

\label{cohomology-definition-strictly-perfect}
Soit $(X, \mathcal{O}_X)$ un espace annelé et soit $\mathcal{E}^\bullet$ un complexe de $\mathcal{O}_X$-modules. On dit que $\mathcal{E}^\bullet$ est {\it strictement parfait} si $\mathcal{E}^i$ est nul sauf pour un nombre fini d'indices $i$ et si $\mathcal{E}^i$ est un facteur direct d'un $\mathcal{O}_X$-module libre de type fini pour tout $i$.
\end{definition}

\noindent
Attention : puisque nous ne supposons pas que $X$ est un espace localement annelé, un facteur direct d'un $\mathcal{O}_X$-module libre de type fini n'est pas nécessairement localement libre de type fini.

\begin{lemma}

\label{cohomology-lemma-cone}

Le cône sur un morphisme de complexes strictement parfaits est strictement parfait.

\end{lemma}

\begin{proof}
Cela résulte immédiatement des définitions.
\end{proof}

\begin{lemma}

\label{cohomology-lemma-tensor}

Le complexe total associé au produit tensoriel de deux complexes strictement parfaits est strictement parfait.

\end{lemma}

\begin{proof}
Démonstration omise.
\end{proof}

\begin{lemma}

\label{cohomology-lemma-strictly-perfect-pullback}

Soit $f : (X, \mathcal{O}_X) \to (Y, \mathcal{O}_Y)$
un morphisme d'espaces annelés. Si $\mathcal{F}^\bullet$ est un complexe strictement parfait de $\mathcal{O}_Y$-modules, alors
$f^*\mathcal{F}^\bullet$ est un complexe strictement parfait de
$\mathcal{O}_X$-modules.

\end{lemma}

\begin{proof}

L'image réciproque d'un module libre de type fini est libre de type fini.
Le foncteur $f^*$ est additif, donc préserve les facteurs directs. Le lemme en résulte.

\end{proof}

\begin{lemma}

\label{cohomology-lemma-local-lift-map}
Soit $(X, \mathcal{O}_X)$ un espace annelé. Étant donné un diagramme de $\mathcal{O}_X$-modules dont les flèches pleines sont fixées $$
\xymatrix{
\mathcal{E} \ar@{..>}[dr] \ar[r] & \mathcal{F} \\
& \mathcal{G} \ar[u]_p
}
$$ où $\mathcal{E}$ est un facteur direct d'un $\mathcal{O}_X$-module libre de type fini et où $p$ est surjectif, il existe localement sur $X$ une flèche pointillée rendant le diagramme commutatif.
\end{lemma}

\begin{proof}

Nous pouvons supposer $\mathcal{E} = \mathcal{O}_X^{\oplus n}$ pour un certain $n$. Trouver la flèche pointillée revient alors à relever dans $\Gamma(X, \mathcal{F})$ les images des éléments de base. C'est possible localement par la caractérisation des morphismes surjectifs de faisceaux (Faisceaux, section \ref{sheaves-section-exactness-points}).

\end{proof}

\begin{lemma}

\label{cohomology-lemma-local-homotopy}

Soit $(X, \mathcal{O}_X)$ un espace annelé.

\begin{enumerate}

\item Soit $\alpha : \mathcal{E}^\bullet \to \mathcal{F}^\bullet$
un morphisme de complexes de $\mathcal{O}_X$-modules, avec $\mathcal{E}^\bullet$ strictement parfait et $\mathcal{F}^\bullet$
acyclique. Alors $\alpha$ est localement homotope à zéro sur $X$.
\item Soit $\alpha : \mathcal{E}^\bullet \to \mathcal{F}^\bullet$
être un morphisme de complexes de $\mathcal{O}_X$-modules avec $\mathcal{E}^\bullet$ strictement parfait, $\mathcal{E}^i = 0$
pour $i < a$ et $H^i(\mathcal{F}^\bullet) = 0$ pour $i \geq a$. Alors $\alpha$ est localement homotope à zéro sur $X$.

\end{enumerate}

\end{lemma}

\begin{proof}

La première assertion résulte de la seconde ; il suffit donc de démontrer (2). Procédons par récurrence sur la longueur du complexe
$\mathcal{E}^\bullet$. Si $\mathcal{E}^\bullet \cong \mathcal{E}[-n]$
pour un facteur direct $\mathcal{E}$ d'un
$\mathcal{O}_X$-module libre de type fini et un entier $n \geq a$, le résultat découle du lemme \ref{cohomology-lemma-local-lift-map} et du fait que
$\mathcal{F}^{n - 1} \to \Ker(\mathcal{F}^n \to \mathcal{F}^{n + 1})$
est surjective par l'annulation supposée de $H^n(\mathcal{F}^\bullet)$. Si $\mathcal{E}^i$ est nul sauf pour $i \in [a, b]$, nous disposons de la suite exacte scindée de complexes
$$
0 \to \mathcal{E}^b[-b] \to \mathcal{E}^\bullet \to
\sigma_{\leq b - 1}\mathcal{E}^\bullet \to 0
$$
qui détermine un triangle distingué dans
$K(\mathcal{O}_X)$, et donc une suite exacte
$$
\Hom_{K(\mathcal{O}_X)}(
\sigma_{\leq b - 1}\mathcal{E}^\bullet, \mathcal{F}^\bullet)
\to
\Hom_{K(\mathcal{O}_X)}(\mathcal{E}^\bullet, \mathcal{F}^\bullet)
\to
\Hom_{K(\mathcal{O}_X)}(\mathcal{E}^b[-b], \mathcal{F}^\bullet)
$$
par les axiomes des catégories triangulées. La composée
$\mathcal{E}^b[-b] \to \mathcal{F}^\bullet$ est localement homotope à zéro ; nous pouvons donc supposer que notre morphisme provient d'un élément du membre de gauche de la suite exacte affichée ci-dessus. Cet élément est localement nul par l'hypothèse de récurrence.

\end{proof}

\begin{lemma}

\label{cohomology-lemma-lift-through-quasi-isomorphism}
Soit $(X, \mathcal{O}_X)$ un espace annelé. Étant donné un diagramme de complexes de $\mathcal{O}_X$-modules dont les flèches pleines sont fixées $$
\xymatrix{
\mathcal{E}^\bullet \ar@{..>}[dr] \ar[r]_\alpha & \mathcal{F}^\bullet \\
& \mathcal{G}^\bullet \ar[u]_f
}
$$ où $\mathcal{E}^\bullet$ est strictement parfait, $\mathcal{E}^j = 0$ pour $j < a$, et où $H^j(f)$ est un isomorphisme pour $j > a$ et une surjection pour $j = a$, il existe localement sur $X$ une flèche pointillée rendant le diagramme commutatif à homotopie près.
\end{lemma}

\begin{proof}

Nos hypothèses sur $f$ impliquent que les faisceaux de cohomologie du cône $C(f)^\bullet$ sont nuls en degrés $\geq a$. Le lemme \ref{cohomology-lemma-local-homotopy} garantit donc l'existence d'un recouvrement ouvert $X = \bigcup U_i$ tel que la composée
$\mathcal{E}^\bullet \to \mathcal{F}^\bullet \to C(f)^\bullet$
soit homotope à zéro sur $U_i$. Comme
$$
\mathcal{G}^\bullet \to \mathcal{F}^\bullet \to C(f)^\bullet \to
\mathcal{G}^\bullet[1]
$$
se restreint en un triangle distingué dans $K(\mathcal{O}_{U_i})$,
nous pouvons relever $\alpha|_{U_i}$, à homotopie près, en un morphisme
$\alpha_i : \mathcal{E}^\bullet|_{U_i} \to \mathcal{G}^\bullet|_{U_i}$
comme désiré.

\end{proof}

\begin{lemma}

\label{cohomology-lemma-local-actual}

Soit $(X, \mathcal{O}_X)$ un espace annelé. Soient $\mathcal{E}^\bullet$, $\mathcal{F}^\bullet$ des complexes de $\mathcal{O}_X$-modules, avec $\mathcal{E}^\bullet$ strictement parfait.

\begin{enumerate}

\item Pour tout élément
$\alpha \in \Hom_{D(\mathcal{O}_X)}(\mathcal{E}^\bullet, \mathcal{F}^\bullet)$
il existe un recouvrement ouvert $X = \bigcup U_i$ tel que
$\alpha|_{U_i}$ soit représenté par un morphisme de complexes
$\alpha_i : \mathcal{E}^\bullet|_{U_i} \to \mathcal{F}^\bullet|_{U_i}$.
\item Étant donné un morphisme de complexes
$\alpha : \mathcal{E}^\bullet \to \mathcal{F}^\bullet$
dont l'image dans le groupe
$\Hom_{D(\mathcal{O}_X)}(\mathcal{E}^\bullet, \mathcal{F}^\bullet)$
est nulle, il existe un recouvrement ouvert $X = \bigcup U_i$ tel que
$\alpha|_{U_i}$ soit homotope à zéro.

\end{enumerate}

\end{lemma}

\begin{proof}
Démonstration de (1). Par construction de la catégorie dérivée, il existe un quasi-isomorphisme $f : \mathcal{F}^\bullet \to \mathcal{G}^\bullet$ et un morphisme de complexes $\beta : \mathcal{E}^\bullet \to \mathcal{G}^\bullet$ tels que $\alpha = f^{-1}\beta$. Le résultat découle donc du lemme \ref{cohomology-lemma-lift-through-quasi-isomorphism}. Nous omettons la démonstration de (2).
\end{proof}

\begin{lemma}

\label{cohomology-lemma-Rhom-strictly-perfect}
Soit $(X, \mathcal{O}_X)$ un espace annelé. Soient $\mathcal{E}^\bullet$, $\mathcal{F}^\bullet$ des complexes de $\mathcal{O}_X$-modules, avec $\mathcal{E}^\bullet$ strictement parfait. Alors le Hom interne $R\SheafHom(\mathcal{E}^\bullet, \mathcal{F}^\bullet)$ est représenté par le complexe $\mathcal{H}^\bullet$ de termes $$
\mathcal{H}^n =
\bigoplus\nolimits_{n = p + q}
\SheafHom_{\mathcal{O}_X}(\mathcal{E}^{-q}, \mathcal{F}^p)
$$ et de différentielle décrite dans la section \ref{cohomology-section-hom-complexes}.
\end{lemma}

\begin{proof}

Choisissons un quasi-isomorphisme $\mathcal{F}^\bullet \to \mathcal{I}^\bullet$
vers un complexe K-injectif. Soit $(\mathcal{H}')^\bullet$ le complexe de termes
$$
(\mathcal{H}')^n =
\prod\nolimits_{n = p + q}
\SheafHom_{\mathcal{O}_X}(\mathcal{E}^{-q}, \mathcal{I}^p)
$$
qui représente $R\SheafHom(\mathcal{E}^\bullet, \mathcal{F}^\bullet)$
par la construction de la section \ref{cohomology-section-internal-hom}. Il suffit de montrer que le morphisme
$$
\mathcal{H}^\bullet \longrightarrow (\mathcal{H}')^\bullet
$$
est un quasi-isomorphisme. Pour tout ouvert $U \subset X$, on constate que
$$
H^0(\mathcal{H}^\bullet(U)) =
\Hom_{K(\mathcal{O}_U)}(\mathcal{E}^\bullet|_U, \mathcal{I}^\bullet|_U)
\to
H^0((\mathcal{H}')^\bullet(U)) =
\Hom_{D(\mathcal{O}_U)}(\mathcal{E}^\bullet|_U, \mathcal{I}^\bullet|_U)
$$
Par le lemme \ref{cohomology-lemma-local-actual}, le faisceau associé au préfaisceau
$U \mapsto H^0(\mathcal{H}^\bullet(U))$
est égal au faisceau associé au préfaisceau
$U \mapsto H^0((\mathcal{H}')^\bullet(U))$. Le même argument s'applique aux autres faisceaux de cohomologie. Ainsi, $\mathcal{H}^\bullet$
est quasi-isomorphe à $(\mathcal{H}')^\bullet$, ce qui démontre le lemme.

\end{proof}

\begin{lemma}

\label{cohomology-lemma-Rhom-strictly-perfect-K-flat}
Dans la situation du lemme \ref{cohomology-lemma-Rhom-strictly-perfect}, si $\mathcal{F}^\bullet$ est K-plat, alors $\mathcal{H}^\bullet$ est K-plat.
\end{lemma}

\begin{proof}

Notons que $\mathcal{H}^\bullet$ est simplement le complexe de Hom
$\SheafHom^\bullet(\mathcal{E}^\bullet, \mathcal{F}^\bullet)$
car la bornitude du complexe strictement parfait
$\mathcal{E}^\bullet$ assure que les produits intervenant dans la définition du complexe de Hom sont des sommes directes. Soit $\mathcal{K}^\bullet$ un complexe acyclique de
$\mathcal{O}_X$-modules. Considérons le morphisme
$$
\gamma :
\text{Tot}(\mathcal{K}^\bullet \otimes
\SheafHom^\bullet(\mathcal{E}^\bullet, \mathcal{F}^\bullet))
\longrightarrow
\SheafHom^\bullet(\mathcal{E}^\bullet,
\text{Tot}(\mathcal{K}^\bullet \otimes \mathcal{F}^\bullet))
$$
du lemme \ref{cohomology-lemma-diagonal-better}. Comme $\mathcal{F}^\bullet$ est K-plat, le complexe $\text{Tot}(\mathcal{K}^\bullet \otimes \mathcal{F}^\bullet)$
est acyclique ; par le lemme \ref{cohomology-lemma-local-actual}
(ou, si l'on préfère, le lemme \ref{cohomology-lemma-Rhom-strictly-perfect}), le but de $\gamma$ est donc lui aussi acyclique. Pour démontrer le lemme, il suffit alors de montrer que $\gamma$ est un isomorphisme de complexes. Procédons par récurrence sur la longueur du complexe $\mathcal{E}^\bullet$. Si cette longueur est $\leq 1$, alors $\mathcal{E}^\bullet$
est un facteur direct de $\mathcal{O}_X^{\oplus n}[k]$ pour certains
$n \geq 0$ et $k \in \mathbf{Z}$, et le résultat est immédiat. Si la longueur est $> 1$, on la réduit en considérant la suite exacte courte de complexes, scindée terme à terme,
$$
0 \to \sigma_{\geq a + 1} \mathcal{E}^\bullet \to
\mathcal{E}^\bullet \to
\sigma_{\leq a} \mathcal{E}^\bullet \to 0
$$
pour un $a \in \mathbf{Z}$ convenable ; voir Homologie, section \ref{homology-section-truncations}. Alors $\gamma$ s'insère dans un morphisme de suites exactes courtes de complexes, scindées terme à terme. Par récurrence, $\gamma$ est un isomorphisme pour
$\sigma_{\geq a + 1} \mathcal{E}^\bullet$
et $\sigma_{\leq a} \mathcal{E}^\bullet$ ; le résultat en découle pour
$\mathcal{E}^\bullet$. Nous omettons quelques détails.

\end{proof}

\begin{lemma}

\label{cohomology-lemma-Rhom-complex-of-direct-summands-finite-free}
Soit $(X, \mathcal{O}_X)$ un espace annelé. Soient $\mathcal{E}^\bullet$, $\mathcal{F}^\bullet$ des complexes de $\mathcal{O}_X$-modules tels que
\begin{enumerate}
\item $\mathcal{F}^n = 0$ pour $n \ll 0$,
\item $\mathcal{E}^n = 0$ pour $n \gg 0$, et
\item $\mathcal{E}^n$ est isomorphe à un facteur direct d'un $\mathcal{O}_X$-module libre de type fini.
\end{enumerate}
Alors le Hom interne $R\SheafHom(\mathcal{E}^\bullet, \mathcal{F}^\bullet)$ est représenté par le complexe $\mathcal{H}^\bullet$ de termes $$
\mathcal{H}^n =
\bigoplus\nolimits_{n = p + q}
\SheafHom_{\mathcal{O}_X}(\mathcal{E}^{-q}, \mathcal{F}^p)
$$ et de différentielle décrite dans la section \ref{cohomology-section-internal-hom}.
\end{lemma}

\begin{proof}

Choisissons un quasi-isomorphisme $\mathcal{F}^\bullet \to \mathcal{I}^\bullet$
où $\mathcal{I}^\bullet$ est un complexe borné inférieurement d'objets injectifs. Notons que $\mathcal{I}^\bullet$ est K-injectif (Catégories dérivées, lemme
\ref{derived-lemma-bounded-below-injectives-K-injective}). La construction de la section \ref{cohomology-section-internal-hom}
montre que
$R\SheafHom(\mathcal{E}^\bullet, \mathcal{F}^\bullet)$ est représenté par le complexe $(\mathcal{H}')^\bullet$ avec des termes
$$
(\mathcal{H}')^n =
\prod\nolimits_{n = p + q}
\SheafHom_{\mathcal{O}_X}(\mathcal{E}^{-q}, \mathcal{I}^p) =
\bigoplus\nolimits_{n = p + q}
\SheafHom_{\mathcal{O}_X}(\mathcal{E}^{-q}, \mathcal{I}^p)
$$
(l'égalité venant de ce qu'il n'y a qu'un nombre fini de termes non nuls). Notons que $\mathcal{H}^\bullet$ est le complexe total associé au bicomplexe de termes
$\SheafHom_{\mathcal{O}_X}(\mathcal{E}^{-q}, \mathcal{F}^p)$
et de même pour $(\mathcal{H}')^\bullet$. Le morphisme naturel $(\mathcal{H}')^\bullet \to \mathcal{H}^\bullet$ provient d'un morphisme de bicomplexes.
Pour montrer qu'il s'agit d'un quasi-isomorphisme, nous pouvons donc utiliser la suite spectrale d'un bicomplexe (Homologie, lemme \ref{homology-lemma-first-quadrant-ss})
$$
{}'E_1^{p, q} =
H^p(\SheafHom_{\mathcal{O}_X}(\mathcal{E}^{-q}, \mathcal{F}^\bullet))
$$
convergeant vers $H^{p + q}(\mathcal{H}^\bullet)$, et de même pour
$(\mathcal{H}')^\bullet$. Pour achever la démonstration, il suffit de montrer que $\mathcal{F}^\bullet \to \mathcal{I}^\bullet$
induit un isomorphisme
$$
H^p(\SheafHom_{\mathcal{O}_X}(\mathcal{E}, \mathcal{F}^\bullet))
\longrightarrow
H^p(\SheafHom_{\mathcal{O}_X}(\mathcal{E}, \mathcal{I}^\bullet))
$$
sur les faisceaux de cohomologie dès que $\mathcal{E}$ est un facteur direct d'un $\mathcal{O}_X$-module libre de type fini. Comme c'est clair lorsque $\mathcal{E}$
est libre de type fini, le résultat en découle.

\end{proof}





\section{Modules pseudo-cohérents}

\label{cohomology-section-pseudo-coherent}

\noindent
Dans cette section, nous discutons des complexes pseudo-cohérents.

\begin{definition}

\label{cohomology-definition-pseudo-coherent}

Soit $(X, \mathcal{O}_X)$ un espace annelé. Soit $\mathcal{E}^\bullet$
un complexe de $\mathcal{O}_X$-modules et soit $m \in \mathbf{Z}$.

\begin{enumerate}

\item On dit que $\mathcal{E}^\bullet$ est {\it $m$-pseudo-cohérent}
s'il existe un recouvrement ouvert $X = \bigcup U_i$ et, pour tout $i$,
un morphisme de complexes
$\alpha_i : \mathcal{E}_i^\bullet \to \mathcal{E}^\bullet|_{U_i}$
où $\mathcal{E}_i^\bullet$ est strictement parfait sur $U_i$ et
$H^j(\alpha_i)$ est un isomorphisme pour $j > m$ et $H^m(\alpha_i)$
est surjectif.
\item On dit que $\mathcal{E}^\bullet$ est {\it pseudo-cohérent}
s'il est $m$-pseudo-cohérent pour tout $m$.
\item On dit qu'un objet $E$ de $D(\mathcal{O}_X)$ est
{\it $m$-pseudo-cohérent} (resp.\ {\it pseudo-cohérent}) s'il peut être représenté par un complexe $m$-pseudo-cohérent (resp.\ pseudo-cohérent) de $\mathcal{O}_X$-modules.

\end{enumerate}

\end{definition}

\noindent
Si $X$ est quasi-compact, tout objet $m$-pseudo-cohérent de $D(\mathcal{O}_X)$ appartient à $D^-(\mathcal{O}_X)$. Ce n'est plus nécessairement vrai lorsque $X$ n'est pas quasi-compact.

\begin{lemma}

\label{cohomology-lemma-pseudo-coherent-independent-representative}

Soit $(X, \mathcal{O}_X)$ un espace annelé et soit $E$ un objet de $D(\mathcal{O}_X)$.

\begin{enumerate}

\item S'il existe un recouvrement ouvert $X = \bigcup U_i$, des complexes strictement parfaits $\mathcal{E}_i^\bullet$ sur $U_i$ et des morphismes $\alpha_i : \mathcal{E}_i^\bullet \to E|_{U_i}$ dans
$D(\mathcal{O}_{U_i})$ avec $H^j(\alpha_i)$ un isomorphisme pour $j > m$
et $H^m(\alpha_i)$ surjectif, alors $E$ est $m$-pseudo-cohérent.
\item Si $E$ est $m$-pseudo-cohérent, alors tout complexe représentant
$E$ est $m$-pseudo-cohérent.

\end{enumerate}

\end{lemma}

\begin{proof}

Soit $\mathcal{F}^\bullet$ un complexe quelconque représentant $E$,
et soient $X = \bigcup U_i$ et
$\alpha_i : \mathcal{E}_i^\bullet \to E|_{U_i}$ comme dans (1). Montrons que $\mathcal{F}^\bullet$ est $m$-pseudo-cohérent comme complexe ; cela démontrera simultanément (1) et (2). Par le lemme \ref{cohomology-lemma-local-actual},
après avoir raffiné le recouvrement ouvert $X = \bigcup U_i$, nous pouvons
représenter les morphismes $\alpha_i$ par des morphismes de complexes
$\alpha_i : \mathcal{E}_i^\bullet \to \mathcal{F}^\bullet|_{U_i}$. Par hypothèse
les $H^j(\alpha_i)$ sont des isomorphismes pour $j > m$, et $H^m(\alpha_i)$
est surjectif ; ainsi $\mathcal{F}^\bullet$ est
$m$-pseudo-cohérent.

\end{proof}

\begin{lemma}

\label{cohomology-lemma-pseudo-coherent-pullback}
Soit $f : (X, \mathcal{O}_X) \to (Y, \mathcal{O}_Y)$ un morphisme d'espaces annelés. Soit $E$ un objet de $D(\mathcal{O}_Y)$. Si $E$ est $m$-pseudo-cohérent, alors $Lf^*E$ est $m$-pseudo-cohérent.
\end{lemma}

\begin{proof}

Représentons $E$ par un complexe $\mathcal{E}^\bullet$ de $\mathcal{O}_Y$-modules et choisissons un recouvrement ouvert $Y = \bigcup V_i$ ainsi que des morphismes
$\alpha_i : \mathcal{E}_i^\bullet \to \mathcal{E}^\bullet|_{V_i}$
comme dans la définition \ref{cohomology-definition-pseudo-coherent}. Posons $U_i = f^{-1}(V_i)$. Par le lemme \ref{cohomology-lemma-pseudo-coherent-independent-representative},
il suffit de montrer que $Lf^*\mathcal{E}^\bullet|_{U_i}$ est
$m$-pseudo-cohérent. Choisissons un triangle distingué
$$
\mathcal{E}_i^\bullet \to
\mathcal{E}^\bullet|_{V_i} \to
C \to
\mathcal{E}_i^\bullet[1]
$$
L'hypothèse sur $\alpha_i$ signifie exactement que les faisceaux de cohomologie
$H^j(C)$ sont nuls pour tout $j \geq m$. Notons $f_i : U_i \to V_i$
la restriction de $f$. Notons que $Lf^*\mathcal{E}^\bullet|_{U_i} = 
Lf_i^*(\mathcal{E}|_{V_i})$. En appliquant $Lf_i^*$, nous obtenons le triangle distingué
$$
Lf_i^*\mathcal{E}_i^\bullet \to
Lf_i^*\mathcal{E}|_{V_i} \to
Lf_i^*C \to
Lf_i^*\mathcal{E}_i^\bullet[1]
$$
La construction de $Lf_i^*$ comme foncteur dérivé à gauche montre que
$H^j(Lf_i^*C) = 0$ pour $j \geq m$ (par Catégories dérivées, lemme
\ref{derived-lemma-negative-vanishing}). Par conséquent, $H^j(Lf_i^*\alpha_i)$ est un isomorphisme pour $j > m$ et $H^m(Lf^*\alpha_i)$ est surjectif.
D'autre part, $Lf_i^*\mathcal{E}_i^\bullet = f_i^*\mathcal{E}_i^\bullet$, qui est strictement parfait par le lemme \ref{cohomology-lemma-strictly-perfect-pullback}. Cela conclut la démonstration.

\end{proof}

\begin{lemma}

\label{cohomology-lemma-cone-pseudo-coherent}

Soit $(X, \mathcal{O}_X)$ un espace annelé et soit $m \in \mathbf{Z}$. Soit $(K, L, M, f, g, h)$ un triangle distingué dans $D(\mathcal{O}_X)$.

\begin{enumerate}

\item Si $K$ est $(m + 1)$-pseudo-cohérent et $L$ est $m$-pseudo-cohérent, alors $M$ est $m$-pseudo-cohérent.
\item Si $K$ et $M$ sont $m$-pseudo-cohérents, alors $L$ est $m$-pseudo-cohérent.
\item Si $L$ est $(m + 1)$-pseudo-cohérent et $M$
est $m$-pseudo-cohérent, alors $K$ est $(m + 1)$-pseudo-cohérent.

\end{enumerate}

\end{lemma}

\begin{proof}

Démonstration de (1). Choisissons un recouvrement ouvert $X = \bigcup U_i$ et des morphismes $\alpha_i : \mathcal{K}_i^\bullet \to K|_{U_i}$ dans $D(\mathcal{O}_{U_i})$
avec $\mathcal{K}_i^\bullet$ strictement parfait et $H^j(\alpha_i)$
isomorphes pour $j > m + 1$ et surjectifs pour $j = m + 1$. Nous pouvons remplacer $\mathcal{K}_i^\bullet$ par
$\sigma_{\geq m + 1}\mathcal{K}_i^\bullet$
et donc nous pouvons supposer que $\mathcal{K}_i^j = 0$
pour $j < m + 1$. Après avoir raffiné le recouvrement ouvert, choisissons des morphismes $\beta_i : \mathcal{L}_i^\bullet \to L|_{U_i}$ dans $D(\mathcal{O}_{U_i})$
avec $\mathcal{L}_i^\bullet$ strictement parfait, tels que
$H^j(\beta)$ soit un isomorphisme pour $j > m$ et surjectif pour $j = m$. Par le lemme \ref{cohomology-lemma-lift-through-quasi-isomorphism},
nous pouvons, après un nouveau raffinement, trouver des morphismes de complexes
$\gamma_i : \mathcal{K}^\bullet \to \mathcal{L}^\bullet$
tels que les diagrammes
$$
\xymatrix{
K|_{U_i} \ar[r] & L|_{U_i} \\
\mathcal{K}_i^\bullet \ar[u]^{\alpha_i} \ar[r]^{\gamma_i} &
\mathcal{L}_i^\bullet \ar[u]_{\beta_i}
}
$$
sont commutatifs dans $D(\mathcal{O}_{U_i})$ (il faut pour cela représenter les morphismes $\alpha_i$, $\beta_i$ et $K|_{U_i} \to L|_{U_i}$
par de véritables morphismes de complexes ; nous omettons quelques détails). Le cône $C(\gamma_i)^\bullet$ est strictement parfait (lemme \ref{cohomology-lemma-cone}). La commutativité du diagramme implique l'existence d'un morphisme de triangles distingués
$$
(\mathcal{K}_i^\bullet, \mathcal{L}_i^\bullet, C(\gamma_i)^\bullet)
\longrightarrow
(K|_{U_i}, L|_{U_i}, M|_{U_i}).
$$
Il résulte du morphisme induit entre les suites exactes longues de cohomologie et de Homologie, lemmes \ref{homology-lemma-four-lemma} et
\ref{homology-lemma-five-lemma}
que $C(\gamma_i)^\bullet \to M|_{U_i}$ induit un isomorphisme en cohomologie dans les degrés $> m$ et une surjection en degré $m$. Ainsi, $M$ est $m$-pseudo-cohérent par le lemme \ref{cohomology-lemma-pseudo-coherent-independent-representative}.

\medskip\noindent

Les assertions (2) et (3) suivent de (1) par rotation du triangle distingué.

\end{proof}

\begin{lemma}

\label{cohomology-lemma-tensor-pseudo-coherent}

Soit $(X, \mathcal{O}_X)$ un espace annelé et soient $K, L$ des objets de $D(\mathcal{O}_X)$.

\begin{enumerate}

\item Si $K$ est $n$-pseudo-cohérent et $H^i(K) = 0$ pour $i > a$
et si $L$ est $m$-pseudo-cohérent et $H^j(L) = 0$ pour $j > b$, alors
$K \otimes_{\mathcal{O}_X}^\mathbf{L} L$ est $t$-pseudo-cohérent avec $t = \max(m + a, n + b)$.
\item Si $K$ et $L$ sont pseudo-cohérents, alors
$K \otimes_{\mathcal{O}_X}^\mathbf{L} L$ est pseudo-cohérent.

\end{enumerate}

\end{lemma}

\begin{proof}
Démonstration de (1). En remplaçant $X$ par les membres d'un recouvrement ouvert, nous pouvons supposer qu'il existe des complexes strictement parfaits $\mathcal{K}^\bullet$ et $\mathcal{L}^\bullet$ et des morphismes $\alpha : \mathcal{K}^\bullet \to K$ et $\beta : \mathcal{L}^\bullet \to L$, tels que $H^i(\alpha)$ soit un isomorphisme pour $i > n$ et surjectif pour $i = n$, et que $H^i(\beta)$ soit un isomorphisme pour $i > m$ et surjectif pour $i = m$. Alors le morphisme $$
\alpha \otimes^\mathbf{L} \beta :
\text{Tot}(\mathcal{K}^\bullet \otimes_{\mathcal{O}_X} \mathcal{L}^\bullet)
\to K \otimes_{\mathcal{O}_X}^\mathbf{L} L
$$ induit des isomorphismes sur les faisceaux de cohomologie en degré $i$ pour $i > t$ et une surjection pour $i = t$. Cela découle de la suite spectrale des Tor (détails omis).

\medskip\noindent

Démonstration de (2). En remplaçant d'abord $X$ par les membres d'un recouvrement ouvert, nous nous ramenons au cas où $K$ et $L$ sont bornés supérieurement. L'énoncé découle alors immédiatement de (1).

\end{proof}

\begin{lemma}

\label{cohomology-lemma-summands-pseudo-coherent}
Soit $(X, \mathcal{O}_X)$ un espace annelé et soit $m \in \mathbf{Z}$. Si $K \oplus L$ est $m$-pseudo-cohérent (resp.\ pseudo-cohérent) dans $D(\mathcal{O}_X)$, alors $K$ et $L$ le sont aussi.
\end{lemma}

\begin{proof}

Supposons que $K \oplus L$ soit $m$-pseudo-cohérent. Après avoir remplacé $X$ par les membres d'un recouvrement ouvert, nous pouvons supposer que $K \oplus L \in D^-(\mathcal{O}_X)$, donc que
$L \in D^-(\mathcal{O}_X)$. Notons qu'il existe un triangle distingué
$$
(K \oplus L, K \oplus L, L \oplus L[1]) =
(K, K, 0) \oplus (L, L, L \oplus L[1])
$$
Voir Catégories dérivées, lemme \ref{derived-lemma-direct-sum-triangles}. Par le lemme \ref{cohomology-lemma-cone-pseudo-coherent},
nous voyons que $L \oplus L[1]$ est $m$-pseudo-cohérent. Ainsi, $L[1] \oplus L[2]$ est lui aussi $m$-pseudo-cohérent. Par récurrence, $L[n] \oplus L[n + 1]$ est $m$-pseudo-cohérent. Comme $L$ est borné supérieurement, $L[n]$ est $m$-pseudo-cohérent pour $n$ assez grand. En remontant alors les triangles distingués
$$
(L[n], L[n] \oplus L[n - 1], L[n - 1])
$$
nous concluons que $L[n - 1], L[n - 2], \ldots, L$ sont $m$-pseudo-cohérents, comme souhaité.

\end{proof}

\begin{lemma}

\label{cohomology-lemma-complex-pseudo-coherent-modules}
Soit $(X, \mathcal{O}_X)$ un espace annelé et soit $m \in \mathbf{Z}$. Soit $\mathcal{F}^\bullet$ un complexe de $\mathcal{O}_X$-modules (localement) borné supérieurement, tel que $\mathcal{F}^i$ soit $(m - i)$-pseudo-cohérent pour tout $i$. Alors $\mathcal{F}^\bullet$ est $m$-pseudo-cohérent.
\end{lemma}

\begin{proof}
Démonstration omise. Indication : utiliser le lemme \ref{cohomology-lemma-cone-pseudo-coherent} et les troncatures comme dans la démonstration de Compléments sur l'algèbre, lemme \ref{more-algebra-lemma-complex-pseudo-coherent-modules}.
\end{proof}

\begin{lemma}

\label{cohomology-lemma-cohomology-pseudo-coherent}
Soit $(X, \mathcal{O}_X)$ un espace annelé, soit $m \in \mathbf{Z}$ et soit $E$ un objet de $D(\mathcal{O}_X)$. Si $E$ est (localement) borné supérieurement et si $H^i(E)$ est $(m - i)$-pseudo-cohérent pour tout $i$, alors $E$ est $m$-pseudo-cohérent.
\end{lemma}

\begin{proof}
Démonstration omise. Indication : utiliser le lemme \ref{cohomology-lemma-cone-pseudo-coherent} et les troncatures comme dans la démonstration de Compléments sur l'algèbre, lemme \ref{more-algebra-lemma-cohomology-pseudo-coherent}.
\end{proof}

\begin{lemma}

\label{cohomology-lemma-finite-cohomology}

Soit $(X, \mathcal{O}_X)$ un espace annelé, soit $K$ un objet de $D(\mathcal{O}_X)$ et soit $m \in \mathbf{Z}$.

\begin{enumerate}

\item Si $K$ est $m$-pseudo-cohérent et $H^i(K) = 0$
pour $i > m$, alors $H^m(K)$ est un $\mathcal{O}_X$-module de type fini.
\item Si $K$ est $m$-pseudo-cohérent et $H^i(K) = 0$
pour $i > m + 1$, alors $H^{m + 1}(K)$ est un
$\mathcal{O}_X$-module de présentation finie.

\end{enumerate}

\end{lemma}

\begin{proof}

(1). Nous pouvons travailler localement sur $X$. Supposons donc qu'il existe un complexe strictement parfait $\mathcal{E}^\bullet$ et un morphisme
$\alpha : \mathcal{E}^\bullet \to K$ induisant un isomorphisme en cohomologie dans les degrés $> m$ et une surjection en degré $m$. Il suffit de démontrer le résultat pour $\mathcal{E}^\bullet$. Soit $n$ le plus grand entier tel que $\mathcal{E}^n \not = 0$. Si $n = m$, alors $H^m(\mathcal{E}^\bullet)$ est un quotient de
$\mathcal{E}^n$ et le résultat est clair. Si $n > m$, alors $\mathcal{E}^{n - 1} \to \mathcal{E}^n$ est surjectif puisque
$H^n(E^\bullet) = 0$. Par le lemme \ref{cohomology-lemma-local-lift-map},
nous pouvons localement trouver une section de cette surjection et écrire
$\mathcal{E}^{n - 1} = \mathcal{E}' \oplus \mathcal{E}^n$. Il suffit donc de prouver le résultat pour le complexe
$(\mathcal{E}')^\bullet$ qui est le même que $\mathcal{E}^\bullet$
à ceci près qu'il contient $\mathcal{E}'$ en degré $n - 1$ et $0$ en degré $n$. On conclut par récurrence sur $n$.

\medskip\noindent
Démonstration de (2). Nous pouvons travailler localement sur $X$. Supposons donc qu'il existe un complexe strictement parfait $\mathcal{E}^\bullet$ et un morphisme $\alpha : \mathcal{E}^\bullet \to K$ induisant un isomorphisme en cohomologie dans les degrés $> m$ et une surjection en degré $m$. Comme dans la démonstration de (1), nous pouvons nous ramener au cas où $\mathcal{E}^i = 0$ pour $i > m + 1$. Alors $H^{m + 1}(K) \cong H^{m + 1}(\mathcal{E}^\bullet) =
\Coker(\mathcal{E}^m \to \mathcal{E}^{m + 1})$, qui est de présentation finie.
\end{proof}

\begin{lemma}

\label{cohomology-lemma-n-pseudo-module}
Soit $(X, \mathcal{O}_X)$ un espace annelé et soit $\mathcal{F}$ un faisceau de $\mathcal{O}_X$-modules.
\begin{enumerate}
\item $\mathcal{F}$, considéré comme objet de $D(\mathcal{O}_X)$, est $0$-pseudo-cohérent si et seulement si $\mathcal{F}$ est un $\mathcal{O}_X$-module de type fini, et
\item $\mathcal{F}$, considéré comme objet de $D(\mathcal{O}_X)$, est $(-1)$-pseudo-cohérent si et seulement si $\mathcal{F}$ est un $\mathcal{O}_X$-module de présentation finie.
\end{enumerate}

\end{lemma}

\begin{proof}

Utiliser le lemme \ref{cohomology-lemma-finite-cohomology}
pour démontrer les implications dans un sens, et le lemme \ref{cohomology-lemma-cohomology-pseudo-coherent} pour la réciproque.

\end{proof}





\section{Dimension de Tor}

\label{cohomology-section-tor}

\noindent
Dans cette section, nous examinons de plus près les résolutions par modules plats.

\begin{definition}

\label{cohomology-definition-tor-amplitude}

Soit $(X, \mathcal{O}_X)$ un espace annelé. Soit $E$ un objet de $D(\mathcal{O}_X)$ et soient $a, b \in \mathbf{Z}$ avec $a \leq b$.

\begin{enumerate}

\item On dit que $E$ a une {\it amplitude de Tor dans $[a, b]$}
si $H^i(E \otimes_{\mathcal{O}_X}^\mathbf{L} \mathcal{F}) = 0$
pour tous $\mathcal{O}_X$-modules $\mathcal{F}$ et tous $i \not \in [a, b]$.
\item On dit que $E$ a une {\it dimension de Tor finie}
s'il a une amplitude de Tor dans $[a, b]$ pour certains $a, b$.
\item On dit que $E$ a {\it localement une dimension de Tor finie}
s'il existe un recouvrement ouvert $X = \bigcup U_i$ tel que
$E|_{U_i}$ ait une dimension de Tor finie pour tout $i$.

\end{enumerate}

Un $\mathcal{O}_X$-module $\mathcal{F}$ a une {\it dimension de Tor $\leq d$}
si $\mathcal{F}[0]$, considéré comme un objet de $D(\mathcal{O}_X)$, a une amplitude de Tor dans $[-d, 0]$.

\end{definition}

\noindent
Notons que si $E$, comme dans la définition, a une dimension de Tor finie, alors $E$ appartient à $D^b(\mathcal{O}_X)$, comme on le voit en prenant $\mathcal{F} = \mathcal{O}_X$ dans la définition ci-dessus.

\begin{lemma}

\label{cohomology-lemma-last-one-flat}
Soit $(X, \mathcal{O}_X)$ un espace annelé. Soit $\mathcal{E}^\bullet$ un complexe borné supérieurement de $\mathcal{O}_X$-modules plats, d'amplitude de Tor contenue dans $[a, b]$. Alors $\Coker(d_{\mathcal{E}^\bullet}^{a - 1})$ est un $\mathcal{O}_X$-module plat.
\end{lemma}

\begin{proof}

Comme $\mathcal{E}^\bullet$ est un complexe borné supérieurement de modules plats, nous avons
$\mathcal{E}^\bullet \otimes_{\mathcal{O}_X} \mathcal{F} =
\mathcal{E}^\bullet \otimes_{\mathcal{O}_X}^{\mathbf{L}} \mathcal{F}$
pour tout $\mathcal{O}_X$-module $\mathcal{F}$. Par conséquent, pour tout $\mathcal{O}_X$-module $\mathcal{F}$, la suite
$$
\mathcal{E}^{a - 2} \otimes_{\mathcal{O}_X} \mathcal{F} \to
\mathcal{E}^{a - 1} \otimes_{\mathcal{O}_X} \mathcal{F} \to
\mathcal{E}^a \otimes_{\mathcal{O}_X} \mathcal{F}
$$
est exacte au milieu. Comme
$\mathcal{E}^{a - 2} \to \mathcal{E}^{a - 1} \to \mathcal{E}^a \to
\Coker(d^{a - 1}) \to 0$
est une résolution plate, il en résulte que
$\text{Tor}_1^{\mathcal{O}_X}(\Coker(d^{a - 1}), \mathcal{F}) = 0$
pour tout $\mathcal{O}_X$-module $\mathcal{F}$. Cela signifie que
$\Coker(d^{a - 1})$ est plat ; voir le lemme \ref{cohomology-lemma-flat-tor-zero}.

\end{proof}

\begin{lemma}

\label{cohomology-lemma-tor-amplitude}
Soit $(X, \mathcal{O}_X)$ un espace annelé, soit $E$ un objet de $D(\mathcal{O}_X)$ et soient $a, b \in \mathbf{Z}$ avec $a \leq b$. Les conditions suivantes sont équivalentes :
\begin{enumerate}
\item $E$ a une amplitude de Tor dans $[a, b]$.
\item $E$ est représenté par un complexe $\mathcal{E}^\bullet$ de $\mathcal{O}_X$-modules plats avec $\mathcal{E}^i = 0$ pour $i \not \in [a, b]$.
\end{enumerate}

\end{lemma}

\begin{proof}
Si (2) est vérifiée, nous pouvons calculer $E \otimes_{\mathcal{O}_X}^\mathbf{L} \mathcal{F} =
\mathcal{E}^\bullet \otimes_{\mathcal{O}_X} \mathcal{F}$, et (1) est immédiate.

\medskip\noindent

Supposons (1). Nous pouvons représenter $E$ par un complexe borné supérieurement de $\mathcal{O}_X$-modules plats $\mathcal{K}^\bullet$ ; voir la section \ref{cohomology-section-flat}. Soit $n$ le plus grand entier tel que $\mathcal{K}^n \not = 0$. Si $n > b$, alors $\mathcal{K}^{n - 1} \to \mathcal{K}^n$ est surjectif puisque
$H^n(\mathcal{K}^\bullet) = 0$. Comme $\mathcal{K}^n$ est plat,
$\Ker(\mathcal{K}^{n - 1} \to \mathcal{K}^n)$ est plat (Modules, lemme \ref{modules-lemma-flat-ses}). Nous pouvons donc remplacer $\mathcal{K}^\bullet$ par
$\tau_{\leq n - 1}\mathcal{K}^\bullet$. Par récurrence sur $n$, nous nous ramenons ainsi au cas où $K^\bullet$ est un complexe de
$\mathcal{O}_X$-modules avec $\mathcal{K}^i = 0$ pour $i > b$.

\medskip\noindent

Posons $\mathcal{E}^\bullet = \tau_{\geq a}\mathcal{K}^\bullet$. Tout est clair, sauf la platitude de $\mathcal{E}^a$, qui résulte immédiatement du lemme \ref{cohomology-lemma-last-one-flat}
et des définitions.

\end{proof}

\begin{lemma}

\label{cohomology-lemma-tor-amplitude-pullback}
Soit $f : (X, \mathcal{O}_X) \to (Y, \mathcal{O}_Y)$ un morphisme d'espaces annelés et soit $E$ un objet de $D(\mathcal{O}_Y)$. Si $E$ a une amplitude de Tor dans $[a, b]$, alors $Lf^*E$ a une amplitude de Tor dans $[a, b]$.
\end{lemma}

\begin{proof}
Supposons que $E$ ait une amplitude de Tor dans $[a, b]$. Par le lemme \ref{cohomology-lemma-tor-amplitude}, nous pouvons représenter $E$ par un complexe $\mathcal{E}^\bullet$ de $\mathcal{O}$-modules plats avec $\mathcal{E}^i = 0$ pour $i \not \in [a, b]$. Alors $Lf^*E$ est représenté par $f^*\mathcal{E}^\bullet$. Par Modules, lemme \ref{modules-lemma-base-change-flat}, les modules $f^*\mathcal{E}^i$ sont plats. Le lemme \ref{cohomology-lemma-tor-amplitude} montre donc que $Lf^*E$ a une amplitude de Tor dans $[a, b]$.
\end{proof}

\begin{lemma}

\label{cohomology-lemma-tor-amplitude-stalk}
Soit $(X, \mathcal{O}_X)$ un espace annelé, soit $E$ un objet de $D(\mathcal{O}_X)$ et soient $a, b \in \mathbf{Z}$ avec $a \leq b$. Les conditions suivantes sont équivalentes :
\begin{enumerate}
\item $E$ a une amplitude de Tor dans $[a, b]$.
\item pour tout $x \in X$, l'objet $E_x$ de $D(\mathcal{O}_{X, x})$ a une amplitude de Tor dans $[a, b]$.
\end{enumerate}

\end{lemma}

\begin{proof}
Prendre le germe en $x$ revient à prendre l'image réciproque par le morphisme d'espaces annelés $(x, \mathcal{O}_{X, x}) \to (X, \mathcal{O}_X)$. L'implication (1) $\Rightarrow$ (2) résulte donc du lemme \ref{cohomology-lemma-tor-amplitude-pullback}. Réciproquement, le passage aux germes commute aux produits tensoriels (Modules, lemme \ref{modules-lemma-stalk-tensor-product}). Ainsi, $$
(E \otimes_{\mathcal{O}_X}^\mathbf{L} \mathcal{F})_x =
E_x \otimes_{\mathcal{O}_{X, x}}^\mathbf{L} \mathcal{F}_x
$$ D'autre part, le passage aux germes est exact, donc $$
H^i(E \otimes_{\mathcal{O}_X}^\mathbf{L} \mathcal{F})_x =
H^i((E \otimes_{\mathcal{O}_X}^\mathbf{L} \mathcal{F})_x) =
H^i(E_x \otimes_{\mathcal{O}_{X, x}}^\mathbf{L} \mathcal{F}_x)
$$ Or un faisceau $H^i(E \otimes_{\mathcal{O}_X}^\mathbf{L} \mathcal{F})$ est nul si et seulement si tous ses germes sont nuls (Modules, lemme \ref{modules-lemma-abelian}). Ainsi, (2) implique (1).
\end{proof}
\begin{lemma}

\label{cohomology-lemma-cone-tor-amplitude}

Soit $(X, \mathcal{O}_X)$ un espace annelé. Soit $(K, L, M, f, g, h)$ un triangle distingué dans $D(\mathcal{O}_X)$ et soient $a, b \in \mathbf{Z}$.

\begin{enumerate}

\item Si $K$ a une amplitude de Tor dans $[a + 1, b + 1]$ et
$L$ une amplitude de Tor dans $[a, b]$, alors $M$ a une amplitude de Tor dans $[a, b]$.
\item Si $K$ et $M$ ont une amplitude de Tor dans $[a, b]$, alors
$L$ a une amplitude de Tor dans $[a, b]$.
\item Si $L$ a une amplitude de Tor dans $[a + 1, b + 1]$
et $M$ une amplitude de Tor dans $[a, b]$, alors
$K$ a une amplitude de Tor dans $[a + 1, b + 1]$.

\end{enumerate}

\end{lemma}

\begin{proof}

Démonstration omise. Indication : cela résulte de la suite exacte longue de cohomologie associée à un triangle distingué et du fait que
$- \otimes_{\mathcal{O}_X}^{\mathbf{L}} \mathcal{F}$
préserve les triangles distingués. L'assertion (2) est la plus facile à démontrer ; les autres s'en déduisent par translation.

\end{proof}

\begin{lemma}

\label{cohomology-lemma-tensor-tor-amplitude}
Soit $(X, \mathcal{O}_X)$ un espace annelé et soient $K, L$ des objets de $D(\mathcal{O}_X)$. Si $K$ a une amplitude de Tor dans $[a, b]$ et $L$ une amplitude de Tor dans $[c, d]$, alors $K \otimes_{\mathcal{O}_X}^\mathbf{L} L$ a une amplitude de Tor dans $[a + c, b + d]$.
\end{lemma}

\begin{proof}

Démonstration omise. Indication : utiliser la suite spectrale des Tor.

\end{proof}

\begin{lemma}

\label{cohomology-lemma-summands-tor-amplitude}
Soit $(X, \mathcal{O}_X)$ un espace annelé et soient $a, b \in \mathbf{Z}$. Pour des objets $K$, $L$ de $D(\mathcal{O}_X)$, si $K \oplus L$ a une amplitude de Tor dans $[a, b]$, alors $K$ et $L$ l'ont également.
\end{lemma}

\begin{proof}

Cela résulte immédiatement de l'additivité des foncteurs Tor.

\end{proof}





\section{Complexes parfaits}

\label{cohomology-section-perfect}

\noindent
Dans cette section, nous discutons des propriétés des complexes parfaits sur les espaces annélés.

\begin{definition}

\label{cohomology-definition-perfect}
Soit $(X, \mathcal{O}_X)$ un espace annelé et soit $\mathcal{E}^\bullet$ un complexe de $\mathcal{O}_X$-modules. On dit que $\mathcal{E}^\bullet$ est {\it parfait} s'il existe un recouvrement ouvert $X = \bigcup U_i$ tel que, pour tout $i$, il existe un quasi-isomorphisme de complexes $\mathcal{E}_i^\bullet \to \mathcal{E}^\bullet|_{U_i}$, où $\mathcal{E}_i^\bullet$ est un complexe strictement parfait de $\mathcal{O}_{U_i}$-modules. Un objet $E$ de $D(\mathcal{O}_X)$ est {\it parfait} s'il peut être représenté par un complexe parfait de $\mathcal{O}_X$-modules.
\end{definition}

\noindent
Si $X$ est quasi-compact, tout objet parfait de $D(\mathcal{O}_X)$ appartient à $D^b(\mathcal{O}_X)$. Ce n'est plus nécessairement vrai lorsque $X$ n'est pas quasi-compact.

\begin{lemma}

\label{cohomology-lemma-perfect-independent-representative}
Soit $(X, \mathcal{O}_X)$ un espace annelé et soit $E$ un objet de $D(\mathcal{O}_X)$.
\begin{enumerate}
\item S'il existe un recouvrement ouvert $X = \bigcup U_i$ et des complexes strictement parfaits $\mathcal{E}_i^\bullet$ sur $U_i$ tels que $\mathcal{E}_i^\bullet$ représente $E|_{U_i}$ dans $D(\mathcal{O}_{U_i})$, alors $E$ est parfait.
\item Si $E$ est parfait, alors tout complexe représentant $E$ est parfait.
\end{enumerate}

\end{lemma}

\begin{proof}
La démonstration est identique à celle du lemme \ref{cohomology-lemma-pseudo-coherent-independent-representative}.
\end{proof}

\begin{lemma}

\label{cohomology-lemma-perfect-on-locally-ringed}

Soit $(X, \mathcal{O}_X)$ un espace annelé et soit $E$ un objet de
$D(\mathcal{O}_X)$. Supposons que tous les germes $\mathcal{O}_{X, x}$
soient des anneaux locaux. Alors les conditions suivantes sont équivalentes :

\begin{enumerate}

\item  $E$ est parfait,
\item il existe un recouvrement ouvert $X = \bigcup U_i$ tel que
$E|_{U_i}$ puisse être représenté par un complexe borné de $\mathcal{O}_{U_i}$-modules localement libres de type fini, et
\item il existe un recouvrement ouvert $X = \bigcup U_i$ tel que
$E|_{U_i}$ puisse être représenté par un complexe borné de $\mathcal{O}_{U_i}$-modules libres de type fini.

\end{enumerate}

\end{lemma}

\begin{proof}
Cela résulte du lemme \ref{cohomology-lemma-perfect-independent-representative} et du fait que, sur $X$, tout facteur direct d'un module libre de type fini est localement libre de type fini. Voir Modules, lemme \ref{modules-lemma-direct-summand-of-locally-free-is-locally-free}.
\end{proof}

\begin{lemma}

\label{cohomology-lemma-perfect-precise}
Soit $(X, \mathcal{O}_X)$ un espace annelé, soit $E$ un objet de $D(\mathcal{O}_X)$ et soient $a \leq b$ des entiers. Si $E$ a une amplitude de Tor dans $[a, b]$ et est $(a - 1)$-pseudo-cohérent, alors $E$ est parfait.
\end{lemma}

\begin{proof}

Après avoir remplacé $X$ par les membres d'un recouvrement ouvert, nous pouvons supposer qu'il existe un complexe strictement parfait $\mathcal{E}^\bullet$ et un morphisme
$\alpha : \mathcal{E}^\bullet \to E$ tel que $H^i(\alpha)$ soit un isomorphisme pour $i \geq a$. Remplaçons $\mathcal{E}^\bullet$ par
$\sigma_{\geq a - 1}\mathcal{E}^\bullet$ et choisissons un triangle distingué
$$
\mathcal{E}^\bullet \to E \to C \to \mathcal{E}^\bullet[1]
$$
L'annulation des faisceaux de cohomologie de $E$ et de $\mathcal{E}^\bullet$,
ainsi que l'hypothèse sur $\alpha$, donnent $C \cong \mathcal{K}[2 - a]$ avec
$\mathcal{K} = \Ker(\mathcal{E}^{a - 1} \to \mathcal{E}^a)$. Soit $\mathcal{F}$ un $\mathcal{O}_X$-module. En appliquant $- \otimes_{\mathcal{O}_X}^\mathbf{L} \mathcal{F}$,
l'hypothèse selon laquelle $E$ a une amplitude de Tor dans $[a, b]$
implique que $\mathcal{K} \otimes_{\mathcal{O}_X} \mathcal{F} \to
\mathcal{E}^{a - 1} \otimes_{\mathcal{O}_X} \mathcal{F}$ a l'image
$\Ker(\mathcal{E}^{a - 1} \otimes_{\mathcal{O}_X} \mathcal{F}
\to \mathcal{E}^a \otimes_{\mathcal{O}_X} \mathcal{F})$. Il s'ensuit que $\text{Tor}_1^{\mathcal{O}_X}(\mathcal{E}', \mathcal{F}) = 0$
où $\mathcal{E}' = \Coker(\mathcal{E}^{a - 1} \to \mathcal{E}^a)$. Ainsi, $\mathcal{E}'$ est plat (lemme \ref{cohomology-lemma-flat-tor-zero}). Le module $\mathcal{E}'$ est donc localement facteur direct d'un module libre de type fini par Modules, lemme \ref{modules-lemma-flat-locally-finite-presentation}. Localement, le complexe
$$
\mathcal{E}' \to \mathcal{E}^{a + 1} \to \ldots \to \mathcal{E}^b
$$
est quasi-isomorphe à $E$ ; ainsi, $E$ est parfait.

\end{proof}

\begin{lemma}

\label{cohomology-lemma-perfect}
Soit $(X, \mathcal{O}_X)$ un espace annelé et soit $E$ un objet de $D(\mathcal{O}_X)$. Les conditions suivantes sont équivalentes :
\begin{enumerate}
\item $E$ est parfait, et
\item $E$ est pseudo-cohérent et a localement une dimension finie de Tor.
\end{enumerate}

\end{lemma}

\begin{proof}
Supposons (1). Par définition, il existe un recouvrement ouvert $X = \bigcup U_i$ tel que $E|_{U_i}$ soit représenté par un complexe strictement parfait. Ainsi, $E$ est pseudo-cohérent (c'est-à-dire $m$-pseudo-cohérent pour tout $m$) par le lemme \ref{cohomology-lemma-pseudo-coherent-independent-representative}. De plus, tout facteur direct d'un module libre de type fini est plat ; le lemme \ref{cohomology-lemma-tor-amplitude} montre donc que $E|_{U_i}$ a une dimension de Tor finie. Ainsi, (2) est vérifiée.

\medskip\noindent
Supposons (2). Après avoir remplacé $X$ par les membres d'un recouvrement ouvert, nous pouvons supposer qu'il existe des entiers $a \leq b$ tels que $E$ ait une amplitude de Tor dans $[a, b]$. Puisque $E$ est $m$-pseudo-cohérent pour tout $m$, nous concluons par le lemme \ref{cohomology-lemma-perfect-precise}.
\end{proof}

\begin{lemma}

\label{cohomology-lemma-perfect-pullback}
Soit $f : (X, \mathcal{O}_X) \to (Y, \mathcal{O}_Y)$ un morphisme d'espaces annelés et soit $E$ un objet de $D(\mathcal{O}_Y)$. Si $E$ est parfait dans $D(\mathcal{O}_Y)$, alors $Lf^*E$ est parfait dans $D(\mathcal{O}_X)$.
\end{lemma}

\begin{proof}

Cela résulte des lemmes \ref{cohomology-lemma-perfect},
\ref{cohomology-lemma-tor-amplitude-pullback},
\ref{cohomology-lemma-pseudo-coherent-pullback}. (Une autre démonstration consiste à reprendre celle du lemme \ref{cohomology-lemma-pseudo-coherent-pullback}.)

\end{proof}

\begin{lemma}

\label{cohomology-lemma-two-out-of-three-perfect}

Soit $(X, \mathcal{O}_X)$ un espace annelé et soit $(K, L, M, f, g, h)$
un triangle distingué dans $D(\mathcal{O}_X)$. Si deux des trois objets
$K, L, M$ sont parfaits, alors le troisième l'est également.

\end{lemma}

\begin{proof}

Première démonstration : combiner les lemmes \ref{cohomology-lemma-perfect}, \ref{cohomology-lemma-cone-pseudo-coherent} et
\ref{cohomology-lemma-cone-tor-amplitude}. Deuxième démonstration (esquisse) : supposons $K$ et $L$ parfaits. Après avoir remplacé
$X$ par les membres d'un recouvrement ouvert, nous pouvons supposer que $K$ et $L$
sont représentés par des complexes strictement parfaits $\mathcal{K}^\bullet$
et $\mathcal{L}^\bullet$. Après un nouveau remplacement de $X$ par les membres d'un recouvrement ouvert, nous pouvons supposer que le morphisme $K \to L$ est représenté par un morphisme de complexes $\alpha : \mathcal{K}^\bullet \to \mathcal{L}^\bullet$ ; voir le lemme \ref{cohomology-lemma-local-actual}. Alors $M$ est isomorphe au cône de $\alpha$, qui est strictement parfait par le lemme \ref{cohomology-lemma-cone}.

\end{proof}

\begin{lemma}

\label{cohomology-lemma-tensor-perfect}
Soit $(X, \mathcal{O}_X)$ un espace annelé. Si $K, L$ sont des objets parfaits de $D(\mathcal{O}_X)$, alors $K \otimes_{\mathcal{O}_X}^\mathbf{L} L$ l'est également.
\end{lemma}

\begin{proof}
Cela résulte des lemmes \ref{cohomology-lemma-perfect}, \ref{cohomology-lemma-tensor-pseudo-coherent} et \ref{cohomology-lemma-tensor-tor-amplitude}.
\end{proof}

\begin{lemma}

\label{cohomology-lemma-summands-perfect}
Soit $(X, \mathcal{O}_X)$ un espace annelé. Si $K \oplus L$ est un objet parfait de $D(\mathcal{O}_X)$, alors $K$ et $L$ le sont également.
\end{lemma}

\begin{proof}
Cela résulte des lemmes \ref{cohomology-lemma-perfect}, \ref{cohomology-lemma-summands-pseudo-coherent} et \ref{cohomology-lemma-summands-tor-amplitude}.
\end{proof}

\begin{lemma}

\label{cohomology-lemma-pushforward-perfect}
Soit $(X, \mathcal{O}_X)$ un espace annelé, soit $j : U \to X$ un sous-espace ouvert et soit $E$ un objet parfait de $D(\mathcal{O}_U)$ dont les faisceaux de cohomologie sont supportés sur une partie fermée $T \subset U$, avec $j(T)$ fermé dans $X$. Alors $Rj_*E$ est un objet parfait de $D(\mathcal{O}_X)$.
\end{lemma}

\begin{proof}
La perfection d'un complexe est une propriété locale sur $X$. Il suffit donc de vérifier que les restrictions de $Rj_*E$ à $U$ et à $V = X \setminus j(T)$ sont parfaites. Or $Rj_*E|_U = E$ est parfait, tandis que $Rj_*E|_V = 0$ puisque $E|_{U \setminus T} = 0$.
\end{proof}

\begin{lemma}

\label{cohomology-lemma-perfect-max-open-coh-loc-free}
Soit $(X, \mathcal{O}_X)$ un espace annelé et soit $E$ un objet parfait de $D(\mathcal{O}_X)$. Supposons que tous les germes $\mathcal{O}_{X, x}$ soient des anneaux locaux. Alors l'ensemble $$
U =
\{x \in X \mid
H^i(E)_x\text{ est un }
\mathcal{O}_{X, x}\text{-module libre de type fini pour tout }i\in \mathbf{Z}\}
$$ est ouvert dans $X$ et constitue le plus grand ouvert $U \subset X$ tel que $H^i(E)|_U$ soit localement libre de type fini pour tout $i \in \mathbf{Z}$.
\end{lemma}

\begin{proof}

Notons que si $V \subset X$ est un ouvert tel que $H^i(E)|_V$
soit localement libre de type fini pour tout $i \in \mathbf{Z}$, alors $V \subset U$. Soit $x \in U$. Montrons qu'un voisinage ouvert de $x$
est contenu dans $U$ et que $H^i(E)$ y est localement libre de type fini pour tout $i$. Cela achèvera la démonstration. Au cours de celle-ci, nous pouvons remplacer un nombre fini de fois
$X$ par un voisinage ouvert de $x$. Nous pouvons donc supposer que $E$ est représenté par un complexe strictement parfait $\mathcal{E}^\bullet$. Supposons que
$\mathcal{E}^i = 0$ pour $i \not \in [a, b]$. Procédons par récurrence sur $b - a$. Le module
$H^b(E) = \Coker(d^{b - 1} : \mathcal{E}^{b - 1} \to \mathcal{E}^b)$
est de présentation finie. Comme $H^b(E)_x$ est libre de type fini, le module $H^b(E)$ est libre de type fini sur un voisinage ouvert de $x$, par Modules, lemme \ref{modules-lemma-finite-presentation-stalk-free}. Après avoir remplacé $X$ par un voisinage ouvert éventuellement plus petit, nous pouvons donc supposer que nous disposons d'une décomposition en somme directe
$\mathcal{E}^b = \Im(d^{b - 1}) \oplus H^b(E)$ et
$H^b(E)$ est libre de type fini ; voir le lemme \ref{cohomology-lemma-local-lift-map}. En répétant le même argument, nous pouvons supposer que $\mathcal{E}^{b - 1} = \Ker(d^{b - 1}) \oplus \Im(d^{b - 1})$. Le complexe
$\mathcal{E}^a \to \ldots \to \mathcal{E}^{b - 2} \to \Ker(d^{b - 1})$
est strictement parfait et représente un objet parfait $E'$ tel que $H^i(E) = H^i(E')$ pour $i \not = b$. L'hypothèse de récurrence permet de conclure.

\end{proof}





\section{Duaux}

\label{cohomology-section-duals}

\noindent
Dans cette section, nous caractérisons les objets dualisables de la catégorie des complexes et de la catégorie dérivée. En particulier, nous verrons qu'un objet de $D(\mathcal{O}_X)$ admet un dual si et seulement s'il est parfait (cela résulte de l'exemple \ref{cohomology-example-dual-derived} et du lemme \ref{cohomology-lemma-left-dual-derived}).

\begin{lemma}

\label{cohomology-lemma-symmetric-monoidal-cat-complexes}

Soit $(X, \mathcal{O}_X)$ un espace annelé. La catégorie des complexes de $\mathcal{O}_X$-modules, munie du produit tensoriel défini par
$\mathcal{F}^\bullet \otimes \mathcal{G}^\bullet =
\text{Tot}(\mathcal{F}^\bullet \otimes_{\mathcal{O}_X} \mathcal{G}^\bullet)$
est une catégorie monoïdale symétrique (pour les règles de signes, voir Compléments sur l'algèbre, section \ref{more-algebra-section-sign-rules}).

\end{lemma}

\begin{proof}
Démonstration omise. Indications : comme unité $\mathbf{1}$, on prend le complexe égal à $\mathcal{O}_X$ en degré $0$ et nul dans les autres degrés, avec les isomorphismes évidents $\text{Tot}(\mathbf{1} \otimes_{\mathcal{O}_X} \mathcal{G}^\bullet) =
\mathcal{G}^\bullet$ et $\text{Tot}(\mathcal{F}^\bullet \otimes_{\mathcal{O}_X} \mathbf{1}) =
\mathcal{F}^\bullet$. Pour démontrer le lemme, il faut vérifier la commutativité de divers diagrammes ; voir Catégories, définitions \ref{categories-definition-monoidal-category} et \ref{categories-definition-symmetric-monoidal-category}. Dans chaque cas, les vérifications sont immédiates.
\end{proof}

\begin{example}

\label{cohomology-example-dual}

Soit $(X, \mathcal{O}_X)$ un espace annelé. Soit $\mathcal{F}^\bullet$
un complexe localement borné de $\mathcal{O}_X$-modules tel que chaque
$\mathcal{F}^n$ soit localement un facteur direct d'un $\mathcal{O}_X$-module libre de type fini. Autrement dit, il existe un recouvrement ouvert
$X = \bigcup U_i$ tel que $\mathcal{F}^\bullet|_{U_i}$ soit un complexe strictement parfait. Considérons le complexe
$$
\mathcal{G}^\bullet = \SheafHom^\bullet(\mathcal{F}^\bullet, \mathcal{O}_X)
$$
comme dans la section \ref{cohomology-section-hom-complexes}. Soient
$$
\eta :
\mathcal{O}_X
\to
\text{Tot}(\mathcal{F}^\bullet \otimes_{\mathcal{O}_X} \mathcal{G}^\bullet)
\quad\text{et}\quad
\epsilon :
\text{Tot}(\mathcal{G}^\bullet \otimes_{\mathcal{O}_X} \mathcal{F}^\bullet)
\to
\mathcal{O}_X
$$
les morphismes $\eta = \sum \eta_n$ et $\epsilon = \sum \epsilon_n$,
où $\eta_n : \mathcal{O}_X \to
\mathcal{F}^n \otimes_{\mathcal{O}_X} \mathcal{G}^{-n}$
et
$\epsilon_n : \mathcal{G}^{-n} \otimes_{\mathcal{O}_X} \mathcal{F}^n
\to \mathcal{O}_X$ sont comme dans Modules, exemple \ref{modules-example-dual}. Alors le triplet $\mathcal{G}^\bullet, \eta, \epsilon$
est un dual à gauche de $\mathcal{F}^\bullet$ au sens de Catégories, définition \ref{categories-definition-dual}. Nous omettons la vérification des égalités
$(1 \otimes \epsilon) \circ (\eta \otimes 1) = \text{id}_{\mathcal{F}^\bullet}$
et
$(\epsilon \otimes 1) \circ (1 \otimes \eta) =
\text{id}_{\mathcal{G}^\bullet}$. Comparer avec Compléments sur l'algèbre, lemme \ref{more-algebra-lemma-left-dual-complex}.

\end{example}

\begin{lemma}

\label{cohomology-lemma-left-dual-complex}
Soit $(X, \mathcal{O}_X)$ un espace annelé et soit $\mathcal{F}^\bullet$ un complexe de $\mathcal{O}_X$-modules. Si $\mathcal{F}^\bullet$ admet un dual à gauche dans la catégorie monoïdale des complexes de $\mathcal{O}_X$-modules (Catégories, définition \ref{categories-definition-dual}), alors $\mathcal{F}^\bullet$ est un complexe localement borné dont les termes sont localement des facteurs directs de $\mathcal{O}_X$-modules libres de type fini, et son dual à gauche est celui construit dans l'exemple \ref{cohomology-example-dual}.
\end{lemma}

\begin{proof}

Par l'unicité des duaux à gauche (Catégories, remarque \ref{categories-remark-left-dual-adjoint}), l'assertion finale résultera de ce que $\mathcal{F}^\bullet$
est bien de la forme annoncée. Soit le triplet $\mathcal{G}^\bullet, \eta, \epsilon$ un dual à gauche. Écrivons $\eta = \sum \eta_n$ et $\epsilon = \sum \epsilon_n$,
où $\eta_n : \mathcal{O}_X \to
\mathcal{F}^n \otimes_{\mathcal{O}_X} \mathcal{G}^{-n}$
et
$\epsilon_n : \mathcal{G}^{-n} \otimes_{\mathcal{O}_X} \mathcal{F}^n
\to \mathcal{O}_X$. Comme
$(1 \otimes \epsilon) \circ (\eta \otimes 1) = \text{id}_{\mathcal{F}^\bullet}$
et
$(\epsilon \otimes 1) \circ (1 \otimes \eta) = \text{id}_{\mathcal{G}^\bullet}$
par Catégories, définition \ref{categories-definition-dual}, nous obtenons immédiatement
$(1 \otimes \epsilon_n) \circ (\eta_n \otimes 1) = \text{id}_{\mathcal{F}^n}$
et
$(\epsilon_n \otimes 1) \circ (1 \otimes \eta_n) =
\text{id}_{\mathcal{G}^{-n}}$. Ainsi, $\mathcal{F}^n$ est localement un facteur direct d'un $\mathcal{O}_X$-module libre de type fini, par Modules, lemme \ref{modules-lemma-left-dual-module}. Comme la somme $\eta = \sum \eta_n$ est localement finie,
$\mathcal{F}^\bullet$ est localement borné.

\end{proof}

\begin{lemma}

\label{cohomology-lemma-internal-hom-evaluate-isom}
Soit $(X, \mathcal{O}_X)$ un espace annelé et soient $K, L, M \in D(\mathcal{O}_X)$. Si $K$ est parfait, alors le morphisme $$
R\SheafHom(L, M) \otimes_{\mathcal{O}_X}^\mathbf{L} K
\longrightarrow
R\SheafHom(R\SheafHom(K, L), M)
$$ du lemme \ref{cohomology-lemma-internal-hom-evaluate} est un isomorphisme.
\end{lemma}

\begin{proof}
Comme ce morphisme est défini globalement et que la formation de chacun de ses deux membres commute à la localisation (voir le lemme \ref{cohomology-lemma-restriction-RHom-to-U}), nous pouvons travailler localement sur $X$. Nous pouvons donc supposer que $K$ est représenté par un complexe strictement parfait $\mathcal{E}^\bullet$.

\medskip\noindent
Si $K_1 \to K_2 \to K_3$ est un triangle distingué dans $D(\mathcal{O}_X)$, nous obtenons des triangles distingués $$
R\SheafHom(L, M) \otimes_{\mathcal{O}_X}^\mathbf{L} K_1 \to
R\SheafHom(L, M) \otimes_{\mathcal{O}_X}^\mathbf{L} K_2 \to
R\SheafHom(L, M) \otimes_{\mathcal{O}_X}^\mathbf{L} K_3
$$ et $$
R\SheafHom(R\SheafHom(K_1, L), M) \to
R\SheafHom(R\SheafHom(K_2, L), M) \to
R\SheafHom(R\SheafHom(K_3, L), M)
$$ Voir la section \ref{cohomology-section-flat} et le lemme \ref{cohomology-lemma-RHom-triangulated}. La flèche du lemme \ref{cohomology-lemma-internal-hom-evaluate} est fonctorielle en $K$ ; nous obtenons donc un morphisme entre ces triangles distingués. Ainsi, si le résultat vaut pour $K_1$ et $K_3$, il vaut pour $K_2$ par Catégories dérivées, lemme \ref{derived-lemma-third-isomorphism-triangle}.

\medskip\noindent
En combinant les remarques ci-dessus avec les triangles distingués $$
\sigma_{\geq n}\mathcal{E}^\bullet \to \mathcal{E}^\bullet \to
\sigma_{\leq n - 1}\mathcal{E}^\bullet
$$ associés aux troncatures brutales, nous nous ramenons au cas où $K$ est un facteur direct d'un $\mathcal{O}_X$-module libre de type fini placé en un certain degré. La compatibilité évidente du problème avec les sommes directes (comme ci-dessus) et les décalages nous ramène au cas où $K = \mathcal{O}_X^{\oplus n}$ pour un certain entier $n$. Ce cas est clair.
\end{proof}

\begin{lemma}

\label{cohomology-lemma-dual-perfect-complex}
Soit $(X, \mathcal{O}_X)$ un espace annelé et soit $K$ un objet parfait de $D(\mathcal{O}_X)$. Alors $K^\vee = R\SheafHom(K, \mathcal{O}_X)$ est également parfait et $(K^\vee)^\vee \cong K$. Il existe des isomorphismes fonctoriels $$
M \otimes^\mathbf{L}_{\mathcal{O}_X} K^\vee = R\SheafHom(K, M)
$$ et $$
H^0(X, M \otimes^\mathbf{L}_{\mathcal{O}_X} K^\vee) =
\Hom_{D(\mathcal{O}_X)}(K, M)
$$ pour tout $M$ dans $D(\mathcal{O}_X)$.
\end{lemma}

\begin{proof}
Par le lemme \ref{cohomology-lemma-internal-hom-evaluate}, il existe une application canonique $$
K = R\SheafHom(\mathcal{O}_X, \mathcal{O}_X)
\otimes_{\mathcal{O}_X}^\mathbf{L} K \longrightarrow
R\SheafHom(R\SheafHom(K, \mathcal{O}_X), \mathcal{O}_X) =
(K^\vee)^\vee
$$ qui est un isomorphisme par le lemme \ref{cohomology-lemma-internal-hom-evaluate-isom}. Pour vérifier les autres assertions, nous utiliserons sans autre mention le fait que la formation du Hom interne commute à la restriction aux ouverts (lemme \ref{cohomology-lemma-restriction-RHom-to-U}). La perfection de $K^\vee$ peut se vérifier localement sur $X$. Par le lemme \ref{cohomology-lemma-dual}, il suffit, pour l'assertion finale, de vérifier que le morphisme (\ref{cohomology-equation-eval}) $$
M \otimes^\mathbf{L}_{\mathcal{O}_X} K^\vee
\longrightarrow
R\SheafHom(K, M)
$$ est un isomorphisme. Cette propriété est elle aussi locale sur $X$. Il suffit donc de démontrer ces deux assertions lorsque $K$ est représenté par un complexe strictement parfait.

\medskip\noindent

Supposons que $K$ soit représenté par le complexe strictement parfait
$\mathcal{E}^\bullet$. Il résulte alors du lemme \ref{cohomology-lemma-Rhom-strictly-perfect}
que $K^\vee$ est représenté par le complexe dont les termes sont
$(\mathcal{E}^{-n})^\vee =
\SheafHom_{\mathcal{O}_X}(\mathcal{E}^{-n}, \mathcal{O}_X)$
en degré $n$. Comme $\mathcal{E}^{-n}$ est un facteur direct d'un $\mathcal{O}_X$-module libre de type fini, il en va de même de $(\mathcal{E}^{-n})^\vee$. Ainsi, $K^\vee$ est lui aussi représenté par un complexe strictement parfait, et nous voyons que $K^\vee$ est parfait. Pour voir que (\ref{cohomology-equation-eval}) est un isomorphisme, représentons
$M$ par un complexe $\mathcal{F}^\bullet$. Par le lemme \ref{cohomology-lemma-Rhom-strictly-perfect}, le complexe
$R\SheafHom(K, M)$ est représenté par le complexe avec des termes
$$
\bigoplus\nolimits_{n = p + q}
\SheafHom_{\mathcal{O}_X}(\mathcal{E}^{-q}, \mathcal{F}^p)
$$
D'autre part, l'objet $M \otimes^\mathbf{L}_{\mathcal{O}_X} K^\vee$
est représenté par le complexe avec des termes
$$
\bigoplus\nolimits_{n = p + q}
\mathcal{F}^p \otimes_{\mathcal{O}_X} (\mathcal{E}^{-q})^\vee
$$
L'assertion selon laquelle (\ref{cohomology-equation-eval}) est un isomorphisme se ramène donc à celle selon laquelle le morphisme canonique
$$
\mathcal{F}
\otimes_{\mathcal{O}_X}
\SheafHom_{\mathcal{O}_X}(\mathcal{E}, \mathcal{O}_X)
\longrightarrow
\SheafHom_{\mathcal{O}_X}(\mathcal{E}, \mathcal{F})
$$
est un isomorphisme lorsque $\mathcal{E}$ est un facteur direct d'un $\mathcal{O}_X$-module libre de type fini et $\mathcal{F}$ un $\mathcal{O}_X$-module quelconque. Cela résulte immédiatement de l'énoncé correspondant lorsque
$\mathcal{E}$ est libre de type fini.

\end{proof}

\begin{lemma}

\label{cohomology-lemma-symmetric-monoidal-derived}

Soit $(X, \mathcal{O}_X)$ un espace annelé. La catégorie dérivée
$D(\mathcal{O}_X)$ est une catégorie monoïdale symétrique dont le produit tensoriel est le produit tensoriel dérivé, muni des contraintes usuelles d'associativité et de commutativité (pour les règles de signes, voir Compléments sur l'algèbre, section \ref{more-algebra-section-sign-rules}).

\end{lemma}

\begin{proof}
Démonstration omise. Comparer avec le lemme \ref{cohomology-lemma-symmetric-monoidal-cat-complexes}.
\end{proof}

\begin{example}

\label{cohomology-example-dual-derived}

Soit $(X, \mathcal{O}_X)$ un espace annelé et soit $K$ un objet parfait de $D(\mathcal{O}_X)$. Posons $K^\vee = R\SheafHom(K, \mathcal{O}_X)$
comme dans le lemme \ref{cohomology-lemma-dual-perfect-complex}. Alors le morphisme
$$
K \otimes_{\mathcal{O}_X}^\mathbf{L} K^\vee \longrightarrow R\SheafHom(K, K)
$$
est un isomorphisme (par ce lemme). Notons
$$
\eta :
\mathcal{O}_X
\longrightarrow
K \otimes_{\mathcal{O}_X}^\mathbf{L} K^\vee
$$
le morphisme envoyant $1$ sur la section correspondant à
$\text{id}_K$ sous l'isomorphisme ci-dessus.
$$
\epsilon : 
K^\vee
\otimes_{\mathcal{O}_X}^\mathbf{L} K
\longrightarrow
\mathcal{O}_X
$$
le morphisme d'évaluation (on peut par exemple le construire à l'aide du lemme \ref{cohomology-lemma-internal-hom-composition}). Alors
le triplet $K^\vee, \eta, \epsilon$ est un dual à gauche de $K$ au sens de Catégories, définition \ref{categories-definition-dual}. Nous omettons de vérifier que
$(1 \otimes \epsilon) \circ (\eta \otimes 1) = \text{id}_K$
et
$(\epsilon \otimes 1) \circ (1 \otimes \eta) =
\text{id}_{K^\vee}$.

\end{example}

\begin{lemma}

\label{cohomology-lemma-left-dual-derived}
Soit $(X, \mathcal{O}_X)$ un espace annelé et soit $M$ un objet de $D(\mathcal{O}_X)$. Si $M$ admet un dual à gauche dans la catégorie monoïdale $D(\mathcal{O}_X)$ (Catégories, définition \ref{categories-definition-dual}), alors $M$ est parfait et son dual à gauche est celui construit dans l'exemple \ref{cohomology-example-dual-derived}.
\end{lemma}

\begin{proof}
Soit $x \in X$. Il suffit de trouver un voisinage ouvert $U$ de $x$ tel que la restriction de $M$ à $U$ soit un complexe parfait. Au cours de la démonstration, nous pouvons donc remplacer un nombre fini de fois $X$ par un voisinage ouvert de $x$. Soit le triplet $N, \eta, \epsilon$ un dual à gauche.

\medskip\noindent

Nous utiliserons plusieurs fois l'argument suivant. Choisissons un complexe $\mathcal{M}^\bullet$
de $\mathcal{O}_X$-modules représentant $M$, ainsi qu'un complexe K-plat
$\mathcal{N}^\bullet$ représentant $N$ et dont les termes sont des
$\mathcal{O}_X$-modules plats ; voir le lemme \ref{cohomology-lemma-K-flat-resolution}. Considérons le morphisme
$$
\eta : \mathcal{O}_X \to
\text{Tot}(\mathcal{M}^\bullet \otimes_{\mathcal{O}_X} \mathcal{N}^\bullet)
$$
Après avoir rétréci $X$, on peut trouver un entier $N$ et, pour
$i = 1, \ldots, N$, des entiers $n_i \in \mathbf{Z}$ et des sections
$f_i$ et $g_i$ de $\mathcal{M}^{n_i}$ et $\mathcal{N}^{-n_i}$
tels que
$$
\eta(1) = \sum\nolimits_i f_i \otimes g_i
$$
Soit $\mathcal{K}^\bullet \subset \mathcal{M}^\bullet$ un sous-complexe quelconque de $\mathcal{O}_X$-modules contenant les sections $f_i$
pour $i = 1, \ldots, N$. Comme
$\text{Tot}(\mathcal{K}^\bullet \otimes_{\mathcal{O}_X} \mathcal{N}^\bullet)
\subset
\text{Tot}(\mathcal{M}^\bullet \otimes_{\mathcal{O}_X} \mathcal{N}^\bullet)$
par platitude des modules $\mathcal{N}^n$, le morphisme $\eta$ se factorise par
$$
\tilde \eta :
\mathcal{O}_X \to
\text{Tot}(\mathcal{K}^\bullet \otimes_{\mathcal{O}_X} \mathcal{N}^\bullet)
$$
En notant $K$ l'objet de $D(\mathcal{O}_X)$ représenté par
$\mathcal{K}^\bullet$, nous obtenons un diagramme commutatif
$$
\xymatrix{
M \ar[rr]_-{\eta \otimes 1} \ar[rrd]_{\tilde \eta \otimes 1} & &
M \otimes^\mathbf{L} N \otimes^\mathbf{L} M
\ar[r]_-{1 \otimes \epsilon} &
M \\
& &
K \otimes^\mathbf{L} N \otimes^\mathbf{L} M
\ar[u] \ar[r]^-{1 \otimes \epsilon} &
K \ar[u]
}
$$
Comme la composée de la ligne supérieure est l'identité de $M$,
nous concluons que $M$ est un facteur direct de $K$ dans $D(\mathcal{O}_X)$.

\medskip\noindent
Comme première application de cet argument, choisissons le sous-complexe $\mathcal{K}^\bullet = \sigma_{\geq a} \tau_{\leq b}\mathcal{M}^\bullet$ avec $a < n_i < b$ pour $i = 1, \ldots, N$. Ainsi, $M$ est un facteur direct, dans $D(\mathcal{O}_X)$, d'un complexe borné ; nous pouvons donc supposer que $M$ appartient à $D^b(\mathcal{O}_X)$. (Rappelons que la procédure ci-dessus impose de rétrécir $X$.)

\medskip\noindent
Puisque $M$ appartient à $D^b(\mathcal{O}_X)$, nous pouvons choisir pour $\mathcal{M}^\bullet$ un complexe borné supérieurement de modules plats (par Modules, lemme \ref{modules-lemma-module-quotient-flat}, et Catégories dérivées, lemme \ref{derived-lemma-subcategory-left-resolution}). Dans l'argument ci-dessus, nous pouvons alors prendre $\mathcal{K}^\bullet = \sigma_{\geq a}\mathcal{M}^\bullet$ avec $a < n_i$ pour $i = 1, \ldots, N$. Nous pouvons donc supposer que $M$ est un facteur direct, dans $D(\mathcal{O}_X)$, d'un complexe borné de modules plats. En particulier, $M$ a une amplitude de Tor finie.

\medskip\noindent

Supposons que $M$ ait une amplitude de Tor dans $[a, b]$. Si $M$ est $m$-pseudo-cohérent, montrons qu'après avoir rétréci $X$ nous pouvons supposer que $M$ est
$(m - 1)$-pseudo-cohérent. Le lemme \ref{cohomology-lemma-perfect-precise} achèvera alors la démonstration, puisque
$M$ est de toute façon $(b + 1)$-pseudo-cohérent. Après avoir rétréci $X$, supposons qu'il existe un complexe strictement parfait $\mathcal{E}^\bullet$ et un morphisme $\alpha : \mathcal{E}^\bullet \to M$
dans $D(\mathcal{O}_X)$, tel que $H^i(\alpha)$ soit un isomorphisme pour
$i > m$ et surjectif pour $i = m$. Nous pouvons et allons supposer que $\mathcal{E}^i = 0$ pour $i < m$. Choisissons un triangle distingué
$$
\mathcal{E}^\bullet \to M \to L \to \mathcal{E}^\bullet[1]
$$
Notons que $H^i(L) = 0$ pour $i \geq m$. Nous pouvons donc représenter
$L$ par un complexe $\mathcal{L}^\bullet$ avec $\mathcal{L}^i = 0$
pour $i \geq m$. Le morphisme $L \to \mathcal{E}^\bullet[1]$
est représenté par un morphisme de complexes
$\mathcal{L}^\bullet \to \mathcal{E}^\bullet[1]$
qui est nul dans tous les degrés sauf en degré $m - 1$,
où nous obtenons un morphisme $\mathcal{L}^{m - 1} \to \mathcal{E}^m$ ; voir Catégories dérivées, lemme \ref{derived-lemma-negative-exts}. Alors $M$ est représenté par le complexe
$$
\mathcal{M}^\bullet :
\ldots \to
\mathcal{L}^{m - 2} \to
\mathcal{L}^{m - 1} \to
\mathcal{E}^m \to
\mathcal{E}^{m + 1} \to \ldots
$$
Appliquons à ce complexe la discussion du deuxième paragraphe afin d'obtenir des sections $f_i$ de $\mathcal{M}^{n_i}$ pour $i = 1, \ldots, N$. Pour $n < m$, soit $\mathcal{K}^n \subset \mathcal{L}^n$
le $\mathcal{O}_X$-sous-module engendré par les sections
$f_i$ telles que $n_i = n$ et par les $d(f_i)$ telles que $n_i = n - 1$. Pour $n \geq m$, posons $\mathcal{K}^n = \mathcal{E}^n$. Nous disposons manifestement d'un morphisme de triangles distingués
$$
\xymatrix{
\mathcal{E}^\bullet \ar[r] &
\mathcal{M}^\bullet \ar[r] &
\mathcal{L}^\bullet \ar[r] &
\mathcal{E}^\bullet[1] \\
\mathcal{E}^\bullet \ar[r] \ar[u] &
\mathcal{K}^\bullet \ar[r] \ar[u] &
\sigma_{\leq m - 1}\mathcal{K}^\bullet \ar[r] \ar[u] &
\mathcal{E}^\bullet[1] \ar[u]
}
$$
où tous les morphismes sont ceux indiqués ci-dessus. Notons $K$ l'objet de $D(\mathcal{O}_X)$ correspondant au complexe
$\mathcal{K}^\bullet$. Par les arguments du deuxième paragraphe de la démonstration, nous obtenons un morphisme $s : M \to K$ dans $D(\mathcal{O}_X)$ tel que la composée
$M \to K \to M$ est l'identité sur $M$. Nous ne savons pas que le diagramme
$$
\xymatrix{
\mathcal{E}^\bullet \ar[r] &
\mathcal{K}^\bullet \ar@{=}[r] &
K \\
\mathcal{E}^\bullet \ar[u]^{\text{id}} \ar[r]^i &
\mathcal{M}^\bullet \ar@{=}[r] &
M \ar[u]_s
}
$$
soit commutatif, mais nous savons qu'il le devient après composition avec le morphisme $K \to M$. Par le lemme \ref{cohomology-lemma-local-actual}, après avoir rétréci $X$, nous pouvons supposer que $s \circ i$ est représenté par un morphisme de complexes
$\sigma : \mathcal{E}^\bullet \to \mathcal{K}^\bullet$. Par le même lemme, nous pouvons supposer que la composée de $\sigma$
avec l'inclusion $\mathcal{K}^\bullet \subset \mathcal{M}^\bullet$
est homotope à zéro par une homotopie
$\{h^i : \mathcal{E}^i \to \mathcal{M}^{i - 1}\}$. Ainsi, après avoir remplacé $\mathcal{K}^{m - 1}$ par
$\mathcal{K}^{m - 1} + \Im(h^m)$ (notons qu'après cette opération, $\mathcal{K}^{m - 1}$ reste engendré par un nombre fini de sections globales), nous voyons que
$\sigma$ lui-même est homotope à zéro ! Il existe donc un diagramme commutatif dont les flèches pleines sont fixées
$$
\xymatrix{
\mathcal{E}^\bullet \ar[r] &
M \ar[r] &
\mathcal{L}^\bullet \ar[r] &
\mathcal{E}^\bullet[1] \\
\mathcal{E}^\bullet \ar[r] \ar[u] &
K \ar[r] \ar[u] &
\sigma_{\leq m - 1}\mathcal{K}^\bullet \ar[r] \ar[u] &
\mathcal{E}^\bullet[1] \ar[u] \\
\mathcal{E}^\bullet \ar[r] \ar[u] &
M \ar[r] \ar[u]^s &
\mathcal{L}^\bullet \ar[r] \ar@{..>}[u] &
\mathcal{E}^\bullet[1] \ar[u]
}
$$
Les axiomes des catégories triangulées fournissent une flèche pointillée complétant le diagramme. En considérant les faisceaux de cohomologie en degré $m - 1$, nous obtenons
$$
\xymatrix{
H^{m - 1}(M) \ar[r] &
H^{m - 1}(\mathcal{L}^\bullet) \ar[r] &
H^m(\mathcal{E}^\bullet) \\
H^{m - 1}(K) \ar[r] \ar[u] &
H^{m - 1}(\sigma_{\leq m - 1}\mathcal{K}^\bullet) \ar[r] \ar[u] &
H^m(\mathcal{E}^\bullet) \ar[u] \\
H^{m - 1}(M) \ar[r] \ar[u] &
H^{m - 1}(\mathcal{L}^\bullet) \ar[r] \ar[u] &
H^m(\mathcal{E}^\bullet) \ar[u]
}
$$
Comme les composées verticales sont l'identité dans les colonnes de gauche et de droite, nous concluons que la composée verticale
$H^{m - 1}(\mathcal{L}^\bullet) \to
H^{m - 1}(\sigma_{\leq m - 1}\mathcal{K}^\bullet) \to
H^{m - 1}(\mathcal{L}^\bullet)$ au milieu est surjective ! En particulier, $H^{m - 1}(\sigma_{\leq m - 1}\mathcal{K}^\bullet) \to
H^{m - 1}(\mathcal{L}^\bullet)$ est surjectif. En utilisant le morphisme induit entre les suites exactes longues de faisceaux de cohomologie provenant du morphisme de triangles ci-dessus, une chasse au diagramme montre que $H^i(K) \to H^i(M)$ est un isomorphisme pour $i \geq m$ et surjectif pour $i = m - 1$. Par construction, nous pouvons choisir $r \geq 0$ et une surjection
$\mathcal{O}_X^{\oplus r} \to \mathcal{K}^{m - 1}$. Alors la composée
$$
(\mathcal{O}_X^{\oplus r} \to \mathcal{E}^m \to
\mathcal{E}^{m + 1} \to \ldots ) \longrightarrow
K \longrightarrow M
$$
induit un isomorphisme sur les faisceaux de cohomologie en degrés $\geq m$ et une surjection en degré $m - 1$, ce qui achève la démonstration.

\end{proof}





\section{Divers}

\label{cohomology-section-misc}

\noindent
Quelques résultats qui ne trouvent pas leur place ailleurs.

\begin{lemma}

\label{cohomology-lemma-colim-and-lim-of-duals}
Soit $(X, \mathcal{O}_X)$ un espace annelé. Soit $(K_n)_{n \in \mathbf{N}}$ un système d'objets parfaits de $D(\mathcal{O}_X)$. Soit $K = \text{hocolim} K_n$ la colimite dérivée (Catégories dérivées, définition \ref{derived-definition-derived-colimit}). Alors, pour tout objet $E$ de $D(\mathcal{O}_X)$, nous avons $$
R\SheafHom(K, E) = R\lim E \otimes^\mathbf{L}_{\mathcal{O}_X} K_n^\vee
$$ où $(K_n^\vee)$ est le système projectif des duaux des complexes parfaits.
\end{lemma}

\begin{proof}
Par le lemme \ref{cohomology-lemma-dual-perfect-complex}, nous avons $R\lim E \otimes^\mathbf{L}_{\mathcal{O}_X} K_n^\vee =
R\lim R\SheafHom(K_n, E)$, qui s'insère dans le triangle distingué $$
R\lim R\SheafHom(K_n, E) \to
\prod R\SheafHom(K_n, E) \to
\prod R\SheafHom(K_n, E)
$$ Comme $K$ s'insère de même dans le triangle distingué $\bigoplus K_n \to \bigoplus K_n \to K$, il suffit de montrer que $\prod R\SheafHom(K_n, E) = R\SheafHom(\bigoplus K_n, E)$. Cela résulte formellement de (\ref{cohomology-equation-internal-hom}) et du fait que le produit tensoriel dérivé commute aux sommes directes.
\end{proof}

\begin{lemma}

\label{cohomology-lemma-ext-composition-is-cup}
Soit $(X, \mathcal{O}_X)$ un espace annelé. Soient $K$ et $E$ des objets de $D(\mathcal{O}_X)$, avec $E$ parfait. Le diagramme $$
\xymatrix{
H^0(X, K \otimes_{\mathcal{O}_X}^\mathbf{L} E^\vee) \times H^0(X, E)
\ar[r] \ar[d] &
H^0(X, K \otimes_{\mathcal{O}_X}^\mathbf{L} E^\vee
\otimes_{\mathcal{O}_X}^\mathbf{L} E) \ar[d] \\
\Hom_X(E, K) \times H^0(X, E) \ar[r] &
H^0(X, K)
}
$$ est commutatif ; la flèche horizontale supérieure est le cup-produit, la flèche verticale de droite utilise $\epsilon : E^\vee \otimes_{\mathcal{O}_X}^\mathbf{L} E \to \mathcal{O}_X$ (exemple \ref{cohomology-example-dual-derived}), la flèche verticale de gauche utilise le lemme \ref{cohomology-lemma-dual-perfect-complex}, et la flèche horizontale inférieure est la flèche évidente.
\end{lemma}

\begin{proof}
Nous abrégerons $\otimes = \otimes_{\mathcal{O}_X}^\mathbf{L}$ et $\mathcal{O} = \mathcal{O}_X$. Nous identifierons $E$ et $K$ à $R\SheafHom(\mathcal{O}, E)$ et $R\SheafHom(\mathcal{O}, K)$, et $E^\vee$ à $R\SheafHom(E, \mathcal{O})$.

\medskip\noindent
Soient $\xi \in H^0(X, K \otimes E^\vee)$ et $\eta \in H^0(X, E)$. Notons $\tilde \xi : \mathcal{O} \to K \otimes E^\vee$ et $\tilde \eta : \mathcal{O} \to E$ les morphismes correspondants dans $D(\mathcal{O})$. Par le lemme \ref{cohomology-lemma-second-cup-equals-first}, le cup-produit $\xi \cup \eta$ correspond à $\tilde \xi \otimes \tilde \eta : \mathcal{O} \to
K \otimes E^\vee \otimes E$.

\medskip\noindent

Nous affirmons que le morphisme $\xi' : E \to K$ correspondant à $\xi$ par le lemme \ref{cohomology-lemma-dual-perfect-complex} est la composée
$$
E = \mathcal{O} \otimes E
\xrightarrow{\tilde \xi \otimes 1_E}
K \otimes E^\vee \otimes E
\xrightarrow{1_K \otimes \epsilon}
K
$$
La construction du lemme \ref{cohomology-lemma-dual-perfect-complex}
utilise le morphisme d'évaluation (\ref{cohomology-equation-eval}), lui-même construit à l'aide de l'identification de $E$ à
$R\SheafHom(\mathcal{O}, E)$ et de la composition
$\underline{\circ}$ construite dans le lemme \ref{cohomology-lemma-internal-hom-composition}. Ainsi, $\xi'$ est la composée

\begin{align*}
E = \mathcal{O} \otimes
R\SheafHom(\mathcal{O}, E)
& \xrightarrow{\tilde \xi \otimes 1}
R\SheafHom(\mathcal{O}, K) \otimes
R\SheafHom(E, \mathcal{O}) \otimes
R\SheafHom(\mathcal{O}, E) \\
& \xrightarrow{\underline{\circ} \otimes 1}
R\SheafHom(E, K) \otimes R\SheafHom(\mathcal{O}, E) \\
& \xrightarrow{\underline{\circ}}
R\SheafHom(\mathcal{O}, K) = K
\end{align*}

L'affirmation résulte immédiatement de ce qui précède, de l'associativité de la composition $\underline{\circ}$ construite dans le lemme \ref{cohomology-lemma-internal-hom-composition} (insérer ici une référence future), et du fait que $\epsilon$
est définie comme la composition
$\underline{\circ} : E^\vee \otimes E \to \mathcal{O}$ dans l'exemple \ref{cohomology-example-dual-derived}.

\medskip\noindent

Au vu des deux paragraphes précédents, l'énoncé du lemme revient à dire que
$(1_K \otimes \epsilon) \circ (\tilde \xi \otimes \tilde \eta)$
est égal à
$(1_K \otimes \epsilon) \circ (\tilde \xi \otimes 1_E)
\circ (1_\mathcal{O} \otimes \tilde \eta)$
ce qui est immédiat.

\end{proof}

\begin{lemma}

\label{cohomology-lemma-pullback-internal-hom}

Soit $h : X \to Y$ un morphisme d'espaces annelés. Soient $K, M$ des objets de $D(\mathcal{O}_Y)$. Le morphisme canonique
$$
Lh^*R\SheafHom(K, M) \longrightarrow R\SheafHom(Lh^*K, Lh^*M)
$$
de la remarque \ref{cohomology-remark-prepare-fancy-base-change}
est un isomorphisme dans les cas suivants :

\begin{enumerate}

\item  $K$ est parfait,
\item  $h$ est plat, $K$ est pseudo-cohérent, et $M$ est (localement) borné inférieurement,
\item  $\mathcal{O}_X$ a dimension finie de Tor sur $h^{-1}\mathcal{O}_Y$,
$K$ est pseudo-cohérent, et $M$ est (localement) borné inférieurement,

\end{enumerate}

\end{lemma}

\begin{proof}

(1). La question est locale sur $Y$ ; nous pouvons donc supposer que
$K$ est représenté par un complexe strictement parfait
$\mathcal{E}^\bullet$ ; voir la section \ref{cohomology-section-perfect}. Choisissons un complexe K-plat $\mathcal{F}^\bullet$ représentant $M$. Le lemme \ref{cohomology-lemma-Rhom-strictly-perfect} montre que
$R\SheafHom(K, L)$ est représenté par le complexe
$\mathcal{H}^\bullet =
\SheafHom^\bullet(\mathcal{E}^\bullet, \mathcal{F}^\bullet)$
avec des termes $\mathcal{H}^n = \bigoplus\nolimits_{n = p + q}
\SheafHom_{\mathcal{O}_X}(\mathcal{E}^{-q}, \mathcal{F}^p)$. Par la construction de $Lh^*$ dans la section \ref{cohomology-section-derived-pullback},
nous voyons que $Lh^*K$ est représenté par le complexe strictement parfait
$h^*\mathcal{E}^\bullet$ (lemme \ref{cohomology-lemma-strictly-perfect-pullback}). De même, l'objet $Lh^*M$ est représenté par le complexe $h^*\mathcal{F}^\bullet$. Enfin, l'objet $Lh^*R\SheafHom(K, M)$
est représenté par $h^*\mathcal{H}^\bullet$, puisque
$\mathcal{H}^\bullet$ est K-plat par le lemme \ref{cohomology-lemma-Rhom-strictly-perfect-K-flat}. Pour achever la démonstration, il suffit donc de montrer que
$h^*\mathcal{H}^\bullet =
\SheafHom^\bullet(h^*\mathcal{E}^\bullet, h^*\mathcal{F}^\bullet)$. Pour cela, il suffit de noter que
$h^*\SheafHom(\mathcal{E}, \mathcal{F}) =
\SheafHom(h^*\mathcal{E}, \mathcal{F})$
dès que $\mathcal{E}$ est un facteur direct d'un
$\mathcal{O}_X$-module libre de type fini.

\medskip\noindent

(2). Puisque $h$ est plat, on peut calculer $Lh^*$
en appliquant simplement $h^*$ à tout complexe de $\mathcal{O}_Y$-modules. En particulier, $H^i(Lh^*K) = h^*H^i(K)$ pour tout $i \in \mathbf{Z}$. Supposons que $H^i(M) = 0$ pour $i < a$. Soit $K' \to K$ un morphisme de $D(\mathcal{O}_Y)$ induisant un isomorphisme
$H^i(K') \to H^i(K)$ pour tout $i \geq b$. Alors les morphismes correspondants
$$
R\SheafHom(K, M) \to R\SheafHom(K', M)
$$
et
$$
R\SheafHom(Lh^*K, Lh^*M) \to R\SheafHom(Lh^*K', Lh^*M)
$$
sont des isomorphismes sur les faisceaux de cohomologie en degrés $< a - b$ (détails omis). Ainsi, pour montrer que le morphisme de l'énoncé du lemme induit un isomorphisme sur les faisceaux de cohomologie en degrés $< a - b$, il suffit de démontrer le résultat pour $K'$ dans ces degrés. Comme dans la démonstration de (1), la question est locale sur $Y$. Nous pouvons donc supposer que $K$ est représenté par un complexe strictement parfait ; voir la section \ref{cohomology-section-pseudo-coherent}. Nous sommes ainsi ramenés au cas (1).

\medskip\noindent

(3). La démonstration est la même que celle de (2), sauf que l'on utilise le fait que $Lh^*$ est de dimension cohomologique bornée afin d'obtenir l'annulation voulue. Nous omettons les détails.

\end{proof}

\begin{lemma}

\label{cohomology-lemma-stalk-internal-hom}
Soit $X$ un espace annelé. Soient $K, M$ des objets de $D(\mathcal{O}_X)$. Soit $x \in X$. Le morphisme canonique $$
R\SheafHom(K, M)_x \longrightarrow
R\Hom_{\mathcal{O}_{X, x}}(K_x, M_x)
$$ est un isomorphisme dans les cas suivants
\begin{enumerate}
\item $K$ est parfait,
\item $K$ est pseudo-cohérent et $M$ est (localement) borné inférieurement.
\end{enumerate}

\end{lemma}

\begin{proof}
Soit $Y = \{x\}$ l'espace annelé réduit à un point, muni du faisceau structural donné par $\mathcal{O}_{X, x}$. On applique alors le lemme \ref{cohomology-lemma-pullback-internal-hom} au morphisme d'inclusion plat $Y \to X$.
\end{proof}







\section{Objets inversibles dans la catégorie dérivée}
\label{cohomology-section-invertible-D-or-R}

\noindent
Dans cette section, nous caractérisons les objets inversibles de la catégorie dérivée d'un espace annelé, aussi bien lorsque les germes du faisceau structural sont des anneaux locaux que lorsqu'ils ne le sont pas.

\begin{lemma}

\label{cohomology-lemma-category-summands-finite-free}
Soit $(X, \mathcal{O}_X)$ un espace annelé. Posons $R = \Gamma(X, \mathcal{O}_X)$. La catégorie des $\mathcal{O}_X$-modules qui sont des facteurs directs de $\mathcal{O}_X$-modules libres de type fini est équivalente à la catégorie des $R$-modules projectifs de type fini.
\end{lemma}

\begin{proof}
Observons qu'un $R$-module projectif de type fini est la même chose qu'un facteur direct d'un $R$-module libre de type fini. L'équivalence est donnée par le foncteur $\mathcal{E} \mapsto
\Gamma(X, \mathcal{E})$. Le foncteur inverse est donné par la construction de Modules, lemme \ref{modules-lemma-construct-quasi-coherent-sheaves}.
\end{proof}

\begin{lemma}

\label{cohomology-lemma-invertible-derived}

Soit $(X, \mathcal{O}_X)$ un espace annelé. Soit $M$ un objet de $D(\mathcal{O}_X)$. Les conditions suivantes sont équivalentes :

\begin{enumerate}

\item  $M$ est inversible dans $D(\mathcal{O}_X)$, voir Catégories, définition \ref{categories-definition-invertible},
\item il existe une décomposition localement finie en produit direct
$$
\mathcal{O}_X = \prod\nolimits_{n \in \mathbf{Z}} \mathcal{O}_n
$$
et, pour chaque $n$, il existe un $\mathcal{O}_n$-module inversible
$\mathcal{H}^n$ (Modules, définition \ref{modules-definition-invertible}) et $M = \bigoplus \mathcal{H}^n[-n]$ en $D(\mathcal{O}_X)$.

\end{enumerate}

Si (1) et (2) sont vérifiées, $M$ est un objet parfait de $D(\mathcal{O}_X)$. Si
$\mathcal{O}_{X, x}$ est un anneau local pour tout $x \in X$, ces conditions sont également équivalentes à :

\begin{enumerate}

\item[(3)] il existe un recouvrement ouvert $X = \bigcup U_i$
et, pour chaque $i$, un entier $n_i$ tel que $M|_{U_i}$
soit représenté par un $\mathcal{O}_{U_i}$-module inversible placé en degré $n_i$.

\end{enumerate}

\end{lemma}

\begin{proof}

Supposons (2). Considérons l'objet $R\SheafHom(M, \mathcal{O}_X)$
et le morphisme de composition
$$
R\SheafHom(M, \mathcal{O}_X) \otimes_{\mathcal{O}_X}^\mathbf{L} M \to
\mathcal{O}_X
$$
Pour montrer qu'il s'agit d'un isomorphisme, nous pouvons travailler localement. Nous pouvons donc supposer que $\mathcal{O}_X = \prod_{a \leq n \leq b} \mathcal{O}_n$
et $M = \bigoplus_{a \leq n \leq b} \mathcal{H}^n[-n]$. Alors il suffit de montrer que
$$
R\SheafHom(\mathcal{H}^m, \mathcal{O}_X)
\otimes_{\mathcal{O}_X}^\mathbf{L} \mathcal{H}^n
$$
est nul si $n \not = m$ et égal à $\mathcal{O}_n$ si $n = m$. Le cas $n \not = m$ résulte du fait que $\mathcal{O}_n$ et
$\mathcal{O}_m$ sont des $\mathcal{O}_X$-algèbres plates telles que
$\mathcal{O}_n \otimes_{\mathcal{O}_X} \mathcal{O}_m = 0$. En utilisant la structure locale des $\mathcal{O}_X$-modules inversibles (Modules, lemme \ref{modules-lemma-invertible}) et en travaillant localement, l'isomorphisme dans le cas $n = m$ s'obtient immédiatement ; nous omettons les détails. Puisque $D(\mathcal{O}_X)$ est monoïdale symétrique, nous concluons que $M$ est inversible.

\medskip\noindent
Supposons (1). La description donnée en (2) fournit un candidat pour $\mathcal{O}_n$, à savoir $\SheafHom_{\mathcal{O}_X}(H^n(M), H^n(M))$. Si ces faisceaux d'anneaux forment une famille localement finie et si $\mathcal{O}_X = \prod \mathcal{O}_n$, nous obtenons immédiatement la décomposition en somme directe $M = \bigoplus H^n(M)[-n]$ à l'aide des idempotents de $\mathcal{O}_X$ issus de la décomposition en produit. Cela montre que, pour démontrer (2), nous pouvons travailler localement sur $X$.

\medskip\noindent

Choisissons un objet $N$ de $D(\mathcal{O}_X)$
et un isomorphisme
$M \otimes_{\mathcal{O}_X}^\mathbf{L} N \cong \mathcal{O}_X$. Soit $x \in X$. Alors $N$ est un dual à gauche de $M$ dans la catégorie monoïdale
$D(\mathcal{O}_X)$, et nous concluons que $M$ est parfait par le lemme \ref{cohomology-lemma-left-dual-derived}. Par symétrie,
$N$ est parfait. Après avoir remplacé $X$ par un voisinage ouvert de $x$, nous pouvons supposer que $M$ et $N$ sont représentés par des complexes strictement parfaits $\mathcal{E}^\bullet$ et $\mathcal{F}^\bullet$. Alors $M \otimes_{\mathcal{O}_X}^\mathbf{L} N$ est représenté par
$\text{Tot}(\mathcal{E}^\bullet \otimes_{\mathcal{O}_X} \mathcal{F}^\bullet)$. Après avoir de nouveau rétréci $X$, nous pouvons supposer que les isomorphismes réciproques
$\mathcal{O}_X \to M \otimes_{\mathcal{O}_X}^\mathbf{L} N$ et
$M \otimes_{\mathcal{O}_X}^\mathbf{L} N \to \mathcal{O}_X$
sont donnés par des morphismes de complexes
$$
\alpha : \mathcal{O}_X \to
\text{Tot}(\mathcal{E}^\bullet \otimes_{\mathcal{O}_X} \mathcal{F}^\bullet)
\quad\text{et}\quad
\beta :
\text{Tot}(\mathcal{E}^\bullet \otimes_{\mathcal{O}_X} \mathcal{F}^\bullet)
\to \mathcal{O}_X
$$
Voir le lemme \ref{cohomology-lemma-local-actual}. Alors $\beta \circ \alpha = 1$
comme morphismes de complexes et $\alpha \circ \beta = 1$ comme morphisme dans $D(\mathcal{O}_X)$. Après avoir rétréci $X$,
nous pouvons supposer que la composée $\alpha \circ \beta$ est homotope à $1$
par une homotopie $\theta$ de composantes
$$
\theta^n :
\text{Tot}^n(\mathcal{E}^\bullet \otimes_{\mathcal{O}_X} \mathcal{F}^\bullet)
\to
\text{Tot}^{n - 1}(
\mathcal{E}^\bullet \otimes_{\mathcal{O}_X} \mathcal{F}^\bullet)
$$
par le même lemme que précédemment. Posons $R = \Gamma(X, \mathcal{O}_X)$. Par le lemme \ref{cohomology-lemma-category-summands-finite-free},
nous obtenons :

\begin{enumerate}

\item  $M^\bullet = \Gamma(X, \mathcal{E}^\bullet)$ est un complexe borné de $R$-modules projectifs de type fini,
\item  $N^\bullet = \Gamma(X, \mathcal{F}^\bullet)$ est un complexe borné de $R$-modules projectifs de type fini,
\item  $\alpha$ et $\beta$ correspondent aux morphismes de complexes
$a : R \to \text{Tot}(M^\bullet \otimes_R N^\bullet)$ et
$b : \text{Tot}(M^\bullet \otimes_R N^\bullet) \to R$,
\item  $\theta^n$ correspond à un morphisme
$h^n : \text{Tot}^n(M^\bullet \otimes_R N^\bullet) \to
\text{Tot}^{n - 1}(M^\bullet \otimes_R N^\bullet)$,
\item  $b \circ a = 1$ et $b \circ a - 1 = dh + hd$,

\end{enumerate}

Il en résulte que $M^\bullet$ et $N^\bullet$ définissent des objets réciproquement inverses de $D(R)$. Par Compléments sur l'algèbre, lemme \ref{more-algebra-lemma-invertible-derived},
nous obtenons une décomposition en produit $R = \prod_{a \leq n \leq b} R_n$
et des $R_n$-modules inversibles $H^n$ tels que $M^\bullet \cong \bigoplus_{a \leq n \leq b} H^n[-n]$. Cet isomorphisme dans $D(R)$ se relève en un morphisme
$$
\bigoplus H^n[-n] \longrightarrow M^\bullet
$$
de complexes, car chaque $H^n$ est projectif comme $R$-module. De même, en appliquant une nouvelle fois le lemme \ref{cohomology-lemma-category-summands-finite-free}, nous obtenons un morphisme
$$
\bigoplus H^n \otimes_R \mathcal{O}_X[-n] \to \mathcal{E}^\bullet
$$
qui est un isomorphisme dans $D(\mathcal{O}_X)$. En posant
$\mathcal{O}_n = R_n \otimes_R \mathcal{O}_X$, nous concluons que (2) est vérifiée.

\medskip\noindent

Si tous les germes de $\mathcal{O}_X$ sont locaux, l'équivalence de (2) et (3) se démontre immédiatement. Nous omettons les détails.

\end{proof}





\section{Objets compacts}

\label{cohomology-section-compact}

\noindent
Dans cette section, nous étudions les objets compacts de la catégorie dérivée des modules sur un espace annelé. Rappelons que les objets compacts sont définis dans Catégories dérivées, définition \ref{derived-definition-compact-object}. Sur les espaces annelés convenables, les objets parfaits sont compacts.

\begin{lemma}

\label{cohomology-lemma-when-jshriek-compact}
Soit $X$ un espace annelé. Soit $j : U \to X$ l'inclusion d'un ouvert. Le $\mathcal{O}_X$-module $j_!\mathcal{O}_U$ est un objet compact de $D(\mathcal{O}_X)$ s'il existe un entier $d$ tel que
\begin{enumerate}
\item $H^p(U, \mathcal{F}) = 0$ pour tous les $p > d$, et
\item les foncteurs $\mathcal{F} \mapsto H^p(U, \mathcal{F})$
commutent avec les sommes directes.
\end{enumerate}

\end{lemma}

\begin{proof}

Supposons (1) et (2). Comme
$\Hom(j_!\mathcal{O}_U, \mathcal{F}) = \mathcal{F}(U)$
par Faisceaux, lemme \ref{sheaves-lemma-j-shriek-modules},
nous avons $\Hom(j_!\mathcal{O}_U, K) = R\Gamma(U, K)$ pour
$K$ dans $D(\mathcal{O}_X)$. Il suffit donc de montrer que $R\Gamma(U, -)$ commute aux sommes directes.
La première hypothèse signifie que le foncteur
$F = H^0(U, -)$ est de dimension cohomologique finie. En outre, la seconde implique que toute somme directe de modules injectifs est $F$-acyclique. Soit la famille $K_i$ d'objets de $D(\mathcal{O}_X)$. Choisissons des représentants K-injectifs $I_i^\bullet$ à termes injectifs pour les $K_i$ ; voir Injectifs, théorème
\ref{injectives-theorem-K-injective-embedding-grothendieck}. Puisque nous pouvons calculer $RF$ en appliquant $F$ à tout complexe d'objets acycliques (Catégories dérivées, lemme \ref{derived-lemma-unbounded-right-derived}) et que $\bigoplus K_i$ est représenté par $\bigoplus I_i^\bullet$
(Injectifs, lemme \ref{injectives-lemma-derived-products}), nous concluons que $R\Gamma(U, \bigoplus K_i)$ est représenté par
$\bigoplus H^0(U, I_i^\bullet)$. Ainsi, $R\Gamma(U, -)$ commute aux sommes directes, comme souhaité.

\end{proof}

\begin{lemma}

\label{cohomology-lemma-perfect-is-compact}

Soit $X$ un espace annelé. Supposons que l'espace topologique sous-jacent à $X$ possède les propriétés suivantes :

\begin{enumerate}

\item  $X$ est quasi-compact,
\item il existe une base d'ouverts quasi-compacts, et
\item l'intersection de deux ouverts quasi-compacts est quasi-compact.

\end{enumerate}

Soit $K$ un objet parfait de $D(\mathcal{O}_X)$. Alors :

\begin{enumerate}

\item[(a)] $K$ est un objet compact de $D^+(\mathcal{O}_X)$
au sens suivant : si $M = \bigoplus_{i \in I} M_i$ est borné inférieurement, alors $\Hom(K, M) = \bigoplus_{i \in I} \Hom(K, M_i)$.
\item[b)] Si $X$ est de dimension cohomologique finie, c'est-à-dire s'il existe un $d$ tel que $H^i(X, \mathcal{F}) = 0$ pour $i > d$, alors
$K$ est un objet compact de $D(\mathcal{O}_X)$.

\end{enumerate}

\end{lemma}

\begin{proof}
Soit $K^\vee$ le dual de $K$ ; voir le lemme \ref{cohomology-lemma-dual-perfect-complex}. Nous avons alors $$
\Hom_{D(\mathcal{O}_X)}(K, M) =
H^0(X, K^\vee \otimes_{\mathcal{O}_X}^\mathbf{L} M)
$$ fonctoriellement en $M$ dans $D(\mathcal{O}_X)$. Puisque $K^\vee \otimes_{\mathcal{O}_X}^\mathbf{L} -$ commute aux sommes directes, il suffit de montrer que $R\Gamma(X, -)$ commute aux sommes directes considérées.

\medskip\noindent
Démonstration de (b). Comme $R\Gamma(X, K) = R\Hom(\mathcal{O}_X, K)$ et que $H^p(X, -)$ commute aux sommes directes par le lemme \ref{cohomology-lemma-quasi-separated-cohomology-colimit}, il s'agit d'un cas particulier du lemme \ref{cohomology-lemma-when-jshriek-compact}.

\medskip\noindent

Démonstration de (a). Soient $\mathcal{I}_i$, $i \in I$, une famille de
$\mathcal{O}_X$-modules injectifs. Par le lemme \ref{cohomology-lemma-quasi-separated-cohomology-colimit},
nous avons
$$
H^p(X, \bigoplus\nolimits_{i \in I} \mathcal{I}_i) =
\bigoplus\nolimits_{i \in I} H^p(X, \mathcal{I}_i) = 0
$$
pour tout $p$. Maintenant, si $M = \bigoplus M_i$ est comme en (a), il existe un $a \in \mathbf{Z}$ tel que $H^n(M_i) = 0$
pour $n < a$. Nous pouvons donc choisir des complexes de $\mathcal{O}_X$-modules injectifs
$\mathcal{I}_i^\bullet$ représentant les $M_i$
avec $\mathcal{I}_i^n = 0$ pour $n < a$ ; voir Catégories dérivées, lemme \ref{derived-lemma-injective-resolutions-exist}. Par Injectifs, lemme \ref{injectives-lemma-derived-products},
le complexe somme directe $\bigoplus \mathcal{I}_i^\bullet$
représente $M$. Par acyclicité de Leray (Catégories dérivées, lemme \ref{derived-lemma-leray-acyclicity}), nous voyons que
$$
R\Gamma(X, M) = \Gamma(X, \bigoplus \mathcal{I}_i^\bullet) =
\bigoplus \Gamma(X, \bigoplus \mathcal{I}_i^\bullet) =
\bigoplus R\Gamma(X, M_i)
$$
comme souhaité.

\end{proof}

















\section{Formule de projection}

\label{cohomology-section-projection-formula}

\noindent
Dans cette section, nous rassemblons plusieurs variantes de la formule de projection. La version la plus élémentaire est le lemme \ref{cohomology-lemma-projection-formula}. Après l'avoir énoncée et démontrée, nous étudions une version plus générale faisant intervenir des complexes parfaits.

\begin{lemma}

\label{cohomology-lemma-injective-tensor-finite-locally-free}
Soit $X$ un espace annelé. Soit $\mathcal{I}$ un $\mathcal{O}_X$-module injectif. Soit $\mathcal{E}$ un $\mathcal{O}_X$-module. Supposons que $\mathcal{E}$ soit localement libre de type fini sur $X$ ; voir Modules, définition \ref{modules-definition-locally-free}. Alors $\mathcal{E} \otimes_{\mathcal{O}_X} \mathcal{I}$ est un $\mathcal{O}_X$-module injectif.
\end{lemma}

\begin{proof}
Sous les hypothèses du lemme, nous avons en effet $$
\Hom_{\mathcal{O}_X}(\mathcal{F},
\mathcal{E} \otimes_{\mathcal{O}_X} \mathcal{I})
=
\Hom_{\mathcal{O}_X}(
\mathcal{F} \otimes_{\mathcal{O}_X} \mathcal{E}^\vee, \mathcal{I})
$$ où $\mathcal{E}^\vee = \SheafHom_{\mathcal{O}_X}(\mathcal{E}, \mathcal{O}_X)$ est le dual de $\mathcal{E}$, lui aussi localement libre de type fini. Puisque le produit tensoriel par un module localement libre de type fini est un foncteur exact, on conclut par Homologie, lemme \ref{homology-lemma-characterize-injectives}.
\end{proof}

\begin{lemma}

\label{cohomology-lemma-projection-formula}
Soit $f : X \to Y$ un morphisme d'espaces annelés. Soit $\mathcal{F}$ un $\mathcal{O}_X$-module. Soit $\mathcal{E}$ un $\mathcal{O}_Y$-module. Supposons que $\mathcal{E}$ soit localement libre de type fini sur $Y$ ; voir Modules, définition \ref{modules-definition-locally-free}. Alors il existe des isomorphismes $$
\mathcal{E} \otimes_{\mathcal{O}_Y} R^qf_*\mathcal{F}
\longrightarrow
R^qf_*(f^*\mathcal{E} \otimes_{\mathcal{O}_X} \mathcal{F})
$$ pour tout $q \geq 0$. En fait, il existe un isomorphisme $$
\mathcal{E} \otimes_{\mathcal{O}_Y} Rf_*\mathcal{F}
\longrightarrow
Rf_*(f^*\mathcal{E} \otimes_{\mathcal{O}_X} \mathcal{F})
$$ dans $D^{+}(Y)$, fonctoriel en $\mathcal{F}$.
\end{lemma}

\begin{proof}
Choisissons une résolution injective $\mathcal{F} \to \mathcal{I}^\bullet$ sur $X$. Notons que $f^*\mathcal{E}$ est lui aussi localement libre de type fini ; nous obtenons donc une résolution $$
f^*\mathcal{E} \otimes_{\mathcal{O}_X} \mathcal{F}
\longrightarrow
f^*\mathcal{E} \otimes_{\mathcal{O}_X} \mathcal{I}^\bullet
$$ qui est injective par le lemme \ref{cohomology-lemma-injective-tensor-finite-locally-free}. En appliquant $f_*$, nous obtenons $$
Rf_*(f^*\mathcal{E} \otimes_{\mathcal{O}_X} \mathcal{F})
=
f_*(f^*\mathcal{E} \otimes_{\mathcal{O}_X} \mathcal{I}^\bullet).
$$ Le lemme en résulte donc s'il est possible de montrer que $f_*(f^*\mathcal{E} \otimes_{\mathcal{O}_X} \mathcal{F}) =
\mathcal{E} \otimes_{\mathcal{O}_Y} f_*(\mathcal{F})$ fonctoriellement en le $\mathcal{O}_X$-module $\mathcal{F}$. C'est clair lorsque $\mathcal{E} = \mathcal{O}_Y^{\oplus n}$, et le cas général s'en déduit en travaillant localement sur $Y$. Nous omettons les détails.
\end{proof}
\noindent

Soit $f : X \to Y$ un morphisme d'espaces annelés. Soient $E \in D(\mathcal{O}_X)$ et $K \in D(\mathcal{O}_Y)$. Sans aucune hypothèse supplémentaire, il existe un morphisme

\begin{equation}
\label{cohomology-equation-projection-formula-map}
Rf_*E \otimes^\mathbf{L}_{\mathcal{O}_Y} K
\longrightarrow
Rf_*(E \otimes^\mathbf{L}_{\mathcal{O}_X} Lf^*K)
\end{equation}

Il s'agit du morphisme adjoint au morphisme canonique
$$
Lf^*(Rf_*E \otimes^\mathbf{L}_{\mathcal{O}_Y} K) =
Lf^*Rf_*E \otimes^\mathbf{L}_{\mathcal{O}_X} Lf^*K
\longrightarrow
E \otimes^\mathbf{L}_{\mathcal{O}_X} Lf^*K
$$
provenant du morphisme $Lf^*Rf_*E \to E$ et des lemmes
\ref{cohomology-lemma-pullback-tensor-product} et \ref{cohomology-lemma-adjoint}. Une version raisonnablement générale de la formule de projection est la suivante.

\begin{lemma}

\label{cohomology-lemma-projection-formula-perfect}
Soit $f : X \to Y$ un morphisme d'espaces annelés. Soient $E \in D(\mathcal{O}_X)$ et $K \in D(\mathcal{O}_Y)$. Si $K$ est parfait, alors $$
Rf_*E \otimes^\mathbf{L}_{\mathcal{O}_Y} K =
Rf_*(E \otimes^\mathbf{L}_{\mathcal{O}_X} Lf^*K)
$$ dans $D(\mathcal{O}_Y)$.
\end{lemma}

\begin{proof}

Pour vérifier que (\ref{cohomology-equation-projection-formula-map}) est un isomorphisme, nous pouvons travailler localement sur $Y$, c'est-à-dire chercher un recouvrement ouvert
$\{V_j \to Y\}$
tel que le morphisme se restreigne en un isomorphisme sur chaque $V_j$. Par définition des objets parfaits, nous pouvons donc supposer que $K$ est représenté par un complexe strictement parfait de $\mathcal{O}_Y$-modules. Notons que, de manière tout à fait générale, l'énoncé vaut pour
$K = K_1 \oplus K_2$ si et seulement s'il vaut pour
$K_1$ et $K_2$. Nous pouvons donc supposer que $K$ est un complexe fini de $\mathcal{O}_Y$-modules libres de type fini. Dans ce cas, un argument simple faisant intervenir les troncatures brutales ramène l'énoncé au cas où $K$ est représenté par un $\mathcal{O}_Y$-module libre de type fini. Comme l'énoncé est invariant par passage aux facteurs directs finis en la variable $K$, il suffit de le démontrer pour $K = \mathcal{O}_Y[n]$,
cas dans lequel il est trivial.

\end{proof}

\noindent
Voici un cas où la formule de projection est vraie en généralité complète.

\begin{lemma}

\label{cohomology-lemma-projection-formula-closed-immersion}
Soit $f : X \to Y$ un morphisme d'espaces annelés tel que $f$ soit un homéomorphisme sur une partie fermée. Alors (\ref{cohomology-equation-projection-formula-map}) est toujours un isomorphisme.
\end{lemma}

\begin{proof}

Puisque $f$ est un homéomorphisme sur une partie fermée, le foncteur $f_*$
est exact (Modules, lemme \ref{modules-lemma-i-star-exact}). Ainsi,
$Rf_*$ se calcule en appliquant $f_*$ à n'importe quel complexe représentant l'objet considéré. Choisissons un complexe K-plat $\mathcal{K}^\bullet$ de $\mathcal{O}_Y$-modules représentant $K$, ainsi qu'un complexe quelconque $\mathcal{E}^\bullet$
de $\mathcal{O}_X$-modules représentant $E$. Alors
$Lf^*K$ est représenté par $f^*\mathcal{K}^\bullet$, qui est un complexe K-plat de $\mathcal{O}_X$-modules (lemme \ref{cohomology-lemma-pullback-K-flat}). Ainsi, le membre de droite de (\ref{cohomology-equation-projection-formula-map}) est représenté par
$$
f_*\text{Tot}(\mathcal{E}^\bullet
\otimes_{\mathcal{O}_X} f^*\mathcal{K}^\bullet)
$$
Par le même raisonnement, le membre de gauche est représenté par
$$
\text{Tot}(f_*\mathcal{E}^\bullet \otimes_{\mathcal{O}_Y} \mathcal{K}^\bullet)
$$
Puisque $f_*$ commute aux sommes directes (Modules, lemme \ref{modules-lemma-i-star-right-adjoint}), il suffit de montrer que
$$
f_*(\mathcal{E} \otimes_{\mathcal{O}_X} f^*\mathcal{K}) =
f_*\mathcal{E} \otimes_{\mathcal{O}_Y} \mathcal{K}
$$
pour tout $\mathcal{O}_X$-module $\mathcal{E}$ et $\mathcal{O}_Y$-module
$\mathcal{K}$. Vérifions cette égalité sur les germes. Soit $y \in Y$. Si $y \not \in f(X)$, les germes des deux membres sont nuls. Si $y = f(x)$, il faut montrer que
$$
\mathcal{E}_x \otimes_{\mathcal{O}_{X, x}}
(\mathcal{O}_{X, x} \otimes_{\mathcal{O}_{Y, y}} \mathcal{F}_y) =
\mathcal{E}_x \otimes_{\mathcal{O}_{Y, y}} \mathcal{F}_y
$$
(en utilisant Faisceaux, lemme \ref{sheaves-lemma-stalks-closed-pushforward}
et lemme \ref{sheaves-lemma-stalk-pullback-modules}). Cette égalité est vraie, ce qui démontre le lemme.

\end{proof}

\begin{remark}

\label{cohomology-remark-compatible-with-diagram}

Le morphisme (\ref{cohomology-equation-projection-formula-map}) est compatible avec le morphisme de changement de base de la remarque \ref{cohomology-remark-base-change} au sens suivant. Supposons que
$$
\xymatrix{
X' \ar[r]_{g'} \ar[d]_{f'} &
X \ar[d]^f \\
Y' \ar[r]^g &
Y
}
$$
soit un diagramme commutatif d'espaces annelés. Soient $E \in D(\mathcal{O}_X)$ et $K \in D(\mathcal{O}_Y)$. Alors le diagramme
$$
\xymatrix{
Lg^*(Rf_*E \otimes^\mathbf{L}_{\mathcal{O}_Y} K) \ar[r]_p \ar[d]_t &
Lg^*Rf_*(E \otimes^\mathbf{L}_{\mathcal{O}_X} Lf^*K) \ar[d]_b \\
Lg^*Rf_*E \otimes^\mathbf{L}_{\mathcal{O}_{Y'}} Lg^*K \ar[d]_b &
Rf'_*L(g')^*(E \otimes^\mathbf{L}_{\mathcal{O}_X} Lf^*K) \ar[d]_t \\
Rf'_*L(g')^*E \otimes^\mathbf{L}_{\mathcal{O}_{Y'}} Lg^*K \ar[rd]_p &
Rf'_*(L(g')^*E \otimes^\mathbf{L}_{\mathcal{O}_{Y'}} L(g')^*Lf^*K) \ar[d]_c \\
& Rf'_*(L(g')^*E \otimes^\mathbf{L}_{\mathcal{O}_{Y'}} L(f')^*Lg^*K)
}
$$
est commutatif. Ici, les flèches étiquetées $t$ sont obtenues en appliquant le lemme \ref{cohomology-lemma-pullback-tensor-product}, celles étiquetées $b$ en appliquant la remarque \ref{cohomology-remark-base-change}, celles étiquetées $p$
en appliquant (\ref{cohomology-equation-projection-formula-map}), et
$c$ vient de $L(g')^* \circ Lf^* = L(f')^* \circ Lg^*$. Nous omettons la vérification.

\end{remark}















\section{Un opérateur introduit par Berthelot et Ogus}

\label{cohomology-section-eta}

\noindent
Cette section continue la discussion commencée dans Compléments sur l'algèbre, section \ref{more-algebra-section-eta}. Nous encourageons vivement le lecteur à lire cette section en premier.

\begin{lemma}

\label{cohomology-lemma-invertible-ideal-sheaf}

Soit $(X, \mathcal{O}_X)$ un espace annelé.
Soit $\mathcal{I} \subset \mathcal{O}_X$ un faisceau d'idéaux. Considérons les deux conditions suivantes :

\begin{enumerate}

\item pour tout $x \in X$, il existe un voisinage ouvert
$U \subset X$ de $x$ et un élément $f \in \mathcal{I}(U)$ tels que
$\mathcal{I}|_U = \mathcal{O}_U \cdot f$ et
$f : \mathcal{O}_U \to \mathcal{O}_U$ est injectif, et
\item  $\mathcal{I}$ est inversible comme $\mathcal{O}_X$-module.

\end{enumerate}

Alors (1) implique (2), et la réciproque est vraie si tous les germes
$\mathcal{O}_{X, x}$ du faisceau structural sont des anneaux locaux.

\end{lemma}

\begin{proof}
Démonstration omise. Indication : utiliser Modules, lemme \ref{modules-lemma-invertible-is-locally-free-rank-1}.
\end{proof}

\begin{situation}

\label{cohomology-situation-eta}
Soit $(X, \mathcal{O}_X)$ un espace annelé. Soit $\mathcal{I} \subset \mathcal{O}_X$ un faisceau d'idéaux satisfaisant à la condition (1) du lemme \ref{cohomology-lemma-invertible-ideal-sheaf}\footnote{La discussion de cette section se généralise au cas où l'on suppose seulement que $\mathcal{I}$ est un $\mathcal{O}_X$-module inversible au sens de Modules, section \ref{modules-section-invertible}.}.
\end{situation}

\begin{lemma}

\label{cohomology-lemma-I-torsion-free}

Dans la situation \ref{cohomology-situation-eta},
soit $\mathcal{F}$ un $\mathcal{O}_X$-module. Les conditions suivantes sont équivalentes :

\begin{enumerate}

\item le sous-faisceau $\mathcal{F}[\mathcal{I}] \subset \mathcal{F}$
des sections annulées par $\mathcal{I}$ est nul,
\item le sous-faisceau $\mathcal{F}[\mathcal{I}^n]$ est nul pour tout $n \geq 1$,
\item l'application de multiplication
$\mathcal{I} \otimes_{\mathcal{O}_X} \mathcal{F} \to \mathcal{F}$
est injectif,
\item pour tout ouvert $U \subset X$ tel que
$\mathcal{I}|_U = \mathcal{O}_U \cdot f$
pour un certain $f \in \mathcal{I}(U)$,
le morphisme $f : \mathcal{F}|_U \to \mathcal{F}|_U$ est injectif,
\item pour tout $x \in X$ et tout générateur $f$ de l'idéal
$\mathcal{I}_x \subset \mathcal{O}_{X, x}$, l'élément $f$
est un non-diviseur de zéro sur le germe $\mathcal{F}_x$.

\end{enumerate}

\end{lemma}

\begin{proof}
Démonstration omise.
\end{proof}

\noindent
Dans la situation \ref{cohomology-situation-eta}, soit $\mathcal{F}$ un $\mathcal{O}_X$-module. Si les conditions équivalentes du lemme \ref{cohomology-lemma-I-torsion-free} sont vérifiées, nous dirons que $\mathcal{F}$ est {\it sans $\mathcal{I}$-torsion}. Dans ce cas, pour tout $i \in \mathbf{Z}$, nous noterons $$
\mathcal{I}^i\mathcal{F} =
\mathcal{I}^{\otimes i} \otimes_{\mathcal{O}_X} \mathcal{F}
$$ de sorte que nous ayons des inclusions $$
\ldots \subset
\mathcal{I}^{i + 1}\mathcal{F} \subset
\mathcal{I}^i\mathcal{F} \subset
\mathcal{I}^{i - 1}\mathcal{F} \subset \ldots
$$ Les modules $\mathcal{I}^i\mathcal{F}$ sont localement, mais pas nécessairement globalement, isomorphes à $\mathcal{F}$ comme $\mathcal{O}_X$-modules.

\medskip\noindent

Soit $\mathcal{F}^\bullet$ un complexe sans $\mathcal{I}$-torsion de $\mathcal{O}_X$-modules, de différentielles $d^i : \mathcal{F}^i \to \mathcal{F}^{i + 1}$. Nous définissons alors $\eta_\mathcal{I}\mathcal{F}^\bullet$
comme le complexe dont les termes sont

\begin{align*}
(\eta_\mathcal{I}\mathcal{F})^i
& =
\Ker\left(
d^i, -1 :
\mathcal{I}^i\mathcal{F}^i \oplus \mathcal{I}^{i + 1}\mathcal{F}^{i + 1}
\to
\mathcal{I}^i\mathcal{F}^{i + 1}
\right) \\
& =
\Ker\left(d^i :
\mathcal{I}^i\mathcal{F}^i
\to
\mathcal{I}^i\mathcal{F}^{i + 1}/
\mathcal{I}^{i + 1}\mathcal{F}^{i + 1}
\right)
\end{align*}

et dont la différentielle est induite par $d^i$. Autrement dit, une section locale
$s$ de $(\eta_\mathcal{I}\mathcal{F})^i$ est la même chose qu'une section locale
$s$ de $\mathcal{I}^i\mathcal{F}^i$ telle que son image $d^i(s)$
dans $\mathcal{I}^i\mathcal{F}^{i + 1}$ appartienne au sous-faisceau
$\mathcal{I}^{i + 1}\mathcal{F}^{i + 1}$. Observons que $\eta_\mathcal{I}\mathcal{F}^\bullet$
est encore un complexe de modules sans $\mathcal{I}$-torsion.

\medskip\noindent
Soit $a^\bullet : \mathcal{F}^\bullet \to \mathcal{G}^\bullet$ un morphisme de complexes sans $\mathcal{I}$-torsion de $\mathcal{O}_X$-modules. Nous obtenons alors un morphisme de complexes $$
\eta_\mathcal{I} a^\bullet :
\eta_\mathcal{I}\mathcal{F}^\bullet
\longrightarrow
\eta_\mathcal{I}\mathcal{G}^\bullet
$$ induit par les morphismes $\mathcal{I}^i\mathcal{F}^i \to \mathcal{I}^i\mathcal{G}^i$. Le lecteur vérifiera que l'on obtient ainsi un endofoncteur de la catégorie des complexes sans $\mathcal{I}$-torsion de $\mathcal{O}_X$-modules.

\medskip\noindent

Si $a^\bullet, b^\bullet : \mathcal{F}^\bullet \to \mathcal{G}^\bullet$
sont deux morphismes de complexes sans $\mathcal{I}$-torsion de $\mathcal{O}_X$-modules et si $h = \{h^i : \mathcal{F}^i \to \mathcal{G}^{i - 1}\}$ est une homotopie entre $a^\bullet$ et $b^\bullet$, nous définissons
$\eta_\mathcal{I}h$ comme la famille de morphismes
$(\eta_\mathcal{I}h)^i : (\eta_\mathcal{I}\mathcal{F})^i \to
(\eta_\mathcal{I}\mathcal{G})^{i - 1}$
qui envoie $x$ sur $h^i(x)$ ; cette définition a un sens, car si
$x$ est une section locale de $\mathcal{I}^i\mathcal{F}^i$,
alors $h^i(x)$ est une section locale de $\mathcal{I}^i\mathcal{G}^{i - 1}$,
qui est bien contenue dans $(\eta_\mathcal{I}\mathcal{G})^{i - 1}$. Le lecteur vérifiera que $\eta_\mathcal{I}h$ est une homotopie entre $\eta_\mathcal{I}a^\bullet$ et $\eta_\mathcal{I}b^\bullet$. En définitive, nous obtenons un foncteur
$$
\eta_f :
K(\mathcal{I}\text{-sans torsion }\mathcal{O}_X\text{-modules})
\longrightarrow
K(\mathcal{I}\text{-sans torsion }\mathcal{O}_X\text{-modules})
$$
sur la catégorie d'homotopie (Catégories dérivées, section \ref{derived-section-homotopy}) de la catégorie additive des modules sans
$\mathcal{I}$-torsion sur $\mathcal{O}_X$. Le foncteur $\eta_\mathcal{I}$ n'est en aucun sens un foncteur exact de catégories triangulées ; comparer avec Compléments sur l'algèbre, exemple \ref{more-algebra-example-eta-not-distinguished}.

\begin{lemma}

\label{cohomology-lemma-eta-stalks}
Dans la situation \ref{cohomology-situation-eta}, soit $\mathcal{F}^\bullet$ un complexe sans $\mathcal{I}$-torsion de $\mathcal{O}_X$-modules. Pour $x \in X$, choisissons un générateur $f \in \mathcal{I}_x$. Alors le germe $(\eta_\mathcal{I}\mathcal{F}^\bullet)_x$ est canoniquement isomorphe au complexe $\eta_f\mathcal{F}^\bullet_x$ construit dans Compléments sur l'algèbre, section \ref{more-algebra-section-eta}.
\end{lemma}

\begin{proof}
Démonstration omise.
\end{proof}

\begin{lemma}

\label{cohomology-lemma-eta-first-property}

Dans la situation \ref{cohomology-situation-eta},
soit $\mathcal{F}^\bullet$ un complexe sans
$\mathcal{I}$-torsion de $\mathcal{O}_X$-modules. Il existe un isomorphisme canonique
$$
\mathcal{I}^{\otimes i} \otimes_{\mathcal{O}_X}
\left(
H^i(\mathcal{F}^\bullet)/H^i(\mathcal{F}^\bullet)[\mathcal{I}]
\right)
\longrightarrow H^i(\eta_\mathcal{I}\mathcal{F}^\bullet)
$$
des faisceaux de cohomologie.

\end{lemma}

\begin{proof}

Nous définissons un morphisme
$$
\mathcal{I}^{\otimes i} \otimes_{\mathcal{O}_X} H^i(\mathcal{F}^\bullet)
\longrightarrow
H^i(\eta_\mathcal{I}\mathcal{F}^\bullet)
$$
comme suit. Soit $g$ une section locale de $\mathcal{I}^{\otimes i}$
et soit $\overline{s}$ une section locale de $H^i(\mathcal{F}^\bullet)$. Alors $\overline{s}$ est localement la classe d'une section locale $s$ de
$\Ker(d^i : \mathcal{F}^i \to \mathcal{F}^{i + 1})$. Nous envoyons $g \otimes \overline{s}$ sur la section locale
$gs$ de $(\eta_\mathcal{I}\mathcal{F})^i \subset \mathcal{I}^i\mathcal{F}$. Bien entendu, $gs$ appartient au noyau de $d^i$ sur
$\eta_\mathcal{I}\mathcal{F}^\bullet$ et définit donc une section locale de $H^i(\eta_\mathcal{I}\mathcal{F}^\bullet)$. Il est immédiat que cette définition ne dépend pas des choix. Nous affirmons que ce morphisme se factorise par l'isomorphisme donné dans le lemme. Nous pouvons le vérifier sur les germes ; le lemme \ref{cohomology-lemma-eta-stalks} ramène alors l'assertion au résultat de Compléments sur l'algèbre, lemme \ref{more-algebra-lemma-eta-first-property}.

\end{proof}

\begin{lemma}

\label{cohomology-lemma-eta-qis}
Dans la situation \ref{cohomology-situation-eta}, soit $\mathcal{F}^\bullet \to \mathcal{G}^\bullet$ un morphisme de complexes sans $\mathcal{I}$-torsion de $\mathcal{O}_X$-modules. Alors le morphisme induit $\eta_\mathcal{I}\mathcal{F}^\bullet \to \eta_\mathcal{I}\mathcal{G}^\bullet$ est lui aussi un quasi-isomorphisme.
\end{lemma}

\begin{proof}

Cela résulte de la compatibilité des isomorphismes du lemme \ref{cohomology-lemma-eta-first-property}
avec les morphismes de complexes.

\end{proof}

\begin{lemma}

\label{cohomology-lemma-Leta}
Dans la situation \ref{cohomology-situation-eta}, il existe un foncteur additif\footnote{Attention : ce foncteur n'est pas exact, c'est-à-dire qu'il ne transforme pas les triangles distingués en triangles distingués.} $L\eta_\mathcal{I} : D(\mathcal{O}_X) \to D(\mathcal{O}_X)$ tel que, si $M$ dans $D(\mathcal{O}_X)$ est représenté par un complexe $\mathcal{F}^\bullet$ sans $\mathcal{I}$-torsion de $\mathcal{O}_X$-modules, alors $L\eta_\mathcal{I}M = \eta_\mathcal{I}\mathcal{F}^\bullet$. Il en va de même pour les morphismes.
\end{lemma}

\begin{proof}

Notons $\mathcal{T} \subset \textit{Mod}(\mathcal{O}_X)$ la sous-catégorie pleine des modules sans $\mathcal{I}$-torsion sur $\mathcal{O}_X$. Nous avons l'inclusion correspondante
$$
K(\mathcal{T})
\quad\subset\quad
K(\textit{Mod}(\mathcal{O}_X)) = K(\mathcal{O}_X)
$$
de $K(\mathcal{T})$ comme sous-catégorie triangulée pleine de $K(\mathcal{O}_X)$. Soit $S \subset \text{Arrows}(K(\mathcal{T}))$ l'ensemble des quasi-isomorphismes. Nous allons appliquer Catégories dérivées, lemme \ref{derived-lemma-localization-subcategory},
pour montrer que l'application
$$
S^{-1}K(\mathcal{T}) \longrightarrow D(\mathcal{O}_X)
$$
est une équivalence de catégories triangulées. D'après ce lemme, il suffit de démontrer que, pour tout complexe $\mathcal{G}^\bullet$ de
$\mathcal{O}_X$-modules, il existe un quasi-isomorphisme $\mathcal{F}^\bullet \to \mathcal{G}^\bullet$
avec $\mathcal{F}^\bullet$
un complexe sans $\mathcal{I}$-torsion de $\mathcal{O}_X$-modules. Par le lemme \ref{cohomology-lemma-K-flat-resolution}, nous pouvons trouver un quasi-isomorphisme
$\mathcal{F}^\bullet \to \mathcal{G}^\bullet$ de sorte que le complexe
$\mathcal{F}^\bullet$ soit K-plat (ce dont nous ne nous servirons pas) et formé de $\mathcal{O}_X$-modules plats $\mathcal{F}^i$. Par la troisième caractérisation du lemme \ref{cohomology-lemma-I-torsion-free},
un $\mathcal{O}_X$-module plat est un module sans
$\mathcal{I}$-torsion parmi les $\mathcal{O}_X$-modules, ce qui achève cette étape.

\medskip\noindent

Ces préliminaires étant établis, nous pouvons définir $L\eta_f$. En effet, la discussion qui suit le lemme \ref{cohomology-lemma-I-torsion-free}
nous a déjà fourni un foncteur bien défini
$$
K(\mathcal{T}) \xrightarrow{\eta_f} K(\mathcal{T}) \to
K(\mathcal{O}_X) \to D(\mathcal{O}_X)
$$
qui, d'après le lemme \ref{cohomology-lemma-eta-qis}, envoie les quasi-isomorphismes sur des quasi-isomorphismes. Ce foncteur se factorise donc par
$S^{-1}K(\mathcal{T}) = D(\mathcal{O}_X)$, en vertu de Catégories, lemme \ref{categories-lemma-properties-left-localization}.

\end{proof}

\noindent

Dans la situation \ref{cohomology-situation-eta}, construisons les opérateurs de Bockstein. Observons d'abord qu'il existe un diagramme commutatif
$$
\xymatrix{
0 \ar[r] &
\mathcal{I}^{i + 1} \ar[r] \ar[d] &
\mathcal{I}^i \ar[r] \ar[d] &
\mathcal{I}^i/\mathcal{I}^{i + 1}
\ar[r] \ar@{=}[d] &
0 \\
0 \ar[r] &
\mathcal{I}^{i + 1}/\mathcal{I}^{i + 2} \ar[r] &
\mathcal{I}^i/\mathcal{I}^{i + 2} \ar[r] &
\mathcal{I}^i/\mathcal{I}^{i + 1} \ar[r] &
0
}
$$
dont les lignes sont des suites exactes courtes de $\mathcal{O}_X$-modules. Soit $M$ un objet de $D(\mathcal{O}_X)$. En tensorisant le diagramme ci-dessus par $M$, on obtient un morphisme
$$
\xymatrix{
M \otimes^\mathbf{L} \mathcal{I}^{i + 1} \ar[r] \ar[d] &
M \otimes^\mathbf{L} \mathcal{I}^i \ar[r] \ar[d] &
M \otimes^\mathbf{L} \mathcal{I}^i/\mathcal{I}^{i + 1} \ar[d]^{\text{id}} \\
M \otimes^\mathbf{L} \mathcal{I}^{i + 1}/\mathcal{I}^{i + 2} \ar[r] &
M \otimes^\mathbf{L} \mathcal{I}^i/\mathcal{I}^{i + 2} \ar[r] &
M \otimes^\mathbf{L} \mathcal{I}^i/\mathcal{I}^{i + 1}
}
$$
de triangles distingués. La suite exacte longue des faisceaux de cohomologie associée au triangle inférieur détermine en particulier un
{\it opérateur de Bockstein}

$$
\beta = \beta^i :
H^i(M \otimes^\mathbf{L} \mathcal{I}^i/\mathcal{I}^{i + 1})
\longrightarrow
H^{i + 1}(M \otimes^\mathbf{L} \mathcal{I}^{i + 1}/\mathcal{I}^{i + 2})
$$
pour tout $i \in \mathbf{Z}$. Pour un usage ultérieur, notons que le diagramme commutatif ci-dessus fournit une factorisation

\begin{equation}
\label{cohomology-equation-factorization-bockstein}
\vcenter{
\xymatrix{
H^i(M \otimes^\mathbf{L} \mathcal{I}^i/\mathcal{I}^{i + 1})
\ar[r]_\delta \ar[rd]_\beta &
H^{i + 1}(M \otimes^\mathbf{L} \mathcal{I}^{i + 1}) \ar[d] \\
&
H^{i + 1}(M \otimes^\mathbf{L} \mathcal{I}^{i + 1}/\mathcal{I}^{i + 2})
}
}
\end{equation}

de l'opérateur de Bockstein, où $\delta$ est l'opérateur de bord provenant du triangle distingué supérieur du diagramme commutatif ci-dessus. Nous obtenons un complexe

\begin{equation}
\label{cohomology-equation-complex-bocksteins}
H^\bullet(M/\mathcal{I}) =
\left[
\begin{matrix}
\ldots \\
\downarrow \\
H^{i - 1}(M \otimes^\mathbf{L} \mathcal{I}^{i - 1}/\mathcal{I}^i) \\
\downarrow \beta \\
H^i(M \otimes^\mathbf{L} \mathcal{I}^i/\mathcal{I}^{i + 1}) \\
\downarrow \beta \\
H^{i + 1}(M \otimes^\mathbf{L} \mathcal{I}^{i + 1}/\mathcal{I}^{i + 2}) \\
\downarrow \\
\ldots
\end{matrix}
\right]
\end{equation}

c'est-à-dire que $\beta \circ \beta = 0$. En effet, nous pouvons le vérifier sur les germes et le déduire alors du résultat algébrique correspondant démontré dans Compléments sur l'algèbre, section \ref{more-algebra-section-eta}. Autre démonstration : les suites exactes courtes
$0 \to \mathcal{I}^{i + 1}/\mathcal{I}^{i + 2}
\to \mathcal{I}^i/\mathcal{I}^{i + 2}
\to \mathcal{I}^i/\mathcal{I}^{i + 1} \to 0$
définissent des morphismes
$b^i : \mathcal{I}^i/\mathcal{I}^{i + 1} \to
(\mathcal{I}^{i + 1}/\mathcal{I}^{i + 2})[1]$
en $D(\mathcal{O}_X)$
qui induisent les morphismes $\beta$ ci-dessus en tensorisant par $M$
et en prenant les faisceaux de cohomologie. On montre ensuite que la composée
$b^{i + 1}[1] \circ b^i : \mathcal{I}^i/\mathcal{I}^{i + 1} \to
(\mathcal{I}^{i + 1}/\mathcal{I}^{i + 2})[1] \to
(\mathcal{I}^{i + 2}/\mathcal{I}^{i + 3})[2]$
est nulle dans $D(\mathcal{O}_X)$ en appliquant le critère de Catégories dérivées, lemme \ref{derived-lemma-cup-ext-1-zero},
et en utilisant le fait que le module $\mathcal{I}^i/\mathcal{I}^{i + 3}$
est une extension de $\mathcal{I}^{i + 1}/\mathcal{I}^{i + 3}$
par $\mathcal{I}^i/\mathcal{I}^{i + 1}$.

\begin{lemma}

\label{cohomology-lemma-eta-second-property}

Dans la situation \ref{cohomology-situation-eta}, soit $M$ un objet de
$D(\mathcal{O}_X)$. Il existe un isomorphisme canonique
$$
L\eta_\mathcal{I}M \otimes^\mathbf{L} \mathcal{O}_X/\mathcal{I}
\longrightarrow
H^\bullet(M/\mathcal{I})
$$
dans $D(\mathcal{O}_X)$, où le membre de droite est le complexe (\ref{cohomology-equation-complex-bocksteins}).

\end{lemma}

\begin{proof}
Par la construction de $L\eta_\mathcal{I}$ dans le lemme \ref{cohomology-lemma-eta-qis}, nous pouvons supposer que $M$ est représenté par un complexe sans $\mathcal{I}$-torsion de $\mathcal{O}_X$-modules $\mathcal{F}^\bullet$. Alors $L\eta_\mathcal{I}M$ est représenté par le complexe $\eta_\mathcal{I}\mathcal{F}^\bullet$, lui aussi sans $\mathcal{I}$-torsion comme complexe de $\mathcal{O}_X$-modules. Ainsi, $L\eta_\mathcal{I}M \otimes^\mathbf{L} \mathcal{O}_X/\mathcal{I}$ est représenté par le complexe $\eta_\mathcal{I}\mathcal{F}^\bullet \otimes \mathcal{O}_X/\mathcal{I}$. De même, le complexe $H^\bullet(M/\mathcal{I})$ a pour termes $H^i(\mathcal{F}^\bullet \otimes \mathcal{I}^i/\mathcal{I}^{i + 1})$.

\medskip\noindent

Soit $f$ un générateur local de $\mathcal{I}$. Soit $s$ une section locale de $(\eta_\mathcal{I}\mathcal{F})^i$. Nous pouvons écrire $s = f^is'$ pour une section locale $s'$ de
$\mathcal{F}^i$ et, de même, $d^i(s) = f^{i + 1}t$ pour une section locale $t$ de $\mathcal{F}^{i + 1}$. Ainsi, $d^i$ envoie $f^is'$
sur zéro dans $\mathcal{F}^{i + 1} \otimes \mathcal{I}^i/\mathcal{I}^{i + 1}$. Nous pouvons donc envoyer $s$ sur la classe de $f^is'$ dans
$H^i(\mathcal{F}^\bullet \otimes \mathcal{I}^i/\mathcal{I}^{i + 1})$. Cette règle définit un morphisme
$$
(\eta_\mathcal{I}\mathcal{F})^i \otimes \mathcal{O}_X/\mathcal{I}
\longrightarrow
H^i(\mathcal{F}^\bullet \otimes \mathcal{I}^i/\mathcal{I}^{i + 1})
$$
de $\mathcal{O}_X$-modules. Un calcul montre que ces morphismes sont compatibles aux différentielles (essentiellement parce que $\beta$
envoie la classe de $f^is'$ sur celle de $f^{i + 1}t$), d'où un morphisme de complexes représentant la flèche de l'énoncé du lemme.

\medskip\noindent

Pour achever la démonstration, observons que la construction du paragraphe précédent coïncide sur les germes avec les morphismes construits dans Compléments sur l'algèbre, lemme
\ref{more-algebra-lemma-eta-second-property}
ce qui permet de conclure.

\end{proof}

\begin{lemma}

\label{cohomology-lemma-eta-tensor-invertible}
Dans la situation \ref{cohomology-situation-eta}, soit $\mathcal{F}^\bullet$ un complexe sans $\mathcal{I}$-torsion de $\mathcal{O}_X$-modules. Soit $\mathcal{L}$ un $\mathcal{O}_X$-module inversible. Alors $\eta_\mathcal{I}(\mathcal{F}^\bullet \otimes \mathcal{L}) =
(\eta_\mathcal{I}\mathcal{F}^\bullet) \otimes \mathcal{L}$.
\end{lemma}

\begin{proof}
C'est immédiat par construction.
\end{proof}

\begin{lemma}

\label{cohomology-lemma-eta-cohomology-locally-free}
Dans la situation \ref{cohomology-situation-eta}, soit $M$ un objet de $D(\mathcal{O}_X)$. Soit $x \in X$ tel que $\mathcal{O}_{X, x}$ soit non nul. Si $H^i(M)_x$ est libre de type fini sur $\mathcal{O}_{X, x}$, alors $H^i(L\eta_\mathcal{I}M)_x$ est libre de type fini sur $\mathcal{O}_{X, x}$ et de même rang.
\end{lemma}

\begin{proof}
Soit en effet $f \in \mathcal{O}_{X, x}$ un générateur du germe $\mathcal{I}_x$. Alors $f$ est un non-diviseur de zéro dans $\mathcal{O}_{X, x}$, d'où $H^i(M)_x[f] = 0$. Par le lemme \ref{cohomology-lemma-eta-first-property}, $H^i(L\eta_\mathcal{I}M)_x$ est donc isomorphe à $\mathcal{I}^i_x \otimes_{\mathcal{O}_{X, x}} H^i(M)_x$, qui est libre de même rang, comme souhaité.
\end{proof}













% OK, start here.
%

\chapter{Cohomologie sur les sites}



\phantomsection
\label{sites-cohomology-section-phantom}


\section{Introduction}
\label{sites-cohomology-section-introduction}

\noindent
Dans ce document, nous développons quelques thèmes relatifs à la cohomologie
des faisceaux. Nous étudions le cas des faisceaux sur des sites,
bien que nous reprenions souvent simplement les raisonnements,
les constructions et les démonstrations du cas
topologique.
Les références de base sont \cite{SGA4}, \cite{Godement} et \cite{Iversen}.




\section{Cohomologie des faisceaux}
\label{sites-cohomology-section-cohomology-sheaves}

\noindent
Soit $\mathcal{C}$ un site, voir
Sites, définition \ref{sites-definition-site}.
Soit $\mathcal{F}$ un faisceau abélien sur $\mathcal{C}$.
Nous savons que la catégorie des faisceaux abéliens sur $\mathcal{C}$
possède suffisamment d'injectifs, voir
Injectifs, théorème \ref{injectives-theorem-sheaves-injectives}.
Nous pouvons donc choisir une résolution injective
$\mathcal{F}[0] \to \mathcal{I}^\bullet$.
Pour tout objet $U$ du site $\mathcal{C}$, nous définissons
\begin{equation}
\label{sites-cohomology-equation-cohomology-object-site}
H^i(U, \mathcal{F}) = H^i(\Gamma(U, \mathcal{I}^\bullet))
\end{equation}
comme le {\it $i$-ième groupe de cohomologie du faisceau abélien
$\mathcal{F}$ sur l'objet $U$}. Autrement dit, ce sont les
foncteurs dérivés à droite du foncteur $\mathcal{F} \mapsto \mathcal{F}(U)$.
La famille de foncteurs $H^i(U, -)$ forme un $\delta$-foncteur universel
$\textit{Ab}(\mathcal{C}) \to \textit{Ab}$.

\medskip\noindent
Il arrive parfois que
le site $\mathcal{C}$ ne possède pas d'objet final. Dans ce
cas, nous définissons les {\it sections globales} d'un préfaisceau
d'ensembles $\mathcal{F}$ sur $\mathcal{C}$ comme l'ensemble
\begin{equation}
\label{sites-cohomology-equation-global-sections}
\Gamma(\mathcal{C}, \mathcal{F}) =
\Mor_{\textit{PSh}(\mathcal{C})}(e, \mathcal{F})
\end{equation}
où $e$ est un objet final de la catégorie des préfaisceaux sur $\mathcal{C}$.
Dans ce cas, étant donné un faisceau abélien $\mathcal{F}$ sur $\mathcal{C}$,
nous définissons le {\it $i$-ième groupe de cohomologie de $\mathcal{F}$ sur $\mathcal{C}$}
comme suit :
\begin{equation}
\label{sites-cohomology-equation-cohomology}
H^i(\mathcal{C}, \mathcal{F}) = H^i(\Gamma(\mathcal{C}, \mathcal{I}^\bullet))
\end{equation}
autrement dit, il s'agit du $i$-ième foncteur dérivé à droite du
foncteur des sections globales.
La famille de foncteurs $H^i(\mathcal{C}, -)$ forme un $\delta$-foncteur universel
$\textit{Ab}(\mathcal{C}) \to \textit{Ab}$.

\medskip\noindent
Soit $f : \Sh(\mathcal{C}) \to \Sh(\mathcal{D})$ un morphisme de topos, voir
Sites, définition \ref{sites-definition-topos}.
Avec $\mathcal{F}[0] \to \mathcal{I}^\bullet$ comme ci-dessus,
nous définissons
\begin{equation}
\label{sites-cohomology-equation-higher-direct-image}
R^if_*\mathcal{F} = H^i(f_*\mathcal{I}^\bullet)
\end{equation}
comme la {\it $i$-ième image directe supérieure de $\mathcal{F}$}.
Ce sont les foncteurs dérivés à droite de $f_*$.
La famille de foncteurs $R^if_*$ forme un $\delta$-foncteur universel
de type $\textit{Ab}(\mathcal{C}) \to \textit{Ab}(\mathcal{D})$.

\medskip\noindent
Soit $(\mathcal{C}, \mathcal{O})$ un site annelé, voir
Modules sur les sites, définition \ref{sites-modules-definition-ringed-site}.
Soit $\mathcal{F}$ un $\mathcal{O}$-module.
Nous savons que la catégorie des $\mathcal{O}$-modules
possède suffisamment d'injectifs, voir
Injectifs, théorème \ref{injectives-theorem-sheaves-modules-injectives}.
Nous pouvons donc choisir une résolution injective
$\mathcal{F}[0] \to \mathcal{I}^\bullet$.
Pour tout objet $U$ du site $\mathcal{C}$, nous définissons
\begin{equation}
\label{sites-cohomology-equation-cohomology-object-site-modules}
H^i(U, \mathcal{F}) = H^i(\Gamma(U, \mathcal{I}^\bullet))
\end{equation}
comme le {\it $i$-ième groupe de cohomologie de $\mathcal{F}$ sur $U$}.
La famille de foncteurs $H^i(U, -)$ forme un $\delta$-foncteur universel
$\textit{Mod}(\mathcal{O}) \to \text{Mod}_{\mathcal{O}(U)}$. De même,
\begin{equation}
\label{sites-cohomology-equation-cohomology-modules}
H^i(\mathcal{C}, \mathcal{F}) = H^i(\Gamma(\mathcal{C}, \mathcal{I}^\bullet))
\end{equation}
est le {\it $i$-ième groupe de cohomologie de $\mathcal{F}$ sur $\mathcal{C}$}.
La famille de foncteurs $H^i(\mathcal{C}, -)$ forme un
$\delta$-foncteur universel
$\textit{Mod}(\mathcal{O}) \to \text{Mod}_{\Gamma(\mathcal{C}, \mathcal{O})}$.

\medskip\noindent
Soit $f : (\Sh(\mathcal{C}), \mathcal{O}) \to (\Sh(\mathcal{D}), \mathcal{O}')$
un morphisme de topos annelés, voir
Modules sur les sites, définition \ref{sites-modules-definition-ringed-topos}.
Avec $\mathcal{F}[0] \to \mathcal{I}^\bullet$ comme ci-dessus,
nous définissons
\begin{equation}
\label{sites-cohomology-equation-higher-direct-image-modules}
R^if_*\mathcal{F} = H^i(f_*\mathcal{I}^\bullet)
\end{equation}
comme la {\it $i$-ième image directe supérieure de $\mathcal{F}$}.
Ce sont les foncteurs dérivés à droite de $f_*$.
La famille de foncteurs $R^if_*$ forme un $\delta$-foncteur universel
de type $\textit{Mod}(\mathcal{O}) \to \textit{Mod}(\mathcal{O}')$.








\section{Foncteurs dérivés}
\label{sites-cohomology-section-derived-functors}

\noindent
Nous exposons brièvement une approche des foncteurs dérivés à droite au moyen
de foncteurs de résolution. Supposons donc que $(\mathcal{C}, \mathcal{O})$ soit un site annelé.
Dans ce chapitre, nous écrirons
$$
K(\mathcal{O}) = K(\textit{Mod}(\mathcal{O}))
\quad
\text{et}
\quad
D(\mathcal{O}) = D(\textit{Mod}(\mathcal{O}))
$$
et de même pour les versions bornées des catégories triangulées
introduites dans
Catégories dérivées, définition \ref{derived-definition-complexes-notation} et
définition \ref{derived-definition-unbounded-derived-category}.
D'après
Catégories dérivées, remarque \ref{derived-remark-big-abelian-category},
il existe un foncteur de résolution
$$
j = j_{(\mathcal{C}, \mathcal{O})} :
K^{+}(\textit{Mod}(\mathcal{O}))
\longrightarrow
K^{+}(\mathcal{I})
$$
où $\mathcal{I}$ est la sous-catégorie additive strictement pleine de
$\textit{Mod}(\mathcal{O})$ constituée des $\mathcal{O}$-modules injectifs.
Pour tout foncteur exact à gauche $F : \textit{Mod}(\mathcal{O}) \to \mathcal{B}$
à valeurs dans une catégorie abélienne $\mathcal{B}$, nous noterons $RF$ le
foncteur dérivé à droite de
Catégories dérivées, section \ref{derived-section-right-derived-functor},
construit à l'aide du foncteur de résolution $j$ que nous venons de décrire :
\begin{equation}
\label{sites-cohomology-equation-RF}
RF = F \circ j' : D^{+}(\mathcal{O}) \longrightarrow D^{+}(\mathcal{B})
\end{equation}
voir
Catégories dérivées, lemme \ref{derived-lemma-right-derived-functor},
pour les notations. Remarquons que nous pouvons considérer $RF$ comme défini sur
$\textit{Mod}(\mathcal{O})$, $\text{Comp}^{+}(\textit{Mod}(\mathcal{O}))$, ou
$K^{+}(\mathcal{O})$, selon la situation. D'après
Catégories dérivées, définition \ref{derived-definition-higher-derived-functors},
nous obtenons le $i$-ième foncteur dérivé à droite
\begin{equation}
\label{sites-cohomology-equation-RFi}
R^iF = H^i \circ RF : \textit{Mod}(\mathcal{O}) \longrightarrow \mathcal{B}
\end{equation}
de sorte que $R^0F = F$ et que $\{R^iF, \delta\}_{i \geq 0}$ est un
$\delta$-foncteur universel, voir
Catégories dérivées, lemme \ref{derived-lemma-higher-derived-functors}.

\medskip\noindent
Voici deux cas particuliers de cette construction. Pour un anneau $R$, nous écrivons
$K(R) = K(\text{Mod}_R)$ et $D(R) = D(\text{Mod}_R)$, et de même pour les
versions bornées. Pour tout objet $U$ de $\mathcal{C}$, nous avons un foncteur exact à gauche
$
\Gamma(U, -) :
\textit{Mod}(\mathcal{O})
\longrightarrow
\text{Mod}_{\mathcal{O}(U)}
$
qui donne lieu à
$$
R\Gamma(U, -) :
D^{+}(\mathcal{O})
\longrightarrow
D^{+}(\mathcal{O}(U))
$$
d'après la discussion ci-dessus. Remarquons que $H^i(U, -) = R^i\Gamma(U, -)$
est compatible avec (\ref{sites-cohomology-equation-cohomology-object-site-modules}) ci-dessus.
De même, nous avons
$$
R\Gamma(\mathcal{C}, -) :
D^{+}(\mathcal{O})
\longrightarrow
D^{+}(\Gamma(\mathcal{C}, \mathcal{O}))
$$
compatible avec (\ref{sites-cohomology-equation-cohomology-modules}). Si
$f : (\Sh(\mathcal{C}), \mathcal{O}) \to (\Sh(\mathcal{D}), \mathcal{O}')$
est un morphisme de topos annelés, nous obtenons alors un foncteur exact à gauche
$f_* : \textit{Mod}(\mathcal{O}) \to \textit{Mod}(\mathcal{O}')$
qui donne lieu à l'{\it image directe dérivée}
$$
Rf_* : D^{+}(\mathcal{O}) \to D^+(\mathcal{O}')
$$
Le $i$-ième faisceau de cohomologie de $Rf_*\mathcal{F}^\bullet$ est noté
$R^if_*\mathcal{F}^\bullet$ et appelé la $i$-ième {\it image directe supérieure},
conformément à (\ref{sites-cohomology-equation-higher-direct-image-modules}).
Les foncteurs affichés ci-dessus sont des foncteurs exacts
de catégories dérivées.







\section{Première cohomologie et torseurs}
\label{sites-cohomology-section-h1-torsors}

\begin{definition}
\label{sites-cohomology-definition-torsor}
Soit $\mathcal{C}$ un site.
Soit $\mathcal{G}$ un faisceau de groupes (éventuellement non commutatifs)
sur $\mathcal{C}$.
Un {\it pseudo-torseur}, ou plus précisément un
{\it pseudo-$\mathcal{G}$-torseur}, est un faisceau
d'ensembles $\mathcal{F}$ sur $\mathcal{C}$ muni d'une action
$\mathcal{G} \times \mathcal{F} \to \mathcal{F}$ telle que
\begin{enumerate}
\item dès que $\mathcal{F}(U)$ est non vide, l'action
$\mathcal{G}(U) \times \mathcal{F}(U) \to \mathcal{F}(U)$
est simplement transitive.
\end{enumerate}
Un {\it morphisme de pseudo-$\mathcal{G}$-torseurs}
$\mathcal{F} \to \mathcal{F}'$
est un morphisme de faisceaux d'ensembles compatible aux
actions de $\mathcal{G}$.
Un {\it torseur}, ou plus précisément un
{\it $\mathcal{G}$-torseur}, est un pseudo-$\mathcal{G}$-torseur qui vérifie
en outre
\begin{enumerate}
\item[(2)] pour tout $U \in \Ob(\mathcal{C})$, il existe
un recouvrement $\{U_i \to U\}_{i \in I}$ de $U$
tel que $\mathcal{F}(U_i)$ soit non vide pour tout $i \in I$.
\end{enumerate}
Un {\it morphisme de $\mathcal{G}$-torseurs} est un morphisme de
pseudo-$\mathcal{G}$-torseurs.
Le {\it $\mathcal{G}$-torseur trivial}
est le faisceau $\mathcal{G}$ muni de l'action à gauche
évidente de $\mathcal{G}$.
\end{definition}

\noindent
Il est clair qu'un morphisme de torseurs est automatiquement un isomorphisme.

\begin{lemma}
\label{sites-cohomology-lemma-trivial-torsor}
Soit $\mathcal{C}$ un site.
Soit $\mathcal{G}$ un faisceau de groupes (éventuellement non commutatifs)
sur $\mathcal{C}$.
Un $\mathcal{G}$-torseur $\mathcal{F}$ est trivial si et seulement si
$\Gamma(\mathcal{C}, \mathcal{F}) \not = \emptyset$.
\end{lemma}

\begin{proof}
Omis.
\end{proof}

\begin{lemma}
\label{sites-cohomology-lemma-torsors-h1}
Soit $\mathcal{C}$ un site.
Soit $\mathcal{H}$ un faisceau abélien sur $\mathcal{C}$.
Il existe une bijection canonique entre l'ensemble des classes
d'isomorphisme de $\mathcal{H}$-torseurs et $H^1(\mathcal{C}, \mathcal{H})$.
\end{lemma}

\begin{proof}
Soit $\mathcal{F}$ un $\mathcal{H}$-torseur.
Considérons le faisceau abélien libre $\mathbf{Z}[\mathcal{F}]$
sur $\mathcal{F}$. C'est la faisceautisation de la règle
qui associe à $U \in \Ob(\mathcal{C})$ l'ensemble des sommes formelles
finies $\sum n_i[s_i]$, où $n_i \in \mathbf{Z}$
et $s_i \in \mathcal{F}(U)$. Il existe un morphisme naturel
$$
\sigma : \mathbf{Z}[\mathcal{F}] \longrightarrow \underline{\mathbf{Z}}
$$
qui, à une section locale $\sum n_i[s_i]$, associe $\sum n_i$.
Le noyau de $\sigma$ est engendré par les sections de la forme
$[s] - [s']$. Il existe un morphisme canonique
$a : \Ker(\sigma) \to \mathcal{H}$
qui envoie $[s] - [s'] \mapsto h$, où $h$ est la section locale de
$\mathcal{H}$ telle que $h \cdot s' = s$. Considérons le diagramme cocartésien
$$
\xymatrix{
0 \ar[r] &
\Ker(\sigma) \ar[r] \ar[d]^a &
\mathbf{Z}[\mathcal{F}] \ar[r] \ar[d] &
\underline{\mathbf{Z}} \ar[r] \ar[d] &
0 \\
0 \ar[r] &
\mathcal{H} \ar[r] &
\mathcal{E} \ar[r] &
\underline{\mathbf{Z}} \ar[r] &
0
}
$$
Ici, $\mathcal{E}$ est l'extension obtenue par somme amalgamée.
À partir de la suite exacte longue de cohomologie associée à la
suite exacte courte inférieure, nous obtenons un élément
$\xi = \xi_\mathcal{F} \in H^1(\mathcal{C}, \mathcal{H})$
en appliquant le morphisme de connexion à
$1 \in H^0(\mathcal{C}, \underline{\mathbf{Z}})$.

\medskip\noindent
Réciproquement, étant donné $\xi \in H^1(\mathcal{C}, \mathcal{H})$, nous
pouvons associer à $\xi$ un torseur comme suit. Choisissons un plongement
$\mathcal{H} \to \mathcal{I}$
de $\mathcal{H}$ dans un faisceau abélien injectif $\mathcal{I}$. Posons
$\mathcal{Q} = \mathcal{I}/\mathcal{H}$, de sorte que nous avons une suite
exacte courte
$$
\xymatrix{
0 \ar[r] &
\mathcal{H} \ar[r] &
\mathcal{I} \ar[r] &
\mathcal{Q} \ar[r] &
0
}
$$
L'élément $\xi$ est l'image d'une section globale
$q \in H^0(\mathcal{C}, \mathcal{Q})$
car $H^1(\mathcal{C}, \mathcal{I}) = 0$ (voir
Catégories dérivées, lemme \ref{derived-lemma-higher-derived-functors}).
Soit $\mathcal{F} \subset \mathcal{I}$ le sous-faisceau (d'ensembles) des sections
qui s'envoient sur $q$ dans le faisceau $\mathcal{Q}$. Il est facile de vérifier que
$\mathcal{F}$ est un $\mathcal{H}$-torseur.

\medskip\noindent
Nous omettons la vérification que les deux constructions données
ci-dessus sont réciproques l'une de l'autre.
\end{proof}






\section{Première cohomologie et extensions}
\label{sites-cohomology-section-h1-extensions}

\begin{lemma}
\label{sites-cohomology-lemma-h1-extensions}
Soit $(\mathcal{C}, \mathcal{O})$ un site annelé.
Soit $\mathcal{F}$ un faisceau de $\mathcal{O}$-modules sur $\mathcal{C}$.
Il existe une bijection canonique
$$
\Ext^1_{\textit{Mod}(\mathcal{O})}(\mathcal{O}, \mathcal{F})
\longrightarrow
H^1(\mathcal{C}, \mathcal{F})
$$
qui associe à l'extension
$$
0 \to \mathcal{F} \to \mathcal{E} \to \mathcal{O} \to 0
$$
l'image de $1 \in \Gamma(\mathcal{C}, \mathcal{O})$ dans
$H^1(\mathcal{C}, \mathcal{F})$.
\end{lemma}

\begin{proof}
Construisons l'inverse de l'application donnée dans le lemme.
Soit $\xi \in H^1(\mathcal{C}, \mathcal{F})$.
Choisissons une injection $\mathcal{F} \subset \mathcal{I}$, avec
$\mathcal{I}$ injectif dans $\textit{Mod}(\mathcal{O})$.
Posons $\mathcal{Q} = \mathcal{I}/\mathcal{F}$.
La suite exacte longue de cohomologie montre que
$\xi$ est l'image d'une section
$\tilde \xi \in \Gamma(\mathcal{C}, \mathcal{Q}) =
\Hom_\mathcal{O}(\mathcal{O}, \mathcal{Q})$.
Il suffit alors de former le produit fibré
$$
\xymatrix{
0 \ar[r] &
\mathcal{F} \ar[r] \ar@{=}[d] &
\mathcal{E} \ar[r] \ar[d] &
\mathcal{O} \ar[r] \ar[d]^{\tilde \xi} &
0 \\
0 \ar[r] &
\mathcal{F} \ar[r] &
\mathcal{I} \ar[r] &
\mathcal{Q} \ar[r] &
0
}
$$
voir Homologie, section \ref{homology-section-extensions}.
\end{proof}

\noindent
Le lemme suivant sera remplacé par le lemme plus général
\ref{sites-cohomology-lemma-cohomology-modules-abelian-agree}.

\begin{lemma}
\label{sites-cohomology-lemma-h1-mod-ab-agree}
Soit $(\mathcal{C}, \mathcal{O})$ un site annelé.
Soit $\mathcal{F}$ un faisceau de $\mathcal{O}$-modules sur $\mathcal{C}$.
Notons $\mathcal{F}_{ab}$ le faisceau de groupes abéliens
sous-jacent. Il existe alors un isomorphisme fonctoriel
$$
H^1(\mathcal{C}, \mathcal{F}_{ab})
=
H^1(\mathcal{C}, \mathcal{F})
$$
où le membre de gauche est la cohomologie calculée dans
$\textit{Ab}(\mathcal{C})$ et le membre de droite
est la cohomologie calculée dans $\textit{Mod}(\mathcal{O})$.
\end{lemma}

\begin{proof}
Notons $\underline{\mathbf{Z}}$ le faisceau constant
$\mathbf{Z}$. Comme
$\textit{Ab}(\mathcal{C}) = \textit{Mod}(\underline{\mathbf{Z}})$
nous pouvons appliquer deux fois le
lemme \ref{sites-cohomology-lemma-h1-extensions},
et il suffit donc de montrer
$$
\Ext^1_{\textit{Mod}(\mathcal{O})}(\mathcal{O}, \mathcal{F})
=
\Ext^1_{\textit{Mod}(\underline{\mathbf{Z}})}(
\underline{\mathbf{Z}}, \mathcal{F}_{ab}).
$$
Supposons que $0 \to \mathcal{F} \to \mathcal{E} \to \mathcal{O} \to 0$
soit une extension dans $\textit{Mod}(\mathcal{O})$. Nous pouvons alors utiliser
le morphisme évident de faisceaux abéliens
$1 : \underline{\mathbf{Z}} \to \mathcal{O}$
et former le produit fibré pour obtenir une extension $\mathcal{E}_{ab}$, comme suit :
$$
\xymatrix{
0 \ar[r] &
\mathcal{F}_{ab} \ar[r] \ar@{=}[d] &
\mathcal{E}_{ab} \ar[r] \ar[d] &
\underline{\mathbf{Z}} \ar[r] \ar[d]^{1} &
0 \\
0 \ar[r] &
\mathcal{F} \ar[r] &
\mathcal{E} \ar[r] &
\mathcal{O} \ar[r] &
0
}
$$
La réciproque est un peu plus plaisante. Supposons que
$0 \to \mathcal{F}_{ab} \to \mathcal{E}_{ab} \to \underline{\mathbf{Z}} \to 0$
soit une extension dans $\textit{Mod}(\underline{\mathbf{Z}})$.
Puisque $\underline{\mathbf{Z}}$ est un $\underline{\mathbf{Z}}$-module plat,
la suite
$$
0 \to \mathcal{F}_{ab} \otimes_{\underline{\mathbf{Z}}} \mathcal{O}
\to \mathcal{E}_{ab} \otimes_{\underline{\mathbf{Z}}} \mathcal{O}
\to \underline{\mathbf{Z}} \otimes_{\underline{\mathbf{Z}}} \mathcal{O}
\to 0
$$
est exacte, voir
Modules sur les sites, lemme \ref{sites-modules-lemma-flat-tor-zero}.
Bien entendu,
$\underline{\mathbf{Z}} \otimes_{\underline{\mathbf{Z}}} \mathcal{O}
= \mathcal{O}$.
Nous pouvons donc former la somme amalgamée par le morphisme de multiplication ($\mathcal{O}$-linéaire)
$\mu : \mathcal{F} \otimes_{\underline{\mathbf{Z}}} \mathcal{O}
\to \mathcal{F}$ afin d'obtenir une extension de $\mathcal{O}$ par
$\mathcal{F}$, comme suit :
$$
\xymatrix{
0 \ar[r] &
\mathcal{F}_{ab} \otimes_{\underline{\mathbf{Z}}} \mathcal{O}
\ar[r] \ar[d]^\mu &
\mathcal{E}_{ab} \otimes_{\underline{\mathbf{Z}}} \mathcal{O}
\ar[r] \ar[d] &
\mathcal{O} \ar[r] \ar@{=}[d] &
0 \\
0 \ar[r] &
\mathcal{F} \ar[r] &
\mathcal{E} \ar[r] &
\mathcal{O} \ar[r] &
0
}
$$
qui est l'extension cherchée. Nous omettons de vérifier que ces
constructions sont réciproques l'une de l'autre.
\end{proof}





\section{Première cohomologie et faisceaux inversibles}
\label{sites-cohomology-section-invertible-sheaves}

\noindent
Le groupe de Picard d'un site annelé est défini dans
Modules sur les sites, section \ref{sites-modules-section-invertible}.

\begin{lemma}
\label{sites-cohomology-lemma-h1-invertible}
Soit $(\mathcal{C}, \mathcal{O})$ un site localement annelé.
Il existe un isomorphisme canonique de groupes abéliens
$$
H^1(\mathcal{C}, \mathcal{O}^*) = \Pic(\mathcal{O}).
$$

\end{lemma}

\begin{proof}
Soit $\mathcal{L}$ un $\mathcal{O}$-module inversible.
Considérons le préfaisceau $\mathcal{L}^*$ défini par la règle
$$
U \longmapsto \{s \in \mathcal{L}(U)
\text{ tel que } \mathcal{O}_U \xrightarrow{s \cdot -} \mathcal{L}_U
\text{ is an isomorphism}\}
$$
Ce préfaisceau vérifie la condition de faisceau. De plus, si
$f \in \mathcal{O}^*(U)$ et $s \in \mathcal{L}^*(U)$, alors il est clair que
$fs \in \mathcal{L}^*(U)$. De même, si $s, s' \in \mathcal{L}^*(U)$,
il existe un unique $f \in \mathcal{O}^*(U)$ tel que
$fs = s'$. En outre, le faisceau $\mathcal{L}^*$ possède localement des sections
d'après Modules sur les sites, lemme
\ref{sites-modules-lemma-invertible-is-locally-free-rank-1}.
Autrement dit, nous
voyons que $\mathcal{L}^*$ est un $\mathcal{O}^*$-torseur. Nous obtenons ainsi
une application
$$
\begin{matrix}
\text{set of invertible sheaves on }(\mathcal{C}, \mathcal{O}) \\
\text{ up to isomorphism}
\end{matrix}
\longrightarrow
\begin{matrix}
\text{set of }\mathcal{O}^*\text{-torsors} \\
\text{ up to isomorphism}
\end{matrix}
$$
Nous omettons de vérifier qu'il s'agit d'un morphisme de groupes abéliens.
D'après le
lemme \ref{sites-cohomology-lemma-torsors-h1},
le membre de droite est canoniquement
en bijection avec $H^1(\mathcal{C}, \mathcal{O}^*)$.
Il reste donc à montrer que cette application est injective et surjective.

\medskip\noindent
Injectivité. Si le torseur $\mathcal{L}^*$ est trivial, cela signifie, d'après le
lemme \ref{sites-cohomology-lemma-trivial-torsor},
que $\mathcal{L}^*$ possède une section globale.
Cela signifie donc précisément que $\mathcal{L} \cong \mathcal{O}$ est
l'élément neutre de $\Pic(\mathcal{O})$.

\medskip\noindent
Surjectivité. Soit $\mathcal{F}$ un $\mathcal{O}^*$-torseur.
Considérons le préfaisceau d'ensembles
$$
\mathcal{L}_1 : U \longmapsto
(\mathcal{F}(U) \times \mathcal{O}(U))/\mathcal{O}^*(U)
$$
où l'action de $f \in \mathcal{O}^*(U)$ sur
$(s, g)$ est donnée par $(fs, f^{-1}g)$. Alors $\mathcal{L}_1$ est un préfaisceau
de $\mathcal{O}$-modules si l'on pose
$(s, g) + (s', g') = (s, g + (s'/s)g')$, où $s'/s$ est la section locale
$f$ de $\mathcal{O}^*$ telle que $fs = s'$, et
$h(s, g) = (s, hg)$ pour toute section locale $h$ de $\mathcal{O}$.
Nous omettons de vérifier que la faisceautisation
$\mathcal{L} = \mathcal{L}_1^\#$ est un $\mathcal{O}$-module inversible
dont le $\mathcal{O}^*$-torseur associé $\mathcal{L}^*$ est isomorphe
à $\mathcal{F}$.
\end{proof}









\section{Localité de la cohomologie}
\label{sites-cohomology-section-locality}

\noindent
Le lemme suivant affirme qu'il n'y a aucune ambiguïté dans la définition de la cohomologie
d'un faisceau $\mathcal{F}$ sur un objet du site.

\begin{lemma}
\label{sites-cohomology-lemma-cohomology-of-open}
Soit $(\mathcal{C}, \mathcal{O})$ un site annelé.
Soit $U$ un objet de $\mathcal{C}$.
\begin{enumerate}
\item Si $\mathcal{I}$ est un $\mathcal{O}$-module injectif,
alors $\mathcal{I}|_U$ est un $\mathcal{O}_U$-module injectif.
\item Pour tout faisceau de $\mathcal{O}$-modules $\mathcal{F}$, nous avons
$H^p(U, \mathcal{F}) = H^p(\mathcal{C}/U, \mathcal{F}|_U)$.
\end{enumerate}
\end{lemma}

\begin{proof}
Rappelons que le foncteur $j_U^{-1}$ de restriction à $U$ est adjoint à droite
du foncteur $j_{U!}$ de prolongement par $0$, voir
Modules sur les sites, section
\ref{sites-modules-section-localize}.
De plus, $j_{U!}$ est exact. Ainsi, (1) découle de
Homologie, lemme \ref{homology-lemma-adjoint-preserve-injectives}.

\medskip\noindent
Par définition, $H^p(U, \mathcal{F}) = H^p(\mathcal{I}^\bullet(U))$,
où $\mathcal{F} \to \mathcal{I}^\bullet$ est une résolution injective
dans $\textit{Mod}(\mathcal{O})$.
D'après ce qui précède, $\mathcal{F}|_U \to \mathcal{I}^\bullet|_U$
est une résolution injective dans $\textit{Mod}(\mathcal{O}_U)$.
Par conséquent, $H^p(U, \mathcal{F}|_U)$ est égal à
$H^p(\mathcal{I}^\bullet|_U(U))$.
Bien entendu, $\mathcal{F}(U) = \mathcal{F}|_U(U)$ pour
tout faisceau $\mathcal{F}$ sur $\mathcal{C}$.
D'où l'égalité de (2).
\end{proof}

\noindent
Le lemme suivant servira à étudier ce qui se passe lorsque l'on change
d'univers partiel, ou à comparer la cohomologie des petits et grands sites
\'etales.

\begin{lemma}
\label{sites-cohomology-lemma-cohomology-bigger-site}
Soient $\mathcal{C}$ et $\mathcal{D}$ des sites.
Soit $u : \mathcal{C} \to \mathcal{D}$ un foncteur.
Supposons que $u$ vérifie les hypothèses de
Sites, lemme \ref{sites-lemma-bigger-site}.
Soit $g : \Sh(\mathcal{C}) \to \Sh(\mathcal{D})$
le morphisme de topos associé.
Pour tout faisceau abélien $\mathcal{F}$ sur $\mathcal{D}$, nous avons des
isomorphismes
$$
R\Gamma(\mathcal{C}, g^{-1}\mathcal{F}) = R\Gamma(\mathcal{D}, \mathcal{F}),
$$
en particulier
$H^p(\mathcal{C}, g^{-1}\mathcal{F}) = H^p(\mathcal{D}, \mathcal{F})$
et, pour tout $U \in \Ob(\mathcal{C})$, nous avons des isomorphismes
$$
R\Gamma(U, g^{-1}\mathcal{F}) = R\Gamma(u(U), \mathcal{F}),
$$
en particulier
$H^p(U, g^{-1}\mathcal{F}) = H^p(u(U), \mathcal{F})$. Tous ces
isomorphismes sont fonctoriels en $\mathcal{F}$.
\end{lemma}

\begin{proof}
Puisqu'il est clair que
$\Gamma(\mathcal{C}, g^{-1}\mathcal{F}) = \Gamma(\mathcal{D}, \mathcal{F})$
d'après l'hypothèse (e), il suffit de montrer que $g^{-1}$ envoie les faisceaux
abéliens injectifs sur des faisceaux abéliens injectifs. Comme d'habitude, nous utilisons
Homologie, lemme \ref{homology-lemma-adjoint-preserve-injectives},
pour le voir. L'adjoint à gauche de $g^{-1}$ est $g_! = f^{-1}$, avec les
notations de
Sites, lemme \ref{sites-lemma-bigger-site},
et ce foncteur est exact. Le lemme s'applique donc bien.
\end{proof}

\noindent
Soit $(\mathcal{C}, \mathcal{O})$ un site annelé.
Soit $\mathcal{F}$ un faisceau de $\mathcal{O}$-modules.
Soit $\varphi : U \to V$ un morphisme de $\mathcal{C}$.
Il existe alors un {\it morphisme de restriction} canonique
\begin{equation}
\label{sites-cohomology-equation-restriction-mapping}
H^n(V, \mathcal{F})
\longrightarrow
H^n(U, \mathcal{F}), \quad
\xi \longmapsto \xi|_U
\end{equation}
fonctoriel en $\mathcal{F}$. En effet, choisissons une résolution
injective quelconque $\mathcal{F} \to \mathcal{I}^\bullet$. Les morphismes de
restriction des faisceaux $\mathcal{I}^p$ donnent un morphisme de complexes
$$
\Gamma(V, \mathcal{I}^\bullet)
\longrightarrow
\Gamma(U, \mathcal{I}^\bullet)
$$
Le membre de gauche est un complexe représentant $R\Gamma(V, \mathcal{F})$
et celui de droite un complexe représentant $R\Gamma(U, \mathcal{F})$.
Nous obtenons le morphisme sur les groupes de cohomologie en appliquant le foncteur $H^n$.
Comme indiqué, nous noterons ce morphisme $\xi \mapsto \xi|_U$.
Ainsi, la règle $U \mapsto H^n(U, \mathcal{F})$ est un préfaisceau de
$\mathcal{O}$-modules. Ce préfaisceau est habituellement noté
$\underline{H}^n(\mathcal{F})$. Nous donnerons une autre interprétation
de ce préfaisceau dans le lemme \ref{sites-cohomology-lemma-include}.

\medskip\noindent
Le lemme suivant affirme qu'il est possible d'annuler les classes de cohomologie
supérieure en passant à un recouvrement.

\begin{lemma}
\label{sites-cohomology-lemma-kill-cohomology-class-on-covering}
Soit $(\mathcal{C}, \mathcal{O})$ un site annelé.
Soit $\mathcal{F}$ un faisceau de $\mathcal{O}$-modules.
Soit $U$ un objet de $\mathcal{C}$.
Soient $n > 0$ et $\xi \in H^n(U, \mathcal{F})$.
Il existe alors un recouvrement $\{U_i \to U\}$ de $\mathcal{C}$
tel que $\xi|_{U_i} = 0$ pour tout $i \in I$.
\end{lemma}

\begin{proof}
Soit $\mathcal{F} \to \mathcal{I}^\bullet$ une résolution injective.
Alors
$$
H^n(U, \mathcal{F}) =
\frac{\Ker(\mathcal{I}^n(U) \to \mathcal{I}^{n + 1}(U))}
{\Im(\mathcal{I}^{n - 1}(U) \to \mathcal{I}^n(U))}.
$$
Choisissons un élément $\tilde \xi \in \mathcal{I}^n(U)$ représentant la
classe de cohomologie dans la présentation ci-dessus. Puisque $\mathcal{I}^\bullet$
est une résolution injective de $\mathcal{F}$ et que $n > 0$, nous voyons que
le complexe $\mathcal{I}^\bullet$ est exact en degré $n$. Ainsi,
$\Im(\mathcal{I}^{n - 1} \to \mathcal{I}^n) =
\Ker(\mathcal{I}^n \to \mathcal{I}^{n + 1})$ comme faisceaux.
Puisque $\tilde \xi$ est une section sur $U$ du faisceau noyau,
nous en déduisons qu'il existe un recouvrement $\{U_i \to U\}$ du site
tel que $\tilde \xi|_{U_i}$ soit l'image par $d$ d'une section
$\xi_i \in \mathcal{I}^{n - 1}(U_i)$. Par définition, la
restriction $\xi|_{U_i}$ correspond à la classe de
$\tilde \xi|_{U_i}$, ce qui permet de conclure.
\end{proof}

\begin{lemma}
\label{sites-cohomology-lemma-higher-direct-images}
Soit $f : (\mathcal{C}, \mathcal{O}_\mathcal{C}) \to
(\mathcal{D}, \mathcal{O}_\mathcal{D})$ un morphisme de sites annelés
correspondant au foncteur continu $u : \mathcal{D} \to \mathcal{C}$.
Pour tout $\mathcal{F} \in \Ob(\textit{Mod}(\mathcal{O}_\mathcal{C}))$,
le faisceau $R^if_*\mathcal{F}$ est le faisceau associé au
préfaisceau
$$
V \longmapsto H^i(u(V), \mathcal{F})
$$
\end{lemma}

\begin{proof}
Soit $\mathcal{F} \to \mathcal{I}^\bullet$ une résolution injective.
Alors $R^if_*\mathcal{F}$ est par définition le $i$-ième faisceau de cohomologie
du complexe
$$
f_*\mathcal{I}^0 \to f_*\mathcal{I}^1 \to f_*\mathcal{I}^2 \to \ldots
$$
Par définition de la structure de catégorie abélienne sur les
$\mathcal{O}_\mathcal{D}$-modules,
ce faisceau de cohomologie est le faisceau associé au préfaisceau
$$
V
\longmapsto
\frac{\Ker(f_*\mathcal{I}^i(V) \to f_*\mathcal{I}^{i + 1}(V))}
{\Im(f_*\mathcal{I}^{i - 1}(V) \to f_*\mathcal{I}^i(V))}
$$
et ceci est manifestement égal à
$$
\frac{\Ker(\mathcal{I}^i(u(V)) \to \mathcal{I}^{i + 1}(u(V)))}
{\Im(\mathcal{I}^{i - 1}(u(V)) \to \mathcal{I}^i(u(V)))}
$$
qui est égal à $H^i(u(V), \mathcal{F})$,
ce qui conclut.
\end{proof}






\section{Le complexe de {\v C}ech et la cohomologie de {\v C}ech}
\label{sites-cohomology-section-cech}

\noindent
Soit $\mathcal{C}$ une catégorie. Soit $\mathcal{U} = \{U_i \to U\}_{i \in I}$
une famille de morphismes de but fixé, voir
Sites, définition \ref{sites-definition-family-morphisms-fixed-target}.
Supposons que tous les produits fibrés $U_{i_0} \times_U \ldots \times_U U_{i_p}$
existent dans $\mathcal{C}$. Soit $\mathcal{F}$ un préfaisceau abélien sur
$\mathcal{C}$. Posons
$$
\check{\mathcal{C}}^p(\mathcal{U}, \mathcal{F})
=
\prod\nolimits_{(i_0, \ldots, i_p) \in I^{p + 1}}
\mathcal{F}(U_{i_0} \times_U \ldots \times_U U_{i_p}).
$$
C'est un groupe abélien. Pour
$s \in \check{\mathcal{C}}^p(\mathcal{U}, \mathcal{F})$, nous notons
$s_{i_0\ldots i_p}$ sa composante dans le facteur
$\mathcal{F}(U_{i_0} \times_U \ldots \times_U U_{i_p})$.
Nous définissons
$$
d : \check{\mathcal{C}}^p(\mathcal{U}, \mathcal{F})
\longrightarrow
\check{\mathcal{C}}^{p + 1}(\mathcal{U}, \mathcal{F})
$$
par la formule
\begin{equation}
\label{sites-cohomology-equation-d-cech}
d(s)_{i_0\ldots i_{p + 1}} =
\sum\nolimits_{j = 0}^{p + 1}
(-1)^j s_{i_0\ldots \hat i_j \ldots i_{p + 1}}
|_{U_{i_0} \times_U \ldots \times_U U_{i_{p + 1}}}
\end{equation}
où la restriction est induite par la projection
$$
U_{i_0} \times_U \ldots \times_U U_{i_{p + 1}} \longrightarrow
U_{i_0} \times_U \ldots \times_U \widehat{U_{i_j}} \times_U
\ldots \times_U U_{i_{p + 1}}.
$$
On vérifie immédiatement que $d \circ d = 0$. Autrement dit,
$\check{\mathcal{C}}^\bullet(\mathcal{U}, \mathcal{F})$ est un complexe.

\begin{definition}
\label{sites-cohomology-definition-cech-complex}
Soit $\mathcal{C}$ une catégorie. Soit $\mathcal{U} = \{U_i \to U\}_{i \in I}$
une famille de morphismes de but fixé telle que tous les produits fibrés
$U_{i_0} \times_U \ldots \times_U U_{i_p}$ existent dans $\mathcal{C}$.
Soit $\mathcal{F}$ un préfaisceau abélien sur $\mathcal{C}$.
Le complexe $\check{\mathcal{C}}^\bullet(\mathcal{U}, \mathcal{F})$
est le {\it complexe de {\v C}ech} associé à $\mathcal{F}$ et à la
famille $\mathcal{U}$. Ses groupes de cohomologie
$H^i(\check{\mathcal{C}}^\bullet(\mathcal{U}, \mathcal{F}))$ sont
appelés les {\it groupes de cohomologie de {\v C}ech} de $\mathcal{F}$ relativement
à $\mathcal{U}$. Ils sont notés $\check H^i(\mathcal{U}, \mathcal{F})$.
\end{definition}

\noindent
Observons que tout recouvrement $\{U_i \to U\}$ d'un site $\mathcal{C}$
est une famille de morphismes de but fixé à laquelle la définition s'applique.

\begin{lemma}
\label{sites-cohomology-lemma-cech-h0}
Soit $\mathcal{C}$ un site.
Soit $\mathcal{F}$ un préfaisceau abélien sur $\mathcal{C}$.
Les assertions suivantes sont équivalentes :
\begin{enumerate}
\item $\mathcal{F}$ est un faisceau abélien sur $\mathcal{C}$ ;
\item pour tout recouvrement $\mathcal{U} = \{U_i \to U\}_{i \in I}$
du site $\mathcal{C}$, l'application naturelle
$$
\mathcal{F}(U) \to \check{H}^0(\mathcal{U}, \mathcal{F})
$$
(voir Sites, section \ref{sites-section-sheafification}) est bijective.
\end{enumerate}
\end{lemma}

\begin{proof}
Cela résulte du fait que la condition de faisceau exige précisément que
$\mathcal{F}(U) \to \check{H}^0(\mathcal{U}, \mathcal{F})$
soit bijective pour tout recouvrement de $\mathcal{C}$.
\end{proof}

\noindent
Soit $\mathcal{C}$ une catégorie. Soit $\mathcal{U} = \{U_i \to U\}_{i\in I}$
une famille de morphismes de $\mathcal{C}$ de but fixé telle que
tous les produits fibrés $U_{i_0} \times_U \ldots \times_U U_{i_p}$ existent dans
$\mathcal{C}$. Soit $\mathcal{V} = \{V_j \to V\}_{j\in J}$ une autre telle famille.
Les données $f : U \to V$, $\alpha : I \to J$ et $f_i : U_i \to V_{\alpha(i)}$
forment un morphisme de familles de morphismes de but fixé, voir
Sites, section \ref{sites-section-refinements}.
Dans ce cas, nous obtenons un morphisme de complexes de {\v C}ech
\begin{equation}
\label{sites-cohomology-equation-map-cech-complexes}
\varphi : \check{\mathcal{C}}^\bullet(\mathcal{V}, \mathcal{F})
\longrightarrow
\check{\mathcal{C}}^\bullet(\mathcal{U}, \mathcal{F})
\end{equation}
qui, en degré $p$, est donné par
$$
\varphi(s)_{i_0 \ldots i_p} =
(f_{i_0} \times \ldots \times f_{i_p})^*s_{\alpha(i_0) \ldots \alpha(i_p)}
$$


\section{La cohomologie de {\v C}ech comme foncteur sur les préfaisceaux}
\sectionmark{Cohomologie de {\v C}ech des préfaisceaux}
\label{sites-cohomology-section-cech-functor}

\noindent
Attention : dans cette section, nous travaillons exclusivement avec des préfaisceaux abéliens
sur une catégorie. Les résultats sont complètement faux dans le
cadre des faisceaux et des catégories de faisceaux !

\medskip\noindent
Soit $\mathcal{C}$ une catégorie. Soit $\mathcal{U} = \{U_i \to U\}_{i \in I}$
une famille de morphismes de but fixé telle que tous les produits fibrés
$U_{i_0} \times_U \ldots \times_U U_{i_p}$ existent dans $\mathcal{C}$.
Soit $\mathcal{F}$ un préfaisceau abélien sur $\mathcal{C}$.
La construction
$$
\mathcal{F} \longmapsto \check{\mathcal{C}}^\bullet(\mathcal{U}, \mathcal{F})
$$
est fonctorielle en $\mathcal{F}$. En fait, c'est un foncteur
\begin{equation}
\label{sites-cohomology-equation-cech-functor}
\check{\mathcal{C}}^\bullet(\mathcal{U}, -) :
\textit{PAb}(\mathcal{C})
\longrightarrow
\text{Comp}^{+}(\textit{Ab})
\end{equation}
voir
Catégories dérivées, définition \ref{derived-definition-complexes-notation},
pour les notations. Rappelons que la catégorie des complexes bornés inférieurement
dans une catégorie abélienne est une catégorie abélienne, voir
Homologie, lemme \ref{homology-lemma-cat-cochain-abelian}.

\begin{lemma}
\label{sites-cohomology-lemma-cech-exact-presheaves}
Le foncteur donné par l'équation (\ref{sites-cohomology-equation-cech-functor})
est exact (voir Homologie, lemme \ref{homology-lemma-exact-functor}).
\end{lemma}

\begin{proof}
Pour tout objet $W$ de $\mathcal{C}$, le foncteur
$\mathcal{F} \mapsto \mathcal{F}(W)$ est un foncteur additif exact
de $\textit{PAb}(\mathcal{C})$ dans $\textit{Ab}$.
Les foncteurs donnant les termes $\check{\mathcal{C}}^p(\mathcal{U}, \mathcal{F})$
du complexe sont des produits de ces foncteurs exacts et sont donc exacts.
De plus, une suite de complexes est exacte si et seulement si la suite
des termes en chaque degré est exacte. Le lemme en découle.
\end{proof}

\begin{lemma}
\label{sites-cohomology-lemma-cech-cohomology-delta-functor-presheaves}
Soit $\mathcal{C}$ une catégorie.
Soit $\mathcal{U} = \{U_i \to U\}_{i \in I}$ une famille de morphismes
de but fixé telle que tous les produits fibrés
$U_{i_0} \times_U \ldots \times_U U_{i_p}$ existent dans $\mathcal{C}$.
Les foncteurs $\mathcal{F} \mapsto \check{H}^n(\mathcal{U}, \mathcal{F})$
forment un $\delta$-foncteur de la catégorie abélienne $\textit{PAb}(\mathcal{C})$
vers la catégorie des $\mathbf{Z}$-modules (voir
Homologie, définition \ref{homology-definition-cohomological-delta-functor}).
\end{lemma}

\begin{proof}
D'après le
lemme \ref{sites-cohomology-lemma-cech-exact-presheaves},
une suite exacte courte de préfaisceaux abéliens
$0 \to \mathcal{F}_1 \to \mathcal{F}_2 \to \mathcal{F}_3 \to 0$
est transformée en une suite exacte courte de complexes de
$\mathbf{Z}$-modules. Nous pouvons donc utiliser
Homologie, lemme \ref{homology-lemma-long-exact-sequence-cochain},
pour obtenir les morphismes de connexion
$\delta_{\mathcal{F}_1 \to \mathcal{F}_2 \to \mathcal{F}_3} :
\check{H}^n(\mathcal{U}, \mathcal{F}_3) \to
\check{H}^{n + 1}(\mathcal{U}, \mathcal{F}_1)$
et une suite exacte longue correspondante. Nous omettons de vérifier
que ces morphismes sont compatibles avec les morphismes entre suites exactes
courtes de préfaisceaux.
\end{proof}

\begin{lemma}
\label{sites-cohomology-lemma-cech-map-into}
Soit $\mathcal{C}$ une catégorie. Soit $\mathcal{U} = \{U_i \to U\}_{i \in I}$
une famille de morphismes de but fixé telle que tous les produits fibrés
$U_{i_0} \times_U \ldots \times_U U_{i_p}$ existent dans $\mathcal{C}$.
Considérons le complexe de chaînes $\mathbf{Z}_{\mathcal{U}, \bullet}$
de préfaisceaux abéliens
$$
\ldots
\to
\bigoplus_{i_0i_1i_2} \mathbf{Z}_{U_{i_0} \times_U U_{i_1} \times_U U_{i_2}}
\to
\bigoplus_{i_0i_1} \mathbf{Z}_{U_{i_0} \times_U U_{i_1}}
\to
\bigoplus_{i_0} \mathbf{Z}_{U_{i_0}}
\to 0 \to \ldots
$$
où le dernier terme non nul est placé en degré $0$
et où le morphisme
$$
\mathbf{Z}_{U_{i_0} \times_U \ldots \times_U U_{i_{p + 1}}}
\longrightarrow
\mathbf{Z}_{U_{i_0} \times_U
\ldots \widehat{U_{i_j}} \ldots \times_U U_{i_{p + 1}}}
$$
est égal au morphisme canonique multiplié par $(-1)^j$.
Il existe alors un isomorphisme
$$
\Hom_{\textit{PAb}(\mathcal{C})}(\mathbf{Z}_{\mathcal{U}, \bullet}, \mathcal{F})
=
\check{\mathcal{C}}^\bullet(\mathcal{U}, \mathcal{F})
$$
fonctoriel en $\mathcal{F} \in \Ob(\textit{PAb}(\mathcal{C}))$.
\end{lemma}

\begin{proof}
C'est une tautologie qui repose sur le fait que
\begin{align*}
\Hom_{\textit{PAb}(\mathcal{C})}(
\bigoplus_{i_0 \ldots i_p}
\mathbf{Z}_{U_{i_0} \times_U \ldots \times_U U_{i_p}},
\mathcal{F})
& =
\prod_{i_0 \ldots i_p}
\Hom_{\textit{PAb}(\mathcal{C})}(
\mathbf{Z}_{U_{i_0} \times_U \ldots \times_U U_{i_p}},
\mathcal{F}) \\
& =
\prod_{i_0 \ldots i_p}
\mathcal{F}(U_{i_0} \times_U \ldots \times_U U_{i_p})
\end{align*}
voir Modules sur les sites, lemme \ref{sites-modules-lemma-obvious-adjointness}.
\end{proof}

\begin{lemma}
\label{sites-cohomology-lemma-homology-complex}
Soit $\mathcal{C}$ une catégorie. Soit
$\mathcal{U} = \{f_i : U_i \to U\}_{i \in I}$ une famille de morphismes
de but fixé telle que tous les produits fibrés
$U_{i_0} \times_U \ldots \times_U U_{i_p}$ existent dans $\mathcal{C}$.
Le complexe de chaînes de préfaisceaux $\mathbf{Z}_{\mathcal{U}, \bullet}$
du lemme \ref{sites-cohomology-lemma-cech-map-into} ci-dessus est exact en degrés
strictement positifs, c'est-à-dire que les préfaisceaux d'homologie
$H_i(\mathbf{Z}_{\mathcal{U}, \bullet})$ sont nuls pour $i > 0$.
\end{lemma}

\begin{proof}
Soit $V$ un objet de $\mathcal{C}$. Nous devons montrer que le complexe de chaînes
de groupes abéliens $\mathbf{Z}_{\mathcal{U}, \bullet}(V)$ est exact en
degrés $> 0$. Il s'agit du complexe
$$
\xymatrix{
\ldots \ar[d] \\
\bigoplus_{i_0i_1i_2}
\mathbf{Z}[
\Mor_\mathcal{C}(V, U_{i_0} \times_U U_{i_1} \times_U U_{i_2})
]
\ar[d] \\
\bigoplus_{i_0i_1}
\mathbf{Z}[
\Mor_\mathcal{C}(V, U_{i_0} \times_U U_{i_1})
]
\ar[d] \\
\bigoplus_{i_0}
\mathbf{Z}[
\Mor_\mathcal{C}(V, U_{i_0})
] \ar[d] \\
0
}
$$
Pour tout morphisme $\varphi : V \to U$, notons
$\Mor_\varphi(V, U_i) = \{\varphi_i : V \to U_i \mid
f_i \circ \varphi_i = \varphi\}$. Nous utiliserons une notation analogue
pour $\Mor_\varphi(V, U_{i_0} \times_U \ldots \times_U U_{i_p})$.
Remarquons que la composition avec les différentes projections entre les
produits fibrés $U_{i_0} \times_U \ldots \times_U U_{i_p}$ préserve
ces ensembles de morphismes. Nous voyons donc que le complexe ci-dessus
est le même que le complexe
$$
\xymatrix{
\ldots \ar[d] \\
\bigoplus_\varphi
\bigoplus_{i_0i_1i_2}
\mathbf{Z}[
\Mor_\varphi(V, U_{i_0} \times_U U_{i_1} \times_U U_{i_2})
]
\ar[d] \\
\bigoplus_\varphi
\bigoplus_{i_0i_1}
\mathbf{Z}[
\Mor_\varphi(V, U_{i_0} \times_U U_{i_1})
]
\ar[d] \\
\bigoplus_\varphi
\bigoplus_{i_0}
\mathbf{Z}[
\Mor_\varphi(V, U_{i_0})
] \ar[d] \\
0
}
$$
Ensuite, remarquons que nous avons
$$
\Mor_\varphi(V, U_{i_0} \times_U \ldots \times_U U_{i_p})
=
\Mor_\varphi(V, U_{i_0}) \times \ldots
\times \Mor_\varphi(V, U_{i_p})
$$
En utilisant ceci et le fait que $\mathbf{Z}[A] \oplus \mathbf{Z}[B] =
\mathbf{Z}[A \amalg B]$, nous voyons que le complexe devient
$$
\xymatrix{
\ldots \ar[d] \\
\bigoplus_\varphi
\mathbf{Z}\left[
\coprod_{i_0i_1i_2}
\Mor_\varphi(V, U_{i_0}) \times \Mor_\varphi(V, U_{i_1}) \times
\Mor_\varphi(V, U_{i_2})
\right]
\ar[d] \\
\bigoplus_\varphi
\mathbf{Z}\left[
\coprod_{i_0i_1}
\Mor_\varphi(V, U_{i_0}) \times \Mor_\varphi(V, U_{i_1})
\right]
\ar[d] \\
\bigoplus_\varphi
\mathbf{Z}\left[
\coprod_{i_0}
\Mor_\varphi(V, U_{i_0})
\right] \ar[d] \\
0
}
$$
Enfin, en posant $S_\varphi = \coprod_{i \in I} \Mor_\varphi(V, U_i)$,
nous obtenons
$$
\bigoplus\nolimits_\varphi \left(\ldots \to
\mathbf{Z}[S_\varphi \times S_\varphi \times S_\varphi] \to
\mathbf{Z}[S_\varphi \times S_\varphi] \to
\mathbf{Z}[S_\varphi] \to 0 \to \ldots
\right)
$$
Nous avons ainsi simplifié notre tâche. Il suffit en effet de montrer que,
pour tout ensemble non vide $S$, le complexe (augmenté) de groupes abéliens libres
$$
\ldots \to
\mathbf{Z}[S \times S \times S] \to
\mathbf{Z}[S \times S] \to
\mathbf{Z}[S] \xrightarrow{\Sigma} \mathbf{Z} \to 0 \to \ldots
$$
est exact en tout degré. Pour le voir, fixons un élément $s \in S$ et
utilisons l'homotopie
$$
n_{(s_0, \ldots, s_p)} \longmapsto n_{(s, s_0, \ldots, s_p)}
$$
avec les notations évidentes.
\end{proof}

\begin{lemma}
\label{sites-cohomology-lemma-complex-tensored-still-exact}
\begin{slogan}
Le complexe de préfaisceaux de {\v C}ech à coefficients entiers est une résolution plate du
préfaisceau constant des entiers.
\end{slogan}
Soit $\mathcal{C}$ une catégorie. Soit
$\mathcal{U} = \{f_i : U_i \to U\}_{i \in I}$ une famille de morphismes
de but fixé telle que tous les produits fibrés
$U_{i_0} \times_U \ldots \times_U U_{i_p}$ existent dans $\mathcal{C}$.
Soit $\mathcal{O}$ un préfaisceau d'anneaux sur $\mathcal{C}$.
Le complexe de chaînes
$$
\mathbf{Z}_{\mathcal{U}, \bullet}
\otimes_{p, \mathbf{Z}}
\mathcal{O}
$$
est exact en degrés strictement positifs. Ici, $\mathbf{Z}_{\mathcal{U}, \bullet}$
est le complexe de chaînes du lemme \ref{sites-cohomology-lemma-cech-map-into}, et
le produit tensoriel est pris sur le préfaisceau constant d'anneaux
de valeur $\mathbf{Z}$.
\end{lemma}

\begin{proof}
Soit $V$ un objet de $\mathcal{C}$.
Dans la démonstration du lemme \ref{sites-cohomology-lemma-homology-complex}, nous avons vu que
$\mathbf{Z}_{\mathcal{U}, \bullet}(V)$ est isomorphe, comme complexe,
à une somme directe de complexes homotopes à $\mathbf{Z}$
placé en degré zéro. Par conséquent,
$\mathbf{Z}_{\mathcal{U}, \bullet}(V) \otimes_\mathbf{Z} \mathcal{O}(V)$
est également isomorphe, comme complexe, à une somme directe de complexes homotopes
à $\mathcal{O}(V)$ placé en degré zéro.
On peut aussi utiliser
Modules sur les sites, lemme \ref{sites-modules-lemma-flat-resolution-of-flat},
qui s'applique puisque les préfaisceaux $\mathbf{Z}_{\mathcal{U}, i}$ sont plats,
et la démonstration du lemme \ref{sites-cohomology-lemma-homology-complex} montre que
$H_0(\mathbf{Z}_{\mathcal{U}, \bullet})$ est lui aussi un préfaisceau plat.
\end{proof}

\begin{lemma}
\label{sites-cohomology-lemma-cech-cohomology-derived-presheaves}
Soit $\mathcal{C}$ une catégorie. Soit
$\mathcal{U} = \{f_i : U_i \to U\}_{i \in I}$ une famille de morphismes
de but fixé telle que tous les produits fibrés
$U_{i_0} \times_U \ldots \times_U U_{i_p}$ existent dans $\mathcal{C}$.
Les foncteurs de cohomologie de {\v C}ech $\check{H}^p(\mathcal{U}, -)$
sont canoniquement isomorphes, en tant que $\delta$-foncteur, aux
foncteurs dérivés à droite du foncteur
$$
\check{H}^0(\mathcal{U}, -) :
\textit{PAb}(\mathcal{C})
\longrightarrow
\textit{Ab}.
$$
De plus, il existe un quasi-isomorphisme fonctoriel
$$
\check{\mathcal{C}}^\bullet(\mathcal{U}, \mathcal{F})
\longrightarrow
R\check{H}^0(\mathcal{U}, \mathcal{F})
$$
où le membre de droite désigne le foncteur dérivé
$$
R\check{H}^0(\mathcal{U}, -) :
D^{+}(\textit{PAb}(\mathcal{C}))
\longrightarrow
D^{+}(\mathbf{Z})
$$
du foncteur exact à gauche $\check{H}^0(\mathcal{U}, -)$.
\end{lemma}

\begin{proof}
Remarquons que la catégorie des préfaisceaux abéliens possède suffisamment d'injectifs, voir
Injectifs, proposition \ref{injectives-proposition-presheaves-injectives}.
Remarquons que $\check{H}^0(\mathcal{U}, -)$ est un foncteur exact à gauche
de la catégorie des préfaisceaux abéliens
vers la catégorie des $\mathbf{Z}$-modules.
Le foncteur dérivé et le foncteur dérivé à droite existent donc, voir
Catégories dérivées, section \ref{derived-section-right-derived-functor}.

\medskip\noindent
Soit $\mathcal{I}$ un préfaisceau abélien injectif.
Dans ce cas, le foncteur
$\Hom_{\textit{PAb}(\mathcal{C})}(-, \mathcal{I})$
est exact sur $\textit{PAb}(\mathcal{C})$. D'après le
lemme \ref{sites-cohomology-lemma-cech-map-into}, nous avons
$$
\Hom_{\textit{PAb}(\mathcal{C})}(
\mathbf{Z}_{\mathcal{U}, \bullet}, \mathcal{I})
=
\check{\mathcal{C}}^\bullet(\mathcal{U}, \mathcal{I}).
$$
D'après le lemme \ref{sites-cohomology-lemma-homology-complex}, le complexe
$\mathbf{Z}_{\mathcal{U}, \bullet}$ est exact en degrés strictement positifs.
L'exactitude du foncteur Hom à valeurs dans $\mathcal{I}$ mentionnée ci-dessus montre donc
que $\check{H}^i(\mathcal{U}, \mathcal{I}) = 0$ pour tout
$i > 0$. Ainsi, le $\delta$-foncteur $(\check{H}^n, \delta)$
(voir le lemme \ref{sites-cohomology-lemma-cech-cohomology-delta-functor-presheaves})
vérifie les hypothèses de
Homologie, lemme \ref{homology-lemma-efface-implies-universal},
et est donc un $\delta$-foncteur universel.

\medskip\noindent
D'après
Catégories dérivées, lemme \ref{derived-lemma-higher-derived-functors},
la famille $R^i\check{H}^0(\mathcal{U}, -)$
forme elle aussi un $\delta$-foncteur universel. Par unicité des
$\delta$-foncteurs universels, voir
Homologie, lemme \ref{homology-lemma-uniqueness-universal-delta-functor},
nous concluons que
$R^i\check{H}^0(\mathcal{U}, -) = \check{H}^i(\mathcal{U}, -)$.
Cela suffit pour la plupart des applications,
et nous suggérons au lecteur de passer la suite de la démonstration.

\medskip\noindent
Soit $\mathcal{F}$ un préfaisceau abélien quelconque sur $\mathcal{C}$.
Choisissons une résolution injective $\mathcal{F} \to \mathcal{I}^\bullet$
dans la catégorie $\textit{PAb}(\mathcal{C})$.
Considérons le complexe double
$\check{\mathcal{C}}^\bullet(\mathcal{U}, \mathcal{I}^\bullet)$
de termes $\check{\mathcal{C}}^p(\mathcal{U}, \mathcal{I}^q)$.
Considérons ensuite le complexe total
$\text{Tot}(\check{\mathcal{C}}^\bullet(\mathcal{U}, \mathcal{I}^\bullet))$
associé à ce complexe double, voir
Homologie, section \ref{homology-section-double-complexes}.
Il existe un morphisme de complexes
$$
\check{\mathcal{C}}^\bullet(\mathcal{U}, \mathcal{F})
\longrightarrow
\text{Tot}(\check{\mathcal{C}}^\bullet(\mathcal{U}, \mathcal{I}^\bullet))
$$
provenant des morphismes
$\check{\mathcal{C}}^p(\mathcal{U}, \mathcal{F})
\to \check{\mathcal{C}}^p(\mathcal{U}, \mathcal{I}^0)$
et il existe un morphisme de complexes
$$
\check{H}^0(\mathcal{U}, \mathcal{I}^\bullet)
\longrightarrow
\text{Tot}(\check{\mathcal{C}}^\bullet(\mathcal{U}, \mathcal{I}^\bullet))
$$
provenant des morphismes
$\check{H}^0(\mathcal{U}, \mathcal{I}^q) \to
\check{\mathcal{C}}^0(\mathcal{U}, \mathcal{I}^q)$.
Ces deux morphismes sont des quasi-isomorphismes par application de
Homologie, lemme \ref{homology-lemma-double-complex-gives-resolution}.
En effet, les colonnes du complexe double sont exactes en degrés strictement positifs
parce que le complexe de {\v C}ech est exact en tant que foncteur
(lemme \ref{sites-cohomology-lemma-cech-exact-presheaves}),
et les lignes du complexe double sont exactes en degrés strictement positifs
puisque, comme nous venons de le voir, les groupes de cohomologie supérieure de {\v C}ech des
préfaisceaux injectifs $\mathcal{I}^q$ sont nuls.
Puisque les quasi-isomorphismes deviennent inversibles
dans $D^{+}(\mathbf{Z})$, ceci donne le dernier morphisme affiché
du lemme. Nous omettons de vérifier que ce morphisme est
fonctoriel.
\end{proof}





\section{Cohomologie de {\v C}ech et cohomologie}
\label{sites-cohomology-section-cech-cohomology-cohomology}

\noindent
La relation entre cohomologie et cohomologie de {\v C}ech vient de ce que la
cohomologie de {\v C}ech d'un faisceau abélien injectif est nulle. Pour le voir,
on remarque qu'un faisceau abélien injectif est un préfaisceau abélien injectif,
puis on applique les résultats de cohomologie de {\v C}ech établis dans la
section précédente.

\begin{lemma}
\label{sites-cohomology-lemma-injective-abelian-sheaf-injective-presheaf}
Soit $\mathcal{C}$ un site. Tout faisceau abélien injectif est également
injectif en tant qu'objet de la catégorie $\textit{PAb}(\mathcal{C})$.
\end{lemma}

\begin{proof}
Appliquons le lemme \ref{homology-lemma-adjoint-preserve-injectives}
d'Homologie aux catégories $\mathcal{A} = \textit{Ab}(\mathcal{C})$ et
$\mathcal{B} = \textit{PAb}(\mathcal{C})$, au foncteur d'inclusion et à la
faisceautisation. (Voir Modules sur les sites, section
\ref{sites-modules-section-abelian-sheaves}, pour constater que toutes les
hypothèses du lemme sont satisfaites.)
\end{proof}

\begin{lemma}
\label{sites-cohomology-lemma-injective-trivial-cech}
Soit $\mathcal{C}$ un site.
Soit $\mathcal{U} = \{U_i \to U\}_{i \in I}$ un recouvrement de
$\mathcal{C}$. Soit $\mathcal{I}$ un faisceau abélien injectif, c'est-à-dire
un objet injectif de $\textit{Ab}(\mathcal{C})$. Alors
$$
\check{H}^p(\mathcal{U}, \mathcal{I}) =
\left\{
\begin{matrix}
\mathcal{I}(U) & \text{if} & p = 0 \\
0 & \text{if} & p > 0
\end{matrix}
\right.
$$
\end{lemma}

\begin{proof}
D'après le lemme \ref{sites-cohomology-lemma-injective-abelian-sheaf-injective-presheaf},
$\mathcal{I}$ est un objet injectif de $\textit{PAb}(\mathcal{C})$.
Nous pouvons donc appliquer le
lemme \ref{sites-cohomology-lemma-cech-cohomology-derived-presheaves} (ou sa preuve) pour
obtenir l'annulation des groupes de cohomologie de {\v C}ech supérieurs.
Pour le degré zéro, voir le lemme \ref{sites-cohomology-lemma-cech-h0}.
\end{proof}

\begin{lemma}
\label{sites-cohomology-lemma-cech-cohomology}
Soit $\mathcal{C}$ un site.
Soit $\mathcal{U} = \{U_i \to U\}_{i \in I}$ un recouvrement de
$\mathcal{C}$. Il existe une transformation
$$
\check{\mathcal{C}}^\bullet(\mathcal{U}, -)
\longrightarrow
R\Gamma(U, -)
$$
de foncteurs $\textit{Ab}(\mathcal{C}) \to D^{+}(\mathbf{Z})$.
En particulier, elle fournit une transformation de foncteurs
$\check{H}^p(\mathcal{U}, \mathcal{F}) \to H^p(U, \mathcal{F})$, lorsque
$\mathcal{F}$ parcourt $\textit{Ab}(\mathcal{C})$.
\end{lemma}

\begin{proof}
Soit $\mathcal{F}$ un faisceau abélien. Choisissons une résolution injective
$\mathcal{F} \to \mathcal{I}^\bullet$. Considérons le bicomplexe
$\check{\mathcal{C}}^\bullet(\mathcal{U}, \mathcal{I}^\bullet)$
dont les termes sont $\check{\mathcal{C}}^p(\mathcal{U}, \mathcal{I}^q)$.
Considérons ensuite le complexe total associé
$\text{Tot}(\check{\mathcal{C}}^\bullet(\mathcal{U}, \mathcal{I}^\bullet))$,
voir Homologie, définition
\ref{homology-definition-associated-simple-complex}. Il existe un morphisme
de complexes
$$
\alpha :
\Gamma(U, \mathcal{I}^\bullet)
\longrightarrow
\text{Tot}(\check{\mathcal{C}}^\bullet(\mathcal{U}, \mathcal{I}^\bullet))
$$
provenant des morphismes
$\mathcal{I}^q(U) \to \check{H}^0(\mathcal{U}, \mathcal{I}^q)$
et un morphisme de complexes
$$
\beta :
\check{\mathcal{C}}^\bullet(\mathcal{U}, \mathcal{F})
\longrightarrow
\text{Tot}(\check{\mathcal{C}}^\bullet(\mathcal{U}, \mathcal{I}^\bullet))
$$
provenant du morphisme $\mathcal{F} \to \mathcal{I}^0$.
Le lemme \ref{homology-lemma-double-complex-gives-resolution} d'Homologie
montre que $\alpha$ est un quasi-isomorphisme.
En effet, le lemme \ref{sites-cohomology-lemma-injective-trivial-cech} implique que la
$q$-ième ligne du bicomplexe
$\check{\mathcal{C}}^\bullet(\mathcal{U}, \mathcal{I}^\bullet)$ est une
résolution de $\Gamma(U, \mathcal{I}^q)$.
Ainsi, $\alpha$ devient inversible dans $D^{+}(\mathbf{Z})$, et la
transformation du lemme est la composée de $\beta$ avec l'inverse de
$\alpha$. Nous omettons la vérification de la fonctorialité.
\end{proof}

\begin{lemma}
\label{sites-cohomology-lemma-cech-h1}
Soit $\mathcal{C}$ un site. Soit $\mathcal{G}$ un faisceau abélien sur
$\mathcal{C}$. Soit $\mathcal{U} = \{U_i \to U\}_{i \in I}$ un recouvrement
de $\mathcal{C}$. Le morphisme
$$
\check{H}^1(\mathcal{U}, \mathcal{G})
\longrightarrow
H^1(U, \mathcal{G})
$$
est injectif et, par la bijection du lemme \ref{sites-cohomology-lemma-torsors-h1}, identifie
$\check{H}^1(\mathcal{U}, \mathcal{G})$ à l'ensemble des classes
d'isomorphisme de $\mathcal{G}|_U$-torseurs dont les restrictions à chaque
$U_i$ sont des torseurs triviaux.
\end{lemma}

\begin{proof}
Pour le voir, construisons un morphisme inverse. Soit $\mathcal{F}$ un
$\mathcal{G}|_U$-torseur sur $\mathcal{C}/U$ dont la restriction à
$\mathcal{C}/U_i$ est triviale. D'après le lemme
\ref{sites-cohomology-lemma-trivial-torsor}, cela signifie qu'il existe une section
$s_i \in \mathcal{F}(U_i)$. Sur $U_{i_0} \times_U U_{i_1}$, il existe une
unique section $s_{i_0i_1}$ de $\mathcal{G}$ telle que
$s_{i_0i_1} \cdot s_{i_0}|_{U_{i_0} \times_U U_{i_1}} =
s_{i_1}|_{U_{i_0} \times_U U_{i_1}}$. Un calcul immédiat montre que
$s_{i_0i_1}$ est un cocycle de {\v C}ech et que sa classe est bien définie
(c'est-à-dire indépendante du choix des sections $s_i$).
L'inverse envoie la classe d'isomorphisme de $\mathcal{F}$ sur la classe de
cohomologie du cocycle $(s_{i_0i_1})$. Nous omettons de vérifier que ce
morphisme est bien un inverse.
\end{proof}

\begin{lemma}
\label{sites-cohomology-lemma-include}
Soit $\mathcal{C}$ un site.
Considérons le foncteur
$i : \textit{Ab}(\mathcal{C}) \to \textit{PAb}(\mathcal{C})$.
C'est un foncteur exact à gauche dont les foncteurs dérivés à droite sont
donnés par
$$
R^pi(\mathcal{F}) = \underline{H}^p(\mathcal{F}) :
U \longmapsto H^p(U, \mathcal{F})
$$
voir la discussion de la section \ref{sites-cohomology-section-locality}.
\end{lemma}

\begin{proof}
Il est clair que $i$ est exact à gauche.
Choisissons une résolution injective $\mathcal{F} \to \mathcal{I}^\bullet$.
Par définition, $R^pi$ est le {\it préfaisceau} de cohomologie de degré $p$
du complexe $\mathcal{I}^\bullet$. Autrement dit, les sections de
$R^pi(\mathcal{F})$ sur un objet $U$ de $\mathcal{C}$ sont données par
$$
\frac{\Ker(\mathcal{I}^n(U) \to \mathcal{I}^{n + 1}(U))}
{\Im(\mathcal{I}^{n - 1}(U) \to \mathcal{I}^n(U))}.
$$
ce qui est la définition de $H^p(U, \mathcal{F})$.
\end{proof}

\begin{lemma}
\label{sites-cohomology-lemma-cech-spectral-sequence}
Soit $\mathcal{C}$ un site. Soit $\mathcal{U} = \{U_i \to U\}_{i \in I}$
un recouvrement de $\mathcal{C}$. Pour tout faisceau abélien $\mathcal{F}$,
il existe une suite spectrale $(E_r, d_r)_{r \geq 0}$ telle que
$$
E_2^{p, q} = \check{H}^p(\mathcal{U}, \underline{H}^q(\mathcal{F}))
$$
qui converge vers $H^{p + q}(U, \mathcal{F})$.
Cette suite spectrale est fonctorielle en $\mathcal{F}$.
\end{lemma}

\begin{proof}
C'est la suite spectrale de Grothendieck (voir Catégories dérivées,
lemme \ref{derived-lemma-grothendieck-spectral-sequence}) associée aux
foncteurs
$$
i :  \textit{Ab}(\mathcal{C}) \to \textit{PAb}(\mathcal{C})
\quad\text{et}\quad
\check{H}^0(\mathcal{U}, - ) : \textit{PAb}(\mathcal{C})
\to \textit{Ab}.
$$
En effet, $\check{H}^0(\mathcal{U}, i(\mathcal{F})) = \mathcal{F}(U)$ d'après
le lemme \ref{sites-cohomology-lemma-cech-h0}. De plus, $i(\mathcal{I})$ est acyclique pour la
cohomologie de {\v C}ech d'après le
lemme \ref{sites-cohomology-lemma-injective-trivial-cech}. Enfin,
$\check{H}^p(\mathcal{U}, -) = R^p\check{H}^0(\mathcal{U}, -)$ comme foncteurs
sur $\textit{PAb}(\mathcal{C})$, d'après le
lemme \ref{sites-cohomology-lemma-cech-cohomology-derived-presheaves}.
La réunion de ces observations prouve le lemme.
\end{proof}

\begin{lemma}
\label{sites-cohomology-lemma-cech-spectral-sequence-application}
Soit $\mathcal{C}$ un site.
Soit $\mathcal{U} = \{U_i \to U\}_{i \in I}$ un recouvrement.
Soit $\mathcal{F} \in \Ob(\textit{Ab}(\mathcal{C}))$.
Supposons que $H^i(U_{i_0} \times_U \ldots \times_U U_{i_p}, \mathcal{F}) = 0$
pour tous $i > 0$, $p \geq 0$ et $i_0, \ldots, i_p \in I$.
Alors $\check{H}^p(\mathcal{U}, \mathcal{F}) = H^p(U, \mathcal{F})$.
\end{lemma}

\begin{proof}
Nous utilisons la suite spectrale du
lemme \ref{sites-cohomology-lemma-cech-spectral-sequence}.
Les hypothèses signifient que $E_2^{p, q} = 0$ pour tout $(p, q)$ tel que
$q \not = 0$. La suite spectrale dégénère donc en $E_2$, d'où le résultat.
\end{proof}

\begin{lemma}
\label{sites-cohomology-lemma-ses-cech-h1}
Soit $\mathcal{C}$ un site.
Soit
$$
0 \to \mathcal{F} \to \mathcal{G} \to \mathcal{H} \to 0
$$
une suite exacte courte de faisceaux abéliens sur $\mathcal{C}$.
Soit $U$ un objet de $\mathcal{C}$. S'il existe un système cofinal de
recouvrements $\mathcal{U}$ de $U$ tels que
$\check{H}^1(\mathcal{U}, \mathcal{F}) = 0$,
alors le morphisme $\mathcal{G}(U) \to \mathcal{H}(U)$ est surjectif.
\end{lemma}

\begin{proof}
Prenons un élément $s \in \mathcal{H}(U)$. Choisissons un recouvrement
$\mathcal{U} = \{U_i \to U\}_{i \in I}$ tel que
(a) $\check{H}^1(\mathcal{U}, \mathcal{F}) = 0$ et (b) $s|_{U_i}$ soit
l'image d'une section $s_i \in \mathcal{G}(U_i)$.
Comme nous pouvons certainement trouver un recouvrement vérifiant (b), les
hypothèses du lemme permettent d'en trouver un qui vérifie à la fois (a) et
(b). Considérons les sections
$$
s_{i_0i_1} =
s_{i_1}|_{U_{i_0} \times_U U_{i_1}} - s_{i_0}|_{U_{i_0} \times_U U_{i_1}}.
$$
Puisque $s_i$ relève $s$, nous avons
$s_{i_0i_1} \in \mathcal{F}(U_{i_0} \times_U U_{i_1})$.
L'annulation de $\check{H}^1(\mathcal{U}, \mathcal{F})$ fournit des sections
$t_i \in \mathcal{F}(U_i)$ telles que
$$
s_{i_0i_1} =
t_{i_1}|_{U_{i_0} \times_U U_{i_1}} - t_{i_0}|_{U_{i_0} \times_U U_{i_1}}.
$$
Les sections $s_i - t_i$ satisfont alors manifestement à la condition de
faisceau et se recollent en une section de $\mathcal{G}$ sur $U$ dont l'image
est $s$. Ceci conclut la preuve.
\end{proof}

\begin{lemma}
\label{sites-cohomology-lemma-cech-vanish-collection}
(Variante de Cohomologie, lemme \ref{cohomology-lemma-cech-vanish}.)
Soit $\mathcal{C}$ un site. Soit $\text{Cov}_\mathcal{C}$ l'ensemble des
recouvrements de $\mathcal{C}$ (voir Sites, définition
\ref{sites-definition-site}). Soient $\mathcal{B} \subset \Ob(\mathcal{C})$
et $\text{Cov} \subset \text{Cov}_\mathcal{C}$ des sous-ensembles. Soit
$\mathcal{F}$ un faisceau abélien sur $\mathcal{C}$. Supposons que
\begin{enumerate}
\item pour tout $\mathcal{U} \in \text{Cov}$,
$\mathcal{U} = \{U_i \to U\}_{i \in I}$, on ait
$U, U_i \in \mathcal{B}$ et tout
$U_{i_0} \times_U \ldots \times_U U_{i_p} \in \mathcal{B}$.
\item pour tout $U \in \mathcal{B}$, les recouvrements de $U$ appartenant à
$\text{Cov}$ forment un système cofinal de recouvrements de $U$ ;
\item pour tout $\mathcal{U} \in \text{Cov}$, on ait
$\check{H}^p(\mathcal{U}, \mathcal{F}) = 0$ pour tout $p > 0$.
\end{enumerate}
Alors $H^p(U, \mathcal{F}) = 0$ pour tout $p > 0$ et tout
$U \in \mathcal{B}$.
\end{lemma}

\begin{proof}
Soient $\mathcal{F}$ et $\text{Cov}$ comme dans le lemme. Nous exprimerons
cette situation en disant que « la cohomologie de {\v C}ech supérieure de
$\mathcal{F}$ s'annule pour tout $\mathcal{U} \in \text{Cov}$ ».
Choisissons un plongement $\mathcal{F} \to \mathcal{I}$ dans un faisceau
abélien injectif. D'après le lemme \ref{sites-cohomology-lemma-injective-trivial-cech}, la
cohomologie de {\v C}ech supérieure de $\mathcal{I}$ s'annule pour tout
$\mathcal{U} \in \text{Cov}$. Posons $\mathcal{Q} = \mathcal{I}/\mathcal{F}$,
de sorte que nous ayons une suite exacte courte
$$
0 \to \mathcal{F} \to \mathcal{I} \to \mathcal{Q} \to 0.
$$
D'après le lemme \ref{sites-cohomology-lemma-ses-cech-h1} et notre hypothèse (2), cette suite
donne naissance à une suite exacte
$$
0 \to \mathcal{F}(U) \to \mathcal{I}(U) \to \mathcal{Q}(U) \to 0.
$$
pour tout $U \in \mathcal{B}$. Par conséquent, pour tout
$\mathcal{U} \in \text{Cov}$, nous obtenons une suite exacte courte de
complexes de {\v C}ech
$$
0 \to
\check{\mathcal{C}}^\bullet(\mathcal{U}, \mathcal{F}) \to
\check{\mathcal{C}}^\bullet(\mathcal{U}, \mathcal{I}) \to
\check{\mathcal{C}}^\bullet(\mathcal{U}, \mathcal{Q}) \to 0
$$
puisque, d'après l'hypothèse (1), chaque terme du complexe de {\v C}ech est
formé d'un produit de valeurs sur des éléments de $\mathcal{B}$.
En particulier, pour tout recouvrement $\mathcal{U} \in \text{Cov}$, nous
avons une suite exacte longue de groupes de cohomologie de {\v C}ech.
Il s'ensuit que $\mathcal{Q}$ est lui aussi un faisceau abélien dont la
cohomologie de {\v C}ech supérieure s'annule pour tout
$\mathcal{U} \in \text{Cov}$.

\medskip\noindent
Considérons ensuite la suite exacte longue de cohomologie
$$
\xymatrix{
0 \ar[r] &
H^0(U, \mathcal{F}) \ar[r] &
H^0(U, \mathcal{I}) \ar[r] &
H^0(U, \mathcal{Q}) \ar[lld] \\
&
H^1(U, \mathcal{F}) \ar[r] &
H^1(U, \mathcal{I}) \ar[r] &
H^1(U, \mathcal{Q}) \ar[lld] \\
&
\ldots & \ldots & \ldots \\
}
$$
pour tout $U \in \mathcal{B}$. Puisque $\mathcal{I}$ est injectif, nous avons
$H^n(U, \mathcal{I}) = 0$ pour $n > 0$ (voir Catégories dérivées,
lemme \ref{derived-lemma-higher-derived-functors}).
D'après ce qui précède, $H^0(U, \mathcal{I}) \to H^0(U, \mathcal{Q})$ est
surjectif, et donc $H^1(U, \mathcal{F}) = 0$.
Comme $\mathcal{F}$ était un faisceau abélien arbitraire dont la cohomologie de
{\v C}ech supérieure s'annule pour tout $\mathcal{U} \in \text{Cov}$, nous
concluons aussi que $H^1(U, \mathcal{Q}) = 0$, puisque $\mathcal{Q}$ est un
autre faisceau de ce type (voir ci-dessus). La suite exacte longue implique à
son tour que $H^2(U, \mathcal{F}) = 0$, et ainsi de suite.
\end{proof}














\section{Deuxième cohomologie et gerbes}
\label{sites-cohomology-section-gerbes}

\noindent
Soit $p : \mathcal{S} \to \mathcal{C}$ une gerbe sur un site dont tous les
groupes d'automorphismes sont commutatifs. Dans cette situation, les premier
et deuxième groupes de cohomologie du faisceau des automorphismes (Champs,
lemme \ref{stacks-lemma-gerbe-abelian-auts}) régissent l'existence d'objets.

\medskip\noindent
Le lemme suivant deviendra obsolète lorsqu'une étude plus complète de cette
relation sera ajoutée.

\begin{lemma}
\label{sites-cohomology-lemma-existence}
Soit $\mathcal{C}$ un site. Soit $p : \mathcal{S} \to \mathcal{C}$ une gerbe
sur un site dont les faisceaux d'automorphismes sont abéliens.
Soit $\mathcal{G}$ le faisceau de groupes abéliens construit dans Champs,
lemme \ref{stacks-lemma-gerbe-abelian-auts}.
Soit $U$ un objet de $\mathcal{C}$ tel que
\begin{enumerate}
\item il existe un système cofinal de recouvrements $\{U_i \to U\}$ de $U$
dans $\mathcal{C}$ tel que $H^1(U_i, \mathcal{G}) = 0$ et
$H^1(U_i \times_U U_j, \mathcal{G}) = 0$
pour tous $i, j$ ; et
\item $H^2(U, \mathcal{G}) = 0$.
\end{enumerate}
Alors il existe un objet de $\mathcal{S}$ au-dessus de $U$.
\end{lemma}

\begin{proof}
D'après Champs, définition \ref{stacks-definition-gerbe}, il existe un
recouvrement $\mathcal{U} = \{U_i \to U\}$ et des $x_i$ dans $\mathcal{S}$
au-dessus de $U_i$. Posons $U_{ij} = U_i \times_U U_j$. En raffinant le
recouvrement conformément à (1), nous pouvons supposer que
$H^1(U_i, \mathcal{G}) = 0$ et $H^1(U_{ij}, \mathcal{G}) = 0$.
Considérons le faisceau
$$
\mathcal{F}_{ij} =
\mathit{Isom}(x_i|_{U_{ij}}, x_j|_{U_{ij}})
$$
sur $\mathcal{C}/U_{ij}$. Puisque
$\mathcal{G}|_{U_{ij}} = \mathit{Aut}(x_i|_{U_{ij}})$, il existe une action
$$
\mathcal{G}|_{U_{ij}} \times \mathcal{F}_{ij} \to \mathcal{F}_{ij}
$$
par précomposition. Il est clair que $\mathcal{F}_{ij}$ est un
pseudo-$\mathcal{G}|_{U_{ij}}$-torseur, et même un torseur, car deux objets
quelconques d'une gerbe sont localement isomorphes.
Par notre choix du recouvrement et le lemme \ref{sites-cohomology-lemma-torsors-h1}, ces
torseurs sont triviaux (et admettent donc des sections globales d'après le
lemme \ref{sites-cohomology-lemma-trivial-torsor}). Autrement dit, nous pouvons choisir des
isomorphismes
$$
\varphi_{ij} : x_i|_{U_{ij}} \longrightarrow x_j|_{U_{ij}}
$$
Pour obtenir un objet $x$ au-dessus de $U$, nous allons modifier notre choix
des $\varphi_{ij}$ afin d'obtenir une donnée de descente (nécessairement
effective puisque $p : \mathcal{S} \to \mathcal{C}$ est un champ).
L'obstruction à ce que ces données forment une donnée de descente est le
possible défaut de la condition de cocycle. Posons
$U_{ijk} = U_i \times_U U_j \times_U U_k$. Nous pouvons considérer
$$
g_{ijk} = \varphi_{ik}^{-1}|_{U_{ijk}} \circ \varphi_{jk}|_{U_{ijk}} \circ
\varphi_{ij}|_{U_{ijk}}
$$
qui est un automorphisme de $x_i$ au-dessus de $U_{ijk}$. Nous considérons
donc $g_{ijk}$ comme une section de $\mathcal{G}$ sur $U_{ijk}$.
Un calcul, que nous omettons, montre que $(g_{i_0i_1i_2})$ est un $2$-cocycle
du complexe de {\v C}ech ${\check C}^\bullet(\mathcal{U}, \mathcal{G})$ de
$\mathcal{G}$ relativement au recouvrement $\mathcal{U}$.
D'après la suite spectrale du lemme \ref{sites-cohomology-lemma-cech-spectral-sequence} et
puisque $H^1(U_i, \mathcal{G}) = 0$ pour tout $i$, le morphisme
${\check H}^2(\mathcal{U}, \mathcal{G}) \to H^2(U, \mathcal{G})$ est injectif.
Par conséquent, $(g_{i_0i_1i_2})$ est un cobord, puisque nous supposons que
$H^2(U, \mathcal{G}) = 0$. Il existe donc des sections
$g_{ij} \in \mathcal{G}(U_{ij})$ telles que
$g_{ik}^{-1}|_{U_{ijk}} g_{jk}|_{U_{ijk}}  g_{ij}|_{U_{ijk}} = g_{ijk}$
pour tous $i, j, k$.
Après avoir remplacé $\varphi_{ij}$ par $\varphi_{ij}g_{ij}^{-1}$, les
$\varphi_{ij}$ définissent une donnée de descente sur les objets $x_i$
au-dessus de $U_i$, ce qui achève la preuve.
\end{proof}












\section{Cohomologie des modules}
\label{sites-cohomology-section-cohomology-modules}

\noindent
Tout ce qui a été dit pour la cohomologie des faisceaux abéliens vaut pour la
cohomologie des modules, puisque les deux coïncident.

\begin{lemma}
\label{sites-cohomology-lemma-injective-module-injective-presheaf}
Soit $(\mathcal{C}, \mathcal{O})$ un site annelé.
Tout faisceau injectif de modules est également injectif en tant qu'objet de
la catégorie $\textit{PMod}(\mathcal{O})$.
\end{lemma}

\begin{proof}
Appliquons le lemme \ref{homology-lemma-adjoint-preserve-injectives}
d'Homologie aux catégories $\mathcal{A} = \textit{Mod}(\mathcal{O})$ et
$\mathcal{B} = \textit{PMod}(\mathcal{O})$, au foncteur d'inclusion et à la
faisceautisation. (Voir Modules sur les sites, section
\ref{sites-modules-section-sheafification-presheaves-modules}, pour constater
que toutes les hypothèses du lemme sont satisfaites.)
\end{proof}

\begin{lemma}
\label{sites-cohomology-lemma-include-modules}
Soit $(\mathcal{C}, \mathcal{O})$ un site annelé.
Considérons le foncteur
$i : \textit{Mod}(\mathcal{C}) \to \textit{PMod}(\mathcal{C})$.
C'est un foncteur exact à gauche dont les foncteurs dérivés à droite sont
donnés par
$$
R^pi(\mathcal{F}) = \underline{H}^p(\mathcal{F}) :
U \longmapsto H^p(U, \mathcal{F})
$$
voir la discussion de la section \ref{sites-cohomology-section-locality}.
\end{lemma}

\begin{proof}
Il est clair que $i$ est exact à gauche.
Choisissons une résolution injective $\mathcal{F} \to \mathcal{I}^\bullet$
dans $\textit{Mod}(\mathcal{O})$.
Par définition, $R^pi$ est le {\it préfaisceau} de cohomologie de degré $p$
du complexe $\mathcal{I}^\bullet$. Autrement dit, les sections de
$R^pi(\mathcal{F})$ sur un objet $U$ de $\mathcal{C}$ sont données par
$$
\frac{\Ker(\mathcal{I}^n(U) \to \mathcal{I}^{n + 1}(U))}
{\Im(\mathcal{I}^{n - 1}(U) \to \mathcal{I}^n(U))}.
$$
ce qui est la définition de $H^p(U, \mathcal{F})$.
\end{proof}

\begin{lemma}
\label{sites-cohomology-lemma-injective-module-trivial-cech}
Soit $(\mathcal{C}, \mathcal{O})$ un site annelé.
Soit $\mathcal{U} = \{U_i \to U\}_{i \in I}$ un recouvrement de
$\mathcal{C}$. Soit $\mathcal{I}$ un $\mathcal{O}$-module injectif,
c'est-à-dire un objet injectif de $\textit{Mod}(\mathcal{O})$. Alors
$$
\check{H}^p(\mathcal{U}, \mathcal{I}) =
\left\{
\begin{matrix}
\mathcal{I}(U) & \text{if} & p = 0 \\
0 & \text{if} & p > 0
\end{matrix}
\right.
$$
\end{lemma}

\begin{proof}
Le lemme \ref{sites-cohomology-lemma-cech-map-into} donne la première égalité de la chaîne
d'égalités suivante :
\begin{align*}
\check{\mathcal{C}}^\bullet(\mathcal{U}, \mathcal{I})
& =
\Mor_{\textit{PAb}(\mathcal{C})}(
\mathbf{Z}_{\mathcal{U}, \bullet}, \mathcal{I}) \\
& =
\Mor_{\textit{PMod}(\mathbf{Z})}(
\mathbf{Z}_{\mathcal{U}, \bullet}, \mathcal{I}) \\
& =
\Mor_{\textit{PMod}(\mathcal{O})}(
\mathbf{Z}_{\mathcal{U}, \bullet} \otimes_{p, \mathbf{Z}} \mathcal{O},
\mathcal{I})
\end{align*}
La troisième égalité découle de Modules sur les sites,
lemme \ref{sites-modules-lemma-adjointness-tensor-restrict-presheaves}.
D'après le lemme \ref{sites-cohomology-lemma-injective-module-injective-presheaf},
$\mathcal{I}$ est un objet injectif de $\textit{PMod}(\mathcal{O})$.
Par conséquent, $\Hom_{\textit{PMod}(\mathcal{O})}(-, \mathcal{I})$ est un
foncteur exact. Le lemme \ref{sites-cohomology-lemma-complex-tensored-still-exact} donne alors
l'annulation des groupes de cohomologie de {\v C}ech supérieurs.
Pour le degré zéro, voir le lemme \ref{sites-cohomology-lemma-cech-h0}.
\end{proof}

\begin{lemma}
\label{sites-cohomology-lemma-cohomology-modules-abelian-agree}
Soit $\mathcal{C}$ un site.
Soit $\mathcal{O}$ un faisceau d'anneaux sur $\mathcal{C}$.
Soit $\mathcal{F}$ un $\mathcal{O}$-module, et notons $\mathcal{F}_{ab}$ le
faisceau de groupes abéliens sous-jacent. Alors
$$
H^i(\mathcal{C}, \mathcal{F}_{ab})
=
H^i(\mathcal{C}, \mathcal{F})
$$
et, pour tout objet $U$ de $\mathcal{C}$, nous avons également
$$
H^i(U, \mathcal{F}_{ab})
=
H^i(U, \mathcal{F}).
$$
Ici, le membre de gauche est la cohomologie calculée dans
$\textit{Ab}(\mathcal{C})$, et le membre de droite celle calculée dans
$\textit{Mod}(\mathcal{O})$.
\end{lemma}

\begin{proof}
D'après Catégories dérivées,
lemme \ref{derived-lemma-higher-derived-functors}, le $\delta$-foncteur
$(\mathcal{F} \mapsto H^p(U, \mathcal{F}))_{p \geq 0}$ est universel. Le
foncteur
$\textit{Mod}(\mathcal{O}) \to \textit{Ab}(\mathcal{C})$,
$\mathcal{F} \mapsto \mathcal{F}_{ab}$ est exact. Ainsi,
$(\mathcal{F} \mapsto H^p(U, \mathcal{F}_{ab}))_{p \geq 0}$
est également un $\delta$-foncteur. Supposons que nous montrions que
$(\mathcal{F} \mapsto H^p(U, \mathcal{F}_{ab}))_{p \geq 0}$
est lui aussi universel. La deuxième assertion du lemme découlera alors de
l'unicité des $\delta$-foncteurs universels ; voir Homologie,
lemme \ref{homology-lemma-uniqueness-universal-delta-functor}.
Puisque $\textit{Mod}(\mathcal{O})$ possède assez d'injectifs, il suffit de
montrer que $H^i(U, \mathcal{I}_{ab}) = 0$ pour tout objet injectif
$\mathcal{I}$ de $\textit{Mod}(\mathcal{O})$ ; voir Homologie,
lemme \ref{homology-lemma-efface-implies-universal}.

\medskip\noindent
Soit $\mathcal{I}$ un objet injectif de $\textit{Mod}(\mathcal{O})$.
Appliquons le lemme \ref{sites-cohomology-lemma-cech-vanish-collection} avec
$\mathcal{F} = \mathcal{I}$, $\mathcal{B} = \mathcal{C}$ et
$\text{Cov} = \text{Cov}_\mathcal{C}$.
L'hypothèse (3) de ce lemme est satisfaite d'après le
lemme \ref{sites-cohomology-lemma-injective-module-trivial-cech}.
Ainsi, $H^i(U, \mathcal{I}_{ab}) = 0$ pour tout objet $U$ de $\mathcal{C}$.

\medskip\noindent
Si $\mathcal{C}$ possède un objet final, cela implique également la première
égalité. Sinon, d'après Sites, lemme \ref{sites-lemma-topos-good-site}, le
topos annelé $(\Sh(\mathcal{C}), \mathcal{O})$ est équivalent à un topos annelé
dont le site sous-jacent possède un objet final. Le lemme en découle.
\end{proof}

\begin{lemma}
\label{sites-cohomology-lemma-cohomology-products}
Soit $\mathcal{C}$ un site. Soit $I$ un ensemble. Pour $i \in I$, soit
$\mathcal{F}_i$ un faisceau abélien sur $\mathcal{C}$. Soit
$U \in \Ob(\mathcal{C})$. Le morphisme canonique
$$
H^p(U, \prod\nolimits_{i \in I} \mathcal{F}_i)
\longrightarrow
\prod\nolimits_{i \in I} H^p(U, \mathcal{F}_i)
$$
est un isomorphisme pour $p = 0$ et est injectif pour $p = 1$.
\end{lemma}

\begin{proof}
L'assertion pour $p = 0$ est vraie parce que le produit des faisceaux est égal
au produit des préfaisceaux sous-jacents ; voir Sites,
lemme \ref{sites-lemma-limit-sheaf}.
Traitons $p = 1$. Posons $\mathcal{F} = \prod \mathcal{F}_i$.
Soit $\xi \in H^1(U, \mathcal{F})$ dont l'image dans
$\prod H^1(U, \mathcal{F}_i)$ est nulle. Par localité de la cohomologie, voir
le lemme \ref{sites-cohomology-lemma-kill-cohomology-class-on-covering}, il existe un
recouvrement $\mathcal{U} = \{U_j \to U\}$ tel que $\xi|_{U_j} = 0$ pour tout
$j$. D'après le lemme \ref{sites-cohomology-lemma-cech-h1}, cela signifie que $\xi$ provient
d'un élément
$\check \xi \in \check H^1(\mathcal{U}, \mathcal{F})$.
Puisque les morphismes
$\check H^1(\mathcal{U}, \mathcal{F}_i) \to H^1(U, \mathcal{F}_i)$
sont injectifs pour tout $i$ (d'après le lemme \ref{sites-cohomology-lemma-cech-h1}) et puisque
l'image de $\xi$ dans $\prod H^1(U, \mathcal{F}_i)$ est nulle, nous voyons
que l'image
$\check \xi_i = 0$ dans $\check H^1(\mathcal{U}, \mathcal{F}_i)$.
Toutefois, comme $\mathcal{F} = \prod \mathcal{F}_i$, le complexe
$\check{\mathcal{C}}^\bullet(\mathcal{U}, \mathcal{F})$ est le produit des
complexes $\check{\mathcal{C}}^\bullet(\mathcal{U}, \mathcal{F}_i)$.
D'après Homologie,
lemme \ref{homology-lemma-product-abelian-groups-exact}, nous concluons que
$\check \xi = 0$, comme voulu.
\end{proof}

\begin{lemma}
\label{sites-cohomology-lemma-restriction-along-monomorphism-surjective}
Soit $(\mathcal{C}, \mathcal{O})$ un site annelé. Soit $a : U' \to U$ un
monomorphisme dans $\mathcal{C}$. Alors, pour tout $\mathcal{O}$-module
injectif $\mathcal{I}$, le morphisme de restriction
$\mathcal{I}(U) \to \mathcal{I}(U')$ est surjectif.
\end{lemma}

\begin{proof}
Soient $j : \mathcal{C}/U \to \mathcal{C}$ et
$j' : \mathcal{C}/U' \to \mathcal{C}$ les morphismes de localisation
(Modules sur les sites, section \ref{sites-modules-section-localize}).
Comme $j_!$ est adjoint à gauche de la restriction, nous avons, pour tout
faisceau $\mathcal{F}$ de $\mathcal{O}$-modules,
$$
\Hom_\mathcal{O}(j_!\mathcal{O}_U, \mathcal{F})
=
\Hom_{\mathcal{O}_U}(\mathcal{O}_U, \mathcal{F}|_U)
=
\mathcal{F}(U)
$$
De même, le faisceau $j'_!\mathcal{O}_{U'}$ représente le foncteur
$\mathcal{F} \mapsto \mathcal{F}(U')$.
Nous décrivons ci-dessous un morphisme canonique de $\mathcal{O}$-modules
$$
j'_!\mathcal{O}_{U'} \longrightarrow j_!\mathcal{O}_U
$$
qui correspond, par le lemme de Yoneda, au morphisme de restriction
$\mathcal{F}(U) \to \mathcal{F}(U')$ (Catégories,
lemme \ref{categories-lemma-yoneda}).
Il suffit de prouver que le morphisme de modules affiché est injectif ; voir
Homologie, lemme \ref{homology-lemma-characterize-injectives}.

\medskip\noindent
Pour construire notre morphisme, il suffit de construire un morphisme entre
les préfaisceaux qui associent à un objet $V$ de $\mathcal{C}$ le
$\mathcal{O}(V)$-module
$$
\bigoplus\nolimits_{\varphi' \in \Mor_\mathcal{C}(V, U')} \mathcal{O}(V)
\quad\text{et}\quad
\bigoplus\nolimits_{\varphi \in \Mor_\mathcal{C}(V, U)} \mathcal{O}(V)
$$
voir Modules sur les sites,
lemme \ref{sites-modules-lemma-extension-by-zero}.
Nous prenons le morphisme qui envoie le facteur correspondant à $\varphi'$
sur celui correspondant à $\varphi = a \circ \varphi'$ par l'identité de
$\mathcal{O}(V)$. Comme $a$ est un monomorphisme, ce morphisme est injectif.
La faisceautisation étant exacte, le résultat en découle.
\end{proof}






\section{Faisceaux totalement acycliques}
\label{sites-cohomology-section-limp}

\noindent
Soit $(\mathcal{C}, \mathcal{O})$ un site annelé.
Soit $K$ un préfaisceau d'ensembles sur $\mathcal{C}$ (nous employons
intentionnellement ici une capitale romaine pour le distinguer des faisceaux
abéliens). Étant donné un faisceau de $\mathcal{O}$-modules $\mathcal{F}$,
posons
$$
\mathcal{F}(K) =
\Mor_{\textit{PSh}(\mathcal{C})}(K, \mathcal{F}) =
\Mor_{\Sh(\mathcal{C})}(K^\#, \mathcal{F})
$$
Le foncteur $\mathcal{F} \mapsto \mathcal{F}(K)$ est un foncteur exact à
gauche $\textit{Mod}(\mathcal{O}) \to \textit{Ab}$ ; il possède donc des
foncteurs dérivés à droite. Nous les noterons $H^p(K, \mathcal{F})$, de sorte
que $H^0(K, \mathcal{F}) = \mathcal{F}(K)$.

\medskip\noindent
Voici quelques observations :
\begin{enumerate}
\item Puisque $\mathcal{F}(K) = \mathcal{F}(K^\#)$, nous avons
$H^p(K, \mathcal{F}) = H^p(K^\#, \mathcal{F})$.
Autoriser $K$ à être un préfaisceau dans la définition ci-dessus n'est qu'une
commodité de notation.
\item Supposons que $K = h_U$ ou $K = h_U^\#$ pour un objet $U$ de
$\mathcal{C}$. Alors $H^p(K, \mathcal{F}) = H^p(U, \mathcal{F})$, car
$\Mor_{\Sh(\mathcal{C})}(h_U^\#, \mathcal{F}) = \mathcal{F}(U)$ ; voir
Sites, section \ref{sites-section-representable-sheaves}.
\item Si $\mathcal{O} = \mathbf{Z}$ (le faisceau constant), alors les groupes
de cohomologie sont des foncteurs
$H^p(K, - ) : \textit{Ab}(\mathcal{C}) \to \textit{Ab}$
car $\textit{Mod}(\mathcal{O}) = \textit{Ab}(\mathcal{C})$ dans ce cas.
\end{enumerate}
Nous pouvons reformuler dans ce langage certains résultats déjà démontrés.

\begin{lemma}
\label{sites-cohomology-lemma-compute-cohomology-on-sheaf-sets}
Soit $(\mathcal{C}, \mathcal{O})$ un site annelé.
Soit $K$ un préfaisceau d'ensembles sur $\mathcal{C}$.
Soit $\mathcal{F}$ un $\mathcal{O}$-module, et notons $\mathcal{F}_{ab}$ le
faisceau de groupes abéliens sous-jacent. Alors
$H^p(K, \mathcal{F}) = H^p(K, \mathcal{F}_{ab})$.
\end{lemma}

\begin{proof}
Nous pouvons remplacer $K$ par son faisceautisé et supposer que $K$ est un
faisceau. Remarquons que $H^p(K, \mathcal{F})$ et
$H^p(K, \mathcal{F}_{ab})$ ne dépendent que du topos, et non du site
sous-jacent. D'après Sites, lemme \ref{sites-lemma-topos-good-site}, nous
pouvons donc remplacer $\mathcal{C}$ par un site « plus grand » tel que
$K = h_U$ pour un objet $U$ de $\mathcal{C}$. Dans ce cas, le résultat découle
du lemme \ref{sites-cohomology-lemma-cohomology-modules-abelian-agree}.
\end{proof}

\begin{lemma}
\label{sites-cohomology-lemma-cech-to-cohomology-sheaf-sets}
Soit $\mathcal{C}$ un site. Soit $K' \to K$ un morphisme de préfaisceaux
d'ensembles sur $\mathcal{C}$ dont le faisceautisé est surjectif. Posons
$K'_p = K' \times_K \ldots \times_K K'$ ($p + 1$ facteurs).
Pour tout faisceau abélien $\mathcal{F}$, il existe une suite spectrale
vérifiant $E_1^{p, q} = H^q(K'_p, \mathcal{F})$ et convergeant vers
$H^{p + q}(K, \mathcal{F})$.
\end{lemma}

\begin{proof}
Puisque la faisceautisation est exacte, nous voyons que
$(K_p')^\#$ est égal à
$(K')^\# \times_{K^\#} \ldots \times_{K^\#} (K')^\#$
($p + 1$ facteurs). Nous pouvons donc remplacer $K$ et $K'$ par leurs
faisceautisés et supposer que $K' \to K$ est un morphisme surjectif de
faisceaux. Après avoir remplacé $\mathcal{C}$ par un site « plus grand » comme
dans Sites, lemme \ref{sites-lemma-topos-good-site}, nous pouvons supposer que
$K, K'$ sont des objets de $\mathcal{C}$ et que
$\mathcal{U} = \{K' \to K\}$ est un recouvrement. Nous disposons alors de la
suite spectrale de la cohomologie de {\v C}ech vers la cohomologie du
lemme \ref{sites-cohomology-lemma-cech-spectral-sequence}, dont la page $E_1$ est celle indiquée
dans l'énoncé.
\end{proof}

\begin{lemma}
\label{sites-cohomology-lemma-cohomology-on-sheaf-sets}
Soit $\mathcal{C}$ un site. Soit $K$ un faisceau d'ensembles sur
$\mathcal{C}$. Considérons le morphisme de topos
$j : \Sh(\mathcal{C}/K) \to \Sh(\mathcal{C})$ ; voir
Sites, lemme \ref{sites-lemma-localize-topos-site}.
Alors $j^{-1}$ préserve les injectifs et
$H^p(K, \mathcal{F}) = H^p(\mathcal{C}/K, j^{-1}\mathcal{F})$
pour tout faisceau abélien $\mathcal{F}$ sur $\mathcal{C}$.
\end{lemma}

\begin{proof}
D'après Sites, lemmes \ref{sites-lemma-localize-topos} et
\ref{sites-lemma-localize-topos-site}, le morphisme de topos $j$ est équivalent
à une localisation. Le résultat découle donc du
lemme \ref{sites-cohomology-lemma-cohomology-of-open}.
\end{proof}

\noindent
Compte tenu du lemme \ref{sites-cohomology-lemma-compute-cohomology-on-sheaf-sets}, nous voyons
que la définition suivante est également la « bonne » pour les faisceaux de
modules sur les sites annelés.

\begin{definition}
\label{sites-cohomology-definition-limp}
Soit $\mathcal{C}$ un site.
On dit qu'un faisceau abélien $\mathcal{F}$ est
{\it totalement acyclique}\footnote{Bien que cette terminologie soit employée
dans \cite[Vbis, Proposition 1.3.10]{SGA4}, elle n'est probablement pas
standard. Dans \cite[V, définition 4.1]{SGA4}, cette propriété est appelée
« flasque », mais nous ne pouvons pas employer ce terme car il entrerait en
conflit avec notre définition des faisceaux flasques sur les espaces
topologiques. Veuillez écrire à
\href{mailto:stacks.project@gmail.com}{stacks.project@gmail.com}
si vous avez une meilleure suggestion.} si, pour tout faisceau d'ensembles
$K$, nous avons $H^p(K, \mathcal{F}) = 0$ pour tout $p \geq 1$.
\end{definition}

\noindent
Il est clair que la propriété d'être totalement acyclique est intrinsèque,
c'est-à-dire préservée par les équivalences de topos.
La cohomologie supérieure d'un faisceau totalement acyclique s'annule sur tous
les objets du site, mais, en général, la condition d'acyclicité totale est
strictement plus forte. Voici une caractérisation parfois utile des faisceaux
totalement acycliques.

\begin{lemma}
\label{sites-cohomology-lemma-characterize-limp}
Soit $\mathcal{C}$ un site. Soit $\mathcal{F}$ un faisceau abélien. Si
\begin{enumerate}
\item $H^p(U, \mathcal{F}) = 0$ pour $p > 0$ et
$U \in \Ob(\mathcal{C})$ ; et
\item pour toute surjection $K' \to K$ de faisceaux d'ensembles, le complexe
de {\v C}ech augmenté
$$
0 \to H^0(K, \mathcal{F}) \to H^0(K', \mathcal{F}) \to
H^0(K' \times_K K', \mathcal{F}) \to \ldots
$$
soit exact,
\end{enumerate}
alors $\mathcal{F}$ est totalement acyclique (et la réciproque est vraie).
\end{lemma}

\begin{proof}
D'après l'hypothèse (1), nous avons
$H^p(h_U^\#, \mathcal{F}) = 0$ pour tout $p > 0$ et tout objet $U$ de
$\mathcal{C}$. Remarquons que si $K = \coprod K_i$ est un coproduit de
faisceaux d'ensembles sur $\mathcal{C}$, alors
$H^p(K, \mathcal{F}) = \prod H^p(K_i, \mathcal{F})$.
Pour tout faisceau d'ensembles $K$, il existe une surjection
$$
K' = \coprod h_{U_i}^\# \longrightarrow K
$$
voir Sites, lemme \ref{sites-lemma-sheaf-coequalizer-representable}.
Nous en concluons que : (*) pour tout faisceau d'ensembles $K$, il existe une
surjection $K' \to K$ de faisceaux d'ensembles telle que
$H^p(K', \mathcal{F}) = 0$ pour $p > 0$. Nous affirmons que (*) et la
condition (2) impliquent que $\mathcal{F}$ est totalement acyclique.
Remarquons que les conditions (*) et (2) ne dépendent que de $\mathcal{F}$ en
tant qu'objet du topos $\Sh(\mathcal{C})$, et non du site sous-jacent.
(Nous n'utiliserons plus la propriété (1) dans la suite de la preuve.)

\medskip\noindent
Nous allons montrer par récurrence sur $n \geq 0$ que (*) et (2) impliquent
l'hypothèse de récurrence suivante $IH_n$ :
$H^p(K, \mathcal{F}) = 0$ pour tout $0 < p \leq n$ et tout faisceau
d'ensembles $K$. Remarquons que $IH_0$ est vraie. Supposons $IH_n$.
Choisissons un faisceau d'ensembles $K$, puis une surjection $K' \to K$ telle
que $H^p(K', \mathcal{F}) = 0$ pour tout $p > 0$. Nous disposons d'une suite
spectrale vérifiant
$$
E_1^{p, q} = H^q(K'_p, \mathcal{F})
$$
et convergeant vers $H^{p + q}(K, \mathcal{F})$ ; voir le
lemme \ref{sites-cohomology-lemma-cech-to-cohomology-sheaf-sets}.
D'après $IH_n$, $E_1^{p, q} = 0$ pour $0 < q \leq n$, et l'hypothèse (2)
donne $E_2^{p, 0} = 0$ pour $p > 0$. Enfin, $E_1^{0, q} = 0$ pour $q > 0$,
car $H^q(K', \mathcal{F}) = 0$ par choix de $K'$. Nous concluons donc que
$H^{n + 1}(K, \mathcal{F}) = 0$, puisque tous les termes $E_2^{p, q}$ tels
que $p + q = n + 1$ sont nuls.
\end{proof}







\section{La suite spectrale de Leray}
\label{sites-cohomology-section-leray}

\noindent
Le lemme suivant est la clé de la preuve de l'existence de la suite spectrale
de Leray.

\begin{lemma}
\label{sites-cohomology-lemma-direct-image-injective-sheaf}
Soit $f : (\Sh(\mathcal{C}), \mathcal{O}_\mathcal{C}) \to
(\Sh(\mathcal{D}), \mathcal{O}_\mathcal{D})$ un morphisme de topos annelés.
Alors, pour tout objet injectif $\mathcal{I}$ de
$\textit{Mod}(\mathcal{O}_\mathcal{C})$, l'image directe
$f_*\mathcal{I}$ est totalement acyclique.
\end{lemma}

\begin{proof}
Soit $K$ un faisceau d'ensembles sur $\mathcal{D}$.
D'après Modules sur les sites, lemme
\ref{sites-modules-lemma-morphism-ringed-topoi-comes-from-morphism-ringed-sites}
nous pouvons remplacer $\mathcal{C}$ et $\mathcal{D}$ par des sites « plus
grands » de sorte que $f$ provienne d'un morphisme de sites annelés induit par
un foncteur continu $u : \mathcal{D} \to \mathcal{C}$ et que $K = h_V$ pour
un objet $V$ de $\mathcal{D}$.

\medskip\noindent
Il nous suffit donc de montrer que $H^q(V, f_*\mathcal{I})$ est nul pour
$q > 0$ et tout objet $V$ de $\mathcal{D}$ lorsque $f$ est donné par un
morphisme de sites annelés. Soit $\mathcal{V} = \{V_j \to V\}$ un
recouvrement quelconque de $\mathcal{D}$. Comme $u$ est continu,
$\mathcal{U} = \{u(V_j) \to u(V)\}$ est un recouvrement de $\mathcal{C}$.
Nous avons alors une égalité de complexes de {\v C}ech
$$
\check{\mathcal{C}}^\bullet(\mathcal{V}, f_*\mathcal{I})
=
\check{\mathcal{C}}^\bullet(\mathcal{U}, \mathcal{I})
$$
par définition de $f_*$. D'après le
lemme \ref{sites-cohomology-lemma-injective-module-trivial-cech}, la cohomologie de ce complexe
est nulle en degrés positifs. Le résultat découle alors du
lemme \ref{sites-cohomology-lemma-cech-vanish-collection}.
\end{proof}

\noindent
Pour les morphismes plats, le foncteur $f_*$ préserve les modules injectifs.
En particulier, le foncteur
$f_* : \textit{Ab}(\mathcal{C}) \to \textit{Ab}(\mathcal{D})$ transforme
toujours les faisceaux abéliens injectifs en faisceaux abéliens injectifs.

\begin{lemma}
\label{sites-cohomology-lemma-pushforward-injective-flat}
Soit $f : (\Sh(\mathcal{C}), \mathcal{O}_\mathcal{C}) \to
(\Sh(\mathcal{D}), \mathcal{O}_\mathcal{D})$ un morphisme de topos annelés.
Si $f$ est plat, alors $f_*\mathcal{I}$ est un
$\mathcal{O}_\mathcal{D}$-module injectif pour tout
$\mathcal{O}_\mathcal{C}$-module injectif $\mathcal{I}$.
\end{lemma}

\begin{proof}
Dans ce cas, le foncteur $f^*$ est exact ; voir Modules sur les sites,
lemme \ref{sites-modules-lemma-flat-pullback-exact}.
Le résultat découle donc d'Homologie,
lemme \ref{homology-lemma-adjoint-preserve-injectives}.
\end{proof}

\begin{lemma}
\label{sites-cohomology-lemma-limp-acyclic}
Soit $(\Sh(\mathcal{C}), \mathcal{O}_\mathcal{C})$ un topos annelé.
Un faisceau totalement acyclique est acyclique à droite pour les foncteurs
suivants :
\begin{enumerate}
\item le foncteur $H^0(U, -)$ pour tout objet $U$ de $\mathcal{C}$ ;
\item le foncteur $\mathcal{F} \mapsto \mathcal{F}(K)$ pour tout préfaisceau
d'ensembles $K$ ;
\item le foncteur de sections globales $\Gamma(\mathcal{C}, -)$ ;
\item le foncteur $f_*$ pour tout morphisme
$f : (\Sh(\mathcal{C}), \mathcal{O}_\mathcal{C}) \to
(\Sh(\mathcal{D}), \mathcal{O}_\mathcal{D})$ de topos annelés.
\end{enumerate}
\end{lemma}

\begin{proof}
La partie (2) est la définition d'un faisceau totalement acyclique.
La partie (1) découle de (2), comme indiqué dans la discussion qui suit la
définition des faisceaux totalement acycliques.
La partie (3) est le cas particulier de (2) où $K = e$ est l'objet final de
$\Sh(\mathcal{C})$.

\medskip\noindent
Pour prouver (4), nous pouvons supposer, d'après Modules sur les sites, lemme
\ref{sites-modules-lemma-morphism-ringed-topoi-comes-from-morphism-ringed-sites}
que $f$ est donné par un morphisme de sites. Dans ce cas, la description des
images directes supérieures du lemme \ref{sites-cohomology-lemma-higher-direct-images} montre
que les $R^if_*$, $i > 0$, d'un faisceau totalement acyclique sont nuls.
\end{proof}

\begin{remark}
\label{sites-cohomology-remark-before-Leray}
Les résultats précédents montrent que Catégories dérivées,
lemme \ref{derived-lemma-compose-derived-functors}, s'applique dans plusieurs
situations. Par exemple, pour tout morphisme
$f : (\Sh(\mathcal{C}), \mathcal{O}_\mathcal{C}) \to
(\Sh(\mathcal{D}), \mathcal{O}_\mathcal{D})$ de topos annelés, nous avons
$$
R\Gamma(\mathcal{D}, Rf_*\mathcal{F}) = R\Gamma(\mathcal{C}, \mathcal{F})
$$
pour tout faisceau de $\mathcal{O}_\mathcal{C}$-modules $\mathcal{F}$.
En effet, pour un $\mathcal{O}_\mathcal{C}$-module injectif $\mathcal{I}$, le
$\mathcal{O}_\mathcal{D}$-module $f_*\mathcal{I}$ est totalement acyclique
d'après le lemme \ref{sites-cohomology-lemma-direct-image-injective-sheaf}, et un faisceau
totalement acyclique est acyclique pour $\Gamma(\mathcal{D}, -)$ d'après le
lemme \ref{sites-cohomology-lemma-limp-acyclic}.
\end{remark}

\begin{lemma}[Suite spectrale de Leray]
\label{sites-cohomology-lemma-Leray}
Soit $f : (\Sh(\mathcal{C}), \mathcal{O}_\mathcal{C}) \to
(\Sh(\mathcal{D}), \mathcal{O}_\mathcal{D})$ un morphisme de topos annelés.
Soit $\mathcal{F}^\bullet$ un complexe borné inférieurement de
$\mathcal{O}_\mathcal{C}$-modules. Il existe une suite spectrale
$$
E_2^{p, q} = H^p(\mathcal{D}, R^qf_*(\mathcal{F}^\bullet))
$$
qui converge vers $H^{p + q}(\mathcal{C}, \mathcal{F}^\bullet)$.
\end{lemma}

\begin{proof}
Il s'agit simplement de la suite spectrale de Grothendieck de Catégories
dérivées, lemme \ref{derived-lemma-grothendieck-spectral-sequence}, provenant
de la composition de foncteurs
$\Gamma(\mathcal{C}, -) = \Gamma(\mathcal{D}, -) \circ f_*$.
Pour vérifier que les hypothèses de Catégories dérivées,
lemme \ref{derived-lemma-grothendieck-spectral-sequence}, sont satisfaites,
voir les lemmes \ref{sites-cohomology-lemma-direct-image-injective-sheaf} et
\ref{sites-cohomology-lemma-limp-acyclic}.
\end{proof}

\begin{lemma}
\label{sites-cohomology-lemma-apply-Leray}
Soit $f : (\Sh(\mathcal{C}), \mathcal{O}_\mathcal{C}) \to
(\Sh(\mathcal{D}), \mathcal{O}_\mathcal{D})$ un morphisme de topos annelés.
Soit $\mathcal{F}$ un $\mathcal{O}_\mathcal{C}$-module.
\begin{enumerate}
\item Si $R^qf_*\mathcal{F} = 0$ pour $q > 0$, alors
$H^p(\mathcal{C}, \mathcal{F}) = H^p(\mathcal{D}, f_*\mathcal{F})$ pour tout
$p$.
\item Si $H^p(\mathcal{D}, R^qf_*\mathcal{F}) = 0$ pour tout $q$ et tout
$p > 0$, alors
$H^q(\mathcal{C}, \mathcal{F}) = H^0(\mathcal{D}, R^qf_*\mathcal{F})$ pour
tout $q$.
\end{enumerate}
\end{lemma}

\begin{proof}
Ce sont deux conditions simples qui forcent la suite spectrale de Leray à
dégénérer. On peut aussi démontrer ces faits directement, sans utiliser la
suite spectrale ; c'est un bon exercice de cohomologie des faisceaux.
\end{proof}

\begin{lemma}[Suite spectrale de Leray relative]
\label{sites-cohomology-lemma-relative-Leray}
Soient
$f : (\Sh(\mathcal{C}), \mathcal{O}_\mathcal{C}) \to
(\Sh(\mathcal{D}), \mathcal{O}_\mathcal{D})$
et
$g : (\Sh(\mathcal{D}), \mathcal{O}_\mathcal{D}) \to
(\Sh(\mathcal{E}), \mathcal{O}_\mathcal{E})$
soient des morphismes de topos annelés.
Soit $\mathcal{F}$ un $\mathcal{O}_\mathcal{C}$-module.
Il existe une suite spectrale telle que
$$
E_2^{p, q} = R^pg_*(R^qf_*\mathcal{F})
$$
qui converge vers $R^{p + q}(g \circ f)_*\mathcal{F}$.
Cette suite spectrale est fonctorielle en $\mathcal{F}$, et il en existe une
version pour les complexes bornés inférieurement de
$\mathcal{O}_\mathcal{C}$-modules.
\end{lemma}

\begin{proof}
C'est une suite spectrale de Grothendieck pour la composition de foncteurs ;
voir Catégories dérivées,
lemme \ref{derived-lemma-grothendieck-spectral-sequence}, ainsi que les lemmes
\ref{sites-cohomology-lemma-direct-image-injective-sheaf} et \ref{sites-cohomology-lemma-limp-acyclic}.
\end{proof}







\section{Le morphisme de changement de base}
\label{sites-cohomology-section-base-change-map}

\noindent
Dans cette section, nous construisons le morphisme de changement de base dans
certains cas ; le cas général est traité dans la remarque
\ref{sites-cohomology-remark-base-change}. Afin d'éviter ici le recours à l'image inverse
dérivée, nous nous restreignons au cas d'un changement de base par un morphisme
plat de sites annelés. Avant d'énoncer le résultat, examinons l'image inverse
plate sur la catégorie dérivée. Supposons que
$g : (\Sh(\mathcal{C}), \mathcal{O}_\mathcal{C})
\to (\Sh(\mathcal{D}), \mathcal{O}_\mathcal{D})$
soit un morphisme plat de topos annelés. D'après Modules sur les sites,
lemme \ref{sites-modules-lemma-flat-pullback-exact}, le foncteur
$g^* : \textit{Mod}(\mathcal{O}_\mathcal{D}) \to
\textit{Mod}(\mathcal{O}_\mathcal{C})$ est exact.
Il admet donc un foncteur dérivé
$$
g^* : D(\mathcal{O}_\mathcal{D}) \to D(\mathcal{O}_\mathcal{C})
$$
qui se calcule simplement en prenant l'image inverse d'un représentant d'un
objet donné de $D(\mathcal{O}_\mathcal{D})$ ; voir Catégories dérivées,
lemme \ref{derived-lemma-right-derived-exact-functor}.
Il préserve les sous-catégories bornées (supérieurement ou inférieurement).
Ainsi, comme annoncé, nous notons ce foncteur $g^*$ plutôt que $Lg^*$.

\begin{lemma}
\label{sites-cohomology-lemma-base-change-map-flat-case}
Soit
$$
\xymatrix{
(\Sh(\mathcal{C}'), \mathcal{O}_{\mathcal{C}'})
\ar[r]_{g'} \ar[d]_{f'} &
(\Sh(\mathcal{C}), \mathcal{O}_\mathcal{C}) \ar[d]^f \\
(\Sh(\mathcal{D}'), \mathcal{O}_{\mathcal{D}'})
\ar[r]^g &
(\Sh(\mathcal{D}), \mathcal{O}_\mathcal{D})
}
$$
un diagramme commutatif de topos annelés.
Soit $\mathcal{F}^\bullet$ un complexe borné inférieurement de
$\mathcal{O}_\mathcal{C}$-modules.
Supposons que $g$ et $g'$ soient plats.
Il existe alors un morphisme canonique de changement de base
$$
g^*Rf_*\mathcal{F}^\bullet
\longrightarrow
R(f')_*(g')^*\mathcal{F}^\bullet
$$
dans $D^{+}(\mathcal{O}_{\mathcal{D}'})$.
\end{lemma}

\begin{proof}
Choisissons des résolutions injectives $\mathcal{F}^\bullet \to \mathcal{I}^\bullet$
et $(g')^*\mathcal{F}^\bullet \to \mathcal{J}^\bullet$.
D'après le lemme \ref{sites-cohomology-lemma-pushforward-injective-flat},
$(g')_*\mathcal{J}^\bullet$ est un complexe d'injectifs qui représente
$R(g')_*(g')^*\mathcal{F}^\bullet$. Par Catégories dérivées, lemmes
\ref{derived-lemma-morphisms-lift}
et \ref{derived-lemma-morphisms-equal-up-to-homotopy},
le morphisme $\beta$ du diagramme
$$
\xymatrix{
(g')_*(g')^*\mathcal{F}^\bullet \ar[r] &
(g')_*\mathcal{J}^\bullet \\
\mathcal{F}^\bullet \ar[u]^{adjunction} \ar[r] &
\mathcal{I}^\bullet \ar[u]_\beta
}
$$
existe et est unique à homotopie près.
En prenant l'image directe vers $\mathcal{D}$, nous obtenons
$$
f_*\beta :
f_*\mathcal{I}^\bullet
\longrightarrow
f_*(g')_*\mathcal{J}^\bullet
=
g_*(f')_*\mathcal{J}^\bullet
$$
Par adjonction entre $g^*$ et $g_*$, nous obtenons un morphisme de complexes
$g^*f_*\mathcal{I}^\bullet \to (f')_*\mathcal{J}^\bullet$.
Ce morphisme est unique à homotopie près, puisque le seul choix effectué au
cours de la construction est celui du morphisme $\beta$ et que tout a été fait
au niveau des complexes.
\end{proof}








\section{Cohomologie et colimites}
\label{sites-cohomology-section-limits}

\noindent
Soit $(\mathcal{C}, \mathcal{O})$ un site annelé.
Soit $\mathcal{I} \to \textit{Mod}(\mathcal{O})$, $i \mapsto \mathcal{F}_i$,
un diagramme indexé par la catégorie $\mathcal{I}$ ; voir
Catégories, section \ref{categories-section-limits}.
Pour tout $i$, il existe un morphisme canonique
$\mathcal{F}_i \to \colim_i \mathcal{F}_i$ qui induit
un morphisme en cohomologie. Nous obtenons donc un morphisme canonique
$$
\colim_i H^p(U, \mathcal{F}_i)
\longrightarrow
H^p(U, \colim_i \mathcal{F}_i)
$$
pour tout $p \geq 0$ et tout objet $U$ de $\mathcal{C}$.
En général, ces morphismes ne sont pas des isomorphismes, même pour $p = 0$.

\medskip\noindent
Le lemme suivant est l'analogue, pour la cohomologie, du
lemme \ref{sites-lemma-directed-colimits-sections} de Sites.

\begin{lemma}
\label{sites-cohomology-lemma-colim-works-over-collection}
Soit $\mathcal{C}$ un site. Soit $\text{Cov}_\mathcal{C}$ l'ensemble
des recouvrements de $\mathcal{C}$ (voir
Sites, définition \ref{sites-definition-site}). Soient
$\mathcal{B} \subset \Ob(\mathcal{C})$ et
$\text{Cov} \subset \text{Cov}_\mathcal{C}$
des sous-ensembles. Supposons que
\begin{enumerate}
\item pour tout $\mathcal{U} \in \text{Cov}$, on ait
$\mathcal{U} = \{U_i \to U\}_{i \in I}$ avec $I$ fini,
$U, U_i \in \mathcal{B}$ et que tout
$U_{i_0} \times_U \ldots \times_U U_{i_p} \in \mathcal{B}$.
\item pour tout $U \in \mathcal{B}$, les recouvrements de $U$
appartenant à $\text{Cov}$ forment un système cofinal de recouvrements de $U$.
\end{enumerate}
Alors le morphisme
$$
\colim_i H^p(U, \mathcal{F}_i)
\longrightarrow
H^p(U, \colim_i \mathcal{F}_i)
$$
est un isomorphisme pour tout $p \geq 0$, tout $U \in \mathcal{B}$ et
tout diagramme filtrant $\mathcal{I} \to \textit{Ab}(\mathcal{C})$.
\end{lemma}

\begin{proof}
Nous raisonnons par récurrence sur $p$.
Remarquons qu'en (1), nous exigeons que les recouvrements
$\mathcal{U} \in \text{Cov}$ soient finis, de sorte que tous les éléments de
$\mathcal{B}$ soient quasi-compacts. Ainsi, (2) et (1) impliquent que tout
$U \in \mathcal{B}$ satisfait à l'hypothèse (4) du lemme
\ref{sites-lemma-directed-colimits-sections} de Sites.
Le résultat est donc vrai pour $p = 0$. Supposons-le vrai pour $p$ et
démontrons-le pour $p + 1$.

\medskip\noindent
Choisissons un diagramme filtrant
$\mathcal{F} : \mathcal{I} \to \textit{Ab}(\mathcal{C})$,
$i \mapsto \mathcal{F}_i$.
Puisque $\textit{Ab}(\mathcal{C})$ possède des plongements injectifs
fonctoriels, voir Injectifs, théorème
\ref{injectives-theorem-sheaves-injectives}, nous pouvons trouver un
morphisme de diagrammes filtrants
$\mathcal{F} \to \mathcal{I}$
tel que chaque $\mathcal{F}_i \to \mathcal{I}_i$ soit un monomorphisme de
faisceaux abéliens dans un faisceau abélien injectif. Notons $\mathcal{Q}_i$
son conoyau ; nous avons ainsi des suites exactes courtes
$$
0 \to
\mathcal{F}_i \to
\mathcal{I}_i \to
\mathcal{Q}_i \to 0.
$$
Comme les colimites de faisceaux sont les faisceautisés des colimites au
niveau des préfaisceaux, comme la faisceautisation est exacte et comme les
colimites filtrantes de groupes abéliens sont exactes
(voir Algèbre, lemme \ref{algebra-lemma-directed-colimit-exact}),
la suite
$$
0 \to
\colim_i \mathcal{F}_i \to
\colim_i \mathcal{I}_i \to
\colim_i \mathcal{Q}_i \to 0.
$$
est elle aussi exacte courte. Nous affirmons que
$H^q(U, \colim_i \mathcal{I}_i) = 0$ pour tout $U \in \mathcal{B}$
pour tout $q \geq 1$. Admettons provisoirement cette affirmation et
considérons le diagramme
$$
\xymatrix{
\colim_i H^p(U, \mathcal{I}_i) \ar[d] \ar[r] &
\colim_i H^p(U, \mathcal{Q}_i) \ar[d] \ar[r] &
\colim_i H^{p + 1}(U, \mathcal{F}_i) \ar[d] \ar[r] &
0 \ar[d] \\
H^p(U, \colim_i \mathcal{I}_i) \ar[r] &
H^p(U, \colim_i \mathcal{Q}_i) \ar[r] &
H^{p + 1}(U, \colim_i \mathcal{F}_i) \ar[r] &
0
}
$$
Le zéro situé en bas à droite provient de l'affirmation, et celui situé en
haut à droite du fait que les faisceaux $\mathcal{I}_i$ sont injectifs.
La ligne du haut est exacte par application du
lemme \ref{algebra-lemma-directed-colimit-exact} d'Algèbre.
Le lemme du serpent donne alors le résultat pour $p + 1$.

\medskip\noindent
Il reste à démontrer l'affirmation. Nous utiliserons le
lemme \ref{sites-cohomology-lemma-cech-vanish-collection}.
D'après le cas $p = 0$, pour $\mathcal{U} \in \text{Cov}$, nous avons
$$
\check{\mathcal{C}}^\bullet(\mathcal{U}, \colim_i \mathcal{I}_i)
=
\colim_i \check{\mathcal{C}}^\bullet(\mathcal{U}, \mathcal{I}_i)
$$
car tous les $U_{j_0} \times_U \ldots \times_U U_{j_p}$
appartiennent à $\mathcal{B}$. D'après le
lemme \ref{sites-cohomology-lemma-injective-trivial-cech}, chacun des complexes de la colimite
de complexes de {\v C}ech est acyclique en degrés $\geq 1$. Par le
lemme \ref{algebra-lemma-directed-colimit-exact} d'Algèbre, le complexe de
{\v C}ech
$\check{\mathcal{C}}^\bullet(\mathcal{U}, \colim_i \mathcal{I}_i)$
est donc lui aussi acyclique en degrés $\geq 1$. Autrement dit,
$\check{H}^p(\mathcal{U}, \colim_i \mathcal{I}_i) = 0$
pour tout $p \geq 1$. Les hypothèses du
lemme \ref{sites-cohomology-lemma-cech-vanish-collection} sont donc satisfaites, d'où
l'affirmation.
\end{proof}

\begin{lemma}
\label{sites-cohomology-lemma-colim-global}
Soit $\mathcal{C}$ un site. Soit $S \subset \Ob(\Sh(\mathcal{C}))$
un sous-ensemble. Notons $*$ l'objet final de $\Sh(\mathcal{C})$. Supposons que
\begin{enumerate}
\item pour un certain $K \in S$, le morphisme $K \to *$ soit surjectif ;
\item pour tout morphisme surjectif de faisceaux $\mathcal{F} \to K$ avec
$K \in S$, il existe $K' \in S$ et un morphisme $K' \to \mathcal{F}$ tels
que la composée $K' \to K$ soit surjective ;
\item pour tous $K, K' \in S$, il existe une surjection
$K'' \to K \times K'$ avec $K'' \in S$ ;
\item pour tous $a, b : K \to K'$ avec $K, K' \in S$, il existe une
surjection $K'' \to \text{Equalizer}(a, b)$ avec $K'' \in S$ ; et
\item tout $K \in S$ soit quasi-compact
(Sites, définition \ref{sites-definition-quasi-compact-topos}).
\end{enumerate}
Alors, pour tout $p \geq 0$, le morphisme
$$
\colim_\lambda H^p(\mathcal{C}, \mathcal{F}_\lambda)
\longrightarrow
H^p(\mathcal{C}, \colim_\lambda \mathcal{F}_\lambda)
$$
est un isomorphisme pour tout diagramme filtrant
$\Lambda \to \textit{Ab}(\mathcal{C})$, $\lambda \mapsto \mathcal{F}_\lambda$.
\end{lemma}

\begin{proof}
Nous raisonnons par récurrence sur $p$. Le cas initial $p = 0$ découle de la
partie (4) du lemme \ref{sites-lemma-directed-colimits-global-sections} de
Sites. Vérifions les hypothèses, bien que nous invitions le lecteur à omettre
cette partie. Supposons que $\mathcal{F} \to *$ soit surjectif.
Choisissons $K \in S$ et une surjection $K \to *$ comme en (1). Alors
$\mathcal{F} \times K \to K$ est surjectif. Comme en (2), choisissons
$K' \to \mathcal{F} \times K$ avec $K' \in S$ et tel que $K' \to K$
soit surjectif.
Il existe alors un morphisme $K' \to \mathcal{F}$, et $K' \to *$ est
surjectif. L'hypothèse (4)(a) du lemme
\ref{sites-lemma-directed-colimits-global-sections} de Sites est donc
satisfaite. D'après le lemme \ref{sites-lemma-image-of-quasi-compact} de Sites
et les hypothèses (3) et (5), $K \times K$ est quasi-compact pour tout
$K \in S$. L'hypothèse (4)(b) du lemme
\ref{sites-lemma-directed-colimits-global-sections} de Sites est donc
satisfaite. Ceci achève la preuve du cas initial.

\medskip\noindent
Étape de récurrence. Supposons le résultat vrai pour $H^p$, lorsque
$p \leq p_0$, et pour tout topos $\Sh(\mathcal{C})$ pour lequel il existe un
ensemble $S \subset \Ob(\Sh(\mathcal{C}))$ satisfaisant à (1) -- (5).
En raisonnant exactement comme dans la preuve du
lemme \ref{sites-cohomology-lemma-colim-works-over-collection}, il suffit de montrer que, pour
toute colimite filtrante $\mathcal{I} = \colim \mathcal{I}_\lambda$ de
faisceaux abéliens injectifs $\mathcal{I}_\lambda$, on a
$H^{p_0 + 1}(\mathcal{C}, \mathcal{I}) = 0$.
Comme en (1), choisissons une surjection $K \to *$ avec $K \in S$.
Notons $K^n$ le produit de $n$ copies de $K$.
Considérons la suite spectrale
$$
E_1^{p, q} = H^q(K^{p + 1}, \mathcal{I}) \Rightarrow
H^{p + q}(*, \mathcal{I}) = H^{p + q}(\mathcal{C}, \mathcal{I})
$$
du lemme \ref{sites-cohomology-lemma-cech-to-cohomology-sheaf-sets}. Rappelons que
$H^q(K^{p + 1}, \mathcal{F}) = H^q(\mathcal{C}/K^{p + 1}, j^{-1}\mathcal{F})$,
pour tout faisceau abélien sur $\mathcal{C}$ ; voir le
lemme \ref{sites-cohomology-lemma-cohomology-on-sheaf-sets}. Nous avons
$j^{-1}\mathcal{I} = \colim j^{-1}\mathcal{I}_\lambda$
car $j^{-1}$ commute aux colimites. D'après le
lemme \ref{sites-cohomology-lemma-cohomology-of-open}, les restrictions
$j^{-1}\mathcal{I}_\lambda$ sont des faisceaux abéliens injectifs sur
$\mathcal{C}/K^{p + 1}$. Nous montrerons ci-dessous que l'hypothèse de
récurrence s'applique à $\mathcal{C}/K^{p + 1}$ ; par conséquent,
$H^q(K^{p + 1}, \mathcal{I}) = \colim H^q(K^{p + 1}, \mathcal{I}_\lambda) = 0$
pour $q < p_0 + 1$ (l'annulation provenant de l'injectivité de
$\mathcal{I}_\lambda$). Il s'ensuit que
$$
H^{p_0 + 1}(\mathcal{C}, \mathcal{I}) =
H^{p_0 + 1}\left(\ldots \to
H^0(K^{p_0}, \mathcal{I}) \to
H^0(K^{p_0 + 1}, \mathcal{I}) \to
H^0(K^{p_0 + 2}, \mathcal{I}) \to \ldots\right)
$$
En utilisant de nouveau l'hypothèse de récurrence, le complexe figurant au
membre de droite est la colimite, sur $\Lambda$, des complexes
$$
\ldots \to
H^0(K^{p_0}, \mathcal{I}_\lambda) \to
H^0(K^{p_0 + 1}, \mathcal{I}_\lambda) \to
H^0(K^{p_0 + 2}, \mathcal{I}_\lambda) \to \ldots
$$
Ces complexes sont exacts puisque $\mathcal{I}_\lambda$ est un faisceau
abélien injectif (cela découle par exemple de la suite spectrale). Comme les
colimites filtrantes sont exactes dans la catégorie des groupes abéliens, nous
obtenons l'annulation voulue.

\medskip\noindent
Il reste à montrer que l'hypothèse de récurrence s'applique au site
$\mathcal{C}/K^n$ pour tout $n \geq 1$. Rappelons que
$\Sh(\mathcal{C}/K^n) = \Sh(\mathcal{C})/K^n$ ; voir Sites,
lemme \ref{sites-lemma-localize-topos-site}.
Nous pouvons donc travailler dans $\Sh(\mathcal{C})/K^n$.
Notons $S_n \subset \Ob(\Sh(\mathcal{C}/K^n))$ l'ensemble des objets de la
forme $K' \to K^n$. Vérifions successivement chaque propriété :
\begin{enumerate}
\item Par (3) et récurrence, il existe une surjection $K' \to K^n$ avec
$K' \in S$. Alors $(K' \to K^n) \to (K^n \to K^n)$ est une surjection dans
$\Sh(\mathcal{C})/K^n$, et $K^n \to K^n$ est l'objet final de
$\Sh(\mathcal{C})/K^n$. Ainsi, (1) est vérifiée pour $S_n$.
\item La propriété (2) pour $S_n$ découle immédiatement de (2) pour $S$.
\item Soient $a : K_1 \to K^n$ et $b : K_2 \to K^n$ dans $S_n$.
Alors $(K_1 \to K^n) \times (K_2 \to K^n)$ est l'objet
$K_1 \times_{K^n} K_2 \to K^n$ de $\Sh(\mathcal{C})/K^n$.
Le sous-faisceau $K_1 \times_{K^n} K_2 \subset K_1 \times K_2$ est
l'égalisateur de $a \circ \text{pr}_1$ et $b \circ \text{pr}_2$.
Écrivons $a = (a_1, \ldots, a_n)$ et $b = (b_1, \ldots, b_n)$.
Choisissons une surjection $K_3 \to K_1 \times K_2$ avec $K_3 \in S$ ;
c'est possible par l'hypothèse (3) pour $\mathcal{C}$.
Choisissons une surjection
$$
K_4
\longrightarrow
\text{Equalizer}(K_3 \to K_1 \times K_2 \xrightarrow{a_1, b_1} K)
$$
avec $K_4 \in S$.
C'est possible par l'hypothèse (4) pour $\mathcal{C}$.
Choisissons une surjection
$$
K_5
\longrightarrow
\text{Equalizer}(K_4 \to K_1 \times K_2 \xrightarrow{a_2, b_2} K)
$$
avec $K_5 \in S$.
C'est encore possible. Poursuivons ainsi jusqu'à obtenir une surjection
$$
K_{3 + n}
\longrightarrow
\text{Equalizer}(K_{2 + n} \to K_1 \times K_2 \xrightarrow{a_n, b_n} K)
$$
avec $K_{3 + n} \in S$. Par construction,
$K_{3 + n} \to K_1 \times_{K^n} K_2$
est surjectif. Ainsi, $(K_{3 + n} \to K^n)$ appartient à $S_n$ et se
surjecte sur le produit $(K_1 \to K^n) \times (K_2 \to K^n)$.
La propriété (3) est donc vérifiée pour $S_n$.
\item La propriété (4) pour $S_n$ découle immédiatement de la propriété
(4) pour $S$.
\item La propriété (5) pour $S_n$ découle de la propriété (5) pour $S$.
En effet, un objet $\mathcal{F} \to K^n$ de $\Sh(\mathcal{C})/K^n$
correspond à un objet quasi-compact de $\Sh(\mathcal{C}/K^n)$ si et seulement
si $\mathcal{F}$ est un objet quasi-compact de $\Sh(\mathcal{C})$.
\end{enumerate}
Ceci achève la preuve du lemme.
\end{proof}

\begin{remark}
\label{sites-cohomology-remark-colim-global}
Soit $\mathcal{C}$ un site. Soit $\mathcal{B} \subset \Ob(\mathcal{C})$
un sous-ensemble. Soit $S \subset \Ob(\Sh(\mathcal{C}))$
l'ensemble des faisceaux $K$ de la forme
$$
K = \coprod\nolimits_{i = 1, \ldots, n} h_{U_i}^\#
$$
où $U_1, \ldots, U_n \in \mathcal{B}$. On peut alors se demander quand cet
ensemble satisfait aux hypothèses du lemme \ref{sites-cohomology-lemma-colim-global}. Une
condition suffisante est que
\begin{enumerate}
\item pour certains $n \geq 0$ et $U_1, \ldots, U_n \in \mathcal{B}$,
le morphisme $\coprod_{i = 1, \ldots, n} h_{U_i}^\# \to *$ soit surjectif ;
\item tout recouvrement de $U \in \mathcal{B}$ admette un raffinement par un
recouvrement de la forme $\{U_i \to U\}_{i = 1, \ldots, n}$ avec
$U_i \in \mathcal{B}$ ;
\item étant donnés $U, U' \in \mathcal{B}$, il existe $n \geq 0$,
$U_1, \ldots, U_n \in \mathcal{B}$ et des morphismes $U_i \to U$ et
$U_i \to U'$ tels que $\coprod_{i = 1, \ldots, n} h_{U_i}^\# \to
h_U^\# \times h_{U'}^\#$ soit surjectif ;
\item étant donnés des morphismes $a, b : U \to U'$ dans $\mathcal{C}$ avec
$U, U' \in \mathcal{B}$, il existe $U_1, \ldots, U_n \in \mathcal{B}$ et
des morphismes $U_i \to U$ qui égalisent $a, b$, tels que
$\coprod_{i = 1, \ldots, n} h_{U_i}^\# \to
\text{Equalizer}(h_a^\#, h_b^\# : h_U^\# \to h_{U'}^\#)$
soit surjectif.
\end{enumerate}
Nous omettons la vérification détaillée ; signalons seulement que la partie
(2) ci-dessus assure que tout élément de $\mathcal{B}$ est quasi-compact et
donc, d'après le lemme \ref{sites-lemma-coproduct-quasi-compact} de Sites,
que tout $K \in S$ est également quasi-compact.
\end{remark}

\begin{lemma}
\label{sites-cohomology-lemma-higher-direct-image-colimit}
Soit $f : \mathcal{C} \to \mathcal{D}$ un morphisme de sites correspondant
au foncteur continu $u : \mathcal{D} \to \mathcal{C}$.
Soit $(\mathcal{F}_i, \varphi_{ii'})$ un système de faisceaux abéliens sur
$\mathcal{C}$. Posons $\mathcal{F} = \colim \mathcal{F}_i$.
Soit $p \geq 0$ un entier. Notons $\mathcal{B}$ l'ensemble des
$V \in \Ob(\mathcal{D})$ tels que
$H^p(u(V), \mathcal{F}) = \colim H^p(u(V), \mathcal{F}_i)$.
Si tout objet de $\mathcal{D}$ possède un recouvrement par des éléments de
$\mathcal{B}$, alors $R^pf_*\mathcal{F} = \colim R^pf_*\mathcal{F}_i$.
\end{lemma}

\begin{proof}
Rappelons que $R^pf_*\mathcal{F}$ est le faisceautisé du préfaisceau
$\mathcal{G}$ qui envoie $V$ sur $H^p(u(V), \mathcal{F})$ ; voir le
lemme \ref{sites-cohomology-lemma-higher-direct-images}. De même,
$R^pf_*\mathcal{F}_i$ est le faisceautisé du préfaisceau
$\mathcal{G}_i$ qui envoie $V$ sur $H^p(u(V), \mathcal{F}_i)$.
Rappelons que la faisceautisation est adjointe à gauche de l'inclusion des
faisceaux dans les préfaisceaux ; voir Sites, section
\ref{sites-section-sheafification}. Elle commute donc aux colimites ; voir
Catégories, lemme \ref{categories-lemma-adjoint-exact}.
Il suffit par conséquent de montrer que le morphisme de préfaisceaux (la
colimite étant prise dans la catégorie des préfaisceaux)
$$
\colim \mathcal{G}_i \longrightarrow \mathcal{G}
$$
induit un isomorphisme sur les faisceautisés. Cela découle du
lemme \ref{sites-lemma-isomorphism-sheafifications} de Sites et de notre
hypothèse sur $\mathcal{B}$.
\end{proof}

\begin{lemma}
\label{sites-cohomology-lemma-colim-sites-injective}
Soit $\mathcal{I}$ une catégorie d'indices cofiltrante, et soit
$(\mathcal{C}_i, f_a)$ un système projectif de sites indexé par $\mathcal{I}$,
comme dans Sites, situation \ref{sites-situation-inverse-limit-sites}.
Posons $\mathcal{C} = \colim \mathcal{C}_i$ comme dans Sites, lemmes
\ref{sites-lemma-colimit-sites} et
\ref{sites-lemma-compute-pullback-to-limit}.
Supposons de plus donnés
\begin{enumerate}
\item un faisceau abélien $\mathcal{F}_i$ sur $\mathcal{C}_i$ pour tout
$i \in \Ob(\mathcal{I})$ ;
\item pour $a : j \to i$, un morphisme
$\varphi_a : f_a^{-1}\mathcal{F}_i \to \mathcal{F}_j$
de faisceaux abéliens sur $\mathcal{C}_j$,
\end{enumerate}
de sorte que $\varphi_c = \varphi_b \circ f_b^{-1}\varphi_a$ dès que
$c = a \circ b$. Il existe alors un morphisme de systèmes
$(\mathcal{F}_i, \varphi_a) \to (\mathcal{G}_i, \psi_a)$
tel que $\mathcal{F}_i \to \mathcal{G}_i$ soit injectif et que
$\mathcal{G}_i$ soit un faisceau abélien injectif.
\end{lemma}

\begin{proof}
Pour chaque $i$, choisissons une injection $\mathcal{F}_i \to \mathcal{A}_i$,
où $\mathcal{A}_i$ est un faisceau abélien injectif sur $\mathcal{C}_i$.
Nous pouvons alors considérer la famille de morphismes
$$
\gamma_i :
\mathcal{F}_i
\longrightarrow
\prod\nolimits_{b : k \to i} f_{b, *}\mathcal{A}_k = \mathcal{G}_i
$$
dont les composantes sont les morphismes adjoints aux morphismes
$f_b^{-1}\mathcal{F}_i \to \mathcal{F}_k \to \mathcal{A}_k$.
Pour $a : j \to i$ dans $\mathcal{I}$, il existe un morphisme canonique
$$
\psi_a : f_a^{-1}\mathcal{G}_i \to \mathcal{G}_j
$$
dont les composantes sont les morphismes canoniques
$f_b^{-1}f_{a \circ b, *}\mathcal{A}_k \to f_{b, *}\mathcal{A}_k$
pour $b : k \to j$. Nous obtenons ainsi une injection
$(\gamma_i) : (\mathcal{F}_i, \varphi_a) \to (\mathcal{G}_i, \psi_a)$
de systèmes de faisceaux abéliens. Remarquons que $\mathcal{G}_i$ est un
faisceau injectif de groupes abéliens sur $\mathcal{C}_i$ ; voir le
lemme \ref{sites-cohomology-lemma-pushforward-injective-flat} et Homologie,
lemme \ref{homology-lemma-product-injectives}.
Ceci achève la construction.
\end{proof}

\begin{lemma}
\label{sites-cohomology-lemma-colimit}
Dans la situation du lemme \ref{sites-cohomology-lemma-colim-sites-injective}, posons
$\mathcal{F} = \colim f_i^{-1}\mathcal{F}_i$.
Soient $i \in \Ob(\mathcal{I})$ et $X_i \in \text{Ob}(\mathcal{C}_i)$. Alors
$$
\colim_{a : j \to i} H^p(u_a(X_i), \mathcal{F}_j) =
H^p(u_i(X_i), \mathcal{F})
$$
pour tout $p \geq 0$.
\end{lemma}

\begin{proof}
Le cas $p = 0$ est le lemme \ref{sites-lemma-colimit} de Sites.

\medskip\noindent
Choisissons $(\mathcal{F}_i, \varphi_a) \to (\mathcal{G}_i, \psi_a)$
comme dans le lemme \ref{sites-cohomology-lemma-colim-sites-injective}.
En raisonnant exactement comme dans la preuve du
lemme \ref{sites-cohomology-lemma-colim-works-over-collection}, il suffit de montrer que
$H^p(X, \colim f_i^{-1}\mathcal{G}_i) = 0$ pour $p > 0$.

\medskip\noindent
Posons $\mathcal{G} = \colim f_i^{-1}\mathcal{G}_i$.
Pour montrer l'annulation de la cohomologie de $\mathcal{G}$ sur tout objet de
$\mathcal{C}$, nous montrons que la cohomologie de {\v C}ech de $\mathcal{G}$
est nulle pour tout recouvrement $\mathcal{U}$ de $\mathcal{C}$
(lemme \ref{sites-cohomology-lemma-cech-vanish-collection}).
Le recouvrement $\mathcal{U}$ provient d'un recouvrement
$\mathcal{U}_i$ de $\mathcal{C}_i$ pour un certain $i$. Nous avons
$$
\check{\mathcal{C}}^\bullet(\mathcal{U}, \mathcal{G}) =
\colim_{a : j \to i}
\check{\mathcal{C}}^\bullet(u_a(\mathcal{U}_i), \mathcal{G}_j)
$$
d'après le cas $p = 0$. Le membre de droite est acyclique en degrés positifs,
car il est une colimite filtrante de complexes acycliques, d'après le
lemme \ref{sites-cohomology-lemma-injective-trivial-cech}. Voir Algèbre,
lemme \ref{algebra-lemma-directed-colimit-exact}.
\end{proof}















\section{Résolutions plates}
\label{sites-cohomology-section-flat}

\noindent
Dans cette section, nous reprenons les arguments de Cohomologie,
section \ref{cohomology-section-flat}, dans le cadre des sites annelés et des
topos annelés.

\begin{lemma}
\label{sites-cohomology-lemma-derived-tor-exact}
Soit $(\mathcal{C}, \mathcal{O})$ un site annelé.
Soit $\mathcal{G}^\bullet$ un complexe de $\mathcal{O}$-modules.
Les foncteurs
$$
K(\textit{Mod}(\mathcal{O}))
\longrightarrow
K(\textit{Mod}(\mathcal{O})),
\quad
\mathcal{F}^\bullet \longmapsto
\text{Tot}(\mathcal{G}^\bullet \otimes_\mathcal{O} \mathcal{F}^\bullet)
$$
et
$$
K(\textit{Mod}(\mathcal{O}))
\longrightarrow
K(\textit{Mod}(\mathcal{O})),
\quad
\mathcal{F}^\bullet \longmapsto
\text{Tot}(\mathcal{F}^\bullet \otimes_\mathcal{O} \mathcal{G}^\bullet)
$$
sont des foncteurs exacts de catégories triangulées.
\end{lemma}

\begin{proof}
Cela résulte de Catégories dérivées, remarque
\ref{derived-remark-double-complex-as-tensor-product-of}.
\end{proof}

\begin{definition}
\label{sites-cohomology-definition-K-flat}
Soit $(\mathcal{C}, \mathcal{O})$ un site annelé.
Un complexe $\mathcal{K}^\bullet$ de $\mathcal{O}$-modules est dit
{\it K-plat} si, pour tout complexe acyclique $\mathcal{F}^\bullet$ de
$\mathcal{O}$-modules, le complexe
$$
\text{Tot}(\mathcal{F}^\bullet \otimes_\mathcal{O} \mathcal{K}^\bullet)
$$
est acyclique.
\end{definition}

\begin{lemma}
\label{sites-cohomology-lemma-K-flat-quasi-isomorphism}
Soit $(\mathcal{C}, \mathcal{O})$ un site annelé.
Soit $\mathcal{K}^\bullet$ un complexe K-plat.
Alors le foncteur
$$
K(\textit{Mod}(\mathcal{O}))
\longrightarrow
K(\textit{Mod}(\mathcal{O})), \quad
\mathcal{F}^\bullet
\longmapsto
\text{Tot}(\mathcal{F}^\bullet \otimes_\mathcal{O} \mathcal{K}^\bullet)
$$
transforme les quasi-isomorphismes en quasi-isomorphismes.
\end{lemma}

\begin{proof}
Cela résulte du lemme \ref{sites-cohomology-lemma-derived-tor-exact} et du fait que les
quasi-isomorphismes se caractérisent par l'acyclicité de leurs cônes.
\end{proof}

\begin{lemma}
\label{sites-cohomology-lemma-restriction-K-flat}
Soit $(\mathcal{C}, \mathcal{O})$ un site annelé.
Soit $U$ un objet de $\mathcal{C}$.
Si $\mathcal{K}^\bullet$ est un complexe K-plat de $\mathcal{O}$-modules,
alors $\mathcal{K}^\bullet|_U$ est un complexe K-plat de
$\mathcal{O}_U$-modules.
\end{lemma}

\begin{proof}
Soit $\mathcal{G}^\bullet$ un complexe exact de $\mathcal{O}_U$-modules.
Puisque $j_{U!}$ est exact (Modules sur les sites,
lemme \ref{sites-modules-lemma-extension-by-zero-exact}) et que
$\mathcal{K}^\bullet$ est un complexe K-plat de $\mathcal{O}$-modules, le
complexe
$$
j_{U!}(\text{Tot}(\mathcal{G}^\bullet \otimes_{\mathcal{O}_U}
\mathcal{K}^\bullet|_U)) =
\text{Tot}(j_{U!}\mathcal{G}^\bullet \otimes_\mathcal{O} \mathcal{K}^\bullet)
$$
est exact. Ici, l'égalité découle de Modules sur les sites,
lemme \ref{sites-modules-lemma-j-shriek-and-tensor}, et du fait que $j_{U!}$
commute aux sommes directes (en tant qu'adjoint à gauche).
Nous concluons puisque $j_{U!}$ reflète l'exactitude, d'après Modules sur les
sites, lemme \ref{sites-modules-lemma-j-shriek-reflects-exactness}.
\end{proof}

\begin{lemma}
\label{sites-cohomology-lemma-tensor-product-K-flat}
Soit $(\mathcal{C}, \mathcal{O})$ un site annelé.
Si $\mathcal{K}^\bullet$ et $\mathcal{L}^\bullet$ sont des complexes K-plats
de $\mathcal{O}$-modules, alors
$\text{Tot}(\mathcal{K}^\bullet \otimes_\mathcal{O} \mathcal{L}^\bullet)$
est un complexe K-plat de $\mathcal{O}$-modules.
\end{lemma}

\begin{proof}
Cela résulte de l'isomorphisme
$$
\text{Tot}(\mathcal{M}^\bullet \otimes_\mathcal{O}
\text{Tot}(\mathcal{K}^\bullet \otimes_\mathcal{O} \mathcal{L}^\bullet))
=
\text{Tot}(\text{Tot}(\mathcal{M}^\bullet \otimes_\mathcal{O}
\mathcal{K}^\bullet) \otimes_\mathcal{O} \mathcal{L}^\bullet)
$$
et de la définition.
\end{proof}

\begin{lemma}
\label{sites-cohomology-lemma-K-flat-two-out-of-three}
Soit $(\mathcal{C}, \mathcal{O})$ un site annelé.
Soit $(\mathcal{K}_1^\bullet, \mathcal{K}_2^\bullet, \mathcal{K}_3^\bullet)$
un triangle distingué dans $K(\textit{Mod}(\mathcal{O}))$.
Si deux des trois complexes $\mathcal{K}_i^\bullet$ sont K-plats, le
troisième l'est aussi.
\end{lemma}

\begin{proof}
Cela résulte du lemme \ref{sites-cohomology-lemma-derived-tor-exact} et du fait que, dans un
triangle distingué de $K(\textit{Mod}(\mathcal{O}))$, si deux des trois objets
sont acycliques, le troisième l'est aussi.
\end{proof}

\begin{lemma}
\label{sites-cohomology-lemma-K-flat-two-out-of-three-ses}
Soit $(\mathcal{C}, \mathcal{O})$ un site annelé. Soit
$0 \to \mathcal{K}_1^\bullet \to \mathcal{K}_2^\bullet \to
\mathcal{K}_3^\bullet \to 0$ une suite exacte courte de complexes telle que
les termes de $\mathcal{K}_3^\bullet$ soient des $\mathcal{O}$-modules plats.
Si deux des trois complexes $\mathcal{K}_i^\bullet$ sont K-plats, le troisième
l'est aussi.
\end{lemma}

\begin{proof}
D'après Modules sur les sites,
lemme \ref{sites-modules-lemma-flat-tor-zero}, pour tout complexe
$\mathcal{L}^\bullet$, nous obtenons une suite exacte courte
$$
0 \to
\text{Tot}(\mathcal{L}^\bullet \otimes_\mathcal{O} \mathcal{K}_1^\bullet) \to
\text{Tot}(\mathcal{L}^\bullet \otimes_\mathcal{O} \mathcal{K}_2^\bullet) \to
\text{Tot}(\mathcal{L}^\bullet \otimes_\mathcal{O} \mathcal{K}_3^\bullet) \to 0
$$
de complexes. Le lemme résulte donc de la suite exacte longue des faisceaux de
cohomologie et de la définition des complexes K-plats.
\end{proof}

\begin{lemma}
\label{sites-cohomology-lemma-bounded-flat-K-flat}
Soit $(\mathcal{C}, \mathcal{O})$ un site annelé. Tout complexe borné
supérieurement de $\mathcal{O}$-modules plats est K-plat.
\end{lemma}

\begin{proof}
Soit $\mathcal{K}^\bullet$ un complexe borné supérieurement de
$\mathcal{O}$-modules plats. Soit $\mathcal{L}^\bullet$ un complexe acyclique
de $\mathcal{O}$-modules. Remarquons que
$\mathcal{L}^\bullet = \colim_m \tau_{\leq m}\mathcal{L}^\bullet$
où les colimites sont prises terme à terme. Ainsi,
$$
\text{Tot}(\mathcal{K}^\bullet \otimes_\mathcal{O} \mathcal{L}^\bullet)
=
\colim_m \text{Tot}(
\mathcal{K}^\bullet \otimes_\mathcal{O} \tau_{\leq m}\mathcal{L}^\bullet)
$$
terme à terme. Pour prouver que le complexe de gauche est acyclique, il suffit
donc de montrer que chacun des complexes de droite l'est. Comme
$\tau_{\leq m}\mathcal{L}^\bullet$ est acyclique, nous sommes ramenés au cas
où $\mathcal{L}^\bullet$ est borné supérieurement.
Dans ce cas, la suite spectrale du
lemme \ref{homology-lemma-first-quadrant-ss} d'Homologie vérifie
$$
{}'E_1^{p, q} = H^p(\mathcal{L}^\bullet \otimes_\mathcal{O} \mathcal{K}^q)
$$
qui est nul puisque $\mathcal{K}^q$ est plat et $\mathcal{L}^\bullet$
acyclique. Ceci conclut la preuve.
\end{proof}

\begin{lemma}
\label{sites-cohomology-lemma-colimit-K-flat}
Soit $(\mathcal{C}, \mathcal{O})$ un site annelé.
Soit $\mathcal{K}_1^\bullet \to \mathcal{K}_2^\bullet \to \ldots$
un système de complexes K-plats.
Alors $\colim_i \mathcal{K}_i^\bullet$ est K-plat.
\end{lemma}

\begin{proof}
Puisque nous prenons les colimites terme à terme, il est clair que
$$
\colim_i \text{Tot}(
\mathcal{F}^\bullet \otimes_\mathcal{O} \mathcal{K}_i^\bullet)
=
\text{Tot}(\mathcal{F}^\bullet \otimes_\mathcal{O}
\colim_i \mathcal{K}_i^\bullet)
$$
Le lemme découle donc de l'exactitude des colimites filtrantes.
\end{proof}

\begin{lemma}
\label{sites-cohomology-lemma-resolution-by-direct-sums-extensions-by-zero}
Soit $(\mathcal{C}, \mathcal{O})$ un site annelé.
Pour tout complexe $\mathcal{G}^\bullet$ de $\mathcal{O}$-modules, il existe
un diagramme commutatif de complexes de $\mathcal{O}$-modules
$$
\xymatrix{
\mathcal{K}_1^\bullet \ar[d] \ar[r] &
\mathcal{K}_2^\bullet \ar[d] \ar[r] & \ldots \\
\tau_{\leq 1}\mathcal{G}^\bullet \ar[r] &
\tau_{\leq 2}\mathcal{G}^\bullet \ar[r] & \ldots
}
$$
ayant les propriétés suivantes : (1) les flèches verticales sont des
quasi-isomorphismes et sont surjectives terme à terme ;
(2) chaque $\mathcal{K}_n^\bullet$ est un complexe borné supérieurement dont
les termes sont des sommes directes de $\mathcal{O}$-modules de la forme
$j_{U!}\mathcal{O}_U$ ; et (3) les morphismes
$\mathcal{K}_n^\bullet \to \mathcal{K}_{n + 1}^\bullet$ sont des injections
scindées terme à terme dont les conoyaux sont des sommes directes de
$\mathcal{O}$-modules de la forme $j_{U!}\mathcal{O}_U$. De plus, le morphisme
$\colim \mathcal{K}_n^\bullet \to \mathcal{G}^\bullet$ est un quasi-isomorphisme.
\end{lemma}

\begin{proof}
L'existence du diagramme et les propriétés (1), (2), (3) découlent
immédiatement de Modules sur les sites,
lemme \ref{sites-modules-lemma-module-quotient-flat}, et de Catégories
dérivées, lemme \ref{derived-lemma-special-direct-system}.
Le morphisme induit
$\colim \mathcal{K}_n^\bullet \to \mathcal{G}^\bullet$
est un quasi-isomorphisme parce que les colimites filtrantes sont exactes.
\end{proof}

\begin{lemma}
\label{sites-cohomology-lemma-K-flat-resolution}
Soit $(\mathcal{C}, \mathcal{O})$ un site annelé. Pour tout complexe
$\mathcal{G}^\bullet$, il existe un complexe $K$-plat $\mathcal{K}^\bullet$
dont les termes sont des $\mathcal{O}$-modules plats, ainsi qu'un
quasi-isomorphisme $\mathcal{K}^\bullet \to \mathcal{G}^\bullet$ surjectif
terme à terme.
\end{lemma}

\begin{proof}
Choisissons un diagramme comme dans le
lemme \ref{sites-cohomology-lemma-resolution-by-direct-sums-extensions-by-zero}.
Chaque complexe $\mathcal{K}_n^\bullet$ est un complexe borné supérieurement
de modules plats ; voir Modules sur les sites,
lemme \ref{sites-modules-lemma-j-shriek-flat}.
Ainsi, $\mathcal{K}_n^\bullet$ est K-plat d'après le
lemme \ref{sites-cohomology-lemma-bounded-flat-K-flat}.
Par conséquent, $\colim \mathcal{K}_n^\bullet$ est K-plat d'après le
lemme \ref{sites-cohomology-lemma-colimit-K-flat}.
Le morphisme induit
$\colim \mathcal{K}_n^\bullet \to \mathcal{G}^\bullet$
est, par construction, un quasi-isomorphisme surjectif terme à terme.
La propriété (3) du
lemme \ref{sites-cohomology-lemma-resolution-by-direct-sums-extensions-by-zero} montre que
$\colim \mathcal{K}_n^m$ est une somme directe de modules plats, donc est plat,
ce qui prouve la dernière assertion.
\end{proof}

\begin{lemma}
\label{sites-cohomology-lemma-derived-tor-quasi-isomorphism-other-side}
Soit $(\mathcal{C}, \mathcal{O})$ un site annelé. Soit
$\alpha : \mathcal{P}^\bullet \to \mathcal{Q}^\bullet$ un quasi-isomorphisme
de complexes K-plats de $\mathcal{O}$-modules.
Pour tout complexe $\mathcal{F}^\bullet$ de $\mathcal{O}$-modules, le
morphisme induit
$$
\text{Tot}(\text{id}_{\mathcal{F}^\bullet} \otimes \alpha) :
\text{Tot}(\mathcal{F}^\bullet \otimes_\mathcal{O} \mathcal{P}^\bullet)
\longrightarrow
\text{Tot}(\mathcal{F}^\bullet \otimes_\mathcal{O} \mathcal{Q}^\bullet)
$$
est un quasi-isomorphisme.
\end{lemma}

\begin{proof}
Choisissons un quasi-isomorphisme $\mathcal{K}^\bullet \to \mathcal{F}^\bullet$,
où $\mathcal{K}^\bullet$ est un complexe K-plat ; voir le
lemme \ref{sites-cohomology-lemma-K-flat-resolution}.
Considérons le diagramme commutatif
$$
\xymatrix{
\text{Tot}(\mathcal{K}^\bullet
\otimes_\mathcal{O} \mathcal{P}^\bullet) \ar[r] \ar[d] &
\text{Tot}(\mathcal{K}^\bullet
\otimes_\mathcal{O} \mathcal{Q}^\bullet) \ar[d] \\
\text{Tot}(\mathcal{F}^\bullet
\otimes_\mathcal{O} \mathcal{P}^\bullet) \ar[r] &
\text{Tot}(\mathcal{F}^\bullet
\otimes_\mathcal{O} \mathcal{Q}^\bullet)
}
$$
Le résultat découle du lemme \ref{sites-cohomology-lemma-K-flat-quasi-isomorphism}, puisque les
flèches verticales et la flèche horizontale supérieure sont des
quasi-isomorphismes.
\end{proof}

\noindent
Soit $(\mathcal{C}, \mathcal{O})$ un site annelé.
Soit $\mathcal{F}^\bullet$ un objet de $D(\mathcal{O})$.
Choisissons une résolution K-plate
$\mathcal{K}^\bullet \to \mathcal{F}^\bullet$ ; voir le
lemme \ref{sites-cohomology-lemma-K-flat-resolution}.
D'après le lemme \ref{sites-cohomology-lemma-derived-tor-exact}, nous obtenons un foncteur exact
de catégories triangulées
$$
K(\mathcal{O})
\longrightarrow
K(\mathcal{O}),
\quad
\mathcal{G}^\bullet
\longmapsto
\text{Tot}(\mathcal{G}^\bullet \otimes_\mathcal{O} \mathcal{K}^\bullet)
$$
D'après le lemme \ref{sites-cohomology-lemma-K-flat-quasi-isomorphism}, ce foncteur induit un
foncteur $D(\mathcal{O}) \to D(\mathcal{O})$, puisque $D(\mathcal{O})$ est la
localisation de $K(\mathcal{O})$ relativement aux quasi-isomorphismes.
Comme la catégorie des résolutions $K$-plates de $\mathcal{F}^\bullet$ est
cofiltrante et compte tenu du
lemme \ref{sites-cohomology-lemma-derived-tor-quasi-isomorphism-other-side}, le foncteur obtenu,
à isomorphisme près, ne dépend pas du choix de $\mathcal{K}^\bullet$.

\begin{definition}
\label{sites-cohomology-definition-derived-tor}
Soit $(\mathcal{C}, \mathcal{O})$ un site annelé.
Soit $\mathcal{F}^\bullet$ un objet de $D(\mathcal{O})$.
Le {\it produit tensoriel dérivé}
$$
- \otimes_\mathcal{O}^{\mathbf{L}} \mathcal{F}^\bullet :
D(\mathcal{O})
\longrightarrow
D(\mathcal{O})
$$
est le foncteur exact de catégories triangulées décrit ci-dessus.
\end{definition}

\noindent
Il ressort de nos constructions explicites qu'il existe un isomorphisme
canonique
$$
\mathcal{F}^\bullet \otimes_\mathcal{O}^{\mathbf{L}} \mathcal{G}^\bullet
\cong
\mathcal{G}^\bullet \otimes_\mathcal{O}^{\mathbf{L}} \mathcal{F}^\bullet
$$
pour $\mathcal{G}^\bullet$ et $\mathcal{F}^\bullet$ dans $D(\mathcal{O})$.
Nous utilisons ici les règles de signes données dans Compléments sur l'algèbre,
section \ref{more-algebra-section-sign-rules}. Ainsi, lorsque nous écrivons
$\mathcal{F}^\bullet \otimes_\mathcal{O}^{\mathbf{L}} \mathcal{G}^\bullet$
nous ne préciserons généralement pas par rapport à laquelle des deux variables
nous définissons le produit tensoriel dérivé.

\begin{definition}
\label{sites-cohomology-definition-tor}
Soit $(\mathcal{C}, \mathcal{O})$ un site annelé.
Soient $\mathcal{F}$ et $\mathcal{G}$ des $\mathcal{O}$-modules.
Les {\it Tor} de $\mathcal{F}$ et $\mathcal{G}$ sont définis par la formule
$$
\text{Tor}_p^\mathcal{O}(\mathcal{F}, \mathcal{G}) =
H^{-p}(\mathcal{F} \otimes_\mathcal{O}^\mathbf{L} \mathcal{G})
$$
où le produit tensoriel dérivé est celui défini ci-dessus.
\end{definition}

\noindent
Cette définition implique que, pour toute suite exacte courte de
$\mathcal{O}$-modules
$0 \to \mathcal{F}_1 \to \mathcal{F}_2 \to \mathcal{F}_3 \to 0$
nous avons une suite exacte longue de cohomologie
$$
\xymatrix{
\mathcal{F}_1 \otimes_\mathcal{O} \mathcal{G} \ar[r] &
\mathcal{F}_2 \otimes_\mathcal{O} \mathcal{G} \ar[r] &
\mathcal{F}_3 \otimes_\mathcal{O} \mathcal{G} \ar[r] & 0 \\
\text{Tor}_1^\mathcal{O}(\mathcal{F}_1, \mathcal{G}) \ar[r] &
\text{Tor}_1^\mathcal{O}(\mathcal{F}_2, \mathcal{G}) \ar[r] &
\text{Tor}_1^\mathcal{O}(\mathcal{F}_3, \mathcal{G}) \ar[ull]
}
$$
pour tout $\mathcal{O}$-module $\mathcal{G}$. On l'appelle la suite exacte
longue de $\text{Tor}$ associée à la situation.

\begin{lemma}
\label{sites-cohomology-lemma-flat-tor-zero}
Soit $(\mathcal{C}, \mathcal{O})$ un site annelé.
Soit $\mathcal{F}$ un $\mathcal{O}$-module.
Les conditions suivantes sont équivalentes :
\begin{enumerate}
\item $\mathcal{F}$ est un $\mathcal{O}$-module plat ;
\item $\text{Tor}_1^\mathcal{O}(\mathcal{F}, \mathcal{G}) = 0$
pour tout $\mathcal{O}$-module $\mathcal{G}$.
\end{enumerate}
\end{lemma}

\begin{proof}
Si $\mathcal{F}$ est plat, alors $\mathcal{F} \otimes_\mathcal{O} -$ est un
foncteur exact et ses foncteurs dérivés s'annulent. Réciproquement, supposons
(2). Si $\mathcal{G} \to \mathcal{H}$ est injectif, de conoyau $\mathcal{Q}$,
la suite exacte longue de $\text{Tor}$ montre que le noyau de
$\mathcal{F} \otimes_\mathcal{O} \mathcal{G} \to
\mathcal{F} \otimes_\mathcal{O} \mathcal{H}$
est un quotient de
$\text{Tor}_1^\mathcal{O}(\mathcal{F}, \mathcal{Q})$
qui est nul par hypothèse. Par conséquent, $\mathcal{F}$ est plat.
\end{proof}

\begin{lemma}
\label{sites-cohomology-lemma-K-flat-flat-acyclic}
Soit $(\mathcal{C}, \mathcal{O})$ un site annelé. Soit $\mathcal{K}^\bullet$
un complexe K-plat, acyclique et à termes plats. Alors
$\mathcal{F} = \Ker(\mathcal{K}^n \to \mathcal{K}^{n + 1})$
est un $\mathcal{O}$-module plat.
\end{lemma}

\begin{proof}
Observons que
$$
\ldots \to \mathcal{K}^{n - 2} \to \mathcal{K}^{n - 1} \to
\mathcal{F} \to 0
$$
est une résolution plate de notre module $\mathcal{F}$. Comme un complexe
borné supérieurement de modules plats est K-plat
(lemme \ref{sites-cohomology-lemma-bounded-flat-K-flat}), nous pouvons utiliser cette résolution
pour calculer $\text{Tor}_i(\mathcal{F}, \mathcal{G})$ pour tout
$\mathcal{O}$-module $\mathcal{G}$. D'une part,
$\mathcal{K}^\bullet \otimes_\mathcal{O}^\mathbf{L} \mathcal{G}$ est nul dans
$D(\mathcal{O})$ parce que $\mathcal{K}^\bullet$ est acyclique et, d'autre
part, il est représenté par
$\mathcal{K}^\bullet \otimes_\mathcal{O} \mathcal{G}$.
Nous voyons donc que
$$
\mathcal{K}^{n - 3} \otimes_\mathcal{O} \mathcal{G} \to
\mathcal{K}^{n - 2} \otimes_\mathcal{O} \mathcal{G} \to
\mathcal{K}^{n - 1} \otimes_\mathcal{O} \mathcal{G}
$$
est exacte. Ainsi, $\text{Tor}_1(\mathcal{F}, \mathcal{G}) = 0$, et nous
concluons par le lemme \ref{sites-cohomology-lemma-flat-tor-zero}.
\end{proof}

\begin{lemma}
\label{sites-cohomology-lemma-factor-through-K-flat}
Soit $(\mathcal{C}, \mathcal{O})$ un site annelé.
Soit $a : \mathcal{K}^\bullet \to \mathcal{L}^\bullet$ un morphisme de
complexes de $\mathcal{O}$-modules. Si $\mathcal{K}^\bullet$ est K-plat, il
existe un complexe $\mathcal{N}^\bullet$ et des morphismes de complexes
$b : \mathcal{K}^\bullet \to \mathcal{N}^\bullet$
et $c : \mathcal{N}^\bullet \to \mathcal{L}^\bullet$ tels que
\begin{enumerate}
\item $\mathcal{N}^\bullet$ est K-plat ;
\item $c$ est un quasi-isomorphisme ;
\item $a$ est homotope à $c \circ b$.
\end{enumerate}
Si les termes de $\mathcal{K}^\bullet$ sont plats, nous pouvons choisir
$\mathcal{N}^\bullet$, $b$ et $c$ de sorte qu'il en soit de même pour
$\mathcal{N}^\bullet$.
\end{lemma}

\begin{proof}
Nous utiliserons le fait que la catégorie d'homotopie
$K(\textit{Mod}(\mathcal{O}))$ est une catégorie triangulée ; voir Catégories
dérivées, proposition
\ref{derived-proposition-homotopy-category-triangulated}.
Choisissons un triangle distingué
$\mathcal{K}^\bullet \to \mathcal{L}^\bullet \to
\mathcal{C}^\bullet \to \mathcal{K}^\bullet[1]$.
Choisissons un quasi-isomorphisme
$\mathcal{M}^\bullet \to \mathcal{C}^\bullet$, où $\mathcal{M}^\bullet$ est
K-plat à termes plats ; voir le lemme \ref{sites-cohomology-lemma-K-flat-resolution}.
Par les axiomes des catégories triangulées, nous pouvons inscrire la composée
$\mathcal{M}^\bullet \to \mathcal{C}^\bullet \to \mathcal{K}^\bullet[1]$
dans un triangle distingué
$\mathcal{K}^\bullet \to \mathcal{N}^\bullet \to
\mathcal{M}^\bullet \to \mathcal{K}^\bullet[1]$.
D'après le lemme \ref{sites-cohomology-lemma-K-flat-two-out-of-three}, $\mathcal{N}^\bullet$ est
K-plat. En utilisant de nouveau les axiomes des catégories triangulées, nous
pouvons choisir un morphisme $\mathcal{N}^\bullet \to \mathcal{L}^\bullet$
qui s'insère dans le morphisme suivant de triangles distingués
$$
\xymatrix{
\mathcal{K}^\bullet \ar[r] \ar[d] &
\mathcal{N}^\bullet \ar[r] \ar[d] &
\mathcal{M}^\bullet \ar[r] \ar[d] &
\mathcal{K}^\bullet[1] \ar[d] \\
\mathcal{K}^\bullet \ar[r] &
\mathcal{L}^\bullet \ar[r] &
\mathcal{C}^\bullet \ar[r] &
\mathcal{K}^\bullet[1]
}
$$
Puisque deux des trois flèches sont des quasi-isomorphismes, la troisième
$\mathcal{N}^\bullet \to \mathcal{L}^\bullet$ l'est aussi, d'après les suites
exactes longues de cohomologie associées à ces triangles distingués (on peut
aussi considérer l'image de ce diagramme dans $D(\mathcal{O})$ et utiliser
Catégories dérivées, lemme \ref{derived-lemma-third-isomorphism-triangle}).
Ceci achève la preuve de (1), (2) et (3).
Pour prouver la dernière assertion, nous pouvons choisir
$\mathcal{N}^\bullet$ de sorte que
$\mathcal{N}^n \cong \mathcal{M}^n \oplus \mathcal{K}^n$ ; voir
Catégories dérivées, lemme
\ref{derived-lemma-improve-distinguished-triangle-homotopy}.
Nous obtenons ainsi la planéité voulue si les termes de
$\mathcal{K}^\bullet$ sont plats.
\end{proof}
\section{Image inverse dérivée}
\label{sites-cohomology-section-derived-pullback}

\noindent
Soit
$f : (\Sh(\mathcal{C}), \mathcal{O}) \to
(\Sh(\mathcal{C}'), \mathcal{O}')$
un morphisme de topos annelés. Nous pouvons utiliser des résolutions
K-plates pour définir un foncteur d'image inverse dérivée
$$
Lf^* : D(\mathcal{O}') \to D(\mathcal{O})
$$

\begin{lemma}
\label{sites-cohomology-lemma-pullback-K-flat}
Soit $f : (\Sh(\mathcal{C}'), \mathcal{O}') \to (\Sh(\mathcal{C}), \mathcal{O})$
un morphisme de topos annelés. Soit $\mathcal{K}^\bullet$ un complexe K-plat
de $\mathcal{O}$-modules dont les termes sont des $\mathcal{O}$-modules plats.
Alors $f^*\mathcal{K}^\bullet$ est un complexe K-plat de
$\mathcal{O}'$-modules dont les termes sont des $\mathcal{O}'$-modules plats.
\end{lemma}

\begin{proof}
Les termes $f^*\mathcal{K}^n$ sont des $\mathcal{O}'$-modules plats d'après
Modules sur les sites, Lemme \ref{sites-modules-lemma-pullback-flat}.
Choisissons un diagramme
$$
\xymatrix{
\mathcal{K}_1^\bullet \ar[d] \ar[r] &
\mathcal{K}_2^\bullet \ar[d] \ar[r] & \ldots \\
\tau_{\leq 1}\mathcal{K}^\bullet \ar[r] &
\tau_{\leq 2}\mathcal{K}^\bullet \ar[r] & \ldots
}
$$
comme dans le Lemme \ref{sites-cohomology-lemma-resolution-by-direct-sums-extensions-by-zero}.
Nous utiliserons sans autre mention toutes les propriétés énoncées dans ce
lemme. Chaque $\mathcal{K}_n^\bullet$ est un complexe borné supérieurement
de modules plats, voir
Modules sur les sites, Lemme \ref{sites-modules-lemma-j-shriek-flat}.
Considérons la suite exacte courte de complexes
$$
0 \to \mathcal{M}^\bullet \to
\colim \mathcal{K}_n^\bullet \to
\mathcal{K}^\bullet \to 0
$$
qui définit $\mathcal{M}^\bullet$. D'après les Lemmes
\ref{sites-cohomology-lemma-bounded-flat-K-flat} et \ref{sites-cohomology-lemma-colimit-K-flat}, le complexe
$\colim \mathcal{K}_n^\bullet$ est K-plat et, d'après
Modules sur les sites, Lemme \ref{sites-modules-lemma-colimits-flat},
ses termes sont plats. D'après Modules sur les sites, Lemme
\ref{sites-modules-lemma-flat-ses}, les termes de $\mathcal{M}^\bullet$
sont plats ; d'après le Lemme \ref{sites-cohomology-lemma-K-flat-two-out-of-three-ses},
$\mathcal{M}^\bullet$ est K-plat ; enfin, d'après la suite exacte longue
de cohomologie, $\mathcal{M}^\bullet$ est acyclique (car la seconde flèche
est un quasi-isomorphisme). L'image inverse
$f^*(\colim \mathcal{K}_n^\bullet) = \colim f^*\mathcal{K}_n^\bullet$
est une colimite de complexes bornés supérieurement de
$\mathcal{O}'$-modules plats, donc elle est K-plate (d'après les mêmes
lemmes que ci-dessus). L'image inverse de notre suite exacte courte
$$
0 \to f^*\mathcal{M}^\bullet \to
f^*(\colim \mathcal{K}_n^\bullet) \to
f^*\mathcal{K}^\bullet \to 0
$$
est une suite exacte courte de complexes d'après
Modules sur les sites, Lemme \ref{sites-modules-lemma-pullback-ses}.
Ainsi, d'après le Lemme \ref{sites-cohomology-lemma-K-flat-two-out-of-three-ses}, il suffit
de montrer que $f^*\mathcal{M}^\bullet$ est K-plat. Cela nous ramène au cas
traité dans le paragraphe suivant.

\medskip\noindent
Supposons que $\mathcal{K}^\bullet$ soit acyclique et K-plat, et que ses
termes soient plats. Le Lemme \ref{sites-cohomology-lemma-K-flat-flat-acyclic} assure alors
que tous les termes de $\tau_{\leq n}\mathcal{K}^\bullet$ sont des
$\mathcal{O}$-modules plats. Choisissons un diagramme comme ci-dessus et
utilisons toutes les propriétés démontrées plus haut pour ce diagramme.
Notons $\mathcal{M}_n^\bullet$ le noyau du morphisme de complexes
$\mathcal{K}_n^\bullet \to \tau_{\leq n}\mathcal{K}^\bullet$
de sorte que nous disposons de suites exactes courtes de complexes
$$
0 \to \mathcal{M}_n^\bullet \to \mathcal{K}_n^\bullet \to
\tau_{\leq n}\mathcal{K}^\bullet \to 0
$$
D'après Modules sur les sites, Lemme \ref{sites-modules-lemma-flat-ses},
les termes du complexe $\mathcal{M}_n^\bullet$ sont plats. Ainsi, dans ce
cas, $\mathcal{M} = \colim \mathcal{M}_n^\bullet$ est une colimite filtrante
de complexes bornés supérieurement de modules plats. Par conséquent,
$f^*\mathcal{M}^\bullet$ est K-plat (par le même argument que ci-dessus),
ce qui conclut.
\end{proof}

\begin{lemma}
\label{sites-cohomology-lemma-derived-base-change}
Soit $f : (\Sh(\mathcal{C}), \mathcal{O}) \to (\Sh(\mathcal{C}'), \mathcal{O}')$
un morphisme de topos annelés. Il existe un foncteur exact
$$
Lf^* : D(\mathcal{O}') \longrightarrow D(\mathcal{O})
$$
entre catégories triangulées tel que
$Lf^*\mathcal{K}^\bullet = f^*\mathcal{K}^\bullet$ pour tout complexe
K-plat $\mathcal{K}^\bullet$ à termes plats et, en particulier, pour tout
complexe borné supérieurement de $\mathcal{O}'$-modules plats.
\end{lemma}

\begin{proof}
Pour le voir, utilisons la théorie générale développée dans
Catégories dérivées, section \ref{derived-section-derived-functors}.
Posons $\mathcal{D} = K(\mathcal{O}')$ et $\mathcal{D}' = D(\mathcal{O})$.
Notons $F : \mathcal{D} \to \mathcal{D}'$ le foncteur exact de catégories
triangulées défini par la règle
$F(\mathcal{G}^\bullet) = f^*\mathcal{G}^\bullet$.
Soit $S$ l'ensemble des quasi-isomorphismes de
$\mathcal{D} = K(\mathcal{O}')$. Nous obtenons ainsi une situation comme
celle de Catégories dérivées, Situation
\ref{derived-situation-derived-functor}, si bien que Catégories dérivées,
Définition \ref{derived-definition-right-derived-functor-defined},
s'applique. Nous affirmons que $LF$ est défini partout. Cela résulte de
Catégories dérivées, Lemme \ref{derived-lemma-find-existence-computes},
en prenant pour $\mathcal{P} \subset \Ob(\mathcal{D})$ la collection des
complexes K-plats $\mathcal{K}^\bullet$ à termes plats. En effet, (1)
résulte du Lemme \ref{sites-cohomology-lemma-K-flat-resolution} et, pour établir (2), nous
devons montrer que, pour tout quasi-isomorphisme
$\mathcal{K}_1^\bullet  \to \mathcal{K}_2^\bullet$ entre objets de
$\mathcal{P}$, le morphisme
$f^*\mathcal{K}_1^\bullet  \to f^*\mathcal{K}_2^\bullet$ est un
quasi-isomorphisme. Écrivons ce morphisme sous la forme
$$
f^{-1}\mathcal{K}_1^\bullet \otimes_{f^{-1}\mathcal{O}'} \mathcal{O}
\longrightarrow
f^{-1}\mathcal{K}_2^\bullet \otimes_{f^{-1}\mathcal{O}'} \mathcal{O}
$$
Le foncteur $f^{-1}$ est exact ; le morphisme
$f^{-1}\mathcal{K}_1^\bullet  \to f^{-1}\mathcal{K}_2^\bullet$ est donc un
quasi-isomorphisme. Les complexes
 $f^{-1}\mathcal{K}_1^\bullet$ et $f^{-1}\mathcal{K}_2^\bullet$
de $f^{-1}\mathcal{O}'$-modules sont K-plats d'après le Lemme
\ref{sites-cohomology-lemma-pullback-K-flat}, car nous pouvons considérer le morphisme de
topos annelés
$(\Sh(\mathcal{C}), f^{-1}\mathcal{O}') \to
(\Sh(\mathcal{C}'), \mathcal{O}')$. Le Lemme
\ref{sites-cohomology-lemma-derived-tor-quasi-isomorphism-other-side} assure donc que le
morphisme affiché est un quasi-isomorphisme. Nous obtenons ainsi un
foncteur dérivé
$$
LF :
D(\mathcal{O}') = S^{-1}\mathcal{D}
\longrightarrow
\mathcal{D}' = D(\mathcal{O})
$$
voir Catégories dérivées, Équation (\ref{derived-equation-everywhere}).
Enfin, Catégories dérivées, Lemme
\ref{derived-lemma-find-existence-computes}, assure également que
$LF(\mathcal{K}^\bullet) = F(\mathcal{K}^\bullet) = f^*\mathcal{K}^\bullet$
lorsque $\mathcal{K}^\bullet$ appartient à $\mathcal{P}$.
La démonstration s'achève en observant que les complexes bornés
supérieurement de modules plats appartiennent à $\mathcal{P}$ d'après le
Lemme \ref{sites-cohomology-lemma-bounded-flat-K-flat}.
\end{proof}

\begin{lemma}
\label{sites-cohomology-lemma-derived-pullback-composition}
Considérons des morphismes de topos annelés
$f : (\Sh(\mathcal{C}), \mathcal{O}_\mathcal{C}) \to
(\Sh(\mathcal{D}), \mathcal{O}_\mathcal{D})$
et
$g : (\Sh(\mathcal{D}), \mathcal{O}_\mathcal{D}) \to
(\Sh(\mathcal{E}), \mathcal{O}_\mathcal{E})$.
Alors $Lf^* \circ Lg^* = L(g \circ f)^*$ comme foncteurs
$D(\mathcal{O}_\mathcal{E}) \to D(\mathcal{O}_\mathcal{C})$.
\end{lemma}

\begin{proof}
Soit $E$ un objet de $D(\mathcal{O}_\mathcal{E})$. Nous pouvons représenter
$E$ par un complexe K-plat $\mathcal{K}^\bullet$ à termes plats, voir le
Lemme \ref{sites-cohomology-lemma-K-flat-resolution}. Par construction, $Lg^*E$ se calcule
par $g^*\mathcal{K}^\bullet$, voir le Lemme
\ref{sites-cohomology-lemma-derived-base-change}. D'après le Lemme
\ref{sites-cohomology-lemma-pullback-K-flat}, le complexe $g^*\mathcal{K}^\bullet$ est
K-plat à termes plats. Ainsi $Lf^*Lg^*E$ est représenté par
$f^*g^*\mathcal{K}^\bullet$. Comme $L(g \circ f)^*E$ est lui aussi
représenté par
$(g \circ f)^*\mathcal{K}^\bullet = f^*g^*\mathcal{K}^\bullet$
nous concluons.
\end{proof}

\begin{lemma}
\label{sites-cohomology-lemma-pullback-tensor-product}
Soit $f : (\Sh(\mathcal{C}), \mathcal{O}) \to (\Sh(\mathcal{D}), \mathcal{O}')$
un morphisme de topos annelés. Il existe un isomorphisme bifonctoriel
canonique
$$
Lf^*(
\mathcal{F}^\bullet \otimes_{\mathcal{O}'}^{\mathbf{L}} \mathcal{G}^\bullet
) =
Lf^*\mathcal{F}^\bullet 
\otimes_{\mathcal{O}}^{\mathbf{L}}
Lf^*\mathcal{G}^\bullet 
$$
pour $\mathcal{F}^\bullet, \mathcal{G}^\bullet \in \Ob(D(\mathcal{O}'))$.
\end{lemma}

\begin{proof}
Grâce à notre construction de l'image inverse dérivée dans le Lemme
\ref{sites-cohomology-lemma-derived-base-change} et à l'existence de résolutions donnée par
le Lemme \ref{sites-cohomology-lemma-K-flat-resolution}, nous pouvons remplacer
$\mathcal{F}^\bullet$ et $\mathcal{G}^\bullet$ par des complexes de
$\mathcal{O}'$-modules qui sont K-plats et dont les termes sont plats.
Dans ce cas,
$\mathcal{F}^\bullet \otimes_{\mathcal{O}'}^{\mathbf{L}} \mathcal{G}^\bullet$
est simplement le complexe total associé au bicomplexe
$\mathcal{F}^\bullet \otimes_{\mathcal{O}'} \mathcal{G}^\bullet$.
Le complexe
$\text{Tot}(\mathcal{F}^\bullet \otimes_{\mathcal{O}'} \mathcal{G}^\bullet)$
est K-plat à termes plats d'après le Lemme
\ref{sites-cohomology-lemma-tensor-product-K-flat} et Modules sur les sites, Lemme
\ref{sites-modules-lemma-tensor-flats}. Ainsi, l'isomorphisme du lemme
provient de l'isomorphisme
$$
\text{Tot}(f^*\mathcal{F}^\bullet \otimes_{\mathcal{O}}
f^*\mathcal{G}^\bullet)
\longrightarrow
f^*\text{Tot}(\mathcal{F}^\bullet \otimes_{\mathcal{O}'} \mathcal{G}^\bullet)
$$
dont les composantes sont les isomorphismes
$f^*\mathcal{F}^p \otimes_{\mathcal{O}} f^*\mathcal{G}^q \to
f^*(\mathcal{F}^p \otimes_{\mathcal{O}'} \mathcal{G}^q)$ fournis par
Modules sur les sites, Lemme
\ref{sites-modules-lemma-tensor-product-pullback}.
\end{proof}

\begin{lemma}
\label{sites-cohomology-lemma-variant-derived-pullback}
Soit $f : (\Sh(\mathcal{C}), \mathcal{O}) \to (\Sh(\mathcal{C}'), \mathcal{O}')$
un morphisme de topos annelés. Il existe un isomorphisme bifonctoriel
canonique
$$
\mathcal{F}^\bullet
\otimes_\mathcal{O}^{\mathbf{L}}
Lf^*\mathcal{G}^\bullet
=
\mathcal{F}^\bullet 
\otimes_{f^{-1}\mathcal{O}'}^{\mathbf{L}}
f^{-1}\mathcal{G}^\bullet 
$$
pour $\mathcal{F}^\bullet$ dans $D(\mathcal{O})$ et
$\mathcal{G}^\bullet$ dans $D(\mathcal{O}')$.
\end{lemma}

\begin{proof}
Soient $\mathcal{F}$ un $\mathcal{O}$-module et $\mathcal{G}$ un
$\mathcal{O}'$-module. Alors
$\mathcal{F} \otimes_{\mathcal{O}} f^*\mathcal{G} =
\mathcal{F} \otimes_{f^{-1}\mathcal{O}'} f^{-1}\mathcal{G}$
car
$f^*\mathcal{G} =
\mathcal{O} \otimes_{f^{-1}\mathcal{O}'} f^{-1}\mathcal{G}$.
Le lemme résulte de cette égalité et des définitions.
\end{proof}











\begin{lemma}
\label{sites-cohomology-lemma-check-K-flat-stalks}
Soit $(\mathcal{C}, \mathcal{O})$ un site annelé.
Soit $\mathcal{K}^\bullet$ un complexe de $\mathcal{O}$-modules.
\begin{enumerate}
\item Si $\mathcal{K}^\bullet$ est K-plat, alors, pour tout point $p$ du
site $\mathcal{C}$, le complexe de $\mathcal{O}_p$-modules
$\mathcal{K}_p^\bullet$ est K-plat au sens de
Compléments d'algèbre, Définition \ref{more-algebra-definition-K-flat}.
\item Si $\mathcal{C}$ possède suffisamment de points, la réciproque est
vraie.
\end{enumerate}
\end{lemma}

\begin{proof}
Démonstration de (2). Supposons que $\mathcal{C}$ possède suffisamment de
points et que $\mathcal{K}_p^\bullet$ soit K-plat pour tout point $p$ de
$\mathcal{C}$. Alors $\mathcal{K}^\bullet$ est K-plat, car $\otimes$ et les
sommes directes commutent à la prise des fibres, et l'exactitude peut se
vérifier sur les fibres ; voir Modules sur les sites, Lemme
\ref{sites-modules-lemma-check-exactness-stalks}.

\medskip\noindent
Démonstration de (1). Supposons $\mathcal{K}^\bullet$ K-plat. Choisissons
un quasi-isomorphisme $a : \mathcal{L}^\bullet \to \mathcal{K}^\bullet$ tel
que $\mathcal{L}^\bullet$ soit K-plat à termes plats, voir le Lemme
\ref{sites-cohomology-lemma-K-flat-resolution}. Toute image inverse de
$\mathcal{L}^\bullet$ est K-plate, voir le Lemme
\ref{sites-cohomology-lemma-pullback-K-flat}. En particulier, la fibre
$\mathcal{L}_p^\bullet$ est un complexe K-plat de
$\mathcal{O}_p$-modules. Le cône $C(a)$ de $a$ est donc un complexe
acyclique K-plat de $\mathcal{O}$-modules (Lemme
\ref{sites-cohomology-lemma-K-flat-two-out-of-three}), et il suffit de montrer que la fibre
de $C(a)$ est K-plate (d'après Compléments d'algèbre, Lemme
\ref{more-algebra-lemma-K-flat-two-out-of-three}). Nous pouvons donc
supposer que $\mathcal{K}^\bullet$ est K-plat et acyclique.

\medskip\noindent
Supposons $\mathcal{K}^\bullet$ acyclique et K-plat. Avant de poursuivre,
remplaçons le site $\mathcal{C}$ par un autre comme dans Sites, Lemme
\ref{sites-lemma-topos-good-site}, afin que $\mathcal{C}$ admette toutes
les limites finies. Il en résulte que la catégorie des voisinages de $p$
est filtrante (Sites, Lemme \ref{sites-lemma-neighbourhoods-directed}) et
que la colimite définissant la fibre d'un faisceau est filtrante. Soit $M$
un $\mathcal{O}_p$-module de présentation finie. Il suffit de montrer que
$\mathcal{K}^\bullet \otimes_{\mathcal{O}_p} M$ est acyclique ; voir
Compléments d'algèbre, Lemme
\ref{more-algebra-lemma-universally-acyclic-K-flat}. Comme
$\mathcal{O}_p$ est la colimite filtrante des $\mathcal{O}(U)$, où $U$
parcourt les voisinages de $p$, nous pouvons trouver un voisinage $(U, x)$
de $p$ et un $\mathcal{O}(U)$-module de présentation finie $M'$ dont le
changement de base à $\mathcal{O}_p$ est $M$ ; voir Algèbre, Lemme
\ref{algebra-lemma-colimit-category-fp-modules}. D'après le Lemme
\ref{sites-cohomology-lemma-restriction-K-flat}, nous pouvons remplacer
$\mathcal{C}, \mathcal{O}, \mathcal{K}^\bullet$ par
$\mathcal{C}/U, \mathcal{O}_U, \mathcal{K}^\bullet|_U$. Nous en déduisons
que nous pouvons supposer qu'il existe un $\mathcal{O}$-module
$\mathcal{F}$ tel que $M \cong \mathcal{F}_p$. Puisque
$\mathcal{K}^\bullet$ est K-plat et acyclique, le complexe
$\mathcal{K}^\bullet \otimes_\mathcal{O} \mathcal{F}$ est acyclique (car
il calcule, par définition, le produit tensoriel dérivé). La prise des
fibres étant un foncteur exact, il en résulte que
$\mathcal{K}^\bullet \otimes_{\mathcal{O}_p} M$ est acyclique, comme voulu.
\end{proof}

\begin{lemma}
\label{sites-cohomology-lemma-pullback-K-flat-points}
Soit $f : (\Sh(\mathcal{C}), \mathcal{O}) \to (\Sh(\mathcal{C}'), \mathcal{O}')$
un morphisme de topos annelés. Si $\mathcal{C}$ possède suffisamment de
points, l'image inverse d'un complexe K-plat de $\mathcal{O}'$-modules est
un complexe K-plat de $\mathcal{O}$-modules.
\end{lemma}

\begin{proof}
Cela résulte du Lemme \ref{sites-cohomology-lemma-check-K-flat-stalks}, de Modules sur les
sites, Lemme \ref{sites-modules-lemma-pullback-stalk}, et de Compléments
d'algèbre, Lemme \ref{more-algebra-lemma-base-change-K-flat}.
\end{proof}

\begin{lemma}
\label{sites-cohomology-lemma-tensor-pull-compatibility}
Soit $f : (\Sh(\mathcal{C}), \mathcal{O}_\mathcal{C}) \to
(\Sh(\mathcal{D}), \mathcal{O}_\mathcal{D})$ un morphisme de topos annelés.
Soient $\mathcal{K}^\bullet$ et $\mathcal{M}^\bullet$ des complexes de
$\mathcal{O}_\mathcal{D}$-modules. Le diagramme
$$
\xymatrix{
Lf^*(\mathcal{K}^\bullet
\otimes_{\mathcal{O}_\mathcal{D}}^\mathbf{L}
\mathcal{M}^\bullet) \ar[r] \ar[d] &
Lf^*\text{Tot}(\mathcal{K}^\bullet
\otimes_{\mathcal{O}_\mathcal{D}}
\mathcal{M}^\bullet) \ar[d] \\
Lf^*\mathcal{K}^\bullet \otimes_{\mathcal{O}_\mathcal{C}}^\mathbf{L}
Lf^*\mathcal{M}^\bullet \ar[d] &
f^*\text{Tot}(\mathcal{K}^\bullet
\otimes_{\mathcal{O}_\mathcal{D}}
\mathcal{M}^\bullet) \ar[d] \\
f^*\mathcal{K}^\bullet \otimes_{\mathcal{O}_\mathcal{C}}^\mathbf{L}
f^*\mathcal{M}^\bullet \ar[r] &
\text{Tot}(f^*\mathcal{K}^\bullet \otimes_{\mathcal{O}_\mathcal{C}}
f^*\mathcal{M}^\bullet)
}
$$
est commutatif.
\end{lemma}

\begin{proof}
Nous utiliserons l'existence de résolutions K-plates à termes plats (Lemme
\ref{sites-cohomology-lemma-K-flat-resolution}), le fait que l'image inverse dérivée se
calcule à l'aide de tels complexes (Lemme
\ref{sites-cohomology-lemma-derived-base-change}), ainsi que la préservation de ces
propriétés par image inverse (Lemme \ref{sites-cohomology-lemma-pullback-K-flat}). Si nous
choisissons de telles résolutions
$\mathcal{P}^\bullet \to \mathcal{K}^\bullet$ et
$\mathcal{Q}^\bullet \to \mathcal{M}^\bullet$, nous voyons alors que
$$
\xymatrix{
Lf^*\text{Tot}(\mathcal{P}^\bullet
\otimes_{\mathcal{O}_\mathcal{D}}
\mathcal{Q}^\bullet) \ar[r] \ar[d] &
Lf^*\text{Tot}(\mathcal{K}^\bullet
\otimes_{\mathcal{O}_\mathcal{D}}
\mathcal{M}^\bullet) \ar[d] \\
f^*\text{Tot}(\mathcal{P}^\bullet
\otimes_{\mathcal{O}_\mathcal{D}}
\mathcal{Q}^\bullet) \ar[d] \ar[r] &
f^*\text{Tot}(\mathcal{K}^\bullet
\otimes_{\mathcal{O}_\mathcal{D}}
\mathcal{M}^\bullet) \ar[d] \\
\text{Tot}(f^*\mathcal{P}^\bullet \otimes_{\mathcal{O}_\mathcal{C}}
f^*\mathcal{Q}^\bullet) \ar[r] &
\text{Tot}(f^*\mathcal{K}^\bullet \otimes_{\mathcal{O}_\mathcal{C}}
f^*\mathcal{M}^\bullet)
}
$$
est commutatif. Or, le côté gauche de ce diagramme est maintenant le côté
gauche du diagramme de l'énoncé, par notre choix de $\mathcal{P}^\bullet$
et $\mathcal{Q}^\bullet$ et par le Lemme
\ref{sites-cohomology-lemma-tensor-product-K-flat}.
\end{proof}









\section{Cohomologie des complexes non bornés}
\label{sites-cohomology-section-unbounded}

\noindent
Soit $(\mathcal{C}, \mathcal{O})$ un site annelé. La catégorie
$\textit{Mod}(\mathcal{O})$ est une catégorie abélienne de Grothendieck :
elle admet toutes les colimites, les colimites filtrantes y sont exactes,
et elle possède un générateur, à savoir
$$
\bigoplus\nolimits_{U \in \Ob(\mathcal{C})} j_{U!}\mathcal{O}_U,
$$
voir Modules sur les sites, section \ref{sites-modules-section-kernels}, et
Lemmes \ref{sites-modules-lemma-j-shriek-flat} et
\ref{sites-modules-lemma-module-quotient-flat}. D'après Injectifs, Théorème
\ref{injectives-theorem-K-injective-embedding-grothendieck}, pour tout
complexe $\mathcal{F}^\bullet$ de $\mathcal{O}$-modules, il existe un
quasi-isomorphisme injectif
$\mathcal{F}^\bullet \to \mathcal{I}^\bullet$ vers un complexe K-injectif
de $\mathcal{O}$-modules ; de plus, cette inclusion peut être choisie
fonctoriellement en $\mathcal{F}^\bullet$. Il résulte de Catégories
dérivées, Lemme \ref{derived-lemma-enough-K-injectives-implies}, que
\begin{enumerate}
\item tout foncteur exact $F : K(\textit{Mod}(\mathcal{O})) \to \mathcal{D}$
à valeurs dans une catégorie triangulée $\mathcal{D}$ admet un foncteur
dérivé à droite
$RF : D(\mathcal{O}) \to \mathcal{D}$,
\item pour tout foncteur additif
$F : \textit{Mod}(\mathcal{O}) \to \mathcal{A}$
à valeurs dans une catégorie abélienne $\mathcal{A}$, nous considérons le
foncteur exact $F : K(\textit{Mod}(\mathcal{O})) \to K(\mathcal{A})$
induit par $F$ et obtenons un foncteur dérivé à droite
$RF : D(\mathcal{O}) \to D(\mathcal{A})$.
\end{enumerate}
Par construction, $RF(\mathcal{F}^\bullet) = F(\mathcal{I}^\bullet)$, où
$\mathcal{F}^\bullet \to \mathcal{I}^\bullet$ est comme ci-dessus.

\medskip\noindent
Voici quelques exemples de ce qui précède :
\begin{enumerate}
\item Le foncteur $\Gamma(\mathcal{C}, -) : \textit{Mod}(\mathcal{O}) \to
\text{Mod}_{\Gamma(\mathcal{C}, \mathcal{O})}$ donne naissance à
$$
R\Gamma(\mathcal{C}, -) :
D(\mathcal{O})
\longrightarrow
D(\Gamma(\mathcal{C}, \mathcal{O}))
$$
Nous noterons
$H^i(\mathcal{C}, K) = H^i(R\Gamma(\mathcal{C}, K))$ la cohomologie.
\item Pour un objet $U$ de $\mathcal{C}$, considérons le foncteur
$\Gamma(U, -) : \textit{Mod}(\mathcal{O}) \to
\text{Mod}_{\Gamma(U, \mathcal{O})}$. Il donne naissance à
$$
R\Gamma(U, -) : D(\mathcal{O}) \to D(\Gamma(U, \mathcal{O}))
$$
Nous noterons $H^i(U, K) = H^i(R\Gamma(U, K))$ la cohomologie.
\item Pour un morphisme de topos annelés
$f : (\Sh(\mathcal{C}), \mathcal{O}) \to (\Sh(\mathcal{D}), \mathcal{O}')$
considérons le foncteur
$f_* : \textit{Mod}(\mathcal{O}) \to \textit{Mod}(\mathcal{O}')$
qui donne naissance à l'image directe dérivée totale
$$
Rf_* : D(\mathcal{O}) \longrightarrow D(\mathcal{O}')
$$
sur les catégories dérivées non bornées.
\end{enumerate}

\begin{lemma}
\label{sites-cohomology-lemma-adjoint}
Soit $f : (\Sh(\mathcal{C}), \mathcal{O}) \to (\Sh(\mathcal{D}), \mathcal{O}')$
un morphisme de topos annelés. Le foncteur $Rf_*$ défini ci-dessus et le
foncteur $Lf^*$ défini dans le Lemme \ref{sites-cohomology-lemma-derived-base-change} sont
adjoints :
$$
\Hom_{D(\mathcal{O})}(Lf^*\mathcal{G}^\bullet, \mathcal{F}^\bullet)
=
\Hom_{D(\mathcal{O}')}(\mathcal{G}^\bullet, Rf_*\mathcal{F}^\bullet)
$$
bifonctoriellement en $\mathcal{F}^\bullet \in \Ob(D(\mathcal{O}))$ et
$\mathcal{G}^\bullet \in \Ob(D(\mathcal{O}'))$.
\end{lemma}

\begin{proof}
Cela résulte formellement de l'existence de $Rf_*$ et de $Lf^*$ ; voir
Catégories dérivées, Lemme
\ref{derived-lemma-derived-adjoint-functors}.
\end{proof}

\begin{lemma}
\label{sites-cohomology-lemma-derived-pushforward-composition}
Soient
$f : (\Sh(\mathcal{C}), \mathcal{O}_\mathcal{C}) \to
(\Sh(\mathcal{D}), \mathcal{O}_\mathcal{D})$
et
$g : (\Sh(\mathcal{D}), \mathcal{O}_\mathcal{D}) \to
(\Sh(\mathcal{E}), \mathcal{O}_\mathcal{E})$
des morphismes de topos annelés. Alors
$Rg_* \circ Rf_* = R(g \circ f)_*$ comme foncteurs
$D(\mathcal{O}_\mathcal{C}) \to D(\mathcal{O}_\mathcal{E})$.
\end{lemma}

\begin{proof}
D'après le Lemme \ref{sites-cohomology-lemma-adjoint}, $Rg_* \circ Rf_*$ est adjoint à
$Lf^* \circ Lg^*$. Or $Lf^* \circ Lg^* = L(g \circ f)^*$ d'après le
Lemme \ref{sites-cohomology-lemma-derived-pullback-composition} ; l'unicité des foncteurs
adjoints donne donc $Rg_* \circ Rf_* = R(g \circ f)_*$.
\end{proof}

\begin{remark}
\label{sites-cohomology-remark-base-change}
La construction des foncteurs dérivés non bornés $Lf^*$ et $Rf_*$ permet
de construire le morphisme de changement de base en toute généralité.
Supposons en effet que
$$
\xymatrix{
(\Sh(\mathcal{C}'), \mathcal{O}_{\mathcal{C}'})
\ar[r]_{g'} \ar[d]_{f'} &
(\Sh(\mathcal{C}), \mathcal{O}_\mathcal{C}) \ar[d]^f \\
(\Sh(\mathcal{D}'), \mathcal{O}_{\mathcal{D}'})
\ar[r]^g &
(\Sh(\mathcal{D}), \mathcal{O}_\mathcal{D})
}
$$
soit un diagramme commutatif de topos annelés. Soit $K$ un objet de
$D(\mathcal{O}_\mathcal{C})$. Il existe alors un morphisme canonique de
changement de base
$$
Lg^*Rf_*K \longrightarrow R(f')_*L(g')^*K
$$
dans $D(\mathcal{O}_{\mathcal{D}'})$. Ce morphisme est en effet adjoint à
un morphisme $L(f')^*Lg^*Rf_*K \to L(g')^*K$. Puisque
$L(f')^* \circ Lg^* = L(g')^* \circ Lf^*$, cela revient à donner un
morphisme $L(g')^*Lf^*Rf_*K \to L(g')^*K$, que l'on peut prendre égal à
l'image par $L(g')^*$ du morphisme d'adjonction $Lf^*Rf_*K \to K$.
\end{remark}

\begin{remark}
\label{sites-cohomology-remark-compose-base-change}
Considérons un diagramme commutatif
$$
\xymatrix{
(\Sh(\mathcal{B}'), \mathcal{O}_{\mathcal{B}'})
\ar[r]_k \ar[d]_{f'} &
(\Sh(\mathcal{B}), \mathcal{O}_\mathcal{B}) \ar[d]^f \\
(\Sh(\mathcal{C}'), \mathcal{O}_{\mathcal{C}'})
\ar[r]^l \ar[d]_{g'} &
(\Sh(\mathcal{C}), \mathcal{O}_\mathcal{C}) \ar[d]^g \\
(\Sh(\mathcal{D}'), \mathcal{O}_{\mathcal{D}'})
\ar[r]^m &
(\Sh(\mathcal{D}), \mathcal{O}_\mathcal{D}) \\
}
$$
de topos annelés. Les morphismes de changement de base de la Remarque
\ref{sites-cohomology-remark-base-change}, pour les deux carrés, se composent alors pour
donner le morphisme de changement de base du rectangle extérieur. Plus
précisément, la composée
\begin{align*}
Lm^* \circ R(g \circ f)_*
& =
Lm^* \circ Rg_* \circ Rf_* \\
& \to Rg'_* \circ Ll^* \circ Rf_* \\
& \to Rg'_* \circ Rf'_* \circ Lk^* \\
& = R(g' \circ f')_* \circ Lk^*
\end{align*}
est le morphisme de changement de base du rectangle. Nous omettons la
vérification.
\end{remark}

\begin{remark}
\label{sites-cohomology-remark-compose-base-change-horizontal}
Considérons un diagramme commutatif
$$
\xymatrix{
(\Sh(\mathcal{C}''), \mathcal{O}_{\mathcal{C}''})
\ar[r]_{g'} \ar[d]_{f''} &
(\Sh(\mathcal{C}'), \mathcal{O}_{\mathcal{C}'})
\ar[r]_g \ar[d]_{f'} &
(\Sh(\mathcal{C}), \mathcal{O}_\mathcal{C}) \ar[d]^f \\
(\Sh(\mathcal{D}''), \mathcal{O}_{\mathcal{D}''})
\ar[r]^{h'} &
(\Sh(\mathcal{D}'), \mathcal{O}_{\mathcal{D}'})
\ar[r]^h &
(\Sh(\mathcal{D}), \mathcal{O}_\mathcal{D})
}
$$
de topos annelés. Les morphismes de changement de base de la Remarque
\ref{sites-cohomology-remark-base-change}, pour les deux carrés, se composent alors pour
donner le morphisme de changement de base du rectangle extérieur. Plus
précisément, la composée
\begin{align*}
L(h \circ h')^* \circ Rf_*
& =
L(h')^* \circ Lh^* \circ Rf_* \\
& \to L(h')^* \circ Rf'_* \circ Lg^* \\
& \to Rf''_* \circ L(g')^* \circ Lg^* \\
& = Rf''_* \circ L(g \circ g')^*
\end{align*}
est le morphisme de changement de base du rectangle. Nous omettons la
vérification.
\end{remark}



\begin{lemma}
\label{sites-cohomology-lemma-adjoints-push-pull-compatibility}
Soit $f : (\Sh(\mathcal{C}), \mathcal{O}_\mathcal{C}) \to
(\Sh(\mathcal{D}), \mathcal{O}_\mathcal{D})$ un morphisme de topos annelés.
Soit $\mathcal{K}^\bullet$ un complexe de
$\mathcal{O}_\mathcal{C}$-modules. Le diagramme
$$
\xymatrix{
Lf^*f_*\mathcal{K}^\bullet \ar[r] \ar[d] &
f^*f_*\mathcal{K}^\bullet \ar[d] \\
Lf^*Rf_*\mathcal{K}^\bullet \ar[r] &
\mathcal{K}^\bullet
}
$$
provenant de $Lf^* \to f^*$ sur les complexes, de $f_* \to Rf_*$ sur les
complexes et de l'adjonction $Lf^* \circ Rf_* \to \text{id}$ est
commutatif dans $D(\mathcal{O}_\mathcal{C})$.
\end{lemma}

\begin{proof}
Nous utiliserons l'existence de résolutions K-plates et de résolutions
K-injectives ; voir les Lemmes \ref{sites-cohomology-lemma-K-flat-resolution},
\ref{sites-cohomology-lemma-derived-base-change} et \ref{sites-cohomology-lemma-pullback-K-flat}, ainsi que
la discussion ci-dessus. Choisissons un quasi-isomorphisme
$\mathcal{K}^\bullet \to \mathcal{I}^\bullet$, où $\mathcal{I}^\bullet$
est K-injectif comme complexe de $\mathcal{O}_\mathcal{C}$-modules.
Choisissons un quasi-isomorphisme
$\mathcal{Q}^\bullet \to f_*\mathcal{I}^\bullet$, où
$\mathcal{Q}^\bullet$ est un complexe K-plat à termes plats de
$\mathcal{O}_\mathcal{D}$-modules. Nous pouvons choisir un complexe K-plat de
$\mathcal{O}_\mathcal{D}$-modules, noté $\mathcal{P}^\bullet$, à termes plats,
ainsi qu'un diagramme de morphismes de complexes
$$
\xymatrix{
\mathcal{P}^\bullet \ar[r] \ar[d] &
f_*\mathcal{K}^\bullet \ar[d] \\
\mathcal{Q}^\bullet \ar[r] & f_*\mathcal{I}^\bullet
}
$$
commutatif à homotopie près, dans lequel la flèche horizontale supérieure
est un quasi-isomorphisme. En effet, nous pouvons d'abord choisir un tel
diagramme pour un certain complexe $\mathcal{P}^\bullet$, car les
quasi-isomorphismes forment un système multiplicatif dans la catégorie
homotopique des complexes, puis choisir une résolution de
$\mathcal{P}^\bullet$ par un complexe K-plat à termes plats. En prenant les
images inverses, nous obtenons un diagramme de morphismes de complexes
$$
\xymatrix{
f^*\mathcal{P}^\bullet \ar[r] \ar[d] &
f^*f_*\mathcal{K}^\bullet \ar[d] \ar[r] &
\mathcal{K}^\bullet \ar[d] \\
f^*\mathcal{Q}^\bullet \ar[r] &
f^*f_*\mathcal{I}^\bullet \ar[r] &
\mathcal{I}^\bullet
}
$$
commutatif à homotopie près. Le rectangle extérieur établit l'assertion
du lemme.
\end{proof}



\begin{remark}
\label{sites-cohomology-remark-cup-product}
Soit $f : (\Sh(\mathcal{C}), \mathcal{O}_\mathcal{C}) \to
(\Sh(\mathcal{D}), \mathcal{O}_\mathcal{D})$ un morphisme de topos
annelés. L'adjonction de $Lf^*$ et $Rf_*$ permet de construire un cup-produit
relatif
$$
Rf_*K \otimes_{\mathcal{O}_\mathcal{D}}^\mathbf{L} Rf_*L
\longrightarrow
Rf_*(K \otimes_{\mathcal{O}_\mathcal{C}}^\mathbf{L} L)
$$
dans $D(\mathcal{O}_\mathcal{D})$, pour tous $K, L$ dans
$D(\mathcal{O}_\mathcal{C})$. Ce morphisme est en effet adjoint à un
morphisme
$Lf^*(Rf_*K \otimes_{\mathcal{O}_\mathcal{D}}^\mathbf{L} Rf_*L) \to
K \otimes_{\mathcal{O}_\mathcal{C}}^\mathbf{L} L$, que l'on peut prendre
égal à la composée de l'isomorphisme
$Lf^*(Rf_*K \otimes_{\mathcal{O}_\mathcal{D}}^\mathbf{L} Rf_*L) =
Lf^*Rf_*K \otimes_{\mathcal{O}_\mathcal{C}}^\mathbf{L} Lf^*Rf_*L$
(Lemme \ref{sites-cohomology-lemma-pullback-tensor-product}) avec le morphisme
$Lf^*Rf_*K \otimes_{\mathcal{O}_\mathcal{C}}^\mathbf{L} Lf^*Rf_*L
\to K \otimes_{\mathcal{O}_\mathcal{C}}^\mathbf{L} L$
provenant de la co-unité $Lf^* \circ Rf_* \to \text{id}$.
\end{remark}

\begin{lemma}
\label{sites-cohomology-lemma-torsion}
Soit $\mathcal{C}$ un site. Notons
$\mathcal{A} \subset \textit{Ab}(\mathcal{C})$ la sous-catégorie de Serre
constituée des faisceaux abéliens de torsion. Alors le foncteur
$D(\mathcal{A}) \to D_\mathcal{A}(\mathcal{C})$ est une équivalence.
\end{lemma}

\begin{proof}
Une observation essentielle est qu'un faisceau abélien injectif
$\mathcal{I}$ est divisible. En effet, si $s \in \mathcal{I}(U)$ est une
section locale, interprétons $s$ comme un morphisme
$s : j_{U!}\mathbf{Z} \to \mathcal{I}$ et appliquons la propriété
définissant un objet injectif au monomorphisme de faisceaux
$n : j_{U!}\mathbf{Z} \to j_{U!}\mathbf{Z}$. Il existe alors
$s' \in \mathcal{I}(U)$ tel que $ns' = s$.

\medskip\noindent
Pour un faisceau $\mathcal{F}$, notons $\mathcal{F}_{tor}$ son sous-faisceau
de torsion. Nous affirmons que, si $\mathcal{I}^\bullet$ est un complexe
de faisceaux abéliens injectifs dont les faisceaux de cohomologie sont de
torsion, alors
$$
\mathcal{I}^\bullet_{tor} \to \mathcal{I}^\bullet
$$
est un quasi-isomorphisme. En effet, la platitude de $\mathbf{Q}$ sur
$\mathbf{Z}$ donne
$$
H^p(\mathcal{I}^\bullet) \otimes_\mathbf{Z} \mathbf{Q} =
H^p(\mathcal{I}^\bullet \otimes_\mathbf{Z} \mathbf{Q})
$$
qui est nul car les faisceaux de cohomologie sont de torsion. Par
divisibilité (établie ci-dessus), le morphisme
$\mathcal{I}^\bullet \to \mathcal{I}^\bullet \otimes_\mathbf{Z} \mathbf{Q}$
est surjectif de noyau $\mathcal{I}^\bullet_{tor}$. L'affirmation résulte
alors de la suite exacte longue des faisceaux de cohomologie associée à la
suite exacte courte ainsi obtenue.

\medskip\noindent
Pour démontrer le lemme, construisons un adjoint à droite
$T : D(\mathcal{C}) \to D(\mathcal{A})$. Étant donné $K$ dans
$D(\mathcal{C})$, nous pouvons représenter $K$ par un complexe K-injectif
$\mathcal{I}^\bullet$ dont les faisceaux de cohomologie sont injectifs ;
voir Injectifs, Théorème
\ref{injectives-theorem-K-injective-embedding-grothendieck}. Posons alors
$T(K) = \mathcal{I}^\bullet_{tor}$ ; autrement dit, $T$ est le foncteur
dérivé à droite du foncteur de torsion. Le foncteur $T$ est adjoint à
droite de $i : D(\mathcal{A}) \to D_\mathcal{A}(\mathcal{C})$. Cela résulte
immédiatement de l'observation suivante : si $\mathcal{F}^\bullet$ est un
complexe de faisceaux de torsion, alors
$$
\Hom_{K(\mathcal{A})}(\mathcal{F}^\bullet, I^\bullet_{tor}) =
\Hom_{K(\textit{Ab}(\mathcal{C}))}(\mathcal{F}^\bullet, I^\bullet)
$$
en particulier, $\mathcal{I}^\bullet_{tor}$ est un complexe K-injectif de
$\mathcal{A}$. Nous omettons quelques détails ; cela résulte aussi, en cas
de doute, du résultat plus général Catégories dérivées, Lemme
\ref{derived-lemma-derived-adjoint-functors}. Notre affirmation ci-dessus
donne $L = T(i(L))$ pour $L$ dans $D(\mathcal{A})$, et $i(T(K)) = K$ si
$K$ appartient à $D_\mathcal{A}(\mathcal{C})$. Le résultat découle alors
de Catégories, Lemme \ref{categories-lemma-adjoint-fully-faithful}.
\end{proof}




\section{Quelques propriétés des complexes K-injectifs}
\label{sites-cohomology-section-properties-K-injective}

\noindent
Soit $(\mathcal{C}, \mathcal{O})$ un site annelé et soit $U$ un objet de
$\mathcal{C}$. Notons
$j : (\Sh(\mathcal{C}/U), \mathcal{O}_U) \to (\Sh(\mathcal{C}), \mathcal{O})$
le morphisme de localisation correspondant. Le foncteur image inverse
$j^*$ est exact, puisqu'il n'est autre que le foncteur de restriction.
Ainsi, sur tout complexe, l'image inverse dérivée $Lj^*$ se calcule par
simple restriction du complexe. Nous notons souvent simplement le foncteur
correspondant
$$
D(\mathcal{O}) \to D(\mathcal{O}_U),
\quad
E \mapsto j^*E = E|_U
$$
De même, le prolongement par zéro
$j_! : \textit{Mod}(\mathcal{O}_U) \to \textit{Mod}(\mathcal{O})$ (voir
Modules sur les sites, Définition
\ref{sites-modules-definition-localize-ringed-site}) est un foncteur exact
(Modules sur les sites, Lemme
\ref{sites-modules-lemma-extension-by-zero-exact}). Il induit donc un
foncteur
$$
j_! : D(\mathcal{O}_U) \to D(\mathcal{O}), \quad
F \mapsto j_!F
$$
en appliquant simplement $j_!$ à tout complexe représentant l'objet $F$.

\begin{lemma}
\label{sites-cohomology-lemma-restrict-K-injective-to-open}
Soit $(\mathcal{C}, \mathcal{O})$ un site annelé et soit $U$ un objet de
$\mathcal{C}$. La restriction d'un complexe K-injectif de
$\mathcal{O}$-modules à $\mathcal{C}/U$ est un complexe K-injectif de
$\mathcal{O}_U$-modules.
\end{lemma}

\begin{proof}
Cela résulte immédiatement de Catégories dérivées, Lemme
\ref{derived-lemma-adjoint-preserve-K-injectives}, et du fait que le
foncteur de restriction admet l'adjoint à gauche exact $j_!$. Voir la
discussion ci-dessus.
\end{proof}

\begin{lemma}
\label{sites-cohomology-lemma-unbounded-cohomology-of-open}
Soit $(\mathcal{C}, \mathcal{O})$ un site annelé. Soit
$U \in \Ob(\mathcal{C})$. Pour $K$ dans $D(\mathcal{O})$, nous avons
$H^p(U, K) = H^p(\mathcal{C}/U, K|_{\mathcal{C}/U})$.
\end{lemma}

\begin{proof}
Soit $\mathcal{I}^\bullet$ un complexe K-injectif de $\mathcal{O}$-modules
représentant $K$. Alors
$$
H^q(U, K) = H^q(\Gamma(U, \mathcal{I}^\bullet)) =
H^q(\Gamma(\mathcal{C}/U, \mathcal{I}^\bullet|_{\mathcal{C}/U}))
$$
par construction de la cohomologie. D'après le Lemme
\ref{sites-cohomology-lemma-restrict-K-injective-to-open}, le complexe
$\mathcal{I}^\bullet|_{\mathcal{C}/U}$ est un complexe K-injectif
représentant $K|_{\mathcal{C}/U}$, d'où le lemme.
\end{proof}

\begin{lemma}
\label{sites-cohomology-lemma-sheafification-cohomology}
Soit $(\mathcal{C}, \mathcal{O})$ un site annelé. Soit $K$ un objet de
$D(\mathcal{O})$. Le faisceau associé au préfaisceau
$$
U \mapsto H^q(U, K) = H^q(\mathcal{C}/U, K|_{\mathcal{C}/U})
$$
est le $q$-ième faisceau de cohomologie $H^q(K)$ de $K$.
\end{lemma}

\begin{proof}
L'égalité $H^q(U, K) = H^q(\mathcal{C}/U, K|_{\mathcal{C}/U})$ résulte du
Lemme \ref{sites-cohomology-lemma-unbounded-cohomology-of-open}. Choisissons un complexe
K-injectif $\mathcal{I}^\bullet$ représentant $K$. Alors
$$
H^q(U, K) =
\frac{\Ker(\mathcal{I}^q(U) \to \mathcal{I}^{q + 1}(U))}
{\Im(\mathcal{I}^{q - 1}(U) \to \mathcal{I}^q(U))}.
$$
par notre construction de la cohomologie. Comme
$H^q(K) = \Ker(\mathcal{I}^q \to \mathcal{I}^{q + 1})/
\Im(\mathcal{I}^{q - 1} \to \mathcal{I}^q)$, le résultat est clair.
\end{proof}

\begin{lemma}
\label{sites-cohomology-lemma-restrict-direct-image-open}
Soit $f : (\mathcal{C}, \mathcal{O}_\mathcal{C}) \to
(\mathcal{D}, \mathcal{O}_\mathcal{D})$ un morphisme de sites annelés
correspondant au foncteur continu $u : \mathcal{D} \to \mathcal{C}$.
Pour $V \in \mathcal{D}$, posons $U = u(V)$ et notons
$g : (\mathcal{C}/U, \mathcal{O}_U) \to (\mathcal{D}/V, \mathcal{O}_V)$
le morphisme induit de sites annelés (Modules sur les sites, Lemme
\ref{sites-modules-lemma-localize-morphism-ringed-sites}).
Alors $(Rf_*E)|_{\mathcal{D}/V} = Rg_*(E|_{\mathcal{C}/U})$ pour $E$ dans
$D(\mathcal{O}_\mathcal{C})$.
\end{lemma}

\begin{proof}
Représentons $E$ par un complexe K-injectif $\mathcal{I}^\bullet$ de
$\mathcal{O}_\mathcal{C}$-modules. Alors
$Rf_*(E) = f_*\mathcal{I}^\bullet$
et $Rg_*(E|_{\mathcal{C}/U}) = g_*(\mathcal{I}^\bullet|_{\mathcal{C}/U})$
d'après le Lemme \ref{sites-cohomology-lemma-restrict-K-injective-to-open}. Comme il est
clair que
$(f_*\mathcal{F})|_{\mathcal{D}/V} = g_*(\mathcal{F}|_{\mathcal{C}/U})$
pour tout faisceau $\mathcal{F}$ sur $\mathcal{C}$ (voir Modules sur les
sites, Lemme \ref{sites-modules-lemma-localize-morphism-ringed-sites}, ou
le résultat plus élémentaire Sites, Lemme
\ref{sites-lemma-localize-morphism}), le résultat s'ensuit.
\end{proof}

\begin{lemma}
\label{sites-cohomology-lemma-Leray-unbounded}
Soit $f : (\mathcal{C}, \mathcal{O}_\mathcal{C}) \to
(\mathcal{D}, \mathcal{O}_\mathcal{D})$ un morphisme de sites annelés
correspondant au foncteur continu $u : \mathcal{D} \to \mathcal{C}$.
Alors $R\Gamma(\mathcal{D}, -) \circ Rf_* = R\Gamma(\mathcal{C}, -)$ comme
foncteurs $D(\mathcal{O}_\mathcal{C}) \to D(\Gamma(\mathcal{O}_\mathcal{D}))$.
Plus généralement, pour $V \in \mathcal{D}$ et $U = u(V)$, nous avons
$R\Gamma(U, -) = R\Gamma(V, -) \circ Rf_*$.
\end{lemma}

\begin{proof}
Considérons le topos ponctuel $pt$ muni de $\mathcal{O}_{pt}$ défini par
l'anneau $\Gamma(\mathcal{O}_\mathcal{D})$. Il existe un morphisme canonique
$(\mathcal{D}, \mathcal{O}_\mathcal{D}) \to (pt, \mathcal{O}_{pt})$
de topos annelés induisant l'identification sur les sections globales des
faisceaux structuraux. Alors
$D(\mathcal{O}_{pt}) = D(\Gamma(\mathcal{O}_\mathcal{D}))$.
L'assertion
$R\Gamma(\mathcal{D}, -) \circ Rf_* = R\Gamma(\mathcal{C}, -)$
résulte du Lemme \ref{sites-cohomology-lemma-derived-pushforward-composition} appliqué à
$$
(\mathcal{C}, \mathcal{O}_\mathcal{C}) \to
(\mathcal{D}, \mathcal{O}_\mathcal{D}) \to (pt, \mathcal{O}_{pt})
$$
La seconde assertion, plus générale, résulte de la première après
application du Lemme \ref{sites-cohomology-lemma-restrict-direct-image-open}.
\end{proof}

\begin{lemma}
\label{sites-cohomology-lemma-unbounded-describe-higher-direct-images}
Soit $f : (\mathcal{C}, \mathcal{O}_\mathcal{C}) \to
(\mathcal{D}, \mathcal{O}_\mathcal{D})$ un morphisme de sites annelés
correspondant au foncteur continu $u : \mathcal{D} \to \mathcal{C}$.
Soit $K$ dans $D(\mathcal{O}_\mathcal{C})$. Alors $H^i(Rf_*K)$ est le
faisceau associé au préfaisceau
$$
V \mapsto H^i(u(V), K) = H^i(V, Rf_*K)
$$
\end{lemma}

\begin{proof}
L'égalité $H^i(u(V), K) = H^i(V, Rf_*K)$ s'obtient en prenant la
cohomologie dans la seconde assertion du Lemme
\ref{sites-cohomology-lemma-Leray-unbounded}. L'assertion concernant le faisceau associé
résulte alors du Lemme \ref{sites-cohomology-lemma-sheafification-cohomology}.
\end{proof}

\begin{lemma}
\label{sites-cohomology-lemma-modules-abelian-unbounded}
Soit $(\mathcal{C}, \mathcal{O}_\mathcal{C})$ un site annelé. Soit $K$ un
objet de $D(\mathcal{O}_\mathcal{C})$ et notons $K_{ab}$ son image dans
$D(\underline{\mathbf{Z}}_\mathcal{C})$.
\begin{enumerate}
\item Il existe un morphisme canonique
$R\Gamma(\mathcal{C}, K) \to R\Gamma(\mathcal{C}, K_{ab})$
qui est un isomorphisme dans $D(\textit{Ab})$.
\item Pour tout $U \in \mathcal{C}$, il existe un morphisme canonique
$R\Gamma(U, K) \to R\Gamma(U, K_{ab})$
qui est un isomorphisme dans $D(\textit{Ab})$.
\item Soit $f : (\mathcal{C}, \mathcal{O}_\mathcal{C}) \to
(\mathcal{D}, \mathcal{O}_\mathcal{D})$ un morphisme de sites annelés.
Il existe un morphisme canonique $Rf_*K \to Rf_*(K_{ab})$ qui est un
isomorphisme dans $D(\underline{\mathbf{Z}}_\mathcal{D})$.
\end{enumerate}
\end{lemma}

\begin{proof}
Le morphisme se construit comme suit. Choisissons un complexe K-injectif
$\mathcal{I}^\bullet$ représentant $K$. Choisissons un quasi-isomorphisme
$\mathcal{I}^\bullet \to \mathcal{J}^\bullet$, où $\mathcal{J}^\bullet$
est un complexe K-injectif de groupes abéliens. Le morphisme de (1) est
alors donné par
$\Gamma(\mathcal{C}, \mathcal{I}^\bullet) \to
\Gamma(\mathcal{C}, \mathcal{J}^\bullet)$
(2) est donné par
$\Gamma(U, \mathcal{I}^\bullet) \to \Gamma(U, \mathcal{J}^\bullet)$
et le morphisme de (3) est donné par
$f_*\mathcal{I}^\bullet \to f_*\mathcal{J}^\bullet$.
Pour montrer que ces morphismes sont des isomorphismes, il suffit de
prouver qu'ils induisent des isomorphismes sur les groupes de cohomologie
et les faisceaux de cohomologie. D'après les Lemmes
\ref{sites-cohomology-lemma-unbounded-cohomology-of-open} et
\ref{sites-cohomology-lemma-unbounded-describe-higher-direct-images}, il suffit de montrer
que le morphisme
$$
H^0(\mathcal{C}, K) \longrightarrow H^0(\mathcal{C}, K_{ab})
$$
est un isomorphisme. Observons que
$$
H^0(\mathcal{C}, K) =
\Hom_{D(\mathcal{O}_\mathcal{C})}(\mathcal{O}_\mathcal{C}, K)
$$
et de même pour l'autre groupe. Choisissons un complexe quelconque
$\mathcal{K}^\bullet$ de $\mathcal{O}_\mathcal{C}$-modules représentant
$K$. Par construction de la catégorie dérivée comme localisation, nous
avons
$$
\Hom_{D(\mathcal{O}_\mathcal{C})}(\mathcal{O}_\mathcal{C}, K) =
\colim_{s : \mathcal{F}^\bullet \to \mathcal{O}_\mathcal{C}}
\Hom_{K(\mathcal{O}_\mathcal{C})}(\mathcal{F}^\bullet, \mathcal{K}^\bullet)
$$
où la colimite porte sur les quasi-isomorphismes $s$ de complexes de
$\mathcal{O}_\mathcal{C}$-modules. De même, nous avons
$$
\Hom_{D(\underline{\mathbf{Z}}_\mathcal{C})}
(\underline{\mathbf{Z}}_\mathcal{C}, K) =
\colim_{s : \mathcal{G}^\bullet \to \underline{\mathbf{Z}}_\mathcal{C}}
\Hom_{K(\underline{\mathbf{Z}}_\mathcal{C})}
(\mathcal{G}^\bullet, \mathcal{K}^\bullet)
$$
Observons ensuite que les quasi-isomorphismes
$s : \mathcal{G}^\bullet \to \underline{\mathbf{Z}}_\mathcal{C}$
pour lesquels $\mathcal{G}^\bullet$ est un complexe borné supérieurement
de $\underline{\mathbf{Z}}_\mathcal{C}$-modules plats forment un système
cofinal. (Cela résulte de Modules sur les sites, Lemme
\ref{sites-modules-lemma-module-quotient-flat}, et de Catégories dérivées,
Lemme \ref{derived-lemma-subcategory-left-resolution} ; voir la discussion
de la section \ref{sites-cohomology-section-flat}.) Nous pouvons donc construire un inverse
du morphisme
$H^0(\mathcal{C}, K) \longrightarrow H^0(\mathcal{C}, K_{ab})$
en représentant un élément $\xi \in H^0(\mathcal{C}, K_{ab})$ par une paire
$$
(s : \mathcal{G}^\bullet \to \underline{\mathbf{Z}}_\mathcal{C},
a : \mathcal{G}^\bullet \to \mathcal{K}^\bullet)
$$
où $\mathcal{G}^\bullet$ est un complexe borné supérieurement de
$\underline{\mathbf{Z}}_\mathcal{C}$-modules plats, et en envoyant cette
paire sur
$$
(\mathcal{G}^\bullet \otimes_{\underline{\mathbf{Z}}_\mathcal{C}}
\mathcal{O}_\mathcal{C}
\to \mathcal{O}_\mathcal{C},
\mathcal{G}^\bullet  \otimes_{\underline{\mathbf{Z}}_\mathcal{C}}
\mathcal{O}_\mathcal{C}
\to \mathcal{K}^\bullet)
$$
Il suffit ici de remarquer que la première flèche est un
quasi-isomorphisme d'après les Lemmes
\ref{sites-cohomology-lemma-derived-tor-quasi-isomorphism-other-side} et
\ref{sites-cohomology-lemma-bounded-flat-K-flat}. Nous omettons la vérification détaillée du
fait que cette construction donne bien un inverse.
\end{proof}

\begin{lemma}
\label{sites-cohomology-lemma-adjoint-lower-shriek-restrict}
Soit $(\mathcal{C}, \mathcal{O})$ un site annelé et soit $U$ un objet de
$\mathcal{C}$. Notons
$j : (\Sh(\mathcal{C}/U), \mathcal{O}_U) \to (\Sh(\mathcal{C}), \mathcal{O})$
le morphisme de localisation correspondant. Le foncteur de restriction
$D(\mathcal{O}) \to D(\mathcal{O}_U)$ est adjoint à droite du prolongement
par zéro $j_! : D(\mathcal{O}_U) \to D(\mathcal{O})$.
\end{lemma}

\begin{proof}
Nous devons montrer que
$$
\Hom_{D(\mathcal{O})}(j_!E, F) = \Hom_{D(\mathcal{O}_U)}(E, F|_U)
$$
Choisissons un complexe $\mathcal{E}^\bullet$ de $\mathcal{O}_U$-modules
représentant $E$ et un complexe K-injectif $\mathcal{I}^\bullet$
représentant $F$. D'après le Lemme
\ref{sites-cohomology-lemma-restrict-K-injective-to-open}, le complexe
$\mathcal{I}^\bullet|_U$ est lui aussi K-injectif. La formule ci-dessus
devient donc
$$
\Hom_{D(\mathcal{O})}(j_!\mathcal{E}^\bullet, \mathcal{I}^\bullet) =
\Hom_{D(\mathcal{O}_U)}(\mathcal{E}^\bullet, \mathcal{I}^\bullet|_U)
$$
ce qui est vrai puisque $|_U$ et $j_!$ sont des foncteurs adjoints
(Modules sur les sites, Lemme
\ref{sites-modules-lemma-extension-by-zero}) et d'après Catégories dérivées,
Lemme \ref{derived-lemma-K-injective}.
\end{proof}

\begin{lemma}
\label{sites-cohomology-lemma-j-shriek-and-tensor}
Soit $(\mathcal{C}, \mathcal{O})$ un site annelé. Soit
$U \in \Ob(\mathcal{C})$. Pour $L$ dans $D(\mathcal{O}_U)$ et $K$ dans
$D(\mathcal{O})$, nous avons
$j_!L \otimes_\mathcal{O}^\mathbf{L} K =
j_!(L \otimes_{\mathcal{O}_U}^\mathbf{L} K|_U)$.
\end{lemma}

\begin{proof}
Représentons $L$ par un complexe de $\mathcal{O}_U$-modules et $K$ par un
complexe K-plat de $\mathcal{O}$-modules, puis appliquons Modules sur les
sites, Lemme \ref{sites-modules-lemma-j-shriek-and-tensor}. Nous omettons
les détails.
\end{proof}

\begin{lemma}
\label{sites-cohomology-lemma-K-injective-flat}
Soit $f : (\Sh(\mathcal{C}), \mathcal{O}_\mathcal{C}) \to
(\Sh(\mathcal{D}), \mathcal{O}_\mathcal{D})$ un morphisme plat de topos
annelés. Si $\mathcal{I}^\bullet$ est un complexe K-injectif de
$\mathcal{O}_\mathcal{C}$-modules, alors $f_*\mathcal{I}^\bullet$ est
K-injectif comme complexe de $\mathcal{O}_\mathcal{D}$-modules.
\end{lemma}

\begin{proof}
Cela est vrai car
$$
\Hom_{K(\mathcal{O}_\mathcal{D})}(\mathcal{F}^\bullet, f_*\mathcal{I}^\bullet)
=
\Hom_{K(\mathcal{O}_\mathcal{C})}(f^*\mathcal{F}^\bullet, \mathcal{I}^\bullet)
$$
d'après Modules sur les sites, Lemme
\ref{sites-modules-lemma-adjoint-pullback-pushforward-modules}, et car
$f^*$ est exact puisque $f$ est supposé plat.
\end{proof}

\begin{lemma}
\label{sites-cohomology-lemma-hom-K-injective}
Soit $\mathcal{C}$ un site. Soit $\mathcal{O} \to \mathcal{O}'$ un
morphisme de faisceaux d'anneaux. Si $\mathcal{I}^\bullet$ est un complexe
K-injectif de $\mathcal{O}$-modules, alors
$\SheafHom_\mathcal{O}(\mathcal{O}', \mathcal{I}^\bullet)$
est un complexe K-injectif de $\mathcal{O}'$-modules.
\end{lemma}

\begin{proof}
Cela est vrai car
$\Hom_{K(\mathcal{O}')}(\mathcal{G}^\bullet,
\Hom_\mathcal{O}(\mathcal{O}', \mathcal{I}^\bullet)) =
\Hom_{K(\mathcal{O})}(\mathcal{G}^\bullet, \mathcal{I}^\bullet)$
d'après Modules sur les sites, Lemme
\ref{sites-modules-lemma-adjoint-hom-restrict}.
\end{proof}






\section{Localisation et cohomologie}
\label{sites-cohomology-section-localization}

\noindent
Soit $\mathcal{C}$ un site et soit $f : X \to Y$ un morphisme de
$\mathcal{C}$. Nous obtenons alors un morphisme de topos
$$
j_{X/Y} : \Sh(\mathcal{C}/X) \longrightarrow \Sh(\mathcal{C}/Y)
$$
Voir Sites, sections \ref{sites-section-localize} et
\ref{sites-section-more-localize}. Certaines questions de cohomologie sont
plus simples pour ce type de morphismes de topos. Voici un exemple où l'on
obtient une forme élémentaire du théorème de changement de base.

\begin{lemma}
\label{sites-cohomology-lemma-localize-cartesian-square}
Soit $\mathcal{C}$ un site. Soit
$$
\xymatrix{
X' \ar[d] \ar[r] & X \ar[d] \\
Y' \ar[r] & Y
}
$$
un diagramme cartésien de $\mathcal{C}$. Alors nous avons
$j_{Y'/Y}^{-1} \circ Rj_{X/Y, *} = Rj_{X'/Y', *} \circ j_{X'/X}^{-1}$
comme foncteurs $D(\mathcal{C}/X) \to D(\mathcal{C}/Y')$.
\end{lemma}

\begin{proof}
Soit $E \in D(\mathcal{C}/X)$. Choisissons un complexe K-injectif
$\mathcal{I}^\bullet$ de faisceaux abéliens sur $\mathcal{C}/X$
représentant $E$. D'après le Lemme
\ref{sites-cohomology-lemma-restrict-K-injective-to-open},
$j_{X'/X}^{-1}\mathcal{I}^\bullet$ est lui aussi K-injectif. Nous pouvons
donc calculer $Rj_{X'/Y'}(j_{X'/X}^{-1}E)$ par
$j_{X'/Y', *}j_{X'/X}^{-1}\mathcal{I}^\bullet$. L'égalité résulte alors de
Sites, Lemme \ref{sites-lemma-localize-cartesian-square}.
\end{proof}

\noindent
Si nous avons un site annelé $(\mathcal{C}, \mathcal{O})$ et un morphisme
$f : X \to Y$ de $\mathcal{C}$, alors $j_{X/Y}$ devient un morphisme de
topos annelés
$$
j_{X/Y} :
(\Sh(\mathcal{C}/X), \mathcal{O}_X)
\longrightarrow
(\Sh(\mathcal{C}/Y), \mathcal{O}_Y)
$$
Voir Modules sur les sites, Lemme
\ref{sites-modules-lemma-relocalize}.

\begin{lemma}
\label{sites-cohomology-lemma-localize-cartesian-square-modules}
Soit $(\mathcal{C}, \mathcal{O})$ un site annelé. Soit
$$
\xymatrix{
X' \ar[d] \ar[r] & X \ar[d] \\
Y' \ar[r] & Y
}
$$
un diagramme cartésien de $\mathcal{C}$. Alors nous avons
$j_{Y'/Y}^* \circ Rj_{X/Y, *} = Rj_{X'/Y', *} \circ j_{X'/X}^*$
comme foncteurs
$D(\mathcal{O}_X) \to D(\mathcal{O}_{Y'})$.
\end{lemma}

\begin{proof}
Puisque $j_{Y'/Y}^{-1}\mathcal{O}_Y = \mathcal{O}_{Y'}$, nous avons
$j_{Y'/Y}^* = Lj_{Y'/Y}^* = j_{Y'/Y}^{-1}$. De même,
$j_{X'/X}^* = Lj_{X'/X}^* = j_{X'/X}^{-1}$. D'après le Lemme
\ref{sites-cohomology-lemma-modules-abelian-unbounded}, il suffit donc de démontrer le
résultat dans les catégories dérivées de faisceaux abéliens, ce que nous
avons fait dans le Lemme \ref{sites-cohomology-lemma-localize-cartesian-square}.
\end{proof}












\section{Systèmes projectifs et cohomologie}
\label{sites-cohomology-section-inverse-systems}

\noindent
Nous démontrons quelques résultats sur les systèmes projectifs de faisceaux
de modules.

\begin{lemma}
\label{sites-cohomology-lemma-ML-general}
Soit $I$ un idéal d'un anneau $A$. Soit $\mathcal{C}$ un site. Soit
$$
\ldots \to \mathcal{F}_3 \to \mathcal{F}_2 \to \mathcal{F}_1
$$
un système projectif de faisceaux de $A$-modules sur $\mathcal{C}$ tel que
$\mathcal{F}_n = \mathcal{F}_{n + 1}/I^n\mathcal{F}_{n + 1}$. Soit
$p \geq 0$. Supposons que
$$
\bigoplus\nolimits_{n \geq 0} H^{p + 1}(\mathcal{C}, I^n\mathcal{F}_{n + 1})
$$
satisfasse la condition de chaîne ascendante comme
$\bigoplus_{n \geq 0} I^n/I^{n + 1}$-module gradué. Alors le système
projectif $M_n = H^p(\mathcal{C}, \mathcal{F}_n)$ satisfait la condition de
Mittag-Leffler\footnote{En fait, il existe $c \geq 0$ tel que
$\Im(M_n \to M_{n - c})$ soit l'image stable pour tout $n \geq c$.}.
\end{lemma}

\begin{proof}
Posons $N_n = H^{p + 1}(\mathcal{C}, I^n\mathcal{F}_{n + 1})$ et soit
$\delta_n : M_n \to N_n$ le morphisme de connexion en cohomologie provenant
de la suite exacte courte
$0 \to I^n\mathcal{F}_{n + 1} \to \mathcal{F}_{n + 1} \to \mathcal{F}_n \to 0$.
Alors $\bigoplus \Im(\delta_n) \subset \bigoplus N_n$ est un sous-module
gradué. En effet, si $s \in M_n$ et $f \in I^m$, nous avons un diagramme
commutatif
$$
\xymatrix{
0 \ar[r] &
I^n\mathcal{F}_{n + 1} \ar[d]_f \ar[r] &
\mathcal{F}_{n + 1} \ar[d]_f \ar[r] &
\mathcal{F}_n \ar[d]_f \ar[r] & 0 \\
0 \ar[r] &
I^{n + m}\mathcal{F}_{n + m + 1} \ar[r] &
\mathcal{F}_{n + m + 1} \ar[r] &
\mathcal{F}_{n + m} \ar[r] & 0
}
$$
Le morphisme vertical médian est défini en relevant une section locale de
$\mathcal{F}_{n + 1}$ en une section de $\mathcal{F}_{n + m + 1}$, puis en
la multipliant par $f$ ; les autres flèches verticales se définissent de
façon analogue. Nous en déduisons que
$\delta_{n + m}(fs) = f \delta_n(s)$. Par hypothèse, nous pouvons trouver
$s_j \in M_{n_j}$, $j = 1, \ldots, N$, tels que les
$\delta_{n_j}(s_j)$ engendrent $\bigoplus \Im(\delta_n)$ comme module
gradué. Soit $n > c = \max(n_j)$ et soit $s \in M_n$. Il existe alors
$f_j \in I^{n - n_j}$ tels que
$\delta_n(s) = \sum f_j \delta_{n_j}(s_j)$. Nous en déduisons que
$\delta(s - \sum f_j s_j) = 0$, c'est-à-dire qu'il existe
$s' \in M_{n + 1}$ dont l'image est $s - \sum f_js_j$ dans $M_n$. Il
s'ensuit que
$$
\Im(M_{n + 1} \to M_{n - c}) = \Im(M_n \to M_{n - c})
$$
En effet, les éléments $f_js_j$ ont une image nulle dans $M_{n - c}$.
Cela démontre le lemme.
\end{proof}

\begin{lemma}
\label{sites-cohomology-lemma-ML-general-better}
Soit $I$ un idéal d'un anneau $A$. Soit $\mathcal{C}$ un site. Soit
$$
\ldots \to \mathcal{F}_3 \to \mathcal{F}_2 \to \mathcal{F}_1
$$
un système projectif de $A$-modules sur $\mathcal{C}$ tel que
$\mathcal{F}_n = \mathcal{F}_{n + 1}/I^n\mathcal{F}_{n + 1}$. Soit
$p \geq 0$. Pour tout $n$, posons
$$
N_n =
\bigcap\nolimits_{m \geq n}
\Im\left(
H^{p + 1}(\mathcal{C}, I^n\mathcal{F}_{m + 1}) \to
H^{p + 1}(\mathcal{C}, I^n\mathcal{F}_{n + 1})
\right)
$$
Si $\bigoplus N_n$ satisfait la condition de chaîne ascendante comme
$\bigoplus_{n \geq 0} I^n/I^{n + 1}$-module gradué, alors le système
projectif $M_n = H^p(\mathcal{C}, \mathcal{F}_n)$ satisfait la condition de
Mittag-Leffler\footnote{En fait, il existe $c \geq 0$ tel que
$\Im(M_n \to M_{n - c})$ soit l'image stable pour tout $n \geq c$.}.
\end{lemma}

\begin{proof}
La démonstration est exactement la même que celle du Lemme
\ref{sites-cohomology-lemma-ML-general}. En effet, les arguments qui y sont donnés
établiront le résultat dès que nous aurons montré que $\bigoplus N_n$ est
un sous-$\bigoplus_{n \geq 0} I^n/I^{n + 1}$-module gradué de
$\bigoplus H^{p + 1}(\mathcal{C}, I^n\mathcal{F}_{n + 1})$ et que l'image
des morphismes de connexion
$\delta_n : M_n \to H^{p + 1}(\mathcal{C}, I^n\mathcal{F}_{n + 1})$ est
contenue dans $N_n$.

\medskip\noindent
Supposons que $\xi \in N_n$ et $f \in I^k$. Choisissons
$m \gg n + k$, puis un élément
$\xi' \in H^{p + 1}(\mathcal{C}, I^n\mathcal{F}_{m + 1})$ relevant
$\xi$. Considérons le diagramme
$$
\xymatrix{
0 \ar[r] &
I^n\mathcal{F}_{m + 1} \ar[d]_f \ar[r] &
\mathcal{F}_{m + 1} \ar[d]_f \ar[r] &
\mathcal{F}_n \ar[d]_f \ar[r] & 0 \\
0 \ar[r] &
I^{n + k}\mathcal{F}_{m + 1} \ar[r] &
\mathcal{F}_{m + 1} \ar[r] &
\mathcal{F}_{n + k} \ar[r] & 0
}
$$
construit comme dans la démonstration du Lemme \ref{sites-cohomology-lemma-ML-general}.
Nous obtenons un morphisme induit en cohomologie et voyons que
$f \xi' \in H^{p + 1}(\mathcal{C}, I^{n + k}\mathcal{F}_{m + 1})$ a pour
image $f \xi$. Comme cela vaut pour tout $m \gg n + k$, nous voyons que
$f\xi$ appartient à $N_{n + k}$, comme voulu.

\medskip\noindent
Pour voir que l'image des morphismes de connexion $\delta_n$ est contenue
dans $N_n$, considérons les diagrammes
$$
\xymatrix{
0 \ar[r] &
I^n\mathcal{F}_{m + 1} \ar[d] \ar[r] &
\mathcal{F}_{m + 1} \ar[d] \ar[r] &
\mathcal{F}_n \ar[d] \ar[r] & 0 \\
0 \ar[r] &
I^n\mathcal{F}_{n + 1} \ar[r] &
\mathcal{F}_{n + 1} \ar[r] &
\mathcal{F}_n \ar[r] & 0
}
$$
pour $m \geq n$. L'examen des morphismes induits en cohomologie permet de
conclure.
\end{proof}

\begin{lemma}
\label{sites-cohomology-lemma-topology-I-adic-general}
Soit $I$ un idéal d'un anneau $A$. Soit $\mathcal{C}$ un site. Soit
$$
\ldots \to \mathcal{F}_3 \to \mathcal{F}_2 \to \mathcal{F}_1
$$
un système projectif de faisceaux de $A$-modules sur $\mathcal{C}$ tel que
$\mathcal{F}_n = \mathcal{F}_{n + 1}/I^n\mathcal{F}_{n + 1}$. Soit
$p \geq 0$. Supposons que
$$
\bigoplus\nolimits_{n \geq 0} H^p(\mathcal{C}, I^n\mathcal{F}_{n + 1})
$$
satisfasse la condition de chaîne ascendante comme
$\bigoplus_{n \geq 0} I^n/I^{n + 1}$-module gradué. Alors la topologie
limite sur $M = \lim H^p(\mathcal{C}, \mathcal{F}_n)$ est la topologie
$I$-adique.
\end{lemma}

\begin{proof}
Posons $F^n = \Ker(M \to H^p(\mathcal{C}, \mathcal{F}_n))$ pour
$n \geq 1$ et $F^0 = M$. Observons que
$I F^n \subset F^{n + 1}$. En particulier, $I^n M \subset F^n$. La
topologie $I$-adique est donc plus fine que la topologie limite. Pour la
réciproque, nous montrerons que, pour tout $n$, il existe $m \geq n$ tel
que $F^m \subset I^nM$\footnote{En fait, il existe $c \geq 0$ tel que
$F^{n + c} \subset I^nM$ pour tout $n$.}. Nous avons des morphismes
injectifs
$$
F^n/F^{n + 1} \longrightarrow H^p(\mathcal{C}, \mathcal{F}_{n + 1})
$$
dont l'image est contenue dans celle de
$H^p(\mathcal{C}, I^n\mathcal{F}_{n + 1}) \to
H^p(\mathcal{C}, \mathcal{F}_{n + 1})$.
Notons
$$
E_n \subset H^p(\mathcal{C}, I^n\mathcal{F}_{n + 1})
$$
l'image inverse de $F^n/F^{n + 1}$. Alors $\bigoplus E_n$ est un
sous-$\bigoplus I^n/I^{n + 1}$-module gradué de
$\bigoplus H^p(\mathcal{C}, I^n\mathcal{F}_{n + 1})$, et
$\bigoplus E_n \to \bigoplus F^n/F^{n + 1}$ est un homomorphisme de
modules gradués ; nous omettons les détails. Par hypothèse,
$\bigoplus E_n$ est engendré par un nombre fini d'éléments homogènes sur
$\bigoplus I^n/I^{n + 1}$. Puisque $E_n \to F^n/F^{n + 1}$ est surjectif,
il en va de même de $\bigoplus F^n/F^{n + 1}$. Nous pouvons donc trouver
$r$, des entiers $c_1, \ldots, c_r \geq 0$ et des $a_i \in F^{c_i}$ dont
les images engendrent $\bigoplus F^n/F^{n + 1}$. Posons
$c = \max(c_i)$.

\medskip\noindent
Pour $n \geq c$, nous affirmons que $I F^n = F^{n + 1}$. Cette
affirmation donne $F^{n + c} = I^nF^c \subset I^nM$, comme voulu. Pour la
démontrer, supposons $a \in F^{n + 1}$. L'image de $a$ dans
$F^{n + 1}/F^{n + 2}$ est une combinaison linéaire des $a_i$. Ainsi,
$a  - \sum f_i a_i \in F^{n + 2}$ pour certains
$f_i \in I^{n + 1 - c_i}$. Comme
$I^{n + 1 - c_i} = I \cdot I^{n - c_i}$ puisque $n \geq c_i$, nous
pouvons écrire $f_i = \sum g_{i, j} h_{i, j}$ avec $g_{i, j} \in I$ et
$h_{i, j}a_i \in F^n$. Nous voyons donc que
$F^{n + 1} = F^{n + 2} + IF^n$. Une récurrence immédiate donne
$F^{n + 1} = F^{n + e} + IF^n$ pour tout $e > 0$. Il en résulte que
$IF^n$ est dense dans $F^{n + 1}$. Choisissons des générateurs
$k_1, \ldots, k_r$ de $I$ et considérons l'application continue
$$
u : (F^n)^{\oplus r} \longrightarrow F^{n + 1},\quad
(x_1, \ldots, x_r) \mapsto \sum k_i x_i
$$
(pour la topologie limite). D'après ce qui précède, l'image de
$(F^m)^{\oplus r}$ par $u$ est dense dans $F^{m + 1}$ pour tout
$m \geq n$. Le lemme de l'application ouverte (Compléments d'algèbre,
Lemme \ref{more-algebra-lemma-open-mapping}) montre que $u$ est ouverte.
Ainsi, $u$ est surjective. Par conséquent, $IF^n = F^{n + 1}$ pour
$n \geq c$, ce
qui achève la démonstration.
\end{proof}

\begin{lemma}
\label{sites-cohomology-lemma-topology-I-adic-general-better}
Soit $I$ un idéal d'un anneau $A$. Soit $\mathcal{C}$ un site. Soit
$$
\ldots \to \mathcal{F}_3 \to \mathcal{F}_2 \to \mathcal{F}_1
$$
un système projectif de faisceaux de $A$-modules sur $\mathcal{C}$ tel que
$\mathcal{F}_n = \mathcal{F}_{n + 1}/I^n\mathcal{F}_{n + 1}$. Soit
$p \geq 0$. Pour tout $n$, posons
$$
N_n =
\bigcap\nolimits_{m \geq n}
\Im\left(
H^p(\mathcal{C}, I^n\mathcal{F}_{m + 1}) \to
H^p(\mathcal{C}, I^n\mathcal{F}_{n + 1})
\right)
$$
Si $\bigoplus N_n$ satisfait la condition de chaîne ascendante comme
$\bigoplus_{n \geq 0} I^n/I^{n + 1}$-module gradué, alors la topologie
limite sur $M = \lim H^p(\mathcal{C}, \mathcal{F}_n)$ est la topologie
$I$-adique.
\end{lemma}

\begin{proof}
La démonstration est exactement la même que celle du Lemme
\ref{sites-cohomology-lemma-topology-I-adic-general}. En effet, les arguments qui y sont
donnés établiront le résultat dès que nous aurons montré que
$\bigoplus N_n$ est un sous-$\bigoplus_{n \geq 0} I^n/I^{n + 1}$-module
gradué de $\bigoplus H^{p + 1}(\mathcal{C}, I^n\mathcal{F}_{n + 1})$ et que
$F^n/F^{n + 1} \subset H^p(\mathcal{C}, \mathcal{F}_{n + 1})$ est contenu
dans l'image de
$N_n \to H^p(\mathcal{C}, \mathcal{F}_{n + 1})$. Dans la démonstration du
Lemme \ref{sites-cohomology-lemma-ML-general-better}, nous avons établi l'assertion sur la
structure de module.

\medskip\noindent
Soit $t \in F^n$. Choisissons un élément
$s \in H^p(\mathcal{C}, I^n\mathcal{F}_{n + 1})$
dont l'image est l'image de $t$ dans
$H^p(\mathcal{C}, \mathcal{F}_{n + 1})$. Nous devons montrer que $s$
appartient à $N_n$. Or $F^n$ est le noyau du morphisme
$M \to H^p(\mathcal{C}, \mathcal{F}_n)$ ; pour tout $m \geq n$, nous
pouvons donc envoyer $t$ sur un élément
$t_m \in H^p(\mathcal{C}, \mathcal{F}_{m + 1})$ dont l'image dans
$H^p(\mathcal{C}, \mathcal{F}_n)$ est nulle. Considérons la suite de
cohomologie
$$
H^{p - 1}(\mathcal{C}, \mathcal{F}_n) \to
H^p(\mathcal{C}, I^n\mathcal{F}_{m + 1}) \to
H^p(\mathcal{C}, \mathcal{F}_{m + 1}) \to
H^p(\mathcal{C}, \mathcal{F}_n)
$$
provenant de la suite exacte courte
$0 \to I^n\mathcal{F}_{m + 1} \to \mathcal{F}_{m + 1} \to \mathcal{F}_n \to 0$.
Nous pouvons choisir
$s_m \in H^p(\mathcal{C}, I^n\mathcal{F}_{m + 1})$ d'image $t_m$. En
comparant la suite ci-dessus à celle obtenue pour $m = n$, nous voyons que
l'image de $s_m$ est égale à $s$ à un élément de l'image de
$H^{p - 1}(\mathcal{C}, \mathcal{F}_n) \to
H^p(\mathcal{C}, I^n\mathcal{F}_{n + 1})$.
près. Cependant, ce morphisme se factorise par le morphisme
$H^p(\mathcal{C}, I^n\mathcal{F}_{m + 1}) \to
H^p(\mathcal{C}, I^n\mathcal{F}_{n + 1})$
et nous voyons que $s$ appartient à l'image, comme voulu.
\end{proof}












\section{Limites dérivées et homotopiques}
\label{sites-cohomology-section-derived-limits}

\noindent
Soit $\mathcal{C}$ un site. Considérons la catégorie
$\mathcal{C} \times \mathbf{N}$ pour laquelle
$\Mor((U, n), (V, m)) = \emptyset$ si $n > m$ et
$\Mor((U, n), (V, m)) = \Mor(U, V)$ sinon. Munissons-la d'une structure
de site en prenant pour recouvrements les familles
$\{(U_i, n) \to (U, n)\}$ telles que $\{U_i \to U\}$ soit un recouvrement
de $\mathcal{C}$. Le lecteur vérifie alors immédiatement que les faisceaux
sur $\mathcal{C} \times \mathbf{N}$ s'identifient aux systèmes projectifs
de faisceaux sur $\mathcal{C}$. En particulier,
$\textit{Ab}(\mathcal{C} \times \mathbf{N})$ est la catégorie des systèmes
projectifs de faisceaux abéliens sur $\mathcal{C}$. Considérons maintenant
le foncteur
$$
\lim : \textit{Ab}(\mathcal{C} \times \mathbf{N}) \to \textit{Ab}(\mathcal{C})
$$
qui associe à un système projectif sa limite. Ce n'est autre que $g_*$, où
$g : \Sh(\mathcal{C} \times \mathbf{N}) \to \Sh(\mathcal{C})$ est le
morphisme de topos associé au foncteur continu et cocontinu
$\mathcal{C} \times \mathbf{N} \to \mathcal{C}$. (Observons que $g^{-1}$
associe à un faisceau sur $\mathcal{C}$ le système projectif constant
correspondant.)

\medskip\noindent
Le formalisme général exposé ci-dessus fournit un foncteur dérivé
$$
R\lim = Rg_* : D(\mathcal{C} \times \mathbf{N}) \to D(\mathcal{C}).
$$
Comme indiqué, ce foncteur est souvent noté $R\lim$.

\medskip\noindent
Par ailleurs, les foncteurs continus et cocontinus
$\mathcal{C} \to \mathcal{C} \times \mathbf{N}$,
$U \mapsto (U, n)$ définissent des morphismes de topos
$i_n : \Sh(\mathcal{C}) \to \Sh(\mathcal{C} \times \mathbf{N})$.
Bien entendu, $i_n^{-1}$ est le foncteur qui extrait le $n$-ième terme du
système projectif. Il existe donc des transformations de foncteurs
$i_{n + 1}^{-1} \to i_n^{-1}$. Ainsi, pour
$K \in D(\mathcal{C} \times \mathbf{N})$, nous obtenons
$K_n = i_n^{-1}K \in D(\mathcal{C})$ et des morphismes
$K_{n + 1} \to K_n$. Dans Catégories dérivées, Définition
\ref{derived-definition-derived-limit}, nous avons défini la notion de
limite homotopique
$$
R\lim K_n \in D(\mathcal{C})
$$
Nous affirmons que les deux notions coïncident (dans la mesure où cette
assertion a un sens).

\begin{lemma}
\label{sites-cohomology-lemma-derived-limit-is-ok}
Soit $\mathcal{C}$ un site. Soit $K$ un objet de
$D(\mathcal{C} \times \mathbf{N})$. Posons $K_n = i_n^{-1}K$ comme
ci-dessus. Alors
$$
R\lim K \cong R\lim K_n
$$
dans $D(\mathcal{C})$.
\end{lemma}

\begin{proof}
Pour calculer $R\lim$ sur un objet $K$ de
$D(\mathcal{C} \times \mathbf{N})$, choisissons un représentant K-injectif
$\mathcal{I}^\bullet$ dont les termes sont des objets injectifs de
$\textit{Ab}(\mathcal{C} \times \mathbf{N})$ ; voir Injectifs, Théorème
\ref{injectives-theorem-K-injective-embedding-grothendieck}. Nous pouvons,
et allons, considérer $\mathcal{I}^\bullet$ comme un système projectif de
complexes $(\mathcal{I}_n^\bullet)$. Nous voyons alors que
$$
R\lim K = \lim \mathcal{I}_n^\bullet
$$
où le membre de droite est la limite projective prise terme à terme.

\medskip\noindent
Soit $\mathcal{J} = (\mathcal{J}_n)$ un objet injectif de
$\textit{Ab}(\mathcal{C} \times \mathbf{N})$. Les morphismes
$(U, n) \to (U, n + 1)$ sont des monomorphismes de
$\mathcal{C} \times \mathbf{N}$ ; par conséquent,
$\mathcal{J}(U, n + 1) \to \mathcal{J}(U, n)$ est surjectif (Lemme
\ref{sites-cohomology-lemma-restriction-along-monomorphism-surjective}). Il s'ensuit que
$\mathcal{J}_{n + 1} \to \mathcal{J}_n$ est surjectif comme morphisme de
préfaisceaux.

\medskip\noindent
Notons que le foncteur $i_n^{-1}$ admet un adjoint à gauche exact
$i_{n, !}$. En effet, $i_{n, !}\mathcal{F}$ est le système projectif
$\ldots 0 \to 0 \to \mathcal{F} \to \ldots \to \mathcal{F}$. Les
complexes $i_n^{-1}\mathcal{I}^\bullet = \mathcal{I}_n^\bullet$ sont donc
K-injectifs d'après Catégories dérivées, Lemme
\ref{derived-lemma-adjoint-preserve-K-injectives}.

\medskip\noindent
Comme nous avons choisi notre complexe K-injectif à termes injectifs, nous
en déduisons que
$$
0 \to  \lim \mathcal{I}_n^\bullet \to \prod \mathcal{I}_n^\bullet
\to \prod \mathcal{I}_n^\bullet \to 0
$$
est une suite exacte courte de complexes de faisceaux abéliens, puisqu'elle
est une suite exacte courte de complexes de préfaisceaux abéliens. De plus,
les produits du milieu et de droite représentent les produits dans
$D(\mathcal{C})$ ; voir Injectifs, Lemme
\ref{injectives-lemma-derived-products}, et sa démonstration (c'est ici que
nous utilisons la K-injectivité de $\mathcal{I}_n^\bullet$). Ainsi,
$R\lim K$ est une limite homotopique du système projectif $(K_n)$, par
définition des limites homotopiques dans les catégories triangulées.
\end{proof}

\begin{lemma}
\label{sites-cohomology-lemma-RGamma-commutes-with-Rlim}
Soit $(\mathcal{C}, \mathcal{O})$ un site annelé. Les foncteurs
$R\Gamma(\mathcal{C}, -)$ et $R\Gamma(U, -)$, pour
$U \in \Ob(\mathcal{C})$, commutent à $R\lim$. De plus, il existe des
suites exactes courtes
$$
0 \to
R^1\lim H^{m - 1}(U, K_n) \to H^m(U, R\lim K_n) \to
\lim H^m(U, K_n) \to 0
$$
pour tout système projectif $(K_n)$ dans $D(\mathcal{O})$ et tout
$m \in \mathbf{Z}$. Il en va de même pour
$H^m(\mathcal{C}, R\lim K_n)$.
\end{lemma}

\begin{proof}
La première assertion résulte d'Injectifs, Lemme
\ref{injectives-lemma-RF-commutes-with-Rlim}. Nous pouvons ensuite appliquer
Compléments d'algèbre, Remarque
\ref{more-algebra-remark-compare-derived-limit}, à
$R\lim R\Gamma(U, K_n) = R\Gamma(U, R\lim K_n)$ pour obtenir les suites
exactes courtes.
\end{proof}

\begin{lemma}
\label{sites-cohomology-lemma-Rf-commutes-with-Rlim}
Soit $f : (\Sh(\mathcal{C}), \mathcal{O}) \to (\Sh(\mathcal{C}'), \mathcal{O}')$
un morphisme de topos annelés. Alors $Rf_*$ commute à $R\lim$, c'est-à-dire
que $Rf_*$ commute aux limites dérivées.
\end{lemma}

\begin{proof}
Soit $(K_n)$ un système projectif d'objets de $D(\mathcal{O})$. Par
récurrence sur $n$, nous pouvons choisir des complexes
$\mathcal{K}_n^\bullet$ de $\mathcal{O}$-modules et des morphismes de
complexes $\mathcal{K}_{n + 1}^\bullet \to \mathcal{K}_n^\bullet$
représentant les morphismes $K_{n + 1} \to K_n$ dans $D(\mathcal{O})$.
Autrement dit, il existe un objet $K$ de
$D(\mathcal{C} \times \mathbf{N})$ dont le système projectif associé est
celui qui nous est donné. Considérons ensuite le diagramme commutatif
$$
\xymatrix{
\Sh(\mathcal{C} \times \mathbf{N}) \ar[r]_g \ar[d]_{f \times 1} &
\Sh(\mathcal{C}) \ar[d]_f \\
\Sh(\mathcal{C}' \times \mathbf{N}) \ar[r]^{g'} &
\Sh(\mathcal{C}')
}
$$
de morphismes de topos. Il en résulte que
$R\lim R(f \times 1)_*K = Rf_* R\lim K$. En explicitant les définitions et
en utilisant le Lemme \ref{sites-cohomology-lemma-derived-limit-is-ok}, nous obtenons
$R\lim (Rf_*K_n) = Rf_*(R\lim K_n)$.

\medskip\noindent
Autre démonstration lorsque $\mathcal{C}$ possède suffisamment de points.
Considérons le triangle distingué qui définit
$$
R\lim K_n \to \prod K_n \to \prod K_n
$$
dans $D(\mathcal{O})$. En appliquant le foncteur exact $Rf_*$, nous obtenons
le triangle distingué
$$
Rf_*(R\lim K_n) \to Rf_*\left(\prod K_n\right) \to Rf_*\left(\prod K_n\right)
$$
dans $D(\mathcal{O}')$. Il suffit donc de prouver que $Rf_*$ commute aux
produits dans la catégorie dérivée (lesquels ne sont pas simplement donnés
par les produits de complexes ; voir Injectifs, Lemme
\ref{injectives-lemma-derived-products}). Or, comme $Rf_*$ est un adjoint à
droite d'après le Lemme \ref{sites-cohomology-lemma-adjoint}, cela résulte formellement de
Catégories, Lemme \ref{categories-lemma-adjoint-exact}. Attention : nous ne
pouvons pas appliquer directement Catégories, Lemme
\ref{categories-lemma-adjoint-exact}, car $R\lim K_n$ n'est pas une limite
dans $D(\mathcal{O})$.
\end{proof}

\begin{remark}
\label{sites-cohomology-remark-discuss-derived-limit}
Soit $(\mathcal{C}, \mathcal{O})$ un site annelé. Soit $(K_n)$ un système
projectif dans $D(\mathcal{O})$. Posons $K = R\lim K_n$. Pour chaque $n$
et $m$, soit $\mathcal{H}^m_n = H^m(K_n)$ le $m$-ième faisceau de
cohomologie de $K_n$, et posons de même $\mathcal{H}^m = H^m(K)$. Notons
$\underline{\mathcal{H}}^m_n$ le préfaisceau
$$
U \longmapsto \underline{\mathcal{H}}^m_n(U) = H^m(U, K_n)
$$
De même, posons $\underline{\mathcal{H}}^m(U) = H^m(U, K)$. D'après le
Lemme \ref{sites-cohomology-lemma-sheafification-cohomology}, $\mathcal{H}^m_n$ est le
faisceau associé à $\underline{\mathcal{H}}^m_n$, et $\mathcal{H}^m$ est
le faisceau associé à $\underline{\mathcal{H}}^m$. Nous avons le diagramme
$$
\xymatrix{
K \ar@{=}[d] &
\underline{\mathcal{H}}^m \ar[d] \ar[r] & 
\mathcal{H}^m \ar[d] \\
R\lim K_n &
\lim \underline{\mathcal{H}}^m_n \ar[r] & 
\lim \mathcal{H}^m_n
}
$$
En général, $\lim \mathcal{H}^m_n$ n'est pas nécessairement le faisceau
associé à $\lim \underline{\mathcal{H}}^m_n$. Si
$U \in \mathcal{C}$, nous avons les suites exactes courtes
\begin{equation}
\label{sites-cohomology-equation-ses-Rlim-over-U}
0 \to
R^1\lim \underline{\mathcal{H}}^{m - 1}_n(U) \to
\underline{\mathcal{H}}^m(U) \to
\lim \underline{\mathcal{H}}^m_n(U) \to 0
\end{equation}
d'après le Lemme \ref{sites-cohomology-lemma-RGamma-commutes-with-Rlim}.
\end{remark}

\noindent
Le lemme suivant s'applique à un système projectif de modules
quasi-cohérents dont les morphismes de transition sont surjectifs sur un
espace algébrique ou un champ algébrique.

\begin{lemma}
\label{sites-cohomology-lemma-inverse-limit-is-derived-limit}
Soit $(\mathcal{C}, \mathcal{O})$ un site annelé. Soit $(\mathcal{F}_n)$ un
système projectif de $\mathcal{O}$-modules. Soit
$\mathcal{B} \subset \Ob(\mathcal{C})$ une partie. Supposons que
\begin{enumerate}
\item tout objet de $\mathcal{C}$ possède un recouvrement dont les membres
appartiennent à $\mathcal{B}$,
\item $H^p(U, \mathcal{F}_n) = 0$ pour $p > 0$ et $U \in \mathcal{B}$,
\item le système projectif $\mathcal{F}_n(U)$ a un $R^1\lim$ nul pour
$U \in \mathcal{B}$.
\end{enumerate}
Alors $R\lim \mathcal{F}_n = \lim \mathcal{F}_n$ et nous avons
$H^p(U, \lim \mathcal{F}_n) = 0$ pour $p > 0$ et $U \in \mathcal{B}$.
\end{lemma}

\begin{proof}
Posons $K_n = \mathcal{F}_n$ et $K = R\lim \mathcal{F}_n$. Avec les
notations de la Remarque \ref{sites-cohomology-remark-discuss-derived-limit}, l'hypothèse
(2) donne, pour $U \in \mathcal{B}$,
$\underline{\mathcal{H}}_n^m(U) = 0$ si $m \not = 0$ et
$\underline{\mathcal{H}}_n^0(U) = \mathcal{F}_n(U)$. L'Équation
(\ref{sites-cohomology-equation-ses-Rlim-over-U}) et l'hypothèse (3) donnent
$\underline{\mathcal{H}}^m(U) = 0$ si $m \not = 0$, et ce groupe est égal à
$\lim \mathcal{F}_n(U)$ si $m = 0$. En prenant les faisceaux associés et en
utilisant (1), nous obtenons $\mathcal{H}^m = 0$ si $m \not = 0$, et ce
faisceau est égal à $\lim \mathcal{F}_n$ si $m = 0$. Ainsi
$K = \lim \mathcal{F}_n$. Comme
$H^m(U, K) = \underline{\mathcal{H}}^m(U) = 0$ pour $m > 0$ (voir
ci-dessus), la seconde assertion est satisfaite.
\end{proof}

\begin{lemma}
\label{sites-cohomology-lemma-cohomology-derived-limit-injective}
Soit $(\mathcal{C}, \mathcal{O})$ un site annelé. Soit $(K_n)$ un système
projectif dans $D(\mathcal{O})$. Soient $V \in \Ob(\mathcal{C})$ et
$m \in \mathbf{Z}$. Supposons qu'il existe un entier $n(V)$ et un système
cofinal $\text{Cov}_V$ de recouvrements de $V$ tels que, pour
$\{V_i \to V\} \in \text{Cov}_V$,
\begin{enumerate}
\item $R^1\lim H^{m - 1}(V_i, K_n) = 0$, et
\item $H^m(V_i, K_n) \to H^m(V_i, K_{n(V)})$ soit injectif pour
$n \geq n(V)$.
\end{enumerate}
Alors le morphisme sur les sections
$H^m(R\lim K_n)(V) \to H^m(K_{n(V)})(V)$ est injectif.
\end{lemma}

\begin{proof}
Soit $\gamma \in H^m(R\lim K_n)(V)$ dont l'image dans
$H^m(K_{n(V)})(V)$ est nulle. Comme $H^m(R\lim K_n)$ est le faisceau
associé au préfaisceau $U \mapsto H^m(U, R\lim K_n)$ (d'après le Lemme
\ref{sites-cohomology-lemma-sheafification-cohomology}), nous pouvons choisir
$\{V_i \to V\} \in \text{Cov}_V$ et des éléments
$\tilde\gamma_i \in H^m(V_i, R\lim K_n)$ dont les images sont les
$\gamma|_{V_i}$. L'image de $\tilde\gamma_i$ est alors un élément
$\tilde\gamma_{i, n(V)} \in H^m(V_i, K_{n(V)})$. Puisque
$H^m(K_{n(V)})$ est le faisceau associé à
$U \mapsto H^m(U, K_{n(V)})$ (à nouveau d'après le Lemme
\ref{sites-cohomology-lemma-sheafification-cohomology}), nous pouvons, après avoir remplacé
$\{V_i \to V\}$ par un raffinement, supposer que
$\tilde\gamma_{i, n(V)} = 0$ pour tout $i$. Pour ce recouvrement,
considérons les suites exactes courtes
$$
0 \to
R^1\lim H^{m - 1}(V_i, K_n) \to H^m(V_i, R\lim K_n) \to
\lim H^m(V_i, K_n) \to 0
$$
du Lemme \ref{sites-cohomology-lemma-RGamma-commutes-with-Rlim}. D'après l'hypothèse (1),
le groupe de gauche est nul et, d'après l'hypothèse (2), le groupe de
droite s'injecte dans $H^m(V_i, K_{n(V)})$. Nous en déduisons que
$\tilde\gamma_i = 0$, puis que $\gamma = 0$, comme voulu.
\end{proof}

\begin{lemma}
\label{sites-cohomology-lemma-is-limit-per-object}
Soit $(\mathcal{C}, \mathcal{O})$ un site annelé. Soit
$E \in D(\mathcal{O})$. Soit $\mathcal{B} \subset \Ob(\mathcal{C})$ une
partie. Supposons que
\begin{enumerate}
\item tout objet de $\mathcal{C}$ possède un recouvrement dont les membres
appartiennent à $\mathcal{B}$, et
\item pour tout $V \in \mathcal{B}$, il existe une fonction
$p(V, -) : \mathbf{Z} \to \mathbf{Z}$ et un système cofinal
$\text{Cov}_V$ de recouvrements de $V$ tels que
$$
H^p(V_i, H^{m - p}(E)) = 0
$$
pour tout $\{V_i \to V\} \in \text{Cov}_V$ et tous entiers $p, m$
vérifiant $p > p(V, m)$.
\end{enumerate}
Alors le morphisme $E \to R\lim \tau_{\geq -n} E$ de Catégories dérivées,
Remarque
\ref{derived-remark-map-into-derived-limit-truncations}
est un isomorphisme dans $D(\mathcal{O})$.
\end{lemma}

\begin{proof}
Posons $K_n = \tau_{\geq -n}E$ et $K = R\lim K_n$. Le morphisme canonique
$E \to K$ provient des morphismes canoniques
$E \to K_n = \tau_{\geq -n}E$. Nous devons montrer que $E \to K$ induit
un isomorphisme $H^m(E) \to H^m(K)$ de faisceaux de cohomologie. Dans la
suite de la démonstration, fixons $m$. Si $n \geq -m$, le morphisme
$E \to \tau_{\geq -n}E = K_n$ induit un isomorphisme
$H^m(E) \to H^m(K_n)$. Pour achever la démonstration, il suffit de montrer
que, pour tout $V \in \mathcal{B}$, il existe un entier $n(V) \geq -m$ tel
que le morphisme $H^m(K)(V) \to H^m(K_{n(V)})(V)$ soit injectif. En effet,
la composée
$$
H^m(E)(V) \to H^m(K)(V) \to H^m(K_{n(V)})(V)
$$
est alors une bijection et la seconde flèche est injective ; la première
flèche est donc bijective. La propriété (1) implique alors que
$H^m(E) \to H^m(K)$ est un isomorphisme. Posons
$$
n(V) = 1 + \max\{-m, p(V, m - 1) - m, -1 + p(V, m) - m, -2 + p(V, m + 1) - m\}.
$$
de sorte que, dans tous les cas, $n(V) \geq -m$. Affirmation : les
morphismes
$$
H^{m - 1}(V_i, K_{n + 1}) \to H^{m - 1}(V_i, K_n)
\quad\text{et}\quad
H^m(V_i, K_{n + 1}) \to H^m(V_i, K_n)
$$
sont des isomorphismes pour $n \geq n(V)$ et
$\{V_i \to V\} \in \text{Cov}_V$. L'affirmation implique que les
conditions (1) et (2) du Lemme
\ref{sites-cohomology-lemma-cohomology-derived-limit-injective} sont satisfaites, et donc
l'injectivité voulue. Rappelons (Catégories dérivées, Remarque
\ref{derived-remark-truncation-distinguished-triangle})
que nous avons des triangles distingués
$$
H^{-n - 1}(E)[n + 1] \to
K_{n + 1} \to K_n \to H^{-n - 1}(E)[n + 2]
$$
En examinant la suite exacte longue de cohomologie associée, l'affirmation
résulte de l'annulation de
$$
H^{m + n}(V_i, H^{-n - 1}(E)),\quad
H^{m + n + 1}(V_i, H^{-n - 1}(E)),\quad
H^{m + n + 2}(V_i, H^{-n - 1}(E))
$$
pour $n \geq n(V)$ et $\{V_i \to V\} \in \text{Cov}_V$. Cela découle de
notre choix de $n(V)$ et de l'hypothèse du lemme.
\end{proof}

\begin{lemma}
\label{sites-cohomology-lemma-is-limit-spaltenstein}
Soit $(\mathcal{C}, \mathcal{O})$ un site annelé. Soit
$E \in D(\mathcal{O})$. Soit $\mathcal{B} \subset \Ob(\mathcal{C})$ une
partie. Supposons que
\begin{enumerate}
\item tout objet de $\mathcal{C}$ possède un recouvrement dont les membres
appartiennent à $\mathcal{B}$, et
\item pour tout $V \in \mathcal{B}$, il existe un entier $d_V \geq 0$ et un
système cofinal $\text{Cov}_V$ de recouvrements de $V$ tels que
$$
H^p(V_i, H^q(E)) = 0 \text{ pour }
\{V_i \to V\} \in \text{Cov}_V,\ p > d_V, \text{ et }q < 0
$$
\end{enumerate}
Alors le morphisme $E \to R\lim \tau_{\geq -n} E$ de Catégories dérivées,
Remarque
\ref{derived-remark-map-into-derived-limit-truncations}
est un isomorphisme dans $D(\mathcal{O})$.
\end{lemma}

\begin{proof}
Cela résulte du Lemme \ref{sites-cohomology-lemma-is-limit-per-object}, avec
$p(V, m) = d_V + \max(0, m)$.
\end{proof}

\begin{lemma}
\label{sites-cohomology-lemma-is-limit}
Soit $(\mathcal{C}, \mathcal{O})$ un site annelé. Soit
$E \in D(\mathcal{O})$. Supposons qu'il existe une fonction
$p(-) : \mathbf{Z} \to \mathbf{Z}$ et une partie
$\mathcal{B} \subset \Ob(\mathcal{C})$ telles que
\begin{enumerate}
\item tout objet de $\mathcal{C}$ possède un recouvrement dont les membres
appartiennent à $\mathcal{B}$,
\item $H^p(V, H^{m - p}(E)) = 0$ pour $p > p(m)$ et
$V \in \mathcal{B}$.
\end{enumerate}
Alors le morphisme $E \to R\lim \tau_{\geq -n} E$ de Catégories dérivées,
Remarque
\ref{derived-remark-map-into-derived-limit-truncations}
est un isomorphisme dans $D(\mathcal{O})$.
\end{lemma}

\begin{proof}
Appliquons le Lemme \ref{sites-cohomology-lemma-is-limit-per-object} avec
$p(V, m) = p(m)$ et $\text{Cov}_V$ égal à l'ensemble des recouvrements
$\{V_i \to V\}$ tels que $V_i \in \mathcal{B}$ pour tout $i$.
\end{proof}

\begin{lemma}
\label{sites-cohomology-lemma-is-limit-dimension}
Soit $(\mathcal{C}, \mathcal{O})$ un site annelé. Soit
$E \in D(\mathcal{O})$. Supposons qu'il existe un entier $d \geq 0$ et
une partie $\mathcal{B} \subset \Ob(\mathcal{C})$ tels que
\begin{enumerate}
\item tout objet de $\mathcal{C}$ possède un recouvrement dont les membres
appartiennent à $\mathcal{B}$,
\item $H^p(V, H^q(E)) = 0$ pour $p > d$, $q < 0$ et
$V \in \mathcal{B}$.
\end{enumerate}
Alors le morphisme $E \to R\lim \tau_{\geq -n} E$ de Catégories dérivées,
Remarque
\ref{derived-remark-map-into-derived-limit-truncations}
est un isomorphisme dans $D(\mathcal{O})$.
\end{lemma}

\begin{proof}
Appliquons le Lemme \ref{sites-cohomology-lemma-is-limit-spaltenstein} avec $d_V = d$ et
$\text{Cov}_V$ égal à l'ensemble des recouvrements $\{V_i \to V\}$ tels
que $V_i \in \mathcal{B}$ pour tout $i$.
\end{proof}

\noindent
Les lemmes précédents permettent de calculer la cohomologie dans certaines
situations.

\begin{lemma}
\label{sites-cohomology-lemma-cohomology-over-U-trivial}
Soit $(\mathcal{C}, \mathcal{O})$ un site annelé. Soit $K$ un objet de
$D(\mathcal{O})$. Soit $\mathcal{B} \subset \Ob(\mathcal{C})$ une partie.
Supposons que
\begin{enumerate}
\item tout objet de $\mathcal{C}$ possède un recouvrement dont les membres
appartiennent à $\mathcal{B}$,
\item $H^p(U, H^q(K)) = 0$ pour tous $p > 0$, $q \in \mathbf{Z}$ et
$U \in \mathcal{B}$.
\end{enumerate}
Alors $H^q(U, K) = H^0(U, H^q(K))$ pour $q \in \mathbf{Z}$ et
$U \in \mathcal{B}$.
\end{lemma}

\begin{proof}
Observons que $K = R\lim \tau_{\geq -n} K$ d'après le Lemme
\ref{sites-cohomology-lemma-is-limit-dimension}, avec $d = 0$. Soit
$U \in \mathcal{B}$. L'Équation (\ref{sites-cohomology-equation-ses-Rlim-over-U}) donne une
suite exacte courte
$$
0 \to R^1\lim H^{q - 1}(U, \tau_{\geq -n}K) \to
H^q(U, K) \to \lim H^q(U, \tau_{\geq -n}K) \to 0
$$
La condition (2) implique
$H^q(U, \tau_{\geq -n} K) = H^0(U, H^q(\tau_{\geq -n} K))$
pour tout $q$, en utilisant la suite spectrale de Catégories dérivées,
Lemme \ref{derived-lemma-two-ss-complex-functor}. Cette suite spectrale
converge car $\tau_{\geq -n}K$ est borné inférieurement. Si $n > -q$,
nous avons $H^q(\tau_{\geq -n}K) = H^q(K)$. Ainsi, les systèmes à gauche
et à droite dans la suite exacte courte affichée sont stationnaires, de
valeurs respectives $H^0(U, H^{q - 1}(K))$ et $H^0(U, H^q(K))$, d'où le
lemme.
\end{proof}

\noindent
Voici un autre cas où nous pouvons décrire la limite dérivée.

\begin{lemma}
\label{sites-cohomology-lemma-derived-limit-suitable-system}
Soit $(\mathcal{C}, \mathcal{O})$ un site annelé. Soit $(K_n)$ un système
projectif d'objets de $D(\mathcal{O})$. Soit
$\mathcal{B} \subset \Ob(\mathcal{C})$ une partie. Supposons que
\begin{enumerate}
\item tout objet de $\mathcal{C}$ possède un recouvrement dont les membres
appartiennent à $\mathcal{B}$,
\item pour tous $U \in \mathcal{B}$ et $q \in \mathbf{Z}$, nous avons
\begin{enumerate}
\item $H^p(U, H^q(K_n)) = 0$ pour $p > 0$,
\item le système projectif $H^0(U, H^q(K_n))$ a un $R^1\lim$ nul.
\end{enumerate}
\end{enumerate}
Alors $H^q(R\lim K_n) = \lim H^q(K_n)$ pour $q \in \mathbf{Z}$.
\end{lemma}

\begin{proof}
Posons $K = R\lim K_n$. Nous utiliserons les notations de la Remarque
\ref{sites-cohomology-remark-discuss-derived-limit}. Soit $U \in \mathcal{B}$. D'après le
Lemme \ref{sites-cohomology-lemma-cohomology-over-U-trivial} et (2)(a), nous avons
$H^q(U, K_n) = H^0(U, H^q(K_n))$. Comme le foncteur $R\Gamma(U, -)$
commute aux limites dérivées, nous avons
$$
H^q(U, K) = H^q(R\lim R\Gamma(U, K_n)) = \lim H^0(U, H^q(K_n))
$$
où la dernière égalité résulte de Compléments d'algèbre, Remarque
\ref{more-algebra-remark-compare-derived-limit}, et de l'hypothèse (2)(b).
Ainsi, $H^q(U, K)$ est la limite projective des sections des faisceaux
$H^q(K_n)$ sur $U$. Comme $\lim H^q(K_n)$ est un faisceau, l'hypothèse
(1) montre que $H^q(K)$, qui est le faisceau associé au préfaisceau
$U \mapsto H^q(U, K)$, est égal à $\lim H^q(K_n)$. Cela démontre le lemme.
\end{proof}







\section{Construction de résolutions K-injectives}
\label{sites-cohomology-section-K-injective}

\noindent
Soit $(\mathcal{C}, \mathcal{O})$ un site annelé. Soit
$\mathcal{F}^\bullet$ un complexe de $\mathcal{O}$-modules. La catégorie
$\textit{Mod}(\mathcal{O})$ possède suffisamment d'injectifs ; nous pouvons
donc utiliser Catégories dérivées, Lemme
\ref{derived-lemma-special-inverse-system}, pour construire un diagramme
$$
\xymatrix{
\ldots \ar[r] &
\tau_{\geq -2}\mathcal{F}^\bullet \ar[r] \ar[d] &
\tau_{\geq -1}\mathcal{F}^\bullet \ar[d] \\
\ldots \ar[r] & \mathcal{I}_2^\bullet \ar[r] & \mathcal{I}_1^\bullet
}
$$
dans la catégorie des complexes de $\mathcal{O}$-modules, tel que
\begin{enumerate}
\item les flèches verticales soient des quasi-isomorphismes,
\item $\mathcal{I}_n^\bullet$ soit un complexe borné inférieurement
d'injectifs,
\item les flèches $\mathcal{I}_{n + 1}^\bullet \to \mathcal{I}_n^\bullet$
soient des surjections scindées terme à terme.
\end{enumerate}
La catégorie des $\mathcal{O}$-modules admet des limites (elles se calculent
au niveau des préfaisceaux) ; nous pouvons donc former la limite terme à
terme $\mathcal{I}^\bullet = \lim_n \mathcal{I}_n^\bullet$. D'après
Catégories dérivées, Lemmes
\ref{derived-lemma-bounded-below-injectives-K-injective} et
\ref{derived-lemma-limit-K-injectives}
il s'agit d'un complexe K-injectif. En général, le morphisme canonique
\begin{equation}
\label{sites-cohomology-equation-into-candidate-K-injective}
\mathcal{F}^\bullet \to \mathcal{I}^\bullet
\end{equation}
n'est pas nécessairement un quasi-isomorphisme. Le lemme suivant donne des
conditions qui assurent qu'il en est un.

\begin{lemma}
\label{sites-cohomology-lemma-K-injective}
Plaçons-nous dans la situation décrite ci-dessus. Notons
$\mathcal{H}^m = H^m(\mathcal{F}^\bullet)$ le $m$-ième faisceau de
cohomologie. Soit $\mathcal{B} \subset \Ob(\mathcal{C})$ une partie et
soit $d \in \mathbf{N}$. Supposons que
\begin{enumerate}
\item tout objet de $\mathcal{C}$ possède un recouvrement dont les membres
appartiennent à $\mathcal{B}$,
\item pour tout $U \in \mathcal{B}$, nous ayons
$H^p(U, \mathcal{H}^q) = 0$ pour $p > d$ et $q < 0$\footnote{Il suffit
que $\forall m$, $\exists p(m)$,
$H^p(U, \mathcal{H}^{m - p}) = 0$ pour $p > p(m)$ ; voir le Lemme
\ref{sites-cohomology-lemma-is-limit}.}.
\end{enumerate}
Alors (\ref{sites-cohomology-equation-into-candidate-K-injective}) est un
quasi-isomorphisme.
\end{lemma}

\begin{proof}
Il suffit, d'après Catégories dérivées, Lemme
\ref{derived-lemma-difficulty-K-injectives}, de montrer que le morphisme
$\mathcal{F}^\bullet \to R\lim \tau_{\geq -n} \mathcal{F}^\bullet$
est un isomorphisme. Cela résulte du Lemme
\ref{sites-cohomology-lemma-is-limit-dimension}.
\end{proof}

\noindent
Voici un lemme technique sur les faisceaux de cohomologie des limites terme
à terme de systèmes projectifs de complexes de modules. Il convient
d'éviter autant que possible d'utiliser ce lemme et de privilégier les
arguments faisant intervenir des limites projectives dérivées.

\begin{lemma}
\label{sites-cohomology-lemma-inverse-limit-complexes}
Soit $(\mathcal{C}, \mathcal{O})$ un site annelé. Soit
$(\mathcal{F}_n^\bullet)$ un système projectif de complexes de
$\mathcal{O}$-modules. Soit $m \in \mathbf{Z}$. Supposons donnés une partie
$\mathcal{B} \subset \Ob(\mathcal{C})$ et un entier $n_0$ tels que
\begin{enumerate}
\item tout objet de $\mathcal{C}$ possède un recouvrement dont les membres
appartiennent à $\mathcal{B}$,
\item pour tout $U \in \mathcal{B}$,
\begin{enumerate}
\item les systèmes projectifs de groupes abéliens
$\mathcal{F}_n^{m - 2}(U)$ et $\mathcal{F}_n^{m - 1}(U)$
aient un $R^1\lim$ nul (par exemple, qu'ils satisfassent la condition de
Mittag-Leffler),
\item le système projectif de groupes abéliens
$H^{m - 1}(\mathcal{F}_n^\bullet(U))$ ait un $R^1\lim$ nul (par exemple,
qu'il satisfasse la condition de Mittag-Leffler), et
\item nous ayons
$H^m(\mathcal{F}_n^\bullet(U)) = H^m(\mathcal{F}_{n_0}^\bullet(U))$
pour tout $n \geq n_0$.
\end{enumerate}
\end{enumerate}
Alors les morphismes
$H^m(\mathcal{F}^\bullet) \to \lim H^m(\mathcal{F}_n^\bullet)
\to H^m(\mathcal{F}_{n_0}^\bullet)$ sont des isomorphismes de faisceaux,
où $\mathcal{F}^\bullet = \lim \mathcal{F}_n^\bullet$ est la limite
projective terme à terme.
\end{lemma}

\begin{proof}
Soit $U \in \mathcal{B}$. Notons que
$H^m(\mathcal{F}^\bullet(U))$ est la cohomologie de
$$
\lim_n \mathcal{F}_n^{m - 2}(U) \to
\lim_n \mathcal{F}_n^{m - 1}(U) \to
\lim_n \mathcal{F}_n^m(U) \to
\lim_n \mathcal{F}_n^{m + 1}(U)
$$
au troisième terme en partant de la gauche. D'après les hypothèses (2)(a)
et (2)(b), nous pouvons appliquer Compléments d'algèbre, Lemme
\ref{more-algebra-lemma-apply-Mittag-Leffler-again}, pour conclure que
$$
H^m(\mathcal{F}^\bullet(U)) = \lim H^m(\mathcal{F}_n^\bullet(U))
$$
D'après l'hypothèse (2)(c), nous en déduisons
$$
H^m(\mathcal{F}^\bullet(U)) = H^m(\mathcal{F}_n^\bullet(U))
$$
pour tout $n \geq n_0$. L'hypothèse (1) montre que le faisceau associé à
$U \mapsto H^m(\mathcal{F}^\bullet(U))$ est égal au faisceau associé à
$U \mapsto H^m(\mathcal{F}_n^\bullet(U))$ pour tout $n \geq n_0$. Le
système projectif de faisceaux $H^m(\mathcal{F}_n^\bullet)$ est donc
constant pour $n \geq n_0$, de valeur $H^m(\mathcal{F}^\bullet)$, ce qui
démontre le lemme.
\end{proof}







\section{Dimension cohomologique bornée}
\label{sites-cohomology-section-bounded}

\noindent
Dans cette section, nous cherchons à déterminer quand un foncteur $Rf_*$
est de dimension cohomologique bornée. Cette question est assez délicate
pour les complexes non bornés.

\begin{situation}
\label{sites-cohomology-situation-olsson-laszlo}
Soit $\mathcal{C}$ un site. Soit $\mathcal{O}$ un faisceau d'anneaux sur
$\mathcal{C}$. Soit $\mathcal{A} \subset \textit{Mod}(\mathcal{O})$ une
sous-catégorie de Serre faible. Nous faisons l'hypothèse suivante : il
existe une partie $\mathcal{B} \subset \Ob(\mathcal{C})$ telle que
\begin{enumerate}
\item tout objet de $\mathcal{C}$ possède un recouvrement dont les membres
appartiennent à $\mathcal{B}$, et
\item pour tout $V \in \mathcal{B}$, il existe un entier $d_V$ et un
système cofinal $\text{Cov}_V$ de recouvrements de $V$ tels que
$$
H^p(V_i, \mathcal{F}) = 0 \text{ pour }
\{V_i \to V\} \in \text{Cov}_V,\ p > d_V, \text{ et }
\mathcal{F} \in \Ob(\mathcal{A})
$$
\end{enumerate}
\end{situation}

\begin{lemma}
\label{sites-cohomology-lemma-olsson-laszlo}
\begin{reference}
C'est \cite[Proposition 2.1.4]{six-I}, avec des hypothèses légèrement
modifiées ; il s'agit de l'analogue pour les sites de
\cite[Proposition 3.13]{Spaltenstein}.
\end{reference}
Dans la Situation \ref{sites-cohomology-situation-olsson-laszlo}, pour tout
$E \in D_\mathcal{A}(\mathcal{O})$, le morphisme
$E \to R\lim \tau_{\geq -n} E$
de Catégories dérivées, Remarque
\ref{derived-remark-map-into-derived-limit-truncations}
est un isomorphisme dans $D(\mathcal{O})$.
\end{lemma}

\begin{proof}
Cela résulte immédiatement du Lemme \ref{sites-cohomology-lemma-is-limit-spaltenstein}.
\end{proof}

\begin{lemma}
\label{sites-cohomology-lemma-olsson-laszlo-modified}
Dans la Situation \ref{sites-cohomology-situation-olsson-laszlo}, soit $(K_n)$ un système
projectif dans $D_\mathcal{A}^+(\mathcal{O})$. Supposons que, pour tout
$j$, le système projectif $(H^j(K_n))$ dans $\mathcal{A}$ soit
stationnaire, de valeur $\mathcal{H}^j$. Alors
$H^j(R\lim K_n) = \mathcal{H}^j$ pour tout $j$.
\end{lemma}

\begin{proof}
Soit $V \in \mathcal{B}$. Soit $\{V_i \to V\}$ un élément de l'ensemble
$\text{Cov}_V$ de la Situation \ref{sites-cohomology-situation-olsson-laszlo}. Comme $K_n$
est borné inférieurement, il existe une suite spectrale
$$
E_2^{p, q} = H^p(V_i, H^q(K_n))
$$
convergeant vers $H^{p + q}(V_i, K_n)$. Voir Catégories dérivées, Lemme
\ref{derived-lemma-two-ss-complex-functor}. Observons que, par hypothèse,
$E_2^{p, q} = 0$ pour $p > d_V$. Choisissons $n_0$ tel que
$$
\begin{matrix}
\mathcal{H}^{j + 1} & = & H^{j + 1}(K_n), \\
\mathcal{H}^j & = & H^j(K_n), \\
\ldots, \\
\mathcal{H}^{j - d_V - 2} & = & H^{j - d_V - 2}(K_n)
\end{matrix}
$$
pour tout $n \geq n_0$. En comparant les suites spectrales ci-dessus pour
$K_n$ et $K_{n_0}$, nous voyons que, pour $n \geq n_0$, les groupes de
cohomologie $H^{j - 1}(V_i, K_n)$ et $H^j(V_i, K_n)$ sont indépendants de
$n$. Il s'ensuit que le morphisme sur les sections
$H^j(R\lim K_n)(V) \to H^j(K_n)(V)$ est injectif pour $n$ assez grand
(en fonction de $V$) ; voir le Lemme
\ref{sites-cohomology-lemma-cohomology-derived-limit-injective}. Comme tout objet de
$\mathcal{C}$ peut être recouvert par des éléments de $\mathcal{B}$, nous
en déduisons que le morphisme
$H^j(R\lim K_n) \to \mathcal{H}^j$ est injectif.

\medskip\noindent
La surjectivité se démontre de manière analogue. Choisissons en effet
$U \in \Ob(\mathcal{C})$ et $\gamma \in \mathcal{H}^j(U)$. Nous voulons
relever $\gamma$ en une section de $H^j(R\lim K_n)$ après avoir remplacé
$U$ par les membres d'un recouvrement. D'après la propriété (1) de la
Situation \ref{sites-cohomology-situation-olsson-laszlo}, nous pouvons donc supposer
$U = V \in \mathcal{B}$. Choisissons $n_0$ tel que
$$
\begin{matrix}
\mathcal{H}^{j + 1} & = & H^{j + 1}(K_n), \\
\mathcal{H}^j & = & H^j(K_n), \\
\ldots, \\
\mathcal{H}^{j - d_V - 2} & = & H^{j - d_V - 2}(K_n)
\end{matrix}
$$
pour tout $n \geq n_0$. Choisissons un élément $\{V_i \to V\}$ de
$\text{Cov}_V$ tel que
$\gamma|_{V_i} \in \mathcal{H}^j(V_i) = H^j(K_{n_0})(V_i)$
se relève en un élément $\gamma_{n_0, i} \in H^j(V_i, K_{n_0})$. Cela est
possible car $H^j(K_{n_0})$ est le faisceau associé au préfaisceau
$U \mapsto H^j(U, K_{n_0})$ d'après le Lemme
\ref{sites-cohomology-lemma-sheafification-cohomology}. La discussion du premier paragraphe
de la démonstration montre que $H^{j - 1}(V_i, K_n)$ et
$H^j(V_i, K_n)$ sont indépendants de $n \geq n_0$. Ainsi,
$\gamma_{n_0, i}$ se relève en un élément
$\gamma_i \in H^j(V_i, R\lim K_n)$ d'après le Lemme
\ref{sites-cohomology-lemma-RGamma-commutes-with-Rlim}. Cela achève la démonstration.
\end{proof}

\begin{lemma}
\label{sites-cohomology-lemma-olsson-laszlo-map-version-one}
\begin{reference}
Il s'agit d'une version de \cite[Lemma 2.1.10]{six-I} avec des hypothèses
légèrement modifiées.
\end{reference}
Soit $f : (\Sh(\mathcal{C}), \mathcal{O}) \to (\Sh(\mathcal{C}'), \mathcal{O}')$
un morphisme de topos annelés. Soient
$\mathcal{A} \subset \textit{Mod}(\mathcal{O})$ et
$\mathcal{A}' \subset \textit{Mod}(\mathcal{O}')$ des sous-catégories de
Serre faibles. Supposons qu'il existe un entier $N$ tel que
\begin{enumerate}
\item $\mathcal{C}, \mathcal{O}, \mathcal{A}$ satisfassent l'hypothèse de
la Situation \ref{sites-cohomology-situation-olsson-laszlo},
\item $\mathcal{C}', \mathcal{O}', \mathcal{A}'$ satisfassent l'hypothèse
de la Situation \ref{sites-cohomology-situation-olsson-laszlo},
\item $R^pf_*\mathcal{F} \in \Ob(\mathcal{A}')$ pour
$p \geq 0$ et $\mathcal{F} \in \Ob(\mathcal{A})$,
\item $R^pf_*\mathcal{F} = 0$ pour
$p > N$ et $\mathcal{F} \in \Ob(\mathcal{A})$,
\end{enumerate}
Alors, pour $K$ dans $D_\mathcal{A}(\mathcal{O})$, nous avons
\begin{enumerate}
\item[(a)] $Rf_*K$ appartient à $D_{\mathcal{A}'}(\mathcal{O}')$,
\item[(b)] le morphisme
$H^j(Rf_*K) \to H^j(Rf_*(\tau_{\geq -n}K))$ est un isomorphisme pour
$j \geq N - n$.
\end{enumerate}
\end{lemma}

\begin{proof}
D'après le Lemme \ref{sites-cohomology-lemma-olsson-laszlo}, nous avons
$K = R\lim \tau_{\geq -n}K$. D'après le Lemme
\ref{sites-cohomology-lemma-Rf-commutes-with-Rlim}, nous avons
$Rf_*K = R\lim Rf_*\tau_{\geq -n}K$. Les complexes
$Rf_*\tau_{\geq -n}K$ sont bornés inférieurement. La suite spectrale
$$
E_2^{p, q} = R^pf_*H^q(\tau_{\geq -n}K)
$$
convergeant vers $H^{p + q}(Rf_*\tau_{\geq -n}K)$ (Catégories dérivées,
Lemme \ref{derived-lemma-two-ss-complex-functor}) et l'hypothèse (3)
montrent que $Rf_*\tau_{\geq -n}K$ appartient à
$D^+_{\mathcal{A}'}(\mathcal{O}')$ ; voir Homologie, Lemme
\ref{homology-lemma-biregular-ss-converges}. Observons que, pour
$m \geq n$, le morphisme
$$
Rf_*(\tau_{\geq -m}K) \longrightarrow Rf_*(\tau_{\geq -n}K)
$$
induit un isomorphisme sur les faisceaux de cohomologie en degrés
$j \geq -n + N$, d'après les suites spectrales ci-dessus. Nous pouvons
donc conclure en appliquant le Lemme
\ref{sites-cohomology-lemma-olsson-laszlo-modified}.
\end{proof}

\noindent
Il s'avère que nous avons parfois besoin d'une variante du lemme précédent
où les hypothèses sont légèrement différentes.

\begin{situation}
\label{sites-cohomology-situation-olsson-laszlo-prime}
Soit
$f : (\mathcal{C}, \mathcal{O}) \to (\mathcal{C}', \mathcal{O}')$
un morphisme de sites annelés. Soit $u : \mathcal{C}' \to \mathcal{C}$ le
foncteur continu de sites correspondant. Soit
$\mathcal{A} \subset \textit{Mod}(\mathcal{O})$ une sous-catégorie de Serre
faible. Nous faisons l'hypothèse suivante : il existe une partie
$\mathcal{B}' \subset \Ob(\mathcal{C}')$ telle que
\begin{enumerate}
\item tout objet de $\mathcal{C}'$ possède un recouvrement dont les membres
appartiennent à $\mathcal{B}'$, et
\item pour tout $V' \in \mathcal{B}'$, il existe un entier $d_{V'}$ et un
système cofinal $\text{Cov}_{V'}$ de recouvrements de $V'$ tels que
$$
H^p(u(V'_i), \mathcal{F}) = 0 \text{ pour }
\{V'_i \to V'\} \in \text{Cov}_{V'},\ p > d_{V'}, \text{ et }
\mathcal{F} \in \Ob(\mathcal{A})
$$
\end{enumerate}
\end{situation}

\begin{lemma}
\label{sites-cohomology-lemma-olsson-laszlo-map-version-two}
\begin{reference}
Il s'agit d'une version de \cite[Lemma 2.1.10]{six-I} avec des hypothèses
légèrement modifiées.
\end{reference}
Soit $f : (\mathcal{C}, \mathcal{O}) \to (\mathcal{C}', \mathcal{O}')$ un
morphisme de sites annelés. Supposons de plus qu'il existe un entier $N$
tel que
\begin{enumerate}
\item $\mathcal{C}, \mathcal{O}, \mathcal{A}$ satisfassent l'hypothèse de
la Situation \ref{sites-cohomology-situation-olsson-laszlo},
\item $f : (\mathcal{C}, \mathcal{O}) \to (\mathcal{C}', \mathcal{O}')$
et $\mathcal{A}$ satisfassent l'hypothèse de la Situation
\ref{sites-cohomology-situation-olsson-laszlo-prime},
\item $R^pf_*\mathcal{F} = 0$ pour
$p > N$ et $\mathcal{F} \in \Ob(\mathcal{A})$,
\end{enumerate}
Alors, pour $K$ dans $D_\mathcal{A}(\mathcal{O})$, le morphisme
$H^j(Rf_*K) \to H^j(Rf_*(\tau_{\geq -n}K))$ est un isomorphisme pour
$j \geq N - n$.
\end{lemma}

\begin{proof}
Soit $K$ dans $D_\mathcal{A}(\mathcal{O})$. D'après le Lemme
\ref{sites-cohomology-lemma-olsson-laszlo}, nous avons
$K = R\lim \tau_{\geq -n}K$. D'après le Lemme
\ref{sites-cohomology-lemma-Rf-commutes-with-Rlim}, nous avons
$Rf_*K = R\lim Rf_*(\tau_{\geq -n}K)$. Soit $V' \in \mathcal{B}'$ et soit
$\{V'_i \to V'\}$ un élément de $\text{Cov}_{V'}$. Considérons
$$
H^j(V'_i, Rf_*K) = H^j(u(V'_i), K)
\quad\text{et}\quad
H^j(V'_i, Rf_*(\tau_{\geq -n}K)) = H^j(u(V'_i), \tau_{\geq -n}K)
$$
L'hypothèse de la Situation \ref{sites-cohomology-situation-olsson-laszlo-prime} implique
que le dernier groupe est indépendant de $n$ pour $n$ assez grand, en
fonction de $j$ et de $d_{V'}$. Nous omettons certains détails. En
appliquant ceci à $j$ et $j - 1$, puis en utilisant le Lemme
\ref{sites-cohomology-lemma-RGamma-commutes-with-Rlim}, nous obtenons
$$
H^j(V'_i, Rf_*K) = \lim H^j(V'_i, Rf_*(\tau_{\geq -n} K))
$$
et le système de droite est constant pour $n$ supérieur à une constante ne
dépendant que de $d_{V'}$ et de $j$. Le Lemme
\ref{sites-cohomology-lemma-cohomology-derived-limit-injective} implique donc que
$$
H^j(Rf_*K)(V') \longrightarrow
\left(\lim H^j(Rf_*(\tau_{\geq -n}K))\right)(V')
$$
est injectif. Comme les éléments $V' \in \mathcal{B}'$ recouvrent tout
objet de $\mathcal{C}'$, nous en déduisons que le morphisme
$H^j(Rf_*K) \to \lim H^j(Rf_*(\tau_{\geq -n}K))$ est injectif. La suite
spectrale
$$
E_2^{p, q} = R^pf_*H^q(\tau_{\geq -n}K)
$$
convergeant vers $H^{p + q}(Rf_*(\tau_{\geq -n}K))$ (Catégories dérivées,
Lemme \ref{derived-lemma-two-ss-complex-functor}) et l'hypothèse (3)
montrent que $H^j(Rf_*(\tau_{\geq -n}K))$ est constant pour
$n \geq N - j$. Ainsi,
$H^j(Rf_*K) \to H^j(Rf_*(\tau_{\geq -n}K))$ est injectif pour
$j \geq N - n$.

\medskip\noindent
Nous avons ainsi démontré le lemme en remplaçant « isomorphisme » par
« injectif » dans sa dernière ligne. Choisissons maintenant $j$ et $n$
tels que $j \geq N - n$, puis considérons le triangle distingué
$$
\tau_{\leq -n - 1}K \to K \to \tau_{\geq -n}K \to (\tau_{\leq -n - 1}K)[1]
$$
Voir Catégories dérivées, Remarque
\ref{derived-remark-truncation-distinguished-triangle}.
Comme $\tau_{\geq -n}\tau_{\leq -n -1}K = 0$, l'injectivité déjà démontrée
pour $\tau_{-n - 1}K$ implique
$$
0 = H^j(Rf_*(\tau_{\leq -n - 1}K)) =
H^{j + 1}(Rf_*(\tau_{\leq -n - 1}K)) =
H^{j + 2}(Rf_*(\tau_{\leq -n - 1}K)) = \ldots
$$
La suite exacte longue de cohomologie associée au triangle distingué
$$
Rf_*(\tau_{\leq -n - 1}K) \to Rf_*K \to Rf_*(\tau_{\geq -n}K) \to
Rf_*(\tau_{\leq -n - 1}K)[1]
$$
implique alors que $H^j(Rf_*K) \to H^j(Rf_*(\tau_{\geq -n}K))$ est un
isomorphisme.
\end{proof}








\section{Mayer--Vietoris}
\label{sites-cohomology-section-mayer-vietoris}

\noindent
Pour l'énoncé et la démonstration usuels de Mayer--Vietoris, voir
Cohomologie, section \ref{cohomology-section-mayer-vietoris}.

\medskip\noindent
Soit $(\mathcal{C}, \mathcal{O})$ un site annelé. Considérons un diagramme
commutatif
$$
\xymatrix{
E \ar[d] \ar[r] & Y \ar[d] \\
Z \ar[r] & X
}
$$
dans la catégorie $\mathcal{C}$. Dans cette situation, étant donné un objet
$K$ de $D(\mathcal{O})$, nous obtenons ce qui ressemble au début d'un
triangle distingué
$$
R\Gamma(X, K) \to
R\Gamma(Z, K) \oplus
R\Gamma(Y, K) \to
R\Gamma(E, K)
$$
Le lemme suivant précise cette observation.

\begin{lemma}
\label{sites-cohomology-lemma-c-square}
Dans la situation ci-dessus, choisissons un complexe K-injectif
$\mathcal{I}^\bullet$ de $\mathcal{O}$-modules représentant $K$. En
multipliant par $-1$ le morphisme canonique pour l'une des quatre flèches,
nous obtenons des morphismes de complexes
$$
\mathcal{I}^\bullet(X) \xrightarrow{\alpha}
\mathcal{I}^\bullet(Z) \oplus
\mathcal{I}^\bullet(Y) \xrightarrow{\beta}
\mathcal{I}^\bullet(E)
$$
avec $\beta \circ \alpha = 0$. Nous obtenons donc un morphisme canonique
$$
c^K_{X, Z, Y, E} :
\mathcal{I}^\bullet(X)
\longrightarrow
C(\beta)^\bullet[-1]
$$
Ce morphisme est canonique au sens où un autre choix de complexe
K-injectif représentant $K$ détermine une flèche isomorphe dans la
catégorie dérivée des groupes abéliens. Si $c^K_{X, Z, Y, E}$ est un
isomorphisme, son inverse fournit un triangle distingué canonique
$$
R\Gamma(X, K) \to
R\Gamma(Z, K) \oplus
R\Gamma(Y, K) \to
R\Gamma(E, K) \to
R\Gamma(X, K)[1]
$$
Toutes ces constructions sont fonctorielles en $K$.
\end{lemma}

\begin{proof}
Le lemme se démontre de lui-même. Par exemple, si
$\mathcal{J}^\bullet$ est un second complexe K-injectif représentant $K$,
nous pouvons choisir un quasi-isomorphisme
$\mathcal{I}^\bullet \to \mathcal{J}^\bullet$ qui détermine des
quasi-isomorphismes entre tous les complexes en jeu. Nous omettons les
détails. Pour la construction des cônes et leur relation avec les triangles
distingués, voir Catégories dérivées, sections
\ref{derived-section-cones} et
\ref{derived-section-homotopy-triangulated}.
\end{proof}

\begin{lemma}
\label{sites-cohomology-lemma-two-out-of-three-blow-up-square}
Dans la situation ci-dessus, soit $K_1 \to K_2 \to K_3 \to K_1[1]$ un
triangle distingué dans $D(\mathcal{O})$. Si $c^{K_i}_{X, Z, Y, E}$ est un
quasi-isomorphisme pour deux valeurs de $i$ parmi $\{1, 2, 3\}$, il en est
de même pour la troisième valeur de $i$.
\end{lemma}

\begin{proof}
En faisant tourner le triangle, nous pouvons supposer que
$c^{K_1}_{X, Z, Y, E}$ et $c^{K_2}_{X, Z, Y, E}$ sont des
quasi-isomorphismes. Choisissons un morphisme
$f : \mathcal{I}^\bullet_1 \to \mathcal{I}^\bullet_2$ de complexes
K-injectifs de $\mathcal{O}$-modules représentant $K_1 \to K_2$. Alors
$K_3$ est représenté par le complexe K-injectif $C(f)^\bullet$ ; voir
Catégories dérivées, Lemme \ref{derived-lemma-triangle-K-injective}. Le
morphisme $c^{K_3}_{X, Z, Y, E}$ est donc un isomorphisme, car il est le
troisième côté d'un morphisme de triangles distingués dans
$K(\textit{Ab})$ dont les deux autres côtés sont des quasi-isomorphismes.
Nous omettons certains détails ; utiliser Catégories dérivées, Lemme
\ref{derived-lemma-third-isomorphism-triangle}.
\end{proof}

\noindent
Donnons un critère assurant que cette construction produit effectivement
un triangle distingué.

\begin{lemma}
\label{sites-cohomology-lemma-square-triangle}
Dans la situation ci-dessus, supposons que
\begin{enumerate}
\item $h_X^\# = h_Y^\# \amalg_{h_E^\#} h_Z^\#$, et
\item $h_E^\# \to h_Y^\#$ soit injectif.
\end{enumerate}
Alors la construction du Lemme \ref{sites-cohomology-lemma-c-square} produit un triangle
distingué
$$
R\Gamma(X, K) \to
R\Gamma(Z, K) \oplus
R\Gamma(Y, K) \to
R\Gamma(E, K) \to R\Gamma(X, K)[1]
$$
fonctoriel en $K$ dans $D(\mathcal{C})$.
\end{lemma}

\begin{proof}
Nous pouvons représenter $K$ par un complexe K-injectif dont les termes
sont des faisceaux abéliens injectifs ; voir la section
\ref{sites-cohomology-section-unbounded}. Il suffit donc de montrer que, si $\mathcal{I}$
est un faisceau abélien injectif, alors
$$
0 \to \mathcal{I}(X) \to
\mathcal{I}(Z) \oplus \mathcal{I}(Y) \to
\mathcal{I}(E) \to 0
$$
est une suite exacte courte. La première flèche est injective car, d'après
la condition (1), le morphisme $h_Y \amalg h_Z \to h_X$ devient surjectif
après faisceautisation, ce qui signifie que
$\{Y \to X, Z \to X\}$ peut être raffiné par un recouvrement de $X$. La
dernière flèche est surjective car
$\mathcal{I}(Y) \to \mathcal{I}(E)$ est surjectif. En effet, nous avons
$\mathcal{I}(E) = \Hom(\mathbf{Z}_E^\#, \mathcal{I})$,
$\mathcal{I}(Y) = \Hom(\mathbf{Z}_Y^\#, \mathcal{I})$,
le morphisme $\mathbf{Z}_E^\# \to \mathbf{Z}_Y^\#$ est injectif d'après
(2), et $\mathcal{I}$ est un faisceau abélien injectif. Comparer à Modules
sur les sites, section
\ref{sites-modules-section-free-abelian-sheaf}.
Enfin, supposons donnés $s \in \mathcal{I}(Y)$ et
$t \in \mathcal{F}(Z)$ ayant la même image dans $\mathcal{I}(E)$. Alors
$s$ et $t$ définissent un morphisme
$$
s \amalg t : h_Y^\# \amalg h_Z^\# \longrightarrow \mathcal{I}
$$
qui, par hypothèse, se factorise par
$h_Y^\# \amalg_{h_E^\#} h_Z^\#$. L'hypothèse (1) fournit donc un unique
morphisme $h_X^\# \to \mathcal{I}$, correspondant à un élément de
$\mathcal{I}(X)$ qui se restreint à $s$ sur $Y$ et à $t$ sur $Z$.
\end{proof}

\begin{lemma}
\label{sites-cohomology-lemma-square-triangle-general}
Soit $\mathcal{C}$ un site. Considérons un diagramme commutatif
$$
\xymatrix{
\mathcal{D} \ar[r] \ar[d] & \mathcal{F} \ar[d] \\
\mathcal{E} \ar[r] & \mathcal{G}
}
$$
de préfaisceaux d'ensembles sur $\mathcal{C}$ et supposons que
\begin{enumerate}
\item $\mathcal{G}^\# =
\mathcal{E}^\# \amalg_{\mathcal{D}^\#} \mathcal{F}^\#$, et
\item $\mathcal{D}^\# \to \mathcal{F}^\#$ soit injectif.
\end{enumerate}
Alors il existe un triangle distingué canonique
$$
R\Gamma(\mathcal{G}, K) \to
R\Gamma(\mathcal{E}, K) \oplus
R\Gamma(\mathcal{F}, K) \to
R\Gamma(\mathcal{D}, K) \to
R\Gamma(\mathcal{G}, K)[1]
$$
fonctoriel en $K \in D(\mathcal{C})$, où $R\Gamma(\mathcal{G}, -)$ désigne
la cohomologie étudiée dans la section \ref{sites-cohomology-section-limp}.
\end{lemma}

\begin{proof}
Comme la faisceautisation est exacte et que
$R\Gamma(\mathcal{G}, -) = R\Gamma(\mathcal{G}^\#, -)$, nous pouvons
supposer que $\mathcal{D}, \mathcal{E}, \mathcal{F}, \mathcal{G}$ sont des
faisceaux d'ensembles. De plus, la cohomologie
$R\Gamma(\mathcal{G}, -)$ ne dépend que du topos, et non du site qui le
présente. D'après Sites, Lemme \ref{sites-lemma-topos-good-site}, nous
pouvons donc remplacer $\mathcal{C}$ par un site « plus grand », muni d'une
topologie sous-canonique, tel que $\mathcal{G} = h_X$,
$\mathcal{F} = h_Y$, $\mathcal{E} = h_Z$ et $\mathcal{D} = h_E$ pour
certains objets $X, Y, Z, E$ de $\mathcal{C}$. Dans ce cas, le résultat
découle du Lemme \ref{sites-cohomology-lemma-square-triangle}.
\end{proof}







\section{Comparaison de deux topologies}
\label{sites-cohomology-section-compare}

\noindent
Soit $\mathcal{C}$ une catégorie. Soient
$\text{Cov}(\mathcal{C}) \supset \text{Cov}'(\mathcal{C})$
deux manières de munir $\mathcal{C}$ d'une structure de site. Notons
$\tau$ la topologie correspondant à $\text{Cov}(\mathcal{C})$ et $\tau'$
celle correspondant à $\text{Cov}'(\mathcal{C})$. Le foncteur identité de
$\mathcal{C}$ définit alors un morphisme de sites
$$
\epsilon : \mathcal{C}_\tau \longrightarrow \mathcal{C}_{\tau'}
$$
où $\epsilon_*$ est le foncteur identité sur les préfaisceaux sous-jacents,
et où $\epsilon^{-1}$ est la faisceautisation pour $\tau$ d'un
$\tau'$-faisceau. Voir Sites, Exemples
\ref{sites-example-finer-topology} et
\ref{sites-example-finer-topology-bis}. Dans la situation ci-dessus, nous
avons les propriétés suivantes :
\begin{enumerate}
\item $\epsilon_* : \Sh(\mathcal{C}_\tau) \to \Sh(\mathcal{C}_{\tau'})$
est pleinement fidèle et $\epsilon^{-1} \circ \epsilon_* = \text{id}$,
\item $\epsilon_* : \textit{Ab}(\mathcal{C}_\tau) \to
\textit{Ab}(\mathcal{C}_{\tau'})$
est pleinement fidèle et $\epsilon^{-1} \circ \epsilon_* = \text{id}$,
\item $R\epsilon_* : D(\mathcal{C}_\tau) \to D(\mathcal{C}_{\tau'})$
est pleinement fidèle et $\epsilon^{-1} \circ R\epsilon_* = \text{id}$,
\item si $\mathcal{O}$ est un faisceau d'anneaux pour la topologie $\tau$,
alors $\mathcal{O}$ est également un faisceau pour la topologie $\tau'$ et
$\epsilon$ devient un morphisme plat de sites annelés
$$
\epsilon :
(\mathcal{C}_\tau, \mathcal{O}_\tau)
\longrightarrow
(\mathcal{C}_{\tau'}, \mathcal{O}_{\tau'})
$$
\item $\epsilon_* : \textit{Mod}(\mathcal{O}_\tau) \to
\textit{Mod}(\mathcal{O}_{\tau'})$ est pleinement fidèle et
$\epsilon^* \circ \epsilon_* = \text{id}$,
\item $R\epsilon_* : D(\mathcal{O}_\tau) \to D(\mathcal{O}_{\tau'})$
est pleinement fidèle et $\epsilon^* \circ R\epsilon_* = \text{id}$.
\end{enumerate}
Voici quelques explications.

\medskip\noindent
Pour (1). Soit $\mathcal{F}$ un faisceau d'ensembles pour la topologie
$\tau$. Alors $\epsilon_*\mathcal{F}$ n'est autre que $\mathcal{F}$
considéré comme faisceau pour la topologie $\tau'$. Appliquer
$\epsilon^{-1}$ revient à prendre le $\tau$-faisceau associé à
$\mathcal{F}$, ce qui ne change rien puisque $\mathcal{F}$ est déjà un
$\tau$-faisceau. Ainsi,
$\epsilon^{-1}(\epsilon_*\mathcal{F})) = \mathcal{F}$. La pleine fidélité
résulte de Catégories, Lemme
\ref{categories-lemma-adjoint-fully-faithful}.

\medskip\noindent
Pour (2). C'est une conséquence de (1), car prendre l'image inverse ou
l'image directe de faisceaux abéliens revient à effectuer ces opérations
sur les faisceaux d'ensembles sous-jacents.

\medskip\noindent
Pour (3). Soit $K$ un objet de $D(\mathcal{C}_\tau)$. Pour calculer
$R\epsilon_*K$, choisissons un complexe K-injectif
$\mathcal{I}^\bullet$ représentant $K$ et posons
$R\epsilon_*K = \epsilon_*\mathcal{I}^\bullet$. Comme
$\epsilon^{-1} : D(\mathcal{C}_{\tau'}) \to D(\mathcal{C}_\tau)$ se
calcule sur un objet $L$ en appliquant le foncteur exact $\epsilon^{-1}$ à
n'importe quel complexe de faisceaux abéliens représentant $L$, nous
voyons que $\epsilon^{-1}R\epsilon_*K$ est représenté par
$\epsilon^{-1}\epsilon_*\mathcal{I}^\bullet$. D'après (1), nous avons
$\mathcal{I}^\bullet = \epsilon^{-1}\epsilon_*\mathcal{I}^\bullet$.
Autrement dit, $\epsilon^{-1} \circ R\epsilon_* = \text{id}$, et nous
concluons comme précédemment.

\medskip\noindent
Pour (4). Observons que
$\epsilon^{-1}\mathcal{O}_{\tau'} = \mathcal{O}_\tau$ ; voir la discussion
de (1). Ainsi, $\epsilon$ est un morphisme plat de sites annelés ; voir
Modules sur les sites, Définition
\ref{sites-modules-definition-flat-morphism}. De plus, il est clair que
$\epsilon^* = \epsilon^{-1}$ sur les $\mathcal{O}_{\tau'}$-modules
(l'image inverse comme module a le même faisceau abélien sous-jacent que
l'image inverse du faisceau abélien sous-jacent).

\medskip\noindent
Pour (5). Cela résulte clairement de (2) et de ce que nous avons dit en (4).

\medskip\noindent
Pour (6). L'argument est analogue à celui de (3). Nous omettons les détails.












\section{Formalismes de la descente cohomologique}
\label{sites-cohomology-section-formal-cohomological-descent}

\noindent
Dans cette section, nous examinons seulement dans quelle mesure un
morphisme de topos annelés détermine un plongement de la catégorie dérivée
de la base dans celle de la source. Voici un résultat typique.

\begin{lemma}
\label{sites-cohomology-lemma-downstairs}
Soit $f : (\Sh(\mathcal{C}), \mathcal{O}_\mathcal{C}) \to
(\Sh(\mathcal{D}), \mathcal{O}_\mathcal{D})$ un morphisme de topos annelés.
Considérons la sous-catégorie pleine
$D' \subset D(\mathcal{O}_\mathcal{D})$ constituée des objets $K$ tels que
$$
K \longrightarrow Rf_*Lf^*K
$$
soit un isomorphisme. Alors $D'$ est une sous-catégorie strictement pleine,
triangulée et saturée de $D(\mathcal{O}_\mathcal{D})$, et le foncteur
$Lf^* : D' \to D(\mathcal{O}_\mathcal{C})$ est pleinement fidèle.
\end{lemma}

\begin{proof}
Voir Catégories dérivées, Définition \ref{derived-definition-saturated},
pour la définition de « saturée » dans ce contexte, et Catégories dérivées,
Lemme \ref{derived-lemma-triangulated-subcategory}, pour une discussion des
sous-catégories triangulées. Le morphisme canonique du lemme est l'unité
du couple de foncteurs adjoints $(Lf^*, Rf_*)$ ; voir le Lemme
\ref{sites-cohomology-lemma-adjoint}. Cela étant dit, nous omettons la démonstration du fait
que $D'$ est une sous-catégorie triangulée saturée : elle résulte
formellement de l'exactitude des foncteurs $Lf^*$ et $Rf_*$ entre
catégories triangulées. La dernière assertion résulte formellement de
l'adjonction de $Lf^*$ et $Rf_*$ ; comparer à Catégories, Lemme
\ref{categories-lemma-adjoint-fully-faithful}.
\end{proof}

\begin{lemma}
\label{sites-cohomology-lemma-upstairs}
Soit $f : (\Sh(\mathcal{C}), \mathcal{O}_\mathcal{C}) \to
(\Sh(\mathcal{D}), \mathcal{O}_\mathcal{D})$ un morphisme de topos annelés.
Considérons la sous-catégorie pleine
$D' \subset D(\mathcal{O}_\mathcal{C})$ constituée des objets $K$ tels que
$$
Lf^*Rf_*K \longrightarrow K
$$
soit un isomorphisme. Alors $D'$ est une sous-catégorie strictement pleine,
triangulée et saturée de $D(\mathcal{O}_\mathcal{C})$, et le foncteur
$Rf_* : D' \to D(\mathcal{O}_\mathcal{D})$ est pleinement fidèle.
\end{lemma}

\begin{proof}
Voir Catégories dérivées, Définition \ref{derived-definition-saturated},
pour la définition de « saturée » dans ce contexte, et Catégories dérivées,
Lemme \ref{derived-lemma-triangulated-subcategory}, pour une discussion des
sous-catégories triangulées. Le morphisme canonique du lemme est la
co-unité du couple de foncteurs adjoints $(Lf^*, Rf_*)$ ; voir le Lemme
\ref{sites-cohomology-lemma-adjoint}. Cela étant dit, nous omettons la démonstration du fait
que $D'$ est une sous-catégorie triangulée saturée : elle résulte
formellement de l'exactitude des foncteurs $Lf^*$ et $Rf_*$ entre
catégories triangulées. La dernière assertion résulte formellement de
l'adjonction de $Lf^*$ et $Rf_*$ ; comparer à Catégories, Lemme
\ref{categories-lemma-adjoint-fully-faithful}.
\end{proof}

\begin{lemma}
\label{sites-cohomology-lemma-bounded-in-image-upstairs}
Soit $f : (\Sh(\mathcal{C}), \mathcal{O}_\mathcal{C}) \to
(\Sh(\mathcal{D}), \mathcal{O}_\mathcal{D})$ un morphisme de topos annelés.
Soit $K$ un objet de $D(\mathcal{O}_\mathcal{C})$. Supposons que
\begin{enumerate}
\item $f$ soit plat,
\item $K$ soit borné inférieurement,
\item $f^*Rf_*H^q(K) \to H^q(K)$ soit un isomorphisme.
\end{enumerate}
Alors $f^*Rf_*K \to K$ est un isomorphisme.
\end{lemma}

\begin{proof}
Observons que $f^*Rf_*K \to K$ est un isomorphisme si et seulement s'il
induit un isomorphisme sur les faisceaux de cohomologie $H^j$. Observons
que
$H^j(f^*Rf_*K) = f^*H^j(Rf_*K) = f^*H^j(Rf_*\tau_{\leq j}K) =
H^j(f^*Rf_*\tau_{\leq j}K)$.
Nous pouvons donc supposer que $K$ est borné. La propriété (3) montre alors
que les faisceaux de cohomologie appartiennent à la sous-catégorie
triangulée $D' \subset D(\mathcal{O}_\mathcal{C})$ du Lemme
\ref{sites-cohomology-lemma-upstairs}. Par conséquent, $K$ y appartient également.
\end{proof}

\begin{lemma}
\label{sites-cohomology-lemma-bounded-in-image-downstairs}
Soit $f : (\Sh(\mathcal{C}), \mathcal{O}_\mathcal{C}) \to
(\Sh(\mathcal{D}), \mathcal{O}_\mathcal{D})$ un morphisme de topos annelés.
Soit $K$ un objet de $D(\mathcal{O}_\mathcal{D})$. Supposons que
\begin{enumerate}
\item $f$ soit plat,
\item $K$ soit borné inférieurement,
\item $H^q(K) \to Rf_*f^*H^q(K)$ soit un isomorphisme.
\end{enumerate}
Alors $K \to Rf_*f^*K$ est un isomorphisme.
\end{lemma}

\begin{proof}
Observons que $K \to Rf_*f^*K$ est un isomorphisme si et seulement s'il
induit un isomorphisme sur les faisceaux de cohomologie $H^j$. Observons
que
$H^j(Rf_*f^*K) = H^j(Rf_*\tau_{\leq j}f^*K) = H^j(Rf_*f^*\tau_{\leq j}K)$.
Nous pouvons donc supposer que $K$ est borné. La propriété (3) montre alors
que les faisceaux de cohomologie appartiennent à la sous-catégorie
triangulée $D' \subset D(\mathcal{O}_\mathcal{D})$ du Lemme
\ref{sites-cohomology-lemma-downstairs}. Par conséquent, $K$ y appartient également.
\end{proof}

\begin{lemma}
\label{sites-cohomology-lemma-equivalence-bounded}
Soit $f : (\Sh(\mathcal{C}), \mathcal{O}) \to (\Sh(\mathcal{C}'), \mathcal{O}')$
un morphisme de topos annelés. Soient
$\mathcal{A} \subset \textit{Mod}(\mathcal{O})$ et
$\mathcal{A}' \subset \textit{Mod}(\mathcal{O}')$ des sous-catégories de
Serre faibles. Supposons que
\begin{enumerate}
\item $f$ soit plat,
\item $f^*$ induise une équivalence de catégories
$\mathcal{A}' \to \mathcal{A}$,
\item $\mathcal{F}' \to Rf_*f^*\mathcal{F}'$ soit un isomorphisme pour
$\mathcal{F}' \in \Ob(\mathcal{A}')$.
\end{enumerate}
Alors
$f^* : D_{\mathcal{A}'}^+(\mathcal{O}') \to D_\mathcal{A}^+(\mathcal{O})$
est une équivalence de catégories dont un quasi-inverse est
$Rf_* : D_\mathcal{A}^+(\mathcal{O}) \to D_{\mathcal{A}'}^+(\mathcal{O}')$.
\end{lemma}

\begin{proof}
D'après les hypothèses (2) et (3) et les Lemmes
\ref{sites-cohomology-lemma-bounded-in-image-downstairs} et \ref{sites-cohomology-lemma-downstairs}, nous
voyons que
$f^* : D_{\mathcal{A}'}^+(\mathcal{O}') \to D_\mathcal{A}^+(\mathcal{O})$
est pleinement fidèle. Soit $\mathcal{F} \in \Ob(\mathcal{A})$. Nous
pouvons écrire $\mathcal{F} = f^*\mathcal{F}'$. Alors
$Rf_*\mathcal{F} = Rf_* f^*\mathcal{F}' = \mathcal{F}'$.
En particulier, $R^pf_*\mathcal{F} = 0$ pour $p > 0$ et
$f_*\mathcal{F} \in \Ob(\mathcal{A}')$. Ainsi, pour tout
$K \in D^+_\mathcal{A}(\mathcal{O})$, la suite spectrale
$E_2^{p, q} = R^pf_*H^q(K)$ convergeant vers $R^{p + q}f_*K$ montre que
$Rf_*K$ appartient à $D^+_{\mathcal{A}'}(\mathcal{O}')$. Les Lemmes
\ref{sites-cohomology-lemma-bounded-in-image-upstairs} et \ref{sites-cohomology-lemma-upstairs} montrent
également que
$Rf_* : D_\mathcal{A}^+(\mathcal{O}) \to D_{\mathcal{A}'}^+(\mathcal{O}')$
est pleinement fidèle. Comme $f^*$ et $Rf_*$ sont adjoints, nous obtenons
alors le résultat du lemme, par exemple à l'aide de Catégories, Lemme
\ref{categories-lemma-adjoint-fully-faithful}.
\end{proof}

\begin{lemma}
\label{sites-cohomology-lemma-equivalence-unbounded-one}
\begin{reference}
C'est l'analogue de \cite[Theorem 2.2.3]{six-I}.
\end{reference}
Soit $f : (\Sh(\mathcal{C}), \mathcal{O}) \to (\Sh(\mathcal{C}'), \mathcal{O}')$
un morphisme de topos annelés. Soient
$\mathcal{A} \subset \textit{Mod}(\mathcal{O})$ et
$\mathcal{A}' \subset \textit{Mod}(\mathcal{O}')$ des sous-catégories de
Serre faibles. Supposons que
\begin{enumerate}
\item $f$ soit plat,
\item $f^*$ induise une équivalence de catégories
$\mathcal{A}' \to \mathcal{A}$,
\item $\mathcal{F}' \to Rf_*f^*\mathcal{F}'$ soit un isomorphisme pour
$\mathcal{F}' \in \Ob(\mathcal{A}')$,
\item $\mathcal{C}, \mathcal{O}, \mathcal{A}$ satisfassent l'hypothèse de
la Situation \ref{sites-cohomology-situation-olsson-laszlo},
\item $\mathcal{C}', \mathcal{O}', \mathcal{A}'$ satisfassent l'hypothèse
de la Situation \ref{sites-cohomology-situation-olsson-laszlo}.
\end{enumerate}
Alors $f^* : D_{\mathcal{A}'}(\mathcal{O}') \to D_\mathcal{A}(\mathcal{O})$
est une équivalence de catégories dont un quasi-inverse est
$Rf_* : D_\mathcal{A}(\mathcal{O}) \to D_{\mathcal{A}'}(\mathcal{O}')$.
\end{lemma}

\begin{proof}
Comme $f^*$ est exact, il est clair que $f^*$ définit un foncteur
$f^* : D_{\mathcal{A}'}(\mathcal{O}') \to D_\mathcal{A}(\mathcal{O})$
comme dans l'énoncé du lemme et que ce foncteur commute aux foncteurs de
troncation $\tau_{\geq -n}$. Nous savons déjà que $f^*$ et $Rf_*$ sont des
équivalences quasi-inverses sur les catégories bornées inférieurement
correspondantes ; voir le Lemme \ref{sites-cohomology-lemma-equivalence-bounded}. Le Lemme
\ref{sites-cohomology-lemma-olsson-laszlo-map-version-one}, avec $N = 0$, montre que $Rf_*$
définit bien un foncteur
$Rf_* : D_\mathcal{A}(\mathcal{O}) \to D_{\mathcal{A}'}(\mathcal{O}')$
et que ce foncteur commute en outre aux foncteurs de troncation
$\tau_{\geq -n}$. Ainsi, pour $K$ dans $D_\mathcal{A}(\mathcal{O})$, le
morphisme $f^*Rf_*K \to K$ est un isomorphisme, puisque c'est le cas sur
les troncations. De même, pour $K'$ dans
$D_{\mathcal{A}'}(\mathcal{O}')$, le morphisme $K' \to Rf_*f^*K'$ est un
isomorphisme, puisque c'est le cas sur les troncations. Cela achève la
démonstration.
\end{proof}

\begin{lemma}
\label{sites-cohomology-lemma-equivalence-unbounded-two}
\begin{reference}
C'est l'analogue de \cite[Theorem 2.2.3]{six-I}.
\end{reference}
Soit $f : (\mathcal{C}, \mathcal{O}) \to (\mathcal{C}', \mathcal{O}')$ un
morphisme de sites annelés. Soient
$\mathcal{A} \subset \textit{Mod}(\mathcal{O})$ et
$\mathcal{A}' \subset \textit{Mod}(\mathcal{O}')$ des sous-catégories de
Serre faibles. Supposons que
\begin{enumerate}
\item $f$ soit plat,
\item $f^*$ induise une équivalence de catégories
$\mathcal{A}' \to \mathcal{A}$,
\item $\mathcal{F}' \to Rf_*f^*\mathcal{F}'$ soit un isomorphisme pour
$\mathcal{F}' \in \Ob(\mathcal{A}')$,
\item $\mathcal{C}, \mathcal{O}, \mathcal{A}$ satisfassent l'hypothèse de
la Situation \ref{sites-cohomology-situation-olsson-laszlo},
\item $f : (\mathcal{C}, \mathcal{O}) \to (\mathcal{C}', \mathcal{O}')$
et $\mathcal{A}$ satisfassent l'hypothèse de la Situation
\ref{sites-cohomology-situation-olsson-laszlo-prime}.
\end{enumerate}
Alors $f^* : D_{\mathcal{A}'}(\mathcal{O}') \to D_\mathcal{A}(\mathcal{O})$
est une équivalence de catégories dont un quasi-inverse est
$Rf_* : D_\mathcal{A}(\mathcal{O}) \to D_{\mathcal{A}'}(\mathcal{O}')$.
\end{lemma}

\begin{proof}
La démonstration de ce lemme est exactement la même que celle du Lemme
\ref{sites-cohomology-lemma-equivalence-unbounded-one}, à ceci près que la référence au
Lemme \ref{sites-cohomology-lemma-olsson-laszlo-map-version-one} est remplacée par une
référence au Lemme \ref{sites-cohomology-lemma-olsson-laszlo-map-version-two}.
\end{proof}








\section{Comparaison de deux topologies, II}
\label{sites-cohomology-section-compare-II}

\noindent
Soit $\mathcal{C}$ une catégorie. Soient
$\text{Cov}(\mathcal{C}) \supset \text{Cov}'(\mathcal{C})$
deux manières de munir $\mathcal{C}$ d'une structure de site. Notons
$\tau$ la topologie correspondant à $\text{Cov}(\mathcal{C})$ et $\tau'$
celle correspondant à $\text{Cov}'(\mathcal{C})$. Le foncteur identité de
$\mathcal{C}$ définit alors un morphisme de sites
$$
\epsilon : \mathcal{C}_\tau \longrightarrow \mathcal{C}_{\tau'}
$$
où $\epsilon_*$ est le foncteur identité sur les préfaisceaux sous-jacents,
et où $\epsilon^{-1}$ est la faisceautisation pour $\tau$ d'un
$\tau'$-faisceau (elle est donc manifestement exacte). Soit
$\mathcal{O}$ un faisceau d'anneaux pour la topologie $\tau$. Alors
$\mathcal{O}$ est également un faisceau pour la topologie $\tau'$ et
$\epsilon$ devient un morphisme de sites annelés
$$
\epsilon :
(\mathcal{C}_\tau, \mathcal{O}_\tau)
\longrightarrow
(\mathcal{C}_{\tau'}, \mathcal{O}_{\tau'})
$$
Pour plus de détails, voir la section \ref{sites-cohomology-section-compare}.

\begin{lemma}
\label{sites-cohomology-lemma-compare-topologies-derived-adequate-modules}
Soit $\epsilon : (\mathcal{C}_\tau, \mathcal{O}_\tau) \to
(\mathcal{C}_{\tau'}, \mathcal{O}_{\tau'})$ comme ci-dessus. Soit
$\mathcal{B} \subset \Ob(\mathcal{C})$ une partie. Soit
$\mathcal{A} \subset \textit{PMod}(\mathcal{O})$ une sous-catégorie
pleine. Supposons que
\begin{enumerate}
\item tout objet de $\mathcal{A}$ soit un faisceau pour la topologie $\tau$,
\item $\mathcal{A}$ soit une sous-catégorie de Serre faible de
$\textit{Mod}(\mathcal{O}_\tau)$,
\item tout objet de $\mathcal{C}$ possède un $\tau'$-recouvrement dont les
membres appartiennent à $\mathcal{B}$, et
\item pour tout $U \in \mathcal{B}$, nous ayons
$H^p_\tau(U, \mathcal{F}) = 0$ pour $p > 0$ et tout
$\mathcal{F} \in \mathcal{A}$.
\end{enumerate}
Alors $\mathcal{A}$ est une sous-catégorie de Serre faible de
$\textit{Mod}(\mathcal{O}_{\tau'})$ et il existe une équivalence de
catégories triangulées
$D_\mathcal{A}(\mathcal{O}_\tau) = D_\mathcal{A}(\mathcal{O}_{\tau'})$
donnée par $\epsilon^*$ et $R\epsilon_*$.
\end{lemma}

\begin{proof}
Comme $\epsilon^{-1}\mathcal{O}_{\tau'} = \mathcal{O}_\tau$, nous voyons
que $\epsilon$ est un morphisme plat de sites annelés et que, de fait,
$\epsilon^{-1} = \epsilon^*$ sur les faisceaux de modules. La propriété
(1) permet de considérer tout objet de $\mathcal{A}$ à la fois comme un
faisceau de $\mathcal{O}_\tau$-modules et comme un faisceau de
$\mathcal{O}_{\tau'}$-modules. Autrement dit, nous avons des foncteurs
d'inclusion pleinement fidèles
$$
\mathcal{A} \to \textit{Mod}(\mathcal{O}_\tau) \to
\textit{Mod}(\mathcal{O}_{\tau'})
$$
Pour éviter toute confusion, nous noterons
$\mathcal{A}' \subset \textit{Mod}(\mathcal{O}_{\tau'})$
les images des objets de $\mathcal{A}$. Il est alors clair que
$\epsilon_* : \mathcal{A} \to \mathcal{A}'$ et
$\epsilon^* : \mathcal{A}' \to \mathcal{A}$ sont
des équivalences quasi-inverses (voir la discussion précédant le lemme et
utiliser le fait que les objets de $\mathcal{A}'$ sont des faisceaux pour
la topologie $\tau$).

\medskip\noindent
Les conditions (3) et (4) impliquent que
$R^p\epsilon_*\mathcal{F} = 0$ pour $p > 0$ et
$\mathcal{F} \in \Ob(\mathcal{A})$. En effet, $R^p\epsilon_*$ est le
faisceau associé au préfaisceau $U \mapsto H^p_\tau(U, \mathcal{F})$ ;
voir le Lemme \ref{sites-cohomology-lemma-higher-direct-images}. Ainsi, tout complexe exact
dans $\mathcal{A}$ (ce qui revient à un complexe exact dans
$\textit{Mod}(\mathcal{O}_\tau)$ dont les termes appartiennent à
$\mathcal{A}$ ; voir Homologie, Lemme
\ref{homology-lemma-characterize-weak-serre-subcategory}) reste exact après
application du foncteur $\epsilon_*$.

\medskip\noindent
Considérons une suite exacte
$$
\mathcal{F}'_0 \to \mathcal{F}'_1 \to
\mathcal{F}'_2 \to \mathcal{F}'_3 \to
\mathcal{F}'_4
$$
dans $\textit{Mod}(\mathcal{O}_{\tau'})$, où
$\mathcal{F}'_0, \mathcal{F}'_1, \mathcal{F}'_3, \mathcal{F}'_4$
appartiennent à $\mathcal{A}'$. L'application du foncteur exact
$\epsilon^*$ donne une suite exacte
$$
\epsilon^*\mathcal{F}'_0 \to \epsilon^*\mathcal{F}'_1 \to
\epsilon^*\mathcal{F}'_2 \to \epsilon^*\mathcal{F}'_3 \to
\epsilon^*\mathcal{F}'_4
$$
dans $\textit{Mod}(\mathcal{O}_\tau)$. Comme $\mathcal{A}$ est une
sous-catégorie de Serre faible et que
$\epsilon^*\mathcal{F}'_0, \epsilon^*\mathcal{F}'_1,
\epsilon^*\mathcal{F}'_3, \epsilon^*\mathcal{F}'_4$
appartiennent à $\mathcal{A}$, nous déduisons d'Homologie, Définition
\ref{homology-definition-serre-subcategory}, que
$\epsilon^*\mathcal{F}'_2$ appartient à $\mathcal{A}$. Considérons le
morphisme de suites
$$
\xymatrix{
\mathcal{F}'_0 \ar[r] \ar[d] &
\mathcal{F}'_1 \ar[r] \ar[d] &
\mathcal{F}'_2 \ar[r] \ar[d] &
\mathcal{F}'_3 \ar[r] \ar[d] &
\mathcal{F}'_4 \ar[d] \\
\epsilon_*\epsilon^*\mathcal{F}'_0 \ar[r] &
\epsilon_*\epsilon^*\mathcal{F}'_1 \ar[r] &
\epsilon_*\epsilon^*\mathcal{F}'_2 \ar[r] &
\epsilon_*\epsilon^*\mathcal{F}'_3 \ar[r] &
\epsilon_*\epsilon^*\mathcal{F}'_4
}
$$
La ligne inférieure est exacte d'après la discussion du paragraphe
précédent. Les flèches verticales d'indices $0$, $1$, $3$, $4$ sont des
isomorphismes d'après la discussion du premier paragraphe. Le lemme des $5$
(Homologie, Lemme \ref{homology-lemma-five-lemma}) donne
$\mathcal{F}'_2 \cong \epsilon_*\epsilon^*\mathcal{F}'_2$, et donc
$\mathcal{F}'_2$ appartient à $\mathcal{A}'$. Nous voyons ainsi que
$\mathcal{A}'$ est une sous-catégorie de Serre faible de
$\textit{Mod}(\mathcal{O}_{\tau'})$ ; voir Homologie, Définition
\ref{homology-definition-serre-subcategory}.

\medskip\noindent
À ce stade, il est légitime de parler des catégories dérivées
$D_\mathcal{A}(\mathcal{O}_\tau)$ et
$D_{\mathcal{A}'}(\mathcal{O}_{\tau'})$ ; voir Catégories dérivées,
section \ref{derived-section-triangulated-sub}. Pour achever la
démonstration, montrons que les conditions (1) -- (5) du Lemme
\ref{sites-cohomology-lemma-equivalence-unbounded-two} s'appliquent. Nous avons déjà établi
(1), (2) et (3). Comme tout objet possède un $\tau'$-recouvrement par des
objets de $\mathcal{B}$, il possède a fortiori un $\tau$-recouvrement par
des objets de $\mathcal{B}$. La condition (4) du Lemme
\ref{sites-cohomology-lemma-equivalence-unbounded-two} est donc satisfaite. La condition (5)
l'est de même.
\end{proof}

\begin{lemma}
\label{sites-cohomology-lemma-descent-squares}
Soit $\epsilon : (\mathcal{C}_\tau, \mathcal{O}_\tau) \to
(\mathcal{C}_{\tau'}, \mathcal{O}_{\tau'})$ comme ci-dessus. Soit $A$ un
ensemble et, pour $\alpha \in A$, soit
$$
\xymatrix{
E_\alpha \ar[d] \ar[r] & Y_\alpha \ar[d] \\
Z_\alpha \ar[r] & X_\alpha
}
$$
un diagramme commutatif dans la catégorie $\mathcal{C}$. Supposons que
\begin{enumerate}
\item un $\tau'$-faisceau $\mathcal{F}'$ soit un $\tau$-faisceau dès que
$\mathcal{F}'(X_\alpha) =
\mathcal{F}'(Z_\alpha) \times_{\mathcal{F}'(E_\alpha)} 
\mathcal{F}'(Y_\alpha)$ pour tout $\alpha$,
\item pour $K'$ dans $D(\mathcal{O}_{\tau'})$ appartenant à l'image
essentielle de $R\epsilon_*$, les morphismes
$c^{K'}_{X_\alpha, Z_\alpha, Y_\alpha, E_\alpha}$ du Lemme
\ref{sites-cohomology-lemma-c-square} soient des isomorphismes pour tout $\alpha$.
\end{enumerate}
Alors $K' \in D^+(\mathcal{O}_{\tau'})$ appartient à l'image essentielle de
$R\epsilon_*$ si et seulement si les morphismes
$c^{K'}_{X_\alpha, Z_\alpha, Y_\alpha, E_\alpha}$ sont des isomorphismes
pour tout $\alpha$.
\end{lemma}

\begin{proof}
L'implication directe résulte de l'hypothèse (2). D'autre part, si $K'$
n'a qu'un seul faisceau de cohomologie non nul, l'implication réciproque
résulte de l'hypothèse (1). Dans le cas général, nous allons démontrer
l'implication réciproque par récurrence. Disons qu'un objet $K'$ de
$D^+(\mathcal{O}_{\tau'})$ satisfait (P) si les morphismes
$c^{K'}_{X_\alpha, Z_\alpha, Y_\alpha, E_\alpha}$ sont des isomorphismes
pour tout $\alpha \in A$.

\medskip\noindent
Soit donc $K'$ un objet de $D^+(\mathcal{O}_{\tau'})$ satisfaisant (P).
Choisissons un triangle distingué
$$
K' \to R\epsilon_*\epsilon^{-1}K' \to M' \to K'[1]
$$
dans $D^+(\mathcal{O}_{\tau'})$, où la première flèche est le morphisme
d'adjonction. D'après (2) et le Lemme
\ref{sites-cohomology-lemma-two-out-of-three-blow-up-square}, $M'$ satisfait (P). D'autre
part, en appliquant $\epsilon^{-1}$ et en utilisant l'égalité
$\epsilon^{-1}R\epsilon_* = \text{id}$ de la section
\ref{sites-cohomology-section-compare}, nous obtenons $\epsilon^{-1}M' = 0$. Nous allons
montrer dans le paragraphe suivant que $M' = 0$, ce qui achèvera la
démonstration.

\medskip\noindent
Soit $K'$ un objet de $D^+(\mathcal{O}_{\tau'})$ satisfaisant (P) et tel
que $\epsilon^{-1}K' = 0$. Montrons que $K' = 0$. Étant donné
$n \in \mathbf{Z}$ tel que $H^i(K') = 0$ pour $i < n$, nous allons montrer
que $H^n(K') = 0$. Pour $\alpha \in A$, nous avons un triangle distingué
$$
R\Gamma_{\tau'}(X_\alpha, K') \to
R\Gamma_{\tau'}(Z_\alpha, K') \oplus
R\Gamma_{\tau'}(Y_\alpha, K') \to
R\Gamma_{\tau'}(E_\alpha, K') \to
R\Gamma_{\tau'}(X_\alpha, K')[1]
$$
d'après le Lemme \ref{sites-cohomology-lemma-c-square}. En prenant la cohomologie en degré
$n$ et en utilisant l'annulation supposée des faisceaux de cohomologie de
$K'$, nous obtenons une suite exacte
$$
0 \to
H^n_{\tau'}(X_\alpha, K') \to
H^n_{\tau'}(Z_\alpha, K') \oplus
H^n_{\tau'}(Y_\alpha, K') \to
H^n_{\tau'}(E_\alpha, K')
$$
qui n'est autre que la suite exacte
$$
0 \to
\Gamma(X_\alpha, H^n(K')) \to
\Gamma(Z_\alpha, H^n(K')) \oplus
\Gamma(Y_\alpha, H^n(K')) \to
\Gamma(E_\alpha, H^n(K'))
$$
L'hypothèse (1) montre que $H^n(K')$ est un $\tau$-faisceau. Or son
$\tau$-faisceau associé
$\epsilon^{-1}H^n(K') = H^n(\epsilon^{-1}K')$ est égal à $0$ puisque
$\epsilon^{-1}K' = 0$. Nous en déduisons que $H^n(K') = 0$, comme voulu.
\end{proof}

\begin{lemma}
\label{sites-cohomology-lemma-descent-squares-helper}
Soit $\epsilon : (\mathcal{C}_\tau, \mathcal{O}_\tau) \to
(\mathcal{C}_{\tau'}, \mathcal{O}_{\tau'})$ comme ci-dessus. Soit
$$
\xymatrix{
E \ar[d] \ar[r] & Y \ar[d] \\
Z \ar[r] & X
}
$$
un diagramme commutatif dans la catégorie $\mathcal{C}$ tel que
\begin{enumerate}
\item $h_X^\# = h_Y^\# \amalg_{h_E^\#} h_Z^\#$, et
\item $h_E^\# \to h_Y^\#$ soit injectif
\end{enumerate}
où ${}^\#$ désigne la $\tau$-faisceautisation. Alors, pour
$K' \in D(\mathcal{O}_{\tau'})$ appartenant à l'image essentielle de
$R\epsilon_*$, le morphisme $c^{K'}_{X, Z, Y, E}$ du Lemme
\ref{sites-cohomology-lemma-c-square} (construit avec la topologie $\tau'$) est un
isomorphisme.
\end{lemma}

\begin{proof}
Ce lemme auxiliaire est une conséquence presque immédiate du Lemme
\ref{sites-cohomology-lemma-square-triangle}, et nous recommandons vivement au lecteur d'en
omettre la démonstration. Écrivons $K' = R\epsilon_*K$. Choisissons un
complexe K-injectif de $\mathcal{O}_\tau$-modules, noté
$\mathcal{J}^\bullet$, représentant $K$. Alors
$\epsilon_*\mathcal{J}^\bullet$ est un complexe K-injectif de
$\mathcal{O}_{\tau'}$-modules représentant $K'$ ; voir le Lemme
\ref{sites-cohomology-lemma-K-injective-flat}. Ensuite,
$$
0 \to
\mathcal{J}^\bullet(X) \xrightarrow{\alpha}
\mathcal{J}^\bullet(Z) \oplus
\mathcal{J}^\bullet(Y) \xrightarrow{\beta}
\mathcal{J}^\bullet(E)
\to 0
$$
est une suite exacte courte de complexes de groupes abéliens ; voir le
Lemme \ref{sites-cohomology-lemma-square-triangle} et sa démonstration. Comme il s'agit de
la même suite de complexes de groupes abéliens que celle utilisée pour
définir $c^{K'}_{X, Z, Y, E}$, nous concluons.
\end{proof}
\section{Comparaison des cohomologies}
\label{sites-cohomology-section-compare-general}

\noindent
Nous développons une théorie générale qui nous aidera à comparer
les cohomologies pour différentes topologies. Étant donnés
$\mathcal{C}$, $\tau$ et $\tau'$ comme à la section \ref{sites-cohomology-section-compare},
ainsi qu'un morphisme $f : X \to Y$ de $\mathcal{C}$, nous obtenons un
diagramme commutatif de morphismes de topos
\begin{equation}
\label{sites-cohomology-equation-commutative-epsilon}
\vcenter{
\xymatrix{
\Sh(\mathcal{C}_\tau/X) \ar[r]_{f_\tau} \ar[d]_{\epsilon_X} &
\Sh(\mathcal{C}_\tau/Y) \ar[d]^{\epsilon_Y} \\
\Sh(\mathcal{C}_{\tau'}/X) \ar[r]^{f_{\tau'}} &
\Sh(\mathcal{C}_{\tau'}/Y)
}
}
\end{equation}
Ici, le morphisme $\epsilon_X$, resp.\ $\epsilon_Y$, est le morphisme de
comparaison de la section \ref{sites-cohomology-section-compare} pour la catégorie
$\mathcal{C}/X$ munie des deux topologies $\tau$ et $\tau'$.
Les morphismes $f_\tau$ et $f_{\tau'}$ sont des morphismes de
«~relocalisation~» (Sites, lemme \ref{sites-lemma-relocalize}).
La commutativité du diagramme est un cas particulier de
Sites, lemme \ref{sites-lemma-localize-morphism}
(appliqué avec $\mathcal{C} = \mathcal{C}_\tau/Y$,
$\mathcal{D} = \mathcal{C}_{\tau'}/Y$,
$u = \text{id}$, $U = X$ et $V = X$). Nous obtenons aussi
$\epsilon_{X, *} \circ f_\tau^{-1} = f_{\tau'}^{-1} \circ \epsilon_{Y, *}$
soit par ce lemme, soit directement.

\begin{situation}
\label{sites-cohomology-situation-compare}
On se place avec $\mathcal{C}$, $\tau$ et $\tau'$ comme à la section
\ref{sites-cohomology-section-compare}. Supposons donnés un sous-ensemble
$\mathcal{P} \subset \text{Flèches}(\mathcal{C})$ et, pour tout objet $X$ de
$\mathcal{C}$, une sous-catégorie de Serre faible
$\mathcal{A}'_X \subset \textit{Ab}(\mathcal{C}_{\tau'}/X)$.
Nous faisons les hypothèses suivantes :
\begin{enumerate}
\item
\label{sites-cohomology-item-base-change-P}
si $f : X \to Y$ appartient à $\mathcal{P}$ et si $Y' \to Y$ est quelconque,
alors $X \times_Y Y'$ existe et $X \times_Y Y' \to Y'$ appartient à $\mathcal{P}$,
\item
\label{sites-cohomology-item-restriction-A}
$f_{\tau'}^{-1}$ envoie $\mathcal{A}'_Y$ dans $\mathcal{A}'_X$
pour tout morphisme $f : X \to Y$ de $\mathcal{C}$,
\item
\label{sites-cohomology-item-A-sheaf}
si $X$ est dans $\mathcal{C}$ et $\mathcal{F}'$ dans $\mathcal{A}'_X$, alors
$\mathcal{F}'$ satisfait la condition de faisceau pour les recouvrements
$\tau$, c'est-à-dire
$\mathcal{F}' = \epsilon_{X, *}\epsilon_X^{-1}\mathcal{F}'$,
\item
\label{sites-cohomology-item-A-and-P}
si $f : X \to Y$ appartient à $\mathcal{P}$ et si
$\mathcal{F}' \in \Ob(\mathcal{A}'_X)$, alors
$R^if_{\tau', *}\mathcal{F}' \in \Ob(\mathcal{A}'_Y)$
pour $i \geq 0$.
\item
\label{sites-cohomology-item-refine-tau-by-P}
si $\{U_i \to U\}_{i \in I}$ est un recouvrement $\tau$, alors il existe
\begin{enumerate}
\item un recouvrement $\tau'$ $\{V_j \to U\}_{j \in J}$,
\item un recouvrement $\tau$ $\{f_j : W_j \to V_j\}$ réduit à un seul
morphisme $f_j \in \mathcal{P}$, et
\item un recouvrement $\tau'$ $\{W_{jk} \to W_j\}_{k \in K_j}$
\end{enumerate}
tels que $\{W_{jk} \to U\}_{j \in J, k \in K_j}$ soit un raffinement de
$\{U_i \to U\}_{i \in I}$.
\end{enumerate}
\end{situation}

\begin{lemma}
\label{sites-cohomology-lemma-A}
Dans la situation \ref{sites-cohomology-situation-compare}, pour $X$ dans $\mathcal{C}$,
notons $\mathcal{A}_X$ la classe des objets de
$\textit{Ab}(\mathcal{C}_\tau/X)$ de la forme
$\epsilon_X^{-1}\mathcal{F}'$, avec $\mathcal{F}'$ dans $\mathcal{A}'_X$.
Alors
\begin{enumerate}
\item pour $\mathcal{F}$ dans $\textit{Ab}(\mathcal{C}_\tau/X)$,
nous avons $\mathcal{F} \in \mathcal{A}_X \Leftrightarrow
\epsilon_{X, *}\mathcal{F} \in \mathcal{A}'_X$, et
\item $f_\tau^{-1}$ envoie $\mathcal{A}_Y$ dans $\mathcal{A}_X$
pour tout morphisme $f : X \to Y$ de $\mathcal{C}$.
\end{enumerate}
\end{lemma}

\begin{proof}
L'assertion (1) résulte de (\ref{sites-cohomology-item-A-sheaf}) et l'assertion (2)
résulte de (\ref{sites-cohomology-item-restriction-A}) et de la commutativité de
(\ref{sites-cohomology-equation-commutative-epsilon}), qui donne
$\epsilon_X^{-1} \circ f_{\tau'}^{-1} = f_\tau^{-1} \circ \epsilon_Y^{-1}$.
\end{proof}

\noindent
Notre prochain objectif est de démontrer les lemmes
\ref{sites-cohomology-lemma-compare-cohomology-general} et
\ref{sites-cohomology-lemma-cohomological-descent-general}. Nous procéderons par récurrence
en utilisant l'hypothèse de récurrence suivante.

\medskip\noindent
$(V_n)$ Pour $X$ dans $\mathcal{C}$ et $\mathcal{F}$ dans $\mathcal{A}_X$,
nous avons $R^i\epsilon_{X, *}\mathcal{F} = 0$ pour $1 \leq i \leq n$.

\begin{lemma}
\label{sites-cohomology-lemma-V-implies-C-general}
Dans la situation \ref{sites-cohomology-situation-compare}, supposons $(V_n)$ vérifiée.
Pour $f : X \to Y$ dans $\mathcal{P}$ et $\mathcal{F}$ dans $\mathcal{A}_X$,
nous avons $R^if_{\tau', *}\epsilon_{X, *}\mathcal{F} =
\epsilon_{Y, *}R^if_{\tau, *}\mathcal{F}$ pour $i \leq n$.
\end{lemma}

\begin{proof}
Nous utiliserons le diagramme commutatif
(\ref{sites-cohomology-equation-commutative-epsilon}) sans plus le mentionner. En particulier,
nous avons
$$
Rf_{\tau', *}R\epsilon_{X, *}\mathcal{F} =
R\epsilon_{Y, *}Rf_{\tau, *}\mathcal{F}
$$
L'hypothèse $(V_n)$ nous dit que
$\epsilon_{X, *}\mathcal{F} \to R\epsilon_{X, *}\mathcal{F}$
est un isomorphisme en degrés $\leq n$. Par conséquent,
$Rf_{\tau', *}\epsilon_{X, *}\mathcal{F} \to
Rf_{\tau', *}R\epsilon_{X, *}\mathcal{F}$ est un isomorphisme
en degrés $\leq n$. Nous en concluons que
$$
R^if_{\tau', *}\epsilon_{X, *}\mathcal{F} \to
H^i(R\epsilon_{Y, *}Rf_{\tau, *}\mathcal{F})
$$
est un isomorphisme pour $i \leq n$. Nous démontrerons le lemme en examinant
la deuxième page de la suite spectrale du
lemme \ref{sites-cohomology-lemma-relative-Leray} pour
$R\epsilon_{Y, *}Rf_{\tau, *}\mathcal{F}$. En voici une représentation :
$$
\begin{matrix}
\ldots &
\ldots &
\ldots &
\ldots \\
\epsilon_{Y, *}R^2f_{\tau, *}\mathcal{F} &
R^1\epsilon_{Y, *}R^2f_{\tau, *}\mathcal{F} &
R^2\epsilon_{Y, *}R^2f_{\tau, *}\mathcal{F} &
\ldots \\
\epsilon_{Y, *}R^1f_{\tau, *}\mathcal{F} &
R^1\epsilon_{Y, *}R^1f_{\tau, *}\mathcal{F} &
R^2\epsilon_{Y, *}R^1f_{\tau, *}\mathcal{F} &
\ldots \\
\epsilon_{Y, *}f_{\tau, *}\mathcal{F} &
R^1\epsilon_{Y, *}f_{\tau, *}\mathcal{F} &
R^2\epsilon_{Y, *}f_{\tau, *}\mathcal{F} &
\ldots
\end{matrix}
$$
Notons $(C_m)$ l'hypothèse : $R^if_{\tau', *}\epsilon_{X, *}\mathcal{F} =
\epsilon_{Y, *}R^if_{\tau, *}\mathcal{F}$ pour $i \leq m$. Observons que
$(C_0)$ est vérifiée. Nous allons montrer que
$(C_{m - 1}) \Rightarrow (C_m)$ pour $m < n$. En effet, si
$(C_{m - 1})$ est vérifiée, alors, pour
$n \geq p > 0$ et $q \leq m - 1$, nous avons
\begin{align*}
R^p\epsilon_{Y, *}R^qf_{\tau, *}\mathcal{F}
& =
R^p\epsilon_{Y, *}
\epsilon_Y^{-1} \epsilon_{Y, *} R^qf_{\tau, *}\mathcal{F} \\
& =
R^p\epsilon_{Y, *}
\epsilon_Y^{-1}R^qf_{\tau', *}\epsilon_{X, *}\mathcal{F} = 0
\end{align*}
La première égalité vient de
$\epsilon_Y^{-1}\epsilon_{Y, *} = \text{id}$, la deuxième de
$(C_{m - 1})$, et la dernière de $(V_n)$, car
$\epsilon_Y^{-1}R^qf_{\tau', *}\epsilon_{X, *}\mathcal{F}$
appartient à $\mathcal{A}_Y$ par (\ref{sites-cohomology-item-A-and-P}).
En examinant la suite spectrale, nous voyons que
$E_2^{0, m} = \epsilon_{Y, *}R^mf_{\tau, *}\mathcal{F}$
est le seul terme $E_2^{p, q}$ non nul avec $p + q = m$.
Rappelons que $\text{d}_r^{p, q} : E_r^{p, q} \to E_r^{p + r, q - r + 1}$.
Il n'y a donc aucune différentielle non nulle
$\text{d}_r^{p, q}$, $r \geq 2$, partant de cette position ou y arrivant.
Nous en concluons que
$H^m(R\epsilon_{Y, *}Rf_{\tau, *}\mathcal{F}) =
\epsilon_{Y, *}R^mf_{\tau, *}\mathcal{F}$, ce qui implique
$(C_m)$ d'après la discussion précédente.

\medskip\noindent
Enfin, supposons $(C_{n - 1})$. La même analyse montre que
$E_2^{0, n} = \epsilon_{Y, *}R^nf_{\tau, *}\mathcal{F}$
est le seul terme $E_2^{p, q}$ non nul avec $p + q = n$.
Il n'y a toujours aucune différentielle non nulle arrivant à cette position,
mais une différentielle non nulle peut en partir, à savoir le morphisme
$d_{n + 1}^{0, n} : \epsilon_{Y, *}R^nf_{\tau, *}\mathcal{F} \to
R^{n + 1}\epsilon_{Y, *}f_{\tau, *}\mathcal{F}$.
Nous en concluons qu'il existe une suite exacte
$$
0 \to R^nf_{\tau', *}\epsilon_{X, *}\mathcal{F} \to 
\epsilon_{Y, *}R^nf_{\tau, *}\mathcal{F} \to
R^{n + 1}\epsilon_{Y, *}f_{\tau, *}\mathcal{F}
$$
Par (\ref{sites-cohomology-item-A-and-P}) et (\ref{sites-cohomology-item-A-sheaf}), le faisceau
$R^nf_{\tau', *}\epsilon_{X, *}\mathcal{F}$
satisfait la condition de faisceau pour les recouvrements $\tau$, tout comme
$\epsilon_{Y, *}R^nf_{\tau, *}\mathcal{F}$
(utiliser la description de $\epsilon_*$ à la section \ref{sites-cohomology-section-compare}).
Cependant, le faisceautisé pour $\tau$ du faisceau pour $\tau'$
$R^{n + 1}\epsilon_{Y, *}f_{\tau, *}\mathcal{F}$
est nul (par localité de la cohomologie ; utiliser les
lemmes \ref{sites-cohomology-lemma-kill-cohomology-class-on-covering} et
\ref{sites-cohomology-lemma-higher-direct-images}).
Ainsi, $R^nf_{\tau', *}\epsilon_{X, *}\mathcal{F} \to 
\epsilon_{Y, *}R^nf_{\tau, *}\mathcal{F}$
doit donc être un isomorphisme, ce qui achève la démonstration.
\end{proof}

\noindent
Si $E'$, resp.\ $E$, est un objet de
$D(\mathcal{C}_{\tau'}/X)$, resp.\ $D(\mathcal{C}_\tau/X)$,
nous noterons
$H^n_{\tau'}(U, E')$, resp.\ $H^n_\tau(U, E)$,
la cohomologie de $E'$, resp.\ $E$,
au-dessus d'un objet $U$ de $\mathcal{C}/X$.

\begin{lemma}
\label{sites-cohomology-lemma-V-implies-cohomology-general}
Dans la situation \ref{sites-cohomology-situation-compare}, si $(V_n)$ est vérifiée, alors,
pour $X$ dans $\mathcal{C}$ et $L \in D(\mathcal{C}_{\tau'}/X)$
tels que $H^i(L) = 0$ pour $i < 0$ et que $H^i(L)$ appartienne à
$\mathcal{A}'_X$ pour $0 \leq i \leq n$, nous avons
$H^n_{\tau'}(X, L) = H^n_\tau(X, \epsilon_X^{-1}L)$.
\end{lemma}

\begin{proof}
Par le lemme \ref{sites-cohomology-lemma-Leray-unbounded}, nous avons
$H^n_\tau(X, \epsilon_X^{-1}L) =
H^n_{\tau'}(X, R\epsilon_{X, *}\epsilon_X^{-1}L)$.
Il existe une suite spectrale
$$
E_2^{p, q} = R^p\epsilon_{X, *}\epsilon_X^{-1}H^q(L)
$$
qui converge vers $H^{p + q}(R\epsilon_{X, *}\epsilon_X^{-1}L)$.
Par $(V_n)$, le terme $E_2^{p, q}$ est nul pour
$0 < p \leq n$ et $0 \leq q \leq n$. Ainsi, les termes
$E_2^{0, q} = \epsilon_{X, *}\epsilon_X^{-1}H^q(L) = H^q(L)$
sont les seuls termes $E_2^{p, q}$ non nuls avec $p + q \leq n$.
Il s'ensuit que le morphisme
$$
L \longrightarrow R\epsilon_{X, *}\epsilon_X^{-1}L
$$
est un isomorphisme en degrés $\leq n$ (on omet un détail mineur).
Nous obtenons donc
$H^i_{\tau'}(X, L) = H^i_{\tau'}(X, R\epsilon_{X, *}\epsilon_X^{-1}L)$
pour $i \leq n$. Le lemme est ainsi démontré.
\end{proof}

\begin{lemma}
\label{sites-cohomology-lemma-V-implies-cohomology-extra-general}
Dans la situation \ref{sites-cohomology-situation-compare}, si $(V_n)$ est vérifiée, alors,
pour $X$ dans $\mathcal{C}$ et $\mathcal{F}$ dans $\mathcal{A}_X$, le morphisme
$H^{n + 1}_{\tau'}(X, \epsilon_{X, *}\mathcal{F}) \to
H^{n + 1}_\tau(X, \mathcal{F})$
est injectif et son image est formée des classes qui deviennent nulles sur
un recouvrement $\tau'$ de $X$.
\end{lemma}

\begin{proof}
Rappelons que $\epsilon_X^{-1}\epsilon_{X, *}\mathcal{F} = \mathcal{F}$ ;
le morphisme est donc donné par l'image inverse des classes de cohomologie
par $\epsilon_X$. La suite spectrale de Leray
(lemme \ref{sites-cohomology-lemma-Leray})
$$
E_2^{p, q} = H^p_{\tau'}(X, R^q\epsilon_{X, *}\mathcal{F})
\Rightarrow
H^{p + q}_\tau(X, \mathcal{F})
$$
combinée à l'annulation supposée donne une suite exacte
$$
0 \to
H^{n + 1}_{\tau'}(X, \epsilon_{X, *}\mathcal{F}) \to
H^{n + 1}_\tau(X, \mathcal{F}) \to
H^0_{\tau'}(X, R^{n + 1}\epsilon_{X, *}\mathcal{F})
$$
C'est une reformulation du lemme.
\end{proof}

\begin{lemma}
\label{sites-cohomology-lemma-make-class-zero-general}
Dans la situation \ref{sites-cohomology-situation-compare}, soit $f : X \to Y$
un élément de $\mathcal{P}$ tel que $\{X \to Y\}$ soit un recouvrement
$\tau$. Soit $\mathcal{F}'$ dans $\mathcal{A}'_Y$. Si $n \geq 0$ et
$$
\theta \in
\text{Égaliseur}\left(
\xymatrix{
H^{n + 1}_{\tau'}(X, \mathcal{F}')
\ar@<1ex>[r] \ar@<-1ex>[r] &
H^{n + 1}_{\tau'}(X \times_Y X, \mathcal{F}')
}
\right)
$$
alors il existe un recouvrement $\tau'$ $\{Y_i \to Y\}$
tel que la restriction de $\theta$ soit nulle dans
$H^{n + 1}_{\tau'}(Y_i \times_Y X, \mathcal{F}')$.
\end{lemma}

\begin{proof}
Observons que $X \times_Y X$ existe d'après (\ref{sites-cohomology-item-base-change-P}).
Pour $Z$ dans $\mathcal{C}/Y$, notons $\mathcal{F}'|_Z$
la restriction de $\mathcal{F}'$ à $\mathcal{C}_{\tau'}/Z$.
Rappelons que $H^{n + 1}_{\tau'}(X, \mathcal{F}') =
H^{n + 1}(\mathcal{C}_{\tau'}/X, \mathcal{F}'|_X)$, voir le
lemme \ref{sites-cohomology-lemma-cohomology-of-open}.
Le lemme affirme que l'image
$\overline{\theta} \in H^0(Y, R^{n + 1}f_{\tau', *}\mathcal{F}'|_X)$
de $\theta$ est nulle. Considérons le diagramme cartésien
$$
\xymatrix{
X \times_Y X \ar[d]_{\text{pr}_1} \ar[r]_{\text{pr}_2} &
X \ar[d]^f \\
X \ar[r]^f & Y
}
$$
Par changement de base trivial (lemme \ref{sites-cohomology-lemma-localize-cartesian-square}),
nous avons
$$
f_{\tau'}^{-1}R^{n + 1}f_{\tau', *}(\mathcal{F}'|_X) =
R^{n + 1}\text{pr}_{1, \tau', *}(\mathcal{F}'|_{X \times_Y X})
$$
Si $\text{pr}_1^{-1}\theta = \text{pr}_2^{-1}\theta$,
alors la section $f_{\tau'}^{-1}\overline{\theta}$ de
$f_{\tau'}^{-1}R^{n + 1}f_{\tau', *}(\mathcal{F}'|_X)$ est nulle,
car il est clair que $\text{pr}_1^{-1}\theta$ s'envoie sur l'élément nul de
$H^0(X, R^{n + 1}\text{pr}_{1, \tau', *}(\mathcal{F}'|_{X \times_Y X}))$.
D'après (\ref{sites-cohomology-item-restriction-A}), $\mathcal{F}'|_X$ appartient à
$\mathcal{A}'_X$. Ainsi,
$\mathcal{G}' = R^{n + 1}f_{\tau', *}(\mathcal{F}'|_X)$
est un objet de $\mathcal{A}'_Y$ par (\ref{sites-cohomology-item-A-and-P}).
Par conséquent, $\mathcal{G}'$ satisfait la condition de faisceau pour les
recouvrements $\tau$ d'après (\ref{sites-cohomology-item-A-sheaf}).
Comme $\{X \to Y\}$ est un recouvrement $\tau$, l'application de restriction
$\mathcal{G}'(Y) \to \mathcal{G}'(X)$ est injective.
Il s'ensuit que $\overline{\theta}$ est nulle.
\end{proof}

\begin{lemma}
\label{sites-cohomology-lemma-induction-step-V-C-general}
Dans la situation \ref{sites-cohomology-situation-compare}, nous avons
$(V_n) \Rightarrow (V_{n + 1})$.
\end{lemma}

\begin{proof}
Soient $X$ dans $\mathcal{C}$ et $\mathcal{F}$ dans $\mathcal{A}_X$.
Soit $\xi \in H^{n + 1}_\tau(U, \mathcal{F})$ pour un certain $U/X$.
Nous devons montrer que la restriction de $\xi$ est nulle sur les membres
d'un recouvrement $\tau'$ de $U$, voir le
lemme \ref{sites-cohomology-lemma-higher-direct-images}.
Il en résulte que nous pouvons remplacer $U$ par les membres d'un
recouvrement $\tau'$ de $U$.

\medskip\noindent
Par localité de la cohomologie
(lemme \ref{sites-cohomology-lemma-kill-cohomology-class-on-covering}),
nous pouvons choisir un recouvrement $\tau$ $\{U_i \to U\}$
tel que la restriction de $\xi$ à $U_i$ soit nulle.
Choisissons $\{V_j \to V\}$, $\{f_j : W_j \to V_j\}$ et
$\{W_{jk} \to W_j\}$ comme dans (\ref{sites-cohomology-item-refine-tau-by-P}).
En remplaçant $U$ par $V_j$ et $\mathcal{F}$ par sa restriction à
$\mathcal{C}_\tau/V_j$, ce qui est permis par
(\ref{sites-cohomology-item-base-change-P}), nous nous ramenons au cas étudié au paragraphe
suivant.

\medskip\noindent
Ici, $f : X \to Y$ est un élément de $\mathcal{P}$ tel que
$\{X \to Y\}$ soit un recouvrement $\tau$,
$\mathcal{F}$ est un objet de $\mathcal{A}_Y$, et
$\xi \in H^{n + 1}_\tau(Y, \mathcal{F})$ est tel qu'il existe un
recouvrement $\tau'$ $\{X_i \to X\}_{i \in I}$ tel que $\xi$ se restreigne
à zéro sur $X_i$ pour tout $i \in I$.
Il s'agit de montrer que la restriction de $\xi$ est nulle sur un
recouvrement $\tau'$ de $Y$.

\medskip\noindent
Par le lemme \ref{sites-cohomology-lemma-V-implies-cohomology-extra-general}, il existe une
unique classe de cohomologie $\tau'$
$\theta \in H^{n + 1}_{\tau'}(X, \epsilon_{X, *}\mathcal{F})$
dont l'image est $\xi|_X$.
Comme les images inverses de $\xi|_X$ par les deux projections donnent la
même classe sur $X \times_Y X$, il en est de même pour $\theta$
(par unicité).
D'après le lemme \ref{sites-cohomology-lemma-make-class-zero-general}, après avoir remplacé
$Y$ par les membres d'un recouvrement $\tau'$, nous pouvons supposer que
$\theta = 0$. Nous pouvons donc supposer que $\xi|_X$ est nulle.

\medskip\noindent
Soit $f : X \to Y$ un élément de $\mathcal{P}$ tel que
$\{X \to Y\}$ soit un recouvrement $\tau$,
soit $\mathcal{F}$ un objet de $\mathcal{A}_Y$, et soit
$\xi \in H^{n + 1}_\tau(Y, \mathcal{F})$ dont l'image dans
$H^{n + 1}_\tau(X, \mathcal{F})$ est nulle.
Il faut montrer que la restriction de $\xi$ est nulle sur un recouvrement
$\tau'$ de $Y$.

\medskip\noindent
Les hypothèses nous disent que l'image de $\xi$ par le morphisme
$$
\mathcal{F} \longrightarrow Rf_{\tau, *}f_\tau^{-1}\mathcal{F}
$$
est nulle ; on utilise ici le lemme \ref{sites-cohomology-lemma-Leray-unbounded}.
Un argument simple fondé sur le triangle distingué de troncation
(Catégories dérivées, remarque
\ref{derived-remark-truncation-distinguished-triangle}) montre que
l'image de $\xi$ par le morphisme
$$
\mathcal{F} \longrightarrow \tau_{\leq n}Rf_{\tau, *}f_\tau^{-1}\mathcal{F}
$$
est nulle. Nous allons comparer celui-ci au morphisme
$\epsilon_{Y, *}\mathcal{F} \to K$, où
$$
K = \tau_{\leq n}Rf_{\tau', *}f_{\tau'}^{-1}\epsilon_{Y, *}\mathcal{F} =
\tau_{\leq n}Rf_{\tau', *}\epsilon_{X, *}f_{\tau}^{-1}\mathcal{F}
$$
L'égalité
$\epsilon_{X, *} f_\tau^{-1} = f_{\tau'}^{-1} \epsilon_{Y, *}$
est une propriété de (\ref{sites-cohomology-equation-commutative-epsilon}). Considérons le
morphisme
$$
Rf_{\tau', *}\epsilon_{X, *}f_{\tau}^{-1}\mathcal{F} \longrightarrow
Rf_{\tau', *}R\epsilon_{X, *}f_{\tau}^{-1}\mathcal{F} =
R\epsilon_{Y, *}Rf_{\tau, *}f_\tau^{-1}\mathcal{F}
$$
utilisé dans la démonstration du lemme \ref{sites-cohomology-lemma-V-implies-C-general} ;
il induit par adjonction un morphisme
$$
\epsilon_Y^{-1} Rf_{\tau', *}\epsilon_{X, *}f_{\tau}^{-1}\mathcal{F} \to
Rf_{\tau, *}f_\tau^{-1}\mathcal{F}
$$
En prenant les troncations, nous obtenons un morphisme
$$
\epsilon_Y^{-1}K
\longrightarrow
\tau_{\leq n}Rf_{\tau, *}f_\tau^{-1}\mathcal{F}
$$
qui est un isomorphisme par le lemme \ref{sites-cohomology-lemma-V-implies-C-general} ;
ce lemme s'applique, car $f_\tau^{-1}\mathcal{F}$ appartient à
$\mathcal{A}_X$ d'après le lemme \ref{sites-cohomology-lemma-A}. Choisissons un triangle
distingué
$$
\epsilon_{Y, *}\mathcal{F} \to K \to L \to \epsilon_{Y, *}\mathcal{F}[1]
$$
Le morphisme $\mathcal{F} \to f_{\tau, *}f_\tau^{-1}\mathcal{F}$
est injectif puisque $\{X \to Y\}$ est un recouvrement $\tau$. Ainsi,
$\epsilon_{Y, *}\mathcal{F} \to
\epsilon_{Y, *}f_{\tau, *}f_\tau^{-1}\mathcal{F} =
f_{\tau', *}f_{\tau'}^{-1}\epsilon_{Y, *}\mathcal{F}$
est lui aussi injectif. Par conséquent, les seuls faisceaux de cohomologie
non nuls de $L$ sont en degrés $0, \ldots, n$.
Comme $f_{\tau', *}f_{\tau'}^{-1}\epsilon_{Y, *}\mathcal{F}$
appartient à $\mathcal{A}'_Y$ d'après (\ref{sites-cohomology-item-restriction-A}) et
(\ref{sites-cohomology-item-A-and-P}), nous en concluons que
$$
H^0(L) = \Coker(\epsilon_{Y, *}\mathcal{F} \to
f_{\tau', *}f_{\tau'}^{-1}\epsilon_{Y, *}\mathcal{F})
$$
appartient à la sous-catégorie de Serre faible $\mathcal{A}'_Y$. Pour
$1 \leq i \leq n$, nous voyons que
$H^i(L) = R^if_{\tau', *}f_{\tau'}^{-1}\epsilon_{Y, *}\mathcal{F}$
appartient à $\mathcal{A}'_Y$ d'après (\ref{sites-cohomology-item-restriction-A}) et
(\ref{sites-cohomology-item-A-and-P}). En tirant le triangle distingué ci-dessus en arrière
par $\epsilon_Y$, nous obtenons le triangle distingué
$$
\mathcal{F} \to \tau_{\leq n}Rf_{\tau, *}f_\tau^{-1}\mathcal{F}
\to \epsilon_Y^{-1}L \to \mathcal{F}[1]
$$
Puisque l'image de $\xi$ dans le terme du milieu est nulle, nous voyons que
$\xi$ est l'image d'un élément
$\xi' \in H^n_\tau(Y, \epsilon_Y^{-1}L)$.
D'après le lemme \ref{sites-cohomology-lemma-V-implies-cohomology-general}, nous avons
$$
H^n_{\tau'}(Y, L) = H^n_\tau(Y, \epsilon_Y^{-1}L),
$$
Nous pouvons donc relever $\xi'$ en un élément de $H^n_{\tau'}(Y, L)$,
puis prendre son bord dans
$H^{n + 1}_{\tau'}(Y, \epsilon_{Y, *}\mathcal{F})$
pour voir que $\xi$ appartient à l'image du morphisme canonique
$H^{n + 1}_{\tau'}(Y, \epsilon_{Y, *}\mathcal{F}) \to
H^{n + 1}_\tau(Y, \mathcal{F})$.
La localité de la cohomologie pour
$H^{n + 1}_{\tau'}(Y,\epsilon_{Y, *}\mathcal{F})$, voir le
lemme \ref{sites-cohomology-lemma-kill-cohomology-class-on-covering}, permet de conclure.
\end{proof}

\begin{lemma}
\label{sites-cohomology-lemma-V-C-all-n-general}
Dans la situation \ref{sites-cohomology-situation-compare}, la propriété $(V_n)$ est vraie
pour tout $n$. De plus :
\begin{enumerate}
\item Pour $X$ dans $\mathcal{C}$ et
$K' \in D^+_{\mathcal{A}'_X}(\mathcal{C}_{\tau'}/X)$, le morphisme
$K' \to R\epsilon_{X, *}(\epsilon_X^{-1}K')$ est un isomorphisme.
\item Pour $f : X \to Y$ dans $\mathcal{P}$ et
$K' \in D^+_{\mathcal{A}'_X}(\mathcal{C}_{\tau'}/X)$, nous avons
$Rf_{\tau', *}K' \in D^+_{\mathcal{A}'_Y}(\mathcal{C}_{\tau'}/Y)$ et
$\epsilon_Y^{-1}(Rf_{\tau', *}K') = Rf_{\tau, *}(\epsilon_X^{-1}K')$.
\end{enumerate}
\end{lemma}

\begin{proof}
Observons que $(V_0)$ est vérifiée, car c'est la condition vide.
Nous obtenons alors $(V_n)$ pour tout $n$ par le
lemme \ref{sites-cohomology-lemma-induction-step-V-C-general}.

\medskip\noindent
Démonstration de (1). L'objet $K = \epsilon_X^{-1}K'$ a pour faisceaux de
cohomologie $H^i(K) = \epsilon_X^{-1}H^i(K')$, qui appartiennent à
$\mathcal{A}_X$. La suite spectrale
$$
E_2^{p, q} = R^p\epsilon_{X, *} H^q(K) \Rightarrow
H^{p + q}(R\epsilon_{X, *}K)
$$
dégénère donc, par $(V_n)$ pour tout $n$, et nous obtenons
$$
H^n(R\epsilon_{X, *}K) = \epsilon_{X, *}H^n(K) =
\epsilon_{X, *}\epsilon_X^{-1}H^n(K') = H^n(K').
$$
La dernière égalité vient encore de ce que $H^n(K')$ appartient à
$\mathcal{A}'_X$. Le morphisme canonique
$K' \to R\epsilon_{X, *}(\epsilon_X^{-1}K')$ est donc un isomorphisme.

\medskip\noindent
Démonstration de (2). En utilisant la suite spectrale
$$
E_2^{p, q} = R^pf_{\tau', *}H^q(K') \Rightarrow R^{p + q}f_{\tau', *}K'
$$
le fait que $R^pf_{\tau', *}H^q(K')$ appartient à
$\mathcal{A}'_Y$ par (\ref{sites-cohomology-item-A-and-P}), le fait que $\mathcal{A}'_Y$ est
une sous-catégorie de Serre faible de
$\textit{Ab}(\mathcal{C}_{\tau'}/Y)$, et Homologie, lemme
\ref{homology-lemma-biregular-ss-converges}, nous concluons que
$Rf_{\tau', *}K' \in D^+_{\mathcal{A}'_Y}(\mathcal{C}_{\tau'}/Y)$.
Pour achever la démonstration, il faut montrer que le morphisme de changement
de base
$$
\epsilon_Y^{-1}(Rf_{\tau', *}K')
\longrightarrow
Rf_{\tau, *}(\epsilon_X^{-1}K')
$$
est un isomorphisme. En comparant la suite spectrale ci-dessus à la suite
spectrale
$$
E_2^{p, q} = R^pf_{\tau, *}H^q(\epsilon_X^{-1}K')
\Rightarrow R^{p + q}f_{\tau, *}\epsilon_X^{-1}K'
$$
nous nous ramenons au cas où $K'$ possède un unique faisceau de cohomologie
non nul $\mathcal{F}'$, appartenant à $\mathcal{A}'_X$ ; nous omettons les
détails. Le lemme \ref{sites-cohomology-lemma-V-implies-C-general} donne alors
$\epsilon_Y^{-1}R^if_{\tau', *}\mathcal{F}' =
R^if_{\tau, *}\epsilon_X^{-1}\mathcal{F}'$ pour tout $i$,
ce qui achève la démonstration.
\end{proof}

\begin{lemma}
\label{sites-cohomology-lemma-cohomological-descent-general}
Dans la situation \ref{sites-cohomology-situation-compare}, pour tout $X$ dans
$\mathcal{C}$, la catégorie
$\mathcal{A}_X \subset \textit{Ab}(\mathcal{C}_\tau/X)$
est une sous-catégorie de Serre faible, et le foncteur
$$
R\epsilon_{X, *} :
D^+_{\mathcal{A}_X}(\mathcal{C}_\tau/X)
\longrightarrow
D^+_{\mathcal{A}'_X}(\mathcal{C}_{\tau'}/X)
$$
est une équivalence dont un quasi-inverse est donné par $\epsilon_X^{-1}$.
\end{lemma}

\begin{proof}
Nous devons vérifier pour $\mathcal{A}_X$ les conditions énumérées dans
Homologie, lemme \ref{homology-lemma-characterize-weak-serre-subcategory}.
Si $\varphi : \mathcal{F} \to \mathcal{G}$ est un morphisme dans
$\mathcal{A}_X$, alors $\epsilon_{X, *}\varphi :
\epsilon_{X, *}\mathcal{F} \to \epsilon_{X, *}\mathcal{G}$
est un morphisme dans $\mathcal{A}'_X$. Par conséquent,
$\Ker(\epsilon_{X, *}\varphi)$ et $\Coker(\epsilon_{X, *}\varphi)$ sont des
objets de $\mathcal{A}'_X$, puisque celle-ci est une sous-catégorie de Serre
faible de $\textit{Ab}(\mathcal{C}_{\tau'}/X)$. En appliquant
$\epsilon_X^{-1}$, nous obtenons une suite exacte
$$
0 \to
\epsilon_X^{-1}\Ker(\epsilon_{X, *}\varphi) \to
\mathcal{F} \to \mathcal{G} \to
\epsilon_X^{-1}\Coker(\epsilon_{X, *}\varphi) \to 0
$$
et nous voyons que $\Ker(\varphi)$ et $\Coker(\varphi)$ appartiennent à
$\mathcal{A}_X$. Enfin, supposons que
$$
0 \to \mathcal{F}_1 \to \mathcal{F}_2 \to \mathcal{F}_3 \to 0
$$
soit une suite exacte courte dans $\textit{Ab}(\mathcal{C}_\tau/X)$,
avec $\mathcal{F}_1$ et $\mathcal{F}_3$ dans $\mathcal{A}_X$.
En appliquant $\epsilon_{X, *}$, nous obtenons une suite exacte
$$
0 \to 
\epsilon_{X, *}\mathcal{F}_1 \to
\epsilon_{X, *}\mathcal{F}_2 \to
\epsilon_{X, *}\mathcal{F}_3 \to
R^1\epsilon_{X, *}\mathcal{F}_1 = 0
$$
L'annulation vient du lemme \ref{sites-cohomology-lemma-V-C-all-n-general}.
Ainsi, $\epsilon_{X, *}\mathcal{F}_2$ appartient à $\mathcal{A}'_X$, car
celle-ci est une sous-catégorie de Serre faible de
$\textit{Ab}(\mathcal{C}_{\tau'}/X)$. En tirant en arrière par $\epsilon_X$,
nous concluons que $\mathcal{F}_2$ appartient à $\mathcal{A}_X$.

\medskip\noindent
Ainsi, $\mathcal{A}_X$ est une sous-catégorie de Serre faible de
$\textit{Ab}(\mathcal{C}_\tau/X)$, et il est légitime de considérer la
catégorie $D^+_{\mathcal{A}_X}(\mathcal{C}_\tau/X)$.
Observons que $\epsilon_X^{-1} : \mathcal{A}'_X \to \mathcal{A}_X$
est une équivalence et que
$\mathcal{F}' \to R\epsilon_{X, *}\epsilon_X^{-1}\mathcal{F}'$
est un isomorphisme pour $\mathcal{F}'$ dans $\mathcal{A}'_X$, puisque
$(V_n)$ est vérifiée pour tout $n$ d'après le
lemme \ref{sites-cohomology-lemma-V-C-all-n-general}. Nous concluons donc par le
lemme \ref{sites-cohomology-lemma-equivalence-bounded}.
\end{proof}

\begin{lemma}
\label{sites-cohomology-lemma-compare-cohomology-general}
Dans la situation \ref{sites-cohomology-situation-compare}, soit $X$ dans $\mathcal{C}$.
\begin{enumerate}
\item pour $\mathcal{F}'$ dans $\mathcal{A}'_X$, nous avons
$H^n_{\tau'}(X, \mathcal{F}') = H^n_\tau(X, \epsilon_X^{-1}\mathcal{F}')$,
\item pour $K' \in D^+_{\mathcal{A}'_X}(\mathcal{C}_{\tau'}/X)$,
nous avons $H^n_{\tau'}(X, K') = H^n_\tau(X, \epsilon_X^{-1}K')$.
\end{enumerate}
\end{lemma}

\begin{proof}
Cela résulte du lemme \ref{sites-cohomology-lemma-V-C-all-n-general} et de la
remarque \ref{sites-cohomology-remark-before-Leray}.
\end{proof}
























\section{Cohomologie sur les espaces séparés et localement quasi-compacts}
\label{sites-cohomology-section-cohomology-LC}

\noindent
Nous continuons à suivre notre convention consistant à dire
«~séparé et localement quasi-compact~» plutôt que «~localement compact~»,
comme on le fait souvent dans la littérature. Notons $\textit{LC}$ la
catégorie dont les objets sont les espaces topologiques séparés et localement
quasi-compacts et dont les morphismes sont les applications continues.

\begin{lemma}
\label{sites-cohomology-lemma-LC-basic}
La catégorie $\textit{LC}$ admet des produits fibrés et un objet final, et
admet donc toutes les limites finies. Étant donnés des morphismes $X \to Z$
et $Y \to Z$ dans $\textit{LC}$ tels que $X$ et $Y$ soient quasi-compacts,
$X \times_Z Y$ est quasi-compact.
\end{lemma}

\begin{proof}
L'objet final est l'espace réduit à un point. Étant donnés des morphismes
$X \to Z$ et $Y \to Z$ de $\textit{LC}$, le produit fibré $X \times_Z Y$
est un sous-espace de $X \times Y$. Ainsi, $X \times_Z Y$ est séparé puisque
$X \times Y$ est séparé d'après
Topologie, section \ref{topology-section-Hausdorff}.

\medskip\noindent
Si $X$ et $Y$ sont quasi-compacts, alors $X \times Y$ est quasi-compact
d'après Topologie, théorème \ref{topology-theorem-tychonov}.
Comme $X \times_Z Y$ est une partie fermée de $X \times Y$
(Topologie, lemme \ref{topology-lemma-fibre-product-closed}),
on voit que $X \times_Z Y$ est quasi-compact d'après
Topologie, lemme \ref{topology-lemma-closed-in-quasi-compact}.

\medskip\noindent
Enfin, revenons au cas général. Si $x \in X$ et $y \in Y$, on peut choisir
des voisinages quasi-compacts $x \in E \subset X$ et $y \in F \subset Y$ ;
le résultat précédent montre alors que $E \times_Z F$ est un voisinage
quasi-compact de $(x, y)$. Ainsi, $X \times_Z Y$ est un objet de
$\textit{LC}$ d'après
Topologie, lemme \ref{topology-lemma-locally-quasi-compact-Hausdorff}.
\end{proof}

\noindent
On peut munir $\textit{LC}$ d'une topologie plus fine que la topologie usuelle.

\begin{definition}
\label{sites-cohomology-definition-covering-LC}
Soit $\{f_i : X_i \to X\}$ une famille de morphismes de but fixé dans la
catégorie $\textit{LC}$. On dit que cette famille est un
{\it recouvrement qc}\footnote{Cette notation n'est pas standard.
Nous l'avons choisie pour rappeler au lecteur les recouvrements fpqc
de schémas.}
si, pour tout $x \in X$, il existe $i_1, \ldots, i_n \in I$ et des parties
quasi-compactes $E_j \subset X_{i_j}$ telles que
$\bigcup f_{i_j}(E_j)$ soit un voisinage de $x$.
\end{definition}

\noindent
Remarquons qu'un recouvrement ouvert $X = \bigcup U_i$ d'un objet de
$\textit{LC}$ fournit un recouvrement qc $\{U_i \to X\}$, puisque $X$ est
localement quasi-compact. Commençons par le lemme obligé.

\begin{lemma}
\label{sites-cohomology-lemma-qc}
Soit $X$ un espace séparé et localement quasi-compact, autrement dit un
objet de $\textit{LC}$.
\begin{enumerate}
\item Si $X' \to X$ est un isomorphisme dans $\textit{LC}$, alors
$\{X' \to X\}$ est un recouvrement qc.
\item Si $\{f_i : X_i \to X\}_{i\in I}$ est un recouvrement qc et si, pour
chaque $i$, on a un recouvrement qc
$\{g_{ij} : X_{ij} \to X_i\}_{j\in J_i}$, alors
$\{X_{ij} \to X\}_{i \in I, j\in J_i}$ est un recouvrement qc.
\item Si $\{X_i \to X\}_{i\in I}$ est un recouvrement qc et si
$X' \to X$ est un morphisme de $\textit{LC}$, alors
$\{X' \times_X X_i \to X'\}_{i\in I}$ est un recouvrement qc.
\end{enumerate}
\end{lemma}

\begin{proof}
Le point (1) résulte de la remarque précédente selon laquelle les
recouvrements ouverts sont des recouvrements qc.

\medskip\noindent
Démonstration de (2). Soit $x \in X$. Choisissons
$i_1, \ldots, i_n \in I$ et des parties quasi-compactes
$E_a \subset X_{i_a}$ telles que $\bigcup f_{i_a}(E_a)$ soit un voisinage
de $x$. Pour tout $e \in E_a$, on peut trouver une partie finie
$J_e \subset J_{i_a}$ et des parties quasi-compactes
$F_{e, j} \subset X_{i_a j}$, $j \in J_e$, telles que
$\bigcup g_{i_a j}(F_{e, j})$ soit un voisinage de $e$. Comme $E_a$ est
quasi-compact, on trouve une famille finie
$e_1, \ldots, e_{m_a}$ telle que
$$
E_a \subset
\bigcup\nolimits_{k = 1, \ldots, m_a}
\bigcup\nolimits_{j \in J_{e_k}} g_{i_a j}(F_{e_k, j})
$$
On voit alors que
$$
\bigcup\nolimits_{a = 1, \ldots, n}
\bigcup\nolimits_{k = 1, \ldots, m_a}
\bigcup\nolimits_{j \in J_{e_k}} f_{i_a}(g_{i_a j}(F_{e_k, j}))
$$
est un voisinage de $x$.

\medskip\noindent
Démonstration de (3). Soit $x' \in X'$ un point, et soit $x \in X$ son
image. Choisissons $i_1, \ldots, i_n \in I$ et des parties quasi-compactes
$E_j \subset X_{i_j}$ telles que $\bigcup f_{i_j}(E_j)$ soit un voisinage
de $x$. Choisissons un voisinage quasi-compact $F \subset X'$ de $x'$ dont
l'image soit contenue dans le voisinage quasi-compact
$\bigcup f_{i_j}(E_j)$ de $x$. Alors
$F \times_X E_j \subset X' \times_X X_{i_j}$ est une partie quasi-compacte,
et $F$ est l'image de l'application
$\coprod F \times_X E_j \to F$. Le changement de base est donc un
recouvrement qc, ce qui achève la démonstration.
\end{proof}

\noindent
Puisque tous les objets de $\textit{LC}$ sont séparés, tout morphisme
$f : X \to Y$ de $\textit{LC}$ est une application continue séparée
d'espaces topologiques. Ainsi, $f$ est une application propre d'espaces
topologiques si et seulement si $f$ est universellement fermée. Voir la
discussion dans Topologie, section \ref{topology-section-proper}.

\begin{lemma}
\label{sites-cohomology-lemma-proper-surjective-is-qc-covering}
Soit $f : X \to Y$ un morphisme de $\textit{LC}$. Si $f$ est propre et
surjectif, alors $\{f : X \to Y\}$ est un recouvrement qc.
\end{lemma}

\begin{proof}
Soit $y \in Y$ un point. Pour tout $x \in X_y$, choisissons un voisinage
quasi-compact $E_x \subset X$, puis un ouvert $x \in U_x \subset E_x$.
Comme $f$ est propre, la fibre $X_y$ est quasi-compacte ; on trouve donc
$x_1, \ldots, x_n \in X_y$ tels que
$X_y \subset U_{x_1} \cup \ldots \cup U_{x_n}$.
Nous affirmons que $f(E_{x_1}) \cup \ldots \cup f(E_{x_n})$ est un voisinage
de $y$. En effet, comme $f$ est fermée
(Topologie, théorème \ref{topology-theorem-characterize-proper}),
on voit que
$Z = f(X \setminus (U_{x_1} \cup \ldots \cup U_{x_n}))$
est une partie fermée de $Y$ qui ne contient pas $y$. Comme $f$ est
surjective, on voit que $Y \setminus Z$ est contenu dans
$f(E_{x_1}) \cup \ldots \cup f(E_{x_n})$, comme voulu.
\end{proof}

\noindent
Mis à part quelques questions ensemblistes, le lemme \ref{sites-cohomology-lemma-qc} montre
que $\textit{LC}$, munie de la collection des recouvrements qc, forme un
site. Nous noterons $\textit{LC}_{qc}$ ce site (convenablement modifié pour
résoudre les questions ensemblistes).

\begin{remark}[Questions ensemblistes]
\label{sites-cohomology-remark-set-theoretic-LC}
La catégorie $\textit{LC}$ est une «~grande~» catégorie, car ses objets
forment une classe propre. De même, les recouvrements forment une classe
propre. Définissons la {\it taille} d'un espace topologique $X$ comme le
cardinal de l'ensemble des points de $X$. Choisissons une fonction $Bound$ sur
les cardinaux, par exemple comme dans
Ensembles, équation (\ref{sets-equation-bound}).
Enfin, soit $S_0$ un ensemble initial d'objets de $\textit{LC}$, par exemple
$S_0 = \{(\mathbf{R}, \text{topologie euclidienne})\}$.
Exactement comme dans Ensembles, lemme \ref{sets-lemma-construct-category},
on peut choisir un ordinal limite $\alpha$ tel que
$\textit{LC}_\alpha = \textit{LC} \cap V_\alpha$ contienne $S_0$ et soit
stable par toutes les limites et colimites dénombrables qui existent dans
$\textit{LC}$. De plus, si $X \in \textit{LC}_\alpha$, si
$Y \in \textit{LC}$ et si
$\text{size}(Y) \leq Bound(\text{size}(X))$, alors $Y$ est isomorphe à un
objet de $\textit{LC}_\alpha$.
Appliquons ensuite Ensembles, lemme \ref{sets-lemma-coverings-site}, pour
choisir un ensemble $\text{Cov}$ de recouvrements qc sur
$\textit{LC}_\alpha$ tel que tout recouvrement qc dans
$\textit{LC}_\alpha$ soit combinatoirement équivalent à un recouvrement
appartenant à cet ensemble. On obtient ainsi un site
$(\textit{LC}_\alpha, \text{Cov})$, que l'on notera $\textit{LC}_{qc}$.
\end{remark}

\noindent
Il existe une seconde topologie sur le site $\textit{LC}_{qc}$ de la
remarque \ref{sites-cohomology-remark-set-theoretic-LC}. Pour un objet $X$, on considère tous
les recouvrements $\{X_i \to X\}$ de $\textit{LC}_{qc}$ tels que
$X_i \to X$ soit une immersion ouverte. Nous notons ce site
$\textit{LC}_{Zar}$. Le foncteur identité
$\textit{LC}_{Zar} \to \textit{LC}_{qc}$ est continu et définit un morphisme
de sites
$$
\epsilon : \textit{LC}_{qc} \longrightarrow \textit{LC}_{Zar}
$$
Voir la section \ref{sites-cohomology-section-compare}.
Pour un espace topologique séparé et localement quasi-compact $X$, plus
précisément pour $X \in \Ob(\textit{LC}_{qc})$, nous notons le morphisme
induit
$$
\epsilon_X : \textit{LC}_{qc}/X \longrightarrow \textit{LC}_{Zar}/X
$$
(voir Sites, lemme \ref{sites-lemma-localize-morphism}).
Soit $X_{Zar}$ le site dont les objets sont les ouverts de $X$, voir
Sites, exemple \ref{sites-example-site-topological}.
Il existe un morphisme de sites
$$
\pi_X : \textit{LC}_{Zar}/X \longrightarrow X_{Zar}
$$
défini par le foncteur continu
$X_{Zar} \to \textit{LC}_{Zar}/X$, $U \mapsto U$.
En effet, $X_{Zar}$ possède des produits fibrés et un objet final, le
foncteur ci-dessus y commute, et Sites, proposition
\ref{sites-proposition-get-morphism}, s'applique.
Nous considérons souvent $\pi$ comme un morphisme de topos
$$
\pi_X : \Sh(\textit{LC}_{Zar}/X) \longrightarrow \Sh(X)
$$
en utilisant l'égalité $\Sh(X_{Zar}) = \Sh(X)$.

\begin{lemma}
\label{sites-cohomology-lemma-describe-pullback-pi}
Soit $X$ un objet de $\textit{LC}_{qc}$ et soit $\mathcal{F}$ un faisceau
sur $X$. La règle
$$
\textit{LC}_{qc}/X \longrightarrow \textit{Ensembles},\quad
(f : Y \to X) \longmapsto \Gamma(Y, f^{-1}\mathcal{F})
$$
est un faisceau et, a fortiori, aussi un faisceau sur
$\textit{LC}_{Zar}/X$. Ce faisceau est égal à
$\pi_X^{-1}\mathcal{F}$ sur $\textit{LC}_{Zar}/X$ et à
$\epsilon_X^{-1}\pi_X^{-1}\mathcal{F}$ sur $\textit{LC}_{qc}/X$.
\end{lemma}

\begin{proof}
Notons $\mathcal{G}$ le préfaisceau donné par la formule du lemme.
Bien entendu, l'image inverse $f^{-1}$ figurant dans la formule désigne
l'image inverse usuelle des faisceaux sur les espaces topologiques.
Il découle immédiatement des définitions que $\mathcal{G}$ est un faisceau
pour la topologie Zar.

\medskip\noindent
Soit $Y \to X$ un morphisme de $\textit{LC}_{qc}$ et soit
$\mathcal{V} = \{g_i : Y_i \to Y\}_{i \in I}$ un recouvrement qc.
Pour montrer que $\mathcal{G}$ est un faisceau pour la topologie qc, il
suffit de montrer que
$\mathcal{G}(Y) \to H^0(\mathcal{V}, \mathcal{G})$
est un isomorphisme, voir Sites, section
\ref{sites-section-sheafification}. Remarquons d'abord que le morphisme est
injectif, puisqu'un recouvrement qc est surjectif et que l'égalité des
sections se détecte sur les germes (utiliser Faisceaux, lemmes
\ref{sheaves-lemma-sheaf-subset-stalks} et
\ref{sheaves-lemma-stalk-pullback-presheaf}). Ainsi, $\mathcal{G}$ est un
préfaisceau séparé sur $\textit{LC}_{qc}$ ; il suffit donc de montrer que
tout élément
$(s_i) \in H^0(\mathcal{V}, \mathcal{G})$
s'envoie sur un élément de l'image de $\mathcal{G}(Y)$ après remplacement de
$\mathcal{V}$ par un raffinement
(Sites, théorème \ref{sites-theorem-plus}).

\medskip\noindent
En identifiant les faisceaux sur $Y_{i, Zar}$ aux faisceaux sur $Y_i$, nous
voyons que $\mathcal{G}|_{Y_{i, Zar}}$ est l'image inverse de
$f^{-1}\mathcal{F}$ par l'application continue $g_i : Y_i \to Y$. Nous
pouvons donc choisir un recouvrement ouvert $Y_i = \bigcup V_{ij}$ tel que,
pour chaque $j$, il existe un ouvert $W_{ij} \subset Y$ et une section
$t_{ij} \in \mathcal{G}(W_{ij})$, avec $V_{ij}$ envoyé dans $W_{ij}$ et
$s|_{V_{ij}}$ image inverse de $t_{ij}$. Autrement dit, après avoir raffiné
le recouvrement $\{Y_i \to Y\}$, nous pouvons supposer qu'il existe des
ouverts $W_i \subset Y$ tels que $Y_i \to Y$ se factorise par $W_i$, et des
sections $t_i$ de $\mathcal{G}$ sur $W_i$ dont les restrictions sont les
sections données $s_i$. De plus, si $y \in Y$ appartient aux images de
$Y_i \to Y$ et de $Y_j \to Y$, alors les images $t_{i, y}$ et $t_{j, y}$
dans le germe $f^{-1}\mathcal{F}_y$ sont égales (car $s_i$ et $s_j$ sont
égales sur $Y_i \times_Y Y_j$). Ainsi, pour $y \in Y$, il existe un élément
bien défini $t_y$ de $f^{-1}\mathcal{F}_y$, égal à $t_{i, y}$ dès que $y$
appartient à l'image de $Y_i \to Y$.
Nous allons montrer que l'élément $(t_y)$ provient d'une section globale de
$f^{-1}\mathcal{F}$ sur $Y$, ce qui achèvera la démonstration du lemme.

\medskip\noindent
Il suffit de le montrer localement sur $Y$, voir Faisceaux, section
\ref{sheaves-section-sheafification}. Soit $y_0 \in Y$. Choisissons
$i_1, \ldots, i_n \in I$ et des parties quasi-compactes
$E_j \subset Y_{i_j}$ telles que $\bigcup g_{i_j}(E_j)$ soit un voisinage de
$y_0$. Soit $V \subset Y$ un voisinage ouvert de $y_0$, contenu dans
$\bigcup g_{i_j}(E_j)$ et dans $W_{i_1} \cap \ldots \cap W_{i_n}$.
Comme $t_{i_1, y_0} = \ldots = t_{i_n, y_0}$, après avoir rétréci $V$, nous
pouvons supposer que les sections $t_{i_j}|_V$, $j = 1, \ldots, n$, de
$f^{-1}\mathcal{F}$ sont égales. Puisque
$V \subset \bigcup g_{i_j}(E_j)$, nous voyons que $(t_y)_{y \in V}$ provient
de cette section.

\medskip\noindent
Il reste à montrer que $\mathcal{G}$ est égal à
$\epsilon_X^{-1}\pi_X^{-1}\mathcal{F}$ sur $\textit{LC}_{qc}$,
resp.\ à $\pi_X^{-1}\mathcal{F}$ sur $\textit{LC}_{Zar}$.
Dans les deux cas, l'image inverse est définie en prenant le préfaisceau
$$
(f : Y \to X)
\longmapsto
\colim_{f(Y) \subset U \subset X} \mathcal{F}(U)
$$
puis en le faisceautisant. La faisceautisation pour la topologie Zar produit
exactement notre faisceau $\mathcal{G}$ ; le fait que $\mathcal{G}$ soit un
faisceau qc montre que cela fonctionne aussi pour la topologie qc.
\end{proof}

\noindent
Soit $X \in \Ob(\textit{LC}_{Zar})$ et soit $\mathcal{H}$ un faisceau
abélien sur $\textit{LC}_{Zar}/X$. Nous noterons
$H^n_{Zar}(U, \mathcal{H})$ la cohomologie de $\mathcal{H}$ au-dessus d'un
objet $U$ de $\textit{LC}_{Zar}/X$.

\begin{lemma}
\label{sites-cohomology-lemma-collect-true-things-Zar}
Soit $X$ un objet de $\textit{LC}_{Zar}$. Alors :
\begin{enumerate}
\item pour $\mathcal{F} \in \textit{Ab}(X)$, nous avons
$H^n_{Zar}(X, \pi_X^{-1}\mathcal{F}) = H^n(X, \mathcal{F})$,
\item $\pi_{X, *} : \textit{Ab}(\textit{LC}_{Zar}/X) \to \textit{Ab}(X)$
est exact,
\item l'unité $\text{id} \to \pi_{X, *} \circ \pi_X^{-1}$ de l'adjonction
est un isomorphisme, et
\item pour $K \in D(X)$, le morphisme canonique
$K \to R\pi_{X, *} \pi_X^{-1}K$ est un isomorphisme.
\end{enumerate}
Soit $f : X \to Y$ un morphisme de $\textit{LC}_{Zar}$. Alors :
\begin{enumerate}
\item[(5)] il existe un diagramme commutatif
$$
\xymatrix{
\Sh(\textit{LC}_{Zar}/X) \ar[r]_{f_{Zar}} \ar[d]_{\pi_X} &
\Sh(\textit{LC}_{Zar}/Y) \ar[d]^{\pi_Y} \\
\Sh(X_{Zar}) \ar[r]^f &
\Sh(Y_{Zar})
}
$$
de topos,
\item[(6)] pour $L \in D^+(Y)$, nous avons
$H^n_{Zar}(X, \pi_Y^{-1}L) = H^n(X, f^{-1}L)$,
\item[(7)] si $f$ est propre, alors :
\begin{enumerate}
\item $\pi_Y^{-1} \circ f_* = f_{Zar, *} \circ \pi_X^{-1}$ comme foncteurs
$\Sh(X) \to \Sh(\textit{LC}_{Zar}/Y)$,
\item $\pi_Y^{-1} \circ Rf_* = Rf_{Zar, *} \circ \pi_X^{-1}$ comme
foncteurs $D^+(X) \to D^+(\textit{LC}_{Zar}/Y)$.
\end{enumerate}
\end{enumerate}
\end{lemma}

\begin{proof}
Démonstration de (1).
L'égalité $H^n_{Zar}(X, \pi_X^{-1}\mathcal{F}) = H^n(X, \mathcal{F})$
est un fait général provenant de l'observation immédiate que les recouvrements
de $X$ dans $\textit{LC}_{Zar}$ sont exactement les recouvrements ouverts de
$X$. Pour une démonstration détaillée, on peut appliquer le
lemme \ref{sites-cohomology-lemma-cohomology-bigger-site} au foncteur
$X_{Zar} \to \textit{LC}_{Zar}$.

\medskip\noindent
Démonstration de (2). Cela est vrai parce que
$\pi_{X, *} = \tau_X^{-1}$ pour un certain morphisme de topos
$\tau_X : \Sh(X_{Zar}) \to \Sh(\textit{LC}_{Zar})$, comme il résulte de
Sites, lemme \ref{sites-lemma-bigger-site}, appliqué au foncteur
$X_{Zar} \to \textit{LC}_{Zar}/X$ utilisé pour définir $\pi_X$.

\medskip\noindent
Démonstration de (3). Cela est vrai parce que
$\tau_X^{-1} \circ \pi_X^{-1}$ est le foncteur identité d'après Sites,
lemme \ref{sites-lemma-bigger-site}. On peut aussi le déduire de la
description explicite de $\pi_X^{-1}$ dans le
lemme \ref{sites-cohomology-lemma-describe-pullback-pi}.

\medskip\noindent
Démonstration de (4). Appliquer (3) à un complexe de faisceaux abéliens
représentant $K$.

\medskip\noindent
Démonstration de (5). Le morphisme de topos $f_{Zar}$ provient d'une
application de Sites, lemme \ref{sites-lemma-relocalize}, et, dans notre cas,
du foncteur continu
$Z/Y \mapsto Z \times_Y X/X$ d'après Sites, lemme
\ref{sites-lemma-relocalize-given-fibre-products}.
Le diagramme commute simplement parce que les foncteurs continus
correspondants se composent correctement
(voir Sites, lemme \ref{sites-lemma-composition-morphisms-sites}).

\medskip\noindent
Démonstration de (6). Nous avons
$H^n_{Zar}(X, \pi_Y^{-1}\mathcal{G}) =
H^n_{Zar}(X, f_{Zar}^{-1}\pi_Y^{-1}\mathcal{G})$
pour $\mathcal{G}$ dans $\textit{Ab}(Y)$, voir le
lemme \ref{sites-cohomology-lemma-cohomology-of-open}. Ceci est égal à
$H^n_{Zar}(X, \pi_X^{-1}f^{-1}\mathcal{G})$ par commutativité du diagramme de
(5). Nous concluons donc par (1) lorsque $L$ est réduit à un seul faisceau en
degré $0$. Le cas général s'obtient en représentant $L$ par un complexe de
faisceaux abéliens borné inférieurement.

\medskip\noindent
Démonstration de (7a). Soit $\mathcal{F}$ un faisceau sur $X$.
Soit $g : Z \to Y$ un objet de $\textit{LC}_{Zar}/Y$. Considérons le produit
fibré
$$
\xymatrix{
Z' \ar[r]_{f'} \ar[d]_{g'} & Z \ar[d]^g \\
X \ar[r]^f & Y
}
$$
Nous avons alors
$$
(f_{Zar, *}\pi_X^{-1}\mathcal{F})(Z/Y) =
(\pi_X^{-1}\mathcal{F})(Z'/X)  =
\Gamma(Z', (g')^{-1}\mathcal{F})  =
\Gamma(Z, f'_*(g')^{-1}\mathcal{F})
$$
la deuxième égalité venant du lemme \ref{sites-cohomology-lemma-describe-pullback-pi}.
D'autre part,
$$
(\pi_Y^{-1}f_*\mathcal{F})(Z/Y) = \Gamma(Z, g^{-1}f_*\mathcal{F})
$$
encore d'après le lemme \ref{sites-cohomology-lemma-describe-pullback-pi}.
Le changement de base propre pour les faisceaux d'ensembles
(Cohomologie, lemme
\ref{cohomology-lemma-proper-base-change-sheaves-of-sets}) montre donc que
les deux ensembles sont canoniquement isomorphes. Cet isomorphisme est
compatible aux applications de restriction et définit un isomorphisme
$\pi_Y^{-1}f_*\mathcal{F} = f_{Zar, *}\pi_X^{-1}\mathcal{F}$.
et donc un isomorphisme de foncteurs
$\pi_Y^{-1} \circ f_* = f_{Zar, *} \circ \pi_X^{-1}$.

\medskip\noindent
Démonstration de (7b). Soit $K \in D^+(X)$. D'après le
lemme \ref{sites-cohomology-lemma-unbounded-describe-higher-direct-images}, le $n$-ième
faisceau de cohomologie de $Rf_{Zar, *}\pi_X^{-1}K$ est le faisceau associé
au préfaisceau
$$
(g : Z \to Y) \longmapsto H^n_{Zar}(Z', \pi_X^{-1}K)
$$
avec les notations ci-dessus. Observons que
\begin{align*}
H^n_{Zar}(Z', \pi_X^{-1}K)
& =
H^n(Z', (g')^{-1}K) \\
& =
H^n(Z, Rf'_*(g')^{-1}K) \\
& =
H^n(Z, g^{-1}Rf_*K) \\
& =
H^n_{Zar}(Z, \pi_Y^{-1}Rf_*K)
\end{align*}
La première égalité est (6) appliqué à $K$ et à $g' : Z' \to X$.
La deuxième est Leray pour $f' : Z' \to Z$
(Cohomologie, lemme \ref{cohomology-lemma-before-Leray}).
La troisième est le théorème de changement de base propre
(Cohomologie, théorème \ref{cohomology-theorem-proper-base-change}).
La quatrième est (6) appliqué à $g : Z \to Y$ et à $Rf_*K$.
Ainsi, $Rf_{Zar, *}\pi_X^{-1}K$ et $\pi_Y^{-1}Rf_*K$ ont les mêmes faisceaux
de cohomologie. Nous omettons la vérification que le morphisme canonique de
changement de base $\pi_Y^{-1}Rf_*K \to Rf_{Zar, *}\pi_X^{-1}K$ induit cet
isomorphisme.
\end{proof}

\noindent
Dans la situation du lemme \ref{sites-cohomology-lemma-describe-pullback-pi}, la composition
de $\epsilon$ et $\pi$ et l'égalité $\Sh(X) = \Sh(X_{Zar})$ déterminent un
morphisme de topos
$$
a_X : \Sh(\textit{LC}_{qc}/X) \longrightarrow \Sh(X)
$$

\begin{lemma}
\label{sites-cohomology-lemma-push-pull-LC}
Soit $f : X \to Y$ un morphisme de $\textit{LC}_{qc}$. Il existe alors des
diagrammes commutatifs de topos
$$
\vcenter{
\xymatrix{
\Sh(\textit{LC}_{qc}/X) \ar[r]_{f_{qc}} \ar[d]_{\epsilon_X} &
\Sh(\textit{LC}_{qc}/Y) \ar[d]^{\epsilon_Y} \\
\Sh(\textit{LC}_{Zar}/X) \ar[r]^{f_{Zar}} &
\Sh(\textit{LC}_{Zar}/Y)
}
}
\quad\text{et}\quad
\vcenter{
\xymatrix{
\Sh(\textit{LC}_{qc}/X) \ar[r]_{f_{qc}} \ar[d]_{a_X} &
\Sh(\textit{LC}_{qc}/Y) \ar[d]^{a_Y} \\
\Sh(X) \ar[r]^f &
\Sh(Y)
}
}
$$
avec $a_X = \pi_X \circ \epsilon_X$, $a_Y = \pi_Y \circ \epsilon_Y$.
Si $f$ est propre, alors $a_Y^{-1} \circ f_* = f_{qc, *} \circ a_X^{-1}$.
\end{lemma}

\begin{proof}
Le morphisme de topos $f_{qc}$ est celui de Sites, lemme
\ref{sites-lemma-relocalize}, qui, dans notre cas, provient du foncteur
continu $Z/Y \mapsto Z \times_Y X/X$, voir Sites, lemme
\ref{sites-lemma-relocalize-given-fibre-products}.
Le diagramme de gauche commute parce que les foncteurs continus
correspondants se composent correctement
(voir Sites, lemme \ref{sites-lemma-composition-morphisms-sites}).
Le diagramme de droite commute parce que celui de gauche commute et à cause
du point (5) du lemme \ref{sites-cohomology-lemma-collect-true-things-Zar}.

\medskip\noindent
Démonstration de la dernière assertion. On pourrait répéter la démonstration
du point (7a) du lemme \ref{sites-cohomology-lemma-collect-true-things-Zar} ; nous allons
plutôt l'en déduire. Comme $\epsilon_{Y, *}$ est le foncteur identité sur les
préfaisceaux sous-jacents, il reflète les isomorphismes. La description du
lemme \ref{sites-cohomology-lemma-describe-pullback-pi} montre que
$\epsilon_{Y, *} \circ a_Y^{-1} = \pi_Y^{-1}$, et de même pour $X$.
Pour montrer que le morphisme canonique
$a_Y^{-1}f_*\mathcal{F} \to f_{qc, *}a_X^{-1}\mathcal{F}$
est un isomorphisme, il suffit de montrer que
$$
\pi_Y^{-1}f_*\mathcal{F} =
\epsilon_{Y, *}a_Y^{-1}f_*\mathcal{F} \to
\epsilon_{Y, *}f_{qc, *}a_X^{-1}\mathcal{F} =
f_{Zar, *}\epsilon_{X, *}a_X^{-1}\mathcal{F} =
f_{Zar, *}\pi_X^{-1}\mathcal{F}
$$
est un isomorphisme. C'est le point (7a) du
lemme \ref{sites-cohomology-lemma-collect-true-things-Zar}.
\end{proof}

\begin{lemma}
\label{sites-cohomology-lemma-compare-qc-zar}
Considérons le morphisme de comparaison
$\epsilon : \textit{LC}_{qc} \to \textit{LC}_{Zar}$.
Notons $\mathcal{P}$ la classe des applications propres d'espaces
topologiques. Pour $X$ dans $\textit{LC}_{Zar}$, notons
$\mathcal{A}'_X \subset \textit{Ab}(\textit{LC}_{Zar}/X)$
la sous-catégorie pleine formée des faisceaux de la forme
$\pi_X^{-1}\mathcal{F}$ avec $\mathcal{F}$ dans $\textit{Ab}(X)$.
Alors les propriétés
(\ref{sites-cohomology-item-base-change-P}),
(\ref{sites-cohomology-item-restriction-A}),
(\ref{sites-cohomology-item-A-sheaf}),
(\ref{sites-cohomology-item-A-and-P}) et
(\ref{sites-cohomology-item-refine-tau-by-P})
de la situation \ref{sites-cohomology-situation-compare} sont vérifiées.
\end{lemma}

\begin{proof}
Nous montrons d'abord que
$\mathcal{A}'_X \subset \textit{Ab}(\textit{LC}_{Zar}/X)$ est une
sous-catégorie de Serre faible en vérifiant les conditions (1), (2), (3) et
(4) de Homologie, lemme
\ref{homology-lemma-characterize-weak-serre-subcategory}.
Les points (1), (2) et (3) sont immédiats, car $\pi_X^{-1}$ est exact et
pleinement fidèle d'après le point (3) du
lemme \ref{sites-cohomology-lemma-collect-true-things-Zar}. Si
$0 \to \pi_X^{-1}\mathcal{F} \to \mathcal{G} \to \pi_X^{-1}\mathcal{F}' \to 0$
est une suite exacte courte dans $\textit{Ab}(\textit{LC}_{Zar}/X)$, alors
$0 \to \mathcal{F} \to \pi_{X, *}\mathcal{G} \to \mathcal{F}' \to 0$ est
exacte d'après le point (2) du lemme
\ref{sites-cohomology-lemma-collect-true-things-Zar}. Par conséquent,
$\mathcal{G} = \pi_X^{-1}\pi_{X, *}\mathcal{G}$ appartient à
$\mathcal{A}'_X$, ce qui vérifie la dernière condition.

\medskip\noindent
La propriété (\ref{sites-cohomology-item-base-change-P}) résulte du lemme
\ref{sites-cohomology-lemma-LC-basic} et du fait que le changement de base d'une application
propre est propre (voir Topologie, théorème
\ref{topology-theorem-characterize-proper} et lemme
\ref{topology-lemma-base-change-separated}).

\medskip\noindent
La propriété (\ref{sites-cohomology-item-restriction-A}) résulte du diagramme commutatif (5)
du lemme \ref{sites-cohomology-lemma-collect-true-things-Zar}.

\medskip\noindent
La propriété (\ref{sites-cohomology-item-A-sheaf}) est le
lemme \ref{sites-cohomology-lemma-describe-pullback-pi}.

\medskip\noindent
La propriété (\ref{sites-cohomology-item-A-and-P}) est le point (7)(b) du
lemme \ref{sites-cohomology-lemma-collect-true-things-Zar}.

\medskip\noindent
Démonstration de (\ref{sites-cohomology-item-refine-tau-by-P}). Soit donné un recouvrement qc
$\{U_i \to U\}$. Pour $u \in U$, choisissons $i_1, \ldots, i_m \in I$ et
des parties quasi-compactes $E_j \subset U_{i_j}$ telles que
$\bigcup f_{i_j}(E_j)$ soit un voisinage de $u$.
Observons que $Y = \coprod_{j = 1, \ldots, m} E_j \to U$ est propre comme
application continue entre espaces séparés quasi-compacts
(Topologie, lemme \ref{topology-lemma-closed-map}).
Choisissons un voisinage ouvert $u \in V$ contenu dans
$\bigcup f_{i_j}(E_j)$. Alors $Y \times_U V \to V$ est un morphisme propre
surjectif, donc un recouvrement $qc$ d'après le
lemme \ref{sites-cohomology-lemma-proper-surjective-is-qc-covering}.
Comme nous pouvons le faire pour tout $u \in U$, la propriété
(\ref{sites-cohomology-item-refine-tau-by-P}) est vérifiée.
\end{proof}

\begin{lemma}
\label{sites-cohomology-lemma-V-C-all-n}
Avec les notations ci-dessus :
\begin{enumerate}
\item Pour $X \in \Ob(\textit{LC}_{qc})$ et un faisceau abélien
$\mathcal{F}$ sur $X$, nous avons
$\epsilon_{X, *}a_X^{-1}\mathcal{F} = \pi_X^{-1}\mathcal{F}$ et
$R^i\epsilon_{X, *}(a_X^{-1}\mathcal{F}) = 0$ pour $i > 0$.
\item Pour un morphisme propre $f : X \to Y$ de $\textit{LC}_{qc}$ et un
faisceau abélien $\mathcal{F}$ sur $X$, nous avons
$a_Y^{-1}(R^if_*\mathcal{F}) = R^if_{qc, *}(a_X^{-1}\mathcal{F})$
pour tout $i$.
\item Pour $X \in \Ob(\textit{LC}_{qc})$ et $K$ dans $D^+(X)$, le
morphisme $\pi_X^{-1}K \to R\epsilon_{X, *}(a_X^{-1}K)$ est un isomorphisme.
\item Pour un morphisme propre $f : X \to Y$ de $\textit{LC}_{qc}$ et $K$
dans $D^+(X)$, nous avons
$a_Y^{-1}(Rf_*K) = Rf_{qc, *}(a_X^{-1}K)$.
\end{enumerate}
\end{lemma}

\begin{proof}
Par le lemme \ref{sites-cohomology-lemma-compare-qc-zar}, tous les lemmes de la
section \ref{sites-cohomology-section-compare-general} s'appliquent dans notre situation.
Pour traduire les résultats, observons que la catégorie $\mathcal{A}_X$ du
lemme \ref{sites-cohomology-lemma-A} est l'image essentielle de
$a_X^{-1} : \textit{Ab}(X) \to \textit{Ab}(\textit{LC}_{qc}/X)$.

\medskip\noindent
Le point (1) équivaut à $(V_n)$ pour tout $n$, ce qui est vérifié d'après le
lemme \ref{sites-cohomology-lemma-V-C-all-n-general}.

\medskip\noindent
Le point (2) résulte de l'application de $\epsilon_Y^{-1}$ à la conclusion
du lemme \ref{sites-cohomology-lemma-V-implies-C-general}.

\medskip\noindent
Le point (3) résulte du point (1) du
lemme \ref{sites-cohomology-lemma-V-C-all-n-general}, car $\pi_X^{-1}K$ appartient à
$D^+_{\mathcal{A}'_X}(\textit{LC}_{Zar}/X)$ et
$a_X^{-1} = \epsilon_X^{-1} \circ \pi_X^{-1}$.

\medskip\noindent
Le point (4) résulte du point (2) du
lemme \ref{sites-cohomology-lemma-V-C-all-n-general} pour la même raison.
\end{proof}

\begin{lemma}
\label{sites-cohomology-lemma-cohomological-descent-LC}
Soit $X$ un objet de $\textit{LC}_{qc}$. Pour $K \in D^+(X)$, le morphisme
$$
K \longrightarrow Ra_{X, *}a_X^{-1}K
$$
est un isomorphisme, où
$a_X : \Sh(\textit{LC}_{qc}/X) \to \Sh(X)$ est comme ci-dessus.
\end{lemma}

\begin{proof}
Nous nous ramenons d'abord au cas où $K$ est donné par un seul faisceau
abélien. Représentons en effet $K$ par un complexe borné inférieurement
$\mathcal{F}^\bullet$. Le cas d'un faisceau montre que
$\mathcal{F}^n = a_{X, *} a_X^{-1} \mathcal{F}^n$
et que les faisceaux $R^qa_{X, *}a_X^{-1}\mathcal{F}^n$ sont nuls pour
$q > 0$. En appliquant le lemme d'acyclicité de Leray
(Catégories dérivées, lemme \ref{derived-lemma-leray-acyclicity}) à
$a_X^{-1}\mathcal{F}^\bullet$ et au foncteur $a_{X, *}$, nous concluons.
Désormais, supposons $K = \mathcal{F}$.

\medskip\noindent
Par le lemme \ref{sites-cohomology-lemma-describe-pullback-pi}, nous avons
$a_{X, *}a_X^{-1}\mathcal{F} = \mathcal{F}$. Il suffit donc de montrer que
$R^qa_{X, *}a_X^{-1}\mathcal{F} = 0$ pour $q > 0$.
Pour cela, nous pouvons utiliser $a_X = \pi_X \circ \epsilon_X$ et la suite
spectrale de Leray du lemme \ref{sites-cohomology-lemma-relative-Leray}.
Par le lemme \ref{sites-cohomology-lemma-V-C-all-n}, nous avons
$R^i\epsilon_{X, *}(a_X^{-1}\mathcal{F}) = 0$ pour $i > 0$ et
$\epsilon_{X, *}a_X^{-1}\mathcal{F} = \pi_X^{-1}\mathcal{F}$.
Par le lemme \ref{sites-cohomology-lemma-collect-true-things-Zar}, nous avons
$R^j\pi_{X, *}(\pi_X^{-1}\mathcal{F}) = 0$ pour $j > 0$.
Cela achève la démonstration.
\end{proof}

\begin{lemma}
\label{sites-cohomology-lemma-compare-cohomology-LC}
Avec $X \in \Ob(\textit{LC}_{qc})$ et
$a_X : \Sh(\textit{LC}_{qc}/X) \to \Sh(X)$ comme ci-dessus :
\begin{enumerate}
\item pour un faisceau abélien $\mathcal{F}$ sur $X$, nous avons
$H^n(X, \mathcal{F}) = H^n_{qc}(X, a_X^{-1}\mathcal{F})$,
\item pour $K \in D^+(X)$, nous avons
$H^n(X, K) = H^n_{qc}(X, a_X^{-1}K)$.
\end{enumerate}
Par exemple, si $A$ est un groupe abélien, alors nous avons
$H^n(X, \underline{A}) = H^n_{qc}(X, \underline{A})$.
\end{lemma}

\begin{proof}
Cela résulte du lemme \ref{sites-cohomology-lemma-cohomological-descent-LC} et de la
remarque \ref{sites-cohomology-remark-before-Leray}.
\end{proof}











\section{Suites spectrales pour Ext}
\label{sites-cohomology-section-spectral-sequence-ext}

\noindent
Dans cette section, nous rassemblons diverses suites spectrales qui interviennent
lorsque l'on considère les foncteurs Ext. Pour toute paire de complexes
$\mathcal{G}^\bullet, \mathcal{F}^\bullet$ de modules sur un site annelé
$(\mathcal{C}, \mathcal{O})$, nous notons
$$
\Ext^n_\mathcal{O}(\mathcal{G}^\bullet, \mathcal{F}^\bullet)
=
\Hom_{D(\mathcal{O})}(\mathcal{G}^\bullet, \mathcal{F}^\bullet[n])
$$
conformément à nos conventions générales de
Catégories dérivées, section \ref{derived-section-ext}.

\begin{example}
\label{sites-cohomology-example-hom-complex-into-sheaf}
Soit $(\mathcal{C}, \mathcal{O})$ un site annelé.
Soit $\mathcal{K}^\bullet$ un complexe de $\mathcal{O}$-modules borné supérieurement.
Soit $\mathcal{F}$ un $\mathcal{O}$-module. Il existe alors une
suite spectrale dont la page $E_2$ est
$$
E_2^{i, j} =
\Ext_\mathcal{O}^i(H^{-j}(\mathcal{K}^\bullet), \mathcal{F})
\Rightarrow
\Ext_\mathcal{O}^{i + j}(\mathcal{K}^\bullet, \mathcal{F})
$$
et une autre suite spectrale dont la page $E_1$ est
$$
E_1^{i, j} =
\Ext_\mathcal{O}^j(\mathcal{K}^{-i}, \mathcal{F})
\Rightarrow
\Ext_\mathcal{O}^{i + j}(\mathcal{K}^\bullet, \mathcal{F}).
$$
Pour construire ces suites spectrales, choisissons une résolution injective
$\mathcal{F} \to \mathcal{I}^\bullet$ et considérons les deux suites
spectrales provenant du complexe double
$\Hom_\mathcal{O}(\mathcal{K}^\bullet, \mathcal{I}^\bullet)$ ; voir
Homologie, section \ref{homology-section-double-complex}.
\end{example}








\section{Cup-produit}
\label{sites-cohomology-section-cup-product}

\noindent
Soit $(\mathcal{C}, \mathcal{O})$ un site annelé. Soient
$K, M$ des objets de $D(\mathcal{O})$. Posons
$A = \Gamma(\mathcal{C}, \mathcal{O})$. Dans ce cadre, le cup-produit (global)
est un morphisme
$$
R\Gamma(\mathcal{C}, K) \otimes_A^\mathbf{L}
R\Gamma(\mathcal{C}, M)
\longrightarrow
R\Gamma(\mathcal{C}, K \otimes_\mathcal{O}^\mathbf{L} M)
$$
dans $D(A)$. Nous le définissons comme le cup-produit relatif associé
au morphisme de topos annelés
$(\Sh(\mathcal{C}), \mathcal{O}) \to (\Sh(pt), A)$
comme dans la remarque \ref{sites-cohomology-remark-cup-product}.

\medskip\noindent
Formulons et démontrons une compatibilité naturelle du
cup-produit relatif. Supposons donc que l'on dispose d'un morphisme
$f : (\Sh(\mathcal{C}), \mathcal{O}_\mathcal{C}) \to
(\Sh(\mathcal{D}), \mathcal{O}_\mathcal{D})$ de topos annelés.
Soient $\mathcal{K}^\bullet$ et $\mathcal{M}^\bullet$
des complexes de $\mathcal{O}_\mathcal{C}$-modules.
Il existe un cup-produit naïf
$$
\text{Tot}(
f_*\mathcal{K}^\bullet
\otimes_{\mathcal{O}_\mathcal{D}}
f_*\mathcal{M}^\bullet)
\longrightarrow
f_*\text{Tot}(\mathcal{K}^\bullet
\otimes_{\mathcal{O}_\mathcal{C}}
\mathcal{M}^\bullet)
$$
Nous affirmons que celui-ci est lié au cup-produit relatif.

\begin{lemma}
\label{sites-cohomology-lemma-cup-compatible-with-naive}
Dans la situation ci-dessus, le diagramme suivant est commutatif
$$
\xymatrix{
f_*\mathcal{K}^\bullet
\otimes_{\mathcal{O}_\mathcal{D}}^\mathbf{L}
f_*\mathcal{M}^\bullet \ar[r] \ar[d]
&
Rf_*\mathcal{K}^\bullet
\otimes_{\mathcal{O}_\mathcal{D}}^\mathbf{L}
Rf_*\mathcal{M}^\bullet \ar[d]^{\text{Remarque \ref{sites-cohomology-remark-cup-product}}} \\
\text{Tot}(
f_*\mathcal{K}^\bullet
\otimes_{\mathcal{O}_\mathcal{D}}
f_*\mathcal{M}^\bullet) \ar[d]_{\text{cup-produit naïf}} &
Rf_*(\mathcal{K}^\bullet
\otimes_{\mathcal{O}_\mathcal{C}}^\mathbf{L}
\mathcal{M}^\bullet) \ar[d] \\
f_*\text{Tot}(\mathcal{K}^\bullet
\otimes_{\mathcal{O}_\mathcal{C}}
\mathcal{M}^\bullet) \ar[r] &
Rf_*\text{Tot}(\mathcal{K}^\bullet
\otimes_{\mathcal{O}_\mathcal{C}}
\mathcal{M}^\bullet)
}
$$
\end{lemma}

\begin{proof}
D'après la construction de la remarque \ref{sites-cohomology-remark-cup-product}, on voit qu'en
parcourant le diagramme dans le sens des aiguilles d'une montre, le morphisme
$$
f_*\mathcal{K}^\bullet
\otimes_{\mathcal{O}_\mathcal{D}}^\mathbf{L}
f_*\mathcal{M}^\bullet 
\longrightarrow
Rf_*\text{Tot}(\mathcal{K}^\bullet
\otimes_{\mathcal{O}_\mathcal{C}}
\mathcal{M}^\bullet)
$$
est adjoint au morphisme
\begin{align*}
Lf^*(f_*\mathcal{K}^\bullet
\otimes_{\mathcal{O}_\mathcal{D}}^\mathbf{L}
f_*\mathcal{M}^\bullet)
& =
Lf^*f_*\mathcal{K}^\bullet
\otimes_{\mathcal{O}_\mathcal{C}}^\mathbf{L}
Lf^*f_*\mathcal{M}^\bullet \\
& \to
Lf^*Rf_*\mathcal{K}^\bullet
\otimes_{\mathcal{O}_\mathcal{C}}^\mathbf{L}
Lf^*Rf_*\mathcal{M}^\bullet \\
& \to
\mathcal{K}^\bullet
\otimes_{\mathcal{O}_\mathcal{C}}^\mathbf{L}
\mathcal{M}^\bullet \\
& \to
\text{Tot}(\mathcal{K}^\bullet
\otimes_{\mathcal{O}_\mathcal{C}}
\mathcal{M}^\bullet)
\end{align*}
D'après le lemme \ref{sites-cohomology-lemma-adjoints-push-pull-compatibility}, ce morphisme est aussi égal à
\begin{align*}
Lf^*(f_*\mathcal{K}^\bullet
\otimes_{\mathcal{O}_\mathcal{D}}^\mathbf{L}
f_*\mathcal{M}^\bullet)
& =
Lf^*f_*\mathcal{K}^\bullet
\otimes_{\mathcal{O}_\mathcal{C}}^\mathbf{L}
Lf^*f_*\mathcal{M}^\bullet \\
& \to
f^*f_*\mathcal{K}^\bullet
\otimes_{\mathcal{O}_\mathcal{C}}^\mathbf{L}
f^*f_*\mathcal{M}^\bullet \\
& \to
\mathcal{K}^\bullet
\otimes_{\mathcal{O}_\mathcal{C}}^\mathbf{L}
\mathcal{M}^\bullet \\
& \to
\text{Tot}(\mathcal{K}^\bullet
\otimes_{\mathcal{O}_\mathcal{C}}
\mathcal{M}^\bullet)
\end{align*}
En parcourant le diagramme dans le sens inverse des aiguilles d'une montre,
on obtient le morphisme adjoint au morphisme
\begin{align*}
Lf^*(f_*\mathcal{K}^\bullet
\otimes_{\mathcal{O}_\mathcal{D}}^\mathbf{L}
f_*\mathcal{M}^\bullet)
& \to
Lf^*\text{Tot}(
f_*\mathcal{K}^\bullet
\otimes_{\mathcal{O}_\mathcal{D}}
f_*\mathcal{M}^\bullet) \\
& \to
Lf^*f_*\text{Tot}(\mathcal{K}^\bullet
\otimes_{\mathcal{O}_\mathcal{C}}
\mathcal{M}^\bullet) \\
& \to
Lf^*Rf_*\text{Tot}(\mathcal{K}^\bullet
\otimes_{\mathcal{O}_\mathcal{C}}
\mathcal{M}^\bullet) \\
& \to
\text{Tot}(\mathcal{K}^\bullet
\otimes_{\mathcal{O}_\mathcal{C}}
\mathcal{M}^\bullet)
\end{align*}
D'après le lemme \ref{sites-cohomology-lemma-adjoints-push-pull-compatibility}, ce morphisme est aussi égal à
\begin{align*}
Lf^*(f_*\mathcal{K}^\bullet
\otimes_{\mathcal{O}_\mathcal{D}}^\mathbf{L}
f_*\mathcal{M}^\bullet)
& \to
Lf^*\text{Tot}(
f_*\mathcal{K}^\bullet
\otimes_{\mathcal{O}_\mathcal{D}}
f_*\mathcal{M}^\bullet) \\
& \to
Lf^*f_*\text{Tot}(\mathcal{K}^\bullet
\otimes_{\mathcal{O}_\mathcal{C}}
\mathcal{M}^\bullet) \\
& \to
f^*f_*\text{Tot}(\mathcal{K}^\bullet
\otimes_{\mathcal{O}_\mathcal{C}}
\mathcal{M}^\bullet) \\
& \to
\text{Tot}(\mathcal{K}^\bullet
\otimes_{\mathcal{O}_\mathcal{C}}
\mathcal{M}^\bullet)
\end{align*}
Pour achever la preuve, il suffit maintenant de considérer le diagramme
$$
\xymatrix{
Lf^*(f_*\mathcal{K}^\bullet
\otimes_{\mathcal{O}_\mathcal{D}}^\mathbf{L}
f_*\mathcal{M}^\bullet) \ar[d] \ar[rr] & &
Lf^*f_*\mathcal{K}^\bullet \otimes_{\mathcal{O}_\mathcal{C}}^\mathbf{L}
Lf^*f_*\mathcal{M}^\bullet \ar[d] \\
Lf^*\text{Tot}(
f_*\mathcal{K}^\bullet
\otimes_{\mathcal{O}_\mathcal{D}}
f_*\mathcal{M}^\bullet) \ar[d]_{naive} \ar[r] &
f^*\text{Tot}(
f_*\mathcal{K}^\bullet
\otimes_{\mathcal{O}_\mathcal{D}}
f_*\mathcal{M}^\bullet) \ar[ldd]^{naive} \ar[dd] &
f^*f_*\mathcal{K}^\bullet \otimes_{\mathcal{O}_\mathcal{C}}^\mathbf{L}
f^*f_*\mathcal{M}^\bullet \ar[dd] \ar[ldd] \\
Lf^*f_*\text{Tot}(\mathcal{K}^\bullet
\otimes_{\mathcal{O}_\mathcal{C}}
\mathcal{M}^\bullet) \ar[d] \\
f^*f_*\text{Tot}(\mathcal{K}^\bullet \otimes_{\mathcal{O}_\mathcal{C}}
\mathcal{M}^\bullet) \ar[rd] &
\text{Tot}(f^*f_*\mathcal{K}^\bullet \otimes_{\mathcal{O}_\mathcal{C}}
f^*f_*\mathcal{M}^\bullet) \ar[d] &
\mathcal{K}^\bullet \otimes_{\mathcal{O}_\mathcal{C}}^\mathbf{L}
\mathcal{M}^\bullet \ar[ld] \\
& \text{Tot}(\mathcal{K}^\bullet
\otimes_{\mathcal{O}_\mathcal{C}}
\mathcal{M}^\bullet)
}
$$
Tous les polygones de ce diagramme sont commutatifs. Celui du haut l'est
d'après le lemme \ref{sites-cohomology-lemma-tensor-pull-compatibility}.
Le carré comportant les deux cup-produits naïfs est commutatif parce que
$Lf^* \to f^*$ est fonctoriel en le complexe de modules.
Il en va de même du carré faisant intervenir les deux morphismes
$\mathcal{A}^\bullet \otimes^\mathbf{L} \mathcal{B}^\bullet \to
\text{Tot}(\mathcal{A}^\bullet \otimes \mathcal{B}^\bullet)$.
Enfin, le dernier carré est commutatif au niveau des complexes et cela peut
être considéré comme la définition du cup-produit naïf (par l'adjonction
entre $f^*$ et $f_*$). La preuve est achevée, car les deux parcours
du bord extérieur du diagramme donnent les deux morphismes
décrits ci-dessus.
\end{proof}

\begin{lemma}
\label{sites-cohomology-lemma-cup-product-associative}
Soit $f : (\Sh(\mathcal{C}), \mathcal{O}) \to (\Sh(\mathcal{C}'), \mathcal{O}')$
un morphisme de topos annelés. Le cup-produit relatif de la
remarque \ref{sites-cohomology-remark-cup-product} est associatif au sens où
le diagramme
$$
\xymatrix{
Rf_*K \otimes_{\mathcal{O}'}^\mathbf{L}
Rf_*L \otimes_{\mathcal{O}'}^\mathbf{L}
Rf_*M \ar[r] \ar[d] &
Rf_*(K \otimes_\mathcal{O}^\mathbf{L} L)
\otimes_{\mathcal{O}'}^\mathbf{L} Rf_*M \ar[d] \\
Rf_*K \otimes_{\mathcal{O}'}^\mathbf{L}
Rf_*(L \otimes_\mathcal{O}^\mathbf{L} M) \ar[r] &
Rf_*(K \otimes_\mathcal{O}^\mathbf{L} 
L \otimes_\mathcal{O}^\mathbf{L} M)
}
$$
est commutatif dans $D(\mathcal{O}')$ pour tous $K, L, M$ dans $D(\mathcal{O})$.
\end{lemma}

\begin{proof}
En suivant l'un ou l'autre chemin, on obtient le morphisme adjoint au morphisme évident
\begin{align*}
Lf^*(Rf_*K \otimes_{\mathcal{O}'}^\mathbf{L}
Rf_*L \otimes_{\mathcal{O}'}^\mathbf{L}
Rf_*M) & =
Lf^*(Rf_*K) \otimes_\mathcal{O}^\mathbf{L}
Lf^*(Rf_*L) \otimes_\mathcal{O}^\mathbf{L}
Lf^*(Rf_*M) \\
& \to
K \otimes_\mathcal{O}^\mathbf{L} 
L \otimes_\mathcal{O}^\mathbf{L} M
\end{align*}
dans $D(\mathcal{O})$.
\end{proof}

\begin{lemma}
\label{sites-cohomology-lemma-cup-product-commutative}
Soit $f : (\Sh(\mathcal{C}), \mathcal{O}) \to (\Sh(\mathcal{C}'), \mathcal{O}')$
un morphisme de topos annelés. Le cup-produit relatif de la
remarque \ref{sites-cohomology-remark-cup-product} est commutatif au sens où
le diagramme
$$
\xymatrix{
Rf_*K \otimes_{\mathcal{O}'}^\mathbf{L} Rf_*L \ar[r] \ar[d]_\psi &
Rf_*(K \otimes_\mathcal{O}^\mathbf{L} L) \ar[d]^{Rf_*\psi} \\
Rf_*L \otimes_{\mathcal{O}'}^\mathbf{L} Rf_*K \ar[r] &
Rf_*(L \otimes_\mathcal{O}^\mathbf{L} K)
}
$$
est commutatif dans $D(\mathcal{O}')$ pour tous $K, L$ dans $D(\mathcal{O})$.
Ici, $\psi$ est la contrainte de commutativité de la catégorie dérivée
(lemme \ref{sites-cohomology-lemma-symmetric-monoidal-derived}).
\end{lemma}

\begin{proof}
La démonstration est omise.
\end{proof}

\begin{lemma}
\label{sites-cohomology-lemma-compose-cup-product}
Soient $f : (\Sh(\mathcal{C}), \mathcal{O}) \to (\Sh(\mathcal{C}'), \mathcal{O}')$
et $f' : (\Sh(\mathcal{C}'), \mathcal{O}') \to
(\Sh(\mathcal{C}''), \mathcal{O}'')$
des morphismes de topos annelés. Le cup-produit relatif de la
remarque \ref{sites-cohomology-remark-cup-product} est compatible aux compositions
au sens où le diagramme
$$
\xymatrix{
R(f' \circ f)_*K \otimes_{\mathcal{O}''}^\mathbf{L} R(f' \circ f)_*L
\ar@{=}[rr] \ar[d] & &
Rf'_*Rf_*K \otimes_{\mathcal{O}''}^\mathbf{L} Rf'_*Rf_*L \ar[d] \\
R(f' \circ f)_*(K \otimes_\mathcal{O}^\mathbf{L} L) \ar@{=}[r] &
Rf'_*Rf_*(K \otimes_\mathcal{O}^\mathbf{L} L) &
Rf'_*(Rf_*K \otimes_{\mathcal{O}'}^\mathbf{L}  Rf_*L) \ar[l]
}
$$
est commutatif dans $D(\mathcal{O}'')$ pour tous $K, L$ dans $D(\mathcal{O})$.
\end{lemma}

\begin{proof}
Cela résulte du fait qu'en suivant l'un ou l'autre chemin dans le diagramme,
on obtient le morphisme adjoint au morphisme
\begin{align*}
& L(f' \circ f)^*\left(R(f' \circ f)_*K
\otimes_{\mathcal{O}''}^\mathbf{L}
R(f' \circ f)_*L\right) \\
& =
L(f' \circ f)^*R(f' \circ f)_*K
\otimes_\mathcal{O}^\mathbf{L}
L(f' \circ f)^*R(f' \circ f)_*L \\
& \to
K \otimes_\mathcal{O}^\mathbf{L} L
\end{align*}
dans $D(\mathcal{O})$. Pour le voir, on utilise le fait que la composée
des co-unités
$$
L(f' \circ f)^*R(f' \circ f)_* =
Lf^* L(f')^* Rf'_* Rf_*  \to
Lf^* Rf_* \to \text{id}
$$
est la co-unité de l'adjonction entre $L(f' \circ f)^*$ et
$R(f' \circ f)_*$. Voir Catégories, lemme \ref{categories-lemma-compose-counits}.
\end{proof}

\begin{lemma}
\label{sites-cohomology-lemma-base-change-cup-product}
Considérons un carré commutatif
$$
\xymatrix{
(\Sh(\mathcal{C}'), \mathcal{O}_{\mathcal{C}'}) \ar[r]_{g'} \ar[d]_{f'} &
(\Sh(\mathcal{C}), \mathcal{O}_{\mathcal{C}}) \ar[d]^f \\
(\Sh(\mathcal{D}'), \mathcal{O}_{\mathcal{D}'}) \ar[r]^g &
(\Sh(\mathcal{D}), \mathcal{O}_{\mathcal{D}}) 
}
$$
de topos annelés. Soient $K, L$ dans $D(\mathcal{O}_{\mathcal{C}})$.
Le cup-produit relatif est compatible avec ce carré
au sens où le diagramme
$$
\xymatrix{
Lg^*(Rf_*K \otimes_{\mathcal{O}_{\mathcal{D}}}^\mathbf{L} Rf_*L)
\ar[r] \ar@{=}[d] &
Lg^*(Rf_*(K \otimes_{\mathcal{O}_{\mathcal{C}}}^\mathbf{L} L)) \ar[d] \\
Lg^*Rf_*K \otimes_{\mathcal{O}_{\mathcal{D}'}}^\mathbf{L} Lg^*Rf_*L \ar[d] &
R(f')_*L(g')^*(K \otimes_{\mathcal{O}_{\mathcal{C}}}^\mathbf{L} L) \ar@{=}[d] \\
R(f')_*L(g')^*K \otimes_{\mathcal{O}_{\mathcal{D}'}}^\mathbf{L} R(f')_*L(g')^*L \ar[r] &
R(f')_*(L(g')^*K \otimes_{\mathcal{O}_{\mathcal{C}'}}^\mathbf{L} L(g')^*L)
}
$$
est commutatif dans $D(\mathcal{O}_{\mathcal{D}'})$.
Les flèches horizontales sont données
par le cup-produit relatif (remarque \ref{sites-cohomology-remark-cup-product})
et les flèches verticales sont données
par le morphisme de changement de base (remarque \ref{sites-cohomology-remark-base-change})
et le lemme \ref{sites-cohomology-lemma-pullback-tensor-product}.
\end{lemma}

\begin{proof}
La démonstration est omise.
\end{proof}






\section{Complexes de Hom}
\label{sites-cohomology-section-hom-complexes}

\noindent
Soit $(\mathcal{C}, \mathcal{O})$ un site annelé. Soient
$\mathcal{L}^\bullet$ et $\mathcal{M}^\bullet$ deux complexes
de $\mathcal{O}$-modules. Nous construisons un complexe
de $\mathcal{O}$-modules
$\SheafHom^\bullet(\mathcal{L}^\bullet, \mathcal{M}^\bullet)$.
Plus précisément, pour tout $n$, nous posons
$$
\SheafHom^n(\mathcal{L}^\bullet, \mathcal{M}^\bullet) =
\prod\nolimits_{n = p + q}
\SheafHom_\mathcal{O}(\mathcal{L}^{-q}, \mathcal{M}^p)
$$
On peut voir $\SheafHom^n$ comme le faisceau de
$\mathcal{O}$-modules formé de toutes les applications $\mathcal{O}$-linéaires
de $\mathcal{L}^\bullet$ dans $\mathcal{M}^\bullet$
(considérés comme des $\mathcal{O}$-modules gradués) qui sont homogènes
de degré $n$. Dans cette terminologie, nous définissons la différentielle par la formule
$$
\text{d}(f) =
\text{d}_\mathcal{M} \circ f - (-1)^n f \circ \text{d}_\mathcal{L}
$$
pour
$f \in \SheafHom^n_\mathcal{O}(\mathcal{L}^\bullet, \mathcal{M}^\bullet)$.
Nous omettons la vérification que $\text{d}^2 = 0$.
Cette construction est un cas particulier d'
Algèbre différentielle graduée, exemple \ref{dga-example-category-complexes}.
Il résulte immédiatement de la construction que l'on a
\begin{equation}
\label{sites-cohomology-equation-cohomology-hom-complex}
H^n(\Gamma(U, \SheafHom^\bullet(\mathcal{L}^\bullet, \mathcal{M}^\bullet))) =
\Hom_{K(\mathcal{O}_U)}(\mathcal{L}^\bullet|_U, \mathcal{M}^\bullet[n]|_U)
\end{equation}
pour tout $n \in \mathbf{Z}$ et tout $U \in \Ob(\mathcal{C})$. De même,
on a
\begin{equation}
\label{sites-cohomology-equation-global-cohomology-hom-complex}
H^n(\Gamma(\mathcal{C},
\SheafHom^\bullet(\mathcal{L}^\bullet, \mathcal{M}^\bullet))) =
\Hom_{K(\mathcal{O})}(\mathcal{L}^\bullet, \mathcal{M}^\bullet[n])
\end{equation}
pour le complexe des sections globales.

\begin{lemma}
\label{sites-cohomology-lemma-compose}
Soit $(\mathcal{C}, \mathcal{O})$ un site annelé.
Étant donnés des complexes $\mathcal{K}^\bullet, \mathcal{L}^\bullet, \mathcal{M}^\bullet$
de $\mathcal{O}$-modules, il existe un isomorphisme
$$
\SheafHom^\bullet(\mathcal{K}^\bullet,
\SheafHom^\bullet(\mathcal{L}^\bullet, \mathcal{M}^\bullet))
=
\SheafHom^\bullet(\text{Tot}(\mathcal{K}^\bullet \otimes_\mathcal{O}
\mathcal{L}^\bullet), \mathcal{M}^\bullet)
$$
de complexes de $\mathcal{O}$-modules, fonctoriel en
$\mathcal{K}^\bullet, \mathcal{L}^\bullet, \mathcal{M}^\bullet$.
\end{lemma}

\begin{proof}
Démonstration omise. Indication : cela se démontre exactement comme dans
Compléments d'algèbre, lemme \ref{more-algebra-lemma-compose}.
\end{proof}

\begin{lemma}
\label{sites-cohomology-lemma-composition}
Soit $(\mathcal{C}, \mathcal{O})$ un site annelé. Étant donnés des complexes
$\mathcal{K}^\bullet, \mathcal{L}^\bullet, \mathcal{M}^\bullet$
de $\mathcal{O}$-modules, il existe un morphisme canonique
$$
\text{Tot}\left(
\SheafHom^\bullet(\mathcal{L}^\bullet, \mathcal{M}^\bullet)
\otimes_\mathcal{O}
\SheafHom^\bullet(\mathcal{K}^\bullet, \mathcal{L}^\bullet)
\right)
\longrightarrow
\SheafHom^\bullet(\mathcal{K}^\bullet, \mathcal{M}^\bullet)
$$
de complexes de $\mathcal{O}$-modules.
\end{lemma}

\begin{proof}
Démonstration omise. Indication : cela se démontre exactement comme dans
Compléments d'algèbre, lemme \ref{more-algebra-lemma-composition}.
\end{proof}

\begin{lemma}
\label{sites-cohomology-lemma-diagonal-better}
Soit $(\mathcal{C}, \mathcal{O})$ un site annelé. Étant donnés des complexes
$\mathcal{K}^\bullet, \mathcal{L}^\bullet, \mathcal{M}^\bullet$
de $\mathcal{O}$-modules, il existe un morphisme canonique
$$
\text{Tot}\left(
\mathcal{K}^\bullet \otimes_\mathcal{O}
\SheafHom^\bullet(\mathcal{M}^\bullet, \mathcal{L}^\bullet)
\right)
\longrightarrow
\SheafHom^\bullet(\mathcal{M}^\bullet,
\text{Tot}(\mathcal{K}^\bullet \otimes_\mathcal{O} \mathcal{L}^\bullet))
$$
de complexes de $\mathcal{O}$-modules, fonctoriel par rapport aux trois complexes.
\end{lemma}

\begin{proof}
Démonstration omise. Indication : cela se démontre exactement comme dans
Compléments d'algèbre, lemme \ref{more-algebra-lemma-diagonal-better}.
\end{proof}

\begin{lemma}
\label{sites-cohomology-lemma-diagonal}
Soit $(\mathcal{C}, \mathcal{O})$ un site annelé. Étant donnés des complexes
$\mathcal{K}^\bullet, \mathcal{L}^\bullet$
de $\mathcal{O}$-modules, il existe un morphisme canonique
$$
\mathcal{K}^\bullet
\longrightarrow
\SheafHom^\bullet(\mathcal{L}^\bullet,
\text{Tot}(\mathcal{K}^\bullet \otimes_\mathcal{O} \mathcal{L}^\bullet))
$$
de complexes de $\mathcal{O}$-modules, fonctoriel par rapport aux deux complexes.
\end{lemma}

\begin{proof}
Démonstration omise. Indication : cela se démontre exactement comme dans
Compléments d'algèbre, lemme \ref{more-algebra-lemma-diagonal}.
\end{proof}

\begin{lemma}
\label{sites-cohomology-lemma-evaluate-and-more}
Soit $(\mathcal{C}, \mathcal{O})$ un site annelé. Étant donnés des complexes
$\mathcal{K}^\bullet, \mathcal{L}^\bullet, \mathcal{M}^\bullet$
de $\mathcal{O}$-modules, il existe un morphisme canonique
$$
\text{Tot}(\SheafHom^\bullet(\mathcal{L}^\bullet,
\mathcal{M}^\bullet) \otimes_\mathcal{O} \mathcal{K}^\bullet)
\longrightarrow
\SheafHom^\bullet(\SheafHom^\bullet(\mathcal{K}^\bullet,
\mathcal{L}^\bullet), \mathcal{M}^\bullet)
$$
de complexes de $\mathcal{O}$-modules, fonctoriel par rapport aux trois complexes.
\end{lemma}

\begin{proof}
Démonstration omise. Indication : cela se démontre exactement comme dans
Compléments d'algèbre, lemme \ref{more-algebra-lemma-evaluate-and-more}.
\end{proof}

\begin{lemma}
\label{sites-cohomology-lemma-RHom-into-K-injective}
Soit $(\mathcal{C}, \mathcal{O})$ un site annelé. Soient $L$ et $M$
des objets de $D(\mathcal{O})$. Soit $\mathcal{I}^\bullet$
un complexe K-injectif de $\mathcal{O}$-modules représentant $M$. Soit
$\mathcal{L}^\bullet$ un complexe de $\mathcal{O}$-modules
représentant $L$.
Alors
$$
H^0(\Gamma(U, \SheafHom^\bullet(\mathcal{L}^\bullet, \mathcal{I}^\bullet))) =
\Hom_{D(\mathcal{O}_U)}(L|_U, M|_U)
$$
pour tout $U \in \Ob(\mathcal{C})$. De même,
$H^0(\Gamma(\mathcal{C},
\SheafHom^\bullet(\mathcal{L}^\bullet, \mathcal{I}^\bullet))) =
\Hom_{D(\mathcal{O})}(L, M)$.
\end{lemma}

\begin{proof}
On a
\begin{align*}
H^0(\Gamma(U, \SheafHom^\bullet(\mathcal{L}^\bullet, \mathcal{I}^\bullet)))
& =
\Hom_{K(\mathcal{O}_U)}(\mathcal{L}^\bullet|_U, \mathcal{I}^\bullet|_U) \\
& =
\Hom_{D(\mathcal{O}_U)}(L|_U, M|_U)
\end{align*}
La première égalité est (\ref{sites-cohomology-equation-cohomology-hom-complex}).
La seconde est vraie parce que $\mathcal{I}^\bullet|_U$
est K-injectif d'après le lemme \ref{sites-cohomology-lemma-restrict-K-injective-to-open}.
La démonstration de l'égalité portant sur les sections globales est analogue, mais utilise
(\ref{sites-cohomology-equation-global-cohomology-hom-complex}).
\end{proof}

\begin{lemma}
\label{sites-cohomology-lemma-RHom-well-defined}
Soit $(\mathcal{C}, \mathcal{O})$ un site annelé. Soit
$(\mathcal{I}')^\bullet \to \mathcal{I}^\bullet$
un quasi-isomorphisme de complexes K-injectifs de $\mathcal{O}$-modules.
Soit $(\mathcal{L}')^\bullet \to \mathcal{L}^\bullet$
un quasi-isomorphisme de complexes de $\mathcal{O}$-modules.
Alors
$$
\SheafHom^\bullet(\mathcal{L}^\bullet, (\mathcal{I}')^\bullet)
\longrightarrow
\SheafHom^\bullet((\mathcal{L}')^\bullet, \mathcal{I}^\bullet)
$$
est un quasi-isomorphisme.
\end{lemma}

\begin{proof}
Soit $M$ l'objet de $D(\mathcal{O})$ représenté par
$\mathcal{I}^\bullet$ et $(\mathcal{I}')^\bullet$.
Soit $L$ l'objet de $D(\mathcal{O})$ représenté par
$\mathcal{L}^\bullet$ et $(\mathcal{L}')^\bullet$.
D'après le lemme \ref{sites-cohomology-lemma-RHom-into-K-injective},
les faisceaux
$$
H^0(\SheafHom^\bullet(\mathcal{L}^\bullet, (\mathcal{I}')^\bullet))
\quad\text{et}\quad
H^0(\SheafHom^\bullet((\mathcal{L}')^\bullet, \mathcal{I}^\bullet))
$$
sont tous deux égaux au faisceau associé au préfaisceau
$$
U \longmapsto \Hom_{D(\mathcal{O}_U)}(L|_U, M|_U)
$$
Par fonctorialité, le morphisme induit l'identité de ce faisceau.
Le même argument appliqué après décalage montre qu'il induit un
isomorphisme sur tous les faisceaux de cohomologie. C'est donc un quasi-isomorphisme.
\end{proof}

\begin{lemma}
\label{sites-cohomology-lemma-RHom-from-K-flat-into-K-injective}
Soit $(\mathcal{C}, \mathcal{O})$ un site annelé. Soit $\mathcal{I}^\bullet$
un complexe K-injectif de $\mathcal{O}$-modules. Soit
$\mathcal{L}^\bullet$ un complexe K-plat de $\mathcal{O}$-modules.
Alors $\SheafHom^\bullet(\mathcal{L}^\bullet, \mathcal{I}^\bullet)$
est un complexe K-injectif de $\mathcal{O}$-modules.
\end{lemma}

\begin{proof}
En effet, si $\mathcal{K}^\bullet$ est un complexe acyclique de
$\mathcal{O}$-modules, alors
\begin{align*}
\Hom_{K(\mathcal{O})}(\mathcal{K}^\bullet,
\SheafHom^\bullet(\mathcal{L}^\bullet, \mathcal{I}^\bullet))
& =
H^0(\Gamma(\mathcal{C},
\SheafHom^\bullet(\mathcal{K}^\bullet,
\SheafHom^\bullet(\mathcal{L}^\bullet, \mathcal{I}^\bullet)))) \\
& =
H^0(\Gamma(\mathcal{C}, \SheafHom^\bullet(\text{Tot}(
\mathcal{K}^\bullet \otimes_\mathcal{O} \mathcal{L}^\bullet),
\mathcal{I}^\bullet))) \\
& =
\Hom_{K(\mathcal{O})}(
\text{Tot}(\mathcal{K}^\bullet \otimes_\mathcal{O} \mathcal{L}^\bullet),
\mathcal{I}^\bullet) \\
& =
0
\end{align*}
La première égalité résulte de (\ref{sites-cohomology-equation-global-cohomology-hom-complex}).
La deuxième résulte du lemme \ref{sites-cohomology-lemma-compose}.
La troisième résulte de (\ref{sites-cohomology-equation-global-cohomology-hom-complex}).
Enfin,
$\text{Tot}(\mathcal{K}^\bullet \otimes_\mathcal{O} \mathcal{L}^\bullet)$
est acyclique puisque $\mathcal{L}^\bullet$ est K-plat
(définition \ref{sites-cohomology-definition-K-flat}) ; la dernière égalité résulte donc du fait que
$\mathcal{I}^\bullet$ est K-injectif.
\end{proof}








\section{Hom interne dans la catégorie dérivée}
\label{sites-cohomology-section-internal-hom}

\noindent
Soit $(\mathcal{C}, \mathcal{O})$ un site annelé. Soient $L, M$ des objets
de $D(\mathcal{O})$. Nous souhaitons construire un objet
$R\SheafHom(L, M)$ de $D(\mathcal{O})$ tel que, pour tout troisième
objet $K$ de $D(\mathcal{O})$, il existe une bijection canonique
\begin{equation}
\label{sites-cohomology-equation-internal-hom}
\Hom_{D(\mathcal{O})}(K, R\SheafHom(L, M))
=
\Hom_{D(\mathcal{O})}(K \otimes_\mathcal{O}^\mathbf{L} L, M)
\end{equation}
Observons que cette formule définit $R\SheafHom(L, M)$ à isomorphisme
unique près, par le lemme de Yoneda
(Catégories, lemme \ref{categories-lemma-yoneda}).

\medskip\noindent
Pour construire un tel objet, choisissons un complexe K-injectif de
$\mathcal{O}$-modules $\mathcal{I}^\bullet$ représentant $M$ et un
complexe quelconque de $\mathcal{O}$-modules $\mathcal{L}^\bullet$ représentant $L$.
Posons alors
$$
R\SheafHom(L, M) = \SheafHom^\bullet(\mathcal{L}^\bullet, \mathcal{I}^\bullet)
$$
où le membre de droite est le complexe de $\mathcal{O}$-modules
construit à la section \ref{sites-cohomology-section-hom-complexes}.
Cette définition est bien fondée par le lemme \ref{sites-cohomology-lemma-RHom-well-defined}.
On obtient un foncteur
$$
D(\mathcal{O})^{opp} \times D(\mathcal{O}) \longrightarrow D(\mathcal{O}),
\quad
(K, L) \longmapsto R\SheafHom(K, L)
$$
Avant de démontrer (\ref{sites-cohomology-equation-internal-hom}),
calculons les groupes de cohomologie de $R\SheafHom(K, L)$.

\begin{lemma}
\label{sites-cohomology-lemma-section-RHom-over-U}
Soit $(\mathcal{C}, \mathcal{O})$ un site annelé. Soient $L, M$ des objets
de $D(\mathcal{O})$. Pour tout objet $U$ de $\mathcal{C}$, on a
$$
H^0(U, R\SheafHom(L, M)) =
\Hom_{D(\mathcal{O}_U)}(L|_U, M|_U)
$$
et l'on a $H^0(\mathcal{C}, R\SheafHom(L, M)) =
\Hom_{D(\mathcal{O})}(L, M)$.
\end{lemma}

\begin{proof}
Choisissons un complexe K-injectif $\mathcal{I}^\bullet$ de
$\mathcal{O}$-modules représentant $M$ et un complexe K-plat
$\mathcal{L}^\bullet$ représentant $L$. Alors
$\SheafHom^\bullet(\mathcal{L}^\bullet, \mathcal{I}^\bullet)$
est K-injectif par le lemme \ref{sites-cohomology-lemma-RHom-from-K-flat-into-K-injective}.
On peut donc calculer la cohomologie sur $U$ en prenant simplement les sections sur $U$,
et le résultat découle du lemme \ref{sites-cohomology-lemma-RHom-into-K-injective}.
\end{proof}

\begin{lemma}
\label{sites-cohomology-lemma-internal-hom}
Soit $(\mathcal{C}, \mathcal{O})$ un site annelé. Soient $K, L, M$ des objets
de $D(\mathcal{O})$. Pour la construction décrite ci-dessus,
il existe un isomorphisme canonique
$$
R\SheafHom(K, R\SheafHom(L, M)) =
R\SheafHom(K \otimes_\mathcal{O}^\mathbf{L} L, M)
$$
dans $D(\mathcal{O})$, fonctoriel en $K, L, M$,
qui redonne (\ref{sites-cohomology-equation-internal-hom}) après application de $H^0(\mathcal{C}, -)$.
\end{lemma}

\begin{proof}
Choisissons un complexe K-injectif $\mathcal{I}^\bullet$ représentant
$M$ et un complexe K-plat de $\mathcal{O}$-modules $\mathcal{L}^\bullet$
représentant $L$.
Pour tout complexe de $\mathcal{O}$-modules $\mathcal{K}^\bullet$,
on a
$$
\SheafHom^\bullet(\mathcal{K}^\bullet,
\SheafHom^\bullet(\mathcal{L}^\bullet, \mathcal{I}^\bullet))
=
\SheafHom^\bullet(
\text{Tot}(\mathcal{K}^\bullet \otimes_\mathcal{O} \mathcal{L}^\bullet),
\mathcal{I}^\bullet)
$$
par le lemme \ref{sites-cohomology-lemma-compose}.
Remarquons que le membre de gauche représente
$R\SheafHom(K, R\SheafHom(L, M))$ (utiliser le
lemme \ref{sites-cohomology-lemma-RHom-from-K-flat-into-K-injective})
et que le membre de droite représente
$R\SheafHom(K \otimes_\mathcal{O}^\mathbf{L} L, M)$.
Cela démontre la formule affichée du lemme.
En prenant les sections globales et en utilisant le lemme \ref{sites-cohomology-lemma-section-RHom-over-U},
on obtient (\ref{sites-cohomology-equation-internal-hom}).
\end{proof}

\begin{lemma}
\label{sites-cohomology-lemma-restriction-RHom-to-U}
Soit $(\mathcal{C}, \mathcal{O})$ un site annelé. Soient $K, L$ des objets
de $D(\mathcal{O})$. La construction de $R\SheafHom(K, L)$
commute aux restrictions, c'est-à-dire que,
pour tout objet $U$ de $\mathcal{C}$, on a
$R\SheafHom(K|_U, L|_U) = R\SheafHom(K, L)|_U$.
\end{lemma}

\begin{proof}
Cela résulte immédiatement de la construction et du
lemme \ref{sites-cohomology-lemma-restrict-K-injective-to-open}.
\end{proof}

\begin{lemma}
\label{sites-cohomology-lemma-RHom-triangulated}
Soit $(\mathcal{C}, \mathcal{O})$ un site annelé. Le bifoncteur
$R\SheafHom(- , -)$ transforme les triangles distingués en
triangles distingués dans chacune des deux variables.
\end{lemma}

\begin{proof}
Cela découle du fait que l'application
$$
(\mathcal{L}^\bullet, \mathcal{M}^\bullet) \longmapsto
\SheafHom^\bullet(\mathcal{L}^\bullet, \mathcal{M}^\bullet)
$$
transforme, dans chacune des variables, toute suite exacte courte de complexes
scindée terme à terme en une suite exacte courte scindée terme à terme. Détails omis.
\end{proof}

\begin{lemma}
\label{sites-cohomology-lemma-internal-hom-evaluate}
Soit $(\mathcal{C}, \mathcal{O})$ un site annelé. Soient $K, L, M$ des objets
de $D(\mathcal{O})$. Il existe un morphisme canonique
$$
R\SheafHom(L, M) \otimes_\mathcal{O}^\mathbf{L} K
\longrightarrow
R\SheafHom(R\SheafHom(K, L), M)
$$
dans $D(\mathcal{O})$, fonctoriel en $K, L, M$.
\end{lemma}

\begin{proof}
Choisissons un complexe K-injectif $\mathcal{I}^\bullet$ représentant $M$,
un complexe K-injectif $\mathcal{J}^\bullet$ représentant $L$, et un complexe
K-plat $\mathcal{K}^\bullet$ représentant $K$. Le morphisme est défini à
l'aide du morphisme
$$
\text{Tot}(\SheafHom^\bullet(\mathcal{J}^\bullet,
\mathcal{I}^\bullet) \otimes_\mathcal{O} \mathcal{K}^\bullet)
\longrightarrow
\SheafHom^\bullet(\SheafHom^\bullet(\mathcal{K}^\bullet,
\mathcal{J}^\bullet), \mathcal{I}^\bullet)
$$
du lemme \ref{sites-cohomology-lemma-evaluate-and-more}. Avec notre choix particulier de
complexes, le membre de gauche représente
$R\SheafHom(L, M) \otimes_\mathcal{O}^\mathbf{L} K$
et le membre de droite représente
$R\SheafHom(R\SheafHom(K, L), M)$. Nous omettons la démonstration de la
fonctorialité par rapport aux trois objets de $D(\mathcal{O})$.
\end{proof}

\begin{lemma}
\label{sites-cohomology-lemma-internal-hom-composition}
\begin{slogan}
Composition sur RSheafHom.
\end{slogan}
Soit $(\mathcal{C}, \mathcal{O})$ un site annelé. Pour $K, L, M$ dans
$D(\mathcal{O})$, il existe un morphisme canonique
$$
R\SheafHom(L, M) \otimes_\mathcal{O}^\mathbf{L} R\SheafHom(K, L)
\longrightarrow R\SheafHom(K, M)
$$
dans $D(\mathcal{O})$.
\end{lemma}

\begin{proof}
Choisissons un complexe K-injectif $\mathcal{I}^\bullet$ représentant $M$,
un complexe K-injectif $\mathcal{J}^\bullet$ représentant $L$, et un complexe
quelconque de $\mathcal{O}$-modules $\mathcal{K}^\bullet$ représentant $K$.
Par le lemme \ref{sites-cohomology-lemma-composition}, il existe un morphisme de complexes
$$
\text{Tot}\left(
\SheafHom^\bullet(\mathcal{J}^\bullet, \mathcal{I}^\bullet)
\otimes_\mathcal{O}
\SheafHom^\bullet(\mathcal{K}^\bullet, \mathcal{J}^\bullet)
\right)
\longrightarrow
\SheafHom^\bullet(\mathcal{K}^\bullet, \mathcal{I}^\bullet)
$$
Les complexes de $\mathcal{O}$-modules
$\SheafHom^\bullet(\mathcal{J}^\bullet, \mathcal{I}^\bullet)$,
$\SheafHom^\bullet(\mathcal{K}^\bullet, \mathcal{J}^\bullet)$ et
$\SheafHom^\bullet(\mathcal{K}^\bullet, \mathcal{I}^\bullet)$
représentent $R\SheafHom(L, M)$, $R\SheafHom(K, L)$ et $R\SheafHom(K, M)$.
Si nous choisissons un complexe K-plat $\mathcal{H}^\bullet$ et un
quasi-isomorphisme
$\mathcal{H}^\bullet \to
\SheafHom^\bullet(\mathcal{K}^\bullet, \mathcal{J}^\bullet)$,
alors il existe un morphisme
$$
\text{Tot}\left(
\SheafHom^\bullet(\mathcal{J}^\bullet, \mathcal{I}^\bullet)
\otimes_\mathcal{O} \mathcal{H}^\bullet
\right)
\longrightarrow
\text{Tot}\left(
\SheafHom^\bullet(\mathcal{J}^\bullet, \mathcal{I}^\bullet)
\otimes_\mathcal{O}
\SheafHom^\bullet(\mathcal{K}^\bullet, \mathcal{J}^\bullet)
\right)
$$
dont la source représente
$R\SheafHom(L, M) \otimes_\mathcal{O}^\mathbf{L} R\SheafHom(K, L)$.
En composant les deux flèches affichées, on obtient le morphisme recherché.
Nous omettons la démonstration de la fonctorialité de la construction.
\end{proof}

\begin{lemma}
\label{sites-cohomology-lemma-internal-hom-diagonal-better}
Soit $(\mathcal{C}, \mathcal{O})$ un site annelé. Pour $K, L, M$ dans
$D(\mathcal{O})$, il existe un morphisme canonique
$$
K \otimes_\mathcal{O}^\mathbf{L} R\SheafHom(M, L)
\longrightarrow
R\SheafHom(M, K \otimes_\mathcal{O}^\mathbf{L} L)
$$
dans $D(\mathcal{O})$, fonctoriel en $K, L, M$.
\end{lemma}

\begin{proof}
Choisissons un complexe K-plat $\mathcal{K}^\bullet$ représentant $K$, un
complexe K-injectif $\mathcal{I}^\bullet$ représentant $L$, et un complexe
quelconque de $\mathcal{O}$-modules $\mathcal{M}^\bullet$ représentant $M$.
Choisissons un quasi-isomorphisme
$\text{Tot}(\mathcal{K}^\bullet \otimes_{\mathcal{O}_X} \mathcal{I}^\bullet)
\to \mathcal{J}^\bullet$
où $\mathcal{J}^\bullet$ est K-injectif. Nous utilisons alors le morphisme
$$
\text{Tot}\left(
\mathcal{K}^\bullet \otimes_\mathcal{O}
\SheafHom^\bullet(\mathcal{M}^\bullet, \mathcal{I}^\bullet)
\right)
\to
\SheafHom^\bullet(\mathcal{M}^\bullet,
\text{Tot}(\mathcal{K}^\bullet \otimes_\mathcal{O} \mathcal{I}^\bullet))
\to
\SheafHom^\bullet(\mathcal{M}^\bullet, \mathcal{J}^\bullet)
$$
où le premier morphisme est celui du lemme \ref{sites-cohomology-lemma-diagonal-better}.
\end{proof}

\begin{lemma}
\label{sites-cohomology-lemma-internal-hom-diagonal}
Soit $(\mathcal{C}, \mathcal{O})$ un site annelé. Pour $K, L$ dans
$D(\mathcal{O})$, il existe un morphisme canonique
$$
K \longrightarrow R\SheafHom(L, K \otimes_\mathcal{O}^\mathbf{L} L)
$$
dans $D(\mathcal{O})$, fonctoriel en $K$ et $L$.
\end{lemma}

\begin{proof}
Choisissons un complexe K-plat $\mathcal{K}^\bullet$ représentant $K$ et un
complexe quelconque de $\mathcal{O}$-modules $\mathcal{L}^\bullet$
représentant $L$. Choisissons un complexe K-injectif
$\mathcal{J}^\bullet$ et un quasi-isomorphisme
$\text{Tot}(\mathcal{K}^\bullet \otimes_\mathcal{O} \mathcal{L}^\bullet)
\to \mathcal{J}^\bullet$. Nous utilisons alors
$$
\mathcal{K}^\bullet \to
\SheafHom^\bullet(\mathcal{L}^\bullet,
\text{Tot}(\mathcal{K}^\bullet \otimes_\mathcal{O} \mathcal{L}^\bullet))
\to
\SheafHom^\bullet(\mathcal{L}^\bullet, \mathcal{J}^\bullet)
$$
où le premier morphisme provient du lemme \ref{sites-cohomology-lemma-diagonal}.
\end{proof}

\begin{lemma}
\label{sites-cohomology-lemma-dual}
Soit $(\mathcal{C}, \mathcal{O})$ un site annelé et soit $L$ un objet de
$D(\mathcal{O})$. Posons $L^\vee = R\SheafHom(L, \mathcal{O})$.
Pour $M$ dans $D(\mathcal{O})$, il existe un morphisme canonique
\begin{equation}
\label{sites-cohomology-equation-eval}
M \otimes^\mathbf{L}_\mathcal{O} L^\vee \longrightarrow R\SheafHom(L, M)
\end{equation}
qui induit un morphisme canonique
$$
H^0(\mathcal{C}, M \otimes_\mathcal{O}^\mathbf{L} L^\vee)
\longrightarrow
\Hom_{D(\mathcal{O})}(L, M)
$$
fonctoriel en $M$ dans $D(\mathcal{O})$.
\end{lemma}

\begin{proof}
Le morphisme (\ref{sites-cohomology-equation-eval}) est un cas particulier du
lemme \ref{sites-cohomology-lemma-internal-hom-composition}, en utilisant l'identification
$M = R\SheafHom(\mathcal{O}, M)$.
\end{proof}

\begin{remark}
\label{sites-cohomology-remark-projection-formula-for-internal-hom}
Soit $f : (\Sh(\mathcal{C}), \mathcal{O}_\mathcal{C}) \to
(\Sh(\mathcal{D}), \mathcal{O}_\mathcal{D})$ un morphisme de topos annelés.
Soient $K, L$ des objets de $D(\mathcal{O}_\mathcal{C})$. Nous affirmons
qu'il existe un morphisme canonique
$$
Rf_*R\SheafHom(L, K) \longrightarrow R\SheafHom(Rf_*L, Rf_*K)
$$
En effet, d'après (\ref{sites-cohomology-equation-internal-hom}), cela revient à donner un
morphisme
$Rf_*R\SheafHom(L, K) \otimes_{\mathcal{O}_\mathcal{D}}^\mathbf{L} Rf_*L
\to Rf_*K$.
Pour cela, nous pouvons utiliser la composition
$$
Rf_*R\SheafHom(L, K) \otimes_{\mathcal{O}_\mathcal{D}}^\mathbf{L} Rf_*L \to
Rf_*(R\SheafHom(L, K) \otimes_{\mathcal{O}_\mathcal{C}}^\mathbf{L} L) \to
Rf_*K
$$
où la première flèche est le cup-produit relatif
(remarque \ref{sites-cohomology-remark-cup-product}) et la deuxième est $Rf_*$ appliqué au
morphisme canonique
$R\SheafHom(L, K) \otimes_{\mathcal{O}_\mathcal{C}}^\mathbf{L} L \to K$
provenant du lemme \ref{sites-cohomology-lemma-internal-hom-composition}
(avec $\mathcal{O}_\mathcal{C}$ à l'une des places).
\end{remark}

\begin{remark}
\label{sites-cohomology-remark-prepare-fancy-base-change}
Soit $h : (\Sh(\mathcal{C}), \mathcal{O}) \to
(\Sh(\mathcal{C}'), \mathcal{O}')$ un morphisme de topos annelés.
Soient $K, L$ des objets de $D(\mathcal{O}')$. Nous affirmons qu'il existe
un morphisme canonique
$$
Lh^*R\SheafHom(K, L) \longrightarrow R\SheafHom(Lh^*K, Lh^*L)
$$
dans $D(\mathcal{O})$. En effet, d'après (\ref{sites-cohomology-equation-internal-hom}),
démontré dans le lemme \ref{sites-cohomology-lemma-internal-hom}, donner un tel morphisme
revient à donner un morphisme
$$
Lh^*R\SheafHom(K, L) \otimes^\mathbf{L} Lh^*K \longrightarrow Lh^*L
$$
La source de cette flèche est
$Lh^*(\SheafHom(K, L) \otimes^\mathbf{L} K)$ d'après le
lemme \ref{sites-cohomology-lemma-pullback-tensor-product}. Il suffit donc de construire un
morphisme canonique
$$
R\SheafHom(K, L) \otimes^\mathbf{L} K \longrightarrow L.
$$
Pour cela, nous prenons la flèche correspondant à
$$
\text{id} :
R\SheafHom(K, L)
\longrightarrow
R\SheafHom(K, L)
$$
par (\ref{sites-cohomology-equation-internal-hom}).
\end{remark}

\begin{remark}
\label{sites-cohomology-remark-fancy-base-change}
Supposons que
$$
\xymatrix{
(\Sh(\mathcal{C}'), \mathcal{O}_{\mathcal{C}'})
\ar[r]_h \ar[d]_{f'} &
(\Sh(\mathcal{C}), \mathcal{O}_\mathcal{C}) \ar[d]^f \\
(\Sh(\mathcal{D}'), \mathcal{O}_{\mathcal{D}'})
\ar[r]^g &
(\Sh(\mathcal{D}), \mathcal{O}_\mathcal{D})
}
$$
soit un diagramme commutatif de topos annelés. Soient $K, L$ des objets de
$D(\mathcal{O}_\mathcal{C})$. Nous affirmons qu'il existe un morphisme
canonique de changement de base
$$
Lg^*Rf_*R\SheafHom(K, L)
\longrightarrow
R(f')_*R\SheafHom(Lh^*K, Lh^*L)
$$
dans $D(\mathcal{O}_{\mathcal{D}'})$. En effet, nous prenons le morphisme
adjoint à la composition
\begin{align*}
L(f')^*Lg^*Rf_*R\SheafHom(K, L)
& =
Lh^*Lf^*Rf_*R\SheafHom(K, L) \\
& \to
Lh^*R\SheafHom(K, L) \\
& \to
R\SheafHom(Lh^*K, Lh^*L)
\end{align*}
où la première flèche utilise le morphisme d'adjonction
$Lf^*Rf_* \to \text{id}$, et la deuxième est le morphisme canonique construit
dans la remarque \ref{sites-cohomology-remark-prepare-fancy-base-change}.
\end{remark}




\section{Hom dérivé global}
\label{sites-cohomology-section-global-RHom}

\noindent
Soit $(\Sh(\mathcal{C}), \mathcal{O})$ un topos annelé.
Soient $K, L \in D(\mathcal{O})$.
En utilisant la construction du Hom interne dans la catégorie dérivée, nous
obtenons un objet bien défini
$$
R\Hom_\mathcal{O}(K, L) = R\Gamma(\mathcal{C}, R\SheafHom(K, L))
$$
de $D(\Gamma(\mathcal{C}, \mathcal{O}))$. D'après le
lemme \ref{sites-cohomology-lemma-section-RHom-over-U}, nous avons
$$
H^0(R\Hom_\mathcal{O}(K, L)) = \Hom_{D(\mathcal{O})}(K, L)
$$
et
$$
H^p(R\Hom_\mathcal{O}(K, L)) = \Ext_{D(\mathcal{O})}^p(K, L)
$$
Si $f : (\mathcal{C}', \mathcal{O}') \to (\mathcal{C}, \mathcal{O})$ est un
morphisme de topos annelés, alors il existe un morphisme canonique
$$
R\Hom_\mathcal{O}(K, L) \longrightarrow R\Hom_{\mathcal{O}'}(Lf^*K, Lf^*L)
$$
dans $D(\Gamma(\mathcal{O}))$, obtenu en prenant les sections globales du
morphisme défini dans la remarque \ref{sites-cohomology-remark-prepare-fancy-base-change}.









\section{Image directe exceptionnelle dérivée}
\label{sites-cohomology-section-derived-lower-shriek}

\noindent
Dans cette section, nous étudions des morphismes $g$ de topos annelés pour
lesquels, outre $Lg^*$ et $Rg_*$, il existe aussi un foncteur dérivé $Lg_!$.

\begin{lemma}
\label{sites-cohomology-lemma-pullback-injective-pre-limp}
Soit $u : \mathcal{C} \to \mathcal{D}$ un foncteur continu et cocontinu de
sites, et soit $g : \Sh(\mathcal{C}) \to \Sh(\mathcal{D})$ le morphisme de
topos correspondant. Soit $\mathcal{O}_\mathcal{D}$ un faisceau d'anneaux et
soit $\mathcal{I}$ un $\mathcal{O}_\mathcal{D}$-module injectif. Alors
$H^p(U, g^{-1}\mathcal{I}) = 0$ pour tout $p > 0$ et tout
$U \in \Ob(\mathcal{C})$.
\end{lemma}

\begin{proof}
L'annulation du lemme résulte du lemme
\ref{sites-cohomology-lemma-cech-vanish-collection} si nous pouvons démontrer l'annulation de
tous les groupes de cohomologie de {\v C}ech supérieurs
$\check H^p(\mathcal{U}, g^{-1}\mathcal{I})$
pour tout recouvrement $\mathcal{U} = \{U_i \to U\}$ de $\mathcal{C}$.
Comme $u$ est continu, $u(\mathcal{U}) = \{u(U_i) \to u(U)\}$ est un
recouvrement de $\mathcal{D}$, et
$u(U_{i_0} \times_U \ldots \times_U U_{i_n}) =
u(U_{i_0}) \times_{u(U)} \ldots \times_{u(U)} u(U_{i_n})$.
Ainsi, nous avons
$$
\check H^p(\mathcal{U}, g^{-1}\mathcal{I}) =
\check H^p(u(\mathcal{U}), \mathcal{I})
$$
car $g^{-1} = u^p$ d'après Sites, lemme \ref{sites-lemma-when-shriek}.
Comme $\mathcal{I}$ est un $\mathcal{O}_\mathcal{D}$-module injectif, ces
groupes de cohomologie de {\v C}ech sont nuls, voir le
lemme \ref{sites-cohomology-lemma-injective-module-trivial-cech}.
\end{proof}

\begin{lemma}
\label{sites-cohomology-lemma-existence-derived-lower-shriek}
Soit $u : \mathcal{C} \to \mathcal{D}$ un foncteur continu et cocontinu de
sites, et soit $g : \Sh(\mathcal{C}) \to \Sh(\mathcal{D})$ le morphisme de
topos correspondant. Soit $\mathcal{O}_\mathcal{D}$ un faisceau d'anneaux et
posons
$\mathcal{O}_\mathcal{C} = g^{-1}\mathcal{O}_\mathcal{D}$.
Le foncteur $g_! : \textit{Mod}(\mathcal{O}_\mathcal{C}) \to
\textit{Mod}(\mathcal{O}_\mathcal{D})$
(voir Modules sur les sites, lemme
\ref{sites-modules-lemma-lower-shriek-modules}) possède un foncteur dérivé à
gauche
$$
Lg_! : D(\mathcal{O}_\mathcal{C}) \longrightarrow D(\mathcal{O}_\mathcal{D})
$$
qui est adjoint à gauche de $g^*$. De plus, pour
$U \in \Ob(\mathcal{C})$, nous avons
$$
Lg_!(j_{U!}\mathcal{O}_U) =
g_!j_{U!}\mathcal{O}_U =
j_{u(U)!} \mathcal{O}_{u(U)}.
$$
où $j_{U!}$ et $j_{u(U)!}$ sont les prolongements par zéro associés aux
morphismes de localisation $j_U : \mathcal{C}/U \to \mathcal{C}$ et
$j_{u(U)} : \mathcal{D}/u(U) \to \mathcal{D}$.
\end{lemma}

\begin{proof}
Nous allons utiliser Catégories dérivées, proposition
\ref{derived-proposition-left-derived-exists}, pour construire $Lg_!$.
Il faut pour cela vérifier les hypothèses (1), (2), (3), (4) et (5) de cette
proposition. Tout d'abord, puisque $g_!$ est un adjoint à gauche, il est exact
à droite et commute à toutes les colimites ; la condition (5) est donc
vérifiée. Les conditions (3) et (4) le sont parce que la catégorie des modules
sur un site annelé est une catégorie abélienne de Grothendieck.
Soit $\mathcal{P} \subset \Ob(\textit{Mod}(\mathcal{O}_\mathcal{C}))$ la
collection des $\mathcal{O}_\mathcal{C}$-modules qui sont des sommes directes
de modules de la forme $j_{U!}\mathcal{O}_U$. Notons que
$g_!j_{U!}\mathcal{O}_U = j_{u(U)!} \mathcal{O}_{u(U)}$, voir la
démonstration de Modules sur les sites, lemme
\ref{sites-modules-lemma-lower-shriek-modules}. Tout
$\mathcal{O}_\mathcal{C}$-module est quotient d'un objet de $\mathcal{P}$,
voir Modules sur les sites, lemme
\ref{sites-modules-lemma-module-quotient-flat}. La condition (1) est donc
vérifiée. Il reste à démontrer (2).
Soit $\mathcal{K}^\bullet$ un complexe acyclique borné supérieurement de
$\mathcal{O}_\mathcal{C}$-modules, avec $\mathcal{K}^n \in \mathcal{P}$ pour
tout $n$. Nous devons montrer que $g_!\mathcal{K}^\bullet$ est exact. Pour
cela, il suffit de montrer, pour tout $\mathcal{O}_\mathcal{D}$-module
injectif $\mathcal{I}$, que
$$
\Hom_{D(\mathcal{O}_\mathcal{D})}(
g_!\mathcal{K}^\bullet, \mathcal{I}[n]) = 0
$$
pour tout $n \in \mathbf{Z}$. Comme $\mathcal{I}$ est injectif, nous avons
\begin{align*}
\Hom_{D(\mathcal{O}_\mathcal{D})}(
g_!\mathcal{K}^\bullet, \mathcal{I}[n])
& =
\Hom_{K(\mathcal{O}_\mathcal{D})}(
g_!\mathcal{K}^\bullet, \mathcal{I}[n]) \\
& =
H^n(\Hom_{\mathcal{O}_\mathcal{D}}(
g_!\mathcal{K}^\bullet, \mathcal{I})) \\
& =
H^n(\Hom_{\mathcal{O}_\mathcal{C}}(
\mathcal{K}^\bullet, g^{-1}\mathcal{I}))
\end{align*}
la dernière égalité venant de l'adjonction entre $g_!$ et $g^{-1}$.

\medskip\noindent
L'annulation de ce groupe serait claire si $g^{-1}\mathcal{I}$ était un
$\mathcal{O}_\mathcal{C}$-module injectif. Mais $g^{-1}\mathcal{I}$ n'est
pas nécessairement un $\mathcal{O}_\mathcal{C}$-module injectif, puisque
$g_!$ n'est en général pas exact. Nous
savons néanmoins que
$$
\Ext^p_{\mathcal{O}_\mathcal{C}}(
j_{U!}\mathcal{O}_U, g^{-1}\mathcal{I}) =
H^p(U, g^{-1}\mathcal{I}) = 0 \text{ pour }p \geq 1
$$
La première égalité vient ici de
$\Hom_{\mathcal{O}_\mathcal{C}}(j_{U!}\mathcal{O}_U, \mathcal{H}) =
\mathcal{H}(U)$ par passage aux foncteurs dérivés, et l'annulation de
$H^p(U, g^{-1}\mathcal{I})$ pour $p > 0$ et
$U \in \Ob(\mathcal{C})$ résulte du
lemme \ref{sites-cohomology-lemma-pullback-injective-pre-limp}.
Comme chaque $\mathcal{K}^{-q}$ est une somme directe de modules de la forme
$j_{U!}\mathcal{O}_U$, nous voyons que
$$
\Ext^p_{\mathcal{O}_\mathcal{C}}(\mathcal{K}^{-q}, g^{-1}\mathcal{I}) = 0
\text{ pour }p \geq 1\text{ et tout }q
$$
Utilisons la suite spectrale (voir l'exemple
\ref{sites-cohomology-example-hom-complex-into-sheaf})
$$
E_1^{p, q} = \Ext^q_{\mathcal{O}_\mathcal{C}}(
\mathcal{K}^{-p}, g^{-1}\mathcal{I})
\Rightarrow
\Ext^{p + q}_{\mathcal{O}_\mathcal{C}}(
\mathcal{K}^\bullet, g^{-1}\mathcal{I}) = 0.
$$
Notons que la suite spectrale aboutit à zéro puisque
$\mathcal{K}^\bullet$ est acyclique (il est donc nul dans la catégorie
dérivée et produit des groupes Ext nuls). Par l'annulation des Ext supérieurs
démontrée ci-dessus, les seuls termes non nuls de la page $E_1$ sont
$E_1^{p, 0} = \Hom_{\mathcal{O}_\mathcal{C}}(
\mathcal{K}^{-p}, g^{-1}\mathcal{I})$.
Nous concluons que le complexe
$\Hom_{\mathcal{O}_\mathcal{C}}(
\mathcal{K}^\bullet, g^{-1}\mathcal{I})$
est acyclique, comme voulu.

\medskip\noindent
Ainsi, le foncteur dérivé à gauche $Lg_!$ existe. Il est adjoint à gauche de
$g^{-1} = g^* = Rg^* = Lg^*$, c'est-à-dire que nous avons
\begin{equation}
\label{sites-cohomology-equation-to-prove}
\Hom_{D(\mathcal{O}_\mathcal{C})}(K, g^*L) =
\Hom_{D(\mathcal{O}_\mathcal{D})}(Lg_!K, L)
\end{equation}
d'après Catégories dérivées, lemme
\ref{derived-lemma-derived-adjoint-functors}. Cela achève la démonstration.
\end{proof}

\begin{remark}
\label{sites-cohomology-remark-when-derived-shriek-equal}
Attention ! Soient $u : \mathcal{C} \to \mathcal{D}$, $g$,
$\mathcal{O}_\mathcal{D}$ et $\mathcal{O}_\mathcal{C}$ comme dans le
lemme \ref{sites-cohomology-lemma-existence-derived-lower-shriek}.
En général, le diagramme
$$
\xymatrix{
D(\mathcal{O}_\mathcal{C}) \ar[r]_{Lg_!} \ar[d]_{forget} &
D(\mathcal{O}_\mathcal{D}) \ar[d]^{forget} \\
D(\mathcal{C}) \ar[r]^{Lg^{Ab}_!} &
D(\mathcal{D})
}
$$
ne commute {\bf pas}, où $Lg_!^{Ab}$ est le foncteur construit dans le
lemme \ref{sites-cohomology-lemma-existence-derived-lower-shriek}, mais en utilisant le
faisceau constant $\mathbf{Z}$ comme faisceau structural sur
$\mathcal{C}$ et $\mathcal{D}$. En général, on n'a même pas
$g_! = g_!^{Ab}$ (voir Modules sur les sites, remarque
\ref{sites-modules-remark-when-shriek-equal}), mais ce phénomène
{\bf peut se produire même si $g_! = g_!^{Ab}$} ! En effet, la construction
de $Lg_!$ dans la démonstration du
lemme \ref{sites-cohomology-lemma-existence-derived-lower-shriek} montre que $Lg_!$ coïncide
avec $Lg_!^{\textit{Ab}}$ si et seulement si les morphismes canoniques
$$
Lg^{Ab}_!j_{U!}\mathcal{O}_U \longrightarrow j_{u(U)!}\mathcal{O}_{u(U)}
$$
sont des isomorphismes dans $D(\mathcal{D})$ pour tout objet $U$ de
$\mathcal{C}$. En général, tout ce que l'on peut dire est qu'il existe une
transformation naturelle
$$
Lg_!^{Ab} \circ forget \longrightarrow forget \circ Lg_!
$$
\end{remark}

\begin{lemma}
\label{sites-cohomology-lemma-pullback-injective-limp}
Soit $u : \mathcal{C} \to \mathcal{D}$ un foncteur continu et cocontinu de
sites, et soit $g : \Sh(\mathcal{C}) \to \Sh(\mathcal{D})$ le morphisme de
topos correspondant. Soit $\mathcal{O}_\mathcal{D}$ un faisceau d'anneaux et
soit $\mathcal{I}$ un $\mathcal{O}_\mathcal{D}$-module injectif. Si
$g_!^{Sh} : \Sh(\mathcal{C}) \to \Sh(\mathcal{D})$
commute aux produits fibrés\footnote{C'est le cas si $\mathcal{C}$ possède
des limites connexes finies et si $u$ y commute, voir Sites, lemme
\ref{sites-lemma-preserve-equalizers}.}, alors $g^{-1}\mathcal{I}$ est
totalement acyclique.
\end{lemma}

\begin{proof}
Nous utiliserons le critère du lemme \ref{sites-cohomology-lemma-characterize-limp}.
La condition (1) est vérifiée d'après le
lemme \ref{sites-cohomology-lemma-pullback-injective-pre-limp}.
Soit $K' \to K$ un morphisme surjectif de faisceaux d'ensembles sur
$\mathcal{C}$. Comme $g_!^{Sh}$ est un adjoint à gauche, le morphisme
$g_!^{Sh}K' \to g_!^{Sh}K$ est surjectif. Observons que
\begin{align*}
H^0(K' \times_K \ldots \times_K K', g^{-1}\mathcal{I}) 
& =
H^0(g_!^{Sh}(K' \times_K \ldots \times_K K'), \mathcal{I}) \\
& =
H^0(g_!^{Sh}K' \times_{g_!^{Sh}K} \ldots \times_{g_!^{Sh}K} g_!^{Sh}K',
\mathcal{I})
\end{align*}
d'après notre hypothèse sur $g_!^{Sh}$. Comme $\mathcal{I}$ est un module
injectif, il est totalement acyclique par le
lemme \ref{sites-cohomology-lemma-direct-image-injective-sheaf} (appliqué à l'identité).
Nous pouvons donc utiliser la réciproque du lemme
\ref{sites-cohomology-lemma-characterize-limp} pour voir que le complexe
$$
0 \to H^0(K, g^{-1}\mathcal{I}) \to H^0(K', g^{-1}\mathcal{I}) \to
H^0(K' \times_K K', g^{-1}\mathcal{I}) \to \ldots
$$
est exact, comme voulu.
\end{proof}

\begin{lemma}
\label{sites-cohomology-lemma-pullback-same-cohomology}
Soit $u : \mathcal{C} \to \mathcal{D}$ un foncteur continu et cocontinu de
sites, et soit $g : \Sh(\mathcal{C}) \to \Sh(\mathcal{D})$ le morphisme de
topos correspondant. Soit $U \in \Ob(\mathcal{C})$.
\begin{enumerate}
\item Pour $M$ dans $D(\mathcal{D})$, nous avons
$R\Gamma(U, g^{-1}M) = R\Gamma(u(U), M)$.
\item Si $\mathcal{O}_\mathcal{D}$ est un faisceau d'anneaux et
$\mathcal{O}_\mathcal{C} = g^{-1}\mathcal{O}_\mathcal{D}$, alors, pour $M$
dans $D(\mathcal{O}_\mathcal{D})$, nous avons
$R\Gamma(U, g^*M) = R\Gamma(u(U), M)$.
\end{enumerate}
\end{lemma}

\begin{proof}
Dans le cas borné inférieurement, on obtient (1) et (2) en représentant $K$
par un complexe borné inférieurement d'injectifs, puis en utilisant le
lemme \ref{sites-cohomology-lemma-pullback-injective-pre-limp} ainsi que le lemme d'acyclicité
de Leray. Dans le cas non borné, remarquons d'abord que (1) est un cas
particulier de (2). Pour (2), nous pouvons utiliser
$$
R\Gamma(U, g^*M) =
R\Hom_{\mathcal{O}_\mathcal{C}}(j_{U!}\mathcal{O}_U, g^*M) =
R\Hom_{\mathcal{O}_\mathcal{D}}(j_{u(U)!}\mathcal{O}_{u(U)}, M) =
R\Gamma(u(U), M)
$$
où l'égalité du milieu résulte du
lemme \ref{sites-cohomology-lemma-existence-derived-lower-shriek}.
\end{proof}

\begin{lemma}
\label{sites-cohomology-lemma-special-square-cocontinuous}
Supposons donné un diagramme commutatif
$$
\xymatrix{
(\Sh(\mathcal{C}'), \mathcal{O}_{\mathcal{C}'})
\ar[r]_{(g', (g')^\sharp)} \ar[d]_{(f', (f')^\sharp)} &
(\Sh(\mathcal{C}), \mathcal{O}_\mathcal{C}) \ar[d]^{(f, f^\sharp)} \\
(\Sh(\mathcal{D}'), \mathcal{O}_{\mathcal{D}'}) \ar[r]^{(g, g^\sharp)} &
(\Sh(\mathcal{D}), \mathcal{O}_\mathcal{D})
}
$$
de topos annelés. Supposons que :
\begin{enumerate}
\item $f$, $f'$, $g$ et $g'$ correspondent aux foncteurs cocontinus
$u$, $u'$, $v$ et $v'$ comme dans Sites, lemme
\ref{sites-lemma-cocontinuous-morphism-topoi},
\item $v \circ u' = u \circ v'$,
\item $v$ et $v'$ sont continus et cocontinus,
\item pour tout objet $V'$ de $\mathcal{D}'$, le foncteur
${}^{u'}_{V'}\mathcal{I} \to {}^{\ \ \ u}_{v(V')}\mathcal{I}$
défini par $v$ est cofinal,
\item $g^{-1}\mathcal{O}_{\mathcal{D}} = \mathcal{O}_{\mathcal{D}'}$
et $(g')^{-1}\mathcal{O}_{\mathcal{C}} = \mathcal{O}_{\mathcal{C}'}$, et
\item $g'_! : \textit{Ab}(\mathcal{C}') \to \textit{Ab}(\mathcal{C})$
est exact\footnote{C'est le cas si les produits fibrés et les égaliseurs
existent dans $\mathcal{C}'$ et si $v'$ y commute, voir Modules sur les
sites, lemme \ref{sites-modules-lemma-exactness-lower-shriek}.}.
\end{enumerate}
Alors $Rf'_* \circ (g')^* = g^* \circ Rf_*$ comme foncteurs
$D(\mathcal{O}_\mathcal{C}) \to D(\mathcal{O}_{\mathcal{D}'})$.
\end{lemma}

\begin{proof}
Nous avons $g^* = Lg^* = g^{-1}$ et
$(g')^* = L(g')^* = (g')^{-1}$ d'après la condition (5).
Par le lemme \ref{sites-cohomology-lemma-modules-abelian-unbounded}, il suffit de démontrer le
résultat dans la catégorie dérivée $D(\mathcal{C})$ des faisceaux abéliens.
Choisissons un objet $K \in D(\mathcal{C})$ et soit
$\mathcal{I}^\bullet$ un complexe K-injectif de faisceaux abéliens sur
$\mathcal{C}$ représentant $K$. D'après Catégories dérivées, lemme
\ref{derived-lemma-adjoint-preserve-K-injectives}, et l'hypothèse (6),
$(g')^{-1}\mathcal{I}^\bullet$ est un complexe K-injectif de faisceaux
abéliens sur $\mathcal{C}'$. Par Modules sur les sites, lemme
\ref{sites-modules-lemma-special-square-cocontinuous}
nous obtenons
$f'_*(g')^{-1}\mathcal{I}^\bullet = g^{-1}f_*\mathcal{I}^\bullet$.
Comme $f_*\mathcal{I}^\bullet$ représente $Rf_*K$ et que
$f'_*(g')^{-1}\mathcal{I}^\bullet$ représente $Rf'_*(g')^{-1}K$, nous
concluons.
\end{proof}

\begin{lemma}
\label{sites-cohomology-lemma-special-square-continuous}
Considérons un diagramme commutatif
$$
\xymatrix{
(\Sh(\mathcal{C}'), \mathcal{O}_{\mathcal{C}'})
\ar[r]_{(g', (g')^\sharp)} \ar[d]_{(f', (f')^\sharp)} &
(\Sh(\mathcal{C}), \mathcal{O}_\mathcal{C}) \ar[d]^{(f, f^\sharp)} \\
(\Sh(\mathcal{D}'), \mathcal{O}_{\mathcal{D}'}) \ar[r]^{(g, g^\sharp)} &
(\Sh(\mathcal{D}), \mathcal{O}_\mathcal{D})
}
$$
de topos annelés et supposons donnés des foncteurs
$$
\xymatrix{
\mathcal{C}' \ar[r]_{v'} &
\mathcal{C} \\
\mathcal{D}' \ar[r]^v \ar[u]^{u'} &
\mathcal{D} \ar[u]_u
}
$$
tels que (avec les notations de Sites, sections
\ref{sites-section-morphism-sites} et
\ref{sites-section-cocontinuous-morphism-topoi}) :
\begin{enumerate}
\item $u$ et $u'$ sont continus et donnent naissance aux morphismes
$f$ et $f'$,
\item $v$ et $v'$ sont cocontinus et donnent naissance aux morphismes
$g$ et $g'$,
\item $u \circ v = v' \circ u'$,
\item $v$ et $v'$ sont continus et cocontinus, et
\item $g^{-1}\mathcal{O}_{\mathcal{D}} = \mathcal{O}_{\mathcal{D}'}$
et $(g')^{-1}\mathcal{O}_{\mathcal{C}} = \mathcal{O}_{\mathcal{C}'}$.
\end{enumerate}
Alors $Rf'_* \circ (g')^* = g^* \circ Rf_*$ comme foncteurs
$D^+(\mathcal{O}_\mathcal{C}) \to D^+(\mathcal{O}_{\mathcal{D}'})$.
Si, de plus,
\begin{enumerate}
\item[(6)] $g'_! : \textit{Ab}(\mathcal{C}') \to \textit{Ab}(\mathcal{C})$
est exact\footnote{C'est le cas si les produits fibrés et les égaliseurs
existent dans $\mathcal{C}'$ et si $v'$ y commute, voir Modules sur les
sites, lemme \ref{sites-modules-lemma-exactness-lower-shriek}.},
\end{enumerate}
alors $Rf'_* \circ (g')^* = g^* \circ Rf_*$ comme foncteurs
$D(\mathcal{O}_\mathcal{C}) \to D(\mathcal{O}_{\mathcal{D}'})$.
\end{lemma}

\begin{proof}
Nous avons $g^* = Lg^* = g^{-1}$ et
$(g')^* = L(g')^* = (g')^{-1}$ d'après la condition (5).
Par le lemme \ref{sites-cohomology-lemma-modules-abelian-unbounded}, il suffit de démontrer le
résultat dans la catégorie dérivée $D^+(\mathcal{C})$ ou $D(\mathcal{C})$
des faisceaux abéliens.

\medskip\noindent
Choisissons un objet $K \in D^+(\mathcal{C})$. Soit
$\mathcal{I}^\bullet$ un complexe borné inférieurement de faisceaux abéliens
injectifs sur $\mathcal{C}$ représentant $K$. D'après le
lemme \ref{sites-cohomology-lemma-pullback-injective-pre-limp}, nous avons
$H^p(U', (g')^{-1}\mathcal{I}^q) = 0$ pour tout $p > 0$, tout $q$ et tout
$U' \in \Ob(\mathcal{C}')$. Rappelons que
$R^pf'_*(g')^{-1}\mathcal{I}^q$ est le faisceau associé au préfaisceau
$V' \mapsto H^p(u'(V'), (g')^{-1}\mathcal{I}^q)$, voir le
lemme \ref{sites-cohomology-lemma-higher-direct-images}. Ainsi,
$(g')^{-1}\mathcal{I}^q$ est acyclique à droite pour le foncteur $f'_*$.
Le lemme d'acyclicité de Leray (Catégories dérivées, lemme
\ref{derived-lemma-leray-acyclicity}) montre alors que
$f'_*(g')^*\mathcal{I}^\bullet$ représente $Rf'_*(g')^{-1}K$.
Par Modules sur les sites, lemme
\ref{sites-modules-lemma-special-square-continuous}
nous obtenons
$f'_*(g')^{-1}\mathcal{I}^\bullet = g^{-1}f_*\mathcal{I}^\bullet$.
Comme $g^{-1}f_*\mathcal{I}^\bullet$ représente $g^{-1}Rf_*K$, nous
concluons.

\medskip\noindent
Choisissons un objet $K \in D(\mathcal{C})$. Soit $\mathcal{I}^\bullet$ un
complexe K-injectif de faisceaux abéliens sur $\mathcal{C}$ représentant
$K$. D'après Catégories dérivées, lemme
\ref{derived-lemma-adjoint-preserve-K-injectives}, et l'hypothèse (6),
$(g')^{-1}\mathcal{I}^\bullet$ est un complexe K-injectif de faisceaux
abéliens sur $\mathcal{C}'$. Par Modules sur les sites, lemme
\ref{sites-modules-lemma-special-square-continuous}
nous obtenons
$f'_*(g')^{-1}\mathcal{I}^\bullet = g^{-1}f_*\mathcal{I}^\bullet$.
Comme $f_*\mathcal{I}^\bullet$ représente $Rf_*K$ et que
$f'_*(g')^{-1}\mathcal{I}^\bullet$ représente $Rf'_*(g')^{-1}K$, nous
concluons.
\end{proof}








\section{Image directe exceptionnelle dérivée pour les catégories fibrées}
\label{sites-cohomology-section-derived-lower-shriek-fibred}

\noindent
Dans cette section, nous étudions certains cas particuliers de la situation
examinée dans la section \ref{sites-cohomology-section-derived-lower-shriek}.
Nous nous assurons que les foncteurs image directe exceptionnelle sur les modules
et sur les faisceaux de groupes abéliens coïncident. Nous conseillons au lecteur de sauter
cette section en première lecture.

\begin{situation}
\label{sites-cohomology-situation-fibred-category}
Soient $(\mathcal{D}, \mathcal{O}_\mathcal{D})$ un site annelé
et $p : \mathcal{C} \to \mathcal{D}$ une catégorie fibrée. Nous munissons
$\mathcal{C}$ de la topologie héritée de $\mathcal{D}$
(Champs, section \ref{stacks-section-topology}). Nous notons
$\pi : \Sh(\mathcal{C}) \to \Sh(\mathcal{D})$ le morphisme de
topos associé à $p$
(Champs, lemme \ref{stacks-lemma-topology-inherited-functorial}).
Nous posons $\mathcal{O}_\mathcal{C} = \pi^{-1}\mathcal{O}_\mathcal{D}$,
de sorte que nous obtenons un morphisme de topos annelés
$$
\pi :
(\Sh(\mathcal{C}), \mathcal{O}_\mathcal{C})
\longrightarrow
(\Sh(\mathcal{D}), \mathcal{O}_\mathcal{D})
$$
\end{situation}

\begin{lemma}
\label{sites-cohomology-lemma-fibred-category-with-object}
On reprend les hypothèses et les notations de la situation \ref{sites-cohomology-situation-fibred-category}.
Pour $U \in \Ob(\mathcal{C})$, considérons le morphisme induit
de topos
$$
\pi_U : \Sh(\mathcal{C}/U) \longrightarrow \Sh(\mathcal{D}/p(U))
$$
Alors il existe un morphisme de topos
$$
\sigma : \Sh(\mathcal{D}/p(U)) \to \Sh(\mathcal{C}/U)
$$
tel que $\pi_U \circ \sigma = \text{id}$ et $\sigma^{-1} = \pi_{U, *}$.
\end{lemma}

\begin{proof}
Remarquons que $\pi_U$ est la restriction de $\pi$ aux localisations, voir
Sites, lemme \ref{sites-lemma-localize-cocontinuous}.
Pour un objet $V \to p(U)$ de $\mathcal{D}/p(U)$, notons
$V \times_{p(U)} U \to U$ le morphisme fortement cartésien de $\mathcal{C}$
au-dessus de $\mathcal{D}$, qui existe puisque $p$ est une catégorie fibrée.
Le foncteur
$$
v : \mathcal{D}/p(U) \to \mathcal{C}/U,\quad
V/p(U) \mapsto V \times_{p(U)} U/U
$$
est continu par définition de la topologie sur $\mathcal{C}$.
De plus, il est adjoint à droite de $p$ par définition des morphismes fortement
cartésiens. Nous sommes donc dans la situation examinée dans
Sites, section \ref{sites-section-cocontinuous-adjoint},
et nous voyons que le faisceau $\pi_{U, *}\mathcal{F}$
est égal à $V \mapsto \mathcal{F}(V \times_{p(U)} U)$
(voir en particulier Sites, lemme
\ref{sites-lemma-have-functor-other-way-morphism}).

\medskip\noindent
Mais nous disposons ici d'un résultat plus fort. En effet, le foncteur $v$
est aussi cocontinu (car tous les morphismes des recouvrements de $\mathcal{C}$ 
sont fortement cartésiens). Ainsi, $v$ définit un morphisme $\sigma$ comme
indiqué dans le lemme. L'égalité $\sigma^{-1} = \pi_{U, *}$
résulte immédiatement de la définition. Puisque $\pi_U^{-1}\mathcal{G}$
est donné par la règle $U'/U \mapsto \mathcal{G}(p(U')/p(U))$,
il s'ensuit que $\sigma^{-1} \circ \pi_U^{-1} = \text{id}$,
ce qui prouve l'égalité
$\pi_U \circ \sigma = \text{id}$.
\end{proof}

\begin{situation}
\label{sites-cohomology-situation-morphism-fibred-categories}
Soit $(\mathcal{D}, \mathcal{O}_\mathcal{D})$ un site annelé.
Soit $u : \mathcal{C}' \to \mathcal{C}$ un $1$-morphisme de catégories
fibrées au-dessus de $\mathcal{D}$
(Catégories, définition \ref{categories-definition-fibred-categories-over-C}).
Munissons $\mathcal{C}$ et $\mathcal{C}'$ de leurs topologies héritées
(Champs, définition \ref{stacks-definition-topology-inherited})
et soient
$\pi : \Sh(\mathcal{C}) \to \Sh(\mathcal{D})$,
$\pi' : \Sh(\mathcal{C}') \to \Sh(\mathcal{D})$, et
$g : \Sh(\mathcal{C}') \to \Sh(\mathcal{C})$
les morphismes de topos correspondants
(Champs, lemme \ref{stacks-lemma-topology-inherited-functorial}).
Posons $\mathcal{O}_\mathcal{C} = \pi^{-1}\mathcal{O}_\mathcal{D}$
et $\mathcal{O}_{\mathcal{C}'} = (\pi')^{-1}\mathcal{O}_\mathcal{D}$.
Remarquons que $g^{-1}\mathcal{O}_\mathcal{C} = \mathcal{O}_{\mathcal{C}'}$,
de sorte que
$$
\xymatrix{
(\Sh(\mathcal{C}'), \mathcal{O}_{\mathcal{C}'}) \ar[rd]_{\pi'} \ar[rr]_g & &
(\Sh(\mathcal{C}), \mathcal{O}_\mathcal{C}) \ar[ld]^\pi \\
& (\Sh(\mathcal{D}), \mathcal{O}_\mathcal{D})
}
$$
est un diagramme commutatif de morphismes de topos annelés.
\end{situation}

\begin{lemma}
\label{sites-cohomology-lemma-morphism-fibred-categories-with-object}
On reprend les hypothèses et les notations de la
situation \ref{sites-cohomology-situation-morphism-fibred-categories}.
Pour $U' \in \Ob(\mathcal{C}')$, posons $U = u(U')$ et $V = p'(U')$, puis
considérons les morphismes induits de topos annelés
$$
\xymatrix{
(\Sh(\mathcal{C}'/U'), \mathcal{O}_{U'}) \ar[rd]_{\pi'_{U'}} \ar[rr]_{g'} & &
(\Sh(\mathcal{C}/U), \mathcal{O}_U) \ar[ld]^{\pi_U} \\
& (\Sh(\mathcal{D}/V), \mathcal{O}_V)
}
$$
Alors il existe un morphisme de topos
$$
\sigma' : \Sh(\mathcal{D}/V) \to \Sh(\mathcal{C}'/U'),
$$
tel que, en posant $\sigma = g' \circ \sigma'$, nous ayons
$\pi'_{U'} \circ \sigma' = \text{id}$, $\pi_U \circ \sigma = \text{id}$,
$(\sigma')^{-1} = \pi'_{U', *}$, et $\sigma^{-1} = \pi_{U, *}$.
\end{lemma}

\begin{proof}
Soit $v' : \mathcal{D}/V \to \mathcal{C}'/U'$ le foncteur construit
dans la démonstration du lemme \ref{sites-cohomology-lemma-fibred-category-with-object} à partir
de $p' : \mathcal{C}' \to \mathcal{D}$ et de l'objet $U'$.
Comme $u$ est un $1$-morphisme de catégories fibrées au-dessus de $\mathcal{D}$,
il transforme les morphismes fortement cartésiens en morphismes fortement cartésiens ;
ainsi, le foncteur $v = u \circ v'$ est le foncteur de
la démonstration du lemme \ref{sites-cohomology-lemma-fibred-category-with-object}
relatif à $p : \mathcal{C} \to \mathcal{D}$ et à $U$. Le présent lemme
résulte donc de ce lemme.
\end{proof}

\begin{lemma}
\label{sites-cohomology-lemma-properties-lower-shriek-fibred-category}
On reprend les hypothèses et les notations de la
situation \ref{sites-cohomology-situation-morphism-fibred-categories}.
\begin{enumerate}
\item Il existe des adjoints à gauche
$g_! : \textit{Mod}(\mathcal{O}_{\mathcal{C}'}) \to
\textit{Mod}(\mathcal{O}_\mathcal{C})$ et
$g_!^{\textit{Ab}} : \textit{Ab}(\mathcal{C}') \to \textit{Ab}(\mathcal{C})$
de $g^* = g^{-1}$ sur les modules et sur les faisceaux abéliens.
\item Le diagramme
$$
\xymatrix{
\textit{Mod}(\mathcal{O}_{\mathcal{C}'}) \ar[d] \ar[r]_{g_!} &
\textit{Mod}(\mathcal{O}_\mathcal{C}) \ar[d] \\
\textit{Ab}(\mathcal{C}') \ar[r]^{g_!^{\textit{Ab}}} &
\textit{Ab}(\mathcal{C})
}
$$
est commutatif.
\item Il existe des adjoints à gauche
$Lg_! : D(\mathcal{O}_{\mathcal{C}'}) \to D(\mathcal{O}_\mathcal{C})$
et
$Lg_!^{\textit{Ab}} : D(\mathcal{C}') \to D(\mathcal{C})$
de $g^* = g^{-1}$ sur les catégories dérivées de modules et de faisceaux abéliens.
\item Le diagramme
$$
\xymatrix{
D(\mathcal{O}_{\mathcal{C}'}) \ar[d] \ar[r]_{Lg_!} &
D(\mathcal{O}_\mathcal{C}) \ar[d] \\
D(\mathcal{C}') \ar[r]^{Lg_!^{\textit{Ab}}} &
D(\mathcal{C})
}
$$
est commutatif.
\end{enumerate}
\end{lemma}

\begin{proof}
Le foncteur $u$ est continu et cocontinu
(Champs, lemme \ref{stacks-lemma-topology-inherited-functorial}).
L'existence des foncteurs $g_!$, $g_!^{\textit{Ab}}$,
$Lg_!$ et $Lg_!^{\textit{Ab}}$ résulte donc de
Modules sur les sites, sections
\ref{sites-modules-section-exactness-lower-shriek} et
\ref{sites-modules-section-lower-shriek-modules}
et de la
section \ref{sites-cohomology-section-derived-lower-shriek}.

\medskip\noindent
Pour démontrer (2), il suffit de montrer que le morphisme canonique
$$
g_!^{\textit{Ab}}j_{U'!}\mathcal{O}_{U'} \to j_{u(U')!}\mathcal{O}_{u(U')}
$$
est un isomorphisme pour tout objet $U'$ de $\mathcal{C}'$, voir
Modules sur les sites, remarque \ref{sites-modules-remark-when-shriek-equal}.
De même, pour démontrer (4), il suffit de montrer que le morphisme canonique
$$
Lg_!^{\textit{Ab}}j_{U'!}\mathcal{O}_{U'} \to j_{u(U')!}\mathcal{O}_{u(U')}
$$
est un isomorphisme dans $D(\mathcal{C})$ pour tout objet $U'$ de
$\mathcal{C}'$, voir la remarque \ref{sites-cohomology-remark-when-derived-shriek-equal}.
Cela entraînera aussi la formule précédente ; c'est donc ce que nous allons montrer.

\medskip\noindent
Nous utiliserons le fait que, pour un morphisme de localisation $j$, les
foncteurs $j_!$ et $j_!^{\textit{Ab}}$ coïncident (voir
Modules sur les sites, remarque \ref{sites-modules-remark-localize-shriek-equal})
et que $j_!$ est exact
(Modules sur les sites, lemme \ref{sites-modules-lemma-extension-by-zero-exact}).
Adoptons les notations du
lemme \ref{sites-cohomology-lemma-morphism-fibred-categories-with-object}.
Puisque $Lg_!^{\textit{Ab}} \circ j_{U'!} = j_{U!} \circ L(g')^{\textit{Ab}}_!$
(par la commutativité du diagramme de Sites, lemme \ref{sites-lemma-localize-cocontinuous},
et l'unicité des foncteurs adjoints), il suffit de prouver que
$L(g')^{\textit{Ab}}_!\mathcal{O}_{U'} = \mathcal{O}_U$. D'après les
résultats du
lemme \ref{sites-cohomology-lemma-morphism-fibred-categories-with-object},
nous avons, pour tout objet $E$ de $D(\mathcal{C}/u(U'))$, la suite
d'égalités suivante
\begin{align*}
\Hom_{D(\mathcal{C}/U)}(L(g')_!^{\textit{Ab}}\mathcal{O}_{U'}, E)
& =
\Hom_{D(\mathcal{C}'/U')}(\mathcal{O}_{U'}, (g')^{-1}E) \\
& =
\Hom_{D(\mathcal{C}'/U')}((\pi'_{U'})^{-1}\mathcal{O}_V, (g')^{-1}E) \\
& =
\Hom_{D(\mathcal{D}/V)}(\mathcal{O}_V, R\pi'_{U', *}(g')^{-1}E) \\
& =
\Hom_{D(\mathcal{D}/V)}(\mathcal{O}_V, (\sigma')^{-1}(g')^{-1}E) \\
& =
\Hom_{D(\mathcal{D}/V)}(\mathcal{O}_V, \sigma^{-1}E) \\
& =
\Hom_{D(\mathcal{D}/V)}(\mathcal{O}_V, \pi_{U, *}E) \\
& =
\Hom_{D(\mathcal{C}/U)}(\pi_U^{-1}\mathcal{O}_V, E) \\
& =
\Hom_{D(\mathcal{C}/U)}(\mathcal{O}_U, E)
\end{align*}
Le lemme de Yoneda permet de conclure.
\end{proof}

\begin{remark}
\label{sites-cohomology-remark-fibred-category}
On reprend les hypothèses et les notations de la
situation \ref{sites-cohomology-situation-fibred-category}. Remarquons qu'en posant
$\mathcal{C}' = \mathcal{D}$ et en prenant pour $u$ le foncteur structural
de $\mathcal{C}$, on obtient une situation comme dans la
situation \ref{sites-cohomology-situation-morphism-fibred-categories}. Le
lemme \ref{sites-cohomology-lemma-properties-lower-shriek-fibred-category} nous donne donc les
foncteurs $\pi_!$, $\pi_!^{\textit{Ab}}$, $L\pi_!$ et
$L\pi_!^{\textit{Ab}}$ tels que
$forget \circ \pi_! = \pi_!^{\textit{Ab}} \circ forget$ et
$forget \circ L\pi_! = L\pi_!^{\textit{Ab}} \circ forget$.
\end{remark}

\begin{remark}
\label{sites-cohomology-remark-morphism-fibred-categories}
On reprend les hypothèses et les notations de la
situation \ref{sites-cohomology-situation-morphism-fibred-categories}.
Soit $\mathcal{F}$ un faisceau abélien sur $\mathcal{C}$, soit
$\mathcal{F}'$ un faisceau abélien sur $\mathcal{C}'$, et soit
$t : \mathcal{F}' \to g^{-1}\mathcal{F}$ un morphisme.
Nous obtenons alors un morphisme canonique
$$
L\pi'_!(\mathcal{F}') \longrightarrow L\pi_!(\mathcal{F})
$$
en utilisant l'adjoint $g_!\mathcal{F}' \to \mathcal{F}$ de $t$, le
morphisme $Lg_!(\mathcal{F}') \to g_!\mathcal{F}'$ et l'égalité
$L\pi'_! = L\pi_! \circ Lg_!$.
\end{remark}

\begin{lemma}
\label{sites-cohomology-lemma-compute-pi-shriek}
On reprend les hypothèses et les notations de la
situation \ref{sites-cohomology-situation-fibred-category}.
Pour $\mathcal{F}$ dans $\textit{Ab}(\mathcal{C})$, le faisceau
$\pi_!\mathcal{F}$ est le faisceau associé au préfaisceau
$$
V \longmapsto \colim_{\mathcal{C}_V^{opp}} \mathcal{F}|_{\mathcal{C}_V}
$$
avec les applications de restriction indiquées dans la démonstration.
\end{lemma}

\begin{proof}
Notons $\mathcal{H}$ la règle du lemme. Pour un morphisme
$h : V' \to V$ de $\mathcal{D}$, il existe un foncteur d'image inverse
$h^* : \mathcal{C}_V \to \mathcal{C}_{V'}$ des catégories fibres
(Catégories, définition
\ref{categories-definition-pullback-functor-fibred-category}).
De plus, pour $U \in \Ob(\mathcal{C}_V)$, il existe un morphisme fortement
cartésien $h^*U \to U$ au-dessus de $h$. La restriction le long de ces
morphismes fortement cartésiens définit une transformation de foncteurs
$$
\mathcal{F}|_{\mathcal{C}_V}
\longrightarrow
\mathcal{F}|_{\mathcal{C}_{V'}} \circ h^*.
$$
et donc un morphisme $\mathcal{H}(V) \to \mathcal{H}(V')$ entre les
colimites, voir Catégories, lemme
\ref{categories-lemma-functorial-colimit}.

\medskip\noindent
Pour démontrer le lemme, nous montrons que
$$
\Mor_{\textit{PSh}(\mathcal{D})}(\mathcal{H}, \mathcal{G}) =
\Mor_{\Sh(\mathcal{C})}(\mathcal{F}, \pi^{-1}\mathcal{G})
$$
pour tout faisceau $\mathcal{G}$ sur $\mathcal{C}$. Un élément du membre de
gauche est un système compatible de morphismes
$\mathcal{F}(U) \to \mathcal{G}(p(U))$ pour tout $U$ de $\mathcal{C}$.
Comme $\pi^{-1}\mathcal{G}(U) = \mathcal{G}(p(U))$ par notre choix de la
topologie sur $\mathcal{C}$, il en est de même pour le membre de droite, ce
qui conclut.
\end{proof}





\section{Homologie d'une catégorie}
\label{sites-cohomology-section-homology}

\noindent
Dans le cas d'une catégorie au-dessus d'un point, nous appellerons les
foncteurs dérivés à gauche d'image directe exceptionnelle les foncteurs
d'homologie.

\begin{example}[Catégorie au-dessus d'un point]
\label{sites-cohomology-example-category-to-point}
Soit $\mathcal{C}$ une catégorie. Munissons $\mathcal{C}$ de la topologie
chaotique (Sites, exemple \ref{sites-example-indiscrete}). Les préfaisceaux
et les faisceaux coïncident alors sur $\mathcal{C}$.
Le foncteur $p : \mathcal{C} \to *$, où $*$ est la catégorie avec un seul
objet et un seul morphisme, est cocontinu et continu. Soit
$\pi : \Sh(\mathcal{C}) \to \Sh(*)$ le morphisme de topos correspondant.
Soit $B$ un anneau. Nous munissons $*$ du faisceau d'anneaux $B$ et
$\mathcal{C}$ de $\mathcal{O}_\mathcal{C} = \pi^{-1}B$, que nous noterons
$\underline{B}$. De cette manière,
$$
\pi : (\Sh(\mathcal{C}), \underline{B}) \to (\Sh(*), B)
$$
est un exemple de la situation \ref{sites-cohomology-situation-fibred-category}.
Par la remarque \ref{sites-cohomology-remark-fibred-category}, il n'est pas nécessaire de
distinguer $\pi_!$ sur les modules et sur les faisceaux abéliens. Le
lemme \ref{sites-cohomology-lemma-compute-pi-shriek} montre que
$\pi_!\mathcal{F} = \colim_{\mathcal{C}^{opp}} \mathcal{F}$.
Ainsi, $L_n\pi_!$ est le $n$-ième foncteur dérivé à gauche du foncteur de
colimite. Dans la suite, nous écrivons
$$
H_n(\mathcal{C}, \mathcal{F}) = L_n\pi_!(\mathcal{F})
$$
et nous appelons ceci le {\it $n$-ième groupe d'homologie de $\mathcal{F}$}
sur $\mathcal{C}$.
\end{example}

\begin{example}[Calcul de l'homologie]
\label{sites-cohomology-example-left-derived-colimits}
Dans l'exemple \ref{sites-cohomology-example-category-to-point}, on peut calculer les
foncteurs $H_n(\mathcal{C}, -)$ comme suit. Soit
$\mathcal{F} \in \Ob(\textit{Ab}(\mathcal{C}))$. Considérons le complexe de
chaînes
$$
K_\bullet(\mathcal{F}) :
\ \ldots \to
\bigoplus\nolimits_{U_2 \to U_1 \to U_0} \mathcal{F}(U_0)
\to
\bigoplus\nolimits_{U_1 \to U_0} \mathcal{F}(U_0)
\to
\bigoplus\nolimits_{U_0} \mathcal{F}(U_0)
$$
où les morphismes de transition sont donnés par
$$
(U_2 \to U_1 \to U_0, s)
\longmapsto
(U_1 \to U_0, s) - (U_2 \to U_0, s) + (U_2 \to U_1, s|_{U_1})
$$
et de même dans les autres degrés. Par construction,
$$
H_0(\mathcal{C}, \mathcal{F}) =
\colim_{\mathcal{C}^{opp}} \mathcal{F} =
H_0(K_\bullet(\mathcal{F})),
$$
voir Catégories, lemme
\ref{categories-lemma-colimits-coproducts-coequalizers}.
La construction de $K_\bullet(\mathcal{F})$ est fonctorielle en
$\mathcal{F}$ et transforme les suites exactes courtes de
$\textit{Ab}(\mathcal{C})$ en suites exactes courtes de complexes. La suite
de foncteurs $\mathcal{F} \mapsto H_n(K_\bullet(\mathcal{F}))$ forme donc un
$\delta$-foncteur, voir Homologie, définition
\ref{homology-definition-cohomological-delta-functor} et lemme
\ref{homology-lemma-long-exact-sequence-cochain}.
Pour $\mathcal{F} = j_{U!}\mathbf{Z}_U$, le complexe
$K_\bullet(\mathcal{F})$ est celui associé au $\mathbf{Z}$-module libre sur
l'ensemble simplicial $X_\bullet$ dont les termes sont
$$
X_n = \coprod\nolimits_{U_n \to \ldots \to U_1 \to U_0}
\Mor_\mathcal{C}(U_0, U)
$$
Cet ensemble simplicial est homotopiquement équivalent à l'ensemble
simplicial constant sur un singleton $\{*\}$. En effet, le morphisme
$X_\bullet \to \{*\}$ est évident, le morphisme $\{*\} \to X_n$ envoie $*$
sur $(U \to \ldots \to U, \text{id}_U)$, et les morphismes
$$
h_{n, i} : X_n \longrightarrow X_n
$$
(Simplicial, lemme \ref{simplicial-lemma-relations-homotopy})
qui définissent l'homotopie entre les deux morphismes
$X_\bullet \to X_\bullet$ sont donnés par la règle
$$
h_{n, i} :
(U_n \to \ldots \to U_0, f)
\longmapsto
(U_n \to \ldots \to U_i \to U \to \ldots \to U, \text{id})
$$
pour $i > 0$, avec $h_{n, 0} = \text{id}$. Vérifications omises.
Il en résulte que $K_\bullet(j_{U!}\mathbf{Z}_U)$ a une homologie nulle en
degrés négatifs (par la fonctorialité de Simplicial, remarque
\ref{simplicial-remark-homotopy-better}, et le résultat de Simplicial, lemme
\ref{simplicial-lemma-homotopy-s-N}).
Ainsi, $K_\bullet(\mathcal{F})$ calcule les foncteurs dérivés à gauche
$H_n(\mathcal{C}, -)$ de $H_0(\mathcal{C}, -)$, par exemple d'après les
énoncés duaux de Homologie, lemme
\ref{homology-lemma-efface-implies-universal}, et Catégories dérivées, lemme
\ref{derived-lemma-right-derived-delta-functor}.
\end{example}

\begin{example}
\label{sites-cohomology-example-morphism-categories}
Soit $u : \mathcal{C}' \to \mathcal{C}$ un foncteur. Munissons
$\mathcal{C}'$ et $\mathcal{C}$ de la topologie chaotique comme dans
l'exemple \ref{sites-cohomology-example-category-to-point}. Les foncteurs $u$,
$\mathcal{C}' \to *$ et $\mathcal{C} \to *$, où $*$ est la catégorie ayant
un seul objet et un seul morphisme, sont cocontinus et continus. Soient
$g : \Sh(\mathcal{C}') \to \Sh(\mathcal{C})$,
$\pi' : \Sh(\mathcal{C}') \to \Sh(*)$, et
$\pi : \Sh(\mathcal{C}) \to \Sh(*)$,
les morphismes de topos correspondants. Soit $B$ un anneau. Nous munissons
$*$ du faisceau d'anneaux $B$, et $\mathcal{C}'$ et $\mathcal{C}$ du
faisceau constant $\underline{B}$. De cette manière,
$$
\xymatrix{
(\Sh(\mathcal{C}'), \underline{B}) \ar[rd]_{\pi'} \ar[rr]_g & &
(\Sh(\mathcal{C}), \underline{B}) \ar[ld]^\pi \\
& (\Sh(*), B)
}
$$
est un exemple de la situation
\ref{sites-cohomology-situation-morphism-fibred-categories}. Ainsi, le
lemme \ref{sites-cohomology-lemma-properties-lower-shriek-fibred-category} s'applique à $g$,
et il n'est pas nécessaire de distinguer $g_!$ sur les modules et sur les
faisceaux abéliens. En particulier, la
remarque \ref{sites-cohomology-remark-morphism-fibred-categories} produit des morphismes
canoniques
$$
H_n(\mathcal{C}', \mathcal{F}')
\longrightarrow
H_n(\mathcal{C}, \mathcal{F})
$$
lorsque $\mathcal{F}$ appartient à $\textit{Ab}(\mathcal{C})$,
$\mathcal{F}'$ à $\textit{Ab}(\mathcal{C}')$, et qu'un morphisme
$t : \mathcal{F}' \to g^{-1}\mathcal{F}$ est donné. En termes du calcul de
l'homologie de l'exemple \ref{sites-cohomology-example-left-derived-colimits}, ces morphismes
proviennent d'un morphisme de complexes
$$
K_\bullet(\mathcal{F}') \longrightarrow K_\bullet(\mathcal{F})
$$
donné par la règle
$$
(U'_n \to \ldots \to U'_0, s') \longmapsto
(u(U'_n) \to \ldots \to u(U'_0), t(s'))
$$
avec les notations évidentes.
\end{example}

\begin{remark}
\label{sites-cohomology-remark-map-evaluation-to-derived}
On reprend les notations et les hypothèses de
l'exemple \ref{sites-cohomology-example-category-to-point}. Soit $\mathcal{F}^\bullet$ un
complexe borné de faisceaux abéliens sur $\mathcal{C}$. Pour tout objet $U$
de $\mathcal{C}$, il existe un morphisme canonique
$$
\mathcal{F}^\bullet(U) \longrightarrow L\pi_!(\mathcal{F}^\bullet)
$$
dans $D(\textit{Ab})$. Si $\mathcal{F}^\bullet$ est un complexe de
$\underline{B}$-modules, ce morphisme appartient à $D(B)$. Pour le montrer,
remarquons que nous calculons $L\pi_!(\mathcal{F}^\bullet)$ en prenant un
quasi-isomorphisme $\mathcal{P}^\bullet \to \mathcal{F}^\bullet$, où
$\mathcal{P}^\bullet$ est un complexe de projectifs. Comme la topologie est
chaotique, $\mathcal{P}^\bullet(U) \to \mathcal{F}^\bullet(U)$ est un
quasi-isomorphisme, donc peut être inversé dans $D(\textit{Ab})$, resp.\
$D(B)$. En le composant avec le morphisme canonique
$\mathcal{P}^\bullet(U) \to \pi_!(\mathcal{P}^\bullet)$ provenant du calcul
de $\pi_!$ comme colimite, nous obtenons la flèche recherchée.
\end{remark}

\begin{lemma}
\label{sites-cohomology-lemma-initial-final}
On reprend les notations et les hypothèses de
l'exemple \ref{sites-cohomology-example-category-to-point}. Si $\mathcal{C}$ possède un objet
initial ou un objet final, alors
$L\pi_! \circ \pi^{-1} = \text{id}$ sur $D(\textit{Ab})$, resp.\ $D(B)$.
\end{lemma}

\begin{proof}
Si $\mathcal{C}$ possède un objet initial, alors $\pi_!$ se calcule par
évaluation sur cet objet et l'énoncé est clair. Si $\mathcal{C}$ possède un
objet final, alors $R\pi_*$ se calcule par évaluation sur cet objet, donc
$R\pi_* \circ \pi^{-1} \cong \text{id}$ sur $D(\textit{Ab})$, resp.\
$D(B)$. Il en résulte que
$\pi^{-1} : D(\textit{Ab}) \to D(\mathcal{C})$,
resp.\ $\pi^{-1} : D(B) \to D(\underline{B})$, est pleinement fidèle, voir
Catégories, lemme \ref{categories-lemma-adjoint-fully-faithful}.
Le même lemme implique alors que $L\pi_! \circ \pi^{-1} = \text{id}$,
comme voulu.
\end{proof}

\begin{lemma}
\label{sites-cohomology-lemma-change-of-rings}
On reprend les notations et les hypothèses de
l'exemple \ref{sites-cohomology-example-category-to-point}. Soit $B \to B'$ un morphisme
d'anneaux. Considérons le diagramme commutatif de topos annelés
$$
\xymatrix{
(\Sh(\mathcal{C}), \underline{B}) \ar[d]_\pi &
(\Sh(\mathcal{C}), \underline{B'}) \ar[d]^{\pi'} \ar[l]^h \\
(*, B) & (*, B') \ar[l]_f
}
$$
Alors $L\pi'_! \circ Lh^* = Lf^* \circ L\pi_!$.
\end{lemma}

\begin{proof}
Les deux foncteurs sont adjoints à gauche du foncteur évident
$D(B') \to D(\underline{B})$.
\end{proof}

\begin{lemma}
\label{sites-cohomology-lemma-compute-by-cosimplicial-resolution}
On reprend les notations et les hypothèses de
l'exemple \ref{sites-cohomology-example-category-to-point}. Soit $U_\bullet$ un objet
cosimplicial de $\mathcal{C}$ tel que, pour tout
$U \in \Ob(\mathcal{C})$, l'ensemble simplicial
$\Mor_\mathcal{C}(U_\bullet, U)$
soit homotopiquement équivalent à l'ensemble simplicial constant sur un
singleton. Alors
$$
L\pi_!(\mathcal{F}) = \mathcal{F}(U_\bullet)
$$
dans $D(\textit{Ab})$, resp.\ $D(B)$, fonctoriellement en $\mathcal{F}$ dans
$\textit{Ab}(\mathcal{C})$, resp.\ $\textit{Mod}(\underline{B})$.
\end{lemma}

\begin{proof}
Comme $L\pi_!$ coïncide pour les modules et les faisceaux abéliens d'après
le lemme \ref{sites-cohomology-lemma-properties-lower-shriek-fibred-category}, il suffit de
démontrer l'énoncé lorsque $\mathcal{F}$ est un faisceau abélien.
Pour $U \in \Ob(\mathcal{C})$, le faisceau abélien
$j_{U!}\mathbf{Z}_U$ est un objet projectif de
$\textit{Ab}(\mathcal{C})$, car
$\Hom(j_{U!}\mathbf{Z}_U, \mathcal{F}) = \mathcal{F}(U)$
et le foncteur des sections est exact puisque la topologie est chaotique.
Tout faisceau abélien est quotient d'une somme directe de
$j_{U!}\mathbf{Z}_U$ d'après Modules sur les sites, lemme
\ref{sites-modules-lemma-module-quotient-flat}. Nous pouvons donc calculer
$L\pi_!(\mathcal{F})$ en choisissant une résolution
$$
\ldots \to \mathcal{G}^{-1} \to \mathcal{G}^0 \to \mathcal{F} \to 0
$$
dont les termes sont des sommes directes de faisceaux de la forme ci-dessus,
puis en prenant $L\pi_!(\mathcal{F}) = \pi_!(\mathcal{G}^\bullet)$.
Considérons le bicomplexe
$A^{\bullet, \bullet} = \mathcal{G}^\bullet(U_\bullet)$.
Le morphisme $\mathcal{G}^0 \to \mathcal{F}$ donne un morphisme de complexes
$A^{0, \bullet} \to \mathcal{F}(U_\bullet)$.
Comme $\pi_!$ se calcule en prenant la colimite sur
$\mathcal{C}^{opp}$ (lemme \ref{sites-cohomology-lemma-compute-pi-shriek}), nous voyons que les
deux compositions
$\mathcal{G}^m(U_1) \to \mathcal{G}^m(U_0) \to \pi_!\mathcal{G}^m$
sont égales. Nous obtenons donc un morphisme canonique de complexes
$$
\text{Tot}(A^{\bullet, \bullet})
\longrightarrow
\pi_!(\mathcal{G}^\bullet) = L\pi_!(\mathcal{F})
$$
Pour démontrer le lemme, il suffit de montrer que les complexes
$$
\ldots \to \mathcal{G}^m(U_1) \to \mathcal{G}^m(U_0) \to
\pi_!\mathcal{G}^m \to 0
$$
sont exacts, voir Homologie, lemme
\ref{homology-lemma-double-complex-gives-resolution}.
Comme les faisceaux $\mathcal{G}^m$ sont des sommes directes des faisceaux
$j_{U!}\mathbf{Z}_U$, nous nous ramenons à
$\mathcal{G} = j_{U!}\mathbf{Z}_U$. Le complexe
$j_{U!}\mathbf{Z}_U(U_\bullet)$ est le complexe de groupes abéliens associé
au $\mathbf{Z}$-module libre sur l'ensemble simplicial
$\Mor_\mathcal{C}(U_\bullet, U)$, que nous avons supposé homotopiquement
équivalent à un singleton. Nous en concluons que
$$
j_{U!}\mathbf{Z}_U(U_\bullet) \to \mathbf{Z}
$$
est une équivalence d'homotopie de groupes abéliens, donc un
quasi-isomorphisme (Simplicial, remarque
\ref{simplicial-remark-homotopy-better} et lemme
\ref{simplicial-lemma-homotopy-s-N}). Cela achève la démonstration, puisque
$\pi_!j_{U!}\mathbf{Z}_U = \mathbf{Z}$, comme cela a été montré dans la
démonstration du lemme
\ref{sites-cohomology-lemma-properties-lower-shriek-fibred-category}.
\end{proof}

\begin{lemma}
\label{sites-cohomology-lemma-get-it-now}
On reprend les notations et les hypothèses de
l'exemple \ref{sites-cohomology-example-morphism-categories}. S'il existe un objet
cosimplicial $U'_\bullet$ de $\mathcal{C}'$ tel que le
lemme \ref{sites-cohomology-lemma-compute-by-cosimplicial-resolution} s'applique à la fois à
$U'_\bullet$ dans $\mathcal{C}'$ et à $u(U'_\bullet)$ dans $\mathcal{C}$,
alors $L\pi'_! \circ g^{-1} = L\pi_!$ comme foncteurs
$D(\mathcal{C}) \to D(\textit{Ab})$,
resp.\ $D(\mathcal{C}, \underline{B}) \to D(B)$.
\end{lemma}

\begin{proof}
Cela résulte immédiatement du
lemme \ref{sites-cohomology-lemma-compute-by-cosimplicial-resolution} et du fait que
$g^{-1}$ est donné par précomposition avec $u$.
\end{proof}

\begin{lemma}
\label{sites-cohomology-lemma-product-categories}
Soient $\mathcal{C}_i$, $i = 1, 2$, des catégories. Soient
$u_i : \mathcal{C}_1 \times \mathcal{C}_2 \to \mathcal{C}_i$ les foncteurs
de projection. Soit $B$ un anneau. Soient
$g_i : (\Sh(\mathcal{C}_1 \times \mathcal{C}_2), \underline{B}) \to
(\Sh(\mathcal{C}_i), \underline{B})$ les morphismes de topos annelés
correspondants, voir l'exemple \ref{sites-cohomology-example-morphism-categories}. Pour
$K_i \in D(\mathcal{C}_i, B)$, nous avons
$$
L(\pi_1 \times \pi_2)_!(
g_1^{-1}K_1 \otimes_{\underline{B}}^\mathbf{L} g_2^{-1}K_2)
=
L\pi_{1, !}(K_1) \otimes_B^\mathbf{L} L\pi_{2, !}(K_2)
$$
dans $D(B)$, avec les notations évidentes.
\end{lemma}

\begin{proof}
Comme les deux membres commutent aux colimites, il suffit de le démontrer pour
$K_1 = j_{U!}\underline{B}_U$ et $K_2 = j_{V!}\underline{B}_V$
avec $U \in \Ob(\mathcal{C}_1)$ et $V \in \Ob(\mathcal{C}_2)$.
Voir la construction de $L\pi_!$ dans le
lemme \ref{sites-cohomology-lemma-existence-derived-lower-shriek}. Dans ce cas,
$$
g_1^{-1}K_1 \otimes_{\underline{B}}^\mathbf{L} g_2^{-1}K_2 =
g_1^{-1}K_1 \otimes_{\underline{B}} g_2^{-1}K_2 =
j_{(U, V)!}\underline{B}_{(U, V)}
$$
Vérification omise. Le résultat s'ensuit, car les membres de gauche et de
droite de la formule du lemme s'évaluent tous deux en $B$, voir la
construction de $L\pi_!$ dans le
lemme \ref{sites-cohomology-lemma-existence-derived-lower-shriek}.
\end{proof}

\begin{lemma}
\label{sites-cohomology-lemma-eilenberg-zilber}
On reprend les notations et les hypothèses de
l'exemple \ref{sites-cohomology-example-category-to-point}. S'il existe un objet cosimplicial
$U_\bullet$ de $\mathcal{C}$ auquel le
lemme \ref{sites-cohomology-lemma-compute-by-cosimplicial-resolution} s'applique, alors
$$
L\pi_!(K_1 \otimes^\mathbf{L}_{\underline{B}} K_2) =
L\pi_!(K_1) \otimes^\mathbf{L}_B L\pi_!(K_2)
$$
pour tous $K_i \in D(\underline{B})$.
\end{lemma}

\begin{proof}
Considérons le diagramme de catégories et de foncteurs
$$
\xymatrix{
& & \mathcal{C} \\
\mathcal{C} \ar[r]^-u &
\mathcal{C} \times \mathcal{C} \ar[rd]^{u_2} \ar[ru]_{u_1} \\
& & \mathcal{C}
}
$$
où $u$ est le foncteur diagonal et les $u_i$ sont les foncteurs de projection.
Cela donne des morphismes de topos annelés $g$, $g_1$, $g_2$.
Pour tout objet $(U_1, U_2)$ de $\mathcal{C}$, nous avons
$$
\Mor_{\mathcal{C} \times \mathcal{C}}(u(U_\bullet), (U_1, U_2)) =
\Mor_\mathcal{C}(U_\bullet, U_1) \times \Mor_\mathcal{C}(U_\bullet, U_2)
$$
qui est homotopiquement équivalent à un point d'après Simplicial, lemme
\ref{simplicial-lemma-products-homotopy}. Ainsi, le
lemme \ref{sites-cohomology-lemma-get-it-now} donne
$L\pi_!(g^{-1}K) = L(\pi \times \pi)_!(K)$ pour tout $K$ dans
$D(\mathcal{C} \times \mathcal{C}, B)$. Prenons
$K = g_1^{-1}K_1 \otimes_B^\mathbf{L} g_2^{-1}K_2$. Alors
$g^{-1}K = K_1 \otimes^\mathbf{L}_{\underline{B}} K_2$, car
$g^{-1} = g^* = Lg^*$ commute au produit tensoriel dérivé
(lemme \ref{sites-cohomology-lemma-pullback-tensor-product}). Pour finir, nous appliquons le
lemme \ref{sites-cohomology-lemma-product-categories}.
\end{proof}

\begin{remark}[Modules simpliciaux]
\label{sites-cohomology-remark-simplicial-modules}
Soit $\mathcal{C} = \Delta$ et soit $B$ un anneau quelconque. C'est un cas
particulier de l'exemple \ref{sites-cohomology-example-category-to-point} où les hypothèses du
lemme \ref{sites-cohomology-lemma-compute-by-cosimplicial-resolution} sont vérifiées.
En effet, soit $U_\bullet$ l'objet cosimplicial de $\Delta$ donné par le
foncteur identité. Pour vérifier la condition, il faut montrer que, pour
$[m] \in \Ob(\Delta)$, l'ensemble simplicial
$\Delta[m] : n \mapsto \Mor_\Delta([n], [m])$ est homotopiquement équivalent
à un point. Cela est expliqué dans Simplicial, exemple
\ref{simplicial-example-simplex-contractible}.

\medskip\noindent
Dans cette situation, la catégorie $\textit{Mod}(\underline{B})$ est
simplement celle des $B$-modules simpliciaux, et le foncteur $L\pi_!$ envoie
un $B$-module simplicial $M_\bullet$ sur son complexe associé
$s(M_\bullet)$ de $B$-modules. Les résultats ci-dessus peuvent donc être
réinterprétés comme des résultats sur les modules simpliciaux. Par exemple,
un cas particulier du lemme \ref{sites-cohomology-lemma-eilenberg-zilber} est le suivant : si
$M_\bullet$, $M'_\bullet$ sont des $B$-modules simpliciaux plats, alors le
complexe $s(M_\bullet \otimes_B M'_\bullet)$ est quasi-isomorphe au complexe
total associé au bicomplexe $s(M_\bullet) \otimes_B s(M'_\bullet)$.
(Indication : utiliser la platitude pour passer des produits tensoriels
dérivés aux produits tensoriels usuels.) C'est un cas particulier du théorème
d'Eilenberg--Zilber que l'on peut trouver dans \cite{Eilenberg-Zilber}.
\end{remark}

\begin{lemma}
\label{sites-cohomology-lemma-O-homology-qis}
Soit $\mathcal{C}$ une catégorie (munie de la topologie chaotique).
Soit $\mathcal{O} \to \mathcal{O}'$ un morphisme de faisceaux d'anneaux sur
$\mathcal{C}$. Supposons que :
\begin{enumerate}
\item il existe un objet cosimplicial $U_\bullet$ de $\mathcal{C}$ comme
dans le lemme \ref{sites-cohomology-lemma-compute-by-cosimplicial-resolution}, et
\item $L\pi_!\mathcal{O} \to L\pi_!\mathcal{O}'$ est un isomorphisme.
\end{enumerate}
Pour $K$ dans $D(\mathcal{O})$, nous avons
$$
L\pi_!(K) = L\pi_!(K \otimes_\mathcal{O}^\mathbf{L} \mathcal{O}')
$$
dans $D(\textit{Ab})$.
\end{lemma}

\begin{proof}
Remarque : dans cette démonstration, $L\pi_!$ désigne le foncteur dérivé à
gauche de $\pi_!$ sur les faisceaux abéliens.
Comme $L\pi_!$ commute aux colimites, il suffit de démontrer l'énoncé pour
les complexes bornés supérieurement de $\mathcal{O}$-modules (comparer avec
l'argument de Catégories dérivées, proposition
\ref{derived-proposition-left-derived-exists}, ou se limiter simplement aux
complexes bornés supérieurement). Tout tel complexe est quasi-isomorphe à un
complexe borné supérieurement dont les termes sont des sommes directes de
$j_{U!}\mathcal{O}_U$, avec $U \in \Ob(\mathcal{C})$, voir Modules sur les
sites, lemme \ref{sites-modules-lemma-module-quotient-flat}.
Il suffit donc de démontrer le lemme pour $j_{U!}\mathcal{O}_U$. Par
hypothèse,
$$
S_\bullet = \Mor_\mathcal{C}(U_\bullet, U)
$$
est un ensemble simplicial homotopiquement équivalent à l'ensemble simplicial
constant sur un singleton. Posons $P_n = \mathcal{O}(U_n)$ et
$P'_n = \mathcal{O}'(U_n)$. Observons que le complexe associé au groupe
abélien simplicial
$$
X_\bullet : n \longmapsto \bigoplus\nolimits_{s \in S_n} P_n
$$
calcule $L\pi_!(j_{U!}\mathcal{O}_U)$ d'après le
lemme \ref{sites-cohomology-lemma-compute-by-cosimplicial-resolution}.
Comme $j_{U!}\mathcal{O}_U$ est un $\mathcal{O}$-module plat, nous avons
$j_{U!}\mathcal{O}_U \otimes^\mathbf{L}_\mathcal{O} \mathcal{O}' =
j_{U!}\mathcal{O}'_U$, et son image par $L\pi_!$ est calculée par le complexe
associé au groupe abélien simplicial
$$
X'_\bullet : n \longmapsto \bigoplus\nolimits_{s \in S_n} P'_n
$$
Comme la règle qui associe à un ensemble simplicial $T_\bullet$ le groupe
abélien simplicial de termes $\bigoplus_{t \in T_n} P_n$ est un foncteur,
$X_\bullet \to P_\bullet$ est une équivalence d'homotopie de groupes
abéliens simpliciaux. De même, la règle qui associe à $T_\bullet$ le groupe
abélien simplicial de termes $\bigoplus_{t \in T_n} P'_n$ est un foncteur.
Ainsi, $X'_\bullet \to P'_\bullet$ est une équivalence d'homotopie de
groupes abéliens simpliciaux.
Par hypothèse, $P_\bullet \to P'_\bullet$ est un quasi-isomorphisme
(puisque $P_\bullet$, resp.\ $P'_\bullet$, calcule $L\pi_!\mathcal{O}$,
resp.\ $L\pi_!\mathcal{O}'$, d'après le
lemme \ref{sites-cohomology-lemma-compute-by-cosimplicial-resolution}).
Nous concluons que $X_\bullet$ et $X'_\bullet$ sont quasi-isomorphes, comme
voulu.
\end{proof}

\begin{remark}
\label{sites-cohomology-remark-O-homology-B-homology-general}
Soient $\mathcal{C}$ et $B$ comme dans
l'exemple \ref{sites-cohomology-example-category-to-point}. Supposons qu'il existe un objet
cosimplicial comme dans le
lemme \ref{sites-cohomology-lemma-compute-by-cosimplicial-resolution}.
Soit $\mathcal{O} \to \underline{B}$ un morphisme de faisceaux d'anneaux sur
$\mathcal{C}$ qui induit un isomorphisme
$L\pi_!\mathcal{O} \to L\pi_!\underline{B}$.
Dans ce cas, nous obtenons un foncteur exact de catégories triangulées
$$
L\pi_! : D(\mathcal{O}) \longrightarrow D(B)
$$
En effet, pour tout objet $K$ de $D(\mathcal{O})$, nous avons
$L\pi^{\textit{Ab}}_!(K) =
L\pi^{\textit{Ab}}_!(K \otimes_{\mathcal{O}}^\mathbf{L} \underline{B})$
par le lemme \ref{sites-cohomology-lemma-O-homology-qis}.
Nous pouvons donc définir le foncteur affiché comme la composition de
$- \otimes^\mathbf{L}_\mathcal{O} \underline{B}$ avec le foncteur
$L\pi_! : D(\underline{B}) \to D(B)$.
Autrement dit, nous obtenons une structure de $B$-module sur $L\pi_!(K)$
provenant de l'identification (canonique et fonctorielle) de $L\pi_!(K)$ à
$L\pi_!(K \otimes_\mathcal{O}^\mathbf{L} \underline{B})$ donnée par le
lemme.
\end{remark}







\section{Calcul de l'image directe exceptionnelle dérivée}
\label{sites-cohomology-section-calculate}

\noindent
Dans cette section, nous appliquons les résultats de la
section \ref{sites-cohomology-section-homology}
pour calculer
$L\pi_!$ dans la situation \ref{sites-cohomology-situation-fibred-category} et
$Lg_!$ dans la situation \ref{sites-cohomology-situation-morphism-fibred-categories}.

\begin{lemma}
\label{sites-cohomology-lemma-compute-left-derived-pi-shriek-pre}
On reprend les hypothèses et les notations de la situation \ref{sites-cohomology-situation-fibred-category}.
Pour $\mathcal{F}$ dans $\textit{PAb}(\mathcal{C})$ et $n \geq 0$,
considérons le faisceau abélien $L_n(\mathcal{F})$ sur $\mathcal{D}$
qui est le faisceau associé au préfaisceau
$$
V \longmapsto H_n(\mathcal{C}_V, \mathcal{F}|_{\mathcal{C}_V})
$$
muni des applications de restriction indiquées dans la démonstration. Alors
$L_n(\mathcal{F}) = L_n(\mathcal{F}^\#)$.
\end{lemma}

\begin{proof}
Pour tout morphisme $h : V' \to V$ de $\mathcal{D}$, il existe un
foncteur de changement de base $h^* : \mathcal{C}_V \to \mathcal{C}_{V'}$ entre les
catégories fibres (Catégories, définition
\ref{categories-definition-pullback-functor-fibred-category}).
De plus, pour $U \in \Ob(\mathcal{C}_V)$, il existe un
morphisme fortement cartésien $h^*U \to U$ au-dessus de $h$.
La restriction le long de ces morphismes fortement cartésiens définit une
transformation de foncteurs
$$
\mathcal{F}|_{\mathcal{C}_V}
\longrightarrow
\mathcal{F}|_{\mathcal{C}_{V'}} \circ h^*.
$$
D'après l'exemple \ref{sites-cohomology-example-morphism-categories},
nous obtenons l'application de restriction voulue
$$
H_n(\mathcal{C}_V, \mathcal{F}|_{\mathcal{C}_V})
\longrightarrow
H_n(\mathcal{C}_{V'}, \mathcal{F}|_{\mathcal{C}_{V'}})
$$
Notons $L_{n, p}(\mathcal{F})$ ce préfaisceau, de sorte que
$L_n(\mathcal{F}) = L_{n, p}(\mathcal{F})^\#$.
Le morphisme canonique $\gamma : \mathcal{F} \to \mathcal{F}^+$
(Sites, théorème \ref{sites-theorem-plus})
définit un morphisme
canonique $L_{n, p}(\mathcal{F}) \to L_{n, p}(\mathcal{F}^+)$.
Il faut montrer que ce morphisme devient un isomorphisme après faisceautisation.

\medskip\noindent
Utilisons le calcul de l'homologie donné dans
l'exemple \ref{sites-cohomology-example-left-derived-colimits}. Notons
$K_\bullet(\mathcal{F}|_{\mathcal{C}_V})$ le complexe associé à
la restriction de $\mathcal{F}$ à la catégorie fibre $\mathcal{C}_V$.
D'après les remarques ci-dessus, nous obtenons un préfaisceau $K_\bullet(\mathcal{F})$
de complexes
$$
V \longmapsto K_\bullet(\mathcal{F}|_{\mathcal{C}_V})
$$
dont les préfaisceaux d'homologie sont les préfaisceaux $L_{n, p}(\mathcal{F})$.
Il suffit donc de montrer que
$$
K_\bullet(\mathcal{F}) \longrightarrow K_\bullet(\mathcal{F}^+)
$$
devient un isomorphisme après faisceautisation.

\medskip\noindent
Injectivité. Soit $V$ un objet de $\mathcal{D}$ et soit
$\xi \in K_n(\mathcal{F})(V)$ un élément dont l'image
est nulle dans $K_n(\mathcal{F}^+)(V)$. Il faut montrer qu'il existe un
recouvrement $\{V_j \to V\}$ tel que $\xi|_{V_j}$ soit nul dans
$K_n(\mathcal{F})(V_j)$. Écrivons
$$
\xi = \sum (U_{i, n} \to \ldots \to U_{i, 0}, \sigma_i)
$$
avec $\sigma_i \in \mathcal{F}(U_{i, 0})$. Nous pouvons faire en sorte que
chaque suite de morphismes $U_n \to \ldots \to U_0$ de $\mathcal{C}_V$
apparaisse au plus une fois. Comme les sommes intervenant dans la définition
du complexe $K_\bullet$ sont des sommes directes, l'image de cet élément ne peut
être nulle dans $K_n(\mathcal{F}^+)(V)$ que si tous les $\sigma_i$ ont une image
nulle dans $\mathcal{F}^+(U_{i, 0})$. Par construction de
$\mathcal{F}^+$, il existe des recouvrements $\{U_{i, 0, j} \to U_{i, 0}\}$
tels que $\sigma_i|_{U_{i, 0, j}}$ soit nul. Par notre construction de
la topologie sur $\mathcal{C}$, on peut écrire $U_{i, 0, j} \to U_{i, 0}$
comme le changement de base (Catégories, définition
\ref{categories-definition-pullback-functor-fibred-category})
de certains morphismes $V_{i, j} \to V$ et, de plus, chaque
$\{V_{i, j} \to V\}$ est un recouvrement. Choisissons un recouvrement
$\{V_j \to V\}$ qui raffine chacun des recouvrements $\{V_{i, j} \to V\}$.
Il est alors clair que $\xi|_{V_j} = 0$.

\medskip\noindent
Surjectivité. Démonstration omise. Indication : raisonner comme dans la preuve de
l'injectivité.
\end{proof}

\begin{lemma}
\label{sites-cohomology-lemma-compute-left-derived-pi-shriek}
On reprend les hypothèses et les notations de la situation \ref{sites-cohomology-situation-fibred-category}.
Pour $\mathcal{F}$ dans $\textit{Ab}(\mathcal{C})$ et $n \geq 0$,
le faisceau $L_n\pi_!(\mathcal{F})$ est égal au faisceau
$L_n(\mathcal{F})$ construit dans le
lemme \ref{sites-cohomology-lemma-compute-left-derived-pi-shriek-pre}.
\end{lemma}

\begin{proof}
Considérons la suite de foncteurs $\mathcal{F} \mapsto L_n(\mathcal{F})$
de $\textit{PAb}(\mathcal{C}) \to \textit{Ab}(\mathcal{D})$.
Puisque, pour tout $V \in \Ob(\mathcal{D})$, la suite de foncteurs
$H_n(\mathcal{C}_V, - )$ forme un $\delta$-foncteur,
il en va de même des foncteurs $\mathcal{F} \mapsto L_n(\mathcal{F})$.
Notre but est de montrer que ceux-ci forment un $\delta$-foncteur universel.
Pour cela, nous construisons des préfaisceaux abéliens
sur lesquels ces foncteurs s'annulent.

\medskip\noindent
Pour $U' \in \Ob(\mathcal{C})$, considérons le préfaisceau abélien
$\mathcal{F}_{U'} = j_{U'!}^{\textit{PAb}}\mathbf{Z}_{U'}$
(Modules sur les sites, remarque \ref{sites-modules-remark-localize-presheaves}).
Rappelons que
$$
\mathcal{F}_{U'}(U) =
\bigoplus\nolimits_{\Mor_\mathcal{C}(U, U')} \mathbf{Z}
$$
Si $U$ est au-dessus de $V = p(U)$ dans $\mathcal{D}$ et si $U'$ est au-dessus de $V' = p(U')$,
alors tout morphisme $a : U \to U'$ se factorise de manière unique sous la forme $U \to h^*U' \to U'$,
où $h = p(a) : V \to V'$ (voir
Catégories, définition
\ref{categories-definition-pullback-functor-fibred-category}).
Nous voyons donc que
$$
\mathcal{F}_{U'}|_{\mathcal{C}_V}
=
\bigoplus\nolimits_{h \in \Mor_\mathcal{D}(V, V')}
j_{h^*U'!}\mathbf{Z}_{h^*U'}
$$
où $j_{h^*U'} : \Sh(\mathcal{C}_V/h^*U') \to \Sh(\mathcal{C}_V)$
est le morphisme de localisation. Les faisceaux $j_{h^*U'!}\mathbf{Z}_{h^*U'}$
ont leurs groupes d'homologie supérieure nuls (voir
l'exemple \ref{sites-cohomology-example-left-derived-colimits}).
Nous en concluons que $L_n(\mathcal{F}_{U'}) = 0$ pour tout $n > 0$ et tout $U'$.
Il s'ensuit que tout préfaisceau abélien $\mathcal{F}$ est un quotient
d'un préfaisceau abélien $\mathcal{G}$ tel que $L_n(\mathcal{G}) = 0$ pour
tout $n > 0$ (Modules sur les sites, lemme
\ref{sites-modules-lemma-module-quotient-flat}).
Comme $L_n(\mathcal{F}) = L_n(\mathcal{F}^\#)$, nous voyons
qu'il en va de même pour les faisceaux abéliens. Ainsi,
la suite de foncteurs $L_n(-)$ est un delta-foncteur universel
sur $\textit{Ab}(\mathcal{C})$
(Homologie, lemme \ref{homology-lemma-efface-implies-universal}).
Comme elle coïncide avec
$H^{-n}(L\pi_!(-))$ pour $n = 0$ d'après le
lemme \ref{sites-cohomology-lemma-compute-pi-shriek},
nous concluons par l'unicité des $\delta$-foncteurs universels
(Homologie, lemme \ref{homology-lemma-uniqueness-universal-delta-functor})
et par
Catégories dérivées, lemme \ref{derived-lemma-right-derived-delta-functor}.
\end{proof}

\begin{lemma}
\label{sites-cohomology-lemma-compute-left-derived-g-shriek}
On reprend les hypothèses et les notations de la
situation \ref{sites-cohomology-situation-morphism-fibred-categories}.
Pour un faisceau abélien $\mathcal{F}'$ sur $\mathcal{C}'$, le faisceau
$L_ng_!(\mathcal{F}')$ est le faisceau associé au préfaisceau
$$
U \longmapsto H_n(\mathcal{I}_U, \mathcal{F}'_U)
$$
Pour les notations et les applications de restriction, voir la démonstration.
\end{lemma}

\begin{proof}
Posons $p(U) = V$. La catégorie $\mathcal{I}_U$ est la catégorie des couples
$(U', \varphi)$ où $\varphi : U \to u(U')$ est un morphisme de $\mathcal{C}$
tel que $p(\varphi) = \text{id}_V$, c'est-à-dire que $\varphi$ est un morphisme de la
catégorie fibre $\mathcal{C}_V$. Les morphismes
$(U'_1, \varphi_1) \to (U'_2, \varphi_2)$ sont donnés par les morphismes
$a : U'_1 \to U'_2$ de la catégorie fibre $\mathcal{C}'_V$ tels que
$\varphi_2 = u(a) \circ \varphi_1$. Le préfaisceau $\mathcal{F}'_U$ associe
à $(U', \varphi)$ le groupe $\mathcal{F}'(U')$.
Nous allons construire ci-dessous les applications de restriction.

\medskip\noindent
Choisissons une factorisation
$$
\xymatrix{
\mathcal{C}' \ar@<1ex>[r]^{u'} &
\mathcal{C}'' \ar[r]^{u''} \ar@<1ex>[l]^w & \mathcal{C}
}
$$
de $u$ comme dans
Catégories, lemme \ref{categories-lemma-ameliorate-morphism-fibred-categories}.
Alors $g_! = g''_! \circ g'_!$, et il en va de même pour les foncteurs dérivés.
D'autre part, le foncteur $g'_!$ est exact, voir
Modules sur les sites, lemme \ref{sites-modules-lemma-have-left-adjoint}.
Nous obtenons donc $Lg_!(\mathcal{F}') = Lg''_!(\mathcal{F}'')$, où
$\mathcal{F}'' = g'_!\mathcal{F}'$. Remarquons que
$\mathcal{F}'' = h^{-1}\mathcal{F}'$, où
$h : \Sh(\mathcal{C}'') \to \Sh(\mathcal{C}')$ est le morphisme de topos
associé à $w$, voir
Sites, lemme \ref{sites-lemma-have-left-adjoint}.
Le foncteur $u''$ fait de $\mathcal{C}''$ une catégorie fibrée
sur $\mathcal{C}$ ; le
lemme \ref{sites-cohomology-lemma-compute-left-derived-pi-shriek}
s'applique donc au calcul de $L_ng''_!$. Le résultat s'ensuit, car la
construction de $\mathcal{C}''$ dans la preuve de
Catégories, lemme \ref{categories-lemma-ameliorate-morphism-fibred-categories}
montre que la catégorie fibre $\mathcal{C}''_U$ est égale à
$\mathcal{I}_U$. De plus, $h^{-1}\mathcal{F}'|_{\mathcal{C}''_U}$
est donné par la règle décrite ci-dessus
(comme $w$ est continu et cocontinu d'après
Champs, lemme \ref{stacks-lemma-topology-inherited-functorial},
nous pouvons appliquer
Sites, lemme \ref{sites-lemma-when-shriek}).
\end{proof}
\section{Modules simpliciaux}
\label{sites-cohomology-section-simplicial-modules}

\noindent
Soit $A_\bullet$ un anneau simplicial. Rappelons que l'on peut voir
$A_\bullet$ comme un faisceau sur $\Delta$ (muni de la topologie chaotique),
voir Simplicial, section
\ref{simplicial-section-simplicial-presheaves}. Un module simplicial
$M_\bullet$ sur $A_\bullet$ n'est alors autre qu'un faisceau de
$A_\bullet$-modules sur $\Delta$. Autrement dit, pour tout $n \geq 0$, on a
un $A_n$-module $M_n$ et, pour toute application
$\varphi : [n] \to [m]$, une application correspondante
$$
M_\bullet(\varphi) : M_m \longrightarrow M_n
$$
qui est $A_\bullet(\varphi)$-linéaire, ces applications se composant de la
manière usuelle.

\medskip\noindent
Soit $\mathcal{C}$ un site. Un {\it faisceau simplicial d'anneaux}
$\mathcal{A}_\bullet$ sur $\mathcal{C}$ est un objet simplicial de la
catégorie des faisceaux d'anneaux sur $\mathcal{C}$. Dans ce cas,
l'affectation $U \mapsto \mathcal{A}_\bullet(U)$ est un faisceau d'anneaux
simpliciaux et, en fait, les deux notions sont équivalentes. Il en va de même
pour les faisceaux abéliens simpliciaux, les faisceaux simpliciaux d'algèbres
de Lie, et ainsi de suite.

\medskip\noindent
Cependant, comme dans le cas des anneaux simpliciaux ci-dessus, on peut
envisager autrement les faisceaux simpliciaux. Considérons en effet la
projection
$$
p : \Delta \times \mathcal{C} \longrightarrow \mathcal{C}
$$
Elle définit une catégorie fibrée dont les morphismes fortement cartésiens
sont exactement ceux de la forme $([n], U) \to ([n], V)$. On munit la
catégorie $\Delta \times \mathcal{C}$ de la topologie héritée de
$\mathcal{C}$ (voir Champs, section \ref{stacks-section-topology}). La
description simple des recouvrements de $\Delta \times \mathcal{C}$
(Champs, lemme \ref{stacks-lemma-topology-inherited}) montre immédiatement
qu'un faisceau simplicial d'anneaux sur $\mathcal{C}$ revient au même qu'un
faisceau d'anneaux sur $\Delta \times \mathcal{C}$.

\medskip\noindent
Par analogie avec les modules simpliciaux sur un anneau simplicial, on
définit comme suit les modules simpliciaux sur un faisceau simplicial
d'anneaux.

\begin{definition}
\label{sites-cohomology-definition-simplicial-module}
Soit $\mathcal{C}$ un site et soit $\mathcal{A}_\bullet$ un faisceau
simplicial d'anneaux sur $\mathcal{C}$. Un
{\it $\mathcal{A}_\bullet$-module simplicial} $\mathcal{F}_\bullet$
(parfois appelé {\it faisceau simplicial de
$\mathcal{A}_\bullet$-modules}) est un faisceau de modules sur le faisceau
d'anneaux sur $\Delta \times \mathcal{C}$ associé à
$\mathcal{A}_\bullet$.
\end{definition}

\noindent
On obtient une catégorie $\textit{Mod}(\mathcal{A}_\bullet)$ de modules
simpliciaux et la catégorie dérivée correspondante
$D(\mathcal{A}_\bullet)$. Un morphisme
$\mathcal{A}_\bullet \to \mathcal{B}_\bullet$ de faisceaux simpliciaux
d'anneaux donne un foncteur
$$
- \otimes^\mathbf{L}_{\mathcal{A}_\bullet} \mathcal{B}_\bullet :
D(\mathcal{A}_\bullet)
\longrightarrow
D(\mathcal{B}_\bullet)
$$
De plus, les résultats des sections précédentes déterminent un foncteur
$$
L\pi_! : D(\mathcal{A}_\bullet) \longrightarrow D(\mathcal{C})
$$
Pour un module simplicial $\mathcal{F}_\bullet$, l'objet
$L\pi_!(\mathcal{F}_\bullet)$ est représenté par le complexe de chaînes
associé $s(\mathcal{F}_\bullet)$ (Simplicial, section
\ref{simplicial-section-complexes}). Cela résulte des lemmes
\ref{sites-cohomology-lemma-compute-left-derived-pi-shriek} et
\ref{sites-cohomology-lemma-compute-by-cosimplicial-resolution}.

\begin{lemma}
\label{sites-cohomology-lemma-base-change-by-qis}
Soit $\mathcal{C}$ un site et soit
$\mathcal{A}_\bullet \to \mathcal{B}_\bullet$ un homomorphisme de faisceaux
simpliciaux d'anneaux sur $\mathcal{C}$. Si
$L\pi_!\mathcal{A}_\bullet \to L\pi_!\mathcal{B}_\bullet$ est un
isomorphisme dans $D(\mathcal{C})$, alors
$$
L\pi_!(K) =
L\pi_!(K \otimes^\mathbf{L}_{\mathcal{A}_\bullet} \mathcal{B}_\bullet)
$$
pour tout $K$ dans $D(\mathcal{A}_\bullet)$.
\end{lemma}

\begin{proof}
Soit $([n], U)$ un objet de $\Delta \times \mathcal{C}$. Comme $L\pi_!$
commute aux colimites, il suffit d'établir l'assertion pour les complexes
bornés supérieurement de $\mathcal{O}$-modules (comparer à l'argument de
Catégories dérivées, proposition
\ref{derived-proposition-left-derived-exists}, ou se limiter simplement aux
complexes bornés supérieurement). Tout tel complexe est quasi-isomorphe à un
complexe borné supérieurement dont les termes sont des modules plats, voir
Modules sur les sites, lemme \ref{sites-modules-lemma-module-quotient-flat}. Il
suffit donc de démontrer le lemme pour un
$\mathcal{A}_\bullet$-module plat $\mathcal{F}$. Dans ce cas, le produit
tensoriel dérivé est le produit tensoriel usuel, qui est encore un faisceau.
Le lemme \ref{sites-cohomology-lemma-compute-left-derived-pi-shriek} permet donc de calculer
les faisceaux de cohomologie des deux membres de l'égalité par le procédé du
lemme \ref{sites-cohomology-lemma-compute-left-derived-pi-shriek-pre}. Il suffit ainsi
d'établir le résultat pour la restriction de $\mathcal{F}$ aux catégories
fibres (c'est-à-dire à $\Delta \times U$). Dans ce cas, le résultat résulte du
lemme \ref{sites-cohomology-lemma-O-homology-qis}.
\end{proof}

\begin{remark}
\label{sites-cohomology-remark-homology-augmentation}
Soit $\mathcal{C}$ un site et soit
$\epsilon : \mathcal{A}_\bullet \to \mathcal{O}$ une augmentation
(Simplicial, définition \ref{simplicial-definition-augmentation}) dans la
catégorie des faisceaux d'anneaux. Supposons que $\epsilon$ induise un
quasi-isomorphisme $s(\mathcal{A}_\bullet) \to \mathcal{O}$. On obtient alors
un foncteur exact de catégories triangulées
$$
L\pi_! : D(\mathcal{A}_\bullet) \longrightarrow D(\mathcal{O})
$$
En effet, pour tout objet $K$ de $D(\mathcal{A}_\bullet)$, on a
$L\pi_!(K) = L\pi_!(K \otimes_{\mathcal{A}_\bullet}^\mathbf{L} \mathcal{O})$
d'après le lemme \ref{sites-cohomology-lemma-base-change-by-qis}. On peut donc définir le
foncteur affiché comme la composée de
$- \otimes^\mathbf{L}_{\mathcal{A}_\bullet} \mathcal{O}$ avec le foncteur
$L\pi_! : D(\Delta \times \mathcal{C}, \pi^{-1}\mathcal{O}) \to
D(\mathcal{O})$ de la remarque \ref{sites-cohomology-remark-fibred-category}. Autrement dit,
on obtient une structure de $\mathcal{O}$-module sur $L\pi_!(K)$ provenant de
l'identification (canonique et fonctorielle) de $L\pi_!(K)$ avec
$L\pi_!(K \otimes_{\mathcal{A}_\bullet}^\mathbf{L} \mathcal{O})$ donnée par
le lemme.
\end{remark}






\section{Cohomologie sur une catégorie}
\label{sites-cohomology-section-cohomology}

\noindent
Dans la situation de l'exemple \ref{sites-cohomology-example-category-to-point}, outre le
foncteur dérivé $L\pi_!$, on dispose également du foncteur $R\pi_*$. Pour un
faisceau abélien $\mathcal{F}$ sur $\mathcal{C}$, on a
$H_n(\mathcal{C}, \mathcal{F}) = H^{-n}(L\pi_!\mathcal{F})$ et
$H^n(\mathcal{C}, \mathcal{F}) = H^n(R\pi_*\mathcal{F})$.

\begin{example}[Calcul de la cohomologie]
\label{sites-cohomology-example-right-derived-limits}
Dans l'exemple \ref{sites-cohomology-example-category-to-point}, on peut calculer les
foncteurs $H^n(\mathcal{C}, -)$ comme suit. Soit
$\mathcal{F} \in \Ob(\textit{Ab}(\mathcal{C}))$. Considérons le complexe de
cochaînes
$$
K^\bullet(\mathcal{F}) :
\prod\nolimits_{U_0} \mathcal{F}(U_0)
\to
\prod\nolimits_{U_0 \to U_1} \mathcal{F}(U_0)
\to
\prod\nolimits_{U_0 \to U_1 \to U_2} \mathcal{F}(U_0)
\to \ldots
$$
dont les applications de transition sont données par
$$
(s_{U_0 \to U_1})
\longmapsto
((U_0 \to U_1 \to U_2) \mapsto s_{U_0 \to U_1} - s_{U_0 \to U_2}
+ s_{U_1 \to U_2}|_{U_0})
$$
et de même aux autres degrés. Par construction,
$$
H^0(\mathcal{C}, \mathcal{F}) =
\lim_{\mathcal{C}^{opp}} \mathcal{F} =
H^0(K^\bullet(\mathcal{F})),
$$
voir Catégories, lemme
\ref{categories-lemma-limits-products-equalizers}. La construction de
$K^\bullet(\mathcal{F})$ est fonctorielle en $\mathcal{F}$ et transforme les
suites exactes courtes de $\textit{Ab}(\mathcal{C})$ en suites exactes
courtes de complexes. La famille de foncteurs
$\mathcal{F} \mapsto H^n(K^\bullet(\mathcal{F}))$ forme donc un
$\delta$-foncteur, voir Homology, définition
\ref{homology-definition-cohomological-delta-functor} et lemme
\ref{homology-lemma-long-exact-sequence-cochain}. Pour un objet $U$ de
$\mathcal{C}$, notons $p_U : \Sh(*) \to \Sh(\mathcal{C})$ le point
correspondant, pour lequel $p_U^{-1}$ est l'évaluation en $U$, voir Sites,
exemple \ref{sites-example-indiscrete-points}. Soit $A$ un groupe abélien et
posons $\mathcal{F} = p_{U, *}A$. Le complexe $K^\bullet(\mathcal{F})$ a alors
pour termes $\text{Map}(X_n, A)$, où
$$
X_n = \coprod\nolimits_{U_0 \to \ldots \to U_{n - 1} \to U_n}
\Mor_\mathcal{C}(U, U_0)
$$
Cet ensemble simplicial est homotopiquement équivalent à l'ensemble
simplicial constant réduit à un point $\{*\}$. En effet, l'application
$X_\bullet \to \{*\}$ est évidente, l'application $\{*\} \to X_n$ envoie
$*$ sur $(U \to \ldots \to U, \text{id}_U)$, et les applications
$$
h_{n, i} : X_n \longrightarrow X_n
$$
(Simplicial, lemme \ref{simplicial-lemma-relations-homotopy}) qui définissent
l'homotopie entre les deux applications
$X_\bullet \to X_\bullet$ sont données par la règle
$$
h_{n, i} :
(U_0 \to \ldots \to U_n, f)
\longmapsto
(U \to \ldots \to U \to U_i \to \ldots \to U_n, \text{id})
$$
pour $i > 0$, et $h_{n, 0} = \text{id}$. Nous omettons les vérifications.
Comme $\text{Map}(-, A)$ est un foncteur contravariant, il s'ensuit que
$K^\bullet(p_{U, *}A)$ a une cohomologie nulle en degrés strictement positifs
(par la fonctorialité de Simplicial, remarque
\ref{simplicial-remark-homotopy-better}, et le résultat de Simplicial, lemme
\ref{simplicial-lemma-homotopy-s-Q}). Il en résulte que
$K^\bullet(\mathcal{F})$ est encore acyclique en degrés strictement positifs
lorsque $\mathcal{F}$ est un produit de faisceaux de la forme $p_{U, *}A$.
Comme tout faisceau abélien sur $\mathcal{C}$ se plonge dans un tel produit,
on conclut que $K^\bullet(\mathcal{F})$ calcule les foncteurs dérivés à droite
$H^n(\mathcal{C}, -)$ de $H^0(\mathcal{C}, -)$, par exemple d'après
Homology, lemme \ref{homology-lemma-efface-implies-universal}, et Derived
Catégories, lemme \ref{derived-lemma-right-derived-delta-functor}.
\end{example}

\begin{example}[Calcul des Ext]
\label{sites-cohomology-example-computing-exts}
Dans l'exemple \ref{sites-cohomology-example-category-to-point}, supposons en outre donné un
faisceau d'anneaux $\mathcal{O}$ sur $\mathcal{C}$. Soient $\mathcal{F}$ et
$\mathcal{G}$ des $\mathcal{O}$-modules. Considérons le complexe
$K^\bullet(\mathcal{G}, \mathcal{F})$ dont le terme de degré $n$ est
$$
\prod\nolimits_{U_0 \to U_1 \to \ldots \to U_n}
\Hom_{\mathcal{O}(U_n)}(\mathcal{G}(U_n), \mathcal{F}(U_0))
$$
et dont l'application de transition est donnée par
$$
(\varphi_{U_0 \to U_1})
\longmapsto
((U_0 \to U_1 \to U_2) \mapsto
\varphi_{U_0 \to U_1} \circ \rho^{U_2}_{U_1}
- \varphi_{U_0 \to U_2}
+ \rho^{U_1}_{U_0} \circ \varphi_{U_1 \to U_2}
$$
et de même aux autres degrés. Les $\rho$ désignent ici les applications de
restriction. Par construction,
$$
\Hom_\mathcal{O}(\mathcal{G}, \mathcal{F}) =
H^0(K^\bullet(\mathcal{G}, \mathcal{F}))
$$
pour toute paire de $\mathcal{O}$-modules $\mathcal{F}, \mathcal{G}$.
L'affectation
$(\mathcal{G}, \mathcal{F}) \mapsto K^\bullet(\mathcal{G}, \mathcal{F})$
est un bifoncteur qui transforme les sommes directes en la première variable
en produits et qui commute aux produits en la seconde variable. Nous
affirmons que
$$
\Ext^i_\mathcal{O}(\mathcal{G}, \mathcal{F}) =
H^i(K^\bullet(\mathcal{G}, \mathcal{F}))
$$
pour $i \geq 0$, pourvu que l'une des conditions suivantes soit satisfaite :
\begin{enumerate}
\item $\mathcal{G}(U)$ est un $\mathcal{O}(U)$-module projectif pour tout
$U \in \Ob(\mathcal{C})$ ;
\item $\mathcal{F}(U)$ est un $\mathcal{O}(U)$-module injectif pour tout
$U \in \Ob(\mathcal{C})$.
\end{enumerate}
Dans le cas (1), le foncteur $K^\bullet(\mathcal{G}, -)$ est un foncteur exact
de la catégorie des $\mathcal{O}$-modules vers celle des complexes de
cochaînes de groupes abéliens. En raisonnant comme dans l'exemple
\ref{sites-cohomology-example-right-derived-limits}, il suffit donc de montrer que
$K^\bullet(\mathcal{G}, \mathcal{F})$ est acyclique en degrés strictement
positifs lorsque $\mathcal{F}$ est de la forme $p_{U, *}A$ pour un
$\mathcal{O}(U)$-module $A$.
Choisissons une suite exacte courte
\begin{equation}
\label{sites-cohomology-equation-split}
0 \to \mathcal{G}' \to \bigoplus j_{U_i!}\mathcal{O}_{U_i} \to
\mathcal{G} \to 0
\end{equation}
voir Modules sur les sites, lemme
\ref{sites-modules-lemma-module-quotient-flat}. Comme (1) vaut pour les
faisceaux du milieu et de droite, il vaut aussi pour $\mathcal{G}'$, et
l'évaluation de (\ref{sites-cohomology-equation-split}) en un objet de $\mathcal{C}$ donne une
suite exacte scindée de modules. On obtient une suite exacte courte de
complexes
$$
0 \to
K^\bullet(\mathcal{G}, \mathcal{F}) \to
\prod K^\bullet(j_{U_i!}\mathcal{O}_{U_i}, \mathcal{F}) \to
K^\bullet(\mathcal{G}', \mathcal{F}) \to 0
$$
pour tout $\mathcal{F}$, en particulier pour
$\mathcal{F} = p_{U, *}A$. En degré $H^0$, on obtient
$$
0 \to \Hom(\mathcal{G}, p_{U, *}A) \to
\Hom(\bigoplus j_{U_i!}\mathcal{O}_{U_i}, p_{U, *}A) \to
\Hom(\mathcal{G}', p_{U, *}A) \to 0
$$
qui est exacte, car
$\Hom(\mathcal{H}, p_{U, *}A) = \Hom_{\mathcal{O}(U)}(\mathcal{H}(U), A)$
et la suite des sections de (\ref{sites-cohomology-equation-split}) sur $U$ est exacte
scindée. Par décalage dimensionnel, il suffit donc de montrer que
$K^\bullet(j_{U'!}\mathcal{O}_{U'}, p_{U, *}A)$ est acyclique en degrés
strictement positifs pour tous $U, U' \in \Ob(\mathcal{C})$. Dans ce cas,
$K^n(j_{U'!}\mathcal{O}_{U'}, p_{U, *}A)$ est égal à
$$
\prod\nolimits_{U \to U_0 \to U_1 \to \ldots \to U_n \to U'} A
$$
Autrement dit, $K^\bullet(j_{U'!}\mathcal{O}_{U'}, p_{U, *}A)$ est le
complexe de termes $\text{Map}(X_\bullet, A)$, où
$$
X_n = \coprod\nolimits_{U_0 \to \ldots \to U_{n - 1} \to U_n}
\Mor_\mathcal{C}(U, U_0) \times \Mor_\mathcal{C}(U_n, U')
$$
Cet ensemble simplicial est homotopiquement équivalent à l'ensemble
simplicial constant $\Mor_\mathcal{C}(U, U')$ : cela se démontre exactement
comme l'assertion correspondante de l'exemple
\ref{sites-cohomology-example-right-derived-limits}. Ceci achève la démonstration de
l'affirmation.

\medskip\noindent
L'argument dans le cas (2) est analogue, par dualité.
\end{example}






\section{Modules sur une catégorie}
\label{sites-cohomology-section-modules-cohomology}

\noindent
Les résultats de cette section serviront à définir une variante de la
catégorie dérivée des modules quasi-cohérents sur un champ en groupoïdes
au-dessus de la catégorie des schémas. Voir Faisceaux on Champs, section
\ref{stacks-sheaves-section-QC}.

\medskip\noindent
Soit $\mathcal{C}$ une catégorie. Considérons $\mathcal{C}$ comme un site muni
de la topologie chaotique. Comme dans l'exemple
\ref{sites-cohomology-example-computing-exts}, soit
$\mathcal{O}$ un faisceau d'anneaux sur $\mathcal{C}$. Autrement dit,
$\mathcal{O}$ est un préfaisceau d'anneaux sur la catégorie $\mathcal{C}$,
voir Catégories, définition \ref{categories-definition-presheaf}.

\begin{definition}
\label{sites-cohomology-definition-cartesian}
Dans la situation ci-dessus, on note
$\mathit{QC}(\mathcal{C}, \mathcal{O})$, ou simplement
{\it $\mathit{QC}(\mathcal{O})$}, la sous-catégorie pleine de
$D(\mathcal{O}) = D(\mathcal{C}, \mathcal{O})$ formée des objets $K$ tels que,
pour tout $U \to V$ dans $\mathcal{C}$, l'application canonique
$$
R\Gamma(V, K) \otimes_{\mathcal{O}(V)}^\mathbf{L} \mathcal{O}(U)
\longrightarrow
R\Gamma(U, K)
$$
soit un isomorphisme dans $D(\mathcal{O}(U))$.
\end{definition}

\begin{lemma}
\label{sites-cohomology-lemma-cartesian-good-sub}
Dans la situation ci-dessus, $\mathit{QC}(\mathcal{O})$ est une
sous-catégorie strictement pleine, saturée et triangulée de
$D(\mathcal{O})$, stable par sommes directes arbitraires.
\end{lemma}

\begin{proof}
Soit $U$ un objet de $\mathcal{C}$. Comme la topologie de $\mathcal{C}$ est
chaotique, le foncteur $\mathcal{F} \mapsto \mathcal{F}(U)$ est exact et
commute aux sommes directes. Le foncteur exact
$K \mapsto R\Gamma(U, K)$ se calcule donc en représentant $K$ par un complexe
quelconque $\mathcal{F}^\bullet$ de $\mathcal{O}$-modules et en prenant
$\mathcal{F}^\bullet(U)$. Ainsi, $R\Gamma(U, -)$ commute aux sommes directes,
voir Injectifs, lemme \ref{injectives-lemma-derived-products}. De même, pour
un morphisme $U \to V$ de $\mathcal{C}$, le foncteur produit tensoriel dérivé
$- \otimes_{\mathcal{O}(V)}^\mathbf{L} \mathcal{O}(U) :
D(\mathcal{O}(V)) \to D(\mathcal{O}(U))$
est exact et commute aux sommes directes. Le lemme résulte immédiatement de
ces observations ; nous omettons les détails.
\end{proof}

\begin{lemma}
\label{sites-cohomology-lemma-QC-bounded-above}
Dans la situation ci-dessus, supposons que $M$ soit un objet de
$\mathit{QC}(\mathcal{O})$ et que $b \in \mathbf{Z}$ vérifie $H^i(M) = 0$ pour
tout $i > b$. Alors $H^b(M)$ est un module quasi-cohérent sur
$(\mathcal{C}, \mathcal{O})$ au sens de Modules sur les sites, définition
\ref{sites-modules-definition-site-local}.
\end{lemma}

\begin{proof}
D'après Modules sur les sites, lemme
\ref{sites-modules-lemma-chaotic-quasi-coherent}, il suffit de montrer que,
pour tout morphisme $U \to V$ de $\mathcal{C}$, l'application
$$
H^b(M)(V) \otimes_{\mathcal{O}(V)} \mathcal{O}(U) \to H^b(M)(U)
$$
est un isomorphisme. Par hypothèse, l'application
$$
R\Gamma(V, M) \otimes_{\mathcal{O}(V)}^\mathbf{L} \mathcal{O}(U)
\to R\Gamma(U, M)
$$
est un isomorphisme. Le résultat découle alors, par exemple, de la suite
spectrale de Tor. Nous omettons les détails.
\end{proof}

\begin{lemma}
\label{sites-cohomology-lemma-QC-final-object}
Dans la situation ci-dessus, supposons que $\mathcal{C}$ possède un objet
final $X$. Posons $R = \mathcal{O}(X)$ et notons
$f : (\mathcal{C}, \mathcal{O}) \to (pt, R)$ le morphisme évident de sites.
Alors $\mathit{QC}(\mathcal{O}) = D(R)$ via $Lf^*$ et $Rf_*$.
\end{lemma}

\begin{proof}
Omis.
\end{proof}

\begin{lemma}
\label{sites-cohomology-lemma-QC-hom-out-of}
Dans la situation ci-dessus, supposons que $K$ soit un objet de
$\mathit{QC}(\mathcal{O})$ et que $M$ soit un objet quelconque de
$D(\mathcal{O})$. Pour tout objet $U$ de $\mathcal{C}$, on a
$$
\Hom_{D(\mathcal{O}_U)}(K|_U, M|_U) =
R\Hom_{\mathcal{O}(U)}(R\Gamma(U, K), R\Gamma(U, M))
$$
\end{lemma}

\begin{proof}
On peut remplacer $\mathcal{C}$ par $\mathcal{C}/U$. On peut donc supposer
que $U = X$ est un objet final de $\mathcal{C}$. D'après le lemme
\ref{sites-cohomology-lemma-QC-final-object}, on a $K = Lf^*P$, où
$P = R\Gamma(U, K) = R\Gamma(X, K) = Rf_*K$. Le résultat s'ensuit, puisque
$Lf^*$ est adjoint à gauche de $Rf_*(-) = R\Gamma(U, -)$.
\end{proof}

\noindent
Soit $(\mathcal{C}, \mathcal{O})$ comme ci-dessus. Pour un complexe
$\mathcal{F}^\bullet$ de $\mathcal{O}$-modules, on définit la
{\it taille} $|\mathcal{F}^\bullet|$ de $\mathcal{F}^\bullet$ par
$$
|\mathcal{F}^\bullet| =
\left|
\coprod\nolimits_{i \in \mathbf{Z},\ U \in \Ob(\mathcal{C})} \mathcal{F}^i(U)
\right|
$$
Pour un objet $K$ de $D(\mathcal{O})$, on appelle {\it taille} $|K|$ de $K$
le cardinal
$$
|K| = \min
\left\{
\left|
\mathcal{F}^\bullet
\right|
\text{ où }\mathcal{F}^\bullet\text{ représente }K
\right\}
$$
Les propriétés des cardinaux assurent l'existence de ce minimum.

\begin{lemma}
\label{sites-cohomology-lemma-build-from}
Dans la situation ci-dessus, il existe un cardinal $\kappa$ ayant la
propriété suivante : étant donnés un complexe $\mathcal{F}^\bullet$ de
$\mathcal{O}$-modules et des sous-ensembles
$\Omega^i_U \subset \mathcal{F}^i(U)$, il existe un sous-complexe
$\mathcal{H}^\bullet \subset \mathcal{F}^\bullet$ tel que
$\Omega^i_U \subset \mathcal{H}^i(U)$ et
$|\mathcal{H}^\bullet| \leq \max(\kappa, |\bigcup \Omega^i_U|)$.
\end{lemma}

\begin{proof}
Définissons $\mathcal{H}^i(U)$ comme le sous-$\mathcal{O}(U)$-module de
$\mathcal{F}^i(U)$ engendré par les images des $\Omega^i_V$ et de
$\text{d}(\Omega^{i - 1}_U)$ par restriction le long de tout morphisme
$f : U \to V$. Le cardinal de $\mathcal{H}^i(U)$ est majoré par
le maximum de $\aleph_0$, du cardinal de $\mathcal{O}(U)$, du cardinal de
$\text{Arrows}(\mathcal{C})$ et de $|\bigcup \Omega^i_U|$. Nous omettons les
détails.
\end{proof}

\begin{lemma}
\label{sites-cohomology-lemma-qis-small}
Dans la situation ci-dessus, il existe un cardinal $\kappa$ ayant la
propriété suivante : pour tout complexe $\mathcal{F}^\bullet$ de
$\mathcal{O}$-modules représentant un objet $K$ de $D(\mathcal{O})$, il existe
un sous-complexe $\mathcal{H}^\bullet \subset \mathcal{F}^\bullet$ tel que
$\mathcal{H}^\bullet$ représente $K$ et tel que
$|\mathcal{H}^\bullet| \leq \max(\kappa, |K|)$.
\end{lemma}

\begin{proof}
Pour commencer, choisissons, pour tout $i$ et tout $U$, un sous-ensemble
$\Omega^i_U \subset \Ker(\text{d} :
\mathcal{F}^i(U) \to \mathcal{F}^{i + 1}(U))$
qui s'applique bijectivement sur
$H^i(K)(U) = H^i(\mathcal{F}^\bullet(U))$.
On a donc $|\Omega^i_U| \leq |K|$, puisque l'on peut représenter $K$ par un
complexe de taille $|K|$. En appliquant le lemme \ref{sites-cohomology-lemma-build-from}, on
obtient un sous-complexe
$\mathcal{S}^\bullet \subset \mathcal{F}^\bullet$ de taille au plus
$\max(\kappa, |K|)$ qui contient les $\Omega^i_U$ ; par conséquent,
$H^i(\mathcal{S}^\bullet) \to H^i(\mathcal{F}^\bullet)$ est une surjection de
faisceaux.

\medskip\noindent
Nous allons construire par récurrence des sous-complexes
$$
\mathcal{S}^\bullet =
\mathcal{S}_0^\bullet \subset
\mathcal{S}_1^\bullet \subset
\mathcal{S}_2^\bullet \subset \ldots \subset \mathcal{F}^\bullet
$$
de taille $\leq \max(\kappa, |K|)$, de sorte que le noyau de
$H^i(\mathcal{S}_n^\bullet) \to H^i(\mathcal{F}^\bullet)$ coïncide avec le
noyau de
$H^i(\mathcal{S}_n^\bullet) \to H^i(\mathcal{S}_{n + 1}^\bullet)$.
Une fois ceci fait, on pourra prendre
$\mathcal{H}^\bullet = \bigcup \mathcal{S}_n^\bullet$.

\medskip\noindent
Construisons $\mathcal{S}_{n + 1}^\bullet$ à partir de
$\mathcal{S}_n^\bullet$. Pour tout $U$ et tout $i$, soit
$\Omega^{i - 1}_U \subset \mathcal{F}^{i - 1}(U)$ un sous-ensemble que
$\text{d} : \mathcal{F}^{i - 1}(U) \to \mathcal{F}^i(U)$ applique
$\Omega^{i - 1}_U$ bijectivement sur
$$
\mathcal{S}_n^i(U) \cap
\text{Im}(\text{d} : \mathcal{F}^{i - 1}(U)\to \mathcal{F}^i(U))
$$
Remarquons que $|\Omega^i_U| \leq \max(\kappa,|K|)$, car
$\mathcal{S}_n^i(U)$ est ainsi majoré. On obtient alors
$\mathcal{S}_{n + 1}^\bullet$ en appliquant le lemme
\ref{sites-cohomology-lemma-build-from} aux sous-ensembles
$$
\mathcal{S}_n^i(U) \cup \Omega^i_U \subset \mathcal{F}^i(U)
$$
ce qui achève la démonstration.
\end{proof}

\begin{lemma}
\label{sites-cohomology-lemma-cartesian-colimit-kappa}
Dans la situation ci-dessus, il existe un cardinal $\kappa$ ayant les
propriétés suivantes :
\begin{enumerate}
\item pour tout objet non nul $K$ de $\mathit{QC}(\mathcal{O})$, il existe un
morphisme non nul $E \to K$ dans $\mathit{QC}(\mathcal{O})$ tel que
$|E| \leq \kappa$ ;
\item pour tout morphisme $\alpha : E \to \bigoplus_n K_n$ dans
$\mathit{QC}(\mathcal{O})$ tel que $|E| \leq \kappa$, il existe des
morphismes $E_n \to K_n$ dans $\mathit{QC}(\mathcal{O})$, avec
$|E_n| \leq \kappa$, tels que $\alpha$ se factorise par
$\bigoplus E_n \to \bigoplus K_n$.
\end{enumerate}
\end{lemma}

\begin{proof}
Soit $\kappa$ un majorant de l'ensemble de cardinaux suivant :
\begin{enumerate}
\item $|\coprod_V j_{U!}\mathcal{O}_U(V)|$ pour tout
$U \in \Ob(\mathcal{C})$ ;
\item les cardinaux $\kappa(\mathcal{O}(V) \to \mathcal{O}(U))$ obtenus dans
Compléments d'algèbre, lemme \ref{more-algebra-lemma-enlarge}, pour tout morphisme
$U \to V$ de $\mathcal{C}$ ;
\item le cardinal obtenu dans le lemme \ref{sites-cohomology-lemma-qis-small}.
\end{enumerate}
Nous affirmons que, pour tout complexe $\mathcal{F}^\bullet$ représentant un
objet de $\mathit{QC}(\mathcal{O})$ et tout sous-complexe
$\mathcal{S}^\bullet \subset \mathcal{F}^\bullet$ vérifiant
$|\mathcal{S}^\bullet| \leq \kappa$, il existe un sous-complexe
$\mathcal{H}^\bullet$ de $\mathcal{F}^\bullet$, contenant
$\mathcal{S}^\bullet$, tel que $\mathcal{H}^\bullet$ représente un objet de
$\mathit{QC}(\mathcal{O})$ et satisfait
$|\mathcal{H}^\bullet| \leq \kappa$. Dans les deux paragraphes suivants,
nous montrons que cette affirmation entraîne le lemme.

\medskip\noindent
Pour établir (1), soit $K$ un objet non nul de
$\mathit{QC}(\mathcal{O})$. Supposons que $K$ soit représenté par le complexe
de $\mathcal{O}$-modules $\mathcal{F}^\bullet$. Alors
$H^i(\mathcal{F}^\bullet)$ est non nul pour un certain $i$. Il existe donc un
objet $U$ de $\mathcal{C}$ et une section
$s \in \mathcal{F}^i(U)$ telle que $d(s) = 0$ et qui détermine une section non
nulle de $H^i(\mathcal{F}^\bullet)$ sur $U$. L'image de
$s : j_{U!}\mathcal{O}_U[-i] \to \mathcal{F}^\bullet$ est alors un
sous-complexe $\mathcal{S}^\bullet \subset \mathcal{F}^\bullet$ tel que
$|\mathcal{S}^\bullet| \leq \kappa$. En appliquant l'affirmation, on obtient
un morphisme non nul $\mathcal{H}^\bullet \to \mathcal{F}^\bullet$ dans
$\mathit{QC}(\mathcal{O})$, avec
$|\mathcal{H}^\bullet| \leq \kappa$. Ceci établit (1).

\medskip\noindent
Soit $\alpha : E \to \bigoplus K_n$ comme dans (2). Choisissons des complexes
$\mathcal{K}_n^\bullet$ représentant les $K_n$. Alors
$\bigoplus \mathcal{K}_n^\bullet$ représente $\bigoplus K_n$. Par la
construction de la catégorie dérivée, on peut représenter $E$ par un complexe
$\mathcal{E}^\bullet$ de sorte que $\alpha$ soit représenté par un morphisme
de complexes
$a : \mathcal{E}^\bullet \to \bigoplus \mathcal{K}_n^\bullet$.
D'après le lemme \ref{sites-cohomology-lemma-qis-small} et le choix de $\kappa$ ci-dessus, on
peut supposer $|\mathcal{E}^\bullet| \leq \kappa$. L'affirmation fournit des
sous-complexes $\mathcal{E}_n^\bullet \subset \mathcal{K}_n^\bullet$
représentant des objets $E_n$ de $\mathit{QC}(\mathcal{O})$ tels que
$|E_n| \leq \kappa$ et contenant l'image de
$a_n : \mathcal{E}^\bullet \to \mathcal{K}_n^\bullet$, comme voulu.

\medskip\noindent
Démonstration de l'affirmation. Soit $\mathcal{F}^\bullet$ un complexe
représentant un objet de $\mathit{QC}(\mathcal{O})$ et soit
$\mathcal{S}^\bullet \subset \mathcal{F}^\bullet$ un sous-complexe de taille
$\leq \kappa$. Nous allons construire par récurrence des sous-complexes
$$
\mathcal{S}^\bullet = \mathcal{S}_0^\bullet \subset
\mathcal{S}_1^\bullet \subset
\mathcal{S}_2^\bullet \subset \ldots \subset \mathcal{F}^\bullet
$$
de taille $\leq \kappa$ tels que, pour tout morphisme $f : U \to V$ de
$\mathcal{C}$ et tout $i \in \mathbf{Z}$,
\begin{enumerate}
\item le noyau de la flèche
$H^i(\mathcal{S}_n^\bullet(V)
\otimes_{\mathcal{O}(V)}^\mathbf{L} \mathcal{O}(U)) \to
H^i(\mathcal{S}_n^\bullet(U))$
s'envoie sur zéro dans
$H^i(\mathcal{S}_{n + 1}^\bullet(V)
\otimes_{\mathcal{O}(V)}^\mathbf{L} \mathcal{O}(U))$,
\item l'image de la flèche
$H^i(\mathcal{S}_n^\bullet(U)) \to H^i(\mathcal{S}_{n + 1}^\bullet(U))$
est contenue dans l'image de
$H^i(\mathcal{S}_{n + 1}^\bullet(V)
\otimes_{\mathcal{O}(V)}^\mathbf{L} \mathcal{O}(U)) \to
H^i(\mathcal{S}_{n + 1}^\bullet(U))$,
\end{enumerate}
Une fois ceci fait, on pourra poser
$\mathcal{H}^\bullet = \bigcup \mathcal{S}_n^\bullet$. En effet, le produit
tensoriel dérivé et la prise de cohomologie des complexes de modules sur des
anneaux commutent aux colimites filtrantes ; les conditions (1) et (2)
garantiront donc ensemble que
$$
\mathcal{H}^\bullet(V) \otimes_{\mathcal{O}(V)}^\mathbf{L} \mathcal{O}(U)
\longrightarrow
\mathcal{H}^\bullet(U)
$$
est un isomorphisme en cohomologie à tout degré, donc un isomorphisme dans
$D(\mathcal{O}(U))$, pour tout $f : U \to V$ dans $\mathcal{C}$. Ainsi,
$\mathcal{H}^\bullet$ représente un objet de $\mathit{QC}(\mathcal{O})$,
comme voulu.

\medskip\noindent
Construisons $\mathcal{S}_{n + 1}$ à partir de $\mathcal{S}_n$. Pour tout
morphisme $f : U \to V$ de $\mathcal{C}$, considérons le diagramme commutatif
$$
\xymatrix{
\mathcal{S}_n^\bullet(V)
\ar[r] \ar[d] &
\mathcal{S}_n^\bullet(U) \ar[d] \\
\mathcal{F}^\bullet(V)
\ar[r] &
\mathcal{F}^\bullet(U)
}
$$
C'est un diagramme du type considéré dans Compléments d'algèbre, lemme
\ref{more-algebra-lemma-enlarge}, pour l'homomorphisme d'anneaux
$\mathcal{O}(V) \to \mathcal{O}(U)$ ; autrement dit, la ligne inférieure
induit un isomorphisme
$$
\mathcal{F}^\bullet(V) \otimes_{\mathcal{O}(V)}^\mathbf{L} \mathcal{O}(U)
\longrightarrow
\mathcal{F}^\bullet(U)
$$
dans $D(\mathcal{O}(U))$. On peut donc choisir des sous-complexes
$$
\mathcal{S}_n^\bullet(V) \subset M^\bullet_f \subset \mathcal{F}^\bullet(V)
\quad\text{et}\quad
\mathcal{S}_n^\bullet(U) \subset N^\bullet_f \subset \mathcal{F}^\bullet(U)
$$
comme dans Compléments d'algèbre, lemme \ref{more-algebra-lemma-enlarge} ; en
particulier, $|N^i_f|, |M^i_f| \leq \kappa$. Appliquons ensuite le lemme
\ref{sites-cohomology-lemma-build-from} aux sous-ensembles
$$
\mathcal{S}_n^i(U) \amalg \coprod\nolimits_{f : U \to V} N^i_f
\amalg \coprod\nolimits_{g : W \to U} M^i_g
\subset
\mathcal{F}^i(U)
$$
afin d'obtenir un sous-complexe
$$
\mathcal{S}_n^\bullet \subset
\mathcal{S}_{n + 1}^\bullet \subset
\mathcal{F}^\bullet
$$
qui contient ces sous-ensembles et vérifie
$|\mathcal{S}_{n + 1}^\bullet| \leq \kappa$. Les conditions (1) et (2) sont
satisfaites, car les assertions correspondantes valent pour
$\mathcal{S}_n^\bullet(V) \subset M^\bullet_f$ et
$\mathcal{S}_n^\bullet(U) \subset N^\bullet_f$ par la construction de More
on Algèbre, lemme \ref{more-algebra-lemma-enlarge}. La démonstration est donc
achevée.
\end{proof}

\begin{proposition}
\label{sites-cohomology-proposition-cartesian-brown}
Soit $\mathcal{C}$ une catégorie considérée comme un site muni de la
topologie chaotique et soit $\mathcal{O}$ un faisceau d'anneaux sur
$\mathcal{C}$. Pour $\mathit{QC}(\mathcal{O})$ défini comme dans la définition
\ref{sites-cohomology-definition-cartesian}, on a :
\begin{enumerate}
\item $\mathit{QC}(\mathcal{O})$ est une sous-catégorie strictement pleine,
saturée et triangulée de $D(\mathcal{O})$, stable par sommes directes
arbitraires ;
\item tout foncteur cohomologique contravariant
$H : \mathit{QC}(\mathcal{O}) \to \textit{Ab}$
qui transforme les sommes directes en produits est représentable ;
\item tout foncteur exact $F : \mathit{QC}(\mathcal{O}) \to \mathcal{D}$ de
catégories triangulées qui transforme les sommes directes en sommes directes
admet un adjoint à droite exact ;
\item le foncteur d'inclusion
$\mathit{QC}(\mathcal{O}) \to D(\mathcal{O})$ admet un adjoint à droite exact.
\end{enumerate}
\end{proposition}

\begin{proof}
L'assertion (1) est le lemme \ref{sites-cohomology-lemma-cartesian-good-sub}. L'assertion (2)
résulte du lemme \ref{sites-cohomology-lemma-cartesian-colimit-kappa} et de Catégories dérivées,
lemme \ref{derived-lemma-brown-bis}. L'assertion (3) résulte du lemme
\ref{sites-cohomology-lemma-cartesian-colimit-kappa} et de Catégories dérivées, proposition
\ref{derived-proposition-brown-bis}. Enfin, (4) est un cas particulier de
(3).
\end{proof}

\noindent
Soit $u : \mathcal{C}' \to \mathcal{C}$ un foncteur entre catégories. Si
l'on considère $\mathcal{C}$ et $\mathcal{C}'$ comme des sites munis de la
topologie chaotique, alors $u$ est un foncteur continu et cocontinu. On obtient
donc un morphisme de topos
$g : \Sh(\mathcal{C}') \to \Sh(\mathcal{C})$, voir Sites, lemme
\ref{sites-lemma-cocontinuous-morphism-topoi}. Supposons en outre donnés des
faisceaux d'anneaux $\mathcal{O}$ sur $\mathcal{C}$ et $\mathcal{O}'$ sur
$\mathcal{C}'$, ainsi qu'un morphisme
$g^\sharp : g^{-1}\mathcal{O} \to \mathcal{O}'$. On note simplement le
morphisme correspondant de topos annelés par
$g : (\Sh(\mathcal{C}'), \mathcal{O}') \to (\Sh(\mathcal{C}), \mathcal{O})$,
voir Modules sur les sites, section \ref{sites-modules-section-ringed-topoi}.

\begin{lemma}
\label{sites-cohomology-lemma-cartesian-functorial}
Soit
$g : (\Sh(\mathcal{C}'), \mathcal{O}') \to (\Sh(\mathcal{C}), \mathcal{O})$
comme ci-dessus. Alors le foncteur
$Lg^* : D(\mathcal{O}) \to D(\mathcal{O}')$ envoie
$\mathit{QC}(\mathcal{O})$ dans $\mathit{QC}(\mathcal{O}')$.
\end{lemma}

\begin{proof}
Soit $U' \in \Ob(\mathcal{C}')$, d'image $U = u(U')$ dans $\mathcal{C}$. Notons
$pt$ la catégorie ayant un seul objet et un seul morphisme, et notons
$(\Sh(pt), \mathcal{O}'(U'))$ et $(\Sh(pt), \mathcal{O}(U))$ les topos annelés
indiqués. On identifie bien entendu les catégories dérivées de modules sur
ces topos annelés à $D(\mathcal{O}'(U'))$ et $D(\mathcal{O}(U))$. On dispose
alors d'un diagramme commutatif de topos annelés
$$
\xymatrix{
(\Sh(pt), \mathcal{O}'(U')) \ar[rr]_{U'} \ar[d] & &
(\Sh(\mathcal{C}'), \mathcal{O}') \ar[d]^g \\
(\Sh(pt), \mathcal{O}(U)) \ar[rr]^U & &
(\Sh(\mathcal{C}), \mathcal{O})
}
$$
L'image inverse le long du morphisme horizontal inférieur envoie $K$ dans
$D(\mathcal{O})$ sur $R\Gamma(U, K)$. Celle le long de la flèche verticale de
gauche envoie $M$ sur
$M \otimes_{\mathcal{O}(U)}^\mathbf{L} \mathcal{O}'(U')$. Les deux parcours du
diagramme donnent le même résultat (lemme
\ref{sites-cohomology-lemma-derived-pullback-composition}) ; par conséquent,
$$
R\Gamma(U', Lg^*K) =
R\Gamma(U, K) \otimes_{\mathcal{O}(U)}^\mathbf{L} \mathcal{O}'(U')
$$
Enfin, soit $f' : U' \to V'$ un morphisme de $\mathcal{C}'$ et notons
$f = u(f') : U = u(U') \to V = u(V')$ son image dans $\mathcal{C}$. Si $K$ est
dans $\mathit{QC}(\mathcal{O})$, alors
\begin{align*}
R\Gamma(V', Lg^*K) \otimes_{\mathcal{O}'(V')}^\mathbf{L} \mathcal{O}'(U')
& =
R\Gamma(V, K) \otimes_{\mathcal{O}(V)}^\mathbf{L} \mathcal{O}'(V')
\otimes_{\mathcal{O}'(V')}^\mathbf{L} \mathcal{O}'(U') \\
& =
R\Gamma(V, K) \otimes_{\mathcal{O}(V)}^\mathbf{L} \mathcal{O}'(U') \\
& =
R\Gamma(V, K) \otimes_{\mathcal{O}(V)}^\mathbf{L} \mathcal{O}(U)
\otimes_{\mathcal{O}(U)}^\mathbf{L} \mathcal{O}'(U') \\
& =
R\Gamma(U, K) \otimes_{\mathcal{O}(U)}^\mathbf{L} \mathcal{O}'(U') \\
& =
R\Gamma(U', Lg^*K)
\end{align*}
comme voulu. Nous avons utilisé l'observation précédente à la fois pour $U'$
et pour $V'$.
\end{proof}

\begin{lemma}
\label{sites-cohomology-lemma-cartesian-flat-quasi-coherent}
Soit $\mathcal{C}$ une catégorie considérée comme un site muni de la
topologie chaotique et soit $\mathcal{O}$ un faisceau d'anneaux sur
$\mathcal{C}$. Supposons que, pour tout $U \to V$ dans $\mathcal{C}$,
l'homomorphisme de restriction
$\mathcal{O}(V) \to \mathcal{O}(U)$ soit plat. Alors
$\mathit{QC}(\mathcal{O})$ coïncide avec la sous-catégorie
$D_\QCoh(\mathcal{O}) \subset D(\mathcal{O})$ des complexes dont les
faisceaux de cohomologie sont quasi-cohérents.
\end{lemma}

\begin{proof}
Rappelons que, sous nos hypothèses,
$\QCoh(\mathcal{O}) \subset \textit{Mod}(\mathcal{O})$ est une sous-catégorie
de Serre faible, voir Modules sur les sites, lemme
\ref{sites-modules-lemma-chaotic-flat-quasi-coherent}. Il est donc légitime de
considérer la sous-catégorie pleine
$$
D_\QCoh(\mathcal{O}) = D_{\QCoh(\mathcal{O})}(\textit{Mod}(\mathcal{O}))
$$
de $D(\mathcal{O})$, voir Catégories dérivées, section
\ref{derived-section-triangulated-sub}. (À strictement parler, ceci n'est pas
nécessaire dans la démonstration du lemme.)

\medskip\noindent
Soit $M$ un objet de $\mathit{QC}(\mathcal{O})$. Pour tout morphisme
$U \to V$ de $\mathcal{C}$, l'homomorphisme de restriction
$\mathcal{O}(V) \to \mathcal{O}(U)$ est plat ; on a donc
\begin{align*}
H^i(M)(U)
& =
H^i(R\Gamma(U, M)) \\
& =
H^i(R\Gamma(V, M) \otimes_{\mathcal{O}(V)}^\mathbf{L} \mathcal{O}(U)) \\
& =
H^i(R\Gamma(V, M)) \otimes_{\mathcal{O}(V)} \mathcal{O}(U) \\
& =
H^i(M)(V) \otimes_{\mathcal{O}(V)} \mathcal{O}(U)
\end{align*}
et $H^i(M)$ est donc quasi-cohérent d'après Modules sur les sites, lemme
\ref{sites-modules-lemma-chaotic-quasi-coherent}. La première et la dernière
égalité ci-dessus résultent du fait que la prise des sections sur un objet de
$\mathcal{C}$ est un foncteur exact, puisque la topologie sur $\mathcal{C}$
est chaotique.

\medskip\noindent
Réciproquement, si $M$ est un objet de $D_\QCoh(\mathcal{O})$, alors Modules
on Sites, lemme \ref{sites-modules-lemma-chaotic-quasi-coherent}, montre que
l'application $R\Gamma(V, M) \to R\Gamma(U, M)$ induit des isomorphismes
$H^i(M)(V) \otimes_{\mathcal{O}(V)} \mathcal{O}(U) \to H^i(M)(U)$.
Par conséquent,
$R\Gamma(V, M) \otimes_{\mathcal{O}(V)}^\mathbf{L} \mathcal{O}(U)
\to R\Gamma(U, M)$ est un isomorphisme dans $D(\mathcal{O}(U))$, par platitude
de $\mathcal{O}(V) \to \mathcal{O}(U)$, et l'on conclut que $M$ appartient à
$\mathit{QC}(\mathcal{O})$.
\end{proof}

\begin{lemma}
\label{sites-cohomology-lemma-cartesion-plus-topology}
Soit $\epsilon : (\mathcal{C}_\tau, \mathcal{O}_\tau) \to
(\mathcal{C}_{\tau'}, \mathcal{O}_{\tau'})$ comme dans la section
\ref{sites-cohomology-section-compare}. Supposons que
\begin{enumerate}
\item $\tau'$ soit la topologie chaotique sur la catégorie $\mathcal{C}$ ;
\item pour tout $U \in \Ob(\mathcal{C})$ et tout complexe K-plat de
$\mathcal{O}(U)$-modules $M^\bullet$, l'application
$$
M^\bullet \longrightarrow
R\Gamma((\mathcal{C}/U)_\tau,
(M^\bullet \otimes_{\mathcal{O}(U)} \mathcal{O}_U)^\#)
$$
soit un quasi-isomorphisme (voir la démonstration pour une explication).
\end{enumerate}
Alors $\epsilon^*$ et $R\epsilon_*$ définissent des équivalences quasi-inverses
l'une de l'autre entre $\mathit{QC}(\mathcal{O})$ et la sous-catégorie pleine
de $D(\mathcal{C}_\tau, \mathcal{O}_\tau)$ formée des objets $K$ tels que
$R\epsilon_*K$ appartienne à $\mathit{QC}(\mathcal{O})$\footnote{Cela signifie
que
$R\Gamma(V, K) \otimes_{\mathcal{O}(V)}^\mathbf{L} \mathcal{O}(U)
\to R\Gamma(U, K)$ est un isomorphisme pour tout $U \to V$ dans
$\mathcal{C}$.}.
\end{lemma}

\begin{proof}
Nous utiliserons sans autre mention les observations de la section
\ref{sites-cohomology-section-compare}. Comme $R\epsilon_*$ est pleinement fidèle et
$\epsilon^* \circ R\epsilon_* = \text{id}$, il suffit, pour démontrer le
lemme, d'établir que, pour $M$ dans $\mathit{QC}(\mathcal{O})$, on a
$R\epsilon_*(\epsilon^*M) = M$. La
condition (2) est précisément celle qui permet de le voir. Choisissons en
effet un complexe K-plat $\mathcal{M}^\bullet$ de $\mathcal{O}$-modules, à
termes plats, qui représente $M$. Alors $\epsilon^*M$ est représenté par la
$\tau$-faisceautisation $(\mathcal{M}^\bullet)^\#$ de
$\mathcal{M}^\bullet$. Soit $U \in \Ob(\mathcal{C})$. Par Leray, on obtient
$$
R\Gamma(U, R\epsilon_*(\epsilon^*M)) =
R\Gamma((\mathcal{C}/U)_\tau, (\mathcal{M}^\bullet)^\#|_{\mathcal{C}/U}) =
R\Gamma((\mathcal{C}/U)_\tau, (\mathcal{M}^\bullet|_{\mathcal{C}/U})^\#)
$$
La dernière égalité résulte du fait que la faisceautisation commute à la
restriction à $\mathcal{C}/U$. Comme d'habitude, notons $\mathcal{O}_U$ la
restriction de $\mathcal{O}$ à $\mathcal{C}/U$. Considérons l'application
$$
\mathcal{M}^\bullet(U) \otimes_{\mathcal{O}(U)} \mathcal{O}_U
\longrightarrow
\mathcal{M}^\bullet|_{\mathcal{C}/U}
$$
de complexes de $\mathcal{O}_U$-modules (pour la topologie $\tau'$). Par le
choix de $\mathcal{M}^\bullet$, le complexe $\mathcal{M}^\bullet(U)$ est un
complexe K-plat de $\mathcal{O}(U)$-modules ; voir le lemme
\ref{sites-cohomology-lemma-pullback-K-flat} et utiliser que l'inclusion de $U$ dans
$\mathcal{C}$ définit un morphisme de topos annelés
$(\Sh(pt), \mathcal{O}(U)) \to (\Sh(\mathcal{C}_{\tau'}), \mathcal{O})$.
Comme $M$ appartient à $\mathit{QC}(\mathcal{O})$, la flèche affichée est un
quasi-isomorphisme. La faisceautisation étant exacte, elle reste un
quasi-isomorphisme après faisceautisation. Par conséquent,
$$
R\Gamma(U, R\epsilon_*(\epsilon^*M)) =
R\Gamma((\mathcal{C}/U)_\tau,
(\mathcal{M}^\bullet(U) \otimes_{\mathcal{O}(U)} \mathcal{O}_U)^\#)
$$
et l'hypothèse (2) montre que ceci est égal à
$R\Gamma(U, M) = \mathcal{M}^\bullet(U)$, comme voulu.
\end{proof}

\begin{lemma}
\label{sites-cohomology-lemma-QC-hom-out-of-plus-topology}
Conservons les notations et les hypothèses du lemme
\ref{sites-cohomology-lemma-cartesion-plus-topology}. Supposons que $K$ soit un objet de
$\mathit{QC}(\mathcal{O})$ et que $M$ soit un objet quelconque de
$D(\mathcal{O}_\tau)$. Pour tout objet $U$ de $\mathcal{C}$, on a
$$
\Hom_{D((\mathcal{O}_U)_\tau)}(\epsilon^*K|_U, M|_U) =
R\Hom_{\mathcal{O}(U)}(R\Gamma(U, K), R\Gamma(U, M))
$$
\end{lemma}

\begin{proof}
On a
$$
\Hom_{D((\mathcal{O}_U)_\tau)}(\epsilon^*K|_U, M|_U) =
\Hom_{D((\mathcal{O}_U)_{\tau'})}(K|_U, R\epsilon_*M|_U)
$$
par adjonction. Le résultat découle donc du lemme
\ref{sites-cohomology-lemma-QC-hom-out-of} et du fait que
$$
R\Gamma(U, M) = R\Gamma(U, R\epsilon_*M)
$$
par Leray.
\end{proof}







\section{Complexes strictement parfaits}
\label{sites-cohomology-section-strictly-perfect}

\noindent
Cette section est l'analogue de Cohomologie, section
\ref{cohomology-section-strictly-perfect}.

\begin{definition}
\label{sites-cohomology-definition-strictly-perfect}
Soit $(\mathcal{C}, \mathcal{O})$ un site annelé et soit
$\mathcal{E}^\bullet$ un complexe de $\mathcal{O}$-modules. On dit que
$\mathcal{E}^\bullet$ est {\it strictement parfait} si
$\mathcal{E}^i$ est nul sauf pour un nombre fini de $i$ et si
$\mathcal{E}^i$ est facteur direct d'un $\mathcal{O}$-module libre de type fini
pour tout $i$.
\end{definition}

\noindent
Soit $U$ un objet de $\mathcal{C}$. Nous dirons souvent « Soit
$\mathcal{E}^\bullet$ un complexe strictement parfait de
$\mathcal{O}_U$-modules » pour signifier que $\mathcal{E}^\bullet$ est un
complexe strictement parfait de modules sur le site annelé
$(\mathcal{C}/U, \mathcal{O}_U)$, voir Modules sur les sites, définition
\ref{sites-modules-definition-localize-ringed-site}.

\begin{lemma}
\label{sites-cohomology-lemma-cone}
Le cône d'un morphisme de complexes strictement parfaits est strictement
parfait.
\end{lemma}

\begin{proof}
C'est immédiat d'après les définitions.
\end{proof}

\begin{lemma}
\label{sites-cohomology-lemma-tensor}
Le complexe total associé au produit tensoriel de deux complexes strictement
parfaits est strictement parfait.
\end{lemma}

\begin{proof}
Omis.
\end{proof}

\begin{lemma}
\label{sites-cohomology-lemma-strictly-perfect-pullback}
Soit $(f, f^\sharp) : (\mathcal{C}, \mathcal{O}_\mathcal{C}) \to
(\mathcal{D}, \mathcal{O}_\mathcal{D})$
un morphisme de topos annelés. Si $\mathcal{F}^\bullet$ est un complexe
strictement parfait de $\mathcal{O}_\mathcal{D}$-modules, alors
$f^*\mathcal{F}^\bullet$ est un complexe strictement parfait de
$\mathcal{O}_\mathcal{C}$-modules.
\end{lemma}

\begin{proof}
Nous avons vu dans Modules sur les sites, lemme
\ref{sites-modules-lemma-global-pullback}, que l'image inverse d'un module
libre de type fini est libre de type fini. Le foncteur $f^*$ est additif et
préserve donc les facteurs directs. Le lemme en résulte.
\end{proof}

\begin{lemma}
\label{sites-cohomology-lemma-local-lift-map}
Soit $(\mathcal{C}, \mathcal{O})$ un site annelé et soit $U$ un objet de
$\mathcal{C}$. Étant donné un diagramme de $\mathcal{O}_U$-modules dont les
flèches pleines sont
$$
\xymatrix{
\mathcal{E} \ar@{..>}[dr] \ar[r] & \mathcal{F} \\
& \mathcal{G} \ar[u]_p
}
$$
où $\mathcal{E}$ est facteur direct d'un $\mathcal{O}_U$-module libre de type
fini et $p$ est surjectif, il existe un recouvrement $\{U_i \to U\}$ tel que,
sur chaque $U_i$, on puisse tracer une flèche pointillée rendant le diagramme
commutatif.
\end{lemma}

\begin{proof}
On peut supposer $\mathcal{E} = \mathcal{O}_U^{\oplus n}$ pour un certain $n$.
Trouver la flèche pointillée revient alors à relever les images des éléments
de la base dans $\Gamma(U, \mathcal{F})$. Cela est possible localement d'après
la caractérisation des morphismes surjectifs de faisceaux (Sites, section
\ref{sites-section-sheaves-injective}).
\end{proof}

\begin{lemma}
\label{sites-cohomology-lemma-local-homotopy}
Soit $(\mathcal{C}, \mathcal{O})$ un site annelé et soit $U$ un objet de
$\mathcal{C}$.
\begin{enumerate}
\item Soit $\alpha : \mathcal{E}^\bullet \to \mathcal{F}^\bullet$ un
morphisme de complexes de $\mathcal{O}_U$-modules, avec
$\mathcal{E}^\bullet$ strictement parfait et $\mathcal{F}^\bullet$ acyclique.
Il existe un recouvrement $\{U_i \to U\}$ tel que chaque
$\alpha|_{U_i}$ soit homotope à zéro.
\item Soit $\alpha : \mathcal{E}^\bullet \to \mathcal{F}^\bullet$ un
morphisme de complexes de $\mathcal{O}_U$-modules, avec
$\mathcal{E}^\bullet$ strictement parfait, $\mathcal{E}^i = 0$ pour $i < a$, et
$H^i(\mathcal{F}^\bullet) = 0$ pour $i \geq a$. Il existe un recouvrement
$\{U_i \to U\}$ tel que chaque $\alpha|_{U_i}$ soit homotope à zéro.
\end{enumerate}
\end{lemma}

\begin{proof}
La première assertion résulte de la seconde ; il suffit donc de démontrer
(2). Procédons par récurrence sur la longueur du complexe
$\mathcal{E}^\bullet$. Si
$\mathcal{E}^\bullet \cong \mathcal{E}[-n]$, où $\mathcal{E}$ est facteur
direct d'un $\mathcal{O}$-module libre de type fini et $n \geq a$, alors le
résultat découle du lemme \ref{sites-cohomology-lemma-local-lift-map} et du fait que
$\mathcal{F}^{n - 1} \to \Ker(\mathcal{F}^n \to \mathcal{F}^{n + 1})$
est surjective par l'annulation supposée de
$H^n(\mathcal{F}^\bullet)$. Si $\mathcal{E}^i$ est nul sauf pour
$i \in [a, b]$, on dispose d'une suite exacte scindée de complexes
$$
0 \to \mathcal{E}^b[-b] \to \mathcal{E}^\bullet \to
\sigma_{\leq b - 1}\mathcal{E}^\bullet \to 0
$$
qui détermine un triangle distingué dans $K(\mathcal{O}_U)$, donc une suite
exacte
$$
\Hom_{K(\mathcal{O}_U)}(
\sigma_{\leq b - 1}\mathcal{E}^\bullet, \mathcal{F}^\bullet)
\to
\Hom_{K(\mathcal{O}_U)}(\mathcal{E}^\bullet, \mathcal{F}^\bullet)
\to
\Hom_{K(\mathcal{O}_U)}(\mathcal{E}^b[-b], \mathcal{F}^\bullet)
$$
d'après les axiomes des catégories triangulées. D'après ce qui précède, la
composée $\mathcal{E}^b[-b] \to \mathcal{F}^\bullet$ est homotope à zéro sur
les membres d'un recouvrement de $U$. On peut donc supposer que notre
morphisme provient d'un élément du membre de gauche de la suite exacte
affichée. Par l'hypothèse de récurrence, cet élément est nul sur les membres
d'un recouvrement de $U$.
\end{proof}

\begin{lemma}
\label{sites-cohomology-lemma-lift-through-quasi-isomorphism}
Soit $(\mathcal{C}, \mathcal{O})$ un site annelé et soit $U$ un objet de
$\mathcal{C}$. Étant donné un diagramme de complexes de
$\mathcal{O}_U$-modules dont les flèches pleines sont
$$
\xymatrix{
\mathcal{E}^\bullet \ar@{..>}[dr] \ar[r]_\alpha & \mathcal{F}^\bullet \\
& \mathcal{G}^\bullet \ar[u]_f
}
$$
où $\mathcal{E}^\bullet$ est strictement parfait,
$\mathcal{E}^j = 0$ pour $j < a$, et $H^j(f)$ est un isomorphisme pour $j > a$
et une surjection pour $j = a$, il existe un recouvrement $\{U_i \to U\}$ et,
pour
chaque $i$, une flèche pointillée sur $U_i$ qui rend le diagramme commutatif à
homotopie près.
\end{lemma}

\begin{proof}
Les hypothèses sur $f$ impliquent que les faisceaux de cohomologie du cône
$C(f)^\bullet$ s'annulent aux degrés $\geq a$. Le lemme
\ref{sites-cohomology-lemma-local-homotopy} fournit donc un recouvrement $\{U_i \to U\}$ tel que
la composée
$\mathcal{E}^\bullet \to \mathcal{F}^\bullet \to C(f)^\bullet$
soit homotope à zéro sur $U_i$. Comme
$$
\mathcal{G}^\bullet \to \mathcal{F}^\bullet \to C(f)^\bullet \to
\mathcal{G}^\bullet[1]
$$
se restreint en un triangle distingué dans $K(\mathcal{O}_{U_i})$, on peut
relever $\alpha|_{U_i}$, à homotopie près, en un morphisme
$\alpha_i : \mathcal{E}^\bullet|_{U_i} \to \mathcal{G}^\bullet|_{U_i}$
comme voulu.
\end{proof}

\begin{lemma}
\label{sites-cohomology-lemma-local-actual}
Soit $(\mathcal{C}, \mathcal{O})$ un site annelé, soit $U$ un objet de
$\mathcal{C}$, et soient $\mathcal{E}^\bullet$, $\mathcal{F}^\bullet$ des
complexes de $\mathcal{O}_U$-modules, avec $\mathcal{E}^\bullet$ strictement
parfait.
\begin{enumerate}
\item Pour tout élément
$\alpha \in \Hom_{D(\mathcal{O}_U)}(\mathcal{E}^\bullet, \mathcal{F}^\bullet)$
il existe un recouvrement $\{U_i \to U\}$ tel que $\alpha|_{U_i}$ soit donné
par un morphisme de complexes
$\alpha_i : \mathcal{E}^\bullet|_{U_i} \to \mathcal{F}^\bullet|_{U_i}$.
\item Soit un morphisme de complexes
$\alpha : \mathcal{E}^\bullet \to \mathcal{F}^\bullet$ dont l'image dans le
groupe
$\Hom_{D(\mathcal{O}_U)}(\mathcal{E}^\bullet, \mathcal{F}^\bullet)$
est nulle. Il existe un recouvrement $\{U_i \to U\}$ tel que
$\alpha|_{U_i}$ soit homotope à zéro.
\end{enumerate}
\end{lemma}

\begin{proof}
Démontrons (1). Par la construction de la catégorie dérivée, on peut trouver
un quasi-isomorphisme
$f : \mathcal{F}^\bullet \to \mathcal{G}^\bullet$ et un morphisme de
complexes $\beta : \mathcal{E}^\bullet \to \mathcal{G}^\bullet$ tels que
$\alpha = f^{-1}\beta$. Le résultat découle alors du lemme
\ref{sites-cohomology-lemma-lift-through-quasi-isomorphism}. Nous omettons la démonstration de
(2).
\end{proof}

\begin{lemma}
\label{sites-cohomology-lemma-Rhom-strictly-perfect}
Soit $(\mathcal{C}, \mathcal{O})$ un site annelé. Soient
$\mathcal{E}^\bullet$, $\mathcal{F}^\bullet$ des complexes de
$\mathcal{O}$-modules, avec $\mathcal{E}^\bullet$ strictement parfait. Alors le
Hom interne $R\SheafHom(\mathcal{E}^\bullet, \mathcal{F}^\bullet)$ est
représenté par le complexe $\mathcal{H}^\bullet$ dont les termes sont
$$
\mathcal{H}^n =
\bigoplus\nolimits_{n = p + q}
\SheafHom_\mathcal{O}(\mathcal{E}^{-q}, \mathcal{F}^p)
$$
et dont la différentielle est décrite dans la section
\ref{sites-cohomology-section-internal-hom}.
\end{lemma}

\begin{proof}
Choisissons un quasi-isomorphisme
$\mathcal{F}^\bullet \to \mathcal{I}^\bullet$ vers un complexe K-injectif.
Soit $(\mathcal{H}')^\bullet$ le complexe dont les termes sont
$$
(\mathcal{H}')^n =
\prod\nolimits_{n = p + q}
\SheafHom_\mathcal{O}(\mathcal{E}^{-q}, \mathcal{I}^p)
$$
Il représente $R\SheafHom(\mathcal{E}^\bullet, \mathcal{F}^\bullet)$ par la
construction de la section \ref{sites-cohomology-section-internal-hom}. Il suffit de montrer
que l'application
$$
\mathcal{H}^\bullet \longrightarrow (\mathcal{H}')^\bullet
$$
est un quasi-isomorphisme. Pour un objet $U$ de $\mathcal{C}$, on constate que
$$
H^0(\mathcal{H}^\bullet(U)) =
\Hom_{K(\mathcal{O}_U)}(\mathcal{E}^\bullet|_U, \mathcal{I}^\bullet|_U)
\to
H^0((\mathcal{H}')^\bullet(U)) =
\Hom_{D(\mathcal{O}_U)}(\mathcal{E}^\bullet|_U, \mathcal{I}^\bullet|_U)
$$
D'après le lemme \ref{sites-cohomology-lemma-local-actual}, la faisceautisation de
$U \mapsto H^0(\mathcal{H}^\bullet(U))$ est égale à celle de
$U \mapsto H^0((\mathcal{H}')^\bullet(U))$. Le même argument vaut pour les
autres faisceaux de cohomologie. Ainsi, $\mathcal{H}^\bullet$ est
quasi-isomorphe à $(\mathcal{H}')^\bullet$, ce qui démontre le lemme.
\end{proof}

\begin{lemma}
\label{sites-cohomology-lemma-Rhom-complex-of-direct-summands-finite-free}
Soit $(\mathcal{C}, \mathcal{O})$ un site annelé. Soient
$\mathcal{E}^\bullet$, $\mathcal{F}^\bullet$ des complexes de
$\mathcal{O}$-modules tels que
\begin{enumerate}
\item $\mathcal{F}^n = 0$ pour $n \ll 0$ ;
\item $\mathcal{E}^n = 0$ pour $n \gg 0$ ;
\item $\mathcal{E}^n$ soit isomorphe à un facteur direct d'un
$\mathcal{O}$-module libre de type fini.
\end{enumerate}
Alors le Hom interne
$R\SheafHom(\mathcal{E}^\bullet, \mathcal{F}^\bullet)$ est représenté par le
complexe $\mathcal{H}^\bullet$ dont les termes sont
$$
\mathcal{H}^n =
\bigoplus\nolimits_{n = p + q}
\SheafHom_\mathcal{O}(\mathcal{E}^{-q}, \mathcal{F}^p)
$$
et dont la différentielle est décrite dans la section
\ref{sites-cohomology-section-internal-hom}.
\end{lemma}

\begin{proof}
Choisissons un quasi-isomorphisme
$\mathcal{F}^\bullet \to \mathcal{I}^\bullet$, où $\mathcal{I}^\bullet$ est un
complexe d'injectifs borné inférieurement. Remarquons que
$\mathcal{I}^\bullet$ est K-injectif (Catégories dérivées, lemme
\ref{derived-lemma-bounded-below-injectives-K-injective}). La construction de
la section \ref{sites-cohomology-section-internal-hom} montre donc que
$R\SheafHom(\mathcal{E}^\bullet, \mathcal{F}^\bullet)$ est représenté par le
complexe $(\mathcal{H}')^\bullet$ dont les termes sont
$$
(\mathcal{H}')^n =
\prod\nolimits_{n = p + q}
\SheafHom_\mathcal{O}(\mathcal{E}^{-q}, \mathcal{I}^p) =
\bigoplus\nolimits_{n = p + q}
\SheafHom_\mathcal{O}(\mathcal{E}^{-q}, \mathcal{I}^p)
$$
Cette égalité vaut parce qu'il n'y a qu'un nombre fini de termes non nuls.
Remarquons que $\mathcal{H}^\bullet$ est le complexe total associé au complexe
double de termes
$\SheafHom_\mathcal{O}(\mathcal{E}^{-q}, \mathcal{F}^p)$
et de même pour $(\mathcal{H}')^\bullet$. L'application naturelle
$\mathcal{H}^\bullet \to (\mathcal{H}')^\bullet$ provient d'un morphisme de
complexes doubles. Pour montrer qu'elle est un quasi-isomorphisme, on peut
donc utiliser la suite spectrale d'un complexe double (Homology, lemme
\ref{homology-lemma-first-quadrant-ss})
$$
{}'E_1^{p, q} =
H^p(\SheafHom_\mathcal{O}(\mathcal{E}^{-q}, \mathcal{F}^\bullet))
$$
qui converge vers $H^{p + q}(\mathcal{H}^\bullet)$, et de même pour
$(\mathcal{H}')^\bullet$. Pour achever la démonstration, il suffit de montrer
que $\mathcal{F}^\bullet \to \mathcal{I}^\bullet$ induit un isomorphisme
$$
H^p(\SheafHom_\mathcal{O}(\mathcal{E}, \mathcal{F}^\bullet))
\longrightarrow
H^p(\SheafHom_\mathcal{O}(\mathcal{E}, \mathcal{I}^\bullet))
$$
sur les faisceaux de cohomologie lorsque $\mathcal{E}$ est facteur direct
d'un $\mathcal{O}$-module libre de type fini. C'est clair lorsque
$\mathcal{E}$ est libre de type fini, d'où le résultat.
\end{proof}









\section{Modules pseudo-cohérents}
\label{sites-cohomology-section-pseudo-coherent}

\noindent
Dans cette section, nous étudions les complexes pseudo-cohérents.

\begin{definition}
\label{sites-cohomology-definition-pseudo-coherent}
Soit $(\mathcal{C}, \mathcal{O})$ un site annelé, soit
$\mathcal{E}^\bullet$ un complexe de $\mathcal{O}$-modules et soit
$m \in \mathbf{Z}$.
\begin{enumerate}
\item On dit que $\mathcal{E}^\bullet$ est {\it $m$-pseudo-cohérent} si, pour
tout objet $U$ de $\mathcal{C}$, il existe un recouvrement $\{U_i \to U\}$ et,
pour chaque $i$, un morphisme de complexes
$\alpha_i : \mathcal{E}_i^\bullet \to \mathcal{E}^\bullet|_{U_i}$
où $\mathcal{E}_i^\bullet$ est un complexe strictement parfait de
$\mathcal{O}_{U_i}$-modules, $H^j(\alpha_i)$ est un isomorphisme pour $j > m$ et
$H^m(\alpha_i)$ est surjectif.
\item On dit que $\mathcal{E}^\bullet$ est {\it pseudo-cohérent} s'il est
$m$-pseudo-cohérent pour tout $m$.
\item On dit qu'un objet $E$ de $D(\mathcal{O})$ est
{\it $m$-pseudo-cohérent} (resp.\ {\it pseudo-cohérent}) si et seulement s'il
peut être représenté par un complexe $m$-pseudo-cohérent
(resp.\ pseudo-cohérent) de $\mathcal{O}$-modules.
\end{enumerate}
\end{definition}

\noindent
Si $\mathcal{C}$ possède un objet final quasi-compact $X$ (par exemple, si
tout recouvrement de $X$ admet un raffinement fini), alors tout objet
$m$-pseudo-cohérent de $D(\mathcal{O})$ appartient à $D^-(\mathcal{O})$.
Ce n'est toutefois pas nécessairement le cas en général.

\begin{lemma}
\label{sites-cohomology-lemma-pseudo-coherent-independent-representative}
Soit $(\mathcal{C}, \mathcal{O})$ un site annelé.
Soit $E$ un objet de $D(\mathcal{O})$.
\begin{enumerate}
\item Si $\mathcal{C}$ possède un objet final $X$ et s'il existe un recouvrement
$\{U_i \to X\}$, des complexes strictement parfaits $\mathcal{E}_i^\bullet$
de $\mathcal{O}_{U_i}$-modules et des morphismes
$\alpha_i : \mathcal{E}_i^\bullet \to E|_{U_i}$ dans
$D(\mathcal{O}_{U_i})$ tels que $H^j(\alpha_i)$ soit un isomorphisme pour
$j > m$ et $H^m(\alpha_i)$ soit surjectif, alors $E$ est
$m$-pseudo-cohérent.
\item Si $E$ est $m$-pseudo-cohérent, alors tout complexe de
$\mathcal{O}$-modules représentant $E$ est $m$-pseudo-cohérent.
\item Si, pour tout objet $U$ de $\mathcal{C}$, il existe un recouvrement
$\{U_i \to U\}$ tel que $E|_{U_i}$ soit $m$-pseudo-cohérent, alors
$E$ est $m$-pseudo-cohérent.
\end{enumerate}
\end{lemma}

\begin{proof}
Soit $\mathcal{F}^\bullet$ un complexe quelconque représentant $E$, et
soient $X$, $\{U_i \to X\}$ et
$\alpha_i : \mathcal{E}_i \to E|_{U_i}$ comme en (1). Nous allons montrer
que $\mathcal{F}^\bullet$ est $m$-pseudo-cohérent comme complexe, ce qui
prouvera (1) et (2) lorsque $\mathcal{C}$ possède un objet final. Par le lemme
\ref{sites-cohomology-lemma-local-actual}, après avoir raffiné le recouvrement
$\{U_i \to X\}$, nous pouvons représenter les morphismes $\alpha_i$ par des
morphismes de complexes
$\alpha_i : \mathcal{E}_i^\bullet \to \mathcal{F}^\bullet|_{U_i}$.
Par hypothèse, les $H^j(\alpha_i)$ sont des isomorphismes pour $j > m$ et
$H^m(\alpha_i)$ est surjectif ; ainsi, $\mathcal{F}^\bullet$ est
$m$-pseudo-cohérent.

\medskip\noindent
Démonstration de (2). D'après ce qui précède,
$\mathcal{F}^\bullet|_U$ est $m$-pseudo-cohérent comme complexe de
$\mathcal{O}_U$-modules pour tout objet $U$ de $\mathcal{C}$. Il résulte
formellement des définitions que $\mathcal{F}^\bullet$ est
$m$-pseudo-cohérent.

\medskip\noindent
Démonstration de (3). Cela résulte des définitions et de Sites,
définition \ref{sites-definition-site}, partie (2).
\end{proof}

\begin{lemma}
\label{sites-cohomology-lemma-pseudo-coherent-pullback}
Soit $(f, f^\sharp) : (\mathcal{C}, \mathcal{O}_\mathcal{C}) \to
(\mathcal{D}, \mathcal{O}_\mathcal{D})$
un morphisme de sites annelés. Soit $E$ un objet de
$D(\mathcal{O}_\mathcal{C})$. Si $E$ est $m$-pseudo-cohérent,
alors $Lf^*E$ est $m$-pseudo-cohérent.
\end{lemma}

\begin{proof}
Supposons $f$ donné par le foncteur $u : \mathcal{D} \to \mathcal{C}$.
Soit $U$ un objet de $\mathcal{C}$. Par Sites, lemme
\ref{sites-lemma-morphism-of-sites-covering}, nous pouvons trouver un
recouvrement $\{U_i \to U\}$ et, pour chaque $i$, un morphisme
$U_i \to u(V_i)$ pour un certain objet $V_i$ de $\mathcal{D}$.
Par le lemme \ref{sites-cohomology-lemma-pseudo-coherent-independent-representative},
il suffit de montrer que $Lf^*E|_{U_i}$ est $m$-pseudo-cohérent.
Pour cela, il suffit de montrer que $Lf^*E|_{u(V_i)}$ est
$m$-pseudo-cohérent, puisque $Lf^*E|_{U_i}$ est la restriction de
$Lf^*E|_{u(V_i)}$ à $\mathcal{C}/U_i$ (via Modules sur les sites, lemme
\ref{sites-modules-lemma-relocalize}).
Par le diagramme commutatif de Modules sur les sites, lemme
\ref{sites-modules-lemma-localize-morphism-ringed-sites}, il suffit de
démontrer le lemme pour le morphisme de sites annelés
$(\mathcal{C}/u(V_i), \mathcal{O}_{u(V_i)}) \to
(\mathcal{D}/V_i, \mathcal{O}_{V_i})$.
Nous pouvons donc supposer que $\mathcal{D}$ possède un objet final $Y$ tel
que $X = u(Y)$ soit un objet final de $\mathcal{C}$.

\medskip\noindent
Soit $\{V_i \to Y\}$ un recouvrement tel que, pour chaque $i$, il existe
un complexe strictement parfait $\mathcal{F}_i^\bullet$ de
$\mathcal{O}_{V_i}$-modules et un morphisme
$\alpha_i : \mathcal{F}_i^\bullet \to E|_{V_i}$ de $D(\mathcal{O}_{V_i})$
tel que $H^j(\alpha_i)$ soit un isomorphisme pour $j > m$ et
$H^m(\alpha_i)$ soit surjectif. En raisonnant comme ci-dessus, il suffit de
démontrer le résultat pour
$(\mathcal{C}/u(V_i), \mathcal{O}_{u(V_i)}) \to
(\mathcal{D}/V_i, \mathcal{O}_{V_i})$. Nous pouvons donc supposer qu'il
existe un complexe strictement parfait $\mathcal{F}^\bullet$ de
$\mathcal{O}_\mathcal{D}$-modules et un morphisme
$\alpha : \mathcal{F}^\bullet \to E$ de $D(\mathcal{O}_\mathcal{D})$ tel
que $H^j(\alpha)$ soit un isomorphisme pour $j > m$ et que
$H^m(\alpha)$ soit surjectif. Dans ce cas, choisissons un triangle distingué
$$
\mathcal{F}^\bullet \to E \to C \to \mathcal{F}^\bullet[1]
$$
L'hypothèse sur $\alpha$ signifie exactement que les faisceaux de cohomologie
$H^j(C)$ sont nuls pour tout $j \geq m$. En appliquant $Lf^*$, nous obtenons
le triangle distingué
$$
Lf^*\mathcal{F}^\bullet \to Lf^*E \to Lf^*C \to Lf^*\mathcal{F}^\bullet[1]
$$
Par la construction de $Lf^*$ comme foncteur dérivé à gauche, on voit que
$H^j(Lf^*C) = 0$ pour $j \geq m$ (par l'énoncé dual de Catégories dérivées,
lemme \ref{derived-lemma-negative-vanishing}). Ainsi $H^j(Lf^*\alpha)$ est
un isomorphisme pour $j > m$ et $H^m(Lf^*\alpha)$ est surjectif.
D'autre part, puisque $\mathcal{F}^\bullet$ est un complexe borné supérieurement
de $\mathcal{O}_\mathcal{D}$-modules plats, on voit que
$Lf^*\mathcal{F}^\bullet = f^*\mathcal{F}^\bullet$.
On conclut en appliquant le lemme \ref{sites-cohomology-lemma-strictly-perfect-pullback}.
\end{proof}

\begin{lemma}
\label{sites-cohomology-lemma-cone-pseudo-coherent}
Soit $(\mathcal{C}, \mathcal{O})$ un site annelé et $m \in \mathbf{Z}$.
Soit $(K, L, M, f, g, h)$ un triangle distingué dans $D(\mathcal{O})$.
\begin{enumerate}
\item Si $K$ est $(m + 1)$-pseudo-cohérent et $L$ est $m$-pseudo-cohérent,
alors $M$ est $m$-pseudo-cohérent.
\item Si $K$ et $M$ sont $m$-pseudo-cohérents, alors $L$ est
$m$-pseudo-cohérent.
\item Si $L$ est $(m + 1)$-pseudo-cohérent et $M$ est
$m$-pseudo-cohérent, alors $K$ est $(m + 1)$-pseudo-cohérent.
\end{enumerate}
\end{lemma}

\begin{proof}
Démonstration de (1). Soit $U$ un objet de $\mathcal{C}$. Choisissons un
recouvrement $\{U_i \to U\}$ et des morphismes
$\alpha_i : \mathcal{K}_i^\bullet \to K|_{U_i}$ dans
$D(\mathcal{O}_{U_i})$, où $\mathcal{K}_i^\bullet$ est strictement parfait,
tels que les $H^j(\alpha_i)$ soient des isomorphismes pour $j > m + 1$ et
soient surjectifs pour $j = m + 1$. Nous pouvons remplacer
$\mathcal{K}_i^\bullet$ par
$\sigma_{\geq m + 1}\mathcal{K}_i^\bullet$ et donc supposer que
$\mathcal{K}_i^j = 0$ pour $j < m + 1$. Après avoir raffiné le recouvrement,
nous pouvons choisir des morphismes
$\beta_i : \mathcal{L}_i^\bullet \to L|_{U_i}$ dans
$D(\mathcal{O}_{U_i})$, où $\mathcal{L}_i^\bullet$ est strictement parfait,
tels que $H^j(\beta_i)$ soit un isomorphisme pour $j > m$ et surjectif pour
$j = m$. Par le lemme \ref{sites-cohomology-lemma-lift-through-quasi-isomorphism}, après avoir
raffiné le recouvrement, nous pouvons trouver des morphismes de complexes
$\gamma_i : \mathcal{K}_i^\bullet \to \mathcal{L}_i^\bullet$ tels que les
diagrammes
$$
\xymatrix{
K|_{U_i} \ar[r] & L|_{U_i} \\
\mathcal{K}_i^\bullet \ar[u]^{\alpha_i} \ar[r]^{\gamma_i} &
\mathcal{L}_i^\bullet \ar[u]_{\beta_i}
}
$$
sont commutatifs dans $D(\mathcal{O}_{U_i})$ (cela exige de représenter les
morphismes $\alpha_i$, $\beta_i$ et $K|_{U_i} \to L|_{U_i}$ par de véritables
morphismes de complexes ; certains détails sont omis). Le cône
$C(\gamma_i)^\bullet$ est strictement parfait (lemme \ref{sites-cohomology-lemma-cone}).
La commutativité du diagramme implique qu'il existe un morphisme de triangles
distingués
$$
(\mathcal{K}_i^\bullet, \mathcal{L}_i^\bullet, C(\gamma_i)^\bullet)
\longrightarrow
(K|_{U_i}, L|_{U_i}, M|_{U_i}).
$$
Il résulte du morphisme induit sur les suites exactes longues de cohomologie
et de Homology, lemmes \ref{homology-lemma-four-lemma} et
\ref{homology-lemma-five-lemma}, que
$C(\gamma_i)^\bullet \to M|_{U_i}$ induit un isomorphisme en cohomologie aux
degrés $> m$ et une surjection au degré $m$. Ainsi, $M$ est
$m$-pseudo-cohérent par le lemme
\ref{sites-cohomology-lemma-pseudo-coherent-independent-representative}.

\medskip\noindent
Les assertions (2) et (3) résultent de (1) en faisant tourner le triangle
distingué.
\end{proof}

\begin{lemma}
\label{sites-cohomology-lemma-tensor-pseudo-coherent}
Soit $(\mathcal{C}, \mathcal{O})$ un site annelé. Soient $K, L$ des objets
de $D(\mathcal{O})$.
\begin{enumerate}
\item Si $K$ est $n$-pseudo-cohérent et $H^i(K) = 0$ pour $i > a$, et si
$L$ est $m$-pseudo-cohérent et $H^j(L) = 0$ pour $j > b$, alors
$K \otimes_\mathcal{O}^\mathbf{L} L$ est $t$-pseudo-cohérent, où
$t = \max(m + a, n + b)$.
\item Si $K$ et $L$ sont pseudo-cohérents, alors
$K \otimes_\mathcal{O}^\mathbf{L} L$ est pseudo-cohérent.
\end{enumerate}
\end{lemma}

\begin{proof}
Démonstration de (1). Soit $U$ un objet de $\mathcal{C}$. En remplaçant $U$
par les membres d'un recouvrement et $\mathcal{C}$ par la localisation
$\mathcal{C}/U$, nous pouvons supposer qu'il existe des complexes strictement
parfaits $\mathcal{K}^\bullet$ et $\mathcal{L}^\bullet$, ainsi que des
morphismes $\alpha : \mathcal{K}^\bullet \to K$ et
$\beta : \mathcal{L}^\bullet \to L$, tels que $H^i(\alpha)$ soit un
isomorphisme pour $i > n$ et surjectif pour $i = n$, et que $H^i(\beta)$ soit
un isomorphisme pour $i > m$ et surjectif pour $i = m$. Alors le morphisme
$$
\alpha \otimes^\mathbf{L} \beta :
\text{Tot}(\mathcal{K}^\bullet \otimes_\mathcal{O} \mathcal{L}^\bullet)
\to K \otimes_\mathcal{O}^\mathbf{L} L
$$
induit des isomorphismes sur les faisceaux de cohomologie au degré $i$ pour
$i > t$ et une surjection pour $i = t$. Cela résulte de la suite spectrale
des Tor (détails omis).

\medskip\noindent
Démonstration de (2). Soit $U$ un objet de $\mathcal{C}$. Nous pouvons d'abord
remplacer $U$ par les membres d'un recouvrement et $\mathcal{C}$ par la
localisation $\mathcal{C}/U$ afin de nous ramener au cas où $K$ et $L$ sont
bornés supérieurement. L'assertion résulte alors immédiatement du cas (1).
\end{proof}

\begin{lemma}
\label{sites-cohomology-lemma-summands-pseudo-coherent}
Soit $(\mathcal{C}, \mathcal{O})$ un site annelé. Soit $m \in \mathbf{Z}$.
Si $K \oplus L$ est $m$-pseudo-cohérent (resp.\ pseudo-cohérent) dans
$D(\mathcal{O})$, il en est de même de $K$ et de $L$.
\end{lemma}

\begin{proof}
Supposons que $K \oplus L$ soit $m$-pseudo-cohérent. Soit $U$ un objet de
$\mathcal{C}$. Après avoir remplacé $U$ par les membres d'un recouvrement,
nous pouvons supposer que $K \oplus L \in D^-(\mathcal{O}_U)$, donc que
$L \in D^-(\mathcal{O}_U)$. Notons qu'il existe un triangle distingué
$$
(K \oplus L, K \oplus L, L \oplus L[1]) =
(K, K, 0) \oplus (L, L, L \oplus L[1])
$$
voir Catégories dérivées, lemme
\ref{derived-lemma-direct-sum-triangles}. Par le lemme
\ref{sites-cohomology-lemma-cone-pseudo-coherent}, on voit que $L \oplus L[1]$ est
$m$-pseudo-cohérent. Par conséquent, $L[1] \oplus L[2]$ est également
$m$-pseudo-cohérent. Par récurrence, $L[n] \oplus L[n + 1]$ est
$m$-pseudo-cohérent. Puisque $L$ est borné supérieurement, on voit que
$L[n]$ est $m$-pseudo-cohérent pour $n$ assez grand. En remontant alors et en
utilisant les triangles distingués
$$
(L[n], L[n] \oplus L[n - 1], L[n - 1])
$$
nous concluons que $L[n - 1], L[n - 2], \ldots, L$ sont
$m$-pseudo-cohérents, comme souhaité.
\end{proof}

\begin{lemma}
\label{sites-cohomology-lemma-finite-cohomology}
Soit $(\mathcal{C}, \mathcal{O})$ un site annelé. Soit $K$ un objet de
$D(\mathcal{O})$. Soit $m \in \mathbf{Z}$.
\begin{enumerate}
\item Si $K$ est $m$-pseudo-cohérent et $H^i(K) = 0$ pour $i > m$, alors
$H^m(K)$ est un $\mathcal{O}$-module de type fini.
\item Si $K$ est $m$-pseudo-cohérent et $H^i(K) = 0$ pour $i > m + 1$, alors
$H^{m + 1}(K)$ est un $\mathcal{O}$-module de présentation finie.
\end{enumerate}
\end{lemma}

\begin{proof}
Démonstration de (1). Soit $U$ un objet de $\mathcal{C}$. Nous devons montrer
que $H^m(K)$ peut être engendré par un nombre fini de sections sur les membres
d'un recouvrement de $U$ (voir Modules sur les sites, définition
\ref{sites-modules-definition-site-local}). Ainsi, au cours de la
démonstration, nous pouvons (un nombre fini de fois) choisir un recouvrement
$\{U_i \to U\}$, remplacer $\mathcal{C}$ par $\mathcal{C}/U_i$ et $U$ par
$U_i$. En particulier, d'après nos définitions, nous pouvons supposer qu'il
existe un complexe strictement parfait $\mathcal{E}^\bullet$ et un morphisme
$\alpha : \mathcal{E}^\bullet \to K$ qui induit un isomorphisme en cohomologie
aux degrés $> m$ et une surjection au degré $m$. Il suffit de démontrer le
résultat pour $\mathcal{E}^\bullet$. Soit $n$ le plus grand entier tel que
$\mathcal{E}^n \not = 0$. Si $n = m$, alors
$H^m(\mathcal{E}^\bullet)$ est un quotient de $\mathcal{E}^n$ et le résultat
est clair. Si $n > m$, alors $\mathcal{E}^{n - 1} \to \mathcal{E}^n$ est
surjectif puisque $H^n(\mathcal{E}^\bullet) = 0$. Par le lemme
\ref{sites-cohomology-lemma-local-lift-map}, nous pouvons (après avoir remplacé $U$ par les
membres d'un recouvrement) trouver une section de cette surjection et écrire
$\mathcal{E}^{n - 1} = \mathcal{E}' \oplus \mathcal{E}^n$. Il suffit donc de
démontrer le résultat pour le complexe $(\mathcal{E}')^\bullet$, qui est le
même que $\mathcal{E}^\bullet$, sauf qu'il a $\mathcal{E}'$ au degré $n - 1$
et $0$ au degré $n$. On conclut par récurrence sur $n$.

\medskip\noindent
Démonstration de (2). Choisissons un objet $U$ de $\mathcal{C}$. Comme dans
la démonstration de (1), nous pouvons travailler localement sur $U$. Nous
pouvons donc supposer qu'il existe un complexe strictement parfait
$\mathcal{E}^\bullet$ et un morphisme
$\alpha : \mathcal{E}^\bullet \to K$ qui induit un isomorphisme en cohomologie
aux degrés $> m$ et une surjection au degré $m$. Comme dans la démonstration
de (1), nous pouvons nous ramener au cas où $\mathcal{E}^i = 0$ pour
$i > m + 1$. Alors on voit que
$H^{m + 1}(K) \cong H^{m + 1}(\mathcal{E}^\bullet) =
\Coker(\mathcal{E}^m \to \mathcal{E}^{m + 1})$
est de présentation finie.
\end{proof}








\section{Dimension de Tor}
\label{sites-cohomology-section-tor}

\noindent
Dans cette section, nous étudions plus en détail les résolutions par des
modules plats.

\begin{definition}
\label{sites-cohomology-definition-tor-amplitude}
Soit $(\mathcal{C}, \mathcal{O})$ un site annelé.
Soit $E$ un objet de $D(\mathcal{O})$.
Soient $a, b \in \mathbf{Z}$ avec $a \leq b$.
\begin{enumerate}
\item On dit que $E$ a une {\it amplitude de Tor dans $[a, b]$} si
$H^i(E \otimes_\mathcal{O}^\mathbf{L} \mathcal{F}) = 0$ pour tout
$\mathcal{O}$-module $\mathcal{F}$ et tout $i \not \in [a, b]$.
\item On dit que $E$ a une {\it dimension de Tor finie} s'il a une amplitude
de Tor dans $[a, b]$ pour certains $a, b$.
\item On dit que $E$ {\it est localement de dimension de Tor finie} si, pour
tout objet $U$ de $\mathcal{C}$, il existe un recouvrement $\{U_i \to U\}$
tel que $E|_{U_i}$ soit de dimension de Tor finie pour tout $i$.
\end{enumerate}
Un $\mathcal{O}$-module $\mathcal{F}$ a une {\it dimension de Tor $\leq d$}
si $\mathcal{F}[0]$, considéré comme un objet de $D(\mathcal{O})$, a une
amplitude de Tor dans $[-d, 0]$.
\end{definition}

\noindent
Notons que si $E$, comme dans la définition, est de dimension de Tor finie,
alors $E$ est un objet de $D^b(\mathcal{O})$, comme on le voit en prenant
$\mathcal{F} = \mathcal{O}$ dans la définition ci-dessus.

\begin{lemma}
\label{sites-cohomology-lemma-last-one-flat}
Soit $(\mathcal{C}, \mathcal{O})$ un site annelé. Soit
$\mathcal{E}^\bullet$ un complexe borné supérieurement de
$\mathcal{O}$-modules plats, d'amplitude de Tor dans $[a, b]$. Alors
$\Coker(d_{\mathcal{E}^\bullet}^{a - 1})$ est un $\mathcal{O}$-module plat.
\end{lemma}

\begin{proof}
Comme $\mathcal{E}^\bullet$ est un complexe borné supérieurement de modules
plats, on voit que
$\mathcal{E}^\bullet \otimes_\mathcal{O} \mathcal{F} =
\mathcal{E}^\bullet \otimes_\mathcal{O}^{\mathbf{L}} \mathcal{F}$
pour tout $\mathcal{O}$-module $\mathcal{F}$. Par conséquent, pour tout
$\mathcal{O}$-module $\mathcal{F}$, la suite
$$
\mathcal{E}^{a - 2} \otimes_\mathcal{O} \mathcal{F} \to
\mathcal{E}^{a - 1} \otimes_\mathcal{O} \mathcal{F} \to
\mathcal{E}^a \otimes_\mathcal{O} \mathcal{F}
$$
est exacte au milieu. Puisque
$\mathcal{E}^{a - 2} \to \mathcal{E}^{a - 1} \to \mathcal{E}^a \to
\Coker(d^{a - 1}) \to 0$
est une résolution plate, cela implique que
$\text{Tor}_1^\mathcal{O}(\Coker(d^{a - 1}), \mathcal{F}) = 0$
pour tout $\mathcal{O}$-module $\mathcal{F}$. Cela signifie que
$\Coker(d^{a - 1})$ est plat, voir le lemme \ref{sites-cohomology-lemma-flat-tor-zero}.
\end{proof}

\begin{lemma}
\label{sites-cohomology-lemma-tor-amplitude}
Soit $(\mathcal{C}, \mathcal{O})$ un site annelé. Soit $E$ un objet de
$D(\mathcal{O})$. Soient $a, b \in \mathbf{Z}$ avec $a \leq b$. Les
conditions suivantes sont équivalentes :
\begin{enumerate}
\item $E$ a une amplitude de Tor dans $[a, b]$.
\item $E$ est représenté par un complexe $\mathcal{E}^\bullet$ de
$\mathcal{O}$-modules plats tel que $\mathcal{E}^i = 0$ pour
$i \not \in [a, b]$.
\end{enumerate}
\end{lemma}

\begin{proof}
Si (2) est satisfaite, alors nous pouvons calculer
$E \otimes_\mathcal{O}^\mathbf{L} \mathcal{F} =
\mathcal{E}^\bullet \otimes_\mathcal{O} \mathcal{F}$
et il est clair que (1) est satisfaite.

\medskip\noindent
Supposons (1) satisfaite. Nous pouvons représenter $E$ par un complexe borné
supérieurement de $\mathcal{O}$-modules plats, noté $\mathcal{K}^\bullet$, voir
la section \ref{sites-cohomology-section-flat}. Soit $n$ le plus grand entier tel que
$\mathcal{K}^n \not = 0$. Si $n > b$, alors
$\mathcal{K}^{n - 1} \to \mathcal{K}^n$ est surjectif puisque
$H^n(\mathcal{K}^\bullet) = 0$. Comme $\mathcal{K}^n$ est plat, on voit que
$\Ker(\mathcal{K}^{n - 1} \to \mathcal{K}^n)$ est plat (Modules sur les sites,
lemme \ref{sites-modules-lemma-flat-ses}). Nous pouvons donc remplacer
$\mathcal{K}^\bullet$ par $\tau_{\leq n - 1}\mathcal{K}^\bullet$. Ainsi, par
récurrence sur $n$, nous nous ramenons au cas où $\mathcal{K}^\bullet$ est un complexe
de $\mathcal{O}$-modules plats tel que $\mathcal{K}^i = 0$ pour $i > b$.

\medskip\noindent
Posons $\mathcal{E}^\bullet = \tau_{\geq a}\mathcal{K}^\bullet$.
Tout est clair, sauf le fait que $\mathcal{E}^a$ soit plat, lequel résulte
immédiatement du lemme \ref{sites-cohomology-lemma-last-one-flat} et des définitions.
\end{proof}

\begin{lemma}
\label{sites-cohomology-lemma-bounded-below-tor-amplitude}
Soit $(\mathcal{C}, \mathcal{O})$ un site annelé. Soit $E$ un objet de
$D(\mathcal{O})$. Soit $a \in \mathbf{Z}$. Les conditions suivantes sont
équivalentes :
\begin{enumerate}
\item $E$ a une amplitude de Tor dans $[a, \infty]$.
\item $E$ peut être représenté par un complexe K-plat
$\mathcal{E}^\bullet$ de $\mathcal{O}$-modules plats tel que
$\mathcal{E}^i = 0$ pour $i \not \in [a, \infty]$.
\end{enumerate}
De plus, nous pouvons choisir $\mathcal{E}^\bullet$ de sorte que toute image
inverse par un morphisme de sites annelés soit un complexe K-plat à termes
plats.
\end{lemma}

\begin{proof}
L'implication (2) $\Rightarrow$ (1) est immédiate. Supposons (1) satisfaite.
Choisissons d'abord un complexe K-plat $\mathcal{K}^\bullet$ à termes plats
représentant $E$, voir le lemme \ref{sites-cohomology-lemma-K-flat-resolution}. Pour tout
$\mathcal{O}$-module $\mathcal{M}$, la cohomologie de
$$
\mathcal{K}^{n - 1} \otimes_\mathcal{O} \mathcal{M} \to
\mathcal{K}^n \otimes_\mathcal{O} \mathcal{M} \to
\mathcal{K}^{n + 1} \otimes_\mathcal{O} \mathcal{M}
$$
calcule $H^n(E \otimes_\mathcal{O}^\mathbf{L} \mathcal{M})$. Celui-ci est
toujours nul pour $n < a$. Par conséquent, si nous appliquons le lemme
\ref{sites-cohomology-lemma-last-one-flat} au complexe
$\ldots \to \mathcal{K}^{a - 1} \to \mathcal{K}^a \to \mathcal{K}^{a + 1}$
nous concluons que
$\mathcal{N} = \Coker(\mathcal{K}^{a - 1} \to \mathcal{K}^a)$ est un
$\mathcal{O}$-module plat. Posons
$$
\mathcal{E}^\bullet = \tau_{\geq a}\mathcal{K}^\bullet =
(\ldots \to 0 \to \mathcal{N} \to \mathcal{K}^{a + 1} \to \ldots )
$$
Le noyau $\mathcal{L}^\bullet$ de
$\mathcal{K}^\bullet \to \mathcal{E}^\bullet$ est le complexe
$$
\mathcal{L}^\bullet = (\ldots \to \mathcal{K}^{a - 1} \to
\mathcal{I} \to 0 \to \ldots)
$$
où $\mathcal{I} \subset \mathcal{K}^a$ est l'image de
$\mathcal{K}^{a - 1} \to \mathcal{K}^a$. Puisque nous avons la suite exacte
courte
$0 \to \mathcal{I} \to \mathcal{K}^a \to \mathcal{N} \to 0$
on voit que $\mathcal{I}$ est un $\mathcal{O}$-module plat. Ainsi,
$\mathcal{L}^\bullet$ est un complexe borné supérieurement de modules plats,
donc K-plat par le lemme \ref{sites-cohomology-lemma-bounded-flat-K-flat}. Il s'ensuit que
$\mathcal{E}^\bullet$ est K-plat par le lemme
\ref{sites-cohomology-lemma-K-flat-two-out-of-three-ses}.

\medskip\noindent
Démonstration de l'assertion finale. Soit
$f : (\mathcal{C}', \mathcal{O}') \to (\mathcal{C}, \mathcal{O})$
un morphisme de sites annelés. Par le lemme \ref{sites-cohomology-lemma-pullback-K-flat}, le
complexe $f^*\mathcal{K}^\bullet$ est K-plat à termes plats. Le complexe
$f^*\mathcal{L}^\bullet$ est K-plat, car c'est un complexe borné
supérieurement de $\mathcal{O}'$-modules plats. Nous avons une suite exacte
courte de complexes de $\mathcal{O}'$-modules
$$
0 \to f^*\mathcal{L}^\bullet \to f^*\mathcal{K}^\bullet \to
f^*\mathcal{E}^\bullet \to 0
$$
car la suite exacte courte
$0 \to \mathcal{I} \to \mathcal{K}^a \to \mathcal{N} \to 0$
de modules plats se tire en arrière en une suite exacte courte. Par le lemme
\ref{sites-cohomology-lemma-K-flat-two-out-of-three-ses}, le complexe
$f^*\mathcal{E}^\bullet$ est K-plat, ce qui achève la démonstration.
\end{proof}

\begin{lemma}
\label{sites-cohomology-lemma-tor-amplitude-pullback}
Soit $(f, f^\sharp) : (\mathcal{C}, \mathcal{O}_\mathcal{C}) \to
(\mathcal{D}, \mathcal{O}_\mathcal{D})$
un morphisme de sites annelés. Soit $E$ un objet de
$D(\mathcal{O}_\mathcal{D})$. Si $E$ a une amplitude de Tor dans $[a, b]$,
alors $Lf^*E$ a une amplitude de Tor dans $[a, b]$.
\end{lemma}

\begin{proof}
Supposons que $E$ ait une amplitude de Tor dans $[a, b]$. Par le lemme
\ref{sites-cohomology-lemma-tor-amplitude}, nous pouvons représenter $E$ par un complexe
$\mathcal{E}^\bullet$ de $\mathcal{O}$-modules plats tel que
$\mathcal{E}^i = 0$ pour $i \not \in [a, b]$. Alors $Lf^*E$ est représenté
par $f^*\mathcal{E}^\bullet$. Par Modules sur les sites, lemme
\ref{sites-modules-lemma-pullback-flat}, les modules
$f^*\mathcal{E}^i$ sont plats. Par le lemme \ref{sites-cohomology-lemma-tor-amplitude}, nous
concluons donc que $Lf^*E$ a une amplitude de Tor dans $[a, b]$.
\end{proof}

\begin{lemma}
\label{sites-cohomology-lemma-cone-tor-amplitude}
Soit $(\mathcal{C}, \mathcal{O})$ un site annelé. Soit
$(K, L, M, f, g, h)$ un triangle distingué dans $D(\mathcal{O})$.
Soient $a, b \in \mathbf{Z}$.
\begin{enumerate}
\item Si $K$ a une amplitude de Tor dans $[a + 1, b + 1]$ et $L$ une
amplitude de Tor dans $[a, b]$, alors $M$ a une amplitude de Tor dans
$[a, b]$.
\item Si $K$ et $M$ ont une amplitude de Tor dans $[a, b]$, alors $L$ a une
amplitude de Tor dans $[a, b]$.
\item Si $L$ a une amplitude de Tor dans $[a + 1, b + 1]$ et $M$ une
amplitude de Tor dans $[a, b]$, alors $K$ a une amplitude de Tor dans
$[a + 1, b + 1]$.
\end{enumerate}
\end{lemma}

\begin{proof}
Démonstration omise. Indication : cela résulte simplement de la suite exacte
longue de cohomologie associée à un triangle distingué et du fait que
$- \otimes_\mathcal{O}^{\mathbf{L}} \mathcal{F}$
préserve les triangles distingués. L'assertion la plus facile à démontrer est
(2), et les autres s'en déduisent par translation.
\end{proof}

\begin{lemma}
\label{sites-cohomology-lemma-tensor-tor-amplitude}
Soit $(\mathcal{C}, \mathcal{O})$ un site annelé. Soient $K, L$ des objets
de $D(\mathcal{O})$. Si $K$ a une amplitude de Tor dans $[a, b]$ et $L$ une
amplitude de Tor dans $[c, d]$, alors $K \otimes_\mathcal{O}^\mathbf{L} L$
a une amplitude de Tor dans $[a + c, b + d]$.
\end{lemma}

\begin{proof}
Démonstration omise. Indication : utiliser la suite spectrale des Tor.
\end{proof}

\begin{lemma}
\label{sites-cohomology-lemma-summands-tor-amplitude}
Soit $(\mathcal{C}, \mathcal{O})$ un site annelé. Soient
$a, b \in \mathbf{Z}$. Si $K$, $L$ sont des objets de $D(\mathcal{O})$ et si
$K \oplus L$ a une amplitude de Tor dans $[a, b]$, il en est de même de $K$
et de $L$.
\end{lemma}

\begin{proof}
C'est clair puisque les foncteurs Tor sont additifs.
\end{proof}

\begin{lemma}
\label{sites-cohomology-lemma-bounded}
Soit $(\mathcal{C}, \mathcal{O})$ un site annelé. Soit
$\mathcal{I} \subset \mathcal{O}$ un faisceau d'idéaux. Soit $K$ un objet
de $D(\mathcal{O})$.
\begin{enumerate}
\item Si $K \otimes_\mathcal{O}^\mathbf{L} \mathcal{O}/\mathcal{I}$ est
borné supérieurement, alors
$K \otimes_\mathcal{O}^\mathbf{L} \mathcal{O}/\mathcal{I}^n$
est uniformément borné supérieurement pour tout $n$.
\item Si $K \otimes_\mathcal{O}^\mathbf{L} \mathcal{O}/\mathcal{I}$,
considéré comme un objet de $D(\mathcal{O}/\mathcal{I})$, a une amplitude de
Tor dans $[a, b]$, alors
$K \otimes_\mathcal{O}^\mathbf{L} \mathcal{O}/\mathcal{I}^n$, considéré
comme un objet de $D(\mathcal{O}/\mathcal{I}^n)$, a une amplitude de Tor dans
$[a, b]$ pour tout $n$.
\end{enumerate}
\end{lemma}

\begin{proof}
Démonstration de (1). Supposons que
$K \otimes_\mathcal{O}^\mathbf{L} \mathcal{O}/\mathcal{I}$
soit borné supérieurement, disons que
$H^i(K \otimes_\mathcal{O}^\mathbf{L} \mathcal{O}/\mathcal{I}) = 0$
pour $i > b$. Notons que nous avons les triangles distingués
$$
K \otimes_\mathcal{O}^\mathbf{L}
\mathcal{I}^n/\mathcal{I}^{n + 1} \to
K \otimes_\mathcal{O}^\mathbf{L}
\mathcal{O}/\mathcal{I}^{n + 1} \to
K \otimes_\mathcal{O}^\mathbf{L}
\mathcal{O}/\mathcal{I}^n \to
K \otimes_\mathcal{O}^\mathbf{L}
\mathcal{I}^n/\mathcal{I}^{n + 1}[1]
$$
et que
$$
K \otimes_\mathcal{O}^\mathbf{L}
\mathcal{I}^n/\mathcal{I}^{n + 1} =
\left(
K \otimes_\mathcal{O}^\mathbf{L}
\mathcal{O}/\mathcal{I}\right)
\otimes_{\mathcal{O}/\mathcal{I}}^\mathbf{L}
\mathcal{I}^n/\mathcal{I}^{n + 1}
$$
Par récurrence, nous concluons que
$H^i(K \otimes_\mathcal{O}^\mathbf{L} \mathcal{O}/\mathcal{I}^n) = 0$
pour $i > b$ et pour tout $n$.

\medskip\noindent
Démonstration de (2). Supposons que
$K \otimes_\mathcal{O}^\mathbf{L} \mathcal{O}/\mathcal{I}$, considéré comme
un objet de $D(\mathcal{O}/\mathcal{I})$, ait une amplitude de Tor dans
$[a, b]$. Soit $\mathcal{F}$ un faisceau de
$\mathcal{O}/\mathcal{I}^n$-modules. Nous avons alors une filtration finie
$$
0 \subset \mathcal{I}^{n - 1}\mathcal{F} \subset \ldots
\subset \mathcal{I}\mathcal{F} \subset \mathcal{F}
$$
dont les quotients successifs sont des faisceaux de
$\mathcal{O}/\mathcal{I}$-modules. Ainsi, pour démontrer que
$K \otimes_\mathcal{O}^\mathbf{L} \mathcal{O}/\mathcal{I}^n$ a une amplitude
de Tor dans $[a, b]$, il suffit de montrer que
$H^i(K \otimes_\mathcal{O}^\mathbf{L} \mathcal{O}/\mathcal{I}^n
\otimes_{\mathcal{O}/\mathcal{I}^n}^\mathbf{L} \mathcal{G})$
est nul pour $i \not \in [a, b]$ et pour tout
$\mathcal{O}/\mathcal{I}$-module $\mathcal{G}$. Puisque
$$
\left(K \otimes_\mathcal{O}^\mathbf{L} \mathcal{O}/\mathcal{I}^n\right)
\otimes_{\mathcal{O}/\mathcal{I}^n}^\mathbf{L} \mathcal{G}
=
\left(K \otimes_\mathcal{O}^\mathbf{L} \mathcal{O}/\mathcal{I}\right)
\otimes_{\mathcal{O}/\mathcal{I}}^\mathbf{L} \mathcal{G}
$$
pour tout faisceau de $\mathcal{O}/\mathcal{I}$-modules $\mathcal{G}$, le
résultat s'ensuit.
\end{proof}

\begin{lemma}
\label{sites-cohomology-lemma-tor-amplitude-stalk}
Soit $(\mathcal{C}, \mathcal{O})$ un site annelé. Soit $E$ un objet de
$D(\mathcal{O})$. Soient $a, b \in \mathbf{Z}$.
\begin{enumerate}
\item Si $E$ a une amplitude de Tor dans $[a, b]$, alors, pour tout point
$p$ du site $\mathcal{C}$, l'objet $E_p$ de $D(\mathcal{O}_p)$ a une amplitude
de Tor dans $[a, b]$.
\item Si $\mathcal{C}$ possède suffisamment de points, la réciproque est
vraie.
\end{enumerate}
\end{lemma}

\begin{proof}
Démonstration de (1). Cela résulte du fait que prendre les germes en $p$
revient à tirer en arrière par le morphisme de sites annelés
$(p, \mathcal{O}_p) \to (\mathcal{C}, \mathcal{O})$ ; nous pouvons donc
appliquer le lemme \ref{sites-cohomology-lemma-tor-amplitude-pullback}.

\medskip\noindent
Démonstration de (2). Si $\mathcal{C}$ possède suffisamment de points, alors
nous pouvons vérifier l'annulation de
$H^i(E \otimes_\mathcal{O}^\mathbf{L} \mathcal{F})$
sur les germes, voir Modules sur les sites, lemme
\ref{sites-modules-lemma-check-exactness-stalks}. Puisque
$H^i(E \otimes_\mathcal{O}^\mathbf{L} \mathcal{F})_p =
H^i(E_p \otimes_{\mathcal{O}_p}^\mathbf{L} \mathcal{F}_p)$, nous concluons.
\end{proof}










\section{Complexes parfaits}
\label{sites-cohomology-section-perfect}

\noindent
Dans cette section, nous étudions les propriétés des complexes parfaits sur
les sites annelés.

\begin{definition}
\label{sites-cohomology-definition-perfect}
Soit $(\mathcal{C}, \mathcal{O})$ un site annelé. Soit
$\mathcal{E}^\bullet$ un complexe de $\mathcal{O}$-modules. On dit que
$\mathcal{E}^\bullet$ est {\it parfait} si, pour tout objet $U$ de
$\mathcal{C}$, il existe un recouvrement $\{U_i \to U\}$ tel que, pour tout
$i$, il existe un morphisme de complexes
$\mathcal{E}_i^\bullet \to \mathcal{E}^\bullet|_{U_i}$
qui soit un quasi-isomorphisme, où $\mathcal{E}_i^\bullet$ est strictement
parfait. Un objet $E$ de $D(\mathcal{O})$ est {\it parfait} s'il peut être
représenté par un complexe parfait de $\mathcal{O}$-modules.
\end{definition}

\noindent
Si $\Sh(\mathcal{C})$ est quasi-compact (Sites, section
\ref{sites-section-quasi-compact}), alors tout objet parfait de
$D(\mathcal{O})$ appartient à $D^b(\mathcal{O})$. Ce n'est toutefois pas
nécessairement le cas sinon.

\begin{lemma}
\label{sites-cohomology-lemma-perfect-independent-representative}
Soit $(\mathcal{C}, \mathcal{O})$ un site annelé.
Soit $E$ un objet de $D(\mathcal{O})$.
\begin{enumerate}
\item Si $\mathcal{C}$ possède un objet final $X$ et s'il existe un
recouvrement $\{U_i \to X\}$, des complexes strictement parfaits
$\mathcal{E}_i^\bullet$ de $\mathcal{O}_{U_i}$-modules et des isomorphismes
 $\alpha_i : \mathcal{E}_i^\bullet \to E|_{U_i}$ dans
$D(\mathcal{O}_{U_i})$, alors $E$ est parfait.
\item Si $E$ est parfait, alors tout complexe représentant $E$ est parfait.
\end{enumerate}
\end{lemma}

\begin{proof}
Identique à la démonstration du lemme
\ref{sites-cohomology-lemma-pseudo-coherent-independent-representative}.
\end{proof}

\begin{lemma}
\label{sites-cohomology-lemma-perfect-precise}
Soit $(\mathcal{C}, \mathcal{O})$ un site annelé. Soit $E$ un objet de
$D(\mathcal{O})$. Soient $a \leq b$ des entiers. Si $E$ a une amplitude de
Tor dans $[a, b]$ et est $(a - 1)$-pseudo-cohérent, alors $E$ est parfait.
\end{lemma}

\begin{proof}
Soit $U$ un objet de $\mathcal{C}$. Après avoir remplacé $U$ par les membres
d'un recouvrement et $\mathcal{C}$ par la localisation $\mathcal{C}/U$, nous
pouvons supposer qu'il existe un complexe strictement parfait
$\mathcal{E}^\bullet$ et un morphisme
$\alpha : \mathcal{E}^\bullet \to E$ tel que $H^i(\alpha)$ soit un
isomorphisme pour $i \geq a$. Nous pouvons remplacer, et remplaçons,
$\mathcal{E}^\bullet$ par $\sigma_{\geq a - 1}\mathcal{E}^\bullet$.
Choisissons un triangle distingué
$$
\mathcal{E}^\bullet \to E \to C \to \mathcal{E}^\bullet[1]
$$
De l'annulation des faisceaux de cohomologie de $E$ et de
$\mathcal{E}^\bullet$, ainsi que de l'hypothèse sur $\alpha$, nous obtenons
$C \cong \mathcal{K}[2 - a]$, où
$\mathcal{K} = \Ker(\mathcal{E}^{a - 1} \to \mathcal{E}^a)$. Soit
$\mathcal{F}$ un $\mathcal{O}$-module. En appliquant
$- \otimes_\mathcal{O}^\mathbf{L} \mathcal{F}$, l'hypothèse selon laquelle
$E$ a une amplitude de Tor dans $[a, b]$ implique que le morphisme
$\mathcal{K} \otimes_\mathcal{O} \mathcal{F} \to
\mathcal{E}^{a - 1} \otimes_\mathcal{O} \mathcal{F}$ a pour image
$\Ker(\mathcal{E}^{a - 1} \otimes_\mathcal{O} \mathcal{F}
\to \mathcal{E}^a \otimes_\mathcal{O} \mathcal{F})$.
Il s'ensuit que
$\text{Tor}_1^\mathcal{O}(\mathcal{E}', \mathcal{F}) = 0$, où
$\mathcal{E}' = \Coker(\mathcal{E}^{a - 1} \to \mathcal{E}^a)$. Ainsi,
$\mathcal{E}'$ est plat (lemme \ref{sites-cohomology-lemma-flat-tor-zero}). Par Modules on
Sites, lemme \ref{sites-modules-lemma-flat-locally-finite-presentation}, il
existe donc un recouvrement $\{U_i \to U\}$ tel que
$\mathcal{E}'|_{U_i}$ soit facteur direct d'un module libre de type fini.
Ainsi, le complexe
$$
\mathcal{E}'|_{U_i} \to \mathcal{E}^{a + 1}|_{U_i} \to \ldots \to
\mathcal{E}^b|_{U_i}
$$
est quasi-isomorphe à $E|_{U_i}$ et $E$ est parfait.
\end{proof}

\begin{lemma}
\label{sites-cohomology-lemma-perfect}
Soit $(\mathcal{C}, \mathcal{O})$ un site annelé. Soit $E$ un objet de
$D(\mathcal{O})$. Les conditions suivantes sont équivalentes :
\begin{enumerate}
\item $E$ est parfait, et
\item $E$ est pseudo-cohérent et est localement de dimension de Tor finie.
\end{enumerate}
\end{lemma}

\begin{proof}
Supposons (1). Soit $U$ un objet de $\mathcal{C}$. Par définition, il existe
un recouvrement $\{U_i \to U\}$ tel que $E|_{U_i}$ soit représenté par un
complexe strictement parfait. Ainsi, $E$ est pseudo-cohérent (c'est-à-dire
$m$-pseudo-cohérent pour tout $m$) par le lemme
\ref{sites-cohomology-lemma-pseudo-coherent-independent-representative}. De plus, un facteur
direct d'un module libre de type fini est plat ; $E|_{U_i}$ est donc de
dimension de Tor finie par le lemme \ref{sites-cohomology-lemma-tor-amplitude}. Ainsi, (2) est
satisfaite.

\medskip\noindent
Supposons (2). Soit $U$ un objet de $\mathcal{C}$. Après avoir remplacé $U$
par les membres d'un recouvrement, nous pouvons supposer qu'il existe des
entiers $a \leq b$ tels que $E|_U$ ait une amplitude de Tor dans $[a, b]$.
Puisque $E|_U$ est $m$-pseudo-cohérent pour tout $m$, nous concluons à l'aide
du lemme \ref{sites-cohomology-lemma-perfect-precise}.
\end{proof}

\begin{lemma}
\label{sites-cohomology-lemma-perfect-pullback}
Soit $(f, f^\sharp) : (\mathcal{C}, \mathcal{O}_\mathcal{C}) \to
(\mathcal{D}, \mathcal{O}_\mathcal{D})$
un morphisme de sites annelés. Soit $E$ un objet de
$D(\mathcal{O}_\mathcal{D})$. Si $E$ est parfait dans
$D(\mathcal{O}_\mathcal{D})$, alors $Lf^*E$ est parfait dans
$D(\mathcal{O}_\mathcal{C})$.
\end{lemma}

\begin{proof}
Cela résulte des lemmes \ref{sites-cohomology-lemma-perfect},
\ref{sites-cohomology-lemma-tor-amplitude-pullback} et
\ref{sites-cohomology-lemma-pseudo-coherent-pullback}.
\end{proof}

\begin{lemma}
\label{sites-cohomology-lemma-two-out-of-three-perfect}
Soit $(\mathcal{C}, \mathcal{O})$ un site annelé. Soit $(K, L, M, f, g, h)$
un triangle distingué dans $D(\mathcal{O})$. Si deux des trois objets
$K, L, M$ sont parfaits, alors le troisième l'est également.
\end{lemma}

\begin{proof}
Première démonstration : combiner les lemmes \ref{sites-cohomology-lemma-perfect},
\ref{sites-cohomology-lemma-cone-pseudo-coherent} et \ref{sites-cohomology-lemma-cone-tor-amplitude}.
Deuxième démonstration (esquisse) : supposons $K$ et $L$ parfaits. Soit $U$
un objet de $\mathcal{C}$. Après avoir remplacé $U$ par les membres d'un
recouvrement, nous pouvons supposer que $K|_U$ et $L|_U$ sont représentés
par des complexes strictement parfaits $\mathcal{K}^\bullet$ et
$\mathcal{L}^\bullet$. Après avoir à nouveau remplacé $U$ par les membres
d'un recouvrement, nous pouvons supposer que le morphisme
$K|_U \to L|_U$ est donné par un morphisme de complexes
$\alpha : \mathcal{K}^\bullet \to \mathcal{L}^\bullet$, voir le lemme
\ref{sites-cohomology-lemma-local-actual}. Alors $M|_U$ est isomorphe au cône de $\alpha$, qui
est strictement parfait par le lemme \ref{sites-cohomology-lemma-cone}.
\end{proof}

\begin{lemma}
\label{sites-cohomology-lemma-tensor-perfect}
Soit $(\mathcal{C}, \mathcal{O})$ un site annelé. Si $K, L$ sont des objets
parfaits de $D(\mathcal{O})$, alors
$K \otimes_\mathcal{O}^\mathbf{L} L$ l'est également.
\end{lemma}

\begin{proof}
Cela résulte des lemmes \ref{sites-cohomology-lemma-perfect},
\ref{sites-cohomology-lemma-tensor-pseudo-coherent} et
\ref{sites-cohomology-lemma-tensor-tor-amplitude}.
\end{proof}

\begin{lemma}
\label{sites-cohomology-lemma-summands-perfect}
Soit $(\mathcal{C}, \mathcal{O})$ un site annelé. Si $K \oplus L$ est un
objet parfait de $D(\mathcal{O})$, alors $K$ et $L$ le sont également.
\end{lemma}

\begin{proof}
Cela résulte des lemmes \ref{sites-cohomology-lemma-perfect},
\ref{sites-cohomology-lemma-summands-pseudo-coherent} et
\ref{sites-cohomology-lemma-summands-tor-amplitude}.
\end{proof}














\section{Duaux}
\label{sites-cohomology-section-duals}

\noindent
Dans cette section, nous caractérisons les objets dualisables de la catégorie
des complexes et de la catégorie dérivée. En particulier, nous verrons qu'un
objet de $D(\mathcal{O})$ possède un dual si et seulement s'il est parfait
(cela résulte de l'exemple \ref{sites-cohomology-example-dual-derived} et du lemme
\ref{sites-cohomology-lemma-left-dual-derived}).

\begin{lemma}
\label{sites-cohomology-lemma-symmetric-monoidal-cat-complexes}
Soit $(\mathcal{C}, \mathcal{O})$ un site annelé. La catégorie des complexes
de $\mathcal{O}$-modules, munie du produit tensoriel défini par
$\mathcal{F}^\bullet \otimes \mathcal{G}^\bullet =
\text{Tot}(\mathcal{F}^\bullet \otimes_\mathcal{O} \mathcal{G}^\bullet)$
est une catégorie monoïdale symétrique.
\end{lemma}

\begin{proof}
Démonstration omise. Indications : comme unité $\mathbf{1}$, on prend le
complexe ayant $\mathcal{O}$ au degré $0$ et zéro aux autres degrés, avec les
isomorphismes évidents
$\text{Tot}(\mathbf{1} \otimes_\mathcal{O} \mathcal{G}^\bullet) =
\mathcal{G}^\bullet$ et
$\text{Tot}(\mathcal{F}^\bullet \otimes_\mathcal{O} \mathbf{1}) =
\mathcal{F}^\bullet$. Pour démontrer le lemme, il faut vérifier la
commutativité de divers diagrammes, voir Catégories, définitions
\ref{categories-definition-monoidal-category} et
\ref{categories-definition-symmetric-monoidal-category}.
Dans chaque cas, les vérifications sont immédiates.
\end{proof}

\begin{example}
\label{sites-cohomology-example-dual}
Soit $(\mathcal{C}, \mathcal{O})$ un site annelé. Soit
$\mathcal{F}^\bullet$ un complexe de $\mathcal{O}$-modules tel que, pour tout
$U \in \Ob(\mathcal{C})$, il existe un recouvrement $\{U_i \to U\}$ tel que
$\mathcal{F}^\bullet|_{U_i}$ soit strictement parfait. Considérons le complexe
$$
\mathcal{G}^\bullet = \SheafHom^\bullet(\mathcal{F}^\bullet, \mathcal{O})
$$
comme dans la section \ref{sites-cohomology-section-hom-complexes}. Soient
$$
\eta :
\mathcal{O}
\to
\text{Tot}(\mathcal{F}^\bullet \otimes_\mathcal{O} \mathcal{G}^\bullet)
\quad\text{et}\quad
\epsilon :
\text{Tot}(\mathcal{G}^\bullet \otimes_\mathcal{O} \mathcal{F}^\bullet)
\to
\mathcal{O}
$$
définis par $\eta = \sum \eta_n$ et $\epsilon = \sum \epsilon_n$, où
$\eta_n : \mathcal{O} \to
\mathcal{F}^n \otimes_\mathcal{O} \mathcal{G}^{-n}$
et
$\epsilon_n : \mathcal{G}^{-n} \otimes_\mathcal{O} \mathcal{F}^n
\to \mathcal{O}$ sont comme dans Modules sur les sites, exemple
\ref{sites-modules-example-dual}. Alors
$\mathcal{G}^\bullet, \eta, \epsilon$ est un dual à gauche de
$\mathcal{F}^\bullet$ au sens de Catégories, définition
\ref{categories-definition-dual}. Nous omettons la vérification de
$(1 \otimes \epsilon) \circ (\eta \otimes 1) = \text{id}_{\mathcal{F}^\bullet}$
et de
$(\epsilon \otimes 1) \circ (1 \otimes \eta) =
\text{id}_{\mathcal{G}^\bullet}$. Comparer avec Compléments d'algèbre, lemme
\ref{more-algebra-lemma-left-dual-complex}.
\end{example}

\begin{lemma}
\label{sites-cohomology-lemma-left-dual-complex}
Soit $(\mathcal{C}, \mathcal{O})$ un site annelé. Soit
$\mathcal{F}^\bullet$ un complexe de $\mathcal{O}$-modules. Si
$\mathcal{F}^\bullet$ possède un dual à gauche dans la catégorie monoïdale
des complexes de $\mathcal{O}$-modules (Catégories, définition
\ref{categories-definition-dual}), alors, pour tout objet $U$ de
$\mathcal{C}$, il existe un recouvrement $\{U_i \to U\}$ tel que
$\mathcal{F}^\bullet|_{U_i}$ soit strictement parfait et que le dual à gauche
soit celui construit dans l'exemple \ref{sites-cohomology-example-dual}.
\end{lemma}

\begin{proof}
Par unicité des duaux à gauche (Catégories, remarque
\ref{categories-remark-left-dual-adjoint}), nous obtenons l'assertion finale
dès que nous avons montré que $\mathcal{F}^\bullet$ est comme annoncé. Soit
$\mathcal{G}^\bullet, \eta, \epsilon$ un dual à gauche. Écrivons
$\eta = \sum \eta_n$ et $\epsilon = \sum \epsilon_n$, où
$\eta_n : \mathcal{O} \to
\mathcal{F}^n \otimes_\mathcal{O} \mathcal{G}^{-n}$
et
$\epsilon_n : \mathcal{G}^{-n} \otimes_\mathcal{O} \mathcal{F}^n
\to \mathcal{O}$. Puisque
$(1 \otimes \epsilon) \circ (\eta \otimes 1) = \text{id}_{\mathcal{F}^\bullet}$
et
$(\epsilon \otimes 1) \circ (1 \otimes \eta) = \text{id}_{\mathcal{G}^\bullet}$
par Catégories, définition \ref{categories-definition-dual}, on voit
immédiatement que
$(1 \otimes \epsilon_n) \circ (\eta_n \otimes 1) = \text{id}_{\mathcal{F}^n}$
et
$(\epsilon_n \otimes 1) \circ (1 \otimes \eta_n) =
\text{id}_{\mathcal{G}^{-n}}$.
Autrement dit, $\mathcal{G}^{-n}$ est un dual à gauche de
$\mathcal{F}^n$, et Modules sur les sites, lemme
\ref{sites-modules-lemma-left-dual-module}, s'applique à chaque
$\mathcal{F}^n$. Soit $U$ un objet de $\mathcal{C}$. Il existe un recouvrement
$\{U_i \to U\}$ tel que, pour tout $i$, seul un nombre fini des
$\eta_n|_{U_i}$ soient non nuls. Ainsi, après avoir remplacé $U$ par $U_i$,
nous pouvons supposer que seul un nombre fini des $\eta_n|_U$ sont non nuls ;
par le lemme cité, cela implique que seul un nombre fini des
$\mathcal{F}^n|_U$ sont non nuls. En utilisant à nouveau le lemme, nous
pouvons alors trouver un recouvrement $\{U_i \to U\}$ tel que chaque
$\mathcal{F}^n|_{U_i}$ soit facteur direct d'un $\mathcal{O}$-module libre de
type fini, ce qui achève la démonstration.
\end{proof}

\begin{lemma}
\label{sites-cohomology-lemma-dual-perfect-complex}
Soit $(\mathcal{C}, \mathcal{O})$ un site annelé. Soit $K$ un objet parfait
de $D(\mathcal{O})$. Alors $K^\vee = R\SheafHom(K, \mathcal{O})$ est également
un objet parfait et $(K^\vee)^\vee \cong K$. Il existe des isomorphismes
fonctoriels
$$
M \otimes^\mathbf{L}_\mathcal{O} K^\vee = R\SheafHom_\mathcal{O}(K, M)
$$
et
$$
H^0(\mathcal{C}, M \otimes^\mathbf{L}_\mathcal{O} K^\vee) =
\Hom_{D(\mathcal{O})}(K, M)
$$
pour $M$ dans $D(\mathcal{O})$.
\end{lemma}

\begin{proof}
Nous utiliserons sans autre mention le fait que la formation du Hom interne
commute à la restriction (lemme \ref{sites-cohomology-lemma-restriction-RHom-to-U}). Soit $U$
un objet quelconque de $\mathcal{C}$. Pour vérifier que $K^\vee$ est parfait,
il suffit de montrer qu'il existe un recouvrement $\{U_i \to U\}$ tel que
$K^\vee|_{U_i}$ soit parfait pour tout $i$. Il existe un morphisme canonique
$$
K = R\SheafHom(\mathcal{O}, \mathcal{O})
\otimes_{\mathcal{O}}^\mathbf{L} K \longrightarrow
R\SheafHom(R\SheafHom(K, \mathcal{O}), \mathcal{O}) =
(K^\vee)^\vee
$$
voir le lemme \ref{sites-cohomology-lemma-internal-hom-evaluate}. Il suffit de démontrer qu'il
existe un recouvrement $\{U_i \to U\}$ tel que la restriction de ce morphisme
à $\mathcal{C}/U_i$ soit un isomorphisme pour tout $i$. Par le lemme
\ref{sites-cohomology-lemma-dual}, pour obtenir l'assertion finale, il suffit de vérifier que
le morphisme (\ref{sites-cohomology-equation-eval})
$$
M \otimes^\mathbf{L}_\mathcal{O} K^\vee
\longrightarrow
R\SheafHom(K, M)
$$
est un isomorphisme. C'est également une question locale (au sens ci-dessus).
Il suffit donc de démontrer le lemme lorsque $K$ est représenté par un
complexe strictement parfait.

\medskip\noindent
Supposons $K$ représenté par le complexe strictement parfait
$\mathcal{E}^\bullet$. Il résulte alors du lemme
\ref{sites-cohomology-lemma-Rhom-strictly-perfect} que $K^\vee$ est représenté par le complexe
dont les termes sont
$(\mathcal{E}^n)^\vee =
\SheafHom_\mathcal{O}(\mathcal{E}^n, \mathcal{O})$
au degré $-n$. Puisque $\mathcal{E}^n$ est facteur direct d'un
$\mathcal{O}$-module libre de type fini, il en est de même de
$(\mathcal{E}^n)^\vee$. Ainsi, $K^\vee$ est lui aussi représenté par un
complexe strictement parfait, et l'on voit que $K^\vee$ est parfait. Le
morphisme $K \to (K^\vee)^\vee$ est un isomorphisme, car il est donné, au
signe près, par les morphismes d'évaluation
$\mathcal{E}^n \to ((\mathcal{E}^n)^\vee)^\vee$, qui sont des isomorphismes.
Pour voir que (\ref{sites-cohomology-equation-eval}) est un isomorphisme, représentons $M$ par
un complexe K-plat $\mathcal{F}^\bullet$. Par le lemme
\ref{sites-cohomology-lemma-Rhom-strictly-perfect}, le complexe $R\SheafHom(K, M)$ est
représenté par le complexe de termes
$$
\bigoplus\nolimits_{n = p + q}
\SheafHom_\mathcal{O}(\mathcal{E}^{-q}, \mathcal{F}^p)
$$
D'autre part, l'objet $M \otimes^\mathbf{L}_\mathcal{O} K^\vee$ est
représenté par le complexe de termes
$$
\bigoplus\nolimits_{n = p + q}
\mathcal{F}^p \otimes_\mathcal{O} (\mathcal{E}^{-q})^\vee
$$
Ainsi, l'assertion selon laquelle (\ref{sites-cohomology-equation-eval}) est un isomorphisme se
ramène à l'assertion selon laquelle le morphisme canonique
$$
\mathcal{F}
\otimes_\mathcal{O}
\SheafHom_\mathcal{O}(\mathcal{E}, \mathcal{O})
\longrightarrow
\SheafHom_\mathcal{O}(\mathcal{E}, \mathcal{F})
$$
est un isomorphisme lorsque $\mathcal{E}$ est facteur direct d'un
$\mathcal{O}$-module libre de type fini et $\mathcal{F}$ est un
$\mathcal{O}$-module quelconque. Cela résulte immédiatement de l'assertion
correspondante lorsque $\mathcal{E}$ est libre de type fini.
\end{proof}

\begin{lemma}
\label{sites-cohomology-lemma-symmetric-monoidal-derived}
Soit $(\mathcal{C}, \mathcal{O})$ un site annelé. La catégorie dérivée
$D(\mathcal{O})$ est une catégorie monoïdale symétrique, dont le produit
tensoriel est le produit tensoriel dérivé, avec les contraintes usuelles
d'associativité et de commutativité (pour les règles de signes, voir More on
Algèbre, section \ref{more-algebra-section-sign-rules}).
\end{lemma}

\begin{proof}
Démonstration omise. Comparer avec le lemme
\ref{sites-cohomology-lemma-symmetric-monoidal-cat-complexes}.
\end{proof}

\begin{example}
\label{sites-cohomology-example-dual-derived}
Soit $(\mathcal{C}, \mathcal{O})$ un site annelé. Soit $K$ un objet parfait
de $D(\mathcal{O})$. Posons $K^\vee = R\SheafHom(K, \mathcal{O})$ comme dans
le lemme \ref{sites-cohomology-lemma-dual-perfect-complex}. Alors le morphisme
$$
K \otimes_\mathcal{O}^\mathbf{L} K^\vee \longrightarrow R\SheafHom(K, K)
$$
est un isomorphisme (par le lemme). Notons
$$
\eta :
\mathcal{O}
\longrightarrow
K \otimes_\mathcal{O}^\mathbf{L} K^\vee
$$
le morphisme qui envoie $1$ sur la section correspondant à $\text{id}_K$
par l'isomorphisme ci-dessus. Notons
$$
\epsilon : 
K^\vee
\otimes_\mathcal{O}^\mathbf{L} K
\longrightarrow
\mathcal{O}
$$
le morphisme d'évaluation (pour le construire, on peut par exemple utiliser
le lemme \ref{sites-cohomology-lemma-internal-hom-composition}). Alors
$K^\vee, \eta, \epsilon$ est un dual à gauche de $K$ au sens de Catégories,
définition \ref{categories-definition-dual}. Nous omettons la vérification de
$(1 \otimes \epsilon) \circ (\eta \otimes 1) = \text{id}_K$
et de
$(\epsilon \otimes 1) \circ (1 \otimes \eta) =
\text{id}_{K^\vee}$.
\end{example}

\begin{lemma}
\label{sites-cohomology-lemma-left-dual-derived}
Soit $(\mathcal{C}, \mathcal{O})$ un site annelé. Soit $M$ un objet de
$D(\mathcal{O})$. Si $M$ possède un dual à gauche dans la catégorie monoïdale
$D(\mathcal{O})$ (Catégories, définition
\ref{categories-definition-dual}), alors $M$ est parfait et le dual à gauche
est celui construit dans l'exemple \ref{sites-cohomology-example-dual-derived}.
\end{lemma}

\begin{proof}
Soit $N, \eta, \epsilon$ un dual à gauche. Observons que, pour tout objet $U$
de $\mathcal{C}$, la restriction $N|_U, \eta|_U, \epsilon|_U$ est un dual à
gauche de $M|_U$.

\medskip\noindent
Soit $U$ un objet de $\mathcal{C}$. Il suffit de trouver un recouvrement
$\{U_i \to U\}_{i \in I}$ de $\mathcal{C}$ tel que $M|_{U_i}$ soit un objet
parfait de $D(\mathcal{O}_{U_i})$. Nous pouvons donc remplacer
$\mathcal{C}, \mathcal{O}, M, N, \eta, \epsilon$ par
$\mathcal{C}/U, \mathcal{O}_U, M|_U, N|_U, \eta|_U, \epsilon|_U$
et supposer que $\mathcal{C}$ possède un objet final $X$. De plus, au cours
de la démonstration, nous pouvons (un nombre fini de fois) remplacer $X$ par
les membres d'un recouvrement $\{U_i \to X\}$ de $X$.

\medskip\noindent
Nous allons utiliser plusieurs fois l'argument suivant. Choisissons un
complexe quelconque $\mathcal{M}^\bullet$ de $\mathcal{O}$-modules
représentant $M$. Choisissons un complexe K-plat $\mathcal{N}^\bullet$
représentant $N$ et dont les termes sont des $\mathcal{O}$-modules plats,
voir le lemme \ref{sites-cohomology-lemma-K-flat-resolution}. Considérons le morphisme
$$
\eta : \mathcal{O} \to
\text{Tot}(\mathcal{M}^\bullet \otimes_\mathcal{O} \mathcal{N}^\bullet)
$$
Après avoir remplacé $X$ par les membres d'un recouvrement, nous pouvons
trouver un entier $N$ et, pour $i = 1, \ldots, N$, des entiers
$n_i \in \mathbf{Z}$ ainsi que des sections $f_i$ et $g_i$ de
$\mathcal{M}^{n_i}$ et de $\mathcal{N}^{-n_i}$ telles que
$$
\eta(1) = \sum\nolimits_i f_i \otimes g_i
$$
Soit $\mathcal{K}^\bullet \subset \mathcal{M}^\bullet$ un sous-complexe
quelconque de $\mathcal{O}$-modules contenant les sections $f_i$ pour
$i = 1, \ldots, N$. Puisque
$\text{Tot}(\mathcal{K}^\bullet \otimes_\mathcal{O} \mathcal{N}^\bullet)
\subset
\text{Tot}(\mathcal{M}^\bullet \otimes_\mathcal{O} \mathcal{N}^\bullet)$
par platitude des modules $\mathcal{N}^n$, on voit que $\eta$ se factorise par
$$
\tilde \eta :
\mathcal{O} \to
\text{Tot}(\mathcal{K}^\bullet \otimes_\mathcal{O} \mathcal{N}^\bullet)
$$
En notant $K$ l'objet de $D(\mathcal{O})$ représenté par
$\mathcal{K}^\bullet$, nous obtenons un diagramme commutatif
$$
\xymatrix{
M \ar[rr]_-{\eta \otimes 1} \ar[rrd]_{\tilde \eta \otimes 1} & &
M \otimes^\mathbf{L} N \otimes^\mathbf{L} M
\ar[r]_-{1 \otimes \epsilon} &
M \\
& &
K \otimes^\mathbf{L} N \otimes^\mathbf{L} M
\ar[u] \ar[r]^-{1 \otimes \epsilon} &
K \ar[u]
}
$$
Puisque la composée de la ligne supérieure est l'identité de $M$, nous
concluons que $M$ est facteur direct de $K$ dans $D(\mathcal{O})$.

\medskip\noindent
Comme première application de l'argument ci-dessus, nous pouvons choisir le
sous-complexe
$\mathcal{K}^\bullet = \sigma_{\geq a} \tau_{\leq b}\mathcal{M}^\bullet$
avec $a < n_i < b$ pour $i = 1, \ldots, N$. Ainsi, $M$ est facteur direct
dans $D(\mathcal{O})$ d'un complexe borné, et nous concluons que nous pouvons
supposer $M$ dans $D^b(\mathcal{O})$. (Rappelons que le procédé ci-dessus
comporte le remplacement de $X$ par les membres d'un recouvrement.)

\medskip\noindent
Puisque $M$ appartient à $D^b(\mathcal{O})$, nous pouvons choisir
$\mathcal{M}^\bullet$ comme un complexe borné supérieurement de modules plats
(par Modules, lemme \ref{modules-lemma-module-quotient-flat}, et Derived
Catégories, lemme \ref{derived-lemma-subcategory-left-resolution}). Dans
l'argument ci-dessus, nous pouvons alors choisir
$\mathcal{K}^\bullet = \sigma_{\geq a}\mathcal{M}^\bullet$ avec $a < n_i$
pour $i = 1, \ldots, N$. Ainsi, nous pouvons supposer que $M$ est facteur
direct dans $D(\mathcal{O})$ d'un complexe borné de modules plats. En
particulier, $M$ a une amplitude de Tor finie.

\medskip\noindent
Disons que $M$ a une amplitude de Tor dans $[a, b]$. En supposant $M$
$m$-pseudo-cohérent, nous allons montrer qu'après avoir remplacé $X$ par les
membres d'un recouvrement, nous pouvons supposer $M$
$(m - 1)$-pseudo-cohérent. Cela achèvera la démonstration par le lemme
\ref{sites-cohomology-lemma-perfect-precise} et le fait que $M$ est de toute façon
$(b + 1)$-pseudo-cohérent. Après avoir remplacé $X$ par les membres d'un
recouvrement, nous pouvons supposer qu'il existe un complexe strictement
parfait $\mathcal{E}^\bullet$ et un morphisme
$\alpha : \mathcal{E}^\bullet \to M$ dans $D(\mathcal{O})$ tel que
$H^i(\alpha)$ soit un isomorphisme pour $i > m$ et surjectif pour $i = m$.
Nous pouvons supposer, et supposons, que $\mathcal{E}^i = 0$ pour $i < m$.
Choisissons un triangle distingué
$$
\mathcal{E}^\bullet \to M \to L \to \mathcal{E}^\bullet[1]
$$
Observons que $H^i(L) = 0$ pour $i \geq m$. Nous pouvons donc représenter $L$
par un complexe $\mathcal{L}^\bullet$ tel que $\mathcal{L}^i = 0$ pour
$i \geq m$. Le morphisme $L \to \mathcal{E}^\bullet[1]$ est donné par un
morphisme de complexes $\mathcal{L}^\bullet \to \mathcal{E}^\bullet[1]$ qui
est nul à tous les degrés sauf au degré $m - 1$, où nous obtenons un morphisme
$\mathcal{L}^{m - 1} \to \mathcal{E}^m$, voir Catégories dérivées, lemme
\ref{derived-lemma-negative-exts}. Alors $M$ est représenté par le complexe
$$
\mathcal{M}^\bullet :
\ldots \to
\mathcal{L}^{m - 2} \to
\mathcal{L}^{m - 1} \to
\mathcal{E}^m \to
\mathcal{E}^{m + 1} \to \ldots
$$
Appliquons à ce complexe la discussion du deuxième paragraphe pour obtenir
des sections $f_i$ de $\mathcal{M}^{n_i}$ pour $i = 1, \ldots, N$. Pour
$n < m$, soit $\mathcal{K}^n \subset \mathcal{L}^n$ le
$\mathcal{O}$-sous-module engendré par les sections $f_i$ pour $n_i = n$ et
par les $d(f_i)$ pour $n_i = n - 1$. Pour $n \geq m$, posons
$\mathcal{K}^n = \mathcal{E}^n$. Nous avons clairement un morphisme de
triangles distingués
$$
\xymatrix{
\mathcal{E}^\bullet \ar[r] &
\mathcal{M}^\bullet \ar[r] &
\mathcal{L}^\bullet \ar[r] &
\mathcal{E}^\bullet[1] \\
\mathcal{E}^\bullet \ar[r] \ar[u] &
\mathcal{K}^\bullet \ar[r] \ar[u] &
\sigma_{\leq m - 1}\mathcal{K}^\bullet \ar[r] \ar[u] &
\mathcal{E}^\bullet[1] \ar[u]
}
$$
où tous les morphismes sont ceux indiqués ci-dessus. Notons $K$ l'objet de
$D(\mathcal{O})$ correspondant au complexe $\mathcal{K}^\bullet$. Par les
arguments du deuxième paragraphe de la démonstration, nous obtenons un
morphisme $s : M \to K$ dans $D(\mathcal{O})$ tel que la composée
$M \to K \to M$ soit l'identité de $M$. Nous ne savons pas si le diagramme
$$
\xymatrix{
\mathcal{E}^\bullet \ar[r] &
\mathcal{K}^\bullet \ar@{=}[r] &
K \\
\mathcal{E}^\bullet \ar[u]^{\text{id}} \ar[r]^i &
\mathcal{M}^\bullet \ar@{=}[r] &
M \ar[u]_s
}
$$
est commutatif, mais nous savons qu'il le devient après composition avec le
morphisme $K \to M$. Par le lemme \ref{sites-cohomology-lemma-local-actual}, après avoir
remplacé $X$ par les membres d'un recouvrement, nous pouvons supposer que
$s \circ i$ est donné par un morphisme de complexes
$\sigma : \mathcal{E}^\bullet \to \mathcal{K}^\bullet$. Par le même lemme,
nous pouvons supposer que la composée de $\sigma$ avec l'inclusion
$\mathcal{K}^\bullet \subset \mathcal{M}^\bullet$ est homotope à zéro par une
certaine homotopie
$\{h^i : \mathcal{E}^i \to \mathcal{M}^{i - 1}\}$. Ainsi, après avoir
remplacé $\mathcal{K}^{m - 1}$ par
$\mathcal{K}^{m - 1} + \Im(h^m)$ (notons qu'après cette opération,
$\mathcal{K}^{m - 1}$ est encore engendré par un nombre fini de sections
globales), on voit que $\sigma$ lui-même est homotope à zéro ! Cela signifie
que nous avons un diagramme commutatif en traits pleins
$$
\xymatrix{
\mathcal{E}^\bullet \ar[r] &
M \ar[r] &
\mathcal{L}^\bullet \ar[r] &
\mathcal{E}^\bullet[1] \\
\mathcal{E}^\bullet \ar[r] \ar[u] &
K \ar[r] \ar[u] &
\sigma_{\leq m - 1}\mathcal{K}^\bullet \ar[r] \ar[u] &
\mathcal{E}^\bullet[1] \ar[u] \\
\mathcal{E}^\bullet \ar[r] \ar[u] &
M \ar[r] \ar[u]^s &
\mathcal{L}^\bullet \ar[r] \ar@{..>}[u] &
\mathcal{E}^\bullet[1] \ar[u]
}
$$
Par les axiomes des catégories triangulées, nous obtenons une flèche en
pointillés complétant le diagramme. En considérant les faisceaux de
cohomologie au degré $m - 1$, nous obtenons
$$
\xymatrix{
H^{m - 1}(M) \ar[r] &
H^{m - 1}(\mathcal{L}^\bullet) \ar[r] &
H^m(\mathcal{E}^\bullet) \\
H^{m - 1}(K) \ar[r] \ar[u] &
H^{m - 1}(\sigma_{\leq m - 1}\mathcal{K}^\bullet) \ar[r] \ar[u] &
H^m(\mathcal{E}^\bullet) \ar[u] \\
H^{m - 1}(M) \ar[r] \ar[u] &
H^{m - 1}(\mathcal{L}^\bullet) \ar[r] \ar[u] &
H^m(\mathcal{E}^\bullet) \ar[u]
}
$$
Puisque les composées verticales sont l'identité dans les colonnes de gauche
et de droite, nous concluons que la composée verticale
$H^{m - 1}(\mathcal{L}^\bullet) \to
H^{m - 1}(\sigma_{\leq m - 1}\mathcal{K}^\bullet) \to
H^{m - 1}(\mathcal{L}^\bullet)$ au milieu est surjective ! En particulier,
$H^{m - 1}(\sigma_{\leq m - 1}\mathcal{K}^\bullet) \to
H^{m - 1}(\mathcal{L}^\bullet)$ est surjectif. En utilisant le morphisme
induit entre les suites exactes longues de faisceaux de cohomologie par le
morphisme de triangles ci-dessus, une chasse au diagramme montre que
$H^i(K) \to H^i(M)$ est un isomorphisme pour $i \geq m$ et est surjectif pour
$i = m - 1$. Par construction, nous pouvons choisir un $r \geq 0$ et une
surjection $\mathcal{O}^{\oplus r} \to \mathcal{K}^{m - 1}$. Alors la composée
$$
(\mathcal{O}^{\oplus r} \to \mathcal{E}^m \to
\mathcal{E}^{m + 1} \to \ldots ) \longrightarrow
K \longrightarrow M
$$
induit un isomorphisme sur les faisceaux de cohomologie aux degrés $\geq m$
et une surjection au degré $m - 1$, ce qui achève la démonstration.
\end{proof}

\begin{lemma}
\label{sites-cohomology-lemma-colim-and-lim-of-duals}
\begin{slogan}
Dualité triviale pour les systèmes d'objets parfaits.
\end{slogan}
Soit $(\mathcal{C}, \mathcal{O})$ un site annelé. Soit
$(K_n)_{n \in \mathbf{N}}$ un système d'objets parfaits de $D(\mathcal{O})$.
Soit $K = \text{hocolim} K_n$ la colimite dérivée (Catégories dérivées,
définition \ref{derived-definition-derived-colimit}). Alors, pour tout objet
$E$ de $D(\mathcal{O})$, nous avons
$$
R\SheafHom(K, E) = R\lim E \otimes^\mathbf{L}_\mathcal{O} K_n^\vee
$$
où $(K_n^\vee)$ est le système inverse de complexes parfaits duaux.
\end{lemma}

\begin{proof}
Par le lemme \ref{sites-cohomology-lemma-dual-perfect-complex}, nous avons
$R\lim E \otimes^\mathbf{L}_\mathcal{O} K_n^\vee =
R\lim R\SheafHom(K_n, E)$
qui s'insère dans le triangle distingué
$$
R\lim R\SheafHom(K_n, E) \to
\prod R\SheafHom(K_n, E) \to
\prod R\SheafHom(K_n, E)
$$
Puisque $K$ s'insère de même dans le triangle distingué
$\bigoplus K_n \to \bigoplus K_n \to K$, il suffit de montrer que
$\prod R\SheafHom(K_n, E) = R\SheafHom(\bigoplus K_n, E)$. C'est une
conséquence formelle de (\ref{sites-cohomology-equation-internal-hom}) et du fait que le
produit tensoriel dérivé commute aux sommes directes.
\end{proof}






\section{Objets inversibles dans la catégorie dérivée}
\label{sites-cohomology-section-invertible-D-or-R}

\noindent
Nous caractérisons les objets inversibles dans la catégorie dérivée d'un site
annelé (à la fois dans le cas d'un topos localement annelé et dans le cas
général).

\begin{lemma}
\label{sites-cohomology-lemma-category-summands-finite-free}
Soit $(\mathcal{C}, \mathcal{O})$ un site annelé. Posons
$R = \Gamma(\mathcal{C}, \mathcal{O})$. La catégorie des
$\mathcal{O}$-modules qui sont facteurs directs de $\mathcal{O}$-modules
libres de type fini est équivalente à la catégorie des $R$-modules projectifs
de type fini.
\end{lemma}

\begin{proof}
Observons qu'un $R$-module projectif de type fini est la même chose qu'un
facteur direct d'un $R$-module libre de type fini. L'équivalence est donnée
par le foncteur
$\mathcal{E} \mapsto \Gamma(\mathcal{C}, \mathcal{E})$. Le foncteur inverse
est donné par la construction suivante. Considérons le morphisme de topos
$f : \Sh(\mathcal{C}) \to \Sh(\text{pt})$, où $f_*$ est donné par la prise
des sections globales et $f^{-1}$ envoie un ensemble $S$, c'est-à-dire un
objet de $\Sh(\text{pt})$, sur le faisceau constant de valeur $S$. Nous
obtenons un morphisme
$(f, f^\sharp) : (\Sh(\mathcal{C}), \mathcal{O}) \to (\Sh(\text{pt}), R)$
de topos annelés en utilisant le morphisme identité $R \to f_*\mathcal{O}$.
Le foncteur inverse est alors donné par $f^*$.
\end{proof}

\begin{lemma}
\label{sites-cohomology-lemma-invertible-derived}
Soit $(\mathcal{C}, \mathcal{O})$ un site annelé. Soit $M$ un objet de
$D(\mathcal{O})$. Les conditions suivantes sont équivalentes :
\begin{enumerate}
\item $M$ est inversible dans $D(\mathcal{O})$, voir Catégories, définition
\ref{categories-definition-invertible}, et
\item il existe une décomposition en produit direct localement
finie\footnote{Cela signifie que, pour tout objet $U$ de $\mathcal{C}$, il
existe un recouvrement $\{U_i \to U\}$ tel que, pour tout $i$, le faisceau
$\mathcal{O}_n|_{U_i}$ ne soit non nul que pour un nombre fini de $n$.}
$$
\mathcal{O} = \prod\nolimits_{n \in \mathbf{Z}} \mathcal{O}_n
$$
et, pour chaque $n$, il existe un $\mathcal{O}_n$-module inversible
$\mathcal{H}^n$ (Modules sur les sites, définition
\ref{sites-modules-definition-invertible-sheaf}), et
$M = \bigoplus \mathcal{H}^n[-n]$ dans $D(\mathcal{O})$.
\end{enumerate}
Si (1) et (2) sont satisfaites, alors $M$ est un objet parfait de
$D(\mathcal{O})$. Si $(\mathcal{C}, \mathcal{O})$ est un site localement
annelé, ces conditions sont également équivalentes à :
\begin{enumerate}
\item[(3)] pour tout objet $U$ de $\mathcal{C}$, il existe un recouvrement
$\{U_i \to U\}$ et, pour chaque $i$, un entier $n_i$ tel que $M|_{U_i}$ soit
représenté par un $\mathcal{O}_{U_i}$-module inversible placé au degré $n_i$.
\end{enumerate}
\end{lemma}

\begin{proof}
Supposons (2). Considérons l'objet $R\SheafHom(M, \mathcal{O})$ et le
morphisme de composition
$$
R\SheafHom(M, \mathcal{O}) \otimes_\mathcal{O}^\mathbf{L} M \to \mathcal{O}
$$
Pour montrer que c'est un isomorphisme, nous pouvons travailler localement.
Nous pouvons donc supposer que
$\mathcal{O} = \prod_{a \leq n \leq b} \mathcal{O}_n$ et
$M = \bigoplus_{a \leq n \leq b} \mathcal{H}^n[-n]$. Il suffit alors de
montrer que
$$
R\SheafHom(\mathcal{H}^m, \mathcal{O})
\otimes_\mathcal{O}^\mathbf{L} \mathcal{H}^n
$$
est nul si $n \not = m$ et égal à $\mathcal{O}_n$ si $n = m$. Le cas
$n \not = m$ résulte du fait que $\mathcal{O}_n$ et $\mathcal{O}_m$ sont des
$\mathcal{O}$-algèbres plates telles que
$\mathcal{O}_n \otimes_\mathcal{O} \mathcal{O}_m = 0$. En utilisant la
structure locale des $\mathcal{O}$-modules inversibles (Modules sur les sites,
lemme \ref{sites-modules-lemma-invertible}) et en travaillant localement,
l'isomorphisme dans le cas $n = m$ s'obtient immédiatement ; nous omettons
les détails. Puisque $D(\mathcal{O})$ est monoïdale symétrique, nous concluons
que $M$ est inversible.

\medskip\noindent
Supposons (1). La description de (2) montre que nous avons un candidat pour
$\mathcal{O}_n$, à savoir
$\SheafHom_\mathcal{O}(H^n(M), H^n(M))$. Si cette famille de faisceaux
d'anneaux est localement finie et si
$\mathcal{O} = \prod \mathcal{O}_n$, alors nous obtenons immédiatement la
décomposition en somme directe $M = \bigoplus H^n(M)[-n]$ en utilisant les
idempotents de $\mathcal{O}$ provenant de la décomposition en produit. Cela
montre que, pour démontrer (2), nous pouvons travailler localement au sens
suivant. Soit $U$ un objet de $\mathcal{C}$. Nous devons montrer qu'il existe
un recouvrement $\{U_i \to U\}$ de $U$ tel qu'avec les $\mathcal{O}_n$
ci-dessus, les assertions précédentes et celles de (2) soient satisfaites
après restriction à $\mathcal{C}/U_i$. Nous pouvons donc supposer que
$\mathcal{C}$ possède un objet final $X$ et, au cours de la démonstration de
(2), remplacer un nombre fini de fois $X$ par les membres d'un recouvrement
de $X$.

\medskip\noindent
Choisissons un objet $N$ de $D(\mathcal{O})$ et un isomorphisme
$M \otimes_\mathcal{O}^\mathbf{L} N \cong \mathcal{O}$. Alors $N$ est un dual
à gauche de $M$ dans la catégorie monoïdale $D(\mathcal{O})$, et nous
concluons que $M$ est parfait par le lemme \ref{sites-cohomology-lemma-left-dual-derived}. Par
symétrie, $N$ est parfait. Après avoir remplacé $X$ par les membres d'un
recouvrement, nous pouvons supposer que $M$ et $N$ sont représentés par des
complexes strictement parfaits $\mathcal{E}^\bullet$ et
$\mathcal{F}^\bullet$. Alors $M \otimes_\mathcal{O}^\mathbf{L} N$ est
représenté par
$\text{Tot}(\mathcal{E}^\bullet \otimes_\mathcal{O} \mathcal{F}^\bullet)$.
Après avoir remplacé $X$ par les membres d'un recouvrement de $X$, nous
pouvons supposer que les isomorphismes réciproques
$\mathcal{O} \to M \otimes_\mathcal{O}^\mathbf{L} N$ et
$M \otimes_\mathcal{O}^\mathbf{L} N \to \mathcal{O}$ sont donnés par des
morphismes de complexes
$$
\alpha : \mathcal{O} \to
\text{Tot}(\mathcal{E}^\bullet \otimes_\mathcal{O} \mathcal{F}^\bullet)
\quad\text{et}\quad
\beta :
\text{Tot}(\mathcal{E}^\bullet \otimes_\mathcal{O} \mathcal{F}^\bullet)
\to \mathcal{O}
$$
voir le lemme \ref{sites-cohomology-lemma-local-actual}. Alors $\beta \circ \alpha = 1$ comme
morphismes de complexes et $\alpha \circ \beta = 1$ comme morphisme dans
$D(\mathcal{O})$. Après avoir remplacé $X$ par les membres d'un recouvrement
de $X$, nous pouvons supposer que la composée $\alpha \circ \beta$ est
homotope à $1$ par une certaine homotopie $\theta$ de composantes
$$
\theta^n :
\text{Tot}^n(\mathcal{E}^\bullet \otimes_\mathcal{O} \mathcal{F}^\bullet)
\to
\text{Tot}^{n - 1}(
\mathcal{E}^\bullet \otimes_\mathcal{O} \mathcal{F}^\bullet)
$$
par le même lemme que précédemment. Posons
$R = \Gamma(\mathcal{C}, \mathcal{O})$. Par le lemme
\ref{sites-cohomology-lemma-category-summands-finite-free}, nous obtenons :
\begin{enumerate}
\item $M^\bullet = \Gamma(X, \mathcal{E}^\bullet)$
est un complexe borné de $R$-modules projectifs de type fini,
\item $N^\bullet = \Gamma(X, \mathcal{F}^\bullet)$
est un complexe borné de $R$-modules projectifs de type fini,
\item $\alpha$ et $\beta$ correspondent aux morphismes de complexes
$a : R \to \text{Tot}(M^\bullet \otimes_R N^\bullet)$ et
$b : \text{Tot}(M^\bullet \otimes_R N^\bullet) \to R$,
\item $\theta^n$ correspond au morphisme
$h^n : \text{Tot}^n(M^\bullet \otimes_R N^\bullet) \to
\text{Tot}^{n - 1}(M^\bullet \otimes_R N^\bullet)$, et
\item $b \circ a = 1$ et $a \circ b - 1 = dh + hd$,
\end{enumerate}
Il s'ensuit que $M^\bullet$ et $N^\bullet$ définissent des objets inverses
l'un de l'autre dans $D(R)$. Par Compléments d'algèbre, lemme
\ref{more-algebra-lemma-invertible-derived}, nous obtenons une décomposition
en produit $R = \prod_{a \leq n \leq b} R_n$ et des $R_n$-modules inversibles
$H^n$ tels que
$M^\bullet \cong \bigoplus_{a \leq n \leq b} H^n[-n]$. Cet isomorphisme dans
$D(R)$ peut être relevé en un morphisme
$$
\bigoplus H^n[-n] \longrightarrow M^\bullet
$$
de complexes, car chaque $H^n$ est projectif comme $R$-module. De manière
correspondante, en utilisant à nouveau le lemme
\ref{sites-cohomology-lemma-category-summands-finite-free}, nous obtenons un morphisme
$$
\bigoplus H^n \otimes_R \mathcal{O}[-n] \to \mathcal{E}^\bullet
$$
qui est un isomorphisme dans $D(\mathcal{O})$. Ici,
$M \otimes_R \mathcal{O}$ désigne le foncteur des $R$-modules projectifs de
type fini vers les $\mathcal{O}$-modules construit dans la démonstration du
lemme \ref{sites-cohomology-lemma-category-summands-finite-free}. En posant
$\mathcal{O}_n = R_n \otimes_R \mathcal{O}$, nous concluons que (2) est vraie.

\medskip\noindent
Si $(\mathcal{C}, \mathcal{O})$ est un site localement annelé, étant donnés
un objet $U$ et une décomposition finie en produit
$\mathcal{O}|_U = \prod_{a \leq n \leq b} \mathcal{O}_n|_U$, nous pouvons
trouver un recouvrement $\{U_i \to U\}$ tel que, pour tout $i$, il existe au
plus un $n$ pour lequel $\mathcal{O}_n|_{U_i}$ est non nul. Cela résulte
aisément de la partie (2) de Modules sur les sites, lemme
\ref{sites-modules-lemma-locally-ringed}, et de la définition des sites
localement annelés donnée dans Modules sur les sites, définition
\ref{sites-modules-definition-locally-ringed}. L'implication
(2) $\Rightarrow$ (3) s'en déduit facilement. L'implication
(3) $\Rightarrow$ (2) est vraie sans aucune hypothèse sur le site annelé.
Nous omettons les détails.
\end{proof}










\section{Formule de projection}
\label{sites-cohomology-section-projection-formula}

\noindent
Soit $f : (\Sh(\mathcal{C}), \mathcal{O}_\mathcal{C}) \to
(\Sh(\mathcal{D}), \mathcal{O}_\mathcal{D})$ un morphisme de topos annelés.
Soient $E \in D(\mathcal{O}_\mathcal{C})$ et
$K \in D(\mathcal{O}_\mathcal{D})$. Sans aucune hypothèse supplémentaire, il
existe un morphisme
\begin{equation}
\label{sites-cohomology-equation-projection-formula-map}
Rf_*E \otimes^\mathbf{L}_{\mathcal{O}_\mathcal{D}} K
\longrightarrow
Rf_*(E \otimes^\mathbf{L}_{\mathcal{O}_\mathcal{C}} Lf^*K)
\end{equation}
À savoir, c'est l'adjoint du morphisme canonique
$$
Lf^*(Rf_*E \otimes^\mathbf{L}_{\mathcal{O}_\mathcal{D}} K) =
Lf^*Rf_*E \otimes^\mathbf{L}_{\mathcal{O}_\mathcal{C}} Lf^*K
\longrightarrow
E \otimes^\mathbf{L}_{\mathcal{O}_\mathcal{C}} Lf^*K
$$
provenant du morphisme $Lf^*Rf_*E \to E$ et des lemmes
\ref{sites-cohomology-lemma-pullback-tensor-product} et \ref{sites-cohomology-lemma-adjoint}. Une version assez
générale de la formule de projection est la suivante.

\begin{lemma}
\label{sites-cohomology-lemma-projection-formula}
Soit $f : (\Sh(\mathcal{C}), \mathcal{O}_\mathcal{C}) \to
(\Sh(\mathcal{D}), \mathcal{O}_\mathcal{D})$ un morphisme de topos annelés.
Soient $E \in D(\mathcal{O}_\mathcal{C})$ et
$K \in D(\mathcal{O}_\mathcal{D})$. Si $K$ est parfait, alors
$$
Rf_*E \otimes^\mathbf{L}_{\mathcal{O}_\mathcal{D}} K =
Rf_*(E \otimes^\mathbf{L}_{\mathcal{O}_\mathcal{C}} Lf^*K)
$$
dans $D(\mathcal{O}_\mathcal{D})$.
\end{lemma}

\begin{proof}
Pour vérifier que (\ref{sites-cohomology-equation-projection-formula-map}) est un isomorphisme,
nous pouvons travailler localement sur $\mathcal{D}$, c'est-à-dire que, pour
tout objet $V$ de $\mathcal{D}$, nous devons trouver un recouvrement
$\{V_j \to V\}$ tel que la restriction du morphisme à $V_j$ soit un
isomorphisme. Par définition des objets parfaits, cela signifie que nous
pouvons supposer $K$ représenté par un complexe strictement parfait de
$\mathcal{O}_\mathcal{D}$-modules. Notons que, tout à fait généralement,
l'assertion est vraie pour $K = K_1 \oplus K_2$ si et seulement si elle est
vraie pour $K_1$ et $K_2$. Nous pouvons donc supposer que $K$ est un complexe
fini de $\mathcal{O}_\mathcal{D}$-modules libres de type fini. Dans ce cas, un
argument simple utilisant les troncations bêtes ramène l'assertion au cas où
$K$ est représenté par un $\mathcal{O}_\mathcal{D}$-module libre de type fini.
Puisque l'assertion est invariante par passage aux facteurs directs finis
dans la variable $K$, nous concluons qu'il suffit de la démontrer pour
$K = \mathcal{O}_\mathcal{D}[n]$, cas où elle est triviale.
\end{proof}

\begin{remark}
\label{sites-cohomology-remark-compatible-with-diagram}
Le morphisme (\ref{sites-cohomology-equation-projection-formula-map}) est compatible avec le
morphisme de changement de base de la remarque \ref{sites-cohomology-remark-base-change} au
sens suivant. Supposons en effet que
$$
\xymatrix{
(\Sh(\mathcal{C}'), \mathcal{O}_{\mathcal{C}'})
\ar[r]_{g'} \ar[d]_{f'} &
(\Sh(\mathcal{C}), \mathcal{O}_\mathcal{C}) \ar[d]^f \\
(\Sh(\mathcal{D}'), \mathcal{O}_{\mathcal{D}'})
\ar[r]^g &
(\Sh(\mathcal{D}), \mathcal{O}_\mathcal{D})
}
$$
soit un diagramme commutatif de topos annelés. Soient
$E \in D(\mathcal{O}_\mathcal{C})$ et $K \in D(\mathcal{O}_\mathcal{D})$.
Alors le diagramme
$$
\xymatrix{
Lg^*(Rf_*E \otimes^\mathbf{L}_{\mathcal{O}_\mathcal{D}} K)
\ar[r]_p \ar[d]_t &
Lg^*Rf_*(E \otimes^\mathbf{L}_{\mathcal{O}_\mathcal{C}} Lf^*K)
\ar[d]_b \\
Lg^*Rf_*E \otimes^\mathbf{L}_{\mathcal{O}_{\mathcal{D}'}}
Lg^*K \ar[d]_b &
Rf'_*L(g')^*(E \otimes^\mathbf{L}_{\mathcal{O}_\mathcal{C}}
Lf^*K) \ar[d]_t \\
Rf'_*L(g')^*E \otimes^\mathbf{L}_{\mathcal{O}_{\mathcal{D}'}}
Lg^*K \ar[rd]_p &
Rf'_*(L(g')^*E \otimes^\mathbf{L}_{\mathcal{O}_{\mathcal{D}'}}
L(g')^*Lf^*K) \ar[d]_c \\
& Rf'_*(L(g')^*E \otimes^\mathbf{L}_{\mathcal{O}_{\mathcal{D}'}}
L(f')^*Lg^*K)
}
$$
est commutatif. Ici, les flèches étiquetées $t$ sont obtenues par une
application du lemme \ref{sites-cohomology-lemma-pullback-tensor-product}, celles étiquetées
$b$ par une application de la remarque \ref{sites-cohomology-remark-base-change}, celles
étiquetées $p$ par une application de
(\ref{sites-cohomology-equation-projection-formula-map}), et $c$ provient de
$L(g')^* \circ Lf^* = L(f')^* \circ Lg^*$. Nous omettons la vérification.
\end{remark}






\section{Objets faiblement contractiles}
\label{sites-cohomology-section-w-contractible}

\noindent
Un objet $U$ d'un site est {\it faiblement contractile} si toute surjection
$\mathcal{F} \to \mathcal{G}$ de faisceaux d'ensembles donne lieu à une
surjection $\mathcal{F}(U) \to \mathcal{G}(U)$, voir Sites, définition
\ref{sites-definition-w-contractible}.

\begin{lemma}
\label{sites-cohomology-lemma-w-contractible}
Soit $\mathcal{C}$ un site. Soit $U$ un objet faiblement contractile de
$\mathcal{C}$. Alors :
\begin{enumerate}
\item le foncteur $\mathcal{F} \mapsto \mathcal{F}(U)$ est un foncteur exact
$\textit{Ab}(\mathcal{C}) \to \textit{Ab}$,
\item $H^p(U, \mathcal{F}) = 0$ pour tout faisceau abélien $\mathcal{F}$ et
tout $p \geq 1$, et
\item pour tout faisceau de groupes $\mathcal{G}$, tout
$\mathcal{G}$-torseur possède une section sur $U$.
\end{enumerate}
\end{lemma}

\begin{proof}
La première assertion résulte immédiatement de la définition (voir aussi
Homology, section \ref{homology-section-functors}). Les foncteurs dérivés
supérieurs s'annulent par Catégories dérivées, lemme
\ref{derived-lemma-right-derived-exact-functor}. Soit $\mathcal{F}$ un
$\mathcal{G}$-torseur. Alors $\mathcal{F} \to *$ est un morphisme surjectif de
faisceaux. Ainsi, (3) résulte également de la définition.
\end{proof}

\noindent
Il est utile d'énumérer ici quelques conséquences de l'existence de
suffisamment d'objets faiblement contractiles.

\begin{proposition}
\label{sites-cohomology-proposition-enough-weakly-contractibles}
Soit $\mathcal{C}$ un site. Soit $\mathcal{B} \subset \Ob(\mathcal{C})$ tel
que tout $U \in \mathcal{B}$ soit faiblement contractile et que tout objet de
$\mathcal{C}$ possède un recouvrement par des éléments de $\mathcal{B}$. Soit
$\mathcal{O}$ un faisceau d'anneaux sur $\mathcal{C}$. Alors :
\begin{enumerate}
\item Un complexe $\mathcal{F}_1 \to \mathcal{F}_2 \to \mathcal{F}_3$ de
$\mathcal{O}$-modules est exact si et seulement si
$\mathcal{F}_1(U) \to \mathcal{F}_2(U) \to \mathcal{F}_3(U)$
est exact pour tout $U \in \mathcal{B}$.
\item Tout objet $K$ de $D(\mathcal{O})$ est la limite dérivée de ses
troncations canoniques : $K = R\lim \tau_{\geq -n} K$.
\item Étant donné un système inverse
$\ldots \to \mathcal{F}_3 \to \mathcal{F}_2 \to \mathcal{F}_1$
dont les morphismes de transition sont surjectifs, la projection
$\lim \mathcal{F}_n \to \mathcal{F}_1$ est surjective.
\item Les produits sont exacts dans $\textit{Mod}(\mathcal{O})$.
\item Les produits dans $D(\mathcal{O})$ peuvent être calculés en prenant les
produits de complexes représentants quelconques.
\item Si $(\mathcal{F}_n)$ est un système inverse de $\mathcal{O}$-modules,
alors $R^p\lim \mathcal{F}_n = 0$ pour tout $p > 1$ et
$$
R^1\lim \mathcal{F}_n  =
\Coker(\prod \mathcal{F}_n \to \prod \mathcal{F}_n)
$$
où le morphisme est $(x_n) \mapsto (x_n - f(x_{n + 1}))$.
\item Si $(K_n)$ est un système inverse d'objets de $D(\mathcal{O})$, alors il
existe des suites exactes courtes
$$
0 \to R^1\lim H^{p - 1}(K_n) \to H^p(R\lim K_n) \to
\lim H^p(K_n) \to 0
$$
\end{enumerate}
\end{proposition}

\begin{proof}
Démonstration de (1). Si la suite est exacte, son évaluation en tout élément
faiblement contractile de $\mathcal{C}$ donne une suite exacte par le lemme
\ref{sites-cohomology-lemma-w-contractible}. Réciproquement, supposons que
$\mathcal{F}_1(U) \to \mathcal{F}_2(U) \to \mathcal{F}_3(U)$
soit exact pour tout $U \in \mathcal{B}$. Soit $V$ un objet de
$\mathcal{C}$ et soit $s \in \mathcal{F}_2(V)$ un élément du noyau de
$\mathcal{F}_2 \to \mathcal{F}_3$. Par hypothèse, il existe un recouvrement
$\{U_i \to V\}$ avec $U_i \in \mathcal{B}$. Alors $s|_{U_i}$ se relève en
une section $s_i \in \mathcal{F}_1(U_i)$. Ainsi, $s$ est une section du
faisceau image $\Im(\mathcal{F}_1 \to \mathcal{F}_2)$. Autrement dit, la
suite $\mathcal{F}_1 \to \mathcal{F}_2 \to \mathcal{F}_3$ est exacte.

\medskip\noindent
Démonstration de (2). Cela résulte du lemme
\ref{sites-cohomology-lemma-is-limit-dimension} avec $d = 0$.

\medskip\noindent
Démonstration de (3). Soit $(\mathcal{F}_n)$ un système comme en (3), et
posons $\mathcal{F} = \lim \mathcal{F}_n$. Si $U \in \mathcal{B}$, alors
$\mathcal{F}(U) = \lim \mathcal{F}_n(U)$
se surjecte sur $\mathcal{F}_1(U)$ puisque tous les morphismes de transition
$\mathcal{F}_{n + 1}(U) \to \mathcal{F}_n(U)$ sont surjectifs. Ainsi,
$\mathcal{F} \to \mathcal{F}_1$ est surjectif par Sites, définition
\ref{sites-definition-sheaves-injective-surjective}, et par l'hypothèse que
tout objet possède un recouvrement par des éléments de $\mathcal{B}$.

\medskip\noindent
Démonstration de (4). Soit
$\mathcal{F}_{i, 1} \to \mathcal{F}_{i, 2} \to \mathcal{F}_{i, 3}$
une famille de suites exactes de $\mathcal{O}$-modules. Nous voulons montrer
que
$\prod \mathcal{F}_{i, 1} \to \prod \mathcal{F}_{i, 2} \to
\prod \mathcal{F}_{i, 3}$ est exacte. Nous utilisons le critère de (1).
Soit $U \in \mathcal{B}$. Alors
$$
(\prod \mathcal{F}_{i, 1})(U) \to
(\prod \mathcal{F}_{i, 2})(U) \to
(\prod \mathcal{F}_{i, 3})(U)
$$
est la même suite que
$$
\prod \mathcal{F}_{i, 1}(U) \to
\prod \mathcal{F}_{i, 2}(U) \to
\prod \mathcal{F}_{i, 3}(U)
$$
Chacune des suites
$\mathcal{F}_{i, 1}(U) \to \mathcal{F}_{i, 2}(U) \to \mathcal{F}_{i, 3}(U)$
est exacte par (1). Ainsi, les suites affichées sont exactes par Homology,
lemme \ref{homology-lemma-product-abelian-groups-exact}. Nous concluons en
appliquant à nouveau (1).

\medskip\noindent
Démonstration de (5). Cela résulte de (4) et de Catégories dérivées, lemme
\ref{derived-lemma-products} (légèrement généralisé).

\medskip\noindent
Démonstration de (6) et (7). Nous renvoyons à la section
\ref{sites-cohomology-section-derived-limits} pour une discussion des limites dérivées, des
limites homotopiques et de leur relation. Par Catégories dérivées, définition
\ref{derived-definition-derived-limit}, nous avons un triangle distingué
$$
R\lim K_n \to \prod K_n \to \prod K_n \to R\lim K_n[1]
$$
En prenant la suite exacte longue des faisceaux de cohomologie, nous obtenons
$$
H^{p - 1}(\prod K_n) \to H^{p - 1}(\prod K_n) \to
H^p(R\lim K_n) \to H^p(\prod K_n) \to H^p(\prod K_n)
$$
Puisque les produits sont exacts par (4), cela devient
$$
\prod H^{p - 1}(K_n) \to \prod H^{p - 1}(K_n) \to
H^p(R\lim K_n) \to \prod H^p(K_n) \to \prod H^p(K_n)
$$
Appliquons d'abord ceci au cas $K_n = \mathcal{F}_n[0]$, où
$(\mathcal{F}_n)$ est comme en (6). Nous concluons que (6) est satisfaite.
Appliquons-le ensuite à $(K_n)$ comme en (7) et nous concluons que (7) est
satisfaite.
\end{proof}






\section{Objets compacts}
\label{sites-cohomology-section-compact}

\noindent
Dans cette section, nous étudions les objets compacts de la catégorie dérivée
des modules sur un site annelé. Rappelons que les objets compacts sont définis
dans Catégories dérivées, définition \ref{derived-definition-compact-object}.

\begin{lemma}
\label{sites-cohomology-lemma-compact-in-terms-of-generators}
Soit $\mathcal{A}$ une catégorie abélienne de Grothendieck. Soit
$S \subset \Ob(\mathcal{A})$ un ensemble d'objets tel que :
\begin{enumerate}
\item tout objet de $\mathcal{A}$ est un quotient d'une somme directe
d'éléments de $S$, et
\item pour tout $E \in S$, le foncteur $\Hom_\mathcal{A}(E, -)$ commute aux
sommes directes.
\end{enumerate}
Alors tout objet compact de $D(\mathcal{A})$ est facteur direct dans
$D(\mathcal{A})$ d'un complexe fini de sommes directes finies d'éléments de
$S$.
\end{lemma}

\begin{proof}
Supposons que $K \in D(\mathcal{A})$ soit un objet compact. Représentons $K$
par un complexe $K^\bullet$ et considérons le morphisme
$$
K^\bullet
\longrightarrow
\bigoplus\nolimits_{n \geq 0} \tau_{\geq n} K^\bullet
$$
où nous avons utilisé les troncations canoniques, voir Homology, section
\ref{homology-section-truncations}. Cela a un sens, car, à chaque degré, la
somme directe de droite est finie. Par hypothèse, ce morphisme se factorise
par une somme directe finie. Nous concluons que
$K \to \tau_{\geq n} K$ est nul pour au moins un $n$, c'est-à-dire que $K$
appartient à $D^{-}(\mathcal{A})$.

\medskip\noindent
Nous pouvons représenter $K$ par un complexe borné supérieurement
$K^\bullet$, dont chaque terme est une somme directe d'objets de $S$, voir
Catégories dérivées, lemme \ref{derived-lemma-subcategory-left-resolution}.
Notons que
$$
K^\bullet = \bigcup\nolimits_{n \leq 0} \sigma_{\geq n}K^\bullet
$$
où nous avons utilisé les troncations bêtes, voir Homology, section
\ref{homology-section-truncations}. Par Catégories dérivées, lemmes
\ref{derived-lemma-colim-hocolim} et
\ref{derived-lemma-commutes-with-countable-sums}, nous voyons donc que
$1 : K^\bullet \to K^\bullet$ se factorise par
$\sigma_{\geq n}K^\bullet \to K^\bullet$ dans $D(\mathcal{A})$. Ainsi,
$1 : K^\bullet \to K^\bullet$ se factorise sous la forme
$$
K^\bullet \xrightarrow{\varphi} L^\bullet \xrightarrow{\psi} K^\bullet
$$
dans $D(\mathcal{A})$, pour un certain complexe borné $L^\bullet$ dont les
termes sont des sommes directes d'éléments de $S$. Disons que $L^i$ est nul
pour $i \not \in [a, b]$. Soit $c$ le plus grand entier $\leq b + 1$ tel que
$L^i$ soit une somme directe finie d'éléments de $S$ pour $i < c$.
Affirmation : si $c < b + 1$, alors nous pouvons modifier $L^\bullet$ pour
augmenter $c$. Par récurrence, cette affirmation montrera que nous avons une
factorisation de $1_K$ sous la forme
$$
K \xrightarrow{\varphi} L \xrightarrow{\psi} K
$$
dans $D(\mathcal{A})$, où $L$ peut être représenté par un complexe fini de
sommes directes finies d'éléments de $S$. Notons que
$e = \varphi \circ \psi \in \text{End}_{D(\mathcal{A})}(L)$ est un
idempotent. Par Catégories dérivées, lemme
\ref{derived-lemma-projectors-have-images-triangulated}, on voit que
$L = \Ker(e) \oplus \Ker(1 - e)$. Le morphisme $\varphi : K \to L$ induit un
isomorphisme avec $\Ker(1 - e)$ dans $D(\mathcal{A})$, et nous concluons.

\medskip\noindent
Démonstration de l'affirmation. Écrivons
$L^c = \bigoplus_{\lambda \in \Lambda} E_\lambda$. Puisque $L^{c - 1}$ est
une somme directe finie d'éléments de $S$, l'hypothèse (2) nous permet de
trouver un sous-ensemble fini $\Lambda' \subset \Lambda$ tel que
$L^{c - 1} \to L^c$ se factorise par
$\bigoplus_{\lambda \in \Lambda'} E_\lambda \subset L^c$. Considérons le
morphisme de complexes
$$
\pi :
L^\bullet
\longrightarrow
(\bigoplus\nolimits_{\lambda \in \Lambda \setminus \Lambda'} E_\lambda)[-c]
$$
donné par la projection sur les facteurs correspondant à
$\Lambda \setminus \Lambda'$ au degré $c$. Par notre hypothèse sur $K$, on
voit qu'après avoir éventuellement remplacé $\Lambda'$ par un sous-ensemble
fini plus grand, nous pouvons supposer que $\pi \circ \varphi = 0$ dans
$D(\mathcal{A})$. Soit $(L')^\bullet \subset L^\bullet$ le noyau de $\pi$.
Comme $\pi$ est surjectif, nous obtenons une suite exacte courte de complexes,
qui donne un triangle distingué dans $D(\mathcal{A})$ (voir Derived
Catégories, lemme \ref{derived-lemma-derived-canonical-delta-functor}).
Puisque $\Hom_{D(\mathcal{A})}(K, -)$ est homologique (voir Derived
Catégories, lemme \ref{derived-lemma-representable-homological}) et que
$\pi \circ \varphi = 0$, nous pouvons trouver un morphisme
$\varphi' : K^\bullet \to (L')^\bullet$ dans $D(\mathcal{A})$ dont la
composée avec $(L')^\bullet \to L^\bullet$ est $\varphi$. En prenant pour
$\psi'$ la composée de $\psi$ avec $(L')^\bullet \to L^\bullet$, nous
obtenons une nouvelle factorisation. Puisque $(L')^\bullet$ coïncide avec
$L^\bullet$ sauf au degré $c$ et que
$(L')^c = \bigoplus_{\lambda \in \Lambda'} E_\lambda$, l'affirmation est
démontrée.
\end{proof}

\begin{lemma}
\label{sites-cohomology-lemma-compact-objects-if-enough-qc}
Soit $(\mathcal{C}, \mathcal{O})$ un site annelé. Supposons que tout objet de
$\mathcal{C}$ possède un recouvrement par des objets quasi-compacts. Alors
tout objet compact de $D(\mathcal{O})$ est facteur direct dans
$D(\mathcal{O})$ d'un complexe fini dont les termes sont des sommes directes
finies de $\mathcal{O}$-modules de la forme $j_!\mathcal{O}_U$, où $U$ est un
objet quasi-compact de $\mathcal{C}$.
\end{lemma}

\begin{proof}
Appliquons le lemme \ref{sites-cohomology-lemma-compact-in-terms-of-generators}, où
$S \subset \Ob(\textit{Mod}(\mathcal{O}))$ est l'ensemble des modules de la
forme $j_!\mathcal{O}_U$ avec $U \in \Ob(\mathcal{C})$ quasi-compact.
L'hypothèse (1) est satisfaite par Modules sur les sites, lemme
\ref{sites-modules-lemma-module-quotient-flat}, et par l'hypothèse que tout
$U$ peut être recouvert par des objets quasi-compacts. L'hypothèse (2) résulte
de
$$
\Hom_\mathcal{O}(j_!\mathcal{O}_U, \mathcal{F}) = \mathcal{F}(U)
$$
qui commute aux sommes directes par Sites, lemme
\ref{sites-lemma-directed-colimits-sections}.
\end{proof}

\noindent
Dans la situation du lemme ci-dessus, il n'est pas toujours vrai que les
modules $j_!\mathcal{O}_U$ soient des objets compacts de $D(\mathcal{O})$
(même si $U$ est un objet quasi-compact de $\mathcal{C}$). Voici deux lemmes
qui traitent cette question.


\begin{lemma}
\label{sites-cohomology-lemma-when-jshriek-lower-compact}
Soit $(\mathcal{C}, \mathcal{O})$ un site annelé. Soit $U$ un objet de
$\mathcal{C}$. Supposons que les foncteurs
$\mathcal{F} \mapsto H^p(U, \mathcal{F})$ commutent aux sommes directes. Alors
le $\mathcal{O}$-module $j_!\mathcal{O}_U$ est un objet compact de
$D^+(\mathcal{O})$ au sens suivant : si
$M = \bigoplus_{i \in I} M_i$ dans $D(\mathcal{O})$ est borné inférieurement,
alors $\Hom(j_{U!}\mathcal{O}_U, M) =
\bigoplus_{i \in I} \Hom(j_{U!}\mathcal{O}_U, M_i)$.
\end{lemma}

\begin{proof}
Puisque $\Hom(j_{U!}\mathcal{O}_U, -)$ est le même foncteur que
$\mathcal{F} \mapsto \mathcal{F}(U)$ par Modules sur les sites, équation
(\ref{sites-modules-equation-map-lower-shriek-OU-into-module}), il suffit de
montrer que $H^p(U, M) = \bigoplus H^p(U, M_i)$. Soit $\mathcal{I}_i$,
$i \in I$, une famille de $\mathcal{O}$-modules injectifs. Par hypothèse,
nous avons
$$
H^p(U, \bigoplus\nolimits_{i \in I} \mathcal{I}_i) =
\bigoplus\nolimits_{i \in I} H^p(U, \mathcal{I}_i) = 0
$$
pour tout $p > 0$. Puisque $M = \bigoplus M_i$ est borné inférieurement, il
existe un $a \in \mathbf{Z}$ tel que $H^n(M_i) = 0$ pour $n < a$. Nous
pouvons donc choisir des complexes de $\mathcal{O}$-modules injectifs
$\mathcal{I}_i^\bullet$ représentant $M_i$ avec $\mathcal{I}_i^n = 0$ pour
$n < a$, voir Catégories dérivées, lemme
\ref{derived-lemma-injective-resolutions-exist}. Par Injectifs, lemme
\ref{injectives-lemma-derived-products}, on voit que le complexe somme directe
$\bigoplus \mathcal{I}_i^\bullet$ représente $M$. Par acyclicité de Leray
(Catégories dérivées, lemme \ref{derived-lemma-leray-acyclicity}), on voit que
$$
R\Gamma(U, M) = \Gamma(U, \bigoplus \mathcal{I}_i^\bullet) =
\bigoplus \Gamma(U, \mathcal{I}_i^\bullet) =
\bigoplus R\Gamma(U, M_i)
$$
comme souhaité.
\end{proof}

\begin{lemma}
\label{sites-cohomology-lemma-when-jshriek-lower-compact-worked-out}
Soit $(\mathcal{C}, \mathcal{O})$ un site annelé dont l'ensemble des
recouvrements est $\text{Cov}_\mathcal{C}$. Soient
$\mathcal{B} \subset \Ob(\mathcal{C})$ et
$\text{Cov} \subset \text{Cov}_\mathcal{C}$
des sous-ensembles. Supposons que :
\begin{enumerate}
\item Pour tout $\mathcal{U} \in \text{Cov}$, nous avons
$\mathcal{U} = \{U_i \to U\}_{i \in I}$ avec $I$ fini,
$U, U_i \in \mathcal{B}$ et tout
$U_{i_0} \times_U \ldots \times_U U_{i_p} \in \mathcal{B}$.
\item Pour tout $U \in \mathcal{B}$, les recouvrements de $U$ appartenant à
$\text{Cov}$ forment un système cofinal de recouvrements de $U$.
\end{enumerate}
Alors, pour $U \in \mathcal{B}$, l'objet $j_{U!}\mathcal{O}_U$ est un objet
compact de $D^+(\mathcal{O})$ au sens suivant : si
$M = \bigoplus_{i \in I} M_i$ dans $D(\mathcal{O})$ est borné inférieurement,
alors $\Hom(j_{U!}\mathcal{O}_U, M) =
\bigoplus_{i \in I} \Hom(j_{U!}\mathcal{O}_U, M_i)$.
\end{lemma}

\begin{proof}
Cela résulte des lemmes \ref{sites-cohomology-lemma-when-jshriek-lower-compact} et
\ref{sites-cohomology-lemma-colim-works-over-collection}.
\end{proof}

\begin{lemma}
\label{sites-cohomology-lemma-when-jshriek-compact}
Soit $(\mathcal{C}, \mathcal{O})$ un site annelé. Soit $U$ un objet de
$\mathcal{C}$. Le $\mathcal{O}$-module $j_!\mathcal{O}_U$ est un objet
compact de $D(\mathcal{O})$ s'il existe un entier $d$ tel que :
\begin{enumerate}
\item $H^p(U, \mathcal{F}) = 0$ pour tout $p > d$, et
\item les foncteurs $\mathcal{F} \mapsto H^p(U, \mathcal{F})$ commutent aux
sommes directes.
\end{enumerate}
\end{lemma}

\begin{proof}
Supposons (1) et (2). Rappelons que
$\Hom(j_!\mathcal{O}_U, K) = R\Gamma(U, K)$ pour $K$ dans $D(\mathcal{O})$.
Nous devons donc montrer que $R\Gamma(U, -)$ commute aux sommes directes. La
première hypothèse signifie que le foncteur $F = H^0(U, -)$ est de dimension
cohomologique finie. De plus, la seconde hypothèse implique que toute somme
directe de modules injectifs est acyclique pour $F$. Soit $K_i$ une famille
d'objets de $D(\mathcal{O})$. Choisissons des représentants K-injectifs
$I_i^\bullet$ à termes injectifs représentant $K_i$, voir Injectifs,
théorème \ref{injectives-theorem-K-injective-embedding-grothendieck}. Puisque
nous pouvons calculer $RF$ en appliquant $F$ à tout complexe d'acycliques
(Catégories dérivées, lemme \ref{derived-lemma-unbounded-right-derived}) et
que $\bigoplus K_i$ est représenté par $\bigoplus I_i^\bullet$ (Injectifs,
lemme \ref{injectives-lemma-derived-products}), nous concluons que
$R\Gamma(U, \bigoplus K_i)$ est représenté par
$\bigoplus H^0(U, I_i^\bullet)$. Ainsi, $R\Gamma(U, -)$ commute aux sommes
directes, comme souhaité.
\end{proof}

\begin{lemma}
\label{sites-cohomology-lemma-quasi-compact-weakly-contractible-compact}
Soit $(\mathcal{C}, \mathcal{O})$ un site annelé. Soit $U$ un objet de
$\mathcal{C}$ qui est quasi-compact et faiblement contractile. Alors
$j_!\mathcal{O}_U$ est un objet compact de $D(\mathcal{O})$.
\end{lemma}

\begin{proof}
Combiner les lemmes \ref{sites-cohomology-lemma-when-jshriek-compact} et
\ref{sites-cohomology-lemma-w-contractible} avec Modules sur les sites, lemme
\ref{sites-modules-lemma-sections-over-quasi-compact}.
\end{proof}

\begin{lemma}
\label{sites-cohomology-lemma-perfect-is-compact}
Soit $(\mathcal{C}, \mathcal{O})$ un site annelé. Supposons que
$\mathcal{C}$ possède les propriétés suivantes :
\begin{enumerate}
\item $\mathcal{C}$ possède un objet final quasi-compact $X$,
\item tout objet quasi-compact de $\mathcal{C}$ possède un système cofinal de
recouvrements finis constitués d'objets quasi-compacts,
\item pour un recouvrement fini $\{U_i \to U\}_{i \in I}$, où $U$ et les
$U_i$ sont quasi-compacts, les produits fibrés $U_i \times_U U_j$ sont
quasi-compacts.
\end{enumerate}
Soit $K$ un objet parfait de $D(\mathcal{O})$. Alors :
\begin{enumerate}
\item[(a)] $K$ est un objet compact de $D^+(\mathcal{O})$ au sens suivant :
si $M = \bigoplus_{i \in I} M_i$ est borné inférieurement, alors
$\Hom(K, M) = \bigoplus_{i \in I} \Hom(K, M_i)$.
\item[(b)] Si $(\mathcal{C}, \mathcal{O})$ est de dimension cohomologique
finie, c'est-à-dire s'il existe un $d$ tel que
$H^i(X, \mathcal{F}) = 0$ pour $i > d$ et pour tout $\mathcal{O}$-module
$\mathcal{F}$, alors $K$ est un objet compact de $D(\mathcal{O})$.
\end{enumerate}
\end{lemma}

\begin{proof}
Soit $K^\vee$ le dual de $K$, voir le lemme
\ref{sites-cohomology-lemma-dual-perfect-complex}. Alors nous avons
$$
\Hom_{D(\mathcal{O})}(K, M) =
H^0(X, K^\vee \otimes_\mathcal{O}^\mathbf{L} M)
$$
fonctoriellement en $M$ dans $D(\mathcal{O})$. Puisque
$K^\vee \otimes_\mathcal{O}^\mathbf{L} -$ commute aux sommes directes, il
suffit de montrer que $R\Gamma(X, -)$ commute aux sommes directes concernées.

\medskip\noindent
Démonstration de (a). Après reformulation comme ci-dessus, c'est un cas
particulier du lemme \ref{sites-cohomology-lemma-when-jshriek-lower-compact-worked-out} avec
$U = X$.

\medskip\noindent
Démonstration de (b). Puisque $R\Gamma(X, K) = R\Hom(\mathcal{O}, K)$ et que
$H^p(X, -)$ commute aux sommes directes par le lemme
\ref{sites-cohomology-lemma-colim-works-over-collection}, c'est un cas particulier du lemme
\ref{sites-cohomology-lemma-when-jshriek-compact}.
\end{proof}






\section{Complexes à faisceaux de cohomologie localement constants}
\sectionmark{Complexes à cohomologie localement constante}
\label{sites-cohomology-section-locally-constant}

\noindent
Les faisceaux localement constants sont introduits dans Modules sur les sites,
section \ref{sites-modules-section-locally-constant}. Soit $\mathcal{C}$ un
site. Soit $\Lambda$ un anneau. Nous notons $D(\mathcal{C}, \Lambda)$ la
catégorie dérivée de la catégorie abélienne des
$\underline{\Lambda}$-modules sur $\mathcal{C}$.

\begin{lemma}
\label{sites-cohomology-lemma-locally-constant}
Soit $\mathcal{C}$ un site possédant un objet final $X$. Soit $\Lambda$ un
anneau noethérien. Soit $K \in D^b(\mathcal{C}, \Lambda)$ tel que les
$H^i(K)$ soient des faisceaux localement constants de $\Lambda$-modules de
type fini. Alors il existe un recouvrement $\{U_i \to X\}$ tel que chaque
$K|_{U_i}$ soit représenté par un complexe de faisceaux localement constants
de $\Lambda$-modules de type fini.
\end{lemma}

\begin{proof}
Soient $a \leq b$ tels que $H^i(K) = 0$ pour $i \not \in [a, b]$. Par
récurrence sur $b - a$, nous allons démontrer qu'il existe un recouvrement
$\{U_i \to X\}$ tel que $K|_{U_i}$ puisse être représenté par un complexe
$\underline{M^\bullet}_{U_i}$, où $M^p$ est un $\Lambda$-module de type fini
et $M^p = 0$ pour $p \not \in [a, b]$. Si $b = a$, c'est clair. En général,
nous pouvons remplacer $X$ par les membres d'un recouvrement et supposer que
$H^b(K)$ est constant, disons $H^b(K) = \underline{M}$. Par Modules sur les sites,
lemme \ref{sites-modules-lemma-locally-constant-finite-type}, le module $M$
est un $\Lambda$-module de type fini. Choisissons une surjection
$\Lambda^{\oplus r} \to M$ donnée par des générateurs $x_1, \ldots, x_r$ de
$M$.

\medskip\noindent
Par une légère généralisation du lemme
\ref{sites-cohomology-lemma-kill-cohomology-class-on-covering} (détails omis), il existe un
recouvrement $\{U_i \to X\}$ tel que, pour tout $i$ et tout $j = 1, \ldots, r$,
la classe $x_j \in H^0(X, H^b(K))$ se relève en un élément de $H^b(U_i, K)$.
Ainsi, après avoir remplacé $X$ par les $U_i$, nous
parvenons à une situation où il existe un morphisme
$\underline{\Lambda^{\oplus r}}[-b] \to K$ induisant une surjection sur les
faisceaux de cohomologie au degré $b$. Choisissons un triangle distingué
$$
\underline{\Lambda^{\oplus r}}[-b] \to K \to L \to
\underline{\Lambda^{\oplus r}}[-b + 1]
$$
Les faisceaux de cohomologie de $L$ ne sont alors non nuls que dans
l'intervalle $[a, b - 1]$, coïncident avec ceux de $K$ dans l'intervalle
$[a, b - 2]$, et il existe une suite exacte courte
$$
0 \to H^{b - 1}(K) \to H^{b - 1}(L) \to
\underline{\Ker(\Lambda^{\oplus r} \to M)} \to 0
$$
au degré $b - 1$. Par Modules sur les sites, lemme
\ref{sites-modules-lemma-kernel-finite-locally-constant}, on voit que
$H^{b - 1}(L)$ est localement constant de type fini. Par l'hypothèse de
récurrence, nous obtenons un isomorphisme
$\underline{M^\bullet} \to L$ dans $D(\mathcal{C}, \Lambda)$,
où $M^p$ est un $\Lambda$-module de type fini et $M^p = 0$ pour
$p \not \in [a, b - 1]$. Le morphisme
$L \to \underline{\Lambda^{\oplus r}}[-b + 1]$ donne un morphisme
$\underline{M^{b - 1}} \to \underline{\Lambda^{\oplus r}}$ qui est
localement constant (Modules sur les sites, lemme
\ref{sites-modules-lemma-morphism-locally-constant}). Nous pouvons donc
supposer qu'il est donné par un morphisme
$M^{b - 1} \to \Lambda^{\oplus r}$. Le triangle distingué montre que la
composée $M^{b - 2} \to M^{b - 1} \to \Lambda^{\oplus r}$ est nulle, et les
axiomes des catégories triangulées produisent un isomorphisme
$$
\underline{M^a \to \ldots \to M^{b - 1} \to \Lambda^{\oplus r}}
\longrightarrow K
$$
dans $D(\mathcal{C}, \Lambda)$.
\end{proof}

\noindent
Soit $\mathcal{C}$ un site. Soit $\Lambda$ un anneau. En utilisant le
morphisme $\Sh(\mathcal{C}) \to \Sh(pt)$, on voit qu'il existe un foncteur
$D(\Lambda) \to D(\mathcal{C}, \Lambda)$, $K \mapsto \underline{K}$.

\begin{lemma}
\label{sites-cohomology-lemma-map-out-of-locally-constant}
Soit $\mathcal{C}$ un site possédant un objet final $X$. Soit $\Lambda$ un
anneau. Soient :
\begin{enumerate}
\item $K$ un objet parfait de $D(\Lambda)$,
\item $K^\bullet$ un complexe fini de $\Lambda$-modules projectifs de type
fini représentant $K$,
\item $\mathcal{L}^\bullet$ un complexe de faisceaux de $\Lambda$-modules, et
\item $\varphi : \underline{K} \to \mathcal{L}^\bullet$ un morphisme dans
$D(\mathcal{C}, \Lambda)$.
\end{enumerate}
Alors il existe un recouvrement $\{U_i \to X\}$ et des morphismes de complexes
$\alpha_i : \underline{K}^\bullet|_{U_i} \to \mathcal{L}^\bullet|_{U_i}$
représentant $\varphi|_{U_i}$.
\end{lemma}

\begin{proof}
Cela résulte immédiatement du lemme \ref{sites-cohomology-lemma-local-actual}.
\end{proof}

\begin{lemma}
\label{sites-cohomology-lemma-locally-constant-map}
Soit $\mathcal{C}$ un site possédant un objet final $X$. Soit $\Lambda$ un
anneau. Soient $K, L$ des objets de $D(\Lambda)$, avec $K$ parfait. Soit
$\varphi : \underline{K} \to \underline{L}$ un morphisme dans
$D(\mathcal{C}, \Lambda)$. Il existe un recouvrement $\{U_i \to X\}$ tel que
$\varphi|_{U_i}$ soit égal à $\underline{\alpha_i}$ pour un certain morphisme
$\alpha_i : K \to L$ dans $D(\Lambda)$.
\end{lemma}

\begin{proof}
Cela résulte du lemme \ref{sites-cohomology-lemma-map-out-of-locally-constant} et de Modules on
Sites, lemme \ref{sites-modules-lemma-morphism-locally-constant}.
\end{proof}

\begin{lemma}
\label{sites-cohomology-lemma-locally-constant-tensor-product}
Soit $\mathcal{C}$ un site. Soit $\Lambda$ un anneau noethérien. Soient
$K, L \in D^-(\mathcal{C}, \Lambda)$. Si les faisceaux de cohomologie de $K$
et de $L$ sont des faisceaux localement constants de $\Lambda$-modules de
type fini, alors les faisceaux de cohomologie de
$K \otimes_\Lambda^\mathbf{L} L$ sont des faisceaux localement constants de
$\Lambda$-modules de type fini.
\end{lemma}

\begin{proof}
Nous allons le démontrer comme application du lemme
\ref{sites-cohomology-lemma-locally-constant}. Notons que
$H^i(K \otimes_\Lambda^\mathbf{L} L)$ est le même que
$H^i(\tau_{\geq i - 1}K \otimes_\Lambda^\mathbf{L} \tau_{\geq i - 1}L)$.
Nous pouvons donc supposer que $K$ et $L$ sont bornés. Par le lemme
\ref{sites-cohomology-lemma-locally-constant}, nous pouvons supposer que $K$ et $L$ sont
représentés par des complexes de faisceaux localement constants de
$\Lambda$-modules de type fini. Nous pouvons alors remplacer ces complexes
par des complexes bornés supérieurement de $\Lambda$-modules libres de type
fini. Dans ce cas, le résultat est clair.
\end{proof}

\begin{lemma}
\label{sites-cohomology-lemma-locally-constant-bounded}
Soit $\mathcal{C}$ un site. Soit $\Lambda$ un anneau noethérien. Soit
$I \subset \Lambda$ un idéal. Soit $K \in D^-(\mathcal{C}, \Lambda)$. Si les
faisceaux de cohomologie de
$K \otimes_\Lambda^\mathbf{L} \underline{\Lambda/I}$ sont des faisceaux
localement constants de $\Lambda/I$-modules de type fini, alors les faisceaux
de cohomologie de
$K \otimes_\Lambda^\mathbf{L} \underline{\Lambda/I^n}$ sont des faisceaux
localement constants de $\Lambda/I^n$-modules de type fini pour tout
$n \geq 1$.
\end{lemma}

\begin{proof}
Rappelons que les faisceaux localement constants de $\Lambda$-modules de type
fini forment une sous-catégorie de Serre faible de tous les
$\underline{\Lambda}$-modules, voir Modules sur les sites, lemme
\ref{sites-modules-lemma-kernel-finite-locally-constant}. Ainsi, la
sous-catégorie de $D(\mathcal{C}, \Lambda)$ constituée des complexes dont les
faisceaux de cohomologie sont des faisceaux localement constants de
$\Lambda$-modules de type fini forme une sous-catégorie triangulée strictement
pleine et saturée de $D(\mathcal{C}, \Lambda)$, voir Catégories dérivées,
lemme \ref{derived-lemma-cohomology-in-serre-subcategory}. Considérons ensuite
les triangles distingués
$$
K \otimes_\Lambda^\mathbf{L} \underline{I^n/I^{n + 1}} \to
K \otimes_\Lambda^\mathbf{L} \underline{\Lambda/I^{n + 1}} \to
K \otimes_\Lambda^\mathbf{L} \underline{\Lambda/I^n} \to
K \otimes_\Lambda^\mathbf{L} \underline{I^n/I^{n + 1}}[1]
$$
et les isomorphismes
$$
K \otimes_\Lambda^\mathbf{L} \underline{I^n/I^{n + 1}}
=
\left(K \otimes_\Lambda^\mathbf{L} \underline{\Lambda/I}\right)
\otimes_{\Lambda/I}^\mathbf{L} \underline{I^n/I^{n + 1}}
$$
En les combinant avec le lemme \ref{sites-cohomology-lemma-locally-constant-tensor-product},
nous obtenons le résultat.
\end{proof}










% OK, commencer ici.
%

\chapter{Algèbre différentielle graduée}



\phantomsection
\label{dga-section-phantom}


\section{Introduction}
\label{dga-section-introduction}

\noindent
Dans ce chapitre, nous traitons des algèbres différentielles graduées, des modules,
des catégories, etc. Une référence de base est \cite{Keller-Deriving}.
Un article de synthèse est \cite{Keller-survey}.

\medskip\noindent
Comme la longueur des développements ne nous préoccupe pas dans le projet Champs,
nous développons d'abord la matière dans le cadre des catégories de modules
différentiels gradués. Nous reprenons ensuite les constructions dans le cadre des
modules différentiels gradués sur des catégories différentielles graduées.



\section{Conventions}
\label{dga-section-conventions}

\noindent
Dans ce chapitre, nous conservons la convention selon laquelle {\it anneau} signifie
anneau commutatif avec unité $1$. Si $R$ est un anneau, une {\it $R$-algèbre $A$}
sera un $R$-module $A$ muni d'une application $R$-bilinéaire $A \times A \to A$
(la multiplication), telle que celle-ci soit associative et possède une unité.
Autrement dit, ce sont des $R$-algèbres associatives avec unité
dont le morphisme structural $R \to A$ prend ses valeurs dans le centre de $A$.

\medskip\noindent
{\bf Règles de signe.} Dans ce chapitre, nous travaillerons avec des algèbres graduées
et des modules gradués, souvent munis de différentielles. Les règles de signe des
complexes sous-jacents seront toujours celles introduites dans
Compléments d'algèbre, section \ref{more-algebra-section-sign-rules} (ou des règles compatibles avec celles-ci).
Il en résultera parfois que la structure multiplicative se trouve tordue d'une manière
inattendue, surtout lorsque l'on considère des modules à gauche
ou le lien entre modules à gauche et modules à droite.





\section{Algèbres différentielles graduées}
\label{dga-section-dga}


\noindent
Nous donnons seulement les définitions.

\begin{definition}
\label{dga-definition-dga}
Soit $R$ un anneau commutatif. Une {\it algèbre différentielle graduée sur $R$}
est l'une ou l'autre des données suivantes :
\begin{enumerate}
\item un complexe de chaînes $A_\bullet$ de $R$-modules muni d'applications
$R$-bilinéaires $A_n \times A_m \to A_{n + m}$,
$(a, b) \mapsto ab$ telles que
$$
\text{d}_{n + m}(ab) = \text{d}_n(a)b + (-1)^n a\text{d}_m(b)
$$
et tel que $\bigoplus A_n$ soit une $R$-algèbre associative
avec unité, ou
\item un complexe de cochaînes $A^\bullet$ de $R$-modules muni d'applications
$R$-bilinéaires $A^n \times A^m \to A^{n + m}$, $(a, b) \mapsto ab$
telles que
$$
\text{d}^{n + m}(ab) = \text{d}^n(a)b + (-1)^n a\text{d}^m(b)
$$
et tel que $\bigoplus A^n$ soit une $R$-algèbre associative avec unité.
\end{enumerate}
\end{definition}

\noindent
Nous écrivons souvent simplement $A = \bigoplus A_n$ ou $A = \bigoplus A^n$ et
voyons cette somme comme une $R$-algèbre associative avec unité, munie d'une
graduation par $\mathbf{Z}$ et d'un opérateur $R$-linéaire $\text{d}$ dont le carré
est nul et qui satisfait la règle de Leibniz décrite ci-dessus. Dans ce cas,
nous disons souvent : « Soit $(A, \text{d})$ une algèbre différentielle graduée. »

\medskip\noindent
La règle de Leibniz reliant les différentielles et la multiplication dans une
$R$-algèbre différentielle graduée $A$ signifie exactement que l'application de multiplication
définit un morphisme de complexes de cochaînes
$$
\text{Tot}(A^\bullet \otimes_R A^\bullet) \to A^\bullet
$$
Ici, $A^\bullet$ désigne le complexe de cochaînes sous-jacent à $A$.

\begin{definition}
\label{dga-definition-homomorphism-dga}
Un {\it homomorphisme d'algèbres différentielles graduées}
$f : (A, \text{d}) \to (B, \text{d})$ est un homomorphisme d'algèbres $f : A \to B$
compatible avec les graduations et avec $\text{d}$.
\end{definition}

\begin{definition}
\label{dga-definition-cdga}
Une algèbre différentielle graduée $(A, \text{d})$ est {\it commutative} si
$ab = (-1)^{nm}ba$ pour $a$ de degré $n$ et $b$ de degré $m$.
Nous disons que $A$ est {\it strictement commutative} si, en outre, $a^2 = 0$
lorsque $\deg(a)$ est impair.
\end{definition}

\noindent
La définition suivante a un sens en général, mais n'est peut-être
« correcte » que pour le produit tensoriel d'algèbres différentielles graduées
commutatives.

\begin{definition}
\label{dga-definition-tensor-product}
Soit $R$ un anneau.
Soient $(A, \text{d})$, $(B, \text{d})$ des algèbres différentielles graduées sur $R$.
Le {\it produit tensoriel} de $A$ et $B$ comme algèbres différentielles graduées
est l'algèbre $A \otimes_R B$ dont la multiplication est définie par
$$
(a \otimes b)(a' \otimes b') = (-1)^{\deg(a')\deg(b)} aa' \otimes bb'
$$
munie de la différentielle $\text{d}$ définie par la règle
$\text{d}(a \otimes b) = \text{d}(a) \otimes b + (-1)^m a \otimes \text{d}(b)$
où $m = \deg(a)$.
\end{definition}

\begin{lemma}
\label{dga-lemma-total-complex-tensor-product}
Soit $R$ un anneau.
Soient $(A, \text{d})$, $(B, \text{d})$ des algèbres différentielles graduées sur $R$.
Désignons par $A^\bullet$, $B^\bullet$ les complexes de cochaînes sous-jacents.
En tant que complexes de cochaînes de $R$-modules, nous avons
$$
(A \otimes_R B)^\bullet = \text{Tot}(A^\bullet \otimes_R B^\bullet).
$$
\end{lemma}

\begin{proof}
Rappelons que la différentielle du complexe total est donnée par
$\text{d}_1^{p, q} + (-1)^p \text{d}_2^{p, q}$ sur $A^p \otimes_R B^q$.
Or c'est exactement la règle qui définit la différentielle
sur $A \otimes_R B$ dans la
définition \ref{dga-definition-tensor-product}.
\end{proof}






\section{Modules différentiels gradués}
\label{dga-section-modules}

\noindent
Par défaut dans ce chapitre, nous considérons les modules à
droite ; nous discutons des modules à gauche dans la section \ref{dga-section-left-modules}.

\begin{definition}
\label{dga-definition-dgm}
Soit $R$ un anneau.
Soit $(A, \text{d})$ une algèbre différentielle graduée sur $R$.
Un {\it module différentiel gradué} $M$ à droite sur $A$ est un $A$-module à droite
$M$ qui possède une graduation $M = \bigoplus M^n$ et une différentielle
$\text{d}$ tels que $M^n A^m \subset M^{n + m}$, tels que
$\text{d}(M^n) \subset M^{n + 1}$, et tels que
$$
\text{d}(ma) = \text{d}(m)a + (-1)^n m\text{d}(a)
$$
pour $a \in A$ et $m \in M^n$. Un
{\it homomorphisme de modules différentiels gradués} $f : M \to N$ est une
application de $A$-modules compatible avec les graduations et les
différentielles. La catégorie des $A$-modules différentiels gradués (droits) est
notée $\text{Mod}_{(A, \text{d})}$.
\end{definition}

\noindent
Notons que nous pouvons considérer $M$ comme un complexe de cochaînes
$M^\bullet$ de $R$-modules (droits). À savoir, pour $r \in R$ nous avons $\text{d}(r)
= 0$ et $r$ s'envoie sur un élément de degré $0$ de $A$,
donc $\text{d}(mr) = \text{d}(m)r$.

\medskip\noindent
La règle de Leibniz reliant les différentielles et la multiplication sur un $R$-module
différentiel gradué $M$ sur une $R$-algèbre différentielle graduée $A$ signifie
exactement que l'application de multiplication définit un morphisme de complexes de cochaînes
$$
\text{Tot}(M^\bullet \otimes_R A^\bullet) \to M^\bullet
$$
Ici $A^\bullet$ et $M^\bullet$ désignent les complexes de cochaînes
sous-jacents de $A$ et $M$.

\begin{lemma}
\label{dga-lemma-dgm-abelian}
Soit $(A, d)$ une algèbre différentielle graduée. La catégorie
$\text{Mod}_{(A, \text{d})}$ est abélienne et possède des limites et colimites arbitraires.
\end{lemma}

\begin{proof}
Les noyaux et conoyaux commutent avec la prise des $A$-modules
sous-jacents. Il en va de même pour les sommes directes et les colimites. En d'autres
termes, ces opérations dans $\text{Mod}_{(A, \text{d})}$ commutent avec le foncteur
d'oubli vers la catégorie des $A$-modules. Ce n'est pas le cas pour les
produits et limites. À savoir, si $N_i$, $i \in I$ est
une famille de $A$-modules différentiels gradués, alors le produit
$\prod N_i$ dans $\text{Mod}_{(A, \text{d})}$ est donné en posant $(\prod N_i)^n = \prod
N_i^n$ et $\prod N_i = \bigoplus_n (\prod N_i)^n$. Ainsi, nous voyons que le produit
commute avec le foncteur d'oubli vers la catégorie des $A$-modules gradués. Une
catégorie avec produits et égalisateurs possède des limites, voir
Catégories, lemme \ref{categories-lemma-limits-products-equalizers}.
\end{proof}

\noindent
Ainsi, si $(A, \text{d})$ est une algèbre
différentielle graduée sur $R$, alors il existe un foncteur exact
$$
\text{Mod}_{(A, \text{d})} \longrightarrow \text{Comp}(R)
$$
de catégories abéliennes. Pour un module différentiel gradué $M$,
les groupes de cohomologie $H^n(M)$ sont définis comme la cohomologie
du complexe correspondant de $R$-modules. Par conséquent, une suite
exacte courte $0 \to K \to L \to M \to 0$ de modules différentiels gradués
donne lieu à une longue suite exacte
\begin{equation}
\label{dga-equation-les}
H^n(K) \to H^n(L) \to H^n(M) \to H^{n + 1}(K)
\end{equation}
de modules de cohomologie,
voir Homologie, lemme \ref{homology-lemma-long-exact-sequence-cochain}.

\medskip\noindent
De plus, à partir de maintenant, nous empruntons toute la
terminologie utilisée pour les complexes de modules. Par exemple, nous disons
qu'un $A$-module différentiel gradué $M$ est {\it acyclique} si
$H^k(M) = 0$ pour tout $k \in \mathbf{Z}$. Nous disons qu'un homomorphisme
$M \to N$ de modules différentiels gradués sur $A$ est un {\it
quasi-isomorphisme} s'il induit des isomorphismes $H^k(M) \to H^k(N)$ pour tout $k \in
\mathbf{Z}$. Et ainsi de suite.

\begin{definition}
\label{dga-definition-shift}
Soit $(A, \text{d})$ une algèbre différentielle graduée.
Soit $M$ un module différentiel gradué dont le complexe sous-jacent de
$R$-modules est $M^\bullet$. Pour tout $k \in \mathbf{Z}$, nous
définissons le {\it module décalé de $k$} $M[k]$ comme suit
\begin{enumerate}
\item le complexe de $R$-modules sous-jacent à $M[k]$ est $M^\bullet[k]$,
c'est-à-dire que nous avons $M[k]^n = M^{n + k}$
et $\text{d}_{M[k]} = (-1)^k\text{d}_M$ et
\item la structure de $A$-module est donnée par la multiplication
$$
(M[k])^n \times A^m \longrightarrow (M[k])^{n + m}
$$
est égale à la multiplication donnée $M^{n + k} \times A^m \to M^{n + k + m}$.
\end{enumerate}
Pour un morphisme $f : M \to N$ de $A$-modules différentiels gradués,
nous notons $f[k] : M[k] \to N[k]$ l'application égale à $f$ sur les
$A$-modules sous-jacents. Cela définit un
foncteur $[k] : \text{Mod}_{(A, \text{d})} \to \text{Mod}_{(A, \text{d})}$.
\end{definition}

\noindent
Vérifions qu'avec ce choix la règle de Leibniz est satisfaite.
Soient $x \in M[k]^n = M^{n + k}$ et $a \in A^m$. En notant
$\cdot_{M[k]}$ le produit dans $M[k]$, on obtient
\begin{align*}
\text{d}_{M[k]}(x \cdot_{M[k]} a)
& =
(-1)^k \text{d}_M(xa) \\
& =
(-1)^k \text{d}_M(x) a + (-1)^{k + n + k} x \text{d}(a) \\ & =
\text{d}_{M[k]}(x)
a + (-1)^n x \text{d}(a) \\ & =
\text{d}_{M[k]}(x)
\cdot_{M[k]} a + (-1)^n x \cdot_{M[k]} \text{d}(a)
\end{align*}
C'est ce que nous voulons puisque $x$ a un degré $n$ en tant qu'élément homogène de $M[k]$. Nous
observons également qu'avec ces choix, nous pouvons considérer
l'application de multiplication comme un morphisme de complexes
$$
\text{Tot}(M^\bullet[k] \otimes _R A^\bullet) \to
\text{Tot}(M^\bullet \otimes _R A^\bullet)[k] \to
M^\bullet[k]
$$
où la première flèche est
celle de Compléments d'algèbre, section
\ref{more-algebra-section-sign-rules} (\ref{more-algebra-item-shift-tensor}) qui dans ce
cas n'implique pas de signe. (En fait, nous aurions pu déduire que la
règle de Leibniz vaut à partir de cette observation.)

\medskip\noindent
Les remarques dans Homologie, section \ref{homology-section-homotopy-shift}
s'appliquent. En particulier, nous identifierons les groupes de cohomologie de tous les
décalages $M[k]$ sans l'intervention de signes.

\medskip\noindent
À ce stade, nous avons suffisamment de structure pour parler des {\it triangles},
voir Catégories dérivées, définition \ref{derived-definition-triangle}. En fait,
notre prochain objectif est de développer suffisamment de théorie pour pouvoir
énoncer et prouver que la catégorie homotopique des modules différentiels gradués
est une catégorie triangulée. Tout d'abord, nous définissons la catégorie homotopique.






\section{La catégorie homotopique}
\label{dga-section-homotopy}

\noindent
Nos homotopies tiennent compte de la structure de module $A$ et de la
graduation, mais pas de la différentielle (bien sûr).

\begin{definition}
\label{dga-definition-homotopy}
Soit $(A, \text{d})$ une algèbre différentielle graduée. Soient
$f, g : M \to N$ des homomorphismes de $A$-modules différentiels gradués. Une
{\it homotopie entre $f$ et $g$} est une application de $A$-modules $h : M \to N$
telle que
\begin{enumerate}
\item $h(M^n) \subset N^{n - 1}$ pour tout $n$, et
\item $f(x) - g(x) = \text{d}_N(h(x)) + h(\text{d}_M(x))$ pour
tout $x \in M$.
\end{enumerate}
Si une homotopie existe, alors nous disons que $f$ et $g$ sont {\it homotopes}.
\end{definition}

\noindent
Ainsi, $h$ est compatible avec la structure de $A$-module et la graduation
mais pas avec la différentielle. Si $f = g$ et si $h$ est une homotopie
au sens de la définition, alors $h$ définit un morphisme $h : M \to N[-1]$
dans $\text{Mod}_{(A, \text{d})}$.

\begin{lemma}
\label{dga-lemma-compose-homotopy}
Soit $(A, \text{d})$ une algèbre différentielle graduée.
Soient $f, g : L \to M$ des homomorphismes de $A$-modules différentiels gradués.
Supposons donnés en outre des homomorphismes $a : K \to L$ et $c : M \to N$.
Si $h : L \to M$ est une application de $A$-modules qui définit une homotopie
entre $f$ et $g$, alors $c \circ h \circ a$ définit une homotopie
entre $c \circ f \circ a$ et $c \circ g \circ a$.
\end{lemma}

\begin{proof}
Résulte immédiatement d'Homologie, lemme \ref{homology-lemma-compose-homotopy-cochain}.
\end{proof}

\noindent
Ce lemme nous permet de définir la catégorie homotopique comme suit.

\begin{definition}
\label{dga-definition-complexes-notation}
Soit $(A, \text{d})$ une algèbre différentielle graduée.
La {\it catégorie homotopique}, notée $K(\text{Mod}_{(A, \text{d})})$, est la
catégorie dont les objets sont les objets de
$\text{Mod}_{(A, \text{d})}$ et dont les morphismes sont les classes d'homotopie
des homomorphismes de $A$-modules différentiels gradués.
\end{definition}

\noindent
La notation $K(\text{Mod}_{(A, \text{d})})$ n'est pas standard mais elle est au
moins cohérente avec l'utilisation de $K(-)$ dans d'autres parties du projet Champs.

\begin{lemma}
\label{dga-lemma-homotopy-direct-sums}
Soit $(A, \text{d})$ une algèbre différentielle graduée. La
catégorie homotopique $K(\text{Mod}_{(A, \text{d})})$ possède des
sommes directes et des produits.
\end{lemma}

\begin{proof}
Démonstration omise. Indication : utiliser les sommes directes et les
produits comme dans le lemme \ref{dga-lemma-dgm-abelian}. Cela fonctionne
parce que nous avons vu que ces foncteurs commutent avec le foncteur d'oubli vers la
catégorie des $A$-modules gradués et parce que $\prod$ est un foncteur
exact sur la catégorie des familles de groupes abéliens.
\end{proof}







\section{Cônes}
\label{dga-section-cones}

\noindent
Nous introduisons les cônes pour la catégorie des modules différentiels gradués.

\begin{definition}
\label{dga-definition-cone}
Soit $(A, \text{d})$ une algèbre différentielle graduée.
Soit $f : K \to L$ un homomorphisme de $A$-modules différentiels gradués. Le
{\it cône} de $f$ est le $A$-module différentiel gradué $C(f)$
donné par $C(f) = L \oplus K$ avec la graduation
$C(f)^n = L^n \oplus K^{n + 1}$ et la
différentielle
$$
d_{C(f)} =
\left(
\begin{matrix}
\text{d}_L & f \\
0 & -\text{d}_K
\end{matrix}
\right)
$$
Il est muni de morphismes canoniques de complexes $i : L \to C(f)$ et $p :
C(f) \to K[1]$ induits par les applications évidentes $L \to C(f)$
et $C(f) \to K$.
\end{definition}

\noindent
La formation du triangle du cône est fonctorielle dans le sens suivant.

\begin{lemma}
\label{dga-lemma-functorial-cone}
Soit $(A, \text{d})$ une algèbre différentielle graduée.
Supposons que
$$
\xymatrix{
K_1 \ar[r]_{f_1} \ar[d]_a & L_1 \ar[d]^b \\
K_2 \ar[r]^{f_2} & L_2
}
$$
soit un diagramme d'homomorphismes de $A$-modules différentiels gradués
qui est commutatif à homotopie
près. Alors il existe un morphisme $c : C(f_1) \to C(f_2)$ qui induit un
morphisme de triangles
$$
(a, b, c) : (K_1, L_1, C(f_1), f_1, i_1, p_1) \to
(K_1, L_1, C(f_1), f_2, i_2, p_2)
$$
dans $K(\text{Mod}_{(A, \text{d})})$.
\end{lemma}

\begin{proof}
Soit $h : K_1 \to L_2$ une homotopie entre $f_2 \circ a$ et $b \circ f_1$.
Définissons $c$ par la matrice
$$
c =
\left(
\begin{matrix}
b & h \\
0 & a
\end{matrix}
\right) :
L_1 \oplus K_1 \to L_2 \oplus K_2
$$
Un calcul matriciel montre que $c$ est un morphisme de modules
différentiels gradués. Il est trivial que $c \circ i_1 = i_2 \circ b$, et il est
également trivial de vérifier que $p_2 \circ c = a \circ p_1$.
\end{proof}











\section{Suites exactes courtes admissibles}
\label{dga-section-admissible}

\noindent
Une suite exacte courte admissible est l'analogue des suites exactes scindées
terme à terme dans le cadre des modules différentiels gradués.

\begin{definition}
\label{dga-definition-admissible-ses}
Soit $(A, \text{d})$ une algèbre différentielle graduée.
\begin{enumerate}
\item Un homomorphisme $K \to L$ de modules différentiels gradués sur $A$
est un {\it monomorphisme admissible} s'il existe un morphisme de $A$-modules
gradués $L \to K$ qui soit inverse à gauche de $K \to L$.
\item Un homomorphisme $L \to M$ de $A$-modules différentiels gradués
est un {\it épimorphisme admissible} s'il existe un morphisme de $A$-modules
gradués $M \to L$ qui soit un inverse à droite de $L \to M$.
\item Une suite exacte courte $0 \to K \to L \to M \to 0$ de $A$-modules
différentiels gradués est une {\it suite exacte courte admissible} si
elle se scinde en tant que suite de $A$-modules gradués.
\end{enumerate}
\end{definition}

\noindent
Ainsi, les scindages sont compatibles avec toutes les données sauf les
différentielles. Étant donnée une suite exacte courte admissible, nous
obtenons un triangle ; c'est la raison pour laquelle nous exigeons que nos
scindages soient compatibles avec la structure de $A$-module.

\begin{lemma}
\label{dga-lemma-admissible-ses}
Soit $(A, \text{d})$ une algèbre différentielle graduée.
Soit $0 \to K \to L \to M \to 0$ une suite exacte courte admissible de
$A$-modules différentiels gradués. Soient $s : M \to L$ et $\pi : L \to K$ des
scindages tels que $\Ker(\pi) = \Im(s)$. Alors
nous obtenons un morphisme
$$
\delta = \pi \circ \text{d}_L \circ s : M \to K[1]
$$
de $\text{Mod}_{(A, \text{d})}$ qui induit les homomorphismes de
liaison dans la suite exacte longue de cohomologie (\ref{dga-equation-les}).
\end{lemma}

\begin{proof}
L'application $\pi \circ \text{d}_L \circ s$ est compatible par construction
avec la structure de $A$-module et les graduations. Elle est compatible
avec les différentielles d'après
Homologie, lemme
\ref{homology-lemma-ses-termwise-split-cochain}. Soit $R$ l'anneau sur lequel $A$ est une algèbre différentielle
graduée. L'égalité des applications est une assertion sur les $R$-modules.
Elle découle donc d'Homologie,
lemmes \ref{homology-lemma-ses-termwise-split-cochain} et
\ref{homology-lemma-ses-termwise-split-long-cochain}.
\end{proof}

\begin{lemma}
\label{dga-lemma-make-commute-map}
Soit $(A, \text{d})$ une algèbre différentielle graduée. Soit
$$
\xymatrix{
K \ar[r]_f \ar[d]_a & L \ar[d]^b \\ M
\ar[r]^g & N
}
$$
un diagramme d'homomorphismes de $A$-modules différentiels gradués
qui commute à homotopie près.
\begin{enumerate}
\item Si $f$ est un monomorphisme admissible, alors $b$ est homotope à un
homomorphisme qui rend le diagramme commutatif.
\item Si $g$ est un épimorphisme admissible, alors $a$ est homotope à
un morphisme qui fait commuter le diagramme.
\end{enumerate}
\end{lemma}

\begin{proof}
Soit $h : K \to N$ une homotopie entre $bf$ et $ga$, c'est-à-dire $bf
- ga = \text{d}h + h\text{d}$. Supposons que $\pi : L \to K$ soit une
application de $A$-modules gradués, inverse à gauche de $f$.
Prenons $b' = b - \text{d}h\pi - h\pi \text{d}$.
Supposons que $s : N \to M$ soit une application de $A$-modules gradués, inverse à
droite de $g$. Prenons $a' = a + \text{d}sh +
sh\text{d}$. Calculs omis.
\end{proof}

\begin{lemma}
\label{dga-lemma-make-injective}
Soit $(A, \text{d})$ une algèbre différentielle
graduée. Soit $\alpha : K \to L$ un homomorphisme de $A$-modules
différentiels gradués. Il existe une factorisation
$$
\xymatrix{
K \ar[r]^{\tilde \alpha} \ar@/_1pc/[rr]_\alpha &
\tilde L \ar[r]^\pi & L
}
$$
dans $\text{Mod}_{(A, \text{d})}$ telle que
\begin{enumerate}
\item $\tilde \alpha$ soit un monomorphisme admissible (voir
la définition \ref{dga-definition-admissible-ses}),
\item il existe un morphisme $s : L \to \tilde
L$ tel que $\pi \circ s = \text{id}_L$ et tel que
$s \circ \pi$ soit homotope à $\text{id}_{\tilde L}$.
\end{enumerate}
\end{lemma}

\begin{proof}
La démonstration est identique à celle de
Catégories dérivées, lemme \ref{derived-lemma-make-injective}. En effet,
nous posons $\tilde L = L \oplus C(1_K)$ et nous utilisons des propriétés
élémentaires de la construction du cône.
\end{proof}

\begin{lemma}
\label{dga-lemma-sequence-maps-split}
Soit $(A, \text{d})$ une algèbre différentielle
graduée. Soit $L_1 \to L_2 \to \ldots \to
L_n$ une suite d'homomorphismes composables de
$A$-modules différentiels gradués.
Il existe un diagramme commutatif
$$
\xymatrix{
L_1 \ar[r] &
L_2 \ar[r] &
\ldots \ar[r] &
L_n \\ M_1
\ar[r] \ar[u] & M_2
\ar[r] \ar[u] & \ldots
\ar[r] & M_n
\ar[u]
}
$$
dans $\text{Mod}_{(A, \text{d})}$ tel que chaque $M_i \to M_{i + 1}$ soit
un monomorphisme admissible et que chaque $M_i \to L_i$ soit
une équivalence d'homotopie.
\end{lemma}

\begin{proof}
Le cas $n = 1$ est immédiat. Le lemme
\ref{dga-lemma-make-injective} traite le cas $n = 2$.
Supposons le diagramme construit à l'exception
de $M_n$. Appliquons le lemme \ref{dga-lemma-make-injective} à la
composition $M_{n - 1} \to L_{n - 1} \to L_n$. Le
résultat est une factorisation $M_{n - 1} \to M_n \to L_n$ comme
souhaité.
\end{proof}



\begin{lemma}
\label{dga-lemma-nilpotent}
Soit $(A, \text{d})$ une algèbre différentielle graduée.
Soient $0 \to K_i \to L_i \to M_i \to 0$, $i = 1, 2, 3$
des suites exactes courtes admissibles de $A$-modules différentiels gradués.
Soient $b : L_1 \to L_2$ et $b' : L_2 \to L_3$
des homomorphismes de modules différentiels gradués tels que
$$
\vcenter{
\xymatrix{
K_1 \ar[d]_0 \ar[r] &
L_1 \ar[r] \ar[d]_b &
M_1 \ar[d]_0 \\ K_2
\ar[r] & L_2 \ar[r] & M_2 } }
\quad\text{et}\quad
\vcenter{
\xymatrix{
K_2
\ar[d]^0
\ar[r] & L_2 \ar[r]
\ar[d]^{b'} & M_2 \ar[d]^0
\\ K_3 \ar[r] &
L_3 \ar[r] & M_3
}
}
$$
commutent à homotopie près. Alors $b' \circ b$ est homotope à $0$.
\end{lemma}

\begin{proof}
Par le lemme \ref{dga-lemma-make-commute-map}, nous pouvons remplacer $b$ et $b'$
par des applications homotopes telles que le carré droit du diagramme de gauche
commute et le carré gauche du diagramme de droite commute. En d'autres termes, nous avons
$\Im(b) \subset \Im(K_2 \to L_2)$ et
$\Ker(b') \supset \Im(K_2 \to L_2)$. Ensuite, $b'
\circ b = 0$ en tant qu'homomorphisme de modules.
\end{proof}
















\section{Triangles distingués}
\label{dga-section-distinguished}

\noindent
Le lemme suivant produit nos triangles distingués.

\begin{lemma}
\label{dga-lemma-triangle-independent-splittings}
Soit $(A, \text{d})$ une algèbre différentielle graduée.
Soit $0 \to K \to L \to M \to 0$ une suite exacte courte admissible
de $A$-modules différentiels gradués. Le triangle
\begin{equation}
\label{dga-equation-triangle-associated-to-admissible-ses}
K \to L \to M \xrightarrow{\delta} K[1]
\end{equation}
avec $\delta$ comme dans le lemme \ref{dga-lemma-admissible-ses} est, à
isomorphisme canonique près dans $K(\text{Mod}_{(A, \text{d})})$, indépendant des choix
effectués dans le lemme \ref{dga-lemma-admissible-ses}.
\end{lemma}

\begin{proof}
En effet, soit $(s', \pi')$ un second choix de scindages comme
dans le lemme \ref{dga-lemma-admissible-ses}. Alors nous affirmons que $\delta$ et
$\delta'$ sont homotopes. Écrivons $s' = s + \alpha \circ h$ et
$\pi' = \pi + g \circ \beta$ pour certains homomorphismes uniques
de $A$-modules $h : M \to K$ et $g : M \to K$ de degré $-1$. Alors
$g = -h$ et $g$ est une homotopie entre $\delta$ et $\delta'$. Les
calculs sont effectués dans la démonstration
d'Homologie, lemme \ref{homology-lemma-ses-termwise-split-homotopy-cochain}.
\end{proof}

\begin{definition}
\label{dga-definition-distinguished-triangle}
Soit $(A, \text{d})$ une algèbre différentielle graduée.
\begin{enumerate}
\item Si $0 \to K \to L \to M \to 0$ est une suite exacte courte
admissible de $A$-modules différentiels gradués, alors le {\it triangle associé
à $0 \to K \to L \to M \to 0$} est le triangle
(\ref{dga-equation-triangle-associated-to-admissible-ses}) de
$K(\text{Mod}_{(A, \text{d})})$.
\item Un triangle de $K(\text{Mod}_{(A, \text{d})})$ est
appelé {\it triangle distingué} s'il est isomorphe à un triangle
associé à une suite exacte courte admissible de
$A$-modules différentiels gradués.
\end{enumerate}
\end{definition}









\section{Cônes et triangles distingués}
\label{dga-section-cones-and-triangles}

\noindent
Soit $(A, \text{d})$ une algèbre différentielle graduée.
Soit $f : K \to L$ un homomorphisme de $A$-modules différentiels gradués.
Alors $(K, L, C(f), f, i, p)$ forme un triangle :
$$
K \to L \to C(f) \to K[1]
$$
dans $\text{Mod}_{(A, \text{d})}$ et donc dans $K(\text{Mod}_{(A, \text{d})})$.
Les cônes ne sont {\bf pas} des triangles distingués en général, mais la
différence tient à un signe ou à une rotation, au choix. Voici deux énoncés précis.

\begin{lemma}
\label{dga-lemma-rotate-cone}
Soit $(A, \text{d})$ une algèbre différentielle
graduée. Soit $f : K \to L$ un homomorphisme de modules différentiels
gradués. Le triangle $(L, C(f), K[1], i, p,
f[1])$ est le triangle associé à la suite exacte courte admissible
$$
0 \to L \to C(f) \to K[1] \to 0
$$
provenant de la définition du cône de $f$.
\end{lemma}

\begin{proof}
Immédiat d'après les définitions.
\end{proof}

\begin{lemma}
\label{dga-lemma-rotate-triangle}
Soit $(A, \text{d})$ une algèbre différentielle graduée.
Soient $\alpha : K \to L$ et $\beta : L \to M$
définissant une suite exacte courte admissible
$$
0 \to K \to L \to M \to 0
$$
de modules différentiels gradués
sur $A$. Soit $(K, L, M, \alpha, \beta,
\delta)$ le triangle associé. Alors les triangles
$$
(M[-1], K, L, \delta[-1], \alpha, \beta)
\quad\text{et}\quad
(M[-1], K, C(\delta[-1]), \delta[-1], i, p)
$$
sont isomorphes.
\end{lemma}

\begin{proof}
En utilisant un choix de scindages, nous écrivons $L = K \oplus M$ et nous
identifions $\alpha$ et $\beta$ avec les applications d'inclusion et de projection
naturelles. Par construction de $\delta$, nous avons
$$
d_B =
\left(
\begin{matrix}
d_K & \delta \\
0 & d_M
\end{matrix}
\right)
$$
D'autre part, le cône de $\delta[-1] : M[-1] \to K$ est
donné par $C(\delta[-1]) = K \oplus M$ avec une différentielle identique
à la matrice ci-dessus ! D'où le lemme.
\end{proof}

\begin{lemma}
\label{dga-lemma-third-isomorphism}
Soit $(A, \text{d})$ une algèbre différentielle graduée.
Soient $f_1 : K_1 \to L_1$ et $f_2 : K_2 \to L_2$ des homomorphismes de
modules différentiels gradués sur $A$. Soit
$$
(a, b, c) :
(K_1, L_1, C(f_1), f_1, i_1, p_1)
\longrightarrow
(K_1, L_1, C(f_1), f_2, i_2, p_2)
$$
un morphisme quelconque de triangles de $K(\text{Mod}_{(A, \text{d})})$. Si $a$ et
$b$ sont des équivalences par homotopie, alors il en est de même pour $c$.
\end{lemma}

\begin{proof}
Soit $a^{-1} : K_2 \to K_1$ un homomorphisme de $A$-modules différentiels gradués qui
est l'inverse de $a$ dans $K(\text{Mod}_{(A, \text{d})})$. Soit
$b^{-1} : L_2 \to L_1$ un homomorphisme de $A$-modules différentiels gradués qui est
l'inverse de $b$ dans $K(\text{Mod}_{(A, \text{d})})$. Soit $c' :
C(f_2) \to C(f_1)$ le morphisme issu du lemme
\ref{dga-lemma-functorial-cone} appliqué à $f_1 \circ a^{-1} =
b^{-1} \circ f_2$. Si nous pouvons montrer
que $c \circ c'$ et $c' \circ c$ sont des isomorphismes dans
$K(\text{Mod}_{(A, \text{d})})$, la démonstration
sera achevée. Il suffit donc de prouver ce qui suit : étant donné un
morphisme de triangles
$(1, 1, c) : (K, L, C(f), f, i, p)$ dans
$K(\text{Mod}_{(A, \text{d})})$, le morphisme $c$ est un isomorphisme dans
$K(\text{Mod}_{(A, \text{d})})$. Par
hypothèse, les deux carrés du diagramme
$$
\xymatrix{
L \ar[r] \ar[d]_1 &
C(f) \ar[r] \ar[d]_c &
K[1] \ar[d]_1 \\ L
\ar[r] & C(f)
\ar[r] &
K[1]
}
$$
commutent à homotopie près. Par construction de $C(f)$, les lignes
forment des suites exactes courtes admissibles. Ainsi, nous voyons
que $(c - 1)^2 = 0$ dans $K(\text{Mod}_{(A, \text{d})})$ par
le lemme \ref{dga-lemma-nilpotent}. Par
conséquent, $c$ est un isomorphisme dans $K(\text{Mod}_{(A, \text{d})})$
avec inverse $2 - c$.
\end{proof}

\noindent
Le lemme suivant montre que les triangles de la catégorie homotopique
donnés par les cônes et les triangles distingués coïncident à
isomorphisme près, du moins à un signe près !

\begin{lemma}
\label{dga-lemma-the-same-up-to-isomorphisms}
Soit $(A, \text{d})$ une algèbre différentielle graduée.
\begin{enumerate}
\item Étant donnée une suite exacte courte
admissible $0 \to K \xrightarrow{\alpha} L \to M \to 0$
de $A$-modules différentiels gradués, il existe une équivalence d'homotopie
$C(\alpha) \to M$ telle que le diagramme
$$
\xymatrix{
K \ar[r] \ar[d] & L \ar[d] \ar[r] &
C(\alpha) \ar[r]_{-p} \ar[d] & K[1] \ar[d] \\ K
\ar[r]^\alpha & L \ar[r]^\beta & M
\ar[r]^\delta & K[1]
}
$$
définisse un isomorphisme de triangles dans $K(\text{Mod}_{(A, \text{d})})$.
\item Étant donné un morphisme de complexes $f : K
\to L$, il existe un isomorphisme de triangles
$$
\xymatrix{
K \ar[r] \ar[d] & \tilde L \ar[d] \ar[r] &
M \ar[r]_{\delta} \ar[d] & K[1] \ar[d] \\ K
\ar[r] & L \ar[r] &
C(f) \ar[r]^{-p} & K[1]
}
$$
où le triangle supérieur est le triangle associé à une
suite exacte courte admissible $K \to \tilde L \to M$.
\end{enumerate}
\end{lemma}

\begin{proof}
Démonstration de (1). Nous avons $C(\alpha) = L \oplus K$ et nous
définissons simplement $C(\alpha) \to M$ via la projection sur $L$ suivie de
$\beta$. Cela définit un morphisme de modules différentiels gradués car les
compositions $K^{n + 1} \to L^{n + 1} \to M^{n + 1}$ sont nulles.
Choisissons des scindages $s : M \to L$ et $\pi : L \to K$ avec
$\Ker(\pi) = \Im(s)$ et posons $\delta
= \pi \circ \text{d}_L \circ s$ comme d'habitude. Pour
obtenir un inverse par homotopie,
nous prenons $M \to C(\alpha)$ donné par $(s , -\delta)$. Cela est compatible
avec les différentielles car $\delta^n$ peut être caractérisé comme
l'application unique $M^n \to K^{n + 1}$
telle que $\text{d} \circ s^n - s^{n + 1} \circ \text{d} = \alpha \circ \delta^n$,
voir la
démonstration de Homologie, lemme
\ref{homology-lemma-ses-termwise-split-cochain}. La composition $M \to C(f) \to M$ est l'identité. La
composition $C(f) \to M \to C(f)$ est égale au morphisme
$$
\left(
\begin{matrix}
s \circ \beta & 0 \\
-\delta \circ \beta & 0
\end{matrix}
\right)
$$
Pour voir que cela est homotope à l'application
identité, utilisons l'homotopie $h : C(\alpha) \to C(\alpha)$
donnée par la matrice
$$
\left(
\begin{matrix}
0 & 0 \\
\pi & 0
\end{matrix}
\right) :
C(\alpha) = L \oplus K
\to
L \oplus K = C(\alpha)
$$
Il est trivial de vérifier que
$$
\left(
\begin{matrix}
1 & 0 \\
0 & 1
\end{matrix}
\right)
-
\left(
\begin{matrix}
s \\
-\delta
\end{matrix}
\right)
\left(
\begin{matrix}
\beta & 0
\end{matrix}
\right)
=
\left(
\begin{matrix}
\text{d} & \alpha \\
0 & -\text{d}
\end{matrix}
\right)
\left(
\begin{matrix}
0 & 0 \\
\pi & 0
\end{matrix}
\right)
+
\left(
\begin{matrix}
0 & 0 \\
\pi & 0
\end{matrix}
\right)
\left(
\begin{matrix}
\text{d} & \alpha \\
0 & -\text{d}
\end{matrix}
\right)
$$
Pour terminer la démonstration de (1), nous devons montrer que les
morphismes $-p : C(\alpha) \to K[1]$ (voir
définition \ref{dga-definition-cone}) et
$C(\alpha) \to M \to K[1]$ concordent à homotopie
près. Cela est clair d'après ce qui précède. En effet, nous pouvons utiliser l'inverse
d'homotopie $(s, -\delta) : M \to C(\alpha)$ et
vérifier à la place que les deux applications
$M \to K[1]$ concordent. Notons en outre que $p
\circ (s, -\delta) = -\delta$ comme souhaité.

\medskip\noindent
Démonstration de (2). Nous notons $\tilde f : K \to
\tilde L$, $s : L \to \tilde
L$ et $\pi : L \to L$ comme dans
le lemme \ref{dga-lemma-make-injective}.
D'après les lemmes \ref{dga-lemma-functorial-cone} et \ref{dga-lemma-third-isomorphism},
les triangles $(K, L, C(f), i, p)$ et $(K,
\tilde L, C(\tilde f), \tilde i, \tilde p)$ sont isomorphes.
Notons que nous pouvons composer les isomorphismes de triangles.
Ainsi, nous pouvons remplacer $L$ par
$\tilde L$ et $f$ par $\tilde f$. En d'autres termes, nous
pouvons supposer que $f$ est un monomorphisme admissible. Dans ce
cas, le résultat découle de la partie (1).
\end{proof}







\section{La catégorie homotopique est triangulée}
\label{dga-section-homotopy-triangulated}

\noindent
Montrons d'abord qu'elle est prétriangulée.

\begin{lemma}
\label{dga-lemma-homotopy-category-pre-triangulated}
Soit $(A, \text{d})$ une algèbre différentielle graduée.
La catégorie homotopique $K(\text{Mod}_{(A, \text{d})})$
avec ses foncteurs de translation naturels et ses triangles distingués
est une catégorie pré-triangulée.
\end{lemma}

\begin{proof}
Démonstration de TR1. Par définition, tout triangle isomorphe à un triangle
distingué est distingué. De plus, tout triangle $(K, K, 0, 1, 0,
0)$ est distingué puisque $0 \to K \to K \to 0 \to 0$ est
une suite exacte courte admissible. Enfin, étant donné tout homomorphisme
$f : K \to L$ de $A$-modules différentiels gradués, le triangle
$(K, L, C(f), f, i, -p)$ est distingué d'après le
lemme \ref{dga-lemma-the-same-up-to-isomorphisms}.

\medskip\noindent
Démonstration de TR2. Soit $(X, Y, Z, f, g, h)$ un
triangle. Supposons que $(Y, Z, X[1], g, h, -f[1])$ soit
distingué. Alors il existe une suite exacte courte admissible
$0 \to K \to L \to M \to 0$ telle que le triangle
associé $(K, L, M, \alpha, \beta, \delta)$ soit
isomorphe à $(Y, Z, X[1], g, h, -f[1])$. En effectuant une rotation inverse,
nous voyons que $(X, Y, Z, f, g, h)$ est
isomorphe à $(M[-1], K, L, -\delta[-1], \alpha,
\beta)$. Il découle du lemme \ref{dga-lemma-rotate-triangle} que le triangle
$(M[-1], K, L, \delta[-1], \alpha, \beta)$ est
isomorphe à $(M[-1],
K, C(\delta[-1]), \delta[-1], i, p)$. En
précomposant l'isomorphisme précédent des triangles avec $-1$ sur $Y$, il s'ensuit
que $(X, Y, Z, f, g, h)$ est isomorphe à $(M[-1], K,
C(\delta[-1]), \delta[-1], i, -p)$. Par conséquent, il
est distingué par le lemme
\ref{dga-lemma-the-same-up-to-isomorphisms}. D'autre part,
supposons que $(X, Y, Z, f, g, h)$ soit distingué. Par le lemme
\ref{dga-lemma-the-same-up-to-isomorphisms}, cela signifie qu'il est isomorphe à un
triangle de la forme $(K, L, C(f), f, i, -p)$
pour un certain morphisme $f$ de $\text{Mod}_{(A,
\text{d})}$. Alors, le triangle obtenu par rotation $(Y, Z,
X[1], g, h, -f[1])$ est isomorphe à
$(L, C(f), K[1], i, -p, -f[1])$, qui est isomorphe au
triangle $(L, C(f), K[1], i, p,
f[1])$. Par le lemme
\ref{dga-lemma-rotate-cone}, ce triangle est distingué. Par conséquent, $(Y, Z, X[1], g,
h, -f[1])$ est distingué comme souhaité.

\medskip\noindent
Démonstration de TR3. Soient $(X, Y, Z, f, g, h)$ et $(X', Y', Z', f', g',
h')$ des triangles distingués de $K(\mathcal{A})$ et soient $a : X \to X'$
et $b : Y \to Y'$ des morphismes tels que $f' \circ a = b \circ f$. D'après
le lemme \ref{dga-lemma-the-same-up-to-isomorphisms}, nous pouvons supposer
que $(X, Y, Z, f, g, h) = (X, Y, C(f), f, i, -p)$ et
$(X', Y', Z', f', g', h') = (X', Y', C(f'), f', i', -p')$. À ce
stade, nous appliquons simplement le lemme \ref{dga-lemma-functorial-cone}
au diagramme commutatif donné par $f, f', a, b$.
\end{proof}

\noindent
Avant de prouver TR4 en général, nous le prouvons dans un cas particulier.

\begin{lemma}
\label{dga-lemma-two-split-injections}
Soit $(A, \text{d})$ une algèbre différentielle graduée. Supposons que
$\alpha : K \to L$ et $\beta : L \to M$ soient des monomorphismes admissibles
de modules différentiels gradués sur $A$. Alors il existe des triangles distingués
$(K, L, Q_1, \alpha, p_1, d_1)$, $(K, M, Q_2, \beta \circ \alpha, p_2, d_2)$ et $(L,
M, Q_3, \beta, p_3, d_3)$ pour lesquels TR4 est vérifié.
\end{lemma}

\begin{proof}
Soient $\pi_1 : L \to K$ et $\pi_3 : M \to L$ des homomorphismes de
$A$-modules gradués qui sont des inverses à gauche de $\alpha$ et $\beta$.
Alors $K \to M$ est également un monomorphisme admissible avec un
inverse à gauche $\pi_2 = \pi_1 \circ \pi_3$. Notons $Q_1$, $Q_2$ et
$Q_3$ pour les conoyaux de $K \to L$, $K \to M$ et $L \to M$. Alors nous
obtenons les identifications (en tant que $A$-modules gradués)
$Q_1 = \Ker(\pi_1)$, $Q_3 = \Ker(\pi_3)$ et $Q_2 =
\Ker(\pi_2)$. Alors $L = K \oplus Q_1$ et $M = L \oplus
Q_3$ en tant que $A$-modules gradués. Cela implique $M = K
\oplus Q_1 \oplus Q_3$. Notons que $\pi_2 = \pi_1 \circ \pi_3$ est nul à la
fois sur $Q_1$ et $Q_3$. Par conséquent $Q_2 = Q_1 \oplus Q_3$.
Considérons le diagramme commutatif
$$
\begin{matrix}
0 & \to & K & \to & L & \to & Q_1 & \to &
0 \\ & & \downarrow & & \downarrow & & \downarrow & \\ 0 &
\to & K & \to & M & \to & Q_2 & \to & 0
\\ & & \downarrow & & \downarrow & & \downarrow & \\ 0 & \to
& L & \to & M & \to & Q_3 & \to & 0
\end{matrix}
$$
Les lignes de ce diagramme sont des suites exactes courtes admissibles, et
déterminent donc par définition des triangles distingués. De plus, les
flèches descendantes dans le diagramme ci-dessus sont compatibles avec les
scindages choisis et définissent donc des morphismes de triangles
$$
(K \to L \to Q_1 \to K[1])
\longrightarrow
(K \to M \to Q_2 \to K[1])
$$
et
$$
(K \to M \to Q_2 \to K[1])
\longrightarrow
(L \to M \to Q_3 \to L[1]).
$$
Notons que les scindages $Q_3 \to M$ de la suite inférieure dans le
diagramme fournissent un scindage pour la suite
scindée $0 \to Q_1 \to Q_2 \to Q_3 \to 0$ lorsqu'on les compose avec $M
\to Q_2$. Il en découle facilement que le morphisme $\delta : Q_3 \to
Q_1[1]$ dans le triangle distingué correspondant
$$
(Q_1 \to Q_2 \to Q_3 \to Q_1[1])
$$
est égal à la composition $Q_3 \to L[1] \to Q_1[1]$. Ainsi,
nous obtenons une structure comme dans la conclusion de l'axiome TR4.
\end{proof}

\noindent
Voici le résultat final.

\begin{proposition}
\label{dga-proposition-homotopy-category-triangulated}
Soit $(A, \text{d})$ une algèbre différentielle graduée. La catégorie homotopique
$K(\text{Mod}_{(A, \text{d})})$ des $A$-modules différentiels gradués avec ses
foncteurs de translation naturels et ses triangles distingués est une catégorie
triangulée.
\end{proposition}

\begin{proof}
Nous savons que $K(\text{Mod}_{(A, \text{d})})$ est une catégorie pré-triangulée. Par
conséquent, il suffit de prouver TR4 et pour le prouver, nous
pouvons utiliser Catégories dérivées, lemme
\ref{derived-lemma-easier-axiom-four}. Soient $K \to L$ et $L \to M$ des morphismes composables
de $K(\text{Mod}_{(A, \text{d})})$.
D'après le lemme \ref{dga-lemma-sequence-maps-split}, nous pouvons
supposer que $K \to L$ et $L \to M$ sont des monomorphismes
admissibles. Dans ce cas, le résultat découle du
lemme \ref{dga-lemma-two-split-injections}.
\end{proof}











\section{Modules à gauche}
\label{dga-section-left-modules}

\noindent
Tout ce que nous avons dit jusqu'à présent a un analogue dans le cadre
des modules différentiels gradués à gauche, sauf qu'il faut faire
attention à certaines règles de signe.

\medskip\noindent
Soit $(A, \text{d})$ une $R$-algèbre différentielle graduée.
Par analogie avec les modules à droite, nous
définissons un {\it module différentiel gradué à gauche sur $A$} $M$
comme un $A$-module à gauche $M$ muni d'une graduation $M =
\bigoplus M^n$ et d'une différentielle $\text{d}$ telles que $A^n M^m
\subset M^{n + m}$, que $\text{d}(M^n) \subset M^{n + 1}$, et que
$$
\text{d}(am) = \text{d}(a) m + (-1)^{\deg(a)}a \text{d}(m)
$$
pour les éléments homogènes $a \in A$ et $m \in M$. Comme auparavant, cette
règle de Leibniz signifie exactement que la multiplication définit un
morphisme de complexes
$$
\text{Tot}(A^\bullet \otimes_R M^\bullet) \to M^\bullet
$$
Ici, $A^\bullet$ et $M^\bullet$ désignent les complexes de $R$-modules
sous-jacents à $A$ et $M$.

\begin{definition}
\label{dga-definition-opposite-dga}
Soit $R$ un anneau. Soit $(A, \text{d})$ une algèbre différentielle graduée
sur $R$. L'{\it algèbre différentielle graduée opposée} est l'algèbre
différentielle graduée $(A^{opp}, \text{d})$ sur $R$ où $A^{opp} = A$ en tant
que $R$-module gradué, $\text{d} = \text{d}$, et la multiplication est
donnée par
$$
a \cdot_{opp} b = (-1)^{\deg(a)\deg(b)} b a
$$
pour les éléments homogènes $a, b \in A$.
\end{definition}

\noindent
Cela a du sens parce que
\begin{align*}
\text{d}(a \cdot_{opp} b)
& =
(-1)^{\deg(a)\deg(b)} \text{d}(b a) \\ &
=
(-1)^{\deg(a)\deg(b)} \text{d}(b) a +
(-1)^{\deg(a)\deg(b) + \deg(b)}b\text{d}(a) \\ & =
(-1)^{\deg(a)}a
\cdot_{opp} \text{d}(b) + \text{d}(a) \cdot_{opp} b
\end{align*}
comme souhaité. En termes de complexes sous-jacents de $R$-modules,
cela signifie que le diagramme
$$
\xymatrix{
\text{Tot}(A^\bullet \otimes_R A^\bullet)
\ar[rrr]_-{\text{multiplication de }A^{opp}}
\ar[d]_{\text{contrainte de commutativité}} & & &
A^\bullet \ar[d]^{\text{id}} \\
\text{Tot}(A^\bullet \otimes_R A^\bullet)
\ar[rrr]^-{\text{multiplication de }A} & & &
A^\bullet
}
$$
commute. Ici, la contrainte de commutativité sur la catégorie monoïdale
symétrique des complexes de $R$-modules est donnée dans
Compléments d'algèbre, section \ref{more-algebra-section-sign-rules}.

\medskip\noindent
Soit $(A, \text{d})$ une algèbre différentielle graduée sur $R$.
Soit $M$ un $A$-module différentiel gradué à gauche. Nous noterons
$M^{opp}$ le module $M$ considéré comme un $A^{opp}$-module à droite avec
multiplication $\cdot_{opp}$ définie par la règle
$$
m \cdot_{opp} a = (-1)^{\deg(a)\deg(m)} a m
$$
pour $a$ et $m$ homogènes. Cela est compatible avec les différentielles
car nous aurions pu utiliser le diagramme
$$
\xymatrix{
\text{Tot}(M^\bullet \otimes_R A^\bullet)
\ar[rrr]_-{\text{multiplication sur }M^{opp}}
\ar[d]_{\text{contrainte de commutativité}} & & &
M^\bullet \ar[d]^{\text{id}} \\
\text{Tot}(A^\bullet \otimes_R M^\bullet)
\ar[rrr]^-{\text{multiplication sur }M} & & &
M^\bullet
}
$$
pour définir la multiplication $\cdot_{opp}$ sur $M^{opp}$.
Pour voir qu'il s'agit d'une multiplication associative, nous calculons
pour les éléments homogènes $a, b \in A$ et $m \in M$ que
\begin{align*}
m \cdot_{opp} (a \cdot_{opp} b)
& =
(-1)^{\deg(a)\deg(b)} m \cdot_{opp} (ba) \\ &
=
(-1)^{\deg(a)\deg(b) + \deg(ab)\deg(m)} bam \\ & =
(-1)^{\deg(a)\deg(b)
+ \deg(ab)\deg(m) + \deg(b)\deg(am)} (am)
\cdot_{opp} b \\ & =
(-1)^{\deg(a)\deg(b)
+ \deg(ab)\deg(m) + \deg(b)\deg(am) + \deg(a)\deg(m)} (m
\cdot_{opp} a) \cdot_{opp} b \\ & = (m
\cdot_{opp}
a) \cdot_{opp} b
\end{align*}
Bien sûr, nous aurions pu le montrer en utilisant la compatibilité entre les
contraintes d'associativité et de commutativité dans la catégorie monoïdale
symétrique des complexes de $R$-modules également.

\begin{lemma}
\label{dga-lemma-left-right}
Soit $(A, \text{d})$ une $R$-algèbre différentielle graduée.
Le foncteur $M \mapsto M^{opp}$ de la catégorie des
$A$-modules différentiels gradués à gauche vers la catégorie des
$A^{opp}$-modules différentiels gradués à droite est une équivalence.
\end{lemma}

\begin{proof}
Démonstration omise.
\end{proof}

\noindent
Passons maintenant aux décalages. Soit $(A, \text{d})$ une algèbre différentielle
graduée. Soit $M$ un $A$-module différentiel gradué à gauche dont le complexe
sous-jacent de $R$-modules est noté $M^\bullet$.
Pour tout $k \in \mathbf{Z}$ nous définissons le {\it module décalé de $k$}
$M[k]$ comme suit
\begin{enumerate}
\item le complexe de $R$-modules sous-jacent à $M[k]$ est $M^\bullet[k]$
\item la structure de $A$-module est donnée par la multiplication
$$
A^n \times (M[k])^m \longrightarrow (M[k])^{n + m}
$$
est égale à $(-1)^{nk}$ multiplié par la multiplication
donnée $A^n \times M^{m + k} \to M^{n + m + k}$.
\end{enumerate}
Vérifions qu'avec ce choix la règle de Leibniz est satisfaite.
Soient $a \in A^n$ et $x \in M[k]^m = M^{m + k}$. En notant
$\cdot_{M[k]}$ le produit dans $M[k]$, on obtient
\begin{align*}
\text{d}_{M[k]}(a \cdot_{M[k]} x)
& =
(-1)^{k + nk} \text{d}_M(ax) \\
& =
(-1)^{k + nk} \text{d}(a) x + (-1)^{k + nk + n} a \text{d}_M(x) \\ &
=
\text{d}(a) \cdot_{M[k]} x + (-1)^{nk + n} a \text{d}_{M[k]}(x) \\ & =
\text{d}(a)
\cdot_{M[k]} x + (-1)^n a \cdot_{M[k]} \text{d}_{M[k]}(x)
\end{align*}
C'est ce que nous voulons puisque $a$ a un degré $n$ en tant qu'élément homogène de $A$. Nous
observons également qu'avec ces choix, nous pouvons considérer
l'application de multiplication comme un morphisme de complexes
$$
\text{Tot}(A^\bullet \otimes_R M^\bullet[k]) \to
\text{Tot}(A^\bullet \otimes_R M^\bullet)[k] \to
M^\bullet[k]
$$
où la première flèche est
celle de Compléments d'algèbre, section \ref{more-algebra-section-sign-rules}
(\ref{more-algebra-item-shift-tensor}), ce qui dans ce cas implique
exactement le signe que nous avons choisi ci-dessus. (En fait, nous aurions pu déduire que la
règle de Leibniz est respectée à partir de cette observation.)

\medskip\noindent
Avec la règle ci-dessus, nous avons des identifications canoniques
$$
(M[k])^{opp} = M^{opp}[k]
$$
de modules différentiels gradués à droite sur
$A^{opp}$ définis sans l'intervention de signes, en d'autres termes, l'équivalence du
lemme \ref{dga-lemma-left-right} est compatible avec les foncteurs de décalage.

\medskip\noindent
Notre choix ci-dessus nécessite la définition suivante.

\begin{definition}
\label{dga-definition-shift-graded-module}
Soit $R$ un anneau. Soit $A$ une algèbre graduée par $\mathbf{Z}$ sur $R$.
\begin{enumerate}
\item Étant donné un $A$-module gradué à droite $M$, nous définissons le
{\it décalé d'ordre $k$, en tant que $A$-module,} $M[k]$ comme étant
identique en tant que $A$-module à droite mais avec la graduation $(M[k])^n = M^{n + k}$.
\item Étant donné un $A$-module gradué à gauche $M$, nous
définissons le {\it module décalé d'ordre $k$ sur $A$} $M[k]$
comme le module de graduation $(M[k])^n = M^{n + k}$, dont la
multiplication $A^n \times (M[k])^m \to
(M[k])^{n + m}$ est égale à $(-1)^{nk}$ fois la multiplication
donnée $A^n \times M^{m + k} \to M^{n + m + k}$.
\end{enumerate}
\end{definition}

\noindent
Soit $(A, \text{d})$ une algèbre différentielle graduée. Soient
$f, g : M \to N$ des homomorphismes de $A$-modules différentiels gradués à gauche. Une
{\it homotopie entre $f$ et $g$} est une application de $A$-modules
gradués $h : M \to N[-1]$ (observez le décalage !) telle que
$$
f(x) - g(x) = \text{d}_N(h(x)) + h(\text{d}_M(x))
$$
pour tout $x \in M$. S'il existe une homotopie, nous disons alors que $f$ et
$g$ sont {\it homotopes}. Ainsi, $h$ est compatible avec la structure de
$A$-module (avec celle décalée sur $N$) et la graduation (avec la graduation décalée sur
$N$) mais pas avec la différentielle. Si $f = g$ et $h$ est une homotopie, alors
$h$ définit un morphisme $h : M \to N[-1]$ de $A$-modules différentiels
gradués à gauche.

\medskip\noindent
Avec la règle ci-dessus, nous constatons que $f, g : M \to N$ sont homotopes
si et seulement si les morphismes induits
$f^{opp}, g^{opp} : M^{opp} \to N^{opp}$ sont
homotopes en tant qu'homomorphismes de modules différentiels gradués à droite sur
$A^{opp}$ (avec la même homotopie).

\medskip\noindent
La catégorie homotopique, les cônes, les suites exactes courtes
admissibles, les triangles distingués sont tous définis exactement de la même manière
que pour les modules différentiels gradués à droite (et tout concorde sur
les complexes sous-jacents de $R$-modules avec les constructions pour les
complexes de $R$-modules). De cette manière, nous obtenons l'analogue de la
proposition \ref{dga-proposition-homotopy-category-triangulated} pour les
modules à gauche également, ou nous pouvons la déduire en travaillant
avec des modules à droite sur l'algèbre opposée.



\section{Produit tensoriel}
\label{dga-section-tensor-product}

\noindent
Soit $R$ un anneau. Soit $A$ une $R$-algèbre (voir la
section \ref{dga-section-conventions}). Étant donné un $A$-module à droite $M$ et
un $A$-module à gauche $N$, il existe un {\it produit tensoriel}
$$
M \otimes_A N
$$
Ce produit tensoriel est un module sur $R$. Comme $R$-module, $M \otimes_A
N$ est engendré par des symboles $x \otimes y$ avec $x \in M$ et $y \in N$
soumis aux relations
$$
\begin{matrix}
(x_1 + x_2) \otimes y - x_1 \otimes y - x_2 \otimes y, \\
x \otimes (y_1 + y_2) - x \otimes y_1 - x \otimes y_2, \\
xa \otimes y - x \otimes ay
\end{matrix}
$$
pour $a \in A$, $x, x_1, x_2 \in M$ et $y, y_1, y_2 \in N$. Nous
énumérons certaines propriétés du produit tensoriel

\medskip\noindent
Dans chaque variable, le produit tensoriel est exact à droite, en fait il commute
avec les sommes directes et les colimites arbitraires.

\medskip\noindent
Le produit tensoriel $M \otimes_A N$ est le réceptacle de l'application
$A$-bilinéaire universelle $M \times N \to M \otimes_A N$, $(x, y) \mapsto x \otimes y$.
Autrement dit,
$$
\text{Bilinear}_A(M \times N, Q) = \Hom_R(M \otimes_A N, Q)
$$
pour tout $R$-module $Q$.

\medskip\noindent
Si $A$ est une algèbre graduée par $\mathbf{Z}$ et si $M$, $N$ sont des
$A$-modules gradués, alors $M \otimes_A N$ est un $R$-module gradué.
La composante de degré $n$, notée $(M \otimes_A N)^n$, de $M \otimes_A N$
est égale à
$$
\Coker\left(
\bigoplus\nolimits_{r + t + s = n}
M^r \otimes_R A^t \otimes_R N^s \to
\bigoplus\nolimits_{p + q = n} M^p \otimes_R N^q
\right)
$$
où l'application envoie $x \otimes a \otimes y$
vers $x \otimes ay - xa \otimes y$ pour
$x \in M^r$, $y \in N^s$ et $a \in A^t$ avec $r + s + t = n$. Dans
ce cas, l'application $M \times N \to M \otimes_A N$ est $A$-bilinéaire et
compatible avec les graduations et universelle dans le sens que
$$
\text{GradedBilinear}_A(M \times N, Q) =
\Hom_{\text{}R\text{-modules gradués}}(M \otimes_A N, Q)
$$
pour tout $R$-module gradué $Q$ avec une notion évidente d'application
bilinéaire graduée.

\medskip\noindent
Si $(A, \text{d})$ est une algèbre différentielle graduée et
$M$ et $N$ sont des $A$-modules différentiels gradués à gauche et à droite,
alors $M \otimes_A N$ est un $R$-module différentiel gradué avec une différentielle
$$
\text{d}(x \otimes y) =
\text{d}(x) \otimes y + (-1)^{\deg(x)}x \otimes \text{d}(y)
$$
pour $x \in M$ et $y \in N$ homogènes. Dans ce cas, l'application
$M \times N \to M \otimes_A N$ est $A$-bilinéaire, compatible avec les graduations,
et compatible avec les différentielles et universelle dans le sens que
$$
\text{DifferentialGradedBilinear}_A(M \times N, Q) =
\Hom_{\text{Comp}(R)}(M \otimes_A N, Q)
$$
pour tout $R$-module différentiel gradué $Q$ avec une notion évidente
d'application bilinéaire différentielle graduée.






\section{Complexes de Hom et modules différentiels gradués}
\label{dga-section-hom-complexes}

\noindent
Nous encourageons le lecteur à sauter cette section.

\medskip\noindent
Soit $R$ un anneau et soit $M^\bullet$ un complexe de $R$-modules.
Considérons le complexe de $R$-modules
$$
E^\bullet = \Hom^\bullet(M^\bullet, M^\bullet)
$$
introduit dans
Compléments d'algèbre, section \ref{more-algebra-section-hom-complexes}.
D'après Compléments d'algèbre, lemme \ref{more-algebra-lemma-composition} il
existe une loi de composition canonique
$$
\text{Tot}(E^\bullet \otimes_R E^\bullet) \to E^\bullet
$$
qui est un morphisme de complexes. Ainsi, nous voyons que $E^\bullet$ avec
cette multiplication est une $R$-algèbre différentielle graduée que nous
noterons $(E, \text{d})$. De plus, en considérant $M^\bullet$ comme
$\Hom^\bullet(R, M^\bullet)$, nous voyons que la composition définit une multiplication
$$
\text{Tot}(E^\bullet \otimes_R M^\bullet) \to M^\bullet
$$
qui munit $M^\bullet$ d'une structure de $E$-module différentiel gradué à {\bf
gauche}, que nous noterons $M$.

\begin{lemma}
\label{dga-lemma-left-module-structure}
Dans la situation ci-dessus, soit $A$ une $R$-algèbre différentielle graduée.
Donner une structure de $A$-module à gauche sur $M$ revient au même que
donner un homomorphisme $A \to E$ d'algèbres différentielles graduées sur $R$.
\end{lemma}

\begin{proof}
Démonstration omise. Observez qu'aucun signe n'intervient dans cette correspondance.
\end{proof}

\noindent
Nous poursuivons avec la discussion ci-dessus et nous supposons donné un
autre complexe $N^\bullet$ de $R$-modules. Considérons le complexe
de $R$-modules $\Hom^\bullet(M^\bullet, N^\bullet)$ introduit dans
Compléments d'algèbre, section \ref{more-algebra-section-hom-complexes}. Comme
ci-dessus, nous voyons que la composition
$$
\text{Tot}(\Hom^\bullet(M^\bullet, N^\bullet) \otimes_R E^\bullet)
\to \Hom^\bullet(M^\bullet, N^\bullet)
$$
définit une multiplication qui transforme $\Hom^\bullet(M^\bullet, N^\bullet)$
en un $E$-module différentiel gradué à {\bf droite}. En utilisant
le lemme \ref{dga-lemma-left-module-structure},
nous concluons que, pour tout $A$-module différentiel gradué à gauche $M$
et tout complexe de $R$-modules $N^\bullet$, il existe un
$A$-module différentiel gradué canonique à droite dont le complexe
sous-jacent de $R$-modules est $\Hom^\bullet(M^\bullet, N^\bullet)$ et où
la multiplication
$$
\Hom^n(M^\bullet, N^\bullet) \times A^m \longrightarrow
\Hom^{n + m}(M^\bullet, N^\bullet)
$$
envoie $f = (f_{p, q})_{p + q = n}$ avec $f_{p, q} \in \Hom(M^{-q}, N^p)$
et $a \in A^m$ vers l'élément $f \cdot a = (f_{p, q} \circ a)$ où $f_{p,
q} \circ a$ est l'application
$$
M^{-q - m} \xrightarrow{a} M^{-q} \xrightarrow{f_{p, q}} N^p, \quad
x \longmapsto f_{p, q}(ax)
$$
sans l'intervention des signes. Utilisons la notation
$\Hom(M, N^\bullet)$ pour désigner ce $A$-module différentiel gradué à droite.

\begin{lemma}
\label{dga-lemma-characterize-hom}
Soit $R$ un anneau. Soit $(A, \text{d})$ une $R$-algèbre différentielle graduée.
Soit $M'$ un $A$-module différentiel gradué à droite et soit
$M$ un $A$-module différentiel gradué à gauche.
Soit $N^\bullet$ un complexe de $R$-modules. Alors nous avons
$$
\Hom_{\text{Mod}_{(A, d)}}(M', \Hom(M, N^\bullet)) =
\Hom_{\text{Comp}(R)}(M' \otimes_A M, N^\bullet)
$$
où $M \otimes_A M$ est considéré comme un complexe de $R$-modules
comme dans la section \ref{dga-section-tensor-product}.
\end{lemma}

\begin{proof}
Montrons que les deux côtés correspondent à des applications $A$-bilinéaires graduées
$$
M' \times M \longrightarrow N^\bullet
$$
compatibles avec les différentielles. Nous l'avons établi pour le membre de
droite dans la section \ref{dga-section-tensor-product}. Étant donné un élément
$g$ du membre de gauche, l'égalité de
Compléments d'algèbre, lemme \ref{more-algebra-lemma-compose}, détermine
un morphisme de complexes de $R$-modules $g' :
\text{Tot}(M' \otimes_R M) \to N^\bullet$. En d'autres
termes, nous obtenons un morphisme $R$-bilinéaire
gradué $g'' : M' \times M \to N^\bullet$ compatible avec les différentielles. La
$A$-linéarité de $g$ se traduit immédiatement par la
$A$-bilinéarité de $g''$.
\end{proof}

\noindent
Soient $R$, $M^\bullet$, $E^\bullet$, $E$ et $M$ comme ci-dessus.
Cependant, supposons maintenant données une $R$-algèbre différentielle graduée $A$ et
une structure à {\bf droite} de $A$-module différentiel gradué sur $M$. Alors
nous pouvons considérer le morphisme
$$
\text{Tot}(A^\bullet \otimes_R M^\bullet)
\xrightarrow{\psi}
\text{Tot}(A^\bullet \otimes_R M^\bullet)
\to
M^\bullet
$$
où la première flèche est la contrainte de commutativité sur la
catégorie différentielle graduée des complexes de $R$-modules. Cela
correspond à une application
$$
\tau : A^\bullet \longrightarrow E^\bullet
$$
de complexes de $R$-modules. Rappelons que
$E^n = \prod_{p + q = n} \Hom_R(M^{-q}, M^p)$ et
écrivons $\tau(a) = (\tau_{p, q}(a))_{p + q = n}$ pour $a \in A^n$. Alors
nous voyons
$$
\tau_{p, q}(a) : M^{-q} \longrightarrow M^p,\quad
x \longmapsto (-1)^{\deg(a)\deg(x)}x a = (-1)^{-nq}xa
$$
Cela n'est pas compatible avec le produit sur $A$ comme le lecteur devrait s'y
attendre d'après la discussion de la section \ref{dga-section-left-modules}. À
savoir, nous avons
$$
\tau(a a') = (-1)^{\deg(a)\deg(a')}\tau(a') \tau(a)
$$
Nous concluons que le lemme suivant est vrai

\begin{lemma}
\label{dga-lemma-right-module-structure}
Dans la situation ci-dessus, soit $A$ une $R$-algèbre différentielle graduée.
Donner une structure de $A$-module à droite sur $M$ revient au même que
donner un homomorphisme $\tau : A \to E^{opp}$
d'algèbres différentielles graduées sur $R$.
\end{lemma}

\begin{proof}
Voir la discussion ci-dessus et noter que la construction de $\tau$ à partir
de l'application de multiplication $M^n \times A^m \to M^{n + m}$
utilise des signes.
\end{proof}

\noindent
Soient $R$, $M^\bullet$, $E^\bullet$, $E$, $A$ et $M$ comme ci-dessus et
soit une structure de $A$-module différentiel gradué à droite sur $M$
donnée comme dans le lemme. Dans ce cas, il existe un $A$-module différentiel
gradué à gauche canonique dont le complexe sous-jacent de $R$-modules est
$\Hom^\bullet(M^\bullet, N^\bullet)$. À savoir, pour la multiplication, nous pouvons
utiliser
\begin{align*}
\text{Tot}(A^\bullet \otimes_R \Hom^\bullet(M^\bullet, N^\bullet))
& \xrightarrow{\psi}
\text{Tot}(\Hom^\bullet(M^\bullet, N^\bullet) \otimes_R A^\bullet) \\ &
\xrightarrow{\tau}
\text{Tot}(\Hom^\bullet(M^\bullet, N^\bullet) \otimes_R
\Hom^\bullet(M^\bullet, M^\bullet)) \\ &
\to
\text{Tot}(\Hom^\bullet(M^\bullet, N^\bullet)
\end{align*}
La première flèche utilise la contrainte de commutativité sur la catégorie
des complexes de $R$-modules, la deuxième flèche est décrite ci-dessus, et la
troisième flèche est la loi de composition pour le complexe Hom. Chacune de ces
applications est un morphisme de complexes ; leur composé est donc un morphisme
de complexes. En fait, cette construction transforme $\Hom^\bullet(M^\bullet,
N^\bullet)$ en un $A$-module différentiel gradué à gauche (l'associativité de la multiplication
peut être démontrée en utilisant la structure monoïdale symétrique ou par un calcul direct en
utilisant les formules ci-dessous). Explicitons la multiplication
$$
A^n \times \Hom^m(M^\bullet, N^\bullet) \longrightarrow
\Hom^{n + m}(M^\bullet, N^\bullet)
$$
Elle envoie $a \in A^n$
et $f = (f_{p, q})_{p + q = m}$ avec $f_{p, q} \in \Hom(M^{-q}, N^p)$ sur
l'élément $a \cdot f$ dont les composantes sont
$$
(-1)^{nm}f_{p, q} \circ \tau_{-q, q + n}(a) =
(-1)^{nm - n(q + n)}f_{p, q} \circ a =
(-1)^{np + n} f_{p, q} \circ a
$$
dans $\Hom_R(M^{-q - n}, N^p)$ où $f_{p, q} \circ a$ est l'application
$$
M^{-q - n} \xrightarrow{a} M^{-q} \xrightarrow{f_{p, q}} N^p,\quad
x \longmapsto f_{p, q}(xa)
$$
Ici intervient un signe $(-1)^{np + n}$. Utilisons la notation
$\Hom(M, N^\bullet)$ pour désigner ce $A$-module différentiel gradué à gauche.

\begin{lemma}
\label{dga-lemma-characterize-hom-other-side}
Soit $R$ un anneau. Soit $(A, \text{d})$ une $R$-algèbre différentielle graduée.
Soit $M$ un $A$-module différentiel gradué à droite et soit
$M'$ un $A$-module différentiel gradué à gauche.
Soit $N^\bullet$ un complexe de $R$-modules. Alors nous avons
$$
\Hom_{\text{}A\text{-modules différentiels gradués à gauche}}(M', \Hom(M,
N^\bullet)) = \Hom_{\text{Comp}(R)}(M \otimes_A M', N^\bullet)
$$
où $M \otimes_A M'$ est vu comme un complexe de $R$-modules
comme dans la section \ref{dga-section-tensor-product}.
\end{lemma}

\begin{proof}
Montrons que les deux membres correspondent à des applications bilinéaires graduées sur $A$
$$
M \times M' \longrightarrow N^\bullet
$$
compatibles avec les différentielles. Nous avons vu que cela est vrai pour le
membre de droite dans la section \ref{dga-section-tensor-product}. Étant donné un
élément $g$ du membre de gauche, l'égalité de
Compléments d'algèbre, lemme \ref{more-algebra-lemma-compose},
détermine un morphisme de complexes
$g' : \text{Tot}(M' \otimes_R M) \to N^\bullet$. Nous
précomposons avec la contrainte de commutativité pour obtenir
$$
\text{Tot}(M \otimes_R M') \xrightarrow{\psi}
\text{Tot}(M' \otimes_R M) \xrightarrow{g'}
N^\bullet
$$
Ce composé correspond à une application
$R$-bilinéaire graduée $g'' : M \times M' \to N^\bullet$ compatible avec les
différentielles. La $A$-linéarité de $g$ se traduit immédiatement par la $A$-bilinéarité de
$g''$. En effet, soient $x \in M^e$ et $x' \in (M')^{e'}$ et $a \in A^n$.
Alors d'une part nous avons
\begin{align*}
g''(x, ax')
& =
(-1)^{e(n + e')} g'(ax' \otimes x) \\ &
=
(-1)^{e(n + e')} g(ax')(x) \\ &
=
(-1)^{e(n + e')} (a \cdot g(x'))(x) \\ & =
(-1)^{e(n
+ e') + n(n + e + e') + n} g(x')(xa)
\end{align*}
et d'autre part nous avons
$$
g''(xa, x') = (-1)^{(e + n)e'} g'(x' \otimes xa) =
(-1)^{(e + n)e'} g(x')(xa)
$$
Ces deux expressions sont égales, comme le montre un calcul élémentaire modulo $2$ sur les exposants.
\end{proof}

\begin{remark}
\label{dga-remark-evaluation-map-left}
Soit $R$ un anneau. Soit $A$ une $R$-algèbre différentielle graduée.
Soit $M$ un $A$-module différentiel gradué à gauche. Soit
$N^\bullet$ un complexe de $R$-modules. Les constructions ci-dessus produisent
un $A$-module différentiel gradué à droite $\Hom(M, N^\bullet)$, puis un
$A$-module différentiel gradué à gauche $\Hom(\Hom(M,
N^\bullet), N^\bullet)$. Nous affirmons qu'il existe une application
d'évaluation
$$
ev : M \longrightarrow \Hom(\Hom(M, N^\bullet), N^\bullet)
$$
dans la catégorie des $A$-modules différentiels gradués à gauche. Pour la définir,
d'après le lemme \ref{dga-lemma-characterize-hom}, il suffit de construire un
couplage $A$-bilinéaire
$$
\Hom(M, N^\bullet) \times M \longrightarrow N^\bullet
$$
compatible avec la graduation et les différentielles. Pour cela nous prenons
$$
(f, x) \longmapsto f(x)
$$
Nous laissons au lecteur le soin de vérifier que ceci est compatible avec la graduation,
les différentielles, et bilinéaire sur $A$. L'application $ev$ sur les complexes de
$R$-modules sous-jacents est celle de Compléments d'algèbre, point (\ref{more-algebra-item-evaluation}).
\end{remark}

\begin{remark}
\label{dga-remark-evaluation-map-right}
Soit $R$ un anneau. Soit $A$ une $R$-algèbre différentielle graduée.
Soit $M$ un $A$-module différentiel gradué à droite. Soit
$N^\bullet$ un complexe de $R$-modules. Les constructions ci-dessus produisent
un $A$-module différentiel gradué à gauche $\Hom(M, N^\bullet)$, puis un
$A$-module différentiel gradué à droite $\Hom(\Hom(M,
N^\bullet), N^\bullet)$. Nous affirmons qu'il existe une application d'évaluation
$$
ev : M \longrightarrow \Hom(\Hom(M, N^\bullet), N^\bullet)
$$
dans la catégorie des $A$-modules différentiels gradués à droite. Pour la définir,
d'après le lemme \ref{dga-lemma-characterize-hom}, il suffit de construire un
couplage $A$-bilinéaire
$$
M \times \Hom(M, N^\bullet) \longrightarrow N^\bullet
$$
compatible avec la graduation et les différentielles. Pour cela, nous prenons
$$
(x, f) \longmapsto (-1)^{\deg(x)\deg(f)}f(x)
$$
Nous laissons au lecteur le soin de vérifier que cela est compatible avec la
graduation, les différentielles et bilinéaire sur $A$. L'application $ev$ sur les complexes de
$R$-modules sous-jacents est celle de Compléments d'algèbre, point (\ref{more-algebra-item-evaluation}).
\end{remark}

\begin{remark}
\label{dga-remark-shift-dual}
Soit $R$ un anneau. Soit $A$ une algèbre différentielle graduée sur
$R$. Soient $M^\bullet$ et $N^\bullet$ des complexes de
$R$-modules. Soit $k \in \mathbf{Z}$ et considérons l'isomorphisme
$$
\Hom^\bullet(M^\bullet, N^\bullet)[-k]
\longrightarrow
\Hom^\bullet(M^\bullet[k], N^\bullet)
$$
de complexes de $R$-modules défini dans
Compléments d'algèbre, point (\ref{more-algebra-item-shift-hom}). Si
$M^\bullet$ a la structure d'un module différentiel gradué à gauche, resp.\ à droite,
sur $A$, alors il s'agit d'un homomorphisme de modules
différentiels gradués à droite, resp.\ à gauche, sur $A$ (avec les
structures de module telles que définies dans cette section).
Nous omettons la vérification ; nous avertissons le lecteur que la
structure de $A$-module sur le décalage d'un $A$-module gradué à gauche est
définie en utilisant un signe,
voir définition \ref{dga-definition-shift-graded-module}.
\end{remark}






\section{Modules projectifs sur les algèbres}
\label{dga-section-projectives-over-algebras}

\noindent
Dans cette section, nous étudions les modules projectifs sur des algèbres, par
analogie avec Algèbre, section
\ref{algebra-section-projective}. Cette section devrait probablement être déplacée ailleurs.

\medskip\noindent
Soit $R$ un anneau et soit $A$ une $R$-algèbre, voir la
section \ref{dga-section-conventions} pour nos conventions. Il
est clair que $A$ est un $A$-module à droite projectif puisque
$\Hom_A(A, M) = M$ pour tout $A$-module à droite $M$ (et donc $\Hom_A(A, -)$ est
exact). Inversement, soit $P$ un $A$-module à droite projectif. Alors nous
pouvons choisir une surjection
$\bigoplus_{i \in I} A \to P$ en choisissant un ensemble $\{p_i\}_{i \in I}$
de générateurs de $P$ sur $A$. Puisque $P$ est projectif, il existe un
inverse à gauche de la surjection, et nous trouvons que $P$ est isomorphe à
un facteur direct d'un module libre, exactement comme dans le cas commutatif
(Algèbre, lemme \ref{algebra-lemma-characterize-projective}).

\medskip\noindent
Nous concluons
\begin{enumerate}
\item la catégorie des $A$-modules a suffisamment de projectifs,
\item $A$ est un $A$-module projectif.
\item chaque $A$-module est un quotient d'une somme directe de copies de $A$,
\item chaque $A$-module projectif est un facteur direct d'une somme
directe de copies de $A$.
\end{enumerate}



\section{Modules projectifs sur les algèbres graduées}
\label{dga-section-projectives-over-graded-algebras}

\noindent
Dans cette section, nous étudions les modules gradués projectifs sur des algèbres
graduées, par analogie avec Algèbre, section \ref{algebra-section-projective}.

\medskip\noindent
Soit $R$ un anneau. Soit $A$ une algèbre $\mathbf{Z}$-graduée sur $R$.
Voir la section \ref{dga-section-conventions} pour nos
conventions. Soit $\text{Mod}_A$ la catégorie des $A$-modules gradués à droite.
Pour un entier $k$, soit $A[k]$ le décalage de $A$. Pour un
$A$-module gradué à droite, nous avons
$$
\Hom_{\text{Mod}_A}(A[k], M) = M^{-k}
$$
Comme le foncteur $M \mapsto M^{-k}$ est exact sur $\text{Mod}_A$, nous
en concluons que $A[k]$ est un objet projectif de $\text{Mod}_A$.
Inversement, supposons que $P$ soit un objet projectif de $\text{Mod}_A$. En
choisissant un ensemble de générateurs homogènes de $P$ comme $A$-module, nous
pouvons trouver une surjection
$$
\bigoplus\nolimits_{i \in I} A[k_i] \longrightarrow P
$$
Ainsi, nous concluons qu'un objet projectif de $\text{Mod}_A$ est
un facteur direct d'une somme directe des décalages $A[k]$.

\medskip\noindent
Nous concluons
\begin{enumerate}
\item la catégorie des $A$-modules gradués a suffisamment d'objets projectifs,
\item $A[k]$ est un $A$-module projectif pour chaque $k \in \mathbf{Z}$,
\item chaque $A$-module gradué est un quotient d'une somme directe
de copies des modules $A[k]$ pour différents $k$.
\item tout $A$-module projectif est un facteur direct d'une somme
directe de copies des modules $A[k]$ pour des $k$ variables.
\end{enumerate}




\section{Modules projectifs et algèbres différentielles graduées}
\label{dga-section-projective-over-differential-graded}

\noindent
Si $(A, \text{d})$ est une algèbre différentielle graduée et $P$ est
un objet de $\text{Mod}_{(A, \text{d})}$ alors nous disons
{\it $P$ est projectif en tant que $A$-module gradué} ou parfois
{\it $P$ est projectif gradué} pour signifier que $P$
est un objet projectif de la catégorie abélienne $\text{Mod}_A$ des
$A$-modules gradués comme dans
la section \ref{dga-section-projectives-over-graded-algebras}.

\begin{lemma}
\label{dga-lemma-target-graded-projective}
Soit $(A, \text{d})$ une algèbre différentielle graduée.
Soit $M \to P$ un homomorphisme surjectif de $A$-modules différentiels
gradués. Si $P$ est projectif en tant que $A$-module gradué, alors
$M \to P$ est un épimorphisme admissible.
\end{lemma}

\begin{proof}
Cela résulte immédiatement des définitions.
\end{proof}

\begin{lemma}
\label{dga-lemma-hom-from-shift-free}
Soit $(A, d)$ une algèbre différentielle graduée. Alors nous avons
$$
\Hom_{\text{Mod}_{(A, \text{d})}}(A[k], M) =
\Ker(\text{d} : M^{-k} \to M^{-k + 1})
$$
et
$$
\Hom_{K(\text{Mod}_{(A, \text{d})})}(A[k], M) = H^{-k}(M)
$$
pour tout $A$-module différentiel gradué $M$.
\end{lemma}

\begin{proof}
Immédiat d'après les définitions.
\end{proof}







\section{Modules injectifs sur les algèbres}
\label{dga-section-modules-noncommutative}

\noindent
Dans cette section, nous étudions les modules injectifs sur des
algèbres, par
analogie avec Compléments d'algèbre, section
\ref{more-algebra-section-injectives-modules}. Cette section devrait probablement être déplacée ailleurs.

\medskip\noindent
Soit $R$ un anneau et soit $A$ une
$R$-algèbre, voir la section \ref{dga-section-conventions} pour nos conventions. Pour un
$A$-module à droite $M$, nous posons
$$
M^\vee = \Hom_\mathbf{Z}(M, \mathbf{Q}/\mathbf{Z})
$$
que nous considérons comme un $A$-module à gauche par la
multiplication $(a f)(x) = f(xa)$. À savoir, $((ab)f)(x) = f(xab) = (bf)(xa) = (a(bf))(x)$.
Inversement, si $M$ est un $A$-module à gauche, alors $M^\vee$ est un
$A$-module à droite. Puisque $\mathbf{Q}/\mathbf{Z}$ est un groupe abélien
injectif (Compléments d'algèbre, lemme \ref{more-algebra-lemma-injective-abelian}), le
foncteur $M \mapsto M^\vee$ est exact
(Compléments d'algèbre, lemme \ref{more-algebra-lemma-vee-exact}). De plus,
l'application d'évaluation $M \to (M^\vee)^\vee$ est
injective pour tous les modules $M$
(Compléments d'algèbre, lemme \ref{more-algebra-lemma-ev-injective}).

\medskip\noindent
Nous affirmons que $A^\vee$ est un $A$-module à droite injectif. À savoir, étant donné
un $A$-module à droite $N$ nous avons
$$
\Hom_A(N, A^\vee) =
\Hom_A(N, \Hom_\mathbf{Z}(A, \mathbf{Q}/\mathbf{Z})) = N^\vee
$$
et nous concluons parce que le foncteur $N \mapsto N^\vee$ est
exact. La seconde égalité tient parce que
$$
\Hom_\mathbf{Z}(N, \Hom_\mathbf{Z}(A, \mathbf{Q}/\mathbf{Z})) =
\Hom_\mathbf{Z}(N \otimes_\mathbf{Z} A, \mathbf{Q}/\mathbf{Z})
$$
d'après Algèbre, lemme \ref{algebra-lemma-hom-from-tensor-product}. À
l'intérieur de ce module, la $A$-linéarité sélectionne exactement les applications
bilinéaires $\varphi : N \times A \to \mathbf{Q}/\mathbf{Z}$ qui
ont la même valeur sur $x \otimes a$ et $xa \otimes 1$, c'est-à-dire qui
proviennent d'éléments de $N^\vee$.

\medskip\noindent
Enfin, pour tout $A$-module à droite $M$, nous pouvons choisir une
surjection $\bigoplus_{i \in I} A \to M^\vee$ pour obtenir une
injection $M \to (M^\vee)^\vee \to \prod_{i \in I} A^\vee$.

\medskip\noindent
Nous concluons
\begin{enumerate}
\item que la catégorie des $A$-modules possède suffisamment d'injectifs,
\item $A^\vee$ est un $A$-module injectif, et
\item tout $A$-module s'injecte dans un produit de copies de $A^\vee$.
\end{enumerate}





\section{Modules injectifs sur les algèbres graduées}
\label{dga-section-modules-noncommutative-graded}

\noindent
Dans cette section, nous discutons des modules gradués injectifs sur les algèbres
graduées de manière
analogue à Compléments d'algèbre, section \ref{more-algebra-section-injectives-modules}.

\medskip\noindent
Soit $R$ un anneau. Soit $A$ une algèbre graduée par $\mathbf{Z}$ sur $R$. Voir
la section \ref{dga-section-conventions} pour nos conventions. Si
$M$ est un $R$-module gradué, nous posons
$$
M^\vee =
\bigoplus\nolimits_{n \in \mathbf{Z}}
\Hom_\mathbf{Z}(M^{-n}, \mathbf{Q}/\mathbf{Z}) =
\bigoplus\nolimits_{n \in \mathbf{Z}} (M^{-n})^\vee
$$
comme un $R$-module gradué (aucun signe dans les actions de $R$ sur
les parties homogènes). Si $M$ a la structure d'un $A$-module à gauche
gradué, alors nous définissons une structure de $A$-module gradué à droite
sur $M^\vee$ en laissant $a \in A^m$ agir par
$$
(M^{-n})^\vee \to (M^{-n - m})^\vee, \quad
f \mapsto f \circ a
$$
comme dans la section
\ref{dga-section-hom-complexes}. Si $M$ a la structure d'un $A$-module à droite
gradué, alors nous définissons une structure de $A$-module gradué à
gauche sur $M^\vee$ en laissant $a \in A^n$ agir par
$$
(M^{-m})^\vee \to (M^{-m - n})^\vee, \quad
f \mapsto (-1)^{nm}f \circ a
$$
comme dans la section \ref{dga-section-hom-complexes} (le signe nous est imposé
parce que nous voulons utiliser la même formule pour le cas où
l'on travaille avec des modules différentiels gradués --- si l'on ne
s'intéresse qu'aux modules gradués, on peut omettre le signe ici). Sur la
catégorie des $A$-modules gradués (à gauche ou à droite), le foncteur $M
\mapsto M^\vee$ est exact (vérifier sur les composantes graduées). De
plus, il existe une application d'évaluation injective
$$
ev : M \longrightarrow (M^\vee)^\vee, \quad
ev^n = (-1)^n \text{ l'application d'évaluation }M^n \to ((M^n)^\vee)^\vee
$$
de modules gradués sur $R$,
voir Compléments d'algèbre, point (\ref{more-algebra-item-evaluation}).
Cette application d'évaluation est un homomorphisme de $A$-modules à gauche, resp.\ à droite,
si $M$ est un $A$-module à gauche, resp.\ à
droite ; voir les remarques \ref{dga-remark-evaluation-map-left} et
\ref{dga-remark-evaluation-map-right}. Enfin, étant donné $k \in \mathbf{Z}$, il existe un isomorphisme canonique
$$
M^\vee[-k] \longrightarrow (M[k])^\vee
$$
de $R$-modules gradués qui utilise un signe et qui, si $M$
est un $A$-module à gauche, resp. \ à droite, est un isomorphisme de
$A$-modules à droite, resp. \ à gauche. Voir la remarque \ref{dga-remark-shift-dual}.

\medskip\noindent
Nous affirmons que $A^\vee$ est un objet injectif de la catégorie
$\text{Mod}_A$ des $A$-modules gradués à droite. En effet, étant donné un
$A$-module gradué à droite $N$ nous avons
$$
\Hom_{\text{Mod}_A}(N, A^\vee) =
\Hom_{\text{Comp}(\mathbf{Z})}(N \otimes_A A, \mathbf{Q}/\mathbf{Z})) =
(N^0)^\vee
$$
selon le lemme
\ref{dga-lemma-characterize-hom} (appliqué au cas où toutes les différentielles sont nulles).
Nous en concluons parce que le foncteur $N \mapsto (N^0)^\vee = (N^\vee)^0$
est exact.

\medskip\noindent
Enfin, pour chaque $A$-module gradué à droite $M$ nous pouvons choisir une surjection
de $A$-modules gradués à gauche
$$
\bigoplus\nolimits_{i \in I} A[k_i] \to M^\vee
$$
où $A[k_i]$ désigne le décalage de $A$ par $k_i \in \mathbf{Z}$. Nous
faisons cela en choisissant des générateurs homogènes pour $M^\vee$. De cette
manière, nous obtenons une injection
$$
M \to (M^\vee)^\vee \to \prod A[k_i]^\vee = \prod A^\vee[-k_i]
$$
Remarquez que les produits dans la formule ci-dessus sont des produits
dans la catégorie des modules gradués (autrement dit, prendre les produits degré
par degré, puis la somme directe des composantes).

\medskip\noindent
Nous concluons que
\begin{enumerate}
\item la catégorie des $A$-modules gradués a assez d'injectifs,
\item pour chaque $k \in \mathbf{Z}$ le module $A^\vee[k]$ est injectif, et
\item chaque $A$-module s'injecte dans un produit dans la catégorie des
modules gradués de copies de décalages $A^\vee[k]$.
\end{enumerate}





\section{Modules injectifs et algèbres différentielles graduées}
\label{dga-section-modules-noncommutative-differential-graded}

\noindent
Si $(A, \text{d})$ est une algèbre différentielle graduée et $I$ est
un objet de $\text{Mod}_{(A, \text{d})}$ alors nous disons
que {\it $I$ est injectif en tant que $A$-module gradué} ou
parfois {\it $I$ est injectif gradué} pour
signifier que $I$ est un objet injectif de la catégorie abélienne $\text{Mod}_A$ des
$A$-modules gradués.

\begin{lemma}
\label{dga-lemma-source-graded-injective}
Soit $(A, \text{d})$ une algèbre différentielle
graduée. Soit $I \to M$ un homomorphisme injectif de $A$-modules
différentiels gradués. Si $I$ est injectif gradué, alors $I
\to M$ est un monomorphisme admissible.
\end{lemma}

\begin{proof}
Cela résulte immédiatement des définitions.
\end{proof}

\noindent
Soit $(A, \text{d})$ une algèbre différentielle graduée. Si $M$ est un
$A$-module différentiel gradué à gauche, resp.\ à droite, alors
$$
M^\vee = \Hom^\bullet(M^\bullet, \mathbf{Q}/\mathbf{Z})
$$
avec la structure de $A$-module construite
dans la section \ref{dga-section-modules-noncommutative-graded} est un
$A$-module différentiel gradué à droite, resp.\ à gauche,
selon la discussion dans la section
\ref{dga-section-hom-complexes}. Selon les
remarques \ref{dga-remark-evaluation-map-left} et \ref{dga-remark-evaluation-map-right},
l'application d'évaluation de la section \ref{dga-section-modules-noncommutative-graded}
$$
M \longrightarrow (M^\vee)^\vee
$$
est un homomorphisme de $A$-modules différentiels gradués à gauche, resp.\ à droite.

\begin{lemma}
\label{dga-lemma-map-into-dual}
Soit $(A, \text{d})$ une algèbre différentielle graduée. Si
$M$ est un $A$-module différentiel gradué à gauche et $N$ est un
$A$-module différentiel gradué à droite, alors
\begin{align*}
\Hom_{\text{Mod}_{(A, \text{d})}}(N, M^\vee)
& =
\Hom_{\text{Comp}(\mathbf{Z})}(N \otimes_A M, \mathbf{Q}/\mathbf{Z}) \\ &
=
\text{DifferentialGradedBilinear}_A(N \times M, \mathbf{Q}/\mathbf{Z})
\end{align*}
\end{lemma}

\begin{proof}
La première égalité découle du lemme
\ref{dga-lemma-characterize-hom} et la deuxième égalité a été démontrée dans la section \ref{dga-section-tensor-product}.
\end{proof}

\begin{lemma}
\label{dga-lemma-hom-into-shift-dual-free}
Soit $(A, \text{d})$ une algèbre différentielle graduée. Alors nous avons
$$
\Hom_{\text{Mod}_{(A, \text{d})}}(M, A^\vee[k]) =
\Ker(\text{d} : (M^\vee)^k \to (M^\vee)^{k + 1})
$$
et
$$
\Hom_{K(\text{Mod}_{(A, \text{d})})}(M, A^\vee[k]) = H^k(M^\vee)
$$
comme foncteurs du $A$-module différentiel gradué $M$.
\end{lemma}

\begin{proof}
Cela est clair à partir de la discussion ci-dessus.
\end{proof}















\section{P-résolutions}
\label{dga-section-P-resolutions}

\noindent
Cette section est l'analogue de
Catégories dérivées, section \ref{derived-section-unbounded}.

\medskip\noindent
Soit $(A, \text{d})$ une algèbre différentielle graduée.
Soit $P$ un $A$-module différentiel gradué. Nous disons
que $P$ {\it a la propriété (P)} s'il existe une filtration
$$
0 = F_{-1}P \subset F_0P \subset F_1P \subset \ldots \subset P
$$
par des sous-modules différentiels gradués, satisfaisant aux conditions suivantes :
\begin{enumerate}
\item $P = \bigcup F_pP$,
\item les inclusions $F_iP \to F_{i + 1}P$ sont des
monomorphismes admissibles,
\item les quotients $F_{i + 1}P/F_iP$ sont isomorphes en tant que modules
différentiels gradués sur $A$ à une somme directe de $A[k]$.
\end{enumerate}
En fait, la condition (2) est une conséquence de la condition (3),
voir le lemme \ref{dga-lemma-target-graded-projective}. De plus, le lecteur
peut vérifier que, en tant que $A$-module gradué, $P$ sera isomorphe à une
somme directe de décalages de $A$.

\begin{lemma}
\label{dga-lemma-property-P-sequence}
Soit $(A, \text{d})$ une algèbre différentielle graduée.
Soit $P$ un module différentiel gradué sur $A$. Si $F_\bullet$ est une
filtration comme dans la propriété (P), alors nous obtenons
une suite exacte courte admissible
$$
0 \to
\bigoplus\nolimits F_iP \to
\bigoplus\nolimits F_iP \to P \to 0
$$
de modules différentiels gradués sur $A$.
\end{lemma}

\begin{proof}
La deuxième application est la somme directe des applications
d'inclusion. La première application sur le facteur direct $F_iP$ de la source
est la somme de l'identité $F_iP \to F_iP$ et de l'opposée de l'application
d'inclusion $F_iP \to F_{i + 1}P$. Choisissons des homomorphismes $s_i : F_{i + 1}P \to F_iP$ de
modules gradués sur $A$ qui sont des inverses à gauche des applications
d'inclusion. La composition donne des applications $s_{j, i} : F_jP \to F_iP$ pour tout
$j > i$. Ensuite, un inverse à gauche de la première flèche envoie $x \in
F_jP$ sur $(s_{j, 0}(x), s_{j, 1}(x), \ldots, s_{j, j - 1}(x), 0, \ldots)$ dans
$\bigoplus F_iP$.
\end{proof}

\noindent
Le lemme suivant montre que les modules différentiels gradués ayant la
propriété (P) sont la notion duale des modules K-injectifs
(c'est-à-dire qu'ils sont K-projectifs dans un certain
sens). Voir Catégories dérivées, définition \ref{derived-definition-K-injective}.

\begin{lemma}
\label{dga-lemma-property-P-K-projective}
Soit $(A, \text{d})$ une algèbre différentielle graduée.
Soit $P$ un $A$-module différentiel gradué ayant la propriété (P).
Alors
$$
\Hom_{K(\text{Mod}_{(A, \text{d})})}(P, N) = 0
$$
pour tous les $A$-modules différentiels gradués acycliques $N$.
\end{lemma}

\begin{proof}
Nous allons utiliser que $K(\text{Mod}_{(A, \text{d})})$ est une
catégorie triangulée (proposition
\ref{dga-proposition-homotopy-category-triangulated}). Soit $F_\bullet$ une filtration sur $P$ comme dans la
propriété (P). La suite exacte courte du lemme
\ref{dga-lemma-property-P-sequence} produit un triangle distingué. Donc, d'après
Catégories dérivées, lemme \ref{derived-lemma-representable-homological}, il
suffit de montrer que
$$
\Hom_{K(\text{Mod}_{(A, \text{d})})}(F_iP, N) = 0
$$
pour tous les $A$-modules différentiels gradués acycliques $N$ et tous les
$i$. Chacun des modules différentiels gradués $F_iP$ possède une filtration finie par des
monomorphismes admissibles, dont les quotients successifs sont des sommes
directes de décalages $A[k]$. Ainsi, il suffit de prouver que
$$
\Hom_{K(\text{Mod}_{(A, \text{d})})}(A[k], N) = 0
$$
pour tous les $A$-modules différentiels gradués acycliques $N$ et tous
les $k$. Cela découle du lemme \ref{dga-lemma-hom-from-shift-free}.
\end{proof}

\begin{lemma}
\label{dga-lemma-good-quotient}
Soit $(A, \text{d})$ une algèbre différentielle
graduée. Soit $M$ un $A$-module différentiel gradué. Il existe un homomorphisme
$P \to M$ de $A$-modules différentiels gradués avec les propriétés
suivantes
\begin{enumerate}
\item $P \to M$ est surjectif,
\item $\Ker(\text{d}_P) \to \Ker(\text{d}_M)$ est surjectif, et
\item $P$ s'insère dans une suite exacte courte admissible
$0 \to P' \to P \to P'' \to 0$ où $P'$, $P''$ sont des sommes directes
de décalages de $A$.
\end{enumerate}
\end{lemma}

\begin{proof}
Soit $P_k$ le $A$-module libre avec des générateurs $x, y$ en degrés $k$ et
$k + 1$. Définissons la structure de $A$-module différentiel gradué sur
$P_k$ en posant $\text{d}(x) = y$ et $\text{d}(y) = 0$. Pour chaque élément $m
\in M^k$, il existe un homomorphisme $P_k \to M$ envoyant $x$
vers $m$ et $y$ vers $\text{d}(m)$. Ainsi, nous voyons qu'il existe
une surjection d'une somme directe de copies de $P_k$ vers $M$.
Cela produit clairement $P \to M$ ayant les propriétés (1) et (3).
Pour obtenir la propriété (2), remarquons que si $m \in
\Ker(\text{d}_M)$ a le degré $k$, alors il existe une application $A[k] \to M$ envoyant
$1$ vers $m$. Par conséquent, nous pouvons atteindre (2) en ajoutant une somme
directe de copies de décalages de $A$.
\end{proof}

\begin{lemma}
\label{dga-lemma-resolve}
Soit $(A, \text{d})$ une algèbre différentielle
graduée. Soit $M$ un $A$-module différentiel gradué. Il existe un homomorphisme
$P \to M$ de $A$-modules différentiels gradués tel que
\begin{enumerate}
\item $P \to M$ soit un quasi-isomorphisme, et
\item $P$ ait la propriété (P).
\end{enumerate}
\end{lemma}

\begin{proof}
Posons $M = M_0$. Nous choisissons par récurrence des suites exactes courtes
$$
0 \to M_{i + 1} \to P_i \to M_i \to 0
$$
où les applications $P_i \to M_i$ sont choisies comme dans le lemme
\ref{dga-lemma-good-quotient}. Cela donne une « résolution »
$$
\ldots \to P_2 \xrightarrow{f_2} P_1 \xrightarrow{f_1} P_0 \to M \to 0
$$
Ensuite, nous posons
$$
P = \bigoplus\nolimits_{i \geq 0} P_i
$$
comme $A$-module, avec une graduation
donnée par $P^n = \bigoplus_{a + b = n} P_{-a}^b$
et une différentielle (comme dans la construction du complexe total associé à
un double complexe) par
$$
\text{d}_P(x) = f_{-a}(x) + (-1)^a \text{d}_{P_{-a}}(x)
$$
pour $x \in P_{-a}^b$. Avec ces conventions, $P$ est en effet un $A$-module
différentiel gradué. En rappelant que chaque $P_i$ a une filtration à deux crans $0
\to P_i' \to P_i \to P_i'' \to 0$, nous posons
$$
F_{2i}P = \bigoplus\nolimits_{i \geq j \geq 0} P_j
\subset
\bigoplus\nolimits_{i \geq 0} P_i = P
$$
et nous ajoutons $P'_{i + 1}$ à $F_{2i}P$ pour obtenir $F_{2i +
1}$. Ce sont des sous-modules différentiels gradués et les quotients successifs
sont des sommes directes de décalages de
$A$. D'après le lemme \ref{dga-lemma-target-graded-projective},
nous voyons que les inclusions $F_iP \to F_{i + 1}P$ sont des monomorphismes
admissibles. Enfin, nous devons montrer que l'application $P \to M$
(donnée par l'augmentation $P_0 \to M$) est un quasi-isomorphisme. Cela découle
d'Homologie, lemme \ref{homology-lemma-good-resolution-gives-qis}.
\end{proof}





\section{I-résolutions}
\label{dga-section-I-resolutions}

\noindent
Cette section est le dual de la section sur les P-résolutions.

\medskip\noindent
Soit $(A, \text{d})$ une algèbre différentielle graduée.
Soit $I$ un $A$-module différentiel gradué. Nous disons
que $I$ {\it a la propriété (I)} s'il existe une filtration
$$
I = F_0I \supset F_1I \supset F_2I \supset \ldots \supset 0
$$
par sous-modules différentiels gradués telle que
\begin{enumerate}
\item $I = \lim I/F_pI$,
\item les applications $I/F_{i + 1}I \to I/F_iI$ soient des épimorphismes admissibles,
\item les quotients $F_iI/F_{i + 1}I$ soient isomorphes comme
$A$-modules différentiels gradués à des produits des modules $A^\vee[k]$ construits
à la section \ref{dga-section-modules-noncommutative-differential-graded}.
\end{enumerate}
En fait, la condition (2) est une conséquence de la condition (3),
voir le lemme \ref{dga-lemma-source-graded-injective}. Le lecteur peut vérifier
qu'en tant que module gradué, $I$ sera isomorphe à un produit de $A^\vee[k]$.

\begin{lemma}
\label{dga-lemma-property-I-sequence}
Soit $(A, \text{d})$ une algèbre différentielle graduée.
Soit $I$ un $A$-module différentiel gradué. Si $F_\bullet$ est une
filtration comme dans la propriété (I), alors on obtient
une suite exacte courte admissible
$$
0 \to I \to
\prod\nolimits I/F_iI \to
\prod\nolimits I/F_iI \to 0
$$
de $A$-modules différentiels gradués.
\end{lemma}

\begin{proof}
Démonstration omise. Indication : c'est le dual du lemme \ref{dga-lemma-property-P-sequence}.
\end{proof}

\noindent
Le lemme suivant montre que les modules différentiels gradués avec
la propriété (I) sont l'analogue des modules K-injectifs.
Voir Catégories dérivées, définition \ref{derived-definition-K-injective}.

\begin{lemma}
\label{dga-lemma-property-I-K-injective}
Soit $(A, \text{d})$ une algèbre différentielle graduée.
Soit $I$ un $A$-module différentiel gradué avec la propriété (I).
Alors
$$
\Hom_{K(\text{Mod}_{(A, \text{d})})}(N, I) = 0
$$
pour tous les $A$-modules différentiels gradués acycliques $N$.
\end{lemma}

\begin{proof}
Nous utiliserons que $K(\text{Mod}_{(A, \text{d})})$ est une
catégorie triangulée (proposition \ref{dga-proposition-homotopy-category-triangulated}).
Soit $F_\bullet$ une filtration sur $I$ comme dans la
propriété (I). La suite exacte courte du lemme
\ref{dga-lemma-property-I-sequence} produit un triangle distingué. Ainsi, par
Catégories dérivées, lemme \ref{derived-lemma-representable-homological}, il
suffit de montrer que
$$
\Hom_{K(\text{Mod}_{(A, \text{d})})}(N, I/F_iI) = 0
$$
pour tous les $A$-modules différentiels gradués acycliques $N$ et tous les
$i$. Chacun des modules différentiels gradués $I/F_iI$ possède une filtration finie
par monomorphismes admissibles, dont les quotients successifs
sont des produits de $A^\vee[k]$. Ainsi il suffit de prouver que
$$
\Hom_{K(\text{Mod}_{(A, \text{d})})}(N, A^\vee[k]) = 0
$$
pour tous les $A$-modules différentiels gradués acycliques $N$ et
tous les $k$. Cela découle du lemme
\ref{dga-lemma-hom-into-shift-dual-free} et du fait que $(-)^\vee$ est un foncteur exact.
\end{proof}

\begin{lemma}
\label{dga-lemma-good-sub}
Soit $(A, \text{d})$ une algèbre différentielle
graduée. Soit $M$ un $A$-module différentiel gradué. Il existe un homomorphisme
$M \to I$ de $A$-modules différentiels gradués avec les propriétés
suivantes
\begin{enumerate}
\item $M \to I$ est injectif,
\item $\Coker(\text{d}_M) \to \Coker(\text{d}_I)$ est injectif,
et
\item $I$ s'insère dans une suite exacte courte
admissible $0 \to I' \to I \to I'' \to 0$ où $I'$, $I''$ sont des produits
de décalages de $A^\vee$.
\end{enumerate}
\end{lemma}

\begin{proof}
Nous utiliserons les foncteurs $N \mapsto N^\vee$ (des modules
différentiels gradués à gauche vers ceux à droite, et des modules différentiels
gradués à droite vers ceux à gauche)
construits dans la section
\ref{dga-section-modules-noncommutative-differential-graded} et toutes leurs propriétés.
Pour chaque $k \in \mathbf{Z}$, soit $Q_k$ le $A$-module libre à gauche avec des
générateurs $x, y$ dans les degrés $k$ et $k + 1$. Définissons la structure de $A$-module
différentiel gradué à gauche sur $Q_k$ en posant $\text{d}(x) = y$ et $\text{d}(y) =
0$. En raisonnant exactement comme dans la démonstration du lemme
\ref{dga-lemma-good-quotient}, nous trouvons une surjection
$$
\bigoplus\nolimits_{i \in I} Q_{k_i} \longrightarrow M^\vee
$$
de modules différentiels gradués à gauche sur $A$. Ensuite, nous pouvons considérer l'injection
$$
M \to (M^\vee)^\vee \to (\bigoplus\nolimits_{i \in I} Q_{k_i})^\vee =
\prod\nolimits_{i \in I} I_{k_i}
$$
où $I_k = Q_{-k}^\vee$ est le $A$-module différentiel gradué à droite « dual ».
De plus, la suite exacte courte $0 \to A[-k - 1] \to Q_k \to A[-k] \to 0$
produit une suite exacte courte
$0 \to A^\vee[k] \to I_k \to A^\vee[k + 1] \to 0$.

\medskip\noindent
Le résultat du paragraphe précédent produit $M \to I$ ayant
les propriétés (1) et (3). Pour obtenir la propriété (2), supposons que
$\overline{m} \in \Coker(\text{d}_M)$ soit un élément non nul de degré
$k$. Choisissons une application $\lambda : M^k \to \mathbf{Q}/\mathbf{Z}$
qui s'annule sur $\Im(M^{k - 1} \to M^k)$ mais pas sur $m$. Selon le
lemme \ref{dga-lemma-hom-into-shift-dual-free}, cela
correspond à un homomorphisme $M \to A^\vee[k]$ de $A$-modules
différentiels gradués qui ne s'annule pas sur $m$. Par conséquent,
nous pouvons obtenir (2) en ajoutant un
produit de copies de décalages de $A^\vee$.
\end{proof}

\begin{lemma}
\label{dga-lemma-right-resolution}
Soit $(A, \text{d})$ une algèbre différentielle
graduée. Soit $M$ un $A$-module différentiel gradué. Il existe un homomorphisme
$M \to I$ de $A$-modules différentiels gradués tel que
\begin{enumerate}
\item $M \to I$ soit un quasi-isomorphisme, et
\item $I$ possède la propriété (I).
\end{enumerate}
\end{lemma}

\begin{proof}
Posons $M = M_0$. Nous choisissons par récurrence des suites exactes courtes
$$
0 \to M_i \to I_i \to M_{i + 1} \to 0
$$
où les applications $M_i \to I_i$ sont choisies comme dans le lemme
\ref{dga-lemma-good-sub}. Cela donne une « résolution »
$$
0 \to M \to I_0 \xrightarrow{f_0} I_1 \xrightarrow{f_1} I_1 \to \ldots
$$
Notons $I$ le $A$-module différentiel gradué dont les composantes graduées sont
$$
I^n = \prod\nolimits_{i \geq 0} I^{n - i}_i
$$
et dont la différentielle est définie par
$$
\text{d}_I(x) = f_i(x) + (-1)^i \text{d}_{I_i}(x)
$$
pour $x \in I_i^{n - i}$. Avec ces conventions, $I$ est en effet un $A$-module
différentiel gradué. En rappelant que chaque $I_i$ possède une filtration à deux crans
$0 \to I_i' \to I_i \to I_i'' \to 0$, nous posons
$$
F_{2i}I^n = \prod\nolimits_{j \geq i} I^{n - j}_j
\subset
\prod\nolimits_{i \geq 0} I^{n - i}_i = I^n
$$
et nous ajoutons un facteur $I'_{i + 1}$ à $F_{2i}I$ pour obtenir $F_{2i +
1}I$. Ce sont des sous-modules différentiels gradués et les quotients
successifs sont des produits de décalages de
$A^\vee$. D'après le lemme \ref{dga-lemma-source-graded-injective}, nous
voyons que les inclusions $F_{i + 1}I \to F_iI$ sont des monomorphismes
admissibles. Enfin, nous devons montrer que l'application $M \to I$
(donnée par l'augmentation $M \to I_0$) est un quasi-isomorphisme. Cela découle
d'Homologie, lemme \ref{homology-lemma-good-right-resolution-gives-qis}.
\end{proof}






\section{La catégorie dérivée}
\label{dga-section-derived}

\noindent
Rappelons que les notions de modules différentiels gradués
acycliques et de quasi-isomorphisme de modules différentiels gradués ont un
sens (voir section \ref{dga-section-modules}).

\begin{lemma}
\label{dga-lemma-acyclic}
Soit $(A, \text{d})$ une algèbre différentielle
graduée. La sous-catégorie pleine $\text{Ac}$ de $K(\text{Mod}_{(A,
\text{d})})$ composée de modules acycliques est une sous-catégorie triangulée
strictement pleine et saturée de $K(\text{Mod}_{(A,
\text{d})})$. Le système multiplicatif saturé correspondant
(voir Catégories dérivées, lemme \ref{derived-lemma-operations})
de $K(\text{Mod}_{(A, \text{d})})$ est la classe $\text{Qis}$ des
quasi-isomorphismes. En particulier, le noyau du foncteur de
localisation
$$
Q : K(\text{Mod}_{(A, \text{d})}) \to
\text{Qis}^{-1}K(\text{Mod}_{(A, \text{d})})
$$
est $\text{Ac}$. De plus, le foncteur $H^0$ se factorise par $Q$.
\end{lemma}

\begin{proof}
Nous savons que $H^0$ est un foncteur homologique par la suite
exacte longue d'homologie (\ref{dga-equation-les}). Le
noyau de $H^0$ est la sous-catégorie des objets acycliques et les
flèches qui induisent des isomorphismes sur tous les $H^i$ sont
les quasi-isomorphismes. Ainsi, ce lemme est un cas particulier
de Catégories dérivées, lemme \ref{derived-lemma-acyclic-general}.

\medskip\noindent
Remarque de théorie des ensembles. La construction de la localisation
donnée dans Catégories dérivées,
proposition
\ref{derived-proposition-construct-localization}, exige que la catégorie triangulée donnée soit « petite », c'est-à-dire
que la collection de ses objets forme un ensemble. Soit $V_\alpha$ un
univers partiel (comme dans Ensembles, section \ref{sets-section-sets-hierarchy})
contenant $(A, \text{d})$ et où la cofinalité de $\alpha$ est supérieure à
$\aleph_0$ (voir Ensembles,
proposition \ref{sets-proposition-exist-ordinals-large-cofinality}). Alors nous pouvons
considérer la catégorie $\text{Mod}_{(A, \text{d}), \alpha}$ des $A$-modules
différentiels gradués contenus dans $V_\alpha$. Une vérification simple
montre que toutes les constructions utilisées dans la démonstration de la
proposition \ref{dga-proposition-homotopy-category-triangulated} fonctionnent à
l'intérieur de $\text{Mod}_{(A, \text{d}), \alpha}$ (car au pire
nous prenons des sommes directes finies de modules différentiels gradués). Ainsi, nous
obtenons une catégorie triangulée
$\text{Qis}_\alpha^{-1}K(\text{Mod}_{(A, \text{d}), \alpha})$. Nous verrons
ci-dessous que si $\beta > \alpha$, alors les foncteurs de transition
$$
\text{Qis}_\alpha^{-1}K(\text{Mod}_{(A, \text{d}), \alpha})
\longrightarrow
\text{Qis}_\beta^{-1}K(\text{Mod}_{(A, \text{d}), \beta})
$$
sont pleinement fidèles, puisque les ensembles de morphismes dans les
catégories quotients se calculent à l'aide de morphismes entre P-résolutions dans les
catégories homotopiques (la construction d'une P-résolution dans la démonstration du lemme
\ref{dga-lemma-resolve} prend des sommes directes dénombrables ainsi que des sommes directes
indexées sur des sous-ensembles du module donné). Le lecteur doit donc considérer la
catégorie du lemme comme l'union de ces sous-catégories.
\end{proof}

\noindent
En tenant compte de la remarque de théorie des ensembles à la fin de la
démonstration du lemme précédent, nous définissons la catégorie dérivée comme suit.

\begin{definition}
\label{dga-definition-unbounded-derived-category}
Soit $(A, \text{d})$ une algèbre différentielle graduée.
Soient $\text{Ac}$ et $\text{Qis}$ comme dans le lemme
\ref{dga-lemma-acyclic}. La {\it catégorie dérivée de $(A, \text{d})$} est la catégorie
triangulée
$$
D(A, \text{d}) =
K(\text{Mod}_{(A, \text{d})})/\text{Ac} =
\text{Qis}^{-1}K(\text{Mod}_{(A, \text{d})}).
$$
Nous notons $H^0 : D(A, \text{d}) \to \text{Mod}_R$ le foncteur unique
dont la composition avec le foncteur quotient redonne le foncteur
$H^0$ défini ci-dessus.
\end{definition}

\noindent
Voici le lemme promis calculant les ensembles de morphismes dans
la catégorie dérivée.

\begin{lemma}
\label{dga-lemma-hom-derived}
Soit $(A, \text{d})$ une algèbre différentielle graduée.
Soient $M$ et $N$ des $A$-modules différentiels gradués.
\begin{enumerate}
\item Soit $P \to M$ une P-résolution comme dans
le lemme \ref{dga-lemma-resolve}. Alors
$$
\Hom_{D(A, \text{d})}(M, N) =
\Hom_{K(\text{Mod}_{(A, \text{d})})}(P, N)
$$
\item Soit $N \to I$ une I-résolution comme dans
le lemme \ref{dga-lemma-right-resolution}. Alors
$$
\Hom_{D(A, \text{d})}(M, N) =
\Hom_{K(\text{Mod}_{(A, \text{d})})}(M, I)
$$
\end{enumerate}
\end{lemma}

\begin{proof}
Soit $P \to M$ comme en (1). Puisque $P \to M$ est un quasi-isomorphisme, nous voyons que
$$
\Hom_{D(A, \text{d})}(P, N) = \Hom_{D(A, \text{d})}(M, N)
$$
par définition de la catégorie dérivée. Un
morphisme $f : P \to N$ dans $D(A, \text{d})$ est égal à
$s^{-1}f'$ où $f' : P \to N'$ est un morphisme et où
$s : N \to N'$ est un quasi-isomorphisme. Choisissons un triangle distingué
$$
N \to N' \to Q \to N[1]
$$
Comme $s$ est un quasi-isomorphisme, nous voyons que $Q$ est acyclique.
Ainsi $\Hom_{K(\text{Mod}_{(A, \text{d})})}(P, Q[k]) = 0$ pour tout $k$ par le
lemme \ref{dga-lemma-property-P-K-projective}. Puisque
$\Hom_{K(\text{Mod}_{(A, \text{d})})}(P, -)$ est
cohomologique, nous concluons que nous pouvons relever $f' : P \to N'$ de manière
unique en un morphisme $f : P \to N$. Cela termine la démonstration.

\medskip\noindent
La démonstration de (2) est duale à celle de (1) en
utilisant le lemme \ref{dga-lemma-property-I-K-injective} au
lieu du lemme \ref{dga-lemma-property-P-K-projective}.
\end{proof}

\begin{lemma}
\label{dga-lemma-derived-products}
Soit $(A, \text{d})$ une algèbre différentielle graduée. Alors
\begin{enumerate}
\item $D(A, \text{d})$ a à la fois des sommes directes et des produits,
\item Les sommes directes sont obtenues en prenant des sommes directes de modules
différentiels gradués,
\item Les produits sont obtenus en prenant les produits de modules
différentiels gradués.
\end{enumerate}
\end{lemma}

\begin{proof}
Nous utiliserons que $\text{Mod}_{(A, \text{d})}$ est une catégorie abélienne
avec des sommes directes et des produits arbitraires, et qu'ils donnent
lieu à des sommes directes et des produits dans $K(\text{Mod}_{(A,
\text{d})})$. Voir les lemmes \ref{dga-lemma-dgm-abelian} et \ref{dga-lemma-homotopy-direct-sums}.

\medskip\noindent
Soit $M_j$ une famille de $A$-modules différentiels gradués.
Considérons la somme directe graduée $M = \bigoplus M_j$ qui est un
$A$-module différentiel gradué de manière évidente. Pour un
$A$-module différentiel gradué $N$, choisissons un quasi-isomorphisme $N \to
I$ où $I$ est un $A$-module différentiel gradué avec la propriété (I). Voir le
lemme \ref{dga-lemma-right-resolution}. En utilisant
le lemme \ref{dga-lemma-hom-derived}, nous avons
\begin{align*}
\Hom_{D(A, \text{d})}(M, N)
& =
\Hom_{K(A, \text{d})}(M, I) \\ &
=
\prod \Hom_{K(A, \text{d})}(M_j, I) \\ & =
\prod
\Hom_{D(A, \text{d})}(M_j, N)
\end{align*}
d'où l'existence de sommes directes dans $D(A, \text{d})$ comme indiqué
dans la partie (2) du lemme.

\medskip\noindent
Soit $M_j$ une famille de $A$-modules différentiels gradués.
Considérons le produit $M = \prod M_j$ de $A$-modules différentiels gradués. Pour un
$A$-module différentiel gradué $N$, choisissons un quasi-isomorphisme $P \to
N$ où $P$ est un $A$-module différentiel gradué ayant la propriété (P). Voir le
lemme \ref{dga-lemma-resolve}. En utilisant
le lemme \ref{dga-lemma-hom-derived}, nous avons
\begin{align*}
\Hom_{D(A, \text{d})}(N, M)
& =
\Hom_{K(A, \text{d})}(P, M) \\ &
=
\prod \Hom_{K(A, \text{d})}(P, M_j) \\ & =
\prod
\Hom_{D(A, \text{d})}(N, M_j)
\end{align*}
d'où l'existence de sommes directes dans $D(A, \text{d})$ comme indiqué
dans la partie (3) du lemme.
\end{proof}

\begin{remark}
\label{dga-remark-P-resolution}
Soit $R$ un anneau. Soit $(A, \text{d})$ une $R$-algèbre différentielle graduée. En
utilisant des P-résolutions, nous pouvons parfois réduire les énoncés concernant
des objets généraux de $D(A, \text{d})$ à des énoncés concernant $A[k]$. À
savoir, soit $T$ une propriété d'objets de $D(A, \text{d})$ et supposons que
\begin{enumerate}
\item si $K_i$, $i \in I$ est une famille d'objets de $D(A, \text{d})$
et $T(K_i)$ vaut pour tous les $i \in I$, alors $T(\bigoplus K_i)$,
\item si $K \to L \to M \to K[1]$ est un triangle distingué de
$D(A, \text{d})$ et $T$ vaut pour deux, alors $T$
vaut pour le troisième objet, et
\item $T(A[k])$ vaut pour tous les $k \in \mathbf{Z}$.
\end{enumerate}
Alors $T$ vaut pour tous les objets de $D(A, \text{d})$. Cela découle
clairement des lemmes \ref{dga-lemma-property-P-sequence} et \ref{dga-lemma-resolve}.
\end{remark}







\section{Le foncteur delta canonique}
\label{dga-section-canonical-delta-functor}

\noindent
Soit $(A, \text{d})$ une algèbre différentielle
graduée. Considérons le
foncteur $\text{Mod}_{(A, \text{d})} \to K(\text{Mod}_{(A,
\text{d})})$. Ce foncteur n'est {\bf pas} un $\delta$-foncteur en
général. Cependant, il s'avère que le
foncteur $\text{Mod}_{(A, \text{d})} \to D(A, \text{d})$ est
un $\delta$-foncteur. Pour voir cela, nous devons définir
les morphismes $\delta$ associés à une suite exacte courte
$$
0 \to K \xrightarrow{a} L \xrightarrow{b} M \to 0
$$
dans la catégorie abélienne $\text{Mod}_{(A, \text{d})}$.
Considérons le cône $C(a)$ du morphisme $a$. Nous avons $C(a) = L \oplus K$ et
nous définissons $q : C(a) \to M$ via la projection sur $L$ suivie de $b$.
On obtient donc un homomorphisme de $A$-modules différentiels gradués
$$
q : C(a) \longrightarrow M.
$$
Il est clair que $q \circ i = b$ où $i$ est comme
dans la définition
\ref{dga-definition-cone}. Notons que, comme $a$ est injective, le noyau de $q$ est identifié avec le
cône de $\text{id}_K$ qui est acyclique. Ainsi, nous voyons
que $q$ est un quasi-isomorphisme. Selon le
lemme \ref{dga-lemma-the-same-up-to-isomorphisms}, le
triangle
$$
(K, L, C(a), a, i, -p)
$$
est un triangle distingué dans $K(\text{Mod}_{(A, \text{d})})$. Comme
le foncteur de localisation
$K(\text{Mod}_{(A, \text{d})}) \to D(A, \text{d})$ est exact,
nous voyons que $(K, L, C(a), a, i, -p)$ est un triangle distingué
dans $D(A, \text{d})$. Puisque $q$ est un quasi-isomorphisme, nous
voyons que $q$ est un isomorphisme dans $D(A, \text{d})$. Ainsi,
nous déduisons que
$$
(K, L, M, a, b, -p \circ q^{-1})
$$
est un triangle distingué de $D(A,
\text{d})$. Cela suggère le lemme suivant.

\begin{lemma}
\label{dga-lemma-derived-canonical-delta-functor}
Soit $(A, \text{d})$ une algèbre différentielle graduée. Le foncteur
$\text{Mod}_{(A, \text{d})} \to D(A, \text{d})$
défini possède la structure naturelle d'un $\delta$-foncteur, avec
$$
\delta_{K \to L \to M} = - p \circ q^{-1}
$$
avec $p$ et $q$ comme expliqué ci-dessus.
\end{lemma}

\begin{proof}
Nous avons déjà vu que ce choix conduit à un triangle distingué
chaque fois qu'une suite exacte courte de complexes est donnée.
Nous devons montrer la fonctorialité de cette construction,
voir Catégories dérivées, définition \ref{derived-definition-delta-functor}.
Cela découle du lemme \ref{dga-lemma-functorial-cone} avec un peu de
travail. Comparer avec
Catégories dérivées, lemme \ref{derived-lemma-derived-canonical-delta-functor}.
\end{proof}

\begin{lemma}
\label{dga-lemma-homotopy-colimit}
Soit $(A, \text{d})$ une algèbre différentielle graduée.
Soit $M_n$ un système de modules différentiels gradués. Alors la colimite
dérivée $\text{hocolim} M_n$ dans $D(A, \text{d})$ est représentée
par le module différentiel gradué $\colim M_n$.
\end{lemma}

\begin{proof}
Posons $M = \colim M_n$. Nous avons une suite exacte de modules différentiels gradués
$$
0 \to \bigoplus M_n \to \bigoplus M_n \to M \to 0
$$
par Catégories dérivées, lemme \ref{derived-lemma-compute-colimit}
(appliqué aux complexes sous-jacents de groupes abéliens).
Les sommes directes sont des sommes directes dans
$D(\mathcal{A})$ par le lemme
\ref{dga-lemma-derived-products}. Ainsi, le résultat découle de la
définition des colimites dérivées dans
Catégories dérivées, définition
\ref{derived-definition-derived-colimit} et du fait qu'une suite exacte courte de
complexes donne un triangle distingué
(lemme \ref{dga-lemma-derived-canonical-delta-functor}).
\end{proof}








\section{Catégories linéaires}
\label{dga-section-linear}

\noindent
Juste les définitions.

\begin{definition}
\label{dga-definition-linear-category}
Soit $R$ un anneau. Une {\it catégorie $R$-linéaire $\mathcal{A}$} est une
catégorie où chaque ensemble de morphismes est muni d'une structure de $R$-module
et où, pour $x, y, z \in \Ob(\mathcal{A})$, la loi de composition
$$
\Hom_\mathcal{A}(y, z) \times \Hom_\mathcal{A}(x, y)
\longrightarrow
\Hom_\mathcal{A}(x, z)
$$
est $R$-bilinéaire.
\end{definition}

\noindent
Ainsi la composition détermine une application $R$-linéaire
$$
\Hom_\mathcal{A}(y, z) \otimes_R \Hom_\mathcal{A}(x, y)
\longrightarrow
\Hom_\mathcal{A}(x, z)
$$
de $R$-modules. Notons que nous ne supposons pas que les catégories $R$-linéaires
soient additives.

\begin{definition}
\label{dga-definition-functor-linear-categories}
Soit $R$ un anneau. Un {\it foncteur de catégories $R$-linéaires}, ou
un {\it foncteur $R$-linéaire} est un foncteur $F : \mathcal{A} \to
\mathcal{B}$ où pour tous objets $x, y$ de $\mathcal{A}$
l'application $F : \Hom_\mathcal{A}(x, y) \to \Hom_\mathcal{B}(F(x), F(y))$ est
un homomorphisme de $R$-modules.
\end{definition}







\section{Catégories graduées}
\label{dga-section-graded}

\noindent
Juste quelques définitions.

\begin{definition}
\label{dga-definition-graded-category}
Soit $R$ un anneau. Une {\it catégorie graduée $\mathcal{A}$
sur $R$} est une catégorie où chaque ensemble de morphismes est doté de la
structure d'un $R$-module gradué et où pour
$x, y, z \in \Ob(\mathcal{A})$ la composition est $R$-bilinéaire et induit un
homomorphisme
$$
\Hom_\mathcal{A}(y, z) \otimes_R \Hom_\mathcal{A}(x, y)
\longrightarrow
\Hom_\mathcal{A}(x, z)
$$
de modules gradués sur $R$ (c'est-à-dire préservant les degrés).
\end{definition}

\noindent
Dans cette situation, nous notons $\Hom_\mathcal{A}^i(x, y)$ la composante de
degré $i$ de l'objet gradué $\Hom_\mathcal{A}(x, y)$, de sorte que
$$
\Hom_\mathcal{A}(x, y) =
\bigoplus\nolimits_{i \in \mathbf{Z}} \Hom_\mathcal{A}^i(x, y)
$$
est la décomposition en somme directe en parties graduées.

\begin{definition}
\label{dga-definition-functor-graded-categories}
Soit $R$ un anneau. Un {\it foncteur de catégories graduées sur $R$}, ou un
{\it foncteur gradué}
est un foncteur $F : \mathcal{A} \to \mathcal{B}$ où, pour tous les objets
$x, y$ de $\mathcal{A}$,
l'application $F : \Hom_\mathcal{A}(x, y) \to \Hom_\mathcal{A}(F(x), F(y))$ est
un homomorphisme de $R$-modules gradués.
\end{definition}

\noindent
Étant donnée une catégorie graduée, nous nous intéressons
souvent à la catégorie « usuelle » des applications de degré $0$.
Voici une définition formelle.

\begin{definition}
\label{dga-definition-H0-of-graded-category}
Soit $R$ un anneau. Soit $\mathcal{A}$ une catégorie graduée
sur $R$. On considère {\it $\mathcal{A}^0$} comme la catégorie ayant
les mêmes objets que $\mathcal{A}$ et avec
$$
\Hom_{\mathcal{A}^0}(x, y) = \Hom^0_\mathcal{A}(x, y)
$$
la composante de degré $0$ du module gradué des morphismes de
$\mathcal{A}$.
\end{definition}

\begin{definition}
\label{dga-definition-graded-direct-sum}
Soit $R$ un anneau. Soit $\mathcal{A}$ une catégorie graduée sur $R$. Une
somme directe $(x, y, z, i, j, p, q)$ dans $\mathcal{A}$ (notation comme
dans Homologie, remarque \ref{homology-remark-direct-sum})
est une {\it somme directe graduée} si $i, j, p, q$ sont homogènes
de degré $0$.
\end{definition}

\begin{example}[Catégorie graduée des objets gradués]
\label{dga-example-graded-category-graded-objects}
Soit $\mathcal{B}$ une catégorie additive. Rappelons que nous avons défini
la catégorie $\text{Gr}(\mathcal{B})$ des objets gradués de $\mathcal{B}$ en
Homologie, définition \ref{homology-definition-graded}. Dans cet
exemple, nous allons construire une catégorie graduée
$\text{Gr}^{gr}(\mathcal{B})$ sur $R = \mathbf{Z}$ dont la
catégorie associée $\text{Gr}^{gr}(\mathcal{B})^0$ redonne
$\text{Gr}(\mathcal{B})$. Comme objets de
$\text{Gr}^{gr}(\mathcal{B})$, nous prenons des
objets gradués de $\mathcal{B}$. Alors, étant donnés les objets gradués $A =
(A^i)$ et $B = (B^i)$ de $\mathcal{B}$ nous posons
$$
\Hom_{\text{Gr}^{gr}(\mathcal{B})}(A, B) =
\bigoplus\nolimits_{n \in \mathbf{Z}} \Hom^n(A, B)
$$
où la composante graduée de degré $n$ est le groupe abélien des applications
homogènes de degré $n$ de $A$ vers $B$. Explicitement, nous avons
$$
\Hom^n(A, B) = \prod\nolimits_{p + q = n} \Hom_\mathcal{B}(A^{-q}, B^p)
$$
(observez l'inversion des indices et observez que nous avons un produit ici et non
une somme directe). En d'autres termes, un morphisme de degré $n$
$f$ de $A$ vers $B$ peut être vu comme un système $f = (f_{p, q})$ où
$p, q \in \mathbf{Z}$, $p + q = n$, et où $f_{p, q} :
A^{-q} \to B^p$ est un morphisme de $\mathcal{B}$. Étant donnés
des objets gradués $A$, $B$, $C$ de $\mathcal{B}$, la
composition des morphismes dans $\text{Gr}^{gr}(\mathcal{B})$ est définie via les
applications
$$
\Hom^m(B, C) \times \Hom^n(A, B) \longrightarrow \Hom^{n + m}(A, C)
$$
par simple composition $(g, f) \mapsto g \circ f$ de morphismes
homogènes d'objets gradués. En termes de composantes, nous avons
$$
(g \circ f)_{p, r} = g_{p, q} \circ f_{-q, r}
$$
où $q$ est tel que $p + q = m$ et $-q + r = n$.
\end{example}

\begin{example}[Catégorie graduée des modules gradués]
\label{dga-example-gm-gr-cat}
Soit $A$ une algèbre graduée par $\mathbf{Z}$ sur un anneau $R$. Nous allons
construire une catégorie graduée $\text{Mod}^{gr}_A$ sur $R$ dont la catégorie associée
$(\text{Mod}^{gr}_A)^0$ est la catégorie des $A$-modules gradués. Comme objets de
$\text{Mod}^{gr}_A$, nous prenons des $A$-modules droits gradués (voir section
\ref{dga-section-projectives-over-algebras}). Étant donnés des $A$-modules
gradués $L$ et $M$, nous posons
$$
\Hom_{\text{Mod}^{gr}_A}(L, M) =
\bigoplus\nolimits_{n \in \mathbf{Z}} \Hom^n(L, M)
$$
où $\Hom^n(L, M)$ est l'ensemble des applications de $A$-modules à droite $L \to
M$ qui sont homogènes de degré $n$, c'est-à-dire $f(L^i) \subset M^{i + n}$
pour tout $i \in \mathbf{Z}$. En termes de composantes, nous avons que
$$
\Hom^n(L, M) \subset \prod\nolimits_{p + q = n} \Hom_R(L^{-q}, M^p)
$$
(observer l'inversion des indices) est le sous-ensemble constitué de
ceux $f = (f_{p, q})$ tels que
$$
f_{p, q}(m a) = f_{p - i, q + i}(m)a
$$
pour $a \in A^i$ et $m \in L^{-q - i}$. Pour les $A$-modules gradués $K$,
$L$, $M$, nous définissons la composition dans $\text{Mod}^{gr}_A$ via les
applications
$$
\Hom^m(L, M) \times \Hom^n(K, L) \longrightarrow \Hom^{n + m}(K, M)
$$
par simple composition d'applications de modules à droite sur $A$ : $(g, f) \mapsto g \circ f$.
\end{example}

\begin{remark}
\label{dga-remark-graded-shift-functors}
Soit $R$ un anneau. Soit $\mathcal{D}$ une catégorie $R$-linéaire dotée d'une
collection de foncteurs $R$-linéaires $[n] : \mathcal{D} \to \mathcal{D}$, $x
\mapsto x[n]$ indexés par $n \in \mathbf{Z}$ tels que $[n]
\circ [m] = [n + m]$ et $[0] = \text{id}_\mathcal{D}$ (égalité comme
foncteurs). Cela nous permet de construire une catégorie graduée $\mathcal{D}^{gr}$ sur
$R$ avec les mêmes objets que $\mathcal{D}$ en posant
$$
\Hom_{\mathcal{D}^{gr}}(x, y) =
\bigoplus\nolimits_{n \in \mathbf{Z}} \Hom_\mathcal{D}(x, y[n])
$$
pour $x, y$ dans $\mathcal{D}$. Observez que $(\mathcal{D}^{gr})^0 = \mathcal{D}$
(voir la définition \ref{dga-definition-H0-of-graded-category}). De plus, la catégorie
graduée $\mathcal{D}^{gr}$ hérite de foncteurs $R$-linéaires gradués
$[n]$ satisfaisant $[n] \circ [m] = [n + m]$ et $[0] = \text{id}_{\mathcal{D}^{gr}}$
avec la propriété que
$$
\Hom_{\mathcal{D}^{gr}}(x, y[n]) = \Hom_{\mathcal{D}^{gr}}(x, y)[n]
$$
comme $R$-modules gradués, de manière compatible avec la composition des morphismes.

\medskip\noindent
Inversement, supposons donnée une catégorie graduée $\mathcal{A}$ sur $R$
dotée d'une collection de foncteurs gradués $R$-linéaires
$[n]$ satisfaisant $[n] \circ [m] = [n + m]$ et $[0] =
\text{id}_\mathcal{A}$, qui sont de plus équipés d'isomorphismes
$$
\Hom_\mathcal{A}(x, y[n]) = \Hom_\mathcal{A}(x, y)[n]
$$
comme $R$-modules gradués, de manière compatible avec la composition des
morphismes. Ensuite, le lecteur montre facilement que $\mathcal{A} = (\mathcal{A}^0)^{gr}$.

\medskip\noindent
Voici deux exemples de la relation
$\mathcal{D} \leftrightarrow \mathcal{A}$ que nous avons établie ci-dessus :
\begin{enumerate}
\item Soit $\mathcal{B}$ une catégorie additive. Si
$\mathcal{D} = \text{Gr}(\mathcal{B})$, alors
$\mathcal{A} = \text{Gr}^{gr}(\mathcal{B})$ comme dans
l'exemple \ref{dga-example-graded-category-graded-objects}.
\item Si $A$ est un anneau gradué et $\mathcal{D} = \text{Mod}_A$
est la catégorie des modules gradués à droite sur $A$,
alors $\mathcal{A} = \text{Mod}^{gr}_A$, voir exemple \ref{dga-example-gm-gr-cat}.
\end{enumerate}
\end{remark}






\section{Catégories différentielles graduées}
\label{dga-section-dga-categories}

\noindent
Notons que si $R$ est un anneau, alors $R$ est une algèbre
différentielle graduée sur lui-même (avec $R = R^0$ bien sûr). Dans ce cas, un
$R$-module différentiel gradué est la même chose qu'un complexe de
$R$-modules. En particulier, étant donnés deux $R$-modules différentiels
gradués $M$ et $N$, nous notons $M \otimes_R N$ le $R$-module
différentiel gradué correspondant au complexe total associé au
double complexe obtenu par le produit tensoriel des complexes de $R$-modules
associés à $M$ et $N$.

\begin{definition}
\label{dga-definition-dga-category}
Soit $R$ un anneau. Une {\it catégorie différentielle graduée
$\mathcal{A}$ sur $R$} est une catégorie où chaque ensemble de morphismes est doté de la
structure d'un $R$-module différentiel gradué et où pour
$x, y, z \in \Ob(\mathcal{A})$ la composition est $R$-bilinéaire et induit un
homomorphisme
$$
\Hom_\mathcal{A}(y, z) \otimes_R \Hom_\mathcal{A}(x, y)
\longrightarrow
\Hom_\mathcal{A}(x, z)
$$
de $R$-modules différentiels gradués.
\end{definition}

\noindent
La condition finale de la définition signifie ce qui suit :
si $f \in \Hom_\mathcal{A}^n(x, y)$ et
$g \in \Hom_\mathcal{A}^m(y, z)$ sont homogènes
de degrés $n$ et $m$, alors
$$
\text{d}(g \circ f) = \text{d}(g) \circ f + (-1)^mg \circ \text{d}(f)
$$
dans $\Hom_\mathcal{A}^{n + m + 1}(x, z)$. Cela découle de la règle du
signe pour la différentielle sur le complexe total d'un double complexe,
voir Homologie, définition \ref{homology-definition-associated-simple-complex}.

\begin{definition}
\label{dga-definition-functor-dga-categories}
Soit $R$ un anneau. Un {\it foncteur de catégories différentielles graduées sur $R$}
est un foncteur $F : \mathcal{A} \to \mathcal{B}$ où pour tous les objets
$x, y$ de $\mathcal{A}$,
l'application $F : \Hom_\mathcal{A}(x, y) \to \Hom_\mathcal{A}(F(x), F(y))$
est un homomorphisme de $R$-modules différentiels gradués.
\end{definition}

\noindent
Étant donnée une catégorie différentielle graduée, nous nous intéressons souvent
aux catégories correspondantes de complexes et à la catégorie
homotopique. Voici une définition formelle.

\begin{definition}
\label{dga-definition-homotopy-category-of-dga-category}
Soit $R$ un anneau. Soit $\mathcal{A}$ une catégorie différentielle
graduée sur $R$. On définit
\begin{enumerate}
\item la {\it catégorie des complexes de $\mathcal{A}$}\footnote{Cette
terminologie n'est peut-être pas usuelle.} comme la
catégorie $\text{Comp}(\mathcal{A})$ dont les objets sont les mêmes que les
objets de $\mathcal{A}$ et avec
$$
\Hom_{\text{Comp}(\mathcal{A})}(x, y) =
\Ker(d : \Hom^0_\mathcal{A}(x, y) \to \Hom^1_\mathcal{A}(x, y))
$$
\item la {\it catégorie homotopique de $\mathcal{A}$} comme la
catégorie $K(\mathcal{A})$ dont les objets sont les mêmes que les
objets de $\mathcal{A}$ et avec
$$
\Hom_{K(\mathcal{A})}(x, y) = H^0(\Hom_\mathcal{A}(x, y))
$$
\end{enumerate}
\end{definition}

\noindent
Notre utilisation du symbole $K(\mathcal{A})$ est non standard, mais au moins
elle est compatible avec l'utilisation de $K(-)$ dans d'autres chapitres du projet Champs.

\begin{definition}
\label{dga-definition-dg-direct-sum}
Soit $R$ un anneau. Soit $\mathcal{A}$ une catégorie différentielle graduée sur
$R$. Une somme directe $(x, y, z, i, j, p, q)$ dans $\mathcal{A}$ (notation comme dans
Homologie, remarque \ref{homology-remark-direct-sum})
est une {\it somme directe différentielle graduée} si $i, j, p, q$ sont homogènes
de degré $0$ et fermés, c'est-à-dire $\text{d}(i) = 0$, etc.
\end{definition}

\begin{lemma}
\label{dga-lemma-functorial}
Soit $R$ un anneau. Un foncteur $F : \mathcal{A} \to \mathcal{B}$
de catégories différentielles graduées sur $R$ induit des
foncteurs $\text{Comp}(\mathcal{A}) \to
\text{Comp}(\mathcal{B})$ et $K(\mathcal{A}) \to K(\mathcal{B})$.
\end{lemma}

\begin{proof}
Démonstration omise.
\end{proof}

\begin{example}[Catégorie différentielle graduée des complexes]
\label{dga-example-category-complexes}
Soit $\mathcal{B}$ une catégorie additive. Nous construirons une
catégorie différentielle graduée $\text{Comp}^{dg}(\mathcal{B})$ sur $R
= \mathbf{Z}$ dont la catégorie associée de complexes est
$\text{Comp}(\mathcal{B})$ et dont la catégorie homotopique associée est
$K(\mathcal{B})$. Comme objets de $\text{Comp}^{dg}(\mathcal{B})$, nous prenons des
complexes de $\mathcal{B}$. Étant donnés les complexes
$A^\bullet$ et $B^\bullet$ de $\mathcal{B}$, nous désignons parfois également par
$A^\bullet$ et $B^\bullet$ les objets gradués correspondants de
$\mathcal{B}$ (c'est-à-dire en oubliant la différentielle). En utilisant
cet abus de notation, nous posons
$$
\Hom_{\text{Comp}^{dg}(\mathcal{B})}(A^\bullet, B^\bullet) =
\Hom_{\text{Gr}^{gr}(\mathcal{B})}(A^\bullet, B^\bullet) =
\bigoplus\nolimits_{n \in \mathbf{Z}} \Hom^n(A, B)
$$
comme $\mathbf{Z}$-module gradué, avec les notations et définitions
de l'exemple \ref{dga-example-graded-category-graded-objects}. En
d'autres termes, la composante de degré $n$
est le groupe abélien des morphismes homogènes de degré $n$ d'objets gradués
$$
\Hom^n(A^\bullet, B^\bullet) =
\prod\nolimits_{p + q = n} \Hom_\mathcal{B}(A^{-q}, B^p)
$$
Notons l'inversion des indices et remarquons que nous avons un produit
direct et non une somme directe. Pour un
élément $f \in \Hom^n(A^\bullet, B^\bullet)$ de degré $n$, nous posons
$$
\text{d}(f) = \text{d}_B \circ f - (-1)^n f \circ \text{d}_A
$$
Le signe est exactement
comme dans Plus sur l'algèbre, section
\ref{more-algebra-section-sign-rules}. Pour comprendre cela, nous pensons à $\text{d}_B$ et $\text{d}_A$
comme des morphismes d'objets gradués de $\mathcal{B}$ homogènes de degré $1$ et
nous utilisons la composition dans la catégorie $\text{Gr}^{gr}(\mathcal{B})$
dans le membre de droite. En termes de composantes, si $f = (f_{p, q})$ avec
$f_{p, q} : A^{-q} \to B^p$, nous avons
\begin{equation}
\label{dga-equation-differential-hom-complex}
\text{d}(f_{p, q}) =
\text{d}_B \circ f_{p, q} - (-1)^{p + q} f_{p, q} \circ \text{d}_A
\end{equation}
Notons que le premier terme de cette expression se trouve
dans $\Hom_\mathcal{B}(A^{-q}, B^{p + 1})$ et que le deuxième terme se
trouve dans $\Hom_\mathcal{B}(A^{-q - 1}, B^p)$. Le lecteur vérifie que
\begin{enumerate}
\item $\text{d}$ est de carré nul,
\item un élément $f$ dans $\Hom^n(A^\bullet, B^\bullet)$
a $\text{d}(f) = 0$ si et seulement si le morphisme
$f : A^\bullet \to B^\bullet[n]$ d'objets gradués de $\mathcal{B}$ est en
fait un morphisme de complexes,
\item en particulier, la catégorie des complexes
de $\text{Comp}^{dg}(\mathcal{B})$ est égale à $\text{Comp}(\mathcal{B})$,
\item le morphisme de complexes défini par $f$ comme dans
(2) est homotope à zéro si et seulement si $f = \text{d}(g)$ pour
un certain $g \in \Hom^{n - 1}(A^\bullet, B^\bullet)$.
\item en particulier, nous obtenons un isomorphisme canonique
$$
\Hom_{K(\mathcal{B})}(A^\bullet, B^\bullet)
\longrightarrow
H^0(\Hom_{\text{Comp}^{dg}(\mathcal{B})}(A^\bullet, B^\bullet))
$$
et la catégorie homotopique de $\text{Comp}^{dg}(\mathcal{B})$ est égale à
$K(\mathcal{B})$.
\end{enumerate}
Étant donnés les complexes $A^\bullet$, $B^\bullet$, $C^\bullet$, nous définissons
la composition
$$
\Hom^m(B^\bullet, C^\bullet) \times \Hom^n(A^\bullet, B^\bullet)
\longrightarrow
\Hom^{n + m}(A^\bullet, C^\bullet)
$$
par composition $(g, f) \mapsto g \circ f$ dans la catégorie graduée
$\text{Gr}^{gr}(\mathcal{B})$, voir
exemple \ref{dga-example-graded-category-graded-objects}. Cela
définit un homomorphisme de modules différentiels gradués
$$
\Hom_{\text{Comp}^{dg}(\mathcal{B})}(B^\bullet, C^\bullet)
\otimes_R
\Hom_{\text{Comp}^{dg}(\mathcal{B})}(A^\bullet, B^\bullet)
\longrightarrow
\Hom_{\text{Comp}^{dg}(\mathcal{B})}(A^\bullet, C^\bullet)
$$
comme requis dans la définition \ref{dga-definition-dga-category}
parce que
\begin{align*}
\text{d}(g \circ f) & =
\text{d}_C \circ g \circ f - (-1)^{n + m} g \circ f \circ \text{d}_A \\
& =
\left(\text{d}_C \circ g - (-1)^m g \circ \text{d}_B\right) \circ f +
(-1)^m g \circ \left(\text{d}_B \circ f - (-1)^n f \circ \text{d}_A\right) \\ & =
\text{d}(g)
\circ f + (-1)^m g \circ \text{d}(f)
\end{align*}
comme voulu.
\end{example}

\begin{lemma}
\label{dga-lemma-additive-functor-induces-dga-functor}
Soit $F : \mathcal{B} \to \mathcal{B}'$ un foncteur additif entre des
catégories additives. Alors $F$ induit un foncteur de catégories
différentielles graduées
$$
F : \text{Comp}^{dg}(\mathcal{B}) \to \text{Comp}^{dg}(\mathcal{B}')
$$
de l'exemple \ref{dga-example-category-complexes}
induisant les foncteurs usuels sur la catégorie des complexes et les
catégories d'homotopie.
\end{lemma}

\begin{proof}
Démonstration omise.
\end{proof}

\begin{example}[Catégorie différentielle graduée des modules différentiels gradués]
\label{dga-example-dgm-dg-cat}
Soit $(A, \text{d})$ une algèbre différentielle graduée sur un anneau $R$. Nous allons
construire une catégorie différentielle graduée $\text{Mod}^{dg}_{(A, \text{d})}$ sur
$R$ dont la catégorie des complexes est $\text{Mod}_{(A, \text{d})}$ et dont la
catégorie homotopique est $K(\text{Mod}_{(A, \text{d})})$. Comme objets de
$\text{Mod}^{dg}_{(A, \text{d})}$, nous prenons les
$A$-modules différentiels gradués. Étant donnés des $A$-modules
différentiels gradués $L$ et $M$, nous définissons
$$
\Hom_{\text{Mod}^{dg}_{(A, \text{d})}}(L, M) =
\Hom_{\text{Mod}^{gr}_A}(L, M) = \bigoplus \Hom^n(L, M)
$$
en tant que $R$-module gradué où le membre de droite est défini comme dans
l'exemple \ref{dga-example-gm-gr-cat}. En d'autres termes, la composante de degré $n$
$\Hom^n(L, M)$ est le $R$-module des morphismes de $A$-modules à droite homogènes
de degré $n$. Pour un élément $f \in \Hom^n(L, M)$, nous posons
$$
\text{d}(f) = \text{d}_M \circ f - (-1)^n f \circ \text{d}_L
$$
Pour comprendre cela, nous considérons $\text{d}_M$ et $\text{d}_L$ comme des
morphismes de $R$-modules gradués et nous utilisons la composition de
morphismes de $R$-modules gradués. Il est clair que $\text{d}(f)$ est homogène de degré
$n + 1$ en tant que morphisme de $R$-modules gradués, et il est $A$-linéaire
parce que
\begin{align*}
\text{d}(f)(xa)
& =
\text{d}_M(f(x) a) - (-1)^n f (\text{d}_L(xa)) \\
& =
\text{d}_M(f(x)) a + (-1)^{\deg(x) + n} f(x) \text{d}(a) -
(-1)^n f(\text{d}_L(x)) a - (-1)^{n + \deg(x)} f(x) \text{d}(a) \\ & =
\text{d}(f)(x) a
\end{align*}
comme souhaité (observez que ce calcul ne fonctionnerait pas sans le signe
dans la définition de notre différentielle sur $\Hom$). Des formules similaires
à celles de l'exemple \ref{dga-example-category-complexes} valent pour la
différentielle de $f$ en termes de composantes. Le
lecteur vérifie (de la même manière que dans
l'exemple \ref{dga-example-category-complexes}) que
\begin{enumerate}
\item $\text{d}$ est de carré nul,
\item un élément $f$ dans $\Hom^n(L, M)$ a $\text{d}(f) = 0$ si et seulement si
$f : L \to M[n]$ est un homomorphisme de $A$-modules différentiels gradués,
\item en particulier, la catégorie des complexes
de $\text{Mod}^{dg}_{(A, \text{d})}$ est $\text{Mod}_{(A, \text{d})}$,
\item l'homomorphisme défini par $f$ comme dans (2) est homotope à zéro
si et seulement si $f = \text{d}(g)$ pour un certain
$g \in \Hom^{n - 1}(L, M)$.
\item en particulier, nous obtenons un isomorphisme canonique
$$
\Hom_{K(\text{Mod}_{(A, \text{d})})}(L, M)
\longrightarrow
H^0(\Hom_{\text{Mod}^{dg}_{(A, \text{d})}}(L, M))
$$
et la catégorie homotopique de $\text{Mod}^{dg}_{(A, \text{d})}$
est $K(\text{Mod}_{(A, \text{d})})$.
\end{enumerate}
Étant donnés des $A$-modules différentiels gradués $K$, $L$, $M$, nous définissons
la composition
$$
\Hom^m(L, M) \times \Hom^n(K, L) \longrightarrow \Hom^{n + m}(K, M)
$$
par composition de morphismes homogènes de $A$-modules à droite $(g, f) \mapsto g \circ f$.
Cela définit un homomorphisme de modules différentiels gradués
$$
\Hom_{\text{Mod}^{dg}_{(A, \text{d})}}(L, M) \otimes_R
\Hom_{\text{Mod}^{dg}_{(A, \text{d})}}(K, L) \longrightarrow
\Hom_{\text{Mod}^{dg}_{(A, \text{d})}}(K, M)
$$
comme requis
dans la définition \ref{dga-definition-dga-category}
parce que
\begin{align*}
\text{d}(g \circ f) & =
\text{d}_M \circ g \circ f - (-1)^{n + m} g \circ f \circ \text{d}_K \\
& =
\left(\text{d}_M \circ g - (-1)^m g \circ \text{d}_L\right) \circ f +
(-1)^m g \circ \left(\text{d}_L \circ f - (-1)^n f \circ \text{d}_K\right) \\ & =
\text{d}(g)
\circ f + (-1)^m g \circ \text{d}(f)
\end{align*}
comme voulu.
\end{example}

\begin{lemma}
\label{dga-lemma-homomorphism-induces-dga-functor}
Soit $\varphi : (A, \text{d}) \to (E, \text{d})$ un homomorphisme
d'algèbres différentielles graduées. Alors $\varphi$ induit un foncteur de catégories
différentielles graduées
$$
F :
\text{Mod}^{dg}_{(E, \text{d})}
\longrightarrow
\text{Mod}^{dg}_{(A, \text{d})}
$$
de l'exemple \ref{dga-example-dgm-dg-cat} induisant des foncteurs de restriction
évidents sur les catégories de modules différentiels gradués et les catégories d'homotopie.
\end{lemma}

\begin{proof}
Démonstration omise.
\end{proof}

\begin{lemma}
\label{dga-lemma-construction}
Soit $R$ un anneau. Soit $\mathcal{A}$ une catégorie différentielle
graduée sur $R$. Soit $x$ un objet de $\mathcal{A}$. Soit
$$
(E, \text{d}) = \Hom_\mathcal{A}(x, x)
$$
l'algèbre différentielle graduée sur $R$ des endomorphismes de $x$.
Nous obtenons un foncteur
$$
\mathcal{A} \longrightarrow \text{Mod}^{dg}_{(E, \text{d})},\quad
y \longmapsto \Hom_\mathcal{A}(x, y)
$$
de catégories différentielles graduées en laissant $E$ agir
sur $\Hom_\mathcal{A}(x, y)$ via la composition dans $\mathcal{A}$.
Ce foncteur induit des foncteurs
$$
\text{Comp}(\mathcal{A}) \to \text{Mod}_{(A, \text{d})}
\quad\text{et}\quad
K(\mathcal{A}) \to K(\text{Mod}_{(A, \text{d})})
$$
par une application du lemme \ref{dga-lemma-functorial}.
\end{lemma}

\begin{proof}
La démonstration est immédiate.
\end{proof}









\section{Construction de catégories triangulées}
\label{dga-section-review}

\noindent
Dans cette section, nous examinons le cadre le plus général auquel s'appliquent
les démonstrations de Catégories dérivées,
proposition
\ref{derived-proposition-homotopy-category-triangulated} et de la proposition \ref{dga-proposition-homotopy-category-triangulated}.

\medskip\noindent
Soit $R$ un anneau. Soit $\mathcal{A}$ une catégorie différentielle graduée sur
$R$. Pour que notre argument fonctionne, nous imposons quelques axiomes sur $\mathcal{A}$ :
\begin{enumerate}
\item[(A)] $\mathcal{A}$ possède un objet nul et des sommes
directes différentielles graduées de deux
objets (comme dans la définition
\ref{dga-definition-dg-direct-sum}). \item[(B)] il existe des foncteurs $[n] : \mathcal{A} \longrightarrow
\mathcal{A}$ de catégories différentielles graduées tels que
$[0] = \text{id}_\mathcal{A}$ et $[n + m] = [n] \circ [m]$ et des
isomorphismes donnés
$$
\Hom_\mathcal{A}(x, y[n]) = \Hom_\mathcal{A}(x, y)[n]
$$
de $R$-modules différentiels gradués compatibles avec la composition.
\end{enumerate}

\noindent
Étant donnée notre catégorie différentielle graduée $\mathcal{A}$, nous disons que
\begin{enumerate}
\item une suite $x \to y \to z$ de morphismes de $\text{Comp}(\mathcal{A})$
est une {\it suite exacte courte admissible} s'il existe un
isomorphisme $y \cong x \oplus z$ dans la catégorie graduée sous-jacente tel
que $x \to z$ et $y \to z$ soient des (co)projections.
\item un morphisme $x\to y$ de $\text{Comp}(\mathcal{A})$ est un
{\it monomorphisme admissible} s'il se prolonge en
une suite exacte courte admissible $x\to y\to z$.
\item un morphisme $y\to z$ de $\text{Comp}(\mathcal{A})$ est un
{\it épimorphisme admissible} s'il se prolonge en
une suite exacte courte admissible $x\to y\to z$.
\end{enumerate}
Le lemme suivant nous indique qu'une suite exacte courte admissible donne un
triangle, à condition que nous ayons les axiomes (A) et (B).

\begin{lemma}
\label{dga-lemma-get-triangle}
Soit $\mathcal{A}$ une catégorie différentielle graduée satisfaisant aux
axiomes (A) et (B). Étant donnée une suite exacte courte admissible $x \to
y \to z$, on obtient (voir la démonstration) un triangle
$$
x \to y \to z \to x[1]
$$
dans $\text{Comp}(\mathcal{A})$ avec la propriété que toute composée de deux flèches
consécutives de $z[-1] \to x \to y \to z \to x[1]$ est nulle dans $K(\mathcal{A})$.
\end{lemma}

\begin{proof}
Choisissons un diagramme
$$
\xymatrix{
x \ar[rr]_1 \ar[rd]_a & & x \\
& y \ar[ru]_\pi \ar[rd]^b & \\ z
\ar[rr]^1 \ar[ru]^s & & z
}
$$
donnant l'isomorphisme des objets gradués $y \cong x \oplus z$ comme dans la
définition d'une suite exacte courte admissible. Voici quelques équations qui sont
valables dans cette situation
\begin{enumerate}
\item $1 = \pi a$ et donc $\text{d}(\pi) a = 0$,
\item $1 = b s$ et donc $b \text{d}(s) = 0$,
\item $1 = a \pi + s b$ et donc $a \text{d}(\pi) + \text{d}(s) b = 0$,
\item $\pi s = 0$ et donc $\text{d}(\pi)s + \pi \text{d}(s) = 0$,
\item $\text{d}(s) = a \pi \text{d}(s)$ parce
que $\text{d}(s) = (a \pi + s b)\text{d}(s)$ et $b\text{d}(s) = 0$,
\item $\text{d}(\pi) = \text{d}(\pi) s b$ parce
que $\text{d}(\pi) = \text{d}(\pi)(a \pi + s b)$ et $\text{d}(\pi)a = 0$,
\item $\text{d}(\pi \text{d}(s)) = 0$ parce que si nous le
postcomposons avec le monomorphisme $a$ nous
obtenons $\text{d}(a\pi \text{d}(s)) = \text{d}(\text{d}(s)) = 0$, et
\item $\text{d}(\text{d}(\pi)s) = 0$ car, d'après (4), c'est
l'opposé de $\text{d}(\pi\text{d}(s))$ qui est $0$ par (7).
\end{enumerate}
Nous avons utilisé à plusieurs reprises que $\text{d}(a) = 0$, $\text{d}(b)
= 0$, et que $\text{d}(1) = 0$. Par (7) nous voyons que
$$
\delta = \pi \text{d}(s) = - \text{d}(\pi) s : z \to x[1]
$$
est un morphisme dans $\text{Comp}(\mathcal{A})$. D'après (5), nous
voyons que la composition $a \delta = a \pi \text{d}(s) =
\text{d}(s)$ est homotope à zéro. D'après (6), nous voyons que la
composition $\delta b = - \text{d}(\pi)sb = \text{d}(-\pi)$ est homotope à zéro.
\end{proof}

\noindent
Outre les axiomes (A) et (B), nous avons besoin d'un axiome concernant
l'existence des cônes. Nous formalisons tout comme suit.

\begin{situation}
\label{dga-situation-ABC}
Ici $R$ est un anneau et $\mathcal{A}$ est une catégorie différentielle graduée
sur $R$ satisfaisant aux axiomes (A), (B), et
\begin{enumerate}
\item[(C)] étant donnée une flèche $f : x \to y$ de degré $0$
avec $\text{d}(f) = 0$, il existe une suite exacte courte admissible
$y \to c(f) \to x[1]$ dans $\text{Comp}(\mathcal{A})$ telle que
l'application $x[1] \to y[1]$ du lemme \ref{dga-lemma-get-triangle} soit égale à $f[1]$.
\end{enumerate}
\end{situation}

\noindent
Nous appellerons $c(f)$ un {\it cône} du morphisme $f$. Si (A),
(B) et (C) sont satisfaits, alors les
cônes sont fonctoriels dans un sens faible.

\begin{lemma}
\label{dga-lemma-cone}
\begin{slogan}
La catégorie homotopique est une catégorie
triangulée. Ce lemme prouve une partie des axiomes d'une catégorie triangulée.
\end{slogan}
Dans la situation \ref{dga-situation-ABC}, supposons que
$$
\xymatrix{
x_1 \ar[r]_{f_1} \ar[d]_a & y_1 \ar[d]^b \\
x_2 \ar[r]^{f_2} & y_2
}
$$
soit un diagramme de $\text{Comp}(\mathcal{A})$ commutatif à homotopie
près. Alors il existe un morphisme $c : c(f_1) \to c(f_2)$ qui donne lieu à un
morphisme de triangles
$$
(a, b, c) : (x_1, y_1, c(f_1)) \to (x_1, y_1, c(f_1))
$$
dans $K(\mathcal{A})$.
\end{lemma}

\begin{proof}
L'hypothèse signifie qu'il existe un morphisme $h : x_1 \to y_2$ de degré
$-1$ tel que $\text{d}(h) = b f_1 - f_2 a$. Choisissons des
isomorphismes $c(f_i) = y_i \oplus x_i[1]$ d'objets gradués compatibles avec les
morphismes $y_i \to c(f_i) \to x_i[1]$. Désignons par
$a_i : y_i \to c(f_i)$, $b_i : c(f_i) \to x_i[1]$, $s_i : x_i[1] \to c(f_i)$ et
$\pi_i : c(f_i) \to y_i$ les morphismes donnés. Rappelons que $x_i[1] \to
y_i[1]$ est donné par $\pi_i \text{d}(s_i)$. Selon l'axiome (C), cela
signifie que
$$
f_i = \pi_i \text{d}(s_i) = - \text{d}(\pi_i) s_i
$$
(nous identifions $\Hom(x_i, y_i)$ avec $\Hom(x_i[1], y_i[1])$ en
utilisant le foncteur de décalage $[1]$).
Posons $c = a_2 b \pi_1 + s_2 a b_1 + a_2hb$. Puis, en utilisant les
égalités trouvées dans la démonstration du lemme \ref{dga-lemma-get-triangle}, nous
obtenons
\begin{align*}
\text{d}(c)
& =
a_2 b \text{d}(\pi_1) + \text{d}(s_2) a b_1 + a_2 \text{d}(h) b_1 \\ &
= -
a_2 b f_1 b_1 + a_2 f_2 a b_1 + a_2 (b f_1 - f_2 a) b_1 \\ & =
0
\end{align*}
(où nous avons utilisé en particulier
que $\text{d}(\pi_1) = \text{d}(\pi_1) s_1 b_1 = f_1 b_1$ et
$\text{d}(s_2) = a_2 \pi_2 \text{d}(s_2) = a_2 f_2$).
Ainsi $c$ est un morphisme de degré $0$ $c : c(f_1) \to c(f_2)$ de $\mathcal{A}$
compatible avec les morphismes donnés $y_i \to c(f_i) \to x_i[1]$.
\end{proof}

\noindent
Dans la situation \ref{dga-situation-ABC}, nous disons qu'un
triangle $(x, y, z, f, g, h)$ dans
$K(\mathcal{A})$ est un {\it triangle distingué} s'il existe une suite
exacte courte admissible $x' \to y' \to z'$ telle que $(x,
y, z, f, g, h)$ soit isomorphe comme triangle dans $K(\mathcal{A})$ au
triangle $(x', y', z', x' \to y', y' \to z', \delta)$ construit
dans le lemme \ref{dga-lemma-get-triangle}. Nous montrerons ci-dessous que
$$
\boxed{
K(\mathcal{A})\text{ est une catégorie triangulée}
}
$$
Ce résultat, bien qu'il ne soit pas aussi général qu'on pourrait le penser,
s'applique à un certain nombre de généralisations naturelles des cas couverts jusqu'à présent
dans le projet Champs. Voici quelques exemples :
\begin{enumerate}
\item Soit $(X, \mathcal{O}_X)$ un espace annelé. Soit $(A, d)$ un
faisceau d'algèbres différentielles graduées sur $\mathcal{O}_X$. Soit
$\mathcal{A}$ la catégorie différentielle graduée des modules
différentiels gradués sur $A$. Alors $K(\mathcal{A})$ est une catégorie triangulée.
\item Soit $(\mathcal{C}, \mathcal{O})$ un site annelé. Soit $(A, d)$ un
faisceau d'algèbres différentielles graduées sur $\mathcal{O}$.
Soit $\mathcal{A}$ la catégorie différentielle graduée des $A$-modules
différentiels gradués. Alors $K(\mathcal{A})$ est une catégorie triangulée.
Voir Faisceaux Différentiels Gradués, proposition
\ref{sdga-proposition-homotopy-category-triangulated}.
\item On trouvera deux exemples d'une autre nature dans Exemples, section
\ref{examples-section-nongraded-differential-graded}.
\end{enumerate}

\noindent
Le lemme élémentaire suivant joue un rôle essentiel dans la construction.

\begin{lemma}
\label{dga-lemma-id-cone-null}
Dans la situation
\ref{dga-situation-ABC}, étant donné tout objet $x$ de $\mathcal{A}$, et le cône $C(1_x)$
du morphisme identité $1_x : x \to x$, le morphisme identité sur
$C(1_x)$ est homotope à zéro.
\end{lemma}

\begin{proof}
Considérons la suite exacte courte admissible donnée par l'axiome (C).
$$
\xymatrix{
x \ar@<0.5ex>[r]^a &
C(1_x) \ar@<0.5ex>[l]^{\pi} \ar@<0.5ex>[r]^b &
x[1]\ar@<0.5ex>[l]^s
}
$$
Alors, d'après le lemme \ref{dga-lemma-get-triangle}, en identifiant les ensembles
de morphismes par décalage, nous avons $1_x=\pi d(s)=-d(\pi)s$ où $s$ est
considéré comme un morphisme dans $\Hom_{\mathcal{A}}^{-1}(x,C(1_x))$. Par
conséquent, $a=a\pi d(s)=d(s)$ en utilisant la formule (5) du lemme
\ref{dga-lemma-get-triangle}, et $b=-d(\pi)sb=-d(\pi)$ par la formule (6) du lemme \ref{dga-lemma-get-triangle}.
Ainsi
$$
1_{C(1_x)} = a\pi + sb = d(s)\pi - sd(\pi) = d(s\pi)
$$
puisque $s$ est de degré $-1$.
\end{proof}

\noindent
Une version plus générale du lemme ci-dessus apparaîtra
dans le lemme \ref{dga-lemma-cone-homotopy}. Le lemme suivant est
l'analogue du lemme \ref{dga-lemma-make-commute-map}.

\begin{lemma}
\label{dga-lemma-homo-change}
Dans la situation \ref{dga-situation-ABC}, soit un diagramme
$$
\xymatrix{x\ar[r]^f\ar[d]_a & y\ar[d]^b\\
z\ar[r]^g & w}
$$
dans $\text{Comp}(\mathcal{A})$ commutatif à homotopie près. Alors
\begin{enumerate}
\item Si $f$ est un monomorphisme admissible, alors $b$ est
homotope à un morphisme $b'$ qui fait commuter le diagramme.
\item Si $g$ est un épimorphisme admissible, alors $a$ est
homotope à un morphisme $a'$ qui fait commuter le diagramme.
\end{enumerate}
\end{lemma}

\begin{proof}
Pour prouver (1), on observe que l'hypothèse implique qu'il existe un
$h\in\Hom_{\mathcal{A}}(x,w)$ de degré $-1$ tel que $bf-ga=d(h)$.
Puisque $f$ est un monomorphisme admissible, il existe un morphisme
$\pi : y \to x$ dans la catégorie $\mathcal{A}$ de degré $0$.
Soit $b' = b - d(h\pi)$. Alors
\begin{align*}
b'f = bf - d(h\pi)f
= &
bf - d(h\pi f) \quad (\text{puisque }d(f) = 0) \\ = &
bf-d(h)
\\ =
&
ga
\end{align*}
comme voulu. La démonstration pour (2) est omise.
\end{proof}

\noindent
Le lemme suivant est l'analogue du lemme \ref{dga-lemma-make-injective}.

\begin{lemma}
\label{dga-lemma-factor}
Dans la situation \ref{dga-situation-ABC}, soit $\alpha : x \to
y$ un morphisme dans $\text{Comp}(\mathcal{A})$. Alors il existe
une factorisation dans $\text{Comp}(\mathcal{A})$ :
$$
\xymatrix{
x \ar[r]^{\tilde{\alpha}} &
\tilde{y} \ar@<0.5ex>[r]^{\pi} &
y\ar@<0.5ex>[l]^s
}
$$
avec les propriétés suivantes :
\begin{enumerate}
\item $\tilde{\alpha}$ est un monomorphisme admissible, et
$\pi\tilde{\alpha}=\alpha$.
\item Il existe un morphisme
$s:y\to\tilde{y}$ dans $\text{Comp}(\mathcal{A})$
tel que $\pi s=1_y$ et que $s\pi$ soit homotope à $1_{\tilde{y}}$.
\end{enumerate}
\end{lemma}

\begin{proof}
D'après l'axiome (A), nous pouvons prendre pour $\tilde{y}$ la somme directe
différentielle graduée de $y$ et $C(1_x)$ ; il existe donc un diagramme
$$
\xymatrix@C=3pc{
y \ar@<0.5ex>[r]^s &
y\oplus C(1_x) \ar@<0.5ex>[l]^{\pi} \ar@<0.5ex>[r]^{p} &
C(1_x)\ar@<0.5ex>[l]^t
}
$$
où tous les morphismes sont de degré zéro, et dans
$\text{Comp}(\mathcal{A})$. Soit $\tilde{y} = y \oplus C(1_x)$. Alors
$1_{\tilde{y}} = s\pi + tp$. Considérons maintenant le diagramme
$$
\xymatrix{
x \ar[r]^{\tilde{\alpha}} &
\tilde{y} \ar@<0.5ex>[r]^{\pi} &
y\ar@<0.5ex>[l]^s
}
$$
où $\tilde{\alpha}$ est induit par le morphisme $x\xrightarrow{\alpha}y$
et le morphisme naturel $x\to C(1_x)$ s'insérant dans la suite
exacte courte admissible
$$
\xymatrix{
x \ar@<0.5ex>[r] &
C(1_x) \ar@<0.5ex>[l] \ar@<0.5ex>[r] &
x[1]\ar@<0.5ex>[l]
}
$$
Ainsi, le morphisme $C(1_x)\to x$ de degré 0 dans ce
diagramme, ainsi que le morphisme nul $y\to x$, induisent un morphisme de
degré 0 $\beta : \tilde{y} \to x$. Ensuite, $\tilde{\alpha}$ est un
monomorphisme admissible puisqu'il s'intègre dans la suite exacte
courte admissible
$$
\xymatrix{
x\ar[r]^{\tilde{\alpha}} &
\tilde{y} \ar[r] &
x[1]
}
$$
De plus, $\pi\tilde{\alpha} = \alpha$ par la construction de
$\tilde{\alpha}$, et $\pi s = 1_y$ par le premier diagramme. Il reste à montrer
que $s\pi$ est homotope à $1_{\tilde{y}}$. Écrivons $1_x$ sous la forme
$d(h)$ pour un certain morphisme de degré $-1$. Ensuite, notre
dernière affirmation découle de
\begin{align*}
1_{\tilde{y}} - s\pi =
& tp
\\ = &
t(dh)p\quad\text{(d'après
le lemme \ref{dga-lemma-id-cone-null})} \\ =
&
d(thp)
\end{align*}
puisque $dt = dp = 0$, et $t$ est de degré zéro.
\end{proof}

\noindent
Le lemme suivant est l'analogue du lemme \ref{dga-lemma-sequence-maps-split}.

\begin{lemma}
\label{dga-lemma-analogue-sequence-maps-split}
Dans la situation
\ref{dga-situation-ABC}, soit $x_1 \to x_2 \to \ldots \to
x_n$ une suite de morphismes composables dans $\text{Comp}(\mathcal{A})$.
Alors il existe un diagramme commutatif dans $\text{Comp}(\mathcal{A})$ :
$$
\xymatrix{x_1\ar[r] & x_2\ar[r] & \ldots\ar[r] & x_n\\
y_1\ar[r]\ar[u] & y_2\ar[r]\ar[u] & \ldots\ar[r] & y_n\ar[u]}
$$
de telle sorte que chaque $y_i\to y_{i+1}$ soit un monomorphisme
admissible et chaque $y_i\to x_i$ soit une équivalence d'homotopie.
\end{lemma}

\begin{proof}
Le cas de $n=1$ est trivial : on prend simplement $y_1 = x_1$ et le
morphisme identité sur $x_1$ est en particulier une équivalence d'homotopie. Le
cas $n = 2$ est donné par le lemme \ref{dga-lemma-factor}. Supposons que nous
ayons construit le diagramme jusqu'à $x_{n - 1}$. Nous
appliquons le lemme \ref{dga-lemma-factor} à la composition
$y_{n - 1} \to x_{n-1} \to x_n$ pour obtenir $y_n$.
Ensuite, $y_{n - 1} \to y_n$ sera un monomorphisme admissible, et $y_n
\to x_n$ une équivalence d'homotopie.
\end{proof}

\noindent
Le lemme suivant est l'analogue du lemme \ref{dga-lemma-nilpotent}.

\begin{lemma}
\label{dga-lemma-triseq}
Dans la situation \ref{dga-situation-ABC}, soient $x_i \to y_i \to
z_i$ des morphismes dans $\mathcal{A}$ ($i=1,2,3$) tels
que $x_2 \to y_2\to z_2$ soit une suite exacte courte admissible.
Soient $b : y_1 \to y_2$ et $b' : y_2\to y_3$ des morphismes
dans $\text{Comp}(\mathcal{A})$ tels que
$$
\vcenter{
\xymatrix{
x_1 \ar[d]_0 \ar[r] &
y_1 \ar[r] \ar[d]_b &
z_1 \ar[d]_0 \\ x_2
\ar[r] & y_2 \ar[r] & z_2 } }
\quad\text{et}\quad
\vcenter{
\xymatrix{
x_2
\ar[d]^0
\ar[r] & y_2 \ar[r]
\ar[d]^{b'} & z_2 \ar[d]^0
\\ x_3 \ar[r] &
y_3 \ar[r] & z_3
}
}
$$
commutent à homotopie près. Alors $b'\circ b$ est homotope à $0$.
\end{lemma}

\begin{proof}
Par le lemme \ref{dga-lemma-homo-change}, nous pouvons remplacer $b$ et
$b'$ par des morphismes homotopes $\tilde{b}$ et $\tilde{b}'$, de telle sorte que
le carré de droite du diagramme de gauche commute et que le carré de gauche du
diagramme de droite commute. Notons $b = \tilde{b} + d(h)$ et $b'=\tilde{b}'+d(h')$
pour des morphismes de degré $-1$ $h$ et $h'$ dans $\mathcal{A}$. D'où
$$
b'b = \tilde{b}'\tilde{b} + d(\tilde{b}'h + h'\tilde{b} + h'd(h))
$$
puisque $d(\tilde{b})=d(\tilde{b}')=0$, c'est-à-dire que $b'b$ est
homotope à $\tilde{b}'\tilde{b}$. Nous voulons maintenant montrer que
$\tilde{b}'\tilde{b}=0$. Comme $x_2\xrightarrow{f} y_2\xrightarrow{g} z_2$ est une suite
exacte courte admissible, il existe des morphismes de degré
$0$ $\pi : y_2 \to x_2$ et $s : z_2 \to y_2$ tels que
$\text{id}_{y_2} = f\pi + sg$. Par conséquent
$$
\tilde{b}'\tilde{b} = \tilde{b}'(f\pi + sg)\tilde{b} = 0
$$
puisque $g\tilde{b} = 0$ et $\tilde{b}'f = 0$ d'après
les deux carrés commutatifs.
\end{proof}

\noindent
Le lemme suivant est l'analogue du
lemme \ref{dga-lemma-triangle-independent-splittings}.

\begin{lemma}
\label{dga-lemma-analogue-triangle-independent-splittings}
Dans la situation
\ref{dga-situation-ABC}, soit $0 \to x \to y \to z \to 0$ une suite exacte
courte admissible dans $\text{Comp}(\mathcal{A})$. Le triangle
$$
\xymatrix{x\ar[r] & y\ar[r] & z\ar[r]^{\delta} & x[1]}
$$
avec $\delta : z \to x[1]$ tel que défini dans le lemme
\ref{dga-lemma-get-triangle} est, à isomorphisme canonique près dans $K(\mathcal{A})$, indépendant des
choix faits dans le lemme \ref{dga-lemma-get-triangle}.
\end{lemma}

\begin{proof}
Supposons que $\delta$ soit défini par le scindage
$$
\xymatrix{
x \ar@<0.5ex>[r]^{a} &
y \ar@<0.5ex>[r]^b\ar@<0.5ex>[l]^{\pi} & z
\ar@<0.5ex>[l]^s
}
$$
et que $\delta'$ soit défini par le scindage avec $\pi',s'$ à
la place de $\pi,s$. Ensuite
$$
s'-s = (a\pi + sb)(s'-s) = a\pi s'
$$
puisque $bs' = bs = 1_z$ et $\pi s = 0$. De même,
$$
\pi' - \pi = (\pi' - \pi)(a\pi + sb) = \pi'sb
$$
Puisque $\delta = \pi d(s)$ et $\delta' = \pi'd(s')$
d'après la construction du lemme \ref{dga-lemma-get-triangle}, nous pouvons calculer
$$
\delta' = \pi'd(s') = (\pi + \pi'sb)d(s + a\pi s') = \delta + d(\pi s')
$$
en utilisant $\pi a = 1_x$, $ba = 0$ et $\pi'sbd(s') = \pi'sba\pi d(s') = 0$
selon la formule (5) dans le lemme \ref{dga-lemma-get-triangle}.
\end{proof}

\noindent
Le lemme suivant est l'analogue du lemme \ref{dga-lemma-rotate-cone}.

\begin{lemma}
\label{dga-lemma-restate-axiom-c}
Dans la situation
\ref{dga-situation-ABC}, soit $f: x \to y$ un morphisme dans
$\text{Comp}(\mathcal{A})$. Le triangle $(y, c(f), x[1], i, p, f[1])$ est le triangle
associé à la suite exacte courte admissible
$$
\xymatrix{y\ar[r] & c(f) \ar[r] & x[1]}
$$
où le cône $c(f)$ est défini comme dans le lemme \ref{dga-lemma-get-triangle}.
\end{lemma}

\begin{proof}
Cela découle de l'axiome (C).
\end{proof}

\noindent
Le lemme suivant est l'analogue du lemme \ref{dga-lemma-rotate-triangle}.

\begin{lemma}
\label{dga-lemma-cone-rotate-isom}
Dans la situation \ref{dga-situation-ABC}, supposons que $\alpha : x \to y$ et $\beta : y \to
z$ définissent une suite exacte courte admissible
$$
\xymatrix{
x \ar[r] &
y\ar[r] &
z
}
$$
dans $\text{Comp}(\mathcal{A})$. Soit $(x, y, z, \alpha, \beta,
\delta)$ le triangle associé dans $K(\mathcal{A})$. Alors les triangles
$$
(z[-1], x, y, \delta[-1], \alpha, \beta)
\quad\text{et}\quad
(z[-1], x, c(\delta[-1]), \delta[-1], i, p)
$$
sont isomorphes.
\end{lemma}

\begin{proof}
Nous avons un diagramme de la forme
$$
\xymatrix{
z[-1]\ar[r]^{\delta[-1]}\ar[d]^1 &
x\ar@<0.5ex>[r]^{\alpha}\ar[d]^1 &
y\ar@<0.5ex>[r]^{\beta}\ar@{.>}[d]\ar@<0.5ex>[l]^{\tilde{\alpha}} &
z\ar[d]^1\ar@<0.5ex>[l]^{\tilde\beta} \\
z[-1] \ar[r]^{\delta[-1]} &
x\ar@<0.5ex>[r]^i &
c(\delta[-1]) \ar@<0.5ex>[r]^p\ar@<0.5ex>[l]^{\tilde i} &
z\ar@<0.5ex>[l]^{\tilde p}
}
$$
avec les scindages de $\alpha, \beta, i$, et $p$ données
par $\tilde{\alpha}, \tilde{\beta}, \tilde{i},$ et $\tilde{p}$ respectivement.
Définissons un morphisme $y \to c(\delta[-1])$ par
$i\tilde{\alpha} + \tilde{p}\beta$ et un morphisme
$c(\delta[-1]) \to y$ par $\alpha \tilde{i} + \tilde{\beta} p$.
Vérifions d'abord que ceux-ci définissent des morphismes dans $\text{Comp}(\mathcal{A})$.
Nous remarquons que, d'après les identités du lemme
\ref{dga-lemma-get-triangle}, nous avons la
relation $\delta[-1] = \tilde{\alpha}d(\tilde{\beta}) = -d(\tilde{\alpha})\tilde{\beta}$ et
la relation $\delta[-1] = \tilde{i}d(\tilde{p})$. Ensuite
\begin{align*}
d(\tilde{\alpha})
& =
d(\tilde{\alpha})\tilde{\beta}\beta \\ &
=
-\delta[-1]\beta
\end{align*}
où nous avons utilisé l'équation (6)
du lemme \ref{dga-lemma-get-triangle} pour la première égalité et
la remarque précédente pour la seconde. De même, nous obtenons
$d(\tilde{p}) = i\delta[-1]$. Ainsi
\begin{align*}
d(i\tilde{\alpha} + \tilde{p}\beta)
& =
d(i)\tilde{\alpha} + id(\tilde{\alpha}) +
d(\tilde{p})\beta + \tilde{p}d(\beta) \\ & =
id(\tilde{\alpha})
+ d(\tilde{p})\beta \\ & =
-i\delta[-1]\beta
+ i\delta[-1]\beta \\ &
=
0
\end{align*}
donc $i\tilde{\alpha} + \tilde{p}\beta$ est en effet un morphisme
de $\text{Comp}(\mathcal{A})$. Par un calcul similaire,
$\alpha \tilde{i} + \tilde{\beta} p$ est également un morphisme de
$\text{Comp}(\mathcal{A})$. Il est immédiat que ces morphismes s'insèrent
dans le diagramme commutatif. Nous calculons :
\begin{align*}
(i\tilde{\alpha} + \tilde{p}\beta)(\alpha \tilde{i} + \tilde{\beta} p)
& =
i\tilde{\alpha}\alpha\tilde{i} + i\tilde{\alpha}\tilde{\beta}p +
\tilde{p}\beta\alpha\tilde{i} + \tilde{p}\beta\tilde{\beta}p \\ & =
i\tilde{i}
+ \tilde{p}p \\ &
=
1_{c(\delta[-1])}
\end{align*}
où nous avons librement utilisé les identités
du lemme \ref{dga-lemma-get-triangle}. De même, nous calculons
$(\alpha \tilde{i} + \tilde{\beta} p)(i\tilde{\alpha} + \tilde{p}\beta) = 1_y$, donc nous
concluons $y \cong c(\delta[-1])$. Par conséquent, les deux triangles en question
sont isomorphes.
\end{proof}

\noindent
Le lemme suivant est l'analogue du
lemme \ref{dga-lemma-third-isomorphism}.

\begin{lemma}
\label{dga-lemma-analogue-third-isomorphism}
Dans la situation \ref{dga-situation-ABC}, soient $f_1 : x_1 \to y_1$
et $f_2 : x_2 \to y_2$ des morphismes dans $\text{Comp}(\mathcal{A})$. Soit
$$
(a,b,c): (x_1,y_1,c(f_1), f_1, i_1, p_1) \to (x_2,y_2, c(f_2), f_2, i_1, p_1)
$$
un morphisme quelconque de triangles dans $K(\mathcal{A})$. Si
$a$ et $b$ sont des équivalences d'homotopie, alors il en est de même pour $c$.
\end{lemma}

\begin{proof}
Puisque $a$ et $b$ sont des équivalences d'homotopie, elles sont
inversibles dans $K(\mathcal{A})$, notons donc $a^{-1}$ et $b^{-1}$ leurs inverses dans
$K(\mathcal{A})$, ce qui nous donne un diagramme commutatif
$$
\xymatrix{
x_2\ar[d]^{a^{-1}}\ar[r]^{f_2} &
y_2\ar[d]^{b^{-1}}\ar[r]^{i_2} &
c(f_2)\ar[d]^{c'} \\
x_1\ar[r]^{f_1} & y_1
\ar[r]^{i_1} &
c(f_1)
}
$$
où l'application $c'$ est définie via le lemme \ref{dga-lemma-cone} appliqué au carré
commutatif de gauche du diagramme ci-dessus. Puisque le diagramme commute
dans $K(\mathcal{A})$, il suffit, d'après le lemme \ref{dga-lemma-triseq}, de
prouver ce qui suit : étant donné un morphisme de triangles
$(1,1,c): (x,y,c(f),f,i,p)\to (x,y,c(f),f,i,p)$ dans
$K(\mathcal{A})$, l'application $c$ est un isomorphisme dans
$K(\mathcal{A})$. Nous avons les diagrammes commutatifs dans $K(\mathcal{A})$:
$$
\vcenter{
\xymatrix{
y\ar[d]^{1}\ar[r] &
c(f)\ar[d]^{c}\ar[r] &
x[1]\ar[d]^{1} \\
y\ar[r] & c(f)
\ar[r] & x[1] }
}
\quad\Rightarrow\quad
\vcenter{
\xymatrix{
y\ar[d]^{0}\ar[r]
&
c(f)\ar[d]^{c-1}\ar[r]
& x[1]\ar[d]^{0}
\\ y\ar[r] & c(f)
\ar[r]
&
x[1]
}
}
$$
Comme les lignes sont des suites exactes courtes admissibles, nous
obtenons l'identité $(c-1)^2 = 0$ par le lemme \ref{dga-lemma-triseq},
dont nous concluons que $2-c$ est l'inverse de $c$ dans $K(\mathcal{A})$ de
sorte que $c$ est un isomorphisme.
\end{proof}

\noindent
Le lemme suivant est l'analogue du
lemme \ref{dga-lemma-the-same-up-to-isomorphisms}.

\begin{lemma}
\label{dga-lemma-cone-homotopy}
Dans la situation \ref{dga-situation-ABC}.
\begin{enumerate}
\item Soit une suite exacte courte admissible
$x\xrightarrow{\alpha} y\xrightarrow{\beta} z$. Il
existe alors une équivalence d'homotopie
$e:C(\alpha)\to z$ telle que le diagramme
\begin{equation}
\label{dga-equation-cone-isom-triangle}
\vcenter{
\xymatrix{
x\ar[r]^{\alpha}\ar[d] &
y\ar[r]^{b}\ar[d] &
C(\alpha)\ar[r]^{-c}\ar@{.>}[d]^{e} &
x[1]\ar[d] \\
x\ar[r]^{\alpha} &
y\ar[r]^{\beta} &
z\ar[r]^{\delta} & x[1]
}
}
\end{equation}
définit un isomorphisme de triangles dans $K(\mathcal{A})$. Ici,
$y\xrightarrow{b}C(\alpha)\xrightarrow{c}x[1]$ est
la suite exacte courte admissible donnée comme dans l'axiome (C).
\item Étant donné un
morphisme $\alpha : x \to y$ dans
$\text{Comp}(\mathcal{A})$, soit $x \xrightarrow{\tilde{\alpha}} \tilde{y} \to y$
la factorisation donnée comme dans le lemme \ref{dga-lemma-factor}, où le
monomorphisme admissible $x \xrightarrow{\tilde{\alpha}} y$ s'étend à la
suite exacte courte admissible
$$
\xymatrix{
x \ar[r]^{\tilde{\alpha}} &
\tilde{y} \ar[r] & z
}
$$
Alors il existe un isomorphisme de triangles
$$
\xymatrix{
x \ar[r]^{\tilde{\alpha}} \ar[d] &
\tilde{y} \ar[r] \ar[d] & z
\ar[r]^{\delta} \ar@{.>}[d]^{e} & x[1]
\ar[d] \\ x
\ar[r]^{\alpha} & y
\ar[r] &
C(\alpha) \ar[r]^{-c} &
x[1]
}
$$
où le triangle supérieur est le
triangle associé à la suite $x
\xrightarrow{\tilde{\alpha}} \tilde{y} \to z$.
\end{enumerate}
\end{lemma}

\begin{proof}
Pour (1), nous considérons le diagramme plus complet, \emph{sans} le
changement de signe sur $c$ :
$$
\xymatrix{
x\ar@<0.5ex>[r]^{\alpha} \ar[d] &
y\ar@<0.5ex>[l]^{\pi} \ar@<0.5ex>[r]^{b}\ar[d] &
C(\alpha)\ar@<0.5ex>[l]^{p} \ar@<0.5ex>[r]^{c}\ar@{.>}@<0.5ex>[d]^{e} &
x[1]\ar@<0.5ex>[l]^{\sigma} \ar[d]\ar@<0.5ex>[r]^{\alpha} &
y[1]\ar@<0.5ex>[l]^{\pi} \\
x\ar@<0.5ex>[r]^{\alpha} &
y\ar@<0.5ex>[r]^{\beta} \ar@<0.5ex>[l]^{\pi} &
z\ar[r]^{\delta}\ar@<0.5ex>[l]^{s} \ar@{.>}@<0.5ex>[u]^{f} &
x[1]
}
$$
où la suite exacte courte admissible
$x\xrightarrow{\alpha} y\xrightarrow{\beta} z$ est
munie du scindage $\pi$, $s$, et la suite exacte courte admissible
$y\xrightarrow{b}C(\alpha)\xrightarrow{c}x[1]$ est munie du scindage $p$, $\sigma$.
Notons que (en identifiant les ensembles de morphismes par décalage)
$$
\alpha = pd(\sigma) = -d(p)\sigma,\quad
\delta = \pi d(s) = -d(\pi)s
$$
par la construction dans le lemme \ref{dga-lemma-get-triangle}.

\medskip\noindent
Nous définissons $e=\beta p$ et $f=bs-\sigma\delta$. Nous vérifions d'abord qu'ils
sont des morphismes dans $\text{Comp}(\mathcal{A})$. Pour montrer que $d(e)=\beta
d(p)$ s'annule, il suffit de montrer que $\beta d(p)b$ et $\beta d(p)\sigma$
s'annulent tous les deux ; or
$$
\beta d(p)b = \beta d(pb) = \beta d(1_y) = 0,\quad
\beta d(p)\sigma = -\beta\alpha = 0
$$
De même, pour vérifier que $d(f)=bd(s)-d(\sigma)\delta$ s'annule, il
suffit de vérifier que les post-compositions par $p$ et $c$ s'annulent toutes
deux ; or
\begin{align*}
pbd(s) - pd(\sigma)\delta
= &
d(s)-\alpha\delta = d(s)-\alpha\pi d(s) = 0 \\
cbd(s)-cd(\sigma)\delta =
&
-cd(\sigma)\delta = -d(c\sigma)\delta = 0
\end{align*}
La commutativité des deux carrés de gauche du
diagramme \ref{dga-equation-cone-isom-triangle} découle directement de la définition.
Avant de prouver la commutativité du carré de droite (à homotopie près), nous
vérifions d'abord que $e$ est une équivalence d'homotopie. Clairement,
$$
ef=\beta p (bs-\sigma\delta)=\beta s=1_z
$$
Pour vérifier que $fe$ est homotope à $1_{C(\alpha)}$, nous observons d'abord
$$
b\alpha = bpd(\alpha) = d(\sigma),\quad
\alpha c = -d(p)\sigma c = -d(p),\quad
d(\pi)p = d(\pi)s\beta p = -\delta\beta p
$$
En utilisant ces identités, nous calculons
\begin{align*}
1_{C(\alpha)} = &
bp + \sigma c
\quad (\text{d'après }y \xrightarrow{b} C(\alpha) \xrightarrow{c} x[1]) \\ = &
b(\alpha\pi
+ s\beta)p + \sigma(\pi\alpha)c \quad
(\text{d'après }x \xrightarrow{\alpha} y \xrightarrow{\beta} z) \\ = &
d(\sigma)\pi
p + bs\beta p - \sigma\pi d(p) \quad
(\text{d'après les deux premières identités ci-dessus}) \\ = &
d(\sigma)\pi
p + bs\beta p - \sigma\delta\beta p +
\sigma\delta\beta p - \sigma\pi d(p) \\ = & (bs -
\sigma\delta)\beta
p + d(\sigma)\pi p - \sigma d(\pi)p -
\sigma\pi d(p)\quad (\text{d'après la
troisième identité ci-dessus}) \\ = & fe +
d(\sigma
\pi p)
\end{align*}
puisque $\sigma \in \Hom^{-1}(x, C(\alpha))$
(cf. démonstration du lemme \ref{dga-lemma-id-cone-null}).
Par conséquent, $e$ et $f$ sont des inverses
homotopiques. Enfin, pour vérifier que le carré
droit du diagramme \ref{dga-equation-cone-isom-triangle} commute à homotopie
près, il suffit de vérifier que $-cf=\delta$. Ceci découle de
$$
-cf = -c(bs-\sigma\delta) = c\sigma\delta = \delta
$$
puisque $cb=0$.

\medskip\noindent
Pour (2), considérons la factorisation
$x\xrightarrow{\tilde{\alpha}}\tilde{y}\to y$
donnée comme dans le lemme \ref{dga-lemma-factor}, donc le deuxième
morphisme est une équivalence d'homotopie. D'après les lemmes
\ref{dga-lemma-cone} et \ref{dga-lemma-analogue-third-isomorphism},
il existe un isomorphisme de triangles entre
$$
x \xrightarrow{\alpha} y \to C(\alpha) \to x[1]
\quad\text{et}\quad
x \xrightarrow{\tilde{\alpha}} \tilde{y} \to C(\tilde{\alpha}) \to x[1]
$$
Puisque nous pouvons composer des isomorphismes de triangles, en
remplaçant $\alpha$ par $\tilde{\alpha}$, $y$ par $\tilde{y}$, et $C(\alpha)$
par $C(\tilde{\alpha})$, nous pouvons supposer que $\alpha$ est un monomorphisme
admissible. Dans ce cas, le résultat découle de (1).
\end{proof}

\noindent
Le lemme suivant est l'analogue du
lemme \ref{dga-lemma-homotopy-category-pre-triangulated}.

\begin{lemma}
\label{dga-lemma-analogue-homotopy-category-pre-triangulated}
Dans la situation \ref{dga-situation-ABC}, la catégorie homotopique
$K(\mathcal{A})$ avec ses foncteurs de translation naturels et ses triangles distingués
est une catégorie pré-triangulée.
\end{lemma}

\begin{proof}
Nous allons vérifier les axiomes TR1, TR2 et TR3.

\medskip\noindent
Démonstration de TR1. Par définition, tout triangle isomorphe à un triangle
distingué est distingué. Puisque
$$
\xymatrix{x\ar[r]^{1_x} & x\ar[r] & 0}
$$
est une suite exacte courte admissible, $(x, x, 0, 1_x, 0, 0)$
est un triangle distingué. De plus, étant donné un morphisme
$\alpha : x \to y$ dans $\text{Comp}(\mathcal{A})$, le triangle
donné par $(x, y, c(\alpha), \alpha, i, -p)$ est distingué d'après le
lemme \ref{dga-lemma-cone-homotopy}.

\medskip\noindent
Démonstration de TR2. Soit $(x,y,z,\alpha,\beta,\gamma)$ un triangle
et supposons que $(y,z,x[1],\beta,\gamma,-\alpha[1])$ soit
distingué. Alors il existe une suite exacte courte admissible
$0 \to x' \to y' \to z' \to 0$ telle que le triangle associé
$(x',y',z',\alpha',\beta',\gamma')$ soit isomorphe à
$(y,z,x[1],\beta,\gamma,-\alpha[1])$. Après rotation, nous concluons que
$(x,y,z,\alpha,\beta,\gamma)$ est isomorphe à $(z'[-1],x',y',
\gamma'[-1], \alpha',\beta')$. D'après le lemme
\ref{dga-lemma-cone-rotate-isom}, nous déduisons
que $(z'[-1],x',y', \gamma'[-1], \alpha',\beta')$ est isomorphe à
$(z'[-1],x',c(\gamma'[-1]), \gamma'[-1], i, p)$. En composant les
deux isomorphismes avec les changements de signe indiqués dans le
diagramme suivant :
$$
\xymatrix@C=3pc{
x\ar[r]^{\alpha}\ar[d] &
y\ar[r]^{\beta}\ar[d] &
z\ar[r]^{\gamma}\ar[d] &
x[1]\ar[d] \\
z'[-1]\ar[r]^{-\gamma'[-1]}\ar[d]_{-1_{z'[-1]}} & x
\ar[r]^{\alpha'}\ar@{=}[d] & y'
\ar[r]^{\beta'} \ar[d] &
z'\ar[d]^{-1_{z'}} \\
z'[-1]\ar[r]^{\gamma'[-1]} & x
\ar[r]^{\alpha'} &
c(\gamma'[-1]) \ar[r]^{-p} &
z'
}
$$
Nous concluons que $(x,y,z,\alpha,\beta,\gamma)$ est distingué par
le lemme \ref{dga-lemma-cone-homotopy} (2). Inversement, supposons que
$(x,y,z,\alpha,\beta,\gamma)$ soit distingué, de sorte que
selon le lemme \ref{dga-lemma-cone-homotopy} (1), il soit isomorphe
à un triangle de la forme $(x',y', c(\alpha'), \alpha', i,
-p)$ pour un certain morphisme $\alpha': x' \to y'$ dans
$\text{Comp}(\mathcal{A})$. Le triangle obtenu par rotation
$(y,z,x[1],\beta,\gamma, -\alpha[1])$ est isomorphe au triangle $(y',c(\alpha'), x'[1], i, -p,
-\alpha[1])$ qui est isomorphe à $(y',c(\alpha'), x'[1], i, p,
\alpha[1])$. D'après le lemme \ref{dga-lemma-restate-axiom-c}, ce triangle est
distingué, d'où il suit que $(y,z,x[1], \beta,\gamma, -\alpha[1])$ est
distingué.

\medskip\noindent
Démonstration de TR3 : Supposons que $(x,y,z,
\alpha,\beta,\gamma)$ et $(x',y',z',\alpha',\beta',\gamma')$ soient des triangles
distingués de $\text{Comp}(\mathcal{A})$ et que $f: x \to x'$ et $g: y
\to y'$ soient des morphismes tels que
$\alpha' \circ f = g \circ \alpha$. D'après le
lemme \ref{dga-lemma-cone-homotopy}, nous pouvons supposer que
$(x,y,z,\alpha,\beta,\gamma)= (x,y,c(\alpha),\alpha, i, -p)$ et
$(x',y',z', \alpha',\beta',\gamma')= (x',y',c(\alpha'), \alpha',i',-p')$. Appliquons maintenant
le lemme \ref{dga-lemma-cone} et nous
avons terminé.
\end{proof}

\noindent
Le lemme suivant est l'analogue du lemme \ref{dga-lemma-two-split-injections}.

\begin{lemma}
\label{dga-lemma-dgc-analogue-tr4}
Dans la situation \ref{dga-situation-ABC}, étant donnés des monomorphismes
admissibles $x \xrightarrow{\alpha} y$, $y \xrightarrow{\beta} z$ dans
$\mathcal{A}$, il existe des triangles
distingués $(x,y,q_1,\alpha,p_1,\delta_1)$, $(x,z,q_2,\beta\alpha,p_2,\delta_2)$ et
$(y,z,q_3,\beta,p_3,\delta_3)$ pour lesquels TR4 est vérifié.
\end{lemma}

\begin{proof}
Étant donnés les monomorphismes admissibles $x\xrightarrow{\alpha}
y$ et $y\xrightarrow{\beta}z$, nous pouvons trouver des triangles
distingués, en les prolongeant en suites exactes courtes admissibles,
$$
\xymatrix{
x\ar@<0.5ex>[r]^{\alpha} &
y\ar@<0.5ex>[l]^{\pi_1} \ar@<0.5ex>[r]^{p_1} &
q_1 \ar[r]^{\delta_1} \ar@<0.5ex>[l]^{s_1} &
x[1]
}
$$
$$
\xymatrix{
x\ar@<0.5ex>[r]^{\beta\alpha} &
z\ar@<0.5ex>[l]^{\pi_1\pi_3} \ar@<0.5ex>[r]^{p_2} &
q_2 \ar[r]^{\delta_2} \ar@<0.5ex>[l]^{s_2} &
x[1]
}
$$
$$
\xymatrix{
y\ar@<0.5ex>[r]^{\beta} &
z\ar@<0.5ex>[l]^{\pi_3} \ar@<0.5ex>[r]^{p_3} &
q_3 \ar[r]^{\delta_3} \ar@<0.5ex>[l]^{s_3} &
x[1]
}
$$
Dans ces diagrammes, les applications $\delta_i$ sont définies comme
$\delta_i = \pi_i d(s_i)$ de manière analogue aux applications définies dans le lemme
\ref{dga-lemma-get-triangle}. Elles s'insèrent dans
le diagramme commutatif suivant, dont les flèches sont en traits pleins
$$
\xymatrix@C=5pc@R=3pc{
x\ar@<0.5ex>[r]^{\alpha} \ar@<0.5ex>[dr]^{\beta\alpha} &
y\ar@<0.5ex>[d]^{\beta} \ar@<0.5ex>[l]^{\pi_1} \ar@<0.5ex>[r]^{p_1} &
q_1 \ar[r]^{\delta_1} \ar@<0.5ex>[l]^{s_1} \ar@{.>}[dd]^{p_2\beta s_1} &
x[1] \\ & z
\ar@<0.5ex>[u]^{\pi_3}\ar@<0.5ex>[d]^{p_3}
\ar@<0.5ex>[dr]^{p_2}
\ar@<0.5ex>[ul]^{\pi_1\pi_3} & & \\ &
q_3\ar@<0.5ex>[u]^{s_3}
\ar[d]^{\delta_3} & q_2
\ar@{.>}[l]^{p_3s_2} \ar@<0.5ex>[ul]^{s_2} \ar[dr]^{\delta_2} \\ & y[1] &
&
x[1]}
$$
où nous avons défini les flèches en pointillés comme indiqué.
Clairement, leur composition $p_3s_2p_2\beta s_1 = 0$ puisque
$s_2p_2 = 0$. Nous affirmons qu'elles sont toutes deux des morphismes de
$\text{Comp}(\mathcal{A})$. Nous pouvons vérifier cela en utilisant les équations
du lemme \ref{dga-lemma-get-triangle} :
$$
d(p_2\beta s_1) = p_2\beta d(s_1) = p_2\beta\alpha\pi_1 d(s_1) = 0
$$
puisque $p_2\beta\alpha = 0$, et
$$
d(p_3s_2) = p_3d(s_2) = p_3\beta\alpha\pi_1\pi_3 d(s_2) = 0
$$
puisque $p_3\beta = 0$. Pour vérifier que $q_1\to q_2\to
q_3$ est une suite exacte courte admissible, il reste à montrer
que, dans la catégorie graduée sous-jacente, $q_2 = q_1\oplus q_3$
avec les deux morphismes ci-dessus comme coprojection et projection.
Pour ce faire, observons que dans la catégorie graduée
sous-jacente $\mathcal{C}$, il y a
$$
y = x\oplus q_1,\quad
z = y\oplus q_3 = x\oplus q_1\oplus q_3
$$
où $\pi_1\pi_3$ donne le morphisme de projection sur le premier
facteur : $x\oplus q_1\oplus q_3\to z$. Par l'axiome (A) sur
$\mathcal{A}$, $\mathcal{C}$ est une catégorie additive, par conséquent
nous pouvons
appliquer Homologie, lemme \ref{homology-lemma-additive-cat-biproduct-kernel}
et conclure que
$$
\Ker(\pi_1\pi_3) = q_1\oplus q_3
$$
dans $\mathcal{C}$. Une autre application
d'Homologie, lemme \ref{homology-lemma-additive-cat-biproduct-kernel} à $z =
x\oplus q_2$ donne $\Ker(\pi_1\pi_3) = q_2$. Par conséquent
$q_2\cong q_1\oplus q_3$ dans $\mathcal{C}$. Il est clair
que les morphismes en pointillés définis ci-dessus donnent la
coprojection et la projection.

\medskip\noindent
Enfin, nous devons vérifier que le morphisme
$\delta : q_3 \to q_1[1]$ induit par la suite exacte
courte admissible $q_1\to q_2\to q_3$ coïncide avec
$p_1\delta_3$. Selon la construction du
lemme \ref{dga-lemma-get-triangle}, le morphisme $\delta$ est donné par
\begin{align*}
p_1\pi_3s_2d(p_2s_3)
= &
p_1\pi_3s_2p_2d(s_3) \\ =
&
p_1\pi_3(1-\beta\alpha\pi_1\pi_3)d(s_3) \\ = &
p_1\pi_3d(s_3)\quad
(\text{puisque }\pi_3\beta = 0) \\ =
&
p_1\delta_3
\end{align*}
comme souhaité. La démonstration est complète.
\end{proof}

\noindent
En rassemblant tout, nous obtenons finalement l'analogue
de la proposition \ref{dga-proposition-homotopy-category-triangulated}.

\begin{proposition}
\label{dga-proposition-ABC-homotopy-category-triangulated}
Dans la situation \ref{dga-situation-ABC}, la catégorie homotopique
$K(\mathcal{A})$ avec ses foncteurs de translation naturels et ses triangles distingués est
une catégorie triangulée.
\end{proposition}

\begin{proof}
D'après le lemme \ref{dga-lemma-analogue-homotopy-category-pre-triangulated}, nous
savons que $K(\mathcal{A})$ est pré-triangulée. En
combinant les lemmes \ref{dga-lemma-analogue-sequence-maps-split}
et \ref{dga-lemma-dgc-analogue-tr4}
avec Catégories dérivées, lemme \ref{derived-lemma-easier-axiom-four},
nous concluons que $K(\mathcal{A})$ est une catégorie triangulée.
\end{proof}

\begin{lemma}
\label{dga-lemma-functor-between-ABC}
Soit $R$ un anneau. Soit $F : \mathcal{A} \to \mathcal{B}$ un foncteur
entre des catégories différentielles graduées sur $R$ satisfaisant aux
axiomes (A), (B) et (C), tel que $F(x[1]) =
F(x)[1]$. Alors $F$ induit un foncteur
exact $K(\mathcal{A}) \to K(\mathcal{B})$ de catégories triangulées.
\end{lemma}

\begin{proof}
En effet, si $x \to y \to z$ est une suite exacte courte
admissible dans $\text{Comp}(\mathcal{A})$, alors $F(x) \to F(y) \to
F(z)$ est une suite exacte courte admissible dans
$\text{Comp}(\mathcal{B})$. De plus, le morphisme de liaison $\delta = \pi\text{d}(s) : z \to x[1]$
construit dans le lemme \ref{dga-lemma-get-triangle} produit le morphisme
$F(\delta) : F(z) \to F(x[1]) = F(x)[1]$ qui est égal au morphisme de
liaison $F(\pi) \text{d}(F(s))$ pour la suite exacte courte admissible $F(x)
\to F(y) \to F(z)$.
\end{proof}





\section{Bimodules}
\label{dga-section-bimodules}

\noindent
Nous continuons la discussion commencée dans la section \ref{dga-section-tensor-product}.

\begin{definition}
\label{dga-definition-bimodule}
Bimodules. Soit $R$ un anneau.
\begin{enumerate}
\item Soient $A$ et $B$ des $R$-algèbres. Un {\it $(A, B)$-bimodule}
est un $R$-module $M$ muni d'applications $R$-bilinéaires
$$
A \times M \to M, (a, x) \mapsto ax
\quad\text{et}\quad
M \times B \to M, (x, b) \mapsto xb
$$
vérifiant les propriétés suivantes
\begin{enumerate}
\item $a'(ax) = (a'a)x$ et $(xb)b' = x(bb')$,
\item $a(xb) = (ax)b$, et
\item $1 x = x = x 1$.
\end{enumerate}
\item Soient $A$ et $B$ des algèbres graduées par $\mathbf{Z}$ sur
$R$. Un {\it $(A, B)$-bimodule gradué} est un $(A, B)$-bimodule $M$ qui
possède une graduation $M = \bigoplus M^n$
telle que $A^n M^m \subset M^{n + m}$ et $M^n B^m \subset M^{n + m}$.
\item Soient $A$ et $B$ des algèbres différentielles graduées
sur $R$. Un {\it $(A, B)$-bimodule différentiel gradué} est un $(A,
B)$-bimodule gradué qui est muni d'une
différentielle $\text{d} : M \to M$ homogène de degré $1$
telle que $\text{d}(ax) = \text{d}(a)x + (-1)^{\deg(a)}a\text{d}(x)$ et
$\text{d}(xb) = \text{d}(x)b + (-1)^{\deg(x)}x\text{d}(b)$ pour
des éléments homogènes $a \in A$, $x \in M$, $b \in B$.
\end{enumerate}
\end{definition}

\noindent
Remarquons qu'un $(A, B)$-bimodule différentiel gradué $M$
est la même chose qu'un module différentiel gradué à droite
sur $B$ qui est également un module différentiel gradué à
gauche sur $A$, avec des graduations et des différentielles qui
coïncident, et tel que la structure de module sur $A$ commute
avec la structure de module sur $B$. Voici un énoncé précis.

\begin{lemma}
\label{dga-lemma-what-makes-a-bimodule-dg}
Soit $R$ un anneau. Soient $(A, \text{d})$ et $(B, \text{d})$ des
algèbres différentielles graduées sur $R$. Soit $M$ un $B$-module
différentiel gradué droit. Il existe une correspondance $1$ à $1$ entre
les structures de bimodule $(A, B)$ sur $M$ compatibles avec la structure
de $B$-module différentiel gradué donnée et les homomorphismes
$$
A
\longrightarrow
\Hom_{\text{Mod}^{dg}_{(B, \text{d})}}(M, M)
$$
d'algèbres différentielles graduées sur $R$.
\end{lemma}

\begin{proof}
Soit $\mu : A \times M \to M$ une application définissant une structure de
$A$-module différentiel gradué à gauche sur le complexe de $R$-modules $M^\bullet$
sous-jacent à $M$. D'après le lemme \ref{dga-lemma-left-module-structure}, la structure $\mu$
correspond à une application $\gamma : A \to \Hom^\bullet(M^\bullet,
M^\bullet)$ d'algèbres différentielles graduées sur $R$. L'assertion du lemme est
simplement que $\mu$ commute avec l'action de $B$ si et seulement si $\gamma$ est à
valeurs dans
$$
\Hom_{\text{Mod}^{dg}_{(B, \text{d})}}(M, M) \subset
\Hom^\bullet(M^\bullet, M^\bullet)
$$
Nous omettons le calcul détaillé.
\end{proof}

\noindent
Soit $M$ un $(A, B)$-bimodule différentiel gradué. Rappelons d'après la
section \ref{dga-section-left-modules} que la structure de module différentiel
gradué à gauche sur $A$ correspond à une structure de module différentiel
gradué à droite sur $A^{opp}$. Comme les structures de module sur $A$ et sur $B$
commutent, cela donne à $M$ la structure d'un module différentiel
gradué sur $A^{opp} \otimes_R B$ :
$$
x \cdot (a \otimes b) = (-1)^{\deg(a)\deg(x)} axb
$$
Inversement, si nous avons un $A^{opp} \otimes_R B$-module différentiel gradué $M$,
alors nous pouvons utiliser la formule ci-dessus pour obtenir un $(A, B)$-bimodule
différentiel gradué.

\begin{lemma}
\label{dga-lemma-bimodule-over-tensor}
Soit $R$ un anneau. Soient $(A, \text{d})$ et $(B, \text{d})$
des algèbres différentielles graduées sur $R$. La construction
ci-dessus définit une équivalence de catégories
$$
\begin{matrix}
\text{bimodules différentiels
gradués sur}\\ (A, B)\text{}
\end{matrix}
\longleftrightarrow
\begin{matrix}
\text{modules différentiels gradués à
droite sur }\\ A^{opp} \otimes_R B\text{}
\end{matrix}
$$
\end{lemma}

\begin{proof}
Cela résulte immédiatement de la discussion ci-dessus.
\end{proof}

\noindent
Soit $R$ un anneau. Soient $(A, \text{d})$ et $(B, \text{d})$ des
$R$-algèbres différentielles
graduées. Soit $P$ un bimodule différentiel gradué sur $(A, B)$. On dit
que $P$ {\it a la propriété (P)} s'il existe une filtration
$$
0 = F_{-1}P \subset F_0P \subset F_1P \subset \ldots \subset P
$$
par des $(A, B)$-bimodules différentiels gradués, telle que
\begin{enumerate}
\item $P = \bigcup F_pP$,
\item les inclusions $F_iP \to F_{i + 1}P$ sont scindées comme morphismes
de $(A, B)$-bimodules gradués,
\item les quotients $F_{i + 1}P/F_iP$ sont isomorphes comme $(A,
B)$-bimodules différentiels gradués à une somme directe de $(A \otimes_R B)[k]$.
\end{enumerate}

\begin{lemma}
\label{dga-lemma-bimodule-resolve}
Soit $R$ un anneau. Soient $(A, \text{d})$ et $(B, \text{d})$ des
$R$-algèbres différentielles graduées. Soit $M$ un bimodule différentiel
gradué sur $(A, B)$. Il existe un homomorphisme $P \to M$ de
bimodules différentiels gradués sur $(A, B)$ qui est un quasi-isomorphisme tel que
$P$ possède la propriété (P) comme définie ci-dessus.
\end{lemma}

\begin{proof}
Cela résulte immédiatement des lemmes
\ref{dga-lemma-bimodule-over-tensor} et \ref{dga-lemma-resolve}.
\end{proof}

\begin{lemma}
\label{dga-lemma-bimodule-property-P-sequence}
Soit $R$ un anneau. Soient $(A, \text{d})$ et $(B, \text{d})$ des
$R$-algèbres différentielles graduées. Soit $P$ un
$(A, B)$-bimodule différentiel gradué possédant la propriété (P)
avec la filtration correspondante $F_\bullet$, alors nous obtenons une
suite exacte courte
$$
0 \to
\bigoplus\nolimits F_iP \to
\bigoplus\nolimits F_iP \to P \to 0
$$
de $(A, B)$-bimodules différentiels gradués, qui est scindée comme suite
de $(A, B)$-bimodules gradués.
\end{lemma}

\begin{proof}
Cela résulte immédiatement des lemmes
\ref{dga-lemma-bimodule-over-tensor} et \ref{dga-lemma-property-P-sequence}.
\end{proof}















\section{Bimodules et produit tensoriel}
\label{dga-section-bimodules-tensor}

\noindent
Soit $R$ un anneau. Soient $A$ et $B$ des $R$-algèbres. Soit $M$ un module à
droite sur $A$. Soit $N$ un bimodule sur $(A, B)$. Alors $M
\otimes_A N$ est un module à droite sur $B$.

\medskip\noindent
Si, dans la situation du paragraphe précédent, $A$
et $B$ sont des algèbres graduées par $\mathbf{Z}$,
$M$ est un module gradué sur $A$, et $N$ est un $(A, B)$-bimodule gradué,
alors $M \otimes_A N$ est un module gradué à droite sur $B$. La construction
est fonctorielle en $M$ et définit un foncteur
$$
- \otimes_A N :
\text{Mod}^{gr}_A
\longrightarrow
\text{Mod}^{gr}_B
$$
de catégories graduées comme dans l'exemple \ref{dga-example-gm-gr-cat}. En
effet, si $M$ et $M'$ sont des $A$-modules gradués et que $f : M \to M'$ est un
homomorphisme de $A$-modules homogène de degré $n$,
alors $f \otimes \text{id}_N : M \otimes_A N \to M' \otimes_A N$
est un homomorphisme de $B$-modules homogène de degré $n$.

\medskip\noindent
Si dans la situation du paragraphe précédent $(A,
\text{d})$ et $(B, \text{d})$ sont des algèbres différentielles graduées, $M$
est un module différentiel gradué sur $A$, et $N$ est un bimodule
différentiel gradué sur $(A, B)$, alors $M \otimes_A N$ est un module
différentiel gradué à droite sur $B$.

\begin{lemma}
\label{dga-lemma-tensor}
Soit $R$ un anneau. Soient $(A, \text{d})$ et $(B, \text{d})$
des algèbres différentielles graduées sur $R$. Soit $N$ un
$(A, B)$-bimodule différentiel gradué. Alors
$M \mapsto M \otimes_A N$ définit un foncteur
$$
- \otimes_A N :
\text{Mod}^{dg}_{(A, \text{d})}
\longrightarrow
\text{Mod}^{dg}_{(B, \text{d})}
$$
de catégories différentielles graduées. Ce foncteur induit des foncteurs
$$
\text{Mod}_{(A, \text{d})} \to \text{Mod}_{(B, \text{d})}
\quad\text{et}\quad
K(\text{Mod}_{(A, \text{d})}) \to K(\text{Mod}_{(B, \text{d})})
$$
par une application du lemme \ref{dga-lemma-functorial}.
\end{lemma}

\begin{proof}
Ci-dessus, nous avons vu comment la construction définit un foncteur des
catégories graduées sous-jacentes. Ainsi, il suffit de montrer que la
construction est compatible avec les différentielles. Soient $M$ et $M'$ des
$A$-modules différentiels gradués et soit $f : M \to M'$ un homomorphisme de
$A$-modules homogène de degré $n$. Alors nous avons
$$
\text{d}(f) = \text{d}_{M'} \circ f - (-1)^n f \circ \text{d}_M
$$
D'un autre côté, nous avons
$$
\text{d}(f \otimes \text{id}_N) =
\text{d}_{M' \otimes_A N} \circ
(f \otimes \text{id}_N) -
(-1)^n
(f \otimes \text{id}_N) \circ \text{d}_{M \otimes_A N}
$$
En appliquant ceci à un élément $x \otimes y$ avec $x \in M$ et $y
\in N$ homogènes, nous obtenons
\begin{align*}
\text{d}(f \otimes \text{id}_N)(x \otimes y)
= &
\text{d}_{M'}(f(x)) \otimes y + (-1)^{n + \deg(x)}f(x) \otimes \text{d}_N(y) \\ &
- (-1)^n f(\text{d}_M(x)) \otimes y -
(-1)^{n + \deg(x)}f(x) \otimes \text{d}_N(y) \\ = &
\text{d}(f)
(x \otimes y)
\end{align*}
Ainsi, nous voyons que $\text{d}(f) \otimes \text{id}_N =
\text{d}(f \otimes \text{id}_N)$ et la démonstration est complète.
\end{proof}

\begin{remark}
\label{dga-remark-shift-tensor-no-sign}
Soit $R$ un anneau. Soient $(A, \text{d})$ et $(B, \text{d})$
des algèbres différentielles graduées sur $R$. Soit $N$ un
bimodule différentiel gradué sur $(A, B)$. Soit $M$ un module différentiel gradué
à droite sur $A$. Alors pour chaque $k \in \mathbf{Z}$ il existe
un isomorphisme
$$
(M \otimes_A N)[k] \longrightarrow
M[k] \otimes_A N
$$
de modules différentiels gradués à droite sur $B$, défini sans intervention
de signes, voir Compléments d'algèbre, section \ref{more-algebra-section-sign-rules}.
\end{remark}

\noindent
Si nous avons un anneau $R$ et des $R$-algèbres $A$, $B$, et
$C$, un $A$-module à droite $M$, un $(A, B)$-bimodule $N$, et un
$(B, C)$-bimodule $N'$, alors
$N \otimes_B N'$ est un $(A, C)$-bimodule et nous avons
$$
(M \otimes_A N) \otimes_B N' = M \otimes_A (N \otimes_B N')
$$
Cette égalité reste valable dans le cas gradué et dans le cas différentiel
gradué. Voir Compléments d'algèbre, section \ref{more-algebra-section-sign-rules} pour les
règles des signes.






\section{Bimodules et Hom interne}
\label{dga-section-bimodules-hom}

\noindent
Soit $R$ un anneau. Si $A$ est une $R$-algèbre (voir nos conventions à la
section \ref{dga-section-conventions}) et que $M$, $M'$ sont des $A$-modules à
droite, alors nous définissons
$$
\Hom_A(M, M') = \{f : M \to M' \mid f \text{ est }A\text{-linéaire}\}
$$
comme d'habitude.

\medskip\noindent
Soit $R$ un anneau. Soient $A$ et $B$ des $R$-algèbres. Soit
$N$ un $(A, B)$-bimodule. Soit $N'$ un module à droite sur $B$.
Dans cette situation, nous considérerons
$$
\Hom_B(N, N')
$$
comme $A$-module à droite, par précomposition.

\medskip\noindent
Soit $R$ un anneau. Soient $A$ et $B$ des algèbres graduées par $\mathbf{Z}$ sur
$R$. Soit $N$ un $(A, B)$-bimodule gradué. Soit $N'$ un module gradué à droite sur
$B$. Dans cette situation, nous considérerons le $R$-module gradué
$$
\Hom_{\text{Mod}^{gr}_B}(N, N')
$$
défini dans l'exemple \ref{dga-example-gm-gr-cat} comme un $A$-module gradué à droite
en utilisant la précomposition. La construction est fonctorielle en $N'$ et définit
un foncteur
$$
\Hom_{\text{Mod}^{gr}_B}(N, -) :
\text{Mod}^{gr}_B
\longrightarrow
\text{Mod}^{gr}_A
$$
de catégories graduées comme dans l'exemple \ref{dga-example-gm-gr-cat}. En
effet, si $N_1$ et $N_2$ sont des $B$-modules gradués et que $f : N_1 \to N_2$ est un
homomorphisme de $B$-modules homogène de degré $n$, alors l'application
induite $\Hom_{\text{Mod}^{gr}_B}(N, N_1) \to \Hom_{\text{Mod}^{gr}_B}(N, N_2)$
est un homomorphisme de $A$-modules homogène de degré $n$.

\medskip\noindent
Soit $R$ un anneau. Soient $A$ et $B$ des algèbres différentielles graduées
par $\mathbf{Z}$ sur $R$. Soit $N$ un bimodule différentiel gradué sur $(A,
B)$. Soit $N'$ un module différentiel gradué à droite sur $B$. Dans cette
situation, nous considérerons le $R$-module différentiel gradué
$$
\Hom_{\text{Mod}^{dg}_{(B, \text{d})}}(N, N')
$$
défini dans l'exemple \ref{dga-example-dgm-dg-cat} comme un $A$-module
différentiel gradué à droite en utilisant la précomposition comme dans le cas gradué.
Ceci est compatible avec les différentielles car la multiplication est la
composition
$$
\Hom_{\text{Mod}^{dg}_B}(N, N') \otimes_R A \to
\Hom_{\text{Mod}^{dg}_B}(N, N') \otimes_R
\Hom_{\text{Mod}^{dg}_B}(N, N) \to
\Hom_{\text{Mod}^{dg}_B}(N, N')
$$
La première flèche utilise
l'application du lemme \ref{dga-lemma-what-makes-a-bimodule-dg} et la
deuxième flèche est la composition dans la catégorie différentielle
graduée $\text{Mod}^{dg}_{(B, \text{d})}$.

\begin{lemma}
\label{dga-lemma-hom}
Soit $R$ un anneau. Soient $(A, \text{d})$ et $(B, \text{d})$
des algèbres différentielles graduées sur $R$. Soit $N$ un
$(A, B)$-bimodule différentiel gradué. La construction ci-dessus
définit un foncteur
$$
\Hom_{\text{Mod}^{dg}_{(B, \text{d})}}(N, -) :
\text{Mod}^{dg}_{(B, \text{d})}
\longrightarrow
\text{Mod}^{dg}_{(A, \text{d})}
$$
de catégories différentielles graduées. Ce foncteur induit des foncteurs
$$
\text{Mod}_{(B, \text{d})} \to \text{Mod}_{(A, \text{d})}
\quad\text{et}\quad
K(\text{Mod}_{(B, \text{d})}) \to K(\text{Mod}_{(A, \text{d})})
$$
par une application du lemme \ref{dga-lemma-functorial}.
\end{lemma}

\begin{proof}
Ci-dessus, nous avons vu comment la construction définit un foncteur des
catégories graduées sous-jacentes. Il suffit donc de montrer que la construction est
compatible avec les différentielles. Soient $N_1$ et $N_2$ des $B$-modules
différentiels gradués. Écrivons
$$
H_{12} = \Hom_{\text{Mod}^{dg}_{(B, \text{d})}}(N_1, N_2),\quad
H_1 = \Hom_{\text{Mod}^{dg}_{(B, \text{d})}}(N, N_1),\quad
H_2 = \Hom_{\text{Mod}^{dg}_{(B, \text{d})}}(N, N_2)
$$
Considérons la composition
$$
c : H_{12} \otimes_R H_1 \longrightarrow H_2
$$
dans la catégorie différentielle graduée $\text{Mod}^{dg}_{(B,
\text{d})}$. Soit $f : N_1 \to N_2$ un homomorphisme de $B$-modules qui est
homogène de degré $n$, en d'autres termes, $f \in
H_{12}^n$. Le foncteur dans le lemme envoie $f$ sur $c_f : H_1 \to H_2$, $g \mapsto c(f,
g)$. De même pour $\text{d}(f)$. D'autre part, la différentielle sur
$$
\Hom_{\text{Mod}^{dg}_{(A, \text{d})}}(H_1, H_2)
$$
envoie $c_f$ sur $\text{d}_{H_2} \circ c_f - (-1)^n c_f \circ \text{d}_{H_1}$. Comme
$c$ est un morphisme de complexes de $R$-modules, nous avons
$\text{d} c(f, g) = c(\text{d}f, g) + (-1)^n c(f, \text{d}g)$. Par
conséquent, nous voyons que
\begin{align*}
(\text{d}c_f)(g)
& =
\text{d}c(f,g) - (-1)^n c(f, \text{d}g) \\
& =
c(\text{d}f, g) + (-1)^n c(f, \text{d}g) - (-1)^n c(f, \text{d}g) \\ & =
c(\text{d}f,
g) = c_{\text{d}f}(g)
\end{align*}
et la démonstration est complète.
\end{proof}

\begin{remark}
\label{dga-remark-shift-hom-no-sign}
Soit $R$ un anneau. Soient $(A, \text{d})$ et $(B, \text{d})$
des algèbres différentielles graduées sur $R$. Soit $N$ un
bimodule différentiel gradué sur $(A, B)$. Soit $N'$ un module différentiel gradué
à droite sur $B$. Alors pour chaque $k \in \mathbf{Z}$ il existe
un isomorphisme
$$
\Hom_{\text{Mod}^{gr}_B}(N, N')[k]
\longrightarrow
\Hom_{\text{Mod}^{gr}_B}(N, N'[k])
$$
de modules différentiels gradués à droite sur $A$, défini sans intervention
de signes, voir Compléments d'algèbre, section \ref{more-algebra-section-sign-rules}.
\end{remark}

\begin{lemma}
\label{dga-lemma-tensor-hom-adjunction}
Soit $R$ un anneau. Soient $A$ et $B$ des $R$-algèbres.
Soit $M$ un module à droite sur $A$, $N$ un bimodule sur $(A, B)$, et $N'$
un module à droite sur $B$. Alors nous avons un isomorphisme canonique
$$
\Hom_B(M \otimes_A N, N') = \Hom_A(M, \Hom_B(N, N'))
$$
de $R$-modules.
Si $A$, $B$, $M$, $N$, $N'$ sont compatiblement gradués, alors nous avons un
isomorphisme canonique
$$
\Hom_{\text{Mod}_B^{gr}}(M \otimes_A N, N') =
\Hom_{\text{Mod}_A^{gr}}(M, \Hom_{\text{Mod}_B^{gr}}(N, N'))
$$
de $R$-modules gradués. Si
$A$, $B$, $M$, $N$, $N'$ sont munis de structures différentielles graduées compatibles,
alors nous avons un isomorphisme canonique
$$
\Hom_{\text{Mod}^{dg}_{(B, \text{d})}}(M \otimes_A N, N') =
\Hom_{\text{Mod}^{dg}_{(A, \text{d})}}(M,
\Hom_{\text{Mod}^{dg}_{(B, \text{d})}}(N, N'))
$$
de complexes de $R$-modules.
\end{lemma}

\begin{proof}
Démonstration omise. Indication : dans le cas non gradué, interpréter les deux membres
comme des applications $A$-bilinéaires $\psi : M \times N \to N'$ qui sont
$B$-linéaires à droite. Dans le cas (différentiel) gradué, utiliser
l'isomorphisme de Compléments d'algèbre, lemme
\ref{more-algebra-lemma-compose} et vérifier qu'il est compatible avec les structures
de modules. On peut aussi utiliser l'isomorphisme du lemme
\ref{dga-lemma-characterize-hom} et montrer qu'il est compatible avec les structures de $B$-modules.
\end{proof}







\section{Hom dérivé}
\label{dga-section-restriction}

\noindent
Cette section est analogue à
Compléments d'algèbre, section \ref{more-algebra-section-RHom}.

\medskip\noindent
Soit $R$ un anneau. Soient $(A, \text{d})$ et $(B, \text{d})$
des algèbres différentielles graduées sur $R$. Soit $N$ un
bimodule différentiel gradué $(A, B)$. Considérons le foncteur
\begin{equation}
\label{dga-equation-restriction}
\Hom_{\text{Mod}^{dg}_{(B, \text{d})}}(N, -) :
\text{Mod}_{(B, \text{d})}
\longrightarrow
\text{Mod}_{(A, \text{d})}
\end{equation}
de la section \ref{dga-section-bimodules-hom}.

\begin{lemma}
\label{dga-lemma-restriction-homotopy}
Le foncteur (\ref{dga-equation-restriction}) définit un foncteur
exact $K(\text{Mod}_{(B, \text{d})}) \to K(\text{Mod}_{(A, \text{d})})$
de catégories triangulées.
\end{lemma}

\begin{proof}
Via le lemme \ref{dga-lemma-hom}
et la remarque
\ref{dga-remark-shift-hom-no-sign}, cela découle du principe général du
lemme \ref{dga-lemma-functor-between-ABC}.
\end{proof}

\noindent
Rappelons que nous avons un foncteur
exact de catégories triangulées
$$
\Hom_{\text{Mod}^{dg}_{(B, \text{d})}}(N, -) :
K(\text{Mod}_{(B, \text{d})}) \to K(\text{Mod}_{(A, \text{d})})
$$
voir le lemme \ref{dga-lemma-restriction-homotopy}. Considérons le diagramme
$$
\xymatrix{
K(\text{Mod}_{(B, \text{d})}) \ar[d] \ar[rr]_{\text{voir ci-dessus}} \ar[rrd]_F & &
K(\text{Mod}_{(A, \text{d})}) \ar[d] \\
D(B, \text{d}) \ar@{..>}[rr] & &
D(A, \text{d})
}
$$
Nous aimerions construire une flèche pointillée en
tant que {\it foncteur dérivé à droite} de la composition
$F$. ({\it Avertissement} : dans la plupart des cas intéressants, le diagramme ne
commutera pas.) Plus précisément, dans le
cadre général de Catégories dérivées, section
\ref{derived-section-derived-functors}, nous voulons
calculer le foncteur dérivé à droite de $F$ par rapport au système multiplicatif des
quasi-isomorphismes dans $K(\text{Mod}_{(A, \text{d})})$.

\begin{lemma}
\label{dga-lemma-derived-restriction}
Dans la situation ci-dessus, le foncteur dérivé à droite de $F$
existe. Nous le notons $R\Hom(N, -) : D(B, \text{d}) \to D(A, \text{d})$.
\end{lemma}

\begin{proof}
Nous
utiliserons Catégories dérivées, lemme \ref{derived-lemma-find-existence-computes} pour
le prouver. Comme collection $\mathcal{I}$ d'objets,
nous prendrons les objets ayant la propriété (I). La propriété
(1) a été démontrée dans le lemme \ref{dga-lemma-right-resolution}. La
propriété (2) est vraie parce que si $s : I \to I'$ est un quasi-isomorphisme de
modules avec la propriété (I), alors $s$ est une équivalence d'homotopie selon
le lemme \ref{dga-lemma-hom-derived}.
\end{proof}

\begin{lemma}
\label{dga-lemma-functoriality-derived-restriction}
Soit $R$ un anneau. Soient $(A, \text{d})$ et $(B, \text{d})$ des
$R$-algèbres différentielles graduées. Soit $f : N \to N'$ un
homomorphisme de $(A, B)$-bimodules différentiels gradués.
Alors $f$ induit un morphisme de foncteurs
$$
- \circ f : R\Hom(N', -) \longrightarrow R\Hom(N, -)
$$
Si $f$ est un quasi-isomorphisme, alors $f \circ -$ est un isomorphisme de
foncteurs.
\end{lemma}

\begin{proof}
Notons $\mathcal{B} = \text{Mod}^{dg}_{(B, \text{d})}$ la
catégorie différentielle graduée des $B$-modules différentiels gradués, voir
l'exemple \ref{dga-example-dgm-dg-cat}. Soit $I$
un $B$-module différentiel gradué avec la propriété (I). Alors $f \circ
- : \Hom_\mathcal{B}(N', I) \to \Hom_\mathcal{B}(N, I)$ est une application de
$A$-modules différentiels gradués. De plus, ceci est fonctoriel par rapport à $I$.
Puisque les foncteurs $ R\Hom(N', -)$ et
$R\Hom(N, -)$ sont calculés en prenant pour
deuxième argument de $\Hom_\mathcal{B}$ des objets possédant la propriété (I)
(lemme \ref{dga-lemma-derived-restriction}), nous obtenons une transformation de foncteurs comme
indiqué.

\medskip\noindent
Supposons que $f$ soit un quasi-isomorphisme. Soit $F_\bullet$ la
filtration donnée sur $I$. Puisque $I = \lim I/F_pI$, nous voyons que
$\Hom_\mathcal{B}(N', I) = \lim \Hom_\mathcal{B}(N', I/F_pI)$ et
$\Hom_\mathcal{B}(N, I) = \lim \Hom_\mathcal{B}(N, I/F_pI)$. Puisque les
applications de transition dans le système $I/F_pI$ sont scindées comme
morphismes de modules gradués, nous voyons que les applications de
transition dans les systèmes $\Hom_\mathcal{B}(N', I/F_pI)$ et $\Hom_\mathcal{B}(N,
I/F_pI)$ sont surjectives. Ainsi, $\Hom_\mathcal{B}(N', I)$, resp. $\Hom_\mathcal{B}(N, I)$
vu comme un complexe de groupes abéliens calcule $R\lim$ du système de
complexes
$\Hom_\mathcal{B}(N', I/F_pI)$, resp. $\Hom_\mathcal{B}(N, I/F_pI)$. Voir plus sur
l'algèbre, lemme \ref{more-algebra-lemma-compute-Rlim}. Ainsi, il suffit de
prouver que chaque morphisme
$$
\Hom_\mathcal{B}(N', I/F_pI) \to \Hom_\mathcal{B}(N, I/F_pI)
$$
est un quasi-isomorphisme. Puisque les surjections $I/F_{p + 1}I \to I/F_pI$ sont
scindées en tant que morphismes de $B$-modules gradués, nous voyons que
$$
0 \to \Hom_\mathcal{B}(N', F_pI/F_{p + 1}I) \to
\Hom_\mathcal{B}(N', I/F_{p + 1}I) \to
\Hom_\mathcal{B}(N', I/F_pI) \to 0
$$
est une suite exacte courte de $A$-modules différentiels gradués. Il
existe une suite analogue pour $N$ et $f$ induit un morphisme
de suites exactes courtes. Ainsi, par récurrence sur $p$ (en commençant par $p = 0$
lorsque $I/F_0I = 0$), nous concluons qu'il suffit de montrer que
l'application
$\Hom_\mathcal{B}(N', F_pI/F_{p + 1}I) \to \Hom_\mathcal{B}(N, F_pI/F_{p + 1}I)$ est un
quasi-isomorphisme. Puisque $F_pI/F_{p + 1}I$ est un produit de décalages de
$A^\vee$, il suffit de prouver que
$\Hom_\mathcal{B}(N', B^\vee[k]) \to \Hom_\mathcal{B}(N, B^\vee[k])$ est un
quasi-isomorphisme. D'après le lemme \ref{dga-lemma-hom-into-shift-dual-free}, il
suffit de montrer que $(N')^\vee \to N^\vee$ est un quasi-isomorphisme. Ceci est
vrai parce que $f$ est un quasi-isomorphisme et $(\ )^\vee$ est un
foncteur exact.
\end{proof}

\begin{lemma}
\label{dga-lemma-derived-restriction-exts}
Soient $(A, \text{d})$ et $(B, \text{d})$ des algèbres différentielles
graduées sur un anneau $R$. Soit $N$ un $(A, B)$-bimodule différentiel gradué.
Alors, pour tout $n \in \mathbf{Z}$, il existe des isomorphismes
$$
H^n(R\Hom(N, M)) = \Ext^n_{D(B, \text{d})}(N, M)
$$
de $R$-modules, fonctoriels en $M$. Ils sont également fonctoriels
en $N$ par rapport à l'opération décrite dans
le lemme \ref{dga-lemma-functoriality-derived-restriction}.
\end{lemma}

\begin{proof}
Dans la démonstration du lemme
\ref{dga-lemma-derived-restriction}, nous avons vu
$$
R\Hom(N, M) =
\Hom_{\text{Mod}^{dg}_{(B, \text{d})}}(N, I)
$$
comme $A$-module différentiel gradué
où $M \to I$ est un quasi-isomorphisme de $M$ vers un $B$-module
différentiel gradué ayant la propriété (I). Par conséquent, ce complexe a les
bons modules de cohomologie selon le lemme \ref{dga-lemma-hom-derived}.
Nous omettons une discussion de la nature fonctorielle de ces
identifications.
\end{proof}

\begin{lemma}
\label{dga-lemma-compute-derived-restriction}
Soit $R$ un anneau. Soient $(A, \text{d})$ et $(B, \text{d})$ des
$R$-algèbres différentielles graduées. Soit $N$ un $(A,
B)$-bimodule différentiel gradué. Si
$\Hom_{D(B, \text{d})}(N, N') = \Hom_{K(\text{Mod}_{(B, \text{d})})}(N, N')$ pour
tout $N' \in K(B, \text{d})$, par exemple si $N$ a la
propriété (P) en tant que $B$-module différentiel gradué, alors
$$
R\Hom(N, M) = \Hom_{\text{Mod}^{dg}_{(B, \text{d})}}(N, M)
$$
de manière fonctorielle en $M$ dans $D(B, \text{d})$.
\end{lemma}

\begin{proof}
Par construction (lemme \ref{dga-lemma-derived-restriction}) pour
trouver $R\Hom(N, M)$ nous choisissons un
quasi-isomorphisme $M \to I$ où $I$ est un $B$-module différentiel gradué avec
la propriété (I) et nous posons
$R\Hom(N, M) = \Hom_{\text{Mod}^{dg}_{(B, \text{d})}}(N, I)$. Par
hypothèse, l'application
$$
\Hom_{\text{Mod}^{dg}_{(B, \text{d})}}(N, M) \longrightarrow
\Hom_{\text{Mod}^{dg}_{(B, \text{d})}}(N, I)
$$
induite par $M \to I$ est un quasi-isomorphisme, voir discussion dans
l'exemple \ref{dga-example-dgm-dg-cat}. Cela prouve le lemme. Si
$N$ a la propriété (P) en tant que $B$-module, alors nous voyons que
l'hypothèse est satisfaite par le lemme \ref{dga-lemma-hom-derived}.
\end{proof}




\section{Variante du Hom dérivé}
\label{dga-section-variant}

\noindent
Soit $\mathcal{A}$ une catégorie abélienne. Considérons la catégorie
différentielle graduée $\text{Comp}^{dg}(\mathcal{A})$ des complexes de $\mathcal{A}$, voir
l'exemple \ref{dga-example-category-complexes}.
Soit $K^\bullet$ un complexe de $\mathcal{A}$. Posons
$$
(E, \text{d}) = \Hom_{\text{Comp}^{dg}(\mathcal{A})}(K^\bullet, K^\bullet)
$$
et considérons le foncteur de catégories différentielles graduées
$$
\text{Comp}^{dg}(\mathcal{A}) \longrightarrow \text{Mod}^{dg}_{(E, \text{d})},
\quad
X^\bullet
\longmapsto
\Hom_{\text{Comp}^{dg}(\mathcal{A})}(K^\bullet, X^\bullet)
$$
du lemme \ref{dga-lemma-construction}.

\begin{lemma}
\label{dga-lemma-existence-of-derived}
Dans la situation ci-dessus, si le foncteur dérivé à droite $R\Hom(K^\bullet,
-)$ de $\Hom(K^\bullet, -) : K(\mathcal{A}) \to D(\textit{Ab})$
est défini partout sur $D(\mathcal{A})$, alors nous obtenons un foncteur exact
canonique
$$
R\Hom(K^\bullet, -) : D(\mathcal{A}) \longrightarrow D(E, \text{d})
$$
de catégories triangulées qui se réduit au foncteur usuel lorsqu'on prend
les complexes associés de groupes abéliens.
\end{lemma}

\begin{proof}
Notons que nous avons un foncteur associé
$K(\mathcal{A}) \to K(\text{Mod}_{(E, \text{d})})$ d'après
le lemme \ref{dga-lemma-construction}.
Nous affirmons que ce foncteur est un foncteur exact de catégories
triangulées. En effet, soit $f : A^\bullet \to B^\bullet$ un morphisme de
complexes de $\mathcal{A}$. Alors un calcul montre que
$$
\Hom_{\text{Comp}^{dg}(\mathcal{A})}(K^\bullet, C(f)^\bullet)
=
C\left(
\Hom_{\text{Comp}^{dg}(\mathcal{A})}(K^\bullet, A^\bullet) \to
\Hom_{\text{Comp}^{dg}(\mathcal{A})}(K^\bullet, B^\bullet)
\right)
$$
où le membre de droite est le cône dans $\text{Mod}_{(E, \text{d})}$
défini plus tôt dans ce chapitre. Cela
montre que notre foncteur est compatible avec les cônes, donc avec
les triangles distingués. Soit $X^\bullet$ un objet de $K(\mathcal{A})$.
Considérons la catégorie des quasi-isomorphismes $s : X^\bullet \to Y^\bullet$. Par
hypothèse, le foncteur $(s :
X^\bullet \to Y^\bullet) \mapsto \Hom_\mathcal{A}(K^\bullet, Y^\bullet)$ est
essentiellement constant lorsqu'il est vu dans $D(\textit{Ab})$. Mais
puisque le foncteur d'oubli $D(E, \text{d}) \to D(\textit{Ab})$ est
compatible avec la prise de cohomologie, la même chose est vraie dans $D(E,
\text{d})$. Cela prouve le lemme.
\end{proof}

\noindent
{\bf Avertissement :} Bien que le lemme soit vrai et puisse être
utile sous cette forme, l'algèbre différentielle $E$ n'est la « bonne » que si
$H^n(E) = \Ext^n_{D(\mathcal{A})}(K^\bullet, K^\bullet)$
pour tout $n \in \mathbf{Z}$.





\section{Produit tensoriel dérivé}
\label{dga-section-base-change}

\noindent
Cette section est analogue à Compléments d'algèbre,
section \ref{more-algebra-section-derived-base-change}.

\medskip\noindent
Soit $R$ un anneau. Soient $(A, \text{d})$ et $(B, \text{d})$ des
algèbres différentielles graduées sur $R$. Soit $N$ un
$(A, B)$-bimodule différentiel gradué. Considérons le foncteur
\begin{equation}
\label{dga-equation-bc}
\text{Mod}_{(A, \text{d})}
\longrightarrow
\text{Mod}_{(B, \text{d})},\quad M
\longmapsto M \otimes_A N
\end{equation}
défini dans la section \ref{dga-section-bimodules-tensor}.

\begin{lemma}
\label{dga-lemma-bc-homotopy}
Le foncteur (\ref{dga-equation-bc}) définit un foncteur
exact de catégories triangulées
$K(\text{Mod}_{(A, \text{d})}) \to K(\text{Mod}_{(B, \text{d})})$.
\end{lemma}

\begin{proof}
Via le lemme \ref{dga-lemma-tensor}
et la remarque
\ref{dga-remark-shift-tensor-no-sign}, cela découle du principe général du
lemme \ref{dga-lemma-functor-between-ABC}.
\end{proof}

\noindent
À ce stade, nous pouvons considérer le diagramme
$$
\xymatrix{
K(\text{Mod}_{(A, \text{d})}) \ar[d] \ar[rr]_{- \otimes_A N} \ar[rrd]_F & &
K(\text{Mod}_{(B, \text{d})}) \ar[d] \\
D(A, \text{d}) \ar@{..>}[rr] & &
D(B, \text{d})
}
$$
La flèche en pointillés que nous allons construire ci-dessous
sera le {\it foncteur dérivé à gauche} de la composition
$F$. ({\it Attention} : le diagramme ne commute
pas.) Plus précisément, dans le cadre
général de Catégories dérivées, section
\ref{derived-section-derived-functors}, nous voulons
calculer le foncteur dérivé à gauche de $F$ par rapport au système multiplicatif de
quasi-isomorphismes dans $K(\text{Mod}_{(A, \text{d})})$.

\begin{lemma}
\label{dga-lemma-derived-bc}
Dans la situation ci-dessus, le foncteur dérivé à gauche de $F$
existe. Nous le
notons $- \otimes_A^\mathbf{L} N : D(A, \text{d}) \to D(B, \text{d})$.
\end{lemma}

\begin{proof}
Nous
utiliserons Catégories dérivées, lemme \ref{derived-lemma-find-existence-computes} pour
le prouver. Comme collection $\mathcal{P}$ d'objets,
nous prendrons les objets ayant la propriété (P). La propriété
(1) a été démontrée dans le lemme \ref{dga-lemma-resolve}. La
propriété (2) est vraie parce que si $s : P \to P'$ est un quasi-isomorphisme de
modules avec la propriété (P), alors $s$ est une équivalence d'homotopie selon
le lemme \ref{dga-lemma-hom-derived}.
\end{proof}

\begin{lemma}
\label{dga-lemma-functoriality-bc}
Soit $R$ un anneau. Soient $(A, \text{d})$ et $(B, \text{d})$ des
$R$-algèbres différentielles graduées. Soit $f : N \to N'$ un
homomorphisme de bimodules différentiels gradués sur $(A, B)$.
Alors $f$ induit un morphisme de foncteurs
$$
1\otimes f :
- \otimes_A^\mathbf{L} N
\longrightarrow -
\otimes_A^\mathbf{L} N'
$$
Si $f$ est un quasi-isomorphisme, alors $1 \otimes f$ est un isomorphisme de
foncteurs.
\end{lemma}

\begin{proof}
Soit $M$ un $A$-module différentiel gradué avec la propriété (P).
Alors $1 \otimes f : M \otimes_A N \to M \otimes_A N'$ est une
application de $B$-modules différentiels gradués. De plus, cela est fonctoriel par
rapport à $M$. Étant donné que les foncteurs $-
\otimes_A^\mathbf{L} N$ et $- \otimes_A^\mathbf{L} N'$ sont calculés par
le produit tensoriel des objets avec la propriété (P)
(lemme \ref{dga-lemma-derived-bc}), nous obtenons une transformation de foncteurs comme
indiqué.

\medskip\noindent
Supposons que $f$ soit un quasi-isomorphisme. Soit $F_\bullet$ la
filtration donnée sur $M$. Observons que $M
\otimes_A N = \colim F_i(M) \otimes_A N$ et $M \otimes_A N'
= \colim F_i(M) \otimes_A N'$. Par conséquent, il
suffit de montrer que $F_n(M) \otimes_A
N \to F_n(M) \otimes_A N'$ est un
quasi-isomorphisme (les limites inductives filtrantes sont exactes, voir Algèbre,
lemme \ref{algebra-lemma-directed-colimit-exact}). Comme les
inclusions $F_n(M) \to F_{n + 1}(M)$ sont scindées en tant que
morphismes de $A$-modules gradués, nous voyons que
$$
0 \to F_n(M) \otimes_A N \to F_{n + 1}(M) \otimes_A N \to
F_{n + 1}(M)/F_n(M) \otimes_A N \to 0
$$
est une suite exacte courte de $B$-modules différentiels gradués. Il
existe une suite analogue pour $N'$ et $f$ induit un morphisme de
suites exactes courtes. Ainsi, par récurrence sur $n$ (en commençant par $n = -1$
lorsque $F_{-1}(M) = 0$), nous concluons qu'il suffit de montrer que
l'application $F_{n + 1}(M)/F_n(M) \otimes_A N \to F_{n + 1}(M)/F_n(M) \otimes_A N'$ est un
quasi-isomorphisme. C'est vrai parce que $F_{n + 1}(M)/F_n(M)$ est une
somme directe de décalages de $A$ et le résultat est vrai pour $A[k]$ puisque
$f : N \to N'$ est un quasi-isomorphisme.
\end{proof}

\begin{lemma}
\label{dga-lemma-compute-bc}
Soit $R$ un anneau. Soient $(A, \text{d})$ et $(B, \text{d})$ des
$R$-algèbres différentielles graduées. Soit $N$ un $(A, B)$-bimodule
différentiel gradué qui possède la propriété (P) en tant que $A$-module différentiel
gradué à gauche. Alors $M \otimes_A^\mathbf{L} N$ est calculé
par $M \otimes_A N$ pour tout $A$-module différentiel gradué $M$.
\end{lemma}

\begin{proof}
Soit $f : M \to M'$ un homomorphisme de $A$-modules différentiels gradués qui
est un quasi-isomorphisme. Nous affirmons que $f \otimes \text{id} : M
\otimes_A N \to M' \otimes_A N$ est un quasi-isomorphisme. Si cela est
vrai, alors par la construction du produit tensoriel dérivé dans la
démonstration du lemme \ref{dga-lemma-derived-bc}, nous obtenons le résultat souhaité. La
construction de l'application $f \otimes \text{id}$ ne dépend que de la
structure de $A$-module différentiel gradué à gauche de $N$. De
plus, nous avons $M \otimes_A N = N \otimes_{A^{opp}} M = N
\otimes_{A^{opp}}^\mathbf{L} M$ parce que $N$ possède la propriété (P) comme
$A^{opp}$-module différentiel gradué. Par conséquent, l'énoncé découle du
lemme \ref{dga-lemma-functoriality-bc}.
\end{proof}

\begin{lemma}
\label{dga-lemma-tensor-hom-adjoint}
Soit $R$ un anneau.
Soient $(A, \text{d})$ et $(B, \text{d})$ des $R$-algèbres différentielles
graduées. Soit $N$ un $(A, B)$-bimodule différentiel gradué.
Alors le foncteur
$$
- \otimes_A^\mathbf{L} N : D(A, \text{d}) \longrightarrow D(B, \text{d})
$$
du lemme \ref{dga-lemma-derived-bc} est un adjoint à gauche du foncteur
$$
R\Hom(N, -) : D(B, \text{d}) \longrightarrow D(A, \text{d})
$$
du lemme \ref{dga-lemma-derived-restriction}.
\end{lemma}

\begin{proof}
Cela découle de Catégories dérivées, lemme
\ref{derived-lemma-pre-derived-adjoint-functors-general} et
du fait que $- \otimes_A N$ et
$\Hom_{\text{Mod}^{dg}_{(B, \text{d})}}(N, -)$ sont adjoints d'après le
lemme \ref{dga-lemma-tensor-hom-adjunction}.
\end{proof}

\begin{example}
\label{dga-example-map-hom-tensor}
Soit $R$ un anneau. Soit $(A, \text{d}) \to (B, \text{d})$ un
homomorphisme d'algèbres différentielles graduées sur $R$. Alors nous pouvons
considérer $B$ comme un $(A, B)$-bimodule différentiel gradué et nous obtenons un foncteur
$$
- \otimes_A B : D(A, \text{d}) \longrightarrow D(B, \text{d})
$$
D'après le lemme \ref{dga-lemma-tensor-hom-adjoint}, le foncteur adjoint à
gauche de celui-ci est le foncteur $R\Hom(B, -)$. Pour un $B$-module
différentiel gradué, notons $N_A$ le $A$-module différentiel gradué obtenu à partir
de $N$ par restriction via $A \to B$. Alors nous avons clairement
un isomorphisme canonique
$$
\Hom_{\text{Mod}^{dg}_{(B, \text{d})}}(B, N) \longrightarrow N_A,\quad
f \longmapsto f(1)
$$
fonctoriel par rapport au $B$-module $N$. Ainsi, nous voyons
que $R\Hom(B, -)$ est le foncteur de restriction et nous obtenons
$$
\Hom_{D(A, \text{d})}(M, N_A) =
\Hom_{D(B, \text{d})}(M \otimes^\mathbf{L}_A B, N)
$$
bifonctoriellement en $M$ et $N$ exactement comme dans le cas des anneaux
commutatifs. Enfin, observons que la restriction est également un foncteur
tensoriel, puisque $N_A = N \otimes_B {}_BB_A = N \otimes_B^\mathbf{L}
{}_BB_A$ où ${}_BB_A$ est $B$ vu comme un $(B, A)$-bimodule différentiel gradué.
\end{example}

\begin{lemma}
\label{dga-lemma-tensor-with-compact-fully-faithful}
Avec les notations et sous les hypothèses du lemme \ref{dga-lemma-tensor-hom-adjoint}, faisons les hypothèses
suivantes :
\begin{enumerate}
\item $N$ définit un objet compact de $D(B, \text{d})$, et
\item l'application $H^k(A) \to \Hom_{D(B, \text{d})}(N, N[k])$ est
un isomorphisme pour tout $k \in \mathbf{Z}$.
\end{enumerate}
Alors le foncteur $-\otimes_A^\mathbf{L} N$ est pleinement fidèle.
\end{lemma}

\begin{proof}
Notre foncteur a un adjoint à gauche donné
par $R\Hom(N, -)$ d'après le lemme
\ref{dga-lemma-tensor-hom-adjoint}. D'après Catégories, lemme \ref{categories-lemma-adjoint-fully-faithful},
il suffit de montrer que pour un $A$-module différentiel gradué $M$,
l'application
$$
M \longrightarrow R\Hom(N, M \otimes_A^\mathbf{L} N)
$$
est un isomorphisme dans $D(A, \text{d})$. Pour cela, il suffit de montrer que
$$
H^n(M) \longrightarrow
\text{Ext}^n_{D(B, \text{d})}(N, M \otimes_A^\mathbf{L} N)
$$
est un isomorphisme, voir le lemme \ref{dga-lemma-derived-restriction-exts}.
Puisque $N$ est un objet compact, le membre de droite commute avec les
sommes directes. Ainsi, d'après la remarque \ref{dga-remark-P-resolution}, il
suffit de prouver que cette application est un isomorphisme pour $M = A[k]$.
Puisque $(A[k] \otimes_A^\mathbf{L} N) = N[k]$ d'après la
remarque \ref{dga-remark-shift-tensor-no-sign},
l'hypothèse (2) sur $N$ affirme précisément le résultat dans ces cas.
\end{proof}

\begin{lemma}
\label{dga-lemma-base-change-K-flat}
Soit $R \to R'$ un homomorphisme d'anneaux. Soit $(A, \text{d})$ une $R$-algèbre
différentielle graduée. Soit $(A', \text{d})$ le changement de base,
c'est-à-dire $A' = A \otimes_R R'$. Si $A$ est K-plat en tant que complexe de $R$-modules,
alors
\begin{enumerate}
\item $- \otimes_A^\mathbf{L} A' : D(A, \text{d}) \to D(A', \text{d})$
est égal au foncteur dérivé à droite de
$$
K(A, \text{d}) \longrightarrow K(A', \text{d}),\quad
M \longmapsto M \otimes_R R'
$$
\item le diagramme
$$
\xymatrix{
D(A, \text{d}) \ar[r]_{- \otimes_A^\mathbf{L} A'} \ar[d]_{restriction} &
D(A', \text{d}) \ar[d]^{restriction} \\
D(R) \ar[r]^{- \otimes_R^\mathbf{L} R'} & D(R')
}
$$
commute, et
\item si $M$ est K-plat en tant que complexe de $R$-modules,
alors le $A'$-module différentiel gradué $M \otimes_R R'$
représente $M \otimes_A^\mathbf{L} A'$.
\end{enumerate}
\end{lemma}

\begin{proof}
Pour tout $A$-module différentiel gradué $M$, il existe une application canonique
$$
c_M : M \otimes_R R' \longrightarrow M \otimes_A A'
$$
Soit $P$ un $A$-module différentiel gradué avec la propriété (P).
Nous affirmons que $c_P$ est un isomorphisme et que $P$ est K-plat en
tant que complexe de $R$-modules. Cela prouvera tous les résultats
énoncés dans le lemme par des arguments formels utilisant la
définition du produit tensoriel dérivé dans le lemme \ref{dga-lemma-derived-bc}
et Compléments d'algèbre, section \ref{more-algebra-section-derived-tensor-product}.

\medskip\noindent
Soit $F_\bullet$ la filtration sur $P$ montrant que $P$ a la propriété (P). Notons que
$c_A$ est un isomorphisme et que $A$ est K-plat en tant que complexe de
$R$-modules par hypothèse. Par conséquent, il en va de même pour les sommes
directes de décalages de $A$ (vous pouvez utiliser
Compléments d'algèbre, lemme \ref{more-algebra-lemma-colimit-K-flat} pour
traiter les sommes directes si vous le souhaitez).
Par conséquent, cela vaut pour les complexes $F_{p + 1}P/F_pP$.
Comme les suites exactes courtes
$$
0 \to F_pP \to F_{p + 1}P \to F_{p + 1}P/F_pP \to 0
$$
sont scindées comme suites de modules gradués, nous pouvons montrer
par récurrence que $c_{F_pP}$ est un isomorphisme pour tout $p$
et que $F_pP$ est K-plat en tant que complexe de $R$-modules (voir
Compléments d'algèbre, lemme \ref{more-algebra-lemma-K-flat-two-out-of-three}).
Enfin, en utilisant que $P = \colim F_pP$, nous concluons que
$c_P$ est un isomorphisme et que $P$ est K-plat en tant
que complexe de $R$-modules (voir
Compléments d'algèbre, lemme \ref{more-algebra-lemma-colimit-K-flat}).
\end{proof}

\begin{lemma}
\label{dga-lemma-RHom-is-tensor}
Soit $R$ un anneau.
Soient $(A, \text{d})$ et $(B, \text{d})$ des $R$-algèbres différentielles graduées. Soit $T$ un
$(A, B)$-bimodule différentiel gradué. Faisons les hypothèses
suivantes :
\begin{enumerate}
\item $T$ définit un objet compact de $D(B, \text{d})$, et
\item $S = \Hom_{\text{Mod}^{dg}_{(B, \text{d})}}(T, B)$
représente $R\Hom(T, B)$ dans $D(A, \text{d})$.
\end{enumerate}
Alors $S$ a une structure de $(B, A)$-bimodule différentiel gradué
et il existe un isomorphisme
$$
N \otimes_B^\mathbf{L} S \longrightarrow R\Hom(T, N)
$$
fonctorielle en $N$ dans $D(B, \text{d})$.
\end{lemma}

\begin{proof}
Notons $\mathcal{B} = \text{Mod}^{dg}_{(B, \text{d})}$. La
structure de $A$-module à droite de $S$ provient de l'application
$A \to \Hom_\mathcal{B}(T, T)$ et de la composition
$\Hom_\mathcal{B}(T, B) \otimes \Hom_\mathcal{B}(T, T) \to
\Hom_\mathcal{B}(T, B)$ définie dans l'exemple \ref{dga-example-dgm-dg-cat}. En utilisant cette
multiplication une seconde fois, on obtient une application
$$
c_N :
N \otimes_B S = \Hom_\mathcal{B}(B, N) \otimes_B \Hom_\mathcal{B}(T, B)
\longrightarrow
\Hom_\mathcal{B}(T, N)
$$
fonctoriel en $N$. Étant donné $N$, nous pouvons choisir des
quasi-isomorphismes $P \to N \to I$ où $P$, resp.\ $I$ est un $B$-module différentiel gradué avec
la propriété (P), resp.\ (I). Ensuite, en utilisant $c_N$, nous obtenons
une application $P \otimes_B S \to \Hom_\mathcal{B}(T, I)$ entre les
objets représentant $S \otimes_B^\mathbf{L} N$ et $R\Hom(T, N)$. De toute
évidence, cela définit une transformation de foncteurs $c$ comme dans le lemme.

\medskip\noindent
Pour prouver que $c$ est un isomorphisme de foncteurs, nous
pouvons supposer que $N$ est un $B$-module différentiel gradué qui
possède la propriété (P). Puisque $T$ définit un objet compact
dans $D(B, \text{d})$ et puisque les deux côtés de la flèche définissent
des foncteurs exacts de catégories triangulées, nous nous ramenons, à
l'aide du lemme
\ref{dga-lemma-property-P-sequence}, au cas où $N$ possède une filtration finie dont les
quotients successifs sont des sommes directes de $B[k]$. En utilisant à
nouveau que les deux côtés de la flèche sont des foncteurs exacts de
catégories triangulées et la compacité de $T$, nous nous
ramenons au cas $N = B[k]$. L'hypothèse (2) est exactement
l'hypothèse que $c$ est un isomorphisme dans ce cas.
\end{proof}






\section{Composition des produits tensoriels dérivés}
\label{dga-section-compose-tensor-functors}

\noindent
Nous encourageons le lecteur à passer cette section.

\medskip\noindent
Soit $R$ un anneau. Soient $(A, \text{d})$, $(B, \text{d})$, et
$(C, \text{d})$ des $R$-algèbres différentielles graduées.
Soit $N$ un $(A, B)$-bimodule différentiel gradué. Soit $N'$
un $(B, C)$-module différentiel gradué. Nous notons
$N_B$ le bimodule $N$ considéré comme un $B$-module différentiel gradué (en
oubliant la structure de module sur $A$). Il existe une application canonique
\begin{equation}
\label{dga-equation-plain-versus-derived}
N_B \otimes_B^\mathbf{L} N'
\longrightarrow
(N \otimes_B N')_C
\end{equation}
dans $D(C, \text{d})$. Ici, $(N \otimes_B N')_C$ désigne le $(A,
C)$-bimodule $N \otimes_B N'$ considéré comme un $C$-module
différentiel gradué. En effet, cette application découle du
fait qu'il existe toujours un morphisme du produit tensoriel dérivé vers le produit
tensoriel ordinaire (comme il s'agit d'un foncteur dérivé à gauche).

\begin{lemma}
\label{dga-lemma-compose-tensor-functors-general}
Soit $R$ un anneau. Soient $(A, \text{d})$, $(B, \text{d})$ et
$(C, \text{d})$ des algèbres différentielles graduées sur
$R$. Soit $N$ un $(A, B)$-bimodule différentiel gradué.
Soit $N'$ un module différentiel gradué $(B, C)$.
Supposons que (\ref{dga-equation-plain-versus-derived}) soit un isomorphisme.
Alors la composition
$$
\xymatrix{
D(A, \text{d}) \ar[rr]^{- \otimes_A^\mathbf{L} N} & &
D(B, \text{d}) \ar[rr]^{- \otimes_B^\mathbf{L} N'} & &
D(C, \text{d})
}
$$
est isomorphe à $- \otimes_A^\mathbf{L} N''$ avec
$N'' = N \otimes_B N'$ considéré comme $(A, C)$-bimodule.
\end{lemma}

\begin{proof}
Définissons une transformation de foncteurs
$$
(- \otimes_A^\mathbf{L} N) \otimes_B^\mathbf{L} N'
\longrightarrow
- \otimes_A^\mathbf{L} N''
$$
Pour ce faire,
soit $M$ un $A$-module différentiel gradué avec la propriété
(P). Selon la construction du foncteur $- \otimes_A^\mathbf{L} N''$ de la
démonstration du lemme \ref{dga-lemma-derived-bc}, le produit tensoriel
ordinaire $M \otimes_A N''$ représente $M \otimes_A^\mathbf{L} N''$ dans $D(C,
\text{d})$. Ensuite, nous écrivons
$$
M \otimes_A N'' =
M \otimes_A (N \otimes_B N') =
(M \otimes_A N) \otimes_B N'
$$
Le module $M \otimes_A N$ représente $M \otimes_A^\mathbf{L} N$ dans
$D(B, \text{d})$. Choisissons un quasi-isomorphisme $Q \to M \otimes_A N$ où $Q$
est un $B$-module différentiel gradué avec la propriété (P). Alors $Q
\otimes_B N'$ représente $(M
\otimes_A^\mathbf{L} N) \otimes_B^\mathbf{L} N'$ dans $D(C, \text{d})$.
Ainsi, nous pouvons
définir notre application via
$$
(M \otimes_A^\mathbf{L} N) \otimes_B^\mathbf{L} N' =
Q \otimes_B N' \to
M \otimes_A N \otimes_B N' = M
\otimes_A^\mathbf{L} N''
$$
La construction de cette application est fonctorielle en $M$ et compatible
avec les triangles distingués et les sommes directes ; nous omettons les
détails. Considérons la propriété $T$ des objets $M$ de $D(A, \text{d})$
exprimant que cette application est un isomorphisme. Alors
\begin{enumerate}
\item si $T$ est vraie pour $M_i$, alors $T$ est vraie pour $\bigoplus M_i$,
\item si $T$ est vraie pour $2$ des $3$ objets d'un triangle
distingué, alors elle est vraie pour le troisième, et
\item $T$ vaut pour $A[k]$ car ici nous obtenons un décalage de
l'application (\ref{dga-equation-plain-versus-derived}) que nous
avons supposée être un isomorphisme.
\end{enumerate}
Ainsi, d'après la remarque \ref{dga-remark-P-resolution}, la propriété $T$
est toujours vraie et la démonstration est complète.
\end{proof}

\noindent
Soit $R$ un anneau. Soient $(A, \text{d})$, $(B, \text{d})$ et
$(C, \text{d})$ des $R$-algèbres différentielles graduées.
Nous notons temporairement $(A \otimes_R B)_B$ l'algèbre
différentielle graduée $A \otimes_R B$ vue comme un module différentiel gradué à
droite sur $B$, et ${}_B(B \otimes_R C)_C$ l'algèbre différentielle
graduée $B \otimes_R C$ vue comme un bimodule différentiel gradué sur
$(B, C)$. Alors il existe une application canonique
\begin{equation}
\label{dga-equation-plain-versus-derived-algebras}
(A \otimes_R B)_B \otimes_B^\mathbf{L} {}_B(B \otimes_R C)_C
\longrightarrow
(A \otimes_R B \otimes_R C)_C
\end{equation}
dans $D(C, \text{d})$ où $(A \otimes_R B \otimes_R C)_C$
désigne la $R$-algèbre
différentielle graduée $A \otimes_R B \otimes_R C$ considérée comme un
$C$-module différentiel gradué (à droite). En effet, cette
application provient de l'identification
$$
(A \otimes_R B)_B \otimes_B {}_B(B \otimes_R C)_C =
(A \otimes_R B \otimes_R C)_C
$$
et du fait que le produit tensoriel dérivé s'envoie toujours dans le
produit tensoriel ordinaire (puisque c'est un foncteur dérivé à gauche).

\begin{lemma}
\label{dga-lemma-compose-tensor-functors-general-algebra}
Soit $R$ un anneau. Soient $(A, \text{d})$, $(B, \text{d})$ et
$(C, \text{d})$ des $R$-algèbres différentielles graduées.
Supposons que (\ref{dga-equation-plain-versus-derived-algebras}) soit un
isomorphisme. Soit $N$ un $(A, B)$-bimodule différentiel gradué.
Soit $N'$ un $(B, C)$-bimodule différentiel gradué.
Alors la composition
$$
\xymatrix{
D(A, \text{d}) \ar[rr]^{- \otimes_A^\mathbf{L} N} & &
D(B, \text{d}) \ar[rr]^{- \otimes_B^\mathbf{L} N'} & &
D(C, \text{d})
}
$$
est isomorphe à $- \otimes_A^\mathbf{L} N''$ pour un $(A, C)$-bimodule
différentiel gradué $N''$ décrit dans la démonstration.
\end{lemma}

\begin{proof}
Selon le lemme \ref{dga-lemma-functoriality-bc}, nous pouvons remplacer $N$ et $N'$ par
des bimodules quasi-isomorphes. Ainsi, nous pouvons supposer que $N$, resp.\ $N'$,
possède la propriété (P) en tant que bimodule différentiel gradué $(A, B)$,
resp.\ $(B, C)$,
voir le lemme \ref{dga-lemma-bimodule-resolve}. Nous affirmons que le
lemme est valable avec le $(A, C)$-bimodule $N'' = N \otimes_B
N'$. Pour le prouver, il suffit de montrer que
$$
N_B \otimes_B^\mathbf{L} N' \longrightarrow (N \otimes_B N')_C
$$
est un isomorphisme dans $D(C, \text{d})$,
voir le lemme \ref{dga-lemma-compose-tensor-functors-general}.

\medskip\noindent
Soit $F_\bullet$ la filtration sur $N$ comme dans la propriété (P) pour les bimodules.
D'après le lemme
\ref{dga-lemma-bimodule-property-P-sequence}, il existe une suite exacte courte
$$
0 \to
\bigoplus\nolimits F_iN \to
\bigoplus\nolimits F_iN \to N \to 0
$$
de $(A, B)$-bimodules différentiels gradués, qui est scindée comme suite
de $(A, B)$-bimodules gradués. A fortiori, il s'agit d'une suite exacte
courte admissible de $B$-modules différentiels gradués et cela produit un triangle
distingué
$$
\bigoplus\nolimits F_iN_B \to
\bigoplus\nolimits F_iN_B \to N_B \to
\bigoplus\nolimits F_iN_B[1]
$$
dans $D(B, \text{d})$. En utilisant que $- \otimes_B^\mathbf{L} N'$
est un foncteur exact de catégories triangulées et commute avec les
sommes directes et en utilisant que $- \otimes_B N'$ transforme les
suites exactes admissibles en suites exactes admissibles et commute
avec les sommes directes, nous nous ramenons à prouver
que
$$
(F_pN)_B \otimes_B^\mathbf{L} N' \longrightarrow (F_pN)_B \otimes_B N'
$$
est un quasi-isomorphisme pour tout $p$. En répétant
l'argument avec les suites exactes courtes de $(A, B)$-bimodules
$$
0 \to F_pN \to F_{p + 1}N \to F_{p + 1}N/F_pN \to 0
$$
qui sont scindées comme suites de $(A, B)$-bimodules
gradués, nous nous ramenons à montrer la même assertion pour $F_{p +
1}N/F_pN$. Comme ces modules sont des sommes directes de décalages de $(A \otimes_R B)_B$,
nous nous ramenons à montrer que
$$
(A \otimes_R B)_B \otimes_B^\mathbf{L} N'
\longrightarrow
(A \otimes_R B)_B \otimes_B N'
$$
est un quasi-isomorphisme.

\medskip\noindent
Choisissons une filtration $F_\bullet$ sur $N'$ comme dans la propriété (P) pour
les bimodules. Choisissons un quasi-isomorphisme $P \to (A
\otimes_R B)_B$ de $B$-modules différentiels gradués où $P$ a la propriété
(P). Nous devons montrer
que $P \otimes_B N' \to (A \otimes_R B)_B \otimes_B N'$ est un
quasi-isomorphisme parce que $P \otimes_B N'$ représente $(A
\otimes_R B)_B \otimes_B^\mathbf{L} N'$ dans $D(C, \text{d})$ par la
construction dans le lemme \ref{dga-lemma-derived-bc}. Puisque $N' =
\colim F_pN'$, nous constatons qu'il
suffit de montrer que $P \otimes_B
F_pN' \to (A \otimes_R B)_B \otimes_B F_pN'$ est un
quasi-isomorphisme. En utilisant les suites exactes courtes $0 \to F_pN' \to
F_{p + 1}N' \to F_{p + 1}N'/F_pN' \to 0$ qui sont scindées comme
suites de $(B, C)$-bimodules gradués, nous nous ramenons à montrer que $P
\otimes_B F_{p + 1}N'/F_pN' \to (A
\otimes_R B)_B \otimes_B F_{p + 1}N'/F_pN'$ est un
quasi-isomorphisme pour tout $p$. Ensuite, en
utilisant enfin que $F_{p + 1}N'/F_pN'$ est une somme
directe de décalages de ${}_B(B \otimes_R C)_C$, nous
concluons qu'il suffit de montrer que
$$
P \otimes_B {}_B(B \otimes_R C)_C \to
(A \otimes_R B)_B \otimes_B {}_B(B \otimes_R C)_C
$$
est un quasi-isomorphisme. Puisque $P \to (A \otimes_R B)_B$
est une résolution par un module satisfaisant à la propriété
(P), cette application de $C$-modules différentiels
gradués représente le morphisme (\ref{dga-equation-plain-versus-derived-algebras}) dans
$D(C, \text{d})$ et la démonstration est terminée.
\end{proof}

\begin{lemma}
\label{dga-lemma-compose-tensor-functors}
Soit $R$ un anneau. Soient $(A, \text{d})$, $(B, \text{d})$ et
$(C, \text{d})$ des $R$-algèbres différentielles
graduées. Si $C$ est K-plat comme complexe de $R$-modules,
alors (\ref{dga-equation-plain-versus-derived-algebras})
est un isomorphisme et la conclusion du
lemme \ref{dga-lemma-compose-tensor-functors-general-algebra} est valable.
\end{lemma}

\begin{proof}
Choisissons un quasi-isomorphisme $P \to (A \otimes_R B)_B$ de $B$-modules
différentiels gradués, où $P$ a la propriété (P). Ensuite, nous devons montrer
que
$$
P \otimes_B (B \otimes_R C) \longrightarrow
(A \otimes_R B) \otimes_B (B \otimes_R C)
$$
est un quasi-isomorphisme. De manière équivalente, il s'agit du morphisme
$$
P \otimes_R C \longrightarrow
A \otimes_R B \otimes_R C
$$
C'est un quasi-isomorphisme si $C$ est K-plat en tant que complexe de $R$-modules
selon Compléments d'algèbre, lemme \ref{more-algebra-lemma-K-flat-quasi-isomorphism}.
\end{proof}





\section{Variante du produit tensoriel dérivé}
\label{dga-section-variant-base-change}

\noindent
Soit $(\mathcal{C}, \mathcal{O})$ un site annelé. Alors nous avons les foncteurs
$$
\text{Comp}(\mathcal{O}) \to K(\mathcal{O}) \to D(\mathcal{O})
$$
et comme nous l'avons vu ci-dessus, nous avons un enrichissement
différentiel gradué $\text{Comp}^{dg}(\mathcal{O})$. Plus précisément, il s'agit
de la catégorie différentielle graduée de l'exemple
\ref{dga-example-category-complexes} associée à la catégorie abélienne
$\textit{Mod}(\mathcal{O})$. Soit $K^\bullet$ un complexe de $\mathcal{O}$-modules, en d'autres
termes, un objet de $\text{Comp}^{dg}(\mathcal{O})$. Posons
$$
(E, \text{d}) =
\Hom_{\text{Comp}^{dg}(\mathcal{O})}(K^\bullet, K^\bullet)
$$
Ceci est une $\mathbf{Z}$-algèbre différentielle graduée. Nous affirmons qu'il
existe un analogue du changement de base dérivé dans cette situation.

\begin{lemma}
\label{dga-lemma-tensor-with-complex}
Dans la situation ci-dessus, il y a un foncteur
$$
- \otimes_E K^\bullet :
\text{Mod}^{dg}_{(E, \text{d})}
\longrightarrow
\text{Comp}^{dg}(\mathcal{O})
$$
de catégories différentielles graduées. Ce foncteur envoie $E$ sur $K^\bullet$ et
commute avec les sommes directes.
\end{lemma}

\begin{proof}
Soit $M$ un $E$-module différentiel gradué. Pour chaque objet $U$ de
$\mathcal{C}$, le complexe $K^\bullet(U)$ est un module différentiel gradué à
gauche sur $E$, ainsi qu'un module à droite sur $\mathcal{O}(U)$. Les
actions commutent, donc nous avons un bimodule.
Ainsi, d'après les constructions des
sections \ref{dga-section-tensor-product} et \ref{dga-section-bimodules}, nous
pouvons former le produit tensoriel
$$
M \otimes_E K^\bullet(U)
$$
qui est un $\mathcal{O}(U)$-module différentiel gradué, c'est-à-dire un complexe
de $\mathcal{O}(U)$-modules. Cette construction est fonctorielle par rapport à $U$,
nous pouvons donc faisceautiser pour obtenir un complexe de $\mathcal{O}$-modules
que nous notons
$$
M \otimes_E K^\bullet
$$
De plus, pour chaque $U$, la construction détermine un foncteur
$\text{Mod}^{dg}_{(E, \text{d})} \to \text{Comp}^{dg}(\mathcal{O}(U))$ de
catégories différentielles graduées selon le lemme \ref{dga-lemma-tensor}.
Il est donc clair que nous obtenons un foncteur tel qu'énoncé dans le lemme.
\end{proof}

\begin{lemma}
\label{dga-lemma-tensor-with-complex-homotopy}
Le foncteur du lemme \ref{dga-lemma-tensor-with-complex} définit un foncteur exact
de catégories triangulées
$K(\text{Mod}_{(E, \text{d})}) \to K(\mathcal{O})$.
\end{lemma}

\begin{proof}
Le foncteur induit un foncteur entre les catégories d'homotopie
selon le lemme
\ref{dga-lemma-functorial}. Nous devons montrer que $- \otimes_E K^\bullet$ transforme les triangles
distingués en triangles distingués. Supposons
que $0 \to K \to L \to M \to 0$ soit une suite exacte courte admissible de
$E$-modules différentiels gradués. Soit $s : M \to L$ un homomorphisme de
$E$-modules gradués qui est l'inverse à gauche de $L \to M$. Alors $s$ définit une
application $M \otimes_E K^\bullet \to L \otimes_E K^\bullet$ de
$\mathcal{O}$-modules gradués (c'est-à-dire respectant la structure de $\mathcal{O}$-module et
la graduation, mais pas les différentielles) qui est
l'inverse à gauche de $L \otimes_E K^\bullet \to M \otimes_E K^\bullet$. Ainsi nous
voyons que
$$
0 \to K \otimes_E K^\bullet \to L \otimes_E K^\bullet \to
M \otimes_E K^\bullet \to 0
$$
est une suite exacte courte scindée terme à terme de complexes,
c'est-à-dire qu'elle définit un triangle distingué dans $K(\mathcal{O})$.
\end{proof}

\begin{lemma}
\label{dga-lemma-tensor-with-complex-derived}
Le foncteur $K(\text{Mod}_{(E, \text{d})}) \to K(\mathcal{O})$ du
lemme \ref{dga-lemma-tensor-with-complex-homotopy} a une version dérivée à
gauche définie sur l'ensemble de $D(E, \text{d})$. Nous la
notons $- \otimes_E^\mathbf{L} K^\bullet : D(E, \text{d}) \to D(\mathcal{O})$.
\end{lemma}

\begin{proof}
Nous
utiliserons Catégories dérivées, lemme \ref{derived-lemma-find-existence-computes} pour
le prouver. Comme collection $\mathcal{P}$ d'objets,
nous prendrons les objets avec la propriété (P). La propriété (1)
a été démontrée dans le lemme \ref{dga-lemma-resolve}. La
propriété (2) est vraie parce que si $s : P \to P'$ est un quasi-isomorphisme de
modules avec la propriété (P), alors $s$ est une équivalence d'homotopie selon
le lemme \ref{dga-lemma-hom-derived}.
\end{proof}

\begin{lemma}
\label{dga-lemma-upgrade-tensor-with-complex-derived}
Soit $R$ un anneau. Soit $\mathcal{C}$ un site. Soit $\mathcal{O}$ un
faisceau de $R$-algèbres commutatives. Soit $K^\bullet$ un
complexe de $\mathcal{O}$-modules. Le
foncteur du lemme
\ref{dga-lemma-tensor-with-complex-derived} possède la propriété suivante :
Pour chaque $M$, $N$ dans $D(E, \text{d})$, il existe une flèche
canonique
$$
R\Hom(M, N)
\longrightarrow
R\Hom_\mathcal{O}(M \otimes_E^\mathbf{L} K^\bullet, N
\otimes_E^\mathbf{L} K^\bullet)
$$
dans $D(R)$ qui sur les modules de cohomologie donne les flèches
$$
\Ext^n_{D(E, \text{d})}(M, N) \to
\Ext^n_{D(\mathcal{O})}
(M \otimes_E^\mathbf{L} K^\bullet, N \otimes_E^\mathbf{L} K^\bullet)
$$
induites par le foncteur $- \otimes_E^\mathbf{L} K^\bullet$.
\end{lemma}

\begin{proof}
Le membre de droite est le Hom dérivé global introduit dans Cohomologie
sur les sites, section \ref{sites-cohomology-section-global-RHom} qui a les bons
modules de cohomologie. Pour le membre de
gauche, nous considérons $M$ comme un $(R, A)$-bimodule et nous avons le
$\Hom$ dérivé introduit dans la section \ref{dga-section-restriction}, qui possède
également les modules de cohomologie corrects. Pour
prouver le lemme, nous pouvons supposer que $M$ et $N$ sont des
$E$-modules différentiels gradués avec la propriété (P) ; cela ne change pas le
membre de gauche selon
le lemme \ref{dga-lemma-functoriality-derived-restriction}.
D'après le lemme
\ref{dga-lemma-compute-derived-restriction}, cela signifie que le membre de gauche devient
$\Hom_{\text{Mod}^{dg}_{(B, \text{d})}}(M, N)$. Dans les
lemmes \ref{dga-lemma-tensor-with-complex},
\ref{dga-lemma-tensor-with-complex-homotopy}, et
\ref{dga-lemma-tensor-with-complex-derived} nous avons
construit un foncteur
$$
- \otimes_E K^\bullet :
\text{Mod}^{dg}_{(E, \text{d})}
\longrightarrow
\text{Comp}^{dg}(\mathcal{O})
$$
de catégories différentielles graduées
et nous avons montré que $- \otimes_E^\mathbf{L} K^\bullet$ est calculé en évaluant
ce foncteur sur des $E$-modules
différentiels gradués ayant la propriété (P). Par conséquent,
nous obtenons un morphisme de complexes de $R$-modules
$$
\Hom_{\text{Mod}^{dg}_{(B, \text{d})}}(M, N)
\longrightarrow
\Hom_{\text{Comp}^{dg}(\mathcal{O})}
(M \otimes_E K^\bullet, N \otimes_E K^\bullet)
$$
Pour tous les complexes de $\mathcal{O}$-modules
$\mathcal{F}^\bullet$ et $\mathcal{G}^\bullet$, il existe une
application canonique
$$
\Hom_{\text{Comp}^{dg}(\mathcal{O})}
(\mathcal{F}^\bullet, \mathcal{G}^\bullet) =
\Gamma(\mathcal{C},
\SheafHom^\bullet(\mathcal{F}^\bullet, \mathcal{G}^\bullet))
\longrightarrow
R\Hom_\mathcal{O}(\mathcal{F}^\bullet, \mathcal{G}^\bullet).
$$
En combinant ces
applications, nous obtenons l'application désirée du lemme.
\end{proof}

\begin{lemma}
\label{dga-lemma-tensor-with-complex-hom-adjoint}
Soit $(\mathcal{C}, \mathcal{O})$ un site annelé.
Soit $K^\bullet$ un complexe de $\mathcal{O}$-modules.
Alors le foncteur
$$
- \otimes_E^\mathbf{L} K^\bullet :
D(E, \text{d})
\longrightarrow
D(\mathcal{O})
$$
du lemme \ref{dga-lemma-tensor-with-complex-derived} est un adjoint à
gauche du foncteur
$$
R\Hom(K^\bullet, -) : D(\mathcal{O}) \longrightarrow D(E, \text{d})
$$
du lemme \ref{dga-lemma-existence-of-derived}.
\end{lemma}

\begin{proof}
L'énoncé signifie que nous avons
$$
\Hom_{D(E, \text{d})}(M, R\Hom(K^\bullet, L^\bullet)) =
\Hom_{D(\mathcal{O})}(M \otimes^\mathbf{L}_E K^\bullet, L^\bullet)
$$
bifonctoriellement en $M$ et $L^\bullet$. Pour voir cela, nous pouvons
remplacer $M$ par un $E$-module différentiel gradué $P$ ayant la
propriété (P). Nous pouvons également remplacer $L^\bullet$ par un
complexe K-injectif de $\mathcal{O}$-modules $I^\bullet$.
Le calcul des foncteurs dérivés donné dans les lemmes mentionnés dans l'énoncé,
combiné avec le lemme \ref{dga-lemma-hom-derived}, traduit ce qui précède en
$$
\Hom_{K(\text{Mod}_{(E, \text{d})})}
(P, \Hom_\mathcal{B}(K^\bullet, I^\bullet)) =
\Hom_{K(\mathcal{O})}(P \otimes_E K^\bullet, I^\bullet)
$$
où $\mathcal{B} = \text{Comp}^{dg}(\mathcal{O})$. Il existe
une application d'évaluation de droite à gauche fonctorielle en
$P$ et $I^\bullet$ (détails omis). Choisissons
une filtration $F_\bullet$ sur $P$ comme dans la définition de la propriété (P). D'après le
lemme \ref{dga-lemma-property-P-sequence} et le fait que les deux
membres de l'égalité sont des foncteurs homologiques en $P$ sur
$K(\text{Mod}_{(E, \text{d})})$, nous nous
ramenons au cas où $P$ est remplacé par le $E$-module
différentiel gradué $\bigoplus F_iP$. Puisque les deux membres
transforment les sommes directes dans la variable $P$ en
produits directs, nous nous ramenons au cas où $P$ est l'un des $E$-modules
différentiels gradués $F_iP$. Puisque chaque $F_iP$
possède une filtration finie (donnée par des monomorphismes admissibles)
dont les quotients successifs sont des $E$-modules projectifs gradués, nous
nous ramenons au cas où $P$ est un $E$-module projectif gradué. Dans ce cas,
nous avons clairement
$$
\Hom_{\text{Mod}^{dg}_{(E, \text{d})}}
(P, \Hom_\mathcal{B}(K^\bullet, I^\bullet)) =
\Hom_{\text{Comp}^{dg}(\mathcal{O})}(P \otimes_E K^\bullet, I^\bullet)
$$
comme $\mathbf{Z}$-modules gradués (car cette affirmation se réduit au cas $P =
E[k]$ où c'est évident). Comme l'isomorphisme est compatible avec les différentielles,
nous obtenons la conclusion.
\end{proof}

\begin{lemma}
\label{dga-lemma-fully-faithful-in-compact-case}
Soit $(\mathcal{C}, \mathcal{O})$ un site annelé. Soit
$K^\bullet$ un complexe de $\mathcal{O}$-modules. Faisons les hypothèses
suivantes :
\begin{enumerate}
\item $K^\bullet$ représente un objet compact de $D(\mathcal{O})$, et
\item $E = \Hom_{\text{Comp}^{dg}(\mathcal{O})}(K^\bullet, K^\bullet)$
calcule les groupes Ext de $K^\bullet$ dans $D(\mathcal{O})$.
\end{enumerate}
Alors le foncteur
$$
- \otimes_E^\mathbf{L} K^\bullet :
D(E, \text{d})
\longrightarrow
D(\mathcal{O})
$$
du lemme \ref{dga-lemma-tensor-with-complex-derived} est pleinement fidèle.
\end{lemma}

\begin{proof}
Comme notre foncteur a un adjoint à gauche donné par
$R\Hom(K^\bullet, -)$ selon le lemme \ref{dga-lemma-tensor-with-complex-hom-adjoint},
il suffit de montrer pour un $E$-module différentiel gradué $M$ que l'application
$$
H^0(M) \longrightarrow
\Hom_{D(\mathcal{O})}(K^\bullet, M \otimes_E^\mathbf{L} K^\bullet)
$$
est un isomorphisme. Nous pouvons supposer que $M = P$ est un $E$-module
différentiel gradué qui a la propriété (P). Puisque $K^\bullet$ définit un objet
compact, nous nous ramenons, à l'aide du
lemme \ref{dga-lemma-property-P-sequence}, au
cas où $P$ a une filtration finie dont les quotients successifs sont des sommes
directes de $E[k]$. Encore une fois, en utilisant la compacité, nous
nous ramenons au cas $P = E[k]$. L'hypothèse sur $K^\bullet$ est que le
résultat est vrai dans ces cas.
\end{proof}











\section{Caractérisation des objets compacts}
\label{dga-section-compact}

\noindent
Les objets compacts des catégories additives sont définis
dans Catégories dérivées, définition \ref{derived-definition-compact-object}. Dans
cette section, nous caractérisons les objets compacts de la catégorie
dérivée d'une algèbre différentielle graduée.

\begin{remark}
\label{dga-remark-source-graded-projective}
Soit $(A, \text{d})$ une algèbre différentielle graduée. Existe-t-il
une caractérisation des $A$-modules différentiels gradués $P$ pour
lesquels nous avons
$$
\Hom_{K(A, \text{d})}(P, M) = \Hom_{D(A, \text{d})}(P, M)
$$
pour tous les $A$-modules différentiels gradués $M$ ?
Soit $\mathcal{D} \subset K(A, \text{d})$ la sous-catégorie pleine
dont les objets sont les objets $P$ satisfaisant ce qui précède. Alors
$\mathcal{D}$ est une sous-catégorie triangulée strictement pleine et saturée de $K(A,
\text{d})$. Si $P$ est projectif en tant que $A$-module gradué, alors pour voir si
$P$ est un objet de $\mathcal{D}$, il suffit de vérifier que
$\Hom_{K(A, \text{d})}(P, M) = 0$ chaque fois que $M$ est
acyclique. Cependant, en général, il ne suffit pas de supposer que $P$ est projectif en
tant que $A$-module gradué. Exemple : prenons $A = R = k[\epsilon]$ où $k$ est
un corps et $k[\epsilon] = k[x]/(x^2)$ est l'anneau des nombres duaux. Soit
$P$ l'objet avec $P^n = R$ pour tous les $n \in \mathbf{Z}$ et de
différentielle donnée par multiplication par $\epsilon$. Alors
$\text{id}_P \in \Hom_{K(A, \text{d})}(P, P)$ est un élément non nul mais $P$ est
acyclique.
\end{remark}

\begin{remark}
\label{dga-remark-graded-projective-is-compact}
Soit $(A, \text{d})$ une algèbre différentielle graduée. Disons qu'un
$A$-module différentiel gradué $M$ est {\it de type fini} si $M$ est engendré, en tant que
module à droite sur $A$, par un nombre fini d'éléments. Si $P$ est un
$A$-module différentiel gradué qui est projectif gradué de type fini, alors nous
pouvons nous demander : Est-ce que $P$ donne un objet compact de $D(A, \text{d})$ ? Nous
nous attendons à ce que ce ne soit pas vrai en général, mais nous ne
connaissons pas de contre-exemple. Cependant, si $P$ est également un objet de la
catégorie $\mathcal{D}$ de la remarque \ref{dga-remark-source-graded-projective},
alors c'est le cas (cela découle du fait que les sommes directes dans $D(A,
\text{d})$ sont données par des sommes directes de modules ; les détails sont omis).
\end{remark}

\begin{lemma}
\label{dga-lemma-factor-through-nicer}
Soit $(A, \text{d})$ une algèbre différentielle graduée. Soit $E$ un objet
compact de $D(A, \text{d})$. Soit $P$ un $A$-module différentiel gradué qui
possède une filtration finie
$$
0 = F_{-1}P \subset F_0P \subset F_1P \subset \ldots \subset F_nP = P
$$
par des sous-modules différentiels gradués, telle que
$$
F_{i + 1}P/F_iP \cong \bigoplus\nolimits_{j \in J_i} A[k_{i, j}]
$$
comme $A$-modules différentiels gradués pour certains ensembles $J_i$ et entiers
$k_{i, j}$. Soit $E \to P$ un morphisme de $D(A,
\text{d})$. Alors il existe un sous-module différentiel gradué $P' \subset P$ tel que
$F_{i + 1}P \cap P'/(F_iP \cap P')$ soit égal à
$\bigoplus_{j \in J'_i} A[k_{i, j}]$ pour certains sous-ensembles
finis $J'_i \subset J_i$ et tel que $E \to P$ se factorise par $P'$.
\end{lemma}

\begin{proof}
Nous allons démontrer par récurrence sur $-1 \leq m \leq n$ qu'il
existe un sous-module différentiel gradué $P' \subset P$ tel que
\begin{enumerate}
\item $F_mP \subset P'$,
\item pour $i \geq m$ le quotient $F_{i + 1}P \cap P'/(F_iP \cap P')$ est
isomorphe à $\bigoplus_{j \in J'_i} A[k_{i, j}]$ pour certains sous-ensembles
finis $J'_i \subset J_i$, et
\item $E \to P$ se factorise par $P'$.
\end{enumerate}
Le cas de base est $m = n$ où nous pouvons prendre $P' = P$.

\medskip\noindent
Étape de récurrence. Supposons que $P'$
convienne pour $m$. Pour $i \geq m$ et $j \in J'_i$, soit $x_{i, j} \in F_{i +
1}P \cap P'$ un élément homogène de degré $k_{i, j}$ dont l'image
dans $F_{i + 1}P \cap P'/(F_iP \cap P')$ est le générateur du
facteur correspondant à $j \in J_i$. Les $x_{i, j}$
engendrent $P'/F_mP$ en tant que $A$-module. Écrivons
$$
\text{d}(x_{i, j}) = \sum x_{i', j'} a_{i, j}^{i', j'} + y_{i, j}
$$
avec $y_{i, j} \in F_mP$ et $a_{i, j}^{i', j'} \in A$. Il
existe un sous-ensemble fini
$J'_{m - 1} \subset J_{m - 1}$ tel que l'image de chaque $y_{i, j}$
appartienne au sous-module $\bigoplus_{j \in J'_{m - 1}} A[k_{m - 1, j}]$ de
$F_mP/F_{m - 1}P$. Soit $P'' \subset F_mP$ l'image inverse de
$\bigoplus_{j \in J'_{m - 1}} A[k_{m - 1, j}]$ sous l'application
$F_mP \to F_mP/F_{m - 1}P$. Alors nous voyons que le $A$-sous-module
$$
P'' + \sum x_{i, j}A
$$
est un sous-module différentiel gradué du type que nous recherchons. De plus
$$
P'/(P'' + \sum x_{i, j}A) =
\bigoplus\nolimits_{j \in J_{m - 1} \setminus J'_{m - 1}} A[k_{m - 1, j}]
$$
Puisque $E$ est compact, la composition de l'application donnée $E \to P'$ avec
l'application quotient se factorise par la somme directe d'un nombre fini des
facteurs du module ci-dessus. Par conséquent, après avoir agrandi $J'_{m - 1}$,
nous pouvons supposer que $E \to P'$ se factorise
par $P'' + \sum x_{i, j}A$ comme souhaité.
\end{proof}

\noindent
Il n'est pas vrai que tout objet compact de $D(A, \text{d})$ provienne
d'un $A$-module différentiel gradué qui soit projectif gradué de type
fini, voir Exemples, section \ref{examples-section-interesting-compact}.

\begin{proposition}
\label{dga-proposition-compact}
Soit $(A, \text{d})$ une algèbre différentielle graduée. Soit $E$ un objet
de $D(A, \text{d})$. Alors les assertions suivantes sont équivalentes
\begin{enumerate}
\item $E$ est un objet compact,
\item $E$ est un facteur direct d'un objet de $D(A, \text{d})$ qui est
représenté par un module différentiel gradué $P$ qui possède une
filtration finie $F_\bullet$ par des sous-modules différentiels gradués, telle que
les $F_iP/F_{i - 1}P$ soient des sommes directes finies de décalages de $A$.
\end{enumerate}
\end{proposition}

\begin{proof}
Supposons que $E$ soit compact. D'après le lemme \ref{dga-lemma-resolve}, nous pouvons
supposer que $E$ est représenté par un $A$-module différentiel gradué $P$ avec la
propriété (P). Considérons le triangle distingué
$$
\bigoplus F_iP \to \bigoplus F_iP \to P
\xrightarrow{\delta} \bigoplus F_iP[1]
$$
provenant de la suite exacte courte admissible du lemme
\ref{dga-lemma-property-P-sequence}. Puisque $E$ est compact, nous avons
$\delta = \sum_{i = 1, \ldots, n} \delta_i$ pour certains
$\delta_i : P \to F_iP[1]$. Puisque la composition de $\delta$ avec
l'application $\bigoplus F_iP[1] \to \bigoplus F_iP[1]$ est nulle (Catégories
dérivées, lemme \ref{derived-lemma-composition-zero}), il s'ensuit que
$\delta = 0$ (car $\bigoplus F_iP \to \bigoplus F_iP$ envoie le facteur direct $F_iP$
via la différence de $\text{id}$ et de l'application d'inclusion dans $F_{i - 1}P$).
Ainsi, nous voyons que
l'identité sur $E$ se factorise à travers $\bigoplus F_iP$ dans
$D(A, \text{d})$ (d'après Catégories
dérivées, lemme \ref{derived-lemma-split}). Ensuite, nous utilisons
à nouveau la compacité de $P$ pour voir que l'application $E \to
\bigoplus F_iP$ se factorise par $\bigoplus_{i = 1, \ldots, n} F_iP$ pour un
certain $n$. En d'autres termes, l'identité sur $E$ se factorise par
$\bigoplus_{i = 1, \ldots, n} F_iP$. Par le lemme
\ref{dga-lemma-factor-through-nicer}, nous voyons
que l'identité de $E$ se factorise comme $E \to P \to E$ où $P$ est
comme dans la partie (2) de l'énoncé du lemme. En d'autres termes,
nous avons prouvé que (1) implique (2).

\medskip\noindent
Supposons (2).
D'après Catégories dérivées, lemme
\ref{derived-lemma-compact-objects-subcategory}, il suffit de montrer que $P$ donne un objet compact. Remarquons que
$P$ a la propriété (P), donc nous avons
$$
\Hom_{D(A, \text{d})}(P, M) = \Hom_{K(A, \text{d})}(P, M)
$$
pour tout module différentiel gradué $M$ d'après le lemme \ref{dga-lemma-hom-derived}.
Comme les sommes directes dans $D(A, \text{d})$ sont données par des sommes
directes de modules gradués (lemme \ref{dga-lemma-derived-products}), nous
nous ramenons à montrer que $\Hom_{K(A, \text{d})}(P, M)$ commute avec les sommes
directes. En utilisant le fait que $K(A, \text{d})$ est une catégorie
triangulée, que $\Hom$ est un foncteur cohomologique dans la
première variable, et la filtration sur $P$, nous nous ramenons au cas où $P$
est une somme directe finie de décalages de $A$. Ainsi, nous nous ramenons
au cas $P = A[k]$, ce qui est clair.
\end{proof}

\begin{lemma}
\label{dga-lemma-compact-implies-bounded}
Soit $(A, \text{d})$ une algèbre différentielle
graduée. Pour tout objet compact $E$ de $D(A, \text{d})$, il
existe des entiers $a \leq b$ tels que $\Hom_{D(A, \text{d})}(E, M) = 0$
si $H^i(M) = 0$ pour $i \in [a, b]$.
\end{lemma}

\begin{proof}
Observons que la collection d'objets de $D(A, \text{d})$ pour lesquels un
tel couple d'entiers existe est une sous-catégorie triangulée saturée et strictement
pleine de $D(A, \text{d})$. Ainsi,
d'après la proposition \ref{dga-proposition-compact}, il suffit de prouver cela
lorsque $E$ est représenté par un module différentiel gradué $P$ qui possède
une filtration finie $F_\bullet$ par des sous-modules différentiels gradués tels
que $F_iP/F_{i - 1}P$ soient des sommes directes finies de décalages de $A$.
En utilisant la compatibilité avec les triangles, nous voyons qu'il
suffit de le prouver pour $P = A$. Dans ce cas, $\Hom_{D(A, \text{d})}(A, M) = H^0(M)$ et
le résultat vaut pour $a = b = 0$.
\end{proof}

\noindent
Si $(A, \text{d})$ est simplement une algèbre placée en degré $0$
avec une différentielle nulle ou, plus généralement,
est concentrée en un nombre fini de degrés, alors nous obtenons
la description plus précise des objets compacts.

\begin{lemma}
\label{dga-lemma-compact}
Soit $(A, \text{d})$ une algèbre différentielle graduée. Supposons que $A^n = 0$
pour $|n| \gg 0$. Soit $E$ un objet de $D(A, \text{d})$. Les
assertions suivantes sont équivalentes
\begin{enumerate}
\item $E$ est un objet compact, et
\item $E$ peut être représenté par un $A$-module différentiel gradué $P$
qui est projectif de type fini en tant que $A$-module gradué et
satisfait $\Hom_{K(A, \text{d})}(P, M) = \Hom_{D(A, \text{d})}(P,
M)$ pour tout $A$-module différentiel gradué $M$.
\end{enumerate}
\end{lemma}

\begin{proof}
Soit $\mathcal{D} \subset K(A, \text{d})$ la sous-catégorie triangulée
discutée dans la remarque \ref{dga-remark-source-graded-projective}.
Soit $P$ un objet de $\mathcal{D}$ qui est projectif de type fini comme
$A$-module gradué. Alors $P$ représente un objet compact de $D(A,
\text{d})$ d'après la remarque \ref{dga-remark-graded-projective-is-compact}.

\medskip\noindent
Pour prouver la réciproque, soit $E$ un objet compact de $D(A, \text{d})$.
Fixons $a \leq b$ comme dans le lemme
\ref{dga-lemma-compact-implies-bounded}. Après avoir diminué $a$ et augmenté $b$ si nécessaire, nous pouvons
également supposer que $H^i(E) = 0$ pour $i \not \in [a, b]$ (cela
découle de la proposition \ref{dga-proposition-compact} et de notre hypothèse sur
$A$). De plus, fixons un entier $c > 0$ tel que $A^n = 0$ si $|n| \geq c$.

\medskip\noindent
Par la proposition \ref{dga-proposition-compact} nous voyons que $E$ est un
facteur direct, dans $D(A, \text{d})$, d'un $A$-module différentiel gradué $P$ qui
possède une filtration finie $F_\bullet$ par des sous-modules
différentiels gradués telle que les $F_iP/F_{i - 1}P$ soient des sommes directes
finies de décalages de $A$. En particulier, $P$ a la propriété (P) et nous
avons $\Hom_{D(A, \text{d})}(P, M) = \Hom_{K(A, \text{d})}(P, M)$ pour tout
module différentiel gradué $M$ selon le lemme \ref{dga-lemma-hom-derived}. En
d'autres termes, $P$ est un objet de la sous-catégorie
triangulée $\mathcal{D} \subset K(A, \text{d})$ discutée dans la remarque
\ref{dga-remark-source-graded-projective}. Notons que
$P$ est libre de type fini comme $A$-module gradué.

\medskip\noindent
Choisissons $n > 0$ tel que $b + 4c - n <
a$. Représentons le projecteur sur $E$ par un endomorphisme $\varphi : P \to P$
de $A$-modules différentiels gradués. Considérons le triangle distingué
$$
P \xrightarrow{1 - \varphi} P \to C \to P[1]
$$
dans $K(A, \text{d})$ où $C$ est le cône de la première flèche. Ensuite,
$C$ est un objet de $\mathcal{D}$, nous
avons $C \cong E \oplus E[1]$ dans $D(A, \text{d})$, et $C$
est un $A$-module gradué libre de type fini.
Ensuite, considérons un triangle distingué
$$
C[1] \to C \to C' \to C[2]
$$
dans $K(A, \text{d})$ où $C'$ est le cône sur un morphisme $C[1] \to
C$ représentant la composition
$$
C[1] \cong E[1] \oplus E[2] \to E[1] \to E \oplus E[1] \cong C
$$
dans $D(A, \text{d})$. Ensuite, nous voyons que $C'$ représente $E \oplus E[2]$.
En continuant de cette manière, nous voyons que nous pouvons trouver un
$A$-module différentiel gradué $P$ qui est un objet de $\mathcal{D}$,
est libre de type fini en tant que $A$-module gradué, et représente $E \oplus E[n]$.

\medskip\noindent
Choisissons une base $x_i$, $i \in I$ d'éléments homogènes pour $P$ en
tant que $A$-module. Posons $d_i =
\deg(x_i)$. Soit $P_1$ le $A$-sous-module de $P$ engendré par $x_i$
et $\text{d}(x_i)$ pour $d_i \leq a - c - 1$.
Soit $P_2$ le $A$-sous-module de $P$ engendré par $x_i$ et
$\text{d}(x_i)$ pour $d_i \geq b - n + c$. Nous
observons
\begin{enumerate}
\item $P_1$ et $P_2$ sont des sous-modules différentiels gradués de $P$,
\item $P_1^t = 0$ pour $t \geq a$,
\item $P_1^t = P^t$ pour $t \leq a - 2c$,
\item $P_2^t = 0$ pour $t \leq b - n$,
\item $P_2^t = P^t$ pour $t \geq b - n + 2c$.
\end{enumerate}
Comme $b - n + 2c \geq a - 2c$ par notre choix de $n$
nous obtenons une suite exacte courte de $A$-modules différentiels gradués
$$
0 \to P_1 \cap P_2 \to P_1 \oplus P_2 \xrightarrow{\pi} P \to 0
$$
Puisque $P$ est projectif en tant que $A$-module gradué, il s'agit d'une suite
exacte courte admissible (lemme \ref{dga-lemma-target-graded-projective}).
Ainsi, nous obtenons un morphisme de
liaison $\delta : P \to (P_1 \cap P_2)[1]$ dans $K(A, \text{d})$, voir
le lemme \ref{dga-lemma-admissible-ses}.
Puisque $P = E \oplus E[n]$ et puisque $P_1 \cap P_2$ est concentré en
degrés $(b - n, a)$, nous constatons que
$\Hom_{D(A, \text{d})}(E \oplus E[n], (P_1 \cap P_2)[1])$ est nul. Par
conséquent, $\delta = 0$ en tant que morphisme dans $K(A, \text{d})$
puisque $P$ est un objet de $\mathcal{D}$. Par
Catégories dérivées, le lemme \ref{derived-lemma-split} nous
permet de trouver une application $s : P \to P_1 \oplus P_2$
telle que $\pi \circ s = \text{id}_P + \text{d}h + h\text{d}$ pour un certain $h : P \to
P$ de degré $-1$. Puisque $P_1 \oplus P_2 \to P$ est surjectif et que $P$ est
projectif comme $A$-module gradué, nous pouvons choisir un relèvement homogène
$\tilde h : P \to P_1 \oplus P_2$ de $h$. Puis nous changeons $s$ en $s +
\text{d} \tilde h + \tilde h \text{d}$ pour obtenir $\pi \circ s =
\text{id}_P$. Cela signifie que nous obtenons une décomposition en
somme directe $P = s^{-1}(P_1) \oplus s^{-1}(P_2)$. Puisque
$s^{-1}(P_2)$ est égal à $P$ en degrés $\geq b - n + 2c$, nous voyons que
$s^{-1}(P_2) \to P \to E$ est un quasi-isomorphisme, c'est-à-dire un
isomorphisme dans $D(A, \text{d})$. Cela termine la démonstration.
\end{proof}








\section{Équivalences de catégories dérivées}
\label{dga-section-equivalence}

\noindent
Soit $R$ un anneau. Soient $(A, \text{d})$ et $(B, \text{d})$ des $R$-algèbres
différentielles graduées. Une question naturelle est de savoir ce que
signifie l'équivalence de $D(A, \text{d})$ avec $D(B, \text{d})$ en tant que
catégorie triangulée linéaire sur $R$. C'est une question plutôt subtile et il
s'avérera qu'elle n'est pas toujours la bonne question à poser.
Néanmoins, dans cette section nous recueillons certaines conditions qui
garantissent que c'est le cas.

\medskip\noindent
Nous recommandons vivement au lecteur de consulter l'article
fondateur \cite{Rickard} sur ce sujet.

\begin{lemma}
\label{dga-lemma-qis-equivalence}
Soit $R$ un anneau. Soit $(A, \text{d}) \to (B, \text{d})$ un
homomorphisme d'algèbres différentielles graduées sur $R$, qui induit un
isomorphisme sur les algèbres de cohomologie. Alors
$$
- \otimes_A^\mathbf{L} B : D(A, \text{d}) \to D(B, \text{d})
$$
donne une équivalence $R$-linéaire de catégories triangulées avec
quasi-inverse le foncteur de restriction $N \mapsto N_A$.
\end{lemma}

\begin{proof}
D'après le lemme
\ref{dga-lemma-tensor-with-compact-fully-faithful}, le foncteur $M \longmapsto M \otimes_A^\mathbf{L} B$ est
pleinement fidèle. D'après le lemme
\ref{dga-lemma-tensor-hom-adjoint}, le foncteur $N \longmapsto R\Hom(B, N) = N_A$ est un adjoint à droite,
voir l'exemple
\ref{dga-example-map-hom-tensor}. Il est clair que le noyau de $R\Hom(B, -)$ est nul. Par
conséquent, le résultat découle de
Catégories dérivées, lemme
\ref{derived-lemma-fully-faithful-adjoint-kernel-zero}.
\end{proof}

\noindent
Lorsque nous analysons la démonstration ci-dessus, nous voyons que nous obtenons sans
travail supplémentaire la généralisation suivante.

\begin{lemma}
\label{dga-lemma-tilting-equivalence}
Soit $R$ un anneau. Soient $(A, \text{d})$ et $(B, \text{d})$ des algèbres
différentielles graduées sur $R$. Soit $N$ un bimodule
différentiel gradué sur $(A, B)$. Faisons les hypothèses suivantes :
\begin{enumerate}
\item $N$ définit un objet compact de $D(B, \text{d})$,
\item si $N' \in D(B, \text{d})$ et
$\Hom_{D(B, \text{d})}(N, N'[n]) = 0$ pour $n \in \mathbf{Z}$,
alors $N' = 0$, et
\item l'application $H^k(A) \to \Hom_{D(B, \text{d})}(N, N[k])$ est
un isomorphisme pour tout $k \in \mathbf{Z}$.
\end{enumerate}
Alors
$$
- \otimes_A^\mathbf{L} N : D(A, \text{d}) \to D(B, \text{d})
$$
donne une équivalence $R$-linéaire de catégories triangulées.
\end{lemma}

\begin{proof}
D'après le lemme
\ref{dga-lemma-tensor-with-compact-fully-faithful} le foncteur $M \longmapsto M \otimes_A^\mathbf{L} N$ est
pleinement fidèle. D'après le lemme
\ref{dga-lemma-tensor-hom-adjoint} le foncteur $N' \longmapsto R\Hom(N, N')$ est un adjoint à droite.
D'après l'hypothèse (3) le noyau de $R\Hom(N, -)$ est nul. Par
conséquent, le résultat découle
de Catégories dérivées, lemme
\ref{derived-lemma-fully-faithful-adjoint-kernel-zero}.
\end{proof}

\begin{remark}
\label{dga-remark-tilting-equivalence}
Dans le lemme \ref{dga-lemma-tilting-equivalence}, nous pouvons
remplacer la condition (2) par la condition selon laquelle $N$ est un
générateur classique pour $D_{compact}(B,
d)$, voir Catégories dérivées,
proposition \ref{derived-proposition-generator-versus-classical-generator}. De plus,
si nous savions que $R\Hom(N, B)$ est un objet compact de $D(A,
\text{d})$, alors il suffirait de vérifier que $N$ est un générateur
faible pour $D_{compact}(B, \text{d})$. Nous omettons la
démonstration ; nous l'ajouterons ici si nous en avons besoin un
jour dans le projet Champs.
\end{remark}

\noindent
Le $B$-module $P$ du lemme ci-dessous est parfois appelé
« complexe $(A, B)$-basculant ».

\begin{lemma}
\label{dga-lemma-rickard}
Soit $R$ un anneau. Soient $(A, \text{d})$ et $(B, \text{d})$ des
$R$-algèbres différentielles graduées. Supposons que $A = H^0(A)$. Les
assertions suivantes sont équivalentes
\begin{enumerate}
\item $D(A, \text{d})$ et $D(B, \text{d})$ sont équivalentes en tant que
catégories triangulées $R$-linéaires, et
\item il existe un objet $P$ de $D(B, \text{d})$ tel que
\begin{enumerate}
\item $P$ est un objet compact de $D(B, \text{d})$,
\item si $N \in D(B, \text{d})$ avec $\Hom_{D(B, \text{d})}(P, N[i]) = 0$
pour $i \in \mathbf{Z}$, alors $N = 0$,
\item $\Hom_{D(B, \text{d})}(P, P[i]) = 0$ pour $i \not = 0$ et
est égal à $A$ pour $i = 0$.
\end{enumerate}
\end{enumerate}
L'équivalence $D(A, \text{d}) \to D(B, \text{d})$
construite dans (2) envoie $A$ sur $P$.
\end{lemma}

\begin{proof}
Soit $F : D(A, \text{d}) \to D(B, \text{d})$ une équivalence. Alors $F$
envoie les objets compacts vers des objets compacts. Par conséquent, $P = F(A)$ est
compact, c'est-à-dire que (2)(a) est vérifié. Les conditions (2)(b) et (2)(c) découlent
immédiatement du fait que $F$ est une équivalence.

\medskip\noindent
Soit $P$ un objet comme en (2). Représentons $P$ par un
module différentiel gradué avec la propriété (P). Posons
$$
(E, \text{d}) = \Hom_{\text{Mod}^{dg}_{(B, \text{d})}}(P, P)
$$
Alors $H^0(E) = A$ et $H^k(E) = 0$ pour $k \not = 0$ par le
lemme \ref{dga-lemma-hom-derived} et l'hypothèse (2)(c). En
considérant $P$ comme un $(E, B)$-bimodule et en utilisant
le lemme \ref{dga-lemma-tilting-equivalence} et l'hypothèse (2)(b), nous
obtenons une équivalence
$$
D(E, \text{d}) \to D(B, \text{d})
$$
envoyant $E$ sur $P$.
Soit $E' \subset E$ la sous-algèbre différentielle graduée sur $R$
avec
$$
(E')^i = \left\{
\begin{matrix}
E^i & \text{si }i < 0 \\
\Ker(E^0 \to E^1) & \text{si }i = 0 \\
0 & \text{si }i > 0
\end{matrix}
\right.
$$
Il y a alors des quasi-isomorphismes d'algèbres
différentielles graduées $(A, \text{d}) \leftarrow (E', \text{d}) \rightarrow (E, \text{d})$. Ainsi,
nous obtenons des équivalences
$$
D(A, \text{d}) \leftarrow D(E', \text{d}) \rightarrow D(E, \text{d})
\rightarrow D(B, \text{d})
$$
par le lemme \ref{dga-lemma-qis-equivalence}.
\end{proof}

\begin{remark}
\label{dga-remark-lift-equivalence-to-dga}
Soit $R$ un anneau. Soient $(A, \text{d})$ et $(B, \text{d})$ des $R$-algèbres
différentielles graduées. Supposons donnée une équivalence $R$-linéaire
$$
F : D(A, \text{d}) \longrightarrow D(B, \text{d})
$$
de catégories triangulées. Posons $N = F(A)$. Alors $N$ est un $B$-module
différentiel gradué. Puisque $F$ est une équivalence et que $A$ est un objet
compact de $D(A, \text{d})$, nous concluons que $N$ est un objet compact de
$D(B, \text{d})$. Puisque $A$ engendre $D(A, \text{d})$ et que $F$ est
une équivalence, nous voyons que $N$ engendre $D(B, \text{d})$. Enfin,
$H^k(A) = \Hom_{D(A, \text{d})}(A, A[k])$ et comme $F$ est une équivalence, nous voyons
que $F$ induit un isomorphisme $H^k(A) =
\Hom_{D(B, \text{d})}(N, N[k])$ pour tous les $k$. Afin de
conclure qu'il existe une équivalence $D(A, \text{d})
\longrightarrow D(B, \text{d})$ qui découle de la construction dans
le lemme
\ref{dga-lemma-tilting-equivalence} tout ce dont nous avons besoin est
une structure de $A$-module à gauche sur $N$ compatible
avec la différentielle et commutant avec la
structure de $B$-module à droite donnée. En fait, il suffit de
le faire après avoir remplacé $N$ par un $B$-module différentiel
gradué quasi-isomorphe. La structure du
module peut être construite dans certains cas. Par exemple, si
nous supposons que $F$ peut être relevé à un foncteur
différentiel gradué
$$
F^{dg} :
\text{Mod}^{dg}_{(A, \text{d})}
\longrightarrow
\text{Mod}^{dg}_{(B, \text{d})}
$$
(pour la notation voir exemple \ref{dga-example-dgm-dg-cat})
entre les catégories différentielles graduées associées,
alors cela est vrai. Un autre cas est discuté dans la proposition ci-dessous.
\end{remark}

\begin{proposition}
\label{dga-proposition-rickard}
Soit $R$ un anneau. Soient $(A, \text{d})$ et $(B, \text{d})$ des
$R$-algèbres différentielles graduées. Soit $F : D(A, \text{d}) \to D(B, \text{d})$ une
équivalence $R$-linéaire de catégories triangulées. Faisons les hypothèses suivantes :
\begin{enumerate}
\item $A = H^0(A)$, et
\item $B$ est K-plat en tant que complexe de $R$-modules.
\end{enumerate}
Alors il existe un $(A, B)$-bimodule $N$ comme dans
le lemme \ref{dga-lemma-tilting-equivalence}.
\end{proposition}

\begin{proof}
Comme dans la remarque \ref{dga-remark-lift-equivalence-to-dga} ci-dessus, nous posons $N =
F(A)$ dans $D(B, \text{d})$. Nous pouvons supposer que $N$ est un $B$-module
différentiel gradué avec la propriété (P). Posons
$$
(E, \text{d}) = \Hom_{\text{Mod}^{dg}_{(B, \text{d})}}(N, N)
$$
Puis $H^0(E) = A$ et $H^k(E) = 0$ pour $k \not = 0$ par le
lemme \ref{dga-lemma-hom-derived}. De
plus, d'après la discussion dans la remarque \ref{dga-remark-lift-equivalence-to-dga} et
par le lemme \ref{dga-lemma-tilting-equivalence},
nous voyons que $N$ en tant que $(E, B)$-bimodule induit
une équivalence $- \otimes_E^\mathbf{L} N : D(E, \text{d}) \to D(B, \text{d})$. Soit
$E' \subset E$ la sous-algèbre différentielle graduée sur $R$
avec
$$
(E')^i = \left\{
\begin{matrix}
E^i & \text{si }i < 0 \\
\Ker(E^0 \to E^1) & \text{si }i = 0 \\
0 & \text{si }i > 0
\end{matrix}
\right.
$$
Ensuite, il y a des quasi-isomorphismes d'algèbres
différentielles graduées $(A, \text{d}) \leftarrow (E', \text{d}) \rightarrow (E, \text{d})$. Ainsi,
nous obtenons des équivalences
$$
D(A, \text{d}) \leftarrow D(E', \text{d}) \rightarrow D(E, \text{d})
\rightarrow D(B, \text{d})
$$
d'après le lemme
\ref{dga-lemma-qis-equivalence}. Notons que le quasi-inverse $D(A, \text{d}) \to D(E', \text{d})$
de la flèche verticale gauche est
donné par $M \mapsto M \otimes_A^\mathbf{L} A$ où $A$ est considéré comme
un bimodule $(A, E')$, voir l'exemple \ref{dga-example-map-hom-tensor}.
D'autre part, le foncteur $D(E', \text{d}) \to D(B, \text{d})$ est donné par $M
\mapsto M \otimes_{E'}^\mathbf{L} N$ où $N$ est comme ci-dessus. Nous
concluons par le lemme \ref{dga-lemma-compose-tensor-functors}.
\end{proof}

\begin{remark}
\label{dga-remark-rickard}
Soient $A, B, F, N$ comme dans la proposition \ref{dga-proposition-rickard}.
Il n'est pas clair que $F$ et le foncteur
$G(-) = - \otimes_A^\mathbf{L} N$ soient isomorphes. Par
construction, il existe un isomorphisme $N =
G(A) \to F(A)$ dans $D(B, \text{d})$. Il est
simple d'étendre cela à un isomorphisme fonctoriel $G(M) \to F(M)$ lorsque
$M$ est un $A$-module différentiel gradué qui est projectif gradué (par
exemple, une somme de décalages de $A$). On peut alors
conclure que $G(M) \cong F(M)$ lorsque $M$ est un cône d'une application entre
de tels modules. Nous ne savons pas si l'on peut en dire davantage en
général.
\end{remark}

\begin{lemma}
\label{dga-lemma-rickard-rings}
Soit $R$ un anneau.
Soient $A$ et $B$ des $R$-algèbres. Les assertions suivantes sont équivalentes
\begin{enumerate}
\item il existe une équivalence $R$-linéaire $D(A) \to
D(B)$ de catégories triangulées,
\item il existe un objet $P$ de $D(B)$ tel que
\begin{enumerate}
\item $P$ peut être représenté par un complexe borné
de $B$-modules projectifs de type fini,
\item si $K \in D(B)$ avec $\Ext^i_B(P, K) = 0$ pour
$i \in \mathbf{Z}$, alors $K = 0$, et
\item $\Ext^i_B(P, P) = 0$ pour $i \not = 0$ et
est égal à $A$ pour $i= 0$.
\end{enumerate}
\end{enumerate}
De plus, si $B$ est plat en tant que $R$-module, alors cela
équivaut également à
\begin{enumerate}
\item[(3)] il existe un $(A, B)$-bimodule $N$ tel que $-
\otimes_A^\mathbf{L} N : D(A) \to D(B)$ soit une équivalence.
\end{enumerate}
\end{lemma}

\begin{proof}
L'équivalence de (1) et (2) est un cas particulier
du lemme \ref{dga-lemma-rickard} combiné avec le résultat du
lemme \ref{dga-lemma-compact} caractérisant les objets compacts de
$D(B)$ (petit détail omis).
L'équivalence avec (3) si $B$ est $R$-plat découle de la
proposition \ref{dga-proposition-rickard}.
\end{proof}

\begin{remark}
\label{dga-remark-centers}
Soit $R$ un anneau. Soient $A$ et $B$ des $R$-algèbres. Si
$D(A)$ et $D(B)$ sont équivalentes comme catégories triangulées
$R$-linéaires, alors les centres de $A$ et $B$ sont isomorphes comme
$R$-algèbres. En particulier, si $A$ et $B$ sont commutatifs, alors $A \cong B$.
La démonstration plutôt délicate peut être trouvée dans
\cite[Proposition 9.2]{Rickard} ou \cite[Proposition 6.3.2]{KZ}. Une autre approche
pourrait être d'utiliser la cohomologie de Hochschild (voir remarque
ci-dessous).
\end{remark}

\begin{remark}
\label{dga-remark-hochschild-cohomology}
Soit $R$ un anneau. Soient $(A, \text{d})$ et $(B, \text{d})$ des $R$-algèbres
différentielles graduées dont les catégories dérivées sont équivalentes, c'est-à-dire
telles qu'il existe une équivalence $R$-linéaire $D(A, \text{d}) \to D(B,
\text{d})$ de catégories triangulées. Nous souhaitons montrer que certains invariants
de $(A, \text{d})$ et $(B, \text{d})$ coïncident. Dans de nombreuses
situations, on a un meilleur contrôle de la situation. Par exemple, il peut arriver
qu'il existe une équivalence de la forme
$$
- \otimes_A \Omega : D(A, \text{d}) \longrightarrow D(B, \text{d})
$$
pour un certain $(A, B)$-bimodule différentiel gradué
$\Omega$ (cela se produit dans la situation de la
proposition \ref{dga-proposition-rickard} et est souvent vrai si
l'équivalence provient d'une construction géométrique). Supposons de
plus que le quasi-inverse de notre foncteur soit donné par
$$
- \otimes_B^\mathbf{L} \Omega' : D(B, \text{d}) \longrightarrow D(A, \text{d})
$$
pour un $(B, A)$-bimodule différentiel gradué $\Omega'$ (et
comme précédemment un tel module $\Omega'$ existe souvent en pratique).
Dans ce cas, nous pouvons considérer le foncteur
$$
D(A^{opp} \otimes_R A, \text{d})
\longrightarrow
D(B^{opp} \otimes_R B, \text{d}),\quad
M \longmapsto \Omega' \otimes^\mathbf{L}_A M \otimes_A^\mathbf{L} \Omega
$$
sur les catégories dérivées de bimodules
(utilisez le lemme \ref{dga-lemma-bimodule-over-tensor} pour transformer les
bimodules en
modules à droite). Observons que ce foncteur envoie le $(A, A)$-bimodule
$A$ sur le $(B, B)$-bimodule $B$. Sous des conditions adéquates
(par exemple, la platitude de $A$, $B$, $\Omega$ sur $R$,
etc.) ce foncteur sera également une équivalence. Si
tel est le cas, il s'ensuit que nous avons des isomorphismes de groupes
de cohomologie de Hochschild
$$
HH^i(A, \text{d}) =
\Hom_{D(A^{opp} \otimes_R A, \text{d})}(A, A[i])
\longrightarrow
\Hom_{D(B^{opp} \otimes_R B, \text{d})}(B, B[i]) =
HH^i(B, \text{d}).
$$
Par exemple, si $A = H^0(A)$, alors $HH^0(A, \text{d})$ est
égal au centre de $A$, et cela donne une démonstration conceptuelle du
résultat mentionné dans la remarque \ref{dga-remark-centers}. Si jamais nous
avons besoin de cette remarque, nous fournirons ici un énoncé précis avec
une démonstration détaillée.
\end{remark}



\section{Résolutions d'algèbres différentielles graduées}
\label{dga-section-resolution-dgas}

\noindent
Soit $R$ un anneau. Sous nos hypothèses, la $R$-algèbre libre
$R\langle S \rangle$ sur un ensemble $S$ est l'algèbre ayant pour $R$-base
les expressions
$$
s_1 s_2 \ldots s_n
$$
où $n \geq 0$ et où $s_1, \ldots, s_n \in S$ forment une suite
d'éléments de $S$. La multiplication est donnée par concaténation
$$
(s_1 s_2 \ldots s_n) \cdot (s'_1 s'_2 \ldots s'_m) =
s_1 \ldots s_n s'_1 \ldots s'_m
$$
Cette algèbre est caractérisée par le fait que l'application
$$
\Mor_{R\text{-alg}}(R\langle S \rangle, A) \to
\text{Map}(S, A),\quad
\varphi \longmapsto (s \mapsto \varphi(s))
$$
est une bijection pour chaque $R$-algèbre $A$.

\medskip\noindent
Dans la catégorie des $R$-algèbres graduées, notre ensemble $S$ devrait être
muni d'une graduation, que nous considérons comme une application $\deg : S \to
\mathbf{Z}$. Alors $R\langle S\rangle$ possède une graduation telle que les monômes
aient un degré
$$
\deg(s_1 s_2 \ldots s_n) = \deg(s_1) + \ldots + \deg(s_n)
$$
Dans ce contexte, l'application
$$
\Mor_{\text{}R\text{-alg graduées}}(R\langle S \rangle, A) \to
\text{Map}_{\text{ensembles gradués}}(S,
A),\quad \varphi \longmapsto (s \mapsto \varphi(s))
$$
est une bijection pour chaque $R$-algèbre graduée $A$.

\medskip\noindent
Si $A$ est une $R$-algèbre graduée et que $S$ est un ensemble gradué,
alors nous pouvons de manière similaire former $A\langle S
\rangle$. Les éléments de $A\langle S \rangle$ sont
des sommes d'éléments de la forme
$$
a_0 s_1 a_1 s_2 \ldots a_{n - 1} s_n a_n
$$
avec $a_i \in A$ modulo les relations selon lesquelles ces
expressions sont $R$-multilinéaires en $(a_0, \ldots,
a_n)$. Ainsi, pour chaque suite $s_1, \ldots, s_n$ d'éléments de $S$,
il y a une inclusion
$$
A \otimes_R \ldots \otimes_R A \subset A\langle S \rangle
$$
et l'algèbre est la somme directe de celles-ci. Avec cette définition, on
vérifie que l'application
$$
\Mor_{\text{}R\text{-alg graduées}}(A\langle S \rangle, B) \to
\Mor_{\text{}R\text{-alg graduées}}(A, B) \times
\text{Map}_{\text{ensembles gradués}}(S, B),
$$
qui envoie $\varphi$ sur $(\varphi|_A, (s \mapsto \varphi(s)))$ est
une bijection pour chaque $R$-algèbre graduée $B$. Nous
observons que si $A$ était une $R$-algèbre graduée libre,
alors $A\langle S \rangle$ l'est aussi.

\medskip\noindent
Supposons que $A$ soit une $R$-algèbre différentielle graduée
et que $S$ soit un ensemble gradué. Supposons de plus que pour
chaque $s \in S$, on nous donne un élément homogène $f_s \in A$ avec $\deg(f_s) = \deg(s)
+ 1$ et $\text{d}f_s = 0$. Alors il existe une structure unique
d'algèbre différentielle graduée sur $A\langle S \rangle$ avec
$\text{d}(s) = f_s$. Par exemple, étant donnés $a, b, c \in A$ et
$s, t \in S$, on définirait
\begin{align*}
\text{d}(asbtc)
& =
\text{d}(a)sbtc + (-1)^{\deg(a)}a f_s b t c +
(-1)^{\deg(a) + \deg(s)} as\text{d}(b)tc \\ &
+ (-1)^{\deg(a) + \deg(s) + \deg(b)} asb f_t c +
(-1)^{\deg(a) + \deg(s) + \deg(b) + \deg(t)} asbt\text{d}(c)
\end{align*}
Nous omettons les détails.

\begin{lemma}
\label{dga-lemma-K-flat-resolution}
Soit $R$ un anneau. Soit $(B, \text{d})$ une $R$-algèbre différentielle
graduée. Il existe un quasi-isomorphisme $(A, \text{d}) \to (B, \text{d})$ de
$R$-algèbres différentielles graduées avec les propriétés suivantes
\begin{enumerate}
\item $A$ est K-plat en tant que complexe de $R$-modules,
\item $A$ est une $R$-algèbre graduée libre.
\end{enumerate}
\end{lemma}

\begin{proof}
Tout d'abord, nous affirmons que nous pouvons trouver $(A_0, \text{d}) \to
(B, \text{d})$ vérifiant (1) et (2) et induisant une surjection
en cohomologie. En effet, prenons un ensemble gradué $S$ et pour
chaque $s \in S$ un élément homogène $b_s \in \Ker(d : B \to B)$ de degré $\deg(s)$
tel que les classes $\overline{b}_s$ dans $H^*(B)$
engendrent $H^*(B)$ en tant que $R$-module.
Ensuite, nous pouvons définir $A_0 = R\langle S \rangle$ avec une différentielle nulle
et $A_0 \to B$ défini en envoyant $s$ sur $b_s$.

\medskip\noindent
Étant donné $A_0 \to B$ induisant une surjection en cohomologie, nous
construisons une suite
$$
A_0 \to A_1 \to A_2 \to \ldots B
$$
par récurrence. Étant donné $A_n \to B$, prenons un ensemble gradué
$S_n$ et, pour chaque $s \in S_n$, soit $a_s \in \Ker(\text{d} : A_n \to A_n)$
un élément homogène de degré $\deg(s) + 1$ dont l'image
est une classe $\overline{a}_s$ dans $H^*(A_n)$ qui
s'envoie sur zéro dans $H^*(B)$. Nous choisissons $S_n$ assez grand pour
que les éléments $\overline{a}_s$ engendrent $\Ker(H^*(A_n) \to H^*(B))$ comme
$R$-module. Ensuite, nous posons
$$
A_{n + 1} = A_n\langle S_n \rangle
$$
avec la différentielle donnée par $\text{d}(s) = a_s$ ; voir la discussion
ci-dessus. Alors chaque $(A_n, \text{d})$ satisfait (1) et (2), nous omettons les détails.
L'application $H^*(A_n) \to H^*(B)$ est surjective comme c'était le cas pour $n = 0$.

\medskip\noindent
Il est clair que $A = \colim A_n$ est une $R$-algèbre graduée libre.
Elle est K-plate d'après Compléments d'algèbre, lemme \ref{more-algebra-lemma-colimit-K-flat}.
L'application $H^*(A) \to H^*(B)$ est un isomorphisme car elle est surjective et
injective : tout élément de $H^*(A)$ provient d'un élément de $H^*(A_n)$
pour un certain $n$ et s'il devient nul dans $H^*(B)$, alors il devient nul
dans $H^*(A_{n + 1})$, donc dans $H^*(A)$.
\end{proof}

\noindent
Comme application, nous prouvons la version « correcte »
du lemme \ref{dga-lemma-compose-tensor-functors-general-algebra}.

\begin{lemma}
\label{dga-lemma-compose-tensor-functors-tor}
Soit $R$ un anneau. Soient $(A, \text{d})$, $(B, \text{d})$, et
$(C, \text{d})$ des $R$-algèbres différentielles graduées.
Supposons que $A \otimes_R C$ représente $A \otimes^\mathbf{L}_R C$ dans
$D(R)$. Soit $N$ un $(A, B)$-bimodule différentiel
gradué. Soit $N'$ un $(B, C)$-bimodule différentiel gradué.
Alors la composition
$$
\xymatrix{
D(A, \text{d}) \ar[rr]^{- \otimes_A^\mathbf{L} N} & &
D(B, \text{d}) \ar[rr]^{- \otimes_B^\mathbf{L} N'} & &
D(C, \text{d})
}
$$
est isomorphe à $- \otimes_A^\mathbf{L} N''$ pour un certain $(A, C)$-bimodule
différentiel gradué $N''$.
\end{lemma}

\begin{proof}
En utilisant le lemme
\ref{dga-lemma-K-flat-resolution}, nous choisissons un quasi-isomorphisme $(B', \text{d}) \to (B,
\text{d})$ avec $B'$ K-plat en tant que complexe de
$R$-modules. D'après le lemme
\ref{dga-lemma-qis-equivalence}, le foncteur $-\otimes^\mathbf{L}_{B'} B : D(B', \text{d}) \to D(B,
\text{d})$ est une équivalence dont le quasi-inverse est donné par la
restriction. Notons que la restriction est canoniquement isomorphe au
foncteur $- \otimes^\mathbf{L}_B B : D(B, \text{d}) \to D(B', \text{d})$ où
$B$ est considéré comme un bimodule sur $(B, B')$.
Il suffit donc de prouver le lemme pour les compositions
$$
D(A) \to D(B) \to D(B'),\quad
D(B') \to D(B) \to D(C),\quad
D(A) \to D(B') \to D(C).
$$
Le premier cas résulte du lemme
\ref{dga-lemma-compose-tensor-functors} parce que $B'$ est K-plat en tant que complexe de
$R$-modules. Le second cas est vrai parce
que $B \otimes_B^\mathbf{L} N' = N' = B \otimes_B N'$ et
donc le lemme \ref{dga-lemma-compose-tensor-functors-general} s'applique. Ainsi,
nous nous ramenons au cas où $B$ est K-plat en tant que complexe de
$R$-modules.

\medskip\noindent
Supposons que $B$ soit K-plat en tant que complexe de $R$-modules. Il
suffit de montrer que (\ref{dga-equation-plain-versus-derived-algebras}) est
un isomorphisme,
voir le lemme \ref{dga-lemma-compose-tensor-functors-general-algebra}.
Choisissons un quasi-isomorphisme $L \to A$ où $L$ est un $R$-module
différentiel gradué qui a la propriété (P). Alors il est clair que $P = L
\otimes_R B$ a la propriété (P) en tant que $B$-module différentiel gradué. Par
conséquent, nous devons montrer que $P \to A \otimes_R B$
induit un quasi-isomorphisme
$$
P \otimes_B (B \otimes_R C)
\longrightarrow
(A \otimes_R B) \otimes_B (B \otimes_R C)
$$
Nous pouvons réécrire ceci comme
$$
P \otimes_R B \otimes_R C \longrightarrow A \otimes_R B \otimes_R C
$$
Puisque $B$ est K-plat en tant que complexe de
$R$-modules, il découle
de Compléments d'algèbre, lemme
\ref{more-algebra-lemma-K-flat-quasi-isomorphism} qu'il suffit de
montrer que
$$
P \otimes_R C \to A \otimes_R C
$$
est un quasi-isomorphisme, ce qui est exactement notre hypothèse.
\end{proof}

\noindent
Le lemme suivant n'appartient pas vraiment à cette section, mais il ne
semble pas y avoir de place naturelle appropriée pour lui.

\begin{lemma}
\label{dga-lemma-countable}
Soit $(A, \text{d})$ une algèbre différentielle graduée avec
$H^i(A)$ dénombrable pour chaque $i$. Soit $M$ un objet de $D(A, \text{d})$. Alors
les énoncés suivants sont équivalents
\begin{enumerate}
\item $M = \text{hocolim} E_n$ avec $E_n$ compact dans $D(A, \text{d})$, et
\item $H^i(M)$ est dénombrable pour chaque $i$.
\end{enumerate}
\end{lemma}

\begin{proof}
Supposons que (1) soit vrai. Ensuite, nous avons $H^i(M) = \colim
H^i(E_n)$ par Catégories dérivées, lemme
\ref{derived-lemma-cohomology-of-hocolim}. Ainsi, il suffit de prouver que $H^i(E_n)$ est dénombrable pour chaque $n$.
Par la proposition \ref{dga-proposition-compact}, nous voyons que
$E_n$ est isomorphe dans $D(A, \text{d})$ à un facteur direct d'un
module différentiel gradué $P$ qui possède une filtration finie
$F_\bullet$ par des sous-modules différentiels gradués telle
que les $F_jP/F_{j - 1}P$ soient des sommes directes finies de décalages de $A$.
Par hypothèse, les groupes $H^i(F_jP/F_{j - 1}P)$ sont dénombrables.
En raisonnant par récurrence sur la longueur de la filtration et en
utilisant la suite exacte longue de cohomologie, nous concluons que (2) est
vrai. L'implication intéressante est l'autre.

\medskip\noindent
Nous affirmons qu'il existe une sous-algèbre différentielle
graduée dénombrable $A' \subset A$ telle que l'application d'inclusion
$A' \to A$ définisse un isomorphisme sur la cohomologie. Pour
construire $A'$, nous choisissons des sous-algèbres différentielles graduées
dénombrables
$$
A_1 \subset A_2 \subset A_3 \subset \ldots
$$
de telle sorte que (a) $H^i(A_1) \to H^i(A)$ soit surjectif, et
(b) pour $n > 1$ le noyau de l'application $H^i(A_{n - 1}) \to
H^i(A_n)$ soit le même que le noyau de l'application $H^i(A_{n - 1}) \to
H^i(A)$. Pour construire $A_1$, prenons une collection dénombrable de
cochaînes $S \subset A$ engendrant la cohomologie de $A$ (comme anneau ou
comme groupe abélien gradué) et soit
$A_1$ la sous-algèbre différentielle graduée de $A$ engendrée par $S$. Pour
construire $A_n$ à partir de $A_{n - 1}$ pour chaque cochaîne $a \in A_{n - 1}^i$
qui s'envoie sur zéro dans $H^i(A)$, choisissons $s_a \in A^{i
- 1}$ avec $\text{d}(s_a) = a$ et soit $A_n$ la sous-algèbre
différentielle graduée de $A$ engendrée par $A_{n - 1}$ et les éléments $s_a$. Enfin,
prenons $A' = \bigcup A_n$.

\medskip\noindent
D'après le lemme
\ref{dga-lemma-qis-equivalence}, l'application de restriction $D(A, \text{d}) \to D(A',
\text{d})$, $M \mapsto M_{A'}$ est une équivalence. Puisque les
groupes de cohomologie de $M$ et $M_{A'}$ sont les mêmes, nous voyons
qu'il suffit de prouver l'implication (2) $\Rightarrow$ (1)
pour $(A', \text{d})$.

\medskip\noindent
Supposons que $A$ soit dénombrable. Par le même type d'argument que celui
donné ci-dessus, nous voyons que pour $M$ dans $D(A,
\text{d})$, les énoncés suivants sont équivalents : $H^i(M)$ est dénombrable pour
chaque $i$ et $M$ peut être représenté par un module différentiel gradué dénombrable.
Par conséquent, afin de prouver l'implication (2) $\Rightarrow$ (1),
nous nous ramenons à la situation décrite dans le paragraphe suivant.

\medskip\noindent
Supposons que $A$ soit dénombrable et que $M$ soit un module différentiel gradué
dénombrable sur $A$. Nous affirmons qu'il existe un homomorphisme $P \to M$ de
$A$-modules différentiels gradués vérifiant les propriétés suivantes :
\begin{enumerate}
\item $P \to M$ est un quasi-isomorphisme,
\item $P$ possède la propriété (P), et
\item $P$ est dénombrable.
\end{enumerate}
En regardant la démonstration de la construction des P-résolutions dans
le lemme \ref{dga-lemma-resolve}, nous voyons qu'il suffit de montrer que nous
pouvons prouver le lemme \ref{dga-lemma-good-quotient}
dans le cadre des modules différentiels gradués dénombrables. Cela
découle immédiatement de la démonstration.

\medskip\noindent
Supposons que $A$ soit dénombrable et que $M$ soit un module
différentiel gradué dénombrable avec la propriété (P). Choisissons une filtration
$$
0 = F_{-1}P \subset F_0P \subset F_1P \subset \ldots \subset P
$$
par des sous-modules différentiels gradués, telle que l'on ait
\begin{enumerate}
\item $P = \bigcup F_pP$,
\item $F_iP \to F_{i + 1}P$ est un monomorphisme admissible,
\item des isomorphismes de modules différentiels
gradués $F_iP/F_{i - 1}P \to \bigoplus_{j \in J_i} A[k_j]$ pour
certains ensembles $J_i$ et entiers $k_j$.
\end{enumerate}
Bien sûr, $J_i$ est dénombrable pour chaque $i$. Pour chaque $i$
et $j \in J_i$, choisissons $x_{i, j} \in F_iP$ de degré $k_j$ dont
l'image dans $F_iP/F_{i - 1}P$ engendre le facteur direct correspondant
à $j$.

\medskip\noindent
Affirmation : Étant donné $n$ et des sous-ensembles finis $S_i \subset J_i$, $i = 1,
\ldots, n$, il existe des sous-ensembles finis $S_i \subset T_i \subset J_i$, $i = 1,
\ldots, n$ tels que $P' = \bigoplus_{i \leq n} \bigoplus_{j \in T_i} Ax_{i, j}$ soit un
sous-module différentiel gradué de $P$. Cela a été montré dans la
démonstration du lemme \ref{dga-lemma-factor-through-nicer} mais cela se montre aussi
facilement directement : les éléments $x_{i, j}$ engendrent librement $P$ en
tant que module à droite sur $A$. La structure de $P$ montre que
$$
\text{d}(x_{i, j}) = \sum\nolimits_{i' < i} x_{i', j'}a_{i', j'}
$$
où bien sûr la somme est finie.
Ainsi, étant donnés $S_0, \ldots, S_n$, nous pouvons d'abord
choisir $S_0 \subset S'_0, \ldots, S_{n - 1} \subset S'_{n - 1}$
avec $\text{d}(x_{n, j}) \in \bigoplus_{i' < n, j' \in S'_{i'}} x_{i', j'}A$ pour
tous $j \in S_n$. Ensuite, par récurrence sur $n$, nous pouvons
choisir $S'_0 \subset T_0, \ldots, S'_{n - 1} \subset T_{n - 1}$
pour s'assurer que $\bigoplus_{i' < n, j' \in T_{i'}} x_{i', j'}A$ soit
un $A$-sous-module différentiel gradué. En définissant $T_n = S_n$, nous trouvons
que $P' = \bigoplus_{i \leq n, j \in T_i} x_{i, j}A$ convient.

\medskip\noindent
D'après l'assertion, il est clair que $P = \bigcup P'_n$
est une union croissante dénombrable de $P'_n$ comme
indiqué ci-dessus. Par construction, chaque $P'_n$ est un module différentiel
gradué avec la propriété (P) tel que la filtration soit finie et que les
quotients successifs sont des sommes directes finies de décalages de $A$.
Ainsi, $P'_n$ définit un objet compact de $D(A, \text{d})$, voir par
exemple la proposition \ref{dga-proposition-compact}.
Puisque $P = \text{hocolim} P'_n$ dans $D(A, \text{d})$ par
le lemme \ref{dga-lemma-homotopy-colimit}, la
démonstration de l'implication (2) $\Rightarrow$ (1) est complète.
\end{proof}








% OK, start here.
%

\chapter{Algèbre à puissances divisées}



\phantomsection
\label{dpa-section-phantom}


\section{Introduction}
\label{dpa-section-introduction}

\noindent
Dans ce chapitre, nous étudions les algèbres à puissances divisées et les
constructions qu'elles permettent. Une référence est l'ouvrage \cite{Berthelot}.





\section{Puissances divisées}
\label{dpa-section-divided-powers}

\noindent
Dans cette section, nous rassemblons quelques résultats sur les anneaux à puissances divisées.
Nous adoptons la convention $0! = 1$ (un produit vide devant valoir $1$).

\begin{definition}
\label{dpa-definition-divided-powers}
Soit $A$ un anneau. Soit $I$ un idéal de $A$. Une famille d'applications
$\gamma_n : I \to I$, $n > 0$, est appelée une {\it structure d'idéal à puissances divisées}
sur $I$ si, pour tous $n \geq 0$, $m > 0$, $x, y \in I$ et $a \in A$, on a
\begin{enumerate}
\item $\gamma_1(x) = x$ ; on pose aussi $\gamma_0(x) = 1$,
\item $\gamma_n(x)\gamma_m(x) = \frac{(n + m)!}{n! m!} \gamma_{n + m}(x)$,
\item $\gamma_n(ax) = a^n \gamma_n(x)$,
\item $\gamma_n(x + y) = \sum_{i = 0, \ldots, n} \gamma_i(x)\gamma_{n - i}(y)$,
\item $\gamma_n(\gamma_m(x)) = \frac{(nm)!}{n! (m!)^n} \gamma_{nm}(x)$.
\end{enumerate}
\end{definition}

\noindent
Remarquons que les nombres rationnels $\frac{(n + m)!}{n! m!}$
et $\frac{(nm)!}{n! (m!)^n}$ qui interviennent dans la définition sont en réalité des
entiers : le premier est le nombre de façons de choisir $n$ éléments parmi $n + m$,
et le second compte les façons de répartir $nm$
objets en $n$ groupes de $m$ objets.
Voici quelques remarques sur la définition qui montrent que
$\gamma_n(x)$ remplace $x^n/n!$ dans $I$.

\begin{lemma}
\label{dpa-lemma-silly}
Soit $A$ un anneau. Soit $I$ un idéal de $A$.
\begin{enumerate}
\item Si $\gamma$ est une structure de puissances divisées\footnote{Ici
et dans la suite, $\gamma$ désigne une suite
d'applications $\gamma_1, \gamma_2, \gamma_3, \ldots$ de $I$ dans $I$.}
sur $I$, alors
$n! \gamma_n(x) = x^n$ pour $n \geq 1$, $x \in I$.
\end{enumerate}
Supposons que $A$ soit sans torsion en tant que $\mathbf{Z}$-module.
\begin{enumerate}
\item[(2)] Une structure de puissances divisées sur $I$, si elle existe, est unique.
\item[(3)] Si $\gamma_n : I \to I$ sont des applications, alors
$$
\gamma\text{ est une structure de puissances divisées}
\Leftrightarrow
n! \gamma_n(x) = x^n\ \forall x \in I, n \geq 1.
$$
\item[(4)] L'idéal $I$ admet une structure de
puissances divisées si et seulement
s'il existe des éléments $x_i$ qui engendrent $I$ comme idéal et tels que,
pour tout $n \geq 1$, on ait $x_i^n \in (n!)I$.
\end{enumerate}
\end{lemma}

\begin{proof}
Démonstration de (1). Si $\gamma$ est une structure de puissances divisées, alors la
condition (2) (appliquée à $1$ et $n-1$ à la place de $n$ et $m$)
entraîne que $n \gamma_n(x) = \gamma_1(x)\gamma_{n - 1}(x)$. Par
récurrence et grâce à la condition (1), on obtient donc $n! \gamma_n(x) = x^n$.

\medskip\noindent
Supposons que $A$ soit sans torsion en tant que $\mathbf{Z}$-module.
Démonstration de (2). Cela résulte immédiatement de (1).

\medskip\noindent
Démonstration de (3). Supposons que $n! \gamma_n(x) = x^n$ pour tout $x \in I$ et tout
$n \geq 1$. Comme $A \subset A \otimes_{\mathbf{Z}} \mathbf{Q}$, il suffit
de démontrer les axiomes (1) -- (5) de la définition
\ref{dpa-definition-divided-powers} lorsque $A$ est une $\mathbf{Q}$-algèbre.
Dans ce cas, $\gamma_n(x) = x^n/n!$, et la vérification de
(1) -- (5) est immédiate ; par exemple, (4) correspond à la formule du
binôme
$$
(x + y)^n = \sum_{i = 0, \ldots, n} \frac{n!}{i!(n - i)!} x^iy^{n - i}
$$
Nous laissons au lecteur le soin d'effectuer ces
vérifications afin de s'assurer que les coefficients sont corrects.

\medskip\noindent
Démonstration de (4). Supposons que des éléments $x_i$ engendrent $I$ comme idéal
et vérifient $x_i^n \in (n!)I$ pour tout $n \geq 1$. Nous affirmons que,
pour tout $x \in I$, on a $x^n \in (n!)I$. Si cette assertion est vraie,
on peut poser $\gamma_n(x) = x^n/n!$, ce qui définit une structure de puissances divisées d'après
(3). Pour démontrer l'assertion, remarquons qu'elle est vraie pour $x = ax_i$. Elle est
donc vraie pour un ensemble d'éléments qui engendre $I$ comme groupe abélien.
Par récurrence sur la longueur d'une expression en fonction de ces générateurs, il
suffit de la vérifier pour $x + y$ lorsqu'elle est vraie pour $x$ et $y$. Cela
découle immédiatement de la formule du binôme.
\end{proof}

\begin{example}
\label{dpa-example-ideal-generated-by-p}
Soit $p$ un nombre premier.
Soit $A$ un anneau dans lequel tout entier $n$ non divisible par $p$
est inversible, autrement dit, $A$ est une $\mathbf{Z}_{(p)}$-algèbre. Alors
$I = pA$ admet une structure canonique de puissances divisées. Plus précisément, pour
$x = pa \in I$, on pose
$$
\gamma_n(x) = \frac{p^n}{n!} a^n
$$
On vérifie immédiatement que $p^n/n! \in p\mathbf{Z}_{(p)}$
pour $n \geq 1$ (cela découle par exemple du fait
que l'exposant de $p$ dans la décomposition en facteurs premiers de $n!$ est
$\left\lfloor n/p \right\rfloor + \left\lfloor n/p^2 \right\rfloor
+ \left\lfloor n/p^3 \right\rfloor + \ldots$),
de sorte que la définition a un sens et fournit une suite
d'applications $\gamma_n : I \to I$. Il est immédiat de vérifier que
les conditions (1) -- (5) de la
définition \ref{dpa-definition-divided-powers} sont satisfaites.
On peut aussi observer que la définition convient manifestement à
$A_0 = \mathbf{Z}_{(p)}$, puis appliquer le
lemme \ref{dpa-lemma-gamma-extends}.
\end{example}

\noindent
On remarque que $\gamma_n\left(0\right) = 0$ pour tout idéal $I$ de
$A$ et toute structure de puissances divisées $\gamma$ sur $I$. (Cela résulte
de l'axiome (3) de la définition \ref{dpa-definition-divided-powers},
en prenant $a=0$.)

\begin{lemma}
\label{dpa-lemma-check-on-generators}
Soit $A$ un anneau. Soit $I$ un idéal de $A$. Soit $\gamma_n : I \to I$,
$n \geq 1$, une suite d'applications. Supposons que
\begin{enumerate}
\item[(a)] les propriétés (1), (3) et (4) de la définition \ref{dpa-definition-divided-powers} soient
vérifiées pour tous $x, y \in I$, et
\item[(b)] les propriétés (2) et (5) soient vérifiées pour tout $x$ appartenant
à un ensemble d'éléments qui engendre $I$ comme idéal.
\end{enumerate}
Alors $\gamma$ est une structure de puissances divisées sur $I$.
\end{lemma}

\begin{proof}
Dans cette démonstration, les numéros (1), (2), (3), (4), (5) renvoient aux
conditions énoncées dans la définition \ref{dpa-definition-divided-powers}. Il
résulte de (3) que, si (2) et (5) sont vérifiées pour $x$, elles le
sont pour $ax$, quel que soit $a \in A$. L'hypothèse (b) entraîne donc que (2)
et (5) sont vérifiées pour un ensemble d'éléments qui
engendre $I$ comme groupe abélien. Par récurrence sur la longueur d'une
expression en fonction de ces générateurs, il suffit par conséquent de montrer que, si
$x, y \in I$ et si (2) et (5) sont vérifiées pour $x$ et $y$, alors elles le sont
pour $x + y$.

\medskip\noindent
Démonstration de (2) pour $x + y$. D'après (4), on a
$$
\gamma_n(x + y)\gamma_m(x + y) =
\sum\nolimits_{i + j = n,\ k + l = m}
\gamma_i(x)\gamma_k(x)\gamma_j(y)\gamma_l(y)
$$
En utilisant (2) pour $x$ et $y$, ceci vaut
$$
\sum\frac{(i + k)!}{i!k!}\frac{(j + l)!}{j!l!}
\gamma_{i + k}(x)\gamma_{j + l}(y)
$$
En comparant avec le développement
$$
\gamma_{n + m}(x + y) = \sum\gamma_a(x)\gamma_b(y)
$$
on voit qu'il faut démontrer que, pour $a + b = n + m$, on a
$$
\sum\nolimits_{i + k = a,\ j + l = b,\ i + j = n,\ k + l = m}
\frac{(i + k)!}{i!k!}\frac{(j + l)!}{j!l!}
=
\frac{(n + m)!}{n!m!}.
$$
Plutôt que de démontrer directement cette identité, remarquons qu'elle est
vraie pour l'idéal $I = (x, y)$ de l'anneau de polynômes $\mathbf{Q}[x, y]$,
car $\gamma_n(f) = f^n/n!$, $f \in I$, définit une structure de
puissances divisées sur $I$. L'identité entre nombres rationnels ci-dessus est donc vraie.

\medskip\noindent
Démonstration de (5) pour $x + y$, sachant que (1) -- (4) sont vérifiées et que
(5) l'est pour $x$ et $y$. Nous nous ramenons de nouveau à une identité
entre nombres rationnels. En effet, (4) permet d'écrire
$\gamma_n(\gamma_m(x + y)) = \gamma_n(\sum\gamma_i(x)\gamma_j(y))$.
Une nouvelle application de (4) permet d'écrire
$\gamma_n(\gamma_m(x + y))$ comme une somme de termes qui sont des produits
de facteurs de la forme $\gamma_k(\gamma_i(x)\gamma_j(y))$.
Si $i > 0$, alors
\begin{align*}
\gamma_k(\gamma_i(x)\gamma_j(y)) & =
\gamma_j(y)^k\gamma_k(\gamma_i(x)) \\
& = \frac{(ki)!}{k!(i!)^k} \gamma_j(y)^k \gamma_{ki}(x) \\
& =
\frac{(ki)!}{k!(i!)^k} \frac{(kj)!}{(j!)^k} \gamma_{ki}(x) \gamma_{kj}(y)
\end{align*}
où l'on a utilisé (3) dans la première égalité, (5) pour $x$ dans la deuxième et
(2) exactement $k$ fois dans la troisième. En utilisant (5) pour $y$, on voit
que la même égalité vaut lorsque $i = 0$. En poursuivant ainsi et en
utilisant tous les axiomes sauf (5), on obtient
$$
\gamma_n(\gamma_m(x + y)) =
\sum\nolimits_{i + j = nm} c_{ij}\gamma_i(x)\gamma_j(y)
$$
où les $c_{ij} \in\mathbf{Z}$ sont certaines constantes
universelles. Comme précédemment, le fait que l'égalité soit valable dans l'anneau de polynômes $\mathbf{Q}[x, y]$
implique que tous les coefficients $c_{ij}$ sont égaux à $(nm)!/n!(m!)^n$,
comme souhaité.
\end{proof}

\begin{lemma}
\label{dpa-lemma-two-ideals}
Soit $A$ un anneau muni de deux idéaux $I, J \subset A$.
Soit $\gamma$ une structure de puissances divisées sur $I$, et soit
$\delta$ une structure de puissances divisées sur $J$.
Alors
\begin{enumerate}
\item $\gamma$ et $\delta$ coïncident sur $IJ$,
\item si $\gamma$ et $\delta$ coïncident sur $I \cap J$, alors elles
sont les restrictions d'une unique structure de puissances divisées $\epsilon$
sur $I + J$.
\end{enumerate}
\end{lemma}

\begin{proof}
Soient $x \in I$ et $y \in J$. Alors
$$
\gamma_n(xy) = y^n\gamma_n(x) = n! \delta_n(y) \gamma_n(x) =
\delta_n(y) x^n = \delta_n(xy).
$$
Par conséquent, $\gamma$ et $\delta$ coïncident sur un ensemble de générateurs (additifs)
de $IJ$. D'après la propriété (4) de la définition \ref{dpa-definition-divided-powers},
elles coïncident donc sur tout $IJ$.

\medskip\noindent
Supposons que $\gamma$ et $\delta$ coïncident sur $I \cap J$.
Soit $z \in I + J$. Écrivons $z = x + y$ avec $x \in I$ et $y \in J$.
On pose alors
$$
\epsilon_n(z) = \sum\gamma_i(x)\delta_{n - i}(y)
$$
pour tout $n \geq 1$.
Pour vérifier que cette définition ne dépend pas du choix, supposons que $z = x' + y'$ soit une
autre écriture avec $x' \in I$ et $y' \in J$. Alors
$w = x - x' = y' - y \in I \cap J$. Par conséquent,
\begin{align*}
\sum\nolimits_{i + j = n} \gamma_i(x)\delta_j(y)
& =
\sum\nolimits_{i + j = n} \gamma_i(x' + w)\delta_j(y) \\
& =
\sum\nolimits_{i' + l + j = n} \gamma_{i'}(x')\gamma_l(w)\delta_j(y) \\
& =
\sum\nolimits_{i' + l + j = n} \gamma_{i'}(x')\delta_l(w)\delta_j(y) \\
& =
\sum\nolimits_{i' + j' = n} \gamma_{i'}(x')\delta_{j'}(y + w) \\
& =
\sum\nolimits_{i' + j' = n} \gamma_{i'}(x')\delta_{j'}(y')
\end{align*}
comme souhaité. Nous avons donc défini des applications
$\epsilon_n : I + J \to I + J$ pour tout $n \geq 1$; il est
immédiat que $\epsilon_n \mid_{I} = \gamma_n$ et
$\epsilon_n \mid_{J} = \delta_n$.
Montrons ensuite les conditions (1) -- (5) de
la définition \ref{dpa-definition-divided-powers} pour la famille
d'applications $\epsilon_n$.
Les propriétés (1) et (3) sont immédiates. Pour vérifier (4),
supposons que $z = x + y$ et $z' = x' + y'$ avec $x, x' \in I$ et $y, y' \in J$,
et calculons
\begin{align*}
\epsilon_n(z + z') & =
\sum\nolimits_{a + b = n} \gamma_a(x + x')\delta_b(y + y') \\
& =
\sum\nolimits_{i + i' + j + j' = n}
\gamma_i(x) \gamma_{i'}(x')\delta_j(y)\delta_{j'}(y') \\
& =
\sum\nolimits_{k = 0, \ldots, n}
\sum\nolimits_{i+j=k} \gamma_i(x)\delta_j(y)
\sum\nolimits_{i'+j'=n-k} \gamma_{i'}(x')\delta_{j'}(y') \\
& =
\sum\nolimits_{k = 0, \ldots, n}\epsilon_k(z)\epsilon_{n-k}(z')
\end{align*}
ce qui donne le résultat voulu. Il suffit maintenant de vérifier (2) et (5) sur
les éléments de $I$ ou de $J$, voir le lemme \ref{dpa-lemma-check-on-generators}.
Cela est immédiat puisque $\gamma$ et $\delta$ sont des structures de
puissances divisées.

\medskip\noindent
Nous avons ainsi démontré l'existence d'une structure de puissances divisées $\epsilon$ sur $I+J$
dont les restrictions à $I$ et à $J$ sont $\gamma$ et $\delta$ ;
son unicité est immédiate.
\end{proof}

\begin{lemma}
\label{dpa-lemma-nil}
Soit $p$ un nombre premier. Soit $A$ un anneau, soit $I \subset A$ un idéal,
et soit $\gamma$ une structure de puissances divisées sur $I$. Supposons que $p$ soit nilpotent
dans $A/I$. Alors $I$ est localement nilpotent si et seulement si $p$ est nilpotent dans $A$.
\end{lemma}

\begin{proof}
Si $p^N = 0$ dans $A$, alors, pour $x \in I$, on a
$x^{pN} = (pN)!\gamma_{pN}(x) = 0$, car $(pN)!$ est
divisible par $p^N$. Réciproquement, supposons que $I$ soit localement nilpotent.
Comme nous avons aussi supposé que $p$ est nilpotent dans $A/I$, on a
$p^r \in I$ pour un certain $r$; ainsi $p^r$ est nilpotent, et donc $p$ est nilpotent.
\end{proof}








\section{Anneaux à puissances divisées}
\label{dpa-section-divided-power-rings}

\noindent
Les anneaux à puissances divisées forment une
catégorie. En voici la définition.

\begin{definition}
\label{dpa-definition-divided-power-ring}
Un {\it anneau à puissances divisées} est un triplet $(A, I, \gamma)$, où
$A$ est un anneau, $I \subset A$ est un idéal et $\gamma = (\gamma_n)_{n \geq 1}$
est une structure de puissances divisées sur $I$.
Un {\it homomorphisme d'anneaux à puissances divisées}
$\varphi : (A, I, \gamma) \to (B, J, \delta)$ est un homomorphisme d'anneaux
$\varphi : A \to B$ tel que $\varphi(I) \subset J$ et
$\delta_n(\varphi(x)) = \varphi(\gamma_n(x))$ pour tout $x \in I$ et tout
$n \geq 1$.
\end{definition}

\noindent
Nous dirons parfois « soit $(B, J, \delta)$ une algèbre à puissances divisées sur
$(A, I, \gamma)$ » pour indiquer que $(B, J, \delta)$ est un anneau à
puissances divisées muni d'un homomorphisme d'anneaux à puissances divisées
$(A, I, \gamma) \to (B, J, \delta)$.

\begin{lemma}
\label{dpa-lemma-limits} La
catégorie des anneaux à puissances divisées admet toutes les limites, et celles-ci coïncident avec les
limites calculées dans la catégorie des anneaux.
\end{lemma}

\begin{proof}
La limite du diagramme vide est l'anneau nul (ce qui est étrange, mais nécessaire
ici). Le produit d'une famille d'anneaux à puissances divisées $(A_t, I_t, \gamma_t)$,
$t \in T$, est $(\prod A_t, \prod I_t, \gamma)$, où
$\gamma_n((x_t)) = (\gamma_{t, n}(x_t))$.
L'égalisateur de $\alpha, \beta : (A, I, \gamma) \to (B, J, \delta)$
est simplement $C = \{a \in A \mid\alpha(a) = \beta(a)\}$, muni de l'idéal $C \cap I$
et des puissances divisées induites. Il en résulte que toutes les limites existent ;
voir Catégories, lemme \ref{categories-lemma-limits-products-equalizers}.
\end{proof}

\noindent
Le lemme suivant illustre, dans le cas des algèbres à puissances
divisées, un phénomène très général de théorie des catégories.

\begin{lemma}
\label{dpa-lemma-a-version-of-brown}
Soit $\mathcal{C}$ la catégorie des anneaux à puissances divisées. Soit
$F : \mathcal{C} \to\textit{Ensembles}$ un foncteur.
Supposons que
\begin{enumerate}
\item il existe un cardinal $\kappa$ tel que, pour tout
$f \in F(A, I, \gamma)$, il existe un morphisme
$(A', I', \gamma') \to (A, I, \gamma)$ de $\mathcal{C}$ tel que $f$
soit l'image d'un élément $f' \in F(A', I', \gamma')$ et que $|A'| \leq\kappa$;
\item $F$ commute aux limites.
\end{enumerate}
Alors $F$ est représentable, c'est-à-dire qu'il existe un objet $(B, J, \delta)$
de $\mathcal{C}$ tel que
$$
F(A, I, \gamma) = \Hom_\mathcal{C}((B, J, \delta), (A, I, \gamma))
$$
fonctoriellement en $(A, I, \gamma)$.
\end{lemma}

\begin{proof}
C'est un cas particulier de
Catégories, lemme \ref{categories-lemma-a-version-of-brown}.
\end{proof}

\begin{lemma}
\label{dpa-lemma-colimits} La
catégorie des anneaux à puissances divisées admet toutes les colimites.
\end{lemma}

\begin{proof} La
colimite du diagramme vide est $\mathbf{Z}$, muni de l'idéal à puissances divisées $(0)$.
Considérons maintenant une colimite quelconque. Soit $\mathcal{C}$ une catégorie et soit
$c \mapsto (A_c, I_c, \gamma_c)$ un diagramme. Considérons le foncteur
$$
F(B, J, \delta) = \lim_{c \in\mathcal{C}}
Hom((A_c, I_c, \gamma_c), (B, J, \delta))
$$
Tout élément $f = (f_c)_{c \in C} \in F(B, J, \delta)$ possède la
propriété suivante : les images $f_c(A_c)$ engendrent un sous-anneau $B'$ de $B$ dont le cardinal est
borné par un cardinal $\kappa$, et les images $f_c(I_c)$ engendrent
un idéal à puissances divisées $J'$ de $B'$. On obtient ainsi une factorisation de
$f$ par un élément $f'$ de $F(B')$, suivie de l'inclusion $B' \to B$. En outre,
$F$ commute aux limites. Nous pouvons donc appliquer
le lemme \ref{dpa-lemma-a-version-of-brown}
pour conclure que $F$ est représentable.
\end{proof}

\begin{remark}
\label{dpa-remark-pushout}
Soit
$$
\xymatrix{
(B, J, \delta) \ar[r] & (B'', J'', \delta'') \\
(A, I, \gamma) \ar[r] \ar[u] & (B', J', \delta') \ar[u]
}
$$
un carré cocartésien dans la catégorie des anneaux à puissances divisées (les sommes amalgamées existent
d'après le lemme \ref{dpa-lemma-colimits}). Alors
\begin{enumerate}
\item $B''/J'' = B/J \otimes_{A/I} B'/J'$,
\item l'application $\varphi : B \otimes_A B' \to B''$ est surjective, et
\item l'application $J \otimes_A B' + B \otimes_A J' \to J''$ induite par $\varphi$
est surjective.
\end{enumerate} Pour
démontrer (1), considérons les applications $(B'', J'', \delta'') \to (C, (0), \emptyset)$
et appliquons la définition d'un carré cocartésien. Pour démontrer (2), considérons
l'image $\tilde B = \varphi(B \otimes_A B')$ et l'idéal
$\tilde J = \varphi(J \otimes_A B' + B \otimes_A J')$. Alors il
est clair que $\delta''_n(\tilde J) \subset\tilde J$ pour tous $n \geq 1$
par compatibilité de $\delta''$ avec $\delta$ et $\delta'$.
On a donc $\tilde B = B''$ et $\tilde J = J''$
par la propriété universelle du carré cocartésien.
\end{remark}

\begin{remark}
\label{dpa-remark-forgetful}
Le foncteur d'oubli $(A, I, \gamma) \mapsto A$ ne commute pas aux
colimites. Par exemple, soit
$$
\xymatrix{
(B, J, \delta) \ar[r] & (B'', J'', \delta'') \\
(A, I, \gamma) \ar[r] \ar[u] & (B', J', \delta') \ar[u]
}
$$
un carré cocartésien dans la catégorie des anneaux à puissances divisées, comme dans
la remarque \ref{dpa-remark-pushout}. En général, l'application
$B \otimes_A B' \to B''$ n'est pas un
isomorphisme. Un exemple explicite est fourni par
$(A, I, \gamma) = (\mathbf{Z}, (0), \emptyset)$,
$(B, J, \delta) = (\mathbf{Z}/4\mathbf{Z}, 2\mathbf{Z}/4\mathbf{Z}, \delta)$,
et
$(B', J', \delta') =
(\mathbf{Z}/4\mathbf{Z}, 2\mathbf{Z}/4\mathbf{Z}, \delta')$
où $\delta_2(2) = 2$ et $\delta'_2(2) = 0$.
Plus précisément, en appliquant le lemme \ref{dpa-lemma-need-only-gamma-p}, on
prend pour $\delta$, respectivement\ $\delta'$, l'unique
structure de puissances divisées sur $J$, respectivement\ $J'$, telle que
$\delta_2 : J \to J$, respectivement\ $\delta'_2 : J' \to J'$,
soit l'application $0 \mapsto 0, 2 \mapsto 2$, respectivement\ $0 \mapsto 0, 2 \mapsto 0$.
Alors $(B'', J'', \delta'') = (\mathbf{F}_2, (0), \emptyset)$,
ce qui ne coïncide pas avec le produit tensoriel.
\end{remark}


\section{Prolongement des puissances divisées}
\label{dpa-section-extend}

\noindent
Voici la définition.

\begin{definition}
\label{dpa-definition-extends} Étant
donnés un anneau à puissances divisées $(A, I, \gamma)$ et un homomorphisme d'anneaux
$A \to B$, nous disons que $\gamma$ {\it se prolonge} à $B$ s'il existe une
structure de puissances divisées $\bar\gamma$ sur $IB$ telle que
$(A, I, \gamma) \to (B, IB, \bar\gamma)$ soit un homomorphisme
d'anneaux à puissances divisées.
\end{definition}

\begin{lemma}
\label{dpa-lemma-gamma-extends}
Soit $(A, I, \gamma)$ un anneau à puissances divisées.
Soit $A \to B$ un homomorphisme
d'anneaux. Si $\gamma$ se prolonge à $B$, ce prolongement est unique.
Supposons que l'une au moins des conditions suivantes soit vérifiée :
\begin{enumerate}
\item $IB = 0$,
\item $I$ est principal, ou
\item $A \to B$ est plat.
\end{enumerate}
Alors $\gamma$ se prolonge à $B$.
\end{lemma}

\begin{proof}
Tout élément de $IB$ s'écrit comme une somme finie
$\sum\nolimits_{i=1}^t b_ix_i$, avec
$b_i \in B$ et $x_i \in I$. Si $\gamma$ admet un prolongement $\bar\gamma$ à $IB$,
alors $\bar\gamma_n(x_i) = \gamma_n(x_i)$.
Les conditions (3) et (4) de la
définition \ref{dpa-definition-divided-powers} impliquent donc que
$$
\bar\gamma_n(\sum\nolimits_{i=1}^t b_ix_i) =
\sum\nolimits_{n_1 + \ldots + n_t = n}
\prod\nolimits_{i = 1}^t b_i^{n_i}\gamma_{n_i}(x_i)
$$
On voit ainsi que $\bar\gamma$, si elle existe, est unique.

\medskip\noindent
Si $IB = 0$, il suffit de poser $\bar\gamma_n(0) = 0$. Si $I = (x)$,
nous définissons $\bar\gamma_n(bx) = b^n\gamma_n(x)$. Cette définition ne dépend pas du coefficient choisi
: si $b'x = bx$, c'est-à-dire si $(b - b')x = 0$, alors
\begin{align*}
b^n\gamma_n(x) - (b')^n\gamma_n(x)
& =
(b^n - (b')^n)\gamma_n(x) \\
& =
(b^{n - 1} + \ldots + (b')^{n - 1})(b - b')\gamma_n(x) = 0
\end{align*}
car $\gamma_n(x)$ est divisible par $x$ (puisque
$\gamma_n(I) \subset I$) et est donc annulé par $b - b'$.
Vérifions maintenant les conditions (1) -- (5) de
la définition \ref{dpa-definition-divided-powers}. Les
conditions (1), (2), (3) et (5) découlent immédiatement de la
construction. Pour (4), soient $y, z \in IB$, disons $y = bx$ et $z = cx$. Alors
$y + z = (b + c)x$. Par conséquent,
\begin{align*}
\bar\gamma_n(y + z)
& =
(b + c)^n\gamma_n(x) \\
& =
\sum\frac{n!}{i!(n - i)!}b^ic^{n -i}\gamma_n(x) \\
& =
\sum b^ic^{n - i}\gamma_i(x)\gamma_{n - i}(x) \\
& =
\sum\bar\gamma_i(y)\bar\gamma_{n -i}(z)
\end{align*}
comme voulu.

\medskip\noindent
Supposons enfin que $A \to B$ soit plat. Soient $b_1, \ldots, b_r \in B$ et
$x_1, \ldots, x_r \in I$. Alors
$$
\bar\gamma_n(\sum b_ix_i) =
\sum b_1^{e_1} \ldots b_r^{e_r} \gamma_{e_1}(x_1) \ldots\gamma_{e_r}(x_r)
$$
où la somme porte sur les $e_1 + \ldots + e_r = n$,
si $\bar\gamma_n$ existe. Supposons ensuite donnés $c_1, \ldots, c_s \in B$
et $a_{ij} \in A$ tels que $b_i = \sum a_{ij}c_j$.
En posant $y_j = \sum a_{ij}x_i$, nous affirmons que
$$
\sum b_1^{e_1} \ldots b_r^{e_r} \gamma_{e_1}(x_1) \ldots\gamma_{e_r}(x_r) =
\sum c_1^{d_1} \ldots c_s^{d_s} \gamma_{d_1}(y_1) \ldots\gamma_{d_s}(y_s)
$$
dans $B$, où, au membre de droite, la somme porte sur les
$d_1 + \ldots + d_s = n$. En effet, les axiomes d'une structure de puissances
divisées permettent de développer les deux membres en sommes, à coefficients
dans $\mathbf{Z}[a_{ij}]$, de termes de la forme
$c_1^{d_1} \ldots c_s^{d_s}\gamma_{e_1}(x_1) \ldots\gamma_{e_r}(x_r)$.
Pour vérifier l'égalité des coefficients, il suffit de remarquer que le résultat est vrai
dans $\mathbf{Q}[x_1, \ldots, x_r, c_1, \ldots, c_s, a_{ij}]$, muni de l'unique structure de puissances divisées
$\gamma$ sur $(x_1, \ldots, x_r)$.
D'après le théorème de Lazard (Algèbre, théorème \ref{algebra-theorem-lazard}),
$B$ est une colimite filtrante de $A$-modules libres
de type fini. En particulier, si $z \in IB$ s'écrit à la fois $z = \sum x_ib_i$ et
$z = \sum x'_{i'}b'_{i'}$, alors nous pouvons trouver $c_1, \ldots, c_s \in B$
et $a_{ij}, a'_{i'j} \in A$ tels que $b_i = \sum a_{ij}c_j$
et $b'_{i'} = \sum a'_{i'j}c_j$, et tels que
$y_j = \sum x_ia_{ij} = \sum x'_{i'}a'_{i'j}$\footnote{On
peut également le démontrer sans
recourir au théorème \ref{algebra-theorem-lazard} d'Algèbre. En effet, si
$z = \sum x_ib_i$ et $z = \sum x'_{i'}b'_{i'}$, alors
$\sum x_ib_i - \sum x'_{i'}b'_{i'} = 0$ est une relation dans le
$A$-module $B$. Le lemme \ref{algebra-lemma-flat-eq}
d'Algèbre (appliqué aux $x_i$ et $x'_{i'}$ à la place des $f_i$,
et aux $b_i$ et $b'_{i'}$ à la place des $x_i$)
fournit alors les éléments $c_1, \ldots, c_s \in B$
et $a_{ij}, a'_{i'j} \in A$ voulus.}. La
procédure ci-dessus définit donc une application $\bar\gamma_n$
sur $IB$, indépendante de l'écriture choisie. Par construction, $\bar\gamma$ satisfait les conditions (1), (3)
et (4). De plus, pour $x \in I$, on a $\bar\gamma_n(x) = \gamma_n(x)$.
Le lemme \ref{dpa-lemma-check-on-generators} montre alors que $\bar\gamma$
est une structure de puissances divisées sur $IB$.
\end{proof}

\begin{lemma}
\label{dpa-lemma-kernel}
Soit $(A, I, \gamma)$ un anneau à puissances divisées.
\begin{enumerate}
\item Si $\varphi : (A, I, \gamma) \to (B, J, \delta)$ est
un homomorphisme d'anneaux à puissances divisées, alors $\Ker(\varphi) \cap I$
est préservé par $\gamma_n$ pour tous $n \geq 1$.
\item Soit $\mathfrak a \subset A$ un idéal et posons
$I' = I \cap\mathfrak a$. Les conditions suivantes sont équivalentes :
\begin{enumerate}
\item $I'$ est préservé par $\gamma_n$ pour tous $n > 0$,
\item $\gamma$ se prolonge à $A/\mathfrak a$,
\item il existe une famille de générateurs $x_i$ de l'idéal $I'$ telle
que $\gamma_n(x_i) \in I'$ pour tout $n > 0$.
\end{enumerate}
\end{enumerate}
\end{lemma}

\begin{proof}
Démonstration de (1). C'est immédiat. Supposons (2)(a). Posons
$\bar\gamma_n(x \bmod I') = \gamma_n(x) \bmod I'$ pour $x \in I$.
Cette définition est indépendante du représentant, car on a $\gamma_n(x + y) = \gamma_n(x) \bmod I'$
pour $y \in I'$ d'après la définition \ref{dpa-definition-divided-powers}, (4),
et l'hypothèse $\gamma_j(y) \in I'$. Il est clair que
$\bar\gamma$ est une structure de puissances divisées puisque $\gamma$ en est
une. Ainsi, (2)(b) est vérifiée. Réciproquement, (2)(b) implique (2)(a)
d'après (1). Il est clair que (2)(a) implique (2)(c). Supposons (2)(c).
Notons que $\gamma_n(x) = a^n\gamma_n(x_i) \in I'$ lorsque $x = ax_i$.
Ainsi, $\gamma_n(x) \in I'$ pour une famille de générateurs de $I'$
comme groupe abélien. Par récurrence sur la longueur d'une expression en
ces générateurs, il suffit de montrer que $\forall n : \gamma_n(x + y) \in I'$
dès que $\forall n : \gamma_n(x), \gamma_n(y) \in I'$. Cela
résulte immédiatement du quatrième axiome d'une structure de puissances divisées.
\end{proof}

\begin{lemma}
\label{dpa-lemma-sub-dp-ideal}
Soit $(A, I, \gamma)$ un anneau à puissances divisées et
soit $E \subset I$ une
partie. Le plus petit idéal $J \subset I$ stable par $\gamma$
et contenant tous les $f \in E$ est l'idéal $J$ engendré par les
$\gamma_n(f)$, pour $n \geq 1$ et $f \in E$.
\end{lemma}

\begin{proof}
Cela résulte immédiatement du lemme \ref{dpa-lemma-kernel}.
\end{proof}

\begin{lemma}
\label{dpa-lemma-extend-to-completion}
Soit $(A, I, \gamma)$ un anneau à puissances divisées. Soit $p$ un nombre
premier. Si $p$ est nilpotent dans $A/I$, alors
\begin{enumerate}
\item l'application canonique du complété $p$-adique $A^\wedge = \lim_e A/p^eA$ vers $A/I$ est surjective ;
\item le noyau de cette application est le complété $p$-adique $I^\wedge$ de $I$;
\item chaque application $\gamma_n$ est continue pour la topologie $p$-adique et se prolonge en une
application $\gamma_n^\wedge : I^\wedge\to I^\wedge$, les applications ainsi obtenues définissant une
structure de puissances divisées sur $I^\wedge$.
\end{enumerate}
Si, de plus, $A$ est une $\mathbf{Z}_{(p)}$-algèbre, alors
\begin{enumerate}
\item[(4)] pour $e$ assez grand, l'idéal $p^eA \subset I$ est stable par la
structure de puissances divisées $\gamma$, et
$$
(A^\wedge, I^\wedge, \gamma^\wedge) = \lim_e (A/p^eA, I/p^eA, \bar\gamma)
$$
dans la catégorie des anneaux à puissances divisées.
\end{enumerate}
\end{lemma}

\begin{proof}
Soit $t \geq 1$ un entier tel que $p^tA/I = 0$, c'est-à-dire $p^tA \subset I$.
L'application $A^\wedge\to A/I$ est la composée $A^\wedge\to A/p^tA \to A/I$,
qui est surjective (par exemple
d'après Algèbre, lemme \ref{algebra-lemma-completion-generalities}).
Comme $p^eI \subset p^eA \cap I \subset p^{e - t}I$ pour $e \geq t$,
le noyau de la composée $A^\wedge\to A/I$ est le complété $p$-adique
de $I$. L'application $\gamma_n$ est continue, car
$$
\gamma_n(x + p^ey) =
\sum\nolimits_{i + j = n} p^{je}\gamma_i(x)\gamma_j(y) =
\gamma_n(x) \bmod p^eI
$$
par les axiomes d'une structure de puissances
divisées. Il est clair que les applications $\gamma_n^\wedge$
satisfont aux mêmes axiomes que les applications $\gamma_n$. Enfin, pour démontrer la dernière assertion, prenons $e > t$.
Alors $p^eA \subset I$ et $\gamma_1(p^eA) \subset p^eA$ ; pour $n > 1$,
on a
$$
\gamma_n(p^ea) = p^n \gamma_n(p^{e - 1}a) = \frac{p^n}{n!} p^{n(e - 1)}a^n
\in p^e A
$$
car $p^n/n! \in\mathbf{Z}_{(p)}$ et, puisque $n \geq 2$ et $e \geq 2$,
$n(e - 1) \geq e$.
Cela prouve que $\gamma$ se prolonge à $A/p^eA$; voir le lemme \ref{dpa-lemma-kernel}.
L'assertion relative aux limites découle immédiatement de la construction des limites donnée dans
la démonstration du lemme \ref{dpa-lemma-limits}.
\end{proof}




\section{Algèbres polynomiales à puissances divisées}
\label{dpa-section-divided-power-polynomial-ring}

\noindent
Un exemple très utile est celui de l'{\it algèbre polynomiale à puissances divisées}.
Soit $A$ un anneau. Soit $t \geq 1$. On note
$A\langle x_1, \ldots, x_t \rangle$ la $A$-algèbre suivante :
en tant que $A$-module, on pose
$$
A\langle x_1, \ldots, x_t \rangle =
\bigoplus\nolimits_{n_1, \ldots, n_t \geq 0} A x_1^{[n_1]} \ldots x_t^{[n_t]}
$$
et la multiplication est donnée par
$$
x_i^{[n]}x_i^{[m]} = \frac{(n + m)!}{n!m!}x_i^{[n + m]}.
$$
On pose aussi $x_i = x_i^{[1]}$. Remarquons que
$1 = x_1^{[0]} \ldots x_t^{[0]}$. Une construction analogue donne
l'algèbre polynomiale à puissances divisées en une infinité de variables.
Il existe un morphisme canonique de $A$-algèbres
$A\langle x_1, \ldots, x_t \rangle\to A$ qui envoie $x_i^{[n]}$ sur zéro
pour $n > 0$. Le noyau de ce morphisme est noté
$A\langle x_1, \ldots, x_t \rangle_{+}$.

\begin{lemma}
\label{dpa-lemma-divided-power-polynomial-algebra}
Soit $(A, I, \gamma)$ un anneau à puissances
divisées. Il existe une unique structure de puissances divisées $\delta$ sur
$$
J = IA\langle x_1, \ldots, x_t \rangle + A\langle x_1, \ldots, x_t \rangle_{+}
$$
telle que
\begin{enumerate}
\item $\delta_n(x_i) = x_i^{[n]}$, et
\item $(A, I, \gamma) \to (A\langle x_1, \ldots, x_t \rangle, J, \delta)$
soit un homomorphisme d'anneaux à puissances divisées.
\end{enumerate}
De plus, $(A\langle x_1, \ldots, x_t \rangle, J, \delta)$ possède
la propriété universelle suivante : se donner un homomorphisme d'anneaux à puissances divisées
$\varphi : (A\langle x_1, \ldots, x_t \rangle, J, \delta) \to
(C, K, \epsilon)$ revient
à se donner un homomorphisme d'anneaux à puissances divisées
$A \to C$ et des éléments $k_1, \ldots, k_t \in K$.
\end{lemma}

\begin{proof}
Nous démontrons le lemme dans le cas d'une algèbre polynomiale à
puissances divisées en une variable. Le cas général se traite exactement de la même
manière ; on peut aussi remarquer que $A\langle x_1, \ldots, x_t\rangle$ est
isomorphe à l'anneau obtenu en adjoignant successivement les variables à puissances divisées
$x_1, \ldots, x_t$.

\medskip\noindent
Soit $A\langle x \rangle_{+}$ l'idéal engendré par
$x, x^{[2]}, x^{[3]}, \ldots$.
Remarquons que $J = IA\langle x \rangle + A\langle x \rangle_{+}$
et que
$$
IA\langle x \rangle\cap A\langle x \rangle_{+} =
IA\langle x \rangle\cdot A\langle x \rangle_{+}
$$
Par conséquent, d'après le lemme \ref{dpa-lemma-two-ideals}, il suffit de montrer qu'il existe
des structures de puissances divisées sur les idéaux $IA\langle x \rangle$ et
$A\langle x \rangle_{+}$. L'existence de la première résulte du
lemme \ref{dpa-lemma-gamma-extends}, puisque $A \to A\langle x \rangle$ est
plat. Pour la seconde, si $A$ est sans torsion, le
lemme \ref{dpa-lemma-silly} (4) montre que $\delta$ existe. En effet, en
choisissant les éléments $x^{[m]}$ comme générateurs, on constate que
$(x^{[m]})^n = \frac{(nm)!}{(m!)^n} x^{[nm]}$
et que $n!$ divise l'entier $\frac{(nm)!}{(m!)^n}$.
Dans le cas général, écrivons $A = R/\mathfrak a$, où $R$
est un anneau sans torsion (par exemple un anneau de polynômes sur $\mathbf{Z}$). Le noyau de
$R\langle x \rangle\to A\langle x \rangle$ est
$\bigoplus\mathfrak a x^{[m]}$. En appliquant le critère (2)(c) du
lemme \ref{dpa-lemma-kernel}, on voit que la structure de puissances divisées
sur $R\langle x \rangle_{+}$ se transporte à $A\langle x \rangle$,
comme souhaité.

\medskip\noindent
Démontrons la propriété universelle. Étant donnés un homomorphisme $\varphi : A \to C$
d'anneaux à puissances divisées et des éléments $k_1, \ldots, k_t \in K$, considérons
$$
A\langle x_1, \ldots, x_t \rangle\to C,\quad
x_1^{[n_1]} \ldots x_t^{[n_t]} \longmapsto
\epsilon_{n_1}(k_1) \ldots\epsilon_{n_t}(k_t)
$$
où l'on utilise $\varphi$ sur les coefficients. Il reste seulement à vérifier qu'il s'agit d'un
homomorphisme de $A$-algèbres (détails omis). La construction
inverse est claire.
\end{proof}

\begin{remark}
\label{dpa-remark-divided-power-polynomial-algebra}
Soit $(A, I, \gamma)$ un anneau à puissances
divisées. Il existe une variante du lemme \ref{dpa-lemma-divided-power-polynomial-algebra}
pour une infinité de variables. Remarquons d'abord que, si $s < t$, il existe
un morphisme canonique
$$
A\langle x_1, \ldots, x_s \rangle\to A\langle x_1, \ldots, x_t\rangle
$$
Ainsi, si $W$ est un ensemble quelconque, on pose
$$
A\langle x_w: w \in W\rangle =
\colim_{E \subset W} A\langle x_e:e \in E\rangle
$$
(la colimite étant prise sur les sous-ensembles finis $E$ de $W$),
avec les morphismes de transition ci-dessus. La définition d'une
colimite montre que la propriété universelle de $A\langle x_w: w \in W\rangle$ est
entièrement analogue à celle énoncée dans le
lemme \ref{dpa-lemma-divided-power-polynomial-algebra}.
\end{remark}

\noindent
Le lemme suivant se trouve dans \cite{BO}.

\begin{lemma}
\label{dpa-lemma-need-only-gamma-p}
Soit $p$ un nombre premier. Soit $A$ un anneau dans lequel tout entier $n$
non divisible par $p$ est inversible, autrement dit $A$ est une $\mathbf{Z}_{(p)}$-algèbre.
Soit $I \subset A$ un idéal. Deux structures de puissances divisées
$\gamma, \gamma'$ sur $I$ sont égales si et seulement si $\gamma_p = \gamma'_p$.
De plus, étant donnée une application $\delta : I \to I$ telle que
\begin{enumerate}
\item $p!\delta(x) = x^p$ pour tout $x \in I$\footnote{Cette condition
résulte des deux autres ; voir la remarque \ref{dpa-remark-ryo-suzuki}.},
\item $\delta(ax) = a^p\delta(x)$ pour tous $a \in A$, $x \in I$, et
\item
$\delta(x + y) =
\delta(x) +
\sum\nolimits_{i + j = p, i,j \geq 1} \frac{1}{i!j!} x^i y^j +
\delta(y) $ pour tous $ x, y \in I$,
\end{enumerate}
il existe une unique structure de puissances divisées $\gamma$ sur $I$ telle
que $\gamma_p = \delta$.
\end{lemma}

\begin{proof}
Si $n$ n'est pas divisible par $p$, alors $\gamma_n(x) = c x \gamma_{n - 1}(x)$,
où $c$ est une unité de $\mathbf{Z}_{(p)}$. De plus,
$$
\gamma_{pm}(x) = c \gamma_m(\gamma_p(x))
$$
où $c$ est une unité de $\mathbf{Z}_{(p)}$. La première assertion est donc
claire. Pour la seconde, on peut remonter ces égalités afin de
définir tous les $\gamma_n$. Plus précisément, si
$n = a_0 + a_1p + \ldots + a_e p^e$ avec $a_i \in\{0, \ldots, p - 1\}$,
on pose
$$
\gamma_n(x) = c_n x^{a_0} \delta(x)^{a_1} \ldots\delta^e(x)^{a_e}
$$
où $c_n \in\mathbf{Z}_{(p)}$ est défini par
$$
c_n =
{(p!)^{a_1 + a_2(1 + p) + \ldots + a_e(1 + \ldots + p^{e - 1})}}/{n!}.
$$
Il faut maintenant vérifier les axiomes (1) -- (5) d'une structure de puissances divisées ; voir
la définition \ref{dpa-definition-divided-powers}. Les axiomes (1) et (3) sont
immédiats. Les axiomes (2) et (5) se vérifient par un calcul direct que nous
omettons. Soient $x, y \in I$. Nous affirmons qu'il existe un homomorphisme d'anneaux
$$
\varphi : \mathbf{Z}_{(p)}\langle u, v \rangle\longrightarrow A
$$
qui envoie $u^{[n]}$ sur $\gamma_n(x)$ et $v^{[n]}$ sur $\gamma_n(y)$.
D'après la construction de $\mathbf{Z}_{(p)}\langle u, v \rangle$, il faut
pour cela vérifier que
$$
\gamma_n(x)\gamma_m(x) = \frac{(n + m)!}{n!m!} \gamma_{n + m}(x)
$$
dans $A$, et de même pour $y$. C'est vrai, puisque l'axiome (2) est vérifié par $\gamma$.
Notons $\epsilon$ la structure de puissances divisées sur
l'idéal $\mathbf{Z}_{(p)}\langle u, v\rangle_{+}$ de
$\mathbf{Z}_{(p)}\langle u, v\rangle$.
Nous affirmons ensuite que $\varphi(\epsilon_n(f)) = \gamma_n(\varphi(f))$
pour $f \in\mathbf{Z}_{(p)}\langle u, v\rangle_{+}$ et tout $n$.
C'est clair pour $n = 0, 1, \ldots, p - 1$. Pour $n = p$, il suffit
de le vérifier sur un système de générateurs de l'idéal
$\mathbf{Z}_{(p)}\langle u, v\rangle_{+}$, car $\epsilon_p$
et $\gamma_p = \delta$ vérifient toutes deux les propriétés (1) et (3)
du lemme. Il suffit donc d'établir
$\gamma_p(\gamma_n(x)) = \frac{(pn)!}{p!(n!)^p}\gamma_{pn}(x)$,
et de même pour $y$; cela résulte de l'axiome (5) pour $\gamma$.
Maintenant, si $n = a_0 + a_1p + \ldots + a_e p^e$
est un entier quelconque écrit sous la forme de son développement $p$-adique ci-dessus, alors
$$
\epsilon_n(f) =
c_n f^{a_0} \gamma_p(f)^{a_1} \ldots\gamma_p^e(f)^{a_e}
$$
car $\epsilon$ est une structure de puissances divisées. On en déduit que
$\varphi(\epsilon_n(f)) = \gamma_n(\varphi(f))$ pour tout $n$.
En appliquant ceci à $f = u + v$, on voit que l'axiome (4) pour $\gamma$
découle du fait que $\epsilon$ est une structure de puissances divisées.
\end{proof}

\begin{remark}
\label{dpa-remark-ryo-suzuki} Dans la
situation $A \supset I$ du lemme \ref{dpa-lemma-need-only-gamma-p},
soit $\delta : I \to I$ une application vérifiant (2) et (3) du lemme. Alors (1)
est également vérifiée.
Montrons d'abord que $\delta(nx) = n\delta(x) + \frac{n^p - n}{p!}x^p$
pour $n \geq 1$, par récurrence sur $n$. Le cas $n = 1$ est
clair. Supposons le résultat établi pour un certain $n$. Alors
$\delta((n + 1)x) = \delta(nx) + \delta(x) +
(\sum_{i = 1}^{p - 1} \frac{n^i}{i!(p - i)!})x^p
= (n + 1)\delta(x) + \left( \frac{n^p - n}{p!} + \frac{(n + 1)^p}{p!}
- \frac{n^p + 1}{p!} \right)x^p =
(n + 1)\delta(x) + \frac{(n + 1)^p - (n + 1)}{p!}x^p$,
ce qui établit le résultat pour $n + 1$. D'autre part,
$\delta(nx) = n^p\delta(x)$, donc
$(n^p - n)\delta(x) = \frac{n^p - n}{p!}x^p$.
Or, pour tout $p$, il existe un entier $n$ tel que $p^2 \not\mid n^p - n$.
En effet,
$(\mathbf{Z}/p^2\mathbf{Z})^\times\cong\mathbf{Z}/p(p-1)\mathbf{Z}$.
Il existe donc $n$ tel que $n^{p-1} \not\equiv 1 \mod p^2$.
On obtient ainsi $p!\delta(x)=x^p$.
\end{remark}









\section{Résolutions de Tate}
\label{dpa-section-tate}

\noindent
Dans cette section, nous présentons brièvement les résolutions construites dans
\cite{Tate-homology} et \cite{AH},
qui combinent les structures de puissances
divisées et les algèbres différentielles
graduées. Nous utiliserons la {\it notation homologique} pour
les algèbres différentielles
graduées. Celles-ci seront concentrées en degrés homologiques positifs ou
nuls. Ainsi, nos algèbres différentielles graduées $(A, \text{d})$
seront des complexes de chaînes
$$
\ldots\to A_2 \to A_1 \to A_0 \to 0 \to\ldots
$$
munis d'une multiplication.

\medskip\noindent
Soit $R$ un anneau (commutatif, comme
toujours). Dans cette section, nous considérerons souvent des
$R$-algèbres graduées $A = \bigoplus_{d \geq 0} A_d$ dont les composantes
de degré négatif sont nulles. On pose $A_+ = \bigoplus_{d > 0} A_d$.
On notera $A_{even} = \bigoplus_{d \geq 0} A_{2d}$ et
$A_{odd} = \bigoplus_{d \geq 0} A_{2d + 1}$.
Rappelons que $A$ est commutative au sens gradué si
$x y = (-1)^{\deg(x)\deg(y)} y x$ pour tous éléments homogènes $x, y$.
Rappelons que $A$ est strictement commutative au sens gradué si, de plus,
$x^2 = 0$ pour tout élément homogène $x$ de degré impair. Enfin, pour comprendre
la définition suivante, gardons à l'esprit que $\gamma_n(x) = x^n/n!$
lorsque $A$ est une $\mathbf{Q}$-algèbre.

\begin{definition}
\label{dpa-definition-divided-powers-graded}
Soit $R$ un anneau. Soit $A = \bigoplus_{d \geq 0} A_d$ une
$R$-algèbre graduée strictement commutative au sens gradué. Une famille d'applications
$\gamma_n : A_{even, +} \to A_{even, +}$, définies pour tout $n > 0$, est appelée
une {\it structure de puissances divisées} sur $A$ si les conditions suivantes sont satisfaites :
\begin{enumerate}
\item $\gamma_n(x) \in A_{2nd}$ si $x \in A_{2d}$,
\item $\gamma_1(x) = x$ pour tout $x$; on pose aussi $\gamma_0(x) = 1$,
\item $\gamma_n(x)\gamma_m(x) = \frac{(n + m)!}{n! m!} \gamma_{n + m}(x)$,
\item $\gamma_n(xy) = x^n \gamma_n(y)$ pour tous $x \in A_{even}$ et
$y \in A_{even, +}$,
\item $\gamma_n(xy) = 0$ si $x, y \in A_{odd}$ sont homogènes et $n > 1$,
\item si $x, y \in A_{even, +}$, alors
$\gamma_n(x + y) = \sum_{i = 0, \ldots, n} \gamma_i(x)\gamma_{n - i}(y)$,
\item $\gamma_n(\gamma_m(x)) =
\frac{(nm)!}{n! (m!)^n} \gamma_{nm}(x) $ pour $ x \in A_{even, +}$.
\end{enumerate}
\end{definition}

\noindent
Remarquons que les conditions (2), (3), (4), (6) et (7) impliquent que
$\gamma$ est une structure de puissances divisées « usuelle » sur l'idéal
$A_{even, +}$ de l'anneau commutatif $A_{even}$; voir
les sections \ref{dpa-section-divided-powers},
\ref{dpa-section-divided-power-rings},
\ref{dpa-section-extend} et
\ref{dpa-section-divided-power-polynomial-ring}.
En particulier, $n! \gamma_n(x) = x^n$ pour tout $x \in A_{even, +}$.
La condition (1) exprime la compatibilité de $\gamma$ avec la graduation, et la condition
(5) dit que $\gamma_n$, pour $n > 1$, s'annule sur les
produits d'éléments homogènes de degré impair. Il peut toutefois arriver
que
$$
\gamma_2(z_1 z_2 + z_3 z_4) = z_1z_2z_3z_4
$$
soit non nul lorsque $z_1, z_2, z_3, z_4$ sont des éléments homogènes de degré impair.

\begin{example}[Adjonction d'une variable impaire]
\label{dpa-example-adjoining-odd}
Soit $R$ un anneau. Soit $(A, \gamma)$ une
$R$-algèbre graduée strictement commutative au sens gradué, munie d'une structure de puissances divisées
comme dans la définition ci-dessus. Soit $d > 0$ un entier impair.
Dans ce cas, on peut adjoindre une variable $T$ de degré $d$ à $A$.
Plus précisément, on pose
$$
A\langle T \rangle = A \oplus AT
$$
avec la graduation donnée par $A\langle T \rangle_m = A_m \oplus A_{m - d}T$.
Il existe une unique structure de puissances divisées sur
$A\langle T \rangle$ compatible avec la structure
donnée sur $A$. Elle est définie par
$$
\gamma_n(x + yT) = \gamma_n(x) + \gamma_{n - 1}(x)yT
$$
pour $x \in A_{even, +}$ et $y \in A_{odd}$.
\end{example}

\begin{example}[Adjonction d'une variable paire]
\label{dpa-example-adjoining-even}
Soit $R$ un anneau. Soit $(A, \gamma)$ une
$R$-algèbre graduée strictement commutative au sens gradué, munie d'une structure de puissances divisées
comme dans la définition ci-dessus. Soit $d > 0$ un entier pair.
Dans ce cas, on peut adjoindre une variable $T$ de degré $d$ à $A$.
Plus précisément, on pose
$$
A\langle T \rangle = A \oplus AT \oplus AT^{(2)} \oplus AT^{(3)} \oplus\ldots
$$
avec la multiplication donnée par
$$
T^{(n)} T^{(m)} = \frac{(n + m)!}{n!m!} T^{(n + m)}
$$
et la graduation donnée par
$$
A\langle T \rangle_m =
A_m \oplus A_{m - d}T \oplus A_{m - 2d}T^{(2)} \oplus\ldots
$$
Il existe une unique structure de puissances divisées sur
$A\langle T \rangle$, compatible avec celle
donnée sur $A$, telle que $\gamma_n(T^{(i)}) = T^{(ni)}$.
Pour la définir, on pose d'abord
$$
\gamma_n\left(\sum\nolimits_{i > 0} x_i T^{(i)}\right) =
\sum\prod\nolimits_{n = \sum e_i} x_i^{e_i} T^{(ie_i)}
$$
si $x_i$ appartient à $A_{even}$. Si $x_0 \in A_{even, +}$,
on pose ensuite
$$
\gamma_n\left(\sum\nolimits_{i \geq 0} x_i T^{(i)}\right) =
\sum\nolimits_{a + b = n}
\gamma_a(x_0)\gamma_b\left(\sum\nolimits_{i > 0} x_iT^{(i)}\right)
$$
où $\gamma_b$ est défini comme ci-dessus.
\end{example}

\begin{remark}
\label{dpa-remark-adjoining-set-of-variables} On
peut aussi adjoindre un ensemble, éventuellement infini, de générateurs extérieurs ou
à puissances divisées dans un degré donné $d > 0$, plutôt qu'un seul comme
dans les exemples \ref{dpa-example-adjoining-odd}
et \ref{dpa-example-adjoining-even}. Plus
précisément, suivant la remarque \ref{dpa-remark-divided-power-polynomial-algebra},
pour $(A,\gamma)$
comme ci-dessus et un ensemble $J$, soit $A\langle
T_j:j\in J\rangle$ la colimite filtrante des algèbres
$A\langle T_j:j\in S\rangle$, où $S$
parcourt les sous-ensembles finis de $J$. Il est immédiat que cette algèbre possède une unique structure de
puissances divisées compatible avec la structure donnée sur $A$ et avec
celle de chaque générateur $T_j$.
\end{remark}

\noindent
Nous relions maintenant la notion de structure de puissances divisées à
celle de différentielle. Pour comprendre la définition, remarquons que
$\text{d}(x^n/n!) = \text{d}(x) x^{n - 1}/(n - 1)!$ si $A$
est une $\mathbf{Q}$-algèbre et si $x \in A_{even, +}$.

\begin{definition}
\label{dpa-definition-divided-powers-dga}
Soit $R$ un anneau. Soit $A = \bigoplus_{d \geq 0} A_d$ une
$R$-algèbre différentielle graduée strictement commutative au sens
gradué. Une structure de puissances divisées $\gamma$ sur $A$ est {\it compatible avec
la structure différentielle graduée} si
$\text{d}(\gamma_n(x)) = \text{d}(x) \gamma_{n - 1}(x)$ pour
tout $x \in A_{even, +}$.
\end{definition}

\noindent
Attention : soit $(A, \text{d}, \gamma)$ comme dans
la définition \ref{dpa-definition-divided-powers-dga}.
Il se peut que $\gamma_n(x)$ ne soit pas un bord, même lorsque
$x$ est un bord. Ainsi, en général, $\gamma$ n'induit pas de structure
de puissances divisées sur l'algèbre d'homologie $H(A)$.
Dans certains articles, les auteurs imposent une condition de
compatibilité supplémentaire afin que ce soit le cas ; nous choisissons de ne
pas le faire.

\begin{lemma}
\label{dpa-lemma-dpdga-good}
Soient $(A, \text{d}, \gamma)$ et $(B, \text{d}, \gamma)$ comme dans
la définition \ref{dpa-definition-divided-powers-dga}. Soit $f : A \to B$
un morphisme d'algèbres différentielles graduées compatible avec les structures de
puissances divisées. Supposons que
\begin{enumerate}
\item $H_k(A) = 0$ pour $k > 0$, et
\item $f$ soit surjectif.
\end{enumerate}
Alors $\gamma$ induit une structure de puissances divisées sur la
$R$-algèbre graduée $H(B)$.
\end{lemma}

\begin{proof}
Supposons que $x$ et $x'$ soient homogènes de même degré $2d$
et représentent la même classe d'homologie dans $H(B)$. Écrivons $x' - x = \text{d}(w)$.
Choisissons un relèvement $y \in A_{2d}$ de $x$ et un relèvement $z \in A_{2d + 1}$
de $w$. Alors $y' = y + \text{d}(z)$ est un relèvement de $x'$.
Par conséquent,
$$
\gamma_n(y') = \sum\gamma_i(y) \gamma_{n - i}(\text{d}(z))
= \gamma_n(y) +
\sum\nolimits_{i < n} \gamma_i(y) \gamma_{n - i}(\text{d}(z))
$$
Comme $A$ est acyclique en degrés positifs et que
$\text{d}(\gamma_j(\text{d}(z))) = 0$ pour tout $j$, on peut écrire cette
égalité sous la forme
$$
\gamma_n(y') = \gamma_n(y) +
\sum\nolimits_{i < n} \gamma_i(y) \text{d}(z_i)
$$
pour certains $z_i$ dans $A$. De plus, pour $0 < i < n$, on a
$$
\text{d}(\gamma_i(y) z_i) =
\text{d}(\gamma_i(y))z_i + \gamma_i(y)\text{d}(z_i) =
\text{d}(y) \gamma_{i - 1}(y) z_i + \gamma_i(y)\text{d}(z_i)
$$
et le premier terme a une image nulle dans $B$, car $\text{d}(y)$ a une image nulle dans $B$.
Ainsi, $\gamma_n(x')$ et $\gamma_n(x)$ ont la même image dans $H(B)$.
On obtient donc une application bien définie $\gamma_n : H_{2d}(B) \to H_{2nd}(B)$
pour tous $d > 0$ et $n > 0$. Nous omettons la vérification du fait qu'elle
définit une structure de puissances divisées sur $H(B)$.
\end{proof}

\begin{lemma}
\label{dpa-lemma-base-change-div}
Soit $(A, \text{d}, \gamma)$ comme dans
la définition \ref{dpa-definition-divided-powers-dga}.
Soit $R \to R'$ un homomorphisme
d'anneaux. Alors $\text{d}$ et $\gamma$ induisent des structures analogues sur
$A' = A \otimes_R R'$ telles que $(A', \text{d}, \gamma)$ soit comme
dans la définition \ref{dpa-definition-divided-powers-dga}.
\end{lemma}

\begin{proof}
Remarquons que $A'_{even} = A_{even} \otimes_R R'$ et
$A'_{even, +} = A_{even, +} \otimes_R R'$. Il s'agit donc de
montrer que les puissances divisées $\gamma$ se prolongent à $A'_{even}$
(au sens de la définition \ref{dpa-definition-extends}).
Une fois établi que $\gamma$ se prolonge, il est facile de voir que ce
prolongement possède toutes les propriétés voulues.

\medskip\noindent
Choisissons une $R$-algèbre polynomiale $P$ sur un ensemble quelconque de
générateurs, ainsi qu'une surjection de $R$-algèbres
$P \to R'$. L'homomorphisme d'anneaux $A_{even} \to A_{even} \otimes_R P$ est plat
; les puissances divisées $\gamma$ se prolongent donc de manière unique à $A_{even} \otimes_R P$
par le lemme \ref{dpa-lemma-gamma-extends}.
Soit $J = \Ker(P \to R')$. Pour montrer que $\gamma$ se
prolonge à $A \otimes_R R'$, il suffit de montrer que
$I' = \Ker(A_{even, +} \otimes_R P \to A_{even, +} \otimes_R R')$
est engendré par des éléments $z$ tels que $\gamma_n(z) \in I'$
pour tout $n > 0$. C'est clair, car $I'$ est engendré par les
éléments de la forme $x \otimes f$, avec
$x \in A_{even, +}$ et $f \in\Ker(P \to R')$.
\end{proof}

\begin{lemma}
\label{dpa-lemma-extend-differential}
Soit $(A, \text{d}, \gamma)$ comme dans
la définition \ref{dpa-definition-divided-powers-dga}.
Soit $d \geq 1$ un entier.
Soit $A\langle T \rangle$ l'algèbre polynomiale graduée à puissances divisées
en $T$, avec $\deg(T) = d$,
construite dans l'exemple \ref{dpa-example-adjoining-odd} ou
\ref{dpa-example-adjoining-even}.
Soit $f \in A_{d - 1}$ un élément tel que $\text{d}(f) = 0$.
Il existe une unique différentielle $\text{d}$
sur $A\langle T\rangle$ telle que $\text{d}(T) = f$ et
telle que $\text{d}$ soit compatible avec la structure de
puissances divisées sur $A\langle T \rangle$.
\end{lemma}

\begin{proof}
Cela résulte d'un calcul direct que nous omettons.
\end{proof}

\noindent
Dans le lemme \ref{dpa-lemma-compute-cohomology-adjoin-variable},
nous calculerons l'homologie de $A\langle T \rangle$ dans certains cas
particuliers. Voici la construction de Tate, telle qu'elle a été
étendue par Avramov et Halperin.

\begin{lemma}
\label{dpa-lemma-tate-resolution}
Soit $R \to S$ un homomorphisme d'anneaux commutatifs.
Il existe une factorisation
$$
R \to A \to S
$$
ayant les propriétés suivantes :
\begin{enumerate}
\item $(A, \text{d}, \gamma)$ est comme
dans la définition \ref{dpa-definition-divided-powers-dga},
\item $A \to S$ est un quasi-isomorphisme (si l'on munit $S$ de
la différentielle nulle),
\item $A_0 = R[x_j: j\in J] \to S$ est une surjection quelconque d'un anneau de
polynômes vers $S$, et
\item $A$ est une algèbre polynomiale graduée à puissances divisées sur $R$.
\end{enumerate}
La dernière condition signifie que $A$ est construite à partir de $A_0$ en
adjoignant successivement un ensemble de variables $T$ dans chaque degré $> 0$, comme
dans l'exemple \ref{dpa-example-adjoining-odd} ou \ref{dpa-example-adjoining-even}.
De plus, si $R$ est noethérien et si $R\to S$ est de type fini, on peut
choisir $A$ avec seulement un nombre fini de générateurs dans
chaque degré.
\end{lemma}

\begin{proof}
Détaillons la construction lorsque $R$ est
noethérien et $R\to S$ est de type fini. Sans ces hypothèses, la démonstration
est la même, sauf qu'il faut utiliser dans chaque degré un
ensemble de générateurs éventuellement infini.

\medskip\noindent
Début de la construction : soit $A(0) = R[x_1, \ldots, x_n]$ un
anneau de polynômes usuel, et soit $A(0) \to S$ une surjection. Pour
la graduation, on pose $A(0)_0 = A(0)$ et $A(0)_d = 0$ pour $d \not = 0$.
Ainsi, $\text{d} = 0$ et $\gamma_n$ est également nulle pour $n > 0$.

\medskip\noindent
Choisissons des générateurs $f_1, \ldots, f_m \in R[x_1, \ldots, x_n]$
du noyau de l'homomorphisme donné $A(0) = R[x_1, \ldots, x_n] \to S$.
En appliquant l'exemple \ref{dpa-example-adjoining-odd} $m$ fois, on obtient
$$
A(1) = A(0)\langle T_1, \ldots, T_m\rangle
$$
avec $\deg(T_i) = 1$, en tant qu'algèbre polynomiale graduée à puissances
divisées. Posons $\text{d}(T_i) = f_i$. Comme $A(1)$ est une algèbre polynomiale à
puissances divisées sur $A(0)$ et que $\text{d}(f_i) = 0$,
ceci se prolonge de manière unique en une différentielle sur $A(1)$,
d'après le lemme \ref{dpa-lemma-extend-differential}.

\medskip\noindent
Hypothèse de récurrence : supposons données des factorisations
$$
R \to A(0) \to A(1) \to\ldots\to A(m) \to S
$$
où $A(0)$ et $A(1)$ sont comme ci-dessus, et telles que chaque $R \to A(m') \to S$,
pour $2 \leq m' \leq m$, vérifie les propriétés (1) et (4)
de l'énoncé du lemme, la propriété (2) étant remplacée par la condition que
$H_i(A(m')) \to H_i(S)$ est un isomorphisme pour
$m' > i \geq 0$. Le cas initial est $m = 1$.

\medskip\noindent
Étape de récurrence : supposons donné $R \to A(m) \to S$
comme dans l'hypothèse de récurrence. Considérons le
groupe $H_m(A(m))$. C'est un module sur $H_0(A(m)) = S$.
En fait, c'est un sous-quotient de $A(m)_m$, qui est un
module de type fini sur $A(m)_0 = R[x_1, \ldots, x_n]$.
On peut donc choisir un nombre fini d'éléments
$$
e_1, \ldots, e_t \in\Ker(\text{d} : A(m)_m \to A(m)_{m - 1})
$$
dont les images engendrent ce module. En appliquant
l'exemple \ref{dpa-example-adjoining-odd} ou
\ref{dpa-example-adjoining-even} $t$ fois, on obtient
$$
A(m + 1) = A(m)\langle T_1, \ldots, T_t\rangle
$$
avec $\deg(T_i) = m + 1$, en tant qu'algèbre graduée à puissances divisées. Posons
$\text{d}(T_i) = e_i$. Comme $A(m+1)$ est une algèbre polynomiale à
puissances divisées sur $A(m)$ et que $\text{d}(e_i) = 0$,
ceci se prolonge de manière unique en une différentielle sur $A(m + 1)$
compatible avec la structure de puissances divisées.
Puisque nous n'avons ajouté des éléments qu'en degré $m + 1$ et au-dessus,
on a $H_i(A(m + 1)) = H_i(A(m))$ pour $i < m$. De
plus, $H_m(A(m + 1)) = 0$ par construction.

\medskip\noindent
Pour achever la démonstration, remarquons que nous avons construit
une suite de morphismes
$$
R \to A(0) \to A(1) \to\ldots\to A(m) \to A(m + 1) \to\ldots\to S
$$
et posons $A = \colim A(m)$.
\end{proof}

\begin{lemma}
\label{dpa-lemma-tate-resoluton-pseudo-coherent-ring-map}
Soit $R \to S$ un homomorphisme pseudo-cohérent d'anneaux (Compléments sur l'algèbre, définition
\ref{more-algebra-definition-pseudo-coherent-perfect}). Alors le
lemme \ref{dpa-lemma-tate-resolution} vaut avec une résolution $A$ de $S$
possédant un nombre fini de générateurs dans chaque degré.
\end{lemma}

\begin{proof}
La démonstration est exactement la même que celle du lemme \ref{dpa-lemma-tate-resolution}. Le
seul point supplémentaire est que, étant donné $A(m) \to S$, il faut
montrer que $H_m = H_m(A(m))$ est un $R[x_1, \ldots, x_m]$-module de
type fini (de sorte qu'à l'étape suivante il ne soit nécessaire d'ajouter qu'un nombre fini de
variables). Considérons le complexe
$$
\ldots\to A(m)_{m - 1} \to A(m)_m \to A(m)_{m - 1} \to
\ldots\to A(m)_0 \to S \to 0
$$
Comme $S$ est un $R[x_1, \ldots, x_n]$-module
pseudo-cohérent et que $A(m)_i$ est un $R[x_1, \ldots, x_n]$-module
libre de type fini, ce complexe est pseudo-cohérent ; voir
Compléments sur l'algèbre, lemme \ref{more-algebra-lemma-complex-pseudo-coherent-modules}.
Comme le complexe est exact en degrés homologiques $> m$,
on conclut que $H_m$ est un $R$-module de type
fini, d'après Compléments sur l'algèbre, lemme \ref{more-algebra-lemma-finite-cohomology}.
\end{proof}

\begin{lemma}
\label{dpa-lemma-uniqueness-tate-resolution}
Soit $R$ un anneau commutatif. Supposons que $(A, \text{d}, \gamma)$ et
$(B, \text{d}, \gamma)$ soient comme
dans la définition \ref{dpa-definition-divided-powers-dga}.
Soit $\overline{\varphi} : H_0(A) \to H_0(B)$ un homomorphisme de $R$-algèbres.
Supposons que
\begin{enumerate}
\item $A$ soit une algèbre polynomiale graduée à puissances divisées sur $R$.
\item $H_k(B) = 0$ pour $k > 0$.
\end{enumerate}
Il existe alors un morphisme $\varphi : A \to B$ de
$R$-algèbres différentielles graduées, compatible avec les puissances
divisées, qui relève $\overline{\varphi}$.
\end{lemma}

\begin{proof}
L'hypothèse signifie que $A$ s'obtient à partir de $R$ en adjoignant successivement
un ensemble de générateurs polynomiaux en degré zéro, des générateurs extérieurs en
degrés impairs positifs et des générateurs à puissances divisées
en degrés pairs positifs. On dispose donc d'une filtration
$R \subset A(0) \subset A(1) \subset\ldots$
de $A$ telle que $A(m + 1)$ s'obtienne à partir de $A(m)$ en adjoignant
des générateurs du type approprié (que nous appellerons simplement «
générateurs à puissances divisées ») en degré $m + 1$.
En particulier, $A(0) \to H_0(A)$ est un morphisme surjectif d'une algèbre
polynomiale usuelle sur $R$ vers $H_0(A)$. On peut donc relever $\overline{\varphi}$
en un homomorphisme de $R$-algèbres $\varphi(0) : A(0) \to B_0$.

\medskip\noindent
Écrivons $A(1) = A(0)\langle T_j:j\in J\rangle$ pour un
ensemble $J$ de variables à puissances divisées $T_j$ de degré $1$. Soit $f_j \in B_0$
défini par $f_j = \varphi(0)(\text{d}(T_j))$. Remarquons que l'image de $f_j$
dans $H_0(B)$ est nulle, puisque l'image de $\text{d}T_j$ dans $H_0(A)$ est
nulle. Il existe donc $b_j \in B_1$ tel que $\text{d}(b_j) = f_j$.
D'après la propriété universelle des algèbres polynomiales à puissances divisées du
lemme \ref{dpa-lemma-divided-power-polynomial-algebra}, il
existe un relèvement $\varphi(1) : A(1) \to B$ de $\varphi(0)$
qui envoie $T_j$ sur $f_j$.

\medskip\noindent
Après avoir construit $\varphi(m)$ pour un certain $m \geq 1$, on construit
$\varphi(m + 1) : A(m + 1) \to B$ exactement de la même manière.
Nous omettons les détails.
\end{proof}

\begin{lemma}
\label{dpa-lemma-divided-powers-on-tor}
Soit $R$ un anneau commutatif. Soient $S$ et $T$ des $R$-algèbres
commutatives. Il existe alors
une structure canonique de $R$-algèbre strictement commutative au sens gradué et munie de puissances divisées sur
$$
\operatorname{Tor}_*^R(S, T).
$$
\end{lemma}

\begin{proof}
Choisissons une factorisation $R \to A \to S$ comme ci-dessus. Comme $A \to S$
est un quasi-isomorphisme et que $A_d$ est un $R$-module
libre, l'algèbre différentielle graduée $B = A \otimes_R T$ calcule
les groupes de Tor affichés dans le lemme. Choisissons une surjection
$R[y_j:j\in J] \to T$. Alors
$B$ est un quotient de l'algèbre différentielle graduée
$A[y_j:j\in J]$, dont l'homologie est concentrée en degré $0$ (elle est
égale à $S[y_j:j\in J]$).
D'après le lemme \ref{dpa-lemma-base-change-div}, les algèbres différentielles graduées $B$ et
$A[y_j:j\in J]$ possèdent des structures de puissances divisées
compatibles avec les différentielles. On obtient donc la structure de
puissances divisées cherchée sur $H(B)$ par le lemme \ref{dpa-lemma-dpdga-good}.

\medskip\noindent
La structure d'algèbre à puissances divisées ainsi construite ne dépend
pas du choix de $A$. En effet, si $A'$ est un second choix, le
lemme \ref{dpa-lemma-uniqueness-tate-resolution}
fournit un morphisme $A \to A'$ qui préserve toutes les structures ainsi
que les augmentations vers $S$. Le morphisme induit
$B = A \otimes_R T \to A' \otimes_R T' = B'$ préserve alors
lui aussi toutes les
structures et est un quasi-isomorphisme. L'isomorphisme induit entre
les algèbres de Tor est donc compatible avec les produits et
les puissances divisées.
\end{proof}





\section{Application aux intersections complètes}
\label{dpa-section-application-ci}

\noindent
Soit $R$ un anneau. Soit $(A, \text{d}, \gamma)$ comme dans
la définition \ref{dpa-definition-divided-powers-dga}.
Une {\it dérivation} de degré $2$ est une application $R$-linéaire
$\theta : A \to A$ possédant les
propriétés suivantes :
\begin{enumerate}
\item $\theta(A_d) \subset A_{d - 2}$,
\item $\theta(xy) = \theta(x)y + x\theta(y)$,
\item $\theta$ commute avec $\text{d}$,
\item $\theta(\gamma_n(x)) = \theta(x) \gamma_{n - 1}(x)$
pour tout $x \in A_{2d}$ et tout $d$.
\end{enumerate}
Dans le lemme suivant, nous construisons une dérivation.

\begin{lemma}
\label{dpa-lemma-get-derivation}
Soit $R$ un anneau. Soit $(A, \text{d}, \gamma)$ comme dans
la définition \ref{dpa-definition-divided-powers-dga}.
Soit $R' \to R$ une surjection d'anneaux dont le
noyau est de carré nul et engendré par un élément $f$.
Si $A$ est une algèbre polynomiale graduée à puissances divisées sur $R$,
ayant un nombre fini de variables en chaque degré,
alors on obtient une dérivation
$\theta : A/IA \to A/IA$, où $I$ est l'annulateur
de $f$ dans $R$.
\end{lemma}

\begin{proof}
Comme $A$ est une algèbre polynomiale à puissances divisées, il existe une algèbre
polynomiale à puissances divisées $A'$ sur $R'$ telle que $A = A' \otimes_{R'} R$.
De plus, on peut relever $\text{d}$ en un opérateur $R$-linéaire
$\text{d}$ sur $A'$ tel que
\begin{enumerate}
\item $\text{d}(xy) = \text{d}(x)y + (-1)^{\deg(x)}x \text{d}(y)$
pour $x, y \in A'$ homogènes, et
\item $\text{d}(\gamma_n(x)) = \text{d}(x) \gamma_{n - 1}(x)$ pour
$x \in A'_{even, +}$.
\end{enumerate}
Nous omettons les détails (indication : procéder une variable
à la fois). Il se peut toutefois que $\text{d}^2$
ne soit pas nul sur $A'$. Il est clair que $\text{d}^2$ envoie $A'$ dans
$fA' \cong A/IA$. Par conséquent, $\text{d}^2$ annule $fA'$ et se factorise en
une application $A \to A/IA$. Comme $\text{d}^2$ est $R$-linéaire, on obtient
l'application voulue $\theta : A/IA \to A/IA$. La vérification des propriétés
d'une dérivation est immédiate.
\end{proof}

\begin{lemma}
\label{dpa-lemma-compute-theta} On
conserve les hypothèses et les notations du lemme \ref{dpa-lemma-get-derivation}.
Supposons que $S = H_0(A)$ soit isomorphe à
$R[x_1, \ldots, x_n]/(f_1, \ldots, f_m)$
pour certains $n$, $m$ et $f_j \in R[x_1, \ldots, x_n]$.
Supposons en outre donnée une relation
$$
\sum r_j f_j = 0
$$
avec $r_j \in R[x_1, \ldots, x_n]$.
Choisissons des relèvements $r'_j, f'_j \in R'[x_1, \ldots, x_n]$ de $r_j, f_j$.
Écrivons $\sum r'_j f'_j = gf$ pour un certain $g \in R/I[x_1, \ldots, x_n]$.
Si $H_1(A) = 0$ et si tous les coefficients de chaque $r_j$ appartiennent à $I$,
alors il existe un élément $\xi\in H_2(A/IA)$ tel que
$\theta(\xi) = g$ dans $S/IS$.
\end{lemma}

\begin{proof}
Soit $A(0) \subset A(1) \subset A(2) \subset\ldots$ la filtration
de $A$ telle que $A(m)$ s'obtienne à partir de $A(m - 1)$ en adjoignant des
variables à puissances divisées de degré $m$. Alors $A(0)$ est une algèbre
polynomiale sur $R$, munie d'une surjection de $R$-algèbres $A(0) \to S$.
On peut donc choisir une application
$$
\varphi : R[x_1, \ldots, x_n] \to A(0)
$$
qui relève les applications d'augmentation vers $S$.
Ensuite, $A(1) = A(0)\langle T_1, \ldots, T_t \rangle$, où
les $T_i$ sont des variables à puissances divisées de degré $1$. Comme $H_0(A) = S$, on
peut choisir des éléments $\xi_j \in\sum A(0)T_i$ tels que $\text{d}(\xi_j) = \varphi(f_j)$.
Alors
$$
\text{d}\left(\sum\varphi(r_j) \xi_j\right) =
\sum\varphi(r_j) \varphi(f_j) = \sum\varphi(r_jf_j) = 0
$$
Comme $H_1(A) = 0$, on peut choisir $\xi\in A_2$ tel que
$\text{d}(\xi) = \sum\varphi(r_j) \xi_j$.
Si les coefficients de $r_j$ appartiennent à $I$,
alors $\varphi(r_j)$ possède aussi cette propriété. Dans ce cas,
$\text{d}(\xi)$ s'annule dans $A_1/IA_1$ et
$\xi$ définit donc une classe dans $H_2(A/IA)$.

\medskip\noindent
La construction de $\theta$ dans la démonstration du lemme \ref{dpa-lemma-get-derivation}
consiste à relever successivement $A(i)$ en $A'(i)$ ainsi qu'à relever
la différentielle $\text{d}$. On relève $\varphi$
en $\varphi' : R'[x_1, \ldots, x_n] \to A'(0)$.
Ensuite, on a $A'(1) = A'(0)\langle T_1, \ldots, T_t\rangle$.
De plus, on peut relever $\xi_j$ en $\xi'_j \in\sum A'(0)T_i$.
Alors $\text{d}(\xi'_j) = \varphi'(f'_j) + f a_j$ pour un certain
$a_j \in A'(0)$.
Considérons un relèvement $\xi' \in A'_2$ de $\xi$.
On sait alors que
$$
\text{d}(\xi') = \sum\varphi'(r'_j)\xi'_j + \sum fb_iT_i
$$
pour certains $b_i \in A(0)$. En appliquant de nouveau $\text{d}$, on obtient
$$
\theta(\xi) = \sum\varphi'(r'_j)\varphi'(f'_j) +
\sum f \varphi'(r'_j) a_j + \sum fb_i \text{d}(T_i)
$$
Le premier terme donne le résultat voulu. Le deuxième est
nul, car les coefficients de $r_j$ appartiennent à $I$ et sont
donc annulés par $f$. Le troisième a une image nulle dans $H_0$,
car l'image de $\text{d}(T_i)$ y est nulle.
\end{proof}

\noindent
La méthode de démonstration du lemme suivant semble due à Gulliksen.

\begin{lemma}
\label{dpa-lemma-not-finite-pd}
Soit $R' \to R$ une surjection d'anneaux noethériens dont le noyau est de
carré nul et engendré par un élément $f$. Soit
$S = R[x_1, \ldots, x_n]/(f_1, \ldots, f_m)$.
Soit $\sum r_j f_j = 0$ une relation dans $R[x_1, \ldots, x_n]$.
Supposons que
\begin{enumerate}
\item chaque $r_j$ ait ses coefficients dans l'annulateur $I$ de $f$ dans $R$,
\item pour certains relèvements $r'_j, f'_j \in R'[x_1, \ldots, x_n]$, on ait
$\sum r'_j f'_j = gf$, où $g$ n'est pas nilpotent dans $S/IS$.
\end{enumerate}
Alors $S$ n'est pas de dimension de Tor finie sur $R$ (autrement dit, $S$ n'est
pas une $R$-algèbre parfaite).
\end{lemma}

\begin{proof}
Choisissons une résolution de Tate $R \to A \to S$ comme
dans le lemme \ref{dpa-lemma-tate-resolution}.
Soient $\xi\in H_2(A/IA)$ et $\theta : A/IA \to A/IA$ l'élément et la
dérivation obtenus dans les lemmes \ref{dpa-lemma-get-derivation} et
\ref{dpa-lemma-compute-theta}.
On observe que
$$
\theta^n(\gamma_n(\xi)) = g^n
$$
dans $H_0(A/IA) = S/IS$.
Ainsi, si $g$ n'est pas nilpotent dans $S/IS$, alors $\xi^n$ est non nul dans
$H_{2n}(A/IA)$ pour tout $n > 0$. Comme
$H_{2n}(A/IA) = \text{Tor}^R_{2n}(S, R/I)$, on conclut.
\end{proof}

\noindent
Le résultat suivant se trouve dans \cite{Rodicio}.

\begin{lemma}
\label{dpa-lemma-injective}
Soit $(A, \mathfrak m)$ un anneau local noethérien. Soient
$I \subset J \subset A$ des idéaux propres. Si $A/J$ est de
dimension de Tor finie sur $A/I$, alors l'application $I/\mathfrak m I \to J/\mathfrak m J$
est injective.
\end{lemma}

\begin{proof}
Soit $f \in I$ un élément dont l'image dans $I/\mathfrak m I$
est non nulle et dont l'image dans $J/\mathfrak mJ$ est nulle. On peut choisir un idéal $I'$
tel que $\mathfrak mI \subset I' \subset I$ et que $I/I'$ soit engendré
par l'image de $f$. Posons $R = A/I$ et $R' = A/I'$. Écrivons $J = (a_1, \ldots, a_m)$
avec $a_j \in A$. Alors $f = \sum b_j a_j$ pour certains $b_j \in\mathfrak m$.
Soient $r_j, f_j \in R$ resp.\ $r'_j, f'_j \in R'$ les images de $b_j, a_j$.
On se trouve alors
dans la situation du lemme \ref{dpa-lemma-not-finite-pd}
(l'idéal $I$ de ce lemme étant égal à $\mathfrak m_R$),
ce qui démontre le lemme.
\end{proof}

\begin{lemma}
\label{dpa-lemma-regular-sequence}
Soit $(A, \mathfrak m)$ un anneau local noethérien. Soient
$I \subset J \subset A$ des idéaux propres. Supposons que
\begin{enumerate}
\item $A/J$ soit de dimension de Tor finie sur $A/I$, et
\item $J$ soit engendré par une suite régulière.
\end{enumerate}
Alors $I$ est engendré par une suite régulière et $J/I$
est engendré par une suite régulière.
\end{lemma}

\begin{proof}
Le lemme \ref{dpa-lemma-injective} montre que l'application
$I/\mathfrak m I \to J/\mathfrak m J$
est injective. Il existe donc $s \leq r$ et un système minimal de
générateurs $f_1, \ldots, f_r$ de $J$ tels que $f_1, \ldots, f_s$ appartiennent à $I$
et forment un système minimal de générateurs de $I$.
Le lemme en résulte, car tout système minimal de générateurs de $J$
est une suite régulière d'après
Compléments sur l'algèbre, lemmes
\ref{more-algebra-lemma-independence-of-generators} et
\ref{more-algebra-lemma-noetherian-finite-all-equivalent}.
\end{proof}

\begin{lemma}
\label{dpa-lemma-perfect-map-ci}
Soit $R \to S$ un homomorphisme local d'anneaux locaux noethériens.
Soient $I \subset R$ et $J \subset S$ des idéaux tels que
$IS \subset J$. Si $R \to S$ est plat et si $S/\mathfrak m_RS$ est
régulier, alors les propriétés suivantes sont équivalentes :
\begin{enumerate}
\item $J$ est engendré par une suite régulière et
$S/J$ est de dimension de Tor finie comme module sur $R/I$,
\item $J$ est engendré par une suite régulière et
$\text{Tor}^{R/I}_p(S/J, R/\mathfrak m_R)$ n'est non nul
que pour un nombre fini d'entiers $p$,
\item $I$ est engendré par une suite
régulière et $J/IS$ est engendré par une suite régulière dans $S/IS$.
\end{enumerate}
\end{lemma}

\begin{proof}
Si (3) est vérifiée, alors $J$ est engendré par une suite régulière ; voir, par exemple,
Compléments sur l'algèbre, lemmes
\ref{more-algebra-lemma-join-koszul-regular-sequences} et
\ref{more-algebra-lemma-noetherian-finite-all-equivalent}. De
plus, sous cette même hypothèse, $S/J = (S/I)/(J/I)$
est de dimension projective finie sur $S/IS$, car le complexe
de Koszul fournit une résolution libre finie de $S/J$ sur $S/IS$.
Comme $R/I \to S/IS$ est plat, il s'ensuit que $S/J$ est de
dimension de Tor finie sur $R/I$,
d'après Compléments sur l'algèbre, lemme \ref{more-algebra-lemma-flat-push-tor-amplitude}.
Ainsi, (3) implique (1).

\medskip\noindent
L'implication (1) $\Rightarrow$ (2) est immédiate.
Supposons (2). D'après
Compléments sur l'algèbre, lemme \ref{more-algebra-lemma-perfect-over-regular-local-ring},
$S/J$ est de dimension de Tor finie sur $S/IS$.
On peut donc appliquer le lemme \ref{dpa-lemma-regular-sequence}
pour conclure que $IS$ et $J/IS$ sont engendrés par des suites régulières.
Soit $f_1, \ldots, f_r \in I$ un système minimal de générateurs de $I$.
Comme $R \to S$ est plat, les éléments $f_1, \ldots, f_r$ forment un
système minimal de générateurs de $IS$ dans $S$. Ainsi, $f_1, \ldots, f_r \in R$
est une suite d'éléments dont les images dans $S$ forment une suite régulière,
d'après Compléments sur l'algèbre, lemmes
\ref{more-algebra-lemma-independence-of-generators} et
\ref{more-algebra-lemma-noetherian-finite-all-equivalent}. Par
conséquent, $f_1, \ldots, f_r$ est une suite régulière dans $R$,
d'après Algèbre, lemme \ref{algebra-lemma-flat-increases-depth}.
\end{proof}





\section{Anneaux d'intersection complète locale}
\label{dpa-section-lci}

\noindent
Soit $(A, \mathfrak m)$ un anneau local noethérien complet.
D'après le théorème de structure de Cohen (voir
Algèbre, théorème \ref{algebra-theorem-cohen-structure-theorem}),
on peut écrire $A$ comme quotient d'un anneau local
noethérien régulier et complet $R$. On dit que $A$ est une
{\it intersection
complète} s'il existe une surjection $R \to A$,
où $R$ est un anneau local régulier, dont le
noyau est engendré par une suite
régulière. Le lemme suivant montre que cette notion ne dépend
pas du choix de la surjection.

\begin{lemma}
\label{dpa-lemma-ci-well-defined}
Soit $(A, \mathfrak m)$ un anneau local noethérien complet. Les
propriétés suivantes sont équivalentes :
\begin{enumerate}
\item pour toute surjection d'anneaux locaux $R \to A$, où $R$
est un anneau local régulier, le noyau de $R \to A$ est engendré
par une suite régulière, et
\item pour une certaine surjection d'anneaux locaux $R \to A$, où $R$
est un anneau local régulier, le noyau de $R \to A$ est engendré
par une suite régulière.
\end{enumerate}
\end{lemma}

\begin{proof}
Soit $k$ le corps résiduel de $A$. Si la caractéristique de
$k$ est $p > 0$, on désigne par $\Lambda$ un anneau de
Cohen (Algèbre, définition \ref{algebra-definition-cohen-ring})
de corps résiduel $k$ (Algèbre, lemme \ref{algebra-lemma-cohen-rings-exist}).
Si la caractéristique de $k$ est $0$, on pose $\Lambda = k$.
Rappelons que $\Lambda[[x_1, \ldots, x_n]]$, pour tout $n$,
est formellement lisse sur $\mathbf{Z}$, resp.\ $\mathbf{Q}$,
pour la topologie $\mathfrak m$-adique ; voir
Compléments sur l'algèbre, lemme
\ref{more-algebra-lemma-power-series-ring-over-Cohen-fs}.
Fixons une surjection $\Lambda[[x_1, \ldots, x_n]] \to A$ donnée
par le théorème de structure de
Cohen (Algèbre, théorème \ref{algebra-theorem-cohen-structure-theorem}).

\medskip\noindent
Soit $R \to A$ une surjection, où $R$ est un
anneau local régulier. Soit $f_1, \ldots, f_r$ un système minimal de générateurs
de $\Ker(R \to A)$. Nous utiliserons sans autre mention le
fait qu'un idéal d'un anneau local noethérien est engendré par une suite
régulière si et seulement si tout système minimal de générateurs est
une suite régulière. On observe que $f_1, \ldots, f_r$
est une suite régulière dans $R$ si et seulement si $f_1, \ldots, f_r$
est une suite régulière dans le complété $R^\wedge$,
d'après Algèbre, lemmes \ref{algebra-lemma-flat-increases-depth} et
\ref{algebra-lemma-completion-flat}.
De plus,
$$
R^\wedge/(f_1, \ldots, f_r)R^\wedge =
(R/(f_1, \ldots, f_n))^\wedge = A^\wedge = A
$$
car $A$ est complet pour la topologie $\mathfrak m_A$-adique (la première égalité résulte
d'Algèbre, lemme \ref{algebra-lemma-completion-tensor}). Enfin,
l'anneau $R^\wedge$ est régulier puisque $R$ est régulier
(Compléments sur l'algèbre, lemme \ref{more-algebra-lemma-completion-regular}). On
peut donc supposer que $R$ est complet.

\medskip\noindent
Si $R$ est complet, on peut choisir une application
$\Lambda[[x_1, \ldots, x_n]] \to R$ relevant l'application donnée
$\Lambda[[x_1, \ldots, x_n]] \to A$; voir
Compléments sur l'algèbre, lemme \ref{more-algebra-lemma-lift-continuous}.
En ajoutant des variables $y_1, \ldots, y_m$ envoyées
sur des générateurs du noyau de $R \to A$, on peut supposer que
$\Lambda[[x_1, \ldots, x_n, y_1, \ldots, y_m]] \to R$ est surjective
(nous omettons quelques détails). On peut alors considérer le diagramme commutatif
$$
\xymatrix{
\Lambda[[x_1, \ldots, x_n, y_1, \ldots, y_m]] \ar[r] \ar[d] & R \ar[d] \\
\Lambda[[x_1, \ldots, x_n]] \ar[r] & A
}
$$
Le lemme \ref{algebra-lemma-ci-well-defined} d'Algèbre montre que
la condition relative à $R \to A$ équivaut à celle relative à
l'application fixée
$\Lambda[[x_1, \ldots, x_n]] \to A$. Cela achève la démonstration du lemme.
\end{proof}

\noindent
Les deux lemmes suivants constituent des vérifications de cohérence de la définition donnée ci-dessus.

\begin{lemma}
\label{dpa-lemma-quotient-regular-ring-by-regular-sequence}
Soit $R$ un anneau régulier. Soit $\mathfrak p \subset R$ un idéal premier.
Soit $f_1, \ldots, f_r \in\mathfrak p$ une suite régulière.
Alors le complété de
$$
A = (R/(f_1, \ldots, f_r))_\mathfrak p =
R_\mathfrak p/(f_1, \ldots, f_r)R_\mathfrak p
$$
est une intersection complète au sens défini ci-dessus.
\end{lemma}

\begin{proof}
Le complété de $A$ est égal à
$A^\wedge = R_\mathfrak p^\wedge/(f_1, \ldots, f_r)R_\mathfrak p^\wedge$,
car le foncteur de complétion est exact sur les modules finis sur l'anneau noethérien
$R_\mathfrak p$ (Algèbre,
lemme \ref{algebra-lemma-completion-tensor}).
L'image de la suite $f_1, \ldots, f_r$ dans $R_\mathfrak p$
est une suite régulière
d'après Algèbre, lemmes \ref{algebra-lemma-completion-flat} et
\ref{algebra-lemma-flat-increases-depth}.
De plus, $R_\mathfrak p^\wedge$ est un anneau local régulier d'après
Compléments sur l'algèbre, lemme \ref{more-algebra-lemma-completion-regular}. Le
résultat découle donc de notre définition d'une intersection
complète pour les anneaux locaux complets.
\end{proof}

\noindent
Le lemme suivant est l'analogue d'Algèbre, lemme \ref{algebra-lemma-lci}.

\begin{lemma}
\label{dpa-lemma-quotient-regular-ring}
Soit $R$ un anneau régulier. Soit $\mathfrak p \subset R$ un idéal premier.
Soit $I \subset\mathfrak p$ un idéal.
Posons $A = (R/I)_\mathfrak p = R_\mathfrak p/I_\mathfrak p$.
Les propriétés suivantes sont équivalentes :
\begin{enumerate}
\item le complété de $A$
est une intersection complète au sens ci-dessus,
\item $I_\mathfrak p \subset R_\mathfrak p$ est engendré
par une suite régulière,
\item le module $(I/I^2)_\mathfrak p$ peut être engendré par
$\dim(R_\mathfrak p) - \dim(A)$ éléments,
\item ajouter quelque chose ici.
\end{enumerate}
\end{lemma}

\begin{proof} On
peut remplacer, et l'on remplace, $R$ par son localisé en $\mathfrak p$.
Alors $\mathfrak p = \mathfrak m$ est l'idéal maximal de $R$
et $A = R/I$. Soit $f_1, \ldots, f_r \in I$ un système minimal
de générateurs. Le complété de $A$ est égal à
$A^\wedge = R^\wedge/(f_1, \ldots, f_r)R^\wedge$,
car le foncteur de complétion est exact sur les modules finis sur l'anneau noethérien
$R_\mathfrak p$ (Algèbre,
lemme \ref{algebra-lemma-completion-tensor}).

\medskip\noindent
Si (1) est vérifiée, alors l'image de la suite $f_1, \ldots, f_r$ dans $R^\wedge$
est une suite régulière par hypothèse. C'est donc une suite régulière
dans $R$ d'après Algèbre, lemmes \ref{algebra-lemma-completion-flat} et
\ref{algebra-lemma-flat-increases-depth}. Ainsi, (1) implique (2).

\medskip\noindent
Supposons (3). Posons $c = \dim(R) - \dim(A)$ et choisissons $f_1, \ldots, f_c \in I$
dont les images engendrent $I/I^2$. Le lemme de Nakayama
(Algèbre, lemme \ref{algebra-lemma-NAK})
montre que $I = (f_1, \ldots, f_c)$. Comme $R$ est régulier, donc de
Cohen--Macaulay (Algèbre, proposition \ref{algebra-proposition-CM-module}),
$f_1, \ldots, f_c$ est une suite régulière
d'après Algèbre, proposition \ref{algebra-proposition-CM-module}.
Ainsi, (3) implique (2).
Enfin, (2) implique (1) d'après le
lemme \ref{dpa-lemma-quotient-regular-ring-by-regular-sequence}.
\end{proof}

\noindent
Le résultat suivant est dû à Avramov ; voir \cite{Avramov}.

\begin{proposition}
\label{dpa-proposition-avramov}
Soit $A \to B$ un homomorphisme local plat d'anneaux locaux noethériens. Les
propriétés suivantes sont équivalentes :
\begin{enumerate}
\item $B^\wedge$ est une intersection complète,
\item $A^\wedge$ et $(B/\mathfrak m_A B)^\wedge$ sont des intersections complètes.
\end{enumerate}
\end{proposition}

\begin{proof}
Considérons le diagramme
$$
\xymatrix{
B \ar[r] & B^\wedge\\
A \ar[u] \ar[r] & A^\wedge\ar[u]
}
$$
Comme les applications horizontales sont fidèlement
plates (Algèbre, lemme \ref{algebra-lemma-completion-faithfully-flat}),
la flèche verticale de droite est plate (par
exemple, d'après Algèbre, lemme
\ref{algebra-lemma-criterion-flatness-fibre-Noetherian}).
De plus,
$(B/\mathfrak m_A B)^\wedge = B^\wedge/\mathfrak m_{A^\wedge} B^\wedge$
d'après Algèbre, lemme \ref{algebra-lemma-completion-tensor}. On
peut donc supposer que $A$ et $B$ sont des anneaux locaux noethériens complets.

\medskip\noindent
Supposons que $A$ et $B$ soient des anneaux locaux noethériens complets.
Choisissons un diagramme
$$
\xymatrix{
S \ar[r] & B \\
R \ar[u] \ar[r] & A \ar[u]
}
$$
comme dans Compléments sur l'algèbre, lemme
\ref{more-algebra-lemma-embed-map-Noetherian-complete-local-rings}.
Posons $I = \Ker(R \to A)$ et $J = \Ker(S \to B)$.
Comme $R/I = A \to B = S/J$ est plat, l'application
$J/IS \otimes_R R/\mathfrak m_R \to J/J \cap\mathfrak m_R S$
est un isomorphisme. Par conséquent, un système minimal de générateurs de $J/IS$
s'envoie sur un système minimal de générateurs de
$\Ker(S/\mathfrak m_R S \to B/\mathfrak m_A B)$.
Enfin, $S/\mathfrak m_R S$ est un anneau local régulier.

\medskip\noindent
Supposons (1), c'est-à-dire que $J$ soit engendré par une suite régulière.
Comme $A = R/I \to B = S/J$ est plat, le
lemme \ref{dpa-lemma-perfect-map-ci} s'applique et montre
que $I$ et $J/IS$ sont engendrés par des suites
régulières. On a $\dim(B) = \dim(A) + \dim(B/\mathfrak m_A B)$ et
$\dim(S/IS) = \dim(A) + \dim(S/\mathfrak m_R S)$
(Algèbre, lemme \ref{algebra-lemma-dimension-base-fibre-equals-total}).
Ainsi, $J/IS$ est engendré par
$$
\dim(S/IS) - \dim(S/J) = \dim(S/\mathfrak m_R S) - \dim(B/\mathfrak m_A B)
$$
éléments (Algèbre, lemme \ref{algebra-lemma-one-equation}).
Il s'ensuit que $\Ker(S/\mathfrak m_R S \to B/\mathfrak m_A B)$
est engendré par le même nombre d'éléments (voir ci-dessus). Par
conséquent, $\Ker(S/\mathfrak m_R S \to B/\mathfrak m_A B)$
est engendré par une suite régulière ; voir, par exemple, le
lemme \ref{dpa-lemma-quotient-regular-ring}.
On obtient ainsi (2).

\medskip\noindent
Si (2) est vérifiée, alors $I$ et $J/J \cap\mathfrak m_RS$
sont engendrés par des suites régulières. En relevant ces générateurs (voir
ci-dessus), puis en utilisant la platitude de $R/I \to S/IS$ et
le lemme de Grothendieck
(Algèbre, lemme \ref{algebra-lemma-grothendieck-regular-sequence}),
on voit que $J/IS$ est engendré par une suite régulière dans $S/IS$.
Le lemme \ref{dpa-lemma-perfect-map-ci} montre alors que $J$
est engendré par une suite régulière, d'où (1).
\end{proof}

\begin{definition}
\label{dpa-definition-lci}
Soit $A$ un anneau noethérien.
\begin{enumerate}
\item Si $A$ est local, on dit que $A$ est une {\it intersection complète}
si son complété est une intersection complète au sens ci-dessus.
\item En général, on dit que $A$ est une {\it intersection complète locale}
si tous ses anneaux locaux sont des intersections complètes.
\end{enumerate}
\end{definition}

\noindent
Nous vérifierons plus loin que cela ne contredit pas la terminologie
introduite dans
Algèbre, définitions \ref{algebra-definition-lci-field} et
\ref{algebra-definition-lci-local-ring}.
Mais montrons d'abord que cette définition est
cohérente : si $A$ est un anneau local
noethérien qui est une intersection complète, alors $A$ est une intersection complète locale,
c'est-à-dire que tous ses anneaux locaux sont des intersections complètes.

\begin{lemma}
\label{dpa-lemma-ci-good}
Soit $(A, \mathfrak m)$ un anneau local noethérien. Soit
$\mathfrak p \subset A$ un idéal premier. Si $A$ est une
intersection complète, alors $A_\mathfrak p$ est aussi une intersection complète.
\end{lemma}

\begin{proof}
Choisissons un idéal premier $\mathfrak q$ de $A^\wedge$ au-dessus de $\mathfrak p$
(cela est possible puisque $A \to A^\wedge$ est fidèlement plat
d'après Algèbre, lemme \ref{algebra-lemma-completion-faithfully-flat}).
Alors $A_\mathfrak p \to (A^\wedge)_\mathfrak q$ est un
homomorphisme local plat. La proposition \ref{dpa-proposition-avramov}
montre donc que $A_\mathfrak p$ est une intersection complète si et seulement si
$(A^\wedge)_\mathfrak q$ est une intersection complète. Il suffit ainsi
de démontrer le lemme lorsque $A$ est complet (c'est l'étape essentielle de
la démonstration).

\medskip\noindent
Supposons $A$ complet. Par définition, on peut écrire
$A = R/(f_1, \ldots, f_r)$ pour une suite régulière
$f_1, \ldots, f_r$ dans un anneau local régulier $R$.
Soit $\mathfrak q \subset R$ l'idéal premier correspondant à $\mathfrak p$.
On observe que $f_1, \ldots, f_r \in\mathfrak q$ et que
$A_\mathfrak p = R_\mathfrak q/(f_1, \ldots, f_r)R_\mathfrak q$.
Ainsi, $A_\mathfrak p$ est une intersection complète d'après
le lemme \ref{dpa-lemma-quotient-regular-ring-by-regular-sequence}.
\end{proof}

\begin{lemma}
\label{dpa-lemma-check-lci-at-maximal-ideals}
Soit $A$ un anneau noethérien. Alors $A$ est une intersection complète locale
si et seulement si $A_\mathfrak m$ est une intersection complète pour
tout idéal maximal $\mathfrak m$ de $A$.
\end{lemma}

\begin{proof}
Cela résulte immédiatement du lemme \ref{dpa-lemma-ci-good} et des définitions.
\end{proof}

\begin{lemma}
\label{dpa-lemma-check-lci-agrees}
Soit $S$ une algèbre de type fini sur un corps $k$.
\begin{enumerate}
\item pour un idéal premier $\mathfrak q \subset S$, l'anneau local $S_\mathfrak q$
est une intersection complète au sens
de la définition \ref{algebra-definition-lci-local-ring} d'Algèbre,
si et seulement si $S_\mathfrak q$ est une
intersection complète au sens de la définition \ref{dpa-definition-lci}, et
\item $S$ est une intersection complète locale au
sens de la définition \ref{algebra-definition-lci-field}
d'Algèbre, si et seulement si $S$ est une
intersection complète locale au sens de la définition \ref{dpa-definition-lci}.
\end{enumerate}
\end{lemma}

\begin{proof}
Démontrons (1). Soit $k[x_1, \ldots, x_n] \to S$ une surjection.
Soit $\mathfrak p \subset k[x_1, \ldots, x_n]$ l'idéal premier
correspondant à $\mathfrak q$.
Soit $I \subset k[x_1, \ldots, x_n]$ le noyau de cette surjection.
L'application $k[x_1, \ldots, x_n]_\mathfrak p \to S_\mathfrak q$
est surjective, de noyau $I_\mathfrak p$. Or
$k[x_1, \ldots, x_n]$ est un anneau régulier
d'après Algèbre, proposition \ref{algebra-proposition-finite-gl-dim-polynomial-ring}.
L'équivalence des deux notions de (1) résulte donc de la
combinaison du
lemme \ref{dpa-lemma-quotient-regular-ring}
avec Algèbre, lemme \ref{algebra-lemma-lci-local}.

\medskip\noindent
Après avoir démontré (1), l'équivalence de (2) résulte de la
définition et d'Algèbre, lemme \ref{algebra-lemma-lci-global}.
\end{proof}

\begin{lemma}
\label{dpa-lemma-avramov}
Soit $A \to B$ un homomorphisme local plat d'anneaux locaux noethériens. Les
propriétés suivantes sont équivalentes :
\begin{enumerate}
\item $B$ est une intersection complète,
\item $A$ et $B/\mathfrak m_A B$ sont des intersections complètes.
\end{enumerate}
\end{lemma}

\begin{proof}
Maintenant que la définition est cohérente, il s'agit d'une reformulation
immédiate de la proposition (non triviale) \ref{dpa-proposition-avramov}.
\end{proof}









\section{Homomorphismes d'intersection complète locale}
\label{dpa-section-lci-homomorphisms}

\noindent
Soit $A \to B$ un homomorphisme local d'anneaux locaux noethériens complets. Une
conséquence du théorème de structure de Cohen est que l'on peut trouver un
diagramme commutatif
$$
\xymatrix{
S \ar[r] & B \\
& A \ar[lu] \ar[u]
}
$$
d'anneaux locaux noethériens complets tel que
$S \to B$ soit surjectif, $A \to S$ soit plat et que
$S/\mathfrak m_A S$ soit un anneau local régulier. Cela résulte de
Compléments sur l'algèbre, lemme
\ref{more-algebra-lemma-embed-map-Noetherian-complete-local-rings}.
Disons (provisoirement) que $A \to S \to B$ est une {\it bonne factorisation}
de $A \to B$ si $S$ est un anneau local noethérien,
$A \to S \to B$ sont des homomorphismes locaux d'anneaux,
$S \to B$ est surjectif, $A \to S$ est plat et $S/\mathfrak m_AS$ est
régulier. Disons que $A \to B$ est un
{\it homomorphisme d'intersection
complète} s'il existe une bonne factorisation $A \to S \to B$
telle que le noyau de $S \to B$ soit engendré par une suite
régulière. Le lemme suivant montre que cette notion ne dépend
pas du choix du diagramme.

\begin{lemma}
\label{dpa-lemma-ci-map-well-defined}
Soit $A \to B$ un homomorphisme local d'anneaux locaux noethériens complets. Les
assertions suivantes sont équivalentes :
\begin{enumerate}
\item pour une certaine bonne factorisation $A \to S \to B$, le noyau de
$S \to B$ est engendré par une suite régulière, et
\item pour toute bonne factorisation $A \to S \to B$, le noyau de
$S \to B$ est engendré par une suite régulière.
\end{enumerate}
\end{lemma}

\begin{proof}
Soit $A \to S \to B$ une bonne factorisation.
Comme $B$ est complet, on obtient une factorisation
$A \to S^\wedge\to B$ où $S^\wedge$ est le complété de $S$.
Remarquons qu'il s'agit encore d'une bonne factorisation :
l'homomorphisme d'anneaux $S \to S^\wedge$ est
plat (Algèbre, lemme \ref{algebra-lemma-completion-flat}),
donc $A \to S^\wedge$ est plat.
L'anneau $S^\wedge/\mathfrak m_A S^\wedge = (S/\mathfrak m_A S)^\wedge$
est régulier puisque $S/\mathfrak m_A S$ est régulier
(Compléments sur l'algèbre, lemme \ref{more-algebra-lemma-completion-regular}).
Soit $f_1, \ldots, f_r$ un système minimal de générateurs
de $\Ker(S \to B)$. Nous utiliserons sans autre mention le
fait qu'un idéal d'un anneau local noethérien est engendré par une suite
régulière si et seulement si tout système minimal de générateurs est
une suite régulière. Remarquons que $f_1, \ldots, f_r$
est une suite régulière dans $S$ si et seulement si $f_1, \ldots, f_r$
est une suite régulière dans le complété $S^\wedge$,
d'après Algèbre, lemme \ref{algebra-lemma-flat-increases-depth}.
De plus, on a
$$
S^\wedge/(f_1, \ldots, f_r)R^\wedge =
(S/(f_1, \ldots, f_n))^\wedge = B^\wedge = B
$$
car $B$ est complet pour la topologie $\mathfrak m_B$-adique (la première égalité résulte
d'Algèbre, lemme \ref{algebra-lemma-completion-tensor}).
Ainsi, le noyau de $S \to B$ est engendré par une suite régulière
si et seulement si le noyau de $S^\wedge\to B$ est engendré par une
suite régulière. Il
suffit donc de considérer les bonnes factorisations pour lesquelles $S$ est complet.

\medskip\noindent
Supposons données deux factorisations $A \to S \to B$ et
$A \to S' \to B$, avec $S$ et $S'$ complets. D'après
Compléments sur l'algèbre, lemme \ref{more-algebra-lemma-dominate-two-surjections},
l'anneau $S \times_B S'$ est un anneau local noethérien complet. Par
conséquent, en appliquant Compléments sur l'algèbre, lemme
\ref{more-algebra-lemma-embed-map-Noetherian-complete-local-rings}, on
peut choisir une bonne factorisation $A \to S'' \to S \times_B S'$
avec $S''$ complet. Il suffit donc de montrer ceci :
si $A \to S' \to S \to B$ sont de bonnes factorisations comparables,
alors $\Ker(S \to B)$ est engendré par une suite régulière
si et seulement si $\Ker(S' \to B)$ est engendré par une suite régulière.

\medskip\noindent
Soient $A \to S' \to S \to B$ de bonnes factorisations comparables. Tout
d'abord, puisque $S'/\mathfrak m_R S' \to S/\mathfrak m_R S$ est
une surjection d'anneaux locaux réguliers, son noyau est engendré
par une suite régulière
$\overline{x}_1, \ldots, \overline{x}_c \in
\mathfrak m_{S'}/\mathfrak m_R S'$
que l'on peut compléter en un système régulier de paramètres de
l'anneau local régulier $S'/\mathfrak m_R S'$;
voir (Algèbre, lemme \ref{algebra-lemma-regular-quotient-regular}).
Posons $I = \Ker(S' \to S)$. Comme $S$ est plat sur $R$, on a
$$
I/\mathfrak m_R I =
\Ker(S'/\mathfrak m_R S' \to S/\mathfrak m_R S) =
(\overline{x}_1, \ldots, \overline{x}_c).
$$
Choisissons des relèvements $x_1, \ldots, x_c \in I$. Ceux-ci forment une suite régulière
qui engendre $I$, puisque $S'$ est plat sur $R$;
voir Algèbre, lemme \ref{algebra-lemma-grothendieck-regular-sequence}.

\medskip\noindent
On en déduit que, si $\Ker(S \to B)$ est lui aussi engendré
par une suite régulière, il en va de même de $\Ker(S' \to B)$; voir
Compléments sur l'algèbre, lemmes
\ref{more-algebra-lemma-join-koszul-regular-sequences} et
\ref{more-algebra-lemma-noetherian-finite-all-equivalent}.

\medskip\noindent
Réciproquement, supposons que $J = \Ker(S' \to B)$ soit
engendré par une suite régulière. Comme les générateurs $x_1, \ldots, x_c$
de $I$ s'envoient sur des éléments linéairement indépendants de
$\mathfrak m_{S'}/\mathfrak m_{S'}^2$, on voit que
$I/\mathfrak m_{S'}I \to J/\mathfrak m_{S'}J$ est injectif.
Il existe donc un système minimal de générateurs
$x_1, \ldots, x_c, y_1, \ldots, y_d$ de $J$.
Alors $x_1, \ldots, x_c, y_1, \ldots, y_d$ est une suite
régulière, et il s'ensuit que les images de $y_1, \ldots, y_d$ dans $S$
forment une suite régulière qui engendre $\Ker(S \to B)$.
Ceci achève la démonstration du lemme.
\end{proof}

\noindent
Dans la proposition suivante, remarquons que la condition d'annulation des
Tor est notamment satisfaite si $B$ est de dimension de Tor finie sur $A$,
et donc, en particulier, si $B$ est plat sur $A$.

\begin{proposition}
\label{dpa-proposition-avramov-map}
Soit $A \to B$ un homomorphisme local d'anneaux locaux noethériens. Les
assertions suivantes sont alors équivalentes :
\begin{enumerate}
\item $B$ est une intersection complète et
$\text{Tor}^A_p(B, A/\mathfrak m_A)$ n'est non nul que pour un nombre fini de $p$,
\item $A$ est une intersection complète et
$A^\wedge\to B^\wedge$ est un homomorphisme d'intersection
complète au sens défini ci-dessus.
\end{enumerate}
\end{proposition}

\begin{proof}
Soit $F_\bullet\to A/\mathfrak m_A$ une résolution par
des $A$-modules libres de type fini. Remarquons que
$\text{Tor}^A_p(B, A/\mathfrak m_A)$
est la $p$-ième homologie du complexe $F_\bullet\otimes_A B$.
Soit $F_\bullet^\wedge = F_\bullet\otimes_A A^\wedge$ le complété.
Alors $F_\bullet^\wedge$ est une résolution de $A^\wedge/\mathfrak m_{A^\wedge}$
par des $A^\wedge$-modules libres de type fini (car $A \to A^\wedge$ est plat et le foncteur de
complétion est exact sur les modules de type fini ;
voir Algèbre, lemmes \ref{algebra-lemma-completion-tensor} et
\ref{algebra-lemma-completion-flat}).
Il s'ensuit que
$$
F_\bullet^\wedge\otimes_{A^\wedge} B^\wedge =
F_\bullet\otimes_A B \otimes_B B^\wedge
$$
La platitude de $B \to B^\wedge$ donne alors
$$
\text{Tor}^{A^\wedge}_p(B^\wedge, A^\wedge/\mathfrak m_{A^\wedge}) =
\text{Tor}^A_p(B, A/\mathfrak m_A) \otimes_B B^\wedge
$$
On voit ainsi que la condition de (1) portant sur l'homomorphisme local $A \to B$
est équivalente à la même condition pour l'homomorphisme local
$A^\wedge\to B^\wedge$.
On peut donc supposer que $A$ et $B$ sont des anneaux locaux noethériens
complets (puisque, de toute façon, les autres conditions sont formulées en termes
des complétés).

\medskip\noindent
Supposons que $A$ et $B$ soient des anneaux locaux noethériens complets.
Choisissons un diagramme
$$
\xymatrix{
S \ar[r] & B \\
R \ar[u] \ar[r] & A \ar[u]
}
$$
comme dans Compléments sur l'algèbre, lemme
\ref{more-algebra-lemma-embed-map-Noetherian-complete-local-rings}.
Posons $I = \Ker(R \to A)$ et $J = \Ker(S \to B)$.
La proposition résulte alors du lemme \ref{dpa-lemma-perfect-map-ci}.
\end{proof}

\begin{remark}
\label{dpa-remark-no-good-ci-map} Il
semble difficile de définir une bonne notion d'« homomorphismes
d'intersection complète locale » pour les homomorphismes entre anneaux noethériens généraux.
La raison en est que, pour un anneau local noethérien $A$, les fibres de
$A \to A^\wedge$ ne sont pas des anneaux d'intersection complète
locale. Ainsi, si $A \to B$ est un homomorphisme local d'anneaux locaux noethériens et si
l'homomorphisme entre les complétés $A^\wedge\to B^\wedge$ est
un homomorphisme d'intersection complète au sens défini ci-dessus,
alors $(A_\mathfrak p)^\wedge\to (B_\mathfrak q)^\wedge$ n'est en général
{\bf pas} un homomorphisme d'intersection complète au sens défini
ci-dessus. Une solution consiste à travailler exclusivement avec des anneaux
noethériens excellents. Plus généralement, on pourrait travailler avec les
anneaux noethériens dont les fibres formelles sont des
intersections complètes ; voir \cite{Rodicio-ci}.
Nous développerons cette théorie
dans Complexes dualisants, section \ref{dualizing-section-formal-fibres}.
\end{remark}

\noindent
Pour terminer cette section, comparons la notion définie ci-dessus
à celle introduite dans
Compléments sur l'algèbre, section \ref{dpa-section-lci}.

\begin{lemma}
\label{dpa-lemma-well-defined-if-you-can-find-good-factorization}
Considérons un diagramme commutatif
$$
\xymatrix{
S \ar[r] & B \\
& A \ar[lu] \ar[u]
}
$$
d'anneaux locaux noethériens tel que $S \to B$ soit surjectif, $A \to S$ soit plat et que
$S/\mathfrak m_A S$ soit un anneau local régulier. Les assertions suivantes sont équivalentes :
\begin{enumerate}
\item $\Ker(S \to B)$ est engendré par une suite régulière, et
\item $A^\wedge\to B^\wedge$ est un homomorphisme d'intersection complète au
sens défini ci-dessus.
\end{enumerate}
\end{lemma}

\begin{proof}
Démonstration omise. Indication : la démonstration est identique à l'argument donné
dans le premier paragraphe de la démonstration du lemme \ref{dpa-lemma-ci-map-well-defined}.
\end{proof}

\begin{lemma}
\label{dpa-lemma-finite-type-lci-map}
Soit $A$ un anneau noethérien.
Soit $A \to B$ un homomorphisme d'anneaux de type fini.
Les assertions suivantes sont équivalentes :
\begin{enumerate}
\item $A \to B$ est une intersection complète locale au sens de
Compléments sur l'algèbre, définition
\ref{more-algebra-definition-local-complete-intersection},
\item pour tout idéal premier $\mathfrak q \subset B$, en posant
$\mathfrak p = A \cap\mathfrak q$, l'homomorphisme d'anneaux
$(A_\mathfrak p)^\wedge\to (B_\mathfrak q)^\wedge$ est
un homomorphisme d'intersection complète au sens défini ci-dessus.
\end{enumerate}
\end{lemma}

\begin{proof}
Choisissons une surjection $R = A[x_1, \ldots, x_n] \to B$.
Remarquons que $A \to R$ est plat à fibres régulières.
Soit $I$ le noyau de $R \to B$.
Supposons (2). On voit alors que
$I$ est localement engendré par une suite régulière,
d'après le
lemme \ref{dpa-lemma-well-defined-if-you-can-find-good-factorization}
et
Algèbre, lemme \ref{algebra-lemma-regular-sequence-in-neighbourhood}.
Autrement dit, (1) est vérifiée.
Réciproquement, supposons (1). Après localisation sur $R$ et $B$,
on peut alors supposer que $I$ est engendré par une suite régulière au sens de Koszul.
D'après Compléments sur l'algèbre, lemme
\ref{more-algebra-lemma-noetherian-finite-all-equivalent},
on trouve que $I$ est localement engendré par une suite régulière. Ainsi,
(2) est vérifiée d'après le
lemme \ref{dpa-lemma-well-defined-if-you-can-find-good-factorization}.
Certains détails sont omis.
\end{proof}

\begin{lemma}
\label{dpa-lemma-avramov-map-finite-type}
Soit $A$ un anneau noethérien. Soit $A \to B$ un homomorphisme d'anneaux de
type fini tel que l'image de $\Spec(B) \to\Spec(A)$ contienne tous les
points fermés de $\Spec(A)$. Les assertions suivantes sont alors équivalentes :
\begin{enumerate}
\item $B$ est une intersection complète et $A \to B$ est de
dimension de Tor finie,
\item $A$ est une intersection complète et $A \to B$ est une intersection
complète locale au sens de Compléments sur l'algèbre, définition
\ref{more-algebra-definition-local-complete-intersection}.
\end{enumerate}
\end{lemma}

\begin{proof}
Il s'agit d'une reformulation de la proposition \ref{dpa-proposition-avramov-map} au
moyen du lemme \ref{dpa-lemma-finite-type-lci-map}.
Nous omettons les détails.
\end{proof}








\section{Homomorphismes d'anneaux lisses et diagonales}
\label{dpa-section-smooth-diagonal-perfect}

\noindent
Dans cette section, nous utilisons les résultats précédents pour caractériser les homomorphismes d'anneaux
lisses comme ceux dont la diagonale est parfaite.

\begin{lemma}
\label{dpa-lemma-local-perfect-diagonal}
Soit $A \to B$ un homomorphisme local d'anneaux locaux noethériens tel que
$B$ soit plat et essentiellement de type fini sur $A$. Si
$$
B \otimes_A B \longrightarrow B
$$
est un homomorphisme d'anneaux parfait, c'est-à-dire si $B$ est de dimension de Tor finie sur
$B \otimes_A B$, alors $B$ est la localisation d'une $A$-algèbre lisse.
\end{lemma}

\begin{proof}
Comme $B$ est essentiellement de type fini sur $A$, il en va de même de $B \otimes_A B$,
et en particulier $B \otimes_A B$ est noethérien. Par conséquent, le quotient $B$ de
$B \otimes_A B$ est pseudo-cohérent sur $B \otimes_A B$
(Compléments sur l'algèbre, lemme \ref{more-algebra-lemma-Noetherian-pseudo-coherent}), ce qui
explique pourquoi la perfection de l'homomorphisme d'anneaux (Compléments sur l'algèbre, définition
\ref{more-algebra-definition-pseudo-coherent-perfect}) équivaut à la
condition de dimension de Tor finie.

\medskip\noindent
On peut écrire $B = R/K$, où $R$ est la localisation de $A[x_1, \ldots, x_n]$
en un idéal premier et où $K \subset R$ est un idéal. Notons
$\mathfrak m \subset R \otimes_A R$ l'idéal maximal qui est l'image
réciproque de l'idéal maximal de $B$ par la surjection
$R \otimes_A R \to B \otimes_A B \to B$. On a alors des surjections
$$
(R \otimes_A R)_\mathfrak m \to (B \otimes_A B)_\mathfrak m \to B
$$
et donc des idéaux $I \subset J \subset (R \otimes_A R)_\mathfrak m$
comme dans le lemme \ref{dpa-lemma-injective}. On en déduit que
$I/\mathfrak m I \to J/\mathfrak m J$ est injectif.

\medskip\noindent
Écrivons $K = (f_1, \ldots, f_r)$ avec $r$ minimal. On peut supposer, et on le fera, que
$f_i \in R$ est l'image d'un élément de $A[x_1, \ldots, x_n]$ que l'on
note également $f_i$. Remarquons que $I$ est engendré
par $f_1 \otimes 1, \ldots, f_r \otimes 1$ et
$1 \otimes f_1, \ldots, 1 \otimes f_r$. Nous affirmons que c'est un système
minimal de générateurs de $I$. En effet, si $\kappa$ est le corps résiduel commun
de $R$, $B$, $(R \otimes_A R)_\mathfrak m$ et $(B \otimes_A B)_\mathfrak m$,
alors on dispose d'un homomorphisme
$R \otimes_A R \to R \otimes_A \kappa\oplus\kappa\otimes_A R$
qui se factorise par $(R \otimes_A R)_\mathfrak m$. Comme $B$ est
plat sur $A$ et que l'on a la suite exacte courte
$0 \to K \to R \to B \to 0$, on voit que
$K \otimes_A \kappa\subset R \otimes_A \kappa$;
voir Algèbre, lemme \ref{algebra-lemma-flat-tor-zero}. En
restreignant ainsi l'homomorphisme
$(R \otimes_A R)_\mathfrak m \to R \otimes_A \kappa\oplus\kappa\otimes_A R$
à $I$, on obtient un homomorphisme
$$
I \to K \otimes_A \kappa\oplus\kappa\otimes_A K \to
K \otimes_B \kappa\oplus\kappa\otimes_B K.
$$
Les éléments
$f_1 \otimes 1, \ldots, f_r \otimes 1, 1 \otimes f_1, \ldots, 1 \otimes f_r$
s'envoient sur une base du but de cet homomorphisme, puisque, d'après le lemme de
Nakayama (Algèbre, lemme \ref{algebra-lemma-NAK}),
$f_1, \ldots, f_r$ s'envoient sur une base de $K \otimes_B \kappa$.
Ceci démontre notre assertion.

\medskip\noindent
L'idéal $J$ est engendré par $f_1 \otimes 1, \ldots, f_r \otimes 1$
et par les éléments $x_1 \otimes 1 - 1 \otimes x_1, \ldots,
x_n \otimes 1 - 1 \otimes x_n$ (pour la démonstration, il
suffit de voir que ces éléments appartiennent à l'idéal $J$).
On peut alors écrire
$$
f_i \otimes 1 - 1 \otimes f_i =
\sum g_{ij} (x_j \otimes 1 - 1 \otimes x_j)
$$
pour certains $g_{ij}$ dans $(R \otimes_A R)_\mathfrak m$. C'est un fait général
concernant les éléments de $A[x_1, \ldots, x_n]$, dont nous omettons la
démonstration. Notons $a_{ij} \in\kappa$ l'image de $g_{ij}$. Un autre calcul
montre que $a_{ij}$ est l'image de $\partial f_i / \partial x_j$ dans $\kappa$.
L'injectivité de $I/\mathfrak m I \to J/\mathfrak m J$ et les remarques
précédentes imposent à la matrice $(a_{ij})$ d'être de rang maximal $r$.
Posons
$$
C = A[x_1, \ldots, x_n]/(f_1, \ldots, f_r)
$$
et considérons le complexe cotangent naïf
$\NL_{C/A} \cong (C^{\oplus r} \to C^{\oplus n})$
où l'homomorphisme est donné par la matrice des dérivées partielles.
Ainsi, $\NL_{C/A} \otimes_A B$
est isomorphe à un $B$-module libre de rang $n - r$ placé en degré $0$.
Ainsi, $C_g$ est lisse sur $A$ pour un certain $g \in C$ dont l'image
dans $B$ est une unité ; voir Algèbre, lemme \ref{algebra-lemma-smooth-at-point}.
Ceci achève la démonstration.
\end{proof}

\begin{lemma}
\label{dpa-lemma-perfect-diagonal}
Soit $A \to B$ un homomorphisme d'anneaux noethériens, plat et de type fini. Si
$$
B \otimes_A B \longrightarrow B
$$
est un homomorphisme d'anneaux parfait, c'est-à-dire si $B$ est de dimension de Tor finie sur
$B \otimes_A B$, alors $B$ est une $A$-algèbre lisse.
\end{lemma}

\begin{proof}
Cela résulte du lemme \ref{dpa-lemma-local-perfect-diagonal} et des
propriétés générales des homomorphismes d'anneaux lisses ;
voir Algèbre, lemmes \ref{algebra-lemma-smooth-at-point} et
\ref{algebra-lemma-locally-smooth}.
On peut aussi modifier légèrement la démonstration du
lemme \ref{dpa-lemma-local-perfect-diagonal} pour démontrer le
présent lemme.
\end{proof}






\section{Liberté du module conormal}
\label{dpa-section-freeness-conormal}

\noindent
Les résolutions de Tate et les dérivations sur celles-ci permettent de
démontrer des versions (plus fortes) des résultats de cette section ; voir \cite{Iyengar}.
Deux références plus élémentaires sont
\cite{Vasconcelos} et \cite{Ferrand-lci}.

\begin{lemma}
\label{dpa-lemma-free-summand-in-ideal-finite-proj-dim}
\begin{reference}
\cite{Vasconcelos}
\end{reference}
Soit $R$ un anneau local noethérien. Soit $I \subset R$ un idéal
de dimension projective finie sur $R$. Si $F \subset I/I^2$ est
un facteur direct isomorphe à $R/I$, alors il existe un élément non diviseur de zéro
$x \in I$ tel que l'image de $x$ dans $I/I^2$ engendre $F$.
\end{lemma}

\begin{proof}
Par hypothèse, on peut choisir une résolution libre finie
$$
0 \to R^{\oplus n_e} \to R^{\oplus n_{e-1}} \to\ldots\to
R^{\oplus n_1} \to R \to R/I \to 0
$$
Alors $\varphi_1 : R^{\oplus n_1} \to R$ est de rang $1$,
et on voit que $I$ contient un élément non diviseur de zéro $y$,
d'après Algèbre, proposition \ref{algebra-proposition-what-exact}.
Soient $\mathfrak p_1, \ldots, \mathfrak p_n$ les idéaux premiers
associés à $R$; voir Algèbre, lemme \ref{algebra-lemma-finite-ass}.
Soit $I^2 \subset J \subset I$ un idéal tel que $J/I^2 = F$.
Alors $J \not\subset\mathfrak p_i$ pour tout $i$,
car $y^2 \in J$ et $y^2 \not\in\mathfrak p_i$;
voir Algèbre, lemme \ref{algebra-lemma-ass-zero-divisors}.
D'après le lemme de Nakayama (Algèbre, lemme \ref{algebra-lemma-NAK}),
on a $J \not\subset\mathfrak m J + I^2$.
D'après Algèbre, lemme \ref{algebra-lemma-silly},
on peut choisir $x \in J$, $x \not\in\mathfrak m J + I^2$ et
$x \not\in\mathfrak p_i$ pour $i = 1, \ldots, n$.
Alors $x$ est un élément non diviseur de zéro, et l'image
de $x$ dans $I/I^2$ engendre (d'après le lemme de Nakayama) le
facteur direct $J/I^2 \cong R/I$.
\end{proof}

\begin{lemma}
\label{dpa-lemma-vasconcelos}
\begin{reference}
Version locale de \cite[Theorem 1.1]{Vasconcelos}
\end{reference}
Soit $R$ un anneau local noethérien. Soit $I \subset R$ un idéal
de dimension projective finie sur $R$. Si $F \subset I/I^2$
est un facteur direct libre de rang $r$, alors il existe une suite régulière
$x_1, \ldots, x_r \in I$ telle que $x_1 \bmod I^2, \ldots, x_r \bmod I^2$
engendrent $F$.
\end{lemma}

\begin{proof}
Si $r = 0$, il n'y a rien à démontrer. Supposons $r > 0$. On peut choisir
$x \in I$ tel que $x$ soit un élément non diviseur de zéro et que $x \bmod I^2$
engendre un facteur direct de $F$ isomorphe à $R/I$; voir
le lemme \ref{dpa-lemma-free-summand-in-ideal-finite-proj-dim}.
Considérons l'anneau $R' = R/(x)$ et l'idéal $I' = I/(x)$.
Bien entendu, $R'/I' = R/I$. La suite exacte courte
$$
0 \to R/I \xrightarrow{x} I/xI \to I' \to 0
$$
est scindée, car l'homomorphisme $I/xI \to I/I^2$ envoie $xR/xI$
sur un facteur direct. Or $I/xI = I \otimes_R^\mathbf{L} R'$ est
de dimension projective finie sur $R'$; voir
Compléments sur l'algèbre, lemmes \ref{more-algebra-lemma-perfect-module} et
\ref{more-algebra-lemma-pull-perfect}. Par
conséquent, le facteur direct $I'$ est de dimension projective finie sur $R'$.
D'autre part, on a la suite exacte courte
$0 \to xR/xI \to I/I^2 \to I'/(I')^2 \to 0$, et on en déduit que
$I'/(I')^2$ possède comme facteur direct libre $F' = F/(R/I \cdot x)$
de rang $r - 1$. Par récurrence sur $r$, on
peut choisir une suite régulière $x'_2, \ldots, x'_r \in I'$
dont les classes engendrent librement $F'$.
Si $x_1 = x$ et si $x_2, \ldots, x_r$ sont des relèvements quelconques de
$x'_1, \ldots, x'_r$ dans $R$, on voit alors que le lemme est vérifié.
\end{proof}

\begin{proposition}
\label{dpa-proposition-regular-ideal}
\begin{reference}
Variante de \cite[Corollary 1]{Vasconcelos}. Voir aussi
\cite{Iyengar} et \cite{Ferrand-lci}.
\end{reference}
Soit $R$ un anneau noethérien. Soit $I \subset R$ un idéal
de dimension projective finie tel que $I/I^2$ soit
localement libre de type fini sur $R/I$. Alors $I$ est un idéal régulier
(Compléments sur l'algèbre, définition \ref{more-algebra-definition-regular-ideal}).
\end{proposition}

\begin{proof}
D'après Algèbre, lemme \ref{algebra-lemma-regular-sequence-in-neighbourhood},
il suffit de montrer que $I_\mathfrak p \subset R_\mathfrak p$ est engendré
par une suite régulière pour tout $\mathfrak p \supset I$. On peut donc
supposer que $R$ est local. Si $I/I^2$ est de rang $r$, alors le
lemme \ref{dpa-lemma-vasconcelos} fournit une suite régulière
$x_1, \ldots, x_r \in I$ qui engendre $I/I^2$.
D'après le lemme de Nakayama (Algèbre, lemme \ref{algebra-lemma-NAK}),
on en déduit que $I$ est engendré par $x_1, \ldots, x_r$.
\end{proof}

\noindent
Pour tout homomorphisme d'anneaux d'intersection complète locale $A \to B$,
le complexe cotangent naïf $\NL_{B/A}$ est parfait
d'amplitude de Tor dans $[-1, 0]$; voir
Compléments sur l'algèbre, lemme \ref{more-algebra-lemma-lci-NL}. À
l'aide de ce qui précède, on peut montrer que cela
caractérise parfois les homomorphismes d'intersection complète locale.

\begin{lemma}
\label{dpa-lemma-perfect-NL-lci}
Soit $A \to B$ un homomorphisme d'anneaux parfait (Compléments sur l'algèbre, définition
\ref{more-algebra-definition-pseudo-coherent-perfect})
entre anneaux noethériens. Les assertions suivantes sont alors équivalentes :
\begin{enumerate}
\item $\NL_{B/A}$ est d'amplitude de Tor dans $[-1, 0]$,
\item $\NL_{B/A}$ est un objet parfait de $D(B)$
d'amplitude de Tor dans $[-1, 0]$, et
\item $A \to B$ est une intersection complète locale
(Compléments sur l'algèbre, définition
\ref{more-algebra-definition-local-complete-intersection}).
\end{enumerate}
\end{lemma}

\begin{proof}
Écrivons $B = A[x_1, \ldots, x_n]/I$. Alors $\NL_{B/A}$ est représenté par
le complexe
$$
I/I^2 \longrightarrow\bigoplus B \text{d}x_i
$$
de $B$-modules dans lequel $I/I^2$ est placé en degré $-1$. Puisque le terme
de degré $0$ est libre de type fini, ce complexe est d'amplitude de Tor dans $[-1, 0]$ si et
seulement si $I/I^2$ est un $B$-module plat ;
voir Compléments sur l'algèbre, lemme \ref{more-algebra-lemma-last-one-flat}.
Comme $I/I^2$ est un $B$-module fini et que $B$ est noethérien, cela
équivaut à ce que $I/I^2$ soit un $B$-module
localement libre de type fini (Algèbre, lemme \ref{algebra-lemma-finite-projective}).
L'équivalence de (1) et (2) est donc claire. En outre, l'équivalence
de (1) et (3) résulte également de la
proposition \ref{dpa-proposition-regular-ideal} (et
du fait qu'un idéal régulier est à la fois un idéal régulier au sens
de Koszul et un idéal quasi-régulier ; voir
Compléments sur l'algèbre, section \ref{more-algebra-section-ideals}).
\end{proof}

\begin{lemma}
\label{dpa-lemma-flat-fp-NL-lci}
Soit $A \to B$ un homomorphisme d'anneaux plat et de présentation finie.
Les assertions suivantes sont alors équivalentes :
\begin{enumerate}
\item $\NL_{B/A}$ est d'amplitude de Tor dans $[-1, 0]$,
\item $\NL_{B/A}$ est un objet parfait de $D(B)$
d'amplitude de Tor dans $[-1, 0]$,
\item $A \to B$ est
syntomique (Algèbre, définition \ref{algebra-definition-lci}), et
\item $A \to B$ est une intersection complète locale
(Compléments sur l'algèbre, définition
\ref{more-algebra-definition-local-complete-intersection}).
\end{enumerate}
\end{lemma}

\begin{proof}
L'équivalence de (3) et (4) est établie dans Compléments sur l'algèbre, lemme
\ref{more-algebra-lemma-syntomic-lci}.

\medskip\noindent
Si $A \to B$ est syntomique, on peut trouver un diagramme cocartésien
$$
\xymatrix{
B_0 \ar[r] & B \\
A_0 \ar[r] \ar[u] & A \ar[u]
}
$$
tel que $A_0 \to B_0$ soit syntomique et que $A_0$ soit noethérien ;
voir Algèbre, lemmes \ref{algebra-lemma-limit-module-finite-presentation} et
\ref{algebra-lemma-colimit-lci}. D'après le lemme \ref{dpa-lemma-perfect-NL-lci},
on voit que $\NL_{B_0/A_0}$ est parfait d'amplitude de Tor dans $[-1, 0]$.
D'après Compléments sur l'algèbre, lemme \ref{more-algebra-lemma-base-change-NL-flat},
on en déduit qu'il en va de même de
$\NL_{B/A} = \NL_{B_0/A_0} \otimes_{B_0}^\mathbf{L} B$ (voir
aussi Compléments sur l'algèbre, lemmes \ref{more-algebra-lemma-pull-tor-amplitude} et
\ref{more-algebra-lemma-pull-perfect}).
Ceci montre que (3) implique (2).

\medskip\noindent
Supposons (1). D'après Compléments sur l'algèbre, lemme
\ref{more-algebra-lemma-base-change-NL-flat}, pour tout
homomorphisme d'anneaux $A \to k$, où
$k$ est un corps, on voit que $\NL_{B \otimes_A k/k}$ est
d'amplitude de Tor dans $[-1, 0]$ (voir
Compléments sur l'algèbre, lemme \ref{more-algebra-lemma-pull-tor-amplitude}).
Le lemme \ref{dpa-lemma-perfect-NL-lci} montre donc que $k \to B \otimes_A k$ est
un homomorphisme d'intersection complète locale. Ainsi, $A \to B$
est syntomique par définition. Ceci montre que (1) implique (3)
et achève la démonstration.
\end{proof}





\section{Complexes de Koszul et résolutions de Tate}
\label{dpa-section-koszul-vs-tate}

\noindent
Dans cette section, nous « relevons » le résultat de
Compléments sur l'algèbre, lemme \ref{more-algebra-lemma-sequence-Koszul-complexes}, à la
catégorie des algèbres différentielles graduées munies de puissances divisées
compatibles avec la structure différentielle graduée (attention : nous représentons ici
les complexes de Koszul comme des complexes de chaînes, tandis que dans la référence
citée ils sont considérés comme des complexes de cochaînes).

\medskip\noindent
Soit $R$ un anneau. Soit $I \subset R$ l'idéal engendré
par $f_1, \ldots, f_r \in R$. Pour $n \geq 1$, nous notons
$$
K_n = K_{n, \bullet} = R\langle\xi_1, \ldots, \xi_r\rangle
$$
l'algèbre différentielle graduée de Koszul, où $\xi_i$ est de degré $1$ et
$\text{d}(\xi_i) = f_i^n$. Elle porte une unique structure de puissances divisées (au
sens de la définition \ref{dpa-definition-divided-powers-dga}) ; voir
l'exemple \ref{dpa-example-adjoining-odd}. Pour $m > n$, l'application de transition
$K_m \to K_n$ est l'homomorphisme d'algèbres différentielles graduées compatible
avec les puissances divisées qui envoie $\xi_i$ sur $f_i^{m - n}\xi_i$.

\begin{lemma}
\label{dpa-lemma-lift-tate-to-koszul}
Dans la situation ci-dessus, si $R$ est noethérien,
alors, pour tout $n$, il existe un entier $N \geq n$ et des applications
$$
K_N \to A \to R/(f_1^N, \ldots, f_r^N)\quad\text{et}\quad A \to K_n
$$
avec les propriétés suivantes
\begin{enumerate}
\item $(A, \text{d}, \gamma)$ est comme
dans la définition \ref{dpa-definition-divided-powers-dga},
\item $A \to R/(f_1^N, \ldots, f_r^N)$ est un quasi-isomorphisme,
\item la composition $K_N \to A \to R/(f_1^N, \ldots, f_r^N)$
est l'application canonique,
\item la composition $K_N \to A \to K_n$ est l'application de transition,
\item $A_0 = R \to R/(f_1^N, \ldots, f_r^N)$ est la
surjection canonique,
\item $A$ est une algèbre polynomiale graduée à puissances divisées sur $R$
ayant un nombre fini de générateurs dans chaque degré ;
\item $A \to K_n$ est un homomorphisme d'algèbres différentielles graduées sur $R$,
compatible avec les puissances divisées, qui induit l'application canonique
$R/(f_1^N, \ldots, f_r^N) \to R/(f_1^n, \ldots, f_r^n)$ en
homologie de degré $0$.
\end{enumerate}
La condition (6) signifie que $A$ est construite à partir de $A_0$ en
adjoignant successivement une famille finie de variables $T$ dans chaque degré
$> 0$, comme dans l'exemple \ref{dpa-example-adjoining-odd} ou \ref{dpa-example-adjoining-even}.
\end{lemma}

\begin{proof}
Fixons $n$. Si $r = 0$, nous pouvons simplement prendre $A = R$. Supposons $r > 0$.
D'après Compléments sur l'algèbre, lemme \ref{more-algebra-lemma-sequence-Koszul-complexes}
(traduit dans le langage des complexes de chaînes), nous pouvons choisir
$$
n_{r} > n_{r - 1} > \ldots > n_1 > n_0 = n
$$
de telle sorte que les applications de transition $K_{n_{i + 1}} \to K_{n_i}$ entre
algèbres de Koszul (définies ci-dessus) induisent l'application nulle en homologie dans les degrés $> 0$.
Nous démontrerons le lemme en prenant $N = n_r$.

\medskip\noindent
Nous construisons $A$ exactement comme dans l'énoncé et la démonstration
du lemme \ref{dpa-lemma-tate-resolution}. Nous aurons donc
$$
A = \colim A(m),\quad\text{et}\quad
A(0) \to A(1) \to A(2) \to\ldots\to R/(f_1^N, \ldots, f_r^N)
$$
On obtient immédiatement les propriétés (1), (2), (5) et (6). Pour achever la
démonstration, nous construirons les homomorphismes de $R$-algèbres
$K_N \to A \to K_n$. Pour ce faire, nous allons construire
\begin{enumerate}
\item un isomorphisme $A(1) \to K_N = K_{n_r}$,
\item une application $A(2) \to K_{n_{r - 1}}$,
\item $\ldots$
\item une application $A(r) \to K_{n_1}$,
\item une application $A(r + 1) \to K_{n_0} = K_n$,
\item une application $A \to K_n$.
\end{enumerate}
À chacune de ces étapes, l'application construite sera un
homomorphisme entre algèbres différentielles graduées munies de puissances
divisées compatibles ; elle sera compatible avec la
précédente par l'intermédiaire des applications de transition entre
algèbres de Koszul et induira l'application canonique évidente en
homologie de degré $0$.

\medskip\noindent
Rappelons que $A(0) = R$. Pour $m = 1$, la démonstration
du lemme \ref{dpa-lemma-tate-resolution}
prend $A(1) = R\langle T_1, \ldots, T_r\rangle$, où
$T_i$ est de degré $1$ et $\text{d}(T_i) = f_i^N$.
En effet, les $f_i^N$ engendrent le noyau de
$A(0) \to R/(f_1^N, \ldots, f_r^N)$.
Pour définir $A(1) \to K_N = K_{n_r}$, nous utilisons donc l'application
$$
\varphi_1 : A(1) \longrightarrow K_{n_r},\quad T_i \longmapsto\xi_i
$$
qui est un isomorphisme.

\medskip\noindent
Pour $m = 2$, la construction dans la démonstration du lemme \ref{dpa-lemma-tate-resolution}
choisit des générateurs $e_1, \ldots, e_t \in\Ker(\text{d} : A(1)_1 \to A(1)_0)$.
Elle prend
$A(2) = A(1)\langle T_1, \ldots, T_t\rangle$
comme algèbre polynomiale à puissances divisées, avec $T_i$ de degré $2$
et $\text{d}(T_i) = e_i$.
Puisque $\varphi_1(e_i)$ est un cycle de $K_{n_r}$,
son image dans $K_{n_{r - 1}}$ est un bord, par
le choix de $n_r$ et $n_{r - 1}$ effectué
ci-dessus. Nous pouvons donc construire le diagramme commutatif
$$
\xymatrix{
A(1) \ar[d] \ar[r]_{\varphi_1} & K_{n_r} \ar[d] \\
A(2) \ar[r]^{\varphi_2} & K_{n_{r - 1}}
}
$$
en envoyant $T_i$ sur un élément de degré $2$ dont le bord est
l'image de $\varphi_1(e_i)$. L'application $\varphi_2$ existe et est compatible avec
la différentielle et les puissances divisées grâce à la propriété
universelle de l'algèbre polynomiale à puissances divisées.

\medskip\noindent
L'algèbre $A(m)$ et l'application $\varphi_m : A(m) \to K_{n_{r + 1 - m}}$
se construisent de la même manière pour $m = 2, \ldots, r$.

\medskip\noindent
Étant donnée l'application $A(r) \to K_{n_1}$, la composée
$H_r(A(r)) \to H_r(K_{n_1}) \to H_r(K_{n_0}) \subset (K_{n_0})_r$
est nulle. On peut donc la prolonger en une application $A(r + 1) \to K_{n_0} = K_n$
en envoyant sur zéro les nouveaux générateurs polynomiaux de $A(r + 1)$.

\medskip\noindent
Une fois construite l'application $A(r + 1) \to K_{n_0} = K_n$, on la prolonge de
l'unique manière possible à $A(r + 2), A(r + 3), \ldots$, en envoyant les
nouveaux générateurs polynomiaux sur zéro. Ceci
achève la démonstration.
\end{proof}

\begin{remark}
\label{dpa-remark-pro-system-koszul}
Dans la situation ci-dessus, si $R$ est noethérien,
nous pouvons choisir par récurrence une
suite d'entiers $1 = n_0 < n_1 < n_2 < \ldots$ telle que,
pour $i = 1, 2, 3, \ldots$, il existe des applications
$K_{n_i} \to A_i \to R/(f_1^{n_i}, \ldots, f_r^{n_i})$
et $A_i \to K_{n_{i - 1}}$ comme dans le lemme \ref{dpa-lemma-lift-tate-to-koszul}.
Notons $A_{i + 1} \to A_i$ la composée $A_{i + 1} \to K_{n_i} \to A_i$.
Le diagramme
$$
\xymatrix{
K_{n_1} \ar[d] &
K_{n_2} \ar[d] \ar[l] &
K_{n_3} \ar[d] \ar[l] &
\ldots\ar[l] \\
A_1 \ar[d] &
A_2 \ar[l] \ar[d] &
A_3 \ar[l] \ar[d] &
\ldots\ar[l] \\
K_1 &
K_{n_1} \ar[l] &
K_{n_2} \ar[l] &
\ldots\ar[l]
}
$$
est alors commutatif. On voit ainsi que les systèmes projectifs
$(K_n)$ et $(A_n)$ sont pro-isomorphes dans la catégorie des algèbres
différentielles graduées sur $R$ munies de puissances divisées compatibles.
\end{remark}

\begin{lemma}
\label{dpa-lemma-compute-cohomology-adjoin-variable}
Soient $(A, \text{d}, \gamma)$, $d \geq 1$, $f \in A_{d - 1}$ et
$A\langle T \rangle$ comme dans le lemme \ref{dpa-lemma-extend-differential}.
\begin{enumerate}
\item Si $d = 1$, il existe une suite exacte longue
$$
\ldots\to H_0(A) \xrightarrow{f} H_0(A) \to H_0(A\langle T \rangle) \to 0
$$
\item Pour $d = 2$, il existe une suite spectrale bornée
$(E_1)_{i, j} = H_{j - i}(A) \cdot T^{[i]}$
qui converge vers $H_{i + j}(A\langle T \rangle)$. La différentielle
$(d_1)_{i, j} : H_{j - i}(A) \cdot T^{[i]} \to
H_{j - i + 1}(A) \cdot T^{[i - 1]}$
envoie $\xi\cdot T^{[i]}$ sur la classe de $f \xi\cdot T^{[i - 1]}$.
\item On pourra compléter cet énoncé pour les autres degrés, si nécessaire.
\end{enumerate}
\end{lemma}

\begin{proof}
Pour $d = 1$, nous avons une suite exacte courte de complexes
$$
0 \to A \to A\langle T \rangle\to A \cdot T \to 0
$$
et (1) en résulte immédiatement. Pour $d = 2$, nous considérons
$A\langle T \rangle$ comme un complexe de chaînes filtré par les sous-complexes
$$
F^pA\langle T \rangle = \bigoplus\nolimits_{i \leq p} A \cdot T^{[i]}
$$
En appliquant la suite spectrale de
Homologie, section \ref{homology-section-filtered-complex}
(traduite en complexes de chaînes), on obtient (2).
\end{proof}

\noindent
Le lemme suivant sera nécessaire plus tard.

\begin{lemma}
\label{dpa-lemma-construct-some-maps}
Dans la situation ci-dessus, pour tous $n \geq t \geq 1$ il existe un $N > n$
et une application
$$
K_t \longrightarrow K_n \otimes_R K_t
$$
dans la catégorie dérivée des $K_N$-modules
différentiels gradués à gauche, dont la composée avec l'application de multiplication est l'application de
transition (dans l'un ou l'autre sens).
\end{lemma}

\begin{proof}
Commençons par le cas $r = 1$. Posons $f = f_1$.
Écrivons $K_t = R\langle x \rangle$,
$K_n = R\langle y \rangle$, $K_N = R\langle z \rangle$
où $x$, $y$ et $z$ sont de degré $1$, et
$\text{d}(x) = f^t$, $\text{d}(y) = f^n$,
$\text{d}(z) = f^N$. Pour tout $N > t$, nous affirmons qu'il existe un quasi-isomorphisme
$$
B_{N, t} = R\langle x, z, u \rangle
\longrightarrow
K_t,\quad
x \mapsto x,\quad
z \mapsto f^{N - t}x,\quad
u \mapsto 0
$$
Ici, le membre de gauche désigne l'algèbre polynomiale à puissances divisées
aux variables $x$ et $z$ de degré $1$ et $u$ de degré $2$, munie de
$\text{d}(x) = f^t$, $\text{d}(z) = f^N$,
$\text{d}(u) = z - f^{N - t}x$. Pour démontrer
l'assertion, remarquons que les trois sous-modules suivants de
$H_*(R\langle x, z\rangle)$ sont les mêmes
\begin{enumerate}
\item le noyau de $H_*(R\langle x, z\rangle) \to H_*(K_t)$,
\item l'image de
$z - f^{N - t}x : H_*(R\langle x, z\rangle) \to H_*(R\langle x, z\rangle)$,
\item le noyau de
$z - f^{N - t}x : H_*(R\langle x, z\rangle) \to H_*(R\langle x, z\rangle)$.
\end{enumerate}
Cette observation résulte d'un calcul direct\footnote{Indication : en posant
$z' = z - f^{N - t}x$, on voit que
$R\langle x, z\rangle = R\langle x, z'\rangle$ avec $\text{d}(z') = 0$
et que l'application $R\langle x, z'\rangle\to K_t$ est
l'application qui annule $z'$.}, détail que nous omettons. On peut alors
appliquer le lemme \ref{dpa-lemma-compute-cohomology-adjoin-variable},
partie (2), pour conclure.

\medskip\noindent
Au moyen de l'homomorphisme $K_N \to B_{N, t}$ d'algèbres différentielles graduées sur $R$ qui
envoie $z$ sur $z$, nous pouvons considérer $B_{N, t} \to K_t$ comme un
quasi-isomorphisme de $K_N$-modules différentiels gradués à gauche. Pour définir la
flèche de l'énoncé, nous utilisons l'homomorphisme
$$
B_{N, t} = R\langle x, z, u \rangle\to K_n \otimes_R K_t,\quad
x \mapsto 1 \otimes x,\quad
z \mapsto f^{N - n}y \otimes 1,\quad
u \mapsto - f^{N - n - t}y \otimes x
$$
Cette définition a un sens dès que $N \geq n + t$. Un calcul
simple montre que, dans l'un ou l'autre ordre de composition, la composée avec
l'application de multiplication est l'application de transition.

\medskip\noindent
Pour $r > 1$, on écrit chacune des algèbres de Koszul comme un
produit tensoriel d'algèbres de Koszul en $1$ variable, puis on
applique la construction précédente. Autrement dit,
$$
K_t = R\langle x_1, \ldots, x_r\rangle =
R\langle x_1\rangle\otimes_R \ldots\otimes_R R\langle x_r\rangle
$$
où $x_i$ est de degré $1$ et $\text{d}(x_i) = f_i^t$.
Lorsque $r > 1$, nous utilisons alors
$$
B_{N, t} =
R\langle x_1, z_1, u_1 \rangle
\otimes_R \ldots\otimes_R
R\langle x_r, z_r, u_r \rangle
$$
où $x_i, z_i$ sont de degré $1$ et $u_i$ est de degré $2$, avec
$\text{d}(x_i) = f_i^t$, $\text{d}(z_i) = f_i^N$,
$\text{d}(u_i) = z_i - f_i^{N - t}x_i$.
L'application tensorielle $B_{N, t} \to K_t$ est un
quasi-isomorphisme, car c'est un produit tensoriel de quasi-isomorphismes
entre complexes de $R$-modules libres bornés inférieurement. Enfin, nous définissons l'application
$$
B_{N, t}
\to
K_n \otimes_R K_t =
R\langle y_1, \ldots, y_r\rangle\otimes_R
R\langle x_1, \ldots, x_r\rangle
$$
comme le produit tensoriel des applications construites dans le cas $r = 1$,
ou directement par les règles $x_i \mapsto 1 \otimes x_i$,
$z_i \mapsto f_i^{N - n}y_i \otimes 1$,
$u_i \mapsto - f_i^{N - n - t}y_i \otimes x_i$, ce qui a
un sens dès que $N \geq n + t$. Nous omettons les détails.
\end{proof}












% OK, start here.
%

\chapter{Faisceaux différentiels gradués}



\phantomsection
\label{sdga-section-phantom}


\section{Introduction}
\label{sdga-section-introduction}

\noindent
Ce chapitre poursuit la discussion commencée dans
Algèbre différentielle graduée, section \ref{dga-section-introduction}.
On trouvera un article de synthèse dans \cite{Keller-survey}.




\section{Conventions}
\label{sdga-section-conventions}

\noindent
Dans ce chapitre, nous conservons la convention selon laquelle {\it anneau} désigne un
anneau commutatif possédant un élément unité $1$. Si $R$ est un anneau, une {\it $R$-algèbre $A$}
sera un $R$-module $A$ muni d'une application $R$-bilinéaire $A \times A \to A$
(la multiplication) telle que la multiplication soit associative et admette un
élément unité.
Autrement dit, ce sont des $R$-algèbres associatives unitaires
telles que l'homomorphisme $R \to A$ définissant la structure d'algèbre a son image dans le centre de $A$.








\section{Faisceaux d'algèbres graduées}
\label{sdga-section-ga}

\noindent
On pourra omettre cette section.

\begin{definition}
\label{sdga-definition-ga}
Soit $(\mathcal{C}, \mathcal{O})$ un site annelé. Un
{\it faisceau d'$\mathcal{O}$-algèbres graduées}
ou un {\it faisceau d'algèbres graduées} sur $(\mathcal{C}, \mathcal{O})$
est la donnée d'une famille $\mathcal{A}^n$ indexée par $n \in \mathbf{Z}$
d'$\mathcal{O}$-modules munis d'applications $\mathcal{O}$-bilinéaires
$$
\mathcal{A}^n \times \mathcal{A}^m \to \mathcal{A}^{n + m},\quad
(a, b) \longmapsto ab
$$
appelées applications de multiplication et vérifiant les propriétés suivantes :
\begin{enumerate}
\item la multiplication est associative, et
\item il existe une section globale $1$ de $\mathcal{A}^0$
qui est un élément unité bilatère pour la multiplication.
\end{enumerate}
Nous notons souvent une telle structure $\mathcal{A}$.
Un {\it homomorphisme d'$\mathcal{O}$-algèbres graduées}
$f : \mathcal{A} \to \mathcal{B}$ est une famille d'homomorphismes
$f^n : \mathcal{A}^n \to \mathcal{B}^n$
d'$\mathcal{O}$-modules compatibles avec les applications de multiplication.
\end{definition}

\noindent
Étant donnés une $\mathcal{O}$-algèbre graduée $\mathcal{A}$
et un objet $U \in \Ob(\mathcal{C})$, nous employons la
notation $$
\mathcal{A}(U)
= \Gamma(U, \mathcal{A})
= \bigoplus\nolimits_{n \in \mathbf{Z}} \mathcal{A}^n(U)
$$
C'est une $\mathcal{O}(U)$-algèbre graduée.

\begin{remark}
\label{sdga-remark-functoriality-ga}
Soit $(f, f^\sharp) : (\Sh(\mathcal{C}), \mathcal{O}_\mathcal{C})
\to (\Sh(\mathcal{D}), \mathcal{O}_\mathcal{D})$ un
morphisme de topos annelés. Nous avons :
\begin{enumerate}
\item Soit $\mathcal{A}$ une $\mathcal{O}_\mathcal{C}$-algèbre graduée. Les
applications de multiplication de $\mathcal{A}$ induisent des applications
$f_*\mathcal{A}^n \times f_*\mathcal{A}^m \to f_*\mathcal{A}^{n + m}$ que,
grâce à $f^\sharp$, nous pouvons considérer comme $\mathcal{O}_\mathcal{D}$-bilinéaires.
Nous noterons $f_*\mathcal{A}$ la
$\mathcal{O}_\mathcal{D}$-algèbre graduée ainsi obtenue.
\item Soit $\mathcal{B}$ une
$\mathcal{O}_\mathcal{D}$-algèbre graduée. Les
applications de multiplication de $\mathcal{B}$ induisent des applications
$f^*\mathcal{B}^n \times f^*\mathcal{B}^m \to f^*\mathcal{B}^{n + m}$ que,
grâce à $f^\sharp$, nous pouvons considérer comme $\mathcal{O}_\mathcal{C}$-bilinéaires.
Nous noterons $f^*\mathcal{B}$
la $\mathcal{O}_\mathcal{C}$-algèbre graduée ainsi obtenue.
\item L'ensemble des homomorphismes $f^*\mathcal{B} \to \mathcal{A}$
de $\mathcal{O}_\mathcal{C}$-algèbres graduées est en
correspondance biunivoque ($1$ à $1$) avec l'ensemble des homomorphismes
$\mathcal{B} \to f_*\mathcal{A}$ de $\mathcal{O}_\mathcal{C}$-algèbres graduées.
\end{enumerate} La
partie (3) résulte immédiatement de l'adjonction usuelle entre $f^*$ et $f_*$ sur les
faisceaux de modules.
\end{remark}




\section{Faisceaux de modules gradués}
\label{sdga-section-graded-modules}

\noindent
On pourra omettre cette section.

\begin{definition}
\label{sdga-definition-gm} Soit
$(\mathcal{C}, \mathcal{O})$ un site annelé. Soit
$\mathcal{A}$ un faisceau d'algèbres graduées sur
$(\mathcal{C}, \mathcal{O})$. Un
{\it $\mathcal{A}$-module gradué (à droite)}, ou un
{\it module gradué (à droite)} sur $\mathcal{A}$, est la
donnée d'une famille $\mathcal{M}^n$, indexée par $n \in \mathbf{Z}$,
de $\mathcal{O}$-modules munis d'applications
$\mathcal{O}$-bilinéaires
$$
\mathcal{M}^n \times \mathcal{A}^m \to \mathcal{M}^{n + m},\quad (x,
a) \longmapsto xa
$$
appelées applications de multiplication et vérifiant les propriétés
suivantes : \begin{enumerate}
\item la multiplication vérifie $(xa)a' = x(aa')$,
\item la section unité $1$ de $\mathcal{A}^0$
agit comme l'identité sur $\mathcal{M}^n$ pour tout $n$.
\end{enumerate}
Nous dirons souvent « soit $\mathcal{M}$ un $\mathcal{A}$-module gradué »
pour désigner cette situation.
Un {\it homomorphisme de $\mathcal{A}$-modules gradués}
$f : \mathcal{M} \to \mathcal{N}$ est une famille d'homomorphismes
$f^n : \mathcal{M}^n \to \mathcal{N}^n$
de $\mathcal{O}$-modules, compatible avec les applications de multiplication.
La catégorie des $\mathcal{A}$-modules gradués (à
droite) est notée $\textit{Mod}(\mathcal{A})$.
\end{definition}

\noindent On
définit exactement de la même manière les {\it modules gradués à gauche}, mais, dans
ce chapitre, les modules seront par défaut des modules à droite.

\medskip\noindent
Étant donnés un $\mathcal{A}$-module gradué $\mathcal{M}$
et un objet $U \in \Ob(\mathcal{C})$, nous employons la
notation $$
\mathcal{M}(U)
= \Gamma(U, \mathcal{M})
= \bigoplus\nolimits_{n \in \mathbf{Z}} \mathcal{M}^n(U)
$$
C'est un $\mathcal{A}(U)$-module gradué (à droite).

\begin{lemma}
\label{sdga-lemma-gm-abelian} Soit
$(\mathcal{C}, \mathcal{O})$ un site annelé. Soit
$\mathcal{A}$ une $\mathcal{O}$-algèbre graduée. La catégorie
$\textit{Mod}(\mathcal{A})$ est une catégorie abélienne possédant les
propriétés suivantes :
\begin{enumerate}
\item $\textit{Mod}(\mathcal{A})$ admet les sommes directes arbitraires,
\item $\textit{Mod}(\mathcal{A})$ admet les limites inductives arbitraires,
\item les limites inductives filtrantes dans $\textit{Mod}(\mathcal{A})$ sont exactes,
\item $\textit{Mod}(\mathcal{A})$ admet les produits arbitraires,
\item $\textit{Mod}(\mathcal{A})$ admet les limites arbitraires.
\end{enumerate}
Le foncteur
$$
\textit{Mod}(\mathcal{A}) \longrightarrow \textit{Mod}(\mathcal{O}),\quad
\mathcal{M} \longmapsto \mathcal{M}^n
$$ qui
envoie un $\mathcal{A}$-module gradué sur son terme de degré $n$ commute à
toutes les limites et limites inductives.
\end{lemma}

\noindent
Le lemme affirme que les limites et limites inductives se calculent terme à terme. Il
affirme aussi (ou l'implique, si l'on préfère) que le foncteur d'oubli
$$
\textit{Mod}(\mathcal{A})
\longrightarrow
\text{les }\mathcal{O}\text{-modules gradués}
$$
commute à toutes les limites et limites inductives.

\begin{proof}
Notons
$\text{gr}^n : \textit{Mod}(\mathcal{A}) \to \textit{Mod}(\mathcal{O})$ le foncteur
de l'énoncé. Considérons un homomorphisme
$f : \mathcal{M} \to \mathcal{N}$ de
$\mathcal{A}$-modules gradués. Le noyau et le conoyau de
$f$, calculés comme morphismes de $\mathcal{O}$-modules gradués, sont eux
aussi munis d'applications de multiplication comme dans la
définition \ref{sdga-definition-gm}. Ce sont donc également
le noyau et le conoyau dans $\textit{Mod}(\mathcal{A})$.
Ainsi, $\textit{Mod}(\mathcal{A})$ est une catégorie abélienne,
et la prise des noyaux et conoyaux commute à $\text{gr}^n$.

\medskip\noindent
Pour démontrer l'existence des limites et limites inductives, il suffit de
démontrer l'existence des produits et des sommes directes, voir
Catégories, lemmes \ref{categories-lemma-limits-products-equalizers} et
\ref{categories-lemma-colimits-coproducts-coequalizers}. Les mêmes
lemmes montrent que, pour
établir que $\text{gr}^n$ commute aux limites et limites inductives, il suffit
d'établir que $\text{gr}^n$ commute aux sommes directes et aux produits.

\medskip\noindent
Soit $\mathcal{M}_t$, $t \in T$, une famille de $\mathcal{A}$-modules gradués.
Considérons le $\mathcal{A}$-module gradué dont le terme de degré $n$ est
$\bigoplus_{t \in T} \mathcal{M}_t^n$ (muni des applications de multiplication évidentes).
Le lecteur vérifie aisément qu'il s'agit d'une somme directe dans
$\textit{Mod}(\mathcal{A})$. Il en va de même des produits.

\medskip\noindent
Observons que $\text{gr}^n$ est un foncteur exact pour tout $n$ et qu'un
complexe $\mathcal{M}_1 \to \mathcal{M}_2 \to \mathcal{M}_3$ de
$\textit{Mod}(\mathcal{A})$ est exact si et seulement si
$\text{gr}^n\mathcal{M}_1 \to \text{gr}^n\mathcal{M}_2 \to
\text{gr}^n\mathcal{M}_3$ est exact dans $\textit{Mod}(\mathcal{O})$ pour
tout $n$. On en déduit (3), puisque les limites inductives filtrantes sont
exactes dans $\textit{Mod}(\mathcal{O})$ ; c'est une
catégorie abélienne de Grothendieck, voir
Cohomologie sur les sites, section \ref{sites-cohomology-section-unbounded}.
\end{proof}






\section{La catégorie graduée des faisceaux de modules gradués}
\label{sdga-section-gm-gr-cat}

\noindent
On pourra omettre cette section. Cette section est
l'analogue de Algèbre différentielle graduée, exemple \ref{dga-example-gm-gr-cat}.
Pour nos conventions sur les catégories graduées,
voir Algèbre différentielle graduée, section \ref{dga-section-graded}.

\medskip\noindent
Soit $(\mathcal{C}, \mathcal{O})$ un site annelé. Soit
$\mathcal{A}$ un faisceau d'algèbres graduées sur
$(\mathcal{C}, \mathcal{O})$. Nous allons construire une
catégorie graduée $\textit{Mod}^{gr}(\mathcal{A})$ sur
$R = \Gamma(\mathcal{C}, \mathcal{O})$ dont
la catégorie associée $(\textit{Mod}^{gr}(\mathcal{A}))^0$ est la
catégorie des $\mathcal{A}$-modules gradués. Les objets de
$\textit{Mod}^{gr}(\mathcal{A})$ sont les
$\mathcal{A}$-modules gradués à droite (voir la section
\ref{sdga-section-graded-modules}). Étant donnés des
$\mathcal{A}$-modules gradués $\mathcal{L}$ et $\mathcal{M}$, posons
$$
\Hom_{\textit{Mod}^{gr}(\mathcal{A})}(\mathcal{L}, \mathcal{M}) =
\bigoplus\nolimits_{n \in \mathbf{Z}} \Hom^n(\mathcal{L}, \mathcal{M})
$$
où
$\Hom^n(\mathcal{L}, \mathcal{M})$ est l'ensemble
des homomorphismes de $\mathcal{A}$-modules à droite
$f : \mathcal{L} \to \mathcal{M}$ homogènes
de degré $n$. Plus précisément, $f$ est la donnée d'une
famille d'homomorphismes $f : \mathcal{L}^i \to \mathcal{M}^{i + n}$, pour
$i \in \mathbf{Z}$, compatible avec les applications de multiplication.
En termes de composantes,
$$
\Hom^n(\mathcal{L}, \mathcal{M})
\subset
\prod\nolimits_{p + q = n}
\Hom_\mathcal{O}(\mathcal{L}^{-q}, \mathcal{M}^p)
$$ (remarquer l'inversion
des indices) est le sous-ensemble des
$f = (f_{p, q})$ tels que
$$
f_{p, q}(m a) = f_{p - i, q + i}(m)a
$$ pour toutes sections locales
$a$ de $\mathcal{A}^i$ et
$m$ de $\mathcal{L}^{-q - i}$. Pour des $\mathcal{A}$-modules gradués
$\mathcal{K}$, $\mathcal{L}$ et $\mathcal{M}$, nous définissons la
composition dans $\textit{Mod}^{gr}(\mathcal{A})$ au moyen des
applications
$$
\Hom^m(\mathcal{L}, \mathcal{M}) \times
\Hom^n(\mathcal{K}, \mathcal{L}) \longrightarrow
\Hom^{n + m}(\mathcal{K}, \mathcal{M})
$$ par
simple composition des applications
de $\mathcal{A}$-modules : $(g, f) \mapsto g \circ f$.





\section{Produit tensoriel de faisceaux de modules gradués}
\label{sdga-section-tensor-product}

\noindent
On pourra omettre cette section. Cette section est analogue à une partie
d'Algèbre différentielle graduée, section \ref{dga-section-tensor-product}.

\medskip\noindent
Soit $(\mathcal{C}, \mathcal{O})$ un site annelé. Soit $\mathcal{A}$ un
faisceau d'algèbres graduées sur $(\mathcal{C}, \mathcal{O})$.
Soit $\mathcal{M}$ un $\mathcal{A}$-module gradué à droite
et soit $\mathcal{N}$ un $\mathcal{A}$-module gradué à gauche.
Nous définissons alors le {\it produit tensoriel}
$\mathcal{M} \otimes_\mathcal{A} \mathcal{N}$
comme le $\mathcal{O}$-module gradué dont le terme de degré $n$
est $$
(\mathcal{M} \otimes_\mathcal{A} \mathcal{N})^n
= \Coker\left(
\bigoplus\nolimits_{r + s + t = n} \mathcal{M}^r \otimes_\mathcal{O}
\mathcal{A}^s \otimes_\mathcal{O} \mathcal{N}^t
\longrightarrow
\bigoplus\nolimits_{p + q = n} \mathcal{M}^p \otimes_\mathcal{O} \mathcal{N}^q
\right)
$$
où l'application envoie la section locale $x \otimes a \otimes y$
de $\mathcal{M}^r \otimes_\mathcal{O} \mathcal{A}^s
\otimes_\mathcal{O} \mathcal{N}^t$
sur $xa \otimes y - x \otimes ay$.
Avec cette définition,
$(\mathcal{M} \otimes_\mathcal{A} \mathcal{N})^n$ est
le faisceau associé au préfaisceau
$U \mapsto (\mathcal{M}(U) \otimes_{\mathcal{A}(U)} \mathcal{N}(U))^n$, où le
produit tensoriel des modules gradués est celui défini dans
Algèbre différentielle graduée, section \ref{dga-section-tensor-product}.

\medskip\noindent
Si l'on fixe le $\mathcal{A}$-module gradué à gauche $\mathcal{N}$,
on obtient un foncteur
$$
- \otimes_\mathcal{A} \mathcal{N} :
\textit{Mod}(\mathcal{A})
\longrightarrow
\text{Gr}(\textit{Mod}(\mathcal{O})) =
\text{les }\mathcal{O}\text{-modules gradués}
$$
Pour la notation $\text{Gr}(-)$, voir
Homologie, définition \ref{homology-definition-graded}.
La catégorie graduée des $\mathcal{O}$-modules gradués est notée
$\text{Gr}^{gr}(\textit{Mod}(\mathcal{O}))$, voir
Algèbre différentielle graduée, exemple
\ref{dga-example-graded-category-graded-objects}.
Le foncteur ci-dessus se relève en un foncteur de catégories graduées
$$
- \otimes_\mathcal{A} \mathcal{N} :
\textit{Mod}^{gr}(\mathcal{A})
\longrightarrow
\text{Gr}^{gr}(\textit{Mod}(\mathcal{O}))
$$
en envoyant les homomorphismes de degré $n$ de $\mathcal{M} \to \mathcal{M}'$
sur l'application induite de degré $n$ de
$\mathcal{M} \otimes_\mathcal{A} \mathcal{N}$ à
$\mathcal{M}' \otimes_\mathcal{A} \mathcal{N}$.






\section{Hom interne de faisceaux de modules gradués}
\label{sdga-section-internal-hom-graded}

\noindent Nous
recommandons vivement au lecteur d'omettre cette section.

\medskip\noindent
Nous aurons besoin de la version faisceautisée de la construction de
la section \ref{sdga-section-gm-gr-cat}. Soient
$(\mathcal{C}, \mathcal{O})$, $\mathcal{A}$,
$\mathcal{M}$ et $\mathcal{L}$ comme dans la section \ref{sdga-section-gm-gr-cat}. Nous
définissons
$$
\SheafHom^{gr}_\mathcal{A}(\mathcal{M}, \mathcal{L})
$$ comme
le $\mathcal{O}$-module gradué dont le terme de degré $n$
$$
\SheafHom^n_\mathcal{A}(\mathcal{M}, \mathcal{L})
\subset
\prod\nolimits_{p + q = n}
\SheafHom_\mathcal{O}(\mathcal{L}^{-q}, \mathcal{M}^p)
$$ est le
sous-faisceau constitué des sections locales $f = (f_{p, q})$ telles que
$$
f_{p, q}(m a) = f_{p - i, q + i}(m)a
$$ pour toutes
sections locales $a$ de $\mathcal{A}^i$ et
$m$ de $\mathcal{L}^{-q - i}$. Comme dans la section \ref{sdga-section-gm-gr-cat}, il
existe une application de
composition $$
\SheafHom^{gr}_\mathcal{A}(\mathcal{L}, \mathcal{M}) \otimes_\mathcal{O}
\SheafHom^{gr}_\mathcal{A}(\mathcal{K}, \mathcal{L})
\longrightarrow
\SheafHom^{gr}_\mathcal{A}(\mathcal{K}, \mathcal{M})
$$ où
le membre de gauche est le produit tensoriel de $\mathcal{O}$-modules gradués défini
dans la section \ref{sdga-section-tensor-product}. Cette
application est donnée par l'application de
composition $$
\SheafHom^m_\mathcal{A}(\mathcal{L}, \mathcal{M}) \otimes_\mathcal{O}
\SheafHom^n_\mathcal{A}(\mathcal{K}, \mathcal{L}) \longrightarrow
\SheafHom^{n + m}_\mathcal{A}(\mathcal{K}, \mathcal{M})
$$
définie par simple composition (localement).

\medskip\noindent
Avec ces définitions, nous avons
$$
\Hom_{\textit{Mod}^{gr}(\mathcal{A})}(\mathcal{L}, \mathcal{M}) =
\Gamma(\mathcal{C}, \SheafHom^{gr}_\mathcal{A}(\mathcal{L}, \mathcal{M}))
$$ comme modules
gradués sur $R$, de manière compatible à la composition.







\section{Faisceaux de bimodules gradués et adjonction tensor-Hom}
\label{sdga-section-graded-bimodules}

\noindent
On pourra omettre cette section.

\begin{definition}
\label{sdga-definition-bimodule}
Soit $(\mathcal{C}, \mathcal{O})$ un site annelé. Soient $\mathcal{A}$ et
$\mathcal{B}$ des faisceaux d'algèbres graduées sur
$(\mathcal{C}, \mathcal{O})$. Un
{\it $(\mathcal{A}, \mathcal{B})$-bimodule gradué} est la
donnée d'une famille $\mathcal{M}^n$, indexée par $n \in \mathbf{Z}$, de
$\mathcal{O}$-modules munis d'applications $\mathcal{O}$-bilinéaires
$$
\mathcal{M}^n \times \mathcal{B}^m \to \mathcal{M}^{n + m},\quad (x,
b) \longmapsto xb
$$
et
$$
\mathcal{A}^n \times \mathcal{M}^m \to \mathcal{M}^{n + m},\quad (a, x)
\longmapsto ax
$$ appelées
applications de multiplication et vérifiant les propriétés
suivantes : \begin{enumerate}
\item la multiplication vérifie $a(a'x) = (aa')x$
et $(xb)b' = x(bb')$,
\item $(ax)b = a(xb)$,
\item la section unité $1$ de $\mathcal{A}^0$ agit comme
l'identité par multiplication, et
\item la section unité $1$ de
$\mathcal{B}^0$ agit comme l'identité par multiplication.
\end{enumerate} Nous notons
souvent une telle structure $\mathcal{M}$. Un
{\it homomorphisme de $(\mathcal{A}, \mathcal{B})$-bimodules gradués}
$f : \mathcal{M} \to \mathcal{N}$ est une famille d'homomorphismes
$f^n : \mathcal{M}^n \to \mathcal{N}^n$ de
$\mathcal{O}$-modules, compatible avec les applications de multiplication.
\end{definition}

\noindent
Étant donnés un $(\mathcal{A}, \mathcal{B})$-bimodule gradué $\mathcal{M}$
et un objet $U \in \Ob(\mathcal{C})$, nous employons la
notation $$
\mathcal{M}(U)
= \Gamma(U, \mathcal{M})
= \bigoplus\nolimits_{n \in \mathbf{Z}} \mathcal{M}^n(U)
$$
C'est un $(\mathcal{A}(U), \mathcal{B}(U))$-bimodule gradué.

\medskip\noindent
Soit $(\mathcal{C}, \mathcal{O})$ un site annelé. Soient $\mathcal{A}$ et
$\mathcal{B}$ des faisceaux d'algèbres graduées sur
$(\mathcal{C}, \mathcal{O})$. Soit $\mathcal{M}$ un
$\mathcal{A}$-module gradué à droite et soit $\mathcal{N}$ un
$(\mathcal{A}, \mathcal{B})$-bimodule gradué. Dans ce cas, le produit tensoriel
gradué défini dans la section \ref{sdga-section-tensor-product}
$$
\mathcal{M} \otimes_\mathcal{A} \mathcal{N}
$$ est un
$\mathcal{B}$-module gradué à droite, muni des applications de multiplication évidentes.
Cette construction définit un foncteur et un foncteur de catégories
graduées $$
\otimes_\mathcal{A} \mathcal{N}
: \textit{Mod}(\mathcal{A})
\longrightarrow
\textit{Mod}(\mathcal{B})
\quad\text{et}\quad
\otimes_\mathcal{A} \mathcal{N} :
\textit{Mod}^{gr}(\mathcal{A})
\longrightarrow
\textit{Mod}^{gr}(\mathcal{B})
$$
en envoyant les homomorphismes de degré $n$ de $\mathcal{M} \to \mathcal{M}'$ sur
les applications induites de degré $n$
de $\mathcal{M} \otimes_\mathcal{A} \mathcal{N}$
vers $\mathcal{M}' \otimes_\mathcal{A} \mathcal{N}$.

\medskip\noindent
Soit $(\mathcal{C}, \mathcal{O})$ un site annelé. Soient $\mathcal{A}$ et
$\mathcal{B}$ des faisceaux d'algèbres graduées sur
$(\mathcal{C}, \mathcal{O})$. Soit $\mathcal{N}$ un
$(\mathcal{A}, \mathcal{B})$-bimodule gradué. Soit
$\mathcal{L}$ un $\mathcal{B}$-module gradué à droite. Dans ce
cas, le Hom interne gradué défini dans
la section \ref{sdga-section-internal-hom-graded}
$$
\SheafHom_\mathcal{B}^{gr}(\mathcal{N}, \mathcal{L})
$$ est un
$\mathcal{A}$-module gradué à droite, muni des
applications de multiplication\footnote{Selon nos conventions,
aucun signe n'intervient ici.}
$$
\SheafHom^n_\mathcal{B}(\mathcal{N}, \mathcal{L})
\times \mathcal{A}^m
\longrightarrow
\SheafHom^{n + m}_\mathcal{B}(\mathcal{N}, \mathcal{L})
$$ qui envoient
une section $f = (f_{p,q})$ de
$\SheafHom^n_\mathcal{B}(\mathcal{N}, \mathcal{L})$ sur $U$ et une
section $a$ de $\mathcal{A}^m$ sur $U$ sur la section
$f a$ de $\SheafHom^{n + m}_\mathcal{B}(\mathcal{N}, \mathcal{L})$ sur $U$ définie par la
famille d'applications
$$
\mathcal{N}^{-q - m}|_U \xrightarrow{a \cdot -}
\mathcal{N}^{-q}|_U \xrightarrow{f_{p, q}}
\mathcal{M}^p|_U
$$ Nous omettons
la vérification que cette définition est correcte. Cette construction
définit un foncteur et un foncteur de catégories graduées
$$
\SheafHom_\mathcal{B}^{gr}(\mathcal{N}, -)
: \textit{Mod}(\mathcal{B})
\longrightarrow
\textit{Mod}(\mathcal{A})
\quad\text{et}\quad
\SheafHom_\mathcal{B}^{gr}(\mathcal{N}, -) :
\textit{Mod}^{gr}(\mathcal{B})
\longrightarrow
\textit{Mod}^{gr}(\mathcal{A})
$$ en
envoyant les homomorphismes de degré $n$ de $\mathcal{L} \to \mathcal{L}'$ sur
les applications induites de degré $n$
de $\SheafHom_\mathcal{B}^{gr}(\mathcal{N}, \mathcal{L})$
vers $\SheafHom_\mathcal{B}^{gr}(\mathcal{N}, \mathcal{L}')$.

\begin{lemma}
\label{sdga-lemma-tensor-hom-adjunction-gr}
Soit $(\mathcal{C}, \mathcal{O})$ un site annelé. Soient $\mathcal{A}$ et
$\mathcal{B}$ des faisceaux d'algèbres graduées sur
$(\mathcal{C}, \mathcal{O})$. Soit $\mathcal{M}$ un
$\mathcal{A}$-module gradué à droite. Soit $\mathcal{N}$ un
$(\mathcal{A}, \mathcal{B})$-bimodule gradué. Soit $\mathcal{L}$ un
$\mathcal{B}$-module gradué à droite. Avec les conventions précédentes,
nous
avons $$
\Hom_{\textit{Mod}^{gr}(\mathcal{B})}(
\mathcal{M} \otimes_\mathcal{A} \mathcal{N}, \mathcal{L})
= \Hom_{\textit{Mod}^{gr}(\mathcal{A})}(
\mathcal{M}, \SheafHom_\mathcal{B}^{gr}(\mathcal{N}, \mathcal{L}))
$$
et
$$
\SheafHom_\mathcal{B}^{gr}(
\mathcal{M} \otimes_\mathcal{A} \mathcal{N}, \mathcal{L}) =
\SheafHom_\mathcal{A}^{gr}(
\mathcal{M}, \SheafHom_\mathcal{B}^{gr}(\mathcal{N}, \mathcal{L}))
$$
fonctoriellement en $\mathcal{M}$, $\mathcal{N}$ et $\mathcal{L}$.
\end{lemma}

\begin{proof}
Omis. Indication : cela résulte de l'interprétation des deux membres
comme applications graduées $\mathcal{A}$-bilinéaires
$\psi : \mathcal{M} \times \mathcal{N} \to \mathcal{L}$,
qui sont $\mathcal{B}$-linéaires à droite.
\end{proof}

\noindent
Soit $(\mathcal{C}, \mathcal{O})$ un site annelé. Soient $\mathcal{A}$ et
$\mathcal{B}$ des faisceaux d'algèbres graduées sur
$(\mathcal{C}, \mathcal{O})$.
Comme cas particulier de ce qui précède, supposons donné
un homomorphisme $\varphi : \mathcal{A} \to \mathcal{B}$
de $\mathcal{O}$-algèbres graduées. On obtient alors un foncteur
et un foncteur de catégories
graduées $$
\otimes_{\mathcal{A}, \varphi} \mathcal{B}
: \textit{Mod}(\mathcal{A})
\longrightarrow
\textit{Mod}(\mathcal{B})
\quad\text{et}\quad
\otimes_{\mathcal{A}, \varphi} \mathcal{B} :
\textit{Mod}^{gr}(\mathcal{A})
\longrightarrow
\textit{Mod}^{gr}(\mathcal{B})
$$
D'autre part, nous avons les foncteurs de restriction
$$
res_\varphi :
\textit{Mod}(\mathcal{B})
\longrightarrow
\textit{Mod}(\mathcal{A})
\quad\text{et}\quad
res_\varphi :
\textit{Mod}^{gr}(\mathcal{B})
\longrightarrow
\textit{Mod}^{gr}(\mathcal{A})
$$ Le
lemme précédent permet de montrer que ces foncteurs sont adjoints
l'un à l'autre (comme d'habitude pour la restriction et le changement
de base). Notons ${}_\mathcal{A}\mathcal{B}_\mathcal{B}$
le faisceau $\mathcal{B}$ considéré comme $(\mathcal{A}, \mathcal{B})$-bimodule
gradué. Pour tout $\mathcal{B}$-module gradué à droite $\mathcal{L}$,
nous
avons alors $$
\SheafHom_\mathcal{B}^{gr}({}_\mathcal{A}\mathcal{B}_\mathcal{B}, \mathcal{L})
= res_\varphi(\mathcal{L})
$$
comme $\mathcal{A}$-modules
gradués à droite. Le lemme \ref{sdga-lemma-tensor-hom-adjunction-gr}
fournit donc un isomorphisme
fonctoriel $$
\Hom_{\textit{Mod}^{gr}(\mathcal{B})}(
\mathcal{M} \otimes_{\mathcal{A}, \varphi} \mathcal{B}, \mathcal{L})
= \Hom_{\textit{Mod}^{gr}(\mathcal{A})}(
\mathcal{M}, res_\varphi(\mathcal{L}))
$$
Lorsque le contexte est clair, nous omettons généralement $\varphi$
dans cette formule. De la même manière, nous
obtenons l'égalité
$$
\SheafHom^{gr}_\mathcal{B}(
\mathcal{M} \otimes_\mathcal{A} \mathcal{B}, \mathcal{L})
= \SheafHom_\mathcal{A}^{gr}(\mathcal{M}, \mathcal{L})
$$
de $\mathcal{O}$-modules gradués.




\section{Image inverse et image directe des faisceaux de modules gradués}
\label{sdga-section-functoriality-graded}

\noindent
Nous conseillons au lecteur d'omettre cette section.

\medskip\noindent
Soit $(f, f^\sharp) : (\Sh(\mathcal{C}), \mathcal{O}_\mathcal{C})
\to (\Sh(\mathcal{D}), \mathcal{O}_\mathcal{D})$ un morphisme
de topos annelés. Soit $\mathcal{A}$ une
$\mathcal{O}_\mathcal{C}$-algèbre graduée. Soit $\mathcal{B}$ une
$\mathcal{O}_\mathcal{D}$-algèbre graduée. Supposons qu'on nous donne
une application
$$
\varphi : f^{-1}\mathcal{B} \to \mathcal{A}
$$ d'algèbres
graduées sur $f^{-1}\mathcal{O}_\mathcal{D}$. Par l'adjonction
entre restriction et extension des scalaires, cela revient à se
donner un morphisme $\varphi : f^*\mathcal{B} \to \mathcal{A}$ de
$\mathcal{O}_\mathcal{C}$-algèbres graduées ou, de manière équivalente,
$\varphi$ peut être vu comme un morphisme
$$
\varphi : \mathcal{B} \to f_*\mathcal{A}
$$
d'algèbres graduées sur $\mathcal{O}_\mathcal{D}$.
Voir la remarque \ref{sdga-remark-functoriality-ga}.

\medskip\noindent
Définissons maintenant un
foncteur $$
f_*
: \textit{Mod}(\mathcal{A})
\longrightarrow
\textit{Mod}(\mathcal{B})
$$
Étant donné un $\mathcal{A}$-module gradué $\mathcal{M}$, nous
définissons $f_*\mathcal{M}$ comme le $\mathcal{B}$-module
gradué dont le terme de degré $n$ est $f_*\mathcal{M}^n$. Comme
multiplication,
nous utilisons $$
f_*\mathcal{M}^n \times \mathcal{B}^m
\xrightarrow{(\text{id}, \varphi^m)}
f_*\mathcal{M}^n \times f_*\mathcal{A}^m
\xrightarrow{f_*\mu_{n, m}}
f_*\mathcal{M}^{n + m}
$$
où $\mu_{n, m} : \mathcal{M}^n \times \mathcal{A}^m
\to \mathcal{M}^{n + m}$ est l'application de multiplication pour $\mathcal{M}$ sur
$\mathcal{A}$. Nous utilisons ici le fait que $f_*$ commute aux produits.
La construction est clairement fonctorielle en
$\mathcal{M}$ et nous obtenons notre foncteur.

\medskip\noindent
Définissons maintenant un
foncteur $$
f^*
: \textit{Mod}(\mathcal{B})
\longrightarrow
\textit{Mod}(\mathcal{A})
$$
Nous définirons ce foncteur comme un composé de foncteurs
$$
\textit{Mod}(\mathcal{B})
\xrightarrow{f^{-1}}
\textit{Mod}(f^{-1}\mathcal{B})
\xrightarrow{ - \otimes_{f^{-1}\mathcal{B}} \mathcal{A}}
\textit{Mod}(\mathcal{A})
$$ Tout d'abord,
étant donné un $\mathcal{B}$-module gradué $\mathcal{N}$, nous définissons
$f^{-1}\mathcal{N}$ comme le $f^{-1}\mathcal{B}$-module gradué
dont le terme de degré $n$ est $f^{-1}\mathcal{N}^n$. Comme multiplication,
nous
utilisons $$
f^{-1}\nu_{n, m}
: f^{-1}\mathcal{N}^n \times f^{-1}\mathcal{B}^m
\longrightarrow
f^{-1}\mathcal{N}^{n + m}
$$
où $\nu_{n, m} : \mathcal{N}^n \times \mathcal{B}^m
\to \mathcal{N}^{n + m}$ est l'application de multiplication pour $\mathcal{N}$ sur
$\mathcal{B}$. Nous utilisons ici le fait que $f^{-1}$ commute aux produits.
La construction est clairement fonctorielle en
$\mathcal{N}$ et nous obtenons notre foncteur $f^{-1}$.
Ceci étant dit, nous pouvons utiliser la
discussion sur le produit tensoriel dans la section \ref{sdga-section-graded-bimodules}
pour définir
le foncteur $$
- \otimes_{f^{-1}\mathcal{B}} \mathcal{A}
: \textit{Mod}(f^{-1}\mathcal{B})
\longrightarrow
\textit{Mod}(\mathcal{A})
$$
Enfin, nous
posons $$
f^*\mathcal{N}
= f^{-1}\mathcal{N} \otimes_{f^{-1}\mathcal{B}, \varphi} \mathcal{A}
$$
comme déjà prédit ci-dessus.

\medskip\noindent
Les foncteurs $f_*$ et $f^*$ se relèvent naturellement en des
foncteurs de catégories graduées
$$
f_*
: \textit{Mod}^{gr}(\mathcal{A})
\longrightarrow
\textit{Mod}^{gr}(\mathcal{B})
\quad\text{et}\quad
f^* :
\textit{Mod}^{gr}(\mathcal{B})
\longrightarrow
\textit{Mod}^{gr}(\mathcal{A})
$$
qui font la même chose sur les objets sous-jacents et sont définis
par la fonctorialité des constructions sur les morphismes homogènes
de degré $n$.

\begin{lemma}
\label{sdga-lemma-adjunction-push-pull-gr} Dans
la situation ci-dessus nous
avons $$
\Hom_{\textit{Mod}^{gr}(\mathcal{B})}(
\mathcal{N}, f_*\mathcal{M})
= \Hom_{\textit{Mod}^{gr}(\mathcal{A})}(
f^*\mathcal{N}, \mathcal{M})
$$
\end{lemma}

\begin{proof}
Omis. Indications : on commence par prouver que $f^{-1}$ et $f_*$ sont adjoints
en tant que foncteurs entre $\textit{Mod}(\mathcal{B})$
et $\textit{Mod}(f^{-1}\mathcal{B})$ en utilisant l'adjonction entre
$f^{-1}$ et $f_*$ sur les faisceaux de groupes abéliens.
On utilise ensuite l'adjonction entre le changement de base et la restriction
donnée à la section \ref{sdga-section-graded-bimodules}.
\end{proof}





\section{Localisation et faisceaux de modules gradués}
\label{sdga-section-localize-graded}

\noindent
Nous conseillons au lecteur d'omettre cette section.

\medskip\noindent
Soit $(\mathcal{C}, \mathcal{O})$ un site annelé. Soit
$U \in \Ob(\mathcal{C})$ et désignons
$$ j
:
(\Sh(\mathcal{C}/U), \mathcal{O}_U)
\longrightarrow
(\Sh(\mathcal{C}), \mathcal{O})
$$ le
morphisme de localisation correspondant (Modules
sur les sites, section \ref{sites-modules-section-localize}).
Ci-dessous, nous utiliserons le fait suivant : pour les $\mathcal{O}_U$-modules
$\mathcal{M}_i$, $i = 1, 2$ et un $\mathcal{O}$-module $\mathcal{A}$,
il existe une application
canonique $$
j_!
: \Hom_{\mathcal{O}_U}(
\mathcal{M}_1 \otimes_{\mathcal{O}_U} \mathcal{A}|_U, \mathcal{M}_2)
\longrightarrow
\Hom_\mathcal{O}(
j_!\mathcal{M}_1 \otimes_\mathcal{O} \mathcal{A}, j_!\mathcal{M}_2)
$$
À savoir, nous
avons $j_!(\mathcal{M}_1 \otimes_{\mathcal{O}_U} \mathcal{A}|_U)
= j_!\mathcal{M}_1 \otimes_\mathcal{O} \mathcal{A}$ par
Modules sur les sites, lemme \ref{sites-modules-lemma-j-shriek-and-tensor}.

\medskip\noindent
Soit $\mathcal{A}$ une $\mathcal{O}$-algèbre graduée. Nous
noterons $\mathcal{A}_U$ la restriction de $\mathcal{A}$ à
$\mathcal{C}/U$, en d'autres termes, nous avons
$\mathcal{A}_U = j^*\mathcal{A} = j^{-1}\mathcal{A}$. Dans la section
\ref{sdga-section-functoriality-graded} nous avons construit
des foncteurs adjoints
$$ j_*
:
\textit{Mod}^{gr}(\mathcal{A}_U)
\longrightarrow
\textit{Mod}^{gr}(\mathcal{A})
\quad\text{et}\quad j^*
:
\textit{Mod}^{gr}(\mathcal{A})
\longrightarrow
\textit{Mod}^{gr}(\mathcal{A}_U)
$$ avec
$j^*$ adjoint à gauche à $j_*$. Nous affirmons qu'il existe en outre un foncteur
exact
$$ j_!
:
\textit{Mod}^{gr}(\mathcal{A}_U)
\longrightarrow
\textit{Mod}^{gr}(\mathcal{A})
$$ adjoint à gauche à
$j^*$. À savoir, étant donné un $\mathcal{A}_U$-module gradué
$\mathcal{M}$, nous définissons $j_!\mathcal{M}$ comme le $\mathcal{A}$-module gradué dont le terme de degré
$n$ est $j_!\mathcal{M}^n$. Comme application de multiplication, nous
utilisons
$$
j_!\mu_{n, m} :
j_!\mathcal{M}^n \times \mathcal{A}^m \to
j_!\mathcal{M}^{n + m}
$$ où
$\mu_{m, n} : \mathcal{M}^n \times \mathcal{A}^m \to \mathcal{M}^{n + m}$ est l'application
de multiplication donnée. Étant donnée une application homogène
$f : \mathcal{M} \to \mathcal{M}'$ de degré $n$ de
$\mathcal{A}_U$-modules gradués, nous obtenons une application homogène
$j_!f : j_!\mathcal{M} \to j_!\mathcal{M}'$ de degré $n$. Ainsi
nous obtenons notre foncteur.

\begin{lemma}
\label{sdga-lemma-extension-by-zero-graded} Dans
la situation ci-dessus nous
avons $$
\Hom_{\textit{Mod}^{gr}(\mathcal{A})}(
j_!\mathcal{M}, \mathcal{N})
= \Hom_{\textit{Mod}^{gr}(\mathcal{A}_U)}(
\mathcal{M}, j^*\mathcal{N})
$$
\end{lemma}

\begin{proof}
D'après la discussion
dans Modules sur les sites, section \ref{sites-modules-section-localize},
les foncteurs $j_!$ et $j^*$ sur les $\mathcal{O}$-modules sont
adjoints. Ainsi, en ne considérant que les structures de $\mathcal{O}$-modules,
nous
savons que $$
\Hom_{\text{Gr}^{gr}(\textit{Mod}(\mathcal{O}))}(
j_!\mathcal{M}, \mathcal{N})
= \Hom_{\text{Gr}^{gr}(\textit{Mod}(\mathcal{O}_U))}(
\mathcal{M}, j^*\mathcal{N})
$$
(Rappelons que $\text{Gr}^{gr}(\textit{Mod}(\mathcal{O}))$
désigne la catégorie graduée des $\mathcal{O}$-modules
gradués.) Il faut ensuite vérifier que ces identifications
envoient les homomorphismes de $\mathcal{A}$-modules du côté
gauche vers les homomorphismes de $\mathcal{A}_U$-modules du
côté droit. Pour le vérifier, étant donnés des morphismes $\mathcal{O}_U$-linéaires
$f^n : \mathcal{M}^n \to j^*\mathcal{N}^{n + d}$
correspondant à des morphismes $\mathcal{O}$-linéaires
$g^n : j_!\mathcal{M}^n \to \mathcal{N}^{n + d}$,
il
suffit de montrer que $$
\xymatrix{
\mathcal{M}^n \otimes_{\mathcal{O}_U} \mathcal{A}_U^m
\ar[r]_{f^n \otimes 1} \ar[d]
& j^*\mathcal{N}^{n + d} \otimes_{\mathcal{O}_U} \mathcal{A}_U^m \ar[d] \\
\mathcal{M}^{n + m} \ar[r]^{f^{n + m}}
& j^*\mathcal{N}^{n + m + d}
}
$$
commute
si et seulement si $$
\xymatrix{
j_!\mathcal{M}^n \otimes_\mathcal{O} \mathcal{A}^m
\ar[r]_{g^n \otimes 1} \ar[d]
& \mathcal{N}^{n + d} \otimes_\mathcal{O} \mathcal{A}_U^m \ar[d] \\
j_!\mathcal{M}^{n + m} \ar[r]^{g^{n + m}}
& \mathcal{N}^{n + m + d}
}
$$
est commutatif.
Cependant, nous savons que \begin{align*}
\Hom_{\mathcal{O}_U}(\mathcal{M}^n \otimes_{\mathcal{O}_U} \mathcal{A}_U^m,
j^*\mathcal{N}^{n + d + m})
&
= \Hom_\mathcal{O}(j_!(\mathcal{M}^n \otimes_{\mathcal{O}_U} \mathcal{A}_U^m),
\mathcal{N}^{n + d + m}) \\
&
= \Hom_\mathcal{O}(j_!\mathcal{M}^n \otimes_\mathcal{O} \mathcal{A}^m,
\mathcal{N}^{n + d + m})
\end{align*}
d'après Modules
sur les sites, lemme \ref{sites-modules-lemma-j-shriek-and-tensor}, déjà
utilisé. Nous omettons la vérification montrant que l'obstruction à la
commutativité du premier diagramme dans le premier groupe est
envoyée sur l'obstruction à la commutativité du second diagramme dans
le dernier groupe.
\end{proof}

\begin{lemma}
\label{sdga-lemma-tensor-with-extension-by-zero} Dans
la situation ci-dessus, soit $\mathcal{M}$ un
$\mathcal{A}_U$-module gradué à droite et soit $\mathcal{N}$ un
$\mathcal{A}$-module gradué à gauche. Alors
$$
j_!\mathcal{M} \otimes_\mathcal{A} \mathcal{N} =
j_!(\mathcal{M} \otimes_{\mathcal{A}_U} \mathcal{N}|_U)
$$ en tant
que $\mathcal{O}$-modules gradués, fonctoriellement en $\mathcal{M}$
et $\mathcal{N}$.
\end{lemma}

\begin{proof}
Rappelons que la composante de degré $n$ de
$j_!\mathcal{M} \otimes_\mathcal{A} \mathcal{N}$ est le conoyau de
l'application canonique
$$
\bigoplus\nolimits_{r + s + t = n}
j_!\mathcal{M}^r \otimes_\mathcal{O}
\mathcal{A}^s \otimes_\mathcal{O}
\mathcal{N}^t
\longrightarrow
\bigoplus\nolimits_{p + q = n}
j_!\mathcal{M}^p \otimes_\mathcal{O} \mathcal{N}^q
$$ Voir
la section \ref{sdga-section-tensor-product}.
Par Modules sur les sites, lemme \ref{sites-modules-lemma-j-shriek-and-tensor}
ceci est la même chose que le conoyau de
$$
\bigoplus\nolimits_{r + s + t = n}
j_!(\mathcal{M}^r \otimes_{\mathcal{O}_U}
\mathcal{A}^s|_U \otimes_{\mathcal{O}_U}
\mathcal{N}^t|_U)
\longrightarrow
\bigoplus\nolimits_{p + q = n}
j_!(\mathcal{M}^p \otimes_{\mathcal{O}_U} \mathcal{N}^q|_U)
$$ ce qui
conclut. Une preuve alternative serait de refaire l'argument de Yoneda
donné dans la démonstration du lemme cité ci-dessus.
\end{proof}








\section{Foncteurs de décalage sur les faisceaux de modules gradués}
\label{sdga-section-shift}

\noindent
Nous recommandons vivement au lecteur d'omettre cette section. Il s'avère que les
faisceaux de modules gradués sur une algèbre graduée sont un exemple du phénomène
discuté dans
Algèbre différentielle graduée, remarque \ref{dga-remark-graded-shift-functors}.

\medskip\noindent
Soit $(\mathcal{C}, \mathcal{O})$ un site annelé. Soit
$\mathcal{A}$ un faisceau d'algèbres graduées sur
$(\mathcal{C}, \mathcal{O})$. Soit
$\mathcal{M}$ un $\mathcal{A}$-module gradué. Soit $k \in \mathbf{Z}$. Nous définissons
le {\it $k$-ième décalage de} $\mathcal{M}$, noté $\mathcal{M}[k]$, comme le
$\mathcal{A}$-module gradué dont la composante de degré $n$ est donnée par
$$
(\mathcal{M}[k])^n = \mathcal{M}^{n + k}
$$
c'est-à-dire la composante de degré $(n + k)$ de $\mathcal{M}$. Pour les
applications de multiplication $$
(\mathcal{M}[k])^n \times \mathcal{A}^m
\longrightarrow
(\mathcal{M}[k])^{n + m}
$$
nous utilisons
simplement les applications de multiplication $$
\mathcal{M}^{n + k} \times \mathcal{A}^m
\longrightarrow
\mathcal{M}^{n + m + k}
$$
de $\mathcal{M}$. Il est clair que nous avons défini un foncteur $[k]$,
que nous avons $[k + l] = [k] \circ [l]$,
et que nous avons $$
\Hom_{\textit{Mod}^{gr}(\mathcal{A})}(\mathcal{L}, \mathcal{M}[k])
= \Hom_{\textit{Mod}^{gr}(\mathcal{A})}(\mathcal{L}, \mathcal{M})[k]
$$
sans intervention
de signes, fonctoriellement en $\mathcal{M}$ et $\mathcal{L}$.
Ainsi, la catégorie graduée des $\mathcal{A}$-modules gradués se retrouve
bien à partir de la catégorie ordinaire des $\mathcal{A}$-modules gradués et
des foncteurs de décalage, comme expliqué dans
Algèbre différentielle graduée, remarque \ref{dga-remark-graded-shift-functors}.

\begin{lemma}
\label{sdga-lemma-gm-grothendieck-abelian}
Soit $(\mathcal{C}, \mathcal{O})$ un site annelé. Soit
$\mathcal{A}$ une $\mathcal{O}$-algèbre graduée. La
catégorie $\textit{Mod}(\mathcal{A})$ est une catégorie abélienne de Grothendieck.
\end{lemma}

\begin{proof}
D'après le lemme \ref{sdga-lemma-gm-abelian} et la définition d'une catégorie
abélienne de Grothendieck
(Injectifs, définition \ref{injectives-definition-grothendieck-conditions}),
il suffit de
montrer que $\textit{Mod}(\mathcal{A})$ possède un générateur. Nous affirmons
que $$
\mathcal{G} = \bigoplus\nolimits_{k, U} j_{U!}\mathcal{A}_U[k]
$$
est un générateur où la somme est effectuée sur tous les objets $U$ de $\mathcal{C}$
et $k \in \mathbf{Z}$. En effet, étant donné un $\mathcal{A}$-module
gradué $\mathcal{M}$, s'il n'existe aucun morphisme non nul de $\mathcal{G}$ vers $\mathcal{M}$,
alors, pour tout $k$ et tout $U$,
on a $$
\Hom_{\textit{Mod}(\mathcal{A})}(j_{U!}\mathcal{A}_U[k], \mathcal{M})
= \Hom_{\textit{Mod}(\mathcal{A}_U)}(\mathcal{A}_U[k], \mathcal{M}|_U) =
\Gamma(U, \mathcal{M}^{-k})
$$ et
ce groupe est nul. Ainsi, $\mathcal{M}$ est nul.
\end{proof}









\section{Faisceaux d'algèbres différentielles graduées}
\label{sdga-section-dga}

\noindent
Cette section est l'analogue
d'Algèbre différentielle graduée, section \ref{dga-section-dga}.

\begin{definition}
\label{sdga-definition-dga} Soit
$(\mathcal{C}, \mathcal{O})$ un site annelé. Un
{\it faisceau de $\mathcal{O}$-algèbres différentielles graduées}, ou
un {\it faisceau d'algèbres différentielles graduées},
sur $(\mathcal{C}, \mathcal{O})$ est un complexe de
cochaînes $\mathcal{A}^\bullet$ de $\mathcal{O}$-modules
muni d'applications $\mathcal{O}$-bilinéaires
$$
\mathcal{A}^n \times \mathcal{A}^m \to \mathcal{A}^{n + m},\quad
(a, b) \longmapsto ab
$$
appelées applications de multiplication et vérifiant les propriétés
suivantes : \begin{enumerate}
\item la multiplication est
associative, \item il existe une section globale $1$ de $\mathcal{A}^0$
qui est une unité bilatère pour la
multiplication, \item pour $U \in \Ob(\mathcal{C})$, $a \in \mathcal{A}^n(U)$
et $b \in \mathcal{A}^m(U)$, on a
$$
\text{d}^{n + m}(ab) = \text{d}^n(a)b + (-1)^n a\text{d}^m(b)
$$
\end{enumerate} Nous
notons souvent une telle structure $(\mathcal{A}, \text{d})$. Un
{\it homomorphisme d'algèbres différentielles graduées sur $\mathcal{O}$}
de $(\mathcal{A}, \text{d})$ vers $(\mathcal{B}, \text{d})$ est un morphisme
$f : \mathcal{A}^\bullet \to \mathcal{B}^\bullet$ de complexes de
$\mathcal{O}$-modules compatible avec les applications de multiplication.
\end{definition}

\noindent
Si l'on donne une $\mathcal{O}$-algèbre différentielle graduée $(\mathcal{A}, \text{d})$
et un objet $U \in \Ob(\mathcal{C})$, on utilise la notation
$$
\mathcal{A}(U)
= \Gamma(U, \mathcal{A}) =
\bigoplus\nolimits_{n \in \mathbf{Z}} \mathcal{A}^n(U)
$$ Il
s'agit d'une $\mathcal{O}(U)$-algèbre différentielle graduée.

\medskip\noindent
Dans la mesure du possible, nous considérerons une
$\mathcal{O}$-algèbre différentielle graduée $(\mathcal{A}, \text{d})$
comme une $\mathcal{O}$-algèbre graduée $\mathcal{A}$ munie
de l'opérateur $\text{d} : \mathcal{A} \to \mathcal{A}$ de degré $1$
(où $\mathcal{A}$ est considérée comme un $\mathcal{O}$-module gradué),
satisfaisant à la règle de Leibniz de la définition.

\begin{remark}
\label{sdga-remark-functoriality-dga} Soit
$(f, f^\sharp) : (\Sh(\mathcal{C}), \mathcal{O}_\mathcal{C})
\to (\Sh(\mathcal{D}), \mathcal{O}_\mathcal{D})$ un morphisme
de topos annelés.
\begin{enumerate}
\item Soit $(\mathcal{A}, \text{d})$ une
$\mathcal{O}_\mathcal{C}$-algèbre différentielle graduée. L'image directe est la
$\mathcal{O}_\mathcal{D}$-algèbre différentielle graduée
$(f_*\mathcal{A}, \text{d})$, où $f_*\mathcal{A}$ est comme dans la remarque
\ref{sdga-remark-functoriality-ga} et
$\text{d} = f_*\text{d}$ comme morphisme $f_*\mathcal{A}^n \to f_*\mathcal{A}^{n + 1}$. Nous omettons
la vérification que la règle de Leibniz est satisfaite.
\item Soit $\mathcal{B}$ une
$\mathcal{O}_\mathcal{D}$-algèbre différentielle graduée. L'image inverse est la
$\mathcal{O}_\mathcal{C}$-algèbre différentielle graduée
$(f^*\mathcal{B}, \text{d})$, où $f^*\mathcal{B}$ est comme dans la remarque
\ref{sdga-remark-functoriality-ga} et
$\text{d} = f^*\text{d}$ comme morphisme $f^*\mathcal{B}^n \to f^*\mathcal{B}^{n + 1}$. Nous omettons
la vérification que la règle de Leibniz est satisfaite.
\item L'ensemble des homomorphismes $f^*\mathcal{B} \to \mathcal{A}$ de
$\mathcal{O}_\mathcal{C}$-algèbres différentielles graduées est en
correspondance biunivoque ($1$ à $1$) avec l'ensemble des homomorphismes
$\mathcal{B} \to f_*\mathcal{A}$ de
$\mathcal{O}_\mathcal{D}$-algèbres différentielles graduées.
\end{enumerate} La
partie (3) découle immédiatement de l'adjonction usuelle entre $f^*$ et $f_*$ sur les
faisceaux de modules.
\end{remark}
















\section{Faisceaux de modules différentiels gradués}
\label{sdga-section-modules}

\noindent
Cette section est l'analogue de
l'Algèbre différentielle graduée, section \ref{dga-section-modules}.

\begin{definition}
\label{sdga-definition-dgm}
Soit $(\mathcal{C}, \mathcal{O})$ un site annelé. Soit
$(\mathcal{A}, \text{d})$ un faisceau d'algèbres différentielles graduées sur
$(\mathcal{C}, \mathcal{O})$. Un
{\it $\mathcal{A}$-module différentiel gradué à droite}, ou
{\it module différentiel gradué à droite} sur $\mathcal{A}$, est un
complexe de cochaînes $\mathcal{M}^\bullet$ muni d'applications
$\mathcal{O}$-bilinéaires
$$
\mathcal{M}^n \times \mathcal{A}^m \to \mathcal{M}^{n + m},\quad (x,
a) \longmapsto xa
$$
appelées applications de multiplication, satisfaisant aux propriétés
suivantes : \begin{enumerate}
\item la multiplication vérifie $(xa)a' = x(aa')$,
\item la section unité $1$ de $\mathcal{A}^0$
agit comme l'identité sur $\mathcal{M}^n$ pour tout $n$,
\item pour $U \in \Ob(\mathcal{C})$, $x \in \mathcal{M}^n(U)$ et
$a \in \mathcal{A}^m(U)$, on a
$$
\text{d}^{n + m}(xa) = \text{d}^n(x)a + (-1)^n x\text{d}^m(a)
$$
\end{enumerate} Nous
disons souvent « soit $\mathcal{M}$ un
$\mathcal{A}$-module différentiel gradué » pour indiquer cette situation. Un
{\it homomorphisme de $\mathcal{A}$-modules différentiels gradués} de
$\mathcal{M}$ vers $\mathcal{N}$ est un morphisme
$f : \mathcal{M}^\bullet \to \mathcal{N}^\bullet$ de complexes de
$\mathcal{O}$-modules compatible avec les applications de multiplication.
La catégorie des $\mathcal{A}$-modules différentiels gradués (à
droite) est notée $\textit{Mod}(\mathcal{A}, \text{d})$.
\end{definition}

\noindent
On définit de même les {\it modules différentiels gradués à gauche}, mais, sauf mention
contraire, les modules de ce chapitre seront des modules à droite.

\medskip\noindent
Étant donnés un $\mathcal{A}$-module différentiel gradué $\mathcal{M}$
et un objet $U \in \Ob(\mathcal{C})$, nous utilisons la notation
$$
\mathcal{M}(U)
= \Gamma(U, \mathcal{M})
= \bigoplus\nolimits_{n \in \mathbf{Z}} \mathcal{M}^n(U)
$$
Il s'agit d'un $\mathcal{A}(U)$-module différentiel gradué à droite.

\begin{lemma}
\label{sdga-lemma-dgm-abelian} Soit
$(\mathcal{C}, \mathcal{O})$ un site annelé. Soit
$(\mathcal{A}, \text{d})$ une $\mathcal{O}$-algèbre différentielle graduée. La catégorie
$\textit{Mod}(\mathcal{A}, \text{d})$ est une catégorie abélienne qui possède les
propriétés suivantes :
\begin{enumerate}
\item $\textit{Mod}(\mathcal{A}, \text{d})$ possède des sommes directes arbitraires,
\item $\textit{Mod}(\mathcal{A}, \text{d})$ possède des limites inductives arbitraires,
\item les limites inductives filtrantes dans $\textit{Mod}(\mathcal{A}, \text{d})$ sont exactes,
\item $\textit{Mod}(\mathcal{A}, \text{d})$ possède des produits
arbitraires, \item $\textit{Mod}(\mathcal{A}, \text{d})$ possède des limites
arbitraires. \end{enumerate}
Le foncteur
d'oubli $$
\textit{Mod}(\mathcal{A}, \text{d})
\longrightarrow
\textit{Mod}(\mathcal{A})
$$
qui envoie un $\mathcal{A}$-module différentiel gradué sur son module
gradué sous-jacent commute à toutes les limites et limites inductives.
\end{lemma}

\begin{proof}
Notons
$F : \textit{Mod}(\mathcal{A}, \text{d}) \to \textit{Mod}(\mathcal{A})$
le foncteur de l'énoncé. La catégorie
$\textit{Mod}(\mathcal{A})$ possède les propriétés (1) -- (5) ; voir le lemme
\ref{sdga-lemma-gm-abelian}.

\medskip\noindent
Considérons un homomorphisme $f : \mathcal{M} \to \mathcal{N}$
de $\mathcal{A}$-modules gradués. Le noyau et
le conoyau de $f$, comme morphisme de $\mathcal{A}$-modules
gradués, sont en outre munis de différentiels comme dans
la définition \ref{sdga-definition-dgm}. Ce sont donc aussi le
noyau et le conoyau dans $\textit{Mod}(\mathcal{A}, \text{d})$. Ainsi,
$\textit{Mod}(\mathcal{A}, \text{d})$ est une catégorie abélienne, et
la formation des noyaux et conoyaux commute avec $F$.

\medskip\noindent
Pour prouver l'existence de limites et de limites inductives, il suffit de
prouver l'existence de produits et de sommes directes, voir
Catégories, lemmes \ref{categories-lemma-limits-products-equalizers} et
\ref{categories-lemma-colimits-coproducts-coequalizers}. Les mêmes lemmes
montrent que, pour prouver
que $F$ commute aux limites et aux limites inductives, il suffit de
prouver que $F$ commute aux sommes directes et aux produits.

\medskip\noindent
Soit $\mathcal{M}_t$, $t \in T$, une famille de
$\mathcal{A}$-modules différentiels gradués. On peut considérer la somme directe
$\bigoplus \mathcal{M}_t$ comme un $\mathcal{A}$-module gradué.
Puisque la somme directe des modules gradués se calcule terme à terme,
$\bigoplus \mathcal{M}_t$ est naturellement muni d'un différentiel. Le lecteur
vérifie facilement que c'est une somme directe en
$\textit{Mod}(\mathcal{A}, \text{d})$. Il en va de même pour les produits.

\medskip\noindent
Observons que $F$ est exact et qu'un complexe
$\mathcal{M}_1 \to \mathcal{M}_2 \to \mathcal{M}_3$ de
$\textit{Mod}(\mathcal{A}, \text{d})$ est exact si et seulement si
$F(\mathcal{M}_1) \to F(\mathcal{M}_2) \to F(\mathcal{M}_3)$ est
exact dans $\textit{Mod}(\mathcal{A})$. On en déduit (3), puisque les
limites inductives filtrantes sont exactes dans
$\textit{Mod}(\mathcal{A})$.
\end{proof}

\noindent
En combinant les lemmes \ref{sdga-lemma-dgm-abelian} et \ref{sdga-lemma-gm-abelian},
on obtient un foncteur exact et fidèle
$$
\textit{Mod}(\mathcal{A}, \text{d}) \longrightarrow \text{Comp}(\mathcal{O})
$$ entre catégories
abéliennes. Pour un $\mathcal{A}$-module différentiel gradué
$\mathcal{M}$, les $\mathcal{O}$-modules de cohomologie, notés $H^i(\mathcal{M})$, sont
définis comme les cohomologies du complexe sous-jacent
de $\mathcal{O}$-modules associé à $\mathcal{M}$. Par
conséquent, une suite exacte courte
$0 \to \mathcal{K} \to \mathcal{L} \to \mathcal{M} \to 0$ de
$\mathcal{A}$-modules différentiels gradués
donne lieu à une suite exacte longue
\begin{equation}
\label{sdga-equation-les}
H^n(\mathcal{K}) \to H^n(\mathcal{L}) \to H^n(\mathcal{M}) \to
H^{n + 1}(\mathcal{K})
\end{equation} des
modules de cohomologie, voir
Homologie, lemme \ref{homology-lemma-long-exact-sequence-cochain}.

\medskip\noindent En
outre, à partir de maintenant, nous empruntons toute la terminologie
utilisée pour les complexes de modules. Par exemple, nous disons
qu'un $\mathcal{A}$-module différentiel gradué $\mathcal{M}$ est {\it acyclique}
si $H^k(\mathcal{M}) = 0$ pour tous les $k \in \mathbf{Z}$.
Un homomorphisme $\mathcal{M} \to \mathcal{N}$
de $\mathcal{A}$-modules différentiels gradués est un {\it quasi-isomorphisme}
s'il induit des isomorphismes $H^k(\mathcal{M}) \to H^k(\mathcal{N})$
pour tout $k \in \mathbf{Z}$. Et ainsi de suite.


















\section{La catégorie différentielle graduée des modules}
\label{sdga-section-dgm-dg-cat}

\noindent
Cette section est l'analogue
d'Algèbre différentielle graduée, exemple \ref{dga-example-dgm-dg-cat}.
Pour nos conventions sur les catégories différentielles graduées,
voir Algèbre différentielle graduée, section \ref{dga-section-dga-categories}.

\medskip\noindent
Soit $(\mathcal{C}, \mathcal{O})$ un site annelé. Soit
$(\mathcal{A}, \text{d})$ un faisceau d'algèbres différentielles graduées sur
$(\mathcal{C}, \mathcal{O})$. Nous allons construire une
catégorie différentielle graduée
$$
\textit{Mod}^{dg}(\mathcal{A}, \text{d})
$$
sur $R = \Gamma(\mathcal{C}, \mathcal{O})$, dont la catégorie de complexes
associée est la catégorie des $\mathcal{A}$-modules différentiels gradués :
$$
\textit{Mod}(\mathcal{A}, \text{d}) =
\text{Comp}(\textit{Mod}^{dg}(\mathcal{A}, \text{d}))
$$ Comme objets
de $\textit{Mod}^{dg}(\mathcal{A}, \text{d})$, nous prenons les
$\mathcal{A}$-modules différentiels gradués à droite, voir
la section \ref{sdga-section-modules}. Étant donnés des
$\mathcal{A}$-modules différentiels gradués $\mathcal{L}$ et $\mathcal{M}$,
posons $$
\Hom_{\textit{Mod}^{dg}(\mathcal{A}, \text{d})}(\mathcal{L}, \mathcal{M})
= \Hom_{\textit{Mod}^{gr}(\mathcal{A})}(\mathcal{L}, \mathcal{M}) =
\bigoplus\nolimits_{n \in \mathbf{Z}} \Hom^n(\mathcal{L}, \mathcal{M})
$$ comme
$R$-module gradué, voir la section \ref{sdga-section-gm-gr-cat}. Autrement
dit, la composante graduée de degré $n$,
$\Hom^n(\mathcal{L}, \mathcal{M})$, est le $R$-module des morphismes
de $\mathcal{A}$-modules à droite homogènes de degré $n$.
Pour $f \in \Hom^n(\mathcal{L}, \mathcal{M})$, on pose
$$
\text{d}(f) =
\text{d}_\mathcal{M} \circ f - (-1)^n f \circ \text{d}_\mathcal{L}
$$ Pour donner
un sens à cette formule, on considère $\text{d}_\mathcal{M}$ et
$\text{d}_\mathcal{L}$ comme des morphismes gradués de $\mathcal{O}$-modules
et on utilise la composition de morphismes gradués de $\mathcal{O}$-modules.
Il est clair que $\text{d}(f)$
est homogène de degré $n + 1$ comme morphisme gradué de $\mathcal{O}$-modules,
et qu'il est $\mathcal{A}$-linéaire, car, pour des sections locales homogènes $x$ et $a$
de $\mathcal{M}$ et $\mathcal{A}$,
on a \begin{align*}
\text{d}(f)(xa)
&
= \text{d}_\mathcal{M}(f(x) a) - (-1)^n f (\text{d}_\mathcal{L}(xa)) \\
&
= \text{d}_\mathcal{M}(f(x)) a + (-1)^{\deg(x) + n} f(x) \text{d}(a)
- (-1)^n f(\text{d}_\mathcal{L}(x)) a - (-1)^{n + \deg(x)} f(x) \text{d}(a) \\
& = \text{d}(f)(x)
a \end{align*}
comme souhaité (observons que ce calcul ne fonctionnerait pas sans le
signe dans la définition de notre différentiel sur $\Hom$).

\medskip\noindent
Pour des $\mathcal{A}$-modules différentiels
gradués $\mathcal{K}$, $\mathcal{L}$ et $\mathcal{M}$,
nous avons déjà défini la composition
$$
\Hom^m(\mathcal{L}, \mathcal{M}) \times
\Hom^n(\mathcal{K}, \mathcal{L}) \longrightarrow
\Hom^{n + m}(\mathcal{K}, \mathcal{M})
$$
dans la section \ref{sdga-section-gm-gr-cat} par composition usuelle
de morphismes de faisceaux. Cela définit un morphisme de modules différentiels
gradués $$
\Hom_{\textit{Mod}^{dg}(\mathcal{A}, \text{d})}(\mathcal{L}, \mathcal{M})
\otimes_R
\Hom_{\textit{Mod}^{dg}(\mathcal{A}, \text{d})}(\mathcal{K}, \mathcal{L})
\longrightarrow
\Hom_{\textit{Mod}^{dg}(\mathcal{A}, \text{d})}(\mathcal{K}, \mathcal{M})
$$
comme
l'exige l'Algèbre différentielle graduée, définition \ref{dga-definition-dga-category},
comme le montre le calcul suivant :
\begin{align*}
\text{d}(g \circ
f) & = \text{d}_\mathcal{M} \circ g \circ
f - (-1)^{n + m} g \circ f \circ \text{d}_\mathcal{K} \\
&
= \left(\text{d}_\mathcal{M} \circ g - (-1)^m g \circ \text{d}_L\right) \circ
f + (-1)^m g \circ \left(\text{d}_\mathcal{L} \circ
f - (-1)^n f \circ \text{d}_\mathcal{K}\right) \\
& =
\text{d}(g) \circ f + (-1)^m g \circ \text{d}(f)
\end{align*}
si $f$ est de degré $n$ et $g$ de degré $m$, comme souhaité.






\section{Produit tensoriel de faisceaux de modules différentiels gradués}
\label{sdga-section-tensor-product-dg}

\noindent
Cette section est l'analogue d'une partie de
l'Algèbre différentielle graduée, section \ref{dga-section-tensor-product}.

\medskip\noindent
Soit $(\mathcal{C}, \mathcal{O})$ un site annelé. Soit
$(\mathcal{A}, \text{d})$ un
faisceau d'algèbres différentielles graduées sur $(\mathcal{C}, \mathcal{O})$.
Soient $\mathcal{M}$ un $\mathcal{A}$-module différentiel gradué à
droite et $\mathcal{N}$ un $\mathcal{A}$-module différentiel gradué à gauche.
Dans cette situation, nous définissons le {\it produit tensoriel}
$\mathcal{M} \otimes_\mathcal{A} \mathcal{N}$ comme
suit. Comme $\mathcal{O}$-module gradué, il est donné par la
construction de la section \ref{sdga-section-tensor-product}. Il est
doté d'un différentiel
$$
\text{d}_{\mathcal{M} \otimes_\mathcal{A} \mathcal{N}}
: (\mathcal{M} \otimes_\mathcal{A} \mathcal{N})^n
\longrightarrow
(\mathcal{M} \otimes_\mathcal{A} \mathcal{N})^{n + 1}
$$
défini par la règle
$$
\text{d}_{\mathcal{M} \otimes_\mathcal{A} \mathcal{N}}(x \otimes y) =
\text{d}_\mathcal{M}(x) \otimes y +
(-1)^{\deg(x)}x \otimes \text{d}_\mathcal{N}(y)
$$ pour
les sections locales homogènes $x$ et $y$ de $\mathcal{M}$ et $\mathcal{N}$. Pour
vérifier que cette définition est correcte, montrons
que $\text{d}_{\mathcal{M} \otimes_\mathcal{A} \mathcal{N}}$ annule les
éléments de la forme $xa \otimes y - x \otimes ay$, pour des sections
locales homogènes $x$, $a$, $y$ de $\mathcal{M}$, $\mathcal{A}$, $\mathcal{N}$.
Nous
calculons \begin{align*}
&
\text{d}_{\mathcal{M} \otimes_\mathcal{A} \mathcal{N}}(
xa \otimes y - x \otimes ay) \\
&
= \text{d}_\mathcal{M}(xa) \otimes y + (-1)^{\deg(x) + \deg(a)}
xa \otimes \text{d}_\mathcal{N}(y)
-\text{d}_\mathcal{M}(x) \otimes ay - (-1)^{\deg(x)}
x \otimes \text{d}_\mathcal{N}(ay) \\
&
= \text{d}_\mathcal{M}(x)a \otimes y + (-1)^{\deg(x)}x\text{d}(a) \otimes
y + (-1)^{\deg(x) + \deg(a)} xa \otimes \text{d}_\mathcal{N}(y) \\
&
-\text{d}_\mathcal{M}(x) \otimes ay - (-1)^{\deg(x)}
x \otimes \text{d}(a)y - (-1)^{\deg(x) +
\deg(a)} x\otimes a\text{d}_\mathcal{N}(y)
\end{align*}
Nous observons que les éléments
$$
\text{d}_\mathcal{M}(x)a \otimes y - \text{d}_\mathcal{M}(x) \otimes ay,\quad
x\text{d}(a) \otimes y - x \otimes \text{d}(a)y,\quad
\text{et}\quad xa
\otimes \text{d}_\mathcal{N}(y) - x\otimes a\text{d}_\mathcal{N}(y)
$$ ont une
image nulle dans $\mathcal{M} \otimes_\mathcal{A} \mathcal{N}$,
ce qui conclut. Nous omettons la vérification
que $\text{d}_{\mathcal{M} \otimes_\mathcal{A} \mathcal{N}} \circ
\text{d}_{\mathcal{M} \otimes_\mathcal{A} \mathcal{N}} = 0$.

\medskip\noindent
Si nous fixons le $\mathcal{A}$-module différentiel gradué à gauche $\mathcal{N}$,
nous obtenons un
foncteur $$
- \otimes_\mathcal{A} \mathcal{N}
: \textit{Mod}(\mathcal{A}, \text{d})
\longrightarrow
\text{Comp}(\mathcal{O})
$$
où le membre de droite est la catégorie des complexes de
$\mathcal{O}$-modules. Ce foncteur se relève en un foncteur de
catégories différentielles graduées
$$
- \otimes_\mathcal{A} \mathcal{N} :
\textit{Mod}^{dg}(\mathcal{A}, \text{d})
\longrightarrow
\text{Comp}^{dg}(\mathcal{O})
$$ Sur
les objets gradués sous-jacents, nous
envoyons un homomorphisme $f : \mathcal{M} \to \mathcal{M}'$ de degré $n$
sur l'application de degré $n$,
$f \otimes \text{id}_\mathcal{N}
: \mathcal{M} \otimes_\mathcal{A} \mathcal{N} \to
\mathcal{M}' \otimes_\mathcal{A} \mathcal{N}$,
conformément à la construction de la section \ref{sdga-section-tensor-product}.
Pour montrer que cela fonctionne, nous devons
vérifier que l'application $$
\Hom_{\textit{Mod}^{dg}(\mathcal{A}, \text{d})}(\mathcal{M}, \mathcal{M}')
\longrightarrow
\Hom_{\text{Comp}^{dg}(\mathcal{O})}(
\mathcal{M} \otimes_\mathcal{A} \mathcal{N},
\mathcal{M}' \otimes_\mathcal{A} \mathcal{N})
$$
est compatible avec les différentiels. Pour voir cela pour $f$
comme
ci-dessus nous devons montrer que $$
(\text{d}_{\mathcal{M}'} \circ f - (-1)^n f \circ \text{d}_\mathcal{M})
\otimes \text{id}_\mathcal{N}
$$
est
égal à $$
\text{d}_{\mathcal{M}' \otimes_\mathcal{A} \mathcal{N}}
\circ (f \otimes \text{id}_\mathcal{N})
- (-1)^n (f \otimes \text{id}_\mathcal{N}) \circ
\text{d}_{\mathcal{M} \otimes_\mathcal{A} \mathcal{N}}
$$
Calculons donc l'effet de ces opérateurs sur
une section locale de la forme $x \otimes y$, où $x$ et $y$ sont des sections
locales homogènes de $\mathcal{M}$ et $\mathcal{N}$. Pour le
premier, nous obtenons $$
(\text{d}_{\mathcal{M}'}(f(x)) - (-1)^n f(\text{d}_\mathcal{M}(x))) \otimes
y $$
et pour
le second nous obtenons \begin{align*}
&\text{d}_{\mathcal{M}' \otimes_\mathcal{A} \mathcal{N}}(f(x) \otimes
y) - (-1)^n (f \otimes \text{id}_\mathcal{N})(
\text{d}_{\mathcal{M} \otimes_\mathcal{A} \mathcal{N}}(x \otimes y) \\
&
= \text{d}_{\mathcal{M}'}(f(x)) \otimes
y + (-1)^{\deg(x) + n}f(x) \otimes \text{d}_\mathcal{N}(y) \\
&
-(-1)^n f(\text{d}_\mathcal{M}(x)) \otimes
y -(-1)^n (-1)^{\deg(x)}f(x) \otimes \text{d}_\mathcal{N}(y)
\end{align*}
qui est bien la même section locale.










\section{Hom interne de faisceaux de modules différentiels gradués}
\label{sdga-section-internal-hom-dg}

\noindent
Nous aurons besoin de la version faisceautisée de la construction de la
section \ref{sdga-section-dgm-dg-cat}. Soient
$(\mathcal{C}, \mathcal{O})$, $\mathcal{A}$,
$\mathcal{M}$ et $\mathcal{L}$ comme dans la section \ref{sdga-section-dgm-dg-cat}. Nous
définissons
$$
\SheafHom^{dg}_\mathcal{A}(\mathcal{M}, \mathcal{L}) =
\SheafHom^{gr}_\mathcal{A}(\mathcal{M}, \mathcal{L}) =
\bigoplus\nolimits_{n \in \mathbf{Z}}
\SheafHom^n_\mathcal{A}(\mathcal{M}, \mathcal{L})
$$ comme
$\mathcal{O}$-module gradué, voir la section
\ref{sdga-section-internal-hom-graded}. Autrement dit, une section
$f$ de la composante graduée de degré $n$,
$\SheafHom^n_\mathcal{A}(\mathcal{L}, \mathcal{M})$, sur $U$ est un morphisme de
$\mathcal{A}_U$-modules à droite
$\mathcal{L}|_U \to \mathcal{M}|_U$, homogène de degré $n$. Pour un tel
$f$, on pose
$$
\text{d}(f) =
\text{d}_\mathcal{M}|_U \circ f - (-1)^n f \circ \text{d}_\mathcal{L}|_U
$$ Pour donner un sens à cette formule,
on considère $\text{d}_\mathcal{M}|_U$ et
$\text{d}_\mathcal{L}|_U$ comme des morphismes gradués de $\mathcal{O}_U$-modules et on utilise la
composition de morphismes gradués de $\mathcal{O}_U$-modules. Il est
clair que $\text{d}(f)$ est homogène de degré
$n + 1$ comme morphisme gradué de $\mathcal{O}_U$-modules. En
utilisant exactement le même calcul que dans la section \ref{sdga-section-dgm-dg-cat},
nous voyons que $\text{d}(f)$ est $\mathcal{A}_U$-linéaire.

\medskip\noindent
Comme dans la section \ref{sdga-section-dgm-dg-cat}, il y a une application de
composition $$
\SheafHom^{dg}_\mathcal{A}(\mathcal{L}, \mathcal{M}) \otimes_\mathcal{O}
\SheafHom^{dg}_\mathcal{A}(\mathcal{K}, \mathcal{L})
\longrightarrow
\SheafHom^{dg}_\mathcal{A}(\mathcal{K}, \mathcal{M})
$$ où le
côté gauche est le produit tensoriel de
$\mathcal{O}$-modules différentiels gradués défini à la section
\ref{sdga-section-tensor-product-dg}. Cette application est donnée
par l'application de composition
$$
\SheafHom^m(\mathcal{L}, \mathcal{M}) \otimes_\mathcal{O}
\SheafHom^n(\mathcal{K}, \mathcal{L}) \longrightarrow
\SheafHom^{n + m}(\mathcal{K}, \mathcal{M})
$$ définie par simple
composition (localement). En appliquant exactement le même calcul
que dans la section \ref{sdga-section-dgm-dg-cat} aux sections locales, nous voyons
que l'application de composition est un morphisme de
$\mathcal{O}$-modules différentiels gradués.

\medskip\noindent
Avec ces définitions, nous avons
$$
\Hom_{\textit{Mod}^{dg}(\mathcal{A})}(\mathcal{L}, \mathcal{M}) =
\Gamma(\mathcal{C}, \SheafHom^{dg}_\mathcal{A}(\mathcal{L}, \mathcal{M}))
$$ comme
$R$-modules gradués, compatiblement à la composition.









\section{Faisceaux de bimodules différentiels gradués et adjonction tensor-Hom}
\label{sdga-section-dg-bimodules}

\noindent
Cette section est l'analogue d'une partie de
l'Algèbre différentielle graduée, section \ref{dga-section-tensor-product}.

\begin{definition}
\label{sdga-definition-dg-bimodule}
Soit $(\mathcal{C}, \mathcal{O})$ un site annelé. Soient $\mathcal{A}$ et
$\mathcal{B}$ des faisceaux d'algèbres différentielles graduées sur
$(\mathcal{C}, \mathcal{O})$. Un
{\it $(\mathcal{A}, \mathcal{B})$-bimodule différentiel gradué} est
la donnée d'un complexe $\mathcal{M}^\bullet$
de $\mathcal{O}$-modules muni d'applications $\mathcal{O}$-bilinéaires
$$
\mathcal{M}^n \times \mathcal{B}^m \to \mathcal{M}^{n + m},\quad
(x, b) \longmapsto xb
$$
et
$$
\mathcal{A}^n \times \mathcal{M}^m \to \mathcal{M}^{n + m},\quad (a,
x) \longmapsto ax
$$
appelées applications de multiplication, avec les propriétés suivantes
\begin{enumerate}
\item la multiplication satisfait $a(a'x) = (aa')x$
et $(xb)b' = x(bb')$,
\item $(ax)b = a(xb)$,
\item $\text{d}(ax) = \text{d}(a) x + (-1)^{\deg(a)}a \text{d}(x)$ et
$\text{d}(xb) = \text{d}(x) b + (-1)^{\deg(x)}x \text{d}(b)$,
\item la section unité $1$ de $\mathcal{A}^0$ agit comme l'identité par
multiplication, et
\item la section unité $1$ de
$\mathcal{B}^0$ agit comme l'identité par multiplication.
\end{enumerate} Nous désignons
souvent une telle structure $\mathcal{M}$ et parfois nous écrivons
${}_\mathcal{A}\mathcal{M}_\mathcal{B}$. Un
{\it homomorphisme de
$(\mathcal{A}, \mathcal{B})$-bimodules différentiels gradués}
$f : \mathcal{M} \to \mathcal{N}$ est un morphisme de complexes
$f : \mathcal{M}^\bullet \to \mathcal{N}^\bullet$ de
$\mathcal{O}$-modules compatible avec les applications de multiplication.
\end{definition}

\noindent
Étant donnés un $(\mathcal{A}, \mathcal{B})$-bimodule différentiel gradué $\mathcal{M}$
et un objet $U \in \Ob(\mathcal{C})$, nous utilisons la
notation $$
\mathcal{M}(U)
= \Gamma(U, \mathcal{M})
= \bigoplus\nolimits_{n \in \mathbf{Z}} \mathcal{M}^n(U)
$$
C'est un $(\mathcal{A}(U), \mathcal{B}(U))$-bimodule différentiel gradué.

\medskip\noindent
Observons qu'un $(\mathcal{A}, \mathcal{B})$-bimodule différentiel gradué
$\mathcal{M}$ revient à se donner un
$\mathcal{B}$-module différentiel gradué à droite qui est aussi un
$\mathcal{A}$-module différentiel gradué à gauche, de sorte que les graduations et les
différentiels coïncident et que les structures de $\mathcal{A}$-module
et de $\mathcal{B}$-module commutent. Voici un énoncé précis.

\begin{lemma}
\label{sdga-lemma-what-makes-a-bimodule-dg} Soit
$(\mathcal{C}, \mathcal{O})$ un site annelé. Soient $\mathcal{A}$ et
$\mathcal{B}$ des faisceaux d'algèbres différentielles graduées sur
$(\mathcal{C}, \mathcal{O})$. Soit $\mathcal{N}$ un
$\mathcal{B}$-module différentiel gradué à droite. Il existe une correspondance biunivoque ($1$ à $1$) entre les
structures de $(\mathcal{A}, \mathcal{B})$-bimodule sur
$\mathcal{N}$ compatibles avec la structure
donnée de $\mathcal{B}$-module différentiel gradué et les homomorphismes
$$
\mathcal{A}
\longrightarrow
\SheafHom^{dg}_\mathcal{B}(\mathcal{N}, \mathcal{N})
$$ de
$\mathcal{O}$-algèbres différentielles graduées.
\end{lemma}

\begin{proof}
Démonstration
omise. \end{proof}

\noindent
Soit $(\mathcal{C}, \mathcal{O})$ un site annelé. Soient $\mathcal{A}$ et
$\mathcal{B}$ des faisceaux d'algèbres différentielles graduées sur
$(\mathcal{C}, \mathcal{O})$. Soient $\mathcal{M}$ un
$\mathcal{A}$-module différentiel gradué à droite et $\mathcal{N}$ un
$(\mathcal{A}, \mathcal{B})$-bimodule différentiel gradué. Dans ce cas, le produit
tensoriel différentiel gradué défini à la section
\ref{sdga-section-tensor-product-dg}
$$
\mathcal{M} \otimes_\mathcal{A} \mathcal{N}
$$ est un
$\mathcal{B}$-module différentiel gradué à droite, avec les
applications de multiplication de la section \ref{sdga-section-graded-bimodules}.
Cette construction définit un foncteur et un foncteur de catégories
graduées $$
\otimes_\mathcal{A} \mathcal{N}
: \textit{Mod}(\mathcal{A}, \text{d})
\longrightarrow
\textit{Mod}(\mathcal{B}, \text{d})
\quad\text{et}\quad
\otimes_\mathcal{A} \mathcal{N}
: \textit{Mod}^{dg}(\mathcal{A}, \text{d})
\longrightarrow
\textit{Mod}^{dg}(\mathcal{B}, \text{d})
$$
en envoyant les homomorphismes de degré $n$ de $\mathcal{M} \to \mathcal{M}'$
sur l'application induite de degré $n$
de $\mathcal{M} \otimes_\mathcal{A} \mathcal{N}$ vers
$\mathcal{M}' \otimes_\mathcal{A} \mathcal{N}$.

\medskip\noindent
Soit $(\mathcal{C}, \mathcal{O})$ un site annelé. Soient $\mathcal{A}$ et
$\mathcal{B}$ des faisceaux d'algèbres différentielles graduées sur
$(\mathcal{C}, \mathcal{O})$. Soit $\mathcal{N}$ un
$(\mathcal{A}, \mathcal{B})$-bimodule différentiel gradué et soit
$\mathcal{L}$ un $\mathcal{B}$-module différentiel gradué à droite. Dans
ce cas, le Hom interne différentiel gradué défini à
la section \ref{sdga-section-internal-hom-dg}
$$
\SheafHom_\mathcal{B}^{dg}(\mathcal{N}, \mathcal{L})
$$ est un
$\mathcal{A}$-module différentiel gradué à droite dont la structure
graduée de $\mathcal{A}$-module à droite est celle de la section
\ref{sdga-section-graded-bimodules}. Une autre définition de la multiplication
consiste à utiliser la composition
$$
\SheafHom_\mathcal{B}^{dg}(\mathcal{N}, \mathcal{L})
\otimes_\mathcal{O} \mathcal{A}
\to
\SheafHom_\mathcal{B}^{dg}(\mathcal{N}, \mathcal{L})
\otimes_\mathcal{O} \SheafHom_\mathcal{B}^{dg}(\mathcal{N}, \mathcal{N})
\to
\SheafHom_\mathcal{B}^{dg}(\mathcal{N}, \mathcal{L})
$$ où la première
flèche provient du lemme \ref{sdga-lemma-what-makes-a-bimodule-dg} et la seconde est
l'application de composition de la section
\ref{sdga-section-internal-hom-dg}. Puisque ces deux flèches sont
compatibles avec les différentiels, nous concluons que nous obtenons
en effet un $\mathcal{A}$-module différentiel gradué.
Cette construction
définit un foncteur et un foncteur de catégories différentielles
graduées $$
\SheafHom_\mathcal{B}^{dg}(\mathcal{N},
-) : \textit{Mod}(\mathcal{B}, \text{d})
\longrightarrow
\textit{Mod}(\mathcal{A})
\quad\text{et}\quad
\SheafHom_\mathcal{B}^{dg}(\mathcal{N}, -)
: \textit{Mod}^{dg}(\mathcal{B}, \text{d})
\longrightarrow
\textit{Mod}^{dg}(\mathcal{A}, \text{d})
$$
en envoyant les homomorphismes de degré $n$ de $\mathcal{L} \to \mathcal{L}'$
sur l'application induite de degré $n$
de $\SheafHom_\mathcal{B}^{dg}(\mathcal{N}, \mathcal{L})$ vers
$\SheafHom_\mathcal{B}^{dg}(\mathcal{N}, \mathcal{L}')$.

\begin{lemma}
\label{sdga-lemma-tensor-hom-adjunction-dg}
Soit $(\mathcal{C}, \mathcal{O})$ un site annelé. Soient $\mathcal{A}$
et $\mathcal{B}$ des faisceaux d'algèbres différentielles graduées
sur $(\mathcal{C}, \mathcal{O})$. Soient $\mathcal{M}$ un
$\mathcal{A}$-module différentiel gradué à droite, $\mathcal{N}$ un
$(\mathcal{A}, \mathcal{B})$-bimodule différentiel gradué et $\mathcal{L}$ un
$\mathcal{B}$-module différentiel gradué à droite. Avec les conventions ci-dessus, nous
avons
$$
\Hom_{\textit{Mod}^{dg}(\mathcal{B}, \text{d})}(
\mathcal{M} \otimes_\mathcal{A} \mathcal{N}, \mathcal{L}) =
\Hom_{\textit{Mod}^{dg}(\mathcal{A}, \text{d})}(
\mathcal{M}, \SheafHom_\mathcal{B}^{dg}(\mathcal{N}, \mathcal{L}))
$$
et
$$
\SheafHom_\mathcal{B}^{dg}(
\mathcal{M} \otimes_\mathcal{A} \mathcal{N}, \mathcal{L}) =
\SheafHom_\mathcal{A}^{dg}(
\mathcal{M}, \SheafHom_\mathcal{B}^{dg}(\mathcal{N}, \mathcal{L}))
$$
fonctoriellement en $\mathcal{M}$, $\mathcal{N}$ et $\mathcal{L}$.
\end{lemma}

\begin{proof}
Omis. Indication : au niveau gradué, nous avons établi ce résultat
dans le lemme \ref{sdga-lemma-tensor-hom-adjunction-gr}. Il suffit
donc de vérifier que les isomorphismes sont compatibles avec les
différentiels, ce qui se fait par un calcul sur les sections locales.
\end{proof}

\noindent
Soit $(\mathcal{C}, \mathcal{O})$ un site annelé. Soient $\mathcal{A}$
et $\mathcal{B}$ des faisceaux d'algèbres différentielles graduées
sur $(\mathcal{C}, \mathcal{O})$.
Comme cas particulier de ce qui précède, supposons
donné un homomorphisme $\varphi : \mathcal{A} \to \mathcal{B}$
de $\mathcal{O}$-algèbres différentielles graduées. On obtient alors
un foncteur et un foncteur de catégories différentielles
graduées $$
\otimes_{\mathcal{A}, \varphi} \mathcal{B}
: \textit{Mod}(\mathcal{A}, \text{d})
\longrightarrow
\textit{Mod}(\mathcal{B}, \text{d})
\quad\text{et}\quad
\otimes_{\mathcal{A}, \varphi} \mathcal{B}
: \textit{Mod}^{dg}(\mathcal{A}, \text{d})
\longrightarrow
\textit{Mod}^{dg}(\mathcal{B}, \text{d})
$$
D'autre part, nous avons les foncteurs de
restriction $$
res_\varphi
: \textit{Mod}(\mathcal{B}, \text{d})
\longrightarrow
\textit{Mod}(\mathcal{A}, \text{d})
\quad\text{et}\quad
res_\varphi
: \textit{Mod}^{dg}(\mathcal{B}, \text{d})
\longrightarrow
\textit{Mod}^{dg}(\mathcal{A}, \text{d})
$$
Le lemme ci-dessus montre que ces foncteurs sont adjoints
(comme d'habitude pour la restriction et le changement de base).
À savoir, écrivons ${}_\mathcal{A}\mathcal{B}_\mathcal{B}$
pour $\mathcal{B}$ considéré comme un
$(\mathcal{A}, \mathcal{B})$-bimodule différentiel
gradué. Alors, pour tout $\mathcal{B}$-module différentiel gradué $\mathcal{L}$,
nous
avons $$
\SheafHom_\mathcal{B}^{dg}({}_\mathcal{A}\mathcal{B}_\mathcal{B}, \mathcal{L})
= res_\varphi(\mathcal{L})
$$
comme $\mathcal{A}$-modules différentiels
gradués. Ainsi, le lemme \ref{sdga-lemma-tensor-hom-adjunction-gr}
fournit un isomorphisme
fonctoriel $$
\Hom_{\textit{Mod}^{dg}(\mathcal{B}, \text{d})}(
\mathcal{M} \otimes_{\mathcal{A}, \varphi} \mathcal{B}, \mathcal{L})
= \Hom_{\textit{Mod}^{dg}(\mathcal{A}, \text{d})}(
\mathcal{M}, res_\varphi(\mathcal{L}))
$$
Lorsque le contexte est clair, nous omettons généralement $\varphi$
dans cette formule. De la même manière, nous obtenons
l'égalité
$$
\SheafHom^{dg}_\mathcal{B}(
\mathcal{M} \otimes_\mathcal{A} \mathcal{B}, \mathcal{L}) =
\SheafHom_\mathcal{A}^{dg}(\mathcal{M}, \mathcal{L})
$$
de $\mathcal{O}$-modules gradués.












\section{Image inverse et image directe des faisceaux de modules différentiels gradués}
\label{sdga-section-functoriality-dg}

\noindent
Soit $(f, f^\sharp) : (\Sh(\mathcal{C}), \mathcal{O}_\mathcal{C})
\to (\Sh(\mathcal{D}), \mathcal{O}_\mathcal{D})$ un morphisme
de topos annelés. Soient $\mathcal{A}$ une
$\mathcal{O}_\mathcal{C}$-algèbre différentielle graduée et $\mathcal{B}$ une
$\mathcal{O}_\mathcal{D}$-algèbre différentielle graduée. Supposons qu'on nous donne
une application
$$
\varphi : f^{-1}\mathcal{B} \to \mathcal{A}
$$ de
$f^{-1}\mathcal{O}_\mathcal{D}$-algèbres différentielles graduées. Par l'adjonction
entre restriction et extension des scalaires, cela revient à se donner
un morphisme $\varphi : f^*\mathcal{B} \to \mathcal{A}$ de
$\mathcal{O}_\mathcal{C}$-algèbres différentielles graduées ou, de manière équivalente,
$\varphi$ peut être considéré comme un morphisme
$$
\varphi : \mathcal{B} \to f_*\mathcal{A}
$$
de $\mathcal{O}_\mathcal{D}$-algèbres différentielles graduées.
Voir la remarque \ref{sdga-remark-functoriality-dga}.

\medskip\noindent
Définissons maintenant un
foncteur $$
f_*
: \textit{Mod}(\mathcal{A}, \text{d})
\longrightarrow
\textit{Mod}(\mathcal{B}, \text{d})
$$
Étant donné un $\mathcal{A}$-module différentiel gradué $\mathcal{M}$, nous
définissons $f_*\mathcal{M}$ comme le $\mathcal{B}$-module gradué construit dans
la section \ref{sdga-section-functoriality-graded}, muni du différentiel
donné par les morphismes $f_*d : f_*\mathcal{M}^n \to f_*\mathcal{M}^{n + 1}$. La
construction est clairement fonctorielle en
$\mathcal{M}$ et nous obtenons notre foncteur.

\medskip\noindent
Définissons maintenant un
foncteur $$
f^*
: \textit{Mod}(\mathcal{B}, \text{d})
\longrightarrow
\textit{Mod}(\mathcal{A}, \text{d})
$$
Étant donné un $\mathcal{B}$-module différentiel gradué $\mathcal{N}$,
nous définissons $f^*\mathcal{N}$ comme le $\mathcal{A}$-module
gradué construit dans la section \ref{sdga-section-functoriality-graded}.
Rappelons
que $$
f^*\mathcal{N} = f^{-1}\mathcal{N} \otimes_{f^{-1}\mathcal{B}} \mathcal{A}
$$
Puisque $f^{-1}\mathcal{N}$ est muni des applications
$f^{-1}\text{d} : f^{-1}\mathcal{N}^n \to f^{-1}\mathcal{N}^{n + 1}$, nous
pouvons considérer ce produit tensoriel comme un exemple
du produit tensoriel étudié dans la section \ref{sdga-section-dg-bimodules},
qui le munit d'un différentiel. La
construction est clairement fonctorielle en
$\mathcal{N}$ et nous obtenons notre foncteur $f^*$.

\medskip\noindent
Les foncteurs $f_*$ et $f^*$ se relèvent naturellement en des
foncteurs de catégories différentielles graduées
$$
f_*
: \textit{Mod}^{dg}(\mathcal{A}, \text{d})
\longrightarrow
\textit{Mod}^{dg}(\mathcal{B}, \text{d})
\quad\text{et}\quad f^*
:
\textit{Mod}^{dg}(\mathcal{B}, \text{d})
\longrightarrow
\textit{Mod}^{dg}(\mathcal{A}, \text{d})
$$ qui
agissent de la même façon sur les objets sous-jacents et sont définis
sur les morphismes homogènes de degré
$n$ par la fonctorialité des constructions.

\begin{lemma}
\label{sdga-lemma-adjunction-push-pull-dg} Dans
la situation ci-dessus, nous
avons $$
\Hom_{\textit{Mod}^{dg}(\mathcal{B}, \text{d})}(
\mathcal{N}, f_*\mathcal{M})
= \Hom_{\textit{Mod}^{dg}(\mathcal{A}, \text{d})}(
f^*\mathcal{N}, \mathcal{M})
$$
\end{lemma}

\begin{proof}
Démonstration omise. Indication : l'énoncé est vrai pour les catégories graduées
sous-jacentes par le lemme \ref{sdga-lemma-adjunction-push-pull-gr}. Un calcul
montre que ces isomorphismes sont compatibles avec les différentiels.
\end{proof}















\section{Localisation et faisceaux de modules différentiels gradués}
\label{sdga-section-localize-dg}

\noindent
Soit $(\mathcal{C}, \mathcal{O})$ un site annelé. Soit
$U \in \Ob(\mathcal{C})$ et notons
$$ j
:
(\Sh(\mathcal{C}/U), \mathcal{O}_U)
\longrightarrow
(\Sh(\mathcal{C}), \mathcal{O})
$$ le
morphisme de localisation correspondant (Modules
sur les sites, section \ref{sites-modules-section-localize}). Nous
utiliserons le fait suivant : pour des $\mathcal{O}_U$-modules
$\mathcal{M}_i$, $i = 1, 2$, et un $\mathcal{O}$-module $\mathcal{A}$, il
existe une application canonique
$$
j_!
: \Hom_{\mathcal{O}_U}(
\mathcal{M}_1 \otimes_{\mathcal{O}_U} \mathcal{A}|_U, \mathcal{M}_2)
\longrightarrow
\Hom_\mathcal{O}(
j_!\mathcal{M}_1 \otimes_\mathcal{O} \mathcal{A}, j_!\mathcal{M}_2)
$$
Nous avons notamment
$j_!(\mathcal{M}_1 \otimes_{\mathcal{O}_U} \mathcal{A}|_U) =
j_!\mathcal{M}_1 \otimes_\mathcal{O} \mathcal{A}$ par
Modules sur les sites, lemme \ref{sites-modules-lemma-j-shriek-and-tensor}.

\medskip\noindent
Soit $\mathcal{A}$ une $\mathcal{O}$-algèbre différentielle graduée.
Notons $\mathcal{A}_U$ la restriction de $\mathcal{A}$ à
$\mathcal{C}/U$, en d'autres termes, nous avons
$\mathcal{A}_U = j^*\mathcal{A} = j^{-1}\mathcal{A}$. Dans la section
\ref{sdga-section-functoriality-dg}, nous avons
construit des foncteurs adjoints
$$ j_*
:
\textit{Mod}^{dg}(\mathcal{A}_U, \text{d})
\longrightarrow
\textit{Mod}^{dg}(\mathcal{A}, \text{d})
\quad\text{et}\quad j^*
:
\textit{Mod}^{dg}(\mathcal{A}, \text{d})
\longrightarrow
\textit{Mod}^{dg}(\mathcal{A}_U, \text{d})
$$ avec
$j^*$ adjoint à gauche à $j_*$. Nous affirmons qu'il existe de plus un foncteur
exact
$$ j_!
:
\textit{Mod}^{dg}(\mathcal{A}_U, \text{d})
\longrightarrow
\textit{Mod}^{dg}(\mathcal{A}, \text{d})
$$ adjoint à droite
de $j_*$. À savoir, étant donné un
$\mathcal{A}_U$-module différentiel gradué $\mathcal{M}$, nous définissons $j_!\mathcal{M}$ comme
le $\mathcal{A}$-module gradué construit dans la section
\ref{sdga-section-localize-graded} avec les
différentiels
$j_!\text{d} : j_!\mathcal{M}^n \to j_!\mathcal{M}^{n + 1}$. Étant donnée
une application homogène
$f : \mathcal{M} \to \mathcal{M}'$ de degré $n$ de
$\mathcal{A}_U$-modules différentiels gradués, nous obtenons une application
homogène $j_!f : j_!\mathcal{M} \to j_!\mathcal{M}'$ de degré $n$
de $\mathcal{A}$-modules différentiels gradués. Nous omettons
la simple vérification que cette
construction est compatible avec les différentiels. Ainsi
nous obtenons notre foncteur.

\begin{lemma}
\label{sdga-lemma-extension-by-zero-dg}
Dans la situation ci-dessus, nous
avons $$
\Hom_{\textit{Mod}^{dg}(\mathcal{A}, \text{d})}(
j_!\mathcal{M}, \mathcal{N})
= \Hom_{\textit{Mod}^{dg}(\mathcal{A}_U, \text{d})}(
\mathcal{M}, j^*\mathcal{N})
$$
\end{lemma}

\begin{proof}
Démonstration omise. Nous avons vu dans le lemme \ref{sdga-lemma-extension-by-zero-graded}
que le lemme est vrai au niveau gradué. Ainsi, tout ce qui doit être
vérifié est que l'isomorphisme résultant est compatible avec les différentiels.
\end{proof}

\begin{lemma}
\label{sdga-lemma-tensor-with-extension-by-zero-dg} Dans
la situation ci-dessus, soit $\mathcal{M}$ un
$\mathcal{A}_U$-module différentiel gradué à droite et soit $\mathcal{N}$ un
$\mathcal{A}$-module différentiel gradué à gauche. Alors
$$
j_!\mathcal{M} \otimes_\mathcal{A} \mathcal{N} =
j_!(\mathcal{M} \otimes_{\mathcal{A}_U} \mathcal{N}|_U)
$$ comme
complexes de $\mathcal{O}$-modules,
fonctoriellement en $\mathcal{M}$ et $\mathcal{N}$.
\end{lemma}

\begin{proof} Au
niveau des modules gradués, cela découle
du lemme \ref{sdga-lemma-tensor-with-extension-by-zero}.
Nous omettons la vérification que cet isomorphisme
est compatible avec les différentiels.
\end{proof}








\section{Foncteurs de décalage sur les faisceaux de modules différentiels gradués}
\label{sdga-section-shift-dg}

\noindent
Soit $(\mathcal{C}, \mathcal{O})$ un site annelé. Soit
$\mathcal{A}$ un faisceau d'algèbres différentielles graduées
sur $(\mathcal{C}, \mathcal{O})$.
Soit $\mathcal{M}$ un $\mathcal{A}$-module différentiel gradué.
Soit $k \in \mathbf{Z}$. Nous définissons le {\it $k$-ième décalage de} $\mathcal{M}$,
noté $\mathcal{M}[k]$, comme
suit : \begin{enumerate}
\item le $\mathcal{A}$-module gradué $\mathcal{M}[k]$ est
celui défini dans la section \ref{sdga-section-shift},
\item le différentiel
$d_{\mathcal{M}[k]} : (\mathcal{M}[k])^n \to (\mathcal{M}[k])^{n + 1}$
est donné par
$(-1)^k\text{d}_\mathcal{M} : \mathcal{M}^{n + k} \to \mathcal{M}^{n + k + 1}$.
\end{enumerate} Pour
un homomorphisme $f : \mathcal{L} \to \mathcal{M}$ de $\mathcal{A}$-modules homogène
de degré $n$, définissons $f[k] : \mathcal{L}[k] \to \mathcal{M}[k]$ par les
mêmes applications composantes que $f$. Alors $f[k]$ est
un homomorphisme de $\mathcal{A}$-modules homogène de degré $n$.
Cela définit
une application $$
\Hom_{\textit{Mod}^{dg}(\mathcal{A}, \text{d})}(\mathcal{L}, \mathcal{M})
\longrightarrow
\Hom_{\textit{Mod}^{dg}(\mathcal{A}, \text{d})}
(\mathcal{L}[k], \mathcal{M}[k])
$$
compatible avec les différentiels, car les
différentiels de $\mathcal{L}$ et
$\mathcal{M}$ sont multipliés par le même signe. Ces choix sont compatibles avec le choix
en Algèbre différentielle graduée,
définition \ref{dga-definition-shift}. Il est clair que nous avons
défini un foncteur
$$ [k]
:
\textit{Mod}^{dg}(\mathcal{A}, \text{d})
\longrightarrow
\textit{Mod}^{dg}(\mathcal{A}, \text{d})
$$ Ce foncteur est un foncteur
de catégories différentielles graduées et l'on a
$[k + l] = [k] \circ [l]$.

\medskip\noindent
Nous affirmons que l'isomorphisme
$$
\Hom_{\textit{Mod}^{dg}(\mathcal{A}, \text{d})}(\mathcal{L}, \mathcal{M}[k])
=
\Hom_{\textit{Mod}^{dg}(\mathcal{A}, \text{d})}(\mathcal{L}, \mathcal{M})[k]
$$ défini à la
section \ref{sdga-section-shift} sur les modules gradués sous-jacents est
compatible avec les
différentiels. Pour le voir, soit une application $\mathcal{A}$-linéaire
$f : \mathcal{L} \to \mathcal{M}[k]$ homogène de degré $n$ ; c'est un élément de degré
$n$ du membre de gauche. Notons
$f' : \mathcal{L} \to \mathcal{M}$ l'homomorphisme homogène de $\mathcal{A}$-modules de degré
$n + k$ dont les applications composantes sont les {\bf mêmes} que celles de $f$.
Selon nos conventions, c'est l'élément correspondant de degré $n$ du membre
de droite.
Par définition du différentiel du membre de gauche, nous
obtenons $$
\text{d}_{LHS}(f)
= \text{d}_{\mathcal{M}[k]} \circ f - (-1)^n f \circ \text{d}_\mathcal{L}
= (-1)^k\text{d}_\mathcal{M} \circ f - (-1)^n f \circ \text{d}_\mathcal{L}
$$
et, pour le différentiel du membre de droite, nous
obtenons $$
\text{d}_{RHS}(f')
= (-1)^k\left(
\text{d}_\mathcal{M} \circ f' - (-1)^{n + k} f' \circ \text{d}_\mathcal{L}
\right)
= (-1)^k\text{d}_\mathcal{M} \circ f' - (-1)^n f' \circ \text{d}_\mathcal{L}
$$ Les
applications composantes sont les mêmes, ce qui achève la démonstration.










\section{La catégorie homotopique}
\label{sdga-section-homotopy}

\noindent
Cette section est l'analogue de
l'Algèbre différentielle graduée, section \ref{dga-section-homotopy}.

\begin{definition}
\label{sdga-definition-homotopy} Soit
$(\mathcal{C}, \mathcal{O})$ un site annelé. Soit
$\mathcal{A}$ un faisceau d'algèbres différentielles graduées sur
$(\mathcal{C}, \mathcal{O})$. Soient
$f, g : \mathcal{M} \to \mathcal{N}$ des
homomorphismes de modules différentiels gradués sur $\mathcal{A}$. Une
{\it homotopie entre $f$ et $g$} est un morphisme gradué de $\mathcal{A}$-modules
$h : \mathcal{M} \to \mathcal{N}$, homogène de degré $-1$, tel
que
$$ f - g
= \text{d}_\mathcal{N} \circ h + h \circ \text{d}_\mathcal{M}
$$ S'il existe
une telle homotopie, nous disons que $f$ et $g$ sont {\it homotopes}.
\end{definition}

\noindent
Dans la situation de la définition, si nous avons des
applications $a : \mathcal{K} \to \mathcal{M}$ et $c : \mathcal{N} \to \mathcal{L}$
alors nous voyons
que $$
\begin{matrix}
h\text{ est une homotopie} \\
\text{ entre }f\text{ et } g
\end{matrix}
\quad
\Rightarrow
\quad
\begin{matrix}
c \circ h \circ a\text{ est une homotopie}\\
\text{entre } c
\circ f \circ a\text{ et } c\circ g \circ a
\end{matrix}
$$ Ainsi,
nous pouvons définir la composition des classes d'homotopie des morphismes
dans $\textit{Mod}(\mathcal{A}, \text{d})$.

\begin{definition}
\label{sdga-definition-complexes-notation}
Soit $(\mathcal{C}, \mathcal{O})$ un site annelé. Soit
$\mathcal{A}$ un faisceau d'algèbres différentielles graduées
sur $(\mathcal{C}, \mathcal{O})$.
La {\it catégorie homotopique}, notée $K(\textit{Mod}(\mathcal{A}, \text{d}))$, est
la catégorie dont les objets sont ceux de
$\textit{Mod}(\mathcal{A}, \text{d})$ et dont les morphismes sont les classes d'homotopie
d'homomorphismes de $\mathcal{A}$-modules différentiels gradués.
\end{definition}

\noindent
La notation $K(\textit{Mod}(\mathcal{A}, \text{d}))$
n'est pas standard, mais elle
est compatible avec les autres usages de $K(-)$ dans le projet Champs.

\medskip\noindent
Dans Algèbre différentielle graduée, définition
\ref{dga-definition-homotopy-category-of-dga-category},
nous avons défini la catégorie des complexes
$\text{Comp}(\mathcal{S})$ et la catégorie
homotopique $K(\mathcal{S})$ d'une catégorie
différentielle graduée $\mathcal{S}$. En appliquant
cela à la catégorie différentielle graduée
$\textit{Mod}^{dg}(\mathcal{A}, \text{d})$, nous obtenons
$$
\textit{Mod}(\mathcal{A}, \text{d}) =
\text{Comp}(\textit{Mod}^{dg}(\mathcal{A}, \text{d}))
$$ (voir la
discussion de la section \ref{sdga-section-dgm-dg-cat}) et nous obtenons
$$
K(\textit{Mod}(\mathcal{A}, \text{d})) =
K(\textit{Mod}^{dg}(\mathcal{A}, \text{d}))
$$
Pour vérifier cette dernière égalité, remarquons que nous avons
$$
\text{d}_{\Hom_{\textit{Mod}^{dg}(\mathcal{A}, \text{d})}
(\mathcal{M}, \mathcal{N})}(h) =
\text{d}_\mathcal{N} \circ h + h \circ \text{d}_\mathcal{M}
$$ lorsque
$h$ est comme dans la définition \ref{sdga-definition-homotopy}. Nous omettons les
détails.

\begin{lemma}
\label{sdga-lemma-homotopy-direct-sums} Soit
$(\mathcal{C}, \mathcal{O})$ un site annelé. Soit
$\mathcal{A}$ un faisceau d'algèbres différentielles graduées sur
$(\mathcal{C}, \mathcal{O})$. La catégorie
homotopique $K(\textit{Mod}(\mathcal{A}, \text{d}))$ possède des
sommes directes et des produits.
\end{lemma}

\begin{proof}
Démonstration omise. Indication : utiliser les sommes directes et
les produits du lemme \ref{sdga-lemma-dgm-abelian}. Cela fonctionne car nous
avons vu que ces foncteurs commutent au foncteur d'oubli vers la catégorie
des modules gradués sur $\mathcal{A}$, et car $\prod$ et $\bigoplus$ sont
des foncteurs exacts sur la catégorie des familles de groupes abéliens.
\end{proof}










\section{Cônes et triangles}
\label{sdga-section-conclude-triangulated}

\noindent
Dans cette section, nous utilisons les
résultats d'Algèbre différentielle graduée, section \ref{dga-section-review},
pour conclure que la catégorie homotopique
de la catégorie des $\mathcal{A}$-modules différentiels
gradués est triangulée.

\begin{lemma}
\label{sdga-lemma-axioms-AB}
Soit $(\mathcal{C}, \mathcal{O})$ un site annelé. Soit
$\mathcal{A}$ un faisceau d'algèbres différentielles graduées
sur $(\mathcal{C}, \mathcal{O})$. La
catégorie différentielle graduée
$\textit{Mod}^{dg}(\mathcal{A}, \text{d})$
satisfait aux axiomes (A) et (B)
d'Algèbre différentielle graduée, section \ref{dga-section-review}.
\end{lemma}

\begin{proof}
Supposons donnés des $\mathcal{A}$-modules différentiels
gradués $\mathcal{M}$ et $\mathcal{N}$. Considérons
le $\mathcal{A}$-module différentiel gradué $\mathcal{M} \oplus \mathcal{N}$
défini de manière évidente. Les coprojections
$i : \mathcal{M} \to \mathcal{M} \oplus \mathcal{N}$ et
$j : \mathcal{N} \to \mathcal{M} \oplus \mathcal{N}$ ainsi que les
projections
$p : \mathcal{M} \oplus \mathcal{N} \to \mathcal{N}$ et
$q : \mathcal{M} \oplus \mathcal{N} \to \mathcal{M}$ sont des morphismes
de modules différentiels gradués sur $\mathcal{A}$. Ainsi,
$i, j, p, q$ sont homogènes de degré
$0$ et fermés, c'est-à-dire que $\text{d}(i) = 0$, etc. Ainsi,
cette somme directe est une somme différentielle graduée au sens de
l'Algèbre différentielle graduée, définition \ref{dga-definition-dg-direct-sum}.
Cela prouve l'axiome (A).

\medskip\noindent
L'axiome (B) a été vérifié dans la section \ref{sdga-section-shift-dg}.
\end{proof}

\noindent
Soit $(\mathcal{C}, \mathcal{O})$ un site annelé. Soit
$\mathcal{A}$ un faisceau d'algèbres différentielles graduées sur
$(\mathcal{C}, \mathcal{O})$.
Rappelons qu'une suite
$$ 0
\to \mathcal{K} \to \mathcal{L} \to \mathcal{N} \to 0
$$ dans
$\textit{Mod}(\mathcal{A}, \text{d})$ est appelée
une {\it suite exacte courte admissible} (voir
Algèbre différentielle graduée, section \ref{dga-section-review}) si elle
est scindée dans $\textit{Mod}(\mathcal{A})$. Autrement dit, elle est scindée comme
suite de $\mathcal{A}$-modules gradués. Notons
$s : \mathcal{N} \to \mathcal{L}$ et
$\pi : \mathcal{L} \to \mathcal{K}$ les scindages
gradués $\mathcal{A}$-linéaires. En combinant
le lemme \ref{sdga-lemma-axioms-AB} et Algèbre
différentielle graduée, lemme \ref{dga-lemma-get-triangle}, nous obtenons
un triangle
$$
\mathcal{K} \to \mathcal{L} \to \mathcal{N} \to \mathcal{K}[1]
$$ où la flèche
$\mathcal{N} \to \mathcal{K}[1]$ dans la preuve d'Algèbre
différentielle graduée, lemme \ref{dga-lemma-get-triangle} est construite
comme
$$
\delta = \pi \circ
\text{d}_{\Hom_{\textit{Mod}^{dg}(\mathcal{A}, \text{d})}
(\mathcal{L}, \mathcal{M})}(s) =
\pi \circ \text{d}_\mathcal{L} \circ s -
\pi \circ s \circ \text{d}_\mathcal{N} =
\pi \circ \text{d}_\mathcal{L} \circ s
$$ avec nos excuses pour cette
notation peu lisible. Quoi qu'il en soit, nous voyons que, dans notre
situation, l'application de bord $\delta$ construite dans Algèbre
différentielle graduée, lemme \ref{dga-lemma-get-triangle}, coïncide
sur les complexes sous-jacents de $\mathcal{O}$-modules avec
l'application de bord usuelle employée dans le projet Champs pour
les suites exactes courtes de complexes scindées terme à terme ;
voir Catégories dérivées, définition \ref{derived-definition-split-ses}.

\begin{definition}
\label{sdga-definition-cone}
Soit $(\mathcal{C}, \mathcal{O})$ un site annelé. Soit
$\mathcal{A}$ un faisceau d'algèbres différentielles graduées sur
$(\mathcal{C}, \mathcal{O})$. Soit
$f : \mathcal{K} \to \mathcal{L}$ un
homomorphisme de modules différentiels gradués sur $\mathcal{A}$. Le
{\it cône} de $f$ est le $\mathcal{A}$-module différentiel gradué
$C(f)$ défini comme suit :
\begin{enumerate}
\item le complexe sous-jacent de $\mathcal{O}$-modules est le cône
du morphisme correspondant
$f : \mathcal{K}^\bullet \to \mathcal{L}^\bullet$ de complexes de
$\mathcal{A}$-modules ; ainsi,
$C(f)^n = \mathcal{L}^n \oplus \mathcal{K}^{n + 1}$ et le différentiel
est
$$
d_{C(f)} =
\left(
\begin{matrix}
\text{d}_\mathcal{L} & f \\ 0 &
-\text{d}_\mathcal{K}
\end{matrix}
\right)
$$
\item l'application de multiplication
$$ C(f)^n
\times \mathcal{A}^m \to C(f)^{n + m}
$$ est la somme
directe de l'application de multiplication
$\mathcal{L}^n \times \mathcal{A}^m \to \mathcal{L}^{n + m}$ et
de l'application de multiplication
$\mathcal{K}^{n + 1} \times \mathcal{A}^m \to \mathcal{K}^{n + 1 + m}$.
\end{enumerate}
Il est équipé d'homomorphismes canoniques de
$\mathcal{A}$-modules différentiels gradués $i : \mathcal{L} \to C(f)$ et
$p : C(f) \to \mathcal{K}[1]$ induits par les applications évidentes.
\end{definition}

\noindent
Remarquons que, dans la situation de la définition, la suite
$$
0 \to \mathcal{L} \to C(f) \to \mathcal{K}[1] \to 0
$$
est une suite exacte courte admissible.

\begin{lemma}
\label{sdga-lemma-axiom-C}
Soit $(\mathcal{C}, \mathcal{O})$ un site annelé. Soit
$\mathcal{A}$ un faisceau d'algèbres différentielles graduées sur
$(\mathcal{C}, \mathcal{O})$. La
catégorie différentielle graduée
$\textit{Mod}^{dg}(\mathcal{A}, \text{d})$
satisfait à l'axiome (C) formulé dans
Algèbre différentielle graduée, situation \ref{dga-situation-ABC}.
\end{lemma}

\begin{proof}
Soit $f : \mathcal{K} \to \mathcal{L}$ un
homomorphisme de modules différentiels gradués sur $\mathcal{A}$. D'après
ce qui précède, nous avons une suite exacte courte admissible
$$
0 \to \mathcal{L} \to C(f) \to \mathcal{K}[1] \to 0
$$ Pour
achever la démonstration, nous devons vérifier que l'application de
bord $$
\delta : \mathcal{K}[1] \to \mathcal{L}[1]
$$
qui lui est associée (voir la discussion ci-dessus) est égale à $f[1]$.
Pour la section $s : \mathcal{K}[1] \to C(f)$, nous utilisons en
degré $n$ le plongement $\mathcal{K}^{n + 1} \to C(f)^n$. En degré
$n$, l'application $\pi$ est donnée par les projections
$C(f)^n \to \mathcal{L}^n$. Alors finalement nous devons calculer
$$
\delta = \pi \circ \text{d}_{C(f)} \circ s
$$ (voir
discussion ci-dessus). En notation matricielle, cette expression est
égale à $$
\left(
\begin{matrix}
1
& 0 \end{matrix}
\right)
\left(
\begin{matrix}
\text{d}_\mathcal{L} & f \\
0 & -\text{d}_\mathcal{K}
\end{matrix}
\right)
\left(
\begin{matrix} 0
\\
1
\end{matrix}
\right) = f
$$
comme souhaité.
\end{proof}

\noindent À
ce stade, nous savons que tous les lemmes prouvés en
Algèbre différentielle graduée, section \ref{dga-section-review} sont
valables pour la catégorie différentielle graduée
$\textit{Mod}^{dg}(\mathcal{A}, \text{d})$. En particulier,
nous obtenons la proposition suivante.

\begin{proposition}
\label{sdga-proposition-homotopy-category-triangulated}
Soit $(\mathcal{C}, \mathcal{O})$ un site annelé. Soit
$\mathcal{A}$ un faisceau d'algèbres différentielles graduées
sur $(\mathcal{C}, \mathcal{O})$.
La catégorie homotopique $K(\textit{Mod}(\mathcal{A}, \text{d}))$
est une catégorie triangulée
où \begin{enumerate}
\item les foncteurs de décalage sont ceux de la
section \ref{sdga-section-shift-dg},
\item les triangles distingués sont les triangles de
$K(\textit{Mod}(\mathcal{A}, \text{d}))$ isomorphes, en
tant que triangles, à un triangle
$$
\mathcal{K} \to \mathcal{L} \to \mathcal{N}
\xrightarrow{\delta} \mathcal{K}[1],\quad\quad
\delta = \pi \circ \text{d}_\mathcal{L} \circ s
$$ construit
à partir d'une suite exacte courte admissible
$0 \to \mathcal{K} \to \mathcal{L} \to \mathcal{N} \to 0$ dans
$\textit{Mod}(\mathcal{A}, \text{d})$ ci-dessus.
\end{enumerate}
\end{proposition}

\begin{proof}
Rappelons que $K(\textit{Mod}(\mathcal{A}, \text{d}))
= K(\textit{Mod}^{dg}(\mathcal{A}, \text{d}))$,
voir la section \ref{sdga-section-homotopy}.
La proposition découle
alors des lemmes \ref{sdga-lemma-axioms-AB} et \ref{sdga-lemma-axiom-C} et
d'Algèbre
différentielle graduée, proposition
\ref{dga-proposition-ABC-homotopy-category-triangulated}.
\end{proof}

\begin{remark}
\label{sdga-remark-cone-identity}
Soit $(\mathcal{C}, \mathcal{O})$ un site annelé. Soit
$\mathcal{A}$ un faisceau d'algèbres différentielles graduées
sur $(\mathcal{C}, \mathcal{O})$. Soit $C = C(\text{id}_\mathcal{A})$
le cône de l'application identité $\mathcal{A} \to \mathcal{A}$, considérée
comme un morphisme de $\mathcal{A}$-modules différentiels
gradués.
Alors $$
\Hom_{\textit{Mod}(\mathcal{A}, \text{d})}(C, \mathcal{M})
= \{(x, y) \in
\Gamma(\mathcal{C}, \mathcal{M}^0) \times
\Gamma(\mathcal{C}, \mathcal{M}^{-1}) \mid
x = \text{d}(y)\}
$$
où l'application de gauche à droite envoie $f$ à la paire $(x, y)$
où $x$ est l'image de la section globale $(0, 1)$
de $C^{-1} = \mathcal{A}^{-1} \oplus \mathcal{A}^0$ et
où $y$ est l'image de la section globale $(1, 0)$ de
$C^0 = \mathcal{A}^0 \oplus \mathcal{A}^1$.
\end{remark}

\begin{lemma}
\label{sdga-lemma-dgm-grothendieck-abelian}
Soit $(\mathcal{C}, \mathcal{O})$ un site annelé. Soit
$(\mathcal{A}, \text{d})$ une $\mathcal{O}$-algèbre différentielle graduée. La
catégorie $\textit{Mod}(\mathcal{A}, \text{d})$ est une
catégorie abélienne de Grothendieck.
\end{lemma}

\begin{proof}
Par le lemme \ref{sdga-lemma-dgm-abelian} et la définition d'une catégorie abélienne
de Grothendieck
(Injectifs, définition \ref{injectives-definition-grothendieck-conditions}), il
suffit de montrer
que $\textit{Mod}(\mathcal{A}, \text{d})$ possède
un générateur. Pour chaque objet $U$ de $\mathcal{C}$, notons
$C_U$ le cône du morphisme identité $\mathcal{A}_U \to \mathcal{A}_U$, comme
dans la remarque \ref{sdga-remark-cone-identity}. Nous affirmons
que $$
\mathcal{G} = \bigoplus\nolimits_{k, U} j_{U!}C_U[k]
$$
est un générateur, où la somme porte sur tous les objets $U$ de $\mathcal{C}$
et tous les $k \in \mathbf{Z}$. En
effet, étant donné un $\mathcal{A}$-module différentiel gradué $\mathcal{M}$,
s'il n'existe aucun morphisme non nul de $\mathcal{G}$ vers $\mathcal{M}$,
alors, pour tous $k$ et $U$,
le groupe \begin{align*}
\Hom_{\textit{Mod}(\mathcal{A})}(j_{U!}C_U[k], \mathcal{M}) \\
&
= \Hom_{\textit{Mod}(\mathcal{A}_U)}(C_U[k], \mathcal{M}|_U) \\
&
= \{(x, y) \in \mathcal{M}^{-k}(U) \times \mathcal{M}^{-k - 1}(U) \mid
x = \text{d}(y)\}
\end{align*}
est égal à zéro. Ainsi, $\mathcal{M}$ est nul.
\end{proof}











\section{Résolutions plates}
\label{sdga-section-P-resolutions}

\noindent
Cette section est l'analogue de
l'Algèbre différentielle graduée, section \ref{dga-section-P-resolutions}.

\medskip\noindent
Soit $(\mathcal{C}, \mathcal{O})$ un site annelé. Soit
$\mathcal{A}$ un faisceau d'algèbres différentielles graduées sur
$(\mathcal{C}, \mathcal{O})$. Nous dirons qu'un
$\mathcal{A}$-module différentiel gradué à droite $\mathcal{P}$ est {\it bon}
si
\begin{enumerate}
\item le foncteur
$\mathcal{N} \mapsto \mathcal{P} \otimes_\mathcal{A} \mathcal{N}$ est
exact sur la catégorie des $\mathcal{A}$-modules gradués à gauche,
\item si $\mathcal{N}$ est un
$\mathcal{A}$-module différentiel gradué acyclique à gauche, alors
$\mathcal{P} \otimes_\mathcal{A} \mathcal{N}$ est acyclique,
\item pour tout morphisme $(f, f^\sharp) : (\Sh(\mathcal{C}'), \mathcal{O}')
\to (\Sh(\mathcal{C}), \mathcal{O})$ de topos
annelés, toute $\mathcal{O}'$-algèbre différentielle graduée
$\mathcal{A}'$ et tout morphisme $\varphi : f^{-1}\mathcal{A} \to \mathcal{A}'$ d'algèbres
différentielles graduées sur $f^{-1}\mathcal{O}_\mathcal{D}$, les propriétés
(1) et (2) valent pour l'image inverse $f^*\mathcal{P}$
(section \ref{sdga-section-functoriality-dg}),
considérée comme un $\mathcal{A}'$-module différentiel gradué.
\end{enumerate}
La première condition signifie que $\mathcal{P}$ est plat comme module gradué
à droite sur $\mathcal{A}$ ; la deuxième signifie que $\mathcal{P}$ est
K-plat au sens de Spaltenstein (voir Cohomologie
sur les sites, section \ref{sites-cohomology-section-flat}) ; la troisième
exige que ces propriétés restent vraies après tout changement de base.

\medskip\noindent Il
est peut-être surprenant qu'il y ait beaucoup de bons modules.

\begin{lemma}
\label{sdga-lemma-supply-good}
Soit $(\mathcal{C}, \mathcal{O})$ un site annelé. Soit
$\mathcal{A}$ un faisceau d'algèbres différentielles graduées
sur $(\mathcal{C}, \mathcal{O})$. Soit $U \in \Ob(\mathcal{C})$.
Alors $j_!\mathcal{A}_U$ est un bon
$\mathcal{A}$-module différentiel gradué.
\end{lemma}

\begin{proof}
Soit $\mathcal{N}$ un module gradué à gauche sur $\mathcal{A}$.
Par le lemme \ref{sdga-lemma-tensor-with-extension-by-zero} nous avons
$$
j_!\mathcal{A}_U \otimes_\mathcal{A} \mathcal{N} =
j_!(\mathcal{A}_U \otimes_{\mathcal{A}_U} \mathcal{N}|_U) =
j_!(\mathcal{N}_U)
$$ comme modules
gradués. Puisque la
restriction à $U$ et le foncteur $j_!$ sont tous deux exacts, cela
prouve (1). Le même argument est valable pour (2) en utilisant
le lemme \ref{sdga-lemma-tensor-with-extension-by-zero-dg}.

\medskip\noindent
Considérons un morphisme $(f, f^\sharp) : (\Sh(\mathcal{C}'), \mathcal{O}')
\to (\Sh(\mathcal{C}), \mathcal{O})$ de
topos annelés, une $\mathcal{O}'$-algèbre différentielle graduée
$\mathcal{A}'$ et un morphisme $\varphi : f^{-1}\mathcal{A} \to \mathcal{A}'$
de $f^{-1}\mathcal{O}$-algèbres différentielles graduées.
Nous devons montrer
que $$
f^*j_!\mathcal{A}_U
= f^{-1}j_!\mathcal{A}_U \otimes_{f^{-1}\mathcal{A}} \mathcal{A}'
$$
satisfait aux points (1) et (2) pour le
topos annelé $(\Sh(\mathcal{C}'), \mathcal{O}')$ doté
du faisceau de $\mathcal{O}'$-algèbres différentielles
graduées $\mathcal{A}'$. Pour le prouver, nous
pouvons remplacer $(\Sh(\mathcal{C}), \mathcal{O})$
et $(\Sh(\mathcal{C}'), \mathcal{O}')$ par des topos
annelés
équivalents. Ainsi, par
Modules sur les sites, lemme \ref{sites-modules-lemma-morphism-ringed-topoi-comes-from-morphism-ringed-sites},
nous pouvons supposer que $f$ provient
d'un morphisme de sites $f : \mathcal{C} \to \mathcal{C}'$
donné par un foncteur continu $u : \mathcal{C} \to \mathcal{C}'$.
Dans ce cas, posons $U' = u(U)$
et notons $j' : \Sh(\mathcal{C}'/U') \to \Sh(\mathcal{C}')$
le morphisme de localisation
correspondant. On obtient un carré commutatif de morphismes
de topos annelés $$
\xymatrix{
(\Sh(\mathcal{C}'/U'), \mathcal{O}'_{U'})
\ar[rr]_{(j', (j')^\sharp)} \ar[d]_{(f', (f')^\sharp)}
& & (\Sh(\mathcal{C}'), \mathcal{O}')
\ar[d]^{(f, f^\sharp)} \\
(\Sh(\mathcal{C}/U), \mathcal{O}_U)
\ar[rr]^{(j, j^\sharp)}
& & (\Sh(\mathcal{C}), \mathcal{O}).
}
$$
et nous avons $f'_*(j')^{-1} = j^{-1}f_*$.
Voir Modules sur
les sites, lemme \ref{sites-modules-lemma-localize-morphism-ringed-sites}.
Par l'unicité des adjoints,
on obtient $f^{-1}j_! = j'_!(f')^{-1}$.
Ainsi nous obtenons \begin{align*}
f^*j_!\mathcal{A}_U
&
= f^{-1}j_!\mathcal{A}_U \otimes_{f^{-1}\mathcal{A}} \mathcal{A}' \\
&
= j'_!(f')^{-1}\mathcal{A}_U \otimes_{f^{-1}\mathcal{A}} \mathcal{A}' \\
&
= j'_!\left(
(f')^{-1}\mathcal{A}_U \otimes_{f^{-1}\mathcal{A}|_{U'}}
\mathcal{A}'|_{U'}\right) \\
&
= j'_!\mathcal{A}'_{U'}
\end{align*}
La première équation est la définition de l'image inverse de $j_!\mathcal{A}_U$
comme module différentiel gradué sur $\mathcal{A}'$.
La deuxième égalité découle de $f^{-1}j_! = j'_!(f')^{-1}$.
La troisième découle du lemme \ref{sdga-lemma-tensor-with-extension-by-zero-dg},
appliqué au site annelé $(\mathcal{C}', f^{-1}\mathcal{O})$ muni
du faisceau d'algèbres différentielles graduées $f^{-1}\mathcal{A}$,
avec les modules différentiels gradués $(f')^{-1}\mathcal{A}_U$ sur $\mathcal{C}'/U'$
et $\mathcal{A}'$ sur $\mathcal{C}'$.
La quatrième égalité résulte de
$(f')^{-1}\mathcal{A}_U = f^{-1}\mathcal{A}|_{U'}$. Nous voyons donc que
l'image inverse est encore un module du même type, et nous
avons démontré ci-dessus les conditions (1) et (2) pour celui-ci.
\end{proof}

\begin{lemma}
\label{sdga-lemma-good-admissible-ses}
Soit $(\mathcal{C}, \mathcal{O})$ un site annelé.
Soit $\mathcal{A}$ un faisceau d'algèbres différentielles
graduées sur $(\mathcal{C}, \mathcal{O})$.
Soit $0 \to \mathcal{P} \to \mathcal{P}' \to \mathcal{P}'' \to 0$
une suite exacte courte admissible de
$\mathcal{A}$-modules différentiels gradués. Si deux modules sur trois sont
bons, le troisième aussi.
\end{lemma}

\begin{proof}
La condition (1) est immédiate, car la suite est scindée au niveau gradué.
Pour la condition (2), remarquons que, pour tout module
différentiel gradué à gauche sur $\mathcal{A}$, la suite
$$ 0
\to
\mathcal{P} \otimes_\mathcal{A} \mathcal{N} \to
\mathcal{P}' \otimes_\mathcal{A} \mathcal{N} \to
\mathcal{P}'' \otimes_\mathcal{A} \mathcal{N} \to 0
$$ est une
suite exacte courte admissible de
$\mathcal{O}$-modules différentiels gradués (puisque, après oubli du différentiel, le
produit tensoriel se calcule dans la catégorie des modules gradués). Ainsi,
si deux des trois sont exacts comme complexes de $\mathcal{O}$-modules, le
troisième aussi. Enfin, le même argument montre que, pour un
morphisme $(f, f^\sharp) : (\Sh(\mathcal{C}'), \mathcal{O}')
\to (\Sh(\mathcal{C}), \mathcal{O})$ de topos
annelés, une $\mathcal{O}'$-algèbre différentielle graduée
$\mathcal{A}'$ et un morphisme $\varphi : f^{-1}\mathcal{A} \to \mathcal{A}'$ de
$f^{-1}\mathcal{O}$-algèbres différentielles graduées,
la suite
$$ 0
\to f^*\mathcal{P} \to f^*\mathcal{P}' \to f^*\mathcal{P}'' \to 0
$$ est
une suite exacte courte admissible de
$\mathcal{A}'$-modules différentiels gradués et le même argument que ci-dessus s'applique ici.
\end{proof}

\begin{lemma}
\label{sdga-lemma-good-direct-sum} Soit
$(\mathcal{C}, \mathcal{O})$ un site annelé. Soit
$\mathcal{A}$ un faisceau d'algèbres différentielles graduées sur
$(\mathcal{C}, \mathcal{O})$. Une somme directe
arbitraire de bons $\mathcal{A}$-modules différentiels gradués est
bonne. Une limite inductive filtrante de bons
$\mathcal{A}$-modules différentiels gradués est bonne.
\end{lemma}

\begin{proof}
Démonstration omise. Indication : les sommes directes et les limites inductives
filtrantes commutent aux produits tensoriels et aux images inverses.
\end{proof}

\begin{lemma}
\label{sdga-lemma-good-quotient}
Soit $(\mathcal{C}, \mathcal{O})$ un site annelé.
Soit $\mathcal{A}$ un faisceau d'algèbres différentielles graduées
sur $(\mathcal{C}, \mathcal{O})$. Soit $\mathcal{M}$
un $\mathcal{A}$-module différentiel gradué. Il existe un homomorphisme
$\mathcal{P} \to \mathcal{M}$ de $\mathcal{A}$-modules différentiels gradués
avec les propriétés suivantes
\begin{enumerate}
\item $\mathcal{P} \to \mathcal{M}$ est surjectif,
\item $\Ker(\text{d}_\mathcal{P}) \to \Ker(\text{d}_\mathcal{M})$
est surjectif, et
\item $\mathcal{P}$ est bon.
\end{enumerate}
\end{lemma}

\begin{proof}
Considérons les triples $(U, k, x)$ où $U$ est un objet de $\mathcal{C}$,
$k \in \mathbf{Z}$, et $x$ est une section de $\mathcal{M}^k$ sur $U$ avec
$\text{d}_\mathcal{M}(x) = 0$. Nous obtenons alors un unique morphisme
de $\mathcal{A}_U$-modules différentiels gradués
$\varphi_x : \mathcal{A}_U[-k] \to \mathcal{M}|_U$
envoyant $1$ sur $x$. Ceci est adjoint à un morphisme
$\psi_x : j_{U!}\mathcal{A}_U[-k] \to \mathcal{M}$.
Remarquons que $1 \in \mathcal{A}_U(U)$ correspond à une
section $1 \in j_{U!}\mathcal{A}_U[-k](U)$ de degré $k$, dont le
différentiel est nul et qui est envoyée sur $x$ par $\psi_x$. Ainsi, si
nous considérons l'application
$$
\bigoplus\nolimits_{(U, k, x)} j_{U!}\mathcal{A}_U[-k]
\longrightarrow
\mathcal{M}
$$
les conditions (2) et (3) seront satisfaites. À savoir, les objets
$j_{U!}\mathcal{A}_U[-k]$ sont bons (lemme \ref{sdga-lemma-supply-good})
et toute somme directe de bons objets est bonne (lemme \ref{sdga-lemma-good-direct-sum}).

\medskip\noindent
Considérons ensuite les triples $(U, k, x)$ où $U$ est un objet de $\mathcal{C}$,
$k \in \mathbf{Z}$ et $x$ une section de $\mathcal{M}^k$, pas nécessairement annulée par
le différentiel. Ensuite, nous pouvons considérer le cône
$C_U$ de l'application identité $\mathcal{A}_U \to \mathcal{A}_U$ comme dans
la remarque \ref{sdga-remark-cone-identity}. L'élément $x$ déterminera une application
$\varphi_x : C_U[-k - 1] \to \mathcal{A}_U$, voir remarque
\ref{sdga-remark-cone-identity}. Maintenant, puisque nous
avons une suite exacte courte admissible
$$
0 \to \mathcal{A}_U \to C_U \to \mathcal{A}_U[1] \to 0
$$
nous concluons que $j_{U!}C_U$ est un bon module par le
lemme \ref{sdga-lemma-good-admissible-ses} et le
lemme \ref{sdga-lemma-supply-good} déjà utilisé. Comme
ci-dessus, nous concluons que la somme directe des applications
$\psi_x : j_{U!}C_U \to \mathcal{M}$ adjointes aux $\varphi_x$
$$
\bigoplus\nolimits_{(U, k, x)} j_{U!}C_U \longrightarrow \mathcal{M}
$$ est
surjective. En prenant la somme directe avec l'application produite
au premier paragraphe, nous concluons.
\end{proof}

\begin{remark}
\label{sdga-remark-sheaf-graded-sets}
Soit $(\mathcal{C}, \mathcal{O})$ un site annelé. Un
{\it faisceau d'ensembles gradués} sur $\mathcal{C}$ est un faisceau
d'ensembles $\mathcal{S}$ muni d'un morphisme
$\deg : \mathcal{S} \to \underline{\mathbf{Z}}$
de faisceaux d'ensembles. Notons $\mathcal{O}[\mathcal{S}]$
le $\mathcal{O}$-module gradué qui est le
$\mathcal{O}$-module libre sur le faisceau d'ensembles gradués $\mathcal{S}$.
Plus précisément, la composante graduée de degré $n$ de
$\mathcal{O}[\mathcal{S}]$ est le faisceau associé au préfaisceau
$$
U \longmapsto
\bigoplus\nolimits_{s \in \mathcal{S}(U),\ \deg(s) = n} s \cdot \mathcal{O}(U)
$$
Avec le différentiel nul, nous pouvons aussi le considérer comme un
$\mathcal{O}$-module différentiel gradué.
Soit $\mathcal{A}$ un faisceau d'$\mathcal{O}$-algèbres graduées.
Nous définissons de même $\mathcal{A}[\mathcal{S}]$ comme le
$\mathcal{A}$-module gradué dont la composante graduée de degré $n$ est le
faisceau associé au préfaisceau
$$
U \longmapsto
\bigoplus\nolimits_{s \in \mathcal{S}(U)} s \cdot \mathcal{A}^{n - \deg(s)}(U)
$$
Si $\mathcal{A}$ est une $\mathcal{O}$-algèbre différentielle graduée,
nous en faisons un $\mathcal{O}$-module différentiel gradué
en posant $\text{d}(s) = 0$ pour tout $s \in \mathcal{S}(U)$,
puis en prenant le faisceau associé.
\end{remark}

\begin{lemma}
\label{sdga-lemma-free-graded-module-good}
Soit $(\mathcal{C}, \mathcal{O})$ un site annelé. Soit
$\mathcal{A}$ une $\mathcal{A}$-algèbre différentielle graduée. Soit
$\mathcal{S}$ un faisceau d'ensembles gradués sur $\mathcal{C}$.
Alors le module gradué libre $\mathcal{A}[\mathcal{S}]$
sur $\mathcal{S}$, muni du différentiel de la remarque
\ref{sdga-remark-sheaf-graded-sets}, est
un bon module différentiel gradué sur $\mathcal{A}$.
\end{lemma}

\begin{proof}
Soit $\mathcal{N}$ un module gradué à gauche sur $\mathcal{A}$.
Alors nous
avons $$
\mathcal{A}[\mathcal{S}] \otimes_\mathcal{A} \mathcal{N} =
\mathcal{O}[\mathcal{S}] \otimes_\mathcal{O} \mathcal{N} =
\mathcal{N}[\mathcal{S}]
$$
où $\mathcal{N}[\mathcal{S}$ est le module
gradué sur $\mathcal{O}$ dont la composante de degré $n$
est le faisceau associé au préfaisceau
$$ U
\longmapsto
\bigoplus\nolimits_{s \in \mathcal{S}(U)} s \cdot \mathcal{N}^{n - \deg(s)}(U)
$$ Il est
clair que $\mathcal{N} \to \mathcal{N}[\mathcal{S}]$ est un foncteur
exact ; ainsi, $\mathcal{A}[\mathcal{S}$ est plat comme module
gradué sur $\mathcal{A}$. Supposons maintenant que $\mathcal{N}$ soit un
module différentiel gradué à gauche sur $\mathcal{A}$. Alors nous avons
$$
H^*(\mathcal{A}[\mathcal{S}] \otimes_\mathcal{A} \mathcal{N}) =
H^*(\mathcal{O}[\mathcal{S}] \otimes_\mathcal{O} \mathcal{N})
$$
comme faisceaux gradués de $\mathcal{O}$-modules, ce qui, par platitude
(sur $\mathcal{O})$ est égal à
$$
H^*(\mathcal{N})[\mathcal{S}]
$$ comme module
gradué sur $\mathcal{O}$. Par conséquent, si $\mathcal{N}$ est acyclique,
alors $\mathcal{A}[\mathcal{S}] \otimes_\mathcal{A} \mathcal{N}$ est
acyclique.

\medskip\noindent
Enfin, considérons un morphisme
$(f, f^\sharp) : (\Sh(\mathcal{C}'), \mathcal{O}')
\to (\Sh(\mathcal{C}), \mathcal{O})$ de
topos annelés, une $\mathcal{O}'$-algèbre différentielle graduée
$\mathcal{A}'$, et une application $\varphi : f^{-1}\mathcal{A} \to \mathcal{A}'$
de $f^{-1}\mathcal{O}$-algèbres différentielles
graduées. Il est donc facile de voir que
$$
f^*\mathcal{A}[\mathcal{S}] = \mathcal{A}'[f^{-1}\mathcal{S}]
$$ ce
qui achève la démonstration que notre module est bon.
\end{proof}

\begin{lemma}
\label{sdga-lemma-resolve}
Soit $(\mathcal{C}, \mathcal{O})$ un site annelé.
Soit $\mathcal{A}$ un faisceau d'algèbres différentielles graduées
sur $(\mathcal{C}, \mathcal{O})$. Soit $\mathcal{M}$
un $\mathcal{A}$-module différentiel gradué. Il existe un homomorphisme
$\mathcal{P} \to \mathcal{M}$ de $\mathcal{A}$-modules différentiels gradués
avec les propriétés suivantes
\begin{enumerate}
\item $\mathcal{P} \to \mathcal{M}$ est un quasi-isomorphisme, et
\item $\mathcal{P}$ est bon.
\end{enumerate}
\end{lemma}

\begin{proof}[Première démonstration]
Soit $\mathcal{S}_0$ le faisceau d'ensembles
gradués (remarque \ref{sdga-remark-sheaf-graded-sets})
dont la composante de degré $n$ est $\Ker(\text{d}_\mathcal{M}^n)$.
Considérons l'homomorphisme de modules différentiels gradués
$$
\mathcal{P}_0 = \mathcal{A}[\mathcal{S}_0] \longrightarrow \mathcal{M}
$$ où
le membre de gauche est défini dans la remarque \ref{sdga-remark-sheaf-graded-sets} et
le morphisme envoie une section locale $s$ de $\mathcal{S}_0$ sur
la section locale correspondante de $\mathcal{M}^{\deg(s)}$ ; celle-ci
appartient au noyau du différentiel, ce qui assure que nous avons
bien un morphisme de modules différentiels gradués. Par construction, les
applications induites sur les faisceaux de cohomologie $H^n(\mathcal{P}_0) \to H^n(\mathcal{M})$
sont surjectives. Nous allons construire par récurrence
des applications $$
\mathcal{P}_0 \to \mathcal{P}_1 \to \mathcal{P}_2 \to \ldots \to \mathcal{M}
$$
Notons bien sûr que $H^*(\mathcal{P}_i) \to H^*(\mathcal{M})$ sera
surjectif pour tous les $i$. Étant
donné $\mathcal{P}_i \to \mathcal{M}$, notons $\mathcal{S}_{i + 1}$ le faisceau
d'ensembles gradués dont la composante de degré $n$
est $$
\Ker(\text{d}_{\mathcal{P}_i}^{n + 1})
\times_{\mathcal{M}^{n + 1}, \text{d}}
\mathcal{M}^n
$$ Posons
alors
$$
\mathcal{P}_{i + 1} = \mathcal{P}_i \oplus \mathcal{A}[\mathcal{S}_{i + 1}]
$$ comme module
gradué sur $\mathcal{A}$, avec le différentiel et le morphisme vers
$\mathcal{M}$ définis comme suit :
\begin{enumerate}
\item sur les sections locales de $\mathcal{P}_i$, on utilise le différentiel de
$\mathcal{P}_i$ et le morphisme donné vers $\mathcal{M}$,
\item pour une section locale $s = (p, m)$ de $\mathcal{S}_{i + 1}$, on pose
$\text{d}(s)$ égal à $p$, vu comme une section de $\mathcal{P}_i$ de degré
$\deg(s) + 1$, et on envoie $s$ sur $m$ dans $\mathcal{M}$, et
\item on prolonge le différentiel de manière unique de façon que la règle de Leibniz soit satisfaite.
\end{enumerate}
Cette construction est bien définie, car $\text{d}(m)$ est l'image de $p$
et $\text{d}(p) = 0$.
Enfin, posons $\mathcal{P} = \colim \mathcal{P}_i$, muni du
morphisme induit vers $\mathcal{M}$.

\medskip\noindent
L'application $\mathcal{P} \to \mathcal{M}$ est un quasi-isomorphisme :
nous avons $H^n(\mathcal{P}) = \colim H^n(\mathcal{P}_i)$
et pour chaque $i$ l'application $H^n(\mathcal{P}_i) \to H^n(\mathcal{M})$
est surjective avec le noyau anéanti par l'application
$H^n(\mathcal{P}_i) \to H^n(\mathcal{P}_{i + 1})$ par construction.
Chaque $\mathcal{P}_i$ est bon parce que $\mathcal{P}_0$ est bon
par le lemme \ref{sdga-lemma-free-graded-module-good} et chaque $\mathcal{P}_{i + 1}$
est au milieu de la suite exacte courte admissible
$0 \to \mathcal{P}_i \to \mathcal{P}_{i + 1} \to
\mathcal{A}[\mathcal{S}_{i + 1}] \to 0$ dont les termes
extérieurs sont bons par induction. Par conséquent,
$\mathcal{P}_{i + 1}$ est bon par
le lemme \ref{sdga-lemma-good-admissible-ses}.
Enfin, nous concluons que $\mathcal{P}$ est bon par
le lemme \ref{sdga-lemma-good-direct-sum}.
\end{proof}

\begin{proof}[Seconde démonstration]
Nous invitons le lecteur à lire la preuve d'Algèbre différentielle
graduée, lemme \ref{dga-lemma-resolve} avant de lire cette preuve.
Posons $\mathcal{M} = \mathcal{M}_0$.
Nous choisissons par récurrence des suites
exactes courtes $$
0 \to \mathcal{M}_{i + 1} \to \mathcal{P}_i \to \mathcal{M}_i \to
0 $$
où
les applications $\mathcal{P}_i \to \mathcal{M}_i$ sont
choisies comme dans le lemme \ref{sdga-lemma-good-quotient}.
Cela
fournit une « résolution » $$
\ldots \to \mathcal{P}_2 \xrightarrow{f_2}
\mathcal{P}_1 \xrightarrow{f_1} \mathcal{P}_0 \to \mathcal{M} \to
0 $$
Définissons alors $\mathcal{P}$ comme le
module différentiel gradué sur $\mathcal{A}$
suivant \begin{enumerate}
\item comme $\mathcal{A}$-module gradué,
posons $\mathcal{P} = \bigoplus_{a \leq 0} \mathcal{P}_{-a}[-a]$ ; sa
composante de degré $n$ est donc
$\mathcal{P}^n = \bigoplus\nolimits_{a + b = n} \mathcal{P}_{-a}^b$,
\item le différentiel sur $\mathcal{P}$ est, comme dans la construction du complexe
total associé à un double complexe, donné par
$$
\text{d}_\mathcal{P}(x) = f_{-a}(x) + (-1)^a \text{d}_{\mathcal{P}_{-a}}(x)
$$ pour
$x$ une section locale de $\mathcal{P}_{-a}^b$.
\end{enumerate} Avec ces conventions
$\mathcal{P}$ est en effet un
$\mathcal{A}$-module différentiel gradué ; nous omettons les détails. Il existe un morphisme
$\mathcal{P} \to \mathcal{M}$ de
$\mathcal{A}$-modules différentiels gradués qui est nul sur les facteurs directs $\mathcal{P}_{-a}[-a]$ pour
$a < 0$ et égal au morphisme donné $\mathcal{P}_0 \to \mathcal{M}$ pour $a = 0$. Notons que
nous avons
$$
\mathcal{P} = \colim_i F_i\mathcal{P}
$$ où
$F_i\mathcal{P} \subset \mathcal{P}$ est le sous-module différentiel gradué sur
$\mathcal{A}$ dont le module gradué sous-jacent sur $\mathcal{A}$
est
$$
F_i\mathcal{P} =
\bigoplus\nolimits_{i \geq -a \geq 0} \mathcal{P}_{-a}[-a]
$$ Il est immédiat
que les applications
$$ 0
\to F_1\mathcal{P} \to F_2\mathcal{P} \to F_3\mathcal{P} \to
\ldots \to \mathcal{P}
$$ sont toutes des
monomorphismes admissibles, et nous avons des suites exactes courtes admissibles
$$ 0
\to F_i\mathcal{P} \to
F_{i + 1}\mathcal{P} \to
\mathcal{P}_{i + 1}[i + 1] \to 0
$$ Par
récurrence et par le lemme \ref{sdga-lemma-good-admissible-ses}, nous
obtenons que $F_i\mathcal{P}$ est un bon module différentiel
gradué sur $\mathcal{A}$. Puisque $\mathcal{P} = \colim F_i\mathcal{P}$ nous
trouvons que $\mathcal{P}$ est bon par le lemme \ref{sdga-lemma-good-direct-sum}.

\medskip\noindent
Enfin, nous devons démontrer que $\mathcal{P} \to \mathcal{M}$ est
un quasi-isomorphisme. Si $\mathcal{C}$ a suffisamment de points, cela
découle du lemme élémentaire d'Homologie
\ref{homology-lemma-good-resolution-gives-qis}, en vérifiant sur
les fibres. En général, on peut raisonner comme suit (cette
démonstration est beaucoup trop longue ; on peut aussi utiliser les sections
locales, comme dans la démonstration élémentaire, mais cet argument est
également assez long).
Puisque les limites inductives filtrantes sont exactes dans la catégorie des
faisceaux abéliens,
nous avons $$
H^d(\mathcal{P}) = \colim H^d(F_i\mathcal{P})
$$
Nous affirmons que, pour tout $i \geq 0$ et tout $d \in \mathbf{Z}$,
nous avons (a) une
suite exacte courte $$
0 \to H^d(\mathcal{M}_{i + 1}[i]) \to
H^d(F_i\mathcal{P}) \to H^d(\mathcal{M}) \to
0 $$
où la deuxième flèche
provient de $F_i\mathcal{P} \to \mathcal{P} \to \mathcal{M}$
et b) la composition
$$
H^d(\mathcal{M}_{i + 1}[i]) \to
H^d(F_i\mathcal{P}) \to
H^d(F_{i + 1}\mathcal{P})
$$ est
nulle. Cette assertion suffit manifestement à achever la démonstration.

\medskip\noindent
Démontrons l'assertion. Pour tout $i \geq 0$, il existe une
application $\mathcal{M}_{i + 1}[i] \to F_i\mathcal{P}$
provenant de l'inclusion de $\mathcal{M}_{i + 1}$ dans $\mathcal{P}_i$
en tant que noyau de $f_i$. Considérons la suite exacte
courte $$
0 \to \mathcal{M}_{i + 1}[i] \to F_i\mathcal{P} \to C_i \to
0 $$
de complexes de modules sur $\mathcal{O}$ définissant $C_i$.
Remarquons que $C_0 = \mathcal{M}_0 = \mathcal{M}$.
Observons de plus que $C_i$ est le complexe total associé
au double complexe $C_i^{\bullet, \bullet}$ avec des
colonnes $$
\mathcal{M}_i = \mathcal{P}_i/\mathcal{M}_{i + 1},
\mathcal{P}_{i - 1}, \ldots, \mathcal{P}_0
$$
en degrés $-i, -i + 1, \ldots, 0$. Il existe un morphisme de doubles
complexes $C_i^{\bullet, \bullet} \to C_{i - 1}^{\bullet, \bullet}$ qui
est $0$ sur la colonne en degré $-i$, est la surjection
$\mathcal{P}_{i - 1} \to \mathcal{M}_{i - 1}$ en degré $-i + 1$, et
est l'identité sur les autres colonnes. Nous
obtenons donc des morphismes de
complexes $$
C_i \longrightarrow C_{i - 1}
$$
Ces morphismes sont des quasi-isomorphismes surjectifs : leur
noyau est le complexe total du double complexe ayant pour
colonnes $\mathcal{M}_i, \mathcal{M}_i$ en degrés $-i, -i + 1$,
avec l'identité entre ces deux colonnes. En utilisant
les identifications qui en
résultent, $H^d(C_i) = H^d(C_{i - 1} = \ldots = H^d(\mathcal{M})$,
nous obtenons déjà une suite exacte longue
$$
H^d(\mathcal{M}_{i + 1}[i]) \to
H^d(F_i\mathcal{P}) \to H^d(\mathcal{M}) \to
H^{d + 1}(\mathcal{M}_{i + 1}[i])
$$
de la suite exacte courte de complexes ci-dessus.
Cependant, nous avons aussi le diagramme commutatif
$$
\xymatrix{
\mathcal{M}_{i + 2}[i + 1] \ar[r]_a & T_{i + 1} \ar[r] &
F_{i + 1}\mathcal{P} \ar[r] & C_{i + 1} \ar[d] \\ &
\mathcal{M}_{i + 1}[i] \ar[r] \ar[u]^b &
F_i\mathcal{P} \ar[u] \ar[r] &
C_i
}
$$ où
$T_{i + 1}$ est le complexe total du double complexe avec les
colonnes $\mathcal{P}_{i + 1}, \mathcal{M}_{i + 1}$ placées en degrés
$-i - 1$ et $-i$. Autrement dit, $T_{i + 1}$ est un décalage du
cône du morphisme
$\mathcal{P}_{i + 1} \to \mathcal{M}_{i + 1}$ ; nous voyons que
$a$ est un quasi-isomorphisme et que $a^{-1} \circ b$ est un décalage du troisième morphisme du
triangle distingué de $D(\mathcal{O})$ associé à la suite exacte
courte
$$ 0
\to \mathcal{M}_{i + 2} \to \mathcal{P}_{i + 1} \to \mathcal{M}_{i + 1} \to 0
$$ L'application
$H^d(\mathcal{P}_{i + 1}) \to H^d(\mathcal{M}_{i + 1})$ est surjective parce que nous
avons choisi nos applications de telle sorte que
$\Ker(\text{d}_{\mathcal{P}_{i + 1}}) \to
\Ker(\text{d}_{\mathcal{M}_{i + 1}})$ est surjective. Ainsi,
$a^{-1} \circ b$ induit l'application nulle sur les faisceaux de cohomologie.
Cela prouve la partie b) de
l'assertion. Puisque $T_{i + 1}$ est le noyau du morphisme surjectif de
complexes $F_{i + 1}\mathcal{P} \to C_i$, nous obtenons un
morphisme de suites exactes
longues de cohomologie $$
\xymatrix{
H^d(T_{i + 1}) \ar[r]
& H^d(F_{i + 1}\mathcal{P}) \ar[r]
& H^d(\mathcal{M}) \ar[r]
& H^{d + 1}(T_{i + 1}) \\
H^d(\mathcal{M}_{i + 1}[i]) \ar[r] \ar[u]
& H^d(F_i\mathcal{P}) \ar[r] \ar[u]
& H^d(\mathcal{M}) \ar[r] \ar[u]
& H^{d + 1}(\mathcal{M}_{i + 1}[i]) \ar[u]
}
$$
D'après la discussion ci-dessus, les flèches
verticales extérieures sont nulles. Ainsi, les
applications $H^d(F_{i + 1}\mathcal{P}) \to H^d(\mathcal{M})$
sont surjectives, ce qui prouve la partie
(a) de l'assertion. Plus précisément, cela prouve l'assertion pour $i > 0$ ;
nous laissons le cas $i = 0$ au lecteur (il suffit en
fait de traiter tous les $i \gg 0$ pour obtenir le lemme).
\end{proof}

\begin{lemma}
\label{sdga-lemma-acyclic-good}
Soit $(\mathcal{C}, \mathcal{O})$ un site annelé. Soit
$\mathcal{A}$ un faisceau d'algèbres différentielles graduées sur
$(\mathcal{C}, \mathcal{O})$. Soit $\mathcal{P}$ un bon module
différentiel gradué acyclique à droite sur $\mathcal{A}$.
\begin{enumerate}
\item pour tout module différentiel gradué à gauche sur $\mathcal{A}$,
$\mathcal{N}$, le produit tensoriel
$\mathcal{P} \otimes_\mathcal{A} \mathcal{N}$ est acyclique,
\item pour tout morphisme $(f, f^\sharp) : (\Sh(\mathcal{C}'), \mathcal{O}')
\to (\Sh(\mathcal{C}), \mathcal{O})$ de
topos annelés, toute $\mathcal{O}'$-algèbre différentielle graduée
$\mathcal{A}'$ et tout morphisme $\varphi : f^{-1}\mathcal{A} \to \mathcal{A}'$ de
$f^{-1}\mathcal{O}$-algèbres différentielles graduées,
l'image inverse $f^*\mathcal{P}$ est acyclique et bonne.
\end{enumerate}
\end{lemma}

\begin{proof}
Démontrons (1). Par le lemme \ref{sdga-lemma-resolve} nous pouvons choisir
un bon module différentiel gradué à gauche $\mathcal{Q}$ et un
quasi-isomorphisme $\mathcal{Q} \to \mathcal{N}$.
Alors $\mathcal{P} \otimes_\mathcal{A} \mathcal{Q}$
est acyclique parce que $\mathcal{Q}$
est bon. Soit $\mathcal{N}'$ le cône du
morphisme $\mathcal{Q} \to \mathcal{N}$.
Ensuite, $\mathcal{P} \otimes_\mathcal{A} \mathcal{N}'$ est acyclique
parce que $\mathcal{P}$ est bon et parce que $\mathcal{N}'$ est
acyclique (comme le cône sur un quasi-isomorphisme). Nous avons un
triangle
distingué $$
\mathcal{Q} \to \mathcal{N} \to \mathcal{N}' \to \mathcal{Q}[1]
$$
dans $K(\textit{Mod}(\mathcal{A}, \text{d}))$ par notre
construction de la structure triangulée. Puisque
$\mathcal{P} \otimes_\mathcal{A} -$ envoie les triangles
distingués sur des triangles distingués, nous obtenons un triangle
distingué
$$
\mathcal{P} \otimes_\mathcal{A} \mathcal{Q} \to
\mathcal{P} \otimes_\mathcal{A} \mathcal{N} \to
\mathcal{P} \otimes_\mathcal{A} \mathcal{N}' \to
\mathcal{P} \otimes_\mathcal{A} \mathcal{Q}[1]
$$ dans
$K(\textit{Mod}(\mathcal{O}))$. C'est ainsi que nous concluons.

\medskip\noindent
Démontrons (2). Notons que $f^*\mathcal{P}$ est bon selon notre définition de
bons modules. Rappelons
que $f^*\mathcal{P} = f^{-1}\mathcal{P} \otimes_{f^{-1}\mathcal{A}} \mathcal{A}'$.
Alors $f^{-1}\mathcal{P}$ est un bon module différentiel gradué acyclique (car $f^{-1}$
est exact) sur $f^{-1}\mathcal{A}$. On voit donc
que $f^*\mathcal{P}$ est acyclique par la partie (1).
\end{proof}






\section{L'enveloppe différentielle graduée d'un module gradué}
\label{sdga-section-dg-hull}

\noindent
L'enveloppe différentielle graduée d'un module gradué $\mathcal{N}$
est le résultat de l'application du foncteur $G$ du lemme suivant.

\begin{lemma}
\label{sdga-lemma-dg-hull}
Soit $(\mathcal{C}, \mathcal{O})$ un site annelé.
Soit $\mathcal{A}$ un faisceau d'algèbres différentielles graduées
sur $(\mathcal{C}, \mathcal{O})$. Le foncteur d'oubli
$F : \textit{Mod}(\mathcal{A}, \text{d}) \to \textit{Mod}(\mathcal{A})$ a
un adjoint à gauche $G : \textit{Mod}(\mathcal{A}) \to
\textit{Mod}(\mathcal{A}, \text{d})$.
\end{lemma}

\begin{proof}
Pour prouver l'existence de $G$, nous pouvons utiliser le théorème du foncteur
adjoint, voir Catégories, théorème \ref{categories-theorem-adjoint-functor} (noter que
nous avons interverti les rôles de $F$ et $G$). Les conditions d'exactitude sur
$F$ sont satisfaites par le lemme \ref{sdga-lemma-dgm-abelian}. La condition
ensembliste se vérifie comme suit : supposons donné un $\mathcal{A}$-module gradué
$\mathcal{N}$. Pour toute application
$$
\varphi : \mathcal{N} \longrightarrow F(\mathcal{M})
$$ nous
pouvons considérer le plus petit sous-module différentiel gradué sur $\mathcal{A}$,
$\mathcal{M}' \subset \mathcal{M}$, tel
que $\Im(\varphi) \subset F(\mathcal{M}')$.
Il est clair que $\mathcal{M}'$ est l'image de l'application
de $\mathcal{A}$-modules gradués
$$
\mathcal{N} \oplus
\mathcal{N}[-1] \otimes_\mathcal{O} \mathcal{A}
\longrightarrow
\mathcal{M}
$$ définie
par
$$ (n,
\sum n_i \otimes a_i) \longmapsto
\varphi(n) + \sum \text{d}(\varphi(n_i)) a_i
$$ car on vérifie
aisément que l'image de cette application est un sous-module différentiel
gradué de $\mathcal{M}$. Le nombre
de classes d'isomorphisme possibles de ces $\mathcal{M}'$ est donc borné, ce
qui conclut la démonstration.
\end{proof}

\noindent
Soit $(\mathcal{C}, \mathcal{O})$ un site annelé. Soit
$\mathcal{A}$ un faisceau d'algèbres différentielles graduées sur
$(\mathcal{C}, \mathcal{O})$. Soit $\mathcal{M}$ un
$\mathcal{A}$-module différentiel gradué et supposons donnée une suite exacte courte
$$ 0
\to \mathcal{N} \to F(\mathcal{M}) \to \mathcal{N}' \to 0
$$ dans
$\textit{Mod}(\mathcal{A})$. Nous obtenons alors un homomorphisme canonique
de modules gradués sur $\mathcal{A}$
$$
\overline{\text{d}} : \mathcal{N} \to \mathcal{N}'[1]
$$ de la manière
suivante : pour une section locale $x$ de $\mathcal{N}$, notons
$\overline{\text{d}}(x)$ l'image dans $\mathcal{N}'$ de
$\text{d}_\mathcal{M}(x)$, où $x$ est considéré comme une section locale
de $\mathcal{M}$.

\begin{lemma}
\label{sdga-lemma-dg-hull-acyclic}
Les foncteurs $F, G$ du lemme \ref{sdga-lemma-dg-hull} ont les
propriétés suivantes. Pour un $\mathcal{A}$-module gradué
$\mathcal{N}$, nous avons
\begin{enumerate}
\item l'unité $\mathcal{N} \to F(G(\mathcal{N}))$ est injective,
\item l'application $\overline{\text{d}} : \mathcal{N} \to
\Coker(\mathcal{N} \to F(G(\mathcal{N})))[1]$ est un isomorphisme, et
\item $G(\mathcal{N})$ est un $\mathcal{A}$-module différentiel gradué acyclique.
\end{enumerate}
\end{lemma}

\begin{proof}
On observe que la propriété (3) est une conséquence des propriétés (1) et (2).
En effet, si $s$ est une section locale non nulle de $F(G(\mathcal{N}))$
telle que $\text{d}(s) = 0$, alors $s$ ne peut pas
appartenir à l'image de $\mathcal{N} \to F(G(\mathcal{N}))$. Ainsi, nous
pouvons écrire l'image $\overline{s}$ de $s$
dans le conoyau comme $\overline{\text{d}}(s')$ pour une section locale $s'$ de $\mathcal{N}$.
Alors $s = \text{d}(s')$, car la différence
$s - \text{d}(s')$ appartient encore au noyau de $\text{d}$
et est contenue dans l'image de l'unité.

\medskip\noindent
Notons provisoirement $\mathcal{A}_{gr}$, respectivement $\mathcal{A}_{dg}$, le
faisceau $\mathcal{A}$ considéré comme un module gradué à droite sur
lui-même, respectivement comme un module différentiel gradué à droite
sur lui-même. Le cas le plus important du lemme est de comprendre
ce qu'est $G(\mathcal{A}_{gr})$. Bien sûr, $G(\mathcal{A}_{gr})$ est
l'objet de $\textit{Mod}(\mathcal{A}, \text{d})$ représentant
le
foncteur $$
\mathcal{M} \longmapsto
\Hom_{\textit{Mod}(\mathcal{A})}(\mathcal{A}_{gr}, F(\mathcal{M}))
= \Gamma(\mathcal{C}, \mathcal{M})
$$
Par la remarque \ref{sdga-remark-cone-identity}, nous voyons que ce
foncteur est représenté par $C[-1]$ où $C$ est le cône sur l'identité de $\mathcal{A}_{dg}$.
Nous avons une
suite exacte courte $$
0 \to \mathcal{A}_{dg}[-1] \to C[-1] \to \mathcal{A}_{dg} \to
0 $$
dans $\textit{Mod}(\mathcal{A}, \text{d})$,
scindée par l'unité $\mathcal{A}_{gr} \to F(C[-1])$ dans $\textit{Mod}(\mathcal{A})$.
Ainsi, $G(\mathcal{A}_{gr})$ vérifie (1) et (2).

\medskip\noindent
Soit $U$ un objet de $\mathcal{C}$. Notons
$j_U : \mathcal{C}/U \to \mathcal{C}$ le morphisme de localisation.
Notons $\mathcal{A}_U$ la restriction de $\mathcal{A}$ à $U$.
Nous noterons $\mathcal{A}_{U, gr}$ le faisceau
$\mathcal{A}_U$ considéré comme un module gradué sur $\mathcal{A}_U$.
Notons $F_U : \textit{Mod}(\mathcal{A}_U, \text{d}) \to
\textit{Mod}(\mathcal{A}_U)$ le foncteur d'oubli et
$G_U$ son adjoint. Nous avons alors les diagrammes commutatifs
$$
\vcenter{
\xymatrix{
\textit{Mod}(\mathcal{A}, \text{d}) \ar[d]_{j_U^*} \ar[r]_F &
\textit{Mod}(\mathcal{A}) \ar[d]^{j_U^*} \\
\textit{Mod}(\mathcal{A}_U, \text{d}) \ar[r]^{F_U} &
\textit{Mod}(\mathcal{A}_U)
}
}
\quad\text{et}\quad
\vcenter{
\xymatrix{
\textit{Mod}(\mathcal{A}_U, \text{d}) \ar[r]_{F_U} \ar[d]_{j_{U!}}
& \textit{Mod}(\mathcal{A}_U) \ar[d]^{j_{U!}} \\
\textit{Mod}(\mathcal{A}, \text{d}) \ar[r]^F
& \textit{Mod}(\mathcal{A})
}
}
$$
par la construction de $j^*_U$ et $j_{U!}$
dans les sections \ref{sdga-section-functoriality-graded},
\ref{sdga-section-functoriality-dg},
\ref{sdga-section-localize-graded}
et \ref{sdga-section-localize-dg}.
Par l'unicité des adjoints, on obtient $j_{U!} \circ G_U = G \circ j_{U!}$.
Puisque $j_{U!}$ est exact, les propriétés (1)
et (2) pour l'unité
$\mathcal{A}_{U, gr} \to F_U(G_U(\mathcal{A}_{U, gr}))$, établies
dans la partie précédente de la démonstration,
impliquent les mêmes propriétés pour l'unité
$j_{U!}\mathcal{A}_{U, gr} \to F(G(j_{U!}\mathcal{A}_{U, gr})) =
j_{U!}F_U(G_U(\mathcal{A}_{U, gr}))$.

\medskip\noindent Dans
la démonstration du lemme \ref{sdga-lemma-gm-grothendieck-abelian},
nous avons vu que tout objet de
$\textit{Mod}(\mathcal{A})$ est un quotient d'une somme directe de copies
de $j_{U!}\mathcal{A}_{U, gr}$. Puisque $G$ est un adjoint à gauche,
nous voyons que $G$ commute aux sommes directes. Ainsi, les
propriétés (1) et (2) sont vérifiées pour les sommes directes
d'objets qui les vérifient. Par conséquent, tout objet $\mathcal{N}$ de
$\textit{Mod}(\mathcal{A})$ s'insère dans une suite exacte
$$
\mathcal{N}_1 \to \mathcal{N}_0 \to \mathcal{N} \to 0
$$ telle que les propriétés
(1) et (2) soient vérifiées pour $\mathcal{N}_1$ et $\mathcal{N}_0$. Nous
laissons au lecteur le soin de déduire (1) et (2) pour
$\mathcal{N}$ en utilisant que $G$ est exact à droite.
\end{proof}









\section{Modules différentiels gradués K-injectifs}
\label{sdga-section-K-injective}

\noindent
Cette section reprend les résultats
d'Injectifs, section \ref{injectives-section-K-injective},
dans le cadre des faisceaux de modules différentiels gradués sur
un faisceau d'algèbres différentielles graduées.

\begin{lemma}
\label{sdga-lemma-characterize-injectives} Soit
$(\mathcal{C}, \mathcal{O})$ un site annelé. Soit $\mathcal{A}$ un faisceau
d'algèbres graduées sur $(\mathcal{C}, \mathcal{O})$. Il existe un
ensemble $T$ et, pour chaque $t \in T$, une application injective
$\mathcal{N}_t \to \mathcal{N}'_t$ de $\mathcal{A}$-modules gradués, tels qu'un
objet $\mathcal{I}$ de $\textit{Mod}(\mathcal{A})$ soit injectif si et
seulement si, pour tout diagramme en traits pleins
$$
\xymatrix{
\mathcal{N}_t \ar[r] \ar[d] & \mathcal{I} \\
\mathcal{N}'_t \ar@{..>}[ru]
}
$$ il existe
une flèche pointillée dans $\textit{Mod}(\mathcal{A})$ rendant le diagramme commutatif.
\end{lemma}

\begin{proof}
Cet énoncé vaut dans toute catégorie abélienne de Grothendieck ;
voir Injectifs, lemme \ref{injectives-lemma-characterize-injective}.
Par le lemme \ref{sdga-lemma-gm-grothendieck-abelian}, la catégorie
$\textit{Mod}(\mathcal{A})$ est une catégorie abélienne de Grothendieck.
\end{proof}

\begin{definition}
\label{sdga-definition-graded-injective}
Soit $(\mathcal{C}, \mathcal{O})$ un site annelé. Soit
$(\mathcal{A}, \text{d})$ un faisceau d'algèbres différentielles graduées sur
$(\mathcal{C}, \mathcal{O})$. Un $\mathcal{A}$-module différentiel gradué
$\mathcal{I}$ est dit {\it injectif gradué}\footnote{Cette terminologie n'est
peut-être pas standard.} si $\mathcal{M}$, considéré comme un
$\mathcal{A}$-module gradué, est un objet injectif de la catégorie
$\textit{Mod}(\mathcal{A})$ des $\mathcal{A}$-modules gradués.
\end{definition}

\begin{remark}
\label{sdga-remark-why-graded-injective}
Soit $(\mathcal{C}, \mathcal{O})$ un site annelé. Soit
$(\mathcal{A}, \text{d})$ un faisceau d'algèbres différentielles graduées sur
$(\mathcal{C}, \mathcal{O})$. Soit $\mathcal{I}$ un
$\mathcal{A}$-module différentiel gradué injectif gradué. Soit
$$ 0
\to \mathcal{M}_1 \to \mathcal{M}_2 \to \mathcal{M}_3 \to 0
$$ une
suite exacte courte de $\mathcal{A}$-modules différentiels gradués. Puisque
$\mathcal{I}$ est injectif gradué, nous
obtenons une suite exacte courte de complexes
$$ 0
\to
\Hom_{\textit{Mod}^{dg}(\mathcal{A}, \text{d})}(\mathcal{M}_3, \mathcal{I})
\to
\Hom_{\textit{Mod}^{dg}(\mathcal{A}, \text{d})}(\mathcal{M}_2, \mathcal{I})
\to
\Hom_{\textit{Mod}^{dg}(\mathcal{A}, \text{d})}(\mathcal{M}_1, \mathcal{I})
\to 0
$$ de modules
sur $\Gamma(\mathcal{C}, \mathcal{O})$. En passant à la cohomologie,
nous obtenons une suite exacte
longue $$
\xymatrix{
\Hom_{K(\textit{Mod}(\mathcal{A}, \text{d}))}(\mathcal{M}_3, \mathcal{I})
\ar[d]
& \Hom_{K(\textit{Mod}(\mathcal{A}, \text{d}))}(\mathcal{M}_3, \mathcal{I})[1]
\ar[d] \\
\Hom_{K(\textit{Mod}(\mathcal{A}, \text{d}))}(\mathcal{M}_2, \mathcal{I})
\ar[d] &
\Hom_{K(\textit{Mod}(\mathcal{A}, \text{d}))}(\mathcal{M}_2, \mathcal{I})[1]
\ar[d] \\
\Hom_{K(\textit{Mod}(\mathcal{A}, \text{d}))}(\mathcal{M}_1, \mathcal{I})
\ar[ruu]
&
\Hom_{K(\textit{Mod}(\mathcal{A}, \text{d}))}(\mathcal{M}_1, \mathcal{I})[1]
}
$$ de
groupes d'homomorphismes dans la catégorie homotopique. Le point est que nous
obtenons cela même si nous n'avons pas supposé que notre suite exacte courte
est admissible (donc la suite exacte courte en général ne définit pas
un triangle distingué dans la catégorie homotopique).
\end{remark}

\begin{lemma}
\label{sdga-lemma-product-graded-injective}
Soit $(\mathcal{C}, \mathcal{O})$ un site annelé. Soit
$(\mathcal{A}, \text{d})$ un faisceau d'algèbres différentielles graduées sur
$(\mathcal{C}, \mathcal{O})$. Soit $T$ un ensemble et, pour
chaque $t \in T$, soit $\mathcal{I}_t$ un module différentiel
gradué injectif gradué sur $\mathcal{A}$. Alors
$\prod \mathcal{I}_t$ est un
$\mathcal{A}$-module différentiel gradué injectif gradué.
\end{lemma}

\begin{proof}
Cela résulte du fait que les produits d'objets injectifs sont injectifs,
voir Homologie, lemme \ref{homology-lemma-product-injectives},
et que les produits
dans $\textit{Mod}(\mathcal{A}, \text{d})$ correspondent
aux produits dans $\textit{Mod}(\mathcal{A})$ par le foncteur d'oubli.
\end{proof}

\begin{lemma}
\label{sdga-lemma-characterize-graded-injectives-in-dg} Soit
$(\mathcal{C}, \mathcal{O})$ un site annelé. Soit
$(\mathcal{A}, \text{d})$ un faisceau
d'algèbres différentielles graduées sur $(\mathcal{C}, \mathcal{O})$. Il existe un ensemble
$T$ et, pour chaque $t \in T$, une application injective
$\mathcal{M}_t \to \mathcal{M}'_t$ de
$\mathcal{A}$-modules différentiels gradués acycliques, tels que, pour
un objet $\mathcal{I}$ de $\textit{Mod}(\mathcal{A}, \text{d})$, les propriétés suivantes
soient équivalentes :
\begin{enumerate}
\item $\mathcal{I}$ est injectif gradué, et
\item pour chaque diagramme en traits pleins
$$
\xymatrix{
\mathcal{M}_t \ar[r] \ar[d] & \mathcal{I} \\
\mathcal{M}'_t \ar@{..>}[ru]
}
$$ il
existe une flèche pointillée dans $\textit{Mod}(\mathcal{A}, \text{d})$
rendant le diagramme commutatif.
\end{enumerate}
\end{lemma}

\begin{proof}
Soient $T$ et $\mathcal{N}_t \to \mathcal{N}'_t$ comme dans
le lemme \ref{sdga-lemma-characterize-injectives}.
Notons $F : \textit{Mod}(\mathcal{A}, \text{d}) \to \textit{Mod}(\mathcal{A})$ le
foncteur d'oubli. Soit
$G$ l'adjoint à gauche de
$F$ donné par le lemme \ref{sdga-lemma-dg-hull}. Posons
$$
\mathcal{M}_t = G(\mathcal{N}_t) \to
G(\mathcal{N}'_t) = \mathcal{M}'_t
$$ Il s'agit d'un
homomorphisme injectif de modules différentiels gradués acycliques
sur $\mathcal{A}$, par le lemme \ref{sdga-lemma-dg-hull-acyclic}.
Puisque $G$ est l'adjoint à gauche de $F$ nous voyons
qu'il existe une flèche pointillée dans le
diagramme $$
\xymatrix{
\mathcal{M}_t \ar[r] \ar[d] & \mathcal{I} \\
\mathcal{M}'_t \ar@{..>}[ru]
}
$$
si et seulement s'il existe une flèche pointillée dans
le diagramme $$
\xymatrix{
\mathcal{N}_t \ar[r] \ar[d] & F(\mathcal{I}) \\
\mathcal{N}'_t \ar@{..>}[ru]
}
$$
L'assertion résulte donc du choix de
notre collection de flèches $\mathcal{N}_t \to \mathcal{N}_t'$.
\end{proof}


\begin{lemma}
\label{sdga-lemma-small-acyclics}
Soit $(\mathcal{C}, \mathcal{O})$ un site annelé. Soit
$(\mathcal{A}, \text{d})$ un faisceau d'algèbres différentielles graduées sur
$(\mathcal{C}, \mathcal{O})$. Il existe un ensemble $S$ et, pour chaque $s$,
un $\mathcal{A}$-module différentiel gradué acyclique $\mathcal{M}_s$, tels
que, pour tout $\mathcal{A}$-module différentiel gradué acyclique non nul
$\mathcal{M}$, il existe $s \in S$ et une application injective
$\mathcal{M}_s \to \mathcal{M}$ dans $\textit{Mod}(\mathcal{A}, \text{d})$.
\end{lemma}

\begin{proof}
Avant de commencer, rappelons que nos conventions garantissent que le site $\mathcal{C}$
a un ensemble d'objets et de morphismes et un ensemble $\text{Cov}(\mathcal{C})$
de
couvertures. Si $\mathcal{F}$
est un $\mathcal{A}$-module différentiel gradué, définissons $|\mathcal{F}|$
comme le cardinal de
$$
\coprod\nolimits_{(U, n)} \mathcal{F}^n(U)
$$
lorsque $U$ parcourt les objets de $\mathcal{C}$ et $n \in \mathbf{Z}$.
Choisissons un cardinal infini $\kappa$ plus grand
que les cardinaux $|\Ob(\mathcal{C})|$, $|\text{Arrows}(\mathcal{C})|$,
$|\text{Cov}(\mathcal{C})|$, $\sup |I|$
pour $\{U_i \to U\}_{i \in I} \in \text{Cov}(\mathcal{C})$
et $|\mathcal{A}|$.

\medskip\noindent
Soit $\mathcal{F} \subset \mathcal{M}$ une inclusion de modules
différentiels gradués sur $\mathcal{A}$.
Supposons donnés un ensemble $K$ et, pour chaque $k \in K$, un
triplet $(U_k, n_k, x_k)$ constitué d'un objet $U_k$ de $\mathcal{C}$,
d'un entier $n_k$ et d'une section $x_k \in \mathcal{M}^{n_k}(U_k)$.
Nous pouvons alors considérer le plus petit sous-module
différentiel gradué sur $\mathcal{A}$, $\mathcal{F}' \subset \mathcal{M}$,
contenant $\mathcal{F}$ et les sections $x_k$ pour $k \in K$.
On peut
décrire $$
(\mathcal{F}')^n(U) \subset \mathcal{M}^n(U)
$$
comme l'ensemble des éléments $x \in \mathcal{M}^n(U)$ pour lesquels
il existe un recouvrement $\{f_i : U_i \to U\}_{i \in I} \in \text{Cov}(\mathcal{C})$
tel que, pour tout $i \in I$, il existe un ensemble fini $T_i$
et des morphismes $g_t : U_i \to U_{k_t}$
vérifiant $$
f_i^*x
= y_i + \sum\nolimits_{t \in T_i} a_{it}g_t^*x_{k_t} + b_{it}g_t^*\text{d}(x_{k_t})
$$
pour une section $y_i \in \mathcal{F}^n(U)$
et des sections $a_{it} \in \mathcal{A}^{n - n_{k_t}}(U_i)$
et $b_{it} \in \mathcal{A}^{n - n_{k_t} - 1}(U_i)$.
(Détails omis ; indication : ces sections appartiennent à $\mathcal{F}'$
et, réciproquement, on vérifie que cette règle définit
un sous-module différentiel gradué sur $\mathcal{A}$.)
Cette description
donne $|\mathcal{F}'| \leq \max(|\mathcal{F}|, |K|, \kappa)$.

\medskip\noindent
Soit $\mathcal{M}$ un
$\mathcal{A}$-module différentiel gradué acyclique non nul. Alors il existe un entier $n$ et une section
non nulle $x$ de $\mathcal{M}^n$ sur un objet $U$ de
$\mathcal{C}$. Soit
$$
\mathcal{F}_0 \subset \mathcal{M}
$$ le plus
petit sous-module différentiel gradué sur $\mathcal{A}$ contenant
$x$. Par le paragraphe précédent, nous avons
$|\mathcal{F}_0| \leq \kappa$. Par récurrence, étant donnés
$\mathcal{F}_0, \ldots, \mathcal{F}_n$, définissons
$\mathcal{F}_{n + 1}$ comme suit. Considérons
l'ensemble $$
L = \{(U, n, x)\}
\{U_i \to U\}_{i \in I}, (x_i)_{i \in I})\}
$$
des triples où $U$ est un objet de $\mathcal{C}$, $n \in \mathbf{Z}$
et $x \in \mathcal{F}_n(U)$ avec $\text{d}(x) = 0$. Puisque
$\mathcal{M}$ est acyclique pour chaque triple $l = (U_l, n_l, x_l) \in L$
nous pouvons
choisir $\{(U_{l, i} \to U_l\}_{i \in I_l} \in \text{Cov}(\mathcal{C})$
et $x_{l, i} \in \mathcal{M}^{n_l - 1}(U_{l, i})$ tels
que $\text{d}(x_{l, i}) = x|_{U_{l, i}}$. Posons alors
$$
K = \{(U_{l, i}, n_l - 1, x_{l, i}) \mid l \in L, i \in I_l\}
$$
et prenons pour $\mathcal{F}_{n + 1}$ le plus petit sous-module
différentiel gradué sur $\mathcal{A}$ de $\mathcal{M}$ contenant
$\mathcal{F}_n$ et les sections $x_{l, i}$.
Puisque $|K| \leq \max(\kappa, |\mathcal{F}_n|)$
nous concluons que $|\mathcal{F}_{n + 1}| \leq \kappa$ par induction.

\medskip\noindent
Par construction, l'inclusion $\mathcal{F}_n \to \mathcal{F}_{n + 1}$ induit
l'application nulle sur les faisceaux de cohomologie. Ainsi,
$\mathcal{F} = \bigcup \mathcal{F}_n$ est un sous-module acyclique non nul tel
que $|\mathcal{F}| \leq \kappa$. Puisque les classes d'isomorphisme
de $\mathcal{A}$-modules différentiels gradués
$\mathcal{F}$ dont $|\mathcal{F}|$ est borné forment un ensemble, nous concluons.
\end{proof}

\begin{definition}
\label{sdga-definition-K-injective} Soit
$(\mathcal{C}, \mathcal{O})$ un site annelé. Soit
$(\mathcal{A}, \text{d})$ un faisceau d'algèbres différentielles graduées sur
$(\mathcal{C}, \mathcal{O})$. Un $\mathcal{A}$-module différentiel gradué
$\mathcal{I}$ est dit {\it K-injectif} si, pour tout module différentiel
gradué acyclique $\mathcal{M}$, nous
avons $$
\Hom_{K(\textit{Mod}(\mathcal{A}, \text{d}))}(\mathcal{M}, \mathcal{I}) = 0
$$
\end{definition}

\noindent
Notons la similitude avec
Catégories dérivées, définition \ref{derived-definition-K-injective}.

\begin{lemma}
\label{sdga-lemma-product-K-injective}
Soit $(\mathcal{C}, \mathcal{O})$ un site annelé. Soit
$(\mathcal{A}, \text{d})$ un faisceau d'algèbres différentielles graduées sur
$(\mathcal{C}, \mathcal{O})$. Soit $T$ un ensemble et, pour
chaque $t \in T$, soit $\mathcal{I}_t$ un module différentiel
gradué K-injectif sur $\mathcal{A}$. Alors
$\prod \mathcal{I}_t$ est un
$\mathcal{A}$-module différentiel gradué K-injectif.
\end{lemma}

\begin{proof}
Soit $\mathcal{K}$ un module différentiel gradué acyclique sur $\mathcal{A}$. Alors
nous avons
$$
\Hom_{\textit{Mod}^{dg}(\mathcal{A}, \text{d})}(\mathcal{K},
\prod\nolimits_{t \in T} \mathcal{I}_t)
=
\prod\nolimits_{t \in T}
\Hom_{\textit{Mod}^{dg}(\mathcal{A}, \text{d})}(\mathcal{K}, \mathcal{I}_t)
$$ car la formation
des produits dans $\textit{Mod}(\mathcal{A}, \text{d})$ commute au
foncteur d'oubli vers les modules gradués sur $\mathcal{A}$. Puisque la
prise de produits est un foncteur exact sur la catégorie des groupes abéliens, nous
en concluons.
\end{proof}

\begin{lemma}
\label{sdga-lemma-first-property-dg-injective} Soit
$(\mathcal{C}, \mathcal{O})$ un site annelé. Soit
$(\mathcal{A}, \text{d})$ un
faisceau d'algèbres différentielles graduées sur $(\mathcal{C}, \mathcal{O})$. Soit
$\mathcal{I}$ un objet K-injectif et injectif gradué de
$\textit{Mod}(\mathcal{A}, \text{d})$. Pour chaque diagramme
en traits pleins dans $\textit{Mod}(\mathcal{A}, \text{d})$
$$
\xymatrix{
\mathcal{M} \ar[r]_a \ar[d]_b & \mathcal{I} \\
\mathcal{M}' \ar@{..>}[ru]
}
$$ où
$b$ est injectif et $\mathcal{M}$ acyclique, il existe une
flèche pointillée rendant le diagramme commutatif.
\end{lemma}

\begin{proof}
Puisque $\mathcal{M}$ est acyclique et $\mathcal{I}$ est K-injectif, il
existe un homomorphisme de $\mathcal{A}$-modules gradués
$h : \mathcal{M} \to \mathcal{I}$ de degré $-1$ tel
que $a = \text{d}(h)$. Puisque $\mathcal{I}$ est injectif gradué et
$b$ est injectif, il existe un homomorphisme de $\mathcal{A}$-modules gradués
$h' : \mathcal{M}' \to \mathcal{I}$ de degré $-1$ tel que
$h = h' \circ b$. Ensuite, nous pouvons prendre $a' = \text{d}(h')$ comme
la flèche pointillée.
\end{proof}

\begin{lemma}
\label{sdga-lemma-second-property-dg-injective} Soit
$(\mathcal{C}, \mathcal{O})$ un site annelé. Soit
$(\mathcal{A}, \text{d})$ un faisceau d'algèbres différentielles graduées sur
$(\mathcal{C}, \mathcal{O})$. Soit $\mathcal{I}$ un objet K-injectif
et injectif gradué de
$\textit{Mod}(\mathcal{A}, \text{d})$. Pour chaque diagramme
en traits pleins dans $\textit{Mod}(\mathcal{A}, \text{d})$
$$
\xymatrix{
\mathcal{M} \ar[r]_a \ar[d]_b & \mathcal{I} \\
\mathcal{M}' \ar@{..>}[ru]
}
$$ où
$b$ est un quasi-isomorphisme, il existe une flèche pointillée rendant le diagramme
commutatif à homotopie près.
\end{lemma}

\begin{proof}
Après avoir remplacé $\mathcal{M}'$ par la somme directe de $\mathcal{M}'$
et du cône de l'identité de $\mathcal{M}$ (qui est acyclique), nous
pouvons supposer que $b$ est également injectif. Alors le conoyau $\mathcal{Q}$
de $b$ est acyclique. Ainsi,
nous voyons que $$
\Hom_{K(\textit{Mod}(\mathcal{A}, \text{d}))}(\mathcal{Q}, \mathcal{I})
= \Hom_{K(\textit{Mod}(\mathcal{A}, \text{d}))}(\mathcal{Q}, \mathcal{I})[1]
= 0 $$
comme $\mathcal{I}$ est K-injectif. Puisque $\mathcal{I}$ est injectif
gradué, la remarque \ref{sdga-remark-why-graded-injective}
montre
que $$
\Hom_{K(\textit{Mod}(\mathcal{A}, \text{d}))}(\mathcal{M}', \mathcal{I})
\longrightarrow
\Hom_{K(\textit{Mod}(\mathcal{A}, \text{d}))}(\mathcal{M}, \mathcal{I})
$$
est bijectif et la preuve est complète.
\end{proof}

\begin{lemma}
\label{sdga-lemma-better-set-of-monos} Soit
$(\mathcal{C}, \mathcal{O})$ un site annelé. Soit
$(\mathcal{A}, \text{d})$ un faisceau
d'algèbres différentielles graduées sur $(\mathcal{C}, \mathcal{O})$. Il existe un
ensemble $R$ et, pour chaque $r \in R$, une application injective
$\mathcal{M}_r \to \mathcal{M}'_r$ de
$\mathcal{A}$-modules différentiels gradués acycliques, tels que,
pour un objet $\mathcal{I}$ de $\textit{Mod}(\mathcal{A}, \text{d})$, les propriétés
suivantes soient équivalentes :
\begin{enumerate}
\item $\mathcal{I}$ est K-injectif et injectif gradué, et
\item pour chaque diagramme en traits pleins
$$
\xymatrix{
\mathcal{M}_r \ar[r] \ar[d] & \mathcal{I} \\
\mathcal{M}'_r \ar@{..>}[ru]
}
$$ il
existe une flèche pointillée dans $\textit{Mod}(\mathcal{A}, \text{d})$
rendant le diagramme commutatif.
\end{enumerate}
\end{lemma}

\begin{proof}
Soient $T$ et $\mathcal{M}_t \to \mathcal{M}'_t$ comme dans le
lemme \ref{sdga-lemma-characterize-graded-injectives-in-dg}. Soient
$S$ et $\mathcal{M}_s$ comme dans le lemme
\ref{sdga-lemma-small-acyclics}. Choisissons
une application injective $\mathcal{M}_s \to \mathcal{M}'_s$
de modules différentiels gradués acycliques sur $\mathcal{A}$
qui soit homotope à zéro. Cela est possible parce que nous
pouvons considérer que $\mathcal{M}'_s$ est le cône de l'identité ;
dans ce cas, il est même vrai que l'identité sur
$\mathcal{M}'_s$ est homotope à zéro, voir Algèbre
différentielle graduée, lemme \ref{dga-lemma-id-cone-null}, qui s'applique
par la discussion de la section \ref{sdga-section-conclude-triangulated}. Nous affirmons
que $R = T \coprod S$ avec les applications données fonctionne.

\medskip\noindent
L'implication (1) $\Rightarrow$ (2) résulte du
lemme \ref{sdga-lemma-first-property-dg-injective}.

\medskip\noindent
Supposons (2). Tout d'abord, par le lemme \ref{sdga-lemma-characterize-graded-injectives-in-dg}
nous voyons que $\mathcal{I}$ est injectif gradué.
Soit maintenant $\mathcal{M}$ un module différentiel gradué acyclique sur $\mathcal{A}$.
Nous
devons montrer que $$
\Hom_{K(\textit{Mod}(\mathcal{A}, \text{d}))}(\mathcal{M}, \mathcal{I})
= 0 $$
La démonstration sera identique à celle
d'Injectifs, lemme \ref{injectives-lemma-characterize-K-injective}.

\medskip\noindent
Nous allons construire par récurrence transfinie sur l'ordinal $\alpha$
un sous-module différentiel gradué acyclique
$\mathcal{K}_\alpha \subset \mathcal{M}$ comme suit.
Pour $\alpha = 0$, posons $\mathcal{K}_0 = 0$. Pour $\alpha > 0$,
on procède comme suit :
\begin{enumerate}
\item Si $\alpha = \beta + 1$ et $\mathcal{K}_\beta = \mathcal{M}$,
alors nous choisissons $\mathcal{K}_\alpha = \mathcal{K}_\beta$.
\item Si $\alpha = \beta + 1$ et $\mathcal{K}_\beta \not = \mathcal{M}$,
alors $\mathcal{M}/\mathcal{K}_\beta$ est un module différentiel gradué
acyclique non nul sur $\mathcal{A}$.
Nous choisissons un sous-$\mathcal{A}$-module différentiel gradué
$\mathcal{N}_\alpha \subset \mathcal{M}/\mathcal{K}_\beta$
isomorphe à $\mathcal{M}_s$ pour un certain $s \in S$ ; voir le
lemme \ref{sdga-lemma-small-acyclics}.
Enfin, nous définissons $\mathcal{K}_\alpha \subset \mathcal{M}$
comme l'image inverse de $\mathcal{N}_\alpha$.
\item Si $\alpha$ est un ordinal limite, posons
$\mathcal{K}_\beta = \colim \mathcal{K}_\alpha$.
\end{enumerate}
Il est clair que $\mathcal{M} = \mathcal{K}_\alpha$ pour un
ordinal $\alpha$ suffisamment grand. Nous allons montrer que
$$
\Hom_{K(\textit{Mod}(\mathcal{A}, \text{d}))}(\mathcal{K}_\alpha, \mathcal{I})
$$
est nul par récurrence transfinie sur $\alpha$. C'est vrai pour $\alpha = 0$,
car $\mathcal{K}_0$ est nul. Supposons-le vrai pour $\beta$ et
$\alpha = \beta + 1$. Dans le cas (1) de la liste ci-dessus, le résultat est clair.
Dans le cas (2), il existe une suite exacte courte
$$
0 \to \mathcal{K}_\beta \to \mathcal{K}_\alpha \to \mathcal{N}_\alpha \to 0
$$
D'après la remarque \ref{sdga-remark-why-graded-injective}
et puisque nous avons vu que $\mathcal{I}$ est injectif
gradué, nous obtenons une suite exacte
$$
\Hom_{K(\textit{Mod}(\mathcal{A}, \text{d}))}(\mathcal{K}_\beta, \mathcal{I})
\to
\Hom_{K(\textit{Mod}(\mathcal{A}, \text{d}))}(\mathcal{K}_\alpha, \mathcal{I})
\to
\Hom_{K(\textit{Mod}(\mathcal{A}, \text{d}))}(\mathcal{N}_\alpha, \mathcal{I})
$$
Par récurrence, le terme de gauche est nul. Par l'hypothèse (2),
celui de droite est nul : toute application $\mathcal{M}_s \to \mathcal{I}$
se factorise par $\mathcal{M}'_s$ et est donc homotope à zéro.
Le groupe du milieu est donc nul lui aussi. Enfin, supposons que $\alpha$ soit un
ordinal limite. Puisque nous avons aussi
$\mathcal{K}_\alpha = \colim \mathcal{K}_\alpha$ comme
$\mathcal{A}$-modules gradués, nous voyons que
$$
\Hom_{\textit{Mod}^{dg}(\mathcal{A}, \text{d})}
(\mathcal{K}_\alpha, \mathcal{I})
= \lim_{\beta < \alpha}
\Hom_{\textit{Mod}^{dg}(\mathcal{A}, \text{d})}
(\mathcal{K}_\beta, \mathcal{I})
$$
comme complexes de groupes abéliens. Les groupes de cohomologie de ces
complexes calculent les morphismes dans $K(\textit{Mod}(\mathcal{A}, \text{d}))$
entre décalages. Les morphismes de transition du système de complexes
sont surjectifs d'après la remarque \ref{sdga-remark-why-graded-injective},
car $\mathcal{I}$ est injectif gradué.
De plus, pour tout ordinal limite $\beta \leq \alpha$,
la limite coïncide avec la valeur. Nous pouvons donc appliquer le
lemme \ref{homology-lemma-ML-over-ordinals} d'Homologie
pour conclure.
\end{proof}

\begin{lemma}
\label{sdga-lemma-functor-set-of-monos} Soit
$(\mathcal{C}, \mathcal{O})$ un site annelé. Soit
$(\mathcal{A}, \text{d})$ un faisceau
d'algèbres différentielles graduées sur $(\mathcal{C}, \mathcal{O})$. Soit
$R$ un ensemble et, pour chaque $r \in R$, soit donnée une application
injective $\mathcal{M}_r \to \mathcal{M}'_r$
de $\mathcal{A}$-modules différentiels gradués acycliques.
Il existe un foncteur $M : \textit{Mod}(\mathcal{A}, \text{d}) \to
\textit{Mod}(\mathcal{A}, \text{d})$ et une transformation naturelle
$j : \text{id} \to M$ telle que
\begin{enumerate}
\item $j_\mathcal{M} : \mathcal{M} \to M(\mathcal{M})$ est injectif et est un
quasi-isomorphisme,
\item pour tout diagramme en traits pleins
$$
\xymatrix{
\mathcal{M}_r \ar[r] \ar[d] & \mathcal{M} \ar[d]^{j_\mathcal{M}} \\
\mathcal{M}'_r \ar@{..>}[r] & M(\mathcal{M})
}
$$ il existe
une flèche pointillée dans $\textit{Mod}(\mathcal{A}, \text{d})$
rendant le diagramme commutatif.
\end{enumerate}
\end{lemma}

\begin{proof}
Définissons $M(\mathcal{M})$ comme la somme amalgamée du diagramme suivant
$$
\xymatrix{
\bigoplus_{(r, \varphi)} \mathcal{M}_r \ar[r] \ar[d] &
\mathcal{M} \ar[d] \\
\bigoplus_{(r, \varphi)} \mathcal{M}'_r \ar[r] &
M(\mathcal{M})
}
$$ où la somme
directe porte sur toutes les paires $(r, \varphi)$ avec
$r \in R$ et $\varphi \in
\Hom_{\textit{Mod}(\mathcal{A}, \text{d})}(\mathcal{M}_r, \mathcal{M})$. Puisque tout changement de cobase
d'un morphisme injectif est injectif, nous voyons que
$\mathcal{M} \to M(\mathcal{M})$ est injectif. Puisque le
conoyau de la flèche verticale gauche est acyclique, nous voyons
que le conoyau, qui lui est isomorphe, de $\mathcal{M} \to M(\mathcal{M})$ est
acyclique ; ainsi $\mathcal{M} \to M(\mathcal{M})$ est un
quasi-isomorphisme. La propriété (2) est vérifiée par construction.
Nous omettons la vérification que
cette procédure peut être transformée en foncteur.
\end{proof}

\begin{theorem}
\label{sdga-theorem-qis-into-dg-injective}
Soit $(\mathcal{C}, \mathcal{O})$ un site annelé. Soit
$(\mathcal{A}, \text{d})$ un
faisceau d'algèbres différentielles graduées sur $(\mathcal{C}, \mathcal{O})$.
Pour tout $\mathcal{A}$-module différentiel gradué $\mathcal{M}$, il
existe un quasi-isomorphisme $\mathcal{M} \to \mathcal{I}$
où $\mathcal{I}$ est un module différentiel gradué injectif
gradué et K-injectif sur $\mathcal{A}$. En outre, la
construction est fonctorielle en $\mathcal{M}$.
\end{theorem}

\begin{proof}
Soit $R$ l'ensemble indexant les morphismes $\mathcal{M}_r \to \mathcal{M}'_r$ de
$\textit{Mod}(\mathcal{A}, \text{d})$
considérés dans le lemme \ref{sdga-lemma-better-set-of-monos}. Soit
$M$, muni de la transformation $\text{id} \to M$, le
foncteur construit au lemme \ref{sdga-lemma-functor-set-of-monos} à
partir de $R$ et des $\mathcal{M}_r \to \mathcal{M}'_r$.
Par récurrence transfinie, nous définissons une suite de foncteurs
$M_\alpha$ et des transformations naturelles $M_\beta \to M_\alpha$
pour $\alpha < \beta$, en posant
\begin{enumerate}
\item $M_0 = \text{id}$,
\item $M_{\alpha + 1} = M \circ M_\alpha$, avec la transformation naturelle
$M_\beta \to M_{\alpha + 1}$ pour $\beta < \alpha + 1$
provenant des transformations déjà construites $M_\beta \to M_\alpha$ et des
morphismes $M_\alpha \to M \circ M_\alpha$ induits par $\text{id} \to M$, et
\item $M_\alpha = \colim_{\beta < \alpha} M_\beta$ si $\alpha$ est un
ordinal limite, avec les coprojections comme transformations
$M_\beta \to M_\alpha$ pour $\alpha < \beta$.
\end{enumerate} Notons que
pour chaque $\mathcal{A}$-module différentiel gradué les applications
$\mathcal{M} \to M_\beta(\mathcal{M}) \to M_\alpha(\mathcal{M})$ sont des
quasi-isomorphismes injectifs (comme les limites inductives filtrantes sont exactes).

\medskip\noindent
Rappelons que $\textit{Mod}(\mathcal{A}, \text{d})$ est une catégorie
abélienne de Grothendieck. Ainsi,
par Injectifs, proposition \ref{injectives-proposition-objects-are-small}
(appliquée à la somme directe des $\mathcal{M}_r$ pour tous les $r \in R$),
il existe un ordinal limite $\alpha$ tel que $\mathcal{M}_r$ soit $\alpha$-petit
par rapport aux injections pour tout $r \in R$.
Nous affirmons que $\mathcal{M} \to M_\alpha(\mathcal{M})$
est le plongement fonctoriel souhaité de $\mathcal{M}$ dans un
module injectif gradué et K-injectif.

\medskip\noindent
En effet, tout morphisme $\mathcal{M}_r \to M_\alpha(\mathcal{M})$
se factorise par $M_\beta(\mathcal{M})$ pour un certain $\beta < \alpha$.
Par la construction de $M$, cela signifie que
$\mathcal{M}_r \to M_{\beta + 1}(\mathcal{M}) = M(M_\beta(\mathcal{M}))$ se factorise
par $\mathcal{M}'_r$. Puisque
$M_\beta(\mathcal{M}) \subset M_{\beta + 1}(\mathcal{M})
\subset M_\alpha(\mathcal{M})$, on obtient la factorisation souhaitée dans
$M_\alpha(\mathcal{M})$. Nous concluons par notre choix de
$R$ et $\mathcal{M}_r \to \mathcal{M}'_r$ dans le
lemme \ref{sdga-lemma-better-set-of-monos}.
\end{proof}










\section{La catégorie dérivée}
\label{sdga-section-derived}

\noindent
Cette section est l'analogue de l'Algèbre différentielle graduée, section
\ref{dga-section-derived}.

\medskip\noindent
Soit $(\mathcal{C}, \mathcal{O})$ un site annelé. Soit
$(\mathcal{A}, \text{d})$ un faisceau d'algèbres différentielles graduées
sur $(\mathcal{C}, \mathcal{O})$. Nous construirons la catégorie
dérivée $D(\mathcal{A}, \text{d})$ en inversant les quasi-isomorphismes
dans $K(\textit{Mod}(\mathcal{A}, \text{d}))$.

\begin{lemma}
\label{sdga-lemma-cohomology-homological}
Soit $(\mathcal{C}, \mathcal{O})$ un site annelé. Soit
$(\mathcal{A}, \text{d})$ un faisceau d'algèbres différentielles graduées
sur $(\mathcal{C}, \mathcal{O})$. Le foncteur
$H^0 : \textit{Mod}(\mathcal{A}, \text{d}) \to \textit{Mod}(\mathcal{O})$
de la section \ref{sdga-section-modules} se factorise par
un
foncteur $$
H^0 : K(\textit{Mod}(\mathcal{A}, \text{d})) \to \textit{Mod}(\mathcal{O})
$$
qui est homologique au sens des
Catégories dérivées, définition \ref{derived-definition-homological}.
\end{lemma}

\begin{proof}
Il découle immédiatement des définitions qu'il existe un
diagramme commutatif
$$
\xymatrix{
\textit{Mod}(\mathcal{A}, \text{d}) \ar[r] \ar[d] &
K(\textit{Mod}(\mathcal{A}, \text{d})) \ar[d] \\
\text{Comp}(\mathcal{O}) \ar[r] &
K(\textit{Mod}(\mathcal{O}))
}
$$
Puisque $H^0(\mathcal{M})$ est le faisceau de cohomologie de
degré zéro du complexe de $\mathcal{O}$-modules sous-jacent à $\mathcal{M}$,
le lemme découle du cas des complexes de $\mathcal{O}$-modules, qui
est un cas particulier de
Catégories dérivées, lemme \ref{derived-lemma-cohomology-homological}.
\end{proof}

\begin{lemma}
\label{sdga-lemma-acyclics}
Soit $(\mathcal{C}, \mathcal{O})$ un site annelé. Soit
$(\mathcal{A}, \text{d})$ un faisceau d'algèbres différentielles graduées
sur $(\mathcal{C}, \mathcal{O})$. La sous-catégorie pleine $\text{Ac}$
de la catégorie homotopique $K(\textit{Mod}(\mathcal{A}, \text{d}))$
composée de modules acycliques est une sous-catégorie triangulée
strictement pleine et saturée de $K(\textit{Mod}(\mathcal{A}, \text{d}))$.
\end{lemma}

\begin{proof}
Bien sûr, un objet $\mathcal{M}$ de $K(\textit{Mod}(\mathcal{A}, \text{d}))$ est
dans $\text{Ac}$ si et seulement si $H^i(\mathcal{M}) = H^0(\mathcal{M}[i])$ est
nul pour tout $i$. Le lemme résulte de cette
observation, du lemme \ref{sdga-lemma-cohomology-homological}
et de Catégories dérivées, lemme \ref{derived-lemma-homological-functor-kernel}.
Voir aussi Catégories dérivées, définitions \ref{derived-definition-saturated}
et \ref{derived-definition-triangulated-subcategory}
et lemme \ref{derived-lemma-triangulated-subcategory}.
\end{proof}

\begin{lemma}
\label{sdga-lemma-qis}
Soit $(\mathcal{C}, \mathcal{O})$ un site annelé. Soit
$(\mathcal{A}, \text{d})$ un faisceau d'algèbres différentielles graduées
sur $(\mathcal{C}, \mathcal{O})$.
Considérons la
sous-classe $\text{Qis} \subset \text{Arrows}(K(\textit{Mod}(\mathcal{A}, \text{d})))$
composée de quasi-isomorphismes. Il s'agit d'un système multiplicatif
saturé compatible avec la structure triangulée sur
$K(\textit{Mod}(\mathcal{A}, \text{d}))$.
\end{lemma}

\begin{proof}
Observons que si $f , g : \mathcal{M} \to \mathcal{N}$ sont des morphismes
de $\textit{Mod}(\mathcal{A}, \text{d})$ qui sont homotopes,
alors $f$ est un quasi-isomorphisme si et seulement si $g$ est un quasi-isomorphisme.
À savoir, les applications $H^i(f) = H^0(f[i])$ et $H^i(g) = H^0(g[i])$
sont les mêmes par le lemme \ref{sdga-lemma-cohomology-homological}.
Il est donc sans ambiguïté de dire qu'un morphisme de la
catégorie homotopique $K(\textit{Mod}(\mathcal{A}, \text{d}))$ est un
quasi-isomorphisme. Pour les définitions de « système multiplicatif »,
« saturé » et « compatible avec la structure triangulée
», voir Catégories dérivées, définition \ref{derived-definition-localization},
et
Catégories, définitions \ref{categories-definition-multiplicative-system}
et \ref{categories-definition-saturated-multiplicative-system}.

\medskip\noindent
Pour démontrer le lemme, considérons la composition de
foncteurs exacts de catégories triangulées
$$
K(\textit{Mod}(\mathcal{A}, \text{d}))
\longrightarrow
K(\textit{Mod}(\mathcal{O}))
\longrightarrow
D(\mathcal{O})
$$ et
remarquons qu'un morphisme $f : \mathcal{M} \to \mathcal{N}$ de
$K(\textit{Mod}(\mathcal{A}, \text{d}))$ appartient à $\text{Qis}$ si
et seulement si son image est un isomorphisme dans $D(\mathcal{O})$.
Ainsi, le lemme découle de Catégories dérivées, lemme
\ref{derived-lemma-triangle-functor-localize}.
\end{proof}

\noindent
Dans le cas du lemme \ref{sdga-lemma-qis}, nous pouvons appliquer
Catégories dérivées, proposition
\ref{derived-proposition-construct-localization} pour
obtenir un foncteur exact de catégories triangulées
$$ Q
:
K(\textit{Mod}(\mathcal{A}, \text{d}))
\longrightarrow
\text{Qis}^{-1}K(\textit{Mod}(\mathcal{A}, \text{d}))
$$ Cependant,
comme $\textit{Mod}(\mathcal{A}, \text{d})$ est une « grande »
catégorie, c'est-à-dire que ses objets forment une classe propre, il n'est pas
immédiatement clair que, pour $\mathcal{M}$ et $\mathcal{N}$ donnés, la
construction de $\text{Qis}^{-1}K(\textit{Mod}(\mathcal{A}, \text{d}))$
produise un {\it ensemble}
$$
\Mor_{\text{Qis}^{-1}K(\textit{Mod}(\mathcal{A}, \text{d}))}
(\mathcal{M}, \mathcal{N})
$$
de morphismes. La construction de complexes K-injectifs permet
toutefois de le démontrer. En effet, par le théorème \ref{sdga-theorem-qis-into-dg-injective},
choisissons un quasi-isomorphisme $s_0 : \mathcal{N} \to \mathcal{I}$ où
$\mathcal{I}$ est un module différentiel gradué injectif
gradué et K-injectif sur $\mathcal{A}$. Rappelons ensuite que les éléments
de l'ensemble affiché sont des classes d'équivalence des
paires $(f : \mathcal{M} \to \mathcal{N}', s : \mathcal{N} \to \mathcal{N}')$
où $f$ est un morphisme arbitraire de
$K(\textit{Mod}(\mathcal{A}, \text{d}))$ et $s$ est un quasi-isomorphisme, voir la
description du calcul gauche des fractions dans Catégories,
section \ref{categories-section-localization}. Par le lemme
\ref{sdga-lemma-second-property-dg-injective} nous pouvons
choisir la flèche pointillée
$$
\xymatrix{
\mathcal{M} \ar[rd]^f & &
\mathcal{N} \ar[ld]_s \ar[rd]^{s_0} \\ &
\mathcal{N}' \ar@{..>}[rr]^{s'} & & \mathcal{I}
}
$$ rendant
le diagramme commutatif (dans la catégorie homotopique).
Ainsi, la paire $(f, s)$ est équivalente à la paire $(s' \circ f, s_0)$ et
nous trouvons que la collection de classes d'équivalence forme un ensemble.

\begin{definition}
\label{sdga-definition-derived-category} Soit
$(\mathcal{C}, \mathcal{O})$ un site annelé. Soit
$(\mathcal{A}, \text{d})$ un faisceau d'algèbres différentielles graduées sur
$(\mathcal{C}, \mathcal{O})$. Soit $\text{Qis}$ comme dans le lemme
\ref{sdga-lemma-qis}. La
{\it catégorie dérivée de $(\mathcal{A}, \text{d})$} est la catégorie
triangulée
$$
D(\mathcal{A}, \text{d}) =
\text{Qis}^{-1}K(\textit{Mod}(\mathcal{A}, \text{d}))
$$ dont la construction a été
discutée plus en détail ci-dessus.
\end{definition}

\noindent
Nous prouvons quelques faits sur cette construction.

\begin{lemma}
\label{sdga-lemma-kernel-localization}
Dans la définition \ref{sdga-definition-derived-category},
le noyau du foncteur de localisation
$Q : K(\textit{Mod}(\mathcal{A}, \text{d})) \to D(\mathcal{A}, \text{d})$
est la catégorie $\text{Ac}$ du lemme \ref{sdga-lemma-acyclics}.
\end{lemma}

\begin{proof}
Cela découle immédiatement
de Catégories dérivées, lemme \ref{derived-lemma-kernel-localization},
et du fait que $0 \to \mathcal{M}$ est un quasi-isomorphisme si
et seulement si $\mathcal{M}$ est acyclique.
\end{proof}

\begin{lemma}
\label{sdga-lemma-H0-over-D}
Dans la définition \ref{sdga-definition-derived-category}, le foncteur
$H^0 : K(\textit{Mod}(\mathcal{A}, \text{d})) \to
\textit{Mod}(\mathcal{O})$ se factorise par un foncteur homologique
$H^0 : D(\mathcal{A}, \text{d}) \to \textit{Mod}(\mathcal{O})$.
\end{lemma}

\begin{proof}
Cela découle immédiatement
de Catégories dérivées, lemme \ref{derived-lemma-universal-property-localization}.
\end{proof}

\noindent
Voici le lemme annoncé qui calcule les ensembles de morphismes dans la
catégorie dérivée.

\begin{lemma}
\label{sdga-lemma-hom-derived}
Soit $(\mathcal{C}, \mathcal{O})$ un site annelé. Soit
$(\mathcal{A}, \text{d})$ un faisceau d'algèbres différentielles graduées
sur $(\mathcal{C}, \mathcal{O})$.
Soient $\mathcal{M}$ et $\mathcal{N}$ des
$\mathcal{A}$-modules différentiels gradués. Soit $\mathcal{N} \to \mathcal{I}$
un quasi-isomorphisme, où $\mathcal{I}$ est un
$\mathcal{A}$-module différentiel gradué injectif gradué et K-injectif. Alors
$$
\Hom_{D(\mathcal{A}, \text{d})}(\mathcal{M}, \mathcal{N}) =
\Hom_{K(\textit{Mod}(\mathcal{A}, \text{d}))}(\mathcal{M}, \mathcal{I})
$$
\end{lemma}

\begin{proof}
Puisque $\mathcal{N} \to \mathcal{I}$ est un quasi-isomorphisme
nous voyons
que $$
\Hom_{D(\mathcal{A}, \text{d})}(\mathcal{M}, \mathcal{N})
= \Hom_{D(\mathcal{A}, \text{d})}(\mathcal{M}, \mathcal{I})
$$
Dans la discussion précédant la définition \ref{sdga-definition-derived-category},
nous avons constaté, en utilisant le lemme \ref{sdga-lemma-second-property-dg-injective},
que tout morphisme $\mathcal{M} \to \mathcal{I}$
dans $D(\mathcal{A}, \text{d})$ peut être représenté
par un morphisme $f : \mathcal{M} \to \mathcal{I}$
dans $K(\textit{Mod}(\mathcal{A}, \text{d}))$.
Si $f, f' :  \mathcal{M} \to \mathcal{I}$ sont deux
morphismes de $K(\textit{Mod}(\mathcal{A}, \text{d}))$, ils
définissent le même morphisme dans $D(\mathcal{A}, \text{d})$ si et
seulement s'il existe un quasi-isomorphisme $g : \mathcal{I} \to \mathcal{K}$
dans $K(\textit{Mod}(\mathcal{A}, \text{d}))$
tel que $g \circ f = g \circ f'$ ; voir
Catégories, lemme \ref{categories-lemma-equality-morphisms-left-localization}.
Cependant, par le lemme \ref{sdga-lemma-second-property-dg-injective}, il
existe une
application $h : \mathcal{K} \to \mathcal{I}$
telle que $h \circ g = \text{id}_\mathcal{I}$
dans $K(\textit{Mod}(\mathcal{A}, \text{d}))$.
Ainsi, $g \circ f = g \circ f'$ signifie $f = f'$ et
la preuve est complète.
\end{proof}

\begin{lemma}
\label{sdga-lemma-derived-products} Soit
$(\mathcal{C}, \mathcal{O})$ un site annelé. Soit
$(\mathcal{A}, \text{d})$ un faisceau d'algèbres différentielles graduées sur
$(\mathcal{C}, \mathcal{O})$. Alors
\begin{enumerate}
\item $D(\mathcal{A}, \text{d})$ possède des sommes directes et des produits,
\item les sommes directes se calculent en prenant les sommes directes de
$\mathcal{A}$-modules différentiels gradués,
\item les produits se calculent en prenant les produits
de modules différentiels gradués K-injectifs.
\end{enumerate}
\end{lemma}

\begin{proof}
Nous utiliserons que $\textit{Mod}(\mathcal{A}, \text{d})$ est une catégorie abélienne
avec des sommes directes et des produits arbitraires, et que ceux-ci donnent
lieu à des sommes directes et des produits dans $K(\textit{Mod}(\mathcal{A}, \text{d}))$.
Voir lemmes \ref{sdga-lemma-dgm-abelian} et \ref{sdga-lemma-homotopy-direct-sums}.

\medskip\noindent
Soit $\mathcal{M}_j$ une famille de $\mathcal{A}$-modules différentiels gradués.
Considérons la somme directe $\mathcal{M} = \bigoplus \mathcal{M}_j$
comme un $\mathcal{A}$-module différentiel gradué.
Pour un $\mathcal{A}$-module différentiel gradué
$\mathcal{N}$, choisissons un quasi-isomorphisme
$\mathcal{N} \to \mathcal{I}$
où $\mathcal{I}$ est injectif gradué et K-injectif comme
$\mathcal{A}$-module différentiel gradué. Voir le
théorème \ref{sdga-theorem-qis-into-dg-injective}.
En utilisant le lemme \ref{sdga-lemma-hom-derived},
nous avons \begin{align*}
\Hom_{D(\mathcal{A}, \text{d})}(\mathcal{M}, \mathcal{N})
&
= \Hom_{K(\mathcal{A}, \text{d})}(\mathcal{M}, \mathcal{I}) \\
&
= \prod \Hom_{K(\mathcal{A}, \text{d})}(\mathcal{M}_j, \mathcal{I}) \\
&
= \prod \Hom_{D(\mathcal{A}, \text{d})}(\mathcal{M}_j, \mathcal{I})
\end{align*}
ce qui prouve l'existence des sommes directes dans $D(A, \text{d})$,
comme indiqué en (2).

\medskip\noindent
Soit $\mathcal{M}_j$ une famille de $\mathcal{A}$-modules différentiels gradués.
Pour chaque $j$, choisissons un
quasi-isomorphisme $\mathcal{M} \to \mathcal{I}_j$
où $\mathcal{I}_j$ est injectif gradué et K-injectif
comme $\mathcal{A}$-module différentiel
gradué. Considérons le produit $\mathcal{I} = \prod \mathcal{I}_j$
de modules différentiels gradués sur $\mathcal{A}$.
Par les lemmes \ref{sdga-lemma-product-K-injective}
et \ref{sdga-lemma-product-graded-injective},
$\mathcal{I}$ est injectif gradué et K-injectif
comme $\mathcal{A}$-module différentiel
gradué. Pour un $\mathcal{A}$-module différentiel
gradué $\mathcal{N}$, le lemme \ref{sdga-lemma-hom-derived} donne
\begin{align*}
\Hom_{D(\mathcal{A}, \text{d})}(\mathcal{N}, \mathcal{I})
&
= \Hom_{K(\mathcal{A}, \text{d})}(\mathcal{N}, \mathcal{I}) \\
&
= \prod \Hom_{K(\mathcal{A}, \text{d})}(\mathcal{N}, \mathcal{I}_j) \\
&
= \prod \Hom_{D(\mathcal{A}, \text{d})}(\mathcal{N}, \mathcal{M}_j)
\end{align*}
ce qui prouve l'existence des produits dans $D(\mathcal{A}, \text{d})$,
comme indiqué en (3).
\end{proof}







\section{Le foncteur delta canonique}
\label{sdga-section-canonical-delta-functor}

\noindent
Soit $(\mathcal{C}, \mathcal{O})$ un site annelé. Soit
$(\mathcal{A}, \text{d})$ un faisceau d'algèbres différentielles graduées
sur $(\mathcal{C}, \mathcal{O})$.
Considérons le foncteur
$\textit{Mod}(\mathcal{A}, \text{d}) \to
K(\textit{Mod}(\mathcal{A}, \text{d}))$.
Ce foncteur n'est {\bf pas} un foncteur $\delta$ en
général. Cependant, il s'avère que le foncteur
$\textit{Mod}(\mathcal{A}, \text{d}) \to D(A, \text{d})$ est
un foncteur $\delta$. Pour voir cela, nous devons définir
les morphismes $\delta$ associés à une suite exacte courte
$$
0 \to \mathcal{K} \xrightarrow{a} \mathcal{L} \xrightarrow{b} \mathcal{M} \to 0
$$
dans la catégorie abélienne $\textit{Mod}(\mathcal{A}, \text{d})$.
Considérons le cône $C(a)$ du morphisme $a$ avec ses morphismes
canoniques $i : \mathcal{L} \to C(a)$ et
$p : C(a) \to \mathcal{K}[1]$, voir définition \ref{sdga-definition-cone}. Il existe un
homomorphisme de modules différentiels gradués sur $\mathcal{A}$
$$ q :
C(a) \longrightarrow \mathcal{M}
$$ par
Algèbre différentielle graduée, lemme \ref{dga-lemma-cone}
(que nous pouvons utiliser dans la discussion
de la section \ref{sdga-section-conclude-triangulated})
appliqué au
diagramme $$
\xymatrix{
\mathcal{K} \ar[r]_a \ar[d]
& \mathcal{L} \ar[d]^b \\
0 \ar[r]
& \mathcal{M}
}
$$
Le morphisme $q$ est un quasi-isomorphisme, par exemple parce
que c'est le cas pour les complexes sous-jacents de $\mathcal{O}$-modules ;
voir
Catégories dérivées, section \ref{derived-section-canonical-delta-functor}.
D'après Algèbre différentielle graduée, lemme \ref{dga-lemma-cone-homotopy} (applicable
par la discussion de la section
\ref{sdga-section-conclude-triangulated}), le
triangle
$$
(\mathcal{K}, \mathcal{L}, C(a), a, i, -p)
$$ est un
triangle distingué en $K(\textit{Mod}(\mathcal{A}, \text{d}))$. Comme le foncteur
de localisation
$K(\textit{Mod}(\mathcal{A}, \text{d})) \to D(\mathcal{A}, \text{d})$ est exact,
nous voyons que $(\mathcal{K}, \mathcal{L}, C(a), a, i, -p)$ est un
triangle distingué
dans $D(\mathcal{A}, \text{d})$. Puisque $q$ est un quasi-isomorphisme,
nous voyons que $q$ est un isomorphisme dans $D(\mathcal{A}, \text{d})$.
Nous déduisons
donc que $$
(\mathcal{K}, \mathcal{L}, \mathcal{M}, a, b, -p \circ q^{-1})
$$
est un triangle distingué de $D(\mathcal{A}, \text{d})$.
Cela conduit au lemme suivant.

\begin{lemma}
\label{sdga-lemma-derived-canonical-delta-functor}
Soit $(\mathcal{C}, \mathcal{O})$ un site annelé. Soit
$(\mathcal{A}, \text{d})$ un faisceau d'algèbres différentielles graduées
sur $(\mathcal{C}, \mathcal{O})$. Le foncteur de localisation
$\textit{Mod}(\mathcal{A}, \text{d}) \to D(\mathcal{A}, \text{d})$ a
la structure naturelle d'un foncteur $\delta$, avec
$$
\delta_{\mathcal{K} \to \mathcal{L} \to \mathcal{M}} = - p \circ q^{-1}
$$ avec
$p$ et $q$ comme expliqué ci-dessus.
\end{lemma}

\begin{proof}
Nous avons déjà vu que ce choix conduit à un triangle distingué
chaque fois qu'une suite exacte courte de complexes est donnée. Il
reste à démontrer la fonctorialité de cette construction ;
voir Catégories dérivées, définition \ref{derived-definition-delta-functor}.
Cela découle, avec quelques vérifications, d'Algèbre différentielle graduée, lemme \ref{dga-lemma-cone},
applicable
par la discussion de la section \ref{sdga-section-conclude-triangulated}.
Comparer avec
Catégories dérivées, lemme \ref{derived-lemma-derived-canonical-delta-functor}.
\end{proof}

\begin{lemma}
\label{sdga-lemma-homotopy-colimit}
Soit $(\mathcal{C}, \mathcal{O})$ un site annelé. Soit
$(\mathcal{A}, \text{d})$ un faisceau d'algèbres différentielles graduées sur
$(\mathcal{C}, \mathcal{O})$. Soit
$\mathcal{M}_n$ un système de modules différentiels gradués sur $\mathcal{A}$.
Ensuite, la limite inductive dérivée $\text{hocolim} \mathcal{M}_n$
dans $D(\mathcal{A}, \text{d})$ est représentée
par le module différentiel gradué $\colim \mathcal{M}_n$.
\end{lemma}

\begin{proof}
Posons $\mathcal{M} = \colim \mathcal{M}_n$.
Nous avons une suite exacte de $\mathcal{A}$-modules différentiels gradués
$$
0 \to \bigoplus \mathcal{M}_n \to \bigoplus \mathcal{M}_n \to \mathcal{M} \to 0
$$
par Catégories dérivées, lemme \ref{derived-lemma-compute-colimit},
appliqué aux complexes sous-jacents de $\mathcal{O}$-modules. Les
sommes directes sont les sommes directes en $D(\mathcal{A}, \text{d})$
par le lemme \ref{sdga-lemma-derived-products}.
Ainsi, le résultat résulte de la
définition des limites inductives dérivées dans Catégories
dérivées, définition \ref{derived-definition-derived-colimit}
et du fait qu'une suite exacte courte de
complexes donne un triangle distingué
(lemme \ref{sdga-lemma-derived-canonical-delta-functor}).
\end{proof}






\section{Image inverse dérivée}
\label{sdga-section-derived-pullback}

\noindent
Soit $(f, f^\sharp) : (\Sh(\mathcal{C}), \mathcal{O}_\mathcal{C})
\to (\Sh(\mathcal{D}), \mathcal{O}_\mathcal{D})$ un morphisme
de topos annelés. Soit $\mathcal{A}$ une
$\mathcal{O}_\mathcal{C}$-algèbre différentielle graduée. Soit $\mathcal{B}$ une
$\mathcal{O}_\mathcal{D}$-algèbre différentielle graduée. Supposons donné un
morphisme
$$
\varphi : f^{-1}\mathcal{B} \to \mathcal{A}
$$ de
$f^{-1}\mathcal{O}_\mathcal{D}$-algèbres différentielles graduées. Par l'adjonction
entre restriction et extension des scalaires, cela revient à se donner
un morphisme $\varphi : f^*\mathcal{B} \to \mathcal{A}$ de
$\mathcal{O}_\mathcal{C}$-algèbres différentielles graduées ou, de manière équivalente,
$\varphi$ peut être considéré comme un morphisme
$$
\varphi : \mathcal{B} \to f_*\mathcal{A}
$$
de $\mathcal{O}_\mathcal{D}$-algèbres différentielles graduées.
Voir la remarque \ref{sdga-remark-functoriality-dga}.

\medskip\noindent
Supposons de plus que $\mathcal{A}'$ soit une seconde
$\mathcal{O}_\mathcal{C}$-algèbre différentielle graduée et que $\mathcal{N}$ soit un
$(\mathcal{A}, \mathcal{A}')$-bimodule différentiel gradué. Dans
ce contexte, nous pouvons considérer le foncteur
$$
\textit{Mod}(\mathcal{B}, \text{d})
\longrightarrow
\textit{Mod}(\mathcal{A}', \text{d}),\quad
\mathcal{M} \longmapsto f^*\mathcal{M} \otimes_{\mathcal{A}} \mathcal{N}
$$
Notons que cela s'étend à un foncteur
$$
\textit{Mod}^{dg}(\mathcal{B}, \text{d})
\longrightarrow
\textit{Mod}^{dg}(\mathcal{A}', \text{d}),\quad
\mathcal{M} \longmapsto f^*\mathcal{M} \otimes_{\mathcal{A}} \mathcal{N}
$$ des
catégories différentielles graduées par la discussion des
sections \ref{sdga-section-functoriality-dg} et \ref{sdga-section-dg-bimodules}. Il
s'ensuit formellement que nous obtenons également un foncteur
exact \begin{equation}
\label{sdga-equation-pullback}
K(\textit{Mod}(\mathcal{B}, \text{d}))
\longrightarrow
K(\textit{Mod}(\mathcal{A}', \text{d})),\quad
\mathcal{M} \longmapsto f^*\mathcal{M} \otimes_{\mathcal{A}} \mathcal{N}
\end{equation}
de catégories triangulées.

\begin{lemma}
\label{sdga-lemma-derived-tensor-product} Dans
la situation ci-dessus, le foncteur (\ref{sdga-equation-pullback}),
composé avec le foncteur de localisation
$K(\textit{Mod}(\mathcal{A}', \text{d})) \to D(\mathcal{A}', \text{d})$, possède
une extension dérivée à
gauche $D(\mathcal{B}, \text{d}) \to D(\mathcal{A}', \text{d})$ dont
la valeur sur un bon $\mathcal{B}$-module différentiel gradué à droite
$\mathcal{P}$ est $f^*\mathcal{P} \otimes_\mathcal{A} \mathcal{N}$.
\end{lemma}

\begin{proof}
Rappelons que pour tout $\mathcal{B}$-module différentiel gradué à
droite $\mathcal{M}$ il existe un quasi-isomorphisme $\mathcal{P} \to \mathcal{M}$
avec $\mathcal{P}$ un bon $\mathcal{B}$-module différentiel gradué. Voir
lemme \ref{sdga-lemma-resolve}. Par
conséquent, d'après Catégories dérivées, lemme \ref{derived-lemma-find-existence-computes},
il suffit de montrer que, pour tout
quasi-isomorphisme $\mathcal{P} \to \mathcal{P}'$ de bons
$\mathcal{B}$-modules différentiels gradués, l'application
induite $$
f^*\mathcal{P} \otimes_\mathcal{A} \mathcal{N}
\longrightarrow
f^*\mathcal{P}' \otimes_\mathcal{A} \mathcal{N}
$$
est un quasi-isomorphisme. Le cône $\mathcal{P}''$
sur $\mathcal{P} \to \mathcal{P}'$ est un bon
$\mathcal{A}$-module différentiel gradué par le lemme \ref{sdga-lemma-good-admissible-ses}.
Puisque nous avons un triangle
distingué $$
\mathcal{P} \to \mathcal{P}' \to \mathcal{P}'' \to \mathcal{P}[1]
$$
dans $K(\textit{Mod}(\mathcal{B}, \text{d}))$ nous obtenons un
triangle
distingué $$
f^*\mathcal{P} \otimes_\mathcal{A} \mathcal{N} \to
f^*\mathcal{P}' \otimes_\mathcal{A} \mathcal{N} \to
f^*\mathcal{P}'' \otimes_\mathcal{A} \mathcal{N} \to
f^*\mathcal{P}[1] \otimes_\mathcal{A} \mathcal{N}
$$
dans $K(\textit{Mod}(\mathcal{A}', \text{d}))$. Par
le lemme \ref{sdga-lemma-acyclic-good}
le module différentiel gradué
$f^*\mathcal{P}'' \otimes_\mathcal{A} \mathcal{N}$
est acyclique et la preuve est complète.
\end{proof}

\begin{definition}
\label{sdga-definition-pullback}
Produit tensoriel dérivé et image inverse dérivée.
\begin{enumerate}
\item Soit $(\mathcal{C}, \mathcal{O})$ un site annelé.
Soient $\mathcal{A}$ et $\mathcal{B}$ des $\mathcal{O}$-algèbres différentielles
graduées. Soit $\mathcal{N}$ un
$(\mathcal{A}, \mathcal{B})$-bimodule différentiel
gradué. Le foncteur $D(\mathcal{A}, \text{d}) \to D(\mathcal{B}, \text{d})$
construit dans le lemme \ref{sdga-lemma-derived-tensor-product}
est appelé le {\it produit tensoriel dérivé} et noté
$- \otimes_\mathcal{A}^\mathbf{L} \mathcal{N}$.
\item Soit $(f, f^\sharp) : (\Sh(\mathcal{C}), \mathcal{O}_\mathcal{C})
\to (\Sh(\mathcal{D}), \mathcal{O}_\mathcal{D})$ un
morphisme de topos annelés. Soit $\mathcal{A}$ une
$\mathcal{O}_\mathcal{C}$-algèbre différentielle graduée. Soit $\mathcal{B}$ une
$\mathcal{O}_\mathcal{D}$-algèbre différentielle graduée. Soit
$\varphi : \mathcal{B} \to f_*\mathcal{A}$ un homomorphisme de
$\mathcal{O}_\mathcal{D}$-algèbres différentielles graduées. Le
foncteur $D(\mathcal{B}, \text{d}) \to D(\mathcal{A}, \text{d})$
construit dans le lemme \ref{sdga-lemma-derived-tensor-product} est
appelé {\it image inverse dérivée}
et noté $Lf^*$.
\end{enumerate}
\end{definition}

\noindent
Avec ce langage en place, nous pouvons exprimer certaines compatibilités évidentes.

\begin{lemma}
\label{sdga-lemma-compose-pullback-tensor}
Dans le lemme \ref{sdga-lemma-derived-tensor-product}, le foncteur
$D(\mathcal{B}, \text{d}) \to D(\mathcal{A}', \text{d})$ est égal
à $\mathcal{M} \mapsto
Lf^*\mathcal{M} \otimes_\mathcal{A}^\mathbf{L} \mathcal{N}$.
\end{lemma}

\begin{proof}
Cela résulte immédiatement du fait que nous pouvons calculer ces
foncteurs en représentant les objets par de bons modules différentiels
gradués et parce que $f^*\mathcal{P}$ est un bon
$\mathcal{A}$-module différentiel gradué si $\mathcal{P}$ est un bon
$\mathcal{B}$-module différentiel gradué.
\end{proof}

\begin{lemma}
\label{sdga-lemma-compose-pullback}
Soient $(f, f^\sharp) : (\Sh(\mathcal{C}), \mathcal{O})
\to (\Sh(\mathcal{C}'), \mathcal{O}')$ et
$(g, g^\sharp) : (\Sh(\mathcal{C}'), \mathcal{O}')
\to (\Sh(\mathcal{C}''), \mathcal{O}'')$ des
morphismes de topos annelés. Soient $\mathcal{A}$, $\mathcal{A}'$ et
$\mathcal{A}''$ des algèbres différentielles graduées respectivement sur $\mathcal{O}$,
$\mathcal{O}'$ et $\mathcal{O}''$. Soient
$\varphi : \mathcal{A}' \to f_*\mathcal{A}$ et
$\varphi' : \mathcal{A}'' \to g_*\mathcal{A}'$ des
homomorphismes d'algèbres différentielles graduées respectivement sur $\mathcal{O}'$
et $\mathcal{O}''$.
Alors nous avons $L(g \circ f)^* = Lf^* \circ Lg^*
: D(\mathcal{A}'', \text{d}) \to D(\mathcal{A}, \text{d})$.
\end{lemma}

\begin{proof}
Cela résulte immédiatement du calcul de ces foncteurs au
moyen de bons modules différentiels gradués, et du fait
que $f^*\mathcal{P}$ est un bon
$\mathcal{A}'$-module différentiel gradué si $\mathcal{P}$ est un bon
$\mathcal{A}$-module différentiel gradué.
\end{proof}

\noindent
Soit $(\mathcal{C}, \mathcal{O})$ un site annelé. Soient
$\mathcal{A}$ et $\mathcal{B}$ des $\mathcal{O}$-algèbres différentielles graduées.
Soit $\mathcal{N} \to \mathcal{N}'$ un homomorphisme de bimodules différentiels
gradués sur $(\mathcal{A}, \mathcal{B})$. On obtient alors des applications
canoniques
$$
t
: \mathcal{M} \otimes_\mathcal{A}^\mathbf{L} \mathcal{N}
\longrightarrow
\mathcal{M} \otimes_\mathcal{A}^\mathbf{L} \mathcal{N}'
$$
fonctorielles en $\mathcal{M}$ dans $D(\mathcal{A}, \text{d})$,
qui définissent une transformation naturelle entre les
foncteurs exacts $D(\mathcal{A}, \text{d}) \to D(\mathcal{B}, \text{d})$
de catégories triangulées. La valeur de $t$ sur un bon
$\mathcal{A}$-module différentiel gradué $\mathcal{P}$ est l'application
évidente $$
\mathcal{P} \otimes_\mathcal{A}^\mathbf{L} \mathcal{N}
= \mathcal{P} \otimes_\mathcal{A} \mathcal{N}
\longrightarrow
\mathcal{P} \otimes_\mathcal{A} \mathcal{N}'
= \mathcal{P} \otimes_\mathcal{A}^\mathbf{L} \mathcal{N}'
$$

\begin{lemma}
\label{sdga-lemma-tensor-symmetry} Dans
la situation ci-dessus, si $\mathcal{N} \to \mathcal{N}'$ est un isomorphisme sur
les faisceaux de cohomologie, alors $t$ est un isomorphisme des
foncteurs $(- \otimes_\mathcal{A}^\mathbf{L} \mathcal{N}) \to
(- \otimes_\mathcal{A}^\mathbf{L} \mathcal{N}')$.
\end{lemma}

\begin{proof}
Il suffit de démontrer que
$\mathcal{P} \otimes_\mathcal{A} \mathcal{N} \to
\mathcal{P} \otimes_\mathcal{A} \mathcal{N}'$ est
un isomorphisme sur les faisceaux de cohomologie pour tout bon
$\mathcal{A}$-module différentiel gradué $\mathcal{P}$.
Pour cela, soit $\mathcal{N}''$ le cône du morphisme
$\mathcal{N} \to \mathcal{N}'$ comme
$\mathcal{A}$-module différentiel gradué à gauche, voir la définition \ref{sdga-definition-cone}.
(Bien sûr, $\mathcal{N}''$ est aussi un bimodule, mais
nous n'en avons pas besoin.) Par fonctorialité de la construction
tensorielle (c'est un foncteur de catégories
différentielles graduées), $\mathcal{P} \otimes_\mathcal{A} \mathcal{N}''$
est le cône, comme complexe de $\mathcal{O}$-modules, du
morphisme $\mathcal{P} \otimes_\mathcal{A} \mathcal{N} \to
\mathcal{P} \otimes_\mathcal{A} \mathcal{N}'$.
Il suffit donc de démontrer
que $\mathcal{P} \otimes_\mathcal{A} \mathcal{N}''$
est acyclique. Cela découle du fait que $\mathcal{P}$ est bon
et du fait que $\mathcal{N}''$ est acyclique en tant
que cône sur un quasi-isomorphisme.
\end{proof}

\begin{lemma}
\label{sdga-lemma-good-on-other-side}
Soit $(\mathcal{C}, \mathcal{O})$ un site annelé. Soient
$\mathcal{A}$ et $\mathcal{B}$ des $\mathcal{O}$-algèbres différentielles
graduées. Soit $\mathcal{N}$ un bimodule différentiel
gradué sur $(\mathcal{A}, \mathcal{B})$. Si $\mathcal{N}$ est
bon comme $\mathcal{A}$-module différentiel gradué à
gauche, alors $\mathcal{M} \otimes_\mathcal{A}^\mathbf{L} \mathcal{N} =
\mathcal{M} \otimes_\mathcal{A} \mathcal{N}$ pour tout
$\mathcal{A}$-module différentiel gradué $\mathcal{M}$.
\end{lemma}

\begin{proof}
Soit $\mathcal{P} \to \mathcal{M}$ un quasi-isomorphisme où
$\mathcal{P}$ est un bon module différentiel gradué à droite sur $\mathcal{A}$. Pour
prouver le lemme, nous devons montrer que
$\mathcal{P} \otimes_\mathcal{A} \mathcal{N} \to
\mathcal{M} \otimes_\mathcal{A} \mathcal{N}$ est un
quasi-isomorphisme. Le cône $C$ du morphisme
$\mathcal{P} \to \mathcal{M}$ est un module différentiel gradué acyclique
à droite sur $\mathcal{A}$. Ainsi,
$C \otimes_\mathcal{A} \mathcal{N}$ est acyclique, car $\mathcal{N}$
est supposé bon comme module différentiel gradué à gauche sur $\mathcal{A}$.
Puisque $C \otimes_\mathcal{A} \mathcal{N}$ est le cône du
morphisme $\mathcal{P} \otimes_\mathcal{A} \mathcal{N} \to
\mathcal{M} \otimes_\mathcal{A} \mathcal{N}$ comme complexe de
$\mathcal{O}$-modules, nous concluons.
\end{proof}

\begin{lemma}
\label{sdga-lemma-compose-tensor} Soit
$(\mathcal{C}, \mathcal{O})$ un site annelé. Soient
$\mathcal{A}$, $\mathcal{A}'$ et $\mathcal{A}''$ des
$\mathcal{O}$-algèbres différentielles graduées. Soient $\mathcal{N}$ et $\mathcal{N}'$ des bimodules
différentiels gradués respectivement sur $(\mathcal{A}, \mathcal{A}')$
et $(\mathcal{A}', \mathcal{A}'')$. Supposons
que l'application canonique
$$
\mathcal{N} \otimes_{\mathcal{A}'}^\mathbf{L} \mathcal{N}'
\longrightarrow
\mathcal{N} \otimes_{\mathcal{A}'} \mathcal{N}'
$$
dans $D(\mathcal{A}'', \text{d})$ soit un quasi-isomorphisme.
Alors nous
avons $$
(\mathcal{M}
\otimes_\mathcal{A}^\mathbf{L} \mathcal{N})
\otimes_{\mathcal{A}'}^\mathbf{L} \mathcal{N}'
=
\mathcal{M}
\otimes_\mathcal{A}^\mathbf{L}
(\mathcal{N} \otimes_{\mathcal{A}'} \mathcal{N}')
$$ en
tant que foncteurs $D(\mathcal{A}, \text{d}) \to D(\mathcal{A}'', \text{d})$.
\end{lemma}

\begin{proof}
Choisissons un bon $\mathcal{A}$-module différentiel gradué
$\mathcal{P}$ et un quasi-isomorphisme $\mathcal{P} \to \mathcal{M}$, voir
lemme \ref{sdga-lemma-resolve}. Alors
$$
\mathcal{M}
\otimes_\mathcal{A}^\mathbf{L}
(\mathcal{N} \otimes_{\mathcal{A}'} \mathcal{N}') =
\mathcal{P} \otimes_\mathcal{A} \mathcal{N}
\otimes_{\mathcal{A}'} \mathcal{N}'
$$ et nous
avons
$$
(\mathcal{M}
\otimes_\mathcal{A}^\mathbf{L} \mathcal{N})
\otimes_{\mathcal{A}'}^\mathbf{L} \mathcal{N}' =
(\mathcal{P} \otimes_\mathcal{A} \mathcal{N})
\otimes_{\mathcal{A}'}^\mathbf{L} \mathcal{N}'
$$ Ainsi, nous
devons montrer que l'application canonique
$$
(\mathcal{P} \otimes_\mathcal{A} \mathcal{N})
\otimes_{\mathcal{A}'}^\mathbf{L} \mathcal{N}'
\longrightarrow
\mathcal{P} \otimes_\mathcal{A} \mathcal{N}
\otimes_{\mathcal{A}'} \mathcal{N}'
$$ est un
quasi-isomorphisme. Choisissons un quasi-isomorphisme
$\mathcal{Q} \to \mathcal{N}'$ où $\mathcal{Q}$ est un
bon module différentiel gradué à gauche sur $\mathcal{A}'$
(lemme \ref{sdga-lemma-resolve}). Par
le lemme \ref{sdga-lemma-good-on-other-side}, le
morphisme ci-dessus, considéré dans la catégorie dérivée des $\mathcal{O}$-modules,
est le morphisme $$
\mathcal{P} \otimes_\mathcal{A} \mathcal{N}
\otimes_{\mathcal{A}'} \mathcal{Q}
\longrightarrow
\mathcal{P} \otimes_\mathcal{A} \mathcal{N}
\otimes_{\mathcal{A}'} \mathcal{N}'
$$
Puisque $\mathcal{N} \otimes_{\mathcal{A}'} \mathcal{Q} \to
\mathcal{N} \otimes_{\mathcal{A}'} \mathcal{N}'$ est un quasi-isomorphisme
par hypothèse et $\mathcal{P}$ est un bon
$\mathcal{A}$-module différentiel gradué, cette application est un quasi-isomorphisme
par le lemme \ref{sdga-lemma-tensor-symmetry} (le côté gauche et
droit calculent $\mathcal{P} \otimes_\mathcal{A}^\mathbf{L} (\mathcal{N}
\otimes_{\mathcal{A}'} \mathcal{Q})$
et $\mathcal{P} \otimes_\mathcal{A}^\mathbf{L} (\mathcal{N}
\otimes_{\mathcal{A}'} \mathcal{N}')$ ou l'on peut simplement
répéter l'argument dans la preuve du lemme).
\end{proof}






\section{Image directe dérivée}
\label{sdga-section-derived-pushforward}

\noindent
L'existence de suffisamment d'objets K-injectifs garantit que nous pouvons prendre le foncteur
dérivé à droite de tout foncteur exact défini sur la catégorie homotopique.

\begin{lemma}
\label{sdga-lemma-right-derived}
Soit $(\mathcal{C}, \mathcal{O})$ un site annelé. Soit
$(\mathcal{A}, \text{d})$ un faisceau d'algèbres différentielles graduées sur
$(\mathcal{C}, \mathcal{O})$. Alors tout foncteur exact
$$
T : K(\textit{Mod}(\mathcal{A}, \text{d})) \longrightarrow \mathcal{D}
$$ de
catégories triangulées possède une extension dérivée à droite
$RT : D(\mathcal{A}, \text{d}) \to \mathcal{D}$
dont la valeur sur un
$\mathcal{A}$-module différentiel gradué $\mathcal{I}$ injectif gradué et K-injectif
est $T(\mathcal{I})$.
\end{lemma}

\begin{proof}
D'après le théorème \ref{sdga-theorem-qis-into-dg-injective},
pour tout module différentiel gradué à droite sur $\mathcal{A}$,
$\mathcal{M}$, il existe un
quasi-isomorphisme $\mathcal{M} \to \mathcal{I}$ où $\mathcal{I}$
est un module différentiel
gradué injectif gradué et K-injectif sur $\mathcal{A}$.
Par conséquent, d'après Catégories dérivées, lemme \ref{derived-lemma-find-existence-computes},
il suffit de montrer
que, pour tout quasi-isomorphisme $\mathcal{I} \to \mathcal{I}'$
de $\mathcal{A}$-modules
différentiels gradués tous deux
injectifs gradués et K-injectifs, $T(\mathcal{I}) \to T(\mathcal{I}')$
est un isomorphisme. En effet, $\mathcal{I} \to \mathcal{I}'$
est un isomorphisme dans $K(\textit{Mod}(\mathcal{A}, \text{d}))$,
comme on le voit par exemple par le lemme \ref{sdga-lemma-hom-derived},
ou par le lemme \ref{sdga-lemma-second-property-dg-injective}.
\end{proof}

\noindent
Nous avons déjà vu un certain nombre de foncteurs auxquels cela s'applique.
Voici deux exemples.

\begin{definition}
\label{sdga-definition-pushforward}
Hom interne dérivé et image directe dérivée.
\begin{enumerate}
\item Soit $(\mathcal{C}, \mathcal{O})$ un site annelé. Soient
$\mathcal{A}$ et $\mathcal{B}$ des $\mathcal{O}$-algèbres différentielles graduées.
Soit $\mathcal{N}$ un
$(\mathcal{A}, \mathcal{B})$-bimodule différentiel gradué. L'extension dérivée à droite
$$
R\SheafHom_\mathcal{B}(\mathcal{N}, -) :
D(\mathcal{B}, \text{d})
\longrightarrow
D(\mathcal{A}, \text{d})
$$
du foncteur Hom interne $\SheafHom_\mathcal{B}^{dg}(\mathcal{N}, -)$ est
appelée {\it Hom interne dérivé}.
\item Soit $(f, f^\sharp) : (\Sh(\mathcal{C}), \mathcal{O}_\mathcal{C})
\to (\Sh(\mathcal{D}), \mathcal{O}_\mathcal{D})$ un morphisme
de topos annelés. Soit $\mathcal{A}$ une
$\mathcal{O}_\mathcal{C}$-algèbre différentielle graduée. Soit $\mathcal{B}$ une
$\mathcal{O}_\mathcal{D}$-algèbre différentielle graduée. Soit
$\varphi : \mathcal{B} \to f_*\mathcal{A}$ un homomorphisme de
$\mathcal{O}_\mathcal{D}$-algèbres différentielles graduées. L'extension
dérivée à droite
$$ Rf_*
:
D(\mathcal{A}, \text{d})
\longrightarrow
D(\mathcal{B}, \text{d})
$$ de
l'image directe $f_*$ est appelée {\it image directe dérivée}.
\end{enumerate}
\end{definition}

\noindent
Le foncteur
$Rf_* : D(\mathcal{A}, \text{d}) \to D(\mathcal{B}, \text{d})$ coïncide
avec l'image directe dérivée des complexes sous-jacents de
$\mathcal{O}$-modules ; voir le lemme \ref{sdga-lemma-pushforward-agrees}.

\medskip\noindent Ces
foncteurs sont les adjoints de l'image inverse dérivée et du
produit tensoriel dérivé.

\begin{lemma}
\label{sdga-lemma-derived-adjoint-tensor-hom}
Soit $(\mathcal{C}, \mathcal{O})$ un site annelé. Soient
$\mathcal{A}$ et $\mathcal{B}$ des $\mathcal{O}$-algèbres différentielles graduées.
Soit $\mathcal{N}$ un
$(\mathcal{A}, \mathcal{B})$-bimodule différentiel gradué. Alors
$$
R\SheafHom_\mathcal{B}(\mathcal{N}, -) :
D(\mathcal{B}, \text{d})
\longrightarrow
D(\mathcal{A}, \text{d})
$$ est
adjoint à droite à
$$
- \otimes_\mathcal{A}^\mathbf{L} \mathcal{N} :
D(\mathcal{A}, \text{d})
\longrightarrow
D(\mathcal{B}, \text{d})
$$
\end{lemma}

\begin{proof}
Cela découle de Catégories dérivées, lemme
\ref{derived-lemma-pre-derived-adjoint-functors-general}
et lemme \ref{sdga-lemma-tensor-hom-adjunction-dg}.
\end{proof}

\begin{lemma}
\label{sdga-lemma-derived-adjoint-push-pull}
Soit $(f, f^\sharp) : (\Sh(\mathcal{C}), \mathcal{O}_\mathcal{C})
\to (\Sh(\mathcal{D}), \mathcal{O}_\mathcal{D})$ un
morphisme de topos annelés. Soit $\mathcal{A}$ une
$\mathcal{O}_\mathcal{C}$-algèbre différentielle graduée. Soit $\mathcal{B}$ une
$\mathcal{O}_\mathcal{D}$-algèbre différentielle graduée. Soit
$\varphi : \mathcal{B} \to f_*\mathcal{A}$ un homomorphisme de
$\mathcal{O}_\mathcal{D}$-algèbres différentielles graduées.
Alors
$$ Rf_*
:
D(\mathcal{A}, \text{d})
\longrightarrow
D(\mathcal{B}, \text{d})
$$ est
adjoint à droite à
$$
Lf^* :
D(\mathcal{B}, \text{d})
\longrightarrow
D(\mathcal{A}, \text{d})
$$
\end{lemma}

\begin{proof}
Cela découle de Catégories dérivées, lemme
\ref{derived-lemma-pre-derived-adjoint-functors-general}
et lemme \ref{sdga-lemma-adjunction-push-pull-dg}.
\end{proof}

\noindent
Ensuite, nous discuterons de ce qui se passe dans la situation considérée
à la section \ref{sdga-section-derived-pullback}.

\medskip\noindent
Soit $(f, f^\sharp) : (\Sh(\mathcal{C}), \mathcal{O}_\mathcal{C})
\to (\Sh(\mathcal{D}), \mathcal{O}_\mathcal{D})$ un morphisme
de topos annelés. Soit $\mathcal{A}$ une
$\mathcal{O}_\mathcal{C}$-algèbre différentielle graduée. Soit $\mathcal{B}$ une
$\mathcal{O}_\mathcal{D}$-algèbre différentielle graduée. Supposons donné un
morphisme
$$
\varphi : f^{-1}\mathcal{B} \to \mathcal{A}
$$ de
$f^{-1}\mathcal{O}_\mathcal{D}$-algèbres différentielles graduées. Par l'adjonction
entre restriction et extension des scalaires, cela revient à se donner
un morphisme $\varphi : f^*\mathcal{B} \to \mathcal{A}$ de
$\mathcal{O}_\mathcal{C}$-algèbres différentielles graduées ou, de manière équivalente,
$\varphi$ peut être considéré comme un morphisme
$$
\varphi : \mathcal{B} \to f_*\mathcal{A}
$$
de $\mathcal{O}_\mathcal{D}$-algèbres différentielles graduées.
Voir la remarque \ref{sdga-remark-functoriality-dga}.

\medskip\noindent
Supposons de plus que $\mathcal{A}'$ soit une seconde
$\mathcal{O}_\mathcal{C}$-algèbre différentielle graduée et que $\mathcal{N}$ soit un
$(\mathcal{A}, \mathcal{A}')$-bimodule différentiel gradué. Dans ce
contexte, nous pouvons considérer le foncteur
$$
\textit{Mod}(\mathcal{A}', \text{d})
\longrightarrow
\textit{Mod}(\mathcal{B}, \text{d}),\quad
\mathcal{M} \longmapsto
f_*\SheafHom^{dg}_{\mathcal{A}'}(\mathcal{N}, \mathcal{M})
$$
Notons que cela s'étend à un foncteur
$$
\textit{Mod}^{dg}(\mathcal{A}', \text{d})
\longrightarrow
\textit{Mod}^{dg}(\mathcal{B}, \text{d}),\quad
\mathcal{M} \longmapsto
f_*\SheafHom^{dg}_{\mathcal{A}'}(\mathcal{N}, \mathcal{M})
$$ des
catégories différentielles graduées par la discussion des
sections \ref{sdga-section-functoriality-dg} et \ref{sdga-section-dg-bimodules}. Il
s'ensuit formellement que nous obtenons également un foncteur
exact \begin{equation}
\label{sdga-equation-pushforward}
K(\textit{Mod}(\mathcal{A}', \text{d}))
\longrightarrow
K(\textit{Mod}(\mathcal{B}, \text{d})),\quad
\mathcal{M} \longmapsto
f_*\SheafHom^{dg}_{\mathcal{A}'}(\mathcal{N}, \mathcal{M})
\end{equation}
de catégories triangulées.

\begin{lemma}
\label{sdga-lemma-compose-pushforward-hom} Dans
la situation ci-dessus, notons $RT : D(\mathcal{A}', \text{d}) \to
D(\mathcal{B}, \text{d})$ l'extension dérivée à droite de
(\ref{sdga-equation-pushforward}). Alors nous avons
$$
RT(\mathcal{M}) = Rf_* R\SheafHom(\mathcal{N}, \mathcal{M})
$$
fonctoriellement en $\mathcal{M}$.
\end{lemma}

\begin{proof}
Par les lemmes \ref{sdga-lemma-tensor-hom-adjunction-dg}
et \ref{sdga-lemma-adjunction-push-pull-dg}, le foncteur
(\ref{sdga-equation-pushforward}) est adjoint à droite
du foncteur (\ref{sdga-equation-pullback}). Par Catégories dérivées, lemme
\ref{derived-lemma-pre-derived-adjoint-functors-general},
le foncteur $RT$ est adjoint à droite au foncteur
du lemme \ref{sdga-lemma-derived-tensor-product}, lequel est égal à
$Lf^*(-) \otimes_\mathcal{A}^\mathbf{L} \mathcal{N}$ par le
lemme \ref{sdga-lemma-compose-pullback-tensor}.
Par les lemmes \ref{sdga-lemma-derived-adjoint-tensor-hom}
et \ref{sdga-lemma-derived-adjoint-push-pull}, le foncteur
$Lf^*(-) \otimes_\mathcal{A}^\mathbf{L} \mathcal{N}$
est adjoint à gauche
à $Rf_* R\SheafHom(\mathcal{N}, -)$.
Nous concluons par unicité des adjoints.
\end{proof}

\begin{lemma}
\label{sdga-lemma-compose-pushforward}
Soient $(f, f^\sharp) : (\Sh(\mathcal{C}), \mathcal{O})
\to (\Sh(\mathcal{C}'), \mathcal{O}')$ et
$(g, g^\sharp) : (\Sh(\mathcal{C}'), \mathcal{O}')
\to (\Sh(\mathcal{C}''), \mathcal{O}'')$ des
morphismes de topos annelés. Soient $\mathcal{A}$, $\mathcal{A}'$ et
$\mathcal{A}''$ des algèbres différentielles graduées respectivement sur $\mathcal{O}$,
$\mathcal{O}'$ et $\mathcal{O}''$. Soient
$\varphi : \mathcal{A}' \to f_*\mathcal{A}$ et
$\varphi' : \mathcal{A}'' \to g_*\mathcal{A}'$ des
homomorphismes d'algèbres différentielles graduées respectivement sur $\mathcal{O}'$
et $\mathcal{O}''$.
Alors nous avons $R(g \circ f)_* = Rg_* \circ Rf_*
: D(\mathcal{A}, \text{d}) \to D(\mathcal{A}'', \text{d})$.
\end{lemma}

\begin{proof}
Cela découle des lemmes \ref{sdga-lemma-compose-pullback} et
\ref{sdga-lemma-derived-adjoint-push-pull} et
de l'unicité des adjoints.
\end{proof}

\begin{lemma}
\label{sdga-lemma-compose-hom} Soit
$(\mathcal{C}, \mathcal{O})$ un site annelé. Soient
$\mathcal{A}$, $\mathcal{A}'$ et $\mathcal{A}''$ des
$\mathcal{O}$-algèbres différentielles graduées. Soient $\mathcal{N}$ et $\mathcal{N}'$ des bimodules
différentiels gradués respectivement sur $(\mathcal{A}, \mathcal{A}')$
et $(\mathcal{A}', \mathcal{A}'')$. Supposons
que l'application canonique
$$
\mathcal{N} \otimes_{\mathcal{A}'}^\mathbf{L} \mathcal{N}'
\longrightarrow
\mathcal{N} \otimes_{\mathcal{A}'} \mathcal{N}'
$$
dans $D(\mathcal{A}'', \text{d})$ soit un quasi-isomorphisme.
Alors nous
avons $$
R\SheafHom_{\mathcal{A}''}
(\mathcal{N} \otimes_{\mathcal{A}'} \mathcal{N}', -)
=
R\SheafHom_{\mathcal{A}'}(\mathcal{N},
R\SheafHom_{\mathcal{A}''}(\mathcal{N}', -))
$$ en
tant que foncteurs $D(\mathcal{A}'', \text{d}) \to D(\mathcal{A}, \text{d})$.
\end{lemma}

\begin{proof}
Cela découle des lemmes \ref{sdga-lemma-compose-tensor} et
\ref{sdga-lemma-derived-adjoint-tensor-hom} et
de l'unicité des adjoints.
\end{proof}

\begin{lemma}
\label{sdga-lemma-pushforward-agrees}
Soit $(f, f^\sharp) : (\Sh(\mathcal{C}), \mathcal{O}_\mathcal{C})
\to (\Sh(\mathcal{D}), \mathcal{O}_\mathcal{D})$ un
morphisme de topos annelés. Soit $\mathcal{A}$ une
$\mathcal{O}_\mathcal{C}$-algèbre différentielle graduée. Soit $\mathcal{B}$ une
$\mathcal{O}_\mathcal{D}$-algèbre différentielle graduée. Soit
$\varphi : \mathcal{B} \to f_*\mathcal{A}$ un homomorphisme de
$\mathcal{O}_\mathcal{D}$-algèbres différentielles graduées. Le
diagramme
$$
\xymatrix{
D(\mathcal{A}, \text{d}) \ar[d]_{Rf_*} \ar[rr]_{forget} & &
D(\mathcal{O}_\mathcal{C}) \ar[d]^{Rf_*} \\
D(\mathcal{B}, \text{d}) \ar[rr]^{forget} & &
D(\mathcal{O}_\mathcal{D})
}
$$
est commutatif.
\end{lemma}

\begin{proof}
Une fois certaines catégories identifiées, ce lemme découle immédiatement
du lemme \ref{sdga-lemma-compose-pushforward}.

\medskip\noindent
Nous pouvons considérer $(\mathcal{O}_\mathcal{C}, 0)$ comme une
$\mathcal{O}_\mathcal{C}$-algèbre différentielle graduée en plaçant $\mathcal{O}_\mathcal{C}$
dans le degré $0$ et en lui conférant le différentiel zéro.
Il est clair
que nous avons $$
\textit{Mod}(\mathcal{O}_\mathcal{C},
0) = \text{Comp}(\mathcal{O}_\mathcal{C})
\quad\text{et}\quad
D(\mathcal{O}_\mathcal{C}, 0) = D(\mathcal{O}_\mathcal{C})
$$
Sous cette identification, le foncteur
d'oubli $\textit{Mod}(\mathcal{A}, \text{d}) \to
\text{Comp}(\mathcal{O}_\mathcal{C})$
est l'« image directe » $\text{id}_{\mathcal{C}, *}$
définie à la section \ref{sdga-section-functoriality-dg},
correspondant au morphisme identité
$\text{id}_\mathcal{C} : (\mathcal{C}, \mathcal{O}_\mathcal{C}) \to
(\mathcal{C}, \mathcal{O}_\mathcal{C})$ de topos annelés et au morphisme
$(\mathcal{O}_\mathcal{C}, 0) \to (\mathcal{A}, \text{d})$ de
$\mathcal{O}_\mathcal{C}$-algèbres différentielles graduées.
Puisque $\text{id}_{\mathcal{C}, *}$ est exact, nous voyons immédiatement
que $$
R\text{id}_{\mathcal{C}, *} = forget
: D(\mathcal{A}, \text{d}) \longrightarrow
D(\mathcal{O}_\mathcal{C}, 0) = D(\mathcal{O}_\mathcal{C})
$$
Le même raisonnement montre que
$$
R\text{id}_{\mathcal{D}, *} = forget :
D(\mathcal{B}, \text{d}) \longrightarrow
D(\mathcal{O}_\mathcal{D}, 0) = D(\mathcal{O}_\mathcal{D})
$$ En outre,
la construction de
$Rf_* : D(\mathcal{O}_\mathcal{C}) \to D(\mathcal{O}_\mathcal{D})$ de Cohomologie
sur les sites, section \ref{sites-cohomology-section-unbounded} concorde avec
la construction de
$Rf_* : D(\mathcal{O}_\mathcal{C}, 0) \to D(\mathcal{O}_\mathcal{D}, 0)$ dans la définition
\ref{sdga-definition-pushforward}, car les deux
foncteurs sont définis comme l'extension dérivée droite de l'image
directe sur les complexes sous-jacents de modules.
Par le lemme \ref{sdga-lemma-compose-pushforward}, nous voyons
que $Rf_* \circ R\text{id}_{\mathcal{C}, *}$
et $R\text{id}_{\mathcal{D}, *} \circ Rf_*$ sont les foncteurs dérivés
de $f_* \circ forget = forget \circ f_*$ et sont donc égaux
par unicité des adjoints.
\end{proof}

\begin{lemma}
\label{sdga-lemma-cohomology-ext}
Soit $(\mathcal{C}, \mathcal{O})$ un site annelé. Soit
$\mathcal{A}$ une $\mathcal{O}$-algèbre différentielle graduée. Soit
$\mathcal{M}$ un $\mathcal{A}$-module différentiel gradué. Soit
$n \in \mathbf{Z}$. Nous avons
$$
H^n(\mathcal{C}, \mathcal{M}) =
\Hom_{D(\mathcal{A}, \text{d})}(\mathcal{A}, \mathcal{M}[n])
$$ où
sur le côté gauche nous avons la cohomologie de $\mathcal{M}$ considérée
comme un complexe de $\mathcal{O}$-modules.
\end{lemma}

\begin{proof}
Pour prouver la formule, observons que
$$
R\Gamma(\mathcal{C}, \mathcal{M}) = \Gamma(\mathcal{C}, \mathcal{I})
$$ où
$\mathcal{M} \to \mathcal{I}$ est un quasi-isomorphisme vers un
module différentiel gradué injectif gradué et K-injectif sur
$\mathcal{A}$, noté $\mathcal{I}$ (en combinant
les lemmes \ref{sdga-lemma-right-derived} et \ref{sdga-lemma-pushforward-agrees}).
Par le lemme \ref{sdga-lemma-hom-derived} nous
avons $$
\Hom_{D(\mathcal{A}, \text{d})}(\mathcal{A}, \mathcal{M}[n])
= \Hom_{K(\textit{Mod}(\mathcal{A}, \text{d}))}(\mathcal{M}, \mathcal{I}[n]) =
H^0(\Gamma(\mathcal{C}, \mathcal{I}[n])) =
H^n(\Gamma(\mathcal{C}, \mathcal{I}))
$$ En
combinant ces deux résultats, nous obtenons notre égalité.
\end{proof}








\section{Équivalences de catégories dérivées}
\label{sdga-section-equivalence}

\noindent
Cette section est l'analogue de
l'Algèbre différentielle graduée, section \ref{dga-section-equivalence}.

\begin{lemma}
\label{sdga-lemma-qis-equivalence}
Soit $(\mathcal{C}, \mathcal{O})$ un site annelé. Si
$\varphi : \mathcal{A} \to \mathcal{B}$ est un homomorphisme de
$\mathcal{O}$-algèbres différentielles graduées qui induit un isomorphisme
sur les faisceaux de cohomologie,
alors $$
D(\mathcal{A}, \text{d}) \longrightarrow D(\mathcal{B}, \text{d}), \quad
\mathcal{M}
\longmapsto
\mathcal{M} \otimes_\mathcal{A}^\mathbf{L} \mathcal{B}
$$
est une équivalence de catégories.
\end{lemma}

\begin{proof}
Rappelons que le foncteur de restriction
$$
\textit{Mod}^{dg}(\mathcal{B}, \text{d}) \to
\textit{Mod}^{dg}(\mathcal{A}, \text{d}),\quad
\mathcal{N} \mapsto res_\varphi \mathcal{N}
$$ est
un adjoint à droite à
$$
\textit{Mod}^{dg}(\mathcal{A}, \text{d}) \to
\textit{Mod}^{dg}(\mathcal{B}, \text{d}),\quad
\mathcal{M} \mapsto \mathcal{M} \otimes_\mathcal{A} \mathcal{B}
$$ Voir la
section \ref{sdga-section-dg-bimodules}. Puisque la restriction
envoie des quasi-isomorphismes sur des quasi-isomorphismes, nous
voyons qu'elle a triviellement une extension dérivée à
gauche (donnée par restriction). Ce foncteur sera adjoint
à droite à $- \otimes_\mathcal{A}^\mathbf{L} \mathcal{B}$
par Catégories dérivées,
lemme \ref{derived-lemma-pre-derived-adjoint-functors-general}.
L'application
de l'adjonction $$
\mathcal{M} \to
res_\varphi(\mathcal{M} \otimes_\mathcal{A}^\mathbf{L} \mathcal{B})
$$
est un isomorphisme dans $D(\mathcal{A}, \text{d})$ par
notre hypothèse que $\mathcal{A} \to \mathcal{B}$ est un
quasi-isomorphisme de $\mathcal{A}$-modules différentiels gradués à gauche.
En particulier, le foncteur du lemme est pleinement
fidèle, voir Catégories, lemme \ref{categories-lemma-adjoint-fully-faithful}.
Il est clair que le noyau du foncteur de restriction
$D(\mathcal{B}, \text{d}) \to D(\mathcal{A}, \text{d})$
est nul. Nous concluons ainsi par Catégories dérivées,
lemme \ref{derived-lemma-fully-faithful-adjoint-kernel-zero}.
\end{proof}













\section{Résolutions d'algèbres différentielles graduées}
\label{sdga-section-resolution-dgas}

\noindent
Cette section est l'analogue de
l'Algèbre différentielle graduée, section \ref{dga-section-resolution-dgas}.

\medskip\noindent
Soit $(\mathcal{C}, \mathcal{O})$ un site annelé. Comme dans la
remarque \ref{sdga-remark-sheaf-graded-sets}, considérons
un faisceau d'ensembles gradués $\mathcal{S}$ sur $\mathcal{C}$.
Considérons le produit de $r$ copies
$\mathcal{S} \times \ldots \times \mathcal{S}$ comme
un faisceau d'ensembles gradués avec la
règle $\deg(s_1 \cdot \ldots \cdot s_r) = \sum \deg(s_i)$.
Pour des sections locales $s_i \in \mathcal{S}(U)$, $i = 1, \ldots, r$, nous
utilisons $s_1 \cdot \ldots \cdot s_r$
pour désigner la section correspondante
de $\mathcal{S} \times \ldots \times \mathcal{S}$
sur $U$. Notons $\mathcal{O}\langle \mathcal{S} \rangle$ la
$\mathcal{O}$-algèbre graduée libre sur $\mathcal{S}$. Plus précisément,
posons
$$
\mathcal{O}\langle \mathcal{S} \rangle =
\mathcal{O} \oplus
\bigoplus\nolimits_{r \geq 1}
\mathcal{O}[\mathcal{S} \times \ldots \times \mathcal{S}]
$$ avec une notation comme dans
la remarque \ref{sdga-remark-sheaf-graded-sets}. Cette construction
devient un faisceau de $\mathcal{O}$-algèbres graduées par
concaténation
$$
(s_1 \cdot \ldots \cdot s_r) (s'_1 \cdot \ldots \cdot s'_{r'}) = s_1
\cdot \ldots s_r \cdot s'_1 \cdot \ldots \cdot s'_{r'}
$$ Nous pouvons
munir $\mathcal{O}\langle \mathcal{S} \rangle$ d'un
différentiel en posant $\text{d}(s) = 0$ pour toute section
locale $s$ de $\mathcal{S}$, puis en prolongeant
de manière unique par la règle de Leibniz ; il
est toutefois important de considérer aussi d'autres différentiels.

\medskip\noindent En
effet, supposons qu'on nous donne un système du genre suivant :
\begin{enumerate}
\item pour $i = 0, 1, 2, \ldots$, des faisceaux d'ensembles gradués $\mathcal{S}_i$,
\item pour $i = 0, 1, 2, \ldots$, des
applications $$
\delta_{i + 1} : \mathcal{S}_{i + 1}
\longrightarrow
\mathcal{A}_i
= \mathcal{O}\langle \mathcal{S}_0 \amalg \ldots \amalg \mathcal{S}_i\rangle
$$
de faisceaux d'ensembles gradués de degré $1$ dont l'image est
contenue dans le noyau du différentiel défini par induction sur la cible.
\end{enumerate} Plus
précisément, posons d'abord
$\mathcal{A}_0 = \mathcal{O}\langle \mathcal{S}_0 \rangle$ et
munissons-la de l'unique différentiel satisfaisant à la règle de Leibniz et tel que
$\text{d}(s) = 0$ pour toute section locale $s$ de $\mathcal{S}$. Par
récurrence, supposons qu'un différentiel $\text{d}$ soit donné
sur $\mathcal{A}_i$. Nous le prolongeons alors en
l'unique différentiel sur $\mathcal{A}_{i + 1}$ satisfaisant à la règle de
Leibniz
et tel que $$
\text{d}(s) = \delta(s)
$$
où $\delta(s) = \delta_j(s)$ si $s$ appartient à la composante $\mathcal{S}_j$
de $\mathcal{S}_0 \amalg \ldots \amalg \mathcal{S}_{i + 1}$.
Cela a du sens précisément parce que $\delta(s)$ est dans le
noyau du différentiel défini par induction.

\begin{lemma}
\label{sdga-lemma-special-good}
Dans la situation ci-dessus, la $\mathcal{O}$-algèbre différentielle graduée
$$
\mathcal{A} = \colim \mathcal{A}_i
$$
possède la propriété suivante : pour tout morphisme
$(f, f^\sharp) : (\Sh(\mathcal{C}'), \mathcal{O}')
\to (\Sh(\mathcal{C}), \mathcal{O})$
de topos annelés, l'image inverse $f^*\mathcal{A}$
est plate comme module gradué sur $\mathcal{O}'$
et K-plate comme complexe de $\mathcal{O}'$-modules.
\end{lemma}

\begin{proof}
Notons que $f^*\mathcal{A} = \colim f^*\mathcal{A}_i$
et que
$$
f^*\mathcal{A}_i = \mathcal{O}'\langle
f^{-1}\mathcal{S}_0 \amalg \ldots \amalg f^{-1}\mathcal{S}_i\rangle
$$
avec un différentiel donné par la procédure inductive ci-dessus en
utilisant $f^{-1}\delta_{i + 1}$. Il suffit donc de prouver que $\mathcal{A}$
est plate comme $\mathcal{O}$-module gradué
et K-plate comme complexe de $\mathcal{O}$-modules.
Il suffit pour cela de montrer que chaque $\mathcal{A}_i$
est plate comme module gradué sur $\mathcal{O}$
et K-plate comme complexe de $\mathcal{O}$-modules ; comparer au
lemme \ref{sdga-lemma-good-direct-sum}.


\medskip\noindent
Pour $i \geq 1$,
écrivons $\mathcal{S} = \mathcal{S}_0 \amalg \ldots \amalg \mathcal{S}_i$,
de sorte que $\mathcal{A}_i = \mathcal{O}\langle \mathcal{S} \rangle$
comme $\mathcal{O}$-algèbre graduée. Nous allons filtrer cette algèbre
par des sous-modules différentiels gradués sur $\mathcal{O}$.

\medskip\noindent
Munissons $W = \mathbf{Z}_{\geq 0}^{i + 1}$ de l'ordre lexicographique.
Autrement dit, étant donnés $w = (w_0, \ldots w_i)$ et
$w' = (w'_0, \ldots, w'_i)$ dans $W$, nous posons
$$
w > w' \Leftrightarrow
\exists j,\ 0 \leq j \leq i :
w_i = w'_i,\ w_{i - 1} = w'_{i - 1},\ \ldots ,
\ w_{j + 1} = w'_{j + 1},\ w_j > w'_j
$$
et ainsi de suite. Supposons donnée une section
$s = s_1 \cdot \ldots \cdot s_r$ de
$\mathcal{S} \times \ldots \times \mathcal{S}$
sur $U$. Nous disons que le {\it poids de $s$ est défini}
si $s_a \in \mathcal{S}_{j_a}(U)$ pour un unique
$0 \leq j_a \leq i$. Dans ce cas, nous définissons le poids
$$
w(s) = (w_0(s), \ldots, w_i(s)) \in W,\quad
w_j(s) = |\{a \mid j_a = j\}|
$$
Le poids de toute section de $\mathcal{S} \times \ldots \times \mathcal{S}$
est défini localement. On vérifie aisément que l'on obtient une décomposition
en union disjointe
$$
\mathcal{S} \times \ldots \times \mathcal{S} =
\coprod\nolimits_{w \in W} \left(
\mathcal{S} \times \ldots \times \mathcal{S}\right)_w
$$
en sous-faisceaux de sections de poids donné. Bien entendu,
seuls les $w \in W$ tels que $\sum_{0 \leq j \leq i} w_j = r$ apparaissent pour un $r$ donné.
Nous obtenons en conséquence une décomposition
$$
\mathcal{A}_i = \mathcal{O} \oplus
\bigoplus\nolimits_{r \geq 1}
\bigoplus\nolimits_{w \in W}
\mathcal{O}[\left(\mathcal{S} \times \ldots \times \mathcal{S}\right)_w]
$$
Le reste de la preuve repose sur l'observation immédiate suivante :
étant donnés $r$, $w$ et une section locale $s = s_1 \cdot \ldots \cdot s_r$ de
$\left(\mathcal{S} \times \ldots \times \mathcal{S}\right)_w$,
nous avons
$$
\text{d}(s) \text{ est une section locale de } \mathcal{O} \oplus
\bigoplus\nolimits_{r' \geq 1}
\bigoplus\nolimits_{w' \in W,\ w' < w}
\mathcal{O}[\left(\mathcal{S} \times \ldots \times \mathcal{S}\right)_{w'}]
$$
En effet, dans chacune des expressions
$$
(-1)^{\deg(s_1) + \ldots + \deg(s_{a - 1})}
s_1 \cdot \ldots s_{a - 1} \cdot \delta(s_a) \cdot
s_{a + 1} \cdot \ldots \cdot s_r
$$
dont la somme donne l'élément $\text{d}(s)$, l'élément $\delta(s_a)$
est localement une combinaison $\mathcal{O}$-linéaire d'éléments
$s'_1 \cdot \ldots \cdot s'_{r'}$, avec $s'_{a'}$ dans
$\mathcal{S}_{j'_a}$ pour un certain $0 \leq j'_{a'} < j_a$, où $j_a$ est tel que
$s_a$ soit une section de $\mathcal{S}_{j_a}$.

\medskip\noindent
Cela signifie ce qui suit. Pour $w \in W$, posons
$$
F_w \mathcal{A}_i = \mathcal{O} \oplus \bigoplus\nolimits_{r \geq 1}
\bigoplus\nolimits_{w' \in W,\ w' \leq w}
\mathcal{O}[\left(\mathcal{S} \times \ldots \times \mathcal{S}\right)_{w'}]
$$
D'après l'observation ci-dessus, c'est un sous-$\mathcal{O}$-module différentiel gradué.
Nous obtenons des suites exactes courtes admissibles
$$
0 \to \colim_{w' < w} F_{w'}\mathcal{A}_i \to
F_w\mathcal{A}_i \to
\bigoplus\nolimits_{r \geq 1}
\mathcal{O}[\left(\mathcal{S} \times \ldots \times \mathcal{S}\right)_w]
\to 0
$$
de $\mathcal{A}$-modules différentiels gradués dont le différentiel
du membre de droite est nul.

\medskip\noindent
Maintenant, nous terminons la preuve par récurrence transfinie sur
l'ensemble ordonné $W$. Le complexe différentiel gradué $F_0\mathcal{A}_0$
est le facteur direct $\mathcal{O}$, qui est K-plat et plat comme module gradué.
Pour $w \in W$, si le résultat est vrai pour $F_{w'}\mathcal{A}_i$
pour $w' < w$, alors par les lemmes \ref{sdga-lemma-good-direct-sum},
\ref{sdga-lemma-good-admissible-ses} et \ref{sdga-lemma-free-graded-module-good}
nous obtenons le résultat pour $w$. Enfin, nous avons
$\mathcal{A}_i = \colim_{w \in W} F_w\mathcal{A}_i$ et
nous concluons.
\end{proof}

\begin{lemma}
\label{sdga-lemma-good-dga} Soit
$(\mathcal{C}, \mathcal{O})$ un site annelé. Soit
$(\mathcal{B}, \text{d})$ une $\mathcal{O}$-algèbre différentielle graduée. Il existe
un quasi-isomorphisme de $\mathcal{O}$-algèbres différentielles graduées
$(\mathcal{A}, \text{d}) \to (\mathcal{B}, \text{d})$ tel que
$\mathcal{A}$ soit plate comme module gradué et K-plate comme complexe de $\mathcal{O}$-modules, et que ces
propriétés restent vraies après image inverse par tout morphisme de
topos annelés.
\end{lemma}

\begin{proof} La
démonstration est identique à la première démonstration du
lemme \ref{sdga-lemma-resolve}, en utilisant cette fois des algèbres
graduées libres au lieu de modules gradués libres.

\medskip\noindent
Nous construirons $\mathcal{A} = \colim \mathcal{A}_i$ comme dans
le lemme \ref{sdga-lemma-special-good}, en construisant
$$
\mathcal{A}_0 \to \mathcal{A}_1 \to \mathcal{A}_2 \to \ldots \to \mathcal{B}
$$
Soit $\mathcal{S}_0$ le faisceau d'ensembles gradués
(remarque \ref{sdga-remark-sheaf-graded-sets})
dont la composante de degré $n$ est $\Ker(\text{d}_\mathcal{B}^n)$.
Considérons l'homomorphisme de
modules différentiels
gradués $$
\mathcal{A}_0 = \mathcal{O}\langle \mathcal{S}_0 \rangle
\longrightarrow
\mathcal{B}
$$
où le morphisme envoie une section locale $s$ de $\mathcal{S}_0$
sur la section locale correspondante de $\mathcal{A}^{\deg(s)}$ ;
celle-ci appartient au noyau du différentiel,
ce qui assure la compatibilité du morphisme avec les différentiels.
Par construction, les applications induites en cohomologie $H^n(\mathcal{A}_0) \to H^n(\mathcal{B})$
sont surjectives et le resteront pour tout $i$.

\medskip\noindent
Étape d'induction de la construction. Étant donné $\mathcal{A}_i \to \mathcal{B}$,
notons $\mathcal{S}_{i + 1}$ le faisceau d'ensembles gradués dont la composante de degré $n$
est $$
\Ker(\text{d}_{\mathcal{A}_i}^{n + 1})
\times_{\mathcal{B}^{n + 1}, \text{d}}
\mathcal{B}^n
$$
Ce faisceau est muni d'une application
canonique $$
\delta_{i + 1} : \mathcal{S}_{i + 1} \longrightarrow
\mathcal{A}_i
$$
dont l'image est contenue dans le noyau de $\text{d}_{\mathcal{A}_i}$
par construction. Par conséquent, $\mathcal{A}_{i + 1} = \mathcal{O}\langle
\mathcal{S}_0 \amalg \ldots \mathcal{S}_{i + 1}\rangle$ a un différentiel
étendant le différentiel sur $\mathcal{A}_i$, voir la discussion au début
de cette section. Le morphisme de $\mathcal{A}_{i + 1}$ vers
$\mathcal{B}$ est l'unique morphisme d'algèbres graduées qui
coïncide avec le morphisme donné sur $\mathcal{A}_i$ et envoie une
section locale $s = (a, b)$ de $\mathcal{S}_{i + 1}$ sur
$b$ dans $\mathcal{B}$. Cela est compatible avec les différentiels
précisément parce que $\text{d}(b)$ est l'image de $a$ dans $\mathcal{B}$.

\medskip\noindent
L'application $\mathcal{A} \to \mathcal{B}$ est un quasi-isomorphisme :
nous avons $H^n(\mathcal{A}) = \colim H^n(\mathcal{A}_i)$
et pour chaque $i$ l'application $H^n(\mathcal{A}_i) \to H^n(\mathcal{B})$
est surjective avec le noyau anéanti par l'application
$H^n(\mathcal{A}_i) \to H^n(\mathcal{A}_{i + 1})$ par construction.
Enfin, la condition de platitude pour $\mathcal{A}$ a été établie dans
le lemme \ref{sdga-lemma-special-good}.
\end{proof}






\section{Compléments}
\label{sdga-section-misc}

\noindent
Soit $(f, f^\sharp) : (\Sh(\mathcal{C}), \mathcal{O})
\to (\Sh(\mathcal{C}'), \mathcal{O}')$ un morphisme de topos
annelés. Soit $\mathcal{A}$ un faisceau de
$\mathcal{O}$-algèbres différentielles graduées. En utilisant la composition\footnote{Il serait
plus précis d'écrire
$F(\mathcal{A}) \otimes_\mathcal{O}^\mathbf{L} F(\mathcal{A})
\to F(\mathcal{A} \otimes_\mathcal{O} \mathcal{A}) \to F(\mathcal{A})$ si
$F$ désigne le foncteur d'oubli pour les complexes de modules sur $\mathcal{O}$. Il
convient également de noter que $\mathcal{A} \otimes_\mathcal{O} \mathcal{A}$
indique le produit tensoriel de la section \ref{sdga-section-tensor-product-dg},
de sorte que $F(\mathcal{A} \otimes_\mathcal{O} \mathcal{A})
= \text{Tot}(F(\mathcal{A}) \otimes_\mathcal{O} F(\mathcal{A}))$.
La première flèche est l'application canonique du produit tensoriel
dérivé de deux complexes de $\mathcal{O}$-modules vers leur produit tensoriel
usuel comme complexes de $\mathcal{O}$-modules.}
$$
\mathcal{A} \otimes_\mathcal{O}^\mathbf{L} \mathcal{A}
\longrightarrow
\mathcal{A} \otimes_\mathcal{O} \mathcal{A}
\longrightarrow
\mathcal{A}
$$
et le cup-produit relatif (voir Cohomologie sur les sites, remarque
\ref{sites-cohomology-remark-cup-product} et
section \ref{sites-cohomology-section-cup-product}),
nous obtenons une multiplication\footnote{Ici et dans
la suite, $Rf_* : D(\mathcal{O}) \to D(\mathcal{O}')$ est le
foncteur dérivé étudié
dans Cohomologie sur les sites, section \ref{sites-cohomology-section-unbounded} et suivantes.}
$$
\mu
: Rf_*\mathcal{A} \otimes_{\mathcal{O}'}^\mathbf{L} Rf_*\mathcal{A}
\longrightarrow
Rf_*\mathcal{A}
$$
dans $D(\mathcal{O}')$. Cette multiplication est
associative dans le
sens où le diagramme $$
\xymatrix{
Rf_*\mathcal{A} \otimes_{\mathcal{O}'}^\mathbf{L} Rf_*\mathcal{A}
\otimes_{\mathcal{O}'}^\mathbf{L} Rf_*\mathcal{A}
\ar[rr]_-{\mu \otimes 1} \ar[d]_{1 \otimes \mu}
& & Rf_*\mathcal{A} \otimes_{\mathcal{O}'}^\mathbf{L} Rf_*\mathcal{A}
\ar[d]^\mu \\
Rf_*\mathcal{A} \otimes_{\mathcal{O}'}^\mathbf{L} Rf_*\mathcal{A}
\ar[rr]^-\mu
& & Rf_*\mathcal{A}
}
$$
est commutatif dans $D(\mathcal{O}')$ ; cela
résulte de Cohomologie sur les sites, lemme \ref{sites-cohomology-lemma-cup-product-associative}.
De même, étant donné
un $\mathcal{A}$-module différentiel gradué à droite $\mathcal{M}$,
nous obtenons une
multiplication $$
\mu_\mathcal{M}
: Rf_*\mathcal{M} \otimes_{\mathcal{O}'}^\mathbf{L} Rf_*\mathcal{A}
\longrightarrow
Rf_*\mathcal{M}
$$
dans $D(\mathcal{O}')$. Cette multiplication est compatible avec $\mu$
ci-dessus au sens où le
diagramme $$
\xymatrix{
Rf_*\mathcal{M} \otimes_{\mathcal{O}'}^\mathbf{L} Rf_*\mathcal{A}
\otimes_{\mathcal{O}'}^\mathbf{L} Rf_*\mathcal{A}
\ar[rr]_-{\mu_\mathcal{M} \otimes 1} \ar[d]_{1 \otimes \mu} & &
Rf_*\mathcal{M} \otimes_{\mathcal{O}'}^\mathbf{L} Rf_*\mathcal{A}
\ar[d]^{\mu_\mathcal{M}} \\
Rf_*\mathcal{M} \otimes_{\mathcal{O}'}^\mathbf{L} Rf_*\mathcal{A}
\ar[rr]^-{\mu_\mathcal{M}} & &
Rf_*\mathcal{M}
}
$$ est
commutatif dans $D(\mathcal{O}')$ ; cela résulte encore de
Cohomologie sur les sites, lemme \ref{sites-cohomology-lemma-cup-product-associative}.

\medskip\noindent
Un cas particulier s'obtient en prenant pour $f$ le morphisme vers le topos
ponctuel $\Sh(pt)$. Dans ce cas, $\mu$ est simplement l'application du
cup-produit
$$
R\Gamma(\mathcal{C}, \mathcal{A})
\otimes_{\Gamma(\mathcal{C}, \mathcal{O})}^\mathbf{L}
R\Gamma(\mathcal{C}, \mathcal{A})
\longrightarrow
R\Gamma(\mathcal{C}, \mathcal{A}),
\quad \eta \otimes \theta \mapsto \eta \cup \theta
$$ et, de même,
$\mu_\mathcal{M}$ est l'application du cup-produit
$$
R\Gamma(\mathcal{C}, \mathcal{M})
\otimes_{\Gamma(\mathcal{C}, \mathcal{O})}^\mathbf{L}
R\Gamma(\mathcal{C}, \mathcal{A})
\longrightarrow
R\Gamma(\mathcal{C}, \mathcal{M}),
\quad \eta \otimes \theta \mapsto \eta \cup \theta
$$ En général,
via l'identification
$$
R\Gamma(\mathcal{C}, \mathcal{A}) =
R\Gamma(\mathcal{C}', Rf_*\mathcal{A})
\quad\text{et}\quad
R\Gamma(\mathcal{C}, \mathcal{M}) =
R\Gamma(\mathcal{C}', Rf_*\mathcal{M})
$$ de Cohomologie sur
les sites, remarque \ref{sites-cohomology-remark-before-Leray}, l'application
$\mu_\mathcal{M}$ induit le cup-produit en cohomologie. Pour le voir,
utilisons
Cohomologie sur les sites, lemme \ref{sites-cohomology-lemma-compose-cup-product},
avec, pour second morphisme de topos, le morphisme de $\Sh(\mathcal{C}')$ vers le
topos ponctuel comme ci-dessus.

\medskip\noindent Si
$\mathcal{M}_1 \to \mathcal{M}_2$ est un homomorphisme de
$\mathcal{A}$-modules différentiels gradués, alors le diagramme
$$
\xymatrix{
Rf_*\mathcal{M}_1 \otimes_{\mathcal{O}'}^\mathbf{L} Rf_*\mathcal{A}
\ar[rr]_-{\mu_{\mathcal{M}_1}} \ar[d] & &
Rf_*\mathcal{M}_1 \ar[d] \\
Rf_*\mathcal{M}_2 \otimes_{\mathcal{O}'}^\mathbf{L} Rf_*\mathcal{A}
\ar[rr]^-{\mu_{\mathcal{M}_2}} & &
Rf_*\mathcal{M}_2
}
$$ est commutatif dans
$D(\mathcal{O}')$ ; cela résulte de la fonctorialité du cup-produit relatif. Supposons
que nous ayons une suite exacte courte
$$ 0
\to \mathcal{M}_1 \xrightarrow{a} \mathcal{M}_2 \to \mathcal{M}_3 \to 0
$$ de modules différentiels gradués à
droite sur $\mathcal{A}$. Alors, nous affirmons que le
diagramme
$$
\xymatrix{
Rf_*\mathcal{M}_3 \otimes_{\mathcal{O}'}^\mathbf{L} Rf_*\mathcal{A}
\ar[rr]_-{\mu_{\mathcal{M}_3}} \ar[d]_{Rf_*\delta \otimes \text{id}} & &
Rf_*\mathcal{M}_3 \ar[d]^{Rf_*\delta} \\
Rf_*\mathcal{M}_1[1] \otimes_{\mathcal{O}'}^\mathbf{L} Rf_*\mathcal{A}
\ar[rr]^-{\mu_{\mathcal{M}_1[1]}} & &
Rf_*\mathcal{M}_1[1]
}
$$ est commutatif dans
$D(\mathcal{O}')$, où
$\delta : \mathcal{M}_3 \to \mathcal{M}_1[1]$ est le morphisme de
$D(\mathcal{O})$ provenant de la suite exacte courte donnée (voir Catégories
dérivées, section
\ref{derived-section-canonical-delta-functor}). Cela est clair si la suite est scindée comme
suite de modules gradués à droite sur
$\mathcal{A}$ : dans ce cas, $\delta$ peut être représenté par un homomorphisme de modules à droite sur
$\mathcal{A}$ et la discussion ci-dessus s'applique. En général, nous argumentons en utilisant
le cône sur $a$ et le diagramme
$$
\xymatrix{
\mathcal{M}_1 \ar[r]_a \ar[d] &
\mathcal{M}_2 \ar[r]_i \ar[d] & C(a)
\ar[r]_{-p} \ar[d]^q &
\mathcal{M}_1[1] \ar[d] \\
\mathcal{M}_1 \ar[r] &
\mathcal{M}_2 \ar[r] &
\mathcal{M}_3 \ar[r]^\delta &
\mathcal{M}_1[1]
}
$$ où le carré de droite est commutatif
dans $D(\mathcal{O})$, par la définition de
$\delta$ dans Catégories dérivées, lemme
\ref{derived-lemma-derived-canonical-delta-functor}. Le cône
$C(a)$ est muni d'une structure de
$\mathcal{A}$-module différentiel gradué à droite telle que $i$, $p$, $q$ soient des homomorphismes de
$\mathcal{A}$-modules différentiels gradués à droite ; voir la définition
\ref{sdga-definition-cone}. Nous savons donc, par ce qui précède,
que les diagrammes correspondants sont commutatifs
pour les morphismes $q$ et $-p$. Puisque
$q$ est un isomorphisme dans $D(\mathcal{O})$, nous concluons que la même chose
est vraie pour $\delta$ comme souhaité.

\medskip\noindent
Dans la situation ci-dessus, étant donné un $\mathcal{A}$-module différentiel gradué
à droite $\mathcal{M}$,
soit $$
\xi \in H^n(\mathcal{C}, \mathcal{M})
$$
Autrement dit, $\xi$ est une classe de cohomologie de degré $n$ de
$\mathcal{M}$, considéré comme un complexe de $\mathcal{O}$-modules.
Par le lemme \ref{sdga-lemma-cohomology-ext}, nous pouvons
construire des morphismes $$
x : \mathcal{A} \rightarrow \mathcal{M}'[n]
\quad\text{et}\quad
s : \mathcal{M} \to \mathcal{M}'
$$
de $\mathcal{A}$-modules différentiels gradués où $s$
est un quasi-isomorphisme et tels que $\xi$ soit l'image
de $1 \in H^0(\mathcal{C}, \mathcal{A})$ par le morphisme
$s[n]^{-1} \circ x$ dans la catégorie dérivée
$D(\mathcal{A}, \text{d})$ et a fortiori dans la catégorie
dérivée $D(\mathcal{O})$. Il s'ensuit que l'application correspondante
$$
\xi' = (s[n])^{-1} \circ x : \mathcal{A} \longrightarrow \mathcal{M}[n]
$$
dans $D(\mathcal{O})$ est uniquement caractérisée par les deux
propriétés suivantes :
\begin{enumerate}
\item $\xi'$ se relève en un morphisme dans $D(\mathcal{A}, \text{d})$, et
\item $\xi = \xi'(1)$ dans
$H^0(\mathcal{C}, \mathcal{M}[n]) = H^n(\mathcal{C}, \mathcal{M})$.
\end{enumerate} En
utilisant les compatibilités de $x$ et $s$ avec le cup-produit relatif
considéré ci-dessus, il s'ensuit que, pour tout\footnote{Par exemple, le
morphisme identité.} morphisme de topos annelés
$(f, f^\sharp) : (\Sh(\mathcal{C}), \mathcal{O})
\to (\Sh(\mathcal{C}'), \mathcal{O}')$, l'image directe dérivée
$$
Rf_*\xi' : Rf_*\mathcal{A} \longrightarrow Rf_*\mathcal{M}[n]
$$ de
$\xi'$ est compatible avec les applications $\mu$ et $\mu_{\mathcal{M}[n]}$ construites
ci-dessus dans le sens où le diagramme
$$
\xymatrix{
Rf_*\mathcal{A} \otimes_{\mathcal{O}'}^\mathbf{L} Rf_*\mathcal{A}
\ar[rr]_-\mu \ar[d]_{Rf_*\xi' \otimes \text{id}} & &
Rf_*\mathcal{A} \ar[d]^{Rf_*\xi'} \\
Rf_*\mathcal{M}[n] \otimes_{\mathcal{O}'}^\mathbf{L} Rf_*\mathcal{A}
\ar[rr]^-{\mu_{\mathcal{M}[n]}} & &
Rf_*\mathcal{M}[n]
}
$$ est commutatif dans
$D(\mathcal{O}')$. En utilisant cette compatibilité pour l'application
aux topos ponctuels, nous voyons en particulier que
$$
\xymatrix{
R\Gamma(\mathcal{C}, \mathcal{A})
\otimes_{\Gamma(\mathcal{C}, \mathcal{O})}^\mathbf{L}
R\Gamma(\mathcal{C}, \mathcal{A})
\ar[d]_{\xi' \otimes \text{id}} \ar[r] &
R\Gamma(\mathcal{C}, \mathcal{A}) \ar[d]^{\xi'} \\
R\Gamma(\mathcal{C}, \mathcal{M}[n])
\otimes_{\Gamma(\mathcal{C}, \mathcal{O})}^\mathbf{L}
R\Gamma(\mathcal{C}, \mathcal{A})
\ar[r] &
R\Gamma(\mathcal{C}, \mathcal{M}[n])
}
$$ est commutatif.
Combiné à $\xi'(1) = \xi$, cela implique que l'application
induite sur la cohomologie
$$
\xi' : R\Gamma(\mathcal{C}, \mathcal{A}) \to
R\Gamma(\mathcal{C}, \mathcal{M}[n]), \quad \eta \mapsto \xi \cup \eta
$$ est donnée
par le cup-produit à gauche par $\xi$, comme indiqué.














\section{Modules différentiels gradués sur une catégorie}
\label{sdga-section-modules-cohomology}

\noindent
Cette section prolonge Cohomologie
sur les sites, section \ref{sites-cohomology-section-modules-cohomology}.

\medskip\noindent
Soit $\mathcal{C}$ une catégorie. Nous pensons à $\mathcal{C}$ comme un site avec
la topologie chaotique. Soit $\mathcal{O}$ un faisceau d'anneaux sur $\mathcal{C}$.
Soit $(\mathcal{A}, \text{d})$ un faisceau de
$\mathcal{O}$-algèbres différentielles graduées. En d'autres termes, $\mathcal{O}$ est
un préfaisceau d'anneaux de la catégorie $\mathcal{C}$ et
$(\mathcal{A}, \text{d})$ est un préfaisceau de
$\mathcal{O}$-algèbres différentielles graduées sur $\mathcal{C}$,
voir Catégories, définition \ref{categories-definition-presheaf}.

\begin{definition}
\label{sdga-definition-cartesian} Dans
la situation ci-dessus, nous notons
{\it $\mathit{QC}(\mathcal{A}, \text{d})$} la
sous-catégorie pleine de $D(\mathcal{A}, \text{d})$
formée des objets $M$ tels que, pour tout
morphisme $U \to V$ de $\mathcal{C}$, l'application
canonique $$
R\Gamma(V, M) \otimes_{\mathcal{A}(V)}^\mathbf{L} \mathcal{A}(U)
\longrightarrow
R\Gamma(U, M)
$$
est un isomorphisme dans $D(\mathcal{A}(U), \text{d})$.
\end{definition}

\begin{lemma}
\label{sdga-lemma-cartesian-good-sub} Dans
la situation ci-dessus, la sous-catégorie $\mathit{QC}(\mathcal{A}, \text{d})$ est
une sous-catégorie strictement pleine, saturée et triangulée de
$D(\mathcal{A}, \text{d})$, stable par sommes directes arbitraires.
\end{lemma}

\begin{proof}
Soit $U$ un objet de $\mathcal{C}$. Puisque la topologie sur
$\mathcal{C}$ est chaotique, le foncteur $\mathcal{F} \mapsto \mathcal{F}(U)$ est
exact et commute aux sommes directes. Par conséquent, le foncteur
exact $M \mapsto R\Gamma(U, M)$ est calculé en représentant $K$ par
un $\mathcal{A}$-module différentiel
gradué quelconque $\mathcal{M}$ et en prenant $\mathcal{M}(U)$.
Ainsi, $R\Gamma(U, -)$ commute aux sommes
directes, voir le lemme \ref{sdga-lemma-derived-products}.
De même, pour un morphisme $U \to V$ de $\mathcal{C}$,
le foncteur de produit
tensoriel dérivé $- \otimes_{\mathcal{O}(A)}^\mathbf{L} \mathcal{A}(U)
: D(\mathcal{A}(V)) \to D(\mathcal{A}(U))$
est exact et commute aux sommes directes. Le
lemme découle directement de ces observations ; les détails sont omis.
\end{proof}

\begin{remark}
\label{sdga-remark-cartesian-brown}
Comme ci-dessus, soient $\mathcal{C}$ une catégorie considérée comme
un site muni de la topologie chaotique, $\mathcal{O}$ un faisceau d'anneaux sur $\mathcal{C}$
et $(\mathcal{A}, \text{d})$ un faisceau de
$\mathcal{O}$-algèbres différentielles graduées. Alors l'analogue
de Cohomologie sur les sites,
proposition \ref{sites-cohomology-proposition-cartesian-brown},
vaut pour $\mathit{QC}(\mathcal{A}, \text{d})$ avec
presque exactement la même
démonstration : \begin{enumerate}
\item tout foncteur cohomologique
contravariant $H : \mathit{QC}(\mathcal{A}, \text{d}) \to \textit{Ab}$
qui transforme les sommes directes en produits est
représentable, \item tout
foncteur exact $F : \mathit{QC}(\mathcal{A}, \text{d}) \to \mathcal{D}$
de catégories triangulées qui transforme les sommes directes en
sommes directes possède un adjoint
à droite exact, et \item le
foncteur d'inclusion $\mathit{QC}(\mathcal{A}, \text{d}) \to D(\mathcal{A}, \text{d})$
possède un
adjoint à droite exact. \end{enumerate}
Si nous en avons besoin, nous le démontrerons ici.
\end{remark}

\noindent
Soit $u : \mathcal{C}' \to \mathcal{C}$ un foncteur entre catégories. Si nous
considérons $\mathcal{C}$ et $\mathcal{C}'$ comme des sites avec la
topologie chaotique, alors $u$ est un foncteur continu et cocontinu. On
obtient donc un morphisme $g : \Sh(\mathcal{C}') \to \Sh(\mathcal{C})$ de
topos, voir Sites, lemme \ref{sites-lemma-cocontinuous-morphism-topoi}. En
outre, supposons donnés des faisceaux d'anneaux $\mathcal{O}$ sur $\mathcal{C}$ et
$\mathcal{O}'$ sur $\mathcal{C}'$ et une application
$g^\sharp : g^{-1}\mathcal{O} \to \mathcal{O}'$.
Notons simplement le morphisme correspondant de topos annelés
$g : (\Sh(\mathcal{C}'), \mathcal{O}') \to (\Sh(\mathcal{C}), \mathcal{O})$, voir
Modules sur les sites, section \ref{sites-modules-section-ringed-topoi}.
Enfin, supposons que $(\mathcal{A}, \text{d})$ soit un faisceau de
$\mathcal{O}$-algèbres différentielles graduées, que
$(\mathcal{A}', \text{d})$ soit un faisceau de
$\mathcal{O}'$-algèbres différentielles graduées, et que nous ayons un morphisme
$\varphi : g^*\mathcal{A} \to \mathcal{A}'$ d'algèbres différentielles
graduées sur $\mathcal{O}'$ (voir la section \ref{sdga-section-functoriality-dg}).

\begin{lemma}
\label{sdga-lemma-cartesian-functorial}
Soient $g : (\Sh(\mathcal{C}'), \mathcal{O}') \to (\Sh(\mathcal{C}), \mathcal{O})$
et $\varphi : g^*\mathcal{A} \to \mathcal{A}'$ comme
ci-dessus. Le foncteur
$Lg^* : D(\mathcal{A}, \text{d}) \to D(\mathcal{A}', \text{d})$ envoie
alors $\mathit{QC}(\mathcal{A}, \text{d})$ dans
$\mathit{QC}(\mathcal{A}', \text{d})$.
\end{lemma}

\begin{proof}
Soit $U' \in \Ob(\mathcal{C}')$, d'image $U = u(U')$ dans $\mathcal{C}$.
Notons $pt$ la catégorie ayant un seul objet et un seul morphisme. Considérons
les topos annelés $(\Sh(pt), \mathcal{O}'(U'))$ et $(\Sh(pt), \mathcal{O}(U))$,
munis des algèbres différentielles graduées
$\mathcal{A}'(U)$ et $\mathcal{A}(U)$. Bien sûr, nous identifions
la catégorie dérivée des modules différentiels gradués sur ceux-ci avec
$D(\mathcal{A}'(U'), \text{d})$ et
$D(\mathcal{A}(U), \text{d})$. Ensuite, nous
avons un diagramme commutatif de topos annelés
$$
\xymatrix{
(\Sh(pt), \mathcal{O}'(U')) \ar[rr]_{U'} \ar[d] & &
(\Sh(\mathcal{C}'), \mathcal{O}') \ar[d]^g \\
(\Sh(pt), \mathcal{O}(U)) \ar[rr]^U & &
(\Sh(\mathcal{C}), \mathcal{O})
}
$$ Chacun
est doté d'algèbres différentielles graduées correspondantes.
L'image inverse par le morphisme horizontal inférieur envoie $M$
dans $D(\mathcal{A}, \text{d})$ sur $R\Gamma(U, K)$,
considéré comme un objet de $D(\mathcal{A}(U), \text{d})$.
L'image inverse par la flèche verticale gauche envoie $M$
sur $M \otimes_{\mathcal{A}(U)}^\mathbf{L} \mathcal{A}'(U')$.
Les deux parcours du diagramme donnent le
même résultat (lemme \ref{sdga-lemma-compose-pullback}), d'où
$$
R\Gamma(U', Lg^*K) =
R\Gamma(U, K) \otimes_{\mathcal{A}(U)}^\mathbf{L} \mathcal{A}'(U')
$$
Enfin, soit $f' : U' \to V'$ un morphisme de $\mathcal{C}'$, et notons
$f = u(f') : U = u(U') \to V = u(V')$ son image dans $\mathcal{C}$. Si
$K$ est dans $\mathit{QC}(\mathcal{A}, \text{d})$ alors nous avons
\begin{align*}
R\Gamma(V', Lg^*K) \otimes_{\mathcal{A}'(V')}^\mathbf{L} \mathcal{A}'(U') &
=
R\Gamma(V, K) \otimes_{\mathcal{A}(V)}^\mathbf{L} \mathcal{A}'(V')
\otimes_{\mathcal{A}'(V')}^\mathbf{L} \mathcal{A}'(U') \\ &
=
R\Gamma(V, K) \otimes_{\mathcal{A}(V)}^\mathbf{L} \mathcal{A}'(U') \\ &
=
R\Gamma(V, K) \otimes_{\mathcal{A}(V)}^\mathbf{L} \mathcal{A}(U)
\otimes_{\mathcal{A}(U)}^\mathbf{L} \mathcal{A}'(U') \\ &
=
R\Gamma(U, K) \otimes_{\mathcal{A}(U)}^\mathbf{L} \mathcal{A}'(U') \\ &
=
R\Gamma(U', Lg^*K)
\end{align*} comme
souhaité. Ici, nous avons utilisé l'observation ci-dessus pour $U'$ et $V'$.
\end{proof}






\section{Modules différentiels gradués sur une catégorie, bis}
\label{sdga-section-modules-cohomology-bis}

\noindent
Nous développons quelques résultats supplémentaires sur la notion de modules
quasi cohérents introduite dans la section \ref{sdga-section-modules-cohomology}.

\begin{lemma}
\label{sdga-lemma-cartesian-cofinal-subcategory}
Soient $\mathcal{C}, \mathcal{O}, \mathcal{A}$ comme
dans la section \ref{sdga-section-modules-cohomology}.
Soit $\mathcal{C}' \subset \mathcal{C}$ une sous-catégorie
pleine telle que, pour tout $U \in \Ob(\mathcal{C})$, la catégorie
$U/\mathcal{C}'$ des flèches $U \to U'$ soit cofiltrée. Notons
$\mathcal{O}', \mathcal{A}'$ les restrictions de
$\mathcal{O}, \mathcal{A}$ à $\mathcal{C}'$. Les
restrictions induisent alors une équivalence
$\mathit{QC}(\mathcal{A}, \text{d}) \to \mathit{QC}(\mathcal{A}', \text{d})$.
\end{lemma}

\begin{proof}
Construisons un quasi-inverse de ce foncteur. Soit
$M'$ un objet de $\mathit{QC}(\mathcal{A}', \text{d})$. Nous pouvons représenter
$M'$ par un bon module différentiel gradué $\mathcal{M}'$, voir lemme
\ref{sdga-lemma-resolve}. Ensuite, pour chaque $U' \in \Ob(\mathcal{C}')$, le
$\mathcal{A}'(U')$-module différentiel gradué $\mathcal{M}'(U)$ est K-plat et plat comme module
gradué et, pour chaque morphisme $U'_1 \to U'_2$ de
$\mathcal{C}'$, l'application
$$
\mathcal{M}'(U'_2) \otimes_{\mathcal{A}'(U'_2)} \mathcal{A}'(U'_1)
\longrightarrow
\mathcal{M}'(U'_1)
$$ est un
quasi-isomorphisme (puisque la source représente le produit tensoriel dérivé). Considérons
le
$\mathcal{A}$-module différentiel gradué $\mathcal{M}$ défini par
$$
\mathcal{M}(U) =
\colim_{U \to U' \in U/\mathcal{C}'}
\mathcal{M}'(U') \otimes_{\mathcal{A}'(U')} \mathcal{A}(U)
$$ Il s'agit d'une limite inductive
filtrante de complexes d'après l'hypothèse du lemme. Puisque
$M'$ est dans $\mathit{QC}(\mathcal{A}', \text{d})$ tous les morphismes de transition
dans le système sont des quasi-isomorphismes. Puisque les limites inductives
filtrantes sont exactes, nous voyons
que $\mathcal{M}(U)$ dans $D(\mathcal{A}(U), \text{d})$
est isomorphe à $\mathcal{M}'(U') \otimes_{\mathcal{A}'(U')} \mathcal{A}(U)$
pour tout morphisme $U \to U'$ avec $U' \in \Ob(\mathcal{C}')$.

\medskip\noindent
Nous affirmons que $\mathcal{M}$ est dans
$\mathit{QC}(\mathcal{A}, \text{d})$ : à savoir, étant donné $U \to V$
dans $\mathcal{C}$, nous choisissons une application $V \to V'$ avec $V' \in \Ob(\mathcal{C}')$.
Par ce qui précède, nous voyons que l'application $\mathcal{M}(V) \to \mathcal{M}(U)$
est
identifiée avec l'application $$
\mathcal{M}'(V') \otimes_{\mathcal{A}'(V')} \mathcal{A}(V)
\longrightarrow
\mathcal{M}'(V') \otimes_{\mathcal{A}'(V')} \mathcal{A}(U)
$$
Puisque $\mathcal{M'}(V')$ est
K-plat comme $\mathcal{A}'(V')$-module différentiel gradué, nous
obtenons l'affirmation.

\medskip\noindent
Le morphisme naturel $\mathcal{M}|_{\mathcal{C}'} \to \mathcal{M}'$
est un isomorphisme dans $D(\mathcal{A}', d)$, d'après ce qui précède.

\medskip\noindent À
l'inverse, si nous avons un objet $E$ de $\mathit{QC}(\mathcal{A}, \text{d})$, alors
nous le représentons par un bon module différentiel gradué $\mathcal{E}$. En
définissant $\mathcal{M}' = \mathcal{E}|_{\mathcal{C}'}$ (c'est un autre bon module
différentiel gradué) nous voyons qu'il y a une application
$$
\mathcal{E} \to \mathcal{M}
$$ qui,
sur $U$ dans $\mathcal{C}$, est donnée par le morphisme
$$
\mathcal{E}(U)
\longrightarrow
\colim_{U \to U' \in U/\mathcal{C}'}
\mathcal{E}(U') \otimes_{\mathcal{A}'(U')} \mathcal{A}(U)
$$ qui est
un quasi-isomorphisme pour la même raison. Ainsi, la restriction et la
construction ci-dessus sont des foncteurs quasi inverses, comme souhaité.
\end{proof}

\begin{lemma}
\label{sdga-lemma-cartesian-qis-equivalence}
Soient $\mathcal{C}, \mathcal{O}$ comme dans
la section \ref{sdga-section-modules-cohomology}. Si
$\varphi : \mathcal{A} \to \mathcal{B}$ est un homomorphisme de
$\mathcal{O}$-algèbres différentielles graduées qui induit un
isomorphisme sur les faisceaux de cohomologie, alors
l'équivalence $D(\mathcal{A}, \text{d}) \to D(\mathcal{B}, \text{d})$
du lemme \ref{sdga-lemma-qis-equivalence} induit une équivalence
$\mathit{QC}(\mathcal{A}, \text{d}) \to
\mathit{QC}(\mathcal{B}, \text{d})$.
\end{lemma}

\begin{proof}
Il suffit de montrer que, pour un morphisme $U \to V$ de $\mathcal{C}$ et un objet
$M$ de $D(\mathcal{A}, \text{d})$, les conditions suivantes sont équivalentes :
\begin{enumerate}
\item $R\Gamma(V, M) \otimes_{\mathcal{A}(V)}^\mathbf{L} \mathcal{A}(U)
\to \Gamma(U, M)$ est un isomorphisme dans $D(\mathcal{A}(U), \text{d})$,
et \item $R\Gamma(V, M \otimes_\mathcal{A}^\mathbf{L} \mathcal{B})
\otimes_{\mathcal{B}(V)}^\mathbf{L} \mathcal{B}(U)
\to \Gamma(U, M \otimes_\mathcal{A}^\mathbf{L} \mathcal{B})$ est un
isomorphisme dans $D(\mathcal{B}(U), \text{d})$.
\end{enumerate} Comme
la topologie sur $\mathcal{C}$ est chaotique, cela se résume simplement au
fait que $\mathcal{A}(U) \to \mathcal{B}(U)$ et
$\mathcal{A}(V) \to \mathcal{B}(V)$ sont des quasi-isomorphismes. Les
détails sont omis.
\end{proof}






\section{Systèmes projectifs d'algèbres différentielles graduées}
\label{sdga-section-qc-inverse-system}

\noindent
Dans cette section, nous examinons le cas particulier suivant de la situation
décrite dans la section \ref{sdga-section-modules-cohomology} :
\begin{enumerate}
\item $\mathcal{C}$ est la catégorie $\mathbf{N}$ avec
un morphisme unique $i \to j$ si et seulement si $i \leq j$,
\item $\mathcal{O}$ est le (pré)faisceau constant d'anneaux de valeur l'anneau
donné $R$.
\end{enumerate} Dans
ce cadre, un faisceau $\mathcal{A}$ de
$\mathcal{O}$-algèbres différentielles graduées est la même chose qu'un système projectif
$(A_n)$ de $R$-algèbres différentielles graduées. Un
faisceau $\mathcal{M}$ de $\mathcal{A}$-modules différentiels gradués est
la même chose qu'un système projectif $(M_n)$ où $M_n$ est un
$A_n$-module différentiel gradué et les morphismes de transition
$M_{n + 1} \to M_n$ sont des morphismes de $A_{n + 1}$-modules.
Dans cette section, nous omettons le différentiel de la
notation et écrivons $\mathit{QC}(\mathcal{A})$ pour la catégorie
notée $\mathit{QC}(\mathcal{A}, \text{d})$ dans
la définition \ref{sdga-definition-cartesian}.

\medskip\noindent
Supposons que $\mathcal{B} = (B_n)$ soit un deuxième système projectif de
$R$-algèbres différentielles graduées. Étant donné un morphisme $\varphi : (A_n) \to (B_n)$
de pro-objets, construisons un foncteur
exact de $\mathit{QC}(\mathcal{A})$ vers $\mathit{QC}(\mathcal{B})$.
D'après Catégories, exemple
\ref{categories-example-pro-morphism-inverse-systems}, le
morphisme $\varphi$ est donné par une suite d'entiers
$\ldots \geq m(3) \geq m(2) \geq m(1)$ et un
diagramme commutatif
$$
\xymatrix{
\ldots \ar[r] &
A_{m(3)} \ar[d]^{\varphi_3} \ar[r] &
A_{m(2)} \ar[d]^{\varphi_2} \ar[r] &
A_{m(1)} \ar[d]^{\varphi_1} \\
\ldots \ar[r] & B_3
\ar[r] & B_2
\ar[r] &
B_1
}
$$ de
$R$-algèbres différentielles graduées. Étant donné un bon faisceau de
$\mathcal{A}$-modules différentiels gradués $\mathcal{M} = (M_n)$
représentant un objet de $\mathit{QC}(\mathcal{A})$,
posons
$$ N_n
= M_{m(n)} \otimes_{A_{m(n)}} B_n
$$ Ce
système projectif détermine un objet de
$\mathit{QC}(\mathcal{B})$ parce que les
$A_{m(n)}$-modules $M_{m(n)}$ sont K-plats ; détails omis. Nous
laissons également au lecteur le soin de montrer que le foncteur
résultant est indépendant des choix faits dans sa construction.

\begin{lemma}
\label{sdga-lemma-pro-isomorphic} Dans
la situation ci-dessus, supposons que $\mathcal{A} = (A_n)$ et
$\mathcal{B} = (B_n)$ soient des systèmes projectifs de
$R$-algèbres différentielles graduées. Si $\varphi : (A_n) \to (B_n)$ est un isomorphisme
de pro-objets, alors le foncteur $\mathit{QC}(\mathcal{A}) \to
\mathit{QC}(\mathcal{B})$ construit ci-dessus est
une équivalence.
\end{lemma}

\begin{proof}
Soit $\psi : (B_n) \to (A_n)$ un morphisme de pro-objets inverse de
$\varphi$. D'après la discussion de Catégories,
exemple
\ref{categories-example-pro-morphism-inverse-systems}, nous pouvons supposer que
$\varphi$ est donné par un système de morphismes comme ci-dessus,
tandis que $\psi$ est donné par une suite $n(1) < n(2) < \ldots$
et un diagramme
commutatif $$
\xymatrix{
\ldots \ar[r]
& B_{n(3)} \ar[d]^{\psi_3} \ar[r]
& B_{n(2)} \ar[d]^{\psi_2} \ar[r]
& B_{n(1)} \ar[d]^{\psi_1} \\
\ldots \ar[r] &
A_3 \ar[r] &
A_2 \ar[r] &
A_1
}
$$
de $R$-algèbres différentielles graduées.
Puisque $\varphi \circ \psi = \text{id}$, nous pouvons, après avoir
éventuellement augmenté les valeurs des fonctions $n(\cdot)$
et $m(\cdot)$, supposer que $B_{n(m(i))} \to A_{m(i)} \to B_i$
est l'application de transition. Il s'ensuit que la composition
des
foncteurs $$
\mathit{QC}(\mathcal{B}) \to
\mathit{QC}(\mathcal{A}) \to
\mathit{QC}(\mathcal{B})
$$
envoie un bon faisceau
de $\mathcal{B}$-modules différentiels gradués $\mathcal{N} = (N_n)$
au système projectif $\mathcal{N}' = (N'_i)$ à valeurs
$$
N'_i = N_{n(m(i))} \otimes_{B_{n(m(i))}} B_i
$$ qui
est canoniquement quasi-isomorphe à $\mathcal{N}$ précisément
parce que $\mathcal{N}$ est un objet de
$\mathit{QC}(\mathcal{B})$ et que
$N_j$ est un module différentiel gradué K-plat pour
tout $j$. Le même argument s'applique à la composition
dans l'autre sens, ce qui conclut la démonstration.
\end{proof}

\noindent
Soient $\mathcal{C} = \mathbf{N}$ et $\mathcal{O}$ le faisceau constant de
valeur un anneau $R$ ; soit $\mathcal{A}$ donné par un système projectif
$(A_n)$ de $R$-algèbres différentielles graduées. Supposons donnés deux
modules différentiels gradués à gauche sur $\mathcal{A}$, $\mathcal{N}$
et $\mathcal{N}'$, par des systèmes projectifs $(N_n)$ et $(N'_n)$.
Ainsi, chaque $N_n$ et $N'_n$ est un module différentiel gradué à gauche sur $A_n$.
Disons provisoirement que $(N_n)$ et $(N'_n)$ sont
{\it pro-isomorphes dans la catégorie dérivée}
s'il existe une suite d'entiers
$$ 1 = n_0
< n_1 < n_2 < n_3 < \ldots
$$ et des
morphismes
$$
N_{n_{2i}} \to N'_{n_{2i - 1}}
\quad\text{dans}\quad
D(A_{n_{2i}}^{opp}, \text{d})
$$
et
$$
N'_{n_{2i + 1}} \to N'_{n_{2i}}
\quad\text{dans}\quad
D(A_{n_{2i + 1}}^{opp}, \text{d})
$$ de
sorte que les compositions $N_{n_{2i}} \to N_{n_{2i - 2}}$ et
$N'_{n_{2i + 1}} \to N'_{2i - 1}$ sont données par les morphismes de
transition des systèmes respectifs.

\begin{lemma}
\label{sdga-lemma-pro-isomorphic-systems-tensor} Si
$(N_n)$ et $(N'_n)$ sont pro-isomorphes dans la catégorie dérivée au sens
défini ci-dessus,
alors pour chaque objet $(M_n)$ de $D(\mathbf{N}, \mathcal{A})$ nous avons
$$
R\lim (M_n \otimes_{A_n}^\mathbf{L} N_n) =
R\lim (M_n \otimes_{A_n}^\mathbf{L} N'_n)
$$
dans $D(R)$.
\end{lemma}

\begin{proof}
L'hypothèse implique que le système projectif
$(M_n \otimes_{A_n}^\mathbf{L} N_n)$ de $D(R)$ est pro-isomorphe (au sens
habituel) au système projectif
$(M_n \otimes_{A_n}^\mathbf{L} N'_n)$ de $D(R)$. Le résultat découle
donc de Compléments
sur l'algèbre, lemme \ref{more-algebra-lemma-pro-isomorphism-Rlim}.
\end{proof}

\begin{lemma}
\label{sdga-lemma-different-tensors}
\begin{reference} Il
s'agit d'une variante de \cite[Lemma 3.5.4]{BS}.
\end{reference}
Soit $R$ un anneau. Soient $f_1, \ldots, f_r \in R$.
Soit $K_n$ le complexe de Koszul sur $f_1^n, \ldots, f_r^n$,
considéré comme une $R$-algèbre différentielle
graduée. Soit $(M_n)$ un objet de $D(\mathbf{N}, (K_n))$.
Alors, pour tout $t \geq 1$, on
a $$
R\lim (M_n \otimes_R^\mathbf{L} K_t) =
R\lim (M_n \otimes_{K_n}^\mathbf{L} K_t)
$$
dans $D(R)$.
\end{lemma}

\begin{proof}
Fixons $t \geq 1$. Pour $n \geq t$, notons ${}_nK_t$
la $R$-algèbre différentielle graduée $K_t$ considérée comme un module différentiel gradué
à gauche sur $K_n$. Observons que
$$
M_n \otimes_R^\mathbf{L} K_t =
M_n \otimes_{K_n}^\mathbf{L} (K_n \otimes_R^\mathbf{L} K_t) =
M_n \otimes_{K_n}^\mathbf{L} (K_n \otimes_R K_t)
$$
Ainsi, d'après le lemme \ref{sdga-lemma-pro-isomorphic-systems-tensor},
il suffit de montrer que
$({}_nK_t)$ et $(K_n \otimes_R K_t)$ sont pro-isomorphes dans la
catégorie dérivée. Les applications de multiplication
$$
K_n \otimes_R K_t \longrightarrow {}_nK_t
$$
sont des morphismes de modules différentiels gradués à gauche sur $K_n$. Ainsi, pour conclure,
il suffit de montrer que, pour tout $n \geq 1$, il existe $N > n$
et un morphisme
$$
{}_NK_t \longrightarrow {}_NK_n \otimes_R K_t
$$
dans $D(K_N^{opp}, \text{d})$ dont la composée avec l'application de multiplication
est l'application de transition (dans l'un ou l'autre sens). Cela est construit dans
Algèbre à puissances divisées, lemme \ref{dpa-lemma-construct-some-maps},
par une construction explicite.
\end{proof}

\begin{proposition}
\label{sdga-proposition-derived-complete} Soit
$R$ un anneau noethérien. Soit $I \subset R$ un idéal. Les trois
catégories suivantes sont canoniquement équivalentes :
\begin{enumerate}
\item Soit $\mathcal{A}$ le faisceau de $R$-algèbres sur $\mathbf{N}$ correspondant
au système projectif de $R$-algèbres $A_n = R/I^n$. La catégorie
$\mathit{QC}(\mathcal{A})$.
\item Choisissons des générateurs $f_1, \ldots, f_r$ de $I$. Soit
$\mathcal{B}$ le faisceau de $R$-algèbres différentielles graduées sur
$\mathbf{N}$ correspondant au système projectif des algèbres de
Koszul sur $f_1^n, \ldots, f_r^n$. La catégorie
$\mathit{QC}(\mathcal{B})$.
\item La sous-catégorie pleine $D_{comp}(R, I) \subset D(R)$ des objets
complets au sens dérivé, voir Compléments sur l'algèbre, définition
\ref{more-algebra-definition-derived-complete} et le texte qui suit.
\end{enumerate}
\end{proposition}

\begin{proof}
Considérons le morphisme évident
$f : (\Sh(\mathbf{N}), \mathcal{A}) \to (\Sh(pt), R)$ de
topos annelés
et considérons les foncteurs adjoints $Lf^*$ et $Rf_*$. Le premier
se restreint à un foncteur
$$ F
: D_{comp}(R, I) \longrightarrow \mathit{QC}(\mathcal{A})
$$ qui
envoie un objet $K$ de $D_{comp}(R, I)$ représenté par un complexe
K-plat $K^\bullet$ à l'objet $(K^\bullet \otimes_R R/I^n)$ de
$\mathit{QC}(\mathcal{A})$. Le second se restreint à un foncteur
$$ G :
\mathit{QC}(\mathcal{A}) \longrightarrow D_{comp}(R, I)
$$ qui envoie un objet
$(M_n^\bullet)$ de $\mathit{QC}(\mathcal{A})$ à
$R\lim M_n^\bullet$. L'objet obtenu est complet au sens dérivé, par exemple d'après Compléments
sur l'algèbre, lemme \ref{more-algebra-lemma-naive-derived-completion}. En
outre, il résulte de Compléments sur l'algèbre, proposition
\ref{more-algebra-proposition-noetherian-naive-completion-is-completion},
que $G \circ F = \text{id}$. Ainsi, pour voir que $F$ et $G$
sont des équivalences quasi-inverses, il suffit de voir que le noyau
de $G$ est nul (voir Catégories dérivées, lemme
\ref{derived-lemma-fully-faithful-adjoint-kernel-zero}).
Cependant, il ne semble pas facile de le montrer directement !

\medskip\noindent
Montrons dans ce paragraphe que $\mathit{QC}(\mathcal{A})$ et
$\mathit{QC}(\mathcal{B})$ sont équivalentes. Écrivons
$\mathcal{B} = (B_n)$, où $B_n$ est le complexe de Koszul considéré
comme un complexe de cochaînes en degrés $-r, -r + 1, \ldots, 0$. Par
D'après Algèbre à puissances divisées, remarque \ref{dpa-remark-pro-system-koszul}
(mais avec des complexes
de chaînes transformés en complexes
de cochaînes) nous pouvons trouver $1 < n_1 < n_2 < \ldots$
et des applications de $R$-algèbres
différentielles graduées $B_{n_i} \to E_i \to R/(f_1^{n_i}, \ldots, f_r^{n_i})$
et $E_i \to B_{n_{i - 1}}$
telles que $$
\xymatrix{
B_{n_1} \ar[d]
& B_{n_2} \ar[d] \ar[l]
& B_{n_3} \ar[d] \ar[l]
& \ldots \ar[l] \\
E_1 \ar[d]
& E_2 \ar[l] \ar[d]
& E_3 \ar[l] \ar[d]
& \ldots \ar[l] \\
B_1
& B_{n_1} \ar[l]
& B_{n_2} \ar[l]
& \ldots \ar[l]
}
$$
soit un diagramme commutatif de $R$-algèbres
différentielles graduées et que $E_i \to R/(f_1^{n_i}, \ldots, f_r^{n_i})$
soit un quasi-isomorphisme.
Nous concluons \begin{enumerate}
\item il existe une équivalence entre $\mathit{QC}(\mathcal{B})$
et $\mathit{QC}((E_i))$,
\item il existe une
équivalence entre $\mathit{QC}((E_i))$ et $\mathit{QC}((R/(f_1^{n_i}, \ldots, f_r^{n_i})))$,
\item il existe une équivalence entre
$\mathit{QC}((R/(f_1^{n_i}, \ldots, f_r^{n_i})))$ et
$\mathit{QC}(\mathcal{A})$.
\end{enumerate} À
savoir, pour (1) nous pouvons appliquer le lemme \ref{sdga-lemma-pro-isomorphic}
au diagramme ci-dessus qui montre que $(E_i)$ et $(B_n)$ sont pro-isomorphes.
Pour (2) nous pouvons appliquer le lemme \ref{sdga-lemma-cartesian-qis-equivalence}
au système projectif des quasi-isomorphismes $E_i \to R/(f_1^{n_i}, \ldots, f_r^{n_i})$.
Pour (3) nous pouvons appliquer le lemme \ref{sdga-lemma-pro-isomorphic}
et le fait élémentaire que les systèmes projectifs $(R/I^n)$
et $(R/(f_1^{n_i}, \ldots, f_r^{n_i})$ sont pro-isomorphes.

\medskip\noindent
Comme dans le premier alinéa de la preuve, nous pouvons définir les
foncteurs adjoints\footnote{Il peut être démontré que ces foncteurs
sont, par les équivalences ci-dessus, compatibles avec $F$ et $G$ définis précédemment.}
$$
F' : D_{comp}(R, I) \longrightarrow \mathit{QC}(\mathcal{B})
\quad\text{et}\quad
G' : \mathit{QC}(\mathcal{B}) \longrightarrow D_{comp}(R, I).
$$
Le premier envoie un objet $K$ de $D_{comp}(R, I)$ représenté par un
complexe K-plat $K^\bullet$ à l'objet $(K^\bullet \otimes_R B_n)$ de
$\mathit{QC}(\mathcal{B})$. Le second envoie un objet
$(M_n)$ de $\mathit{QC}(\mathcal{B})$ à $R\lim M_n$. En raisonnant
comme ci-dessus, il suffit de montrer que le noyau de $G'$ est nul. Soit donc
$\mathcal{M} = (M_n)$ un bon faisceau de modules différentiels gradués
sur $\mathcal{B}$ représentant un objet de
$\mathit{QC}(\mathcal{B})$ dans le noyau de $G'$. Alors
$$ 0 =
R\lim M_n \Rightarrow 0 =
(R\lim M_n) \otimes_R^\mathbf{L} B_t =
R\lim (M_n \otimes_R^\mathbf{L} B_t)
$$ Par le lemme
\ref{sdga-lemma-different-tensors} nous avons
$R\lim (M_n \otimes_R^\mathbf{L} B_t) =
R\lim (M_n \otimes_{B_n}^\mathbf{L} B_t)$. Puisque
$(M_n)$ est un objet de $\mathit{QC}(\mathcal{B})$, nous voyons que le système
projectif $M_n \otimes_{B_n}^\mathbf{L} B_t$ est constant à partir d'un
certain rang, de valeur $M_t$. Par conséquent, $M_t = 0$, comme voulu.
\end{proof}

\begin{remark}
\label{sdga-remark-proregular} Soit
$R$ un anneau et soit $f_1, \ldots, f_r \in R$ une suite d'éléments engendrant un idéal
$I$. Soit $K_n$ le complexe de Koszul sur
$f_1^n, \ldots, f_r^n$, considéré comme une $R$-algèbre différentielle graduée. Nous disons que
$f_1, \ldots, f_r$ est une {\it suite faiblement prorégulière} si, pour tout
$n$, il existe $m > n$ tel que $K_m \to K_n$ induise l'application
nulle en cohomologie, sauf en degré $0$. Si tel est le cas, les arguments de la
preuve de la proposition \ref{sdga-proposition-derived-complete} continuent à fonctionner
même lorsque $R$ n'est pas noethérien. En particulier, nous
voyons que $\mathit{QC}(\{R/I^n\})$ est équivalente,
comme catégorie triangulée $R$-linéaire, à la catégorie
$D_{comp}(R, I)$ des objets complets au sens dérivé, à condition que
$I$ puisse être engendré par une suite faiblement prorégulière. Si cela s'avère
nécessaire, nous en donnerons ici un énoncé et une démonstration précis.
\end{remark}












% OK, commencer ici.
%

\chapter{Hyperrecouvrements}



\phantomsection
\label{hypercovering-section-phantom}


\section{Introduction}
\label{hypercovering-section-introduction}

\noindent
Soit $\mathcal{C}$ un site ; voir Sites, Définition \ref{sites-definition-site}.
Soit $X$ un objet de $\mathcal{C}$.
Étant donné un faisceau abélien $\mathcal{F}$
sur $\mathcal{C}$, nous voudrions calculer
ses groupes de cohomologie
$$
H^i(X, \mathcal{F}).
$$
D'après nos définitions générales (Cohomologie sur les sites, section
\ref{sites-cohomology-section-cohomology-sheaves})
ce groupe de cohomologie se calcule en
choisissant une résolution injective
$
0 \to \mathcal{F} \to \mathcal{I}^0 \to \mathcal{I}^1 \to \ldots
$
et en posant
$$
H^i(X, \mathcal{F})
=
H^i(
\Gamma(X, \mathcal{I}^0) \to
\Gamma(X, \mathcal{I}^1) \to
\Gamma(X, \mathcal{I}^2)\to \ldots)
$$
Le but de ce chapitre est de montrer que l'on peut aussi calculer ces
groupes de cohomologie sans choisir de résolution injective
(lorsque $\mathcal{C}$ admet des produits fibrés). Pour cela,
nous utiliserons des hyperrecouvrements.

\medskip\noindent
Un hyperrecouvrement dans un site est une généralisation d'un recouvrement ; voir
\cite[Expos\'e V, Sec. 7]{SGA4}. Étant donné un hyperrecouvrement $K$ d'un objet
$X$, il existe une suite spectrale de {\v C}ech à la cohomologie
qui exprime la cohomologie d'un faisceau abélien $\mathcal{F}$
sur $X$ en fonction de la cohomologie du faisceau sur les
composantes $K_n$ de $K$. Il se trouve qu'il existe toujours
assez d'hyperrecouvrements ; en prenant donc la limite inductive sur tous les hyperrecouvrements,
la suite spectrale dégénère et la cohomologie de $\mathcal{F}$
sur $X$ se calcule par la limite inductive des groupes de cohomologie de {\v C}ech.

\medskip\noindent
Un dispositif plus général que l'on peut considérer est une augmentation simpliciale pour laquelle
on a descente cohomologique ; voir \cite[Expos\'e Vbis]{SGA4}. Le texte de
Brian Conrad constitue un bel exposé sur la descente cohomologique ; voir
\url{https://math.stanford.edu/~conrad/papers/hypercover.pdf}.
Nous reviendrons sur ces questions dans le chapitre consacré aux espaces simpliciaux,
où nous montrerons, par exemple, que les hyperrecouvrements propres d'espaces
topologiques ``localement compacts'' sont de descente cohomologique
(Espaces simpliciaux, section
\ref{spaces-simplicial-section-proper-hypercovering}).
Notre méthode consistera à ramener cet énoncé à la suite spectrale de {\v C}ech
à la cohomologie construite dans ce chapitre.






























\section{Objets semi-représentables}
\label{hypercovering-section-semi-representable}

\noindent
Pour commencer, nous posons la définition suivante.
Les lettres « SR » signifient semi-représentable.

\begin{definition}
\label{hypercovering-definition-SR}
Soit $\mathcal{C}$ une catégorie. On note $\text{SR}(\mathcal{C})$
la catégorie des {\it objets semi-représentables}, définie comme suit :
\begin{enumerate}
\item les objets sont les familles d'objets $\{U_i\}_{i \in I}$, et
\item les morphismes $\{U_i\}_{i \in I} \to \{V_j\}_{j \in J}$ sont donnés par
une application $\alpha : I \to J$ et, pour chaque $i \in I$,
un morphisme $f_i : U_i \to V_{\alpha(i)}$ de $\mathcal{C}$.
\end{enumerate}
Soit $X \in \Ob(\mathcal{C})$ un objet de $\mathcal{C}$.
La catégorie des {\it objets semi-représentables au-dessus de $X$}
est la catégorie
$\text{SR}(\mathcal{C}, X) = \text{SR}(\mathcal{C}/X)$.
\end{definition}

\noindent
Cette définition est essentiellement équivalente à celle de
\cite[Expos\'e V, sous-section 7.3.0]{SGA4}. Notons qu'il s'agit
d'une « grande » catégorie. Nous bornerons plus loin la taille des ensembles
d'indices $I$ dont nous avons besoin pour les hyperrecouvrements de $X$. Nous pourrons alors redéfinir
$\text{SR}(\mathcal{C}, X)$ de manière à obtenir une catégorie. Décrivons explicitement
les objets et les morphismes de $\text{SR}(\mathcal{C}, X)$ :
\begin{enumerate}
\item les objets sont les familles de morphismes
$\{U_i \to X\}_{i \in I}$, et
\item les morphismes $\{U_i \to X\}_{i \in I} \to
\{V_j \to X\}_{j \in J}$ sont donnés par
une application $\alpha : I \to J$ et, pour chaque $i \in I$,
un morphisme $f_i : U_i \to V_{\alpha(i)}$ au-dessus de $X$.
\end{enumerate}
Il existe un foncteur d'oubli
$\text{SR}(\mathcal{C}, X) \to \text{SR}(\mathcal{C})$.

\begin{definition}
\label{hypercovering-definition-SR-F}
Soit $\mathcal{C}$ une catégorie.
On note $F$ le foncteur {\it qui associe un préfaisceau à un
objet semi-représentable}. En formule,
\begin{eqnarray*}
F : \text{SR}(\mathcal{C}) & \longrightarrow & \textit{PSh}(\mathcal{C}) \\
\{U_i\}_{i \in I} & \longmapsto & \amalg_{i\in I} h_{U_i}
\end{eqnarray*}
où $h_U$ désigne le préfaisceau représentable associé à
l'objet $U$.
\end{definition}

\noindent
À un morphisme $U \to X$ est associé un morphisme $h_U \to h_X$ de préfaisceaux
représentables. On considère donc souvent $F$ sur $\text{SR}(\mathcal{C}, X)$
comme un foncteur à valeurs dans la catégorie des préfaisceaux d'ensembles au-dessus de $h_X$,
à savoir $\textit{PSh}(\mathcal{C})/h_X$. Voici le diagramme correspondant :
$$
\xymatrix{
\text{SR}(\mathcal{C}, X) \ar[r]_F \ar[d] &
\textit{PSh}(\mathcal{C})/h_X \ar[d] \\
\text{SR}(\mathcal{C}) \ar[r]^F &
\textit{PSh}(\mathcal{C})
}
$$
Étudions maintenant l'existence de limites dans la catégorie des objets semi-
représentables.

\begin{lemma}
\label{hypercovering-lemma-coprod-prod-SR}
Soit $\mathcal{C}$ une catégorie.
\begin{enumerate}
\item la catégorie $\text{SR}(\mathcal{C})$ admet des coproduits
et $F$ commute à ceux-ci,
\item le foncteur $F : \text{SR}(\mathcal{C}) \to \textit{PSh}(\mathcal{C})$
commute aux limites,
\item si $\mathcal{C}$ admet des produits fibrés, alors $\text{SR}(\mathcal{C})$
admet des produits fibrés,
\item si $\mathcal{C}$ admet les produits de deux objets, alors
$\text{SR}(\mathcal{C})$ admet les produits de deux objets,
\item si $\mathcal{C}$ admet des égalisateurs, il en est de même de $\text{SR}(\mathcal{C})$, et
\item si $\mathcal{C}$ admet un objet final, il en est de même de $\text{SR}(\mathcal{C})$.
\end{enumerate}
Soit $X \in \Ob(\mathcal{C})$.
\begin{enumerate}
\item la catégorie $\text{SR}(\mathcal{C}, X)$ admet des coproduits
et $F$ commute à ceux-ci,
\item si $\mathcal{C}$ admet des produits fibrés, alors $\text{SR}(\mathcal{C}, X)$
admet des limites finies et
$F : \text{SR}(\mathcal{C}, X) \to \textit{PSh}(\mathcal{C})/h_X$
commute à celles-ci.
\end{enumerate}
\end{lemma}

\begin{proof}
Démonstration des assertions sur $\text{SR}(\mathcal{C})$.
Démonstration de (1). Le coproduit de $\{U_i\}_{i \in I}$ et $\{V_j\}_{j \in J}$ est
$\{U_i\}_{i \in I} \amalg \{V_j\}_{j \in J}$, autrement dit la famille
d'objets dont l'ensemble d'indices est $I \amalg J$ et qui, pour un élément
$k \in I \amalg J$, donne $U_i$ si $k = i \in I$ et donne $V_j$ si
$k = j \in J$. Il en va de même pour les coproduits
de familles d'objets. Il est clair que $F$ commute à ceux-ci.

\medskip\noindent
Démonstration de (2). Pour $U$ dans $\Ob(\mathcal{C})$, considérons l'objet $\{U\}$ de
$\text{SR}(\mathcal{C})$. Il est clair que
$\Mor_{\text{SR}(\mathcal{C})}(\{U\}, K)) = F(K)(U)$
pour $K \in \Ob(\text{SR}(\mathcal{C}))$. Comme les limites de préfaisceaux
se calculent au niveau des sections
(Sites, section \ref{sites-section-limits-colimits-PSh}),
on conclut que $F$ commute aux limites.

\medskip\noindent
Démonstration de (3). Supposons donnés un morphisme
$(\alpha, f_i) : \{U_i\}_{i \in I} \to \{V_j\}_{j \in J}$
et un morphisme
$(\beta, g_k) : \{W_k\}_{k \in K} \to \{V_j\}_{j \in J}$.
Le produit fibré de ces morphismes est donné par
$$
\{ U_i \times_{f_i, V_j, g_k} W_k\}_{(i, j, k) \in I \times J \times K
\text{ tels que } j = \alpha(i) = \beta(k)}
$$
Les produits fibrés existent si $\mathcal{C}$ admet des produits fibrés.

\medskip\noindent
Démonstration de (4). Le produit de $\{U_i\}_{i \in I}$ et $\{V_j\}_{j \in J}$ est
$\{U_i \times V_j\}_{i \in I, j \in J}$. Les produits existent si
$\mathcal{C}$ admet des produits.

\medskip\noindent
Démonstration de (5). L'égalisateur de deux applications
$(\alpha, f_i), (\alpha', f'_i) : \{U_i\}_{i \in I} \to \{V_j\}_{j \in J}$
est
$$
\{
\text{Eq}(f_i, f'_i : U_i \to V_{\alpha(i)})
\}_{i \in I,\ \alpha(i) = \alpha'(i)}
$$
Les égalisateurs existent si $\mathcal{C}$ admet des égalisateurs.

\medskip\noindent
Démonstration de (6). Si $X$ est un objet final de $\mathcal{C}$, alors
$\{X\}$ est un objet final de $\text{SR}(\mathcal{C})$.

\medskip\noindent
Démonstration des assertions sur $\text{SR}(\mathcal{C}, X)$.
Elles résultent des résultats précédents appliqués à la catégorie
$\mathcal{C}/X$, compte tenu de
$\text{SR}(\mathcal{C}/X) = \text{SR}(\mathcal{C}, X)$ et de
$\textit{PSh}(\mathcal{C}/X) = \textit{PSh}(\mathcal{C})/h_X$
(Sites, Lemme \ref{sites-lemma-essential-image-j-shriek} appliqué
à $\mathcal{C}$ munie de la topologie chaotique). Toutefois,
donnons aussi l'argument direct suivant.
Il est clair que le coproduit de
$\{U_i \to X\}_{i \in I}$ et $\{V_j \to X\}_{j \in J}$
est $\{U_i \to X\}_{i \in I} \amalg \{V_j \to X\}_{j \in J}$,
et de même pour les coproduits de
familles de familles de morphismes de but $X$.
L'objet $\{X \to X\}$ est un objet final
de $\text{SR}(\mathcal{C}, X)$.
Supposons donnés un morphisme
$(\alpha, f_i) : \{U_i \to X\}_{i \in I} \to \{V_j \to X\}_{j \in J}$
et un morphisme
$(\beta, g_k) : \{W_k \to X\}_{k \in K} \to \{V_j \to X\}_{j \in J}$.
Le produit fibré de ces morphismes est donné par
$$
\{ U_i \times_{f_i, V_j, g_k} W_k \to X \}_{(i, j, k) \in I \times J \times K
\text{ tels que } j = \alpha(i) = \beta(k)}
$$
Les produits fibrés existent d'après l'hypothèse selon laquelle
$\mathcal{C}$ admet des produits fibrés.
Ainsi $\text{SR}(\mathcal{C}, X)$ admet des limites finies ;
voir Catégories, Lemme \ref{categories-lemma-finite-limits-exist}.
Nous omettons la vérification des assertions sur le foncteur $F$ dans ce cas.
\end{proof}





\section{Hyperrecouvrements}
\label{hypercovering-section-hypercoverings}

\noindent
Si nous supposons que notre catégorie est un site, nous pouvons poser la
définition suivante.

\begin{definition}
\label{hypercovering-definition-covering-SR}
Soit $\mathcal{C}$ un site. Soit
$f = (\alpha, f_i) : \{U_i\}_{i \in I} \to \{V_j\}_{j \in J}$
un morphisme de la catégorie $\text{SR}(\mathcal{C})$.
On dit que $f$ est un {\it recouvrement} si, pour tout $j \in J$, la
famille de morphismes $\{U_i \to V_j\}_{i \in I, \alpha(i) = j}$
est couvrante pour le site $\mathcal{C}$.
Soit $X$ un objet de $\mathcal{C}$.
Un morphisme $K \to L$ de $\text{SR}(\mathcal{C}, X)$ est
un {\it recouvrement} si son image dans $\text{SR}(\mathcal{C})$ est
un recouvrement.
\end{definition}

\begin{lemma}
\label{hypercovering-lemma-covering-permanence}
Soit $\mathcal{C}$ un site.
\begin{enumerate}
\item Un composé de recouvrements dans $\text{SR}(\mathcal{C})$
est un recouvrement.
\item Si $K \to L$ est un recouvrement dans $\text{SR}(\mathcal{C})$
et $L' \to L$ est un morphisme, alors $L' \times_L K$ existe
et $L' \times_L K \to L'$ est un recouvrement.
\item Si $\mathcal{C}$ admet les produits de deux objets, et si
$A \to B$ et $K \to L$ sont des recouvrements dans $\text{SR}(\mathcal{C})$,
alors $A \times K \to B \times L$ est un recouvrement.
\end{enumerate}
Soit $X \in \Ob(\mathcal{C})$. Alors (1) et (2) valent pour
$\text{SR}(\mathcal{C}, X)$, et (3) vaut si $\mathcal{C}$
admet des produits fibrés.
\end{lemma}

\begin{proof}
L'assertion (1) résulte immédiatement des axiomes d'un site.
L'assertion (2) résulte de la construction des produits fibrés
dans $\text{SR}(\mathcal{C})$ donnée dans la démonstration du
Lemme \ref{hypercovering-lemma-coprod-prod-SR}
et de l'exigence que les morphismes d'un recouvrement
de $\mathcal{C}$ soient représentables.
L'assertion (3) s'obtient en considérant $A \times K \to B \times L$
comme le composé $A \times K \to B \times K \to B \times L$,
donc comme un composé de changements de base de recouvrements.
La dernière assertion résulte de l'égalité $\text{SR}(\mathcal{C}, X) =
\text{SR}(\mathcal{C}/X)$.
\end{proof}

\noindent
D'après le Lemme \ref{hypercovering-lemma-coprod-prod-SR} et le
Lemme \ref{simplicial-lemma-existence-cosk} de Simplicial,
le cosquelette d'un objet simplicial tronqué de
$\text{SR}(\mathcal{C}, X)$ existe si $\mathcal{C}$ admet des produits fibrés.
La définition suivante a donc un sens.

\begin{definition}
\label{hypercovering-definition-hypercovering}
Soit $\mathcal{C}$ un site. Supposons que $\mathcal{C}$ admette des produits fibrés.
Soit $X \in \Ob(\mathcal{C})$ un objet de $\mathcal{C}$.
Un {\it hyperrecouvrement de $X$} est un objet simplicial
$K$ de $\text{SR}(\mathcal{C}, X)$ tel que
\begin{enumerate}
\item L'objet $K_0$ est un recouvrement de $X$ pour le site $\mathcal{C}$.
\item Pour tout $n \geq 0$, le morphisme canonique
$$
K_{n + 1} \longrightarrow (\text{cosk}_n \text{sk}_n K)_{n + 1}
$$
est un recouvrement au sens défini ci-dessus.
\end{enumerate}
\end{definition}

\noindent
La condition (1) a un sens puisque tout objet de
$\text{SR}(\mathcal{C}, X)$ est, après tout, une famille
de morphismes de but $X$. On pourrait aussi la
formuler en disant que le morphisme de $K_0$ vers
l'objet final de $\text{SR}(\mathcal{C}, X)$
est un recouvrement.

\begin{example}[Hyperrecouvrements de {\v C}ech]
\label{hypercovering-example-cech}
Soit $\mathcal{C}$ un site admettant des produits fibrés.
Soit $\{U_i \to X\}_{i \in I}$ un recouvrement de $\mathcal{C}$.
Posons $K_0 = \{U_i \to X\}_{i \in I}$.
Alors $K_0$ est un objet simplicial $0$-tronqué de
$\text{SR}(\mathcal{C}, X)$. Nous pouvons donc former
$$
K = \text{cosk}_0 K_0.
$$
Il est clair que $K$ vérifie la condition (1) de la Définition \ref{hypercovering-definition-hypercovering}.
Comme tous les morphismes $K_{n + 1} \to (\text{cosk}_n \text{sk}_n K)_{n + 1}$
sont des isomorphismes d'après le
Lemme \ref{simplicial-lemma-cosk-up} de Simplicial,
il vérifie aussi la condition (2). Remarquons que
les termes $K_n$ sont les objets usuels
$$
K_n = \{
U_{i_0} \times_X U_{i_1} \times_X \ldots \times_X U_{i_n} \to X
\}_{(i_0, i_1, \ldots, i_n) \in I^{n + 1}}
$$
Un hyperrecouvrement de $X$ de cette forme est appelé un
{\it hyperrecouvrement de {\v C}ech} de $X$.
\end{example}

\begin{example}[Hyperrecouvrement par un objet simplicial du site]
\label{hypercovering-example-hypercovering-in-C}
Soit $\mathcal{C}$ un site admettant des produits fibrés. Soit
$X \in \Ob(\mathcal{C})$. Soit $U$ un objet simplicial de $\mathcal{C}$.
Comme d'habitude, on note $U_n = U([n])$. Enfin, supposons donnée une augmentation
$$
a : U \to X
$$
Dans cette situation, on peut considérer l'objet simplicial $K$
de $\text{SR}(\mathcal{C}, X)$ dont les termes sont $K_n = \{U_n \to X\}$.
Alors $K$ est un hyperrecouvrement de $X$ au sens de la
Définition \ref{hypercovering-definition-hypercovering}
si et seulement si les trois
conditions suivantes\footnote{Comme $\mathcal{C}$ admet des produits fibrés, la
catégorie $\mathcal{C}/X$ admet toutes les limites finies.
Les cosquelettes requis existent donc d'après le
Lemme \ref{simplicial-lemma-existence-cosk} de Simplicial.} sont vérifiées :
\begin{enumerate}
\item $\{U_0 \to X\}$ est un recouvrement de $\mathcal{C}$,
\item $\{U_1 \to U_0 \times_X U_0\}$ est un recouvrement de $\mathcal{C}$,
\item $\{U_{n + 1} \to (\text{cosk}_n\text{sk}_n U)_{n + 1}\}$
est un recouvrement de $\mathcal{C}$ pour $n \geq 1$.
\end{enumerate}
Nous omettons la vérification immédiate.
\end{example}

\begin{example}[Hyperrecouvrement de {\v C}ech associé à un recouvrement]
\label{hypercovering-example-cech-cover}
Soit $\mathcal{C}$ un site admettant des produits fibrés. Soit $U \to X$ un
morphisme de $\mathcal{C}$ tel que $\{U \to X\}$ soit un recouvrement de
$\mathcal{C}$\footnote{Un morphisme de $\mathcal{C}$ ayant cette propriété
est parfois appelé un « recouvrement ».}. Considérons l'objet simplicial $K$ de
$\text{SR}(\mathcal{C}, X)$ dont les termes sont
$$
K_n = \{U \times_X U \times_X \ldots \times_X U \to X\}
\quad (n + 1 \text{ facteurs})
$$
Alors $K$ est un hyperrecouvrement de $X$. Cet exemple est un cas particulier à la fois de
l'Exemple \ref{hypercovering-example-cech} et de
l'Exemple \ref{hypercovering-example-hypercovering-in-C}.
\end{example}

\begin{lemma}
\label{hypercovering-lemma-hypercoverings-set}
Soit $\mathcal{C}$ un site admettant des produits fibrés.
Soit $X \in \Ob(\mathcal{C})$ un objet de $\mathcal{C}$.
La collection de tous les hyperrecouvrements de $X$ forme un ensemble.
\end{lemma}

\begin{proof}
Puisque $\mathcal{C}$ est un site, l'ensemble de tous les recouvrements de
$X$ forme un ensemble. On voit donc que la collection
des $K_0$ possibles forme un ensemble. Supposons avoir montré que
la collection de tous les $K_0, \ldots, K_n$ possibles forme
un ensemble. Il suffit alors de montrer qu'étant donnés
$K_0, \ldots, K_n$, la collection de tous les
$K_{n + 1}$ possibles forme un ensemble. Cela est clair, puisque
il faut choisir $K_{n + 1}$ parmi tous les recouvrements possibles
de $(\text{cosk}_n \text{sk}_n K)_{n + 1}$.
\end{proof}

\begin{remark}
\label{hypercovering-remark-hypercoverings-really-set}
Le lemme ne dit pas seulement qu'il existe un système cofinal
de choix d'hyperrecouvrements qui soit un ensemble,
mais que les hyperrecouvrements eux-mêmes forment effectivement un ensemble.
\end{remark}

\noindent
La catégorie des préfaisceaux sur $\mathcal{C}$ admet des
(co)limites finies. Les foncteurs $\text{cosk}_n$
existent donc pour les préfaisceaux d'ensembles.

\begin{lemma}
\label{hypercovering-lemma-hypercovering-F}
Soit $\mathcal{C}$ un site admettant des produits fibrés.
Soit $X \in \Ob(\mathcal{C})$ un objet de $\mathcal{C}$.
Soit $K$ un hyperrecouvrement de $X$.
Considérons l'objet simplicial $F(K)$ de $\textit{PSh}(\mathcal{C})$,
muni de son augmentation vers le préfaisceau simplicial constant $h_X$.
\begin{enumerate}
\item Le morphisme de préfaisceaux $F(K)_0 \to h_X$ devient
surjectif après passage aux faisceaux associés.
\item Le morphisme
$$
(d^1_0, d^1_1) :
F(K)_1
\longrightarrow
F(K)_0 \times_{h_X} F(K)_0
$$
devient surjectif après passage aux faisceaux associés.
\item Pour tout $n \geq 1$, le morphisme
$$
F(K)_{n + 1} \longrightarrow (\text{cosk}_n \text{sk}_n F(K))_{n + 1}
$$
devient surjectif après passage aux faisceaux associés.
\end{enumerate}
\end{lemma}

\begin{proof}
Nous utiliserons le fait que si
$\{U_i \to U\}_{i \in I}$ est un recouvrement du site
$\mathcal{C}$, alors le morphisme
$$
\amalg_{i \in I} h_{U_i} \to h_U
$$
devient surjectif après passage aux faisceaux associés ; voir
Sites, Lemme \ref{sites-lemma-covering-surjective-after-sheafification}.
La première assertion en résulte immédiatement.

\medskip\noindent
Pour la deuxième assertion, remarquons que, d'après
Simplicial, Exemple \ref{simplicial-example-cosk0},
l'objet simplicial $\text{cosk}_0 \text{sk}_0 K$
a pour termes $K_0 \times \ldots \times K_0$. Ainsi,
d'après la définition d'un hyperrecouvrement,
on voit que $(d^1_0, d^1_1) : K_1 \to K_0 \times K_0$ est un
recouvrement. L'assertion (2) résulte donc de l'observation précédente
et du fait que $F$ transforme les produits en produits
fibrés au-dessus de $h_X$.

\medskip\noindent
Pour la troisième assertion, nous affirmons que
$\text{cosk}_n \text{sk}_n F(K) =
F(\text{cosk}_n \text{sk}_n K)$ pour $n \geq 1$.
Pour le montrer, notons temporairement $F'$ le foncteur
$\text{SR}(\mathcal{C}, X) \to \textit{PSh}(\mathcal{C})/h_X$.
D'après le Lemme \ref{hypercovering-lemma-coprod-prod-SR}, le foncteur
$F'$ commute aux limites finies.
La description du foncteur $\text{cosk}_n$ donnée dans
Simplicial, section \ref{simplicial-section-skeleton},
montre que $\text{cosk}_n \text{sk}_n F'(K) =
F'(\text{cosk}_n \text{sk}_n K)$.
Rappelons que la catégorie intervenant dans la description de
$(\text{cosk}_n U)_m$ du
Lemme \ref{simplicial-lemma-existence-cosk} de Simplicial
est la catégorie $(\Delta/[m])^{opp}_{\leq n}$. C'est un
exercice plaisant que de montrer que $(\Delta/[m])_{\leq n}$ est
une catégorie connexe (voir
Catégories, Définition \ref{categories-definition-category-connected})
dès que $n \geq 1$. Par conséquent, le
Lemme \ref{categories-lemma-connected-limit-over-X} de Catégories
montre que $\text{cosk}_n \text{sk}_n F'(K) =
\text{cosk}_n \text{sk}_n F(K)$. D'où l'assertion.
L'assertion (3) en résulte, car on voit maintenant que
le morphisme de (3) s'obtient en appliquant le
foncteur $F$ à un recouvrement comme dans la Définition \ref{hypercovering-definition-covering-SR},
et le résultat découle du premier fait mentionné
dans cette démonstration.
\end{proof}



\section{Acyclicité}
\label{hypercovering-section-acyclicity}

\noindent
Soit $\mathcal{C}$ un site.
Pour un préfaisceau d'ensembles $\mathcal{F}$, on note $\mathbf{Z}_\mathcal{F}$
le préfaisceau de groupes abéliens défini par la règle
$$
\mathbf{Z}_\mathcal{F}(U) = \text{groupe abélien libre sur }\mathcal{F}(U).
$$
Nous l'appellerons parfois le {\it préfaisceau abélien libre sur $\mathcal{F}$}.
Bien entendu, la construction $\mathcal{F} \mapsto \mathbf{Z}_\mathcal{F}$
est un foncteur, adjoint à gauche du foncteur d'oubli
$\textit{PAb}(\mathcal{C}) \to \textit{PSh}(\mathcal{C})$.
Bien entendu, le faisceau associé $\mathbf{Z}_\mathcal{F}^\#$ est
un faisceau de groupes abéliens, et le foncteur
$\mathcal{F} \mapsto \mathbf{Z}_\mathcal{F}^\#$ est lui aussi
un adjoint à gauche. Nous appelons parfois $\mathbf{Z}_\mathcal{F}^\#$
le {\it faisceau abélien libre sur $\mathcal{F}$}.

\medskip\noindent
Pour un objet $X$ du site $\mathcal{C}$, on note
$\mathbf{Z}_X$ le préfaisceau abélien libre sur $h_X$, et
on note $\mathbf{Z}_X^\#$ son faisceau associé.

\begin{definition}
\label{hypercovering-definition-homology}
Soit $\mathcal{C}$ un site.
Soit $K$ un objet simplicial de $\textit{PSh}(\mathcal{C})$.
Par ce qui précède, on obtient un objet simplicial $\mathbf{Z}_K^\#$ de
$\textit{Ab}(\mathcal{C})$. On peut prendre son
complexe associé de préfaisceaux abéliens $s(\mathbf{Z}_K^\#)$ ; voir
Simplicial, section \ref{simplicial-section-complexes}.
L'{\it homologie de $K$} est l'homologie du
complexe de faisceaux abéliens $s(\mathbf{Z}_K^\#)$.
\end{definition}

\noindent
Autrement dit, la {\it $i$-ième homologie $H_i(K)$ de $K$}
est le faisceau de groupes abéliens $H_i(K) = H_i(s(\mathbf{Z}_K^\#))$.
Dans cette section, nous étudions l'homologie lorsque $K$
est un hyperrecouvrement d'un objet $X$ de $\mathcal{C}$.

\begin{lemma}
\label{hypercovering-lemma-compare-cosk0}
Soit $\mathcal{C}$ un site.
Soit $\mathcal{F} \to \mathcal{G}$ un morphisme
de préfaisceaux d'ensembles. Notons $K$ l'objet
simplicial de $\textit{PSh}(\mathcal{C})$ dont le $n$-ième
terme est le $(n + 1)$-ième produit fibré de $\mathcal{F}$
au-dessus de $\mathcal{G}$ ; voir
Simplicial, Exemple \ref{simplicial-example-fibre-products-simplicial-object}.
Alors, si $\mathcal{F} \to \mathcal{G}$ devient surjectif après
passage aux faisceaux associés, on a
$$
H_i(K) =
\left\{
\begin{matrix}
0 & \text{si} & i > 0\\
\mathbf{Z}_\mathcal{G}^\# & \text{si} & i = 0
\end{matrix}
\right.
$$
L'isomorphisme en degré $0$ est donné par le
morphisme $H_0(K) \to \mathbf{Z}_\mathcal{G}^\#$
provenant de l'application $(\mathbf{Z}_K^\#)_0 =
\mathbf{Z}_\mathcal{F}^\# \to \mathbf{Z}_\mathcal{G}^\#$.
\end{lemma}

\begin{proof}
Soit $\mathcal{G}' \subset \mathcal{G}$ l'image du
morphisme $\mathcal{F} \to \mathcal{G}$.
Soit $U \in \Ob(\mathcal{C})$. Posons
$A = \mathcal{F}(U)$ et $B = \mathcal{G}'(U)$.
Alors l'ensemble simplicial $K(U)$ est égal à l'ensemble
simplicial dont les $n$-simplexes sont donnés par
$$
A \times_B A \times_B \ldots \times_B A\ (n + 1 \text{ facteurs)}.
$$
D'après Simplicial, Lemme \ref{simplicial-lemma-cosk-minus-one-equivalence},
le morphisme $K(U) \to B$ est une fibration de Kan triviale.
C'est donc une équivalence d'homotopie
(Simplicial, Lemme \ref{simplicial-lemma-trivial-kan-homotopy}).
En appliquant le foncteur « groupe abélien libre sur » à ce morphisme,
on en déduit que
$$
\mathbf{Z}_K(U) \longrightarrow \mathbf{Z}_B
$$
est une équivalence d'homotopie. Remarquons que $s(\mathbf{Z}_B)$ est
le complexe
$$
\ldots \to
\bigoplus\nolimits_{b \in B}\mathbf{Z} \xrightarrow{0}
\bigoplus\nolimits_{b \in B}\mathbf{Z} \xrightarrow{1}
\bigoplus\nolimits_{b \in B}\mathbf{Z} \xrightarrow{0}
\bigoplus\nolimits_{b \in B}\mathbf{Z} \to 0
$$
voir Simplicial, Lemme \ref{simplicial-lemma-homology-eilenberg-maclane}.
On voit donc que
$H_i(s(\mathbf{Z}_K(U))) = 0$ pour $i > 0$, et
$H_0(s(\mathbf{Z}_K(U))) = \bigoplus_{b \in B}\mathbf{Z}
= \bigoplus_{s \in \mathcal{G}'(U)} \mathbf{Z}$.
Ces identifications sont compatibles aux morphismes de
restriction.

\medskip\noindent
On conclut que $H_i(s(\mathbf{Z}_K)) = 0$ pour $i > 0$ et que
$H_0(s(\mathbf{Z}_K)) = \mathbf{Z}_{\mathcal{G}'}$, où l'on
calcule ici les groupes d'homologie dans $\textit{PAb}(\mathcal{C})$. Comme
le foncteur « faisceau associé » est exact, on en déduit le résultat
du lemme. En effet, l'exactitude implique
que $H_0(s(\mathbf{Z}_K))^\# = H_0(s(\mathbf{Z}_K^\#))$,
et de même pour les autres indices.
\end{proof}

\begin{lemma}
\label{hypercovering-lemma-acyclicity}
Soit $\mathcal{C}$ un site.
Soit $f : L \to K$ un morphisme d'objets
simpliciaux de $\textit{PSh}(\mathcal{C})$.
Soit $n \geq 0$ un entier.
Supposons que
\begin{enumerate}
\item Pour $i < n$, le morphisme $L_i \to K_i$ est un isomorphisme.
\item Le morphisme $L_n \to K_n$ devient surjectif après passage aux faisceaux associés.
\item L'application canonique $L \to \text{cosk}_n \text{sk}_n L$ est un isomorphisme.
\item L'application canonique $K \to \text{cosk}_n \text{sk}_n K$ est un isomorphisme.
\end{enumerate}
Alors $H_i(f) : H_i(L) \to H_i(K)$ est un isomorphisme.
\end{lemma}

\begin{proof}
Cette démonstration est exactement la même que celle du
Lemme \ref{hypercovering-lemma-compare-cosk0} ci-dessus. En effet,
on commence par noter $K_n' \subset K_n$ le sous-préfaisceau
qui est l'image de l'application $L_n \to K_n$. L'hypothèse
(2) signifie que le faisceau associé à $K_n'$ est égal au
faisceau associé à $K_n$. De plus, puisque $L_i = K_i$
pour tout $i < n$, on voit que l'on obtient un préfaisceau simplicial
$n$-tronqué $U$ en prenant
$U_0 = L_0 = K_0, \ldots, U_{n - 1} = L_{n - 1} = K_{n - 1}, U_n = K'_n$.
Notons $K' = \text{cosk}_n U$, qui est un préfaisceau simplicial.
Comme on peut construire $K'_m$ par une limite finie, et
que le foncteur « faisceau associé » est exact, on voit que
$(K'_m)^\# = K_m$. Autrement dit, $(K')^\# = K^\#$.
On conclut, en utilisant encore une fois l'exactitude du faisceau associé,
que $H_i(K) = H_i(K')$. Il suffit donc de démontrer le lemme
pour le morphisme $L \to K'$ ; autrement dit, on peut
supposer que $L_n \to K_n$ est un morphisme surjectif
de {\it préfaisceaux} !

\medskip\noindent
Dans ce cas, pour tout objet $U$ de $\mathcal{C}$,
on voit que le morphisme d'ensembles simpliciaux
$$
L(U) \longrightarrow K(U)
$$
satisfait toutes les hypothèses du
Lemme \ref{simplicial-lemma-section} de Simplicial.
C'est donc une fibration de Kan triviale. En particulier, c'est
une équivalence d'homotopie
(Simplicial, Lemme \ref{simplicial-lemma-trivial-kan-homotopy}).
Ainsi
$$
\mathbf{Z}_L(U) \longrightarrow \mathbf{Z}_K(U)
$$
est elle aussi une équivalence d'homotopie, et ce pour tout $U$.
Le résultat en découle.
\end{proof}

\begin{lemma}
\label{hypercovering-lemma-acyclic-hypercover-sheaves}
Soit $\mathcal{C}$ un site.
Soit $K$ un préfaisceau simplicial.
Soit $\mathcal{G}$ un préfaisceau.
Soit $K \to \mathcal{G}$ une augmentation de $K$
vers $\mathcal{G}$. Supposons que
\begin{enumerate}
\item Le morphisme de préfaisceaux $K_0 \to \mathcal{G}$ devienne
surjectif après passage aux faisceaux associés.
\item Le morphisme
$$
(d^1_0, d^1_1) :
K_1
\longrightarrow
K_0 \times_\mathcal{G} K_0
$$
devienne surjectif après passage aux faisceaux associés.
\item Pour tout $n \geq 1$, le morphisme
$$
K_{n + 1} \longrightarrow (\text{cosk}_n \text{sk}_n K)_{n + 1}
$$
devienne surjectif après passage aux faisceaux associés.
\end{enumerate}
Alors $H_i(K) = 0$ pour $i > 0$ et
$H_0(K) = \mathbf{Z}_\mathcal{G}^\#$.
\end{lemma}

\begin{proof}
Notons $K^n = \text{cosk}_n \text{sk}_n K$ pour $n \geq 1$.
Définissons $K^0$ comme l'objet simplicial dont les termes
$(K^0)_n$ sont égaux au produit fibré itéré $(n + 1)$ fois
$K_0 \times_\mathcal{G} \ldots \times_\mathcal{G} K_0$ ;
voir Simplicial,
Exemple \ref{simplicial-example-fibre-products-simplicial-object}.
On a des morphismes
$$
K \longrightarrow \ldots \to K^n \to K^{n - 1} \to \ldots \to K^1 \to K^0.
$$
Les morphismes $K \to K^i$, $K^j \to K^i$ pour $j \geq i \geq 1$ proviennent
des propriétés universelles des foncteurs $\text{cosk}_n$.
Le morphisme $K^1 \to K^0$ est le morphisme canonique
de
Simplicial, Remarque \ref{simplicial-remark-augmentation}.
Rappelons aussi que $K^0 \to \text{cosk}_1 \text{sk}_1 K^0$
est un isomorphisme ; voir
Simplicial, Lemme \ref{simplicial-lemma-cosk-minus-one}.

\medskip\noindent
D'après le Lemme \ref{hypercovering-lemma-compare-cosk0}, on voit que
$H_i(K^0) = 0$ pour $i > 0$ et $H_0(K^0) = \mathbf{Z}_\mathcal{G}^\#$.

\medskip\noindent
Fixons $n \geq 1$. Considérons le morphisme $K^n \to K^{n - 1}$.
C'est un isomorphisme sur les termes de degré $< n$.
Remarquons que $K^n \to \text{cosk}_n \text{sk}_n K^n$ et
$K^{n - 1} \to \text{cosk}_n \text{sk}_n K^{n - 1}$
sont des isomorphismes. Remarquons que $(K^n)_n = K_n$ et
que $(K^{n - 1})_n = (\text{cosk}_{n - 1} \text{sk}_{n - 1} K)_n$.
D'après l'hypothèse, on a donc que $(K^n)_n \to (K^{n - 1})_n$
est un morphisme de préfaisceaux qui devient surjectif après
passage aux faisceaux associés. Le Lemme \ref{hypercovering-lemma-acyclicity} donne alors
$H_i(K^n) = H_i(K^{n - 1})$.
Combiné avec ce qui précède, ceci démontre le lemme.
\end{proof}

\begin{lemma}
\label{hypercovering-lemma-hypercovering-acyclic}
Soit $\mathcal{C}$ un site admettant des produits fibrés.
Soit $X$ un objet de $\mathcal{C}$.
Soit $K$ un hyperrecouvrement de $X$.
L'homologie du préfaisceau simplicial $F(K)$ est
$0$ dans les degrés $> 0$ et égale à $\mathbf{Z}_X^\#$
en degré $0$.
\end{lemma}

\begin{proof}
Combiner les Lemmes \ref{hypercovering-lemma-acyclic-hypercover-sheaves}
et \ref{hypercovering-lemma-hypercovering-F}.
\end{proof}












\section{Cohomologie de {\v C}ech et hyperrecouvrements}
\label{hypercovering-section-hyper-cech}

\noindent
Soit $\mathcal{C}$ un site. Considérons un préfaisceau de
groupes abéliens $\mathcal{F}$ sur le site $\mathcal{C}$.
Il définit un foncteur
\begin{eqnarray*}
\mathcal{F} : \text{SR}(\mathcal{C})^{opp}
& \longrightarrow &
\textit{Ab} \\
\{U_i\}_{i \in I} &
\longmapsto &
\prod\nolimits_{i \in I} \mathcal{F}(U_i)
\end{eqnarray*}
Ainsi, un objet simplicial $K$ de $\text{SR}(\mathcal{C})$
donne un objet cosimplicial $\mathcal{F}(K)$ de $\textit{Ab}$.
Le complexe de cochaînes $s(\mathcal{F}(K))$ associé à $\mathcal{F}(K)$
(Simplicial, section
\ref{simplicial-section-dold-kan-cosimplicial})
est appelé le complexe de {\v C}ech de $\mathcal{F}$ relativement
à l'objet simplicial $K$. On pose
$$
\check{H}^i(K, \mathcal{F})
=
H^i(s(\mathcal{F}(K))).
$$
et on l'appelle le $i$-ième groupe de cohomologie de {\v C}ech
de $\mathcal{F}$ relativement à $K$.
Dans cette section, nous démontrons des analogues de certains résultats sur la
cohomologie de {\v C}ech des recouvrements ouverts établis dans
Cohomologie, sections \ref{cohomology-section-cech},
\ref{cohomology-section-cech-functor} et
\ref{cohomology-section-cech-cohomology-cohomology}.

\begin{lemma}
\label{hypercovering-lemma-h0-cech}
Soit $\mathcal{C}$ un site admettant des produits fibrés.
Soit $X$ un objet de $\mathcal{C}$.
Soit $K$ un hyperrecouvrement de $X$.
Soit $\mathcal{F}$ un faisceau de groupes abéliens sur $\mathcal{C}$.
Alors $\check{H}^0(K, \mathcal{F}) = \mathcal{F}(X)$.
\end{lemma}

\begin{proof}
On a
$$
\check{H}^0(K, \mathcal{F})
=
\Ker(\mathcal{F}(K_0) \longrightarrow \mathcal{F}(K_1))
$$
Écrivons $K_0 = \{U_i \to X\}$. C'est un recouvrement dans le site
$\mathcal{C}$. De même, $K_1 \to K_0 \times K_0$
est un recouvrement dans $\text{SR}(\mathcal{C}, X)$. On peut donc
écrire $K_1 = \amalg_{i_0, i_1 \in I} \{V_{i_0i_1j} \to X\}$
de telle sorte que le morphisme $K_1 \to K_0 \times K_0$ soit donné
par les recouvrements $\{V_{i_0i_1j} \to U_{i_0} \times_X U_{i_1}\}$
du site $\mathcal{C}$. On peut alors encore identifier
$$
\check{H}^0(K, \mathcal{F})
=
\Ker(
\prod\nolimits_i \mathcal{F}(U_i)
\longrightarrow
\prod\nolimits_{i_0i_1 j} \mathcal{F}(V_{i_0i_1j})
)
$$
où l'application est évidente. La propriété de faisceau de $\mathcal{F}$
implique que $\check{H}^0(K, \mathcal{F}) = H^0(X, \mathcal{F})$.
\end{proof}

\noindent
Cette propriété caractérise bien sûr les faisceaux abéliens parmi tous les
préfaisceaux abéliens sur $\mathcal{C}$.
L'analogue de Cohomologie, Lemme \ref{hypercovering-lemma-injective-trivial-cech},
est ici le suivant.

\begin{lemma}
\label{hypercovering-lemma-injective-trivial-cech}
Soit $\mathcal{C}$ un site admettant des produits fibrés.
Soit $X$ un objet de $\mathcal{C}$.
Soit $K$ un hyperrecouvrement de $X$.
Soit $\mathcal{I}$ un faisceau injectif de groupes abéliens sur $\mathcal{C}$.
Alors
$$
\check{H}^p(K, \mathcal{I}) =
\left\{
\begin{matrix}
\mathcal{I}(X) & \text{si} & p = 0 \\
0 & \text{si} & p > 0
\end{matrix}
\right.
$$
\end{lemma}

\begin{proof}
Observons que, pour tout objet $Z = \{U_i \to X\}$ de
$\text{SR}(\mathcal{C}, X)$ et tout faisceau abélien
$\mathcal{F}$ sur $\mathcal{C}$, on a
\begin{eqnarray*}
\mathcal{F}(Z)
& = &
\prod \mathcal{F}(U_i) \\
& = &
\prod \Mor_{\textit{PSh}(\mathcal{C})}(h_{U_i}, \mathcal{F})\\
& = &
\Mor_{\textit{PSh}(\mathcal{C})}(F(Z), \mathcal{F})\\
& = &
\Mor_{\textit{PAb}(\mathcal{C})}(\mathbf{Z}_{F(Z)}, \mathcal{F}) \\
& = &
\Mor_{\textit{Ab}(\mathcal{C})}(\mathbf{Z}_{F(Z)}^\#, \mathcal{F})
\end{eqnarray*}
On voit donc que, pour tout objet simplicial $K$ de
$\text{SR}(\mathcal{C}, X)$, on a
\begin{equation}
\label{hypercovering-equation-identify-cech}
s(\mathcal{F}(K))
=
\Hom_{\textit{Ab}(\mathcal{C})}(s(\mathbf{Z}_{F(K)}^\#), \mathcal{F})
\end{equation}
voir la Définition \ref{hypercovering-definition-homology} pour les notations.
Le complexe de faisceaux $s(\mathbf{Z}_{F(K)}^\#)$ est quasi-isomorphe
à $\mathbf{Z}_X^\#$ si $K$ est un hyperrecouvrement ; voir le
Lemme \ref{hypercovering-lemma-hypercovering-acyclic}. On en conclut
que si $\mathcal{I}$ est un faisceau abélien injectif et si
$K$ est un hyperrecouvrement, alors le complexe $s(\mathcal{I}(K))$
est acyclique, sauf peut-être en degré $0$.
Autrement dit, on a
$$
\check{H}^i(K, \mathcal{I}) = 0
$$
pour $i > 0$. Combiné au Lemme \ref{hypercovering-lemma-h0-cech}, ceci démontre le lemme.
\end{proof}

\noindent
Venons-en à l'analogue de Cohomologie sur les sites, Lemme
\ref{sites-cohomology-lemma-cech-spectral-sequence}.
Soit $\mathcal{C}$ un site.
Soit $\mathcal{F}$ un faisceau de groupes abéliens sur $\mathcal{C}$.
Rappelons que $\underline{H}^i(\mathcal{F})$ désigne le préfaisceau
de groupes abéliens sur $\mathcal{C}$ défini par la
règle $\underline{H}^i(\mathcal{F}) : U \longmapsto H^i(U, \mathcal{F})$.
On l'étend à $\text{SR}(\mathcal{C})$ comme dans l'introduction
de cette section.

\begin{lemma}
\label{hypercovering-lemma-cech-spectral-sequence}
Soit $\mathcal{C}$ un site admettant des produits fibrés.
Soit $X$ un objet de $\mathcal{C}$.
Soit $K$ un hyperrecouvrement de $X$.
Soit $\mathcal{F}$ un faisceau de groupes abéliens sur $\mathcal{C}$.
Il existe un morphisme
$$
s(\mathcal{F}(K))
\longrightarrow
R\Gamma(X, \mathcal{F})
$$
dans $D^{+}(\textit{Ab})$, fonctoriel en $\mathcal{F}$, qui induit
des transformations naturelles
$$
\check{H}^i(K, -) \longrightarrow H^i(X, -)
$$
entre foncteurs $\textit{Ab}(\mathcal{C}) \to \textit{Ab}$. De plus,
il existe une suite spectrale $(E_r, d_r)_{r \geq 0}$ avec
$$
E_2^{p, q} = \check{H}^p(K, \underline{H}^q(\mathcal{F}))
$$
convergeant vers $H^{p + q}(X, \mathcal{F})$.
Cette suite spectrale est fonctorielle en $\mathcal{F}$ et
en l'hyperrecouvrement $K$.
\end{lemma}

\begin{proof}
On pourrait le démontrer par la même méthode que celle employée dans le lemme
correspondant du chapitre consacré à la cohomologie. Donnons plutôt une
démonstration par un complexe double.

\medskip\noindent
Choisissons une résolution injective $\mathcal{F} \to \mathcal{I}^\bullet$
dans la catégorie des faisceaux abéliens sur $\mathcal{C}$. Considérons le
complexe double $A^{\bullet, \bullet}$ de termes
$$
A^{p, q} = \mathcal{I}^q(K_p)
$$
où la différentielle $d_1^{p, q} : A^{p, q} \to A^{p + 1, q}$ est celle
qui provient de la différentielle du complexe $s(\mathcal{I}^q(K))$
associé au groupe abélien cosimplicial $\mathcal{I}^p(K)$,
et où la différentielle $d_2^{p, q} : A^{p, q} \to A^{p, q + 1}$ est celle
qui provient de la différentielle $\mathcal{I}^q \to \mathcal{I}^{q + 1}$.
Notons $\text{Tot}(A^{\bullet, \bullet})$ le complexe total associé au
complexe double $A^{\bullet, \bullet}$ ; voir
Homologie, section \ref{homology-section-double-complexes}.
Nous utiliserons les deux suites
spectrales $({}'E_r, {}'d_r)$ et $({}''E_r, {}''d_r)$
associées à ce complexe double ; voir
Homologie, section \ref{homology-section-double-complex}.

\medskip\noindent
D'après le Lemme \ref{hypercovering-lemma-injective-trivial-cech},
les complexes $s(\mathcal{I}^q(K))$ sont acycliques en
degrés positifs et ont pour $H^0$ le groupe $\mathcal{I}^q(X)$.
Par conséquent, d'après
Homologie, Lemme \ref{homology-lemma-double-complex-gives-resolution},
le morphisme naturel
$$
\mathcal{I}^\bullet(X) \longrightarrow \text{Tot}(A^{\bullet, \bullet})
$$
est un quasi-isomorphisme de complexes de groupes abéliens. En particulier,
on conclut que $H^n(\text{Tot}(A^{\bullet, \bullet})) = H^n(X, \mathcal{F})$.

\medskip\noindent
Le morphisme $s(\mathcal{F}(K)) \longrightarrow R\Gamma(X, \mathcal{F})$ du
lemme est le composé du morphisme
$s(\mathcal{F}(K)) \to \text{Tot}(A^{\bullet, \bullet})$
avec l'inverse
du quasi-isomorphisme affiché ci-dessus. Cela convient puisque
$\mathcal{I}^\bullet(X)$ représente $R\Gamma(X, \mathcal{F})$.

\medskip\noindent
Considérons la suite spectrale $({}'E_r, {}'d_r)_{r \geq 0}$. D'après
Homologie, Lemme \ref{homology-lemma-ss-double-complex},
on voit que
$$
{}'E_2^{p, q} = H^p_I(H^q_{II}(A^{\bullet, \bullet}))
$$
Autrement dit, on calcule d'abord la cohomologie relativement à
$d_2$, ce qui donne les groupes
${}'E_1^{p, q} = \underline{H}^q(\mathcal{F})(K_p)$.
Il en résulte bien (d'après la description de la différentielle
${}'d_1$) que
${}'E_2^{p, q} = \check{H}^p(K, \underline{H}^q(\mathcal{F}))$.
De ce qui précède et de Homologie, Lemme \ref{homology-lemma-first-quadrant-ss},
on déduit que cette suite converge vers $H^n(X, \mathcal{F})$, comme voulu.

\medskip\noindent
Nous omettons la démonstration des assertions relatives à la fonctorialité des
constructions précédentes en le faisceau abélien $\mathcal{F}$ et en
l'hyperrecouvrement $K$.
\end{proof}









\section{Hyperrecouvrements à la Verdier}
\label{hypercovering-section-hypercoverings-verdier}

\noindent
Le lecteur attentif aura remarqué que, pour obtenir
la suite spectrale de {\v C}ech à la cohomologie associée à un
hyperrecouvrement d'un objet $X$, il suffit de disposer de la
conclusion du Lemme \ref{hypercovering-lemma-hypercovering-F}.
La définition suivante a donc un sens.

\begin{definition}
\label{hypercovering-definition-hypercovering-variant}
Soit $\mathcal{C}$ un site. Supposons que $\mathcal{C}$ admette des égalisateurs
et des produits fibrés. Soit $\mathcal{G}$ un préfaisceau d'ensembles.
Un {\it hyperrecouvrement de $\mathcal{G}$} est un objet simplicial
$K$ de $\text{SR}(\mathcal{C})$, muni d'une augmentation
$F(K) \to \mathcal{G}$ telle que
\begin{enumerate}
\item $F(K_0) \to \mathcal{G}$ devienne surjectif
après passage aux faisceaux associés,
\item $F(K_1) \to F(K_0) \times_\mathcal{G} F(K_0)$
devienne surjectif après passage aux faisceaux associés, et
\item $F(K_{n + 1}) \longrightarrow F((\text{cosk}_n \text{sk}_n K)_{n + 1})$
devienne surjectif après passage aux faisceaux associés pour $n \geq 1$.
\end{enumerate}
On dit qu'un objet simplicial $K$ de $\text{SR}(\mathcal{C})$
est un {\it hyperrecouvrement} si $K$ est un hyperrecouvrement de l'objet
final $*$ de $\textit{PSh}(\mathcal{C})$.
\end{definition}

\noindent
L'hypothèse selon laquelle $\mathcal{C}$ admet des produits fibrés et des égalisateurs
garantit que $\text{SR}(\mathcal{C})$ admet des produits fibrés
et des égalisateurs, et que $F$ commute à ceux-ci
(Lemme \ref{hypercovering-lemma-coprod-prod-SR}), ce qui suffit
pour définir les foncteurs cosquelette utilisés (voir
Simplicial, Remarque \ref{simplicial-remark-existence-cosk} et
Catégories, Lemme \ref{categories-lemma-fibre-products-equalizers-exist}).
Si $\mathcal{C}$ est quelconque, on peut remplacer la condition (3) par la
condition selon laquelle
$F(K_{n + 1}) \longrightarrow ((\text{cosk}_n \text{sk}_n F(K))_{n + 1})$
devient surjectif après passage aux faisceaux associés pour $n \geq 1$, et les
résultats de cette section restent valables.

\medskip\noindent
Soit $\mathcal{F}$ un faisceau abélien sur $\mathcal{C}$.
Dans la section précédente, nous avons défini le complexe de {\v C}ech de $\mathcal{F}$
relativement à un objet simplicial $K$ de $\text{SR}(\mathcal{C})$.
Puis, étant donné un préfaisceau $\mathcal{G}$, on pose
$$
H^0(\mathcal{G}, \mathcal{F}) =
\Mor_{\textit{PSh}(\mathcal{C})}(\mathcal{G}, \mathcal{F}) =
\Mor_{\Sh(\mathcal{C})}(\mathcal{G}^\#, \mathcal{F}) =
H^0(\mathcal{G}^\#, \mathcal{F})
$$
avec les notations de
Cohomologie sur les sites, section \ref{sites-cohomology-section-limp}.
C'est un foncteur exact à gauche, et ses foncteurs dérivés supérieurs
(brièvement étudiés dans
Cohomologie sur les sites, section \ref{sites-cohomology-section-limp})
sont notés $H^i(\mathcal{G}, \mathcal{F})$.
Nous montrerons qu'étant donné un hyperrecouvrement $K$ de $\mathcal{G}$,
il existe une suite spectrale de {\v C}ech à la cohomologie convergeant vers la
cohomologie $H^i(\mathcal{G}, \mathcal{F})$.
Remarquons que si $\mathcal{G} = *$, alors
$H^i(*, \mathcal{F}) = H^i(\mathcal{C}, \mathcal{F})$ redonne
la cohomologie de $\mathcal{F}$ sur le site $\mathcal{C}$.

\begin{lemma}
\label{hypercovering-lemma-h0-cech-variant}
Soit $\mathcal{C}$ un site admettant des égalisateurs et des produits fibrés.
Soit $\mathcal{G}$ un préfaisceau sur $\mathcal{C}$.
Soit $K$ un hyperrecouvrement de $\mathcal{G}$.
Soit $\mathcal{F}$ un faisceau de groupes abéliens sur $\mathcal{C}$.
Alors $\check{H}^0(K, \mathcal{F}) = H^0(\mathcal{G}, \mathcal{F})$.
\end{lemma}

\begin{proof}
Cela résulte de la définition de $H^0(\mathcal{G}, \mathcal{F})$
et du fait que
$$
\xymatrix{
F(K_1) \ar@<1ex>[r] \ar@<-1ex>[r] &
F(K_0) \ar[r] & \mathcal{G}
}
$$
devient un diagramme coégalisateur après passage aux faisceaux associés.
\end{proof}

\begin{lemma}
\label{hypercovering-lemma-injective-trivial-cech-variant}
Soit $\mathcal{C}$ un site admettant des égalisateurs et des produits fibrés.
Soit $\mathcal{G}$ un préfaisceau sur $\mathcal{C}$.
Soit $K$ un hyperrecouvrement de $\mathcal{G}$.
Soit $\mathcal{I}$ un faisceau injectif de groupes abéliens sur $\mathcal{C}$.
Alors
$$
\check{H}^p(K, \mathcal{I}) =
\left\{
\begin{matrix}
H^0(\mathcal{G}, \mathcal{I}) & \text{si} & p = 0 \\
0 & \text{si} & p > 0
\end{matrix}
\right.
$$
\end{lemma}

\begin{proof}
D'après (\ref{hypercovering-equation-identify-cech}), on a
$$
s(\mathcal{F}(K))
=
\Hom_{\textit{Ab}(\mathcal{C})}(s(\mathbf{Z}_{F(K)}^\#), \mathcal{F})
$$
Le complexe $s(\mathbf{Z}_{F(K)}^\#)$ est quasi-isomorphe
à $\mathbf{Z}_\mathcal{G}^\#$ ; voir le
Lemme \ref{hypercovering-lemma-acyclic-hypercover-sheaves}. On en conclut
que si $\mathcal{I}$ est un faisceau abélien injectif, alors
le complexe $s(\mathcal{I}(K))$ est acyclique, sauf peut-être en degré $0$.
Autrement dit, on a $\check{H}^i(K, \mathcal{I}) = 0$
pour $i > 0$. Combiné au Lemme \ref{hypercovering-lemma-h0-cech-variant},
ceci démontre le lemme.
\end{proof}

\begin{lemma}
\label{hypercovering-lemma-cech-spectral-sequence-variant}
Soit $\mathcal{C}$ un site admettant des égalisateurs et des produits fibrés.
Soit $\mathcal{G}$ un préfaisceau sur $\mathcal{C}$.
Soit $K$ un hyperrecouvrement de $\mathcal{G}$.
Soit $\mathcal{F}$ un faisceau de groupes abéliens sur $\mathcal{C}$.
Il existe un morphisme
$$
s(\mathcal{F}(K)) \longrightarrow R\Gamma(\mathcal{G}, \mathcal{F})
$$
dans $D^{+}(\textit{Ab})$, fonctoriel en $\mathcal{F}$, qui induit
une transformation naturelle
$$
\check{H}^i(K, -) \longrightarrow H^i(\mathcal{G}, -)
$$
de foncteurs $\textit{Ab}(\mathcal{C}) \to \textit{Ab}$. De plus,
il existe une suite spectrale $(E_r, d_r)_{r \geq 0}$ avec
$$
E_2^{p, q} = \check{H}^p(K, \underline{H}^q(\mathcal{F}))
$$
convergeant vers $H^{p + q}(\mathcal{G}, \mathcal{F})$.
Cette suite spectrale est fonctorielle en $\mathcal{F}$ et
en l'hyperrecouvrement $K$.
\end{lemma}

\begin{proof}
Choisissons une résolution injective $\mathcal{F} \to \mathcal{I}^\bullet$
dans la catégorie des faisceaux abéliens sur $\mathcal{C}$. Considérons le
complexe double $A^{\bullet, \bullet}$ de termes
$$
A^{p, q} = \mathcal{I}^q(K_p)
$$
où la différentielle $d_1^{p, q} : A^{p, q} \to A^{p + 1, q}$
est celle provenant de la différentielle $\mathcal{I}^p \to \mathcal{I}^{p + 1}$,
et où la différentielle $d_2^{p, q} : A^{p, q} \to A^{p, q + 1}$ est
celle provenant de la différentielle du complexe
$s(\mathcal{I}^p(K))$ associé au groupe abélien cosimplicial
$\mathcal{I}^p(K)$, comme expliqué ci-dessus.
Nous utiliserons les deux suites
spectrales $({}'E_r, {}'d_r)$ et $({}''E_r, {}''d_r)$
associées à ce complexe double ; voir
Homologie, section \ref{homology-section-double-complex}.

\medskip\noindent
D'après le Lemme \ref{hypercovering-lemma-injective-trivial-cech-variant}, les complexes
$s(\mathcal{I}^p(K))$ sont acycliques en degrés positifs et ont pour
$H^0$ le groupe $H^0(\mathcal{G}, \mathcal{I}^p)$. Par conséquent, d'après
Homologie, Lemme \ref{homology-lemma-double-complex-gives-resolution},
et sa démonstration, la suite spectrale $({}'E_r, {}'d_r)$ dégénère,
et le morphisme naturel
$$
H^0(\mathcal{G}, \mathcal{I}^\bullet) \longrightarrow
\text{Tot}(A^{\bullet, \bullet})
$$
est un quasi-isomorphisme de complexes de groupes abéliens. Le morphisme
$s(\mathcal{F}(K)) \longrightarrow R\Gamma(\mathcal{G}, \mathcal{F})$
du lemme est le composé du morphisme naturel
$s(\mathcal{F}(K)) \to \text{Tot}(A^{\bullet, \bullet})$
avec l'inverse du quasi-isomorphisme affiché ci-dessus.
Cela convient parce que $H^0(\mathcal{G}, \mathcal{I}^\bullet)$
représente $R\Gamma(\mathcal{G}, \mathcal{F})$.

\medskip\noindent
Considérons la suite spectrale $({}''E_r, {}''d_r)_{r \geq 0}$. D'après
Homologie, Lemme \ref{homology-lemma-ss-double-complex},
on voit que
$$
{}''E_2^{p, q} = H^p_{II}(H^q_I(A^{\bullet, \bullet}))
$$
Autrement dit, on calcule d'abord la cohomologie relativement à
$d_1$, ce qui donne les groupes
${}''E_1^{p, q} = \underline{H}^p(\mathcal{F})(K_q)$.
Il en résulte bien (d'après la description de la différentielle
${}''d_1$) que
${}''E_2^{p, q} = \check{H}^p(K, \underline{H}^q(\mathcal{F}))$.
Puisque cette suite spectrale converge vers la cohomologie de
$\text{Tot}(A^{\bullet, \bullet})$, la démonstration est terminée.
\end{proof}

\begin{lemma}
\label{hypercovering-lemma-cech-spectral-sequence-verdier}
Soit $\mathcal{C}$ un site admettant des égalisateurs et des produits fibrés.
Soit $K$ un hyperrecouvrement.
Soit $\mathcal{F}$ un faisceau abélien. Il existe une
suite spectrale $(E_r, d_r)_{r \geq 0}$ avec
$$
E_2^{p, q} = \check{H}^p(K, \underline{H}^q(\mathcal{F}))
$$
convergeant vers les groupes de cohomologie globale $H^{p + q}(\mathcal{F})$.
\end{lemma}

\begin{proof}
C'est un cas particulier du Lemme \ref{hypercovering-lemma-cech-spectral-sequence-variant}.
\end{proof}







\section{Recouvrement des hyperrecouvrements}
\label{hypercovering-section-covering}

\noindent
Voici quelques méthodes pour construire des hyperrecouvrements.
Remarquons que, puisque la catégorie
$\text{SR}(\mathcal{C}, X)$ admet des produits fibrés,
la catégorie des objets simpliciaux
de $\text{SR}(\mathcal{C}, X)$ en admet elle aussi ;
voir Simplicial, Lemme \ref{simplicial-lemma-fibre-product}.

\begin{lemma}
\label{hypercovering-lemma-funny-gamma}
Soit $\mathcal{C}$ un site admettant des produits fibrés.
Soit $X$ un objet de $\mathcal{C}$.
Soient $K, L, M$ des objets simpliciaux de $\text{SR}(\mathcal{C}, X)$.
Soient $a : K \to L$, $b : M \to L$ des morphismes.
Supposons que
\begin{enumerate}
\item $K$ soit un hyperrecouvrement de $X$,
\item le morphisme $M_0 \to L_0$ soit un recouvrement, et
\item pour tout $n \geq 0$, dans le diagramme
$$
\xymatrix{
M_{n + 1} \ar[dd] \ar[rr] \ar[rd]^\gamma &
&
(\text{cosk}_n \text{sk}_n M)_{n + 1} \ar[dd] \\
&
L_{n + 1}
\times_{(\text{cosk}_n \text{sk}_n L)_{n + 1}}
(\text{cosk}_n \text{sk}_n M)_{n + 1}
\ar[ld] \ar[ru]
& \\
L_{n + 1} \ar[rr] & & (\text{cosk}_n \text{sk}_n L)_{n + 1}
}
$$
la flèche $\gamma$ soit un recouvrement.
\end{enumerate}
Alors le produit fibré $K \times_L M$ est un hyperrecouvrement de $X$.
\end{lemma}

\begin{proof}
Le morphisme $(K \times_L M)_0 = K_0 \times_{L_0} M_0 \to K_0$
est le changement de base d'un recouvrement d'après (2), donc un recouvrement ; voir
le Lemme \ref{hypercovering-lemma-covering-permanence}. Et $K_0 \to \{X \to X\}$
est un recouvrement d'après (1). Ainsi, $(K \times_L M)_0 \to \{X \to X\}$
est un recouvrement d'après le Lemme \ref{hypercovering-lemma-covering-permanence}. Par conséquent,
$K \times_L M$ satisfait la première condition de la Définition
\ref{hypercovering-definition-hypercovering}.

\medskip\noindent
Il reste à vérifier que
$$
K_{n + 1} \times_{L_{n + 1}} M_{n + 1} = (K \times_L M)_{n + 1}
\longrightarrow
(\text{cosk}_n \text{sk}_n (K \times_L M))_{n + 1}
$$
est un recouvrement pour tout $n \geq 0$. Adoptons les abréviations suivantes :
$A = (\text{cosk}_n \text{sk}_n K)_{n + 1}$,
$B = (\text{cosk}_n \text{sk}_n L)_{n + 1}$, et
$C = (\text{cosk}_n \text{sk}_n M)_{n + 1}$.
Le foncteur $\text{cosk}_n \text{sk}_n$ commute aux produits fibrés ;
voir Simplicial, Lemme \ref{simplicial-lemma-cosk-fibre-product}.
Le membre de droite ci-dessus est donc égal à $A \times_B C$.
Considérons le diagramme commutatif suivant
$$
\xymatrix{
K_{n + 1} \times_{L_{n + 1}} M_{n + 1} \ar[r] \ar[d] &
M_{n + 1} \ar[d] \ar[rd]_\gamma \ar[rrd] &
& \\
K_{n + 1} \ar[r] \ar[rd] &
L_{n + 1} \ar[rrd] &
L_{n + 1} \times_B C \ar[l] \ar[r] &
C \ar[d] \\
&
A \ar[rr] &
&
B
}
$$
Ce diagramme montre que
$$
K_{n + 1} \times_{L_{n + 1}} M_{n + 1}
=
(K_{n + 1} \times_B C)
\times_{(L_{n + 1} \times_B C), \gamma}
M_{n + 1}
$$
Or $K_{n + 1} \times_B C \to A \times_B C$
est le changement de base du recouvrement $K_{n + 1} \to A$
par le morphisme $A \times_B C \to A$ ; c'est donc un
recouvrement. D'après l'hypothèse (3), le morphisme $\gamma$ est un recouvrement.
Par conséquent, le morphisme
$$
(K_{n + 1} \times_B C)
\times_{(L_{n + 1} \times_B C), \gamma}
M_{n + 1}
\longrightarrow
K_{n + 1} \times_B C
$$
est un recouvrement, puisqu'il est le changement de base d'un recouvrement.
Le lemme en découle, car un composé de recouvrements
est un recouvrement.
\end{proof}

\begin{lemma}
\label{hypercovering-lemma-product-hypercoverings}
Soit $\mathcal{C}$ un site admettant des produits fibrés.
Soit $X$ un objet de $\mathcal{C}$.
Si $K, L$ sont des hyperrecouvrements de $X$, alors
$K \times L$ est un hyperrecouvrement de $X$.
\end{lemma}

\begin{proof}
On peut soit le vérifier directement, soit appliquer le
Lemme \ref{hypercovering-lemma-funny-gamma} ci-dessus et vérifier que $L \to \{X \to X\}$
possède la propriété (3).
\end{proof}


\noindent
Soit $\mathcal{C}$ un site admettant des produits fibrés.
Soit $X$ un objet de $\mathcal{C}$.
Puisque la catégorie $\text{SR}(\mathcal{C}, X)$ admet des coproduits et des
limites finies, il est permis de parler des objets
$U \times K$ et $\Hom(U, K)$ pour certains ensembles simpliciaux $U$
(par exemple ceux qui ne possèdent qu'un nombre fini de simplexes non dégénérés)
et tout objet simplicial $K$ de $\text{SR}(\mathcal{C}, X)$.
Voir Simplicial, sections
\ref{simplicial-section-product-with-simplicial-sets} et
\ref{simplicial-section-hom-from-simplicial-sets}.

\begin{lemma}
\label{hypercovering-lemma-covering}
Soit $\mathcal{C}$ un site admettant des produits fibrés.
Soit $X$ un objet de $\mathcal{C}$.
Soit $K$ un hyperrecouvrement de $X$.
Soit $k \geq 0$ un entier.
Soit $u : Z \to K_k$ un recouvrement
dans $\text{SR}(\mathcal{C}, X)$.
Alors il existe un morphisme d'hyperrecouvrements
$f: L \to K$ tel que $L_k \to K_k$
se factorise par $u$.
\end{lemma}

\begin{proof}
Notons $Y = K_k$. Soit $C[k]$ l'ensemble cosimplicial défini dans
Simplicial, Exemple \ref{simplicial-example-simplex-cosimplicial-set}.
Nous utiliserons la description de $\Hom(C[k], Y)$ et $\Hom(C[k], Z)$
donnée dans
Simplicial, Lemme \ref{simplicial-lemma-morphism-into-product}.
Il existe un morphisme canonique
$K \to \Hom(C[k], Y)$ correspondant à $\text{id} : K_k = Y \to Y$.
Considérons le morphisme $\Hom(C[k], Z) \to \Hom(C[k], Y)$
qui, sur les termes de degré $n$, est le morphisme
$$
\prod\nolimits_{\alpha : [k] \to [n]} Z
\longrightarrow
\prod\nolimits_{\alpha : [k] \to [n]} Y
$$
obtenu en utilisant le morphisme donné $Z \to Y$ sur chaque facteur. Posons
$$
L = K \times_{\Hom(C[k], Y)} \Hom(C[k], Z).
$$
Le morphisme $L_k \to K_k$ s'insère dans un diagramme commutatif
$$
\xymatrix{
L_k \ar[r] \ar[d] &
\prod_{\alpha : [k] \to [k]} Z \ar[r]^-{\text{pr}_{\text{id}_{[k]}}} \ar[d] &
Z \ar[d] \\
K_k \ar[r] &
\prod_{\alpha : [k] \to [k]} Y \ar[r]^-{\text{pr}_{\text{id}_{[k]}}} &
Y
}
$$
Comme le composé des deux flèches du bas est l'identité,
on en conclut que l'on a la factorisation voulue.

\medskip\noindent
Il reste à montrer que $L$ est un hyperrecouvrement de $X$.
Pour cela, nous appliquerons le Lemme \ref{hypercovering-lemma-funny-gamma}.
La condition (1) est satisfaite par hypothèse.
Pour (2), le morphisme
$$
\Hom(C[k], Z)_0 \to \Hom(C[k], Y)_0
$$
est un recouvrement parce qu'il est isomorphe à $Z \to Y$, puisqu'il
n'existe qu'un seul morphisme $[k] \to [0]$.

\medskip\noindent
Considérons la condition (3) pour $n = 0$. Alors, puisque
$(\text{cosk}_0 T)_1 = T \times T$
(Simplicial, Exemple \ref{simplicial-example-cosk0})
et puisque $\Hom(C[k], Z)_1 = \prod_{\alpha : [k] \to [1]} Z$,
on obtient le diagramme
$$
\xymatrix{
\prod\nolimits_{\alpha : [k] \to [1]} Z \ar[r] \ar[d] &
Z \times Z \ar[d] \\
\prod\nolimits_{\alpha : [k] \to [1]} Y \ar[r] &
Y \times Y
}
$$
où les flèches horizontales correspondent à la projection sur les facteurs
associés aux deux $\alpha$ non surjectifs. La flèche $\gamma$
est donc le morphisme
$$
\prod\nolimits_{\alpha : [k] \to [1]} Z
\longrightarrow
\prod\nolimits_{\alpha : [k] \to [1]\text{ non surjective}} Z
\times
\prod\nolimits_{\alpha : [k] \to [1]\text{ surjective}} Y
$$
qui est un produit de recouvrements, donc un recouvrement d'après le
Lemme \ref{hypercovering-lemma-covering-permanence}.

\medskip\noindent
Considérons la condition (3) pour $n > 0$. Nous affirmons qu'il existe une
application injective $\tau : S' \to S$ d'ensembles finis telle que, pour tout
objet $T$ de $\text{SR}(\mathcal{C}, X)$, le morphisme
\begin{equation}
\label{hypercovering-equation-map}
\Hom(C[k], T)_{n + 1} \to
(\text{cosk}_n \text{sk}_n \Hom(C[k], T))_{n + 1}
\end{equation}
soit isomorphe à la projection $\prod_{s \in S} T \to \prod_{s' \in S'} T$,
fonctoriellement en $T$. Si cela est vrai, on voit, en raisonnant comme au
paragraphe précédent, que la flèche $\gamma$ est le morphisme
$$
\prod\nolimits_{s \in S} Z
\longrightarrow
\prod\nolimits_{s \in S'} Z
\times
\prod\nolimits_{s \not\in \tau(S')} Y
$$
qui est un produit de recouvrements, donc un recouvrement d'après le
Lemme \ref{hypercovering-lemma-covering-permanence}. Par construction, on a
$\Hom(C[k], T)_{n + 1} = \prod_{\alpha : [k] \to [n + 1]} T$
(voir Simplicial, Lemme \ref{simplicial-lemma-morphism-into-product}).
On prend donc $S = \text{Map}([k], [n + 1])$.
D'autre part, Simplicial, Lemme \ref{simplicial-lemma-formula-limit},
donne une description des points de
$(\text{cosk}_n \text{sk}_n \Hom(C[k], T))_{n + 1}$
comme des suites $(f_0, \ldots, f_{n + 1})$ de points de $\Hom(C[k], T)_n$
vérifiant $d^n_{j - 1} f_i = d^n_i f_j$ pour $0 \leq i < j \leq n + 1$.
On peut écrire $f_i = (f_{i, \alpha})$, où $f_{i, \alpha}$ est un point de $T$
et $\alpha \in \text{Map}([k], [n])$. Les conditions se traduisent par
$$
f_{i, \delta^n_{j - 1} \circ \beta} = f_{j, \delta_i^n \circ \beta}
$$
pour tout $0 \leq i < j \leq n + 1$ et tout $\beta : [k] \to [n - 1]$. On
voit donc que
$$
S' = \{0, \ldots, n + 1\} \times \text{Map}([k], [n]) / \sim
$$
où la relation d'équivalence est engendrée par les équivalences
$$
(i, \delta^n_{j - 1} \circ \beta) \sim (j, \delta_i^n \circ \beta)
$$
pour $0 \leq i < j \leq n + 1$ et $\beta : [k] \to [n - 1]$.
Un calcul (omis) montre que le morphisme (\ref{hypercovering-equation-map})
correspond à l'application $S' \to S$ qui envoie $(i, \alpha)$ sur
$\delta^{n + 1}_i \circ \alpha \in S$. (Le lecteur pourra être rassuré
de voir que cette application est bien définie d'après la partie (1) du
Lemme \ref{simplicial-lemma-relations-face-degeneracy} de Simplicial.)
Pour achever la démonstration, il suffit de montrer que si
$\alpha, \alpha' : [k] \to [n]$ et $0 \leq i < j \leq n + 1$
sont tels que
$$
\delta^{n + 1}_i \circ \alpha = \delta^{n + 1}_j \circ \alpha'
$$
alors on a $\alpha = \delta^n_{j - 1} \circ \beta$
et $\alpha' = \delta_i^n \circ \beta$ pour un certain $\beta : [k] \to [n - 1]$.
Cela se voit aisément, et nous l'omettons.
\end{proof}

\begin{lemma}
\label{hypercovering-lemma-covering-sheaf}
Soit $\mathcal{C}$ un site admettant des produits fibrés.
Soit $X$ un objet de $\mathcal{C}$.
Soit $K$ un hyperrecouvrement de $X$.
Soit $n \geq 0$ un entier.
Soit $u : \mathcal{F} \to F(K_n)$ un morphisme
de préfaisceaux qui devient surjectif
après passage aux faisceaux associés.
Alors il existe un morphisme d'hyperrecouvrements
$f: L \to K$ tel que $F(f_n) : F(L_n) \to F(K_n)$
se factorise par $u$.
\end{lemma}

\begin{proof}
Écrivons $K_n = \{U_i \to X\}_{i \in I}$.
Ainsi, l'application $u$ est un morphisme de préfaisceaux d'ensembles
$u : \mathcal{F} \to \amalg h_{u_i}$.
L'hypothèse sur $u$ signifie que, pour tout
$i \in I$, il existe un recouvrement $\{U_{ij} \to U_i\}_{j \in I_i}$
du site $\mathcal{C}$ et un morphisme de préfaisceaux
$t_{ij} : h_{U_{ij}} \to \mathcal{F}$ tel que
$u \circ t_{ij}$ soit l'application $h_{U_{ij}} \to h_{U_i}$
provenant du morphisme $U_{ij} \to U_i$.
Posons $J = \amalg_{i \in I} I_i$, et soit
$\alpha : J \to I$ l'application évidente.
Pour $j \in J$, notons $V_j = U_{\alpha(j)j}$. Posons
$Z = \{V_j \to X\}_{j \in J}$.
Enfin, considérons le morphisme
$u' : Z \to K_n$ donné par $\alpha : J \to I$
et par les morphismes $V_j = U_{\alpha(j)j} \to U_{\alpha(j)}$
ci-dessus. Il est clair que c'est un recouvrement dans la
catégorie $\text{SR}(\mathcal{C}, X)$ et, par
construction, $F(u') : F(Z) \to F(K_n)$ se factorise par $u$.
Le résultat découle donc du Lemme \ref{hypercovering-lemma-covering} ci-dessus.
\end{proof}


\section{Ajout de simplexes}
\label{hypercovering-section-adding-simplices}

\noindent
Dans cette section, nous démontrons quelques lemmes techniques qui serviront plus loin.
Soit $\mathcal{C}$ un site admettant des produits fibrés.
Soit $X$ un objet de $\mathcal{C}$.
Comme nous l'avons signalé dans la section \ref{hypercovering-section-covering} ci-dessus,
les objets $U \times K$ et $\Hom(U, K)$
sont définis pour certains ensembles simpliciaux $U$
et tout objet simplicial $K$ de $\text{SR}(\mathcal{C}, X)$.
Voir Simplicial, sections
\ref{simplicial-section-product-with-simplicial-sets} et
\ref{simplicial-section-hom-from-simplicial-sets}.

\begin{lemma}
\label{hypercovering-lemma-one-more-simplex}
Soit $\mathcal{C}$ un site admettant des produits fibrés.
Soit $X$ un objet de $\mathcal{C}$.
Soit $K$ un hyperrecouvrement de $X$.
Soient $U \subset V$ des ensembles simpliciaux tels que $U_n, V_n$
soient finis non vides pour tout $n$.
Supposons que $U$ n'ait qu'un nombre fini de simplexes non dégénérés.
Supposons que $n \geq 0$ et que $x \in V_n$,
$x \not \in U_n$, soient tels que
\begin{enumerate}
\item $V_i = U_i$ pour $i < n$,
\item $V_n = U_n \cup \{x\}$,
\item tout $z \in V_j$, $z \not \in U_j$, pour $j > n$,
soit dégénéré.
\end{enumerate}
Alors le morphisme
$$
\Hom(V, K)_0
\longrightarrow
\Hom(U, K)_0
$$
de $\text{SR}(\mathcal{C}, X)$ est un recouvrement.
\end{lemma}

\begin{proof}
Si $n = 0$, il résulte aisément que $V = U \amalg \Delta[0]$
(voir ci-dessous). Dans ce cas, $\Hom(V, K)_0 =
\Hom(U, K)_0 \times K_0$. Le résultat découle alors
du Lemme \ref{hypercovering-lemma-covering-permanence}.

\medskip\noindent
Soit $a : \Delta[n] \to V$ le morphisme associé à $x$
comme dans Simplicial, Lemme \ref{simplicial-lemma-simplex-map}.
Écrivons $\partial \Delta[n] = i_{(n-1)!} \text{sk}_{n - 1} \Delta[n]$
pour le $(n - 1)$-squelette de $\Delta[n]$.
Soit $b : \partial \Delta[n] \to U$ la restriction
de $a$ au $(n - 1)$-squelette de $\Delta[n]$. D'après
Simplicial, Lemme \ref{simplicial-lemma-glue-simplex},
on a $V = U \amalg_{\partial \Delta[n]} \Delta[n]$. D'après
Simplicial, Lemme
\ref{simplicial-lemma-hom-from-coprod},
le carré
$$
\xymatrix{
\Hom(V, K)_0 \ar[r] \ar[d] &
\Hom(U, K)_0 \ar[d] \\
\Hom(\Delta[n], K)_0 \ar[r] &
\Hom(\partial \Delta[n], K)_0
}
$$
est cartésien. Il suffit donc de montrer que
la flèche horizontale inférieure est un recouvrement. D'après
Simplicial, Lemme \ref{simplicial-lemma-cosk-shriek},
cette flèche s'identifie à
$$
K_n \to (\text{cosk}_{n - 1} \text{sk}_{n - 1} K)_n
$$
et c'est donc un recouvrement par définition d'un hyperrecouvrement.
\end{proof}

\begin{lemma}
\label{hypercovering-lemma-add-simplices}
Soit $\mathcal{C}$ un site admettant des produits fibrés.
Soit $X$ un objet de $\mathcal{C}$.
Soit $K$ un hyperrecouvrement de $X$.
Soient $U \subset V$ des ensembles simpliciaux tels que $U_n, V_n$
soient finis non vides pour tout $n$.
Supposons que $U$ et $V$ n'aient qu'un nombre fini de simplexes non dégénérés.
Alors le morphisme
$$
\Hom(V, K)_0
\longrightarrow
\Hom(U, K)_0
$$
de $\text{SR}(\mathcal{C}, X)$ est un recouvrement.
\end{lemma}

\begin{proof}
D'après le Lemme \ref{hypercovering-lemma-one-more-simplex}
ci-dessus, il suffit de démontrer un lemme élémentaire
sur les inclusions d'ensembles simpliciaux $U \subset V$ comme dans le
lemme. Or c'est exactement le résultat de
Simplicial, Lemme \ref{simplicial-lemma-add-simplices}.
\end{proof}

\begin{lemma}
\label{hypercovering-lemma-degeneracy-maps-coverings}
Soit $\mathcal{C}$ un site admettant des produits fibrés. Soit $X$ un objet de
$\mathcal{C}$. Soit $K$ un hyperrecouvrement de $X$. Alors
\begin{enumerate}
\item $K_n$ est un recouvrement de $X$ pour tout $n \geq 0$,
\item $d^n_i : K_n \to K_{n - 1}$ est un recouvrement pour tous $n \geq 1$
et $0 \leq i \leq n$.
\end{enumerate}
\end{lemma}

\begin{proof}
Rappelons que $K_0$ est un recouvrement de $X$ d'après la
Définition \ref{hypercovering-definition-hypercovering}
et que cela équivaut à dire que
$K_0 \to \{X \to X\}$ est un recouvrement au sens
de la Définition \ref{hypercovering-definition-covering-SR}.
L'assertion (1) résulte donc de (2), car celle-ci montrera que
le composé
$K_n \to K_{n - 1} \to \ldots \to K_0 \to \{X \to X\}$
est un recouvrement d'après le Lemme \ref{hypercovering-lemma-covering-permanence}.

\medskip\noindent
Démonstration de (2). Observons que
$\Mor(\Delta[n], K)_0 = K_n$ d'après
Simplicial, Lemme \ref{simplicial-lemma-exists-hom-from-simplicial-set-finite}.
L'assertion (2) résulte donc du Lemme \ref{hypercovering-lemma-add-simplices}
appliqué aux $n + 1$ inclusions distinctes $\Delta[n - 1] \to \Delta[n]$.
\end{proof}

\begin{remark}
\label{hypercovering-remark-P-covering}
Un cas particulier utile des Lemmes \ref{hypercovering-lemma-add-simplices} et
\ref{hypercovering-lemma-degeneracy-maps-coverings} est le suivant.
Supposons donnée une catégorie $\mathcal{C}$ admettant des produits fibrés.
Soit $P \subset \text{Arrows}(\mathcal{C})$ un sous-ensemble
stable par changement de base, stable par composition
et contenant tous les isomorphismes. On appelle alors
{\it $P$-hyperrecouvrement} une augmentation $a : U \to X$
issue d'un objet simplicial de $\mathcal{C}$ telle que
\begin{enumerate}
\item $U_0 \to X$ appartienne à $P$,
\item $U_1 \to U_0 \times_X U_0$ appartienne à $P$,
\item $U_{n + 1} \to (\text{cosk}_n\text{sk}_n U)_{n + 1}$
appartienne à $P$ pour $n \geq 1$.
\end{enumerate}
La catégorie $\mathcal{C}/X$ admet toutes les limites finies ; les
cosquelettes employés dans la formulation ci-dessus existent donc
(voir Catégories, Lemme \ref{categories-lemma-finite-limits-exist}).
Nous affirmons alors que les morphismes $U_n \to X$ et $d^n_i : U_n \to U_{n - 1}$
appartiennent à $P$. Cela résulte des lemmes précités
en munissant $\mathcal{C}$ de la structure de site dont les recouvrements
sont les $\{f : V \to U\}$ avec $f \in P$, et en prenant pour $K$ l'objet défini par
$K_n = \{U_n \to X\}$.
\end{remark}


\section{Homotopies}
\label{hypercovering-section-homotopies}

\noindent
Soit $\mathcal{C}$ un site admettant des produits fibrés.
Soit $X$ un objet de $\mathcal{C}$.
Soit $L$ un objet simplicial de $\text{SR}(\mathcal{C}, X)$.
D'après
Simplicial, Lemme \ref{simplicial-lemma-exists-hom-from-simplicial-set-finite},
il existe un objet $\Hom(\Delta[1], L)$
de la catégorie $\text{Simp}(\text{SR}(\mathcal{C}, X))$ qui représente le
foncteur
$$
T
\longmapsto
\Mor_{\text{Simp}(\text{SR}(\mathcal{C}, X))}(\Delta[1] \times T, L)
$$
Il existe un morphisme canonique
$$
\Hom(\Delta[1], L) \to L \times L
$$
provenant de $e_i : \Delta[0] \to \Delta[1]$ et de l'identification
$\Hom(\Delta[0], L) = L$.

\begin{lemma}
\label{hypercovering-lemma-hom-hypercovering}
Soit $\mathcal{C}$ un site admettant des produits fibrés.
Soit $X$ un objet de $\mathcal{C}$.
Soit $L$ un objet simplicial de $\text{SR}(\mathcal{C}, X)$.
Soit $n \geq 0$. Considérons le diagramme commutatif
\begin{equation}
\label{hypercovering-equation-diagram}
\xymatrix{
\Hom(\Delta[1], L)_{n + 1} \ar[r] \ar[d] &
(\text{cosk}_n \text{sk}_n \Hom(\Delta[1], L))_{n + 1} \ar[d] \\
(L \times L)_{n + 1} \ar[r] &
(\text{cosk}_n \text{sk}_n (L \times L))_{n + 1}
}
\end{equation}
provenant du morphisme défini ci-dessus.
On peut identifier les termes de ce diagramme comme suit,
où
$\partial \Delta[n + 1] = i_{n!}\text{sk}_n \Delta[n + 1]$
est le $n$-squelette du $(n + 1)$-simplexe :
\begin{eqnarray*}
\Hom(\Delta[1], L)_{n + 1}
& = &
\Hom(\Delta[1] \times \Delta[n + 1], L)_0 \\
(\text{cosk}_n \text{sk}_n \Hom(\Delta[1], L))_{n + 1}
& = &
\Hom(\Delta[1] \times \partial \Delta[n + 1], L)_0 \\
(L \times L)_{n + 1}
& = &
\Hom(
(\Delta[n + 1] \amalg \Delta[n + 1], L)_0 \\
(\text{cosk}_n \text{sk}_n (L \times L))_{n + 1}
& = &
\Hom(
\partial \Delta[n + 1]
\amalg
\partial \Delta[n + 1], L)_0
\end{eqnarray*}
et les morphismes entre ces objets de $\text{SR}(\mathcal{C}, X)$
proviennent du diagramme commutatif d'ensembles simpliciaux
\begin{equation}
\label{hypercovering-equation-dual-diagram}
\xymatrix{
\Delta[1] \times \Delta[n + 1] &
\Delta[1] \times \partial\Delta[n + 1] \ar[l] \\
\Delta[n + 1] \amalg \Delta[n + 1] \ar[u] &
\partial\Delta[n + 1] \amalg \partial\Delta[n + 1]
\ar[l] \ar[u]
}
\end{equation}
De plus, le produit fibré de la flèche inférieure et de la
flèche de droite dans (\ref{hypercovering-equation-diagram}) est égal à
$$
\Hom(U, L)_0
$$
où $U \subset \Delta[1] \times \Delta[n + 1]$
est le plus petit sous-ensemble simplicial tel que
$\Delta[n + 1] \amalg \Delta[n + 1]$ et
$\Delta[1] \times \partial\Delta[n + 1]$ s'y envoient tous deux.
\end{lemma}

\begin{proof}
La première et la troisième égalité sont données par
Simplicial, Lemme \ref{simplicial-lemma-exists-hom-from-simplicial-set-finite}.
La deuxième et la quatrième résultent du lemme cité, combiné au
Lemme \ref{simplicial-lemma-cosk-shriek} de Simplicial.
La dernière assertion résulte du fait que
$U$ est la somme amalgamée de la flèche inférieure et de la flèche droite du
diagramme (\ref{hypercovering-equation-dual-diagram}), d'après
Simplicial, Lemme \ref{simplicial-lemma-hom-from-coprod}.
Pour voir que $U$ est égal à cette somme amalgamée, il suffit
de voir que l'intersection de
$\Delta[n + 1] \amalg \Delta[n + 1]$ et de
$\Delta[1] \times \partial\Delta[n + 1]$
dans $\Delta[1] \times \Delta[n + 1]$ est égale à
$\partial\Delta[n + 1] \amalg \partial\Delta[n + 1]$.
Nous laissons cette vérification au lecteur.
\end{proof}

\begin{lemma}
\label{hypercovering-lemma-homotopy}
Soit $\mathcal{C}$ un site admettant des produits fibrés.
Soit $X$ un objet de $\mathcal{C}$.
Soient $K, L$ des hyperrecouvrements de $X$.
Soient $a, b : K \to L$ des morphismes d'hyperrecouvrements.
Il existe un morphisme d'hyperrecouvrements
$c : K' \to K$ tel que $a \circ c$ soit homotope
à $b \circ c$.
\end{lemma}

\begin{proof}
Considérons le diagramme commutatif suivant
$$
\xymatrix{
K' \ar@{=}[r]^-{def} \ar[rd]_c &
K \times_{(L \times L)} \Hom(\Delta[1], L)
\ar[r] \ar[d] & \Hom(\Delta[1], L) \ar[d] \\
& K \ar[r]^{(a, b)} & L \times L
}
$$
Par la propriété fonctorielle de $\Hom(\Delta[1], L)$,
le composé des morphismes horizontaux
correspond à un morphisme $K' \times \Delta[1] \to L$ qui
définit une homotopie entre $c \circ a$ et $c \circ b$.
Ainsi, si l'on peut montrer que $K'$ est un
hyperrecouvrement de $X$, on obtient le lemme.
Pour cela, nous appliquerons le Lemme \ref{hypercovering-lemma-funny-gamma}
à la paire de morphismes $K \to L \times L$
et $\Hom(\Delta[1], L) \to L \times L$.
La condition (1) du Lemme \ref{hypercovering-lemma-funny-gamma} est satisfaite.
La condition (2) du Lemme \ref{hypercovering-lemma-funny-gamma} est vérifiée parce que
$\Hom(\Delta[1], L)_0 = L_1$, et que le morphisme
$(d^1_0, d^1_1) : L_1 \to L_0 \times L_0$ est un
recouvrement de $\text{SR}(\mathcal{C}, X)$, puisque
$L$ est par hypothèse un hyperrecouvrement.
Pour démontrer la condition (3) du Lemme \ref{hypercovering-lemma-funny-gamma},
on utilise le Lemme \ref{hypercovering-lemma-hom-hypercovering} ci-dessus. D'après
ce lemme, le morphisme $\gamma$ de la condition (3) du Lemme
\ref{hypercovering-lemma-funny-gamma} est le morphisme
$$
\Hom(\Delta[1] \times \Delta[n + 1], L)_0
\longrightarrow
\Hom(U, L)_0
$$
où $U \subset \Delta[1] \times \Delta[n + 1]$.
D'après le Lemme \ref{hypercovering-lemma-add-simplices},
c'est un recouvrement, ce qui démontre l'assertion.
\end{proof}

\begin{remark}
\label{hypercovering-remark-contractible-category}
Remarquons que le point essentiel de la démonstration est l'emploi du
Lemme \ref{hypercovering-lemma-add-simplices}. Ce lemme
est tout à fait général et ne dépend pas de la
forme exacte des ensembles simpliciaux (pourvu qu'ils
n'aient qu'un nombre fini de simplexes non dégénérés).
Il paraît tout à fait raisonnable d'attendre un résultat
du type suivant :
Étant donné un morphisme $a : K \times \partial \Delta[k]
\to L$, où $K$ et $L$ sont des hyperrecouvrements, il
existe un morphisme d'hyperrecouvrements $c : K' \to K$
et un morphisme $g : K' \times \Delta[k] \to L$
tels que
$g|_{K' \times \partial \Delta[k]} =
a \circ (c \times \text{id}_{\partial \Delta[k]})$.
Autrement dit, la catégorie des hyperrecouvrements est
contractile en un sens convenable.
\end{remark}

















\section{Cohomologie et hyperrecouvrements}
\label{hypercovering-section-cohomology}

\noindent
Soit $\mathcal{C}$ un site admettant des produits fibrés.
Soit $X$ un objet de $\mathcal{C}$.
Soit $\mathcal{F}$ un faisceau de groupes abéliens sur $\mathcal{C}$.
Soient $K, L$ des hyperrecouvrements de $X$.
Si $a, b : K \to L$ sont des applications homotopes,
alors $\mathcal{F}(a), \mathcal{F}(b) : \mathcal{F}(K) \to \mathcal{F}(L)$
sont des applications homotopes ; voir
Simplicial, Lemme \ref{simplicial-lemma-functorial-homotopy}.
Elles ont donc le même effet sur les groupes de cohomologie des complexes
de cochaînes associés ; voir
Simplicial, Lemme \ref{simplicial-lemma-homotopy-s-Q}.
Nous allons nous en servir pour définir la limite inductive sur tous les
hyperrecouvrements.

\medskip\noindent
Notons temporairement $\text{HC}(\mathcal{C}, X)$
la catégorie dont les objets sont les hyperrecouvrements de $X$
et dont les morphismes sont les applications entre hyperrecouvrements de $X$
à homotopie près. Nous avons vu qu'il s'agit
d'une catégorie et non d'une « grande » catégorie ;
voir le Lemme \ref{hypercovering-lemma-hypercoverings-set}.
L'opposée de $\text{HC}(\mathcal{C}, X)$
sera la catégorie d'indices de notre diagramme ; voir
Catégories, section \ref{categories-section-limits}, pour la terminologie.
Considérons le diagramme
$$
\check{H}^i(-, \mathcal{F}) :
\text{HC}(\mathcal{C}, X)^{opp}
\longrightarrow
\textit{Ab}.
$$
D'après les Lemmes \ref{hypercovering-lemma-product-hypercoverings} et \ref{hypercovering-lemma-homotopy},
ainsi que la remarque ci-dessus sur les homotopies, ce diagramme est filtrant ; voir
Catégories, Définition \ref{categories-definition-directed}.
La limite inductive
$$
\check{H}^i_{\text{HC}}(X, \mathcal{F}) =
\colim_{K \in \text{HC}(\mathcal{C}, X)} \check{H}^i(K, \mathcal{F})
$$
admet donc une description particulièrement simple (voir la référence citée).

\begin{theorem}
\label{hypercovering-theorem-cohomology-hypercoverings}
Soit $\mathcal{C}$ un site admettant des produits fibrés.
Soit $X$ un objet de $\mathcal{C}$. Soit $i \geq 0$.
Les foncteurs
\begin{eqnarray*}
\textit{Ab}(\mathcal{C}) & \longrightarrow & \textit{Ab} \\
\mathcal{F} & \longmapsto & H^i(X, \mathcal{F}) \\
\mathcal{F} & \longmapsto & \check{H}^i_{\text{HC}}(X, \mathcal{F})
\end{eqnarray*}
sont canoniquement isomorphes.
\end{theorem}

\begin{proof}[Démonstration à l'aide des suites spectrales.]
Supposons que $\xi \in H^p(X, \mathcal{F})$ pour un certain $p \geq 0$.
Montrons que $\xi$ appartient à l'image de l'application
$\check{H}^p(X, \mathcal{F}) \to H^p(X, \mathcal{F})$ du
Lemme \ref{hypercovering-lemma-cech-spectral-sequence}
pour un certain hyperrecouvrement $K$ de $X$.

\medskip\noindent
C'est vrai si $p = 0$ d'après le Lemme \ref{hypercovering-lemma-h0-cech}.
Si $p = 1$, choisissons un hyperrecouvrement de {\v C}ech $K$ de $X$ comme dans
l'Exemple \ref{hypercovering-example-cech}, en partant d'un recouvrement
$K_0 = \{U_i \to X\}$ du site $\mathcal{C}$ tel que
$\xi|_{U_i} = 0$ ; voir
Cohomologie sur les sites,
Lemme \ref{sites-cohomology-lemma-kill-cohomology-class-on-covering}.
Il résulte immédiatement de la suite spectrale
du Lemme \ref{hypercovering-lemma-cech-spectral-sequence} que, dans ce cas, $\xi$ provient
d'un élément de $\check{H}^1(K, \mathcal{F})$.
Dans le cas général, choisissons un hyperrecouvrement quelconque $K$ de $X$
tel que l'image de $\xi$ soit nulle dans $\underline{H}^p(\mathcal{F})(K_0)$
(en utilisant l'Exemple \ref{hypercovering-example-cech} et
Cohomologie sur les sites,
Lemme \ref{sites-cohomology-lemma-kill-cohomology-class-on-covering}
une fois encore).
D'après la suite spectrale du Lemme \ref{hypercovering-lemma-cech-spectral-sequence},
l'obstruction à ce que $\xi$ provienne d'un élément de
$\check{H}^p(K, \mathcal{F})$ est une suite d'éléments
$\xi_1, \ldots, \xi_{p - 1}$ avec
$\xi_q \in \check{H}^{p - q}(K, \underline{H}^q(\mathcal{F}))$
(plus précisément, les images des $\xi_q$ dans certains sous-quotients
de ces groupes).

\medskip\noindent
On peut remplacer par récurrence l'hyperrecouvrement $K$ par des raffinements
de telle sorte que les restrictions des obstructions $\xi_1, \ldots, \xi_{p - 1}$ soient nulles
(et pas seulement leurs images
dans les sous-quotients : il n'y a donc ici aucune subtilité). En effet, supposons être
déjà parvenus à une situation où
$\xi_{q + 1}, \ldots, \xi_{p - 1}$ sont nuls.
Remarquons que $\xi_q \in \check{H}^{p - q}(K, \underline{H}^q(\mathcal{F}))$
est la classe d'un élément
$$
\tilde \xi_q \in
\underline{H}^q(\mathcal{F})(K_{p - q}) =
\prod H^q(U_i, \mathcal{F})
$$
si $K_{p - q} = \{U_i \to X\}_{i \in I}$. Soit $\xi_{q, i}$
la composante de $\tilde \xi_q$ dans $H^q(U_i, \mathcal{F})$.
Comme $q \geq 1$, on peut appliquer
Cohomologie sur les sites,
Lemme \ref{sites-cohomology-lemma-kill-cohomology-class-on-covering},
une fois encore, afin de choisir des recouvrements $\{U_{i, j} \to U_i\}$
du site tels que chaque restriction $\xi_{q, i}|_{U_{i, j}} = 0$.
Considérons l'objet $Z = \{U_{i, j} \to X\}$ de la catégorie
$\text{SR}(\mathcal{C}, X)$ et son morphisme évident
$u : Z \to K_{p - q}$. Il est clair que $u$ est un recouvrement ; voir
la Définition \ref{hypercovering-definition-covering-SR}. D'après le
Lemme \ref{hypercovering-lemma-covering}, il
existe un morphisme $L \to K$ d'hyperrecouvrements de $X$ tel que
$L_{p - q} \to K_{p - q}$ se factorise par $u$. Alors l'image de
$\xi_q$ dans $\underline{H}^q(\mathcal{F})(L_{p - q})$
est nulle. Puisque la suite spectrale du
Lemme \ref{hypercovering-lemma-cech-spectral-sequence}
est fonctorielle, cela signifie qu'après avoir remplacé $K$ par $L$, on atteint la
situation où $\xi_q, \ldots, \xi_{p - 1}$ sont tous nuls.
En poursuivant ainsi, on obtient un hyperrecouvrement où ils sont tous
nuls, et donc $\xi$ appartient à l'image de l'application
$\check{H}^p(X, \mathcal{F}) \to H^p(X, \mathcal{F})$.

\medskip\noindent
Supposons que $K$ soit un hyperrecouvrement de $X$, que
$\xi \in \check{H}^p(K, \mathcal{F})$ et que l'image de
$\xi$ par l'application
$\check{H}^p(X, \mathcal{F}) \to H^p(X, \mathcal{F})$ du
Lemme \ref{hypercovering-lemma-cech-spectral-sequence}
soit nulle. Pour achever la démonstration du théorème, il faut montrer qu'il
existe un morphisme d'hyperrecouvrements $L \to K$ tel que
la restriction de $\xi$ soit nulle dans $\check{H}^p(L, \mathcal{F})$.
D'après la suite spectrale du Lemme \ref{hypercovering-lemma-cech-spectral-sequence},
l'annulation de l'image de $\xi$ dans $H^p(X, \mathcal{F})$
signifie qu'il existe des éléments $\xi_1, \ldots, \xi_{p - 2}$
avec $\xi_q \in \check{H}^{p - 1 - q}(K, \underline{H}^q(\mathcal{F}))$
(plus précisément, leurs images dans certains sous-quotients)
tels que les images $d_{q + 1}^{p - 1 - q, q}\xi_q$ (dans la suite
spectrale) aient pour somme $\xi$. Par le même procédé que ci-dessus,
on peut donc trouver un morphisme d'hyperrecouvrements $L \to K$ tel que
les restrictions des éléments $\xi_q$, $q = 1, \ldots, p - 2$,
dans $\check{H}^{p - 1 - q}(L, \underline{H}^q(\mathcal{F}))$ soient nulles.
Il en résulte que $\xi$ est nul, puisque le morphisme $L \to K$
induit un morphisme de suites spectrales d'après le
Lemme \ref{hypercovering-lemma-cech-spectral-sequence}.
\end{proof}

\begin{proof}[Démonstration sans suites spectrales.]
Nous avons vu le résultat pour $i = 0$ ; voir le Lemme \ref{hypercovering-lemma-h0-cech}.
Nous savons que les foncteurs $H^i(X, -)$ forment un $\delta$-foncteur universel ; voir
Catégories dérivées, Lemme \ref{derived-lemma-higher-derived-functors}.
Pour démontrer le théorème, il suffit de montrer que
la suite de foncteurs $\check{H}^i_{HC}(X, -)$ forme un
$\delta$-foncteur. En effet, nous savons que la cohomologie de {\v C}ech
est nulle sur les faisceaux injectifs (Lemme \ref{hypercovering-lemma-injective-trivial-cech}),
et on peut alors appliquer
Homologie, Lemme \ref{homology-lemma-efface-implies-universal}.

\medskip\noindent
Soit
$$
0 \to \mathcal{F} \to \mathcal{G} \to \mathcal{H} \to 0
$$
une suite exacte courte de faisceaux abéliens sur $\mathcal{C}$.
Soit $\xi \in \check{H}^p_{HC}(X, \mathcal{H})$. Choisissons un hyperrecouvrement
$K$ de $X$ et un élément $\sigma \in \mathcal{H}(K_p)$ représentant
$\xi$ en cohomologie. Il existe une suite exacte correspondante de
complexes
$$
0 \to s(\mathcal{F}(K)) \to s(\mathcal{G}(K)) \to s(\mathcal{H}(K))
$$
mais rien ne garantit qu'il y ait aussi un zéro à droite, et c'est
la seule chose qui
nous empêche de définir $\delta(\xi)$ par une simple application du
lemme du serpent. Rappelons que
$$
\mathcal{H}(K_p) = \prod \mathcal{H}(U_i)
$$
si $K_p = \{U_i \to X\}$. Écrivons $\sigma =\prod \sigma_i$ avec
$\sigma_i \in \mathcal{H}(U_i)$. Comme $\mathcal{G} \to \mathcal{H}$ est
un morphisme surjectif de faisceaux, il existe des recouvrements
$\{U_{i, j} \to U_i\}$ tels que $\sigma_i|_{U_{i, j}}$ soit
l'image d'un élément $\tau_{i, j} \in \mathcal{G}(U_{i, j})$.
Considérons l'objet $Z = \{U_{i, j} \to X\}$ de la catégorie
$\text{SR}(\mathcal{C}, X)$ et son morphisme évident
$u : Z \to K_p$. Il est clair que $u$ est un recouvrement ; voir
la Définition \ref{hypercovering-definition-covering-SR}. D'après le
Lemme \ref{hypercovering-lemma-covering}, il
existe un morphisme $L \to K$ d'hyperrecouvrements de $X$ tel que
$L_p \to K_p$ se factorise par $u$. Après avoir remplacé $K$ par $L$,
on peut donc supposer que $\sigma$ est l'image d'un
élément $\tau \in \mathcal{G}(K_p)$. Remarquons que $d(\sigma) = 0$,
mais pas nécessairement que $d(\tau) = 0$. Ainsi, $d(\tau) \in \mathcal{F}(K_{p + 1})$
est un cocycle. Dans cette situation, on définit
$\delta(\xi)$ comme la classe du cocycle $d(\tau)$ dans
$\check{H}^{p + 1}_{HC}(X, \mathcal{F})$.

\medskip\noindent
Il faut maintenant vérifier plusieurs points :
(a) $\delta(\xi)$ ne dépend pas du choix de $\tau$,
(b) $\delta(\xi)$ ne dépend pas du choix de l'hyperrecouvrement
$L \to K$ sur lequel $\sigma$ se relève, et
(c) $\delta(\xi)$ ne dépend pas de l'hyperrecouvrement initial ni de
$\sigma$ choisi pour représenter $\xi$. Nous omettons la vérification de
(a), (b) et (c) ; l'indépendance des choix d'hyperrecouvrements
se ramène en fait aux Lemmes \ref{hypercovering-lemma-product-hypercoverings}
et \ref{hypercovering-lemma-homotopy}. Nous omettons aussi de vérifier que
$\delta$ est fonctoriel à l'égard des morphismes de suites exactes
courtes de faisceaux abéliens sur $\mathcal{C}$.

\medskip\noindent
Enfin, il faut vérifier qu'avec cette définition de $\delta$,
la suite exacte courte de faisceaux abéliens ci-dessus fournit une
suite exacte longue de groupes de cohomologie de {\v C}ech.
Montrons d'abord que si $\delta(\xi) = 0$ (avec $\xi$ comme ci-dessus), alors
$\xi$ est l'image d'un élément
$\xi' \in \check{H}^p_{HC}(X, \mathcal{G})$.
En effet, si $\delta(\xi) = 0$, alors, avec les notations précédentes, on
voit que la classe de $d(\tau)$ est nulle dans
$\check{H}^{p + 1}_{HC}(X, \mathcal{F})$. Il existe donc
un morphisme d'hyperrecouvrements $L \to K$ tel que la restriction
de $d(\tau)$ à un élément de $\mathcal{F}(L_{p + 1})$ soit
égale à $d(\upsilon)$ pour un certain $\upsilon \in \mathcal{F}(L_p)$.
Cela implique que $\tau|_{L_p} + \upsilon$ forme un
cocycle et détermine une classe $\xi' \in \check{H}^p(L, \mathcal{G})$
qui s'envoie sur $\xi$, comme voulu.

\medskip\noindent
Nous omettons la démonstration du fait que si $\xi' \in \check{H}^{p + 1}_{HC}(X, \mathcal{F})$
s'envoie sur zéro dans $\check{H}^{p + 1}_{HC}(X, \mathcal{G})$, alors il est
égal à $\delta(\xi)$ pour un certain $\xi \in \check{H}^p_{HC}(X, \mathcal{H})$.
\end{proof}

\noindent
Déduisons maintenant le cas de Verdier du
Théorème \ref{hypercovering-theorem-cohomology-hypercoverings}
par un artifice.

\begin{proposition}
\label{hypercovering-proposition-cohomology-hypercoverings}
Soit $\mathcal{C}$ un site admettant des produits fibrés et les produits de deux objets.
Soit $\mathcal{F}$ un faisceau abélien sur $\mathcal{C}$.
Soit $i \geq 0$. Alors
\begin{enumerate}
\item pour tout $\xi \in H^i(\mathcal{F})$, il existe un hyperrecouvrement
$K$ tel que $\xi$ appartienne à l'image de l'application canonique
$\check{H}^i(K, \mathcal{F}) \to H^i(\mathcal{F})$, et
\item si $K, L$ sont des hyperrecouvrements et $\xi_K \in \check{H}^i(K, \mathcal{F})$,
$\xi_L \in \check{H}^i(L, \mathcal{F})$ des éléments s'envoyant
sur le même élément de $H^i(\mathcal{F})$, alors il existe
un hyperrecouvrement $M$ et des morphismes $M \to K$ et $M \to L$ tels
que $\xi_K$ et $\xi_L$ s'envoient sur le même élément de
$\check{H}^i(M, \mathcal{F})$.
\end{enumerate}
Autrement dit, à des questions ensemblistes près, les groupes de
cohomologie de $\mathcal{F}$ sur $\mathcal{C}$ sont la limite inductive des
groupes de cohomologie de {\v C}ech de $\mathcal{F}$ sur tous les hyperrecouvrements.
\end{proposition}

\begin{proof}
Ce résultat est une conséquence immédiate du
Théorème \ref{hypercovering-theorem-cohomology-hypercoverings}.
En effet, on peut remplacer artificiellement $\mathcal{C}$ par un site
légèrement plus grand $\mathcal{C}'$ tel que
(I) $\mathcal{C}'$ possède un objet final $X$ et (II) les
hyperrecouvrements dans $\mathcal{C}$ soient essentiellement la
même chose que les hyperrecouvrements de $X$ dans $\mathcal{C}'$.
Mais la nature de la construction impose un assez grand nombre de
vérifications.

\medskip\noindent
Appelons une famille de morphismes $\{U_i \to U\}$ de $\mathcal{C}$
de but fixé un {\it recouvrement faible} si l'application
$\coprod_{i \in I} h_{U_i} \to h_U$ devient surjective après passage aux faisceaux associés.
Construisons un nouveau site $\mathcal{C}'$ comme suit :
\begin{enumerate}
\item comme catégorie, on pose $\Ob(\mathcal{C}') = \Ob(\mathcal{C}) \amalg \{X\}$
et on ajoute un unique morphisme vers $X$ depuis tout objet de $\mathcal{C}'$,
\item $\mathcal{C}'$ admet des produits fibrés dès que les produits fibrés et les produits
de deux objets existent dans $\mathcal{C}$,
\item les recouvrements de $\mathcal{C}'$ sont les recouvrements faibles de $\mathcal{C}$,
ainsi que les $\{U_i \to X\}_{i \in I}$ tels que soit $U_i = X$
pour un certain $i$, soit $U_i \not = X$ pour tout $i$ et que l'application
$\coprod h_{U_i} \to *$ de préfaisceaux sur $\mathcal{C}$ devienne
surjective après passage aux faisceaux associés sur $\mathcal{C}$,
\item on applique Ensembles, Lemme \ref{sets-lemma-coverings-site},
pour restreindre les recouvrements et obtenir notre site $\mathcal{C}'$.
\end{enumerate}
Alors $\Sh(\mathcal{C}') = \Sh(\mathcal{C})$, car le foncteur d'inclusion
$\mathcal{C} \to \mathcal{C}'$ est spécialement cocontinu
(voir Sites, Définition \ref{sites-definition-special-cocontinuous-functor}).
Nous omettons les vérifications immédiates.

\medskip\noindent
Choisissons un recouvrement $\{U_i \to X\}$ de $\mathcal{C}'$ tel que $U_i$ soit un
objet de $\mathcal{C}$ pour tout $i$ (cela est possible puisque
$\mathcal{C} \to \mathcal{C}'$ est spécialement cocontinu).
Alors $K_0 = \{U_i \to X\}$ est un recouvrement dans le
site $\mathcal{C}'$ construit ci-dessus. Considérons $K_0$ comme un objet de
$\text{SR}(\mathcal{C}', X)$ et posons $K_{init} = \text{cosk}_0(K_0)$.
Alors $K_{init}$ est un hyperrecouvrement de $X$ ; voir
l'Exemple \ref{hypercovering-example-cech}. Remarquons que chaque $K_{init, n}$ est de la forme
$\{W_j \to X\}$ avec $W_j \in \Ob(\mathcal{C})$.

\medskip\noindent
Démonstration de (1). Choisissons $\xi \in H^i(\mathcal{F}) = H^i(X, \mathcal{F}')$,
où $\mathcal{F}'$ est le faisceau abélien sur $\mathcal{C}'$ correspondant
à $\mathcal{F}$ sur $\mathcal{C}$. D'après le
Théorème \ref{hypercovering-theorem-cohomology-hypercoverings},
il existe un morphisme d'hyperrecouvrements $K' \to K_{init}$
de $X$ dans $\mathcal{C}'$ tel que $\xi$ provienne d'un élément
de $\check{H}^i(K', \mathcal{F})$.
Écrivons $K'_n = \{U_{n, j} \to X\}$. Comme $K'_n$ s'envoie dans
$K_{init, n}$, on voit que $U_{n, j}$ est un objet de $\mathcal{C}$.
On peut donc définir un objet simplicial $K$ de $\text{SR}(\mathcal{C})$
en posant $K_n = \{U_{n, j}\}$. Puisque les recouvrements de
$\mathcal{C}'$ formés de familles de morphismes de $\mathcal{C}$
sont des recouvrements faibles, on voit que $K$ est un hyperrecouvrement au sens
de la Définition \ref{hypercovering-definition-hypercovering-variant}.
Enfin, puisque $\mathcal{F}'$ est l'unique faisceau sur $\mathcal{C}'$
dont la restriction à $\mathcal{C}$ est égale à $\mathcal{F}$,
on voit que les complexes de {\v C}ech $s(\mathcal{F}(K))$
et $s(\mathcal{F}'(K'))$ sont identiques, ce qui démontre (1).
(La compatibilité avec l'application vers les groupes de cohomologie est omise.)

\medskip\noindent
Démonstration de (2). Soient $K$ et $L$ des hyperrecouvrements dans $\mathcal{C}$.
Soient $K'$ et $L'$ les objets simpliciaux de $\text{SR}(\mathcal{C}', X)$
obtenus à partir de $K$ et $L$ par le foncteur
$\text{SR}(\mathcal{C}) \to \text{SR}(\mathcal{C}', X)$,
$\{U_i\} \mapsto \{U_i \to X\}$. Comme précédemment, les complexes de
{\v C}ech sont égaux ; on obtient donc $\xi_{K'}$ et
$\xi_{L'}$ qui s'envoient sur la même classe de cohomologie de $\mathcal{F}'$
sur $\mathcal{C}'$. Après avoir éventuellement élargi notre choix
de recouvrements dans $\mathcal{C}'$ (en raison d'une question ensembliste),
on peut supposer que $K'$ et $L'$ sont des hyperrecouvrements de $X$ dans
$\mathcal{C}'$ ; cela résulte de notre définition des hyperrecouvrements dans la
Définition \ref{hypercovering-definition-hypercovering-variant} et
du fait que les recouvrements faibles de $\mathcal{C}$ donnent des recouvrements dans
$\mathcal{C}'$. D'après le
Théorème \ref{hypercovering-theorem-cohomology-hypercoverings},
il existe un hyperrecouvrement $M'$ de $X$ dans $\mathcal{C}'$
et des morphismes $M' \to K'$, $M' \to L'$ et $M' \to K_{init}$
tels que les restrictions de $\xi_{K'}$ et $\xi_{L'}$ soient le même élément de
$\check{H}^i(M', \mathcal{F})$. En explicitant cet énoncé comme ci-dessus,
on obtient (2).
\end{proof}





\section{Hyperrecouvrements d'espaces}
\label{hypercovering-section-hypercoverings-spaces}

\noindent
La théorie précédente présente déjà quelque intérêt dans le cas des espaces
topologiques. Dans ce cas, on peut expliciter ce qu'est un hyperrecouvrement et voir
ce que dit effectivement le résultat.

\medskip\noindent
Soit $X$ un espace topologique. Considérons le site $X_{Zar}$
de Sites, Exemple \ref{sites-example-site-topological}. Rappelons qu'un
objet de $X_{Zar}$ est simplement un ouvert de $X$ et que les morphismes
de $X_{Zar}$ correspondent simplement aux inclusions. Qu'est donc un
hyperrecouvrement de $X$ pour le site $X_{Zar}$ ?

\medskip\noindent
Explicitons d'abord la Définition \ref{hypercovering-definition-SR}.
Un objet de $\text{SR}(X_{Zar}, X)$ est simplement donné par un ensemble
$I$ et, pour chaque $i \in I$, un ouvert $U_i \subset X$.
Notons-le $\{U_i\}_{i \in I}$, puisqu'aucune
confusion n'est possible au sujet du morphisme $U_i \to X$.
Un morphisme $\{U_i\}_{i \in I} \to \{V_j\}_{j \in J}$
entre deux tels objets est donné par une application d'ensembles
$\alpha : I \to J$ telle que $U_i \subset V_{\alpha(i)}$ pour tout
$i \in I$. Quand un tel morphisme est-il un recouvrement ? C'est le cas
si et seulement si, pour tout $j \in J$, on a
$V_j = \bigcup_{i\in I, \ \alpha(i) = j} U_i$ (et que la famille est
un recouvrement dans le site $X_{Zar}$).

\medskip\noindent
Ce qui précède donne la description suivante d'un hyperrecouvrement
dans le site $X_{Zar}$. Un hyperrecouvrement de $X$ dans $X_{Zar}$
est donné par les données suivantes :
\begin{enumerate}
\item un ensemble simplicial $I$ (voir
Simplicial, section \ref{simplicial-section-simplicial-set}), et
\item pour tout $n \geq 0$ et tout $i \in I_n$, un ouvert $U_i \subset X$.
\end{enumerate}
Nous noterons une telle collection de données $(I, \{U_i\})$.
Pour que ce soit un hyperrecouvrement de $X$, on exige
les propriétés suivantes :
\begin{itemize}
\item pour $i \in I_n$ et $0 \leq a \leq n$,
on a $U_i \subset U_{d^n_a(i)}$,
\item pour $i \in I_n$ et $0 \leq a \leq n$, on a $U_i = U_{s^n_a(i)}$,
\item on a
\begin{equation}
\label{hypercovering-equation-covering-X}
X = \bigcup\nolimits_{i \in I_0} U_i,
\end{equation}
\item pour tous $i_0, i_1 \in I_0$, on a
\begin{equation}
\label{hypercovering-equation-covering-two}
U_{i_0} \cap U_{i_1} =
\bigcup\nolimits_{i \in I_1, \ d^1_0(i) = i_0, \ d^1_1(i) = i_1} U_i,
\end{equation}
\item pour tout $n \geq 1$ et tout
$(i_0, \ldots, i_{n + 1}) \in (I_n)^{n + 2}$ tels que
$d^n_{b - 1}(i_a) = d^n_a(i_b)$ pour tous $0\leq a < b\leq n + 1$,
on a
\begin{equation}
\label{hypercovering-equation-covering-general}
U_{i_0} \cap \ldots \cap U_{i_{n + 1}} =
\bigcup\nolimits_{i \in I_{n + 1},
\ d^{n + 1}_a(i) = i_a, \ a = 0, \ldots, n + 1} U_i,
\end{equation}
\item chacun des recouvrements ouverts (\ref{hypercovering-equation-covering-X}),
(\ref{hypercovering-equation-covering-two}) et (\ref{hypercovering-equation-covering-general})
est un élément de $\text{Cov}(X_{Zar})$
(c'est une condition ensembliste qui borne
la taille des ensembles d'indices des recouvrements).
\end{itemize}
Les conditions (\ref{hypercovering-equation-covering-X}) et
(\ref{hypercovering-equation-covering-two}) devraient être familières, par exemple d'après le
chapitre consacré aux faisceaux sur les espaces, et la condition
(\ref{hypercovering-equation-covering-general}) en est la généralisation naturelle.

\begin{remark}
\label{hypercovering-remark-not-covering-set}
Une particularité de cette description est que, si l'une des intersections
multiples $U_{i_0} \cap \ldots \cap U_{i_{n + 1}}$ est vide, alors
le recouvrement du membre de droite peut être le recouvrement vide.
Ainsi, les applications
$I_{n + 1} \to (\text{cosk}_n\text{sk}_n I)_{n + 1}$ ne sont pas nécessairement surjectives.
Cela signifie que la réalisation géométrique de $I$ peut être un espace
intéressant (non contractile).

\medskip\noindent
En effet, soit $I'_n \subset I_n$ le sous-ensemble
formé des simplexes $i \in I_n$ tels que
$U_i \not = \emptyset$. On vérifie aisément que $I' \subset I$
est un sous-ensemble simplicial et que $(I', \{U_i\})$ est un hyperrecouvrement.
On peut donc toujours raffiner un hyperrecouvrement en un hyperrecouvrement où
aucun des ouverts $U_i$ n'est vide.
\end{remark}

\begin{remark}
\label{hypercovering-remark-repackage-into-simplicial-space}
Reformulons encore une fois ces informations.
Supposons donc que $(I, \{U_i\})$ soit un hyperrecouvrement de
l'espace topologique $X$. À partir de ces données, on peut construire
un espace topologique simplicial $U_\bullet$ en posant
$$
U_n = \coprod\nolimits_{i \in I_n} U_i,
$$
et, pour $\varphi : [n] \to [m]$ donné, on définit
les morphismes $U(\varphi) : U_n \to U_m$ par les
inclusions $U_i \subset U_{\varphi(i)}$
pour $i \in I_n$. Cet espace topologique simplicial est
muni d'une augmentation $\epsilon : U_\bullet \to X$.
Avec ce morphisme, l'espace simplicial $U_\bullet$ devient
un hyperrecouvrement de $X$ le long duquel on a descente cohomologique
au sens de \cite[Expos\'e Vbis]{SGA4}.
Autrement dit, $H^n(U_\bullet, \epsilon^*\mathcal{F}) = H^n(X, \mathcal{F})$.
(Insérer ici une référence future à la cohomologie sur les espaces
simpliciaux et à la descente cohomologique formulée en ces termes.)
Supposons que $\mathcal{F}$ soit un faisceau abélien sur $X$.
Dans ce cas, la suite spectrale du Lemme \ref{hypercovering-lemma-cech-spectral-sequence}
devient la suite spectrale dont le terme $E_1$ est
$$
E_1^{p, q} = H^q(U_p, \epsilon_q^*\mathcal{F})
\Rightarrow
H^{p + q}(U_\bullet, \epsilon^*\mathcal{F}) = H^{p + q}(X, \mathcal{F})
$$
qui compare la cohomologie totale de $\epsilon^*\mathcal{F}$
aux groupes de cohomologie de $\mathcal{F}$ sur les morceaux
de $U_\bullet$. (Insérer une référence future à cette suite spectrale
ici.)
\end{remark}

\noindent
En topologie, on cherche souvent des hyperrecouvrements de $X$ qui
possèdent la propriété que tous les $U_i$ appartiennent à une base donnée de la topologie
de $X$ et que tous les recouvrements
(\ref{hypercovering-equation-covering-two}) et (\ref{hypercovering-equation-covering-general})
appartiennent à une collection cofinale donnée de recouvrements.
Voici deux lemmes qui en fournissent des exemples.

\begin{lemma}
\label{hypercovering-lemma-basis-hypercovering}
Soit $X$ un espace topologique.
Soit $\mathcal{B}$ une base de la topologie de $X$.
Il existe un hyperrecouvrement $(I, \{U_i\})$ de $X$
tel que chaque $U_i$ soit un élément de $\mathcal{B}$.
\end{lemma}

\begin{proof}
Soit $n \geq 0$.
Appelons {\it hyperrecouvrement $n$-tronqué} de $X$ une donnée
formée d'un ensemble simplicial $n$-tronqué $I$ et, pour chaque
$i \in I_a$, $0 \leq a \leq n$, d'un ouvert $U_i$ de $X$, de telle sorte que
les conditions définissant un hyperrecouvrement soient satisfaites lorsqu'elles ont un sens.
Autrement dit, on n'exige les relations d'inclusion et les conditions de
recouvrement que lorsque tous les simplexes qui y interviennent
sont des $a$-simplexes avec $a \leq n$. Le lemme découle de l'assertion
suivante : tout hyperrecouvrement $n$-tronqué $(I, \{U_i\})$ tel que
tous les $U_i \in \mathcal{B}$ se prolonge en un hyperrecouvrement $(n + 1)$-tronqué
sans ajout de $a$-simplexes pour $a \leq n$.
Procédons comme suit. Considérons d'abord l'ensemble simplicial
$(n + 1)$-tronqué $I'$ défini par
$I' = \text{sk}_{n + 1}(\text{cosk}_n I)$.
Rappelons que
$$
I'_{n + 1} =
\left\{
\begin{matrix}
(i_0, \ldots, i_{n + 1}) \in (I_n)^{n + 2} \text{ tels que}\\
d^n_{b - 1}(i_a) = d^n_a(i_b) \text{ pour tous }0\leq a < b\leq n + 1
\end{matrix}
\right\}
$$
Si $i' \in I'_{n + 1}$ est dégénéré, disons $i' = s^n_a(i)$, on pose
$U_{i'} = U_i$ (c'est de toute façon imposé par la deuxième condition).
On pose aussi $J_{i'} = \{i'\}$ dans ce cas.
Si $i' \in I'_{n + 1}$ est non dégénéré, disons
$i' = (i_0, \ldots, i_{n + 1})$, on choisit un ensemble
$J_{i'}$ et un recouvrement ouvert
\begin{equation}
\label{hypercovering-equation-choose-covering}
U_{i_0} \cap \ldots \cap U_{i_{n + 1}} =
\bigcup\nolimits_{i \in J_{i'}} U_i,
\end{equation}
avec $U_i \in \mathcal{B}$ pour $i \in J_{i'}$.
Posons
$$
I_{n + 1} = \coprod\nolimits_{i' \in I'_{n + 1}} J_{i'}
$$
Il existe une application canonique $\pi : I_{n + 1} \to I'_{n + 1}$ qui est
une bijection au-dessus de l'ensemble des simplexes dégénérés de $I'_{n + 1}$ par
construction.
Pour $i \in I_{n + 1}$, on définit $d^{n + 1}_a(i) = d^{n + 1}_a(\pi(i))$.
Pour $i \in I_n$, on définit $s^n_a(i) \in I_{n + 1}$ comme l'unique
simplexe situé au-dessus du simplexe dégénéré $s^n_a(i) \in I'_{n + 1}$.
Nous omettons la vérification que cela définit un hyperrecouvrement
$(n + 1)$-tronqué de $X$.
\end{proof}

\begin{lemma}
\label{hypercovering-lemma-quasi-separated-quasi-compact-hypercovering}
Soit $X$ un espace topologique.
Soit $\mathcal{B}$ une base de la topologie de $X$.
Supposons que
\begin{enumerate}
\item $X$ soit quasi-compact,
\item chaque $U \in \mathcal{B}$ soit un ouvert quasi-compact, et
\item l'intersection de deux ouverts quasi-compacts quelconques de
$X$ soit quasi-compacte.
\end{enumerate}
Alors il existe un hyperrecouvrement $(I, \{U_i\})$ de $X$ ayant les
propriétés suivantes :
\begin{enumerate}
\item chaque $U_i$ est un élément de la base $\mathcal{B}$,
\item chacun des $I_n$ est un ensemble fini, et, en particulier,
\item chacun des recouvrements (\ref{hypercovering-equation-covering-X}),
(\ref{hypercovering-equation-covering-two}) et (\ref{hypercovering-equation-covering-general})
est fini.
\end{enumerate}
\end{lemma}

\begin{proof}
Cela résulte directement de la construction donnée dans la démonstration du
Lemme \ref{hypercovering-lemma-basis-hypercovering}, si l'on choisit des recouvrements finis
par des éléments de $\mathcal{B}$ dans (\ref{hypercovering-equation-choose-covering}).
Nous omettons les détails.
\end{proof}







\section{Construction d'hyperrecouvrements}
\label{hypercovering-section-hypercovering-sites}

\noindent
Soit $\mathcal{C}$ un site. Dans cette section, nous considérerons un
objet simplicial de $\text{SR}(\mathcal{C})$ de la manière suivante.
Comme d'habitude, on pose $K_n = K([n])$ et on note $K(\varphi) : K_n \to K_m$
le morphisme associé à $\varphi : [m] \to [n]$.
On peut écrire $K_n = \{U_{n, i}\}_{i \in I_n}$. Pour
$\varphi : [m] \to [n]$, le morphisme $K(\varphi) : K_n \to K_m$
est donné par une application $\alpha(\varphi) : I_n \to I_m$ et des morphismes
$f_{\varphi, i} : U_{n, i} \to U_{m, \alpha(\varphi)(i)}$
pour $i \in I_n$. Le fait que $K$ soit un objet simplicial de
$\text{SR}(\mathcal{C})$ implique que $(I_n, \alpha(\varphi))$
est un ensemble simplicial
et que $f_{\psi, \alpha(\varphi)(i)} \circ f_{\varphi, i} =
f_{\varphi \circ \psi, i}$ lorsque $\psi : [l] \to [m]$.

\begin{lemma}
\label{hypercovering-lemma-split}
Soit $\mathcal{C}$ un site. Soit $K$ un objet simplicial $r$-tronqué
de $\text{SR}(\mathcal{C})$. Les conditions suivantes sont équivalentes :
\begin{enumerate}
\item $K$ est scindé (Simplicial, Définition \ref{simplicial-definition-split}),
\item $f_{\varphi, i} : U_{n, i} \to U_{m, \alpha(\varphi)(i)}$
est un isomorphisme pour $r \geq n \geq 0$,
$\varphi : [m] \to [n]$ surjective, $i \in I_n$, et
\item $f_{\sigma^n_j, i} : U_{n, i} \to U_{n + 1, \alpha(\sigma^n_j)(i)}$
est un isomorphisme pour $0 \leq j \leq n < r$, $i \in I_n$.
\end{enumerate}
Le même énoncé vaut pour les objets simpliciaux si, dans (2) et (3),
on pose $r = \infty$.
\end{lemma}

\begin{proof}
Le scindage d'un ensemble simplicial est unique et donné par
les indices non dégénérés $N(I_n)$ en chaque degré $n$ ; voir
Simplicial, Lemme \ref{simplicial-lemma-splitting-simplicial-sets}.
Le coproduit de deux objets $\{U_i\}_{i \in I}$ et $\{U_j\}_{j \in J}$
de $\text{SR}(\mathcal{C})$ est donné par $\{U_l\}_{l \in I \amalg J}$,
avec les notations évidentes. Un scindage de $K$ doit donc être donné par
$N(K_n) = \{U_i\}_{i \in N(I_n)}$. L'équivalence de (1) et (2)
résulte alors de l'explicitation des définitions. L'équivalence de (2)
et (3) résulte du fait que toute surjection
$\varphi : [m] \to [n]$ est un composé de morphismes
$\sigma^k_j$ avec $k = n, n + 1, \ldots, m - 1$.
\end{proof}

\begin{lemma}
\label{hypercovering-lemma-hypercovering-object}
Soit $\mathcal{C}$ un site admettant des produits fibrés.
Soit $\mathcal{B} \subset \Ob(\mathcal{C})$ un sous-ensemble.
Supposons que
\begin{enumerate}
\item tout objet $U$ de $\mathcal{C}$ admette un recouvrement
$\{U_j \to U\}_{j \in J}$ avec $U_j \in \mathcal{B}$, et
\item si $\{U_j \to U\}_{j \in J}$ est un recouvrement
avec $U_j \in \mathcal{B}$ et si $\{U' \to U\}$ est un morphisme avec
$U' \in \mathcal{B}$, alors $\{U_j \to U\}_{j \in J} \amalg \{U' \to U\}$
soit un recouvrement.
\end{enumerate}
Alors, pour tout $X$ de $\mathcal{C}$, il existe un hyperrecouvrement $K$
de $X$ tel que $K_n = \{U_{n, i}\}_{i \in I_n}$
avec $U_{n, i} \in \mathcal{B}$ pour tout $i \in I_n$.
\end{lemma}

\begin{proof}
La démonstration du Lemme \ref{hypercovering-lemma-basis-hypercovering}
constitue un échauffement pour celle-ci, et
nous encourageons le lecteur à la lire d'abord.

\medskip\noindent
On commence par remplacer $\mathcal{C}$ par le site $\mathcal{C}/X$.
Après ce remplacement, on peut supposer que $X$ est l'objet final
de $\mathcal{C}$ et que $\mathcal{C}$ admet toutes les limites finies
(Catégories, Lemme \ref{categories-lemma-finite-limits-exist}).

\medskip\noindent
Soit $n \geq 0$. Appelons
{\it hyperrecouvrement $n$-tronqué de type $\mathcal{B}$ de $X$}
un objet simplicial $n$-tronqué $K$
de $\text{SR}(\mathcal{C})$
tel que, pour $i \in I_a$, $0 \leq a \leq n$,
on ait $U_{a, i} \in \mathcal{B}$, que
$K_0$ soit un recouvrement de $X$ et que
$K_{a + 1} \to (\text{cosk}_a \text{sk}_a K)_{a + 1}$
soit un recouvrement pour $a = 0, \ldots, n - 1$,
comme dans la Définition \ref{hypercovering-definition-covering-SR}.

\medskip\noindent
Puisque, par hypothèse, $X$ admet un recouvrement $\{U_{0, i} \to X\}_{i \in I_0}$
avec $U_i \in \mathcal{B}$, on obtient un hyperrecouvrement
$0$-tronqué de type $\mathcal{B}$ de $X$. Observons que tout hyperrecouvrement
$0$-tronqué de type $\mathcal{B}$ de $X$ est scindé ; voir le
Lemme \ref{hypercovering-lemma-split}.

\medskip\noindent
Le lemme découle de l'assertion suivante, pour $n \geq 0$ : étant donné un
hyperrecouvrement $n$-tronqué de type $\mathcal{B}$, noté $K$, de $X$ et scindé,
on peut le prolonger en un
hyperrecouvrement $(n + 1)$-tronqué de type $\mathcal{B}$, scindé, de $X$.

\medskip\noindent
Construction du prolongement. Considérons l'objet simplicial $(n + 1)$-tronqué
$K' = \text{sk}_{n + 1}(\text{cosk}_n K)$ de $\text{SR}(\mathcal{C})$.
Écrivons
$$
K'_{n + 1} = \{U'_{n + 1, i}\}_{i \in I'_{n + 1}}
$$
Puisque $K = \text{sk}_n K'$, on a $K_a = K'_a$ pour $0 \leq a \leq n$.
Pour tout $i' \in I'_{n + 1}$, choisissons un recouvrement
\begin{equation}
\label{hypercovering-equation-choose-covering-B}
\{g_{n + 1, j} : U_{n + 1, j} \to U'_{n + 1, i'}\}_{j \in J_{i'}}
\end{equation}
avec $U_{n + 1, j} \in \mathcal{B}$ pour $j \in J_{i'}$.
C'est possible d'après l'hypothèse du lemme sur $\mathcal{B}$.
Pour $0 \leq m \leq n$, notons $N_m \subset I_m$ le sous-ensemble des
indices non dégénérés. Posons
$$
I_{n + 1} =
\coprod\nolimits_{\varphi : [n + 1] \to [m]\text{ surjective, }0\leq m \leq n}
N_m \amalg
\coprod\nolimits_{i' \in I'_{n + 1}} J_{i'}
$$
Pour $j \in I_{n + 1}$, posons
$$
U_{n + 1, j} =
\left\{
\begin{matrix}
U_{m, i} & \text{si} &
j = (\varphi, i) & \text{où} & \varphi : [n + 1] \to [m], i \in N_m \\
U_{n + 1, j} & \text{si} &
j \in J_{i'} & \text{où} & i' \in I'_{n + 1}
\end{matrix}
\right.
$$
avec les notations évidentes. Posons $K_{n + 1} = \{U_{n + 1, j}\}_{j \in I_{n + 1}}$.
Par construction, $U_{n + 1, j}$ est un élément
de $\mathcal{B}$ pour tout $j \in I_{n + 1}$. Définissons des
applications compatibles
$$
I_{n + 1} \to I'_{n + 1}
\quad\text{et}\quad
K_{n + 1} \to K'_{n + 1}
$$
La première application est donnée par
$(\varphi, i) \mapsto \alpha'(\varphi)(i)$ et
$(j \in J_{i'}) \mapsto i'$.
Pour la deuxième, on utilise les morphismes
$$
f'_{\varphi, i} : U_{m, i} \to U'_{n + 1, \alpha'(\varphi)(i)}
\quad\text{et}\quad
g_{n + 1, j} : U_{n + 1, j} \to U'_{n + 1, i'}
$$
Nous affirmons que le morphisme
$$
K_{n + 1} \to K'_{n + 1} =
(\text{cosk}_n \text{sk}_n K')_{n + 1} =
(\text{cosk}_n K)_{n + 1}
$$
est un recouvrement au sens de la Définition \ref{hypercovering-definition-covering-SR}.
En effet, si $i' \in I'_{n + 1}$, alors soit $i'$ est non dégénéré
et l'image inverse de $i'$ dans $I_{n + 1}$ est égale à $J_{i'}$,
auquel cas on obtient un recouvrement de $U'_{n + 1, i'}$ par notre choix
(\ref{hypercovering-equation-choose-covering-B}), soit $i'$ est dégénéré et
l'image inverse de $i'$ dans $I_{n + 1}$ est
$J_{i'} \amalg \{(\varphi, i)\}$ pour un unique couple $(\varphi, i)$,
auquel cas notre choix (\ref{hypercovering-equation-choose-covering-B})
et l'hypothèse (2) du lemme donnent un recouvrement.

\medskip\noindent
Pour achever la démonstration, il faut définir les morphismes
$K(\varphi) : K_{n + 1} \to K_m$ correspondant aux morphismes
$\varphi : [m] \to [n + 1]$, $0 \leq m \leq n$, ainsi que les morphismes
$K(\varphi) : K_m \to K_{n + 1}$ correspondant aux morphismes
$\varphi : [n + 1] \to [m]$, $0 \leq m \leq n$,
de manière à satisfaire les relations de composition convenables.
Pour le premier type, on utilise le composé
$$
K_{n + 1} \to K'_{n + 1} \xrightarrow{K'(\varphi)} K'_m = K_m
$$
pour définir $K(\varphi) : K_{n + 1} \to K_m$.
Pour le second type, supposons donné $\varphi : [n + 1] \to [m]$,
$0 \leq m \leq n$. On définit le morphisme correspondant
$K(\varphi) : K_m \to K_{n + 1}$ comme suit :
\begin{enumerate}
\item pour $i \in I_m$, il existe une unique application surjective
$\psi : [m] \to [m_0]$ et un unique $i_0 \in I_{m_0}$ non dégénéré
tels que $\alpha(\psi)(i_0) = i$\footnote{Par exemple, si $i$ est
non dégénéré, alors $m = m_0$ et $\psi = \text{id}_{[m]}$.},
\item on pose $\varphi_0 = \psi_0 \circ \varphi : [n + 1] \to [m_0]$
et on envoie
$i \in I_m$ sur $(\varphi_0, i_0) \in I_{n + 1}$ ; autrement dit,
$\alpha(\varphi)(i) = (\varphi_0, i_0)$, et
\item le morphisme
$f_{\varphi, i} : U_{m, i} \to U_{n + 1, \alpha(\varphi)(i)} = U_{m_0, i_0}$
est l'inverse de l'isomorphisme $f_{\psi, i_0} : U_{m_0, i_0} \to U_{m, i}$
(voir le Lemme \ref{hypercovering-lemma-split}).
\end{enumerate}
Nous omettons la vérification immédiate mais laborieuse que cela définit
un hyperrecouvrement $(n + 1)$-tronqué de type $\mathcal{B}$, scindé, de $X$
prolongeant celui qui était donné en degré $n$. Tout est en fait clair
d'après ce qui précède, sauf la vérification que les morphismes
$K(\varphi)$ se composent correctement pour tout $\varphi : [a] \to [b]$
avec $0 \leq a, b \leq n + 1$.
\end{proof}

\begin{lemma}
\label{hypercovering-lemma-hypercovering-site}
Soit $\mathcal{C}$ un site admettant des égalisateurs et des produits fibrés.
Soit $\mathcal{B} \subset \Ob(\mathcal{C})$ un sous-ensemble. Supposons
que tout objet de $\mathcal{C}$ possède un recouvrement
dont les membres sont des éléments de $\mathcal{B}$.
Alors il existe un hyperrecouvrement $K$ tel que
$K_n = \{U_i\}_{i \in I_n}$ avec $U_i \in \mathcal{B}$
pour tout $i \in I_n$.
\end{lemma}

\begin{proof}
Cette démonstration est presque identique à celle du
Lemme \ref{hypercovering-lemma-hypercovering-object}. Nous n'en
expliquerons que les différences.

\medskip\noindent
Soit $n \geq 1$. Appelons
{\it hyperrecouvrement $n$-tronqué de type $\mathcal{B}$}
un objet simplicial
$n$-tronqué $K$ de $\text{SR}(\mathcal{C})$
tel que, pour $i \in I_a$, $0 \leq a \leq n$,
on ait $U_{a, i} \in \mathcal{B}$ et que
\begin{enumerate}
\item $F(K_0)^\# \to *$ soit surjectif,
\item $F(K_1)^\# \to F(K_0)^\# \times F(K_0)^\#$ soit surjectif,
\item $F(K_{a + 1})^\# \to F((\text{cosk}_a \text{sk}_a K)_{a + 1})^\#$
soit surjectif pour $a = 1, \ldots, n - 1$.
\end{enumerate}
Construisons d'abord explicitement un hyperrecouvrement $1$-tronqué de type $\mathcal{B}$ et scindé.

\medskip\noindent
Prenons $I_0 = \mathcal{B}$ et $K_0 = \{U\}_{U \in \mathcal{B}}$.
Alors (1) est vérifiée d'après notre hypothèse sur $\mathcal{B}$. Posons
$$
\Omega =
\{(U, V, W, a, b) \mid U, V, W \in \mathcal{B}, a : U \to V, b : U \to W\}
$$
Puis posons $I_1 = I_0 \amalg \Omega$. Pour $i \in I_1$, on pose
$U_{1, i} = U_{0, i}$ si $i \in I_0$, et $U_{1, i} = U$
si $i = (U, V, W, a, b) \in \Omega$. L'application
$K(\sigma^0_0) : K_0 \to K_1$ correspond à
l'inclusion $\alpha(\sigma^0_0) : I_0 \to I_1$
et à l'identité $f_{\sigma^0_0, i} : U_{0, i} \to U_{1, i}$
sur les objets. Les applications $K(\delta^1_0), K(\delta^1_1) : K_1 \to K_0$
correspondent aux deux applications $I_1 \to I_0$ qui sont
l'identité sur $I_0 \subset I_1$ et envoient $(U, V, W, a, b) \in \Omega \subset I_1$
sur $V$, resp.\ $W$. Les morphismes correspondants
$f_{\delta^1_0, i}, f_{\delta^1_1, i} : U_{1, i} \to U_{0, i}$ sont
l'identité si $i \in I_0$, et $a, b$ si $i = (U, V, W, a, b) \in \Omega$.
La condition (2) est vérifiée parce que toute section de
$F(K_0)^\# \times F(K_0)^\#$ sur un objet $U$ de $\mathcal{C}$
provient, après remplacement de $U$ par les membres d'un recouvrement,
d'une application $U \to F(K_0) \times F(K_0)$.
Cela signifie à son tour qu'on dispose de $V, W \in \mathcal{B}$
et de deux morphismes $U \to V$ et $U \to W$. En remplaçant encore
$U$ par les membres d'un recouvrement, on peut supposer que $U \in \mathcal{B}$,
comme voulu.

\medskip\noindent
Le lemme découle de l'assertion suivante : étant donné un
hyperrecouvrement $n$-tronqué de type $\mathcal{B}$, noté $K$ et scindé, pour $n \geq 1$,
on peut le prolonger en un hyperrecouvrement $(n + 1)$-tronqué de type $\mathcal{B}$ et scindé.
Ici, l'argument procède exactement comme dans la démonstration du
Lemme \ref{hypercovering-lemma-hypercovering-object}.
Nous omettons les détails précis, à l'exception des remarques suivantes.
Premièrement, l'hypothèse (2) de la démonstration du présent
lemme est inutile, car nous n'avons pas besoin que le morphisme
$K_{n + 1} \to (\text{cosk}_n K)_{n + 1}$ soit un recouvrement ;
il suffit qu'il induise une surjection sur les faisceaux d'ensembles associés,
ce qui résulte de
Sites, Lemme \ref{sites-lemma-covering-surjective-after-sheafification}.
Deuxièmement, l'hypothèse selon laquelle $\mathcal{C}$ admet des produits fibrés et des égalisateurs
garantit que $\text{SR}(\mathcal{C})$ admet des produits fibrés
et des égalisateurs, et que $F$ commute à ceux-ci
(Lemme \ref{hypercovering-lemma-coprod-prod-SR}). Cela suffit à
garantir l'existence des foncteurs cosquelette employés (voir
Simplicial, Remarque \ref{simplicial-remark-existence-cosk} et
Catégories, Lemme \ref{categories-lemma-fibre-products-equalizers-exist}).
\end{proof}

\begin{lemma}
\label{hypercovering-lemma-hypercovering-morphism-sites}
Soit $f : \mathcal{C} \to \mathcal{D}$ un morphisme de sites
donné par le foncteur $u : \mathcal{D} \to \mathcal{C}$.
Supposons que $\mathcal{D}$ et $\mathcal{C}$ admettent des égalisateurs et des
produits fibrés, et que $u$ commute à ceux-ci.
Si un objet simplicial $K$ de $\text{SR}(\mathcal{D})$
est un hyperrecouvrement, alors $u(K)$ est un hyperrecouvrement.
\end{lemma}

\begin{proof}
Si l'on écrit $K_n = \{U_{n, i}\}_{i \in I_n}$ comme dans l'introduction
de cette section, alors $u(K)$ est l'objet de $\text{SR}(\mathcal{C})$
donné par $u(K_n) = \{u(U_i)\}_{i \in I_n}$.
D'après Sites, Lemme \ref{sites-lemma-pullback-representable-sheaf},
on a $f^{-1}h_U^\# = h_{u(U)}^\#$ pour $U \in \Ob(\mathcal{D})$.
Cela signifie que $f^{-1}F(K_n)^\# = F(u(K_n))^\#$ pour tout $n$.
Vérifions les conditions (1), (2) et (3) pour que $u(K)$ soit un
hyperrecouvrement selon la Définition \ref{hypercovering-definition-hypercovering-variant}.
Comme $f^{-1}$ est un foncteur exact, on trouve que
$$
F(u(K_0))^\# = f^{-1}F(K_0)^\# \to f^{-1}* = *
$$
est surjectif comme image inverse d'un morphisme surjectif, ce qui donne (1).
De même,
$$
F(u(K_1))^\# = f^{-1}F(K_1)^\# \to
f^{-1} (F(K_0) \times F(K_0))^\# = F(u(K_0))^\# \times F(u(K_0))^\#
$$
est surjectif comme image inverse, ce qui donne (2). Pour la condition (3),
il suffit, afin de conclure par la même méthode, que
$$
F((\text{cosk}_n \text{sk}_n u(K))_{n + 1})^\# =
f^{-1}F((\text{cosk}_n \text{sk}_n K)_{n + 1})^\#
$$
Ce qui précède montre que $f^{-1}F(-) = F(u(-))$. Il suffit donc de montrer
que $u$ commute aux limites utilisées pour définir
$(\text{cosk}_n \text{sk}_n K)_{n + 1}$ pour $n \geq 1$.
D'après Simplicial, Remarque \ref{simplicial-remark-existence-cosk},
ces limites sont des limites finies connexes, auxquelles $u$ commute
par hypothèse.
\end{proof}

\begin{lemma}
\label{hypercovering-lemma-hypercovering-continuous-functor}
Soient $\mathcal{C}$, $\mathcal{D}$ des sites. Soit
$u : \mathcal{D} \to \mathcal{C}$ un foncteur continu.
Supposons que $\mathcal{D}$ et $\mathcal{C}$ admettent des produits fibrés
et que $u$ commute à ceux-ci. Soient $Y \in \mathcal{D}$ et
$K \in \text{SR}(\mathcal{D}, Y)$ un hyperrecouvrement de $Y$.
Alors $u(K)$ est un hyperrecouvrement de $u(Y)$.
\end{lemma}

\begin{proof}
C'est plus facile que la démonstration du Lemme \ref{hypercovering-lemma-hypercovering-morphism-sites},
car la notion d'hyperrecouvrement d'un objet est plus forte ; voir les
Définitions \ref{hypercovering-definition-hypercovering} et \ref{hypercovering-definition-covering-SR}.
En effet, $u$ envoie les recouvrements sur des recouvrements par définition
d'un morphisme de sites. Il suffit de vérifier que $u$ commute aux
limites utilisées pour définir
$(\text{cosk}_n \text{sk}_n K)_{n + 1}$ pour $n \geq 1$.
C'est clair, car le foncteur induit
$\mathcal{D}/Y \to \mathcal{C}/X$ commute à toutes les limites finies
(et la source et le but admettent toutes les limites finies d'après
Catégories, Lemme \ref{categories-lemma-finite-limits-exist}).
\end{proof}

\begin{lemma}
\label{hypercovering-lemma-w-contractible}
Soit $\mathcal{C}$ un site. Soit $\mathcal{B} \subset \Ob(\mathcal{C})$
un sous-ensemble. Supposons que
\begin{enumerate}
\item $\mathcal{C}$ admette des produits fibrés,
\item pour tout $X \in \Ob(\mathcal{C})$, il existe un recouvrement fini
$\{U_i \to X\}_{i \in I}$ avec $U_i \in \mathcal{B}$,
\item si $\{U_i \to X\}_{i \in I}$ est un recouvrement fini avec
$U_i \in \mathcal{B}$ et si $U \to X$ est un morphisme avec $U \in \mathcal{B}$,
alors $\{U_i \to X\}_{i \in I} \amalg \{U \to X\}$ soit un recouvrement.
\end{enumerate}
Alors, pour tout $X$, il existe un hyperrecouvrement $K$ de $X$
tel que chaque $K_n = \{U_{n, i} \to X\}_{i \in I_n}$, où
$I_n$ est fini et $U_{n, i} \in \mathcal{B}$.
\end{lemma}

\begin{proof}
Ce lemme est l'analogue du
Lemme \ref{hypercovering-lemma-quasi-separated-quasi-compact-hypercovering}
pour les sites. Pour le démontrer, on suit exactement la démonstration du
Lemme \ref{hypercovering-lemma-hypercovering-object},
en prêtant attention aux deux points suivants :
\begin{enumerate}
\item[(a)] on choisit le recouvrement initial $\{U_{0, i} \to X\}_{i \in I_0}$
avec $U_{0, i} \in \mathcal{B}$ de telle sorte que l'ensemble d'indices $I_0$ soit fini, et
\item[(b)] lors du choix des recouvrements
(\ref{hypercovering-equation-choose-covering-B}),
on choisit $J_{i'}$ fini.
\end{enumerate}
Le lecteur voit aisément qu'avec ces modifications, on obtient
des ensembles d'indices finis $I_n$ pour tout $n$.
\end{proof}

\begin{remark}
\label{hypercovering-remark-taking-disjoint-unions}
Soit $\mathcal{C}$ un site. Soient
$K$ et $L$ des objets de $\text{SR}(\mathcal{C})$.
Écrivons $K = \{U_i\}_{i \in I}$ et $L = \{V_j\}_{j \in J}$.
Supposons que $U = \coprod_{i \in I} U_i$ et $V = \coprod_{j \in J} V_j$
existent. On obtient alors
$$
\Mor_{\text{SR}(\mathcal{C})}(K, L) \longrightarrow \Mor_\mathcal{C}(U, V)
$$
comme suit. Étant donné $f : K \to L$, défini par $\alpha : I \to J$
et les $f_i : U_i \to V_{\alpha(i)}$, on obtient une transformation de foncteurs
$$
\Mor_\mathcal{C}(V, -) =
\prod\nolimits_{j \in J} \Mor_\mathcal{C}(V_j, -)
\to
\prod\nolimits_{i \in I} \Mor_\mathcal{C}(U_i, -) =
\Mor_\mathcal{C}(U, -)
$$
qui envoie $(g_j)_{j \in J}$ sur
$(g_{\alpha(i)} \circ f_i)_{i \in I}$. Le lemme de Yoneda
produit donc l'application correspondante $U \to V$. Bien entendu, $U \to V$
envoie le facteur $U_i$ dans le facteur $V_{\alpha(i)}$ par
le morphisme $f_i$.
\end{remark}

\begin{remark}
\label{hypercovering-remark-take-unions-hypercovering}
Soit $\mathcal{C}$ un site. Supposons que $\mathcal{C}$ admette
des produits fibrés et des égalisateurs, et soit $K$ un hyperrecouvrement.
Écrivons $K_n = \{U_{n, i}\}_{i \in I_n}$. Supposons que
\begin{enumerate}
\item[(a)] $U_n = \coprod_{i \in I_n} U_{n, i}$ existe, et
\item[(b)] $\coprod_{i \in I_n} h_{U_{n, i}} \to h_{U_n}$ induise
un isomorphisme sur les faisceaux associés.
\end{enumerate}
On obtient alors un autre objet simplicial $L$ de $\text{SR}(\mathcal{C})$
avec $L_n = \{U_n\}$ ; voir
la Remarque \ref{hypercovering-remark-taking-disjoint-unions}.
Nous affirmons que $L$ est un hyperrecouvrement.
Vérifions pour cela les conditions (1), (2) et (3) de la
Définition \ref{hypercovering-definition-hypercovering-variant}.
La condition (1) résulte de (b) et de la condition (1) pour $K$.
La condition (2) se démontre exactement de la même manière.
La condition (3) résulte des égalités
\begin{align*}
F((\text{cosk}_n \text{sk}_n L)_{n + 1})^\#
& =
((\text{cosk}_n \text{sk}_n F(L)^\#)_{n + 1}) \\
& =
((\text{cosk}_n \text{sk}_n F(K)^\#)_{n + 1}) \\
& =
F((\text{cosk}_n \text{sk}_n K)_{n + 1})^\#
\end{align*}
pour $n \geq 1$ ; la condition pour $K$ implique donc celle pour
$L$, exactement comme dans (1) et (2).
Remarquons que $F$ commute aux limites connexes et que le foncteur faisceau associé est exact,
ce qui démontre la première et la dernière égalité ; l'égalité du milieu résulte de
$F(K)^\# = F(L)^\#$ d'après (b).
\end{remark}

\begin{remark}
\label{hypercovering-remark-take-unions-hypercovering-X}
Soit $\mathcal{C}$ un site. Soit $X \in \Ob(\mathcal{C})$.
Supposons que $\mathcal{C}$ admette des produits fibrés, et soit $K$ un hyperrecouvrement de $X$.
Écrivons $K_n = \{U_{n, i}\}_{i \in I_n}$. Supposons que
\begin{enumerate}
\item[(a)] $U_n = \coprod_{i \in I_n} U_{n, i}$ existe,
\item[(b)] étant donnés des morphismes
$(\alpha, f_i) : \{U_i\}_{i \in I} \to \{V_j\}_{j \in J}$ et
$(\beta, g_k) : \{W_k\}_{k \in K} \to \{V_j\}_{j \in J}$
dans $\text{SR}(\mathcal{C})$, tels que
$U = \coprod U_i$, $V = \coprod V_j$ et $W = \coprod W_j$
existent, alors $U \times_V W =
\coprod_{(i, j, k), \alpha(i) = j = \beta(k)} U_i \times_{V_j} W_k$,
\item[(c)] si $(\alpha, f_i) : \{U_i\}_{i \in I} \to \{V_j\}_{j \in J}$
est un recouvrement au sens de la
Définition \ref{hypercovering-definition-covering-SR}
et si $U = \coprod U_i$ et $V = \coprod V_j$ existent,
alors le morphisme correspondant $U \to V$
de la Remarque \ref{hypercovering-remark-taking-disjoint-unions}
est un recouvrement de $\mathcal{C}$.
\end{enumerate}
On obtient alors un autre objet simplicial $L$ de $\text{SR}(\mathcal{C})$
avec $L_n = \{U_n\}$ ; voir
la Remarque \ref{hypercovering-remark-taking-disjoint-unions}.
Nous affirmons que $L$ est un hyperrecouvrement de $X$.
Vérifions pour cela les conditions (1) et (2) de la
Définition \ref{hypercovering-definition-hypercovering}.
La condition (1) résulte de (c) et de la condition (1) pour $K$,
car la condition (1) pour $K$ dit que $K_0 = \{U_{0, i}\}_{i \in I_0}$
est un recouvrement de $\{X\}$ au sens de la
Définition \ref{hypercovering-definition-covering-SR}.
La condition (2) résulte du fait que $\mathcal{C}/X$ admet
toutes les limites finies, donc que $\text{SR}(\mathcal{C}/X)$
admet toutes les limites finies, et du fait que la condition (b) dit que la
construction « prendre les sommes disjointes » commute
à ces limites finies. Ainsi, le morphisme
$$
L_{n + 1} \longrightarrow (\text{cosk}_n \text{sk}_n L)_{n + 1}
$$
est un recouvrement, puisqu'il résulte de l'application de notre
foncteur « prendre les sommes disjointes » au morphisme
$$
K_{n + 1} \longrightarrow (\text{cosk}_n \text{sk}_n K)_{n + 1}
$$
qui est supposé être un recouvrement au sens de la
Définition \ref{hypercovering-definition-covering-SR}, d'après la condition (2) pour $K$.
Cela a un sens parce que la propriété (b) garantit en particulier
que, si l'on part d'un diagramme fini d'objets
semi-représentables au-dessus de $X$
pour lesquels on peut prendre les sommes disjointes, alors
la limite du diagramme dans $\text{SR}(\mathcal{C}/X)$
est encore un objet semi-représentable au-dessus de $X$ pour lequel
on peut prendre les sommes disjointes.
\end{remark}








\bibliography{my}
\bibliographystyle{amsalpha}
\end{document}
